\section{Experimental Method}


\begin{figure}[H]
    \centering
    \resizebox{\columnwidth}{!}{%% Creator: Matplotlib, PGF backend
%%
%% To include the figure in your LaTeX document, write
%%   \input{<filename>.pgf}
%%
%% Make sure the required packages are loaded in your preamble
%%   \usepackage{pgf}
%%
%% Also ensure that all the required font packages are loaded; for instance,
%% the lmodern package is sometimes necessary when using math font.
%%   \usepackage{lmodern}
%%
%% Figures using additional raster images can only be included by \input if
%% they are in the same directory as the main LaTeX file. For loading figures
%% from other directories you can use the `import` package
%%   \usepackage{import}
%%
%% and then include the figures with
%%   \import{<path to file>}{<filename>.pgf}
%%
%% Matplotlib used the following preamble
%%   \def\mathdefault#1{#1}
%%   \everymath=\expandafter{\the\everymath\displaystyle}
%%   \IfFileExists{scrextend.sty}{
%%     \usepackage[fontsize=10.000000pt]{scrextend}
%%   }{
%%     \renewcommand{\normalsize}{\fontsize{10.000000}{12.000000}\selectfont}
%%     \normalsize
%%   }
%%   
%%   \makeatletter\@ifpackageloaded{underscore}{}{\usepackage[strings]{underscore}}\makeatother
%%
\begingroup%
\makeatletter%
\begin{pgfpicture}%
\pgfpathrectangle{\pgfpointorigin}{\pgfqpoint{6.400000in}{4.800000in}}%
\pgfusepath{use as bounding box, clip}%
\begin{pgfscope}%
\pgfsetbuttcap%
\pgfsetmiterjoin%
\definecolor{currentfill}{rgb}{1.000000,1.000000,1.000000}%
\pgfsetfillcolor{currentfill}%
\pgfsetlinewidth{0.000000pt}%
\definecolor{currentstroke}{rgb}{1.000000,1.000000,1.000000}%
\pgfsetstrokecolor{currentstroke}%
\pgfsetdash{}{0pt}%
\pgfpathmoveto{\pgfqpoint{0.000000in}{0.000000in}}%
\pgfpathlineto{\pgfqpoint{6.400000in}{0.000000in}}%
\pgfpathlineto{\pgfqpoint{6.400000in}{4.800000in}}%
\pgfpathlineto{\pgfqpoint{0.000000in}{4.800000in}}%
\pgfpathlineto{\pgfqpoint{0.000000in}{0.000000in}}%
\pgfpathclose%
\pgfusepath{fill}%
\end{pgfscope}%
\begin{pgfscope}%
\pgfsetbuttcap%
\pgfsetmiterjoin%
\definecolor{currentfill}{rgb}{1.000000,1.000000,1.000000}%
\pgfsetfillcolor{currentfill}%
\pgfsetlinewidth{0.000000pt}%
\definecolor{currentstroke}{rgb}{0.000000,0.000000,0.000000}%
\pgfsetstrokecolor{currentstroke}%
\pgfsetstrokeopacity{0.000000}%
\pgfsetdash{}{0pt}%
\pgfpathmoveto{\pgfqpoint{1.420821in}{0.917143in}}%
\pgfpathlineto{\pgfqpoint{6.250000in}{0.917143in}}%
\pgfpathlineto{\pgfqpoint{6.250000in}{4.650000in}}%
\pgfpathlineto{\pgfqpoint{1.420821in}{4.650000in}}%
\pgfpathlineto{\pgfqpoint{1.420821in}{0.917143in}}%
\pgfpathclose%
\pgfusepath{fill}%
\end{pgfscope}%
\begin{pgfscope}%
\pgfpathrectangle{\pgfqpoint{1.420821in}{0.917143in}}{\pgfqpoint{4.829179in}{3.732857in}}%
\pgfusepath{clip}%
\pgfsetbuttcap%
\pgfsetroundjoin%
\definecolor{currentfill}{rgb}{0.121569,0.466667,0.705882}%
\pgfsetfillcolor{currentfill}%
\pgfsetlinewidth{1.003750pt}%
\definecolor{currentstroke}{rgb}{0.121569,0.466667,0.705882}%
\pgfsetstrokecolor{currentstroke}%
\pgfsetdash{}{0pt}%
\pgfsys@defobject{currentmarker}{\pgfqpoint{-0.020833in}{-0.020833in}}{\pgfqpoint{0.020833in}{0.020833in}}{%
\pgfpathmoveto{\pgfqpoint{0.000000in}{-0.020833in}}%
\pgfpathcurveto{\pgfqpoint{0.005525in}{-0.020833in}}{\pgfqpoint{0.010825in}{-0.018638in}}{\pgfqpoint{0.014731in}{-0.014731in}}%
\pgfpathcurveto{\pgfqpoint{0.018638in}{-0.010825in}}{\pgfqpoint{0.020833in}{-0.005525in}}{\pgfqpoint{0.020833in}{0.000000in}}%
\pgfpathcurveto{\pgfqpoint{0.020833in}{0.005525in}}{\pgfqpoint{0.018638in}{0.010825in}}{\pgfqpoint{0.014731in}{0.014731in}}%
\pgfpathcurveto{\pgfqpoint{0.010825in}{0.018638in}}{\pgfqpoint{0.005525in}{0.020833in}}{\pgfqpoint{0.000000in}{0.020833in}}%
\pgfpathcurveto{\pgfqpoint{-0.005525in}{0.020833in}}{\pgfqpoint{-0.010825in}{0.018638in}}{\pgfqpoint{-0.014731in}{0.014731in}}%
\pgfpathcurveto{\pgfqpoint{-0.018638in}{0.010825in}}{\pgfqpoint{-0.020833in}{0.005525in}}{\pgfqpoint{-0.020833in}{0.000000in}}%
\pgfpathcurveto{\pgfqpoint{-0.020833in}{-0.005525in}}{\pgfqpoint{-0.018638in}{-0.010825in}}{\pgfqpoint{-0.014731in}{-0.014731in}}%
\pgfpathcurveto{\pgfqpoint{-0.010825in}{-0.018638in}}{\pgfqpoint{-0.005525in}{-0.020833in}}{\pgfqpoint{0.000000in}{-0.020833in}}%
\pgfpathlineto{\pgfqpoint{0.000000in}{-0.020833in}}%
\pgfpathclose%
\pgfusepath{stroke,fill}%
}%
\begin{pgfscope}%
\pgfsys@transformshift{1.640329in}{1.086819in}%
\pgfsys@useobject{currentmarker}{}%
\end{pgfscope}%
\begin{pgfscope}%
\pgfsys@transformshift{1.953912in}{4.480325in}%
\pgfsys@useobject{currentmarker}{}%
\end{pgfscope}%
\begin{pgfscope}%
\pgfsys@transformshift{2.267495in}{3.014323in}%
\pgfsys@useobject{currentmarker}{}%
\end{pgfscope}%
\begin{pgfscope}%
\pgfsys@transformshift{2.581078in}{3.749634in}%
\pgfsys@useobject{currentmarker}{}%
\end{pgfscope}%
\begin{pgfscope}%
\pgfsys@transformshift{2.894661in}{4.140043in}%
\pgfsys@useobject{currentmarker}{}%
\end{pgfscope}%
\begin{pgfscope}%
\pgfsys@transformshift{3.208244in}{2.755731in}%
\pgfsys@useobject{currentmarker}{}%
\end{pgfscope}%
\begin{pgfscope}%
\pgfsys@transformshift{3.521827in}{3.182171in}%
\pgfsys@useobject{currentmarker}{}%
\end{pgfscope}%
\begin{pgfscope}%
\pgfsys@transformshift{3.835410in}{2.899613in}%
\pgfsys@useobject{currentmarker}{}%
\end{pgfscope}%
\begin{pgfscope}%
\pgfsys@transformshift{4.148993in}{4.245282in}%
\pgfsys@useobject{currentmarker}{}%
\end{pgfscope}%
\begin{pgfscope}%
\pgfsys@transformshift{4.462577in}{2.978491in}%
\pgfsys@useobject{currentmarker}{}%
\end{pgfscope}%
\begin{pgfscope}%
\pgfsys@transformshift{4.776160in}{2.383609in}%
\pgfsys@useobject{currentmarker}{}%
\end{pgfscope}%
\begin{pgfscope}%
\pgfsys@transformshift{5.089743in}{2.652900in}%
\pgfsys@useobject{currentmarker}{}%
\end{pgfscope}%
\begin{pgfscope}%
\pgfsys@transformshift{5.403326in}{3.229834in}%
\pgfsys@useobject{currentmarker}{}%
\end{pgfscope}%
\begin{pgfscope}%
\pgfsys@transformshift{5.716909in}{2.099902in}%
\pgfsys@useobject{currentmarker}{}%
\end{pgfscope}%
\begin{pgfscope}%
\pgfsys@transformshift{6.030492in}{2.634649in}%
\pgfsys@useobject{currentmarker}{}%
\end{pgfscope}%
\end{pgfscope}%
\begin{pgfscope}%
\pgfsetbuttcap%
\pgfsetroundjoin%
\definecolor{currentfill}{rgb}{0.000000,0.000000,0.000000}%
\pgfsetfillcolor{currentfill}%
\pgfsetlinewidth{0.803000pt}%
\definecolor{currentstroke}{rgb}{0.000000,0.000000,0.000000}%
\pgfsetstrokecolor{currentstroke}%
\pgfsetdash{}{0pt}%
\pgfsys@defobject{currentmarker}{\pgfqpoint{0.000000in}{-0.048611in}}{\pgfqpoint{0.000000in}{0.000000in}}{%
\pgfpathmoveto{\pgfqpoint{0.000000in}{0.000000in}}%
\pgfpathlineto{\pgfqpoint{0.000000in}{-0.048611in}}%
\pgfusepath{stroke,fill}%
}%
\begin{pgfscope}%
\pgfsys@transformshift{1.550734in}{0.917143in}%
\pgfsys@useobject{currentmarker}{}%
\end{pgfscope}%
\end{pgfscope}%
\begin{pgfscope}%
\definecolor{textcolor}{rgb}{0.000000,0.000000,0.000000}%
\pgfsetstrokecolor{textcolor}%
\pgfsetfillcolor{textcolor}%
\pgftext[x=1.550734in,y=0.819921in,,top]{\color{textcolor}{\rmfamily\fontsize{24.000000}{28.800000}\selectfont\catcode`\^=\active\def^{\ifmmode\sp\else\^{}\fi}\catcode`\%=\active\def%{\%}$\mathdefault{0}$}}%
\end{pgfscope}%
\begin{pgfscope}%
\pgfsetbuttcap%
\pgfsetroundjoin%
\definecolor{currentfill}{rgb}{0.000000,0.000000,0.000000}%
\pgfsetfillcolor{currentfill}%
\pgfsetlinewidth{0.803000pt}%
\definecolor{currentstroke}{rgb}{0.000000,0.000000,0.000000}%
\pgfsetstrokecolor{currentstroke}%
\pgfsetdash{}{0pt}%
\pgfsys@defobject{currentmarker}{\pgfqpoint{0.000000in}{-0.048611in}}{\pgfqpoint{0.000000in}{0.000000in}}{%
\pgfpathmoveto{\pgfqpoint{0.000000in}{0.000000in}}%
\pgfpathlineto{\pgfqpoint{0.000000in}{-0.048611in}}%
\pgfusepath{stroke,fill}%
}%
\begin{pgfscope}%
\pgfsys@transformshift{2.148035in}{0.917143in}%
\pgfsys@useobject{currentmarker}{}%
\end{pgfscope}%
\end{pgfscope}%
\begin{pgfscope}%
\definecolor{textcolor}{rgb}{0.000000,0.000000,0.000000}%
\pgfsetstrokecolor{textcolor}%
\pgfsetfillcolor{textcolor}%
\pgftext[x=2.148035in,y=0.819921in,,top]{\color{textcolor}{\rmfamily\fontsize{24.000000}{28.800000}\selectfont\catcode`\^=\active\def^{\ifmmode\sp\else\^{}\fi}\catcode`\%=\active\def%{\%}$\mathdefault{2}$}}%
\end{pgfscope}%
\begin{pgfscope}%
\pgfsetbuttcap%
\pgfsetroundjoin%
\definecolor{currentfill}{rgb}{0.000000,0.000000,0.000000}%
\pgfsetfillcolor{currentfill}%
\pgfsetlinewidth{0.803000pt}%
\definecolor{currentstroke}{rgb}{0.000000,0.000000,0.000000}%
\pgfsetstrokecolor{currentstroke}%
\pgfsetdash{}{0pt}%
\pgfsys@defobject{currentmarker}{\pgfqpoint{0.000000in}{-0.048611in}}{\pgfqpoint{0.000000in}{0.000000in}}{%
\pgfpathmoveto{\pgfqpoint{0.000000in}{0.000000in}}%
\pgfpathlineto{\pgfqpoint{0.000000in}{-0.048611in}}%
\pgfusepath{stroke,fill}%
}%
\begin{pgfscope}%
\pgfsys@transformshift{2.745336in}{0.917143in}%
\pgfsys@useobject{currentmarker}{}%
\end{pgfscope}%
\end{pgfscope}%
\begin{pgfscope}%
\definecolor{textcolor}{rgb}{0.000000,0.000000,0.000000}%
\pgfsetstrokecolor{textcolor}%
\pgfsetfillcolor{textcolor}%
\pgftext[x=2.745336in,y=0.819921in,,top]{\color{textcolor}{\rmfamily\fontsize{24.000000}{28.800000}\selectfont\catcode`\^=\active\def^{\ifmmode\sp\else\^{}\fi}\catcode`\%=\active\def%{\%}$\mathdefault{4}$}}%
\end{pgfscope}%
\begin{pgfscope}%
\pgfsetbuttcap%
\pgfsetroundjoin%
\definecolor{currentfill}{rgb}{0.000000,0.000000,0.000000}%
\pgfsetfillcolor{currentfill}%
\pgfsetlinewidth{0.803000pt}%
\definecolor{currentstroke}{rgb}{0.000000,0.000000,0.000000}%
\pgfsetstrokecolor{currentstroke}%
\pgfsetdash{}{0pt}%
\pgfsys@defobject{currentmarker}{\pgfqpoint{0.000000in}{-0.048611in}}{\pgfqpoint{0.000000in}{0.000000in}}{%
\pgfpathmoveto{\pgfqpoint{0.000000in}{0.000000in}}%
\pgfpathlineto{\pgfqpoint{0.000000in}{-0.048611in}}%
\pgfusepath{stroke,fill}%
}%
\begin{pgfscope}%
\pgfsys@transformshift{3.342637in}{0.917143in}%
\pgfsys@useobject{currentmarker}{}%
\end{pgfscope}%
\end{pgfscope}%
\begin{pgfscope}%
\definecolor{textcolor}{rgb}{0.000000,0.000000,0.000000}%
\pgfsetstrokecolor{textcolor}%
\pgfsetfillcolor{textcolor}%
\pgftext[x=3.342637in,y=0.819921in,,top]{\color{textcolor}{\rmfamily\fontsize{24.000000}{28.800000}\selectfont\catcode`\^=\active\def^{\ifmmode\sp\else\^{}\fi}\catcode`\%=\active\def%{\%}$\mathdefault{6}$}}%
\end{pgfscope}%
\begin{pgfscope}%
\pgfsetbuttcap%
\pgfsetroundjoin%
\definecolor{currentfill}{rgb}{0.000000,0.000000,0.000000}%
\pgfsetfillcolor{currentfill}%
\pgfsetlinewidth{0.803000pt}%
\definecolor{currentstroke}{rgb}{0.000000,0.000000,0.000000}%
\pgfsetstrokecolor{currentstroke}%
\pgfsetdash{}{0pt}%
\pgfsys@defobject{currentmarker}{\pgfqpoint{0.000000in}{-0.048611in}}{\pgfqpoint{0.000000in}{0.000000in}}{%
\pgfpathmoveto{\pgfqpoint{0.000000in}{0.000000in}}%
\pgfpathlineto{\pgfqpoint{0.000000in}{-0.048611in}}%
\pgfusepath{stroke,fill}%
}%
\begin{pgfscope}%
\pgfsys@transformshift{3.939938in}{0.917143in}%
\pgfsys@useobject{currentmarker}{}%
\end{pgfscope}%
\end{pgfscope}%
\begin{pgfscope}%
\definecolor{textcolor}{rgb}{0.000000,0.000000,0.000000}%
\pgfsetstrokecolor{textcolor}%
\pgfsetfillcolor{textcolor}%
\pgftext[x=3.939938in,y=0.819921in,,top]{\color{textcolor}{\rmfamily\fontsize{24.000000}{28.800000}\selectfont\catcode`\^=\active\def^{\ifmmode\sp\else\^{}\fi}\catcode`\%=\active\def%{\%}$\mathdefault{8}$}}%
\end{pgfscope}%
\begin{pgfscope}%
\pgfsetbuttcap%
\pgfsetroundjoin%
\definecolor{currentfill}{rgb}{0.000000,0.000000,0.000000}%
\pgfsetfillcolor{currentfill}%
\pgfsetlinewidth{0.803000pt}%
\definecolor{currentstroke}{rgb}{0.000000,0.000000,0.000000}%
\pgfsetstrokecolor{currentstroke}%
\pgfsetdash{}{0pt}%
\pgfsys@defobject{currentmarker}{\pgfqpoint{0.000000in}{-0.048611in}}{\pgfqpoint{0.000000in}{0.000000in}}{%
\pgfpathmoveto{\pgfqpoint{0.000000in}{0.000000in}}%
\pgfpathlineto{\pgfqpoint{0.000000in}{-0.048611in}}%
\pgfusepath{stroke,fill}%
}%
\begin{pgfscope}%
\pgfsys@transformshift{4.537239in}{0.917143in}%
\pgfsys@useobject{currentmarker}{}%
\end{pgfscope}%
\end{pgfscope}%
\begin{pgfscope}%
\definecolor{textcolor}{rgb}{0.000000,0.000000,0.000000}%
\pgfsetstrokecolor{textcolor}%
\pgfsetfillcolor{textcolor}%
\pgftext[x=4.537239in,y=0.819921in,,top]{\color{textcolor}{\rmfamily\fontsize{24.000000}{28.800000}\selectfont\catcode`\^=\active\def^{\ifmmode\sp\else\^{}\fi}\catcode`\%=\active\def%{\%}$\mathdefault{10}$}}%
\end{pgfscope}%
\begin{pgfscope}%
\pgfsetbuttcap%
\pgfsetroundjoin%
\definecolor{currentfill}{rgb}{0.000000,0.000000,0.000000}%
\pgfsetfillcolor{currentfill}%
\pgfsetlinewidth{0.803000pt}%
\definecolor{currentstroke}{rgb}{0.000000,0.000000,0.000000}%
\pgfsetstrokecolor{currentstroke}%
\pgfsetdash{}{0pt}%
\pgfsys@defobject{currentmarker}{\pgfqpoint{0.000000in}{-0.048611in}}{\pgfqpoint{0.000000in}{0.000000in}}{%
\pgfpathmoveto{\pgfqpoint{0.000000in}{0.000000in}}%
\pgfpathlineto{\pgfqpoint{0.000000in}{-0.048611in}}%
\pgfusepath{stroke,fill}%
}%
\begin{pgfscope}%
\pgfsys@transformshift{5.134540in}{0.917143in}%
\pgfsys@useobject{currentmarker}{}%
\end{pgfscope}%
\end{pgfscope}%
\begin{pgfscope}%
\definecolor{textcolor}{rgb}{0.000000,0.000000,0.000000}%
\pgfsetstrokecolor{textcolor}%
\pgfsetfillcolor{textcolor}%
\pgftext[x=5.134540in,y=0.819921in,,top]{\color{textcolor}{\rmfamily\fontsize{24.000000}{28.800000}\selectfont\catcode`\^=\active\def^{\ifmmode\sp\else\^{}\fi}\catcode`\%=\active\def%{\%}$\mathdefault{12}$}}%
\end{pgfscope}%
\begin{pgfscope}%
\pgfsetbuttcap%
\pgfsetroundjoin%
\definecolor{currentfill}{rgb}{0.000000,0.000000,0.000000}%
\pgfsetfillcolor{currentfill}%
\pgfsetlinewidth{0.803000pt}%
\definecolor{currentstroke}{rgb}{0.000000,0.000000,0.000000}%
\pgfsetstrokecolor{currentstroke}%
\pgfsetdash{}{0pt}%
\pgfsys@defobject{currentmarker}{\pgfqpoint{0.000000in}{-0.048611in}}{\pgfqpoint{0.000000in}{0.000000in}}{%
\pgfpathmoveto{\pgfqpoint{0.000000in}{0.000000in}}%
\pgfpathlineto{\pgfqpoint{0.000000in}{-0.048611in}}%
\pgfusepath{stroke,fill}%
}%
\begin{pgfscope}%
\pgfsys@transformshift{5.731841in}{0.917143in}%
\pgfsys@useobject{currentmarker}{}%
\end{pgfscope}%
\end{pgfscope}%
\begin{pgfscope}%
\definecolor{textcolor}{rgb}{0.000000,0.000000,0.000000}%
\pgfsetstrokecolor{textcolor}%
\pgfsetfillcolor{textcolor}%
\pgftext[x=5.731841in,y=0.819921in,,top]{\color{textcolor}{\rmfamily\fontsize{24.000000}{28.800000}\selectfont\catcode`\^=\active\def^{\ifmmode\sp\else\^{}\fi}\catcode`\%=\active\def%{\%}$\mathdefault{14}$}}%
\end{pgfscope}%
\begin{pgfscope}%
\definecolor{textcolor}{rgb}{0.000000,0.000000,0.000000}%
\pgfsetstrokecolor{textcolor}%
\pgfsetfillcolor{textcolor}%
\pgftext[x=3.835410in,y=0.457183in,,top]{\color{textcolor}{\rmfamily\fontsize{24.000000}{28.800000}\selectfont\catcode`\^=\active\def^{\ifmmode\sp\else\^{}\fi}\catcode`\%=\active\def%{\%}Flashing time $\tau$}}%
\end{pgfscope}%
\begin{pgfscope}%
\pgfsetbuttcap%
\pgfsetroundjoin%
\definecolor{currentfill}{rgb}{0.000000,0.000000,0.000000}%
\pgfsetfillcolor{currentfill}%
\pgfsetlinewidth{0.803000pt}%
\definecolor{currentstroke}{rgb}{0.000000,0.000000,0.000000}%
\pgfsetstrokecolor{currentstroke}%
\pgfsetdash{}{0pt}%
\pgfsys@defobject{currentmarker}{\pgfqpoint{-0.048611in}{0.000000in}}{\pgfqpoint{-0.000000in}{0.000000in}}{%
\pgfpathmoveto{\pgfqpoint{-0.000000in}{0.000000in}}%
\pgfpathlineto{\pgfqpoint{-0.048611in}{0.000000in}}%
\pgfusepath{stroke,fill}%
}%
\begin{pgfscope}%
\pgfsys@transformshift{1.420821in}{1.127971in}%
\pgfsys@useobject{currentmarker}{}%
\end{pgfscope}%
\end{pgfscope}%
\begin{pgfscope}%
\definecolor{textcolor}{rgb}{0.000000,0.000000,0.000000}%
\pgfsetstrokecolor{textcolor}%
\pgfsetfillcolor{textcolor}%
\pgftext[x=0.595700in, y=1.007987in, left, base]{\color{textcolor}{\rmfamily\fontsize{24.000000}{28.800000}\selectfont\catcode`\^=\active\def^{\ifmmode\sp\else\^{}\fi}\catcode`\%=\active\def%{\%}$\mathdefault{0.000}$}}%
\end{pgfscope}%
\begin{pgfscope}%
\pgfsetbuttcap%
\pgfsetroundjoin%
\definecolor{currentfill}{rgb}{0.000000,0.000000,0.000000}%
\pgfsetfillcolor{currentfill}%
\pgfsetlinewidth{0.803000pt}%
\definecolor{currentstroke}{rgb}{0.000000,0.000000,0.000000}%
\pgfsetstrokecolor{currentstroke}%
\pgfsetdash{}{0pt}%
\pgfsys@defobject{currentmarker}{\pgfqpoint{-0.048611in}{0.000000in}}{\pgfqpoint{-0.000000in}{0.000000in}}{%
\pgfpathmoveto{\pgfqpoint{-0.000000in}{0.000000in}}%
\pgfpathlineto{\pgfqpoint{-0.048611in}{0.000000in}}%
\pgfusepath{stroke,fill}%
}%
\begin{pgfscope}%
\pgfsys@transformshift{1.420821in}{1.580037in}%
\pgfsys@useobject{currentmarker}{}%
\end{pgfscope}%
\end{pgfscope}%
\begin{pgfscope}%
\definecolor{textcolor}{rgb}{0.000000,0.000000,0.000000}%
\pgfsetstrokecolor{textcolor}%
\pgfsetfillcolor{textcolor}%
\pgftext[x=0.595700in, y=1.460053in, left, base]{\color{textcolor}{\rmfamily\fontsize{24.000000}{28.800000}\selectfont\catcode`\^=\active\def^{\ifmmode\sp\else\^{}\fi}\catcode`\%=\active\def%{\%}$\mathdefault{0.005}$}}%
\end{pgfscope}%
\begin{pgfscope}%
\pgfsetbuttcap%
\pgfsetroundjoin%
\definecolor{currentfill}{rgb}{0.000000,0.000000,0.000000}%
\pgfsetfillcolor{currentfill}%
\pgfsetlinewidth{0.803000pt}%
\definecolor{currentstroke}{rgb}{0.000000,0.000000,0.000000}%
\pgfsetstrokecolor{currentstroke}%
\pgfsetdash{}{0pt}%
\pgfsys@defobject{currentmarker}{\pgfqpoint{-0.048611in}{0.000000in}}{\pgfqpoint{-0.000000in}{0.000000in}}{%
\pgfpathmoveto{\pgfqpoint{-0.000000in}{0.000000in}}%
\pgfpathlineto{\pgfqpoint{-0.048611in}{0.000000in}}%
\pgfusepath{stroke,fill}%
}%
\begin{pgfscope}%
\pgfsys@transformshift{1.420821in}{2.032104in}%
\pgfsys@useobject{currentmarker}{}%
\end{pgfscope}%
\end{pgfscope}%
\begin{pgfscope}%
\definecolor{textcolor}{rgb}{0.000000,0.000000,0.000000}%
\pgfsetstrokecolor{textcolor}%
\pgfsetfillcolor{textcolor}%
\pgftext[x=0.595700in, y=1.912119in, left, base]{\color{textcolor}{\rmfamily\fontsize{24.000000}{28.800000}\selectfont\catcode`\^=\active\def^{\ifmmode\sp\else\^{}\fi}\catcode`\%=\active\def%{\%}$\mathdefault{0.010}$}}%
\end{pgfscope}%
\begin{pgfscope}%
\pgfsetbuttcap%
\pgfsetroundjoin%
\definecolor{currentfill}{rgb}{0.000000,0.000000,0.000000}%
\pgfsetfillcolor{currentfill}%
\pgfsetlinewidth{0.803000pt}%
\definecolor{currentstroke}{rgb}{0.000000,0.000000,0.000000}%
\pgfsetstrokecolor{currentstroke}%
\pgfsetdash{}{0pt}%
\pgfsys@defobject{currentmarker}{\pgfqpoint{-0.048611in}{0.000000in}}{\pgfqpoint{-0.000000in}{0.000000in}}{%
\pgfpathmoveto{\pgfqpoint{-0.000000in}{0.000000in}}%
\pgfpathlineto{\pgfqpoint{-0.048611in}{0.000000in}}%
\pgfusepath{stroke,fill}%
}%
\begin{pgfscope}%
\pgfsys@transformshift{1.420821in}{2.484170in}%
\pgfsys@useobject{currentmarker}{}%
\end{pgfscope}%
\end{pgfscope}%
\begin{pgfscope}%
\definecolor{textcolor}{rgb}{0.000000,0.000000,0.000000}%
\pgfsetstrokecolor{textcolor}%
\pgfsetfillcolor{textcolor}%
\pgftext[x=0.595700in, y=2.364185in, left, base]{\color{textcolor}{\rmfamily\fontsize{24.000000}{28.800000}\selectfont\catcode`\^=\active\def^{\ifmmode\sp\else\^{}\fi}\catcode`\%=\active\def%{\%}$\mathdefault{0.015}$}}%
\end{pgfscope}%
\begin{pgfscope}%
\pgfsetbuttcap%
\pgfsetroundjoin%
\definecolor{currentfill}{rgb}{0.000000,0.000000,0.000000}%
\pgfsetfillcolor{currentfill}%
\pgfsetlinewidth{0.803000pt}%
\definecolor{currentstroke}{rgb}{0.000000,0.000000,0.000000}%
\pgfsetstrokecolor{currentstroke}%
\pgfsetdash{}{0pt}%
\pgfsys@defobject{currentmarker}{\pgfqpoint{-0.048611in}{0.000000in}}{\pgfqpoint{-0.000000in}{0.000000in}}{%
\pgfpathmoveto{\pgfqpoint{-0.000000in}{0.000000in}}%
\pgfpathlineto{\pgfqpoint{-0.048611in}{0.000000in}}%
\pgfusepath{stroke,fill}%
}%
\begin{pgfscope}%
\pgfsys@transformshift{1.420821in}{2.936236in}%
\pgfsys@useobject{currentmarker}{}%
\end{pgfscope}%
\end{pgfscope}%
\begin{pgfscope}%
\definecolor{textcolor}{rgb}{0.000000,0.000000,0.000000}%
\pgfsetstrokecolor{textcolor}%
\pgfsetfillcolor{textcolor}%
\pgftext[x=0.595700in, y=2.816251in, left, base]{\color{textcolor}{\rmfamily\fontsize{24.000000}{28.800000}\selectfont\catcode`\^=\active\def^{\ifmmode\sp\else\^{}\fi}\catcode`\%=\active\def%{\%}$\mathdefault{0.020}$}}%
\end{pgfscope}%
\begin{pgfscope}%
\pgfsetbuttcap%
\pgfsetroundjoin%
\definecolor{currentfill}{rgb}{0.000000,0.000000,0.000000}%
\pgfsetfillcolor{currentfill}%
\pgfsetlinewidth{0.803000pt}%
\definecolor{currentstroke}{rgb}{0.000000,0.000000,0.000000}%
\pgfsetstrokecolor{currentstroke}%
\pgfsetdash{}{0pt}%
\pgfsys@defobject{currentmarker}{\pgfqpoint{-0.048611in}{0.000000in}}{\pgfqpoint{-0.000000in}{0.000000in}}{%
\pgfpathmoveto{\pgfqpoint{-0.000000in}{0.000000in}}%
\pgfpathlineto{\pgfqpoint{-0.048611in}{0.000000in}}%
\pgfusepath{stroke,fill}%
}%
\begin{pgfscope}%
\pgfsys@transformshift{1.420821in}{3.388302in}%
\pgfsys@useobject{currentmarker}{}%
\end{pgfscope}%
\end{pgfscope}%
\begin{pgfscope}%
\definecolor{textcolor}{rgb}{0.000000,0.000000,0.000000}%
\pgfsetstrokecolor{textcolor}%
\pgfsetfillcolor{textcolor}%
\pgftext[x=0.595700in, y=3.268317in, left, base]{\color{textcolor}{\rmfamily\fontsize{24.000000}{28.800000}\selectfont\catcode`\^=\active\def^{\ifmmode\sp\else\^{}\fi}\catcode`\%=\active\def%{\%}$\mathdefault{0.025}$}}%
\end{pgfscope}%
\begin{pgfscope}%
\pgfsetbuttcap%
\pgfsetroundjoin%
\definecolor{currentfill}{rgb}{0.000000,0.000000,0.000000}%
\pgfsetfillcolor{currentfill}%
\pgfsetlinewidth{0.803000pt}%
\definecolor{currentstroke}{rgb}{0.000000,0.000000,0.000000}%
\pgfsetstrokecolor{currentstroke}%
\pgfsetdash{}{0pt}%
\pgfsys@defobject{currentmarker}{\pgfqpoint{-0.048611in}{0.000000in}}{\pgfqpoint{-0.000000in}{0.000000in}}{%
\pgfpathmoveto{\pgfqpoint{-0.000000in}{0.000000in}}%
\pgfpathlineto{\pgfqpoint{-0.048611in}{0.000000in}}%
\pgfusepath{stroke,fill}%
}%
\begin{pgfscope}%
\pgfsys@transformshift{1.420821in}{3.840368in}%
\pgfsys@useobject{currentmarker}{}%
\end{pgfscope}%
\end{pgfscope}%
\begin{pgfscope}%
\definecolor{textcolor}{rgb}{0.000000,0.000000,0.000000}%
\pgfsetstrokecolor{textcolor}%
\pgfsetfillcolor{textcolor}%
\pgftext[x=0.595700in, y=3.720384in, left, base]{\color{textcolor}{\rmfamily\fontsize{24.000000}{28.800000}\selectfont\catcode`\^=\active\def^{\ifmmode\sp\else\^{}\fi}\catcode`\%=\active\def%{\%}$\mathdefault{0.030}$}}%
\end{pgfscope}%
\begin{pgfscope}%
\pgfsetbuttcap%
\pgfsetroundjoin%
\definecolor{currentfill}{rgb}{0.000000,0.000000,0.000000}%
\pgfsetfillcolor{currentfill}%
\pgfsetlinewidth{0.803000pt}%
\definecolor{currentstroke}{rgb}{0.000000,0.000000,0.000000}%
\pgfsetstrokecolor{currentstroke}%
\pgfsetdash{}{0pt}%
\pgfsys@defobject{currentmarker}{\pgfqpoint{-0.048611in}{0.000000in}}{\pgfqpoint{-0.000000in}{0.000000in}}{%
\pgfpathmoveto{\pgfqpoint{-0.000000in}{0.000000in}}%
\pgfpathlineto{\pgfqpoint{-0.048611in}{0.000000in}}%
\pgfusepath{stroke,fill}%
}%
\begin{pgfscope}%
\pgfsys@transformshift{1.420821in}{4.292434in}%
\pgfsys@useobject{currentmarker}{}%
\end{pgfscope}%
\end{pgfscope}%
\begin{pgfscope}%
\definecolor{textcolor}{rgb}{0.000000,0.000000,0.000000}%
\pgfsetstrokecolor{textcolor}%
\pgfsetfillcolor{textcolor}%
\pgftext[x=0.595700in, y=4.172450in, left, base]{\color{textcolor}{\rmfamily\fontsize{24.000000}{28.800000}\selectfont\catcode`\^=\active\def^{\ifmmode\sp\else\^{}\fi}\catcode`\%=\active\def%{\%}$\mathdefault{0.035}$}}%
\end{pgfscope}%
\begin{pgfscope}%
\definecolor{textcolor}{rgb}{0.000000,0.000000,0.000000}%
\pgfsetstrokecolor{textcolor}%
\pgfsetfillcolor{textcolor}%
\pgftext[x=0.540144in,y=2.783572in,,bottom,rotate=90.000000]{\color{textcolor}{\rmfamily\fontsize{24.000000}{28.800000}\selectfont\catcode`\^=\active\def^{\ifmmode\sp\else\^{}\fi}\catcode`\%=\active\def%{\%}Mean drift velocity ($\hat{x} / \hat{t}$)}}%
\end{pgfscope}%
\begin{pgfscope}%
\pgfpathrectangle{\pgfqpoint{1.420821in}{0.917143in}}{\pgfqpoint{4.829179in}{3.732857in}}%
\pgfusepath{clip}%
\pgfsetbuttcap%
\pgfsetroundjoin%
\pgfsetlinewidth{1.505625pt}%
\definecolor{currentstroke}{rgb}{1.000000,0.000000,0.000000}%
\pgfsetstrokecolor{currentstroke}%
\pgfsetdash{{5.550000pt}{2.400000pt}}{0.000000pt}%
\pgfpathmoveto{\pgfqpoint{2.894661in}{0.917143in}}%
\pgfpathlineto{\pgfqpoint{2.894661in}{4.650000in}}%
\pgfusepath{stroke}%
\end{pgfscope}%
\begin{pgfscope}%
\pgfsetrectcap%
\pgfsetmiterjoin%
\pgfsetlinewidth{0.803000pt}%
\definecolor{currentstroke}{rgb}{0.000000,0.000000,0.000000}%
\pgfsetstrokecolor{currentstroke}%
\pgfsetdash{}{0pt}%
\pgfpathmoveto{\pgfqpoint{1.420821in}{0.917143in}}%
\pgfpathlineto{\pgfqpoint{1.420821in}{4.650000in}}%
\pgfusepath{stroke}%
\end{pgfscope}%
\begin{pgfscope}%
\pgfsetrectcap%
\pgfsetmiterjoin%
\pgfsetlinewidth{0.803000pt}%
\definecolor{currentstroke}{rgb}{0.000000,0.000000,0.000000}%
\pgfsetstrokecolor{currentstroke}%
\pgfsetdash{}{0pt}%
\pgfpathmoveto{\pgfqpoint{6.250000in}{0.917143in}}%
\pgfpathlineto{\pgfqpoint{6.250000in}{4.650000in}}%
\pgfusepath{stroke}%
\end{pgfscope}%
\begin{pgfscope}%
\pgfsetrectcap%
\pgfsetmiterjoin%
\pgfsetlinewidth{0.803000pt}%
\definecolor{currentstroke}{rgb}{0.000000,0.000000,0.000000}%
\pgfsetstrokecolor{currentstroke}%
\pgfsetdash{}{0pt}%
\pgfpathmoveto{\pgfqpoint{1.420821in}{0.917143in}}%
\pgfpathlineto{\pgfqpoint{6.250000in}{0.917143in}}%
\pgfusepath{stroke}%
\end{pgfscope}%
\begin{pgfscope}%
\pgfsetrectcap%
\pgfsetmiterjoin%
\pgfsetlinewidth{0.803000pt}%
\definecolor{currentstroke}{rgb}{0.000000,0.000000,0.000000}%
\pgfsetstrokecolor{currentstroke}%
\pgfsetdash{}{0pt}%
\pgfpathmoveto{\pgfqpoint{1.420821in}{4.650000in}}%
\pgfpathlineto{\pgfqpoint{6.250000in}{4.650000in}}%
\pgfusepath{stroke}%
\end{pgfscope}%
\begin{pgfscope}%
\pgfsetbuttcap%
\pgfsetmiterjoin%
\definecolor{currentfill}{rgb}{1.000000,1.000000,1.000000}%
\pgfsetfillcolor{currentfill}%
\pgfsetfillopacity{0.800000}%
\pgfsetlinewidth{1.003750pt}%
\definecolor{currentstroke}{rgb}{0.800000,0.800000,0.800000}%
\pgfsetstrokecolor{currentstroke}%
\pgfsetstrokeopacity{0.800000}%
\pgfsetdash{}{0pt}%
\pgfpathmoveto{\pgfqpoint{1.654154in}{1.083810in}}%
\pgfpathlineto{\pgfqpoint{5.454636in}{1.083810in}}%
\pgfpathquadraticcurveto{\pgfqpoint{5.521303in}{1.083810in}}{\pgfqpoint{5.521303in}{1.150477in}}%
\pgfpathlineto{\pgfqpoint{5.521303in}{2.157719in}}%
\pgfpathquadraticcurveto{\pgfqpoint{5.521303in}{2.224386in}}{\pgfqpoint{5.454636in}{2.224386in}}%
\pgfpathlineto{\pgfqpoint{1.654154in}{2.224386in}}%
\pgfpathquadraticcurveto{\pgfqpoint{1.587487in}{2.224386in}}{\pgfqpoint{1.587487in}{2.157719in}}%
\pgfpathlineto{\pgfqpoint{1.587487in}{1.150477in}}%
\pgfpathquadraticcurveto{\pgfqpoint{1.587487in}{1.083810in}}{\pgfqpoint{1.654154in}{1.083810in}}%
\pgfpathlineto{\pgfqpoint{1.654154in}{1.083810in}}%
\pgfpathclose%
\pgfusepath{stroke,fill}%
\end{pgfscope}%
\begin{pgfscope}%
\pgfsetbuttcap%
\pgfsetroundjoin%
\pgfsetlinewidth{1.505625pt}%
\definecolor{currentstroke}{rgb}{1.000000,0.000000,0.000000}%
\pgfsetstrokecolor{currentstroke}%
\pgfsetdash{{5.550000pt}{2.400000pt}}{0.000000pt}%
\pgfpathmoveto{\pgfqpoint{1.720821in}{1.893759in}}%
\pgfpathlineto{\pgfqpoint{2.054154in}{1.893759in}}%
\pgfpathlineto{\pgfqpoint{2.387488in}{1.893759in}}%
\pgfusepath{stroke}%
\end{pgfscope}%
\begin{pgfscope}%
\definecolor{textcolor}{rgb}{0.000000,0.000000,0.000000}%
\pgfsetstrokecolor{textcolor}%
\pgfsetfillcolor{textcolor}%
\pgftext[x=2.654154in,y=1.777093in,left,base]{\color{textcolor}{\rmfamily\fontsize{24.000000}{28.800000}\selectfont\catcode`\^=\active\def^{\ifmmode\sp\else\^{}\fi}\catcode`\%=\active\def%{\%}Max at ($\hat15.0$)=4.50}}%
\end{pgfscope}%
\begin{pgfscope}%
\pgfsetbuttcap%
\pgfsetroundjoin%
\definecolor{currentfill}{rgb}{0.121569,0.466667,0.705882}%
\pgfsetfillcolor{currentfill}%
\pgfsetlinewidth{1.003750pt}%
\definecolor{currentstroke}{rgb}{0.121569,0.466667,0.705882}%
\pgfsetstrokecolor{currentstroke}%
\pgfsetdash{}{0pt}%
\pgfsys@defobject{currentmarker}{\pgfqpoint{-0.020833in}{-0.020833in}}{\pgfqpoint{0.020833in}{0.020833in}}{%
\pgfpathmoveto{\pgfqpoint{0.000000in}{-0.020833in}}%
\pgfpathcurveto{\pgfqpoint{0.005525in}{-0.020833in}}{\pgfqpoint{0.010825in}{-0.018638in}}{\pgfqpoint{0.014731in}{-0.014731in}}%
\pgfpathcurveto{\pgfqpoint{0.018638in}{-0.010825in}}{\pgfqpoint{0.020833in}{-0.005525in}}{\pgfqpoint{0.020833in}{0.000000in}}%
\pgfpathcurveto{\pgfqpoint{0.020833in}{0.005525in}}{\pgfqpoint{0.018638in}{0.010825in}}{\pgfqpoint{0.014731in}{0.014731in}}%
\pgfpathcurveto{\pgfqpoint{0.010825in}{0.018638in}}{\pgfqpoint{0.005525in}{0.020833in}}{\pgfqpoint{0.000000in}{0.020833in}}%
\pgfpathcurveto{\pgfqpoint{-0.005525in}{0.020833in}}{\pgfqpoint{-0.010825in}{0.018638in}}{\pgfqpoint{-0.014731in}{0.014731in}}%
\pgfpathcurveto{\pgfqpoint{-0.018638in}{0.010825in}}{\pgfqpoint{-0.020833in}{0.005525in}}{\pgfqpoint{-0.020833in}{0.000000in}}%
\pgfpathcurveto{\pgfqpoint{-0.020833in}{-0.005525in}}{\pgfqpoint{-0.018638in}{-0.010825in}}{\pgfqpoint{-0.014731in}{-0.014731in}}%
\pgfpathcurveto{\pgfqpoint{-0.010825in}{-0.018638in}}{\pgfqpoint{-0.005525in}{-0.020833in}}{\pgfqpoint{0.000000in}{-0.020833in}}%
\pgfpathlineto{\pgfqpoint{0.000000in}{-0.020833in}}%
\pgfpathclose%
\pgfusepath{stroke,fill}%
}%
\begin{pgfscope}%
\pgfsys@transformshift{2.054154in}{1.371857in}%
\pgfsys@useobject{currentmarker}{}%
\end{pgfscope}%
\end{pgfscope}%
\begin{pgfscope}%
\definecolor{textcolor}{rgb}{0.000000,0.000000,0.000000}%
\pgfsetstrokecolor{textcolor}%
\pgfsetfillcolor{textcolor}%
\pgftext[x=2.654154in,y=1.284357in,left,base]{\color{textcolor}{\rmfamily\fontsize{24.000000}{28.800000}\selectfont\catcode`\^=\active\def^{\ifmmode\sp\else\^{}\fi}\catcode`\%=\active\def%{\%}Mean drift velocity}}%
\end{pgfscope}%
\end{pgfpicture}%
\makeatother%
\endgroup%
}
    \caption{Drift velocity as a  of flashing time.}
    \label{fig:drift_velocity}
\end{figure}

\begin{figure}[H]
    \centering
    \resizebox{\columnwidth}{!}{%% Creator: Matplotlib, PGF backend
%%
%% To include the figure in your LaTeX document, write
%%   \input{<filename>.pgf}
%%
%% Make sure the required packages are loaded in your preamble
%%   \usepackage{pgf}
%%
%% Also ensure that all the required font packages are loaded; for instance,
%% the lmodern package is sometimes necessary when using math font.
%%   \usepackage{lmodern}
%%
%% Figures using additional raster images can only be included by \input if
%% they are in the same directory as the main LaTeX file. For loading figures
%% from other directories you can use the `import` package
%%   \usepackage{import}
%%
%% and then include the figures with
%%   \import{<path to file>}{<filename>.pgf}
%%
%% Matplotlib used the following preamble
%%   \def\mathdefault#1{#1}
%%   \everymath=\expandafter{\the\everymath\displaystyle}
%%   \IfFileExists{scrextend.sty}{
%%     \usepackage[fontsize=10.000000pt]{scrextend}
%%   }{
%%     \renewcommand{\normalsize}{\fontsize{10.000000}{12.000000}\selectfont}
%%     \normalsize
%%   }
%%   
%%   \makeatletter\@ifpackageloaded{underscore}{}{\usepackage[strings]{underscore}}\makeatother
%%
\begingroup%
\makeatletter%
\begin{pgfpicture}%
\pgfpathrectangle{\pgfpointorigin}{\pgfqpoint{6.400000in}{4.800000in}}%
\pgfusepath{use as bounding box, clip}%
\begin{pgfscope}%
\pgfsetbuttcap%
\pgfsetmiterjoin%
\definecolor{currentfill}{rgb}{1.000000,1.000000,1.000000}%
\pgfsetfillcolor{currentfill}%
\pgfsetlinewidth{0.000000pt}%
\definecolor{currentstroke}{rgb}{1.000000,1.000000,1.000000}%
\pgfsetstrokecolor{currentstroke}%
\pgfsetdash{}{0pt}%
\pgfpathmoveto{\pgfqpoint{0.000000in}{0.000000in}}%
\pgfpathlineto{\pgfqpoint{6.400000in}{0.000000in}}%
\pgfpathlineto{\pgfqpoint{6.400000in}{4.800000in}}%
\pgfpathlineto{\pgfqpoint{0.000000in}{4.800000in}}%
\pgfpathlineto{\pgfqpoint{0.000000in}{0.000000in}}%
\pgfpathclose%
\pgfusepath{fill}%
\end{pgfscope}%
\begin{pgfscope}%
\pgfsetbuttcap%
\pgfsetmiterjoin%
\definecolor{currentfill}{rgb}{1.000000,1.000000,1.000000}%
\pgfsetfillcolor{currentfill}%
\pgfsetlinewidth{0.000000pt}%
\definecolor{currentstroke}{rgb}{0.000000,0.000000,0.000000}%
\pgfsetstrokecolor{currentstroke}%
\pgfsetstrokeopacity{0.000000}%
\pgfsetdash{}{0pt}%
\pgfpathmoveto{\pgfqpoint{1.198472in}{0.823309in}}%
\pgfpathlineto{\pgfqpoint{6.250000in}{0.823309in}}%
\pgfpathlineto{\pgfqpoint{6.250000in}{4.650000in}}%
\pgfpathlineto{\pgfqpoint{1.198472in}{4.650000in}}%
\pgfpathlineto{\pgfqpoint{1.198472in}{0.823309in}}%
\pgfpathclose%
\pgfusepath{fill}%
\end{pgfscope}%
\begin{pgfscope}%
\pgfpathrectangle{\pgfqpoint{1.198472in}{0.823309in}}{\pgfqpoint{5.051528in}{3.826691in}}%
\pgfusepath{clip}%
\pgfsetbuttcap%
\pgfsetroundjoin%
\definecolor{currentfill}{rgb}{0.121569,0.466667,0.705882}%
\pgfsetfillcolor{currentfill}%
\pgfsetlinewidth{1.003750pt}%
\definecolor{currentstroke}{rgb}{0.121569,0.466667,0.705882}%
\pgfsetstrokecolor{currentstroke}%
\pgfsetdash{}{0pt}%
\pgfsys@defobject{currentmarker}{\pgfqpoint{-0.020833in}{-0.020833in}}{\pgfqpoint{0.020833in}{0.020833in}}{%
\pgfpathmoveto{\pgfqpoint{0.000000in}{-0.020833in}}%
\pgfpathcurveto{\pgfqpoint{0.005525in}{-0.020833in}}{\pgfqpoint{0.010825in}{-0.018638in}}{\pgfqpoint{0.014731in}{-0.014731in}}%
\pgfpathcurveto{\pgfqpoint{0.018638in}{-0.010825in}}{\pgfqpoint{0.020833in}{-0.005525in}}{\pgfqpoint{0.020833in}{0.000000in}}%
\pgfpathcurveto{\pgfqpoint{0.020833in}{0.005525in}}{\pgfqpoint{0.018638in}{0.010825in}}{\pgfqpoint{0.014731in}{0.014731in}}%
\pgfpathcurveto{\pgfqpoint{0.010825in}{0.018638in}}{\pgfqpoint{0.005525in}{0.020833in}}{\pgfqpoint{0.000000in}{0.020833in}}%
\pgfpathcurveto{\pgfqpoint{-0.005525in}{0.020833in}}{\pgfqpoint{-0.010825in}{0.018638in}}{\pgfqpoint{-0.014731in}{0.014731in}}%
\pgfpathcurveto{\pgfqpoint{-0.018638in}{0.010825in}}{\pgfqpoint{-0.020833in}{0.005525in}}{\pgfqpoint{-0.020833in}{0.000000in}}%
\pgfpathcurveto{\pgfqpoint{-0.020833in}{-0.005525in}}{\pgfqpoint{-0.018638in}{-0.010825in}}{\pgfqpoint{-0.014731in}{-0.014731in}}%
\pgfpathcurveto{\pgfqpoint{-0.010825in}{-0.018638in}}{\pgfqpoint{-0.005525in}{-0.020833in}}{\pgfqpoint{0.000000in}{-0.020833in}}%
\pgfpathlineto{\pgfqpoint{0.000000in}{-0.020833in}}%
\pgfpathclose%
\pgfusepath{stroke,fill}%
}%
\begin{pgfscope}%
\pgfsys@transformshift{1.428087in}{1.029035in}%
\pgfsys@useobject{currentmarker}{}%
\end{pgfscope}%
\begin{pgfscope}%
\pgfsys@transformshift{1.474473in}{1.031930in}%
\pgfsys@useobject{currentmarker}{}%
\end{pgfscope}%
\begin{pgfscope}%
\pgfsys@transformshift{1.520860in}{1.417404in}%
\pgfsys@useobject{currentmarker}{}%
\end{pgfscope}%
\begin{pgfscope}%
\pgfsys@transformshift{1.567247in}{1.652517in}%
\pgfsys@useobject{currentmarker}{}%
\end{pgfscope}%
\begin{pgfscope}%
\pgfsys@transformshift{1.613634in}{1.902730in}%
\pgfsys@useobject{currentmarker}{}%
\end{pgfscope}%
\begin{pgfscope}%
\pgfsys@transformshift{1.660021in}{1.622878in}%
\pgfsys@useobject{currentmarker}{}%
\end{pgfscope}%
\begin{pgfscope}%
\pgfsys@transformshift{1.706408in}{1.763515in}%
\pgfsys@useobject{currentmarker}{}%
\end{pgfscope}%
\begin{pgfscope}%
\pgfsys@transformshift{1.752795in}{1.739529in}%
\pgfsys@useobject{currentmarker}{}%
\end{pgfscope}%
\begin{pgfscope}%
\pgfsys@transformshift{1.799181in}{2.529799in}%
\pgfsys@useobject{currentmarker}{}%
\end{pgfscope}%
\begin{pgfscope}%
\pgfsys@transformshift{1.845568in}{2.729735in}%
\pgfsys@useobject{currentmarker}{}%
\end{pgfscope}%
\begin{pgfscope}%
\pgfsys@transformshift{1.891955in}{3.119332in}%
\pgfsys@useobject{currentmarker}{}%
\end{pgfscope}%
\begin{pgfscope}%
\pgfsys@transformshift{1.938342in}{3.041700in}%
\pgfsys@useobject{currentmarker}{}%
\end{pgfscope}%
\begin{pgfscope}%
\pgfsys@transformshift{1.984729in}{2.939907in}%
\pgfsys@useobject{currentmarker}{}%
\end{pgfscope}%
\begin{pgfscope}%
\pgfsys@transformshift{2.031116in}{3.247228in}%
\pgfsys@useobject{currentmarker}{}%
\end{pgfscope}%
\begin{pgfscope}%
\pgfsys@transformshift{2.077503in}{3.277809in}%
\pgfsys@useobject{currentmarker}{}%
\end{pgfscope}%
\begin{pgfscope}%
\pgfsys@transformshift{2.123889in}{3.458160in}%
\pgfsys@useobject{currentmarker}{}%
\end{pgfscope}%
\begin{pgfscope}%
\pgfsys@transformshift{2.170276in}{4.168953in}%
\pgfsys@useobject{currentmarker}{}%
\end{pgfscope}%
\begin{pgfscope}%
\pgfsys@transformshift{2.216663in}{3.170995in}%
\pgfsys@useobject{currentmarker}{}%
\end{pgfscope}%
\begin{pgfscope}%
\pgfsys@transformshift{2.263050in}{3.567808in}%
\pgfsys@useobject{currentmarker}{}%
\end{pgfscope}%
\begin{pgfscope}%
\pgfsys@transformshift{2.309437in}{4.133891in}%
\pgfsys@useobject{currentmarker}{}%
\end{pgfscope}%
\begin{pgfscope}%
\pgfsys@transformshift{2.355824in}{4.019726in}%
\pgfsys@useobject{currentmarker}{}%
\end{pgfscope}%
\begin{pgfscope}%
\pgfsys@transformshift{2.402211in}{3.762430in}%
\pgfsys@useobject{currentmarker}{}%
\end{pgfscope}%
\begin{pgfscope}%
\pgfsys@transformshift{2.448597in}{4.289309in}%
\pgfsys@useobject{currentmarker}{}%
\end{pgfscope}%
\begin{pgfscope}%
\pgfsys@transformshift{2.494984in}{3.780983in}%
\pgfsys@useobject{currentmarker}{}%
\end{pgfscope}%
\begin{pgfscope}%
\pgfsys@transformshift{2.541371in}{4.402448in}%
\pgfsys@useobject{currentmarker}{}%
\end{pgfscope}%
\begin{pgfscope}%
\pgfsys@transformshift{2.587758in}{4.346526in}%
\pgfsys@useobject{currentmarker}{}%
\end{pgfscope}%
\begin{pgfscope}%
\pgfsys@transformshift{2.634145in}{3.888282in}%
\pgfsys@useobject{currentmarker}{}%
\end{pgfscope}%
\begin{pgfscope}%
\pgfsys@transformshift{2.680532in}{4.080641in}%
\pgfsys@useobject{currentmarker}{}%
\end{pgfscope}%
\begin{pgfscope}%
\pgfsys@transformshift{2.726918in}{4.319215in}%
\pgfsys@useobject{currentmarker}{}%
\end{pgfscope}%
\begin{pgfscope}%
\pgfsys@transformshift{2.773305in}{4.160212in}%
\pgfsys@useobject{currentmarker}{}%
\end{pgfscope}%
\begin{pgfscope}%
\pgfsys@transformshift{2.819692in}{3.983558in}%
\pgfsys@useobject{currentmarker}{}%
\end{pgfscope}%
\begin{pgfscope}%
\pgfsys@transformshift{2.866079in}{4.029502in}%
\pgfsys@useobject{currentmarker}{}%
\end{pgfscope}%
\begin{pgfscope}%
\pgfsys@transformshift{2.912466in}{4.476060in}%
\pgfsys@useobject{currentmarker}{}%
\end{pgfscope}%
\begin{pgfscope}%
\pgfsys@transformshift{2.958853in}{3.826161in}%
\pgfsys@useobject{currentmarker}{}%
\end{pgfscope}%
\begin{pgfscope}%
\pgfsys@transformshift{3.005240in}{4.057798in}%
\pgfsys@useobject{currentmarker}{}%
\end{pgfscope}%
\begin{pgfscope}%
\pgfsys@transformshift{3.051626in}{4.157806in}%
\pgfsys@useobject{currentmarker}{}%
\end{pgfscope}%
\begin{pgfscope}%
\pgfsys@transformshift{3.098013in}{4.321673in}%
\pgfsys@useobject{currentmarker}{}%
\end{pgfscope}%
\begin{pgfscope}%
\pgfsys@transformshift{3.144400in}{4.194408in}%
\pgfsys@useobject{currentmarker}{}%
\end{pgfscope}%
\begin{pgfscope}%
\pgfsys@transformshift{3.190787in}{4.106303in}%
\pgfsys@useobject{currentmarker}{}%
\end{pgfscope}%
\begin{pgfscope}%
\pgfsys@transformshift{3.237174in}{4.217397in}%
\pgfsys@useobject{currentmarker}{}%
\end{pgfscope}%
\begin{pgfscope}%
\pgfsys@transformshift{3.283561in}{4.059182in}%
\pgfsys@useobject{currentmarker}{}%
\end{pgfscope}%
\begin{pgfscope}%
\pgfsys@transformshift{3.329948in}{3.907826in}%
\pgfsys@useobject{currentmarker}{}%
\end{pgfscope}%
\begin{pgfscope}%
\pgfsys@transformshift{3.376334in}{4.277046in}%
\pgfsys@useobject{currentmarker}{}%
\end{pgfscope}%
\begin{pgfscope}%
\pgfsys@transformshift{3.422721in}{4.098876in}%
\pgfsys@useobject{currentmarker}{}%
\end{pgfscope}%
\begin{pgfscope}%
\pgfsys@transformshift{3.469108in}{3.956603in}%
\pgfsys@useobject{currentmarker}{}%
\end{pgfscope}%
\begin{pgfscope}%
\pgfsys@transformshift{3.515495in}{4.047100in}%
\pgfsys@useobject{currentmarker}{}%
\end{pgfscope}%
\begin{pgfscope}%
\pgfsys@transformshift{3.561882in}{3.977132in}%
\pgfsys@useobject{currentmarker}{}%
\end{pgfscope}%
\begin{pgfscope}%
\pgfsys@transformshift{3.608269in}{4.017708in}%
\pgfsys@useobject{currentmarker}{}%
\end{pgfscope}%
\begin{pgfscope}%
\pgfsys@transformshift{3.654656in}{3.977617in}%
\pgfsys@useobject{currentmarker}{}%
\end{pgfscope}%
\begin{pgfscope}%
\pgfsys@transformshift{3.701042in}{4.068990in}%
\pgfsys@useobject{currentmarker}{}%
\end{pgfscope}%
\begin{pgfscope}%
\pgfsys@transformshift{3.747429in}{3.746862in}%
\pgfsys@useobject{currentmarker}{}%
\end{pgfscope}%
\begin{pgfscope}%
\pgfsys@transformshift{3.793816in}{3.940185in}%
\pgfsys@useobject{currentmarker}{}%
\end{pgfscope}%
\begin{pgfscope}%
\pgfsys@transformshift{3.840203in}{3.723326in}%
\pgfsys@useobject{currentmarker}{}%
\end{pgfscope}%
\begin{pgfscope}%
\pgfsys@transformshift{3.886590in}{3.902024in}%
\pgfsys@useobject{currentmarker}{}%
\end{pgfscope}%
\begin{pgfscope}%
\pgfsys@transformshift{3.932977in}{3.708164in}%
\pgfsys@useobject{currentmarker}{}%
\end{pgfscope}%
\begin{pgfscope}%
\pgfsys@transformshift{3.979364in}{3.826662in}%
\pgfsys@useobject{currentmarker}{}%
\end{pgfscope}%
\begin{pgfscope}%
\pgfsys@transformshift{4.025750in}{3.963841in}%
\pgfsys@useobject{currentmarker}{}%
\end{pgfscope}%
\begin{pgfscope}%
\pgfsys@transformshift{4.072137in}{3.937920in}%
\pgfsys@useobject{currentmarker}{}%
\end{pgfscope}%
\begin{pgfscope}%
\pgfsys@transformshift{4.118524in}{3.757263in}%
\pgfsys@useobject{currentmarker}{}%
\end{pgfscope}%
\begin{pgfscope}%
\pgfsys@transformshift{4.164911in}{3.516342in}%
\pgfsys@useobject{currentmarker}{}%
\end{pgfscope}%
\begin{pgfscope}%
\pgfsys@transformshift{4.211298in}{3.714071in}%
\pgfsys@useobject{currentmarker}{}%
\end{pgfscope}%
\begin{pgfscope}%
\pgfsys@transformshift{4.257685in}{3.700810in}%
\pgfsys@useobject{currentmarker}{}%
\end{pgfscope}%
\begin{pgfscope}%
\pgfsys@transformshift{4.304072in}{3.604658in}%
\pgfsys@useobject{currentmarker}{}%
\end{pgfscope}%
\begin{pgfscope}%
\pgfsys@transformshift{4.350458in}{3.867464in}%
\pgfsys@useobject{currentmarker}{}%
\end{pgfscope}%
\begin{pgfscope}%
\pgfsys@transformshift{4.396845in}{4.000738in}%
\pgfsys@useobject{currentmarker}{}%
\end{pgfscope}%
\begin{pgfscope}%
\pgfsys@transformshift{4.443232in}{3.725003in}%
\pgfsys@useobject{currentmarker}{}%
\end{pgfscope}%
\begin{pgfscope}%
\pgfsys@transformshift{4.489619in}{3.692578in}%
\pgfsys@useobject{currentmarker}{}%
\end{pgfscope}%
\begin{pgfscope}%
\pgfsys@transformshift{4.536006in}{3.783930in}%
\pgfsys@useobject{currentmarker}{}%
\end{pgfscope}%
\begin{pgfscope}%
\pgfsys@transformshift{4.582393in}{3.630938in}%
\pgfsys@useobject{currentmarker}{}%
\end{pgfscope}%
\begin{pgfscope}%
\pgfsys@transformshift{4.628779in}{3.657836in}%
\pgfsys@useobject{currentmarker}{}%
\end{pgfscope}%
\begin{pgfscope}%
\pgfsys@transformshift{4.675166in}{3.554406in}%
\pgfsys@useobject{currentmarker}{}%
\end{pgfscope}%
\begin{pgfscope}%
\pgfsys@transformshift{4.721553in}{3.715771in}%
\pgfsys@useobject{currentmarker}{}%
\end{pgfscope}%
\begin{pgfscope}%
\pgfsys@transformshift{4.767940in}{3.529670in}%
\pgfsys@useobject{currentmarker}{}%
\end{pgfscope}%
\begin{pgfscope}%
\pgfsys@transformshift{4.814327in}{3.622021in}%
\pgfsys@useobject{currentmarker}{}%
\end{pgfscope}%
\begin{pgfscope}%
\pgfsys@transformshift{4.860714in}{3.688223in}%
\pgfsys@useobject{currentmarker}{}%
\end{pgfscope}%
\begin{pgfscope}%
\pgfsys@transformshift{4.907101in}{3.631471in}%
\pgfsys@useobject{currentmarker}{}%
\end{pgfscope}%
\begin{pgfscope}%
\pgfsys@transformshift{4.953487in}{3.678922in}%
\pgfsys@useobject{currentmarker}{}%
\end{pgfscope}%
\begin{pgfscope}%
\pgfsys@transformshift{4.999874in}{3.600914in}%
\pgfsys@useobject{currentmarker}{}%
\end{pgfscope}%
\begin{pgfscope}%
\pgfsys@transformshift{5.046261in}{3.485905in}%
\pgfsys@useobject{currentmarker}{}%
\end{pgfscope}%
\begin{pgfscope}%
\pgfsys@transformshift{5.092648in}{3.362309in}%
\pgfsys@useobject{currentmarker}{}%
\end{pgfscope}%
\begin{pgfscope}%
\pgfsys@transformshift{5.139035in}{3.470716in}%
\pgfsys@useobject{currentmarker}{}%
\end{pgfscope}%
\begin{pgfscope}%
\pgfsys@transformshift{5.185422in}{3.431034in}%
\pgfsys@useobject{currentmarker}{}%
\end{pgfscope}%
\begin{pgfscope}%
\pgfsys@transformshift{5.231809in}{3.690377in}%
\pgfsys@useobject{currentmarker}{}%
\end{pgfscope}%
\begin{pgfscope}%
\pgfsys@transformshift{5.278195in}{3.439220in}%
\pgfsys@useobject{currentmarker}{}%
\end{pgfscope}%
\begin{pgfscope}%
\pgfsys@transformshift{5.324582in}{3.271323in}%
\pgfsys@useobject{currentmarker}{}%
\end{pgfscope}%
\begin{pgfscope}%
\pgfsys@transformshift{5.370969in}{3.605810in}%
\pgfsys@useobject{currentmarker}{}%
\end{pgfscope}%
\begin{pgfscope}%
\pgfsys@transformshift{5.417356in}{3.368594in}%
\pgfsys@useobject{currentmarker}{}%
\end{pgfscope}%
\begin{pgfscope}%
\pgfsys@transformshift{5.463743in}{3.347888in}%
\pgfsys@useobject{currentmarker}{}%
\end{pgfscope}%
\begin{pgfscope}%
\pgfsys@transformshift{5.510130in}{3.322738in}%
\pgfsys@useobject{currentmarker}{}%
\end{pgfscope}%
\begin{pgfscope}%
\pgfsys@transformshift{5.556517in}{3.219312in}%
\pgfsys@useobject{currentmarker}{}%
\end{pgfscope}%
\begin{pgfscope}%
\pgfsys@transformshift{5.602903in}{3.303101in}%
\pgfsys@useobject{currentmarker}{}%
\end{pgfscope}%
\begin{pgfscope}%
\pgfsys@transformshift{5.649290in}{3.217008in}%
\pgfsys@useobject{currentmarker}{}%
\end{pgfscope}%
\begin{pgfscope}%
\pgfsys@transformshift{5.695677in}{3.284293in}%
\pgfsys@useobject{currentmarker}{}%
\end{pgfscope}%
\begin{pgfscope}%
\pgfsys@transformshift{5.742064in}{3.449672in}%
\pgfsys@useobject{currentmarker}{}%
\end{pgfscope}%
\begin{pgfscope}%
\pgfsys@transformshift{5.788451in}{3.397079in}%
\pgfsys@useobject{currentmarker}{}%
\end{pgfscope}%
\begin{pgfscope}%
\pgfsys@transformshift{5.834838in}{3.299247in}%
\pgfsys@useobject{currentmarker}{}%
\end{pgfscope}%
\begin{pgfscope}%
\pgfsys@transformshift{5.881225in}{3.121382in}%
\pgfsys@useobject{currentmarker}{}%
\end{pgfscope}%
\begin{pgfscope}%
\pgfsys@transformshift{5.927611in}{3.276714in}%
\pgfsys@useobject{currentmarker}{}%
\end{pgfscope}%
\begin{pgfscope}%
\pgfsys@transformshift{5.973998in}{3.317536in}%
\pgfsys@useobject{currentmarker}{}%
\end{pgfscope}%
\begin{pgfscope}%
\pgfsys@transformshift{6.020385in}{3.344049in}%
\pgfsys@useobject{currentmarker}{}%
\end{pgfscope}%
\end{pgfscope}%
\begin{pgfscope}%
\pgfpathrectangle{\pgfqpoint{1.198472in}{0.823309in}}{\pgfqpoint{5.051528in}{3.826691in}}%
\pgfusepath{clip}%
\pgfsetbuttcap%
\pgfsetroundjoin%
\pgfsetlinewidth{1.505625pt}%
\definecolor{currentstroke}{rgb}{1.000000,0.498039,0.313725}%
\pgfsetstrokecolor{currentstroke}%
\pgfsetdash{}{0pt}%
\pgfpathmoveto{\pgfqpoint{1.386420in}{1.002147in}}%
\pgfpathlineto{\pgfqpoint{1.469753in}{1.002147in}}%
\pgfpathlineto{\pgfqpoint{1.469753in}{1.085480in}}%
\pgfpathlineto{\pgfqpoint{1.386420in}{1.085480in}}%
\pgfpathlineto{\pgfqpoint{1.386420in}{1.002147in}}%
\pgfpathclose%
\pgfusepath{stroke}%
\end{pgfscope}%
\begin{pgfscope}%
\pgfpathrectangle{\pgfqpoint{1.198472in}{0.823309in}}{\pgfqpoint{5.051528in}{3.826691in}}%
\pgfusepath{clip}%
\pgfsetbuttcap%
\pgfsetroundjoin%
\pgfsetlinewidth{1.505625pt}%
\definecolor{currentstroke}{rgb}{1.000000,0.498039,0.313725}%
\pgfsetstrokecolor{currentstroke}%
\pgfsetdash{}{0pt}%
\pgfpathmoveto{\pgfqpoint{1.432807in}{1.003595in}}%
\pgfpathlineto{\pgfqpoint{1.516140in}{1.003595in}}%
\pgfpathlineto{\pgfqpoint{1.516140in}{1.086928in}}%
\pgfpathlineto{\pgfqpoint{1.432807in}{1.086928in}}%
\pgfpathlineto{\pgfqpoint{1.432807in}{1.003595in}}%
\pgfpathclose%
\pgfusepath{stroke}%
\end{pgfscope}%
\begin{pgfscope}%
\pgfpathrectangle{\pgfqpoint{1.198472in}{0.823309in}}{\pgfqpoint{5.051528in}{3.826691in}}%
\pgfusepath{clip}%
\pgfsetbuttcap%
\pgfsetroundjoin%
\pgfsetlinewidth{1.505625pt}%
\definecolor{currentstroke}{rgb}{1.000000,0.498039,0.313725}%
\pgfsetstrokecolor{currentstroke}%
\pgfsetdash{}{0pt}%
\pgfpathmoveto{\pgfqpoint{1.479194in}{1.196332in}}%
\pgfpathlineto{\pgfqpoint{1.562527in}{1.196332in}}%
\pgfpathlineto{\pgfqpoint{1.562527in}{1.279665in}}%
\pgfpathlineto{\pgfqpoint{1.479194in}{1.279665in}}%
\pgfpathlineto{\pgfqpoint{1.479194in}{1.196332in}}%
\pgfpathclose%
\pgfusepath{stroke}%
\end{pgfscope}%
\begin{pgfscope}%
\pgfpathrectangle{\pgfqpoint{1.198472in}{0.823309in}}{\pgfqpoint{5.051528in}{3.826691in}}%
\pgfusepath{clip}%
\pgfsetbuttcap%
\pgfsetroundjoin%
\pgfsetlinewidth{1.505625pt}%
\definecolor{currentstroke}{rgb}{1.000000,0.498039,0.313725}%
\pgfsetstrokecolor{currentstroke}%
\pgfsetdash{}{0pt}%
\pgfpathmoveto{\pgfqpoint{1.525580in}{1.313888in}}%
\pgfpathlineto{\pgfqpoint{1.608914in}{1.313888in}}%
\pgfpathlineto{\pgfqpoint{1.608914in}{1.397222in}}%
\pgfpathlineto{\pgfqpoint{1.525580in}{1.397222in}}%
\pgfpathlineto{\pgfqpoint{1.525580in}{1.313888in}}%
\pgfpathclose%
\pgfusepath{stroke}%
\end{pgfscope}%
\begin{pgfscope}%
\pgfpathrectangle{\pgfqpoint{1.198472in}{0.823309in}}{\pgfqpoint{5.051528in}{3.826691in}}%
\pgfusepath{clip}%
\pgfsetbuttcap%
\pgfsetroundjoin%
\pgfsetlinewidth{1.505625pt}%
\definecolor{currentstroke}{rgb}{1.000000,0.498039,0.313725}%
\pgfsetstrokecolor{currentstroke}%
\pgfsetdash{}{0pt}%
\pgfpathmoveto{\pgfqpoint{1.571967in}{1.438995in}}%
\pgfpathlineto{\pgfqpoint{1.655301in}{1.438995in}}%
\pgfpathlineto{\pgfqpoint{1.655301in}{1.522328in}}%
\pgfpathlineto{\pgfqpoint{1.571967in}{1.522328in}}%
\pgfpathlineto{\pgfqpoint{1.571967in}{1.438995in}}%
\pgfpathclose%
\pgfusepath{stroke}%
\end{pgfscope}%
\begin{pgfscope}%
\pgfpathrectangle{\pgfqpoint{1.198472in}{0.823309in}}{\pgfqpoint{5.051528in}{3.826691in}}%
\pgfusepath{clip}%
\pgfsetbuttcap%
\pgfsetroundjoin%
\pgfsetlinewidth{1.505625pt}%
\definecolor{currentstroke}{rgb}{1.000000,0.498039,0.313725}%
\pgfsetstrokecolor{currentstroke}%
\pgfsetdash{}{0pt}%
\pgfpathmoveto{\pgfqpoint{1.618354in}{1.299069in}}%
\pgfpathlineto{\pgfqpoint{1.701688in}{1.299069in}}%
\pgfpathlineto{\pgfqpoint{1.701688in}{1.382402in}}%
\pgfpathlineto{\pgfqpoint{1.618354in}{1.382402in}}%
\pgfpathlineto{\pgfqpoint{1.618354in}{1.299069in}}%
\pgfpathclose%
\pgfusepath{stroke}%
\end{pgfscope}%
\begin{pgfscope}%
\pgfpathrectangle{\pgfqpoint{1.198472in}{0.823309in}}{\pgfqpoint{5.051528in}{3.826691in}}%
\pgfusepath{clip}%
\pgfsetbuttcap%
\pgfsetroundjoin%
\pgfsetlinewidth{1.505625pt}%
\definecolor{currentstroke}{rgb}{1.000000,0.498039,0.313725}%
\pgfsetstrokecolor{currentstroke}%
\pgfsetdash{}{0pt}%
\pgfpathmoveto{\pgfqpoint{1.664741in}{1.369387in}}%
\pgfpathlineto{\pgfqpoint{1.748074in}{1.369387in}}%
\pgfpathlineto{\pgfqpoint{1.748074in}{1.452721in}}%
\pgfpathlineto{\pgfqpoint{1.664741in}{1.452721in}}%
\pgfpathlineto{\pgfqpoint{1.664741in}{1.369387in}}%
\pgfpathclose%
\pgfusepath{stroke}%
\end{pgfscope}%
\begin{pgfscope}%
\pgfpathrectangle{\pgfqpoint{1.198472in}{0.823309in}}{\pgfqpoint{5.051528in}{3.826691in}}%
\pgfusepath{clip}%
\pgfsetbuttcap%
\pgfsetroundjoin%
\pgfsetlinewidth{1.505625pt}%
\definecolor{currentstroke}{rgb}{1.000000,0.498039,0.313725}%
\pgfsetstrokecolor{currentstroke}%
\pgfsetdash{}{0pt}%
\pgfpathmoveto{\pgfqpoint{1.711128in}{1.357394in}}%
\pgfpathlineto{\pgfqpoint{1.794461in}{1.357394in}}%
\pgfpathlineto{\pgfqpoint{1.794461in}{1.440728in}}%
\pgfpathlineto{\pgfqpoint{1.711128in}{1.440728in}}%
\pgfpathlineto{\pgfqpoint{1.711128in}{1.357394in}}%
\pgfpathclose%
\pgfusepath{stroke}%
\end{pgfscope}%
\begin{pgfscope}%
\pgfpathrectangle{\pgfqpoint{1.198472in}{0.823309in}}{\pgfqpoint{5.051528in}{3.826691in}}%
\pgfusepath{clip}%
\pgfsetbuttcap%
\pgfsetroundjoin%
\pgfsetlinewidth{1.505625pt}%
\definecolor{currentstroke}{rgb}{1.000000,0.498039,0.313725}%
\pgfsetstrokecolor{currentstroke}%
\pgfsetdash{}{0pt}%
\pgfpathmoveto{\pgfqpoint{1.757515in}{1.752529in}}%
\pgfpathlineto{\pgfqpoint{1.840848in}{1.752529in}}%
\pgfpathlineto{\pgfqpoint{1.840848in}{1.835863in}}%
\pgfpathlineto{\pgfqpoint{1.757515in}{1.835863in}}%
\pgfpathlineto{\pgfqpoint{1.757515in}{1.752529in}}%
\pgfpathclose%
\pgfusepath{stroke}%
\end{pgfscope}%
\begin{pgfscope}%
\pgfpathrectangle{\pgfqpoint{1.198472in}{0.823309in}}{\pgfqpoint{5.051528in}{3.826691in}}%
\pgfusepath{clip}%
\pgfsetbuttcap%
\pgfsetroundjoin%
\pgfsetlinewidth{1.505625pt}%
\definecolor{currentstroke}{rgb}{1.000000,0.498039,0.313725}%
\pgfsetstrokecolor{currentstroke}%
\pgfsetdash{}{0pt}%
\pgfpathmoveto{\pgfqpoint{1.803902in}{1.852497in}}%
\pgfpathlineto{\pgfqpoint{1.887235in}{1.852497in}}%
\pgfpathlineto{\pgfqpoint{1.887235in}{1.935830in}}%
\pgfpathlineto{\pgfqpoint{1.803902in}{1.935830in}}%
\pgfpathlineto{\pgfqpoint{1.803902in}{1.852497in}}%
\pgfpathclose%
\pgfusepath{stroke}%
\end{pgfscope}%
\begin{pgfscope}%
\pgfpathrectangle{\pgfqpoint{1.198472in}{0.823309in}}{\pgfqpoint{5.051528in}{3.826691in}}%
\pgfusepath{clip}%
\pgfsetbuttcap%
\pgfsetroundjoin%
\pgfsetlinewidth{1.505625pt}%
\definecolor{currentstroke}{rgb}{1.000000,0.498039,0.313725}%
\pgfsetstrokecolor{currentstroke}%
\pgfsetdash{}{0pt}%
\pgfpathmoveto{\pgfqpoint{1.850288in}{2.047296in}}%
\pgfpathlineto{\pgfqpoint{1.933622in}{2.047296in}}%
\pgfpathlineto{\pgfqpoint{1.933622in}{2.130629in}}%
\pgfpathlineto{\pgfqpoint{1.850288in}{2.130629in}}%
\pgfpathlineto{\pgfqpoint{1.850288in}{2.047296in}}%
\pgfpathclose%
\pgfusepath{stroke}%
\end{pgfscope}%
\begin{pgfscope}%
\pgfpathrectangle{\pgfqpoint{1.198472in}{0.823309in}}{\pgfqpoint{5.051528in}{3.826691in}}%
\pgfusepath{clip}%
\pgfsetbuttcap%
\pgfsetroundjoin%
\pgfsetlinewidth{1.505625pt}%
\definecolor{currentstroke}{rgb}{1.000000,0.498039,0.313725}%
\pgfsetstrokecolor{currentstroke}%
\pgfsetdash{}{0pt}%
\pgfpathmoveto{\pgfqpoint{1.896675in}{2.008479in}}%
\pgfpathlineto{\pgfqpoint{1.980009in}{2.008479in}}%
\pgfpathlineto{\pgfqpoint{1.980009in}{2.091813in}}%
\pgfpathlineto{\pgfqpoint{1.896675in}{2.091813in}}%
\pgfpathlineto{\pgfqpoint{1.896675in}{2.008479in}}%
\pgfpathclose%
\pgfusepath{stroke}%
\end{pgfscope}%
\begin{pgfscope}%
\pgfpathrectangle{\pgfqpoint{1.198472in}{0.823309in}}{\pgfqpoint{5.051528in}{3.826691in}}%
\pgfusepath{clip}%
\pgfsetbuttcap%
\pgfsetroundjoin%
\pgfsetlinewidth{1.505625pt}%
\definecolor{currentstroke}{rgb}{1.000000,0.498039,0.313725}%
\pgfsetstrokecolor{currentstroke}%
\pgfsetdash{}{0pt}%
\pgfpathmoveto{\pgfqpoint{1.943062in}{1.957583in}}%
\pgfpathlineto{\pgfqpoint{2.026395in}{1.957583in}}%
\pgfpathlineto{\pgfqpoint{2.026395in}{2.040917in}}%
\pgfpathlineto{\pgfqpoint{1.943062in}{2.040917in}}%
\pgfpathlineto{\pgfqpoint{1.943062in}{1.957583in}}%
\pgfpathclose%
\pgfusepath{stroke}%
\end{pgfscope}%
\begin{pgfscope}%
\pgfpathrectangle{\pgfqpoint{1.198472in}{0.823309in}}{\pgfqpoint{5.051528in}{3.826691in}}%
\pgfusepath{clip}%
\pgfsetbuttcap%
\pgfsetroundjoin%
\pgfsetlinewidth{1.505625pt}%
\definecolor{currentstroke}{rgb}{1.000000,0.498039,0.313725}%
\pgfsetstrokecolor{currentstroke}%
\pgfsetdash{}{0pt}%
\pgfpathmoveto{\pgfqpoint{1.989449in}{2.111244in}}%
\pgfpathlineto{\pgfqpoint{2.072782in}{2.111244in}}%
\pgfpathlineto{\pgfqpoint{2.072782in}{2.194577in}}%
\pgfpathlineto{\pgfqpoint{1.989449in}{2.194577in}}%
\pgfpathlineto{\pgfqpoint{1.989449in}{2.111244in}}%
\pgfpathclose%
\pgfusepath{stroke}%
\end{pgfscope}%
\begin{pgfscope}%
\pgfpathrectangle{\pgfqpoint{1.198472in}{0.823309in}}{\pgfqpoint{5.051528in}{3.826691in}}%
\pgfusepath{clip}%
\pgfsetbuttcap%
\pgfsetroundjoin%
\pgfsetlinewidth{1.505625pt}%
\definecolor{currentstroke}{rgb}{1.000000,0.498039,0.313725}%
\pgfsetstrokecolor{currentstroke}%
\pgfsetdash{}{0pt}%
\pgfpathmoveto{\pgfqpoint{2.035836in}{2.126534in}}%
\pgfpathlineto{\pgfqpoint{2.119169in}{2.126534in}}%
\pgfpathlineto{\pgfqpoint{2.119169in}{2.209868in}}%
\pgfpathlineto{\pgfqpoint{2.035836in}{2.209868in}}%
\pgfpathlineto{\pgfqpoint{2.035836in}{2.126534in}}%
\pgfpathclose%
\pgfusepath{stroke}%
\end{pgfscope}%
\begin{pgfscope}%
\pgfpathrectangle{\pgfqpoint{1.198472in}{0.823309in}}{\pgfqpoint{5.051528in}{3.826691in}}%
\pgfusepath{clip}%
\pgfsetbuttcap%
\pgfsetroundjoin%
\pgfsetlinewidth{1.505625pt}%
\definecolor{currentstroke}{rgb}{1.000000,0.498039,0.313725}%
\pgfsetstrokecolor{currentstroke}%
\pgfsetdash{}{0pt}%
\pgfpathmoveto{\pgfqpoint{2.082223in}{2.216710in}}%
\pgfpathlineto{\pgfqpoint{2.165556in}{2.216710in}}%
\pgfpathlineto{\pgfqpoint{2.165556in}{2.300043in}}%
\pgfpathlineto{\pgfqpoint{2.082223in}{2.300043in}}%
\pgfpathlineto{\pgfqpoint{2.082223in}{2.216710in}}%
\pgfpathclose%
\pgfusepath{stroke}%
\end{pgfscope}%
\begin{pgfscope}%
\pgfpathrectangle{\pgfqpoint{1.198472in}{0.823309in}}{\pgfqpoint{5.051528in}{3.826691in}}%
\pgfusepath{clip}%
\pgfsetbuttcap%
\pgfsetroundjoin%
\pgfsetlinewidth{1.505625pt}%
\definecolor{currentstroke}{rgb}{1.000000,0.498039,0.313725}%
\pgfsetstrokecolor{currentstroke}%
\pgfsetdash{}{0pt}%
\pgfpathmoveto{\pgfqpoint{2.128610in}{2.572106in}}%
\pgfpathlineto{\pgfqpoint{2.211943in}{2.572106in}}%
\pgfpathlineto{\pgfqpoint{2.211943in}{2.655439in}}%
\pgfpathlineto{\pgfqpoint{2.128610in}{2.655439in}}%
\pgfpathlineto{\pgfqpoint{2.128610in}{2.572106in}}%
\pgfpathclose%
\pgfusepath{stroke}%
\end{pgfscope}%
\begin{pgfscope}%
\pgfpathrectangle{\pgfqpoint{1.198472in}{0.823309in}}{\pgfqpoint{5.051528in}{3.826691in}}%
\pgfusepath{clip}%
\pgfsetbuttcap%
\pgfsetroundjoin%
\pgfsetlinewidth{1.505625pt}%
\definecolor{currentstroke}{rgb}{1.000000,0.498039,0.313725}%
\pgfsetstrokecolor{currentstroke}%
\pgfsetdash{}{0pt}%
\pgfpathmoveto{\pgfqpoint{2.174996in}{2.073127in}}%
\pgfpathlineto{\pgfqpoint{2.258330in}{2.073127in}}%
\pgfpathlineto{\pgfqpoint{2.258330in}{2.156460in}}%
\pgfpathlineto{\pgfqpoint{2.174996in}{2.156460in}}%
\pgfpathlineto{\pgfqpoint{2.174996in}{2.073127in}}%
\pgfpathclose%
\pgfusepath{stroke}%
\end{pgfscope}%
\begin{pgfscope}%
\pgfpathrectangle{\pgfqpoint{1.198472in}{0.823309in}}{\pgfqpoint{5.051528in}{3.826691in}}%
\pgfusepath{clip}%
\pgfsetbuttcap%
\pgfsetroundjoin%
\pgfsetlinewidth{1.505625pt}%
\definecolor{currentstroke}{rgb}{1.000000,0.498039,0.313725}%
\pgfsetstrokecolor{currentstroke}%
\pgfsetdash{}{0pt}%
\pgfpathmoveto{\pgfqpoint{2.221383in}{2.271534in}}%
\pgfpathlineto{\pgfqpoint{2.304717in}{2.271534in}}%
\pgfpathlineto{\pgfqpoint{2.304717in}{2.354867in}}%
\pgfpathlineto{\pgfqpoint{2.221383in}{2.354867in}}%
\pgfpathlineto{\pgfqpoint{2.221383in}{2.271534in}}%
\pgfpathclose%
\pgfusepath{stroke}%
\end{pgfscope}%
\begin{pgfscope}%
\pgfpathrectangle{\pgfqpoint{1.198472in}{0.823309in}}{\pgfqpoint{5.051528in}{3.826691in}}%
\pgfusepath{clip}%
\pgfsetbuttcap%
\pgfsetroundjoin%
\pgfsetlinewidth{1.505625pt}%
\definecolor{currentstroke}{rgb}{1.000000,0.498039,0.313725}%
\pgfsetstrokecolor{currentstroke}%
\pgfsetdash{}{0pt}%
\pgfpathmoveto{\pgfqpoint{2.267770in}{2.554575in}}%
\pgfpathlineto{\pgfqpoint{2.351103in}{2.554575in}}%
\pgfpathlineto{\pgfqpoint{2.351103in}{2.637908in}}%
\pgfpathlineto{\pgfqpoint{2.267770in}{2.637908in}}%
\pgfpathlineto{\pgfqpoint{2.267770in}{2.554575in}}%
\pgfpathclose%
\pgfusepath{stroke}%
\end{pgfscope}%
\begin{pgfscope}%
\pgfpathrectangle{\pgfqpoint{1.198472in}{0.823309in}}{\pgfqpoint{5.051528in}{3.826691in}}%
\pgfusepath{clip}%
\pgfsetbuttcap%
\pgfsetroundjoin%
\pgfsetlinewidth{1.505625pt}%
\definecolor{currentstroke}{rgb}{1.000000,0.498039,0.313725}%
\pgfsetstrokecolor{currentstroke}%
\pgfsetdash{}{0pt}%
\pgfpathmoveto{\pgfqpoint{2.314157in}{2.497493in}}%
\pgfpathlineto{\pgfqpoint{2.397490in}{2.497493in}}%
\pgfpathlineto{\pgfqpoint{2.397490in}{2.580826in}}%
\pgfpathlineto{\pgfqpoint{2.314157in}{2.580826in}}%
\pgfpathlineto{\pgfqpoint{2.314157in}{2.497493in}}%
\pgfpathclose%
\pgfusepath{stroke}%
\end{pgfscope}%
\begin{pgfscope}%
\pgfpathrectangle{\pgfqpoint{1.198472in}{0.823309in}}{\pgfqpoint{5.051528in}{3.826691in}}%
\pgfusepath{clip}%
\pgfsetbuttcap%
\pgfsetroundjoin%
\pgfsetlinewidth{1.505625pt}%
\definecolor{currentstroke}{rgb}{1.000000,0.498039,0.313725}%
\pgfsetstrokecolor{currentstroke}%
\pgfsetdash{}{0pt}%
\pgfpathmoveto{\pgfqpoint{2.360544in}{2.368845in}}%
\pgfpathlineto{\pgfqpoint{2.443877in}{2.368845in}}%
\pgfpathlineto{\pgfqpoint{2.443877in}{2.452178in}}%
\pgfpathlineto{\pgfqpoint{2.360544in}{2.452178in}}%
\pgfpathlineto{\pgfqpoint{2.360544in}{2.368845in}}%
\pgfpathclose%
\pgfusepath{stroke}%
\end{pgfscope}%
\begin{pgfscope}%
\pgfpathrectangle{\pgfqpoint{1.198472in}{0.823309in}}{\pgfqpoint{5.051528in}{3.826691in}}%
\pgfusepath{clip}%
\pgfsetbuttcap%
\pgfsetroundjoin%
\pgfsetlinewidth{1.505625pt}%
\definecolor{currentstroke}{rgb}{1.000000,0.498039,0.313725}%
\pgfsetstrokecolor{currentstroke}%
\pgfsetdash{}{0pt}%
\pgfpathmoveto{\pgfqpoint{2.406931in}{2.632284in}}%
\pgfpathlineto{\pgfqpoint{2.490264in}{2.632284in}}%
\pgfpathlineto{\pgfqpoint{2.490264in}{2.715617in}}%
\pgfpathlineto{\pgfqpoint{2.406931in}{2.715617in}}%
\pgfpathlineto{\pgfqpoint{2.406931in}{2.632284in}}%
\pgfpathclose%
\pgfusepath{stroke}%
\end{pgfscope}%
\begin{pgfscope}%
\pgfpathrectangle{\pgfqpoint{1.198472in}{0.823309in}}{\pgfqpoint{5.051528in}{3.826691in}}%
\pgfusepath{clip}%
\pgfsetbuttcap%
\pgfsetroundjoin%
\pgfsetlinewidth{1.505625pt}%
\definecolor{currentstroke}{rgb}{1.000000,0.498039,0.313725}%
\pgfsetstrokecolor{currentstroke}%
\pgfsetdash{}{0pt}%
\pgfpathmoveto{\pgfqpoint{2.453318in}{2.378121in}}%
\pgfpathlineto{\pgfqpoint{2.536651in}{2.378121in}}%
\pgfpathlineto{\pgfqpoint{2.536651in}{2.461455in}}%
\pgfpathlineto{\pgfqpoint{2.453318in}{2.461455in}}%
\pgfpathlineto{\pgfqpoint{2.453318in}{2.378121in}}%
\pgfpathclose%
\pgfusepath{stroke}%
\end{pgfscope}%
\begin{pgfscope}%
\pgfpathrectangle{\pgfqpoint{1.198472in}{0.823309in}}{\pgfqpoint{5.051528in}{3.826691in}}%
\pgfusepath{clip}%
\pgfsetbuttcap%
\pgfsetroundjoin%
\pgfsetlinewidth{1.505625pt}%
\definecolor{currentstroke}{rgb}{1.000000,0.498039,0.313725}%
\pgfsetstrokecolor{currentstroke}%
\pgfsetdash{}{0pt}%
\pgfpathmoveto{\pgfqpoint{2.499704in}{2.688854in}}%
\pgfpathlineto{\pgfqpoint{2.583038in}{2.688854in}}%
\pgfpathlineto{\pgfqpoint{2.583038in}{2.772187in}}%
\pgfpathlineto{\pgfqpoint{2.499704in}{2.772187in}}%
\pgfpathlineto{\pgfqpoint{2.499704in}{2.688854in}}%
\pgfpathclose%
\pgfusepath{stroke}%
\end{pgfscope}%
\begin{pgfscope}%
\pgfpathrectangle{\pgfqpoint{1.198472in}{0.823309in}}{\pgfqpoint{5.051528in}{3.826691in}}%
\pgfusepath{clip}%
\pgfsetbuttcap%
\pgfsetroundjoin%
\pgfsetlinewidth{1.505625pt}%
\definecolor{currentstroke}{rgb}{1.000000,0.498039,0.313725}%
\pgfsetstrokecolor{currentstroke}%
\pgfsetdash{}{0pt}%
\pgfpathmoveto{\pgfqpoint{2.546091in}{2.660893in}}%
\pgfpathlineto{\pgfqpoint{2.629425in}{2.660893in}}%
\pgfpathlineto{\pgfqpoint{2.629425in}{2.744226in}}%
\pgfpathlineto{\pgfqpoint{2.546091in}{2.744226in}}%
\pgfpathlineto{\pgfqpoint{2.546091in}{2.660893in}}%
\pgfpathclose%
\pgfusepath{stroke}%
\end{pgfscope}%
\begin{pgfscope}%
\pgfpathrectangle{\pgfqpoint{1.198472in}{0.823309in}}{\pgfqpoint{5.051528in}{3.826691in}}%
\pgfusepath{clip}%
\pgfsetbuttcap%
\pgfsetroundjoin%
\pgfsetlinewidth{1.505625pt}%
\definecolor{currentstroke}{rgb}{1.000000,0.498039,0.313725}%
\pgfsetstrokecolor{currentstroke}%
\pgfsetdash{}{0pt}%
\pgfpathmoveto{\pgfqpoint{2.592478in}{2.431771in}}%
\pgfpathlineto{\pgfqpoint{2.675811in}{2.431771in}}%
\pgfpathlineto{\pgfqpoint{2.675811in}{2.515104in}}%
\pgfpathlineto{\pgfqpoint{2.592478in}{2.515104in}}%
\pgfpathlineto{\pgfqpoint{2.592478in}{2.431771in}}%
\pgfpathclose%
\pgfusepath{stroke}%
\end{pgfscope}%
\begin{pgfscope}%
\pgfpathrectangle{\pgfqpoint{1.198472in}{0.823309in}}{\pgfqpoint{5.051528in}{3.826691in}}%
\pgfusepath{clip}%
\pgfsetbuttcap%
\pgfsetroundjoin%
\pgfsetlinewidth{1.505625pt}%
\definecolor{currentstroke}{rgb}{1.000000,0.498039,0.313725}%
\pgfsetstrokecolor{currentstroke}%
\pgfsetdash{}{0pt}%
\pgfpathmoveto{\pgfqpoint{2.638865in}{2.527950in}}%
\pgfpathlineto{\pgfqpoint{2.722198in}{2.527950in}}%
\pgfpathlineto{\pgfqpoint{2.722198in}{2.611283in}}%
\pgfpathlineto{\pgfqpoint{2.638865in}{2.611283in}}%
\pgfpathlineto{\pgfqpoint{2.638865in}{2.527950in}}%
\pgfpathclose%
\pgfusepath{stroke}%
\end{pgfscope}%
\begin{pgfscope}%
\pgfpathrectangle{\pgfqpoint{1.198472in}{0.823309in}}{\pgfqpoint{5.051528in}{3.826691in}}%
\pgfusepath{clip}%
\pgfsetbuttcap%
\pgfsetroundjoin%
\pgfsetlinewidth{1.505625pt}%
\definecolor{currentstroke}{rgb}{1.000000,0.498039,0.313725}%
\pgfsetstrokecolor{currentstroke}%
\pgfsetdash{}{0pt}%
\pgfpathmoveto{\pgfqpoint{2.685252in}{2.647237in}}%
\pgfpathlineto{\pgfqpoint{2.768585in}{2.647237in}}%
\pgfpathlineto{\pgfqpoint{2.768585in}{2.730570in}}%
\pgfpathlineto{\pgfqpoint{2.685252in}{2.730570in}}%
\pgfpathlineto{\pgfqpoint{2.685252in}{2.647237in}}%
\pgfpathclose%
\pgfusepath{stroke}%
\end{pgfscope}%
\begin{pgfscope}%
\pgfpathrectangle{\pgfqpoint{1.198472in}{0.823309in}}{\pgfqpoint{5.051528in}{3.826691in}}%
\pgfusepath{clip}%
\pgfsetbuttcap%
\pgfsetroundjoin%
\pgfsetlinewidth{1.505625pt}%
\definecolor{currentstroke}{rgb}{1.000000,0.498039,0.313725}%
\pgfsetstrokecolor{currentstroke}%
\pgfsetdash{}{0pt}%
\pgfpathmoveto{\pgfqpoint{2.731639in}{2.567736in}}%
\pgfpathlineto{\pgfqpoint{2.814972in}{2.567736in}}%
\pgfpathlineto{\pgfqpoint{2.814972in}{2.651069in}}%
\pgfpathlineto{\pgfqpoint{2.731639in}{2.651069in}}%
\pgfpathlineto{\pgfqpoint{2.731639in}{2.567736in}}%
\pgfpathclose%
\pgfusepath{stroke}%
\end{pgfscope}%
\begin{pgfscope}%
\pgfpathrectangle{\pgfqpoint{1.198472in}{0.823309in}}{\pgfqpoint{5.051528in}{3.826691in}}%
\pgfusepath{clip}%
\pgfsetbuttcap%
\pgfsetroundjoin%
\pgfsetlinewidth{1.505625pt}%
\definecolor{currentstroke}{rgb}{1.000000,0.498039,0.313725}%
\pgfsetstrokecolor{currentstroke}%
\pgfsetdash{}{0pt}%
\pgfpathmoveto{\pgfqpoint{2.778026in}{2.479409in}}%
\pgfpathlineto{\pgfqpoint{2.861359in}{2.479409in}}%
\pgfpathlineto{\pgfqpoint{2.861359in}{2.562742in}}%
\pgfpathlineto{\pgfqpoint{2.778026in}{2.562742in}}%
\pgfpathlineto{\pgfqpoint{2.778026in}{2.479409in}}%
\pgfpathclose%
\pgfusepath{stroke}%
\end{pgfscope}%
\begin{pgfscope}%
\pgfpathrectangle{\pgfqpoint{1.198472in}{0.823309in}}{\pgfqpoint{5.051528in}{3.826691in}}%
\pgfusepath{clip}%
\pgfsetbuttcap%
\pgfsetroundjoin%
\pgfsetlinewidth{1.505625pt}%
\definecolor{currentstroke}{rgb}{1.000000,0.498039,0.313725}%
\pgfsetstrokecolor{currentstroke}%
\pgfsetdash{}{0pt}%
\pgfpathmoveto{\pgfqpoint{2.824412in}{2.502381in}}%
\pgfpathlineto{\pgfqpoint{2.907746in}{2.502381in}}%
\pgfpathlineto{\pgfqpoint{2.907746in}{2.585714in}}%
\pgfpathlineto{\pgfqpoint{2.824412in}{2.585714in}}%
\pgfpathlineto{\pgfqpoint{2.824412in}{2.502381in}}%
\pgfpathclose%
\pgfusepath{stroke}%
\end{pgfscope}%
\begin{pgfscope}%
\pgfpathrectangle{\pgfqpoint{1.198472in}{0.823309in}}{\pgfqpoint{5.051528in}{3.826691in}}%
\pgfusepath{clip}%
\pgfsetbuttcap%
\pgfsetroundjoin%
\pgfsetlinewidth{1.505625pt}%
\definecolor{currentstroke}{rgb}{1.000000,0.498039,0.313725}%
\pgfsetstrokecolor{currentstroke}%
\pgfsetdash{}{0pt}%
\pgfpathmoveto{\pgfqpoint{2.870799in}{2.725659in}}%
\pgfpathlineto{\pgfqpoint{2.954133in}{2.725659in}}%
\pgfpathlineto{\pgfqpoint{2.954133in}{2.808993in}}%
\pgfpathlineto{\pgfqpoint{2.870799in}{2.808993in}}%
\pgfpathlineto{\pgfqpoint{2.870799in}{2.725659in}}%
\pgfpathclose%
\pgfusepath{stroke}%
\end{pgfscope}%
\begin{pgfscope}%
\pgfpathrectangle{\pgfqpoint{1.198472in}{0.823309in}}{\pgfqpoint{5.051528in}{3.826691in}}%
\pgfusepath{clip}%
\pgfsetbuttcap%
\pgfsetroundjoin%
\pgfsetlinewidth{1.505625pt}%
\definecolor{currentstroke}{rgb}{1.000000,0.498039,0.313725}%
\pgfsetstrokecolor{currentstroke}%
\pgfsetdash{}{0pt}%
\pgfpathmoveto{\pgfqpoint{2.917186in}{2.400710in}}%
\pgfpathlineto{\pgfqpoint{3.000519in}{2.400710in}}%
\pgfpathlineto{\pgfqpoint{3.000519in}{2.484044in}}%
\pgfpathlineto{\pgfqpoint{2.917186in}{2.484044in}}%
\pgfpathlineto{\pgfqpoint{2.917186in}{2.400710in}}%
\pgfpathclose%
\pgfusepath{stroke}%
\end{pgfscope}%
\begin{pgfscope}%
\pgfpathrectangle{\pgfqpoint{1.198472in}{0.823309in}}{\pgfqpoint{5.051528in}{3.826691in}}%
\pgfusepath{clip}%
\pgfsetbuttcap%
\pgfsetroundjoin%
\pgfsetlinewidth{1.505625pt}%
\definecolor{currentstroke}{rgb}{1.000000,0.498039,0.313725}%
\pgfsetstrokecolor{currentstroke}%
\pgfsetdash{}{0pt}%
\pgfpathmoveto{\pgfqpoint{2.963573in}{2.516529in}}%
\pgfpathlineto{\pgfqpoint{3.046906in}{2.516529in}}%
\pgfpathlineto{\pgfqpoint{3.046906in}{2.599862in}}%
\pgfpathlineto{\pgfqpoint{2.963573in}{2.599862in}}%
\pgfpathlineto{\pgfqpoint{2.963573in}{2.516529in}}%
\pgfpathclose%
\pgfusepath{stroke}%
\end{pgfscope}%
\begin{pgfscope}%
\pgfpathrectangle{\pgfqpoint{1.198472in}{0.823309in}}{\pgfqpoint{5.051528in}{3.826691in}}%
\pgfusepath{clip}%
\pgfsetbuttcap%
\pgfsetroundjoin%
\pgfsetlinewidth{1.505625pt}%
\definecolor{currentstroke}{rgb}{1.000000,0.498039,0.313725}%
\pgfsetstrokecolor{currentstroke}%
\pgfsetdash{}{0pt}%
\pgfpathmoveto{\pgfqpoint{3.009960in}{2.566533in}}%
\pgfpathlineto{\pgfqpoint{3.093293in}{2.566533in}}%
\pgfpathlineto{\pgfqpoint{3.093293in}{2.649866in}}%
\pgfpathlineto{\pgfqpoint{3.009960in}{2.649866in}}%
\pgfpathlineto{\pgfqpoint{3.009960in}{2.566533in}}%
\pgfpathclose%
\pgfusepath{stroke}%
\end{pgfscope}%
\begin{pgfscope}%
\pgfpathrectangle{\pgfqpoint{1.198472in}{0.823309in}}{\pgfqpoint{5.051528in}{3.826691in}}%
\pgfusepath{clip}%
\pgfsetbuttcap%
\pgfsetroundjoin%
\pgfsetlinewidth{1.505625pt}%
\definecolor{currentstroke}{rgb}{1.000000,0.498039,0.313725}%
\pgfsetstrokecolor{currentstroke}%
\pgfsetdash{}{0pt}%
\pgfpathmoveto{\pgfqpoint{3.056347in}{2.648466in}}%
\pgfpathlineto{\pgfqpoint{3.139680in}{2.648466in}}%
\pgfpathlineto{\pgfqpoint{3.139680in}{2.731799in}}%
\pgfpathlineto{\pgfqpoint{3.056347in}{2.731799in}}%
\pgfpathlineto{\pgfqpoint{3.056347in}{2.648466in}}%
\pgfpathclose%
\pgfusepath{stroke}%
\end{pgfscope}%
\begin{pgfscope}%
\pgfpathrectangle{\pgfqpoint{1.198472in}{0.823309in}}{\pgfqpoint{5.051528in}{3.826691in}}%
\pgfusepath{clip}%
\pgfsetbuttcap%
\pgfsetroundjoin%
\pgfsetlinewidth{1.505625pt}%
\definecolor{currentstroke}{rgb}{1.000000,0.498039,0.313725}%
\pgfsetstrokecolor{currentstroke}%
\pgfsetdash{}{0pt}%
\pgfpathmoveto{\pgfqpoint{3.102733in}{2.584834in}}%
\pgfpathlineto{\pgfqpoint{3.186067in}{2.584834in}}%
\pgfpathlineto{\pgfqpoint{3.186067in}{2.668167in}}%
\pgfpathlineto{\pgfqpoint{3.102733in}{2.668167in}}%
\pgfpathlineto{\pgfqpoint{3.102733in}{2.584834in}}%
\pgfpathclose%
\pgfusepath{stroke}%
\end{pgfscope}%
\begin{pgfscope}%
\pgfpathrectangle{\pgfqpoint{1.198472in}{0.823309in}}{\pgfqpoint{5.051528in}{3.826691in}}%
\pgfusepath{clip}%
\pgfsetbuttcap%
\pgfsetroundjoin%
\pgfsetlinewidth{1.505625pt}%
\definecolor{currentstroke}{rgb}{1.000000,0.498039,0.313725}%
\pgfsetstrokecolor{currentstroke}%
\pgfsetdash{}{0pt}%
\pgfpathmoveto{\pgfqpoint{3.149120in}{2.540781in}}%
\pgfpathlineto{\pgfqpoint{3.232454in}{2.540781in}}%
\pgfpathlineto{\pgfqpoint{3.232454in}{2.624114in}}%
\pgfpathlineto{\pgfqpoint{3.149120in}{2.624114in}}%
\pgfpathlineto{\pgfqpoint{3.149120in}{2.540781in}}%
\pgfpathclose%
\pgfusepath{stroke}%
\end{pgfscope}%
\begin{pgfscope}%
\pgfpathrectangle{\pgfqpoint{1.198472in}{0.823309in}}{\pgfqpoint{5.051528in}{3.826691in}}%
\pgfusepath{clip}%
\pgfsetbuttcap%
\pgfsetroundjoin%
\pgfsetlinewidth{1.505625pt}%
\definecolor{currentstroke}{rgb}{1.000000,0.498039,0.313725}%
\pgfsetstrokecolor{currentstroke}%
\pgfsetdash{}{0pt}%
\pgfpathmoveto{\pgfqpoint{3.195507in}{2.596328in}}%
\pgfpathlineto{\pgfqpoint{3.278841in}{2.596328in}}%
\pgfpathlineto{\pgfqpoint{3.278841in}{2.679661in}}%
\pgfpathlineto{\pgfqpoint{3.195507in}{2.679661in}}%
\pgfpathlineto{\pgfqpoint{3.195507in}{2.596328in}}%
\pgfpathclose%
\pgfusepath{stroke}%
\end{pgfscope}%
\begin{pgfscope}%
\pgfpathrectangle{\pgfqpoint{1.198472in}{0.823309in}}{\pgfqpoint{5.051528in}{3.826691in}}%
\pgfusepath{clip}%
\pgfsetbuttcap%
\pgfsetroundjoin%
\pgfsetlinewidth{1.505625pt}%
\definecolor{currentstroke}{rgb}{1.000000,0.498039,0.313725}%
\pgfsetstrokecolor{currentstroke}%
\pgfsetdash{}{0pt}%
\pgfpathmoveto{\pgfqpoint{3.241894in}{2.517221in}}%
\pgfpathlineto{\pgfqpoint{3.325227in}{2.517221in}}%
\pgfpathlineto{\pgfqpoint{3.325227in}{2.600554in}}%
\pgfpathlineto{\pgfqpoint{3.241894in}{2.600554in}}%
\pgfpathlineto{\pgfqpoint{3.241894in}{2.517221in}}%
\pgfpathclose%
\pgfusepath{stroke}%
\end{pgfscope}%
\begin{pgfscope}%
\pgfpathrectangle{\pgfqpoint{1.198472in}{0.823309in}}{\pgfqpoint{5.051528in}{3.826691in}}%
\pgfusepath{clip}%
\pgfsetbuttcap%
\pgfsetroundjoin%
\pgfsetlinewidth{1.505625pt}%
\definecolor{currentstroke}{rgb}{1.000000,0.498039,0.313725}%
\pgfsetstrokecolor{currentstroke}%
\pgfsetdash{}{0pt}%
\pgfpathmoveto{\pgfqpoint{3.288281in}{2.441543in}}%
\pgfpathlineto{\pgfqpoint{3.371614in}{2.441543in}}%
\pgfpathlineto{\pgfqpoint{3.371614in}{2.524876in}}%
\pgfpathlineto{\pgfqpoint{3.288281in}{2.524876in}}%
\pgfpathlineto{\pgfqpoint{3.288281in}{2.441543in}}%
\pgfpathclose%
\pgfusepath{stroke}%
\end{pgfscope}%
\begin{pgfscope}%
\pgfpathrectangle{\pgfqpoint{1.198472in}{0.823309in}}{\pgfqpoint{5.051528in}{3.826691in}}%
\pgfusepath{clip}%
\pgfsetbuttcap%
\pgfsetroundjoin%
\pgfsetlinewidth{1.505625pt}%
\definecolor{currentstroke}{rgb}{1.000000,0.498039,0.313725}%
\pgfsetstrokecolor{currentstroke}%
\pgfsetdash{}{0pt}%
\pgfpathmoveto{\pgfqpoint{3.334668in}{2.626153in}}%
\pgfpathlineto{\pgfqpoint{3.418001in}{2.626153in}}%
\pgfpathlineto{\pgfqpoint{3.418001in}{2.709486in}}%
\pgfpathlineto{\pgfqpoint{3.334668in}{2.709486in}}%
\pgfpathlineto{\pgfqpoint{3.334668in}{2.626153in}}%
\pgfpathclose%
\pgfusepath{stroke}%
\end{pgfscope}%
\begin{pgfscope}%
\pgfpathrectangle{\pgfqpoint{1.198472in}{0.823309in}}{\pgfqpoint{5.051528in}{3.826691in}}%
\pgfusepath{clip}%
\pgfsetbuttcap%
\pgfsetroundjoin%
\pgfsetlinewidth{1.505625pt}%
\definecolor{currentstroke}{rgb}{1.000000,0.498039,0.313725}%
\pgfsetstrokecolor{currentstroke}%
\pgfsetdash{}{0pt}%
\pgfpathmoveto{\pgfqpoint{3.381055in}{2.537068in}}%
\pgfpathlineto{\pgfqpoint{3.464388in}{2.537068in}}%
\pgfpathlineto{\pgfqpoint{3.464388in}{2.620401in}}%
\pgfpathlineto{\pgfqpoint{3.381055in}{2.620401in}}%
\pgfpathlineto{\pgfqpoint{3.381055in}{2.537068in}}%
\pgfpathclose%
\pgfusepath{stroke}%
\end{pgfscope}%
\begin{pgfscope}%
\pgfpathrectangle{\pgfqpoint{1.198472in}{0.823309in}}{\pgfqpoint{5.051528in}{3.826691in}}%
\pgfusepath{clip}%
\pgfsetbuttcap%
\pgfsetroundjoin%
\pgfsetlinewidth{1.505625pt}%
\definecolor{currentstroke}{rgb}{1.000000,0.498039,0.313725}%
\pgfsetstrokecolor{currentstroke}%
\pgfsetdash{}{0pt}%
\pgfpathmoveto{\pgfqpoint{3.427441in}{2.465931in}}%
\pgfpathlineto{\pgfqpoint{3.510775in}{2.465931in}}%
\pgfpathlineto{\pgfqpoint{3.510775in}{2.549264in}}%
\pgfpathlineto{\pgfqpoint{3.427441in}{2.549264in}}%
\pgfpathlineto{\pgfqpoint{3.427441in}{2.465931in}}%
\pgfpathclose%
\pgfusepath{stroke}%
\end{pgfscope}%
\begin{pgfscope}%
\pgfpathrectangle{\pgfqpoint{1.198472in}{0.823309in}}{\pgfqpoint{5.051528in}{3.826691in}}%
\pgfusepath{clip}%
\pgfsetbuttcap%
\pgfsetroundjoin%
\pgfsetlinewidth{1.505625pt}%
\definecolor{currentstroke}{rgb}{1.000000,0.498039,0.313725}%
\pgfsetstrokecolor{currentstroke}%
\pgfsetdash{}{0pt}%
\pgfpathmoveto{\pgfqpoint{3.473828in}{2.511180in}}%
\pgfpathlineto{\pgfqpoint{3.557162in}{2.511180in}}%
\pgfpathlineto{\pgfqpoint{3.557162in}{2.594513in}}%
\pgfpathlineto{\pgfqpoint{3.473828in}{2.594513in}}%
\pgfpathlineto{\pgfqpoint{3.473828in}{2.511180in}}%
\pgfpathclose%
\pgfusepath{stroke}%
\end{pgfscope}%
\begin{pgfscope}%
\pgfpathrectangle{\pgfqpoint{1.198472in}{0.823309in}}{\pgfqpoint{5.051528in}{3.826691in}}%
\pgfusepath{clip}%
\pgfsetbuttcap%
\pgfsetroundjoin%
\pgfsetlinewidth{1.505625pt}%
\definecolor{currentstroke}{rgb}{1.000000,0.498039,0.313725}%
\pgfsetstrokecolor{currentstroke}%
\pgfsetdash{}{0pt}%
\pgfpathmoveto{\pgfqpoint{3.520215in}{2.476196in}}%
\pgfpathlineto{\pgfqpoint{3.603549in}{2.476196in}}%
\pgfpathlineto{\pgfqpoint{3.603549in}{2.559529in}}%
\pgfpathlineto{\pgfqpoint{3.520215in}{2.559529in}}%
\pgfpathlineto{\pgfqpoint{3.520215in}{2.476196in}}%
\pgfpathclose%
\pgfusepath{stroke}%
\end{pgfscope}%
\begin{pgfscope}%
\pgfpathrectangle{\pgfqpoint{1.198472in}{0.823309in}}{\pgfqpoint{5.051528in}{3.826691in}}%
\pgfusepath{clip}%
\pgfsetbuttcap%
\pgfsetroundjoin%
\pgfsetlinewidth{1.505625pt}%
\definecolor{currentstroke}{rgb}{1.000000,0.498039,0.313725}%
\pgfsetstrokecolor{currentstroke}%
\pgfsetdash{}{0pt}%
\pgfpathmoveto{\pgfqpoint{3.566602in}{2.496484in}}%
\pgfpathlineto{\pgfqpoint{3.649935in}{2.496484in}}%
\pgfpathlineto{\pgfqpoint{3.649935in}{2.579817in}}%
\pgfpathlineto{\pgfqpoint{3.566602in}{2.579817in}}%
\pgfpathlineto{\pgfqpoint{3.566602in}{2.496484in}}%
\pgfpathclose%
\pgfusepath{stroke}%
\end{pgfscope}%
\begin{pgfscope}%
\pgfpathrectangle{\pgfqpoint{1.198472in}{0.823309in}}{\pgfqpoint{5.051528in}{3.826691in}}%
\pgfusepath{clip}%
\pgfsetbuttcap%
\pgfsetroundjoin%
\pgfsetlinewidth{1.505625pt}%
\definecolor{currentstroke}{rgb}{1.000000,0.498039,0.313725}%
\pgfsetstrokecolor{currentstroke}%
\pgfsetdash{}{0pt}%
\pgfpathmoveto{\pgfqpoint{3.612989in}{2.476438in}}%
\pgfpathlineto{\pgfqpoint{3.696322in}{2.476438in}}%
\pgfpathlineto{\pgfqpoint{3.696322in}{2.559772in}}%
\pgfpathlineto{\pgfqpoint{3.612989in}{2.559772in}}%
\pgfpathlineto{\pgfqpoint{3.612989in}{2.476438in}}%
\pgfpathclose%
\pgfusepath{stroke}%
\end{pgfscope}%
\begin{pgfscope}%
\pgfpathrectangle{\pgfqpoint{1.198472in}{0.823309in}}{\pgfqpoint{5.051528in}{3.826691in}}%
\pgfusepath{clip}%
\pgfsetbuttcap%
\pgfsetroundjoin%
\pgfsetlinewidth{1.505625pt}%
\definecolor{currentstroke}{rgb}{1.000000,0.498039,0.313725}%
\pgfsetstrokecolor{currentstroke}%
\pgfsetdash{}{0pt}%
\pgfpathmoveto{\pgfqpoint{3.659376in}{2.522125in}}%
\pgfpathlineto{\pgfqpoint{3.742709in}{2.522125in}}%
\pgfpathlineto{\pgfqpoint{3.742709in}{2.605458in}}%
\pgfpathlineto{\pgfqpoint{3.659376in}{2.605458in}}%
\pgfpathlineto{\pgfqpoint{3.659376in}{2.522125in}}%
\pgfpathclose%
\pgfusepath{stroke}%
\end{pgfscope}%
\begin{pgfscope}%
\pgfpathrectangle{\pgfqpoint{1.198472in}{0.823309in}}{\pgfqpoint{5.051528in}{3.826691in}}%
\pgfusepath{clip}%
\pgfsetbuttcap%
\pgfsetroundjoin%
\pgfsetlinewidth{1.505625pt}%
\definecolor{currentstroke}{rgb}{1.000000,0.498039,0.313725}%
\pgfsetstrokecolor{currentstroke}%
\pgfsetdash{}{0pt}%
\pgfpathmoveto{\pgfqpoint{3.705763in}{2.361061in}}%
\pgfpathlineto{\pgfqpoint{3.789096in}{2.361061in}}%
\pgfpathlineto{\pgfqpoint{3.789096in}{2.444394in}}%
\pgfpathlineto{\pgfqpoint{3.705763in}{2.444394in}}%
\pgfpathlineto{\pgfqpoint{3.705763in}{2.361061in}}%
\pgfpathclose%
\pgfusepath{stroke}%
\end{pgfscope}%
\begin{pgfscope}%
\pgfpathrectangle{\pgfqpoint{1.198472in}{0.823309in}}{\pgfqpoint{5.051528in}{3.826691in}}%
\pgfusepath{clip}%
\pgfsetbuttcap%
\pgfsetroundjoin%
\pgfsetlinewidth{1.505625pt}%
\definecolor{currentstroke}{rgb}{1.000000,0.498039,0.313725}%
\pgfsetstrokecolor{currentstroke}%
\pgfsetdash{}{0pt}%
\pgfpathmoveto{\pgfqpoint{3.752149in}{2.457722in}}%
\pgfpathlineto{\pgfqpoint{3.835483in}{2.457722in}}%
\pgfpathlineto{\pgfqpoint{3.835483in}{2.541056in}}%
\pgfpathlineto{\pgfqpoint{3.752149in}{2.541056in}}%
\pgfpathlineto{\pgfqpoint{3.752149in}{2.457722in}}%
\pgfpathclose%
\pgfusepath{stroke}%
\end{pgfscope}%
\begin{pgfscope}%
\pgfpathrectangle{\pgfqpoint{1.198472in}{0.823309in}}{\pgfqpoint{5.051528in}{3.826691in}}%
\pgfusepath{clip}%
\pgfsetbuttcap%
\pgfsetroundjoin%
\pgfsetlinewidth{1.505625pt}%
\definecolor{currentstroke}{rgb}{1.000000,0.498039,0.313725}%
\pgfsetstrokecolor{currentstroke}%
\pgfsetdash{}{0pt}%
\pgfpathmoveto{\pgfqpoint{3.798536in}{2.349293in}}%
\pgfpathlineto{\pgfqpoint{3.881870in}{2.349293in}}%
\pgfpathlineto{\pgfqpoint{3.881870in}{2.432626in}}%
\pgfpathlineto{\pgfqpoint{3.798536in}{2.432626in}}%
\pgfpathlineto{\pgfqpoint{3.798536in}{2.349293in}}%
\pgfpathclose%
\pgfusepath{stroke}%
\end{pgfscope}%
\begin{pgfscope}%
\pgfpathrectangle{\pgfqpoint{1.198472in}{0.823309in}}{\pgfqpoint{5.051528in}{3.826691in}}%
\pgfusepath{clip}%
\pgfsetbuttcap%
\pgfsetroundjoin%
\pgfsetlinewidth{1.505625pt}%
\definecolor{currentstroke}{rgb}{1.000000,0.498039,0.313725}%
\pgfsetstrokecolor{currentstroke}%
\pgfsetdash{}{0pt}%
\pgfpathmoveto{\pgfqpoint{3.844923in}{2.438642in}}%
\pgfpathlineto{\pgfqpoint{3.928256in}{2.438642in}}%
\pgfpathlineto{\pgfqpoint{3.928256in}{2.521975in}}%
\pgfpathlineto{\pgfqpoint{3.844923in}{2.521975in}}%
\pgfpathlineto{\pgfqpoint{3.844923in}{2.438642in}}%
\pgfpathclose%
\pgfusepath{stroke}%
\end{pgfscope}%
\begin{pgfscope}%
\pgfpathrectangle{\pgfqpoint{1.198472in}{0.823309in}}{\pgfqpoint{5.051528in}{3.826691in}}%
\pgfusepath{clip}%
\pgfsetbuttcap%
\pgfsetroundjoin%
\pgfsetlinewidth{1.505625pt}%
\definecolor{currentstroke}{rgb}{1.000000,0.498039,0.313725}%
\pgfsetstrokecolor{currentstroke}%
\pgfsetdash{}{0pt}%
\pgfpathmoveto{\pgfqpoint{3.891310in}{2.341711in}}%
\pgfpathlineto{\pgfqpoint{3.974643in}{2.341711in}}%
\pgfpathlineto{\pgfqpoint{3.974643in}{2.425045in}}%
\pgfpathlineto{\pgfqpoint{3.891310in}{2.425045in}}%
\pgfpathlineto{\pgfqpoint{3.891310in}{2.341711in}}%
\pgfpathclose%
\pgfusepath{stroke}%
\end{pgfscope}%
\begin{pgfscope}%
\pgfpathrectangle{\pgfqpoint{1.198472in}{0.823309in}}{\pgfqpoint{5.051528in}{3.826691in}}%
\pgfusepath{clip}%
\pgfsetbuttcap%
\pgfsetroundjoin%
\pgfsetlinewidth{1.505625pt}%
\definecolor{currentstroke}{rgb}{1.000000,0.498039,0.313725}%
\pgfsetstrokecolor{currentstroke}%
\pgfsetdash{}{0pt}%
\pgfpathmoveto{\pgfqpoint{3.937697in}{2.400961in}}%
\pgfpathlineto{\pgfqpoint{4.021030in}{2.400961in}}%
\pgfpathlineto{\pgfqpoint{4.021030in}{2.484294in}}%
\pgfpathlineto{\pgfqpoint{3.937697in}{2.484294in}}%
\pgfpathlineto{\pgfqpoint{3.937697in}{2.400961in}}%
\pgfpathclose%
\pgfusepath{stroke}%
\end{pgfscope}%
\begin{pgfscope}%
\pgfpathrectangle{\pgfqpoint{1.198472in}{0.823309in}}{\pgfqpoint{5.051528in}{3.826691in}}%
\pgfusepath{clip}%
\pgfsetbuttcap%
\pgfsetroundjoin%
\pgfsetlinewidth{1.505625pt}%
\definecolor{currentstroke}{rgb}{1.000000,0.498039,0.313725}%
\pgfsetstrokecolor{currentstroke}%
\pgfsetdash{}{0pt}%
\pgfpathmoveto{\pgfqpoint{3.984084in}{2.469550in}}%
\pgfpathlineto{\pgfqpoint{4.067417in}{2.469550in}}%
\pgfpathlineto{\pgfqpoint{4.067417in}{2.552884in}}%
\pgfpathlineto{\pgfqpoint{3.984084in}{2.552884in}}%
\pgfpathlineto{\pgfqpoint{3.984084in}{2.469550in}}%
\pgfpathclose%
\pgfusepath{stroke}%
\end{pgfscope}%
\begin{pgfscope}%
\pgfpathrectangle{\pgfqpoint{1.198472in}{0.823309in}}{\pgfqpoint{5.051528in}{3.826691in}}%
\pgfusepath{clip}%
\pgfsetbuttcap%
\pgfsetroundjoin%
\pgfsetlinewidth{1.505625pt}%
\definecolor{currentstroke}{rgb}{1.000000,0.498039,0.313725}%
\pgfsetstrokecolor{currentstroke}%
\pgfsetdash{}{0pt}%
\pgfpathmoveto{\pgfqpoint{4.030471in}{2.456590in}}%
\pgfpathlineto{\pgfqpoint{4.113804in}{2.456590in}}%
\pgfpathlineto{\pgfqpoint{4.113804in}{2.539923in}}%
\pgfpathlineto{\pgfqpoint{4.030471in}{2.539923in}}%
\pgfpathlineto{\pgfqpoint{4.030471in}{2.456590in}}%
\pgfpathclose%
\pgfusepath{stroke}%
\end{pgfscope}%
\begin{pgfscope}%
\pgfpathrectangle{\pgfqpoint{1.198472in}{0.823309in}}{\pgfqpoint{5.051528in}{3.826691in}}%
\pgfusepath{clip}%
\pgfsetbuttcap%
\pgfsetroundjoin%
\pgfsetlinewidth{1.505625pt}%
\definecolor{currentstroke}{rgb}{1.000000,0.498039,0.313725}%
\pgfsetstrokecolor{currentstroke}%
\pgfsetdash{}{0pt}%
\pgfpathmoveto{\pgfqpoint{4.076857in}{2.366261in}}%
\pgfpathlineto{\pgfqpoint{4.160191in}{2.366261in}}%
\pgfpathlineto{\pgfqpoint{4.160191in}{2.449594in}}%
\pgfpathlineto{\pgfqpoint{4.076857in}{2.449594in}}%
\pgfpathlineto{\pgfqpoint{4.076857in}{2.366261in}}%
\pgfpathclose%
\pgfusepath{stroke}%
\end{pgfscope}%
\begin{pgfscope}%
\pgfpathrectangle{\pgfqpoint{1.198472in}{0.823309in}}{\pgfqpoint{5.051528in}{3.826691in}}%
\pgfusepath{clip}%
\pgfsetbuttcap%
\pgfsetroundjoin%
\pgfsetlinewidth{1.505625pt}%
\definecolor{currentstroke}{rgb}{1.000000,0.498039,0.313725}%
\pgfsetstrokecolor{currentstroke}%
\pgfsetdash{}{0pt}%
\pgfpathmoveto{\pgfqpoint{4.123244in}{2.245801in}}%
\pgfpathlineto{\pgfqpoint{4.206578in}{2.245801in}}%
\pgfpathlineto{\pgfqpoint{4.206578in}{2.329134in}}%
\pgfpathlineto{\pgfqpoint{4.123244in}{2.329134in}}%
\pgfpathlineto{\pgfqpoint{4.123244in}{2.245801in}}%
\pgfpathclose%
\pgfusepath{stroke}%
\end{pgfscope}%
\begin{pgfscope}%
\pgfpathrectangle{\pgfqpoint{1.198472in}{0.823309in}}{\pgfqpoint{5.051528in}{3.826691in}}%
\pgfusepath{clip}%
\pgfsetbuttcap%
\pgfsetroundjoin%
\pgfsetlinewidth{1.505625pt}%
\definecolor{currentstroke}{rgb}{1.000000,0.498039,0.313725}%
\pgfsetstrokecolor{currentstroke}%
\pgfsetdash{}{0pt}%
\pgfpathmoveto{\pgfqpoint{4.169631in}{2.344665in}}%
\pgfpathlineto{\pgfqpoint{4.252964in}{2.344665in}}%
\pgfpathlineto{\pgfqpoint{4.252964in}{2.427999in}}%
\pgfpathlineto{\pgfqpoint{4.169631in}{2.427999in}}%
\pgfpathlineto{\pgfqpoint{4.169631in}{2.344665in}}%
\pgfpathclose%
\pgfusepath{stroke}%
\end{pgfscope}%
\begin{pgfscope}%
\pgfpathrectangle{\pgfqpoint{1.198472in}{0.823309in}}{\pgfqpoint{5.051528in}{3.826691in}}%
\pgfusepath{clip}%
\pgfsetbuttcap%
\pgfsetroundjoin%
\pgfsetlinewidth{1.505625pt}%
\definecolor{currentstroke}{rgb}{1.000000,0.498039,0.313725}%
\pgfsetstrokecolor{currentstroke}%
\pgfsetdash{}{0pt}%
\pgfpathmoveto{\pgfqpoint{4.216018in}{2.338035in}}%
\pgfpathlineto{\pgfqpoint{4.299351in}{2.338035in}}%
\pgfpathlineto{\pgfqpoint{4.299351in}{2.421368in}}%
\pgfpathlineto{\pgfqpoint{4.216018in}{2.421368in}}%
\pgfpathlineto{\pgfqpoint{4.216018in}{2.338035in}}%
\pgfpathclose%
\pgfusepath{stroke}%
\end{pgfscope}%
\begin{pgfscope}%
\pgfpathrectangle{\pgfqpoint{1.198472in}{0.823309in}}{\pgfqpoint{5.051528in}{3.826691in}}%
\pgfusepath{clip}%
\pgfsetbuttcap%
\pgfsetroundjoin%
\pgfsetlinewidth{1.505625pt}%
\definecolor{currentstroke}{rgb}{1.000000,0.498039,0.313725}%
\pgfsetstrokecolor{currentstroke}%
\pgfsetdash{}{0pt}%
\pgfpathmoveto{\pgfqpoint{4.262405in}{2.289958in}}%
\pgfpathlineto{\pgfqpoint{4.345738in}{2.289958in}}%
\pgfpathlineto{\pgfqpoint{4.345738in}{2.373292in}}%
\pgfpathlineto{\pgfqpoint{4.262405in}{2.373292in}}%
\pgfpathlineto{\pgfqpoint{4.262405in}{2.289958in}}%
\pgfpathclose%
\pgfusepath{stroke}%
\end{pgfscope}%
\begin{pgfscope}%
\pgfpathrectangle{\pgfqpoint{1.198472in}{0.823309in}}{\pgfqpoint{5.051528in}{3.826691in}}%
\pgfusepath{clip}%
\pgfsetbuttcap%
\pgfsetroundjoin%
\pgfsetlinewidth{1.505625pt}%
\definecolor{currentstroke}{rgb}{1.000000,0.498039,0.313725}%
\pgfsetstrokecolor{currentstroke}%
\pgfsetdash{}{0pt}%
\pgfpathmoveto{\pgfqpoint{4.308792in}{2.421362in}}%
\pgfpathlineto{\pgfqpoint{4.392125in}{2.421362in}}%
\pgfpathlineto{\pgfqpoint{4.392125in}{2.504695in}}%
\pgfpathlineto{\pgfqpoint{4.308792in}{2.504695in}}%
\pgfpathlineto{\pgfqpoint{4.308792in}{2.421362in}}%
\pgfpathclose%
\pgfusepath{stroke}%
\end{pgfscope}%
\begin{pgfscope}%
\pgfpathrectangle{\pgfqpoint{1.198472in}{0.823309in}}{\pgfqpoint{5.051528in}{3.826691in}}%
\pgfusepath{clip}%
\pgfsetbuttcap%
\pgfsetroundjoin%
\pgfsetlinewidth{1.505625pt}%
\definecolor{currentstroke}{rgb}{1.000000,0.498039,0.313725}%
\pgfsetstrokecolor{currentstroke}%
\pgfsetdash{}{0pt}%
\pgfpathmoveto{\pgfqpoint{4.355179in}{2.487999in}}%
\pgfpathlineto{\pgfqpoint{4.438512in}{2.487999in}}%
\pgfpathlineto{\pgfqpoint{4.438512in}{2.571332in}}%
\pgfpathlineto{\pgfqpoint{4.355179in}{2.571332in}}%
\pgfpathlineto{\pgfqpoint{4.355179in}{2.487999in}}%
\pgfpathclose%
\pgfusepath{stroke}%
\end{pgfscope}%
\begin{pgfscope}%
\pgfpathrectangle{\pgfqpoint{1.198472in}{0.823309in}}{\pgfqpoint{5.051528in}{3.826691in}}%
\pgfusepath{clip}%
\pgfsetbuttcap%
\pgfsetroundjoin%
\pgfsetlinewidth{1.505625pt}%
\definecolor{currentstroke}{rgb}{1.000000,0.498039,0.313725}%
\pgfsetstrokecolor{currentstroke}%
\pgfsetdash{}{0pt}%
\pgfpathmoveto{\pgfqpoint{4.401565in}{2.350131in}}%
\pgfpathlineto{\pgfqpoint{4.484899in}{2.350131in}}%
\pgfpathlineto{\pgfqpoint{4.484899in}{2.433465in}}%
\pgfpathlineto{\pgfqpoint{4.401565in}{2.433465in}}%
\pgfpathlineto{\pgfqpoint{4.401565in}{2.350131in}}%
\pgfpathclose%
\pgfusepath{stroke}%
\end{pgfscope}%
\begin{pgfscope}%
\pgfpathrectangle{\pgfqpoint{1.198472in}{0.823309in}}{\pgfqpoint{5.051528in}{3.826691in}}%
\pgfusepath{clip}%
\pgfsetbuttcap%
\pgfsetroundjoin%
\pgfsetlinewidth{1.505625pt}%
\definecolor{currentstroke}{rgb}{1.000000,0.498039,0.313725}%
\pgfsetstrokecolor{currentstroke}%
\pgfsetdash{}{0pt}%
\pgfpathmoveto{\pgfqpoint{4.447952in}{2.333919in}}%
\pgfpathlineto{\pgfqpoint{4.531286in}{2.333919in}}%
\pgfpathlineto{\pgfqpoint{4.531286in}{2.417252in}}%
\pgfpathlineto{\pgfqpoint{4.447952in}{2.417252in}}%
\pgfpathlineto{\pgfqpoint{4.447952in}{2.333919in}}%
\pgfpathclose%
\pgfusepath{stroke}%
\end{pgfscope}%
\begin{pgfscope}%
\pgfpathrectangle{\pgfqpoint{1.198472in}{0.823309in}}{\pgfqpoint{5.051528in}{3.826691in}}%
\pgfusepath{clip}%
\pgfsetbuttcap%
\pgfsetroundjoin%
\pgfsetlinewidth{1.505625pt}%
\definecolor{currentstroke}{rgb}{1.000000,0.498039,0.313725}%
\pgfsetstrokecolor{currentstroke}%
\pgfsetdash{}{0pt}%
\pgfpathmoveto{\pgfqpoint{4.494339in}{2.379595in}}%
\pgfpathlineto{\pgfqpoint{4.577672in}{2.379595in}}%
\pgfpathlineto{\pgfqpoint{4.577672in}{2.462928in}}%
\pgfpathlineto{\pgfqpoint{4.494339in}{2.462928in}}%
\pgfpathlineto{\pgfqpoint{4.494339in}{2.379595in}}%
\pgfpathclose%
\pgfusepath{stroke}%
\end{pgfscope}%
\begin{pgfscope}%
\pgfpathrectangle{\pgfqpoint{1.198472in}{0.823309in}}{\pgfqpoint{5.051528in}{3.826691in}}%
\pgfusepath{clip}%
\pgfsetbuttcap%
\pgfsetroundjoin%
\pgfsetlinewidth{1.505625pt}%
\definecolor{currentstroke}{rgb}{1.000000,0.498039,0.313725}%
\pgfsetstrokecolor{currentstroke}%
\pgfsetdash{}{0pt}%
\pgfpathmoveto{\pgfqpoint{4.540726in}{2.303099in}}%
\pgfpathlineto{\pgfqpoint{4.624059in}{2.303099in}}%
\pgfpathlineto{\pgfqpoint{4.624059in}{2.386432in}}%
\pgfpathlineto{\pgfqpoint{4.540726in}{2.386432in}}%
\pgfpathlineto{\pgfqpoint{4.540726in}{2.303099in}}%
\pgfpathclose%
\pgfusepath{stroke}%
\end{pgfscope}%
\begin{pgfscope}%
\pgfpathrectangle{\pgfqpoint{1.198472in}{0.823309in}}{\pgfqpoint{5.051528in}{3.826691in}}%
\pgfusepath{clip}%
\pgfsetbuttcap%
\pgfsetroundjoin%
\pgfsetlinewidth{1.505625pt}%
\definecolor{currentstroke}{rgb}{1.000000,0.498039,0.313725}%
\pgfsetstrokecolor{currentstroke}%
\pgfsetdash{}{0pt}%
\pgfpathmoveto{\pgfqpoint{4.587113in}{2.316548in}}%
\pgfpathlineto{\pgfqpoint{4.670446in}{2.316548in}}%
\pgfpathlineto{\pgfqpoint{4.670446in}{2.399881in}}%
\pgfpathlineto{\pgfqpoint{4.587113in}{2.399881in}}%
\pgfpathlineto{\pgfqpoint{4.587113in}{2.316548in}}%
\pgfpathclose%
\pgfusepath{stroke}%
\end{pgfscope}%
\begin{pgfscope}%
\pgfpathrectangle{\pgfqpoint{1.198472in}{0.823309in}}{\pgfqpoint{5.051528in}{3.826691in}}%
\pgfusepath{clip}%
\pgfsetbuttcap%
\pgfsetroundjoin%
\pgfsetlinewidth{1.505625pt}%
\definecolor{currentstroke}{rgb}{1.000000,0.498039,0.313725}%
\pgfsetstrokecolor{currentstroke}%
\pgfsetdash{}{0pt}%
\pgfpathmoveto{\pgfqpoint{4.633500in}{2.264833in}}%
\pgfpathlineto{\pgfqpoint{4.716833in}{2.264833in}}%
\pgfpathlineto{\pgfqpoint{4.716833in}{2.348166in}}%
\pgfpathlineto{\pgfqpoint{4.633500in}{2.348166in}}%
\pgfpathlineto{\pgfqpoint{4.633500in}{2.264833in}}%
\pgfpathclose%
\pgfusepath{stroke}%
\end{pgfscope}%
\begin{pgfscope}%
\pgfpathrectangle{\pgfqpoint{1.198472in}{0.823309in}}{\pgfqpoint{5.051528in}{3.826691in}}%
\pgfusepath{clip}%
\pgfsetbuttcap%
\pgfsetroundjoin%
\pgfsetlinewidth{1.505625pt}%
\definecolor{currentstroke}{rgb}{1.000000,0.498039,0.313725}%
\pgfsetstrokecolor{currentstroke}%
\pgfsetdash{}{0pt}%
\pgfpathmoveto{\pgfqpoint{4.679887in}{2.345515in}}%
\pgfpathlineto{\pgfqpoint{4.763220in}{2.345515in}}%
\pgfpathlineto{\pgfqpoint{4.763220in}{2.428849in}}%
\pgfpathlineto{\pgfqpoint{4.679887in}{2.428849in}}%
\pgfpathlineto{\pgfqpoint{4.679887in}{2.345515in}}%
\pgfpathclose%
\pgfusepath{stroke}%
\end{pgfscope}%
\begin{pgfscope}%
\pgfpathrectangle{\pgfqpoint{1.198472in}{0.823309in}}{\pgfqpoint{5.051528in}{3.826691in}}%
\pgfusepath{clip}%
\pgfsetbuttcap%
\pgfsetroundjoin%
\pgfsetlinewidth{1.505625pt}%
\definecolor{currentstroke}{rgb}{1.000000,0.498039,0.313725}%
\pgfsetstrokecolor{currentstroke}%
\pgfsetdash{}{0pt}%
\pgfpathmoveto{\pgfqpoint{4.726273in}{2.252465in}}%
\pgfpathlineto{\pgfqpoint{4.809607in}{2.252465in}}%
\pgfpathlineto{\pgfqpoint{4.809607in}{2.335798in}}%
\pgfpathlineto{\pgfqpoint{4.726273in}{2.335798in}}%
\pgfpathlineto{\pgfqpoint{4.726273in}{2.252465in}}%
\pgfpathclose%
\pgfusepath{stroke}%
\end{pgfscope}%
\begin{pgfscope}%
\pgfpathrectangle{\pgfqpoint{1.198472in}{0.823309in}}{\pgfqpoint{5.051528in}{3.826691in}}%
\pgfusepath{clip}%
\pgfsetbuttcap%
\pgfsetroundjoin%
\pgfsetlinewidth{1.505625pt}%
\definecolor{currentstroke}{rgb}{1.000000,0.498039,0.313725}%
\pgfsetstrokecolor{currentstroke}%
\pgfsetdash{}{0pt}%
\pgfpathmoveto{\pgfqpoint{4.772660in}{2.298640in}}%
\pgfpathlineto{\pgfqpoint{4.855994in}{2.298640in}}%
\pgfpathlineto{\pgfqpoint{4.855994in}{2.381974in}}%
\pgfpathlineto{\pgfqpoint{4.772660in}{2.381974in}}%
\pgfpathlineto{\pgfqpoint{4.772660in}{2.298640in}}%
\pgfpathclose%
\pgfusepath{stroke}%
\end{pgfscope}%
\begin{pgfscope}%
\pgfpathrectangle{\pgfqpoint{1.198472in}{0.823309in}}{\pgfqpoint{5.051528in}{3.826691in}}%
\pgfusepath{clip}%
\pgfsetbuttcap%
\pgfsetroundjoin%
\pgfsetlinewidth{1.505625pt}%
\definecolor{currentstroke}{rgb}{1.000000,0.498039,0.313725}%
\pgfsetstrokecolor{currentstroke}%
\pgfsetdash{}{0pt}%
\pgfpathmoveto{\pgfqpoint{4.819047in}{2.331741in}}%
\pgfpathlineto{\pgfqpoint{4.902380in}{2.331741in}}%
\pgfpathlineto{\pgfqpoint{4.902380in}{2.415075in}}%
\pgfpathlineto{\pgfqpoint{4.819047in}{2.415075in}}%
\pgfpathlineto{\pgfqpoint{4.819047in}{2.331741in}}%
\pgfpathclose%
\pgfusepath{stroke}%
\end{pgfscope}%
\begin{pgfscope}%
\pgfpathrectangle{\pgfqpoint{1.198472in}{0.823309in}}{\pgfqpoint{5.051528in}{3.826691in}}%
\pgfusepath{clip}%
\pgfsetbuttcap%
\pgfsetroundjoin%
\pgfsetlinewidth{1.505625pt}%
\definecolor{currentstroke}{rgb}{1.000000,0.498039,0.313725}%
\pgfsetstrokecolor{currentstroke}%
\pgfsetdash{}{0pt}%
\pgfpathmoveto{\pgfqpoint{4.865434in}{2.303365in}}%
\pgfpathlineto{\pgfqpoint{4.948767in}{2.303365in}}%
\pgfpathlineto{\pgfqpoint{4.948767in}{2.386699in}}%
\pgfpathlineto{\pgfqpoint{4.865434in}{2.386699in}}%
\pgfpathlineto{\pgfqpoint{4.865434in}{2.303365in}}%
\pgfpathclose%
\pgfusepath{stroke}%
\end{pgfscope}%
\begin{pgfscope}%
\pgfpathrectangle{\pgfqpoint{1.198472in}{0.823309in}}{\pgfqpoint{5.051528in}{3.826691in}}%
\pgfusepath{clip}%
\pgfsetbuttcap%
\pgfsetroundjoin%
\pgfsetlinewidth{1.505625pt}%
\definecolor{currentstroke}{rgb}{1.000000,0.498039,0.313725}%
\pgfsetstrokecolor{currentstroke}%
\pgfsetdash{}{0pt}%
\pgfpathmoveto{\pgfqpoint{4.911821in}{2.327091in}}%
\pgfpathlineto{\pgfqpoint{4.995154in}{2.327091in}}%
\pgfpathlineto{\pgfqpoint{4.995154in}{2.410424in}}%
\pgfpathlineto{\pgfqpoint{4.911821in}{2.410424in}}%
\pgfpathlineto{\pgfqpoint{4.911821in}{2.327091in}}%
\pgfpathclose%
\pgfusepath{stroke}%
\end{pgfscope}%
\begin{pgfscope}%
\pgfpathrectangle{\pgfqpoint{1.198472in}{0.823309in}}{\pgfqpoint{5.051528in}{3.826691in}}%
\pgfusepath{clip}%
\pgfsetbuttcap%
\pgfsetroundjoin%
\pgfsetlinewidth{1.505625pt}%
\definecolor{currentstroke}{rgb}{1.000000,0.498039,0.313725}%
\pgfsetstrokecolor{currentstroke}%
\pgfsetdash{}{0pt}%
\pgfpathmoveto{\pgfqpoint{4.958208in}{2.288087in}}%
\pgfpathlineto{\pgfqpoint{5.041541in}{2.288087in}}%
\pgfpathlineto{\pgfqpoint{5.041541in}{2.371420in}}%
\pgfpathlineto{\pgfqpoint{4.958208in}{2.371420in}}%
\pgfpathlineto{\pgfqpoint{4.958208in}{2.288087in}}%
\pgfpathclose%
\pgfusepath{stroke}%
\end{pgfscope}%
\begin{pgfscope}%
\pgfpathrectangle{\pgfqpoint{1.198472in}{0.823309in}}{\pgfqpoint{5.051528in}{3.826691in}}%
\pgfusepath{clip}%
\pgfsetbuttcap%
\pgfsetroundjoin%
\pgfsetlinewidth{1.505625pt}%
\definecolor{currentstroke}{rgb}{1.000000,0.498039,0.313725}%
\pgfsetstrokecolor{currentstroke}%
\pgfsetdash{}{0pt}%
\pgfpathmoveto{\pgfqpoint{5.004594in}{2.230582in}}%
\pgfpathlineto{\pgfqpoint{5.087928in}{2.230582in}}%
\pgfpathlineto{\pgfqpoint{5.087928in}{2.313916in}}%
\pgfpathlineto{\pgfqpoint{5.004594in}{2.313916in}}%
\pgfpathlineto{\pgfqpoint{5.004594in}{2.230582in}}%
\pgfpathclose%
\pgfusepath{stroke}%
\end{pgfscope}%
\begin{pgfscope}%
\pgfpathrectangle{\pgfqpoint{1.198472in}{0.823309in}}{\pgfqpoint{5.051528in}{3.826691in}}%
\pgfusepath{clip}%
\pgfsetbuttcap%
\pgfsetroundjoin%
\pgfsetlinewidth{1.505625pt}%
\definecolor{currentstroke}{rgb}{1.000000,0.498039,0.313725}%
\pgfsetstrokecolor{currentstroke}%
\pgfsetdash{}{0pt}%
\pgfpathmoveto{\pgfqpoint{5.050981in}{2.168784in}}%
\pgfpathlineto{\pgfqpoint{5.134315in}{2.168784in}}%
\pgfpathlineto{\pgfqpoint{5.134315in}{2.252118in}}%
\pgfpathlineto{\pgfqpoint{5.050981in}{2.252118in}}%
\pgfpathlineto{\pgfqpoint{5.050981in}{2.168784in}}%
\pgfpathclose%
\pgfusepath{stroke}%
\end{pgfscope}%
\begin{pgfscope}%
\pgfpathrectangle{\pgfqpoint{1.198472in}{0.823309in}}{\pgfqpoint{5.051528in}{3.826691in}}%
\pgfusepath{clip}%
\pgfsetbuttcap%
\pgfsetroundjoin%
\pgfsetlinewidth{1.505625pt}%
\definecolor{currentstroke}{rgb}{1.000000,0.498039,0.313725}%
\pgfsetstrokecolor{currentstroke}%
\pgfsetdash{}{0pt}%
\pgfpathmoveto{\pgfqpoint{5.097368in}{2.222988in}}%
\pgfpathlineto{\pgfqpoint{5.180702in}{2.222988in}}%
\pgfpathlineto{\pgfqpoint{5.180702in}{2.306321in}}%
\pgfpathlineto{\pgfqpoint{5.097368in}{2.306321in}}%
\pgfpathlineto{\pgfqpoint{5.097368in}{2.222988in}}%
\pgfpathclose%
\pgfusepath{stroke}%
\end{pgfscope}%
\begin{pgfscope}%
\pgfpathrectangle{\pgfqpoint{1.198472in}{0.823309in}}{\pgfqpoint{5.051528in}{3.826691in}}%
\pgfusepath{clip}%
\pgfsetbuttcap%
\pgfsetroundjoin%
\pgfsetlinewidth{1.505625pt}%
\definecolor{currentstroke}{rgb}{1.000000,0.498039,0.313725}%
\pgfsetstrokecolor{currentstroke}%
\pgfsetdash{}{0pt}%
\pgfpathmoveto{\pgfqpoint{5.143755in}{2.203147in}}%
\pgfpathlineto{\pgfqpoint{5.227088in}{2.203147in}}%
\pgfpathlineto{\pgfqpoint{5.227088in}{2.286480in}}%
\pgfpathlineto{\pgfqpoint{5.143755in}{2.286480in}}%
\pgfpathlineto{\pgfqpoint{5.143755in}{2.203147in}}%
\pgfpathclose%
\pgfusepath{stroke}%
\end{pgfscope}%
\begin{pgfscope}%
\pgfpathrectangle{\pgfqpoint{1.198472in}{0.823309in}}{\pgfqpoint{5.051528in}{3.826691in}}%
\pgfusepath{clip}%
\pgfsetbuttcap%
\pgfsetroundjoin%
\pgfsetlinewidth{1.505625pt}%
\definecolor{currentstroke}{rgb}{1.000000,0.498039,0.313725}%
\pgfsetstrokecolor{currentstroke}%
\pgfsetdash{}{0pt}%
\pgfpathmoveto{\pgfqpoint{5.190142in}{2.332818in}}%
\pgfpathlineto{\pgfqpoint{5.273475in}{2.332818in}}%
\pgfpathlineto{\pgfqpoint{5.273475in}{2.416152in}}%
\pgfpathlineto{\pgfqpoint{5.190142in}{2.416152in}}%
\pgfpathlineto{\pgfqpoint{5.190142in}{2.332818in}}%
\pgfpathclose%
\pgfusepath{stroke}%
\end{pgfscope}%
\begin{pgfscope}%
\pgfpathrectangle{\pgfqpoint{1.198472in}{0.823309in}}{\pgfqpoint{5.051528in}{3.826691in}}%
\pgfusepath{clip}%
\pgfsetbuttcap%
\pgfsetroundjoin%
\pgfsetlinewidth{1.505625pt}%
\definecolor{currentstroke}{rgb}{1.000000,0.498039,0.313725}%
\pgfsetstrokecolor{currentstroke}%
\pgfsetdash{}{0pt}%
\pgfpathmoveto{\pgfqpoint{5.236529in}{2.207239in}}%
\pgfpathlineto{\pgfqpoint{5.319862in}{2.207239in}}%
\pgfpathlineto{\pgfqpoint{5.319862in}{2.290573in}}%
\pgfpathlineto{\pgfqpoint{5.236529in}{2.290573in}}%
\pgfpathlineto{\pgfqpoint{5.236529in}{2.207239in}}%
\pgfpathclose%
\pgfusepath{stroke}%
\end{pgfscope}%
\begin{pgfscope}%
\pgfpathrectangle{\pgfqpoint{1.198472in}{0.823309in}}{\pgfqpoint{5.051528in}{3.826691in}}%
\pgfusepath{clip}%
\pgfsetbuttcap%
\pgfsetroundjoin%
\pgfsetlinewidth{1.505625pt}%
\definecolor{currentstroke}{rgb}{1.000000,0.498039,0.313725}%
\pgfsetstrokecolor{currentstroke}%
\pgfsetdash{}{0pt}%
\pgfpathmoveto{\pgfqpoint{5.282916in}{2.123291in}}%
\pgfpathlineto{\pgfqpoint{5.366249in}{2.123291in}}%
\pgfpathlineto{\pgfqpoint{5.366249in}{2.206624in}}%
\pgfpathlineto{\pgfqpoint{5.282916in}{2.206624in}}%
\pgfpathlineto{\pgfqpoint{5.282916in}{2.123291in}}%
\pgfpathclose%
\pgfusepath{stroke}%
\end{pgfscope}%
\begin{pgfscope}%
\pgfpathrectangle{\pgfqpoint{1.198472in}{0.823309in}}{\pgfqpoint{5.051528in}{3.826691in}}%
\pgfusepath{clip}%
\pgfsetbuttcap%
\pgfsetroundjoin%
\pgfsetlinewidth{1.505625pt}%
\definecolor{currentstroke}{rgb}{1.000000,0.498039,0.313725}%
\pgfsetstrokecolor{currentstroke}%
\pgfsetdash{}{0pt}%
\pgfpathmoveto{\pgfqpoint{5.329302in}{2.290535in}}%
\pgfpathlineto{\pgfqpoint{5.412636in}{2.290535in}}%
\pgfpathlineto{\pgfqpoint{5.412636in}{2.373868in}}%
\pgfpathlineto{\pgfqpoint{5.329302in}{2.373868in}}%
\pgfpathlineto{\pgfqpoint{5.329302in}{2.290535in}}%
\pgfpathclose%
\pgfusepath{stroke}%
\end{pgfscope}%
\begin{pgfscope}%
\pgfpathrectangle{\pgfqpoint{1.198472in}{0.823309in}}{\pgfqpoint{5.051528in}{3.826691in}}%
\pgfusepath{clip}%
\pgfsetbuttcap%
\pgfsetroundjoin%
\pgfsetlinewidth{1.505625pt}%
\definecolor{currentstroke}{rgb}{1.000000,0.498039,0.313725}%
\pgfsetstrokecolor{currentstroke}%
\pgfsetdash{}{0pt}%
\pgfpathmoveto{\pgfqpoint{5.375689in}{2.171927in}}%
\pgfpathlineto{\pgfqpoint{5.459023in}{2.171927in}}%
\pgfpathlineto{\pgfqpoint{5.459023in}{2.255260in}}%
\pgfpathlineto{\pgfqpoint{5.375689in}{2.255260in}}%
\pgfpathlineto{\pgfqpoint{5.375689in}{2.171927in}}%
\pgfpathclose%
\pgfusepath{stroke}%
\end{pgfscope}%
\begin{pgfscope}%
\pgfpathrectangle{\pgfqpoint{1.198472in}{0.823309in}}{\pgfqpoint{5.051528in}{3.826691in}}%
\pgfusepath{clip}%
\pgfsetbuttcap%
\pgfsetroundjoin%
\pgfsetlinewidth{1.505625pt}%
\definecolor{currentstroke}{rgb}{1.000000,0.498039,0.313725}%
\pgfsetstrokecolor{currentstroke}%
\pgfsetdash{}{0pt}%
\pgfpathmoveto{\pgfqpoint{5.422076in}{2.161574in}}%
\pgfpathlineto{\pgfqpoint{5.505410in}{2.161574in}}%
\pgfpathlineto{\pgfqpoint{5.505410in}{2.244907in}}%
\pgfpathlineto{\pgfqpoint{5.422076in}{2.244907in}}%
\pgfpathlineto{\pgfqpoint{5.422076in}{2.161574in}}%
\pgfpathclose%
\pgfusepath{stroke}%
\end{pgfscope}%
\begin{pgfscope}%
\pgfpathrectangle{\pgfqpoint{1.198472in}{0.823309in}}{\pgfqpoint{5.051528in}{3.826691in}}%
\pgfusepath{clip}%
\pgfsetbuttcap%
\pgfsetroundjoin%
\pgfsetlinewidth{1.505625pt}%
\definecolor{currentstroke}{rgb}{1.000000,0.498039,0.313725}%
\pgfsetstrokecolor{currentstroke}%
\pgfsetdash{}{0pt}%
\pgfpathmoveto{\pgfqpoint{5.468463in}{2.148999in}}%
\pgfpathlineto{\pgfqpoint{5.551796in}{2.148999in}}%
\pgfpathlineto{\pgfqpoint{5.551796in}{2.232332in}}%
\pgfpathlineto{\pgfqpoint{5.468463in}{2.232332in}}%
\pgfpathlineto{\pgfqpoint{5.468463in}{2.148999in}}%
\pgfpathclose%
\pgfusepath{stroke}%
\end{pgfscope}%
\begin{pgfscope}%
\pgfpathrectangle{\pgfqpoint{1.198472in}{0.823309in}}{\pgfqpoint{5.051528in}{3.826691in}}%
\pgfusepath{clip}%
\pgfsetbuttcap%
\pgfsetroundjoin%
\pgfsetlinewidth{1.505625pt}%
\definecolor{currentstroke}{rgb}{1.000000,0.498039,0.313725}%
\pgfsetstrokecolor{currentstroke}%
\pgfsetdash{}{0pt}%
\pgfpathmoveto{\pgfqpoint{5.514850in}{2.097286in}}%
\pgfpathlineto{\pgfqpoint{5.598183in}{2.097286in}}%
\pgfpathlineto{\pgfqpoint{5.598183in}{2.180619in}}%
\pgfpathlineto{\pgfqpoint{5.514850in}{2.180619in}}%
\pgfpathlineto{\pgfqpoint{5.514850in}{2.097286in}}%
\pgfpathclose%
\pgfusepath{stroke}%
\end{pgfscope}%
\begin{pgfscope}%
\pgfpathrectangle{\pgfqpoint{1.198472in}{0.823309in}}{\pgfqpoint{5.051528in}{3.826691in}}%
\pgfusepath{clip}%
\pgfsetbuttcap%
\pgfsetroundjoin%
\pgfsetlinewidth{1.505625pt}%
\definecolor{currentstroke}{rgb}{1.000000,0.498039,0.313725}%
\pgfsetstrokecolor{currentstroke}%
\pgfsetdash{}{0pt}%
\pgfpathmoveto{\pgfqpoint{5.561237in}{2.139180in}}%
\pgfpathlineto{\pgfqpoint{5.644570in}{2.139180in}}%
\pgfpathlineto{\pgfqpoint{5.644570in}{2.222513in}}%
\pgfpathlineto{\pgfqpoint{5.561237in}{2.222513in}}%
\pgfpathlineto{\pgfqpoint{5.561237in}{2.139180in}}%
\pgfpathclose%
\pgfusepath{stroke}%
\end{pgfscope}%
\begin{pgfscope}%
\pgfpathrectangle{\pgfqpoint{1.198472in}{0.823309in}}{\pgfqpoint{5.051528in}{3.826691in}}%
\pgfusepath{clip}%
\pgfsetbuttcap%
\pgfsetroundjoin%
\pgfsetlinewidth{1.505625pt}%
\definecolor{currentstroke}{rgb}{1.000000,0.498039,0.313725}%
\pgfsetstrokecolor{currentstroke}%
\pgfsetdash{}{0pt}%
\pgfpathmoveto{\pgfqpoint{5.607624in}{2.096134in}}%
\pgfpathlineto{\pgfqpoint{5.690957in}{2.096134in}}%
\pgfpathlineto{\pgfqpoint{5.690957in}{2.179467in}}%
\pgfpathlineto{\pgfqpoint{5.607624in}{2.179467in}}%
\pgfpathlineto{\pgfqpoint{5.607624in}{2.096134in}}%
\pgfpathclose%
\pgfusepath{stroke}%
\end{pgfscope}%
\begin{pgfscope}%
\pgfpathrectangle{\pgfqpoint{1.198472in}{0.823309in}}{\pgfqpoint{5.051528in}{3.826691in}}%
\pgfusepath{clip}%
\pgfsetbuttcap%
\pgfsetroundjoin%
\pgfsetlinewidth{1.505625pt}%
\definecolor{currentstroke}{rgb}{1.000000,0.498039,0.313725}%
\pgfsetstrokecolor{currentstroke}%
\pgfsetdash{}{0pt}%
\pgfpathmoveto{\pgfqpoint{5.654010in}{2.129776in}}%
\pgfpathlineto{\pgfqpoint{5.737344in}{2.129776in}}%
\pgfpathlineto{\pgfqpoint{5.737344in}{2.213109in}}%
\pgfpathlineto{\pgfqpoint{5.654010in}{2.213109in}}%
\pgfpathlineto{\pgfqpoint{5.654010in}{2.129776in}}%
\pgfpathclose%
\pgfusepath{stroke}%
\end{pgfscope}%
\begin{pgfscope}%
\pgfpathrectangle{\pgfqpoint{1.198472in}{0.823309in}}{\pgfqpoint{5.051528in}{3.826691in}}%
\pgfusepath{clip}%
\pgfsetbuttcap%
\pgfsetroundjoin%
\pgfsetlinewidth{1.505625pt}%
\definecolor{currentstroke}{rgb}{1.000000,0.498039,0.313725}%
\pgfsetstrokecolor{currentstroke}%
\pgfsetdash{}{0pt}%
\pgfpathmoveto{\pgfqpoint{5.700397in}{2.212466in}}%
\pgfpathlineto{\pgfqpoint{5.783731in}{2.212466in}}%
\pgfpathlineto{\pgfqpoint{5.783731in}{2.295799in}}%
\pgfpathlineto{\pgfqpoint{5.700397in}{2.295799in}}%
\pgfpathlineto{\pgfqpoint{5.700397in}{2.212466in}}%
\pgfpathclose%
\pgfusepath{stroke}%
\end{pgfscope}%
\begin{pgfscope}%
\pgfpathrectangle{\pgfqpoint{1.198472in}{0.823309in}}{\pgfqpoint{5.051528in}{3.826691in}}%
\pgfusepath{clip}%
\pgfsetbuttcap%
\pgfsetroundjoin%
\pgfsetlinewidth{1.505625pt}%
\definecolor{currentstroke}{rgb}{1.000000,0.498039,0.313725}%
\pgfsetstrokecolor{currentstroke}%
\pgfsetdash{}{0pt}%
\pgfpathmoveto{\pgfqpoint{5.746784in}{2.186169in}}%
\pgfpathlineto{\pgfqpoint{5.830117in}{2.186169in}}%
\pgfpathlineto{\pgfqpoint{5.830117in}{2.269502in}}%
\pgfpathlineto{\pgfqpoint{5.746784in}{2.269502in}}%
\pgfpathlineto{\pgfqpoint{5.746784in}{2.186169in}}%
\pgfpathclose%
\pgfusepath{stroke}%
\end{pgfscope}%
\begin{pgfscope}%
\pgfpathrectangle{\pgfqpoint{1.198472in}{0.823309in}}{\pgfqpoint{5.051528in}{3.826691in}}%
\pgfusepath{clip}%
\pgfsetbuttcap%
\pgfsetroundjoin%
\pgfsetlinewidth{1.505625pt}%
\definecolor{currentstroke}{rgb}{1.000000,0.498039,0.313725}%
\pgfsetstrokecolor{currentstroke}%
\pgfsetdash{}{0pt}%
\pgfpathmoveto{\pgfqpoint{5.793171in}{2.137253in}}%
\pgfpathlineto{\pgfqpoint{5.876504in}{2.137253in}}%
\pgfpathlineto{\pgfqpoint{5.876504in}{2.220587in}}%
\pgfpathlineto{\pgfqpoint{5.793171in}{2.220587in}}%
\pgfpathlineto{\pgfqpoint{5.793171in}{2.137253in}}%
\pgfpathclose%
\pgfusepath{stroke}%
\end{pgfscope}%
\begin{pgfscope}%
\pgfpathrectangle{\pgfqpoint{1.198472in}{0.823309in}}{\pgfqpoint{5.051528in}{3.826691in}}%
\pgfusepath{clip}%
\pgfsetbuttcap%
\pgfsetroundjoin%
\pgfsetlinewidth{1.505625pt}%
\definecolor{currentstroke}{rgb}{1.000000,0.498039,0.313725}%
\pgfsetstrokecolor{currentstroke}%
\pgfsetdash{}{0pt}%
\pgfpathmoveto{\pgfqpoint{5.839558in}{2.048321in}}%
\pgfpathlineto{\pgfqpoint{5.922891in}{2.048321in}}%
\pgfpathlineto{\pgfqpoint{5.922891in}{2.131654in}}%
\pgfpathlineto{\pgfqpoint{5.839558in}{2.131654in}}%
\pgfpathlineto{\pgfqpoint{5.839558in}{2.048321in}}%
\pgfpathclose%
\pgfusepath{stroke}%
\end{pgfscope}%
\begin{pgfscope}%
\pgfpathrectangle{\pgfqpoint{1.198472in}{0.823309in}}{\pgfqpoint{5.051528in}{3.826691in}}%
\pgfusepath{clip}%
\pgfsetbuttcap%
\pgfsetroundjoin%
\pgfsetlinewidth{1.505625pt}%
\definecolor{currentstroke}{rgb}{1.000000,0.498039,0.313725}%
\pgfsetstrokecolor{currentstroke}%
\pgfsetdash{}{0pt}%
\pgfpathmoveto{\pgfqpoint{5.885945in}{2.125987in}}%
\pgfpathlineto{\pgfqpoint{5.969278in}{2.125987in}}%
\pgfpathlineto{\pgfqpoint{5.969278in}{2.209320in}}%
\pgfpathlineto{\pgfqpoint{5.885945in}{2.209320in}}%
\pgfpathlineto{\pgfqpoint{5.885945in}{2.125987in}}%
\pgfpathclose%
\pgfusepath{stroke}%
\end{pgfscope}%
\begin{pgfscope}%
\pgfpathrectangle{\pgfqpoint{1.198472in}{0.823309in}}{\pgfqpoint{5.051528in}{3.826691in}}%
\pgfusepath{clip}%
\pgfsetbuttcap%
\pgfsetroundjoin%
\pgfsetlinewidth{1.505625pt}%
\definecolor{currentstroke}{rgb}{1.000000,0.498039,0.313725}%
\pgfsetstrokecolor{currentstroke}%
\pgfsetdash{}{0pt}%
\pgfpathmoveto{\pgfqpoint{5.932332in}{2.146398in}}%
\pgfpathlineto{\pgfqpoint{6.015665in}{2.146398in}}%
\pgfpathlineto{\pgfqpoint{6.015665in}{2.229731in}}%
\pgfpathlineto{\pgfqpoint{5.932332in}{2.229731in}}%
\pgfpathlineto{\pgfqpoint{5.932332in}{2.146398in}}%
\pgfpathclose%
\pgfusepath{stroke}%
\end{pgfscope}%
\begin{pgfscope}%
\pgfpathrectangle{\pgfqpoint{1.198472in}{0.823309in}}{\pgfqpoint{5.051528in}{3.826691in}}%
\pgfusepath{clip}%
\pgfsetbuttcap%
\pgfsetroundjoin%
\pgfsetlinewidth{1.505625pt}%
\definecolor{currentstroke}{rgb}{1.000000,0.498039,0.313725}%
\pgfsetstrokecolor{currentstroke}%
\pgfsetdash{}{0pt}%
\pgfpathmoveto{\pgfqpoint{5.978718in}{2.159654in}}%
\pgfpathlineto{\pgfqpoint{6.062052in}{2.159654in}}%
\pgfpathlineto{\pgfqpoint{6.062052in}{2.242987in}}%
\pgfpathlineto{\pgfqpoint{5.978718in}{2.242987in}}%
\pgfpathlineto{\pgfqpoint{5.978718in}{2.159654in}}%
\pgfpathclose%
\pgfusepath{stroke}%
\end{pgfscope}%
\begin{pgfscope}%
\pgfpathrectangle{\pgfqpoint{1.198472in}{0.823309in}}{\pgfqpoint{5.051528in}{3.826691in}}%
\pgfusepath{clip}%
\pgfsetbuttcap%
\pgfsetroundjoin%
\definecolor{currentfill}{rgb}{1.000000,0.498039,0.313725}%
\pgfsetfillcolor{currentfill}%
\pgfsetlinewidth{1.003750pt}%
\definecolor{currentstroke}{rgb}{1.000000,0.498039,0.313725}%
\pgfsetstrokecolor{currentstroke}%
\pgfsetdash{}{0pt}%
\pgfsys@defobject{currentmarker}{\pgfqpoint{-0.020833in}{-0.020833in}}{\pgfqpoint{0.020833in}{0.020833in}}{%
\pgfpathmoveto{\pgfqpoint{0.000000in}{-0.020833in}}%
\pgfpathcurveto{\pgfqpoint{0.005525in}{-0.020833in}}{\pgfqpoint{0.010825in}{-0.018638in}}{\pgfqpoint{0.014731in}{-0.014731in}}%
\pgfpathcurveto{\pgfqpoint{0.018638in}{-0.010825in}}{\pgfqpoint{0.020833in}{-0.005525in}}{\pgfqpoint{0.020833in}{0.000000in}}%
\pgfpathcurveto{\pgfqpoint{0.020833in}{0.005525in}}{\pgfqpoint{0.018638in}{0.010825in}}{\pgfqpoint{0.014731in}{0.014731in}}%
\pgfpathcurveto{\pgfqpoint{0.010825in}{0.018638in}}{\pgfqpoint{0.005525in}{0.020833in}}{\pgfqpoint{0.000000in}{0.020833in}}%
\pgfpathcurveto{\pgfqpoint{-0.005525in}{0.020833in}}{\pgfqpoint{-0.010825in}{0.018638in}}{\pgfqpoint{-0.014731in}{0.014731in}}%
\pgfpathcurveto{\pgfqpoint{-0.018638in}{0.010825in}}{\pgfqpoint{-0.020833in}{0.005525in}}{\pgfqpoint{-0.020833in}{0.000000in}}%
\pgfpathcurveto{\pgfqpoint{-0.020833in}{-0.005525in}}{\pgfqpoint{-0.018638in}{-0.010825in}}{\pgfqpoint{-0.014731in}{-0.014731in}}%
\pgfpathcurveto{\pgfqpoint{-0.010825in}{-0.018638in}}{\pgfqpoint{-0.005525in}{-0.020833in}}{\pgfqpoint{0.000000in}{-0.020833in}}%
\pgfpathlineto{\pgfqpoint{0.000000in}{-0.020833in}}%
\pgfpathclose%
\pgfusepath{stroke,fill}%
}%
\begin{pgfscope}%
\pgfsys@transformshift{1.428087in}{0.997250in}%
\pgfsys@useobject{currentmarker}{}%
\end{pgfscope}%
\begin{pgfscope}%
\pgfsys@transformshift{1.474473in}{1.012481in}%
\pgfsys@useobject{currentmarker}{}%
\end{pgfscope}%
\begin{pgfscope}%
\pgfsys@transformshift{1.520860in}{1.013172in}%
\pgfsys@useobject{currentmarker}{}%
\end{pgfscope}%
\begin{pgfscope}%
\pgfsys@transformshift{1.567247in}{1.026091in}%
\pgfsys@useobject{currentmarker}{}%
\end{pgfscope}%
\begin{pgfscope}%
\pgfsys@transformshift{1.613634in}{1.034707in}%
\pgfsys@useobject{currentmarker}{}%
\end{pgfscope}%
\begin{pgfscope}%
\pgfsys@transformshift{1.660021in}{1.039162in}%
\pgfsys@useobject{currentmarker}{}%
\end{pgfscope}%
\begin{pgfscope}%
\pgfsys@transformshift{1.706408in}{1.139428in}%
\pgfsys@useobject{currentmarker}{}%
\end{pgfscope}%
\begin{pgfscope}%
\pgfsys@transformshift{1.752795in}{1.225150in}%
\pgfsys@useobject{currentmarker}{}%
\end{pgfscope}%
\begin{pgfscope}%
\pgfsys@transformshift{1.799181in}{1.532278in}%
\pgfsys@useobject{currentmarker}{}%
\end{pgfscope}%
\begin{pgfscope}%
\pgfsys@transformshift{1.845568in}{1.349763in}%
\pgfsys@useobject{currentmarker}{}%
\end{pgfscope}%
\begin{pgfscope}%
\pgfsys@transformshift{1.891955in}{1.464789in}%
\pgfsys@useobject{currentmarker}{}%
\end{pgfscope}%
\begin{pgfscope}%
\pgfsys@transformshift{1.938342in}{1.737212in}%
\pgfsys@useobject{currentmarker}{}%
\end{pgfscope}%
\begin{pgfscope}%
\pgfsys@transformshift{1.984729in}{1.535472in}%
\pgfsys@useobject{currentmarker}{}%
\end{pgfscope}%
\begin{pgfscope}%
\pgfsys@transformshift{2.031116in}{1.771574in}%
\pgfsys@useobject{currentmarker}{}%
\end{pgfscope}%
\begin{pgfscope}%
\pgfsys@transformshift{2.077503in}{1.799516in}%
\pgfsys@useobject{currentmarker}{}%
\end{pgfscope}%
\begin{pgfscope}%
\pgfsys@transformshift{2.123889in}{2.099475in}%
\pgfsys@useobject{currentmarker}{}%
\end{pgfscope}%
\begin{pgfscope}%
\pgfsys@transformshift{2.170276in}{2.341339in}%
\pgfsys@useobject{currentmarker}{}%
\end{pgfscope}%
\begin{pgfscope}%
\pgfsys@transformshift{2.216663in}{2.186947in}%
\pgfsys@useobject{currentmarker}{}%
\end{pgfscope}%
\begin{pgfscope}%
\pgfsys@transformshift{2.263050in}{2.632888in}%
\pgfsys@useobject{currentmarker}{}%
\end{pgfscope}%
\begin{pgfscope}%
\pgfsys@transformshift{2.309437in}{2.610777in}%
\pgfsys@useobject{currentmarker}{}%
\end{pgfscope}%
\begin{pgfscope}%
\pgfsys@transformshift{2.355824in}{2.514294in}%
\pgfsys@useobject{currentmarker}{}%
\end{pgfscope}%
\begin{pgfscope}%
\pgfsys@transformshift{2.402211in}{2.562279in}%
\pgfsys@useobject{currentmarker}{}%
\end{pgfscope}%
\begin{pgfscope}%
\pgfsys@transformshift{2.448597in}{2.533866in}%
\pgfsys@useobject{currentmarker}{}%
\end{pgfscope}%
\begin{pgfscope}%
\pgfsys@transformshift{2.494984in}{2.895021in}%
\pgfsys@useobject{currentmarker}{}%
\end{pgfscope}%
\begin{pgfscope}%
\pgfsys@transformshift{2.541371in}{2.633160in}%
\pgfsys@useobject{currentmarker}{}%
\end{pgfscope}%
\begin{pgfscope}%
\pgfsys@transformshift{2.587758in}{2.901049in}%
\pgfsys@useobject{currentmarker}{}%
\end{pgfscope}%
\begin{pgfscope}%
\pgfsys@transformshift{2.634145in}{3.021650in}%
\pgfsys@useobject{currentmarker}{}%
\end{pgfscope}%
\begin{pgfscope}%
\pgfsys@transformshift{2.680532in}{2.835288in}%
\pgfsys@useobject{currentmarker}{}%
\end{pgfscope}%
\begin{pgfscope}%
\pgfsys@transformshift{2.726918in}{3.658374in}%
\pgfsys@useobject{currentmarker}{}%
\end{pgfscope}%
\begin{pgfscope}%
\pgfsys@transformshift{2.773305in}{3.292656in}%
\pgfsys@useobject{currentmarker}{}%
\end{pgfscope}%
\begin{pgfscope}%
\pgfsys@transformshift{2.819692in}{3.527797in}%
\pgfsys@useobject{currentmarker}{}%
\end{pgfscope}%
\begin{pgfscope}%
\pgfsys@transformshift{2.866079in}{3.667257in}%
\pgfsys@useobject{currentmarker}{}%
\end{pgfscope}%
\begin{pgfscope}%
\pgfsys@transformshift{2.912466in}{3.384434in}%
\pgfsys@useobject{currentmarker}{}%
\end{pgfscope}%
\begin{pgfscope}%
\pgfsys@transformshift{2.958853in}{3.270464in}%
\pgfsys@useobject{currentmarker}{}%
\end{pgfscope}%
\begin{pgfscope}%
\pgfsys@transformshift{3.005240in}{3.676297in}%
\pgfsys@useobject{currentmarker}{}%
\end{pgfscope}%
\begin{pgfscope}%
\pgfsys@transformshift{3.051626in}{3.749343in}%
\pgfsys@useobject{currentmarker}{}%
\end{pgfscope}%
\begin{pgfscope}%
\pgfsys@transformshift{3.098013in}{3.703305in}%
\pgfsys@useobject{currentmarker}{}%
\end{pgfscope}%
\begin{pgfscope}%
\pgfsys@transformshift{3.144400in}{4.180441in}%
\pgfsys@useobject{currentmarker}{}%
\end{pgfscope}%
\begin{pgfscope}%
\pgfsys@transformshift{3.190787in}{3.795626in}%
\pgfsys@useobject{currentmarker}{}%
\end{pgfscope}%
\begin{pgfscope}%
\pgfsys@transformshift{3.237174in}{3.685947in}%
\pgfsys@useobject{currentmarker}{}%
\end{pgfscope}%
\begin{pgfscope}%
\pgfsys@transformshift{3.283561in}{3.497847in}%
\pgfsys@useobject{currentmarker}{}%
\end{pgfscope}%
\begin{pgfscope}%
\pgfsys@transformshift{3.329948in}{3.709901in}%
\pgfsys@useobject{currentmarker}{}%
\end{pgfscope}%
\begin{pgfscope}%
\pgfsys@transformshift{3.376334in}{4.075131in}%
\pgfsys@useobject{currentmarker}{}%
\end{pgfscope}%
\begin{pgfscope}%
\pgfsys@transformshift{3.422721in}{3.418288in}%
\pgfsys@useobject{currentmarker}{}%
\end{pgfscope}%
\begin{pgfscope}%
\pgfsys@transformshift{3.469108in}{3.790723in}%
\pgfsys@useobject{currentmarker}{}%
\end{pgfscope}%
\begin{pgfscope}%
\pgfsys@transformshift{3.515495in}{3.677527in}%
\pgfsys@useobject{currentmarker}{}%
\end{pgfscope}%
\begin{pgfscope}%
\pgfsys@transformshift{3.561882in}{4.067330in}%
\pgfsys@useobject{currentmarker}{}%
\end{pgfscope}%
\begin{pgfscope}%
\pgfsys@transformshift{3.608269in}{3.880374in}%
\pgfsys@useobject{currentmarker}{}%
\end{pgfscope}%
\begin{pgfscope}%
\pgfsys@transformshift{3.654656in}{3.913181in}%
\pgfsys@useobject{currentmarker}{}%
\end{pgfscope}%
\begin{pgfscope}%
\pgfsys@transformshift{3.701042in}{3.980950in}%
\pgfsys@useobject{currentmarker}{}%
\end{pgfscope}%
\begin{pgfscope}%
\pgfsys@transformshift{3.747429in}{4.300916in}%
\pgfsys@useobject{currentmarker}{}%
\end{pgfscope}%
\begin{pgfscope}%
\pgfsys@transformshift{3.793816in}{3.888651in}%
\pgfsys@useobject{currentmarker}{}%
\end{pgfscope}%
\begin{pgfscope}%
\pgfsys@transformshift{3.840203in}{4.034117in}%
\pgfsys@useobject{currentmarker}{}%
\end{pgfscope}%
\begin{pgfscope}%
\pgfsys@transformshift{3.886590in}{3.916498in}%
\pgfsys@useobject{currentmarker}{}%
\end{pgfscope}%
\begin{pgfscope}%
\pgfsys@transformshift{3.932977in}{4.215254in}%
\pgfsys@useobject{currentmarker}{}%
\end{pgfscope}%
\begin{pgfscope}%
\pgfsys@transformshift{3.979364in}{4.160332in}%
\pgfsys@useobject{currentmarker}{}%
\end{pgfscope}%
\begin{pgfscope}%
\pgfsys@transformshift{4.025750in}{3.984695in}%
\pgfsys@useobject{currentmarker}{}%
\end{pgfscope}%
\begin{pgfscope}%
\pgfsys@transformshift{4.072137in}{3.846547in}%
\pgfsys@useobject{currentmarker}{}%
\end{pgfscope}%
\begin{pgfscope}%
\pgfsys@transformshift{4.118524in}{4.066563in}%
\pgfsys@useobject{currentmarker}{}%
\end{pgfscope}%
\begin{pgfscope}%
\pgfsys@transformshift{4.164911in}{4.457277in}%
\pgfsys@useobject{currentmarker}{}%
\end{pgfscope}%
\begin{pgfscope}%
\pgfsys@transformshift{4.211298in}{4.186634in}%
\pgfsys@useobject{currentmarker}{}%
\end{pgfscope}%
\begin{pgfscope}%
\pgfsys@transformshift{4.257685in}{3.755723in}%
\pgfsys@useobject{currentmarker}{}%
\end{pgfscope}%
\begin{pgfscope}%
\pgfsys@transformshift{4.304072in}{4.132317in}%
\pgfsys@useobject{currentmarker}{}%
\end{pgfscope}%
\begin{pgfscope}%
\pgfsys@transformshift{4.350458in}{4.153501in}%
\pgfsys@useobject{currentmarker}{}%
\end{pgfscope}%
\begin{pgfscope}%
\pgfsys@transformshift{4.396845in}{4.013123in}%
\pgfsys@useobject{currentmarker}{}%
\end{pgfscope}%
\begin{pgfscope}%
\pgfsys@transformshift{4.443232in}{3.930061in}%
\pgfsys@useobject{currentmarker}{}%
\end{pgfscope}%
\begin{pgfscope}%
\pgfsys@transformshift{4.489619in}{4.189591in}%
\pgfsys@useobject{currentmarker}{}%
\end{pgfscope}%
\begin{pgfscope}%
\pgfsys@transformshift{4.536006in}{4.273468in}%
\pgfsys@useobject{currentmarker}{}%
\end{pgfscope}%
\begin{pgfscope}%
\pgfsys@transformshift{4.582393in}{4.240931in}%
\pgfsys@useobject{currentmarker}{}%
\end{pgfscope}%
\begin{pgfscope}%
\pgfsys@transformshift{4.628779in}{4.385791in}%
\pgfsys@useobject{currentmarker}{}%
\end{pgfscope}%
\begin{pgfscope}%
\pgfsys@transformshift{4.675166in}{4.290889in}%
\pgfsys@useobject{currentmarker}{}%
\end{pgfscope}%
\begin{pgfscope}%
\pgfsys@transformshift{4.721553in}{4.319269in}%
\pgfsys@useobject{currentmarker}{}%
\end{pgfscope}%
\begin{pgfscope}%
\pgfsys@transformshift{4.767940in}{4.046991in}%
\pgfsys@useobject{currentmarker}{}%
\end{pgfscope}%
\begin{pgfscope}%
\pgfsys@transformshift{4.814327in}{4.329021in}%
\pgfsys@useobject{currentmarker}{}%
\end{pgfscope}%
\begin{pgfscope}%
\pgfsys@transformshift{4.860714in}{4.273727in}%
\pgfsys@useobject{currentmarker}{}%
\end{pgfscope}%
\begin{pgfscope}%
\pgfsys@transformshift{4.907101in}{3.769000in}%
\pgfsys@useobject{currentmarker}{}%
\end{pgfscope}%
\begin{pgfscope}%
\pgfsys@transformshift{4.953487in}{4.215926in}%
\pgfsys@useobject{currentmarker}{}%
\end{pgfscope}%
\begin{pgfscope}%
\pgfsys@transformshift{4.999874in}{4.266928in}%
\pgfsys@useobject{currentmarker}{}%
\end{pgfscope}%
\begin{pgfscope}%
\pgfsys@transformshift{5.046261in}{3.959476in}%
\pgfsys@useobject{currentmarker}{}%
\end{pgfscope}%
\begin{pgfscope}%
\pgfsys@transformshift{5.092648in}{4.254477in}%
\pgfsys@useobject{currentmarker}{}%
\end{pgfscope}%
\begin{pgfscope}%
\pgfsys@transformshift{5.139035in}{4.183201in}%
\pgfsys@useobject{currentmarker}{}%
\end{pgfscope}%
\begin{pgfscope}%
\pgfsys@transformshift{5.185422in}{4.382621in}%
\pgfsys@useobject{currentmarker}{}%
\end{pgfscope}%
\begin{pgfscope}%
\pgfsys@transformshift{5.231809in}{4.067223in}%
\pgfsys@useobject{currentmarker}{}%
\end{pgfscope}%
\begin{pgfscope}%
\pgfsys@transformshift{5.278195in}{3.937429in}%
\pgfsys@useobject{currentmarker}{}%
\end{pgfscope}%
\begin{pgfscope}%
\pgfsys@transformshift{5.324582in}{4.060121in}%
\pgfsys@useobject{currentmarker}{}%
\end{pgfscope}%
\begin{pgfscope}%
\pgfsys@transformshift{5.370969in}{4.066628in}%
\pgfsys@useobject{currentmarker}{}%
\end{pgfscope}%
\begin{pgfscope}%
\pgfsys@transformshift{5.417356in}{3.981939in}%
\pgfsys@useobject{currentmarker}{}%
\end{pgfscope}%
\begin{pgfscope}%
\pgfsys@transformshift{5.463743in}{3.950170in}%
\pgfsys@useobject{currentmarker}{}%
\end{pgfscope}%
\begin{pgfscope}%
\pgfsys@transformshift{5.510130in}{4.066919in}%
\pgfsys@useobject{currentmarker}{}%
\end{pgfscope}%
\begin{pgfscope}%
\pgfsys@transformshift{5.556517in}{4.073818in}%
\pgfsys@useobject{currentmarker}{}%
\end{pgfscope}%
\begin{pgfscope}%
\pgfsys@transformshift{5.602903in}{4.031461in}%
\pgfsys@useobject{currentmarker}{}%
\end{pgfscope}%
\begin{pgfscope}%
\pgfsys@transformshift{5.649290in}{4.047410in}%
\pgfsys@useobject{currentmarker}{}%
\end{pgfscope}%
\begin{pgfscope}%
\pgfsys@transformshift{5.695677in}{3.940416in}%
\pgfsys@useobject{currentmarker}{}%
\end{pgfscope}%
\begin{pgfscope}%
\pgfsys@transformshift{5.742064in}{3.900276in}%
\pgfsys@useobject{currentmarker}{}%
\end{pgfscope}%
\begin{pgfscope}%
\pgfsys@transformshift{5.788451in}{3.983201in}%
\pgfsys@useobject{currentmarker}{}%
\end{pgfscope}%
\begin{pgfscope}%
\pgfsys@transformshift{5.834838in}{3.944047in}%
\pgfsys@useobject{currentmarker}{}%
\end{pgfscope}%
\begin{pgfscope}%
\pgfsys@transformshift{5.881225in}{4.051491in}%
\pgfsys@useobject{currentmarker}{}%
\end{pgfscope}%
\begin{pgfscope}%
\pgfsys@transformshift{5.927611in}{4.229590in}%
\pgfsys@useobject{currentmarker}{}%
\end{pgfscope}%
\begin{pgfscope}%
\pgfsys@transformshift{5.973998in}{4.117787in}%
\pgfsys@useobject{currentmarker}{}%
\end{pgfscope}%
\begin{pgfscope}%
\pgfsys@transformshift{6.020385in}{4.247250in}%
\pgfsys@useobject{currentmarker}{}%
\end{pgfscope}%
\end{pgfscope}%
\begin{pgfscope}%
\pgfsetbuttcap%
\pgfsetroundjoin%
\definecolor{currentfill}{rgb}{0.000000,0.000000,0.000000}%
\pgfsetfillcolor{currentfill}%
\pgfsetlinewidth{0.803000pt}%
\definecolor{currentstroke}{rgb}{0.000000,0.000000,0.000000}%
\pgfsetstrokecolor{currentstroke}%
\pgfsetdash{}{0pt}%
\pgfsys@defobject{currentmarker}{\pgfqpoint{0.000000in}{-0.048611in}}{\pgfqpoint{0.000000in}{0.000000in}}{%
\pgfpathmoveto{\pgfqpoint{0.000000in}{0.000000in}}%
\pgfpathlineto{\pgfqpoint{0.000000in}{-0.048611in}}%
\pgfusepath{stroke,fill}%
}%
\begin{pgfscope}%
\pgfsys@transformshift{1.296878in}{0.823309in}%
\pgfsys@useobject{currentmarker}{}%
\end{pgfscope}%
\end{pgfscope}%
\begin{pgfscope}%
\definecolor{textcolor}{rgb}{0.000000,0.000000,0.000000}%
\pgfsetstrokecolor{textcolor}%
\pgfsetfillcolor{textcolor}%
\pgftext[x=1.296878in,y=0.726087in,,top]{\color{textcolor}{\rmfamily\fontsize{16.000000}{19.200000}\selectfont\catcode`\^=\active\def^{\ifmmode\sp\else\^{}\fi}\catcode`\%=\active\def%{\%}$\mathdefault{0.0}$}}%
\end{pgfscope}%
\begin{pgfscope}%
\pgfsetbuttcap%
\pgfsetroundjoin%
\definecolor{currentfill}{rgb}{0.000000,0.000000,0.000000}%
\pgfsetfillcolor{currentfill}%
\pgfsetlinewidth{0.803000pt}%
\definecolor{currentstroke}{rgb}{0.000000,0.000000,0.000000}%
\pgfsetstrokecolor{currentstroke}%
\pgfsetdash{}{0pt}%
\pgfsys@defobject{currentmarker}{\pgfqpoint{0.000000in}{-0.048611in}}{\pgfqpoint{0.000000in}{0.000000in}}{%
\pgfpathmoveto{\pgfqpoint{0.000000in}{0.000000in}}%
\pgfpathlineto{\pgfqpoint{0.000000in}{-0.048611in}}%
\pgfusepath{stroke,fill}%
}%
\begin{pgfscope}%
\pgfsys@transformshift{1.952921in}{0.823309in}%
\pgfsys@useobject{currentmarker}{}%
\end{pgfscope}%
\end{pgfscope}%
\begin{pgfscope}%
\definecolor{textcolor}{rgb}{0.000000,0.000000,0.000000}%
\pgfsetstrokecolor{textcolor}%
\pgfsetfillcolor{textcolor}%
\pgftext[x=1.952921in,y=0.726087in,,top]{\color{textcolor}{\rmfamily\fontsize{16.000000}{19.200000}\selectfont\catcode`\^=\active\def^{\ifmmode\sp\else\^{}\fi}\catcode`\%=\active\def%{\%}$\mathdefault{2.5}$}}%
\end{pgfscope}%
\begin{pgfscope}%
\pgfsetbuttcap%
\pgfsetroundjoin%
\definecolor{currentfill}{rgb}{0.000000,0.000000,0.000000}%
\pgfsetfillcolor{currentfill}%
\pgfsetlinewidth{0.803000pt}%
\definecolor{currentstroke}{rgb}{0.000000,0.000000,0.000000}%
\pgfsetstrokecolor{currentstroke}%
\pgfsetdash{}{0pt}%
\pgfsys@defobject{currentmarker}{\pgfqpoint{0.000000in}{-0.048611in}}{\pgfqpoint{0.000000in}{0.000000in}}{%
\pgfpathmoveto{\pgfqpoint{0.000000in}{0.000000in}}%
\pgfpathlineto{\pgfqpoint{0.000000in}{-0.048611in}}%
\pgfusepath{stroke,fill}%
}%
\begin{pgfscope}%
\pgfsys@transformshift{2.608963in}{0.823309in}%
\pgfsys@useobject{currentmarker}{}%
\end{pgfscope}%
\end{pgfscope}%
\begin{pgfscope}%
\definecolor{textcolor}{rgb}{0.000000,0.000000,0.000000}%
\pgfsetstrokecolor{textcolor}%
\pgfsetfillcolor{textcolor}%
\pgftext[x=2.608963in,y=0.726087in,,top]{\color{textcolor}{\rmfamily\fontsize{16.000000}{19.200000}\selectfont\catcode`\^=\active\def^{\ifmmode\sp\else\^{}\fi}\catcode`\%=\active\def%{\%}$\mathdefault{5.0}$}}%
\end{pgfscope}%
\begin{pgfscope}%
\pgfsetbuttcap%
\pgfsetroundjoin%
\definecolor{currentfill}{rgb}{0.000000,0.000000,0.000000}%
\pgfsetfillcolor{currentfill}%
\pgfsetlinewidth{0.803000pt}%
\definecolor{currentstroke}{rgb}{0.000000,0.000000,0.000000}%
\pgfsetstrokecolor{currentstroke}%
\pgfsetdash{}{0pt}%
\pgfsys@defobject{currentmarker}{\pgfqpoint{0.000000in}{-0.048611in}}{\pgfqpoint{0.000000in}{0.000000in}}{%
\pgfpathmoveto{\pgfqpoint{0.000000in}{0.000000in}}%
\pgfpathlineto{\pgfqpoint{0.000000in}{-0.048611in}}%
\pgfusepath{stroke,fill}%
}%
\begin{pgfscope}%
\pgfsys@transformshift{3.265006in}{0.823309in}%
\pgfsys@useobject{currentmarker}{}%
\end{pgfscope}%
\end{pgfscope}%
\begin{pgfscope}%
\definecolor{textcolor}{rgb}{0.000000,0.000000,0.000000}%
\pgfsetstrokecolor{textcolor}%
\pgfsetfillcolor{textcolor}%
\pgftext[x=3.265006in,y=0.726087in,,top]{\color{textcolor}{\rmfamily\fontsize{16.000000}{19.200000}\selectfont\catcode`\^=\active\def^{\ifmmode\sp\else\^{}\fi}\catcode`\%=\active\def%{\%}$\mathdefault{7.5}$}}%
\end{pgfscope}%
\begin{pgfscope}%
\pgfsetbuttcap%
\pgfsetroundjoin%
\definecolor{currentfill}{rgb}{0.000000,0.000000,0.000000}%
\pgfsetfillcolor{currentfill}%
\pgfsetlinewidth{0.803000pt}%
\definecolor{currentstroke}{rgb}{0.000000,0.000000,0.000000}%
\pgfsetstrokecolor{currentstroke}%
\pgfsetdash{}{0pt}%
\pgfsys@defobject{currentmarker}{\pgfqpoint{0.000000in}{-0.048611in}}{\pgfqpoint{0.000000in}{0.000000in}}{%
\pgfpathmoveto{\pgfqpoint{0.000000in}{0.000000in}}%
\pgfpathlineto{\pgfqpoint{0.000000in}{-0.048611in}}%
\pgfusepath{stroke,fill}%
}%
\begin{pgfscope}%
\pgfsys@transformshift{3.921049in}{0.823309in}%
\pgfsys@useobject{currentmarker}{}%
\end{pgfscope}%
\end{pgfscope}%
\begin{pgfscope}%
\definecolor{textcolor}{rgb}{0.000000,0.000000,0.000000}%
\pgfsetstrokecolor{textcolor}%
\pgfsetfillcolor{textcolor}%
\pgftext[x=3.921049in,y=0.726087in,,top]{\color{textcolor}{\rmfamily\fontsize{16.000000}{19.200000}\selectfont\catcode`\^=\active\def^{\ifmmode\sp\else\^{}\fi}\catcode`\%=\active\def%{\%}$\mathdefault{10.0}$}}%
\end{pgfscope}%
\begin{pgfscope}%
\pgfsetbuttcap%
\pgfsetroundjoin%
\definecolor{currentfill}{rgb}{0.000000,0.000000,0.000000}%
\pgfsetfillcolor{currentfill}%
\pgfsetlinewidth{0.803000pt}%
\definecolor{currentstroke}{rgb}{0.000000,0.000000,0.000000}%
\pgfsetstrokecolor{currentstroke}%
\pgfsetdash{}{0pt}%
\pgfsys@defobject{currentmarker}{\pgfqpoint{0.000000in}{-0.048611in}}{\pgfqpoint{0.000000in}{0.000000in}}{%
\pgfpathmoveto{\pgfqpoint{0.000000in}{0.000000in}}%
\pgfpathlineto{\pgfqpoint{0.000000in}{-0.048611in}}%
\pgfusepath{stroke,fill}%
}%
\begin{pgfscope}%
\pgfsys@transformshift{4.577091in}{0.823309in}%
\pgfsys@useobject{currentmarker}{}%
\end{pgfscope}%
\end{pgfscope}%
\begin{pgfscope}%
\definecolor{textcolor}{rgb}{0.000000,0.000000,0.000000}%
\pgfsetstrokecolor{textcolor}%
\pgfsetfillcolor{textcolor}%
\pgftext[x=4.577091in,y=0.726087in,,top]{\color{textcolor}{\rmfamily\fontsize{16.000000}{19.200000}\selectfont\catcode`\^=\active\def^{\ifmmode\sp\else\^{}\fi}\catcode`\%=\active\def%{\%}$\mathdefault{12.5}$}}%
\end{pgfscope}%
\begin{pgfscope}%
\pgfsetbuttcap%
\pgfsetroundjoin%
\definecolor{currentfill}{rgb}{0.000000,0.000000,0.000000}%
\pgfsetfillcolor{currentfill}%
\pgfsetlinewidth{0.803000pt}%
\definecolor{currentstroke}{rgb}{0.000000,0.000000,0.000000}%
\pgfsetstrokecolor{currentstroke}%
\pgfsetdash{}{0pt}%
\pgfsys@defobject{currentmarker}{\pgfqpoint{0.000000in}{-0.048611in}}{\pgfqpoint{0.000000in}{0.000000in}}{%
\pgfpathmoveto{\pgfqpoint{0.000000in}{0.000000in}}%
\pgfpathlineto{\pgfqpoint{0.000000in}{-0.048611in}}%
\pgfusepath{stroke,fill}%
}%
\begin{pgfscope}%
\pgfsys@transformshift{5.233134in}{0.823309in}%
\pgfsys@useobject{currentmarker}{}%
\end{pgfscope}%
\end{pgfscope}%
\begin{pgfscope}%
\definecolor{textcolor}{rgb}{0.000000,0.000000,0.000000}%
\pgfsetstrokecolor{textcolor}%
\pgfsetfillcolor{textcolor}%
\pgftext[x=5.233134in,y=0.726087in,,top]{\color{textcolor}{\rmfamily\fontsize{16.000000}{19.200000}\selectfont\catcode`\^=\active\def^{\ifmmode\sp\else\^{}\fi}\catcode`\%=\active\def%{\%}$\mathdefault{15.0}$}}%
\end{pgfscope}%
\begin{pgfscope}%
\pgfsetbuttcap%
\pgfsetroundjoin%
\definecolor{currentfill}{rgb}{0.000000,0.000000,0.000000}%
\pgfsetfillcolor{currentfill}%
\pgfsetlinewidth{0.803000pt}%
\definecolor{currentstroke}{rgb}{0.000000,0.000000,0.000000}%
\pgfsetstrokecolor{currentstroke}%
\pgfsetdash{}{0pt}%
\pgfsys@defobject{currentmarker}{\pgfqpoint{0.000000in}{-0.048611in}}{\pgfqpoint{0.000000in}{0.000000in}}{%
\pgfpathmoveto{\pgfqpoint{0.000000in}{0.000000in}}%
\pgfpathlineto{\pgfqpoint{0.000000in}{-0.048611in}}%
\pgfusepath{stroke,fill}%
}%
\begin{pgfscope}%
\pgfsys@transformshift{5.889177in}{0.823309in}%
\pgfsys@useobject{currentmarker}{}%
\end{pgfscope}%
\end{pgfscope}%
\begin{pgfscope}%
\definecolor{textcolor}{rgb}{0.000000,0.000000,0.000000}%
\pgfsetstrokecolor{textcolor}%
\pgfsetfillcolor{textcolor}%
\pgftext[x=5.889177in,y=0.726087in,,top]{\color{textcolor}{\rmfamily\fontsize{16.000000}{19.200000}\selectfont\catcode`\^=\active\def^{\ifmmode\sp\else\^{}\fi}\catcode`\%=\active\def%{\%}$\mathdefault{17.5}$}}%
\end{pgfscope}%
\begin{pgfscope}%
\definecolor{textcolor}{rgb}{0.000000,0.000000,0.000000}%
\pgfsetstrokecolor{textcolor}%
\pgfsetfillcolor{textcolor}%
\pgftext[x=3.724236in,y=0.457183in,,top]{\color{textcolor}{\rmfamily\fontsize{24.000000}{28.800000}\selectfont\catcode`\^=\active\def^{\ifmmode\sp\else\^{}\fi}\catcode`\%=\active\def%{\%}Flashing time $\tau$}}%
\end{pgfscope}%
\begin{pgfscope}%
\pgfsetbuttcap%
\pgfsetroundjoin%
\definecolor{currentfill}{rgb}{0.000000,0.000000,0.000000}%
\pgfsetfillcolor{currentfill}%
\pgfsetlinewidth{0.803000pt}%
\definecolor{currentstroke}{rgb}{0.000000,0.000000,0.000000}%
\pgfsetstrokecolor{currentstroke}%
\pgfsetdash{}{0pt}%
\pgfsys@defobject{currentmarker}{\pgfqpoint{-0.048611in}{0.000000in}}{\pgfqpoint{-0.000000in}{0.000000in}}{%
\pgfpathmoveto{\pgfqpoint{-0.000000in}{0.000000in}}%
\pgfpathlineto{\pgfqpoint{-0.048611in}{0.000000in}}%
\pgfusepath{stroke,fill}%
}%
\begin{pgfscope}%
\pgfsys@transformshift{1.198472in}{1.058593in}%
\pgfsys@useobject{currentmarker}{}%
\end{pgfscope}%
\end{pgfscope}%
\begin{pgfscope}%
\definecolor{textcolor}{rgb}{0.000000,0.000000,0.000000}%
\pgfsetstrokecolor{textcolor}%
\pgfsetfillcolor{textcolor}%
\pgftext[x=0.595700in, y=0.975259in, left, base]{\color{textcolor}{\rmfamily\fontsize{16.000000}{19.200000}\selectfont\catcode`\^=\active\def^{\ifmmode\sp\else\^{}\fi}\catcode`\%=\active\def%{\%}$\mathdefault{0.000}$}}%
\end{pgfscope}%
\begin{pgfscope}%
\pgfsetbuttcap%
\pgfsetroundjoin%
\definecolor{currentfill}{rgb}{0.000000,0.000000,0.000000}%
\pgfsetfillcolor{currentfill}%
\pgfsetlinewidth{0.803000pt}%
\definecolor{currentstroke}{rgb}{0.000000,0.000000,0.000000}%
\pgfsetstrokecolor{currentstroke}%
\pgfsetdash{}{0pt}%
\pgfsys@defobject{currentmarker}{\pgfqpoint{-0.048611in}{0.000000in}}{\pgfqpoint{-0.000000in}{0.000000in}}{%
\pgfpathmoveto{\pgfqpoint{-0.000000in}{0.000000in}}%
\pgfpathlineto{\pgfqpoint{-0.048611in}{0.000000in}}%
\pgfusepath{stroke,fill}%
}%
\begin{pgfscope}%
\pgfsys@transformshift{1.198472in}{1.856170in}%
\pgfsys@useobject{currentmarker}{}%
\end{pgfscope}%
\end{pgfscope}%
\begin{pgfscope}%
\definecolor{textcolor}{rgb}{0.000000,0.000000,0.000000}%
\pgfsetstrokecolor{textcolor}%
\pgfsetfillcolor{textcolor}%
\pgftext[x=0.595700in, y=1.772837in, left, base]{\color{textcolor}{\rmfamily\fontsize{16.000000}{19.200000}\selectfont\catcode`\^=\active\def^{\ifmmode\sp\else\^{}\fi}\catcode`\%=\active\def%{\%}$\mathdefault{0.005}$}}%
\end{pgfscope}%
\begin{pgfscope}%
\pgfsetbuttcap%
\pgfsetroundjoin%
\definecolor{currentfill}{rgb}{0.000000,0.000000,0.000000}%
\pgfsetfillcolor{currentfill}%
\pgfsetlinewidth{0.803000pt}%
\definecolor{currentstroke}{rgb}{0.000000,0.000000,0.000000}%
\pgfsetstrokecolor{currentstroke}%
\pgfsetdash{}{0pt}%
\pgfsys@defobject{currentmarker}{\pgfqpoint{-0.048611in}{0.000000in}}{\pgfqpoint{-0.000000in}{0.000000in}}{%
\pgfpathmoveto{\pgfqpoint{-0.000000in}{0.000000in}}%
\pgfpathlineto{\pgfqpoint{-0.048611in}{0.000000in}}%
\pgfusepath{stroke,fill}%
}%
\begin{pgfscope}%
\pgfsys@transformshift{1.198472in}{2.653748in}%
\pgfsys@useobject{currentmarker}{}%
\end{pgfscope}%
\end{pgfscope}%
\begin{pgfscope}%
\definecolor{textcolor}{rgb}{0.000000,0.000000,0.000000}%
\pgfsetstrokecolor{textcolor}%
\pgfsetfillcolor{textcolor}%
\pgftext[x=0.595700in, y=2.570414in, left, base]{\color{textcolor}{\rmfamily\fontsize{16.000000}{19.200000}\selectfont\catcode`\^=\active\def^{\ifmmode\sp\else\^{}\fi}\catcode`\%=\active\def%{\%}$\mathdefault{0.010}$}}%
\end{pgfscope}%
\begin{pgfscope}%
\pgfsetbuttcap%
\pgfsetroundjoin%
\definecolor{currentfill}{rgb}{0.000000,0.000000,0.000000}%
\pgfsetfillcolor{currentfill}%
\pgfsetlinewidth{0.803000pt}%
\definecolor{currentstroke}{rgb}{0.000000,0.000000,0.000000}%
\pgfsetstrokecolor{currentstroke}%
\pgfsetdash{}{0pt}%
\pgfsys@defobject{currentmarker}{\pgfqpoint{-0.048611in}{0.000000in}}{\pgfqpoint{-0.000000in}{0.000000in}}{%
\pgfpathmoveto{\pgfqpoint{-0.000000in}{0.000000in}}%
\pgfpathlineto{\pgfqpoint{-0.048611in}{0.000000in}}%
\pgfusepath{stroke,fill}%
}%
\begin{pgfscope}%
\pgfsys@transformshift{1.198472in}{3.451325in}%
\pgfsys@useobject{currentmarker}{}%
\end{pgfscope}%
\end{pgfscope}%
\begin{pgfscope}%
\definecolor{textcolor}{rgb}{0.000000,0.000000,0.000000}%
\pgfsetstrokecolor{textcolor}%
\pgfsetfillcolor{textcolor}%
\pgftext[x=0.595700in, y=3.367992in, left, base]{\color{textcolor}{\rmfamily\fontsize{16.000000}{19.200000}\selectfont\catcode`\^=\active\def^{\ifmmode\sp\else\^{}\fi}\catcode`\%=\active\def%{\%}$\mathdefault{0.015}$}}%
\end{pgfscope}%
\begin{pgfscope}%
\pgfsetbuttcap%
\pgfsetroundjoin%
\definecolor{currentfill}{rgb}{0.000000,0.000000,0.000000}%
\pgfsetfillcolor{currentfill}%
\pgfsetlinewidth{0.803000pt}%
\definecolor{currentstroke}{rgb}{0.000000,0.000000,0.000000}%
\pgfsetstrokecolor{currentstroke}%
\pgfsetdash{}{0pt}%
\pgfsys@defobject{currentmarker}{\pgfqpoint{-0.048611in}{0.000000in}}{\pgfqpoint{-0.000000in}{0.000000in}}{%
\pgfpathmoveto{\pgfqpoint{-0.000000in}{0.000000in}}%
\pgfpathlineto{\pgfqpoint{-0.048611in}{0.000000in}}%
\pgfusepath{stroke,fill}%
}%
\begin{pgfscope}%
\pgfsys@transformshift{1.198472in}{4.248902in}%
\pgfsys@useobject{currentmarker}{}%
\end{pgfscope}%
\end{pgfscope}%
\begin{pgfscope}%
\definecolor{textcolor}{rgb}{0.000000,0.000000,0.000000}%
\pgfsetstrokecolor{textcolor}%
\pgfsetfillcolor{textcolor}%
\pgftext[x=0.595700in, y=4.165569in, left, base]{\color{textcolor}{\rmfamily\fontsize{16.000000}{19.200000}\selectfont\catcode`\^=\active\def^{\ifmmode\sp\else\^{}\fi}\catcode`\%=\active\def%{\%}$\mathdefault{0.020}$}}%
\end{pgfscope}%
\begin{pgfscope}%
\definecolor{textcolor}{rgb}{0.000000,0.000000,0.000000}%
\pgfsetstrokecolor{textcolor}%
\pgfsetfillcolor{textcolor}%
\pgftext[x=0.540144in,y=2.736655in,,bottom,rotate=90.000000]{\color{textcolor}{\rmfamily\fontsize{24.000000}{28.800000}\selectfont\catcode`\^=\active\def^{\ifmmode\sp\else\^{}\fi}\catcode`\%=\active\def%{\%}Mean drift velocity ($\hat{x} / \hat{t}$)}}%
\end{pgfscope}%
\begin{pgfscope}%
\pgfpathrectangle{\pgfqpoint{1.198472in}{0.823309in}}{\pgfqpoint{5.051528in}{3.826691in}}%
\pgfusepath{clip}%
\pgfsetbuttcap%
\pgfsetroundjoin%
\pgfsetlinewidth{1.505625pt}%
\definecolor{currentstroke}{rgb}{1.000000,0.000000,0.000000}%
\pgfsetstrokecolor{currentstroke}%
\pgfsetdash{{5.550000pt}{2.400000pt}}{0.000000pt}%
\pgfpathmoveto{\pgfqpoint{3.144400in}{0.823309in}}%
\pgfpathlineto{\pgfqpoint{3.144400in}{4.650000in}}%
\pgfusepath{stroke}%
\end{pgfscope}%
\begin{pgfscope}%
\pgfsetrectcap%
\pgfsetmiterjoin%
\pgfsetlinewidth{0.803000pt}%
\definecolor{currentstroke}{rgb}{0.000000,0.000000,0.000000}%
\pgfsetstrokecolor{currentstroke}%
\pgfsetdash{}{0pt}%
\pgfpathmoveto{\pgfqpoint{1.198472in}{0.823309in}}%
\pgfpathlineto{\pgfqpoint{1.198472in}{4.650000in}}%
\pgfusepath{stroke}%
\end{pgfscope}%
\begin{pgfscope}%
\pgfsetrectcap%
\pgfsetmiterjoin%
\pgfsetlinewidth{0.803000pt}%
\definecolor{currentstroke}{rgb}{0.000000,0.000000,0.000000}%
\pgfsetstrokecolor{currentstroke}%
\pgfsetdash{}{0pt}%
\pgfpathmoveto{\pgfqpoint{6.250000in}{0.823309in}}%
\pgfpathlineto{\pgfqpoint{6.250000in}{4.650000in}}%
\pgfusepath{stroke}%
\end{pgfscope}%
\begin{pgfscope}%
\pgfsetrectcap%
\pgfsetmiterjoin%
\pgfsetlinewidth{0.803000pt}%
\definecolor{currentstroke}{rgb}{0.000000,0.000000,0.000000}%
\pgfsetstrokecolor{currentstroke}%
\pgfsetdash{}{0pt}%
\pgfpathmoveto{\pgfqpoint{1.198472in}{0.823309in}}%
\pgfpathlineto{\pgfqpoint{6.250000in}{0.823309in}}%
\pgfusepath{stroke}%
\end{pgfscope}%
\begin{pgfscope}%
\pgfsetrectcap%
\pgfsetmiterjoin%
\pgfsetlinewidth{0.803000pt}%
\definecolor{currentstroke}{rgb}{0.000000,0.000000,0.000000}%
\pgfsetstrokecolor{currentstroke}%
\pgfsetdash{}{0pt}%
\pgfpathmoveto{\pgfqpoint{1.198472in}{4.650000in}}%
\pgfpathlineto{\pgfqpoint{6.250000in}{4.650000in}}%
\pgfusepath{stroke}%
\end{pgfscope}%
\begin{pgfscope}%
\pgfsetbuttcap%
\pgfsetmiterjoin%
\definecolor{currentfill}{rgb}{1.000000,1.000000,1.000000}%
\pgfsetfillcolor{currentfill}%
\pgfsetfillopacity{0.800000}%
\pgfsetlinewidth{1.003750pt}%
\definecolor{currentstroke}{rgb}{0.800000,0.800000,0.800000}%
\pgfsetstrokecolor{currentstroke}%
\pgfsetstrokeopacity{0.800000}%
\pgfsetdash{}{0pt}%
\pgfpathmoveto{\pgfqpoint{2.318913in}{0.934420in}}%
\pgfpathlineto{\pgfqpoint{5.129558in}{0.934420in}}%
\pgfpathquadraticcurveto{\pgfqpoint{5.174003in}{0.934420in}}{\pgfqpoint{5.174003in}{0.978865in}}%
\pgfpathlineto{\pgfqpoint{5.174003in}{2.254481in}}%
\pgfpathquadraticcurveto{\pgfqpoint{5.174003in}{2.298926in}}{\pgfqpoint{5.129558in}{2.298926in}}%
\pgfpathlineto{\pgfqpoint{2.318913in}{2.298926in}}%
\pgfpathquadraticcurveto{\pgfqpoint{2.274469in}{2.298926in}}{\pgfqpoint{2.274469in}{2.254481in}}%
\pgfpathlineto{\pgfqpoint{2.274469in}{0.978865in}}%
\pgfpathquadraticcurveto{\pgfqpoint{2.274469in}{0.934420in}}{\pgfqpoint{2.318913in}{0.934420in}}%
\pgfpathlineto{\pgfqpoint{2.318913in}{0.934420in}}%
\pgfpathclose%
\pgfusepath{stroke,fill}%
\end{pgfscope}%
\begin{pgfscope}%
\pgfsetbuttcap%
\pgfsetroundjoin%
\pgfsetlinewidth{1.505625pt}%
\definecolor{currentstroke}{rgb}{1.000000,0.000000,0.000000}%
\pgfsetstrokecolor{currentstroke}%
\pgfsetdash{{5.550000pt}{2.400000pt}}{0.000000pt}%
\pgfpathmoveto{\pgfqpoint{2.363358in}{2.121148in}}%
\pgfpathlineto{\pgfqpoint{2.585580in}{2.121148in}}%
\pgfpathlineto{\pgfqpoint{2.807802in}{2.121148in}}%
\pgfusepath{stroke}%
\end{pgfscope}%
\begin{pgfscope}%
\definecolor{textcolor}{rgb}{0.000000,0.000000,0.000000}%
\pgfsetstrokecolor{textcolor}%
\pgfsetfillcolor{textcolor}%
\pgftext[x=2.985580in,y=2.043370in,left,base]{\color{textcolor}{\rmfamily\fontsize{16.000000}{19.200000}\selectfont\catcode`\^=\active\def^{\ifmmode\sp\else\^{}\fi}\catcode`\%=\active\def%{\%}Optimal flashing time}}%
\end{pgfscope}%
\begin{pgfscope}%
\pgfsetbuttcap%
\pgfsetroundjoin%
\definecolor{currentfill}{rgb}{0.121569,0.466667,0.705882}%
\pgfsetfillcolor{currentfill}%
\pgfsetlinewidth{1.003750pt}%
\definecolor{currentstroke}{rgb}{0.121569,0.466667,0.705882}%
\pgfsetstrokecolor{currentstroke}%
\pgfsetdash{}{0pt}%
\pgfsys@defobject{currentmarker}{\pgfqpoint{-0.020833in}{-0.020833in}}{\pgfqpoint{0.020833in}{0.020833in}}{%
\pgfpathmoveto{\pgfqpoint{0.000000in}{-0.020833in}}%
\pgfpathcurveto{\pgfqpoint{0.005525in}{-0.020833in}}{\pgfqpoint{0.010825in}{-0.018638in}}{\pgfqpoint{0.014731in}{-0.014731in}}%
\pgfpathcurveto{\pgfqpoint{0.018638in}{-0.010825in}}{\pgfqpoint{0.020833in}{-0.005525in}}{\pgfqpoint{0.020833in}{0.000000in}}%
\pgfpathcurveto{\pgfqpoint{0.020833in}{0.005525in}}{\pgfqpoint{0.018638in}{0.010825in}}{\pgfqpoint{0.014731in}{0.014731in}}%
\pgfpathcurveto{\pgfqpoint{0.010825in}{0.018638in}}{\pgfqpoint{0.005525in}{0.020833in}}{\pgfqpoint{0.000000in}{0.020833in}}%
\pgfpathcurveto{\pgfqpoint{-0.005525in}{0.020833in}}{\pgfqpoint{-0.010825in}{0.018638in}}{\pgfqpoint{-0.014731in}{0.014731in}}%
\pgfpathcurveto{\pgfqpoint{-0.018638in}{0.010825in}}{\pgfqpoint{-0.020833in}{0.005525in}}{\pgfqpoint{-0.020833in}{0.000000in}}%
\pgfpathcurveto{\pgfqpoint{-0.020833in}{-0.005525in}}{\pgfqpoint{-0.018638in}{-0.010825in}}{\pgfqpoint{-0.014731in}{-0.014731in}}%
\pgfpathcurveto{\pgfqpoint{-0.010825in}{-0.018638in}}{\pgfqpoint{-0.005525in}{-0.020833in}}{\pgfqpoint{0.000000in}{-0.020833in}}%
\pgfpathlineto{\pgfqpoint{0.000000in}{-0.020833in}}%
\pgfpathclose%
\pgfusepath{stroke,fill}%
}%
\begin{pgfscope}%
\pgfsys@transformshift{2.585580in}{1.777244in}%
\pgfsys@useobject{currentmarker}{}%
\end{pgfscope}%
\end{pgfscope}%
\begin{pgfscope}%
\definecolor{textcolor}{rgb}{0.000000,0.000000,0.000000}%
\pgfsetstrokecolor{textcolor}%
\pgfsetfillcolor{textcolor}%
\pgftext[x=2.985580in,y=1.718911in,left,base]{\color{textcolor}{\rmfamily\fontsize{16.000000}{19.200000}\selectfont\catcode`\^=\active\def^{\ifmmode\sp\else\^{}\fi}\catcode`\%=\active\def%{\%}Simulation 12 nm}}%
\end{pgfscope}%
\begin{pgfscope}%
\pgfsetbuttcap%
\pgfsetroundjoin%
\pgfsetlinewidth{1.505625pt}%
\definecolor{currentstroke}{rgb}{1.000000,0.498039,0.313725}%
\pgfsetstrokecolor{currentstroke}%
\pgfsetdash{}{0pt}%
\pgfpathmoveto{\pgfqpoint{2.543913in}{1.411118in}}%
\pgfpathlineto{\pgfqpoint{2.627247in}{1.411118in}}%
\pgfpathlineto{\pgfqpoint{2.627247in}{1.494451in}}%
\pgfpathlineto{\pgfqpoint{2.543913in}{1.494451in}}%
\pgfpathlineto{\pgfqpoint{2.543913in}{1.411118in}}%
\pgfpathclose%
\pgfusepath{stroke}%
\end{pgfscope}%
\begin{pgfscope}%
\definecolor{textcolor}{rgb}{0.000000,0.000000,0.000000}%
\pgfsetstrokecolor{textcolor}%
\pgfsetfillcolor{textcolor}%
\pgftext[x=2.985580in,y=1.394451in,left,base]{\color{textcolor}{\rmfamily\fontsize{16.000000}{19.200000}\selectfont\catcode`\^=\active\def^{\ifmmode\sp\else\^{}\fi}\catcode`\%=\active\def%{\%}Prediction 24 nm}}%
\end{pgfscope}%
\begin{pgfscope}%
\pgfsetbuttcap%
\pgfsetroundjoin%
\definecolor{currentfill}{rgb}{1.000000,0.498039,0.313725}%
\pgfsetfillcolor{currentfill}%
\pgfsetlinewidth{1.003750pt}%
\definecolor{currentstroke}{rgb}{1.000000,0.498039,0.313725}%
\pgfsetstrokecolor{currentstroke}%
\pgfsetdash{}{0pt}%
\pgfsys@defobject{currentmarker}{\pgfqpoint{-0.020833in}{-0.020833in}}{\pgfqpoint{0.020833in}{0.020833in}}{%
\pgfpathmoveto{\pgfqpoint{0.000000in}{-0.020833in}}%
\pgfpathcurveto{\pgfqpoint{0.005525in}{-0.020833in}}{\pgfqpoint{0.010825in}{-0.018638in}}{\pgfqpoint{0.014731in}{-0.014731in}}%
\pgfpathcurveto{\pgfqpoint{0.018638in}{-0.010825in}}{\pgfqpoint{0.020833in}{-0.005525in}}{\pgfqpoint{0.020833in}{0.000000in}}%
\pgfpathcurveto{\pgfqpoint{0.020833in}{0.005525in}}{\pgfqpoint{0.018638in}{0.010825in}}{\pgfqpoint{0.014731in}{0.014731in}}%
\pgfpathcurveto{\pgfqpoint{0.010825in}{0.018638in}}{\pgfqpoint{0.005525in}{0.020833in}}{\pgfqpoint{0.000000in}{0.020833in}}%
\pgfpathcurveto{\pgfqpoint{-0.005525in}{0.020833in}}{\pgfqpoint{-0.010825in}{0.018638in}}{\pgfqpoint{-0.014731in}{0.014731in}}%
\pgfpathcurveto{\pgfqpoint{-0.018638in}{0.010825in}}{\pgfqpoint{-0.020833in}{0.005525in}}{\pgfqpoint{-0.020833in}{0.000000in}}%
\pgfpathcurveto{\pgfqpoint{-0.020833in}{-0.005525in}}{\pgfqpoint{-0.018638in}{-0.010825in}}{\pgfqpoint{-0.014731in}{-0.014731in}}%
\pgfpathcurveto{\pgfqpoint{-0.010825in}{-0.018638in}}{\pgfqpoint{-0.005525in}{-0.020833in}}{\pgfqpoint{0.000000in}{-0.020833in}}%
\pgfpathlineto{\pgfqpoint{0.000000in}{-0.020833in}}%
\pgfpathclose%
\pgfusepath{stroke,fill}%
}%
\begin{pgfscope}%
\pgfsys@transformshift{2.585580in}{1.128324in}%
\pgfsys@useobject{currentmarker}{}%
\end{pgfscope}%
\end{pgfscope}%
\begin{pgfscope}%
\definecolor{textcolor}{rgb}{0.000000,0.000000,0.000000}%
\pgfsetstrokecolor{textcolor}%
\pgfsetfillcolor{textcolor}%
\pgftext[x=2.985580in,y=1.069991in,left,base]{\color{textcolor}{\rmfamily\fontsize{16.000000}{19.200000}\selectfont\catcode`\^=\active\def^{\ifmmode\sp\else\^{}\fi}\catcode`\%=\active\def%{\%}Simulation 24 nm}}%
\end{pgfscope}%
\end{pgfpicture}%
\makeatother%
\endgroup%
}
    \caption{Drift velocity as a function of flashing time.}
    \label{fig:drift_velocitytest}
\end{figure}

\begin{figure}[H]
    \centering
    \resizebox{\columnwidth}{!}{%% Creator: Matplotlib, PGF backend
%%
%% To include the figure in your LaTeX document, write
%%   \input{<filename>.pgf}
%%
%% Make sure the required packages are loaded in your preamble
%%   \usepackage{pgf}
%%
%% Also ensure that all the required font packages are loaded; for instance,
%% the lmodern package is sometimes necessary when using math font.
%%   \usepackage{lmodern}
%%
%% Figures using additional raster images can only be included by \input if
%% they are in the same directory as the main LaTeX file. For loading figures
%% from other directories you can use the `import` package
%%   \usepackage{import}
%%
%% and then include the figures with
%%   \import{<path to file>}{<filename>.pgf}
%%
%% Matplotlib used the following preamble
%%   \def\mathdefault#1{#1}
%%   \everymath=\expandafter{\the\everymath\displaystyle}
%%   \IfFileExists{scrextend.sty}{
%%     \usepackage[fontsize=10.000000pt]{scrextend}
%%   }{
%%     \renewcommand{\normalsize}{\fontsize{10.000000}{12.000000}\selectfont}
%%     \normalsize
%%   }
%%   
%%   \makeatletter\@ifpackageloaded{underscore}{}{\usepackage[strings]{underscore}}\makeatother
%%
\begingroup%
\makeatletter%
\begin{pgfpicture}%
\pgfpathrectangle{\pgfpointorigin}{\pgfqpoint{8.000000in}{9.000000in}}%
\pgfusepath{use as bounding box, clip}%
\begin{pgfscope}%
\pgfsetbuttcap%
\pgfsetmiterjoin%
\definecolor{currentfill}{rgb}{1.000000,1.000000,1.000000}%
\pgfsetfillcolor{currentfill}%
\pgfsetlinewidth{0.000000pt}%
\definecolor{currentstroke}{rgb}{1.000000,1.000000,1.000000}%
\pgfsetstrokecolor{currentstroke}%
\pgfsetdash{}{0pt}%
\pgfpathmoveto{\pgfqpoint{0.000000in}{0.000000in}}%
\pgfpathlineto{\pgfqpoint{8.000000in}{0.000000in}}%
\pgfpathlineto{\pgfqpoint{8.000000in}{9.000000in}}%
\pgfpathlineto{\pgfqpoint{0.000000in}{9.000000in}}%
\pgfpathlineto{\pgfqpoint{0.000000in}{0.000000in}}%
\pgfpathclose%
\pgfusepath{fill}%
\end{pgfscope}%
\begin{pgfscope}%
\pgfsetbuttcap%
\pgfsetmiterjoin%
\definecolor{currentfill}{rgb}{1.000000,1.000000,1.000000}%
\pgfsetfillcolor{currentfill}%
\pgfsetlinewidth{0.000000pt}%
\definecolor{currentstroke}{rgb}{0.000000,0.000000,0.000000}%
\pgfsetstrokecolor{currentstroke}%
\pgfsetstrokeopacity{0.000000}%
\pgfsetdash{}{0pt}%
\pgfpathmoveto{\pgfqpoint{0.787430in}{6.294627in}}%
\pgfpathlineto{\pgfqpoint{7.850000in}{6.294627in}}%
\pgfpathlineto{\pgfqpoint{7.850000in}{8.850000in}}%
\pgfpathlineto{\pgfqpoint{0.787430in}{8.850000in}}%
\pgfpathlineto{\pgfqpoint{0.787430in}{6.294627in}}%
\pgfpathclose%
\pgfusepath{fill}%
\end{pgfscope}%
\begin{pgfscope}%
\pgfpathrectangle{\pgfqpoint{0.787430in}{6.294627in}}{\pgfqpoint{7.062570in}{2.555373in}}%
\pgfusepath{clip}%
\pgfsetbuttcap%
\pgfsetmiterjoin%
\definecolor{currentfill}{rgb}{0.121569,0.466667,0.705882}%
\pgfsetfillcolor{currentfill}%
\pgfsetlinewidth{0.000000pt}%
\definecolor{currentstroke}{rgb}{0.000000,0.000000,0.000000}%
\pgfsetstrokecolor{currentstroke}%
\pgfsetstrokeopacity{0.000000}%
\pgfsetdash{}{0pt}%
\pgfpathmoveto{\pgfqpoint{2.572192in}{6.294627in}}%
\pgfpathlineto{\pgfqpoint{2.613063in}{6.294627in}}%
\pgfpathlineto{\pgfqpoint{2.613063in}{6.509364in}}%
\pgfpathlineto{\pgfqpoint{2.572192in}{6.509364in}}%
\pgfpathlineto{\pgfqpoint{2.572192in}{6.294627in}}%
\pgfpathclose%
\pgfusepath{fill}%
\end{pgfscope}%
\begin{pgfscope}%
\pgfpathrectangle{\pgfqpoint{0.787430in}{6.294627in}}{\pgfqpoint{7.062570in}{2.555373in}}%
\pgfusepath{clip}%
\pgfsetbuttcap%
\pgfsetmiterjoin%
\definecolor{currentfill}{rgb}{0.121569,0.466667,0.705882}%
\pgfsetfillcolor{currentfill}%
\pgfsetlinewidth{0.000000pt}%
\definecolor{currentstroke}{rgb}{0.000000,0.000000,0.000000}%
\pgfsetstrokecolor{currentstroke}%
\pgfsetstrokeopacity{0.000000}%
\pgfsetdash{}{0pt}%
\pgfpathmoveto{\pgfqpoint{2.613063in}{6.294627in}}%
\pgfpathlineto{\pgfqpoint{2.653934in}{6.294627in}}%
\pgfpathlineto{\pgfqpoint{2.653934in}{6.294627in}}%
\pgfpathlineto{\pgfqpoint{2.613063in}{6.294627in}}%
\pgfpathlineto{\pgfqpoint{2.613063in}{6.294627in}}%
\pgfpathclose%
\pgfusepath{fill}%
\end{pgfscope}%
\begin{pgfscope}%
\pgfpathrectangle{\pgfqpoint{0.787430in}{6.294627in}}{\pgfqpoint{7.062570in}{2.555373in}}%
\pgfusepath{clip}%
\pgfsetbuttcap%
\pgfsetmiterjoin%
\definecolor{currentfill}{rgb}{0.121569,0.466667,0.705882}%
\pgfsetfillcolor{currentfill}%
\pgfsetlinewidth{0.000000pt}%
\definecolor{currentstroke}{rgb}{0.000000,0.000000,0.000000}%
\pgfsetstrokecolor{currentstroke}%
\pgfsetstrokeopacity{0.000000}%
\pgfsetdash{}{0pt}%
\pgfpathmoveto{\pgfqpoint{2.653934in}{6.294627in}}%
\pgfpathlineto{\pgfqpoint{2.694805in}{6.294627in}}%
\pgfpathlineto{\pgfqpoint{2.694805in}{6.294627in}}%
\pgfpathlineto{\pgfqpoint{2.653934in}{6.294627in}}%
\pgfpathlineto{\pgfqpoint{2.653934in}{6.294627in}}%
\pgfpathclose%
\pgfusepath{fill}%
\end{pgfscope}%
\begin{pgfscope}%
\pgfpathrectangle{\pgfqpoint{0.787430in}{6.294627in}}{\pgfqpoint{7.062570in}{2.555373in}}%
\pgfusepath{clip}%
\pgfsetbuttcap%
\pgfsetmiterjoin%
\definecolor{currentfill}{rgb}{0.121569,0.466667,0.705882}%
\pgfsetfillcolor{currentfill}%
\pgfsetlinewidth{0.000000pt}%
\definecolor{currentstroke}{rgb}{0.000000,0.000000,0.000000}%
\pgfsetstrokecolor{currentstroke}%
\pgfsetstrokeopacity{0.000000}%
\pgfsetdash{}{0pt}%
\pgfpathmoveto{\pgfqpoint{2.694805in}{6.294627in}}%
\pgfpathlineto{\pgfqpoint{2.735676in}{6.294627in}}%
\pgfpathlineto{\pgfqpoint{2.735676in}{6.294627in}}%
\pgfpathlineto{\pgfqpoint{2.694805in}{6.294627in}}%
\pgfpathlineto{\pgfqpoint{2.694805in}{6.294627in}}%
\pgfpathclose%
\pgfusepath{fill}%
\end{pgfscope}%
\begin{pgfscope}%
\pgfpathrectangle{\pgfqpoint{0.787430in}{6.294627in}}{\pgfqpoint{7.062570in}{2.555373in}}%
\pgfusepath{clip}%
\pgfsetbuttcap%
\pgfsetmiterjoin%
\definecolor{currentfill}{rgb}{0.121569,0.466667,0.705882}%
\pgfsetfillcolor{currentfill}%
\pgfsetlinewidth{0.000000pt}%
\definecolor{currentstroke}{rgb}{0.000000,0.000000,0.000000}%
\pgfsetstrokecolor{currentstroke}%
\pgfsetstrokeopacity{0.000000}%
\pgfsetdash{}{0pt}%
\pgfpathmoveto{\pgfqpoint{2.735676in}{6.294627in}}%
\pgfpathlineto{\pgfqpoint{2.776548in}{6.294627in}}%
\pgfpathlineto{\pgfqpoint{2.776548in}{6.294627in}}%
\pgfpathlineto{\pgfqpoint{2.735676in}{6.294627in}}%
\pgfpathlineto{\pgfqpoint{2.735676in}{6.294627in}}%
\pgfpathclose%
\pgfusepath{fill}%
\end{pgfscope}%
\begin{pgfscope}%
\pgfpathrectangle{\pgfqpoint{0.787430in}{6.294627in}}{\pgfqpoint{7.062570in}{2.555373in}}%
\pgfusepath{clip}%
\pgfsetbuttcap%
\pgfsetmiterjoin%
\definecolor{currentfill}{rgb}{0.121569,0.466667,0.705882}%
\pgfsetfillcolor{currentfill}%
\pgfsetlinewidth{0.000000pt}%
\definecolor{currentstroke}{rgb}{0.000000,0.000000,0.000000}%
\pgfsetstrokecolor{currentstroke}%
\pgfsetstrokeopacity{0.000000}%
\pgfsetdash{}{0pt}%
\pgfpathmoveto{\pgfqpoint{2.776548in}{6.294627in}}%
\pgfpathlineto{\pgfqpoint{2.817419in}{6.294627in}}%
\pgfpathlineto{\pgfqpoint{2.817419in}{6.366206in}}%
\pgfpathlineto{\pgfqpoint{2.776548in}{6.366206in}}%
\pgfpathlineto{\pgfqpoint{2.776548in}{6.294627in}}%
\pgfpathclose%
\pgfusepath{fill}%
\end{pgfscope}%
\begin{pgfscope}%
\pgfpathrectangle{\pgfqpoint{0.787430in}{6.294627in}}{\pgfqpoint{7.062570in}{2.555373in}}%
\pgfusepath{clip}%
\pgfsetbuttcap%
\pgfsetmiterjoin%
\definecolor{currentfill}{rgb}{0.121569,0.466667,0.705882}%
\pgfsetfillcolor{currentfill}%
\pgfsetlinewidth{0.000000pt}%
\definecolor{currentstroke}{rgb}{0.000000,0.000000,0.000000}%
\pgfsetstrokecolor{currentstroke}%
\pgfsetstrokeopacity{0.000000}%
\pgfsetdash{}{0pt}%
\pgfpathmoveto{\pgfqpoint{2.817419in}{6.294627in}}%
\pgfpathlineto{\pgfqpoint{2.858290in}{6.294627in}}%
\pgfpathlineto{\pgfqpoint{2.858290in}{6.437785in}}%
\pgfpathlineto{\pgfqpoint{2.817419in}{6.437785in}}%
\pgfpathlineto{\pgfqpoint{2.817419in}{6.294627in}}%
\pgfpathclose%
\pgfusepath{fill}%
\end{pgfscope}%
\begin{pgfscope}%
\pgfpathrectangle{\pgfqpoint{0.787430in}{6.294627in}}{\pgfqpoint{7.062570in}{2.555373in}}%
\pgfusepath{clip}%
\pgfsetbuttcap%
\pgfsetmiterjoin%
\definecolor{currentfill}{rgb}{0.121569,0.466667,0.705882}%
\pgfsetfillcolor{currentfill}%
\pgfsetlinewidth{0.000000pt}%
\definecolor{currentstroke}{rgb}{0.000000,0.000000,0.000000}%
\pgfsetstrokecolor{currentstroke}%
\pgfsetstrokeopacity{0.000000}%
\pgfsetdash{}{0pt}%
\pgfpathmoveto{\pgfqpoint{2.858290in}{6.294627in}}%
\pgfpathlineto{\pgfqpoint{2.899161in}{6.294627in}}%
\pgfpathlineto{\pgfqpoint{2.899161in}{6.437785in}}%
\pgfpathlineto{\pgfqpoint{2.858290in}{6.437785in}}%
\pgfpathlineto{\pgfqpoint{2.858290in}{6.294627in}}%
\pgfpathclose%
\pgfusepath{fill}%
\end{pgfscope}%
\begin{pgfscope}%
\pgfpathrectangle{\pgfqpoint{0.787430in}{6.294627in}}{\pgfqpoint{7.062570in}{2.555373in}}%
\pgfusepath{clip}%
\pgfsetbuttcap%
\pgfsetmiterjoin%
\definecolor{currentfill}{rgb}{0.121569,0.466667,0.705882}%
\pgfsetfillcolor{currentfill}%
\pgfsetlinewidth{0.000000pt}%
\definecolor{currentstroke}{rgb}{0.000000,0.000000,0.000000}%
\pgfsetstrokecolor{currentstroke}%
\pgfsetstrokeopacity{0.000000}%
\pgfsetdash{}{0pt}%
\pgfpathmoveto{\pgfqpoint{2.899161in}{6.294627in}}%
\pgfpathlineto{\pgfqpoint{2.940032in}{6.294627in}}%
\pgfpathlineto{\pgfqpoint{2.940032in}{6.366206in}}%
\pgfpathlineto{\pgfqpoint{2.899161in}{6.366206in}}%
\pgfpathlineto{\pgfqpoint{2.899161in}{6.294627in}}%
\pgfpathclose%
\pgfusepath{fill}%
\end{pgfscope}%
\begin{pgfscope}%
\pgfpathrectangle{\pgfqpoint{0.787430in}{6.294627in}}{\pgfqpoint{7.062570in}{2.555373in}}%
\pgfusepath{clip}%
\pgfsetbuttcap%
\pgfsetmiterjoin%
\definecolor{currentfill}{rgb}{0.121569,0.466667,0.705882}%
\pgfsetfillcolor{currentfill}%
\pgfsetlinewidth{0.000000pt}%
\definecolor{currentstroke}{rgb}{0.000000,0.000000,0.000000}%
\pgfsetstrokecolor{currentstroke}%
\pgfsetstrokeopacity{0.000000}%
\pgfsetdash{}{0pt}%
\pgfpathmoveto{\pgfqpoint{2.940032in}{6.294627in}}%
\pgfpathlineto{\pgfqpoint{2.980904in}{6.294627in}}%
\pgfpathlineto{\pgfqpoint{2.980904in}{6.366206in}}%
\pgfpathlineto{\pgfqpoint{2.940032in}{6.366206in}}%
\pgfpathlineto{\pgfqpoint{2.940032in}{6.294627in}}%
\pgfpathclose%
\pgfusepath{fill}%
\end{pgfscope}%
\begin{pgfscope}%
\pgfpathrectangle{\pgfqpoint{0.787430in}{6.294627in}}{\pgfqpoint{7.062570in}{2.555373in}}%
\pgfusepath{clip}%
\pgfsetbuttcap%
\pgfsetmiterjoin%
\definecolor{currentfill}{rgb}{0.121569,0.466667,0.705882}%
\pgfsetfillcolor{currentfill}%
\pgfsetlinewidth{0.000000pt}%
\definecolor{currentstroke}{rgb}{0.000000,0.000000,0.000000}%
\pgfsetstrokecolor{currentstroke}%
\pgfsetstrokeopacity{0.000000}%
\pgfsetdash{}{0pt}%
\pgfpathmoveto{\pgfqpoint{2.980904in}{6.294627in}}%
\pgfpathlineto{\pgfqpoint{3.021775in}{6.294627in}}%
\pgfpathlineto{\pgfqpoint{3.021775in}{6.509364in}}%
\pgfpathlineto{\pgfqpoint{2.980904in}{6.509364in}}%
\pgfpathlineto{\pgfqpoint{2.980904in}{6.294627in}}%
\pgfpathclose%
\pgfusepath{fill}%
\end{pgfscope}%
\begin{pgfscope}%
\pgfpathrectangle{\pgfqpoint{0.787430in}{6.294627in}}{\pgfqpoint{7.062570in}{2.555373in}}%
\pgfusepath{clip}%
\pgfsetbuttcap%
\pgfsetmiterjoin%
\definecolor{currentfill}{rgb}{0.121569,0.466667,0.705882}%
\pgfsetfillcolor{currentfill}%
\pgfsetlinewidth{0.000000pt}%
\definecolor{currentstroke}{rgb}{0.000000,0.000000,0.000000}%
\pgfsetstrokecolor{currentstroke}%
\pgfsetstrokeopacity{0.000000}%
\pgfsetdash{}{0pt}%
\pgfpathmoveto{\pgfqpoint{3.021775in}{6.294627in}}%
\pgfpathlineto{\pgfqpoint{3.062646in}{6.294627in}}%
\pgfpathlineto{\pgfqpoint{3.062646in}{6.580943in}}%
\pgfpathlineto{\pgfqpoint{3.021775in}{6.580943in}}%
\pgfpathlineto{\pgfqpoint{3.021775in}{6.294627in}}%
\pgfpathclose%
\pgfusepath{fill}%
\end{pgfscope}%
\begin{pgfscope}%
\pgfpathrectangle{\pgfqpoint{0.787430in}{6.294627in}}{\pgfqpoint{7.062570in}{2.555373in}}%
\pgfusepath{clip}%
\pgfsetbuttcap%
\pgfsetmiterjoin%
\definecolor{currentfill}{rgb}{0.121569,0.466667,0.705882}%
\pgfsetfillcolor{currentfill}%
\pgfsetlinewidth{0.000000pt}%
\definecolor{currentstroke}{rgb}{0.000000,0.000000,0.000000}%
\pgfsetstrokecolor{currentstroke}%
\pgfsetstrokeopacity{0.000000}%
\pgfsetdash{}{0pt}%
\pgfpathmoveto{\pgfqpoint{3.062646in}{6.294627in}}%
\pgfpathlineto{\pgfqpoint{3.103517in}{6.294627in}}%
\pgfpathlineto{\pgfqpoint{3.103517in}{6.366206in}}%
\pgfpathlineto{\pgfqpoint{3.062646in}{6.366206in}}%
\pgfpathlineto{\pgfqpoint{3.062646in}{6.294627in}}%
\pgfpathclose%
\pgfusepath{fill}%
\end{pgfscope}%
\begin{pgfscope}%
\pgfpathrectangle{\pgfqpoint{0.787430in}{6.294627in}}{\pgfqpoint{7.062570in}{2.555373in}}%
\pgfusepath{clip}%
\pgfsetbuttcap%
\pgfsetmiterjoin%
\definecolor{currentfill}{rgb}{0.121569,0.466667,0.705882}%
\pgfsetfillcolor{currentfill}%
\pgfsetlinewidth{0.000000pt}%
\definecolor{currentstroke}{rgb}{0.000000,0.000000,0.000000}%
\pgfsetstrokecolor{currentstroke}%
\pgfsetstrokeopacity{0.000000}%
\pgfsetdash{}{0pt}%
\pgfpathmoveto{\pgfqpoint{3.103517in}{6.294627in}}%
\pgfpathlineto{\pgfqpoint{3.144388in}{6.294627in}}%
\pgfpathlineto{\pgfqpoint{3.144388in}{6.652522in}}%
\pgfpathlineto{\pgfqpoint{3.103517in}{6.652522in}}%
\pgfpathlineto{\pgfqpoint{3.103517in}{6.294627in}}%
\pgfpathclose%
\pgfusepath{fill}%
\end{pgfscope}%
\begin{pgfscope}%
\pgfpathrectangle{\pgfqpoint{0.787430in}{6.294627in}}{\pgfqpoint{7.062570in}{2.555373in}}%
\pgfusepath{clip}%
\pgfsetbuttcap%
\pgfsetmiterjoin%
\definecolor{currentfill}{rgb}{0.121569,0.466667,0.705882}%
\pgfsetfillcolor{currentfill}%
\pgfsetlinewidth{0.000000pt}%
\definecolor{currentstroke}{rgb}{0.000000,0.000000,0.000000}%
\pgfsetstrokecolor{currentstroke}%
\pgfsetstrokeopacity{0.000000}%
\pgfsetdash{}{0pt}%
\pgfpathmoveto{\pgfqpoint{3.144388in}{6.294627in}}%
\pgfpathlineto{\pgfqpoint{3.185260in}{6.294627in}}%
\pgfpathlineto{\pgfqpoint{3.185260in}{6.509364in}}%
\pgfpathlineto{\pgfqpoint{3.144388in}{6.509364in}}%
\pgfpathlineto{\pgfqpoint{3.144388in}{6.294627in}}%
\pgfpathclose%
\pgfusepath{fill}%
\end{pgfscope}%
\begin{pgfscope}%
\pgfpathrectangle{\pgfqpoint{0.787430in}{6.294627in}}{\pgfqpoint{7.062570in}{2.555373in}}%
\pgfusepath{clip}%
\pgfsetbuttcap%
\pgfsetmiterjoin%
\definecolor{currentfill}{rgb}{0.121569,0.466667,0.705882}%
\pgfsetfillcolor{currentfill}%
\pgfsetlinewidth{0.000000pt}%
\definecolor{currentstroke}{rgb}{0.000000,0.000000,0.000000}%
\pgfsetstrokecolor{currentstroke}%
\pgfsetstrokeopacity{0.000000}%
\pgfsetdash{}{0pt}%
\pgfpathmoveto{\pgfqpoint{3.185260in}{6.294627in}}%
\pgfpathlineto{\pgfqpoint{3.226131in}{6.294627in}}%
\pgfpathlineto{\pgfqpoint{3.226131in}{6.652522in}}%
\pgfpathlineto{\pgfqpoint{3.185260in}{6.652522in}}%
\pgfpathlineto{\pgfqpoint{3.185260in}{6.294627in}}%
\pgfpathclose%
\pgfusepath{fill}%
\end{pgfscope}%
\begin{pgfscope}%
\pgfpathrectangle{\pgfqpoint{0.787430in}{6.294627in}}{\pgfqpoint{7.062570in}{2.555373in}}%
\pgfusepath{clip}%
\pgfsetbuttcap%
\pgfsetmiterjoin%
\definecolor{currentfill}{rgb}{0.121569,0.466667,0.705882}%
\pgfsetfillcolor{currentfill}%
\pgfsetlinewidth{0.000000pt}%
\definecolor{currentstroke}{rgb}{0.000000,0.000000,0.000000}%
\pgfsetstrokecolor{currentstroke}%
\pgfsetstrokeopacity{0.000000}%
\pgfsetdash{}{0pt}%
\pgfpathmoveto{\pgfqpoint{3.226131in}{6.294627in}}%
\pgfpathlineto{\pgfqpoint{3.267002in}{6.294627in}}%
\pgfpathlineto{\pgfqpoint{3.267002in}{6.509364in}}%
\pgfpathlineto{\pgfqpoint{3.226131in}{6.509364in}}%
\pgfpathlineto{\pgfqpoint{3.226131in}{6.294627in}}%
\pgfpathclose%
\pgfusepath{fill}%
\end{pgfscope}%
\begin{pgfscope}%
\pgfpathrectangle{\pgfqpoint{0.787430in}{6.294627in}}{\pgfqpoint{7.062570in}{2.555373in}}%
\pgfusepath{clip}%
\pgfsetbuttcap%
\pgfsetmiterjoin%
\definecolor{currentfill}{rgb}{0.121569,0.466667,0.705882}%
\pgfsetfillcolor{currentfill}%
\pgfsetlinewidth{0.000000pt}%
\definecolor{currentstroke}{rgb}{0.000000,0.000000,0.000000}%
\pgfsetstrokecolor{currentstroke}%
\pgfsetstrokeopacity{0.000000}%
\pgfsetdash{}{0pt}%
\pgfpathmoveto{\pgfqpoint{3.267002in}{6.294627in}}%
\pgfpathlineto{\pgfqpoint{3.307873in}{6.294627in}}%
\pgfpathlineto{\pgfqpoint{3.307873in}{6.724101in}}%
\pgfpathlineto{\pgfqpoint{3.267002in}{6.724101in}}%
\pgfpathlineto{\pgfqpoint{3.267002in}{6.294627in}}%
\pgfpathclose%
\pgfusepath{fill}%
\end{pgfscope}%
\begin{pgfscope}%
\pgfpathrectangle{\pgfqpoint{0.787430in}{6.294627in}}{\pgfqpoint{7.062570in}{2.555373in}}%
\pgfusepath{clip}%
\pgfsetbuttcap%
\pgfsetmiterjoin%
\definecolor{currentfill}{rgb}{0.121569,0.466667,0.705882}%
\pgfsetfillcolor{currentfill}%
\pgfsetlinewidth{0.000000pt}%
\definecolor{currentstroke}{rgb}{0.000000,0.000000,0.000000}%
\pgfsetstrokecolor{currentstroke}%
\pgfsetstrokeopacity{0.000000}%
\pgfsetdash{}{0pt}%
\pgfpathmoveto{\pgfqpoint{3.307873in}{6.294627in}}%
\pgfpathlineto{\pgfqpoint{3.348744in}{6.294627in}}%
\pgfpathlineto{\pgfqpoint{3.348744in}{6.795680in}}%
\pgfpathlineto{\pgfqpoint{3.307873in}{6.795680in}}%
\pgfpathlineto{\pgfqpoint{3.307873in}{6.294627in}}%
\pgfpathclose%
\pgfusepath{fill}%
\end{pgfscope}%
\begin{pgfscope}%
\pgfpathrectangle{\pgfqpoint{0.787430in}{6.294627in}}{\pgfqpoint{7.062570in}{2.555373in}}%
\pgfusepath{clip}%
\pgfsetbuttcap%
\pgfsetmiterjoin%
\definecolor{currentfill}{rgb}{0.121569,0.466667,0.705882}%
\pgfsetfillcolor{currentfill}%
\pgfsetlinewidth{0.000000pt}%
\definecolor{currentstroke}{rgb}{0.000000,0.000000,0.000000}%
\pgfsetstrokecolor{currentstroke}%
\pgfsetstrokeopacity{0.000000}%
\pgfsetdash{}{0pt}%
\pgfpathmoveto{\pgfqpoint{3.348744in}{6.294627in}}%
\pgfpathlineto{\pgfqpoint{3.389616in}{6.294627in}}%
\pgfpathlineto{\pgfqpoint{3.389616in}{7.081997in}}%
\pgfpathlineto{\pgfqpoint{3.348744in}{7.081997in}}%
\pgfpathlineto{\pgfqpoint{3.348744in}{6.294627in}}%
\pgfpathclose%
\pgfusepath{fill}%
\end{pgfscope}%
\begin{pgfscope}%
\pgfpathrectangle{\pgfqpoint{0.787430in}{6.294627in}}{\pgfqpoint{7.062570in}{2.555373in}}%
\pgfusepath{clip}%
\pgfsetbuttcap%
\pgfsetmiterjoin%
\definecolor{currentfill}{rgb}{0.121569,0.466667,0.705882}%
\pgfsetfillcolor{currentfill}%
\pgfsetlinewidth{0.000000pt}%
\definecolor{currentstroke}{rgb}{0.000000,0.000000,0.000000}%
\pgfsetstrokecolor{currentstroke}%
\pgfsetstrokeopacity{0.000000}%
\pgfsetdash{}{0pt}%
\pgfpathmoveto{\pgfqpoint{3.389616in}{6.294627in}}%
\pgfpathlineto{\pgfqpoint{3.430487in}{6.294627in}}%
\pgfpathlineto{\pgfqpoint{3.430487in}{6.795680in}}%
\pgfpathlineto{\pgfqpoint{3.389616in}{6.795680in}}%
\pgfpathlineto{\pgfqpoint{3.389616in}{6.294627in}}%
\pgfpathclose%
\pgfusepath{fill}%
\end{pgfscope}%
\begin{pgfscope}%
\pgfpathrectangle{\pgfqpoint{0.787430in}{6.294627in}}{\pgfqpoint{7.062570in}{2.555373in}}%
\pgfusepath{clip}%
\pgfsetbuttcap%
\pgfsetmiterjoin%
\definecolor{currentfill}{rgb}{0.121569,0.466667,0.705882}%
\pgfsetfillcolor{currentfill}%
\pgfsetlinewidth{0.000000pt}%
\definecolor{currentstroke}{rgb}{0.000000,0.000000,0.000000}%
\pgfsetstrokecolor{currentstroke}%
\pgfsetstrokeopacity{0.000000}%
\pgfsetdash{}{0pt}%
\pgfpathmoveto{\pgfqpoint{3.430487in}{6.294627in}}%
\pgfpathlineto{\pgfqpoint{3.471358in}{6.294627in}}%
\pgfpathlineto{\pgfqpoint{3.471358in}{6.724101in}}%
\pgfpathlineto{\pgfqpoint{3.430487in}{6.724101in}}%
\pgfpathlineto{\pgfqpoint{3.430487in}{6.294627in}}%
\pgfpathclose%
\pgfusepath{fill}%
\end{pgfscope}%
\begin{pgfscope}%
\pgfpathrectangle{\pgfqpoint{0.787430in}{6.294627in}}{\pgfqpoint{7.062570in}{2.555373in}}%
\pgfusepath{clip}%
\pgfsetbuttcap%
\pgfsetmiterjoin%
\definecolor{currentfill}{rgb}{0.121569,0.466667,0.705882}%
\pgfsetfillcolor{currentfill}%
\pgfsetlinewidth{0.000000pt}%
\definecolor{currentstroke}{rgb}{0.000000,0.000000,0.000000}%
\pgfsetstrokecolor{currentstroke}%
\pgfsetstrokeopacity{0.000000}%
\pgfsetdash{}{0pt}%
\pgfpathmoveto{\pgfqpoint{3.471358in}{6.294627in}}%
\pgfpathlineto{\pgfqpoint{3.512229in}{6.294627in}}%
\pgfpathlineto{\pgfqpoint{3.512229in}{7.153576in}}%
\pgfpathlineto{\pgfqpoint{3.471358in}{7.153576in}}%
\pgfpathlineto{\pgfqpoint{3.471358in}{6.294627in}}%
\pgfpathclose%
\pgfusepath{fill}%
\end{pgfscope}%
\begin{pgfscope}%
\pgfpathrectangle{\pgfqpoint{0.787430in}{6.294627in}}{\pgfqpoint{7.062570in}{2.555373in}}%
\pgfusepath{clip}%
\pgfsetbuttcap%
\pgfsetmiterjoin%
\definecolor{currentfill}{rgb}{0.121569,0.466667,0.705882}%
\pgfsetfillcolor{currentfill}%
\pgfsetlinewidth{0.000000pt}%
\definecolor{currentstroke}{rgb}{0.000000,0.000000,0.000000}%
\pgfsetstrokecolor{currentstroke}%
\pgfsetstrokeopacity{0.000000}%
\pgfsetdash{}{0pt}%
\pgfpathmoveto{\pgfqpoint{3.512229in}{6.294627in}}%
\pgfpathlineto{\pgfqpoint{3.553100in}{6.294627in}}%
\pgfpathlineto{\pgfqpoint{3.553100in}{7.225155in}}%
\pgfpathlineto{\pgfqpoint{3.512229in}{7.225155in}}%
\pgfpathlineto{\pgfqpoint{3.512229in}{6.294627in}}%
\pgfpathclose%
\pgfusepath{fill}%
\end{pgfscope}%
\begin{pgfscope}%
\pgfpathrectangle{\pgfqpoint{0.787430in}{6.294627in}}{\pgfqpoint{7.062570in}{2.555373in}}%
\pgfusepath{clip}%
\pgfsetbuttcap%
\pgfsetmiterjoin%
\definecolor{currentfill}{rgb}{0.121569,0.466667,0.705882}%
\pgfsetfillcolor{currentfill}%
\pgfsetlinewidth{0.000000pt}%
\definecolor{currentstroke}{rgb}{0.000000,0.000000,0.000000}%
\pgfsetstrokecolor{currentstroke}%
\pgfsetstrokeopacity{0.000000}%
\pgfsetdash{}{0pt}%
\pgfpathmoveto{\pgfqpoint{3.553100in}{6.294627in}}%
\pgfpathlineto{\pgfqpoint{3.593972in}{6.294627in}}%
\pgfpathlineto{\pgfqpoint{3.593972in}{6.938838in}}%
\pgfpathlineto{\pgfqpoint{3.553100in}{6.938838in}}%
\pgfpathlineto{\pgfqpoint{3.553100in}{6.294627in}}%
\pgfpathclose%
\pgfusepath{fill}%
\end{pgfscope}%
\begin{pgfscope}%
\pgfpathrectangle{\pgfqpoint{0.787430in}{6.294627in}}{\pgfqpoint{7.062570in}{2.555373in}}%
\pgfusepath{clip}%
\pgfsetbuttcap%
\pgfsetmiterjoin%
\definecolor{currentfill}{rgb}{0.121569,0.466667,0.705882}%
\pgfsetfillcolor{currentfill}%
\pgfsetlinewidth{0.000000pt}%
\definecolor{currentstroke}{rgb}{0.000000,0.000000,0.000000}%
\pgfsetstrokecolor{currentstroke}%
\pgfsetstrokeopacity{0.000000}%
\pgfsetdash{}{0pt}%
\pgfpathmoveto{\pgfqpoint{3.593972in}{6.294627in}}%
\pgfpathlineto{\pgfqpoint{3.634843in}{6.294627in}}%
\pgfpathlineto{\pgfqpoint{3.634843in}{7.439892in}}%
\pgfpathlineto{\pgfqpoint{3.593972in}{7.439892in}}%
\pgfpathlineto{\pgfqpoint{3.593972in}{6.294627in}}%
\pgfpathclose%
\pgfusepath{fill}%
\end{pgfscope}%
\begin{pgfscope}%
\pgfpathrectangle{\pgfqpoint{0.787430in}{6.294627in}}{\pgfqpoint{7.062570in}{2.555373in}}%
\pgfusepath{clip}%
\pgfsetbuttcap%
\pgfsetmiterjoin%
\definecolor{currentfill}{rgb}{0.121569,0.466667,0.705882}%
\pgfsetfillcolor{currentfill}%
\pgfsetlinewidth{0.000000pt}%
\definecolor{currentstroke}{rgb}{0.000000,0.000000,0.000000}%
\pgfsetstrokecolor{currentstroke}%
\pgfsetstrokeopacity{0.000000}%
\pgfsetdash{}{0pt}%
\pgfpathmoveto{\pgfqpoint{3.634843in}{6.294627in}}%
\pgfpathlineto{\pgfqpoint{3.675714in}{6.294627in}}%
\pgfpathlineto{\pgfqpoint{3.675714in}{7.010418in}}%
\pgfpathlineto{\pgfqpoint{3.634843in}{7.010418in}}%
\pgfpathlineto{\pgfqpoint{3.634843in}{6.294627in}}%
\pgfpathclose%
\pgfusepath{fill}%
\end{pgfscope}%
\begin{pgfscope}%
\pgfpathrectangle{\pgfqpoint{0.787430in}{6.294627in}}{\pgfqpoint{7.062570in}{2.555373in}}%
\pgfusepath{clip}%
\pgfsetbuttcap%
\pgfsetmiterjoin%
\definecolor{currentfill}{rgb}{0.121569,0.466667,0.705882}%
\pgfsetfillcolor{currentfill}%
\pgfsetlinewidth{0.000000pt}%
\definecolor{currentstroke}{rgb}{0.000000,0.000000,0.000000}%
\pgfsetstrokecolor{currentstroke}%
\pgfsetstrokeopacity{0.000000}%
\pgfsetdash{}{0pt}%
\pgfpathmoveto{\pgfqpoint{3.675714in}{6.294627in}}%
\pgfpathlineto{\pgfqpoint{3.716585in}{6.294627in}}%
\pgfpathlineto{\pgfqpoint{3.716585in}{6.867259in}}%
\pgfpathlineto{\pgfqpoint{3.675714in}{6.867259in}}%
\pgfpathlineto{\pgfqpoint{3.675714in}{6.294627in}}%
\pgfpathclose%
\pgfusepath{fill}%
\end{pgfscope}%
\begin{pgfscope}%
\pgfpathrectangle{\pgfqpoint{0.787430in}{6.294627in}}{\pgfqpoint{7.062570in}{2.555373in}}%
\pgfusepath{clip}%
\pgfsetbuttcap%
\pgfsetmiterjoin%
\definecolor{currentfill}{rgb}{0.121569,0.466667,0.705882}%
\pgfsetfillcolor{currentfill}%
\pgfsetlinewidth{0.000000pt}%
\definecolor{currentstroke}{rgb}{0.000000,0.000000,0.000000}%
\pgfsetstrokecolor{currentstroke}%
\pgfsetstrokeopacity{0.000000}%
\pgfsetdash{}{0pt}%
\pgfpathmoveto{\pgfqpoint{3.716585in}{6.294627in}}%
\pgfpathlineto{\pgfqpoint{3.757456in}{6.294627in}}%
\pgfpathlineto{\pgfqpoint{3.757456in}{7.511471in}}%
\pgfpathlineto{\pgfqpoint{3.716585in}{7.511471in}}%
\pgfpathlineto{\pgfqpoint{3.716585in}{6.294627in}}%
\pgfpathclose%
\pgfusepath{fill}%
\end{pgfscope}%
\begin{pgfscope}%
\pgfpathrectangle{\pgfqpoint{0.787430in}{6.294627in}}{\pgfqpoint{7.062570in}{2.555373in}}%
\pgfusepath{clip}%
\pgfsetbuttcap%
\pgfsetmiterjoin%
\definecolor{currentfill}{rgb}{0.121569,0.466667,0.705882}%
\pgfsetfillcolor{currentfill}%
\pgfsetlinewidth{0.000000pt}%
\definecolor{currentstroke}{rgb}{0.000000,0.000000,0.000000}%
\pgfsetstrokecolor{currentstroke}%
\pgfsetstrokeopacity{0.000000}%
\pgfsetdash{}{0pt}%
\pgfpathmoveto{\pgfqpoint{3.757456in}{6.294627in}}%
\pgfpathlineto{\pgfqpoint{3.798328in}{6.294627in}}%
\pgfpathlineto{\pgfqpoint{3.798328in}{7.153576in}}%
\pgfpathlineto{\pgfqpoint{3.757456in}{7.153576in}}%
\pgfpathlineto{\pgfqpoint{3.757456in}{6.294627in}}%
\pgfpathclose%
\pgfusepath{fill}%
\end{pgfscope}%
\begin{pgfscope}%
\pgfpathrectangle{\pgfqpoint{0.787430in}{6.294627in}}{\pgfqpoint{7.062570in}{2.555373in}}%
\pgfusepath{clip}%
\pgfsetbuttcap%
\pgfsetmiterjoin%
\definecolor{currentfill}{rgb}{0.121569,0.466667,0.705882}%
\pgfsetfillcolor{currentfill}%
\pgfsetlinewidth{0.000000pt}%
\definecolor{currentstroke}{rgb}{0.000000,0.000000,0.000000}%
\pgfsetstrokecolor{currentstroke}%
\pgfsetstrokeopacity{0.000000}%
\pgfsetdash{}{0pt}%
\pgfpathmoveto{\pgfqpoint{3.798328in}{6.294627in}}%
\pgfpathlineto{\pgfqpoint{3.839199in}{6.294627in}}%
\pgfpathlineto{\pgfqpoint{3.839199in}{7.153576in}}%
\pgfpathlineto{\pgfqpoint{3.798328in}{7.153576in}}%
\pgfpathlineto{\pgfqpoint{3.798328in}{6.294627in}}%
\pgfpathclose%
\pgfusepath{fill}%
\end{pgfscope}%
\begin{pgfscope}%
\pgfpathrectangle{\pgfqpoint{0.787430in}{6.294627in}}{\pgfqpoint{7.062570in}{2.555373in}}%
\pgfusepath{clip}%
\pgfsetbuttcap%
\pgfsetmiterjoin%
\definecolor{currentfill}{rgb}{0.121569,0.466667,0.705882}%
\pgfsetfillcolor{currentfill}%
\pgfsetlinewidth{0.000000pt}%
\definecolor{currentstroke}{rgb}{0.000000,0.000000,0.000000}%
\pgfsetstrokecolor{currentstroke}%
\pgfsetstrokeopacity{0.000000}%
\pgfsetdash{}{0pt}%
\pgfpathmoveto{\pgfqpoint{3.839199in}{6.294627in}}%
\pgfpathlineto{\pgfqpoint{3.880070in}{6.294627in}}%
\pgfpathlineto{\pgfqpoint{3.880070in}{8.513578in}}%
\pgfpathlineto{\pgfqpoint{3.839199in}{8.513578in}}%
\pgfpathlineto{\pgfqpoint{3.839199in}{6.294627in}}%
\pgfpathclose%
\pgfusepath{fill}%
\end{pgfscope}%
\begin{pgfscope}%
\pgfpathrectangle{\pgfqpoint{0.787430in}{6.294627in}}{\pgfqpoint{7.062570in}{2.555373in}}%
\pgfusepath{clip}%
\pgfsetbuttcap%
\pgfsetmiterjoin%
\definecolor{currentfill}{rgb}{0.121569,0.466667,0.705882}%
\pgfsetfillcolor{currentfill}%
\pgfsetlinewidth{0.000000pt}%
\definecolor{currentstroke}{rgb}{0.000000,0.000000,0.000000}%
\pgfsetstrokecolor{currentstroke}%
\pgfsetstrokeopacity{0.000000}%
\pgfsetdash{}{0pt}%
\pgfpathmoveto{\pgfqpoint{3.880070in}{6.294627in}}%
\pgfpathlineto{\pgfqpoint{3.920941in}{6.294627in}}%
\pgfpathlineto{\pgfqpoint{3.920941in}{7.583050in}}%
\pgfpathlineto{\pgfqpoint{3.880070in}{7.583050in}}%
\pgfpathlineto{\pgfqpoint{3.880070in}{6.294627in}}%
\pgfpathclose%
\pgfusepath{fill}%
\end{pgfscope}%
\begin{pgfscope}%
\pgfpathrectangle{\pgfqpoint{0.787430in}{6.294627in}}{\pgfqpoint{7.062570in}{2.555373in}}%
\pgfusepath{clip}%
\pgfsetbuttcap%
\pgfsetmiterjoin%
\definecolor{currentfill}{rgb}{0.121569,0.466667,0.705882}%
\pgfsetfillcolor{currentfill}%
\pgfsetlinewidth{0.000000pt}%
\definecolor{currentstroke}{rgb}{0.000000,0.000000,0.000000}%
\pgfsetstrokecolor{currentstroke}%
\pgfsetstrokeopacity{0.000000}%
\pgfsetdash{}{0pt}%
\pgfpathmoveto{\pgfqpoint{3.920941in}{6.294627in}}%
\pgfpathlineto{\pgfqpoint{3.961812in}{6.294627in}}%
\pgfpathlineto{\pgfqpoint{3.961812in}{8.585157in}}%
\pgfpathlineto{\pgfqpoint{3.920941in}{8.585157in}}%
\pgfpathlineto{\pgfqpoint{3.920941in}{6.294627in}}%
\pgfpathclose%
\pgfusepath{fill}%
\end{pgfscope}%
\begin{pgfscope}%
\pgfpathrectangle{\pgfqpoint{0.787430in}{6.294627in}}{\pgfqpoint{7.062570in}{2.555373in}}%
\pgfusepath{clip}%
\pgfsetbuttcap%
\pgfsetmiterjoin%
\definecolor{currentfill}{rgb}{0.121569,0.466667,0.705882}%
\pgfsetfillcolor{currentfill}%
\pgfsetlinewidth{0.000000pt}%
\definecolor{currentstroke}{rgb}{0.000000,0.000000,0.000000}%
\pgfsetstrokecolor{currentstroke}%
\pgfsetstrokeopacity{0.000000}%
\pgfsetdash{}{0pt}%
\pgfpathmoveto{\pgfqpoint{3.961812in}{6.294627in}}%
\pgfpathlineto{\pgfqpoint{4.002684in}{6.294627in}}%
\pgfpathlineto{\pgfqpoint{4.002684in}{7.511471in}}%
\pgfpathlineto{\pgfqpoint{3.961812in}{7.511471in}}%
\pgfpathlineto{\pgfqpoint{3.961812in}{6.294627in}}%
\pgfpathclose%
\pgfusepath{fill}%
\end{pgfscope}%
\begin{pgfscope}%
\pgfpathrectangle{\pgfqpoint{0.787430in}{6.294627in}}{\pgfqpoint{7.062570in}{2.555373in}}%
\pgfusepath{clip}%
\pgfsetbuttcap%
\pgfsetmiterjoin%
\definecolor{currentfill}{rgb}{0.121569,0.466667,0.705882}%
\pgfsetfillcolor{currentfill}%
\pgfsetlinewidth{0.000000pt}%
\definecolor{currentstroke}{rgb}{0.000000,0.000000,0.000000}%
\pgfsetstrokecolor{currentstroke}%
\pgfsetstrokeopacity{0.000000}%
\pgfsetdash{}{0pt}%
\pgfpathmoveto{\pgfqpoint{4.002684in}{6.294627in}}%
\pgfpathlineto{\pgfqpoint{4.043555in}{6.294627in}}%
\pgfpathlineto{\pgfqpoint{4.043555in}{7.869367in}}%
\pgfpathlineto{\pgfqpoint{4.002684in}{7.869367in}}%
\pgfpathlineto{\pgfqpoint{4.002684in}{6.294627in}}%
\pgfpathclose%
\pgfusepath{fill}%
\end{pgfscope}%
\begin{pgfscope}%
\pgfpathrectangle{\pgfqpoint{0.787430in}{6.294627in}}{\pgfqpoint{7.062570in}{2.555373in}}%
\pgfusepath{clip}%
\pgfsetbuttcap%
\pgfsetmiterjoin%
\definecolor{currentfill}{rgb}{0.121569,0.466667,0.705882}%
\pgfsetfillcolor{currentfill}%
\pgfsetlinewidth{0.000000pt}%
\definecolor{currentstroke}{rgb}{0.000000,0.000000,0.000000}%
\pgfsetstrokecolor{currentstroke}%
\pgfsetstrokeopacity{0.000000}%
\pgfsetdash{}{0pt}%
\pgfpathmoveto{\pgfqpoint{4.043555in}{6.294627in}}%
\pgfpathlineto{\pgfqpoint{4.084426in}{6.294627in}}%
\pgfpathlineto{\pgfqpoint{4.084426in}{7.726208in}}%
\pgfpathlineto{\pgfqpoint{4.043555in}{7.726208in}}%
\pgfpathlineto{\pgfqpoint{4.043555in}{6.294627in}}%
\pgfpathclose%
\pgfusepath{fill}%
\end{pgfscope}%
\begin{pgfscope}%
\pgfpathrectangle{\pgfqpoint{0.787430in}{6.294627in}}{\pgfqpoint{7.062570in}{2.555373in}}%
\pgfusepath{clip}%
\pgfsetbuttcap%
\pgfsetmiterjoin%
\definecolor{currentfill}{rgb}{0.121569,0.466667,0.705882}%
\pgfsetfillcolor{currentfill}%
\pgfsetlinewidth{0.000000pt}%
\definecolor{currentstroke}{rgb}{0.000000,0.000000,0.000000}%
\pgfsetstrokecolor{currentstroke}%
\pgfsetstrokeopacity{0.000000}%
\pgfsetdash{}{0pt}%
\pgfpathmoveto{\pgfqpoint{4.084426in}{6.294627in}}%
\pgfpathlineto{\pgfqpoint{4.125297in}{6.294627in}}%
\pgfpathlineto{\pgfqpoint{4.125297in}{7.869367in}}%
\pgfpathlineto{\pgfqpoint{4.084426in}{7.869367in}}%
\pgfpathlineto{\pgfqpoint{4.084426in}{6.294627in}}%
\pgfpathclose%
\pgfusepath{fill}%
\end{pgfscope}%
\begin{pgfscope}%
\pgfpathrectangle{\pgfqpoint{0.787430in}{6.294627in}}{\pgfqpoint{7.062570in}{2.555373in}}%
\pgfusepath{clip}%
\pgfsetbuttcap%
\pgfsetmiterjoin%
\definecolor{currentfill}{rgb}{0.121569,0.466667,0.705882}%
\pgfsetfillcolor{currentfill}%
\pgfsetlinewidth{0.000000pt}%
\definecolor{currentstroke}{rgb}{0.000000,0.000000,0.000000}%
\pgfsetstrokecolor{currentstroke}%
\pgfsetstrokeopacity{0.000000}%
\pgfsetdash{}{0pt}%
\pgfpathmoveto{\pgfqpoint{4.125297in}{6.294627in}}%
\pgfpathlineto{\pgfqpoint{4.166168in}{6.294627in}}%
\pgfpathlineto{\pgfqpoint{4.166168in}{7.869367in}}%
\pgfpathlineto{\pgfqpoint{4.125297in}{7.869367in}}%
\pgfpathlineto{\pgfqpoint{4.125297in}{6.294627in}}%
\pgfpathclose%
\pgfusepath{fill}%
\end{pgfscope}%
\begin{pgfscope}%
\pgfpathrectangle{\pgfqpoint{0.787430in}{6.294627in}}{\pgfqpoint{7.062570in}{2.555373in}}%
\pgfusepath{clip}%
\pgfsetbuttcap%
\pgfsetmiterjoin%
\definecolor{currentfill}{rgb}{0.121569,0.466667,0.705882}%
\pgfsetfillcolor{currentfill}%
\pgfsetlinewidth{0.000000pt}%
\definecolor{currentstroke}{rgb}{0.000000,0.000000,0.000000}%
\pgfsetstrokecolor{currentstroke}%
\pgfsetstrokeopacity{0.000000}%
\pgfsetdash{}{0pt}%
\pgfpathmoveto{\pgfqpoint{4.166168in}{6.294627in}}%
\pgfpathlineto{\pgfqpoint{4.207040in}{6.294627in}}%
\pgfpathlineto{\pgfqpoint{4.207040in}{7.439892in}}%
\pgfpathlineto{\pgfqpoint{4.166168in}{7.439892in}}%
\pgfpathlineto{\pgfqpoint{4.166168in}{6.294627in}}%
\pgfpathclose%
\pgfusepath{fill}%
\end{pgfscope}%
\begin{pgfscope}%
\pgfpathrectangle{\pgfqpoint{0.787430in}{6.294627in}}{\pgfqpoint{7.062570in}{2.555373in}}%
\pgfusepath{clip}%
\pgfsetbuttcap%
\pgfsetmiterjoin%
\definecolor{currentfill}{rgb}{0.121569,0.466667,0.705882}%
\pgfsetfillcolor{currentfill}%
\pgfsetlinewidth{0.000000pt}%
\definecolor{currentstroke}{rgb}{0.000000,0.000000,0.000000}%
\pgfsetstrokecolor{currentstroke}%
\pgfsetstrokeopacity{0.000000}%
\pgfsetdash{}{0pt}%
\pgfpathmoveto{\pgfqpoint{4.207040in}{6.294627in}}%
\pgfpathlineto{\pgfqpoint{4.247911in}{6.294627in}}%
\pgfpathlineto{\pgfqpoint{4.247911in}{8.728316in}}%
\pgfpathlineto{\pgfqpoint{4.207040in}{8.728316in}}%
\pgfpathlineto{\pgfqpoint{4.207040in}{6.294627in}}%
\pgfpathclose%
\pgfusepath{fill}%
\end{pgfscope}%
\begin{pgfscope}%
\pgfpathrectangle{\pgfqpoint{0.787430in}{6.294627in}}{\pgfqpoint{7.062570in}{2.555373in}}%
\pgfusepath{clip}%
\pgfsetbuttcap%
\pgfsetmiterjoin%
\definecolor{currentfill}{rgb}{0.121569,0.466667,0.705882}%
\pgfsetfillcolor{currentfill}%
\pgfsetlinewidth{0.000000pt}%
\definecolor{currentstroke}{rgb}{0.000000,0.000000,0.000000}%
\pgfsetstrokecolor{currentstroke}%
\pgfsetstrokeopacity{0.000000}%
\pgfsetdash{}{0pt}%
\pgfpathmoveto{\pgfqpoint{4.247911in}{6.294627in}}%
\pgfpathlineto{\pgfqpoint{4.288782in}{6.294627in}}%
\pgfpathlineto{\pgfqpoint{4.288782in}{7.940946in}}%
\pgfpathlineto{\pgfqpoint{4.247911in}{7.940946in}}%
\pgfpathlineto{\pgfqpoint{4.247911in}{6.294627in}}%
\pgfpathclose%
\pgfusepath{fill}%
\end{pgfscope}%
\begin{pgfscope}%
\pgfpathrectangle{\pgfqpoint{0.787430in}{6.294627in}}{\pgfqpoint{7.062570in}{2.555373in}}%
\pgfusepath{clip}%
\pgfsetbuttcap%
\pgfsetmiterjoin%
\definecolor{currentfill}{rgb}{0.121569,0.466667,0.705882}%
\pgfsetfillcolor{currentfill}%
\pgfsetlinewidth{0.000000pt}%
\definecolor{currentstroke}{rgb}{0.000000,0.000000,0.000000}%
\pgfsetstrokecolor{currentstroke}%
\pgfsetstrokeopacity{0.000000}%
\pgfsetdash{}{0pt}%
\pgfpathmoveto{\pgfqpoint{4.288782in}{6.294627in}}%
\pgfpathlineto{\pgfqpoint{4.329653in}{6.294627in}}%
\pgfpathlineto{\pgfqpoint{4.329653in}{8.227262in}}%
\pgfpathlineto{\pgfqpoint{4.288782in}{8.227262in}}%
\pgfpathlineto{\pgfqpoint{4.288782in}{6.294627in}}%
\pgfpathclose%
\pgfusepath{fill}%
\end{pgfscope}%
\begin{pgfscope}%
\pgfpathrectangle{\pgfqpoint{0.787430in}{6.294627in}}{\pgfqpoint{7.062570in}{2.555373in}}%
\pgfusepath{clip}%
\pgfsetbuttcap%
\pgfsetmiterjoin%
\definecolor{currentfill}{rgb}{0.121569,0.466667,0.705882}%
\pgfsetfillcolor{currentfill}%
\pgfsetlinewidth{0.000000pt}%
\definecolor{currentstroke}{rgb}{0.000000,0.000000,0.000000}%
\pgfsetstrokecolor{currentstroke}%
\pgfsetstrokeopacity{0.000000}%
\pgfsetdash{}{0pt}%
\pgfpathmoveto{\pgfqpoint{4.329653in}{6.294627in}}%
\pgfpathlineto{\pgfqpoint{4.370524in}{6.294627in}}%
\pgfpathlineto{\pgfqpoint{4.370524in}{8.227262in}}%
\pgfpathlineto{\pgfqpoint{4.329653in}{8.227262in}}%
\pgfpathlineto{\pgfqpoint{4.329653in}{6.294627in}}%
\pgfpathclose%
\pgfusepath{fill}%
\end{pgfscope}%
\begin{pgfscope}%
\pgfpathrectangle{\pgfqpoint{0.787430in}{6.294627in}}{\pgfqpoint{7.062570in}{2.555373in}}%
\pgfusepath{clip}%
\pgfsetbuttcap%
\pgfsetmiterjoin%
\definecolor{currentfill}{rgb}{0.121569,0.466667,0.705882}%
\pgfsetfillcolor{currentfill}%
\pgfsetlinewidth{0.000000pt}%
\definecolor{currentstroke}{rgb}{0.000000,0.000000,0.000000}%
\pgfsetstrokecolor{currentstroke}%
\pgfsetstrokeopacity{0.000000}%
\pgfsetdash{}{0pt}%
\pgfpathmoveto{\pgfqpoint{4.370524in}{6.294627in}}%
\pgfpathlineto{\pgfqpoint{4.411396in}{6.294627in}}%
\pgfpathlineto{\pgfqpoint{4.411396in}{7.869367in}}%
\pgfpathlineto{\pgfqpoint{4.370524in}{7.869367in}}%
\pgfpathlineto{\pgfqpoint{4.370524in}{6.294627in}}%
\pgfpathclose%
\pgfusepath{fill}%
\end{pgfscope}%
\begin{pgfscope}%
\pgfpathrectangle{\pgfqpoint{0.787430in}{6.294627in}}{\pgfqpoint{7.062570in}{2.555373in}}%
\pgfusepath{clip}%
\pgfsetbuttcap%
\pgfsetmiterjoin%
\definecolor{currentfill}{rgb}{0.121569,0.466667,0.705882}%
\pgfsetfillcolor{currentfill}%
\pgfsetlinewidth{0.000000pt}%
\definecolor{currentstroke}{rgb}{0.000000,0.000000,0.000000}%
\pgfsetstrokecolor{currentstroke}%
\pgfsetstrokeopacity{0.000000}%
\pgfsetdash{}{0pt}%
\pgfpathmoveto{\pgfqpoint{4.411396in}{6.294627in}}%
\pgfpathlineto{\pgfqpoint{4.452267in}{6.294627in}}%
\pgfpathlineto{\pgfqpoint{4.452267in}{7.869367in}}%
\pgfpathlineto{\pgfqpoint{4.411396in}{7.869367in}}%
\pgfpathlineto{\pgfqpoint{4.411396in}{6.294627in}}%
\pgfpathclose%
\pgfusepath{fill}%
\end{pgfscope}%
\begin{pgfscope}%
\pgfpathrectangle{\pgfqpoint{0.787430in}{6.294627in}}{\pgfqpoint{7.062570in}{2.555373in}}%
\pgfusepath{clip}%
\pgfsetbuttcap%
\pgfsetmiterjoin%
\definecolor{currentfill}{rgb}{0.121569,0.466667,0.705882}%
\pgfsetfillcolor{currentfill}%
\pgfsetlinewidth{0.000000pt}%
\definecolor{currentstroke}{rgb}{0.000000,0.000000,0.000000}%
\pgfsetstrokecolor{currentstroke}%
\pgfsetstrokeopacity{0.000000}%
\pgfsetdash{}{0pt}%
\pgfpathmoveto{\pgfqpoint{4.452267in}{6.294627in}}%
\pgfpathlineto{\pgfqpoint{4.493138in}{6.294627in}}%
\pgfpathlineto{\pgfqpoint{4.493138in}{8.513578in}}%
\pgfpathlineto{\pgfqpoint{4.452267in}{8.513578in}}%
\pgfpathlineto{\pgfqpoint{4.452267in}{6.294627in}}%
\pgfpathclose%
\pgfusepath{fill}%
\end{pgfscope}%
\begin{pgfscope}%
\pgfpathrectangle{\pgfqpoint{0.787430in}{6.294627in}}{\pgfqpoint{7.062570in}{2.555373in}}%
\pgfusepath{clip}%
\pgfsetbuttcap%
\pgfsetmiterjoin%
\definecolor{currentfill}{rgb}{0.121569,0.466667,0.705882}%
\pgfsetfillcolor{currentfill}%
\pgfsetlinewidth{0.000000pt}%
\definecolor{currentstroke}{rgb}{0.000000,0.000000,0.000000}%
\pgfsetstrokecolor{currentstroke}%
\pgfsetstrokeopacity{0.000000}%
\pgfsetdash{}{0pt}%
\pgfpathmoveto{\pgfqpoint{4.493138in}{6.294627in}}%
\pgfpathlineto{\pgfqpoint{4.534009in}{6.294627in}}%
\pgfpathlineto{\pgfqpoint{4.534009in}{8.513578in}}%
\pgfpathlineto{\pgfqpoint{4.493138in}{8.513578in}}%
\pgfpathlineto{\pgfqpoint{4.493138in}{6.294627in}}%
\pgfpathclose%
\pgfusepath{fill}%
\end{pgfscope}%
\begin{pgfscope}%
\pgfpathrectangle{\pgfqpoint{0.787430in}{6.294627in}}{\pgfqpoint{7.062570in}{2.555373in}}%
\pgfusepath{clip}%
\pgfsetbuttcap%
\pgfsetmiterjoin%
\definecolor{currentfill}{rgb}{0.121569,0.466667,0.705882}%
\pgfsetfillcolor{currentfill}%
\pgfsetlinewidth{0.000000pt}%
\definecolor{currentstroke}{rgb}{0.000000,0.000000,0.000000}%
\pgfsetstrokecolor{currentstroke}%
\pgfsetstrokeopacity{0.000000}%
\pgfsetdash{}{0pt}%
\pgfpathmoveto{\pgfqpoint{4.534009in}{6.294627in}}%
\pgfpathlineto{\pgfqpoint{4.574880in}{6.294627in}}%
\pgfpathlineto{\pgfqpoint{4.574880in}{8.513578in}}%
\pgfpathlineto{\pgfqpoint{4.534009in}{8.513578in}}%
\pgfpathlineto{\pgfqpoint{4.534009in}{6.294627in}}%
\pgfpathclose%
\pgfusepath{fill}%
\end{pgfscope}%
\begin{pgfscope}%
\pgfpathrectangle{\pgfqpoint{0.787430in}{6.294627in}}{\pgfqpoint{7.062570in}{2.555373in}}%
\pgfusepath{clip}%
\pgfsetbuttcap%
\pgfsetmiterjoin%
\definecolor{currentfill}{rgb}{0.121569,0.466667,0.705882}%
\pgfsetfillcolor{currentfill}%
\pgfsetlinewidth{0.000000pt}%
\definecolor{currentstroke}{rgb}{0.000000,0.000000,0.000000}%
\pgfsetstrokecolor{currentstroke}%
\pgfsetstrokeopacity{0.000000}%
\pgfsetdash{}{0pt}%
\pgfpathmoveto{\pgfqpoint{4.574880in}{6.294627in}}%
\pgfpathlineto{\pgfqpoint{4.615752in}{6.294627in}}%
\pgfpathlineto{\pgfqpoint{4.615752in}{8.012525in}}%
\pgfpathlineto{\pgfqpoint{4.574880in}{8.012525in}}%
\pgfpathlineto{\pgfqpoint{4.574880in}{6.294627in}}%
\pgfpathclose%
\pgfusepath{fill}%
\end{pgfscope}%
\begin{pgfscope}%
\pgfpathrectangle{\pgfqpoint{0.787430in}{6.294627in}}{\pgfqpoint{7.062570in}{2.555373in}}%
\pgfusepath{clip}%
\pgfsetbuttcap%
\pgfsetmiterjoin%
\definecolor{currentfill}{rgb}{0.121569,0.466667,0.705882}%
\pgfsetfillcolor{currentfill}%
\pgfsetlinewidth{0.000000pt}%
\definecolor{currentstroke}{rgb}{0.000000,0.000000,0.000000}%
\pgfsetstrokecolor{currentstroke}%
\pgfsetstrokeopacity{0.000000}%
\pgfsetdash{}{0pt}%
\pgfpathmoveto{\pgfqpoint{4.615752in}{6.294627in}}%
\pgfpathlineto{\pgfqpoint{4.656623in}{6.294627in}}%
\pgfpathlineto{\pgfqpoint{4.656623in}{7.869367in}}%
\pgfpathlineto{\pgfqpoint{4.615752in}{7.869367in}}%
\pgfpathlineto{\pgfqpoint{4.615752in}{6.294627in}}%
\pgfpathclose%
\pgfusepath{fill}%
\end{pgfscope}%
\begin{pgfscope}%
\pgfpathrectangle{\pgfqpoint{0.787430in}{6.294627in}}{\pgfqpoint{7.062570in}{2.555373in}}%
\pgfusepath{clip}%
\pgfsetbuttcap%
\pgfsetmiterjoin%
\definecolor{currentfill}{rgb}{0.121569,0.466667,0.705882}%
\pgfsetfillcolor{currentfill}%
\pgfsetlinewidth{0.000000pt}%
\definecolor{currentstroke}{rgb}{0.000000,0.000000,0.000000}%
\pgfsetstrokecolor{currentstroke}%
\pgfsetstrokeopacity{0.000000}%
\pgfsetdash{}{0pt}%
\pgfpathmoveto{\pgfqpoint{4.656623in}{6.294627in}}%
\pgfpathlineto{\pgfqpoint{4.697494in}{6.294627in}}%
\pgfpathlineto{\pgfqpoint{4.697494in}{7.654629in}}%
\pgfpathlineto{\pgfqpoint{4.656623in}{7.654629in}}%
\pgfpathlineto{\pgfqpoint{4.656623in}{6.294627in}}%
\pgfpathclose%
\pgfusepath{fill}%
\end{pgfscope}%
\begin{pgfscope}%
\pgfpathrectangle{\pgfqpoint{0.787430in}{6.294627in}}{\pgfqpoint{7.062570in}{2.555373in}}%
\pgfusepath{clip}%
\pgfsetbuttcap%
\pgfsetmiterjoin%
\definecolor{currentfill}{rgb}{0.121569,0.466667,0.705882}%
\pgfsetfillcolor{currentfill}%
\pgfsetlinewidth{0.000000pt}%
\definecolor{currentstroke}{rgb}{0.000000,0.000000,0.000000}%
\pgfsetstrokecolor{currentstroke}%
\pgfsetstrokeopacity{0.000000}%
\pgfsetdash{}{0pt}%
\pgfpathmoveto{\pgfqpoint{4.697494in}{6.294627in}}%
\pgfpathlineto{\pgfqpoint{4.738365in}{6.294627in}}%
\pgfpathlineto{\pgfqpoint{4.738365in}{7.869367in}}%
\pgfpathlineto{\pgfqpoint{4.697494in}{7.869367in}}%
\pgfpathlineto{\pgfqpoint{4.697494in}{6.294627in}}%
\pgfpathclose%
\pgfusepath{fill}%
\end{pgfscope}%
\begin{pgfscope}%
\pgfpathrectangle{\pgfqpoint{0.787430in}{6.294627in}}{\pgfqpoint{7.062570in}{2.555373in}}%
\pgfusepath{clip}%
\pgfsetbuttcap%
\pgfsetmiterjoin%
\definecolor{currentfill}{rgb}{0.121569,0.466667,0.705882}%
\pgfsetfillcolor{currentfill}%
\pgfsetlinewidth{0.000000pt}%
\definecolor{currentstroke}{rgb}{0.000000,0.000000,0.000000}%
\pgfsetstrokecolor{currentstroke}%
\pgfsetstrokeopacity{0.000000}%
\pgfsetdash{}{0pt}%
\pgfpathmoveto{\pgfqpoint{4.738365in}{6.294627in}}%
\pgfpathlineto{\pgfqpoint{4.779236in}{6.294627in}}%
\pgfpathlineto{\pgfqpoint{4.779236in}{8.012525in}}%
\pgfpathlineto{\pgfqpoint{4.738365in}{8.012525in}}%
\pgfpathlineto{\pgfqpoint{4.738365in}{6.294627in}}%
\pgfpathclose%
\pgfusepath{fill}%
\end{pgfscope}%
\begin{pgfscope}%
\pgfpathrectangle{\pgfqpoint{0.787430in}{6.294627in}}{\pgfqpoint{7.062570in}{2.555373in}}%
\pgfusepath{clip}%
\pgfsetbuttcap%
\pgfsetmiterjoin%
\definecolor{currentfill}{rgb}{0.121569,0.466667,0.705882}%
\pgfsetfillcolor{currentfill}%
\pgfsetlinewidth{0.000000pt}%
\definecolor{currentstroke}{rgb}{0.000000,0.000000,0.000000}%
\pgfsetstrokecolor{currentstroke}%
\pgfsetstrokeopacity{0.000000}%
\pgfsetdash{}{0pt}%
\pgfpathmoveto{\pgfqpoint{4.779236in}{6.294627in}}%
\pgfpathlineto{\pgfqpoint{4.820108in}{6.294627in}}%
\pgfpathlineto{\pgfqpoint{4.820108in}{7.368313in}}%
\pgfpathlineto{\pgfqpoint{4.779236in}{7.368313in}}%
\pgfpathlineto{\pgfqpoint{4.779236in}{6.294627in}}%
\pgfpathclose%
\pgfusepath{fill}%
\end{pgfscope}%
\begin{pgfscope}%
\pgfpathrectangle{\pgfqpoint{0.787430in}{6.294627in}}{\pgfqpoint{7.062570in}{2.555373in}}%
\pgfusepath{clip}%
\pgfsetbuttcap%
\pgfsetmiterjoin%
\definecolor{currentfill}{rgb}{0.121569,0.466667,0.705882}%
\pgfsetfillcolor{currentfill}%
\pgfsetlinewidth{0.000000pt}%
\definecolor{currentstroke}{rgb}{0.000000,0.000000,0.000000}%
\pgfsetstrokecolor{currentstroke}%
\pgfsetstrokeopacity{0.000000}%
\pgfsetdash{}{0pt}%
\pgfpathmoveto{\pgfqpoint{4.820108in}{6.294627in}}%
\pgfpathlineto{\pgfqpoint{4.860979in}{6.294627in}}%
\pgfpathlineto{\pgfqpoint{4.860979in}{7.296734in}}%
\pgfpathlineto{\pgfqpoint{4.820108in}{7.296734in}}%
\pgfpathlineto{\pgfqpoint{4.820108in}{6.294627in}}%
\pgfpathclose%
\pgfusepath{fill}%
\end{pgfscope}%
\begin{pgfscope}%
\pgfpathrectangle{\pgfqpoint{0.787430in}{6.294627in}}{\pgfqpoint{7.062570in}{2.555373in}}%
\pgfusepath{clip}%
\pgfsetbuttcap%
\pgfsetmiterjoin%
\definecolor{currentfill}{rgb}{0.121569,0.466667,0.705882}%
\pgfsetfillcolor{currentfill}%
\pgfsetlinewidth{0.000000pt}%
\definecolor{currentstroke}{rgb}{0.000000,0.000000,0.000000}%
\pgfsetstrokecolor{currentstroke}%
\pgfsetstrokeopacity{0.000000}%
\pgfsetdash{}{0pt}%
\pgfpathmoveto{\pgfqpoint{4.860979in}{6.294627in}}%
\pgfpathlineto{\pgfqpoint{4.901850in}{6.294627in}}%
\pgfpathlineto{\pgfqpoint{4.901850in}{8.012525in}}%
\pgfpathlineto{\pgfqpoint{4.860979in}{8.012525in}}%
\pgfpathlineto{\pgfqpoint{4.860979in}{6.294627in}}%
\pgfpathclose%
\pgfusepath{fill}%
\end{pgfscope}%
\begin{pgfscope}%
\pgfpathrectangle{\pgfqpoint{0.787430in}{6.294627in}}{\pgfqpoint{7.062570in}{2.555373in}}%
\pgfusepath{clip}%
\pgfsetbuttcap%
\pgfsetmiterjoin%
\definecolor{currentfill}{rgb}{0.121569,0.466667,0.705882}%
\pgfsetfillcolor{currentfill}%
\pgfsetlinewidth{0.000000pt}%
\definecolor{currentstroke}{rgb}{0.000000,0.000000,0.000000}%
\pgfsetstrokecolor{currentstroke}%
\pgfsetstrokeopacity{0.000000}%
\pgfsetdash{}{0pt}%
\pgfpathmoveto{\pgfqpoint{4.901850in}{6.294627in}}%
\pgfpathlineto{\pgfqpoint{4.942721in}{6.294627in}}%
\pgfpathlineto{\pgfqpoint{4.942721in}{7.940946in}}%
\pgfpathlineto{\pgfqpoint{4.901850in}{7.940946in}}%
\pgfpathlineto{\pgfqpoint{4.901850in}{6.294627in}}%
\pgfpathclose%
\pgfusepath{fill}%
\end{pgfscope}%
\begin{pgfscope}%
\pgfpathrectangle{\pgfqpoint{0.787430in}{6.294627in}}{\pgfqpoint{7.062570in}{2.555373in}}%
\pgfusepath{clip}%
\pgfsetbuttcap%
\pgfsetmiterjoin%
\definecolor{currentfill}{rgb}{0.121569,0.466667,0.705882}%
\pgfsetfillcolor{currentfill}%
\pgfsetlinewidth{0.000000pt}%
\definecolor{currentstroke}{rgb}{0.000000,0.000000,0.000000}%
\pgfsetstrokecolor{currentstroke}%
\pgfsetstrokeopacity{0.000000}%
\pgfsetdash{}{0pt}%
\pgfpathmoveto{\pgfqpoint{4.942721in}{6.294627in}}%
\pgfpathlineto{\pgfqpoint{4.983592in}{6.294627in}}%
\pgfpathlineto{\pgfqpoint{4.983592in}{7.511471in}}%
\pgfpathlineto{\pgfqpoint{4.942721in}{7.511471in}}%
\pgfpathlineto{\pgfqpoint{4.942721in}{6.294627in}}%
\pgfpathclose%
\pgfusepath{fill}%
\end{pgfscope}%
\begin{pgfscope}%
\pgfpathrectangle{\pgfqpoint{0.787430in}{6.294627in}}{\pgfqpoint{7.062570in}{2.555373in}}%
\pgfusepath{clip}%
\pgfsetbuttcap%
\pgfsetmiterjoin%
\definecolor{currentfill}{rgb}{0.121569,0.466667,0.705882}%
\pgfsetfillcolor{currentfill}%
\pgfsetlinewidth{0.000000pt}%
\definecolor{currentstroke}{rgb}{0.000000,0.000000,0.000000}%
\pgfsetstrokecolor{currentstroke}%
\pgfsetstrokeopacity{0.000000}%
\pgfsetdash{}{0pt}%
\pgfpathmoveto{\pgfqpoint{4.983592in}{6.294627in}}%
\pgfpathlineto{\pgfqpoint{5.024464in}{6.294627in}}%
\pgfpathlineto{\pgfqpoint{5.024464in}{7.225155in}}%
\pgfpathlineto{\pgfqpoint{4.983592in}{7.225155in}}%
\pgfpathlineto{\pgfqpoint{4.983592in}{6.294627in}}%
\pgfpathclose%
\pgfusepath{fill}%
\end{pgfscope}%
\begin{pgfscope}%
\pgfpathrectangle{\pgfqpoint{0.787430in}{6.294627in}}{\pgfqpoint{7.062570in}{2.555373in}}%
\pgfusepath{clip}%
\pgfsetbuttcap%
\pgfsetmiterjoin%
\definecolor{currentfill}{rgb}{0.121569,0.466667,0.705882}%
\pgfsetfillcolor{currentfill}%
\pgfsetlinewidth{0.000000pt}%
\definecolor{currentstroke}{rgb}{0.000000,0.000000,0.000000}%
\pgfsetstrokecolor{currentstroke}%
\pgfsetstrokeopacity{0.000000}%
\pgfsetdash{}{0pt}%
\pgfpathmoveto{\pgfqpoint{5.024464in}{6.294627in}}%
\pgfpathlineto{\pgfqpoint{5.065335in}{6.294627in}}%
\pgfpathlineto{\pgfqpoint{5.065335in}{7.296734in}}%
\pgfpathlineto{\pgfqpoint{5.024464in}{7.296734in}}%
\pgfpathlineto{\pgfqpoint{5.024464in}{6.294627in}}%
\pgfpathclose%
\pgfusepath{fill}%
\end{pgfscope}%
\begin{pgfscope}%
\pgfpathrectangle{\pgfqpoint{0.787430in}{6.294627in}}{\pgfqpoint{7.062570in}{2.555373in}}%
\pgfusepath{clip}%
\pgfsetbuttcap%
\pgfsetmiterjoin%
\definecolor{currentfill}{rgb}{0.121569,0.466667,0.705882}%
\pgfsetfillcolor{currentfill}%
\pgfsetlinewidth{0.000000pt}%
\definecolor{currentstroke}{rgb}{0.000000,0.000000,0.000000}%
\pgfsetstrokecolor{currentstroke}%
\pgfsetstrokeopacity{0.000000}%
\pgfsetdash{}{0pt}%
\pgfpathmoveto{\pgfqpoint{5.065335in}{6.294627in}}%
\pgfpathlineto{\pgfqpoint{5.106206in}{6.294627in}}%
\pgfpathlineto{\pgfqpoint{5.106206in}{7.511471in}}%
\pgfpathlineto{\pgfqpoint{5.065335in}{7.511471in}}%
\pgfpathlineto{\pgfqpoint{5.065335in}{6.294627in}}%
\pgfpathclose%
\pgfusepath{fill}%
\end{pgfscope}%
\begin{pgfscope}%
\pgfpathrectangle{\pgfqpoint{0.787430in}{6.294627in}}{\pgfqpoint{7.062570in}{2.555373in}}%
\pgfusepath{clip}%
\pgfsetbuttcap%
\pgfsetmiterjoin%
\definecolor{currentfill}{rgb}{0.121569,0.466667,0.705882}%
\pgfsetfillcolor{currentfill}%
\pgfsetlinewidth{0.000000pt}%
\definecolor{currentstroke}{rgb}{0.000000,0.000000,0.000000}%
\pgfsetstrokecolor{currentstroke}%
\pgfsetstrokeopacity{0.000000}%
\pgfsetdash{}{0pt}%
\pgfpathmoveto{\pgfqpoint{5.106206in}{6.294627in}}%
\pgfpathlineto{\pgfqpoint{5.147077in}{6.294627in}}%
\pgfpathlineto{\pgfqpoint{5.147077in}{6.867259in}}%
\pgfpathlineto{\pgfqpoint{5.106206in}{6.867259in}}%
\pgfpathlineto{\pgfqpoint{5.106206in}{6.294627in}}%
\pgfpathclose%
\pgfusepath{fill}%
\end{pgfscope}%
\begin{pgfscope}%
\pgfpathrectangle{\pgfqpoint{0.787430in}{6.294627in}}{\pgfqpoint{7.062570in}{2.555373in}}%
\pgfusepath{clip}%
\pgfsetbuttcap%
\pgfsetmiterjoin%
\definecolor{currentfill}{rgb}{0.121569,0.466667,0.705882}%
\pgfsetfillcolor{currentfill}%
\pgfsetlinewidth{0.000000pt}%
\definecolor{currentstroke}{rgb}{0.000000,0.000000,0.000000}%
\pgfsetstrokecolor{currentstroke}%
\pgfsetstrokeopacity{0.000000}%
\pgfsetdash{}{0pt}%
\pgfpathmoveto{\pgfqpoint{5.147077in}{6.294627in}}%
\pgfpathlineto{\pgfqpoint{5.187948in}{6.294627in}}%
\pgfpathlineto{\pgfqpoint{5.187948in}{6.938838in}}%
\pgfpathlineto{\pgfqpoint{5.147077in}{6.938838in}}%
\pgfpathlineto{\pgfqpoint{5.147077in}{6.294627in}}%
\pgfpathclose%
\pgfusepath{fill}%
\end{pgfscope}%
\begin{pgfscope}%
\pgfpathrectangle{\pgfqpoint{0.787430in}{6.294627in}}{\pgfqpoint{7.062570in}{2.555373in}}%
\pgfusepath{clip}%
\pgfsetbuttcap%
\pgfsetmiterjoin%
\definecolor{currentfill}{rgb}{0.121569,0.466667,0.705882}%
\pgfsetfillcolor{currentfill}%
\pgfsetlinewidth{0.000000pt}%
\definecolor{currentstroke}{rgb}{0.000000,0.000000,0.000000}%
\pgfsetstrokecolor{currentstroke}%
\pgfsetstrokeopacity{0.000000}%
\pgfsetdash{}{0pt}%
\pgfpathmoveto{\pgfqpoint{5.187948in}{6.294627in}}%
\pgfpathlineto{\pgfqpoint{5.228820in}{6.294627in}}%
\pgfpathlineto{\pgfqpoint{5.228820in}{7.153576in}}%
\pgfpathlineto{\pgfqpoint{5.187948in}{7.153576in}}%
\pgfpathlineto{\pgfqpoint{5.187948in}{6.294627in}}%
\pgfpathclose%
\pgfusepath{fill}%
\end{pgfscope}%
\begin{pgfscope}%
\pgfpathrectangle{\pgfqpoint{0.787430in}{6.294627in}}{\pgfqpoint{7.062570in}{2.555373in}}%
\pgfusepath{clip}%
\pgfsetbuttcap%
\pgfsetmiterjoin%
\definecolor{currentfill}{rgb}{0.121569,0.466667,0.705882}%
\pgfsetfillcolor{currentfill}%
\pgfsetlinewidth{0.000000pt}%
\definecolor{currentstroke}{rgb}{0.000000,0.000000,0.000000}%
\pgfsetstrokecolor{currentstroke}%
\pgfsetstrokeopacity{0.000000}%
\pgfsetdash{}{0pt}%
\pgfpathmoveto{\pgfqpoint{5.228820in}{6.294627in}}%
\pgfpathlineto{\pgfqpoint{5.269691in}{6.294627in}}%
\pgfpathlineto{\pgfqpoint{5.269691in}{7.153576in}}%
\pgfpathlineto{\pgfqpoint{5.228820in}{7.153576in}}%
\pgfpathlineto{\pgfqpoint{5.228820in}{6.294627in}}%
\pgfpathclose%
\pgfusepath{fill}%
\end{pgfscope}%
\begin{pgfscope}%
\pgfpathrectangle{\pgfqpoint{0.787430in}{6.294627in}}{\pgfqpoint{7.062570in}{2.555373in}}%
\pgfusepath{clip}%
\pgfsetbuttcap%
\pgfsetmiterjoin%
\definecolor{currentfill}{rgb}{0.121569,0.466667,0.705882}%
\pgfsetfillcolor{currentfill}%
\pgfsetlinewidth{0.000000pt}%
\definecolor{currentstroke}{rgb}{0.000000,0.000000,0.000000}%
\pgfsetstrokecolor{currentstroke}%
\pgfsetstrokeopacity{0.000000}%
\pgfsetdash{}{0pt}%
\pgfpathmoveto{\pgfqpoint{5.269691in}{6.294627in}}%
\pgfpathlineto{\pgfqpoint{5.310562in}{6.294627in}}%
\pgfpathlineto{\pgfqpoint{5.310562in}{6.867259in}}%
\pgfpathlineto{\pgfqpoint{5.269691in}{6.867259in}}%
\pgfpathlineto{\pgfqpoint{5.269691in}{6.294627in}}%
\pgfpathclose%
\pgfusepath{fill}%
\end{pgfscope}%
\begin{pgfscope}%
\pgfpathrectangle{\pgfqpoint{0.787430in}{6.294627in}}{\pgfqpoint{7.062570in}{2.555373in}}%
\pgfusepath{clip}%
\pgfsetbuttcap%
\pgfsetmiterjoin%
\definecolor{currentfill}{rgb}{0.121569,0.466667,0.705882}%
\pgfsetfillcolor{currentfill}%
\pgfsetlinewidth{0.000000pt}%
\definecolor{currentstroke}{rgb}{0.000000,0.000000,0.000000}%
\pgfsetstrokecolor{currentstroke}%
\pgfsetstrokeopacity{0.000000}%
\pgfsetdash{}{0pt}%
\pgfpathmoveto{\pgfqpoint{5.310562in}{6.294627in}}%
\pgfpathlineto{\pgfqpoint{5.351433in}{6.294627in}}%
\pgfpathlineto{\pgfqpoint{5.351433in}{7.010418in}}%
\pgfpathlineto{\pgfqpoint{5.310562in}{7.010418in}}%
\pgfpathlineto{\pgfqpoint{5.310562in}{6.294627in}}%
\pgfpathclose%
\pgfusepath{fill}%
\end{pgfscope}%
\begin{pgfscope}%
\pgfpathrectangle{\pgfqpoint{0.787430in}{6.294627in}}{\pgfqpoint{7.062570in}{2.555373in}}%
\pgfusepath{clip}%
\pgfsetbuttcap%
\pgfsetmiterjoin%
\definecolor{currentfill}{rgb}{0.121569,0.466667,0.705882}%
\pgfsetfillcolor{currentfill}%
\pgfsetlinewidth{0.000000pt}%
\definecolor{currentstroke}{rgb}{0.000000,0.000000,0.000000}%
\pgfsetstrokecolor{currentstroke}%
\pgfsetstrokeopacity{0.000000}%
\pgfsetdash{}{0pt}%
\pgfpathmoveto{\pgfqpoint{5.351433in}{6.294627in}}%
\pgfpathlineto{\pgfqpoint{5.392304in}{6.294627in}}%
\pgfpathlineto{\pgfqpoint{5.392304in}{6.724101in}}%
\pgfpathlineto{\pgfqpoint{5.351433in}{6.724101in}}%
\pgfpathlineto{\pgfqpoint{5.351433in}{6.294627in}}%
\pgfpathclose%
\pgfusepath{fill}%
\end{pgfscope}%
\begin{pgfscope}%
\pgfpathrectangle{\pgfqpoint{0.787430in}{6.294627in}}{\pgfqpoint{7.062570in}{2.555373in}}%
\pgfusepath{clip}%
\pgfsetbuttcap%
\pgfsetmiterjoin%
\definecolor{currentfill}{rgb}{0.121569,0.466667,0.705882}%
\pgfsetfillcolor{currentfill}%
\pgfsetlinewidth{0.000000pt}%
\definecolor{currentstroke}{rgb}{0.000000,0.000000,0.000000}%
\pgfsetstrokecolor{currentstroke}%
\pgfsetstrokeopacity{0.000000}%
\pgfsetdash{}{0pt}%
\pgfpathmoveto{\pgfqpoint{5.392304in}{6.294627in}}%
\pgfpathlineto{\pgfqpoint{5.433176in}{6.294627in}}%
\pgfpathlineto{\pgfqpoint{5.433176in}{7.010418in}}%
\pgfpathlineto{\pgfqpoint{5.392304in}{7.010418in}}%
\pgfpathlineto{\pgfqpoint{5.392304in}{6.294627in}}%
\pgfpathclose%
\pgfusepath{fill}%
\end{pgfscope}%
\begin{pgfscope}%
\pgfpathrectangle{\pgfqpoint{0.787430in}{6.294627in}}{\pgfqpoint{7.062570in}{2.555373in}}%
\pgfusepath{clip}%
\pgfsetbuttcap%
\pgfsetmiterjoin%
\definecolor{currentfill}{rgb}{0.121569,0.466667,0.705882}%
\pgfsetfillcolor{currentfill}%
\pgfsetlinewidth{0.000000pt}%
\definecolor{currentstroke}{rgb}{0.000000,0.000000,0.000000}%
\pgfsetstrokecolor{currentstroke}%
\pgfsetstrokeopacity{0.000000}%
\pgfsetdash{}{0pt}%
\pgfpathmoveto{\pgfqpoint{5.433176in}{6.294627in}}%
\pgfpathlineto{\pgfqpoint{5.474047in}{6.294627in}}%
\pgfpathlineto{\pgfqpoint{5.474047in}{6.938838in}}%
\pgfpathlineto{\pgfqpoint{5.433176in}{6.938838in}}%
\pgfpathlineto{\pgfqpoint{5.433176in}{6.294627in}}%
\pgfpathclose%
\pgfusepath{fill}%
\end{pgfscope}%
\begin{pgfscope}%
\pgfpathrectangle{\pgfqpoint{0.787430in}{6.294627in}}{\pgfqpoint{7.062570in}{2.555373in}}%
\pgfusepath{clip}%
\pgfsetbuttcap%
\pgfsetmiterjoin%
\definecolor{currentfill}{rgb}{0.121569,0.466667,0.705882}%
\pgfsetfillcolor{currentfill}%
\pgfsetlinewidth{0.000000pt}%
\definecolor{currentstroke}{rgb}{0.000000,0.000000,0.000000}%
\pgfsetstrokecolor{currentstroke}%
\pgfsetstrokeopacity{0.000000}%
\pgfsetdash{}{0pt}%
\pgfpathmoveto{\pgfqpoint{5.474047in}{6.294627in}}%
\pgfpathlineto{\pgfqpoint{5.514918in}{6.294627in}}%
\pgfpathlineto{\pgfqpoint{5.514918in}{7.010418in}}%
\pgfpathlineto{\pgfqpoint{5.474047in}{7.010418in}}%
\pgfpathlineto{\pgfqpoint{5.474047in}{6.294627in}}%
\pgfpathclose%
\pgfusepath{fill}%
\end{pgfscope}%
\begin{pgfscope}%
\pgfpathrectangle{\pgfqpoint{0.787430in}{6.294627in}}{\pgfqpoint{7.062570in}{2.555373in}}%
\pgfusepath{clip}%
\pgfsetbuttcap%
\pgfsetmiterjoin%
\definecolor{currentfill}{rgb}{0.121569,0.466667,0.705882}%
\pgfsetfillcolor{currentfill}%
\pgfsetlinewidth{0.000000pt}%
\definecolor{currentstroke}{rgb}{0.000000,0.000000,0.000000}%
\pgfsetstrokecolor{currentstroke}%
\pgfsetstrokeopacity{0.000000}%
\pgfsetdash{}{0pt}%
\pgfpathmoveto{\pgfqpoint{5.514918in}{6.294627in}}%
\pgfpathlineto{\pgfqpoint{5.555789in}{6.294627in}}%
\pgfpathlineto{\pgfqpoint{5.555789in}{6.724101in}}%
\pgfpathlineto{\pgfqpoint{5.514918in}{6.724101in}}%
\pgfpathlineto{\pgfqpoint{5.514918in}{6.294627in}}%
\pgfpathclose%
\pgfusepath{fill}%
\end{pgfscope}%
\begin{pgfscope}%
\pgfpathrectangle{\pgfqpoint{0.787430in}{6.294627in}}{\pgfqpoint{7.062570in}{2.555373in}}%
\pgfusepath{clip}%
\pgfsetbuttcap%
\pgfsetmiterjoin%
\definecolor{currentfill}{rgb}{0.121569,0.466667,0.705882}%
\pgfsetfillcolor{currentfill}%
\pgfsetlinewidth{0.000000pt}%
\definecolor{currentstroke}{rgb}{0.000000,0.000000,0.000000}%
\pgfsetstrokecolor{currentstroke}%
\pgfsetstrokeopacity{0.000000}%
\pgfsetdash{}{0pt}%
\pgfpathmoveto{\pgfqpoint{5.555789in}{6.294627in}}%
\pgfpathlineto{\pgfqpoint{5.596660in}{6.294627in}}%
\pgfpathlineto{\pgfqpoint{5.596660in}{6.580943in}}%
\pgfpathlineto{\pgfqpoint{5.555789in}{6.580943in}}%
\pgfpathlineto{\pgfqpoint{5.555789in}{6.294627in}}%
\pgfpathclose%
\pgfusepath{fill}%
\end{pgfscope}%
\begin{pgfscope}%
\pgfpathrectangle{\pgfqpoint{0.787430in}{6.294627in}}{\pgfqpoint{7.062570in}{2.555373in}}%
\pgfusepath{clip}%
\pgfsetbuttcap%
\pgfsetmiterjoin%
\definecolor{currentfill}{rgb}{0.121569,0.466667,0.705882}%
\pgfsetfillcolor{currentfill}%
\pgfsetlinewidth{0.000000pt}%
\definecolor{currentstroke}{rgb}{0.000000,0.000000,0.000000}%
\pgfsetstrokecolor{currentstroke}%
\pgfsetstrokeopacity{0.000000}%
\pgfsetdash{}{0pt}%
\pgfpathmoveto{\pgfqpoint{5.596660in}{6.294627in}}%
\pgfpathlineto{\pgfqpoint{5.637532in}{6.294627in}}%
\pgfpathlineto{\pgfqpoint{5.637532in}{6.724101in}}%
\pgfpathlineto{\pgfqpoint{5.596660in}{6.724101in}}%
\pgfpathlineto{\pgfqpoint{5.596660in}{6.294627in}}%
\pgfpathclose%
\pgfusepath{fill}%
\end{pgfscope}%
\begin{pgfscope}%
\pgfpathrectangle{\pgfqpoint{0.787430in}{6.294627in}}{\pgfqpoint{7.062570in}{2.555373in}}%
\pgfusepath{clip}%
\pgfsetbuttcap%
\pgfsetmiterjoin%
\definecolor{currentfill}{rgb}{0.121569,0.466667,0.705882}%
\pgfsetfillcolor{currentfill}%
\pgfsetlinewidth{0.000000pt}%
\definecolor{currentstroke}{rgb}{0.000000,0.000000,0.000000}%
\pgfsetstrokecolor{currentstroke}%
\pgfsetstrokeopacity{0.000000}%
\pgfsetdash{}{0pt}%
\pgfpathmoveto{\pgfqpoint{5.637532in}{6.294627in}}%
\pgfpathlineto{\pgfqpoint{5.678403in}{6.294627in}}%
\pgfpathlineto{\pgfqpoint{5.678403in}{6.509364in}}%
\pgfpathlineto{\pgfqpoint{5.637532in}{6.509364in}}%
\pgfpathlineto{\pgfqpoint{5.637532in}{6.294627in}}%
\pgfpathclose%
\pgfusepath{fill}%
\end{pgfscope}%
\begin{pgfscope}%
\pgfpathrectangle{\pgfqpoint{0.787430in}{6.294627in}}{\pgfqpoint{7.062570in}{2.555373in}}%
\pgfusepath{clip}%
\pgfsetbuttcap%
\pgfsetmiterjoin%
\definecolor{currentfill}{rgb}{0.121569,0.466667,0.705882}%
\pgfsetfillcolor{currentfill}%
\pgfsetlinewidth{0.000000pt}%
\definecolor{currentstroke}{rgb}{0.000000,0.000000,0.000000}%
\pgfsetstrokecolor{currentstroke}%
\pgfsetstrokeopacity{0.000000}%
\pgfsetdash{}{0pt}%
\pgfpathmoveto{\pgfqpoint{5.678403in}{6.294627in}}%
\pgfpathlineto{\pgfqpoint{5.719274in}{6.294627in}}%
\pgfpathlineto{\pgfqpoint{5.719274in}{6.437785in}}%
\pgfpathlineto{\pgfqpoint{5.678403in}{6.437785in}}%
\pgfpathlineto{\pgfqpoint{5.678403in}{6.294627in}}%
\pgfpathclose%
\pgfusepath{fill}%
\end{pgfscope}%
\begin{pgfscope}%
\pgfpathrectangle{\pgfqpoint{0.787430in}{6.294627in}}{\pgfqpoint{7.062570in}{2.555373in}}%
\pgfusepath{clip}%
\pgfsetbuttcap%
\pgfsetmiterjoin%
\definecolor{currentfill}{rgb}{0.121569,0.466667,0.705882}%
\pgfsetfillcolor{currentfill}%
\pgfsetlinewidth{0.000000pt}%
\definecolor{currentstroke}{rgb}{0.000000,0.000000,0.000000}%
\pgfsetstrokecolor{currentstroke}%
\pgfsetstrokeopacity{0.000000}%
\pgfsetdash{}{0pt}%
\pgfpathmoveto{\pgfqpoint{5.719274in}{6.294627in}}%
\pgfpathlineto{\pgfqpoint{5.760145in}{6.294627in}}%
\pgfpathlineto{\pgfqpoint{5.760145in}{6.366206in}}%
\pgfpathlineto{\pgfqpoint{5.719274in}{6.366206in}}%
\pgfpathlineto{\pgfqpoint{5.719274in}{6.294627in}}%
\pgfpathclose%
\pgfusepath{fill}%
\end{pgfscope}%
\begin{pgfscope}%
\pgfpathrectangle{\pgfqpoint{0.787430in}{6.294627in}}{\pgfqpoint{7.062570in}{2.555373in}}%
\pgfusepath{clip}%
\pgfsetbuttcap%
\pgfsetmiterjoin%
\definecolor{currentfill}{rgb}{0.121569,0.466667,0.705882}%
\pgfsetfillcolor{currentfill}%
\pgfsetlinewidth{0.000000pt}%
\definecolor{currentstroke}{rgb}{0.000000,0.000000,0.000000}%
\pgfsetstrokecolor{currentstroke}%
\pgfsetstrokeopacity{0.000000}%
\pgfsetdash{}{0pt}%
\pgfpathmoveto{\pgfqpoint{5.760145in}{6.294627in}}%
\pgfpathlineto{\pgfqpoint{5.801016in}{6.294627in}}%
\pgfpathlineto{\pgfqpoint{5.801016in}{6.437785in}}%
\pgfpathlineto{\pgfqpoint{5.760145in}{6.437785in}}%
\pgfpathlineto{\pgfqpoint{5.760145in}{6.294627in}}%
\pgfpathclose%
\pgfusepath{fill}%
\end{pgfscope}%
\begin{pgfscope}%
\pgfpathrectangle{\pgfqpoint{0.787430in}{6.294627in}}{\pgfqpoint{7.062570in}{2.555373in}}%
\pgfusepath{clip}%
\pgfsetbuttcap%
\pgfsetmiterjoin%
\definecolor{currentfill}{rgb}{0.121569,0.466667,0.705882}%
\pgfsetfillcolor{currentfill}%
\pgfsetlinewidth{0.000000pt}%
\definecolor{currentstroke}{rgb}{0.000000,0.000000,0.000000}%
\pgfsetstrokecolor{currentstroke}%
\pgfsetstrokeopacity{0.000000}%
\pgfsetdash{}{0pt}%
\pgfpathmoveto{\pgfqpoint{5.801016in}{6.294627in}}%
\pgfpathlineto{\pgfqpoint{5.841888in}{6.294627in}}%
\pgfpathlineto{\pgfqpoint{5.841888in}{6.366206in}}%
\pgfpathlineto{\pgfqpoint{5.801016in}{6.366206in}}%
\pgfpathlineto{\pgfqpoint{5.801016in}{6.294627in}}%
\pgfpathclose%
\pgfusepath{fill}%
\end{pgfscope}%
\begin{pgfscope}%
\pgfpathrectangle{\pgfqpoint{0.787430in}{6.294627in}}{\pgfqpoint{7.062570in}{2.555373in}}%
\pgfusepath{clip}%
\pgfsetbuttcap%
\pgfsetmiterjoin%
\definecolor{currentfill}{rgb}{0.121569,0.466667,0.705882}%
\pgfsetfillcolor{currentfill}%
\pgfsetlinewidth{0.000000pt}%
\definecolor{currentstroke}{rgb}{0.000000,0.000000,0.000000}%
\pgfsetstrokecolor{currentstroke}%
\pgfsetstrokeopacity{0.000000}%
\pgfsetdash{}{0pt}%
\pgfpathmoveto{\pgfqpoint{5.841888in}{6.294627in}}%
\pgfpathlineto{\pgfqpoint{5.882759in}{6.294627in}}%
\pgfpathlineto{\pgfqpoint{5.882759in}{6.366206in}}%
\pgfpathlineto{\pgfqpoint{5.841888in}{6.366206in}}%
\pgfpathlineto{\pgfqpoint{5.841888in}{6.294627in}}%
\pgfpathclose%
\pgfusepath{fill}%
\end{pgfscope}%
\begin{pgfscope}%
\pgfpathrectangle{\pgfqpoint{0.787430in}{6.294627in}}{\pgfqpoint{7.062570in}{2.555373in}}%
\pgfusepath{clip}%
\pgfsetbuttcap%
\pgfsetmiterjoin%
\definecolor{currentfill}{rgb}{0.121569,0.466667,0.705882}%
\pgfsetfillcolor{currentfill}%
\pgfsetlinewidth{0.000000pt}%
\definecolor{currentstroke}{rgb}{0.000000,0.000000,0.000000}%
\pgfsetstrokecolor{currentstroke}%
\pgfsetstrokeopacity{0.000000}%
\pgfsetdash{}{0pt}%
\pgfpathmoveto{\pgfqpoint{5.882759in}{6.294627in}}%
\pgfpathlineto{\pgfqpoint{5.923630in}{6.294627in}}%
\pgfpathlineto{\pgfqpoint{5.923630in}{6.294627in}}%
\pgfpathlineto{\pgfqpoint{5.882759in}{6.294627in}}%
\pgfpathlineto{\pgfqpoint{5.882759in}{6.294627in}}%
\pgfpathclose%
\pgfusepath{fill}%
\end{pgfscope}%
\begin{pgfscope}%
\pgfpathrectangle{\pgfqpoint{0.787430in}{6.294627in}}{\pgfqpoint{7.062570in}{2.555373in}}%
\pgfusepath{clip}%
\pgfsetbuttcap%
\pgfsetmiterjoin%
\definecolor{currentfill}{rgb}{0.121569,0.466667,0.705882}%
\pgfsetfillcolor{currentfill}%
\pgfsetlinewidth{0.000000pt}%
\definecolor{currentstroke}{rgb}{0.000000,0.000000,0.000000}%
\pgfsetstrokecolor{currentstroke}%
\pgfsetstrokeopacity{0.000000}%
\pgfsetdash{}{0pt}%
\pgfpathmoveto{\pgfqpoint{5.923630in}{6.294627in}}%
\pgfpathlineto{\pgfqpoint{5.964501in}{6.294627in}}%
\pgfpathlineto{\pgfqpoint{5.964501in}{6.366206in}}%
\pgfpathlineto{\pgfqpoint{5.923630in}{6.366206in}}%
\pgfpathlineto{\pgfqpoint{5.923630in}{6.294627in}}%
\pgfpathclose%
\pgfusepath{fill}%
\end{pgfscope}%
\begin{pgfscope}%
\pgfpathrectangle{\pgfqpoint{0.787430in}{6.294627in}}{\pgfqpoint{7.062570in}{2.555373in}}%
\pgfusepath{clip}%
\pgfsetbuttcap%
\pgfsetmiterjoin%
\definecolor{currentfill}{rgb}{0.121569,0.466667,0.705882}%
\pgfsetfillcolor{currentfill}%
\pgfsetlinewidth{0.000000pt}%
\definecolor{currentstroke}{rgb}{0.000000,0.000000,0.000000}%
\pgfsetstrokecolor{currentstroke}%
\pgfsetstrokeopacity{0.000000}%
\pgfsetdash{}{0pt}%
\pgfpathmoveto{\pgfqpoint{5.964501in}{6.294627in}}%
\pgfpathlineto{\pgfqpoint{6.005372in}{6.294627in}}%
\pgfpathlineto{\pgfqpoint{6.005372in}{6.294627in}}%
\pgfpathlineto{\pgfqpoint{5.964501in}{6.294627in}}%
\pgfpathlineto{\pgfqpoint{5.964501in}{6.294627in}}%
\pgfpathclose%
\pgfusepath{fill}%
\end{pgfscope}%
\begin{pgfscope}%
\pgfpathrectangle{\pgfqpoint{0.787430in}{6.294627in}}{\pgfqpoint{7.062570in}{2.555373in}}%
\pgfusepath{clip}%
\pgfsetbuttcap%
\pgfsetmiterjoin%
\definecolor{currentfill}{rgb}{0.121569,0.466667,0.705882}%
\pgfsetfillcolor{currentfill}%
\pgfsetlinewidth{0.000000pt}%
\definecolor{currentstroke}{rgb}{0.000000,0.000000,0.000000}%
\pgfsetstrokecolor{currentstroke}%
\pgfsetstrokeopacity{0.000000}%
\pgfsetdash{}{0pt}%
\pgfpathmoveto{\pgfqpoint{6.005372in}{6.294627in}}%
\pgfpathlineto{\pgfqpoint{6.046244in}{6.294627in}}%
\pgfpathlineto{\pgfqpoint{6.046244in}{6.294627in}}%
\pgfpathlineto{\pgfqpoint{6.005372in}{6.294627in}}%
\pgfpathlineto{\pgfqpoint{6.005372in}{6.294627in}}%
\pgfpathclose%
\pgfusepath{fill}%
\end{pgfscope}%
\begin{pgfscope}%
\pgfpathrectangle{\pgfqpoint{0.787430in}{6.294627in}}{\pgfqpoint{7.062570in}{2.555373in}}%
\pgfusepath{clip}%
\pgfsetbuttcap%
\pgfsetmiterjoin%
\definecolor{currentfill}{rgb}{0.121569,0.466667,0.705882}%
\pgfsetfillcolor{currentfill}%
\pgfsetlinewidth{0.000000pt}%
\definecolor{currentstroke}{rgb}{0.000000,0.000000,0.000000}%
\pgfsetstrokecolor{currentstroke}%
\pgfsetstrokeopacity{0.000000}%
\pgfsetdash{}{0pt}%
\pgfpathmoveto{\pgfqpoint{6.046244in}{6.294627in}}%
\pgfpathlineto{\pgfqpoint{6.087115in}{6.294627in}}%
\pgfpathlineto{\pgfqpoint{6.087115in}{6.294627in}}%
\pgfpathlineto{\pgfqpoint{6.046244in}{6.294627in}}%
\pgfpathlineto{\pgfqpoint{6.046244in}{6.294627in}}%
\pgfpathclose%
\pgfusepath{fill}%
\end{pgfscope}%
\begin{pgfscope}%
\pgfpathrectangle{\pgfqpoint{0.787430in}{6.294627in}}{\pgfqpoint{7.062570in}{2.555373in}}%
\pgfusepath{clip}%
\pgfsetbuttcap%
\pgfsetmiterjoin%
\definecolor{currentfill}{rgb}{0.121569,0.466667,0.705882}%
\pgfsetfillcolor{currentfill}%
\pgfsetlinewidth{0.000000pt}%
\definecolor{currentstroke}{rgb}{0.000000,0.000000,0.000000}%
\pgfsetstrokecolor{currentstroke}%
\pgfsetstrokeopacity{0.000000}%
\pgfsetdash{}{0pt}%
\pgfpathmoveto{\pgfqpoint{6.087115in}{6.294627in}}%
\pgfpathlineto{\pgfqpoint{6.127986in}{6.294627in}}%
\pgfpathlineto{\pgfqpoint{6.127986in}{6.366206in}}%
\pgfpathlineto{\pgfqpoint{6.087115in}{6.366206in}}%
\pgfpathlineto{\pgfqpoint{6.087115in}{6.294627in}}%
\pgfpathclose%
\pgfusepath{fill}%
\end{pgfscope}%
\begin{pgfscope}%
\pgfpathrectangle{\pgfqpoint{0.787430in}{6.294627in}}{\pgfqpoint{7.062570in}{2.555373in}}%
\pgfusepath{clip}%
\pgfsetbuttcap%
\pgfsetmiterjoin%
\definecolor{currentfill}{rgb}{0.121569,0.466667,0.705882}%
\pgfsetfillcolor{currentfill}%
\pgfsetlinewidth{0.000000pt}%
\definecolor{currentstroke}{rgb}{0.000000,0.000000,0.000000}%
\pgfsetstrokecolor{currentstroke}%
\pgfsetstrokeopacity{0.000000}%
\pgfsetdash{}{0pt}%
\pgfpathmoveto{\pgfqpoint{6.127986in}{6.294627in}}%
\pgfpathlineto{\pgfqpoint{6.168857in}{6.294627in}}%
\pgfpathlineto{\pgfqpoint{6.168857in}{6.294627in}}%
\pgfpathlineto{\pgfqpoint{6.127986in}{6.294627in}}%
\pgfpathlineto{\pgfqpoint{6.127986in}{6.294627in}}%
\pgfpathclose%
\pgfusepath{fill}%
\end{pgfscope}%
\begin{pgfscope}%
\pgfpathrectangle{\pgfqpoint{0.787430in}{6.294627in}}{\pgfqpoint{7.062570in}{2.555373in}}%
\pgfusepath{clip}%
\pgfsetbuttcap%
\pgfsetmiterjoin%
\definecolor{currentfill}{rgb}{0.121569,0.466667,0.705882}%
\pgfsetfillcolor{currentfill}%
\pgfsetlinewidth{0.000000pt}%
\definecolor{currentstroke}{rgb}{0.000000,0.000000,0.000000}%
\pgfsetstrokecolor{currentstroke}%
\pgfsetstrokeopacity{0.000000}%
\pgfsetdash{}{0pt}%
\pgfpathmoveto{\pgfqpoint{6.168857in}{6.294627in}}%
\pgfpathlineto{\pgfqpoint{6.209728in}{6.294627in}}%
\pgfpathlineto{\pgfqpoint{6.209728in}{6.294627in}}%
\pgfpathlineto{\pgfqpoint{6.168857in}{6.294627in}}%
\pgfpathlineto{\pgfqpoint{6.168857in}{6.294627in}}%
\pgfpathclose%
\pgfusepath{fill}%
\end{pgfscope}%
\begin{pgfscope}%
\pgfpathrectangle{\pgfqpoint{0.787430in}{6.294627in}}{\pgfqpoint{7.062570in}{2.555373in}}%
\pgfusepath{clip}%
\pgfsetbuttcap%
\pgfsetmiterjoin%
\definecolor{currentfill}{rgb}{0.121569,0.466667,0.705882}%
\pgfsetfillcolor{currentfill}%
\pgfsetlinewidth{0.000000pt}%
\definecolor{currentstroke}{rgb}{0.000000,0.000000,0.000000}%
\pgfsetstrokecolor{currentstroke}%
\pgfsetstrokeopacity{0.000000}%
\pgfsetdash{}{0pt}%
\pgfpathmoveto{\pgfqpoint{6.209728in}{6.294627in}}%
\pgfpathlineto{\pgfqpoint{6.250600in}{6.294627in}}%
\pgfpathlineto{\pgfqpoint{6.250600in}{6.294627in}}%
\pgfpathlineto{\pgfqpoint{6.209728in}{6.294627in}}%
\pgfpathlineto{\pgfqpoint{6.209728in}{6.294627in}}%
\pgfpathclose%
\pgfusepath{fill}%
\end{pgfscope}%
\begin{pgfscope}%
\pgfpathrectangle{\pgfqpoint{0.787430in}{6.294627in}}{\pgfqpoint{7.062570in}{2.555373in}}%
\pgfusepath{clip}%
\pgfsetbuttcap%
\pgfsetmiterjoin%
\definecolor{currentfill}{rgb}{0.121569,0.466667,0.705882}%
\pgfsetfillcolor{currentfill}%
\pgfsetlinewidth{0.000000pt}%
\definecolor{currentstroke}{rgb}{0.000000,0.000000,0.000000}%
\pgfsetstrokecolor{currentstroke}%
\pgfsetstrokeopacity{0.000000}%
\pgfsetdash{}{0pt}%
\pgfpathmoveto{\pgfqpoint{6.250600in}{6.294627in}}%
\pgfpathlineto{\pgfqpoint{6.291471in}{6.294627in}}%
\pgfpathlineto{\pgfqpoint{6.291471in}{6.366206in}}%
\pgfpathlineto{\pgfqpoint{6.250600in}{6.366206in}}%
\pgfpathlineto{\pgfqpoint{6.250600in}{6.294627in}}%
\pgfpathclose%
\pgfusepath{fill}%
\end{pgfscope}%
\begin{pgfscope}%
\pgfpathrectangle{\pgfqpoint{0.787430in}{6.294627in}}{\pgfqpoint{7.062570in}{2.555373in}}%
\pgfusepath{clip}%
\pgfsetbuttcap%
\pgfsetmiterjoin%
\definecolor{currentfill}{rgb}{0.121569,0.466667,0.705882}%
\pgfsetfillcolor{currentfill}%
\pgfsetlinewidth{0.000000pt}%
\definecolor{currentstroke}{rgb}{0.000000,0.000000,0.000000}%
\pgfsetstrokecolor{currentstroke}%
\pgfsetstrokeopacity{0.000000}%
\pgfsetdash{}{0pt}%
\pgfpathmoveto{\pgfqpoint{6.291471in}{6.294627in}}%
\pgfpathlineto{\pgfqpoint{6.332342in}{6.294627in}}%
\pgfpathlineto{\pgfqpoint{6.332342in}{6.294627in}}%
\pgfpathlineto{\pgfqpoint{6.291471in}{6.294627in}}%
\pgfpathlineto{\pgfqpoint{6.291471in}{6.294627in}}%
\pgfpathclose%
\pgfusepath{fill}%
\end{pgfscope}%
\begin{pgfscope}%
\pgfpathrectangle{\pgfqpoint{0.787430in}{6.294627in}}{\pgfqpoint{7.062570in}{2.555373in}}%
\pgfusepath{clip}%
\pgfsetbuttcap%
\pgfsetmiterjoin%
\definecolor{currentfill}{rgb}{0.121569,0.466667,0.705882}%
\pgfsetfillcolor{currentfill}%
\pgfsetlinewidth{0.000000pt}%
\definecolor{currentstroke}{rgb}{0.000000,0.000000,0.000000}%
\pgfsetstrokecolor{currentstroke}%
\pgfsetstrokeopacity{0.000000}%
\pgfsetdash{}{0pt}%
\pgfpathmoveto{\pgfqpoint{6.332342in}{6.294627in}}%
\pgfpathlineto{\pgfqpoint{6.373213in}{6.294627in}}%
\pgfpathlineto{\pgfqpoint{6.373213in}{6.294627in}}%
\pgfpathlineto{\pgfqpoint{6.332342in}{6.294627in}}%
\pgfpathlineto{\pgfqpoint{6.332342in}{6.294627in}}%
\pgfpathclose%
\pgfusepath{fill}%
\end{pgfscope}%
\begin{pgfscope}%
\pgfpathrectangle{\pgfqpoint{0.787430in}{6.294627in}}{\pgfqpoint{7.062570in}{2.555373in}}%
\pgfusepath{clip}%
\pgfsetbuttcap%
\pgfsetmiterjoin%
\definecolor{currentfill}{rgb}{0.121569,0.466667,0.705882}%
\pgfsetfillcolor{currentfill}%
\pgfsetlinewidth{0.000000pt}%
\definecolor{currentstroke}{rgb}{0.000000,0.000000,0.000000}%
\pgfsetstrokecolor{currentstroke}%
\pgfsetstrokeopacity{0.000000}%
\pgfsetdash{}{0pt}%
\pgfpathmoveto{\pgfqpoint{6.373213in}{6.294627in}}%
\pgfpathlineto{\pgfqpoint{6.414084in}{6.294627in}}%
\pgfpathlineto{\pgfqpoint{6.414084in}{6.294627in}}%
\pgfpathlineto{\pgfqpoint{6.373213in}{6.294627in}}%
\pgfpathlineto{\pgfqpoint{6.373213in}{6.294627in}}%
\pgfpathclose%
\pgfusepath{fill}%
\end{pgfscope}%
\begin{pgfscope}%
\pgfpathrectangle{\pgfqpoint{0.787430in}{6.294627in}}{\pgfqpoint{7.062570in}{2.555373in}}%
\pgfusepath{clip}%
\pgfsetbuttcap%
\pgfsetmiterjoin%
\definecolor{currentfill}{rgb}{0.121569,0.466667,0.705882}%
\pgfsetfillcolor{currentfill}%
\pgfsetlinewidth{0.000000pt}%
\definecolor{currentstroke}{rgb}{0.000000,0.000000,0.000000}%
\pgfsetstrokecolor{currentstroke}%
\pgfsetstrokeopacity{0.000000}%
\pgfsetdash{}{0pt}%
\pgfpathmoveto{\pgfqpoint{6.414084in}{6.294627in}}%
\pgfpathlineto{\pgfqpoint{6.454956in}{6.294627in}}%
\pgfpathlineto{\pgfqpoint{6.454956in}{6.294627in}}%
\pgfpathlineto{\pgfqpoint{6.414084in}{6.294627in}}%
\pgfpathlineto{\pgfqpoint{6.414084in}{6.294627in}}%
\pgfpathclose%
\pgfusepath{fill}%
\end{pgfscope}%
\begin{pgfscope}%
\pgfpathrectangle{\pgfqpoint{0.787430in}{6.294627in}}{\pgfqpoint{7.062570in}{2.555373in}}%
\pgfusepath{clip}%
\pgfsetbuttcap%
\pgfsetmiterjoin%
\definecolor{currentfill}{rgb}{0.121569,0.466667,0.705882}%
\pgfsetfillcolor{currentfill}%
\pgfsetlinewidth{0.000000pt}%
\definecolor{currentstroke}{rgb}{0.000000,0.000000,0.000000}%
\pgfsetstrokecolor{currentstroke}%
\pgfsetstrokeopacity{0.000000}%
\pgfsetdash{}{0pt}%
\pgfpathmoveto{\pgfqpoint{6.454956in}{6.294627in}}%
\pgfpathlineto{\pgfqpoint{6.495827in}{6.294627in}}%
\pgfpathlineto{\pgfqpoint{6.495827in}{6.294627in}}%
\pgfpathlineto{\pgfqpoint{6.454956in}{6.294627in}}%
\pgfpathlineto{\pgfqpoint{6.454956in}{6.294627in}}%
\pgfpathclose%
\pgfusepath{fill}%
\end{pgfscope}%
\begin{pgfscope}%
\pgfpathrectangle{\pgfqpoint{0.787430in}{6.294627in}}{\pgfqpoint{7.062570in}{2.555373in}}%
\pgfusepath{clip}%
\pgfsetbuttcap%
\pgfsetmiterjoin%
\definecolor{currentfill}{rgb}{0.121569,0.466667,0.705882}%
\pgfsetfillcolor{currentfill}%
\pgfsetlinewidth{0.000000pt}%
\definecolor{currentstroke}{rgb}{0.000000,0.000000,0.000000}%
\pgfsetstrokecolor{currentstroke}%
\pgfsetstrokeopacity{0.000000}%
\pgfsetdash{}{0pt}%
\pgfpathmoveto{\pgfqpoint{6.495827in}{6.294627in}}%
\pgfpathlineto{\pgfqpoint{6.536698in}{6.294627in}}%
\pgfpathlineto{\pgfqpoint{6.536698in}{6.294627in}}%
\pgfpathlineto{\pgfqpoint{6.495827in}{6.294627in}}%
\pgfpathlineto{\pgfqpoint{6.495827in}{6.294627in}}%
\pgfpathclose%
\pgfusepath{fill}%
\end{pgfscope}%
\begin{pgfscope}%
\pgfpathrectangle{\pgfqpoint{0.787430in}{6.294627in}}{\pgfqpoint{7.062570in}{2.555373in}}%
\pgfusepath{clip}%
\pgfsetbuttcap%
\pgfsetmiterjoin%
\definecolor{currentfill}{rgb}{0.121569,0.466667,0.705882}%
\pgfsetfillcolor{currentfill}%
\pgfsetlinewidth{0.000000pt}%
\definecolor{currentstroke}{rgb}{0.000000,0.000000,0.000000}%
\pgfsetstrokecolor{currentstroke}%
\pgfsetstrokeopacity{0.000000}%
\pgfsetdash{}{0pt}%
\pgfpathmoveto{\pgfqpoint{6.536698in}{6.294627in}}%
\pgfpathlineto{\pgfqpoint{6.577569in}{6.294627in}}%
\pgfpathlineto{\pgfqpoint{6.577569in}{6.294627in}}%
\pgfpathlineto{\pgfqpoint{6.536698in}{6.294627in}}%
\pgfpathlineto{\pgfqpoint{6.536698in}{6.294627in}}%
\pgfpathclose%
\pgfusepath{fill}%
\end{pgfscope}%
\begin{pgfscope}%
\pgfpathrectangle{\pgfqpoint{0.787430in}{6.294627in}}{\pgfqpoint{7.062570in}{2.555373in}}%
\pgfusepath{clip}%
\pgfsetbuttcap%
\pgfsetmiterjoin%
\definecolor{currentfill}{rgb}{0.121569,0.466667,0.705882}%
\pgfsetfillcolor{currentfill}%
\pgfsetlinewidth{0.000000pt}%
\definecolor{currentstroke}{rgb}{0.000000,0.000000,0.000000}%
\pgfsetstrokecolor{currentstroke}%
\pgfsetstrokeopacity{0.000000}%
\pgfsetdash{}{0pt}%
\pgfpathmoveto{\pgfqpoint{6.577569in}{6.294627in}}%
\pgfpathlineto{\pgfqpoint{6.618440in}{6.294627in}}%
\pgfpathlineto{\pgfqpoint{6.618440in}{6.294627in}}%
\pgfpathlineto{\pgfqpoint{6.577569in}{6.294627in}}%
\pgfpathlineto{\pgfqpoint{6.577569in}{6.294627in}}%
\pgfpathclose%
\pgfusepath{fill}%
\end{pgfscope}%
\begin{pgfscope}%
\pgfpathrectangle{\pgfqpoint{0.787430in}{6.294627in}}{\pgfqpoint{7.062570in}{2.555373in}}%
\pgfusepath{clip}%
\pgfsetbuttcap%
\pgfsetmiterjoin%
\definecolor{currentfill}{rgb}{0.121569,0.466667,0.705882}%
\pgfsetfillcolor{currentfill}%
\pgfsetlinewidth{0.000000pt}%
\definecolor{currentstroke}{rgb}{0.000000,0.000000,0.000000}%
\pgfsetstrokecolor{currentstroke}%
\pgfsetstrokeopacity{0.000000}%
\pgfsetdash{}{0pt}%
\pgfpathmoveto{\pgfqpoint{6.618440in}{6.294627in}}%
\pgfpathlineto{\pgfqpoint{6.659312in}{6.294627in}}%
\pgfpathlineto{\pgfqpoint{6.659312in}{6.366206in}}%
\pgfpathlineto{\pgfqpoint{6.618440in}{6.366206in}}%
\pgfpathlineto{\pgfqpoint{6.618440in}{6.294627in}}%
\pgfpathclose%
\pgfusepath{fill}%
\end{pgfscope}%
\begin{pgfscope}%
\pgfsetbuttcap%
\pgfsetroundjoin%
\definecolor{currentfill}{rgb}{0.000000,0.000000,0.000000}%
\pgfsetfillcolor{currentfill}%
\pgfsetlinewidth{0.803000pt}%
\definecolor{currentstroke}{rgb}{0.000000,0.000000,0.000000}%
\pgfsetstrokecolor{currentstroke}%
\pgfsetdash{}{0pt}%
\pgfsys@defobject{currentmarker}{\pgfqpoint{0.000000in}{-0.048611in}}{\pgfqpoint{0.000000in}{0.000000in}}{%
\pgfpathmoveto{\pgfqpoint{0.000000in}{0.000000in}}%
\pgfpathlineto{\pgfqpoint{0.000000in}{-0.048611in}}%
\pgfusepath{stroke,fill}%
}%
\begin{pgfscope}%
\pgfsys@transformshift{1.923827in}{6.294627in}%
\pgfsys@useobject{currentmarker}{}%
\end{pgfscope}%
\end{pgfscope}%
\begin{pgfscope}%
\pgfsetbuttcap%
\pgfsetroundjoin%
\definecolor{currentfill}{rgb}{0.000000,0.000000,0.000000}%
\pgfsetfillcolor{currentfill}%
\pgfsetlinewidth{0.803000pt}%
\definecolor{currentstroke}{rgb}{0.000000,0.000000,0.000000}%
\pgfsetstrokecolor{currentstroke}%
\pgfsetdash{}{0pt}%
\pgfsys@defobject{currentmarker}{\pgfqpoint{0.000000in}{-0.048611in}}{\pgfqpoint{0.000000in}{0.000000in}}{%
\pgfpathmoveto{\pgfqpoint{0.000000in}{0.000000in}}%
\pgfpathlineto{\pgfqpoint{0.000000in}{-0.048611in}}%
\pgfusepath{stroke,fill}%
}%
\begin{pgfscope}%
\pgfsys@transformshift{3.153546in}{6.294627in}%
\pgfsys@useobject{currentmarker}{}%
\end{pgfscope}%
\end{pgfscope}%
\begin{pgfscope}%
\pgfsetbuttcap%
\pgfsetroundjoin%
\definecolor{currentfill}{rgb}{0.000000,0.000000,0.000000}%
\pgfsetfillcolor{currentfill}%
\pgfsetlinewidth{0.803000pt}%
\definecolor{currentstroke}{rgb}{0.000000,0.000000,0.000000}%
\pgfsetstrokecolor{currentstroke}%
\pgfsetdash{}{0pt}%
\pgfsys@defobject{currentmarker}{\pgfqpoint{0.000000in}{-0.048611in}}{\pgfqpoint{0.000000in}{0.000000in}}{%
\pgfpathmoveto{\pgfqpoint{0.000000in}{0.000000in}}%
\pgfpathlineto{\pgfqpoint{0.000000in}{-0.048611in}}%
\pgfusepath{stroke,fill}%
}%
\begin{pgfscope}%
\pgfsys@transformshift{4.383264in}{6.294627in}%
\pgfsys@useobject{currentmarker}{}%
\end{pgfscope}%
\end{pgfscope}%
\begin{pgfscope}%
\pgfsetbuttcap%
\pgfsetroundjoin%
\definecolor{currentfill}{rgb}{0.000000,0.000000,0.000000}%
\pgfsetfillcolor{currentfill}%
\pgfsetlinewidth{0.803000pt}%
\definecolor{currentstroke}{rgb}{0.000000,0.000000,0.000000}%
\pgfsetstrokecolor{currentstroke}%
\pgfsetdash{}{0pt}%
\pgfsys@defobject{currentmarker}{\pgfqpoint{0.000000in}{-0.048611in}}{\pgfqpoint{0.000000in}{0.000000in}}{%
\pgfpathmoveto{\pgfqpoint{0.000000in}{0.000000in}}%
\pgfpathlineto{\pgfqpoint{0.000000in}{-0.048611in}}%
\pgfusepath{stroke,fill}%
}%
\begin{pgfscope}%
\pgfsys@transformshift{5.612983in}{6.294627in}%
\pgfsys@useobject{currentmarker}{}%
\end{pgfscope}%
\end{pgfscope}%
\begin{pgfscope}%
\pgfsetbuttcap%
\pgfsetroundjoin%
\definecolor{currentfill}{rgb}{0.000000,0.000000,0.000000}%
\pgfsetfillcolor{currentfill}%
\pgfsetlinewidth{0.803000pt}%
\definecolor{currentstroke}{rgb}{0.000000,0.000000,0.000000}%
\pgfsetstrokecolor{currentstroke}%
\pgfsetdash{}{0pt}%
\pgfsys@defobject{currentmarker}{\pgfqpoint{0.000000in}{-0.048611in}}{\pgfqpoint{0.000000in}{0.000000in}}{%
\pgfpathmoveto{\pgfqpoint{0.000000in}{0.000000in}}%
\pgfpathlineto{\pgfqpoint{0.000000in}{-0.048611in}}%
\pgfusepath{stroke,fill}%
}%
\begin{pgfscope}%
\pgfsys@transformshift{6.842701in}{6.294627in}%
\pgfsys@useobject{currentmarker}{}%
\end{pgfscope}%
\end{pgfscope}%
\begin{pgfscope}%
\pgfsetbuttcap%
\pgfsetroundjoin%
\definecolor{currentfill}{rgb}{0.000000,0.000000,0.000000}%
\pgfsetfillcolor{currentfill}%
\pgfsetlinewidth{0.803000pt}%
\definecolor{currentstroke}{rgb}{0.000000,0.000000,0.000000}%
\pgfsetstrokecolor{currentstroke}%
\pgfsetdash{}{0pt}%
\pgfsys@defobject{currentmarker}{\pgfqpoint{-0.048611in}{0.000000in}}{\pgfqpoint{-0.000000in}{0.000000in}}{%
\pgfpathmoveto{\pgfqpoint{-0.000000in}{0.000000in}}%
\pgfpathlineto{\pgfqpoint{-0.048611in}{0.000000in}}%
\pgfusepath{stroke,fill}%
}%
\begin{pgfscope}%
\pgfsys@transformshift{0.787430in}{6.294627in}%
\pgfsys@useobject{currentmarker}{}%
\end{pgfscope}%
\end{pgfscope}%
\begin{pgfscope}%
\definecolor{textcolor}{rgb}{0.000000,0.000000,0.000000}%
\pgfsetstrokecolor{textcolor}%
\pgfsetfillcolor{textcolor}%
\pgftext[x=0.512738in, y=6.246401in, left, base]{\color{textcolor}{\rmfamily\fontsize{10.000000}{12.000000}\selectfont\catcode`\^=\active\def^{\ifmmode\sp\else\^{}\fi}\catcode`\%=\active\def%{\%}$\mathdefault{0.0}$}}%
\end{pgfscope}%
\begin{pgfscope}%
\pgfsetbuttcap%
\pgfsetroundjoin%
\definecolor{currentfill}{rgb}{0.000000,0.000000,0.000000}%
\pgfsetfillcolor{currentfill}%
\pgfsetlinewidth{0.803000pt}%
\definecolor{currentstroke}{rgb}{0.000000,0.000000,0.000000}%
\pgfsetstrokecolor{currentstroke}%
\pgfsetdash{}{0pt}%
\pgfsys@defobject{currentmarker}{\pgfqpoint{-0.048611in}{0.000000in}}{\pgfqpoint{-0.000000in}{0.000000in}}{%
\pgfpathmoveto{\pgfqpoint{-0.000000in}{0.000000in}}%
\pgfpathlineto{\pgfqpoint{-0.048611in}{0.000000in}}%
\pgfusepath{stroke,fill}%
}%
\begin{pgfscope}%
\pgfsys@transformshift{0.787430in}{7.246234in}%
\pgfsys@useobject{currentmarker}{}%
\end{pgfscope}%
\end{pgfscope}%
\begin{pgfscope}%
\definecolor{textcolor}{rgb}{0.000000,0.000000,0.000000}%
\pgfsetstrokecolor{textcolor}%
\pgfsetfillcolor{textcolor}%
\pgftext[x=0.512738in, y=7.198009in, left, base]{\color{textcolor}{\rmfamily\fontsize{10.000000}{12.000000}\selectfont\catcode`\^=\active\def^{\ifmmode\sp\else\^{}\fi}\catcode`\%=\active\def%{\%}$\mathdefault{0.2}$}}%
\end{pgfscope}%
\begin{pgfscope}%
\pgfsetbuttcap%
\pgfsetroundjoin%
\definecolor{currentfill}{rgb}{0.000000,0.000000,0.000000}%
\pgfsetfillcolor{currentfill}%
\pgfsetlinewidth{0.803000pt}%
\definecolor{currentstroke}{rgb}{0.000000,0.000000,0.000000}%
\pgfsetstrokecolor{currentstroke}%
\pgfsetdash{}{0pt}%
\pgfsys@defobject{currentmarker}{\pgfqpoint{-0.048611in}{0.000000in}}{\pgfqpoint{-0.000000in}{0.000000in}}{%
\pgfpathmoveto{\pgfqpoint{-0.000000in}{0.000000in}}%
\pgfpathlineto{\pgfqpoint{-0.048611in}{0.000000in}}%
\pgfusepath{stroke,fill}%
}%
\begin{pgfscope}%
\pgfsys@transformshift{0.787430in}{8.197841in}%
\pgfsys@useobject{currentmarker}{}%
\end{pgfscope}%
\end{pgfscope}%
\begin{pgfscope}%
\definecolor{textcolor}{rgb}{0.000000,0.000000,0.000000}%
\pgfsetstrokecolor{textcolor}%
\pgfsetfillcolor{textcolor}%
\pgftext[x=0.512738in, y=8.149616in, left, base]{\color{textcolor}{\rmfamily\fontsize{10.000000}{12.000000}\selectfont\catcode`\^=\active\def^{\ifmmode\sp\else\^{}\fi}\catcode`\%=\active\def%{\%}$\mathdefault{0.4}$}}%
\end{pgfscope}%
\begin{pgfscope}%
\pgfpathrectangle{\pgfqpoint{0.787430in}{6.294627in}}{\pgfqpoint{7.062570in}{2.555373in}}%
\pgfusepath{clip}%
\pgfsetrectcap%
\pgfsetroundjoin%
\pgfsetlinewidth{2.007500pt}%
\definecolor{currentstroke}{rgb}{1.000000,0.000000,0.000000}%
\pgfsetstrokecolor{currentstroke}%
\pgfsetdash{}{0pt}%
\pgfpathmoveto{\pgfqpoint{1.923827in}{6.295263in}}%
\pgfpathlineto{\pgfqpoint{2.155246in}{6.297300in}}%
\pgfpathlineto{\pgfqpoint{2.298036in}{6.300664in}}%
\pgfpathlineto{\pgfqpoint{2.401436in}{6.305155in}}%
\pgfpathlineto{\pgfqpoint{2.490064in}{6.311208in}}%
\pgfpathlineto{\pgfqpoint{2.563921in}{6.318456in}}%
\pgfpathlineto{\pgfqpoint{2.627930in}{6.326878in}}%
\pgfpathlineto{\pgfqpoint{2.687016in}{6.336863in}}%
\pgfpathlineto{\pgfqpoint{2.741178in}{6.348273in}}%
\pgfpathlineto{\pgfqpoint{2.790416in}{6.360853in}}%
\pgfpathlineto{\pgfqpoint{2.839654in}{6.375860in}}%
\pgfpathlineto{\pgfqpoint{2.883968in}{6.391719in}}%
\pgfpathlineto{\pgfqpoint{2.928282in}{6.410073in}}%
\pgfpathlineto{\pgfqpoint{2.967672in}{6.428696in}}%
\pgfpathlineto{\pgfqpoint{3.007063in}{6.449684in}}%
\pgfpathlineto{\pgfqpoint{3.046453in}{6.473224in}}%
\pgfpathlineto{\pgfqpoint{3.085844in}{6.499495in}}%
\pgfpathlineto{\pgfqpoint{3.125234in}{6.528668in}}%
\pgfpathlineto{\pgfqpoint{3.164624in}{6.560901in}}%
\pgfpathlineto{\pgfqpoint{3.204015in}{6.596331in}}%
\pgfpathlineto{\pgfqpoint{3.243405in}{6.635076in}}%
\pgfpathlineto{\pgfqpoint{3.282796in}{6.677222in}}%
\pgfpathlineto{\pgfqpoint{3.322186in}{6.722825in}}%
\pgfpathlineto{\pgfqpoint{3.361576in}{6.771901in}}%
\pgfpathlineto{\pgfqpoint{3.400967in}{6.824423in}}%
\pgfpathlineto{\pgfqpoint{3.445281in}{6.887534in}}%
\pgfpathlineto{\pgfqpoint{3.489595in}{6.954725in}}%
\pgfpathlineto{\pgfqpoint{3.538833in}{7.033831in}}%
\pgfpathlineto{\pgfqpoint{3.588071in}{7.117126in}}%
\pgfpathlineto{\pgfqpoint{3.647157in}{7.221680in}}%
\pgfpathlineto{\pgfqpoint{3.721014in}{7.357365in}}%
\pgfpathlineto{\pgfqpoint{3.898270in}{7.685320in}}%
\pgfpathlineto{\pgfqpoint{3.952432in}{7.779614in}}%
\pgfpathlineto{\pgfqpoint{3.996746in}{7.852481in}}%
\pgfpathlineto{\pgfqpoint{4.036137in}{7.913176in}}%
\pgfpathlineto{\pgfqpoint{4.070603in}{7.962595in}}%
\pgfpathlineto{\pgfqpoint{4.105070in}{8.008130in}}%
\pgfpathlineto{\pgfqpoint{4.134613in}{8.043768in}}%
\pgfpathlineto{\pgfqpoint{4.164155in}{8.076030in}}%
\pgfpathlineto{\pgfqpoint{4.193698in}{8.104704in}}%
\pgfpathlineto{\pgfqpoint{4.218317in}{8.125719in}}%
\pgfpathlineto{\pgfqpoint{4.242936in}{8.144011in}}%
\pgfpathlineto{\pgfqpoint{4.267555in}{8.159493in}}%
\pgfpathlineto{\pgfqpoint{4.292174in}{8.172092in}}%
\pgfpathlineto{\pgfqpoint{4.316793in}{8.181749in}}%
\pgfpathlineto{\pgfqpoint{4.336488in}{8.187324in}}%
\pgfpathlineto{\pgfqpoint{4.356183in}{8.190969in}}%
\pgfpathlineto{\pgfqpoint{4.375879in}{8.192672in}}%
\pgfpathlineto{\pgfqpoint{4.395574in}{8.192428in}}%
\pgfpathlineto{\pgfqpoint{4.415269in}{8.190239in}}%
\pgfpathlineto{\pgfqpoint{4.434964in}{8.186110in}}%
\pgfpathlineto{\pgfqpoint{4.454659in}{8.180055in}}%
\pgfpathlineto{\pgfqpoint{4.474355in}{8.172092in}}%
\pgfpathlineto{\pgfqpoint{4.498974in}{8.159493in}}%
\pgfpathlineto{\pgfqpoint{4.523593in}{8.144011in}}%
\pgfpathlineto{\pgfqpoint{4.548212in}{8.125719in}}%
\pgfpathlineto{\pgfqpoint{4.572831in}{8.104704in}}%
\pgfpathlineto{\pgfqpoint{4.597450in}{8.081064in}}%
\pgfpathlineto{\pgfqpoint{4.626992in}{8.049386in}}%
\pgfpathlineto{\pgfqpoint{4.656535in}{8.014295in}}%
\pgfpathlineto{\pgfqpoint{4.686078in}{7.976019in}}%
\pgfpathlineto{\pgfqpoint{4.720544in}{7.927668in}}%
\pgfpathlineto{\pgfqpoint{4.759935in}{7.868043in}}%
\pgfpathlineto{\pgfqpoint{4.799325in}{7.804387in}}%
\pgfpathlineto{\pgfqpoint{4.848563in}{7.720175in}}%
\pgfpathlineto{\pgfqpoint{4.907649in}{7.614080in}}%
\pgfpathlineto{\pgfqpoint{4.996277in}{7.449390in}}%
\pgfpathlineto{\pgfqpoint{5.109524in}{7.239506in}}%
\pgfpathlineto{\pgfqpoint{5.173534in}{7.125662in}}%
\pgfpathlineto{\pgfqpoint{5.227696in}{7.033831in}}%
\pgfpathlineto{\pgfqpoint{5.276934in}{6.954725in}}%
\pgfpathlineto{\pgfqpoint{5.321248in}{6.887534in}}%
\pgfpathlineto{\pgfqpoint{5.365562in}{6.824423in}}%
\pgfpathlineto{\pgfqpoint{5.404952in}{6.771901in}}%
\pgfpathlineto{\pgfqpoint{5.444343in}{6.722825in}}%
\pgfpathlineto{\pgfqpoint{5.483733in}{6.677222in}}%
\pgfpathlineto{\pgfqpoint{5.523124in}{6.635076in}}%
\pgfpathlineto{\pgfqpoint{5.562514in}{6.596331in}}%
\pgfpathlineto{\pgfqpoint{5.601904in}{6.560901in}}%
\pgfpathlineto{\pgfqpoint{5.641295in}{6.528668in}}%
\pgfpathlineto{\pgfqpoint{5.680685in}{6.499495in}}%
\pgfpathlineto{\pgfqpoint{5.720075in}{6.473224in}}%
\pgfpathlineto{\pgfqpoint{5.759466in}{6.449684in}}%
\pgfpathlineto{\pgfqpoint{5.798856in}{6.428696in}}%
\pgfpathlineto{\pgfqpoint{5.838247in}{6.410073in}}%
\pgfpathlineto{\pgfqpoint{5.882561in}{6.391719in}}%
\pgfpathlineto{\pgfqpoint{5.926875in}{6.375860in}}%
\pgfpathlineto{\pgfqpoint{5.971189in}{6.362239in}}%
\pgfpathlineto{\pgfqpoint{6.020427in}{6.349431in}}%
\pgfpathlineto{\pgfqpoint{6.074589in}{6.337805in}}%
\pgfpathlineto{\pgfqpoint{6.133674in}{6.327623in}}%
\pgfpathlineto{\pgfqpoint{6.197684in}{6.319027in}}%
\pgfpathlineto{\pgfqpoint{6.266617in}{6.312044in}}%
\pgfpathlineto{\pgfqpoint{6.345398in}{6.306294in}}%
\pgfpathlineto{\pgfqpoint{6.438950in}{6.301724in}}%
\pgfpathlineto{\pgfqpoint{6.557121in}{6.298291in}}%
\pgfpathlineto{\pgfqpoint{6.719606in}{6.296017in}}%
\pgfpathlineto{\pgfqpoint{6.842701in}{6.295263in}}%
\pgfpathlineto{\pgfqpoint{6.842701in}{6.295263in}}%
\pgfusepath{stroke}%
\end{pgfscope}%
\begin{pgfscope}%
\pgfsetrectcap%
\pgfsetmiterjoin%
\pgfsetlinewidth{0.803000pt}%
\definecolor{currentstroke}{rgb}{0.000000,0.000000,0.000000}%
\pgfsetstrokecolor{currentstroke}%
\pgfsetdash{}{0pt}%
\pgfpathmoveto{\pgfqpoint{0.787430in}{6.294627in}}%
\pgfpathlineto{\pgfqpoint{0.787430in}{8.850000in}}%
\pgfusepath{stroke}%
\end{pgfscope}%
\begin{pgfscope}%
\pgfsetrectcap%
\pgfsetmiterjoin%
\pgfsetlinewidth{0.803000pt}%
\definecolor{currentstroke}{rgb}{0.000000,0.000000,0.000000}%
\pgfsetstrokecolor{currentstroke}%
\pgfsetdash{}{0pt}%
\pgfpathmoveto{\pgfqpoint{7.850000in}{6.294627in}}%
\pgfpathlineto{\pgfqpoint{7.850000in}{8.850000in}}%
\pgfusepath{stroke}%
\end{pgfscope}%
\begin{pgfscope}%
\pgfsetrectcap%
\pgfsetmiterjoin%
\pgfsetlinewidth{0.803000pt}%
\definecolor{currentstroke}{rgb}{0.000000,0.000000,0.000000}%
\pgfsetstrokecolor{currentstroke}%
\pgfsetdash{}{0pt}%
\pgfpathmoveto{\pgfqpoint{0.787430in}{6.294627in}}%
\pgfpathlineto{\pgfqpoint{7.850000in}{6.294627in}}%
\pgfusepath{stroke}%
\end{pgfscope}%
\begin{pgfscope}%
\pgfsetrectcap%
\pgfsetmiterjoin%
\pgfsetlinewidth{0.803000pt}%
\definecolor{currentstroke}{rgb}{0.000000,0.000000,0.000000}%
\pgfsetstrokecolor{currentstroke}%
\pgfsetdash{}{0pt}%
\pgfpathmoveto{\pgfqpoint{0.787430in}{8.850000in}}%
\pgfpathlineto{\pgfqpoint{7.850000in}{8.850000in}}%
\pgfusepath{stroke}%
\end{pgfscope}%
\begin{pgfscope}%
\pgfsetbuttcap%
\pgfsetmiterjoin%
\definecolor{currentfill}{rgb}{1.000000,1.000000,1.000000}%
\pgfsetfillcolor{currentfill}%
\pgfsetfillopacity{0.800000}%
\pgfsetlinewidth{1.003750pt}%
\definecolor{currentstroke}{rgb}{0.800000,0.800000,0.800000}%
\pgfsetstrokecolor{currentstroke}%
\pgfsetstrokeopacity{0.800000}%
\pgfsetdash{}{0pt}%
\pgfpathmoveto{\pgfqpoint{6.955323in}{8.351543in}}%
\pgfpathlineto{\pgfqpoint{7.752778in}{8.351543in}}%
\pgfpathquadraticcurveto{\pgfqpoint{7.780556in}{8.351543in}}{\pgfqpoint{7.780556in}{8.379321in}}%
\pgfpathlineto{\pgfqpoint{7.780556in}{8.752778in}}%
\pgfpathquadraticcurveto{\pgfqpoint{7.780556in}{8.780556in}}{\pgfqpoint{7.752778in}{8.780556in}}%
\pgfpathlineto{\pgfqpoint{6.955323in}{8.780556in}}%
\pgfpathquadraticcurveto{\pgfqpoint{6.927545in}{8.780556in}}{\pgfqpoint{6.927545in}{8.752778in}}%
\pgfpathlineto{\pgfqpoint{6.927545in}{8.379321in}}%
\pgfpathquadraticcurveto{\pgfqpoint{6.927545in}{8.351543in}}{\pgfqpoint{6.955323in}{8.351543in}}%
\pgfpathlineto{\pgfqpoint{6.955323in}{8.351543in}}%
\pgfpathclose%
\pgfusepath{stroke,fill}%
\end{pgfscope}%
\begin{pgfscope}%
\pgfsetbuttcap%
\pgfsetmiterjoin%
\definecolor{currentfill}{rgb}{0.121569,0.466667,0.705882}%
\pgfsetfillcolor{currentfill}%
\pgfsetlinewidth{0.000000pt}%
\definecolor{currentstroke}{rgb}{0.000000,0.000000,0.000000}%
\pgfsetstrokecolor{currentstroke}%
\pgfsetstrokeopacity{0.000000}%
\pgfsetdash{}{0pt}%
\pgfpathmoveto{\pgfqpoint{6.983101in}{8.627778in}}%
\pgfpathlineto{\pgfqpoint{7.260879in}{8.627778in}}%
\pgfpathlineto{\pgfqpoint{7.260879in}{8.725000in}}%
\pgfpathlineto{\pgfqpoint{6.983101in}{8.725000in}}%
\pgfpathlineto{\pgfqpoint{6.983101in}{8.627778in}}%
\pgfpathclose%
\pgfusepath{fill}%
\end{pgfscope}%
\begin{pgfscope}%
\definecolor{textcolor}{rgb}{0.000000,0.000000,0.000000}%
\pgfsetstrokecolor{textcolor}%
\pgfsetfillcolor{textcolor}%
\pgftext[x=7.371990in,y=8.627778in,left,base]{\color{textcolor}{\rmfamily\fontsize{10.000000}{12.000000}\selectfont\catcode`\^=\active\def^{\ifmmode\sp\else\^{}\fi}\catcode`\%=\active\def%{\%}RNG}}%
\end{pgfscope}%
\begin{pgfscope}%
\pgfsetrectcap%
\pgfsetroundjoin%
\pgfsetlinewidth{2.007500pt}%
\definecolor{currentstroke}{rgb}{1.000000,0.000000,0.000000}%
\pgfsetstrokecolor{currentstroke}%
\pgfsetdash{}{0pt}%
\pgfpathmoveto{\pgfqpoint{6.983101in}{8.482716in}}%
\pgfpathlineto{\pgfqpoint{7.121990in}{8.482716in}}%
\pgfpathlineto{\pgfqpoint{7.260879in}{8.482716in}}%
\pgfusepath{stroke}%
\end{pgfscope}%
\begin{pgfscope}%
\definecolor{textcolor}{rgb}{0.000000,0.000000,0.000000}%
\pgfsetstrokecolor{textcolor}%
\pgfsetfillcolor{textcolor}%
\pgftext[x=7.371990in,y=8.434105in,left,base]{\color{textcolor}{\rmfamily\fontsize{10.000000}{12.000000}\selectfont\catcode`\^=\active\def^{\ifmmode\sp\else\^{}\fi}\catcode`\%=\active\def%{\%}Exact}}%
\end{pgfscope}%
\begin{pgfscope}%
\pgfsetbuttcap%
\pgfsetmiterjoin%
\definecolor{currentfill}{rgb}{1.000000,1.000000,1.000000}%
\pgfsetfillcolor{currentfill}%
\pgfsetlinewidth{0.000000pt}%
\definecolor{currentstroke}{rgb}{0.000000,0.000000,0.000000}%
\pgfsetstrokecolor{currentstroke}%
\pgfsetstrokeopacity{0.000000}%
\pgfsetdash{}{0pt}%
\pgfpathmoveto{\pgfqpoint{0.787430in}{3.514022in}}%
\pgfpathlineto{\pgfqpoint{7.850000in}{3.514022in}}%
\pgfpathlineto{\pgfqpoint{7.850000in}{6.069395in}}%
\pgfpathlineto{\pgfqpoint{0.787430in}{6.069395in}}%
\pgfpathlineto{\pgfqpoint{0.787430in}{3.514022in}}%
\pgfpathclose%
\pgfusepath{fill}%
\end{pgfscope}%
\begin{pgfscope}%
\pgfpathrectangle{\pgfqpoint{0.787430in}{3.514022in}}{\pgfqpoint{7.062570in}{2.555373in}}%
\pgfusepath{clip}%
\pgfsetbuttcap%
\pgfsetmiterjoin%
\definecolor{currentfill}{rgb}{0.121569,0.466667,0.705882}%
\pgfsetfillcolor{currentfill}%
\pgfsetlinewidth{0.000000pt}%
\definecolor{currentstroke}{rgb}{0.000000,0.000000,0.000000}%
\pgfsetstrokecolor{currentstroke}%
\pgfsetstrokeopacity{0.000000}%
\pgfsetdash{}{0pt}%
\pgfpathmoveto{\pgfqpoint{1.558481in}{3.514022in}}%
\pgfpathlineto{\pgfqpoint{1.612475in}{3.514022in}}%
\pgfpathlineto{\pgfqpoint{1.612475in}{3.514713in}}%
\pgfpathlineto{\pgfqpoint{1.558481in}{3.514713in}}%
\pgfpathlineto{\pgfqpoint{1.558481in}{3.514022in}}%
\pgfpathclose%
\pgfusepath{fill}%
\end{pgfscope}%
\begin{pgfscope}%
\pgfpathrectangle{\pgfqpoint{0.787430in}{3.514022in}}{\pgfqpoint{7.062570in}{2.555373in}}%
\pgfusepath{clip}%
\pgfsetbuttcap%
\pgfsetmiterjoin%
\definecolor{currentfill}{rgb}{0.121569,0.466667,0.705882}%
\pgfsetfillcolor{currentfill}%
\pgfsetlinewidth{0.000000pt}%
\definecolor{currentstroke}{rgb}{0.000000,0.000000,0.000000}%
\pgfsetstrokecolor{currentstroke}%
\pgfsetstrokeopacity{0.000000}%
\pgfsetdash{}{0pt}%
\pgfpathmoveto{\pgfqpoint{1.612475in}{3.514022in}}%
\pgfpathlineto{\pgfqpoint{1.666469in}{3.514022in}}%
\pgfpathlineto{\pgfqpoint{1.666469in}{3.514022in}}%
\pgfpathlineto{\pgfqpoint{1.612475in}{3.514022in}}%
\pgfpathlineto{\pgfqpoint{1.612475in}{3.514022in}}%
\pgfpathclose%
\pgfusepath{fill}%
\end{pgfscope}%
\begin{pgfscope}%
\pgfpathrectangle{\pgfqpoint{0.787430in}{3.514022in}}{\pgfqpoint{7.062570in}{2.555373in}}%
\pgfusepath{clip}%
\pgfsetbuttcap%
\pgfsetmiterjoin%
\definecolor{currentfill}{rgb}{0.121569,0.466667,0.705882}%
\pgfsetfillcolor{currentfill}%
\pgfsetlinewidth{0.000000pt}%
\definecolor{currentstroke}{rgb}{0.000000,0.000000,0.000000}%
\pgfsetstrokecolor{currentstroke}%
\pgfsetstrokeopacity{0.000000}%
\pgfsetdash{}{0pt}%
\pgfpathmoveto{\pgfqpoint{1.666469in}{3.514022in}}%
\pgfpathlineto{\pgfqpoint{1.720464in}{3.514022in}}%
\pgfpathlineto{\pgfqpoint{1.720464in}{3.514022in}}%
\pgfpathlineto{\pgfqpoint{1.666469in}{3.514022in}}%
\pgfpathlineto{\pgfqpoint{1.666469in}{3.514022in}}%
\pgfpathclose%
\pgfusepath{fill}%
\end{pgfscope}%
\begin{pgfscope}%
\pgfpathrectangle{\pgfqpoint{0.787430in}{3.514022in}}{\pgfqpoint{7.062570in}{2.555373in}}%
\pgfusepath{clip}%
\pgfsetbuttcap%
\pgfsetmiterjoin%
\definecolor{currentfill}{rgb}{0.121569,0.466667,0.705882}%
\pgfsetfillcolor{currentfill}%
\pgfsetlinewidth{0.000000pt}%
\definecolor{currentstroke}{rgb}{0.000000,0.000000,0.000000}%
\pgfsetstrokecolor{currentstroke}%
\pgfsetstrokeopacity{0.000000}%
\pgfsetdash{}{0pt}%
\pgfpathmoveto{\pgfqpoint{1.720464in}{3.514022in}}%
\pgfpathlineto{\pgfqpoint{1.774458in}{3.514022in}}%
\pgfpathlineto{\pgfqpoint{1.774458in}{3.514022in}}%
\pgfpathlineto{\pgfqpoint{1.720464in}{3.514022in}}%
\pgfpathlineto{\pgfqpoint{1.720464in}{3.514022in}}%
\pgfpathclose%
\pgfusepath{fill}%
\end{pgfscope}%
\begin{pgfscope}%
\pgfpathrectangle{\pgfqpoint{0.787430in}{3.514022in}}{\pgfqpoint{7.062570in}{2.555373in}}%
\pgfusepath{clip}%
\pgfsetbuttcap%
\pgfsetmiterjoin%
\definecolor{currentfill}{rgb}{0.121569,0.466667,0.705882}%
\pgfsetfillcolor{currentfill}%
\pgfsetlinewidth{0.000000pt}%
\definecolor{currentstroke}{rgb}{0.000000,0.000000,0.000000}%
\pgfsetstrokecolor{currentstroke}%
\pgfsetstrokeopacity{0.000000}%
\pgfsetdash{}{0pt}%
\pgfpathmoveto{\pgfqpoint{1.774458in}{3.514022in}}%
\pgfpathlineto{\pgfqpoint{1.828452in}{3.514022in}}%
\pgfpathlineto{\pgfqpoint{1.828452in}{3.514022in}}%
\pgfpathlineto{\pgfqpoint{1.774458in}{3.514022in}}%
\pgfpathlineto{\pgfqpoint{1.774458in}{3.514022in}}%
\pgfpathclose%
\pgfusepath{fill}%
\end{pgfscope}%
\begin{pgfscope}%
\pgfpathrectangle{\pgfqpoint{0.787430in}{3.514022in}}{\pgfqpoint{7.062570in}{2.555373in}}%
\pgfusepath{clip}%
\pgfsetbuttcap%
\pgfsetmiterjoin%
\definecolor{currentfill}{rgb}{0.121569,0.466667,0.705882}%
\pgfsetfillcolor{currentfill}%
\pgfsetlinewidth{0.000000pt}%
\definecolor{currentstroke}{rgb}{0.000000,0.000000,0.000000}%
\pgfsetstrokecolor{currentstroke}%
\pgfsetstrokeopacity{0.000000}%
\pgfsetdash{}{0pt}%
\pgfpathmoveto{\pgfqpoint{1.828452in}{3.514022in}}%
\pgfpathlineto{\pgfqpoint{1.882446in}{3.514022in}}%
\pgfpathlineto{\pgfqpoint{1.882446in}{3.514022in}}%
\pgfpathlineto{\pgfqpoint{1.828452in}{3.514022in}}%
\pgfpathlineto{\pgfqpoint{1.828452in}{3.514022in}}%
\pgfpathclose%
\pgfusepath{fill}%
\end{pgfscope}%
\begin{pgfscope}%
\pgfpathrectangle{\pgfqpoint{0.787430in}{3.514022in}}{\pgfqpoint{7.062570in}{2.555373in}}%
\pgfusepath{clip}%
\pgfsetbuttcap%
\pgfsetmiterjoin%
\definecolor{currentfill}{rgb}{0.121569,0.466667,0.705882}%
\pgfsetfillcolor{currentfill}%
\pgfsetlinewidth{0.000000pt}%
\definecolor{currentstroke}{rgb}{0.000000,0.000000,0.000000}%
\pgfsetstrokecolor{currentstroke}%
\pgfsetstrokeopacity{0.000000}%
\pgfsetdash{}{0pt}%
\pgfpathmoveto{\pgfqpoint{1.882446in}{3.514022in}}%
\pgfpathlineto{\pgfqpoint{1.936440in}{3.514022in}}%
\pgfpathlineto{\pgfqpoint{1.936440in}{3.515404in}}%
\pgfpathlineto{\pgfqpoint{1.882446in}{3.515404in}}%
\pgfpathlineto{\pgfqpoint{1.882446in}{3.514022in}}%
\pgfpathclose%
\pgfusepath{fill}%
\end{pgfscope}%
\begin{pgfscope}%
\pgfpathrectangle{\pgfqpoint{0.787430in}{3.514022in}}{\pgfqpoint{7.062570in}{2.555373in}}%
\pgfusepath{clip}%
\pgfsetbuttcap%
\pgfsetmiterjoin%
\definecolor{currentfill}{rgb}{0.121569,0.466667,0.705882}%
\pgfsetfillcolor{currentfill}%
\pgfsetlinewidth{0.000000pt}%
\definecolor{currentstroke}{rgb}{0.000000,0.000000,0.000000}%
\pgfsetstrokecolor{currentstroke}%
\pgfsetstrokeopacity{0.000000}%
\pgfsetdash{}{0pt}%
\pgfpathmoveto{\pgfqpoint{1.936440in}{3.514022in}}%
\pgfpathlineto{\pgfqpoint{1.990434in}{3.514022in}}%
\pgfpathlineto{\pgfqpoint{1.990434in}{3.514022in}}%
\pgfpathlineto{\pgfqpoint{1.936440in}{3.514022in}}%
\pgfpathlineto{\pgfqpoint{1.936440in}{3.514022in}}%
\pgfpathclose%
\pgfusepath{fill}%
\end{pgfscope}%
\begin{pgfscope}%
\pgfpathrectangle{\pgfqpoint{0.787430in}{3.514022in}}{\pgfqpoint{7.062570in}{2.555373in}}%
\pgfusepath{clip}%
\pgfsetbuttcap%
\pgfsetmiterjoin%
\definecolor{currentfill}{rgb}{0.121569,0.466667,0.705882}%
\pgfsetfillcolor{currentfill}%
\pgfsetlinewidth{0.000000pt}%
\definecolor{currentstroke}{rgb}{0.000000,0.000000,0.000000}%
\pgfsetstrokecolor{currentstroke}%
\pgfsetstrokeopacity{0.000000}%
\pgfsetdash{}{0pt}%
\pgfpathmoveto{\pgfqpoint{1.990434in}{3.514022in}}%
\pgfpathlineto{\pgfqpoint{2.044428in}{3.514022in}}%
\pgfpathlineto{\pgfqpoint{2.044428in}{3.514713in}}%
\pgfpathlineto{\pgfqpoint{1.990434in}{3.514713in}}%
\pgfpathlineto{\pgfqpoint{1.990434in}{3.514022in}}%
\pgfpathclose%
\pgfusepath{fill}%
\end{pgfscope}%
\begin{pgfscope}%
\pgfpathrectangle{\pgfqpoint{0.787430in}{3.514022in}}{\pgfqpoint{7.062570in}{2.555373in}}%
\pgfusepath{clip}%
\pgfsetbuttcap%
\pgfsetmiterjoin%
\definecolor{currentfill}{rgb}{0.121569,0.466667,0.705882}%
\pgfsetfillcolor{currentfill}%
\pgfsetlinewidth{0.000000pt}%
\definecolor{currentstroke}{rgb}{0.000000,0.000000,0.000000}%
\pgfsetstrokecolor{currentstroke}%
\pgfsetstrokeopacity{0.000000}%
\pgfsetdash{}{0pt}%
\pgfpathmoveto{\pgfqpoint{2.044428in}{3.514022in}}%
\pgfpathlineto{\pgfqpoint{2.098423in}{3.514022in}}%
\pgfpathlineto{\pgfqpoint{2.098423in}{3.516094in}}%
\pgfpathlineto{\pgfqpoint{2.044428in}{3.516094in}}%
\pgfpathlineto{\pgfqpoint{2.044428in}{3.514022in}}%
\pgfpathclose%
\pgfusepath{fill}%
\end{pgfscope}%
\begin{pgfscope}%
\pgfpathrectangle{\pgfqpoint{0.787430in}{3.514022in}}{\pgfqpoint{7.062570in}{2.555373in}}%
\pgfusepath{clip}%
\pgfsetbuttcap%
\pgfsetmiterjoin%
\definecolor{currentfill}{rgb}{0.121569,0.466667,0.705882}%
\pgfsetfillcolor{currentfill}%
\pgfsetlinewidth{0.000000pt}%
\definecolor{currentstroke}{rgb}{0.000000,0.000000,0.000000}%
\pgfsetstrokecolor{currentstroke}%
\pgfsetstrokeopacity{0.000000}%
\pgfsetdash{}{0pt}%
\pgfpathmoveto{\pgfqpoint{2.098423in}{3.514022in}}%
\pgfpathlineto{\pgfqpoint{2.152417in}{3.514022in}}%
\pgfpathlineto{\pgfqpoint{2.152417in}{3.516785in}}%
\pgfpathlineto{\pgfqpoint{2.098423in}{3.516785in}}%
\pgfpathlineto{\pgfqpoint{2.098423in}{3.514022in}}%
\pgfpathclose%
\pgfusepath{fill}%
\end{pgfscope}%
\begin{pgfscope}%
\pgfpathrectangle{\pgfqpoint{0.787430in}{3.514022in}}{\pgfqpoint{7.062570in}{2.555373in}}%
\pgfusepath{clip}%
\pgfsetbuttcap%
\pgfsetmiterjoin%
\definecolor{currentfill}{rgb}{0.121569,0.466667,0.705882}%
\pgfsetfillcolor{currentfill}%
\pgfsetlinewidth{0.000000pt}%
\definecolor{currentstroke}{rgb}{0.000000,0.000000,0.000000}%
\pgfsetstrokecolor{currentstroke}%
\pgfsetstrokeopacity{0.000000}%
\pgfsetdash{}{0pt}%
\pgfpathmoveto{\pgfqpoint{2.152417in}{3.514022in}}%
\pgfpathlineto{\pgfqpoint{2.206411in}{3.514022in}}%
\pgfpathlineto{\pgfqpoint{2.206411in}{3.520930in}}%
\pgfpathlineto{\pgfqpoint{2.152417in}{3.520930in}}%
\pgfpathlineto{\pgfqpoint{2.152417in}{3.514022in}}%
\pgfpathclose%
\pgfusepath{fill}%
\end{pgfscope}%
\begin{pgfscope}%
\pgfpathrectangle{\pgfqpoint{0.787430in}{3.514022in}}{\pgfqpoint{7.062570in}{2.555373in}}%
\pgfusepath{clip}%
\pgfsetbuttcap%
\pgfsetmiterjoin%
\definecolor{currentfill}{rgb}{0.121569,0.466667,0.705882}%
\pgfsetfillcolor{currentfill}%
\pgfsetlinewidth{0.000000pt}%
\definecolor{currentstroke}{rgb}{0.000000,0.000000,0.000000}%
\pgfsetstrokecolor{currentstroke}%
\pgfsetstrokeopacity{0.000000}%
\pgfsetdash{}{0pt}%
\pgfpathmoveto{\pgfqpoint{2.206411in}{3.514022in}}%
\pgfpathlineto{\pgfqpoint{2.260405in}{3.514022in}}%
\pgfpathlineto{\pgfqpoint{2.260405in}{3.517476in}}%
\pgfpathlineto{\pgfqpoint{2.206411in}{3.517476in}}%
\pgfpathlineto{\pgfqpoint{2.206411in}{3.514022in}}%
\pgfpathclose%
\pgfusepath{fill}%
\end{pgfscope}%
\begin{pgfscope}%
\pgfpathrectangle{\pgfqpoint{0.787430in}{3.514022in}}{\pgfqpoint{7.062570in}{2.555373in}}%
\pgfusepath{clip}%
\pgfsetbuttcap%
\pgfsetmiterjoin%
\definecolor{currentfill}{rgb}{0.121569,0.466667,0.705882}%
\pgfsetfillcolor{currentfill}%
\pgfsetlinewidth{0.000000pt}%
\definecolor{currentstroke}{rgb}{0.000000,0.000000,0.000000}%
\pgfsetstrokecolor{currentstroke}%
\pgfsetstrokeopacity{0.000000}%
\pgfsetdash{}{0pt}%
\pgfpathmoveto{\pgfqpoint{2.260405in}{3.514022in}}%
\pgfpathlineto{\pgfqpoint{2.314399in}{3.514022in}}%
\pgfpathlineto{\pgfqpoint{2.314399in}{3.521621in}}%
\pgfpathlineto{\pgfqpoint{2.260405in}{3.521621in}}%
\pgfpathlineto{\pgfqpoint{2.260405in}{3.514022in}}%
\pgfpathclose%
\pgfusepath{fill}%
\end{pgfscope}%
\begin{pgfscope}%
\pgfpathrectangle{\pgfqpoint{0.787430in}{3.514022in}}{\pgfqpoint{7.062570in}{2.555373in}}%
\pgfusepath{clip}%
\pgfsetbuttcap%
\pgfsetmiterjoin%
\definecolor{currentfill}{rgb}{0.121569,0.466667,0.705882}%
\pgfsetfillcolor{currentfill}%
\pgfsetlinewidth{0.000000pt}%
\definecolor{currentstroke}{rgb}{0.000000,0.000000,0.000000}%
\pgfsetstrokecolor{currentstroke}%
\pgfsetstrokeopacity{0.000000}%
\pgfsetdash{}{0pt}%
\pgfpathmoveto{\pgfqpoint{2.314399in}{3.514022in}}%
\pgfpathlineto{\pgfqpoint{2.368393in}{3.514022in}}%
\pgfpathlineto{\pgfqpoint{2.368393in}{3.525766in}}%
\pgfpathlineto{\pgfqpoint{2.314399in}{3.525766in}}%
\pgfpathlineto{\pgfqpoint{2.314399in}{3.514022in}}%
\pgfpathclose%
\pgfusepath{fill}%
\end{pgfscope}%
\begin{pgfscope}%
\pgfpathrectangle{\pgfqpoint{0.787430in}{3.514022in}}{\pgfqpoint{7.062570in}{2.555373in}}%
\pgfusepath{clip}%
\pgfsetbuttcap%
\pgfsetmiterjoin%
\definecolor{currentfill}{rgb}{0.121569,0.466667,0.705882}%
\pgfsetfillcolor{currentfill}%
\pgfsetlinewidth{0.000000pt}%
\definecolor{currentstroke}{rgb}{0.000000,0.000000,0.000000}%
\pgfsetstrokecolor{currentstroke}%
\pgfsetstrokeopacity{0.000000}%
\pgfsetdash{}{0pt}%
\pgfpathmoveto{\pgfqpoint{2.368393in}{3.514022in}}%
\pgfpathlineto{\pgfqpoint{2.422387in}{3.514022in}}%
\pgfpathlineto{\pgfqpoint{2.422387in}{3.527147in}}%
\pgfpathlineto{\pgfqpoint{2.368393in}{3.527147in}}%
\pgfpathlineto{\pgfqpoint{2.368393in}{3.514022in}}%
\pgfpathclose%
\pgfusepath{fill}%
\end{pgfscope}%
\begin{pgfscope}%
\pgfpathrectangle{\pgfqpoint{0.787430in}{3.514022in}}{\pgfqpoint{7.062570in}{2.555373in}}%
\pgfusepath{clip}%
\pgfsetbuttcap%
\pgfsetmiterjoin%
\definecolor{currentfill}{rgb}{0.121569,0.466667,0.705882}%
\pgfsetfillcolor{currentfill}%
\pgfsetlinewidth{0.000000pt}%
\definecolor{currentstroke}{rgb}{0.000000,0.000000,0.000000}%
\pgfsetstrokecolor{currentstroke}%
\pgfsetstrokeopacity{0.000000}%
\pgfsetdash{}{0pt}%
\pgfpathmoveto{\pgfqpoint{2.422387in}{3.514022in}}%
\pgfpathlineto{\pgfqpoint{2.476382in}{3.514022in}}%
\pgfpathlineto{\pgfqpoint{2.476382in}{3.538200in}}%
\pgfpathlineto{\pgfqpoint{2.422387in}{3.538200in}}%
\pgfpathlineto{\pgfqpoint{2.422387in}{3.514022in}}%
\pgfpathclose%
\pgfusepath{fill}%
\end{pgfscope}%
\begin{pgfscope}%
\pgfpathrectangle{\pgfqpoint{0.787430in}{3.514022in}}{\pgfqpoint{7.062570in}{2.555373in}}%
\pgfusepath{clip}%
\pgfsetbuttcap%
\pgfsetmiterjoin%
\definecolor{currentfill}{rgb}{0.121569,0.466667,0.705882}%
\pgfsetfillcolor{currentfill}%
\pgfsetlinewidth{0.000000pt}%
\definecolor{currentstroke}{rgb}{0.000000,0.000000,0.000000}%
\pgfsetstrokecolor{currentstroke}%
\pgfsetstrokeopacity{0.000000}%
\pgfsetdash{}{0pt}%
\pgfpathmoveto{\pgfqpoint{2.476382in}{3.514022in}}%
\pgfpathlineto{\pgfqpoint{2.530376in}{3.514022in}}%
\pgfpathlineto{\pgfqpoint{2.530376in}{3.536128in}}%
\pgfpathlineto{\pgfqpoint{2.476382in}{3.536128in}}%
\pgfpathlineto{\pgfqpoint{2.476382in}{3.514022in}}%
\pgfpathclose%
\pgfusepath{fill}%
\end{pgfscope}%
\begin{pgfscope}%
\pgfpathrectangle{\pgfqpoint{0.787430in}{3.514022in}}{\pgfqpoint{7.062570in}{2.555373in}}%
\pgfusepath{clip}%
\pgfsetbuttcap%
\pgfsetmiterjoin%
\definecolor{currentfill}{rgb}{0.121569,0.466667,0.705882}%
\pgfsetfillcolor{currentfill}%
\pgfsetlinewidth{0.000000pt}%
\definecolor{currentstroke}{rgb}{0.000000,0.000000,0.000000}%
\pgfsetstrokecolor{currentstroke}%
\pgfsetstrokeopacity{0.000000}%
\pgfsetdash{}{0pt}%
\pgfpathmoveto{\pgfqpoint{2.530376in}{3.514022in}}%
\pgfpathlineto{\pgfqpoint{2.584370in}{3.514022in}}%
\pgfpathlineto{\pgfqpoint{2.584370in}{3.545108in}}%
\pgfpathlineto{\pgfqpoint{2.530376in}{3.545108in}}%
\pgfpathlineto{\pgfqpoint{2.530376in}{3.514022in}}%
\pgfpathclose%
\pgfusepath{fill}%
\end{pgfscope}%
\begin{pgfscope}%
\pgfpathrectangle{\pgfqpoint{0.787430in}{3.514022in}}{\pgfqpoint{7.062570in}{2.555373in}}%
\pgfusepath{clip}%
\pgfsetbuttcap%
\pgfsetmiterjoin%
\definecolor{currentfill}{rgb}{0.121569,0.466667,0.705882}%
\pgfsetfillcolor{currentfill}%
\pgfsetlinewidth{0.000000pt}%
\definecolor{currentstroke}{rgb}{0.000000,0.000000,0.000000}%
\pgfsetstrokecolor{currentstroke}%
\pgfsetstrokeopacity{0.000000}%
\pgfsetdash{}{0pt}%
\pgfpathmoveto{\pgfqpoint{2.584370in}{3.514022in}}%
\pgfpathlineto{\pgfqpoint{2.638364in}{3.514022in}}%
\pgfpathlineto{\pgfqpoint{2.638364in}{3.556852in}}%
\pgfpathlineto{\pgfqpoint{2.584370in}{3.556852in}}%
\pgfpathlineto{\pgfqpoint{2.584370in}{3.514022in}}%
\pgfpathclose%
\pgfusepath{fill}%
\end{pgfscope}%
\begin{pgfscope}%
\pgfpathrectangle{\pgfqpoint{0.787430in}{3.514022in}}{\pgfqpoint{7.062570in}{2.555373in}}%
\pgfusepath{clip}%
\pgfsetbuttcap%
\pgfsetmiterjoin%
\definecolor{currentfill}{rgb}{0.121569,0.466667,0.705882}%
\pgfsetfillcolor{currentfill}%
\pgfsetlinewidth{0.000000pt}%
\definecolor{currentstroke}{rgb}{0.000000,0.000000,0.000000}%
\pgfsetstrokecolor{currentstroke}%
\pgfsetstrokeopacity{0.000000}%
\pgfsetdash{}{0pt}%
\pgfpathmoveto{\pgfqpoint{2.638364in}{3.514022in}}%
\pgfpathlineto{\pgfqpoint{2.692358in}{3.514022in}}%
\pgfpathlineto{\pgfqpoint{2.692358in}{3.558924in}}%
\pgfpathlineto{\pgfqpoint{2.638364in}{3.558924in}}%
\pgfpathlineto{\pgfqpoint{2.638364in}{3.514022in}}%
\pgfpathclose%
\pgfusepath{fill}%
\end{pgfscope}%
\begin{pgfscope}%
\pgfpathrectangle{\pgfqpoint{0.787430in}{3.514022in}}{\pgfqpoint{7.062570in}{2.555373in}}%
\pgfusepath{clip}%
\pgfsetbuttcap%
\pgfsetmiterjoin%
\definecolor{currentfill}{rgb}{0.121569,0.466667,0.705882}%
\pgfsetfillcolor{currentfill}%
\pgfsetlinewidth{0.000000pt}%
\definecolor{currentstroke}{rgb}{0.000000,0.000000,0.000000}%
\pgfsetstrokecolor{currentstroke}%
\pgfsetstrokeopacity{0.000000}%
\pgfsetdash{}{0pt}%
\pgfpathmoveto{\pgfqpoint{2.692358in}{3.514022in}}%
\pgfpathlineto{\pgfqpoint{2.746352in}{3.514022in}}%
\pgfpathlineto{\pgfqpoint{2.746352in}{3.578266in}}%
\pgfpathlineto{\pgfqpoint{2.692358in}{3.578266in}}%
\pgfpathlineto{\pgfqpoint{2.692358in}{3.514022in}}%
\pgfpathclose%
\pgfusepath{fill}%
\end{pgfscope}%
\begin{pgfscope}%
\pgfpathrectangle{\pgfqpoint{0.787430in}{3.514022in}}{\pgfqpoint{7.062570in}{2.555373in}}%
\pgfusepath{clip}%
\pgfsetbuttcap%
\pgfsetmiterjoin%
\definecolor{currentfill}{rgb}{0.121569,0.466667,0.705882}%
\pgfsetfillcolor{currentfill}%
\pgfsetlinewidth{0.000000pt}%
\definecolor{currentstroke}{rgb}{0.000000,0.000000,0.000000}%
\pgfsetstrokecolor{currentstroke}%
\pgfsetstrokeopacity{0.000000}%
\pgfsetdash{}{0pt}%
\pgfpathmoveto{\pgfqpoint{2.746352in}{3.514022in}}%
\pgfpathlineto{\pgfqpoint{2.800346in}{3.514022in}}%
\pgfpathlineto{\pgfqpoint{2.800346in}{3.587938in}}%
\pgfpathlineto{\pgfqpoint{2.746352in}{3.587938in}}%
\pgfpathlineto{\pgfqpoint{2.746352in}{3.514022in}}%
\pgfpathclose%
\pgfusepath{fill}%
\end{pgfscope}%
\begin{pgfscope}%
\pgfpathrectangle{\pgfqpoint{0.787430in}{3.514022in}}{\pgfqpoint{7.062570in}{2.555373in}}%
\pgfusepath{clip}%
\pgfsetbuttcap%
\pgfsetmiterjoin%
\definecolor{currentfill}{rgb}{0.121569,0.466667,0.705882}%
\pgfsetfillcolor{currentfill}%
\pgfsetlinewidth{0.000000pt}%
\definecolor{currentstroke}{rgb}{0.000000,0.000000,0.000000}%
\pgfsetstrokecolor{currentstroke}%
\pgfsetstrokeopacity{0.000000}%
\pgfsetdash{}{0pt}%
\pgfpathmoveto{\pgfqpoint{2.800346in}{3.514022in}}%
\pgfpathlineto{\pgfqpoint{2.854341in}{3.514022in}}%
\pgfpathlineto{\pgfqpoint{2.854341in}{3.619714in}}%
\pgfpathlineto{\pgfqpoint{2.800346in}{3.619714in}}%
\pgfpathlineto{\pgfqpoint{2.800346in}{3.514022in}}%
\pgfpathclose%
\pgfusepath{fill}%
\end{pgfscope}%
\begin{pgfscope}%
\pgfpathrectangle{\pgfqpoint{0.787430in}{3.514022in}}{\pgfqpoint{7.062570in}{2.555373in}}%
\pgfusepath{clip}%
\pgfsetbuttcap%
\pgfsetmiterjoin%
\definecolor{currentfill}{rgb}{0.121569,0.466667,0.705882}%
\pgfsetfillcolor{currentfill}%
\pgfsetlinewidth{0.000000pt}%
\definecolor{currentstroke}{rgb}{0.000000,0.000000,0.000000}%
\pgfsetstrokecolor{currentstroke}%
\pgfsetstrokeopacity{0.000000}%
\pgfsetdash{}{0pt}%
\pgfpathmoveto{\pgfqpoint{2.854341in}{3.514022in}}%
\pgfpathlineto{\pgfqpoint{2.908335in}{3.514022in}}%
\pgfpathlineto{\pgfqpoint{2.908335in}{3.647346in}}%
\pgfpathlineto{\pgfqpoint{2.854341in}{3.647346in}}%
\pgfpathlineto{\pgfqpoint{2.854341in}{3.514022in}}%
\pgfpathclose%
\pgfusepath{fill}%
\end{pgfscope}%
\begin{pgfscope}%
\pgfpathrectangle{\pgfqpoint{0.787430in}{3.514022in}}{\pgfqpoint{7.062570in}{2.555373in}}%
\pgfusepath{clip}%
\pgfsetbuttcap%
\pgfsetmiterjoin%
\definecolor{currentfill}{rgb}{0.121569,0.466667,0.705882}%
\pgfsetfillcolor{currentfill}%
\pgfsetlinewidth{0.000000pt}%
\definecolor{currentstroke}{rgb}{0.000000,0.000000,0.000000}%
\pgfsetstrokecolor{currentstroke}%
\pgfsetstrokeopacity{0.000000}%
\pgfsetdash{}{0pt}%
\pgfpathmoveto{\pgfqpoint{2.908335in}{3.514022in}}%
\pgfpathlineto{\pgfqpoint{2.962329in}{3.514022in}}%
\pgfpathlineto{\pgfqpoint{2.962329in}{3.654945in}}%
\pgfpathlineto{\pgfqpoint{2.908335in}{3.654945in}}%
\pgfpathlineto{\pgfqpoint{2.908335in}{3.514022in}}%
\pgfpathclose%
\pgfusepath{fill}%
\end{pgfscope}%
\begin{pgfscope}%
\pgfpathrectangle{\pgfqpoint{0.787430in}{3.514022in}}{\pgfqpoint{7.062570in}{2.555373in}}%
\pgfusepath{clip}%
\pgfsetbuttcap%
\pgfsetmiterjoin%
\definecolor{currentfill}{rgb}{0.121569,0.466667,0.705882}%
\pgfsetfillcolor{currentfill}%
\pgfsetlinewidth{0.000000pt}%
\definecolor{currentstroke}{rgb}{0.000000,0.000000,0.000000}%
\pgfsetstrokecolor{currentstroke}%
\pgfsetstrokeopacity{0.000000}%
\pgfsetdash{}{0pt}%
\pgfpathmoveto{\pgfqpoint{2.962329in}{3.514022in}}%
\pgfpathlineto{\pgfqpoint{3.016323in}{3.514022in}}%
\pgfpathlineto{\pgfqpoint{3.016323in}{3.686722in}}%
\pgfpathlineto{\pgfqpoint{2.962329in}{3.686722in}}%
\pgfpathlineto{\pgfqpoint{2.962329in}{3.514022in}}%
\pgfpathclose%
\pgfusepath{fill}%
\end{pgfscope}%
\begin{pgfscope}%
\pgfpathrectangle{\pgfqpoint{0.787430in}{3.514022in}}{\pgfqpoint{7.062570in}{2.555373in}}%
\pgfusepath{clip}%
\pgfsetbuttcap%
\pgfsetmiterjoin%
\definecolor{currentfill}{rgb}{0.121569,0.466667,0.705882}%
\pgfsetfillcolor{currentfill}%
\pgfsetlinewidth{0.000000pt}%
\definecolor{currentstroke}{rgb}{0.000000,0.000000,0.000000}%
\pgfsetstrokecolor{currentstroke}%
\pgfsetstrokeopacity{0.000000}%
\pgfsetdash{}{0pt}%
\pgfpathmoveto{\pgfqpoint{3.016323in}{3.514022in}}%
\pgfpathlineto{\pgfqpoint{3.070317in}{3.514022in}}%
\pgfpathlineto{\pgfqpoint{3.070317in}{3.759256in}}%
\pgfpathlineto{\pgfqpoint{3.016323in}{3.759256in}}%
\pgfpathlineto{\pgfqpoint{3.016323in}{3.514022in}}%
\pgfpathclose%
\pgfusepath{fill}%
\end{pgfscope}%
\begin{pgfscope}%
\pgfpathrectangle{\pgfqpoint{0.787430in}{3.514022in}}{\pgfqpoint{7.062570in}{2.555373in}}%
\pgfusepath{clip}%
\pgfsetbuttcap%
\pgfsetmiterjoin%
\definecolor{currentfill}{rgb}{0.121569,0.466667,0.705882}%
\pgfsetfillcolor{currentfill}%
\pgfsetlinewidth{0.000000pt}%
\definecolor{currentstroke}{rgb}{0.000000,0.000000,0.000000}%
\pgfsetstrokecolor{currentstroke}%
\pgfsetstrokeopacity{0.000000}%
\pgfsetdash{}{0pt}%
\pgfpathmoveto{\pgfqpoint{3.070317in}{3.514022in}}%
\pgfpathlineto{\pgfqpoint{3.124311in}{3.514022in}}%
\pgfpathlineto{\pgfqpoint{3.124311in}{3.771000in}}%
\pgfpathlineto{\pgfqpoint{3.070317in}{3.771000in}}%
\pgfpathlineto{\pgfqpoint{3.070317in}{3.514022in}}%
\pgfpathclose%
\pgfusepath{fill}%
\end{pgfscope}%
\begin{pgfscope}%
\pgfpathrectangle{\pgfqpoint{0.787430in}{3.514022in}}{\pgfqpoint{7.062570in}{2.555373in}}%
\pgfusepath{clip}%
\pgfsetbuttcap%
\pgfsetmiterjoin%
\definecolor{currentfill}{rgb}{0.121569,0.466667,0.705882}%
\pgfsetfillcolor{currentfill}%
\pgfsetlinewidth{0.000000pt}%
\definecolor{currentstroke}{rgb}{0.000000,0.000000,0.000000}%
\pgfsetstrokecolor{currentstroke}%
\pgfsetstrokeopacity{0.000000}%
\pgfsetdash{}{0pt}%
\pgfpathmoveto{\pgfqpoint{3.124311in}{3.514022in}}%
\pgfpathlineto{\pgfqpoint{3.178305in}{3.514022in}}%
\pgfpathlineto{\pgfqpoint{3.178305in}{3.829718in}}%
\pgfpathlineto{\pgfqpoint{3.124311in}{3.829718in}}%
\pgfpathlineto{\pgfqpoint{3.124311in}{3.514022in}}%
\pgfpathclose%
\pgfusepath{fill}%
\end{pgfscope}%
\begin{pgfscope}%
\pgfpathrectangle{\pgfqpoint{0.787430in}{3.514022in}}{\pgfqpoint{7.062570in}{2.555373in}}%
\pgfusepath{clip}%
\pgfsetbuttcap%
\pgfsetmiterjoin%
\definecolor{currentfill}{rgb}{0.121569,0.466667,0.705882}%
\pgfsetfillcolor{currentfill}%
\pgfsetlinewidth{0.000000pt}%
\definecolor{currentstroke}{rgb}{0.000000,0.000000,0.000000}%
\pgfsetstrokecolor{currentstroke}%
\pgfsetstrokeopacity{0.000000}%
\pgfsetdash{}{0pt}%
\pgfpathmoveto{\pgfqpoint{3.178305in}{3.514022in}}%
\pgfpathlineto{\pgfqpoint{3.232300in}{3.514022in}}%
\pgfpathlineto{\pgfqpoint{3.232300in}{3.926430in}}%
\pgfpathlineto{\pgfqpoint{3.178305in}{3.926430in}}%
\pgfpathlineto{\pgfqpoint{3.178305in}{3.514022in}}%
\pgfpathclose%
\pgfusepath{fill}%
\end{pgfscope}%
\begin{pgfscope}%
\pgfpathrectangle{\pgfqpoint{0.787430in}{3.514022in}}{\pgfqpoint{7.062570in}{2.555373in}}%
\pgfusepath{clip}%
\pgfsetbuttcap%
\pgfsetmiterjoin%
\definecolor{currentfill}{rgb}{0.121569,0.466667,0.705882}%
\pgfsetfillcolor{currentfill}%
\pgfsetlinewidth{0.000000pt}%
\definecolor{currentstroke}{rgb}{0.000000,0.000000,0.000000}%
\pgfsetstrokecolor{currentstroke}%
\pgfsetstrokeopacity{0.000000}%
\pgfsetdash{}{0pt}%
\pgfpathmoveto{\pgfqpoint{3.232300in}{3.514022in}}%
\pgfpathlineto{\pgfqpoint{3.286294in}{3.514022in}}%
\pgfpathlineto{\pgfqpoint{3.286294in}{3.970641in}}%
\pgfpathlineto{\pgfqpoint{3.232300in}{3.970641in}}%
\pgfpathlineto{\pgfqpoint{3.232300in}{3.514022in}}%
\pgfpathclose%
\pgfusepath{fill}%
\end{pgfscope}%
\begin{pgfscope}%
\pgfpathrectangle{\pgfqpoint{0.787430in}{3.514022in}}{\pgfqpoint{7.062570in}{2.555373in}}%
\pgfusepath{clip}%
\pgfsetbuttcap%
\pgfsetmiterjoin%
\definecolor{currentfill}{rgb}{0.121569,0.466667,0.705882}%
\pgfsetfillcolor{currentfill}%
\pgfsetlinewidth{0.000000pt}%
\definecolor{currentstroke}{rgb}{0.000000,0.000000,0.000000}%
\pgfsetstrokecolor{currentstroke}%
\pgfsetstrokeopacity{0.000000}%
\pgfsetdash{}{0pt}%
\pgfpathmoveto{\pgfqpoint{3.286294in}{3.514022in}}%
\pgfpathlineto{\pgfqpoint{3.340288in}{3.514022in}}%
\pgfpathlineto{\pgfqpoint{3.340288in}{4.075642in}}%
\pgfpathlineto{\pgfqpoint{3.286294in}{4.075642in}}%
\pgfpathlineto{\pgfqpoint{3.286294in}{3.514022in}}%
\pgfpathclose%
\pgfusepath{fill}%
\end{pgfscope}%
\begin{pgfscope}%
\pgfpathrectangle{\pgfqpoint{0.787430in}{3.514022in}}{\pgfqpoint{7.062570in}{2.555373in}}%
\pgfusepath{clip}%
\pgfsetbuttcap%
\pgfsetmiterjoin%
\definecolor{currentfill}{rgb}{0.121569,0.466667,0.705882}%
\pgfsetfillcolor{currentfill}%
\pgfsetlinewidth{0.000000pt}%
\definecolor{currentstroke}{rgb}{0.000000,0.000000,0.000000}%
\pgfsetstrokecolor{currentstroke}%
\pgfsetstrokeopacity{0.000000}%
\pgfsetdash{}{0pt}%
\pgfpathmoveto{\pgfqpoint{3.340288in}{3.514022in}}%
\pgfpathlineto{\pgfqpoint{3.394282in}{3.514022in}}%
\pgfpathlineto{\pgfqpoint{3.394282in}{4.143341in}}%
\pgfpathlineto{\pgfqpoint{3.340288in}{4.143341in}}%
\pgfpathlineto{\pgfqpoint{3.340288in}{3.514022in}}%
\pgfpathclose%
\pgfusepath{fill}%
\end{pgfscope}%
\begin{pgfscope}%
\pgfpathrectangle{\pgfqpoint{0.787430in}{3.514022in}}{\pgfqpoint{7.062570in}{2.555373in}}%
\pgfusepath{clip}%
\pgfsetbuttcap%
\pgfsetmiterjoin%
\definecolor{currentfill}{rgb}{0.121569,0.466667,0.705882}%
\pgfsetfillcolor{currentfill}%
\pgfsetlinewidth{0.000000pt}%
\definecolor{currentstroke}{rgb}{0.000000,0.000000,0.000000}%
\pgfsetstrokecolor{currentstroke}%
\pgfsetstrokeopacity{0.000000}%
\pgfsetdash{}{0pt}%
\pgfpathmoveto{\pgfqpoint{3.394282in}{3.514022in}}%
\pgfpathlineto{\pgfqpoint{3.448276in}{3.514022in}}%
\pgfpathlineto{\pgfqpoint{3.448276in}{4.191697in}}%
\pgfpathlineto{\pgfqpoint{3.394282in}{4.191697in}}%
\pgfpathlineto{\pgfqpoint{3.394282in}{3.514022in}}%
\pgfpathclose%
\pgfusepath{fill}%
\end{pgfscope}%
\begin{pgfscope}%
\pgfpathrectangle{\pgfqpoint{0.787430in}{3.514022in}}{\pgfqpoint{7.062570in}{2.555373in}}%
\pgfusepath{clip}%
\pgfsetbuttcap%
\pgfsetmiterjoin%
\definecolor{currentfill}{rgb}{0.121569,0.466667,0.705882}%
\pgfsetfillcolor{currentfill}%
\pgfsetlinewidth{0.000000pt}%
\definecolor{currentstroke}{rgb}{0.000000,0.000000,0.000000}%
\pgfsetstrokecolor{currentstroke}%
\pgfsetstrokeopacity{0.000000}%
\pgfsetdash{}{0pt}%
\pgfpathmoveto{\pgfqpoint{3.448276in}{3.514022in}}%
\pgfpathlineto{\pgfqpoint{3.502270in}{3.514022in}}%
\pgfpathlineto{\pgfqpoint{3.502270in}{4.326403in}}%
\pgfpathlineto{\pgfqpoint{3.448276in}{4.326403in}}%
\pgfpathlineto{\pgfqpoint{3.448276in}{3.514022in}}%
\pgfpathclose%
\pgfusepath{fill}%
\end{pgfscope}%
\begin{pgfscope}%
\pgfpathrectangle{\pgfqpoint{0.787430in}{3.514022in}}{\pgfqpoint{7.062570in}{2.555373in}}%
\pgfusepath{clip}%
\pgfsetbuttcap%
\pgfsetmiterjoin%
\definecolor{currentfill}{rgb}{0.121569,0.466667,0.705882}%
\pgfsetfillcolor{currentfill}%
\pgfsetlinewidth{0.000000pt}%
\definecolor{currentstroke}{rgb}{0.000000,0.000000,0.000000}%
\pgfsetstrokecolor{currentstroke}%
\pgfsetstrokeopacity{0.000000}%
\pgfsetdash{}{0pt}%
\pgfpathmoveto{\pgfqpoint{3.502270in}{3.514022in}}%
\pgfpathlineto{\pgfqpoint{3.556264in}{3.514022in}}%
\pgfpathlineto{\pgfqpoint{3.556264in}{4.414825in}}%
\pgfpathlineto{\pgfqpoint{3.502270in}{4.414825in}}%
\pgfpathlineto{\pgfqpoint{3.502270in}{3.514022in}}%
\pgfpathclose%
\pgfusepath{fill}%
\end{pgfscope}%
\begin{pgfscope}%
\pgfpathrectangle{\pgfqpoint{0.787430in}{3.514022in}}{\pgfqpoint{7.062570in}{2.555373in}}%
\pgfusepath{clip}%
\pgfsetbuttcap%
\pgfsetmiterjoin%
\definecolor{currentfill}{rgb}{0.121569,0.466667,0.705882}%
\pgfsetfillcolor{currentfill}%
\pgfsetlinewidth{0.000000pt}%
\definecolor{currentstroke}{rgb}{0.000000,0.000000,0.000000}%
\pgfsetstrokecolor{currentstroke}%
\pgfsetstrokeopacity{0.000000}%
\pgfsetdash{}{0pt}%
\pgfpathmoveto{\pgfqpoint{3.556264in}{3.514022in}}%
\pgfpathlineto{\pgfqpoint{3.610259in}{3.514022in}}%
\pgfpathlineto{\pgfqpoint{3.610259in}{4.551604in}}%
\pgfpathlineto{\pgfqpoint{3.556264in}{4.551604in}}%
\pgfpathlineto{\pgfqpoint{3.556264in}{3.514022in}}%
\pgfpathclose%
\pgfusepath{fill}%
\end{pgfscope}%
\begin{pgfscope}%
\pgfpathrectangle{\pgfqpoint{0.787430in}{3.514022in}}{\pgfqpoint{7.062570in}{2.555373in}}%
\pgfusepath{clip}%
\pgfsetbuttcap%
\pgfsetmiterjoin%
\definecolor{currentfill}{rgb}{0.121569,0.466667,0.705882}%
\pgfsetfillcolor{currentfill}%
\pgfsetlinewidth{0.000000pt}%
\definecolor{currentstroke}{rgb}{0.000000,0.000000,0.000000}%
\pgfsetstrokecolor{currentstroke}%
\pgfsetstrokeopacity{0.000000}%
\pgfsetdash{}{0pt}%
\pgfpathmoveto{\pgfqpoint{3.610259in}{3.514022in}}%
\pgfpathlineto{\pgfqpoint{3.664253in}{3.514022in}}%
\pgfpathlineto{\pgfqpoint{3.664253in}{4.654533in}}%
\pgfpathlineto{\pgfqpoint{3.610259in}{4.654533in}}%
\pgfpathlineto{\pgfqpoint{3.610259in}{3.514022in}}%
\pgfpathclose%
\pgfusepath{fill}%
\end{pgfscope}%
\begin{pgfscope}%
\pgfpathrectangle{\pgfqpoint{0.787430in}{3.514022in}}{\pgfqpoint{7.062570in}{2.555373in}}%
\pgfusepath{clip}%
\pgfsetbuttcap%
\pgfsetmiterjoin%
\definecolor{currentfill}{rgb}{0.121569,0.466667,0.705882}%
\pgfsetfillcolor{currentfill}%
\pgfsetlinewidth{0.000000pt}%
\definecolor{currentstroke}{rgb}{0.000000,0.000000,0.000000}%
\pgfsetstrokecolor{currentstroke}%
\pgfsetstrokeopacity{0.000000}%
\pgfsetdash{}{0pt}%
\pgfpathmoveto{\pgfqpoint{3.664253in}{3.514022in}}%
\pgfpathlineto{\pgfqpoint{3.718247in}{3.514022in}}%
\pgfpathlineto{\pgfqpoint{3.718247in}{4.804437in}}%
\pgfpathlineto{\pgfqpoint{3.664253in}{4.804437in}}%
\pgfpathlineto{\pgfqpoint{3.664253in}{3.514022in}}%
\pgfpathclose%
\pgfusepath{fill}%
\end{pgfscope}%
\begin{pgfscope}%
\pgfpathrectangle{\pgfqpoint{0.787430in}{3.514022in}}{\pgfqpoint{7.062570in}{2.555373in}}%
\pgfusepath{clip}%
\pgfsetbuttcap%
\pgfsetmiterjoin%
\definecolor{currentfill}{rgb}{0.121569,0.466667,0.705882}%
\pgfsetfillcolor{currentfill}%
\pgfsetlinewidth{0.000000pt}%
\definecolor{currentstroke}{rgb}{0.000000,0.000000,0.000000}%
\pgfsetstrokecolor{currentstroke}%
\pgfsetstrokeopacity{0.000000}%
\pgfsetdash{}{0pt}%
\pgfpathmoveto{\pgfqpoint{3.718247in}{3.514022in}}%
\pgfpathlineto{\pgfqpoint{3.772241in}{3.514022in}}%
\pgfpathlineto{\pgfqpoint{3.772241in}{4.866609in}}%
\pgfpathlineto{\pgfqpoint{3.718247in}{4.866609in}}%
\pgfpathlineto{\pgfqpoint{3.718247in}{3.514022in}}%
\pgfpathclose%
\pgfusepath{fill}%
\end{pgfscope}%
\begin{pgfscope}%
\pgfpathrectangle{\pgfqpoint{0.787430in}{3.514022in}}{\pgfqpoint{7.062570in}{2.555373in}}%
\pgfusepath{clip}%
\pgfsetbuttcap%
\pgfsetmiterjoin%
\definecolor{currentfill}{rgb}{0.121569,0.466667,0.705882}%
\pgfsetfillcolor{currentfill}%
\pgfsetlinewidth{0.000000pt}%
\definecolor{currentstroke}{rgb}{0.000000,0.000000,0.000000}%
\pgfsetstrokecolor{currentstroke}%
\pgfsetstrokeopacity{0.000000}%
\pgfsetdash{}{0pt}%
\pgfpathmoveto{\pgfqpoint{3.772241in}{3.514022in}}%
\pgfpathlineto{\pgfqpoint{3.826235in}{3.514022in}}%
\pgfpathlineto{\pgfqpoint{3.826235in}{5.055888in}}%
\pgfpathlineto{\pgfqpoint{3.772241in}{5.055888in}}%
\pgfpathlineto{\pgfqpoint{3.772241in}{3.514022in}}%
\pgfpathclose%
\pgfusepath{fill}%
\end{pgfscope}%
\begin{pgfscope}%
\pgfpathrectangle{\pgfqpoint{0.787430in}{3.514022in}}{\pgfqpoint{7.062570in}{2.555373in}}%
\pgfusepath{clip}%
\pgfsetbuttcap%
\pgfsetmiterjoin%
\definecolor{currentfill}{rgb}{0.121569,0.466667,0.705882}%
\pgfsetfillcolor{currentfill}%
\pgfsetlinewidth{0.000000pt}%
\definecolor{currentstroke}{rgb}{0.000000,0.000000,0.000000}%
\pgfsetstrokecolor{currentstroke}%
\pgfsetstrokeopacity{0.000000}%
\pgfsetdash{}{0pt}%
\pgfpathmoveto{\pgfqpoint{3.826235in}{3.514022in}}%
\pgfpathlineto{\pgfqpoint{3.880229in}{3.514022in}}%
\pgfpathlineto{\pgfqpoint{3.880229in}{5.232733in}}%
\pgfpathlineto{\pgfqpoint{3.826235in}{5.232733in}}%
\pgfpathlineto{\pgfqpoint{3.826235in}{3.514022in}}%
\pgfpathclose%
\pgfusepath{fill}%
\end{pgfscope}%
\begin{pgfscope}%
\pgfpathrectangle{\pgfqpoint{0.787430in}{3.514022in}}{\pgfqpoint{7.062570in}{2.555373in}}%
\pgfusepath{clip}%
\pgfsetbuttcap%
\pgfsetmiterjoin%
\definecolor{currentfill}{rgb}{0.121569,0.466667,0.705882}%
\pgfsetfillcolor{currentfill}%
\pgfsetlinewidth{0.000000pt}%
\definecolor{currentstroke}{rgb}{0.000000,0.000000,0.000000}%
\pgfsetstrokecolor{currentstroke}%
\pgfsetstrokeopacity{0.000000}%
\pgfsetdash{}{0pt}%
\pgfpathmoveto{\pgfqpoint{3.880229in}{3.514022in}}%
\pgfpathlineto{\pgfqpoint{3.934223in}{3.514022in}}%
\pgfpathlineto{\pgfqpoint{3.934223in}{5.325300in}}%
\pgfpathlineto{\pgfqpoint{3.880229in}{5.325300in}}%
\pgfpathlineto{\pgfqpoint{3.880229in}{3.514022in}}%
\pgfpathclose%
\pgfusepath{fill}%
\end{pgfscope}%
\begin{pgfscope}%
\pgfpathrectangle{\pgfqpoint{0.787430in}{3.514022in}}{\pgfqpoint{7.062570in}{2.555373in}}%
\pgfusepath{clip}%
\pgfsetbuttcap%
\pgfsetmiterjoin%
\definecolor{currentfill}{rgb}{0.121569,0.466667,0.705882}%
\pgfsetfillcolor{currentfill}%
\pgfsetlinewidth{0.000000pt}%
\definecolor{currentstroke}{rgb}{0.000000,0.000000,0.000000}%
\pgfsetstrokecolor{currentstroke}%
\pgfsetstrokeopacity{0.000000}%
\pgfsetdash{}{0pt}%
\pgfpathmoveto{\pgfqpoint{3.934223in}{3.514022in}}%
\pgfpathlineto{\pgfqpoint{3.988218in}{3.514022in}}%
\pgfpathlineto{\pgfqpoint{3.988218in}{5.425466in}}%
\pgfpathlineto{\pgfqpoint{3.934223in}{5.425466in}}%
\pgfpathlineto{\pgfqpoint{3.934223in}{3.514022in}}%
\pgfpathclose%
\pgfusepath{fill}%
\end{pgfscope}%
\begin{pgfscope}%
\pgfpathrectangle{\pgfqpoint{0.787430in}{3.514022in}}{\pgfqpoint{7.062570in}{2.555373in}}%
\pgfusepath{clip}%
\pgfsetbuttcap%
\pgfsetmiterjoin%
\definecolor{currentfill}{rgb}{0.121569,0.466667,0.705882}%
\pgfsetfillcolor{currentfill}%
\pgfsetlinewidth{0.000000pt}%
\definecolor{currentstroke}{rgb}{0.000000,0.000000,0.000000}%
\pgfsetstrokecolor{currentstroke}%
\pgfsetstrokeopacity{0.000000}%
\pgfsetdash{}{0pt}%
\pgfpathmoveto{\pgfqpoint{3.988218in}{3.514022in}}%
\pgfpathlineto{\pgfqpoint{4.042212in}{3.514022in}}%
\pgfpathlineto{\pgfqpoint{4.042212in}{5.536685in}}%
\pgfpathlineto{\pgfqpoint{3.988218in}{5.536685in}}%
\pgfpathlineto{\pgfqpoint{3.988218in}{3.514022in}}%
\pgfpathclose%
\pgfusepath{fill}%
\end{pgfscope}%
\begin{pgfscope}%
\pgfpathrectangle{\pgfqpoint{0.787430in}{3.514022in}}{\pgfqpoint{7.062570in}{2.555373in}}%
\pgfusepath{clip}%
\pgfsetbuttcap%
\pgfsetmiterjoin%
\definecolor{currentfill}{rgb}{0.121569,0.466667,0.705882}%
\pgfsetfillcolor{currentfill}%
\pgfsetlinewidth{0.000000pt}%
\definecolor{currentstroke}{rgb}{0.000000,0.000000,0.000000}%
\pgfsetstrokecolor{currentstroke}%
\pgfsetstrokeopacity{0.000000}%
\pgfsetdash{}{0pt}%
\pgfpathmoveto{\pgfqpoint{4.042212in}{3.514022in}}%
\pgfpathlineto{\pgfqpoint{4.096206in}{3.514022in}}%
\pgfpathlineto{\pgfqpoint{4.096206in}{5.668628in}}%
\pgfpathlineto{\pgfqpoint{4.042212in}{5.668628in}}%
\pgfpathlineto{\pgfqpoint{4.042212in}{3.514022in}}%
\pgfpathclose%
\pgfusepath{fill}%
\end{pgfscope}%
\begin{pgfscope}%
\pgfpathrectangle{\pgfqpoint{0.787430in}{3.514022in}}{\pgfqpoint{7.062570in}{2.555373in}}%
\pgfusepath{clip}%
\pgfsetbuttcap%
\pgfsetmiterjoin%
\definecolor{currentfill}{rgb}{0.121569,0.466667,0.705882}%
\pgfsetfillcolor{currentfill}%
\pgfsetlinewidth{0.000000pt}%
\definecolor{currentstroke}{rgb}{0.000000,0.000000,0.000000}%
\pgfsetstrokecolor{currentstroke}%
\pgfsetstrokeopacity{0.000000}%
\pgfsetdash{}{0pt}%
\pgfpathmoveto{\pgfqpoint{4.096206in}{3.514022in}}%
\pgfpathlineto{\pgfqpoint{4.150200in}{3.514022in}}%
\pgfpathlineto{\pgfqpoint{4.150200in}{5.703858in}}%
\pgfpathlineto{\pgfqpoint{4.096206in}{5.703858in}}%
\pgfpathlineto{\pgfqpoint{4.096206in}{3.514022in}}%
\pgfpathclose%
\pgfusepath{fill}%
\end{pgfscope}%
\begin{pgfscope}%
\pgfpathrectangle{\pgfqpoint{0.787430in}{3.514022in}}{\pgfqpoint{7.062570in}{2.555373in}}%
\pgfusepath{clip}%
\pgfsetbuttcap%
\pgfsetmiterjoin%
\definecolor{currentfill}{rgb}{0.121569,0.466667,0.705882}%
\pgfsetfillcolor{currentfill}%
\pgfsetlinewidth{0.000000pt}%
\definecolor{currentstroke}{rgb}{0.000000,0.000000,0.000000}%
\pgfsetstrokecolor{currentstroke}%
\pgfsetstrokeopacity{0.000000}%
\pgfsetdash{}{0pt}%
\pgfpathmoveto{\pgfqpoint{4.150200in}{3.514022in}}%
\pgfpathlineto{\pgfqpoint{4.204194in}{3.514022in}}%
\pgfpathlineto{\pgfqpoint{4.204194in}{5.857216in}}%
\pgfpathlineto{\pgfqpoint{4.150200in}{5.857216in}}%
\pgfpathlineto{\pgfqpoint{4.150200in}{3.514022in}}%
\pgfpathclose%
\pgfusepath{fill}%
\end{pgfscope}%
\begin{pgfscope}%
\pgfpathrectangle{\pgfqpoint{0.787430in}{3.514022in}}{\pgfqpoint{7.062570in}{2.555373in}}%
\pgfusepath{clip}%
\pgfsetbuttcap%
\pgfsetmiterjoin%
\definecolor{currentfill}{rgb}{0.121569,0.466667,0.705882}%
\pgfsetfillcolor{currentfill}%
\pgfsetlinewidth{0.000000pt}%
\definecolor{currentstroke}{rgb}{0.000000,0.000000,0.000000}%
\pgfsetstrokecolor{currentstroke}%
\pgfsetstrokeopacity{0.000000}%
\pgfsetdash{}{0pt}%
\pgfpathmoveto{\pgfqpoint{4.204194in}{3.514022in}}%
\pgfpathlineto{\pgfqpoint{4.258188in}{3.514022in}}%
\pgfpathlineto{\pgfqpoint{4.258188in}{5.837183in}}%
\pgfpathlineto{\pgfqpoint{4.204194in}{5.837183in}}%
\pgfpathlineto{\pgfqpoint{4.204194in}{3.514022in}}%
\pgfpathclose%
\pgfusepath{fill}%
\end{pgfscope}%
\begin{pgfscope}%
\pgfpathrectangle{\pgfqpoint{0.787430in}{3.514022in}}{\pgfqpoint{7.062570in}{2.555373in}}%
\pgfusepath{clip}%
\pgfsetbuttcap%
\pgfsetmiterjoin%
\definecolor{currentfill}{rgb}{0.121569,0.466667,0.705882}%
\pgfsetfillcolor{currentfill}%
\pgfsetlinewidth{0.000000pt}%
\definecolor{currentstroke}{rgb}{0.000000,0.000000,0.000000}%
\pgfsetstrokecolor{currentstroke}%
\pgfsetstrokeopacity{0.000000}%
\pgfsetdash{}{0pt}%
\pgfpathmoveto{\pgfqpoint{4.258188in}{3.514022in}}%
\pgfpathlineto{\pgfqpoint{4.312182in}{3.514022in}}%
\pgfpathlineto{\pgfqpoint{4.312182in}{5.915934in}}%
\pgfpathlineto{\pgfqpoint{4.258188in}{5.915934in}}%
\pgfpathlineto{\pgfqpoint{4.258188in}{3.514022in}}%
\pgfpathclose%
\pgfusepath{fill}%
\end{pgfscope}%
\begin{pgfscope}%
\pgfpathrectangle{\pgfqpoint{0.787430in}{3.514022in}}{\pgfqpoint{7.062570in}{2.555373in}}%
\pgfusepath{clip}%
\pgfsetbuttcap%
\pgfsetmiterjoin%
\definecolor{currentfill}{rgb}{0.121569,0.466667,0.705882}%
\pgfsetfillcolor{currentfill}%
\pgfsetlinewidth{0.000000pt}%
\definecolor{currentstroke}{rgb}{0.000000,0.000000,0.000000}%
\pgfsetstrokecolor{currentstroke}%
\pgfsetstrokeopacity{0.000000}%
\pgfsetdash{}{0pt}%
\pgfpathmoveto{\pgfqpoint{4.312182in}{3.514022in}}%
\pgfpathlineto{\pgfqpoint{4.366177in}{3.514022in}}%
\pgfpathlineto{\pgfqpoint{4.366177in}{5.938730in}}%
\pgfpathlineto{\pgfqpoint{4.312182in}{5.938730in}}%
\pgfpathlineto{\pgfqpoint{4.312182in}{3.514022in}}%
\pgfpathclose%
\pgfusepath{fill}%
\end{pgfscope}%
\begin{pgfscope}%
\pgfpathrectangle{\pgfqpoint{0.787430in}{3.514022in}}{\pgfqpoint{7.062570in}{2.555373in}}%
\pgfusepath{clip}%
\pgfsetbuttcap%
\pgfsetmiterjoin%
\definecolor{currentfill}{rgb}{0.121569,0.466667,0.705882}%
\pgfsetfillcolor{currentfill}%
\pgfsetlinewidth{0.000000pt}%
\definecolor{currentstroke}{rgb}{0.000000,0.000000,0.000000}%
\pgfsetstrokecolor{currentstroke}%
\pgfsetstrokeopacity{0.000000}%
\pgfsetdash{}{0pt}%
\pgfpathmoveto{\pgfqpoint{4.366177in}{3.514022in}}%
\pgfpathlineto{\pgfqpoint{4.420171in}{3.514022in}}%
\pgfpathlineto{\pgfqpoint{4.420171in}{5.947711in}}%
\pgfpathlineto{\pgfqpoint{4.366177in}{5.947711in}}%
\pgfpathlineto{\pgfqpoint{4.366177in}{3.514022in}}%
\pgfpathclose%
\pgfusepath{fill}%
\end{pgfscope}%
\begin{pgfscope}%
\pgfpathrectangle{\pgfqpoint{0.787430in}{3.514022in}}{\pgfqpoint{7.062570in}{2.555373in}}%
\pgfusepath{clip}%
\pgfsetbuttcap%
\pgfsetmiterjoin%
\definecolor{currentfill}{rgb}{0.121569,0.466667,0.705882}%
\pgfsetfillcolor{currentfill}%
\pgfsetlinewidth{0.000000pt}%
\definecolor{currentstroke}{rgb}{0.000000,0.000000,0.000000}%
\pgfsetstrokecolor{currentstroke}%
\pgfsetstrokeopacity{0.000000}%
\pgfsetdash{}{0pt}%
\pgfpathmoveto{\pgfqpoint{4.420171in}{3.514022in}}%
\pgfpathlineto{\pgfqpoint{4.474165in}{3.514022in}}%
\pgfpathlineto{\pgfqpoint{4.474165in}{5.902809in}}%
\pgfpathlineto{\pgfqpoint{4.420171in}{5.902809in}}%
\pgfpathlineto{\pgfqpoint{4.420171in}{3.514022in}}%
\pgfpathclose%
\pgfusepath{fill}%
\end{pgfscope}%
\begin{pgfscope}%
\pgfpathrectangle{\pgfqpoint{0.787430in}{3.514022in}}{\pgfqpoint{7.062570in}{2.555373in}}%
\pgfusepath{clip}%
\pgfsetbuttcap%
\pgfsetmiterjoin%
\definecolor{currentfill}{rgb}{0.121569,0.466667,0.705882}%
\pgfsetfillcolor{currentfill}%
\pgfsetlinewidth{0.000000pt}%
\definecolor{currentstroke}{rgb}{0.000000,0.000000,0.000000}%
\pgfsetstrokecolor{currentstroke}%
\pgfsetstrokeopacity{0.000000}%
\pgfsetdash{}{0pt}%
\pgfpathmoveto{\pgfqpoint{4.474165in}{3.514022in}}%
\pgfpathlineto{\pgfqpoint{4.528159in}{3.514022in}}%
\pgfpathlineto{\pgfqpoint{4.528159in}{5.924914in}}%
\pgfpathlineto{\pgfqpoint{4.474165in}{5.924914in}}%
\pgfpathlineto{\pgfqpoint{4.474165in}{3.514022in}}%
\pgfpathclose%
\pgfusepath{fill}%
\end{pgfscope}%
\begin{pgfscope}%
\pgfpathrectangle{\pgfqpoint{0.787430in}{3.514022in}}{\pgfqpoint{7.062570in}{2.555373in}}%
\pgfusepath{clip}%
\pgfsetbuttcap%
\pgfsetmiterjoin%
\definecolor{currentfill}{rgb}{0.121569,0.466667,0.705882}%
\pgfsetfillcolor{currentfill}%
\pgfsetlinewidth{0.000000pt}%
\definecolor{currentstroke}{rgb}{0.000000,0.000000,0.000000}%
\pgfsetstrokecolor{currentstroke}%
\pgfsetstrokeopacity{0.000000}%
\pgfsetdash{}{0pt}%
\pgfpathmoveto{\pgfqpoint{4.528159in}{3.514022in}}%
\pgfpathlineto{\pgfqpoint{4.582153in}{3.514022in}}%
\pgfpathlineto{\pgfqpoint{4.582153in}{5.861361in}}%
\pgfpathlineto{\pgfqpoint{4.528159in}{5.861361in}}%
\pgfpathlineto{\pgfqpoint{4.528159in}{3.514022in}}%
\pgfpathclose%
\pgfusepath{fill}%
\end{pgfscope}%
\begin{pgfscope}%
\pgfpathrectangle{\pgfqpoint{0.787430in}{3.514022in}}{\pgfqpoint{7.062570in}{2.555373in}}%
\pgfusepath{clip}%
\pgfsetbuttcap%
\pgfsetmiterjoin%
\definecolor{currentfill}{rgb}{0.121569,0.466667,0.705882}%
\pgfsetfillcolor{currentfill}%
\pgfsetlinewidth{0.000000pt}%
\definecolor{currentstroke}{rgb}{0.000000,0.000000,0.000000}%
\pgfsetstrokecolor{currentstroke}%
\pgfsetstrokeopacity{0.000000}%
\pgfsetdash{}{0pt}%
\pgfpathmoveto{\pgfqpoint{4.582153in}{3.514022in}}%
\pgfpathlineto{\pgfqpoint{4.636147in}{3.514022in}}%
\pgfpathlineto{\pgfqpoint{4.636147in}{5.781919in}}%
\pgfpathlineto{\pgfqpoint{4.582153in}{5.781919in}}%
\pgfpathlineto{\pgfqpoint{4.582153in}{3.514022in}}%
\pgfpathclose%
\pgfusepath{fill}%
\end{pgfscope}%
\begin{pgfscope}%
\pgfpathrectangle{\pgfqpoint{0.787430in}{3.514022in}}{\pgfqpoint{7.062570in}{2.555373in}}%
\pgfusepath{clip}%
\pgfsetbuttcap%
\pgfsetmiterjoin%
\definecolor{currentfill}{rgb}{0.121569,0.466667,0.705882}%
\pgfsetfillcolor{currentfill}%
\pgfsetlinewidth{0.000000pt}%
\definecolor{currentstroke}{rgb}{0.000000,0.000000,0.000000}%
\pgfsetstrokecolor{currentstroke}%
\pgfsetstrokeopacity{0.000000}%
\pgfsetdash{}{0pt}%
\pgfpathmoveto{\pgfqpoint{4.636147in}{3.514022in}}%
\pgfpathlineto{\pgfqpoint{4.690141in}{3.514022in}}%
\pgfpathlineto{\pgfqpoint{4.690141in}{5.634088in}}%
\pgfpathlineto{\pgfqpoint{4.636147in}{5.634088in}}%
\pgfpathlineto{\pgfqpoint{4.636147in}{3.514022in}}%
\pgfpathclose%
\pgfusepath{fill}%
\end{pgfscope}%
\begin{pgfscope}%
\pgfpathrectangle{\pgfqpoint{0.787430in}{3.514022in}}{\pgfqpoint{7.062570in}{2.555373in}}%
\pgfusepath{clip}%
\pgfsetbuttcap%
\pgfsetmiterjoin%
\definecolor{currentfill}{rgb}{0.121569,0.466667,0.705882}%
\pgfsetfillcolor{currentfill}%
\pgfsetlinewidth{0.000000pt}%
\definecolor{currentstroke}{rgb}{0.000000,0.000000,0.000000}%
\pgfsetstrokecolor{currentstroke}%
\pgfsetstrokeopacity{0.000000}%
\pgfsetdash{}{0pt}%
\pgfpathmoveto{\pgfqpoint{4.690141in}{3.514022in}}%
\pgfpathlineto{\pgfqpoint{4.744136in}{3.514022in}}%
\pgfpathlineto{\pgfqpoint{4.744136in}{5.601620in}}%
\pgfpathlineto{\pgfqpoint{4.690141in}{5.601620in}}%
\pgfpathlineto{\pgfqpoint{4.690141in}{3.514022in}}%
\pgfpathclose%
\pgfusepath{fill}%
\end{pgfscope}%
\begin{pgfscope}%
\pgfpathrectangle{\pgfqpoint{0.787430in}{3.514022in}}{\pgfqpoint{7.062570in}{2.555373in}}%
\pgfusepath{clip}%
\pgfsetbuttcap%
\pgfsetmiterjoin%
\definecolor{currentfill}{rgb}{0.121569,0.466667,0.705882}%
\pgfsetfillcolor{currentfill}%
\pgfsetlinewidth{0.000000pt}%
\definecolor{currentstroke}{rgb}{0.000000,0.000000,0.000000}%
\pgfsetstrokecolor{currentstroke}%
\pgfsetstrokeopacity{0.000000}%
\pgfsetdash{}{0pt}%
\pgfpathmoveto{\pgfqpoint{4.744136in}{3.514022in}}%
\pgfpathlineto{\pgfqpoint{4.798130in}{3.514022in}}%
\pgfpathlineto{\pgfqpoint{4.798130in}{5.513888in}}%
\pgfpathlineto{\pgfqpoint{4.744136in}{5.513888in}}%
\pgfpathlineto{\pgfqpoint{4.744136in}{3.514022in}}%
\pgfpathclose%
\pgfusepath{fill}%
\end{pgfscope}%
\begin{pgfscope}%
\pgfpathrectangle{\pgfqpoint{0.787430in}{3.514022in}}{\pgfqpoint{7.062570in}{2.555373in}}%
\pgfusepath{clip}%
\pgfsetbuttcap%
\pgfsetmiterjoin%
\definecolor{currentfill}{rgb}{0.121569,0.466667,0.705882}%
\pgfsetfillcolor{currentfill}%
\pgfsetlinewidth{0.000000pt}%
\definecolor{currentstroke}{rgb}{0.000000,0.000000,0.000000}%
\pgfsetstrokecolor{currentstroke}%
\pgfsetstrokeopacity{0.000000}%
\pgfsetdash{}{0pt}%
\pgfpathmoveto{\pgfqpoint{4.798130in}{3.514022in}}%
\pgfpathlineto{\pgfqpoint{4.852124in}{3.514022in}}%
\pgfpathlineto{\pgfqpoint{4.852124in}{5.346715in}}%
\pgfpathlineto{\pgfqpoint{4.798130in}{5.346715in}}%
\pgfpathlineto{\pgfqpoint{4.798130in}{3.514022in}}%
\pgfpathclose%
\pgfusepath{fill}%
\end{pgfscope}%
\begin{pgfscope}%
\pgfpathrectangle{\pgfqpoint{0.787430in}{3.514022in}}{\pgfqpoint{7.062570in}{2.555373in}}%
\pgfusepath{clip}%
\pgfsetbuttcap%
\pgfsetmiterjoin%
\definecolor{currentfill}{rgb}{0.121569,0.466667,0.705882}%
\pgfsetfillcolor{currentfill}%
\pgfsetlinewidth{0.000000pt}%
\definecolor{currentstroke}{rgb}{0.000000,0.000000,0.000000}%
\pgfsetstrokecolor{currentstroke}%
\pgfsetstrokeopacity{0.000000}%
\pgfsetdash{}{0pt}%
\pgfpathmoveto{\pgfqpoint{4.852124in}{3.514022in}}%
\pgfpathlineto{\pgfqpoint{4.906118in}{3.514022in}}%
\pgfpathlineto{\pgfqpoint{4.906118in}{5.300431in}}%
\pgfpathlineto{\pgfqpoint{4.852124in}{5.300431in}}%
\pgfpathlineto{\pgfqpoint{4.852124in}{3.514022in}}%
\pgfpathclose%
\pgfusepath{fill}%
\end{pgfscope}%
\begin{pgfscope}%
\pgfpathrectangle{\pgfqpoint{0.787430in}{3.514022in}}{\pgfqpoint{7.062570in}{2.555373in}}%
\pgfusepath{clip}%
\pgfsetbuttcap%
\pgfsetmiterjoin%
\definecolor{currentfill}{rgb}{0.121569,0.466667,0.705882}%
\pgfsetfillcolor{currentfill}%
\pgfsetlinewidth{0.000000pt}%
\definecolor{currentstroke}{rgb}{0.000000,0.000000,0.000000}%
\pgfsetstrokecolor{currentstroke}%
\pgfsetstrokeopacity{0.000000}%
\pgfsetdash{}{0pt}%
\pgfpathmoveto{\pgfqpoint{4.906118in}{3.514022in}}%
\pgfpathlineto{\pgfqpoint{4.960112in}{3.514022in}}%
\pgfpathlineto{\pgfqpoint{4.960112in}{5.077303in}}%
\pgfpathlineto{\pgfqpoint{4.906118in}{5.077303in}}%
\pgfpathlineto{\pgfqpoint{4.906118in}{3.514022in}}%
\pgfpathclose%
\pgfusepath{fill}%
\end{pgfscope}%
\begin{pgfscope}%
\pgfpathrectangle{\pgfqpoint{0.787430in}{3.514022in}}{\pgfqpoint{7.062570in}{2.555373in}}%
\pgfusepath{clip}%
\pgfsetbuttcap%
\pgfsetmiterjoin%
\definecolor{currentfill}{rgb}{0.121569,0.466667,0.705882}%
\pgfsetfillcolor{currentfill}%
\pgfsetlinewidth{0.000000pt}%
\definecolor{currentstroke}{rgb}{0.000000,0.000000,0.000000}%
\pgfsetstrokecolor{currentstroke}%
\pgfsetstrokeopacity{0.000000}%
\pgfsetdash{}{0pt}%
\pgfpathmoveto{\pgfqpoint{4.960112in}{3.514022in}}%
\pgfpathlineto{\pgfqpoint{5.014106in}{3.514022in}}%
\pgfpathlineto{\pgfqpoint{5.014106in}{5.006841in}}%
\pgfpathlineto{\pgfqpoint{4.960112in}{5.006841in}}%
\pgfpathlineto{\pgfqpoint{4.960112in}{3.514022in}}%
\pgfpathclose%
\pgfusepath{fill}%
\end{pgfscope}%
\begin{pgfscope}%
\pgfpathrectangle{\pgfqpoint{0.787430in}{3.514022in}}{\pgfqpoint{7.062570in}{2.555373in}}%
\pgfusepath{clip}%
\pgfsetbuttcap%
\pgfsetmiterjoin%
\definecolor{currentfill}{rgb}{0.121569,0.466667,0.705882}%
\pgfsetfillcolor{currentfill}%
\pgfsetlinewidth{0.000000pt}%
\definecolor{currentstroke}{rgb}{0.000000,0.000000,0.000000}%
\pgfsetstrokecolor{currentstroke}%
\pgfsetstrokeopacity{0.000000}%
\pgfsetdash{}{0pt}%
\pgfpathmoveto{\pgfqpoint{5.014106in}{3.514022in}}%
\pgfpathlineto{\pgfqpoint{5.068100in}{3.514022in}}%
\pgfpathlineto{\pgfqpoint{5.068100in}{4.875589in}}%
\pgfpathlineto{\pgfqpoint{5.014106in}{4.875589in}}%
\pgfpathlineto{\pgfqpoint{5.014106in}{3.514022in}}%
\pgfpathclose%
\pgfusepath{fill}%
\end{pgfscope}%
\begin{pgfscope}%
\pgfpathrectangle{\pgfqpoint{0.787430in}{3.514022in}}{\pgfqpoint{7.062570in}{2.555373in}}%
\pgfusepath{clip}%
\pgfsetbuttcap%
\pgfsetmiterjoin%
\definecolor{currentfill}{rgb}{0.121569,0.466667,0.705882}%
\pgfsetfillcolor{currentfill}%
\pgfsetlinewidth{0.000000pt}%
\definecolor{currentstroke}{rgb}{0.000000,0.000000,0.000000}%
\pgfsetstrokecolor{currentstroke}%
\pgfsetstrokeopacity{0.000000}%
\pgfsetdash{}{0pt}%
\pgfpathmoveto{\pgfqpoint{5.068100in}{3.514022in}}%
\pgfpathlineto{\pgfqpoint{5.122095in}{3.514022in}}%
\pgfpathlineto{\pgfqpoint{5.122095in}{4.736047in}}%
\pgfpathlineto{\pgfqpoint{5.068100in}{4.736047in}}%
\pgfpathlineto{\pgfqpoint{5.068100in}{3.514022in}}%
\pgfpathclose%
\pgfusepath{fill}%
\end{pgfscope}%
\begin{pgfscope}%
\pgfpathrectangle{\pgfqpoint{0.787430in}{3.514022in}}{\pgfqpoint{7.062570in}{2.555373in}}%
\pgfusepath{clip}%
\pgfsetbuttcap%
\pgfsetmiterjoin%
\definecolor{currentfill}{rgb}{0.121569,0.466667,0.705882}%
\pgfsetfillcolor{currentfill}%
\pgfsetlinewidth{0.000000pt}%
\definecolor{currentstroke}{rgb}{0.000000,0.000000,0.000000}%
\pgfsetstrokecolor{currentstroke}%
\pgfsetstrokeopacity{0.000000}%
\pgfsetdash{}{0pt}%
\pgfpathmoveto{\pgfqpoint{5.122095in}{3.514022in}}%
\pgfpathlineto{\pgfqpoint{5.176089in}{3.514022in}}%
\pgfpathlineto{\pgfqpoint{5.176089in}{4.662823in}}%
\pgfpathlineto{\pgfqpoint{5.122095in}{4.662823in}}%
\pgfpathlineto{\pgfqpoint{5.122095in}{3.514022in}}%
\pgfpathclose%
\pgfusepath{fill}%
\end{pgfscope}%
\begin{pgfscope}%
\pgfpathrectangle{\pgfqpoint{0.787430in}{3.514022in}}{\pgfqpoint{7.062570in}{2.555373in}}%
\pgfusepath{clip}%
\pgfsetbuttcap%
\pgfsetmiterjoin%
\definecolor{currentfill}{rgb}{0.121569,0.466667,0.705882}%
\pgfsetfillcolor{currentfill}%
\pgfsetlinewidth{0.000000pt}%
\definecolor{currentstroke}{rgb}{0.000000,0.000000,0.000000}%
\pgfsetstrokecolor{currentstroke}%
\pgfsetstrokeopacity{0.000000}%
\pgfsetdash{}{0pt}%
\pgfpathmoveto{\pgfqpoint{5.176089in}{3.514022in}}%
\pgfpathlineto{\pgfqpoint{5.230083in}{3.514022in}}%
\pgfpathlineto{\pgfqpoint{5.230083in}{4.510156in}}%
\pgfpathlineto{\pgfqpoint{5.176089in}{4.510156in}}%
\pgfpathlineto{\pgfqpoint{5.176089in}{3.514022in}}%
\pgfpathclose%
\pgfusepath{fill}%
\end{pgfscope}%
\begin{pgfscope}%
\pgfpathrectangle{\pgfqpoint{0.787430in}{3.514022in}}{\pgfqpoint{7.062570in}{2.555373in}}%
\pgfusepath{clip}%
\pgfsetbuttcap%
\pgfsetmiterjoin%
\definecolor{currentfill}{rgb}{0.121569,0.466667,0.705882}%
\pgfsetfillcolor{currentfill}%
\pgfsetlinewidth{0.000000pt}%
\definecolor{currentstroke}{rgb}{0.000000,0.000000,0.000000}%
\pgfsetstrokecolor{currentstroke}%
\pgfsetstrokeopacity{0.000000}%
\pgfsetdash{}{0pt}%
\pgfpathmoveto{\pgfqpoint{5.230083in}{3.514022in}}%
\pgfpathlineto{\pgfqpoint{5.284077in}{3.514022in}}%
\pgfpathlineto{\pgfqpoint{5.284077in}{4.399628in}}%
\pgfpathlineto{\pgfqpoint{5.230083in}{4.399628in}}%
\pgfpathlineto{\pgfqpoint{5.230083in}{3.514022in}}%
\pgfpathclose%
\pgfusepath{fill}%
\end{pgfscope}%
\begin{pgfscope}%
\pgfpathrectangle{\pgfqpoint{0.787430in}{3.514022in}}{\pgfqpoint{7.062570in}{2.555373in}}%
\pgfusepath{clip}%
\pgfsetbuttcap%
\pgfsetmiterjoin%
\definecolor{currentfill}{rgb}{0.121569,0.466667,0.705882}%
\pgfsetfillcolor{currentfill}%
\pgfsetlinewidth{0.000000pt}%
\definecolor{currentstroke}{rgb}{0.000000,0.000000,0.000000}%
\pgfsetstrokecolor{currentstroke}%
\pgfsetstrokeopacity{0.000000}%
\pgfsetdash{}{0pt}%
\pgfpathmoveto{\pgfqpoint{5.284077in}{3.514022in}}%
\pgfpathlineto{\pgfqpoint{5.338071in}{3.514022in}}%
\pgfpathlineto{\pgfqpoint{5.338071in}{4.260086in}}%
\pgfpathlineto{\pgfqpoint{5.284077in}{4.260086in}}%
\pgfpathlineto{\pgfqpoint{5.284077in}{3.514022in}}%
\pgfpathclose%
\pgfusepath{fill}%
\end{pgfscope}%
\begin{pgfscope}%
\pgfpathrectangle{\pgfqpoint{0.787430in}{3.514022in}}{\pgfqpoint{7.062570in}{2.555373in}}%
\pgfusepath{clip}%
\pgfsetbuttcap%
\pgfsetmiterjoin%
\definecolor{currentfill}{rgb}{0.121569,0.466667,0.705882}%
\pgfsetfillcolor{currentfill}%
\pgfsetlinewidth{0.000000pt}%
\definecolor{currentstroke}{rgb}{0.000000,0.000000,0.000000}%
\pgfsetstrokecolor{currentstroke}%
\pgfsetstrokeopacity{0.000000}%
\pgfsetdash{}{0pt}%
\pgfpathmoveto{\pgfqpoint{5.338071in}{3.514022in}}%
\pgfpathlineto{\pgfqpoint{5.392065in}{3.514022in}}%
\pgfpathlineto{\pgfqpoint{5.392065in}{4.191006in}}%
\pgfpathlineto{\pgfqpoint{5.338071in}{4.191006in}}%
\pgfpathlineto{\pgfqpoint{5.338071in}{3.514022in}}%
\pgfpathclose%
\pgfusepath{fill}%
\end{pgfscope}%
\begin{pgfscope}%
\pgfpathrectangle{\pgfqpoint{0.787430in}{3.514022in}}{\pgfqpoint{7.062570in}{2.555373in}}%
\pgfusepath{clip}%
\pgfsetbuttcap%
\pgfsetmiterjoin%
\definecolor{currentfill}{rgb}{0.121569,0.466667,0.705882}%
\pgfsetfillcolor{currentfill}%
\pgfsetlinewidth{0.000000pt}%
\definecolor{currentstroke}{rgb}{0.000000,0.000000,0.000000}%
\pgfsetstrokecolor{currentstroke}%
\pgfsetstrokeopacity{0.000000}%
\pgfsetdash{}{0pt}%
\pgfpathmoveto{\pgfqpoint{5.392065in}{3.514022in}}%
\pgfpathlineto{\pgfqpoint{5.446059in}{3.514022in}}%
\pgfpathlineto{\pgfqpoint{5.446059in}{4.089458in}}%
\pgfpathlineto{\pgfqpoint{5.392065in}{4.089458in}}%
\pgfpathlineto{\pgfqpoint{5.392065in}{3.514022in}}%
\pgfpathclose%
\pgfusepath{fill}%
\end{pgfscope}%
\begin{pgfscope}%
\pgfpathrectangle{\pgfqpoint{0.787430in}{3.514022in}}{\pgfqpoint{7.062570in}{2.555373in}}%
\pgfusepath{clip}%
\pgfsetbuttcap%
\pgfsetmiterjoin%
\definecolor{currentfill}{rgb}{0.121569,0.466667,0.705882}%
\pgfsetfillcolor{currentfill}%
\pgfsetlinewidth{0.000000pt}%
\definecolor{currentstroke}{rgb}{0.000000,0.000000,0.000000}%
\pgfsetstrokecolor{currentstroke}%
\pgfsetstrokeopacity{0.000000}%
\pgfsetdash{}{0pt}%
\pgfpathmoveto{\pgfqpoint{5.446059in}{3.514022in}}%
\pgfpathlineto{\pgfqpoint{5.500054in}{3.514022in}}%
\pgfpathlineto{\pgfqpoint{5.500054in}{4.017615in}}%
\pgfpathlineto{\pgfqpoint{5.446059in}{4.017615in}}%
\pgfpathlineto{\pgfqpoint{5.446059in}{3.514022in}}%
\pgfpathclose%
\pgfusepath{fill}%
\end{pgfscope}%
\begin{pgfscope}%
\pgfpathrectangle{\pgfqpoint{0.787430in}{3.514022in}}{\pgfqpoint{7.062570in}{2.555373in}}%
\pgfusepath{clip}%
\pgfsetbuttcap%
\pgfsetmiterjoin%
\definecolor{currentfill}{rgb}{0.121569,0.466667,0.705882}%
\pgfsetfillcolor{currentfill}%
\pgfsetlinewidth{0.000000pt}%
\definecolor{currentstroke}{rgb}{0.000000,0.000000,0.000000}%
\pgfsetstrokecolor{currentstroke}%
\pgfsetstrokeopacity{0.000000}%
\pgfsetdash{}{0pt}%
\pgfpathmoveto{\pgfqpoint{5.500054in}{3.514022in}}%
\pgfpathlineto{\pgfqpoint{5.554048in}{3.514022in}}%
\pgfpathlineto{\pgfqpoint{5.554048in}{3.942318in}}%
\pgfpathlineto{\pgfqpoint{5.500054in}{3.942318in}}%
\pgfpathlineto{\pgfqpoint{5.500054in}{3.514022in}}%
\pgfpathclose%
\pgfusepath{fill}%
\end{pgfscope}%
\begin{pgfscope}%
\pgfpathrectangle{\pgfqpoint{0.787430in}{3.514022in}}{\pgfqpoint{7.062570in}{2.555373in}}%
\pgfusepath{clip}%
\pgfsetbuttcap%
\pgfsetmiterjoin%
\definecolor{currentfill}{rgb}{0.121569,0.466667,0.705882}%
\pgfsetfillcolor{currentfill}%
\pgfsetlinewidth{0.000000pt}%
\definecolor{currentstroke}{rgb}{0.000000,0.000000,0.000000}%
\pgfsetstrokecolor{currentstroke}%
\pgfsetstrokeopacity{0.000000}%
\pgfsetdash{}{0pt}%
\pgfpathmoveto{\pgfqpoint{5.554048in}{3.514022in}}%
\pgfpathlineto{\pgfqpoint{5.608042in}{3.514022in}}%
\pgfpathlineto{\pgfqpoint{5.608042in}{3.876692in}}%
\pgfpathlineto{\pgfqpoint{5.554048in}{3.876692in}}%
\pgfpathlineto{\pgfqpoint{5.554048in}{3.514022in}}%
\pgfpathclose%
\pgfusepath{fill}%
\end{pgfscope}%
\begin{pgfscope}%
\pgfpathrectangle{\pgfqpoint{0.787430in}{3.514022in}}{\pgfqpoint{7.062570in}{2.555373in}}%
\pgfusepath{clip}%
\pgfsetbuttcap%
\pgfsetmiterjoin%
\definecolor{currentfill}{rgb}{0.121569,0.466667,0.705882}%
\pgfsetfillcolor{currentfill}%
\pgfsetlinewidth{0.000000pt}%
\definecolor{currentstroke}{rgb}{0.000000,0.000000,0.000000}%
\pgfsetstrokecolor{currentstroke}%
\pgfsetstrokeopacity{0.000000}%
\pgfsetdash{}{0pt}%
\pgfpathmoveto{\pgfqpoint{5.608042in}{3.514022in}}%
\pgfpathlineto{\pgfqpoint{5.662036in}{3.514022in}}%
\pgfpathlineto{\pgfqpoint{5.662036in}{3.833172in}}%
\pgfpathlineto{\pgfqpoint{5.608042in}{3.833172in}}%
\pgfpathlineto{\pgfqpoint{5.608042in}{3.514022in}}%
\pgfpathclose%
\pgfusepath{fill}%
\end{pgfscope}%
\begin{pgfscope}%
\pgfpathrectangle{\pgfqpoint{0.787430in}{3.514022in}}{\pgfqpoint{7.062570in}{2.555373in}}%
\pgfusepath{clip}%
\pgfsetbuttcap%
\pgfsetmiterjoin%
\definecolor{currentfill}{rgb}{0.121569,0.466667,0.705882}%
\pgfsetfillcolor{currentfill}%
\pgfsetlinewidth{0.000000pt}%
\definecolor{currentstroke}{rgb}{0.000000,0.000000,0.000000}%
\pgfsetstrokecolor{currentstroke}%
\pgfsetstrokeopacity{0.000000}%
\pgfsetdash{}{0pt}%
\pgfpathmoveto{\pgfqpoint{5.662036in}{3.514022in}}%
\pgfpathlineto{\pgfqpoint{5.716030in}{3.514022in}}%
\pgfpathlineto{\pgfqpoint{5.716030in}{3.780671in}}%
\pgfpathlineto{\pgfqpoint{5.662036in}{3.780671in}}%
\pgfpathlineto{\pgfqpoint{5.662036in}{3.514022in}}%
\pgfpathclose%
\pgfusepath{fill}%
\end{pgfscope}%
\begin{pgfscope}%
\pgfpathrectangle{\pgfqpoint{0.787430in}{3.514022in}}{\pgfqpoint{7.062570in}{2.555373in}}%
\pgfusepath{clip}%
\pgfsetbuttcap%
\pgfsetmiterjoin%
\definecolor{currentfill}{rgb}{0.121569,0.466667,0.705882}%
\pgfsetfillcolor{currentfill}%
\pgfsetlinewidth{0.000000pt}%
\definecolor{currentstroke}{rgb}{0.000000,0.000000,0.000000}%
\pgfsetstrokecolor{currentstroke}%
\pgfsetstrokeopacity{0.000000}%
\pgfsetdash{}{0pt}%
\pgfpathmoveto{\pgfqpoint{5.716030in}{3.514022in}}%
\pgfpathlineto{\pgfqpoint{5.770024in}{3.514022in}}%
\pgfpathlineto{\pgfqpoint{5.770024in}{3.730933in}}%
\pgfpathlineto{\pgfqpoint{5.716030in}{3.730933in}}%
\pgfpathlineto{\pgfqpoint{5.716030in}{3.514022in}}%
\pgfpathclose%
\pgfusepath{fill}%
\end{pgfscope}%
\begin{pgfscope}%
\pgfpathrectangle{\pgfqpoint{0.787430in}{3.514022in}}{\pgfqpoint{7.062570in}{2.555373in}}%
\pgfusepath{clip}%
\pgfsetbuttcap%
\pgfsetmiterjoin%
\definecolor{currentfill}{rgb}{0.121569,0.466667,0.705882}%
\pgfsetfillcolor{currentfill}%
\pgfsetlinewidth{0.000000pt}%
\definecolor{currentstroke}{rgb}{0.000000,0.000000,0.000000}%
\pgfsetstrokecolor{currentstroke}%
\pgfsetstrokeopacity{0.000000}%
\pgfsetdash{}{0pt}%
\pgfpathmoveto{\pgfqpoint{5.770024in}{3.514022in}}%
\pgfpathlineto{\pgfqpoint{5.824018in}{3.514022in}}%
\pgfpathlineto{\pgfqpoint{5.824018in}{3.688104in}}%
\pgfpathlineto{\pgfqpoint{5.770024in}{3.688104in}}%
\pgfpathlineto{\pgfqpoint{5.770024in}{3.514022in}}%
\pgfpathclose%
\pgfusepath{fill}%
\end{pgfscope}%
\begin{pgfscope}%
\pgfpathrectangle{\pgfqpoint{0.787430in}{3.514022in}}{\pgfqpoint{7.062570in}{2.555373in}}%
\pgfusepath{clip}%
\pgfsetbuttcap%
\pgfsetmiterjoin%
\definecolor{currentfill}{rgb}{0.121569,0.466667,0.705882}%
\pgfsetfillcolor{currentfill}%
\pgfsetlinewidth{0.000000pt}%
\definecolor{currentstroke}{rgb}{0.000000,0.000000,0.000000}%
\pgfsetstrokecolor{currentstroke}%
\pgfsetstrokeopacity{0.000000}%
\pgfsetdash{}{0pt}%
\pgfpathmoveto{\pgfqpoint{5.824018in}{3.514022in}}%
\pgfpathlineto{\pgfqpoint{5.878013in}{3.514022in}}%
\pgfpathlineto{\pgfqpoint{5.878013in}{3.645965in}}%
\pgfpathlineto{\pgfqpoint{5.824018in}{3.645965in}}%
\pgfpathlineto{\pgfqpoint{5.824018in}{3.514022in}}%
\pgfpathclose%
\pgfusepath{fill}%
\end{pgfscope}%
\begin{pgfscope}%
\pgfpathrectangle{\pgfqpoint{0.787430in}{3.514022in}}{\pgfqpoint{7.062570in}{2.555373in}}%
\pgfusepath{clip}%
\pgfsetbuttcap%
\pgfsetmiterjoin%
\definecolor{currentfill}{rgb}{0.121569,0.466667,0.705882}%
\pgfsetfillcolor{currentfill}%
\pgfsetlinewidth{0.000000pt}%
\definecolor{currentstroke}{rgb}{0.000000,0.000000,0.000000}%
\pgfsetstrokecolor{currentstroke}%
\pgfsetstrokeopacity{0.000000}%
\pgfsetdash{}{0pt}%
\pgfpathmoveto{\pgfqpoint{5.878013in}{3.514022in}}%
\pgfpathlineto{\pgfqpoint{5.932007in}{3.514022in}}%
\pgfpathlineto{\pgfqpoint{5.932007in}{3.621096in}}%
\pgfpathlineto{\pgfqpoint{5.878013in}{3.621096in}}%
\pgfpathlineto{\pgfqpoint{5.878013in}{3.514022in}}%
\pgfpathclose%
\pgfusepath{fill}%
\end{pgfscope}%
\begin{pgfscope}%
\pgfpathrectangle{\pgfqpoint{0.787430in}{3.514022in}}{\pgfqpoint{7.062570in}{2.555373in}}%
\pgfusepath{clip}%
\pgfsetbuttcap%
\pgfsetmiterjoin%
\definecolor{currentfill}{rgb}{0.121569,0.466667,0.705882}%
\pgfsetfillcolor{currentfill}%
\pgfsetlinewidth{0.000000pt}%
\definecolor{currentstroke}{rgb}{0.000000,0.000000,0.000000}%
\pgfsetstrokecolor{currentstroke}%
\pgfsetstrokeopacity{0.000000}%
\pgfsetdash{}{0pt}%
\pgfpathmoveto{\pgfqpoint{5.932007in}{3.514022in}}%
\pgfpathlineto{\pgfqpoint{5.986001in}{3.514022in}}%
\pgfpathlineto{\pgfqpoint{5.986001in}{3.593464in}}%
\pgfpathlineto{\pgfqpoint{5.932007in}{3.593464in}}%
\pgfpathlineto{\pgfqpoint{5.932007in}{3.514022in}}%
\pgfpathclose%
\pgfusepath{fill}%
\end{pgfscope}%
\begin{pgfscope}%
\pgfpathrectangle{\pgfqpoint{0.787430in}{3.514022in}}{\pgfqpoint{7.062570in}{2.555373in}}%
\pgfusepath{clip}%
\pgfsetbuttcap%
\pgfsetmiterjoin%
\definecolor{currentfill}{rgb}{0.121569,0.466667,0.705882}%
\pgfsetfillcolor{currentfill}%
\pgfsetlinewidth{0.000000pt}%
\definecolor{currentstroke}{rgb}{0.000000,0.000000,0.000000}%
\pgfsetstrokecolor{currentstroke}%
\pgfsetstrokeopacity{0.000000}%
\pgfsetdash{}{0pt}%
\pgfpathmoveto{\pgfqpoint{5.986001in}{3.514022in}}%
\pgfpathlineto{\pgfqpoint{6.039995in}{3.514022in}}%
\pgfpathlineto{\pgfqpoint{6.039995in}{3.580339in}}%
\pgfpathlineto{\pgfqpoint{5.986001in}{3.580339in}}%
\pgfpathlineto{\pgfqpoint{5.986001in}{3.514022in}}%
\pgfpathclose%
\pgfusepath{fill}%
\end{pgfscope}%
\begin{pgfscope}%
\pgfpathrectangle{\pgfqpoint{0.787430in}{3.514022in}}{\pgfqpoint{7.062570in}{2.555373in}}%
\pgfusepath{clip}%
\pgfsetbuttcap%
\pgfsetmiterjoin%
\definecolor{currentfill}{rgb}{0.121569,0.466667,0.705882}%
\pgfsetfillcolor{currentfill}%
\pgfsetlinewidth{0.000000pt}%
\definecolor{currentstroke}{rgb}{0.000000,0.000000,0.000000}%
\pgfsetstrokecolor{currentstroke}%
\pgfsetstrokeopacity{0.000000}%
\pgfsetdash{}{0pt}%
\pgfpathmoveto{\pgfqpoint{6.039995in}{3.514022in}}%
\pgfpathlineto{\pgfqpoint{6.093989in}{3.514022in}}%
\pgfpathlineto{\pgfqpoint{6.093989in}{3.567214in}}%
\pgfpathlineto{\pgfqpoint{6.039995in}{3.567214in}}%
\pgfpathlineto{\pgfqpoint{6.039995in}{3.514022in}}%
\pgfpathclose%
\pgfusepath{fill}%
\end{pgfscope}%
\begin{pgfscope}%
\pgfpathrectangle{\pgfqpoint{0.787430in}{3.514022in}}{\pgfqpoint{7.062570in}{2.555373in}}%
\pgfusepath{clip}%
\pgfsetbuttcap%
\pgfsetmiterjoin%
\definecolor{currentfill}{rgb}{0.121569,0.466667,0.705882}%
\pgfsetfillcolor{currentfill}%
\pgfsetlinewidth{0.000000pt}%
\definecolor{currentstroke}{rgb}{0.000000,0.000000,0.000000}%
\pgfsetstrokecolor{currentstroke}%
\pgfsetstrokeopacity{0.000000}%
\pgfsetdash{}{0pt}%
\pgfpathmoveto{\pgfqpoint{6.093989in}{3.514022in}}%
\pgfpathlineto{\pgfqpoint{6.147983in}{3.514022in}}%
\pgfpathlineto{\pgfqpoint{6.147983in}{3.565141in}}%
\pgfpathlineto{\pgfqpoint{6.093989in}{3.565141in}}%
\pgfpathlineto{\pgfqpoint{6.093989in}{3.514022in}}%
\pgfpathclose%
\pgfusepath{fill}%
\end{pgfscope}%
\begin{pgfscope}%
\pgfpathrectangle{\pgfqpoint{0.787430in}{3.514022in}}{\pgfqpoint{7.062570in}{2.555373in}}%
\pgfusepath{clip}%
\pgfsetbuttcap%
\pgfsetmiterjoin%
\definecolor{currentfill}{rgb}{0.121569,0.466667,0.705882}%
\pgfsetfillcolor{currentfill}%
\pgfsetlinewidth{0.000000pt}%
\definecolor{currentstroke}{rgb}{0.000000,0.000000,0.000000}%
\pgfsetstrokecolor{currentstroke}%
\pgfsetstrokeopacity{0.000000}%
\pgfsetdash{}{0pt}%
\pgfpathmoveto{\pgfqpoint{6.147983in}{3.514022in}}%
\pgfpathlineto{\pgfqpoint{6.201977in}{3.514022in}}%
\pgfpathlineto{\pgfqpoint{6.201977in}{3.550634in}}%
\pgfpathlineto{\pgfqpoint{6.147983in}{3.550634in}}%
\pgfpathlineto{\pgfqpoint{6.147983in}{3.514022in}}%
\pgfpathclose%
\pgfusepath{fill}%
\end{pgfscope}%
\begin{pgfscope}%
\pgfpathrectangle{\pgfqpoint{0.787430in}{3.514022in}}{\pgfqpoint{7.062570in}{2.555373in}}%
\pgfusepath{clip}%
\pgfsetbuttcap%
\pgfsetmiterjoin%
\definecolor{currentfill}{rgb}{0.121569,0.466667,0.705882}%
\pgfsetfillcolor{currentfill}%
\pgfsetlinewidth{0.000000pt}%
\definecolor{currentstroke}{rgb}{0.000000,0.000000,0.000000}%
\pgfsetstrokecolor{currentstroke}%
\pgfsetstrokeopacity{0.000000}%
\pgfsetdash{}{0pt}%
\pgfpathmoveto{\pgfqpoint{6.201977in}{3.514022in}}%
\pgfpathlineto{\pgfqpoint{6.255972in}{3.514022in}}%
\pgfpathlineto{\pgfqpoint{6.255972in}{3.537509in}}%
\pgfpathlineto{\pgfqpoint{6.201977in}{3.537509in}}%
\pgfpathlineto{\pgfqpoint{6.201977in}{3.514022in}}%
\pgfpathclose%
\pgfusepath{fill}%
\end{pgfscope}%
\begin{pgfscope}%
\pgfpathrectangle{\pgfqpoint{0.787430in}{3.514022in}}{\pgfqpoint{7.062570in}{2.555373in}}%
\pgfusepath{clip}%
\pgfsetbuttcap%
\pgfsetmiterjoin%
\definecolor{currentfill}{rgb}{0.121569,0.466667,0.705882}%
\pgfsetfillcolor{currentfill}%
\pgfsetlinewidth{0.000000pt}%
\definecolor{currentstroke}{rgb}{0.000000,0.000000,0.000000}%
\pgfsetstrokecolor{currentstroke}%
\pgfsetstrokeopacity{0.000000}%
\pgfsetdash{}{0pt}%
\pgfpathmoveto{\pgfqpoint{6.255972in}{3.514022in}}%
\pgfpathlineto{\pgfqpoint{6.309966in}{3.514022in}}%
\pgfpathlineto{\pgfqpoint{6.309966in}{3.530601in}}%
\pgfpathlineto{\pgfqpoint{6.255972in}{3.530601in}}%
\pgfpathlineto{\pgfqpoint{6.255972in}{3.514022in}}%
\pgfpathclose%
\pgfusepath{fill}%
\end{pgfscope}%
\begin{pgfscope}%
\pgfpathrectangle{\pgfqpoint{0.787430in}{3.514022in}}{\pgfqpoint{7.062570in}{2.555373in}}%
\pgfusepath{clip}%
\pgfsetbuttcap%
\pgfsetmiterjoin%
\definecolor{currentfill}{rgb}{0.121569,0.466667,0.705882}%
\pgfsetfillcolor{currentfill}%
\pgfsetlinewidth{0.000000pt}%
\definecolor{currentstroke}{rgb}{0.000000,0.000000,0.000000}%
\pgfsetstrokecolor{currentstroke}%
\pgfsetstrokeopacity{0.000000}%
\pgfsetdash{}{0pt}%
\pgfpathmoveto{\pgfqpoint{6.309966in}{3.514022in}}%
\pgfpathlineto{\pgfqpoint{6.363960in}{3.514022in}}%
\pgfpathlineto{\pgfqpoint{6.363960in}{3.534055in}}%
\pgfpathlineto{\pgfqpoint{6.309966in}{3.534055in}}%
\pgfpathlineto{\pgfqpoint{6.309966in}{3.514022in}}%
\pgfpathclose%
\pgfusepath{fill}%
\end{pgfscope}%
\begin{pgfscope}%
\pgfpathrectangle{\pgfqpoint{0.787430in}{3.514022in}}{\pgfqpoint{7.062570in}{2.555373in}}%
\pgfusepath{clip}%
\pgfsetbuttcap%
\pgfsetmiterjoin%
\definecolor{currentfill}{rgb}{0.121569,0.466667,0.705882}%
\pgfsetfillcolor{currentfill}%
\pgfsetlinewidth{0.000000pt}%
\definecolor{currentstroke}{rgb}{0.000000,0.000000,0.000000}%
\pgfsetstrokecolor{currentstroke}%
\pgfsetstrokeopacity{0.000000}%
\pgfsetdash{}{0pt}%
\pgfpathmoveto{\pgfqpoint{6.363960in}{3.514022in}}%
\pgfpathlineto{\pgfqpoint{6.417954in}{3.514022in}}%
\pgfpathlineto{\pgfqpoint{6.417954in}{3.523693in}}%
\pgfpathlineto{\pgfqpoint{6.363960in}{3.523693in}}%
\pgfpathlineto{\pgfqpoint{6.363960in}{3.514022in}}%
\pgfpathclose%
\pgfusepath{fill}%
\end{pgfscope}%
\begin{pgfscope}%
\pgfpathrectangle{\pgfqpoint{0.787430in}{3.514022in}}{\pgfqpoint{7.062570in}{2.555373in}}%
\pgfusepath{clip}%
\pgfsetbuttcap%
\pgfsetmiterjoin%
\definecolor{currentfill}{rgb}{0.121569,0.466667,0.705882}%
\pgfsetfillcolor{currentfill}%
\pgfsetlinewidth{0.000000pt}%
\definecolor{currentstroke}{rgb}{0.000000,0.000000,0.000000}%
\pgfsetstrokecolor{currentstroke}%
\pgfsetstrokeopacity{0.000000}%
\pgfsetdash{}{0pt}%
\pgfpathmoveto{\pgfqpoint{6.417954in}{3.514022in}}%
\pgfpathlineto{\pgfqpoint{6.471948in}{3.514022in}}%
\pgfpathlineto{\pgfqpoint{6.471948in}{3.522312in}}%
\pgfpathlineto{\pgfqpoint{6.417954in}{3.522312in}}%
\pgfpathlineto{\pgfqpoint{6.417954in}{3.514022in}}%
\pgfpathclose%
\pgfusepath{fill}%
\end{pgfscope}%
\begin{pgfscope}%
\pgfpathrectangle{\pgfqpoint{0.787430in}{3.514022in}}{\pgfqpoint{7.062570in}{2.555373in}}%
\pgfusepath{clip}%
\pgfsetbuttcap%
\pgfsetmiterjoin%
\definecolor{currentfill}{rgb}{0.121569,0.466667,0.705882}%
\pgfsetfillcolor{currentfill}%
\pgfsetlinewidth{0.000000pt}%
\definecolor{currentstroke}{rgb}{0.000000,0.000000,0.000000}%
\pgfsetstrokecolor{currentstroke}%
\pgfsetstrokeopacity{0.000000}%
\pgfsetdash{}{0pt}%
\pgfpathmoveto{\pgfqpoint{6.471948in}{3.514022in}}%
\pgfpathlineto{\pgfqpoint{6.525942in}{3.514022in}}%
\pgfpathlineto{\pgfqpoint{6.525942in}{3.517476in}}%
\pgfpathlineto{\pgfqpoint{6.471948in}{3.517476in}}%
\pgfpathlineto{\pgfqpoint{6.471948in}{3.514022in}}%
\pgfpathclose%
\pgfusepath{fill}%
\end{pgfscope}%
\begin{pgfscope}%
\pgfpathrectangle{\pgfqpoint{0.787430in}{3.514022in}}{\pgfqpoint{7.062570in}{2.555373in}}%
\pgfusepath{clip}%
\pgfsetbuttcap%
\pgfsetmiterjoin%
\definecolor{currentfill}{rgb}{0.121569,0.466667,0.705882}%
\pgfsetfillcolor{currentfill}%
\pgfsetlinewidth{0.000000pt}%
\definecolor{currentstroke}{rgb}{0.000000,0.000000,0.000000}%
\pgfsetstrokecolor{currentstroke}%
\pgfsetstrokeopacity{0.000000}%
\pgfsetdash{}{0pt}%
\pgfpathmoveto{\pgfqpoint{6.525942in}{3.514022in}}%
\pgfpathlineto{\pgfqpoint{6.579936in}{3.514022in}}%
\pgfpathlineto{\pgfqpoint{6.579936in}{3.517476in}}%
\pgfpathlineto{\pgfqpoint{6.525942in}{3.517476in}}%
\pgfpathlineto{\pgfqpoint{6.525942in}{3.514022in}}%
\pgfpathclose%
\pgfusepath{fill}%
\end{pgfscope}%
\begin{pgfscope}%
\pgfpathrectangle{\pgfqpoint{0.787430in}{3.514022in}}{\pgfqpoint{7.062570in}{2.555373in}}%
\pgfusepath{clip}%
\pgfsetbuttcap%
\pgfsetmiterjoin%
\definecolor{currentfill}{rgb}{0.121569,0.466667,0.705882}%
\pgfsetfillcolor{currentfill}%
\pgfsetlinewidth{0.000000pt}%
\definecolor{currentstroke}{rgb}{0.000000,0.000000,0.000000}%
\pgfsetstrokecolor{currentstroke}%
\pgfsetstrokeopacity{0.000000}%
\pgfsetdash{}{0pt}%
\pgfpathmoveto{\pgfqpoint{6.579936in}{3.514022in}}%
\pgfpathlineto{\pgfqpoint{6.633931in}{3.514022in}}%
\pgfpathlineto{\pgfqpoint{6.633931in}{3.518167in}}%
\pgfpathlineto{\pgfqpoint{6.579936in}{3.518167in}}%
\pgfpathlineto{\pgfqpoint{6.579936in}{3.514022in}}%
\pgfpathclose%
\pgfusepath{fill}%
\end{pgfscope}%
\begin{pgfscope}%
\pgfpathrectangle{\pgfqpoint{0.787430in}{3.514022in}}{\pgfqpoint{7.062570in}{2.555373in}}%
\pgfusepath{clip}%
\pgfsetbuttcap%
\pgfsetmiterjoin%
\definecolor{currentfill}{rgb}{0.121569,0.466667,0.705882}%
\pgfsetfillcolor{currentfill}%
\pgfsetlinewidth{0.000000pt}%
\definecolor{currentstroke}{rgb}{0.000000,0.000000,0.000000}%
\pgfsetstrokecolor{currentstroke}%
\pgfsetstrokeopacity{0.000000}%
\pgfsetdash{}{0pt}%
\pgfpathmoveto{\pgfqpoint{6.633931in}{3.514022in}}%
\pgfpathlineto{\pgfqpoint{6.687925in}{3.514022in}}%
\pgfpathlineto{\pgfqpoint{6.687925in}{3.518858in}}%
\pgfpathlineto{\pgfqpoint{6.633931in}{3.518858in}}%
\pgfpathlineto{\pgfqpoint{6.633931in}{3.514022in}}%
\pgfpathclose%
\pgfusepath{fill}%
\end{pgfscope}%
\begin{pgfscope}%
\pgfpathrectangle{\pgfqpoint{0.787430in}{3.514022in}}{\pgfqpoint{7.062570in}{2.555373in}}%
\pgfusepath{clip}%
\pgfsetbuttcap%
\pgfsetmiterjoin%
\definecolor{currentfill}{rgb}{0.121569,0.466667,0.705882}%
\pgfsetfillcolor{currentfill}%
\pgfsetlinewidth{0.000000pt}%
\definecolor{currentstroke}{rgb}{0.000000,0.000000,0.000000}%
\pgfsetstrokecolor{currentstroke}%
\pgfsetstrokeopacity{0.000000}%
\pgfsetdash{}{0pt}%
\pgfpathmoveto{\pgfqpoint{6.687925in}{3.514022in}}%
\pgfpathlineto{\pgfqpoint{6.741919in}{3.514022in}}%
\pgfpathlineto{\pgfqpoint{6.741919in}{3.514022in}}%
\pgfpathlineto{\pgfqpoint{6.687925in}{3.514022in}}%
\pgfpathlineto{\pgfqpoint{6.687925in}{3.514022in}}%
\pgfpathclose%
\pgfusepath{fill}%
\end{pgfscope}%
\begin{pgfscope}%
\pgfpathrectangle{\pgfqpoint{0.787430in}{3.514022in}}{\pgfqpoint{7.062570in}{2.555373in}}%
\pgfusepath{clip}%
\pgfsetbuttcap%
\pgfsetmiterjoin%
\definecolor{currentfill}{rgb}{0.121569,0.466667,0.705882}%
\pgfsetfillcolor{currentfill}%
\pgfsetlinewidth{0.000000pt}%
\definecolor{currentstroke}{rgb}{0.000000,0.000000,0.000000}%
\pgfsetstrokecolor{currentstroke}%
\pgfsetstrokeopacity{0.000000}%
\pgfsetdash{}{0pt}%
\pgfpathmoveto{\pgfqpoint{6.741919in}{3.514022in}}%
\pgfpathlineto{\pgfqpoint{6.795913in}{3.514022in}}%
\pgfpathlineto{\pgfqpoint{6.795913in}{3.514713in}}%
\pgfpathlineto{\pgfqpoint{6.741919in}{3.514713in}}%
\pgfpathlineto{\pgfqpoint{6.741919in}{3.514022in}}%
\pgfpathclose%
\pgfusepath{fill}%
\end{pgfscope}%
\begin{pgfscope}%
\pgfpathrectangle{\pgfqpoint{0.787430in}{3.514022in}}{\pgfqpoint{7.062570in}{2.555373in}}%
\pgfusepath{clip}%
\pgfsetbuttcap%
\pgfsetmiterjoin%
\definecolor{currentfill}{rgb}{0.121569,0.466667,0.705882}%
\pgfsetfillcolor{currentfill}%
\pgfsetlinewidth{0.000000pt}%
\definecolor{currentstroke}{rgb}{0.000000,0.000000,0.000000}%
\pgfsetstrokecolor{currentstroke}%
\pgfsetstrokeopacity{0.000000}%
\pgfsetdash{}{0pt}%
\pgfpathmoveto{\pgfqpoint{6.795913in}{3.514022in}}%
\pgfpathlineto{\pgfqpoint{6.849907in}{3.514022in}}%
\pgfpathlineto{\pgfqpoint{6.849907in}{3.515404in}}%
\pgfpathlineto{\pgfqpoint{6.795913in}{3.515404in}}%
\pgfpathlineto{\pgfqpoint{6.795913in}{3.514022in}}%
\pgfpathclose%
\pgfusepath{fill}%
\end{pgfscope}%
\begin{pgfscope}%
\pgfpathrectangle{\pgfqpoint{0.787430in}{3.514022in}}{\pgfqpoint{7.062570in}{2.555373in}}%
\pgfusepath{clip}%
\pgfsetbuttcap%
\pgfsetmiterjoin%
\definecolor{currentfill}{rgb}{0.121569,0.466667,0.705882}%
\pgfsetfillcolor{currentfill}%
\pgfsetlinewidth{0.000000pt}%
\definecolor{currentstroke}{rgb}{0.000000,0.000000,0.000000}%
\pgfsetstrokecolor{currentstroke}%
\pgfsetstrokeopacity{0.000000}%
\pgfsetdash{}{0pt}%
\pgfpathmoveto{\pgfqpoint{6.849907in}{3.514022in}}%
\pgfpathlineto{\pgfqpoint{6.903901in}{3.514022in}}%
\pgfpathlineto{\pgfqpoint{6.903901in}{3.516785in}}%
\pgfpathlineto{\pgfqpoint{6.849907in}{3.516785in}}%
\pgfpathlineto{\pgfqpoint{6.849907in}{3.514022in}}%
\pgfpathclose%
\pgfusepath{fill}%
\end{pgfscope}%
\begin{pgfscope}%
\pgfpathrectangle{\pgfqpoint{0.787430in}{3.514022in}}{\pgfqpoint{7.062570in}{2.555373in}}%
\pgfusepath{clip}%
\pgfsetbuttcap%
\pgfsetmiterjoin%
\definecolor{currentfill}{rgb}{0.121569,0.466667,0.705882}%
\pgfsetfillcolor{currentfill}%
\pgfsetlinewidth{0.000000pt}%
\definecolor{currentstroke}{rgb}{0.000000,0.000000,0.000000}%
\pgfsetstrokecolor{currentstroke}%
\pgfsetstrokeopacity{0.000000}%
\pgfsetdash{}{0pt}%
\pgfpathmoveto{\pgfqpoint{6.903901in}{3.514022in}}%
\pgfpathlineto{\pgfqpoint{6.957895in}{3.514022in}}%
\pgfpathlineto{\pgfqpoint{6.957895in}{3.514713in}}%
\pgfpathlineto{\pgfqpoint{6.903901in}{3.514713in}}%
\pgfpathlineto{\pgfqpoint{6.903901in}{3.514022in}}%
\pgfpathclose%
\pgfusepath{fill}%
\end{pgfscope}%
\begin{pgfscope}%
\pgfsetbuttcap%
\pgfsetroundjoin%
\definecolor{currentfill}{rgb}{0.000000,0.000000,0.000000}%
\pgfsetfillcolor{currentfill}%
\pgfsetlinewidth{0.803000pt}%
\definecolor{currentstroke}{rgb}{0.000000,0.000000,0.000000}%
\pgfsetstrokecolor{currentstroke}%
\pgfsetdash{}{0pt}%
\pgfsys@defobject{currentmarker}{\pgfqpoint{0.000000in}{-0.048611in}}{\pgfqpoint{0.000000in}{0.000000in}}{%
\pgfpathmoveto{\pgfqpoint{0.000000in}{0.000000in}}%
\pgfpathlineto{\pgfqpoint{0.000000in}{-0.048611in}}%
\pgfusepath{stroke,fill}%
}%
\begin{pgfscope}%
\pgfsys@transformshift{1.923827in}{3.514022in}%
\pgfsys@useobject{currentmarker}{}%
\end{pgfscope}%
\end{pgfscope}%
\begin{pgfscope}%
\pgfsetbuttcap%
\pgfsetroundjoin%
\definecolor{currentfill}{rgb}{0.000000,0.000000,0.000000}%
\pgfsetfillcolor{currentfill}%
\pgfsetlinewidth{0.803000pt}%
\definecolor{currentstroke}{rgb}{0.000000,0.000000,0.000000}%
\pgfsetstrokecolor{currentstroke}%
\pgfsetdash{}{0pt}%
\pgfsys@defobject{currentmarker}{\pgfqpoint{0.000000in}{-0.048611in}}{\pgfqpoint{0.000000in}{0.000000in}}{%
\pgfpathmoveto{\pgfqpoint{0.000000in}{0.000000in}}%
\pgfpathlineto{\pgfqpoint{0.000000in}{-0.048611in}}%
\pgfusepath{stroke,fill}%
}%
\begin{pgfscope}%
\pgfsys@transformshift{3.153546in}{3.514022in}%
\pgfsys@useobject{currentmarker}{}%
\end{pgfscope}%
\end{pgfscope}%
\begin{pgfscope}%
\pgfsetbuttcap%
\pgfsetroundjoin%
\definecolor{currentfill}{rgb}{0.000000,0.000000,0.000000}%
\pgfsetfillcolor{currentfill}%
\pgfsetlinewidth{0.803000pt}%
\definecolor{currentstroke}{rgb}{0.000000,0.000000,0.000000}%
\pgfsetstrokecolor{currentstroke}%
\pgfsetdash{}{0pt}%
\pgfsys@defobject{currentmarker}{\pgfqpoint{0.000000in}{-0.048611in}}{\pgfqpoint{0.000000in}{0.000000in}}{%
\pgfpathmoveto{\pgfqpoint{0.000000in}{0.000000in}}%
\pgfpathlineto{\pgfqpoint{0.000000in}{-0.048611in}}%
\pgfusepath{stroke,fill}%
}%
\begin{pgfscope}%
\pgfsys@transformshift{4.383264in}{3.514022in}%
\pgfsys@useobject{currentmarker}{}%
\end{pgfscope}%
\end{pgfscope}%
\begin{pgfscope}%
\pgfsetbuttcap%
\pgfsetroundjoin%
\definecolor{currentfill}{rgb}{0.000000,0.000000,0.000000}%
\pgfsetfillcolor{currentfill}%
\pgfsetlinewidth{0.803000pt}%
\definecolor{currentstroke}{rgb}{0.000000,0.000000,0.000000}%
\pgfsetstrokecolor{currentstroke}%
\pgfsetdash{}{0pt}%
\pgfsys@defobject{currentmarker}{\pgfqpoint{0.000000in}{-0.048611in}}{\pgfqpoint{0.000000in}{0.000000in}}{%
\pgfpathmoveto{\pgfqpoint{0.000000in}{0.000000in}}%
\pgfpathlineto{\pgfqpoint{0.000000in}{-0.048611in}}%
\pgfusepath{stroke,fill}%
}%
\begin{pgfscope}%
\pgfsys@transformshift{5.612983in}{3.514022in}%
\pgfsys@useobject{currentmarker}{}%
\end{pgfscope}%
\end{pgfscope}%
\begin{pgfscope}%
\pgfsetbuttcap%
\pgfsetroundjoin%
\definecolor{currentfill}{rgb}{0.000000,0.000000,0.000000}%
\pgfsetfillcolor{currentfill}%
\pgfsetlinewidth{0.803000pt}%
\definecolor{currentstroke}{rgb}{0.000000,0.000000,0.000000}%
\pgfsetstrokecolor{currentstroke}%
\pgfsetdash{}{0pt}%
\pgfsys@defobject{currentmarker}{\pgfqpoint{0.000000in}{-0.048611in}}{\pgfqpoint{0.000000in}{0.000000in}}{%
\pgfpathmoveto{\pgfqpoint{0.000000in}{0.000000in}}%
\pgfpathlineto{\pgfqpoint{0.000000in}{-0.048611in}}%
\pgfusepath{stroke,fill}%
}%
\begin{pgfscope}%
\pgfsys@transformshift{6.842701in}{3.514022in}%
\pgfsys@useobject{currentmarker}{}%
\end{pgfscope}%
\end{pgfscope}%
\begin{pgfscope}%
\pgfsetbuttcap%
\pgfsetroundjoin%
\definecolor{currentfill}{rgb}{0.000000,0.000000,0.000000}%
\pgfsetfillcolor{currentfill}%
\pgfsetlinewidth{0.803000pt}%
\definecolor{currentstroke}{rgb}{0.000000,0.000000,0.000000}%
\pgfsetstrokecolor{currentstroke}%
\pgfsetdash{}{0pt}%
\pgfsys@defobject{currentmarker}{\pgfqpoint{-0.048611in}{0.000000in}}{\pgfqpoint{-0.000000in}{0.000000in}}{%
\pgfpathmoveto{\pgfqpoint{-0.000000in}{0.000000in}}%
\pgfpathlineto{\pgfqpoint{-0.048611in}{0.000000in}}%
\pgfusepath{stroke,fill}%
}%
\begin{pgfscope}%
\pgfsys@transformshift{0.787430in}{3.514022in}%
\pgfsys@useobject{currentmarker}{}%
\end{pgfscope}%
\end{pgfscope}%
\begin{pgfscope}%
\definecolor{textcolor}{rgb}{0.000000,0.000000,0.000000}%
\pgfsetstrokecolor{textcolor}%
\pgfsetfillcolor{textcolor}%
\pgftext[x=0.512738in, y=3.465797in, left, base]{\color{textcolor}{\rmfamily\fontsize{10.000000}{12.000000}\selectfont\catcode`\^=\active\def^{\ifmmode\sp\else\^{}\fi}\catcode`\%=\active\def%{\%}$\mathdefault{0.0}$}}%
\end{pgfscope}%
\begin{pgfscope}%
\pgfsetbuttcap%
\pgfsetroundjoin%
\definecolor{currentfill}{rgb}{0.000000,0.000000,0.000000}%
\pgfsetfillcolor{currentfill}%
\pgfsetlinewidth{0.803000pt}%
\definecolor{currentstroke}{rgb}{0.000000,0.000000,0.000000}%
\pgfsetstrokecolor{currentstroke}%
\pgfsetdash{}{0pt}%
\pgfsys@defobject{currentmarker}{\pgfqpoint{-0.048611in}{0.000000in}}{\pgfqpoint{-0.000000in}{0.000000in}}{%
\pgfpathmoveto{\pgfqpoint{-0.000000in}{0.000000in}}%
\pgfpathlineto{\pgfqpoint{-0.048611in}{0.000000in}}%
\pgfusepath{stroke,fill}%
}%
\begin{pgfscope}%
\pgfsys@transformshift{0.787430in}{4.727280in}%
\pgfsys@useobject{currentmarker}{}%
\end{pgfscope}%
\end{pgfscope}%
\begin{pgfscope}%
\definecolor{textcolor}{rgb}{0.000000,0.000000,0.000000}%
\pgfsetstrokecolor{textcolor}%
\pgfsetfillcolor{textcolor}%
\pgftext[x=0.512738in, y=4.679055in, left, base]{\color{textcolor}{\rmfamily\fontsize{10.000000}{12.000000}\selectfont\catcode`\^=\active\def^{\ifmmode\sp\else\^{}\fi}\catcode`\%=\active\def%{\%}$\mathdefault{0.2}$}}%
\end{pgfscope}%
\begin{pgfscope}%
\pgfsetbuttcap%
\pgfsetroundjoin%
\definecolor{currentfill}{rgb}{0.000000,0.000000,0.000000}%
\pgfsetfillcolor{currentfill}%
\pgfsetlinewidth{0.803000pt}%
\definecolor{currentstroke}{rgb}{0.000000,0.000000,0.000000}%
\pgfsetstrokecolor{currentstroke}%
\pgfsetdash{}{0pt}%
\pgfsys@defobject{currentmarker}{\pgfqpoint{-0.048611in}{0.000000in}}{\pgfqpoint{-0.000000in}{0.000000in}}{%
\pgfpathmoveto{\pgfqpoint{-0.000000in}{0.000000in}}%
\pgfpathlineto{\pgfqpoint{-0.048611in}{0.000000in}}%
\pgfusepath{stroke,fill}%
}%
\begin{pgfscope}%
\pgfsys@transformshift{0.787430in}{5.940539in}%
\pgfsys@useobject{currentmarker}{}%
\end{pgfscope}%
\end{pgfscope}%
\begin{pgfscope}%
\definecolor{textcolor}{rgb}{0.000000,0.000000,0.000000}%
\pgfsetstrokecolor{textcolor}%
\pgfsetfillcolor{textcolor}%
\pgftext[x=0.512738in, y=5.892314in, left, base]{\color{textcolor}{\rmfamily\fontsize{10.000000}{12.000000}\selectfont\catcode`\^=\active\def^{\ifmmode\sp\else\^{}\fi}\catcode`\%=\active\def%{\%}$\mathdefault{0.4}$}}%
\end{pgfscope}%
\begin{pgfscope}%
\definecolor{textcolor}{rgb}{0.000000,0.000000,0.000000}%
\pgfsetstrokecolor{textcolor}%
\pgfsetfillcolor{textcolor}%
\pgftext[x=0.457183in,y=4.791709in,,bottom,rotate=90.000000]{\color{textcolor}{\rmfamily\fontsize{24.000000}{28.800000}\selectfont\catcode`\^=\active\def^{\ifmmode\sp\else\^{}\fi}\catcode`\%=\active\def%{\%}$P_x$}}%
\end{pgfscope}%
\begin{pgfscope}%
\pgfpathrectangle{\pgfqpoint{0.787430in}{3.514022in}}{\pgfqpoint{7.062570in}{2.555373in}}%
\pgfusepath{clip}%
\pgfsetrectcap%
\pgfsetroundjoin%
\pgfsetlinewidth{2.007500pt}%
\definecolor{currentstroke}{rgb}{1.000000,0.000000,0.000000}%
\pgfsetstrokecolor{currentstroke}%
\pgfsetdash{}{0pt}%
\pgfpathmoveto{\pgfqpoint{1.923827in}{3.514834in}}%
\pgfpathlineto{\pgfqpoint{2.135551in}{3.517055in}}%
\pgfpathlineto{\pgfqpoint{2.268493in}{3.520554in}}%
\pgfpathlineto{\pgfqpoint{2.366969in}{3.525209in}}%
\pgfpathlineto{\pgfqpoint{2.450674in}{3.531342in}}%
\pgfpathlineto{\pgfqpoint{2.519607in}{3.538505in}}%
\pgfpathlineto{\pgfqpoint{2.583616in}{3.547407in}}%
\pgfpathlineto{\pgfqpoint{2.637778in}{3.557059in}}%
\pgfpathlineto{\pgfqpoint{2.687016in}{3.567871in}}%
\pgfpathlineto{\pgfqpoint{2.736254in}{3.580969in}}%
\pgfpathlineto{\pgfqpoint{2.780568in}{3.595015in}}%
\pgfpathlineto{\pgfqpoint{2.824882in}{3.611500in}}%
\pgfpathlineto{\pgfqpoint{2.864273in}{3.628451in}}%
\pgfpathlineto{\pgfqpoint{2.903663in}{3.647798in}}%
\pgfpathlineto{\pgfqpoint{2.943053in}{3.669776in}}%
\pgfpathlineto{\pgfqpoint{2.982444in}{3.694623in}}%
\pgfpathlineto{\pgfqpoint{3.016910in}{3.718903in}}%
\pgfpathlineto{\pgfqpoint{3.051377in}{3.745718in}}%
\pgfpathlineto{\pgfqpoint{3.085844in}{3.775221in}}%
\pgfpathlineto{\pgfqpoint{3.120310in}{3.807556in}}%
\pgfpathlineto{\pgfqpoint{3.154777in}{3.842860in}}%
\pgfpathlineto{\pgfqpoint{3.189243in}{3.881254in}}%
\pgfpathlineto{\pgfqpoint{3.223710in}{3.922845in}}%
\pgfpathlineto{\pgfqpoint{3.258177in}{3.967717in}}%
\pgfpathlineto{\pgfqpoint{3.292643in}{4.015935in}}%
\pgfpathlineto{\pgfqpoint{3.332033in}{4.075185in}}%
\pgfpathlineto{\pgfqpoint{3.371424in}{4.138858in}}%
\pgfpathlineto{\pgfqpoint{3.410814in}{4.206907in}}%
\pgfpathlineto{\pgfqpoint{3.450205in}{4.279220in}}%
\pgfpathlineto{\pgfqpoint{3.494519in}{4.365444in}}%
\pgfpathlineto{\pgfqpoint{3.538833in}{4.456476in}}%
\pgfpathlineto{\pgfqpoint{3.588071in}{4.562674in}}%
\pgfpathlineto{\pgfqpoint{3.642233in}{4.684660in}}%
\pgfpathlineto{\pgfqpoint{3.711166in}{4.845624in}}%
\pgfpathlineto{\pgfqpoint{3.908118in}{5.309408in}}%
\pgfpathlineto{\pgfqpoint{3.957356in}{5.417909in}}%
\pgfpathlineto{\pgfqpoint{4.001670in}{5.510180in}}%
\pgfpathlineto{\pgfqpoint{4.041060in}{5.586887in}}%
\pgfpathlineto{\pgfqpoint{4.075527in}{5.649219in}}%
\pgfpathlineto{\pgfqpoint{4.105070in}{5.698665in}}%
\pgfpathlineto{\pgfqpoint{4.134613in}{5.744102in}}%
\pgfpathlineto{\pgfqpoint{4.164155in}{5.785235in}}%
\pgfpathlineto{\pgfqpoint{4.188774in}{5.816028in}}%
\pgfpathlineto{\pgfqpoint{4.213393in}{5.843501in}}%
\pgfpathlineto{\pgfqpoint{4.238012in}{5.867526in}}%
\pgfpathlineto{\pgfqpoint{4.262631in}{5.887990in}}%
\pgfpathlineto{\pgfqpoint{4.282326in}{5.901731in}}%
\pgfpathlineto{\pgfqpoint{4.302022in}{5.913088in}}%
\pgfpathlineto{\pgfqpoint{4.321717in}{5.922028in}}%
\pgfpathlineto{\pgfqpoint{4.341412in}{5.928523in}}%
\pgfpathlineto{\pgfqpoint{4.361107in}{5.932552in}}%
\pgfpathlineto{\pgfqpoint{4.380802in}{5.934103in}}%
\pgfpathlineto{\pgfqpoint{4.400498in}{5.933172in}}%
\pgfpathlineto{\pgfqpoint{4.420193in}{5.929762in}}%
\pgfpathlineto{\pgfqpoint{4.439888in}{5.923882in}}%
\pgfpathlineto{\pgfqpoint{4.459583in}{5.915551in}}%
\pgfpathlineto{\pgfqpoint{4.479278in}{5.904795in}}%
\pgfpathlineto{\pgfqpoint{4.498974in}{5.891646in}}%
\pgfpathlineto{\pgfqpoint{4.518669in}{5.876145in}}%
\pgfpathlineto{\pgfqpoint{4.543288in}{5.853532in}}%
\pgfpathlineto{\pgfqpoint{4.567907in}{5.827424in}}%
\pgfpathlineto{\pgfqpoint{4.592526in}{5.797942in}}%
\pgfpathlineto{\pgfqpoint{4.617145in}{5.765224in}}%
\pgfpathlineto{\pgfqpoint{4.646688in}{5.721904in}}%
\pgfpathlineto{\pgfqpoint{4.676230in}{5.674424in}}%
\pgfpathlineto{\pgfqpoint{4.710697in}{5.614182in}}%
\pgfpathlineto{\pgfqpoint{4.745163in}{5.549215in}}%
\pgfpathlineto{\pgfqpoint{4.784554in}{5.469885in}}%
\pgfpathlineto{\pgfqpoint{4.828868in}{5.375175in}}%
\pgfpathlineto{\pgfqpoint{4.878106in}{5.264613in}}%
\pgfpathlineto{\pgfqpoint{4.942115in}{5.115219in}}%
\pgfpathlineto{\pgfqpoint{5.143991in}{4.639751in}}%
\pgfpathlineto{\pgfqpoint{5.198153in}{4.519603in}}%
\pgfpathlineto{\pgfqpoint{5.247391in}{4.415452in}}%
\pgfpathlineto{\pgfqpoint{5.291705in}{4.326505in}}%
\pgfpathlineto{\pgfqpoint{5.336019in}{4.242540in}}%
\pgfpathlineto{\pgfqpoint{5.375410in}{4.172341in}}%
\pgfpathlineto{\pgfqpoint{5.414800in}{4.106470in}}%
\pgfpathlineto{\pgfqpoint{5.454190in}{4.045006in}}%
\pgfpathlineto{\pgfqpoint{5.493581in}{3.987970in}}%
\pgfpathlineto{\pgfqpoint{5.528047in}{3.941670in}}%
\pgfpathlineto{\pgfqpoint{5.562514in}{3.898682in}}%
\pgfpathlineto{\pgfqpoint{5.596980in}{3.858930in}}%
\pgfpathlineto{\pgfqpoint{5.631447in}{3.822316in}}%
\pgfpathlineto{\pgfqpoint{5.665914in}{3.788724in}}%
\pgfpathlineto{\pgfqpoint{5.700380in}{3.758024in}}%
\pgfpathlineto{\pgfqpoint{5.734847in}{3.730076in}}%
\pgfpathlineto{\pgfqpoint{5.769313in}{3.704728in}}%
\pgfpathlineto{\pgfqpoint{5.803780in}{3.681826in}}%
\pgfpathlineto{\pgfqpoint{5.843170in}{3.658444in}}%
\pgfpathlineto{\pgfqpoint{5.882561in}{3.637810in}}%
\pgfpathlineto{\pgfqpoint{5.921951in}{3.619690in}}%
\pgfpathlineto{\pgfqpoint{5.961342in}{3.603853in}}%
\pgfpathlineto{\pgfqpoint{6.005656in}{3.588489in}}%
\pgfpathlineto{\pgfqpoint{6.049970in}{3.575433in}}%
\pgfpathlineto{\pgfqpoint{6.099208in}{3.563292in}}%
\pgfpathlineto{\pgfqpoint{6.153370in}{3.552404in}}%
\pgfpathlineto{\pgfqpoint{6.212455in}{3.542994in}}%
\pgfpathlineto{\pgfqpoint{6.276465in}{3.535163in}}%
\pgfpathlineto{\pgfqpoint{6.350322in}{3.528522in}}%
\pgfpathlineto{\pgfqpoint{6.434026in}{3.523316in}}%
\pgfpathlineto{\pgfqpoint{6.537426in}{3.519251in}}%
\pgfpathlineto{\pgfqpoint{6.670368in}{3.516417in}}%
\pgfpathlineto{\pgfqpoint{6.842701in}{3.514834in}}%
\pgfpathlineto{\pgfqpoint{6.842701in}{3.514834in}}%
\pgfusepath{stroke}%
\end{pgfscope}%
\begin{pgfscope}%
\pgfsetrectcap%
\pgfsetmiterjoin%
\pgfsetlinewidth{0.803000pt}%
\definecolor{currentstroke}{rgb}{0.000000,0.000000,0.000000}%
\pgfsetstrokecolor{currentstroke}%
\pgfsetdash{}{0pt}%
\pgfpathmoveto{\pgfqpoint{0.787430in}{3.514022in}}%
\pgfpathlineto{\pgfqpoint{0.787430in}{6.069395in}}%
\pgfusepath{stroke}%
\end{pgfscope}%
\begin{pgfscope}%
\pgfsetrectcap%
\pgfsetmiterjoin%
\pgfsetlinewidth{0.803000pt}%
\definecolor{currentstroke}{rgb}{0.000000,0.000000,0.000000}%
\pgfsetstrokecolor{currentstroke}%
\pgfsetdash{}{0pt}%
\pgfpathmoveto{\pgfqpoint{7.850000in}{3.514022in}}%
\pgfpathlineto{\pgfqpoint{7.850000in}{6.069395in}}%
\pgfusepath{stroke}%
\end{pgfscope}%
\begin{pgfscope}%
\pgfsetrectcap%
\pgfsetmiterjoin%
\pgfsetlinewidth{0.803000pt}%
\definecolor{currentstroke}{rgb}{0.000000,0.000000,0.000000}%
\pgfsetstrokecolor{currentstroke}%
\pgfsetdash{}{0pt}%
\pgfpathmoveto{\pgfqpoint{0.787430in}{3.514022in}}%
\pgfpathlineto{\pgfqpoint{7.850000in}{3.514022in}}%
\pgfusepath{stroke}%
\end{pgfscope}%
\begin{pgfscope}%
\pgfsetrectcap%
\pgfsetmiterjoin%
\pgfsetlinewidth{0.803000pt}%
\definecolor{currentstroke}{rgb}{0.000000,0.000000,0.000000}%
\pgfsetstrokecolor{currentstroke}%
\pgfsetdash{}{0pt}%
\pgfpathmoveto{\pgfqpoint{0.787430in}{6.069395in}}%
\pgfpathlineto{\pgfqpoint{7.850000in}{6.069395in}}%
\pgfusepath{stroke}%
\end{pgfscope}%
\begin{pgfscope}%
\pgfsetbuttcap%
\pgfsetmiterjoin%
\definecolor{currentfill}{rgb}{1.000000,1.000000,1.000000}%
\pgfsetfillcolor{currentfill}%
\pgfsetfillopacity{0.800000}%
\pgfsetlinewidth{1.003750pt}%
\definecolor{currentstroke}{rgb}{0.800000,0.800000,0.800000}%
\pgfsetstrokecolor{currentstroke}%
\pgfsetstrokeopacity{0.800000}%
\pgfsetdash{}{0pt}%
\pgfpathmoveto{\pgfqpoint{6.955323in}{5.570939in}}%
\pgfpathlineto{\pgfqpoint{7.752778in}{5.570939in}}%
\pgfpathquadraticcurveto{\pgfqpoint{7.780556in}{5.570939in}}{\pgfqpoint{7.780556in}{5.598716in}}%
\pgfpathlineto{\pgfqpoint{7.780556in}{5.972173in}}%
\pgfpathquadraticcurveto{\pgfqpoint{7.780556in}{5.999951in}}{\pgfqpoint{7.752778in}{5.999951in}}%
\pgfpathlineto{\pgfqpoint{6.955323in}{5.999951in}}%
\pgfpathquadraticcurveto{\pgfqpoint{6.927545in}{5.999951in}}{\pgfqpoint{6.927545in}{5.972173in}}%
\pgfpathlineto{\pgfqpoint{6.927545in}{5.598716in}}%
\pgfpathquadraticcurveto{\pgfqpoint{6.927545in}{5.570939in}}{\pgfqpoint{6.955323in}{5.570939in}}%
\pgfpathlineto{\pgfqpoint{6.955323in}{5.570939in}}%
\pgfpathclose%
\pgfusepath{stroke,fill}%
\end{pgfscope}%
\begin{pgfscope}%
\pgfsetbuttcap%
\pgfsetmiterjoin%
\definecolor{currentfill}{rgb}{0.121569,0.466667,0.705882}%
\pgfsetfillcolor{currentfill}%
\pgfsetlinewidth{0.000000pt}%
\definecolor{currentstroke}{rgb}{0.000000,0.000000,0.000000}%
\pgfsetstrokecolor{currentstroke}%
\pgfsetstrokeopacity{0.000000}%
\pgfsetdash{}{0pt}%
\pgfpathmoveto{\pgfqpoint{6.983101in}{5.847173in}}%
\pgfpathlineto{\pgfqpoint{7.260879in}{5.847173in}}%
\pgfpathlineto{\pgfqpoint{7.260879in}{5.944395in}}%
\pgfpathlineto{\pgfqpoint{6.983101in}{5.944395in}}%
\pgfpathlineto{\pgfqpoint{6.983101in}{5.847173in}}%
\pgfpathclose%
\pgfusepath{fill}%
\end{pgfscope}%
\begin{pgfscope}%
\definecolor{textcolor}{rgb}{0.000000,0.000000,0.000000}%
\pgfsetstrokecolor{textcolor}%
\pgfsetfillcolor{textcolor}%
\pgftext[x=7.371990in,y=5.847173in,left,base]{\color{textcolor}{\rmfamily\fontsize{10.000000}{12.000000}\selectfont\catcode`\^=\active\def^{\ifmmode\sp\else\^{}\fi}\catcode`\%=\active\def%{\%}RNG}}%
\end{pgfscope}%
\begin{pgfscope}%
\pgfsetrectcap%
\pgfsetroundjoin%
\pgfsetlinewidth{2.007500pt}%
\definecolor{currentstroke}{rgb}{1.000000,0.000000,0.000000}%
\pgfsetstrokecolor{currentstroke}%
\pgfsetdash{}{0pt}%
\pgfpathmoveto{\pgfqpoint{6.983101in}{5.702111in}}%
\pgfpathlineto{\pgfqpoint{7.121990in}{5.702111in}}%
\pgfpathlineto{\pgfqpoint{7.260879in}{5.702111in}}%
\pgfusepath{stroke}%
\end{pgfscope}%
\begin{pgfscope}%
\definecolor{textcolor}{rgb}{0.000000,0.000000,0.000000}%
\pgfsetstrokecolor{textcolor}%
\pgfsetfillcolor{textcolor}%
\pgftext[x=7.371990in,y=5.653500in,left,base]{\color{textcolor}{\rmfamily\fontsize{10.000000}{12.000000}\selectfont\catcode`\^=\active\def^{\ifmmode\sp\else\^{}\fi}\catcode`\%=\active\def%{\%}Exact}}%
\end{pgfscope}%
\begin{pgfscope}%
\pgfsetbuttcap%
\pgfsetmiterjoin%
\definecolor{currentfill}{rgb}{1.000000,1.000000,1.000000}%
\pgfsetfillcolor{currentfill}%
\pgfsetlinewidth{0.000000pt}%
\definecolor{currentstroke}{rgb}{0.000000,0.000000,0.000000}%
\pgfsetstrokecolor{currentstroke}%
\pgfsetstrokeopacity{0.000000}%
\pgfsetdash{}{0pt}%
\pgfpathmoveto{\pgfqpoint{0.787430in}{0.733417in}}%
\pgfpathlineto{\pgfqpoint{7.850000in}{0.733417in}}%
\pgfpathlineto{\pgfqpoint{7.850000in}{3.288791in}}%
\pgfpathlineto{\pgfqpoint{0.787430in}{3.288791in}}%
\pgfpathlineto{\pgfqpoint{0.787430in}{0.733417in}}%
\pgfpathclose%
\pgfusepath{fill}%
\end{pgfscope}%
\begin{pgfscope}%
\pgfpathrectangle{\pgfqpoint{0.787430in}{0.733417in}}{\pgfqpoint{7.062570in}{2.555373in}}%
\pgfusepath{clip}%
\pgfsetbuttcap%
\pgfsetmiterjoin%
\definecolor{currentfill}{rgb}{0.121569,0.466667,0.705882}%
\pgfsetfillcolor{currentfill}%
\pgfsetlinewidth{0.000000pt}%
\definecolor{currentstroke}{rgb}{0.000000,0.000000,0.000000}%
\pgfsetstrokecolor{currentstroke}%
\pgfsetstrokeopacity{0.000000}%
\pgfsetdash{}{0pt}%
\pgfpathmoveto{\pgfqpoint{1.108456in}{0.733417in}}%
\pgfpathlineto{\pgfqpoint{1.172661in}{0.733417in}}%
\pgfpathlineto{\pgfqpoint{1.172661in}{0.733423in}}%
\pgfpathlineto{\pgfqpoint{1.108456in}{0.733423in}}%
\pgfpathlineto{\pgfqpoint{1.108456in}{0.733417in}}%
\pgfpathclose%
\pgfusepath{fill}%
\end{pgfscope}%
\begin{pgfscope}%
\pgfpathrectangle{\pgfqpoint{0.787430in}{0.733417in}}{\pgfqpoint{7.062570in}{2.555373in}}%
\pgfusepath{clip}%
\pgfsetbuttcap%
\pgfsetmiterjoin%
\definecolor{currentfill}{rgb}{0.121569,0.466667,0.705882}%
\pgfsetfillcolor{currentfill}%
\pgfsetlinewidth{0.000000pt}%
\definecolor{currentstroke}{rgb}{0.000000,0.000000,0.000000}%
\pgfsetstrokecolor{currentstroke}%
\pgfsetstrokeopacity{0.000000}%
\pgfsetdash{}{0pt}%
\pgfpathmoveto{\pgfqpoint{1.172661in}{0.733417in}}%
\pgfpathlineto{\pgfqpoint{1.236867in}{0.733417in}}%
\pgfpathlineto{\pgfqpoint{1.236867in}{0.733417in}}%
\pgfpathlineto{\pgfqpoint{1.172661in}{0.733417in}}%
\pgfpathlineto{\pgfqpoint{1.172661in}{0.733417in}}%
\pgfpathclose%
\pgfusepath{fill}%
\end{pgfscope}%
\begin{pgfscope}%
\pgfpathrectangle{\pgfqpoint{0.787430in}{0.733417in}}{\pgfqpoint{7.062570in}{2.555373in}}%
\pgfusepath{clip}%
\pgfsetbuttcap%
\pgfsetmiterjoin%
\definecolor{currentfill}{rgb}{0.121569,0.466667,0.705882}%
\pgfsetfillcolor{currentfill}%
\pgfsetlinewidth{0.000000pt}%
\definecolor{currentstroke}{rgb}{0.000000,0.000000,0.000000}%
\pgfsetstrokecolor{currentstroke}%
\pgfsetstrokeopacity{0.000000}%
\pgfsetdash{}{0pt}%
\pgfpathmoveto{\pgfqpoint{1.236867in}{0.733417in}}%
\pgfpathlineto{\pgfqpoint{1.301072in}{0.733417in}}%
\pgfpathlineto{\pgfqpoint{1.301072in}{0.733429in}}%
\pgfpathlineto{\pgfqpoint{1.236867in}{0.733429in}}%
\pgfpathlineto{\pgfqpoint{1.236867in}{0.733417in}}%
\pgfpathclose%
\pgfusepath{fill}%
\end{pgfscope}%
\begin{pgfscope}%
\pgfpathrectangle{\pgfqpoint{0.787430in}{0.733417in}}{\pgfqpoint{7.062570in}{2.555373in}}%
\pgfusepath{clip}%
\pgfsetbuttcap%
\pgfsetmiterjoin%
\definecolor{currentfill}{rgb}{0.121569,0.466667,0.705882}%
\pgfsetfillcolor{currentfill}%
\pgfsetlinewidth{0.000000pt}%
\definecolor{currentstroke}{rgb}{0.000000,0.000000,0.000000}%
\pgfsetstrokecolor{currentstroke}%
\pgfsetstrokeopacity{0.000000}%
\pgfsetdash{}{0pt}%
\pgfpathmoveto{\pgfqpoint{1.301072in}{0.733417in}}%
\pgfpathlineto{\pgfqpoint{1.365277in}{0.733417in}}%
\pgfpathlineto{\pgfqpoint{1.365277in}{0.733429in}}%
\pgfpathlineto{\pgfqpoint{1.301072in}{0.733429in}}%
\pgfpathlineto{\pgfqpoint{1.301072in}{0.733417in}}%
\pgfpathclose%
\pgfusepath{fill}%
\end{pgfscope}%
\begin{pgfscope}%
\pgfpathrectangle{\pgfqpoint{0.787430in}{0.733417in}}{\pgfqpoint{7.062570in}{2.555373in}}%
\pgfusepath{clip}%
\pgfsetbuttcap%
\pgfsetmiterjoin%
\definecolor{currentfill}{rgb}{0.121569,0.466667,0.705882}%
\pgfsetfillcolor{currentfill}%
\pgfsetlinewidth{0.000000pt}%
\definecolor{currentstroke}{rgb}{0.000000,0.000000,0.000000}%
\pgfsetstrokecolor{currentstroke}%
\pgfsetstrokeopacity{0.000000}%
\pgfsetdash{}{0pt}%
\pgfpathmoveto{\pgfqpoint{1.365277in}{0.733417in}}%
\pgfpathlineto{\pgfqpoint{1.429482in}{0.733417in}}%
\pgfpathlineto{\pgfqpoint{1.429482in}{0.733446in}}%
\pgfpathlineto{\pgfqpoint{1.365277in}{0.733446in}}%
\pgfpathlineto{\pgfqpoint{1.365277in}{0.733417in}}%
\pgfpathclose%
\pgfusepath{fill}%
\end{pgfscope}%
\begin{pgfscope}%
\pgfpathrectangle{\pgfqpoint{0.787430in}{0.733417in}}{\pgfqpoint{7.062570in}{2.555373in}}%
\pgfusepath{clip}%
\pgfsetbuttcap%
\pgfsetmiterjoin%
\definecolor{currentfill}{rgb}{0.121569,0.466667,0.705882}%
\pgfsetfillcolor{currentfill}%
\pgfsetlinewidth{0.000000pt}%
\definecolor{currentstroke}{rgb}{0.000000,0.000000,0.000000}%
\pgfsetstrokecolor{currentstroke}%
\pgfsetstrokeopacity{0.000000}%
\pgfsetdash{}{0pt}%
\pgfpathmoveto{\pgfqpoint{1.429482in}{0.733417in}}%
\pgfpathlineto{\pgfqpoint{1.493687in}{0.733417in}}%
\pgfpathlineto{\pgfqpoint{1.493687in}{0.733446in}}%
\pgfpathlineto{\pgfqpoint{1.429482in}{0.733446in}}%
\pgfpathlineto{\pgfqpoint{1.429482in}{0.733417in}}%
\pgfpathclose%
\pgfusepath{fill}%
\end{pgfscope}%
\begin{pgfscope}%
\pgfpathrectangle{\pgfqpoint{0.787430in}{0.733417in}}{\pgfqpoint{7.062570in}{2.555373in}}%
\pgfusepath{clip}%
\pgfsetbuttcap%
\pgfsetmiterjoin%
\definecolor{currentfill}{rgb}{0.121569,0.466667,0.705882}%
\pgfsetfillcolor{currentfill}%
\pgfsetlinewidth{0.000000pt}%
\definecolor{currentstroke}{rgb}{0.000000,0.000000,0.000000}%
\pgfsetstrokecolor{currentstroke}%
\pgfsetstrokeopacity{0.000000}%
\pgfsetdash{}{0pt}%
\pgfpathmoveto{\pgfqpoint{1.493687in}{0.733417in}}%
\pgfpathlineto{\pgfqpoint{1.557892in}{0.733417in}}%
\pgfpathlineto{\pgfqpoint{1.557892in}{0.733481in}}%
\pgfpathlineto{\pgfqpoint{1.493687in}{0.733481in}}%
\pgfpathlineto{\pgfqpoint{1.493687in}{0.733417in}}%
\pgfpathclose%
\pgfusepath{fill}%
\end{pgfscope}%
\begin{pgfscope}%
\pgfpathrectangle{\pgfqpoint{0.787430in}{0.733417in}}{\pgfqpoint{7.062570in}{2.555373in}}%
\pgfusepath{clip}%
\pgfsetbuttcap%
\pgfsetmiterjoin%
\definecolor{currentfill}{rgb}{0.121569,0.466667,0.705882}%
\pgfsetfillcolor{currentfill}%
\pgfsetlinewidth{0.000000pt}%
\definecolor{currentstroke}{rgb}{0.000000,0.000000,0.000000}%
\pgfsetstrokecolor{currentstroke}%
\pgfsetstrokeopacity{0.000000}%
\pgfsetdash{}{0pt}%
\pgfpathmoveto{\pgfqpoint{1.557892in}{0.733417in}}%
\pgfpathlineto{\pgfqpoint{1.622098in}{0.733417in}}%
\pgfpathlineto{\pgfqpoint{1.622098in}{0.733481in}}%
\pgfpathlineto{\pgfqpoint{1.557892in}{0.733481in}}%
\pgfpathlineto{\pgfqpoint{1.557892in}{0.733417in}}%
\pgfpathclose%
\pgfusepath{fill}%
\end{pgfscope}%
\begin{pgfscope}%
\pgfpathrectangle{\pgfqpoint{0.787430in}{0.733417in}}{\pgfqpoint{7.062570in}{2.555373in}}%
\pgfusepath{clip}%
\pgfsetbuttcap%
\pgfsetmiterjoin%
\definecolor{currentfill}{rgb}{0.121569,0.466667,0.705882}%
\pgfsetfillcolor{currentfill}%
\pgfsetlinewidth{0.000000pt}%
\definecolor{currentstroke}{rgb}{0.000000,0.000000,0.000000}%
\pgfsetstrokecolor{currentstroke}%
\pgfsetstrokeopacity{0.000000}%
\pgfsetdash{}{0pt}%
\pgfpathmoveto{\pgfqpoint{1.622098in}{0.733417in}}%
\pgfpathlineto{\pgfqpoint{1.686303in}{0.733417in}}%
\pgfpathlineto{\pgfqpoint{1.686303in}{0.733528in}}%
\pgfpathlineto{\pgfqpoint{1.622098in}{0.733528in}}%
\pgfpathlineto{\pgfqpoint{1.622098in}{0.733417in}}%
\pgfpathclose%
\pgfusepath{fill}%
\end{pgfscope}%
\begin{pgfscope}%
\pgfpathrectangle{\pgfqpoint{0.787430in}{0.733417in}}{\pgfqpoint{7.062570in}{2.555373in}}%
\pgfusepath{clip}%
\pgfsetbuttcap%
\pgfsetmiterjoin%
\definecolor{currentfill}{rgb}{0.121569,0.466667,0.705882}%
\pgfsetfillcolor{currentfill}%
\pgfsetlinewidth{0.000000pt}%
\definecolor{currentstroke}{rgb}{0.000000,0.000000,0.000000}%
\pgfsetstrokecolor{currentstroke}%
\pgfsetstrokeopacity{0.000000}%
\pgfsetdash{}{0pt}%
\pgfpathmoveto{\pgfqpoint{1.686303in}{0.733417in}}%
\pgfpathlineto{\pgfqpoint{1.750508in}{0.733417in}}%
\pgfpathlineto{\pgfqpoint{1.750508in}{0.733610in}}%
\pgfpathlineto{\pgfqpoint{1.686303in}{0.733610in}}%
\pgfpathlineto{\pgfqpoint{1.686303in}{0.733417in}}%
\pgfpathclose%
\pgfusepath{fill}%
\end{pgfscope}%
\begin{pgfscope}%
\pgfpathrectangle{\pgfqpoint{0.787430in}{0.733417in}}{\pgfqpoint{7.062570in}{2.555373in}}%
\pgfusepath{clip}%
\pgfsetbuttcap%
\pgfsetmiterjoin%
\definecolor{currentfill}{rgb}{0.121569,0.466667,0.705882}%
\pgfsetfillcolor{currentfill}%
\pgfsetlinewidth{0.000000pt}%
\definecolor{currentstroke}{rgb}{0.000000,0.000000,0.000000}%
\pgfsetstrokecolor{currentstroke}%
\pgfsetstrokeopacity{0.000000}%
\pgfsetdash{}{0pt}%
\pgfpathmoveto{\pgfqpoint{1.750508in}{0.733417in}}%
\pgfpathlineto{\pgfqpoint{1.814713in}{0.733417in}}%
\pgfpathlineto{\pgfqpoint{1.814713in}{0.733727in}}%
\pgfpathlineto{\pgfqpoint{1.750508in}{0.733727in}}%
\pgfpathlineto{\pgfqpoint{1.750508in}{0.733417in}}%
\pgfpathclose%
\pgfusepath{fill}%
\end{pgfscope}%
\begin{pgfscope}%
\pgfpathrectangle{\pgfqpoint{0.787430in}{0.733417in}}{\pgfqpoint{7.062570in}{2.555373in}}%
\pgfusepath{clip}%
\pgfsetbuttcap%
\pgfsetmiterjoin%
\definecolor{currentfill}{rgb}{0.121569,0.466667,0.705882}%
\pgfsetfillcolor{currentfill}%
\pgfsetlinewidth{0.000000pt}%
\definecolor{currentstroke}{rgb}{0.000000,0.000000,0.000000}%
\pgfsetstrokecolor{currentstroke}%
\pgfsetstrokeopacity{0.000000}%
\pgfsetdash{}{0pt}%
\pgfpathmoveto{\pgfqpoint{1.814713in}{0.733417in}}%
\pgfpathlineto{\pgfqpoint{1.878918in}{0.733417in}}%
\pgfpathlineto{\pgfqpoint{1.878918in}{0.733867in}}%
\pgfpathlineto{\pgfqpoint{1.814713in}{0.733867in}}%
\pgfpathlineto{\pgfqpoint{1.814713in}{0.733417in}}%
\pgfpathclose%
\pgfusepath{fill}%
\end{pgfscope}%
\begin{pgfscope}%
\pgfpathrectangle{\pgfqpoint{0.787430in}{0.733417in}}{\pgfqpoint{7.062570in}{2.555373in}}%
\pgfusepath{clip}%
\pgfsetbuttcap%
\pgfsetmiterjoin%
\definecolor{currentfill}{rgb}{0.121569,0.466667,0.705882}%
\pgfsetfillcolor{currentfill}%
\pgfsetlinewidth{0.000000pt}%
\definecolor{currentstroke}{rgb}{0.000000,0.000000,0.000000}%
\pgfsetstrokecolor{currentstroke}%
\pgfsetstrokeopacity{0.000000}%
\pgfsetdash{}{0pt}%
\pgfpathmoveto{\pgfqpoint{1.878918in}{0.733417in}}%
\pgfpathlineto{\pgfqpoint{1.943124in}{0.733417in}}%
\pgfpathlineto{\pgfqpoint{1.943124in}{0.734247in}}%
\pgfpathlineto{\pgfqpoint{1.878918in}{0.734247in}}%
\pgfpathlineto{\pgfqpoint{1.878918in}{0.733417in}}%
\pgfpathclose%
\pgfusepath{fill}%
\end{pgfscope}%
\begin{pgfscope}%
\pgfpathrectangle{\pgfqpoint{0.787430in}{0.733417in}}{\pgfqpoint{7.062570in}{2.555373in}}%
\pgfusepath{clip}%
\pgfsetbuttcap%
\pgfsetmiterjoin%
\definecolor{currentfill}{rgb}{0.121569,0.466667,0.705882}%
\pgfsetfillcolor{currentfill}%
\pgfsetlinewidth{0.000000pt}%
\definecolor{currentstroke}{rgb}{0.000000,0.000000,0.000000}%
\pgfsetstrokecolor{currentstroke}%
\pgfsetstrokeopacity{0.000000}%
\pgfsetdash{}{0pt}%
\pgfpathmoveto{\pgfqpoint{1.943124in}{0.733417in}}%
\pgfpathlineto{\pgfqpoint{2.007329in}{0.733417in}}%
\pgfpathlineto{\pgfqpoint{2.007329in}{0.734632in}}%
\pgfpathlineto{\pgfqpoint{1.943124in}{0.734632in}}%
\pgfpathlineto{\pgfqpoint{1.943124in}{0.733417in}}%
\pgfpathclose%
\pgfusepath{fill}%
\end{pgfscope}%
\begin{pgfscope}%
\pgfpathrectangle{\pgfqpoint{0.787430in}{0.733417in}}{\pgfqpoint{7.062570in}{2.555373in}}%
\pgfusepath{clip}%
\pgfsetbuttcap%
\pgfsetmiterjoin%
\definecolor{currentfill}{rgb}{0.121569,0.466667,0.705882}%
\pgfsetfillcolor{currentfill}%
\pgfsetlinewidth{0.000000pt}%
\definecolor{currentstroke}{rgb}{0.000000,0.000000,0.000000}%
\pgfsetstrokecolor{currentstroke}%
\pgfsetstrokeopacity{0.000000}%
\pgfsetdash{}{0pt}%
\pgfpathmoveto{\pgfqpoint{2.007329in}{0.733417in}}%
\pgfpathlineto{\pgfqpoint{2.071534in}{0.733417in}}%
\pgfpathlineto{\pgfqpoint{2.071534in}{0.735421in}}%
\pgfpathlineto{\pgfqpoint{2.007329in}{0.735421in}}%
\pgfpathlineto{\pgfqpoint{2.007329in}{0.733417in}}%
\pgfpathclose%
\pgfusepath{fill}%
\end{pgfscope}%
\begin{pgfscope}%
\pgfpathrectangle{\pgfqpoint{0.787430in}{0.733417in}}{\pgfqpoint{7.062570in}{2.555373in}}%
\pgfusepath{clip}%
\pgfsetbuttcap%
\pgfsetmiterjoin%
\definecolor{currentfill}{rgb}{0.121569,0.466667,0.705882}%
\pgfsetfillcolor{currentfill}%
\pgfsetlinewidth{0.000000pt}%
\definecolor{currentstroke}{rgb}{0.000000,0.000000,0.000000}%
\pgfsetstrokecolor{currentstroke}%
\pgfsetstrokeopacity{0.000000}%
\pgfsetdash{}{0pt}%
\pgfpathmoveto{\pgfqpoint{2.071534in}{0.733417in}}%
\pgfpathlineto{\pgfqpoint{2.135739in}{0.733417in}}%
\pgfpathlineto{\pgfqpoint{2.135739in}{0.735883in}}%
\pgfpathlineto{\pgfqpoint{2.071534in}{0.735883in}}%
\pgfpathlineto{\pgfqpoint{2.071534in}{0.733417in}}%
\pgfpathclose%
\pgfusepath{fill}%
\end{pgfscope}%
\begin{pgfscope}%
\pgfpathrectangle{\pgfqpoint{0.787430in}{0.733417in}}{\pgfqpoint{7.062570in}{2.555373in}}%
\pgfusepath{clip}%
\pgfsetbuttcap%
\pgfsetmiterjoin%
\definecolor{currentfill}{rgb}{0.121569,0.466667,0.705882}%
\pgfsetfillcolor{currentfill}%
\pgfsetlinewidth{0.000000pt}%
\definecolor{currentstroke}{rgb}{0.000000,0.000000,0.000000}%
\pgfsetstrokecolor{currentstroke}%
\pgfsetstrokeopacity{0.000000}%
\pgfsetdash{}{0pt}%
\pgfpathmoveto{\pgfqpoint{2.135739in}{0.733417in}}%
\pgfpathlineto{\pgfqpoint{2.199944in}{0.733417in}}%
\pgfpathlineto{\pgfqpoint{2.199944in}{0.737238in}}%
\pgfpathlineto{\pgfqpoint{2.135739in}{0.737238in}}%
\pgfpathlineto{\pgfqpoint{2.135739in}{0.733417in}}%
\pgfpathclose%
\pgfusepath{fill}%
\end{pgfscope}%
\begin{pgfscope}%
\pgfpathrectangle{\pgfqpoint{0.787430in}{0.733417in}}{\pgfqpoint{7.062570in}{2.555373in}}%
\pgfusepath{clip}%
\pgfsetbuttcap%
\pgfsetmiterjoin%
\definecolor{currentfill}{rgb}{0.121569,0.466667,0.705882}%
\pgfsetfillcolor{currentfill}%
\pgfsetlinewidth{0.000000pt}%
\definecolor{currentstroke}{rgb}{0.000000,0.000000,0.000000}%
\pgfsetstrokecolor{currentstroke}%
\pgfsetstrokeopacity{0.000000}%
\pgfsetdash{}{0pt}%
\pgfpathmoveto{\pgfqpoint{2.199944in}{0.733417in}}%
\pgfpathlineto{\pgfqpoint{2.264149in}{0.733417in}}%
\pgfpathlineto{\pgfqpoint{2.264149in}{0.738728in}}%
\pgfpathlineto{\pgfqpoint{2.199944in}{0.738728in}}%
\pgfpathlineto{\pgfqpoint{2.199944in}{0.733417in}}%
\pgfpathclose%
\pgfusepath{fill}%
\end{pgfscope}%
\begin{pgfscope}%
\pgfpathrectangle{\pgfqpoint{0.787430in}{0.733417in}}{\pgfqpoint{7.062570in}{2.555373in}}%
\pgfusepath{clip}%
\pgfsetbuttcap%
\pgfsetmiterjoin%
\definecolor{currentfill}{rgb}{0.121569,0.466667,0.705882}%
\pgfsetfillcolor{currentfill}%
\pgfsetlinewidth{0.000000pt}%
\definecolor{currentstroke}{rgb}{0.000000,0.000000,0.000000}%
\pgfsetstrokecolor{currentstroke}%
\pgfsetstrokeopacity{0.000000}%
\pgfsetdash{}{0pt}%
\pgfpathmoveto{\pgfqpoint{2.264149in}{0.733417in}}%
\pgfpathlineto{\pgfqpoint{2.328355in}{0.733417in}}%
\pgfpathlineto{\pgfqpoint{2.328355in}{0.741164in}}%
\pgfpathlineto{\pgfqpoint{2.264149in}{0.741164in}}%
\pgfpathlineto{\pgfqpoint{2.264149in}{0.733417in}}%
\pgfpathclose%
\pgfusepath{fill}%
\end{pgfscope}%
\begin{pgfscope}%
\pgfpathrectangle{\pgfqpoint{0.787430in}{0.733417in}}{\pgfqpoint{7.062570in}{2.555373in}}%
\pgfusepath{clip}%
\pgfsetbuttcap%
\pgfsetmiterjoin%
\definecolor{currentfill}{rgb}{0.121569,0.466667,0.705882}%
\pgfsetfillcolor{currentfill}%
\pgfsetlinewidth{0.000000pt}%
\definecolor{currentstroke}{rgb}{0.000000,0.000000,0.000000}%
\pgfsetstrokecolor{currentstroke}%
\pgfsetstrokeopacity{0.000000}%
\pgfsetdash{}{0pt}%
\pgfpathmoveto{\pgfqpoint{2.328355in}{0.733417in}}%
\pgfpathlineto{\pgfqpoint{2.392560in}{0.733417in}}%
\pgfpathlineto{\pgfqpoint{2.392560in}{0.743810in}}%
\pgfpathlineto{\pgfqpoint{2.328355in}{0.743810in}}%
\pgfpathlineto{\pgfqpoint{2.328355in}{0.733417in}}%
\pgfpathclose%
\pgfusepath{fill}%
\end{pgfscope}%
\begin{pgfscope}%
\pgfpathrectangle{\pgfqpoint{0.787430in}{0.733417in}}{\pgfqpoint{7.062570in}{2.555373in}}%
\pgfusepath{clip}%
\pgfsetbuttcap%
\pgfsetmiterjoin%
\definecolor{currentfill}{rgb}{0.121569,0.466667,0.705882}%
\pgfsetfillcolor{currentfill}%
\pgfsetlinewidth{0.000000pt}%
\definecolor{currentstroke}{rgb}{0.000000,0.000000,0.000000}%
\pgfsetstrokecolor{currentstroke}%
\pgfsetstrokeopacity{0.000000}%
\pgfsetdash{}{0pt}%
\pgfpathmoveto{\pgfqpoint{2.392560in}{0.733417in}}%
\pgfpathlineto{\pgfqpoint{2.456765in}{0.733417in}}%
\pgfpathlineto{\pgfqpoint{2.456765in}{0.748922in}}%
\pgfpathlineto{\pgfqpoint{2.392560in}{0.748922in}}%
\pgfpathlineto{\pgfqpoint{2.392560in}{0.733417in}}%
\pgfpathclose%
\pgfusepath{fill}%
\end{pgfscope}%
\begin{pgfscope}%
\pgfpathrectangle{\pgfqpoint{0.787430in}{0.733417in}}{\pgfqpoint{7.062570in}{2.555373in}}%
\pgfusepath{clip}%
\pgfsetbuttcap%
\pgfsetmiterjoin%
\definecolor{currentfill}{rgb}{0.121569,0.466667,0.705882}%
\pgfsetfillcolor{currentfill}%
\pgfsetlinewidth{0.000000pt}%
\definecolor{currentstroke}{rgb}{0.000000,0.000000,0.000000}%
\pgfsetstrokecolor{currentstroke}%
\pgfsetstrokeopacity{0.000000}%
\pgfsetdash{}{0pt}%
\pgfpathmoveto{\pgfqpoint{2.456765in}{0.733417in}}%
\pgfpathlineto{\pgfqpoint{2.520970in}{0.733417in}}%
\pgfpathlineto{\pgfqpoint{2.520970in}{0.754723in}}%
\pgfpathlineto{\pgfqpoint{2.456765in}{0.754723in}}%
\pgfpathlineto{\pgfqpoint{2.456765in}{0.733417in}}%
\pgfpathclose%
\pgfusepath{fill}%
\end{pgfscope}%
\begin{pgfscope}%
\pgfpathrectangle{\pgfqpoint{0.787430in}{0.733417in}}{\pgfqpoint{7.062570in}{2.555373in}}%
\pgfusepath{clip}%
\pgfsetbuttcap%
\pgfsetmiterjoin%
\definecolor{currentfill}{rgb}{0.121569,0.466667,0.705882}%
\pgfsetfillcolor{currentfill}%
\pgfsetlinewidth{0.000000pt}%
\definecolor{currentstroke}{rgb}{0.000000,0.000000,0.000000}%
\pgfsetstrokecolor{currentstroke}%
\pgfsetstrokeopacity{0.000000}%
\pgfsetdash{}{0pt}%
\pgfpathmoveto{\pgfqpoint{2.520970in}{0.733417in}}%
\pgfpathlineto{\pgfqpoint{2.585175in}{0.733417in}}%
\pgfpathlineto{\pgfqpoint{2.585175in}{0.762984in}}%
\pgfpathlineto{\pgfqpoint{2.520970in}{0.762984in}}%
\pgfpathlineto{\pgfqpoint{2.520970in}{0.733417in}}%
\pgfpathclose%
\pgfusepath{fill}%
\end{pgfscope}%
\begin{pgfscope}%
\pgfpathrectangle{\pgfqpoint{0.787430in}{0.733417in}}{\pgfqpoint{7.062570in}{2.555373in}}%
\pgfusepath{clip}%
\pgfsetbuttcap%
\pgfsetmiterjoin%
\definecolor{currentfill}{rgb}{0.121569,0.466667,0.705882}%
\pgfsetfillcolor{currentfill}%
\pgfsetlinewidth{0.000000pt}%
\definecolor{currentstroke}{rgb}{0.000000,0.000000,0.000000}%
\pgfsetstrokecolor{currentstroke}%
\pgfsetstrokeopacity{0.000000}%
\pgfsetdash{}{0pt}%
\pgfpathmoveto{\pgfqpoint{2.585175in}{0.733417in}}%
\pgfpathlineto{\pgfqpoint{2.649380in}{0.733417in}}%
\pgfpathlineto{\pgfqpoint{2.649380in}{0.772974in}}%
\pgfpathlineto{\pgfqpoint{2.585175in}{0.772974in}}%
\pgfpathlineto{\pgfqpoint{2.585175in}{0.733417in}}%
\pgfpathclose%
\pgfusepath{fill}%
\end{pgfscope}%
\begin{pgfscope}%
\pgfpathrectangle{\pgfqpoint{0.787430in}{0.733417in}}{\pgfqpoint{7.062570in}{2.555373in}}%
\pgfusepath{clip}%
\pgfsetbuttcap%
\pgfsetmiterjoin%
\definecolor{currentfill}{rgb}{0.121569,0.466667,0.705882}%
\pgfsetfillcolor{currentfill}%
\pgfsetlinewidth{0.000000pt}%
\definecolor{currentstroke}{rgb}{0.000000,0.000000,0.000000}%
\pgfsetstrokecolor{currentstroke}%
\pgfsetstrokeopacity{0.000000}%
\pgfsetdash{}{0pt}%
\pgfpathmoveto{\pgfqpoint{2.649380in}{0.733417in}}%
\pgfpathlineto{\pgfqpoint{2.713586in}{0.733417in}}%
\pgfpathlineto{\pgfqpoint{2.713586in}{0.785388in}}%
\pgfpathlineto{\pgfqpoint{2.649380in}{0.785388in}}%
\pgfpathlineto{\pgfqpoint{2.649380in}{0.733417in}}%
\pgfpathclose%
\pgfusepath{fill}%
\end{pgfscope}%
\begin{pgfscope}%
\pgfpathrectangle{\pgfqpoint{0.787430in}{0.733417in}}{\pgfqpoint{7.062570in}{2.555373in}}%
\pgfusepath{clip}%
\pgfsetbuttcap%
\pgfsetmiterjoin%
\definecolor{currentfill}{rgb}{0.121569,0.466667,0.705882}%
\pgfsetfillcolor{currentfill}%
\pgfsetlinewidth{0.000000pt}%
\definecolor{currentstroke}{rgb}{0.000000,0.000000,0.000000}%
\pgfsetstrokecolor{currentstroke}%
\pgfsetstrokeopacity{0.000000}%
\pgfsetdash{}{0pt}%
\pgfpathmoveto{\pgfqpoint{2.713586in}{0.733417in}}%
\pgfpathlineto{\pgfqpoint{2.777791in}{0.733417in}}%
\pgfpathlineto{\pgfqpoint{2.777791in}{0.803323in}}%
\pgfpathlineto{\pgfqpoint{2.713586in}{0.803323in}}%
\pgfpathlineto{\pgfqpoint{2.713586in}{0.733417in}}%
\pgfpathclose%
\pgfusepath{fill}%
\end{pgfscope}%
\begin{pgfscope}%
\pgfpathrectangle{\pgfqpoint{0.787430in}{0.733417in}}{\pgfqpoint{7.062570in}{2.555373in}}%
\pgfusepath{clip}%
\pgfsetbuttcap%
\pgfsetmiterjoin%
\definecolor{currentfill}{rgb}{0.121569,0.466667,0.705882}%
\pgfsetfillcolor{currentfill}%
\pgfsetlinewidth{0.000000pt}%
\definecolor{currentstroke}{rgb}{0.000000,0.000000,0.000000}%
\pgfsetstrokecolor{currentstroke}%
\pgfsetstrokeopacity{0.000000}%
\pgfsetdash{}{0pt}%
\pgfpathmoveto{\pgfqpoint{2.777791in}{0.733417in}}%
\pgfpathlineto{\pgfqpoint{2.841996in}{0.733417in}}%
\pgfpathlineto{\pgfqpoint{2.841996in}{0.826084in}}%
\pgfpathlineto{\pgfqpoint{2.777791in}{0.826084in}}%
\pgfpathlineto{\pgfqpoint{2.777791in}{0.733417in}}%
\pgfpathclose%
\pgfusepath{fill}%
\end{pgfscope}%
\begin{pgfscope}%
\pgfpathrectangle{\pgfqpoint{0.787430in}{0.733417in}}{\pgfqpoint{7.062570in}{2.555373in}}%
\pgfusepath{clip}%
\pgfsetbuttcap%
\pgfsetmiterjoin%
\definecolor{currentfill}{rgb}{0.121569,0.466667,0.705882}%
\pgfsetfillcolor{currentfill}%
\pgfsetlinewidth{0.000000pt}%
\definecolor{currentstroke}{rgb}{0.000000,0.000000,0.000000}%
\pgfsetstrokecolor{currentstroke}%
\pgfsetstrokeopacity{0.000000}%
\pgfsetdash{}{0pt}%
\pgfpathmoveto{\pgfqpoint{2.841996in}{0.733417in}}%
\pgfpathlineto{\pgfqpoint{2.906201in}{0.733417in}}%
\pgfpathlineto{\pgfqpoint{2.906201in}{0.852770in}}%
\pgfpathlineto{\pgfqpoint{2.841996in}{0.852770in}}%
\pgfpathlineto{\pgfqpoint{2.841996in}{0.733417in}}%
\pgfpathclose%
\pgfusepath{fill}%
\end{pgfscope}%
\begin{pgfscope}%
\pgfpathrectangle{\pgfqpoint{0.787430in}{0.733417in}}{\pgfqpoint{7.062570in}{2.555373in}}%
\pgfusepath{clip}%
\pgfsetbuttcap%
\pgfsetmiterjoin%
\definecolor{currentfill}{rgb}{0.121569,0.466667,0.705882}%
\pgfsetfillcolor{currentfill}%
\pgfsetlinewidth{0.000000pt}%
\definecolor{currentstroke}{rgb}{0.000000,0.000000,0.000000}%
\pgfsetstrokecolor{currentstroke}%
\pgfsetstrokeopacity{0.000000}%
\pgfsetdash{}{0pt}%
\pgfpathmoveto{\pgfqpoint{2.906201in}{0.733417in}}%
\pgfpathlineto{\pgfqpoint{2.970406in}{0.733417in}}%
\pgfpathlineto{\pgfqpoint{2.970406in}{0.888424in}}%
\pgfpathlineto{\pgfqpoint{2.906201in}{0.888424in}}%
\pgfpathlineto{\pgfqpoint{2.906201in}{0.733417in}}%
\pgfpathclose%
\pgfusepath{fill}%
\end{pgfscope}%
\begin{pgfscope}%
\pgfpathrectangle{\pgfqpoint{0.787430in}{0.733417in}}{\pgfqpoint{7.062570in}{2.555373in}}%
\pgfusepath{clip}%
\pgfsetbuttcap%
\pgfsetmiterjoin%
\definecolor{currentfill}{rgb}{0.121569,0.466667,0.705882}%
\pgfsetfillcolor{currentfill}%
\pgfsetlinewidth{0.000000pt}%
\definecolor{currentstroke}{rgb}{0.000000,0.000000,0.000000}%
\pgfsetstrokecolor{currentstroke}%
\pgfsetstrokeopacity{0.000000}%
\pgfsetdash{}{0pt}%
\pgfpathmoveto{\pgfqpoint{2.970406in}{0.733417in}}%
\pgfpathlineto{\pgfqpoint{3.034612in}{0.733417in}}%
\pgfpathlineto{\pgfqpoint{3.034612in}{0.931719in}}%
\pgfpathlineto{\pgfqpoint{2.970406in}{0.931719in}}%
\pgfpathlineto{\pgfqpoint{2.970406in}{0.733417in}}%
\pgfpathclose%
\pgfusepath{fill}%
\end{pgfscope}%
\begin{pgfscope}%
\pgfpathrectangle{\pgfqpoint{0.787430in}{0.733417in}}{\pgfqpoint{7.062570in}{2.555373in}}%
\pgfusepath{clip}%
\pgfsetbuttcap%
\pgfsetmiterjoin%
\definecolor{currentfill}{rgb}{0.121569,0.466667,0.705882}%
\pgfsetfillcolor{currentfill}%
\pgfsetlinewidth{0.000000pt}%
\definecolor{currentstroke}{rgb}{0.000000,0.000000,0.000000}%
\pgfsetstrokecolor{currentstroke}%
\pgfsetstrokeopacity{0.000000}%
\pgfsetdash{}{0pt}%
\pgfpathmoveto{\pgfqpoint{3.034612in}{0.733417in}}%
\pgfpathlineto{\pgfqpoint{3.098817in}{0.733417in}}%
\pgfpathlineto{\pgfqpoint{3.098817in}{0.976650in}}%
\pgfpathlineto{\pgfqpoint{3.034612in}{0.976650in}}%
\pgfpathlineto{\pgfqpoint{3.034612in}{0.733417in}}%
\pgfpathclose%
\pgfusepath{fill}%
\end{pgfscope}%
\begin{pgfscope}%
\pgfpathrectangle{\pgfqpoint{0.787430in}{0.733417in}}{\pgfqpoint{7.062570in}{2.555373in}}%
\pgfusepath{clip}%
\pgfsetbuttcap%
\pgfsetmiterjoin%
\definecolor{currentfill}{rgb}{0.121569,0.466667,0.705882}%
\pgfsetfillcolor{currentfill}%
\pgfsetlinewidth{0.000000pt}%
\definecolor{currentstroke}{rgb}{0.000000,0.000000,0.000000}%
\pgfsetstrokecolor{currentstroke}%
\pgfsetstrokeopacity{0.000000}%
\pgfsetdash{}{0pt}%
\pgfpathmoveto{\pgfqpoint{3.098817in}{0.733417in}}%
\pgfpathlineto{\pgfqpoint{3.163022in}{0.733417in}}%
\pgfpathlineto{\pgfqpoint{3.163022in}{1.038173in}}%
\pgfpathlineto{\pgfqpoint{3.098817in}{1.038173in}}%
\pgfpathlineto{\pgfqpoint{3.098817in}{0.733417in}}%
\pgfpathclose%
\pgfusepath{fill}%
\end{pgfscope}%
\begin{pgfscope}%
\pgfpathrectangle{\pgfqpoint{0.787430in}{0.733417in}}{\pgfqpoint{7.062570in}{2.555373in}}%
\pgfusepath{clip}%
\pgfsetbuttcap%
\pgfsetmiterjoin%
\definecolor{currentfill}{rgb}{0.121569,0.466667,0.705882}%
\pgfsetfillcolor{currentfill}%
\pgfsetlinewidth{0.000000pt}%
\definecolor{currentstroke}{rgb}{0.000000,0.000000,0.000000}%
\pgfsetstrokecolor{currentstroke}%
\pgfsetstrokeopacity{0.000000}%
\pgfsetdash{}{0pt}%
\pgfpathmoveto{\pgfqpoint{3.163022in}{0.733417in}}%
\pgfpathlineto{\pgfqpoint{3.227227in}{0.733417in}}%
\pgfpathlineto{\pgfqpoint{3.227227in}{1.110561in}}%
\pgfpathlineto{\pgfqpoint{3.163022in}{1.110561in}}%
\pgfpathlineto{\pgfqpoint{3.163022in}{0.733417in}}%
\pgfpathclose%
\pgfusepath{fill}%
\end{pgfscope}%
\begin{pgfscope}%
\pgfpathrectangle{\pgfqpoint{0.787430in}{0.733417in}}{\pgfqpoint{7.062570in}{2.555373in}}%
\pgfusepath{clip}%
\pgfsetbuttcap%
\pgfsetmiterjoin%
\definecolor{currentfill}{rgb}{0.121569,0.466667,0.705882}%
\pgfsetfillcolor{currentfill}%
\pgfsetlinewidth{0.000000pt}%
\definecolor{currentstroke}{rgb}{0.000000,0.000000,0.000000}%
\pgfsetstrokecolor{currentstroke}%
\pgfsetstrokeopacity{0.000000}%
\pgfsetdash{}{0pt}%
\pgfpathmoveto{\pgfqpoint{3.227227in}{0.733417in}}%
\pgfpathlineto{\pgfqpoint{3.291432in}{0.733417in}}%
\pgfpathlineto{\pgfqpoint{3.291432in}{1.191281in}}%
\pgfpathlineto{\pgfqpoint{3.227227in}{1.191281in}}%
\pgfpathlineto{\pgfqpoint{3.227227in}{0.733417in}}%
\pgfpathclose%
\pgfusepath{fill}%
\end{pgfscope}%
\begin{pgfscope}%
\pgfpathrectangle{\pgfqpoint{0.787430in}{0.733417in}}{\pgfqpoint{7.062570in}{2.555373in}}%
\pgfusepath{clip}%
\pgfsetbuttcap%
\pgfsetmiterjoin%
\definecolor{currentfill}{rgb}{0.121569,0.466667,0.705882}%
\pgfsetfillcolor{currentfill}%
\pgfsetlinewidth{0.000000pt}%
\definecolor{currentstroke}{rgb}{0.000000,0.000000,0.000000}%
\pgfsetstrokecolor{currentstroke}%
\pgfsetstrokeopacity{0.000000}%
\pgfsetdash{}{0pt}%
\pgfpathmoveto{\pgfqpoint{3.291432in}{0.733417in}}%
\pgfpathlineto{\pgfqpoint{3.355637in}{0.733417in}}%
\pgfpathlineto{\pgfqpoint{3.355637in}{1.287336in}}%
\pgfpathlineto{\pgfqpoint{3.291432in}{1.287336in}}%
\pgfpathlineto{\pgfqpoint{3.291432in}{0.733417in}}%
\pgfpathclose%
\pgfusepath{fill}%
\end{pgfscope}%
\begin{pgfscope}%
\pgfpathrectangle{\pgfqpoint{0.787430in}{0.733417in}}{\pgfqpoint{7.062570in}{2.555373in}}%
\pgfusepath{clip}%
\pgfsetbuttcap%
\pgfsetmiterjoin%
\definecolor{currentfill}{rgb}{0.121569,0.466667,0.705882}%
\pgfsetfillcolor{currentfill}%
\pgfsetlinewidth{0.000000pt}%
\definecolor{currentstroke}{rgb}{0.000000,0.000000,0.000000}%
\pgfsetstrokecolor{currentstroke}%
\pgfsetstrokeopacity{0.000000}%
\pgfsetdash{}{0pt}%
\pgfpathmoveto{\pgfqpoint{3.355637in}{0.733417in}}%
\pgfpathlineto{\pgfqpoint{3.419843in}{0.733417in}}%
\pgfpathlineto{\pgfqpoint{3.419843in}{1.391774in}}%
\pgfpathlineto{\pgfqpoint{3.355637in}{1.391774in}}%
\pgfpathlineto{\pgfqpoint{3.355637in}{0.733417in}}%
\pgfpathclose%
\pgfusepath{fill}%
\end{pgfscope}%
\begin{pgfscope}%
\pgfpathrectangle{\pgfqpoint{0.787430in}{0.733417in}}{\pgfqpoint{7.062570in}{2.555373in}}%
\pgfusepath{clip}%
\pgfsetbuttcap%
\pgfsetmiterjoin%
\definecolor{currentfill}{rgb}{0.121569,0.466667,0.705882}%
\pgfsetfillcolor{currentfill}%
\pgfsetlinewidth{0.000000pt}%
\definecolor{currentstroke}{rgb}{0.000000,0.000000,0.000000}%
\pgfsetstrokecolor{currentstroke}%
\pgfsetstrokeopacity{0.000000}%
\pgfsetdash{}{0pt}%
\pgfpathmoveto{\pgfqpoint{3.419843in}{0.733417in}}%
\pgfpathlineto{\pgfqpoint{3.484048in}{0.733417in}}%
\pgfpathlineto{\pgfqpoint{3.484048in}{1.505664in}}%
\pgfpathlineto{\pgfqpoint{3.419843in}{1.505664in}}%
\pgfpathlineto{\pgfqpoint{3.419843in}{0.733417in}}%
\pgfpathclose%
\pgfusepath{fill}%
\end{pgfscope}%
\begin{pgfscope}%
\pgfpathrectangle{\pgfqpoint{0.787430in}{0.733417in}}{\pgfqpoint{7.062570in}{2.555373in}}%
\pgfusepath{clip}%
\pgfsetbuttcap%
\pgfsetmiterjoin%
\definecolor{currentfill}{rgb}{0.121569,0.466667,0.705882}%
\pgfsetfillcolor{currentfill}%
\pgfsetlinewidth{0.000000pt}%
\definecolor{currentstroke}{rgb}{0.000000,0.000000,0.000000}%
\pgfsetstrokecolor{currentstroke}%
\pgfsetstrokeopacity{0.000000}%
\pgfsetdash{}{0pt}%
\pgfpathmoveto{\pgfqpoint{3.484048in}{0.733417in}}%
\pgfpathlineto{\pgfqpoint{3.548253in}{0.733417in}}%
\pgfpathlineto{\pgfqpoint{3.548253in}{1.637554in}}%
\pgfpathlineto{\pgfqpoint{3.484048in}{1.637554in}}%
\pgfpathlineto{\pgfqpoint{3.484048in}{0.733417in}}%
\pgfpathclose%
\pgfusepath{fill}%
\end{pgfscope}%
\begin{pgfscope}%
\pgfpathrectangle{\pgfqpoint{0.787430in}{0.733417in}}{\pgfqpoint{7.062570in}{2.555373in}}%
\pgfusepath{clip}%
\pgfsetbuttcap%
\pgfsetmiterjoin%
\definecolor{currentfill}{rgb}{0.121569,0.466667,0.705882}%
\pgfsetfillcolor{currentfill}%
\pgfsetlinewidth{0.000000pt}%
\definecolor{currentstroke}{rgb}{0.000000,0.000000,0.000000}%
\pgfsetstrokecolor{currentstroke}%
\pgfsetstrokeopacity{0.000000}%
\pgfsetdash{}{0pt}%
\pgfpathmoveto{\pgfqpoint{3.548253in}{0.733417in}}%
\pgfpathlineto{\pgfqpoint{3.612458in}{0.733417in}}%
\pgfpathlineto{\pgfqpoint{3.612458in}{1.773919in}}%
\pgfpathlineto{\pgfqpoint{3.548253in}{1.773919in}}%
\pgfpathlineto{\pgfqpoint{3.548253in}{0.733417in}}%
\pgfpathclose%
\pgfusepath{fill}%
\end{pgfscope}%
\begin{pgfscope}%
\pgfpathrectangle{\pgfqpoint{0.787430in}{0.733417in}}{\pgfqpoint{7.062570in}{2.555373in}}%
\pgfusepath{clip}%
\pgfsetbuttcap%
\pgfsetmiterjoin%
\definecolor{currentfill}{rgb}{0.121569,0.466667,0.705882}%
\pgfsetfillcolor{currentfill}%
\pgfsetlinewidth{0.000000pt}%
\definecolor{currentstroke}{rgb}{0.000000,0.000000,0.000000}%
\pgfsetstrokecolor{currentstroke}%
\pgfsetstrokeopacity{0.000000}%
\pgfsetdash{}{0pt}%
\pgfpathmoveto{\pgfqpoint{3.612458in}{0.733417in}}%
\pgfpathlineto{\pgfqpoint{3.676663in}{0.733417in}}%
\pgfpathlineto{\pgfqpoint{3.676663in}{1.920460in}}%
\pgfpathlineto{\pgfqpoint{3.612458in}{1.920460in}}%
\pgfpathlineto{\pgfqpoint{3.612458in}{0.733417in}}%
\pgfpathclose%
\pgfusepath{fill}%
\end{pgfscope}%
\begin{pgfscope}%
\pgfpathrectangle{\pgfqpoint{0.787430in}{0.733417in}}{\pgfqpoint{7.062570in}{2.555373in}}%
\pgfusepath{clip}%
\pgfsetbuttcap%
\pgfsetmiterjoin%
\definecolor{currentfill}{rgb}{0.121569,0.466667,0.705882}%
\pgfsetfillcolor{currentfill}%
\pgfsetlinewidth{0.000000pt}%
\definecolor{currentstroke}{rgb}{0.000000,0.000000,0.000000}%
\pgfsetstrokecolor{currentstroke}%
\pgfsetstrokeopacity{0.000000}%
\pgfsetdash{}{0pt}%
\pgfpathmoveto{\pgfqpoint{3.676663in}{0.733417in}}%
\pgfpathlineto{\pgfqpoint{3.740869in}{0.733417in}}%
\pgfpathlineto{\pgfqpoint{3.740869in}{2.066096in}}%
\pgfpathlineto{\pgfqpoint{3.676663in}{2.066096in}}%
\pgfpathlineto{\pgfqpoint{3.676663in}{0.733417in}}%
\pgfpathclose%
\pgfusepath{fill}%
\end{pgfscope}%
\begin{pgfscope}%
\pgfpathrectangle{\pgfqpoint{0.787430in}{0.733417in}}{\pgfqpoint{7.062570in}{2.555373in}}%
\pgfusepath{clip}%
\pgfsetbuttcap%
\pgfsetmiterjoin%
\definecolor{currentfill}{rgb}{0.121569,0.466667,0.705882}%
\pgfsetfillcolor{currentfill}%
\pgfsetlinewidth{0.000000pt}%
\definecolor{currentstroke}{rgb}{0.000000,0.000000,0.000000}%
\pgfsetstrokecolor{currentstroke}%
\pgfsetstrokeopacity{0.000000}%
\pgfsetdash{}{0pt}%
\pgfpathmoveto{\pgfqpoint{3.740869in}{0.733417in}}%
\pgfpathlineto{\pgfqpoint{3.805074in}{0.733417in}}%
\pgfpathlineto{\pgfqpoint{3.805074in}{2.218001in}}%
\pgfpathlineto{\pgfqpoint{3.740869in}{2.218001in}}%
\pgfpathlineto{\pgfqpoint{3.740869in}{0.733417in}}%
\pgfpathclose%
\pgfusepath{fill}%
\end{pgfscope}%
\begin{pgfscope}%
\pgfpathrectangle{\pgfqpoint{0.787430in}{0.733417in}}{\pgfqpoint{7.062570in}{2.555373in}}%
\pgfusepath{clip}%
\pgfsetbuttcap%
\pgfsetmiterjoin%
\definecolor{currentfill}{rgb}{0.121569,0.466667,0.705882}%
\pgfsetfillcolor{currentfill}%
\pgfsetlinewidth{0.000000pt}%
\definecolor{currentstroke}{rgb}{0.000000,0.000000,0.000000}%
\pgfsetstrokecolor{currentstroke}%
\pgfsetstrokeopacity{0.000000}%
\pgfsetdash{}{0pt}%
\pgfpathmoveto{\pgfqpoint{3.805074in}{0.733417in}}%
\pgfpathlineto{\pgfqpoint{3.869279in}{0.733417in}}%
\pgfpathlineto{\pgfqpoint{3.869279in}{2.377769in}}%
\pgfpathlineto{\pgfqpoint{3.805074in}{2.377769in}}%
\pgfpathlineto{\pgfqpoint{3.805074in}{0.733417in}}%
\pgfpathclose%
\pgfusepath{fill}%
\end{pgfscope}%
\begin{pgfscope}%
\pgfpathrectangle{\pgfqpoint{0.787430in}{0.733417in}}{\pgfqpoint{7.062570in}{2.555373in}}%
\pgfusepath{clip}%
\pgfsetbuttcap%
\pgfsetmiterjoin%
\definecolor{currentfill}{rgb}{0.121569,0.466667,0.705882}%
\pgfsetfillcolor{currentfill}%
\pgfsetlinewidth{0.000000pt}%
\definecolor{currentstroke}{rgb}{0.000000,0.000000,0.000000}%
\pgfsetstrokecolor{currentstroke}%
\pgfsetstrokeopacity{0.000000}%
\pgfsetdash{}{0pt}%
\pgfpathmoveto{\pgfqpoint{3.869279in}{0.733417in}}%
\pgfpathlineto{\pgfqpoint{3.933484in}{0.733417in}}%
\pgfpathlineto{\pgfqpoint{3.933484in}{2.524246in}}%
\pgfpathlineto{\pgfqpoint{3.869279in}{2.524246in}}%
\pgfpathlineto{\pgfqpoint{3.869279in}{0.733417in}}%
\pgfpathclose%
\pgfusepath{fill}%
\end{pgfscope}%
\begin{pgfscope}%
\pgfpathrectangle{\pgfqpoint{0.787430in}{0.733417in}}{\pgfqpoint{7.062570in}{2.555373in}}%
\pgfusepath{clip}%
\pgfsetbuttcap%
\pgfsetmiterjoin%
\definecolor{currentfill}{rgb}{0.121569,0.466667,0.705882}%
\pgfsetfillcolor{currentfill}%
\pgfsetlinewidth{0.000000pt}%
\definecolor{currentstroke}{rgb}{0.000000,0.000000,0.000000}%
\pgfsetstrokecolor{currentstroke}%
\pgfsetstrokeopacity{0.000000}%
\pgfsetdash{}{0pt}%
\pgfpathmoveto{\pgfqpoint{3.933484in}{0.733417in}}%
\pgfpathlineto{\pgfqpoint{3.997689in}{0.733417in}}%
\pgfpathlineto{\pgfqpoint{3.997689in}{2.663298in}}%
\pgfpathlineto{\pgfqpoint{3.933484in}{2.663298in}}%
\pgfpathlineto{\pgfqpoint{3.933484in}{0.733417in}}%
\pgfpathclose%
\pgfusepath{fill}%
\end{pgfscope}%
\begin{pgfscope}%
\pgfpathrectangle{\pgfqpoint{0.787430in}{0.733417in}}{\pgfqpoint{7.062570in}{2.555373in}}%
\pgfusepath{clip}%
\pgfsetbuttcap%
\pgfsetmiterjoin%
\definecolor{currentfill}{rgb}{0.121569,0.466667,0.705882}%
\pgfsetfillcolor{currentfill}%
\pgfsetlinewidth{0.000000pt}%
\definecolor{currentstroke}{rgb}{0.000000,0.000000,0.000000}%
\pgfsetstrokecolor{currentstroke}%
\pgfsetstrokeopacity{0.000000}%
\pgfsetdash{}{0pt}%
\pgfpathmoveto{\pgfqpoint{3.997689in}{0.733417in}}%
\pgfpathlineto{\pgfqpoint{4.061894in}{0.733417in}}%
\pgfpathlineto{\pgfqpoint{4.061894in}{2.799581in}}%
\pgfpathlineto{\pgfqpoint{3.997689in}{2.799581in}}%
\pgfpathlineto{\pgfqpoint{3.997689in}{0.733417in}}%
\pgfpathclose%
\pgfusepath{fill}%
\end{pgfscope}%
\begin{pgfscope}%
\pgfpathrectangle{\pgfqpoint{0.787430in}{0.733417in}}{\pgfqpoint{7.062570in}{2.555373in}}%
\pgfusepath{clip}%
\pgfsetbuttcap%
\pgfsetmiterjoin%
\definecolor{currentfill}{rgb}{0.121569,0.466667,0.705882}%
\pgfsetfillcolor{currentfill}%
\pgfsetlinewidth{0.000000pt}%
\definecolor{currentstroke}{rgb}{0.000000,0.000000,0.000000}%
\pgfsetstrokecolor{currentstroke}%
\pgfsetstrokeopacity{0.000000}%
\pgfsetdash{}{0pt}%
\pgfpathmoveto{\pgfqpoint{4.061894in}{0.733417in}}%
\pgfpathlineto{\pgfqpoint{4.126100in}{0.733417in}}%
\pgfpathlineto{\pgfqpoint{4.126100in}{2.912373in}}%
\pgfpathlineto{\pgfqpoint{4.061894in}{2.912373in}}%
\pgfpathlineto{\pgfqpoint{4.061894in}{0.733417in}}%
\pgfpathclose%
\pgfusepath{fill}%
\end{pgfscope}%
\begin{pgfscope}%
\pgfpathrectangle{\pgfqpoint{0.787430in}{0.733417in}}{\pgfqpoint{7.062570in}{2.555373in}}%
\pgfusepath{clip}%
\pgfsetbuttcap%
\pgfsetmiterjoin%
\definecolor{currentfill}{rgb}{0.121569,0.466667,0.705882}%
\pgfsetfillcolor{currentfill}%
\pgfsetlinewidth{0.000000pt}%
\definecolor{currentstroke}{rgb}{0.000000,0.000000,0.000000}%
\pgfsetstrokecolor{currentstroke}%
\pgfsetstrokeopacity{0.000000}%
\pgfsetdash{}{0pt}%
\pgfpathmoveto{\pgfqpoint{4.126100in}{0.733417in}}%
\pgfpathlineto{\pgfqpoint{4.190305in}{0.733417in}}%
\pgfpathlineto{\pgfqpoint{4.190305in}{3.006646in}}%
\pgfpathlineto{\pgfqpoint{4.126100in}{3.006646in}}%
\pgfpathlineto{\pgfqpoint{4.126100in}{0.733417in}}%
\pgfpathclose%
\pgfusepath{fill}%
\end{pgfscope}%
\begin{pgfscope}%
\pgfpathrectangle{\pgfqpoint{0.787430in}{0.733417in}}{\pgfqpoint{7.062570in}{2.555373in}}%
\pgfusepath{clip}%
\pgfsetbuttcap%
\pgfsetmiterjoin%
\definecolor{currentfill}{rgb}{0.121569,0.466667,0.705882}%
\pgfsetfillcolor{currentfill}%
\pgfsetlinewidth{0.000000pt}%
\definecolor{currentstroke}{rgb}{0.000000,0.000000,0.000000}%
\pgfsetstrokecolor{currentstroke}%
\pgfsetstrokeopacity{0.000000}%
\pgfsetdash{}{0pt}%
\pgfpathmoveto{\pgfqpoint{4.190305in}{0.733417in}}%
\pgfpathlineto{\pgfqpoint{4.254510in}{0.733417in}}%
\pgfpathlineto{\pgfqpoint{4.254510in}{3.084573in}}%
\pgfpathlineto{\pgfqpoint{4.190305in}{3.084573in}}%
\pgfpathlineto{\pgfqpoint{4.190305in}{0.733417in}}%
\pgfpathclose%
\pgfusepath{fill}%
\end{pgfscope}%
\begin{pgfscope}%
\pgfpathrectangle{\pgfqpoint{0.787430in}{0.733417in}}{\pgfqpoint{7.062570in}{2.555373in}}%
\pgfusepath{clip}%
\pgfsetbuttcap%
\pgfsetmiterjoin%
\definecolor{currentfill}{rgb}{0.121569,0.466667,0.705882}%
\pgfsetfillcolor{currentfill}%
\pgfsetlinewidth{0.000000pt}%
\definecolor{currentstroke}{rgb}{0.000000,0.000000,0.000000}%
\pgfsetstrokecolor{currentstroke}%
\pgfsetstrokeopacity{0.000000}%
\pgfsetdash{}{0pt}%
\pgfpathmoveto{\pgfqpoint{4.254510in}{0.733417in}}%
\pgfpathlineto{\pgfqpoint{4.318715in}{0.733417in}}%
\pgfpathlineto{\pgfqpoint{4.318715in}{3.132268in}}%
\pgfpathlineto{\pgfqpoint{4.254510in}{3.132268in}}%
\pgfpathlineto{\pgfqpoint{4.254510in}{0.733417in}}%
\pgfpathclose%
\pgfusepath{fill}%
\end{pgfscope}%
\begin{pgfscope}%
\pgfpathrectangle{\pgfqpoint{0.787430in}{0.733417in}}{\pgfqpoint{7.062570in}{2.555373in}}%
\pgfusepath{clip}%
\pgfsetbuttcap%
\pgfsetmiterjoin%
\definecolor{currentfill}{rgb}{0.121569,0.466667,0.705882}%
\pgfsetfillcolor{currentfill}%
\pgfsetlinewidth{0.000000pt}%
\definecolor{currentstroke}{rgb}{0.000000,0.000000,0.000000}%
\pgfsetstrokecolor{currentstroke}%
\pgfsetstrokeopacity{0.000000}%
\pgfsetdash{}{0pt}%
\pgfpathmoveto{\pgfqpoint{4.318715in}{0.733417in}}%
\pgfpathlineto{\pgfqpoint{4.382920in}{0.733417in}}%
\pgfpathlineto{\pgfqpoint{4.382920in}{3.160374in}}%
\pgfpathlineto{\pgfqpoint{4.318715in}{3.160374in}}%
\pgfpathlineto{\pgfqpoint{4.318715in}{0.733417in}}%
\pgfpathclose%
\pgfusepath{fill}%
\end{pgfscope}%
\begin{pgfscope}%
\pgfpathrectangle{\pgfqpoint{0.787430in}{0.733417in}}{\pgfqpoint{7.062570in}{2.555373in}}%
\pgfusepath{clip}%
\pgfsetbuttcap%
\pgfsetmiterjoin%
\definecolor{currentfill}{rgb}{0.121569,0.466667,0.705882}%
\pgfsetfillcolor{currentfill}%
\pgfsetlinewidth{0.000000pt}%
\definecolor{currentstroke}{rgb}{0.000000,0.000000,0.000000}%
\pgfsetstrokecolor{currentstroke}%
\pgfsetstrokeopacity{0.000000}%
\pgfsetdash{}{0pt}%
\pgfpathmoveto{\pgfqpoint{4.382920in}{0.733417in}}%
\pgfpathlineto{\pgfqpoint{4.447125in}{0.733417in}}%
\pgfpathlineto{\pgfqpoint{4.447125in}{3.165468in}}%
\pgfpathlineto{\pgfqpoint{4.382920in}{3.165468in}}%
\pgfpathlineto{\pgfqpoint{4.382920in}{0.733417in}}%
\pgfpathclose%
\pgfusepath{fill}%
\end{pgfscope}%
\begin{pgfscope}%
\pgfpathrectangle{\pgfqpoint{0.787430in}{0.733417in}}{\pgfqpoint{7.062570in}{2.555373in}}%
\pgfusepath{clip}%
\pgfsetbuttcap%
\pgfsetmiterjoin%
\definecolor{currentfill}{rgb}{0.121569,0.466667,0.705882}%
\pgfsetfillcolor{currentfill}%
\pgfsetlinewidth{0.000000pt}%
\definecolor{currentstroke}{rgb}{0.000000,0.000000,0.000000}%
\pgfsetstrokecolor{currentstroke}%
\pgfsetstrokeopacity{0.000000}%
\pgfsetdash{}{0pt}%
\pgfpathmoveto{\pgfqpoint{4.447125in}{0.733417in}}%
\pgfpathlineto{\pgfqpoint{4.511331in}{0.733417in}}%
\pgfpathlineto{\pgfqpoint{4.511331in}{3.135954in}}%
\pgfpathlineto{\pgfqpoint{4.447125in}{3.135954in}}%
\pgfpathlineto{\pgfqpoint{4.447125in}{0.733417in}}%
\pgfpathclose%
\pgfusepath{fill}%
\end{pgfscope}%
\begin{pgfscope}%
\pgfpathrectangle{\pgfqpoint{0.787430in}{0.733417in}}{\pgfqpoint{7.062570in}{2.555373in}}%
\pgfusepath{clip}%
\pgfsetbuttcap%
\pgfsetmiterjoin%
\definecolor{currentfill}{rgb}{0.121569,0.466667,0.705882}%
\pgfsetfillcolor{currentfill}%
\pgfsetlinewidth{0.000000pt}%
\definecolor{currentstroke}{rgb}{0.000000,0.000000,0.000000}%
\pgfsetstrokecolor{currentstroke}%
\pgfsetstrokeopacity{0.000000}%
\pgfsetdash{}{0pt}%
\pgfpathmoveto{\pgfqpoint{4.511331in}{0.733417in}}%
\pgfpathlineto{\pgfqpoint{4.575536in}{0.733417in}}%
\pgfpathlineto{\pgfqpoint{4.575536in}{3.085940in}}%
\pgfpathlineto{\pgfqpoint{4.511331in}{3.085940in}}%
\pgfpathlineto{\pgfqpoint{4.511331in}{0.733417in}}%
\pgfpathclose%
\pgfusepath{fill}%
\end{pgfscope}%
\begin{pgfscope}%
\pgfpathrectangle{\pgfqpoint{0.787430in}{0.733417in}}{\pgfqpoint{7.062570in}{2.555373in}}%
\pgfusepath{clip}%
\pgfsetbuttcap%
\pgfsetmiterjoin%
\definecolor{currentfill}{rgb}{0.121569,0.466667,0.705882}%
\pgfsetfillcolor{currentfill}%
\pgfsetlinewidth{0.000000pt}%
\definecolor{currentstroke}{rgb}{0.000000,0.000000,0.000000}%
\pgfsetstrokecolor{currentstroke}%
\pgfsetstrokeopacity{0.000000}%
\pgfsetdash{}{0pt}%
\pgfpathmoveto{\pgfqpoint{4.575536in}{0.733417in}}%
\pgfpathlineto{\pgfqpoint{4.639741in}{0.733417in}}%
\pgfpathlineto{\pgfqpoint{4.639741in}{3.007026in}}%
\pgfpathlineto{\pgfqpoint{4.575536in}{3.007026in}}%
\pgfpathlineto{\pgfqpoint{4.575536in}{0.733417in}}%
\pgfpathclose%
\pgfusepath{fill}%
\end{pgfscope}%
\begin{pgfscope}%
\pgfpathrectangle{\pgfqpoint{0.787430in}{0.733417in}}{\pgfqpoint{7.062570in}{2.555373in}}%
\pgfusepath{clip}%
\pgfsetbuttcap%
\pgfsetmiterjoin%
\definecolor{currentfill}{rgb}{0.121569,0.466667,0.705882}%
\pgfsetfillcolor{currentfill}%
\pgfsetlinewidth{0.000000pt}%
\definecolor{currentstroke}{rgb}{0.000000,0.000000,0.000000}%
\pgfsetstrokecolor{currentstroke}%
\pgfsetstrokeopacity{0.000000}%
\pgfsetdash{}{0pt}%
\pgfpathmoveto{\pgfqpoint{4.639741in}{0.733417in}}%
\pgfpathlineto{\pgfqpoint{4.703946in}{0.733417in}}%
\pgfpathlineto{\pgfqpoint{4.703946in}{2.909283in}}%
\pgfpathlineto{\pgfqpoint{4.639741in}{2.909283in}}%
\pgfpathlineto{\pgfqpoint{4.639741in}{0.733417in}}%
\pgfpathclose%
\pgfusepath{fill}%
\end{pgfscope}%
\begin{pgfscope}%
\pgfpathrectangle{\pgfqpoint{0.787430in}{0.733417in}}{\pgfqpoint{7.062570in}{2.555373in}}%
\pgfusepath{clip}%
\pgfsetbuttcap%
\pgfsetmiterjoin%
\definecolor{currentfill}{rgb}{0.121569,0.466667,0.705882}%
\pgfsetfillcolor{currentfill}%
\pgfsetlinewidth{0.000000pt}%
\definecolor{currentstroke}{rgb}{0.000000,0.000000,0.000000}%
\pgfsetstrokecolor{currentstroke}%
\pgfsetstrokeopacity{0.000000}%
\pgfsetdash{}{0pt}%
\pgfpathmoveto{\pgfqpoint{4.703946in}{0.733417in}}%
\pgfpathlineto{\pgfqpoint{4.768151in}{0.733417in}}%
\pgfpathlineto{\pgfqpoint{4.768151in}{2.794931in}}%
\pgfpathlineto{\pgfqpoint{4.703946in}{2.794931in}}%
\pgfpathlineto{\pgfqpoint{4.703946in}{0.733417in}}%
\pgfpathclose%
\pgfusepath{fill}%
\end{pgfscope}%
\begin{pgfscope}%
\pgfpathrectangle{\pgfqpoint{0.787430in}{0.733417in}}{\pgfqpoint{7.062570in}{2.555373in}}%
\pgfusepath{clip}%
\pgfsetbuttcap%
\pgfsetmiterjoin%
\definecolor{currentfill}{rgb}{0.121569,0.466667,0.705882}%
\pgfsetfillcolor{currentfill}%
\pgfsetlinewidth{0.000000pt}%
\definecolor{currentstroke}{rgb}{0.000000,0.000000,0.000000}%
\pgfsetstrokecolor{currentstroke}%
\pgfsetstrokeopacity{0.000000}%
\pgfsetdash{}{0pt}%
\pgfpathmoveto{\pgfqpoint{4.768151in}{0.733417in}}%
\pgfpathlineto{\pgfqpoint{4.832357in}{0.733417in}}%
\pgfpathlineto{\pgfqpoint{4.832357in}{2.665805in}}%
\pgfpathlineto{\pgfqpoint{4.768151in}{2.665805in}}%
\pgfpathlineto{\pgfqpoint{4.768151in}{0.733417in}}%
\pgfpathclose%
\pgfusepath{fill}%
\end{pgfscope}%
\begin{pgfscope}%
\pgfpathrectangle{\pgfqpoint{0.787430in}{0.733417in}}{\pgfqpoint{7.062570in}{2.555373in}}%
\pgfusepath{clip}%
\pgfsetbuttcap%
\pgfsetmiterjoin%
\definecolor{currentfill}{rgb}{0.121569,0.466667,0.705882}%
\pgfsetfillcolor{currentfill}%
\pgfsetlinewidth{0.000000pt}%
\definecolor{currentstroke}{rgb}{0.000000,0.000000,0.000000}%
\pgfsetstrokecolor{currentstroke}%
\pgfsetstrokeopacity{0.000000}%
\pgfsetdash{}{0pt}%
\pgfpathmoveto{\pgfqpoint{4.832357in}{0.733417in}}%
\pgfpathlineto{\pgfqpoint{4.896562in}{0.733417in}}%
\pgfpathlineto{\pgfqpoint{4.896562in}{2.528634in}}%
\pgfpathlineto{\pgfqpoint{4.832357in}{2.528634in}}%
\pgfpathlineto{\pgfqpoint{4.832357in}{0.733417in}}%
\pgfpathclose%
\pgfusepath{fill}%
\end{pgfscope}%
\begin{pgfscope}%
\pgfpathrectangle{\pgfqpoint{0.787430in}{0.733417in}}{\pgfqpoint{7.062570in}{2.555373in}}%
\pgfusepath{clip}%
\pgfsetbuttcap%
\pgfsetmiterjoin%
\definecolor{currentfill}{rgb}{0.121569,0.466667,0.705882}%
\pgfsetfillcolor{currentfill}%
\pgfsetlinewidth{0.000000pt}%
\definecolor{currentstroke}{rgb}{0.000000,0.000000,0.000000}%
\pgfsetstrokecolor{currentstroke}%
\pgfsetstrokeopacity{0.000000}%
\pgfsetdash{}{0pt}%
\pgfpathmoveto{\pgfqpoint{4.896562in}{0.733417in}}%
\pgfpathlineto{\pgfqpoint{4.960767in}{0.733417in}}%
\pgfpathlineto{\pgfqpoint{4.960767in}{2.382121in}}%
\pgfpathlineto{\pgfqpoint{4.896562in}{2.382121in}}%
\pgfpathlineto{\pgfqpoint{4.896562in}{0.733417in}}%
\pgfpathclose%
\pgfusepath{fill}%
\end{pgfscope}%
\begin{pgfscope}%
\pgfpathrectangle{\pgfqpoint{0.787430in}{0.733417in}}{\pgfqpoint{7.062570in}{2.555373in}}%
\pgfusepath{clip}%
\pgfsetbuttcap%
\pgfsetmiterjoin%
\definecolor{currentfill}{rgb}{0.121569,0.466667,0.705882}%
\pgfsetfillcolor{currentfill}%
\pgfsetlinewidth{0.000000pt}%
\definecolor{currentstroke}{rgb}{0.000000,0.000000,0.000000}%
\pgfsetstrokecolor{currentstroke}%
\pgfsetstrokeopacity{0.000000}%
\pgfsetdash{}{0pt}%
\pgfpathmoveto{\pgfqpoint{4.960767in}{0.733417in}}%
\pgfpathlineto{\pgfqpoint{5.024972in}{0.733417in}}%
\pgfpathlineto{\pgfqpoint{5.024972in}{2.221383in}}%
\pgfpathlineto{\pgfqpoint{4.960767in}{2.221383in}}%
\pgfpathlineto{\pgfqpoint{4.960767in}{0.733417in}}%
\pgfpathclose%
\pgfusepath{fill}%
\end{pgfscope}%
\begin{pgfscope}%
\pgfpathrectangle{\pgfqpoint{0.787430in}{0.733417in}}{\pgfqpoint{7.062570in}{2.555373in}}%
\pgfusepath{clip}%
\pgfsetbuttcap%
\pgfsetmiterjoin%
\definecolor{currentfill}{rgb}{0.121569,0.466667,0.705882}%
\pgfsetfillcolor{currentfill}%
\pgfsetlinewidth{0.000000pt}%
\definecolor{currentstroke}{rgb}{0.000000,0.000000,0.000000}%
\pgfsetstrokecolor{currentstroke}%
\pgfsetstrokeopacity{0.000000}%
\pgfsetdash{}{0pt}%
\pgfpathmoveto{\pgfqpoint{5.024972in}{0.733417in}}%
\pgfpathlineto{\pgfqpoint{5.089177in}{0.733417in}}%
\pgfpathlineto{\pgfqpoint{5.089177in}{2.068509in}}%
\pgfpathlineto{\pgfqpoint{5.024972in}{2.068509in}}%
\pgfpathlineto{\pgfqpoint{5.024972in}{0.733417in}}%
\pgfpathclose%
\pgfusepath{fill}%
\end{pgfscope}%
\begin{pgfscope}%
\pgfpathrectangle{\pgfqpoint{0.787430in}{0.733417in}}{\pgfqpoint{7.062570in}{2.555373in}}%
\pgfusepath{clip}%
\pgfsetbuttcap%
\pgfsetmiterjoin%
\definecolor{currentfill}{rgb}{0.121569,0.466667,0.705882}%
\pgfsetfillcolor{currentfill}%
\pgfsetlinewidth{0.000000pt}%
\definecolor{currentstroke}{rgb}{0.000000,0.000000,0.000000}%
\pgfsetstrokecolor{currentstroke}%
\pgfsetstrokeopacity{0.000000}%
\pgfsetdash{}{0pt}%
\pgfpathmoveto{\pgfqpoint{5.089177in}{0.733417in}}%
\pgfpathlineto{\pgfqpoint{5.153382in}{0.733417in}}%
\pgfpathlineto{\pgfqpoint{5.153382in}{1.916745in}}%
\pgfpathlineto{\pgfqpoint{5.089177in}{1.916745in}}%
\pgfpathlineto{\pgfqpoint{5.089177in}{0.733417in}}%
\pgfpathclose%
\pgfusepath{fill}%
\end{pgfscope}%
\begin{pgfscope}%
\pgfpathrectangle{\pgfqpoint{0.787430in}{0.733417in}}{\pgfqpoint{7.062570in}{2.555373in}}%
\pgfusepath{clip}%
\pgfsetbuttcap%
\pgfsetmiterjoin%
\definecolor{currentfill}{rgb}{0.121569,0.466667,0.705882}%
\pgfsetfillcolor{currentfill}%
\pgfsetlinewidth{0.000000pt}%
\definecolor{currentstroke}{rgb}{0.000000,0.000000,0.000000}%
\pgfsetstrokecolor{currentstroke}%
\pgfsetstrokeopacity{0.000000}%
\pgfsetdash{}{0pt}%
\pgfpathmoveto{\pgfqpoint{5.153382in}{0.733417in}}%
\pgfpathlineto{\pgfqpoint{5.217588in}{0.733417in}}%
\pgfpathlineto{\pgfqpoint{5.217588in}{1.768930in}}%
\pgfpathlineto{\pgfqpoint{5.153382in}{1.768930in}}%
\pgfpathlineto{\pgfqpoint{5.153382in}{0.733417in}}%
\pgfpathclose%
\pgfusepath{fill}%
\end{pgfscope}%
\begin{pgfscope}%
\pgfpathrectangle{\pgfqpoint{0.787430in}{0.733417in}}{\pgfqpoint{7.062570in}{2.555373in}}%
\pgfusepath{clip}%
\pgfsetbuttcap%
\pgfsetmiterjoin%
\definecolor{currentfill}{rgb}{0.121569,0.466667,0.705882}%
\pgfsetfillcolor{currentfill}%
\pgfsetlinewidth{0.000000pt}%
\definecolor{currentstroke}{rgb}{0.000000,0.000000,0.000000}%
\pgfsetstrokecolor{currentstroke}%
\pgfsetstrokeopacity{0.000000}%
\pgfsetdash{}{0pt}%
\pgfpathmoveto{\pgfqpoint{5.217588in}{0.733417in}}%
\pgfpathlineto{\pgfqpoint{5.281793in}{0.733417in}}%
\pgfpathlineto{\pgfqpoint{5.281793in}{1.633593in}}%
\pgfpathlineto{\pgfqpoint{5.217588in}{1.633593in}}%
\pgfpathlineto{\pgfqpoint{5.217588in}{0.733417in}}%
\pgfpathclose%
\pgfusepath{fill}%
\end{pgfscope}%
\begin{pgfscope}%
\pgfpathrectangle{\pgfqpoint{0.787430in}{0.733417in}}{\pgfqpoint{7.062570in}{2.555373in}}%
\pgfusepath{clip}%
\pgfsetbuttcap%
\pgfsetmiterjoin%
\definecolor{currentfill}{rgb}{0.121569,0.466667,0.705882}%
\pgfsetfillcolor{currentfill}%
\pgfsetlinewidth{0.000000pt}%
\definecolor{currentstroke}{rgb}{0.000000,0.000000,0.000000}%
\pgfsetstrokecolor{currentstroke}%
\pgfsetstrokeopacity{0.000000}%
\pgfsetdash{}{0pt}%
\pgfpathmoveto{\pgfqpoint{5.281793in}{0.733417in}}%
\pgfpathlineto{\pgfqpoint{5.345998in}{0.733417in}}%
\pgfpathlineto{\pgfqpoint{5.345998in}{1.503018in}}%
\pgfpathlineto{\pgfqpoint{5.281793in}{1.503018in}}%
\pgfpathlineto{\pgfqpoint{5.281793in}{0.733417in}}%
\pgfpathclose%
\pgfusepath{fill}%
\end{pgfscope}%
\begin{pgfscope}%
\pgfpathrectangle{\pgfqpoint{0.787430in}{0.733417in}}{\pgfqpoint{7.062570in}{2.555373in}}%
\pgfusepath{clip}%
\pgfsetbuttcap%
\pgfsetmiterjoin%
\definecolor{currentfill}{rgb}{0.121569,0.466667,0.705882}%
\pgfsetfillcolor{currentfill}%
\pgfsetlinewidth{0.000000pt}%
\definecolor{currentstroke}{rgb}{0.000000,0.000000,0.000000}%
\pgfsetstrokecolor{currentstroke}%
\pgfsetstrokeopacity{0.000000}%
\pgfsetdash{}{0pt}%
\pgfpathmoveto{\pgfqpoint{5.345998in}{0.733417in}}%
\pgfpathlineto{\pgfqpoint{5.410203in}{0.733417in}}%
\pgfpathlineto{\pgfqpoint{5.410203in}{1.389238in}}%
\pgfpathlineto{\pgfqpoint{5.345998in}{1.389238in}}%
\pgfpathlineto{\pgfqpoint{5.345998in}{0.733417in}}%
\pgfpathclose%
\pgfusepath{fill}%
\end{pgfscope}%
\begin{pgfscope}%
\pgfpathrectangle{\pgfqpoint{0.787430in}{0.733417in}}{\pgfqpoint{7.062570in}{2.555373in}}%
\pgfusepath{clip}%
\pgfsetbuttcap%
\pgfsetmiterjoin%
\definecolor{currentfill}{rgb}{0.121569,0.466667,0.705882}%
\pgfsetfillcolor{currentfill}%
\pgfsetlinewidth{0.000000pt}%
\definecolor{currentstroke}{rgb}{0.000000,0.000000,0.000000}%
\pgfsetstrokecolor{currentstroke}%
\pgfsetstrokeopacity{0.000000}%
\pgfsetdash{}{0pt}%
\pgfpathmoveto{\pgfqpoint{5.410203in}{0.733417in}}%
\pgfpathlineto{\pgfqpoint{5.474408in}{0.733417in}}%
\pgfpathlineto{\pgfqpoint{5.474408in}{1.284566in}}%
\pgfpathlineto{\pgfqpoint{5.410203in}{1.284566in}}%
\pgfpathlineto{\pgfqpoint{5.410203in}{0.733417in}}%
\pgfpathclose%
\pgfusepath{fill}%
\end{pgfscope}%
\begin{pgfscope}%
\pgfpathrectangle{\pgfqpoint{0.787430in}{0.733417in}}{\pgfqpoint{7.062570in}{2.555373in}}%
\pgfusepath{clip}%
\pgfsetbuttcap%
\pgfsetmiterjoin%
\definecolor{currentfill}{rgb}{0.121569,0.466667,0.705882}%
\pgfsetfillcolor{currentfill}%
\pgfsetlinewidth{0.000000pt}%
\definecolor{currentstroke}{rgb}{0.000000,0.000000,0.000000}%
\pgfsetstrokecolor{currentstroke}%
\pgfsetstrokeopacity{0.000000}%
\pgfsetdash{}{0pt}%
\pgfpathmoveto{\pgfqpoint{5.474408in}{0.733417in}}%
\pgfpathlineto{\pgfqpoint{5.538614in}{0.733417in}}%
\pgfpathlineto{\pgfqpoint{5.538614in}{1.193524in}}%
\pgfpathlineto{\pgfqpoint{5.474408in}{1.193524in}}%
\pgfpathlineto{\pgfqpoint{5.474408in}{0.733417in}}%
\pgfpathclose%
\pgfusepath{fill}%
\end{pgfscope}%
\begin{pgfscope}%
\pgfpathrectangle{\pgfqpoint{0.787430in}{0.733417in}}{\pgfqpoint{7.062570in}{2.555373in}}%
\pgfusepath{clip}%
\pgfsetbuttcap%
\pgfsetmiterjoin%
\definecolor{currentfill}{rgb}{0.121569,0.466667,0.705882}%
\pgfsetfillcolor{currentfill}%
\pgfsetlinewidth{0.000000pt}%
\definecolor{currentstroke}{rgb}{0.000000,0.000000,0.000000}%
\pgfsetstrokecolor{currentstroke}%
\pgfsetstrokeopacity{0.000000}%
\pgfsetdash{}{0pt}%
\pgfpathmoveto{\pgfqpoint{5.538614in}{0.733417in}}%
\pgfpathlineto{\pgfqpoint{5.602819in}{0.733417in}}%
\pgfpathlineto{\pgfqpoint{5.602819in}{1.109265in}}%
\pgfpathlineto{\pgfqpoint{5.538614in}{1.109265in}}%
\pgfpathlineto{\pgfqpoint{5.538614in}{0.733417in}}%
\pgfpathclose%
\pgfusepath{fill}%
\end{pgfscope}%
\begin{pgfscope}%
\pgfpathrectangle{\pgfqpoint{0.787430in}{0.733417in}}{\pgfqpoint{7.062570in}{2.555373in}}%
\pgfusepath{clip}%
\pgfsetbuttcap%
\pgfsetmiterjoin%
\definecolor{currentfill}{rgb}{0.121569,0.466667,0.705882}%
\pgfsetfillcolor{currentfill}%
\pgfsetlinewidth{0.000000pt}%
\definecolor{currentstroke}{rgb}{0.000000,0.000000,0.000000}%
\pgfsetstrokecolor{currentstroke}%
\pgfsetstrokeopacity{0.000000}%
\pgfsetdash{}{0pt}%
\pgfpathmoveto{\pgfqpoint{5.602819in}{0.733417in}}%
\pgfpathlineto{\pgfqpoint{5.667024in}{0.733417in}}%
\pgfpathlineto{\pgfqpoint{5.667024in}{1.039949in}}%
\pgfpathlineto{\pgfqpoint{5.602819in}{1.039949in}}%
\pgfpathlineto{\pgfqpoint{5.602819in}{0.733417in}}%
\pgfpathclose%
\pgfusepath{fill}%
\end{pgfscope}%
\begin{pgfscope}%
\pgfpathrectangle{\pgfqpoint{0.787430in}{0.733417in}}{\pgfqpoint{7.062570in}{2.555373in}}%
\pgfusepath{clip}%
\pgfsetbuttcap%
\pgfsetmiterjoin%
\definecolor{currentfill}{rgb}{0.121569,0.466667,0.705882}%
\pgfsetfillcolor{currentfill}%
\pgfsetlinewidth{0.000000pt}%
\definecolor{currentstroke}{rgb}{0.000000,0.000000,0.000000}%
\pgfsetstrokecolor{currentstroke}%
\pgfsetstrokeopacity{0.000000}%
\pgfsetdash{}{0pt}%
\pgfpathmoveto{\pgfqpoint{5.667024in}{0.733417in}}%
\pgfpathlineto{\pgfqpoint{5.731229in}{0.733417in}}%
\pgfpathlineto{\pgfqpoint{5.731229in}{0.981219in}}%
\pgfpathlineto{\pgfqpoint{5.667024in}{0.981219in}}%
\pgfpathlineto{\pgfqpoint{5.667024in}{0.733417in}}%
\pgfpathclose%
\pgfusepath{fill}%
\end{pgfscope}%
\begin{pgfscope}%
\pgfpathrectangle{\pgfqpoint{0.787430in}{0.733417in}}{\pgfqpoint{7.062570in}{2.555373in}}%
\pgfusepath{clip}%
\pgfsetbuttcap%
\pgfsetmiterjoin%
\definecolor{currentfill}{rgb}{0.121569,0.466667,0.705882}%
\pgfsetfillcolor{currentfill}%
\pgfsetlinewidth{0.000000pt}%
\definecolor{currentstroke}{rgb}{0.000000,0.000000,0.000000}%
\pgfsetstrokecolor{currentstroke}%
\pgfsetstrokeopacity{0.000000}%
\pgfsetdash{}{0pt}%
\pgfpathmoveto{\pgfqpoint{5.731229in}{0.733417in}}%
\pgfpathlineto{\pgfqpoint{5.795434in}{0.733417in}}%
\pgfpathlineto{\pgfqpoint{5.795434in}{0.931532in}}%
\pgfpathlineto{\pgfqpoint{5.731229in}{0.931532in}}%
\pgfpathlineto{\pgfqpoint{5.731229in}{0.733417in}}%
\pgfpathclose%
\pgfusepath{fill}%
\end{pgfscope}%
\begin{pgfscope}%
\pgfpathrectangle{\pgfqpoint{0.787430in}{0.733417in}}{\pgfqpoint{7.062570in}{2.555373in}}%
\pgfusepath{clip}%
\pgfsetbuttcap%
\pgfsetmiterjoin%
\definecolor{currentfill}{rgb}{0.121569,0.466667,0.705882}%
\pgfsetfillcolor{currentfill}%
\pgfsetlinewidth{0.000000pt}%
\definecolor{currentstroke}{rgb}{0.000000,0.000000,0.000000}%
\pgfsetstrokecolor{currentstroke}%
\pgfsetstrokeopacity{0.000000}%
\pgfsetdash{}{0pt}%
\pgfpathmoveto{\pgfqpoint{5.795434in}{0.733417in}}%
\pgfpathlineto{\pgfqpoint{5.859639in}{0.733417in}}%
\pgfpathlineto{\pgfqpoint{5.859639in}{0.887910in}}%
\pgfpathlineto{\pgfqpoint{5.795434in}{0.887910in}}%
\pgfpathlineto{\pgfqpoint{5.795434in}{0.733417in}}%
\pgfpathclose%
\pgfusepath{fill}%
\end{pgfscope}%
\begin{pgfscope}%
\pgfpathrectangle{\pgfqpoint{0.787430in}{0.733417in}}{\pgfqpoint{7.062570in}{2.555373in}}%
\pgfusepath{clip}%
\pgfsetbuttcap%
\pgfsetmiterjoin%
\definecolor{currentfill}{rgb}{0.121569,0.466667,0.705882}%
\pgfsetfillcolor{currentfill}%
\pgfsetlinewidth{0.000000pt}%
\definecolor{currentstroke}{rgb}{0.000000,0.000000,0.000000}%
\pgfsetstrokecolor{currentstroke}%
\pgfsetstrokeopacity{0.000000}%
\pgfsetdash{}{0pt}%
\pgfpathmoveto{\pgfqpoint{5.859639in}{0.733417in}}%
\pgfpathlineto{\pgfqpoint{5.923845in}{0.733417in}}%
\pgfpathlineto{\pgfqpoint{5.923845in}{0.854014in}}%
\pgfpathlineto{\pgfqpoint{5.859639in}{0.854014in}}%
\pgfpathlineto{\pgfqpoint{5.859639in}{0.733417in}}%
\pgfpathclose%
\pgfusepath{fill}%
\end{pgfscope}%
\begin{pgfscope}%
\pgfpathrectangle{\pgfqpoint{0.787430in}{0.733417in}}{\pgfqpoint{7.062570in}{2.555373in}}%
\pgfusepath{clip}%
\pgfsetbuttcap%
\pgfsetmiterjoin%
\definecolor{currentfill}{rgb}{0.121569,0.466667,0.705882}%
\pgfsetfillcolor{currentfill}%
\pgfsetlinewidth{0.000000pt}%
\definecolor{currentstroke}{rgb}{0.000000,0.000000,0.000000}%
\pgfsetstrokecolor{currentstroke}%
\pgfsetstrokeopacity{0.000000}%
\pgfsetdash{}{0pt}%
\pgfpathmoveto{\pgfqpoint{5.923845in}{0.733417in}}%
\pgfpathlineto{\pgfqpoint{5.988050in}{0.733417in}}%
\pgfpathlineto{\pgfqpoint{5.988050in}{0.827083in}}%
\pgfpathlineto{\pgfqpoint{5.923845in}{0.827083in}}%
\pgfpathlineto{\pgfqpoint{5.923845in}{0.733417in}}%
\pgfpathclose%
\pgfusepath{fill}%
\end{pgfscope}%
\begin{pgfscope}%
\pgfpathrectangle{\pgfqpoint{0.787430in}{0.733417in}}{\pgfqpoint{7.062570in}{2.555373in}}%
\pgfusepath{clip}%
\pgfsetbuttcap%
\pgfsetmiterjoin%
\definecolor{currentfill}{rgb}{0.121569,0.466667,0.705882}%
\pgfsetfillcolor{currentfill}%
\pgfsetlinewidth{0.000000pt}%
\definecolor{currentstroke}{rgb}{0.000000,0.000000,0.000000}%
\pgfsetstrokecolor{currentstroke}%
\pgfsetstrokeopacity{0.000000}%
\pgfsetdash{}{0pt}%
\pgfpathmoveto{\pgfqpoint{5.988050in}{0.733417in}}%
\pgfpathlineto{\pgfqpoint{6.052255in}{0.733417in}}%
\pgfpathlineto{\pgfqpoint{6.052255in}{0.804404in}}%
\pgfpathlineto{\pgfqpoint{5.988050in}{0.804404in}}%
\pgfpathlineto{\pgfqpoint{5.988050in}{0.733417in}}%
\pgfpathclose%
\pgfusepath{fill}%
\end{pgfscope}%
\begin{pgfscope}%
\pgfpathrectangle{\pgfqpoint{0.787430in}{0.733417in}}{\pgfqpoint{7.062570in}{2.555373in}}%
\pgfusepath{clip}%
\pgfsetbuttcap%
\pgfsetmiterjoin%
\definecolor{currentfill}{rgb}{0.121569,0.466667,0.705882}%
\pgfsetfillcolor{currentfill}%
\pgfsetlinewidth{0.000000pt}%
\definecolor{currentstroke}{rgb}{0.000000,0.000000,0.000000}%
\pgfsetstrokecolor{currentstroke}%
\pgfsetstrokeopacity{0.000000}%
\pgfsetdash{}{0pt}%
\pgfpathmoveto{\pgfqpoint{6.052255in}{0.733417in}}%
\pgfpathlineto{\pgfqpoint{6.116460in}{0.733417in}}%
\pgfpathlineto{\pgfqpoint{6.116460in}{0.786971in}}%
\pgfpathlineto{\pgfqpoint{6.052255in}{0.786971in}}%
\pgfpathlineto{\pgfqpoint{6.052255in}{0.733417in}}%
\pgfpathclose%
\pgfusepath{fill}%
\end{pgfscope}%
\begin{pgfscope}%
\pgfpathrectangle{\pgfqpoint{0.787430in}{0.733417in}}{\pgfqpoint{7.062570in}{2.555373in}}%
\pgfusepath{clip}%
\pgfsetbuttcap%
\pgfsetmiterjoin%
\definecolor{currentfill}{rgb}{0.121569,0.466667,0.705882}%
\pgfsetfillcolor{currentfill}%
\pgfsetlinewidth{0.000000pt}%
\definecolor{currentstroke}{rgb}{0.000000,0.000000,0.000000}%
\pgfsetstrokecolor{currentstroke}%
\pgfsetstrokeopacity{0.000000}%
\pgfsetdash{}{0pt}%
\pgfpathmoveto{\pgfqpoint{6.116460in}{0.733417in}}%
\pgfpathlineto{\pgfqpoint{6.180665in}{0.733417in}}%
\pgfpathlineto{\pgfqpoint{6.180665in}{0.773055in}}%
\pgfpathlineto{\pgfqpoint{6.116460in}{0.773055in}}%
\pgfpathlineto{\pgfqpoint{6.116460in}{0.733417in}}%
\pgfpathclose%
\pgfusepath{fill}%
\end{pgfscope}%
\begin{pgfscope}%
\pgfpathrectangle{\pgfqpoint{0.787430in}{0.733417in}}{\pgfqpoint{7.062570in}{2.555373in}}%
\pgfusepath{clip}%
\pgfsetbuttcap%
\pgfsetmiterjoin%
\definecolor{currentfill}{rgb}{0.121569,0.466667,0.705882}%
\pgfsetfillcolor{currentfill}%
\pgfsetlinewidth{0.000000pt}%
\definecolor{currentstroke}{rgb}{0.000000,0.000000,0.000000}%
\pgfsetstrokecolor{currentstroke}%
\pgfsetstrokeopacity{0.000000}%
\pgfsetdash{}{0pt}%
\pgfpathmoveto{\pgfqpoint{6.180665in}{0.733417in}}%
\pgfpathlineto{\pgfqpoint{6.244871in}{0.733417in}}%
\pgfpathlineto{\pgfqpoint{6.244871in}{0.762481in}}%
\pgfpathlineto{\pgfqpoint{6.180665in}{0.762481in}}%
\pgfpathlineto{\pgfqpoint{6.180665in}{0.733417in}}%
\pgfpathclose%
\pgfusepath{fill}%
\end{pgfscope}%
\begin{pgfscope}%
\pgfpathrectangle{\pgfqpoint{0.787430in}{0.733417in}}{\pgfqpoint{7.062570in}{2.555373in}}%
\pgfusepath{clip}%
\pgfsetbuttcap%
\pgfsetmiterjoin%
\definecolor{currentfill}{rgb}{0.121569,0.466667,0.705882}%
\pgfsetfillcolor{currentfill}%
\pgfsetlinewidth{0.000000pt}%
\definecolor{currentstroke}{rgb}{0.000000,0.000000,0.000000}%
\pgfsetstrokecolor{currentstroke}%
\pgfsetstrokeopacity{0.000000}%
\pgfsetdash{}{0pt}%
\pgfpathmoveto{\pgfqpoint{6.244871in}{0.733417in}}%
\pgfpathlineto{\pgfqpoint{6.309076in}{0.733417in}}%
\pgfpathlineto{\pgfqpoint{6.309076in}{0.754665in}}%
\pgfpathlineto{\pgfqpoint{6.244871in}{0.754665in}}%
\pgfpathlineto{\pgfqpoint{6.244871in}{0.733417in}}%
\pgfpathclose%
\pgfusepath{fill}%
\end{pgfscope}%
\begin{pgfscope}%
\pgfpathrectangle{\pgfqpoint{0.787430in}{0.733417in}}{\pgfqpoint{7.062570in}{2.555373in}}%
\pgfusepath{clip}%
\pgfsetbuttcap%
\pgfsetmiterjoin%
\definecolor{currentfill}{rgb}{0.121569,0.466667,0.705882}%
\pgfsetfillcolor{currentfill}%
\pgfsetlinewidth{0.000000pt}%
\definecolor{currentstroke}{rgb}{0.000000,0.000000,0.000000}%
\pgfsetstrokecolor{currentstroke}%
\pgfsetstrokeopacity{0.000000}%
\pgfsetdash{}{0pt}%
\pgfpathmoveto{\pgfqpoint{6.309076in}{0.733417in}}%
\pgfpathlineto{\pgfqpoint{6.373281in}{0.733417in}}%
\pgfpathlineto{\pgfqpoint{6.373281in}{0.748519in}}%
\pgfpathlineto{\pgfqpoint{6.309076in}{0.748519in}}%
\pgfpathlineto{\pgfqpoint{6.309076in}{0.733417in}}%
\pgfpathclose%
\pgfusepath{fill}%
\end{pgfscope}%
\begin{pgfscope}%
\pgfpathrectangle{\pgfqpoint{0.787430in}{0.733417in}}{\pgfqpoint{7.062570in}{2.555373in}}%
\pgfusepath{clip}%
\pgfsetbuttcap%
\pgfsetmiterjoin%
\definecolor{currentfill}{rgb}{0.121569,0.466667,0.705882}%
\pgfsetfillcolor{currentfill}%
\pgfsetlinewidth{0.000000pt}%
\definecolor{currentstroke}{rgb}{0.000000,0.000000,0.000000}%
\pgfsetstrokecolor{currentstroke}%
\pgfsetstrokeopacity{0.000000}%
\pgfsetdash{}{0pt}%
\pgfpathmoveto{\pgfqpoint{6.373281in}{0.733417in}}%
\pgfpathlineto{\pgfqpoint{6.437486in}{0.733417in}}%
\pgfpathlineto{\pgfqpoint{6.437486in}{0.744348in}}%
\pgfpathlineto{\pgfqpoint{6.373281in}{0.744348in}}%
\pgfpathlineto{\pgfqpoint{6.373281in}{0.733417in}}%
\pgfpathclose%
\pgfusepath{fill}%
\end{pgfscope}%
\begin{pgfscope}%
\pgfpathrectangle{\pgfqpoint{0.787430in}{0.733417in}}{\pgfqpoint{7.062570in}{2.555373in}}%
\pgfusepath{clip}%
\pgfsetbuttcap%
\pgfsetmiterjoin%
\definecolor{currentfill}{rgb}{0.121569,0.466667,0.705882}%
\pgfsetfillcolor{currentfill}%
\pgfsetlinewidth{0.000000pt}%
\definecolor{currentstroke}{rgb}{0.000000,0.000000,0.000000}%
\pgfsetstrokecolor{currentstroke}%
\pgfsetstrokeopacity{0.000000}%
\pgfsetdash{}{0pt}%
\pgfpathmoveto{\pgfqpoint{6.437486in}{0.733417in}}%
\pgfpathlineto{\pgfqpoint{6.501691in}{0.733417in}}%
\pgfpathlineto{\pgfqpoint{6.501691in}{0.741105in}}%
\pgfpathlineto{\pgfqpoint{6.437486in}{0.741105in}}%
\pgfpathlineto{\pgfqpoint{6.437486in}{0.733417in}}%
\pgfpathclose%
\pgfusepath{fill}%
\end{pgfscope}%
\begin{pgfscope}%
\pgfpathrectangle{\pgfqpoint{0.787430in}{0.733417in}}{\pgfqpoint{7.062570in}{2.555373in}}%
\pgfusepath{clip}%
\pgfsetbuttcap%
\pgfsetmiterjoin%
\definecolor{currentfill}{rgb}{0.121569,0.466667,0.705882}%
\pgfsetfillcolor{currentfill}%
\pgfsetlinewidth{0.000000pt}%
\definecolor{currentstroke}{rgb}{0.000000,0.000000,0.000000}%
\pgfsetstrokecolor{currentstroke}%
\pgfsetstrokeopacity{0.000000}%
\pgfsetdash{}{0pt}%
\pgfpathmoveto{\pgfqpoint{6.501691in}{0.733417in}}%
\pgfpathlineto{\pgfqpoint{6.565896in}{0.733417in}}%
\pgfpathlineto{\pgfqpoint{6.565896in}{0.738891in}}%
\pgfpathlineto{\pgfqpoint{6.501691in}{0.738891in}}%
\pgfpathlineto{\pgfqpoint{6.501691in}{0.733417in}}%
\pgfpathclose%
\pgfusepath{fill}%
\end{pgfscope}%
\begin{pgfscope}%
\pgfpathrectangle{\pgfqpoint{0.787430in}{0.733417in}}{\pgfqpoint{7.062570in}{2.555373in}}%
\pgfusepath{clip}%
\pgfsetbuttcap%
\pgfsetmiterjoin%
\definecolor{currentfill}{rgb}{0.121569,0.466667,0.705882}%
\pgfsetfillcolor{currentfill}%
\pgfsetlinewidth{0.000000pt}%
\definecolor{currentstroke}{rgb}{0.000000,0.000000,0.000000}%
\pgfsetstrokecolor{currentstroke}%
\pgfsetstrokeopacity{0.000000}%
\pgfsetdash{}{0pt}%
\pgfpathmoveto{\pgfqpoint{6.565896in}{0.733417in}}%
\pgfpathlineto{\pgfqpoint{6.630102in}{0.733417in}}%
\pgfpathlineto{\pgfqpoint{6.630102in}{0.737261in}}%
\pgfpathlineto{\pgfqpoint{6.565896in}{0.737261in}}%
\pgfpathlineto{\pgfqpoint{6.565896in}{0.733417in}}%
\pgfpathclose%
\pgfusepath{fill}%
\end{pgfscope}%
\begin{pgfscope}%
\pgfpathrectangle{\pgfqpoint{0.787430in}{0.733417in}}{\pgfqpoint{7.062570in}{2.555373in}}%
\pgfusepath{clip}%
\pgfsetbuttcap%
\pgfsetmiterjoin%
\definecolor{currentfill}{rgb}{0.121569,0.466667,0.705882}%
\pgfsetfillcolor{currentfill}%
\pgfsetlinewidth{0.000000pt}%
\definecolor{currentstroke}{rgb}{0.000000,0.000000,0.000000}%
\pgfsetstrokecolor{currentstroke}%
\pgfsetstrokeopacity{0.000000}%
\pgfsetdash{}{0pt}%
\pgfpathmoveto{\pgfqpoint{6.630102in}{0.733417in}}%
\pgfpathlineto{\pgfqpoint{6.694307in}{0.733417in}}%
\pgfpathlineto{\pgfqpoint{6.694307in}{0.736093in}}%
\pgfpathlineto{\pgfqpoint{6.630102in}{0.736093in}}%
\pgfpathlineto{\pgfqpoint{6.630102in}{0.733417in}}%
\pgfpathclose%
\pgfusepath{fill}%
\end{pgfscope}%
\begin{pgfscope}%
\pgfpathrectangle{\pgfqpoint{0.787430in}{0.733417in}}{\pgfqpoint{7.062570in}{2.555373in}}%
\pgfusepath{clip}%
\pgfsetbuttcap%
\pgfsetmiterjoin%
\definecolor{currentfill}{rgb}{0.121569,0.466667,0.705882}%
\pgfsetfillcolor{currentfill}%
\pgfsetlinewidth{0.000000pt}%
\definecolor{currentstroke}{rgb}{0.000000,0.000000,0.000000}%
\pgfsetstrokecolor{currentstroke}%
\pgfsetstrokeopacity{0.000000}%
\pgfsetdash{}{0pt}%
\pgfpathmoveto{\pgfqpoint{6.694307in}{0.733417in}}%
\pgfpathlineto{\pgfqpoint{6.758512in}{0.733417in}}%
\pgfpathlineto{\pgfqpoint{6.758512in}{0.735135in}}%
\pgfpathlineto{\pgfqpoint{6.694307in}{0.735135in}}%
\pgfpathlineto{\pgfqpoint{6.694307in}{0.733417in}}%
\pgfpathclose%
\pgfusepath{fill}%
\end{pgfscope}%
\begin{pgfscope}%
\pgfpathrectangle{\pgfqpoint{0.787430in}{0.733417in}}{\pgfqpoint{7.062570in}{2.555373in}}%
\pgfusepath{clip}%
\pgfsetbuttcap%
\pgfsetmiterjoin%
\definecolor{currentfill}{rgb}{0.121569,0.466667,0.705882}%
\pgfsetfillcolor{currentfill}%
\pgfsetlinewidth{0.000000pt}%
\definecolor{currentstroke}{rgb}{0.000000,0.000000,0.000000}%
\pgfsetstrokecolor{currentstroke}%
\pgfsetstrokeopacity{0.000000}%
\pgfsetdash{}{0pt}%
\pgfpathmoveto{\pgfqpoint{6.758512in}{0.733417in}}%
\pgfpathlineto{\pgfqpoint{6.822717in}{0.733417in}}%
\pgfpathlineto{\pgfqpoint{6.822717in}{0.734416in}}%
\pgfpathlineto{\pgfqpoint{6.758512in}{0.734416in}}%
\pgfpathlineto{\pgfqpoint{6.758512in}{0.733417in}}%
\pgfpathclose%
\pgfusepath{fill}%
\end{pgfscope}%
\begin{pgfscope}%
\pgfpathrectangle{\pgfqpoint{0.787430in}{0.733417in}}{\pgfqpoint{7.062570in}{2.555373in}}%
\pgfusepath{clip}%
\pgfsetbuttcap%
\pgfsetmiterjoin%
\definecolor{currentfill}{rgb}{0.121569,0.466667,0.705882}%
\pgfsetfillcolor{currentfill}%
\pgfsetlinewidth{0.000000pt}%
\definecolor{currentstroke}{rgb}{0.000000,0.000000,0.000000}%
\pgfsetstrokecolor{currentstroke}%
\pgfsetstrokeopacity{0.000000}%
\pgfsetdash{}{0pt}%
\pgfpathmoveto{\pgfqpoint{6.822717in}{0.733417in}}%
\pgfpathlineto{\pgfqpoint{6.886922in}{0.733417in}}%
\pgfpathlineto{\pgfqpoint{6.886922in}{0.734194in}}%
\pgfpathlineto{\pgfqpoint{6.822717in}{0.734194in}}%
\pgfpathlineto{\pgfqpoint{6.822717in}{0.733417in}}%
\pgfpathclose%
\pgfusepath{fill}%
\end{pgfscope}%
\begin{pgfscope}%
\pgfpathrectangle{\pgfqpoint{0.787430in}{0.733417in}}{\pgfqpoint{7.062570in}{2.555373in}}%
\pgfusepath{clip}%
\pgfsetbuttcap%
\pgfsetmiterjoin%
\definecolor{currentfill}{rgb}{0.121569,0.466667,0.705882}%
\pgfsetfillcolor{currentfill}%
\pgfsetlinewidth{0.000000pt}%
\definecolor{currentstroke}{rgb}{0.000000,0.000000,0.000000}%
\pgfsetstrokecolor{currentstroke}%
\pgfsetstrokeopacity{0.000000}%
\pgfsetdash{}{0pt}%
\pgfpathmoveto{\pgfqpoint{6.886922in}{0.733417in}}%
\pgfpathlineto{\pgfqpoint{6.951127in}{0.733417in}}%
\pgfpathlineto{\pgfqpoint{6.951127in}{0.733961in}}%
\pgfpathlineto{\pgfqpoint{6.886922in}{0.733961in}}%
\pgfpathlineto{\pgfqpoint{6.886922in}{0.733417in}}%
\pgfpathclose%
\pgfusepath{fill}%
\end{pgfscope}%
\begin{pgfscope}%
\pgfpathrectangle{\pgfqpoint{0.787430in}{0.733417in}}{\pgfqpoint{7.062570in}{2.555373in}}%
\pgfusepath{clip}%
\pgfsetbuttcap%
\pgfsetmiterjoin%
\definecolor{currentfill}{rgb}{0.121569,0.466667,0.705882}%
\pgfsetfillcolor{currentfill}%
\pgfsetlinewidth{0.000000pt}%
\definecolor{currentstroke}{rgb}{0.000000,0.000000,0.000000}%
\pgfsetstrokecolor{currentstroke}%
\pgfsetstrokeopacity{0.000000}%
\pgfsetdash{}{0pt}%
\pgfpathmoveto{\pgfqpoint{6.951127in}{0.733417in}}%
\pgfpathlineto{\pgfqpoint{7.015333in}{0.733417in}}%
\pgfpathlineto{\pgfqpoint{7.015333in}{0.733756in}}%
\pgfpathlineto{\pgfqpoint{6.951127in}{0.733756in}}%
\pgfpathlineto{\pgfqpoint{6.951127in}{0.733417in}}%
\pgfpathclose%
\pgfusepath{fill}%
\end{pgfscope}%
\begin{pgfscope}%
\pgfpathrectangle{\pgfqpoint{0.787430in}{0.733417in}}{\pgfqpoint{7.062570in}{2.555373in}}%
\pgfusepath{clip}%
\pgfsetbuttcap%
\pgfsetmiterjoin%
\definecolor{currentfill}{rgb}{0.121569,0.466667,0.705882}%
\pgfsetfillcolor{currentfill}%
\pgfsetlinewidth{0.000000pt}%
\definecolor{currentstroke}{rgb}{0.000000,0.000000,0.000000}%
\pgfsetstrokecolor{currentstroke}%
\pgfsetstrokeopacity{0.000000}%
\pgfsetdash{}{0pt}%
\pgfpathmoveto{\pgfqpoint{7.015333in}{0.733417in}}%
\pgfpathlineto{\pgfqpoint{7.079538in}{0.733417in}}%
\pgfpathlineto{\pgfqpoint{7.079538in}{0.733616in}}%
\pgfpathlineto{\pgfqpoint{7.015333in}{0.733616in}}%
\pgfpathlineto{\pgfqpoint{7.015333in}{0.733417in}}%
\pgfpathclose%
\pgfusepath{fill}%
\end{pgfscope}%
\begin{pgfscope}%
\pgfpathrectangle{\pgfqpoint{0.787430in}{0.733417in}}{\pgfqpoint{7.062570in}{2.555373in}}%
\pgfusepath{clip}%
\pgfsetbuttcap%
\pgfsetmiterjoin%
\definecolor{currentfill}{rgb}{0.121569,0.466667,0.705882}%
\pgfsetfillcolor{currentfill}%
\pgfsetlinewidth{0.000000pt}%
\definecolor{currentstroke}{rgb}{0.000000,0.000000,0.000000}%
\pgfsetstrokecolor{currentstroke}%
\pgfsetstrokeopacity{0.000000}%
\pgfsetdash{}{0pt}%
\pgfpathmoveto{\pgfqpoint{7.079538in}{0.733417in}}%
\pgfpathlineto{\pgfqpoint{7.143743in}{0.733417in}}%
\pgfpathlineto{\pgfqpoint{7.143743in}{0.733534in}}%
\pgfpathlineto{\pgfqpoint{7.079538in}{0.733534in}}%
\pgfpathlineto{\pgfqpoint{7.079538in}{0.733417in}}%
\pgfpathclose%
\pgfusepath{fill}%
\end{pgfscope}%
\begin{pgfscope}%
\pgfpathrectangle{\pgfqpoint{0.787430in}{0.733417in}}{\pgfqpoint{7.062570in}{2.555373in}}%
\pgfusepath{clip}%
\pgfsetbuttcap%
\pgfsetmiterjoin%
\definecolor{currentfill}{rgb}{0.121569,0.466667,0.705882}%
\pgfsetfillcolor{currentfill}%
\pgfsetlinewidth{0.000000pt}%
\definecolor{currentstroke}{rgb}{0.000000,0.000000,0.000000}%
\pgfsetstrokecolor{currentstroke}%
\pgfsetstrokeopacity{0.000000}%
\pgfsetdash{}{0pt}%
\pgfpathmoveto{\pgfqpoint{7.143743in}{0.733417in}}%
\pgfpathlineto{\pgfqpoint{7.207948in}{0.733417in}}%
\pgfpathlineto{\pgfqpoint{7.207948in}{0.733464in}}%
\pgfpathlineto{\pgfqpoint{7.143743in}{0.733464in}}%
\pgfpathlineto{\pgfqpoint{7.143743in}{0.733417in}}%
\pgfpathclose%
\pgfusepath{fill}%
\end{pgfscope}%
\begin{pgfscope}%
\pgfpathrectangle{\pgfqpoint{0.787430in}{0.733417in}}{\pgfqpoint{7.062570in}{2.555373in}}%
\pgfusepath{clip}%
\pgfsetbuttcap%
\pgfsetmiterjoin%
\definecolor{currentfill}{rgb}{0.121569,0.466667,0.705882}%
\pgfsetfillcolor{currentfill}%
\pgfsetlinewidth{0.000000pt}%
\definecolor{currentstroke}{rgb}{0.000000,0.000000,0.000000}%
\pgfsetstrokecolor{currentstroke}%
\pgfsetstrokeopacity{0.000000}%
\pgfsetdash{}{0pt}%
\pgfpathmoveto{\pgfqpoint{7.207948in}{0.733417in}}%
\pgfpathlineto{\pgfqpoint{7.272153in}{0.733417in}}%
\pgfpathlineto{\pgfqpoint{7.272153in}{0.733476in}}%
\pgfpathlineto{\pgfqpoint{7.207948in}{0.733476in}}%
\pgfpathlineto{\pgfqpoint{7.207948in}{0.733417in}}%
\pgfpathclose%
\pgfusepath{fill}%
\end{pgfscope}%
\begin{pgfscope}%
\pgfpathrectangle{\pgfqpoint{0.787430in}{0.733417in}}{\pgfqpoint{7.062570in}{2.555373in}}%
\pgfusepath{clip}%
\pgfsetbuttcap%
\pgfsetmiterjoin%
\definecolor{currentfill}{rgb}{0.121569,0.466667,0.705882}%
\pgfsetfillcolor{currentfill}%
\pgfsetlinewidth{0.000000pt}%
\definecolor{currentstroke}{rgb}{0.000000,0.000000,0.000000}%
\pgfsetstrokecolor{currentstroke}%
\pgfsetstrokeopacity{0.000000}%
\pgfsetdash{}{0pt}%
\pgfpathmoveto{\pgfqpoint{7.272153in}{0.733417in}}%
\pgfpathlineto{\pgfqpoint{7.336359in}{0.733417in}}%
\pgfpathlineto{\pgfqpoint{7.336359in}{0.733452in}}%
\pgfpathlineto{\pgfqpoint{7.272153in}{0.733452in}}%
\pgfpathlineto{\pgfqpoint{7.272153in}{0.733417in}}%
\pgfpathclose%
\pgfusepath{fill}%
\end{pgfscope}%
\begin{pgfscope}%
\pgfpathrectangle{\pgfqpoint{0.787430in}{0.733417in}}{\pgfqpoint{7.062570in}{2.555373in}}%
\pgfusepath{clip}%
\pgfsetbuttcap%
\pgfsetmiterjoin%
\definecolor{currentfill}{rgb}{0.121569,0.466667,0.705882}%
\pgfsetfillcolor{currentfill}%
\pgfsetlinewidth{0.000000pt}%
\definecolor{currentstroke}{rgb}{0.000000,0.000000,0.000000}%
\pgfsetstrokecolor{currentstroke}%
\pgfsetstrokeopacity{0.000000}%
\pgfsetdash{}{0pt}%
\pgfpathmoveto{\pgfqpoint{7.336359in}{0.733417in}}%
\pgfpathlineto{\pgfqpoint{7.400564in}{0.733417in}}%
\pgfpathlineto{\pgfqpoint{7.400564in}{0.733452in}}%
\pgfpathlineto{\pgfqpoint{7.336359in}{0.733452in}}%
\pgfpathlineto{\pgfqpoint{7.336359in}{0.733417in}}%
\pgfpathclose%
\pgfusepath{fill}%
\end{pgfscope}%
\begin{pgfscope}%
\pgfpathrectangle{\pgfqpoint{0.787430in}{0.733417in}}{\pgfqpoint{7.062570in}{2.555373in}}%
\pgfusepath{clip}%
\pgfsetbuttcap%
\pgfsetmiterjoin%
\definecolor{currentfill}{rgb}{0.121569,0.466667,0.705882}%
\pgfsetfillcolor{currentfill}%
\pgfsetlinewidth{0.000000pt}%
\definecolor{currentstroke}{rgb}{0.000000,0.000000,0.000000}%
\pgfsetstrokecolor{currentstroke}%
\pgfsetstrokeopacity{0.000000}%
\pgfsetdash{}{0pt}%
\pgfpathmoveto{\pgfqpoint{7.400564in}{0.733417in}}%
\pgfpathlineto{\pgfqpoint{7.464769in}{0.733417in}}%
\pgfpathlineto{\pgfqpoint{7.464769in}{0.733435in}}%
\pgfpathlineto{\pgfqpoint{7.400564in}{0.733435in}}%
\pgfpathlineto{\pgfqpoint{7.400564in}{0.733417in}}%
\pgfpathclose%
\pgfusepath{fill}%
\end{pgfscope}%
\begin{pgfscope}%
\pgfpathrectangle{\pgfqpoint{0.787430in}{0.733417in}}{\pgfqpoint{7.062570in}{2.555373in}}%
\pgfusepath{clip}%
\pgfsetbuttcap%
\pgfsetmiterjoin%
\definecolor{currentfill}{rgb}{0.121569,0.466667,0.705882}%
\pgfsetfillcolor{currentfill}%
\pgfsetlinewidth{0.000000pt}%
\definecolor{currentstroke}{rgb}{0.000000,0.000000,0.000000}%
\pgfsetstrokecolor{currentstroke}%
\pgfsetstrokeopacity{0.000000}%
\pgfsetdash{}{0pt}%
\pgfpathmoveto{\pgfqpoint{7.464769in}{0.733417in}}%
\pgfpathlineto{\pgfqpoint{7.528974in}{0.733417in}}%
\pgfpathlineto{\pgfqpoint{7.528974in}{0.733429in}}%
\pgfpathlineto{\pgfqpoint{7.464769in}{0.733429in}}%
\pgfpathlineto{\pgfqpoint{7.464769in}{0.733417in}}%
\pgfpathclose%
\pgfusepath{fill}%
\end{pgfscope}%
\begin{pgfscope}%
\pgfsetbuttcap%
\pgfsetroundjoin%
\definecolor{currentfill}{rgb}{0.000000,0.000000,0.000000}%
\pgfsetfillcolor{currentfill}%
\pgfsetlinewidth{0.803000pt}%
\definecolor{currentstroke}{rgb}{0.000000,0.000000,0.000000}%
\pgfsetstrokecolor{currentstroke}%
\pgfsetdash{}{0pt}%
\pgfsys@defobject{currentmarker}{\pgfqpoint{0.000000in}{-0.048611in}}{\pgfqpoint{0.000000in}{0.000000in}}{%
\pgfpathmoveto{\pgfqpoint{0.000000in}{0.000000in}}%
\pgfpathlineto{\pgfqpoint{0.000000in}{-0.048611in}}%
\pgfusepath{stroke,fill}%
}%
\begin{pgfscope}%
\pgfsys@transformshift{1.923827in}{0.733417in}%
\pgfsys@useobject{currentmarker}{}%
\end{pgfscope}%
\end{pgfscope}%
\begin{pgfscope}%
\definecolor{textcolor}{rgb}{0.000000,0.000000,0.000000}%
\pgfsetstrokecolor{textcolor}%
\pgfsetfillcolor{textcolor}%
\pgftext[x=1.923827in,y=0.636195in,,top]{\color{textcolor}{\rmfamily\fontsize{10.000000}{12.000000}\selectfont\catcode`\^=\active\def^{\ifmmode\sp\else\^{}\fi}\catcode`\%=\active\def%{\%}$\mathdefault{\ensuremath{-}4}$}}%
\end{pgfscope}%
\begin{pgfscope}%
\pgfsetbuttcap%
\pgfsetroundjoin%
\definecolor{currentfill}{rgb}{0.000000,0.000000,0.000000}%
\pgfsetfillcolor{currentfill}%
\pgfsetlinewidth{0.803000pt}%
\definecolor{currentstroke}{rgb}{0.000000,0.000000,0.000000}%
\pgfsetstrokecolor{currentstroke}%
\pgfsetdash{}{0pt}%
\pgfsys@defobject{currentmarker}{\pgfqpoint{0.000000in}{-0.048611in}}{\pgfqpoint{0.000000in}{0.000000in}}{%
\pgfpathmoveto{\pgfqpoint{0.000000in}{0.000000in}}%
\pgfpathlineto{\pgfqpoint{0.000000in}{-0.048611in}}%
\pgfusepath{stroke,fill}%
}%
\begin{pgfscope}%
\pgfsys@transformshift{3.153546in}{0.733417in}%
\pgfsys@useobject{currentmarker}{}%
\end{pgfscope}%
\end{pgfscope}%
\begin{pgfscope}%
\definecolor{textcolor}{rgb}{0.000000,0.000000,0.000000}%
\pgfsetstrokecolor{textcolor}%
\pgfsetfillcolor{textcolor}%
\pgftext[x=3.153546in,y=0.636195in,,top]{\color{textcolor}{\rmfamily\fontsize{10.000000}{12.000000}\selectfont\catcode`\^=\active\def^{\ifmmode\sp\else\^{}\fi}\catcode`\%=\active\def%{\%}$\mathdefault{\ensuremath{-}2}$}}%
\end{pgfscope}%
\begin{pgfscope}%
\pgfsetbuttcap%
\pgfsetroundjoin%
\definecolor{currentfill}{rgb}{0.000000,0.000000,0.000000}%
\pgfsetfillcolor{currentfill}%
\pgfsetlinewidth{0.803000pt}%
\definecolor{currentstroke}{rgb}{0.000000,0.000000,0.000000}%
\pgfsetstrokecolor{currentstroke}%
\pgfsetdash{}{0pt}%
\pgfsys@defobject{currentmarker}{\pgfqpoint{0.000000in}{-0.048611in}}{\pgfqpoint{0.000000in}{0.000000in}}{%
\pgfpathmoveto{\pgfqpoint{0.000000in}{0.000000in}}%
\pgfpathlineto{\pgfqpoint{0.000000in}{-0.048611in}}%
\pgfusepath{stroke,fill}%
}%
\begin{pgfscope}%
\pgfsys@transformshift{4.383264in}{0.733417in}%
\pgfsys@useobject{currentmarker}{}%
\end{pgfscope}%
\end{pgfscope}%
\begin{pgfscope}%
\definecolor{textcolor}{rgb}{0.000000,0.000000,0.000000}%
\pgfsetstrokecolor{textcolor}%
\pgfsetfillcolor{textcolor}%
\pgftext[x=4.383264in,y=0.636195in,,top]{\color{textcolor}{\rmfamily\fontsize{10.000000}{12.000000}\selectfont\catcode`\^=\active\def^{\ifmmode\sp\else\^{}\fi}\catcode`\%=\active\def%{\%}$\mathdefault{0}$}}%
\end{pgfscope}%
\begin{pgfscope}%
\pgfsetbuttcap%
\pgfsetroundjoin%
\definecolor{currentfill}{rgb}{0.000000,0.000000,0.000000}%
\pgfsetfillcolor{currentfill}%
\pgfsetlinewidth{0.803000pt}%
\definecolor{currentstroke}{rgb}{0.000000,0.000000,0.000000}%
\pgfsetstrokecolor{currentstroke}%
\pgfsetdash{}{0pt}%
\pgfsys@defobject{currentmarker}{\pgfqpoint{0.000000in}{-0.048611in}}{\pgfqpoint{0.000000in}{0.000000in}}{%
\pgfpathmoveto{\pgfqpoint{0.000000in}{0.000000in}}%
\pgfpathlineto{\pgfqpoint{0.000000in}{-0.048611in}}%
\pgfusepath{stroke,fill}%
}%
\begin{pgfscope}%
\pgfsys@transformshift{5.612983in}{0.733417in}%
\pgfsys@useobject{currentmarker}{}%
\end{pgfscope}%
\end{pgfscope}%
\begin{pgfscope}%
\definecolor{textcolor}{rgb}{0.000000,0.000000,0.000000}%
\pgfsetstrokecolor{textcolor}%
\pgfsetfillcolor{textcolor}%
\pgftext[x=5.612983in,y=0.636195in,,top]{\color{textcolor}{\rmfamily\fontsize{10.000000}{12.000000}\selectfont\catcode`\^=\active\def^{\ifmmode\sp\else\^{}\fi}\catcode`\%=\active\def%{\%}$\mathdefault{2}$}}%
\end{pgfscope}%
\begin{pgfscope}%
\pgfsetbuttcap%
\pgfsetroundjoin%
\definecolor{currentfill}{rgb}{0.000000,0.000000,0.000000}%
\pgfsetfillcolor{currentfill}%
\pgfsetlinewidth{0.803000pt}%
\definecolor{currentstroke}{rgb}{0.000000,0.000000,0.000000}%
\pgfsetstrokecolor{currentstroke}%
\pgfsetdash{}{0pt}%
\pgfsys@defobject{currentmarker}{\pgfqpoint{0.000000in}{-0.048611in}}{\pgfqpoint{0.000000in}{0.000000in}}{%
\pgfpathmoveto{\pgfqpoint{0.000000in}{0.000000in}}%
\pgfpathlineto{\pgfqpoint{0.000000in}{-0.048611in}}%
\pgfusepath{stroke,fill}%
}%
\begin{pgfscope}%
\pgfsys@transformshift{6.842701in}{0.733417in}%
\pgfsys@useobject{currentmarker}{}%
\end{pgfscope}%
\end{pgfscope}%
\begin{pgfscope}%
\definecolor{textcolor}{rgb}{0.000000,0.000000,0.000000}%
\pgfsetstrokecolor{textcolor}%
\pgfsetfillcolor{textcolor}%
\pgftext[x=6.842701in,y=0.636195in,,top]{\color{textcolor}{\rmfamily\fontsize{10.000000}{12.000000}\selectfont\catcode`\^=\active\def^{\ifmmode\sp\else\^{}\fi}\catcode`\%=\active\def%{\%}$\mathdefault{4}$}}%
\end{pgfscope}%
\begin{pgfscope}%
\definecolor{textcolor}{rgb}{0.000000,0.000000,0.000000}%
\pgfsetstrokecolor{textcolor}%
\pgfsetfillcolor{textcolor}%
\pgftext[x=4.318715in,y=0.457183in,,top]{\color{textcolor}{\rmfamily\fontsize{24.000000}{28.800000}\selectfont\catcode`\^=\active\def^{\ifmmode\sp\else\^{}\fi}\catcode`\%=\active\def%{\%}$x$}}%
\end{pgfscope}%
\begin{pgfscope}%
\pgfsetbuttcap%
\pgfsetroundjoin%
\definecolor{currentfill}{rgb}{0.000000,0.000000,0.000000}%
\pgfsetfillcolor{currentfill}%
\pgfsetlinewidth{0.803000pt}%
\definecolor{currentstroke}{rgb}{0.000000,0.000000,0.000000}%
\pgfsetstrokecolor{currentstroke}%
\pgfsetdash{}{0pt}%
\pgfsys@defobject{currentmarker}{\pgfqpoint{-0.048611in}{0.000000in}}{\pgfqpoint{-0.000000in}{0.000000in}}{%
\pgfpathmoveto{\pgfqpoint{-0.000000in}{0.000000in}}%
\pgfpathlineto{\pgfqpoint{-0.048611in}{0.000000in}}%
\pgfusepath{stroke,fill}%
}%
\begin{pgfscope}%
\pgfsys@transformshift{0.787430in}{0.733417in}%
\pgfsys@useobject{currentmarker}{}%
\end{pgfscope}%
\end{pgfscope}%
\begin{pgfscope}%
\definecolor{textcolor}{rgb}{0.000000,0.000000,0.000000}%
\pgfsetstrokecolor{textcolor}%
\pgfsetfillcolor{textcolor}%
\pgftext[x=0.512738in, y=0.685192in, left, base]{\color{textcolor}{\rmfamily\fontsize{10.000000}{12.000000}\selectfont\catcode`\^=\active\def^{\ifmmode\sp\else\^{}\fi}\catcode`\%=\active\def%{\%}$\mathdefault{0.0}$}}%
\end{pgfscope}%
\begin{pgfscope}%
\pgfsetbuttcap%
\pgfsetroundjoin%
\definecolor{currentfill}{rgb}{0.000000,0.000000,0.000000}%
\pgfsetfillcolor{currentfill}%
\pgfsetlinewidth{0.803000pt}%
\definecolor{currentstroke}{rgb}{0.000000,0.000000,0.000000}%
\pgfsetstrokecolor{currentstroke}%
\pgfsetdash{}{0pt}%
\pgfsys@defobject{currentmarker}{\pgfqpoint{-0.048611in}{0.000000in}}{\pgfqpoint{-0.000000in}{0.000000in}}{%
\pgfpathmoveto{\pgfqpoint{-0.000000in}{0.000000in}}%
\pgfpathlineto{\pgfqpoint{-0.048611in}{0.000000in}}%
\pgfusepath{stroke,fill}%
}%
\begin{pgfscope}%
\pgfsys@transformshift{0.787430in}{1.953498in}%
\pgfsys@useobject{currentmarker}{}%
\end{pgfscope}%
\end{pgfscope}%
\begin{pgfscope}%
\definecolor{textcolor}{rgb}{0.000000,0.000000,0.000000}%
\pgfsetstrokecolor{textcolor}%
\pgfsetfillcolor{textcolor}%
\pgftext[x=0.512738in, y=1.905272in, left, base]{\color{textcolor}{\rmfamily\fontsize{10.000000}{12.000000}\selectfont\catcode`\^=\active\def^{\ifmmode\sp\else\^{}\fi}\catcode`\%=\active\def%{\%}$\mathdefault{0.2}$}}%
\end{pgfscope}%
\begin{pgfscope}%
\pgfsetbuttcap%
\pgfsetroundjoin%
\definecolor{currentfill}{rgb}{0.000000,0.000000,0.000000}%
\pgfsetfillcolor{currentfill}%
\pgfsetlinewidth{0.803000pt}%
\definecolor{currentstroke}{rgb}{0.000000,0.000000,0.000000}%
\pgfsetstrokecolor{currentstroke}%
\pgfsetdash{}{0pt}%
\pgfsys@defobject{currentmarker}{\pgfqpoint{-0.048611in}{0.000000in}}{\pgfqpoint{-0.000000in}{0.000000in}}{%
\pgfpathmoveto{\pgfqpoint{-0.000000in}{0.000000in}}%
\pgfpathlineto{\pgfqpoint{-0.048611in}{0.000000in}}%
\pgfusepath{stroke,fill}%
}%
\begin{pgfscope}%
\pgfsys@transformshift{0.787430in}{3.173578in}%
\pgfsys@useobject{currentmarker}{}%
\end{pgfscope}%
\end{pgfscope}%
\begin{pgfscope}%
\definecolor{textcolor}{rgb}{0.000000,0.000000,0.000000}%
\pgfsetstrokecolor{textcolor}%
\pgfsetfillcolor{textcolor}%
\pgftext[x=0.512738in, y=3.125353in, left, base]{\color{textcolor}{\rmfamily\fontsize{10.000000}{12.000000}\selectfont\catcode`\^=\active\def^{\ifmmode\sp\else\^{}\fi}\catcode`\%=\active\def%{\%}$\mathdefault{0.4}$}}%
\end{pgfscope}%
\begin{pgfscope}%
\pgfpathrectangle{\pgfqpoint{0.787430in}{0.733417in}}{\pgfqpoint{7.062570in}{2.555373in}}%
\pgfusepath{clip}%
\pgfsetrectcap%
\pgfsetroundjoin%
\pgfsetlinewidth{2.007500pt}%
\definecolor{currentstroke}{rgb}{1.000000,0.000000,0.000000}%
\pgfsetstrokecolor{currentstroke}%
\pgfsetdash{}{0pt}%
\pgfpathmoveto{\pgfqpoint{1.923827in}{0.734234in}}%
\pgfpathlineto{\pgfqpoint{2.135551in}{0.736468in}}%
\pgfpathlineto{\pgfqpoint{2.268493in}{0.739986in}}%
\pgfpathlineto{\pgfqpoint{2.366969in}{0.744667in}}%
\pgfpathlineto{\pgfqpoint{2.450674in}{0.750835in}}%
\pgfpathlineto{\pgfqpoint{2.519607in}{0.758038in}}%
\pgfpathlineto{\pgfqpoint{2.583616in}{0.766990in}}%
\pgfpathlineto{\pgfqpoint{2.637778in}{0.776696in}}%
\pgfpathlineto{\pgfqpoint{2.687016in}{0.787569in}}%
\pgfpathlineto{\pgfqpoint{2.736254in}{0.800741in}}%
\pgfpathlineto{\pgfqpoint{2.780568in}{0.814866in}}%
\pgfpathlineto{\pgfqpoint{2.824882in}{0.831444in}}%
\pgfpathlineto{\pgfqpoint{2.864273in}{0.848489in}}%
\pgfpathlineto{\pgfqpoint{2.903663in}{0.867946in}}%
\pgfpathlineto{\pgfqpoint{2.943053in}{0.890047in}}%
\pgfpathlineto{\pgfqpoint{2.977520in}{0.911744in}}%
\pgfpathlineto{\pgfqpoint{3.011987in}{0.935810in}}%
\pgfpathlineto{\pgfqpoint{3.046453in}{0.962402in}}%
\pgfpathlineto{\pgfqpoint{3.080920in}{0.991675in}}%
\pgfpathlineto{\pgfqpoint{3.115386in}{1.023777in}}%
\pgfpathlineto{\pgfqpoint{3.149853in}{1.058845in}}%
\pgfpathlineto{\pgfqpoint{3.184320in}{1.097004in}}%
\pgfpathlineto{\pgfqpoint{3.218786in}{1.138364in}}%
\pgfpathlineto{\pgfqpoint{3.253253in}{1.183012in}}%
\pgfpathlineto{\pgfqpoint{3.287719in}{1.231018in}}%
\pgfpathlineto{\pgfqpoint{3.327110in}{1.290044in}}%
\pgfpathlineto{\pgfqpoint{3.366500in}{1.353521in}}%
\pgfpathlineto{\pgfqpoint{3.405890in}{1.421407in}}%
\pgfpathlineto{\pgfqpoint{3.445281in}{1.493599in}}%
\pgfpathlineto{\pgfqpoint{3.489595in}{1.579746in}}%
\pgfpathlineto{\pgfqpoint{3.533909in}{1.670774in}}%
\pgfpathlineto{\pgfqpoint{3.583147in}{1.777067in}}%
\pgfpathlineto{\pgfqpoint{3.637309in}{1.899294in}}%
\pgfpathlineto{\pgfqpoint{3.706242in}{2.060794in}}%
\pgfpathlineto{\pgfqpoint{3.913042in}{2.550048in}}%
\pgfpathlineto{\pgfqpoint{3.962280in}{2.658598in}}%
\pgfpathlineto{\pgfqpoint{4.006594in}{2.750736in}}%
\pgfpathlineto{\pgfqpoint{4.045984in}{2.827181in}}%
\pgfpathlineto{\pgfqpoint{4.080451in}{2.889174in}}%
\pgfpathlineto{\pgfqpoint{4.109994in}{2.938248in}}%
\pgfpathlineto{\pgfqpoint{4.139536in}{2.983239in}}%
\pgfpathlineto{\pgfqpoint{4.169079in}{3.023855in}}%
\pgfpathlineto{\pgfqpoint{4.193698in}{3.054164in}}%
\pgfpathlineto{\pgfqpoint{4.218317in}{3.081108in}}%
\pgfpathlineto{\pgfqpoint{4.242936in}{3.104560in}}%
\pgfpathlineto{\pgfqpoint{4.267555in}{3.124410in}}%
\pgfpathlineto{\pgfqpoint{4.287250in}{3.137633in}}%
\pgfpathlineto{\pgfqpoint{4.306945in}{3.148450in}}%
\pgfpathlineto{\pgfqpoint{4.326641in}{3.156827in}}%
\pgfpathlineto{\pgfqpoint{4.346336in}{3.162740in}}%
\pgfpathlineto{\pgfqpoint{4.366031in}{3.166170in}}%
\pgfpathlineto{\pgfqpoint{4.385726in}{3.167106in}}%
\pgfpathlineto{\pgfqpoint{4.405421in}{3.165546in}}%
\pgfpathlineto{\pgfqpoint{4.425117in}{3.161494in}}%
\pgfpathlineto{\pgfqpoint{4.444812in}{3.154963in}}%
\pgfpathlineto{\pgfqpoint{4.464507in}{3.145973in}}%
\pgfpathlineto{\pgfqpoint{4.484202in}{3.134552in}}%
\pgfpathlineto{\pgfqpoint{4.503897in}{3.120733in}}%
\pgfpathlineto{\pgfqpoint{4.523593in}{3.104560in}}%
\pgfpathlineto{\pgfqpoint{4.548212in}{3.081108in}}%
\pgfpathlineto{\pgfqpoint{4.572831in}{3.054164in}}%
\pgfpathlineto{\pgfqpoint{4.597450in}{3.023855in}}%
\pgfpathlineto{\pgfqpoint{4.622069in}{2.990320in}}%
\pgfpathlineto{\pgfqpoint{4.651611in}{2.946038in}}%
\pgfpathlineto{\pgfqpoint{4.681154in}{2.897623in}}%
\pgfpathlineto{\pgfqpoint{4.715621in}{2.836332in}}%
\pgfpathlineto{\pgfqpoint{4.750087in}{2.770365in}}%
\pgfpathlineto{\pgfqpoint{4.789478in}{2.689962in}}%
\pgfpathlineto{\pgfqpoint{4.833792in}{2.594145in}}%
\pgfpathlineto{\pgfqpoint{4.887954in}{2.471083in}}%
\pgfpathlineto{\pgfqpoint{4.956887in}{2.308384in}}%
\pgfpathlineto{\pgfqpoint{5.124296in}{1.910638in}}%
\pgfpathlineto{\pgfqpoint{5.183381in}{1.777067in}}%
\pgfpathlineto{\pgfqpoint{5.232619in}{1.670774in}}%
\pgfpathlineto{\pgfqpoint{5.281857in}{1.569926in}}%
\pgfpathlineto{\pgfqpoint{5.326172in}{1.484345in}}%
\pgfpathlineto{\pgfqpoint{5.365562in}{1.412683in}}%
\pgfpathlineto{\pgfqpoint{5.404952in}{1.345343in}}%
\pgfpathlineto{\pgfqpoint{5.444343in}{1.282422in}}%
\pgfpathlineto{\pgfqpoint{5.483733in}{1.223953in}}%
\pgfpathlineto{\pgfqpoint{5.518200in}{1.176430in}}%
\pgfpathlineto{\pgfqpoint{5.552666in}{1.132256in}}%
\pgfpathlineto{\pgfqpoint{5.587133in}{1.091360in}}%
\pgfpathlineto{\pgfqpoint{5.621599in}{1.053649in}}%
\pgfpathlineto{\pgfqpoint{5.656066in}{1.019013in}}%
\pgfpathlineto{\pgfqpoint{5.690533in}{0.987324in}}%
\pgfpathlineto{\pgfqpoint{5.724999in}{0.958442in}}%
\pgfpathlineto{\pgfqpoint{5.759466in}{0.932221in}}%
\pgfpathlineto{\pgfqpoint{5.793932in}{0.908503in}}%
\pgfpathlineto{\pgfqpoint{5.833323in}{0.884261in}}%
\pgfpathlineto{\pgfqpoint{5.872713in}{0.862843in}}%
\pgfpathlineto{\pgfqpoint{5.912104in}{0.844011in}}%
\pgfpathlineto{\pgfqpoint{5.951494in}{0.827532in}}%
\pgfpathlineto{\pgfqpoint{5.995808in}{0.811526in}}%
\pgfpathlineto{\pgfqpoint{6.040122in}{0.797905in}}%
\pgfpathlineto{\pgfqpoint{6.089360in}{0.785222in}}%
\pgfpathlineto{\pgfqpoint{6.138598in}{0.774767in}}%
\pgfpathlineto{\pgfqpoint{6.192760in}{0.765449in}}%
\pgfpathlineto{\pgfqpoint{6.251846in}{0.757447in}}%
\pgfpathlineto{\pgfqpoint{6.320779in}{0.750402in}}%
\pgfpathlineto{\pgfqpoint{6.399560in}{0.744667in}}%
\pgfpathlineto{\pgfqpoint{6.493112in}{0.740169in}}%
\pgfpathlineto{\pgfqpoint{6.611283in}{0.736845in}}%
\pgfpathlineto{\pgfqpoint{6.773768in}{0.734688in}}%
\pgfpathlineto{\pgfqpoint{6.842701in}{0.734234in}}%
\pgfpathlineto{\pgfqpoint{6.842701in}{0.734234in}}%
\pgfusepath{stroke}%
\end{pgfscope}%
\begin{pgfscope}%
\pgfsetrectcap%
\pgfsetmiterjoin%
\pgfsetlinewidth{0.803000pt}%
\definecolor{currentstroke}{rgb}{0.000000,0.000000,0.000000}%
\pgfsetstrokecolor{currentstroke}%
\pgfsetdash{}{0pt}%
\pgfpathmoveto{\pgfqpoint{0.787430in}{0.733417in}}%
\pgfpathlineto{\pgfqpoint{0.787430in}{3.288791in}}%
\pgfusepath{stroke}%
\end{pgfscope}%
\begin{pgfscope}%
\pgfsetrectcap%
\pgfsetmiterjoin%
\pgfsetlinewidth{0.803000pt}%
\definecolor{currentstroke}{rgb}{0.000000,0.000000,0.000000}%
\pgfsetstrokecolor{currentstroke}%
\pgfsetdash{}{0pt}%
\pgfpathmoveto{\pgfqpoint{7.850000in}{0.733417in}}%
\pgfpathlineto{\pgfqpoint{7.850000in}{3.288791in}}%
\pgfusepath{stroke}%
\end{pgfscope}%
\begin{pgfscope}%
\pgfsetrectcap%
\pgfsetmiterjoin%
\pgfsetlinewidth{0.803000pt}%
\definecolor{currentstroke}{rgb}{0.000000,0.000000,0.000000}%
\pgfsetstrokecolor{currentstroke}%
\pgfsetdash{}{0pt}%
\pgfpathmoveto{\pgfqpoint{0.787430in}{0.733417in}}%
\pgfpathlineto{\pgfqpoint{7.850000in}{0.733417in}}%
\pgfusepath{stroke}%
\end{pgfscope}%
\begin{pgfscope}%
\pgfsetrectcap%
\pgfsetmiterjoin%
\pgfsetlinewidth{0.803000pt}%
\definecolor{currentstroke}{rgb}{0.000000,0.000000,0.000000}%
\pgfsetstrokecolor{currentstroke}%
\pgfsetdash{}{0pt}%
\pgfpathmoveto{\pgfqpoint{0.787430in}{3.288791in}}%
\pgfpathlineto{\pgfqpoint{7.850000in}{3.288791in}}%
\pgfusepath{stroke}%
\end{pgfscope}%
\begin{pgfscope}%
\pgfsetbuttcap%
\pgfsetmiterjoin%
\definecolor{currentfill}{rgb}{1.000000,1.000000,1.000000}%
\pgfsetfillcolor{currentfill}%
\pgfsetfillopacity{0.800000}%
\pgfsetlinewidth{1.003750pt}%
\definecolor{currentstroke}{rgb}{0.800000,0.800000,0.800000}%
\pgfsetstrokecolor{currentstroke}%
\pgfsetstrokeopacity{0.800000}%
\pgfsetdash{}{0pt}%
\pgfpathmoveto{\pgfqpoint{6.955323in}{2.790334in}}%
\pgfpathlineto{\pgfqpoint{7.752778in}{2.790334in}}%
\pgfpathquadraticcurveto{\pgfqpoint{7.780556in}{2.790334in}}{\pgfqpoint{7.780556in}{2.818112in}}%
\pgfpathlineto{\pgfqpoint{7.780556in}{3.191568in}}%
\pgfpathquadraticcurveto{\pgfqpoint{7.780556in}{3.219346in}}{\pgfqpoint{7.752778in}{3.219346in}}%
\pgfpathlineto{\pgfqpoint{6.955323in}{3.219346in}}%
\pgfpathquadraticcurveto{\pgfqpoint{6.927545in}{3.219346in}}{\pgfqpoint{6.927545in}{3.191568in}}%
\pgfpathlineto{\pgfqpoint{6.927545in}{2.818112in}}%
\pgfpathquadraticcurveto{\pgfqpoint{6.927545in}{2.790334in}}{\pgfqpoint{6.955323in}{2.790334in}}%
\pgfpathlineto{\pgfqpoint{6.955323in}{2.790334in}}%
\pgfpathclose%
\pgfusepath{stroke,fill}%
\end{pgfscope}%
\begin{pgfscope}%
\pgfsetbuttcap%
\pgfsetmiterjoin%
\definecolor{currentfill}{rgb}{0.121569,0.466667,0.705882}%
\pgfsetfillcolor{currentfill}%
\pgfsetlinewidth{0.000000pt}%
\definecolor{currentstroke}{rgb}{0.000000,0.000000,0.000000}%
\pgfsetstrokecolor{currentstroke}%
\pgfsetstrokeopacity{0.000000}%
\pgfsetdash{}{0pt}%
\pgfpathmoveto{\pgfqpoint{6.983101in}{3.066568in}}%
\pgfpathlineto{\pgfqpoint{7.260879in}{3.066568in}}%
\pgfpathlineto{\pgfqpoint{7.260879in}{3.163791in}}%
\pgfpathlineto{\pgfqpoint{6.983101in}{3.163791in}}%
\pgfpathlineto{\pgfqpoint{6.983101in}{3.066568in}}%
\pgfpathclose%
\pgfusepath{fill}%
\end{pgfscope}%
\begin{pgfscope}%
\definecolor{textcolor}{rgb}{0.000000,0.000000,0.000000}%
\pgfsetstrokecolor{textcolor}%
\pgfsetfillcolor{textcolor}%
\pgftext[x=7.371990in,y=3.066568in,left,base]{\color{textcolor}{\rmfamily\fontsize{10.000000}{12.000000}\selectfont\catcode`\^=\active\def^{\ifmmode\sp\else\^{}\fi}\catcode`\%=\active\def%{\%}RNG}}%
\end{pgfscope}%
\begin{pgfscope}%
\pgfsetrectcap%
\pgfsetroundjoin%
\pgfsetlinewidth{2.007500pt}%
\definecolor{currentstroke}{rgb}{1.000000,0.000000,0.000000}%
\pgfsetstrokecolor{currentstroke}%
\pgfsetdash{}{0pt}%
\pgfpathmoveto{\pgfqpoint{6.983101in}{2.921507in}}%
\pgfpathlineto{\pgfqpoint{7.121990in}{2.921507in}}%
\pgfpathlineto{\pgfqpoint{7.260879in}{2.921507in}}%
\pgfusepath{stroke}%
\end{pgfscope}%
\begin{pgfscope}%
\definecolor{textcolor}{rgb}{0.000000,0.000000,0.000000}%
\pgfsetstrokecolor{textcolor}%
\pgfsetfillcolor{textcolor}%
\pgftext[x=7.371990in,y=2.872896in,left,base]{\color{textcolor}{\rmfamily\fontsize{10.000000}{12.000000}\selectfont\catcode`\^=\active\def^{\ifmmode\sp\else\^{}\fi}\catcode`\%=\active\def%{\%}Exact}}%
\end{pgfscope}%
\end{pgfpicture}%
\makeatother%
\endgroup%
}
    \caption{Histogram of Gaussian random numbers generated by the Box-Muller algorithm. The red curve represents the probability density function of a Gaussian distribution with the same mean and standard deviation as the generated samples.}
    \label{fig:gaussian}
\end{figure}

\begin{figure}[H]
    \centering
    \resizebox{\columnwidth}{!}{%% Creator: Matplotlib, PGF backend
%%
%% To include the figure in your LaTeX document, write
%%   \input{<filename>.pgf}
%%
%% Make sure the required packages are loaded in your preamble
%%   \usepackage{pgf}
%%
%% Also ensure that all the required font packages are loaded; for instance,
%% the lmodern package is sometimes necessary when using math font.
%%   \usepackage{lmodern}
%%
%% Figures using additional raster images can only be included by \input if
%% they are in the same directory as the main LaTeX file. For loading figures
%% from other directories you can use the `import` package
%%   \usepackage{import}
%%
%% and then include the figures with
%%   \import{<path to file>}{<filename>.pgf}
%%
%% Matplotlib used the following preamble
%%   \def\mathdefault#1{#1}
%%   \everymath=\expandafter{\the\everymath\displaystyle}
%%   \IfFileExists{scrextend.sty}{
%%     \usepackage[fontsize=10.000000pt]{scrextend}
%%   }{
%%     \renewcommand{\normalsize}{\fontsize{10.000000}{12.000000}\selectfont}
%%     \normalsize
%%   }
%%   
%%   \makeatletter\@ifpackageloaded{underscore}{}{\usepackage[strings]{underscore}}\makeatother
%%
\begingroup%
\makeatletter%
\begin{pgfpicture}%
\pgfpathrectangle{\pgfpointorigin}{\pgfqpoint{6.400000in}{4.800000in}}%
\pgfusepath{use as bounding box, clip}%
\begin{pgfscope}%
\pgfsetbuttcap%
\pgfsetmiterjoin%
\definecolor{currentfill}{rgb}{1.000000,1.000000,1.000000}%
\pgfsetfillcolor{currentfill}%
\pgfsetlinewidth{0.000000pt}%
\definecolor{currentstroke}{rgb}{1.000000,1.000000,1.000000}%
\pgfsetstrokecolor{currentstroke}%
\pgfsetdash{}{0pt}%
\pgfpathmoveto{\pgfqpoint{0.000000in}{0.000000in}}%
\pgfpathlineto{\pgfqpoint{6.400000in}{0.000000in}}%
\pgfpathlineto{\pgfqpoint{6.400000in}{4.800000in}}%
\pgfpathlineto{\pgfqpoint{0.000000in}{4.800000in}}%
\pgfpathlineto{\pgfqpoint{0.000000in}{0.000000in}}%
\pgfpathclose%
\pgfusepath{fill}%
\end{pgfscope}%
\begin{pgfscope}%
\pgfsetbuttcap%
\pgfsetmiterjoin%
\definecolor{currentfill}{rgb}{1.000000,1.000000,1.000000}%
\pgfsetfillcolor{currentfill}%
\pgfsetlinewidth{0.000000pt}%
\definecolor{currentstroke}{rgb}{0.000000,0.000000,0.000000}%
\pgfsetstrokecolor{currentstroke}%
\pgfsetstrokeopacity{0.000000}%
\pgfsetdash{}{0pt}%
\pgfpathmoveto{\pgfqpoint{0.611806in}{2.699151in}}%
\pgfpathlineto{\pgfqpoint{6.250000in}{2.699151in}}%
\pgfpathlineto{\pgfqpoint{6.250000in}{4.650000in}}%
\pgfpathlineto{\pgfqpoint{0.611806in}{4.650000in}}%
\pgfpathlineto{\pgfqpoint{0.611806in}{2.699151in}}%
\pgfpathclose%
\pgfusepath{fill}%
\end{pgfscope}%
\begin{pgfscope}%
\pgfsetbuttcap%
\pgfsetroundjoin%
\definecolor{currentfill}{rgb}{0.000000,0.000000,0.000000}%
\pgfsetfillcolor{currentfill}%
\pgfsetlinewidth{0.803000pt}%
\definecolor{currentstroke}{rgb}{0.000000,0.000000,0.000000}%
\pgfsetstrokecolor{currentstroke}%
\pgfsetdash{}{0pt}%
\pgfsys@defobject{currentmarker}{\pgfqpoint{0.000000in}{-0.048611in}}{\pgfqpoint{0.000000in}{0.000000in}}{%
\pgfpathmoveto{\pgfqpoint{0.000000in}{0.000000in}}%
\pgfpathlineto{\pgfqpoint{0.000000in}{-0.048611in}}%
\pgfusepath{stroke,fill}%
}%
\begin{pgfscope}%
\pgfsys@transformshift{0.868088in}{2.699151in}%
\pgfsys@useobject{currentmarker}{}%
\end{pgfscope}%
\end{pgfscope}%
\begin{pgfscope}%
\pgfsetbuttcap%
\pgfsetroundjoin%
\definecolor{currentfill}{rgb}{0.000000,0.000000,0.000000}%
\pgfsetfillcolor{currentfill}%
\pgfsetlinewidth{0.803000pt}%
\definecolor{currentstroke}{rgb}{0.000000,0.000000,0.000000}%
\pgfsetstrokecolor{currentstroke}%
\pgfsetdash{}{0pt}%
\pgfsys@defobject{currentmarker}{\pgfqpoint{0.000000in}{-0.048611in}}{\pgfqpoint{0.000000in}{0.000000in}}{%
\pgfpathmoveto{\pgfqpoint{0.000000in}{0.000000in}}%
\pgfpathlineto{\pgfqpoint{0.000000in}{-0.048611in}}%
\pgfusepath{stroke,fill}%
}%
\begin{pgfscope}%
\pgfsys@transformshift{1.722359in}{2.699151in}%
\pgfsys@useobject{currentmarker}{}%
\end{pgfscope}%
\end{pgfscope}%
\begin{pgfscope}%
\pgfsetbuttcap%
\pgfsetroundjoin%
\definecolor{currentfill}{rgb}{0.000000,0.000000,0.000000}%
\pgfsetfillcolor{currentfill}%
\pgfsetlinewidth{0.803000pt}%
\definecolor{currentstroke}{rgb}{0.000000,0.000000,0.000000}%
\pgfsetstrokecolor{currentstroke}%
\pgfsetdash{}{0pt}%
\pgfsys@defobject{currentmarker}{\pgfqpoint{0.000000in}{-0.048611in}}{\pgfqpoint{0.000000in}{0.000000in}}{%
\pgfpathmoveto{\pgfqpoint{0.000000in}{0.000000in}}%
\pgfpathlineto{\pgfqpoint{0.000000in}{-0.048611in}}%
\pgfusepath{stroke,fill}%
}%
\begin{pgfscope}%
\pgfsys@transformshift{2.576631in}{2.699151in}%
\pgfsys@useobject{currentmarker}{}%
\end{pgfscope}%
\end{pgfscope}%
\begin{pgfscope}%
\pgfsetbuttcap%
\pgfsetroundjoin%
\definecolor{currentfill}{rgb}{0.000000,0.000000,0.000000}%
\pgfsetfillcolor{currentfill}%
\pgfsetlinewidth{0.803000pt}%
\definecolor{currentstroke}{rgb}{0.000000,0.000000,0.000000}%
\pgfsetstrokecolor{currentstroke}%
\pgfsetdash{}{0pt}%
\pgfsys@defobject{currentmarker}{\pgfqpoint{0.000000in}{-0.048611in}}{\pgfqpoint{0.000000in}{0.000000in}}{%
\pgfpathmoveto{\pgfqpoint{0.000000in}{0.000000in}}%
\pgfpathlineto{\pgfqpoint{0.000000in}{-0.048611in}}%
\pgfusepath{stroke,fill}%
}%
\begin{pgfscope}%
\pgfsys@transformshift{3.430903in}{2.699151in}%
\pgfsys@useobject{currentmarker}{}%
\end{pgfscope}%
\end{pgfscope}%
\begin{pgfscope}%
\pgfsetbuttcap%
\pgfsetroundjoin%
\definecolor{currentfill}{rgb}{0.000000,0.000000,0.000000}%
\pgfsetfillcolor{currentfill}%
\pgfsetlinewidth{0.803000pt}%
\definecolor{currentstroke}{rgb}{0.000000,0.000000,0.000000}%
\pgfsetstrokecolor{currentstroke}%
\pgfsetdash{}{0pt}%
\pgfsys@defobject{currentmarker}{\pgfqpoint{0.000000in}{-0.048611in}}{\pgfqpoint{0.000000in}{0.000000in}}{%
\pgfpathmoveto{\pgfqpoint{0.000000in}{0.000000in}}%
\pgfpathlineto{\pgfqpoint{0.000000in}{-0.048611in}}%
\pgfusepath{stroke,fill}%
}%
\begin{pgfscope}%
\pgfsys@transformshift{4.285175in}{2.699151in}%
\pgfsys@useobject{currentmarker}{}%
\end{pgfscope}%
\end{pgfscope}%
\begin{pgfscope}%
\pgfsetbuttcap%
\pgfsetroundjoin%
\definecolor{currentfill}{rgb}{0.000000,0.000000,0.000000}%
\pgfsetfillcolor{currentfill}%
\pgfsetlinewidth{0.803000pt}%
\definecolor{currentstroke}{rgb}{0.000000,0.000000,0.000000}%
\pgfsetstrokecolor{currentstroke}%
\pgfsetdash{}{0pt}%
\pgfsys@defobject{currentmarker}{\pgfqpoint{0.000000in}{-0.048611in}}{\pgfqpoint{0.000000in}{0.000000in}}{%
\pgfpathmoveto{\pgfqpoint{0.000000in}{0.000000in}}%
\pgfpathlineto{\pgfqpoint{0.000000in}{-0.048611in}}%
\pgfusepath{stroke,fill}%
}%
\begin{pgfscope}%
\pgfsys@transformshift{5.139447in}{2.699151in}%
\pgfsys@useobject{currentmarker}{}%
\end{pgfscope}%
\end{pgfscope}%
\begin{pgfscope}%
\pgfsetbuttcap%
\pgfsetroundjoin%
\definecolor{currentfill}{rgb}{0.000000,0.000000,0.000000}%
\pgfsetfillcolor{currentfill}%
\pgfsetlinewidth{0.803000pt}%
\definecolor{currentstroke}{rgb}{0.000000,0.000000,0.000000}%
\pgfsetstrokecolor{currentstroke}%
\pgfsetdash{}{0pt}%
\pgfsys@defobject{currentmarker}{\pgfqpoint{0.000000in}{-0.048611in}}{\pgfqpoint{0.000000in}{0.000000in}}{%
\pgfpathmoveto{\pgfqpoint{0.000000in}{0.000000in}}%
\pgfpathlineto{\pgfqpoint{0.000000in}{-0.048611in}}%
\pgfusepath{stroke,fill}%
}%
\begin{pgfscope}%
\pgfsys@transformshift{5.993718in}{2.699151in}%
\pgfsys@useobject{currentmarker}{}%
\end{pgfscope}%
\end{pgfscope}%
\begin{pgfscope}%
\pgfsetbuttcap%
\pgfsetroundjoin%
\definecolor{currentfill}{rgb}{0.000000,0.000000,0.000000}%
\pgfsetfillcolor{currentfill}%
\pgfsetlinewidth{0.803000pt}%
\definecolor{currentstroke}{rgb}{0.000000,0.000000,0.000000}%
\pgfsetstrokecolor{currentstroke}%
\pgfsetdash{}{0pt}%
\pgfsys@defobject{currentmarker}{\pgfqpoint{-0.048611in}{0.000000in}}{\pgfqpoint{-0.000000in}{0.000000in}}{%
\pgfpathmoveto{\pgfqpoint{-0.000000in}{0.000000in}}%
\pgfpathlineto{\pgfqpoint{-0.048611in}{0.000000in}}%
\pgfusepath{stroke,fill}%
}%
\begin{pgfscope}%
\pgfsys@transformshift{0.611806in}{2.787826in}%
\pgfsys@useobject{currentmarker}{}%
\end{pgfscope}%
\end{pgfscope}%
\begin{pgfscope}%
\definecolor{textcolor}{rgb}{0.000000,0.000000,0.000000}%
\pgfsetstrokecolor{textcolor}%
\pgfsetfillcolor{textcolor}%
\pgftext[x=0.337114in, y=2.739601in, left, base]{\color{textcolor}{\rmfamily\fontsize{10.000000}{12.000000}\selectfont\catcode`\^=\active\def^{\ifmmode\sp\else\^{}\fi}\catcode`\%=\active\def%{\%}$\mathdefault{0.0}$}}%
\end{pgfscope}%
\begin{pgfscope}%
\pgfsetbuttcap%
\pgfsetroundjoin%
\definecolor{currentfill}{rgb}{0.000000,0.000000,0.000000}%
\pgfsetfillcolor{currentfill}%
\pgfsetlinewidth{0.803000pt}%
\definecolor{currentstroke}{rgb}{0.000000,0.000000,0.000000}%
\pgfsetstrokecolor{currentstroke}%
\pgfsetdash{}{0pt}%
\pgfsys@defobject{currentmarker}{\pgfqpoint{-0.048611in}{0.000000in}}{\pgfqpoint{-0.000000in}{0.000000in}}{%
\pgfpathmoveto{\pgfqpoint{-0.000000in}{0.000000in}}%
\pgfpathlineto{\pgfqpoint{-0.048611in}{0.000000in}}%
\pgfusepath{stroke,fill}%
}%
\begin{pgfscope}%
\pgfsys@transformshift{0.611806in}{3.143059in}%
\pgfsys@useobject{currentmarker}{}%
\end{pgfscope}%
\end{pgfscope}%
\begin{pgfscope}%
\definecolor{textcolor}{rgb}{0.000000,0.000000,0.000000}%
\pgfsetstrokecolor{textcolor}%
\pgfsetfillcolor{textcolor}%
\pgftext[x=0.337114in, y=3.094834in, left, base]{\color{textcolor}{\rmfamily\fontsize{10.000000}{12.000000}\selectfont\catcode`\^=\active\def^{\ifmmode\sp\else\^{}\fi}\catcode`\%=\active\def%{\%}$\mathdefault{0.2}$}}%
\end{pgfscope}%
\begin{pgfscope}%
\pgfsetbuttcap%
\pgfsetroundjoin%
\definecolor{currentfill}{rgb}{0.000000,0.000000,0.000000}%
\pgfsetfillcolor{currentfill}%
\pgfsetlinewidth{0.803000pt}%
\definecolor{currentstroke}{rgb}{0.000000,0.000000,0.000000}%
\pgfsetstrokecolor{currentstroke}%
\pgfsetdash{}{0pt}%
\pgfsys@defobject{currentmarker}{\pgfqpoint{-0.048611in}{0.000000in}}{\pgfqpoint{-0.000000in}{0.000000in}}{%
\pgfpathmoveto{\pgfqpoint{-0.000000in}{0.000000in}}%
\pgfpathlineto{\pgfqpoint{-0.048611in}{0.000000in}}%
\pgfusepath{stroke,fill}%
}%
\begin{pgfscope}%
\pgfsys@transformshift{0.611806in}{3.498292in}%
\pgfsys@useobject{currentmarker}{}%
\end{pgfscope}%
\end{pgfscope}%
\begin{pgfscope}%
\definecolor{textcolor}{rgb}{0.000000,0.000000,0.000000}%
\pgfsetstrokecolor{textcolor}%
\pgfsetfillcolor{textcolor}%
\pgftext[x=0.337114in, y=3.450067in, left, base]{\color{textcolor}{\rmfamily\fontsize{10.000000}{12.000000}\selectfont\catcode`\^=\active\def^{\ifmmode\sp\else\^{}\fi}\catcode`\%=\active\def%{\%}$\mathdefault{0.4}$}}%
\end{pgfscope}%
\begin{pgfscope}%
\pgfsetbuttcap%
\pgfsetroundjoin%
\definecolor{currentfill}{rgb}{0.000000,0.000000,0.000000}%
\pgfsetfillcolor{currentfill}%
\pgfsetlinewidth{0.803000pt}%
\definecolor{currentstroke}{rgb}{0.000000,0.000000,0.000000}%
\pgfsetstrokecolor{currentstroke}%
\pgfsetdash{}{0pt}%
\pgfsys@defobject{currentmarker}{\pgfqpoint{-0.048611in}{0.000000in}}{\pgfqpoint{-0.000000in}{0.000000in}}{%
\pgfpathmoveto{\pgfqpoint{-0.000000in}{0.000000in}}%
\pgfpathlineto{\pgfqpoint{-0.048611in}{0.000000in}}%
\pgfusepath{stroke,fill}%
}%
\begin{pgfscope}%
\pgfsys@transformshift{0.611806in}{3.853526in}%
\pgfsys@useobject{currentmarker}{}%
\end{pgfscope}%
\end{pgfscope}%
\begin{pgfscope}%
\definecolor{textcolor}{rgb}{0.000000,0.000000,0.000000}%
\pgfsetstrokecolor{textcolor}%
\pgfsetfillcolor{textcolor}%
\pgftext[x=0.337114in, y=3.805300in, left, base]{\color{textcolor}{\rmfamily\fontsize{10.000000}{12.000000}\selectfont\catcode`\^=\active\def^{\ifmmode\sp\else\^{}\fi}\catcode`\%=\active\def%{\%}$\mathdefault{0.6}$}}%
\end{pgfscope}%
\begin{pgfscope}%
\pgfsetbuttcap%
\pgfsetroundjoin%
\definecolor{currentfill}{rgb}{0.000000,0.000000,0.000000}%
\pgfsetfillcolor{currentfill}%
\pgfsetlinewidth{0.803000pt}%
\definecolor{currentstroke}{rgb}{0.000000,0.000000,0.000000}%
\pgfsetstrokecolor{currentstroke}%
\pgfsetdash{}{0pt}%
\pgfsys@defobject{currentmarker}{\pgfqpoint{-0.048611in}{0.000000in}}{\pgfqpoint{-0.000000in}{0.000000in}}{%
\pgfpathmoveto{\pgfqpoint{-0.000000in}{0.000000in}}%
\pgfpathlineto{\pgfqpoint{-0.048611in}{0.000000in}}%
\pgfusepath{stroke,fill}%
}%
\begin{pgfscope}%
\pgfsys@transformshift{0.611806in}{4.208759in}%
\pgfsys@useobject{currentmarker}{}%
\end{pgfscope}%
\end{pgfscope}%
\begin{pgfscope}%
\definecolor{textcolor}{rgb}{0.000000,0.000000,0.000000}%
\pgfsetstrokecolor{textcolor}%
\pgfsetfillcolor{textcolor}%
\pgftext[x=0.337114in, y=4.160534in, left, base]{\color{textcolor}{\rmfamily\fontsize{10.000000}{12.000000}\selectfont\catcode`\^=\active\def^{\ifmmode\sp\else\^{}\fi}\catcode`\%=\active\def%{\%}$\mathdefault{0.8}$}}%
\end{pgfscope}%
\begin{pgfscope}%
\pgfsetbuttcap%
\pgfsetroundjoin%
\definecolor{currentfill}{rgb}{0.000000,0.000000,0.000000}%
\pgfsetfillcolor{currentfill}%
\pgfsetlinewidth{0.803000pt}%
\definecolor{currentstroke}{rgb}{0.000000,0.000000,0.000000}%
\pgfsetstrokecolor{currentstroke}%
\pgfsetdash{}{0pt}%
\pgfsys@defobject{currentmarker}{\pgfqpoint{-0.048611in}{0.000000in}}{\pgfqpoint{-0.000000in}{0.000000in}}{%
\pgfpathmoveto{\pgfqpoint{-0.000000in}{0.000000in}}%
\pgfpathlineto{\pgfqpoint{-0.048611in}{0.000000in}}%
\pgfusepath{stroke,fill}%
}%
\begin{pgfscope}%
\pgfsys@transformshift{0.611806in}{4.563992in}%
\pgfsys@useobject{currentmarker}{}%
\end{pgfscope}%
\end{pgfscope}%
\begin{pgfscope}%
\definecolor{textcolor}{rgb}{0.000000,0.000000,0.000000}%
\pgfsetstrokecolor{textcolor}%
\pgfsetfillcolor{textcolor}%
\pgftext[x=0.337114in, y=4.515767in, left, base]{\color{textcolor}{\rmfamily\fontsize{10.000000}{12.000000}\selectfont\catcode`\^=\active\def^{\ifmmode\sp\else\^{}\fi}\catcode`\%=\active\def%{\%}$\mathdefault{1.0}$}}%
\end{pgfscope}%
\begin{pgfscope}%
\definecolor{textcolor}{rgb}{0.000000,0.000000,0.000000}%
\pgfsetstrokecolor{textcolor}%
\pgfsetfillcolor{textcolor}%
\pgftext[x=0.281558in,y=3.674576in,,bottom,rotate=90.000000]{\color{textcolor}{\rmfamily\fontsize{10.000000}{12.000000}\selectfont\catcode`\^=\active\def^{\ifmmode\sp\else\^{}\fi}\catcode`\%=\active\def%{\%}$\hat{U}$}}%
\end{pgfscope}%
\begin{pgfscope}%
\pgfpathrectangle{\pgfqpoint{0.611806in}{2.699151in}}{\pgfqpoint{5.638194in}{1.950849in}}%
\pgfusepath{clip}%
\pgfsetrectcap%
\pgfsetroundjoin%
\pgfsetlinewidth{1.505625pt}%
\definecolor{currentstroke}{rgb}{0.121569,0.466667,0.705882}%
\pgfsetstrokecolor{currentstroke}%
\pgfsetdash{}{0pt}%
\pgfpathmoveto{\pgfqpoint{0.868088in}{2.787826in}}%
\pgfpathlineto{\pgfqpoint{1.206718in}{4.547990in}}%
\pgfpathlineto{\pgfqpoint{1.211849in}{4.561325in}}%
\pgfpathlineto{\pgfqpoint{2.576631in}{2.787826in}}%
\pgfpathlineto{\pgfqpoint{2.915261in}{4.547990in}}%
\pgfpathlineto{\pgfqpoint{2.920392in}{4.561325in}}%
\pgfpathlineto{\pgfqpoint{4.285175in}{2.787826in}}%
\pgfpathlineto{\pgfqpoint{4.623805in}{4.547990in}}%
\pgfpathlineto{\pgfqpoint{4.628936in}{4.561325in}}%
\pgfpathlineto{\pgfqpoint{5.993718in}{2.787826in}}%
\pgfpathlineto{\pgfqpoint{5.993718in}{2.787826in}}%
\pgfusepath{stroke}%
\end{pgfscope}%
\begin{pgfscope}%
\pgfsetrectcap%
\pgfsetmiterjoin%
\pgfsetlinewidth{0.803000pt}%
\definecolor{currentstroke}{rgb}{0.000000,0.000000,0.000000}%
\pgfsetstrokecolor{currentstroke}%
\pgfsetdash{}{0pt}%
\pgfpathmoveto{\pgfqpoint{0.611806in}{2.699151in}}%
\pgfpathlineto{\pgfqpoint{0.611806in}{4.650000in}}%
\pgfusepath{stroke}%
\end{pgfscope}%
\begin{pgfscope}%
\pgfsetrectcap%
\pgfsetmiterjoin%
\pgfsetlinewidth{0.803000pt}%
\definecolor{currentstroke}{rgb}{0.000000,0.000000,0.000000}%
\pgfsetstrokecolor{currentstroke}%
\pgfsetdash{}{0pt}%
\pgfpathmoveto{\pgfqpoint{6.250000in}{2.699151in}}%
\pgfpathlineto{\pgfqpoint{6.250000in}{4.650000in}}%
\pgfusepath{stroke}%
\end{pgfscope}%
\begin{pgfscope}%
\pgfsetrectcap%
\pgfsetmiterjoin%
\pgfsetlinewidth{0.803000pt}%
\definecolor{currentstroke}{rgb}{0.000000,0.000000,0.000000}%
\pgfsetstrokecolor{currentstroke}%
\pgfsetdash{}{0pt}%
\pgfpathmoveto{\pgfqpoint{0.611806in}{2.699151in}}%
\pgfpathlineto{\pgfqpoint{6.250000in}{2.699151in}}%
\pgfusepath{stroke}%
\end{pgfscope}%
\begin{pgfscope}%
\pgfsetrectcap%
\pgfsetmiterjoin%
\pgfsetlinewidth{0.803000pt}%
\definecolor{currentstroke}{rgb}{0.000000,0.000000,0.000000}%
\pgfsetstrokecolor{currentstroke}%
\pgfsetdash{}{0pt}%
\pgfpathmoveto{\pgfqpoint{0.611806in}{4.650000in}}%
\pgfpathlineto{\pgfqpoint{6.250000in}{4.650000in}}%
\pgfusepath{stroke}%
\end{pgfscope}%
\begin{pgfscope}%
\pgfsetbuttcap%
\pgfsetmiterjoin%
\definecolor{currentfill}{rgb}{1.000000,1.000000,1.000000}%
\pgfsetfillcolor{currentfill}%
\pgfsetfillopacity{0.800000}%
\pgfsetlinewidth{1.003750pt}%
\definecolor{currentstroke}{rgb}{0.800000,0.800000,0.800000}%
\pgfsetstrokecolor{currentstroke}%
\pgfsetstrokeopacity{0.800000}%
\pgfsetdash{}{0pt}%
\pgfpathmoveto{\pgfqpoint{5.299075in}{4.303164in}}%
\pgfpathlineto{\pgfqpoint{6.152778in}{4.303164in}}%
\pgfpathquadraticcurveto{\pgfqpoint{6.180556in}{4.303164in}}{\pgfqpoint{6.180556in}{4.330942in}}%
\pgfpathlineto{\pgfqpoint{6.180556in}{4.552778in}}%
\pgfpathquadraticcurveto{\pgfqpoint{6.180556in}{4.580556in}}{\pgfqpoint{6.152778in}{4.580556in}}%
\pgfpathlineto{\pgfqpoint{5.299075in}{4.580556in}}%
\pgfpathquadraticcurveto{\pgfqpoint{5.271297in}{4.580556in}}{\pgfqpoint{5.271297in}{4.552778in}}%
\pgfpathlineto{\pgfqpoint{5.271297in}{4.330942in}}%
\pgfpathquadraticcurveto{\pgfqpoint{5.271297in}{4.303164in}}{\pgfqpoint{5.299075in}{4.303164in}}%
\pgfpathlineto{\pgfqpoint{5.299075in}{4.303164in}}%
\pgfpathclose%
\pgfusepath{stroke,fill}%
\end{pgfscope}%
\begin{pgfscope}%
\pgfsetrectcap%
\pgfsetroundjoin%
\pgfsetlinewidth{1.505625pt}%
\definecolor{currentstroke}{rgb}{0.121569,0.466667,0.705882}%
\pgfsetstrokecolor{currentstroke}%
\pgfsetdash{}{0pt}%
\pgfpathmoveto{\pgfqpoint{5.326852in}{4.442053in}}%
\pgfpathlineto{\pgfqpoint{5.465741in}{4.442053in}}%
\pgfpathlineto{\pgfqpoint{5.604630in}{4.442053in}}%
\pgfusepath{stroke}%
\end{pgfscope}%
\begin{pgfscope}%
\definecolor{textcolor}{rgb}{0.000000,0.000000,0.000000}%
\pgfsetstrokecolor{textcolor}%
\pgfsetfillcolor{textcolor}%
\pgftext[x=5.715741in,y=4.393442in,left,base]{\color{textcolor}{\rmfamily\fontsize{10.000000}{12.000000}\selectfont\catcode`\^=\active\def^{\ifmmode\sp\else\^{}\fi}\catcode`\%=\active\def%{\%}$\hat{U}(\hat{x},\hat{t})$}}%
\end{pgfscope}%
\begin{pgfscope}%
\pgfsetbuttcap%
\pgfsetmiterjoin%
\definecolor{currentfill}{rgb}{1.000000,1.000000,1.000000}%
\pgfsetfillcolor{currentfill}%
\pgfsetlinewidth{0.000000pt}%
\definecolor{currentstroke}{rgb}{0.000000,0.000000,0.000000}%
\pgfsetstrokecolor{currentstroke}%
\pgfsetstrokeopacity{0.000000}%
\pgfsetdash{}{0pt}%
\pgfpathmoveto{\pgfqpoint{0.611806in}{0.549691in}}%
\pgfpathlineto{\pgfqpoint{6.250000in}{0.549691in}}%
\pgfpathlineto{\pgfqpoint{6.250000in}{2.500540in}}%
\pgfpathlineto{\pgfqpoint{0.611806in}{2.500540in}}%
\pgfpathlineto{\pgfqpoint{0.611806in}{0.549691in}}%
\pgfpathclose%
\pgfusepath{fill}%
\end{pgfscope}%
\begin{pgfscope}%
\pgfsetbuttcap%
\pgfsetroundjoin%
\definecolor{currentfill}{rgb}{0.000000,0.000000,0.000000}%
\pgfsetfillcolor{currentfill}%
\pgfsetlinewidth{0.803000pt}%
\definecolor{currentstroke}{rgb}{0.000000,0.000000,0.000000}%
\pgfsetstrokecolor{currentstroke}%
\pgfsetdash{}{0pt}%
\pgfsys@defobject{currentmarker}{\pgfqpoint{0.000000in}{-0.048611in}}{\pgfqpoint{0.000000in}{0.000000in}}{%
\pgfpathmoveto{\pgfqpoint{0.000000in}{0.000000in}}%
\pgfpathlineto{\pgfqpoint{0.000000in}{-0.048611in}}%
\pgfusepath{stroke,fill}%
}%
\begin{pgfscope}%
\pgfsys@transformshift{0.868088in}{0.549691in}%
\pgfsys@useobject{currentmarker}{}%
\end{pgfscope}%
\end{pgfscope}%
\begin{pgfscope}%
\definecolor{textcolor}{rgb}{0.000000,0.000000,0.000000}%
\pgfsetstrokecolor{textcolor}%
\pgfsetfillcolor{textcolor}%
\pgftext[x=0.868088in,y=0.452469in,,top]{\color{textcolor}{\rmfamily\fontsize{10.000000}{12.000000}\selectfont\catcode`\^=\active\def^{\ifmmode\sp\else\^{}\fi}\catcode`\%=\active\def%{\%}$\mathdefault{0.0}$}}%
\end{pgfscope}%
\begin{pgfscope}%
\pgfsetbuttcap%
\pgfsetroundjoin%
\definecolor{currentfill}{rgb}{0.000000,0.000000,0.000000}%
\pgfsetfillcolor{currentfill}%
\pgfsetlinewidth{0.803000pt}%
\definecolor{currentstroke}{rgb}{0.000000,0.000000,0.000000}%
\pgfsetstrokecolor{currentstroke}%
\pgfsetdash{}{0pt}%
\pgfsys@defobject{currentmarker}{\pgfqpoint{0.000000in}{-0.048611in}}{\pgfqpoint{0.000000in}{0.000000in}}{%
\pgfpathmoveto{\pgfqpoint{0.000000in}{0.000000in}}%
\pgfpathlineto{\pgfqpoint{0.000000in}{-0.048611in}}%
\pgfusepath{stroke,fill}%
}%
\begin{pgfscope}%
\pgfsys@transformshift{1.722359in}{0.549691in}%
\pgfsys@useobject{currentmarker}{}%
\end{pgfscope}%
\end{pgfscope}%
\begin{pgfscope}%
\definecolor{textcolor}{rgb}{0.000000,0.000000,0.000000}%
\pgfsetstrokecolor{textcolor}%
\pgfsetfillcolor{textcolor}%
\pgftext[x=1.722359in,y=0.452469in,,top]{\color{textcolor}{\rmfamily\fontsize{10.000000}{12.000000}\selectfont\catcode`\^=\active\def^{\ifmmode\sp\else\^{}\fi}\catcode`\%=\active\def%{\%}$\mathdefault{0.5}$}}%
\end{pgfscope}%
\begin{pgfscope}%
\pgfsetbuttcap%
\pgfsetroundjoin%
\definecolor{currentfill}{rgb}{0.000000,0.000000,0.000000}%
\pgfsetfillcolor{currentfill}%
\pgfsetlinewidth{0.803000pt}%
\definecolor{currentstroke}{rgb}{0.000000,0.000000,0.000000}%
\pgfsetstrokecolor{currentstroke}%
\pgfsetdash{}{0pt}%
\pgfsys@defobject{currentmarker}{\pgfqpoint{0.000000in}{-0.048611in}}{\pgfqpoint{0.000000in}{0.000000in}}{%
\pgfpathmoveto{\pgfqpoint{0.000000in}{0.000000in}}%
\pgfpathlineto{\pgfqpoint{0.000000in}{-0.048611in}}%
\pgfusepath{stroke,fill}%
}%
\begin{pgfscope}%
\pgfsys@transformshift{2.576631in}{0.549691in}%
\pgfsys@useobject{currentmarker}{}%
\end{pgfscope}%
\end{pgfscope}%
\begin{pgfscope}%
\definecolor{textcolor}{rgb}{0.000000,0.000000,0.000000}%
\pgfsetstrokecolor{textcolor}%
\pgfsetfillcolor{textcolor}%
\pgftext[x=2.576631in,y=0.452469in,,top]{\color{textcolor}{\rmfamily\fontsize{10.000000}{12.000000}\selectfont\catcode`\^=\active\def^{\ifmmode\sp\else\^{}\fi}\catcode`\%=\active\def%{\%}$\mathdefault{1.0}$}}%
\end{pgfscope}%
\begin{pgfscope}%
\pgfsetbuttcap%
\pgfsetroundjoin%
\definecolor{currentfill}{rgb}{0.000000,0.000000,0.000000}%
\pgfsetfillcolor{currentfill}%
\pgfsetlinewidth{0.803000pt}%
\definecolor{currentstroke}{rgb}{0.000000,0.000000,0.000000}%
\pgfsetstrokecolor{currentstroke}%
\pgfsetdash{}{0pt}%
\pgfsys@defobject{currentmarker}{\pgfqpoint{0.000000in}{-0.048611in}}{\pgfqpoint{0.000000in}{0.000000in}}{%
\pgfpathmoveto{\pgfqpoint{0.000000in}{0.000000in}}%
\pgfpathlineto{\pgfqpoint{0.000000in}{-0.048611in}}%
\pgfusepath{stroke,fill}%
}%
\begin{pgfscope}%
\pgfsys@transformshift{3.430903in}{0.549691in}%
\pgfsys@useobject{currentmarker}{}%
\end{pgfscope}%
\end{pgfscope}%
\begin{pgfscope}%
\definecolor{textcolor}{rgb}{0.000000,0.000000,0.000000}%
\pgfsetstrokecolor{textcolor}%
\pgfsetfillcolor{textcolor}%
\pgftext[x=3.430903in,y=0.452469in,,top]{\color{textcolor}{\rmfamily\fontsize{10.000000}{12.000000}\selectfont\catcode`\^=\active\def^{\ifmmode\sp\else\^{}\fi}\catcode`\%=\active\def%{\%}$\mathdefault{1.5}$}}%
\end{pgfscope}%
\begin{pgfscope}%
\pgfsetbuttcap%
\pgfsetroundjoin%
\definecolor{currentfill}{rgb}{0.000000,0.000000,0.000000}%
\pgfsetfillcolor{currentfill}%
\pgfsetlinewidth{0.803000pt}%
\definecolor{currentstroke}{rgb}{0.000000,0.000000,0.000000}%
\pgfsetstrokecolor{currentstroke}%
\pgfsetdash{}{0pt}%
\pgfsys@defobject{currentmarker}{\pgfqpoint{0.000000in}{-0.048611in}}{\pgfqpoint{0.000000in}{0.000000in}}{%
\pgfpathmoveto{\pgfqpoint{0.000000in}{0.000000in}}%
\pgfpathlineto{\pgfqpoint{0.000000in}{-0.048611in}}%
\pgfusepath{stroke,fill}%
}%
\begin{pgfscope}%
\pgfsys@transformshift{4.285175in}{0.549691in}%
\pgfsys@useobject{currentmarker}{}%
\end{pgfscope}%
\end{pgfscope}%
\begin{pgfscope}%
\definecolor{textcolor}{rgb}{0.000000,0.000000,0.000000}%
\pgfsetstrokecolor{textcolor}%
\pgfsetfillcolor{textcolor}%
\pgftext[x=4.285175in,y=0.452469in,,top]{\color{textcolor}{\rmfamily\fontsize{10.000000}{12.000000}\selectfont\catcode`\^=\active\def^{\ifmmode\sp\else\^{}\fi}\catcode`\%=\active\def%{\%}$\mathdefault{2.0}$}}%
\end{pgfscope}%
\begin{pgfscope}%
\pgfsetbuttcap%
\pgfsetroundjoin%
\definecolor{currentfill}{rgb}{0.000000,0.000000,0.000000}%
\pgfsetfillcolor{currentfill}%
\pgfsetlinewidth{0.803000pt}%
\definecolor{currentstroke}{rgb}{0.000000,0.000000,0.000000}%
\pgfsetstrokecolor{currentstroke}%
\pgfsetdash{}{0pt}%
\pgfsys@defobject{currentmarker}{\pgfqpoint{0.000000in}{-0.048611in}}{\pgfqpoint{0.000000in}{0.000000in}}{%
\pgfpathmoveto{\pgfqpoint{0.000000in}{0.000000in}}%
\pgfpathlineto{\pgfqpoint{0.000000in}{-0.048611in}}%
\pgfusepath{stroke,fill}%
}%
\begin{pgfscope}%
\pgfsys@transformshift{5.139447in}{0.549691in}%
\pgfsys@useobject{currentmarker}{}%
\end{pgfscope}%
\end{pgfscope}%
\begin{pgfscope}%
\definecolor{textcolor}{rgb}{0.000000,0.000000,0.000000}%
\pgfsetstrokecolor{textcolor}%
\pgfsetfillcolor{textcolor}%
\pgftext[x=5.139447in,y=0.452469in,,top]{\color{textcolor}{\rmfamily\fontsize{10.000000}{12.000000}\selectfont\catcode`\^=\active\def^{\ifmmode\sp\else\^{}\fi}\catcode`\%=\active\def%{\%}$\mathdefault{2.5}$}}%
\end{pgfscope}%
\begin{pgfscope}%
\pgfsetbuttcap%
\pgfsetroundjoin%
\definecolor{currentfill}{rgb}{0.000000,0.000000,0.000000}%
\pgfsetfillcolor{currentfill}%
\pgfsetlinewidth{0.803000pt}%
\definecolor{currentstroke}{rgb}{0.000000,0.000000,0.000000}%
\pgfsetstrokecolor{currentstroke}%
\pgfsetdash{}{0pt}%
\pgfsys@defobject{currentmarker}{\pgfqpoint{0.000000in}{-0.048611in}}{\pgfqpoint{0.000000in}{0.000000in}}{%
\pgfpathmoveto{\pgfqpoint{0.000000in}{0.000000in}}%
\pgfpathlineto{\pgfqpoint{0.000000in}{-0.048611in}}%
\pgfusepath{stroke,fill}%
}%
\begin{pgfscope}%
\pgfsys@transformshift{5.993718in}{0.549691in}%
\pgfsys@useobject{currentmarker}{}%
\end{pgfscope}%
\end{pgfscope}%
\begin{pgfscope}%
\definecolor{textcolor}{rgb}{0.000000,0.000000,0.000000}%
\pgfsetstrokecolor{textcolor}%
\pgfsetfillcolor{textcolor}%
\pgftext[x=5.993718in,y=0.452469in,,top]{\color{textcolor}{\rmfamily\fontsize{10.000000}{12.000000}\selectfont\catcode`\^=\active\def^{\ifmmode\sp\else\^{}\fi}\catcode`\%=\active\def%{\%}$\mathdefault{3.0}$}}%
\end{pgfscope}%
\begin{pgfscope}%
\definecolor{textcolor}{rgb}{0.000000,0.000000,0.000000}%
\pgfsetstrokecolor{textcolor}%
\pgfsetfillcolor{textcolor}%
\pgftext[x=3.430903in,y=0.273457in,,top]{\color{textcolor}{\rmfamily\fontsize{10.000000}{12.000000}\selectfont\catcode`\^=\active\def^{\ifmmode\sp\else\^{}\fi}\catcode`\%=\active\def%{\%}$\hat{x}$}}%
\end{pgfscope}%
\begin{pgfscope}%
\pgfsetbuttcap%
\pgfsetroundjoin%
\definecolor{currentfill}{rgb}{0.000000,0.000000,0.000000}%
\pgfsetfillcolor{currentfill}%
\pgfsetlinewidth{0.803000pt}%
\definecolor{currentstroke}{rgb}{0.000000,0.000000,0.000000}%
\pgfsetstrokecolor{currentstroke}%
\pgfsetdash{}{0pt}%
\pgfsys@defobject{currentmarker}{\pgfqpoint{-0.048611in}{0.000000in}}{\pgfqpoint{-0.000000in}{0.000000in}}{%
\pgfpathmoveto{\pgfqpoint{-0.000000in}{0.000000in}}%
\pgfpathlineto{\pgfqpoint{-0.048611in}{0.000000in}}%
\pgfusepath{stroke,fill}%
}%
\begin{pgfscope}%
\pgfsys@transformshift{0.611806in}{0.638366in}%
\pgfsys@useobject{currentmarker}{}%
\end{pgfscope}%
\end{pgfscope}%
\begin{pgfscope}%
\definecolor{textcolor}{rgb}{0.000000,0.000000,0.000000}%
\pgfsetstrokecolor{textcolor}%
\pgfsetfillcolor{textcolor}%
\pgftext[x=0.337114in, y=0.590141in, left, base]{\color{textcolor}{\rmfamily\fontsize{10.000000}{12.000000}\selectfont\catcode`\^=\active\def^{\ifmmode\sp\else\^{}\fi}\catcode`\%=\active\def%{\%}$\mathdefault{\ensuremath{-}5}$}}%
\end{pgfscope}%
\begin{pgfscope}%
\pgfsetbuttcap%
\pgfsetroundjoin%
\definecolor{currentfill}{rgb}{0.000000,0.000000,0.000000}%
\pgfsetfillcolor{currentfill}%
\pgfsetlinewidth{0.803000pt}%
\definecolor{currentstroke}{rgb}{0.000000,0.000000,0.000000}%
\pgfsetstrokecolor{currentstroke}%
\pgfsetdash{}{0pt}%
\pgfsys@defobject{currentmarker}{\pgfqpoint{-0.048611in}{0.000000in}}{\pgfqpoint{-0.000000in}{0.000000in}}{%
\pgfpathmoveto{\pgfqpoint{-0.000000in}{0.000000in}}%
\pgfpathlineto{\pgfqpoint{-0.048611in}{0.000000in}}%
\pgfusepath{stroke,fill}%
}%
\begin{pgfscope}%
\pgfsys@transformshift{0.611806in}{0.922126in}%
\pgfsys@useobject{currentmarker}{}%
\end{pgfscope}%
\end{pgfscope}%
\begin{pgfscope}%
\definecolor{textcolor}{rgb}{0.000000,0.000000,0.000000}%
\pgfsetstrokecolor{textcolor}%
\pgfsetfillcolor{textcolor}%
\pgftext[x=0.337114in, y=0.873901in, left, base]{\color{textcolor}{\rmfamily\fontsize{10.000000}{12.000000}\selectfont\catcode`\^=\active\def^{\ifmmode\sp\else\^{}\fi}\catcode`\%=\active\def%{\%}$\mathdefault{\ensuremath{-}4}$}}%
\end{pgfscope}%
\begin{pgfscope}%
\pgfsetbuttcap%
\pgfsetroundjoin%
\definecolor{currentfill}{rgb}{0.000000,0.000000,0.000000}%
\pgfsetfillcolor{currentfill}%
\pgfsetlinewidth{0.803000pt}%
\definecolor{currentstroke}{rgb}{0.000000,0.000000,0.000000}%
\pgfsetstrokecolor{currentstroke}%
\pgfsetdash{}{0pt}%
\pgfsys@defobject{currentmarker}{\pgfqpoint{-0.048611in}{0.000000in}}{\pgfqpoint{-0.000000in}{0.000000in}}{%
\pgfpathmoveto{\pgfqpoint{-0.000000in}{0.000000in}}%
\pgfpathlineto{\pgfqpoint{-0.048611in}{0.000000in}}%
\pgfusepath{stroke,fill}%
}%
\begin{pgfscope}%
\pgfsys@transformshift{0.611806in}{1.205886in}%
\pgfsys@useobject{currentmarker}{}%
\end{pgfscope}%
\end{pgfscope}%
\begin{pgfscope}%
\definecolor{textcolor}{rgb}{0.000000,0.000000,0.000000}%
\pgfsetstrokecolor{textcolor}%
\pgfsetfillcolor{textcolor}%
\pgftext[x=0.337114in, y=1.157660in, left, base]{\color{textcolor}{\rmfamily\fontsize{10.000000}{12.000000}\selectfont\catcode`\^=\active\def^{\ifmmode\sp\else\^{}\fi}\catcode`\%=\active\def%{\%}$\mathdefault{\ensuremath{-}3}$}}%
\end{pgfscope}%
\begin{pgfscope}%
\pgfsetbuttcap%
\pgfsetroundjoin%
\definecolor{currentfill}{rgb}{0.000000,0.000000,0.000000}%
\pgfsetfillcolor{currentfill}%
\pgfsetlinewidth{0.803000pt}%
\definecolor{currentstroke}{rgb}{0.000000,0.000000,0.000000}%
\pgfsetstrokecolor{currentstroke}%
\pgfsetdash{}{0pt}%
\pgfsys@defobject{currentmarker}{\pgfqpoint{-0.048611in}{0.000000in}}{\pgfqpoint{-0.000000in}{0.000000in}}{%
\pgfpathmoveto{\pgfqpoint{-0.000000in}{0.000000in}}%
\pgfpathlineto{\pgfqpoint{-0.048611in}{0.000000in}}%
\pgfusepath{stroke,fill}%
}%
\begin{pgfscope}%
\pgfsys@transformshift{0.611806in}{1.489646in}%
\pgfsys@useobject{currentmarker}{}%
\end{pgfscope}%
\end{pgfscope}%
\begin{pgfscope}%
\definecolor{textcolor}{rgb}{0.000000,0.000000,0.000000}%
\pgfsetstrokecolor{textcolor}%
\pgfsetfillcolor{textcolor}%
\pgftext[x=0.337114in, y=1.441420in, left, base]{\color{textcolor}{\rmfamily\fontsize{10.000000}{12.000000}\selectfont\catcode`\^=\active\def^{\ifmmode\sp\else\^{}\fi}\catcode`\%=\active\def%{\%}$\mathdefault{\ensuremath{-}2}$}}%
\end{pgfscope}%
\begin{pgfscope}%
\pgfsetbuttcap%
\pgfsetroundjoin%
\definecolor{currentfill}{rgb}{0.000000,0.000000,0.000000}%
\pgfsetfillcolor{currentfill}%
\pgfsetlinewidth{0.803000pt}%
\definecolor{currentstroke}{rgb}{0.000000,0.000000,0.000000}%
\pgfsetstrokecolor{currentstroke}%
\pgfsetdash{}{0pt}%
\pgfsys@defobject{currentmarker}{\pgfqpoint{-0.048611in}{0.000000in}}{\pgfqpoint{-0.000000in}{0.000000in}}{%
\pgfpathmoveto{\pgfqpoint{-0.000000in}{0.000000in}}%
\pgfpathlineto{\pgfqpoint{-0.048611in}{0.000000in}}%
\pgfusepath{stroke,fill}%
}%
\begin{pgfscope}%
\pgfsys@transformshift{0.611806in}{1.773405in}%
\pgfsys@useobject{currentmarker}{}%
\end{pgfscope}%
\end{pgfscope}%
\begin{pgfscope}%
\definecolor{textcolor}{rgb}{0.000000,0.000000,0.000000}%
\pgfsetstrokecolor{textcolor}%
\pgfsetfillcolor{textcolor}%
\pgftext[x=0.337114in, y=1.725180in, left, base]{\color{textcolor}{\rmfamily\fontsize{10.000000}{12.000000}\selectfont\catcode`\^=\active\def^{\ifmmode\sp\else\^{}\fi}\catcode`\%=\active\def%{\%}$\mathdefault{\ensuremath{-}1}$}}%
\end{pgfscope}%
\begin{pgfscope}%
\pgfsetbuttcap%
\pgfsetroundjoin%
\definecolor{currentfill}{rgb}{0.000000,0.000000,0.000000}%
\pgfsetfillcolor{currentfill}%
\pgfsetlinewidth{0.803000pt}%
\definecolor{currentstroke}{rgb}{0.000000,0.000000,0.000000}%
\pgfsetstrokecolor{currentstroke}%
\pgfsetdash{}{0pt}%
\pgfsys@defobject{currentmarker}{\pgfqpoint{-0.048611in}{0.000000in}}{\pgfqpoint{-0.000000in}{0.000000in}}{%
\pgfpathmoveto{\pgfqpoint{-0.000000in}{0.000000in}}%
\pgfpathlineto{\pgfqpoint{-0.048611in}{0.000000in}}%
\pgfusepath{stroke,fill}%
}%
\begin{pgfscope}%
\pgfsys@transformshift{0.611806in}{2.057165in}%
\pgfsys@useobject{currentmarker}{}%
\end{pgfscope}%
\end{pgfscope}%
\begin{pgfscope}%
\definecolor{textcolor}{rgb}{0.000000,0.000000,0.000000}%
\pgfsetstrokecolor{textcolor}%
\pgfsetfillcolor{textcolor}%
\pgftext[x=0.445139in, y=2.008940in, left, base]{\color{textcolor}{\rmfamily\fontsize{10.000000}{12.000000}\selectfont\catcode`\^=\active\def^{\ifmmode\sp\else\^{}\fi}\catcode`\%=\active\def%{\%}$\mathdefault{0}$}}%
\end{pgfscope}%
\begin{pgfscope}%
\pgfsetbuttcap%
\pgfsetroundjoin%
\definecolor{currentfill}{rgb}{0.000000,0.000000,0.000000}%
\pgfsetfillcolor{currentfill}%
\pgfsetlinewidth{0.803000pt}%
\definecolor{currentstroke}{rgb}{0.000000,0.000000,0.000000}%
\pgfsetstrokecolor{currentstroke}%
\pgfsetdash{}{0pt}%
\pgfsys@defobject{currentmarker}{\pgfqpoint{-0.048611in}{0.000000in}}{\pgfqpoint{-0.000000in}{0.000000in}}{%
\pgfpathmoveto{\pgfqpoint{-0.000000in}{0.000000in}}%
\pgfpathlineto{\pgfqpoint{-0.048611in}{0.000000in}}%
\pgfusepath{stroke,fill}%
}%
\begin{pgfscope}%
\pgfsys@transformshift{0.611806in}{2.340925in}%
\pgfsys@useobject{currentmarker}{}%
\end{pgfscope}%
\end{pgfscope}%
\begin{pgfscope}%
\definecolor{textcolor}{rgb}{0.000000,0.000000,0.000000}%
\pgfsetstrokecolor{textcolor}%
\pgfsetfillcolor{textcolor}%
\pgftext[x=0.445139in, y=2.292700in, left, base]{\color{textcolor}{\rmfamily\fontsize{10.000000}{12.000000}\selectfont\catcode`\^=\active\def^{\ifmmode\sp\else\^{}\fi}\catcode`\%=\active\def%{\%}$\mathdefault{1}$}}%
\end{pgfscope}%
\begin{pgfscope}%
\definecolor{textcolor}{rgb}{0.000000,0.000000,0.000000}%
\pgfsetstrokecolor{textcolor}%
\pgfsetfillcolor{textcolor}%
\pgftext[x=0.281558in,y=1.525116in,,bottom,rotate=90.000000]{\color{textcolor}{\rmfamily\fontsize{10.000000}{12.000000}\selectfont\catcode`\^=\active\def^{\ifmmode\sp\else\^{}\fi}\catcode`\%=\active\def%{\%}$\partial\hat{U}\partial \hat{x}$}}%
\end{pgfscope}%
\begin{pgfscope}%
\pgfpathrectangle{\pgfqpoint{0.611806in}{0.549691in}}{\pgfqpoint{5.638194in}{1.950849in}}%
\pgfusepath{clip}%
\pgfsetrectcap%
\pgfsetroundjoin%
\pgfsetlinewidth{1.505625pt}%
\definecolor{currentstroke}{rgb}{1.000000,0.647059,0.000000}%
\pgfsetstrokecolor{currentstroke}%
\pgfsetdash{}{0pt}%
\pgfpathmoveto{\pgfqpoint{0.868088in}{0.638366in}}%
\pgfpathlineto{\pgfqpoint{1.206718in}{0.638366in}}%
\pgfpathlineto{\pgfqpoint{1.211849in}{2.411865in}}%
\pgfpathlineto{\pgfqpoint{2.571500in}{2.411865in}}%
\pgfpathlineto{\pgfqpoint{2.576631in}{0.638366in}}%
\pgfpathlineto{\pgfqpoint{2.915261in}{0.638366in}}%
\pgfpathlineto{\pgfqpoint{2.920392in}{2.411865in}}%
\pgfpathlineto{\pgfqpoint{4.280044in}{2.411865in}}%
\pgfpathlineto{\pgfqpoint{4.285175in}{0.638366in}}%
\pgfpathlineto{\pgfqpoint{4.623805in}{0.638366in}}%
\pgfpathlineto{\pgfqpoint{4.628936in}{2.411865in}}%
\pgfpathlineto{\pgfqpoint{5.988588in}{2.411865in}}%
\pgfpathlineto{\pgfqpoint{5.993718in}{0.638366in}}%
\pgfpathlineto{\pgfqpoint{5.993718in}{0.638366in}}%
\pgfusepath{stroke}%
\end{pgfscope}%
\begin{pgfscope}%
\pgfsetrectcap%
\pgfsetmiterjoin%
\pgfsetlinewidth{0.803000pt}%
\definecolor{currentstroke}{rgb}{0.000000,0.000000,0.000000}%
\pgfsetstrokecolor{currentstroke}%
\pgfsetdash{}{0pt}%
\pgfpathmoveto{\pgfqpoint{0.611806in}{0.549691in}}%
\pgfpathlineto{\pgfqpoint{0.611806in}{2.500540in}}%
\pgfusepath{stroke}%
\end{pgfscope}%
\begin{pgfscope}%
\pgfsetrectcap%
\pgfsetmiterjoin%
\pgfsetlinewidth{0.803000pt}%
\definecolor{currentstroke}{rgb}{0.000000,0.000000,0.000000}%
\pgfsetstrokecolor{currentstroke}%
\pgfsetdash{}{0pt}%
\pgfpathmoveto{\pgfqpoint{6.250000in}{0.549691in}}%
\pgfpathlineto{\pgfqpoint{6.250000in}{2.500540in}}%
\pgfusepath{stroke}%
\end{pgfscope}%
\begin{pgfscope}%
\pgfsetrectcap%
\pgfsetmiterjoin%
\pgfsetlinewidth{0.803000pt}%
\definecolor{currentstroke}{rgb}{0.000000,0.000000,0.000000}%
\pgfsetstrokecolor{currentstroke}%
\pgfsetdash{}{0pt}%
\pgfpathmoveto{\pgfqpoint{0.611806in}{0.549691in}}%
\pgfpathlineto{\pgfqpoint{6.250000in}{0.549691in}}%
\pgfusepath{stroke}%
\end{pgfscope}%
\begin{pgfscope}%
\pgfsetrectcap%
\pgfsetmiterjoin%
\pgfsetlinewidth{0.803000pt}%
\definecolor{currentstroke}{rgb}{0.000000,0.000000,0.000000}%
\pgfsetstrokecolor{currentstroke}%
\pgfsetdash{}{0pt}%
\pgfpathmoveto{\pgfqpoint{0.611806in}{2.500540in}}%
\pgfpathlineto{\pgfqpoint{6.250000in}{2.500540in}}%
\pgfusepath{stroke}%
\end{pgfscope}%
\begin{pgfscope}%
\pgfsetbuttcap%
\pgfsetmiterjoin%
\definecolor{currentfill}{rgb}{1.000000,1.000000,1.000000}%
\pgfsetfillcolor{currentfill}%
\pgfsetfillopacity{0.800000}%
\pgfsetlinewidth{1.003750pt}%
\definecolor{currentstroke}{rgb}{0.800000,0.800000,0.800000}%
\pgfsetstrokecolor{currentstroke}%
\pgfsetstrokeopacity{0.800000}%
\pgfsetdash{}{0pt}%
\pgfpathmoveto{\pgfqpoint{5.356077in}{2.188426in}}%
\pgfpathlineto{\pgfqpoint{6.152778in}{2.188426in}}%
\pgfpathquadraticcurveto{\pgfqpoint{6.180556in}{2.188426in}}{\pgfqpoint{6.180556in}{2.216204in}}%
\pgfpathlineto{\pgfqpoint{6.180556in}{2.403318in}}%
\pgfpathquadraticcurveto{\pgfqpoint{6.180556in}{2.431096in}}{\pgfqpoint{6.152778in}{2.431096in}}%
\pgfpathlineto{\pgfqpoint{5.356077in}{2.431096in}}%
\pgfpathquadraticcurveto{\pgfqpoint{5.328299in}{2.431096in}}{\pgfqpoint{5.328299in}{2.403318in}}%
\pgfpathlineto{\pgfqpoint{5.328299in}{2.216204in}}%
\pgfpathquadraticcurveto{\pgfqpoint{5.328299in}{2.188426in}}{\pgfqpoint{5.356077in}{2.188426in}}%
\pgfpathlineto{\pgfqpoint{5.356077in}{2.188426in}}%
\pgfpathclose%
\pgfusepath{stroke,fill}%
\end{pgfscope}%
\begin{pgfscope}%
\pgfsetrectcap%
\pgfsetroundjoin%
\pgfsetlinewidth{1.505625pt}%
\definecolor{currentstroke}{rgb}{1.000000,0.647059,0.000000}%
\pgfsetstrokecolor{currentstroke}%
\pgfsetdash{}{0pt}%
\pgfpathmoveto{\pgfqpoint{5.383855in}{2.319599in}}%
\pgfpathlineto{\pgfqpoint{5.522744in}{2.319599in}}%
\pgfpathlineto{\pgfqpoint{5.661633in}{2.319599in}}%
\pgfusepath{stroke}%
\end{pgfscope}%
\begin{pgfscope}%
\definecolor{textcolor}{rgb}{0.000000,0.000000,0.000000}%
\pgfsetstrokecolor{textcolor}%
\pgfsetfillcolor{textcolor}%
\pgftext[x=5.772744in,y=2.270988in,left,base]{\color{textcolor}{\rmfamily\fontsize{10.000000}{12.000000}\selectfont\catcode`\^=\active\def^{\ifmmode\sp\else\^{}\fi}\catcode`\%=\active\def%{\%}$\partial\hat{U}\partial \hat{x}$}}%
\end{pgfscope}%
\end{pgfpicture}%
\makeatother%
\endgroup%
}
    \caption{Plotted shape of potential and force}
    \label{fig:potentia_force}
\end{figure}


\begin{figure}[H]
    \centering
    \resizebox{\columnwidth}{!}{%% Creator: Matplotlib, PGF backend
%%
%% To include the figure in your LaTeX document, write
%%   \input{<filename>.pgf}
%%
%% Make sure the required packages are loaded in your preamble
%%   \usepackage{pgf}
%%
%% Also ensure that all the required font packages are loaded; for instance,
%% the lmodern package is sometimes necessary when using math font.
%%   \usepackage{lmodern}
%%
%% Figures using additional raster images can only be included by \input if
%% they are in the same directory as the main LaTeX file. For loading figures
%% from other directories you can use the `import` package
%%   \usepackage{import}
%%
%% and then include the figures with
%%   \import{<path to file>}{<filename>.pgf}
%%
%% Matplotlib used the following preamble
%%   \def\mathdefault#1{#1}
%%   \everymath=\expandafter{\the\everymath\displaystyle}
%%   \IfFileExists{scrextend.sty}{
%%     \usepackage[fontsize=10.000000pt]{scrextend}
%%   }{
%%     \renewcommand{\normalsize}{\fontsize{10.000000}{12.000000}\selectfont}
%%     \normalsize
%%   }
%%   
%%   \makeatletter\@ifpackageloaded{underscore}{}{\usepackage[strings]{underscore}}\makeatother
%%
\begingroup%
\makeatletter%
\begin{pgfpicture}%
\pgfpathrectangle{\pgfpointorigin}{\pgfqpoint{6.400000in}{4.800000in}}%
\pgfusepath{use as bounding box, clip}%
\begin{pgfscope}%
\pgfsetbuttcap%
\pgfsetmiterjoin%
\definecolor{currentfill}{rgb}{1.000000,1.000000,1.000000}%
\pgfsetfillcolor{currentfill}%
\pgfsetlinewidth{0.000000pt}%
\definecolor{currentstroke}{rgb}{1.000000,1.000000,1.000000}%
\pgfsetstrokecolor{currentstroke}%
\pgfsetdash{}{0pt}%
\pgfpathmoveto{\pgfqpoint{0.000000in}{0.000000in}}%
\pgfpathlineto{\pgfqpoint{6.400000in}{0.000000in}}%
\pgfpathlineto{\pgfqpoint{6.400000in}{4.800000in}}%
\pgfpathlineto{\pgfqpoint{0.000000in}{4.800000in}}%
\pgfpathlineto{\pgfqpoint{0.000000in}{0.000000in}}%
\pgfpathclose%
\pgfusepath{fill}%
\end{pgfscope}%
\begin{pgfscope}%
\pgfsetbuttcap%
\pgfsetmiterjoin%
\definecolor{currentfill}{rgb}{1.000000,1.000000,1.000000}%
\pgfsetfillcolor{currentfill}%
\pgfsetlinewidth{0.000000pt}%
\definecolor{currentstroke}{rgb}{0.000000,0.000000,0.000000}%
\pgfsetstrokecolor{currentstroke}%
\pgfsetstrokeopacity{0.000000}%
\pgfsetdash{}{0pt}%
\pgfpathmoveto{\pgfqpoint{0.894648in}{0.771190in}}%
\pgfpathlineto{\pgfqpoint{6.250000in}{0.771190in}}%
\pgfpathlineto{\pgfqpoint{6.250000in}{4.601775in}}%
\pgfpathlineto{\pgfqpoint{0.894648in}{4.601775in}}%
\pgfpathlineto{\pgfqpoint{0.894648in}{0.771190in}}%
\pgfpathclose%
\pgfusepath{fill}%
\end{pgfscope}%
\begin{pgfscope}%
\pgfsetbuttcap%
\pgfsetroundjoin%
\definecolor{currentfill}{rgb}{0.000000,0.000000,0.000000}%
\pgfsetfillcolor{currentfill}%
\pgfsetlinewidth{0.803000pt}%
\definecolor{currentstroke}{rgb}{0.000000,0.000000,0.000000}%
\pgfsetstrokecolor{currentstroke}%
\pgfsetdash{}{0pt}%
\pgfsys@defobject{currentmarker}{\pgfqpoint{0.000000in}{-0.048611in}}{\pgfqpoint{0.000000in}{0.000000in}}{%
\pgfpathmoveto{\pgfqpoint{0.000000in}{0.000000in}}%
\pgfpathlineto{\pgfqpoint{0.000000in}{-0.048611in}}%
\pgfusepath{stroke,fill}%
}%
\begin{pgfscope}%
\pgfsys@transformshift{1.138073in}{0.771190in}%
\pgfsys@useobject{currentmarker}{}%
\end{pgfscope}%
\end{pgfscope}%
\begin{pgfscope}%
\definecolor{textcolor}{rgb}{0.000000,0.000000,0.000000}%
\pgfsetstrokecolor{textcolor}%
\pgfsetfillcolor{textcolor}%
\pgftext[x=1.138073in,y=0.673968in,,top]{\color{textcolor}{\rmfamily\fontsize{10.000000}{12.000000}\selectfont\catcode`\^=\active\def^{\ifmmode\sp\else\^{}\fi}\catcode`\%=\active\def%{\%}$\mathdefault{0}$}}%
\end{pgfscope}%
\begin{pgfscope}%
\pgfsetbuttcap%
\pgfsetroundjoin%
\definecolor{currentfill}{rgb}{0.000000,0.000000,0.000000}%
\pgfsetfillcolor{currentfill}%
\pgfsetlinewidth{0.803000pt}%
\definecolor{currentstroke}{rgb}{0.000000,0.000000,0.000000}%
\pgfsetstrokecolor{currentstroke}%
\pgfsetdash{}{0pt}%
\pgfsys@defobject{currentmarker}{\pgfqpoint{0.000000in}{-0.048611in}}{\pgfqpoint{0.000000in}{0.000000in}}{%
\pgfpathmoveto{\pgfqpoint{0.000000in}{0.000000in}}%
\pgfpathlineto{\pgfqpoint{0.000000in}{-0.048611in}}%
\pgfusepath{stroke,fill}%
}%
\begin{pgfscope}%
\pgfsys@transformshift{2.111773in}{0.771190in}%
\pgfsys@useobject{currentmarker}{}%
\end{pgfscope}%
\end{pgfscope}%
\begin{pgfscope}%
\definecolor{textcolor}{rgb}{0.000000,0.000000,0.000000}%
\pgfsetstrokecolor{textcolor}%
\pgfsetfillcolor{textcolor}%
\pgftext[x=2.111773in,y=0.673968in,,top]{\color{textcolor}{\rmfamily\fontsize{10.000000}{12.000000}\selectfont\catcode`\^=\active\def^{\ifmmode\sp\else\^{}\fi}\catcode`\%=\active\def%{\%}$\mathdefault{20}$}}%
\end{pgfscope}%
\begin{pgfscope}%
\pgfsetbuttcap%
\pgfsetroundjoin%
\definecolor{currentfill}{rgb}{0.000000,0.000000,0.000000}%
\pgfsetfillcolor{currentfill}%
\pgfsetlinewidth{0.803000pt}%
\definecolor{currentstroke}{rgb}{0.000000,0.000000,0.000000}%
\pgfsetstrokecolor{currentstroke}%
\pgfsetdash{}{0pt}%
\pgfsys@defobject{currentmarker}{\pgfqpoint{0.000000in}{-0.048611in}}{\pgfqpoint{0.000000in}{0.000000in}}{%
\pgfpathmoveto{\pgfqpoint{0.000000in}{0.000000in}}%
\pgfpathlineto{\pgfqpoint{0.000000in}{-0.048611in}}%
\pgfusepath{stroke,fill}%
}%
\begin{pgfscope}%
\pgfsys@transformshift{3.085474in}{0.771190in}%
\pgfsys@useobject{currentmarker}{}%
\end{pgfscope}%
\end{pgfscope}%
\begin{pgfscope}%
\definecolor{textcolor}{rgb}{0.000000,0.000000,0.000000}%
\pgfsetstrokecolor{textcolor}%
\pgfsetfillcolor{textcolor}%
\pgftext[x=3.085474in,y=0.673968in,,top]{\color{textcolor}{\rmfamily\fontsize{10.000000}{12.000000}\selectfont\catcode`\^=\active\def^{\ifmmode\sp\else\^{}\fi}\catcode`\%=\active\def%{\%}$\mathdefault{40}$}}%
\end{pgfscope}%
\begin{pgfscope}%
\pgfsetbuttcap%
\pgfsetroundjoin%
\definecolor{currentfill}{rgb}{0.000000,0.000000,0.000000}%
\pgfsetfillcolor{currentfill}%
\pgfsetlinewidth{0.803000pt}%
\definecolor{currentstroke}{rgb}{0.000000,0.000000,0.000000}%
\pgfsetstrokecolor{currentstroke}%
\pgfsetdash{}{0pt}%
\pgfsys@defobject{currentmarker}{\pgfqpoint{0.000000in}{-0.048611in}}{\pgfqpoint{0.000000in}{0.000000in}}{%
\pgfpathmoveto{\pgfqpoint{0.000000in}{0.000000in}}%
\pgfpathlineto{\pgfqpoint{0.000000in}{-0.048611in}}%
\pgfusepath{stroke,fill}%
}%
\begin{pgfscope}%
\pgfsys@transformshift{4.059174in}{0.771190in}%
\pgfsys@useobject{currentmarker}{}%
\end{pgfscope}%
\end{pgfscope}%
\begin{pgfscope}%
\definecolor{textcolor}{rgb}{0.000000,0.000000,0.000000}%
\pgfsetstrokecolor{textcolor}%
\pgfsetfillcolor{textcolor}%
\pgftext[x=4.059174in,y=0.673968in,,top]{\color{textcolor}{\rmfamily\fontsize{10.000000}{12.000000}\selectfont\catcode`\^=\active\def^{\ifmmode\sp\else\^{}\fi}\catcode`\%=\active\def%{\%}$\mathdefault{60}$}}%
\end{pgfscope}%
\begin{pgfscope}%
\pgfsetbuttcap%
\pgfsetroundjoin%
\definecolor{currentfill}{rgb}{0.000000,0.000000,0.000000}%
\pgfsetfillcolor{currentfill}%
\pgfsetlinewidth{0.803000pt}%
\definecolor{currentstroke}{rgb}{0.000000,0.000000,0.000000}%
\pgfsetstrokecolor{currentstroke}%
\pgfsetdash{}{0pt}%
\pgfsys@defobject{currentmarker}{\pgfqpoint{0.000000in}{-0.048611in}}{\pgfqpoint{0.000000in}{0.000000in}}{%
\pgfpathmoveto{\pgfqpoint{0.000000in}{0.000000in}}%
\pgfpathlineto{\pgfqpoint{0.000000in}{-0.048611in}}%
\pgfusepath{stroke,fill}%
}%
\begin{pgfscope}%
\pgfsys@transformshift{5.032874in}{0.771190in}%
\pgfsys@useobject{currentmarker}{}%
\end{pgfscope}%
\end{pgfscope}%
\begin{pgfscope}%
\definecolor{textcolor}{rgb}{0.000000,0.000000,0.000000}%
\pgfsetstrokecolor{textcolor}%
\pgfsetfillcolor{textcolor}%
\pgftext[x=5.032874in,y=0.673968in,,top]{\color{textcolor}{\rmfamily\fontsize{10.000000}{12.000000}\selectfont\catcode`\^=\active\def^{\ifmmode\sp\else\^{}\fi}\catcode`\%=\active\def%{\%}$\mathdefault{80}$}}%
\end{pgfscope}%
\begin{pgfscope}%
\pgfsetbuttcap%
\pgfsetroundjoin%
\definecolor{currentfill}{rgb}{0.000000,0.000000,0.000000}%
\pgfsetfillcolor{currentfill}%
\pgfsetlinewidth{0.803000pt}%
\definecolor{currentstroke}{rgb}{0.000000,0.000000,0.000000}%
\pgfsetstrokecolor{currentstroke}%
\pgfsetdash{}{0pt}%
\pgfsys@defobject{currentmarker}{\pgfqpoint{0.000000in}{-0.048611in}}{\pgfqpoint{0.000000in}{0.000000in}}{%
\pgfpathmoveto{\pgfqpoint{0.000000in}{0.000000in}}%
\pgfpathlineto{\pgfqpoint{0.000000in}{-0.048611in}}%
\pgfusepath{stroke,fill}%
}%
\begin{pgfscope}%
\pgfsys@transformshift{6.006575in}{0.771190in}%
\pgfsys@useobject{currentmarker}{}%
\end{pgfscope}%
\end{pgfscope}%
\begin{pgfscope}%
\definecolor{textcolor}{rgb}{0.000000,0.000000,0.000000}%
\pgfsetstrokecolor{textcolor}%
\pgfsetfillcolor{textcolor}%
\pgftext[x=6.006575in,y=0.673968in,,top]{\color{textcolor}{\rmfamily\fontsize{10.000000}{12.000000}\selectfont\catcode`\^=\active\def^{\ifmmode\sp\else\^{}\fi}\catcode`\%=\active\def%{\%}$\mathdefault{100}$}}%
\end{pgfscope}%
\begin{pgfscope}%
\definecolor{textcolor}{rgb}{0.000000,0.000000,0.000000}%
\pgfsetstrokecolor{textcolor}%
\pgfsetfillcolor{textcolor}%
\pgftext[x=3.572324in,y=0.494955in,,top]{\color{textcolor}{\rmfamily\fontsize{24.000000}{28.800000}\selectfont\catcode`\^=\active\def^{\ifmmode\sp\else\^{}\fi}\catcode`\%=\active\def%{\%}$t (s)$}}%
\end{pgfscope}%
\begin{pgfscope}%
\pgfsetbuttcap%
\pgfsetroundjoin%
\definecolor{currentfill}{rgb}{0.000000,0.000000,0.000000}%
\pgfsetfillcolor{currentfill}%
\pgfsetlinewidth{0.803000pt}%
\definecolor{currentstroke}{rgb}{0.000000,0.000000,0.000000}%
\pgfsetstrokecolor{currentstroke}%
\pgfsetdash{}{0pt}%
\pgfsys@defobject{currentmarker}{\pgfqpoint{-0.048611in}{0.000000in}}{\pgfqpoint{-0.000000in}{0.000000in}}{%
\pgfpathmoveto{\pgfqpoint{-0.000000in}{0.000000in}}%
\pgfpathlineto{\pgfqpoint{-0.048611in}{0.000000in}}%
\pgfusepath{stroke,fill}%
}%
\begin{pgfscope}%
\pgfsys@transformshift{0.894648in}{0.771190in}%
\pgfsys@useobject{currentmarker}{}%
\end{pgfscope}%
\end{pgfscope}%
\begin{pgfscope}%
\definecolor{textcolor}{rgb}{0.000000,0.000000,0.000000}%
\pgfsetstrokecolor{textcolor}%
\pgfsetfillcolor{textcolor}%
\pgftext[x=0.550511in, y=0.722965in, left, base]{\color{textcolor}{\rmfamily\fontsize{10.000000}{12.000000}\selectfont\catcode`\^=\active\def^{\ifmmode\sp\else\^{}\fi}\catcode`\%=\active\def%{\%}$\mathdefault{\ensuremath{-}40}$}}%
\end{pgfscope}%
\begin{pgfscope}%
\pgfsetbuttcap%
\pgfsetroundjoin%
\definecolor{currentfill}{rgb}{0.000000,0.000000,0.000000}%
\pgfsetfillcolor{currentfill}%
\pgfsetlinewidth{0.803000pt}%
\definecolor{currentstroke}{rgb}{0.000000,0.000000,0.000000}%
\pgfsetstrokecolor{currentstroke}%
\pgfsetdash{}{0pt}%
\pgfsys@defobject{currentmarker}{\pgfqpoint{-0.048611in}{0.000000in}}{\pgfqpoint{-0.000000in}{0.000000in}}{%
\pgfpathmoveto{\pgfqpoint{-0.000000in}{0.000000in}}%
\pgfpathlineto{\pgfqpoint{-0.048611in}{0.000000in}}%
\pgfusepath{stroke,fill}%
}%
\begin{pgfscope}%
\pgfsys@transformshift{0.894648in}{1.537307in}%
\pgfsys@useobject{currentmarker}{}%
\end{pgfscope}%
\end{pgfscope}%
\begin{pgfscope}%
\definecolor{textcolor}{rgb}{0.000000,0.000000,0.000000}%
\pgfsetstrokecolor{textcolor}%
\pgfsetfillcolor{textcolor}%
\pgftext[x=0.550511in, y=1.489082in, left, base]{\color{textcolor}{\rmfamily\fontsize{10.000000}{12.000000}\selectfont\catcode`\^=\active\def^{\ifmmode\sp\else\^{}\fi}\catcode`\%=\active\def%{\%}$\mathdefault{\ensuremath{-}20}$}}%
\end{pgfscope}%
\begin{pgfscope}%
\pgfsetbuttcap%
\pgfsetroundjoin%
\definecolor{currentfill}{rgb}{0.000000,0.000000,0.000000}%
\pgfsetfillcolor{currentfill}%
\pgfsetlinewidth{0.803000pt}%
\definecolor{currentstroke}{rgb}{0.000000,0.000000,0.000000}%
\pgfsetstrokecolor{currentstroke}%
\pgfsetdash{}{0pt}%
\pgfsys@defobject{currentmarker}{\pgfqpoint{-0.048611in}{0.000000in}}{\pgfqpoint{-0.000000in}{0.000000in}}{%
\pgfpathmoveto{\pgfqpoint{-0.000000in}{0.000000in}}%
\pgfpathlineto{\pgfqpoint{-0.048611in}{0.000000in}}%
\pgfusepath{stroke,fill}%
}%
\begin{pgfscope}%
\pgfsys@transformshift{0.894648in}{2.303424in}%
\pgfsys@useobject{currentmarker}{}%
\end{pgfscope}%
\end{pgfscope}%
\begin{pgfscope}%
\definecolor{textcolor}{rgb}{0.000000,0.000000,0.000000}%
\pgfsetstrokecolor{textcolor}%
\pgfsetfillcolor{textcolor}%
\pgftext[x=0.727981in, y=2.255199in, left, base]{\color{textcolor}{\rmfamily\fontsize{10.000000}{12.000000}\selectfont\catcode`\^=\active\def^{\ifmmode\sp\else\^{}\fi}\catcode`\%=\active\def%{\%}$\mathdefault{0}$}}%
\end{pgfscope}%
\begin{pgfscope}%
\pgfsetbuttcap%
\pgfsetroundjoin%
\definecolor{currentfill}{rgb}{0.000000,0.000000,0.000000}%
\pgfsetfillcolor{currentfill}%
\pgfsetlinewidth{0.803000pt}%
\definecolor{currentstroke}{rgb}{0.000000,0.000000,0.000000}%
\pgfsetstrokecolor{currentstroke}%
\pgfsetdash{}{0pt}%
\pgfsys@defobject{currentmarker}{\pgfqpoint{-0.048611in}{0.000000in}}{\pgfqpoint{-0.000000in}{0.000000in}}{%
\pgfpathmoveto{\pgfqpoint{-0.000000in}{0.000000in}}%
\pgfpathlineto{\pgfqpoint{-0.048611in}{0.000000in}}%
\pgfusepath{stroke,fill}%
}%
\begin{pgfscope}%
\pgfsys@transformshift{0.894648in}{3.069541in}%
\pgfsys@useobject{currentmarker}{}%
\end{pgfscope}%
\end{pgfscope}%
\begin{pgfscope}%
\definecolor{textcolor}{rgb}{0.000000,0.000000,0.000000}%
\pgfsetstrokecolor{textcolor}%
\pgfsetfillcolor{textcolor}%
\pgftext[x=0.658536in, y=3.021316in, left, base]{\color{textcolor}{\rmfamily\fontsize{10.000000}{12.000000}\selectfont\catcode`\^=\active\def^{\ifmmode\sp\else\^{}\fi}\catcode`\%=\active\def%{\%}$\mathdefault{20}$}}%
\end{pgfscope}%
\begin{pgfscope}%
\pgfsetbuttcap%
\pgfsetroundjoin%
\definecolor{currentfill}{rgb}{0.000000,0.000000,0.000000}%
\pgfsetfillcolor{currentfill}%
\pgfsetlinewidth{0.803000pt}%
\definecolor{currentstroke}{rgb}{0.000000,0.000000,0.000000}%
\pgfsetstrokecolor{currentstroke}%
\pgfsetdash{}{0pt}%
\pgfsys@defobject{currentmarker}{\pgfqpoint{-0.048611in}{0.000000in}}{\pgfqpoint{-0.000000in}{0.000000in}}{%
\pgfpathmoveto{\pgfqpoint{-0.000000in}{0.000000in}}%
\pgfpathlineto{\pgfqpoint{-0.048611in}{0.000000in}}%
\pgfusepath{stroke,fill}%
}%
\begin{pgfscope}%
\pgfsys@transformshift{0.894648in}{3.835658in}%
\pgfsys@useobject{currentmarker}{}%
\end{pgfscope}%
\end{pgfscope}%
\begin{pgfscope}%
\definecolor{textcolor}{rgb}{0.000000,0.000000,0.000000}%
\pgfsetstrokecolor{textcolor}%
\pgfsetfillcolor{textcolor}%
\pgftext[x=0.658536in, y=3.787432in, left, base]{\color{textcolor}{\rmfamily\fontsize{10.000000}{12.000000}\selectfont\catcode`\^=\active\def^{\ifmmode\sp\else\^{}\fi}\catcode`\%=\active\def%{\%}$\mathdefault{40}$}}%
\end{pgfscope}%
\begin{pgfscope}%
\pgfsetbuttcap%
\pgfsetroundjoin%
\definecolor{currentfill}{rgb}{0.000000,0.000000,0.000000}%
\pgfsetfillcolor{currentfill}%
\pgfsetlinewidth{0.803000pt}%
\definecolor{currentstroke}{rgb}{0.000000,0.000000,0.000000}%
\pgfsetstrokecolor{currentstroke}%
\pgfsetdash{}{0pt}%
\pgfsys@defobject{currentmarker}{\pgfqpoint{-0.048611in}{0.000000in}}{\pgfqpoint{-0.000000in}{0.000000in}}{%
\pgfpathmoveto{\pgfqpoint{-0.000000in}{0.000000in}}%
\pgfpathlineto{\pgfqpoint{-0.048611in}{0.000000in}}%
\pgfusepath{stroke,fill}%
}%
\begin{pgfscope}%
\pgfsys@transformshift{0.894648in}{4.601775in}%
\pgfsys@useobject{currentmarker}{}%
\end{pgfscope}%
\end{pgfscope}%
\begin{pgfscope}%
\definecolor{textcolor}{rgb}{0.000000,0.000000,0.000000}%
\pgfsetstrokecolor{textcolor}%
\pgfsetfillcolor{textcolor}%
\pgftext[x=0.658536in, y=4.553549in, left, base]{\color{textcolor}{\rmfamily\fontsize{10.000000}{12.000000}\selectfont\catcode`\^=\active\def^{\ifmmode\sp\else\^{}\fi}\catcode`\%=\active\def%{\%}$\mathdefault{60}$}}%
\end{pgfscope}%
\begin{pgfscope}%
\definecolor{textcolor}{rgb}{0.000000,0.000000,0.000000}%
\pgfsetstrokecolor{textcolor}%
\pgfsetfillcolor{textcolor}%
\pgftext[x=0.494955in,y=2.686482in,,bottom,rotate=90.000000]{\color{textcolor}{\rmfamily\fontsize{24.000000}{28.800000}\selectfont\catcode`\^=\active\def^{\ifmmode\sp\else\^{}\fi}\catcode`\%=\active\def%{\%}$x (\mu m)$}}%
\end{pgfscope}%
\begin{pgfscope}%
\pgfpathrectangle{\pgfqpoint{0.894648in}{0.771190in}}{\pgfqpoint{5.355353in}{3.830585in}}%
\pgfusepath{clip}%
\pgfsetrectcap%
\pgfsetroundjoin%
\pgfsetlinewidth{1.505625pt}%
\definecolor{currentstroke}{rgb}{0.121569,0.466667,0.705882}%
\pgfsetstrokecolor{currentstroke}%
\pgfsetdash{}{0pt}%
\pgfpathmoveto{\pgfqpoint{1.138073in}{2.303424in}}%
\pgfpathlineto{\pgfqpoint{1.138073in}{2.303424in}}%
\pgfpathlineto{\pgfqpoint{1.138102in}{2.300827in}}%
\pgfpathlineto{\pgfqpoint{1.138876in}{2.340053in}}%
\pgfpathlineto{\pgfqpoint{1.139017in}{2.335176in}}%
\pgfpathlineto{\pgfqpoint{1.139962in}{2.363453in}}%
\pgfpathlineto{\pgfqpoint{1.139338in}{2.318073in}}%
\pgfpathlineto{\pgfqpoint{1.140161in}{2.361922in}}%
\pgfpathlineto{\pgfqpoint{1.140390in}{2.368733in}}%
\pgfpathlineto{\pgfqpoint{1.140853in}{2.340110in}}%
\pgfpathlineto{\pgfqpoint{1.141208in}{2.346353in}}%
\pgfpathlineto{\pgfqpoint{1.142201in}{2.331992in}}%
\pgfpathlineto{\pgfqpoint{1.141597in}{2.354594in}}%
\pgfpathlineto{\pgfqpoint{1.142313in}{2.344473in}}%
\pgfpathlineto{\pgfqpoint{1.143019in}{2.385232in}}%
\pgfpathlineto{\pgfqpoint{1.143457in}{2.366991in}}%
\pgfpathlineto{\pgfqpoint{1.144368in}{2.330634in}}%
\pgfpathlineto{\pgfqpoint{1.143486in}{2.375994in}}%
\pgfpathlineto{\pgfqpoint{1.144669in}{2.336104in}}%
\pgfpathlineto{\pgfqpoint{1.145565in}{2.281867in}}%
\pgfpathlineto{\pgfqpoint{1.144699in}{2.336357in}}%
\pgfpathlineto{\pgfqpoint{1.145994in}{2.311388in}}%
\pgfpathlineto{\pgfqpoint{1.147109in}{2.334478in}}%
\pgfpathlineto{\pgfqpoint{1.146568in}{2.287779in}}%
\pgfpathlineto{\pgfqpoint{1.147148in}{2.324441in}}%
\pgfpathlineto{\pgfqpoint{1.148034in}{2.308670in}}%
\pgfpathlineto{\pgfqpoint{1.147479in}{2.338157in}}%
\pgfpathlineto{\pgfqpoint{1.148301in}{2.314487in}}%
\pgfpathlineto{\pgfqpoint{1.148369in}{2.328842in}}%
\pgfpathlineto{\pgfqpoint{1.149251in}{2.270682in}}%
\pgfpathlineto{\pgfqpoint{1.150526in}{2.309280in}}%
\pgfpathlineto{\pgfqpoint{1.149333in}{2.260316in}}%
\pgfpathlineto{\pgfqpoint{1.150672in}{2.293190in}}%
\pgfpathlineto{\pgfqpoint{1.151028in}{2.270229in}}%
\pgfpathlineto{\pgfqpoint{1.151714in}{2.300197in}}%
\pgfpathlineto{\pgfqpoint{1.151797in}{2.282774in}}%
\pgfpathlineto{\pgfqpoint{1.152250in}{2.324778in}}%
\pgfpathlineto{\pgfqpoint{1.151816in}{2.281686in}}%
\pgfpathlineto{\pgfqpoint{1.152951in}{2.287017in}}%
\pgfpathlineto{\pgfqpoint{1.154148in}{2.236819in}}%
\pgfpathlineto{\pgfqpoint{1.152965in}{2.287107in}}%
\pgfpathlineto{\pgfqpoint{1.154187in}{2.237271in}}%
\pgfpathlineto{\pgfqpoint{1.154981in}{2.264154in}}%
\pgfpathlineto{\pgfqpoint{1.155219in}{2.235613in}}%
\pgfpathlineto{\pgfqpoint{1.155297in}{2.241237in}}%
\pgfpathlineto{\pgfqpoint{1.156480in}{2.211260in}}%
\pgfpathlineto{\pgfqpoint{1.155327in}{2.243414in}}%
\pgfpathlineto{\pgfqpoint{1.156553in}{2.219280in}}%
\pgfpathlineto{\pgfqpoint{1.157581in}{2.247651in}}%
\pgfpathlineto{\pgfqpoint{1.156651in}{2.209615in}}%
\pgfpathlineto{\pgfqpoint{1.157736in}{2.228002in}}%
\pgfpathlineto{\pgfqpoint{1.158384in}{2.218816in}}%
\pgfpathlineto{\pgfqpoint{1.158520in}{2.244500in}}%
\pgfpathlineto{\pgfqpoint{1.158788in}{2.236948in}}%
\pgfpathlineto{\pgfqpoint{1.159255in}{2.259943in}}%
\pgfpathlineto{\pgfqpoint{1.159056in}{2.233001in}}%
\pgfpathlineto{\pgfqpoint{1.159918in}{2.241866in}}%
\pgfpathlineto{\pgfqpoint{1.160254in}{2.209890in}}%
\pgfpathlineto{\pgfqpoint{1.161023in}{2.242327in}}%
\pgfpathlineto{\pgfqpoint{1.161028in}{2.241889in}}%
\pgfpathlineto{\pgfqpoint{1.161500in}{2.254987in}}%
\pgfpathlineto{\pgfqpoint{1.161894in}{2.222518in}}%
\pgfpathlineto{\pgfqpoint{1.162069in}{2.223925in}}%
\pgfpathlineto{\pgfqpoint{1.162951in}{2.177545in}}%
\pgfpathlineto{\pgfqpoint{1.162444in}{2.226082in}}%
\pgfpathlineto{\pgfqpoint{1.163340in}{2.182998in}}%
\pgfpathlineto{\pgfqpoint{1.164479in}{2.216985in}}%
\pgfpathlineto{\pgfqpoint{1.163545in}{2.176654in}}%
\pgfpathlineto{\pgfqpoint{1.164509in}{2.213400in}}%
\pgfpathlineto{\pgfqpoint{1.165482in}{2.192615in}}%
\pgfpathlineto{\pgfqpoint{1.165585in}{2.214016in}}%
\pgfpathlineto{\pgfqpoint{1.165623in}{2.217175in}}%
\pgfpathlineto{\pgfqpoint{1.166427in}{2.171964in}}%
\pgfpathlineto{\pgfqpoint{1.166578in}{2.179175in}}%
\pgfpathlineto{\pgfqpoint{1.167108in}{2.166908in}}%
\pgfpathlineto{\pgfqpoint{1.167225in}{2.189279in}}%
\pgfpathlineto{\pgfqpoint{1.167678in}{2.179601in}}%
\pgfpathlineto{\pgfqpoint{1.167717in}{2.188373in}}%
\pgfpathlineto{\pgfqpoint{1.168111in}{2.149279in}}%
\pgfpathlineto{\pgfqpoint{1.168744in}{2.158815in}}%
\pgfpathlineto{\pgfqpoint{1.169582in}{2.131290in}}%
\pgfpathlineto{\pgfqpoint{1.168812in}{2.160489in}}%
\pgfpathlineto{\pgfqpoint{1.169864in}{2.153296in}}%
\pgfpathlineto{\pgfqpoint{1.170005in}{2.170343in}}%
\pgfpathlineto{\pgfqpoint{1.170828in}{2.134320in}}%
\pgfpathlineto{\pgfqpoint{1.170954in}{2.143783in}}%
\pgfpathlineto{\pgfqpoint{1.171086in}{2.136428in}}%
\pgfpathlineto{\pgfqpoint{1.171500in}{2.175412in}}%
\pgfpathlineto{\pgfqpoint{1.172045in}{2.155183in}}%
\pgfpathlineto{\pgfqpoint{1.172780in}{2.112586in}}%
\pgfpathlineto{\pgfqpoint{1.173063in}{2.157302in}}%
\pgfpathlineto{\pgfqpoint{1.173155in}{2.151674in}}%
\pgfpathlineto{\pgfqpoint{1.174255in}{2.173174in}}%
\pgfpathlineto{\pgfqpoint{1.173403in}{2.145238in}}%
\pgfpathlineto{\pgfqpoint{1.174299in}{2.173075in}}%
\pgfpathlineto{\pgfqpoint{1.176091in}{2.124651in}}%
\pgfpathlineto{\pgfqpoint{1.174494in}{2.189100in}}%
\pgfpathlineto{\pgfqpoint{1.176135in}{2.138176in}}%
\pgfpathlineto{\pgfqpoint{1.176227in}{2.146260in}}%
\pgfpathlineto{\pgfqpoint{1.176719in}{2.114701in}}%
\pgfpathlineto{\pgfqpoint{1.177254in}{2.140764in}}%
\pgfpathlineto{\pgfqpoint{1.178447in}{2.168244in}}%
\pgfpathlineto{\pgfqpoint{1.177376in}{2.127607in}}%
\pgfpathlineto{\pgfqpoint{1.178588in}{2.163169in}}%
\pgfpathlineto{\pgfqpoint{1.179689in}{2.121904in}}%
\pgfpathlineto{\pgfqpoint{1.178788in}{2.168579in}}%
\pgfpathlineto{\pgfqpoint{1.179781in}{2.136001in}}%
\pgfpathlineto{\pgfqpoint{1.180711in}{2.153551in}}%
\pgfpathlineto{\pgfqpoint{1.180127in}{2.114056in}}%
\pgfpathlineto{\pgfqpoint{1.180911in}{2.144829in}}%
\pgfpathlineto{\pgfqpoint{1.182225in}{2.193717in}}%
\pgfpathlineto{\pgfqpoint{1.181558in}{2.142310in}}%
\pgfpathlineto{\pgfqpoint{1.182288in}{2.188371in}}%
\pgfpathlineto{\pgfqpoint{1.182342in}{2.176179in}}%
\pgfpathlineto{\pgfqpoint{1.183097in}{2.198168in}}%
\pgfpathlineto{\pgfqpoint{1.183447in}{2.179242in}}%
\pgfpathlineto{\pgfqpoint{1.184533in}{2.212387in}}%
\pgfpathlineto{\pgfqpoint{1.184654in}{2.207896in}}%
\pgfpathlineto{\pgfqpoint{1.185643in}{2.164789in}}%
\pgfpathlineto{\pgfqpoint{1.184718in}{2.209089in}}%
\pgfpathlineto{\pgfqpoint{1.185818in}{2.178471in}}%
\pgfpathlineto{\pgfqpoint{1.186368in}{2.162657in}}%
\pgfpathlineto{\pgfqpoint{1.185969in}{2.191523in}}%
\pgfpathlineto{\pgfqpoint{1.186923in}{2.177607in}}%
\pgfpathlineto{\pgfqpoint{1.187298in}{2.224111in}}%
\pgfpathlineto{\pgfqpoint{1.187858in}{2.174361in}}%
\pgfpathlineto{\pgfqpoint{1.188062in}{2.187105in}}%
\pgfpathlineto{\pgfqpoint{1.188189in}{2.172059in}}%
\pgfpathlineto{\pgfqpoint{1.188447in}{2.205052in}}%
\pgfpathlineto{\pgfqpoint{1.189138in}{2.198187in}}%
\pgfpathlineto{\pgfqpoint{1.189455in}{2.181973in}}%
\pgfpathlineto{\pgfqpoint{1.189197in}{2.206638in}}%
\pgfpathlineto{\pgfqpoint{1.189830in}{2.203055in}}%
\pgfpathlineto{\pgfqpoint{1.190156in}{2.210817in}}%
\pgfpathlineto{\pgfqpoint{1.190725in}{2.187089in}}%
\pgfpathlineto{\pgfqpoint{1.190930in}{2.202472in}}%
\pgfpathlineto{\pgfqpoint{1.191300in}{2.178614in}}%
\pgfpathlineto{\pgfqpoint{1.191758in}{2.228592in}}%
\pgfpathlineto{\pgfqpoint{1.191952in}{2.219682in}}%
\pgfpathlineto{\pgfqpoint{1.192381in}{2.256459in}}%
\pgfpathlineto{\pgfqpoint{1.193072in}{2.227105in}}%
\pgfpathlineto{\pgfqpoint{1.193987in}{2.211629in}}%
\pgfpathlineto{\pgfqpoint{1.193559in}{2.237334in}}%
\pgfpathlineto{\pgfqpoint{1.194177in}{2.228692in}}%
\pgfpathlineto{\pgfqpoint{1.194576in}{2.246576in}}%
\pgfpathlineto{\pgfqpoint{1.195200in}{2.175942in}}%
\pgfpathlineto{\pgfqpoint{1.195759in}{2.168974in}}%
\pgfpathlineto{\pgfqpoint{1.195657in}{2.193278in}}%
\pgfpathlineto{\pgfqpoint{1.196232in}{2.184469in}}%
\pgfpathlineto{\pgfqpoint{1.197313in}{2.203811in}}%
\pgfpathlineto{\pgfqpoint{1.196534in}{2.160506in}}%
\pgfpathlineto{\pgfqpoint{1.197351in}{2.194796in}}%
\pgfpathlineto{\pgfqpoint{1.197425in}{2.188426in}}%
\pgfpathlineto{\pgfqpoint{1.197907in}{2.228761in}}%
\pgfpathlineto{\pgfqpoint{1.198413in}{2.209272in}}%
\pgfpathlineto{\pgfqpoint{1.198710in}{2.222210in}}%
\pgfpathlineto{\pgfqpoint{1.199197in}{2.187204in}}%
\pgfpathlineto{\pgfqpoint{1.199503in}{2.200929in}}%
\pgfpathlineto{\pgfqpoint{1.199513in}{2.195915in}}%
\pgfpathlineto{\pgfqpoint{1.200277in}{2.231871in}}%
\pgfpathlineto{\pgfqpoint{1.200385in}{2.230138in}}%
\pgfpathlineto{\pgfqpoint{1.200424in}{2.236603in}}%
\pgfpathlineto{\pgfqpoint{1.201280in}{2.184582in}}%
\pgfpathlineto{\pgfqpoint{1.202303in}{2.149074in}}%
\pgfpathlineto{\pgfqpoint{1.202473in}{2.160461in}}%
\pgfpathlineto{\pgfqpoint{1.202945in}{2.173886in}}%
\pgfpathlineto{\pgfqpoint{1.203476in}{2.134025in}}%
\pgfpathlineto{\pgfqpoint{1.203569in}{2.147798in}}%
\pgfpathlineto{\pgfqpoint{1.203856in}{2.168917in}}%
\pgfpathlineto{\pgfqpoint{1.204552in}{2.142347in}}%
\pgfpathlineto{\pgfqpoint{1.204562in}{2.143797in}}%
\pgfpathlineto{\pgfqpoint{1.205092in}{2.125762in}}%
\pgfpathlineto{\pgfqpoint{1.205195in}{2.145187in}}%
\pgfpathlineto{\pgfqpoint{1.205677in}{2.138142in}}%
\pgfpathlineto{\pgfqpoint{1.206134in}{2.157571in}}%
\pgfpathlineto{\pgfqpoint{1.206665in}{2.127516in}}%
\pgfpathlineto{\pgfqpoint{1.206787in}{2.140034in}}%
\pgfpathlineto{\pgfqpoint{1.207342in}{2.168562in}}%
\pgfpathlineto{\pgfqpoint{1.206826in}{2.135395in}}%
\pgfpathlineto{\pgfqpoint{1.207843in}{2.141397in}}%
\pgfpathlineto{\pgfqpoint{1.207936in}{2.137016in}}%
\pgfpathlineto{\pgfqpoint{1.208861in}{2.190965in}}%
\pgfpathlineto{\pgfqpoint{1.210019in}{2.244128in}}%
\pgfpathlineto{\pgfqpoint{1.210058in}{2.242454in}}%
\pgfpathlineto{\pgfqpoint{1.210097in}{2.247491in}}%
\pgfpathlineto{\pgfqpoint{1.210350in}{2.221927in}}%
\pgfpathlineto{\pgfqpoint{1.210608in}{2.228573in}}%
\pgfpathlineto{\pgfqpoint{1.210935in}{2.198582in}}%
\pgfpathlineto{\pgfqpoint{1.211499in}{2.237751in}}%
\pgfpathlineto{\pgfqpoint{1.211718in}{2.227362in}}%
\pgfpathlineto{\pgfqpoint{1.211889in}{2.234438in}}%
\pgfpathlineto{\pgfqpoint{1.212239in}{2.197454in}}%
\pgfpathlineto{\pgfqpoint{1.212746in}{2.214121in}}%
\pgfpathlineto{\pgfqpoint{1.213627in}{2.207447in}}%
\pgfpathlineto{\pgfqpoint{1.213086in}{2.243908in}}%
\pgfpathlineto{\pgfqpoint{1.213807in}{2.224461in}}%
\pgfpathlineto{\pgfqpoint{1.213914in}{2.236826in}}%
\pgfpathlineto{\pgfqpoint{1.214868in}{2.206299in}}%
\pgfpathlineto{\pgfqpoint{1.214883in}{2.209926in}}%
\pgfpathlineto{\pgfqpoint{1.216110in}{2.137323in}}%
\pgfpathlineto{\pgfqpoint{1.216154in}{2.144114in}}%
\pgfpathlineto{\pgfqpoint{1.216339in}{2.157844in}}%
\pgfpathlineto{\pgfqpoint{1.217059in}{2.125196in}}%
\pgfpathlineto{\pgfqpoint{1.217254in}{2.136623in}}%
\pgfpathlineto{\pgfqpoint{1.217434in}{2.127224in}}%
\pgfpathlineto{\pgfqpoint{1.218111in}{2.179603in}}%
\pgfpathlineto{\pgfqpoint{1.218228in}{2.170515in}}%
\pgfpathlineto{\pgfqpoint{1.218252in}{2.172470in}}%
\pgfpathlineto{\pgfqpoint{1.218267in}{2.168901in}}%
\pgfpathlineto{\pgfqpoint{1.218938in}{2.126641in}}%
\pgfpathlineto{\pgfqpoint{1.219469in}{2.126870in}}%
\pgfpathlineto{\pgfqpoint{1.220044in}{2.148712in}}%
\pgfpathlineto{\pgfqpoint{1.220501in}{2.089208in}}%
\pgfpathlineto{\pgfqpoint{1.220516in}{2.091768in}}%
\pgfpathlineto{\pgfqpoint{1.221947in}{2.007787in}}%
\pgfpathlineto{\pgfqpoint{1.222093in}{2.011481in}}%
\pgfpathlineto{\pgfqpoint{1.222862in}{2.042996in}}%
\pgfpathlineto{\pgfqpoint{1.223296in}{2.038694in}}%
\pgfpathlineto{\pgfqpoint{1.224586in}{1.963104in}}%
\pgfpathlineto{\pgfqpoint{1.223325in}{2.039956in}}%
\pgfpathlineto{\pgfqpoint{1.224615in}{1.965684in}}%
\pgfpathlineto{\pgfqpoint{1.225452in}{1.949596in}}%
\pgfpathlineto{\pgfqpoint{1.225350in}{1.974754in}}%
\pgfpathlineto{\pgfqpoint{1.225657in}{1.971404in}}%
\pgfpathlineto{\pgfqpoint{1.225662in}{1.972129in}}%
\pgfpathlineto{\pgfqpoint{1.226431in}{1.938871in}}%
\pgfpathlineto{\pgfqpoint{1.226563in}{1.942045in}}%
\pgfpathlineto{\pgfqpoint{1.226825in}{1.927621in}}%
\pgfpathlineto{\pgfqpoint{1.227556in}{1.969005in}}%
\pgfpathlineto{\pgfqpoint{1.227600in}{1.975708in}}%
\pgfpathlineto{\pgfqpoint{1.228607in}{1.932650in}}%
\pgfpathlineto{\pgfqpoint{1.229450in}{1.958106in}}%
\pgfpathlineto{\pgfqpoint{1.229045in}{1.921890in}}%
\pgfpathlineto{\pgfqpoint{1.229863in}{1.952371in}}%
\pgfpathlineto{\pgfqpoint{1.230039in}{1.948598in}}%
\pgfpathlineto{\pgfqpoint{1.230720in}{1.989547in}}%
\pgfpathlineto{\pgfqpoint{1.230788in}{1.989123in}}%
\pgfpathlineto{\pgfqpoint{1.231027in}{2.011045in}}%
\pgfpathlineto{\pgfqpoint{1.231548in}{1.949937in}}%
\pgfpathlineto{\pgfqpoint{1.231747in}{1.951161in}}%
\pgfpathlineto{\pgfqpoint{1.231869in}{1.942567in}}%
\pgfpathlineto{\pgfqpoint{1.232619in}{1.994760in}}%
\pgfpathlineto{\pgfqpoint{1.232716in}{1.991428in}}%
\pgfpathlineto{\pgfqpoint{1.233257in}{2.003422in}}%
\pgfpathlineto{\pgfqpoint{1.233461in}{1.983335in}}%
\pgfpathlineto{\pgfqpoint{1.233666in}{1.987119in}}%
\pgfpathlineto{\pgfqpoint{1.233987in}{1.973243in}}%
\pgfpathlineto{\pgfqpoint{1.234391in}{1.998340in}}%
\pgfpathlineto{\pgfqpoint{1.234776in}{1.986011in}}%
\pgfpathlineto{\pgfqpoint{1.235189in}{2.017755in}}%
\pgfpathlineto{\pgfqpoint{1.234795in}{1.983649in}}%
\pgfpathlineto{\pgfqpoint{1.236080in}{2.002347in}}%
\pgfpathlineto{\pgfqpoint{1.237093in}{1.972593in}}%
\pgfpathlineto{\pgfqpoint{1.236708in}{2.008177in}}%
\pgfpathlineto{\pgfqpoint{1.237229in}{1.981110in}}%
\pgfpathlineto{\pgfqpoint{1.237268in}{1.987110in}}%
\pgfpathlineto{\pgfqpoint{1.238203in}{1.951984in}}%
\pgfpathlineto{\pgfqpoint{1.238310in}{1.956724in}}%
\pgfpathlineto{\pgfqpoint{1.238495in}{1.939222in}}%
\pgfpathlineto{\pgfqpoint{1.238875in}{1.985029in}}%
\pgfpathlineto{\pgfqpoint{1.239138in}{1.974739in}}%
\pgfpathlineto{\pgfqpoint{1.239493in}{1.993561in}}%
\pgfpathlineto{\pgfqpoint{1.240009in}{1.959161in}}%
\pgfpathlineto{\pgfqpoint{1.240238in}{1.969439in}}%
\pgfpathlineto{\pgfqpoint{1.240262in}{1.967714in}}%
\pgfpathlineto{\pgfqpoint{1.240267in}{1.971204in}}%
\pgfpathlineto{\pgfqpoint{1.240783in}{1.998473in}}%
\pgfpathlineto{\pgfqpoint{1.241319in}{1.941399in}}%
\pgfpathlineto{\pgfqpoint{1.241334in}{1.941575in}}%
\pgfpathlineto{\pgfqpoint{1.242721in}{2.011795in}}%
\pgfpathlineto{\pgfqpoint{1.241475in}{1.924225in}}%
\pgfpathlineto{\pgfqpoint{1.242726in}{2.011111in}}%
\pgfpathlineto{\pgfqpoint{1.242755in}{2.004081in}}%
\pgfpathlineto{\pgfqpoint{1.243422in}{2.045459in}}%
\pgfpathlineto{\pgfqpoint{1.243763in}{2.033765in}}%
\pgfpathlineto{\pgfqpoint{1.243919in}{2.036977in}}%
\pgfpathlineto{\pgfqpoint{1.244118in}{1.996866in}}%
\pgfpathlineto{\pgfqpoint{1.244610in}{2.012612in}}%
\pgfpathlineto{\pgfqpoint{1.245589in}{1.998404in}}%
\pgfpathlineto{\pgfqpoint{1.244824in}{2.024759in}}%
\pgfpathlineto{\pgfqpoint{1.245735in}{2.009051in}}%
\pgfpathlineto{\pgfqpoint{1.246426in}{2.028518in}}%
\pgfpathlineto{\pgfqpoint{1.245842in}{2.003724in}}%
\pgfpathlineto{\pgfqpoint{1.246898in}{2.026616in}}%
\pgfpathlineto{\pgfqpoint{1.246903in}{2.025574in}}%
\pgfpathlineto{\pgfqpoint{1.247740in}{2.069828in}}%
\pgfpathlineto{\pgfqpoint{1.247921in}{2.042000in}}%
\pgfpathlineto{\pgfqpoint{1.248816in}{2.092340in}}%
\pgfpathlineto{\pgfqpoint{1.247999in}{2.039761in}}%
\pgfpathlineto{\pgfqpoint{1.249079in}{2.083385in}}%
\pgfpathlineto{\pgfqpoint{1.249201in}{2.092636in}}%
\pgfpathlineto{\pgfqpoint{1.249766in}{2.024520in}}%
\pgfpathlineto{\pgfqpoint{1.250068in}{2.042240in}}%
\pgfpathlineto{\pgfqpoint{1.250598in}{2.013293in}}%
\pgfpathlineto{\pgfqpoint{1.250092in}{2.049474in}}%
\pgfpathlineto{\pgfqpoint{1.251265in}{2.026424in}}%
\pgfpathlineto{\pgfqpoint{1.251679in}{2.055753in}}%
\pgfpathlineto{\pgfqpoint{1.252171in}{2.016464in}}%
\pgfpathlineto{\pgfqpoint{1.252341in}{2.021088in}}%
\pgfpathlineto{\pgfqpoint{1.252356in}{2.016334in}}%
\pgfpathlineto{\pgfqpoint{1.253398in}{2.038488in}}%
\pgfpathlineto{\pgfqpoint{1.253636in}{2.037588in}}%
\pgfpathlineto{\pgfqpoint{1.253962in}{2.061558in}}%
\pgfpathlineto{\pgfqpoint{1.254357in}{2.049904in}}%
\pgfpathlineto{\pgfqpoint{1.254819in}{2.061076in}}%
\pgfpathlineto{\pgfqpoint{1.254552in}{2.039852in}}%
\pgfpathlineto{\pgfqpoint{1.255447in}{2.044088in}}%
\pgfpathlineto{\pgfqpoint{1.255486in}{2.039455in}}%
\pgfpathlineto{\pgfqpoint{1.255774in}{2.056014in}}%
\pgfpathlineto{\pgfqpoint{1.256591in}{2.069719in}}%
\pgfpathlineto{\pgfqpoint{1.256849in}{2.036889in}}%
\pgfpathlineto{\pgfqpoint{1.257229in}{2.026788in}}%
\pgfpathlineto{\pgfqpoint{1.257497in}{2.052739in}}%
\pgfpathlineto{\pgfqpoint{1.258062in}{2.061397in}}%
\pgfpathlineto{\pgfqpoint{1.258568in}{2.038925in}}%
\pgfpathlineto{\pgfqpoint{1.258617in}{2.024047in}}%
\pgfpathlineto{\pgfqpoint{1.259483in}{2.058515in}}%
\pgfpathlineto{\pgfqpoint{1.259620in}{2.057369in}}%
\pgfpathlineto{\pgfqpoint{1.259756in}{2.075936in}}%
\pgfpathlineto{\pgfqpoint{1.260671in}{2.046883in}}%
\pgfpathlineto{\pgfqpoint{1.260705in}{2.052190in}}%
\pgfpathlineto{\pgfqpoint{1.261139in}{2.031788in}}%
\pgfpathlineto{\pgfqpoint{1.261655in}{2.064551in}}%
\pgfpathlineto{\pgfqpoint{1.261786in}{2.057294in}}%
\pgfpathlineto{\pgfqpoint{1.262765in}{2.110850in}}%
\pgfpathlineto{\pgfqpoint{1.261825in}{2.055216in}}%
\pgfpathlineto{\pgfqpoint{1.263008in}{2.106297in}}%
\pgfpathlineto{\pgfqpoint{1.263417in}{2.061151in}}%
\pgfpathlineto{\pgfqpoint{1.263096in}{2.108206in}}%
\pgfpathlineto{\pgfqpoint{1.264201in}{2.068724in}}%
\pgfpathlineto{\pgfqpoint{1.265496in}{2.120371in}}%
\pgfpathlineto{\pgfqpoint{1.265501in}{2.118614in}}%
\pgfpathlineto{\pgfqpoint{1.266333in}{2.152498in}}%
\pgfpathlineto{\pgfqpoint{1.266518in}{2.151613in}}%
\pgfpathlineto{\pgfqpoint{1.266659in}{2.161695in}}%
\pgfpathlineto{\pgfqpoint{1.267317in}{2.117030in}}%
\pgfpathlineto{\pgfqpoint{1.267536in}{2.122492in}}%
\pgfpathlineto{\pgfqpoint{1.267974in}{2.135226in}}%
\pgfpathlineto{\pgfqpoint{1.268563in}{2.109177in}}%
\pgfpathlineto{\pgfqpoint{1.268583in}{2.112585in}}%
\pgfpathlineto{\pgfqpoint{1.269732in}{2.080333in}}%
\pgfpathlineto{\pgfqpoint{1.269975in}{2.090459in}}%
\pgfpathlineto{\pgfqpoint{1.271109in}{2.148847in}}%
\pgfpathlineto{\pgfqpoint{1.270087in}{2.076064in}}%
\pgfpathlineto{\pgfqpoint{1.271299in}{2.142675in}}%
\pgfpathlineto{\pgfqpoint{1.271392in}{2.149158in}}%
\pgfpathlineto{\pgfqpoint{1.271791in}{2.129081in}}%
\pgfpathlineto{\pgfqpoint{1.271825in}{2.130684in}}%
\pgfpathlineto{\pgfqpoint{1.272823in}{2.124017in}}%
\pgfpathlineto{\pgfqpoint{1.272161in}{2.152439in}}%
\pgfpathlineto{\pgfqpoint{1.272940in}{2.128360in}}%
\pgfpathlineto{\pgfqpoint{1.273076in}{2.132652in}}%
\pgfpathlineto{\pgfqpoint{1.273490in}{2.100111in}}%
\pgfpathlineto{\pgfqpoint{1.273991in}{2.117125in}}%
\pgfpathlineto{\pgfqpoint{1.274069in}{2.113337in}}%
\pgfpathlineto{\pgfqpoint{1.274245in}{2.138766in}}%
\pgfpathlineto{\pgfqpoint{1.274605in}{2.131257in}}%
\pgfpathlineto{\pgfqpoint{1.275067in}{2.150073in}}%
\pgfpathlineto{\pgfqpoint{1.275608in}{2.123323in}}%
\pgfpathlineto{\pgfqpoint{1.275705in}{2.129266in}}%
\pgfpathlineto{\pgfqpoint{1.275880in}{2.122361in}}%
\pgfpathlineto{\pgfqpoint{1.276494in}{2.187356in}}%
\pgfpathlineto{\pgfqpoint{1.276586in}{2.184030in}}%
\pgfpathlineto{\pgfqpoint{1.277171in}{2.177358in}}%
\pgfpathlineto{\pgfqpoint{1.277502in}{2.205659in}}%
\pgfpathlineto{\pgfqpoint{1.277526in}{2.204550in}}%
\pgfpathlineto{\pgfqpoint{1.278631in}{2.261027in}}%
\pgfpathlineto{\pgfqpoint{1.278694in}{2.256204in}}%
\pgfpathlineto{\pgfqpoint{1.279060in}{2.248456in}}%
\pgfpathlineto{\pgfqpoint{1.279697in}{2.276056in}}%
\pgfpathlineto{\pgfqpoint{1.279731in}{2.274432in}}%
\pgfpathlineto{\pgfqpoint{1.279790in}{2.282426in}}%
\pgfpathlineto{\pgfqpoint{1.280754in}{2.230944in}}%
\pgfpathlineto{\pgfqpoint{1.280773in}{2.233952in}}%
\pgfpathlineto{\pgfqpoint{1.281226in}{2.267174in}}%
\pgfpathlineto{\pgfqpoint{1.280817in}{2.233296in}}%
\pgfpathlineto{\pgfqpoint{1.281231in}{2.266562in}}%
\pgfpathlineto{\pgfqpoint{1.282102in}{2.300915in}}%
\pgfpathlineto{\pgfqpoint{1.281450in}{2.253446in}}%
\pgfpathlineto{\pgfqpoint{1.282360in}{2.274895in}}%
\pgfpathlineto{\pgfqpoint{1.282628in}{2.249375in}}%
\pgfpathlineto{\pgfqpoint{1.283495in}{2.266447in}}%
\pgfpathlineto{\pgfqpoint{1.284079in}{2.248651in}}%
\pgfpathlineto{\pgfqpoint{1.283655in}{2.274382in}}%
\pgfpathlineto{\pgfqpoint{1.284478in}{2.267009in}}%
\pgfpathlineto{\pgfqpoint{1.284887in}{2.293037in}}%
\pgfpathlineto{\pgfqpoint{1.285116in}{2.260404in}}%
\pgfpathlineto{\pgfqpoint{1.285583in}{2.267619in}}%
\pgfpathlineto{\pgfqpoint{1.286606in}{2.251685in}}%
\pgfpathlineto{\pgfqpoint{1.286211in}{2.283587in}}%
\pgfpathlineto{\pgfqpoint{1.286713in}{2.257600in}}%
\pgfpathlineto{\pgfqpoint{1.287706in}{2.223189in}}%
\pgfpathlineto{\pgfqpoint{1.287175in}{2.270084in}}%
\pgfpathlineto{\pgfqpoint{1.287969in}{2.227923in}}%
\pgfpathlineto{\pgfqpoint{1.288242in}{2.234562in}}%
\pgfpathlineto{\pgfqpoint{1.288646in}{2.208760in}}%
\pgfpathlineto{\pgfqpoint{1.289050in}{2.217960in}}%
\pgfpathlineto{\pgfqpoint{1.289113in}{2.225240in}}%
\pgfpathlineto{\pgfqpoint{1.289989in}{2.201661in}}%
\pgfpathlineto{\pgfqpoint{1.290053in}{2.203898in}}%
\pgfpathlineto{\pgfqpoint{1.290403in}{2.182752in}}%
\pgfpathlineto{\pgfqpoint{1.290958in}{2.225990in}}%
\pgfpathlineto{\pgfqpoint{1.291109in}{2.215900in}}%
\pgfpathlineto{\pgfqpoint{1.291484in}{2.227911in}}%
\pgfpathlineto{\pgfqpoint{1.291757in}{2.206054in}}%
\pgfpathlineto{\pgfqpoint{1.292190in}{2.210752in}}%
\pgfpathlineto{\pgfqpoint{1.292341in}{2.192335in}}%
\pgfpathlineto{\pgfqpoint{1.293120in}{2.249698in}}%
\pgfpathlineto{\pgfqpoint{1.293212in}{2.246128in}}%
\pgfpathlineto{\pgfqpoint{1.294259in}{2.272926in}}%
\pgfpathlineto{\pgfqpoint{1.293251in}{2.242176in}}%
\pgfpathlineto{\pgfqpoint{1.294352in}{2.267964in}}%
\pgfpathlineto{\pgfqpoint{1.294619in}{2.251415in}}%
\pgfpathlineto{\pgfqpoint{1.294400in}{2.268341in}}%
\pgfpathlineto{\pgfqpoint{1.294649in}{2.252025in}}%
\pgfpathlineto{\pgfqpoint{1.295671in}{2.243614in}}%
\pgfpathlineto{\pgfqpoint{1.294809in}{2.269418in}}%
\pgfpathlineto{\pgfqpoint{1.295754in}{2.253929in}}%
\pgfpathlineto{\pgfqpoint{1.295788in}{2.248765in}}%
\pgfpathlineto{\pgfqpoint{1.295968in}{2.292799in}}%
\pgfpathlineto{\pgfqpoint{1.296002in}{2.303250in}}%
\pgfpathlineto{\pgfqpoint{1.297019in}{2.263968in}}%
\pgfpathlineto{\pgfqpoint{1.297034in}{2.265763in}}%
\pgfpathlineto{\pgfqpoint{1.297039in}{2.264748in}}%
\pgfpathlineto{\pgfqpoint{1.297682in}{2.307157in}}%
\pgfpathlineto{\pgfqpoint{1.298027in}{2.290257in}}%
\pgfpathlineto{\pgfqpoint{1.298850in}{2.301572in}}%
\pgfpathlineto{\pgfqpoint{1.298349in}{2.268108in}}%
\pgfpathlineto{\pgfqpoint{1.299128in}{2.289713in}}%
\pgfpathlineto{\pgfqpoint{1.300121in}{2.268988in}}%
\pgfpathlineto{\pgfqpoint{1.299536in}{2.295335in}}%
\pgfpathlineto{\pgfqpoint{1.300252in}{2.282007in}}%
\pgfpathlineto{\pgfqpoint{1.300661in}{2.307637in}}%
\pgfpathlineto{\pgfqpoint{1.301323in}{2.266692in}}%
\pgfpathlineto{\pgfqpoint{1.301338in}{2.271511in}}%
\pgfpathlineto{\pgfqpoint{1.301946in}{2.263272in}}%
\pgfpathlineto{\pgfqpoint{1.302287in}{2.281169in}}%
\pgfpathlineto{\pgfqpoint{1.302409in}{2.279488in}}%
\pgfpathlineto{\pgfqpoint{1.302920in}{2.328110in}}%
\pgfpathlineto{\pgfqpoint{1.302448in}{2.272006in}}%
\pgfpathlineto{\pgfqpoint{1.303616in}{2.303590in}}%
\pgfpathlineto{\pgfqpoint{1.304201in}{2.294503in}}%
\pgfpathlineto{\pgfqpoint{1.304653in}{2.328317in}}%
\pgfpathlineto{\pgfqpoint{1.304697in}{2.323217in}}%
\pgfpathlineto{\pgfqpoint{1.305890in}{2.285430in}}%
\pgfpathlineto{\pgfqpoint{1.305476in}{2.344905in}}%
\pgfpathlineto{\pgfqpoint{1.305904in}{2.286270in}}%
\pgfpathlineto{\pgfqpoint{1.306094in}{2.294544in}}%
\pgfpathlineto{\pgfqpoint{1.306825in}{2.246158in}}%
\pgfpathlineto{\pgfqpoint{1.306868in}{2.251451in}}%
\pgfpathlineto{\pgfqpoint{1.307282in}{2.230576in}}%
\pgfpathlineto{\pgfqpoint{1.307764in}{2.272296in}}%
\pgfpathlineto{\pgfqpoint{1.307896in}{2.262760in}}%
\pgfpathlineto{\pgfqpoint{1.307935in}{2.273253in}}%
\pgfpathlineto{\pgfqpoint{1.308538in}{2.192832in}}%
\pgfpathlineto{\pgfqpoint{1.308918in}{2.214887in}}%
\pgfpathlineto{\pgfqpoint{1.309249in}{2.206127in}}%
\pgfpathlineto{\pgfqpoint{1.309575in}{2.236309in}}%
\pgfpathlineto{\pgfqpoint{1.309970in}{2.228989in}}%
\pgfpathlineto{\pgfqpoint{1.310374in}{2.270257in}}%
\pgfpathlineto{\pgfqpoint{1.309984in}{2.227686in}}%
\pgfpathlineto{\pgfqpoint{1.311377in}{2.258421in}}%
\pgfpathlineto{\pgfqpoint{1.312272in}{2.223048in}}%
\pgfpathlineto{\pgfqpoint{1.312526in}{2.241808in}}%
\pgfpathlineto{\pgfqpoint{1.313665in}{2.269710in}}%
\pgfpathlineto{\pgfqpoint{1.313670in}{2.268593in}}%
\pgfpathlineto{\pgfqpoint{1.314566in}{2.253071in}}%
\pgfpathlineto{\pgfqpoint{1.313894in}{2.277804in}}%
\pgfpathlineto{\pgfqpoint{1.314790in}{2.262366in}}%
\pgfpathlineto{\pgfqpoint{1.315831in}{2.195257in}}%
\pgfpathlineto{\pgfqpoint{1.314858in}{2.268287in}}%
\pgfpathlineto{\pgfqpoint{1.316226in}{2.204874in}}%
\pgfpathlineto{\pgfqpoint{1.316644in}{2.216487in}}%
\pgfpathlineto{\pgfqpoint{1.317073in}{2.178970in}}%
\pgfpathlineto{\pgfqpoint{1.317326in}{2.202177in}}%
\pgfpathlineto{\pgfqpoint{1.318490in}{2.179104in}}%
\pgfpathlineto{\pgfqpoint{1.317604in}{2.225853in}}%
\pgfpathlineto{\pgfqpoint{1.318514in}{2.184364in}}%
\pgfpathlineto{\pgfqpoint{1.318787in}{2.202566in}}%
\pgfpathlineto{\pgfqpoint{1.319356in}{2.174378in}}%
\pgfpathlineto{\pgfqpoint{1.319629in}{2.189396in}}%
\pgfpathlineto{\pgfqpoint{1.320413in}{2.177345in}}%
\pgfpathlineto{\pgfqpoint{1.320695in}{2.207427in}}%
\pgfpathlineto{\pgfqpoint{1.321810in}{2.173179in}}%
\pgfpathlineto{\pgfqpoint{1.320880in}{2.216934in}}%
\pgfpathlineto{\pgfqpoint{1.321907in}{2.192275in}}%
\pgfpathlineto{\pgfqpoint{1.322652in}{2.206550in}}%
\pgfpathlineto{\pgfqpoint{1.322287in}{2.169023in}}%
\pgfpathlineto{\pgfqpoint{1.323012in}{2.189511in}}%
\pgfpathlineto{\pgfqpoint{1.323051in}{2.195903in}}%
\pgfpathlineto{\pgfqpoint{1.323173in}{2.179214in}}%
\pgfpathlineto{\pgfqpoint{1.323329in}{2.179992in}}%
\pgfpathlineto{\pgfqpoint{1.323392in}{2.173799in}}%
\pgfpathlineto{\pgfqpoint{1.323986in}{2.211474in}}%
\pgfpathlineto{\pgfqpoint{1.324419in}{2.184780in}}%
\pgfpathlineto{\pgfqpoint{1.324989in}{2.202320in}}%
\pgfpathlineto{\pgfqpoint{1.324760in}{2.178283in}}%
\pgfpathlineto{\pgfqpoint{1.325534in}{2.187663in}}%
\pgfpathlineto{\pgfqpoint{1.326693in}{2.208379in}}%
\pgfpathlineto{\pgfqpoint{1.325700in}{2.182395in}}%
\pgfpathlineto{\pgfqpoint{1.326698in}{2.208357in}}%
\pgfpathlineto{\pgfqpoint{1.326990in}{2.198746in}}%
\pgfpathlineto{\pgfqpoint{1.327379in}{2.224010in}}%
\pgfpathlineto{\pgfqpoint{1.327433in}{2.221671in}}%
\pgfpathlineto{\pgfqpoint{1.327754in}{2.246026in}}%
\pgfpathlineto{\pgfqpoint{1.328412in}{2.215198in}}%
\pgfpathlineto{\pgfqpoint{1.328553in}{2.227395in}}%
\pgfpathlineto{\pgfqpoint{1.328621in}{2.225296in}}%
\pgfpathlineto{\pgfqpoint{1.328689in}{2.238533in}}%
\pgfpathlineto{\pgfqpoint{1.328723in}{2.238382in}}%
\pgfpathlineto{\pgfqpoint{1.329531in}{2.270427in}}%
\pgfpathlineto{\pgfqpoint{1.328752in}{2.237166in}}%
\pgfpathlineto{\pgfqpoint{1.329916in}{2.263146in}}%
\pgfpathlineto{\pgfqpoint{1.330091in}{2.272489in}}%
\pgfpathlineto{\pgfqpoint{1.330442in}{2.251471in}}%
\pgfpathlineto{\pgfqpoint{1.330500in}{2.254493in}}%
\pgfpathlineto{\pgfqpoint{1.331299in}{2.238019in}}%
\pgfpathlineto{\pgfqpoint{1.331576in}{2.274493in}}%
\pgfpathlineto{\pgfqpoint{1.331586in}{2.273919in}}%
\pgfpathlineto{\pgfqpoint{1.332657in}{2.295944in}}%
\pgfpathlineto{\pgfqpoint{1.332029in}{2.262360in}}%
\pgfpathlineto{\pgfqpoint{1.332730in}{2.287237in}}%
\pgfpathlineto{\pgfqpoint{1.332964in}{2.265328in}}%
\pgfpathlineto{\pgfqpoint{1.333640in}{2.314970in}}%
\pgfpathlineto{\pgfqpoint{1.333816in}{2.306205in}}%
\pgfpathlineto{\pgfqpoint{1.333937in}{2.300890in}}%
\pgfpathlineto{\pgfqpoint{1.334775in}{2.332099in}}%
\pgfpathlineto{\pgfqpoint{1.335773in}{2.341863in}}%
\pgfpathlineto{\pgfqpoint{1.335096in}{2.313312in}}%
\pgfpathlineto{\pgfqpoint{1.335880in}{2.331441in}}%
\pgfpathlineto{\pgfqpoint{1.335992in}{2.323309in}}%
\pgfpathlineto{\pgfqpoint{1.336639in}{2.352582in}}%
\pgfpathlineto{\pgfqpoint{1.336664in}{2.348853in}}%
\pgfpathlineto{\pgfqpoint{1.337808in}{2.391276in}}%
\pgfpathlineto{\pgfqpoint{1.336712in}{2.346345in}}%
\pgfpathlineto{\pgfqpoint{1.337827in}{2.387161in}}%
\pgfpathlineto{\pgfqpoint{1.339210in}{2.426616in}}%
\pgfpathlineto{\pgfqpoint{1.337915in}{2.382503in}}%
\pgfpathlineto{\pgfqpoint{1.339215in}{2.423656in}}%
\pgfpathlineto{\pgfqpoint{1.339336in}{2.408093in}}%
\pgfpathlineto{\pgfqpoint{1.340140in}{2.442249in}}%
\pgfpathlineto{\pgfqpoint{1.340305in}{2.437580in}}%
\pgfpathlineto{\pgfqpoint{1.340792in}{2.466149in}}%
\pgfpathlineto{\pgfqpoint{1.340354in}{2.432604in}}%
\pgfpathlineto{\pgfqpoint{1.341445in}{2.455635in}}%
\pgfpathlineto{\pgfqpoint{1.342759in}{2.518722in}}%
\pgfpathlineto{\pgfqpoint{1.342769in}{2.517346in}}%
\pgfpathlineto{\pgfqpoint{1.343927in}{2.459327in}}%
\pgfpathlineto{\pgfqpoint{1.342993in}{2.520250in}}%
\pgfpathlineto{\pgfqpoint{1.343947in}{2.463985in}}%
\pgfpathlineto{\pgfqpoint{1.344385in}{2.474541in}}%
\pgfpathlineto{\pgfqpoint{1.344891in}{2.449478in}}%
\pgfpathlineto{\pgfqpoint{1.345038in}{2.460734in}}%
\pgfpathlineto{\pgfqpoint{1.345782in}{2.406667in}}%
\pgfpathlineto{\pgfqpoint{1.345101in}{2.464189in}}%
\pgfpathlineto{\pgfqpoint{1.346274in}{2.442708in}}%
\pgfpathlineto{\pgfqpoint{1.346985in}{2.455878in}}%
\pgfpathlineto{\pgfqpoint{1.346401in}{2.427291in}}%
\pgfpathlineto{\pgfqpoint{1.347389in}{2.444225in}}%
\pgfpathlineto{\pgfqpoint{1.347759in}{2.456707in}}%
\pgfpathlineto{\pgfqpoint{1.348387in}{2.437385in}}%
\pgfpathlineto{\pgfqpoint{1.348465in}{2.439731in}}%
\pgfpathlineto{\pgfqpoint{1.348835in}{2.414523in}}%
\pgfpathlineto{\pgfqpoint{1.349473in}{2.442833in}}%
\pgfpathlineto{\pgfqpoint{1.349590in}{2.428136in}}%
\pgfpathlineto{\pgfqpoint{1.349950in}{2.456734in}}%
\pgfpathlineto{\pgfqpoint{1.350670in}{2.414316in}}%
\pgfpathlineto{\pgfqpoint{1.350685in}{2.417443in}}%
\pgfpathlineto{\pgfqpoint{1.351182in}{2.393613in}}%
\pgfpathlineto{\pgfqpoint{1.350700in}{2.419321in}}%
\pgfpathlineto{\pgfqpoint{1.351834in}{2.396234in}}%
\pgfpathlineto{\pgfqpoint{1.352043in}{2.417977in}}%
\pgfpathlineto{\pgfqpoint{1.352744in}{2.389328in}}%
\pgfpathlineto{\pgfqpoint{1.352949in}{2.401988in}}%
\pgfpathlineto{\pgfqpoint{1.353256in}{2.382792in}}%
\pgfpathlineto{\pgfqpoint{1.354025in}{2.414031in}}%
\pgfpathlineto{\pgfqpoint{1.354039in}{2.413258in}}%
\pgfpathlineto{\pgfqpoint{1.355115in}{2.423047in}}%
\pgfpathlineto{\pgfqpoint{1.354244in}{2.402968in}}%
\pgfpathlineto{\pgfqpoint{1.355164in}{2.417785in}}%
\pgfpathlineto{\pgfqpoint{1.355232in}{2.424130in}}%
\pgfpathlineto{\pgfqpoint{1.355393in}{2.406051in}}%
\pgfpathlineto{\pgfqpoint{1.356459in}{2.380700in}}%
\pgfpathlineto{\pgfqpoint{1.355787in}{2.430760in}}%
\pgfpathlineto{\pgfqpoint{1.356552in}{2.383664in}}%
\pgfpathlineto{\pgfqpoint{1.356595in}{2.387949in}}%
\pgfpathlineto{\pgfqpoint{1.357316in}{2.340482in}}%
\pgfpathlineto{\pgfqpoint{1.357584in}{2.363907in}}%
\pgfpathlineto{\pgfqpoint{1.358737in}{2.313190in}}%
\pgfpathlineto{\pgfqpoint{1.358757in}{2.316606in}}%
\pgfpathlineto{\pgfqpoint{1.359273in}{2.293665in}}%
\pgfpathlineto{\pgfqpoint{1.359628in}{2.304887in}}%
\pgfpathlineto{\pgfqpoint{1.360437in}{2.293682in}}%
\pgfpathlineto{\pgfqpoint{1.359736in}{2.316503in}}%
\pgfpathlineto{\pgfqpoint{1.360743in}{2.301336in}}%
\pgfpathlineto{\pgfqpoint{1.360748in}{2.300945in}}%
\pgfpathlineto{\pgfqpoint{1.361396in}{2.337216in}}%
\pgfpathlineto{\pgfqpoint{1.361571in}{2.323866in}}%
\pgfpathlineto{\pgfqpoint{1.361576in}{2.325128in}}%
\pgfpathlineto{\pgfqpoint{1.362520in}{2.284744in}}%
\pgfpathlineto{\pgfqpoint{1.362530in}{2.286315in}}%
\pgfpathlineto{\pgfqpoint{1.363085in}{2.271796in}}%
\pgfpathlineto{\pgfqpoint{1.363562in}{2.295652in}}%
\pgfpathlineto{\pgfqpoint{1.363591in}{2.295184in}}%
\pgfpathlineto{\pgfqpoint{1.363995in}{2.302276in}}%
\pgfpathlineto{\pgfqpoint{1.364385in}{2.265498in}}%
\pgfpathlineto{\pgfqpoint{1.364580in}{2.280619in}}%
\pgfpathlineto{\pgfqpoint{1.365738in}{2.245405in}}%
\pgfpathlineto{\pgfqpoint{1.365807in}{2.253685in}}%
\pgfpathlineto{\pgfqpoint{1.366863in}{2.268398in}}%
\pgfpathlineto{\pgfqpoint{1.366469in}{2.240616in}}%
\pgfpathlineto{\pgfqpoint{1.366975in}{2.265160in}}%
\pgfpathlineto{\pgfqpoint{1.367121in}{2.243803in}}%
\pgfpathlineto{\pgfqpoint{1.367905in}{2.287636in}}%
\pgfpathlineto{\pgfqpoint{1.368075in}{2.272259in}}%
\pgfpathlineto{\pgfqpoint{1.368742in}{2.218644in}}%
\pgfpathlineto{\pgfqpoint{1.368178in}{2.272744in}}%
\pgfpathlineto{\pgfqpoint{1.369292in}{2.222445in}}%
\pgfpathlineto{\pgfqpoint{1.369721in}{2.227349in}}%
\pgfpathlineto{\pgfqpoint{1.370023in}{2.203286in}}%
\pgfpathlineto{\pgfqpoint{1.370232in}{2.212467in}}%
\pgfpathlineto{\pgfqpoint{1.370446in}{2.194418in}}%
\pgfpathlineto{\pgfqpoint{1.370719in}{2.231040in}}%
\pgfpathlineto{\pgfqpoint{1.371362in}{2.197652in}}%
\pgfpathlineto{\pgfqpoint{1.371420in}{2.206491in}}%
\pgfpathlineto{\pgfqpoint{1.372335in}{2.163734in}}%
\pgfpathlineto{\pgfqpoint{1.372345in}{2.163993in}}%
\pgfpathlineto{\pgfqpoint{1.372822in}{2.179907in}}%
\pgfpathlineto{\pgfqpoint{1.373022in}{2.156398in}}%
\pgfpathlineto{\pgfqpoint{1.373494in}{2.171019in}}%
\pgfpathlineto{\pgfqpoint{1.374560in}{2.149690in}}%
\pgfpathlineto{\pgfqpoint{1.374273in}{2.173321in}}%
\pgfpathlineto{\pgfqpoint{1.374648in}{2.152436in}}%
\pgfpathlineto{\pgfqpoint{1.374716in}{2.150049in}}%
\pgfpathlineto{\pgfqpoint{1.375154in}{2.167834in}}%
\pgfpathlineto{\pgfqpoint{1.375451in}{2.178232in}}%
\pgfpathlineto{\pgfqpoint{1.376045in}{2.152077in}}%
\pgfpathlineto{\pgfqpoint{1.376240in}{2.161149in}}%
\pgfpathlineto{\pgfqpoint{1.376469in}{2.152950in}}%
\pgfpathlineto{\pgfqpoint{1.376761in}{2.196947in}}%
\pgfpathlineto{\pgfqpoint{1.377097in}{2.186217in}}%
\pgfpathlineto{\pgfqpoint{1.377228in}{2.210761in}}%
\pgfpathlineto{\pgfqpoint{1.377564in}{2.173859in}}%
\pgfpathlineto{\pgfqpoint{1.378207in}{2.187431in}}%
\pgfpathlineto{\pgfqpoint{1.378674in}{2.193335in}}%
\pgfpathlineto{\pgfqpoint{1.378387in}{2.169214in}}%
\pgfpathlineto{\pgfqpoint{1.378952in}{2.175232in}}%
\pgfpathlineto{\pgfqpoint{1.379872in}{2.143583in}}%
\pgfpathlineto{\pgfqpoint{1.379511in}{2.178460in}}%
\pgfpathlineto{\pgfqpoint{1.380100in}{2.148080in}}%
\pgfpathlineto{\pgfqpoint{1.380310in}{2.169463in}}%
\pgfpathlineto{\pgfqpoint{1.380850in}{2.131872in}}%
\pgfpathlineto{\pgfqpoint{1.381162in}{2.137702in}}%
\pgfpathlineto{\pgfqpoint{1.381347in}{2.129047in}}%
\pgfpathlineto{\pgfqpoint{1.381887in}{2.164765in}}%
\pgfpathlineto{\pgfqpoint{1.382267in}{2.137803in}}%
\pgfpathlineto{\pgfqpoint{1.382291in}{2.143421in}}%
\pgfpathlineto{\pgfqpoint{1.383109in}{2.107172in}}%
\pgfpathlineto{\pgfqpoint{1.383270in}{2.112429in}}%
\pgfpathlineto{\pgfqpoint{1.383859in}{2.093620in}}%
\pgfpathlineto{\pgfqpoint{1.384020in}{2.117216in}}%
\pgfpathlineto{\pgfqpoint{1.384341in}{2.115616in}}%
\pgfpathlineto{\pgfqpoint{1.384536in}{2.132522in}}%
\pgfpathlineto{\pgfqpoint{1.385057in}{2.085640in}}%
\pgfpathlineto{\pgfqpoint{1.385407in}{2.099221in}}%
\pgfpathlineto{\pgfqpoint{1.385879in}{2.068176in}}%
\pgfpathlineto{\pgfqpoint{1.386342in}{2.104751in}}%
\pgfpathlineto{\pgfqpoint{1.386527in}{2.095492in}}%
\pgfpathlineto{\pgfqpoint{1.386615in}{2.101424in}}%
\pgfpathlineto{\pgfqpoint{1.386804in}{2.071531in}}%
\pgfpathlineto{\pgfqpoint{1.387617in}{2.089781in}}%
\pgfpathlineto{\pgfqpoint{1.387832in}{2.080362in}}%
\pgfpathlineto{\pgfqpoint{1.388611in}{2.104589in}}%
\pgfpathlineto{\pgfqpoint{1.388713in}{2.093749in}}%
\pgfpathlineto{\pgfqpoint{1.389828in}{2.114658in}}%
\pgfpathlineto{\pgfqpoint{1.389317in}{2.084621in}}%
\pgfpathlineto{\pgfqpoint{1.389862in}{2.108618in}}%
\pgfpathlineto{\pgfqpoint{1.390855in}{2.086039in}}%
\pgfpathlineto{\pgfqpoint{1.389998in}{2.120124in}}%
\pgfpathlineto{\pgfqpoint{1.391006in}{2.090579in}}%
\pgfpathlineto{\pgfqpoint{1.391478in}{2.070303in}}%
\pgfpathlineto{\pgfqpoint{1.392019in}{2.103743in}}%
\pgfpathlineto{\pgfqpoint{1.392345in}{2.130909in}}%
\pgfpathlineto{\pgfqpoint{1.393153in}{2.112591in}}%
\pgfpathlineto{\pgfqpoint{1.393489in}{2.107914in}}%
\pgfpathlineto{\pgfqpoint{1.393211in}{2.129392in}}%
\pgfpathlineto{\pgfqpoint{1.393830in}{2.127008in}}%
\pgfpathlineto{\pgfqpoint{1.394127in}{2.145266in}}%
\pgfpathlineto{\pgfqpoint{1.394594in}{2.096547in}}%
\pgfpathlineto{\pgfqpoint{1.394891in}{2.109661in}}%
\pgfpathlineto{\pgfqpoint{1.394901in}{2.110564in}}%
\pgfpathlineto{\pgfqpoint{1.395105in}{2.082559in}}%
\pgfpathlineto{\pgfqpoint{1.395222in}{2.085192in}}%
\pgfpathlineto{\pgfqpoint{1.395246in}{2.076582in}}%
\pgfpathlineto{\pgfqpoint{1.396147in}{2.108180in}}%
\pgfpathlineto{\pgfqpoint{1.396293in}{2.102258in}}%
\pgfpathlineto{\pgfqpoint{1.396303in}{2.108618in}}%
\pgfpathlineto{\pgfqpoint{1.396887in}{2.074338in}}%
\pgfpathlineto{\pgfqpoint{1.397389in}{2.096587in}}%
\pgfpathlineto{\pgfqpoint{1.398070in}{2.055056in}}%
\pgfpathlineto{\pgfqpoint{1.398562in}{2.076581in}}%
\pgfpathlineto{\pgfqpoint{1.399813in}{2.114086in}}%
\pgfpathlineto{\pgfqpoint{1.399034in}{2.041049in}}%
\pgfpathlineto{\pgfqpoint{1.399828in}{2.108329in}}%
\pgfpathlineto{\pgfqpoint{1.400232in}{2.116053in}}%
\pgfpathlineto{\pgfqpoint{1.401059in}{2.081132in}}%
\pgfpathlineto{\pgfqpoint{1.402023in}{2.116301in}}%
\pgfpathlineto{\pgfqpoint{1.401541in}{2.076954in}}%
\pgfpathlineto{\pgfqpoint{1.402238in}{2.111063in}}%
\pgfpathlineto{\pgfqpoint{1.402272in}{2.100096in}}%
\pgfpathlineto{\pgfqpoint{1.403250in}{2.131021in}}%
\pgfpathlineto{\pgfqpoint{1.403318in}{2.126102in}}%
\pgfpathlineto{\pgfqpoint{1.403406in}{2.133593in}}%
\pgfpathlineto{\pgfqpoint{1.404170in}{2.069636in}}%
\pgfpathlineto{\pgfqpoint{1.404234in}{2.075087in}}%
\pgfpathlineto{\pgfqpoint{1.404321in}{2.068596in}}%
\pgfpathlineto{\pgfqpoint{1.404706in}{2.092683in}}%
\pgfpathlineto{\pgfqpoint{1.404828in}{2.089231in}}%
\pgfpathlineto{\pgfqpoint{1.405572in}{2.130505in}}%
\pgfpathlineto{\pgfqpoint{1.404871in}{2.081370in}}%
\pgfpathlineto{\pgfqpoint{1.406074in}{2.124111in}}%
\pgfpathlineto{\pgfqpoint{1.407004in}{2.055742in}}%
\pgfpathlineto{\pgfqpoint{1.406162in}{2.125097in}}%
\pgfpathlineto{\pgfqpoint{1.407228in}{2.080982in}}%
\pgfpathlineto{\pgfqpoint{1.407851in}{2.076250in}}%
\pgfpathlineto{\pgfqpoint{1.407486in}{2.096660in}}%
\pgfpathlineto{\pgfqpoint{1.408245in}{2.086001in}}%
\pgfpathlineto{\pgfqpoint{1.408854in}{2.130568in}}%
\pgfpathlineto{\pgfqpoint{1.409389in}{2.119663in}}%
\pgfpathlineto{\pgfqpoint{1.409414in}{2.116009in}}%
\pgfpathlineto{\pgfqpoint{1.409910in}{2.153017in}}%
\pgfpathlineto{\pgfqpoint{1.410368in}{2.127289in}}%
\pgfpathlineto{\pgfqpoint{1.410704in}{2.145225in}}%
\pgfpathlineto{\pgfqpoint{1.411337in}{2.106342in}}%
\pgfpathlineto{\pgfqpoint{1.411385in}{2.107604in}}%
\pgfpathlineto{\pgfqpoint{1.411692in}{2.087521in}}%
\pgfpathlineto{\pgfqpoint{1.412447in}{2.113554in}}%
\pgfpathlineto{\pgfqpoint{1.412452in}{2.110781in}}%
\pgfpathlineto{\pgfqpoint{1.413333in}{2.132818in}}%
\pgfpathlineto{\pgfqpoint{1.412588in}{2.098484in}}%
\pgfpathlineto{\pgfqpoint{1.413567in}{2.113766in}}%
\pgfpathlineto{\pgfqpoint{1.413781in}{2.124142in}}%
\pgfpathlineto{\pgfqpoint{1.414214in}{2.072175in}}%
\pgfpathlineto{\pgfqpoint{1.414604in}{2.091448in}}%
\pgfpathlineto{\pgfqpoint{1.414662in}{2.085518in}}%
\pgfpathlineto{\pgfqpoint{1.415139in}{2.125148in}}%
\pgfpathlineto{\pgfqpoint{1.415655in}{2.099737in}}%
\pgfpathlineto{\pgfqpoint{1.416161in}{2.087614in}}%
\pgfpathlineto{\pgfqpoint{1.416785in}{2.111407in}}%
\pgfpathlineto{\pgfqpoint{1.416824in}{2.108959in}}%
\pgfpathlineto{\pgfqpoint{1.417505in}{2.162900in}}%
\pgfpathlineto{\pgfqpoint{1.417744in}{2.147980in}}%
\pgfpathlineto{\pgfqpoint{1.418756in}{2.170074in}}%
\pgfpathlineto{\pgfqpoint{1.418109in}{2.137843in}}%
\pgfpathlineto{\pgfqpoint{1.418868in}{2.154429in}}%
\pgfpathlineto{\pgfqpoint{1.419175in}{2.146406in}}%
\pgfpathlineto{\pgfqpoint{1.419886in}{2.173842in}}%
\pgfpathlineto{\pgfqpoint{1.419974in}{2.155651in}}%
\pgfpathlineto{\pgfqpoint{1.420227in}{2.169914in}}%
\pgfpathlineto{\pgfqpoint{1.420787in}{2.142195in}}%
\pgfpathlineto{\pgfqpoint{1.420864in}{2.144507in}}%
\pgfpathlineto{\pgfqpoint{1.421166in}{2.129988in}}%
\pgfpathlineto{\pgfqpoint{1.421882in}{2.155039in}}%
\pgfpathlineto{\pgfqpoint{1.421940in}{2.153308in}}%
\pgfpathlineto{\pgfqpoint{1.423080in}{2.170537in}}%
\pgfpathlineto{\pgfqpoint{1.422529in}{2.129419in}}%
\pgfpathlineto{\pgfqpoint{1.423089in}{2.170423in}}%
\pgfpathlineto{\pgfqpoint{1.423119in}{2.172141in}}%
\pgfpathlineto{\pgfqpoint{1.423644in}{2.142374in}}%
\pgfpathlineto{\pgfqpoint{1.424180in}{2.168714in}}%
\pgfpathlineto{\pgfqpoint{1.424292in}{2.161051in}}%
\pgfpathlineto{\pgfqpoint{1.424920in}{2.188418in}}%
\pgfpathlineto{\pgfqpoint{1.425295in}{2.166209in}}%
\pgfpathlineto{\pgfqpoint{1.425300in}{2.167414in}}%
\pgfpathlineto{\pgfqpoint{1.425840in}{2.137399in}}%
\pgfpathlineto{\pgfqpoint{1.426317in}{2.144890in}}%
\pgfpathlineto{\pgfqpoint{1.426819in}{2.136537in}}%
\pgfpathlineto{\pgfqpoint{1.427262in}{2.162852in}}%
\pgfpathlineto{\pgfqpoint{1.427422in}{2.149332in}}%
\pgfpathlineto{\pgfqpoint{1.428172in}{2.131146in}}%
\pgfpathlineto{\pgfqpoint{1.427948in}{2.173940in}}%
\pgfpathlineto{\pgfqpoint{1.428581in}{2.135494in}}%
\pgfpathlineto{\pgfqpoint{1.429501in}{2.172693in}}%
\pgfpathlineto{\pgfqpoint{1.428771in}{2.122982in}}%
\pgfpathlineto{\pgfqpoint{1.429745in}{2.162729in}}%
\pgfpathlineto{\pgfqpoint{1.430217in}{2.135004in}}%
\pgfpathlineto{\pgfqpoint{1.430606in}{2.170654in}}%
\pgfpathlineto{\pgfqpoint{1.430874in}{2.153458in}}%
\pgfpathlineto{\pgfqpoint{1.430957in}{2.151610in}}%
\pgfpathlineto{\pgfqpoint{1.431083in}{2.173221in}}%
\pgfpathlineto{\pgfqpoint{1.431191in}{2.171261in}}%
\pgfpathlineto{\pgfqpoint{1.431736in}{2.193423in}}%
\pgfpathlineto{\pgfqpoint{1.432164in}{2.159338in}}%
\pgfpathlineto{\pgfqpoint{1.432276in}{2.167923in}}%
\pgfpathlineto{\pgfqpoint{1.432330in}{2.157230in}}%
\pgfpathlineto{\pgfqpoint{1.432690in}{2.185583in}}%
\pgfpathlineto{\pgfqpoint{1.433386in}{2.168321in}}%
\pgfpathlineto{\pgfqpoint{1.433469in}{2.159289in}}%
\pgfpathlineto{\pgfqpoint{1.434238in}{2.194255in}}%
\pgfpathlineto{\pgfqpoint{1.434277in}{2.192906in}}%
\pgfpathlineto{\pgfqpoint{1.434457in}{2.201799in}}%
\pgfpathlineto{\pgfqpoint{1.435080in}{2.142957in}}%
\pgfpathlineto{\pgfqpoint{1.436161in}{2.077709in}}%
\pgfpathlineto{\pgfqpoint{1.436337in}{2.082597in}}%
\pgfpathlineto{\pgfqpoint{1.436385in}{2.077397in}}%
\pgfpathlineto{\pgfqpoint{1.437510in}{2.113913in}}%
\pgfpathlineto{\pgfqpoint{1.437929in}{2.095598in}}%
\pgfpathlineto{\pgfqpoint{1.438449in}{2.129297in}}%
\pgfpathlineto{\pgfqpoint{1.438625in}{2.110209in}}%
\pgfpathlineto{\pgfqpoint{1.439107in}{2.134777in}}%
\pgfpathlineto{\pgfqpoint{1.439598in}{2.104400in}}%
\pgfpathlineto{\pgfqpoint{1.439725in}{2.105781in}}%
\pgfpathlineto{\pgfqpoint{1.440431in}{2.087677in}}%
\pgfpathlineto{\pgfqpoint{1.439978in}{2.119030in}}%
\pgfpathlineto{\pgfqpoint{1.440825in}{2.108483in}}%
\pgfpathlineto{\pgfqpoint{1.440932in}{2.113344in}}%
\pgfpathlineto{\pgfqpoint{1.441429in}{2.084358in}}%
\pgfpathlineto{\pgfqpoint{1.441716in}{2.096302in}}%
\pgfpathlineto{\pgfqpoint{1.442013in}{2.076283in}}%
\pgfpathlineto{\pgfqpoint{1.442612in}{2.134131in}}%
\pgfpathlineto{\pgfqpoint{1.442753in}{2.121917in}}%
\pgfpathlineto{\pgfqpoint{1.443372in}{2.144427in}}%
\pgfpathlineto{\pgfqpoint{1.443206in}{2.121286in}}%
\pgfpathlineto{\pgfqpoint{1.443888in}{2.141122in}}%
\pgfpathlineto{\pgfqpoint{1.443951in}{2.135689in}}%
\pgfpathlineto{\pgfqpoint{1.444384in}{2.167419in}}%
\pgfpathlineto{\pgfqpoint{1.444954in}{2.163565in}}%
\pgfpathlineto{\pgfqpoint{1.444973in}{2.174537in}}%
\pgfpathlineto{\pgfqpoint{1.446020in}{2.134044in}}%
\pgfpathlineto{\pgfqpoint{1.446030in}{2.132502in}}%
\pgfpathlineto{\pgfqpoint{1.446741in}{2.179617in}}%
\pgfpathlineto{\pgfqpoint{1.447096in}{2.142718in}}%
\pgfpathlineto{\pgfqpoint{1.447208in}{2.151007in}}%
\pgfpathlineto{\pgfqpoint{1.447836in}{2.120342in}}%
\pgfpathlineto{\pgfqpoint{1.448162in}{2.135214in}}%
\pgfpathlineto{\pgfqpoint{1.448508in}{2.109827in}}%
\pgfpathlineto{\pgfqpoint{1.449223in}{2.146948in}}%
\pgfpathlineto{\pgfqpoint{1.449253in}{2.140175in}}%
\pgfpathlineto{\pgfqpoint{1.449272in}{2.145790in}}%
\pgfpathlineto{\pgfqpoint{1.449978in}{2.095868in}}%
\pgfpathlineto{\pgfqpoint{1.450197in}{2.116381in}}%
\pgfpathlineto{\pgfqpoint{1.450689in}{2.087382in}}%
\pgfpathlineto{\pgfqpoint{1.451127in}{2.117908in}}%
\pgfpathlineto{\pgfqpoint{1.451327in}{2.105348in}}%
\pgfpathlineto{\pgfqpoint{1.452549in}{2.143114in}}%
\pgfpathlineto{\pgfqpoint{1.451341in}{2.099674in}}%
\pgfpathlineto{\pgfqpoint{1.452554in}{2.140691in}}%
\pgfpathlineto{\pgfqpoint{1.452558in}{2.140050in}}%
\pgfpathlineto{\pgfqpoint{1.453177in}{2.173475in}}%
\pgfpathlineto{\pgfqpoint{1.453517in}{2.159154in}}%
\pgfpathlineto{\pgfqpoint{1.454228in}{2.177691in}}%
\pgfpathlineto{\pgfqpoint{1.453941in}{2.147555in}}%
\pgfpathlineto{\pgfqpoint{1.454666in}{2.170389in}}%
\pgfpathlineto{\pgfqpoint{1.455163in}{2.142269in}}%
\pgfpathlineto{\pgfqpoint{1.454842in}{2.174637in}}%
\pgfpathlineto{\pgfqpoint{1.455830in}{2.163556in}}%
\pgfpathlineto{\pgfqpoint{1.456760in}{2.213399in}}%
\pgfpathlineto{\pgfqpoint{1.457023in}{2.201031in}}%
\pgfpathlineto{\pgfqpoint{1.457665in}{2.165991in}}%
\pgfpathlineto{\pgfqpoint{1.457115in}{2.206387in}}%
\pgfpathlineto{\pgfqpoint{1.458162in}{2.179033in}}%
\pgfpathlineto{\pgfqpoint{1.458381in}{2.200759in}}%
\pgfpathlineto{\pgfqpoint{1.459141in}{2.130677in}}%
\pgfpathlineto{\pgfqpoint{1.460358in}{2.073636in}}%
\pgfpathlineto{\pgfqpoint{1.461682in}{2.157769in}}%
\pgfpathlineto{\pgfqpoint{1.461794in}{2.151552in}}%
\pgfpathlineto{\pgfqpoint{1.461886in}{2.141392in}}%
\pgfpathlineto{\pgfqpoint{1.462846in}{2.183402in}}%
\pgfpathlineto{\pgfqpoint{1.462855in}{2.179316in}}%
\pgfpathlineto{\pgfqpoint{1.462909in}{2.184417in}}%
\pgfpathlineto{\pgfqpoint{1.462923in}{2.173625in}}%
\pgfpathlineto{\pgfqpoint{1.462928in}{2.173651in}}%
\pgfpathlineto{\pgfqpoint{1.463459in}{2.136835in}}%
\pgfpathlineto{\pgfqpoint{1.463094in}{2.174121in}}%
\pgfpathlineto{\pgfqpoint{1.464131in}{2.142167in}}%
\pgfpathlineto{\pgfqpoint{1.465226in}{2.162584in}}%
\pgfpathlineto{\pgfqpoint{1.464749in}{2.129651in}}%
\pgfpathlineto{\pgfqpoint{1.465324in}{2.156183in}}%
\pgfpathlineto{\pgfqpoint{1.466127in}{2.115125in}}%
\pgfpathlineto{\pgfqpoint{1.465358in}{2.163385in}}%
\pgfpathlineto{\pgfqpoint{1.466468in}{2.146110in}}%
\pgfpathlineto{\pgfqpoint{1.466993in}{2.155276in}}%
\pgfpathlineto{\pgfqpoint{1.467183in}{2.124114in}}%
\pgfpathlineto{\pgfqpoint{1.467242in}{2.126897in}}%
\pgfpathlineto{\pgfqpoint{1.468250in}{2.100988in}}%
\pgfpathlineto{\pgfqpoint{1.467310in}{2.129037in}}%
\pgfpathlineto{\pgfqpoint{1.468376in}{2.118379in}}%
\pgfpathlineto{\pgfqpoint{1.469136in}{2.099618in}}%
\pgfpathlineto{\pgfqpoint{1.468610in}{2.124227in}}%
\pgfpathlineto{\pgfqpoint{1.469511in}{2.111725in}}%
\pgfpathlineto{\pgfqpoint{1.470562in}{2.080081in}}%
\pgfpathlineto{\pgfqpoint{1.469822in}{2.128874in}}%
\pgfpathlineto{\pgfqpoint{1.470762in}{2.088666in}}%
\pgfpathlineto{\pgfqpoint{1.471117in}{2.099034in}}%
\pgfpathlineto{\pgfqpoint{1.471341in}{2.072279in}}%
\pgfpathlineto{\pgfqpoint{1.471862in}{2.088838in}}%
\pgfpathlineto{\pgfqpoint{1.472709in}{2.063201in}}%
\pgfpathlineto{\pgfqpoint{1.471969in}{2.094685in}}%
\pgfpathlineto{\pgfqpoint{1.472977in}{2.080756in}}%
\pgfpathlineto{\pgfqpoint{1.473055in}{2.084933in}}%
\pgfpathlineto{\pgfqpoint{1.473113in}{2.076305in}}%
\pgfpathlineto{\pgfqpoint{1.473435in}{2.054849in}}%
\pgfpathlineto{\pgfqpoint{1.473785in}{2.083900in}}%
\pgfpathlineto{\pgfqpoint{1.474248in}{2.063517in}}%
\pgfpathlineto{\pgfqpoint{1.475304in}{2.085870in}}%
\pgfpathlineto{\pgfqpoint{1.474496in}{2.047736in}}%
\pgfpathlineto{\pgfqpoint{1.475416in}{2.079823in}}%
\pgfpathlineto{\pgfqpoint{1.475640in}{2.062641in}}%
\pgfpathlineto{\pgfqpoint{1.475932in}{2.092061in}}%
\pgfpathlineto{\pgfqpoint{1.476662in}{2.074053in}}%
\pgfpathlineto{\pgfqpoint{1.476706in}{2.082318in}}%
\pgfpathlineto{\pgfqpoint{1.477514in}{2.049227in}}%
\pgfpathlineto{\pgfqpoint{1.477738in}{2.056821in}}%
\pgfpathlineto{\pgfqpoint{1.477767in}{2.054269in}}%
\pgfpathlineto{\pgfqpoint{1.478527in}{2.080467in}}%
\pgfpathlineto{\pgfqpoint{1.479637in}{2.086925in}}%
\pgfpathlineto{\pgfqpoint{1.479145in}{2.052106in}}%
\pgfpathlineto{\pgfqpoint{1.479642in}{2.085060in}}%
\pgfpathlineto{\pgfqpoint{1.480844in}{2.014000in}}%
\pgfpathlineto{\pgfqpoint{1.480849in}{2.016366in}}%
\pgfpathlineto{\pgfqpoint{1.480937in}{2.035746in}}%
\pgfpathlineto{\pgfqpoint{1.481151in}{2.048344in}}%
\pgfpathlineto{\pgfqpoint{1.481580in}{1.995321in}}%
\pgfpathlineto{\pgfqpoint{1.481993in}{2.014711in}}%
\pgfpathlineto{\pgfqpoint{1.483074in}{1.958854in}}%
\pgfpathlineto{\pgfqpoint{1.483425in}{1.976003in}}%
\pgfpathlineto{\pgfqpoint{1.484267in}{2.013912in}}%
\pgfpathlineto{\pgfqpoint{1.483537in}{1.975283in}}%
\pgfpathlineto{\pgfqpoint{1.484627in}{2.010811in}}%
\pgfpathlineto{\pgfqpoint{1.484632in}{2.008524in}}%
\pgfpathlineto{\pgfqpoint{1.484983in}{2.057428in}}%
\pgfpathlineto{\pgfqpoint{1.485669in}{2.038166in}}%
\pgfpathlineto{\pgfqpoint{1.486808in}{2.066019in}}%
\pgfpathlineto{\pgfqpoint{1.486078in}{2.026734in}}%
\pgfpathlineto{\pgfqpoint{1.486813in}{2.065220in}}%
\pgfpathlineto{\pgfqpoint{1.487641in}{2.021322in}}%
\pgfpathlineto{\pgfqpoint{1.486847in}{2.067106in}}%
\pgfpathlineto{\pgfqpoint{1.487996in}{2.040367in}}%
\pgfpathlineto{\pgfqpoint{1.488050in}{2.051375in}}%
\pgfpathlineto{\pgfqpoint{1.488561in}{2.032834in}}%
\pgfpathlineto{\pgfqpoint{1.489096in}{2.035789in}}%
\pgfpathlineto{\pgfqpoint{1.489165in}{2.032348in}}%
\pgfpathlineto{\pgfqpoint{1.489910in}{2.088435in}}%
\pgfpathlineto{\pgfqpoint{1.490060in}{2.077318in}}%
\pgfpathlineto{\pgfqpoint{1.491244in}{2.119519in}}%
\pgfpathlineto{\pgfqpoint{1.490664in}{2.067588in}}%
\pgfpathlineto{\pgfqpoint{1.491317in}{2.113636in}}%
\pgfpathlineto{\pgfqpoint{1.491949in}{2.080510in}}%
\pgfpathlineto{\pgfqpoint{1.491346in}{2.114100in}}%
\pgfpathlineto{\pgfqpoint{1.492465in}{2.089787in}}%
\pgfpathlineto{\pgfqpoint{1.493483in}{2.112708in}}%
\pgfpathlineto{\pgfqpoint{1.492529in}{2.084794in}}%
\pgfpathlineto{\pgfqpoint{1.493717in}{2.093154in}}%
\pgfpathlineto{\pgfqpoint{1.493761in}{2.088462in}}%
\pgfpathlineto{\pgfqpoint{1.494101in}{2.125705in}}%
\pgfpathlineto{\pgfqpoint{1.494798in}{2.103184in}}%
\pgfpathlineto{\pgfqpoint{1.495630in}{2.068290in}}%
\pgfpathlineto{\pgfqpoint{1.494832in}{2.105688in}}%
\pgfpathlineto{\pgfqpoint{1.496180in}{2.083984in}}%
\pgfpathlineto{\pgfqpoint{1.497008in}{2.101068in}}%
\pgfpathlineto{\pgfqpoint{1.496516in}{2.058748in}}%
\pgfpathlineto{\pgfqpoint{1.497310in}{2.094716in}}%
\pgfpathlineto{\pgfqpoint{1.497543in}{2.085308in}}%
\pgfpathlineto{\pgfqpoint{1.498347in}{2.118292in}}%
\pgfpathlineto{\pgfqpoint{1.498376in}{2.113726in}}%
\pgfpathlineto{\pgfqpoint{1.499247in}{2.136498in}}%
\pgfpathlineto{\pgfqpoint{1.498799in}{2.103395in}}%
\pgfpathlineto{\pgfqpoint{1.499437in}{2.104803in}}%
\pgfpathlineto{\pgfqpoint{1.500499in}{2.057443in}}%
\pgfpathlineto{\pgfqpoint{1.499525in}{2.105559in}}%
\pgfpathlineto{\pgfqpoint{1.500625in}{2.073928in}}%
\pgfpathlineto{\pgfqpoint{1.500679in}{2.083326in}}%
\pgfpathlineto{\pgfqpoint{1.501594in}{2.034713in}}%
\pgfpathlineto{\pgfqpoint{1.501609in}{2.037906in}}%
\pgfpathlineto{\pgfqpoint{1.502456in}{2.004586in}}%
\pgfpathlineto{\pgfqpoint{1.501633in}{2.041172in}}%
\pgfpathlineto{\pgfqpoint{1.502767in}{2.020851in}}%
\pgfpathlineto{\pgfqpoint{1.503244in}{2.027194in}}%
\pgfpathlineto{\pgfqpoint{1.503536in}{2.002013in}}%
\pgfpathlineto{\pgfqpoint{1.503683in}{2.012598in}}%
\pgfpathlineto{\pgfqpoint{1.504651in}{1.934174in}}%
\pgfpathlineto{\pgfqpoint{1.504880in}{1.962353in}}%
\pgfpathlineto{\pgfqpoint{1.505718in}{1.983568in}}%
\pgfpathlineto{\pgfqpoint{1.505279in}{1.934095in}}%
\pgfpathlineto{\pgfqpoint{1.506078in}{1.974170in}}%
\pgfpathlineto{\pgfqpoint{1.506998in}{1.942489in}}%
\pgfpathlineto{\pgfqpoint{1.506540in}{1.986521in}}%
\pgfpathlineto{\pgfqpoint{1.507222in}{1.964411in}}%
\pgfpathlineto{\pgfqpoint{1.507767in}{1.986957in}}%
\pgfpathlineto{\pgfqpoint{1.508006in}{1.961584in}}%
\pgfpathlineto{\pgfqpoint{1.508351in}{1.971389in}}%
\pgfpathlineto{\pgfqpoint{1.509048in}{1.931351in}}%
\pgfpathlineto{\pgfqpoint{1.508424in}{1.981504in}}%
\pgfpathlineto{\pgfqpoint{1.509588in}{1.952103in}}%
\pgfpathlineto{\pgfqpoint{1.510859in}{2.004922in}}%
\pgfpathlineto{\pgfqpoint{1.509603in}{1.951952in}}%
\pgfpathlineto{\pgfqpoint{1.510864in}{2.002394in}}%
\pgfpathlineto{\pgfqpoint{1.512003in}{1.978159in}}%
\pgfpathlineto{\pgfqpoint{1.511039in}{2.006706in}}%
\pgfpathlineto{\pgfqpoint{1.512125in}{1.993141in}}%
\pgfpathlineto{\pgfqpoint{1.512154in}{1.996475in}}%
\pgfpathlineto{\pgfqpoint{1.512728in}{1.934090in}}%
\pgfpathlineto{\pgfqpoint{1.513103in}{1.944390in}}%
\pgfpathlineto{\pgfqpoint{1.513955in}{1.952402in}}%
\pgfpathlineto{\pgfqpoint{1.513366in}{1.930757in}}%
\pgfpathlineto{\pgfqpoint{1.514155in}{1.938648in}}%
\pgfpathlineto{\pgfqpoint{1.515396in}{1.879837in}}%
\pgfpathlineto{\pgfqpoint{1.515411in}{1.880411in}}%
\pgfpathlineto{\pgfqpoint{1.515484in}{1.883932in}}%
\pgfpathlineto{\pgfqpoint{1.516404in}{1.855837in}}%
\pgfpathlineto{\pgfqpoint{1.516501in}{1.870315in}}%
\pgfpathlineto{\pgfqpoint{1.516506in}{1.870202in}}%
\pgfpathlineto{\pgfqpoint{1.516613in}{1.892869in}}%
\pgfpathlineto{\pgfqpoint{1.516618in}{1.892736in}}%
\pgfpathlineto{\pgfqpoint{1.516750in}{1.900093in}}%
\pgfpathlineto{\pgfqpoint{1.517309in}{1.863827in}}%
\pgfpathlineto{\pgfqpoint{1.517689in}{1.871758in}}%
\pgfpathlineto{\pgfqpoint{1.519213in}{1.937356in}}%
\pgfpathlineto{\pgfqpoint{1.517709in}{1.868164in}}%
\pgfpathlineto{\pgfqpoint{1.519218in}{1.937204in}}%
\pgfpathlineto{\pgfqpoint{1.519223in}{1.936968in}}%
\pgfpathlineto{\pgfqpoint{1.519622in}{1.968363in}}%
\pgfpathlineto{\pgfqpoint{1.519661in}{1.967100in}}%
\pgfpathlineto{\pgfqpoint{1.520440in}{2.003564in}}%
\pgfpathlineto{\pgfqpoint{1.519831in}{1.966249in}}%
\pgfpathlineto{\pgfqpoint{1.520825in}{1.992765in}}%
\pgfpathlineto{\pgfqpoint{1.521399in}{1.964278in}}%
\pgfpathlineto{\pgfqpoint{1.521959in}{1.978630in}}%
\pgfpathlineto{\pgfqpoint{1.522889in}{2.035038in}}%
\pgfpathlineto{\pgfqpoint{1.523147in}{2.009331in}}%
\pgfpathlineto{\pgfqpoint{1.523366in}{2.002093in}}%
\pgfpathlineto{\pgfqpoint{1.523502in}{2.020624in}}%
\pgfpathlineto{\pgfqpoint{1.524218in}{2.014469in}}%
\pgfpathlineto{\pgfqpoint{1.524413in}{2.026775in}}%
\pgfpathlineto{\pgfqpoint{1.524997in}{1.995930in}}%
\pgfpathlineto{\pgfqpoint{1.525308in}{2.011095in}}%
\pgfpathlineto{\pgfqpoint{1.525377in}{2.006210in}}%
\pgfpathlineto{\pgfqpoint{1.526292in}{2.093951in}}%
\pgfpathlineto{\pgfqpoint{1.526384in}{2.102457in}}%
\pgfpathlineto{\pgfqpoint{1.527037in}{2.068354in}}%
\pgfpathlineto{\pgfqpoint{1.527085in}{2.070479in}}%
\pgfpathlineto{\pgfqpoint{1.528054in}{2.044114in}}%
\pgfpathlineto{\pgfqpoint{1.527572in}{2.085252in}}%
\pgfpathlineto{\pgfqpoint{1.528230in}{2.048426in}}%
\pgfpathlineto{\pgfqpoint{1.528531in}{2.054741in}}%
\pgfpathlineto{\pgfqpoint{1.528843in}{2.031116in}}%
\pgfpathlineto{\pgfqpoint{1.529325in}{2.041468in}}%
\pgfpathlineto{\pgfqpoint{1.530245in}{2.015117in}}%
\pgfpathlineto{\pgfqpoint{1.529378in}{2.043136in}}%
\pgfpathlineto{\pgfqpoint{1.530498in}{2.022242in}}%
\pgfpathlineto{\pgfqpoint{1.530523in}{2.024316in}}%
\pgfpathlineto{\pgfqpoint{1.531399in}{1.982316in}}%
\pgfpathlineto{\pgfqpoint{1.532645in}{1.955789in}}%
\pgfpathlineto{\pgfqpoint{1.531491in}{1.992474in}}%
\pgfpathlineto{\pgfqpoint{1.532665in}{1.958512in}}%
\pgfpathlineto{\pgfqpoint{1.533799in}{2.000844in}}%
\pgfpathlineto{\pgfqpoint{1.533828in}{1.995093in}}%
\pgfpathlineto{\pgfqpoint{1.533974in}{1.980460in}}%
\pgfpathlineto{\pgfqpoint{1.534641in}{2.020462in}}%
\pgfpathlineto{\pgfqpoint{1.534656in}{2.017019in}}%
\pgfpathlineto{\pgfqpoint{1.535863in}{2.082872in}}%
\pgfpathlineto{\pgfqpoint{1.536769in}{2.043471in}}%
\pgfpathlineto{\pgfqpoint{1.535888in}{2.085579in}}%
\pgfpathlineto{\pgfqpoint{1.537017in}{2.061348in}}%
\pgfpathlineto{\pgfqpoint{1.537674in}{2.045223in}}%
\pgfpathlineto{\pgfqpoint{1.537202in}{2.072701in}}%
\pgfpathlineto{\pgfqpoint{1.538147in}{2.050454in}}%
\pgfpathlineto{\pgfqpoint{1.538497in}{2.078324in}}%
\pgfpathlineto{\pgfqpoint{1.539096in}{2.031717in}}%
\pgfpathlineto{\pgfqpoint{1.539106in}{2.031073in}}%
\pgfpathlineto{\pgfqpoint{1.539344in}{2.062046in}}%
\pgfpathlineto{\pgfqpoint{1.540050in}{2.048291in}}%
\pgfpathlineto{\pgfqpoint{1.541189in}{2.067962in}}%
\pgfpathlineto{\pgfqpoint{1.540469in}{2.040470in}}%
\pgfpathlineto{\pgfqpoint{1.541194in}{2.066555in}}%
\pgfpathlineto{\pgfqpoint{1.542334in}{2.032235in}}%
\pgfpathlineto{\pgfqpoint{1.541409in}{2.080503in}}%
\pgfpathlineto{\pgfqpoint{1.542358in}{2.034459in}}%
\pgfpathlineto{\pgfqpoint{1.542373in}{2.037579in}}%
\pgfpathlineto{\pgfqpoint{1.543371in}{1.994974in}}%
\pgfpathlineto{\pgfqpoint{1.543385in}{1.998841in}}%
\pgfpathlineto{\pgfqpoint{1.544597in}{1.969505in}}%
\pgfpathlineto{\pgfqpoint{1.543668in}{2.011759in}}%
\pgfpathlineto{\pgfqpoint{1.544612in}{1.970009in}}%
\pgfpathlineto{\pgfqpoint{1.545722in}{1.995726in}}%
\pgfpathlineto{\pgfqpoint{1.545187in}{1.964209in}}%
\pgfpathlineto{\pgfqpoint{1.545761in}{1.991700in}}%
\pgfpathlineto{\pgfqpoint{1.545766in}{1.989841in}}%
\pgfpathlineto{\pgfqpoint{1.546559in}{2.040902in}}%
\pgfpathlineto{\pgfqpoint{1.546793in}{2.021571in}}%
\pgfpathlineto{\pgfqpoint{1.547085in}{2.047281in}}%
\pgfpathlineto{\pgfqpoint{1.548078in}{2.032833in}}%
\pgfpathlineto{\pgfqpoint{1.549091in}{2.014738in}}%
\pgfpathlineto{\pgfqpoint{1.548210in}{2.035540in}}%
\pgfpathlineto{\pgfqpoint{1.549208in}{2.015697in}}%
\pgfpathlineto{\pgfqpoint{1.549807in}{1.997873in}}%
\pgfpathlineto{\pgfqpoint{1.550235in}{2.021011in}}%
\pgfpathlineto{\pgfqpoint{1.550289in}{2.017358in}}%
\pgfpathlineto{\pgfqpoint{1.550333in}{2.025127in}}%
\pgfpathlineto{\pgfqpoint{1.550878in}{1.984123in}}%
\pgfpathlineto{\pgfqpoint{1.551365in}{1.995631in}}%
\pgfpathlineto{\pgfqpoint{1.552592in}{1.964674in}}%
\pgfpathlineto{\pgfqpoint{1.552338in}{2.000512in}}%
\pgfpathlineto{\pgfqpoint{1.552669in}{1.981434in}}%
\pgfpathlineto{\pgfqpoint{1.553726in}{2.040446in}}%
\pgfpathlineto{\pgfqpoint{1.552884in}{1.969441in}}%
\pgfpathlineto{\pgfqpoint{1.553906in}{2.031970in}}%
\pgfpathlineto{\pgfqpoint{1.553921in}{2.029474in}}%
\pgfpathlineto{\pgfqpoint{1.554515in}{2.060866in}}%
\pgfpathlineto{\pgfqpoint{1.554948in}{2.045990in}}%
\pgfpathlineto{\pgfqpoint{1.554958in}{2.047251in}}%
\pgfpathlineto{\pgfqpoint{1.555138in}{2.021817in}}%
\pgfpathlineto{\pgfqpoint{1.555269in}{2.005022in}}%
\pgfpathlineto{\pgfqpoint{1.555615in}{2.030744in}}%
\pgfpathlineto{\pgfqpoint{1.556150in}{2.025247in}}%
\pgfpathlineto{\pgfqpoint{1.556257in}{2.039760in}}%
\pgfpathlineto{\pgfqpoint{1.557139in}{1.975341in}}%
\pgfpathlineto{\pgfqpoint{1.557212in}{1.983075in}}%
\pgfpathlineto{\pgfqpoint{1.558565in}{1.931071in}}%
\pgfpathlineto{\pgfqpoint{1.557494in}{1.990722in}}%
\pgfpathlineto{\pgfqpoint{1.558663in}{1.939284in}}%
\pgfpathlineto{\pgfqpoint{1.559738in}{1.952311in}}%
\pgfpathlineto{\pgfqpoint{1.559451in}{1.919642in}}%
\pgfpathlineto{\pgfqpoint{1.559807in}{1.949775in}}%
\pgfpathlineto{\pgfqpoint{1.559885in}{1.945760in}}%
\pgfpathlineto{\pgfqpoint{1.560050in}{1.960202in}}%
\pgfpathlineto{\pgfqpoint{1.560483in}{1.975273in}}%
\pgfpathlineto{\pgfqpoint{1.560956in}{1.942205in}}%
\pgfpathlineto{\pgfqpoint{1.561150in}{1.956658in}}%
\pgfpathlineto{\pgfqpoint{1.562426in}{1.897794in}}%
\pgfpathlineto{\pgfqpoint{1.562431in}{1.897910in}}%
\pgfpathlineto{\pgfqpoint{1.562908in}{1.912314in}}%
\pgfpathlineto{\pgfqpoint{1.563390in}{1.876924in}}%
\pgfpathlineto{\pgfqpoint{1.563395in}{1.878658in}}%
\pgfpathlineto{\pgfqpoint{1.564169in}{1.843852in}}%
\pgfpathlineto{\pgfqpoint{1.563672in}{1.891300in}}%
\pgfpathlineto{\pgfqpoint{1.564544in}{1.857773in}}%
\pgfpathlineto{\pgfqpoint{1.565571in}{1.831211in}}%
\pgfpathlineto{\pgfqpoint{1.564982in}{1.870203in}}%
\pgfpathlineto{\pgfqpoint{1.565736in}{1.839968in}}%
\pgfpathlineto{\pgfqpoint{1.566121in}{1.846278in}}%
\pgfpathlineto{\pgfqpoint{1.566214in}{1.825476in}}%
\pgfpathlineto{\pgfqpoint{1.566671in}{1.835406in}}%
\pgfpathlineto{\pgfqpoint{1.567820in}{1.820423in}}%
\pgfpathlineto{\pgfqpoint{1.567314in}{1.846231in}}%
\pgfpathlineto{\pgfqpoint{1.567830in}{1.821437in}}%
\pgfpathlineto{\pgfqpoint{1.568818in}{1.852793in}}%
\pgfpathlineto{\pgfqpoint{1.568054in}{1.809158in}}%
\pgfpathlineto{\pgfqpoint{1.568993in}{1.847962in}}%
\pgfpathlineto{\pgfqpoint{1.569261in}{1.816788in}}%
\pgfpathlineto{\pgfqpoint{1.569860in}{1.871558in}}%
\pgfpathlineto{\pgfqpoint{1.570079in}{1.860879in}}%
\pgfpathlineto{\pgfqpoint{1.570571in}{1.885850in}}%
\pgfpathlineto{\pgfqpoint{1.570999in}{1.846405in}}%
\pgfpathlineto{\pgfqpoint{1.571175in}{1.856444in}}%
\pgfpathlineto{\pgfqpoint{1.572567in}{1.780023in}}%
\pgfpathlineto{\pgfqpoint{1.571243in}{1.867014in}}%
\pgfpathlineto{\pgfqpoint{1.572572in}{1.780044in}}%
\pgfpathlineto{\pgfqpoint{1.572747in}{1.785312in}}%
\pgfpathlineto{\pgfqpoint{1.573511in}{1.730640in}}%
\pgfpathlineto{\pgfqpoint{1.573604in}{1.732600in}}%
\pgfpathlineto{\pgfqpoint{1.574519in}{1.739077in}}%
\pgfpathlineto{\pgfqpoint{1.574427in}{1.715478in}}%
\pgfpathlineto{\pgfqpoint{1.574690in}{1.729585in}}%
\pgfpathlineto{\pgfqpoint{1.575259in}{1.706951in}}%
\pgfpathlineto{\pgfqpoint{1.575658in}{1.746689in}}%
\pgfpathlineto{\pgfqpoint{1.575663in}{1.745890in}}%
\pgfpathlineto{\pgfqpoint{1.575717in}{1.753859in}}%
\pgfpathlineto{\pgfqpoint{1.576686in}{1.707516in}}%
\pgfpathlineto{\pgfqpoint{1.576715in}{1.709272in}}%
\pgfpathlineto{\pgfqpoint{1.577620in}{1.676140in}}%
\pgfpathlineto{\pgfqpoint{1.576817in}{1.714303in}}%
\pgfpathlineto{\pgfqpoint{1.577893in}{1.688821in}}%
\pgfpathlineto{\pgfqpoint{1.578307in}{1.723138in}}%
\pgfpathlineto{\pgfqpoint{1.577913in}{1.687775in}}%
\pgfpathlineto{\pgfqpoint{1.579042in}{1.721871in}}%
\pgfpathlineto{\pgfqpoint{1.579412in}{1.696335in}}%
\pgfpathlineto{\pgfqpoint{1.580064in}{1.769222in}}%
\pgfpathlineto{\pgfqpoint{1.580196in}{1.775318in}}%
\pgfpathlineto{\pgfqpoint{1.580916in}{1.743360in}}%
\pgfpathlineto{\pgfqpoint{1.581131in}{1.755062in}}%
\pgfpathlineto{\pgfqpoint{1.582255in}{1.749014in}}%
\pgfpathlineto{\pgfqpoint{1.581963in}{1.767879in}}%
\pgfpathlineto{\pgfqpoint{1.582265in}{1.751986in}}%
\pgfpathlineto{\pgfqpoint{1.582440in}{1.762734in}}%
\pgfpathlineto{\pgfqpoint{1.583073in}{1.742038in}}%
\pgfpathlineto{\pgfqpoint{1.583390in}{1.755346in}}%
\pgfpathlineto{\pgfqpoint{1.583633in}{1.760693in}}%
\pgfpathlineto{\pgfqpoint{1.584057in}{1.746723in}}%
\pgfpathlineto{\pgfqpoint{1.584261in}{1.735875in}}%
\pgfpathlineto{\pgfqpoint{1.584441in}{1.770962in}}%
\pgfpathlineto{\pgfqpoint{1.585074in}{1.766614in}}%
\pgfpathlineto{\pgfqpoint{1.585098in}{1.770947in}}%
\pgfpathlineto{\pgfqpoint{1.585585in}{1.751687in}}%
\pgfpathlineto{\pgfqpoint{1.586160in}{1.759469in}}%
\pgfpathlineto{\pgfqpoint{1.586661in}{1.730805in}}%
\pgfpathlineto{\pgfqpoint{1.586501in}{1.767015in}}%
\pgfpathlineto{\pgfqpoint{1.587299in}{1.733399in}}%
\pgfpathlineto{\pgfqpoint{1.587479in}{1.743389in}}%
\pgfpathlineto{\pgfqpoint{1.588093in}{1.699540in}}%
\pgfpathlineto{\pgfqpoint{1.588287in}{1.711174in}}%
\pgfpathlineto{\pgfqpoint{1.588453in}{1.689137in}}%
\pgfpathlineto{\pgfqpoint{1.589305in}{1.732868in}}%
\pgfpathlineto{\pgfqpoint{1.589324in}{1.728097in}}%
\pgfpathlineto{\pgfqpoint{1.589329in}{1.729964in}}%
\pgfpathlineto{\pgfqpoint{1.589904in}{1.694274in}}%
\pgfpathlineto{\pgfqpoint{1.590415in}{1.720476in}}%
\pgfpathlineto{\pgfqpoint{1.591072in}{1.701541in}}%
\pgfpathlineto{\pgfqpoint{1.591574in}{1.710142in}}%
\pgfpathlineto{\pgfqpoint{1.591622in}{1.712069in}}%
\pgfpathlineto{\pgfqpoint{1.592280in}{1.673598in}}%
\pgfpathlineto{\pgfqpoint{1.592309in}{1.679316in}}%
\pgfpathlineto{\pgfqpoint{1.593409in}{1.631842in}}%
\pgfpathlineto{\pgfqpoint{1.592426in}{1.679562in}}%
\pgfpathlineto{\pgfqpoint{1.593463in}{1.633423in}}%
\pgfpathlineto{\pgfqpoint{1.593550in}{1.646720in}}%
\pgfpathlineto{\pgfqpoint{1.594373in}{1.603317in}}%
\pgfpathlineto{\pgfqpoint{1.594539in}{1.607741in}}%
\pgfpathlineto{\pgfqpoint{1.595303in}{1.575213in}}%
\pgfpathlineto{\pgfqpoint{1.594631in}{1.611547in}}%
\pgfpathlineto{\pgfqpoint{1.595892in}{1.590494in}}%
\pgfpathlineto{\pgfqpoint{1.596812in}{1.634710in}}%
\pgfpathlineto{\pgfqpoint{1.597065in}{1.632348in}}%
\pgfpathlineto{\pgfqpoint{1.597981in}{1.579081in}}%
\pgfpathlineto{\pgfqpoint{1.598351in}{1.596816in}}%
\pgfpathlineto{\pgfqpoint{1.598497in}{1.616950in}}%
\pgfpathlineto{\pgfqpoint{1.599100in}{1.582787in}}%
\pgfpathlineto{\pgfqpoint{1.599339in}{1.587161in}}%
\pgfpathlineto{\pgfqpoint{1.599865in}{1.571011in}}%
\pgfpathlineto{\pgfqpoint{1.600386in}{1.604660in}}%
\pgfpathlineto{\pgfqpoint{1.600439in}{1.594933in}}%
\pgfpathlineto{\pgfqpoint{1.600590in}{1.578792in}}%
\pgfpathlineto{\pgfqpoint{1.600829in}{1.602247in}}%
\pgfpathlineto{\pgfqpoint{1.601023in}{1.601342in}}%
\pgfpathlineto{\pgfqpoint{1.602211in}{1.653787in}}%
\pgfpathlineto{\pgfqpoint{1.601223in}{1.596803in}}%
\pgfpathlineto{\pgfqpoint{1.602221in}{1.653462in}}%
\pgfpathlineto{\pgfqpoint{1.602606in}{1.674600in}}%
\pgfpathlineto{\pgfqpoint{1.602990in}{1.648885in}}%
\pgfpathlineto{\pgfqpoint{1.603735in}{1.665631in}}%
\pgfpathlineto{\pgfqpoint{1.604363in}{1.623584in}}%
\pgfpathlineto{\pgfqpoint{1.604865in}{1.646822in}}%
\pgfpathlineto{\pgfqpoint{1.605288in}{1.653295in}}%
\pgfpathlineto{\pgfqpoint{1.605644in}{1.616804in}}%
\pgfpathlineto{\pgfqpoint{1.605668in}{1.619404in}}%
\pgfpathlineto{\pgfqpoint{1.605673in}{1.614967in}}%
\pgfpathlineto{\pgfqpoint{1.606145in}{1.662863in}}%
\pgfpathlineto{\pgfqpoint{1.606773in}{1.624261in}}%
\pgfpathlineto{\pgfqpoint{1.607133in}{1.678864in}}%
\pgfpathlineto{\pgfqpoint{1.606783in}{1.622900in}}%
\pgfpathlineto{\pgfqpoint{1.608141in}{1.671751in}}%
\pgfpathlineto{\pgfqpoint{1.608477in}{1.659383in}}%
\pgfpathlineto{\pgfqpoint{1.609027in}{1.704502in}}%
\pgfpathlineto{\pgfqpoint{1.609198in}{1.703397in}}%
\pgfpathlineto{\pgfqpoint{1.609402in}{1.716582in}}%
\pgfpathlineto{\pgfqpoint{1.609889in}{1.682924in}}%
\pgfpathlineto{\pgfqpoint{1.610244in}{1.690804in}}%
\pgfpathlineto{\pgfqpoint{1.610278in}{1.686610in}}%
\pgfpathlineto{\pgfqpoint{1.610673in}{1.712720in}}%
\pgfpathlineto{\pgfqpoint{1.611291in}{1.696045in}}%
\pgfpathlineto{\pgfqpoint{1.611724in}{1.726248in}}%
\pgfpathlineto{\pgfqpoint{1.612347in}{1.684787in}}%
\pgfpathlineto{\pgfqpoint{1.612396in}{1.695714in}}%
\pgfpathlineto{\pgfqpoint{1.612756in}{1.650414in}}%
\pgfpathlineto{\pgfqpoint{1.613623in}{1.673083in}}%
\pgfpathlineto{\pgfqpoint{1.614090in}{1.676365in}}%
\pgfpathlineto{\pgfqpoint{1.614392in}{1.659511in}}%
\pgfpathlineto{\pgfqpoint{1.615351in}{1.646661in}}%
\pgfpathlineto{\pgfqpoint{1.615166in}{1.676198in}}%
\pgfpathlineto{\pgfqpoint{1.615483in}{1.663972in}}%
\pgfpathlineto{\pgfqpoint{1.616690in}{1.734729in}}%
\pgfpathlineto{\pgfqpoint{1.615770in}{1.656567in}}%
\pgfpathlineto{\pgfqpoint{1.616817in}{1.725312in}}%
\pgfpathlineto{\pgfqpoint{1.617995in}{1.744544in}}%
\pgfpathlineto{\pgfqpoint{1.617235in}{1.704609in}}%
\pgfpathlineto{\pgfqpoint{1.618005in}{1.740789in}}%
\pgfpathlineto{\pgfqpoint{1.618010in}{1.740328in}}%
\pgfpathlineto{\pgfqpoint{1.618623in}{1.785106in}}%
\pgfpathlineto{\pgfqpoint{1.618662in}{1.782914in}}%
\pgfpathlineto{\pgfqpoint{1.619645in}{1.818939in}}%
\pgfpathlineto{\pgfqpoint{1.619188in}{1.776969in}}%
\pgfpathlineto{\pgfqpoint{1.619918in}{1.810830in}}%
\pgfpathlineto{\pgfqpoint{1.620814in}{1.801807in}}%
\pgfpathlineto{\pgfqpoint{1.620006in}{1.824638in}}%
\pgfpathlineto{\pgfqpoint{1.621033in}{1.809077in}}%
\pgfpathlineto{\pgfqpoint{1.621121in}{1.805941in}}%
\pgfpathlineto{\pgfqpoint{1.621934in}{1.835784in}}%
\pgfpathlineto{\pgfqpoint{1.621938in}{1.833330in}}%
\pgfpathlineto{\pgfqpoint{1.622912in}{1.849240in}}%
\pgfpathlineto{\pgfqpoint{1.622318in}{1.813828in}}%
\pgfpathlineto{\pgfqpoint{1.623058in}{1.843877in}}%
\pgfpathlineto{\pgfqpoint{1.623949in}{1.809150in}}%
\pgfpathlineto{\pgfqpoint{1.624202in}{1.826370in}}%
\pgfpathlineto{\pgfqpoint{1.624305in}{1.823794in}}%
\pgfpathlineto{\pgfqpoint{1.624743in}{1.856013in}}%
\pgfpathlineto{\pgfqpoint{1.625293in}{1.828784in}}%
\pgfpathlineto{\pgfqpoint{1.625424in}{1.845501in}}%
\pgfpathlineto{\pgfqpoint{1.625731in}{1.818846in}}%
\pgfpathlineto{\pgfqpoint{1.626393in}{1.827143in}}%
\pgfpathlineto{\pgfqpoint{1.627240in}{1.812093in}}%
\pgfpathlineto{\pgfqpoint{1.626710in}{1.842653in}}%
\pgfpathlineto{\pgfqpoint{1.627450in}{1.827954in}}%
\pgfpathlineto{\pgfqpoint{1.627537in}{1.840673in}}%
\pgfpathlineto{\pgfqpoint{1.627829in}{1.807940in}}%
\pgfpathlineto{\pgfqpoint{1.628550in}{1.821365in}}%
\pgfpathlineto{\pgfqpoint{1.629431in}{1.780632in}}%
\pgfpathlineto{\pgfqpoint{1.628579in}{1.822624in}}%
\pgfpathlineto{\pgfqpoint{1.629723in}{1.805985in}}%
\pgfpathlineto{\pgfqpoint{1.630760in}{1.778324in}}%
\pgfpathlineto{\pgfqpoint{1.629835in}{1.811357in}}%
\pgfpathlineto{\pgfqpoint{1.630950in}{1.784459in}}%
\pgfpathlineto{\pgfqpoint{1.631909in}{1.820224in}}%
\pgfpathlineto{\pgfqpoint{1.631218in}{1.777056in}}%
\pgfpathlineto{\pgfqpoint{1.632109in}{1.804018in}}%
\pgfpathlineto{\pgfqpoint{1.633044in}{1.827306in}}%
\pgfpathlineto{\pgfqpoint{1.632226in}{1.792116in}}%
\pgfpathlineto{\pgfqpoint{1.633258in}{1.815586in}}%
\pgfpathlineto{\pgfqpoint{1.633321in}{1.815032in}}%
\pgfpathlineto{\pgfqpoint{1.634660in}{1.862042in}}%
\pgfpathlineto{\pgfqpoint{1.634976in}{1.833186in}}%
\pgfpathlineto{\pgfqpoint{1.634782in}{1.862324in}}%
\pgfpathlineto{\pgfqpoint{1.635789in}{1.856218in}}%
\pgfpathlineto{\pgfqpoint{1.635862in}{1.863823in}}%
\pgfpathlineto{\pgfqpoint{1.636388in}{1.821974in}}%
\pgfpathlineto{\pgfqpoint{1.636685in}{1.823429in}}%
\pgfpathlineto{\pgfqpoint{1.637041in}{1.790942in}}%
\pgfpathlineto{\pgfqpoint{1.636714in}{1.826515in}}%
\pgfpathlineto{\pgfqpoint{1.637917in}{1.806016in}}%
\pgfpathlineto{\pgfqpoint{1.637936in}{1.809986in}}%
\pgfpathlineto{\pgfqpoint{1.638428in}{1.777609in}}%
\pgfpathlineto{\pgfqpoint{1.638973in}{1.783593in}}%
\pgfpathlineto{\pgfqpoint{1.639421in}{1.809797in}}%
\pgfpathlineto{\pgfqpoint{1.639738in}{1.766155in}}%
\pgfpathlineto{\pgfqpoint{1.639996in}{1.775380in}}%
\pgfpathlineto{\pgfqpoint{1.640171in}{1.764359in}}%
\pgfpathlineto{\pgfqpoint{1.640887in}{1.811269in}}%
\pgfpathlineto{\pgfqpoint{1.641047in}{1.806510in}}%
\pgfpathlineto{\pgfqpoint{1.641310in}{1.793216in}}%
\pgfpathlineto{\pgfqpoint{1.641899in}{1.843371in}}%
\pgfpathlineto{\pgfqpoint{1.641982in}{1.837156in}}%
\pgfpathlineto{\pgfqpoint{1.641992in}{1.836575in}}%
\pgfpathlineto{\pgfqpoint{1.642323in}{1.861759in}}%
\pgfpathlineto{\pgfqpoint{1.642518in}{1.872964in}}%
\pgfpathlineto{\pgfqpoint{1.643058in}{1.842732in}}%
\pgfpathlineto{\pgfqpoint{1.643443in}{1.867519in}}%
\pgfpathlineto{\pgfqpoint{1.643457in}{1.864179in}}%
\pgfpathlineto{\pgfqpoint{1.644139in}{1.929824in}}%
\pgfpathlineto{\pgfqpoint{1.644421in}{1.911168in}}%
\pgfpathlineto{\pgfqpoint{1.644928in}{1.929866in}}%
\pgfpathlineto{\pgfqpoint{1.645317in}{1.880124in}}%
\pgfpathlineto{\pgfqpoint{1.645366in}{1.882199in}}%
\pgfpathlineto{\pgfqpoint{1.645682in}{1.873939in}}%
\pgfpathlineto{\pgfqpoint{1.646403in}{1.917243in}}%
\pgfpathlineto{\pgfqpoint{1.646432in}{1.909864in}}%
\pgfpathlineto{\pgfqpoint{1.646724in}{1.935638in}}%
\pgfpathlineto{\pgfqpoint{1.646446in}{1.907929in}}%
\pgfpathlineto{\pgfqpoint{1.647537in}{1.911652in}}%
\pgfpathlineto{\pgfqpoint{1.647605in}{1.901195in}}%
\pgfpathlineto{\pgfqpoint{1.647834in}{1.927688in}}%
\pgfpathlineto{\pgfqpoint{1.648628in}{1.919241in}}%
\pgfpathlineto{\pgfqpoint{1.649509in}{1.932838in}}%
\pgfpathlineto{\pgfqpoint{1.649187in}{1.904372in}}%
\pgfpathlineto{\pgfqpoint{1.649738in}{1.918739in}}%
\pgfpathlineto{\pgfqpoint{1.649767in}{1.920869in}}%
\pgfpathlineto{\pgfqpoint{1.650297in}{1.888617in}}%
\pgfpathlineto{\pgfqpoint{1.650385in}{1.891421in}}%
\pgfpathlineto{\pgfqpoint{1.651076in}{1.849291in}}%
\pgfpathlineto{\pgfqpoint{1.651519in}{1.879905in}}%
\pgfpathlineto{\pgfqpoint{1.651685in}{1.883800in}}%
\pgfpathlineto{\pgfqpoint{1.652362in}{1.823143in}}%
\pgfpathlineto{\pgfqpoint{1.652508in}{1.827800in}}%
\pgfpathlineto{\pgfqpoint{1.652522in}{1.828758in}}%
\pgfpathlineto{\pgfqpoint{1.652518in}{1.827129in}}%
\pgfpathlineto{\pgfqpoint{1.652527in}{1.827295in}}%
\pgfpathlineto{\pgfqpoint{1.653243in}{1.807440in}}%
\pgfpathlineto{\pgfqpoint{1.653623in}{1.831768in}}%
\pgfpathlineto{\pgfqpoint{1.653628in}{1.831544in}}%
\pgfpathlineto{\pgfqpoint{1.654333in}{1.793194in}}%
\pgfpathlineto{\pgfqpoint{1.653793in}{1.839977in}}%
\pgfpathlineto{\pgfqpoint{1.654772in}{1.819495in}}%
\pgfpathlineto{\pgfqpoint{1.654859in}{1.823767in}}%
\pgfpathlineto{\pgfqpoint{1.655258in}{1.776999in}}%
\pgfpathlineto{\pgfqpoint{1.655268in}{1.777016in}}%
\pgfpathlineto{\pgfqpoint{1.656179in}{1.748579in}}%
\pgfpathlineto{\pgfqpoint{1.655560in}{1.788701in}}%
\pgfpathlineto{\pgfqpoint{1.656461in}{1.754059in}}%
\pgfpathlineto{\pgfqpoint{1.656485in}{1.757263in}}%
\pgfpathlineto{\pgfqpoint{1.657035in}{1.735961in}}%
\pgfpathlineto{\pgfqpoint{1.657065in}{1.736119in}}%
\pgfpathlineto{\pgfqpoint{1.658204in}{1.716444in}}%
\pgfpathlineto{\pgfqpoint{1.657225in}{1.746354in}}%
\pgfpathlineto{\pgfqpoint{1.658219in}{1.716656in}}%
\pgfpathlineto{\pgfqpoint{1.658822in}{1.728732in}}%
\pgfpathlineto{\pgfqpoint{1.658642in}{1.702915in}}%
\pgfpathlineto{\pgfqpoint{1.659329in}{1.717799in}}%
\pgfpathlineto{\pgfqpoint{1.659679in}{1.731674in}}%
\pgfpathlineto{\pgfqpoint{1.659874in}{1.711317in}}%
\pgfpathlineto{\pgfqpoint{1.659898in}{1.712894in}}%
\pgfpathlineto{\pgfqpoint{1.660005in}{1.693779in}}%
\pgfpathlineto{\pgfqpoint{1.660901in}{1.756496in}}%
\pgfpathlineto{\pgfqpoint{1.660964in}{1.749325in}}%
\pgfpathlineto{\pgfqpoint{1.661135in}{1.773998in}}%
\pgfpathlineto{\pgfqpoint{1.660994in}{1.740678in}}%
\pgfpathlineto{\pgfqpoint{1.662045in}{1.742786in}}%
\pgfpathlineto{\pgfqpoint{1.662474in}{1.769181in}}%
\pgfpathlineto{\pgfqpoint{1.663189in}{1.719133in}}%
\pgfpathlineto{\pgfqpoint{1.663199in}{1.721376in}}%
\pgfpathlineto{\pgfqpoint{1.664095in}{1.674054in}}%
\pgfpathlineto{\pgfqpoint{1.664143in}{1.678076in}}%
\pgfpathlineto{\pgfqpoint{1.664689in}{1.647827in}}%
\pgfpathlineto{\pgfqpoint{1.664285in}{1.682150in}}%
\pgfpathlineto{\pgfqpoint{1.665263in}{1.672952in}}%
\pgfpathlineto{\pgfqpoint{1.666388in}{1.715274in}}%
\pgfpathlineto{\pgfqpoint{1.666432in}{1.707599in}}%
\pgfpathlineto{\pgfqpoint{1.666549in}{1.701019in}}%
\pgfpathlineto{\pgfqpoint{1.667031in}{1.721821in}}%
\pgfpathlineto{\pgfqpoint{1.667293in}{1.731380in}}%
\pgfpathlineto{\pgfqpoint{1.667805in}{1.683141in}}%
\pgfpathlineto{\pgfqpoint{1.668102in}{1.705661in}}%
\pgfpathlineto{\pgfqpoint{1.668944in}{1.684894in}}%
\pgfpathlineto{\pgfqpoint{1.669153in}{1.713325in}}%
\pgfpathlineto{\pgfqpoint{1.669163in}{1.716452in}}%
\pgfpathlineto{\pgfqpoint{1.669645in}{1.665462in}}%
\pgfpathlineto{\pgfqpoint{1.670185in}{1.672327in}}%
\pgfpathlineto{\pgfqpoint{1.670322in}{1.638544in}}%
\pgfpathlineto{\pgfqpoint{1.671334in}{1.665768in}}%
\pgfpathlineto{\pgfqpoint{1.671738in}{1.672341in}}%
\pgfpathlineto{\pgfqpoint{1.672172in}{1.641636in}}%
\pgfpathlineto{\pgfqpoint{1.672381in}{1.647427in}}%
\pgfpathlineto{\pgfqpoint{1.672405in}{1.645198in}}%
\pgfpathlineto{\pgfqpoint{1.672639in}{1.670230in}}%
\pgfpathlineto{\pgfqpoint{1.673447in}{1.651733in}}%
\pgfpathlineto{\pgfqpoint{1.674421in}{1.689035in}}%
\pgfpathlineto{\pgfqpoint{1.673627in}{1.651400in}}%
\pgfpathlineto{\pgfqpoint{1.674616in}{1.673824in}}%
\pgfpathlineto{\pgfqpoint{1.674771in}{1.665135in}}%
\pgfpathlineto{\pgfqpoint{1.675151in}{1.709093in}}%
\pgfpathlineto{\pgfqpoint{1.675589in}{1.700942in}}%
\pgfpathlineto{\pgfqpoint{1.675838in}{1.732346in}}%
\pgfpathlineto{\pgfqpoint{1.675643in}{1.693186in}}%
\pgfpathlineto{\pgfqpoint{1.676865in}{1.719522in}}%
\pgfpathlineto{\pgfqpoint{1.676875in}{1.717419in}}%
\pgfpathlineto{\pgfqpoint{1.677780in}{1.753229in}}%
\pgfpathlineto{\pgfqpoint{1.677824in}{1.746768in}}%
\pgfpathlineto{\pgfqpoint{1.677994in}{1.749904in}}%
\pgfpathlineto{\pgfqpoint{1.678842in}{1.724335in}}%
\pgfpathlineto{\pgfqpoint{1.678876in}{1.725978in}}%
\pgfpathlineto{\pgfqpoint{1.679406in}{1.718303in}}%
\pgfpathlineto{\pgfqpoint{1.679693in}{1.740981in}}%
\pgfpathlineto{\pgfqpoint{1.679961in}{1.733113in}}%
\pgfpathlineto{\pgfqpoint{1.681582in}{1.819374in}}%
\pgfpathlineto{\pgfqpoint{1.680117in}{1.729149in}}%
\pgfpathlineto{\pgfqpoint{1.681694in}{1.801619in}}%
\pgfpathlineto{\pgfqpoint{1.682683in}{1.758672in}}%
\pgfpathlineto{\pgfqpoint{1.682994in}{1.771274in}}%
\pgfpathlineto{\pgfqpoint{1.684085in}{1.783146in}}%
\pgfpathlineto{\pgfqpoint{1.683535in}{1.752836in}}%
\pgfpathlineto{\pgfqpoint{1.684114in}{1.778413in}}%
\pgfpathlineto{\pgfqpoint{1.684825in}{1.783674in}}%
\pgfpathlineto{\pgfqpoint{1.684479in}{1.761660in}}%
\pgfpathlineto{\pgfqpoint{1.685190in}{1.774401in}}%
\pgfpathlineto{\pgfqpoint{1.686173in}{1.725545in}}%
\pgfpathlineto{\pgfqpoint{1.685321in}{1.776009in}}%
\pgfpathlineto{\pgfqpoint{1.686373in}{1.749920in}}%
\pgfpathlineto{\pgfqpoint{1.687318in}{1.774566in}}%
\pgfpathlineto{\pgfqpoint{1.686675in}{1.737291in}}%
\pgfpathlineto{\pgfqpoint{1.687488in}{1.756381in}}%
\pgfpathlineto{\pgfqpoint{1.687498in}{1.753926in}}%
\pgfpathlineto{\pgfqpoint{1.688481in}{1.804193in}}%
\pgfpathlineto{\pgfqpoint{1.688491in}{1.802636in}}%
\pgfpathlineto{\pgfqpoint{1.688515in}{1.799720in}}%
\pgfpathlineto{\pgfqpoint{1.689411in}{1.838647in}}%
\pgfpathlineto{\pgfqpoint{1.689611in}{1.843806in}}%
\pgfpathlineto{\pgfqpoint{1.690404in}{1.796656in}}%
\pgfpathlineto{\pgfqpoint{1.690419in}{1.799146in}}%
\pgfpathlineto{\pgfqpoint{1.692133in}{1.729840in}}%
\pgfpathlineto{\pgfqpoint{1.690429in}{1.799190in}}%
\pgfpathlineto{\pgfqpoint{1.692220in}{1.742612in}}%
\pgfpathlineto{\pgfqpoint{1.692619in}{1.750697in}}%
\pgfpathlineto{\pgfqpoint{1.693262in}{1.713650in}}%
\pgfpathlineto{\pgfqpoint{1.693281in}{1.716259in}}%
\pgfpathlineto{\pgfqpoint{1.694469in}{1.776691in}}%
\pgfpathlineto{\pgfqpoint{1.694830in}{1.762848in}}%
\pgfpathlineto{\pgfqpoint{1.695073in}{1.744021in}}%
\pgfpathlineto{\pgfqpoint{1.695857in}{1.801961in}}%
\pgfpathlineto{\pgfqpoint{1.695872in}{1.798830in}}%
\pgfpathlineto{\pgfqpoint{1.696290in}{1.782393in}}%
\pgfpathlineto{\pgfqpoint{1.696782in}{1.816499in}}%
\pgfpathlineto{\pgfqpoint{1.696933in}{1.808653in}}%
\pgfpathlineto{\pgfqpoint{1.697989in}{1.783719in}}%
\pgfpathlineto{\pgfqpoint{1.697541in}{1.821167in}}%
\pgfpathlineto{\pgfqpoint{1.698233in}{1.791492in}}%
\pgfpathlineto{\pgfqpoint{1.699206in}{1.808305in}}%
\pgfpathlineto{\pgfqpoint{1.698968in}{1.772625in}}%
\pgfpathlineto{\pgfqpoint{1.699338in}{1.789148in}}%
\pgfpathlineto{\pgfqpoint{1.700312in}{1.746189in}}%
\pgfpathlineto{\pgfqpoint{1.699353in}{1.789455in}}%
\pgfpathlineto{\pgfqpoint{1.700574in}{1.749520in}}%
\pgfpathlineto{\pgfqpoint{1.700652in}{1.748704in}}%
\pgfpathlineto{\pgfqpoint{1.701378in}{1.790922in}}%
\pgfpathlineto{\pgfqpoint{1.701412in}{1.789172in}}%
\pgfpathlineto{\pgfqpoint{1.702288in}{1.826893in}}%
\pgfpathlineto{\pgfqpoint{1.702561in}{1.819681in}}%
\pgfpathlineto{\pgfqpoint{1.702712in}{1.836717in}}%
\pgfpathlineto{\pgfqpoint{1.703471in}{1.791065in}}%
\pgfpathlineto{\pgfqpoint{1.703617in}{1.801102in}}%
\pgfpathlineto{\pgfqpoint{1.704766in}{1.757546in}}%
\pgfpathlineto{\pgfqpoint{1.703724in}{1.811685in}}%
\pgfpathlineto{\pgfqpoint{1.704903in}{1.774307in}}%
\pgfpathlineto{\pgfqpoint{1.706013in}{1.824365in}}%
\pgfpathlineto{\pgfqpoint{1.705151in}{1.771688in}}%
\pgfpathlineto{\pgfqpoint{1.706100in}{1.819632in}}%
\pgfpathlineto{\pgfqpoint{1.706120in}{1.813259in}}%
\pgfpathlineto{\pgfqpoint{1.706743in}{1.854771in}}%
\pgfpathlineto{\pgfqpoint{1.707201in}{1.827677in}}%
\pgfpathlineto{\pgfqpoint{1.707400in}{1.812011in}}%
\pgfpathlineto{\pgfqpoint{1.708043in}{1.836518in}}%
\pgfpathlineto{\pgfqpoint{1.708325in}{1.821700in}}%
\pgfpathlineto{\pgfqpoint{1.708384in}{1.825142in}}%
\pgfpathlineto{\pgfqpoint{1.709094in}{1.779309in}}%
\pgfpathlineto{\pgfqpoint{1.709406in}{1.814888in}}%
\pgfpathlineto{\pgfqpoint{1.709421in}{1.817925in}}%
\pgfpathlineto{\pgfqpoint{1.709927in}{1.784563in}}%
\pgfpathlineto{\pgfqpoint{1.710229in}{1.787239in}}%
\pgfpathlineto{\pgfqpoint{1.710808in}{1.746111in}}%
\pgfpathlineto{\pgfqpoint{1.710389in}{1.797371in}}%
\pgfpathlineto{\pgfqpoint{1.711353in}{1.771388in}}%
\pgfpathlineto{\pgfqpoint{1.712127in}{1.752816in}}%
\pgfpathlineto{\pgfqpoint{1.712376in}{1.787630in}}%
\pgfpathlineto{\pgfqpoint{1.712390in}{1.785033in}}%
\pgfpathlineto{\pgfqpoint{1.712770in}{1.780542in}}%
\pgfpathlineto{\pgfqpoint{1.712970in}{1.805925in}}%
\pgfpathlineto{\pgfqpoint{1.712979in}{1.805795in}}%
\pgfpathlineto{\pgfqpoint{1.713062in}{1.821875in}}%
\pgfpathlineto{\pgfqpoint{1.713797in}{1.745790in}}%
\pgfpathlineto{\pgfqpoint{1.713890in}{1.754566in}}%
\pgfpathlineto{\pgfqpoint{1.715068in}{1.792381in}}%
\pgfpathlineto{\pgfqpoint{1.714133in}{1.731747in}}%
\pgfpathlineto{\pgfqpoint{1.715224in}{1.781861in}}%
\pgfpathlineto{\pgfqpoint{1.716232in}{1.744355in}}%
\pgfpathlineto{\pgfqpoint{1.715258in}{1.785889in}}%
\pgfpathlineto{\pgfqpoint{1.716426in}{1.757482in}}%
\pgfpathlineto{\pgfqpoint{1.716675in}{1.773661in}}%
\pgfpathlineto{\pgfqpoint{1.717215in}{1.722208in}}%
\pgfpathlineto{\pgfqpoint{1.717493in}{1.740443in}}%
\pgfpathlineto{\pgfqpoint{1.717502in}{1.737370in}}%
\pgfpathlineto{\pgfqpoint{1.718491in}{1.762993in}}%
\pgfpathlineto{\pgfqpoint{1.718534in}{1.760502in}}%
\pgfpathlineto{\pgfqpoint{1.718539in}{1.760564in}}%
\pgfpathlineto{\pgfqpoint{1.718632in}{1.748010in}}%
\pgfpathlineto{\pgfqpoint{1.718676in}{1.748187in}}%
\pgfpathlineto{\pgfqpoint{1.719907in}{1.724961in}}%
\pgfpathlineto{\pgfqpoint{1.719060in}{1.767630in}}%
\pgfpathlineto{\pgfqpoint{1.719932in}{1.732741in}}%
\pgfpathlineto{\pgfqpoint{1.719990in}{1.739503in}}%
\pgfpathlineto{\pgfqpoint{1.720871in}{1.708532in}}%
\pgfpathlineto{\pgfqpoint{1.721027in}{1.724931in}}%
\pgfpathlineto{\pgfqpoint{1.721032in}{1.725121in}}%
\pgfpathlineto{\pgfqpoint{1.721480in}{1.699951in}}%
\pgfpathlineto{\pgfqpoint{1.721509in}{1.704071in}}%
\pgfpathlineto{\pgfqpoint{1.722531in}{1.654149in}}%
\pgfpathlineto{\pgfqpoint{1.722658in}{1.669538in}}%
\pgfpathlineto{\pgfqpoint{1.723447in}{1.646604in}}%
\pgfpathlineto{\pgfqpoint{1.722955in}{1.685820in}}%
\pgfpathlineto{\pgfqpoint{1.723778in}{1.664031in}}%
\pgfpathlineto{\pgfqpoint{1.724343in}{1.690176in}}%
\pgfpathlineto{\pgfqpoint{1.724751in}{1.646373in}}%
\pgfpathlineto{\pgfqpoint{1.724854in}{1.653146in}}%
\pgfpathlineto{\pgfqpoint{1.725058in}{1.649571in}}%
\pgfpathlineto{\pgfqpoint{1.725725in}{1.682188in}}%
\pgfpathlineto{\pgfqpoint{1.725905in}{1.670685in}}%
\pgfpathlineto{\pgfqpoint{1.726787in}{1.701024in}}%
\pgfpathlineto{\pgfqpoint{1.726227in}{1.667246in}}%
\pgfpathlineto{\pgfqpoint{1.727045in}{1.694985in}}%
\pgfpathlineto{\pgfqpoint{1.728125in}{1.677584in}}%
\pgfpathlineto{\pgfqpoint{1.727833in}{1.720203in}}%
\pgfpathlineto{\pgfqpoint{1.728174in}{1.685907in}}%
\pgfpathlineto{\pgfqpoint{1.728612in}{1.698643in}}%
\pgfpathlineto{\pgfqpoint{1.729075in}{1.664809in}}%
\pgfpathlineto{\pgfqpoint{1.729250in}{1.677774in}}%
\pgfpathlineto{\pgfqpoint{1.730034in}{1.632231in}}%
\pgfpathlineto{\pgfqpoint{1.729674in}{1.682946in}}%
\pgfpathlineto{\pgfqpoint{1.730452in}{1.640753in}}%
\pgfpathlineto{\pgfqpoint{1.733052in}{1.721068in}}%
\pgfpathlineto{\pgfqpoint{1.730803in}{1.630654in}}%
\pgfpathlineto{\pgfqpoint{1.733821in}{1.700058in}}%
\pgfpathlineto{\pgfqpoint{1.733831in}{1.694986in}}%
\pgfpathlineto{\pgfqpoint{1.734562in}{1.741660in}}%
\pgfpathlineto{\pgfqpoint{1.734902in}{1.721370in}}%
\pgfpathlineto{\pgfqpoint{1.735379in}{1.736053in}}%
\pgfpathlineto{\pgfqpoint{1.735180in}{1.712899in}}%
\pgfpathlineto{\pgfqpoint{1.736027in}{1.731212in}}%
\pgfpathlineto{\pgfqpoint{1.736908in}{1.710523in}}%
\pgfpathlineto{\pgfqpoint{1.736304in}{1.734169in}}%
\pgfpathlineto{\pgfqpoint{1.737195in}{1.718511in}}%
\pgfpathlineto{\pgfqpoint{1.738135in}{1.767363in}}%
\pgfpathlineto{\pgfqpoint{1.738344in}{1.747312in}}%
\pgfpathlineto{\pgfqpoint{1.739591in}{1.715196in}}%
\pgfpathlineto{\pgfqpoint{1.738408in}{1.748661in}}%
\pgfpathlineto{\pgfqpoint{1.739625in}{1.716915in}}%
\pgfpathlineto{\pgfqpoint{1.740550in}{1.728881in}}%
\pgfpathlineto{\pgfqpoint{1.740399in}{1.707978in}}%
\pgfpathlineto{\pgfqpoint{1.740749in}{1.724343in}}%
\pgfpathlineto{\pgfqpoint{1.741387in}{1.698789in}}%
\pgfpathlineto{\pgfqpoint{1.741738in}{1.737998in}}%
\pgfpathlineto{\pgfqpoint{1.741791in}{1.731449in}}%
\pgfpathlineto{\pgfqpoint{1.741928in}{1.747447in}}%
\pgfpathlineto{\pgfqpoint{1.742828in}{1.706117in}}%
\pgfpathlineto{\pgfqpoint{1.742838in}{1.706787in}}%
\pgfpathlineto{\pgfqpoint{1.744084in}{1.669991in}}%
\pgfpathlineto{\pgfqpoint{1.742984in}{1.715981in}}%
\pgfpathlineto{\pgfqpoint{1.744114in}{1.676078in}}%
\pgfpathlineto{\pgfqpoint{1.744990in}{1.698525in}}%
\pgfpathlineto{\pgfqpoint{1.744591in}{1.671832in}}%
\pgfpathlineto{\pgfqpoint{1.745316in}{1.688413in}}%
\pgfpathlineto{\pgfqpoint{1.745691in}{1.665834in}}%
\pgfpathlineto{\pgfqpoint{1.746407in}{1.696243in}}%
\pgfpathlineto{\pgfqpoint{1.746421in}{1.691127in}}%
\pgfpathlineto{\pgfqpoint{1.746431in}{1.694495in}}%
\pgfpathlineto{\pgfqpoint{1.747385in}{1.624545in}}%
\pgfpathlineto{\pgfqpoint{1.747429in}{1.625641in}}%
\pgfpathlineto{\pgfqpoint{1.747604in}{1.620530in}}%
\pgfpathlineto{\pgfqpoint{1.747901in}{1.645828in}}%
\pgfpathlineto{\pgfqpoint{1.748442in}{1.634781in}}%
\pgfpathlineto{\pgfqpoint{1.749639in}{1.696037in}}%
\pgfpathlineto{\pgfqpoint{1.749732in}{1.683460in}}%
\pgfpathlineto{\pgfqpoint{1.750496in}{1.636725in}}%
\pgfpathlineto{\pgfqpoint{1.750886in}{1.659751in}}%
\pgfpathlineto{\pgfqpoint{1.750963in}{1.666274in}}%
\pgfpathlineto{\pgfqpoint{1.751445in}{1.636978in}}%
\pgfpathlineto{\pgfqpoint{1.751962in}{1.637845in}}%
\pgfpathlineto{\pgfqpoint{1.752682in}{1.679549in}}%
\pgfpathlineto{\pgfqpoint{1.752088in}{1.633001in}}%
\pgfpathlineto{\pgfqpoint{1.753291in}{1.669788in}}%
\pgfpathlineto{\pgfqpoint{1.753670in}{1.650373in}}%
\pgfpathlineto{\pgfqpoint{1.754147in}{1.684911in}}%
\pgfpathlineto{\pgfqpoint{1.754386in}{1.678602in}}%
\pgfpathlineto{\pgfqpoint{1.754464in}{1.685946in}}%
\pgfpathlineto{\pgfqpoint{1.755107in}{1.662550in}}%
\pgfpathlineto{\pgfqpoint{1.755477in}{1.674875in}}%
\pgfpathlineto{\pgfqpoint{1.755584in}{1.669560in}}%
\pgfpathlineto{\pgfqpoint{1.756148in}{1.695758in}}%
\pgfpathlineto{\pgfqpoint{1.756514in}{1.691202in}}%
\pgfpathlineto{\pgfqpoint{1.756815in}{1.708641in}}%
\pgfpathlineto{\pgfqpoint{1.757171in}{1.672635in}}%
\pgfpathlineto{\pgfqpoint{1.757624in}{1.691896in}}%
\pgfpathlineto{\pgfqpoint{1.757687in}{1.683066in}}%
\pgfpathlineto{\pgfqpoint{1.758193in}{1.711917in}}%
\pgfpathlineto{\pgfqpoint{1.758588in}{1.741992in}}%
\pgfpathlineto{\pgfqpoint{1.759264in}{1.706199in}}%
\pgfpathlineto{\pgfqpoint{1.759294in}{1.708558in}}%
\pgfpathlineto{\pgfqpoint{1.760462in}{1.648220in}}%
\pgfpathlineto{\pgfqpoint{1.759342in}{1.709502in}}%
\pgfpathlineto{\pgfqpoint{1.760486in}{1.652230in}}%
\pgfpathlineto{\pgfqpoint{1.761304in}{1.626952in}}%
\pgfpathlineto{\pgfqpoint{1.760530in}{1.656314in}}%
\pgfpathlineto{\pgfqpoint{1.761591in}{1.651689in}}%
\pgfpathlineto{\pgfqpoint{1.762867in}{1.731920in}}%
\pgfpathlineto{\pgfqpoint{1.762877in}{1.729384in}}%
\pgfpathlineto{\pgfqpoint{1.763354in}{1.769873in}}%
\pgfpathlineto{\pgfqpoint{1.763899in}{1.756268in}}%
\pgfpathlineto{\pgfqpoint{1.763904in}{1.758860in}}%
\pgfpathlineto{\pgfqpoint{1.764108in}{1.724124in}}%
\pgfpathlineto{\pgfqpoint{1.764960in}{1.732982in}}%
\pgfpathlineto{\pgfqpoint{1.765958in}{1.665991in}}%
\pgfpathlineto{\pgfqpoint{1.764995in}{1.734005in}}%
\pgfpathlineto{\pgfqpoint{1.766202in}{1.685965in}}%
\pgfpathlineto{\pgfqpoint{1.767502in}{1.735996in}}%
\pgfpathlineto{\pgfqpoint{1.766294in}{1.679275in}}%
\pgfpathlineto{\pgfqpoint{1.768003in}{1.704337in}}%
\pgfpathlineto{\pgfqpoint{1.768096in}{1.692317in}}%
\pgfpathlineto{\pgfqpoint{1.768962in}{1.714602in}}%
\pgfpathlineto{\pgfqpoint{1.769079in}{1.709511in}}%
\pgfpathlineto{\pgfqpoint{1.770530in}{1.765582in}}%
\pgfpathlineto{\pgfqpoint{1.770647in}{1.759426in}}%
\pgfpathlineto{\pgfqpoint{1.771265in}{1.739356in}}%
\pgfpathlineto{\pgfqpoint{1.771401in}{1.767128in}}%
\pgfpathlineto{\pgfqpoint{1.771752in}{1.762050in}}%
\pgfpathlineto{\pgfqpoint{1.772643in}{1.754316in}}%
\pgfpathlineto{\pgfqpoint{1.772078in}{1.777640in}}%
\pgfpathlineto{\pgfqpoint{1.772833in}{1.768586in}}%
\pgfpathlineto{\pgfqpoint{1.773037in}{1.776398in}}%
\pgfpathlineto{\pgfqpoint{1.773753in}{1.729043in}}%
\pgfpathlineto{\pgfqpoint{1.773768in}{1.730290in}}%
\pgfpathlineto{\pgfqpoint{1.775277in}{1.777218in}}%
\pgfpathlineto{\pgfqpoint{1.773821in}{1.727222in}}%
\pgfpathlineto{\pgfqpoint{1.775476in}{1.754871in}}%
\pgfpathlineto{\pgfqpoint{1.776947in}{1.686385in}}%
\pgfpathlineto{\pgfqpoint{1.775695in}{1.756862in}}%
\pgfpathlineto{\pgfqpoint{1.776986in}{1.695163in}}%
\pgfpathlineto{\pgfqpoint{1.777682in}{1.739074in}}%
\pgfpathlineto{\pgfqpoint{1.777112in}{1.683067in}}%
\pgfpathlineto{\pgfqpoint{1.778159in}{1.720640in}}%
\pgfpathlineto{\pgfqpoint{1.778982in}{1.686773in}}%
\pgfpathlineto{\pgfqpoint{1.779298in}{1.698933in}}%
\pgfpathlineto{\pgfqpoint{1.779323in}{1.703834in}}%
\pgfpathlineto{\pgfqpoint{1.780286in}{1.645111in}}%
\pgfpathlineto{\pgfqpoint{1.780296in}{1.646576in}}%
\pgfpathlineto{\pgfqpoint{1.780725in}{1.649452in}}%
\pgfpathlineto{\pgfqpoint{1.781241in}{1.611473in}}%
\pgfpathlineto{\pgfqpoint{1.781319in}{1.618014in}}%
\pgfpathlineto{\pgfqpoint{1.781338in}{1.608747in}}%
\pgfpathlineto{\pgfqpoint{1.782049in}{1.653478in}}%
\pgfpathlineto{\pgfqpoint{1.782404in}{1.636864in}}%
\pgfpathlineto{\pgfqpoint{1.782560in}{1.631262in}}%
\pgfpathlineto{\pgfqpoint{1.783100in}{1.659633in}}%
\pgfpathlineto{\pgfqpoint{1.783475in}{1.645112in}}%
\pgfpathlineto{\pgfqpoint{1.784673in}{1.627695in}}%
\pgfpathlineto{\pgfqpoint{1.784259in}{1.662106in}}%
\pgfpathlineto{\pgfqpoint{1.784692in}{1.630060in}}%
\pgfpathlineto{\pgfqpoint{1.785194in}{1.640038in}}%
\pgfpathlineto{\pgfqpoint{1.785700in}{1.600520in}}%
\pgfpathlineto{\pgfqpoint{1.785720in}{1.600525in}}%
\pgfpathlineto{\pgfqpoint{1.786445in}{1.568444in}}%
\pgfpathlineto{\pgfqpoint{1.786854in}{1.589366in}}%
\pgfpathlineto{\pgfqpoint{1.787906in}{1.605257in}}%
\pgfpathlineto{\pgfqpoint{1.787059in}{1.575856in}}%
\pgfpathlineto{\pgfqpoint{1.787988in}{1.600112in}}%
\pgfpathlineto{\pgfqpoint{1.788320in}{1.578173in}}%
\pgfpathlineto{\pgfqpoint{1.788047in}{1.609776in}}%
\pgfpathlineto{\pgfqpoint{1.789147in}{1.589894in}}%
\pgfpathlineto{\pgfqpoint{1.789172in}{1.588852in}}%
\pgfpathlineto{\pgfqpoint{1.791241in}{1.639175in}}%
\pgfpathlineto{\pgfqpoint{1.792360in}{1.655363in}}%
\pgfpathlineto{\pgfqpoint{1.791825in}{1.623403in}}%
\pgfpathlineto{\pgfqpoint{1.792380in}{1.651967in}}%
\pgfpathlineto{\pgfqpoint{1.793130in}{1.670156in}}%
\pgfpathlineto{\pgfqpoint{1.792842in}{1.639976in}}%
\pgfpathlineto{\pgfqpoint{1.793519in}{1.658281in}}%
\pgfpathlineto{\pgfqpoint{1.793558in}{1.655305in}}%
\pgfpathlineto{\pgfqpoint{1.793889in}{1.687245in}}%
\pgfpathlineto{\pgfqpoint{1.794576in}{1.668481in}}%
\pgfpathlineto{\pgfqpoint{1.794838in}{1.654540in}}%
\pgfpathlineto{\pgfqpoint{1.795486in}{1.693969in}}%
\pgfpathlineto{\pgfqpoint{1.796318in}{1.714442in}}%
\pgfpathlineto{\pgfqpoint{1.795880in}{1.677801in}}%
\pgfpathlineto{\pgfqpoint{1.796611in}{1.699828in}}%
\pgfpathlineto{\pgfqpoint{1.797253in}{1.678543in}}%
\pgfpathlineto{\pgfqpoint{1.797516in}{1.711194in}}%
\pgfpathlineto{\pgfqpoint{1.797721in}{1.693823in}}%
\pgfpathlineto{\pgfqpoint{1.797920in}{1.717173in}}%
\pgfpathlineto{\pgfqpoint{1.798689in}{1.661107in}}%
\pgfpathlineto{\pgfqpoint{1.798792in}{1.664537in}}%
\pgfpathlineto{\pgfqpoint{1.799960in}{1.720597in}}%
\pgfpathlineto{\pgfqpoint{1.798938in}{1.655018in}}%
\pgfpathlineto{\pgfqpoint{1.800004in}{1.714701in}}%
\pgfpathlineto{\pgfqpoint{1.800247in}{1.729889in}}%
\pgfpathlineto{\pgfqpoint{1.801021in}{1.695094in}}%
\pgfpathlineto{\pgfqpoint{1.801026in}{1.696294in}}%
\pgfpathlineto{\pgfqpoint{1.801830in}{1.673840in}}%
\pgfpathlineto{\pgfqpoint{1.801314in}{1.700680in}}%
\pgfpathlineto{\pgfqpoint{1.802170in}{1.685157in}}%
\pgfpathlineto{\pgfqpoint{1.802365in}{1.698639in}}%
\pgfpathlineto{\pgfqpoint{1.802901in}{1.680895in}}%
\pgfpathlineto{\pgfqpoint{1.802930in}{1.681564in}}%
\pgfpathlineto{\pgfqpoint{1.803071in}{1.668825in}}%
\pgfpathlineto{\pgfqpoint{1.803675in}{1.713079in}}%
\pgfpathlineto{\pgfqpoint{1.803967in}{1.697609in}}%
\pgfpathlineto{\pgfqpoint{1.804074in}{1.701321in}}%
\pgfpathlineto{\pgfqpoint{1.804575in}{1.676344in}}%
\pgfpathlineto{\pgfqpoint{1.805043in}{1.683452in}}%
\pgfpathlineto{\pgfqpoint{1.805734in}{1.670512in}}%
\pgfpathlineto{\pgfqpoint{1.806099in}{1.694274in}}%
\pgfpathlineto{\pgfqpoint{1.806128in}{1.687500in}}%
\pgfpathlineto{\pgfqpoint{1.806382in}{1.705083in}}%
\pgfpathlineto{\pgfqpoint{1.806688in}{1.674526in}}%
\pgfpathlineto{\pgfqpoint{1.807229in}{1.681125in}}%
\pgfpathlineto{\pgfqpoint{1.807234in}{1.681135in}}%
\pgfpathlineto{\pgfqpoint{1.807239in}{1.680163in}}%
\pgfpathlineto{\pgfqpoint{1.807297in}{1.681702in}}%
\pgfpathlineto{\pgfqpoint{1.808431in}{1.647111in}}%
\pgfpathlineto{\pgfqpoint{1.808660in}{1.659179in}}%
\pgfpathlineto{\pgfqpoint{1.809171in}{1.623919in}}%
\pgfpathlineto{\pgfqpoint{1.809512in}{1.636333in}}%
\pgfpathlineto{\pgfqpoint{1.809663in}{1.622186in}}%
\pgfpathlineto{\pgfqpoint{1.810155in}{1.656120in}}%
\pgfpathlineto{\pgfqpoint{1.810608in}{1.639176in}}%
\pgfpathlineto{\pgfqpoint{1.811718in}{1.626209in}}%
\pgfpathlineto{\pgfqpoint{1.811070in}{1.656929in}}%
\pgfpathlineto{\pgfqpoint{1.811830in}{1.635221in}}%
\pgfpathlineto{\pgfqpoint{1.812024in}{1.649487in}}%
\pgfpathlineto{\pgfqpoint{1.812701in}{1.603830in}}%
\pgfpathlineto{\pgfqpoint{1.812881in}{1.621104in}}%
\pgfpathlineto{\pgfqpoint{1.814025in}{1.597278in}}%
\pgfpathlineto{\pgfqpoint{1.813129in}{1.640111in}}%
\pgfpathlineto{\pgfqpoint{1.814040in}{1.600468in}}%
\pgfpathlineto{\pgfqpoint{1.815004in}{1.636003in}}%
\pgfpathlineto{\pgfqpoint{1.814054in}{1.600130in}}%
\pgfpathlineto{\pgfqpoint{1.815237in}{1.625049in}}%
\pgfpathlineto{\pgfqpoint{1.815666in}{1.594242in}}%
\pgfpathlineto{\pgfqpoint{1.815267in}{1.630647in}}%
\pgfpathlineto{\pgfqpoint{1.816396in}{1.603614in}}%
\pgfpathlineto{\pgfqpoint{1.817044in}{1.650147in}}%
\pgfpathlineto{\pgfqpoint{1.817565in}{1.644234in}}%
\pgfpathlineto{\pgfqpoint{1.818850in}{1.696960in}}%
\pgfpathlineto{\pgfqpoint{1.817652in}{1.640552in}}%
\pgfpathlineto{\pgfqpoint{1.818884in}{1.694001in}}%
\pgfpathlineto{\pgfqpoint{1.819867in}{1.677594in}}%
\pgfpathlineto{\pgfqpoint{1.819142in}{1.723563in}}%
\pgfpathlineto{\pgfqpoint{1.820018in}{1.680845in}}%
\pgfpathlineto{\pgfqpoint{1.820929in}{1.738905in}}%
\pgfpathlineto{\pgfqpoint{1.821260in}{1.726977in}}%
\pgfpathlineto{\pgfqpoint{1.822496in}{1.764024in}}%
\pgfpathlineto{\pgfqpoint{1.821299in}{1.724864in}}%
\pgfpathlineto{\pgfqpoint{1.822540in}{1.760354in}}%
\pgfpathlineto{\pgfqpoint{1.823149in}{1.744566in}}%
\pgfpathlineto{\pgfqpoint{1.822618in}{1.764347in}}%
\pgfpathlineto{\pgfqpoint{1.823616in}{1.762992in}}%
\pgfpathlineto{\pgfqpoint{1.823952in}{1.793786in}}%
\pgfpathlineto{\pgfqpoint{1.824746in}{1.773261in}}%
\pgfpathlineto{\pgfqpoint{1.824955in}{1.754691in}}%
\pgfpathlineto{\pgfqpoint{1.825568in}{1.790191in}}%
\pgfpathlineto{\pgfqpoint{1.826713in}{1.818293in}}%
\pgfpathlineto{\pgfqpoint{1.826401in}{1.777510in}}%
\pgfpathlineto{\pgfqpoint{1.826732in}{1.811919in}}%
\pgfpathlineto{\pgfqpoint{1.827842in}{1.772442in}}%
\pgfpathlineto{\pgfqpoint{1.827457in}{1.816384in}}%
\pgfpathlineto{\pgfqpoint{1.827915in}{1.777634in}}%
\pgfpathlineto{\pgfqpoint{1.828134in}{1.788387in}}%
\pgfpathlineto{\pgfqpoint{1.828582in}{1.761211in}}%
\pgfpathlineto{\pgfqpoint{1.829006in}{1.766980in}}%
\pgfpathlineto{\pgfqpoint{1.829045in}{1.770528in}}%
\pgfpathlineto{\pgfqpoint{1.829434in}{1.743538in}}%
\pgfpathlineto{\pgfqpoint{1.830111in}{1.765508in}}%
\pgfpathlineto{\pgfqpoint{1.830593in}{1.728579in}}%
\pgfpathlineto{\pgfqpoint{1.831211in}{1.770921in}}%
\pgfpathlineto{\pgfqpoint{1.831216in}{1.767470in}}%
\pgfpathlineto{\pgfqpoint{1.831230in}{1.774152in}}%
\pgfpathlineto{\pgfqpoint{1.831946in}{1.733076in}}%
\pgfpathlineto{\pgfqpoint{1.831951in}{1.734042in}}%
\pgfpathlineto{\pgfqpoint{1.832370in}{1.706955in}}%
\pgfpathlineto{\pgfqpoint{1.833158in}{1.715259in}}%
\pgfpathlineto{\pgfqpoint{1.833971in}{1.703539in}}%
\pgfpathlineto{\pgfqpoint{1.833358in}{1.746522in}}%
\pgfpathlineto{\pgfqpoint{1.834254in}{1.716887in}}%
\pgfpathlineto{\pgfqpoint{1.835378in}{1.773096in}}%
\pgfpathlineto{\pgfqpoint{1.835471in}{1.771929in}}%
\pgfpathlineto{\pgfqpoint{1.835719in}{1.762456in}}%
\pgfpathlineto{\pgfqpoint{1.836435in}{1.799922in}}%
\pgfpathlineto{\pgfqpoint{1.836479in}{1.792337in}}%
\pgfpathlineto{\pgfqpoint{1.836484in}{1.794562in}}%
\pgfpathlineto{\pgfqpoint{1.837087in}{1.753057in}}%
\pgfpathlineto{\pgfqpoint{1.837535in}{1.769712in}}%
\pgfpathlineto{\pgfqpoint{1.837905in}{1.766877in}}%
\pgfpathlineto{\pgfqpoint{1.838309in}{1.805566in}}%
\pgfpathlineto{\pgfqpoint{1.838543in}{1.800972in}}%
\pgfpathlineto{\pgfqpoint{1.838548in}{1.803516in}}%
\pgfpathlineto{\pgfqpoint{1.839405in}{1.754429in}}%
\pgfpathlineto{\pgfqpoint{1.839575in}{1.764057in}}%
\pgfpathlineto{\pgfqpoint{1.840680in}{1.716523in}}%
\pgfpathlineto{\pgfqpoint{1.839604in}{1.766759in}}%
\pgfpathlineto{\pgfqpoint{1.840748in}{1.720538in}}%
\pgfpathlineto{\pgfqpoint{1.841065in}{1.751808in}}%
\pgfpathlineto{\pgfqpoint{1.841342in}{1.712554in}}%
\pgfpathlineto{\pgfqpoint{1.841868in}{1.728484in}}%
\pgfpathlineto{\pgfqpoint{1.841941in}{1.721295in}}%
\pgfpathlineto{\pgfqpoint{1.842467in}{1.755782in}}%
\pgfpathlineto{\pgfqpoint{1.842949in}{1.741479in}}%
\pgfpathlineto{\pgfqpoint{1.843952in}{1.768608in}}%
\pgfpathlineto{\pgfqpoint{1.842964in}{1.740369in}}%
\pgfpathlineto{\pgfqpoint{1.844453in}{1.751953in}}%
\pgfpathlineto{\pgfqpoint{1.844862in}{1.726574in}}%
\pgfpathlineto{\pgfqpoint{1.845305in}{1.765158in}}%
\pgfpathlineto{\pgfqpoint{1.845554in}{1.760162in}}%
\pgfpathlineto{\pgfqpoint{1.845670in}{1.757406in}}%
\pgfpathlineto{\pgfqpoint{1.845607in}{1.762916in}}%
\pgfpathlineto{\pgfqpoint{1.845705in}{1.757478in}}%
\pgfpathlineto{\pgfqpoint{1.845865in}{1.747779in}}%
\pgfpathlineto{\pgfqpoint{1.846673in}{1.792384in}}%
\pgfpathlineto{\pgfqpoint{1.846751in}{1.786944in}}%
\pgfpathlineto{\pgfqpoint{1.847871in}{1.824548in}}%
\pgfpathlineto{\pgfqpoint{1.846761in}{1.784793in}}%
\pgfpathlineto{\pgfqpoint{1.848129in}{1.793778in}}%
\pgfpathlineto{\pgfqpoint{1.848694in}{1.768849in}}%
\pgfpathlineto{\pgfqpoint{1.848202in}{1.800147in}}%
\pgfpathlineto{\pgfqpoint{1.849273in}{1.782376in}}%
\pgfpathlineto{\pgfqpoint{1.850636in}{1.730503in}}%
\pgfpathlineto{\pgfqpoint{1.849351in}{1.788076in}}%
\pgfpathlineto{\pgfqpoint{1.850666in}{1.736195in}}%
\pgfpathlineto{\pgfqpoint{1.852048in}{1.798714in}}%
\pgfpathlineto{\pgfqpoint{1.852053in}{1.797971in}}%
\pgfpathlineto{\pgfqpoint{1.853075in}{1.766126in}}%
\pgfpathlineto{\pgfqpoint{1.853192in}{1.770100in}}%
\pgfpathlineto{\pgfqpoint{1.853611in}{1.794207in}}%
\pgfpathlineto{\pgfqpoint{1.854239in}{1.757828in}}%
\pgfpathlineto{\pgfqpoint{1.854244in}{1.758712in}}%
\pgfpathlineto{\pgfqpoint{1.855359in}{1.745151in}}%
\pgfpathlineto{\pgfqpoint{1.854848in}{1.782900in}}%
\pgfpathlineto{\pgfqpoint{1.855373in}{1.748789in}}%
\pgfpathlineto{\pgfqpoint{1.856191in}{1.788773in}}%
\pgfpathlineto{\pgfqpoint{1.855558in}{1.746677in}}%
\pgfpathlineto{\pgfqpoint{1.856552in}{1.780933in}}%
\pgfpathlineto{\pgfqpoint{1.857223in}{1.766130in}}%
\pgfpathlineto{\pgfqpoint{1.856634in}{1.793224in}}%
\pgfpathlineto{\pgfqpoint{1.857662in}{1.779741in}}%
\pgfpathlineto{\pgfqpoint{1.858119in}{1.796723in}}%
\pgfpathlineto{\pgfqpoint{1.857827in}{1.762819in}}%
\pgfpathlineto{\pgfqpoint{1.858733in}{1.769876in}}%
\pgfpathlineto{\pgfqpoint{1.858738in}{1.769817in}}%
\pgfpathlineto{\pgfqpoint{1.858815in}{1.780510in}}%
\pgfpathlineto{\pgfqpoint{1.858957in}{1.779871in}}%
\pgfpathlineto{\pgfqpoint{1.859555in}{1.810609in}}%
\pgfpathlineto{\pgfqpoint{1.860101in}{1.789463in}}%
\pgfpathlineto{\pgfqpoint{1.860247in}{1.772502in}}%
\pgfpathlineto{\pgfqpoint{1.861016in}{1.815677in}}%
\pgfpathlineto{\pgfqpoint{1.861021in}{1.814923in}}%
\pgfpathlineto{\pgfqpoint{1.862029in}{1.827433in}}%
\pgfpathlineto{\pgfqpoint{1.861581in}{1.801322in}}%
\pgfpathlineto{\pgfqpoint{1.862131in}{1.816077in}}%
\pgfpathlineto{\pgfqpoint{1.862525in}{1.799706in}}%
\pgfpathlineto{\pgfqpoint{1.863134in}{1.843386in}}%
\pgfpathlineto{\pgfqpoint{1.864346in}{1.905132in}}%
\pgfpathlineto{\pgfqpoint{1.864351in}{1.904278in}}%
\pgfpathlineto{\pgfqpoint{1.864998in}{1.873479in}}%
\pgfpathlineto{\pgfqpoint{1.865524in}{1.920899in}}%
\pgfpathlineto{\pgfqpoint{1.866571in}{1.866493in}}%
\pgfpathlineto{\pgfqpoint{1.866683in}{1.867822in}}%
\pgfpathlineto{\pgfqpoint{1.866805in}{1.878223in}}%
\pgfpathlineto{\pgfqpoint{1.867710in}{1.841137in}}%
\pgfpathlineto{\pgfqpoint{1.867744in}{1.842917in}}%
\pgfpathlineto{\pgfqpoint{1.868796in}{1.812194in}}%
\pgfpathlineto{\pgfqpoint{1.868231in}{1.845271in}}%
\pgfpathlineto{\pgfqpoint{1.868888in}{1.817823in}}%
\pgfpathlineto{\pgfqpoint{1.868898in}{1.821143in}}%
\pgfpathlineto{\pgfqpoint{1.869677in}{1.775235in}}%
\pgfpathlineto{\pgfqpoint{1.869828in}{1.787744in}}%
\pgfpathlineto{\pgfqpoint{1.870904in}{1.755687in}}%
\pgfpathlineto{\pgfqpoint{1.870290in}{1.797315in}}%
\pgfpathlineto{\pgfqpoint{1.870992in}{1.765908in}}%
\pgfpathlineto{\pgfqpoint{1.871016in}{1.761942in}}%
\pgfpathlineto{\pgfqpoint{1.871250in}{1.777156in}}%
\pgfpathlineto{\pgfqpoint{1.871610in}{1.756754in}}%
\pgfpathlineto{\pgfqpoint{1.872428in}{1.815440in}}%
\pgfpathlineto{\pgfqpoint{1.873134in}{1.821793in}}%
\pgfpathlineto{\pgfqpoint{1.872876in}{1.792473in}}%
\pgfpathlineto{\pgfqpoint{1.873509in}{1.812669in}}%
\pgfpathlineto{\pgfqpoint{1.873650in}{1.801270in}}%
\pgfpathlineto{\pgfqpoint{1.874463in}{1.840646in}}%
\pgfpathlineto{\pgfqpoint{1.875607in}{1.873272in}}%
\pgfpathlineto{\pgfqpoint{1.874599in}{1.835143in}}%
\pgfpathlineto{\pgfqpoint{1.875621in}{1.870896in}}%
\pgfpathlineto{\pgfqpoint{1.876157in}{1.861971in}}%
\pgfpathlineto{\pgfqpoint{1.875772in}{1.885991in}}%
\pgfpathlineto{\pgfqpoint{1.876303in}{1.881001in}}%
\pgfpathlineto{\pgfqpoint{1.876318in}{1.883055in}}%
\pgfpathlineto{\pgfqpoint{1.876990in}{1.850110in}}%
\pgfpathlineto{\pgfqpoint{1.877204in}{1.866777in}}%
\pgfpathlineto{\pgfqpoint{1.877394in}{1.858465in}}%
\pgfpathlineto{\pgfqpoint{1.877954in}{1.893395in}}%
\pgfpathlineto{\pgfqpoint{1.878309in}{1.869133in}}%
\pgfpathlineto{\pgfqpoint{1.878445in}{1.880593in}}%
\pgfpathlineto{\pgfqpoint{1.879098in}{1.855369in}}%
\pgfpathlineto{\pgfqpoint{1.879341in}{1.845648in}}%
\pgfpathlineto{\pgfqpoint{1.879711in}{1.875131in}}%
\pgfpathlineto{\pgfqpoint{1.880178in}{1.860969in}}%
\pgfpathlineto{\pgfqpoint{1.881702in}{1.790610in}}%
\pgfpathlineto{\pgfqpoint{1.880266in}{1.866510in}}%
\pgfpathlineto{\pgfqpoint{1.881858in}{1.805391in}}%
\pgfpathlineto{\pgfqpoint{1.882871in}{1.835888in}}%
\pgfpathlineto{\pgfqpoint{1.882291in}{1.802765in}}%
\pgfpathlineto{\pgfqpoint{1.883017in}{1.812888in}}%
\pgfpathlineto{\pgfqpoint{1.883450in}{1.811191in}}%
\pgfpathlineto{\pgfqpoint{1.883737in}{1.830716in}}%
\pgfpathlineto{\pgfqpoint{1.883781in}{1.830579in}}%
\pgfpathlineto{\pgfqpoint{1.884205in}{1.823007in}}%
\pgfpathlineto{\pgfqpoint{1.884950in}{1.864457in}}%
\pgfpathlineto{\pgfqpoint{1.886142in}{1.809136in}}%
\pgfpathlineto{\pgfqpoint{1.884984in}{1.870122in}}%
\pgfpathlineto{\pgfqpoint{1.886147in}{1.809558in}}%
\pgfpathlineto{\pgfqpoint{1.887077in}{1.788060in}}%
\pgfpathlineto{\pgfqpoint{1.886546in}{1.826905in}}%
\pgfpathlineto{\pgfqpoint{1.887277in}{1.801836in}}%
\pgfpathlineto{\pgfqpoint{1.887286in}{1.803043in}}%
\pgfpathlineto{\pgfqpoint{1.888216in}{1.769329in}}%
\pgfpathlineto{\pgfqpoint{1.889341in}{1.741211in}}%
\pgfpathlineto{\pgfqpoint{1.888669in}{1.779557in}}%
\pgfpathlineto{\pgfqpoint{1.889380in}{1.752510in}}%
\pgfpathlineto{\pgfqpoint{1.890505in}{1.769779in}}%
\pgfpathlineto{\pgfqpoint{1.889872in}{1.733202in}}%
\pgfpathlineto{\pgfqpoint{1.890534in}{1.767458in}}%
\pgfpathlineto{\pgfqpoint{1.891615in}{1.759165in}}%
\pgfpathlineto{\pgfqpoint{1.891327in}{1.795249in}}%
\pgfpathlineto{\pgfqpoint{1.891673in}{1.761965in}}%
\pgfpathlineto{\pgfqpoint{1.891916in}{1.769021in}}%
\pgfpathlineto{\pgfqpoint{1.892505in}{1.741240in}}%
\pgfpathlineto{\pgfqpoint{1.892574in}{1.744792in}}%
\pgfpathlineto{\pgfqpoint{1.893367in}{1.725828in}}%
\pgfpathlineto{\pgfqpoint{1.893051in}{1.761335in}}%
\pgfpathlineto{\pgfqpoint{1.893723in}{1.735141in}}%
\pgfpathlineto{\pgfqpoint{1.894229in}{1.751904in}}%
\pgfpathlineto{\pgfqpoint{1.894560in}{1.717613in}}%
\pgfpathlineto{\pgfqpoint{1.894828in}{1.735042in}}%
\pgfpathlineto{\pgfqpoint{1.895256in}{1.700593in}}%
\pgfpathlineto{\pgfqpoint{1.894901in}{1.735304in}}%
\pgfpathlineto{\pgfqpoint{1.895952in}{1.727531in}}%
\pgfpathlineto{\pgfqpoint{1.896989in}{1.757942in}}%
\pgfpathlineto{\pgfqpoint{1.896118in}{1.717878in}}%
\pgfpathlineto{\pgfqpoint{1.897135in}{1.747851in}}%
\pgfpathlineto{\pgfqpoint{1.897880in}{1.706308in}}%
\pgfpathlineto{\pgfqpoint{1.897189in}{1.751589in}}%
\pgfpathlineto{\pgfqpoint{1.898353in}{1.717323in}}%
\pgfpathlineto{\pgfqpoint{1.899687in}{1.752090in}}%
\pgfpathlineto{\pgfqpoint{1.899107in}{1.711032in}}%
\pgfpathlineto{\pgfqpoint{1.899774in}{1.745966in}}%
\pgfpathlineto{\pgfqpoint{1.900680in}{1.731061in}}%
\pgfpathlineto{\pgfqpoint{1.900178in}{1.758940in}}%
\pgfpathlineto{\pgfqpoint{1.900889in}{1.739096in}}%
\pgfpathlineto{\pgfqpoint{1.901687in}{1.756382in}}%
\pgfpathlineto{\pgfqpoint{1.901911in}{1.726707in}}%
\pgfpathlineto{\pgfqpoint{1.901989in}{1.734954in}}%
\pgfpathlineto{\pgfqpoint{1.902004in}{1.733729in}}%
\pgfpathlineto{\pgfqpoint{1.902252in}{1.764569in}}%
\pgfpathlineto{\pgfqpoint{1.902973in}{1.751111in}}%
\pgfpathlineto{\pgfqpoint{1.902978in}{1.752652in}}%
\pgfpathlineto{\pgfqpoint{1.903635in}{1.710909in}}%
\pgfpathlineto{\pgfqpoint{1.903995in}{1.725569in}}%
\pgfpathlineto{\pgfqpoint{1.904448in}{1.714935in}}%
\pgfpathlineto{\pgfqpoint{1.904915in}{1.750707in}}%
\pgfpathlineto{\pgfqpoint{1.905066in}{1.741320in}}%
\pgfpathlineto{\pgfqpoint{1.905164in}{1.750937in}}%
\pgfpathlineto{\pgfqpoint{1.905806in}{1.711207in}}%
\pgfpathlineto{\pgfqpoint{1.905826in}{1.714220in}}%
\pgfpathlineto{\pgfqpoint{1.906196in}{1.695209in}}%
\pgfpathlineto{\pgfqpoint{1.906838in}{1.718735in}}%
\pgfpathlineto{\pgfqpoint{1.906931in}{1.715351in}}%
\pgfpathlineto{\pgfqpoint{1.908182in}{1.756157in}}%
\pgfpathlineto{\pgfqpoint{1.907384in}{1.700273in}}%
\pgfpathlineto{\pgfqpoint{1.908192in}{1.755335in}}%
\pgfpathlineto{\pgfqpoint{1.908216in}{1.759047in}}%
\pgfpathlineto{\pgfqpoint{1.908489in}{1.738698in}}%
\pgfpathlineto{\pgfqpoint{1.909151in}{1.743971in}}%
\pgfpathlineto{\pgfqpoint{1.910129in}{1.713604in}}%
\pgfpathlineto{\pgfqpoint{1.909560in}{1.747953in}}%
\pgfpathlineto{\pgfqpoint{1.910280in}{1.729644in}}%
\pgfpathlineto{\pgfqpoint{1.911317in}{1.767415in}}%
\pgfpathlineto{\pgfqpoint{1.910310in}{1.724785in}}%
\pgfpathlineto{\pgfqpoint{1.911502in}{1.759148in}}%
\pgfpathlineto{\pgfqpoint{1.912461in}{1.794683in}}%
\pgfpathlineto{\pgfqpoint{1.911653in}{1.751502in}}%
\pgfpathlineto{\pgfqpoint{1.912749in}{1.775356in}}%
\pgfpathlineto{\pgfqpoint{1.913065in}{1.768867in}}%
\pgfpathlineto{\pgfqpoint{1.913596in}{1.799786in}}%
\pgfpathlineto{\pgfqpoint{1.913844in}{1.786558in}}%
\pgfpathlineto{\pgfqpoint{1.913980in}{1.800275in}}%
\pgfpathlineto{\pgfqpoint{1.914886in}{1.776609in}}%
\pgfpathlineto{\pgfqpoint{1.914944in}{1.785624in}}%
\pgfpathlineto{\pgfqpoint{1.915860in}{1.738548in}}%
\pgfpathlineto{\pgfqpoint{1.915110in}{1.801832in}}%
\pgfpathlineto{\pgfqpoint{1.916108in}{1.752118in}}%
\pgfpathlineto{\pgfqpoint{1.917827in}{1.701496in}}%
\pgfpathlineto{\pgfqpoint{1.917997in}{1.707182in}}%
\pgfpathlineto{\pgfqpoint{1.918372in}{1.718655in}}%
\pgfpathlineto{\pgfqpoint{1.918859in}{1.692830in}}%
\pgfpathlineto{\pgfqpoint{1.919112in}{1.711451in}}%
\pgfpathlineto{\pgfqpoint{1.920271in}{1.752851in}}%
\pgfpathlineto{\pgfqpoint{1.919141in}{1.710095in}}%
\pgfpathlineto{\pgfqpoint{1.920899in}{1.730277in}}%
\pgfpathlineto{\pgfqpoint{1.922077in}{1.688962in}}%
\pgfpathlineto{\pgfqpoint{1.921244in}{1.739603in}}%
\pgfpathlineto{\pgfqpoint{1.922082in}{1.691103in}}%
\pgfpathlineto{\pgfqpoint{1.923985in}{1.740299in}}%
\pgfpathlineto{\pgfqpoint{1.922208in}{1.681276in}}%
\pgfpathlineto{\pgfqpoint{1.924034in}{1.734836in}}%
\pgfpathlineto{\pgfqpoint{1.924180in}{1.724251in}}%
\pgfpathlineto{\pgfqpoint{1.924866in}{1.761588in}}%
\pgfpathlineto{\pgfqpoint{1.925105in}{1.758799in}}%
\pgfpathlineto{\pgfqpoint{1.925124in}{1.765114in}}%
\pgfpathlineto{\pgfqpoint{1.926025in}{1.735350in}}%
\pgfpathlineto{\pgfqpoint{1.926064in}{1.736856in}}%
\pgfpathlineto{\pgfqpoint{1.926084in}{1.732400in}}%
\pgfpathlineto{\pgfqpoint{1.926512in}{1.766635in}}%
\pgfpathlineto{\pgfqpoint{1.927116in}{1.762420in}}%
\pgfpathlineto{\pgfqpoint{1.928294in}{1.809634in}}%
\pgfpathlineto{\pgfqpoint{1.927125in}{1.760468in}}%
\pgfpathlineto{\pgfqpoint{1.928508in}{1.808618in}}%
\pgfpathlineto{\pgfqpoint{1.928528in}{1.806747in}}%
\pgfpathlineto{\pgfqpoint{1.929414in}{1.845053in}}%
\pgfpathlineto{\pgfqpoint{1.929540in}{1.826620in}}%
\pgfpathlineto{\pgfqpoint{1.929545in}{1.830444in}}%
\pgfpathlineto{\pgfqpoint{1.930120in}{1.796573in}}%
\pgfpathlineto{\pgfqpoint{1.930611in}{1.798674in}}%
\pgfpathlineto{\pgfqpoint{1.931653in}{1.745297in}}%
\pgfpathlineto{\pgfqpoint{1.930626in}{1.799390in}}%
\pgfpathlineto{\pgfqpoint{1.931931in}{1.756152in}}%
\pgfpathlineto{\pgfqpoint{1.933201in}{1.809603in}}%
\pgfpathlineto{\pgfqpoint{1.932218in}{1.734728in}}%
\pgfpathlineto{\pgfqpoint{1.933226in}{1.808908in}}%
\pgfpathlineto{\pgfqpoint{1.934029in}{1.799724in}}%
\pgfpathlineto{\pgfqpoint{1.933459in}{1.830051in}}%
\pgfpathlineto{\pgfqpoint{1.934350in}{1.802734in}}%
\pgfpathlineto{\pgfqpoint{1.935280in}{1.823794in}}%
\pgfpathlineto{\pgfqpoint{1.934910in}{1.792771in}}%
\pgfpathlineto{\pgfqpoint{1.935475in}{1.809841in}}%
\pgfpathlineto{\pgfqpoint{1.936010in}{1.775962in}}%
\pgfpathlineto{\pgfqpoint{1.935524in}{1.813633in}}%
\pgfpathlineto{\pgfqpoint{1.936624in}{1.791427in}}%
\pgfpathlineto{\pgfqpoint{1.936906in}{1.797779in}}%
\pgfpathlineto{\pgfqpoint{1.937617in}{1.781130in}}%
\pgfpathlineto{\pgfqpoint{1.937646in}{1.784591in}}%
\pgfpathlineto{\pgfqpoint{1.938586in}{1.754193in}}%
\pgfpathlineto{\pgfqpoint{1.938187in}{1.795114in}}%
\pgfpathlineto{\pgfqpoint{1.938800in}{1.760415in}}%
\pgfpathlineto{\pgfqpoint{1.939516in}{1.796590in}}%
\pgfpathlineto{\pgfqpoint{1.939959in}{1.782128in}}%
\pgfpathlineto{\pgfqpoint{1.940042in}{1.771795in}}%
\pgfpathlineto{\pgfqpoint{1.940222in}{1.796026in}}%
\pgfpathlineto{\pgfqpoint{1.940665in}{1.795433in}}%
\pgfpathlineto{\pgfqpoint{1.941483in}{1.828017in}}%
\pgfpathlineto{\pgfqpoint{1.940728in}{1.790368in}}%
\pgfpathlineto{\pgfqpoint{1.941799in}{1.807222in}}%
\pgfpathlineto{\pgfqpoint{1.941906in}{1.784271in}}%
\pgfpathlineto{\pgfqpoint{1.942145in}{1.815431in}}%
\pgfpathlineto{\pgfqpoint{1.942894in}{1.813116in}}%
\pgfpathlineto{\pgfqpoint{1.943221in}{1.825853in}}%
\pgfpathlineto{\pgfqpoint{1.943902in}{1.770124in}}%
\pgfpathlineto{\pgfqpoint{1.943961in}{1.782588in}}%
\pgfpathlineto{\pgfqpoint{1.944043in}{1.774071in}}%
\pgfpathlineto{\pgfqpoint{1.944925in}{1.829488in}}%
\pgfpathlineto{\pgfqpoint{1.944954in}{1.823968in}}%
\pgfpathlineto{\pgfqpoint{1.945188in}{1.844200in}}%
\pgfpathlineto{\pgfqpoint{1.945382in}{1.815716in}}%
\pgfpathlineto{\pgfqpoint{1.946044in}{1.819822in}}%
\pgfpathlineto{\pgfqpoint{1.946638in}{1.811825in}}%
\pgfpathlineto{\pgfqpoint{1.947101in}{1.836628in}}%
\pgfpathlineto{\pgfqpoint{1.947106in}{1.835438in}}%
\pgfpathlineto{\pgfqpoint{1.947291in}{1.850769in}}%
\pgfpathlineto{\pgfqpoint{1.947919in}{1.816666in}}%
\pgfpathlineto{\pgfqpoint{1.948211in}{1.833581in}}%
\pgfpathlineto{\pgfqpoint{1.948469in}{1.846945in}}%
\pgfpathlineto{\pgfqpoint{1.949082in}{1.819418in}}%
\pgfpathlineto{\pgfqpoint{1.949204in}{1.815793in}}%
\pgfpathlineto{\pgfqpoint{1.950066in}{1.849865in}}%
\pgfpathlineto{\pgfqpoint{1.950158in}{1.844731in}}%
\pgfpathlineto{\pgfqpoint{1.951088in}{1.806933in}}%
\pgfpathlineto{\pgfqpoint{1.950207in}{1.847852in}}%
\pgfpathlineto{\pgfqpoint{1.951677in}{1.830489in}}%
\pgfpathlineto{\pgfqpoint{1.952924in}{1.905407in}}%
\pgfpathlineto{\pgfqpoint{1.951702in}{1.827074in}}%
\pgfpathlineto{\pgfqpoint{1.952933in}{1.904140in}}%
\pgfpathlineto{\pgfqpoint{1.954068in}{1.866264in}}%
\pgfpathlineto{\pgfqpoint{1.953113in}{1.914601in}}%
\pgfpathlineto{\pgfqpoint{1.954121in}{1.870452in}}%
\pgfpathlineto{\pgfqpoint{1.954136in}{1.877434in}}%
\pgfpathlineto{\pgfqpoint{1.954973in}{1.839819in}}%
\pgfpathlineto{\pgfqpoint{1.955222in}{1.864574in}}%
\pgfpathlineto{\pgfqpoint{1.955226in}{1.864686in}}%
\pgfpathlineto{\pgfqpoint{1.955421in}{1.844787in}}%
\pgfpathlineto{\pgfqpoint{1.955548in}{1.848121in}}%
\pgfpathlineto{\pgfqpoint{1.955562in}{1.842991in}}%
\pgfpathlineto{\pgfqpoint{1.956015in}{1.879412in}}%
\pgfpathlineto{\pgfqpoint{1.956624in}{1.874634in}}%
\pgfpathlineto{\pgfqpoint{1.957846in}{1.911940in}}%
\pgfpathlineto{\pgfqpoint{1.958415in}{1.918927in}}%
\pgfpathlineto{\pgfqpoint{1.957982in}{1.896071in}}%
\pgfpathlineto{\pgfqpoint{1.958892in}{1.908751in}}%
\pgfpathlineto{\pgfqpoint{1.959472in}{1.875254in}}%
\pgfpathlineto{\pgfqpoint{1.958912in}{1.910600in}}%
\pgfpathlineto{\pgfqpoint{1.960071in}{1.902349in}}%
\pgfpathlineto{\pgfqpoint{1.962933in}{2.004576in}}%
\pgfpathlineto{\pgfqpoint{1.960153in}{1.893146in}}%
\pgfpathlineto{\pgfqpoint{1.963070in}{1.983602in}}%
\pgfpathlineto{\pgfqpoint{1.964189in}{1.965081in}}%
\pgfpathlineto{\pgfqpoint{1.963683in}{2.000541in}}%
\pgfpathlineto{\pgfqpoint{1.964214in}{1.969671in}}%
\pgfpathlineto{\pgfqpoint{1.965207in}{1.976933in}}%
\pgfpathlineto{\pgfqpoint{1.964652in}{1.939694in}}%
\pgfpathlineto{\pgfqpoint{1.965309in}{1.965769in}}%
\pgfpathlineto{\pgfqpoint{1.966385in}{1.928348in}}%
\pgfpathlineto{\pgfqpoint{1.965382in}{1.977411in}}%
\pgfpathlineto{\pgfqpoint{1.966492in}{1.933480in}}%
\pgfpathlineto{\pgfqpoint{1.966589in}{1.942995in}}%
\pgfpathlineto{\pgfqpoint{1.966755in}{1.920003in}}%
\pgfpathlineto{\pgfqpoint{1.967125in}{1.893589in}}%
\pgfpathlineto{\pgfqpoint{1.967631in}{1.925372in}}%
\pgfpathlineto{\pgfqpoint{1.967855in}{1.923455in}}%
\pgfpathlineto{\pgfqpoint{1.968756in}{1.953550in}}%
\pgfpathlineto{\pgfqpoint{1.968031in}{1.914196in}}%
\pgfpathlineto{\pgfqpoint{1.969082in}{1.943081in}}%
\pgfpathlineto{\pgfqpoint{1.969798in}{1.924116in}}%
\pgfpathlineto{\pgfqpoint{1.969150in}{1.955479in}}%
\pgfpathlineto{\pgfqpoint{1.970187in}{1.942498in}}%
\pgfpathlineto{\pgfqpoint{1.971156in}{1.949960in}}%
\pgfpathlineto{\pgfqpoint{1.970669in}{1.912322in}}%
\pgfpathlineto{\pgfqpoint{1.971297in}{1.945367in}}%
\pgfpathlineto{\pgfqpoint{1.971794in}{1.968141in}}%
\pgfpathlineto{\pgfqpoint{1.972437in}{1.924250in}}%
\pgfpathlineto{\pgfqpoint{1.972719in}{1.936561in}}%
\pgfpathlineto{\pgfqpoint{1.973284in}{1.894438in}}%
\pgfpathlineto{\pgfqpoint{1.973478in}{1.904904in}}%
\pgfpathlineto{\pgfqpoint{1.974194in}{1.860979in}}%
\pgfpathlineto{\pgfqpoint{1.974608in}{1.885745in}}%
\pgfpathlineto{\pgfqpoint{1.975674in}{1.903379in}}%
\pgfpathlineto{\pgfqpoint{1.974744in}{1.877068in}}%
\pgfpathlineto{\pgfqpoint{1.975742in}{1.899139in}}%
\pgfpathlineto{\pgfqpoint{1.976911in}{1.879225in}}%
\pgfpathlineto{\pgfqpoint{1.975840in}{1.906456in}}%
\pgfpathlineto{\pgfqpoint{1.976916in}{1.883111in}}%
\pgfpathlineto{\pgfqpoint{1.977232in}{1.908982in}}%
\pgfpathlineto{\pgfqpoint{1.977733in}{1.875956in}}%
\pgfpathlineto{\pgfqpoint{1.978030in}{1.885054in}}%
\pgfpathlineto{\pgfqpoint{1.978060in}{1.880959in}}%
\pgfpathlineto{\pgfqpoint{1.978069in}{1.885763in}}%
\pgfpathlineto{\pgfqpoint{1.978099in}{1.882057in}}%
\pgfpathlineto{\pgfqpoint{1.979199in}{1.857006in}}%
\pgfpathlineto{\pgfqpoint{1.978230in}{1.889113in}}%
\pgfpathlineto{\pgfqpoint{1.979233in}{1.865473in}}%
\pgfpathlineto{\pgfqpoint{1.979627in}{1.840622in}}%
\pgfpathlineto{\pgfqpoint{1.979929in}{1.874147in}}%
\pgfpathlineto{\pgfqpoint{1.980343in}{1.863975in}}%
\pgfpathlineto{\pgfqpoint{1.980869in}{1.883748in}}%
\pgfpathlineto{\pgfqpoint{1.981331in}{1.845461in}}%
\pgfpathlineto{\pgfqpoint{1.981336in}{1.846236in}}%
\pgfpathlineto{\pgfqpoint{1.981945in}{1.790292in}}%
\pgfpathlineto{\pgfqpoint{1.981356in}{1.847980in}}%
\pgfpathlineto{\pgfqpoint{1.982524in}{1.809557in}}%
\pgfpathlineto{\pgfqpoint{1.983834in}{1.864516in}}%
\pgfpathlineto{\pgfqpoint{1.982578in}{1.806758in}}%
\pgfpathlineto{\pgfqpoint{1.983892in}{1.859901in}}%
\pgfpathlineto{\pgfqpoint{1.985007in}{1.807763in}}%
\pgfpathlineto{\pgfqpoint{1.984077in}{1.864050in}}%
\pgfpathlineto{\pgfqpoint{1.985445in}{1.825106in}}%
\pgfpathlineto{\pgfqpoint{1.985494in}{1.833660in}}%
\pgfpathlineto{\pgfqpoint{1.986433in}{1.803999in}}%
\pgfpathlineto{\pgfqpoint{1.986526in}{1.810887in}}%
\pgfpathlineto{\pgfqpoint{1.986541in}{1.814041in}}%
\pgfpathlineto{\pgfqpoint{1.987475in}{1.783989in}}%
\pgfpathlineto{\pgfqpoint{1.987621in}{1.786715in}}%
\pgfpathlineto{\pgfqpoint{1.988634in}{1.758137in}}%
\pgfpathlineto{\pgfqpoint{1.989764in}{1.849527in}}%
\pgfpathlineto{\pgfqpoint{1.989949in}{1.836365in}}%
\pgfpathlineto{\pgfqpoint{1.990143in}{1.824441in}}%
\pgfpathlineto{\pgfqpoint{1.990849in}{1.852532in}}%
\pgfpathlineto{\pgfqpoint{1.991054in}{1.841180in}}%
\pgfpathlineto{\pgfqpoint{1.992125in}{1.862102in}}%
\pgfpathlineto{\pgfqpoint{1.991463in}{1.833557in}}%
\pgfpathlineto{\pgfqpoint{1.992193in}{1.859706in}}%
\pgfpathlineto{\pgfqpoint{1.993493in}{1.796399in}}%
\pgfpathlineto{\pgfqpoint{1.992246in}{1.865987in}}%
\pgfpathlineto{\pgfqpoint{1.993644in}{1.807336in}}%
\pgfpathlineto{\pgfqpoint{1.994900in}{1.846788in}}%
\pgfpathlineto{\pgfqpoint{1.993751in}{1.795392in}}%
\pgfpathlineto{\pgfqpoint{1.994910in}{1.843080in}}%
\pgfpathlineto{\pgfqpoint{1.995226in}{1.829154in}}%
\pgfpathlineto{\pgfqpoint{1.995927in}{1.855255in}}%
\pgfpathlineto{\pgfqpoint{1.995995in}{1.853524in}}%
\pgfpathlineto{\pgfqpoint{1.996910in}{1.897229in}}%
\pgfpathlineto{\pgfqpoint{1.996200in}{1.852323in}}%
\pgfpathlineto{\pgfqpoint{1.997188in}{1.886358in}}%
\pgfpathlineto{\pgfqpoint{1.997514in}{1.868877in}}%
\pgfpathlineto{\pgfqpoint{1.997767in}{1.891215in}}%
\pgfpathlineto{\pgfqpoint{1.998327in}{1.880064in}}%
\pgfpathlineto{\pgfqpoint{1.998902in}{1.904690in}}%
\pgfpathlineto{\pgfqpoint{1.999403in}{1.870169in}}%
\pgfpathlineto{\pgfqpoint{2.000241in}{1.823251in}}%
\pgfpathlineto{\pgfqpoint{2.000533in}{1.863193in}}%
\pgfpathlineto{\pgfqpoint{2.000581in}{1.870529in}}%
\pgfpathlineto{\pgfqpoint{2.001000in}{1.813963in}}%
\pgfpathlineto{\pgfqpoint{2.001565in}{1.825644in}}%
\pgfpathlineto{\pgfqpoint{2.002203in}{1.804662in}}%
\pgfpathlineto{\pgfqpoint{2.001789in}{1.828739in}}%
\pgfpathlineto{\pgfqpoint{2.002685in}{1.819110in}}%
\pgfpathlineto{\pgfqpoint{2.003118in}{1.851273in}}%
\pgfpathlineto{\pgfqpoint{2.002762in}{1.816446in}}%
\pgfpathlineto{\pgfqpoint{2.003824in}{1.844263in}}%
\pgfpathlineto{\pgfqpoint{2.004812in}{1.825912in}}%
\pgfpathlineto{\pgfqpoint{2.004301in}{1.865393in}}%
\pgfpathlineto{\pgfqpoint{2.004992in}{1.831072in}}%
\pgfpathlineto{\pgfqpoint{2.005075in}{1.841995in}}%
\pgfpathlineto{\pgfqpoint{2.005976in}{1.800058in}}%
\pgfpathlineto{\pgfqpoint{2.006068in}{1.805314in}}%
\pgfpathlineto{\pgfqpoint{2.006443in}{1.822855in}}%
\pgfpathlineto{\pgfqpoint{2.006131in}{1.796526in}}%
\pgfpathlineto{\pgfqpoint{2.007129in}{1.799044in}}%
\pgfpathlineto{\pgfqpoint{2.008152in}{1.781452in}}%
\pgfpathlineto{\pgfqpoint{2.007899in}{1.810567in}}%
\pgfpathlineto{\pgfqpoint{2.008264in}{1.785905in}}%
\pgfpathlineto{\pgfqpoint{2.008605in}{1.775800in}}%
\pgfpathlineto{\pgfqpoint{2.008926in}{1.804186in}}%
\pgfpathlineto{\pgfqpoint{2.009388in}{1.778762in}}%
\pgfpathlineto{\pgfqpoint{2.009651in}{1.814567in}}%
\pgfpathlineto{\pgfqpoint{2.010542in}{1.806677in}}%
\pgfpathlineto{\pgfqpoint{2.010552in}{1.802220in}}%
\pgfpathlineto{\pgfqpoint{2.011462in}{1.828992in}}%
\pgfpathlineto{\pgfqpoint{2.011628in}{1.822760in}}%
\pgfpathlineto{\pgfqpoint{2.011633in}{1.823904in}}%
\pgfpathlineto{\pgfqpoint{2.012392in}{1.771856in}}%
\pgfpathlineto{\pgfqpoint{2.012558in}{1.782712in}}%
\pgfpathlineto{\pgfqpoint{2.013609in}{1.758461in}}%
\pgfpathlineto{\pgfqpoint{2.012626in}{1.789583in}}%
\pgfpathlineto{\pgfqpoint{2.013687in}{1.773254in}}%
\pgfpathlineto{\pgfqpoint{2.014320in}{1.798047in}}%
\pgfpathlineto{\pgfqpoint{2.013887in}{1.771510in}}%
\pgfpathlineto{\pgfqpoint{2.014856in}{1.782440in}}%
\pgfpathlineto{\pgfqpoint{2.015435in}{1.768064in}}%
\pgfpathlineto{\pgfqpoint{2.015883in}{1.798566in}}%
\pgfpathlineto{\pgfqpoint{2.015951in}{1.792796in}}%
\pgfpathlineto{\pgfqpoint{2.016370in}{1.803866in}}%
\pgfpathlineto{\pgfqpoint{2.016019in}{1.786755in}}%
\pgfpathlineto{\pgfqpoint{2.016574in}{1.788604in}}%
\pgfpathlineto{\pgfqpoint{2.017621in}{1.761396in}}%
\pgfpathlineto{\pgfqpoint{2.016764in}{1.804943in}}%
\pgfpathlineto{\pgfqpoint{2.017718in}{1.766213in}}%
\pgfpathlineto{\pgfqpoint{2.018230in}{1.787231in}}%
\pgfpathlineto{\pgfqpoint{2.017855in}{1.758521in}}%
\pgfpathlineto{\pgfqpoint{2.018843in}{1.770864in}}%
\pgfpathlineto{\pgfqpoint{2.020031in}{1.738436in}}%
\pgfpathlineto{\pgfqpoint{2.019179in}{1.784481in}}%
\pgfpathlineto{\pgfqpoint{2.020075in}{1.748344in}}%
\pgfpathlineto{\pgfqpoint{2.020805in}{1.778474in}}%
\pgfpathlineto{\pgfqpoint{2.021204in}{1.759618in}}%
\pgfpathlineto{\pgfqpoint{2.021876in}{1.779945in}}%
\pgfpathlineto{\pgfqpoint{2.021350in}{1.753929in}}%
\pgfpathlineto{\pgfqpoint{2.022095in}{1.757905in}}%
\pgfpathlineto{\pgfqpoint{2.022499in}{1.749207in}}%
\pgfpathlineto{\pgfqpoint{2.023054in}{1.780542in}}%
\pgfpathlineto{\pgfqpoint{2.023176in}{1.775926in}}%
\pgfpathlineto{\pgfqpoint{2.023298in}{1.780161in}}%
\pgfpathlineto{\pgfqpoint{2.023731in}{1.754938in}}%
\pgfpathlineto{\pgfqpoint{2.024276in}{1.774940in}}%
\pgfpathlineto{\pgfqpoint{2.024710in}{1.764799in}}%
\pgfpathlineto{\pgfqpoint{2.024943in}{1.784153in}}%
\pgfpathlineto{\pgfqpoint{2.025401in}{1.770387in}}%
\pgfpathlineto{\pgfqpoint{2.026418in}{1.760204in}}%
\pgfpathlineto{\pgfqpoint{2.025571in}{1.791175in}}%
\pgfpathlineto{\pgfqpoint{2.026530in}{1.765474in}}%
\pgfpathlineto{\pgfqpoint{2.026837in}{1.790987in}}%
\pgfpathlineto{\pgfqpoint{2.027451in}{1.756795in}}%
\pgfpathlineto{\pgfqpoint{2.027704in}{1.783600in}}%
\pgfpathlineto{\pgfqpoint{2.027864in}{1.775982in}}%
\pgfpathlineto{\pgfqpoint{2.028488in}{1.814834in}}%
\pgfpathlineto{\pgfqpoint{2.028809in}{1.785805in}}%
\pgfpathlineto{\pgfqpoint{2.028853in}{1.778405in}}%
\pgfpathlineto{\pgfqpoint{2.028921in}{1.800323in}}%
\pgfpathlineto{\pgfqpoint{2.029914in}{1.785362in}}%
\pgfpathlineto{\pgfqpoint{2.031092in}{1.869990in}}%
\pgfpathlineto{\pgfqpoint{2.031453in}{1.849500in}}%
\pgfpathlineto{\pgfqpoint{2.032256in}{1.808235in}}%
\pgfpathlineto{\pgfqpoint{2.032611in}{1.825184in}}%
\pgfpathlineto{\pgfqpoint{2.033059in}{1.843806in}}%
\pgfpathlineto{\pgfqpoint{2.033371in}{1.819911in}}%
\pgfpathlineto{\pgfqpoint{2.033741in}{1.834603in}}%
\pgfpathlineto{\pgfqpoint{2.034052in}{1.822016in}}%
\pgfpathlineto{\pgfqpoint{2.034505in}{1.867315in}}%
\pgfpathlineto{\pgfqpoint{2.034851in}{1.834000in}}%
\pgfpathlineto{\pgfqpoint{2.035805in}{1.848594in}}%
\pgfpathlineto{\pgfqpoint{2.035182in}{1.818627in}}%
\pgfpathlineto{\pgfqpoint{2.035980in}{1.846540in}}%
\pgfpathlineto{\pgfqpoint{2.036414in}{1.824375in}}%
\pgfpathlineto{\pgfqpoint{2.036949in}{1.857777in}}%
\pgfpathlineto{\pgfqpoint{2.036954in}{1.857580in}}%
\pgfpathlineto{\pgfqpoint{2.038098in}{1.886002in}}%
\pgfpathlineto{\pgfqpoint{2.037100in}{1.851441in}}%
\pgfpathlineto{\pgfqpoint{2.038132in}{1.883507in}}%
\pgfpathlineto{\pgfqpoint{2.039130in}{1.862962in}}%
\pgfpathlineto{\pgfqpoint{2.038225in}{1.893138in}}%
\pgfpathlineto{\pgfqpoint{2.039276in}{1.872204in}}%
\pgfpathlineto{\pgfqpoint{2.040552in}{1.909168in}}%
\pgfpathlineto{\pgfqpoint{2.040557in}{1.908500in}}%
\pgfpathlineto{\pgfqpoint{2.040980in}{1.893112in}}%
\pgfpathlineto{\pgfqpoint{2.041608in}{1.931421in}}%
\pgfpathlineto{\pgfqpoint{2.041681in}{1.938462in}}%
\pgfpathlineto{\pgfqpoint{2.042110in}{1.901687in}}%
\pgfpathlineto{\pgfqpoint{2.042645in}{1.907280in}}%
\pgfpathlineto{\pgfqpoint{2.042670in}{1.914860in}}%
\pgfpathlineto{\pgfqpoint{2.042684in}{1.913748in}}%
\pgfpathlineto{\pgfqpoint{2.043565in}{1.973078in}}%
\pgfpathlineto{\pgfqpoint{2.043935in}{1.968354in}}%
\pgfpathlineto{\pgfqpoint{2.044870in}{1.989896in}}%
\pgfpathlineto{\pgfqpoint{2.044349in}{1.952232in}}%
\pgfpathlineto{\pgfqpoint{2.044977in}{1.965438in}}%
\pgfpathlineto{\pgfqpoint{2.045498in}{1.947407in}}%
\pgfpathlineto{\pgfqpoint{2.045045in}{1.966589in}}%
\pgfpathlineto{\pgfqpoint{2.046092in}{1.960503in}}%
\pgfpathlineto{\pgfqpoint{2.047061in}{1.930859in}}%
\pgfpathlineto{\pgfqpoint{2.046194in}{1.962458in}}%
\pgfpathlineto{\pgfqpoint{2.047265in}{1.949183in}}%
\pgfpathlineto{\pgfqpoint{2.047470in}{1.973042in}}%
\pgfpathlineto{\pgfqpoint{2.048410in}{1.961981in}}%
\pgfpathlineto{\pgfqpoint{2.049349in}{1.940047in}}%
\pgfpathlineto{\pgfqpoint{2.048697in}{1.966276in}}%
\pgfpathlineto{\pgfqpoint{2.049554in}{1.949488in}}%
\pgfpathlineto{\pgfqpoint{2.050016in}{1.994407in}}%
\pgfpathlineto{\pgfqpoint{2.049602in}{1.948907in}}%
\pgfpathlineto{\pgfqpoint{2.050673in}{1.956054in}}%
\pgfpathlineto{\pgfqpoint{2.050883in}{1.937528in}}%
\pgfpathlineto{\pgfqpoint{2.051691in}{1.987626in}}%
\pgfpathlineto{\pgfqpoint{2.051696in}{1.986152in}}%
\pgfpathlineto{\pgfqpoint{2.052475in}{2.039150in}}%
\pgfpathlineto{\pgfqpoint{2.052645in}{2.034173in}}%
\pgfpathlineto{\pgfqpoint{2.053400in}{2.068536in}}%
\pgfpathlineto{\pgfqpoint{2.053702in}{2.031933in}}%
\pgfpathlineto{\pgfqpoint{2.053828in}{2.039763in}}%
\pgfpathlineto{\pgfqpoint{2.053833in}{2.039544in}}%
\pgfpathlineto{\pgfqpoint{2.054179in}{2.065928in}}%
\pgfpathlineto{\pgfqpoint{2.054359in}{2.063925in}}%
\pgfpathlineto{\pgfqpoint{2.054495in}{2.078854in}}%
\pgfpathlineto{\pgfqpoint{2.054977in}{2.028130in}}%
\pgfpathlineto{\pgfqpoint{2.055454in}{2.059290in}}%
\pgfpathlineto{\pgfqpoint{2.055474in}{2.053478in}}%
\pgfpathlineto{\pgfqpoint{2.055810in}{2.076614in}}%
\pgfpathlineto{\pgfqpoint{2.056379in}{2.070911in}}%
\pgfpathlineto{\pgfqpoint{2.057329in}{2.096269in}}%
\pgfpathlineto{\pgfqpoint{2.056433in}{2.070782in}}%
\pgfpathlineto{\pgfqpoint{2.057518in}{2.084212in}}%
\pgfpathlineto{\pgfqpoint{2.057742in}{2.072931in}}%
\pgfpathlineto{\pgfqpoint{2.058054in}{2.098839in}}%
\pgfpathlineto{\pgfqpoint{2.058638in}{2.081896in}}%
\pgfpathlineto{\pgfqpoint{2.058682in}{2.086841in}}%
\pgfpathlineto{\pgfqpoint{2.059592in}{2.057413in}}%
\pgfpathlineto{\pgfqpoint{2.059787in}{2.047245in}}%
\pgfpathlineto{\pgfqpoint{2.060586in}{2.094659in}}%
\pgfpathlineto{\pgfqpoint{2.060615in}{2.092037in}}%
\pgfpathlineto{\pgfqpoint{2.060912in}{2.123586in}}%
\pgfpathlineto{\pgfqpoint{2.061754in}{2.103194in}}%
\pgfpathlineto{\pgfqpoint{2.062942in}{2.070674in}}%
\pgfpathlineto{\pgfqpoint{2.062387in}{2.119562in}}%
\pgfpathlineto{\pgfqpoint{2.062947in}{2.071383in}}%
\pgfpathlineto{\pgfqpoint{2.063779in}{2.088270in}}%
\pgfpathlineto{\pgfqpoint{2.063307in}{2.048536in}}%
\pgfpathlineto{\pgfqpoint{2.064145in}{2.077436in}}%
\pgfpathlineto{\pgfqpoint{2.064441in}{2.067130in}}%
\pgfpathlineto{\pgfqpoint{2.064709in}{2.098686in}}%
\pgfpathlineto{\pgfqpoint{2.065259in}{2.073789in}}%
\pgfpathlineto{\pgfqpoint{2.065362in}{2.077448in}}%
\pgfpathlineto{\pgfqpoint{2.066048in}{2.044898in}}%
\pgfpathlineto{\pgfqpoint{2.066189in}{2.055141in}}%
\pgfpathlineto{\pgfqpoint{2.066700in}{2.029615in}}%
\pgfpathlineto{\pgfqpoint{2.067241in}{2.083592in}}%
\pgfpathlineto{\pgfqpoint{2.068098in}{2.099623in}}%
\pgfpathlineto{\pgfqpoint{2.068229in}{2.077075in}}%
\pgfpathlineto{\pgfqpoint{2.068361in}{2.091242in}}%
\pgfpathlineto{\pgfqpoint{2.069549in}{2.030207in}}%
\pgfpathlineto{\pgfqpoint{2.069563in}{2.032242in}}%
\pgfpathlineto{\pgfqpoint{2.070074in}{2.048266in}}%
\pgfpathlineto{\pgfqpoint{2.070386in}{2.004483in}}%
\pgfpathlineto{\pgfqpoint{2.070659in}{2.024810in}}%
\pgfpathlineto{\pgfqpoint{2.070814in}{2.037582in}}%
\pgfpathlineto{\pgfqpoint{2.071155in}{2.008225in}}%
\pgfpathlineto{\pgfqpoint{2.071550in}{2.018543in}}%
\pgfpathlineto{\pgfqpoint{2.072109in}{1.992996in}}%
\pgfpathlineto{\pgfqpoint{2.072582in}{2.038094in}}%
\pgfpathlineto{\pgfqpoint{2.072640in}{2.028569in}}%
\pgfpathlineto{\pgfqpoint{2.072913in}{2.054399in}}%
\pgfpathlineto{\pgfqpoint{2.073463in}{2.011887in}}%
\pgfpathlineto{\pgfqpoint{2.073682in}{2.013314in}}%
\pgfpathlineto{\pgfqpoint{2.073833in}{1.991776in}}%
\pgfpathlineto{\pgfqpoint{2.074719in}{2.023389in}}%
\pgfpathlineto{\pgfqpoint{2.074782in}{2.019257in}}%
\pgfpathlineto{\pgfqpoint{2.075430in}{1.996532in}}%
\pgfpathlineto{\pgfqpoint{2.075050in}{2.022096in}}%
\pgfpathlineto{\pgfqpoint{2.075907in}{2.015568in}}%
\pgfpathlineto{\pgfqpoint{2.076019in}{2.022602in}}%
\pgfpathlineto{\pgfqpoint{2.076204in}{2.003668in}}%
\pgfpathlineto{\pgfqpoint{2.077022in}{2.018858in}}%
\pgfpathlineto{\pgfqpoint{2.078210in}{1.988049in}}%
\pgfpathlineto{\pgfqpoint{2.077411in}{2.025918in}}%
\pgfpathlineto{\pgfqpoint{2.078214in}{1.988807in}}%
\pgfpathlineto{\pgfqpoint{2.078726in}{2.013194in}}%
\pgfpathlineto{\pgfqpoint{2.078443in}{1.983647in}}%
\pgfpathlineto{\pgfqpoint{2.079344in}{2.003523in}}%
\pgfpathlineto{\pgfqpoint{2.080118in}{1.983987in}}%
\pgfpathlineto{\pgfqpoint{2.079621in}{2.025210in}}%
\pgfpathlineto{\pgfqpoint{2.080478in}{1.997695in}}%
\pgfpathlineto{\pgfqpoint{2.080936in}{2.019024in}}%
\pgfpathlineto{\pgfqpoint{2.081136in}{1.983911in}}%
\pgfpathlineto{\pgfqpoint{2.081583in}{1.998037in}}%
\pgfpathlineto{\pgfqpoint{2.081657in}{1.979864in}}%
\pgfpathlineto{\pgfqpoint{2.081988in}{2.025696in}}%
\pgfpathlineto{\pgfqpoint{2.082694in}{1.994795in}}%
\pgfpathlineto{\pgfqpoint{2.083911in}{2.026197in}}%
\pgfpathlineto{\pgfqpoint{2.083073in}{1.979772in}}%
\pgfpathlineto{\pgfqpoint{2.083930in}{2.023150in}}%
\pgfpathlineto{\pgfqpoint{2.084110in}{2.009204in}}%
\pgfpathlineto{\pgfqpoint{2.084626in}{2.032003in}}%
\pgfpathlineto{\pgfqpoint{2.085011in}{2.029813in}}%
\pgfpathlineto{\pgfqpoint{2.085902in}{2.071311in}}%
\pgfpathlineto{\pgfqpoint{2.085118in}{2.028388in}}%
\pgfpathlineto{\pgfqpoint{2.086165in}{2.063905in}}%
\pgfpathlineto{\pgfqpoint{2.086403in}{2.030178in}}%
\pgfpathlineto{\pgfqpoint{2.087036in}{2.075157in}}%
\pgfpathlineto{\pgfqpoint{2.087246in}{2.071560in}}%
\pgfpathlineto{\pgfqpoint{2.087684in}{2.089112in}}%
\pgfpathlineto{\pgfqpoint{2.088210in}{2.048819in}}%
\pgfpathlineto{\pgfqpoint{2.088248in}{2.051473in}}%
\pgfpathlineto{\pgfqpoint{2.089237in}{2.042146in}}%
\pgfpathlineto{\pgfqpoint{2.088813in}{2.064900in}}%
\pgfpathlineto{\pgfqpoint{2.089378in}{2.043186in}}%
\pgfpathlineto{\pgfqpoint{2.090225in}{2.077352in}}%
\pgfpathlineto{\pgfqpoint{2.089495in}{2.034655in}}%
\pgfpathlineto{\pgfqpoint{2.090624in}{2.070798in}}%
\pgfpathlineto{\pgfqpoint{2.090790in}{2.068635in}}%
\pgfpathlineto{\pgfqpoint{2.091637in}{2.103148in}}%
\pgfpathlineto{\pgfqpoint{2.091661in}{2.101261in}}%
\pgfpathlineto{\pgfqpoint{2.091764in}{2.113120in}}%
\pgfpathlineto{\pgfqpoint{2.091997in}{2.090687in}}%
\pgfpathlineto{\pgfqpoint{2.092606in}{2.094935in}}%
\pgfpathlineto{\pgfqpoint{2.092620in}{2.092685in}}%
\pgfpathlineto{\pgfqpoint{2.093073in}{2.143279in}}%
\pgfpathlineto{\pgfqpoint{2.093443in}{2.143011in}}%
\pgfpathlineto{\pgfqpoint{2.094612in}{2.172696in}}%
\pgfpathlineto{\pgfqpoint{2.093657in}{2.132096in}}%
\pgfpathlineto{\pgfqpoint{2.094626in}{2.170393in}}%
\pgfpathlineto{\pgfqpoint{2.094660in}{2.167267in}}%
\pgfpathlineto{\pgfqpoint{2.095288in}{2.191335in}}%
\pgfpathlineto{\pgfqpoint{2.095293in}{2.190319in}}%
\pgfpathlineto{\pgfqpoint{2.095425in}{2.207977in}}%
\pgfpathlineto{\pgfqpoint{2.095843in}{2.175409in}}%
\pgfpathlineto{\pgfqpoint{2.096398in}{2.188057in}}%
\pgfpathlineto{\pgfqpoint{2.096905in}{2.193778in}}%
\pgfpathlineto{\pgfqpoint{2.096525in}{2.174552in}}%
\pgfpathlineto{\pgfqpoint{2.097401in}{2.175111in}}%
\pgfpathlineto{\pgfqpoint{2.098419in}{2.157618in}}%
\pgfpathlineto{\pgfqpoint{2.097820in}{2.189048in}}%
\pgfpathlineto{\pgfqpoint{2.098526in}{2.162788in}}%
\pgfpathlineto{\pgfqpoint{2.098701in}{2.165199in}}%
\pgfpathlineto{\pgfqpoint{2.099052in}{2.138229in}}%
\pgfpathlineto{\pgfqpoint{2.099582in}{2.147019in}}%
\pgfpathlineto{\pgfqpoint{2.099977in}{2.125979in}}%
\pgfpathlineto{\pgfqpoint{2.100595in}{2.166699in}}%
\pgfpathlineto{\pgfqpoint{2.100634in}{2.163928in}}%
\pgfpathlineto{\pgfqpoint{2.100649in}{2.167763in}}%
\pgfpathlineto{\pgfqpoint{2.100955in}{2.144478in}}%
\pgfpathlineto{\pgfqpoint{2.101398in}{2.135003in}}%
\pgfpathlineto{\pgfqpoint{2.101817in}{2.168560in}}%
\pgfpathlineto{\pgfqpoint{2.102051in}{2.146013in}}%
\pgfpathlineto{\pgfqpoint{2.102231in}{2.151775in}}%
\pgfpathlineto{\pgfqpoint{2.102946in}{2.090313in}}%
\pgfpathlineto{\pgfqpoint{2.102995in}{2.094425in}}%
\pgfpathlineto{\pgfqpoint{2.103721in}{2.115618in}}%
\pgfpathlineto{\pgfqpoint{2.104173in}{2.077340in}}%
\pgfpathlineto{\pgfqpoint{2.104928in}{2.130537in}}%
\pgfpathlineto{\pgfqpoint{2.104324in}{2.071444in}}%
\pgfpathlineto{\pgfqpoint{2.105454in}{2.118641in}}%
\pgfpathlineto{\pgfqpoint{2.106603in}{2.087794in}}%
\pgfpathlineto{\pgfqpoint{2.105668in}{2.134858in}}%
\pgfpathlineto{\pgfqpoint{2.106632in}{2.091005in}}%
\pgfpathlineto{\pgfqpoint{2.107528in}{2.127436in}}%
\pgfpathlineto{\pgfqpoint{2.106739in}{2.087127in}}%
\pgfpathlineto{\pgfqpoint{2.107810in}{2.114939in}}%
\pgfpathlineto{\pgfqpoint{2.108910in}{2.082454in}}%
\pgfpathlineto{\pgfqpoint{2.107951in}{2.133088in}}%
\pgfpathlineto{\pgfqpoint{2.108964in}{2.091176in}}%
\pgfpathlineto{\pgfqpoint{2.109022in}{2.081700in}}%
\pgfpathlineto{\pgfqpoint{2.109714in}{2.106296in}}%
\pgfpathlineto{\pgfqpoint{2.110045in}{2.096393in}}%
\pgfpathlineto{\pgfqpoint{2.110887in}{2.122267in}}%
\pgfpathlineto{\pgfqpoint{2.110527in}{2.094746in}}%
\pgfpathlineto{\pgfqpoint{2.111174in}{2.118825in}}%
\pgfpathlineto{\pgfqpoint{2.112372in}{2.053708in}}%
\pgfpathlineto{\pgfqpoint{2.112406in}{2.057987in}}%
\pgfpathlineto{\pgfqpoint{2.113224in}{2.083797in}}%
\pgfpathlineto{\pgfqpoint{2.112757in}{2.054398in}}%
\pgfpathlineto{\pgfqpoint{2.113516in}{2.058291in}}%
\pgfpathlineto{\pgfqpoint{2.113555in}{2.051082in}}%
\pgfpathlineto{\pgfqpoint{2.113774in}{2.072665in}}%
\pgfpathlineto{\pgfqpoint{2.114047in}{2.069951in}}%
\pgfpathlineto{\pgfqpoint{2.115254in}{2.149225in}}%
\pgfpathlineto{\pgfqpoint{2.114076in}{2.065819in}}%
\pgfpathlineto{\pgfqpoint{2.115293in}{2.144673in}}%
\pgfpathlineto{\pgfqpoint{2.115454in}{2.134362in}}%
\pgfpathlineto{\pgfqpoint{2.116252in}{2.198371in}}%
\pgfpathlineto{\pgfqpoint{2.116257in}{2.197045in}}%
\pgfpathlineto{\pgfqpoint{2.116909in}{2.251642in}}%
\pgfpathlineto{\pgfqpoint{2.117406in}{2.212591in}}%
\pgfpathlineto{\pgfqpoint{2.117883in}{2.239572in}}%
\pgfpathlineto{\pgfqpoint{2.118492in}{2.208398in}}%
\pgfpathlineto{\pgfqpoint{2.118506in}{2.212193in}}%
\pgfpathlineto{\pgfqpoint{2.118613in}{2.197102in}}%
\pgfpathlineto{\pgfqpoint{2.118935in}{2.246889in}}%
\pgfpathlineto{\pgfqpoint{2.119543in}{2.238536in}}%
\pgfpathlineto{\pgfqpoint{2.119748in}{2.247353in}}%
\pgfpathlineto{\pgfqpoint{2.120575in}{2.204115in}}%
\pgfpathlineto{\pgfqpoint{2.120605in}{2.208398in}}%
\pgfpathlineto{\pgfqpoint{2.120945in}{2.199067in}}%
\pgfpathlineto{\pgfqpoint{2.121603in}{2.220457in}}%
\pgfpathlineto{\pgfqpoint{2.121719in}{2.205736in}}%
\pgfpathlineto{\pgfqpoint{2.121929in}{2.195205in}}%
\pgfpathlineto{\pgfqpoint{2.122377in}{2.232554in}}%
\pgfpathlineto{\pgfqpoint{2.122615in}{2.224758in}}%
\pgfpathlineto{\pgfqpoint{2.122693in}{2.234753in}}%
\pgfpathlineto{\pgfqpoint{2.123170in}{2.198937in}}%
\pgfpathlineto{\pgfqpoint{2.123672in}{2.205618in}}%
\pgfpathlineto{\pgfqpoint{2.123779in}{2.203696in}}%
\pgfpathlineto{\pgfqpoint{2.124567in}{2.267578in}}%
\pgfpathlineto{\pgfqpoint{2.124582in}{2.265707in}}%
\pgfpathlineto{\pgfqpoint{2.124611in}{2.278000in}}%
\pgfpathlineto{\pgfqpoint{2.125590in}{2.226548in}}%
\pgfpathlineto{\pgfqpoint{2.125624in}{2.231349in}}%
\pgfpathlineto{\pgfqpoint{2.126096in}{2.189325in}}%
\pgfpathlineto{\pgfqpoint{2.125648in}{2.235879in}}%
\pgfpathlineto{\pgfqpoint{2.126758in}{2.215793in}}%
\pgfpathlineto{\pgfqpoint{2.127683in}{2.234617in}}%
\pgfpathlineto{\pgfqpoint{2.127849in}{2.210554in}}%
\pgfpathlineto{\pgfqpoint{2.127854in}{2.212142in}}%
\pgfpathlineto{\pgfqpoint{2.128010in}{2.197248in}}%
\pgfpathlineto{\pgfqpoint{2.128759in}{2.236798in}}%
\pgfpathlineto{\pgfqpoint{2.128769in}{2.233408in}}%
\pgfpathlineto{\pgfqpoint{2.129903in}{2.272618in}}%
\pgfpathlineto{\pgfqpoint{2.129918in}{2.265159in}}%
\pgfpathlineto{\pgfqpoint{2.130575in}{2.249830in}}%
\pgfpathlineto{\pgfqpoint{2.131082in}{2.288550in}}%
\pgfpathlineto{\pgfqpoint{2.131885in}{2.266315in}}%
\pgfpathlineto{\pgfqpoint{2.131510in}{2.300228in}}%
\pgfpathlineto{\pgfqpoint{2.132177in}{2.291721in}}%
\pgfpathlineto{\pgfqpoint{2.132766in}{2.302704in}}%
\pgfpathlineto{\pgfqpoint{2.133029in}{2.259414in}}%
\pgfpathlineto{\pgfqpoint{2.133044in}{2.261741in}}%
\pgfpathlineto{\pgfqpoint{2.133881in}{2.216522in}}%
\pgfpathlineto{\pgfqpoint{2.133214in}{2.266926in}}%
\pgfpathlineto{\pgfqpoint{2.134231in}{2.222308in}}%
\pgfpathlineto{\pgfqpoint{2.134353in}{2.224383in}}%
\pgfpathlineto{\pgfqpoint{2.134528in}{2.205927in}}%
\pgfpathlineto{\pgfqpoint{2.134533in}{2.203543in}}%
\pgfpathlineto{\pgfqpoint{2.135049in}{2.235624in}}%
\pgfpathlineto{\pgfqpoint{2.135590in}{2.228310in}}%
\pgfpathlineto{\pgfqpoint{2.136734in}{2.246877in}}%
\pgfpathlineto{\pgfqpoint{2.136437in}{2.217299in}}%
\pgfpathlineto{\pgfqpoint{2.136739in}{2.245619in}}%
\pgfpathlineto{\pgfqpoint{2.137226in}{2.258844in}}%
\pgfpathlineto{\pgfqpoint{2.137425in}{2.231647in}}%
\pgfpathlineto{\pgfqpoint{2.137527in}{2.231909in}}%
\pgfpathlineto{\pgfqpoint{2.138506in}{2.219096in}}%
\pgfpathlineto{\pgfqpoint{2.137829in}{2.256290in}}%
\pgfpathlineto{\pgfqpoint{2.138628in}{2.234112in}}%
\pgfpathlineto{\pgfqpoint{2.138876in}{2.250230in}}%
\pgfpathlineto{\pgfqpoint{2.139665in}{2.210865in}}%
\pgfpathlineto{\pgfqpoint{2.140711in}{2.182549in}}%
\pgfpathlineto{\pgfqpoint{2.139679in}{2.211073in}}%
\pgfpathlineto{\pgfqpoint{2.140848in}{2.190200in}}%
\pgfpathlineto{\pgfqpoint{2.141135in}{2.218810in}}%
\pgfpathlineto{\pgfqpoint{2.141817in}{2.187350in}}%
\pgfpathlineto{\pgfqpoint{2.141890in}{2.187376in}}%
\pgfpathlineto{\pgfqpoint{2.141919in}{2.182553in}}%
\pgfpathlineto{\pgfqpoint{2.142415in}{2.223077in}}%
\pgfpathlineto{\pgfqpoint{2.142907in}{2.209299in}}%
\pgfpathlineto{\pgfqpoint{2.143024in}{2.227076in}}%
\pgfpathlineto{\pgfqpoint{2.143676in}{2.196973in}}%
\pgfpathlineto{\pgfqpoint{2.143964in}{2.198584in}}%
\pgfpathlineto{\pgfqpoint{2.144027in}{2.192921in}}%
\pgfpathlineto{\pgfqpoint{2.144592in}{2.222301in}}%
\pgfpathlineto{\pgfqpoint{2.144986in}{2.209534in}}%
\pgfpathlineto{\pgfqpoint{2.145015in}{2.214260in}}%
\pgfpathlineto{\pgfqpoint{2.145848in}{2.179185in}}%
\pgfpathlineto{\pgfqpoint{2.145940in}{2.188168in}}%
\pgfpathlineto{\pgfqpoint{2.145969in}{2.186152in}}%
\pgfpathlineto{\pgfqpoint{2.146529in}{2.222379in}}%
\pgfpathlineto{\pgfqpoint{2.146943in}{2.207499in}}%
\pgfpathlineto{\pgfqpoint{2.146977in}{2.212683in}}%
\pgfpathlineto{\pgfqpoint{2.147936in}{2.173403in}}%
\pgfpathlineto{\pgfqpoint{2.147980in}{2.173882in}}%
\pgfpathlineto{\pgfqpoint{2.148837in}{2.142426in}}%
\pgfpathlineto{\pgfqpoint{2.148141in}{2.180109in}}%
\pgfpathlineto{\pgfqpoint{2.149090in}{2.172813in}}%
\pgfpathlineto{\pgfqpoint{2.150229in}{2.203593in}}%
\pgfpathlineto{\pgfqpoint{2.149158in}{2.162478in}}%
\pgfpathlineto{\pgfqpoint{2.150375in}{2.198787in}}%
\pgfpathlineto{\pgfqpoint{2.151378in}{2.172094in}}%
\pgfpathlineto{\pgfqpoint{2.150531in}{2.200995in}}%
\pgfpathlineto{\pgfqpoint{2.151500in}{2.182935in}}%
\pgfpathlineto{\pgfqpoint{2.151549in}{2.172074in}}%
\pgfpathlineto{\pgfqpoint{2.152075in}{2.217627in}}%
\pgfpathlineto{\pgfqpoint{2.152430in}{2.198838in}}%
\pgfpathlineto{\pgfqpoint{2.152844in}{2.234060in}}%
\pgfpathlineto{\pgfqpoint{2.153559in}{2.220481in}}%
\pgfpathlineto{\pgfqpoint{2.153564in}{2.220458in}}%
\pgfpathlineto{\pgfqpoint{2.153579in}{2.225363in}}%
\pgfpathlineto{\pgfqpoint{2.154460in}{2.252617in}}%
\pgfpathlineto{\pgfqpoint{2.154718in}{2.241179in}}%
\pgfpathlineto{\pgfqpoint{2.155487in}{2.219403in}}%
\pgfpathlineto{\pgfqpoint{2.155142in}{2.243875in}}%
\pgfpathlineto{\pgfqpoint{2.155848in}{2.234421in}}%
\pgfpathlineto{\pgfqpoint{2.155950in}{2.245322in}}%
\pgfpathlineto{\pgfqpoint{2.156753in}{2.206078in}}%
\pgfpathlineto{\pgfqpoint{2.156773in}{2.209842in}}%
\pgfpathlineto{\pgfqpoint{2.156782in}{2.206983in}}%
\pgfpathlineto{\pgfqpoint{2.157527in}{2.245794in}}%
\pgfpathlineto{\pgfqpoint{2.157849in}{2.227354in}}%
\pgfpathlineto{\pgfqpoint{2.158063in}{2.218315in}}%
\pgfpathlineto{\pgfqpoint{2.158496in}{2.247577in}}%
\pgfpathlineto{\pgfqpoint{2.158934in}{2.235648in}}%
\pgfpathlineto{\pgfqpoint{2.159893in}{2.255556in}}%
\pgfpathlineto{\pgfqpoint{2.159348in}{2.228818in}}%
\pgfpathlineto{\pgfqpoint{2.160181in}{2.244827in}}%
\pgfpathlineto{\pgfqpoint{2.160711in}{2.224333in}}%
\pgfpathlineto{\pgfqpoint{2.160307in}{2.260125in}}%
\pgfpathlineto{\pgfqpoint{2.161305in}{2.234545in}}%
\pgfpathlineto{\pgfqpoint{2.161461in}{2.241838in}}%
\pgfpathlineto{\pgfqpoint{2.161758in}{2.217118in}}%
\pgfpathlineto{\pgfqpoint{2.161919in}{2.219923in}}%
\pgfpathlineto{\pgfqpoint{2.162035in}{2.205337in}}%
\pgfpathlineto{\pgfqpoint{2.162878in}{2.262416in}}%
\pgfpathlineto{\pgfqpoint{2.162931in}{2.250973in}}%
\pgfpathlineto{\pgfqpoint{2.163486in}{2.278613in}}%
\pgfpathlineto{\pgfqpoint{2.164061in}{2.266842in}}%
\pgfpathlineto{\pgfqpoint{2.165132in}{2.239208in}}%
\pgfpathlineto{\pgfqpoint{2.164168in}{2.277529in}}%
\pgfpathlineto{\pgfqpoint{2.165205in}{2.248633in}}%
\pgfpathlineto{\pgfqpoint{2.166451in}{2.298773in}}%
\pgfpathlineto{\pgfqpoint{2.165380in}{2.241857in}}%
\pgfpathlineto{\pgfqpoint{2.166515in}{2.292625in}}%
\pgfpathlineto{\pgfqpoint{2.166563in}{2.283098in}}%
\pgfpathlineto{\pgfqpoint{2.167488in}{2.319884in}}%
\pgfpathlineto{\pgfqpoint{2.167615in}{2.299839in}}%
\pgfpathlineto{\pgfqpoint{2.168545in}{2.331881in}}%
\pgfpathlineto{\pgfqpoint{2.168301in}{2.296479in}}%
\pgfpathlineto{\pgfqpoint{2.168803in}{2.309597in}}%
\pgfpathlineto{\pgfqpoint{2.169158in}{2.286360in}}%
\pgfpathlineto{\pgfqpoint{2.169820in}{2.337843in}}%
\pgfpathlineto{\pgfqpoint{2.169869in}{2.331311in}}%
\pgfpathlineto{\pgfqpoint{2.169879in}{2.332845in}}%
\pgfpathlineto{\pgfqpoint{2.170249in}{2.311197in}}%
\pgfpathlineto{\pgfqpoint{2.170268in}{2.311988in}}%
\pgfpathlineto{\pgfqpoint{2.170404in}{2.305225in}}%
\pgfpathlineto{\pgfqpoint{2.171305in}{2.354779in}}%
\pgfpathlineto{\pgfqpoint{2.171310in}{2.357385in}}%
\pgfpathlineto{\pgfqpoint{2.172303in}{2.334687in}}%
\pgfpathlineto{\pgfqpoint{2.172386in}{2.341592in}}%
\pgfpathlineto{\pgfqpoint{2.172517in}{2.331525in}}%
\pgfpathlineto{\pgfqpoint{2.173155in}{2.363274in}}%
\pgfpathlineto{\pgfqpoint{2.173472in}{2.345508in}}%
\pgfpathlineto{\pgfqpoint{2.174650in}{2.374357in}}%
\pgfpathlineto{\pgfqpoint{2.173559in}{2.344513in}}%
\pgfpathlineto{\pgfqpoint{2.174669in}{2.368722in}}%
\pgfpathlineto{\pgfqpoint{2.175215in}{2.342837in}}%
\pgfpathlineto{\pgfqpoint{2.175789in}{2.360719in}}%
\pgfpathlineto{\pgfqpoint{2.175794in}{2.361096in}}%
\pgfpathlineto{\pgfqpoint{2.176295in}{2.317296in}}%
\pgfpathlineto{\pgfqpoint{2.176339in}{2.319182in}}%
\pgfpathlineto{\pgfqpoint{2.176398in}{2.311387in}}%
\pgfpathlineto{\pgfqpoint{2.176918in}{2.342749in}}%
\pgfpathlineto{\pgfqpoint{2.177420in}{2.328979in}}%
\pgfpathlineto{\pgfqpoint{2.178282in}{2.335576in}}%
\pgfpathlineto{\pgfqpoint{2.177970in}{2.310283in}}%
\pgfpathlineto{\pgfqpoint{2.178433in}{2.325091in}}%
\pgfpathlineto{\pgfqpoint{2.179504in}{2.298856in}}%
\pgfpathlineto{\pgfqpoint{2.178671in}{2.343940in}}%
\pgfpathlineto{\pgfqpoint{2.179567in}{2.309312in}}%
\pgfpathlineto{\pgfqpoint{2.179952in}{2.318639in}}%
\pgfpathlineto{\pgfqpoint{2.180531in}{2.296024in}}%
\pgfpathlineto{\pgfqpoint{2.180667in}{2.300966in}}%
\pgfpathlineto{\pgfqpoint{2.181003in}{2.312040in}}%
\pgfpathlineto{\pgfqpoint{2.181738in}{2.288526in}}%
\pgfpathlineto{\pgfqpoint{2.181758in}{2.291829in}}%
\pgfpathlineto{\pgfqpoint{2.181768in}{2.293461in}}%
\pgfpathlineto{\pgfqpoint{2.182147in}{2.257438in}}%
\pgfpathlineto{\pgfqpoint{2.183019in}{2.326010in}}%
\pgfpathlineto{\pgfqpoint{2.183759in}{2.287583in}}%
\pgfpathlineto{\pgfqpoint{2.184158in}{2.303038in}}%
\pgfpathlineto{\pgfqpoint{2.184883in}{2.335613in}}%
\pgfpathlineto{\pgfqpoint{2.184489in}{2.301334in}}%
\pgfpathlineto{\pgfqpoint{2.185326in}{2.331457in}}%
\pgfpathlineto{\pgfqpoint{2.185370in}{2.328547in}}%
\pgfpathlineto{\pgfqpoint{2.185794in}{2.372363in}}%
\pgfpathlineto{\pgfqpoint{2.186110in}{2.363939in}}%
\pgfpathlineto{\pgfqpoint{2.186738in}{2.380845in}}%
\pgfpathlineto{\pgfqpoint{2.186217in}{2.359894in}}%
\pgfpathlineto{\pgfqpoint{2.187220in}{2.365794in}}%
\pgfpathlineto{\pgfqpoint{2.187313in}{2.354312in}}%
\pgfpathlineto{\pgfqpoint{2.188077in}{2.387716in}}%
\pgfpathlineto{\pgfqpoint{2.188311in}{2.371195in}}%
\pgfpathlineto{\pgfqpoint{2.189104in}{2.423753in}}%
\pgfpathlineto{\pgfqpoint{2.189469in}{2.380271in}}%
\pgfpathlineto{\pgfqpoint{2.189796in}{2.363589in}}%
\pgfpathlineto{\pgfqpoint{2.190506in}{2.401063in}}%
\pgfpathlineto{\pgfqpoint{2.190521in}{2.398678in}}%
\pgfpathlineto{\pgfqpoint{2.191047in}{2.425387in}}%
\pgfpathlineto{\pgfqpoint{2.190536in}{2.397026in}}%
\pgfpathlineto{\pgfqpoint{2.191928in}{2.411564in}}%
\pgfpathlineto{\pgfqpoint{2.192668in}{2.397781in}}%
\pgfpathlineto{\pgfqpoint{2.192342in}{2.436296in}}%
\pgfpathlineto{\pgfqpoint{2.192921in}{2.431108in}}%
\pgfpathlineto{\pgfqpoint{2.194523in}{2.486768in}}%
\pgfpathlineto{\pgfqpoint{2.192941in}{2.429775in}}%
\pgfpathlineto{\pgfqpoint{2.194581in}{2.484310in}}%
\pgfpathlineto{\pgfqpoint{2.194586in}{2.481741in}}%
\pgfpathlineto{\pgfqpoint{2.195317in}{2.512465in}}%
\pgfpathlineto{\pgfqpoint{2.195652in}{2.501766in}}%
\pgfpathlineto{\pgfqpoint{2.195988in}{2.541475in}}%
\pgfpathlineto{\pgfqpoint{2.196680in}{2.487179in}}%
\pgfpathlineto{\pgfqpoint{2.196733in}{2.490262in}}%
\pgfpathlineto{\pgfqpoint{2.196831in}{2.491941in}}%
\pgfpathlineto{\pgfqpoint{2.197663in}{2.464998in}}%
\pgfpathlineto{\pgfqpoint{2.197751in}{2.472630in}}%
\pgfpathlineto{\pgfqpoint{2.197756in}{2.472334in}}%
\pgfpathlineto{\pgfqpoint{2.198340in}{2.500861in}}%
\pgfpathlineto{\pgfqpoint{2.198471in}{2.492820in}}%
\pgfpathlineto{\pgfqpoint{2.199606in}{2.521783in}}%
\pgfpathlineto{\pgfqpoint{2.198744in}{2.487872in}}%
\pgfpathlineto{\pgfqpoint{2.199615in}{2.521604in}}%
\pgfpathlineto{\pgfqpoint{2.200701in}{2.500539in}}%
\pgfpathlineto{\pgfqpoint{2.200107in}{2.557924in}}%
\pgfpathlineto{\pgfqpoint{2.200745in}{2.510292in}}%
\pgfpathlineto{\pgfqpoint{2.201772in}{2.561155in}}%
\pgfpathlineto{\pgfqpoint{2.200945in}{2.505662in}}%
\pgfpathlineto{\pgfqpoint{2.201957in}{2.548111in}}%
\pgfpathlineto{\pgfqpoint{2.202308in}{2.540126in}}%
\pgfpathlineto{\pgfqpoint{2.203023in}{2.571193in}}%
\pgfpathlineto{\pgfqpoint{2.203033in}{2.567456in}}%
\pgfpathlineto{\pgfqpoint{2.203038in}{2.570657in}}%
\pgfpathlineto{\pgfqpoint{2.203973in}{2.533210in}}%
\pgfpathlineto{\pgfqpoint{2.204094in}{2.536139in}}%
\pgfpathlineto{\pgfqpoint{2.205093in}{2.494513in}}%
\pgfpathlineto{\pgfqpoint{2.204304in}{2.549362in}}%
\pgfpathlineto{\pgfqpoint{2.205258in}{2.496075in}}%
\pgfpathlineto{\pgfqpoint{2.205876in}{2.470095in}}%
\pgfpathlineto{\pgfqpoint{2.206256in}{2.520948in}}%
\pgfpathlineto{\pgfqpoint{2.206300in}{2.520916in}}%
\pgfpathlineto{\pgfqpoint{2.207434in}{2.555053in}}%
\pgfpathlineto{\pgfqpoint{2.206319in}{2.518740in}}%
\pgfpathlineto{\pgfqpoint{2.207463in}{2.547443in}}%
\pgfpathlineto{\pgfqpoint{2.208457in}{2.540293in}}%
\pgfpathlineto{\pgfqpoint{2.208018in}{2.573943in}}%
\pgfpathlineto{\pgfqpoint{2.208525in}{2.556506in}}%
\pgfpathlineto{\pgfqpoint{2.209664in}{2.586732in}}%
\pgfpathlineto{\pgfqpoint{2.209688in}{2.583479in}}%
\pgfpathlineto{\pgfqpoint{2.210068in}{2.568709in}}%
\pgfpathlineto{\pgfqpoint{2.210448in}{2.599180in}}%
\pgfpathlineto{\pgfqpoint{2.210750in}{2.595942in}}%
\pgfpathlineto{\pgfqpoint{2.211052in}{2.610233in}}%
\pgfpathlineto{\pgfqpoint{2.211577in}{2.578044in}}%
\pgfpathlineto{\pgfqpoint{2.211641in}{2.578834in}}%
\pgfpathlineto{\pgfqpoint{2.212011in}{2.527649in}}%
\pgfpathlineto{\pgfqpoint{2.211660in}{2.581770in}}%
\pgfpathlineto{\pgfqpoint{2.212848in}{2.558683in}}%
\pgfpathlineto{\pgfqpoint{2.212916in}{2.562839in}}%
\pgfpathlineto{\pgfqpoint{2.213685in}{2.504739in}}%
\pgfpathlineto{\pgfqpoint{2.213851in}{2.519641in}}%
\pgfpathlineto{\pgfqpoint{2.214343in}{2.498838in}}%
\pgfpathlineto{\pgfqpoint{2.214888in}{2.528890in}}%
\pgfpathlineto{\pgfqpoint{2.214946in}{2.522737in}}%
\pgfpathlineto{\pgfqpoint{2.215930in}{2.548631in}}%
\pgfpathlineto{\pgfqpoint{2.215414in}{2.518302in}}%
\pgfpathlineto{\pgfqpoint{2.216061in}{2.525938in}}%
\pgfpathlineto{\pgfqpoint{2.216991in}{2.557499in}}%
\pgfpathlineto{\pgfqpoint{2.216568in}{2.520220in}}%
\pgfpathlineto{\pgfqpoint{2.217225in}{2.530744in}}%
\pgfpathlineto{\pgfqpoint{2.217790in}{2.503216in}}%
\pgfpathlineto{\pgfqpoint{2.217327in}{2.536732in}}%
\pgfpathlineto{\pgfqpoint{2.218320in}{2.535425in}}%
\pgfpathlineto{\pgfqpoint{2.219240in}{2.559557in}}%
\pgfpathlineto{\pgfqpoint{2.218471in}{2.527583in}}%
\pgfpathlineto{\pgfqpoint{2.219440in}{2.539584in}}%
\pgfpathlineto{\pgfqpoint{2.219903in}{2.521164in}}%
\pgfpathlineto{\pgfqpoint{2.219547in}{2.548466in}}%
\pgfpathlineto{\pgfqpoint{2.220414in}{2.543979in}}%
\pgfpathlineto{\pgfqpoint{2.220725in}{2.565586in}}%
\pgfpathlineto{\pgfqpoint{2.221441in}{2.519390in}}%
\pgfpathlineto{\pgfqpoint{2.221495in}{2.528070in}}%
\pgfpathlineto{\pgfqpoint{2.222069in}{2.504413in}}%
\pgfpathlineto{\pgfqpoint{2.222566in}{2.541632in}}%
\pgfpathlineto{\pgfqpoint{2.222570in}{2.540609in}}%
\pgfpathlineto{\pgfqpoint{2.222853in}{2.560138in}}%
\pgfpathlineto{\pgfqpoint{2.223281in}{2.503214in}}%
\pgfpathlineto{\pgfqpoint{2.223676in}{2.540331in}}%
\pgfpathlineto{\pgfqpoint{2.223895in}{2.527170in}}%
\pgfpathlineto{\pgfqpoint{2.224474in}{2.568697in}}%
\pgfpathlineto{\pgfqpoint{2.224941in}{2.579281in}}%
\pgfpathlineto{\pgfqpoint{2.225535in}{2.549938in}}%
\pgfpathlineto{\pgfqpoint{2.225545in}{2.553974in}}%
\pgfpathlineto{\pgfqpoint{2.226602in}{2.483077in}}%
\pgfpathlineto{\pgfqpoint{2.226791in}{2.495014in}}%
\pgfpathlineto{\pgfqpoint{2.226903in}{2.473834in}}%
\pgfpathlineto{\pgfqpoint{2.227483in}{2.526218in}}%
\pgfpathlineto{\pgfqpoint{2.227634in}{2.523409in}}%
\pgfpathlineto{\pgfqpoint{2.228788in}{2.607558in}}%
\pgfpathlineto{\pgfqpoint{2.228870in}{2.603181in}}%
\pgfpathlineto{\pgfqpoint{2.229430in}{2.580474in}}%
\pgfpathlineto{\pgfqpoint{2.229849in}{2.625282in}}%
\pgfpathlineto{\pgfqpoint{2.229946in}{2.616225in}}%
\pgfpathlineto{\pgfqpoint{2.230005in}{2.621047in}}%
\pgfpathlineto{\pgfqpoint{2.230886in}{2.566059in}}%
\pgfpathlineto{\pgfqpoint{2.230930in}{2.578421in}}%
\pgfpathlineto{\pgfqpoint{2.230973in}{2.567967in}}%
\pgfpathlineto{\pgfqpoint{2.231378in}{2.598895in}}%
\pgfpathlineto{\pgfqpoint{2.232025in}{2.594193in}}%
\pgfpathlineto{\pgfqpoint{2.232122in}{2.603609in}}%
\pgfpathlineto{\pgfqpoint{2.232833in}{2.566598in}}%
\pgfpathlineto{\pgfqpoint{2.233082in}{2.572654in}}%
\pgfpathlineto{\pgfqpoint{2.233203in}{2.564296in}}%
\pgfpathlineto{\pgfqpoint{2.233417in}{2.595837in}}%
\pgfpathlineto{\pgfqpoint{2.234162in}{2.579596in}}%
\pgfpathlineto{\pgfqpoint{2.235175in}{2.585078in}}%
\pgfpathlineto{\pgfqpoint{2.234747in}{2.556763in}}%
\pgfpathlineto{\pgfqpoint{2.235263in}{2.577224in}}%
\pgfpathlineto{\pgfqpoint{2.236051in}{2.557624in}}%
\pgfpathlineto{\pgfqpoint{2.235443in}{2.586992in}}%
\pgfpathlineto{\pgfqpoint{2.236368in}{2.577261in}}%
\pgfpathlineto{\pgfqpoint{2.237156in}{2.586497in}}%
\pgfpathlineto{\pgfqpoint{2.237424in}{2.557425in}}%
\pgfpathlineto{\pgfqpoint{2.237429in}{2.558330in}}%
\pgfpathlineto{\pgfqpoint{2.237434in}{2.555654in}}%
\pgfpathlineto{\pgfqpoint{2.238145in}{2.588120in}}%
\pgfpathlineto{\pgfqpoint{2.238505in}{2.573556in}}%
\pgfpathlineto{\pgfqpoint{2.239167in}{2.629317in}}%
\pgfpathlineto{\pgfqpoint{2.238607in}{2.573036in}}%
\pgfpathlineto{\pgfqpoint{2.239649in}{2.593504in}}%
\pgfpathlineto{\pgfqpoint{2.240691in}{2.631944in}}%
\pgfpathlineto{\pgfqpoint{2.239854in}{2.585221in}}%
\pgfpathlineto{\pgfqpoint{2.240910in}{2.616412in}}%
\pgfpathlineto{\pgfqpoint{2.241879in}{2.587937in}}%
\pgfpathlineto{\pgfqpoint{2.241149in}{2.625685in}}%
\pgfpathlineto{\pgfqpoint{2.242064in}{2.594520in}}%
\pgfpathlineto{\pgfqpoint{2.242882in}{2.580656in}}%
\pgfpathlineto{\pgfqpoint{2.242283in}{2.606209in}}%
\pgfpathlineto{\pgfqpoint{2.243154in}{2.598362in}}%
\pgfpathlineto{\pgfqpoint{2.243403in}{2.591228in}}%
\pgfpathlineto{\pgfqpoint{2.244352in}{2.644176in}}%
\pgfpathlineto{\pgfqpoint{2.245443in}{2.672612in}}%
\pgfpathlineto{\pgfqpoint{2.245243in}{2.642548in}}%
\pgfpathlineto{\pgfqpoint{2.245589in}{2.663570in}}%
\pgfpathlineto{\pgfqpoint{2.245603in}{2.657128in}}%
\pgfpathlineto{\pgfqpoint{2.245988in}{2.693838in}}%
\pgfpathlineto{\pgfqpoint{2.246694in}{2.667341in}}%
\pgfpathlineto{\pgfqpoint{2.246801in}{2.678465in}}%
\pgfpathlineto{\pgfqpoint{2.247487in}{2.645275in}}%
\pgfpathlineto{\pgfqpoint{2.247901in}{2.616199in}}%
\pgfpathlineto{\pgfqpoint{2.248456in}{2.678015in}}%
\pgfpathlineto{\pgfqpoint{2.248461in}{2.677350in}}%
\pgfpathlineto{\pgfqpoint{2.249045in}{2.707593in}}%
\pgfpathlineto{\pgfqpoint{2.249449in}{2.670603in}}%
\pgfpathlineto{\pgfqpoint{2.249571in}{2.677018in}}%
\pgfpathlineto{\pgfqpoint{2.250521in}{2.628622in}}%
\pgfpathlineto{\pgfqpoint{2.249712in}{2.678833in}}%
\pgfpathlineto{\pgfqpoint{2.250730in}{2.654370in}}%
\pgfpathlineto{\pgfqpoint{2.250788in}{2.658446in}}%
\pgfpathlineto{\pgfqpoint{2.250856in}{2.650752in}}%
\pgfpathlineto{\pgfqpoint{2.251100in}{2.643005in}}%
\pgfpathlineto{\pgfqpoint{2.251889in}{2.696921in}}%
\pgfpathlineto{\pgfqpoint{2.251898in}{2.695985in}}%
\pgfpathlineto{\pgfqpoint{2.252916in}{2.716001in}}%
\pgfpathlineto{\pgfqpoint{2.252453in}{2.672850in}}%
\pgfpathlineto{\pgfqpoint{2.253076in}{2.714931in}}%
\pgfpathlineto{\pgfqpoint{2.254235in}{2.654524in}}%
\pgfpathlineto{\pgfqpoint{2.254260in}{2.655214in}}%
\pgfpathlineto{\pgfqpoint{2.254732in}{2.661707in}}%
\pgfpathlineto{\pgfqpoint{2.255058in}{2.628735in}}%
\pgfpathlineto{\pgfqpoint{2.255331in}{2.638753in}}%
\pgfpathlineto{\pgfqpoint{2.255374in}{2.633236in}}%
\pgfpathlineto{\pgfqpoint{2.255832in}{2.658400in}}%
\pgfpathlineto{\pgfqpoint{2.256411in}{2.645199in}}%
\pgfpathlineto{\pgfqpoint{2.256421in}{2.644291in}}%
\pgfpathlineto{\pgfqpoint{2.256470in}{2.649403in}}%
\pgfpathlineto{\pgfqpoint{2.257375in}{2.666004in}}%
\pgfpathlineto{\pgfqpoint{2.256884in}{2.635961in}}%
\pgfpathlineto{\pgfqpoint{2.257575in}{2.651741in}}%
\pgfpathlineto{\pgfqpoint{2.258330in}{2.628977in}}%
\pgfpathlineto{\pgfqpoint{2.258018in}{2.657513in}}%
\pgfpathlineto{\pgfqpoint{2.258719in}{2.632700in}}%
\pgfpathlineto{\pgfqpoint{2.259912in}{2.673831in}}%
\pgfpathlineto{\pgfqpoint{2.259045in}{2.620162in}}%
\pgfpathlineto{\pgfqpoint{2.259926in}{2.672091in}}%
\pgfpathlineto{\pgfqpoint{2.260929in}{2.630278in}}%
\pgfpathlineto{\pgfqpoint{2.260053in}{2.673499in}}%
\pgfpathlineto{\pgfqpoint{2.261144in}{2.638556in}}%
\pgfpathlineto{\pgfqpoint{2.261577in}{2.654717in}}%
\pgfpathlineto{\pgfqpoint{2.261952in}{2.627434in}}%
\pgfpathlineto{\pgfqpoint{2.262288in}{2.642803in}}%
\pgfpathlineto{\pgfqpoint{2.262828in}{2.617174in}}%
\pgfpathlineto{\pgfqpoint{2.262521in}{2.646651in}}%
\pgfpathlineto{\pgfqpoint{2.263461in}{2.636283in}}%
\pgfpathlineto{\pgfqpoint{2.263958in}{2.649603in}}%
\pgfpathlineto{\pgfqpoint{2.264420in}{2.614877in}}%
\pgfpathlineto{\pgfqpoint{2.264556in}{2.630836in}}%
\pgfpathlineto{\pgfqpoint{2.264595in}{2.634075in}}%
\pgfpathlineto{\pgfqpoint{2.264678in}{2.623080in}}%
\pgfpathlineto{\pgfqpoint{2.264761in}{2.624429in}}%
\pgfpathlineto{\pgfqpoint{2.264834in}{2.616124in}}%
\pgfpathlineto{\pgfqpoint{2.265462in}{2.652411in}}%
\pgfpathlineto{\pgfqpoint{2.265842in}{2.643543in}}%
\pgfpathlineto{\pgfqpoint{2.265851in}{2.646569in}}%
\pgfpathlineto{\pgfqpoint{2.266835in}{2.614816in}}%
\pgfpathlineto{\pgfqpoint{2.268037in}{2.542869in}}%
\pgfpathlineto{\pgfqpoint{2.266903in}{2.619962in}}%
\pgfpathlineto{\pgfqpoint{2.268071in}{2.548031in}}%
\pgfpathlineto{\pgfqpoint{2.268417in}{2.558789in}}%
\pgfpathlineto{\pgfqpoint{2.269143in}{2.513810in}}%
\pgfpathlineto{\pgfqpoint{2.269454in}{2.500097in}}%
\pgfpathlineto{\pgfqpoint{2.270097in}{2.528356in}}%
\pgfpathlineto{\pgfqpoint{2.270262in}{2.510297in}}%
\pgfpathlineto{\pgfqpoint{2.271095in}{2.542929in}}%
\pgfpathlineto{\pgfqpoint{2.270272in}{2.509364in}}%
\pgfpathlineto{\pgfqpoint{2.271411in}{2.528571in}}%
\pgfpathlineto{\pgfqpoint{2.271445in}{2.520320in}}%
\pgfpathlineto{\pgfqpoint{2.271947in}{2.545158in}}%
\pgfpathlineto{\pgfqpoint{2.272249in}{2.538650in}}%
\pgfpathlineto{\pgfqpoint{2.272740in}{2.562614in}}%
\pgfpathlineto{\pgfqpoint{2.273247in}{2.527059in}}%
\pgfpathlineto{\pgfqpoint{2.273329in}{2.530942in}}%
\pgfpathlineto{\pgfqpoint{2.274517in}{2.476886in}}%
\pgfpathlineto{\pgfqpoint{2.273539in}{2.546806in}}%
\pgfpathlineto{\pgfqpoint{2.274697in}{2.485221in}}%
\pgfpathlineto{\pgfqpoint{2.274873in}{2.479538in}}%
\pgfpathlineto{\pgfqpoint{2.275467in}{2.501209in}}%
\pgfpathlineto{\pgfqpoint{2.275584in}{2.492441in}}%
\pgfpathlineto{\pgfqpoint{2.275642in}{2.502416in}}%
\pgfpathlineto{\pgfqpoint{2.276650in}{2.465741in}}%
\pgfpathlineto{\pgfqpoint{2.276669in}{2.469438in}}%
\pgfpathlineto{\pgfqpoint{2.276771in}{2.450497in}}%
\pgfpathlineto{\pgfqpoint{2.277292in}{2.474145in}}%
\pgfpathlineto{\pgfqpoint{2.277891in}{2.459206in}}%
\pgfpathlineto{\pgfqpoint{2.277896in}{2.460776in}}%
\pgfpathlineto{\pgfqpoint{2.278617in}{2.410715in}}%
\pgfpathlineto{\pgfqpoint{2.278914in}{2.434412in}}%
\pgfpathlineto{\pgfqpoint{2.279250in}{2.405991in}}%
\pgfpathlineto{\pgfqpoint{2.279887in}{2.436825in}}%
\pgfpathlineto{\pgfqpoint{2.280072in}{2.408929in}}%
\pgfpathlineto{\pgfqpoint{2.280661in}{2.436681in}}%
\pgfpathlineto{\pgfqpoint{2.281148in}{2.404680in}}%
\pgfpathlineto{\pgfqpoint{2.281192in}{2.417092in}}%
\pgfpathlineto{\pgfqpoint{2.282112in}{2.434427in}}%
\pgfpathlineto{\pgfqpoint{2.281474in}{2.400285in}}%
\pgfpathlineto{\pgfqpoint{2.282322in}{2.429206in}}%
\pgfpathlineto{\pgfqpoint{2.283427in}{2.385137in}}%
\pgfpathlineto{\pgfqpoint{2.283480in}{2.386073in}}%
\pgfpathlineto{\pgfqpoint{2.284624in}{2.414186in}}%
\pgfpathlineto{\pgfqpoint{2.283836in}{2.361353in}}%
\pgfpathlineto{\pgfqpoint{2.284629in}{2.413800in}}%
\pgfpathlineto{\pgfqpoint{2.284668in}{2.411660in}}%
\pgfpathlineto{\pgfqpoint{2.285145in}{2.447607in}}%
\pgfpathlineto{\pgfqpoint{2.285588in}{2.429139in}}%
\pgfpathlineto{\pgfqpoint{2.285832in}{2.455906in}}%
\pgfpathlineto{\pgfqpoint{2.286640in}{2.413544in}}%
\pgfpathlineto{\pgfqpoint{2.286679in}{2.421719in}}%
\pgfpathlineto{\pgfqpoint{2.286693in}{2.417855in}}%
\pgfpathlineto{\pgfqpoint{2.287210in}{2.459968in}}%
\pgfpathlineto{\pgfqpoint{2.287774in}{2.429139in}}%
\pgfpathlineto{\pgfqpoint{2.287779in}{2.430899in}}%
\pgfpathlineto{\pgfqpoint{2.288281in}{2.407658in}}%
\pgfpathlineto{\pgfqpoint{2.288865in}{2.421791in}}%
\pgfpathlineto{\pgfqpoint{2.289921in}{2.461973in}}%
\pgfpathlineto{\pgfqpoint{2.288972in}{2.416670in}}%
\pgfpathlineto{\pgfqpoint{2.290101in}{2.449310in}}%
\pgfpathlineto{\pgfqpoint{2.290569in}{2.439432in}}%
\pgfpathlineto{\pgfqpoint{2.291080in}{2.477458in}}%
\pgfpathlineto{\pgfqpoint{2.292224in}{2.519434in}}%
\pgfpathlineto{\pgfqpoint{2.292234in}{2.517312in}}%
\pgfpathlineto{\pgfqpoint{2.292312in}{2.530055in}}%
\pgfpathlineto{\pgfqpoint{2.292511in}{2.491661in}}%
\pgfpathlineto{\pgfqpoint{2.293320in}{2.514123in}}%
\pgfpathlineto{\pgfqpoint{2.293407in}{2.505470in}}%
\pgfpathlineto{\pgfqpoint{2.293928in}{2.533877in}}%
\pgfpathlineto{\pgfqpoint{2.294395in}{2.533244in}}%
\pgfpathlineto{\pgfqpoint{2.294936in}{2.557406in}}%
\pgfpathlineto{\pgfqpoint{2.295384in}{2.519839in}}%
\pgfpathlineto{\pgfqpoint{2.295462in}{2.521474in}}%
\pgfpathlineto{\pgfqpoint{2.295579in}{2.521846in}}%
\pgfpathlineto{\pgfqpoint{2.295744in}{2.505029in}}%
\pgfpathlineto{\pgfqpoint{2.295768in}{2.507431in}}%
\pgfpathlineto{\pgfqpoint{2.295846in}{2.500878in}}%
\pgfpathlineto{\pgfqpoint{2.296815in}{2.534653in}}%
\pgfpathlineto{\pgfqpoint{2.296839in}{2.525316in}}%
\pgfpathlineto{\pgfqpoint{2.297730in}{2.502585in}}%
\pgfpathlineto{\pgfqpoint{2.296922in}{2.531719in}}%
\pgfpathlineto{\pgfqpoint{2.298037in}{2.518523in}}%
\pgfpathlineto{\pgfqpoint{2.298417in}{2.537311in}}%
\pgfpathlineto{\pgfqpoint{2.299118in}{2.508352in}}%
\pgfpathlineto{\pgfqpoint{2.299133in}{2.510805in}}%
\pgfpathlineto{\pgfqpoint{2.299171in}{2.501893in}}%
\pgfpathlineto{\pgfqpoint{2.299931in}{2.537156in}}%
\pgfpathlineto{\pgfqpoint{2.300140in}{2.535730in}}%
\pgfpathlineto{\pgfqpoint{2.300428in}{2.551167in}}%
\pgfpathlineto{\pgfqpoint{2.300958in}{2.512009in}}%
\pgfpathlineto{\pgfqpoint{2.301236in}{2.524739in}}%
\pgfpathlineto{\pgfqpoint{2.301250in}{2.527596in}}%
\pgfpathlineto{\pgfqpoint{2.302000in}{2.493466in}}%
\pgfpathlineto{\pgfqpoint{2.302278in}{2.504765in}}%
\pgfpathlineto{\pgfqpoint{2.303061in}{2.469008in}}%
\pgfpathlineto{\pgfqpoint{2.303427in}{2.497041in}}%
\pgfpathlineto{\pgfqpoint{2.303456in}{2.501825in}}%
\pgfpathlineto{\pgfqpoint{2.304254in}{2.443721in}}%
\pgfpathlineto{\pgfqpoint{2.304454in}{2.455040in}}%
\pgfpathlineto{\pgfqpoint{2.304551in}{2.451923in}}%
\pgfpathlineto{\pgfqpoint{2.305437in}{2.482072in}}%
\pgfpathlineto{\pgfqpoint{2.305442in}{2.479935in}}%
\pgfpathlineto{\pgfqpoint{2.305525in}{2.489323in}}%
\pgfpathlineto{\pgfqpoint{2.306411in}{2.445590in}}%
\pgfpathlineto{\pgfqpoint{2.306528in}{2.454075in}}%
\pgfpathlineto{\pgfqpoint{2.307404in}{2.436177in}}%
\pgfpathlineto{\pgfqpoint{2.306572in}{2.460180in}}%
\pgfpathlineto{\pgfqpoint{2.307701in}{2.444654in}}%
\pgfpathlineto{\pgfqpoint{2.308271in}{2.474975in}}%
\pgfpathlineto{\pgfqpoint{2.307716in}{2.441477in}}%
\pgfpathlineto{\pgfqpoint{2.308845in}{2.455174in}}%
\pgfpathlineto{\pgfqpoint{2.309960in}{2.442640in}}%
\pgfpathlineto{\pgfqpoint{2.309483in}{2.479573in}}%
\pgfpathlineto{\pgfqpoint{2.310004in}{2.445944in}}%
\pgfpathlineto{\pgfqpoint{2.311172in}{2.506139in}}%
\pgfpathlineto{\pgfqpoint{2.311177in}{2.505462in}}%
\pgfpathlineto{\pgfqpoint{2.312229in}{2.456264in}}%
\pgfpathlineto{\pgfqpoint{2.312365in}{2.461972in}}%
\pgfpathlineto{\pgfqpoint{2.312511in}{2.443820in}}%
\pgfpathlineto{\pgfqpoint{2.313159in}{2.476589in}}%
\pgfpathlineto{\pgfqpoint{2.313509in}{2.450953in}}%
\pgfpathlineto{\pgfqpoint{2.314897in}{2.489653in}}%
\pgfpathlineto{\pgfqpoint{2.313704in}{2.430300in}}%
\pgfpathlineto{\pgfqpoint{2.315121in}{2.474348in}}%
\pgfpathlineto{\pgfqpoint{2.315987in}{2.457679in}}%
\pgfpathlineto{\pgfqpoint{2.315388in}{2.485862in}}%
\pgfpathlineto{\pgfqpoint{2.316211in}{2.477519in}}%
\pgfpathlineto{\pgfqpoint{2.317151in}{2.508398in}}%
\pgfpathlineto{\pgfqpoint{2.317355in}{2.491619in}}%
\pgfpathlineto{\pgfqpoint{2.317375in}{2.493923in}}%
\pgfpathlineto{\pgfqpoint{2.317881in}{2.474074in}}%
\pgfpathlineto{\pgfqpoint{2.317891in}{2.476289in}}%
\pgfpathlineto{\pgfqpoint{2.318811in}{2.420852in}}%
\pgfpathlineto{\pgfqpoint{2.318178in}{2.484099in}}%
\pgfpathlineto{\pgfqpoint{2.319074in}{2.428887in}}%
\pgfpathlineto{\pgfqpoint{2.320140in}{2.435316in}}%
\pgfpathlineto{\pgfqpoint{2.319824in}{2.403365in}}%
\pgfpathlineto{\pgfqpoint{2.320199in}{2.432641in}}%
\pgfpathlineto{\pgfqpoint{2.320364in}{2.422129in}}%
\pgfpathlineto{\pgfqpoint{2.320749in}{2.445688in}}%
\pgfpathlineto{\pgfqpoint{2.321021in}{2.439528in}}%
\pgfpathlineto{\pgfqpoint{2.321274in}{2.462530in}}%
\pgfpathlineto{\pgfqpoint{2.321868in}{2.420173in}}%
\pgfpathlineto{\pgfqpoint{2.322083in}{2.431017in}}%
\pgfpathlineto{\pgfqpoint{2.322462in}{2.419699in}}%
\pgfpathlineto{\pgfqpoint{2.323056in}{2.453544in}}%
\pgfpathlineto{\pgfqpoint{2.323154in}{2.450408in}}%
\pgfpathlineto{\pgfqpoint{2.324361in}{2.425959in}}%
\pgfpathlineto{\pgfqpoint{2.323227in}{2.458735in}}%
\pgfpathlineto{\pgfqpoint{2.324385in}{2.428736in}}%
\pgfpathlineto{\pgfqpoint{2.325043in}{2.450078in}}%
\pgfpathlineto{\pgfqpoint{2.325349in}{2.423460in}}%
\pgfpathlineto{\pgfqpoint{2.325505in}{2.440395in}}%
\pgfpathlineto{\pgfqpoint{2.325515in}{2.438015in}}%
\pgfpathlineto{\pgfqpoint{2.326435in}{2.472394in}}%
\pgfpathlineto{\pgfqpoint{2.326440in}{2.471033in}}%
\pgfpathlineto{\pgfqpoint{2.327652in}{2.513564in}}%
\pgfpathlineto{\pgfqpoint{2.326493in}{2.470728in}}%
\pgfpathlineto{\pgfqpoint{2.327662in}{2.512512in}}%
\pgfpathlineto{\pgfqpoint{2.328718in}{2.468912in}}%
\pgfpathlineto{\pgfqpoint{2.328811in}{2.481102in}}%
\pgfpathlineto{\pgfqpoint{2.328947in}{2.498748in}}%
\pgfpathlineto{\pgfqpoint{2.329176in}{2.466061in}}%
\pgfpathlineto{\pgfqpoint{2.329950in}{2.496371in}}%
\pgfpathlineto{\pgfqpoint{2.331143in}{2.466657in}}%
\pgfpathlineto{\pgfqpoint{2.330111in}{2.502048in}}%
\pgfpathlineto{\pgfqpoint{2.331162in}{2.471281in}}%
\pgfpathlineto{\pgfqpoint{2.331231in}{2.473884in}}%
\pgfpathlineto{\pgfqpoint{2.331766in}{2.445344in}}%
\pgfpathlineto{\pgfqpoint{2.332258in}{2.468200in}}%
\pgfpathlineto{\pgfqpoint{2.332988in}{2.460284in}}%
\pgfpathlineto{\pgfqpoint{2.332861in}{2.483684in}}%
\pgfpathlineto{\pgfqpoint{2.333358in}{2.469920in}}%
\pgfpathlineto{\pgfqpoint{2.333986in}{2.530804in}}%
\pgfpathlineto{\pgfqpoint{2.333499in}{2.462676in}}%
\pgfpathlineto{\pgfqpoint{2.334522in}{2.502836in}}%
\pgfpathlineto{\pgfqpoint{2.334935in}{2.493027in}}%
\pgfpathlineto{\pgfqpoint{2.335568in}{2.527454in}}%
\pgfpathlineto{\pgfqpoint{2.335593in}{2.526885in}}%
\pgfpathlineto{\pgfqpoint{2.335602in}{2.528912in}}%
\pgfpathlineto{\pgfqpoint{2.336221in}{2.494270in}}%
\pgfpathlineto{\pgfqpoint{2.336527in}{2.512345in}}%
\pgfpathlineto{\pgfqpoint{2.337550in}{2.503298in}}%
\pgfpathlineto{\pgfqpoint{2.337204in}{2.537510in}}%
\pgfpathlineto{\pgfqpoint{2.337652in}{2.506892in}}%
\pgfpathlineto{\pgfqpoint{2.338095in}{2.521275in}}%
\pgfpathlineto{\pgfqpoint{2.338577in}{2.480849in}}%
\pgfpathlineto{\pgfqpoint{2.338616in}{2.483497in}}%
\pgfpathlineto{\pgfqpoint{2.339794in}{2.445110in}}%
\pgfpathlineto{\pgfqpoint{2.339098in}{2.501462in}}%
\pgfpathlineto{\pgfqpoint{2.339833in}{2.449999in}}%
\pgfpathlineto{\pgfqpoint{2.340958in}{2.506716in}}%
\pgfpathlineto{\pgfqpoint{2.339892in}{2.444267in}}%
\pgfpathlineto{\pgfqpoint{2.341211in}{2.487097in}}%
\pgfpathlineto{\pgfqpoint{2.342418in}{2.408428in}}%
\pgfpathlineto{\pgfqpoint{2.341240in}{2.494054in}}%
\pgfpathlineto{\pgfqpoint{2.342433in}{2.414374in}}%
\pgfpathlineto{\pgfqpoint{2.342983in}{2.447136in}}%
\pgfpathlineto{\pgfqpoint{2.343548in}{2.417679in}}%
\pgfpathlineto{\pgfqpoint{2.344648in}{2.407511in}}%
\pgfpathlineto{\pgfqpoint{2.343806in}{2.445727in}}%
\pgfpathlineto{\pgfqpoint{2.344668in}{2.414576in}}%
\pgfpathlineto{\pgfqpoint{2.344877in}{2.401994in}}%
\pgfpathlineto{\pgfqpoint{2.345154in}{2.439537in}}%
\pgfpathlineto{\pgfqpoint{2.345534in}{2.434190in}}%
\pgfpathlineto{\pgfqpoint{2.345583in}{2.443587in}}%
\pgfpathlineto{\pgfqpoint{2.346016in}{2.411066in}}%
\pgfpathlineto{\pgfqpoint{2.346634in}{2.430876in}}%
\pgfpathlineto{\pgfqpoint{2.347360in}{2.397473in}}%
\pgfpathlineto{\pgfqpoint{2.347618in}{2.431721in}}%
\pgfpathlineto{\pgfqpoint{2.347730in}{2.427685in}}%
\pgfpathlineto{\pgfqpoint{2.348718in}{2.448985in}}%
\pgfpathlineto{\pgfqpoint{2.348397in}{2.413378in}}%
\pgfpathlineto{\pgfqpoint{2.348855in}{2.441464in}}%
\pgfpathlineto{\pgfqpoint{2.349156in}{2.411370in}}%
\pgfpathlineto{\pgfqpoint{2.349862in}{2.483111in}}%
\pgfpathlineto{\pgfqpoint{2.350924in}{2.521196in}}%
\pgfpathlineto{\pgfqpoint{2.351118in}{2.499972in}}%
\pgfpathlineto{\pgfqpoint{2.351206in}{2.491325in}}%
\pgfpathlineto{\pgfqpoint{2.351756in}{2.515721in}}%
\pgfpathlineto{\pgfqpoint{2.352189in}{2.506297in}}%
\pgfpathlineto{\pgfqpoint{2.352321in}{2.527484in}}%
\pgfpathlineto{\pgfqpoint{2.352973in}{2.483263in}}%
\pgfpathlineto{\pgfqpoint{2.353100in}{2.486877in}}%
\pgfpathlineto{\pgfqpoint{2.353124in}{2.484171in}}%
\pgfpathlineto{\pgfqpoint{2.353309in}{2.512134in}}%
\pgfpathlineto{\pgfqpoint{2.353314in}{2.511080in}}%
\pgfpathlineto{\pgfqpoint{2.353334in}{2.510395in}}%
\pgfpathlineto{\pgfqpoint{2.354580in}{2.595437in}}%
\pgfpathlineto{\pgfqpoint{2.355651in}{2.554269in}}%
\pgfpathlineto{\pgfqpoint{2.355855in}{2.564496in}}%
\pgfpathlineto{\pgfqpoint{2.357063in}{2.601656in}}%
\pgfpathlineto{\pgfqpoint{2.356303in}{2.556264in}}%
\pgfpathlineto{\pgfqpoint{2.357073in}{2.599837in}}%
\pgfpathlineto{\pgfqpoint{2.357282in}{2.596000in}}%
\pgfpathlineto{\pgfqpoint{2.357827in}{2.642996in}}%
\pgfpathlineto{\pgfqpoint{2.357852in}{2.642989in}}%
\pgfpathlineto{\pgfqpoint{2.357978in}{2.657707in}}%
\pgfpathlineto{\pgfqpoint{2.358704in}{2.623936in}}%
\pgfpathlineto{\pgfqpoint{2.358947in}{2.639550in}}%
\pgfpathlineto{\pgfqpoint{2.359093in}{2.627094in}}%
\pgfpathlineto{\pgfqpoint{2.359546in}{2.663611in}}%
\pgfpathlineto{\pgfqpoint{2.360037in}{2.645180in}}%
\pgfpathlineto{\pgfqpoint{2.360125in}{2.655128in}}%
\pgfpathlineto{\pgfqpoint{2.360763in}{2.633540in}}%
\pgfpathlineto{\pgfqpoint{2.361138in}{2.644404in}}%
\pgfpathlineto{\pgfqpoint{2.362331in}{2.613710in}}%
\pgfpathlineto{\pgfqpoint{2.361182in}{2.650785in}}%
\pgfpathlineto{\pgfqpoint{2.362374in}{2.618740in}}%
\pgfpathlineto{\pgfqpoint{2.363041in}{2.642485in}}%
\pgfpathlineto{\pgfqpoint{2.363489in}{2.623365in}}%
\pgfpathlineto{\pgfqpoint{2.364254in}{2.607570in}}%
\pgfpathlineto{\pgfqpoint{2.363669in}{2.633930in}}%
\pgfpathlineto{\pgfqpoint{2.364604in}{2.618427in}}%
\pgfpathlineto{\pgfqpoint{2.364765in}{2.636929in}}%
\pgfpathlineto{\pgfqpoint{2.365222in}{2.592552in}}%
\pgfpathlineto{\pgfqpoint{2.365578in}{2.597974in}}%
\pgfpathlineto{\pgfqpoint{2.366527in}{2.590230in}}%
\pgfpathlineto{\pgfqpoint{2.365836in}{2.619667in}}%
\pgfpathlineto{\pgfqpoint{2.366664in}{2.606165in}}%
\pgfpathlineto{\pgfqpoint{2.368251in}{2.665804in}}%
\pgfpathlineto{\pgfqpoint{2.366751in}{2.599409in}}%
\pgfpathlineto{\pgfqpoint{2.368275in}{2.663824in}}%
\pgfpathlineto{\pgfqpoint{2.368747in}{2.651013in}}%
\pgfpathlineto{\pgfqpoint{2.369122in}{2.686935in}}%
\pgfpathlineto{\pgfqpoint{2.369356in}{2.672074in}}%
\pgfpathlineto{\pgfqpoint{2.369672in}{2.688488in}}%
\pgfpathlineto{\pgfqpoint{2.369434in}{2.667712in}}%
\pgfpathlineto{\pgfqpoint{2.370500in}{2.687896in}}%
\pgfpathlineto{\pgfqpoint{2.370510in}{2.688788in}}%
\pgfpathlineto{\pgfqpoint{2.371625in}{2.642435in}}%
\pgfpathlineto{\pgfqpoint{2.371941in}{2.632422in}}%
\pgfpathlineto{\pgfqpoint{2.372189in}{2.653497in}}%
\pgfpathlineto{\pgfqpoint{2.372730in}{2.643458in}}%
\pgfpathlineto{\pgfqpoint{2.373372in}{2.661046in}}%
\pgfpathlineto{\pgfqpoint{2.373742in}{2.624579in}}%
\pgfpathlineto{\pgfqpoint{2.374794in}{2.579635in}}%
\pgfpathlineto{\pgfqpoint{2.374361in}{2.629215in}}%
\pgfpathlineto{\pgfqpoint{2.375018in}{2.581886in}}%
\pgfpathlineto{\pgfqpoint{2.375845in}{2.594682in}}%
\pgfpathlineto{\pgfqpoint{2.375475in}{2.571292in}}%
\pgfpathlineto{\pgfqpoint{2.376123in}{2.582291in}}%
\pgfpathlineto{\pgfqpoint{2.377223in}{2.537032in}}%
\pgfpathlineto{\pgfqpoint{2.377301in}{2.555420in}}%
\pgfpathlineto{\pgfqpoint{2.378367in}{2.566586in}}%
\pgfpathlineto{\pgfqpoint{2.377885in}{2.541720in}}%
\pgfpathlineto{\pgfqpoint{2.378450in}{2.560066in}}%
\pgfpathlineto{\pgfqpoint{2.378538in}{2.547878in}}%
\pgfpathlineto{\pgfqpoint{2.379434in}{2.584876in}}%
\pgfpathlineto{\pgfqpoint{2.379521in}{2.576031in}}%
\pgfpathlineto{\pgfqpoint{2.380617in}{2.600383in}}%
\pgfpathlineto{\pgfqpoint{2.379925in}{2.561175in}}%
\pgfpathlineto{\pgfqpoint{2.380675in}{2.593482in}}%
\pgfpathlineto{\pgfqpoint{2.380996in}{2.563645in}}%
\pgfpathlineto{\pgfqpoint{2.381595in}{2.603316in}}%
\pgfpathlineto{\pgfqpoint{2.381790in}{2.589704in}}%
\pgfpathlineto{\pgfqpoint{2.381809in}{2.597195in}}%
\pgfpathlineto{\pgfqpoint{2.382369in}{2.561349in}}%
\pgfpathlineto{\pgfqpoint{2.382866in}{2.570694in}}%
\pgfpathlineto{\pgfqpoint{2.382919in}{2.564234in}}%
\pgfpathlineto{\pgfqpoint{2.383431in}{2.596352in}}%
\pgfpathlineto{\pgfqpoint{2.383981in}{2.566713in}}%
\pgfpathlineto{\pgfqpoint{2.385032in}{2.585764in}}%
\pgfpathlineto{\pgfqpoint{2.384283in}{2.535079in}}%
\pgfpathlineto{\pgfqpoint{2.385125in}{2.575645in}}%
\pgfpathlineto{\pgfqpoint{2.386220in}{2.555127in}}%
\pgfpathlineto{\pgfqpoint{2.385617in}{2.576850in}}%
\pgfpathlineto{\pgfqpoint{2.386254in}{2.563648in}}%
\pgfpathlineto{\pgfqpoint{2.386517in}{2.580246in}}%
\pgfpathlineto{\pgfqpoint{2.387165in}{2.539725in}}%
\pgfpathlineto{\pgfqpoint{2.387291in}{2.542748in}}%
\pgfpathlineto{\pgfqpoint{2.387394in}{2.528662in}}%
\pgfpathlineto{\pgfqpoint{2.388231in}{2.571857in}}%
\pgfpathlineto{\pgfqpoint{2.388319in}{2.558155in}}%
\pgfpathlineto{\pgfqpoint{2.388669in}{2.572026in}}%
\pgfpathlineto{\pgfqpoint{2.389107in}{2.547636in}}%
\pgfpathlineto{\pgfqpoint{2.389443in}{2.566782in}}%
\pgfpathlineto{\pgfqpoint{2.389458in}{2.564596in}}%
\pgfpathlineto{\pgfqpoint{2.389526in}{2.574231in}}%
\pgfpathlineto{\pgfqpoint{2.389750in}{2.598306in}}%
\pgfpathlineto{\pgfqpoint{2.390393in}{2.531121in}}%
\pgfpathlineto{\pgfqpoint{2.390539in}{2.531974in}}%
\pgfpathlineto{\pgfqpoint{2.390724in}{2.522045in}}%
\pgfpathlineto{\pgfqpoint{2.391566in}{2.560488in}}%
\pgfpathlineto{\pgfqpoint{2.391595in}{2.556405in}}%
\pgfpathlineto{\pgfqpoint{2.392759in}{2.574215in}}%
\pgfpathlineto{\pgfqpoint{2.392160in}{2.536659in}}%
\pgfpathlineto{\pgfqpoint{2.392778in}{2.574211in}}%
\pgfpathlineto{\pgfqpoint{2.393440in}{2.521137in}}%
\pgfpathlineto{\pgfqpoint{2.393917in}{2.561509in}}%
\pgfpathlineto{\pgfqpoint{2.394555in}{2.569334in}}%
\pgfpathlineto{\pgfqpoint{2.394852in}{2.539714in}}%
\pgfpathlineto{\pgfqpoint{2.394979in}{2.541799in}}%
\pgfpathlineto{\pgfqpoint{2.395023in}{2.539839in}}%
\pgfpathlineto{\pgfqpoint{2.395982in}{2.567634in}}%
\pgfpathlineto{\pgfqpoint{2.396011in}{2.562600in}}%
\pgfpathlineto{\pgfqpoint{2.396025in}{2.568303in}}%
\pgfpathlineto{\pgfqpoint{2.396634in}{2.541516in}}%
\pgfpathlineto{\pgfqpoint{2.397121in}{2.563122in}}%
\pgfpathlineto{\pgfqpoint{2.398416in}{2.615494in}}%
\pgfpathlineto{\pgfqpoint{2.397145in}{2.559281in}}%
\pgfpathlineto{\pgfqpoint{2.398460in}{2.607192in}}%
\pgfpathlineto{\pgfqpoint{2.399278in}{2.573975in}}%
\pgfpathlineto{\pgfqpoint{2.398912in}{2.609444in}}%
\pgfpathlineto{\pgfqpoint{2.399594in}{2.593666in}}%
\pgfpathlineto{\pgfqpoint{2.399643in}{2.583766in}}%
\pgfpathlineto{\pgfqpoint{2.400480in}{2.612008in}}%
\pgfpathlineto{\pgfqpoint{2.401799in}{2.662829in}}%
\pgfpathlineto{\pgfqpoint{2.400719in}{2.599374in}}%
\pgfpathlineto{\pgfqpoint{2.401804in}{2.661484in}}%
\pgfpathlineto{\pgfqpoint{2.403031in}{2.610560in}}%
\pgfpathlineto{\pgfqpoint{2.401907in}{2.665163in}}%
\pgfpathlineto{\pgfqpoint{2.403075in}{2.616417in}}%
\pgfpathlineto{\pgfqpoint{2.403211in}{2.604113in}}%
\pgfpathlineto{\pgfqpoint{2.403800in}{2.654938in}}%
\pgfpathlineto{\pgfqpoint{2.404024in}{2.640091in}}%
\pgfpathlineto{\pgfqpoint{2.404847in}{2.662851in}}%
\pgfpathlineto{\pgfqpoint{2.404531in}{2.636401in}}%
\pgfpathlineto{\pgfqpoint{2.405159in}{2.652976in}}%
\pgfpathlineto{\pgfqpoint{2.406011in}{2.634805in}}%
\pgfpathlineto{\pgfqpoint{2.405548in}{2.660488in}}%
\pgfpathlineto{\pgfqpoint{2.406274in}{2.650067in}}%
\pgfpathlineto{\pgfqpoint{2.407160in}{2.666206in}}%
\pgfpathlineto{\pgfqpoint{2.406872in}{2.635744in}}%
\pgfpathlineto{\pgfqpoint{2.407427in}{2.663583in}}%
\pgfpathlineto{\pgfqpoint{2.407481in}{2.663845in}}%
\pgfpathlineto{\pgfqpoint{2.408659in}{2.608908in}}%
\pgfpathlineto{\pgfqpoint{2.408805in}{2.603127in}}%
\pgfpathlineto{\pgfqpoint{2.409146in}{2.630678in}}%
\pgfpathlineto{\pgfqpoint{2.409399in}{2.629599in}}%
\pgfpathlineto{\pgfqpoint{2.409794in}{2.667057in}}%
\pgfpathlineto{\pgfqpoint{2.409458in}{2.624940in}}%
\pgfpathlineto{\pgfqpoint{2.410611in}{2.654132in}}%
\pgfpathlineto{\pgfqpoint{2.411614in}{2.676665in}}%
\pgfpathlineto{\pgfqpoint{2.411274in}{2.644480in}}%
\pgfpathlineto{\pgfqpoint{2.411790in}{2.668071in}}%
\pgfpathlineto{\pgfqpoint{2.411902in}{2.670470in}}%
\pgfpathlineto{\pgfqpoint{2.413065in}{2.627029in}}%
\pgfpathlineto{\pgfqpoint{2.414122in}{2.668475in}}%
\pgfpathlineto{\pgfqpoint{2.413109in}{2.621846in}}%
\pgfpathlineto{\pgfqpoint{2.414229in}{2.655849in}}%
\pgfpathlineto{\pgfqpoint{2.414448in}{2.640606in}}%
\pgfpathlineto{\pgfqpoint{2.414930in}{2.664432in}}%
\pgfpathlineto{\pgfqpoint{2.415378in}{2.646519in}}%
\pgfpathlineto{\pgfqpoint{2.415431in}{2.632902in}}%
\pgfpathlineto{\pgfqpoint{2.415689in}{2.664855in}}%
\pgfpathlineto{\pgfqpoint{2.416502in}{2.644600in}}%
\pgfpathlineto{\pgfqpoint{2.416580in}{2.636116in}}%
\pgfpathlineto{\pgfqpoint{2.417666in}{2.670156in}}%
\pgfpathlineto{\pgfqpoint{2.417671in}{2.669462in}}%
\pgfpathlineto{\pgfqpoint{2.418104in}{2.699338in}}%
\pgfpathlineto{\pgfqpoint{2.418635in}{2.688875in}}%
\pgfpathlineto{\pgfqpoint{2.419813in}{2.729829in}}%
\pgfpathlineto{\pgfqpoint{2.418990in}{2.652991in}}%
\pgfpathlineto{\pgfqpoint{2.419818in}{2.729393in}}%
\pgfpathlineto{\pgfqpoint{2.420081in}{2.710825in}}%
\pgfpathlineto{\pgfqpoint{2.420865in}{2.750256in}}%
\pgfpathlineto{\pgfqpoint{2.420874in}{2.748788in}}%
\pgfpathlineto{\pgfqpoint{2.421770in}{2.775667in}}%
\pgfpathlineto{\pgfqpoint{2.421205in}{2.735110in}}%
\pgfpathlineto{\pgfqpoint{2.422028in}{2.757617in}}%
\pgfpathlineto{\pgfqpoint{2.422603in}{2.737909in}}%
\pgfpathlineto{\pgfqpoint{2.423055in}{2.770990in}}%
\pgfpathlineto{\pgfqpoint{2.423143in}{2.755464in}}%
\pgfpathlineto{\pgfqpoint{2.423581in}{2.794134in}}%
\pgfpathlineto{\pgfqpoint{2.424078in}{2.752748in}}%
\pgfpathlineto{\pgfqpoint{2.424272in}{2.768591in}}%
\pgfpathlineto{\pgfqpoint{2.424287in}{2.772496in}}%
\pgfpathlineto{\pgfqpoint{2.424774in}{2.735368in}}%
\pgfpathlineto{\pgfqpoint{2.425290in}{2.745060in}}%
\pgfpathlineto{\pgfqpoint{2.425348in}{2.743641in}}%
\pgfpathlineto{\pgfqpoint{2.425592in}{2.773023in}}%
\pgfpathlineto{\pgfqpoint{2.426030in}{2.759458in}}%
\pgfpathlineto{\pgfqpoint{2.426955in}{2.763774in}}%
\pgfpathlineto{\pgfqpoint{2.426200in}{2.739290in}}%
\pgfpathlineto{\pgfqpoint{2.427086in}{2.752840in}}%
\pgfpathlineto{\pgfqpoint{2.427106in}{2.758042in}}%
\pgfpathlineto{\pgfqpoint{2.427179in}{2.747161in}}%
\pgfpathlineto{\pgfqpoint{2.427354in}{2.721406in}}%
\pgfpathlineto{\pgfqpoint{2.427895in}{2.767090in}}%
\pgfpathlineto{\pgfqpoint{2.428284in}{2.749343in}}%
\pgfpathlineto{\pgfqpoint{2.428883in}{2.759709in}}%
\pgfpathlineto{\pgfqpoint{2.428425in}{2.737597in}}%
\pgfpathlineto{\pgfqpoint{2.429355in}{2.746277in}}%
\pgfpathlineto{\pgfqpoint{2.430504in}{2.718091in}}%
\pgfpathlineto{\pgfqpoint{2.429555in}{2.757663in}}%
\pgfpathlineto{\pgfqpoint{2.430533in}{2.723423in}}%
\pgfpathlineto{\pgfqpoint{2.430752in}{2.746889in}}%
\pgfpathlineto{\pgfqpoint{2.431507in}{2.711433in}}%
\pgfpathlineto{\pgfqpoint{2.431648in}{2.728290in}}%
\pgfpathlineto{\pgfqpoint{2.431897in}{2.720071in}}%
\pgfpathlineto{\pgfqpoint{2.432276in}{2.742635in}}%
\pgfpathlineto{\pgfqpoint{2.432573in}{2.765739in}}%
\pgfpathlineto{\pgfqpoint{2.432315in}{2.734541in}}%
\pgfpathlineto{\pgfqpoint{2.433430in}{2.750497in}}%
\pgfpathlineto{\pgfqpoint{2.433620in}{2.741221in}}%
\pgfpathlineto{\pgfqpoint{2.434428in}{2.791673in}}%
\pgfpathlineto{\pgfqpoint{2.434482in}{2.781075in}}%
\pgfpathlineto{\pgfqpoint{2.434487in}{2.781942in}}%
\pgfpathlineto{\pgfqpoint{2.435436in}{2.756891in}}%
\pgfpathlineto{\pgfqpoint{2.435485in}{2.762098in}}%
\pgfpathlineto{\pgfqpoint{2.435548in}{2.754532in}}%
\pgfpathlineto{\pgfqpoint{2.436186in}{2.776771in}}%
\pgfpathlineto{\pgfqpoint{2.436580in}{2.768874in}}%
\pgfpathlineto{\pgfqpoint{2.437710in}{2.810621in}}%
\pgfpathlineto{\pgfqpoint{2.437773in}{2.799923in}}%
\pgfpathlineto{\pgfqpoint{2.438737in}{2.782638in}}%
\pgfpathlineto{\pgfqpoint{2.438503in}{2.808670in}}%
\pgfpathlineto{\pgfqpoint{2.438888in}{2.795342in}}%
\pgfpathlineto{\pgfqpoint{2.439910in}{2.773804in}}%
\pgfpathlineto{\pgfqpoint{2.439044in}{2.800837in}}%
\pgfpathlineto{\pgfqpoint{2.440168in}{2.782383in}}%
\pgfpathlineto{\pgfqpoint{2.440957in}{2.796592in}}%
\pgfpathlineto{\pgfqpoint{2.441147in}{2.769712in}}%
\pgfpathlineto{\pgfqpoint{2.441249in}{2.771062in}}%
\pgfpathlineto{\pgfqpoint{2.441960in}{2.784728in}}%
\pgfpathlineto{\pgfqpoint{2.442266in}{2.747176in}}%
\pgfpathlineto{\pgfqpoint{2.442515in}{2.727893in}}%
\pgfpathlineto{\pgfqpoint{2.442992in}{2.762386in}}%
\pgfpathlineto{\pgfqpoint{2.443386in}{2.745510in}}%
\pgfpathlineto{\pgfqpoint{2.443396in}{2.750064in}}%
\pgfpathlineto{\pgfqpoint{2.443566in}{2.723283in}}%
\pgfpathlineto{\pgfqpoint{2.444501in}{2.748300in}}%
\pgfpathlineto{\pgfqpoint{2.446220in}{2.705493in}}%
\pgfpathlineto{\pgfqpoint{2.444598in}{2.754886in}}%
\pgfpathlineto{\pgfqpoint{2.446361in}{2.716368in}}%
\pgfpathlineto{\pgfqpoint{2.447305in}{2.740321in}}%
\pgfpathlineto{\pgfqpoint{2.446541in}{2.708605in}}%
\pgfpathlineto{\pgfqpoint{2.447471in}{2.717515in}}%
\pgfpathlineto{\pgfqpoint{2.448415in}{2.682913in}}%
\pgfpathlineto{\pgfqpoint{2.448625in}{2.699723in}}%
\pgfpathlineto{\pgfqpoint{2.448931in}{2.716608in}}%
\pgfpathlineto{\pgfqpoint{2.449433in}{2.670887in}}%
\pgfpathlineto{\pgfqpoint{2.449574in}{2.681206in}}%
\pgfpathlineto{\pgfqpoint{2.450290in}{2.650735in}}%
\pgfpathlineto{\pgfqpoint{2.449603in}{2.686978in}}%
\pgfpathlineto{\pgfqpoint{2.450747in}{2.652878in}}%
\pgfpathlineto{\pgfqpoint{2.451375in}{2.698625in}}%
\pgfpathlineto{\pgfqpoint{2.451887in}{2.660446in}}%
\pgfpathlineto{\pgfqpoint{2.451979in}{2.662356in}}%
\pgfpathlineto{\pgfqpoint{2.452651in}{2.612228in}}%
\pgfpathlineto{\pgfqpoint{2.452709in}{2.619619in}}%
\pgfpathlineto{\pgfqpoint{2.452743in}{2.612748in}}%
\pgfpathlineto{\pgfqpoint{2.453732in}{2.670178in}}%
\pgfpathlineto{\pgfqpoint{2.453746in}{2.669451in}}%
\pgfpathlineto{\pgfqpoint{2.454058in}{2.639320in}}%
\pgfpathlineto{\pgfqpoint{2.454735in}{2.694047in}}%
\pgfpathlineto{\pgfqpoint{2.454740in}{2.695666in}}%
\pgfpathlineto{\pgfqpoint{2.455640in}{2.665870in}}%
\pgfpathlineto{\pgfqpoint{2.455781in}{2.673913in}}%
\pgfpathlineto{\pgfqpoint{2.456940in}{2.641171in}}%
\pgfpathlineto{\pgfqpoint{2.456332in}{2.676844in}}%
\pgfpathlineto{\pgfqpoint{2.456999in}{2.642968in}}%
\pgfpathlineto{\pgfqpoint{2.457033in}{2.648802in}}%
\pgfpathlineto{\pgfqpoint{2.457792in}{2.605425in}}%
\pgfpathlineto{\pgfqpoint{2.458070in}{2.595740in}}%
\pgfpathlineto{\pgfqpoint{2.458737in}{2.637206in}}%
\pgfpathlineto{\pgfqpoint{2.459710in}{2.641663in}}%
\pgfpathlineto{\pgfqpoint{2.459038in}{2.624711in}}%
\pgfpathlineto{\pgfqpoint{2.459827in}{2.631697in}}%
\pgfpathlineto{\pgfqpoint{2.460343in}{2.599504in}}%
\pgfpathlineto{\pgfqpoint{2.459973in}{2.640963in}}%
\pgfpathlineto{\pgfqpoint{2.460966in}{2.612248in}}%
\pgfpathlineto{\pgfqpoint{2.461906in}{2.652938in}}%
\pgfpathlineto{\pgfqpoint{2.461468in}{2.594309in}}%
\pgfpathlineto{\pgfqpoint{2.462159in}{2.640397in}}%
\pgfpathlineto{\pgfqpoint{2.463045in}{2.616262in}}%
\pgfpathlineto{\pgfqpoint{2.462578in}{2.651417in}}%
\pgfpathlineto{\pgfqpoint{2.463332in}{2.625478in}}%
\pgfpathlineto{\pgfqpoint{2.464146in}{2.612436in}}%
\pgfpathlineto{\pgfqpoint{2.463780in}{2.636069in}}%
\pgfpathlineto{\pgfqpoint{2.464287in}{2.630824in}}%
\pgfpathlineto{\pgfqpoint{2.464959in}{2.659218in}}%
\pgfpathlineto{\pgfqpoint{2.465416in}{2.648640in}}%
\pgfpathlineto{\pgfqpoint{2.466098in}{2.631658in}}%
\pgfpathlineto{\pgfqpoint{2.466341in}{2.669577in}}%
\pgfpathlineto{\pgfqpoint{2.466497in}{2.656952in}}%
\pgfpathlineto{\pgfqpoint{2.466984in}{2.644450in}}%
\pgfpathlineto{\pgfqpoint{2.466628in}{2.672735in}}%
\pgfpathlineto{\pgfqpoint{2.467115in}{2.662519in}}%
\pgfpathlineto{\pgfqpoint{2.467140in}{2.674506in}}%
\pgfpathlineto{\pgfqpoint{2.468099in}{2.645449in}}%
\pgfpathlineto{\pgfqpoint{2.468211in}{2.651016in}}%
\pgfpathlineto{\pgfqpoint{2.469034in}{2.625884in}}%
\pgfpathlineto{\pgfqpoint{2.468352in}{2.658640in}}%
\pgfpathlineto{\pgfqpoint{2.469340in}{2.639835in}}%
\pgfpathlineto{\pgfqpoint{2.470499in}{2.623118in}}%
\pgfpathlineto{\pgfqpoint{2.469598in}{2.655369in}}%
\pgfpathlineto{\pgfqpoint{2.470509in}{2.625105in}}%
\pgfpathlineto{\pgfqpoint{2.471059in}{2.612008in}}%
\pgfpathlineto{\pgfqpoint{2.470752in}{2.632174in}}%
\pgfpathlineto{\pgfqpoint{2.471156in}{2.631787in}}%
\pgfpathlineto{\pgfqpoint{2.472797in}{2.696713in}}%
\pgfpathlineto{\pgfqpoint{2.472953in}{2.683972in}}%
\pgfpathlineto{\pgfqpoint{2.473196in}{2.677206in}}%
\pgfpathlineto{\pgfqpoint{2.473357in}{2.702156in}}%
\pgfpathlineto{\pgfqpoint{2.474038in}{2.692296in}}%
\pgfpathlineto{\pgfqpoint{2.474272in}{2.709088in}}%
\pgfpathlineto{\pgfqpoint{2.475056in}{2.672487in}}%
\pgfpathlineto{\pgfqpoint{2.475105in}{2.679282in}}%
\pgfpathlineto{\pgfqpoint{2.476322in}{2.638055in}}%
\pgfpathlineto{\pgfqpoint{2.475601in}{2.684209in}}%
\pgfpathlineto{\pgfqpoint{2.476370in}{2.639458in}}%
\pgfpathlineto{\pgfqpoint{2.477553in}{2.708179in}}%
\pgfpathlineto{\pgfqpoint{2.476390in}{2.639043in}}%
\pgfpathlineto{\pgfqpoint{2.477607in}{2.701887in}}%
\pgfpathlineto{\pgfqpoint{2.477797in}{2.690506in}}%
\pgfpathlineto{\pgfqpoint{2.478503in}{2.729080in}}%
\pgfpathlineto{\pgfqpoint{2.478532in}{2.727222in}}%
\pgfpathlineto{\pgfqpoint{2.479720in}{2.765916in}}%
\pgfpathlineto{\pgfqpoint{2.479949in}{2.746634in}}%
\pgfpathlineto{\pgfqpoint{2.479954in}{2.744911in}}%
\pgfpathlineto{\pgfqpoint{2.480869in}{2.809887in}}%
\pgfpathlineto{\pgfqpoint{2.480874in}{2.809029in}}%
\pgfpathlineto{\pgfqpoint{2.480918in}{2.818728in}}%
\pgfpathlineto{\pgfqpoint{2.481862in}{2.748373in}}%
\pgfpathlineto{\pgfqpoint{2.481911in}{2.762640in}}%
\pgfpathlineto{\pgfqpoint{2.482987in}{2.741417in}}%
\pgfpathlineto{\pgfqpoint{2.482091in}{2.769071in}}%
\pgfpathlineto{\pgfqpoint{2.483055in}{2.753138in}}%
\pgfpathlineto{\pgfqpoint{2.483459in}{2.772208in}}%
\pgfpathlineto{\pgfqpoint{2.484063in}{2.735207in}}%
\pgfpathlineto{\pgfqpoint{2.484145in}{2.742297in}}%
\pgfpathlineto{\pgfqpoint{2.484433in}{2.768704in}}%
\pgfpathlineto{\pgfqpoint{2.484257in}{2.739582in}}%
\pgfpathlineto{\pgfqpoint{2.485299in}{2.752892in}}%
\pgfpathlineto{\pgfqpoint{2.485353in}{2.741716in}}%
\pgfpathlineto{\pgfqpoint{2.486361in}{2.778568in}}%
\pgfpathlineto{\pgfqpoint{2.486726in}{2.793283in}}%
\pgfpathlineto{\pgfqpoint{2.487359in}{2.760093in}}%
\pgfpathlineto{\pgfqpoint{2.487363in}{2.760916in}}%
\pgfpathlineto{\pgfqpoint{2.487821in}{2.740551in}}%
\pgfpathlineto{\pgfqpoint{2.487597in}{2.770809in}}%
\pgfpathlineto{\pgfqpoint{2.488478in}{2.758350in}}%
\pgfpathlineto{\pgfqpoint{2.488683in}{2.767405in}}%
\pgfpathlineto{\pgfqpoint{2.489228in}{2.725745in}}%
\pgfpathlineto{\pgfqpoint{2.489569in}{2.749365in}}%
\pgfpathlineto{\pgfqpoint{2.489574in}{2.749432in}}%
\pgfpathlineto{\pgfqpoint{2.489686in}{2.736778in}}%
\pgfpathlineto{\pgfqpoint{2.489691in}{2.737128in}}%
\pgfpathlineto{\pgfqpoint{2.489841in}{2.733710in}}%
\pgfpathlineto{\pgfqpoint{2.490601in}{2.762531in}}%
\pgfpathlineto{\pgfqpoint{2.490737in}{2.757567in}}%
\pgfpathlineto{\pgfqpoint{2.490752in}{2.761192in}}%
\pgfpathlineto{\pgfqpoint{2.491477in}{2.723568in}}%
\pgfpathlineto{\pgfqpoint{2.491794in}{2.731924in}}%
\pgfpathlineto{\pgfqpoint{2.491915in}{2.725200in}}%
\pgfpathlineto{\pgfqpoint{2.492436in}{2.741154in}}%
\pgfpathlineto{\pgfqpoint{2.492909in}{2.726706in}}%
\pgfpathlineto{\pgfqpoint{2.493590in}{2.744389in}}%
\pgfpathlineto{\pgfqpoint{2.493099in}{2.720115in}}%
\pgfpathlineto{\pgfqpoint{2.494038in}{2.741759in}}%
\pgfpathlineto{\pgfqpoint{2.494150in}{2.744074in}}%
\pgfpathlineto{\pgfqpoint{2.494613in}{2.696551in}}%
\pgfpathlineto{\pgfqpoint{2.495640in}{2.656763in}}%
\pgfpathlineto{\pgfqpoint{2.495382in}{2.697646in}}%
\pgfpathlineto{\pgfqpoint{2.495757in}{2.678014in}}%
\pgfpathlineto{\pgfqpoint{2.496351in}{2.686508in}}%
\pgfpathlineto{\pgfqpoint{2.496774in}{2.654699in}}%
\pgfpathlineto{\pgfqpoint{2.496808in}{2.659769in}}%
\pgfpathlineto{\pgfqpoint{2.497529in}{2.694716in}}%
\pgfpathlineto{\pgfqpoint{2.496891in}{2.640723in}}%
\pgfpathlineto{\pgfqpoint{2.498006in}{2.680715in}}%
\pgfpathlineto{\pgfqpoint{2.498055in}{2.674656in}}%
\pgfpathlineto{\pgfqpoint{2.498688in}{2.695162in}}%
\pgfpathlineto{\pgfqpoint{2.499097in}{2.684265in}}%
\pgfpathlineto{\pgfqpoint{2.500114in}{2.730548in}}%
\pgfpathlineto{\pgfqpoint{2.499471in}{2.672374in}}%
\pgfpathlineto{\pgfqpoint{2.500523in}{2.724561in}}%
\pgfpathlineto{\pgfqpoint{2.501599in}{2.706102in}}%
\pgfpathlineto{\pgfqpoint{2.500990in}{2.737444in}}%
\pgfpathlineto{\pgfqpoint{2.501677in}{2.708070in}}%
\pgfpathlineto{\pgfqpoint{2.502621in}{2.736050in}}%
\pgfpathlineto{\pgfqpoint{2.501974in}{2.692350in}}%
\pgfpathlineto{\pgfqpoint{2.502860in}{2.723015in}}%
\pgfpathlineto{\pgfqpoint{2.503512in}{2.747566in}}%
\pgfpathlineto{\pgfqpoint{2.503892in}{2.709731in}}%
\pgfpathlineto{\pgfqpoint{2.504160in}{2.692524in}}%
\pgfpathlineto{\pgfqpoint{2.504885in}{2.730599in}}%
\pgfpathlineto{\pgfqpoint{2.504934in}{2.723679in}}%
\pgfpathlineto{\pgfqpoint{2.505080in}{2.744055in}}%
\pgfpathlineto{\pgfqpoint{2.505640in}{2.701214in}}%
\pgfpathlineto{\pgfqpoint{2.506029in}{2.718251in}}%
\pgfpathlineto{\pgfqpoint{2.506093in}{2.710234in}}%
\pgfpathlineto{\pgfqpoint{2.506964in}{2.768885in}}%
\pgfpathlineto{\pgfqpoint{2.507422in}{2.791474in}}%
\pgfpathlineto{\pgfqpoint{2.507095in}{2.761569in}}%
\pgfpathlineto{\pgfqpoint{2.508132in}{2.779399in}}%
\pgfpathlineto{\pgfqpoint{2.508162in}{2.769285in}}%
\pgfpathlineto{\pgfqpoint{2.508853in}{2.798532in}}%
\pgfpathlineto{\pgfqpoint{2.509194in}{2.792280in}}%
\pgfpathlineto{\pgfqpoint{2.510017in}{2.799564in}}%
\pgfpathlineto{\pgfqpoint{2.509457in}{2.771290in}}%
\pgfpathlineto{\pgfqpoint{2.510265in}{2.772900in}}%
\pgfpathlineto{\pgfqpoint{2.510275in}{2.770549in}}%
\pgfpathlineto{\pgfqpoint{2.510289in}{2.774245in}}%
\pgfpathlineto{\pgfqpoint{2.510294in}{2.774118in}}%
\pgfpathlineto{\pgfqpoint{2.510304in}{2.778431in}}%
\pgfpathlineto{\pgfqpoint{2.510718in}{2.745223in}}%
\pgfpathlineto{\pgfqpoint{2.511365in}{2.762595in}}%
\pgfpathlineto{\pgfqpoint{2.511531in}{2.754160in}}%
\pgfpathlineto{\pgfqpoint{2.512164in}{2.795436in}}%
\pgfpathlineto{\pgfqpoint{2.512353in}{2.787984in}}%
\pgfpathlineto{\pgfqpoint{2.512368in}{2.793228in}}%
\pgfpathlineto{\pgfqpoint{2.512908in}{2.767108in}}%
\pgfpathlineto{\pgfqpoint{2.513449in}{2.781456in}}%
\pgfpathlineto{\pgfqpoint{2.513600in}{2.769014in}}%
\pgfpathlineto{\pgfqpoint{2.514427in}{2.805603in}}%
\pgfpathlineto{\pgfqpoint{2.514515in}{2.800245in}}%
\pgfpathlineto{\pgfqpoint{2.515187in}{2.759632in}}%
\pgfpathlineto{\pgfqpoint{2.515645in}{2.792468in}}%
\pgfpathlineto{\pgfqpoint{2.516453in}{2.816842in}}%
\pgfpathlineto{\pgfqpoint{2.516058in}{2.783808in}}%
\pgfpathlineto{\pgfqpoint{2.516755in}{2.793457in}}%
\pgfpathlineto{\pgfqpoint{2.517553in}{2.754056in}}%
\pgfpathlineto{\pgfqpoint{2.517918in}{2.758766in}}%
\pgfpathlineto{\pgfqpoint{2.518931in}{2.714809in}}%
\pgfpathlineto{\pgfqpoint{2.518045in}{2.769007in}}%
\pgfpathlineto{\pgfqpoint{2.519233in}{2.723647in}}%
\pgfpathlineto{\pgfqpoint{2.519296in}{2.725735in}}%
\pgfpathlineto{\pgfqpoint{2.519870in}{2.679645in}}%
\pgfpathlineto{\pgfqpoint{2.520153in}{2.687052in}}%
\pgfpathlineto{\pgfqpoint{2.521360in}{2.619703in}}%
\pgfpathlineto{\pgfqpoint{2.520265in}{2.700100in}}%
\pgfpathlineto{\pgfqpoint{2.521365in}{2.620926in}}%
\pgfpathlineto{\pgfqpoint{2.521535in}{2.629867in}}%
\pgfpathlineto{\pgfqpoint{2.522412in}{2.605731in}}%
\pgfpathlineto{\pgfqpoint{2.522456in}{2.601597in}}%
\pgfpathlineto{\pgfqpoint{2.522436in}{2.608624in}}%
\pgfpathlineto{\pgfqpoint{2.522460in}{2.603259in}}%
\pgfpathlineto{\pgfqpoint{2.522490in}{2.604028in}}%
\pgfpathlineto{\pgfqpoint{2.523639in}{2.556972in}}%
\pgfpathlineto{\pgfqpoint{2.523858in}{2.550347in}}%
\pgfpathlineto{\pgfqpoint{2.524457in}{2.587519in}}%
\pgfpathlineto{\pgfqpoint{2.524763in}{2.552764in}}%
\pgfpathlineto{\pgfqpoint{2.524768in}{2.553851in}}%
\pgfpathlineto{\pgfqpoint{2.525138in}{2.527840in}}%
\pgfpathlineto{\pgfqpoint{2.525839in}{2.544919in}}%
\pgfpathlineto{\pgfqpoint{2.526024in}{2.527657in}}%
\pgfpathlineto{\pgfqpoint{2.526394in}{2.559751in}}%
\pgfpathlineto{\pgfqpoint{2.526935in}{2.546146in}}%
\pgfpathlineto{\pgfqpoint{2.526983in}{2.549949in}}%
\pgfpathlineto{\pgfqpoint{2.527353in}{2.509447in}}%
\pgfpathlineto{\pgfqpoint{2.527782in}{2.521433in}}%
\pgfpathlineto{\pgfqpoint{2.527801in}{2.514816in}}%
\pgfpathlineto{\pgfqpoint{2.528176in}{2.544503in}}%
\pgfpathlineto{\pgfqpoint{2.528833in}{2.535752in}}%
\pgfpathlineto{\pgfqpoint{2.529495in}{2.510417in}}%
\pgfpathlineto{\pgfqpoint{2.529695in}{2.541027in}}%
\pgfpathlineto{\pgfqpoint{2.530050in}{2.530209in}}%
\pgfpathlineto{\pgfqpoint{2.530055in}{2.530623in}}%
\pgfpathlineto{\pgfqpoint{2.530752in}{2.497525in}}%
\pgfpathlineto{\pgfqpoint{2.531944in}{2.443547in}}%
\pgfpathlineto{\pgfqpoint{2.530937in}{2.509906in}}%
\pgfpathlineto{\pgfqpoint{2.531949in}{2.443874in}}%
\pgfpathlineto{\pgfqpoint{2.533006in}{2.463837in}}%
\pgfpathlineto{\pgfqpoint{2.532382in}{2.438651in}}%
\pgfpathlineto{\pgfqpoint{2.533079in}{2.449727in}}%
\pgfpathlineto{\pgfqpoint{2.533171in}{2.446666in}}%
\pgfpathlineto{\pgfqpoint{2.534096in}{2.470860in}}%
\pgfpathlineto{\pgfqpoint{2.534116in}{2.466593in}}%
\pgfpathlineto{\pgfqpoint{2.535201in}{2.472176in}}%
\pgfpathlineto{\pgfqpoint{2.534452in}{2.447037in}}%
\pgfpathlineto{\pgfqpoint{2.535226in}{2.466217in}}%
\pgfpathlineto{\pgfqpoint{2.535231in}{2.466193in}}%
\pgfpathlineto{\pgfqpoint{2.535250in}{2.470206in}}%
\pgfpathlineto{\pgfqpoint{2.535265in}{2.470078in}}%
\pgfpathlineto{\pgfqpoint{2.535795in}{2.491345in}}%
\pgfpathlineto{\pgfqpoint{2.536117in}{2.449721in}}%
\pgfpathlineto{\pgfqpoint{2.536350in}{2.455171in}}%
\pgfpathlineto{\pgfqpoint{2.536414in}{2.457080in}}%
\pgfpathlineto{\pgfqpoint{2.537514in}{2.416186in}}%
\pgfpathlineto{\pgfqpoint{2.537777in}{2.426425in}}%
\pgfpathlineto{\pgfqpoint{2.538527in}{2.405958in}}%
\pgfpathlineto{\pgfqpoint{2.538624in}{2.419087in}}%
\pgfpathlineto{\pgfqpoint{2.538833in}{2.429829in}}%
\pgfpathlineto{\pgfqpoint{2.538931in}{2.409124in}}%
\pgfpathlineto{\pgfqpoint{2.539018in}{2.411204in}}%
\pgfpathlineto{\pgfqpoint{2.539135in}{2.402476in}}%
\pgfpathlineto{\pgfqpoint{2.539797in}{2.453172in}}%
\pgfpathlineto{\pgfqpoint{2.540007in}{2.451612in}}%
\pgfpathlineto{\pgfqpoint{2.542850in}{2.364444in}}%
\pgfpathlineto{\pgfqpoint{2.540045in}{2.454268in}}%
\pgfpathlineto{\pgfqpoint{2.542942in}{2.381075in}}%
\pgfpathlineto{\pgfqpoint{2.542947in}{2.382521in}}%
\pgfpathlineto{\pgfqpoint{2.543410in}{2.341610in}}%
\pgfpathlineto{\pgfqpoint{2.544008in}{2.370457in}}%
\pgfpathlineto{\pgfqpoint{2.545041in}{2.322025in}}%
\pgfpathlineto{\pgfqpoint{2.544081in}{2.374395in}}%
\pgfpathlineto{\pgfqpoint{2.545211in}{2.326254in}}%
\pgfpathlineto{\pgfqpoint{2.545673in}{2.354850in}}%
\pgfpathlineto{\pgfqpoint{2.546277in}{2.301498in}}%
\pgfpathlineto{\pgfqpoint{2.546297in}{2.307703in}}%
\pgfpathlineto{\pgfqpoint{2.546306in}{2.309839in}}%
\pgfpathlineto{\pgfqpoint{2.547363in}{2.343964in}}%
\pgfpathlineto{\pgfqpoint{2.546920in}{2.298300in}}%
\pgfpathlineto{\pgfqpoint{2.547494in}{2.334099in}}%
\pgfpathlineto{\pgfqpoint{2.548823in}{2.391621in}}%
\pgfpathlineto{\pgfqpoint{2.547626in}{2.326743in}}%
\pgfpathlineto{\pgfqpoint{2.548955in}{2.376514in}}%
\pgfpathlineto{\pgfqpoint{2.549554in}{2.356112in}}%
\pgfpathlineto{\pgfqpoint{2.549091in}{2.391762in}}%
\pgfpathlineto{\pgfqpoint{2.550079in}{2.370466in}}%
\pgfpathlineto{\pgfqpoint{2.551180in}{2.395308in}}%
\pgfpathlineto{\pgfqpoint{2.551233in}{2.389757in}}%
\pgfpathlineto{\pgfqpoint{2.551988in}{2.373491in}}%
\pgfpathlineto{\pgfqpoint{2.551555in}{2.409426in}}%
\pgfpathlineto{\pgfqpoint{2.552353in}{2.384680in}}%
\pgfpathlineto{\pgfqpoint{2.552421in}{2.400807in}}%
\pgfpathlineto{\pgfqpoint{2.552971in}{2.365097in}}%
\pgfpathlineto{\pgfqpoint{2.553458in}{2.380609in}}%
\pgfpathlineto{\pgfqpoint{2.554578in}{2.340661in}}%
\pgfpathlineto{\pgfqpoint{2.554622in}{2.345244in}}%
\pgfpathlineto{\pgfqpoint{2.555717in}{2.370205in}}%
\pgfpathlineto{\pgfqpoint{2.554817in}{2.334731in}}%
\pgfpathlineto{\pgfqpoint{2.555746in}{2.361697in}}%
\pgfpathlineto{\pgfqpoint{2.556910in}{2.337464in}}%
\pgfpathlineto{\pgfqpoint{2.556238in}{2.381093in}}%
\pgfpathlineto{\pgfqpoint{2.557002in}{2.341760in}}%
\pgfpathlineto{\pgfqpoint{2.557562in}{2.380929in}}%
\pgfpathlineto{\pgfqpoint{2.557046in}{2.335058in}}%
\pgfpathlineto{\pgfqpoint{2.558151in}{2.360478in}}%
\pgfpathlineto{\pgfqpoint{2.559257in}{2.356032in}}%
\pgfpathlineto{\pgfqpoint{2.558857in}{2.398910in}}%
\pgfpathlineto{\pgfqpoint{2.559266in}{2.358493in}}%
\pgfpathlineto{\pgfqpoint{2.559276in}{2.359652in}}%
\pgfpathlineto{\pgfqpoint{2.559885in}{2.319116in}}%
\pgfpathlineto{\pgfqpoint{2.559914in}{2.322205in}}%
\pgfpathlineto{\pgfqpoint{2.559977in}{2.306374in}}%
\pgfpathlineto{\pgfqpoint{2.560396in}{2.337120in}}%
\pgfpathlineto{\pgfqpoint{2.561024in}{2.321016in}}%
\pgfpathlineto{\pgfqpoint{2.562236in}{2.360423in}}%
\pgfpathlineto{\pgfqpoint{2.561399in}{2.304302in}}%
\pgfpathlineto{\pgfqpoint{2.562509in}{2.337497in}}%
\pgfpathlineto{\pgfqpoint{2.563341in}{2.317646in}}%
\pgfpathlineto{\pgfqpoint{2.563112in}{2.342718in}}%
\pgfpathlineto{\pgfqpoint{2.563667in}{2.322764in}}%
\pgfpathlineto{\pgfqpoint{2.563716in}{2.330330in}}%
\pgfpathlineto{\pgfqpoint{2.564607in}{2.296190in}}%
\pgfpathlineto{\pgfqpoint{2.565561in}{2.254823in}}%
\pgfpathlineto{\pgfqpoint{2.565756in}{2.268485in}}%
\pgfpathlineto{\pgfqpoint{2.565922in}{2.246465in}}%
\pgfpathlineto{\pgfqpoint{2.566233in}{2.278004in}}%
\pgfpathlineto{\pgfqpoint{2.566881in}{2.262464in}}%
\pgfpathlineto{\pgfqpoint{2.567329in}{2.272097in}}%
\pgfpathlineto{\pgfqpoint{2.567650in}{2.240206in}}%
\pgfpathlineto{\pgfqpoint{2.567923in}{2.251333in}}%
\pgfpathlineto{\pgfqpoint{2.568195in}{2.230131in}}%
\pgfpathlineto{\pgfqpoint{2.568604in}{2.254226in}}%
\pgfpathlineto{\pgfqpoint{2.569047in}{2.240647in}}%
\pgfpathlineto{\pgfqpoint{2.570269in}{2.284078in}}%
\pgfpathlineto{\pgfqpoint{2.569300in}{2.226701in}}%
\pgfpathlineto{\pgfqpoint{2.570293in}{2.279022in}}%
\pgfpathlineto{\pgfqpoint{2.571306in}{2.245369in}}%
\pgfpathlineto{\pgfqpoint{2.570839in}{2.295529in}}%
\pgfpathlineto{\pgfqpoint{2.571496in}{2.249250in}}%
\pgfpathlineto{\pgfqpoint{2.571657in}{2.239012in}}%
\pgfpathlineto{\pgfqpoint{2.572747in}{2.302363in}}%
\pgfpathlineto{\pgfqpoint{2.573804in}{2.295635in}}%
\pgfpathlineto{\pgfqpoint{2.573624in}{2.324470in}}%
\pgfpathlineto{\pgfqpoint{2.573852in}{2.303366in}}%
\pgfpathlineto{\pgfqpoint{2.574359in}{2.309285in}}%
\pgfpathlineto{\pgfqpoint{2.573930in}{2.284886in}}%
\pgfpathlineto{\pgfqpoint{2.574953in}{2.302375in}}%
\pgfpathlineto{\pgfqpoint{2.575444in}{2.296805in}}%
\pgfpathlineto{\pgfqpoint{2.575844in}{2.319988in}}%
\pgfpathlineto{\pgfqpoint{2.576019in}{2.314478in}}%
\pgfpathlineto{\pgfqpoint{2.576043in}{2.324769in}}%
\pgfpathlineto{\pgfqpoint{2.577066in}{2.282583in}}%
\pgfpathlineto{\pgfqpoint{2.577669in}{2.258574in}}%
\pgfpathlineto{\pgfqpoint{2.577436in}{2.286686in}}%
\pgfpathlineto{\pgfqpoint{2.578190in}{2.266157in}}%
\pgfpathlineto{\pgfqpoint{2.579281in}{2.298226in}}%
\pgfpathlineto{\pgfqpoint{2.579364in}{2.285874in}}%
\pgfpathlineto{\pgfqpoint{2.579398in}{2.279545in}}%
\pgfpathlineto{\pgfqpoint{2.579992in}{2.311977in}}%
\pgfpathlineto{\pgfqpoint{2.580391in}{2.306034in}}%
\pgfpathlineto{\pgfqpoint{2.580430in}{2.310588in}}%
\pgfpathlineto{\pgfqpoint{2.581287in}{2.256914in}}%
\pgfpathlineto{\pgfqpoint{2.581399in}{2.267911in}}%
\pgfpathlineto{\pgfqpoint{2.582139in}{2.236287in}}%
\pgfpathlineto{\pgfqpoint{2.581593in}{2.271686in}}%
\pgfpathlineto{\pgfqpoint{2.582543in}{2.247288in}}%
\pgfpathlineto{\pgfqpoint{2.583555in}{2.263492in}}%
\pgfpathlineto{\pgfqpoint{2.583073in}{2.232616in}}%
\pgfpathlineto{\pgfqpoint{2.583692in}{2.255188in}}%
\pgfpathlineto{\pgfqpoint{2.584714in}{2.216792in}}%
\pgfpathlineto{\pgfqpoint{2.584144in}{2.273046in}}%
\pgfpathlineto{\pgfqpoint{2.584894in}{2.218825in}}%
\pgfpathlineto{\pgfqpoint{2.585975in}{2.261584in}}%
\pgfpathlineto{\pgfqpoint{2.585235in}{2.206408in}}%
\pgfpathlineto{\pgfqpoint{2.586121in}{2.250020in}}%
\pgfpathlineto{\pgfqpoint{2.587207in}{2.228486in}}%
\pgfpathlineto{\pgfqpoint{2.587246in}{2.239900in}}%
\pgfpathlineto{\pgfqpoint{2.587362in}{2.230020in}}%
\pgfpathlineto{\pgfqpoint{2.587280in}{2.240303in}}%
\pgfpathlineto{\pgfqpoint{2.587367in}{2.230264in}}%
\pgfpathlineto{\pgfqpoint{2.587684in}{2.203344in}}%
\pgfpathlineto{\pgfqpoint{2.588482in}{2.225369in}}%
\pgfpathlineto{\pgfqpoint{2.588935in}{2.219001in}}%
\pgfpathlineto{\pgfqpoint{2.588662in}{2.248368in}}%
\pgfpathlineto{\pgfqpoint{2.589373in}{2.239154in}}%
\pgfpathlineto{\pgfqpoint{2.590079in}{2.269584in}}%
\pgfpathlineto{\pgfqpoint{2.589509in}{2.239145in}}%
\pgfpathlineto{\pgfqpoint{2.590512in}{2.245190in}}%
\pgfpathlineto{\pgfqpoint{2.592061in}{2.188488in}}%
\pgfpathlineto{\pgfqpoint{2.590527in}{2.249852in}}%
\pgfpathlineto{\pgfqpoint{2.592966in}{2.211988in}}%
\pgfpathlineto{\pgfqpoint{2.594149in}{2.247074in}}%
\pgfpathlineto{\pgfqpoint{2.594198in}{2.242044in}}%
\pgfpathlineto{\pgfqpoint{2.594232in}{2.233814in}}%
\pgfpathlineto{\pgfqpoint{2.595269in}{2.278729in}}%
\pgfpathlineto{\pgfqpoint{2.596681in}{2.340009in}}%
\pgfpathlineto{\pgfqpoint{2.595332in}{2.273429in}}%
\pgfpathlineto{\pgfqpoint{2.596691in}{2.337635in}}%
\pgfpathlineto{\pgfqpoint{2.598171in}{2.288224in}}%
\pgfpathlineto{\pgfqpoint{2.596793in}{2.346298in}}%
\pgfpathlineto{\pgfqpoint{2.598312in}{2.299636in}}%
\pgfpathlineto{\pgfqpoint{2.599120in}{2.338847in}}%
\pgfpathlineto{\pgfqpoint{2.598356in}{2.298535in}}%
\pgfpathlineto{\pgfqpoint{2.599436in}{2.302994in}}%
\pgfpathlineto{\pgfqpoint{2.599456in}{2.297612in}}%
\pgfpathlineto{\pgfqpoint{2.600342in}{2.320032in}}%
\pgfpathlineto{\pgfqpoint{2.600454in}{2.312196in}}%
\pgfpathlineto{\pgfqpoint{2.600921in}{2.323540in}}%
\pgfpathlineto{\pgfqpoint{2.601199in}{2.300492in}}%
\pgfpathlineto{\pgfqpoint{2.601569in}{2.320530in}}%
\pgfpathlineto{\pgfqpoint{2.602718in}{2.331103in}}%
\pgfpathlineto{\pgfqpoint{2.602060in}{2.301838in}}%
\pgfpathlineto{\pgfqpoint{2.602732in}{2.330807in}}%
\pgfpathlineto{\pgfqpoint{2.602800in}{2.313911in}}%
\pgfpathlineto{\pgfqpoint{2.603414in}{2.366409in}}%
\pgfpathlineto{\pgfqpoint{2.603769in}{2.359683in}}%
\pgfpathlineto{\pgfqpoint{2.604577in}{2.371020in}}%
\pgfpathlineto{\pgfqpoint{2.604427in}{2.337182in}}%
\pgfpathlineto{\pgfqpoint{2.604860in}{2.358531in}}%
\pgfpathlineto{\pgfqpoint{2.604918in}{2.355304in}}%
\pgfpathlineto{\pgfqpoint{2.605741in}{2.422028in}}%
\pgfpathlineto{\pgfqpoint{2.605887in}{2.392694in}}%
\pgfpathlineto{\pgfqpoint{2.606476in}{2.405915in}}%
\pgfpathlineto{\pgfqpoint{2.606145in}{2.388616in}}%
\pgfpathlineto{\pgfqpoint{2.606710in}{2.388975in}}%
\pgfpathlineto{\pgfqpoint{2.606983in}{2.367421in}}%
\pgfpathlineto{\pgfqpoint{2.607654in}{2.437623in}}%
\pgfpathlineto{\pgfqpoint{2.608842in}{2.460764in}}%
\pgfpathlineto{\pgfqpoint{2.607873in}{2.431170in}}%
\pgfpathlineto{\pgfqpoint{2.608862in}{2.455454in}}%
\pgfpathlineto{\pgfqpoint{2.609524in}{2.443308in}}%
\pgfpathlineto{\pgfqpoint{2.609738in}{2.470652in}}%
\pgfpathlineto{\pgfqpoint{2.609986in}{2.448587in}}%
\pgfpathlineto{\pgfqpoint{2.610137in}{2.462406in}}%
\pgfpathlineto{\pgfqpoint{2.610571in}{2.435855in}}%
\pgfpathlineto{\pgfqpoint{2.611111in}{2.455673in}}%
\pgfpathlineto{\pgfqpoint{2.611953in}{2.475555in}}%
\pgfpathlineto{\pgfqpoint{2.611539in}{2.427721in}}%
\pgfpathlineto{\pgfqpoint{2.612187in}{2.453392in}}%
\pgfpathlineto{\pgfqpoint{2.613341in}{2.393126in}}%
\pgfpathlineto{\pgfqpoint{2.612503in}{2.474720in}}%
\pgfpathlineto{\pgfqpoint{2.613433in}{2.398532in}}%
\pgfpathlineto{\pgfqpoint{2.613458in}{2.395121in}}%
\pgfpathlineto{\pgfqpoint{2.613988in}{2.438925in}}%
\pgfpathlineto{\pgfqpoint{2.614008in}{2.436158in}}%
\pgfpathlineto{\pgfqpoint{2.614879in}{2.464349in}}%
\pgfpathlineto{\pgfqpoint{2.614115in}{2.430928in}}%
\pgfpathlineto{\pgfqpoint{2.615166in}{2.456853in}}%
\pgfpathlineto{\pgfqpoint{2.616121in}{2.438577in}}%
\pgfpathlineto{\pgfqpoint{2.615843in}{2.469005in}}%
\pgfpathlineto{\pgfqpoint{2.616301in}{2.447238in}}%
\pgfpathlineto{\pgfqpoint{2.616486in}{2.459909in}}%
\pgfpathlineto{\pgfqpoint{2.617202in}{2.416146in}}%
\pgfpathlineto{\pgfqpoint{2.617387in}{2.436015in}}%
\pgfpathlineto{\pgfqpoint{2.617849in}{2.450391in}}%
\pgfpathlineto{\pgfqpoint{2.617528in}{2.425986in}}%
\pgfpathlineto{\pgfqpoint{2.618170in}{2.429786in}}%
\pgfpathlineto{\pgfqpoint{2.618920in}{2.405269in}}%
\pgfpathlineto{\pgfqpoint{2.618195in}{2.434484in}}%
\pgfpathlineto{\pgfqpoint{2.619280in}{2.429980in}}%
\pgfpathlineto{\pgfqpoint{2.619918in}{2.479979in}}%
\pgfpathlineto{\pgfqpoint{2.620507in}{2.471659in}}%
\pgfpathlineto{\pgfqpoint{2.620804in}{2.453155in}}%
\pgfpathlineto{\pgfqpoint{2.621632in}{2.468826in}}%
\pgfpathlineto{\pgfqpoint{2.622708in}{2.482321in}}%
\pgfpathlineto{\pgfqpoint{2.622060in}{2.461466in}}%
\pgfpathlineto{\pgfqpoint{2.622761in}{2.472686in}}%
\pgfpathlineto{\pgfqpoint{2.623910in}{2.442002in}}%
\pgfpathlineto{\pgfqpoint{2.622810in}{2.476584in}}%
\pgfpathlineto{\pgfqpoint{2.623998in}{2.451946in}}%
\pgfpathlineto{\pgfqpoint{2.625152in}{2.511171in}}%
\pgfpathlineto{\pgfqpoint{2.625220in}{2.492425in}}%
\pgfpathlineto{\pgfqpoint{2.625235in}{2.487170in}}%
\pgfpathlineto{\pgfqpoint{2.626164in}{2.538781in}}%
\pgfpathlineto{\pgfqpoint{2.626267in}{2.526554in}}%
\pgfpathlineto{\pgfqpoint{2.626491in}{2.549791in}}%
\pgfpathlineto{\pgfqpoint{2.626398in}{2.516627in}}%
\pgfpathlineto{\pgfqpoint{2.627284in}{2.517762in}}%
\pgfpathlineto{\pgfqpoint{2.627864in}{2.502669in}}%
\pgfpathlineto{\pgfqpoint{2.627596in}{2.527228in}}%
\pgfpathlineto{\pgfqpoint{2.628394in}{2.517085in}}%
\pgfpathlineto{\pgfqpoint{2.628642in}{2.512439in}}%
\pgfpathlineto{\pgfqpoint{2.628925in}{2.545875in}}%
\pgfpathlineto{\pgfqpoint{2.629392in}{2.532810in}}%
\pgfpathlineto{\pgfqpoint{2.629679in}{2.518192in}}%
\pgfpathlineto{\pgfqpoint{2.630512in}{2.539487in}}%
\pgfpathlineto{\pgfqpoint{2.631471in}{2.503406in}}%
\pgfpathlineto{\pgfqpoint{2.630639in}{2.559124in}}%
\pgfpathlineto{\pgfqpoint{2.631680in}{2.508848in}}%
\pgfpathlineto{\pgfqpoint{2.631831in}{2.521509in}}%
\pgfpathlineto{\pgfqpoint{2.632576in}{2.496328in}}%
\pgfpathlineto{\pgfqpoint{2.632713in}{2.501090in}}%
\pgfpathlineto{\pgfqpoint{2.632873in}{2.485669in}}%
\pgfpathlineto{\pgfqpoint{2.633229in}{2.525890in}}%
\pgfpathlineto{\pgfqpoint{2.633808in}{2.511496in}}%
\pgfpathlineto{\pgfqpoint{2.634640in}{2.564049in}}%
\pgfpathlineto{\pgfqpoint{2.633832in}{2.509251in}}%
\pgfpathlineto{\pgfqpoint{2.635249in}{2.533689in}}%
\pgfpathlineto{\pgfqpoint{2.635955in}{2.518051in}}%
\pgfpathlineto{\pgfqpoint{2.636232in}{2.546318in}}%
\pgfpathlineto{\pgfqpoint{2.636349in}{2.536797in}}%
\pgfpathlineto{\pgfqpoint{2.639178in}{2.610053in}}%
\pgfpathlineto{\pgfqpoint{2.636515in}{2.522729in}}%
\pgfpathlineto{\pgfqpoint{2.640035in}{2.587023in}}%
\pgfpathlineto{\pgfqpoint{2.641018in}{2.553441in}}%
\pgfpathlineto{\pgfqpoint{2.640259in}{2.595288in}}%
\pgfpathlineto{\pgfqpoint{2.641257in}{2.564089in}}%
\pgfpathlineto{\pgfqpoint{2.641364in}{2.570200in}}%
\pgfpathlineto{\pgfqpoint{2.642016in}{2.528053in}}%
\pgfpathlineto{\pgfqpoint{2.642026in}{2.524665in}}%
\pgfpathlineto{\pgfqpoint{2.642513in}{2.557063in}}%
\pgfpathlineto{\pgfqpoint{2.643082in}{2.552229in}}%
\pgfpathlineto{\pgfqpoint{2.643233in}{2.557593in}}%
\pgfpathlineto{\pgfqpoint{2.643686in}{2.533505in}}%
\pgfpathlineto{\pgfqpoint{2.643959in}{2.534426in}}%
\pgfpathlineto{\pgfqpoint{2.645010in}{2.512855in}}%
\pgfpathlineto{\pgfqpoint{2.644553in}{2.544440in}}%
\pgfpathlineto{\pgfqpoint{2.645108in}{2.518735in}}%
\pgfpathlineto{\pgfqpoint{2.645395in}{2.534908in}}%
\pgfpathlineto{\pgfqpoint{2.646169in}{2.505193in}}%
\pgfpathlineto{\pgfqpoint{2.646189in}{2.509040in}}%
\pgfpathlineto{\pgfqpoint{2.646899in}{2.477876in}}%
\pgfpathlineto{\pgfqpoint{2.646218in}{2.509439in}}%
\pgfpathlineto{\pgfqpoint{2.647338in}{2.493179in}}%
\pgfpathlineto{\pgfqpoint{2.647941in}{2.523297in}}%
\pgfpathlineto{\pgfqpoint{2.647508in}{2.478355in}}%
\pgfpathlineto{\pgfqpoint{2.648462in}{2.507398in}}%
\pgfpathlineto{\pgfqpoint{2.649489in}{2.464791in}}%
\pgfpathlineto{\pgfqpoint{2.648545in}{2.509175in}}%
\pgfpathlineto{\pgfqpoint{2.649694in}{2.471895in}}%
\pgfpathlineto{\pgfqpoint{2.650385in}{2.526724in}}%
\pgfpathlineto{\pgfqpoint{2.650848in}{2.509705in}}%
\pgfpathlineto{\pgfqpoint{2.651159in}{2.489177in}}%
\pgfpathlineto{\pgfqpoint{2.651597in}{2.545487in}}%
\pgfpathlineto{\pgfqpoint{2.651656in}{2.544634in}}%
\pgfpathlineto{\pgfqpoint{2.651768in}{2.558137in}}%
\pgfpathlineto{\pgfqpoint{2.652206in}{2.528567in}}%
\pgfpathlineto{\pgfqpoint{2.652289in}{2.529850in}}%
\pgfpathlineto{\pgfqpoint{2.653258in}{2.493417in}}%
\pgfpathlineto{\pgfqpoint{2.652328in}{2.530287in}}%
\pgfpathlineto{\pgfqpoint{2.653418in}{2.516216in}}%
\pgfpathlineto{\pgfqpoint{2.654606in}{2.561532in}}%
\pgfpathlineto{\pgfqpoint{2.653433in}{2.514562in}}%
\pgfpathlineto{\pgfqpoint{2.654781in}{2.543216in}}%
\pgfpathlineto{\pgfqpoint{2.655371in}{2.515575in}}%
\pgfpathlineto{\pgfqpoint{2.655833in}{2.549436in}}%
\pgfpathlineto{\pgfqpoint{2.655926in}{2.535018in}}%
\pgfpathlineto{\pgfqpoint{2.657070in}{2.584959in}}%
\pgfpathlineto{\pgfqpoint{2.656081in}{2.529428in}}%
\pgfpathlineto{\pgfqpoint{2.657113in}{2.577452in}}%
\pgfpathlineto{\pgfqpoint{2.657391in}{2.561784in}}%
\pgfpathlineto{\pgfqpoint{2.657922in}{2.621170in}}%
\pgfpathlineto{\pgfqpoint{2.658150in}{2.613321in}}%
\pgfpathlineto{\pgfqpoint{2.658671in}{2.631497in}}%
\pgfpathlineto{\pgfqpoint{2.659158in}{2.586888in}}%
\pgfpathlineto{\pgfqpoint{2.659231in}{2.588881in}}%
\pgfpathlineto{\pgfqpoint{2.659752in}{2.574364in}}%
\pgfpathlineto{\pgfqpoint{2.659981in}{2.617088in}}%
\pgfpathlineto{\pgfqpoint{2.660302in}{2.640846in}}%
\pgfpathlineto{\pgfqpoint{2.660721in}{2.608734in}}%
\pgfpathlineto{\pgfqpoint{2.661120in}{2.635861in}}%
\pgfpathlineto{\pgfqpoint{2.661145in}{2.627006in}}%
\pgfpathlineto{\pgfqpoint{2.661349in}{2.649335in}}%
\pgfpathlineto{\pgfqpoint{2.662221in}{2.639907in}}%
\pgfpathlineto{\pgfqpoint{2.663467in}{2.671112in}}%
\pgfpathlineto{\pgfqpoint{2.662420in}{2.620399in}}%
\pgfpathlineto{\pgfqpoint{2.663559in}{2.662169in}}%
\pgfpathlineto{\pgfqpoint{2.664635in}{2.645063in}}%
\pgfpathlineto{\pgfqpoint{2.663954in}{2.680340in}}%
\pgfpathlineto{\pgfqpoint{2.664694in}{2.653995in}}%
\pgfpathlineto{\pgfqpoint{2.665867in}{2.700526in}}%
\pgfpathlineto{\pgfqpoint{2.666164in}{2.724426in}}%
\pgfpathlineto{\pgfqpoint{2.666378in}{2.692199in}}%
\pgfpathlineto{\pgfqpoint{2.667060in}{2.711912in}}%
\pgfpathlineto{\pgfqpoint{2.667162in}{2.702717in}}%
\pgfpathlineto{\pgfqpoint{2.667814in}{2.741037in}}%
\pgfpathlineto{\pgfqpoint{2.668146in}{2.727727in}}%
\pgfpathlineto{\pgfqpoint{2.668890in}{2.705302in}}%
\pgfpathlineto{\pgfqpoint{2.668243in}{2.743557in}}%
\pgfpathlineto{\pgfqpoint{2.669265in}{2.722680in}}%
\pgfpathlineto{\pgfqpoint{2.670093in}{2.750983in}}%
\pgfpathlineto{\pgfqpoint{2.669333in}{2.719941in}}%
\pgfpathlineto{\pgfqpoint{2.670390in}{2.737599in}}%
\pgfpathlineto{\pgfqpoint{2.670877in}{2.714687in}}%
\pgfpathlineto{\pgfqpoint{2.670672in}{2.739795in}}%
\pgfpathlineto{\pgfqpoint{2.671500in}{2.734146in}}%
\pgfpathlineto{\pgfqpoint{2.672079in}{2.767802in}}%
\pgfpathlineto{\pgfqpoint{2.671797in}{2.732959in}}%
\pgfpathlineto{\pgfqpoint{2.672663in}{2.747477in}}%
\pgfpathlineto{\pgfqpoint{2.673501in}{2.693287in}}%
\pgfpathlineto{\pgfqpoint{2.673876in}{2.705776in}}%
\pgfpathlineto{\pgfqpoint{2.674294in}{2.731426in}}%
\pgfpathlineto{\pgfqpoint{2.674684in}{2.701658in}}%
\pgfpathlineto{\pgfqpoint{2.674986in}{2.709362in}}%
\pgfpathlineto{\pgfqpoint{2.674991in}{2.709613in}}%
\pgfpathlineto{\pgfqpoint{2.675526in}{2.680517in}}%
\pgfpathlineto{\pgfqpoint{2.675555in}{2.681961in}}%
\pgfpathlineto{\pgfqpoint{2.676110in}{2.666549in}}%
\pgfpathlineto{\pgfqpoint{2.676631in}{2.702493in}}%
\pgfpathlineto{\pgfqpoint{2.676641in}{2.700118in}}%
\pgfpathlineto{\pgfqpoint{2.677732in}{2.711228in}}%
\pgfpathlineto{\pgfqpoint{2.677274in}{2.677771in}}%
\pgfpathlineto{\pgfqpoint{2.677761in}{2.705079in}}%
\pgfpathlineto{\pgfqpoint{2.677809in}{2.694670in}}%
\pgfpathlineto{\pgfqpoint{2.678778in}{2.743750in}}%
\pgfpathlineto{\pgfqpoint{2.678788in}{2.738982in}}%
\pgfpathlineto{\pgfqpoint{2.680346in}{2.808807in}}%
\pgfpathlineto{\pgfqpoint{2.680361in}{2.807458in}}%
\pgfpathlineto{\pgfqpoint{2.681471in}{2.776038in}}%
\pgfpathlineto{\pgfqpoint{2.680940in}{2.811768in}}%
\pgfpathlineto{\pgfqpoint{2.681505in}{2.785042in}}%
\pgfpathlineto{\pgfqpoint{2.682561in}{2.793609in}}%
\pgfpathlineto{\pgfqpoint{2.681967in}{2.758050in}}%
\pgfpathlineto{\pgfqpoint{2.682624in}{2.791365in}}%
\pgfpathlineto{\pgfqpoint{2.683681in}{2.742798in}}%
\pgfpathlineto{\pgfqpoint{2.683803in}{2.748778in}}%
\pgfpathlineto{\pgfqpoint{2.685059in}{2.793657in}}%
\pgfpathlineto{\pgfqpoint{2.683944in}{2.740098in}}%
\pgfpathlineto{\pgfqpoint{2.685068in}{2.793198in}}%
\pgfpathlineto{\pgfqpoint{2.685463in}{2.775853in}}%
\pgfpathlineto{\pgfqpoint{2.686071in}{2.811504in}}%
\pgfpathlineto{\pgfqpoint{2.686125in}{2.809537in}}%
\pgfpathlineto{\pgfqpoint{2.686174in}{2.820155in}}%
\pgfpathlineto{\pgfqpoint{2.686558in}{2.789413in}}%
\pgfpathlineto{\pgfqpoint{2.687225in}{2.801466in}}%
\pgfpathlineto{\pgfqpoint{2.687546in}{2.829228in}}%
\pgfpathlineto{\pgfqpoint{2.688364in}{2.812418in}}%
\pgfpathlineto{\pgfqpoint{2.689538in}{2.849118in}}%
\pgfpathlineto{\pgfqpoint{2.688457in}{2.808846in}}%
\pgfpathlineto{\pgfqpoint{2.689674in}{2.839761in}}%
\pgfpathlineto{\pgfqpoint{2.690078in}{2.826843in}}%
\pgfpathlineto{\pgfqpoint{2.690740in}{2.862837in}}%
\pgfpathlineto{\pgfqpoint{2.691295in}{2.874711in}}%
\pgfpathlineto{\pgfqpoint{2.691607in}{2.829967in}}%
\pgfpathlineto{\pgfqpoint{2.691777in}{2.834853in}}%
\pgfpathlineto{\pgfqpoint{2.692288in}{2.812719in}}%
\pgfpathlineto{\pgfqpoint{2.692858in}{2.836651in}}%
\pgfpathlineto{\pgfqpoint{2.692887in}{2.834123in}}%
\pgfpathlineto{\pgfqpoint{2.693418in}{2.839470in}}%
\pgfpathlineto{\pgfqpoint{2.693919in}{2.787405in}}%
\pgfpathlineto{\pgfqpoint{2.693934in}{2.791028in}}%
\pgfpathlineto{\pgfqpoint{2.694026in}{2.800478in}}%
\pgfpathlineto{\pgfqpoint{2.694625in}{2.757284in}}%
\pgfpathlineto{\pgfqpoint{2.694674in}{2.758383in}}%
\pgfpathlineto{\pgfqpoint{2.694723in}{2.753959in}}%
\pgfpathlineto{\pgfqpoint{2.695628in}{2.776306in}}%
\pgfpathlineto{\pgfqpoint{2.695653in}{2.783675in}}%
\pgfpathlineto{\pgfqpoint{2.696368in}{2.756948in}}%
\pgfpathlineto{\pgfqpoint{2.696724in}{2.762741in}}%
\pgfpathlineto{\pgfqpoint{2.697391in}{2.750287in}}%
\pgfpathlineto{\pgfqpoint{2.696821in}{2.771201in}}%
\pgfpathlineto{\pgfqpoint{2.697532in}{2.769821in}}%
\pgfpathlineto{\pgfqpoint{2.698218in}{2.809134in}}%
\pgfpathlineto{\pgfqpoint{2.697585in}{2.760466in}}%
\pgfpathlineto{\pgfqpoint{2.698700in}{2.787478in}}%
\pgfpathlineto{\pgfqpoint{2.699596in}{2.756418in}}%
\pgfpathlineto{\pgfqpoint{2.699119in}{2.789610in}}%
\pgfpathlineto{\pgfqpoint{2.699825in}{2.771905in}}%
\pgfpathlineto{\pgfqpoint{2.700872in}{2.797026in}}%
\pgfpathlineto{\pgfqpoint{2.700409in}{2.764543in}}%
\pgfpathlineto{\pgfqpoint{2.700993in}{2.782657in}}%
\pgfpathlineto{\pgfqpoint{2.701539in}{2.764669in}}%
\pgfpathlineto{\pgfqpoint{2.701977in}{2.795823in}}%
\pgfpathlineto{\pgfqpoint{2.701991in}{2.793657in}}%
\pgfpathlineto{\pgfqpoint{2.703208in}{2.829876in}}%
\pgfpathlineto{\pgfqpoint{2.702069in}{2.791702in}}%
\pgfpathlineto{\pgfqpoint{2.703223in}{2.825911in}}%
\pgfpathlineto{\pgfqpoint{2.703768in}{2.815336in}}%
\pgfpathlineto{\pgfqpoint{2.704245in}{2.840533in}}%
\pgfpathlineto{\pgfqpoint{2.704304in}{2.836931in}}%
\pgfpathlineto{\pgfqpoint{2.705200in}{2.869236in}}%
\pgfpathlineto{\pgfqpoint{2.704474in}{2.827764in}}%
\pgfpathlineto{\pgfqpoint{2.705458in}{2.861805in}}%
\pgfpathlineto{\pgfqpoint{2.705521in}{2.842431in}}%
\pgfpathlineto{\pgfqpoint{2.706207in}{2.884347in}}%
\pgfpathlineto{\pgfqpoint{2.706524in}{2.879828in}}%
\pgfpathlineto{\pgfqpoint{2.706612in}{2.885601in}}%
\pgfpathlineto{\pgfqpoint{2.707386in}{2.862066in}}%
\pgfpathlineto{\pgfqpoint{2.707566in}{2.873030in}}%
\pgfpathlineto{\pgfqpoint{2.708574in}{2.854734in}}%
\pgfpathlineto{\pgfqpoint{2.707975in}{2.876694in}}%
\pgfpathlineto{\pgfqpoint{2.708695in}{2.864993in}}%
\pgfpathlineto{\pgfqpoint{2.708700in}{2.864379in}}%
\pgfpathlineto{\pgfqpoint{2.709314in}{2.898288in}}%
\pgfpathlineto{\pgfqpoint{2.709601in}{2.889186in}}%
\pgfpathlineto{\pgfqpoint{2.709737in}{2.899221in}}%
\pgfpathlineto{\pgfqpoint{2.710579in}{2.836798in}}%
\pgfpathlineto{\pgfqpoint{2.710638in}{2.838599in}}%
\pgfpathlineto{\pgfqpoint{2.710803in}{2.849067in}}%
\pgfpathlineto{\pgfqpoint{2.710735in}{2.831880in}}%
\pgfpathlineto{\pgfqpoint{2.711008in}{2.833346in}}%
\pgfpathlineto{\pgfqpoint{2.711013in}{2.830152in}}%
\pgfpathlineto{\pgfqpoint{2.711407in}{2.861696in}}%
\pgfpathlineto{\pgfqpoint{2.712093in}{2.849218in}}%
\pgfpathlineto{\pgfqpoint{2.713218in}{2.813074in}}%
\pgfpathlineto{\pgfqpoint{2.713272in}{2.821391in}}%
\pgfpathlineto{\pgfqpoint{2.714031in}{2.842431in}}%
\pgfpathlineto{\pgfqpoint{2.713330in}{2.812377in}}%
\pgfpathlineto{\pgfqpoint{2.714416in}{2.832428in}}%
\pgfpathlineto{\pgfqpoint{2.714586in}{2.826278in}}%
\pgfpathlineto{\pgfqpoint{2.714985in}{2.860672in}}%
\pgfpathlineto{\pgfqpoint{2.716047in}{2.894552in}}%
\pgfpathlineto{\pgfqpoint{2.715540in}{2.844851in}}%
\pgfpathlineto{\pgfqpoint{2.716139in}{2.884595in}}%
\pgfpathlineto{\pgfqpoint{2.717074in}{2.876867in}}%
\pgfpathlineto{\pgfqpoint{2.716597in}{2.902666in}}%
\pgfpathlineto{\pgfqpoint{2.717210in}{2.894508in}}%
\pgfpathlineto{\pgfqpoint{2.717692in}{2.937089in}}%
\pgfpathlineto{\pgfqpoint{2.717254in}{2.891098in}}%
\pgfpathlineto{\pgfqpoint{2.718369in}{2.907070in}}%
\pgfpathlineto{\pgfqpoint{2.719304in}{2.874484in}}%
\pgfpathlineto{\pgfqpoint{2.718423in}{2.909246in}}%
\pgfpathlineto{\pgfqpoint{2.719528in}{2.891410in}}%
\pgfpathlineto{\pgfqpoint{2.719732in}{2.901194in}}%
\pgfpathlineto{\pgfqpoint{2.720302in}{2.880383in}}%
\pgfpathlineto{\pgfqpoint{2.720628in}{2.889580in}}%
\pgfpathlineto{\pgfqpoint{2.720647in}{2.884349in}}%
\pgfpathlineto{\pgfqpoint{2.721582in}{2.920989in}}%
\pgfpathlineto{\pgfqpoint{2.721641in}{2.915289in}}%
\pgfpathlineto{\pgfqpoint{2.722746in}{2.965761in}}%
\pgfpathlineto{\pgfqpoint{2.722867in}{2.951848in}}%
\pgfpathlineto{\pgfqpoint{2.722872in}{2.950789in}}%
\pgfpathlineto{\pgfqpoint{2.723797in}{2.986979in}}%
\pgfpathlineto{\pgfqpoint{2.723827in}{2.983307in}}%
\pgfpathlineto{\pgfqpoint{2.724172in}{2.985891in}}%
\pgfpathlineto{\pgfqpoint{2.724362in}{2.962255in}}%
\pgfpathlineto{\pgfqpoint{2.724888in}{2.964082in}}%
\pgfpathlineto{\pgfqpoint{2.725779in}{2.929787in}}%
\pgfpathlineto{\pgfqpoint{2.725170in}{2.989179in}}%
\pgfpathlineto{\pgfqpoint{2.726037in}{2.938960in}}%
\pgfpathlineto{\pgfqpoint{2.726577in}{2.978102in}}%
\pgfpathlineto{\pgfqpoint{2.726090in}{2.936079in}}%
\pgfpathlineto{\pgfqpoint{2.727147in}{2.939109in}}%
\pgfpathlineto{\pgfqpoint{2.727181in}{2.936741in}}%
\pgfpathlineto{\pgfqpoint{2.728052in}{2.977160in}}%
\pgfpathlineto{\pgfqpoint{2.728140in}{2.962716in}}%
\pgfpathlineto{\pgfqpoint{2.728169in}{2.970610in}}%
\pgfpathlineto{\pgfqpoint{2.728953in}{2.878534in}}%
\pgfpathlineto{\pgfqpoint{2.729089in}{2.885987in}}%
\pgfpathlineto{\pgfqpoint{2.729099in}{2.880802in}}%
\pgfpathlineto{\pgfqpoint{2.729440in}{2.907486in}}%
\pgfpathlineto{\pgfqpoint{2.730180in}{2.893276in}}%
\pgfpathlineto{\pgfqpoint{2.730214in}{2.887701in}}%
\pgfpathlineto{\pgfqpoint{2.730764in}{2.923102in}}%
\pgfpathlineto{\pgfqpoint{2.731022in}{2.910109in}}%
\pgfpathlineto{\pgfqpoint{2.731421in}{2.932302in}}%
\pgfpathlineto{\pgfqpoint{2.731894in}{2.886948in}}%
\pgfpathlineto{\pgfqpoint{2.732113in}{2.900681in}}%
\pgfpathlineto{\pgfqpoint{2.732118in}{2.898318in}}%
\pgfpathlineto{\pgfqpoint{2.732741in}{2.930112in}}%
\pgfpathlineto{\pgfqpoint{2.733164in}{2.925047in}}%
\pgfpathlineto{\pgfqpoint{2.733632in}{2.946124in}}%
\pgfpathlineto{\pgfqpoint{2.733846in}{2.922685in}}%
\pgfpathlineto{\pgfqpoint{2.734289in}{2.935196in}}%
\pgfpathlineto{\pgfqpoint{2.734420in}{2.960470in}}%
\pgfpathlineto{\pgfqpoint{2.735326in}{2.912362in}}%
\pgfpathlineto{\pgfqpoint{2.735345in}{2.917446in}}%
\pgfpathlineto{\pgfqpoint{2.735618in}{2.897548in}}%
\pgfpathlineto{\pgfqpoint{2.736246in}{2.950113in}}%
\pgfpathlineto{\pgfqpoint{2.736251in}{2.949457in}}%
\pgfpathlineto{\pgfqpoint{2.736709in}{2.964034in}}%
\pgfpathlineto{\pgfqpoint{2.736348in}{2.940955in}}%
\pgfpathlineto{\pgfqpoint{2.737385in}{2.960217in}}%
\pgfpathlineto{\pgfqpoint{2.738305in}{2.918549in}}%
\pgfpathlineto{\pgfqpoint{2.738622in}{2.935361in}}%
\pgfpathlineto{\pgfqpoint{2.739893in}{3.002503in}}%
\pgfpathlineto{\pgfqpoint{2.739046in}{2.917921in}}%
\pgfpathlineto{\pgfqpoint{2.739946in}{2.996802in}}%
\pgfpathlineto{\pgfqpoint{2.740024in}{2.986391in}}%
\pgfpathlineto{\pgfqpoint{2.740866in}{3.019683in}}%
\pgfpathlineto{\pgfqpoint{2.741022in}{3.008329in}}%
\pgfpathlineto{\pgfqpoint{2.741027in}{3.007167in}}%
\pgfpathlineto{\pgfqpoint{2.741918in}{3.051654in}}%
\pgfpathlineto{\pgfqpoint{2.741967in}{3.047632in}}%
\pgfpathlineto{\pgfqpoint{2.742380in}{3.063886in}}%
\pgfpathlineto{\pgfqpoint{2.742692in}{3.035778in}}%
\pgfpathlineto{\pgfqpoint{2.743062in}{3.044094in}}%
\pgfpathlineto{\pgfqpoint{2.743880in}{3.011868in}}%
\pgfpathlineto{\pgfqpoint{2.744187in}{3.043500in}}%
\pgfpathlineto{\pgfqpoint{2.744425in}{3.049374in}}%
\pgfpathlineto{\pgfqpoint{2.745131in}{3.016818in}}%
\pgfpathlineto{\pgfqpoint{2.745282in}{3.033949in}}%
\pgfpathlineto{\pgfqpoint{2.745652in}{3.022684in}}%
\pgfpathlineto{\pgfqpoint{2.746144in}{3.061844in}}%
\pgfpathlineto{\pgfqpoint{2.746343in}{3.051382in}}%
\pgfpathlineto{\pgfqpoint{2.747059in}{3.055961in}}%
\pgfpathlineto{\pgfqpoint{2.746786in}{3.023169in}}%
\pgfpathlineto{\pgfqpoint{2.747444in}{3.051834in}}%
\pgfpathlineto{\pgfqpoint{2.748612in}{3.005937in}}%
\pgfpathlineto{\pgfqpoint{2.748646in}{3.007105in}}%
\pgfpathlineto{\pgfqpoint{2.749386in}{3.048477in}}%
\pgfpathlineto{\pgfqpoint{2.748661in}{3.006897in}}%
\pgfpathlineto{\pgfqpoint{2.749824in}{3.021506in}}%
\pgfpathlineto{\pgfqpoint{2.751003in}{2.987564in}}%
\pgfpathlineto{\pgfqpoint{2.749834in}{3.021726in}}%
\pgfpathlineto{\pgfqpoint{2.751032in}{2.993291in}}%
\pgfpathlineto{\pgfqpoint{2.751066in}{2.992169in}}%
\pgfpathlineto{\pgfqpoint{2.751095in}{2.998347in}}%
\pgfpathlineto{\pgfqpoint{2.751105in}{2.998041in}}%
\pgfpathlineto{\pgfqpoint{2.752181in}{3.075498in}}%
\pgfpathlineto{\pgfqpoint{2.751226in}{2.993447in}}%
\pgfpathlineto{\pgfqpoint{2.752341in}{3.068050in}}%
\pgfpathlineto{\pgfqpoint{2.752750in}{3.026887in}}%
\pgfpathlineto{\pgfqpoint{2.753325in}{3.070809in}}%
\pgfpathlineto{\pgfqpoint{2.753447in}{3.066891in}}%
\pgfpathlineto{\pgfqpoint{2.754347in}{3.081772in}}%
\pgfpathlineto{\pgfqpoint{2.753554in}{3.053306in}}%
\pgfpathlineto{\pgfqpoint{2.754561in}{3.071753in}}%
\pgfpathlineto{\pgfqpoint{2.755715in}{3.044004in}}%
\pgfpathlineto{\pgfqpoint{2.755248in}{3.084971in}}%
\pgfpathlineto{\pgfqpoint{2.755730in}{3.047087in}}%
\pgfpathlineto{\pgfqpoint{2.756119in}{3.057858in}}%
\pgfpathlineto{\pgfqpoint{2.756606in}{3.022386in}}%
\pgfpathlineto{\pgfqpoint{2.756743in}{3.030220in}}%
\pgfpathlineto{\pgfqpoint{2.757926in}{3.008418in}}%
\pgfpathlineto{\pgfqpoint{2.757619in}{3.055231in}}%
\pgfpathlineto{\pgfqpoint{2.757940in}{3.009941in}}%
\pgfpathlineto{\pgfqpoint{2.757945in}{3.011410in}}%
\pgfpathlineto{\pgfqpoint{2.758826in}{2.945440in}}%
\pgfpathlineto{\pgfqpoint{2.758851in}{2.949974in}}%
\pgfpathlineto{\pgfqpoint{2.759537in}{2.925916in}}%
\pgfpathlineto{\pgfqpoint{2.759927in}{2.952935in}}%
\pgfpathlineto{\pgfqpoint{2.761095in}{3.001480in}}%
\pgfpathlineto{\pgfqpoint{2.760058in}{2.950166in}}%
\pgfpathlineto{\pgfqpoint{2.761163in}{2.990210in}}%
\pgfpathlineto{\pgfqpoint{2.762147in}{2.932785in}}%
\pgfpathlineto{\pgfqpoint{2.762361in}{2.944484in}}%
\pgfpathlineto{\pgfqpoint{2.763451in}{2.972842in}}%
\pgfpathlineto{\pgfqpoint{2.762502in}{2.932238in}}%
\pgfpathlineto{\pgfqpoint{2.763505in}{2.968826in}}%
\pgfpathlineto{\pgfqpoint{2.764279in}{2.930244in}}%
\pgfpathlineto{\pgfqpoint{2.763554in}{2.973463in}}%
\pgfpathlineto{\pgfqpoint{2.764829in}{2.935886in}}%
\pgfpathlineto{\pgfqpoint{2.765491in}{2.972145in}}%
\pgfpathlineto{\pgfqpoint{2.765959in}{2.953957in}}%
\pgfpathlineto{\pgfqpoint{2.766051in}{2.961915in}}%
\pgfpathlineto{\pgfqpoint{2.766596in}{2.926575in}}%
\pgfpathlineto{\pgfqpoint{2.766913in}{2.933649in}}%
\pgfpathlineto{\pgfqpoint{2.768101in}{2.907168in}}%
\pgfpathlineto{\pgfqpoint{2.767405in}{2.942234in}}%
\pgfpathlineto{\pgfqpoint{2.768106in}{2.907957in}}%
\pgfpathlineto{\pgfqpoint{2.769084in}{2.919716in}}%
\pgfpathlineto{\pgfqpoint{2.768529in}{2.900223in}}%
\pgfpathlineto{\pgfqpoint{2.769235in}{2.912931in}}%
\pgfpathlineto{\pgfqpoint{2.770369in}{2.886307in}}%
\pgfpathlineto{\pgfqpoint{2.770442in}{2.891470in}}%
\pgfpathlineto{\pgfqpoint{2.770457in}{2.895751in}}%
\pgfpathlineto{\pgfqpoint{2.771416in}{2.867130in}}%
\pgfpathlineto{\pgfqpoint{2.771445in}{2.872010in}}%
\pgfpathlineto{\pgfqpoint{2.772434in}{2.825235in}}%
\pgfpathlineto{\pgfqpoint{2.771514in}{2.876207in}}%
\pgfpathlineto{\pgfqpoint{2.772619in}{2.832726in}}%
\pgfpathlineto{\pgfqpoint{2.772959in}{2.856419in}}%
\pgfpathlineto{\pgfqpoint{2.772648in}{2.832042in}}%
\pgfpathlineto{\pgfqpoint{2.772964in}{2.855240in}}%
\pgfpathlineto{\pgfqpoint{2.774128in}{2.915380in}}%
\pgfpathlineto{\pgfqpoint{2.773310in}{2.848678in}}%
\pgfpathlineto{\pgfqpoint{2.774143in}{2.909046in}}%
\pgfpathlineto{\pgfqpoint{2.774508in}{2.865645in}}%
\pgfpathlineto{\pgfqpoint{2.775335in}{2.876662in}}%
\pgfpathlineto{\pgfqpoint{2.775632in}{2.909788in}}%
\pgfpathlineto{\pgfqpoint{2.775379in}{2.876557in}}%
\pgfpathlineto{\pgfqpoint{2.776465in}{2.895985in}}%
\pgfpathlineto{\pgfqpoint{2.776674in}{2.884614in}}%
\pgfpathlineto{\pgfqpoint{2.777424in}{2.917444in}}%
\pgfpathlineto{\pgfqpoint{2.777482in}{2.916624in}}%
\pgfpathlineto{\pgfqpoint{2.778315in}{2.956395in}}%
\pgfpathlineto{\pgfqpoint{2.777502in}{2.915772in}}%
\pgfpathlineto{\pgfqpoint{2.778661in}{2.941980in}}%
\pgfpathlineto{\pgfqpoint{2.778870in}{2.927831in}}%
\pgfpathlineto{\pgfqpoint{2.779245in}{2.961631in}}%
\pgfpathlineto{\pgfqpoint{2.779493in}{2.953108in}}%
\pgfpathlineto{\pgfqpoint{2.780540in}{2.989293in}}%
\pgfpathlineto{\pgfqpoint{2.779736in}{2.938332in}}%
\pgfpathlineto{\pgfqpoint{2.780666in}{2.965097in}}%
\pgfpathlineto{\pgfqpoint{2.780691in}{2.960102in}}%
\pgfpathlineto{\pgfqpoint{2.781562in}{3.026795in}}%
\pgfpathlineto{\pgfqpoint{2.781718in}{3.014567in}}%
\pgfpathlineto{\pgfqpoint{2.781723in}{3.018842in}}%
\pgfpathlineto{\pgfqpoint{2.782760in}{2.987099in}}%
\pgfpathlineto{\pgfqpoint{2.782799in}{2.988006in}}%
\pgfpathlineto{\pgfqpoint{2.782901in}{2.969081in}}%
\pgfpathlineto{\pgfqpoint{2.783982in}{3.013126in}}%
\pgfpathlineto{\pgfqpoint{2.784493in}{3.002576in}}%
\pgfpathlineto{\pgfqpoint{2.784318in}{3.026410in}}%
\pgfpathlineto{\pgfqpoint{2.785077in}{3.012621in}}%
\pgfpathlineto{\pgfqpoint{2.785954in}{3.035850in}}%
\pgfpathlineto{\pgfqpoint{2.785340in}{2.984321in}}%
\pgfpathlineto{\pgfqpoint{2.786197in}{3.021519in}}%
\pgfpathlineto{\pgfqpoint{2.786226in}{3.017913in}}%
\pgfpathlineto{\pgfqpoint{2.786650in}{3.043941in}}%
\pgfpathlineto{\pgfqpoint{2.787278in}{3.028598in}}%
\pgfpathlineto{\pgfqpoint{2.787862in}{3.033521in}}%
\pgfpathlineto{\pgfqpoint{2.788315in}{3.008224in}}%
\pgfpathlineto{\pgfqpoint{2.788339in}{3.014725in}}%
\pgfpathlineto{\pgfqpoint{2.788344in}{3.014720in}}%
\pgfpathlineto{\pgfqpoint{2.788349in}{3.014947in}}%
\pgfpathlineto{\pgfqpoint{2.788407in}{3.026111in}}%
\pgfpathlineto{\pgfqpoint{2.789084in}{2.986005in}}%
\pgfpathlineto{\pgfqpoint{2.789435in}{3.006053in}}%
\pgfpathlineto{\pgfqpoint{2.789478in}{2.998807in}}%
\pgfpathlineto{\pgfqpoint{2.790199in}{3.038431in}}%
\pgfpathlineto{\pgfqpoint{2.790394in}{3.028227in}}%
\pgfpathlineto{\pgfqpoint{2.790880in}{3.037290in}}%
\pgfpathlineto{\pgfqpoint{2.791265in}{3.010235in}}%
\pgfpathlineto{\pgfqpoint{2.791494in}{3.022680in}}%
\pgfpathlineto{\pgfqpoint{2.792010in}{3.040155in}}%
\pgfpathlineto{\pgfqpoint{2.791830in}{3.009430in}}%
\pgfpathlineto{\pgfqpoint{2.792584in}{3.015531in}}%
\pgfpathlineto{\pgfqpoint{2.793607in}{3.001684in}}%
\pgfpathlineto{\pgfqpoint{2.792672in}{3.033571in}}%
\pgfpathlineto{\pgfqpoint{2.793802in}{3.010978in}}%
\pgfpathlineto{\pgfqpoint{2.793811in}{3.015542in}}%
\pgfpathlineto{\pgfqpoint{2.794410in}{2.983986in}}%
\pgfpathlineto{\pgfqpoint{2.794858in}{2.990738in}}%
\pgfpathlineto{\pgfqpoint{2.795686in}{2.965518in}}%
\pgfpathlineto{\pgfqpoint{2.795028in}{3.009710in}}%
\pgfpathlineto{\pgfqpoint{2.795953in}{2.992401in}}%
\pgfpathlineto{\pgfqpoint{2.796698in}{3.006630in}}%
\pgfpathlineto{\pgfqpoint{2.796148in}{2.977331in}}%
\pgfpathlineto{\pgfqpoint{2.797059in}{2.991569in}}%
\pgfpathlineto{\pgfqpoint{2.797063in}{2.990609in}}%
\pgfpathlineto{\pgfqpoint{2.797594in}{3.030559in}}%
\pgfpathlineto{\pgfqpoint{2.797993in}{3.026229in}}%
\pgfpathlineto{\pgfqpoint{2.798821in}{3.058199in}}%
\pgfpathlineto{\pgfqpoint{2.798023in}{3.025617in}}%
\pgfpathlineto{\pgfqpoint{2.799128in}{3.032224in}}%
\pgfpathlineto{\pgfqpoint{2.799790in}{3.047964in}}%
\pgfpathlineto{\pgfqpoint{2.799361in}{3.014831in}}%
\pgfpathlineto{\pgfqpoint{2.800028in}{3.020044in}}%
\pgfpathlineto{\pgfqpoint{2.801046in}{2.983678in}}%
\pgfpathlineto{\pgfqpoint{2.800335in}{3.025326in}}%
\pgfpathlineto{\pgfqpoint{2.801187in}{2.990923in}}%
\pgfpathlineto{\pgfqpoint{2.801513in}{2.970287in}}%
\pgfpathlineto{\pgfqpoint{2.802321in}{3.011270in}}%
\pgfpathlineto{\pgfqpoint{2.803709in}{2.949832in}}%
\pgfpathlineto{\pgfqpoint{2.802463in}{3.023349in}}%
\pgfpathlineto{\pgfqpoint{2.804079in}{2.962718in}}%
\pgfpathlineto{\pgfqpoint{2.804629in}{3.017613in}}%
\pgfpathlineto{\pgfqpoint{2.805257in}{2.975851in}}%
\pgfpathlineto{\pgfqpoint{2.805875in}{2.946641in}}%
\pgfpathlineto{\pgfqpoint{2.806187in}{2.979733in}}%
\pgfpathlineto{\pgfqpoint{2.806387in}{2.967581in}}%
\pgfpathlineto{\pgfqpoint{2.807239in}{2.947728in}}%
\pgfpathlineto{\pgfqpoint{2.807394in}{2.973918in}}%
\pgfpathlineto{\pgfqpoint{2.807482in}{2.968746in}}%
\pgfpathlineto{\pgfqpoint{2.808529in}{2.995122in}}%
\pgfpathlineto{\pgfqpoint{2.808091in}{2.965271in}}%
\pgfpathlineto{\pgfqpoint{2.808646in}{2.978485in}}%
\pgfpathlineto{\pgfqpoint{2.809619in}{2.948179in}}%
\pgfpathlineto{\pgfqpoint{2.808738in}{2.980362in}}%
\pgfpathlineto{\pgfqpoint{2.809799in}{2.952613in}}%
\pgfpathlineto{\pgfqpoint{2.809853in}{2.954493in}}%
\pgfpathlineto{\pgfqpoint{2.810788in}{2.914170in}}%
\pgfpathlineto{\pgfqpoint{2.810846in}{2.910960in}}%
\pgfpathlineto{\pgfqpoint{2.811791in}{2.953107in}}%
\pgfpathlineto{\pgfqpoint{2.811805in}{2.951347in}}%
\pgfpathlineto{\pgfqpoint{2.811839in}{2.963429in}}%
\pgfpathlineto{\pgfqpoint{2.812282in}{2.938891in}}%
\pgfpathlineto{\pgfqpoint{2.812935in}{2.954995in}}%
\pgfpathlineto{\pgfqpoint{2.814006in}{2.947616in}}%
\pgfpathlineto{\pgfqpoint{2.813203in}{2.978525in}}%
\pgfpathlineto{\pgfqpoint{2.814040in}{2.951894in}}%
\pgfpathlineto{\pgfqpoint{2.814848in}{2.991721in}}%
\pgfpathlineto{\pgfqpoint{2.815174in}{2.976821in}}%
\pgfpathlineto{\pgfqpoint{2.815247in}{2.968922in}}%
\pgfpathlineto{\pgfqpoint{2.816182in}{3.026681in}}%
\pgfpathlineto{\pgfqpoint{2.817214in}{3.066932in}}%
\pgfpathlineto{\pgfqpoint{2.817360in}{3.045595in}}%
\pgfpathlineto{\pgfqpoint{2.817794in}{3.016207in}}%
\pgfpathlineto{\pgfqpoint{2.817565in}{3.049088in}}%
\pgfpathlineto{\pgfqpoint{2.818480in}{3.038411in}}%
\pgfpathlineto{\pgfqpoint{2.819132in}{3.027313in}}%
\pgfpathlineto{\pgfqpoint{2.819488in}{3.064795in}}%
\pgfpathlineto{\pgfqpoint{2.819507in}{3.062283in}}%
\pgfpathlineto{\pgfqpoint{2.819785in}{3.058553in}}%
\pgfpathlineto{\pgfqpoint{2.819975in}{3.079513in}}%
\pgfpathlineto{\pgfqpoint{2.820150in}{3.090304in}}%
\pgfpathlineto{\pgfqpoint{2.820943in}{3.063975in}}%
\pgfpathlineto{\pgfqpoint{2.821065in}{3.069397in}}%
\pgfpathlineto{\pgfqpoint{2.822185in}{3.063268in}}%
\pgfpathlineto{\pgfqpoint{2.821172in}{3.083139in}}%
\pgfpathlineto{\pgfqpoint{2.822190in}{3.064970in}}%
\pgfpathlineto{\pgfqpoint{2.822677in}{3.054642in}}%
\pgfpathlineto{\pgfqpoint{2.822867in}{3.075717in}}%
\pgfpathlineto{\pgfqpoint{2.822988in}{3.072301in}}%
\pgfpathlineto{\pgfqpoint{2.823305in}{3.083977in}}%
\pgfpathlineto{\pgfqpoint{2.823957in}{3.044849in}}%
\pgfpathlineto{\pgfqpoint{2.824030in}{3.058701in}}%
\pgfpathlineto{\pgfqpoint{2.824142in}{3.050386in}}%
\pgfpathlineto{\pgfqpoint{2.824906in}{3.108495in}}%
\pgfpathlineto{\pgfqpoint{2.824950in}{3.103952in}}%
\pgfpathlineto{\pgfqpoint{2.824975in}{3.113678in}}%
\pgfpathlineto{\pgfqpoint{2.825379in}{3.080876in}}%
\pgfpathlineto{\pgfqpoint{2.826065in}{3.106475in}}%
\pgfpathlineto{\pgfqpoint{2.826664in}{3.099476in}}%
\pgfpathlineto{\pgfqpoint{2.827092in}{3.140074in}}%
\pgfpathlineto{\pgfqpoint{2.827107in}{3.136554in}}%
\pgfpathlineto{\pgfqpoint{2.827380in}{3.159837in}}%
\pgfpathlineto{\pgfqpoint{2.828100in}{3.118656in}}%
\pgfpathlineto{\pgfqpoint{2.828202in}{3.122821in}}%
\pgfpathlineto{\pgfqpoint{2.828222in}{3.126125in}}%
\pgfpathlineto{\pgfqpoint{2.828757in}{3.102152in}}%
\pgfpathlineto{\pgfqpoint{2.829936in}{3.048114in}}%
\pgfpathlineto{\pgfqpoint{2.828957in}{3.103937in}}%
\pgfpathlineto{\pgfqpoint{2.830101in}{3.076811in}}%
\pgfpathlineto{\pgfqpoint{2.830656in}{3.102639in}}%
\pgfpathlineto{\pgfqpoint{2.830111in}{3.074153in}}%
\pgfpathlineto{\pgfqpoint{2.831377in}{3.086931in}}%
\pgfpathlineto{\pgfqpoint{2.832321in}{3.041361in}}%
\pgfpathlineto{\pgfqpoint{2.831678in}{3.098630in}}%
\pgfpathlineto{\pgfqpoint{2.832696in}{3.075241in}}%
\pgfpathlineto{\pgfqpoint{2.832701in}{3.076802in}}%
\pgfpathlineto{\pgfqpoint{2.833597in}{3.020700in}}%
\pgfpathlineto{\pgfqpoint{2.833645in}{3.025812in}}%
\pgfpathlineto{\pgfqpoint{2.834858in}{2.994281in}}%
\pgfpathlineto{\pgfqpoint{2.834877in}{2.996260in}}%
\pgfpathlineto{\pgfqpoint{2.835447in}{2.983217in}}%
\pgfpathlineto{\pgfqpoint{2.835822in}{3.013034in}}%
\pgfpathlineto{\pgfqpoint{2.835992in}{2.995133in}}%
\pgfpathlineto{\pgfqpoint{2.836518in}{3.020994in}}%
\pgfpathlineto{\pgfqpoint{2.836946in}{2.986277in}}%
\pgfpathlineto{\pgfqpoint{2.838163in}{2.874517in}}%
\pgfpathlineto{\pgfqpoint{2.837058in}{2.989592in}}%
\pgfpathlineto{\pgfqpoint{2.838285in}{2.888820in}}%
\pgfpathlineto{\pgfqpoint{2.838358in}{2.892481in}}%
\pgfpathlineto{\pgfqpoint{2.838835in}{2.860980in}}%
\pgfpathlineto{\pgfqpoint{2.839239in}{2.871138in}}%
\pgfpathlineto{\pgfqpoint{2.839249in}{2.870529in}}%
\pgfpathlineto{\pgfqpoint{2.839317in}{2.880109in}}%
\pgfpathlineto{\pgfqpoint{2.839322in}{2.879936in}}%
\pgfpathlineto{\pgfqpoint{2.840145in}{2.905736in}}%
\pgfpathlineto{\pgfqpoint{2.839385in}{2.868301in}}%
\pgfpathlineto{\pgfqpoint{2.840427in}{2.877770in}}%
\pgfpathlineto{\pgfqpoint{2.840680in}{2.862512in}}%
\pgfpathlineto{\pgfqpoint{2.840515in}{2.890038in}}%
\pgfpathlineto{\pgfqpoint{2.841226in}{2.886333in}}%
\pgfpathlineto{\pgfqpoint{2.841727in}{2.904122in}}%
\pgfpathlineto{\pgfqpoint{2.842048in}{2.874915in}}%
\pgfpathlineto{\pgfqpoint{2.842321in}{2.876383in}}%
\pgfpathlineto{\pgfqpoint{2.843003in}{2.880741in}}%
\pgfpathlineto{\pgfqpoint{2.843251in}{2.855865in}}%
\pgfpathlineto{\pgfqpoint{2.843256in}{2.855990in}}%
\pgfpathlineto{\pgfqpoint{2.844390in}{2.838653in}}%
\pgfpathlineto{\pgfqpoint{2.843684in}{2.871519in}}%
\pgfpathlineto{\pgfqpoint{2.844434in}{2.843264in}}%
\pgfpathlineto{\pgfqpoint{2.846863in}{2.770740in}}%
\pgfpathlineto{\pgfqpoint{2.844746in}{2.864212in}}%
\pgfpathlineto{\pgfqpoint{2.846995in}{2.790179in}}%
\pgfpathlineto{\pgfqpoint{2.847107in}{2.803750in}}%
\pgfpathlineto{\pgfqpoint{2.847891in}{2.774691in}}%
\pgfpathlineto{\pgfqpoint{2.848080in}{2.776719in}}%
\pgfpathlineto{\pgfqpoint{2.848699in}{2.767221in}}%
\pgfpathlineto{\pgfqpoint{2.848436in}{2.795785in}}%
\pgfpathlineto{\pgfqpoint{2.848840in}{2.785319in}}%
\pgfpathlineto{\pgfqpoint{2.849103in}{2.806172in}}%
\pgfpathlineto{\pgfqpoint{2.849833in}{2.752208in}}%
\pgfpathlineto{\pgfqpoint{2.850739in}{2.714248in}}%
\pgfpathlineto{\pgfqpoint{2.850958in}{2.740674in}}%
\pgfpathlineto{\pgfqpoint{2.851980in}{2.778713in}}%
\pgfpathlineto{\pgfqpoint{2.851128in}{2.733077in}}%
\pgfpathlineto{\pgfqpoint{2.852355in}{2.762384in}}%
\pgfpathlineto{\pgfqpoint{2.853562in}{2.740130in}}%
\pgfpathlineto{\pgfqpoint{2.852910in}{2.769281in}}%
\pgfpathlineto{\pgfqpoint{2.853606in}{2.746831in}}%
\pgfpathlineto{\pgfqpoint{2.854745in}{2.791912in}}%
\pgfpathlineto{\pgfqpoint{2.854770in}{2.789121in}}%
\pgfpathlineto{\pgfqpoint{2.854809in}{2.777334in}}%
\pgfpathlineto{\pgfqpoint{2.855734in}{2.810083in}}%
\pgfpathlineto{\pgfqpoint{2.855778in}{2.808065in}}%
\pgfpathlineto{\pgfqpoint{2.856581in}{2.818618in}}%
\pgfpathlineto{\pgfqpoint{2.856849in}{2.785171in}}%
\pgfpathlineto{\pgfqpoint{2.856863in}{2.787904in}}%
\pgfpathlineto{\pgfqpoint{2.857034in}{2.772145in}}%
\pgfpathlineto{\pgfqpoint{2.857671in}{2.814224in}}%
\pgfpathlineto{\pgfqpoint{2.857866in}{2.809129in}}%
\pgfpathlineto{\pgfqpoint{2.858942in}{2.832780in}}%
\pgfpathlineto{\pgfqpoint{2.858348in}{2.796691in}}%
\pgfpathlineto{\pgfqpoint{2.859010in}{2.827065in}}%
\pgfpathlineto{\pgfqpoint{2.859239in}{2.814315in}}%
\pgfpathlineto{\pgfqpoint{2.859497in}{2.842093in}}%
\pgfpathlineto{\pgfqpoint{2.860115in}{2.826573in}}%
\pgfpathlineto{\pgfqpoint{2.860787in}{2.839531in}}%
\pgfpathlineto{\pgfqpoint{2.860261in}{2.808700in}}%
\pgfpathlineto{\pgfqpoint{2.861225in}{2.828427in}}%
\pgfpathlineto{\pgfqpoint{2.862418in}{2.798705in}}%
\pgfpathlineto{\pgfqpoint{2.861513in}{2.830449in}}%
\pgfpathlineto{\pgfqpoint{2.862438in}{2.801291in}}%
\pgfpathlineto{\pgfqpoint{2.863144in}{2.833828in}}%
\pgfpathlineto{\pgfqpoint{2.863577in}{2.815936in}}%
\pgfpathlineto{\pgfqpoint{2.864083in}{2.794051in}}%
\pgfpathlineto{\pgfqpoint{2.863752in}{2.818532in}}%
\pgfpathlineto{\pgfqpoint{2.864702in}{2.806828in}}%
\pgfpathlineto{\pgfqpoint{2.864926in}{2.820742in}}%
\pgfpathlineto{\pgfqpoint{2.865213in}{2.798019in}}%
\pgfpathlineto{\pgfqpoint{2.865802in}{2.805251in}}%
\pgfpathlineto{\pgfqpoint{2.866347in}{2.792378in}}%
\pgfpathlineto{\pgfqpoint{2.866819in}{2.820140in}}%
\pgfpathlineto{\pgfqpoint{2.866853in}{2.815694in}}%
\pgfpathlineto{\pgfqpoint{2.866858in}{2.817029in}}%
\pgfpathlineto{\pgfqpoint{2.867404in}{2.793979in}}%
\pgfpathlineto{\pgfqpoint{2.867929in}{2.807308in}}%
\pgfpathlineto{\pgfqpoint{2.867949in}{2.813794in}}%
\pgfpathlineto{\pgfqpoint{2.868110in}{2.796310in}}%
\pgfpathlineto{\pgfqpoint{2.868251in}{2.784106in}}%
\pgfpathlineto{\pgfqpoint{2.868655in}{2.808161in}}%
\pgfpathlineto{\pgfqpoint{2.868859in}{2.805428in}}%
\pgfpathlineto{\pgfqpoint{2.869010in}{2.813411in}}%
\pgfpathlineto{\pgfqpoint{2.869765in}{2.754176in}}%
\pgfpathlineto{\pgfqpoint{2.869867in}{2.755470in}}%
\pgfpathlineto{\pgfqpoint{2.870802in}{2.762453in}}%
\pgfpathlineto{\pgfqpoint{2.870281in}{2.731827in}}%
\pgfpathlineto{\pgfqpoint{2.870962in}{2.755529in}}%
\pgfpathlineto{\pgfqpoint{2.871985in}{2.741076in}}%
\pgfpathlineto{\pgfqpoint{2.871289in}{2.767753in}}%
\pgfpathlineto{\pgfqpoint{2.872087in}{2.748777in}}%
\pgfpathlineto{\pgfqpoint{2.872637in}{2.762393in}}%
\pgfpathlineto{\pgfqpoint{2.872827in}{2.737592in}}%
\pgfpathlineto{\pgfqpoint{2.873066in}{2.742005in}}%
\pgfpathlineto{\pgfqpoint{2.873742in}{2.681956in}}%
\pgfpathlineto{\pgfqpoint{2.873134in}{2.746196in}}%
\pgfpathlineto{\pgfqpoint{2.874297in}{2.702543in}}%
\pgfpathlineto{\pgfqpoint{2.875208in}{2.714861in}}%
\pgfpathlineto{\pgfqpoint{2.874838in}{2.689014in}}%
\pgfpathlineto{\pgfqpoint{2.875422in}{2.705876in}}%
\pgfpathlineto{\pgfqpoint{2.876683in}{2.640620in}}%
\pgfpathlineto{\pgfqpoint{2.875597in}{2.707929in}}%
\pgfpathlineto{\pgfqpoint{2.876766in}{2.650315in}}%
\pgfpathlineto{\pgfqpoint{2.876941in}{2.666832in}}%
\pgfpathlineto{\pgfqpoint{2.877642in}{2.631092in}}%
\pgfpathlineto{\pgfqpoint{2.877710in}{2.635931in}}%
\pgfpathlineto{\pgfqpoint{2.878830in}{2.613544in}}%
\pgfpathlineto{\pgfqpoint{2.877778in}{2.641059in}}%
\pgfpathlineto{\pgfqpoint{2.878859in}{2.616135in}}%
\pgfpathlineto{\pgfqpoint{2.879059in}{2.622812in}}%
\pgfpathlineto{\pgfqpoint{2.879609in}{2.584052in}}%
\pgfpathlineto{\pgfqpoint{2.879935in}{2.591878in}}%
\pgfpathlineto{\pgfqpoint{2.879945in}{2.587169in}}%
\pgfpathlineto{\pgfqpoint{2.880734in}{2.619478in}}%
\pgfpathlineto{\pgfqpoint{2.880792in}{2.615304in}}%
\pgfpathlineto{\pgfqpoint{2.880816in}{2.622436in}}%
\pgfpathlineto{\pgfqpoint{2.881474in}{2.565874in}}%
\pgfpathlineto{\pgfqpoint{2.881858in}{2.595051in}}%
\pgfpathlineto{\pgfqpoint{2.882257in}{2.586217in}}%
\pgfpathlineto{\pgfqpoint{2.882491in}{2.621452in}}%
\pgfpathlineto{\pgfqpoint{2.882944in}{2.605252in}}%
\pgfpathlineto{\pgfqpoint{2.884112in}{2.630288in}}%
\pgfpathlineto{\pgfqpoint{2.883689in}{2.582571in}}%
\pgfpathlineto{\pgfqpoint{2.884132in}{2.629591in}}%
\pgfpathlineto{\pgfqpoint{2.884351in}{2.617002in}}%
\pgfpathlineto{\pgfqpoint{2.884682in}{2.643178in}}%
\pgfpathlineto{\pgfqpoint{2.884726in}{2.641282in}}%
\pgfpathlineto{\pgfqpoint{2.885763in}{2.682980in}}%
\pgfpathlineto{\pgfqpoint{2.885222in}{2.638156in}}%
\pgfpathlineto{\pgfqpoint{2.885884in}{2.674279in}}%
\pgfpathlineto{\pgfqpoint{2.886108in}{2.689991in}}%
\pgfpathlineto{\pgfqpoint{2.886634in}{2.653172in}}%
\pgfpathlineto{\pgfqpoint{2.886946in}{2.657408in}}%
\pgfpathlineto{\pgfqpoint{2.888241in}{2.725137in}}%
\pgfpathlineto{\pgfqpoint{2.886965in}{2.653555in}}%
\pgfpathlineto{\pgfqpoint{2.888694in}{2.706618in}}%
\pgfpathlineto{\pgfqpoint{2.888767in}{2.696941in}}%
\pgfpathlineto{\pgfqpoint{2.889682in}{2.744779in}}%
\pgfpathlineto{\pgfqpoint{2.889769in}{2.724491in}}%
\pgfpathlineto{\pgfqpoint{2.889877in}{2.743620in}}%
\pgfpathlineto{\pgfqpoint{2.890363in}{2.713881in}}%
\pgfpathlineto{\pgfqpoint{2.890831in}{2.716950in}}%
\pgfpathlineto{\pgfqpoint{2.892140in}{2.631632in}}%
\pgfpathlineto{\pgfqpoint{2.890884in}{2.720004in}}%
\pgfpathlineto{\pgfqpoint{2.892179in}{2.642475in}}%
\pgfpathlineto{\pgfqpoint{2.893017in}{2.628236in}}%
\pgfpathlineto{\pgfqpoint{2.892515in}{2.658465in}}%
\pgfpathlineto{\pgfqpoint{2.893221in}{2.654103in}}%
\pgfpathlineto{\pgfqpoint{2.893255in}{2.660540in}}%
\pgfpathlineto{\pgfqpoint{2.893762in}{2.630762in}}%
\pgfpathlineto{\pgfqpoint{2.894317in}{2.640872in}}%
\pgfpathlineto{\pgfqpoint{2.894881in}{2.608870in}}%
\pgfpathlineto{\pgfqpoint{2.894550in}{2.652316in}}%
\pgfpathlineto{\pgfqpoint{2.895563in}{2.627704in}}%
\pgfpathlineto{\pgfqpoint{2.896722in}{2.656744in}}%
\pgfpathlineto{\pgfqpoint{2.896147in}{2.617454in}}%
\pgfpathlineto{\pgfqpoint{2.896741in}{2.652994in}}%
\pgfpathlineto{\pgfqpoint{2.897559in}{2.621269in}}%
\pgfpathlineto{\pgfqpoint{2.896965in}{2.663911in}}%
\pgfpathlineto{\pgfqpoint{2.897895in}{2.638219in}}%
\pgfpathlineto{\pgfqpoint{2.898012in}{2.631894in}}%
\pgfpathlineto{\pgfqpoint{2.898469in}{2.661672in}}%
\pgfpathlineto{\pgfqpoint{2.898859in}{2.661595in}}%
\pgfpathlineto{\pgfqpoint{2.899896in}{2.704197in}}%
\pgfpathlineto{\pgfqpoint{2.898898in}{2.660192in}}%
\pgfpathlineto{\pgfqpoint{2.900027in}{2.691652in}}%
\pgfpathlineto{\pgfqpoint{2.900612in}{2.737490in}}%
\pgfpathlineto{\pgfqpoint{2.901172in}{2.696309in}}%
\pgfpathlineto{\pgfqpoint{2.902408in}{2.735668in}}%
\pgfpathlineto{\pgfqpoint{2.901274in}{2.691482in}}%
\pgfpathlineto{\pgfqpoint{2.902549in}{2.714339in}}%
\pgfpathlineto{\pgfqpoint{2.903810in}{2.658403in}}%
\pgfpathlineto{\pgfqpoint{2.903927in}{2.669780in}}%
\pgfpathlineto{\pgfqpoint{2.904166in}{2.638494in}}%
\pgfpathlineto{\pgfqpoint{2.904959in}{2.666639in}}%
\pgfpathlineto{\pgfqpoint{2.905777in}{2.657345in}}%
\pgfpathlineto{\pgfqpoint{2.905977in}{2.680389in}}%
\pgfpathlineto{\pgfqpoint{2.906055in}{2.671701in}}%
\pgfpathlineto{\pgfqpoint{2.906391in}{2.662001in}}%
\pgfpathlineto{\pgfqpoint{2.906517in}{2.679602in}}%
\pgfpathlineto{\pgfqpoint{2.907428in}{2.691693in}}%
\pgfpathlineto{\pgfqpoint{2.906775in}{2.661926in}}%
\pgfpathlineto{\pgfqpoint{2.907632in}{2.682751in}}%
\pgfpathlineto{\pgfqpoint{2.907720in}{2.689794in}}%
\pgfpathlineto{\pgfqpoint{2.908674in}{2.651085in}}%
\pgfpathlineto{\pgfqpoint{2.908679in}{2.651867in}}%
\pgfpathlineto{\pgfqpoint{2.909579in}{2.613538in}}%
\pgfpathlineto{\pgfqpoint{2.908791in}{2.658037in}}%
\pgfpathlineto{\pgfqpoint{2.909842in}{2.625414in}}%
\pgfpathlineto{\pgfqpoint{2.910110in}{2.642288in}}%
\pgfpathlineto{\pgfqpoint{2.910738in}{2.598707in}}%
\pgfpathlineto{\pgfqpoint{2.910874in}{2.605129in}}%
\pgfpathlineto{\pgfqpoint{2.911371in}{2.620744in}}%
\pgfpathlineto{\pgfqpoint{2.911658in}{2.599524in}}%
\pgfpathlineto{\pgfqpoint{2.911843in}{2.588612in}}%
\pgfpathlineto{\pgfqpoint{2.912598in}{2.615411in}}%
\pgfpathlineto{\pgfqpoint{2.912729in}{2.611107in}}%
\pgfpathlineto{\pgfqpoint{2.912832in}{2.621564in}}%
\pgfpathlineto{\pgfqpoint{2.913231in}{2.588506in}}%
\pgfpathlineto{\pgfqpoint{2.913805in}{2.606719in}}%
\pgfpathlineto{\pgfqpoint{2.914297in}{2.583698in}}%
\pgfpathlineto{\pgfqpoint{2.914633in}{2.632526in}}%
\pgfpathlineto{\pgfqpoint{2.914925in}{2.602097in}}%
\pgfpathlineto{\pgfqpoint{2.914993in}{2.611674in}}%
\pgfpathlineto{\pgfqpoint{2.915860in}{2.562194in}}%
\pgfpathlineto{\pgfqpoint{2.915938in}{2.566560in}}%
\pgfpathlineto{\pgfqpoint{2.916863in}{2.520528in}}%
\pgfpathlineto{\pgfqpoint{2.916025in}{2.579321in}}%
\pgfpathlineto{\pgfqpoint{2.917096in}{2.543905in}}%
\pgfpathlineto{\pgfqpoint{2.917247in}{2.537908in}}%
\pgfpathlineto{\pgfqpoint{2.917875in}{2.568147in}}%
\pgfpathlineto{\pgfqpoint{2.918206in}{2.544191in}}%
\pgfpathlineto{\pgfqpoint{2.918693in}{2.554773in}}%
\pgfpathlineto{\pgfqpoint{2.919127in}{2.524980in}}%
\pgfpathlineto{\pgfqpoint{2.919219in}{2.533327in}}%
\pgfpathlineto{\pgfqpoint{2.920110in}{2.502464in}}%
\pgfpathlineto{\pgfqpoint{2.919467in}{2.535966in}}%
\pgfpathlineto{\pgfqpoint{2.920387in}{2.505528in}}%
\pgfpathlineto{\pgfqpoint{2.921571in}{2.564411in}}%
\pgfpathlineto{\pgfqpoint{2.920616in}{2.488470in}}%
\pgfpathlineto{\pgfqpoint{2.921575in}{2.563952in}}%
\pgfpathlineto{\pgfqpoint{2.922681in}{2.536305in}}%
\pgfpathlineto{\pgfqpoint{2.921833in}{2.574127in}}%
\pgfpathlineto{\pgfqpoint{2.922724in}{2.546823in}}%
\pgfpathlineto{\pgfqpoint{2.923659in}{2.573002in}}%
\pgfpathlineto{\pgfqpoint{2.923133in}{2.531010in}}%
\pgfpathlineto{\pgfqpoint{2.923859in}{2.555906in}}%
\pgfpathlineto{\pgfqpoint{2.923864in}{2.555863in}}%
\pgfpathlineto{\pgfqpoint{2.923995in}{2.564063in}}%
\pgfpathlineto{\pgfqpoint{2.924000in}{2.563281in}}%
\pgfpathlineto{\pgfqpoint{2.924015in}{2.570134in}}%
\pgfpathlineto{\pgfqpoint{2.925003in}{2.511290in}}%
\pgfpathlineto{\pgfqpoint{2.925008in}{2.511259in}}%
\pgfpathlineto{\pgfqpoint{2.925027in}{2.516930in}}%
\pgfpathlineto{\pgfqpoint{2.926054in}{2.542152in}}%
\pgfpathlineto{\pgfqpoint{2.925460in}{2.497313in}}%
\pgfpathlineto{\pgfqpoint{2.926152in}{2.527556in}}%
\pgfpathlineto{\pgfqpoint{2.926235in}{2.552659in}}%
\pgfpathlineto{\pgfqpoint{2.926814in}{2.521999in}}%
\pgfpathlineto{\pgfqpoint{2.927335in}{2.542937in}}%
\pgfpathlineto{\pgfqpoint{2.928046in}{2.539303in}}%
\pgfpathlineto{\pgfqpoint{2.927710in}{2.569648in}}%
\pgfpathlineto{\pgfqpoint{2.928435in}{2.545742in}}%
\pgfpathlineto{\pgfqpoint{2.929053in}{2.579864in}}%
\pgfpathlineto{\pgfqpoint{2.928459in}{2.542066in}}%
\pgfpathlineto{\pgfqpoint{2.929511in}{2.543395in}}%
\pgfpathlineto{\pgfqpoint{2.930553in}{2.535619in}}%
\pgfpathlineto{\pgfqpoint{2.929920in}{2.571504in}}%
\pgfpathlineto{\pgfqpoint{2.930616in}{2.542680in}}%
\pgfpathlineto{\pgfqpoint{2.931322in}{2.559120in}}%
\pgfpathlineto{\pgfqpoint{2.931152in}{2.531921in}}%
\pgfpathlineto{\pgfqpoint{2.931731in}{2.545256in}}%
\pgfpathlineto{\pgfqpoint{2.931780in}{2.535069in}}%
\pgfpathlineto{\pgfqpoint{2.932710in}{2.574325in}}%
\pgfpathlineto{\pgfqpoint{2.932758in}{2.568080in}}%
\pgfpathlineto{\pgfqpoint{2.934574in}{2.615158in}}%
\pgfpathlineto{\pgfqpoint{2.932797in}{2.560578in}}%
\pgfpathlineto{\pgfqpoint{2.935261in}{2.597417in}}%
\pgfpathlineto{\pgfqpoint{2.935538in}{2.576795in}}%
\pgfpathlineto{\pgfqpoint{2.935718in}{2.606841in}}%
\pgfpathlineto{\pgfqpoint{2.936376in}{2.592960in}}%
\pgfpathlineto{\pgfqpoint{2.937456in}{2.549184in}}%
\pgfpathlineto{\pgfqpoint{2.936439in}{2.595857in}}%
\pgfpathlineto{\pgfqpoint{2.937603in}{2.558059in}}%
\pgfpathlineto{\pgfqpoint{2.937612in}{2.559488in}}%
\pgfpathlineto{\pgfqpoint{2.938469in}{2.527035in}}%
\pgfpathlineto{\pgfqpoint{2.938489in}{2.528518in}}%
\pgfpathlineto{\pgfqpoint{2.938688in}{2.508545in}}%
\pgfpathlineto{\pgfqpoint{2.939428in}{2.535462in}}%
\pgfpathlineto{\pgfqpoint{2.939589in}{2.532123in}}%
\pgfpathlineto{\pgfqpoint{2.940884in}{2.569472in}}%
\pgfpathlineto{\pgfqpoint{2.939657in}{2.524144in}}%
\pgfpathlineto{\pgfqpoint{2.940933in}{2.566292in}}%
\pgfpathlineto{\pgfqpoint{2.942018in}{2.541093in}}%
\pgfpathlineto{\pgfqpoint{2.941259in}{2.567704in}}%
\pgfpathlineto{\pgfqpoint{2.942067in}{2.543817in}}%
\pgfpathlineto{\pgfqpoint{2.942885in}{2.535004in}}%
\pgfpathlineto{\pgfqpoint{2.942301in}{2.560328in}}%
\pgfpathlineto{\pgfqpoint{2.943094in}{2.548954in}}%
\pgfpathlineto{\pgfqpoint{2.943532in}{2.553841in}}%
\pgfpathlineto{\pgfqpoint{2.944005in}{2.497573in}}%
\pgfpathlineto{\pgfqpoint{2.944019in}{2.499138in}}%
\pgfpathlineto{\pgfqpoint{2.944949in}{2.474080in}}%
\pgfpathlineto{\pgfqpoint{2.944185in}{2.506193in}}%
\pgfpathlineto{\pgfqpoint{2.945188in}{2.485584in}}%
\pgfpathlineto{\pgfqpoint{2.945197in}{2.486985in}}%
\pgfpathlineto{\pgfqpoint{2.945859in}{2.444943in}}%
\pgfpathlineto{\pgfqpoint{2.946200in}{2.460393in}}%
\pgfpathlineto{\pgfqpoint{2.946843in}{2.422335in}}%
\pgfpathlineto{\pgfqpoint{2.946244in}{2.461794in}}%
\pgfpathlineto{\pgfqpoint{2.947344in}{2.439496in}}%
\pgfpathlineto{\pgfqpoint{2.947851in}{2.495057in}}%
\pgfpathlineto{\pgfqpoint{2.947374in}{2.438394in}}%
\pgfpathlineto{\pgfqpoint{2.948464in}{2.443980in}}%
\pgfpathlineto{\pgfqpoint{2.948503in}{2.439007in}}%
\pgfpathlineto{\pgfqpoint{2.948922in}{2.471723in}}%
\pgfpathlineto{\pgfqpoint{2.949555in}{2.453975in}}%
\pgfpathlineto{\pgfqpoint{2.949560in}{2.452788in}}%
\pgfpathlineto{\pgfqpoint{2.950285in}{2.485171in}}%
\pgfpathlineto{\pgfqpoint{2.950587in}{2.471606in}}%
\pgfpathlineto{\pgfqpoint{2.951390in}{2.491052in}}%
\pgfpathlineto{\pgfqpoint{2.950850in}{2.451633in}}%
\pgfpathlineto{\pgfqpoint{2.951697in}{2.472657in}}%
\pgfpathlineto{\pgfqpoint{2.952505in}{2.445418in}}%
\pgfpathlineto{\pgfqpoint{2.952841in}{2.451532in}}%
\pgfpathlineto{\pgfqpoint{2.953396in}{2.476100in}}%
\pgfpathlineto{\pgfqpoint{2.953523in}{2.449924in}}%
\pgfpathlineto{\pgfqpoint{2.954024in}{2.466822in}}%
\pgfpathlineto{\pgfqpoint{2.955012in}{2.436300in}}%
\pgfpathlineto{\pgfqpoint{2.955168in}{2.448803in}}%
\pgfpathlineto{\pgfqpoint{2.955363in}{2.463919in}}%
\pgfpathlineto{\pgfqpoint{2.956044in}{2.432886in}}%
\pgfpathlineto{\pgfqpoint{2.956278in}{2.450469in}}%
\pgfpathlineto{\pgfqpoint{2.956770in}{2.433964in}}%
\pgfpathlineto{\pgfqpoint{2.957067in}{2.474358in}}%
\pgfpathlineto{\pgfqpoint{2.957349in}{2.465792in}}%
\pgfpathlineto{\pgfqpoint{2.958415in}{2.478832in}}%
\pgfpathlineto{\pgfqpoint{2.958143in}{2.444260in}}%
\pgfpathlineto{\pgfqpoint{2.958469in}{2.473894in}}%
\pgfpathlineto{\pgfqpoint{2.958693in}{2.461814in}}%
\pgfpathlineto{\pgfqpoint{2.959146in}{2.494078in}}%
\pgfpathlineto{\pgfqpoint{2.959564in}{2.481185in}}%
\pgfpathlineto{\pgfqpoint{2.959584in}{2.487933in}}%
\pgfpathlineto{\pgfqpoint{2.960606in}{2.441564in}}%
\pgfpathlineto{\pgfqpoint{2.960621in}{2.444901in}}%
\pgfpathlineto{\pgfqpoint{2.961721in}{2.433043in}}%
\pgfpathlineto{\pgfqpoint{2.960947in}{2.468873in}}%
\pgfpathlineto{\pgfqpoint{2.961745in}{2.441052in}}%
\pgfpathlineto{\pgfqpoint{2.962193in}{2.491088in}}%
\pgfpathlineto{\pgfqpoint{2.962924in}{2.472607in}}%
\pgfpathlineto{\pgfqpoint{2.963021in}{2.476306in}}%
\pgfpathlineto{\pgfqpoint{2.963347in}{2.438683in}}%
\pgfpathlineto{\pgfqpoint{2.963586in}{2.447755in}}%
\pgfpathlineto{\pgfqpoint{2.963737in}{2.428959in}}%
\pgfpathlineto{\pgfqpoint{2.964593in}{2.461780in}}%
\pgfpathlineto{\pgfqpoint{2.964671in}{2.460545in}}%
\pgfpathlineto{\pgfqpoint{2.965898in}{2.407439in}}%
\pgfpathlineto{\pgfqpoint{2.965105in}{2.466413in}}%
\pgfpathlineto{\pgfqpoint{2.965908in}{2.409012in}}%
\pgfpathlineto{\pgfqpoint{2.965947in}{2.412371in}}%
\pgfpathlineto{\pgfqpoint{2.966789in}{2.385670in}}%
\pgfpathlineto{\pgfqpoint{2.966994in}{2.393933in}}%
\pgfpathlineto{\pgfqpoint{2.966999in}{2.395368in}}%
\pgfpathlineto{\pgfqpoint{2.967899in}{2.352474in}}%
\pgfpathlineto{\pgfqpoint{2.967987in}{2.362762in}}%
\pgfpathlineto{\pgfqpoint{2.968698in}{2.311721in}}%
\pgfpathlineto{\pgfqpoint{2.969175in}{2.320001in}}%
\pgfpathlineto{\pgfqpoint{2.969180in}{2.319642in}}%
\pgfpathlineto{\pgfqpoint{2.969744in}{2.358349in}}%
\pgfpathlineto{\pgfqpoint{2.969861in}{2.361514in}}%
\pgfpathlineto{\pgfqpoint{2.970660in}{2.329246in}}%
\pgfpathlineto{\pgfqpoint{2.970806in}{2.339512in}}%
\pgfpathlineto{\pgfqpoint{2.972028in}{2.379510in}}%
\pgfpathlineto{\pgfqpoint{2.970879in}{2.337640in}}%
\pgfpathlineto{\pgfqpoint{2.972081in}{2.371322in}}%
\pgfpathlineto{\pgfqpoint{2.972188in}{2.362200in}}%
\pgfpathlineto{\pgfqpoint{2.972904in}{2.390695in}}%
\pgfpathlineto{\pgfqpoint{2.972933in}{2.388885in}}%
\pgfpathlineto{\pgfqpoint{2.973396in}{2.415795in}}%
\pgfpathlineto{\pgfqpoint{2.973021in}{2.383458in}}%
\pgfpathlineto{\pgfqpoint{2.974053in}{2.396923in}}%
\pgfpathlineto{\pgfqpoint{2.974102in}{2.386817in}}%
\pgfpathlineto{\pgfqpoint{2.974890in}{2.416608in}}%
\pgfpathlineto{\pgfqpoint{2.975148in}{2.404352in}}%
\pgfpathlineto{\pgfqpoint{2.976385in}{2.362284in}}%
\pgfpathlineto{\pgfqpoint{2.975221in}{2.413261in}}%
\pgfpathlineto{\pgfqpoint{2.976448in}{2.366742in}}%
\pgfpathlineto{\pgfqpoint{2.977685in}{2.425277in}}%
\pgfpathlineto{\pgfqpoint{2.976477in}{2.366730in}}%
\pgfpathlineto{\pgfqpoint{2.977719in}{2.420545in}}%
\pgfpathlineto{\pgfqpoint{2.977938in}{2.424272in}}%
\pgfpathlineto{\pgfqpoint{2.978245in}{2.390253in}}%
\pgfpathlineto{\pgfqpoint{2.978454in}{2.396444in}}%
\pgfpathlineto{\pgfqpoint{2.979822in}{2.343605in}}%
\pgfpathlineto{\pgfqpoint{2.980017in}{2.349449in}}%
\pgfpathlineto{\pgfqpoint{2.980153in}{2.359885in}}%
\pgfpathlineto{\pgfqpoint{2.980995in}{2.314216in}}%
\pgfpathlineto{\pgfqpoint{2.981059in}{2.320089in}}%
\pgfpathlineto{\pgfqpoint{2.981122in}{2.314632in}}%
\pgfpathlineto{\pgfqpoint{2.981770in}{2.357059in}}%
\pgfpathlineto{\pgfqpoint{2.982008in}{2.377232in}}%
\pgfpathlineto{\pgfqpoint{2.982466in}{2.330455in}}%
\pgfpathlineto{\pgfqpoint{2.982841in}{2.345349in}}%
\pgfpathlineto{\pgfqpoint{2.982855in}{2.344413in}}%
\pgfpathlineto{\pgfqpoint{2.982860in}{2.345867in}}%
\pgfpathlineto{\pgfqpoint{2.983391in}{2.360043in}}%
\pgfpathlineto{\pgfqpoint{2.983668in}{2.331135in}}%
\pgfpathlineto{\pgfqpoint{2.983956in}{2.344303in}}%
\pgfpathlineto{\pgfqpoint{2.984175in}{2.308974in}}%
\pgfpathlineto{\pgfqpoint{2.983980in}{2.344329in}}%
\pgfpathlineto{\pgfqpoint{2.985070in}{2.340101in}}%
\pgfpathlineto{\pgfqpoint{2.986224in}{2.361314in}}%
\pgfpathlineto{\pgfqpoint{2.985440in}{2.328429in}}%
\pgfpathlineto{\pgfqpoint{2.986229in}{2.360015in}}%
\pgfpathlineto{\pgfqpoint{2.987135in}{2.396958in}}%
\pgfpathlineto{\pgfqpoint{2.986312in}{2.353468in}}%
\pgfpathlineto{\pgfqpoint{2.987490in}{2.372846in}}%
\pgfpathlineto{\pgfqpoint{2.987709in}{2.354840in}}%
\pgfpathlineto{\pgfqpoint{2.988254in}{2.391333in}}%
\pgfpathlineto{\pgfqpoint{2.988542in}{2.387996in}}%
\pgfpathlineto{\pgfqpoint{2.988790in}{2.402944in}}%
\pgfpathlineto{\pgfqpoint{2.989525in}{2.374356in}}%
\pgfpathlineto{\pgfqpoint{2.989608in}{2.380293in}}%
\pgfpathlineto{\pgfqpoint{2.990324in}{2.349749in}}%
\pgfpathlineto{\pgfqpoint{2.990698in}{2.386782in}}%
\pgfpathlineto{\pgfqpoint{2.990713in}{2.383098in}}%
\pgfpathlineto{\pgfqpoint{2.990893in}{2.399528in}}%
\pgfpathlineto{\pgfqpoint{2.991750in}{2.368210in}}%
\pgfpathlineto{\pgfqpoint{2.991779in}{2.368483in}}%
\pgfpathlineto{\pgfqpoint{2.992203in}{2.344550in}}%
\pgfpathlineto{\pgfqpoint{2.992928in}{2.357653in}}%
\pgfpathlineto{\pgfqpoint{2.993240in}{2.382333in}}%
\pgfpathlineto{\pgfqpoint{2.993965in}{2.341725in}}%
\pgfpathlineto{\pgfqpoint{2.993999in}{2.348559in}}%
\pgfpathlineto{\pgfqpoint{2.995104in}{2.326435in}}%
\pgfpathlineto{\pgfqpoint{2.994097in}{2.360942in}}%
\pgfpathlineto{\pgfqpoint{2.995177in}{2.329971in}}%
\pgfpathlineto{\pgfqpoint{2.996127in}{2.363939in}}%
\pgfpathlineto{\pgfqpoint{2.995216in}{2.327654in}}%
\pgfpathlineto{\pgfqpoint{2.996336in}{2.357072in}}%
\pgfpathlineto{\pgfqpoint{2.997514in}{2.332424in}}%
\pgfpathlineto{\pgfqpoint{2.996443in}{2.370615in}}%
\pgfpathlineto{\pgfqpoint{2.997524in}{2.335579in}}%
\pgfpathlineto{\pgfqpoint{2.998731in}{2.296135in}}%
\pgfpathlineto{\pgfqpoint{2.998912in}{2.304464in}}%
\pgfpathlineto{\pgfqpoint{3.000007in}{2.322858in}}%
\pgfpathlineto{\pgfqpoint{2.999209in}{2.282973in}}%
\pgfpathlineto{\pgfqpoint{3.000056in}{2.316555in}}%
\pgfpathlineto{\pgfqpoint{3.000192in}{2.311737in}}%
\pgfpathlineto{\pgfqpoint{3.000932in}{2.360890in}}%
\pgfpathlineto{\pgfqpoint{3.001093in}{2.355487in}}%
\pgfpathlineto{\pgfqpoint{3.001429in}{2.346785in}}%
\pgfpathlineto{\pgfqpoint{3.001979in}{2.372504in}}%
\pgfpathlineto{\pgfqpoint{3.002135in}{2.370413in}}%
\pgfpathlineto{\pgfqpoint{3.002597in}{2.383382in}}%
\pgfpathlineto{\pgfqpoint{3.003108in}{2.344325in}}%
\pgfpathlineto{\pgfqpoint{3.003994in}{2.318088in}}%
\pgfpathlineto{\pgfqpoint{3.003503in}{2.350670in}}%
\pgfpathlineto{\pgfqpoint{3.004262in}{2.330487in}}%
\pgfpathlineto{\pgfqpoint{3.004944in}{2.371482in}}%
\pgfpathlineto{\pgfqpoint{3.005387in}{2.344262in}}%
\pgfpathlineto{\pgfqpoint{3.006472in}{2.294168in}}%
\pgfpathlineto{\pgfqpoint{3.005611in}{2.357991in}}%
\pgfpathlineto{\pgfqpoint{3.006609in}{2.303635in}}%
\pgfpathlineto{\pgfqpoint{3.007281in}{2.332873in}}%
\pgfpathlineto{\pgfqpoint{3.006638in}{2.301264in}}%
\pgfpathlineto{\pgfqpoint{3.007758in}{2.318464in}}%
\pgfpathlineto{\pgfqpoint{3.008269in}{2.286600in}}%
\pgfpathlineto{\pgfqpoint{3.007899in}{2.330804in}}%
\pgfpathlineto{\pgfqpoint{3.008955in}{2.297699in}}%
\pgfpathlineto{\pgfqpoint{3.010265in}{2.362290in}}%
\pgfpathlineto{\pgfqpoint{3.009043in}{2.292432in}}%
\pgfpathlineto{\pgfqpoint{3.010294in}{2.353828in}}%
\pgfpathlineto{\pgfqpoint{3.010976in}{2.329772in}}%
\pgfpathlineto{\pgfqpoint{3.010810in}{2.362196in}}%
\pgfpathlineto{\pgfqpoint{3.011419in}{2.347327in}}%
\pgfpathlineto{\pgfqpoint{3.011769in}{2.354367in}}%
\pgfpathlineto{\pgfqpoint{3.011706in}{2.337847in}}%
\pgfpathlineto{\pgfqpoint{3.012086in}{2.339776in}}%
\pgfpathlineto{\pgfqpoint{3.012095in}{2.334183in}}%
\pgfpathlineto{\pgfqpoint{3.012436in}{2.373813in}}%
\pgfpathlineto{\pgfqpoint{3.013157in}{2.366610in}}%
\pgfpathlineto{\pgfqpoint{3.013736in}{2.376543in}}%
\pgfpathlineto{\pgfqpoint{3.014009in}{2.336091in}}%
\pgfpathlineto{\pgfqpoint{3.014160in}{2.347945in}}%
\pgfpathlineto{\pgfqpoint{3.015260in}{2.305022in}}%
\pgfpathlineto{\pgfqpoint{3.014286in}{2.352897in}}%
\pgfpathlineto{\pgfqpoint{3.015338in}{2.308938in}}%
\pgfpathlineto{\pgfqpoint{3.015479in}{2.319783in}}%
\pgfpathlineto{\pgfqpoint{3.016351in}{2.290601in}}%
\pgfpathlineto{\pgfqpoint{3.016419in}{2.297675in}}%
\pgfpathlineto{\pgfqpoint{3.016458in}{2.302950in}}%
\pgfpathlineto{\pgfqpoint{3.016599in}{2.287444in}}%
\pgfpathlineto{\pgfqpoint{3.016604in}{2.287578in}}%
\pgfpathlineto{\pgfqpoint{3.017490in}{2.278260in}}%
\pgfpathlineto{\pgfqpoint{3.017003in}{2.303533in}}%
\pgfpathlineto{\pgfqpoint{3.017704in}{2.289415in}}%
\pgfpathlineto{\pgfqpoint{3.018887in}{2.270151in}}%
\pgfpathlineto{\pgfqpoint{3.018371in}{2.302393in}}%
\pgfpathlineto{\pgfqpoint{3.018906in}{2.276035in}}%
\pgfpathlineto{\pgfqpoint{3.019978in}{2.285796in}}%
\pgfpathlineto{\pgfqpoint{3.019496in}{2.251863in}}%
\pgfpathlineto{\pgfqpoint{3.020031in}{2.283903in}}%
\pgfpathlineto{\pgfqpoint{3.020163in}{2.263199in}}%
\pgfpathlineto{\pgfqpoint{3.020805in}{2.294888in}}%
\pgfpathlineto{\pgfqpoint{3.021170in}{2.275169in}}%
\pgfpathlineto{\pgfqpoint{3.021175in}{2.276266in}}%
\pgfpathlineto{\pgfqpoint{3.022125in}{2.240558in}}%
\pgfpathlineto{\pgfqpoint{3.022757in}{2.228960in}}%
\pgfpathlineto{\pgfqpoint{3.022378in}{2.254197in}}%
\pgfpathlineto{\pgfqpoint{3.023215in}{2.252098in}}%
\pgfpathlineto{\pgfqpoint{3.023230in}{2.253922in}}%
\pgfpathlineto{\pgfqpoint{3.023897in}{2.197620in}}%
\pgfpathlineto{\pgfqpoint{3.023902in}{2.198856in}}%
\pgfpathlineto{\pgfqpoint{3.024685in}{2.176352in}}%
\pgfpathlineto{\pgfqpoint{3.024242in}{2.222819in}}%
\pgfpathlineto{\pgfqpoint{3.025075in}{2.186639in}}%
\pgfpathlineto{\pgfqpoint{3.025805in}{2.212134in}}%
\pgfpathlineto{\pgfqpoint{3.026151in}{2.175322in}}%
\pgfpathlineto{\pgfqpoint{3.026170in}{2.178486in}}%
\pgfpathlineto{\pgfqpoint{3.026526in}{2.157493in}}%
\pgfpathlineto{\pgfqpoint{3.026808in}{2.194427in}}%
\pgfpathlineto{\pgfqpoint{3.027042in}{2.193743in}}%
\pgfpathlineto{\pgfqpoint{3.027139in}{2.197320in}}%
\pgfpathlineto{\pgfqpoint{3.028020in}{2.148871in}}%
\pgfpathlineto{\pgfqpoint{3.028079in}{2.159240in}}%
\pgfpathlineto{\pgfqpoint{3.028127in}{2.154775in}}%
\pgfpathlineto{\pgfqpoint{3.028532in}{2.186947in}}%
\pgfpathlineto{\pgfqpoint{3.029121in}{2.176904in}}%
\pgfpathlineto{\pgfqpoint{3.029388in}{2.195546in}}%
\pgfpathlineto{\pgfqpoint{3.030007in}{2.139107in}}%
\pgfpathlineto{\pgfqpoint{3.030177in}{2.144383in}}%
\pgfpathlineto{\pgfqpoint{3.031224in}{2.117232in}}%
\pgfpathlineto{\pgfqpoint{3.030425in}{2.150642in}}%
\pgfpathlineto{\pgfqpoint{3.031326in}{2.125888in}}%
\pgfpathlineto{\pgfqpoint{3.032353in}{2.138717in}}%
\pgfpathlineto{\pgfqpoint{3.031876in}{2.112891in}}%
\pgfpathlineto{\pgfqpoint{3.032436in}{2.126910in}}%
\pgfpathlineto{\pgfqpoint{3.033015in}{2.097777in}}%
\pgfpathlineto{\pgfqpoint{3.032465in}{2.128884in}}%
\pgfpathlineto{\pgfqpoint{3.033600in}{2.124002in}}%
\pgfpathlineto{\pgfqpoint{3.033989in}{2.155783in}}%
\pgfpathlineto{\pgfqpoint{3.034666in}{2.115097in}}%
\pgfpathlineto{\pgfqpoint{3.034724in}{2.129313in}}%
\pgfpathlineto{\pgfqpoint{3.034851in}{2.137256in}}%
\pgfpathlineto{\pgfqpoint{3.035070in}{2.113472in}}%
\pgfpathlineto{\pgfqpoint{3.035537in}{2.115043in}}%
\pgfpathlineto{\pgfqpoint{3.036112in}{2.104912in}}%
\pgfpathlineto{\pgfqpoint{3.036599in}{2.142455in}}%
\pgfpathlineto{\pgfqpoint{3.036725in}{2.150911in}}%
\pgfpathlineto{\pgfqpoint{3.036832in}{2.136161in}}%
\pgfpathlineto{\pgfqpoint{3.036862in}{2.133624in}}%
\pgfpathlineto{\pgfqpoint{3.037733in}{2.171332in}}%
\pgfpathlineto{\pgfqpoint{3.037738in}{2.170442in}}%
\pgfpathlineto{\pgfqpoint{3.038897in}{2.190684in}}%
\pgfpathlineto{\pgfqpoint{3.037806in}{2.166927in}}%
\pgfpathlineto{\pgfqpoint{3.038950in}{2.181676in}}%
\pgfpathlineto{\pgfqpoint{3.039734in}{2.121497in}}%
\pgfpathlineto{\pgfqpoint{3.039106in}{2.181999in}}%
\pgfpathlineto{\pgfqpoint{3.040167in}{2.146890in}}%
\pgfpathlineto{\pgfqpoint{3.040362in}{2.132201in}}%
\pgfpathlineto{\pgfqpoint{3.040946in}{2.181896in}}%
\pgfpathlineto{\pgfqpoint{3.041194in}{2.155941in}}%
\pgfpathlineto{\pgfqpoint{3.042334in}{2.193176in}}%
\pgfpathlineto{\pgfqpoint{3.042343in}{2.192604in}}%
\pgfpathlineto{\pgfqpoint{3.043419in}{2.148798in}}%
\pgfpathlineto{\pgfqpoint{3.042504in}{2.198638in}}%
\pgfpathlineto{\pgfqpoint{3.043551in}{2.160191in}}%
\pgfpathlineto{\pgfqpoint{3.043736in}{2.192097in}}%
\pgfpathlineto{\pgfqpoint{3.044247in}{2.149363in}}%
\pgfpathlineto{\pgfqpoint{3.044680in}{2.169157in}}%
\pgfpathlineto{\pgfqpoint{3.045435in}{2.145142in}}%
\pgfpathlineto{\pgfqpoint{3.044724in}{2.172118in}}%
\pgfpathlineto{\pgfqpoint{3.045781in}{2.169970in}}%
\pgfpathlineto{\pgfqpoint{3.046107in}{2.182997in}}%
\pgfpathlineto{\pgfqpoint{3.046613in}{2.155231in}}%
\pgfpathlineto{\pgfqpoint{3.046866in}{2.164516in}}%
\pgfpathlineto{\pgfqpoint{3.047046in}{2.150314in}}%
\pgfpathlineto{\pgfqpoint{3.047587in}{2.175925in}}%
\pgfpathlineto{\pgfqpoint{3.048025in}{2.159238in}}%
\pgfpathlineto{\pgfqpoint{3.048044in}{2.163717in}}%
\pgfpathlineto{\pgfqpoint{3.049096in}{2.134559in}}%
\pgfpathlineto{\pgfqpoint{3.049111in}{2.142229in}}%
\pgfpathlineto{\pgfqpoint{3.050060in}{2.124628in}}%
\pgfpathlineto{\pgfqpoint{3.050201in}{2.145915in}}%
\pgfpathlineto{\pgfqpoint{3.050230in}{2.137154in}}%
\pgfpathlineto{\pgfqpoint{3.050308in}{2.154294in}}%
\pgfpathlineto{\pgfqpoint{3.050552in}{2.126346in}}%
\pgfpathlineto{\pgfqpoint{3.051389in}{2.147127in}}%
\pgfpathlineto{\pgfqpoint{3.051978in}{2.157557in}}%
\pgfpathlineto{\pgfqpoint{3.051662in}{2.136030in}}%
\pgfpathlineto{\pgfqpoint{3.052227in}{2.141402in}}%
\pgfpathlineto{\pgfqpoint{3.052621in}{2.120896in}}%
\pgfpathlineto{\pgfqpoint{3.052421in}{2.154520in}}%
\pgfpathlineto{\pgfqpoint{3.053375in}{2.133193in}}%
\pgfpathlineto{\pgfqpoint{3.054617in}{2.187484in}}%
\pgfpathlineto{\pgfqpoint{3.053517in}{2.126820in}}%
\pgfpathlineto{\pgfqpoint{3.054646in}{2.186453in}}%
\pgfpathlineto{\pgfqpoint{3.054982in}{2.166893in}}%
\pgfpathlineto{\pgfqpoint{3.055800in}{2.176382in}}%
\pgfpathlineto{\pgfqpoint{3.056545in}{2.170214in}}%
\pgfpathlineto{\pgfqpoint{3.056243in}{2.203802in}}%
\pgfpathlineto{\pgfqpoint{3.056818in}{2.190013in}}%
\pgfpathlineto{\pgfqpoint{3.057558in}{2.212966in}}%
\pgfpathlineto{\pgfqpoint{3.057261in}{2.184625in}}%
\pgfpathlineto{\pgfqpoint{3.057962in}{2.204271in}}%
\pgfpathlineto{\pgfqpoint{3.058823in}{2.175001in}}%
\pgfpathlineto{\pgfqpoint{3.058088in}{2.214405in}}%
\pgfpathlineto{\pgfqpoint{3.059057in}{2.210771in}}%
\pgfpathlineto{\pgfqpoint{3.060508in}{2.249477in}}%
\pgfpathlineto{\pgfqpoint{3.059096in}{2.210180in}}%
\pgfpathlineto{\pgfqpoint{3.060581in}{2.245345in}}%
\pgfpathlineto{\pgfqpoint{3.061039in}{2.228079in}}%
\pgfpathlineto{\pgfqpoint{3.061623in}{2.266325in}}%
\pgfpathlineto{\pgfqpoint{3.061667in}{2.258315in}}%
\pgfpathlineto{\pgfqpoint{3.062781in}{2.272808in}}%
\pgfpathlineto{\pgfqpoint{3.062226in}{2.241527in}}%
\pgfpathlineto{\pgfqpoint{3.062791in}{2.266568in}}%
\pgfpathlineto{\pgfqpoint{3.063409in}{2.252931in}}%
\pgfpathlineto{\pgfqpoint{3.063765in}{2.278850in}}%
\pgfpathlineto{\pgfqpoint{3.063857in}{2.276861in}}%
\pgfpathlineto{\pgfqpoint{3.065006in}{2.286370in}}%
\pgfpathlineto{\pgfqpoint{3.064359in}{2.242903in}}%
\pgfpathlineto{\pgfqpoint{3.065031in}{2.285609in}}%
\pgfpathlineto{\pgfqpoint{3.065493in}{2.257357in}}%
\pgfpathlineto{\pgfqpoint{3.065970in}{2.292811in}}%
\pgfpathlineto{\pgfqpoint{3.066160in}{2.273847in}}%
\pgfpathlineto{\pgfqpoint{3.066545in}{2.256180in}}%
\pgfpathlineto{\pgfqpoint{3.066915in}{2.283512in}}%
\pgfpathlineto{\pgfqpoint{3.067051in}{2.280112in}}%
\pgfpathlineto{\pgfqpoint{3.067966in}{2.297646in}}%
\pgfpathlineto{\pgfqpoint{3.067329in}{2.268633in}}%
\pgfpathlineto{\pgfqpoint{3.068171in}{2.286489in}}%
\pgfpathlineto{\pgfqpoint{3.069485in}{2.227965in}}%
\pgfpathlineto{\pgfqpoint{3.068195in}{2.287146in}}%
\pgfpathlineto{\pgfqpoint{3.069505in}{2.230004in}}%
\pgfpathlineto{\pgfqpoint{3.069524in}{2.234656in}}%
\pgfpathlineto{\pgfqpoint{3.070513in}{2.186290in}}%
\pgfpathlineto{\pgfqpoint{3.070542in}{2.190916in}}%
\pgfpathlineto{\pgfqpoint{3.071155in}{2.173921in}}%
\pgfpathlineto{\pgfqpoint{3.070756in}{2.200202in}}%
\pgfpathlineto{\pgfqpoint{3.071666in}{2.188636in}}%
\pgfpathlineto{\pgfqpoint{3.071783in}{2.192096in}}%
\pgfpathlineto{\pgfqpoint{3.072436in}{2.143258in}}%
\pgfpathlineto{\pgfqpoint{3.072708in}{2.170479in}}%
\pgfpathlineto{\pgfqpoint{3.072903in}{2.162503in}}%
\pgfpathlineto{\pgfqpoint{3.073765in}{2.190263in}}%
\pgfpathlineto{\pgfqpoint{3.073775in}{2.187220in}}%
\pgfpathlineto{\pgfqpoint{3.074782in}{2.234508in}}%
\pgfpathlineto{\pgfqpoint{3.073940in}{2.186628in}}%
\pgfpathlineto{\pgfqpoint{3.074933in}{2.223599in}}%
\pgfpathlineto{\pgfqpoint{3.075629in}{2.188652in}}%
\pgfpathlineto{\pgfqpoint{3.075157in}{2.225880in}}%
\pgfpathlineto{\pgfqpoint{3.076092in}{2.208359in}}%
\pgfpathlineto{\pgfqpoint{3.076842in}{2.241651in}}%
\pgfpathlineto{\pgfqpoint{3.076257in}{2.201778in}}%
\pgfpathlineto{\pgfqpoint{3.077226in}{2.215102in}}%
\pgfpathlineto{\pgfqpoint{3.077231in}{2.214930in}}%
\pgfpathlineto{\pgfqpoint{3.077460in}{2.237755in}}%
\pgfpathlineto{\pgfqpoint{3.077733in}{2.233688in}}%
\pgfpathlineto{\pgfqpoint{3.078497in}{2.262898in}}%
\pgfpathlineto{\pgfqpoint{3.078828in}{2.230801in}}%
\pgfpathlineto{\pgfqpoint{3.078847in}{2.236427in}}%
\pgfpathlineto{\pgfqpoint{3.079047in}{2.246443in}}%
\pgfpathlineto{\pgfqpoint{3.079393in}{2.200757in}}%
\pgfpathlineto{\pgfqpoint{3.079953in}{2.234980in}}%
\pgfpathlineto{\pgfqpoint{3.079987in}{2.240631in}}%
\pgfpathlineto{\pgfqpoint{3.080590in}{2.205463in}}%
\pgfpathlineto{\pgfqpoint{3.080605in}{2.206624in}}%
\pgfpathlineto{\pgfqpoint{3.080620in}{2.204050in}}%
\pgfpathlineto{\pgfqpoint{3.081360in}{2.254812in}}%
\pgfpathlineto{\pgfqpoint{3.082353in}{2.291618in}}%
\pgfpathlineto{\pgfqpoint{3.081408in}{2.249771in}}%
\pgfpathlineto{\pgfqpoint{3.082499in}{2.265566in}}%
\pgfpathlineto{\pgfqpoint{3.082606in}{2.272295in}}%
\pgfpathlineto{\pgfqpoint{3.083030in}{2.245150in}}%
\pgfpathlineto{\pgfqpoint{3.083438in}{2.249073in}}%
\pgfpathlineto{\pgfqpoint{3.083633in}{2.216216in}}%
\pgfpathlineto{\pgfqpoint{3.084475in}{2.251639in}}%
\pgfpathlineto{\pgfqpoint{3.084553in}{2.242388in}}%
\pgfpathlineto{\pgfqpoint{3.085668in}{2.287211in}}%
\pgfpathlineto{\pgfqpoint{3.085030in}{2.238039in}}%
\pgfpathlineto{\pgfqpoint{3.086296in}{2.272381in}}%
\pgfpathlineto{\pgfqpoint{3.086905in}{2.262966in}}%
\pgfpathlineto{\pgfqpoint{3.087197in}{2.283977in}}%
\pgfpathlineto{\pgfqpoint{3.087401in}{2.271561in}}%
\pgfpathlineto{\pgfqpoint{3.088526in}{2.296815in}}%
\pgfpathlineto{\pgfqpoint{3.087436in}{2.268042in}}%
\pgfpathlineto{\pgfqpoint{3.088570in}{2.295249in}}%
\pgfpathlineto{\pgfqpoint{3.089641in}{2.281412in}}%
\pgfpathlineto{\pgfqpoint{3.089266in}{2.316413in}}%
\pgfpathlineto{\pgfqpoint{3.089690in}{2.286345in}}%
\pgfpathlineto{\pgfqpoint{3.090668in}{2.313258in}}%
\pgfpathlineto{\pgfqpoint{3.089704in}{2.284952in}}%
\pgfpathlineto{\pgfqpoint{3.090848in}{2.295243in}}%
\pgfpathlineto{\pgfqpoint{3.091472in}{2.300012in}}%
\pgfpathlineto{\pgfqpoint{3.091043in}{2.279325in}}%
\pgfpathlineto{\pgfqpoint{3.091637in}{2.282646in}}%
\pgfpathlineto{\pgfqpoint{3.091842in}{2.288589in}}%
\pgfpathlineto{\pgfqpoint{3.092898in}{2.214797in}}%
\pgfpathlineto{\pgfqpoint{3.093692in}{2.233674in}}%
\pgfpathlineto{\pgfqpoint{3.093005in}{2.199188in}}%
\pgfpathlineto{\pgfqpoint{3.094027in}{2.219776in}}%
\pgfpathlineto{\pgfqpoint{3.094855in}{2.189796in}}%
\pgfpathlineto{\pgfqpoint{3.094232in}{2.229821in}}%
\pgfpathlineto{\pgfqpoint{3.095157in}{2.201165in}}%
\pgfpathlineto{\pgfqpoint{3.095624in}{2.211580in}}%
\pgfpathlineto{\pgfqpoint{3.096136in}{2.173988in}}%
\pgfpathlineto{\pgfqpoint{3.096140in}{2.175112in}}%
\pgfpathlineto{\pgfqpoint{3.096759in}{2.167406in}}%
\pgfpathlineto{\pgfqpoint{3.097041in}{2.189974in}}%
\pgfpathlineto{\pgfqpoint{3.097138in}{2.189418in}}%
\pgfpathlineto{\pgfqpoint{3.098010in}{2.198988in}}%
\pgfpathlineto{\pgfqpoint{3.097236in}{2.172941in}}%
\pgfpathlineto{\pgfqpoint{3.098219in}{2.182495in}}%
\pgfpathlineto{\pgfqpoint{3.098253in}{2.177154in}}%
\pgfpathlineto{\pgfqpoint{3.098662in}{2.203465in}}%
\pgfpathlineto{\pgfqpoint{3.099285in}{2.198836in}}%
\pgfpathlineto{\pgfqpoint{3.099524in}{2.229587in}}%
\pgfpathlineto{\pgfqpoint{3.099310in}{2.197131in}}%
\pgfpathlineto{\pgfqpoint{3.100439in}{2.211101in}}%
\pgfpathlineto{\pgfqpoint{3.101247in}{2.201030in}}%
\pgfpathlineto{\pgfqpoint{3.101432in}{2.227270in}}%
\pgfpathlineto{\pgfqpoint{3.101535in}{2.217146in}}%
\pgfpathlineto{\pgfqpoint{3.101914in}{2.240693in}}%
\pgfpathlineto{\pgfqpoint{3.102280in}{2.212390in}}%
\pgfpathlineto{\pgfqpoint{3.102689in}{2.238730in}}%
\pgfpathlineto{\pgfqpoint{3.103151in}{2.217393in}}%
\pgfpathlineto{\pgfqpoint{3.103813in}{2.233281in}}%
\pgfpathlineto{\pgfqpoint{3.103818in}{2.233182in}}%
\pgfpathlineto{\pgfqpoint{3.104149in}{2.247518in}}%
\pgfpathlineto{\pgfqpoint{3.104169in}{2.252423in}}%
\pgfpathlineto{\pgfqpoint{3.105147in}{2.208686in}}%
\pgfpathlineto{\pgfqpoint{3.105152in}{2.211109in}}%
\pgfpathlineto{\pgfqpoint{3.105663in}{2.220264in}}%
\pgfpathlineto{\pgfqpoint{3.105882in}{2.198066in}}%
\pgfpathlineto{\pgfqpoint{3.106053in}{2.203435in}}%
\pgfpathlineto{\pgfqpoint{3.106291in}{2.178509in}}%
\pgfpathlineto{\pgfqpoint{3.107041in}{2.226777in}}%
\pgfpathlineto{\pgfqpoint{3.107255in}{2.234039in}}%
\pgfpathlineto{\pgfqpoint{3.107713in}{2.208434in}}%
\pgfpathlineto{\pgfqpoint{3.108034in}{2.219268in}}%
\pgfpathlineto{\pgfqpoint{3.108575in}{2.215480in}}%
\pgfpathlineto{\pgfqpoint{3.109042in}{2.237979in}}%
\pgfpathlineto{\pgfqpoint{3.109130in}{2.224519in}}%
\pgfpathlineto{\pgfqpoint{3.109909in}{2.251154in}}%
\pgfpathlineto{\pgfqpoint{3.109504in}{2.219366in}}%
\pgfpathlineto{\pgfqpoint{3.110240in}{2.226321in}}%
\pgfpathlineto{\pgfqpoint{3.110585in}{2.207249in}}%
\pgfpathlineto{\pgfqpoint{3.111340in}{2.230350in}}%
\pgfpathlineto{\pgfqpoint{3.112021in}{2.242724in}}%
\pgfpathlineto{\pgfqpoint{3.111798in}{2.220698in}}%
\pgfpathlineto{\pgfqpoint{3.112377in}{2.227074in}}%
\pgfpathlineto{\pgfqpoint{3.112640in}{2.212489in}}%
\pgfpathlineto{\pgfqpoint{3.113156in}{2.252124in}}%
\pgfpathlineto{\pgfqpoint{3.113497in}{2.223290in}}%
\pgfpathlineto{\pgfqpoint{3.114597in}{2.180773in}}%
\pgfpathlineto{\pgfqpoint{3.113506in}{2.223823in}}%
\pgfpathlineto{\pgfqpoint{3.115084in}{2.191065in}}%
\pgfpathlineto{\pgfqpoint{3.115089in}{2.192524in}}%
\pgfpathlineto{\pgfqpoint{3.115605in}{2.146328in}}%
\pgfpathlineto{\pgfqpoint{3.116082in}{2.158977in}}%
\pgfpathlineto{\pgfqpoint{3.116174in}{2.155721in}}%
\pgfpathlineto{\pgfqpoint{3.116749in}{2.184022in}}%
\pgfpathlineto{\pgfqpoint{3.117133in}{2.179803in}}%
\pgfpathlineto{\pgfqpoint{3.118331in}{2.197117in}}%
\pgfpathlineto{\pgfqpoint{3.117343in}{2.156214in}}%
\pgfpathlineto{\pgfqpoint{3.118428in}{2.186511in}}%
\pgfpathlineto{\pgfqpoint{3.119373in}{2.159495in}}%
\pgfpathlineto{\pgfqpoint{3.118570in}{2.196646in}}%
\pgfpathlineto{\pgfqpoint{3.119577in}{2.163071in}}%
\pgfpathlineto{\pgfqpoint{3.120741in}{2.185519in}}%
\pgfpathlineto{\pgfqpoint{3.120317in}{2.145525in}}%
\pgfpathlineto{\pgfqpoint{3.120746in}{2.185251in}}%
\pgfpathlineto{\pgfqpoint{3.121924in}{2.169744in}}%
\pgfpathlineto{\pgfqpoint{3.121515in}{2.199493in}}%
\pgfpathlineto{\pgfqpoint{3.121934in}{2.172166in}}%
\pgfpathlineto{\pgfqpoint{3.122055in}{2.187400in}}%
\pgfpathlineto{\pgfqpoint{3.122401in}{2.200009in}}%
\pgfpathlineto{\pgfqpoint{3.122698in}{2.174551in}}%
\pgfpathlineto{\pgfqpoint{3.123136in}{2.179338in}}%
\pgfpathlineto{\pgfqpoint{3.123725in}{2.206943in}}%
\pgfpathlineto{\pgfqpoint{3.123268in}{2.161974in}}%
\pgfpathlineto{\pgfqpoint{3.124344in}{2.193984in}}%
\pgfpathlineto{\pgfqpoint{3.124855in}{2.168617in}}%
\pgfpathlineto{\pgfqpoint{3.125473in}{2.184358in}}%
\pgfpathlineto{\pgfqpoint{3.126671in}{2.238479in}}%
\pgfpathlineto{\pgfqpoint{3.125585in}{2.178267in}}%
\pgfpathlineto{\pgfqpoint{3.127055in}{2.213402in}}%
\pgfpathlineto{\pgfqpoint{3.127211in}{2.210881in}}%
\pgfpathlineto{\pgfqpoint{3.127679in}{2.231634in}}%
\pgfpathlineto{\pgfqpoint{3.128126in}{2.228151in}}%
\pgfpathlineto{\pgfqpoint{3.129300in}{2.247976in}}%
\pgfpathlineto{\pgfqpoint{3.128433in}{2.225873in}}%
\pgfpathlineto{\pgfqpoint{3.129310in}{2.243514in}}%
\pgfpathlineto{\pgfqpoint{3.130551in}{2.223111in}}%
\pgfpathlineto{\pgfqpoint{3.130011in}{2.257795in}}%
\pgfpathlineto{\pgfqpoint{3.130590in}{2.226848in}}%
\pgfpathlineto{\pgfqpoint{3.130600in}{2.229460in}}%
\pgfpathlineto{\pgfqpoint{3.131101in}{2.195344in}}%
\pgfpathlineto{\pgfqpoint{3.131612in}{2.199911in}}%
\pgfpathlineto{\pgfqpoint{3.132411in}{2.155659in}}%
\pgfpathlineto{\pgfqpoint{3.131749in}{2.201418in}}%
\pgfpathlineto{\pgfqpoint{3.132752in}{2.175074in}}%
\pgfpathlineto{\pgfqpoint{3.132941in}{2.158287in}}%
\pgfpathlineto{\pgfqpoint{3.133482in}{2.192787in}}%
\pgfpathlineto{\pgfqpoint{3.133531in}{2.188973in}}%
\pgfpathlineto{\pgfqpoint{3.134456in}{2.242761in}}%
\pgfpathlineto{\pgfqpoint{3.133638in}{2.184015in}}%
\pgfpathlineto{\pgfqpoint{3.134689in}{2.233372in}}%
\pgfpathlineto{\pgfqpoint{3.135566in}{2.207694in}}%
\pgfpathlineto{\pgfqpoint{3.135035in}{2.254660in}}%
\pgfpathlineto{\pgfqpoint{3.135799in}{2.233149in}}%
\pgfpathlineto{\pgfqpoint{3.137114in}{2.296929in}}%
\pgfpathlineto{\pgfqpoint{3.136203in}{2.232854in}}%
\pgfpathlineto{\pgfqpoint{3.137119in}{2.294676in}}%
\pgfpathlineto{\pgfqpoint{3.137489in}{2.282742in}}%
\pgfpathlineto{\pgfqpoint{3.137766in}{2.306224in}}%
\pgfpathlineto{\pgfqpoint{3.138195in}{2.297039in}}%
\pgfpathlineto{\pgfqpoint{3.139027in}{2.356714in}}%
\pgfpathlineto{\pgfqpoint{3.138389in}{2.294054in}}%
\pgfpathlineto{\pgfqpoint{3.139368in}{2.345285in}}%
\pgfpathlineto{\pgfqpoint{3.140283in}{2.349431in}}%
\pgfpathlineto{\pgfqpoint{3.139806in}{2.295670in}}%
\pgfpathlineto{\pgfqpoint{3.140429in}{2.340301in}}%
\pgfpathlineto{\pgfqpoint{3.141325in}{2.302658in}}%
\pgfpathlineto{\pgfqpoint{3.140760in}{2.349020in}}%
\pgfpathlineto{\pgfqpoint{3.141564in}{2.317393in}}%
\pgfpathlineto{\pgfqpoint{3.141632in}{2.330128in}}%
\pgfpathlineto{\pgfqpoint{3.142026in}{2.293477in}}%
\pgfpathlineto{\pgfqpoint{3.142644in}{2.297280in}}%
\pgfpathlineto{\pgfqpoint{3.142659in}{2.300209in}}%
\pgfpathlineto{\pgfqpoint{3.142951in}{2.277194in}}%
\pgfpathlineto{\pgfqpoint{3.143618in}{2.285679in}}%
\pgfpathlineto{\pgfqpoint{3.143623in}{2.282560in}}%
\pgfpathlineto{\pgfqpoint{3.144081in}{2.318449in}}%
\pgfpathlineto{\pgfqpoint{3.144684in}{2.313367in}}%
\pgfpathlineto{\pgfqpoint{3.144777in}{2.292682in}}%
\pgfpathlineto{\pgfqpoint{3.145458in}{2.329246in}}%
\pgfpathlineto{\pgfqpoint{3.145463in}{2.326068in}}%
\pgfpathlineto{\pgfqpoint{3.146525in}{2.345547in}}%
\pgfpathlineto{\pgfqpoint{3.145843in}{2.315870in}}%
\pgfpathlineto{\pgfqpoint{3.146593in}{2.342928in}}%
\pgfpathlineto{\pgfqpoint{3.147265in}{2.350568in}}%
\pgfpathlineto{\pgfqpoint{3.147566in}{2.310368in}}%
\pgfpathlineto{\pgfqpoint{3.147581in}{2.313022in}}%
\pgfpathlineto{\pgfqpoint{3.147703in}{2.307785in}}%
\pgfpathlineto{\pgfqpoint{3.148112in}{2.337664in}}%
\pgfpathlineto{\pgfqpoint{3.148657in}{2.328668in}}%
\pgfpathlineto{\pgfqpoint{3.148798in}{2.342777in}}%
\pgfpathlineto{\pgfqpoint{3.149139in}{2.317429in}}%
\pgfpathlineto{\pgfqpoint{3.150638in}{2.252208in}}%
\pgfpathlineto{\pgfqpoint{3.150702in}{2.253698in}}%
\pgfpathlineto{\pgfqpoint{3.151515in}{2.271318in}}%
\pgfpathlineto{\pgfqpoint{3.151374in}{2.253618in}}%
\pgfpathlineto{\pgfqpoint{3.151826in}{2.260705in}}%
\pgfpathlineto{\pgfqpoint{3.152990in}{2.303722in}}%
\pgfpathlineto{\pgfqpoint{3.151963in}{2.260298in}}%
\pgfpathlineto{\pgfqpoint{3.153370in}{2.279368in}}%
\pgfpathlineto{\pgfqpoint{3.154012in}{2.296193in}}%
\pgfpathlineto{\pgfqpoint{3.154489in}{2.265835in}}%
\pgfpathlineto{\pgfqpoint{3.155483in}{2.268247in}}%
\pgfpathlineto{\pgfqpoint{3.155181in}{2.227086in}}%
\pgfpathlineto{\pgfqpoint{3.155556in}{2.260203in}}%
\pgfpathlineto{\pgfqpoint{3.156130in}{2.257595in}}%
\pgfpathlineto{\pgfqpoint{3.155862in}{2.278069in}}%
\pgfpathlineto{\pgfqpoint{3.156568in}{2.268246in}}%
\pgfpathlineto{\pgfqpoint{3.156773in}{2.276271in}}%
\pgfpathlineto{\pgfqpoint{3.157240in}{2.248217in}}%
\pgfpathlineto{\pgfqpoint{3.157659in}{2.257278in}}%
\pgfpathlineto{\pgfqpoint{3.157922in}{2.270344in}}%
\pgfpathlineto{\pgfqpoint{3.158667in}{2.237502in}}%
\pgfpathlineto{\pgfqpoint{3.158676in}{2.238875in}}%
\pgfpathlineto{\pgfqpoint{3.159801in}{2.211908in}}%
\pgfpathlineto{\pgfqpoint{3.159105in}{2.257599in}}%
\pgfpathlineto{\pgfqpoint{3.159820in}{2.212497in}}%
\pgfpathlineto{\pgfqpoint{3.160765in}{2.232881in}}%
\pgfpathlineto{\pgfqpoint{3.160190in}{2.189766in}}%
\pgfpathlineto{\pgfqpoint{3.160955in}{2.223941in}}%
\pgfpathlineto{\pgfqpoint{3.160960in}{2.223678in}}%
\pgfpathlineto{\pgfqpoint{3.161266in}{2.256869in}}%
\pgfpathlineto{\pgfqpoint{3.161520in}{2.251427in}}%
\pgfpathlineto{\pgfqpoint{3.162688in}{2.310708in}}%
\pgfpathlineto{\pgfqpoint{3.163141in}{2.279345in}}%
\pgfpathlineto{\pgfqpoint{3.163589in}{2.326935in}}%
\pgfpathlineto{\pgfqpoint{3.163779in}{2.313902in}}%
\pgfpathlineto{\pgfqpoint{3.164811in}{2.350259in}}%
\pgfpathlineto{\pgfqpoint{3.164236in}{2.304151in}}%
\pgfpathlineto{\pgfqpoint{3.164947in}{2.337489in}}%
\pgfpathlineto{\pgfqpoint{3.165526in}{2.308400in}}%
\pgfpathlineto{\pgfqpoint{3.165093in}{2.346524in}}%
\pgfpathlineto{\pgfqpoint{3.166096in}{2.320341in}}%
\pgfpathlineto{\pgfqpoint{3.166344in}{2.332447in}}%
\pgfpathlineto{\pgfqpoint{3.166802in}{2.296237in}}%
\pgfpathlineto{\pgfqpoint{3.167104in}{2.307598in}}%
\pgfpathlineto{\pgfqpoint{3.167148in}{2.303988in}}%
\pgfpathlineto{\pgfqpoint{3.168073in}{2.328230in}}%
\pgfpathlineto{\pgfqpoint{3.168180in}{2.320076in}}%
\pgfpathlineto{\pgfqpoint{3.168189in}{2.321916in}}%
\pgfpathlineto{\pgfqpoint{3.168199in}{2.320579in}}%
\pgfpathlineto{\pgfqpoint{3.169105in}{2.370367in}}%
\pgfpathlineto{\pgfqpoint{3.169407in}{2.358971in}}%
\pgfpathlineto{\pgfqpoint{3.169509in}{2.348592in}}%
\pgfpathlineto{\pgfqpoint{3.169957in}{2.388582in}}%
\pgfpathlineto{\pgfqpoint{3.170375in}{2.385382in}}%
\pgfpathlineto{\pgfqpoint{3.171072in}{2.422792in}}%
\pgfpathlineto{\pgfqpoint{3.171665in}{2.401874in}}%
\pgfpathlineto{\pgfqpoint{3.171797in}{2.387046in}}%
\pgfpathlineto{\pgfqpoint{3.172673in}{2.440897in}}%
\pgfpathlineto{\pgfqpoint{3.172707in}{2.432834in}}%
\pgfpathlineto{\pgfqpoint{3.173116in}{2.458432in}}%
\pgfpathlineto{\pgfqpoint{3.172761in}{2.421902in}}%
\pgfpathlineto{\pgfqpoint{3.173603in}{2.426319in}}%
\pgfpathlineto{\pgfqpoint{3.174445in}{2.402917in}}%
\pgfpathlineto{\pgfqpoint{3.173939in}{2.449056in}}%
\pgfpathlineto{\pgfqpoint{3.174738in}{2.416062in}}%
\pgfpathlineto{\pgfqpoint{3.175541in}{2.465022in}}%
\pgfpathlineto{\pgfqpoint{3.175921in}{2.430341in}}%
\pgfpathlineto{\pgfqpoint{3.175925in}{2.428195in}}%
\pgfpathlineto{\pgfqpoint{3.176816in}{2.488247in}}%
\pgfpathlineto{\pgfqpoint{3.176923in}{2.476800in}}%
\pgfpathlineto{\pgfqpoint{3.177853in}{2.520135in}}%
\pgfpathlineto{\pgfqpoint{3.178218in}{2.511725in}}%
\pgfpathlineto{\pgfqpoint{3.179338in}{2.469170in}}%
\pgfpathlineto{\pgfqpoint{3.178374in}{2.517422in}}%
\pgfpathlineto{\pgfqpoint{3.179436in}{2.472859in}}%
\pgfpathlineto{\pgfqpoint{3.180570in}{2.507040in}}%
\pgfpathlineto{\pgfqpoint{3.179465in}{2.464293in}}%
\pgfpathlineto{\pgfqpoint{3.180770in}{2.496066in}}%
\pgfpathlineto{\pgfqpoint{3.181909in}{2.455443in}}%
\pgfpathlineto{\pgfqpoint{3.181953in}{2.460862in}}%
\pgfpathlineto{\pgfqpoint{3.182902in}{2.485004in}}%
\pgfpathlineto{\pgfqpoint{3.182069in}{2.448666in}}%
\pgfpathlineto{\pgfqpoint{3.183384in}{2.467765in}}%
\pgfpathlineto{\pgfqpoint{3.184105in}{2.424296in}}%
\pgfpathlineto{\pgfqpoint{3.183403in}{2.472263in}}%
\pgfpathlineto{\pgfqpoint{3.184849in}{2.428540in}}%
\pgfpathlineto{\pgfqpoint{3.185901in}{2.452967in}}%
\pgfpathlineto{\pgfqpoint{3.185132in}{2.417353in}}%
\pgfpathlineto{\pgfqpoint{3.185993in}{2.436310in}}%
\pgfpathlineto{\pgfqpoint{3.186295in}{2.408790in}}%
\pgfpathlineto{\pgfqpoint{3.186826in}{2.476855in}}%
\pgfpathlineto{\pgfqpoint{3.187079in}{2.449892in}}%
\pgfpathlineto{\pgfqpoint{3.187600in}{2.475262in}}%
\pgfpathlineto{\pgfqpoint{3.188116in}{2.435906in}}%
\pgfpathlineto{\pgfqpoint{3.188165in}{2.441066in}}%
\pgfpathlineto{\pgfqpoint{3.188954in}{2.437370in}}%
\pgfpathlineto{\pgfqpoint{3.189104in}{2.463087in}}%
\pgfpathlineto{\pgfqpoint{3.189119in}{2.461141in}}%
\pgfpathlineto{\pgfqpoint{3.189163in}{2.473282in}}%
\pgfpathlineto{\pgfqpoint{3.190214in}{2.456771in}}%
\pgfpathlineto{\pgfqpoint{3.190224in}{2.455172in}}%
\pgfpathlineto{\pgfqpoint{3.190721in}{2.500345in}}%
\pgfpathlineto{\pgfqpoint{3.191169in}{2.484311in}}%
\pgfpathlineto{\pgfqpoint{3.191354in}{2.494408in}}%
\pgfpathlineto{\pgfqpoint{3.192118in}{2.456206in}}%
\pgfpathlineto{\pgfqpoint{3.192201in}{2.463197in}}%
\pgfpathlineto{\pgfqpoint{3.192853in}{2.432205in}}%
\pgfpathlineto{\pgfqpoint{3.193272in}{2.475760in}}%
\pgfpathlineto{\pgfqpoint{3.193632in}{2.497015in}}%
\pgfpathlineto{\pgfqpoint{3.193895in}{2.457451in}}%
\pgfpathlineto{\pgfqpoint{3.194022in}{2.459062in}}%
\pgfpathlineto{\pgfqpoint{3.194090in}{2.453397in}}%
\pgfpathlineto{\pgfqpoint{3.195010in}{2.508443in}}%
\pgfpathlineto{\pgfqpoint{3.195064in}{2.503432in}}%
\pgfpathlineto{\pgfqpoint{3.195068in}{2.504238in}}%
\pgfpathlineto{\pgfqpoint{3.195638in}{2.466164in}}%
\pgfpathlineto{\pgfqpoint{3.196018in}{2.478087in}}%
\pgfpathlineto{\pgfqpoint{3.197693in}{2.379668in}}%
\pgfpathlineto{\pgfqpoint{3.197790in}{2.391553in}}%
\pgfpathlineto{\pgfqpoint{3.198535in}{2.454632in}}%
\pgfpathlineto{\pgfqpoint{3.198939in}{2.420885in}}%
\pgfpathlineto{\pgfqpoint{3.199285in}{2.414499in}}%
\pgfpathlineto{\pgfqpoint{3.199786in}{2.459025in}}%
\pgfpathlineto{\pgfqpoint{3.199966in}{2.445912in}}%
\pgfpathlineto{\pgfqpoint{3.201027in}{2.468301in}}%
\pgfpathlineto{\pgfqpoint{3.200122in}{2.441280in}}%
\pgfpathlineto{\pgfqpoint{3.201173in}{2.460983in}}%
\pgfpathlineto{\pgfqpoint{3.201646in}{2.442702in}}%
\pgfpathlineto{\pgfqpoint{3.201295in}{2.470825in}}%
\pgfpathlineto{\pgfqpoint{3.202181in}{2.469474in}}%
\pgfpathlineto{\pgfqpoint{3.202950in}{2.475758in}}%
\pgfpathlineto{\pgfqpoint{3.202327in}{2.455191in}}%
\pgfpathlineto{\pgfqpoint{3.203233in}{2.465599in}}%
\pgfpathlineto{\pgfqpoint{3.204314in}{2.431432in}}%
\pgfpathlineto{\pgfqpoint{3.204367in}{2.440659in}}%
\pgfpathlineto{\pgfqpoint{3.204377in}{2.438816in}}%
\pgfpathlineto{\pgfqpoint{3.204513in}{2.455398in}}%
\pgfpathlineto{\pgfqpoint{3.204528in}{2.454329in}}%
\pgfpathlineto{\pgfqpoint{3.204586in}{2.457326in}}%
\pgfpathlineto{\pgfqpoint{3.205239in}{2.427195in}}%
\pgfpathlineto{\pgfqpoint{3.205604in}{2.440857in}}%
\pgfpathlineto{\pgfqpoint{3.206500in}{2.388983in}}%
\pgfpathlineto{\pgfqpoint{3.205657in}{2.442041in}}%
\pgfpathlineto{\pgfqpoint{3.206767in}{2.417896in}}%
\pgfpathlineto{\pgfqpoint{3.207050in}{2.447037in}}%
\pgfpathlineto{\pgfqpoint{3.207493in}{2.403563in}}%
\pgfpathlineto{\pgfqpoint{3.207907in}{2.430326in}}%
\pgfpathlineto{\pgfqpoint{3.208350in}{2.375123in}}%
\pgfpathlineto{\pgfqpoint{3.209051in}{2.407053in}}%
\pgfpathlineto{\pgfqpoint{3.210190in}{2.435622in}}%
\pgfpathlineto{\pgfqpoint{3.209168in}{2.396494in}}%
\pgfpathlineto{\pgfqpoint{3.210409in}{2.428395in}}%
\pgfpathlineto{\pgfqpoint{3.211086in}{2.381331in}}%
\pgfpathlineto{\pgfqpoint{3.211689in}{2.385538in}}%
\pgfpathlineto{\pgfqpoint{3.212084in}{2.406978in}}%
\pgfpathlineto{\pgfqpoint{3.212644in}{2.371906in}}%
\pgfpathlineto{\pgfqpoint{3.212780in}{2.373861in}}%
\pgfpathlineto{\pgfqpoint{3.212950in}{2.393535in}}%
\pgfpathlineto{\pgfqpoint{3.213150in}{2.360916in}}%
\pgfpathlineto{\pgfqpoint{3.213359in}{2.364967in}}%
\pgfpathlineto{\pgfqpoint{3.214387in}{2.325845in}}%
\pgfpathlineto{\pgfqpoint{3.213515in}{2.368526in}}%
\pgfpathlineto{\pgfqpoint{3.214508in}{2.333292in}}%
\pgfpathlineto{\pgfqpoint{3.214864in}{2.302027in}}%
\pgfpathlineto{\pgfqpoint{3.214538in}{2.334879in}}%
\pgfpathlineto{\pgfqpoint{3.215652in}{2.321459in}}%
\pgfpathlineto{\pgfqpoint{3.215906in}{2.350620in}}%
\pgfpathlineto{\pgfqpoint{3.216509in}{2.320366in}}%
\pgfpathlineto{\pgfqpoint{3.216801in}{2.343128in}}%
\pgfpathlineto{\pgfqpoint{3.217235in}{2.316601in}}%
\pgfpathlineto{\pgfqpoint{3.217829in}{2.348341in}}%
\pgfpathlineto{\pgfqpoint{3.217941in}{2.327075in}}%
\pgfpathlineto{\pgfqpoint{3.217955in}{2.330610in}}%
\pgfpathlineto{\pgfqpoint{3.218963in}{2.289644in}}%
\pgfpathlineto{\pgfqpoint{3.218968in}{2.290670in}}%
\pgfpathlineto{\pgfqpoint{3.220127in}{2.266588in}}%
\pgfpathlineto{\pgfqpoint{3.219411in}{2.310034in}}%
\pgfpathlineto{\pgfqpoint{3.220175in}{2.272250in}}%
\pgfpathlineto{\pgfqpoint{3.221125in}{2.284700in}}%
\pgfpathlineto{\pgfqpoint{3.220803in}{2.263488in}}%
\pgfpathlineto{\pgfqpoint{3.221290in}{2.277592in}}%
\pgfpathlineto{\pgfqpoint{3.221694in}{2.293109in}}%
\pgfpathlineto{\pgfqpoint{3.222108in}{2.262147in}}%
\pgfpathlineto{\pgfqpoint{3.222313in}{2.267583in}}%
\pgfpathlineto{\pgfqpoint{3.222317in}{2.265026in}}%
\pgfpathlineto{\pgfqpoint{3.222595in}{2.293487in}}%
\pgfpathlineto{\pgfqpoint{3.223418in}{2.270496in}}%
\pgfpathlineto{\pgfqpoint{3.223895in}{2.248782in}}%
\pgfpathlineto{\pgfqpoint{3.223617in}{2.278650in}}%
\pgfpathlineto{\pgfqpoint{3.224552in}{2.262702in}}%
\pgfpathlineto{\pgfqpoint{3.225312in}{2.305284in}}%
\pgfpathlineto{\pgfqpoint{3.224795in}{2.256986in}}%
\pgfpathlineto{\pgfqpoint{3.225711in}{2.286584in}}%
\pgfpathlineto{\pgfqpoint{3.226860in}{2.260073in}}%
\pgfpathlineto{\pgfqpoint{3.226519in}{2.301975in}}%
\pgfpathlineto{\pgfqpoint{3.226899in}{2.264230in}}%
\pgfpathlineto{\pgfqpoint{3.226913in}{2.268565in}}%
\pgfpathlineto{\pgfqpoint{3.227892in}{2.220541in}}%
\pgfpathlineto{\pgfqpoint{3.227911in}{2.222335in}}%
\pgfpathlineto{\pgfqpoint{3.229041in}{2.196268in}}%
\pgfpathlineto{\pgfqpoint{3.229693in}{2.188718in}}%
\pgfpathlineto{\pgfqpoint{3.230024in}{2.215373in}}%
\pgfpathlineto{\pgfqpoint{3.230107in}{2.210051in}}%
\pgfpathlineto{\pgfqpoint{3.230112in}{2.213344in}}%
\pgfpathlineto{\pgfqpoint{3.230974in}{2.139767in}}%
\pgfpathlineto{\pgfqpoint{3.231125in}{2.147148in}}%
\pgfpathlineto{\pgfqpoint{3.231212in}{2.153358in}}%
\pgfpathlineto{\pgfqpoint{3.232127in}{2.106219in}}%
\pgfpathlineto{\pgfqpoint{3.232157in}{2.110963in}}%
\pgfpathlineto{\pgfqpoint{3.233880in}{2.175470in}}%
\pgfpathlineto{\pgfqpoint{3.232205in}{2.110872in}}%
\pgfpathlineto{\pgfqpoint{3.234328in}{2.153154in}}%
\pgfpathlineto{\pgfqpoint{3.234664in}{2.150804in}}%
\pgfpathlineto{\pgfqpoint{3.235307in}{2.185978in}}%
\pgfpathlineto{\pgfqpoint{3.235326in}{2.182101in}}%
\pgfpathlineto{\pgfqpoint{3.235750in}{2.210848in}}%
\pgfpathlineto{\pgfqpoint{3.235380in}{2.178508in}}%
\pgfpathlineto{\pgfqpoint{3.236465in}{2.199491in}}%
\pgfpathlineto{\pgfqpoint{3.237366in}{2.133483in}}%
\pgfpathlineto{\pgfqpoint{3.236553in}{2.201207in}}%
\pgfpathlineto{\pgfqpoint{3.237770in}{2.157810in}}%
\pgfpathlineto{\pgfqpoint{3.238369in}{2.182202in}}%
\pgfpathlineto{\pgfqpoint{3.238077in}{2.147244in}}%
\pgfpathlineto{\pgfqpoint{3.238904in}{2.169902in}}%
\pgfpathlineto{\pgfqpoint{3.240087in}{2.138473in}}%
\pgfpathlineto{\pgfqpoint{3.239109in}{2.171576in}}%
\pgfpathlineto{\pgfqpoint{3.240097in}{2.141006in}}%
\pgfpathlineto{\pgfqpoint{3.240823in}{2.162740in}}%
\pgfpathlineto{\pgfqpoint{3.240423in}{2.140277in}}%
\pgfpathlineto{\pgfqpoint{3.241227in}{2.152555in}}%
\pgfpathlineto{\pgfqpoint{3.241256in}{2.155285in}}%
\pgfpathlineto{\pgfqpoint{3.241402in}{2.120522in}}%
\pgfpathlineto{\pgfqpoint{3.242220in}{2.131745in}}%
\pgfpathlineto{\pgfqpoint{3.243335in}{2.093467in}}%
\pgfpathlineto{\pgfqpoint{3.242682in}{2.141992in}}%
\pgfpathlineto{\pgfqpoint{3.243447in}{2.105374in}}%
\pgfpathlineto{\pgfqpoint{3.243578in}{2.119421in}}%
\pgfpathlineto{\pgfqpoint{3.244060in}{2.077314in}}%
\pgfpathlineto{\pgfqpoint{3.244489in}{2.080891in}}%
\pgfpathlineto{\pgfqpoint{3.244542in}{2.078275in}}%
\pgfpathlineto{\pgfqpoint{3.244664in}{2.106869in}}%
\pgfpathlineto{\pgfqpoint{3.245540in}{2.096619in}}%
\pgfpathlineto{\pgfqpoint{3.246611in}{2.106606in}}%
\pgfpathlineto{\pgfqpoint{3.245969in}{2.082430in}}%
\pgfpathlineto{\pgfqpoint{3.246660in}{2.099169in}}%
\pgfpathlineto{\pgfqpoint{3.246811in}{2.082949in}}%
\pgfpathlineto{\pgfqpoint{3.246689in}{2.100920in}}%
\pgfpathlineto{\pgfqpoint{3.246816in}{2.083195in}}%
\pgfpathlineto{\pgfqpoint{3.246830in}{2.077181in}}%
\pgfpathlineto{\pgfqpoint{3.247789in}{2.120768in}}%
\pgfpathlineto{\pgfqpoint{3.247867in}{2.108342in}}%
\pgfpathlineto{\pgfqpoint{3.247872in}{2.112745in}}%
\pgfpathlineto{\pgfqpoint{3.248933in}{2.089620in}}%
\pgfpathlineto{\pgfqpoint{3.248958in}{2.091745in}}%
\pgfpathlineto{\pgfqpoint{3.249591in}{2.053950in}}%
\pgfpathlineto{\pgfqpoint{3.249016in}{2.097496in}}%
\pgfpathlineto{\pgfqpoint{3.250048in}{2.091628in}}%
\pgfpathlineto{\pgfqpoint{3.250053in}{2.093942in}}%
\pgfpathlineto{\pgfqpoint{3.250876in}{2.009089in}}%
\pgfpathlineto{\pgfqpoint{3.250988in}{2.013180in}}%
\pgfpathlineto{\pgfqpoint{3.251514in}{2.027295in}}%
\pgfpathlineto{\pgfqpoint{3.252137in}{1.977445in}}%
\pgfpathlineto{\pgfqpoint{3.253325in}{1.938734in}}%
\pgfpathlineto{\pgfqpoint{3.252302in}{1.982755in}}%
\pgfpathlineto{\pgfqpoint{3.253666in}{1.953293in}}%
\pgfpathlineto{\pgfqpoint{3.253982in}{1.978665in}}%
\pgfpathlineto{\pgfqpoint{3.254328in}{1.941442in}}%
\pgfpathlineto{\pgfqpoint{3.254824in}{1.963817in}}%
\pgfpathlineto{\pgfqpoint{3.255297in}{1.930235in}}%
\pgfpathlineto{\pgfqpoint{3.255939in}{1.961155in}}%
\pgfpathlineto{\pgfqpoint{3.256343in}{1.972669in}}%
\pgfpathlineto{\pgfqpoint{3.256752in}{1.938926in}}%
\pgfpathlineto{\pgfqpoint{3.256986in}{1.943233in}}%
\pgfpathlineto{\pgfqpoint{3.257371in}{1.935742in}}%
\pgfpathlineto{\pgfqpoint{3.257887in}{1.962938in}}%
\pgfpathlineto{\pgfqpoint{3.258057in}{1.948145in}}%
\pgfpathlineto{\pgfqpoint{3.258310in}{1.978629in}}%
\pgfpathlineto{\pgfqpoint{3.259079in}{1.917914in}}%
\pgfpathlineto{\pgfqpoint{3.259094in}{1.918690in}}%
\pgfpathlineto{\pgfqpoint{3.259240in}{1.907022in}}%
\pgfpathlineto{\pgfqpoint{3.260048in}{1.957490in}}%
\pgfpathlineto{\pgfqpoint{3.260063in}{1.952487in}}%
\pgfpathlineto{\pgfqpoint{3.260077in}{1.958066in}}%
\pgfpathlineto{\pgfqpoint{3.260433in}{1.925120in}}%
\pgfpathlineto{\pgfqpoint{3.261129in}{1.927010in}}%
\pgfpathlineto{\pgfqpoint{3.261207in}{1.937535in}}%
\pgfpathlineto{\pgfqpoint{3.261718in}{1.909183in}}%
\pgfpathlineto{\pgfqpoint{3.262210in}{1.923186in}}%
\pgfpathlineto{\pgfqpoint{3.262755in}{1.904384in}}%
\pgfpathlineto{\pgfqpoint{3.262521in}{1.931501in}}%
\pgfpathlineto{\pgfqpoint{3.263335in}{1.907960in}}%
\pgfpathlineto{\pgfqpoint{3.264162in}{1.940648in}}%
\pgfpathlineto{\pgfqpoint{3.263349in}{1.907240in}}%
\pgfpathlineto{\pgfqpoint{3.264483in}{1.937045in}}%
\pgfpathlineto{\pgfqpoint{3.264498in}{1.934508in}}%
\pgfpathlineto{\pgfqpoint{3.264873in}{1.977116in}}%
\pgfpathlineto{\pgfqpoint{3.265511in}{1.952126in}}%
\pgfpathlineto{\pgfqpoint{3.265520in}{1.957739in}}%
\pgfpathlineto{\pgfqpoint{3.266411in}{1.913898in}}%
\pgfpathlineto{\pgfqpoint{3.266572in}{1.925430in}}%
\pgfpathlineto{\pgfqpoint{3.267609in}{1.943964in}}%
\pgfpathlineto{\pgfqpoint{3.266689in}{1.919022in}}%
\pgfpathlineto{\pgfqpoint{3.267711in}{1.932232in}}%
\pgfpathlineto{\pgfqpoint{3.268028in}{1.923167in}}%
\pgfpathlineto{\pgfqpoint{3.268661in}{1.963687in}}%
\pgfpathlineto{\pgfqpoint{3.268724in}{1.958085in}}%
\pgfpathlineto{\pgfqpoint{3.268851in}{1.972521in}}%
\pgfpathlineto{\pgfqpoint{3.269766in}{1.930135in}}%
\pgfpathlineto{\pgfqpoint{3.269771in}{1.930920in}}%
\pgfpathlineto{\pgfqpoint{3.269892in}{1.923357in}}%
\pgfpathlineto{\pgfqpoint{3.270379in}{1.942576in}}%
\pgfpathlineto{\pgfqpoint{3.270871in}{1.931016in}}%
\pgfpathlineto{\pgfqpoint{3.272098in}{1.963502in}}%
\pgfpathlineto{\pgfqpoint{3.270949in}{1.925905in}}%
\pgfpathlineto{\pgfqpoint{3.272137in}{1.956288in}}%
\pgfpathlineto{\pgfqpoint{3.273354in}{1.937331in}}%
\pgfpathlineto{\pgfqpoint{3.272244in}{1.972987in}}%
\pgfpathlineto{\pgfqpoint{3.273373in}{1.939581in}}%
\pgfpathlineto{\pgfqpoint{3.274498in}{1.955131in}}%
\pgfpathlineto{\pgfqpoint{3.273690in}{1.926087in}}%
\pgfpathlineto{\pgfqpoint{3.274503in}{1.953241in}}%
\pgfpathlineto{\pgfqpoint{3.274566in}{1.943178in}}%
\pgfpathlineto{\pgfqpoint{3.275004in}{1.969726in}}%
\pgfpathlineto{\pgfqpoint{3.275593in}{1.951597in}}%
\pgfpathlineto{\pgfqpoint{3.275866in}{1.965681in}}%
\pgfpathlineto{\pgfqpoint{3.276426in}{1.937260in}}%
\pgfpathlineto{\pgfqpoint{3.276699in}{1.951674in}}%
\pgfpathlineto{\pgfqpoint{3.276708in}{1.948510in}}%
\pgfpathlineto{\pgfqpoint{3.277356in}{1.981484in}}%
\pgfpathlineto{\pgfqpoint{3.277804in}{1.951428in}}%
\pgfpathlineto{\pgfqpoint{3.278996in}{2.035399in}}%
\pgfpathlineto{\pgfqpoint{3.277867in}{1.949363in}}%
\pgfpathlineto{\pgfqpoint{3.279079in}{2.028641in}}%
\pgfpathlineto{\pgfqpoint{3.279152in}{2.023824in}}%
\pgfpathlineto{\pgfqpoint{3.280033in}{2.052256in}}%
\pgfpathlineto{\pgfqpoint{3.280136in}{2.048364in}}%
\pgfpathlineto{\pgfqpoint{3.280141in}{2.049203in}}%
\pgfpathlineto{\pgfqpoint{3.280881in}{1.985549in}}%
\pgfpathlineto{\pgfqpoint{3.280885in}{1.986478in}}%
\pgfpathlineto{\pgfqpoint{3.281265in}{1.996959in}}%
\pgfpathlineto{\pgfqpoint{3.281552in}{1.976872in}}%
\pgfpathlineto{\pgfqpoint{3.282015in}{1.989781in}}%
\pgfpathlineto{\pgfqpoint{3.282896in}{1.955914in}}%
\pgfpathlineto{\pgfqpoint{3.283179in}{1.957344in}}%
\pgfpathlineto{\pgfqpoint{3.284138in}{1.997541in}}%
\pgfpathlineto{\pgfqpoint{3.284396in}{1.994572in}}%
\pgfpathlineto{\pgfqpoint{3.284663in}{1.983273in}}%
\pgfpathlineto{\pgfqpoint{3.285433in}{2.015706in}}%
\pgfpathlineto{\pgfqpoint{3.285442in}{2.013700in}}%
\pgfpathlineto{\pgfqpoint{3.285652in}{2.007693in}}%
\pgfpathlineto{\pgfqpoint{3.285924in}{2.027792in}}%
\pgfpathlineto{\pgfqpoint{3.286903in}{2.054540in}}%
\pgfpathlineto{\pgfqpoint{3.286504in}{2.022331in}}%
\pgfpathlineto{\pgfqpoint{3.287059in}{2.046103in}}%
\pgfpathlineto{\pgfqpoint{3.287361in}{2.068066in}}%
\pgfpathlineto{\pgfqpoint{3.287964in}{2.023646in}}%
\pgfpathlineto{\pgfqpoint{3.288037in}{2.026214in}}%
\pgfpathlineto{\pgfqpoint{3.288052in}{2.024537in}}%
\pgfpathlineto{\pgfqpoint{3.288782in}{2.056751in}}%
\pgfpathlineto{\pgfqpoint{3.289123in}{2.034362in}}%
\pgfpathlineto{\pgfqpoint{3.289133in}{2.032431in}}%
\pgfpathlineto{\pgfqpoint{3.289892in}{2.060314in}}%
\pgfpathlineto{\pgfqpoint{3.289965in}{2.054301in}}%
\pgfpathlineto{\pgfqpoint{3.291109in}{2.072628in}}%
\pgfpathlineto{\pgfqpoint{3.290574in}{2.038377in}}%
\pgfpathlineto{\pgfqpoint{3.291114in}{2.072430in}}%
\pgfpathlineto{\pgfqpoint{3.292171in}{2.034413in}}%
\pgfpathlineto{\pgfqpoint{3.291187in}{2.077114in}}%
\pgfpathlineto{\pgfqpoint{3.292302in}{2.048524in}}%
\pgfpathlineto{\pgfqpoint{3.292964in}{2.072803in}}%
\pgfpathlineto{\pgfqpoint{3.292482in}{2.042493in}}%
\pgfpathlineto{\pgfqpoint{3.293451in}{2.056391in}}%
\pgfpathlineto{\pgfqpoint{3.293597in}{2.045724in}}%
\pgfpathlineto{\pgfqpoint{3.294469in}{2.082975in}}%
\pgfpathlineto{\pgfqpoint{3.294527in}{2.080406in}}%
\pgfpathlineto{\pgfqpoint{3.295637in}{2.085003in}}%
\pgfpathlineto{\pgfqpoint{3.295155in}{2.052792in}}%
\pgfpathlineto{\pgfqpoint{3.295642in}{2.083379in}}%
\pgfpathlineto{\pgfqpoint{3.295700in}{2.087744in}}%
\pgfpathlineto{\pgfqpoint{3.295895in}{2.100784in}}%
\pgfpathlineto{\pgfqpoint{3.296119in}{2.074313in}}%
\pgfpathlineto{\pgfqpoint{3.296801in}{2.081631in}}%
\pgfpathlineto{\pgfqpoint{3.297404in}{2.036181in}}%
\pgfpathlineto{\pgfqpoint{3.297025in}{2.094143in}}%
\pgfpathlineto{\pgfqpoint{3.298023in}{2.036636in}}%
\pgfpathlineto{\pgfqpoint{3.298222in}{2.051805in}}%
\pgfpathlineto{\pgfqpoint{3.298982in}{2.014177in}}%
\pgfpathlineto{\pgfqpoint{3.299030in}{2.021803in}}%
\pgfpathlineto{\pgfqpoint{3.300097in}{1.956508in}}%
\pgfpathlineto{\pgfqpoint{3.299137in}{2.021824in}}%
\pgfpathlineto{\pgfqpoint{3.300530in}{1.976815in}}%
\pgfpathlineto{\pgfqpoint{3.300535in}{1.977597in}}%
\pgfpathlineto{\pgfqpoint{3.301372in}{1.938161in}}%
\pgfpathlineto{\pgfqpoint{3.301421in}{1.942523in}}%
\pgfpathlineto{\pgfqpoint{3.302531in}{1.914401in}}%
\pgfpathlineto{\pgfqpoint{3.302239in}{1.947508in}}%
\pgfpathlineto{\pgfqpoint{3.302750in}{1.927221in}}%
\pgfpathlineto{\pgfqpoint{3.302959in}{1.933283in}}%
\pgfpathlineto{\pgfqpoint{3.303563in}{1.882127in}}%
\pgfpathlineto{\pgfqpoint{3.303913in}{1.861086in}}%
\pgfpathlineto{\pgfqpoint{3.304191in}{1.914335in}}%
\pgfpathlineto{\pgfqpoint{3.304785in}{1.937003in}}%
\pgfpathlineto{\pgfqpoint{3.304405in}{1.908610in}}%
\pgfpathlineto{\pgfqpoint{3.305311in}{1.922919in}}%
\pgfpathlineto{\pgfqpoint{3.305320in}{1.918277in}}%
\pgfpathlineto{\pgfqpoint{3.306299in}{1.955990in}}%
\pgfpathlineto{\pgfqpoint{3.306377in}{1.949538in}}%
\pgfpathlineto{\pgfqpoint{3.306460in}{1.945882in}}%
\pgfpathlineto{\pgfqpoint{3.306562in}{1.967868in}}%
\pgfpathlineto{\pgfqpoint{3.306567in}{1.967814in}}%
\pgfpathlineto{\pgfqpoint{3.306815in}{1.991309in}}%
\pgfpathlineto{\pgfqpoint{3.307579in}{1.940690in}}%
\pgfpathlineto{\pgfqpoint{3.307604in}{1.942145in}}%
\pgfpathlineto{\pgfqpoint{3.308728in}{1.891038in}}%
\pgfpathlineto{\pgfqpoint{3.308816in}{1.903140in}}%
\pgfpathlineto{\pgfqpoint{3.309868in}{1.906367in}}%
\pgfpathlineto{\pgfqpoint{3.309084in}{1.877906in}}%
\pgfpathlineto{\pgfqpoint{3.309911in}{1.902991in}}%
\pgfpathlineto{\pgfqpoint{3.311245in}{1.852546in}}%
\pgfpathlineto{\pgfqpoint{3.310038in}{1.915923in}}%
\pgfpathlineto{\pgfqpoint{3.311270in}{1.853552in}}%
\pgfpathlineto{\pgfqpoint{3.312205in}{1.896702in}}%
\pgfpathlineto{\pgfqpoint{3.311401in}{1.842699in}}%
\pgfpathlineto{\pgfqpoint{3.312570in}{1.884842in}}%
\pgfpathlineto{\pgfqpoint{3.313081in}{1.851416in}}%
\pgfpathlineto{\pgfqpoint{3.312652in}{1.888230in}}%
\pgfpathlineto{\pgfqpoint{3.313685in}{1.878845in}}%
\pgfpathlineto{\pgfqpoint{3.314634in}{1.912524in}}%
\pgfpathlineto{\pgfqpoint{3.314157in}{1.878411in}}%
\pgfpathlineto{\pgfqpoint{3.314853in}{1.899805in}}%
\pgfpathlineto{\pgfqpoint{3.315637in}{1.857237in}}%
\pgfpathlineto{\pgfqpoint{3.316080in}{1.874084in}}%
\pgfpathlineto{\pgfqpoint{3.317141in}{1.898057in}}%
\pgfpathlineto{\pgfqpoint{3.316304in}{1.854823in}}%
\pgfpathlineto{\pgfqpoint{3.317200in}{1.886763in}}%
\pgfpathlineto{\pgfqpoint{3.317618in}{1.886871in}}%
\pgfpathlineto{\pgfqpoint{3.318422in}{1.850960in}}%
\pgfpathlineto{\pgfqpoint{3.319079in}{1.872914in}}%
\pgfpathlineto{\pgfqpoint{3.318704in}{1.833917in}}%
\pgfpathlineto{\pgfqpoint{3.319556in}{1.856017in}}%
\pgfpathlineto{\pgfqpoint{3.319843in}{1.842069in}}%
\pgfpathlineto{\pgfqpoint{3.320573in}{1.887484in}}%
\pgfpathlineto{\pgfqpoint{3.320627in}{1.876476in}}%
\pgfpathlineto{\pgfqpoint{3.320632in}{1.877949in}}%
\pgfpathlineto{\pgfqpoint{3.321659in}{1.854895in}}%
\pgfpathlineto{\pgfqpoint{3.321684in}{1.860784in}}%
\pgfpathlineto{\pgfqpoint{3.323757in}{1.785152in}}%
\pgfpathlineto{\pgfqpoint{3.321912in}{1.866450in}}%
\pgfpathlineto{\pgfqpoint{3.324254in}{1.819121in}}%
\pgfpathlineto{\pgfqpoint{3.324264in}{1.822453in}}%
\pgfpathlineto{\pgfqpoint{3.324926in}{1.777734in}}%
\pgfpathlineto{\pgfqpoint{3.325340in}{1.810346in}}%
\pgfpathlineto{\pgfqpoint{3.325345in}{1.810515in}}%
\pgfpathlineto{\pgfqpoint{3.325656in}{1.785470in}}%
\pgfpathlineto{\pgfqpoint{3.325924in}{1.792838in}}%
\pgfpathlineto{\pgfqpoint{3.326800in}{1.761595in}}%
\pgfpathlineto{\pgfqpoint{3.326377in}{1.804284in}}%
\pgfpathlineto{\pgfqpoint{3.327141in}{1.765433in}}%
\pgfpathlineto{\pgfqpoint{3.327555in}{1.776052in}}%
\pgfpathlineto{\pgfqpoint{3.327964in}{1.745072in}}%
\pgfpathlineto{\pgfqpoint{3.328227in}{1.755419in}}%
\pgfpathlineto{\pgfqpoint{3.328237in}{1.753649in}}%
\pgfpathlineto{\pgfqpoint{3.328373in}{1.765770in}}%
\pgfpathlineto{\pgfqpoint{3.329181in}{1.815180in}}%
\pgfpathlineto{\pgfqpoint{3.328436in}{1.763962in}}%
\pgfpathlineto{\pgfqpoint{3.329507in}{1.799312in}}%
\pgfpathlineto{\pgfqpoint{3.331226in}{1.694094in}}%
\pgfpathlineto{\pgfqpoint{3.329605in}{1.801516in}}%
\pgfpathlineto{\pgfqpoint{3.331630in}{1.728187in}}%
\pgfpathlineto{\pgfqpoint{3.332462in}{1.752383in}}%
\pgfpathlineto{\pgfqpoint{3.331654in}{1.724886in}}%
\pgfpathlineto{\pgfqpoint{3.332837in}{1.741073in}}%
\pgfpathlineto{\pgfqpoint{3.333042in}{1.721723in}}%
\pgfpathlineto{\pgfqpoint{3.333421in}{1.753394in}}%
\pgfpathlineto{\pgfqpoint{3.333850in}{1.745802in}}%
\pgfpathlineto{\pgfqpoint{3.333855in}{1.748001in}}%
\pgfpathlineto{\pgfqpoint{3.334872in}{1.707540in}}%
\pgfpathlineto{\pgfqpoint{3.336274in}{1.656068in}}%
\pgfpathlineto{\pgfqpoint{3.334940in}{1.716350in}}%
\pgfpathlineto{\pgfqpoint{3.336347in}{1.660068in}}%
\pgfpathlineto{\pgfqpoint{3.337375in}{1.699439in}}%
\pgfpathlineto{\pgfqpoint{3.336382in}{1.657952in}}%
\pgfpathlineto{\pgfqpoint{3.337487in}{1.679337in}}%
\pgfpathlineto{\pgfqpoint{3.338509in}{1.634890in}}%
\pgfpathlineto{\pgfqpoint{3.337837in}{1.700851in}}%
\pgfpathlineto{\pgfqpoint{3.338641in}{1.656996in}}%
\pgfpathlineto{\pgfqpoint{3.338689in}{1.664044in}}%
\pgfpathlineto{\pgfqpoint{3.338860in}{1.638189in}}%
\pgfpathlineto{\pgfqpoint{3.339609in}{1.645863in}}%
\pgfpathlineto{\pgfqpoint{3.339960in}{1.619147in}}%
\pgfpathlineto{\pgfqpoint{3.339687in}{1.652256in}}%
\pgfpathlineto{\pgfqpoint{3.340758in}{1.628337in}}%
\pgfpathlineto{\pgfqpoint{3.340987in}{1.633398in}}%
\pgfpathlineto{\pgfqpoint{3.340870in}{1.622424in}}%
\pgfpathlineto{\pgfqpoint{3.341187in}{1.624104in}}%
\pgfpathlineto{\pgfqpoint{3.341688in}{1.606375in}}%
\pgfpathlineto{\pgfqpoint{3.342190in}{1.672057in}}%
\pgfpathlineto{\pgfqpoint{3.342199in}{1.669862in}}%
\pgfpathlineto{\pgfqpoint{3.343066in}{1.659685in}}%
\pgfpathlineto{\pgfqpoint{3.342808in}{1.707903in}}%
\pgfpathlineto{\pgfqpoint{3.343290in}{1.671095in}}%
\pgfpathlineto{\pgfqpoint{3.343343in}{1.685663in}}%
\pgfpathlineto{\pgfqpoint{3.343723in}{1.639188in}}%
\pgfpathlineto{\pgfqpoint{3.344390in}{1.668922in}}%
\pgfpathlineto{\pgfqpoint{3.345281in}{1.634878in}}%
\pgfpathlineto{\pgfqpoint{3.344634in}{1.687965in}}%
\pgfpathlineto{\pgfqpoint{3.345568in}{1.648631in}}%
\pgfpathlineto{\pgfqpoint{3.345612in}{1.643980in}}%
\pgfpathlineto{\pgfqpoint{3.346318in}{1.674914in}}%
\pgfpathlineto{\pgfqpoint{3.346352in}{1.671451in}}%
\pgfpathlineto{\pgfqpoint{3.347496in}{1.699632in}}%
\pgfpathlineto{\pgfqpoint{3.347097in}{1.658899in}}%
\pgfpathlineto{\pgfqpoint{3.347550in}{1.695665in}}%
\pgfpathlineto{\pgfqpoint{3.348382in}{1.685711in}}%
\pgfpathlineto{\pgfqpoint{3.347886in}{1.711750in}}%
\pgfpathlineto{\pgfqpoint{3.348670in}{1.691719in}}%
\pgfpathlineto{\pgfqpoint{3.349848in}{1.648497in}}%
\pgfpathlineto{\pgfqpoint{3.348733in}{1.702797in}}%
\pgfpathlineto{\pgfqpoint{3.350057in}{1.653197in}}%
\pgfpathlineto{\pgfqpoint{3.350067in}{1.655044in}}%
\pgfpathlineto{\pgfqpoint{3.350296in}{1.616860in}}%
\pgfpathlineto{\pgfqpoint{3.350831in}{1.625455in}}%
\pgfpathlineto{\pgfqpoint{3.350870in}{1.615552in}}%
\pgfpathlineto{\pgfqpoint{3.351450in}{1.651523in}}%
\pgfpathlineto{\pgfqpoint{3.351917in}{1.632563in}}%
\pgfpathlineto{\pgfqpoint{3.352676in}{1.648469in}}%
\pgfpathlineto{\pgfqpoint{3.352409in}{1.618089in}}%
\pgfpathlineto{\pgfqpoint{3.353071in}{1.645551in}}%
\pgfpathlineto{\pgfqpoint{3.353402in}{1.647857in}}%
\pgfpathlineto{\pgfqpoint{3.354225in}{1.607928in}}%
\pgfpathlineto{\pgfqpoint{3.355505in}{1.567210in}}%
\pgfpathlineto{\pgfqpoint{3.354877in}{1.616491in}}%
\pgfpathlineto{\pgfqpoint{3.355641in}{1.571614in}}%
\pgfpathlineto{\pgfqpoint{3.355870in}{1.592867in}}%
\pgfpathlineto{\pgfqpoint{3.356328in}{1.563933in}}%
\pgfpathlineto{\pgfqpoint{3.356771in}{1.586551in}}%
\pgfpathlineto{\pgfqpoint{3.357876in}{1.571043in}}%
\pgfpathlineto{\pgfqpoint{3.357579in}{1.615712in}}%
\pgfpathlineto{\pgfqpoint{3.357944in}{1.578221in}}%
\pgfpathlineto{\pgfqpoint{3.357983in}{1.587910in}}%
\pgfpathlineto{\pgfqpoint{3.358889in}{1.546282in}}%
\pgfpathlineto{\pgfqpoint{3.358942in}{1.549824in}}%
\pgfpathlineto{\pgfqpoint{3.359950in}{1.493527in}}%
\pgfpathlineto{\pgfqpoint{3.358976in}{1.556446in}}%
\pgfpathlineto{\pgfqpoint{3.360150in}{1.497271in}}%
\pgfpathlineto{\pgfqpoint{3.360339in}{1.511142in}}%
\pgfpathlineto{\pgfqpoint{3.361167in}{1.450396in}}%
\pgfpathlineto{\pgfqpoint{3.361172in}{1.452437in}}%
\pgfpathlineto{\pgfqpoint{3.362321in}{1.434117in}}%
\pgfpathlineto{\pgfqpoint{3.361630in}{1.481631in}}%
\pgfpathlineto{\pgfqpoint{3.362374in}{1.442788in}}%
\pgfpathlineto{\pgfqpoint{3.363056in}{1.457773in}}%
\pgfpathlineto{\pgfqpoint{3.362769in}{1.436834in}}%
\pgfpathlineto{\pgfqpoint{3.363484in}{1.443022in}}%
\pgfpathlineto{\pgfqpoint{3.364312in}{1.479170in}}%
\pgfpathlineto{\pgfqpoint{3.363611in}{1.433332in}}%
\pgfpathlineto{\pgfqpoint{3.364804in}{1.453374in}}%
\pgfpathlineto{\pgfqpoint{3.365505in}{1.435448in}}%
\pgfpathlineto{\pgfqpoint{3.365865in}{1.460924in}}%
\pgfpathlineto{\pgfqpoint{3.366036in}{1.475428in}}%
\pgfpathlineto{\pgfqpoint{3.366771in}{1.419162in}}%
\pgfpathlineto{\pgfqpoint{3.366892in}{1.426588in}}%
\pgfpathlineto{\pgfqpoint{3.367000in}{1.422219in}}%
\pgfpathlineto{\pgfqpoint{3.367408in}{1.457181in}}%
\pgfpathlineto{\pgfqpoint{3.367959in}{1.437207in}}%
\pgfpathlineto{\pgfqpoint{3.368032in}{1.442289in}}%
\pgfpathlineto{\pgfqpoint{3.369005in}{1.402301in}}%
\pgfpathlineto{\pgfqpoint{3.369015in}{1.403098in}}%
\pgfpathlineto{\pgfqpoint{3.369176in}{1.394510in}}%
\pgfpathlineto{\pgfqpoint{3.369463in}{1.420699in}}%
\pgfpathlineto{\pgfqpoint{3.369468in}{1.420025in}}%
\pgfpathlineto{\pgfqpoint{3.369541in}{1.427431in}}%
\pgfpathlineto{\pgfqpoint{3.370091in}{1.394202in}}%
\pgfpathlineto{\pgfqpoint{3.370554in}{1.404900in}}%
\pgfpathlineto{\pgfqpoint{3.370675in}{1.398522in}}%
\pgfpathlineto{\pgfqpoint{3.370948in}{1.426564in}}%
\pgfpathlineto{\pgfqpoint{3.371664in}{1.403813in}}%
\pgfpathlineto{\pgfqpoint{3.372554in}{1.431245in}}%
\pgfpathlineto{\pgfqpoint{3.371698in}{1.403575in}}%
\pgfpathlineto{\pgfqpoint{3.372788in}{1.410774in}}%
\pgfpathlineto{\pgfqpoint{3.372827in}{1.404446in}}%
\pgfpathlineto{\pgfqpoint{3.373333in}{1.452780in}}%
\pgfpathlineto{\pgfqpoint{3.373445in}{1.444017in}}%
\pgfpathlineto{\pgfqpoint{3.374142in}{1.470850in}}%
\pgfpathlineto{\pgfqpoint{3.373723in}{1.444013in}}%
\pgfpathlineto{\pgfqpoint{3.374575in}{1.461293in}}%
\pgfpathlineto{\pgfqpoint{3.374624in}{1.455086in}}%
\pgfpathlineto{\pgfqpoint{3.375393in}{1.489213in}}%
\pgfpathlineto{\pgfqpoint{3.375515in}{1.482217in}}%
\pgfpathlineto{\pgfqpoint{3.376133in}{1.510205in}}%
\pgfpathlineto{\pgfqpoint{3.376625in}{1.483686in}}%
\pgfpathlineto{\pgfqpoint{3.378382in}{1.564085in}}%
\pgfpathlineto{\pgfqpoint{3.376849in}{1.473729in}}%
\pgfpathlineto{\pgfqpoint{3.378402in}{1.558492in}}%
\pgfpathlineto{\pgfqpoint{3.378757in}{1.556798in}}%
\pgfpathlineto{\pgfqpoint{3.379361in}{1.580406in}}%
\pgfpathlineto{\pgfqpoint{3.379482in}{1.562615in}}%
\pgfpathlineto{\pgfqpoint{3.379565in}{1.570904in}}%
\pgfpathlineto{\pgfqpoint{3.380368in}{1.509629in}}%
\pgfpathlineto{\pgfqpoint{3.380505in}{1.514809in}}%
\pgfpathlineto{\pgfqpoint{3.381371in}{1.469374in}}%
\pgfpathlineto{\pgfqpoint{3.380563in}{1.514941in}}%
\pgfpathlineto{\pgfqpoint{3.381678in}{1.490769in}}%
\pgfpathlineto{\pgfqpoint{3.381688in}{1.493426in}}%
\pgfpathlineto{\pgfqpoint{3.382150in}{1.461481in}}%
\pgfpathlineto{\pgfqpoint{3.382189in}{1.463570in}}%
\pgfpathlineto{\pgfqpoint{3.382433in}{1.451066in}}%
\pgfpathlineto{\pgfqpoint{3.383085in}{1.475498in}}%
\pgfpathlineto{\pgfqpoint{3.383251in}{1.471037in}}%
\pgfpathlineto{\pgfqpoint{3.384419in}{1.485529in}}%
\pgfpathlineto{\pgfqpoint{3.383796in}{1.460432in}}%
\pgfpathlineto{\pgfqpoint{3.384448in}{1.483166in}}%
\pgfpathlineto{\pgfqpoint{3.384604in}{1.467532in}}%
\pgfpathlineto{\pgfqpoint{3.385164in}{1.507844in}}%
\pgfpathlineto{\pgfqpoint{3.385261in}{1.505943in}}%
\pgfpathlineto{\pgfqpoint{3.386313in}{1.527170in}}%
\pgfpathlineto{\pgfqpoint{3.385583in}{1.484731in}}%
\pgfpathlineto{\pgfqpoint{3.386410in}{1.519842in}}%
\pgfpathlineto{\pgfqpoint{3.386834in}{1.510205in}}%
\pgfpathlineto{\pgfqpoint{3.387106in}{1.540820in}}%
\pgfpathlineto{\pgfqpoint{3.387476in}{1.533133in}}%
\pgfpathlineto{\pgfqpoint{3.388338in}{1.577144in}}%
\pgfpathlineto{\pgfqpoint{3.387734in}{1.521731in}}%
\pgfpathlineto{\pgfqpoint{3.388718in}{1.557630in}}%
\pgfpathlineto{\pgfqpoint{3.389190in}{1.542852in}}%
\pgfpathlineto{\pgfqpoint{3.389726in}{1.567000in}}%
\pgfpathlineto{\pgfqpoint{3.389735in}{1.568405in}}%
\pgfpathlineto{\pgfqpoint{3.390602in}{1.536601in}}%
\pgfpathlineto{\pgfqpoint{3.390622in}{1.540750in}}%
\pgfpathlineto{\pgfqpoint{3.391167in}{1.517655in}}%
\pgfpathlineto{\pgfqpoint{3.391736in}{1.536777in}}%
\pgfpathlineto{\pgfqpoint{3.391761in}{1.542999in}}%
\pgfpathlineto{\pgfqpoint{3.392209in}{1.515694in}}%
\pgfpathlineto{\pgfqpoint{3.392725in}{1.522873in}}%
\pgfpathlineto{\pgfqpoint{3.393840in}{1.512073in}}%
\pgfpathlineto{\pgfqpoint{3.393168in}{1.552552in}}%
\pgfpathlineto{\pgfqpoint{3.393844in}{1.513470in}}%
\pgfpathlineto{\pgfqpoint{3.393991in}{1.525882in}}%
\pgfpathlineto{\pgfqpoint{3.394701in}{1.477339in}}%
\pgfpathlineto{\pgfqpoint{3.394808in}{1.483375in}}%
\pgfpathlineto{\pgfqpoint{3.394950in}{1.443889in}}%
\pgfpathlineto{\pgfqpoint{3.395914in}{1.481310in}}%
\pgfpathlineto{\pgfqpoint{3.396498in}{1.500801in}}%
\pgfpathlineto{\pgfqpoint{3.396800in}{1.464027in}}%
\pgfpathlineto{\pgfqpoint{3.397028in}{1.487839in}}%
\pgfpathlineto{\pgfqpoint{3.398143in}{1.444099in}}%
\pgfpathlineto{\pgfqpoint{3.397267in}{1.502380in}}%
\pgfpathlineto{\pgfqpoint{3.398182in}{1.448684in}}%
\pgfpathlineto{\pgfqpoint{3.399214in}{1.511878in}}%
\pgfpathlineto{\pgfqpoint{3.398577in}{1.436065in}}%
\pgfpathlineto{\pgfqpoint{3.399453in}{1.501649in}}%
\pgfpathlineto{\pgfqpoint{3.399696in}{1.518111in}}%
\pgfpathlineto{\pgfqpoint{3.400461in}{1.480331in}}%
\pgfpathlineto{\pgfqpoint{3.401074in}{1.432892in}}%
\pgfpathlineto{\pgfqpoint{3.401615in}{1.450260in}}%
\pgfpathlineto{\pgfqpoint{3.402082in}{1.470881in}}%
\pgfpathlineto{\pgfqpoint{3.402350in}{1.432782in}}%
\pgfpathlineto{\pgfqpoint{3.402715in}{1.447954in}}%
\pgfpathlineto{\pgfqpoint{3.403533in}{1.396317in}}%
\pgfpathlineto{\pgfqpoint{3.403864in}{1.426889in}}%
\pgfpathlineto{\pgfqpoint{3.404253in}{1.407867in}}%
\pgfpathlineto{\pgfqpoint{3.404721in}{1.433063in}}%
\pgfpathlineto{\pgfqpoint{3.404794in}{1.431888in}}%
\pgfpathlineto{\pgfqpoint{3.405845in}{1.468840in}}%
\pgfpathlineto{\pgfqpoint{3.404935in}{1.431082in}}%
\pgfpathlineto{\pgfqpoint{3.406011in}{1.458144in}}%
\pgfpathlineto{\pgfqpoint{3.406025in}{1.455491in}}%
\pgfpathlineto{\pgfqpoint{3.406941in}{1.524716in}}%
\pgfpathlineto{\pgfqpoint{3.406980in}{1.514686in}}%
\pgfpathlineto{\pgfqpoint{3.407043in}{1.524859in}}%
\pgfpathlineto{\pgfqpoint{3.407953in}{1.482236in}}%
\pgfpathlineto{\pgfqpoint{3.408065in}{1.493705in}}%
\pgfpathlineto{\pgfqpoint{3.408533in}{1.482271in}}%
\pgfpathlineto{\pgfqpoint{3.408182in}{1.504261in}}%
\pgfpathlineto{\pgfqpoint{3.409112in}{1.503484in}}%
\pgfpathlineto{\pgfqpoint{3.409141in}{1.509235in}}%
\pgfpathlineto{\pgfqpoint{3.409726in}{1.480575in}}%
\pgfpathlineto{\pgfqpoint{3.410607in}{1.432420in}}%
\pgfpathlineto{\pgfqpoint{3.410909in}{1.437113in}}%
\pgfpathlineto{\pgfqpoint{3.411600in}{1.479502in}}%
\pgfpathlineto{\pgfqpoint{3.411337in}{1.432514in}}%
\pgfpathlineto{\pgfqpoint{3.412038in}{1.449928in}}%
\pgfpathlineto{\pgfqpoint{3.413255in}{1.416266in}}%
\pgfpathlineto{\pgfqpoint{3.412359in}{1.468983in}}%
\pgfpathlineto{\pgfqpoint{3.413260in}{1.417083in}}%
\pgfpathlineto{\pgfqpoint{3.413771in}{1.420834in}}%
\pgfpathlineto{\pgfqpoint{3.414253in}{1.381585in}}%
\pgfpathlineto{\pgfqpoint{3.414317in}{1.387425in}}%
\pgfpathlineto{\pgfqpoint{3.415568in}{1.415478in}}%
\pgfpathlineto{\pgfqpoint{3.414784in}{1.378492in}}%
\pgfpathlineto{\pgfqpoint{3.415689in}{1.402824in}}%
\pgfpathlineto{\pgfqpoint{3.415714in}{1.396819in}}%
\pgfpathlineto{\pgfqpoint{3.415928in}{1.420384in}}%
\pgfpathlineto{\pgfqpoint{3.416780in}{1.415578in}}%
\pgfpathlineto{\pgfqpoint{3.417583in}{1.454036in}}%
\pgfpathlineto{\pgfqpoint{3.416834in}{1.412068in}}%
\pgfpathlineto{\pgfqpoint{3.418007in}{1.425611in}}%
\pgfpathlineto{\pgfqpoint{3.418279in}{1.436602in}}%
\pgfpathlineto{\pgfqpoint{3.419166in}{1.407693in}}%
\pgfpathlineto{\pgfqpoint{3.419297in}{1.423772in}}%
\pgfpathlineto{\pgfqpoint{3.420222in}{1.392716in}}%
\pgfpathlineto{\pgfqpoint{3.420256in}{1.402073in}}%
\pgfpathlineto{\pgfqpoint{3.421322in}{1.348647in}}%
\pgfpathlineto{\pgfqpoint{3.421624in}{1.358878in}}%
\pgfpathlineto{\pgfqpoint{3.422783in}{1.394668in}}%
\pgfpathlineto{\pgfqpoint{3.421848in}{1.348002in}}%
\pgfpathlineto{\pgfqpoint{3.422851in}{1.387571in}}%
\pgfpathlineto{\pgfqpoint{3.423167in}{1.352690in}}%
\pgfpathlineto{\pgfqpoint{3.423907in}{1.389304in}}%
\pgfpathlineto{\pgfqpoint{3.423956in}{1.387996in}}%
\pgfpathlineto{\pgfqpoint{3.425461in}{1.458766in}}%
\pgfpathlineto{\pgfqpoint{3.424273in}{1.376471in}}%
\pgfpathlineto{\pgfqpoint{3.425572in}{1.454865in}}%
\pgfpathlineto{\pgfqpoint{3.426205in}{1.432441in}}%
\pgfpathlineto{\pgfqpoint{3.425855in}{1.460185in}}%
\pgfpathlineto{\pgfqpoint{3.426702in}{1.445036in}}%
\pgfpathlineto{\pgfqpoint{3.427203in}{1.464224in}}%
\pgfpathlineto{\pgfqpoint{3.426746in}{1.434995in}}%
\pgfpathlineto{\pgfqpoint{3.427802in}{1.438962in}}%
\pgfpathlineto{\pgfqpoint{3.427968in}{1.452968in}}%
\pgfpathlineto{\pgfqpoint{3.428776in}{1.427660in}}%
\pgfpathlineto{\pgfqpoint{3.428815in}{1.427830in}}%
\pgfpathlineto{\pgfqpoint{3.429978in}{1.382458in}}%
\pgfpathlineto{\pgfqpoint{3.428873in}{1.430234in}}%
\pgfpathlineto{\pgfqpoint{3.430022in}{1.387401in}}%
\pgfpathlineto{\pgfqpoint{3.430538in}{1.417757in}}%
\pgfpathlineto{\pgfqpoint{3.431191in}{1.406444in}}%
\pgfpathlineto{\pgfqpoint{3.431950in}{1.352622in}}%
\pgfpathlineto{\pgfqpoint{3.431376in}{1.416733in}}%
\pgfpathlineto{\pgfqpoint{3.432384in}{1.362700in}}%
\pgfpathlineto{\pgfqpoint{3.433002in}{1.372774in}}%
\pgfpathlineto{\pgfqpoint{3.432466in}{1.351842in}}%
\pgfpathlineto{\pgfqpoint{3.433484in}{1.362211in}}%
\pgfpathlineto{\pgfqpoint{3.434691in}{1.274782in}}%
\pgfpathlineto{\pgfqpoint{3.433508in}{1.365519in}}%
\pgfpathlineto{\pgfqpoint{3.434793in}{1.280291in}}%
\pgfpathlineto{\pgfqpoint{3.435991in}{1.319928in}}%
\pgfpathlineto{\pgfqpoint{3.434866in}{1.279664in}}%
\pgfpathlineto{\pgfqpoint{3.436050in}{1.310490in}}%
\pgfpathlineto{\pgfqpoint{3.436989in}{1.324129in}}%
\pgfpathlineto{\pgfqpoint{3.436249in}{1.297869in}}%
\pgfpathlineto{\pgfqpoint{3.437247in}{1.317485in}}%
\pgfpathlineto{\pgfqpoint{3.437349in}{1.292526in}}%
\pgfpathlineto{\pgfqpoint{3.437948in}{1.354065in}}%
\pgfpathlineto{\pgfqpoint{3.438289in}{1.345475in}}%
\pgfpathlineto{\pgfqpoint{3.439521in}{1.392629in}}%
\pgfpathlineto{\pgfqpoint{3.438333in}{1.339242in}}%
\pgfpathlineto{\pgfqpoint{3.439608in}{1.383548in}}%
\pgfpathlineto{\pgfqpoint{3.440665in}{1.341843in}}%
\pgfpathlineto{\pgfqpoint{3.439998in}{1.386434in}}%
\pgfpathlineto{\pgfqpoint{3.440767in}{1.353106in}}%
\pgfpathlineto{\pgfqpoint{3.441775in}{1.336975in}}%
\pgfpathlineto{\pgfqpoint{3.440976in}{1.359469in}}%
\pgfpathlineto{\pgfqpoint{3.442038in}{1.346947in}}%
\pgfpathlineto{\pgfqpoint{3.443133in}{1.401803in}}%
\pgfpathlineto{\pgfqpoint{3.442116in}{1.345848in}}%
\pgfpathlineto{\pgfqpoint{3.443196in}{1.392581in}}%
\pgfpathlineto{\pgfqpoint{3.443381in}{1.400204in}}%
\pgfpathlineto{\pgfqpoint{3.444156in}{1.367357in}}%
\pgfpathlineto{\pgfqpoint{3.444209in}{1.369209in}}%
\pgfpathlineto{\pgfqpoint{3.446069in}{1.434776in}}%
\pgfpathlineto{\pgfqpoint{3.444238in}{1.367713in}}%
\pgfpathlineto{\pgfqpoint{3.446215in}{1.410452in}}%
\pgfpathlineto{\pgfqpoint{3.447388in}{1.383258in}}%
\pgfpathlineto{\pgfqpoint{3.446833in}{1.419787in}}%
\pgfpathlineto{\pgfqpoint{3.447422in}{1.396196in}}%
\pgfpathlineto{\pgfqpoint{3.448167in}{1.410455in}}%
\pgfpathlineto{\pgfqpoint{3.447895in}{1.388278in}}%
\pgfpathlineto{\pgfqpoint{3.448523in}{1.396260in}}%
\pgfpathlineto{\pgfqpoint{3.449715in}{1.356140in}}%
\pgfpathlineto{\pgfqpoint{3.449730in}{1.356517in}}%
\pgfpathlineto{\pgfqpoint{3.450825in}{1.371737in}}%
\pgfpathlineto{\pgfqpoint{3.450348in}{1.322712in}}%
\pgfpathlineto{\pgfqpoint{3.450864in}{1.366111in}}%
\pgfpathlineto{\pgfqpoint{3.452437in}{1.301796in}}%
\pgfpathlineto{\pgfqpoint{3.451069in}{1.379302in}}%
\pgfpathlineto{\pgfqpoint{3.452627in}{1.309530in}}%
\pgfpathlineto{\pgfqpoint{3.452846in}{1.320671in}}%
\pgfpathlineto{\pgfqpoint{3.453581in}{1.263498in}}%
\pgfpathlineto{\pgfqpoint{3.453664in}{1.272701in}}%
\pgfpathlineto{\pgfqpoint{3.454345in}{1.265424in}}%
\pgfpathlineto{\pgfqpoint{3.453834in}{1.296562in}}%
\pgfpathlineto{\pgfqpoint{3.454730in}{1.294885in}}%
\pgfpathlineto{\pgfqpoint{3.455762in}{1.308028in}}%
\pgfpathlineto{\pgfqpoint{3.455051in}{1.270745in}}%
\pgfpathlineto{\pgfqpoint{3.455874in}{1.299503in}}%
\pgfpathlineto{\pgfqpoint{3.456156in}{1.317423in}}%
\pgfpathlineto{\pgfqpoint{3.457004in}{1.287525in}}%
\pgfpathlineto{\pgfqpoint{3.457038in}{1.297298in}}%
\pgfpathlineto{\pgfqpoint{3.457636in}{1.251362in}}%
\pgfpathlineto{\pgfqpoint{3.458109in}{1.284472in}}%
\pgfpathlineto{\pgfqpoint{3.458610in}{1.251844in}}%
\pgfpathlineto{\pgfqpoint{3.459048in}{1.307731in}}%
\pgfpathlineto{\pgfqpoint{3.459233in}{1.280711in}}%
\pgfpathlineto{\pgfqpoint{3.459238in}{1.282883in}}%
\pgfpathlineto{\pgfqpoint{3.460076in}{1.246199in}}%
\pgfpathlineto{\pgfqpoint{3.460275in}{1.255251in}}%
\pgfpathlineto{\pgfqpoint{3.460893in}{1.230375in}}%
\pgfpathlineto{\pgfqpoint{3.461127in}{1.255460in}}%
\pgfpathlineto{\pgfqpoint{3.461410in}{1.237504in}}%
\pgfpathlineto{\pgfqpoint{3.462456in}{1.251691in}}%
\pgfpathlineto{\pgfqpoint{3.462042in}{1.208678in}}%
\pgfpathlineto{\pgfqpoint{3.462578in}{1.246964in}}%
\pgfpathlineto{\pgfqpoint{3.463065in}{1.217161in}}%
\pgfpathlineto{\pgfqpoint{3.462666in}{1.253739in}}%
\pgfpathlineto{\pgfqpoint{3.463824in}{1.223107in}}%
\pgfpathlineto{\pgfqpoint{3.464861in}{1.247605in}}%
\pgfpathlineto{\pgfqpoint{3.464039in}{1.209553in}}%
\pgfpathlineto{\pgfqpoint{3.465090in}{1.233032in}}%
\pgfpathlineto{\pgfqpoint{3.465402in}{1.209384in}}%
\pgfpathlineto{\pgfqpoint{3.466117in}{1.235995in}}%
\pgfpathlineto{\pgfqpoint{3.466224in}{1.216676in}}%
\pgfpathlineto{\pgfqpoint{3.466784in}{1.212291in}}%
\pgfpathlineto{\pgfqpoint{3.466375in}{1.234082in}}%
\pgfpathlineto{\pgfqpoint{3.467169in}{1.222349in}}%
\pgfpathlineto{\pgfqpoint{3.467846in}{1.280582in}}%
\pgfpathlineto{\pgfqpoint{3.467354in}{1.221191in}}%
\pgfpathlineto{\pgfqpoint{3.468308in}{1.249431in}}%
\pgfpathlineto{\pgfqpoint{3.469233in}{1.203092in}}%
\pgfpathlineto{\pgfqpoint{3.468425in}{1.250839in}}%
\pgfpathlineto{\pgfqpoint{3.469452in}{1.216883in}}%
\pgfpathlineto{\pgfqpoint{3.469501in}{1.224043in}}%
\pgfpathlineto{\pgfqpoint{3.470348in}{1.190713in}}%
\pgfpathlineto{\pgfqpoint{3.470363in}{1.191293in}}%
\pgfpathlineto{\pgfqpoint{3.471254in}{1.144670in}}%
\pgfpathlineto{\pgfqpoint{3.470450in}{1.193805in}}%
\pgfpathlineto{\pgfqpoint{3.471551in}{1.150418in}}%
\pgfpathlineto{\pgfqpoint{3.471867in}{1.161257in}}%
\pgfpathlineto{\pgfqpoint{3.472364in}{1.122035in}}%
\pgfpathlineto{\pgfqpoint{3.472373in}{1.124144in}}%
\pgfpathlineto{\pgfqpoint{3.472549in}{1.115788in}}%
\pgfpathlineto{\pgfqpoint{3.473070in}{1.145985in}}%
\pgfpathlineto{\pgfqpoint{3.473469in}{1.129793in}}%
\pgfpathlineto{\pgfqpoint{3.474501in}{1.135921in}}%
\pgfpathlineto{\pgfqpoint{3.473785in}{1.102826in}}%
\pgfpathlineto{\pgfqpoint{3.474613in}{1.134956in}}%
\pgfpathlineto{\pgfqpoint{3.474978in}{1.112702in}}%
\pgfpathlineto{\pgfqpoint{3.475489in}{1.160677in}}%
\pgfpathlineto{\pgfqpoint{3.475674in}{1.180452in}}%
\pgfpathlineto{\pgfqpoint{3.476624in}{1.175518in}}%
\pgfpathlineto{\pgfqpoint{3.477042in}{1.188496in}}%
\pgfpathlineto{\pgfqpoint{3.477636in}{1.164728in}}%
\pgfpathlineto{\pgfqpoint{3.477661in}{1.166522in}}%
\pgfpathlineto{\pgfqpoint{3.477675in}{1.160703in}}%
\pgfpathlineto{\pgfqpoint{3.478625in}{1.207412in}}%
\pgfpathlineto{\pgfqpoint{3.479589in}{1.222423in}}%
\pgfpathlineto{\pgfqpoint{3.478873in}{1.193599in}}%
\pgfpathlineto{\pgfqpoint{3.479744in}{1.212765in}}%
\pgfpathlineto{\pgfqpoint{3.480499in}{1.188892in}}%
\pgfpathlineto{\pgfqpoint{3.479817in}{1.225114in}}%
\pgfpathlineto{\pgfqpoint{3.480918in}{1.208712in}}%
\pgfpathlineto{\pgfqpoint{3.480922in}{1.212895in}}%
\pgfpathlineto{\pgfqpoint{3.481901in}{1.153678in}}%
\pgfpathlineto{\pgfqpoint{3.481964in}{1.159687in}}%
\pgfpathlineto{\pgfqpoint{3.482437in}{1.139716in}}%
\pgfpathlineto{\pgfqpoint{3.483011in}{1.188946in}}%
\pgfpathlineto{\pgfqpoint{3.483031in}{1.187809in}}%
\pgfpathlineto{\pgfqpoint{3.483040in}{1.190480in}}%
\pgfpathlineto{\pgfqpoint{3.483892in}{1.164993in}}%
\pgfpathlineto{\pgfqpoint{3.483936in}{1.167770in}}%
\pgfpathlineto{\pgfqpoint{3.484165in}{1.135572in}}%
\pgfpathlineto{\pgfqpoint{3.483946in}{1.167915in}}%
\pgfpathlineto{\pgfqpoint{3.485109in}{1.144253in}}%
\pgfpathlineto{\pgfqpoint{3.485153in}{1.156678in}}%
\pgfpathlineto{\pgfqpoint{3.485957in}{1.123580in}}%
\pgfpathlineto{\pgfqpoint{3.486200in}{1.136792in}}%
\pgfpathlineto{\pgfqpoint{3.487183in}{1.106791in}}%
\pgfpathlineto{\pgfqpoint{3.486312in}{1.138557in}}%
\pgfpathlineto{\pgfqpoint{3.487412in}{1.109479in}}%
\pgfpathlineto{\pgfqpoint{3.488313in}{1.128889in}}%
\pgfpathlineto{\pgfqpoint{3.487617in}{1.091967in}}%
\pgfpathlineto{\pgfqpoint{3.488522in}{1.110751in}}%
\pgfpathlineto{\pgfqpoint{3.489973in}{1.071205in}}%
\pgfpathlineto{\pgfqpoint{3.488839in}{1.140919in}}%
\pgfpathlineto{\pgfqpoint{3.489978in}{1.073095in}}%
\pgfpathlineto{\pgfqpoint{3.490747in}{1.096957in}}%
\pgfpathlineto{\pgfqpoint{3.490158in}{1.065915in}}%
\pgfpathlineto{\pgfqpoint{3.491083in}{1.074388in}}%
\pgfpathlineto{\pgfqpoint{3.492271in}{1.014642in}}%
\pgfpathlineto{\pgfqpoint{3.491151in}{1.078539in}}%
\pgfpathlineto{\pgfqpoint{3.492281in}{1.018995in}}%
\pgfpathlineto{\pgfqpoint{3.492422in}{1.027349in}}%
\pgfpathlineto{\pgfqpoint{3.493040in}{1.000033in}}%
\pgfpathlineto{\pgfqpoint{3.493288in}{1.003331in}}%
\pgfpathlineto{\pgfqpoint{3.493327in}{0.997045in}}%
\pgfpathlineto{\pgfqpoint{3.493668in}{1.032676in}}%
\pgfpathlineto{\pgfqpoint{3.494374in}{1.012909in}}%
\pgfpathlineto{\pgfqpoint{3.494997in}{1.046517in}}%
\pgfpathlineto{\pgfqpoint{3.495324in}{1.009261in}}%
\pgfpathlineto{\pgfqpoint{3.495538in}{1.016744in}}%
\pgfpathlineto{\pgfqpoint{3.496146in}{0.997919in}}%
\pgfpathlineto{\pgfqpoint{3.495606in}{1.021915in}}%
\pgfpathlineto{\pgfqpoint{3.496687in}{1.002115in}}%
\pgfpathlineto{\pgfqpoint{3.496998in}{0.977880in}}%
\pgfpathlineto{\pgfqpoint{3.497534in}{1.006419in}}%
\pgfpathlineto{\pgfqpoint{3.497855in}{0.983098in}}%
\pgfpathlineto{\pgfqpoint{3.498990in}{1.022272in}}%
\pgfpathlineto{\pgfqpoint{3.497928in}{0.980059in}}%
\pgfpathlineto{\pgfqpoint{3.499024in}{1.017503in}}%
\pgfpathlineto{\pgfqpoint{3.499754in}{1.041787in}}%
\pgfpathlineto{\pgfqpoint{3.499121in}{1.008797in}}%
\pgfpathlineto{\pgfqpoint{3.500387in}{1.030834in}}%
\pgfpathlineto{\pgfqpoint{3.501234in}{0.999302in}}%
\pgfpathlineto{\pgfqpoint{3.501482in}{1.032691in}}%
\pgfpathlineto{\pgfqpoint{3.501497in}{1.030604in}}%
\pgfpathlineto{\pgfqpoint{3.501706in}{1.011824in}}%
\pgfpathlineto{\pgfqpoint{3.502709in}{1.043837in}}%
\pgfpathlineto{\pgfqpoint{3.502987in}{1.029382in}}%
\pgfpathlineto{\pgfqpoint{3.503469in}{1.067151in}}%
\pgfpathlineto{\pgfqpoint{3.503804in}{1.051568in}}%
\pgfpathlineto{\pgfqpoint{3.503834in}{1.052858in}}%
\pgfpathlineto{\pgfqpoint{3.504598in}{1.015811in}}%
\pgfpathlineto{\pgfqpoint{3.504783in}{1.022297in}}%
\pgfpathlineto{\pgfqpoint{3.505017in}{0.996372in}}%
\pgfpathlineto{\pgfqpoint{3.504793in}{1.023144in}}%
\pgfpathlineto{\pgfqpoint{3.505937in}{1.018877in}}%
\pgfpathlineto{\pgfqpoint{3.506214in}{1.037872in}}%
\pgfpathlineto{\pgfqpoint{3.506950in}{0.977891in}}%
\pgfpathlineto{\pgfqpoint{3.507135in}{0.954705in}}%
\pgfpathlineto{\pgfqpoint{3.507855in}{1.006754in}}%
\pgfpathlineto{\pgfqpoint{3.507875in}{1.004618in}}%
\pgfpathlineto{\pgfqpoint{3.508498in}{1.025615in}}%
\pgfpathlineto{\pgfqpoint{3.508785in}{0.994467in}}%
\pgfpathlineto{\pgfqpoint{3.509004in}{1.016638in}}%
\pgfpathlineto{\pgfqpoint{3.509574in}{0.993528in}}%
\pgfpathlineto{\pgfqpoint{3.510065in}{1.024280in}}%
\pgfpathlineto{\pgfqpoint{3.510075in}{1.023892in}}%
\pgfpathlineto{\pgfqpoint{3.511263in}{1.064264in}}%
\pgfpathlineto{\pgfqpoint{3.512363in}{1.057305in}}%
\pgfpathlineto{\pgfqpoint{3.511657in}{1.088513in}}%
\pgfpathlineto{\pgfqpoint{3.512388in}{1.058426in}}%
\pgfpathlineto{\pgfqpoint{3.512480in}{1.061870in}}%
\pgfpathlineto{\pgfqpoint{3.512767in}{1.045687in}}%
\pgfpathlineto{\pgfqpoint{3.512918in}{1.028752in}}%
\pgfpathlineto{\pgfqpoint{3.513707in}{1.073318in}}%
\pgfpathlineto{\pgfqpoint{3.513819in}{1.065493in}}%
\pgfpathlineto{\pgfqpoint{3.514160in}{1.084111in}}%
\pgfpathlineto{\pgfqpoint{3.514393in}{1.054277in}}%
\pgfpathlineto{\pgfqpoint{3.514924in}{1.068779in}}%
\pgfpathlineto{\pgfqpoint{3.516054in}{1.043748in}}%
\pgfpathlineto{\pgfqpoint{3.515606in}{1.080542in}}%
\pgfpathlineto{\pgfqpoint{3.516063in}{1.044392in}}%
\pgfpathlineto{\pgfqpoint{3.516901in}{1.042098in}}%
\pgfpathlineto{\pgfqpoint{3.516560in}{1.067009in}}%
\pgfpathlineto{\pgfqpoint{3.517139in}{1.053659in}}%
\pgfpathlineto{\pgfqpoint{3.517947in}{1.083021in}}%
\pgfpathlineto{\pgfqpoint{3.517251in}{1.047874in}}%
\pgfpathlineto{\pgfqpoint{3.518288in}{1.064043in}}%
\pgfpathlineto{\pgfqpoint{3.518916in}{1.019123in}}%
\pgfpathlineto{\pgfqpoint{3.519817in}{1.034392in}}%
\pgfpathlineto{\pgfqpoint{3.520630in}{1.045760in}}%
\pgfpathlineto{\pgfqpoint{3.520284in}{1.015488in}}%
\pgfpathlineto{\pgfqpoint{3.520927in}{1.038077in}}%
\pgfpathlineto{\pgfqpoint{3.520956in}{1.032206in}}%
\pgfpathlineto{\pgfqpoint{3.521808in}{1.093282in}}%
\pgfpathlineto{\pgfqpoint{3.522962in}{1.136820in}}%
\pgfpathlineto{\pgfqpoint{3.522996in}{1.131499in}}%
\pgfpathlineto{\pgfqpoint{3.523205in}{1.106810in}}%
\pgfpathlineto{\pgfqpoint{3.523989in}{1.153762in}}%
\pgfpathlineto{\pgfqpoint{3.524057in}{1.146988in}}%
\pgfpathlineto{\pgfqpoint{3.525216in}{1.197897in}}%
\pgfpathlineto{\pgfqpoint{3.524082in}{1.146093in}}%
\pgfpathlineto{\pgfqpoint{3.525236in}{1.193096in}}%
\pgfpathlineto{\pgfqpoint{3.525966in}{1.185451in}}%
\pgfpathlineto{\pgfqpoint{3.526224in}{1.207873in}}%
\pgfpathlineto{\pgfqpoint{3.526346in}{1.192657in}}%
\pgfpathlineto{\pgfqpoint{3.526979in}{1.222487in}}%
\pgfpathlineto{\pgfqpoint{3.527397in}{1.176539in}}%
\pgfpathlineto{\pgfqpoint{3.527407in}{1.176932in}}%
\pgfpathlineto{\pgfqpoint{3.527422in}{1.171711in}}%
\pgfpathlineto{\pgfqpoint{3.528059in}{1.209104in}}%
\pgfpathlineto{\pgfqpoint{3.528463in}{1.196247in}}%
\pgfpathlineto{\pgfqpoint{3.529573in}{1.170045in}}%
\pgfpathlineto{\pgfqpoint{3.528483in}{1.201086in}}%
\pgfpathlineto{\pgfqpoint{3.529729in}{1.183361in}}%
\pgfpathlineto{\pgfqpoint{3.530323in}{1.226497in}}%
\pgfpathlineto{\pgfqpoint{3.529870in}{1.182197in}}%
\pgfpathlineto{\pgfqpoint{3.530961in}{1.211236in}}%
\pgfpathlineto{\pgfqpoint{3.531190in}{1.197360in}}%
\pgfpathlineto{\pgfqpoint{3.531964in}{1.267264in}}%
\pgfpathlineto{\pgfqpoint{3.532008in}{1.254837in}}%
\pgfpathlineto{\pgfqpoint{3.533249in}{1.316966in}}%
\pgfpathlineto{\pgfqpoint{3.532047in}{1.253478in}}%
\pgfpathlineto{\pgfqpoint{3.533288in}{1.312076in}}%
\pgfpathlineto{\pgfqpoint{3.534320in}{1.303696in}}%
\pgfpathlineto{\pgfqpoint{3.533697in}{1.328632in}}%
\pgfpathlineto{\pgfqpoint{3.534408in}{1.306956in}}%
\pgfpathlineto{\pgfqpoint{3.534700in}{1.308153in}}%
\pgfpathlineto{\pgfqpoint{3.535260in}{1.288996in}}%
\pgfpathlineto{\pgfqpoint{3.535576in}{1.283165in}}%
\pgfpathlineto{\pgfqpoint{3.536199in}{1.319442in}}%
\pgfpathlineto{\pgfqpoint{3.536589in}{1.331430in}}%
\pgfpathlineto{\pgfqpoint{3.536935in}{1.304900in}}%
\pgfpathlineto{\pgfqpoint{3.537290in}{1.311409in}}%
\pgfpathlineto{\pgfqpoint{3.538376in}{1.301947in}}%
\pgfpathlineto{\pgfqpoint{3.537821in}{1.329791in}}%
\pgfpathlineto{\pgfqpoint{3.538449in}{1.305323in}}%
\pgfpathlineto{\pgfqpoint{3.538780in}{1.321313in}}%
\pgfpathlineto{\pgfqpoint{3.539466in}{1.287094in}}%
\pgfpathlineto{\pgfqpoint{3.539530in}{1.289359in}}%
\pgfpathlineto{\pgfqpoint{3.540382in}{1.312007in}}%
\pgfpathlineto{\pgfqpoint{3.539666in}{1.282832in}}%
\pgfpathlineto{\pgfqpoint{3.540713in}{1.294545in}}%
\pgfpathlineto{\pgfqpoint{3.541852in}{1.272405in}}%
\pgfpathlineto{\pgfqpoint{3.541112in}{1.308660in}}%
\pgfpathlineto{\pgfqpoint{3.541891in}{1.275064in}}%
\pgfpathlineto{\pgfqpoint{3.542918in}{1.294933in}}%
\pgfpathlineto{\pgfqpoint{3.542056in}{1.265640in}}%
\pgfpathlineto{\pgfqpoint{3.543010in}{1.280428in}}%
\pgfpathlineto{\pgfqpoint{3.544101in}{1.222469in}}%
\pgfpathlineto{\pgfqpoint{3.543054in}{1.283702in}}%
\pgfpathlineto{\pgfqpoint{3.544286in}{1.238391in}}%
\pgfpathlineto{\pgfqpoint{3.544943in}{1.270396in}}%
\pgfpathlineto{\pgfqpoint{3.544388in}{1.225811in}}%
\pgfpathlineto{\pgfqpoint{3.545396in}{1.239296in}}%
\pgfpathlineto{\pgfqpoint{3.545498in}{1.230853in}}%
\pgfpathlineto{\pgfqpoint{3.546360in}{1.280395in}}%
\pgfpathlineto{\pgfqpoint{3.546384in}{1.277360in}}%
\pgfpathlineto{\pgfqpoint{3.546881in}{1.288083in}}%
\pgfpathlineto{\pgfqpoint{3.546594in}{1.258213in}}%
\pgfpathlineto{\pgfqpoint{3.547485in}{1.276842in}}%
\pgfpathlineto{\pgfqpoint{3.547679in}{1.258837in}}%
\pgfpathlineto{\pgfqpoint{3.548488in}{1.298500in}}%
\pgfpathlineto{\pgfqpoint{3.549617in}{1.315486in}}%
\pgfpathlineto{\pgfqpoint{3.549033in}{1.284852in}}%
\pgfpathlineto{\pgfqpoint{3.549656in}{1.313964in}}%
\pgfpathlineto{\pgfqpoint{3.550868in}{1.287539in}}%
\pgfpathlineto{\pgfqpoint{3.549734in}{1.321209in}}%
\pgfpathlineto{\pgfqpoint{3.550873in}{1.287923in}}%
\pgfpathlineto{\pgfqpoint{3.551560in}{1.316260in}}%
\pgfpathlineto{\pgfqpoint{3.550951in}{1.278853in}}%
\pgfpathlineto{\pgfqpoint{3.551983in}{1.290749in}}%
\pgfpathlineto{\pgfqpoint{3.552295in}{1.286247in}}%
\pgfpathlineto{\pgfqpoint{3.552942in}{1.310768in}}%
\pgfpathlineto{\pgfqpoint{3.553054in}{1.304146in}}%
\pgfpathlineto{\pgfqpoint{3.554101in}{1.335293in}}%
\pgfpathlineto{\pgfqpoint{3.553371in}{1.287751in}}%
\pgfpathlineto{\pgfqpoint{3.554208in}{1.327486in}}%
\pgfpathlineto{\pgfqpoint{3.554413in}{1.311589in}}%
\pgfpathlineto{\pgfqpoint{3.555006in}{1.354259in}}%
\pgfpathlineto{\pgfqpoint{3.555260in}{1.340079in}}%
\pgfpathlineto{\pgfqpoint{3.555927in}{1.311788in}}%
\pgfpathlineto{\pgfqpoint{3.556360in}{1.341079in}}%
\pgfpathlineto{\pgfqpoint{3.556365in}{1.339710in}}%
\pgfpathlineto{\pgfqpoint{3.556793in}{1.356781in}}%
\pgfpathlineto{\pgfqpoint{3.557183in}{1.312999in}}%
\pgfpathlineto{\pgfqpoint{3.557426in}{1.318149in}}%
\pgfpathlineto{\pgfqpoint{3.557460in}{1.316290in}}%
\pgfpathlineto{\pgfqpoint{3.557494in}{1.323748in}}%
\pgfpathlineto{\pgfqpoint{3.557606in}{1.336719in}}%
\pgfpathlineto{\pgfqpoint{3.558346in}{1.288992in}}%
\pgfpathlineto{\pgfqpoint{3.558400in}{1.293785in}}%
\pgfpathlineto{\pgfqpoint{3.559028in}{1.289058in}}%
\pgfpathlineto{\pgfqpoint{3.559456in}{1.321929in}}%
\pgfpathlineto{\pgfqpoint{3.559471in}{1.317032in}}%
\pgfpathlineto{\pgfqpoint{3.560177in}{1.377315in}}%
\pgfpathlineto{\pgfqpoint{3.560605in}{1.339536in}}%
\pgfpathlineto{\pgfqpoint{3.560766in}{1.335819in}}%
\pgfpathlineto{\pgfqpoint{3.561418in}{1.366444in}}%
\pgfpathlineto{\pgfqpoint{3.561598in}{1.355749in}}%
\pgfpathlineto{\pgfqpoint{3.562416in}{1.388706in}}%
\pgfpathlineto{\pgfqpoint{3.562747in}{1.373061in}}%
\pgfpathlineto{\pgfqpoint{3.563575in}{1.348314in}}%
\pgfpathlineto{\pgfqpoint{3.563702in}{1.387973in}}%
\pgfpathlineto{\pgfqpoint{3.563707in}{1.387026in}}%
\pgfpathlineto{\pgfqpoint{3.565036in}{1.436310in}}%
\pgfpathlineto{\pgfqpoint{3.563794in}{1.384925in}}%
\pgfpathlineto{\pgfqpoint{3.565060in}{1.432512in}}%
\pgfpathlineto{\pgfqpoint{3.565927in}{1.405645in}}%
\pgfpathlineto{\pgfqpoint{3.565269in}{1.437727in}}%
\pgfpathlineto{\pgfqpoint{3.566185in}{1.425123in}}%
\pgfpathlineto{\pgfqpoint{3.566238in}{1.433529in}}%
\pgfpathlineto{\pgfqpoint{3.567110in}{1.387863in}}%
\pgfpathlineto{\pgfqpoint{3.567222in}{1.402499in}}%
\pgfpathlineto{\pgfqpoint{3.567840in}{1.424738in}}%
\pgfpathlineto{\pgfqpoint{3.567236in}{1.401499in}}%
\pgfpathlineto{\pgfqpoint{3.568434in}{1.415763in}}%
\pgfpathlineto{\pgfqpoint{3.568590in}{1.400267in}}%
\pgfpathlineto{\pgfqpoint{3.568906in}{1.422898in}}%
\pgfpathlineto{\pgfqpoint{3.569544in}{1.414636in}}%
\pgfpathlineto{\pgfqpoint{3.569549in}{1.414165in}}%
\pgfpathlineto{\pgfqpoint{3.570001in}{1.436951in}}%
\pgfpathlineto{\pgfqpoint{3.570459in}{1.432899in}}%
\pgfpathlineto{\pgfqpoint{3.571749in}{1.473730in}}%
\pgfpathlineto{\pgfqpoint{3.570615in}{1.421912in}}%
\pgfpathlineto{\pgfqpoint{3.571813in}{1.466804in}}%
\pgfpathlineto{\pgfqpoint{3.572996in}{1.424227in}}%
\pgfpathlineto{\pgfqpoint{3.573015in}{1.430359in}}%
\pgfpathlineto{\pgfqpoint{3.573083in}{1.427251in}}%
\pgfpathlineto{\pgfqpoint{3.573531in}{1.461097in}}%
\pgfpathlineto{\pgfqpoint{3.573687in}{1.454121in}}%
\pgfpathlineto{\pgfqpoint{3.573760in}{1.459842in}}%
\pgfpathlineto{\pgfqpoint{3.574451in}{1.386780in}}%
\pgfpathlineto{\pgfqpoint{3.574466in}{1.388792in}}%
\pgfpathlineto{\pgfqpoint{3.574476in}{1.386296in}}%
\pgfpathlineto{\pgfqpoint{3.574709in}{1.416796in}}%
\pgfpathlineto{\pgfqpoint{3.575542in}{1.395800in}}%
\pgfpathlineto{\pgfqpoint{3.575795in}{1.402790in}}%
\pgfpathlineto{\pgfqpoint{3.576301in}{1.378957in}}%
\pgfpathlineto{\pgfqpoint{3.576452in}{1.379546in}}%
\pgfpathlineto{\pgfqpoint{3.576457in}{1.376127in}}%
\pgfpathlineto{\pgfqpoint{3.577441in}{1.417275in}}%
\pgfpathlineto{\pgfqpoint{3.577523in}{1.406947in}}%
\pgfpathlineto{\pgfqpoint{3.578317in}{1.399206in}}%
\pgfpathlineto{\pgfqpoint{3.578044in}{1.423596in}}%
\pgfpathlineto{\pgfqpoint{3.578633in}{1.406124in}}%
\pgfpathlineto{\pgfqpoint{3.579583in}{1.434709in}}%
\pgfpathlineto{\pgfqpoint{3.579768in}{1.414084in}}%
\pgfpathlineto{\pgfqpoint{3.580328in}{1.385664in}}%
\pgfpathlineto{\pgfqpoint{3.579850in}{1.420867in}}%
\pgfpathlineto{\pgfqpoint{3.580887in}{1.407204in}}%
\pgfpathlineto{\pgfqpoint{3.580980in}{1.418340in}}%
\pgfpathlineto{\pgfqpoint{3.581442in}{1.393593in}}%
\pgfpathlineto{\pgfqpoint{3.581983in}{1.400536in}}%
\pgfpathlineto{\pgfqpoint{3.583210in}{1.336370in}}%
\pgfpathlineto{\pgfqpoint{3.582158in}{1.402453in}}%
\pgfpathlineto{\pgfqpoint{3.583215in}{1.337050in}}%
\pgfpathlineto{\pgfqpoint{3.584359in}{1.308240in}}%
\pgfpathlineto{\pgfqpoint{3.583906in}{1.359097in}}%
\pgfpathlineto{\pgfqpoint{3.584505in}{1.324343in}}%
\pgfpathlineto{\pgfqpoint{3.585615in}{1.341340in}}%
\pgfpathlineto{\pgfqpoint{3.584928in}{1.293554in}}%
\pgfpathlineto{\pgfqpoint{3.585639in}{1.335967in}}%
\pgfpathlineto{\pgfqpoint{3.585698in}{1.323864in}}%
\pgfpathlineto{\pgfqpoint{3.586408in}{1.357021in}}%
\pgfpathlineto{\pgfqpoint{3.586696in}{1.348391in}}%
\pgfpathlineto{\pgfqpoint{3.586993in}{1.349614in}}%
\pgfpathlineto{\pgfqpoint{3.587440in}{1.325357in}}%
\pgfpathlineto{\pgfqpoint{3.587548in}{1.326861in}}%
\pgfpathlineto{\pgfqpoint{3.587689in}{1.308181in}}%
\pgfpathlineto{\pgfqpoint{3.588609in}{1.344522in}}%
\pgfpathlineto{\pgfqpoint{3.588808in}{1.350941in}}%
\pgfpathlineto{\pgfqpoint{3.589578in}{1.294958in}}%
\pgfpathlineto{\pgfqpoint{3.589651in}{1.300082in}}%
\pgfpathlineto{\pgfqpoint{3.589880in}{1.311122in}}%
\pgfpathlineto{\pgfqpoint{3.590542in}{1.272513in}}%
\pgfpathlineto{\pgfqpoint{3.590547in}{1.269500in}}%
\pgfpathlineto{\pgfqpoint{3.591569in}{1.314479in}}%
\pgfpathlineto{\pgfqpoint{3.591588in}{1.311842in}}%
\pgfpathlineto{\pgfqpoint{3.591593in}{1.312304in}}%
\pgfpathlineto{\pgfqpoint{3.592007in}{1.271739in}}%
\pgfpathlineto{\pgfqpoint{3.592470in}{1.291232in}}%
\pgfpathlineto{\pgfqpoint{3.592499in}{1.286298in}}%
\pgfpathlineto{\pgfqpoint{3.593463in}{1.323253in}}%
\pgfpathlineto{\pgfqpoint{3.593531in}{1.320689in}}%
\pgfpathlineto{\pgfqpoint{3.593555in}{1.319021in}}%
\pgfpathlineto{\pgfqpoint{3.594344in}{1.345220in}}%
\pgfpathlineto{\pgfqpoint{3.595853in}{1.308222in}}%
\pgfpathlineto{\pgfqpoint{3.595921in}{1.314515in}}%
\pgfpathlineto{\pgfqpoint{3.596807in}{1.345000in}}%
\pgfpathlineto{\pgfqpoint{3.595975in}{1.311140in}}%
\pgfpathlineto{\pgfqpoint{3.597070in}{1.339275in}}%
\pgfpathlineto{\pgfqpoint{3.597348in}{1.330045in}}%
\pgfpathlineto{\pgfqpoint{3.597903in}{1.369683in}}%
\pgfpathlineto{\pgfqpoint{3.597913in}{1.371113in}}%
\pgfpathlineto{\pgfqpoint{3.598210in}{1.347917in}}%
\pgfpathlineto{\pgfqpoint{3.598287in}{1.348189in}}%
\pgfpathlineto{\pgfqpoint{3.598823in}{1.318238in}}%
\pgfpathlineto{\pgfqpoint{3.599407in}{1.341014in}}%
\pgfpathlineto{\pgfqpoint{3.599412in}{1.340280in}}%
\pgfpathlineto{\pgfqpoint{3.600030in}{1.387128in}}%
\pgfpathlineto{\pgfqpoint{3.600254in}{1.380701in}}%
\pgfpathlineto{\pgfqpoint{3.600658in}{1.406921in}}%
\pgfpathlineto{\pgfqpoint{3.601131in}{1.375800in}}%
\pgfpathlineto{\pgfqpoint{3.601379in}{1.385868in}}%
\pgfpathlineto{\pgfqpoint{3.601939in}{1.355136in}}%
\pgfpathlineto{\pgfqpoint{3.602348in}{1.394102in}}%
\pgfpathlineto{\pgfqpoint{3.602489in}{1.383636in}}%
\pgfpathlineto{\pgfqpoint{3.603156in}{1.389005in}}%
\pgfpathlineto{\pgfqpoint{3.603360in}{1.365167in}}%
\pgfpathlineto{\pgfqpoint{3.603555in}{1.371875in}}%
\pgfpathlineto{\pgfqpoint{3.603687in}{1.359141in}}%
\pgfpathlineto{\pgfqpoint{3.604271in}{1.394673in}}%
\pgfpathlineto{\pgfqpoint{3.604626in}{1.380838in}}%
\pgfpathlineto{\pgfqpoint{3.605040in}{1.397794in}}%
\pgfpathlineto{\pgfqpoint{3.604840in}{1.367833in}}%
\pgfpathlineto{\pgfqpoint{3.605761in}{1.392000in}}%
\pgfpathlineto{\pgfqpoint{3.605775in}{1.388809in}}%
\pgfpathlineto{\pgfqpoint{3.606622in}{1.411383in}}%
\pgfpathlineto{\pgfqpoint{3.606700in}{1.406251in}}%
\pgfpathlineto{\pgfqpoint{3.607557in}{1.450106in}}%
\pgfpathlineto{\pgfqpoint{3.606768in}{1.404201in}}%
\pgfpathlineto{\pgfqpoint{3.607854in}{1.434992in}}%
\pgfpathlineto{\pgfqpoint{3.608998in}{1.405190in}}%
\pgfpathlineto{\pgfqpoint{3.607937in}{1.448531in}}%
\pgfpathlineto{\pgfqpoint{3.609023in}{1.409252in}}%
\pgfpathlineto{\pgfqpoint{3.609237in}{1.434701in}}%
\pgfpathlineto{\pgfqpoint{3.609441in}{1.400666in}}%
\pgfpathlineto{\pgfqpoint{3.610123in}{1.407712in}}%
\pgfpathlineto{\pgfqpoint{3.610580in}{1.387033in}}%
\pgfpathlineto{\pgfqpoint{3.610761in}{1.422147in}}%
\pgfpathlineto{\pgfqpoint{3.611218in}{1.413390in}}%
\pgfpathlineto{\pgfqpoint{3.611544in}{1.424729in}}%
\pgfpathlineto{\pgfqpoint{3.611939in}{1.393939in}}%
\pgfpathlineto{\pgfqpoint{3.612226in}{1.398182in}}%
\pgfpathlineto{\pgfqpoint{3.613740in}{1.449193in}}%
\pgfpathlineto{\pgfqpoint{3.612245in}{1.397817in}}%
\pgfpathlineto{\pgfqpoint{3.613910in}{1.435130in}}%
\pgfpathlineto{\pgfqpoint{3.614173in}{1.421296in}}%
\pgfpathlineto{\pgfqpoint{3.614339in}{1.448309in}}%
\pgfpathlineto{\pgfqpoint{3.615006in}{1.439367in}}%
\pgfpathlineto{\pgfqpoint{3.615016in}{1.446207in}}%
\pgfpathlineto{\pgfqpoint{3.615522in}{1.397973in}}%
\pgfpathlineto{\pgfqpoint{3.616121in}{1.444876in}}%
\pgfpathlineto{\pgfqpoint{3.617241in}{1.433603in}}%
\pgfpathlineto{\pgfqpoint{3.616569in}{1.458488in}}%
\pgfpathlineto{\pgfqpoint{3.617250in}{1.434841in}}%
\pgfpathlineto{\pgfqpoint{3.618049in}{1.481751in}}%
\pgfpathlineto{\pgfqpoint{3.617435in}{1.429780in}}%
\pgfpathlineto{\pgfqpoint{3.618414in}{1.445312in}}%
\pgfpathlineto{\pgfqpoint{3.619105in}{1.421015in}}%
\pgfpathlineto{\pgfqpoint{3.619427in}{1.449030in}}%
\pgfpathlineto{\pgfqpoint{3.620653in}{1.480962in}}%
\pgfpathlineto{\pgfqpoint{3.620210in}{1.430255in}}%
\pgfpathlineto{\pgfqpoint{3.620678in}{1.474291in}}%
\pgfpathlineto{\pgfqpoint{3.621106in}{1.451823in}}%
\pgfpathlineto{\pgfqpoint{3.621710in}{1.493382in}}%
\pgfpathlineto{\pgfqpoint{3.621720in}{1.493011in}}%
\pgfpathlineto{\pgfqpoint{3.622888in}{1.505714in}}%
\pgfpathlineto{\pgfqpoint{3.622542in}{1.469550in}}%
\pgfpathlineto{\pgfqpoint{3.622893in}{1.503685in}}%
\pgfpathlineto{\pgfqpoint{3.622932in}{1.498686in}}%
\pgfpathlineto{\pgfqpoint{3.623735in}{1.525938in}}%
\pgfpathlineto{\pgfqpoint{3.623954in}{1.520294in}}%
\pgfpathlineto{\pgfqpoint{3.625235in}{1.580330in}}%
\pgfpathlineto{\pgfqpoint{3.624018in}{1.517646in}}%
\pgfpathlineto{\pgfqpoint{3.625240in}{1.579007in}}%
\pgfpathlineto{\pgfqpoint{3.625244in}{1.577472in}}%
\pgfpathlineto{\pgfqpoint{3.625721in}{1.614359in}}%
\pgfpathlineto{\pgfqpoint{3.626306in}{1.592723in}}%
\pgfpathlineto{\pgfqpoint{3.626793in}{1.607337in}}%
\pgfpathlineto{\pgfqpoint{3.626427in}{1.564345in}}%
\pgfpathlineto{\pgfqpoint{3.627255in}{1.586227in}}%
\pgfpathlineto{\pgfqpoint{3.627576in}{1.577323in}}%
\pgfpathlineto{\pgfqpoint{3.628258in}{1.601521in}}%
\pgfpathlineto{\pgfqpoint{3.628312in}{1.599160in}}%
\pgfpathlineto{\pgfqpoint{3.629728in}{1.660757in}}%
\pgfpathlineto{\pgfqpoint{3.628331in}{1.596683in}}%
\pgfpathlineto{\pgfqpoint{3.629918in}{1.657109in}}%
\pgfpathlineto{\pgfqpoint{3.630166in}{1.652205in}}%
\pgfpathlineto{\pgfqpoint{3.630838in}{1.695126in}}%
\pgfpathlineto{\pgfqpoint{3.630906in}{1.682789in}}%
\pgfpathlineto{\pgfqpoint{3.632226in}{1.734225in}}%
\pgfpathlineto{\pgfqpoint{3.632236in}{1.731814in}}%
\pgfpathlineto{\pgfqpoint{3.632776in}{1.703752in}}%
\pgfpathlineto{\pgfqpoint{3.633360in}{1.720129in}}%
\pgfpathlineto{\pgfqpoint{3.633419in}{1.729908in}}%
\pgfpathlineto{\pgfqpoint{3.633808in}{1.695676in}}%
\pgfpathlineto{\pgfqpoint{3.634465in}{1.717381in}}%
\pgfpathlineto{\pgfqpoint{3.635614in}{1.676797in}}%
\pgfpathlineto{\pgfqpoint{3.634504in}{1.718724in}}%
\pgfpathlineto{\pgfqpoint{3.635634in}{1.678461in}}%
\pgfpathlineto{\pgfqpoint{3.636734in}{1.666899in}}%
\pgfpathlineto{\pgfqpoint{3.636145in}{1.697930in}}%
\pgfpathlineto{\pgfqpoint{3.636778in}{1.671788in}}%
\pgfpathlineto{\pgfqpoint{3.636792in}{1.674582in}}%
\pgfpathlineto{\pgfqpoint{3.637240in}{1.644218in}}%
\pgfpathlineto{\pgfqpoint{3.637596in}{1.646751in}}%
\pgfpathlineto{\pgfqpoint{3.638336in}{1.613554in}}%
\pgfpathlineto{\pgfqpoint{3.638730in}{1.627399in}}%
\pgfpathlineto{\pgfqpoint{3.638998in}{1.615609in}}%
\pgfpathlineto{\pgfqpoint{3.639928in}{1.674608in}}%
\pgfpathlineto{\pgfqpoint{3.641096in}{1.662082in}}%
\pgfpathlineto{\pgfqpoint{3.640215in}{1.694534in}}%
\pgfpathlineto{\pgfqpoint{3.641101in}{1.662518in}}%
\pgfpathlineto{\pgfqpoint{3.641145in}{1.667928in}}%
\pgfpathlineto{\pgfqpoint{3.642007in}{1.592643in}}%
\pgfpathlineto{\pgfqpoint{3.642046in}{1.596682in}}%
\pgfpathlineto{\pgfqpoint{3.643053in}{1.567554in}}%
\pgfpathlineto{\pgfqpoint{3.643180in}{1.579527in}}%
\pgfpathlineto{\pgfqpoint{3.643190in}{1.584045in}}%
\pgfpathlineto{\pgfqpoint{3.643973in}{1.548973in}}%
\pgfpathlineto{\pgfqpoint{3.644270in}{1.564213in}}%
\pgfpathlineto{\pgfqpoint{3.644324in}{1.569039in}}%
\pgfpathlineto{\pgfqpoint{3.644933in}{1.541242in}}%
\pgfpathlineto{\pgfqpoint{3.645098in}{1.528594in}}%
\pgfpathlineto{\pgfqpoint{3.645682in}{1.581306in}}%
\pgfpathlineto{\pgfqpoint{3.645999in}{1.571692in}}%
\pgfpathlineto{\pgfqpoint{3.646140in}{1.578270in}}%
\pgfpathlineto{\pgfqpoint{3.646437in}{1.538448in}}%
\pgfpathlineto{\pgfqpoint{3.647104in}{1.571932in}}%
\pgfpathlineto{\pgfqpoint{3.647566in}{1.549696in}}%
\pgfpathlineto{\pgfqpoint{3.648146in}{1.580709in}}%
\pgfpathlineto{\pgfqpoint{3.648219in}{1.569952in}}%
\pgfpathlineto{\pgfqpoint{3.649295in}{1.640160in}}%
\pgfpathlineto{\pgfqpoint{3.649553in}{1.636016in}}%
\pgfpathlineto{\pgfqpoint{3.649563in}{1.629171in}}%
\pgfpathlineto{\pgfqpoint{3.650268in}{1.665463in}}%
\pgfpathlineto{\pgfqpoint{3.650643in}{1.654392in}}%
\pgfpathlineto{\pgfqpoint{3.650746in}{1.645346in}}%
\pgfpathlineto{\pgfqpoint{3.651096in}{1.671975in}}%
\pgfpathlineto{\pgfqpoint{3.651627in}{1.666847in}}%
\pgfpathlineto{\pgfqpoint{3.651690in}{1.673085in}}%
\pgfpathlineto{\pgfqpoint{3.652562in}{1.629097in}}%
\pgfpathlineto{\pgfqpoint{3.652654in}{1.638346in}}%
\pgfpathlineto{\pgfqpoint{3.652815in}{1.627485in}}%
\pgfpathlineto{\pgfqpoint{3.653199in}{1.662141in}}%
\pgfpathlineto{\pgfqpoint{3.653720in}{1.642969in}}%
\pgfpathlineto{\pgfqpoint{3.654027in}{1.663306in}}%
\pgfpathlineto{\pgfqpoint{3.654733in}{1.628068in}}%
\pgfpathlineto{\pgfqpoint{3.654816in}{1.635012in}}%
\pgfpathlineto{\pgfqpoint{3.655531in}{1.619532in}}%
\pgfpathlineto{\pgfqpoint{3.655819in}{1.655618in}}%
\pgfpathlineto{\pgfqpoint{3.655911in}{1.641896in}}%
\pgfpathlineto{\pgfqpoint{3.656622in}{1.603853in}}%
\pgfpathlineto{\pgfqpoint{3.656086in}{1.652419in}}%
\pgfpathlineto{\pgfqpoint{3.657109in}{1.611895in}}%
\pgfpathlineto{\pgfqpoint{3.657240in}{1.623261in}}%
\pgfpathlineto{\pgfqpoint{3.657912in}{1.577760in}}%
\pgfpathlineto{\pgfqpoint{3.658116in}{1.590325in}}%
\pgfpathlineto{\pgfqpoint{3.659295in}{1.566571in}}%
\pgfpathlineto{\pgfqpoint{3.658233in}{1.597462in}}%
\pgfpathlineto{\pgfqpoint{3.659309in}{1.570031in}}%
\pgfpathlineto{\pgfqpoint{3.659679in}{1.593821in}}%
\pgfpathlineto{\pgfqpoint{3.660390in}{1.559431in}}%
\pgfpathlineto{\pgfqpoint{3.660614in}{1.577343in}}%
\pgfpathlineto{\pgfqpoint{3.661271in}{1.546716in}}%
\pgfpathlineto{\pgfqpoint{3.661442in}{1.548839in}}%
\pgfpathlineto{\pgfqpoint{3.662445in}{1.524652in}}%
\pgfpathlineto{\pgfqpoint{3.661612in}{1.559924in}}%
\pgfpathlineto{\pgfqpoint{3.662586in}{1.537385in}}%
\pgfpathlineto{\pgfqpoint{3.663491in}{1.555349in}}%
\pgfpathlineto{\pgfqpoint{3.662927in}{1.521163in}}%
\pgfpathlineto{\pgfqpoint{3.663720in}{1.543320in}}%
\pgfpathlineto{\pgfqpoint{3.663891in}{1.533185in}}%
\pgfpathlineto{\pgfqpoint{3.664450in}{1.578718in}}%
\pgfpathlineto{\pgfqpoint{3.664816in}{1.548678in}}%
\pgfpathlineto{\pgfqpoint{3.665210in}{1.569418in}}%
\pgfpathlineto{\pgfqpoint{3.665492in}{1.541921in}}%
\pgfpathlineto{\pgfqpoint{3.665794in}{1.545497in}}%
\pgfpathlineto{\pgfqpoint{3.666651in}{1.486953in}}%
\pgfpathlineto{\pgfqpoint{3.666967in}{1.509804in}}%
\pgfpathlineto{\pgfqpoint{3.667240in}{1.491290in}}%
\pgfpathlineto{\pgfqpoint{3.667571in}{1.519042in}}%
\pgfpathlineto{\pgfqpoint{3.668165in}{1.494544in}}%
\pgfpathlineto{\pgfqpoint{3.668404in}{1.487766in}}%
\pgfpathlineto{\pgfqpoint{3.668774in}{1.511525in}}%
\pgfpathlineto{\pgfqpoint{3.669781in}{1.530753in}}%
\pgfpathlineto{\pgfqpoint{3.669290in}{1.503189in}}%
\pgfpathlineto{\pgfqpoint{3.669908in}{1.520467in}}%
\pgfpathlineto{\pgfqpoint{3.669923in}{1.520921in}}%
\pgfpathlineto{\pgfqpoint{3.670083in}{1.502457in}}%
\pgfpathlineto{\pgfqpoint{3.670093in}{1.498413in}}%
\pgfpathlineto{\pgfqpoint{3.671067in}{1.555142in}}%
\pgfpathlineto{\pgfqpoint{3.671125in}{1.552541in}}%
\pgfpathlineto{\pgfqpoint{3.672045in}{1.555303in}}%
\pgfpathlineto{\pgfqpoint{3.671529in}{1.521950in}}%
\pgfpathlineto{\pgfqpoint{3.672084in}{1.545759in}}%
\pgfpathlineto{\pgfqpoint{3.673185in}{1.482223in}}%
\pgfpathlineto{\pgfqpoint{3.672269in}{1.548367in}}%
\pgfpathlineto{\pgfqpoint{3.673384in}{1.489132in}}%
\pgfpathlineto{\pgfqpoint{3.673413in}{1.484417in}}%
\pgfpathlineto{\pgfqpoint{3.674859in}{1.544987in}}%
\pgfpathlineto{\pgfqpoint{3.675848in}{1.521503in}}%
\pgfpathlineto{\pgfqpoint{3.675682in}{1.545789in}}%
\pgfpathlineto{\pgfqpoint{3.675984in}{1.536208in}}%
\pgfpathlineto{\pgfqpoint{3.676130in}{1.528243in}}%
\pgfpathlineto{\pgfqpoint{3.677133in}{1.558646in}}%
\pgfpathlineto{\pgfqpoint{3.677138in}{1.559872in}}%
\pgfpathlineto{\pgfqpoint{3.677995in}{1.495927in}}%
\pgfpathlineto{\pgfqpoint{3.678014in}{1.500813in}}%
\pgfpathlineto{\pgfqpoint{3.678121in}{1.487196in}}%
\pgfpathlineto{\pgfqpoint{3.678744in}{1.512061in}}%
\pgfpathlineto{\pgfqpoint{3.679066in}{1.508216in}}%
\pgfpathlineto{\pgfqpoint{3.679100in}{1.511990in}}%
\pgfpathlineto{\pgfqpoint{3.680059in}{1.484847in}}%
\pgfpathlineto{\pgfqpoint{3.680176in}{1.508893in}}%
\pgfpathlineto{\pgfqpoint{3.680599in}{1.494881in}}%
\pgfpathlineto{\pgfqpoint{3.680960in}{1.519919in}}%
\pgfpathlineto{\pgfqpoint{3.681325in}{1.503390in}}%
\pgfpathlineto{\pgfqpoint{3.681870in}{1.519221in}}%
\pgfpathlineto{\pgfqpoint{3.681700in}{1.491458in}}%
\pgfpathlineto{\pgfqpoint{3.682420in}{1.500734in}}%
\pgfpathlineto{\pgfqpoint{3.683024in}{1.484694in}}%
\pgfpathlineto{\pgfqpoint{3.682712in}{1.509419in}}%
\pgfpathlineto{\pgfqpoint{3.683554in}{1.489266in}}%
\pgfpathlineto{\pgfqpoint{3.684314in}{1.504510in}}%
\pgfpathlineto{\pgfqpoint{3.683788in}{1.476548in}}%
\pgfpathlineto{\pgfqpoint{3.684630in}{1.483295in}}%
\pgfpathlineto{\pgfqpoint{3.685404in}{1.453556in}}%
\pgfpathlineto{\pgfqpoint{3.685127in}{1.496298in}}%
\pgfpathlineto{\pgfqpoint{3.685774in}{1.466206in}}%
\pgfpathlineto{\pgfqpoint{3.687074in}{1.530474in}}%
\pgfpathlineto{\pgfqpoint{3.685955in}{1.465733in}}%
\pgfpathlineto{\pgfqpoint{3.687094in}{1.530241in}}%
\pgfpathlineto{\pgfqpoint{3.687181in}{1.522068in}}%
\pgfpathlineto{\pgfqpoint{3.687809in}{1.555180in}}%
\pgfpathlineto{\pgfqpoint{3.688204in}{1.530390in}}%
\pgfpathlineto{\pgfqpoint{3.689363in}{1.570393in}}%
\pgfpathlineto{\pgfqpoint{3.689445in}{1.559207in}}%
\pgfpathlineto{\pgfqpoint{3.690512in}{1.536829in}}%
\pgfpathlineto{\pgfqpoint{3.689572in}{1.568840in}}%
\pgfpathlineto{\pgfqpoint{3.690585in}{1.545437in}}%
\pgfpathlineto{\pgfqpoint{3.691261in}{1.562358in}}%
\pgfpathlineto{\pgfqpoint{3.691471in}{1.535852in}}%
\pgfpathlineto{\pgfqpoint{3.691724in}{1.559947in}}%
\pgfpathlineto{\pgfqpoint{3.692970in}{1.519435in}}%
\pgfpathlineto{\pgfqpoint{3.691987in}{1.564105in}}%
\pgfpathlineto{\pgfqpoint{3.692975in}{1.520398in}}%
\pgfpathlineto{\pgfqpoint{3.693087in}{1.534700in}}%
\pgfpathlineto{\pgfqpoint{3.693856in}{1.499558in}}%
\pgfpathlineto{\pgfqpoint{3.694051in}{1.501481in}}%
\pgfpathlineto{\pgfqpoint{3.695278in}{1.456675in}}%
\pgfpathlineto{\pgfqpoint{3.694139in}{1.508528in}}%
\pgfpathlineto{\pgfqpoint{3.695297in}{1.457556in}}%
\pgfpathlineto{\pgfqpoint{3.696091in}{1.472033in}}%
\pgfpathlineto{\pgfqpoint{3.696266in}{1.447214in}}%
\pgfpathlineto{\pgfqpoint{3.696510in}{1.432585in}}%
\pgfpathlineto{\pgfqpoint{3.697293in}{1.464233in}}%
\pgfpathlineto{\pgfqpoint{3.697366in}{1.450689in}}%
\pgfpathlineto{\pgfqpoint{3.698583in}{1.508422in}}%
\pgfpathlineto{\pgfqpoint{3.697444in}{1.442200in}}%
\pgfpathlineto{\pgfqpoint{3.698661in}{1.503523in}}%
\pgfpathlineto{\pgfqpoint{3.698744in}{1.496364in}}%
\pgfpathlineto{\pgfqpoint{3.699591in}{1.541005in}}%
\pgfpathlineto{\pgfqpoint{3.699635in}{1.535563in}}%
\pgfpathlineto{\pgfqpoint{3.700808in}{1.571171in}}%
\pgfpathlineto{\pgfqpoint{3.700059in}{1.516308in}}%
\pgfpathlineto{\pgfqpoint{3.700833in}{1.569147in}}%
\pgfpathlineto{\pgfqpoint{3.701027in}{1.548020in}}%
\pgfpathlineto{\pgfqpoint{3.701787in}{1.597540in}}%
\pgfpathlineto{\pgfqpoint{3.701860in}{1.586725in}}%
\pgfpathlineto{\pgfqpoint{3.702775in}{1.618560in}}%
\pgfpathlineto{\pgfqpoint{3.702196in}{1.575578in}}%
\pgfpathlineto{\pgfqpoint{3.702985in}{1.598601in}}%
\pgfpathlineto{\pgfqpoint{3.704056in}{1.548913in}}%
\pgfpathlineto{\pgfqpoint{3.703077in}{1.599660in}}%
\pgfpathlineto{\pgfqpoint{3.704143in}{1.571912in}}%
\pgfpathlineto{\pgfqpoint{3.704713in}{1.600401in}}%
\pgfpathlineto{\pgfqpoint{3.705141in}{1.562663in}}%
\pgfpathlineto{\pgfqpoint{3.705224in}{1.567079in}}%
\pgfpathlineto{\pgfqpoint{3.705789in}{1.548181in}}%
\pgfpathlineto{\pgfqpoint{3.706232in}{1.590612in}}%
\pgfpathlineto{\pgfqpoint{3.706305in}{1.582311in}}%
\pgfpathlineto{\pgfqpoint{3.706728in}{1.599246in}}%
\pgfpathlineto{\pgfqpoint{3.707025in}{1.568504in}}%
\pgfpathlineto{\pgfqpoint{3.707074in}{1.569369in}}%
\pgfpathlineto{\pgfqpoint{3.707415in}{1.561110in}}%
\pgfpathlineto{\pgfqpoint{3.707926in}{1.585385in}}%
\pgfpathlineto{\pgfqpoint{3.708174in}{1.570682in}}%
\pgfpathlineto{\pgfqpoint{3.708652in}{1.564192in}}%
\pgfpathlineto{\pgfqpoint{3.709323in}{1.600320in}}%
\pgfpathlineto{\pgfqpoint{3.710234in}{1.564562in}}%
\pgfpathlineto{\pgfqpoint{3.709635in}{1.620504in}}%
\pgfpathlineto{\pgfqpoint{3.710502in}{1.569810in}}%
\pgfpathlineto{\pgfqpoint{3.710521in}{1.574687in}}%
\pgfpathlineto{\pgfqpoint{3.711519in}{1.532787in}}%
\pgfpathlineto{\pgfqpoint{3.711548in}{1.534462in}}%
\pgfpathlineto{\pgfqpoint{3.711558in}{1.534818in}}%
\pgfpathlineto{\pgfqpoint{3.711977in}{1.507204in}}%
\pgfpathlineto{\pgfqpoint{3.712167in}{1.498331in}}%
\pgfpathlineto{\pgfqpoint{3.712897in}{1.516036in}}%
\pgfpathlineto{\pgfqpoint{3.713092in}{1.499871in}}%
\pgfpathlineto{\pgfqpoint{3.713316in}{1.527600in}}%
\pgfpathlineto{\pgfqpoint{3.713160in}{1.492797in}}%
\pgfpathlineto{\pgfqpoint{3.714226in}{1.508884in}}%
\pgfpathlineto{\pgfqpoint{3.715209in}{1.489255in}}%
\pgfpathlineto{\pgfqpoint{3.714796in}{1.528899in}}%
\pgfpathlineto{\pgfqpoint{3.715424in}{1.499503in}}%
\pgfpathlineto{\pgfqpoint{3.715560in}{1.503595in}}%
\pgfpathlineto{\pgfqpoint{3.715730in}{1.480190in}}%
\pgfpathlineto{\pgfqpoint{3.716436in}{1.495810in}}%
\pgfpathlineto{\pgfqpoint{3.717595in}{1.474994in}}%
\pgfpathlineto{\pgfqpoint{3.716933in}{1.514863in}}%
\pgfpathlineto{\pgfqpoint{3.717610in}{1.475984in}}%
\pgfpathlineto{\pgfqpoint{3.718617in}{1.486560in}}%
\pgfpathlineto{\pgfqpoint{3.718389in}{1.464572in}}%
\pgfpathlineto{\pgfqpoint{3.718734in}{1.485195in}}%
\pgfpathlineto{\pgfqpoint{3.719241in}{1.438093in}}%
\pgfpathlineto{\pgfqpoint{3.718807in}{1.487111in}}%
\pgfpathlineto{\pgfqpoint{3.719883in}{1.467509in}}%
\pgfpathlineto{\pgfqpoint{3.720312in}{1.473615in}}%
\pgfpathlineto{\pgfqpoint{3.720867in}{1.447193in}}%
\pgfpathlineto{\pgfqpoint{3.720886in}{1.448500in}}%
\pgfpathlineto{\pgfqpoint{3.722016in}{1.416446in}}%
\pgfpathlineto{\pgfqpoint{3.722098in}{1.420293in}}%
\pgfpathlineto{\pgfqpoint{3.722595in}{1.439021in}}%
\pgfpathlineto{\pgfqpoint{3.722960in}{1.407190in}}%
\pgfpathlineto{\pgfqpoint{3.723199in}{1.419499in}}%
\pgfpathlineto{\pgfqpoint{3.724425in}{1.357438in}}%
\pgfpathlineto{\pgfqpoint{3.723296in}{1.420540in}}%
\pgfpathlineto{\pgfqpoint{3.724435in}{1.357651in}}%
\pgfpathlineto{\pgfqpoint{3.725058in}{1.390664in}}%
\pgfpathlineto{\pgfqpoint{3.724474in}{1.348466in}}%
\pgfpathlineto{\pgfqpoint{3.725511in}{1.351939in}}%
\pgfpathlineto{\pgfqpoint{3.725526in}{1.346964in}}%
\pgfpathlineto{\pgfqpoint{3.726568in}{1.385156in}}%
\pgfpathlineto{\pgfqpoint{3.726587in}{1.385751in}}%
\pgfpathlineto{\pgfqpoint{3.727288in}{1.354209in}}%
\pgfpathlineto{\pgfqpoint{3.727356in}{1.367445in}}%
\pgfpathlineto{\pgfqpoint{3.727639in}{1.329409in}}%
\pgfpathlineto{\pgfqpoint{3.728184in}{1.343857in}}%
\pgfpathlineto{\pgfqpoint{3.728272in}{1.320812in}}%
\pgfpathlineto{\pgfqpoint{3.728827in}{1.347239in}}%
\pgfpathlineto{\pgfqpoint{3.729313in}{1.335878in}}%
\pgfpathlineto{\pgfqpoint{3.730253in}{1.368007in}}%
\pgfpathlineto{\pgfqpoint{3.729464in}{1.330300in}}%
\pgfpathlineto{\pgfqpoint{3.730477in}{1.355450in}}%
\pgfpathlineto{\pgfqpoint{3.730492in}{1.349787in}}%
\pgfpathlineto{\pgfqpoint{3.731037in}{1.388007in}}%
\pgfpathlineto{\pgfqpoint{3.731582in}{1.360192in}}%
\pgfpathlineto{\pgfqpoint{3.732420in}{1.377307in}}%
\pgfpathlineto{\pgfqpoint{3.732205in}{1.346200in}}%
\pgfpathlineto{\pgfqpoint{3.732687in}{1.358641in}}%
\pgfpathlineto{\pgfqpoint{3.733437in}{1.341983in}}%
\pgfpathlineto{\pgfqpoint{3.732955in}{1.372102in}}%
\pgfpathlineto{\pgfqpoint{3.733792in}{1.360259in}}%
\pgfpathlineto{\pgfqpoint{3.734382in}{1.378938in}}%
\pgfpathlineto{\pgfqpoint{3.734800in}{1.354961in}}%
\pgfpathlineto{\pgfqpoint{3.734805in}{1.355421in}}%
\pgfpathlineto{\pgfqpoint{3.735745in}{1.320828in}}%
\pgfpathlineto{\pgfqpoint{3.735088in}{1.358551in}}%
\pgfpathlineto{\pgfqpoint{3.735944in}{1.332866in}}%
\pgfpathlineto{\pgfqpoint{3.736665in}{1.350891in}}%
\pgfpathlineto{\pgfqpoint{3.736519in}{1.319593in}}%
\pgfpathlineto{\pgfqpoint{3.737069in}{1.341550in}}%
\pgfpathlineto{\pgfqpoint{3.737940in}{1.376226in}}%
\pgfpathlineto{\pgfqpoint{3.737400in}{1.338630in}}%
\pgfpathlineto{\pgfqpoint{3.738262in}{1.350937in}}%
\pgfpathlineto{\pgfqpoint{3.738788in}{1.328361in}}%
\pgfpathlineto{\pgfqpoint{3.738335in}{1.360970in}}%
\pgfpathlineto{\pgfqpoint{3.739386in}{1.343830in}}%
\pgfpathlineto{\pgfqpoint{3.740316in}{1.377828in}}%
\pgfpathlineto{\pgfqpoint{3.739513in}{1.337592in}}%
\pgfpathlineto{\pgfqpoint{3.740555in}{1.376567in}}%
\pgfpathlineto{\pgfqpoint{3.740599in}{1.366872in}}%
\pgfpathlineto{\pgfqpoint{3.741232in}{1.401284in}}%
\pgfpathlineto{\pgfqpoint{3.741485in}{1.395843in}}%
\pgfpathlineto{\pgfqpoint{3.741494in}{1.398050in}}%
\pgfpathlineto{\pgfqpoint{3.741855in}{1.369038in}}%
\pgfpathlineto{\pgfqpoint{3.742536in}{1.373291in}}%
\pgfpathlineto{\pgfqpoint{3.742648in}{1.362628in}}%
\pgfpathlineto{\pgfqpoint{3.743325in}{1.392708in}}%
\pgfpathlineto{\pgfqpoint{3.743622in}{1.376549in}}%
\pgfpathlineto{\pgfqpoint{3.743685in}{1.391099in}}%
\pgfpathlineto{\pgfqpoint{3.744338in}{1.348988in}}%
\pgfpathlineto{\pgfqpoint{3.744713in}{1.361472in}}%
\pgfpathlineto{\pgfqpoint{3.745492in}{1.398702in}}%
\pgfpathlineto{\pgfqpoint{3.745112in}{1.346246in}}%
\pgfpathlineto{\pgfqpoint{3.745988in}{1.384029in}}%
\pgfpathlineto{\pgfqpoint{3.746972in}{1.360181in}}%
\pgfpathlineto{\pgfqpoint{3.746236in}{1.389609in}}%
\pgfpathlineto{\pgfqpoint{3.747127in}{1.368205in}}%
\pgfpathlineto{\pgfqpoint{3.747239in}{1.372881in}}%
\pgfpathlineto{\pgfqpoint{3.747609in}{1.346199in}}%
\pgfpathlineto{\pgfqpoint{3.748174in}{1.356805in}}%
\pgfpathlineto{\pgfqpoint{3.748865in}{1.375606in}}%
\pgfpathlineto{\pgfqpoint{3.749318in}{1.344141in}}%
\pgfpathlineto{\pgfqpoint{3.749922in}{1.387116in}}%
\pgfpathlineto{\pgfqpoint{3.749425in}{1.336967in}}%
\pgfpathlineto{\pgfqpoint{3.750521in}{1.366077in}}%
\pgfpathlineto{\pgfqpoint{3.750662in}{1.383332in}}%
\pgfpathlineto{\pgfqpoint{3.751373in}{1.350774in}}%
\pgfpathlineto{\pgfqpoint{3.751572in}{1.357175in}}%
\pgfpathlineto{\pgfqpoint{3.751587in}{1.350003in}}%
\pgfpathlineto{\pgfqpoint{3.752239in}{1.389171in}}%
\pgfpathlineto{\pgfqpoint{3.752668in}{1.370476in}}%
\pgfpathlineto{\pgfqpoint{3.752721in}{1.366301in}}%
\pgfpathlineto{\pgfqpoint{3.752682in}{1.373496in}}%
\pgfpathlineto{\pgfqpoint{3.752736in}{1.366394in}}%
\pgfpathlineto{\pgfqpoint{3.753870in}{1.343949in}}%
\pgfpathlineto{\pgfqpoint{3.753919in}{1.352218in}}%
\pgfpathlineto{\pgfqpoint{3.754571in}{1.306832in}}%
\pgfpathlineto{\pgfqpoint{3.754863in}{1.321733in}}%
\pgfpathlineto{\pgfqpoint{3.756407in}{1.389686in}}%
\pgfpathlineto{\pgfqpoint{3.755097in}{1.287380in}}%
\pgfpathlineto{\pgfqpoint{3.756524in}{1.378649in}}%
\pgfpathlineto{\pgfqpoint{3.757507in}{1.359156in}}%
\pgfpathlineto{\pgfqpoint{3.757614in}{1.383335in}}%
\pgfpathlineto{\pgfqpoint{3.757619in}{1.382885in}}%
\pgfpathlineto{\pgfqpoint{3.757867in}{1.376608in}}%
\pgfpathlineto{\pgfqpoint{3.758164in}{1.415645in}}%
\pgfpathlineto{\pgfqpoint{3.758242in}{1.414213in}}%
\pgfpathlineto{\pgfqpoint{3.759498in}{1.452788in}}%
\pgfpathlineto{\pgfqpoint{3.758301in}{1.409742in}}%
\pgfpathlineto{\pgfqpoint{3.759503in}{1.450162in}}%
\pgfpathlineto{\pgfqpoint{3.759523in}{1.443786in}}%
\pgfpathlineto{\pgfqpoint{3.760448in}{1.489695in}}%
\pgfpathlineto{\pgfqpoint{3.760574in}{1.479364in}}%
\pgfpathlineto{\pgfqpoint{3.761358in}{1.489405in}}%
\pgfpathlineto{\pgfqpoint{3.760740in}{1.467987in}}%
\pgfpathlineto{\pgfqpoint{3.761660in}{1.475488in}}%
\pgfpathlineto{\pgfqpoint{3.761665in}{1.473694in}}%
\pgfpathlineto{\pgfqpoint{3.762697in}{1.501124in}}%
\pgfpathlineto{\pgfqpoint{3.762702in}{1.500454in}}%
\pgfpathlineto{\pgfqpoint{3.763232in}{1.491573in}}%
\pgfpathlineto{\pgfqpoint{3.762862in}{1.522322in}}%
\pgfpathlineto{\pgfqpoint{3.763505in}{1.512649in}}%
\pgfpathlineto{\pgfqpoint{3.763938in}{1.561857in}}%
\pgfpathlineto{\pgfqpoint{3.764634in}{1.527098in}}%
\pgfpathlineto{\pgfqpoint{3.765740in}{1.518988in}}%
\pgfpathlineto{\pgfqpoint{3.765559in}{1.545207in}}%
\pgfpathlineto{\pgfqpoint{3.765754in}{1.523746in}}%
\pgfpathlineto{\pgfqpoint{3.766256in}{1.545945in}}%
\pgfpathlineto{\pgfqpoint{3.766767in}{1.516899in}}%
\pgfpathlineto{\pgfqpoint{3.768062in}{1.474196in}}%
\pgfpathlineto{\pgfqpoint{3.766889in}{1.529360in}}%
\pgfpathlineto{\pgfqpoint{3.768101in}{1.477867in}}%
\pgfpathlineto{\pgfqpoint{3.768140in}{1.488140in}}%
\pgfpathlineto{\pgfqpoint{3.769123in}{1.446851in}}%
\pgfpathlineto{\pgfqpoint{3.769143in}{1.448775in}}%
\pgfpathlineto{\pgfqpoint{3.770004in}{1.446356in}}%
\pgfpathlineto{\pgfqpoint{3.769425in}{1.470764in}}%
\pgfpathlineto{\pgfqpoint{3.770107in}{1.456022in}}%
\pgfpathlineto{\pgfqpoint{3.770141in}{1.465803in}}%
\pgfpathlineto{\pgfqpoint{3.770579in}{1.434599in}}%
\pgfpathlineto{\pgfqpoint{3.771187in}{1.441317in}}%
\pgfpathlineto{\pgfqpoint{3.771436in}{1.436808in}}%
\pgfpathlineto{\pgfqpoint{3.772117in}{1.479827in}}%
\pgfpathlineto{\pgfqpoint{3.772239in}{1.467984in}}%
\pgfpathlineto{\pgfqpoint{3.772249in}{1.468230in}}%
\pgfpathlineto{\pgfqpoint{3.772293in}{1.456669in}}%
\pgfpathlineto{\pgfqpoint{3.773330in}{1.429219in}}%
\pgfpathlineto{\pgfqpoint{3.772409in}{1.466511in}}%
\pgfpathlineto{\pgfqpoint{3.773442in}{1.440798in}}%
\pgfpathlineto{\pgfqpoint{3.773451in}{1.442090in}}%
\pgfpathlineto{\pgfqpoint{3.774396in}{1.412463in}}%
\pgfpathlineto{\pgfqpoint{3.774459in}{1.421732in}}%
\pgfpathlineto{\pgfqpoint{3.775516in}{1.361870in}}%
\pgfpathlineto{\pgfqpoint{3.775705in}{1.364737in}}%
\pgfpathlineto{\pgfqpoint{3.775813in}{1.383089in}}%
\pgfpathlineto{\pgfqpoint{3.776713in}{1.313609in}}%
\pgfpathlineto{\pgfqpoint{3.776742in}{1.313899in}}%
\pgfpathlineto{\pgfqpoint{3.777244in}{1.305333in}}%
\pgfpathlineto{\pgfqpoint{3.777419in}{1.332239in}}%
\pgfpathlineto{\pgfqpoint{3.777521in}{1.327579in}}%
\pgfpathlineto{\pgfqpoint{3.777862in}{1.340245in}}%
\pgfpathlineto{\pgfqpoint{3.778461in}{1.293396in}}%
\pgfpathlineto{\pgfqpoint{3.778519in}{1.299537in}}%
\pgfpathlineto{\pgfqpoint{3.778894in}{1.264401in}}%
\pgfpathlineto{\pgfqpoint{3.779620in}{1.301644in}}%
\pgfpathlineto{\pgfqpoint{3.779629in}{1.299911in}}%
\pgfpathlineto{\pgfqpoint{3.779956in}{1.308717in}}%
\pgfpathlineto{\pgfqpoint{3.780384in}{1.286930in}}%
\pgfpathlineto{\pgfqpoint{3.780701in}{1.290133in}}%
\pgfpathlineto{\pgfqpoint{3.781786in}{1.274163in}}%
\pgfpathlineto{\pgfqpoint{3.781577in}{1.294753in}}%
\pgfpathlineto{\pgfqpoint{3.781825in}{1.276273in}}%
\pgfpathlineto{\pgfqpoint{3.782799in}{1.297589in}}%
\pgfpathlineto{\pgfqpoint{3.782000in}{1.276227in}}%
\pgfpathlineto{\pgfqpoint{3.782945in}{1.287041in}}%
\pgfpathlineto{\pgfqpoint{3.782955in}{1.284149in}}%
\pgfpathlineto{\pgfqpoint{3.783456in}{1.327721in}}%
\pgfpathlineto{\pgfqpoint{3.783996in}{1.311818in}}%
\pgfpathlineto{\pgfqpoint{3.785160in}{1.327645in}}%
\pgfpathlineto{\pgfqpoint{3.784639in}{1.297599in}}%
\pgfpathlineto{\pgfqpoint{3.785189in}{1.324334in}}%
\pgfpathlineto{\pgfqpoint{3.785194in}{1.322800in}}%
\pgfpathlineto{\pgfqpoint{3.785554in}{1.349656in}}%
\pgfpathlineto{\pgfqpoint{3.786236in}{1.341550in}}%
\pgfpathlineto{\pgfqpoint{3.786265in}{1.346759in}}%
\pgfpathlineto{\pgfqpoint{3.786922in}{1.314412in}}%
\pgfpathlineto{\pgfqpoint{3.787317in}{1.321624in}}%
\pgfpathlineto{\pgfqpoint{3.787964in}{1.311016in}}%
\pgfpathlineto{\pgfqpoint{3.787453in}{1.336186in}}%
\pgfpathlineto{\pgfqpoint{3.788441in}{1.319293in}}%
\pgfpathlineto{\pgfqpoint{3.788597in}{1.344264in}}%
\pgfpathlineto{\pgfqpoint{3.789459in}{1.358308in}}%
\pgfpathlineto{\pgfqpoint{3.788850in}{1.329837in}}%
\pgfpathlineto{\pgfqpoint{3.789693in}{1.339802in}}%
\pgfpathlineto{\pgfqpoint{3.790856in}{1.302304in}}%
\pgfpathlineto{\pgfqpoint{3.789907in}{1.346555in}}%
\pgfpathlineto{\pgfqpoint{3.790905in}{1.307897in}}%
\pgfpathlineto{\pgfqpoint{3.791966in}{1.330608in}}%
\pgfpathlineto{\pgfqpoint{3.791319in}{1.290258in}}%
\pgfpathlineto{\pgfqpoint{3.792151in}{1.314244in}}%
\pgfpathlineto{\pgfqpoint{3.792808in}{1.292771in}}%
\pgfpathlineto{\pgfqpoint{3.793208in}{1.335254in}}%
\pgfpathlineto{\pgfqpoint{3.793222in}{1.331331in}}%
\pgfpathlineto{\pgfqpoint{3.793286in}{1.349700in}}%
\pgfpathlineto{\pgfqpoint{3.793957in}{1.310596in}}%
\pgfpathlineto{\pgfqpoint{3.794332in}{1.334466in}}%
\pgfpathlineto{\pgfqpoint{3.794337in}{1.330940in}}%
\pgfpathlineto{\pgfqpoint{3.795233in}{1.368798in}}%
\pgfpathlineto{\pgfqpoint{3.795403in}{1.359371in}}%
\pgfpathlineto{\pgfqpoint{3.795598in}{1.363946in}}%
\pgfpathlineto{\pgfqpoint{3.796372in}{1.326903in}}%
\pgfpathlineto{\pgfqpoint{3.796377in}{1.327279in}}%
\pgfpathlineto{\pgfqpoint{3.797112in}{1.346615in}}%
\pgfpathlineto{\pgfqpoint{3.797516in}{1.309909in}}%
\pgfpathlineto{\pgfqpoint{3.798086in}{1.361525in}}%
\pgfpathlineto{\pgfqpoint{3.798656in}{1.326092in}}%
\pgfpathlineto{\pgfqpoint{3.799371in}{1.308924in}}%
\pgfpathlineto{\pgfqpoint{3.798918in}{1.347552in}}%
\pgfpathlineto{\pgfqpoint{3.799741in}{1.338401in}}%
\pgfpathlineto{\pgfqpoint{3.800218in}{1.325210in}}%
\pgfpathlineto{\pgfqpoint{3.800087in}{1.345530in}}%
\pgfpathlineto{\pgfqpoint{3.800793in}{1.340711in}}%
\pgfpathlineto{\pgfqpoint{3.801285in}{1.352612in}}%
\pgfpathlineto{\pgfqpoint{3.801591in}{1.326118in}}%
\pgfpathlineto{\pgfqpoint{3.801635in}{1.329126in}}%
\pgfpathlineto{\pgfqpoint{3.801674in}{1.321545in}}%
\pgfpathlineto{\pgfqpoint{3.802682in}{1.362428in}}%
\pgfpathlineto{\pgfqpoint{3.802706in}{1.357988in}}%
\pgfpathlineto{\pgfqpoint{3.802808in}{1.356919in}}%
\pgfpathlineto{\pgfqpoint{3.803461in}{1.377536in}}%
\pgfpathlineto{\pgfqpoint{3.803578in}{1.378347in}}%
\pgfpathlineto{\pgfqpoint{3.804099in}{1.352286in}}%
\pgfpathlineto{\pgfqpoint{3.804420in}{1.356166in}}%
\pgfpathlineto{\pgfqpoint{3.805131in}{1.340926in}}%
\pgfpathlineto{\pgfqpoint{3.804649in}{1.365580in}}%
\pgfpathlineto{\pgfqpoint{3.805544in}{1.348809in}}%
\pgfpathlineto{\pgfqpoint{3.805783in}{1.359607in}}%
\pgfpathlineto{\pgfqpoint{3.806061in}{1.336745in}}%
\pgfpathlineto{\pgfqpoint{3.806635in}{1.344146in}}%
\pgfpathlineto{\pgfqpoint{3.806810in}{1.321466in}}%
\pgfpathlineto{\pgfqpoint{3.807696in}{1.357042in}}%
\pgfpathlineto{\pgfqpoint{3.807706in}{1.354814in}}%
\pgfpathlineto{\pgfqpoint{3.808217in}{1.309471in}}%
\pgfpathlineto{\pgfqpoint{3.808943in}{1.335334in}}%
\pgfpathlineto{\pgfqpoint{3.808996in}{1.344696in}}%
\pgfpathlineto{\pgfqpoint{3.809800in}{1.316216in}}%
\pgfpathlineto{\pgfqpoint{3.810014in}{1.316325in}}%
\pgfpathlineto{\pgfqpoint{3.811168in}{1.291150in}}%
\pgfpathlineto{\pgfqpoint{3.811479in}{1.308487in}}%
\pgfpathlineto{\pgfqpoint{3.811284in}{1.285740in}}%
\pgfpathlineto{\pgfqpoint{3.811523in}{1.307550in}}%
\pgfpathlineto{\pgfqpoint{3.812643in}{1.324392in}}%
\pgfpathlineto{\pgfqpoint{3.812097in}{1.288802in}}%
\pgfpathlineto{\pgfqpoint{3.812657in}{1.323871in}}%
\pgfpathlineto{\pgfqpoint{3.812954in}{1.305653in}}%
\pgfpathlineto{\pgfqpoint{3.812828in}{1.331058in}}%
\pgfpathlineto{\pgfqpoint{3.813753in}{1.328462in}}%
\pgfpathlineto{\pgfqpoint{3.814955in}{1.367088in}}%
\pgfpathlineto{\pgfqpoint{3.813806in}{1.327166in}}%
\pgfpathlineto{\pgfqpoint{3.814960in}{1.365463in}}%
\pgfpathlineto{\pgfqpoint{3.814965in}{1.362773in}}%
\pgfpathlineto{\pgfqpoint{3.815778in}{1.394799in}}%
\pgfpathlineto{\pgfqpoint{3.816060in}{1.369960in}}%
\pgfpathlineto{\pgfqpoint{3.817146in}{1.387532in}}%
\pgfpathlineto{\pgfqpoint{3.816484in}{1.357351in}}%
\pgfpathlineto{\pgfqpoint{3.817200in}{1.383006in}}%
\pgfpathlineto{\pgfqpoint{3.817326in}{1.366696in}}%
\pgfpathlineto{\pgfqpoint{3.818256in}{1.395571in}}%
\pgfpathlineto{\pgfqpoint{3.818276in}{1.391428in}}%
\pgfpathlineto{\pgfqpoint{3.818582in}{1.408437in}}%
\pgfpathlineto{\pgfqpoint{3.819196in}{1.363738in}}%
\pgfpathlineto{\pgfqpoint{3.819327in}{1.373782in}}%
\pgfpathlineto{\pgfqpoint{3.819975in}{1.337239in}}%
\pgfpathlineto{\pgfqpoint{3.819342in}{1.375446in}}%
\pgfpathlineto{\pgfqpoint{3.820462in}{1.370402in}}%
\pgfpathlineto{\pgfqpoint{3.820681in}{1.380621in}}%
\pgfpathlineto{\pgfqpoint{3.820978in}{1.355830in}}%
\pgfpathlineto{\pgfqpoint{3.821508in}{1.358489in}}%
\pgfpathlineto{\pgfqpoint{3.821669in}{1.355035in}}%
\pgfpathlineto{\pgfqpoint{3.822365in}{1.380982in}}%
\pgfpathlineto{\pgfqpoint{3.822380in}{1.379930in}}%
\pgfpathlineto{\pgfqpoint{3.822399in}{1.382523in}}%
\pgfpathlineto{\pgfqpoint{3.823280in}{1.354691in}}%
\pgfpathlineto{\pgfqpoint{3.823461in}{1.367871in}}%
\pgfpathlineto{\pgfqpoint{3.824629in}{1.355940in}}%
\pgfpathlineto{\pgfqpoint{3.823714in}{1.382210in}}%
\pgfpathlineto{\pgfqpoint{3.824644in}{1.361904in}}%
\pgfpathlineto{\pgfqpoint{3.825213in}{1.402149in}}%
\pgfpathlineto{\pgfqpoint{3.824673in}{1.358765in}}%
\pgfpathlineto{\pgfqpoint{3.825822in}{1.378422in}}%
\pgfpathlineto{\pgfqpoint{3.825948in}{1.369895in}}%
\pgfpathlineto{\pgfqpoint{3.826421in}{1.396667in}}%
\pgfpathlineto{\pgfqpoint{3.826917in}{1.385587in}}%
\pgfpathlineto{\pgfqpoint{3.826922in}{1.386353in}}%
\pgfpathlineto{\pgfqpoint{3.827404in}{1.359217in}}%
\pgfpathlineto{\pgfqpoint{3.827915in}{1.368573in}}%
\pgfpathlineto{\pgfqpoint{3.827954in}{1.359910in}}%
\pgfpathlineto{\pgfqpoint{3.828758in}{1.399645in}}%
\pgfpathlineto{\pgfqpoint{3.828801in}{1.398120in}}%
\pgfpathlineto{\pgfqpoint{3.828899in}{1.410114in}}%
\pgfpathlineto{\pgfqpoint{3.829756in}{1.381987in}}%
\pgfpathlineto{\pgfqpoint{3.829902in}{1.390666in}}%
\pgfpathlineto{\pgfqpoint{3.831133in}{1.347722in}}%
\pgfpathlineto{\pgfqpoint{3.829921in}{1.391735in}}%
\pgfpathlineto{\pgfqpoint{3.831430in}{1.380659in}}%
\pgfpathlineto{\pgfqpoint{3.831562in}{1.388221in}}%
\pgfpathlineto{\pgfqpoint{3.832112in}{1.356130in}}%
\pgfpathlineto{\pgfqpoint{3.832526in}{1.377060in}}%
\pgfpathlineto{\pgfqpoint{3.833066in}{1.365648in}}%
\pgfpathlineto{\pgfqpoint{3.833509in}{1.408079in}}%
\pgfpathlineto{\pgfqpoint{3.833514in}{1.406916in}}%
\pgfpathlineto{\pgfqpoint{3.833519in}{1.406939in}}%
\pgfpathlineto{\pgfqpoint{3.833524in}{1.406553in}}%
\pgfpathlineto{\pgfqpoint{3.833981in}{1.361902in}}%
\pgfpathlineto{\pgfqpoint{3.834653in}{1.377030in}}%
\pgfpathlineto{\pgfqpoint{3.835033in}{1.363653in}}%
\pgfpathlineto{\pgfqpoint{3.835403in}{1.409836in}}%
\pgfpathlineto{\pgfqpoint{3.836318in}{1.424289in}}%
\pgfpathlineto{\pgfqpoint{3.836016in}{1.393543in}}%
\pgfpathlineto{\pgfqpoint{3.836523in}{1.412190in}}%
\pgfpathlineto{\pgfqpoint{3.837375in}{1.388322in}}%
\pgfpathlineto{\pgfqpoint{3.836829in}{1.419103in}}%
\pgfpathlineto{\pgfqpoint{3.837633in}{1.409524in}}%
\pgfpathlineto{\pgfqpoint{3.838470in}{1.427104in}}%
\pgfpathlineto{\pgfqpoint{3.838013in}{1.396739in}}%
\pgfpathlineto{\pgfqpoint{3.838728in}{1.402246in}}%
\pgfpathlineto{\pgfqpoint{3.838869in}{1.392407in}}%
\pgfpathlineto{\pgfqpoint{3.839123in}{1.424797in}}%
\pgfpathlineto{\pgfqpoint{3.839833in}{1.404295in}}%
\pgfpathlineto{\pgfqpoint{3.840783in}{1.352886in}}%
\pgfpathlineto{\pgfqpoint{3.839955in}{1.409501in}}%
\pgfpathlineto{\pgfqpoint{3.841075in}{1.362775in}}%
\pgfpathlineto{\pgfqpoint{3.841255in}{1.374315in}}%
\pgfpathlineto{\pgfqpoint{3.842005in}{1.336773in}}%
\pgfpathlineto{\pgfqpoint{3.842180in}{1.360516in}}%
\pgfpathlineto{\pgfqpoint{3.842910in}{1.389108in}}%
\pgfpathlineto{\pgfqpoint{3.843309in}{1.365907in}}%
\pgfpathlineto{\pgfqpoint{3.844030in}{1.356435in}}%
\pgfpathlineto{\pgfqpoint{3.843519in}{1.384320in}}%
\pgfpathlineto{\pgfqpoint{3.844400in}{1.371648in}}%
\pgfpathlineto{\pgfqpoint{3.844551in}{1.384967in}}%
\pgfpathlineto{\pgfqpoint{3.845213in}{1.353094in}}%
\pgfpathlineto{\pgfqpoint{3.845491in}{1.365987in}}%
\pgfpathlineto{\pgfqpoint{3.846440in}{1.382265in}}%
\pgfpathlineto{\pgfqpoint{3.845841in}{1.350954in}}%
\pgfpathlineto{\pgfqpoint{3.846640in}{1.371095in}}%
\pgfpathlineto{\pgfqpoint{3.847102in}{1.353241in}}%
\pgfpathlineto{\pgfqpoint{3.847662in}{1.393195in}}%
\pgfpathlineto{\pgfqpoint{3.847720in}{1.386589in}}%
\pgfpathlineto{\pgfqpoint{3.848937in}{1.462444in}}%
\pgfpathlineto{\pgfqpoint{3.847769in}{1.383452in}}%
\pgfpathlineto{\pgfqpoint{3.848972in}{1.457473in}}%
\pgfpathlineto{\pgfqpoint{3.849108in}{1.431097in}}%
\pgfpathlineto{\pgfqpoint{3.850018in}{1.479055in}}%
\pgfpathlineto{\pgfqpoint{3.850038in}{1.474559in}}%
\pgfpathlineto{\pgfqpoint{3.850520in}{1.489649in}}%
\pgfpathlineto{\pgfqpoint{3.851070in}{1.465852in}}%
\pgfpathlineto{\pgfqpoint{3.851177in}{1.486779in}}%
\pgfpathlineto{\pgfqpoint{3.851829in}{1.440163in}}%
\pgfpathlineto{\pgfqpoint{3.852316in}{1.458187in}}%
\pgfpathlineto{\pgfqpoint{3.852341in}{1.455725in}}%
\pgfpathlineto{\pgfqpoint{3.853056in}{1.492784in}}%
\pgfpathlineto{\pgfqpoint{3.853066in}{1.491434in}}%
\pgfpathlineto{\pgfqpoint{3.853256in}{1.496422in}}%
\pgfpathlineto{\pgfqpoint{3.853908in}{1.454455in}}%
\pgfpathlineto{\pgfqpoint{3.854113in}{1.462310in}}%
\pgfpathlineto{\pgfqpoint{3.854132in}{1.457626in}}%
\pgfpathlineto{\pgfqpoint{3.855106in}{1.494170in}}%
\pgfpathlineto{\pgfqpoint{3.855145in}{1.491048in}}%
\pgfpathlineto{\pgfqpoint{3.855758in}{1.504760in}}%
\pgfpathlineto{\pgfqpoint{3.855490in}{1.482505in}}%
\pgfpathlineto{\pgfqpoint{3.856245in}{1.485114in}}%
\pgfpathlineto{\pgfqpoint{3.856289in}{1.479514in}}%
\pgfpathlineto{\pgfqpoint{3.857243in}{1.561260in}}%
\pgfpathlineto{\pgfqpoint{3.858937in}{1.488406in}}%
\pgfpathlineto{\pgfqpoint{3.857477in}{1.576396in}}%
\pgfpathlineto{\pgfqpoint{3.859317in}{1.522875in}}%
\pgfpathlineto{\pgfqpoint{3.860198in}{1.549626in}}%
\pgfpathlineto{\pgfqpoint{3.859522in}{1.518606in}}%
\pgfpathlineto{\pgfqpoint{3.860466in}{1.542554in}}%
\pgfpathlineto{\pgfqpoint{3.860471in}{1.539057in}}%
\pgfpathlineto{\pgfqpoint{3.861333in}{1.581164in}}%
\pgfpathlineto{\pgfqpoint{3.861532in}{1.574277in}}%
\pgfpathlineto{\pgfqpoint{3.862272in}{1.611978in}}%
\pgfpathlineto{\pgfqpoint{3.862540in}{1.572043in}}%
\pgfpathlineto{\pgfqpoint{3.862740in}{1.589464in}}%
\pgfpathlineto{\pgfqpoint{3.863509in}{1.572126in}}%
\pgfpathlineto{\pgfqpoint{3.863635in}{1.599032in}}%
\pgfpathlineto{\pgfqpoint{3.863708in}{1.597925in}}%
\pgfpathlineto{\pgfqpoint{3.863728in}{1.607811in}}%
\pgfpathlineto{\pgfqpoint{3.864648in}{1.549450in}}%
\pgfpathlineto{\pgfqpoint{3.864755in}{1.552766in}}%
\pgfpathlineto{\pgfqpoint{3.865753in}{1.529408in}}%
\pgfpathlineto{\pgfqpoint{3.865203in}{1.574755in}}%
\pgfpathlineto{\pgfqpoint{3.865909in}{1.535523in}}%
\pgfpathlineto{\pgfqpoint{3.865919in}{1.538687in}}%
\pgfpathlineto{\pgfqpoint{3.866669in}{1.444374in}}%
\pgfpathlineto{\pgfqpoint{3.866844in}{1.448436in}}%
\pgfpathlineto{\pgfqpoint{3.867462in}{1.429684in}}%
\pgfpathlineto{\pgfqpoint{3.866922in}{1.457687in}}%
\pgfpathlineto{\pgfqpoint{3.867944in}{1.449753in}}%
\pgfpathlineto{\pgfqpoint{3.868962in}{1.464618in}}%
\pgfpathlineto{\pgfqpoint{3.868431in}{1.431467in}}%
\pgfpathlineto{\pgfqpoint{3.869083in}{1.455966in}}%
\pgfpathlineto{\pgfqpoint{3.870369in}{1.415681in}}%
\pgfpathlineto{\pgfqpoint{3.869195in}{1.467041in}}%
\pgfpathlineto{\pgfqpoint{3.870388in}{1.419593in}}%
\pgfpathlineto{\pgfqpoint{3.871396in}{1.473830in}}%
\pgfpathlineto{\pgfqpoint{3.871595in}{1.473630in}}%
\pgfpathlineto{\pgfqpoint{3.871980in}{1.462705in}}%
\pgfpathlineto{\pgfqpoint{3.872618in}{1.498517in}}%
\pgfpathlineto{\pgfqpoint{3.872623in}{1.498239in}}%
\pgfpathlineto{\pgfqpoint{3.872725in}{1.501652in}}%
\pgfpathlineto{\pgfqpoint{3.873377in}{1.468337in}}%
\pgfpathlineto{\pgfqpoint{3.873557in}{1.472267in}}%
\pgfpathlineto{\pgfqpoint{3.873596in}{1.462291in}}%
\pgfpathlineto{\pgfqpoint{3.874224in}{1.494811in}}%
\pgfpathlineto{\pgfqpoint{3.874629in}{1.480382in}}%
\pgfpathlineto{\pgfqpoint{3.875768in}{1.528037in}}%
\pgfpathlineto{\pgfqpoint{3.874716in}{1.475420in}}%
\pgfpathlineto{\pgfqpoint{3.875899in}{1.523164in}}%
\pgfpathlineto{\pgfqpoint{3.875953in}{1.507179in}}%
\pgfpathlineto{\pgfqpoint{3.876693in}{1.542012in}}%
\pgfpathlineto{\pgfqpoint{3.876985in}{1.531567in}}%
\pgfpathlineto{\pgfqpoint{3.877111in}{1.549387in}}%
\pgfpathlineto{\pgfqpoint{3.877705in}{1.519914in}}%
\pgfpathlineto{\pgfqpoint{3.878051in}{1.522150in}}%
\pgfpathlineto{\pgfqpoint{3.879336in}{1.503886in}}%
\pgfpathlineto{\pgfqpoint{3.878645in}{1.542289in}}%
\pgfpathlineto{\pgfqpoint{3.879375in}{1.505984in}}%
\pgfpathlineto{\pgfqpoint{3.880563in}{1.540351in}}%
\pgfpathlineto{\pgfqpoint{3.879599in}{1.503554in}}%
\pgfpathlineto{\pgfqpoint{3.880573in}{1.535137in}}%
\pgfpathlineto{\pgfqpoint{3.881401in}{1.504311in}}%
\pgfpathlineto{\pgfqpoint{3.880836in}{1.538169in}}%
\pgfpathlineto{\pgfqpoint{3.881717in}{1.516698in}}%
\pgfpathlineto{\pgfqpoint{3.881917in}{1.533580in}}%
\pgfpathlineto{\pgfqpoint{3.882306in}{1.506611in}}%
\pgfpathlineto{\pgfqpoint{3.882832in}{1.518913in}}%
\pgfpathlineto{\pgfqpoint{3.884901in}{1.416954in}}%
\pgfpathlineto{\pgfqpoint{3.882871in}{1.521874in}}%
\pgfpathlineto{\pgfqpoint{3.885558in}{1.448576in}}%
\pgfpathlineto{\pgfqpoint{3.886006in}{1.472080in}}%
\pgfpathlineto{\pgfqpoint{3.886493in}{1.446379in}}%
\pgfpathlineto{\pgfqpoint{3.886702in}{1.466543in}}%
\pgfpathlineto{\pgfqpoint{3.886717in}{1.460588in}}%
\pgfpathlineto{\pgfqpoint{3.887433in}{1.484989in}}%
\pgfpathlineto{\pgfqpoint{3.887803in}{1.471246in}}%
\pgfpathlineto{\pgfqpoint{3.889044in}{1.402324in}}%
\pgfpathlineto{\pgfqpoint{3.887973in}{1.474223in}}%
\pgfpathlineto{\pgfqpoint{3.889161in}{1.414132in}}%
\pgfpathlineto{\pgfqpoint{3.890096in}{1.446635in}}%
\pgfpathlineto{\pgfqpoint{3.889200in}{1.409057in}}%
\pgfpathlineto{\pgfqpoint{3.890549in}{1.422331in}}%
\pgfpathlineto{\pgfqpoint{3.891610in}{1.390821in}}%
\pgfpathlineto{\pgfqpoint{3.891371in}{1.431748in}}%
\pgfpathlineto{\pgfqpoint{3.891707in}{1.392745in}}%
\pgfpathlineto{\pgfqpoint{3.892150in}{1.416508in}}%
\pgfpathlineto{\pgfqpoint{3.892486in}{1.391534in}}%
\pgfpathlineto{\pgfqpoint{3.892851in}{1.408054in}}%
\pgfpathlineto{\pgfqpoint{3.892983in}{1.396391in}}%
\pgfpathlineto{\pgfqpoint{3.893416in}{1.423735in}}%
\pgfpathlineto{\pgfqpoint{3.893742in}{1.420900in}}%
\pgfpathlineto{\pgfqpoint{3.894351in}{1.445352in}}%
\pgfpathlineto{\pgfqpoint{3.894799in}{1.398144in}}%
\pgfpathlineto{\pgfqpoint{3.894813in}{1.400784in}}%
\pgfpathlineto{\pgfqpoint{3.895154in}{1.382073in}}%
\pgfpathlineto{\pgfqpoint{3.895573in}{1.410402in}}%
\pgfpathlineto{\pgfqpoint{3.896376in}{1.434266in}}%
\pgfpathlineto{\pgfqpoint{3.895656in}{1.405759in}}%
\pgfpathlineto{\pgfqpoint{3.896707in}{1.423166in}}%
\pgfpathlineto{\pgfqpoint{3.897622in}{1.401761in}}%
\pgfpathlineto{\pgfqpoint{3.897228in}{1.442478in}}%
\pgfpathlineto{\pgfqpoint{3.897842in}{1.411249in}}%
\pgfpathlineto{\pgfqpoint{3.898893in}{1.452408in}}%
\pgfpathlineto{\pgfqpoint{3.898085in}{1.402934in}}%
\pgfpathlineto{\pgfqpoint{3.898995in}{1.451093in}}%
\pgfpathlineto{\pgfqpoint{3.899253in}{1.425870in}}%
\pgfpathlineto{\pgfqpoint{3.899993in}{1.467568in}}%
\pgfpathlineto{\pgfqpoint{3.900076in}{1.463680in}}%
\pgfpathlineto{\pgfqpoint{3.900933in}{1.505109in}}%
\pgfpathlineto{\pgfqpoint{3.900232in}{1.453891in}}%
\pgfpathlineto{\pgfqpoint{3.901381in}{1.480577in}}%
\pgfpathlineto{\pgfqpoint{3.901435in}{1.472495in}}%
\pgfpathlineto{\pgfqpoint{3.902345in}{1.559630in}}%
\pgfpathlineto{\pgfqpoint{3.902423in}{1.549251in}}%
\pgfpathlineto{\pgfqpoint{3.903148in}{1.574512in}}%
\pgfpathlineto{\pgfqpoint{3.902510in}{1.543296in}}%
\pgfpathlineto{\pgfqpoint{3.903557in}{1.563976in}}%
\pgfpathlineto{\pgfqpoint{3.904171in}{1.543833in}}%
\pgfpathlineto{\pgfqpoint{3.903791in}{1.576589in}}%
\pgfpathlineto{\pgfqpoint{3.904721in}{1.554700in}}%
\pgfpathlineto{\pgfqpoint{3.905100in}{1.574197in}}%
\pgfpathlineto{\pgfqpoint{3.905733in}{1.547441in}}%
\pgfpathlineto{\pgfqpoint{3.905875in}{1.538194in}}%
\pgfpathlineto{\pgfqpoint{3.906804in}{1.558125in}}%
\pgfpathlineto{\pgfqpoint{3.906819in}{1.556673in}}%
\pgfpathlineto{\pgfqpoint{3.907335in}{1.574950in}}%
\pgfpathlineto{\pgfqpoint{3.907827in}{1.544162in}}%
\pgfpathlineto{\pgfqpoint{3.907837in}{1.546990in}}%
\pgfpathlineto{\pgfqpoint{3.908109in}{1.530865in}}%
\pgfpathlineto{\pgfqpoint{3.908221in}{1.551492in}}%
\pgfpathlineto{\pgfqpoint{3.908951in}{1.544043in}}%
\pgfpathlineto{\pgfqpoint{3.909200in}{1.566986in}}%
\pgfpathlineto{\pgfqpoint{3.909857in}{1.541086in}}%
\pgfpathlineto{\pgfqpoint{3.910120in}{1.552832in}}%
\pgfpathlineto{\pgfqpoint{3.911060in}{1.527677in}}%
\pgfpathlineto{\pgfqpoint{3.910685in}{1.554692in}}%
\pgfpathlineto{\pgfqpoint{3.911444in}{1.539442in}}%
\pgfpathlineto{\pgfqpoint{3.912019in}{1.570049in}}%
\pgfpathlineto{\pgfqpoint{3.912505in}{1.532575in}}%
\pgfpathlineto{\pgfqpoint{3.912544in}{1.534461in}}%
\pgfpathlineto{\pgfqpoint{3.913800in}{1.503920in}}%
\pgfpathlineto{\pgfqpoint{3.912788in}{1.542092in}}%
\pgfpathlineto{\pgfqpoint{3.913820in}{1.506260in}}%
\pgfpathlineto{\pgfqpoint{3.914896in}{1.519601in}}%
\pgfpathlineto{\pgfqpoint{3.914020in}{1.495613in}}%
\pgfpathlineto{\pgfqpoint{3.914954in}{1.516170in}}%
\pgfpathlineto{\pgfqpoint{3.915972in}{1.497773in}}%
\pgfpathlineto{\pgfqpoint{3.915592in}{1.527179in}}%
\pgfpathlineto{\pgfqpoint{3.916098in}{1.506222in}}%
\pgfpathlineto{\pgfqpoint{3.917145in}{1.539053in}}%
\pgfpathlineto{\pgfqpoint{3.916201in}{1.504365in}}%
\pgfpathlineto{\pgfqpoint{3.917277in}{1.536751in}}%
\pgfpathlineto{\pgfqpoint{3.918270in}{1.475486in}}%
\pgfpathlineto{\pgfqpoint{3.917291in}{1.537887in}}%
\pgfpathlineto{\pgfqpoint{3.918465in}{1.492281in}}%
\pgfpathlineto{\pgfqpoint{3.918869in}{1.502545in}}%
\pgfpathlineto{\pgfqpoint{3.918766in}{1.481702in}}%
\pgfpathlineto{\pgfqpoint{3.919560in}{1.483420in}}%
\pgfpathlineto{\pgfqpoint{3.920490in}{1.517834in}}%
\pgfpathlineto{\pgfqpoint{3.919618in}{1.481573in}}%
\pgfpathlineto{\pgfqpoint{3.920670in}{1.484814in}}%
\pgfpathlineto{\pgfqpoint{3.921868in}{1.449795in}}%
\pgfpathlineto{\pgfqpoint{3.920685in}{1.486255in}}%
\pgfpathlineto{\pgfqpoint{3.921911in}{1.461606in}}%
\pgfpathlineto{\pgfqpoint{3.922929in}{1.497893in}}%
\pgfpathlineto{\pgfqpoint{3.923060in}{1.494934in}}%
\pgfpathlineto{\pgfqpoint{3.924239in}{1.443746in}}%
\pgfpathlineto{\pgfqpoint{3.923275in}{1.502690in}}%
\pgfpathlineto{\pgfqpoint{3.924258in}{1.447539in}}%
\pgfpathlineto{\pgfqpoint{3.924443in}{1.458352in}}%
\pgfpathlineto{\pgfqpoint{3.925261in}{1.402406in}}%
\pgfpathlineto{\pgfqpoint{3.925271in}{1.402642in}}%
\pgfpathlineto{\pgfqpoint{3.926069in}{1.432990in}}%
\pgfpathlineto{\pgfqpoint{3.925407in}{1.395996in}}%
\pgfpathlineto{\pgfqpoint{3.926498in}{1.424603in}}%
\pgfpathlineto{\pgfqpoint{3.926507in}{1.424948in}}%
\pgfpathlineto{\pgfqpoint{3.927277in}{1.433200in}}%
\pgfpathlineto{\pgfqpoint{3.926639in}{1.412354in}}%
\pgfpathlineto{\pgfqpoint{3.927622in}{1.426337in}}%
\pgfpathlineto{\pgfqpoint{3.927836in}{1.436411in}}%
\pgfpathlineto{\pgfqpoint{3.928158in}{1.408357in}}%
\pgfpathlineto{\pgfqpoint{3.929803in}{1.353560in}}%
\pgfpathlineto{\pgfqpoint{3.928270in}{1.410577in}}%
\pgfpathlineto{\pgfqpoint{3.929808in}{1.354185in}}%
\pgfpathlineto{\pgfqpoint{3.930008in}{1.369861in}}%
\pgfpathlineto{\pgfqpoint{3.930655in}{1.336259in}}%
\pgfpathlineto{\pgfqpoint{3.930860in}{1.338525in}}%
\pgfpathlineto{\pgfqpoint{3.930865in}{1.336967in}}%
\pgfpathlineto{\pgfqpoint{3.931191in}{1.378966in}}%
\pgfpathlineto{\pgfqpoint{3.931887in}{1.365197in}}%
\pgfpathlineto{\pgfqpoint{3.932427in}{1.357309in}}%
\pgfpathlineto{\pgfqpoint{3.932646in}{1.381171in}}%
\pgfpathlineto{\pgfqpoint{3.932812in}{1.378172in}}%
\pgfpathlineto{\pgfqpoint{3.932992in}{1.367690in}}%
\pgfpathlineto{\pgfqpoint{3.933961in}{1.413208in}}%
\pgfpathlineto{\pgfqpoint{3.934311in}{1.408213in}}%
\pgfpathlineto{\pgfqpoint{3.934803in}{1.452901in}}%
\pgfpathlineto{\pgfqpoint{3.934901in}{1.448564in}}%
\pgfpathlineto{\pgfqpoint{3.935792in}{1.471343in}}%
\pgfpathlineto{\pgfqpoint{3.935480in}{1.433076in}}%
\pgfpathlineto{\pgfqpoint{3.935996in}{1.444931in}}%
\pgfpathlineto{\pgfqpoint{3.936074in}{1.438032in}}%
\pgfpathlineto{\pgfqpoint{3.936911in}{1.508886in}}%
\pgfpathlineto{\pgfqpoint{3.937004in}{1.503698in}}%
\pgfpathlineto{\pgfqpoint{3.937208in}{1.526892in}}%
\pgfpathlineto{\pgfqpoint{3.938065in}{1.493196in}}%
\pgfpathlineto{\pgfqpoint{3.938094in}{1.498703in}}%
\pgfpathlineto{\pgfqpoint{3.939277in}{1.431007in}}%
\pgfpathlineto{\pgfqpoint{3.938133in}{1.504832in}}%
\pgfpathlineto{\pgfqpoint{3.939326in}{1.441967in}}%
\pgfpathlineto{\pgfqpoint{3.939331in}{1.442554in}}%
\pgfpathlineto{\pgfqpoint{3.939706in}{1.416685in}}%
\pgfpathlineto{\pgfqpoint{3.940222in}{1.417479in}}%
\pgfpathlineto{\pgfqpoint{3.940436in}{1.396906in}}%
\pgfpathlineto{\pgfqpoint{3.940246in}{1.419706in}}%
\pgfpathlineto{\pgfqpoint{3.941332in}{1.417382in}}%
\pgfpathlineto{\pgfqpoint{3.942257in}{1.447955in}}%
\pgfpathlineto{\pgfqpoint{3.941863in}{1.409319in}}%
\pgfpathlineto{\pgfqpoint{3.942627in}{1.439251in}}%
\pgfpathlineto{\pgfqpoint{3.943050in}{1.402532in}}%
\pgfpathlineto{\pgfqpoint{3.943761in}{1.419862in}}%
\pgfpathlineto{\pgfqpoint{3.944136in}{1.425578in}}%
\pgfpathlineto{\pgfqpoint{3.944448in}{1.396243in}}%
\pgfpathlineto{\pgfqpoint{3.944789in}{1.400826in}}%
\pgfpathlineto{\pgfqpoint{3.944793in}{1.397732in}}%
\pgfpathlineto{\pgfqpoint{3.945738in}{1.447553in}}%
\pgfpathlineto{\pgfqpoint{3.945835in}{1.441527in}}%
\pgfpathlineto{\pgfqpoint{3.946595in}{1.413031in}}%
\pgfpathlineto{\pgfqpoint{3.946960in}{1.434523in}}%
\pgfpathlineto{\pgfqpoint{3.947603in}{1.455273in}}%
\pgfpathlineto{\pgfqpoint{3.947987in}{1.430833in}}%
\pgfpathlineto{\pgfqpoint{3.948075in}{1.437166in}}%
\pgfpathlineto{\pgfqpoint{3.948080in}{1.434886in}}%
\pgfpathlineto{\pgfqpoint{3.948420in}{1.463691in}}%
\pgfpathlineto{\pgfqpoint{3.949156in}{1.449231in}}%
\pgfpathlineto{\pgfqpoint{3.949793in}{1.474944in}}%
\pgfpathlineto{\pgfqpoint{3.949214in}{1.446665in}}%
\pgfpathlineto{\pgfqpoint{3.950300in}{1.460114in}}%
\pgfpathlineto{\pgfqpoint{3.950446in}{1.447692in}}%
\pgfpathlineto{\pgfqpoint{3.951303in}{1.481473in}}%
\pgfpathlineto{\pgfqpoint{3.951312in}{1.481318in}}%
\pgfpathlineto{\pgfqpoint{3.951682in}{1.502411in}}%
\pgfpathlineto{\pgfqpoint{3.952052in}{1.457983in}}%
\pgfpathlineto{\pgfqpoint{3.952086in}{1.458801in}}%
\pgfpathlineto{\pgfqpoint{3.952184in}{1.448773in}}%
\pgfpathlineto{\pgfqpoint{3.952593in}{1.489584in}}%
\pgfpathlineto{\pgfqpoint{3.953167in}{1.468837in}}%
\pgfpathlineto{\pgfqpoint{3.953323in}{1.477295in}}%
\pgfpathlineto{\pgfqpoint{3.953951in}{1.430183in}}%
\pgfpathlineto{\pgfqpoint{3.953956in}{1.432036in}}%
\pgfpathlineto{\pgfqpoint{3.954978in}{1.379884in}}%
\pgfpathlineto{\pgfqpoint{3.954068in}{1.444943in}}%
\pgfpathlineto{\pgfqpoint{3.955212in}{1.390973in}}%
\pgfpathlineto{\pgfqpoint{3.956220in}{1.349490in}}%
\pgfpathlineto{\pgfqpoint{3.956473in}{1.354361in}}%
\pgfpathlineto{\pgfqpoint{3.956648in}{1.363367in}}%
\pgfpathlineto{\pgfqpoint{3.957252in}{1.339109in}}%
\pgfpathlineto{\pgfqpoint{3.957559in}{1.348240in}}%
\pgfpathlineto{\pgfqpoint{3.958318in}{1.334644in}}%
\pgfpathlineto{\pgfqpoint{3.958050in}{1.355054in}}%
\pgfpathlineto{\pgfqpoint{3.958698in}{1.335911in}}%
\pgfpathlineto{\pgfqpoint{3.959784in}{1.360549in}}%
\pgfpathlineto{\pgfqpoint{3.959248in}{1.325722in}}%
\pgfpathlineto{\pgfqpoint{3.959837in}{1.360184in}}%
\pgfpathlineto{\pgfqpoint{3.961122in}{1.288877in}}%
\pgfpathlineto{\pgfqpoint{3.959949in}{1.360613in}}%
\pgfpathlineto{\pgfqpoint{3.961132in}{1.291502in}}%
\pgfpathlineto{\pgfqpoint{3.962266in}{1.320066in}}%
\pgfpathlineto{\pgfqpoint{3.961682in}{1.278820in}}%
\pgfpathlineto{\pgfqpoint{3.962296in}{1.314425in}}%
\pgfpathlineto{\pgfqpoint{3.962354in}{1.308153in}}%
\pgfpathlineto{\pgfqpoint{3.963050in}{1.340243in}}%
\pgfpathlineto{\pgfqpoint{3.963333in}{1.327141in}}%
\pgfpathlineto{\pgfqpoint{3.964331in}{1.360397in}}%
\pgfpathlineto{\pgfqpoint{3.963790in}{1.321675in}}%
\pgfpathlineto{\pgfqpoint{3.964506in}{1.349239in}}%
\pgfpathlineto{\pgfqpoint{3.965767in}{1.279143in}}%
\pgfpathlineto{\pgfqpoint{3.965845in}{1.286408in}}%
\pgfpathlineto{\pgfqpoint{3.965850in}{1.288490in}}%
\pgfpathlineto{\pgfqpoint{3.966108in}{1.266835in}}%
\pgfpathlineto{\pgfqpoint{3.966930in}{1.276345in}}%
\pgfpathlineto{\pgfqpoint{3.970372in}{1.376294in}}%
\pgfpathlineto{\pgfqpoint{3.966950in}{1.274229in}}%
\pgfpathlineto{\pgfqpoint{3.970626in}{1.350597in}}%
\pgfpathlineto{\pgfqpoint{3.971677in}{1.320372in}}%
\pgfpathlineto{\pgfqpoint{3.971103in}{1.359865in}}%
\pgfpathlineto{\pgfqpoint{3.971814in}{1.335759in}}%
\pgfpathlineto{\pgfqpoint{3.973245in}{1.383136in}}%
\pgfpathlineto{\pgfqpoint{3.972115in}{1.326604in}}%
\pgfpathlineto{\pgfqpoint{3.973298in}{1.367418in}}%
\pgfpathlineto{\pgfqpoint{3.974150in}{1.342859in}}%
\pgfpathlineto{\pgfqpoint{3.973751in}{1.371620in}}%
\pgfpathlineto{\pgfqpoint{3.974447in}{1.353766in}}%
\pgfpathlineto{\pgfqpoint{3.974467in}{1.350713in}}%
\pgfpathlineto{\pgfqpoint{3.974642in}{1.366022in}}%
\pgfpathlineto{\pgfqpoint{3.974681in}{1.365958in}}%
\pgfpathlineto{\pgfqpoint{3.974842in}{1.365540in}}%
\pgfpathlineto{\pgfqpoint{3.975820in}{1.392149in}}%
\pgfpathlineto{\pgfqpoint{3.975913in}{1.372409in}}%
\pgfpathlineto{\pgfqpoint{3.976332in}{1.416635in}}%
\pgfpathlineto{\pgfqpoint{3.976867in}{1.412105in}}%
\pgfpathlineto{\pgfqpoint{3.977962in}{1.435898in}}%
\pgfpathlineto{\pgfqpoint{3.977471in}{1.404034in}}%
\pgfpathlineto{\pgfqpoint{3.978021in}{1.429094in}}%
\pgfpathlineto{\pgfqpoint{3.979087in}{1.385138in}}%
\pgfpathlineto{\pgfqpoint{3.978216in}{1.440717in}}%
\pgfpathlineto{\pgfqpoint{3.979253in}{1.398590in}}%
\pgfpathlineto{\pgfqpoint{3.980533in}{1.470295in}}%
\pgfpathlineto{\pgfqpoint{3.980577in}{1.461503in}}%
\pgfpathlineto{\pgfqpoint{3.981843in}{1.393924in}}%
\pgfpathlineto{\pgfqpoint{3.980655in}{1.463187in}}%
\pgfpathlineto{\pgfqpoint{3.981857in}{1.394591in}}%
\pgfpathlineto{\pgfqpoint{3.981862in}{1.394804in}}%
\pgfpathlineto{\pgfqpoint{3.982344in}{1.368371in}}%
\pgfpathlineto{\pgfqpoint{3.982349in}{1.366064in}}%
\pgfpathlineto{\pgfqpoint{3.983259in}{1.407575in}}%
\pgfpathlineto{\pgfqpoint{3.983415in}{1.385302in}}%
\pgfpathlineto{\pgfqpoint{3.984637in}{1.447286in}}%
\pgfpathlineto{\pgfqpoint{3.983459in}{1.382260in}}%
\pgfpathlineto{\pgfqpoint{3.984764in}{1.437968in}}%
\pgfpathlineto{\pgfqpoint{3.985562in}{1.427031in}}%
\pgfpathlineto{\pgfqpoint{3.985217in}{1.450176in}}%
\pgfpathlineto{\pgfqpoint{3.985815in}{1.448226in}}%
\pgfpathlineto{\pgfqpoint{3.986916in}{1.461895in}}%
\pgfpathlineto{\pgfqpoint{3.986492in}{1.431308in}}%
\pgfpathlineto{\pgfqpoint{3.986945in}{1.455801in}}%
\pgfpathlineto{\pgfqpoint{3.987149in}{1.451431in}}%
\pgfpathlineto{\pgfqpoint{3.987636in}{1.479982in}}%
\pgfpathlineto{\pgfqpoint{3.988031in}{1.463056in}}%
\pgfpathlineto{\pgfqpoint{3.988435in}{1.483032in}}%
\pgfpathlineto{\pgfqpoint{3.988269in}{1.456940in}}%
\pgfpathlineto{\pgfqpoint{3.989106in}{1.457195in}}%
\pgfpathlineto{\pgfqpoint{3.989506in}{1.437799in}}%
\pgfpathlineto{\pgfqpoint{3.990046in}{1.477439in}}%
\pgfpathlineto{\pgfqpoint{3.990187in}{1.468862in}}%
\pgfpathlineto{\pgfqpoint{3.991497in}{1.393238in}}%
\pgfpathlineto{\pgfqpoint{3.991731in}{1.409543in}}%
\pgfpathlineto{\pgfqpoint{3.992811in}{1.434659in}}%
\pgfpathlineto{\pgfqpoint{3.992310in}{1.402624in}}%
\pgfpathlineto{\pgfqpoint{3.992880in}{1.425746in}}%
\pgfpathlineto{\pgfqpoint{3.993975in}{1.416154in}}%
\pgfpathlineto{\pgfqpoint{3.993235in}{1.444569in}}%
\pgfpathlineto{\pgfqpoint{3.993999in}{1.418700in}}%
\pgfpathlineto{\pgfqpoint{3.994072in}{1.433684in}}%
\pgfpathlineto{\pgfqpoint{3.995017in}{1.373042in}}%
\pgfpathlineto{\pgfqpoint{3.995333in}{1.347794in}}%
\pgfpathlineto{\pgfqpoint{3.995178in}{1.384813in}}%
\pgfpathlineto{\pgfqpoint{3.995338in}{1.349275in}}%
\pgfpathlineto{\pgfqpoint{3.995416in}{1.341597in}}%
\pgfpathlineto{\pgfqpoint{3.996278in}{1.390938in}}%
\pgfpathlineto{\pgfqpoint{3.996395in}{1.381585in}}%
\pgfpathlineto{\pgfqpoint{3.996400in}{1.381008in}}%
\pgfpathlineto{\pgfqpoint{3.997037in}{1.417660in}}%
\pgfpathlineto{\pgfqpoint{3.997213in}{1.415716in}}%
\pgfpathlineto{\pgfqpoint{3.997782in}{1.428485in}}%
\pgfpathlineto{\pgfqpoint{3.998259in}{1.398632in}}%
\pgfpathlineto{\pgfqpoint{3.998308in}{1.408473in}}%
\pgfpathlineto{\pgfqpoint{3.999145in}{1.386859in}}%
\pgfpathlineto{\pgfqpoint{3.998483in}{1.424453in}}%
\pgfpathlineto{\pgfqpoint{3.999501in}{1.402774in}}%
\pgfpathlineto{\pgfqpoint{4.000114in}{1.422387in}}%
\pgfpathlineto{\pgfqpoint{3.999574in}{1.391408in}}%
\pgfpathlineto{\pgfqpoint{4.000596in}{1.395615in}}%
\pgfpathlineto{\pgfqpoint{4.001721in}{1.372576in}}%
\pgfpathlineto{\pgfqpoint{4.001399in}{1.413926in}}%
\pgfpathlineto{\pgfqpoint{4.001731in}{1.374808in}}%
\pgfpathlineto{\pgfqpoint{4.002081in}{1.403488in}}%
\pgfpathlineto{\pgfqpoint{4.002860in}{1.388500in}}%
\pgfpathlineto{\pgfqpoint{4.003683in}{1.376054in}}%
\pgfpathlineto{\pgfqpoint{4.003405in}{1.406331in}}%
\pgfpathlineto{\pgfqpoint{4.003970in}{1.386226in}}%
\pgfpathlineto{\pgfqpoint{4.003975in}{1.387742in}}%
\pgfpathlineto{\pgfqpoint{4.004350in}{1.356258in}}%
\pgfpathlineto{\pgfqpoint{4.004978in}{1.358007in}}%
\pgfpathlineto{\pgfqpoint{4.006098in}{1.310493in}}%
\pgfpathlineto{\pgfqpoint{4.006132in}{1.314007in}}%
\pgfpathlineto{\pgfqpoint{4.007251in}{1.322353in}}%
\pgfpathlineto{\pgfqpoint{4.006803in}{1.296591in}}%
\pgfpathlineto{\pgfqpoint{4.007256in}{1.320787in}}%
\pgfpathlineto{\pgfqpoint{4.008858in}{1.358839in}}%
\pgfpathlineto{\pgfqpoint{4.007573in}{1.297577in}}%
\pgfpathlineto{\pgfqpoint{4.009082in}{1.349123in}}%
\pgfpathlineto{\pgfqpoint{4.009476in}{1.327401in}}%
\pgfpathlineto{\pgfqpoint{4.010177in}{1.350936in}}%
\pgfpathlineto{\pgfqpoint{4.010187in}{1.349487in}}%
\pgfpathlineto{\pgfqpoint{4.010216in}{1.356164in}}%
\pgfpathlineto{\pgfqpoint{4.010825in}{1.329261in}}%
\pgfpathlineto{\pgfqpoint{4.011268in}{1.346133in}}%
\pgfpathlineto{\pgfqpoint{4.012135in}{1.293325in}}%
\pgfpathlineto{\pgfqpoint{4.011278in}{1.346517in}}%
\pgfpathlineto{\pgfqpoint{4.012529in}{1.302397in}}%
\pgfpathlineto{\pgfqpoint{4.012933in}{1.316875in}}%
\pgfpathlineto{\pgfqpoint{4.013503in}{1.277628in}}%
\pgfpathlineto{\pgfqpoint{4.013921in}{1.242583in}}%
\pgfpathlineto{\pgfqpoint{4.013712in}{1.278524in}}%
\pgfpathlineto{\pgfqpoint{4.014637in}{1.258988in}}%
\pgfpathlineto{\pgfqpoint{4.015810in}{1.296768in}}%
\pgfpathlineto{\pgfqpoint{4.014729in}{1.253399in}}%
\pgfpathlineto{\pgfqpoint{4.015830in}{1.293579in}}%
\pgfpathlineto{\pgfqpoint{4.016015in}{1.275415in}}%
\pgfpathlineto{\pgfqpoint{4.016837in}{1.315376in}}%
\pgfpathlineto{\pgfqpoint{4.017344in}{1.336220in}}%
\pgfpathlineto{\pgfqpoint{4.017022in}{1.308184in}}%
\pgfpathlineto{\pgfqpoint{4.017938in}{1.312606in}}%
\pgfpathlineto{\pgfqpoint{4.018907in}{1.358229in}}%
\pgfpathlineto{\pgfqpoint{4.019233in}{1.323339in}}%
\pgfpathlineto{\pgfqpoint{4.019622in}{1.288266in}}%
\pgfpathlineto{\pgfqpoint{4.020372in}{1.302956in}}%
\pgfpathlineto{\pgfqpoint{4.020382in}{1.301330in}}%
\pgfpathlineto{\pgfqpoint{4.020873in}{1.328616in}}%
\pgfpathlineto{\pgfqpoint{4.021346in}{1.314243in}}%
\pgfpathlineto{\pgfqpoint{4.021360in}{1.316709in}}%
\pgfpathlineto{\pgfqpoint{4.022120in}{1.274564in}}%
\pgfpathlineto{\pgfqpoint{4.022383in}{1.285798in}}%
\pgfpathlineto{\pgfqpoint{4.022660in}{1.279920in}}%
\pgfpathlineto{\pgfqpoint{4.023429in}{1.311049in}}%
\pgfpathlineto{\pgfqpoint{4.024145in}{1.249034in}}%
\pgfpathlineto{\pgfqpoint{4.023454in}{1.311944in}}%
\pgfpathlineto{\pgfqpoint{4.024861in}{1.267022in}}%
\pgfpathlineto{\pgfqpoint{4.025002in}{1.273477in}}%
\pgfpathlineto{\pgfqpoint{4.025503in}{1.223757in}}%
\pgfpathlineto{\pgfqpoint{4.025922in}{1.250719in}}%
\pgfpathlineto{\pgfqpoint{4.026307in}{1.237037in}}%
\pgfpathlineto{\pgfqpoint{4.026049in}{1.255727in}}%
\pgfpathlineto{\pgfqpoint{4.026969in}{1.254058in}}%
\pgfpathlineto{\pgfqpoint{4.028006in}{1.277333in}}%
\pgfpathlineto{\pgfqpoint{4.027071in}{1.241922in}}%
\pgfpathlineto{\pgfqpoint{4.028103in}{1.270942in}}%
\pgfpathlineto{\pgfqpoint{4.028882in}{1.265154in}}%
\pgfpathlineto{\pgfqpoint{4.029096in}{1.294351in}}%
\pgfpathlineto{\pgfqpoint{4.029106in}{1.293656in}}%
\pgfpathlineto{\pgfqpoint{4.029116in}{1.295624in}}%
\pgfpathlineto{\pgfqpoint{4.030041in}{1.262619in}}%
\pgfpathlineto{\pgfqpoint{4.030114in}{1.265890in}}%
\pgfpathlineto{\pgfqpoint{4.030669in}{1.249482in}}%
\pgfpathlineto{\pgfqpoint{4.031010in}{1.290469in}}%
\pgfpathlineto{\pgfqpoint{4.031058in}{1.285083in}}%
\pgfpathlineto{\pgfqpoint{4.031881in}{1.321171in}}%
\pgfpathlineto{\pgfqpoint{4.031127in}{1.284929in}}%
\pgfpathlineto{\pgfqpoint{4.032183in}{1.296067in}}%
\pgfpathlineto{\pgfqpoint{4.032305in}{1.286202in}}%
\pgfpathlineto{\pgfqpoint{4.033059in}{1.317128in}}%
\pgfpathlineto{\pgfqpoint{4.033210in}{1.312645in}}%
\pgfpathlineto{\pgfqpoint{4.033312in}{1.300074in}}%
\pgfpathlineto{\pgfqpoint{4.034359in}{1.336376in}}%
\pgfpathlineto{\pgfqpoint{4.035299in}{1.321995in}}%
\pgfpathlineto{\pgfqpoint{4.034588in}{1.347778in}}%
\pgfpathlineto{\pgfqpoint{4.035464in}{1.339671in}}%
\pgfpathlineto{\pgfqpoint{4.036253in}{1.330069in}}%
\pgfpathlineto{\pgfqpoint{4.035674in}{1.357493in}}%
\pgfpathlineto{\pgfqpoint{4.036540in}{1.343401in}}%
\pgfpathlineto{\pgfqpoint{4.037538in}{1.359371in}}%
\pgfpathlineto{\pgfqpoint{4.037241in}{1.334167in}}%
\pgfpathlineto{\pgfqpoint{4.037675in}{1.357153in}}%
\pgfpathlineto{\pgfqpoint{4.038678in}{1.327205in}}%
\pgfpathlineto{\pgfqpoint{4.038220in}{1.367118in}}%
\pgfpathlineto{\pgfqpoint{4.038897in}{1.337448in}}%
\pgfpathlineto{\pgfqpoint{4.039330in}{1.357049in}}%
\pgfpathlineto{\pgfqpoint{4.039651in}{1.321721in}}%
\pgfpathlineto{\pgfqpoint{4.039992in}{1.331887in}}%
\pgfpathlineto{\pgfqpoint{4.041073in}{1.343605in}}%
\pgfpathlineto{\pgfqpoint{4.040221in}{1.310883in}}%
\pgfpathlineto{\pgfqpoint{4.041117in}{1.335428in}}%
\pgfpathlineto{\pgfqpoint{4.041569in}{1.317288in}}%
\pgfpathlineto{\pgfqpoint{4.041974in}{1.336524in}}%
\pgfpathlineto{\pgfqpoint{4.042246in}{1.329318in}}%
\pgfpathlineto{\pgfqpoint{4.043536in}{1.294847in}}%
\pgfpathlineto{\pgfqpoint{4.043609in}{1.299720in}}%
\pgfpathlineto{\pgfqpoint{4.043653in}{1.305652in}}%
\pgfpathlineto{\pgfqpoint{4.044583in}{1.249254in}}%
\pgfpathlineto{\pgfqpoint{4.044695in}{1.242560in}}%
\pgfpathlineto{\pgfqpoint{4.045498in}{1.297441in}}%
\pgfpathlineto{\pgfqpoint{4.045571in}{1.293587in}}%
\pgfpathlineto{\pgfqpoint{4.045649in}{1.299423in}}%
\pgfpathlineto{\pgfqpoint{4.046102in}{1.273556in}}%
\pgfpathlineto{\pgfqpoint{4.046438in}{1.284233in}}%
\pgfpathlineto{\pgfqpoint{4.047329in}{1.273504in}}%
\pgfpathlineto{\pgfqpoint{4.047076in}{1.305636in}}%
\pgfpathlineto{\pgfqpoint{4.047558in}{1.278898in}}%
\pgfpathlineto{\pgfqpoint{4.048312in}{1.309714in}}%
\pgfpathlineto{\pgfqpoint{4.048741in}{1.290072in}}%
\pgfpathlineto{\pgfqpoint{4.048819in}{1.292830in}}%
\pgfpathlineto{\pgfqpoint{4.049953in}{1.254934in}}%
\pgfpathlineto{\pgfqpoint{4.050343in}{1.249169in}}%
\pgfpathlineto{\pgfqpoint{4.050187in}{1.264069in}}%
\pgfpathlineto{\pgfqpoint{4.050401in}{1.260835in}}%
\pgfpathlineto{\pgfqpoint{4.050601in}{1.269062in}}%
\pgfpathlineto{\pgfqpoint{4.050795in}{1.247228in}}%
\pgfpathlineto{\pgfqpoint{4.051501in}{1.256292in}}%
\pgfpathlineto{\pgfqpoint{4.053648in}{1.379977in}}%
\pgfpathlineto{\pgfqpoint{4.051555in}{1.252549in}}%
\pgfpathlineto{\pgfqpoint{4.053668in}{1.375733in}}%
\pgfpathlineto{\pgfqpoint{4.054461in}{1.339157in}}%
\pgfpathlineto{\pgfqpoint{4.053897in}{1.381060in}}%
\pgfpathlineto{\pgfqpoint{4.054831in}{1.355339in}}%
\pgfpathlineto{\pgfqpoint{4.056063in}{1.414926in}}%
\pgfpathlineto{\pgfqpoint{4.055075in}{1.337671in}}%
\pgfpathlineto{\pgfqpoint{4.056068in}{1.414246in}}%
\pgfpathlineto{\pgfqpoint{4.056087in}{1.411150in}}%
\pgfpathlineto{\pgfqpoint{4.056384in}{1.432417in}}%
\pgfpathlineto{\pgfqpoint{4.056472in}{1.432018in}}%
\pgfpathlineto{\pgfqpoint{4.056569in}{1.444003in}}%
\pgfpathlineto{\pgfqpoint{4.057110in}{1.407477in}}%
\pgfpathlineto{\pgfqpoint{4.057426in}{1.415256in}}%
\pgfpathlineto{\pgfqpoint{4.057743in}{1.392509in}}%
\pgfpathlineto{\pgfqpoint{4.058307in}{1.423704in}}%
\pgfpathlineto{\pgfqpoint{4.058614in}{1.443052in}}%
\pgfpathlineto{\pgfqpoint{4.058341in}{1.414416in}}%
\pgfpathlineto{\pgfqpoint{4.059393in}{1.420952in}}%
\pgfpathlineto{\pgfqpoint{4.059500in}{1.410352in}}%
\pgfpathlineto{\pgfqpoint{4.059895in}{1.431176in}}%
\pgfpathlineto{\pgfqpoint{4.060513in}{1.410904in}}%
\pgfpathlineto{\pgfqpoint{4.060907in}{1.429048in}}%
\pgfpathlineto{\pgfqpoint{4.061326in}{1.398211in}}%
\pgfpathlineto{\pgfqpoint{4.061345in}{1.398640in}}%
\pgfpathlineto{\pgfqpoint{4.061501in}{1.390621in}}%
\pgfpathlineto{\pgfqpoint{4.061920in}{1.428576in}}%
\pgfpathlineto{\pgfqpoint{4.062382in}{1.423616in}}%
\pgfpathlineto{\pgfqpoint{4.062494in}{1.441905in}}%
\pgfpathlineto{\pgfqpoint{4.063361in}{1.380763in}}%
\pgfpathlineto{\pgfqpoint{4.063439in}{1.382856in}}%
\pgfpathlineto{\pgfqpoint{4.063444in}{1.383970in}}%
\pgfpathlineto{\pgfqpoint{4.064286in}{1.358826in}}%
\pgfpathlineto{\pgfqpoint{4.064451in}{1.360910in}}%
\pgfpathlineto{\pgfqpoint{4.064958in}{1.349775in}}%
\pgfpathlineto{\pgfqpoint{4.065435in}{1.383200in}}%
\pgfpathlineto{\pgfqpoint{4.065488in}{1.379155in}}%
\pgfpathlineto{\pgfqpoint{4.065493in}{1.379255in}}%
\pgfpathlineto{\pgfqpoint{4.065586in}{1.360122in}}%
\pgfpathlineto{\pgfqpoint{4.066637in}{1.336103in}}%
\pgfpathlineto{\pgfqpoint{4.065776in}{1.385475in}}%
\pgfpathlineto{\pgfqpoint{4.066725in}{1.346736in}}%
\pgfpathlineto{\pgfqpoint{4.067085in}{1.356792in}}%
\pgfpathlineto{\pgfqpoint{4.067217in}{1.337118in}}%
\pgfpathlineto{\pgfqpoint{4.067222in}{1.339057in}}%
\pgfpathlineto{\pgfqpoint{4.067962in}{1.327380in}}%
\pgfpathlineto{\pgfqpoint{4.067689in}{1.349891in}}%
\pgfpathlineto{\pgfqpoint{4.068341in}{1.331144in}}%
\pgfpathlineto{\pgfqpoint{4.069510in}{1.279739in}}%
\pgfpathlineto{\pgfqpoint{4.069524in}{1.286093in}}%
\pgfpathlineto{\pgfqpoint{4.070250in}{1.321893in}}%
\pgfpathlineto{\pgfqpoint{4.069666in}{1.282059in}}%
\pgfpathlineto{\pgfqpoint{4.070712in}{1.315316in}}%
\pgfpathlineto{\pgfqpoint{4.071467in}{1.336944in}}%
\pgfpathlineto{\pgfqpoint{4.071891in}{1.298568in}}%
\pgfpathlineto{\pgfqpoint{4.072562in}{1.310113in}}%
\pgfpathlineto{\pgfqpoint{4.072962in}{1.279130in}}%
\pgfpathlineto{\pgfqpoint{4.073624in}{1.259946in}}%
\pgfpathlineto{\pgfqpoint{4.073186in}{1.286336in}}%
\pgfpathlineto{\pgfqpoint{4.074067in}{1.280958in}}%
\pgfpathlineto{\pgfqpoint{4.074232in}{1.292316in}}%
\pgfpathlineto{\pgfqpoint{4.074719in}{1.256552in}}%
\pgfpathlineto{\pgfqpoint{4.075172in}{1.281068in}}%
\pgfpathlineto{\pgfqpoint{4.076024in}{1.264376in}}%
\pgfpathlineto{\pgfqpoint{4.075445in}{1.291523in}}%
\pgfpathlineto{\pgfqpoint{4.076287in}{1.276705in}}%
\pgfpathlineto{\pgfqpoint{4.076345in}{1.282728in}}%
\pgfpathlineto{\pgfqpoint{4.076740in}{1.257147in}}%
\pgfpathlineto{\pgfqpoint{4.077348in}{1.263642in}}%
\pgfpathlineto{\pgfqpoint{4.077971in}{1.287002in}}%
\pgfpathlineto{\pgfqpoint{4.077416in}{1.256511in}}%
\pgfpathlineto{\pgfqpoint{4.078434in}{1.264007in}}%
\pgfpathlineto{\pgfqpoint{4.079320in}{1.219681in}}%
\pgfpathlineto{\pgfqpoint{4.078541in}{1.268124in}}%
\pgfpathlineto{\pgfqpoint{4.079617in}{1.242522in}}%
\pgfpathlineto{\pgfqpoint{4.079826in}{1.222553in}}%
\pgfpathlineto{\pgfqpoint{4.080902in}{1.279227in}}%
\pgfpathlineto{\pgfqpoint{4.081472in}{1.272652in}}%
\pgfpathlineto{\pgfqpoint{4.081949in}{1.317457in}}%
\pgfpathlineto{\pgfqpoint{4.082032in}{1.319855in}}%
\pgfpathlineto{\pgfqpoint{4.082864in}{1.287165in}}%
\pgfpathlineto{\pgfqpoint{4.082971in}{1.290868in}}%
\pgfpathlineto{\pgfqpoint{4.083302in}{1.284444in}}%
\pgfpathlineto{\pgfqpoint{4.083945in}{1.316544in}}%
\pgfpathlineto{\pgfqpoint{4.084052in}{1.304576in}}%
\pgfpathlineto{\pgfqpoint{4.085220in}{1.338438in}}%
\pgfpathlineto{\pgfqpoint{4.084081in}{1.303894in}}%
\pgfpathlineto{\pgfqpoint{4.085230in}{1.338426in}}%
\pgfpathlineto{\pgfqpoint{4.085488in}{1.325413in}}%
\pgfpathlineto{\pgfqpoint{4.085960in}{1.374755in}}%
\pgfpathlineto{\pgfqpoint{4.086277in}{1.354836in}}%
\pgfpathlineto{\pgfqpoint{4.087684in}{1.405606in}}%
\pgfpathlineto{\pgfqpoint{4.086326in}{1.348352in}}%
\pgfpathlineto{\pgfqpoint{4.087825in}{1.400279in}}%
\pgfpathlineto{\pgfqpoint{4.087859in}{1.393836in}}%
\pgfpathlineto{\pgfqpoint{4.088794in}{1.445287in}}%
\pgfpathlineto{\pgfqpoint{4.088964in}{1.438046in}}%
\pgfpathlineto{\pgfqpoint{4.089329in}{1.468224in}}%
\pgfpathlineto{\pgfqpoint{4.089880in}{1.448538in}}%
\pgfpathlineto{\pgfqpoint{4.090610in}{1.474217in}}%
\pgfpathlineto{\pgfqpoint{4.090912in}{1.441475in}}%
\pgfpathlineto{\pgfqpoint{4.090990in}{1.450519in}}%
\pgfpathlineto{\pgfqpoint{4.090999in}{1.448500in}}%
\pgfpathlineto{\pgfqpoint{4.091939in}{1.502088in}}%
\pgfpathlineto{\pgfqpoint{4.091949in}{1.500583in}}%
\pgfpathlineto{\pgfqpoint{4.092095in}{1.506038in}}%
\pgfpathlineto{\pgfqpoint{4.092630in}{1.470439in}}%
\pgfpathlineto{\pgfqpoint{4.092903in}{1.472786in}}%
\pgfpathlineto{\pgfqpoint{4.093794in}{1.416253in}}%
\pgfpathlineto{\pgfqpoint{4.094247in}{1.448415in}}%
\pgfpathlineto{\pgfqpoint{4.095138in}{1.467474in}}%
\pgfpathlineto{\pgfqpoint{4.094631in}{1.427060in}}%
\pgfpathlineto{\pgfqpoint{4.095371in}{1.454363in}}%
\pgfpathlineto{\pgfqpoint{4.095697in}{1.452650in}}%
\pgfpathlineto{\pgfqpoint{4.096335in}{1.479048in}}%
\pgfpathlineto{\pgfqpoint{4.096350in}{1.477786in}}%
\pgfpathlineto{\pgfqpoint{4.096730in}{1.496796in}}%
\pgfpathlineto{\pgfqpoint{4.097002in}{1.476569in}}%
\pgfpathlineto{\pgfqpoint{4.097489in}{1.482765in}}%
\pgfpathlineto{\pgfqpoint{4.098073in}{1.460318in}}%
\pgfpathlineto{\pgfqpoint{4.098307in}{1.501841in}}%
\pgfpathlineto{\pgfqpoint{4.098575in}{1.492147in}}%
\pgfpathlineto{\pgfqpoint{4.100011in}{1.436333in}}%
\pgfpathlineto{\pgfqpoint{4.100055in}{1.442055in}}%
\pgfpathlineto{\pgfqpoint{4.100123in}{1.449527in}}%
\pgfpathlineto{\pgfqpoint{4.101072in}{1.402992in}}%
\pgfpathlineto{\pgfqpoint{4.101092in}{1.403083in}}%
\pgfpathlineto{\pgfqpoint{4.101140in}{1.397301in}}%
\pgfpathlineto{\pgfqpoint{4.102090in}{1.477370in}}%
\pgfpathlineto{\pgfqpoint{4.102104in}{1.475825in}}%
\pgfpathlineto{\pgfqpoint{4.102820in}{1.479968in}}%
\pgfpathlineto{\pgfqpoint{4.102474in}{1.447649in}}%
\pgfpathlineto{\pgfqpoint{4.103190in}{1.465620in}}%
\pgfpathlineto{\pgfqpoint{4.104329in}{1.424886in}}%
\pgfpathlineto{\pgfqpoint{4.103404in}{1.480169in}}%
\pgfpathlineto{\pgfqpoint{4.104427in}{1.433075in}}%
\pgfpathlineto{\pgfqpoint{4.104573in}{1.451111in}}%
\pgfpathlineto{\pgfqpoint{4.105327in}{1.399094in}}%
\pgfpathlineto{\pgfqpoint{4.105415in}{1.407480in}}%
\pgfpathlineto{\pgfqpoint{4.106306in}{1.376797in}}%
\pgfpathlineto{\pgfqpoint{4.106535in}{1.403136in}}%
\pgfpathlineto{\pgfqpoint{4.106632in}{1.407363in}}%
\pgfpathlineto{\pgfqpoint{4.106656in}{1.400235in}}%
\pgfpathlineto{\pgfqpoint{4.107528in}{1.365935in}}%
\pgfpathlineto{\pgfqpoint{4.106715in}{1.403868in}}%
\pgfpathlineto{\pgfqpoint{4.107805in}{1.378206in}}%
\pgfpathlineto{\pgfqpoint{4.108726in}{1.350665in}}%
\pgfpathlineto{\pgfqpoint{4.107922in}{1.385600in}}%
\pgfpathlineto{\pgfqpoint{4.108930in}{1.371081in}}%
\pgfpathlineto{\pgfqpoint{4.110060in}{1.396064in}}%
\pgfpathlineto{\pgfqpoint{4.109086in}{1.357936in}}%
\pgfpathlineto{\pgfqpoint{4.110064in}{1.392299in}}%
\pgfpathlineto{\pgfqpoint{4.110912in}{1.369334in}}%
\pgfpathlineto{\pgfqpoint{4.111131in}{1.397341in}}%
\pgfpathlineto{\pgfqpoint{4.111184in}{1.386960in}}%
\pgfpathlineto{\pgfqpoint{4.112202in}{1.401551in}}%
\pgfpathlineto{\pgfqpoint{4.111486in}{1.357796in}}%
\pgfpathlineto{\pgfqpoint{4.112299in}{1.389253in}}%
\pgfpathlineto{\pgfqpoint{4.112494in}{1.385684in}}%
\pgfpathlineto{\pgfqpoint{4.113044in}{1.417685in}}%
\pgfpathlineto{\pgfqpoint{4.114013in}{1.446223in}}%
\pgfpathlineto{\pgfqpoint{4.113579in}{1.403954in}}%
\pgfpathlineto{\pgfqpoint{4.114178in}{1.435969in}}%
\pgfpathlineto{\pgfqpoint{4.114709in}{1.413514in}}%
\pgfpathlineto{\pgfqpoint{4.114285in}{1.444821in}}%
\pgfpathlineto{\pgfqpoint{4.115288in}{1.431620in}}%
\pgfpathlineto{\pgfqpoint{4.115580in}{1.455666in}}%
\pgfpathlineto{\pgfqpoint{4.116067in}{1.419189in}}%
\pgfpathlineto{\pgfqpoint{4.116432in}{1.452113in}}%
\pgfpathlineto{\pgfqpoint{4.117586in}{1.406499in}}%
\pgfpathlineto{\pgfqpoint{4.117611in}{1.409871in}}%
\pgfpathlineto{\pgfqpoint{4.117947in}{1.432337in}}%
\pgfpathlineto{\pgfqpoint{4.118594in}{1.396346in}}%
\pgfpathlineto{\pgfqpoint{4.119641in}{1.365255in}}%
\pgfpathlineto{\pgfqpoint{4.118876in}{1.410413in}}%
\pgfpathlineto{\pgfqpoint{4.119753in}{1.375436in}}%
\pgfpathlineto{\pgfqpoint{4.120449in}{1.402384in}}%
\pgfpathlineto{\pgfqpoint{4.120098in}{1.371266in}}%
\pgfpathlineto{\pgfqpoint{4.120907in}{1.390713in}}%
\pgfpathlineto{\pgfqpoint{4.121092in}{1.401996in}}%
\pgfpathlineto{\pgfqpoint{4.121598in}{1.368091in}}%
\pgfpathlineto{\pgfqpoint{4.121929in}{1.372083in}}%
\pgfpathlineto{\pgfqpoint{4.122065in}{1.364157in}}%
\pgfpathlineto{\pgfqpoint{4.122523in}{1.387562in}}%
\pgfpathlineto{\pgfqpoint{4.123784in}{1.422219in}}%
\pgfpathlineto{\pgfqpoint{4.123073in}{1.380316in}}%
\pgfpathlineto{\pgfqpoint{4.123949in}{1.397102in}}%
\pgfpathlineto{\pgfqpoint{4.124008in}{1.388934in}}%
\pgfpathlineto{\pgfqpoint{4.124850in}{1.443027in}}%
\pgfpathlineto{\pgfqpoint{4.124938in}{1.429540in}}%
\pgfpathlineto{\pgfqpoint{4.125775in}{1.462331in}}%
\pgfpathlineto{\pgfqpoint{4.124986in}{1.416250in}}%
\pgfpathlineto{\pgfqpoint{4.126165in}{1.439076in}}%
\pgfpathlineto{\pgfqpoint{4.126520in}{1.410206in}}%
\pgfpathlineto{\pgfqpoint{4.126233in}{1.443797in}}%
\pgfpathlineto{\pgfqpoint{4.127309in}{1.423608in}}%
\pgfpathlineto{\pgfqpoint{4.127591in}{1.402776in}}%
\pgfpathlineto{\pgfqpoint{4.127352in}{1.430827in}}%
\pgfpathlineto{\pgfqpoint{4.128312in}{1.426847in}}%
\pgfpathlineto{\pgfqpoint{4.128370in}{1.436105in}}%
\pgfpathlineto{\pgfqpoint{4.128881in}{1.405199in}}%
\pgfpathlineto{\pgfqpoint{4.129368in}{1.412656in}}%
\pgfpathlineto{\pgfqpoint{4.129495in}{1.391769in}}%
\pgfpathlineto{\pgfqpoint{4.129894in}{1.417219in}}%
\pgfpathlineto{\pgfqpoint{4.130498in}{1.406218in}}%
\pgfpathlineto{\pgfqpoint{4.130790in}{1.424478in}}%
\pgfpathlineto{\pgfqpoint{4.131530in}{1.390910in}}%
\pgfpathlineto{\pgfqpoint{4.131588in}{1.392868in}}%
\pgfpathlineto{\pgfqpoint{4.131608in}{1.389796in}}%
\pgfpathlineto{\pgfqpoint{4.132328in}{1.417568in}}%
\pgfpathlineto{\pgfqpoint{4.132645in}{1.407978in}}%
\pgfpathlineto{\pgfqpoint{4.132747in}{1.419532in}}%
\pgfpathlineto{\pgfqpoint{4.133725in}{1.397792in}}%
\pgfpathlineto{\pgfqpoint{4.133730in}{1.398297in}}%
\pgfpathlineto{\pgfqpoint{4.133871in}{1.389645in}}%
\pgfpathlineto{\pgfqpoint{4.134519in}{1.430100in}}%
\pgfpathlineto{\pgfqpoint{4.134811in}{1.416854in}}%
\pgfpathlineto{\pgfqpoint{4.134908in}{1.412748in}}%
\pgfpathlineto{\pgfqpoint{4.135989in}{1.454972in}}%
\pgfpathlineto{\pgfqpoint{4.136749in}{1.435949in}}%
\pgfpathlineto{\pgfqpoint{4.136135in}{1.461402in}}%
\pgfpathlineto{\pgfqpoint{4.137099in}{1.454862in}}%
\pgfpathlineto{\pgfqpoint{4.137226in}{1.451213in}}%
\pgfpathlineto{\pgfqpoint{4.138005in}{1.475304in}}%
\pgfpathlineto{\pgfqpoint{4.138078in}{1.468822in}}%
\pgfpathlineto{\pgfqpoint{4.138428in}{1.467766in}}%
\pgfpathlineto{\pgfqpoint{4.139314in}{1.531721in}}%
\pgfpathlineto{\pgfqpoint{4.139426in}{1.529311in}}%
\pgfpathlineto{\pgfqpoint{4.139495in}{1.543556in}}%
\pgfpathlineto{\pgfqpoint{4.140605in}{1.593812in}}%
\pgfpathlineto{\pgfqpoint{4.139952in}{1.516250in}}%
\pgfpathlineto{\pgfqpoint{4.140682in}{1.583636in}}%
\pgfpathlineto{\pgfqpoint{4.141797in}{1.536963in}}%
\pgfpathlineto{\pgfqpoint{4.141973in}{1.563209in}}%
\pgfpathlineto{\pgfqpoint{4.142046in}{1.567580in}}%
\pgfpathlineto{\pgfqpoint{4.142990in}{1.531701in}}%
\pgfpathlineto{\pgfqpoint{4.143024in}{1.536909in}}%
\pgfpathlineto{\pgfqpoint{4.143813in}{1.509055in}}%
\pgfpathlineto{\pgfqpoint{4.143068in}{1.538034in}}%
\pgfpathlineto{\pgfqpoint{4.144149in}{1.528906in}}%
\pgfpathlineto{\pgfqpoint{4.144660in}{1.552244in}}%
\pgfpathlineto{\pgfqpoint{4.145205in}{1.514653in}}%
\pgfpathlineto{\pgfqpoint{4.145234in}{1.519019in}}%
\pgfpathlineto{\pgfqpoint{4.145244in}{1.517873in}}%
\pgfpathlineto{\pgfqpoint{4.146023in}{1.541388in}}%
\pgfpathlineto{\pgfqpoint{4.146048in}{1.540001in}}%
\pgfpathlineto{\pgfqpoint{4.146573in}{1.558294in}}%
\pgfpathlineto{\pgfqpoint{4.146247in}{1.530632in}}%
\pgfpathlineto{\pgfqpoint{4.147182in}{1.551088in}}%
\pgfpathlineto{\pgfqpoint{4.147201in}{1.548390in}}%
\pgfpathlineto{\pgfqpoint{4.147659in}{1.590330in}}%
\pgfpathlineto{\pgfqpoint{4.148190in}{1.564231in}}%
\pgfpathlineto{\pgfqpoint{4.149344in}{1.601173in}}%
\pgfpathlineto{\pgfqpoint{4.149373in}{1.600603in}}%
\pgfpathlineto{\pgfqpoint{4.150936in}{1.555864in}}%
\pgfpathlineto{\pgfqpoint{4.149928in}{1.608648in}}%
\pgfpathlineto{\pgfqpoint{4.150950in}{1.563001in}}%
\pgfpathlineto{\pgfqpoint{4.151242in}{1.575573in}}%
\pgfpathlineto{\pgfqpoint{4.151856in}{1.545364in}}%
\pgfpathlineto{\pgfqpoint{4.152050in}{1.561346in}}%
\pgfpathlineto{\pgfqpoint{4.152518in}{1.532070in}}%
\pgfpathlineto{\pgfqpoint{4.152245in}{1.562299in}}%
\pgfpathlineto{\pgfqpoint{4.153160in}{1.561663in}}%
\pgfpathlineto{\pgfqpoint{4.154300in}{1.590376in}}%
\pgfpathlineto{\pgfqpoint{4.153399in}{1.546962in}}%
\pgfpathlineto{\pgfqpoint{4.154314in}{1.589429in}}%
\pgfpathlineto{\pgfqpoint{4.155473in}{1.606291in}}%
\pgfpathlineto{\pgfqpoint{4.154416in}{1.581591in}}%
\pgfpathlineto{\pgfqpoint{4.155682in}{1.595782in}}%
\pgfpathlineto{\pgfqpoint{4.156700in}{1.586343in}}%
\pgfpathlineto{\pgfqpoint{4.156354in}{1.627834in}}%
\pgfpathlineto{\pgfqpoint{4.156787in}{1.598854in}}%
\pgfpathlineto{\pgfqpoint{4.157177in}{1.596746in}}%
\pgfpathlineto{\pgfqpoint{4.157435in}{1.648553in}}%
\pgfpathlineto{\pgfqpoint{4.157678in}{1.636024in}}%
\pgfpathlineto{\pgfqpoint{4.157966in}{1.659867in}}%
\pgfpathlineto{\pgfqpoint{4.158706in}{1.634487in}}%
\pgfpathlineto{\pgfqpoint{4.158798in}{1.644114in}}%
\pgfpathlineto{\pgfqpoint{4.158930in}{1.632952in}}%
\pgfpathlineto{\pgfqpoint{4.159514in}{1.683738in}}%
\pgfpathlineto{\pgfqpoint{4.159864in}{1.652325in}}%
\pgfpathlineto{\pgfqpoint{4.159918in}{1.658351in}}%
\pgfpathlineto{\pgfqpoint{4.160653in}{1.607430in}}%
\pgfpathlineto{\pgfqpoint{4.160750in}{1.609009in}}%
\pgfpathlineto{\pgfqpoint{4.161062in}{1.594146in}}%
\pgfpathlineto{\pgfqpoint{4.161695in}{1.627206in}}%
\pgfpathlineto{\pgfqpoint{4.161831in}{1.618196in}}%
\pgfpathlineto{\pgfqpoint{4.164022in}{1.699412in}}%
\pgfpathlineto{\pgfqpoint{4.162002in}{1.617789in}}%
\pgfpathlineto{\pgfqpoint{4.164192in}{1.677872in}}%
\pgfpathlineto{\pgfqpoint{4.164952in}{1.643813in}}%
\pgfpathlineto{\pgfqpoint{4.164270in}{1.685002in}}%
\pgfpathlineto{\pgfqpoint{4.165346in}{1.657751in}}%
\pgfpathlineto{\pgfqpoint{4.165862in}{1.665857in}}%
\pgfpathlineto{\pgfqpoint{4.165765in}{1.644546in}}%
\pgfpathlineto{\pgfqpoint{4.166471in}{1.662376in}}%
\pgfpathlineto{\pgfqpoint{4.166987in}{1.642972in}}%
\pgfpathlineto{\pgfqpoint{4.167396in}{1.673481in}}%
\pgfpathlineto{\pgfqpoint{4.167488in}{1.671457in}}%
\pgfpathlineto{\pgfqpoint{4.167931in}{1.700770in}}%
\pgfpathlineto{\pgfqpoint{4.168511in}{1.661912in}}%
\pgfpathlineto{\pgfqpoint{4.168598in}{1.672159in}}%
\pgfpathlineto{\pgfqpoint{4.169981in}{1.710586in}}%
\pgfpathlineto{\pgfqpoint{4.168939in}{1.661542in}}%
\pgfpathlineto{\pgfqpoint{4.170677in}{1.691962in}}%
\pgfpathlineto{\pgfqpoint{4.170687in}{1.688539in}}%
\pgfpathlineto{\pgfqpoint{4.171709in}{1.729074in}}%
\pgfpathlineto{\pgfqpoint{4.171724in}{1.720943in}}%
\pgfpathlineto{\pgfqpoint{4.172547in}{1.734214in}}%
\pgfpathlineto{\pgfqpoint{4.172698in}{1.707910in}}%
\pgfpathlineto{\pgfqpoint{4.173822in}{1.650457in}}%
\pgfpathlineto{\pgfqpoint{4.173871in}{1.653080in}}%
\pgfpathlineto{\pgfqpoint{4.173993in}{1.657956in}}%
\pgfpathlineto{\pgfqpoint{4.174353in}{1.613266in}}%
\pgfpathlineto{\pgfqpoint{4.174742in}{1.634684in}}%
\pgfpathlineto{\pgfqpoint{4.175190in}{1.625042in}}%
\pgfpathlineto{\pgfqpoint{4.175804in}{1.641893in}}%
\pgfpathlineto{\pgfqpoint{4.175828in}{1.638387in}}%
\pgfpathlineto{\pgfqpoint{4.175999in}{1.648250in}}%
\pgfpathlineto{\pgfqpoint{4.176739in}{1.609854in}}%
\pgfpathlineto{\pgfqpoint{4.176841in}{1.620016in}}%
\pgfpathlineto{\pgfqpoint{4.177688in}{1.583551in}}%
\pgfpathlineto{\pgfqpoint{4.177123in}{1.620079in}}%
\pgfpathlineto{\pgfqpoint{4.177980in}{1.595719in}}%
\pgfpathlineto{\pgfqpoint{4.178598in}{1.611786in}}%
\pgfpathlineto{\pgfqpoint{4.178973in}{1.565935in}}%
\pgfpathlineto{\pgfqpoint{4.179061in}{1.579960in}}%
\pgfpathlineto{\pgfqpoint{4.179553in}{1.554893in}}%
\pgfpathlineto{\pgfqpoint{4.179952in}{1.586422in}}%
\pgfpathlineto{\pgfqpoint{4.180215in}{1.573743in}}%
\pgfpathlineto{\pgfqpoint{4.180234in}{1.576363in}}%
\pgfpathlineto{\pgfqpoint{4.181203in}{1.526058in}}%
\pgfpathlineto{\pgfqpoint{4.181227in}{1.529826in}}%
\pgfpathlineto{\pgfqpoint{4.181773in}{1.498791in}}%
\pgfpathlineto{\pgfqpoint{4.181257in}{1.530047in}}%
\pgfpathlineto{\pgfqpoint{4.182352in}{1.519497in}}%
\pgfpathlineto{\pgfqpoint{4.182819in}{1.533863in}}%
\pgfpathlineto{\pgfqpoint{4.183228in}{1.503109in}}%
\pgfpathlineto{\pgfqpoint{4.183442in}{1.516817in}}%
\pgfpathlineto{\pgfqpoint{4.183837in}{1.485694in}}%
\pgfpathlineto{\pgfqpoint{4.184562in}{1.508066in}}%
\pgfpathlineto{\pgfqpoint{4.184596in}{1.517164in}}%
\pgfpathlineto{\pgfqpoint{4.185010in}{1.469766in}}%
\pgfpathlineto{\pgfqpoint{4.185429in}{1.482737in}}%
\pgfpathlineto{\pgfqpoint{4.185687in}{1.468733in}}%
\pgfpathlineto{\pgfqpoint{4.186344in}{1.493153in}}%
\pgfpathlineto{\pgfqpoint{4.186539in}{1.482826in}}%
\pgfpathlineto{\pgfqpoint{4.186588in}{1.490782in}}%
\pgfpathlineto{\pgfqpoint{4.186875in}{1.453864in}}%
\pgfpathlineto{\pgfqpoint{4.187639in}{1.476991in}}%
\pgfpathlineto{\pgfqpoint{4.187688in}{1.474302in}}%
\pgfpathlineto{\pgfqpoint{4.187887in}{1.497361in}}%
\pgfpathlineto{\pgfqpoint{4.188520in}{1.484850in}}%
\pgfpathlineto{\pgfqpoint{4.188915in}{1.519888in}}%
\pgfpathlineto{\pgfqpoint{4.188662in}{1.477350in}}%
\pgfpathlineto{\pgfqpoint{4.189630in}{1.486029in}}%
\pgfpathlineto{\pgfqpoint{4.189888in}{1.454622in}}%
\pgfpathlineto{\pgfqpoint{4.189655in}{1.488274in}}%
\pgfpathlineto{\pgfqpoint{4.190750in}{1.481422in}}%
\pgfpathlineto{\pgfqpoint{4.191914in}{1.453700in}}%
\pgfpathlineto{\pgfqpoint{4.190838in}{1.486230in}}%
\pgfpathlineto{\pgfqpoint{4.191943in}{1.459293in}}%
\pgfpathlineto{\pgfqpoint{4.192133in}{1.464578in}}%
\pgfpathlineto{\pgfqpoint{4.192259in}{1.442440in}}%
\pgfpathlineto{\pgfqpoint{4.192795in}{1.445604in}}%
\pgfpathlineto{\pgfqpoint{4.193988in}{1.396218in}}%
\pgfpathlineto{\pgfqpoint{4.192917in}{1.450045in}}%
\pgfpathlineto{\pgfqpoint{4.194031in}{1.399942in}}%
\pgfpathlineto{\pgfqpoint{4.194942in}{1.419326in}}%
\pgfpathlineto{\pgfqpoint{4.194518in}{1.396321in}}%
\pgfpathlineto{\pgfqpoint{4.195161in}{1.413534in}}%
\pgfpathlineto{\pgfqpoint{4.195219in}{1.403117in}}%
\pgfpathlineto{\pgfqpoint{4.196110in}{1.457345in}}%
\pgfpathlineto{\pgfqpoint{4.196115in}{1.456280in}}%
\pgfpathlineto{\pgfqpoint{4.196982in}{1.471906in}}%
\pgfpathlineto{\pgfqpoint{4.196261in}{1.446397in}}%
\pgfpathlineto{\pgfqpoint{4.197211in}{1.450748in}}%
\pgfpathlineto{\pgfqpoint{4.198214in}{1.415450in}}%
\pgfpathlineto{\pgfqpoint{4.197298in}{1.455465in}}%
\pgfpathlineto{\pgfqpoint{4.198369in}{1.421831in}}%
\pgfpathlineto{\pgfqpoint{4.198924in}{1.447373in}}%
\pgfpathlineto{\pgfqpoint{4.199255in}{1.407086in}}%
\pgfpathlineto{\pgfqpoint{4.199416in}{1.391056in}}%
\pgfpathlineto{\pgfqpoint{4.199314in}{1.414011in}}%
\pgfpathlineto{\pgfqpoint{4.200375in}{1.401721in}}%
\pgfpathlineto{\pgfqpoint{4.200385in}{1.403466in}}%
\pgfpathlineto{\pgfqpoint{4.200526in}{1.380688in}}%
\pgfpathlineto{\pgfqpoint{4.201461in}{1.395290in}}%
\pgfpathlineto{\pgfqpoint{4.202323in}{1.376063in}}%
\pgfpathlineto{\pgfqpoint{4.201534in}{1.401167in}}%
\pgfpathlineto{\pgfqpoint{4.202610in}{1.378304in}}%
\pgfpathlineto{\pgfqpoint{4.203730in}{1.402638in}}%
\pgfpathlineto{\pgfqpoint{4.202663in}{1.371064in}}%
\pgfpathlineto{\pgfqpoint{4.203754in}{1.399058in}}%
\pgfpathlineto{\pgfqpoint{4.203939in}{1.417340in}}%
\pgfpathlineto{\pgfqpoint{4.204280in}{1.373764in}}%
\pgfpathlineto{\pgfqpoint{4.204781in}{1.385075in}}%
\pgfpathlineto{\pgfqpoint{4.204830in}{1.375115in}}%
\pgfpathlineto{\pgfqpoint{4.205706in}{1.427062in}}%
\pgfpathlineto{\pgfqpoint{4.205828in}{1.411689in}}%
\pgfpathlineto{\pgfqpoint{4.205945in}{1.421972in}}%
\pgfpathlineto{\pgfqpoint{4.206607in}{1.388341in}}%
\pgfpathlineto{\pgfqpoint{4.206928in}{1.406524in}}%
\pgfpathlineto{\pgfqpoint{4.208218in}{1.368998in}}%
\pgfpathlineto{\pgfqpoint{4.207133in}{1.422197in}}%
\pgfpathlineto{\pgfqpoint{4.208316in}{1.387549in}}%
\pgfpathlineto{\pgfqpoint{4.208666in}{1.417926in}}%
\pgfpathlineto{\pgfqpoint{4.208350in}{1.386885in}}%
\pgfpathlineto{\pgfqpoint{4.209460in}{1.398419in}}%
\pgfpathlineto{\pgfqpoint{4.209932in}{1.361584in}}%
\pgfpathlineto{\pgfqpoint{4.209635in}{1.403954in}}%
\pgfpathlineto{\pgfqpoint{4.210584in}{1.392163in}}%
\pgfpathlineto{\pgfqpoint{4.211154in}{1.402047in}}%
\pgfpathlineto{\pgfqpoint{4.211407in}{1.367810in}}%
\pgfpathlineto{\pgfqpoint{4.211660in}{1.380644in}}%
\pgfpathlineto{\pgfqpoint{4.211685in}{1.372010in}}%
\pgfpathlineto{\pgfqpoint{4.212727in}{1.408204in}}%
\pgfpathlineto{\pgfqpoint{4.213087in}{1.385794in}}%
\pgfpathlineto{\pgfqpoint{4.212863in}{1.425122in}}%
\pgfpathlineto{\pgfqpoint{4.213773in}{1.418943in}}%
\pgfpathlineto{\pgfqpoint{4.213924in}{1.433069in}}%
\pgfpathlineto{\pgfqpoint{4.214776in}{1.404056in}}%
\pgfpathlineto{\pgfqpoint{4.214830in}{1.409278in}}%
\pgfpathlineto{\pgfqpoint{4.214839in}{1.406450in}}%
\pgfpathlineto{\pgfqpoint{4.215122in}{1.428470in}}%
\pgfpathlineto{\pgfqpoint{4.215911in}{1.418222in}}%
\pgfpathlineto{\pgfqpoint{4.216242in}{1.426444in}}%
\pgfpathlineto{\pgfqpoint{4.216865in}{1.397200in}}%
\pgfpathlineto{\pgfqpoint{4.216875in}{1.399304in}}%
\pgfpathlineto{\pgfqpoint{4.217692in}{1.372189in}}%
\pgfpathlineto{\pgfqpoint{4.218004in}{1.389740in}}%
\pgfpathlineto{\pgfqpoint{4.219430in}{1.453859in}}%
\pgfpathlineto{\pgfqpoint{4.219484in}{1.447476in}}%
\pgfpathlineto{\pgfqpoint{4.220409in}{1.434751in}}%
\pgfpathlineto{\pgfqpoint{4.219786in}{1.464171in}}%
\pgfpathlineto{\pgfqpoint{4.220579in}{1.448258in}}%
\pgfpathlineto{\pgfqpoint{4.221553in}{1.466258in}}%
\pgfpathlineto{\pgfqpoint{4.221217in}{1.427708in}}%
\pgfpathlineto{\pgfqpoint{4.221709in}{1.465978in}}%
\pgfpathlineto{\pgfqpoint{4.223169in}{1.511678in}}%
\pgfpathlineto{\pgfqpoint{4.222113in}{1.454036in}}%
\pgfpathlineto{\pgfqpoint{4.223189in}{1.506615in}}%
\pgfpathlineto{\pgfqpoint{4.223316in}{1.501192in}}%
\pgfpathlineto{\pgfqpoint{4.224143in}{1.552666in}}%
\pgfpathlineto{\pgfqpoint{4.224197in}{1.545680in}}%
\pgfpathlineto{\pgfqpoint{4.224299in}{1.551892in}}%
\pgfpathlineto{\pgfqpoint{4.224684in}{1.528219in}}%
\pgfpathlineto{\pgfqpoint{4.225273in}{1.537147in}}%
\pgfpathlineto{\pgfqpoint{4.225876in}{1.492856in}}%
\pgfpathlineto{\pgfqpoint{4.226402in}{1.528407in}}%
\pgfpathlineto{\pgfqpoint{4.227561in}{1.505336in}}%
\pgfpathlineto{\pgfqpoint{4.226870in}{1.532847in}}%
\pgfpathlineto{\pgfqpoint{4.227571in}{1.510481in}}%
\pgfpathlineto{\pgfqpoint{4.228627in}{1.530299in}}%
\pgfpathlineto{\pgfqpoint{4.227780in}{1.497383in}}%
\pgfpathlineto{\pgfqpoint{4.228700in}{1.518311in}}%
\pgfpathlineto{\pgfqpoint{4.228720in}{1.514895in}}%
\pgfpathlineto{\pgfqpoint{4.229513in}{1.575071in}}%
\pgfpathlineto{\pgfqpoint{4.229722in}{1.555005in}}%
\pgfpathlineto{\pgfqpoint{4.231280in}{1.511933in}}%
\pgfpathlineto{\pgfqpoint{4.229757in}{1.557268in}}%
\pgfpathlineto{\pgfqpoint{4.231397in}{1.520462in}}%
\pgfpathlineto{\pgfqpoint{4.232566in}{1.581473in}}%
\pgfpathlineto{\pgfqpoint{4.231412in}{1.518699in}}%
\pgfpathlineto{\pgfqpoint{4.232610in}{1.572347in}}%
\pgfpathlineto{\pgfqpoint{4.232741in}{1.562280in}}%
\pgfpathlineto{\pgfqpoint{4.232906in}{1.595893in}}%
\pgfpathlineto{\pgfqpoint{4.233715in}{1.570725in}}%
\pgfpathlineto{\pgfqpoint{4.234859in}{1.602837in}}%
\pgfpathlineto{\pgfqpoint{4.233987in}{1.569903in}}%
\pgfpathlineto{\pgfqpoint{4.234878in}{1.599271in}}%
\pgfpathlineto{\pgfqpoint{4.236154in}{1.634982in}}%
\pgfpathlineto{\pgfqpoint{4.235462in}{1.584540in}}%
\pgfpathlineto{\pgfqpoint{4.236198in}{1.627785in}}%
\pgfpathlineto{\pgfqpoint{4.236407in}{1.617275in}}%
\pgfpathlineto{\pgfqpoint{4.236830in}{1.650869in}}%
\pgfpathlineto{\pgfqpoint{4.237244in}{1.638118in}}%
\pgfpathlineto{\pgfqpoint{4.237371in}{1.649190in}}%
\pgfpathlineto{\pgfqpoint{4.237848in}{1.617533in}}%
\pgfpathlineto{\pgfqpoint{4.238359in}{1.643801in}}%
\pgfpathlineto{\pgfqpoint{4.239450in}{1.633654in}}%
\pgfpathlineto{\pgfqpoint{4.238885in}{1.663994in}}%
\pgfpathlineto{\pgfqpoint{4.239498in}{1.637929in}}%
\pgfpathlineto{\pgfqpoint{4.239503in}{1.639077in}}%
\pgfpathlineto{\pgfqpoint{4.240219in}{1.573777in}}%
\pgfpathlineto{\pgfqpoint{4.240423in}{1.596088in}}%
\pgfpathlineto{\pgfqpoint{4.241047in}{1.582329in}}%
\pgfpathlineto{\pgfqpoint{4.240604in}{1.602326in}}%
\pgfpathlineto{\pgfqpoint{4.241548in}{1.586023in}}%
\pgfpathlineto{\pgfqpoint{4.241553in}{1.588134in}}%
\pgfpathlineto{\pgfqpoint{4.242390in}{1.510167in}}%
\pgfpathlineto{\pgfqpoint{4.242483in}{1.515455in}}%
\pgfpathlineto{\pgfqpoint{4.243544in}{1.490139in}}%
\pgfpathlineto{\pgfqpoint{4.242668in}{1.518885in}}%
\pgfpathlineto{\pgfqpoint{4.243617in}{1.492068in}}%
\pgfpathlineto{\pgfqpoint{4.244747in}{1.542933in}}%
\pgfpathlineto{\pgfqpoint{4.244781in}{1.537238in}}%
\pgfpathlineto{\pgfqpoint{4.245871in}{1.550670in}}%
\pgfpathlineto{\pgfqpoint{4.245681in}{1.525079in}}%
\pgfpathlineto{\pgfqpoint{4.245925in}{1.547096in}}%
\pgfpathlineto{\pgfqpoint{4.247064in}{1.490501in}}%
\pgfpathlineto{\pgfqpoint{4.246042in}{1.548638in}}%
\pgfpathlineto{\pgfqpoint{4.247108in}{1.496138in}}%
\pgfpathlineto{\pgfqpoint{4.247234in}{1.511827in}}%
\pgfpathlineto{\pgfqpoint{4.247931in}{1.467957in}}%
\pgfpathlineto{\pgfqpoint{4.248155in}{1.475892in}}%
\pgfpathlineto{\pgfqpoint{4.249226in}{1.418740in}}%
\pgfpathlineto{\pgfqpoint{4.248247in}{1.483976in}}%
\pgfpathlineto{\pgfqpoint{4.249790in}{1.444872in}}%
\pgfpathlineto{\pgfqpoint{4.249975in}{1.457968in}}%
\pgfpathlineto{\pgfqpoint{4.250565in}{1.414471in}}%
\pgfpathlineto{\pgfqpoint{4.250852in}{1.432616in}}%
\pgfpathlineto{\pgfqpoint{4.251777in}{1.394263in}}%
\pgfpathlineto{\pgfqpoint{4.251154in}{1.436045in}}%
\pgfpathlineto{\pgfqpoint{4.252020in}{1.414237in}}%
\pgfpathlineto{\pgfqpoint{4.253242in}{1.481194in}}%
\pgfpathlineto{\pgfqpoint{4.253247in}{1.481032in}}%
\pgfpathlineto{\pgfqpoint{4.253354in}{1.469367in}}%
\pgfpathlineto{\pgfqpoint{4.253792in}{1.493435in}}%
\pgfpathlineto{\pgfqpoint{4.254342in}{1.486067in}}%
\pgfpathlineto{\pgfqpoint{4.255379in}{1.509768in}}%
\pgfpathlineto{\pgfqpoint{4.254508in}{1.480570in}}%
\pgfpathlineto{\pgfqpoint{4.255482in}{1.505324in}}%
\pgfpathlineto{\pgfqpoint{4.256460in}{1.484782in}}%
\pgfpathlineto{\pgfqpoint{4.256003in}{1.509353in}}%
\pgfpathlineto{\pgfqpoint{4.256626in}{1.490547in}}%
\pgfpathlineto{\pgfqpoint{4.256631in}{1.492108in}}%
\pgfpathlineto{\pgfqpoint{4.257531in}{1.451303in}}%
\pgfpathlineto{\pgfqpoint{4.257629in}{1.456402in}}%
\pgfpathlineto{\pgfqpoint{4.258734in}{1.490075in}}%
\pgfpathlineto{\pgfqpoint{4.257692in}{1.450937in}}%
\pgfpathlineto{\pgfqpoint{4.258890in}{1.479665in}}%
\pgfpathlineto{\pgfqpoint{4.260477in}{1.527401in}}%
\pgfpathlineto{\pgfqpoint{4.258909in}{1.477306in}}%
\pgfpathlineto{\pgfqpoint{4.260530in}{1.519405in}}%
\pgfpathlineto{\pgfqpoint{4.260540in}{1.517218in}}%
\pgfpathlineto{\pgfqpoint{4.261533in}{1.560225in}}%
\pgfpathlineto{\pgfqpoint{4.261874in}{1.566536in}}%
\pgfpathlineto{\pgfqpoint{4.262351in}{1.539298in}}%
\pgfpathlineto{\pgfqpoint{4.262648in}{1.561995in}}%
\pgfpathlineto{\pgfqpoint{4.262663in}{1.555949in}}%
\pgfpathlineto{\pgfqpoint{4.262750in}{1.580372in}}%
\pgfpathlineto{\pgfqpoint{4.263393in}{1.577787in}}%
\pgfpathlineto{\pgfqpoint{4.263739in}{1.609155in}}%
\pgfpathlineto{\pgfqpoint{4.264367in}{1.559010in}}%
\pgfpathlineto{\pgfqpoint{4.264469in}{1.563327in}}%
\pgfpathlineto{\pgfqpoint{4.264751in}{1.528412in}}%
\pgfpathlineto{\pgfqpoint{4.265594in}{1.549482in}}%
\pgfpathlineto{\pgfqpoint{4.265754in}{1.552619in}}%
\pgfpathlineto{\pgfqpoint{4.266149in}{1.523482in}}%
\pgfpathlineto{\pgfqpoint{4.266670in}{1.542460in}}%
\pgfpathlineto{\pgfqpoint{4.267434in}{1.520311in}}%
\pgfpathlineto{\pgfqpoint{4.266893in}{1.550899in}}%
\pgfpathlineto{\pgfqpoint{4.267804in}{1.529485in}}%
\pgfpathlineto{\pgfqpoint{4.268417in}{1.572721in}}%
\pgfpathlineto{\pgfqpoint{4.268958in}{1.560275in}}%
\pgfpathlineto{\pgfqpoint{4.268997in}{1.553248in}}%
\pgfpathlineto{\pgfqpoint{4.269469in}{1.585599in}}%
\pgfpathlineto{\pgfqpoint{4.270000in}{1.571708in}}%
\pgfpathlineto{\pgfqpoint{4.270516in}{1.595701in}}%
\pgfpathlineto{\pgfqpoint{4.270039in}{1.570411in}}%
\pgfpathlineto{\pgfqpoint{4.271129in}{1.576260in}}%
\pgfpathlineto{\pgfqpoint{4.271256in}{1.566142in}}%
\pgfpathlineto{\pgfqpoint{4.272098in}{1.614310in}}%
\pgfpathlineto{\pgfqpoint{4.272186in}{1.610997in}}%
\pgfpathlineto{\pgfqpoint{4.272195in}{1.617149in}}%
\pgfpathlineto{\pgfqpoint{4.273184in}{1.568957in}}%
\pgfpathlineto{\pgfqpoint{4.273223in}{1.569036in}}%
\pgfpathlineto{\pgfqpoint{4.273524in}{1.542321in}}%
\pgfpathlineto{\pgfqpoint{4.274221in}{1.612710in}}%
\pgfpathlineto{\pgfqpoint{4.274298in}{1.610883in}}%
\pgfpathlineto{\pgfqpoint{4.274415in}{1.605535in}}%
\pgfpathlineto{\pgfqpoint{4.275112in}{1.668408in}}%
\pgfpathlineto{\pgfqpoint{4.275189in}{1.659962in}}%
\pgfpathlineto{\pgfqpoint{4.275545in}{1.680687in}}%
\pgfpathlineto{\pgfqpoint{4.275788in}{1.648233in}}%
\pgfpathlineto{\pgfqpoint{4.276226in}{1.650185in}}%
\pgfpathlineto{\pgfqpoint{4.276996in}{1.604916in}}%
\pgfpathlineto{\pgfqpoint{4.277371in}{1.612510in}}%
\pgfpathlineto{\pgfqpoint{4.278461in}{1.575248in}}%
\pgfpathlineto{\pgfqpoint{4.277492in}{1.621003in}}%
\pgfpathlineto{\pgfqpoint{4.278568in}{1.591324in}}%
\pgfpathlineto{\pgfqpoint{4.279435in}{1.603931in}}%
\pgfpathlineto{\pgfqpoint{4.278953in}{1.573276in}}%
\pgfpathlineto{\pgfqpoint{4.279678in}{1.593687in}}%
\pgfpathlineto{\pgfqpoint{4.279970in}{1.581526in}}%
\pgfpathlineto{\pgfqpoint{4.280428in}{1.607645in}}%
\pgfpathlineto{\pgfqpoint{4.280803in}{1.587499in}}%
\pgfpathlineto{\pgfqpoint{4.281407in}{1.590715in}}%
\pgfpathlineto{\pgfqpoint{4.281718in}{1.548126in}}%
\pgfpathlineto{\pgfqpoint{4.281728in}{1.548465in}}%
\pgfpathlineto{\pgfqpoint{4.281772in}{1.537480in}}%
\pgfpathlineto{\pgfqpoint{4.282702in}{1.579563in}}%
\pgfpathlineto{\pgfqpoint{4.282813in}{1.559726in}}%
\pgfpathlineto{\pgfqpoint{4.283972in}{1.532200in}}%
\pgfpathlineto{\pgfqpoint{4.282916in}{1.566626in}}%
\pgfpathlineto{\pgfqpoint{4.284099in}{1.543359in}}%
\pgfpathlineto{\pgfqpoint{4.285233in}{1.585181in}}%
\pgfpathlineto{\pgfqpoint{4.285287in}{1.583817in}}%
\pgfpathlineto{\pgfqpoint{4.285915in}{1.556590in}}%
\pgfpathlineto{\pgfqpoint{4.285710in}{1.587238in}}%
\pgfpathlineto{\pgfqpoint{4.286397in}{1.582869in}}%
\pgfpathlineto{\pgfqpoint{4.287429in}{1.637944in}}%
\pgfpathlineto{\pgfqpoint{4.287779in}{1.617980in}}%
\pgfpathlineto{\pgfqpoint{4.287896in}{1.598282in}}%
\pgfpathlineto{\pgfqpoint{4.288271in}{1.626163in}}%
\pgfpathlineto{\pgfqpoint{4.288933in}{1.605754in}}%
\pgfpathlineto{\pgfqpoint{4.289420in}{1.624786in}}%
\pgfpathlineto{\pgfqpoint{4.289128in}{1.601939in}}%
\pgfpathlineto{\pgfqpoint{4.290058in}{1.614119in}}%
\pgfpathlineto{\pgfqpoint{4.290725in}{1.603282in}}%
\pgfpathlineto{\pgfqpoint{4.290550in}{1.628787in}}%
\pgfpathlineto{\pgfqpoint{4.291158in}{1.617166in}}%
\pgfpathlineto{\pgfqpoint{4.292151in}{1.648590in}}%
\pgfpathlineto{\pgfqpoint{4.291718in}{1.600605in}}%
\pgfpathlineto{\pgfqpoint{4.292302in}{1.627768in}}%
\pgfpathlineto{\pgfqpoint{4.293315in}{1.605704in}}%
\pgfpathlineto{\pgfqpoint{4.293032in}{1.630696in}}%
\pgfpathlineto{\pgfqpoint{4.293417in}{1.623455in}}%
\pgfpathlineto{\pgfqpoint{4.294474in}{1.642983in}}%
\pgfpathlineto{\pgfqpoint{4.294006in}{1.610429in}}%
\pgfpathlineto{\pgfqpoint{4.294542in}{1.632217in}}%
\pgfpathlineto{\pgfqpoint{4.295661in}{1.609803in}}%
\pgfpathlineto{\pgfqpoint{4.295204in}{1.648109in}}%
\pgfpathlineto{\pgfqpoint{4.295739in}{1.618522in}}%
\pgfpathlineto{\pgfqpoint{4.295827in}{1.626329in}}%
\pgfpathlineto{\pgfqpoint{4.296178in}{1.604203in}}%
\pgfpathlineto{\pgfqpoint{4.296820in}{1.608925in}}%
\pgfpathlineto{\pgfqpoint{4.297025in}{1.593376in}}%
\pgfpathlineto{\pgfqpoint{4.297229in}{1.629256in}}%
\pgfpathlineto{\pgfqpoint{4.297804in}{1.626547in}}%
\pgfpathlineto{\pgfqpoint{4.298568in}{1.666409in}}%
\pgfpathlineto{\pgfqpoint{4.298957in}{1.652136in}}%
\pgfpathlineto{\pgfqpoint{4.299571in}{1.647371in}}%
\pgfpathlineto{\pgfqpoint{4.299766in}{1.673710in}}%
\pgfpathlineto{\pgfqpoint{4.299770in}{1.672967in}}%
\pgfpathlineto{\pgfqpoint{4.300053in}{1.685591in}}%
\pgfpathlineto{\pgfqpoint{4.300559in}{1.645223in}}%
\pgfpathlineto{\pgfqpoint{4.300846in}{1.661563in}}%
\pgfpathlineto{\pgfqpoint{4.301572in}{1.631283in}}%
\pgfpathlineto{\pgfqpoint{4.300895in}{1.664644in}}%
\pgfpathlineto{\pgfqpoint{4.302005in}{1.649539in}}%
\pgfpathlineto{\pgfqpoint{4.302667in}{1.669120in}}%
\pgfpathlineto{\pgfqpoint{4.302945in}{1.635912in}}%
\pgfpathlineto{\pgfqpoint{4.302989in}{1.636395in}}%
\pgfpathlineto{\pgfqpoint{4.303217in}{1.646035in}}%
\pgfpathlineto{\pgfqpoint{4.303242in}{1.645939in}}%
\pgfpathlineto{\pgfqpoint{4.303271in}{1.650667in}}%
\pgfpathlineto{\pgfqpoint{4.303889in}{1.600904in}}%
\pgfpathlineto{\pgfqpoint{4.304308in}{1.633553in}}%
\pgfpathlineto{\pgfqpoint{4.305618in}{1.602241in}}%
\pgfpathlineto{\pgfqpoint{4.305072in}{1.643982in}}%
\pgfpathlineto{\pgfqpoint{4.305622in}{1.603669in}}%
\pgfpathlineto{\pgfqpoint{4.305695in}{1.620659in}}%
\pgfpathlineto{\pgfqpoint{4.306679in}{1.571416in}}%
\pgfpathlineto{\pgfqpoint{4.306684in}{1.572780in}}%
\pgfpathlineto{\pgfqpoint{4.306883in}{1.574098in}}%
\pgfpathlineto{\pgfqpoint{4.307351in}{1.533636in}}%
\pgfpathlineto{\pgfqpoint{4.307545in}{1.542627in}}%
\pgfpathlineto{\pgfqpoint{4.308685in}{1.513889in}}%
\pgfpathlineto{\pgfqpoint{4.307560in}{1.548250in}}%
\pgfpathlineto{\pgfqpoint{4.308753in}{1.521434in}}%
\pgfpathlineto{\pgfqpoint{4.309147in}{1.542727in}}%
\pgfpathlineto{\pgfqpoint{4.309809in}{1.513235in}}%
\pgfpathlineto{\pgfqpoint{4.309853in}{1.518067in}}%
\pgfpathlineto{\pgfqpoint{4.309960in}{1.521097in}}%
\pgfpathlineto{\pgfqpoint{4.310510in}{1.496490in}}%
\pgfpathlineto{\pgfqpoint{4.310515in}{1.496645in}}%
\pgfpathlineto{\pgfqpoint{4.311022in}{1.479470in}}%
\pgfpathlineto{\pgfqpoint{4.311543in}{1.509214in}}%
\pgfpathlineto{\pgfqpoint{4.311577in}{1.502928in}}%
\pgfpathlineto{\pgfqpoint{4.312000in}{1.483187in}}%
\pgfpathlineto{\pgfqpoint{4.312711in}{1.515678in}}%
\pgfpathlineto{\pgfqpoint{4.312838in}{1.510255in}}%
\pgfpathlineto{\pgfqpoint{4.313636in}{1.573588in}}%
\pgfpathlineto{\pgfqpoint{4.313651in}{1.570951in}}%
\pgfpathlineto{\pgfqpoint{4.313694in}{1.577438in}}%
\pgfpathlineto{\pgfqpoint{4.314619in}{1.550002in}}%
\pgfpathlineto{\pgfqpoint{4.315243in}{1.520137in}}%
\pgfpathlineto{\pgfqpoint{4.315632in}{1.559057in}}%
\pgfpathlineto{\pgfqpoint{4.315690in}{1.554636in}}%
\pgfpathlineto{\pgfqpoint{4.316839in}{1.594892in}}%
\pgfpathlineto{\pgfqpoint{4.316859in}{1.589599in}}%
\pgfpathlineto{\pgfqpoint{4.317015in}{1.582591in}}%
\pgfpathlineto{\pgfqpoint{4.317282in}{1.602905in}}%
\pgfpathlineto{\pgfqpoint{4.317984in}{1.585807in}}%
\pgfpathlineto{\pgfqpoint{4.318047in}{1.596621in}}%
\pgfpathlineto{\pgfqpoint{4.318996in}{1.568328in}}%
\pgfpathlineto{\pgfqpoint{4.319089in}{1.581987in}}%
\pgfpathlineto{\pgfqpoint{4.319989in}{1.563795in}}%
\pgfpathlineto{\pgfqpoint{4.319449in}{1.588319in}}%
\pgfpathlineto{\pgfqpoint{4.320140in}{1.583353in}}%
\pgfpathlineto{\pgfqpoint{4.320749in}{1.579502in}}%
\pgfpathlineto{\pgfqpoint{4.321343in}{1.639055in}}%
\pgfpathlineto{\pgfqpoint{4.322229in}{1.603203in}}%
\pgfpathlineto{\pgfqpoint{4.321508in}{1.649587in}}%
\pgfpathlineto{\pgfqpoint{4.322652in}{1.629045in}}%
\pgfpathlineto{\pgfqpoint{4.322920in}{1.639385in}}%
\pgfpathlineto{\pgfqpoint{4.323222in}{1.613947in}}%
\pgfpathlineto{\pgfqpoint{4.323767in}{1.635079in}}%
\pgfpathlineto{\pgfqpoint{4.323792in}{1.629706in}}%
\pgfpathlineto{\pgfqpoint{4.324064in}{1.672801in}}%
\pgfpathlineto{\pgfqpoint{4.324824in}{1.642322in}}%
\pgfpathlineto{\pgfqpoint{4.325973in}{1.693257in}}%
\pgfpathlineto{\pgfqpoint{4.326021in}{1.689378in}}%
\pgfpathlineto{\pgfqpoint{4.326202in}{1.679633in}}%
\pgfpathlineto{\pgfqpoint{4.326839in}{1.712213in}}%
\pgfpathlineto{\pgfqpoint{4.326966in}{1.717566in}}%
\pgfpathlineto{\pgfqpoint{4.327901in}{1.690209in}}%
\pgfpathlineto{\pgfqpoint{4.327920in}{1.695781in}}%
\pgfpathlineto{\pgfqpoint{4.327954in}{1.691034in}}%
\pgfpathlineto{\pgfqpoint{4.328534in}{1.727636in}}%
\pgfpathlineto{\pgfqpoint{4.329006in}{1.704603in}}%
\pgfpathlineto{\pgfqpoint{4.330160in}{1.731548in}}%
\pgfpathlineto{\pgfqpoint{4.329016in}{1.703979in}}%
\pgfpathlineto{\pgfqpoint{4.330203in}{1.719421in}}%
\pgfpathlineto{\pgfqpoint{4.330642in}{1.708011in}}%
\pgfpathlineto{\pgfqpoint{4.330817in}{1.725422in}}%
\pgfpathlineto{\pgfqpoint{4.331333in}{1.709065in}}%
\pgfpathlineto{\pgfqpoint{4.331508in}{1.730539in}}%
\pgfpathlineto{\pgfqpoint{4.331372in}{1.708471in}}%
\pgfpathlineto{\pgfqpoint{4.332497in}{1.725015in}}%
\pgfpathlineto{\pgfqpoint{4.332828in}{1.712450in}}%
\pgfpathlineto{\pgfqpoint{4.333159in}{1.736438in}}%
\pgfpathlineto{\pgfqpoint{4.333164in}{1.735636in}}%
\pgfpathlineto{\pgfqpoint{4.333237in}{1.746722in}}%
\pgfpathlineto{\pgfqpoint{4.333675in}{1.722223in}}%
\pgfpathlineto{\pgfqpoint{4.334254in}{1.729750in}}%
\pgfpathlineto{\pgfqpoint{4.334741in}{1.708636in}}%
\pgfpathlineto{\pgfqpoint{4.335111in}{1.761158in}}%
\pgfpathlineto{\pgfqpoint{4.335218in}{1.755591in}}%
\pgfpathlineto{\pgfqpoint{4.336226in}{1.760494in}}%
\pgfpathlineto{\pgfqpoint{4.335914in}{1.732820in}}%
\pgfpathlineto{\pgfqpoint{4.336323in}{1.756659in}}%
\pgfpathlineto{\pgfqpoint{4.336328in}{1.753550in}}%
\pgfpathlineto{\pgfqpoint{4.337034in}{1.818966in}}%
\pgfpathlineto{\pgfqpoint{4.337370in}{1.793728in}}%
\pgfpathlineto{\pgfqpoint{4.337375in}{1.795356in}}%
\pgfpathlineto{\pgfqpoint{4.338017in}{1.747540in}}%
\pgfpathlineto{\pgfqpoint{4.338412in}{1.771092in}}%
\pgfpathlineto{\pgfqpoint{4.338947in}{1.741454in}}%
\pgfpathlineto{\pgfqpoint{4.338441in}{1.776537in}}%
\pgfpathlineto{\pgfqpoint{4.339575in}{1.755739in}}%
\pgfpathlineto{\pgfqpoint{4.340773in}{1.789212in}}%
\pgfpathlineto{\pgfqpoint{4.340057in}{1.748239in}}%
\pgfpathlineto{\pgfqpoint{4.340778in}{1.788458in}}%
\pgfpathlineto{\pgfqpoint{4.340846in}{1.777302in}}%
\pgfpathlineto{\pgfqpoint{4.341620in}{1.809196in}}%
\pgfpathlineto{\pgfqpoint{4.341873in}{1.795966in}}%
\pgfpathlineto{\pgfqpoint{4.341878in}{1.794999in}}%
\pgfpathlineto{\pgfqpoint{4.342604in}{1.846057in}}%
\pgfpathlineto{\pgfqpoint{4.342769in}{1.837651in}}%
\pgfpathlineto{\pgfqpoint{4.343266in}{1.860836in}}%
\pgfpathlineto{\pgfqpoint{4.342896in}{1.826329in}}%
\pgfpathlineto{\pgfqpoint{4.343879in}{1.838952in}}%
\pgfpathlineto{\pgfqpoint{4.344196in}{1.855222in}}%
\pgfpathlineto{\pgfqpoint{4.345082in}{1.815188in}}%
\pgfpathlineto{\pgfqpoint{4.345087in}{1.818497in}}%
\pgfpathlineto{\pgfqpoint{4.346133in}{1.782370in}}%
\pgfpathlineto{\pgfqpoint{4.346158in}{1.792311in}}%
\pgfpathlineto{\pgfqpoint{4.347165in}{1.780554in}}%
\pgfpathlineto{\pgfqpoint{4.346781in}{1.810965in}}%
\pgfpathlineto{\pgfqpoint{4.347248in}{1.794592in}}%
\pgfpathlineto{\pgfqpoint{4.348120in}{1.823867in}}%
\pgfpathlineto{\pgfqpoint{4.347555in}{1.780811in}}%
\pgfpathlineto{\pgfqpoint{4.348368in}{1.798075in}}%
\pgfpathlineto{\pgfqpoint{4.348499in}{1.785954in}}%
\pgfpathlineto{\pgfqpoint{4.349264in}{1.827977in}}%
\pgfpathlineto{\pgfqpoint{4.349327in}{1.823285in}}%
\pgfpathlineto{\pgfqpoint{4.349673in}{1.859124in}}%
\pgfpathlineto{\pgfqpoint{4.350306in}{1.813873in}}%
\pgfpathlineto{\pgfqpoint{4.350447in}{1.830349in}}%
\pgfpathlineto{\pgfqpoint{4.351396in}{1.835309in}}%
\pgfpathlineto{\pgfqpoint{4.351172in}{1.807491in}}%
\pgfpathlineto{\pgfqpoint{4.351489in}{1.815366in}}%
\pgfpathlineto{\pgfqpoint{4.352414in}{1.775912in}}%
\pgfpathlineto{\pgfqpoint{4.352638in}{1.791759in}}%
\pgfpathlineto{\pgfqpoint{4.353441in}{1.806997in}}%
\pgfpathlineto{\pgfqpoint{4.353105in}{1.775302in}}%
\pgfpathlineto{\pgfqpoint{4.353709in}{1.786726in}}%
\pgfpathlineto{\pgfqpoint{4.353850in}{1.773968in}}%
\pgfpathlineto{\pgfqpoint{4.354244in}{1.803549in}}%
\pgfpathlineto{\pgfqpoint{4.354819in}{1.786451in}}%
\pgfpathlineto{\pgfqpoint{4.355267in}{1.807180in}}%
\pgfpathlineto{\pgfqpoint{4.355671in}{1.768709in}}%
\pgfpathlineto{\pgfqpoint{4.355919in}{1.784594in}}%
\pgfpathlineto{\pgfqpoint{4.356708in}{1.741526in}}%
\pgfpathlineto{\pgfqpoint{4.357092in}{1.769107in}}%
\pgfpathlineto{\pgfqpoint{4.357121in}{1.775482in}}%
\pgfpathlineto{\pgfqpoint{4.358066in}{1.752786in}}%
\pgfpathlineto{\pgfqpoint{4.358139in}{1.758296in}}%
\pgfpathlineto{\pgfqpoint{4.358991in}{1.720705in}}%
\pgfpathlineto{\pgfqpoint{4.359278in}{1.746251in}}%
\pgfpathlineto{\pgfqpoint{4.360836in}{1.677796in}}%
\pgfpathlineto{\pgfqpoint{4.361703in}{1.700426in}}%
\pgfpathlineto{\pgfqpoint{4.361917in}{1.706680in}}%
\pgfpathlineto{\pgfqpoint{4.362399in}{1.687154in}}%
\pgfpathlineto{\pgfqpoint{4.362827in}{1.705425in}}%
\pgfpathlineto{\pgfqpoint{4.363192in}{1.693438in}}%
\pgfpathlineto{\pgfqpoint{4.363728in}{1.738561in}}%
\pgfpathlineto{\pgfqpoint{4.363913in}{1.714400in}}%
\pgfpathlineto{\pgfqpoint{4.364215in}{1.737159in}}%
\pgfpathlineto{\pgfqpoint{4.365013in}{1.712394in}}%
\pgfpathlineto{\pgfqpoint{4.365023in}{1.716155in}}%
\pgfpathlineto{\pgfqpoint{4.365719in}{1.665410in}}%
\pgfpathlineto{\pgfqpoint{4.365091in}{1.720854in}}%
\pgfpathlineto{\pgfqpoint{4.366201in}{1.693635in}}%
\pgfpathlineto{\pgfqpoint{4.366552in}{1.697801in}}%
\pgfpathlineto{\pgfqpoint{4.367126in}{1.670818in}}%
\pgfpathlineto{\pgfqpoint{4.367146in}{1.673881in}}%
\pgfpathlineto{\pgfqpoint{4.367394in}{1.664538in}}%
\pgfpathlineto{\pgfqpoint{4.367905in}{1.697787in}}%
\pgfpathlineto{\pgfqpoint{4.368227in}{1.693917in}}%
\pgfpathlineto{\pgfqpoint{4.368869in}{1.681486in}}%
\pgfpathlineto{\pgfqpoint{4.369205in}{1.703059in}}%
\pgfpathlineto{\pgfqpoint{4.369332in}{1.694340in}}%
\pgfpathlineto{\pgfqpoint{4.369361in}{1.701178in}}%
\pgfpathlineto{\pgfqpoint{4.369556in}{1.718134in}}%
\pgfpathlineto{\pgfqpoint{4.370271in}{1.689148in}}%
\pgfpathlineto{\pgfqpoint{4.370437in}{1.689790in}}%
\pgfpathlineto{\pgfqpoint{4.370568in}{1.676462in}}%
\pgfpathlineto{\pgfqpoint{4.371152in}{1.716108in}}%
\pgfpathlineto{\pgfqpoint{4.371518in}{1.704579in}}%
\pgfpathlineto{\pgfqpoint{4.372277in}{1.752136in}}%
\pgfpathlineto{\pgfqpoint{4.371600in}{1.698011in}}%
\pgfpathlineto{\pgfqpoint{4.372876in}{1.741286in}}%
\pgfpathlineto{\pgfqpoint{4.373222in}{1.733757in}}%
\pgfpathlineto{\pgfqpoint{4.373942in}{1.768566in}}%
\pgfpathlineto{\pgfqpoint{4.374010in}{1.776534in}}%
\pgfpathlineto{\pgfqpoint{4.374595in}{1.731043in}}%
\pgfpathlineto{\pgfqpoint{4.374935in}{1.746316in}}%
\pgfpathlineto{\pgfqpoint{4.375855in}{1.723074in}}%
\pgfpathlineto{\pgfqpoint{4.375315in}{1.750109in}}%
\pgfpathlineto{\pgfqpoint{4.376065in}{1.735464in}}%
\pgfpathlineto{\pgfqpoint{4.376323in}{1.763292in}}%
\pgfpathlineto{\pgfqpoint{4.377034in}{1.732261in}}%
\pgfpathlineto{\pgfqpoint{4.377189in}{1.749099in}}%
\pgfpathlineto{\pgfqpoint{4.377764in}{1.753995in}}%
\pgfpathlineto{\pgfqpoint{4.378032in}{1.730876in}}%
\pgfpathlineto{\pgfqpoint{4.378124in}{1.737188in}}%
\pgfpathlineto{\pgfqpoint{4.378153in}{1.732107in}}%
\pgfpathlineto{\pgfqpoint{4.379088in}{1.783746in}}%
\pgfpathlineto{\pgfqpoint{4.379151in}{1.773915in}}%
\pgfpathlineto{\pgfqpoint{4.380154in}{1.803778in}}%
\pgfpathlineto{\pgfqpoint{4.379482in}{1.748909in}}%
\pgfpathlineto{\pgfqpoint{4.380320in}{1.786967in}}%
\pgfpathlineto{\pgfqpoint{4.381425in}{1.781440in}}%
\pgfpathlineto{\pgfqpoint{4.380524in}{1.804230in}}%
\pgfpathlineto{\pgfqpoint{4.381435in}{1.783056in}}%
\pgfpathlineto{\pgfqpoint{4.381537in}{1.794529in}}%
\pgfpathlineto{\pgfqpoint{4.382472in}{1.734880in}}%
\pgfpathlineto{\pgfqpoint{4.382496in}{1.736231in}}%
\pgfpathlineto{\pgfqpoint{4.383533in}{1.714646in}}%
\pgfpathlineto{\pgfqpoint{4.382817in}{1.742061in}}%
\pgfpathlineto{\pgfqpoint{4.383665in}{1.729121in}}%
\pgfpathlineto{\pgfqpoint{4.384570in}{1.747924in}}%
\pgfpathlineto{\pgfqpoint{4.384103in}{1.716848in}}%
\pgfpathlineto{\pgfqpoint{4.384775in}{1.730172in}}%
\pgfpathlineto{\pgfqpoint{4.385597in}{1.743031in}}%
\pgfpathlineto{\pgfqpoint{4.385101in}{1.709337in}}%
\pgfpathlineto{\pgfqpoint{4.385899in}{1.735107in}}%
\pgfpathlineto{\pgfqpoint{4.386001in}{1.727127in}}%
\pgfpathlineto{\pgfqpoint{4.386849in}{1.803520in}}%
\pgfpathlineto{\pgfqpoint{4.386863in}{1.801520in}}%
\pgfpathlineto{\pgfqpoint{4.387540in}{1.842778in}}%
\pgfpathlineto{\pgfqpoint{4.387842in}{1.821045in}}%
\pgfpathlineto{\pgfqpoint{4.387993in}{1.828424in}}%
\pgfpathlineto{\pgfqpoint{4.388786in}{1.776152in}}%
\pgfpathlineto{\pgfqpoint{4.388845in}{1.777515in}}%
\pgfpathlineto{\pgfqpoint{4.389083in}{1.753196in}}%
\pgfpathlineto{\pgfqpoint{4.389366in}{1.791970in}}%
\pgfpathlineto{\pgfqpoint{4.389969in}{1.763699in}}%
\pgfpathlineto{\pgfqpoint{4.391084in}{1.857005in}}%
\pgfpathlineto{\pgfqpoint{4.391206in}{1.844499in}}%
\pgfpathlineto{\pgfqpoint{4.391810in}{1.794635in}}%
\pgfpathlineto{\pgfqpoint{4.391289in}{1.844868in}}%
\pgfpathlineto{\pgfqpoint{4.392335in}{1.832409in}}%
\pgfpathlineto{\pgfqpoint{4.392360in}{1.844408in}}%
\pgfpathlineto{\pgfqpoint{4.392666in}{1.798581in}}%
\pgfpathlineto{\pgfqpoint{4.393426in}{1.831974in}}%
\pgfpathlineto{\pgfqpoint{4.394151in}{1.804270in}}%
\pgfpathlineto{\pgfqpoint{4.393645in}{1.842207in}}%
\pgfpathlineto{\pgfqpoint{4.394536in}{1.827471in}}%
\pgfpathlineto{\pgfqpoint{4.395665in}{1.878168in}}%
\pgfpathlineto{\pgfqpoint{4.394624in}{1.824981in}}%
\pgfpathlineto{\pgfqpoint{4.395768in}{1.873418in}}%
\pgfpathlineto{\pgfqpoint{4.396965in}{1.794723in}}%
\pgfpathlineto{\pgfqpoint{4.397014in}{1.801370in}}%
\pgfpathlineto{\pgfqpoint{4.398061in}{1.823041in}}%
\pgfpathlineto{\pgfqpoint{4.397228in}{1.788396in}}%
\pgfpathlineto{\pgfqpoint{4.398139in}{1.805224in}}%
\pgfpathlineto{\pgfqpoint{4.398470in}{1.781745in}}%
\pgfpathlineto{\pgfqpoint{4.398221in}{1.815402in}}%
\pgfpathlineto{\pgfqpoint{4.399292in}{1.787386in}}%
\pgfpathlineto{\pgfqpoint{4.400164in}{1.797225in}}%
\pgfpathlineto{\pgfqpoint{4.399886in}{1.758218in}}%
\pgfpathlineto{\pgfqpoint{4.400373in}{1.781905in}}%
\pgfpathlineto{\pgfqpoint{4.400553in}{1.797467in}}%
\pgfpathlineto{\pgfqpoint{4.400422in}{1.775944in}}%
\pgfpathlineto{\pgfqpoint{4.401099in}{1.776289in}}%
\pgfpathlineto{\pgfqpoint{4.401376in}{1.748202in}}%
\pgfpathlineto{\pgfqpoint{4.401629in}{1.781488in}}%
\pgfpathlineto{\pgfqpoint{4.402218in}{1.767823in}}%
\pgfpathlineto{\pgfqpoint{4.403348in}{1.792513in}}%
\pgfpathlineto{\pgfqpoint{4.402910in}{1.751468in}}%
\pgfpathlineto{\pgfqpoint{4.403372in}{1.791234in}}%
\pgfpathlineto{\pgfqpoint{4.404161in}{1.779743in}}%
\pgfpathlineto{\pgfqpoint{4.404307in}{1.799213in}}%
\pgfpathlineto{\pgfqpoint{4.404458in}{1.797810in}}%
\pgfpathlineto{\pgfqpoint{4.404468in}{1.801952in}}%
\pgfpathlineto{\pgfqpoint{4.405427in}{1.733565in}}%
\pgfpathlineto{\pgfqpoint{4.406430in}{1.693750in}}%
\pgfpathlineto{\pgfqpoint{4.406581in}{1.702491in}}%
\pgfpathlineto{\pgfqpoint{4.406610in}{1.701119in}}%
\pgfpathlineto{\pgfqpoint{4.407131in}{1.741464in}}%
\pgfpathlineto{\pgfqpoint{4.407549in}{1.734267in}}%
\pgfpathlineto{\pgfqpoint{4.407807in}{1.715163in}}%
\pgfpathlineto{\pgfqpoint{4.408494in}{1.750043in}}%
\pgfpathlineto{\pgfqpoint{4.408669in}{1.728696in}}%
\pgfpathlineto{\pgfqpoint{4.409068in}{1.760747in}}%
\pgfpathlineto{\pgfqpoint{4.409370in}{1.726936in}}%
\pgfpathlineto{\pgfqpoint{4.409784in}{1.735618in}}%
\pgfpathlineto{\pgfqpoint{4.410003in}{1.718803in}}%
\pgfpathlineto{\pgfqpoint{4.410446in}{1.752775in}}%
\pgfpathlineto{\pgfqpoint{4.410904in}{1.733069in}}%
\pgfpathlineto{\pgfqpoint{4.412184in}{1.769746in}}%
\pgfpathlineto{\pgfqpoint{4.411542in}{1.731711in}}%
\pgfpathlineto{\pgfqpoint{4.412257in}{1.753070in}}%
\pgfpathlineto{\pgfqpoint{4.413435in}{1.687243in}}%
\pgfpathlineto{\pgfqpoint{4.413445in}{1.689514in}}%
\pgfpathlineto{\pgfqpoint{4.413723in}{1.646342in}}%
\pgfpathlineto{\pgfqpoint{4.414253in}{1.671008in}}%
\pgfpathlineto{\pgfqpoint{4.414784in}{1.619589in}}%
\pgfpathlineto{\pgfqpoint{4.415490in}{1.636602in}}%
\pgfpathlineto{\pgfqpoint{4.415539in}{1.629793in}}%
\pgfpathlineto{\pgfqpoint{4.416293in}{1.677946in}}%
\pgfpathlineto{\pgfqpoint{4.416366in}{1.669932in}}%
\pgfpathlineto{\pgfqpoint{4.417428in}{1.703924in}}%
\pgfpathlineto{\pgfqpoint{4.417544in}{1.695331in}}%
\pgfpathlineto{\pgfqpoint{4.418430in}{1.689785in}}%
\pgfpathlineto{\pgfqpoint{4.417851in}{1.729427in}}%
\pgfpathlineto{\pgfqpoint{4.418640in}{1.697283in}}%
\pgfpathlineto{\pgfqpoint{4.419156in}{1.744078in}}%
\pgfpathlineto{\pgfqpoint{4.418732in}{1.695825in}}%
\pgfpathlineto{\pgfqpoint{4.419769in}{1.708823in}}%
\pgfpathlineto{\pgfqpoint{4.419779in}{1.707478in}}%
\pgfpathlineto{\pgfqpoint{4.420665in}{1.736347in}}%
\pgfpathlineto{\pgfqpoint{4.420728in}{1.735958in}}%
\pgfpathlineto{\pgfqpoint{4.420952in}{1.709658in}}%
\pgfpathlineto{\pgfqpoint{4.421692in}{1.744494in}}%
\pgfpathlineto{\pgfqpoint{4.421926in}{1.727259in}}%
\pgfpathlineto{\pgfqpoint{4.422953in}{1.732026in}}%
\pgfpathlineto{\pgfqpoint{4.422267in}{1.700442in}}%
\pgfpathlineto{\pgfqpoint{4.422997in}{1.721050in}}%
\pgfpathlineto{\pgfqpoint{4.423903in}{1.678696in}}%
\pgfpathlineto{\pgfqpoint{4.423343in}{1.728454in}}%
\pgfpathlineto{\pgfqpoint{4.424136in}{1.691559in}}%
\pgfpathlineto{\pgfqpoint{4.425280in}{1.743947in}}%
\pgfpathlineto{\pgfqpoint{4.424209in}{1.690664in}}%
\pgfpathlineto{\pgfqpoint{4.425339in}{1.732297in}}%
\pgfpathlineto{\pgfqpoint{4.426152in}{1.725146in}}%
\pgfpathlineto{\pgfqpoint{4.425573in}{1.742809in}}%
\pgfpathlineto{\pgfqpoint{4.426434in}{1.739381in}}%
\pgfpathlineto{\pgfqpoint{4.426819in}{1.716689in}}%
\pgfpathlineto{\pgfqpoint{4.427213in}{1.751784in}}%
\pgfpathlineto{\pgfqpoint{4.427354in}{1.749032in}}%
\pgfpathlineto{\pgfqpoint{4.428440in}{1.774528in}}%
\pgfpathlineto{\pgfqpoint{4.428109in}{1.744010in}}%
\pgfpathlineto{\pgfqpoint{4.428499in}{1.772932in}}%
\pgfpathlineto{\pgfqpoint{4.428503in}{1.773030in}}%
\pgfpathlineto{\pgfqpoint{4.428757in}{1.757735in}}%
\pgfpathlineto{\pgfqpoint{4.428820in}{1.752936in}}%
\pgfpathlineto{\pgfqpoint{4.429122in}{1.784684in}}%
\pgfpathlineto{\pgfqpoint{4.429823in}{1.781478in}}%
\pgfpathlineto{\pgfqpoint{4.430504in}{1.808323in}}%
\pgfpathlineto{\pgfqpoint{4.430066in}{1.771638in}}%
\pgfpathlineto{\pgfqpoint{4.430957in}{1.793398in}}%
\pgfpathlineto{\pgfqpoint{4.430991in}{1.788341in}}%
\pgfpathlineto{\pgfqpoint{4.431892in}{1.827472in}}%
\pgfpathlineto{\pgfqpoint{4.431902in}{1.826016in}}%
\pgfpathlineto{\pgfqpoint{4.432257in}{1.838257in}}%
\pgfpathlineto{\pgfqpoint{4.432471in}{1.822294in}}%
\pgfpathlineto{\pgfqpoint{4.433124in}{1.803196in}}%
\pgfpathlineto{\pgfqpoint{4.433430in}{1.825403in}}%
\pgfpathlineto{\pgfqpoint{4.433645in}{1.806804in}}%
\pgfpathlineto{\pgfqpoint{4.433649in}{1.809812in}}%
\pgfpathlineto{\pgfqpoint{4.434652in}{1.750729in}}%
\pgfpathlineto{\pgfqpoint{4.434667in}{1.751978in}}%
\pgfpathlineto{\pgfqpoint{4.435188in}{1.733247in}}%
\pgfpathlineto{\pgfqpoint{4.435597in}{1.783364in}}%
\pgfpathlineto{\pgfqpoint{4.435733in}{1.772495in}}%
\pgfpathlineto{\pgfqpoint{4.436434in}{1.811107in}}%
\pgfpathlineto{\pgfqpoint{4.435762in}{1.766570in}}%
\pgfpathlineto{\pgfqpoint{4.436950in}{1.796098in}}%
\pgfpathlineto{\pgfqpoint{4.437505in}{1.787504in}}%
\pgfpathlineto{\pgfqpoint{4.437788in}{1.821988in}}%
\pgfpathlineto{\pgfqpoint{4.437817in}{1.821767in}}%
\pgfpathlineto{\pgfqpoint{4.438591in}{1.858893in}}%
\pgfpathlineto{\pgfqpoint{4.439146in}{1.851488in}}%
\pgfpathlineto{\pgfqpoint{4.439170in}{1.849039in}}%
\pgfpathlineto{\pgfqpoint{4.439793in}{1.872950in}}%
\pgfpathlineto{\pgfqpoint{4.440217in}{1.862555in}}%
\pgfpathlineto{\pgfqpoint{4.440563in}{1.895982in}}%
\pgfpathlineto{\pgfqpoint{4.441210in}{1.853596in}}%
\pgfpathlineto{\pgfqpoint{4.441244in}{1.857535in}}%
\pgfpathlineto{\pgfqpoint{4.442242in}{1.822391in}}%
\pgfpathlineto{\pgfqpoint{4.442384in}{1.831466in}}%
\pgfpathlineto{\pgfqpoint{4.442500in}{1.824774in}}%
\pgfpathlineto{\pgfqpoint{4.443615in}{1.903245in}}%
\pgfpathlineto{\pgfqpoint{4.444433in}{1.917483in}}%
\pgfpathlineto{\pgfqpoint{4.444005in}{1.876577in}}%
\pgfpathlineto{\pgfqpoint{4.444711in}{1.902364in}}%
\pgfpathlineto{\pgfqpoint{4.444716in}{1.898368in}}%
\pgfpathlineto{\pgfqpoint{4.445324in}{1.931747in}}%
\pgfpathlineto{\pgfqpoint{4.445796in}{1.919260in}}%
\pgfpathlineto{\pgfqpoint{4.446906in}{1.969419in}}%
\pgfpathlineto{\pgfqpoint{4.445952in}{1.910848in}}%
\pgfpathlineto{\pgfqpoint{4.446950in}{1.960249in}}%
\pgfpathlineto{\pgfqpoint{4.446955in}{1.960091in}}%
\pgfpathlineto{\pgfqpoint{4.447218in}{1.984693in}}%
\pgfpathlineto{\pgfqpoint{4.447233in}{1.984160in}}%
\pgfpathlineto{\pgfqpoint{4.447817in}{2.018418in}}%
\pgfpathlineto{\pgfqpoint{4.448158in}{1.983946in}}%
\pgfpathlineto{\pgfqpoint{4.448372in}{2.011117in}}%
\pgfpathlineto{\pgfqpoint{4.449268in}{2.105190in}}%
\pgfpathlineto{\pgfqpoint{4.448401in}{2.010252in}}%
\pgfpathlineto{\pgfqpoint{4.449696in}{2.076398in}}%
\pgfpathlineto{\pgfqpoint{4.449813in}{2.074779in}}%
\pgfpathlineto{\pgfqpoint{4.450095in}{2.101793in}}%
\pgfpathlineto{\pgfqpoint{4.450115in}{2.100728in}}%
\pgfpathlineto{\pgfqpoint{4.450120in}{2.103992in}}%
\pgfpathlineto{\pgfqpoint{4.450913in}{2.063986in}}%
\pgfpathlineto{\pgfqpoint{4.451210in}{2.093816in}}%
\pgfpathlineto{\pgfqpoint{4.452106in}{2.106883in}}%
\pgfpathlineto{\pgfqpoint{4.452398in}{2.067989in}}%
\pgfpathlineto{\pgfqpoint{4.453542in}{2.114946in}}%
\pgfpathlineto{\pgfqpoint{4.452418in}{2.066309in}}%
\pgfpathlineto{\pgfqpoint{4.453664in}{2.110384in}}%
\pgfpathlineto{\pgfqpoint{4.453834in}{2.095181in}}%
\pgfpathlineto{\pgfqpoint{4.454633in}{2.115502in}}%
\pgfpathlineto{\pgfqpoint{4.454750in}{2.113008in}}%
\pgfpathlineto{\pgfqpoint{4.454759in}{2.114272in}}%
\pgfpathlineto{\pgfqpoint{4.455193in}{2.065236in}}%
\pgfpathlineto{\pgfqpoint{4.455689in}{2.070273in}}%
\pgfpathlineto{\pgfqpoint{4.456025in}{2.056525in}}%
\pgfpathlineto{\pgfqpoint{4.456726in}{2.078884in}}%
\pgfpathlineto{\pgfqpoint{4.456794in}{2.072172in}}%
\pgfpathlineto{\pgfqpoint{4.457598in}{2.081895in}}%
\pgfpathlineto{\pgfqpoint{4.457344in}{2.062751in}}%
\pgfpathlineto{\pgfqpoint{4.457880in}{2.067183in}}%
\pgfpathlineto{\pgfqpoint{4.457885in}{2.066705in}}%
\pgfpathlineto{\pgfqpoint{4.458518in}{2.112571in}}%
\pgfpathlineto{\pgfqpoint{4.458557in}{2.108929in}}%
\pgfpathlineto{\pgfqpoint{4.459258in}{2.144100in}}%
\pgfpathlineto{\pgfqpoint{4.458566in}{2.108159in}}%
\pgfpathlineto{\pgfqpoint{4.459711in}{2.117615in}}%
\pgfpathlineto{\pgfqpoint{4.460864in}{2.083744in}}%
\pgfpathlineto{\pgfqpoint{4.459920in}{2.119988in}}%
\pgfpathlineto{\pgfqpoint{4.460933in}{2.088238in}}%
\pgfpathlineto{\pgfqpoint{4.460967in}{2.084422in}}%
\pgfpathlineto{\pgfqpoint{4.461556in}{2.128315in}}%
\pgfpathlineto{\pgfqpoint{4.461590in}{2.126945in}}%
\pgfpathlineto{\pgfqpoint{4.462320in}{2.124660in}}%
\pgfpathlineto{\pgfqpoint{4.462773in}{2.166163in}}%
\pgfpathlineto{\pgfqpoint{4.464000in}{2.134848in}}%
\pgfpathlineto{\pgfqpoint{4.463055in}{2.186858in}}%
\pgfpathlineto{\pgfqpoint{4.464024in}{2.137850in}}%
\pgfpathlineto{\pgfqpoint{4.464710in}{2.120034in}}%
\pgfpathlineto{\pgfqpoint{4.465149in}{2.152265in}}%
\pgfpathlineto{\pgfqpoint{4.465699in}{2.116274in}}%
\pgfpathlineto{\pgfqpoint{4.465163in}{2.153105in}}%
\pgfpathlineto{\pgfqpoint{4.466351in}{2.130046in}}%
\pgfpathlineto{\pgfqpoint{4.467505in}{2.162734in}}%
\pgfpathlineto{\pgfqpoint{4.466522in}{2.122329in}}%
\pgfpathlineto{\pgfqpoint{4.467563in}{2.159072in}}%
\pgfpathlineto{\pgfqpoint{4.468196in}{2.138184in}}%
\pgfpathlineto{\pgfqpoint{4.468586in}{2.170172in}}%
\pgfpathlineto{\pgfqpoint{4.468654in}{2.165797in}}%
\pgfpathlineto{\pgfqpoint{4.469764in}{2.213841in}}%
\pgfpathlineto{\pgfqpoint{4.468727in}{2.160206in}}%
\pgfpathlineto{\pgfqpoint{4.469861in}{2.208135in}}%
\pgfpathlineto{\pgfqpoint{4.470548in}{2.182858in}}%
\pgfpathlineto{\pgfqpoint{4.470869in}{2.215298in}}%
\pgfpathlineto{\pgfqpoint{4.470976in}{2.205762in}}%
\pgfpathlineto{\pgfqpoint{4.471400in}{2.193641in}}%
\pgfpathlineto{\pgfqpoint{4.471624in}{2.219736in}}%
\pgfpathlineto{\pgfqpoint{4.471745in}{2.242326in}}%
\pgfpathlineto{\pgfqpoint{4.472495in}{2.206474in}}%
\pgfpathlineto{\pgfqpoint{4.472719in}{2.215480in}}%
\pgfpathlineto{\pgfqpoint{4.474214in}{2.272005in}}%
\pgfpathlineto{\pgfqpoint{4.474258in}{2.265809in}}%
\pgfpathlineto{\pgfqpoint{4.475324in}{2.253007in}}%
\pgfpathlineto{\pgfqpoint{4.474457in}{2.282547in}}%
\pgfpathlineto{\pgfqpoint{4.475382in}{2.258665in}}%
\pgfpathlineto{\pgfqpoint{4.476560in}{2.303526in}}%
\pgfpathlineto{\pgfqpoint{4.475747in}{2.249501in}}%
\pgfpathlineto{\pgfqpoint{4.476711in}{2.299081in}}%
\pgfpathlineto{\pgfqpoint{4.477437in}{2.276541in}}%
\pgfpathlineto{\pgfqpoint{4.477208in}{2.308604in}}%
\pgfpathlineto{\pgfqpoint{4.477831in}{2.294127in}}%
\pgfpathlineto{\pgfqpoint{4.477860in}{2.296774in}}%
\pgfpathlineto{\pgfqpoint{4.478581in}{2.243250in}}%
\pgfpathlineto{\pgfqpoint{4.478790in}{2.254233in}}%
\pgfpathlineto{\pgfqpoint{4.479603in}{2.297900in}}%
\pgfpathlineto{\pgfqpoint{4.480475in}{2.277836in}}%
\pgfpathlineto{\pgfqpoint{4.481512in}{2.205622in}}%
\pgfpathlineto{\pgfqpoint{4.481667in}{2.216388in}}%
\pgfpathlineto{\pgfqpoint{4.482417in}{2.249652in}}%
\pgfpathlineto{\pgfqpoint{4.481692in}{2.215828in}}%
\pgfpathlineto{\pgfqpoint{4.482855in}{2.245030in}}%
\pgfpathlineto{\pgfqpoint{4.483965in}{2.205604in}}%
\pgfpathlineto{\pgfqpoint{4.484063in}{2.208636in}}%
\pgfpathlineto{\pgfqpoint{4.484175in}{2.219724in}}%
\pgfpathlineto{\pgfqpoint{4.485007in}{2.173750in}}%
\pgfpathlineto{\pgfqpoint{4.485046in}{2.181058in}}%
\pgfpathlineto{\pgfqpoint{4.486010in}{2.153049in}}%
\pgfpathlineto{\pgfqpoint{4.485231in}{2.191225in}}%
\pgfpathlineto{\pgfqpoint{4.486200in}{2.157712in}}%
\pgfpathlineto{\pgfqpoint{4.487037in}{2.162550in}}%
\pgfpathlineto{\pgfqpoint{4.486779in}{2.143749in}}%
\pgfpathlineto{\pgfqpoint{4.487295in}{2.150335in}}%
\pgfpathlineto{\pgfqpoint{4.487763in}{2.141486in}}%
\pgfpathlineto{\pgfqpoint{4.488323in}{2.168800in}}%
\pgfpathlineto{\pgfqpoint{4.488347in}{2.166773in}}%
\pgfpathlineto{\pgfqpoint{4.489472in}{2.193923in}}%
\pgfpathlineto{\pgfqpoint{4.488922in}{2.166754in}}%
\pgfpathlineto{\pgfqpoint{4.489501in}{2.192355in}}%
\pgfpathlineto{\pgfqpoint{4.489515in}{2.197304in}}%
\pgfpathlineto{\pgfqpoint{4.490392in}{2.155110in}}%
\pgfpathlineto{\pgfqpoint{4.490514in}{2.168264in}}%
\pgfpathlineto{\pgfqpoint{4.491069in}{2.151031in}}%
\pgfpathlineto{\pgfqpoint{4.491555in}{2.182735in}}%
\pgfpathlineto{\pgfqpoint{4.491599in}{2.179179in}}%
\pgfpathlineto{\pgfqpoint{4.492295in}{2.203148in}}%
\pgfpathlineto{\pgfqpoint{4.491624in}{2.176631in}}%
\pgfpathlineto{\pgfqpoint{4.492768in}{2.199992in}}%
\pgfpathlineto{\pgfqpoint{4.493858in}{2.188083in}}%
\pgfpathlineto{\pgfqpoint{4.493245in}{2.213961in}}%
\pgfpathlineto{\pgfqpoint{4.493926in}{2.197438in}}%
\pgfpathlineto{\pgfqpoint{4.495075in}{2.230289in}}%
\pgfpathlineto{\pgfqpoint{4.494063in}{2.194752in}}%
\pgfpathlineto{\pgfqpoint{4.495095in}{2.225364in}}%
\pgfpathlineto{\pgfqpoint{4.496273in}{2.192732in}}%
\pgfpathlineto{\pgfqpoint{4.495143in}{2.231144in}}%
\pgfpathlineto{\pgfqpoint{4.496292in}{2.194918in}}%
\pgfpathlineto{\pgfqpoint{4.497003in}{2.234840in}}%
\pgfpathlineto{\pgfqpoint{4.497437in}{2.199208in}}%
\pgfpathlineto{\pgfqpoint{4.498435in}{2.187712in}}%
\pgfpathlineto{\pgfqpoint{4.497758in}{2.210703in}}%
\pgfpathlineto{\pgfqpoint{4.498537in}{2.203083in}}%
\pgfpathlineto{\pgfqpoint{4.499082in}{2.228446in}}%
\pgfpathlineto{\pgfqpoint{4.499306in}{2.196180in}}%
\pgfpathlineto{\pgfqpoint{4.499681in}{2.205056in}}%
\pgfpathlineto{\pgfqpoint{4.500124in}{2.200970in}}%
\pgfpathlineto{\pgfqpoint{4.500562in}{2.228415in}}%
\pgfpathlineto{\pgfqpoint{4.500747in}{2.263077in}}%
\pgfpathlineto{\pgfqpoint{4.501278in}{2.220061in}}%
\pgfpathlineto{\pgfqpoint{4.501692in}{2.239315in}}%
\pgfpathlineto{\pgfqpoint{4.502300in}{2.258251in}}%
\pgfpathlineto{\pgfqpoint{4.502729in}{2.229580in}}%
\pgfpathlineto{\pgfqpoint{4.502743in}{2.230293in}}%
\pgfpathlineto{\pgfqpoint{4.503425in}{2.219659in}}%
\pgfpathlineto{\pgfqpoint{4.503259in}{2.244857in}}%
\pgfpathlineto{\pgfqpoint{4.503853in}{2.230725in}}%
\pgfpathlineto{\pgfqpoint{4.503858in}{2.233717in}}%
\pgfpathlineto{\pgfqpoint{4.504248in}{2.194474in}}%
\pgfpathlineto{\pgfqpoint{4.504954in}{2.226413in}}%
\pgfpathlineto{\pgfqpoint{4.505835in}{2.200037in}}%
\pgfpathlineto{\pgfqpoint{4.505314in}{2.235439in}}%
\pgfpathlineto{\pgfqpoint{4.506083in}{2.217344in}}%
\pgfpathlineto{\pgfqpoint{4.506088in}{2.219112in}}%
\pgfpathlineto{\pgfqpoint{4.506760in}{2.171777in}}%
\pgfpathlineto{\pgfqpoint{4.507149in}{2.201118in}}%
\pgfpathlineto{\pgfqpoint{4.507154in}{2.198367in}}%
\pgfpathlineto{\pgfqpoint{4.508050in}{2.229591in}}%
\pgfpathlineto{\pgfqpoint{4.508215in}{2.225415in}}%
\pgfpathlineto{\pgfqpoint{4.508371in}{2.237144in}}%
\pgfpathlineto{\pgfqpoint{4.509004in}{2.205224in}}%
\pgfpathlineto{\pgfqpoint{4.509301in}{2.208909in}}%
\pgfpathlineto{\pgfqpoint{4.509364in}{2.216530in}}%
\pgfpathlineto{\pgfqpoint{4.509588in}{2.188454in}}%
\pgfpathlineto{\pgfqpoint{4.509793in}{2.189787in}}%
\pgfpathlineto{\pgfqpoint{4.510431in}{2.148769in}}%
\pgfpathlineto{\pgfqpoint{4.510927in}{2.169405in}}%
\pgfpathlineto{\pgfqpoint{4.511117in}{2.188553in}}%
\pgfpathlineto{\pgfqpoint{4.511623in}{2.162794in}}%
\pgfpathlineto{\pgfqpoint{4.511828in}{2.166532in}}%
\pgfpathlineto{\pgfqpoint{4.512227in}{2.135369in}}%
\pgfpathlineto{\pgfqpoint{4.512987in}{2.138652in}}%
\pgfpathlineto{\pgfqpoint{4.513459in}{2.173364in}}%
\pgfpathlineto{\pgfqpoint{4.513186in}{2.130866in}}%
\pgfpathlineto{\pgfqpoint{4.514101in}{2.142410in}}%
\pgfpathlineto{\pgfqpoint{4.514934in}{2.158962in}}%
\pgfpathlineto{\pgfqpoint{4.514535in}{2.133352in}}%
\pgfpathlineto{\pgfqpoint{4.515246in}{2.153128in}}%
\pgfpathlineto{\pgfqpoint{4.515591in}{2.197034in}}%
\pgfpathlineto{\pgfqpoint{4.516468in}{2.168734in}}%
\pgfpathlineto{\pgfqpoint{4.516959in}{2.158015in}}%
\pgfpathlineto{\pgfqpoint{4.517378in}{2.191409in}}%
\pgfpathlineto{\pgfqpoint{4.517402in}{2.190246in}}%
\pgfpathlineto{\pgfqpoint{4.518702in}{2.220022in}}%
\pgfpathlineto{\pgfqpoint{4.517982in}{2.177447in}}%
\pgfpathlineto{\pgfqpoint{4.518712in}{2.217316in}}%
\pgfpathlineto{\pgfqpoint{4.519890in}{2.191707in}}%
\pgfpathlineto{\pgfqpoint{4.519082in}{2.230725in}}%
\pgfpathlineto{\pgfqpoint{4.519900in}{2.193524in}}%
\pgfpathlineto{\pgfqpoint{4.520844in}{2.180376in}}%
\pgfpathlineto{\pgfqpoint{4.520620in}{2.206873in}}%
\pgfpathlineto{\pgfqpoint{4.521020in}{2.193061in}}%
\pgfpathlineto{\pgfqpoint{4.521097in}{2.204623in}}%
\pgfpathlineto{\pgfqpoint{4.521726in}{2.158833in}}%
\pgfpathlineto{\pgfqpoint{4.522115in}{2.187175in}}%
\pgfpathlineto{\pgfqpoint{4.522694in}{2.134811in}}%
\pgfpathlineto{\pgfqpoint{4.523274in}{2.136425in}}%
\pgfpathlineto{\pgfqpoint{4.523629in}{2.152862in}}%
\pgfpathlineto{\pgfqpoint{4.523371in}{2.128581in}}%
\pgfpathlineto{\pgfqpoint{4.524374in}{2.135193in}}%
\pgfpathlineto{\pgfqpoint{4.524408in}{2.131251in}}%
\pgfpathlineto{\pgfqpoint{4.525168in}{2.176773in}}%
\pgfpathlineto{\pgfqpoint{4.525435in}{2.153493in}}%
\pgfpathlineto{\pgfqpoint{4.526560in}{2.214619in}}%
\pgfpathlineto{\pgfqpoint{4.526789in}{2.210306in}}%
\pgfpathlineto{\pgfqpoint{4.527465in}{2.198992in}}%
\pgfpathlineto{\pgfqpoint{4.527047in}{2.232132in}}%
\pgfpathlineto{\pgfqpoint{4.527899in}{2.208654in}}%
\pgfpathlineto{\pgfqpoint{4.527923in}{2.210501in}}%
\pgfpathlineto{\pgfqpoint{4.528537in}{2.179139in}}%
\pgfpathlineto{\pgfqpoint{4.528799in}{2.192270in}}%
\pgfpathlineto{\pgfqpoint{4.529413in}{2.158684in}}%
\pgfpathlineto{\pgfqpoint{4.529939in}{2.170157in}}%
\pgfpathlineto{\pgfqpoint{4.531068in}{2.142337in}}%
\pgfpathlineto{\pgfqpoint{4.531102in}{2.146266in}}%
\pgfpathlineto{\pgfqpoint{4.531424in}{2.134876in}}%
\pgfpathlineto{\pgfqpoint{4.532115in}{2.171966in}}%
\pgfpathlineto{\pgfqpoint{4.532129in}{2.168882in}}%
\pgfpathlineto{\pgfqpoint{4.533400in}{2.219540in}}%
\pgfpathlineto{\pgfqpoint{4.532232in}{2.159378in}}%
\pgfpathlineto{\pgfqpoint{4.533405in}{2.217255in}}%
\pgfpathlineto{\pgfqpoint{4.533556in}{2.204785in}}%
\pgfpathlineto{\pgfqpoint{4.534072in}{2.229492in}}%
\pgfpathlineto{\pgfqpoint{4.534481in}{2.225716in}}%
\pgfpathlineto{\pgfqpoint{4.534573in}{2.241257in}}%
\pgfpathlineto{\pgfqpoint{4.535270in}{2.184838in}}%
\pgfpathlineto{\pgfqpoint{4.535513in}{2.199017in}}%
\pgfpathlineto{\pgfqpoint{4.536716in}{2.152065in}}%
\pgfpathlineto{\pgfqpoint{4.536832in}{2.160978in}}%
\pgfpathlineto{\pgfqpoint{4.537719in}{2.192127in}}%
\pgfpathlineto{\pgfqpoint{4.537967in}{2.174702in}}%
\pgfpathlineto{\pgfqpoint{4.538906in}{2.150940in}}%
\pgfpathlineto{\pgfqpoint{4.538084in}{2.176836in}}%
\pgfpathlineto{\pgfqpoint{4.539077in}{2.172158in}}%
\pgfpathlineto{\pgfqpoint{4.539646in}{2.183870in}}%
\pgfpathlineto{\pgfqpoint{4.539997in}{2.153482in}}%
\pgfpathlineto{\pgfqpoint{4.540163in}{2.164054in}}%
\pgfpathlineto{\pgfqpoint{4.540187in}{2.162306in}}%
\pgfpathlineto{\pgfqpoint{4.540313in}{2.183803in}}%
\pgfpathlineto{\pgfqpoint{4.541672in}{2.274007in}}%
\pgfpathlineto{\pgfqpoint{4.540377in}{2.180371in}}%
\pgfpathlineto{\pgfqpoint{4.541711in}{2.268248in}}%
\pgfpathlineto{\pgfqpoint{4.542451in}{2.300683in}}%
\pgfpathlineto{\pgfqpoint{4.541823in}{2.266700in}}%
\pgfpathlineto{\pgfqpoint{4.542962in}{2.276070in}}%
\pgfpathlineto{\pgfqpoint{4.542972in}{2.276199in}}%
\pgfpathlineto{\pgfqpoint{4.544228in}{2.217991in}}%
\pgfpathlineto{\pgfqpoint{4.544846in}{2.236139in}}%
\pgfpathlineto{\pgfqpoint{4.544340in}{2.206034in}}%
\pgfpathlineto{\pgfqpoint{4.545338in}{2.217747in}}%
\pgfpathlineto{\pgfqpoint{4.546433in}{2.201858in}}%
\pgfpathlineto{\pgfqpoint{4.545966in}{2.244784in}}%
\pgfpathlineto{\pgfqpoint{4.546462in}{2.206717in}}%
\pgfpathlineto{\pgfqpoint{4.547105in}{2.175082in}}%
\pgfpathlineto{\pgfqpoint{4.546511in}{2.212385in}}%
\pgfpathlineto{\pgfqpoint{4.547577in}{2.205601in}}%
\pgfpathlineto{\pgfqpoint{4.548361in}{2.245315in}}%
\pgfpathlineto{\pgfqpoint{4.547684in}{2.196575in}}%
\pgfpathlineto{\pgfqpoint{4.548809in}{2.236909in}}%
\pgfpathlineto{\pgfqpoint{4.549437in}{2.194495in}}%
\pgfpathlineto{\pgfqpoint{4.549938in}{2.221793in}}%
\pgfpathlineto{\pgfqpoint{4.549953in}{2.221254in}}%
\pgfpathlineto{\pgfqpoint{4.549973in}{2.228516in}}%
\pgfpathlineto{\pgfqpoint{4.550995in}{2.243195in}}%
\pgfpathlineto{\pgfqpoint{4.550148in}{2.217796in}}%
\pgfpathlineto{\pgfqpoint{4.551087in}{2.231807in}}%
\pgfpathlineto{\pgfqpoint{4.551292in}{2.230210in}}%
\pgfpathlineto{\pgfqpoint{4.551930in}{2.264680in}}%
\pgfpathlineto{\pgfqpoint{4.552076in}{2.257476in}}%
\pgfpathlineto{\pgfqpoint{4.553395in}{2.291397in}}%
\pgfpathlineto{\pgfqpoint{4.552197in}{2.247749in}}%
\pgfpathlineto{\pgfqpoint{4.553507in}{2.281848in}}%
\pgfpathlineto{\pgfqpoint{4.554724in}{2.179301in}}%
\pgfpathlineto{\pgfqpoint{4.554768in}{2.181184in}}%
\pgfpathlineto{\pgfqpoint{4.555109in}{2.204365in}}%
\pgfpathlineto{\pgfqpoint{4.555849in}{2.170008in}}%
\pgfpathlineto{\pgfqpoint{4.555932in}{2.171502in}}%
\pgfpathlineto{\pgfqpoint{4.556530in}{2.148165in}}%
\pgfpathlineto{\pgfqpoint{4.556866in}{2.159370in}}%
\pgfpathlineto{\pgfqpoint{4.557154in}{2.120686in}}%
\pgfpathlineto{\pgfqpoint{4.557976in}{2.155615in}}%
\pgfpathlineto{\pgfqpoint{4.558303in}{2.172145in}}%
\pgfpathlineto{\pgfqpoint{4.558541in}{2.145646in}}%
\pgfpathlineto{\pgfqpoint{4.559082in}{2.155472in}}%
\pgfpathlineto{\pgfqpoint{4.559972in}{2.129772in}}%
\pgfpathlineto{\pgfqpoint{4.559164in}{2.163382in}}%
\pgfpathlineto{\pgfqpoint{4.560255in}{2.146862in}}%
\pgfpathlineto{\pgfqpoint{4.561433in}{2.174147in}}%
\pgfpathlineto{\pgfqpoint{4.560484in}{2.141271in}}%
\pgfpathlineto{\pgfqpoint{4.561443in}{2.173519in}}%
\pgfpathlineto{\pgfqpoint{4.561487in}{2.170112in}}%
\pgfpathlineto{\pgfqpoint{4.561822in}{2.201693in}}%
\pgfpathlineto{\pgfqpoint{4.562504in}{2.185391in}}%
\pgfpathlineto{\pgfqpoint{4.562791in}{2.190964in}}%
\pgfpathlineto{\pgfqpoint{4.563181in}{2.153030in}}%
\pgfpathlineto{\pgfqpoint{4.563546in}{2.166268in}}%
\pgfpathlineto{\pgfqpoint{4.564257in}{2.135507in}}%
\pgfpathlineto{\pgfqpoint{4.564666in}{2.153798in}}%
\pgfpathlineto{\pgfqpoint{4.565114in}{2.184893in}}%
\pgfpathlineto{\pgfqpoint{4.565790in}{2.165717in}}%
\pgfpathlineto{\pgfqpoint{4.566964in}{2.194674in}}%
\pgfpathlineto{\pgfqpoint{4.566379in}{2.153574in}}%
\pgfpathlineto{\pgfqpoint{4.567022in}{2.193902in}}%
\pgfpathlineto{\pgfqpoint{4.567309in}{2.178226in}}%
\pgfpathlineto{\pgfqpoint{4.567207in}{2.200078in}}%
\pgfpathlineto{\pgfqpoint{4.568127in}{2.194070in}}%
\pgfpathlineto{\pgfqpoint{4.568775in}{2.222174in}}%
\pgfpathlineto{\pgfqpoint{4.568307in}{2.175945in}}%
\pgfpathlineto{\pgfqpoint{4.569262in}{2.204789in}}%
\pgfpathlineto{\pgfqpoint{4.570011in}{2.182976in}}%
\pgfpathlineto{\pgfqpoint{4.570376in}{2.199247in}}%
\pgfpathlineto{\pgfqpoint{4.571457in}{2.228094in}}%
\pgfpathlineto{\pgfqpoint{4.571116in}{2.189419in}}%
\pgfpathlineto{\pgfqpoint{4.571559in}{2.226572in}}%
\pgfpathlineto{\pgfqpoint{4.571735in}{2.217944in}}%
\pgfpathlineto{\pgfqpoint{4.572402in}{2.261616in}}%
\pgfpathlineto{\pgfqpoint{4.572592in}{2.252556in}}%
\pgfpathlineto{\pgfqpoint{4.572718in}{2.264793in}}%
\pgfpathlineto{\pgfqpoint{4.572962in}{2.230539in}}%
\pgfpathlineto{\pgfqpoint{4.573682in}{2.242296in}}%
\pgfpathlineto{\pgfqpoint{4.573726in}{2.249228in}}%
\pgfpathlineto{\pgfqpoint{4.574544in}{2.225951in}}%
\pgfpathlineto{\pgfqpoint{4.574700in}{2.232083in}}%
\pgfpathlineto{\pgfqpoint{4.574705in}{2.228593in}}%
\pgfpathlineto{\pgfqpoint{4.575644in}{2.261320in}}%
\pgfpathlineto{\pgfqpoint{4.575795in}{2.242610in}}%
\pgfpathlineto{\pgfqpoint{4.576744in}{2.263192in}}%
\pgfpathlineto{\pgfqpoint{4.575873in}{2.237657in}}%
\pgfpathlineto{\pgfqpoint{4.576964in}{2.259015in}}%
\pgfpathlineto{\pgfqpoint{4.578010in}{2.221482in}}%
\pgfpathlineto{\pgfqpoint{4.577455in}{2.261335in}}%
\pgfpathlineto{\pgfqpoint{4.578122in}{2.233716in}}%
\pgfpathlineto{\pgfqpoint{4.578156in}{2.240179in}}%
\pgfpathlineto{\pgfqpoint{4.578882in}{2.173332in}}%
\pgfpathlineto{\pgfqpoint{4.578964in}{2.176656in}}%
\pgfpathlineto{\pgfqpoint{4.579086in}{2.159423in}}%
\pgfpathlineto{\pgfqpoint{4.579281in}{2.178881in}}%
\pgfpathlineto{\pgfqpoint{4.580070in}{2.175682in}}%
\pgfpathlineto{\pgfqpoint{4.580561in}{2.189867in}}%
\pgfpathlineto{\pgfqpoint{4.580824in}{2.164171in}}%
\pgfpathlineto{\pgfqpoint{4.581185in}{2.178614in}}%
\pgfpathlineto{\pgfqpoint{4.581608in}{2.170153in}}%
\pgfpathlineto{\pgfqpoint{4.582022in}{2.203995in}}%
\pgfpathlineto{\pgfqpoint{4.582256in}{2.185318in}}%
\pgfpathlineto{\pgfqpoint{4.582986in}{2.196674in}}%
\pgfpathlineto{\pgfqpoint{4.582616in}{2.157248in}}%
\pgfpathlineto{\pgfqpoint{4.583380in}{2.191171in}}%
\pgfpathlineto{\pgfqpoint{4.584218in}{2.159347in}}%
\pgfpathlineto{\pgfqpoint{4.583560in}{2.204586in}}%
\pgfpathlineto{\pgfqpoint{4.584500in}{2.183740in}}%
\pgfpathlineto{\pgfqpoint{4.585644in}{2.224760in}}%
\pgfpathlineto{\pgfqpoint{4.584690in}{2.171674in}}%
\pgfpathlineto{\pgfqpoint{4.585766in}{2.218827in}}%
\pgfpathlineto{\pgfqpoint{4.587251in}{2.119235in}}%
\pgfpathlineto{\pgfqpoint{4.587275in}{2.119302in}}%
\pgfpathlineto{\pgfqpoint{4.587518in}{2.139012in}}%
\pgfpathlineto{\pgfqpoint{4.588317in}{2.110440in}}%
\pgfpathlineto{\pgfqpoint{4.588346in}{2.114254in}}%
\pgfpathlineto{\pgfqpoint{4.588458in}{2.102806in}}%
\pgfpathlineto{\pgfqpoint{4.588998in}{2.125807in}}%
\pgfpathlineto{\pgfqpoint{4.589471in}{2.111122in}}%
\pgfpathlineto{\pgfqpoint{4.589490in}{2.110643in}}%
\pgfpathlineto{\pgfqpoint{4.589510in}{2.112575in}}%
\pgfpathlineto{\pgfqpoint{4.589519in}{2.111189in}}%
\pgfpathlineto{\pgfqpoint{4.590264in}{2.171578in}}%
\pgfpathlineto{\pgfqpoint{4.589563in}{2.103286in}}%
\pgfpathlineto{\pgfqpoint{4.590688in}{2.138128in}}%
\pgfpathlineto{\pgfqpoint{4.591355in}{2.105993in}}%
\pgfpathlineto{\pgfqpoint{4.590814in}{2.139118in}}%
\pgfpathlineto{\pgfqpoint{4.591822in}{2.124780in}}%
\pgfpathlineto{\pgfqpoint{4.591827in}{2.126472in}}%
\pgfpathlineto{\pgfqpoint{4.592158in}{2.098087in}}%
\pgfpathlineto{\pgfqpoint{4.592913in}{2.119479in}}%
\pgfpathlineto{\pgfqpoint{4.594135in}{2.092328in}}%
\pgfpathlineto{\pgfqpoint{4.593458in}{2.129717in}}%
\pgfpathlineto{\pgfqpoint{4.594140in}{2.094972in}}%
\pgfpathlineto{\pgfqpoint{4.594914in}{2.129768in}}%
\pgfpathlineto{\pgfqpoint{4.594154in}{2.091912in}}%
\pgfpathlineto{\pgfqpoint{4.595391in}{2.101797in}}%
\pgfpathlineto{\pgfqpoint{4.596525in}{2.078375in}}%
\pgfpathlineto{\pgfqpoint{4.596160in}{2.110349in}}%
\pgfpathlineto{\pgfqpoint{4.596579in}{2.081249in}}%
\pgfpathlineto{\pgfqpoint{4.596715in}{2.093092in}}%
\pgfpathlineto{\pgfqpoint{4.596603in}{2.072792in}}%
\pgfpathlineto{\pgfqpoint{4.597616in}{2.079354in}}%
\pgfpathlineto{\pgfqpoint{4.598521in}{2.046407in}}%
\pgfpathlineto{\pgfqpoint{4.598137in}{2.089647in}}%
\pgfpathlineto{\pgfqpoint{4.598740in}{2.070592in}}%
\pgfpathlineto{\pgfqpoint{4.598750in}{2.072742in}}%
\pgfpathlineto{\pgfqpoint{4.599023in}{2.037161in}}%
\pgfpathlineto{\pgfqpoint{4.599826in}{2.060800in}}%
\pgfpathlineto{\pgfqpoint{4.600172in}{2.046584in}}%
\pgfpathlineto{\pgfqpoint{4.599977in}{2.074202in}}%
\pgfpathlineto{\pgfqpoint{4.600605in}{2.066811in}}%
\pgfpathlineto{\pgfqpoint{4.601029in}{2.098102in}}%
\pgfpathlineto{\pgfqpoint{4.601705in}{2.066374in}}%
\pgfpathlineto{\pgfqpoint{4.601715in}{2.068174in}}%
\pgfpathlineto{\pgfqpoint{4.602815in}{2.102457in}}%
\pgfpathlineto{\pgfqpoint{4.601885in}{2.046110in}}%
\pgfpathlineto{\pgfqpoint{4.603005in}{2.085954in}}%
\pgfpathlineto{\pgfqpoint{4.603852in}{2.067616in}}%
\pgfpathlineto{\pgfqpoint{4.603258in}{2.093205in}}%
\pgfpathlineto{\pgfqpoint{4.604135in}{2.076698in}}%
\pgfpathlineto{\pgfqpoint{4.605323in}{2.129482in}}%
\pgfpathlineto{\pgfqpoint{4.605508in}{2.138961in}}%
\pgfpathlineto{\pgfqpoint{4.606515in}{2.155744in}}%
\pgfpathlineto{\pgfqpoint{4.606588in}{2.135304in}}%
\pgfpathlineto{\pgfqpoint{4.606618in}{2.140554in}}%
\pgfpathlineto{\pgfqpoint{4.607844in}{2.047987in}}%
\pgfpathlineto{\pgfqpoint{4.606866in}{2.146227in}}%
\pgfpathlineto{\pgfqpoint{4.607981in}{2.056417in}}%
\pgfpathlineto{\pgfqpoint{4.609052in}{2.088244in}}%
\pgfpathlineto{\pgfqpoint{4.608351in}{2.039623in}}%
\pgfpathlineto{\pgfqpoint{4.609139in}{2.076276in}}%
\pgfpathlineto{\pgfqpoint{4.610269in}{2.052014in}}%
\pgfpathlineto{\pgfqpoint{4.609432in}{2.089413in}}%
\pgfpathlineto{\pgfqpoint{4.610313in}{2.059866in}}%
\pgfpathlineto{\pgfqpoint{4.610712in}{2.066035in}}%
\pgfpathlineto{\pgfqpoint{4.611282in}{2.039499in}}%
\pgfpathlineto{\pgfqpoint{4.611384in}{2.047516in}}%
\pgfpathlineto{\pgfqpoint{4.611394in}{2.044207in}}%
\pgfpathlineto{\pgfqpoint{4.612392in}{2.073845in}}%
\pgfpathlineto{\pgfqpoint{4.612455in}{2.068627in}}%
\pgfpathlineto{\pgfqpoint{4.613307in}{2.087669in}}%
\pgfpathlineto{\pgfqpoint{4.612650in}{2.062220in}}%
\pgfpathlineto{\pgfqpoint{4.613648in}{2.085409in}}%
\pgfpathlineto{\pgfqpoint{4.614617in}{2.052975in}}%
\pgfpathlineto{\pgfqpoint{4.613735in}{2.087873in}}%
\pgfpathlineto{\pgfqpoint{4.614792in}{2.064400in}}%
\pgfpathlineto{\pgfqpoint{4.615264in}{2.044932in}}%
\pgfpathlineto{\pgfqpoint{4.615410in}{2.068562in}}%
\pgfpathlineto{\pgfqpoint{4.615907in}{2.061674in}}%
\pgfpathlineto{\pgfqpoint{4.616788in}{2.078455in}}%
\pgfpathlineto{\pgfqpoint{4.616087in}{2.058630in}}%
\pgfpathlineto{\pgfqpoint{4.617051in}{2.069677in}}%
\pgfpathlineto{\pgfqpoint{4.617085in}{2.057023in}}%
\pgfpathlineto{\pgfqpoint{4.618029in}{2.094931in}}%
\pgfpathlineto{\pgfqpoint{4.618073in}{2.088523in}}%
\pgfpathlineto{\pgfqpoint{4.618248in}{2.098621in}}%
\pgfpathlineto{\pgfqpoint{4.618404in}{2.077777in}}%
\pgfpathlineto{\pgfqpoint{4.619159in}{2.081436in}}%
\pgfpathlineto{\pgfqpoint{4.619188in}{2.082295in}}%
\pgfpathlineto{\pgfqpoint{4.620507in}{2.008680in}}%
\pgfpathlineto{\pgfqpoint{4.621233in}{2.020559in}}%
\pgfpathlineto{\pgfqpoint{4.620941in}{2.003575in}}%
\pgfpathlineto{\pgfqpoint{4.621632in}{2.016008in}}%
\pgfpathlineto{\pgfqpoint{4.621944in}{1.999338in}}%
\pgfpathlineto{\pgfqpoint{4.622572in}{2.055288in}}%
\pgfpathlineto{\pgfqpoint{4.623609in}{2.077958in}}%
\pgfpathlineto{\pgfqpoint{4.622878in}{2.042598in}}%
\pgfpathlineto{\pgfqpoint{4.623711in}{2.066092in}}%
\pgfpathlineto{\pgfqpoint{4.624022in}{2.082718in}}%
\pgfpathlineto{\pgfqpoint{4.624315in}{2.055804in}}%
\pgfpathlineto{\pgfqpoint{4.624806in}{2.062265in}}%
\pgfpathlineto{\pgfqpoint{4.625507in}{2.035171in}}%
\pgfpathlineto{\pgfqpoint{4.625157in}{2.073345in}}%
\pgfpathlineto{\pgfqpoint{4.625926in}{2.054648in}}%
\pgfpathlineto{\pgfqpoint{4.626048in}{2.076969in}}%
\pgfpathlineto{\pgfqpoint{4.626924in}{2.025231in}}%
\pgfpathlineto{\pgfqpoint{4.626929in}{2.025954in}}%
\pgfpathlineto{\pgfqpoint{4.628122in}{2.062790in}}%
\pgfpathlineto{\pgfqpoint{4.627046in}{2.004215in}}%
\pgfpathlineto{\pgfqpoint{4.628180in}{2.060724in}}%
\pgfpathlineto{\pgfqpoint{4.628263in}{2.052911in}}%
\pgfpathlineto{\pgfqpoint{4.629164in}{2.090307in}}%
\pgfpathlineto{\pgfqpoint{4.629227in}{2.086887in}}%
\pgfpathlineto{\pgfqpoint{4.629568in}{2.098469in}}%
\pgfpathlineto{\pgfqpoint{4.629957in}{2.057781in}}%
\pgfpathlineto{\pgfqpoint{4.630283in}{2.067089in}}%
\pgfpathlineto{\pgfqpoint{4.630692in}{2.058821in}}%
\pgfpathlineto{\pgfqpoint{4.631111in}{2.101058in}}%
\pgfpathlineto{\pgfqpoint{4.631340in}{2.081938in}}%
\pgfpathlineto{\pgfqpoint{4.631875in}{2.108688in}}%
\pgfpathlineto{\pgfqpoint{4.632284in}{2.059411in}}%
\pgfpathlineto{\pgfqpoint{4.632426in}{2.072773in}}%
\pgfpathlineto{\pgfqpoint{4.633448in}{2.112301in}}%
\pgfpathlineto{\pgfqpoint{4.632572in}{2.061680in}}%
\pgfpathlineto{\pgfqpoint{4.633662in}{2.090690in}}%
\pgfpathlineto{\pgfqpoint{4.634485in}{2.068922in}}%
\pgfpathlineto{\pgfqpoint{4.634081in}{2.094102in}}%
\pgfpathlineto{\pgfqpoint{4.634777in}{2.088192in}}%
\pgfpathlineto{\pgfqpoint{4.634947in}{2.095427in}}%
\pgfpathlineto{\pgfqpoint{4.635269in}{2.056461in}}%
\pgfpathlineto{\pgfqpoint{4.635814in}{2.061960in}}%
\pgfpathlineto{\pgfqpoint{4.636793in}{2.045108in}}%
\pgfpathlineto{\pgfqpoint{4.636238in}{2.084322in}}%
\pgfpathlineto{\pgfqpoint{4.636939in}{2.054403in}}%
\pgfpathlineto{\pgfqpoint{4.638190in}{2.096035in}}%
\pgfpathlineto{\pgfqpoint{4.636968in}{2.051116in}}%
\pgfpathlineto{\pgfqpoint{4.638219in}{2.092015in}}%
\pgfpathlineto{\pgfqpoint{4.638239in}{2.088932in}}%
\pgfpathlineto{\pgfqpoint{4.638959in}{2.133988in}}%
\pgfpathlineto{\pgfqpoint{4.639275in}{2.108581in}}%
\pgfpathlineto{\pgfqpoint{4.639295in}{2.110289in}}%
\pgfpathlineto{\pgfqpoint{4.639680in}{2.076384in}}%
\pgfpathlineto{\pgfqpoint{4.640278in}{2.089031in}}%
\pgfpathlineto{\pgfqpoint{4.640527in}{2.076181in}}%
\pgfpathlineto{\pgfqpoint{4.641160in}{2.100359in}}%
\pgfpathlineto{\pgfqpoint{4.641374in}{2.094593in}}%
\pgfpathlineto{\pgfqpoint{4.642581in}{2.159656in}}%
\pgfpathlineto{\pgfqpoint{4.641457in}{2.083071in}}%
\pgfpathlineto{\pgfqpoint{4.642630in}{2.157081in}}%
\pgfpathlineto{\pgfqpoint{4.642820in}{2.144281in}}%
\pgfpathlineto{\pgfqpoint{4.643088in}{2.170819in}}%
\pgfpathlineto{\pgfqpoint{4.643389in}{2.167323in}}%
\pgfpathlineto{\pgfqpoint{4.643759in}{2.183061in}}%
\pgfpathlineto{\pgfqpoint{4.644310in}{2.156594in}}%
\pgfpathlineto{\pgfqpoint{4.644490in}{2.163712in}}%
\pgfpathlineto{\pgfqpoint{4.645678in}{2.138790in}}%
\pgfpathlineto{\pgfqpoint{4.644918in}{2.182943in}}%
\pgfpathlineto{\pgfqpoint{4.645682in}{2.141196in}}%
\pgfpathlineto{\pgfqpoint{4.645755in}{2.151667in}}%
\pgfpathlineto{\pgfqpoint{4.646194in}{2.119060in}}%
\pgfpathlineto{\pgfqpoint{4.646792in}{2.141943in}}%
\pgfpathlineto{\pgfqpoint{4.647756in}{2.114717in}}%
\pgfpathlineto{\pgfqpoint{4.647016in}{2.153726in}}%
\pgfpathlineto{\pgfqpoint{4.648165in}{2.130499in}}%
\pgfpathlineto{\pgfqpoint{4.648623in}{2.166274in}}%
\pgfpathlineto{\pgfqpoint{4.649305in}{2.148323in}}%
\pgfpathlineto{\pgfqpoint{4.650424in}{2.164317in}}%
\pgfpathlineto{\pgfqpoint{4.650093in}{2.137879in}}%
\pgfpathlineto{\pgfqpoint{4.650488in}{2.160233in}}%
\pgfpathlineto{\pgfqpoint{4.650828in}{2.147434in}}%
\pgfpathlineto{\pgfqpoint{4.651549in}{2.191160in}}%
\pgfpathlineto{\pgfqpoint{4.652250in}{2.213600in}}%
\pgfpathlineto{\pgfqpoint{4.652518in}{2.180322in}}%
\pgfpathlineto{\pgfqpoint{4.652649in}{2.185312in}}%
\pgfpathlineto{\pgfqpoint{4.653360in}{2.191074in}}%
\pgfpathlineto{\pgfqpoint{4.653633in}{2.154671in}}%
\pgfpathlineto{\pgfqpoint{4.653652in}{2.150576in}}%
\pgfpathlineto{\pgfqpoint{4.654392in}{2.191283in}}%
\pgfpathlineto{\pgfqpoint{4.654660in}{2.182184in}}%
\pgfpathlineto{\pgfqpoint{4.655181in}{2.212451in}}%
\pgfpathlineto{\pgfqpoint{4.654840in}{2.168692in}}%
\pgfpathlineto{\pgfqpoint{4.655765in}{2.179805in}}%
\pgfpathlineto{\pgfqpoint{4.655794in}{2.176494in}}%
\pgfpathlineto{\pgfqpoint{4.656271in}{2.201662in}}%
\pgfpathlineto{\pgfqpoint{4.656836in}{2.188511in}}%
\pgfpathlineto{\pgfqpoint{4.656841in}{2.190937in}}%
\pgfpathlineto{\pgfqpoint{4.657669in}{2.153344in}}%
\pgfpathlineto{\pgfqpoint{4.657898in}{2.166884in}}%
\pgfpathlineto{\pgfqpoint{4.658638in}{2.127157in}}%
\pgfpathlineto{\pgfqpoint{4.658097in}{2.174723in}}%
\pgfpathlineto{\pgfqpoint{4.659139in}{2.134492in}}%
\pgfpathlineto{\pgfqpoint{4.660264in}{2.157953in}}%
\pgfpathlineto{\pgfqpoint{4.659499in}{2.134110in}}%
\pgfpathlineto{\pgfqpoint{4.660283in}{2.151445in}}%
\pgfpathlineto{\pgfqpoint{4.660911in}{2.143844in}}%
\pgfpathlineto{\pgfqpoint{4.661232in}{2.173573in}}%
\pgfpathlineto{\pgfqpoint{4.661242in}{2.169731in}}%
\pgfpathlineto{\pgfqpoint{4.661286in}{2.175493in}}%
\pgfpathlineto{\pgfqpoint{4.661817in}{2.128053in}}%
\pgfpathlineto{\pgfqpoint{4.662318in}{2.140323in}}%
\pgfpathlineto{\pgfqpoint{4.663370in}{2.167168in}}%
\pgfpathlineto{\pgfqpoint{4.662367in}{2.137582in}}%
\pgfpathlineto{\pgfqpoint{4.663521in}{2.163628in}}%
\pgfpathlineto{\pgfqpoint{4.663667in}{2.154651in}}%
\pgfpathlineto{\pgfqpoint{4.664207in}{2.198355in}}%
\pgfpathlineto{\pgfqpoint{4.664582in}{2.186489in}}%
\pgfpathlineto{\pgfqpoint{4.664626in}{2.192985in}}%
\pgfpathlineto{\pgfqpoint{4.664937in}{2.158357in}}%
\pgfpathlineto{\pgfqpoint{4.665677in}{2.178960in}}%
\pgfpathlineto{\pgfqpoint{4.665809in}{2.174737in}}%
\pgfpathlineto{\pgfqpoint{4.666554in}{2.216476in}}%
\pgfpathlineto{\pgfqpoint{4.666724in}{2.209952in}}%
\pgfpathlineto{\pgfqpoint{4.666997in}{2.237755in}}%
\pgfpathlineto{\pgfqpoint{4.667795in}{2.208841in}}%
\pgfpathlineto{\pgfqpoint{4.667839in}{2.211711in}}%
\pgfpathlineto{\pgfqpoint{4.668146in}{2.186031in}}%
\pgfpathlineto{\pgfqpoint{4.667893in}{2.214295in}}%
\pgfpathlineto{\pgfqpoint{4.668973in}{2.199708in}}%
\pgfpathlineto{\pgfqpoint{4.670025in}{2.206102in}}%
\pgfpathlineto{\pgfqpoint{4.669348in}{2.163066in}}%
\pgfpathlineto{\pgfqpoint{4.670079in}{2.198508in}}%
\pgfpathlineto{\pgfqpoint{4.670268in}{2.187082in}}%
\pgfpathlineto{\pgfqpoint{4.670882in}{2.236239in}}%
\pgfpathlineto{\pgfqpoint{4.671047in}{2.229796in}}%
\pgfpathlineto{\pgfqpoint{4.671086in}{2.233676in}}%
\pgfpathlineto{\pgfqpoint{4.671335in}{2.196832in}}%
\pgfpathlineto{\pgfqpoint{4.671977in}{2.211409in}}%
\pgfpathlineto{\pgfqpoint{4.672313in}{2.189331in}}%
\pgfpathlineto{\pgfqpoint{4.672766in}{2.232366in}}%
\pgfpathlineto{\pgfqpoint{4.673043in}{2.223236in}}%
\pgfpathlineto{\pgfqpoint{4.673691in}{2.244083in}}%
\pgfpathlineto{\pgfqpoint{4.673331in}{2.206020in}}%
\pgfpathlineto{\pgfqpoint{4.674158in}{2.227570in}}%
\pgfpathlineto{\pgfqpoint{4.674207in}{2.217801in}}%
\pgfpathlineto{\pgfqpoint{4.675005in}{2.243890in}}%
\pgfpathlineto{\pgfqpoint{4.675263in}{2.230772in}}%
\pgfpathlineto{\pgfqpoint{4.675653in}{2.238234in}}%
\pgfpathlineto{\pgfqpoint{4.676213in}{2.186114in}}%
\pgfpathlineto{\pgfqpoint{4.676252in}{2.191868in}}%
\pgfpathlineto{\pgfqpoint{4.677367in}{2.157103in}}%
\pgfpathlineto{\pgfqpoint{4.676714in}{2.205930in}}%
\pgfpathlineto{\pgfqpoint{4.677401in}{2.162686in}}%
\pgfpathlineto{\pgfqpoint{4.678642in}{2.130724in}}%
\pgfpathlineto{\pgfqpoint{4.677479in}{2.169185in}}%
\pgfpathlineto{\pgfqpoint{4.678667in}{2.137185in}}%
\pgfpathlineto{\pgfqpoint{4.679801in}{2.156726in}}%
\pgfpathlineto{\pgfqpoint{4.679261in}{2.106000in}}%
\pgfpathlineto{\pgfqpoint{4.679811in}{2.155879in}}%
\pgfpathlineto{\pgfqpoint{4.679816in}{2.155895in}}%
\pgfpathlineto{\pgfqpoint{4.679820in}{2.154147in}}%
\pgfpathlineto{\pgfqpoint{4.679825in}{2.154479in}}%
\pgfpathlineto{\pgfqpoint{4.680590in}{2.133672in}}%
\pgfpathlineto{\pgfqpoint{4.680779in}{2.175066in}}%
\pgfpathlineto{\pgfqpoint{4.680935in}{2.151504in}}%
\pgfpathlineto{\pgfqpoint{4.681149in}{2.139785in}}%
\pgfpathlineto{\pgfqpoint{4.682104in}{2.173793in}}%
\pgfpathlineto{\pgfqpoint{4.684178in}{2.085564in}}%
\pgfpathlineto{\pgfqpoint{4.682225in}{2.190349in}}%
\pgfpathlineto{\pgfqpoint{4.684353in}{2.103696in}}%
\pgfpathlineto{\pgfqpoint{4.684450in}{2.118344in}}%
\pgfpathlineto{\pgfqpoint{4.685069in}{2.086783in}}%
\pgfpathlineto{\pgfqpoint{4.685463in}{2.104770in}}%
\pgfpathlineto{\pgfqpoint{4.686286in}{2.073538in}}%
\pgfpathlineto{\pgfqpoint{4.685638in}{2.132159in}}%
\pgfpathlineto{\pgfqpoint{4.686544in}{2.104856in}}%
\pgfpathlineto{\pgfqpoint{4.686948in}{2.115817in}}%
\pgfpathlineto{\pgfqpoint{4.687337in}{2.082507in}}%
\pgfpathlineto{\pgfqpoint{4.687620in}{2.086962in}}%
\pgfpathlineto{\pgfqpoint{4.688715in}{2.127252in}}%
\pgfpathlineto{\pgfqpoint{4.687639in}{2.084733in}}%
\pgfpathlineto{\pgfqpoint{4.688827in}{2.121129in}}%
\pgfpathlineto{\pgfqpoint{4.688837in}{2.118708in}}%
\pgfpathlineto{\pgfqpoint{4.689265in}{2.182566in}}%
\pgfpathlineto{\pgfqpoint{4.689509in}{2.170634in}}%
\pgfpathlineto{\pgfqpoint{4.690658in}{2.186673in}}%
\pgfpathlineto{\pgfqpoint{4.690288in}{2.160721in}}%
\pgfpathlineto{\pgfqpoint{4.690663in}{2.186299in}}%
\pgfpathlineto{\pgfqpoint{4.691578in}{2.146617in}}%
\pgfpathlineto{\pgfqpoint{4.690682in}{2.188101in}}%
\pgfpathlineto{\pgfqpoint{4.691865in}{2.162934in}}%
\pgfpathlineto{\pgfqpoint{4.691880in}{2.161808in}}%
\pgfpathlineto{\pgfqpoint{4.691875in}{2.163704in}}%
\pgfpathlineto{\pgfqpoint{4.691889in}{2.162516in}}%
\pgfpathlineto{\pgfqpoint{4.692503in}{2.194521in}}%
\pgfpathlineto{\pgfqpoint{4.693033in}{2.145177in}}%
\pgfpathlineto{\pgfqpoint{4.693632in}{2.185668in}}%
\pgfpathlineto{\pgfqpoint{4.694202in}{2.172089in}}%
\pgfpathlineto{\pgfqpoint{4.694231in}{2.163795in}}%
\pgfpathlineto{\pgfqpoint{4.694854in}{2.211593in}}%
\pgfpathlineto{\pgfqpoint{4.695288in}{2.185730in}}%
\pgfpathlineto{\pgfqpoint{4.696427in}{2.168982in}}%
\pgfpathlineto{\pgfqpoint{4.696032in}{2.204696in}}%
\pgfpathlineto{\pgfqpoint{4.696441in}{2.170625in}}%
\pgfpathlineto{\pgfqpoint{4.696899in}{2.186972in}}%
\pgfpathlineto{\pgfqpoint{4.697216in}{2.156252in}}%
\pgfpathlineto{\pgfqpoint{4.697551in}{2.175597in}}%
\pgfpathlineto{\pgfqpoint{4.697868in}{2.154109in}}%
\pgfpathlineto{\pgfqpoint{4.698618in}{2.198259in}}%
\pgfpathlineto{\pgfqpoint{4.698627in}{2.197406in}}%
\pgfpathlineto{\pgfqpoint{4.699728in}{2.207789in}}%
\pgfpathlineto{\pgfqpoint{4.698769in}{2.181115in}}%
\pgfpathlineto{\pgfqpoint{4.699767in}{2.205342in}}%
\pgfpathlineto{\pgfqpoint{4.699791in}{2.203581in}}%
\pgfpathlineto{\pgfqpoint{4.700765in}{2.234329in}}%
\pgfpathlineto{\pgfqpoint{4.700774in}{2.238200in}}%
\pgfpathlineto{\pgfqpoint{4.701807in}{2.210006in}}%
\pgfpathlineto{\pgfqpoint{4.701831in}{2.210712in}}%
\pgfpathlineto{\pgfqpoint{4.702576in}{2.199973in}}%
\pgfpathlineto{\pgfqpoint{4.702293in}{2.225341in}}%
\pgfpathlineto{\pgfqpoint{4.702897in}{2.214966in}}%
\pgfpathlineto{\pgfqpoint{4.703374in}{2.262318in}}%
\pgfpathlineto{\pgfqpoint{4.704080in}{2.251951in}}%
\pgfpathlineto{\pgfqpoint{4.704085in}{2.252181in}}%
\pgfpathlineto{\pgfqpoint{4.704659in}{2.227643in}}%
\pgfpathlineto{\pgfqpoint{4.704684in}{2.231104in}}%
\pgfpathlineto{\pgfqpoint{4.704796in}{2.208572in}}%
\pgfpathlineto{\pgfqpoint{4.705799in}{2.222371in}}%
\pgfpathlineto{\pgfqpoint{4.705808in}{2.223159in}}%
\pgfpathlineto{\pgfqpoint{4.705852in}{2.217994in}}%
\pgfpathlineto{\pgfqpoint{4.706003in}{2.204512in}}%
\pgfpathlineto{\pgfqpoint{4.706709in}{2.236607in}}%
\pgfpathlineto{\pgfqpoint{4.706953in}{2.224107in}}%
\pgfpathlineto{\pgfqpoint{4.707483in}{2.199962in}}%
\pgfpathlineto{\pgfqpoint{4.707951in}{2.230273in}}%
\pgfpathlineto{\pgfqpoint{4.708028in}{2.226838in}}%
\pgfpathlineto{\pgfqpoint{4.708876in}{2.241917in}}%
\pgfpathlineto{\pgfqpoint{4.709085in}{2.213401in}}%
\pgfpathlineto{\pgfqpoint{4.709966in}{2.166731in}}%
\pgfpathlineto{\pgfqpoint{4.710258in}{2.175626in}}%
\pgfpathlineto{\pgfqpoint{4.710599in}{2.193947in}}%
\pgfpathlineto{\pgfqpoint{4.711120in}{2.150161in}}%
\pgfpathlineto{\pgfqpoint{4.711212in}{2.155873in}}%
\pgfpathlineto{\pgfqpoint{4.712284in}{2.139241in}}%
\pgfpathlineto{\pgfqpoint{4.711407in}{2.167587in}}%
\pgfpathlineto{\pgfqpoint{4.712366in}{2.145767in}}%
\pgfpathlineto{\pgfqpoint{4.713028in}{2.132393in}}%
\pgfpathlineto{\pgfqpoint{4.713359in}{2.155169in}}%
\pgfpathlineto{\pgfqpoint{4.713364in}{2.154937in}}%
\pgfpathlineto{\pgfqpoint{4.713725in}{2.173855in}}%
\pgfpathlineto{\pgfqpoint{4.714241in}{2.137223in}}%
\pgfpathlineto{\pgfqpoint{4.714314in}{2.141893in}}%
\pgfpathlineto{\pgfqpoint{4.715025in}{2.097957in}}%
\pgfpathlineto{\pgfqpoint{4.715438in}{2.131552in}}%
\pgfpathlineto{\pgfqpoint{4.717025in}{2.095522in}}%
\pgfpathlineto{\pgfqpoint{4.715653in}{2.138462in}}%
\pgfpathlineto{\pgfqpoint{4.717030in}{2.098809in}}%
\pgfpathlineto{\pgfqpoint{4.717230in}{2.104146in}}%
\pgfpathlineto{\pgfqpoint{4.718082in}{2.063111in}}%
\pgfpathlineto{\pgfqpoint{4.718262in}{2.070491in}}%
\pgfpathlineto{\pgfqpoint{4.718759in}{2.035837in}}%
\pgfpathlineto{\pgfqpoint{4.719099in}{2.054296in}}%
\pgfpathlineto{\pgfqpoint{4.719479in}{2.038260in}}%
\pgfpathlineto{\pgfqpoint{4.719966in}{2.064637in}}%
\pgfpathlineto{\pgfqpoint{4.720258in}{2.038347in}}%
\pgfpathlineto{\pgfqpoint{4.721427in}{2.076382in}}%
\pgfpathlineto{\pgfqpoint{4.720360in}{2.029375in}}%
\pgfpathlineto{\pgfqpoint{4.721461in}{2.075338in}}%
\pgfpathlineto{\pgfqpoint{4.722478in}{2.048413in}}%
\pgfpathlineto{\pgfqpoint{4.721777in}{2.090957in}}%
\pgfpathlineto{\pgfqpoint{4.722600in}{2.054933in}}%
\pgfpathlineto{\pgfqpoint{4.723549in}{2.048727in}}%
\pgfpathlineto{\pgfqpoint{4.723043in}{2.077403in}}%
\pgfpathlineto{\pgfqpoint{4.723686in}{2.064422in}}%
\pgfpathlineto{\pgfqpoint{4.723880in}{2.094669in}}%
\pgfpathlineto{\pgfqpoint{4.724396in}{2.059097in}}%
\pgfpathlineto{\pgfqpoint{4.724839in}{2.070765in}}%
\pgfpathlineto{\pgfqpoint{4.725463in}{2.039820in}}%
\pgfpathlineto{\pgfqpoint{4.725974in}{2.050343in}}%
\pgfpathlineto{\pgfqpoint{4.726894in}{2.085592in}}%
\pgfpathlineto{\pgfqpoint{4.727191in}{2.066827in}}%
\pgfpathlineto{\pgfqpoint{4.728486in}{2.021424in}}%
\pgfpathlineto{\pgfqpoint{4.727254in}{2.072757in}}%
\pgfpathlineto{\pgfqpoint{4.728501in}{2.029718in}}%
\pgfpathlineto{\pgfqpoint{4.729416in}{2.044843in}}%
\pgfpathlineto{\pgfqpoint{4.728739in}{2.011316in}}%
\pgfpathlineto{\pgfqpoint{4.729474in}{2.027308in}}%
\pgfpathlineto{\pgfqpoint{4.729542in}{2.020462in}}%
\pgfpathlineto{\pgfqpoint{4.730102in}{2.055598in}}%
\pgfpathlineto{\pgfqpoint{4.730550in}{2.039548in}}%
\pgfpathlineto{\pgfqpoint{4.730686in}{2.061110in}}%
\pgfpathlineto{\pgfqpoint{4.731631in}{2.020924in}}%
\pgfpathlineto{\pgfqpoint{4.731782in}{2.004615in}}%
\pgfpathlineto{\pgfqpoint{4.732639in}{2.041627in}}%
\pgfpathlineto{\pgfqpoint{4.732658in}{2.039663in}}%
\pgfpathlineto{\pgfqpoint{4.733116in}{2.057448in}}%
\pgfpathlineto{\pgfqpoint{4.733724in}{2.025738in}}%
\pgfpathlineto{\pgfqpoint{4.734883in}{1.980210in}}%
\pgfpathlineto{\pgfqpoint{4.733797in}{2.031084in}}%
\pgfpathlineto{\pgfqpoint{4.734907in}{1.984006in}}%
\pgfpathlineto{\pgfqpoint{4.735609in}{1.996294in}}%
\pgfpathlineto{\pgfqpoint{4.735015in}{1.966341in}}%
\pgfpathlineto{\pgfqpoint{4.735998in}{1.978155in}}%
\pgfpathlineto{\pgfqpoint{4.736105in}{1.968122in}}%
\pgfpathlineto{\pgfqpoint{4.736675in}{2.022181in}}%
\pgfpathlineto{\pgfqpoint{4.737059in}{2.004826in}}%
\pgfpathlineto{\pgfqpoint{4.737064in}{2.006172in}}%
\pgfpathlineto{\pgfqpoint{4.737575in}{1.975576in}}%
\pgfpathlineto{\pgfqpoint{4.738087in}{1.981792in}}%
\pgfpathlineto{\pgfqpoint{4.738340in}{2.006590in}}%
\pgfpathlineto{\pgfqpoint{4.739007in}{1.942991in}}%
\pgfpathlineto{\pgfqpoint{4.739012in}{1.943946in}}%
\pgfpathlineto{\pgfqpoint{4.740122in}{1.904484in}}%
\pgfpathlineto{\pgfqpoint{4.740185in}{1.911797in}}%
\pgfpathlineto{\pgfqpoint{4.740190in}{1.912563in}}%
\pgfpathlineto{\pgfqpoint{4.740798in}{1.871591in}}%
\pgfpathlineto{\pgfqpoint{4.741222in}{1.900351in}}%
\pgfpathlineto{\pgfqpoint{4.741490in}{1.904234in}}%
\pgfpathlineto{\pgfqpoint{4.741422in}{1.891548in}}%
\pgfpathlineto{\pgfqpoint{4.741597in}{1.894226in}}%
\pgfpathlineto{\pgfqpoint{4.741607in}{1.891108in}}%
\pgfpathlineto{\pgfqpoint{4.742502in}{1.959904in}}%
\pgfpathlineto{\pgfqpoint{4.742551in}{1.958897in}}%
\pgfpathlineto{\pgfqpoint{4.742663in}{1.963615in}}%
\pgfpathlineto{\pgfqpoint{4.743014in}{1.933529in}}%
\pgfpathlineto{\pgfqpoint{4.743637in}{1.946654in}}%
\pgfpathlineto{\pgfqpoint{4.743661in}{1.945403in}}%
\pgfpathlineto{\pgfqpoint{4.744459in}{1.999706in}}%
\pgfpathlineto{\pgfqpoint{4.744469in}{1.997358in}}%
\pgfpathlineto{\pgfqpoint{4.744903in}{2.003898in}}%
\pgfpathlineto{\pgfqpoint{4.744640in}{1.988514in}}%
\pgfpathlineto{\pgfqpoint{4.745005in}{1.993112in}}%
\pgfpathlineto{\pgfqpoint{4.746017in}{1.968490in}}%
\pgfpathlineto{\pgfqpoint{4.745389in}{2.006662in}}%
\pgfpathlineto{\pgfqpoint{4.746134in}{1.975044in}}%
\pgfpathlineto{\pgfqpoint{4.746149in}{1.973830in}}%
\pgfpathlineto{\pgfqpoint{4.746183in}{1.984667in}}%
\pgfpathlineto{\pgfqpoint{4.746188in}{1.983308in}}%
\pgfpathlineto{\pgfqpoint{4.746397in}{1.994389in}}%
\pgfpathlineto{\pgfqpoint{4.747220in}{1.948680in}}%
\pgfpathlineto{\pgfqpoint{4.747244in}{1.953335in}}%
\pgfpathlineto{\pgfqpoint{4.747502in}{1.937373in}}%
\pgfpathlineto{\pgfqpoint{4.748242in}{1.966103in}}%
\pgfpathlineto{\pgfqpoint{4.748247in}{1.965062in}}%
\pgfpathlineto{\pgfqpoint{4.748530in}{1.979481in}}%
\pgfpathlineto{\pgfqpoint{4.749109in}{1.949217in}}%
\pgfpathlineto{\pgfqpoint{4.749289in}{1.961959in}}%
\pgfpathlineto{\pgfqpoint{4.750282in}{1.913430in}}%
\pgfpathlineto{\pgfqpoint{4.750433in}{1.919927in}}%
\pgfpathlineto{\pgfqpoint{4.750462in}{1.925344in}}%
\pgfpathlineto{\pgfqpoint{4.750803in}{1.897774in}}%
\pgfpathlineto{\pgfqpoint{4.751524in}{1.915275in}}%
\pgfpathlineto{\pgfqpoint{4.752308in}{1.892441in}}%
\pgfpathlineto{\pgfqpoint{4.751665in}{1.926008in}}%
\pgfpathlineto{\pgfqpoint{4.752600in}{1.921429in}}%
\pgfpathlineto{\pgfqpoint{4.753160in}{1.940758in}}%
\pgfpathlineto{\pgfqpoint{4.753622in}{1.913294in}}%
\pgfpathlineto{\pgfqpoint{4.753700in}{1.921036in}}%
\pgfpathlineto{\pgfqpoint{4.753719in}{1.915594in}}%
\pgfpathlineto{\pgfqpoint{4.754669in}{1.944388in}}%
\pgfpathlineto{\pgfqpoint{4.754776in}{1.943174in}}%
\pgfpathlineto{\pgfqpoint{4.755097in}{1.950367in}}%
\pgfpathlineto{\pgfqpoint{4.755779in}{1.916669in}}%
\pgfpathlineto{\pgfqpoint{4.755847in}{1.927881in}}%
\pgfpathlineto{\pgfqpoint{4.756295in}{1.918833in}}%
\pgfpathlineto{\pgfqpoint{4.756139in}{1.941304in}}%
\pgfpathlineto{\pgfqpoint{4.756981in}{1.921434in}}%
\pgfpathlineto{\pgfqpoint{4.757074in}{1.933285in}}%
\pgfpathlineto{\pgfqpoint{4.758028in}{1.894908in}}%
\pgfpathlineto{\pgfqpoint{4.758052in}{1.896211in}}%
\pgfpathlineto{\pgfqpoint{4.758077in}{1.893250in}}%
\pgfpathlineto{\pgfqpoint{4.758948in}{1.927991in}}%
\pgfpathlineto{\pgfqpoint{4.759050in}{1.917695in}}%
\pgfpathlineto{\pgfqpoint{4.759445in}{1.952020in}}%
\pgfpathlineto{\pgfqpoint{4.760121in}{1.908173in}}%
\pgfpathlineto{\pgfqpoint{4.760126in}{1.905658in}}%
\pgfpathlineto{\pgfqpoint{4.760973in}{1.936534in}}%
\pgfpathlineto{\pgfqpoint{4.761188in}{1.930508in}}%
\pgfpathlineto{\pgfqpoint{4.761319in}{1.915612in}}%
\pgfpathlineto{\pgfqpoint{4.761796in}{1.953413in}}%
\pgfpathlineto{\pgfqpoint{4.761918in}{1.974754in}}%
\pgfpathlineto{\pgfqpoint{4.762166in}{1.946256in}}%
\pgfpathlineto{\pgfqpoint{4.762935in}{1.972996in}}%
\pgfpathlineto{\pgfqpoint{4.763281in}{1.993012in}}%
\pgfpathlineto{\pgfqpoint{4.763890in}{1.951411in}}%
\pgfpathlineto{\pgfqpoint{4.764011in}{1.959867in}}%
\pgfpathlineto{\pgfqpoint{4.765399in}{2.007269in}}%
\pgfpathlineto{\pgfqpoint{4.764075in}{1.959034in}}%
\pgfpathlineto{\pgfqpoint{4.765414in}{2.005007in}}%
\pgfpathlineto{\pgfqpoint{4.766285in}{1.993368in}}%
\pgfpathlineto{\pgfqpoint{4.766460in}{2.014411in}}%
\pgfpathlineto{\pgfqpoint{4.766499in}{2.007195in}}%
\pgfpathlineto{\pgfqpoint{4.767599in}{2.040387in}}%
\pgfpathlineto{\pgfqpoint{4.766923in}{2.003623in}}%
\pgfpathlineto{\pgfqpoint{4.767653in}{2.033235in}}%
\pgfpathlineto{\pgfqpoint{4.767697in}{2.042664in}}%
\pgfpathlineto{\pgfqpoint{4.767926in}{2.017754in}}%
\pgfpathlineto{\pgfqpoint{4.768456in}{1.989912in}}%
\pgfpathlineto{\pgfqpoint{4.768763in}{2.021133in}}%
\pgfpathlineto{\pgfqpoint{4.769041in}{2.015095in}}%
\pgfpathlineto{\pgfqpoint{4.769167in}{2.028047in}}%
\pgfpathlineto{\pgfqpoint{4.769771in}{1.980478in}}%
\pgfpathlineto{\pgfqpoint{4.769776in}{1.981140in}}%
\pgfpathlineto{\pgfqpoint{4.770934in}{1.957731in}}%
\pgfpathlineto{\pgfqpoint{4.770180in}{1.996513in}}%
\pgfpathlineto{\pgfqpoint{4.770964in}{1.961751in}}%
\pgfpathlineto{\pgfqpoint{4.770968in}{1.966427in}}%
\pgfpathlineto{\pgfqpoint{4.771762in}{1.905177in}}%
\pgfpathlineto{\pgfqpoint{4.772020in}{1.913546in}}%
\pgfpathlineto{\pgfqpoint{4.773150in}{1.872360in}}%
\pgfpathlineto{\pgfqpoint{4.772035in}{1.915581in}}%
\pgfpathlineto{\pgfqpoint{4.773559in}{1.893098in}}%
\pgfpathlineto{\pgfqpoint{4.773563in}{1.895847in}}%
\pgfpathlineto{\pgfqpoint{4.774440in}{1.829958in}}%
\pgfpathlineto{\pgfqpoint{4.774571in}{1.838759in}}%
\pgfpathlineto{\pgfqpoint{4.775486in}{1.860765in}}%
\pgfpathlineto{\pgfqpoint{4.774693in}{1.822032in}}%
\pgfpathlineto{\pgfqpoint{4.775876in}{1.843661in}}%
\pgfpathlineto{\pgfqpoint{4.777049in}{1.822947in}}%
\pgfpathlineto{\pgfqpoint{4.776348in}{1.850424in}}%
\pgfpathlineto{\pgfqpoint{4.777054in}{1.826032in}}%
\pgfpathlineto{\pgfqpoint{4.778412in}{1.930761in}}%
\pgfpathlineto{\pgfqpoint{4.778422in}{1.928283in}}%
\pgfpathlineto{\pgfqpoint{4.779488in}{1.894256in}}%
\pgfpathlineto{\pgfqpoint{4.778476in}{1.929480in}}%
\pgfpathlineto{\pgfqpoint{4.779620in}{1.908606in}}%
\pgfpathlineto{\pgfqpoint{4.779625in}{1.909472in}}%
\pgfpathlineto{\pgfqpoint{4.780306in}{1.848245in}}%
\pgfpathlineto{\pgfqpoint{4.780394in}{1.849358in}}%
\pgfpathlineto{\pgfqpoint{4.780404in}{1.843400in}}%
\pgfpathlineto{\pgfqpoint{4.781411in}{1.894822in}}%
\pgfpathlineto{\pgfqpoint{4.781465in}{1.885406in}}%
\pgfpathlineto{\pgfqpoint{4.782629in}{1.908220in}}%
\pgfpathlineto{\pgfqpoint{4.781796in}{1.868513in}}%
\pgfpathlineto{\pgfqpoint{4.782643in}{1.907060in}}%
\pgfpathlineto{\pgfqpoint{4.782677in}{1.902696in}}%
\pgfpathlineto{\pgfqpoint{4.783461in}{1.939546in}}%
\pgfpathlineto{\pgfqpoint{4.783607in}{1.934404in}}%
\pgfpathlineto{\pgfqpoint{4.783641in}{1.939888in}}%
\pgfpathlineto{\pgfqpoint{4.783909in}{1.901738in}}%
\pgfpathlineto{\pgfqpoint{4.784707in}{1.929723in}}%
\pgfpathlineto{\pgfqpoint{4.785535in}{1.901123in}}%
\pgfpathlineto{\pgfqpoint{4.785185in}{1.931430in}}%
\pgfpathlineto{\pgfqpoint{4.785842in}{1.915859in}}%
\pgfpathlineto{\pgfqpoint{4.785886in}{1.926000in}}%
\pgfpathlineto{\pgfqpoint{4.786893in}{1.885024in}}%
\pgfpathlineto{\pgfqpoint{4.787458in}{1.875121in}}%
\pgfpathlineto{\pgfqpoint{4.787181in}{1.901043in}}%
\pgfpathlineto{\pgfqpoint{4.787979in}{1.895606in}}%
\pgfpathlineto{\pgfqpoint{4.788276in}{1.903717in}}%
\pgfpathlineto{\pgfqpoint{4.789035in}{1.873568in}}%
\pgfpathlineto{\pgfqpoint{4.789230in}{1.867843in}}%
\pgfpathlineto{\pgfqpoint{4.789600in}{1.907114in}}%
\pgfpathlineto{\pgfqpoint{4.789951in}{1.901617in}}%
\pgfpathlineto{\pgfqpoint{4.790778in}{1.973789in}}%
\pgfpathlineto{\pgfqpoint{4.789985in}{1.896617in}}%
\pgfpathlineto{\pgfqpoint{4.791148in}{1.936699in}}%
\pgfpathlineto{\pgfqpoint{4.792322in}{1.905288in}}%
\pgfpathlineto{\pgfqpoint{4.791674in}{1.955186in}}%
\pgfpathlineto{\pgfqpoint{4.792331in}{1.905420in}}%
\pgfpathlineto{\pgfqpoint{4.792770in}{1.926556in}}%
\pgfpathlineto{\pgfqpoint{4.793149in}{1.897970in}}%
\pgfpathlineto{\pgfqpoint{4.793437in}{1.905899in}}%
\pgfpathlineto{\pgfqpoint{4.793592in}{1.882675in}}%
\pgfpathlineto{\pgfqpoint{4.794050in}{1.926461in}}%
\pgfpathlineto{\pgfqpoint{4.794498in}{1.914793in}}%
\pgfpathlineto{\pgfqpoint{4.794785in}{1.917959in}}%
\pgfpathlineto{\pgfqpoint{4.794712in}{1.896294in}}%
\pgfpathlineto{\pgfqpoint{4.795472in}{1.905677in}}%
\pgfpathlineto{\pgfqpoint{4.796348in}{1.873750in}}%
\pgfpathlineto{\pgfqpoint{4.795530in}{1.906024in}}%
\pgfpathlineto{\pgfqpoint{4.796611in}{1.884514in}}%
\pgfpathlineto{\pgfqpoint{4.797804in}{1.929153in}}%
\pgfpathlineto{\pgfqpoint{4.796689in}{1.883408in}}%
\pgfpathlineto{\pgfqpoint{4.797838in}{1.916635in}}%
\pgfpathlineto{\pgfqpoint{4.797867in}{1.905709in}}%
\pgfpathlineto{\pgfqpoint{4.798631in}{1.962942in}}%
\pgfpathlineto{\pgfqpoint{4.798748in}{1.956491in}}%
\pgfpathlineto{\pgfqpoint{4.799021in}{1.974314in}}%
\pgfpathlineto{\pgfqpoint{4.799586in}{1.949700in}}%
\pgfpathlineto{\pgfqpoint{4.799829in}{1.952707in}}%
\pgfpathlineto{\pgfqpoint{4.800126in}{1.947279in}}%
\pgfpathlineto{\pgfqpoint{4.800394in}{1.987173in}}%
\pgfpathlineto{\pgfqpoint{4.800666in}{1.979896in}}%
\pgfpathlineto{\pgfqpoint{4.802122in}{2.034358in}}%
\pgfpathlineto{\pgfqpoint{4.800817in}{1.971341in}}%
\pgfpathlineto{\pgfqpoint{4.802195in}{2.024615in}}%
\pgfpathlineto{\pgfqpoint{4.802794in}{2.030828in}}%
\pgfpathlineto{\pgfqpoint{4.803383in}{2.007119in}}%
\pgfpathlineto{\pgfqpoint{4.803471in}{2.020681in}}%
\pgfpathlineto{\pgfqpoint{4.804303in}{1.963526in}}%
\pgfpathlineto{\pgfqpoint{4.804405in}{1.969810in}}%
\pgfpathlineto{\pgfqpoint{4.804420in}{1.968681in}}%
\pgfpathlineto{\pgfqpoint{4.804425in}{1.970082in}}%
\pgfpathlineto{\pgfqpoint{4.804659in}{1.944830in}}%
\pgfpathlineto{\pgfqpoint{4.804469in}{1.973847in}}%
\pgfpathlineto{\pgfqpoint{4.805549in}{1.956234in}}%
\pgfpathlineto{\pgfqpoint{4.806737in}{2.019006in}}%
\pgfpathlineto{\pgfqpoint{4.805564in}{1.956113in}}%
\pgfpathlineto{\pgfqpoint{4.806767in}{2.011957in}}%
\pgfpathlineto{\pgfqpoint{4.808086in}{1.950786in}}%
\pgfpathlineto{\pgfqpoint{4.806781in}{2.014611in}}%
\pgfpathlineto{\pgfqpoint{4.808354in}{1.960213in}}%
\pgfpathlineto{\pgfqpoint{4.808529in}{1.965533in}}%
\pgfpathlineto{\pgfqpoint{4.808923in}{1.936861in}}%
\pgfpathlineto{\pgfqpoint{4.809435in}{1.946329in}}%
\pgfpathlineto{\pgfqpoint{4.810953in}{2.025337in}}%
\pgfpathlineto{\pgfqpoint{4.809532in}{1.943646in}}%
\pgfpathlineto{\pgfqpoint{4.811411in}{1.995051in}}%
\pgfpathlineto{\pgfqpoint{4.811762in}{1.985278in}}%
\pgfpathlineto{\pgfqpoint{4.812351in}{2.017238in}}%
\pgfpathlineto{\pgfqpoint{4.812502in}{1.999732in}}%
\pgfpathlineto{\pgfqpoint{4.813675in}{1.967371in}}%
\pgfpathlineto{\pgfqpoint{4.813188in}{2.011606in}}%
\pgfpathlineto{\pgfqpoint{4.813894in}{1.982202in}}%
\pgfpathlineto{\pgfqpoint{4.814693in}{2.013134in}}%
\pgfpathlineto{\pgfqpoint{4.815014in}{1.990397in}}%
\pgfpathlineto{\pgfqpoint{4.815798in}{1.951007in}}%
\pgfpathlineto{\pgfqpoint{4.815369in}{1.995055in}}%
\pgfpathlineto{\pgfqpoint{4.816148in}{1.973564in}}%
\pgfpathlineto{\pgfqpoint{4.816431in}{2.017412in}}%
\pgfpathlineto{\pgfqpoint{4.817492in}{2.011554in}}%
\pgfpathlineto{\pgfqpoint{4.817959in}{1.985267in}}%
\pgfpathlineto{\pgfqpoint{4.818339in}{2.017925in}}%
\pgfpathlineto{\pgfqpoint{4.818602in}{2.010901in}}%
\pgfpathlineto{\pgfqpoint{4.818621in}{2.004583in}}%
\pgfpathlineto{\pgfqpoint{4.819405in}{2.020412in}}%
\pgfpathlineto{\pgfqpoint{4.819644in}{2.016933in}}%
\pgfpathlineto{\pgfqpoint{4.819853in}{2.029308in}}%
\pgfpathlineto{\pgfqpoint{4.820233in}{1.998855in}}%
\pgfpathlineto{\pgfqpoint{4.820725in}{2.005149in}}%
\pgfpathlineto{\pgfqpoint{4.821635in}{1.978925in}}%
\pgfpathlineto{\pgfqpoint{4.820749in}{2.005967in}}%
\pgfpathlineto{\pgfqpoint{4.821869in}{1.992939in}}%
\pgfpathlineto{\pgfqpoint{4.823091in}{2.062651in}}%
\pgfpathlineto{\pgfqpoint{4.823130in}{2.058484in}}%
\pgfpathlineto{\pgfqpoint{4.824415in}{1.974293in}}%
\pgfpathlineto{\pgfqpoint{4.824541in}{1.980608in}}%
\pgfpathlineto{\pgfqpoint{4.824673in}{1.966516in}}%
\pgfpathlineto{\pgfqpoint{4.825174in}{1.999798in}}%
\pgfpathlineto{\pgfqpoint{4.825593in}{1.988027in}}%
\pgfpathlineto{\pgfqpoint{4.826119in}{1.973964in}}%
\pgfpathlineto{\pgfqpoint{4.826742in}{2.016200in}}%
\pgfpathlineto{\pgfqpoint{4.826981in}{2.012420in}}%
\pgfpathlineto{\pgfqpoint{4.827706in}{2.054291in}}%
\pgfpathlineto{\pgfqpoint{4.827750in}{2.049913in}}%
\pgfpathlineto{\pgfqpoint{4.828573in}{2.061423in}}%
\pgfpathlineto{\pgfqpoint{4.828237in}{2.033382in}}%
\pgfpathlineto{\pgfqpoint{4.828840in}{2.034249in}}%
\pgfpathlineto{\pgfqpoint{4.828855in}{2.032705in}}%
\pgfpathlineto{\pgfqpoint{4.829980in}{2.002546in}}%
\pgfpathlineto{\pgfqpoint{4.829259in}{2.043620in}}%
\pgfpathlineto{\pgfqpoint{4.830023in}{2.006689in}}%
\pgfpathlineto{\pgfqpoint{4.830155in}{2.016452in}}%
\pgfpathlineto{\pgfqpoint{4.830661in}{1.986237in}}%
\pgfpathlineto{\pgfqpoint{4.831070in}{1.996259in}}%
\pgfpathlineto{\pgfqpoint{4.832190in}{1.946876in}}%
\pgfpathlineto{\pgfqpoint{4.831455in}{2.005620in}}%
\pgfpathlineto{\pgfqpoint{4.832243in}{1.955842in}}%
\pgfpathlineto{\pgfqpoint{4.833256in}{1.986797in}}%
\pgfpathlineto{\pgfqpoint{4.833456in}{1.984063in}}%
\pgfpathlineto{\pgfqpoint{4.834098in}{1.978904in}}%
\pgfpathlineto{\pgfqpoint{4.834001in}{1.992939in}}%
\pgfpathlineto{\pgfqpoint{4.834566in}{1.983404in}}%
\pgfpathlineto{\pgfqpoint{4.834712in}{1.996958in}}%
\pgfpathlineto{\pgfqpoint{4.835082in}{1.973220in}}%
\pgfpathlineto{\pgfqpoint{4.835681in}{1.986520in}}%
\pgfpathlineto{\pgfqpoint{4.836026in}{1.952482in}}%
\pgfpathlineto{\pgfqpoint{4.835749in}{1.988420in}}%
\pgfpathlineto{\pgfqpoint{4.836844in}{1.960557in}}%
\pgfpathlineto{\pgfqpoint{4.837078in}{1.967880in}}%
\pgfpathlineto{\pgfqpoint{4.837506in}{1.929838in}}%
\pgfpathlineto{\pgfqpoint{4.837935in}{1.945420in}}%
\pgfpathlineto{\pgfqpoint{4.837944in}{1.944689in}}%
\pgfpathlineto{\pgfqpoint{4.838280in}{1.968279in}}%
\pgfpathlineto{\pgfqpoint{4.839196in}{2.004564in}}%
\pgfpathlineto{\pgfqpoint{4.838485in}{1.963950in}}%
\pgfpathlineto{\pgfqpoint{4.839420in}{1.996366in}}%
\pgfpathlineto{\pgfqpoint{4.839478in}{2.001328in}}%
\pgfpathlineto{\pgfqpoint{4.839790in}{1.980685in}}%
\pgfpathlineto{\pgfqpoint{4.839882in}{1.970318in}}%
\pgfpathlineto{\pgfqpoint{4.840817in}{2.030188in}}%
\pgfpathlineto{\pgfqpoint{4.840836in}{2.027755in}}%
\pgfpathlineto{\pgfqpoint{4.842034in}{2.056373in}}%
\pgfpathlineto{\pgfqpoint{4.841158in}{2.013767in}}%
\pgfpathlineto{\pgfqpoint{4.842049in}{2.049936in}}%
\pgfpathlineto{\pgfqpoint{4.843261in}{1.995652in}}%
\pgfpathlineto{\pgfqpoint{4.843417in}{2.006352in}}%
\pgfpathlineto{\pgfqpoint{4.844015in}{2.043684in}}%
\pgfpathlineto{\pgfqpoint{4.843436in}{2.004991in}}%
\pgfpathlineto{\pgfqpoint{4.844575in}{2.027231in}}%
\pgfpathlineto{\pgfqpoint{4.845617in}{2.033576in}}%
\pgfpathlineto{\pgfqpoint{4.844726in}{2.009397in}}%
\pgfpathlineto{\pgfqpoint{4.845651in}{2.025412in}}%
\pgfpathlineto{\pgfqpoint{4.846182in}{1.997545in}}%
\pgfpathlineto{\pgfqpoint{4.845734in}{2.031347in}}%
\pgfpathlineto{\pgfqpoint{4.846766in}{2.022344in}}%
\pgfpathlineto{\pgfqpoint{4.846776in}{2.025956in}}%
\pgfpathlineto{\pgfqpoint{4.847779in}{1.987419in}}%
\pgfpathlineto{\pgfqpoint{4.847818in}{1.990802in}}%
\pgfpathlineto{\pgfqpoint{4.848300in}{1.979363in}}%
\pgfpathlineto{\pgfqpoint{4.848120in}{2.000112in}}%
\pgfpathlineto{\pgfqpoint{4.848928in}{1.988030in}}%
\pgfpathlineto{\pgfqpoint{4.849605in}{2.000208in}}%
\pgfpathlineto{\pgfqpoint{4.849079in}{1.979805in}}%
\pgfpathlineto{\pgfqpoint{4.850033in}{1.987102in}}%
\pgfpathlineto{\pgfqpoint{4.850067in}{1.990847in}}%
\pgfpathlineto{\pgfqpoint{4.850184in}{1.996478in}}%
\pgfpathlineto{\pgfqpoint{4.850851in}{1.949300in}}%
\pgfpathlineto{\pgfqpoint{4.851104in}{1.949740in}}%
\pgfpathlineto{\pgfqpoint{4.852039in}{1.924800in}}%
\pgfpathlineto{\pgfqpoint{4.852253in}{1.940932in}}%
\pgfpathlineto{\pgfqpoint{4.852433in}{1.952379in}}%
\pgfpathlineto{\pgfqpoint{4.853047in}{1.907049in}}%
\pgfpathlineto{\pgfqpoint{4.853309in}{1.911722in}}%
\pgfpathlineto{\pgfqpoint{4.853324in}{1.911708in}}%
\pgfpathlineto{\pgfqpoint{4.853407in}{1.930833in}}%
\pgfpathlineto{\pgfqpoint{4.853417in}{1.934100in}}%
\pgfpathlineto{\pgfqpoint{4.854269in}{1.902679in}}%
\pgfpathlineto{\pgfqpoint{4.854478in}{1.917274in}}%
\pgfpathlineto{\pgfqpoint{4.855213in}{1.901292in}}%
\pgfpathlineto{\pgfqpoint{4.854682in}{1.932451in}}%
\pgfpathlineto{\pgfqpoint{4.855471in}{1.919741in}}%
\pgfpathlineto{\pgfqpoint{4.856532in}{1.966421in}}%
\pgfpathlineto{\pgfqpoint{4.855817in}{1.911260in}}%
\pgfpathlineto{\pgfqpoint{4.856713in}{1.952636in}}%
\pgfpathlineto{\pgfqpoint{4.856893in}{1.942286in}}%
\pgfpathlineto{\pgfqpoint{4.857521in}{1.979122in}}%
\pgfpathlineto{\pgfqpoint{4.857715in}{1.974043in}}%
\pgfpathlineto{\pgfqpoint{4.858533in}{1.988552in}}%
\pgfpathlineto{\pgfqpoint{4.858733in}{1.962390in}}%
\pgfpathlineto{\pgfqpoint{4.858796in}{1.965889in}}%
\pgfpathlineto{\pgfqpoint{4.859239in}{1.930686in}}%
\pgfpathlineto{\pgfqpoint{4.859911in}{1.961854in}}%
\pgfpathlineto{\pgfqpoint{4.859979in}{1.972555in}}%
\pgfpathlineto{\pgfqpoint{4.860705in}{1.930842in}}%
\pgfpathlineto{\pgfqpoint{4.860997in}{1.940957in}}%
\pgfpathlineto{\pgfqpoint{4.861007in}{1.940150in}}%
\pgfpathlineto{\pgfqpoint{4.861498in}{1.975819in}}%
\pgfpathlineto{\pgfqpoint{4.861537in}{1.979382in}}%
\pgfpathlineto{\pgfqpoint{4.861834in}{1.941480in}}%
\pgfpathlineto{\pgfqpoint{4.862453in}{1.943694in}}%
\pgfpathlineto{\pgfqpoint{4.862944in}{1.922133in}}%
\pgfpathlineto{\pgfqpoint{4.863485in}{1.951351in}}%
\pgfpathlineto{\pgfqpoint{4.863558in}{1.944675in}}%
\pgfpathlineto{\pgfqpoint{4.863757in}{1.961613in}}%
\pgfpathlineto{\pgfqpoint{4.864546in}{1.919272in}}%
\pgfpathlineto{\pgfqpoint{4.864561in}{1.920994in}}%
\pgfpathlineto{\pgfqpoint{4.865033in}{1.897725in}}%
\pgfpathlineto{\pgfqpoint{4.864663in}{1.922526in}}%
\pgfpathlineto{\pgfqpoint{4.865675in}{1.915010in}}%
\pgfpathlineto{\pgfqpoint{4.866352in}{1.939953in}}%
\pgfpathlineto{\pgfqpoint{4.866717in}{1.904600in}}%
\pgfpathlineto{\pgfqpoint{4.866751in}{1.909039in}}%
\pgfpathlineto{\pgfqpoint{4.867258in}{1.886035in}}%
\pgfpathlineto{\pgfqpoint{4.867667in}{1.920379in}}%
\pgfpathlineto{\pgfqpoint{4.867852in}{1.910727in}}%
\pgfpathlineto{\pgfqpoint{4.868139in}{1.928662in}}%
\pgfpathlineto{\pgfqpoint{4.867939in}{1.908815in}}%
\pgfpathlineto{\pgfqpoint{4.868971in}{1.927105in}}%
\pgfpathlineto{\pgfqpoint{4.868976in}{1.926030in}}%
\pgfpathlineto{\pgfqpoint{4.869809in}{1.955484in}}%
\pgfpathlineto{\pgfqpoint{4.870043in}{1.935379in}}%
\pgfpathlineto{\pgfqpoint{4.870159in}{1.945849in}}%
\pgfpathlineto{\pgfqpoint{4.870471in}{1.920002in}}%
\pgfpathlineto{\pgfqpoint{4.871045in}{1.928170in}}%
\pgfpathlineto{\pgfqpoint{4.871065in}{1.925582in}}%
\pgfpathlineto{\pgfqpoint{4.871980in}{1.982746in}}%
\pgfpathlineto{\pgfqpoint{4.872034in}{1.980686in}}%
\pgfpathlineto{\pgfqpoint{4.872073in}{1.973548in}}%
\pgfpathlineto{\pgfqpoint{4.872399in}{1.998691in}}%
\pgfpathlineto{\pgfqpoint{4.873119in}{1.986074in}}%
\pgfpathlineto{\pgfqpoint{4.873937in}{2.011758in}}%
\pgfpathlineto{\pgfqpoint{4.873212in}{1.976976in}}%
\pgfpathlineto{\pgfqpoint{4.874244in}{1.992567in}}%
\pgfpathlineto{\pgfqpoint{4.874327in}{1.982289in}}%
\pgfpathlineto{\pgfqpoint{4.874926in}{2.013737in}}%
\pgfpathlineto{\pgfqpoint{4.875354in}{1.989132in}}%
\pgfpathlineto{\pgfqpoint{4.875515in}{2.004105in}}%
\pgfpathlineto{\pgfqpoint{4.876406in}{1.970800in}}%
\pgfpathlineto{\pgfqpoint{4.876430in}{1.974952in}}%
\pgfpathlineto{\pgfqpoint{4.877457in}{1.938923in}}%
\pgfpathlineto{\pgfqpoint{4.876508in}{1.979806in}}%
\pgfpathlineto{\pgfqpoint{4.877613in}{1.942937in}}%
\pgfpathlineto{\pgfqpoint{4.877691in}{1.955232in}}%
\pgfpathlineto{\pgfqpoint{4.878338in}{1.913915in}}%
\pgfpathlineto{\pgfqpoint{4.878694in}{1.935404in}}%
\pgfpathlineto{\pgfqpoint{4.879156in}{1.924296in}}%
\pgfpathlineto{\pgfqpoint{4.879702in}{1.958555in}}%
\pgfpathlineto{\pgfqpoint{4.879760in}{1.951669in}}%
\pgfpathlineto{\pgfqpoint{4.879999in}{1.964541in}}%
\pgfpathlineto{\pgfqpoint{4.880115in}{1.942161in}}%
\pgfpathlineto{\pgfqpoint{4.880967in}{1.957558in}}%
\pgfpathlineto{\pgfqpoint{4.880997in}{1.953726in}}%
\pgfpathlineto{\pgfqpoint{4.881693in}{2.006348in}}%
\pgfpathlineto{\pgfqpoint{4.882019in}{1.986589in}}%
\pgfpathlineto{\pgfqpoint{4.882053in}{1.989910in}}%
\pgfpathlineto{\pgfqpoint{4.882671in}{1.957628in}}%
\pgfpathlineto{\pgfqpoint{4.883095in}{1.974135in}}%
\pgfpathlineto{\pgfqpoint{4.883708in}{1.960079in}}%
\pgfpathlineto{\pgfqpoint{4.884137in}{1.989316in}}%
\pgfpathlineto{\pgfqpoint{4.885067in}{2.040472in}}%
\pgfpathlineto{\pgfqpoint{4.885437in}{2.032371in}}%
\pgfpathlineto{\pgfqpoint{4.886235in}{2.018062in}}%
\pgfpathlineto{\pgfqpoint{4.885987in}{2.040882in}}%
\pgfpathlineto{\pgfqpoint{4.886542in}{2.034443in}}%
\pgfpathlineto{\pgfqpoint{4.887628in}{2.054863in}}%
\pgfpathlineto{\pgfqpoint{4.886926in}{2.005108in}}%
\pgfpathlineto{\pgfqpoint{4.887662in}{2.041898in}}%
\pgfpathlineto{\pgfqpoint{4.888528in}{1.997508in}}%
\pgfpathlineto{\pgfqpoint{4.888806in}{2.025729in}}%
\pgfpathlineto{\pgfqpoint{4.888869in}{2.028995in}}%
\pgfpathlineto{\pgfqpoint{4.889468in}{1.996534in}}%
\pgfpathlineto{\pgfqpoint{4.889916in}{2.027915in}}%
\pgfpathlineto{\pgfqpoint{4.890330in}{2.031273in}}%
\pgfpathlineto{\pgfqpoint{4.891216in}{1.981667in}}%
\pgfpathlineto{\pgfqpoint{4.891990in}{2.012800in}}%
\pgfpathlineto{\pgfqpoint{4.891279in}{1.977610in}}%
\pgfpathlineto{\pgfqpoint{4.892374in}{1.998337in}}%
\pgfpathlineto{\pgfqpoint{4.892408in}{1.992803in}}%
\pgfpathlineto{\pgfqpoint{4.893367in}{2.026167in}}%
\pgfpathlineto{\pgfqpoint{4.893441in}{2.017143in}}%
\pgfpathlineto{\pgfqpoint{4.893514in}{2.012113in}}%
\pgfpathlineto{\pgfqpoint{4.893713in}{2.031213in}}%
\pgfpathlineto{\pgfqpoint{4.894882in}{2.053432in}}%
\pgfpathlineto{\pgfqpoint{4.893849in}{2.017642in}}%
\pgfpathlineto{\pgfqpoint{4.894886in}{2.052315in}}%
\pgfpathlineto{\pgfqpoint{4.894984in}{2.034528in}}%
\pgfpathlineto{\pgfqpoint{4.895553in}{2.059553in}}%
\pgfpathlineto{\pgfqpoint{4.895992in}{2.053914in}}%
\pgfpathlineto{\pgfqpoint{4.896079in}{2.058678in}}%
\pgfpathlineto{\pgfqpoint{4.896683in}{2.002189in}}%
\pgfpathlineto{\pgfqpoint{4.896737in}{2.005208in}}%
\pgfpathlineto{\pgfqpoint{4.896824in}{2.001733in}}%
\pgfpathlineto{\pgfqpoint{4.897584in}{2.033939in}}%
\pgfpathlineto{\pgfqpoint{4.897788in}{2.044505in}}%
\pgfpathlineto{\pgfqpoint{4.898591in}{2.005531in}}%
\pgfpathlineto{\pgfqpoint{4.898664in}{2.014971in}}%
\pgfpathlineto{\pgfqpoint{4.899784in}{2.031807in}}%
\pgfpathlineto{\pgfqpoint{4.898864in}{1.999932in}}%
\pgfpathlineto{\pgfqpoint{4.899813in}{2.026864in}}%
\pgfpathlineto{\pgfqpoint{4.900198in}{2.015189in}}%
\pgfpathlineto{\pgfqpoint{4.899930in}{2.045833in}}%
\pgfpathlineto{\pgfqpoint{4.900875in}{2.036571in}}%
\pgfpathlineto{\pgfqpoint{4.902024in}{2.072983in}}%
\pgfpathlineto{\pgfqpoint{4.901264in}{2.026596in}}%
\pgfpathlineto{\pgfqpoint{4.902038in}{2.071442in}}%
\pgfpathlineto{\pgfqpoint{4.902058in}{2.067224in}}%
\pgfpathlineto{\pgfqpoint{4.902793in}{2.112061in}}%
\pgfpathlineto{\pgfqpoint{4.903119in}{2.080689in}}%
\pgfpathlineto{\pgfqpoint{4.903406in}{2.095557in}}%
\pgfpathlineto{\pgfqpoint{4.903781in}{2.075257in}}%
\pgfpathlineto{\pgfqpoint{4.904234in}{2.083501in}}%
\pgfpathlineto{\pgfqpoint{4.905388in}{2.036424in}}%
\pgfpathlineto{\pgfqpoint{4.904292in}{2.089486in}}%
\pgfpathlineto{\pgfqpoint{4.905422in}{2.042637in}}%
\pgfpathlineto{\pgfqpoint{4.905441in}{2.036754in}}%
\pgfpathlineto{\pgfqpoint{4.905982in}{2.072223in}}%
\pgfpathlineto{\pgfqpoint{4.906459in}{2.059330in}}%
\pgfpathlineto{\pgfqpoint{4.906493in}{2.066168in}}%
\pgfpathlineto{\pgfqpoint{4.906741in}{2.041567in}}%
\pgfpathlineto{\pgfqpoint{4.907579in}{2.063264in}}%
\pgfpathlineto{\pgfqpoint{4.908143in}{2.039006in}}%
\pgfpathlineto{\pgfqpoint{4.908591in}{2.072511in}}%
\pgfpathlineto{\pgfqpoint{4.908596in}{2.071728in}}%
\pgfpathlineto{\pgfqpoint{4.909156in}{2.095360in}}%
\pgfpathlineto{\pgfqpoint{4.909541in}{2.064340in}}%
\pgfpathlineto{\pgfqpoint{4.909721in}{2.074887in}}%
\pgfpathlineto{\pgfqpoint{4.910534in}{2.067294in}}%
\pgfpathlineto{\pgfqpoint{4.910003in}{2.096676in}}%
\pgfpathlineto{\pgfqpoint{4.910806in}{2.084389in}}%
\pgfpathlineto{\pgfqpoint{4.910948in}{2.070337in}}%
\pgfpathlineto{\pgfqpoint{4.911332in}{2.099194in}}%
\pgfpathlineto{\pgfqpoint{4.911396in}{2.096915in}}%
\pgfpathlineto{\pgfqpoint{4.911654in}{2.120872in}}%
\pgfpathlineto{\pgfqpoint{4.912442in}{2.087027in}}%
\pgfpathlineto{\pgfqpoint{4.912452in}{2.087158in}}%
\pgfpathlineto{\pgfqpoint{4.912691in}{2.104953in}}%
\pgfpathlineto{\pgfqpoint{4.913650in}{2.065040in}}%
\pgfpathlineto{\pgfqpoint{4.914779in}{2.091600in}}%
\pgfpathlineto{\pgfqpoint{4.913917in}{2.053619in}}%
\pgfpathlineto{\pgfqpoint{4.914789in}{2.090546in}}%
\pgfpathlineto{\pgfqpoint{4.914808in}{2.086309in}}%
\pgfpathlineto{\pgfqpoint{4.915870in}{2.114130in}}%
\pgfpathlineto{\pgfqpoint{4.916079in}{2.109843in}}%
\pgfpathlineto{\pgfqpoint{4.916858in}{2.158593in}}%
\pgfpathlineto{\pgfqpoint{4.920660in}{2.054755in}}%
\pgfpathlineto{\pgfqpoint{4.917043in}{2.165582in}}%
\pgfpathlineto{\pgfqpoint{4.920850in}{2.065609in}}%
\pgfpathlineto{\pgfqpoint{4.921308in}{2.086271in}}%
\pgfpathlineto{\pgfqpoint{4.921804in}{2.053188in}}%
\pgfpathlineto{\pgfqpoint{4.921975in}{2.068653in}}%
\pgfpathlineto{\pgfqpoint{4.922627in}{2.041502in}}%
\pgfpathlineto{\pgfqpoint{4.923124in}{2.056010in}}%
\pgfpathlineto{\pgfqpoint{4.923888in}{2.104847in}}%
\pgfpathlineto{\pgfqpoint{4.923221in}{2.048879in}}%
\pgfpathlineto{\pgfqpoint{4.924258in}{2.061896in}}%
\pgfpathlineto{\pgfqpoint{4.924589in}{2.038982in}}%
\pgfpathlineto{\pgfqpoint{4.925217in}{2.086558in}}%
\pgfpathlineto{\pgfqpoint{4.925324in}{2.077199in}}%
\pgfpathlineto{\pgfqpoint{4.926381in}{2.127714in}}%
\pgfpathlineto{\pgfqpoint{4.926624in}{2.113066in}}%
\pgfpathlineto{\pgfqpoint{4.927111in}{2.107897in}}%
\pgfpathlineto{\pgfqpoint{4.927432in}{2.137463in}}%
\pgfpathlineto{\pgfqpoint{4.927734in}{2.112321in}}%
\pgfpathlineto{\pgfqpoint{4.928791in}{2.128986in}}%
\pgfpathlineto{\pgfqpoint{4.928119in}{2.087044in}}%
\pgfpathlineto{\pgfqpoint{4.928873in}{2.118394in}}%
\pgfpathlineto{\pgfqpoint{4.929969in}{2.064615in}}%
\pgfpathlineto{\pgfqpoint{4.928898in}{2.118831in}}%
\pgfpathlineto{\pgfqpoint{4.930100in}{2.068530in}}%
\pgfpathlineto{\pgfqpoint{4.930675in}{2.083418in}}%
\pgfpathlineto{\pgfqpoint{4.930280in}{2.059130in}}%
\pgfpathlineto{\pgfqpoint{4.930991in}{2.062403in}}%
\pgfpathlineto{\pgfqpoint{4.931288in}{2.056392in}}%
\pgfpathlineto{\pgfqpoint{4.931887in}{2.077422in}}%
\pgfpathlineto{\pgfqpoint{4.932087in}{2.069901in}}%
\pgfpathlineto{\pgfqpoint{4.932987in}{2.108303in}}%
\pgfpathlineto{\pgfqpoint{4.933279in}{2.104379in}}%
\pgfpathlineto{\pgfqpoint{4.934721in}{2.056757in}}%
\pgfpathlineto{\pgfqpoint{4.933503in}{2.112596in}}%
\pgfpathlineto{\pgfqpoint{4.934730in}{2.058739in}}%
\pgfpathlineto{\pgfqpoint{4.934735in}{2.060959in}}%
\pgfpathlineto{\pgfqpoint{4.935592in}{2.029566in}}%
\pgfpathlineto{\pgfqpoint{4.935806in}{2.043668in}}%
\pgfpathlineto{\pgfqpoint{4.936230in}{2.055652in}}%
\pgfpathlineto{\pgfqpoint{4.936707in}{2.027833in}}%
\pgfpathlineto{\pgfqpoint{4.936843in}{2.037739in}}%
\pgfpathlineto{\pgfqpoint{4.936848in}{2.035102in}}%
\pgfpathlineto{\pgfqpoint{4.937909in}{2.061489in}}%
\pgfpathlineto{\pgfqpoint{4.937997in}{2.069379in}}%
\pgfpathlineto{\pgfqpoint{4.938781in}{2.022081in}}%
\pgfpathlineto{\pgfqpoint{4.938786in}{2.023013in}}%
\pgfpathlineto{\pgfqpoint{4.938985in}{2.006033in}}%
\pgfpathlineto{\pgfqpoint{4.939584in}{2.052616in}}%
\pgfpathlineto{\pgfqpoint{4.939643in}{2.050011in}}%
\pgfpathlineto{\pgfqpoint{4.940816in}{2.081410in}}%
\pgfpathlineto{\pgfqpoint{4.939745in}{2.033631in}}%
\pgfpathlineto{\pgfqpoint{4.940840in}{2.075978in}}%
\pgfpathlineto{\pgfqpoint{4.940981in}{2.061974in}}%
\pgfpathlineto{\pgfqpoint{4.941293in}{2.101952in}}%
\pgfpathlineto{\pgfqpoint{4.941926in}{2.087497in}}%
\pgfpathlineto{\pgfqpoint{4.943124in}{2.148525in}}%
\pgfpathlineto{\pgfqpoint{4.943133in}{2.145157in}}%
\pgfpathlineto{\pgfqpoint{4.943382in}{2.125813in}}%
\pgfpathlineto{\pgfqpoint{4.944190in}{2.157905in}}%
\pgfpathlineto{\pgfqpoint{4.944234in}{2.146881in}}%
\pgfpathlineto{\pgfqpoint{4.944983in}{2.162050in}}%
\pgfpathlineto{\pgfqpoint{4.944268in}{2.138681in}}%
\pgfpathlineto{\pgfqpoint{4.945275in}{2.144674in}}%
\pgfpathlineto{\pgfqpoint{4.945714in}{2.128489in}}%
\pgfpathlineto{\pgfqpoint{4.946283in}{2.178320in}}%
\pgfpathlineto{\pgfqpoint{4.946351in}{2.177091in}}%
\pgfpathlineto{\pgfqpoint{4.948825in}{2.087391in}}%
\pgfpathlineto{\pgfqpoint{4.948976in}{2.101521in}}%
\pgfpathlineto{\pgfqpoint{4.949793in}{2.124393in}}%
\pgfpathlineto{\pgfqpoint{4.949214in}{2.089563in}}%
\pgfpathlineto{\pgfqpoint{4.950105in}{2.104316in}}%
\pgfpathlineto{\pgfqpoint{4.951176in}{2.087658in}}%
\pgfpathlineto{\pgfqpoint{4.950626in}{2.133549in}}%
\pgfpathlineto{\pgfqpoint{4.951239in}{2.090132in}}%
\pgfpathlineto{\pgfqpoint{4.951716in}{2.096689in}}%
\pgfpathlineto{\pgfqpoint{4.952174in}{2.057387in}}%
\pgfpathlineto{\pgfqpoint{4.952281in}{2.061633in}}%
\pgfpathlineto{\pgfqpoint{4.952432in}{2.049822in}}%
\pgfpathlineto{\pgfqpoint{4.953182in}{2.080498in}}%
\pgfpathlineto{\pgfqpoint{4.953289in}{2.074285in}}%
\pgfpathlineto{\pgfqpoint{4.953640in}{2.084590in}}%
\pgfpathlineto{\pgfqpoint{4.954122in}{2.030419in}}%
\pgfpathlineto{\pgfqpoint{4.954263in}{2.039529in}}%
\pgfpathlineto{\pgfqpoint{4.955047in}{2.064547in}}%
\pgfpathlineto{\pgfqpoint{4.954428in}{2.036768in}}%
\pgfpathlineto{\pgfqpoint{4.955382in}{2.040189in}}%
\pgfpathlineto{\pgfqpoint{4.956244in}{2.005661in}}%
\pgfpathlineto{\pgfqpoint{4.955743in}{2.064437in}}%
\pgfpathlineto{\pgfqpoint{4.956561in}{2.016704in}}%
\pgfpathlineto{\pgfqpoint{4.956755in}{2.034978in}}%
\pgfpathlineto{\pgfqpoint{4.957452in}{2.008288in}}%
\pgfpathlineto{\pgfqpoint{4.957656in}{2.010223in}}%
\pgfpathlineto{\pgfqpoint{4.958674in}{1.980224in}}%
\pgfpathlineto{\pgfqpoint{4.957685in}{2.015795in}}%
\pgfpathlineto{\pgfqpoint{4.958898in}{2.007209in}}%
\pgfpathlineto{\pgfqpoint{4.959983in}{2.062104in}}%
\pgfpathlineto{\pgfqpoint{4.959005in}{1.998290in}}%
\pgfpathlineto{\pgfqpoint{4.960085in}{2.049557in}}%
\pgfpathlineto{\pgfqpoint{4.961083in}{2.029553in}}%
\pgfpathlineto{\pgfqpoint{4.960480in}{2.056458in}}%
\pgfpathlineto{\pgfqpoint{4.961283in}{2.043842in}}%
\pgfpathlineto{\pgfqpoint{4.962486in}{2.114954in}}%
\pgfpathlineto{\pgfqpoint{4.962573in}{2.114860in}}%
\pgfpathlineto{\pgfqpoint{4.963002in}{2.103532in}}%
\pgfpathlineto{\pgfqpoint{4.963260in}{2.129928in}}%
\pgfpathlineto{\pgfqpoint{4.963693in}{2.110427in}}%
\pgfpathlineto{\pgfqpoint{4.963859in}{2.130095in}}%
\pgfpathlineto{\pgfqpoint{4.964676in}{2.092965in}}%
\pgfpathlineto{\pgfqpoint{4.964774in}{2.097673in}}%
\pgfpathlineto{\pgfqpoint{4.965733in}{2.082254in}}%
\pgfpathlineto{\pgfqpoint{4.965256in}{2.109248in}}%
\pgfpathlineto{\pgfqpoint{4.965884in}{2.097031in}}%
\pgfpathlineto{\pgfqpoint{4.966030in}{2.113000in}}%
\pgfpathlineto{\pgfqpoint{4.966853in}{2.082921in}}%
\pgfpathlineto{\pgfqpoint{4.966916in}{2.087390in}}%
\pgfpathlineto{\pgfqpoint{4.967369in}{2.067412in}}%
\pgfpathlineto{\pgfqpoint{4.967680in}{2.098328in}}%
\pgfpathlineto{\pgfqpoint{4.968075in}{2.076564in}}%
\pgfpathlineto{\pgfqpoint{4.969209in}{2.103907in}}%
\pgfpathlineto{\pgfqpoint{4.968274in}{2.056258in}}%
\pgfpathlineto{\pgfqpoint{4.969219in}{2.100691in}}%
\pgfpathlineto{\pgfqpoint{4.970256in}{2.093626in}}%
\pgfpathlineto{\pgfqpoint{4.969409in}{2.118045in}}%
\pgfpathlineto{\pgfqpoint{4.970295in}{2.101838in}}%
\pgfpathlineto{\pgfqpoint{4.971750in}{2.163978in}}%
\pgfpathlineto{\pgfqpoint{4.970728in}{2.095027in}}%
\pgfpathlineto{\pgfqpoint{4.971780in}{2.156288in}}%
\pgfpathlineto{\pgfqpoint{4.971911in}{2.144963in}}%
\pgfpathlineto{\pgfqpoint{4.972096in}{2.173150in}}%
\pgfpathlineto{\pgfqpoint{4.972860in}{2.170944in}}%
\pgfpathlineto{\pgfqpoint{4.972875in}{2.171403in}}%
\pgfpathlineto{\pgfqpoint{4.972885in}{2.164507in}}%
\pgfpathlineto{\pgfqpoint{4.973800in}{2.144121in}}%
\pgfpathlineto{\pgfqpoint{4.973557in}{2.164701in}}%
\pgfpathlineto{\pgfqpoint{4.974014in}{2.155260in}}%
\pgfpathlineto{\pgfqpoint{4.975085in}{2.151100in}}%
\pgfpathlineto{\pgfqpoint{4.974715in}{2.186176in}}%
\pgfpathlineto{\pgfqpoint{4.975129in}{2.152162in}}%
\pgfpathlineto{\pgfqpoint{4.976448in}{2.216397in}}%
\pgfpathlineto{\pgfqpoint{4.975387in}{2.143102in}}%
\pgfpathlineto{\pgfqpoint{4.976468in}{2.214728in}}%
\pgfpathlineto{\pgfqpoint{4.976473in}{2.213020in}}%
\pgfpathlineto{\pgfqpoint{4.976999in}{2.254466in}}%
\pgfpathlineto{\pgfqpoint{4.977485in}{2.246705in}}%
\pgfpathlineto{\pgfqpoint{4.978615in}{2.262594in}}%
\pgfpathlineto{\pgfqpoint{4.978235in}{2.243020in}}%
\pgfpathlineto{\pgfqpoint{4.978620in}{2.261264in}}%
\pgfpathlineto{\pgfqpoint{4.979199in}{2.246371in}}%
\pgfpathlineto{\pgfqpoint{4.979394in}{2.268923in}}%
\pgfpathlineto{\pgfqpoint{4.979735in}{2.257008in}}%
\pgfpathlineto{\pgfqpoint{4.980046in}{2.240929in}}%
\pgfpathlineto{\pgfqpoint{4.980606in}{2.282723in}}%
\pgfpathlineto{\pgfqpoint{4.980640in}{2.279109in}}%
\pgfpathlineto{\pgfqpoint{4.981736in}{2.303338in}}%
\pgfpathlineto{\pgfqpoint{4.981205in}{2.256269in}}%
\pgfpathlineto{\pgfqpoint{4.981823in}{2.299878in}}%
\pgfpathlineto{\pgfqpoint{4.981843in}{2.296169in}}%
\pgfpathlineto{\pgfqpoint{4.982622in}{2.332601in}}%
\pgfpathlineto{\pgfqpoint{4.982743in}{2.331269in}}%
\pgfpathlineto{\pgfqpoint{4.983328in}{2.349751in}}%
\pgfpathlineto{\pgfqpoint{4.983576in}{2.316120in}}%
\pgfpathlineto{\pgfqpoint{4.983805in}{2.323996in}}%
\pgfpathlineto{\pgfqpoint{4.983853in}{2.319014in}}%
\pgfpathlineto{\pgfqpoint{4.984832in}{2.369272in}}%
\pgfpathlineto{\pgfqpoint{4.984861in}{2.364872in}}%
\pgfpathlineto{\pgfqpoint{4.984983in}{2.385061in}}%
\pgfpathlineto{\pgfqpoint{4.985889in}{2.357889in}}%
\pgfpathlineto{\pgfqpoint{4.985971in}{2.366020in}}%
\pgfpathlineto{\pgfqpoint{4.986044in}{2.374621in}}%
\pgfpathlineto{\pgfqpoint{4.986984in}{2.349038in}}%
\pgfpathlineto{\pgfqpoint{4.987003in}{2.352089in}}%
\pgfpathlineto{\pgfqpoint{4.987656in}{2.332074in}}%
\pgfpathlineto{\pgfqpoint{4.987296in}{2.364672in}}%
\pgfpathlineto{\pgfqpoint{4.988099in}{2.352986in}}%
\pgfpathlineto{\pgfqpoint{4.989204in}{2.432411in}}%
\pgfpathlineto{\pgfqpoint{4.989418in}{2.412453in}}%
\pgfpathlineto{\pgfqpoint{4.989481in}{2.409465in}}%
\pgfpathlineto{\pgfqpoint{4.989525in}{2.419287in}}%
\pgfpathlineto{\pgfqpoint{4.990212in}{2.425947in}}%
\pgfpathlineto{\pgfqpoint{4.990582in}{2.397427in}}%
\pgfpathlineto{\pgfqpoint{4.990596in}{2.398701in}}%
\pgfpathlineto{\pgfqpoint{4.991813in}{2.344662in}}%
\pgfpathlineto{\pgfqpoint{4.991916in}{2.357339in}}%
\pgfpathlineto{\pgfqpoint{4.992651in}{2.375704in}}%
\pgfpathlineto{\pgfqpoint{4.992271in}{2.337058in}}%
\pgfpathlineto{\pgfqpoint{4.993016in}{2.354115in}}%
\pgfpathlineto{\pgfqpoint{4.994330in}{2.404111in}}%
\pgfpathlineto{\pgfqpoint{4.993074in}{2.346473in}}%
\pgfpathlineto{\pgfqpoint{4.994671in}{2.392574in}}%
\pgfpathlineto{\pgfqpoint{4.995762in}{2.354863in}}%
\pgfpathlineto{\pgfqpoint{4.994710in}{2.396243in}}%
\pgfpathlineto{\pgfqpoint{4.995825in}{2.367756in}}%
\pgfpathlineto{\pgfqpoint{4.996400in}{2.380893in}}%
\pgfpathlineto{\pgfqpoint{4.996721in}{2.356059in}}%
\pgfpathlineto{\pgfqpoint{4.996823in}{2.356489in}}%
\pgfpathlineto{\pgfqpoint{4.997033in}{2.352566in}}%
\pgfpathlineto{\pgfqpoint{4.997797in}{2.382278in}}%
\pgfpathlineto{\pgfqpoint{4.997909in}{2.374072in}}%
\pgfpathlineto{\pgfqpoint{4.997919in}{2.371242in}}%
\pgfpathlineto{\pgfqpoint{4.998668in}{2.402772in}}%
\pgfpathlineto{\pgfqpoint{4.998936in}{2.395844in}}%
\pgfpathlineto{\pgfqpoint{4.999603in}{2.382742in}}%
\pgfpathlineto{\pgfqpoint{4.999004in}{2.410429in}}%
\pgfpathlineto{\pgfqpoint{5.000036in}{2.397013in}}%
\pgfpathlineto{\pgfqpoint{5.000051in}{2.400646in}}%
\pgfpathlineto{\pgfqpoint{5.000440in}{2.357509in}}%
\pgfpathlineto{\pgfqpoint{5.001142in}{2.396649in}}%
\pgfpathlineto{\pgfqpoint{5.001195in}{2.392035in}}%
\pgfpathlineto{\pgfqpoint{5.001604in}{2.432017in}}%
\pgfpathlineto{\pgfqpoint{5.001813in}{2.429840in}}%
\pgfpathlineto{\pgfqpoint{5.003011in}{2.501312in}}%
\pgfpathlineto{\pgfqpoint{5.003133in}{2.492346in}}%
\pgfpathlineto{\pgfqpoint{5.003712in}{2.482294in}}%
\pgfpathlineto{\pgfqpoint{5.003318in}{2.506629in}}%
\pgfpathlineto{\pgfqpoint{5.004223in}{2.501828in}}%
\pgfpathlineto{\pgfqpoint{5.005489in}{2.532172in}}%
\pgfpathlineto{\pgfqpoint{5.004238in}{2.494776in}}%
\pgfpathlineto{\pgfqpoint{5.005499in}{2.531658in}}%
\pgfpathlineto{\pgfqpoint{5.006239in}{2.517082in}}%
\pgfpathlineto{\pgfqpoint{5.005888in}{2.558902in}}%
\pgfpathlineto{\pgfqpoint{5.006643in}{2.519225in}}%
\pgfpathlineto{\pgfqpoint{5.007013in}{2.552327in}}%
\pgfpathlineto{\pgfqpoint{5.007656in}{2.511037in}}%
\pgfpathlineto{\pgfqpoint{5.007680in}{2.511503in}}%
\pgfpathlineto{\pgfqpoint{5.008551in}{2.474764in}}%
\pgfpathlineto{\pgfqpoint{5.007816in}{2.528207in}}%
\pgfpathlineto{\pgfqpoint{5.008819in}{2.487995in}}%
\pgfpathlineto{\pgfqpoint{5.009077in}{2.471565in}}%
\pgfpathlineto{\pgfqpoint{5.009379in}{2.500688in}}%
\pgfpathlineto{\pgfqpoint{5.009885in}{2.496490in}}%
\pgfpathlineto{\pgfqpoint{5.010212in}{2.505546in}}%
\pgfpathlineto{\pgfqpoint{5.010557in}{2.452988in}}%
\pgfpathlineto{\pgfqpoint{5.010694in}{2.436070in}}%
\pgfpathlineto{\pgfqpoint{5.011502in}{2.474130in}}%
\pgfpathlineto{\pgfqpoint{5.011667in}{2.449710in}}%
\pgfpathlineto{\pgfqpoint{5.012680in}{2.468433in}}%
\pgfpathlineto{\pgfqpoint{5.012310in}{2.432804in}}%
\pgfpathlineto{\pgfqpoint{5.012787in}{2.457646in}}%
\pgfpathlineto{\pgfqpoint{5.013756in}{2.407843in}}%
\pgfpathlineto{\pgfqpoint{5.012982in}{2.460452in}}%
\pgfpathlineto{\pgfqpoint{5.013955in}{2.424562in}}%
\pgfpathlineto{\pgfqpoint{5.014671in}{2.407751in}}%
\pgfpathlineto{\pgfqpoint{5.014340in}{2.435040in}}%
\pgfpathlineto{\pgfqpoint{5.014944in}{2.428760in}}%
\pgfpathlineto{\pgfqpoint{5.016307in}{2.503573in}}%
\pgfpathlineto{\pgfqpoint{5.015027in}{2.424028in}}%
\pgfpathlineto{\pgfqpoint{5.016380in}{2.498480in}}%
\pgfpathlineto{\pgfqpoint{5.016414in}{2.492559in}}%
\pgfpathlineto{\pgfqpoint{5.017315in}{2.570069in}}%
\pgfpathlineto{\pgfqpoint{5.017383in}{2.560974in}}%
\pgfpathlineto{\pgfqpoint{5.018527in}{2.597878in}}%
\pgfpathlineto{\pgfqpoint{5.017427in}{2.557018in}}%
\pgfpathlineto{\pgfqpoint{5.018532in}{2.595924in}}%
\pgfpathlineto{\pgfqpoint{5.019048in}{2.576159in}}%
\pgfpathlineto{\pgfqpoint{5.018775in}{2.603086in}}%
\pgfpathlineto{\pgfqpoint{5.019632in}{2.596975in}}%
\pgfpathlineto{\pgfqpoint{5.020031in}{2.572931in}}%
\pgfpathlineto{\pgfqpoint{5.020786in}{2.610792in}}%
\pgfpathlineto{\pgfqpoint{5.020801in}{2.607323in}}%
\pgfpathlineto{\pgfqpoint{5.021570in}{2.640034in}}%
\pgfpathlineto{\pgfqpoint{5.021818in}{2.626219in}}%
\pgfpathlineto{\pgfqpoint{5.022777in}{2.637327in}}%
\pgfpathlineto{\pgfqpoint{5.022490in}{2.607639in}}%
\pgfpathlineto{\pgfqpoint{5.022933in}{2.628059in}}%
\pgfpathlineto{\pgfqpoint{5.023468in}{2.644744in}}%
\pgfpathlineto{\pgfqpoint{5.023585in}{2.626179in}}%
\pgfpathlineto{\pgfqpoint{5.023907in}{2.627883in}}%
\pgfpathlineto{\pgfqpoint{5.024238in}{2.616399in}}%
\pgfpathlineto{\pgfqpoint{5.024958in}{2.642945in}}%
\pgfpathlineto{\pgfqpoint{5.025012in}{2.630627in}}%
\pgfpathlineto{\pgfqpoint{5.025017in}{2.630631in}}%
\pgfpathlineto{\pgfqpoint{5.025693in}{2.605473in}}%
\pgfpathlineto{\pgfqpoint{5.025844in}{2.635386in}}%
\pgfpathlineto{\pgfqpoint{5.026175in}{2.610501in}}%
\pgfpathlineto{\pgfqpoint{5.027125in}{2.635286in}}%
\pgfpathlineto{\pgfqpoint{5.026502in}{2.602105in}}%
\pgfpathlineto{\pgfqpoint{5.027295in}{2.616067in}}%
\pgfpathlineto{\pgfqpoint{5.027300in}{2.615131in}}%
\pgfpathlineto{\pgfqpoint{5.028196in}{2.650947in}}%
\pgfpathlineto{\pgfqpoint{5.028210in}{2.649487in}}%
\pgfpathlineto{\pgfqpoint{5.028566in}{2.664259in}}%
\pgfpathlineto{\pgfqpoint{5.029272in}{2.615685in}}%
\pgfpathlineto{\pgfqpoint{5.029277in}{2.617729in}}%
\pgfpathlineto{\pgfqpoint{5.029832in}{2.632403in}}%
\pgfpathlineto{\pgfqpoint{5.030299in}{2.612413in}}%
\pgfpathlineto{\pgfqpoint{5.030372in}{2.619504in}}%
\pgfpathlineto{\pgfqpoint{5.030416in}{2.611102in}}%
\pgfpathlineto{\pgfqpoint{5.031307in}{2.675455in}}%
\pgfpathlineto{\pgfqpoint{5.031424in}{2.664327in}}%
\pgfpathlineto{\pgfqpoint{5.031613in}{2.684842in}}%
\pgfpathlineto{\pgfqpoint{5.032573in}{2.675791in}}%
\pgfpathlineto{\pgfqpoint{5.032646in}{2.663495in}}%
\pgfpathlineto{\pgfqpoint{5.033210in}{2.693499in}}%
\pgfpathlineto{\pgfqpoint{5.033663in}{2.688491in}}%
\pgfpathlineto{\pgfqpoint{5.033717in}{2.701101in}}%
\pgfpathlineto{\pgfqpoint{5.034481in}{2.641678in}}%
\pgfpathlineto{\pgfqpoint{5.034525in}{2.651838in}}%
\pgfpathlineto{\pgfqpoint{5.035430in}{2.624412in}}%
\pgfpathlineto{\pgfqpoint{5.034603in}{2.658209in}}%
\pgfpathlineto{\pgfqpoint{5.035649in}{2.642570in}}%
\pgfpathlineto{\pgfqpoint{5.036214in}{2.616835in}}%
\pgfpathlineto{\pgfqpoint{5.036686in}{2.663392in}}%
\pgfpathlineto{\pgfqpoint{5.036706in}{2.659716in}}%
\pgfpathlineto{\pgfqpoint{5.037514in}{2.720593in}}%
\pgfpathlineto{\pgfqpoint{5.037991in}{2.692667in}}%
\pgfpathlineto{\pgfqpoint{5.038225in}{2.675542in}}%
\pgfpathlineto{\pgfqpoint{5.038765in}{2.695979in}}%
\pgfpathlineto{\pgfqpoint{5.039101in}{2.688612in}}%
\pgfpathlineto{\pgfqpoint{5.040245in}{2.717106in}}%
\pgfpathlineto{\pgfqpoint{5.039695in}{2.677292in}}%
\pgfpathlineto{\pgfqpoint{5.040265in}{2.713822in}}%
\pgfpathlineto{\pgfqpoint{5.040503in}{2.729074in}}%
\pgfpathlineto{\pgfqpoint{5.041092in}{2.690264in}}%
\pgfpathlineto{\pgfqpoint{5.041234in}{2.670830in}}%
\pgfpathlineto{\pgfqpoint{5.041828in}{2.723816in}}%
\pgfpathlineto{\pgfqpoint{5.042139in}{2.713114in}}%
\pgfpathlineto{\pgfqpoint{5.042607in}{2.698256in}}%
\pgfpathlineto{\pgfqpoint{5.042407in}{2.725567in}}%
\pgfpathlineto{\pgfqpoint{5.043225in}{2.712774in}}%
\pgfpathlineto{\pgfqpoint{5.044169in}{2.765217in}}%
\pgfpathlineto{\pgfqpoint{5.044364in}{2.751233in}}%
\pgfpathlineto{\pgfqpoint{5.045153in}{2.741283in}}%
\pgfpathlineto{\pgfqpoint{5.045382in}{2.766187in}}%
\pgfpathlineto{\pgfqpoint{5.045455in}{2.761352in}}%
\pgfpathlineto{\pgfqpoint{5.045518in}{2.750127in}}%
\pgfpathlineto{\pgfqpoint{5.046477in}{2.782129in}}%
\pgfpathlineto{\pgfqpoint{5.046492in}{2.779111in}}%
\pgfpathlineto{\pgfqpoint{5.046545in}{2.789400in}}%
\pgfpathlineto{\pgfqpoint{5.046808in}{2.763046in}}%
\pgfpathlineto{\pgfqpoint{5.047538in}{2.768645in}}%
\pgfpathlineto{\pgfqpoint{5.048575in}{2.750591in}}%
\pgfpathlineto{\pgfqpoint{5.048113in}{2.790256in}}%
\pgfpathlineto{\pgfqpoint{5.048663in}{2.755260in}}%
\pgfpathlineto{\pgfqpoint{5.049160in}{2.785863in}}%
\pgfpathlineto{\pgfqpoint{5.048702in}{2.745847in}}%
\pgfpathlineto{\pgfqpoint{5.049856in}{2.776307in}}%
\pgfpathlineto{\pgfqpoint{5.050557in}{2.814623in}}%
\pgfpathlineto{\pgfqpoint{5.049953in}{2.774140in}}%
\pgfpathlineto{\pgfqpoint{5.051126in}{2.801828in}}%
\pgfpathlineto{\pgfqpoint{5.052076in}{2.745785in}}%
\pgfpathlineto{\pgfqpoint{5.051433in}{2.802551in}}%
\pgfpathlineto{\pgfqpoint{5.052305in}{2.750877in}}%
\pgfpathlineto{\pgfqpoint{5.052597in}{2.770762in}}%
\pgfpathlineto{\pgfqpoint{5.053269in}{2.718002in}}%
\pgfpathlineto{\pgfqpoint{5.054271in}{2.690005in}}%
\pgfpathlineto{\pgfqpoint{5.053658in}{2.739129in}}%
\pgfpathlineto{\pgfqpoint{5.054486in}{2.697438in}}%
\pgfpathlineto{\pgfqpoint{5.055459in}{2.722806in}}%
\pgfpathlineto{\pgfqpoint{5.054792in}{2.689949in}}%
\pgfpathlineto{\pgfqpoint{5.055625in}{2.703532in}}%
\pgfpathlineto{\pgfqpoint{5.056321in}{2.690593in}}%
\pgfpathlineto{\pgfqpoint{5.056063in}{2.715714in}}%
\pgfpathlineto{\pgfqpoint{5.056730in}{2.703593in}}%
\pgfpathlineto{\pgfqpoint{5.056750in}{2.708560in}}%
\pgfpathlineto{\pgfqpoint{5.057236in}{2.682223in}}%
\pgfpathlineto{\pgfqpoint{5.057826in}{2.701631in}}%
\pgfpathlineto{\pgfqpoint{5.058838in}{2.674511in}}%
\pgfpathlineto{\pgfqpoint{5.057899in}{2.709302in}}%
\pgfpathlineto{\pgfqpoint{5.058960in}{2.685854in}}%
\pgfpathlineto{\pgfqpoint{5.059651in}{2.698522in}}%
\pgfpathlineto{\pgfqpoint{5.059301in}{2.662412in}}%
\pgfpathlineto{\pgfqpoint{5.060119in}{2.694497in}}%
\pgfpathlineto{\pgfqpoint{5.061151in}{2.681459in}}%
\pgfpathlineto{\pgfqpoint{5.060318in}{2.722313in}}%
\pgfpathlineto{\pgfqpoint{5.061229in}{2.693045in}}%
\pgfpathlineto{\pgfqpoint{5.062042in}{2.696977in}}%
\pgfpathlineto{\pgfqpoint{5.061521in}{2.668221in}}%
\pgfpathlineto{\pgfqpoint{5.062227in}{2.681064in}}%
\pgfpathlineto{\pgfqpoint{5.062275in}{2.673079in}}%
\pgfpathlineto{\pgfqpoint{5.062665in}{2.711177in}}%
\pgfpathlineto{\pgfqpoint{5.063317in}{2.693166in}}%
\pgfpathlineto{\pgfqpoint{5.063794in}{2.679534in}}%
\pgfpathlineto{\pgfqpoint{5.064062in}{2.711617in}}%
\pgfpathlineto{\pgfqpoint{5.064452in}{2.684961in}}%
\pgfpathlineto{\pgfqpoint{5.064612in}{2.675370in}}%
\pgfpathlineto{\pgfqpoint{5.065172in}{2.717389in}}%
\pgfpathlineto{\pgfqpoint{5.065391in}{2.708928in}}%
\pgfpathlineto{\pgfqpoint{5.066102in}{2.719303in}}%
\pgfpathlineto{\pgfqpoint{5.065717in}{2.687722in}}%
\pgfpathlineto{\pgfqpoint{5.066443in}{2.703485in}}%
\pgfpathlineto{\pgfqpoint{5.067514in}{2.686557in}}%
\pgfpathlineto{\pgfqpoint{5.067095in}{2.723506in}}%
\pgfpathlineto{\pgfqpoint{5.067611in}{2.694899in}}%
\pgfpathlineto{\pgfqpoint{5.067694in}{2.697110in}}%
\pgfpathlineto{\pgfqpoint{5.067704in}{2.692085in}}%
\pgfpathlineto{\pgfqpoint{5.067718in}{2.688376in}}%
\pgfpathlineto{\pgfqpoint{5.068492in}{2.730254in}}%
\pgfpathlineto{\pgfqpoint{5.068736in}{2.723761in}}%
\pgfpathlineto{\pgfqpoint{5.069320in}{2.737831in}}%
\pgfpathlineto{\pgfqpoint{5.069607in}{2.701210in}}%
\pgfpathlineto{\pgfqpoint{5.069675in}{2.703757in}}%
\pgfpathlineto{\pgfqpoint{5.070245in}{2.675704in}}%
\pgfpathlineto{\pgfqpoint{5.069724in}{2.707640in}}%
\pgfpathlineto{\pgfqpoint{5.070800in}{2.689171in}}%
\pgfpathlineto{\pgfqpoint{5.070834in}{2.690805in}}%
\pgfpathlineto{\pgfqpoint{5.071175in}{2.665265in}}%
\pgfpathlineto{\pgfqpoint{5.071214in}{2.667521in}}%
\pgfpathlineto{\pgfqpoint{5.071949in}{2.635649in}}%
\pgfpathlineto{\pgfqpoint{5.071248in}{2.667912in}}%
\pgfpathlineto{\pgfqpoint{5.072329in}{2.664076in}}%
\pgfpathlineto{\pgfqpoint{5.072426in}{2.650864in}}%
\pgfpathlineto{\pgfqpoint{5.073590in}{2.718379in}}%
\pgfpathlineto{\pgfqpoint{5.074461in}{2.669357in}}%
\pgfpathlineto{\pgfqpoint{5.073599in}{2.720499in}}%
\pgfpathlineto{\pgfqpoint{5.074787in}{2.689938in}}%
\pgfpathlineto{\pgfqpoint{5.075678in}{2.727850in}}%
\pgfpathlineto{\pgfqpoint{5.074885in}{2.677402in}}%
\pgfpathlineto{\pgfqpoint{5.075946in}{2.724365in}}%
\pgfpathlineto{\pgfqpoint{5.076233in}{2.678527in}}%
\pgfpathlineto{\pgfqpoint{5.077110in}{2.698711in}}%
\pgfpathlineto{\pgfqpoint{5.077845in}{2.717212in}}%
\pgfpathlineto{\pgfqpoint{5.077504in}{2.681537in}}%
\pgfpathlineto{\pgfqpoint{5.078234in}{2.712505in}}%
\pgfpathlineto{\pgfqpoint{5.078259in}{2.720299in}}%
\pgfpathlineto{\pgfqpoint{5.078814in}{2.692388in}}%
\pgfpathlineto{\pgfqpoint{5.079286in}{2.712363in}}%
\pgfpathlineto{\pgfqpoint{5.079373in}{2.696930in}}%
\pgfpathlineto{\pgfqpoint{5.079597in}{2.724277in}}%
\pgfpathlineto{\pgfqpoint{5.080391in}{2.712701in}}%
\pgfpathlineto{\pgfqpoint{5.081267in}{2.737693in}}%
\pgfpathlineto{\pgfqpoint{5.080625in}{2.707719in}}%
\pgfpathlineto{\pgfqpoint{5.081530in}{2.729819in}}%
\pgfpathlineto{\pgfqpoint{5.081671in}{2.715924in}}%
\pgfpathlineto{\pgfqpoint{5.082280in}{2.745201in}}%
\pgfpathlineto{\pgfqpoint{5.082631in}{2.731464in}}%
\pgfpathlineto{\pgfqpoint{5.083663in}{2.756353in}}%
\pgfpathlineto{\pgfqpoint{5.082840in}{2.721194in}}%
\pgfpathlineto{\pgfqpoint{5.083857in}{2.747270in}}%
\pgfpathlineto{\pgfqpoint{5.083969in}{2.745020in}}%
\pgfpathlineto{\pgfqpoint{5.084300in}{2.783537in}}%
\pgfpathlineto{\pgfqpoint{5.084890in}{2.816770in}}%
\pgfpathlineto{\pgfqpoint{5.084437in}{2.767424in}}%
\pgfpathlineto{\pgfqpoint{5.085391in}{2.770144in}}%
\pgfpathlineto{\pgfqpoint{5.085410in}{2.762339in}}%
\pgfpathlineto{\pgfqpoint{5.086262in}{2.842054in}}%
\pgfpathlineto{\pgfqpoint{5.086374in}{2.831852in}}%
\pgfpathlineto{\pgfqpoint{5.086934in}{2.809956in}}%
\pgfpathlineto{\pgfqpoint{5.086701in}{2.843708in}}%
\pgfpathlineto{\pgfqpoint{5.087470in}{2.835933in}}%
\pgfpathlineto{\pgfqpoint{5.087582in}{2.848801in}}%
\pgfpathlineto{\pgfqpoint{5.088512in}{2.822272in}}%
\pgfpathlineto{\pgfqpoint{5.088536in}{2.824462in}}%
\pgfpathlineto{\pgfqpoint{5.088711in}{2.822192in}}%
\pgfpathlineto{\pgfqpoint{5.088969in}{2.854614in}}%
\pgfpathlineto{\pgfqpoint{5.089130in}{2.863826in}}%
\pgfpathlineto{\pgfqpoint{5.089919in}{2.814014in}}%
\pgfpathlineto{\pgfqpoint{5.090026in}{2.816854in}}%
\pgfpathlineto{\pgfqpoint{5.091097in}{2.793663in}}%
\pgfpathlineto{\pgfqpoint{5.090323in}{2.832112in}}%
\pgfpathlineto{\pgfqpoint{5.091170in}{2.802387in}}%
\pgfpathlineto{\pgfqpoint{5.092285in}{2.822400in}}%
\pgfpathlineto{\pgfqpoint{5.091584in}{2.779526in}}%
\pgfpathlineto{\pgfqpoint{5.092338in}{2.816677in}}%
\pgfpathlineto{\pgfqpoint{5.092616in}{2.825005in}}%
\pgfpathlineto{\pgfqpoint{5.092426in}{2.802905in}}%
\pgfpathlineto{\pgfqpoint{5.093200in}{2.811467in}}%
\pgfpathlineto{\pgfqpoint{5.093346in}{2.791930in}}%
\pgfpathlineto{\pgfqpoint{5.094072in}{2.839723in}}%
\pgfpathlineto{\pgfqpoint{5.094222in}{2.831877in}}%
\pgfpathlineto{\pgfqpoint{5.095337in}{2.865936in}}%
\pgfpathlineto{\pgfqpoint{5.095513in}{2.845002in}}%
\pgfpathlineto{\pgfqpoint{5.096681in}{2.802275in}}%
\pgfpathlineto{\pgfqpoint{5.095795in}{2.851455in}}%
\pgfpathlineto{\pgfqpoint{5.096691in}{2.807694in}}%
\pgfpathlineto{\pgfqpoint{5.097455in}{2.738409in}}%
\pgfpathlineto{\pgfqpoint{5.096759in}{2.810381in}}%
\pgfpathlineto{\pgfqpoint{5.097893in}{2.764511in}}%
\pgfpathlineto{\pgfqpoint{5.098804in}{2.787650in}}%
\pgfpathlineto{\pgfqpoint{5.098117in}{2.762023in}}%
\pgfpathlineto{\pgfqpoint{5.098945in}{2.762788in}}%
\pgfpathlineto{\pgfqpoint{5.100177in}{2.711211in}}%
\pgfpathlineto{\pgfqpoint{5.100298in}{2.724570in}}%
\pgfpathlineto{\pgfqpoint{5.100980in}{2.685946in}}%
\pgfpathlineto{\pgfqpoint{5.101248in}{2.694988in}}%
\pgfpathlineto{\pgfqpoint{5.101861in}{2.711882in}}%
\pgfpathlineto{\pgfqpoint{5.102070in}{2.681975in}}%
\pgfpathlineto{\pgfqpoint{5.103044in}{2.667002in}}%
\pgfpathlineto{\pgfqpoint{5.102333in}{2.691333in}}%
\pgfpathlineto{\pgfqpoint{5.103190in}{2.674719in}}%
\pgfpathlineto{\pgfqpoint{5.104028in}{2.722071in}}%
\pgfpathlineto{\pgfqpoint{5.104354in}{2.699217in}}%
\pgfpathlineto{\pgfqpoint{5.106131in}{2.756908in}}%
\pgfpathlineto{\pgfqpoint{5.104738in}{2.699007in}}%
\pgfpathlineto{\pgfqpoint{5.106199in}{2.748414in}}%
\pgfpathlineto{\pgfqpoint{5.106340in}{2.739696in}}%
\pgfpathlineto{\pgfqpoint{5.106929in}{2.819337in}}%
\pgfpathlineto{\pgfqpoint{5.107212in}{2.798468in}}%
\pgfpathlineto{\pgfqpoint{5.107216in}{2.798002in}}%
\pgfpathlineto{\pgfqpoint{5.107830in}{2.842615in}}%
\pgfpathlineto{\pgfqpoint{5.107854in}{2.841548in}}%
\pgfpathlineto{\pgfqpoint{5.108146in}{2.855517in}}%
\pgfpathlineto{\pgfqpoint{5.108731in}{2.824779in}}%
\pgfpathlineto{\pgfqpoint{5.108950in}{2.835568in}}%
\pgfpathlineto{\pgfqpoint{5.109490in}{2.829675in}}%
\pgfpathlineto{\pgfqpoint{5.110011in}{2.856914in}}%
\pgfpathlineto{\pgfqpoint{5.110016in}{2.856597in}}%
\pgfpathlineto{\pgfqpoint{5.110035in}{2.853812in}}%
\pgfpathlineto{\pgfqpoint{5.110420in}{2.879827in}}%
\pgfpathlineto{\pgfqpoint{5.111106in}{2.861312in}}%
\pgfpathlineto{\pgfqpoint{5.112319in}{2.917059in}}%
\pgfpathlineto{\pgfqpoint{5.112333in}{2.912251in}}%
\pgfpathlineto{\pgfqpoint{5.112557in}{2.897670in}}%
\pgfpathlineto{\pgfqpoint{5.112762in}{2.923198in}}%
\pgfpathlineto{\pgfqpoint{5.113297in}{2.914630in}}%
\pgfpathlineto{\pgfqpoint{5.113974in}{2.927878in}}%
\pgfpathlineto{\pgfqpoint{5.114130in}{2.902942in}}%
\pgfpathlineto{\pgfqpoint{5.114407in}{2.914826in}}%
\pgfpathlineto{\pgfqpoint{5.115658in}{2.881707in}}%
\pgfpathlineto{\pgfqpoint{5.115030in}{2.930884in}}%
\pgfpathlineto{\pgfqpoint{5.115688in}{2.884348in}}%
\pgfpathlineto{\pgfqpoint{5.116496in}{2.927095in}}%
\pgfpathlineto{\pgfqpoint{5.116866in}{2.909209in}}%
\pgfpathlineto{\pgfqpoint{5.117883in}{2.883435in}}%
\pgfpathlineto{\pgfqpoint{5.116983in}{2.918542in}}%
\pgfpathlineto{\pgfqpoint{5.117986in}{2.905455in}}%
\pgfpathlineto{\pgfqpoint{5.118015in}{2.904169in}}%
\pgfpathlineto{\pgfqpoint{5.118395in}{2.926075in}}%
\pgfpathlineto{\pgfqpoint{5.118404in}{2.925892in}}%
\pgfpathlineto{\pgfqpoint{5.119529in}{2.943396in}}%
\pgfpathlineto{\pgfqpoint{5.119144in}{2.915601in}}%
\pgfpathlineto{\pgfqpoint{5.119539in}{2.940533in}}%
\pgfpathlineto{\pgfqpoint{5.120395in}{2.916288in}}%
\pgfpathlineto{\pgfqpoint{5.120610in}{2.944156in}}%
\pgfpathlineto{\pgfqpoint{5.121647in}{2.962994in}}%
\pgfpathlineto{\pgfqpoint{5.120843in}{2.913075in}}%
\pgfpathlineto{\pgfqpoint{5.121749in}{2.957650in}}%
\pgfpathlineto{\pgfqpoint{5.121812in}{2.953216in}}%
\pgfpathlineto{\pgfqpoint{5.122567in}{2.998139in}}%
\pgfpathlineto{\pgfqpoint{5.122650in}{2.991068in}}%
\pgfpathlineto{\pgfqpoint{5.123852in}{3.022808in}}%
\pgfpathlineto{\pgfqpoint{5.122927in}{2.975018in}}%
\pgfpathlineto{\pgfqpoint{5.123857in}{3.022762in}}%
\pgfpathlineto{\pgfqpoint{5.125001in}{2.988570in}}%
\pgfpathlineto{\pgfqpoint{5.124451in}{3.038115in}}%
\pgfpathlineto{\pgfqpoint{5.125030in}{2.991028in}}%
\pgfpathlineto{\pgfqpoint{5.125702in}{3.021969in}}%
\pgfpathlineto{\pgfqpoint{5.125045in}{2.987097in}}%
\pgfpathlineto{\pgfqpoint{5.126243in}{3.009339in}}%
\pgfpathlineto{\pgfqpoint{5.127421in}{2.973328in}}%
\pgfpathlineto{\pgfqpoint{5.126452in}{3.018030in}}%
\pgfpathlineto{\pgfqpoint{5.127469in}{2.979870in}}%
\pgfpathlineto{\pgfqpoint{5.128146in}{2.998085in}}%
\pgfpathlineto{\pgfqpoint{5.128487in}{2.959312in}}%
\pgfpathlineto{\pgfqpoint{5.128818in}{2.929655in}}%
\pgfpathlineto{\pgfqpoint{5.129543in}{2.961787in}}%
\pgfpathlineto{\pgfqpoint{5.129548in}{2.961580in}}%
\pgfpathlineto{\pgfqpoint{5.129641in}{2.970857in}}%
\pgfpathlineto{\pgfqpoint{5.130098in}{2.939414in}}%
\pgfpathlineto{\pgfqpoint{5.130629in}{2.956537in}}%
\pgfpathlineto{\pgfqpoint{5.131403in}{2.932853in}}%
\pgfpathlineto{\pgfqpoint{5.130668in}{2.962071in}}%
\pgfpathlineto{\pgfqpoint{5.131734in}{2.957295in}}%
\pgfpathlineto{\pgfqpoint{5.132791in}{2.974313in}}%
\pgfpathlineto{\pgfqpoint{5.132124in}{2.933261in}}%
\pgfpathlineto{\pgfqpoint{5.132859in}{2.963215in}}%
\pgfpathlineto{\pgfqpoint{5.133613in}{2.946877in}}%
\pgfpathlineto{\pgfqpoint{5.133906in}{2.980388in}}%
\pgfpathlineto{\pgfqpoint{5.133945in}{2.973648in}}%
\pgfpathlineto{\pgfqpoint{5.134821in}{2.999667in}}%
\pgfpathlineto{\pgfqpoint{5.133974in}{2.972819in}}%
\pgfpathlineto{\pgfqpoint{5.135074in}{2.979437in}}%
\pgfpathlineto{\pgfqpoint{5.136350in}{2.949392in}}%
\pgfpathlineto{\pgfqpoint{5.135123in}{2.988110in}}%
\pgfpathlineto{\pgfqpoint{5.136403in}{2.956282in}}%
\pgfpathlineto{\pgfqpoint{5.136632in}{2.980971in}}%
\pgfpathlineto{\pgfqpoint{5.137090in}{2.922135in}}%
\pgfpathlineto{\pgfqpoint{5.137094in}{2.922303in}}%
\pgfpathlineto{\pgfqpoint{5.138204in}{2.863346in}}%
\pgfpathlineto{\pgfqpoint{5.138365in}{2.882129in}}%
\pgfpathlineto{\pgfqpoint{5.139008in}{2.902340in}}%
\pgfpathlineto{\pgfqpoint{5.139276in}{2.869237in}}%
\pgfpathlineto{\pgfqpoint{5.139470in}{2.882234in}}%
\pgfpathlineto{\pgfqpoint{5.139928in}{2.851618in}}%
\pgfpathlineto{\pgfqpoint{5.140351in}{2.886765in}}%
\pgfpathlineto{\pgfqpoint{5.140629in}{2.858337in}}%
\pgfpathlineto{\pgfqpoint{5.141530in}{2.883544in}}%
\pgfpathlineto{\pgfqpoint{5.140683in}{2.850980in}}%
\pgfpathlineto{\pgfqpoint{5.141885in}{2.879514in}}%
\pgfpathlineto{\pgfqpoint{5.143029in}{2.862542in}}%
\pgfpathlineto{\pgfqpoint{5.142148in}{2.884256in}}%
\pgfpathlineto{\pgfqpoint{5.143034in}{2.862748in}}%
\pgfpathlineto{\pgfqpoint{5.144144in}{2.937274in}}%
\pgfpathlineto{\pgfqpoint{5.143068in}{2.858845in}}%
\pgfpathlineto{\pgfqpoint{5.144358in}{2.934708in}}%
\pgfpathlineto{\pgfqpoint{5.145064in}{2.911197in}}%
\pgfpathlineto{\pgfqpoint{5.145371in}{2.938794in}}%
\pgfpathlineto{\pgfqpoint{5.145459in}{2.936465in}}%
\pgfpathlineto{\pgfqpoint{5.145497in}{2.941678in}}%
\pgfpathlineto{\pgfqpoint{5.146476in}{2.873666in}}%
\pgfpathlineto{\pgfqpoint{5.146486in}{2.876335in}}%
\pgfpathlineto{\pgfqpoint{5.146793in}{2.882580in}}%
\pgfpathlineto{\pgfqpoint{5.147542in}{2.854834in}}%
\pgfpathlineto{\pgfqpoint{5.147951in}{2.873929in}}%
\pgfpathlineto{\pgfqpoint{5.148156in}{2.851627in}}%
\pgfpathlineto{\pgfqpoint{5.148555in}{2.855240in}}%
\pgfpathlineto{\pgfqpoint{5.148901in}{2.845367in}}%
\pgfpathlineto{\pgfqpoint{5.149456in}{2.875515in}}%
\pgfpathlineto{\pgfqpoint{5.149660in}{2.856961in}}%
\pgfpathlineto{\pgfqpoint{5.149928in}{2.845483in}}%
\pgfpathlineto{\pgfqpoint{5.150137in}{2.869410in}}%
\pgfpathlineto{\pgfqpoint{5.151203in}{2.887214in}}%
\pgfpathlineto{\pgfqpoint{5.150965in}{2.864798in}}%
\pgfpathlineto{\pgfqpoint{5.151276in}{2.883295in}}%
\pgfpathlineto{\pgfqpoint{5.152021in}{2.904922in}}%
\pgfpathlineto{\pgfqpoint{5.151345in}{2.875381in}}%
\pgfpathlineto{\pgfqpoint{5.152440in}{2.893693in}}%
\pgfpathlineto{\pgfqpoint{5.153331in}{2.848613in}}%
\pgfpathlineto{\pgfqpoint{5.153589in}{2.860101in}}%
\pgfpathlineto{\pgfqpoint{5.153677in}{2.873688in}}%
\pgfpathlineto{\pgfqpoint{5.154358in}{2.834916in}}%
\pgfpathlineto{\pgfqpoint{5.155497in}{2.783900in}}%
\pgfpathlineto{\pgfqpoint{5.155580in}{2.785644in}}%
\pgfpathlineto{\pgfqpoint{5.155979in}{2.804948in}}%
\pgfpathlineto{\pgfqpoint{5.156452in}{2.772375in}}%
\pgfpathlineto{\pgfqpoint{5.156671in}{2.779244in}}%
\pgfpathlineto{\pgfqpoint{5.156680in}{2.777171in}}%
\pgfpathlineto{\pgfqpoint{5.157552in}{2.809649in}}%
\pgfpathlineto{\pgfqpoint{5.157688in}{2.801903in}}%
\pgfpathlineto{\pgfqpoint{5.157786in}{2.821083in}}%
\pgfpathlineto{\pgfqpoint{5.158394in}{2.780309in}}%
\pgfpathlineto{\pgfqpoint{5.158774in}{2.787531in}}%
\pgfpathlineto{\pgfqpoint{5.158793in}{2.785747in}}%
\pgfpathlineto{\pgfqpoint{5.159275in}{2.836198in}}%
\pgfpathlineto{\pgfqpoint{5.159368in}{2.828751in}}%
\pgfpathlineto{\pgfqpoint{5.160088in}{2.873412in}}%
\pgfpathlineto{\pgfqpoint{5.159417in}{2.821762in}}%
\pgfpathlineto{\pgfqpoint{5.160497in}{2.845099in}}%
\pgfpathlineto{\pgfqpoint{5.161899in}{2.798652in}}%
\pgfpathlineto{\pgfqpoint{5.161354in}{2.855319in}}%
\pgfpathlineto{\pgfqpoint{5.161904in}{2.800367in}}%
\pgfpathlineto{\pgfqpoint{5.162235in}{2.824397in}}%
\pgfpathlineto{\pgfqpoint{5.162990in}{2.793047in}}%
\pgfpathlineto{\pgfqpoint{5.162995in}{2.797051in}}%
\pgfpathlineto{\pgfqpoint{5.163121in}{2.770474in}}%
\pgfpathlineto{\pgfqpoint{5.163939in}{2.800354in}}%
\pgfpathlineto{\pgfqpoint{5.164100in}{2.797925in}}%
\pgfpathlineto{\pgfqpoint{5.164120in}{2.799577in}}%
\pgfpathlineto{\pgfqpoint{5.164485in}{2.770650in}}%
\pgfpathlineto{\pgfqpoint{5.165152in}{2.787238in}}%
\pgfpathlineto{\pgfqpoint{5.165200in}{2.776476in}}%
\pgfpathlineto{\pgfqpoint{5.165984in}{2.808529in}}%
\pgfpathlineto{\pgfqpoint{5.166262in}{2.787075in}}%
\pgfpathlineto{\pgfqpoint{5.167425in}{2.819218in}}%
\pgfpathlineto{\pgfqpoint{5.166958in}{2.780260in}}%
\pgfpathlineto{\pgfqpoint{5.167493in}{2.812274in}}%
\pgfpathlineto{\pgfqpoint{5.167678in}{2.803078in}}%
\pgfpathlineto{\pgfqpoint{5.167863in}{2.819878in}}%
\pgfpathlineto{\pgfqpoint{5.168316in}{2.840758in}}%
\pgfpathlineto{\pgfqpoint{5.168793in}{2.803301in}}%
\pgfpathlineto{\pgfqpoint{5.168973in}{2.820224in}}%
\pgfpathlineto{\pgfqpoint{5.169007in}{2.813599in}}%
\pgfpathlineto{\pgfqpoint{5.169816in}{2.850507in}}%
\pgfpathlineto{\pgfqpoint{5.169947in}{2.835461in}}%
\pgfpathlineto{\pgfqpoint{5.170551in}{2.875565in}}%
\pgfpathlineto{\pgfqpoint{5.170001in}{2.833018in}}%
\pgfpathlineto{\pgfqpoint{5.171125in}{2.847637in}}%
\pgfpathlineto{\pgfqpoint{5.171437in}{2.841931in}}%
\pgfpathlineto{\pgfqpoint{5.171943in}{2.865944in}}%
\pgfpathlineto{\pgfqpoint{5.173151in}{2.904883in}}%
\pgfpathlineto{\pgfqpoint{5.172386in}{2.862569in}}%
\pgfpathlineto{\pgfqpoint{5.173272in}{2.894585in}}%
\pgfpathlineto{\pgfqpoint{5.173413in}{2.886385in}}%
\pgfpathlineto{\pgfqpoint{5.174124in}{2.922244in}}%
\pgfpathlineto{\pgfqpoint{5.174280in}{2.910805in}}%
\pgfpathlineto{\pgfqpoint{5.174840in}{2.929249in}}%
\pgfpathlineto{\pgfqpoint{5.174533in}{2.890946in}}%
\pgfpathlineto{\pgfqpoint{5.175375in}{2.904982in}}%
\pgfpathlineto{\pgfqpoint{5.175380in}{2.904361in}}%
\pgfpathlineto{\pgfqpoint{5.175892in}{2.937782in}}%
\pgfpathlineto{\pgfqpoint{5.176354in}{2.921643in}}%
\pgfpathlineto{\pgfqpoint{5.177440in}{2.966284in}}%
\pgfpathlineto{\pgfqpoint{5.176383in}{2.919542in}}%
\pgfpathlineto{\pgfqpoint{5.177591in}{2.949940in}}%
\pgfpathlineto{\pgfqpoint{5.178345in}{2.931058in}}%
\pgfpathlineto{\pgfqpoint{5.177795in}{2.957736in}}%
\pgfpathlineto{\pgfqpoint{5.178706in}{2.947570in}}%
\pgfpathlineto{\pgfqpoint{5.179606in}{2.982435in}}%
\pgfpathlineto{\pgfqpoint{5.179830in}{2.952410in}}%
\pgfpathlineto{\pgfqpoint{5.181096in}{2.932647in}}%
\pgfpathlineto{\pgfqpoint{5.180020in}{2.960154in}}%
\pgfpathlineto{\pgfqpoint{5.181252in}{2.946906in}}%
\pgfpathlineto{\pgfqpoint{5.181296in}{2.957978in}}%
\pgfpathlineto{\pgfqpoint{5.181865in}{2.937203in}}%
\pgfpathlineto{\pgfqpoint{5.182362in}{2.947507in}}%
\pgfpathlineto{\pgfqpoint{5.183516in}{2.979680in}}%
\pgfpathlineto{\pgfqpoint{5.182474in}{2.941697in}}%
\pgfpathlineto{\pgfqpoint{5.183764in}{2.958260in}}%
\pgfpathlineto{\pgfqpoint{5.183798in}{2.952096in}}%
\pgfpathlineto{\pgfqpoint{5.184265in}{2.982090in}}%
\pgfpathlineto{\pgfqpoint{5.184830in}{2.973436in}}%
\pgfpathlineto{\pgfqpoint{5.186008in}{3.046273in}}%
\pgfpathlineto{\pgfqpoint{5.186052in}{3.038725in}}%
\pgfpathlineto{\pgfqpoint{5.186076in}{3.034978in}}%
\pgfpathlineto{\pgfqpoint{5.187045in}{3.075568in}}%
\pgfpathlineto{\pgfqpoint{5.187469in}{3.105448in}}%
\pgfpathlineto{\pgfqpoint{5.187863in}{3.063585in}}%
\pgfpathlineto{\pgfqpoint{5.188165in}{3.079624in}}%
\pgfpathlineto{\pgfqpoint{5.189260in}{3.050329in}}%
\pgfpathlineto{\pgfqpoint{5.188209in}{3.084490in}}%
\pgfpathlineto{\pgfqpoint{5.189333in}{3.054336in}}%
\pgfpathlineto{\pgfqpoint{5.192313in}{3.152069in}}%
\pgfpathlineto{\pgfqpoint{5.189368in}{3.049806in}}%
\pgfpathlineto{\pgfqpoint{5.192454in}{3.130163in}}%
\pgfpathlineto{\pgfqpoint{5.193379in}{3.098671in}}%
\pgfpathlineto{\pgfqpoint{5.193608in}{3.112638in}}%
\pgfpathlineto{\pgfqpoint{5.193949in}{3.142861in}}%
\pgfpathlineto{\pgfqpoint{5.194674in}{3.097477in}}%
\pgfpathlineto{\pgfqpoint{5.194684in}{3.100409in}}%
\pgfpathlineto{\pgfqpoint{5.195400in}{3.114363in}}%
\pgfpathlineto{\pgfqpoint{5.195551in}{3.091796in}}%
\pgfpathlineto{\pgfqpoint{5.195706in}{3.096576in}}%
\pgfpathlineto{\pgfqpoint{5.196003in}{3.079106in}}%
\pgfpathlineto{\pgfqpoint{5.196592in}{3.116402in}}%
\pgfpathlineto{\pgfqpoint{5.196816in}{3.096794in}}%
\pgfpathlineto{\pgfqpoint{5.197016in}{3.066466in}}%
\pgfpathlineto{\pgfqpoint{5.197795in}{3.124007in}}%
\pgfpathlineto{\pgfqpoint{5.198642in}{3.136443in}}%
\pgfpathlineto{\pgfqpoint{5.198428in}{3.109414in}}%
\pgfpathlineto{\pgfqpoint{5.198871in}{3.117555in}}%
\pgfpathlineto{\pgfqpoint{5.199650in}{3.090737in}}%
\pgfpathlineto{\pgfqpoint{5.198905in}{3.121066in}}%
\pgfpathlineto{\pgfqpoint{5.200073in}{3.096082in}}%
\pgfpathlineto{\pgfqpoint{5.200594in}{3.104544in}}%
\pgfpathlineto{\pgfqpoint{5.201144in}{3.081246in}}%
\pgfpathlineto{\pgfqpoint{5.201154in}{3.085793in}}%
\pgfpathlineto{\pgfqpoint{5.201923in}{3.096061in}}%
\pgfpathlineto{\pgfqpoint{5.201563in}{3.072399in}}%
\pgfpathlineto{\pgfqpoint{5.202225in}{3.081883in}}%
\pgfpathlineto{\pgfqpoint{5.202717in}{3.052934in}}%
\pgfpathlineto{\pgfqpoint{5.202235in}{3.082396in}}%
\pgfpathlineto{\pgfqpoint{5.203394in}{3.069760in}}%
\pgfpathlineto{\pgfqpoint{5.204075in}{3.123228in}}%
\pgfpathlineto{\pgfqpoint{5.203457in}{3.069538in}}%
\pgfpathlineto{\pgfqpoint{5.204557in}{3.096641in}}%
\pgfpathlineto{\pgfqpoint{5.204845in}{3.089629in}}%
\pgfpathlineto{\pgfqpoint{5.205599in}{3.125936in}}%
\pgfpathlineto{\pgfqpoint{5.205604in}{3.124853in}}%
\pgfpathlineto{\pgfqpoint{5.205682in}{3.134844in}}%
\pgfpathlineto{\pgfqpoint{5.206607in}{3.091299in}}%
\pgfpathlineto{\pgfqpoint{5.206675in}{3.097738in}}%
\pgfpathlineto{\pgfqpoint{5.208150in}{3.145675in}}%
\pgfpathlineto{\pgfqpoint{5.206992in}{3.090172in}}%
\pgfpathlineto{\pgfqpoint{5.208184in}{3.143720in}}%
\pgfpathlineto{\pgfqpoint{5.208287in}{3.123203in}}%
\pgfpathlineto{\pgfqpoint{5.208661in}{3.156838in}}%
\pgfpathlineto{\pgfqpoint{5.209246in}{3.156464in}}%
\pgfpathlineto{\pgfqpoint{5.210302in}{3.168324in}}%
\pgfpathlineto{\pgfqpoint{5.209854in}{3.139695in}}%
\pgfpathlineto{\pgfqpoint{5.210365in}{3.160866in}}%
\pgfpathlineto{\pgfqpoint{5.211101in}{3.143947in}}%
\pgfpathlineto{\pgfqpoint{5.211412in}{3.198242in}}%
\pgfpathlineto{\pgfqpoint{5.211422in}{3.197816in}}%
\pgfpathlineto{\pgfqpoint{5.213145in}{3.290155in}}%
\pgfpathlineto{\pgfqpoint{5.211719in}{3.186054in}}%
\pgfpathlineto{\pgfqpoint{5.213754in}{3.256716in}}%
\pgfpathlineto{\pgfqpoint{5.214372in}{3.245210in}}%
\pgfpathlineto{\pgfqpoint{5.214659in}{3.267648in}}%
\pgfpathlineto{\pgfqpoint{5.214859in}{3.259374in}}%
\pgfpathlineto{\pgfqpoint{5.215112in}{3.277929in}}%
\pgfpathlineto{\pgfqpoint{5.215886in}{3.243820in}}%
\pgfpathlineto{\pgfqpoint{5.216991in}{3.227912in}}%
\pgfpathlineto{\pgfqpoint{5.216149in}{3.262999in}}%
\pgfpathlineto{\pgfqpoint{5.217001in}{3.229908in}}%
\pgfpathlineto{\pgfqpoint{5.218179in}{3.268754in}}%
\pgfpathlineto{\pgfqpoint{5.217108in}{3.219830in}}%
\pgfpathlineto{\pgfqpoint{5.218184in}{3.266889in}}%
\pgfpathlineto{\pgfqpoint{5.219411in}{3.239288in}}%
\pgfpathlineto{\pgfqpoint{5.218257in}{3.281778in}}%
\pgfpathlineto{\pgfqpoint{5.219489in}{3.250056in}}%
\pgfpathlineto{\pgfqpoint{5.220356in}{3.269069in}}%
\pgfpathlineto{\pgfqpoint{5.219679in}{3.229370in}}%
\pgfpathlineto{\pgfqpoint{5.220589in}{3.249631in}}%
\pgfpathlineto{\pgfqpoint{5.220818in}{3.243050in}}%
\pgfpathlineto{\pgfqpoint{5.221276in}{3.293391in}}%
\pgfpathlineto{\pgfqpoint{5.221612in}{3.267622in}}%
\pgfpathlineto{\pgfqpoint{5.221831in}{3.256896in}}%
\pgfpathlineto{\pgfqpoint{5.222809in}{3.310685in}}%
\pgfpathlineto{\pgfqpoint{5.223408in}{3.349570in}}%
\pgfpathlineto{\pgfqpoint{5.222843in}{3.307816in}}%
\pgfpathlineto{\pgfqpoint{5.224056in}{3.345240in}}%
\pgfpathlineto{\pgfqpoint{5.224431in}{3.333625in}}%
\pgfpathlineto{\pgfqpoint{5.224864in}{3.390969in}}%
\pgfpathlineto{\pgfqpoint{5.225141in}{3.355628in}}%
\pgfpathlineto{\pgfqpoint{5.226120in}{3.389649in}}%
\pgfpathlineto{\pgfqpoint{5.225521in}{3.327803in}}%
\pgfpathlineto{\pgfqpoint{5.226295in}{3.380881in}}%
\pgfpathlineto{\pgfqpoint{5.227756in}{3.331045in}}%
\pgfpathlineto{\pgfqpoint{5.226461in}{3.399213in}}%
\pgfpathlineto{\pgfqpoint{5.227790in}{3.336519in}}%
\pgfpathlineto{\pgfqpoint{5.228671in}{3.358170in}}%
\pgfpathlineto{\pgfqpoint{5.228277in}{3.332474in}}%
\pgfpathlineto{\pgfqpoint{5.228934in}{3.340206in}}%
\pgfpathlineto{\pgfqpoint{5.229883in}{3.332875in}}%
\pgfpathlineto{\pgfqpoint{5.229143in}{3.365338in}}%
\pgfpathlineto{\pgfqpoint{5.229947in}{3.344496in}}%
\pgfpathlineto{\pgfqpoint{5.230341in}{3.362171in}}%
\pgfpathlineto{\pgfqpoint{5.230901in}{3.327455in}}%
\pgfpathlineto{\pgfqpoint{5.231008in}{3.328678in}}%
\pgfpathlineto{\pgfqpoint{5.231324in}{3.309681in}}%
\pgfpathlineto{\pgfqpoint{5.231125in}{3.344373in}}%
\pgfpathlineto{\pgfqpoint{5.232089in}{3.333203in}}%
\pgfpathlineto{\pgfqpoint{5.232629in}{3.339869in}}%
\pgfpathlineto{\pgfqpoint{5.232429in}{3.313730in}}%
\pgfpathlineto{\pgfqpoint{5.233135in}{3.320495in}}%
\pgfpathlineto{\pgfqpoint{5.234430in}{3.272465in}}%
\pgfpathlineto{\pgfqpoint{5.234435in}{3.274693in}}%
\pgfpathlineto{\pgfqpoint{5.235575in}{3.330817in}}%
\pgfpathlineto{\pgfqpoint{5.234450in}{3.273890in}}%
\pgfpathlineto{\pgfqpoint{5.235662in}{3.323969in}}%
\pgfpathlineto{\pgfqpoint{5.236675in}{3.331629in}}%
\pgfpathlineto{\pgfqpoint{5.236018in}{3.310712in}}%
\pgfpathlineto{\pgfqpoint{5.236758in}{3.322438in}}%
\pgfpathlineto{\pgfqpoint{5.236850in}{3.315489in}}%
\pgfpathlineto{\pgfqpoint{5.237624in}{3.394372in}}%
\pgfpathlineto{\pgfqpoint{5.237765in}{3.380099in}}%
\pgfpathlineto{\pgfqpoint{5.238520in}{3.373628in}}%
\pgfpathlineto{\pgfqpoint{5.237941in}{3.396087in}}%
\pgfpathlineto{\pgfqpoint{5.238856in}{3.381334in}}%
\pgfpathlineto{\pgfqpoint{5.238992in}{3.388875in}}%
\pgfpathlineto{\pgfqpoint{5.239484in}{3.356080in}}%
\pgfpathlineto{\pgfqpoint{5.239888in}{3.367797in}}%
\pgfpathlineto{\pgfqpoint{5.240404in}{3.351044in}}%
\pgfpathlineto{\pgfqpoint{5.240190in}{3.372744in}}%
\pgfpathlineto{\pgfqpoint{5.241003in}{3.362710in}}%
\pgfpathlineto{\pgfqpoint{5.241699in}{3.392824in}}%
\pgfpathlineto{\pgfqpoint{5.241061in}{3.352859in}}%
\pgfpathlineto{\pgfqpoint{5.242113in}{3.362576in}}%
\pgfpathlineto{\pgfqpoint{5.242498in}{3.328768in}}%
\pgfpathlineto{\pgfqpoint{5.242278in}{3.368600in}}%
\pgfpathlineto{\pgfqpoint{5.243330in}{3.347181in}}%
\pgfpathlineto{\pgfqpoint{5.243335in}{3.349605in}}%
\pgfpathlineto{\pgfqpoint{5.244338in}{3.309554in}}%
\pgfpathlineto{\pgfqpoint{5.244372in}{3.315056in}}%
\pgfpathlineto{\pgfqpoint{5.244795in}{3.292253in}}%
\pgfpathlineto{\pgfqpoint{5.245375in}{3.359835in}}%
\pgfpathlineto{\pgfqpoint{5.245394in}{3.358373in}}%
\pgfpathlineto{\pgfqpoint{5.245842in}{3.378153in}}%
\pgfpathlineto{\pgfqpoint{5.246441in}{3.337838in}}%
\pgfpathlineto{\pgfqpoint{5.246475in}{3.339627in}}%
\pgfpathlineto{\pgfqpoint{5.246845in}{3.354179in}}%
\pgfpathlineto{\pgfqpoint{5.247463in}{3.325716in}}%
\pgfpathlineto{\pgfqpoint{5.247551in}{3.327858in}}%
\pgfpathlineto{\pgfqpoint{5.248530in}{3.307777in}}%
\pgfpathlineto{\pgfqpoint{5.248043in}{3.334346in}}%
\pgfpathlineto{\pgfqpoint{5.248719in}{3.311283in}}%
\pgfpathlineto{\pgfqpoint{5.248744in}{3.319050in}}%
\pgfpathlineto{\pgfqpoint{5.249494in}{3.272748in}}%
\pgfpathlineto{\pgfqpoint{5.249747in}{3.279907in}}%
\pgfpathlineto{\pgfqpoint{5.250131in}{3.304851in}}%
\pgfpathlineto{\pgfqpoint{5.250419in}{3.267469in}}%
\pgfpathlineto{\pgfqpoint{5.250857in}{3.280648in}}%
\pgfpathlineto{\pgfqpoint{5.252103in}{3.243386in}}%
\pgfpathlineto{\pgfqpoint{5.251655in}{3.285374in}}%
\pgfpathlineto{\pgfqpoint{5.252113in}{3.246918in}}%
\pgfpathlineto{\pgfqpoint{5.252162in}{3.253744in}}%
\pgfpathlineto{\pgfqpoint{5.253028in}{3.192632in}}%
\pgfpathlineto{\pgfqpoint{5.253043in}{3.195712in}}%
\pgfpathlineto{\pgfqpoint{5.253213in}{3.181342in}}%
\pgfpathlineto{\pgfqpoint{5.253783in}{3.208243in}}%
\pgfpathlineto{\pgfqpoint{5.254109in}{3.199736in}}%
\pgfpathlineto{\pgfqpoint{5.254469in}{3.219717in}}%
\pgfpathlineto{\pgfqpoint{5.254917in}{3.194750in}}%
\pgfpathlineto{\pgfqpoint{5.255983in}{3.169387in}}%
\pgfpathlineto{\pgfqpoint{5.255302in}{3.210753in}}%
\pgfpathlineto{\pgfqpoint{5.256090in}{3.172418in}}%
\pgfpathlineto{\pgfqpoint{5.256344in}{3.187694in}}%
\pgfpathlineto{\pgfqpoint{5.257161in}{3.151588in}}%
\pgfpathlineto{\pgfqpoint{5.257171in}{3.157600in}}%
\pgfpathlineto{\pgfqpoint{5.257312in}{3.145525in}}%
\pgfpathlineto{\pgfqpoint{5.258018in}{3.170410in}}%
\pgfpathlineto{\pgfqpoint{5.258310in}{3.151351in}}%
\pgfpathlineto{\pgfqpoint{5.259313in}{3.189685in}}%
\pgfpathlineto{\pgfqpoint{5.259455in}{3.178321in}}%
\pgfpathlineto{\pgfqpoint{5.259523in}{3.187081in}}%
\pgfpathlineto{\pgfqpoint{5.259756in}{3.162061in}}%
\pgfpathlineto{\pgfqpoint{5.260521in}{3.163100in}}%
\pgfpathlineto{\pgfqpoint{5.261300in}{3.184106in}}%
\pgfpathlineto{\pgfqpoint{5.260633in}{3.156978in}}%
\pgfpathlineto{\pgfqpoint{5.261626in}{3.165169in}}%
\pgfpathlineto{\pgfqpoint{5.262541in}{3.088743in}}%
\pgfpathlineto{\pgfqpoint{5.261675in}{3.169953in}}%
\pgfpathlineto{\pgfqpoint{5.262838in}{3.107869in}}%
\pgfpathlineto{\pgfqpoint{5.263106in}{3.107658in}}%
\pgfpathlineto{\pgfqpoint{5.264021in}{3.161699in}}%
\pgfpathlineto{\pgfqpoint{5.264824in}{3.133340in}}%
\pgfpathlineto{\pgfqpoint{5.264187in}{3.169865in}}%
\pgfpathlineto{\pgfqpoint{5.265151in}{3.144416in}}%
\pgfpathlineto{\pgfqpoint{5.265915in}{3.186648in}}%
\pgfpathlineto{\pgfqpoint{5.265389in}{3.130206in}}%
\pgfpathlineto{\pgfqpoint{5.266305in}{3.181579in}}%
\pgfpathlineto{\pgfqpoint{5.266324in}{3.179341in}}%
\pgfpathlineto{\pgfqpoint{5.266348in}{3.183285in}}%
\pgfpathlineto{\pgfqpoint{5.266392in}{3.182535in}}%
\pgfpathlineto{\pgfqpoint{5.266402in}{3.186711in}}%
\pgfpathlineto{\pgfqpoint{5.267147in}{3.156664in}}%
\pgfpathlineto{\pgfqpoint{5.267458in}{3.157392in}}%
\pgfpathlineto{\pgfqpoint{5.268125in}{3.134331in}}%
\pgfpathlineto{\pgfqpoint{5.267546in}{3.159412in}}%
\pgfpathlineto{\pgfqpoint{5.268627in}{3.140685in}}%
\pgfpathlineto{\pgfqpoint{5.268773in}{3.151616in}}%
\pgfpathlineto{\pgfqpoint{5.269255in}{3.103001in}}%
\pgfpathlineto{\pgfqpoint{5.269493in}{3.105585in}}%
\pgfpathlineto{\pgfqpoint{5.270365in}{3.083009in}}%
\pgfpathlineto{\pgfqpoint{5.270024in}{3.119642in}}%
\pgfpathlineto{\pgfqpoint{5.270613in}{3.095397in}}%
\pgfpathlineto{\pgfqpoint{5.270676in}{3.106026in}}%
\pgfpathlineto{\pgfqpoint{5.271665in}{3.060138in}}%
\pgfpathlineto{\pgfqpoint{5.272931in}{3.031352in}}%
\pgfpathlineto{\pgfqpoint{5.272239in}{3.066133in}}%
\pgfpathlineto{\pgfqpoint{5.272945in}{3.035223in}}%
\pgfpathlineto{\pgfqpoint{5.273140in}{3.056554in}}%
\pgfpathlineto{\pgfqpoint{5.273583in}{3.028965in}}%
\pgfpathlineto{\pgfqpoint{5.274089in}{3.046513in}}%
\pgfpathlineto{\pgfqpoint{5.274109in}{3.043930in}}%
\pgfpathlineto{\pgfqpoint{5.274742in}{3.077243in}}%
\pgfpathlineto{\pgfqpoint{5.275194in}{3.048098in}}%
\pgfpathlineto{\pgfqpoint{5.275472in}{3.034993in}}%
\pgfpathlineto{\pgfqpoint{5.276163in}{3.069608in}}%
\pgfpathlineto{\pgfqpoint{5.277088in}{3.108104in}}%
\pgfpathlineto{\pgfqpoint{5.277332in}{3.106562in}}%
\pgfpathlineto{\pgfqpoint{5.278339in}{3.085162in}}%
\pgfpathlineto{\pgfqpoint{5.277901in}{3.114305in}}%
\pgfpathlineto{\pgfqpoint{5.278451in}{3.098763in}}%
\pgfpathlineto{\pgfqpoint{5.278821in}{3.133824in}}%
\pgfpathlineto{\pgfqpoint{5.279406in}{3.094281in}}%
\pgfpathlineto{\pgfqpoint{5.279566in}{3.102364in}}%
\pgfpathlineto{\pgfqpoint{5.280696in}{3.142054in}}%
\pgfpathlineto{\pgfqpoint{5.280087in}{3.094355in}}%
\pgfpathlineto{\pgfqpoint{5.280774in}{3.129824in}}%
\pgfpathlineto{\pgfqpoint{5.281966in}{3.084927in}}%
\pgfpathlineto{\pgfqpoint{5.281105in}{3.142753in}}%
\pgfpathlineto{\pgfqpoint{5.282156in}{3.100128in}}%
\pgfpathlineto{\pgfqpoint{5.282390in}{3.102244in}}%
\pgfpathlineto{\pgfqpoint{5.282551in}{3.085832in}}%
\pgfpathlineto{\pgfqpoint{5.283539in}{3.059355in}}%
\pgfpathlineto{\pgfqpoint{5.283169in}{3.096044in}}%
\pgfpathlineto{\pgfqpoint{5.283685in}{3.062770in}}%
\pgfpathlineto{\pgfqpoint{5.284313in}{3.042898in}}%
\pgfpathlineto{\pgfqpoint{5.284819in}{3.055486in}}%
\pgfpathlineto{\pgfqpoint{5.285910in}{3.041183in}}%
\pgfpathlineto{\pgfqpoint{5.285277in}{3.076214in}}%
\pgfpathlineto{\pgfqpoint{5.286041in}{3.050059in}}%
\pgfpathlineto{\pgfqpoint{5.286645in}{3.091969in}}%
\pgfpathlineto{\pgfqpoint{5.286090in}{3.048900in}}%
\pgfpathlineto{\pgfqpoint{5.287181in}{3.068445in}}%
\pgfpathlineto{\pgfqpoint{5.287186in}{3.068505in}}%
\pgfpathlineto{\pgfqpoint{5.287288in}{3.056789in}}%
\pgfpathlineto{\pgfqpoint{5.287293in}{3.057115in}}%
\pgfpathlineto{\pgfqpoint{5.287531in}{3.032788in}}%
\pgfpathlineto{\pgfqpoint{5.287857in}{3.063382in}}%
\pgfpathlineto{\pgfqpoint{5.288476in}{3.040213in}}%
\pgfpathlineto{\pgfqpoint{5.288656in}{3.056228in}}%
\pgfpathlineto{\pgfqpoint{5.289109in}{3.017112in}}%
\pgfpathlineto{\pgfqpoint{5.289371in}{3.031878in}}%
\pgfpathlineto{\pgfqpoint{5.290535in}{2.984041in}}%
\pgfpathlineto{\pgfqpoint{5.290564in}{2.986760in}}%
\pgfpathlineto{\pgfqpoint{5.291952in}{3.012410in}}%
\pgfpathlineto{\pgfqpoint{5.291037in}{2.974717in}}%
\pgfpathlineto{\pgfqpoint{5.292035in}{3.008270in}}%
\pgfpathlineto{\pgfqpoint{5.292069in}{3.000855in}}%
\pgfpathlineto{\pgfqpoint{5.292482in}{3.039055in}}%
\pgfpathlineto{\pgfqpoint{5.293115in}{3.028386in}}%
\pgfpathlineto{\pgfqpoint{5.293227in}{3.036113in}}%
\pgfpathlineto{\pgfqpoint{5.294040in}{2.980254in}}%
\pgfpathlineto{\pgfqpoint{5.294099in}{2.982755in}}%
\pgfpathlineto{\pgfqpoint{5.294756in}{3.015716in}}%
\pgfpathlineto{\pgfqpoint{5.295024in}{2.967292in}}%
\pgfpathlineto{\pgfqpoint{5.295150in}{2.979398in}}%
\pgfpathlineto{\pgfqpoint{5.296002in}{2.975142in}}%
\pgfpathlineto{\pgfqpoint{5.295593in}{3.001242in}}%
\pgfpathlineto{\pgfqpoint{5.296221in}{2.997476in}}%
\pgfpathlineto{\pgfqpoint{5.296961in}{3.014847in}}%
\pgfpathlineto{\pgfqpoint{5.296655in}{2.975650in}}%
\pgfpathlineto{\pgfqpoint{5.297346in}{3.005325in}}%
\pgfpathlineto{\pgfqpoint{5.297960in}{3.041478in}}%
\pgfpathlineto{\pgfqpoint{5.298510in}{3.015927in}}%
\pgfpathlineto{\pgfqpoint{5.298943in}{3.000097in}}%
\pgfpathlineto{\pgfqpoint{5.299255in}{3.020541in}}%
\pgfpathlineto{\pgfqpoint{5.299634in}{3.006098in}}%
\pgfpathlineto{\pgfqpoint{5.299980in}{3.039494in}}%
\pgfpathlineto{\pgfqpoint{5.299698in}{3.002777in}}%
\pgfpathlineto{\pgfqpoint{5.300754in}{3.010892in}}%
\pgfpathlineto{\pgfqpoint{5.300808in}{3.004305in}}%
\pgfpathlineto{\pgfqpoint{5.301810in}{3.047146in}}%
\pgfpathlineto{\pgfqpoint{5.301815in}{3.046321in}}%
\pgfpathlineto{\pgfqpoint{5.302083in}{3.067726in}}%
\pgfpathlineto{\pgfqpoint{5.302740in}{3.028180in}}%
\pgfpathlineto{\pgfqpoint{5.302891in}{3.042230in}}%
\pgfpathlineto{\pgfqpoint{5.302930in}{3.032567in}}%
\pgfpathlineto{\pgfqpoint{5.303811in}{3.070557in}}%
\pgfpathlineto{\pgfqpoint{5.303958in}{3.058262in}}%
\pgfpathlineto{\pgfqpoint{5.305082in}{3.102343in}}%
\pgfpathlineto{\pgfqpoint{5.304001in}{3.046985in}}%
\pgfpathlineto{\pgfqpoint{5.305228in}{3.099685in}}%
\pgfpathlineto{\pgfqpoint{5.305403in}{3.081139in}}%
\pgfpathlineto{\pgfqpoint{5.306066in}{3.105543in}}%
\pgfpathlineto{\pgfqpoint{5.306348in}{3.094489in}}%
\pgfpathlineto{\pgfqpoint{5.307502in}{3.127201in}}%
\pgfpathlineto{\pgfqpoint{5.306406in}{3.090294in}}%
\pgfpathlineto{\pgfqpoint{5.307550in}{3.123575in}}%
\pgfpathlineto{\pgfqpoint{5.308398in}{3.080081in}}%
\pgfpathlineto{\pgfqpoint{5.308709in}{3.096266in}}%
\pgfpathlineto{\pgfqpoint{5.309508in}{3.127826in}}%
\pgfpathlineto{\pgfqpoint{5.308855in}{3.088303in}}%
\pgfpathlineto{\pgfqpoint{5.309926in}{3.117674in}}%
\pgfpathlineto{\pgfqpoint{5.311056in}{3.077612in}}%
\pgfpathlineto{\pgfqpoint{5.310141in}{3.128718in}}%
\pgfpathlineto{\pgfqpoint{5.311075in}{3.083520in}}%
\pgfpathlineto{\pgfqpoint{5.311309in}{3.090889in}}%
\pgfpathlineto{\pgfqpoint{5.312122in}{3.063284in}}%
\pgfpathlineto{\pgfqpoint{5.312151in}{3.069448in}}%
\pgfpathlineto{\pgfqpoint{5.312185in}{3.073376in}}%
\pgfpathlineto{\pgfqpoint{5.312881in}{3.016431in}}%
\pgfpathlineto{\pgfqpoint{5.313081in}{3.023254in}}%
\pgfpathlineto{\pgfqpoint{5.314250in}{3.006679in}}%
\pgfpathlineto{\pgfqpoint{5.314006in}{3.035607in}}%
\pgfpathlineto{\pgfqpoint{5.314284in}{3.012487in}}%
\pgfpathlineto{\pgfqpoint{5.315871in}{2.976158in}}%
\pgfpathlineto{\pgfqpoint{5.314527in}{3.022251in}}%
\pgfpathlineto{\pgfqpoint{5.315934in}{2.988313in}}%
\pgfpathlineto{\pgfqpoint{5.317005in}{3.014920in}}%
\pgfpathlineto{\pgfqpoint{5.316791in}{2.978306in}}%
\pgfpathlineto{\pgfqpoint{5.317093in}{3.002572in}}%
\pgfpathlineto{\pgfqpoint{5.317745in}{2.989711in}}%
\pgfpathlineto{\pgfqpoint{5.317390in}{3.012572in}}%
\pgfpathlineto{\pgfqpoint{5.318193in}{3.004711in}}%
\pgfpathlineto{\pgfqpoint{5.318436in}{3.037416in}}%
\pgfpathlineto{\pgfqpoint{5.319357in}{3.035914in}}%
\pgfpathlineto{\pgfqpoint{5.319955in}{3.011246in}}%
\pgfpathlineto{\pgfqpoint{5.319678in}{3.049370in}}%
\pgfpathlineto{\pgfqpoint{5.320413in}{3.039058in}}%
\pgfpathlineto{\pgfqpoint{5.320973in}{3.044377in}}%
\pgfpathlineto{\pgfqpoint{5.320618in}{3.021076in}}%
\pgfpathlineto{\pgfqpoint{5.321460in}{3.027522in}}%
\pgfpathlineto{\pgfqpoint{5.322146in}{3.003431in}}%
\pgfpathlineto{\pgfqpoint{5.321801in}{3.055309in}}%
\pgfpathlineto{\pgfqpoint{5.322628in}{3.008049in}}%
\pgfpathlineto{\pgfqpoint{5.322808in}{3.021669in}}%
\pgfpathlineto{\pgfqpoint{5.323704in}{2.993520in}}%
\pgfpathlineto{\pgfqpoint{5.323709in}{2.993704in}}%
\pgfpathlineto{\pgfqpoint{5.323719in}{2.996280in}}%
\pgfpathlineto{\pgfqpoint{5.323991in}{2.971454in}}%
\pgfpathlineto{\pgfqpoint{5.324488in}{2.974856in}}%
\pgfpathlineto{\pgfqpoint{5.325501in}{2.948623in}}%
\pgfpathlineto{\pgfqpoint{5.324761in}{2.992777in}}%
\pgfpathlineto{\pgfqpoint{5.325661in}{2.953336in}}%
\pgfpathlineto{\pgfqpoint{5.326839in}{3.016943in}}%
\pgfpathlineto{\pgfqpoint{5.325841in}{2.951956in}}%
\pgfpathlineto{\pgfqpoint{5.326951in}{3.002650in}}%
\pgfpathlineto{\pgfqpoint{5.327541in}{2.966800in}}%
\pgfpathlineto{\pgfqpoint{5.327886in}{3.011314in}}%
\pgfpathlineto{\pgfqpoint{5.328057in}{3.005857in}}%
\pgfpathlineto{\pgfqpoint{5.328071in}{3.008500in}}%
\pgfpathlineto{\pgfqpoint{5.328076in}{3.009845in}}%
\pgfpathlineto{\pgfqpoint{5.328777in}{2.942467in}}%
\pgfpathlineto{\pgfqpoint{5.328952in}{2.947802in}}%
\pgfpathlineto{\pgfqpoint{5.329069in}{2.938417in}}%
\pgfpathlineto{\pgfqpoint{5.329829in}{2.975688in}}%
\pgfpathlineto{\pgfqpoint{5.330019in}{2.966140in}}%
\pgfpathlineto{\pgfqpoint{5.330140in}{2.974274in}}%
\pgfpathlineto{\pgfqpoint{5.330744in}{2.929026in}}%
\pgfpathlineto{\pgfqpoint{5.331085in}{2.943221in}}%
\pgfpathlineto{\pgfqpoint{5.331109in}{2.937317in}}%
\pgfpathlineto{\pgfqpoint{5.331168in}{2.949040in}}%
\pgfpathlineto{\pgfqpoint{5.331172in}{2.948904in}}%
\pgfpathlineto{\pgfqpoint{5.331314in}{2.968291in}}%
\pgfpathlineto{\pgfqpoint{5.331873in}{2.930286in}}%
\pgfpathlineto{\pgfqpoint{5.332292in}{2.952786in}}%
\pgfpathlineto{\pgfqpoint{5.333105in}{2.931478in}}%
\pgfpathlineto{\pgfqpoint{5.332648in}{2.969835in}}%
\pgfpathlineto{\pgfqpoint{5.333461in}{2.937837in}}%
\pgfpathlineto{\pgfqpoint{5.333709in}{2.957037in}}%
\pgfpathlineto{\pgfqpoint{5.334089in}{2.909110in}}%
\pgfpathlineto{\pgfqpoint{5.334532in}{2.933267in}}%
\pgfpathlineto{\pgfqpoint{5.335690in}{2.878122in}}%
\pgfpathlineto{\pgfqpoint{5.334610in}{2.933759in}}%
\pgfpathlineto{\pgfqpoint{5.335695in}{2.880923in}}%
\pgfpathlineto{\pgfqpoint{5.336353in}{2.855673in}}%
\pgfpathlineto{\pgfqpoint{5.336752in}{2.891912in}}%
\pgfpathlineto{\pgfqpoint{5.338052in}{2.928181in}}%
\pgfpathlineto{\pgfqpoint{5.338091in}{2.920048in}}%
\pgfpathlineto{\pgfqpoint{5.339240in}{2.891325in}}%
\pgfpathlineto{\pgfqpoint{5.338120in}{2.921489in}}%
\pgfpathlineto{\pgfqpoint{5.339254in}{2.892690in}}%
\pgfpathlineto{\pgfqpoint{5.339580in}{2.868254in}}%
\pgfpathlineto{\pgfqpoint{5.339936in}{2.904413in}}%
\pgfpathlineto{\pgfqpoint{5.340087in}{2.917598in}}%
\pgfpathlineto{\pgfqpoint{5.340165in}{2.900905in}}%
\pgfpathlineto{\pgfqpoint{5.341051in}{2.907885in}}%
\pgfpathlineto{\pgfqpoint{5.341742in}{2.898042in}}%
\pgfpathlineto{\pgfqpoint{5.341601in}{2.926628in}}%
\pgfpathlineto{\pgfqpoint{5.342122in}{2.920835in}}%
\pgfpathlineto{\pgfqpoint{5.342755in}{2.953568in}}%
\pgfpathlineto{\pgfqpoint{5.343251in}{2.935140in}}%
\pgfpathlineto{\pgfqpoint{5.344176in}{2.908126in}}%
\pgfpathlineto{\pgfqpoint{5.343714in}{2.966632in}}%
\pgfpathlineto{\pgfqpoint{5.344381in}{2.930931in}}%
\pgfpathlineto{\pgfqpoint{5.344405in}{2.934453in}}%
\pgfpathlineto{\pgfqpoint{5.344868in}{2.897824in}}%
\pgfpathlineto{\pgfqpoint{5.345028in}{2.899534in}}%
\pgfpathlineto{\pgfqpoint{5.345281in}{2.884519in}}%
\pgfpathlineto{\pgfqpoint{5.345895in}{2.918216in}}%
\pgfpathlineto{\pgfqpoint{5.346060in}{2.916179in}}%
\pgfpathlineto{\pgfqpoint{5.346236in}{2.911944in}}%
\pgfpathlineto{\pgfqpoint{5.347292in}{2.952649in}}%
\pgfpathlineto{\pgfqpoint{5.348056in}{2.931935in}}%
\pgfpathlineto{\pgfqpoint{5.348329in}{2.953535in}}%
\pgfpathlineto{\pgfqpoint{5.348407in}{2.947736in}}%
\pgfpathlineto{\pgfqpoint{5.348860in}{2.969687in}}%
\pgfpathlineto{\pgfqpoint{5.348519in}{2.935777in}}%
\pgfpathlineto{\pgfqpoint{5.349527in}{2.952556in}}%
\pgfpathlineto{\pgfqpoint{5.350466in}{2.910716in}}%
\pgfpathlineto{\pgfqpoint{5.350763in}{2.938188in}}%
\pgfpathlineto{\pgfqpoint{5.351002in}{2.946540in}}%
\pgfpathlineto{\pgfqpoint{5.351265in}{2.924277in}}%
\pgfpathlineto{\pgfqpoint{5.351508in}{2.927082in}}%
\pgfpathlineto{\pgfqpoint{5.352647in}{2.900129in}}%
\pgfpathlineto{\pgfqpoint{5.352652in}{2.903593in}}%
\pgfpathlineto{\pgfqpoint{5.353076in}{2.880767in}}%
\pgfpathlineto{\pgfqpoint{5.353602in}{2.909769in}}%
\pgfpathlineto{\pgfqpoint{5.353796in}{2.884945in}}%
\pgfpathlineto{\pgfqpoint{5.354093in}{2.897814in}}%
\pgfpathlineto{\pgfqpoint{5.354434in}{2.864148in}}%
\pgfpathlineto{\pgfqpoint{5.354897in}{2.880771in}}%
\pgfpathlineto{\pgfqpoint{5.354926in}{2.873290in}}%
\pgfpathlineto{\pgfqpoint{5.355719in}{2.921892in}}%
\pgfpathlineto{\pgfqpoint{5.355929in}{2.912513in}}%
\pgfpathlineto{\pgfqpoint{5.356761in}{2.948627in}}%
\pgfpathlineto{\pgfqpoint{5.356265in}{2.905173in}}%
\pgfpathlineto{\pgfqpoint{5.357102in}{2.940231in}}%
\pgfpathlineto{\pgfqpoint{5.358606in}{2.865068in}}%
\pgfpathlineto{\pgfqpoint{5.358631in}{2.873648in}}%
\pgfpathlineto{\pgfqpoint{5.359609in}{2.893954in}}%
\pgfpathlineto{\pgfqpoint{5.359760in}{2.885379in}}%
\pgfpathlineto{\pgfqpoint{5.360461in}{2.865719in}}%
\pgfpathlineto{\pgfqpoint{5.359799in}{2.890903in}}%
\pgfpathlineto{\pgfqpoint{5.360875in}{2.883629in}}%
\pgfpathlineto{\pgfqpoint{5.361932in}{2.900260in}}%
\pgfpathlineto{\pgfqpoint{5.361240in}{2.864361in}}%
\pgfpathlineto{\pgfqpoint{5.362010in}{2.886742in}}%
\pgfpathlineto{\pgfqpoint{5.362097in}{2.876513in}}%
\pgfpathlineto{\pgfqpoint{5.362998in}{2.943893in}}%
\pgfpathlineto{\pgfqpoint{5.363061in}{2.933735in}}%
\pgfpathlineto{\pgfqpoint{5.363275in}{2.924941in}}%
\pgfpathlineto{\pgfqpoint{5.363694in}{2.948500in}}%
\pgfpathlineto{\pgfqpoint{5.363879in}{2.943457in}}%
\pgfpathlineto{\pgfqpoint{5.364385in}{2.989575in}}%
\pgfpathlineto{\pgfqpoint{5.365125in}{2.976271in}}%
\pgfpathlineto{\pgfqpoint{5.365159in}{2.968968in}}%
\pgfpathlineto{\pgfqpoint{5.365573in}{2.988808in}}%
\pgfpathlineto{\pgfqpoint{5.366216in}{2.984447in}}%
\pgfpathlineto{\pgfqpoint{5.366824in}{3.031293in}}%
\pgfpathlineto{\pgfqpoint{5.367418in}{3.007249in}}%
\pgfpathlineto{\pgfqpoint{5.368305in}{2.965310in}}%
\pgfpathlineto{\pgfqpoint{5.367472in}{3.017090in}}%
\pgfpathlineto{\pgfqpoint{5.368597in}{2.979453in}}%
\pgfpathlineto{\pgfqpoint{5.369483in}{3.012846in}}%
\pgfpathlineto{\pgfqpoint{5.369746in}{2.998024in}}%
\pgfpathlineto{\pgfqpoint{5.369965in}{3.028766in}}%
\pgfpathlineto{\pgfqpoint{5.370641in}{2.995219in}}%
\pgfpathlineto{\pgfqpoint{5.370943in}{3.012522in}}%
\pgfpathlineto{\pgfqpoint{5.372014in}{2.965182in}}%
\pgfpathlineto{\pgfqpoint{5.370963in}{3.013441in}}%
\pgfpathlineto{\pgfqpoint{5.372404in}{2.986171in}}%
\pgfpathlineto{\pgfqpoint{5.372409in}{2.986047in}}%
\pgfpathlineto{\pgfqpoint{5.372579in}{3.010017in}}%
\pgfpathlineto{\pgfqpoint{5.373168in}{3.025264in}}%
\pgfpathlineto{\pgfqpoint{5.372813in}{2.999503in}}%
\pgfpathlineto{\pgfqpoint{5.373475in}{3.008319in}}%
\pgfpathlineto{\pgfqpoint{5.374137in}{2.979290in}}%
\pgfpathlineto{\pgfqpoint{5.374595in}{3.001302in}}%
\pgfpathlineto{\pgfqpoint{5.374599in}{3.005040in}}%
\pgfpathlineto{\pgfqpoint{5.375529in}{2.964116in}}%
\pgfpathlineto{\pgfqpoint{5.375671in}{2.976323in}}%
\pgfpathlineto{\pgfqpoint{5.375690in}{2.980155in}}%
\pgfpathlineto{\pgfqpoint{5.376600in}{2.960062in}}%
\pgfpathlineto{\pgfqpoint{5.377613in}{2.926460in}}%
\pgfpathlineto{\pgfqpoint{5.377058in}{2.968461in}}%
\pgfpathlineto{\pgfqpoint{5.377779in}{2.931245in}}%
\pgfpathlineto{\pgfqpoint{5.378893in}{2.961682in}}%
\pgfpathlineto{\pgfqpoint{5.378163in}{2.926714in}}%
\pgfpathlineto{\pgfqpoint{5.378971in}{2.959892in}}%
\pgfpathlineto{\pgfqpoint{5.379293in}{2.934225in}}%
\pgfpathlineto{\pgfqpoint{5.380062in}{2.976027in}}%
\pgfpathlineto{\pgfqpoint{5.380067in}{2.974866in}}%
\pgfpathlineto{\pgfqpoint{5.381167in}{3.010787in}}%
\pgfpathlineto{\pgfqpoint{5.380266in}{2.956896in}}%
\pgfpathlineto{\pgfqpoint{5.381230in}{3.001350in}}%
\pgfpathlineto{\pgfqpoint{5.382034in}{2.961983in}}%
\pgfpathlineto{\pgfqpoint{5.382360in}{2.977952in}}%
\pgfpathlineto{\pgfqpoint{5.382399in}{2.972023in}}%
\pgfpathlineto{\pgfqpoint{5.382852in}{2.995446in}}%
\pgfpathlineto{\pgfqpoint{5.384161in}{3.024175in}}%
\pgfpathlineto{\pgfqpoint{5.383226in}{2.988071in}}%
\pgfpathlineto{\pgfqpoint{5.384273in}{3.015513in}}%
\pgfpathlineto{\pgfqpoint{5.385135in}{2.988664in}}%
\pgfpathlineto{\pgfqpoint{5.384283in}{3.016172in}}%
\pgfpathlineto{\pgfqpoint{5.385515in}{3.000568in}}%
\pgfpathlineto{\pgfqpoint{5.386221in}{3.017430in}}%
\pgfpathlineto{\pgfqpoint{5.385797in}{2.976949in}}%
\pgfpathlineto{\pgfqpoint{5.386630in}{3.003999in}}%
\pgfpathlineto{\pgfqpoint{5.386664in}{3.006445in}}%
\pgfpathlineto{\pgfqpoint{5.387335in}{2.980128in}}%
\pgfpathlineto{\pgfqpoint{5.387545in}{2.987394in}}%
\pgfpathlineto{\pgfqpoint{5.387550in}{2.986573in}}%
\pgfpathlineto{\pgfqpoint{5.388012in}{3.020194in}}%
\pgfpathlineto{\pgfqpoint{5.388572in}{2.999149in}}%
\pgfpathlineto{\pgfqpoint{5.389575in}{3.013431in}}%
\pgfpathlineto{\pgfqpoint{5.388859in}{2.970670in}}%
\pgfpathlineto{\pgfqpoint{5.389692in}{3.004044in}}%
\pgfpathlineto{\pgfqpoint{5.389702in}{3.000101in}}%
\pgfpathlineto{\pgfqpoint{5.390661in}{3.047566in}}%
\pgfpathlineto{\pgfqpoint{5.390753in}{3.037044in}}%
\pgfpathlineto{\pgfqpoint{5.390797in}{3.033708in}}%
\pgfpathlineto{\pgfqpoint{5.391391in}{3.066353in}}%
\pgfpathlineto{\pgfqpoint{5.391785in}{3.056997in}}%
\pgfpathlineto{\pgfqpoint{5.392360in}{3.072081in}}%
\pgfpathlineto{\pgfqpoint{5.391912in}{3.034427in}}%
\pgfpathlineto{\pgfqpoint{5.392934in}{3.066911in}}%
\pgfpathlineto{\pgfqpoint{5.393192in}{3.062332in}}%
\pgfpathlineto{\pgfqpoint{5.393640in}{3.080974in}}%
\pgfpathlineto{\pgfqpoint{5.394025in}{3.070807in}}%
\pgfpathlineto{\pgfqpoint{5.395149in}{3.083783in}}%
\pgfpathlineto{\pgfqpoint{5.394521in}{3.059175in}}%
\pgfpathlineto{\pgfqpoint{5.395159in}{3.080511in}}%
\pgfpathlineto{\pgfqpoint{5.395193in}{3.077218in}}%
\pgfpathlineto{\pgfqpoint{5.395846in}{3.118790in}}%
\pgfpathlineto{\pgfqpoint{5.396211in}{3.104716in}}%
\pgfpathlineto{\pgfqpoint{5.396537in}{3.123926in}}%
\pgfpathlineto{\pgfqpoint{5.396970in}{3.082581in}}%
\pgfpathlineto{\pgfqpoint{5.397029in}{3.085619in}}%
\pgfpathlineto{\pgfqpoint{5.397589in}{3.074569in}}%
\pgfpathlineto{\pgfqpoint{5.398090in}{3.097534in}}%
\pgfpathlineto{\pgfqpoint{5.399215in}{3.135149in}}%
\pgfpathlineto{\pgfqpoint{5.398655in}{3.085944in}}%
\pgfpathlineto{\pgfqpoint{5.399234in}{3.133243in}}%
\pgfpathlineto{\pgfqpoint{5.399263in}{3.128783in}}%
\pgfpathlineto{\pgfqpoint{5.399731in}{3.169160in}}%
\pgfpathlineto{\pgfqpoint{5.400286in}{3.155327in}}%
\pgfpathlineto{\pgfqpoint{5.400295in}{3.156458in}}%
\pgfpathlineto{\pgfqpoint{5.400996in}{3.126703in}}%
\pgfpathlineto{\pgfqpoint{5.401152in}{3.139513in}}%
\pgfpathlineto{\pgfqpoint{5.401162in}{3.135875in}}%
\pgfpathlineto{\pgfqpoint{5.402068in}{3.178576in}}%
\pgfpathlineto{\pgfqpoint{5.402194in}{3.169582in}}%
\pgfpathlineto{\pgfqpoint{5.402257in}{3.160344in}}%
\pgfpathlineto{\pgfqpoint{5.402559in}{3.199023in}}%
\pgfpathlineto{\pgfqpoint{5.403202in}{3.184666in}}%
\pgfpathlineto{\pgfqpoint{5.403217in}{3.186562in}}%
\pgfpathlineto{\pgfqpoint{5.404088in}{3.157260in}}%
\pgfpathlineto{\pgfqpoint{5.404254in}{3.170695in}}%
\pgfpathlineto{\pgfqpoint{5.404619in}{3.171198in}}%
\pgfpathlineto{\pgfqpoint{5.405432in}{3.140856in}}%
\pgfpathlineto{\pgfqpoint{5.405500in}{3.148305in}}%
\pgfpathlineto{\pgfqpoint{5.406328in}{3.124909in}}%
\pgfpathlineto{\pgfqpoint{5.406532in}{3.136547in}}%
\pgfpathlineto{\pgfqpoint{5.407983in}{3.092290in}}%
\pgfpathlineto{\pgfqpoint{5.407993in}{3.095616in}}%
\pgfpathlineto{\pgfqpoint{5.408148in}{3.100703in}}%
\pgfpathlineto{\pgfqpoint{5.408908in}{3.060498in}}%
\pgfpathlineto{\pgfqpoint{5.409049in}{3.075483in}}%
\pgfpathlineto{\pgfqpoint{5.410393in}{3.103922in}}%
\pgfpathlineto{\pgfqpoint{5.409161in}{3.064731in}}%
\pgfpathlineto{\pgfqpoint{5.410461in}{3.099053in}}%
\pgfpathlineto{\pgfqpoint{5.410514in}{3.088681in}}%
\pgfpathlineto{\pgfqpoint{5.411508in}{3.121299in}}%
\pgfpathlineto{\pgfqpoint{5.411542in}{3.118997in}}%
\pgfpathlineto{\pgfqpoint{5.413168in}{3.177637in}}%
\pgfpathlineto{\pgfqpoint{5.413192in}{3.172647in}}%
\pgfpathlineto{\pgfqpoint{5.413538in}{3.160079in}}%
\pgfpathlineto{\pgfqpoint{5.414039in}{3.227517in}}%
\pgfpathlineto{\pgfqpoint{5.414185in}{3.217270in}}%
\pgfpathlineto{\pgfqpoint{5.414219in}{3.223394in}}%
\pgfpathlineto{\pgfqpoint{5.414862in}{3.181664in}}%
\pgfpathlineto{\pgfqpoint{5.415266in}{3.213365in}}%
\pgfpathlineto{\pgfqpoint{5.415310in}{3.205052in}}%
\pgfpathlineto{\pgfqpoint{5.415962in}{3.237008in}}%
\pgfpathlineto{\pgfqpoint{5.416323in}{3.231023in}}%
\pgfpathlineto{\pgfqpoint{5.416449in}{3.243703in}}%
\pgfpathlineto{\pgfqpoint{5.417223in}{3.211612in}}%
\pgfpathlineto{\pgfqpoint{5.417369in}{3.216034in}}%
\pgfpathlineto{\pgfqpoint{5.417374in}{3.216565in}}%
\pgfpathlineto{\pgfqpoint{5.417622in}{3.181842in}}%
\pgfpathlineto{\pgfqpoint{5.418265in}{3.194829in}}%
\pgfpathlineto{\pgfqpoint{5.419190in}{3.165223in}}%
\pgfpathlineto{\pgfqpoint{5.418869in}{3.205656in}}%
\pgfpathlineto{\pgfqpoint{5.419429in}{3.170047in}}%
\pgfpathlineto{\pgfqpoint{5.419740in}{3.152023in}}%
\pgfpathlineto{\pgfqpoint{5.420266in}{3.178829in}}%
\pgfpathlineto{\pgfqpoint{5.420524in}{3.172596in}}%
\pgfpathlineto{\pgfqpoint{5.420544in}{3.174513in}}%
\pgfpathlineto{\pgfqpoint{5.421376in}{3.149590in}}%
\pgfpathlineto{\pgfqpoint{5.421469in}{3.156771in}}%
\pgfpathlineto{\pgfqpoint{5.421644in}{3.133059in}}%
\pgfpathlineto{\pgfqpoint{5.422301in}{3.179098in}}%
\pgfpathlineto{\pgfqpoint{5.422583in}{3.154345in}}%
\pgfpathlineto{\pgfqpoint{5.422764in}{3.164508in}}%
\pgfpathlineto{\pgfqpoint{5.423737in}{3.139862in}}%
\pgfpathlineto{\pgfqpoint{5.424492in}{3.156111in}}%
\pgfpathlineto{\pgfqpoint{5.424088in}{3.125721in}}%
\pgfpathlineto{\pgfqpoint{5.424740in}{3.138925in}}%
\pgfpathlineto{\pgfqpoint{5.425831in}{3.102912in}}%
\pgfpathlineto{\pgfqpoint{5.425057in}{3.147511in}}%
\pgfpathlineto{\pgfqpoint{5.425923in}{3.110116in}}%
\pgfpathlineto{\pgfqpoint{5.426639in}{3.143831in}}%
\pgfpathlineto{\pgfqpoint{5.426055in}{3.102477in}}%
\pgfpathlineto{\pgfqpoint{5.427140in}{3.126060in}}%
\pgfpathlineto{\pgfqpoint{5.427204in}{3.118056in}}%
\pgfpathlineto{\pgfqpoint{5.427802in}{3.151772in}}%
\pgfpathlineto{\pgfqpoint{5.428197in}{3.143105in}}%
\pgfpathlineto{\pgfqpoint{5.428640in}{3.175554in}}%
\pgfpathlineto{\pgfqpoint{5.428216in}{3.140998in}}%
\pgfpathlineto{\pgfqpoint{5.429390in}{3.158460in}}%
\pgfpathlineto{\pgfqpoint{5.429584in}{3.134051in}}%
\pgfpathlineto{\pgfqpoint{5.430217in}{3.184809in}}%
\pgfpathlineto{\pgfqpoint{5.430397in}{3.183080in}}%
\pgfpathlineto{\pgfqpoint{5.431201in}{3.227227in}}%
\pgfpathlineto{\pgfqpoint{5.430417in}{3.179519in}}%
\pgfpathlineto{\pgfqpoint{5.431541in}{3.216306in}}%
\pgfpathlineto{\pgfqpoint{5.432690in}{3.178525in}}%
\pgfpathlineto{\pgfqpoint{5.432705in}{3.181615in}}%
\pgfpathlineto{\pgfqpoint{5.433104in}{3.204628in}}%
\pgfpathlineto{\pgfqpoint{5.432739in}{3.179265in}}%
\pgfpathlineto{\pgfqpoint{5.433864in}{3.196193in}}%
\pgfpathlineto{\pgfqpoint{5.434487in}{3.187174in}}%
\pgfpathlineto{\pgfqpoint{5.434755in}{3.226011in}}%
\pgfpathlineto{\pgfqpoint{5.435952in}{3.260371in}}%
\pgfpathlineto{\pgfqpoint{5.434813in}{3.224342in}}%
\pgfpathlineto{\pgfqpoint{5.436006in}{3.253070in}}%
\pgfpathlineto{\pgfqpoint{5.436756in}{3.197367in}}%
\pgfpathlineto{\pgfqpoint{5.436094in}{3.265860in}}%
\pgfpathlineto{\pgfqpoint{5.437223in}{3.201122in}}%
\pgfpathlineto{\pgfqpoint{5.438391in}{3.273733in}}%
\pgfpathlineto{\pgfqpoint{5.438421in}{3.269831in}}%
\pgfpathlineto{\pgfqpoint{5.439326in}{3.240592in}}%
\pgfpathlineto{\pgfqpoint{5.438800in}{3.270819in}}%
\pgfpathlineto{\pgfqpoint{5.439599in}{3.246510in}}%
\pgfpathlineto{\pgfqpoint{5.440285in}{3.222161in}}%
\pgfpathlineto{\pgfqpoint{5.439813in}{3.264962in}}%
\pgfpathlineto{\pgfqpoint{5.440738in}{3.241918in}}%
\pgfpathlineto{\pgfqpoint{5.441001in}{3.264560in}}%
\pgfpathlineto{\pgfqpoint{5.441751in}{3.226014in}}%
\pgfpathlineto{\pgfqpoint{5.441824in}{3.230456in}}%
\pgfpathlineto{\pgfqpoint{5.442466in}{3.192700in}}%
\pgfpathlineto{\pgfqpoint{5.441892in}{3.231792in}}%
\pgfpathlineto{\pgfqpoint{5.443036in}{3.212566in}}%
\pgfpathlineto{\pgfqpoint{5.443757in}{3.246050in}}%
\pgfpathlineto{\pgfqpoint{5.443148in}{3.209068in}}%
\pgfpathlineto{\pgfqpoint{5.444224in}{3.239612in}}%
\pgfpathlineto{\pgfqpoint{5.445387in}{3.294273in}}%
\pgfpathlineto{\pgfqpoint{5.445782in}{3.271064in}}%
\pgfpathlineto{\pgfqpoint{5.446882in}{3.232204in}}%
\pgfpathlineto{\pgfqpoint{5.446926in}{3.234297in}}%
\pgfpathlineto{\pgfqpoint{5.447111in}{3.250251in}}%
\pgfpathlineto{\pgfqpoint{5.447413in}{3.217160in}}%
\pgfpathlineto{\pgfqpoint{5.448036in}{3.238047in}}%
\pgfpathlineto{\pgfqpoint{5.448070in}{3.224554in}}%
\pgfpathlineto{\pgfqpoint{5.449058in}{3.261445in}}%
\pgfpathlineto{\pgfqpoint{5.449097in}{3.258989in}}%
\pgfpathlineto{\pgfqpoint{5.450071in}{3.274575in}}%
\pgfpathlineto{\pgfqpoint{5.449604in}{3.242579in}}%
\pgfpathlineto{\pgfqpoint{5.450212in}{3.262650in}}%
\pgfpathlineto{\pgfqpoint{5.450266in}{3.256240in}}%
\pgfpathlineto{\pgfqpoint{5.451137in}{3.311688in}}%
\pgfpathlineto{\pgfqpoint{5.451176in}{3.311412in}}%
\pgfpathlineto{\pgfqpoint{5.451712in}{3.349029in}}%
\pgfpathlineto{\pgfqpoint{5.451225in}{3.309101in}}%
\pgfpathlineto{\pgfqpoint{5.452330in}{3.330435in}}%
\pgfpathlineto{\pgfqpoint{5.453396in}{3.359060in}}%
\pgfpathlineto{\pgfqpoint{5.452866in}{3.316530in}}%
\pgfpathlineto{\pgfqpoint{5.453912in}{3.350927in}}%
\pgfpathlineto{\pgfqpoint{5.454735in}{3.299667in}}%
\pgfpathlineto{\pgfqpoint{5.455061in}{3.322396in}}%
\pgfpathlineto{\pgfqpoint{5.455105in}{3.330066in}}%
\pgfpathlineto{\pgfqpoint{5.455801in}{3.292515in}}%
\pgfpathlineto{\pgfqpoint{5.456093in}{3.298072in}}%
\pgfpathlineto{\pgfqpoint{5.457247in}{3.276415in}}%
\pgfpathlineto{\pgfqpoint{5.456546in}{3.322875in}}%
\pgfpathlineto{\pgfqpoint{5.457267in}{3.277253in}}%
\pgfpathlineto{\pgfqpoint{5.457948in}{3.309632in}}%
\pgfpathlineto{\pgfqpoint{5.458401in}{3.302796in}}%
\pgfpathlineto{\pgfqpoint{5.459126in}{3.352325in}}%
\pgfpathlineto{\pgfqpoint{5.459633in}{3.319507in}}%
\pgfpathlineto{\pgfqpoint{5.460850in}{3.260121in}}%
\pgfpathlineto{\pgfqpoint{5.459891in}{3.321879in}}%
\pgfpathlineto{\pgfqpoint{5.460855in}{3.261820in}}%
\pgfpathlineto{\pgfqpoint{5.460991in}{3.264228in}}%
\pgfpathlineto{\pgfqpoint{5.461113in}{3.245791in}}%
\pgfpathlineto{\pgfqpoint{5.461322in}{3.227942in}}%
\pgfpathlineto{\pgfqpoint{5.462077in}{3.247816in}}%
\pgfpathlineto{\pgfqpoint{5.462237in}{3.239380in}}%
\pgfpathlineto{\pgfqpoint{5.462242in}{3.240550in}}%
\pgfpathlineto{\pgfqpoint{5.463109in}{3.187927in}}%
\pgfpathlineto{\pgfqpoint{5.463153in}{3.194541in}}%
\pgfpathlineto{\pgfqpoint{5.464277in}{3.179381in}}%
\pgfpathlineto{\pgfqpoint{5.463440in}{3.207692in}}%
\pgfpathlineto{\pgfqpoint{5.464282in}{3.180924in}}%
\pgfpathlineto{\pgfqpoint{5.464603in}{3.195596in}}%
\pgfpathlineto{\pgfqpoint{5.465124in}{3.169440in}}%
\pgfpathlineto{\pgfqpoint{5.465397in}{3.182762in}}%
\pgfpathlineto{\pgfqpoint{5.465490in}{3.177451in}}%
\pgfpathlineto{\pgfqpoint{5.466025in}{3.204454in}}%
\pgfpathlineto{\pgfqpoint{5.466502in}{3.184548in}}%
\pgfpathlineto{\pgfqpoint{5.467228in}{3.209701in}}%
\pgfpathlineto{\pgfqpoint{5.466648in}{3.182487in}}%
\pgfpathlineto{\pgfqpoint{5.467622in}{3.190507in}}%
\pgfpathlineto{\pgfqpoint{5.468464in}{3.160424in}}%
\pgfpathlineto{\pgfqpoint{5.467734in}{3.194389in}}%
\pgfpathlineto{\pgfqpoint{5.468781in}{3.162204in}}%
\pgfpathlineto{\pgfqpoint{5.468946in}{3.173872in}}%
\pgfpathlineto{\pgfqpoint{5.469579in}{3.153481in}}%
\pgfpathlineto{\pgfqpoint{5.469891in}{3.164034in}}%
\pgfpathlineto{\pgfqpoint{5.470962in}{3.137411in}}%
\pgfpathlineto{\pgfqpoint{5.470081in}{3.171151in}}%
\pgfpathlineto{\pgfqpoint{5.471059in}{3.138379in}}%
\pgfpathlineto{\pgfqpoint{5.472413in}{3.208396in}}%
\pgfpathlineto{\pgfqpoint{5.471215in}{3.134955in}}%
\pgfpathlineto{\pgfqpoint{5.472515in}{3.207397in}}%
\pgfpathlineto{\pgfqpoint{5.473036in}{3.181010in}}%
\pgfpathlineto{\pgfqpoint{5.473469in}{3.218935in}}%
\pgfpathlineto{\pgfqpoint{5.473635in}{3.201222in}}%
\pgfpathlineto{\pgfqpoint{5.473669in}{3.196986in}}%
\pgfpathlineto{\pgfqpoint{5.473795in}{3.205018in}}%
\pgfpathlineto{\pgfqpoint{5.473805in}{3.203215in}}%
\pgfpathlineto{\pgfqpoint{5.474521in}{3.237112in}}%
\pgfpathlineto{\pgfqpoint{5.473941in}{3.199972in}}%
\pgfpathlineto{\pgfqpoint{5.474934in}{3.219628in}}%
\pgfpathlineto{\pgfqpoint{5.475280in}{3.215387in}}%
\pgfpathlineto{\pgfqpoint{5.476132in}{3.252160in}}%
\pgfpathlineto{\pgfqpoint{5.476702in}{3.226930in}}%
\pgfpathlineto{\pgfqpoint{5.477208in}{3.259472in}}%
\pgfpathlineto{\pgfqpoint{5.477237in}{3.255309in}}%
\pgfpathlineto{\pgfqpoint{5.478084in}{3.237448in}}%
\pgfpathlineto{\pgfqpoint{5.477632in}{3.264435in}}%
\pgfpathlineto{\pgfqpoint{5.478415in}{3.248710in}}%
\pgfpathlineto{\pgfqpoint{5.479443in}{3.275522in}}%
\pgfpathlineto{\pgfqpoint{5.478654in}{3.233784in}}%
\pgfpathlineto{\pgfqpoint{5.479560in}{3.269958in}}%
\pgfpathlineto{\pgfqpoint{5.480017in}{3.240085in}}%
\pgfpathlineto{\pgfqpoint{5.479827in}{3.272603in}}%
\pgfpathlineto{\pgfqpoint{5.480699in}{3.252956in}}%
\pgfpathlineto{\pgfqpoint{5.480903in}{3.233484in}}%
\pgfpathlineto{\pgfqpoint{5.481692in}{3.255616in}}%
\pgfpathlineto{\pgfqpoint{5.481711in}{3.261203in}}%
\pgfpathlineto{\pgfqpoint{5.482520in}{3.194422in}}%
\pgfpathlineto{\pgfqpoint{5.482612in}{3.204918in}}%
\pgfpathlineto{\pgfqpoint{5.483922in}{3.132044in}}%
\pgfpathlineto{\pgfqpoint{5.483931in}{3.134920in}}%
\pgfpathlineto{\pgfqpoint{5.484998in}{3.154378in}}%
\pgfpathlineto{\pgfqpoint{5.484384in}{3.125449in}}%
\pgfpathlineto{\pgfqpoint{5.485085in}{3.153980in}}%
\pgfpathlineto{\pgfqpoint{5.486653in}{3.105101in}}%
\pgfpathlineto{\pgfqpoint{5.485665in}{3.169912in}}%
\pgfpathlineto{\pgfqpoint{5.486750in}{3.117434in}}%
\pgfpathlineto{\pgfqpoint{5.487125in}{3.132842in}}%
\pgfpathlineto{\pgfqpoint{5.486848in}{3.102139in}}%
\pgfpathlineto{\pgfqpoint{5.487870in}{3.121290in}}%
\pgfpathlineto{\pgfqpoint{5.488844in}{3.070516in}}%
\pgfpathlineto{\pgfqpoint{5.487958in}{3.128318in}}%
\pgfpathlineto{\pgfqpoint{5.489063in}{3.079172in}}%
\pgfpathlineto{\pgfqpoint{5.489988in}{3.133458in}}%
\pgfpathlineto{\pgfqpoint{5.489355in}{3.073799in}}%
\pgfpathlineto{\pgfqpoint{5.490455in}{3.113765in}}%
\pgfpathlineto{\pgfqpoint{5.491687in}{3.071924in}}%
\pgfpathlineto{\pgfqpoint{5.491750in}{3.081293in}}%
\pgfpathlineto{\pgfqpoint{5.491921in}{3.100423in}}%
\pgfpathlineto{\pgfqpoint{5.492359in}{3.076718in}}%
\pgfpathlineto{\pgfqpoint{5.492870in}{3.089082in}}%
\pgfpathlineto{\pgfqpoint{5.494063in}{3.065305in}}%
\pgfpathlineto{\pgfqpoint{5.493167in}{3.098830in}}%
\pgfpathlineto{\pgfqpoint{5.494077in}{3.068050in}}%
\pgfpathlineto{\pgfqpoint{5.494141in}{3.067599in}}%
\pgfpathlineto{\pgfqpoint{5.494506in}{3.096802in}}%
\pgfpathlineto{\pgfqpoint{5.494671in}{3.092814in}}%
\pgfpathlineto{\pgfqpoint{5.495411in}{3.119924in}}%
\pgfpathlineto{\pgfqpoint{5.494866in}{3.088883in}}%
\pgfpathlineto{\pgfqpoint{5.495942in}{3.102368in}}%
\pgfpathlineto{\pgfqpoint{5.496293in}{3.094970in}}%
\pgfpathlineto{\pgfqpoint{5.496794in}{3.137007in}}%
\pgfpathlineto{\pgfqpoint{5.496979in}{3.128657in}}%
\pgfpathlineto{\pgfqpoint{5.496984in}{3.130483in}}%
\pgfpathlineto{\pgfqpoint{5.497524in}{3.082258in}}%
\pgfpathlineto{\pgfqpoint{5.498031in}{3.107111in}}%
\pgfpathlineto{\pgfqpoint{5.498883in}{3.092584in}}%
\pgfpathlineto{\pgfqpoint{5.498493in}{3.126650in}}%
\pgfpathlineto{\pgfqpoint{5.499150in}{3.098038in}}%
\pgfpathlineto{\pgfqpoint{5.500100in}{3.112192in}}%
\pgfpathlineto{\pgfqpoint{5.499525in}{3.075835in}}%
\pgfpathlineto{\pgfqpoint{5.500265in}{3.098986in}}%
\pgfpathlineto{\pgfqpoint{5.501117in}{3.065003in}}%
\pgfpathlineto{\pgfqpoint{5.500290in}{3.102343in}}%
\pgfpathlineto{\pgfqpoint{5.501473in}{3.087040in}}%
\pgfpathlineto{\pgfqpoint{5.502578in}{3.124693in}}%
\pgfpathlineto{\pgfqpoint{5.502052in}{3.084734in}}%
\pgfpathlineto{\pgfqpoint{5.502738in}{3.115088in}}%
\pgfpathlineto{\pgfqpoint{5.502743in}{3.112951in}}%
\pgfpathlineto{\pgfqpoint{5.503449in}{3.167049in}}%
\pgfpathlineto{\pgfqpoint{5.503800in}{3.134821in}}%
\pgfpathlineto{\pgfqpoint{5.504924in}{3.143548in}}%
\pgfpathlineto{\pgfqpoint{5.504603in}{3.107722in}}%
\pgfpathlineto{\pgfqpoint{5.504929in}{3.142015in}}%
\pgfpathlineto{\pgfqpoint{5.504973in}{3.137289in}}%
\pgfpathlineto{\pgfqpoint{5.505017in}{3.150464in}}%
\pgfpathlineto{\pgfqpoint{5.506239in}{3.219191in}}%
\pgfpathlineto{\pgfqpoint{5.506254in}{3.214932in}}%
\pgfpathlineto{\pgfqpoint{5.506872in}{3.186366in}}%
\pgfpathlineto{\pgfqpoint{5.507291in}{3.225109in}}%
\pgfpathlineto{\pgfqpoint{5.507334in}{3.220631in}}%
\pgfpathlineto{\pgfqpoint{5.508396in}{3.234007in}}%
\pgfpathlineto{\pgfqpoint{5.508157in}{3.217126in}}%
\pgfpathlineto{\pgfqpoint{5.508439in}{3.222593in}}%
\pgfpathlineto{\pgfqpoint{5.508522in}{3.209799in}}%
\pgfpathlineto{\pgfqpoint{5.509330in}{3.254802in}}%
\pgfpathlineto{\pgfqpoint{5.509530in}{3.237957in}}%
\pgfpathlineto{\pgfqpoint{5.510479in}{3.224230in}}%
\pgfpathlineto{\pgfqpoint{5.510221in}{3.250140in}}%
\pgfpathlineto{\pgfqpoint{5.510616in}{3.246234in}}%
\pgfpathlineto{\pgfqpoint{5.510961in}{3.266650in}}%
\pgfpathlineto{\pgfqpoint{5.511458in}{3.219364in}}%
\pgfpathlineto{\pgfqpoint{5.511701in}{3.229520in}}%
\pgfpathlineto{\pgfqpoint{5.512690in}{3.275985in}}%
\pgfpathlineto{\pgfqpoint{5.511838in}{3.227929in}}%
\pgfpathlineto{\pgfqpoint{5.512992in}{3.257428in}}%
\pgfpathlineto{\pgfqpoint{5.513001in}{3.258065in}}%
\pgfpathlineto{\pgfqpoint{5.513454in}{3.218660in}}%
\pgfpathlineto{\pgfqpoint{5.513941in}{3.243543in}}%
\pgfpathlineto{\pgfqpoint{5.515095in}{3.216965in}}%
\pgfpathlineto{\pgfqpoint{5.514379in}{3.255808in}}%
\pgfpathlineto{\pgfqpoint{5.515134in}{3.222047in}}%
\pgfpathlineto{\pgfqpoint{5.515903in}{3.291170in}}%
\pgfpathlineto{\pgfqpoint{5.516312in}{3.280361in}}%
\pgfpathlineto{\pgfqpoint{5.516950in}{3.264270in}}%
\pgfpathlineto{\pgfqpoint{5.516521in}{3.297931in}}%
\pgfpathlineto{\pgfqpoint{5.517432in}{3.274929in}}%
\pgfpathlineto{\pgfqpoint{5.517972in}{3.288727in}}%
\pgfpathlineto{\pgfqpoint{5.518147in}{3.266600in}}%
\pgfpathlineto{\pgfqpoint{5.518172in}{3.267090in}}%
\pgfpathlineto{\pgfqpoint{5.519321in}{3.222855in}}%
\pgfpathlineto{\pgfqpoint{5.519330in}{3.225344in}}%
\pgfpathlineto{\pgfqpoint{5.520148in}{3.275021in}}%
\pgfpathlineto{\pgfqpoint{5.520679in}{3.254415in}}%
\pgfpathlineto{\pgfqpoint{5.521711in}{3.240135in}}%
\pgfpathlineto{\pgfqpoint{5.521112in}{3.273756in}}%
\pgfpathlineto{\pgfqpoint{5.521818in}{3.246640in}}%
\pgfpathlineto{\pgfqpoint{5.522217in}{3.273835in}}%
\pgfpathlineto{\pgfqpoint{5.521847in}{3.243936in}}%
\pgfpathlineto{\pgfqpoint{5.522977in}{3.259495in}}%
\pgfpathlineto{\pgfqpoint{5.523761in}{3.228454in}}%
\pgfpathlineto{\pgfqpoint{5.524126in}{3.241103in}}%
\pgfpathlineto{\pgfqpoint{5.524286in}{3.245165in}}%
\pgfpathlineto{\pgfqpoint{5.524793in}{3.210202in}}%
\pgfpathlineto{\pgfqpoint{5.524934in}{3.214080in}}%
\pgfpathlineto{\pgfqpoint{5.525163in}{3.197310in}}%
\pgfpathlineto{\pgfqpoint{5.525450in}{3.229063in}}%
\pgfpathlineto{\pgfqpoint{5.526039in}{3.218906in}}%
\pgfpathlineto{\pgfqpoint{5.527042in}{3.158762in}}%
\pgfpathlineto{\pgfqpoint{5.526171in}{3.223839in}}%
\pgfpathlineto{\pgfqpoint{5.527266in}{3.165251in}}%
\pgfpathlineto{\pgfqpoint{5.527339in}{3.176907in}}%
\pgfpathlineto{\pgfqpoint{5.527763in}{3.136698in}}%
\pgfpathlineto{\pgfqpoint{5.528376in}{3.168245in}}%
\pgfpathlineto{\pgfqpoint{5.528761in}{3.161716in}}%
\pgfpathlineto{\pgfqpoint{5.529004in}{3.191024in}}%
\pgfpathlineto{\pgfqpoint{5.529442in}{3.185660in}}%
\pgfpathlineto{\pgfqpoint{5.529452in}{3.187425in}}%
\pgfpathlineto{\pgfqpoint{5.530060in}{3.152458in}}%
\pgfpathlineto{\pgfqpoint{5.530465in}{3.161589in}}%
\pgfpathlineto{\pgfqpoint{5.531093in}{3.114181in}}%
\pgfpathlineto{\pgfqpoint{5.530484in}{3.163803in}}%
\pgfpathlineto{\pgfqpoint{5.531633in}{3.140669in}}%
\pgfpathlineto{\pgfqpoint{5.531638in}{3.144090in}}%
\pgfpathlineto{\pgfqpoint{5.531876in}{3.115939in}}%
\pgfpathlineto{\pgfqpoint{5.532709in}{3.119008in}}%
\pgfpathlineto{\pgfqpoint{5.533405in}{3.096125in}}%
\pgfpathlineto{\pgfqpoint{5.532758in}{3.124340in}}%
\pgfpathlineto{\pgfqpoint{5.533795in}{3.118375in}}%
\pgfpathlineto{\pgfqpoint{5.534564in}{3.132403in}}%
\pgfpathlineto{\pgfqpoint{5.534218in}{3.091833in}}%
\pgfpathlineto{\pgfqpoint{5.534905in}{3.120073in}}%
\pgfpathlineto{\pgfqpoint{5.535499in}{3.144639in}}%
\pgfpathlineto{\pgfqpoint{5.535786in}{3.116094in}}%
\pgfpathlineto{\pgfqpoint{5.535791in}{3.117174in}}%
\pgfpathlineto{\pgfqpoint{5.536881in}{3.073989in}}%
\pgfpathlineto{\pgfqpoint{5.535947in}{3.120134in}}%
\pgfpathlineto{\pgfqpoint{5.536935in}{3.083451in}}%
\pgfpathlineto{\pgfqpoint{5.537261in}{3.108194in}}%
\pgfpathlineto{\pgfqpoint{5.537967in}{3.081265in}}%
\pgfpathlineto{\pgfqpoint{5.538040in}{3.083669in}}%
\pgfpathlineto{\pgfqpoint{5.538420in}{3.068457in}}%
\pgfpathlineto{\pgfqpoint{5.538799in}{3.106988in}}%
\pgfpathlineto{\pgfqpoint{5.539140in}{3.086363in}}%
\pgfpathlineto{\pgfqpoint{5.539330in}{3.070102in}}%
\pgfpathlineto{\pgfqpoint{5.539968in}{3.115457in}}%
\pgfpathlineto{\pgfqpoint{5.541068in}{3.137158in}}%
\pgfpathlineto{\pgfqpoint{5.540382in}{3.102709in}}%
\pgfpathlineto{\pgfqpoint{5.541180in}{3.132775in}}%
\pgfpathlineto{\pgfqpoint{5.541779in}{3.118791in}}%
\pgfpathlineto{\pgfqpoint{5.542105in}{3.153829in}}%
\pgfpathlineto{\pgfqpoint{5.542271in}{3.143647in}}%
\pgfpathlineto{\pgfqpoint{5.542490in}{3.152975in}}%
\pgfpathlineto{\pgfqpoint{5.542446in}{3.139768in}}%
\pgfpathlineto{\pgfqpoint{5.543278in}{3.140100in}}%
\pgfpathlineto{\pgfqpoint{5.543473in}{3.131149in}}%
\pgfpathlineto{\pgfqpoint{5.543770in}{3.163437in}}%
\pgfpathlineto{\pgfqpoint{5.544350in}{3.155656in}}%
\pgfpathlineto{\pgfqpoint{5.545284in}{3.192075in}}%
\pgfpathlineto{\pgfqpoint{5.544520in}{3.146162in}}%
\pgfpathlineto{\pgfqpoint{5.545557in}{3.170022in}}%
\pgfpathlineto{\pgfqpoint{5.546331in}{3.163501in}}%
\pgfpathlineto{\pgfqpoint{5.546501in}{3.192081in}}%
\pgfpathlineto{\pgfqpoint{5.546560in}{3.191420in}}%
\pgfpathlineto{\pgfqpoint{5.547680in}{3.214630in}}%
\pgfpathlineto{\pgfqpoint{5.547714in}{3.207861in}}%
\pgfpathlineto{\pgfqpoint{5.548595in}{3.247113in}}%
\pgfpathlineto{\pgfqpoint{5.547835in}{3.192840in}}%
\pgfpathlineto{\pgfqpoint{5.548960in}{3.225338in}}%
\pgfpathlineto{\pgfqpoint{5.549588in}{3.211415in}}%
\pgfpathlineto{\pgfqpoint{5.550002in}{3.241610in}}%
\pgfpathlineto{\pgfqpoint{5.550007in}{3.241150in}}%
\pgfpathlineto{\pgfqpoint{5.551307in}{3.212114in}}%
\pgfpathlineto{\pgfqpoint{5.550075in}{3.243189in}}%
\pgfpathlineto{\pgfqpoint{5.551389in}{3.223704in}}%
\pgfpathlineto{\pgfqpoint{5.551492in}{3.236590in}}%
\pgfpathlineto{\pgfqpoint{5.552475in}{3.208130in}}%
\pgfpathlineto{\pgfqpoint{5.552762in}{3.205663in}}%
\pgfpathlineto{\pgfqpoint{5.553308in}{3.244516in}}%
\pgfpathlineto{\pgfqpoint{5.553327in}{3.243075in}}%
\pgfpathlineto{\pgfqpoint{5.553721in}{3.261239in}}%
\pgfpathlineto{\pgfqpoint{5.553931in}{3.220939in}}%
\pgfpathlineto{\pgfqpoint{5.554461in}{3.255257in}}%
\pgfpathlineto{\pgfqpoint{5.554481in}{3.252258in}}%
\pgfpathlineto{\pgfqpoint{5.555367in}{3.292527in}}%
\pgfpathlineto{\pgfqpoint{5.555445in}{3.287514in}}%
\pgfpathlineto{\pgfqpoint{5.556404in}{3.310226in}}%
\pgfpathlineto{\pgfqpoint{5.555864in}{3.279334in}}%
\pgfpathlineto{\pgfqpoint{5.556560in}{3.290933in}}%
\pgfpathlineto{\pgfqpoint{5.557392in}{3.332319in}}%
\pgfpathlineto{\pgfqpoint{5.556633in}{3.284332in}}%
\pgfpathlineto{\pgfqpoint{5.557933in}{3.316434in}}%
\pgfpathlineto{\pgfqpoint{5.558984in}{3.282245in}}%
\pgfpathlineto{\pgfqpoint{5.557991in}{3.318432in}}%
\pgfpathlineto{\pgfqpoint{5.559077in}{3.288413in}}%
\pgfpathlineto{\pgfqpoint{5.559895in}{3.323767in}}%
\pgfpathlineto{\pgfqpoint{5.559140in}{3.285941in}}%
\pgfpathlineto{\pgfqpoint{5.560245in}{3.303138in}}%
\pgfpathlineto{\pgfqpoint{5.560756in}{3.326570in}}%
\pgfpathlineto{\pgfqpoint{5.561263in}{3.289369in}}%
\pgfpathlineto{\pgfqpoint{5.561287in}{3.289757in}}%
\pgfpathlineto{\pgfqpoint{5.561662in}{3.282610in}}%
\pgfpathlineto{\pgfqpoint{5.561477in}{3.300066in}}%
\pgfpathlineto{\pgfqpoint{5.562382in}{3.294923in}}%
\pgfpathlineto{\pgfqpoint{5.565586in}{3.380047in}}%
\pgfpathlineto{\pgfqpoint{5.562606in}{3.287247in}}%
\pgfpathlineto{\pgfqpoint{5.566268in}{3.356925in}}%
\pgfpathlineto{\pgfqpoint{5.567348in}{3.331904in}}%
\pgfpathlineto{\pgfqpoint{5.566521in}{3.380621in}}%
\pgfpathlineto{\pgfqpoint{5.567446in}{3.337568in}}%
\pgfpathlineto{\pgfqpoint{5.567524in}{3.346013in}}%
\pgfpathlineto{\pgfqpoint{5.568127in}{3.291368in}}%
\pgfpathlineto{\pgfqpoint{5.568453in}{3.294495in}}%
\pgfpathlineto{\pgfqpoint{5.569359in}{3.268521in}}%
\pgfpathlineto{\pgfqpoint{5.568668in}{3.301330in}}%
\pgfpathlineto{\pgfqpoint{5.569656in}{3.277842in}}%
\pgfpathlineto{\pgfqpoint{5.570430in}{3.271588in}}%
\pgfpathlineto{\pgfqpoint{5.569968in}{3.309399in}}%
\pgfpathlineto{\pgfqpoint{5.570620in}{3.285536in}}%
\pgfpathlineto{\pgfqpoint{5.571501in}{3.305619in}}%
\pgfpathlineto{\pgfqpoint{5.571701in}{3.278987in}}%
\pgfpathlineto{\pgfqpoint{5.572748in}{3.236945in}}%
\pgfpathlineto{\pgfqpoint{5.572056in}{3.285571in}}%
\pgfpathlineto{\pgfqpoint{5.572874in}{3.240329in}}%
\pgfpathlineto{\pgfqpoint{5.573268in}{3.218485in}}%
\pgfpathlineto{\pgfqpoint{5.573872in}{3.249409in}}%
\pgfpathlineto{\pgfqpoint{5.573945in}{3.242096in}}%
\pgfpathlineto{\pgfqpoint{5.575011in}{3.294420in}}%
\pgfpathlineto{\pgfqpoint{5.575162in}{3.284530in}}%
\pgfpathlineto{\pgfqpoint{5.575377in}{3.296926in}}%
\pgfpathlineto{\pgfqpoint{5.575712in}{3.257969in}}%
\pgfpathlineto{\pgfqpoint{5.576190in}{3.266991in}}%
\pgfpathlineto{\pgfqpoint{5.577231in}{3.229588in}}%
\pgfpathlineto{\pgfqpoint{5.576199in}{3.268733in}}%
\pgfpathlineto{\pgfqpoint{5.577412in}{3.244686in}}%
\pgfpathlineto{\pgfqpoint{5.577884in}{3.273985in}}%
\pgfpathlineto{\pgfqpoint{5.578541in}{3.257328in}}%
\pgfpathlineto{\pgfqpoint{5.579490in}{3.234785in}}%
\pgfpathlineto{\pgfqpoint{5.578561in}{3.261162in}}%
\pgfpathlineto{\pgfqpoint{5.579671in}{3.243712in}}%
\pgfpathlineto{\pgfqpoint{5.579748in}{3.253041in}}%
\pgfpathlineto{\pgfqpoint{5.580075in}{3.232776in}}%
\pgfpathlineto{\pgfqpoint{5.580766in}{3.235507in}}%
\pgfpathlineto{\pgfqpoint{5.580975in}{3.229743in}}%
\pgfpathlineto{\pgfqpoint{5.581516in}{3.254744in}}%
\pgfpathlineto{\pgfqpoint{5.582217in}{3.312587in}}%
\pgfpathlineto{\pgfqpoint{5.582655in}{3.274940in}}%
\pgfpathlineto{\pgfqpoint{5.582665in}{3.274691in}}%
\pgfpathlineto{\pgfqpoint{5.582670in}{3.275076in}}%
\pgfpathlineto{\pgfqpoint{5.583614in}{3.315323in}}%
\pgfpathlineto{\pgfqpoint{5.582918in}{3.272953in}}%
\pgfpathlineto{\pgfqpoint{5.583809in}{3.300503in}}%
\pgfpathlineto{\pgfqpoint{5.584159in}{3.279435in}}%
\pgfpathlineto{\pgfqpoint{5.583979in}{3.309949in}}%
\pgfpathlineto{\pgfqpoint{5.584948in}{3.290778in}}%
\pgfpathlineto{\pgfqpoint{5.585693in}{3.334641in}}%
\pgfpathlineto{\pgfqpoint{5.585040in}{3.286399in}}%
\pgfpathlineto{\pgfqpoint{5.586082in}{3.308363in}}%
\pgfpathlineto{\pgfqpoint{5.586657in}{3.288823in}}%
\pgfpathlineto{\pgfqpoint{5.586292in}{3.315510in}}%
\pgfpathlineto{\pgfqpoint{5.587197in}{3.304851in}}%
\pgfpathlineto{\pgfqpoint{5.587329in}{3.321255in}}%
\pgfpathlineto{\pgfqpoint{5.587757in}{3.291506in}}%
\pgfpathlineto{\pgfqpoint{5.588336in}{3.308767in}}%
\pgfpathlineto{\pgfqpoint{5.589631in}{3.241563in}}%
\pgfpathlineto{\pgfqpoint{5.589651in}{3.243926in}}%
\pgfpathlineto{\pgfqpoint{5.589675in}{3.244999in}}%
\pgfpathlineto{\pgfqpoint{5.590435in}{3.201453in}}%
\pgfpathlineto{\pgfqpoint{5.591194in}{3.180593in}}%
\pgfpathlineto{\pgfqpoint{5.590527in}{3.205987in}}%
\pgfpathlineto{\pgfqpoint{5.591564in}{3.189848in}}%
\pgfpathlineto{\pgfqpoint{5.591598in}{3.185220in}}%
\pgfpathlineto{\pgfqpoint{5.592110in}{3.213382in}}%
\pgfpathlineto{\pgfqpoint{5.592363in}{3.207292in}}%
\pgfpathlineto{\pgfqpoint{5.593555in}{3.258210in}}%
\pgfpathlineto{\pgfqpoint{5.592738in}{3.205943in}}%
\pgfpathlineto{\pgfqpoint{5.593565in}{3.255189in}}%
\pgfpathlineto{\pgfqpoint{5.594305in}{3.238361in}}%
\pgfpathlineto{\pgfqpoint{5.594515in}{3.264720in}}%
\pgfpathlineto{\pgfqpoint{5.594661in}{3.261239in}}%
\pgfpathlineto{\pgfqpoint{5.595191in}{3.234560in}}%
\pgfpathlineto{\pgfqpoint{5.594670in}{3.264442in}}%
\pgfpathlineto{\pgfqpoint{5.595654in}{3.264271in}}%
\pgfpathlineto{\pgfqpoint{5.596579in}{3.310566in}}%
\pgfpathlineto{\pgfqpoint{5.595771in}{3.263856in}}%
\pgfpathlineto{\pgfqpoint{5.596813in}{3.307690in}}%
\pgfpathlineto{\pgfqpoint{5.597290in}{3.301789in}}%
\pgfpathlineto{\pgfqpoint{5.597100in}{3.324356in}}%
\pgfpathlineto{\pgfqpoint{5.597708in}{3.314534in}}%
\pgfpathlineto{\pgfqpoint{5.598604in}{3.294167in}}%
\pgfpathlineto{\pgfqpoint{5.598872in}{3.336883in}}%
\pgfpathlineto{\pgfqpoint{5.599398in}{3.318926in}}%
\pgfpathlineto{\pgfqpoint{5.599923in}{3.358817in}}%
\pgfpathlineto{\pgfqpoint{5.599928in}{3.357070in}}%
\pgfpathlineto{\pgfqpoint{5.600610in}{3.386395in}}%
\pgfpathlineto{\pgfqpoint{5.601029in}{3.353198in}}%
\pgfpathlineto{\pgfqpoint{5.601034in}{3.353711in}}%
\pgfpathlineto{\pgfqpoint{5.601350in}{3.337246in}}%
\pgfpathlineto{\pgfqpoint{5.602075in}{3.362494in}}%
\pgfpathlineto{\pgfqpoint{5.602090in}{3.360297in}}%
\pgfpathlineto{\pgfqpoint{5.602275in}{3.371843in}}%
\pgfpathlineto{\pgfqpoint{5.602893in}{3.318032in}}%
\pgfpathlineto{\pgfqpoint{5.603107in}{3.342472in}}%
\pgfpathlineto{\pgfqpoint{5.604203in}{3.331765in}}%
\pgfpathlineto{\pgfqpoint{5.603784in}{3.367954in}}%
\pgfpathlineto{\pgfqpoint{5.604237in}{3.331810in}}%
\pgfpathlineto{\pgfqpoint{5.604247in}{3.334393in}}%
\pgfpathlineto{\pgfqpoint{5.604305in}{3.325284in}}%
\pgfpathlineto{\pgfqpoint{5.604359in}{3.325643in}}%
\pgfpathlineto{\pgfqpoint{5.605566in}{3.261513in}}%
\pgfpathlineto{\pgfqpoint{5.604407in}{3.325913in}}%
\pgfpathlineto{\pgfqpoint{5.605673in}{3.264530in}}%
\pgfpathlineto{\pgfqpoint{5.607032in}{3.302409in}}%
\pgfpathlineto{\pgfqpoint{5.606194in}{3.255620in}}%
\pgfpathlineto{\pgfqpoint{5.607041in}{3.300080in}}%
\pgfpathlineto{\pgfqpoint{5.607158in}{3.293469in}}%
\pgfpathlineto{\pgfqpoint{5.607328in}{3.318003in}}%
\pgfpathlineto{\pgfqpoint{5.608127in}{3.306497in}}%
\pgfpathlineto{\pgfqpoint{5.609938in}{3.380099in}}%
\pgfpathlineto{\pgfqpoint{5.608653in}{3.294236in}}%
\pgfpathlineto{\pgfqpoint{5.609943in}{3.377386in}}%
\pgfpathlineto{\pgfqpoint{5.611004in}{3.355824in}}%
\pgfpathlineto{\pgfqpoint{5.610469in}{3.386174in}}%
\pgfpathlineto{\pgfqpoint{5.611072in}{3.365901in}}%
\pgfpathlineto{\pgfqpoint{5.612046in}{3.353871in}}%
\pgfpathlineto{\pgfqpoint{5.611438in}{3.387920in}}%
\pgfpathlineto{\pgfqpoint{5.612163in}{3.369153in}}%
\pgfpathlineto{\pgfqpoint{5.612586in}{3.387826in}}%
\pgfpathlineto{\pgfqpoint{5.612747in}{3.366347in}}%
\pgfpathlineto{\pgfqpoint{5.613288in}{3.375593in}}%
\pgfpathlineto{\pgfqpoint{5.613687in}{3.358886in}}%
\pgfpathlineto{\pgfqpoint{5.614310in}{3.411955in}}%
\pgfpathlineto{\pgfqpoint{5.614729in}{3.430218in}}%
\pgfpathlineto{\pgfqpoint{5.615016in}{3.395572in}}%
\pgfpathlineto{\pgfqpoint{5.615449in}{3.421045in}}%
\pgfpathlineto{\pgfqpoint{5.615800in}{3.403301in}}%
\pgfpathlineto{\pgfqpoint{5.616506in}{3.428437in}}%
\pgfpathlineto{\pgfqpoint{5.616554in}{3.423083in}}%
\pgfpathlineto{\pgfqpoint{5.617007in}{3.439280in}}%
\pgfpathlineto{\pgfqpoint{5.617465in}{3.371507in}}%
\pgfpathlineto{\pgfqpoint{5.617543in}{3.377255in}}%
\pgfpathlineto{\pgfqpoint{5.617650in}{3.393313in}}%
\pgfpathlineto{\pgfqpoint{5.617991in}{3.349900in}}%
\pgfpathlineto{\pgfqpoint{5.618565in}{3.359777in}}%
\pgfpathlineto{\pgfqpoint{5.619651in}{3.319196in}}%
\pgfpathlineto{\pgfqpoint{5.619792in}{3.329539in}}%
\pgfpathlineto{\pgfqpoint{5.620318in}{3.349933in}}%
\pgfpathlineto{\pgfqpoint{5.620868in}{3.312375in}}%
\pgfpathlineto{\pgfqpoint{5.620878in}{3.315220in}}%
\pgfpathlineto{\pgfqpoint{5.621063in}{3.324172in}}%
\pgfpathlineto{\pgfqpoint{5.622207in}{3.233671in}}%
\pgfpathlineto{\pgfqpoint{5.622893in}{3.248667in}}%
\pgfpathlineto{\pgfqpoint{5.623107in}{3.223813in}}%
\pgfpathlineto{\pgfqpoint{5.623341in}{3.243176in}}%
\pgfpathlineto{\pgfqpoint{5.623687in}{3.236022in}}%
\pgfpathlineto{\pgfqpoint{5.624232in}{3.278029in}}%
\pgfpathlineto{\pgfqpoint{5.624237in}{3.277474in}}%
\pgfpathlineto{\pgfqpoint{5.624344in}{3.264778in}}%
\pgfpathlineto{\pgfqpoint{5.625585in}{3.325384in}}%
\pgfpathlineto{\pgfqpoint{5.626705in}{3.292939in}}%
\pgfpathlineto{\pgfqpoint{5.626082in}{3.336156in}}%
\pgfpathlineto{\pgfqpoint{5.626773in}{3.295730in}}%
\pgfpathlineto{\pgfqpoint{5.627065in}{3.313578in}}%
\pgfpathlineto{\pgfqpoint{5.627757in}{3.277737in}}%
\pgfpathlineto{\pgfqpoint{5.627844in}{3.283055in}}%
\pgfpathlineto{\pgfqpoint{5.629008in}{3.321982in}}%
\pgfpathlineto{\pgfqpoint{5.628044in}{3.273663in}}%
\pgfpathlineto{\pgfqpoint{5.629203in}{3.303764in}}%
\pgfpathlineto{\pgfqpoint{5.629212in}{3.300007in}}%
\pgfpathlineto{\pgfqpoint{5.630040in}{3.345670in}}%
\pgfpathlineto{\pgfqpoint{5.630264in}{3.334578in}}%
\pgfpathlineto{\pgfqpoint{5.630770in}{3.353201in}}%
\pgfpathlineto{\pgfqpoint{5.630332in}{3.328123in}}%
\pgfpathlineto{\pgfqpoint{5.631369in}{3.333695in}}%
\pgfpathlineto{\pgfqpoint{5.632319in}{3.317179in}}%
\pgfpathlineto{\pgfqpoint{5.631778in}{3.353602in}}%
\pgfpathlineto{\pgfqpoint{5.632494in}{3.327059in}}%
\pgfpathlineto{\pgfqpoint{5.633054in}{3.335883in}}%
\pgfpathlineto{\pgfqpoint{5.633385in}{3.294834in}}%
\pgfpathlineto{\pgfqpoint{5.633453in}{3.300274in}}%
\pgfpathlineto{\pgfqpoint{5.634461in}{3.286279in}}%
\pgfpathlineto{\pgfqpoint{5.634091in}{3.320073in}}%
\pgfpathlineto{\pgfqpoint{5.634573in}{3.296866in}}%
\pgfpathlineto{\pgfqpoint{5.635244in}{3.304999in}}%
\pgfpathlineto{\pgfqpoint{5.635444in}{3.281359in}}%
\pgfpathlineto{\pgfqpoint{5.635454in}{3.281911in}}%
\pgfpathlineto{\pgfqpoint{5.635590in}{3.286280in}}%
\pgfpathlineto{\pgfqpoint{5.636598in}{3.256518in}}%
\pgfpathlineto{\pgfqpoint{5.636900in}{3.284824in}}%
\pgfpathlineto{\pgfqpoint{5.637688in}{3.250579in}}%
\pgfpathlineto{\pgfqpoint{5.637703in}{3.253531in}}%
\pgfpathlineto{\pgfqpoint{5.637723in}{3.250842in}}%
\pgfpathlineto{\pgfqpoint{5.638531in}{3.301472in}}%
\pgfpathlineto{\pgfqpoint{5.638652in}{3.293365in}}%
\pgfpathlineto{\pgfqpoint{5.639217in}{3.311686in}}%
\pgfpathlineto{\pgfqpoint{5.639641in}{3.271372in}}%
\pgfpathlineto{\pgfqpoint{5.639655in}{3.274263in}}%
\pgfpathlineto{\pgfqpoint{5.640765in}{3.251406in}}%
\pgfpathlineto{\pgfqpoint{5.640264in}{3.290557in}}%
\pgfpathlineto{\pgfqpoint{5.640824in}{3.262584in}}%
\pgfpathlineto{\pgfqpoint{5.642080in}{3.292798in}}%
\pgfpathlineto{\pgfqpoint{5.641077in}{3.255158in}}%
\pgfpathlineto{\pgfqpoint{5.642109in}{3.289370in}}%
\pgfpathlineto{\pgfqpoint{5.642143in}{3.282021in}}%
\pgfpathlineto{\pgfqpoint{5.642508in}{3.324351in}}%
\pgfpathlineto{\pgfqpoint{5.643170in}{3.307198in}}%
\pgfpathlineto{\pgfqpoint{5.643599in}{3.331702in}}%
\pgfpathlineto{\pgfqpoint{5.644052in}{3.298228in}}%
\pgfpathlineto{\pgfqpoint{5.644271in}{3.305163in}}%
\pgfpathlineto{\pgfqpoint{5.644285in}{3.300800in}}%
\pgfpathlineto{\pgfqpoint{5.644339in}{3.309572in}}%
\pgfpathlineto{\pgfqpoint{5.644353in}{3.312649in}}%
\pgfpathlineto{\pgfqpoint{5.645274in}{3.279694in}}%
\pgfpathlineto{\pgfqpoint{5.645434in}{3.305058in}}%
\pgfpathlineto{\pgfqpoint{5.646569in}{3.278573in}}%
\pgfpathlineto{\pgfqpoint{5.645702in}{3.319352in}}%
\pgfpathlineto{\pgfqpoint{5.646651in}{3.286021in}}%
\pgfpathlineto{\pgfqpoint{5.646661in}{3.286505in}}%
\pgfpathlineto{\pgfqpoint{5.647406in}{3.263564in}}%
\pgfpathlineto{\pgfqpoint{5.647440in}{3.264883in}}%
\pgfpathlineto{\pgfqpoint{5.647786in}{3.258407in}}%
\pgfpathlineto{\pgfqpoint{5.648419in}{3.292019in}}%
\pgfpathlineto{\pgfqpoint{5.648438in}{3.289763in}}%
\pgfpathlineto{\pgfqpoint{5.649237in}{3.317640in}}%
\pgfpathlineto{\pgfqpoint{5.648745in}{3.284780in}}%
\pgfpathlineto{\pgfqpoint{5.649563in}{3.294649in}}%
\pgfpathlineto{\pgfqpoint{5.650434in}{3.264358in}}%
\pgfpathlineto{\pgfqpoint{5.649728in}{3.307130in}}%
\pgfpathlineto{\pgfqpoint{5.650712in}{3.279820in}}%
\pgfpathlineto{\pgfqpoint{5.651569in}{3.311277in}}%
\pgfpathlineto{\pgfqpoint{5.650829in}{3.278420in}}%
\pgfpathlineto{\pgfqpoint{5.651836in}{3.287730in}}%
\pgfpathlineto{\pgfqpoint{5.652659in}{3.260977in}}%
\pgfpathlineto{\pgfqpoint{5.651870in}{3.291438in}}%
\pgfpathlineto{\pgfqpoint{5.653029in}{3.276155in}}%
\pgfpathlineto{\pgfqpoint{5.653049in}{3.283548in}}%
\pgfpathlineto{\pgfqpoint{5.653974in}{3.235720in}}%
\pgfpathlineto{\pgfqpoint{5.654037in}{3.232932in}}%
\pgfpathlineto{\pgfqpoint{5.654592in}{3.260135in}}%
\pgfpathlineto{\pgfqpoint{5.655025in}{3.255866in}}%
\pgfpathlineto{\pgfqpoint{5.655050in}{3.250648in}}%
\pgfpathlineto{\pgfqpoint{5.656603in}{3.314538in}}%
\pgfpathlineto{\pgfqpoint{5.657747in}{3.352668in}}%
\pgfpathlineto{\pgfqpoint{5.657221in}{3.307717in}}%
\pgfpathlineto{\pgfqpoint{5.657761in}{3.348996in}}%
\pgfpathlineto{\pgfqpoint{5.657771in}{3.350812in}}%
\pgfpathlineto{\pgfqpoint{5.657995in}{3.333254in}}%
\pgfpathlineto{\pgfqpoint{5.658204in}{3.334336in}}%
\pgfpathlineto{\pgfqpoint{5.659061in}{3.298150in}}%
\pgfpathlineto{\pgfqpoint{5.658263in}{3.335562in}}%
\pgfpathlineto{\pgfqpoint{5.659558in}{3.315041in}}%
\pgfpathlineto{\pgfqpoint{5.660497in}{3.323181in}}%
\pgfpathlineto{\pgfqpoint{5.660118in}{3.288516in}}%
\pgfpathlineto{\pgfqpoint{5.660639in}{3.308195in}}%
\pgfpathlineto{\pgfqpoint{5.661408in}{3.277917in}}%
\pgfpathlineto{\pgfqpoint{5.661783in}{3.292027in}}%
\pgfpathlineto{\pgfqpoint{5.661880in}{3.302384in}}%
\pgfpathlineto{\pgfqpoint{5.662795in}{3.259790in}}%
\pgfpathlineto{\pgfqpoint{5.662868in}{3.276070in}}%
\pgfpathlineto{\pgfqpoint{5.662883in}{3.270771in}}%
\pgfpathlineto{\pgfqpoint{5.663720in}{3.323806in}}%
\pgfpathlineto{\pgfqpoint{5.663930in}{3.311617in}}%
\pgfpathlineto{\pgfqpoint{5.664699in}{3.343431in}}%
\pgfpathlineto{\pgfqpoint{5.664426in}{3.300072in}}%
\pgfpathlineto{\pgfqpoint{5.665006in}{3.302471in}}%
\pgfpathlineto{\pgfqpoint{5.665346in}{3.283134in}}%
\pgfpathlineto{\pgfqpoint{5.665877in}{3.319680in}}%
\pgfpathlineto{\pgfqpoint{5.666125in}{3.293366in}}%
\pgfpathlineto{\pgfqpoint{5.666213in}{3.298732in}}%
\pgfpathlineto{\pgfqpoint{5.666622in}{3.252118in}}%
\pgfpathlineto{\pgfqpoint{5.667206in}{3.285217in}}%
\pgfpathlineto{\pgfqpoint{5.668360in}{3.247389in}}%
\pgfpathlineto{\pgfqpoint{5.667610in}{3.329705in}}%
\pgfpathlineto{\pgfqpoint{5.668380in}{3.252512in}}%
\pgfpathlineto{\pgfqpoint{5.668535in}{3.237811in}}%
\pgfpathlineto{\pgfqpoint{5.669631in}{3.306398in}}%
\pgfpathlineto{\pgfqpoint{5.669713in}{3.316239in}}%
\pgfpathlineto{\pgfqpoint{5.670828in}{3.258286in}}%
\pgfpathlineto{\pgfqpoint{5.671856in}{3.228343in}}%
\pgfpathlineto{\pgfqpoint{5.671091in}{3.260211in}}%
\pgfpathlineto{\pgfqpoint{5.671982in}{3.239293in}}%
\pgfpathlineto{\pgfqpoint{5.673053in}{3.283495in}}%
\pgfpathlineto{\pgfqpoint{5.672080in}{3.238001in}}%
\pgfpathlineto{\pgfqpoint{5.673151in}{3.276871in}}%
\pgfpathlineto{\pgfqpoint{5.673156in}{3.276868in}}%
\pgfpathlineto{\pgfqpoint{5.674251in}{3.289449in}}%
\pgfpathlineto{\pgfqpoint{5.673482in}{3.264647in}}%
\pgfpathlineto{\pgfqpoint{5.674275in}{3.285032in}}%
\pgfpathlineto{\pgfqpoint{5.674412in}{3.274847in}}%
\pgfpathlineto{\pgfqpoint{5.675210in}{3.302155in}}%
\pgfpathlineto{\pgfqpoint{5.675638in}{3.336183in}}%
\pgfpathlineto{\pgfqpoint{5.675322in}{3.296993in}}%
\pgfpathlineto{\pgfqpoint{5.676354in}{3.314568in}}%
\pgfpathlineto{\pgfqpoint{5.676690in}{3.323152in}}%
\pgfpathlineto{\pgfqpoint{5.677513in}{3.279279in}}%
\pgfpathlineto{\pgfqpoint{5.678428in}{3.260736in}}%
\pgfpathlineto{\pgfqpoint{5.678151in}{3.285378in}}%
\pgfpathlineto{\pgfqpoint{5.678628in}{3.278420in}}%
\pgfpathlineto{\pgfqpoint{5.679533in}{3.298589in}}%
\pgfpathlineto{\pgfqpoint{5.679051in}{3.262585in}}%
\pgfpathlineto{\pgfqpoint{5.679782in}{3.290538in}}%
\pgfpathlineto{\pgfqpoint{5.679967in}{3.275579in}}%
\pgfpathlineto{\pgfqpoint{5.680152in}{3.317883in}}%
\pgfpathlineto{\pgfqpoint{5.680746in}{3.302974in}}%
\pgfpathlineto{\pgfqpoint{5.680926in}{3.309544in}}%
\pgfpathlineto{\pgfqpoint{5.681695in}{3.262456in}}%
\pgfpathlineto{\pgfqpoint{5.681773in}{3.277133in}}%
\pgfpathlineto{\pgfqpoint{5.682333in}{3.256267in}}%
\pgfpathlineto{\pgfqpoint{5.682109in}{3.281675in}}%
\pgfpathlineto{\pgfqpoint{5.682893in}{3.275223in}}%
\pgfpathlineto{\pgfqpoint{5.684110in}{3.318524in}}%
\pgfpathlineto{\pgfqpoint{5.683326in}{3.261661in}}%
\pgfpathlineto{\pgfqpoint{5.684548in}{3.295736in}}%
\pgfpathlineto{\pgfqpoint{5.685653in}{3.270813in}}%
\pgfpathlineto{\pgfqpoint{5.684679in}{3.304565in}}%
\pgfpathlineto{\pgfqpoint{5.685682in}{3.272313in}}%
\pgfpathlineto{\pgfqpoint{5.685818in}{3.294158in}}%
\pgfpathlineto{\pgfqpoint{5.686208in}{3.258592in}}%
\pgfpathlineto{\pgfqpoint{5.686797in}{3.279486in}}%
\pgfpathlineto{\pgfqpoint{5.686807in}{3.274605in}}%
\pgfpathlineto{\pgfqpoint{5.687824in}{3.342474in}}%
\pgfpathlineto{\pgfqpoint{5.687829in}{3.340981in}}%
\pgfpathlineto{\pgfqpoint{5.688647in}{3.368544in}}%
\pgfpathlineto{\pgfqpoint{5.688891in}{3.339758in}}%
\pgfpathlineto{\pgfqpoint{5.688939in}{3.342287in}}%
\pgfpathlineto{\pgfqpoint{5.690044in}{3.289682in}}%
\pgfpathlineto{\pgfqpoint{5.690205in}{3.298194in}}%
\pgfpathlineto{\pgfqpoint{5.690215in}{3.298231in}}%
\pgfpathlineto{\pgfqpoint{5.690229in}{3.296469in}}%
\pgfpathlineto{\pgfqpoint{5.690677in}{3.279877in}}%
\pgfpathlineto{\pgfqpoint{5.690867in}{3.307933in}}%
\pgfpathlineto{\pgfqpoint{5.691359in}{3.280018in}}%
\pgfpathlineto{\pgfqpoint{5.692469in}{3.325329in}}%
\pgfpathlineto{\pgfqpoint{5.692591in}{3.314052in}}%
\pgfpathlineto{\pgfqpoint{5.693335in}{3.278833in}}%
\pgfpathlineto{\pgfqpoint{5.692985in}{3.320550in}}%
\pgfpathlineto{\pgfqpoint{5.693735in}{3.300236in}}%
\pgfpathlineto{\pgfqpoint{5.694845in}{3.325243in}}%
\pgfpathlineto{\pgfqpoint{5.694338in}{3.289787in}}%
\pgfpathlineto{\pgfqpoint{5.694874in}{3.324902in}}%
\pgfpathlineto{\pgfqpoint{5.695132in}{3.349754in}}%
\pgfpathlineto{\pgfqpoint{5.695648in}{3.314940in}}%
\pgfpathlineto{\pgfqpoint{5.696432in}{3.287323in}}%
\pgfpathlineto{\pgfqpoint{5.696081in}{3.318537in}}%
\pgfpathlineto{\pgfqpoint{5.696768in}{3.310435in}}%
\pgfpathlineto{\pgfqpoint{5.697795in}{3.339475in}}%
\pgfpathlineto{\pgfqpoint{5.697026in}{3.300470in}}%
\pgfpathlineto{\pgfqpoint{5.697931in}{3.321990in}}%
\pgfpathlineto{\pgfqpoint{5.698954in}{3.292458in}}%
\pgfpathlineto{\pgfqpoint{5.697956in}{3.325844in}}%
\pgfpathlineto{\pgfqpoint{5.699119in}{3.302122in}}%
\pgfpathlineto{\pgfqpoint{5.700093in}{3.354351in}}%
\pgfpathlineto{\pgfqpoint{5.699144in}{3.302003in}}%
\pgfpathlineto{\pgfqpoint{5.700346in}{3.339084in}}%
\pgfpathlineto{\pgfqpoint{5.700507in}{3.336662in}}%
\pgfpathlineto{\pgfqpoint{5.700784in}{3.356758in}}%
\pgfpathlineto{\pgfqpoint{5.701164in}{3.350433in}}%
\pgfpathlineto{\pgfqpoint{5.701188in}{3.352535in}}%
\pgfpathlineto{\pgfqpoint{5.701558in}{3.324264in}}%
\pgfpathlineto{\pgfqpoint{5.702040in}{3.332786in}}%
\pgfpathlineto{\pgfqpoint{5.702245in}{3.311034in}}%
\pgfpathlineto{\pgfqpoint{5.702070in}{3.336299in}}%
\pgfpathlineto{\pgfqpoint{5.703165in}{3.324664in}}%
\pgfpathlineto{\pgfqpoint{5.703783in}{3.346699in}}%
\pgfpathlineto{\pgfqpoint{5.703228in}{3.317943in}}%
\pgfpathlineto{\pgfqpoint{5.704270in}{3.322520in}}%
\pgfpathlineto{\pgfqpoint{5.704275in}{3.321976in}}%
\pgfpathlineto{\pgfqpoint{5.704864in}{3.353542in}}%
\pgfpathlineto{\pgfqpoint{5.705219in}{3.338835in}}%
\pgfpathlineto{\pgfqpoint{5.705745in}{3.378361in}}%
\pgfpathlineto{\pgfqpoint{5.706364in}{3.360322in}}%
\pgfpathlineto{\pgfqpoint{5.707590in}{3.298761in}}%
\pgfpathlineto{\pgfqpoint{5.706393in}{3.364949in}}%
\pgfpathlineto{\pgfqpoint{5.707634in}{3.303587in}}%
\pgfpathlineto{\pgfqpoint{5.707751in}{3.309615in}}%
\pgfpathlineto{\pgfqpoint{5.708082in}{3.253226in}}%
\pgfpathlineto{\pgfqpoint{5.708715in}{3.285120in}}%
\pgfpathlineto{\pgfqpoint{5.709265in}{3.344998in}}%
\pgfpathlineto{\pgfqpoint{5.710059in}{3.338063in}}%
\pgfpathlineto{\pgfqpoint{5.710122in}{3.324288in}}%
\pgfpathlineto{\pgfqpoint{5.710731in}{3.377445in}}%
\pgfpathlineto{\pgfqpoint{5.710828in}{3.373001in}}%
\pgfpathlineto{\pgfqpoint{5.711354in}{3.379100in}}%
\pgfpathlineto{\pgfqpoint{5.711602in}{3.355569in}}%
\pgfpathlineto{\pgfqpoint{5.711909in}{3.361793in}}%
\pgfpathlineto{\pgfqpoint{5.712551in}{3.320627in}}%
\pgfpathlineto{\pgfqpoint{5.712868in}{3.363611in}}%
\pgfpathlineto{\pgfqpoint{5.713014in}{3.362111in}}%
\pgfpathlineto{\pgfqpoint{5.713369in}{3.383467in}}%
\pgfpathlineto{\pgfqpoint{5.713394in}{3.383451in}}%
\pgfpathlineto{\pgfqpoint{5.713476in}{3.393126in}}%
\pgfpathlineto{\pgfqpoint{5.714372in}{3.357141in}}%
\pgfpathlineto{\pgfqpoint{5.714421in}{3.361142in}}%
\pgfpathlineto{\pgfqpoint{5.714606in}{3.342681in}}%
\pgfpathlineto{\pgfqpoint{5.715346in}{3.378582in}}%
\pgfpathlineto{\pgfqpoint{5.715546in}{3.354411in}}%
\pgfpathlineto{\pgfqpoint{5.716344in}{3.379360in}}%
\pgfpathlineto{\pgfqpoint{5.715687in}{3.351863in}}%
\pgfpathlineto{\pgfqpoint{5.716675in}{3.371102in}}%
\pgfpathlineto{\pgfqpoint{5.716768in}{3.377380in}}%
\pgfpathlineto{\pgfqpoint{5.717444in}{3.342568in}}%
\pgfpathlineto{\pgfqpoint{5.717678in}{3.353041in}}%
\pgfpathlineto{\pgfqpoint{5.717795in}{3.338746in}}%
\pgfpathlineto{\pgfqpoint{5.718355in}{3.407172in}}%
\pgfpathlineto{\pgfqpoint{5.718681in}{3.381540in}}%
\pgfpathlineto{\pgfqpoint{5.718764in}{3.393739in}}%
\pgfpathlineto{\pgfqpoint{5.718988in}{3.372242in}}%
\pgfpathlineto{\pgfqpoint{5.719509in}{3.376578in}}%
\pgfpathlineto{\pgfqpoint{5.720672in}{3.348641in}}%
\pgfpathlineto{\pgfqpoint{5.720687in}{3.350177in}}%
\pgfpathlineto{\pgfqpoint{5.720945in}{3.370346in}}%
\pgfpathlineto{\pgfqpoint{5.721480in}{3.336592in}}%
\pgfpathlineto{\pgfqpoint{5.721490in}{3.337537in}}%
\pgfpathlineto{\pgfqpoint{5.721733in}{3.334197in}}%
\pgfpathlineto{\pgfqpoint{5.722469in}{3.369010in}}%
\pgfpathlineto{\pgfqpoint{5.722478in}{3.367989in}}%
\pgfpathlineto{\pgfqpoint{5.723452in}{3.389723in}}%
\pgfpathlineto{\pgfqpoint{5.722498in}{3.363682in}}%
\pgfpathlineto{\pgfqpoint{5.723871in}{3.376710in}}%
\pgfpathlineto{\pgfqpoint{5.724518in}{3.353952in}}%
\pgfpathlineto{\pgfqpoint{5.724830in}{3.385599in}}%
\pgfpathlineto{\pgfqpoint{5.724981in}{3.373933in}}%
\pgfpathlineto{\pgfqpoint{5.725755in}{3.378476in}}%
\pgfpathlineto{\pgfqpoint{5.725307in}{3.346908in}}%
\pgfpathlineto{\pgfqpoint{5.726052in}{3.366864in}}%
\pgfpathlineto{\pgfqpoint{5.726514in}{3.334720in}}%
\pgfpathlineto{\pgfqpoint{5.726100in}{3.367153in}}%
\pgfpathlineto{\pgfqpoint{5.727181in}{3.362944in}}%
\pgfpathlineto{\pgfqpoint{5.728291in}{3.390082in}}%
\pgfpathlineto{\pgfqpoint{5.727220in}{3.360448in}}%
\pgfpathlineto{\pgfqpoint{5.728330in}{3.387914in}}%
\pgfpathlineto{\pgfqpoint{5.729275in}{3.361808in}}%
\pgfpathlineto{\pgfqpoint{5.728622in}{3.392788in}}%
\pgfpathlineto{\pgfqpoint{5.729450in}{3.382121in}}%
\pgfpathlineto{\pgfqpoint{5.729830in}{3.414614in}}%
\pgfpathlineto{\pgfqpoint{5.730565in}{3.389888in}}%
\pgfpathlineto{\pgfqpoint{5.731694in}{3.356924in}}%
\pgfpathlineto{\pgfqpoint{5.730594in}{3.392810in}}%
\pgfpathlineto{\pgfqpoint{5.731801in}{3.367317in}}%
\pgfpathlineto{\pgfqpoint{5.732342in}{3.380213in}}%
\pgfpathlineto{\pgfqpoint{5.732001in}{3.355446in}}%
\pgfpathlineto{\pgfqpoint{5.732902in}{3.366097in}}%
\pgfpathlineto{\pgfqpoint{5.733087in}{3.334878in}}%
\pgfpathlineto{\pgfqpoint{5.733992in}{3.368291in}}%
\pgfpathlineto{\pgfqpoint{5.734012in}{3.365012in}}%
\pgfpathlineto{\pgfqpoint{5.734226in}{3.350926in}}%
\pgfpathlineto{\pgfqpoint{5.734883in}{3.395554in}}%
\pgfpathlineto{\pgfqpoint{5.734898in}{3.394293in}}%
\pgfpathlineto{\pgfqpoint{5.735093in}{3.409580in}}%
\pgfpathlineto{\pgfqpoint{5.735954in}{3.367983in}}%
\pgfpathlineto{\pgfqpoint{5.735964in}{3.369528in}}%
\pgfpathlineto{\pgfqpoint{5.737011in}{3.304123in}}%
\pgfpathlineto{\pgfqpoint{5.735984in}{3.369606in}}%
\pgfpathlineto{\pgfqpoint{5.737235in}{3.317952in}}%
\pgfpathlineto{\pgfqpoint{5.737391in}{3.323953in}}%
\pgfpathlineto{\pgfqpoint{5.737576in}{3.300954in}}%
\pgfpathlineto{\pgfqpoint{5.738150in}{3.302605in}}%
\pgfpathlineto{\pgfqpoint{5.738432in}{3.264497in}}%
\pgfpathlineto{\pgfqpoint{5.739284in}{3.295330in}}%
\pgfpathlineto{\pgfqpoint{5.740122in}{3.326612in}}%
\pgfpathlineto{\pgfqpoint{5.739693in}{3.292009in}}%
\pgfpathlineto{\pgfqpoint{5.740472in}{3.310441in}}%
\pgfpathlineto{\pgfqpoint{5.740487in}{3.307896in}}%
\pgfpathlineto{\pgfqpoint{5.740881in}{3.270417in}}%
\pgfpathlineto{\pgfqpoint{5.741315in}{3.320414in}}%
\pgfpathlineto{\pgfqpoint{5.741636in}{3.302658in}}%
\pgfpathlineto{\pgfqpoint{5.742639in}{3.339043in}}%
\pgfpathlineto{\pgfqpoint{5.741646in}{3.301465in}}%
\pgfpathlineto{\pgfqpoint{5.742819in}{3.337729in}}%
\pgfpathlineto{\pgfqpoint{5.743564in}{3.297980in}}%
\pgfpathlineto{\pgfqpoint{5.742848in}{3.341177in}}%
\pgfpathlineto{\pgfqpoint{5.743953in}{3.322373in}}%
\pgfpathlineto{\pgfqpoint{5.745278in}{3.396232in}}%
\pgfpathlineto{\pgfqpoint{5.743982in}{3.320911in}}%
\pgfpathlineto{\pgfqpoint{5.745404in}{3.389927in}}%
\pgfpathlineto{\pgfqpoint{5.745560in}{3.363822in}}%
\pgfpathlineto{\pgfqpoint{5.746008in}{3.404696in}}%
\pgfpathlineto{\pgfqpoint{5.746524in}{3.386633in}}%
\pgfpathlineto{\pgfqpoint{5.747371in}{3.373902in}}%
\pgfpathlineto{\pgfqpoint{5.747108in}{3.395901in}}%
\pgfpathlineto{\pgfqpoint{5.747600in}{3.390479in}}%
\pgfpathlineto{\pgfqpoint{5.747751in}{3.378125in}}%
\pgfpathlineto{\pgfqpoint{5.748194in}{3.409712in}}%
\pgfpathlineto{\pgfqpoint{5.748447in}{3.402695in}}%
\pgfpathlineto{\pgfqpoint{5.748457in}{3.404415in}}%
\pgfpathlineto{\pgfqpoint{5.748700in}{3.382240in}}%
\pgfpathlineto{\pgfqpoint{5.749445in}{3.394365in}}%
\pgfpathlineto{\pgfqpoint{5.749864in}{3.370562in}}%
\pgfpathlineto{\pgfqpoint{5.750424in}{3.403535in}}%
\pgfpathlineto{\pgfqpoint{5.750565in}{3.391575in}}%
\pgfpathlineto{\pgfqpoint{5.750574in}{3.396134in}}%
\pgfpathlineto{\pgfqpoint{5.751611in}{3.344063in}}%
\pgfpathlineto{\pgfqpoint{5.751616in}{3.344782in}}%
\pgfpathlineto{\pgfqpoint{5.752527in}{3.319692in}}%
\pgfpathlineto{\pgfqpoint{5.751670in}{3.350957in}}%
\pgfpathlineto{\pgfqpoint{5.752726in}{3.344678in}}%
\pgfpathlineto{\pgfqpoint{5.753286in}{3.323153in}}%
\pgfpathlineto{\pgfqpoint{5.753622in}{3.354515in}}%
\pgfpathlineto{\pgfqpoint{5.753827in}{3.342974in}}%
\pgfpathlineto{\pgfqpoint{5.753890in}{3.351547in}}%
\pgfpathlineto{\pgfqpoint{5.754810in}{3.302565in}}%
\pgfpathlineto{\pgfqpoint{5.754849in}{3.305618in}}%
\pgfpathlineto{\pgfqpoint{5.755589in}{3.267094in}}%
\pgfpathlineto{\pgfqpoint{5.754903in}{3.312033in}}%
\pgfpathlineto{\pgfqpoint{5.755964in}{3.302970in}}%
\pgfpathlineto{\pgfqpoint{5.756475in}{3.287579in}}%
\pgfpathlineto{\pgfqpoint{5.756154in}{3.318083in}}%
\pgfpathlineto{\pgfqpoint{5.757045in}{3.305984in}}%
\pgfpathlineto{\pgfqpoint{5.757059in}{3.307188in}}%
\pgfpathlineto{\pgfqpoint{5.757916in}{3.270805in}}%
\pgfpathlineto{\pgfqpoint{5.758023in}{3.283668in}}%
\pgfpathlineto{\pgfqpoint{5.758121in}{3.277141in}}%
\pgfpathlineto{\pgfqpoint{5.758890in}{3.321910in}}%
\pgfpathlineto{\pgfqpoint{5.758904in}{3.321451in}}%
\pgfpathlineto{\pgfqpoint{5.759917in}{3.344827in}}%
\pgfpathlineto{\pgfqpoint{5.759138in}{3.309868in}}%
\pgfpathlineto{\pgfqpoint{5.760029in}{3.326474in}}%
\pgfpathlineto{\pgfqpoint{5.760365in}{3.299729in}}%
\pgfpathlineto{\pgfqpoint{5.760949in}{3.364852in}}%
\pgfpathlineto{\pgfqpoint{5.761051in}{3.353809in}}%
\pgfpathlineto{\pgfqpoint{5.762259in}{3.382726in}}%
\pgfpathlineto{\pgfqpoint{5.761071in}{3.347450in}}%
\pgfpathlineto{\pgfqpoint{5.762283in}{3.381718in}}%
\pgfpathlineto{\pgfqpoint{5.762298in}{3.378925in}}%
\pgfpathlineto{\pgfqpoint{5.762507in}{3.393649in}}%
\pgfpathlineto{\pgfqpoint{5.762663in}{3.392210in}}%
\pgfpathlineto{\pgfqpoint{5.763023in}{3.396947in}}%
\pgfpathlineto{\pgfqpoint{5.763525in}{3.367768in}}%
\pgfpathlineto{\pgfqpoint{5.763753in}{3.383078in}}%
\pgfpathlineto{\pgfqpoint{5.764235in}{3.370527in}}%
\pgfpathlineto{\pgfqpoint{5.764669in}{3.402294in}}%
\pgfpathlineto{\pgfqpoint{5.764776in}{3.395146in}}%
\pgfpathlineto{\pgfqpoint{5.765506in}{3.430931in}}%
\pgfpathlineto{\pgfqpoint{5.764829in}{3.387541in}}%
\pgfpathlineto{\pgfqpoint{5.765930in}{3.418370in}}%
\pgfpathlineto{\pgfqpoint{5.766115in}{3.407219in}}%
\pgfpathlineto{\pgfqpoint{5.766845in}{3.438928in}}%
\pgfpathlineto{\pgfqpoint{5.766972in}{3.434800in}}%
\pgfpathlineto{\pgfqpoint{5.768121in}{3.478372in}}%
\pgfpathlineto{\pgfqpoint{5.767376in}{3.430449in}}%
\pgfpathlineto{\pgfqpoint{5.768554in}{3.459815in}}%
\pgfpathlineto{\pgfqpoint{5.769742in}{3.421421in}}%
\pgfpathlineto{\pgfqpoint{5.768573in}{3.463112in}}%
\pgfpathlineto{\pgfqpoint{5.769781in}{3.429431in}}%
\pgfpathlineto{\pgfqpoint{5.770272in}{3.439697in}}%
\pgfpathlineto{\pgfqpoint{5.770730in}{3.401360in}}%
\pgfpathlineto{\pgfqpoint{5.770798in}{3.409046in}}%
\pgfpathlineto{\pgfqpoint{5.770803in}{3.408978in}}%
\pgfpathlineto{\pgfqpoint{5.770939in}{3.421990in}}%
\pgfpathlineto{\pgfqpoint{5.772079in}{3.463455in}}%
\pgfpathlineto{\pgfqpoint{5.771003in}{3.419410in}}%
\pgfpathlineto{\pgfqpoint{5.772103in}{3.460983in}}%
\pgfpathlineto{\pgfqpoint{5.772390in}{3.445394in}}%
\pgfpathlineto{\pgfqpoint{5.772838in}{3.474806in}}%
\pgfpathlineto{\pgfqpoint{5.773193in}{3.466485in}}%
\pgfpathlineto{\pgfqpoint{5.773257in}{3.474777in}}%
\pgfpathlineto{\pgfqpoint{5.773739in}{3.440182in}}%
\pgfpathlineto{\pgfqpoint{5.774080in}{3.446893in}}%
\pgfpathlineto{\pgfqpoint{5.774532in}{3.443308in}}%
\pgfpathlineto{\pgfqpoint{5.774863in}{3.469388in}}%
\pgfpathlineto{\pgfqpoint{5.775121in}{3.462303in}}%
\pgfpathlineto{\pgfqpoint{5.776431in}{3.544941in}}%
\pgfpathlineto{\pgfqpoint{5.775141in}{3.461962in}}%
\pgfpathlineto{\pgfqpoint{5.776485in}{3.542238in}}%
\pgfpathlineto{\pgfqpoint{5.777439in}{3.517831in}}%
\pgfpathlineto{\pgfqpoint{5.776923in}{3.547950in}}%
\pgfpathlineto{\pgfqpoint{5.777604in}{3.534804in}}%
\pgfpathlineto{\pgfqpoint{5.778666in}{3.554430in}}%
\pgfpathlineto{\pgfqpoint{5.777853in}{3.525559in}}%
\pgfpathlineto{\pgfqpoint{5.778724in}{3.542752in}}%
\pgfpathlineto{\pgfqpoint{5.779148in}{3.527249in}}%
\pgfpathlineto{\pgfqpoint{5.779586in}{3.560409in}}%
\pgfpathlineto{\pgfqpoint{5.779785in}{3.556157in}}%
\pgfpathlineto{\pgfqpoint{5.780837in}{3.559516in}}%
\pgfpathlineto{\pgfqpoint{5.780053in}{3.523340in}}%
\pgfpathlineto{\pgfqpoint{5.780876in}{3.552098in}}%
\pgfpathlineto{\pgfqpoint{5.781747in}{3.539013in}}%
\pgfpathlineto{\pgfqpoint{5.780910in}{3.559281in}}%
\pgfpathlineto{\pgfqpoint{5.782152in}{3.549720in}}%
\pgfpathlineto{\pgfqpoint{5.783266in}{3.611750in}}%
\pgfpathlineto{\pgfqpoint{5.782341in}{3.539747in}}%
\pgfpathlineto{\pgfqpoint{5.783412in}{3.595686in}}%
\pgfpathlineto{\pgfqpoint{5.784396in}{3.590070in}}%
\pgfpathlineto{\pgfqpoint{5.783671in}{3.618730in}}%
\pgfpathlineto{\pgfqpoint{5.784527in}{3.592259in}}%
\pgfpathlineto{\pgfqpoint{5.785662in}{3.637762in}}%
\pgfpathlineto{\pgfqpoint{5.784815in}{3.573934in}}%
\pgfpathlineto{\pgfqpoint{5.785720in}{3.633240in}}%
\pgfpathlineto{\pgfqpoint{5.786567in}{3.607382in}}%
\pgfpathlineto{\pgfqpoint{5.786879in}{3.616845in}}%
\pgfpathlineto{\pgfqpoint{5.787122in}{3.634644in}}%
\pgfpathlineto{\pgfqpoint{5.787419in}{3.608499in}}%
\pgfpathlineto{\pgfqpoint{5.788013in}{3.633189in}}%
\pgfpathlineto{\pgfqpoint{5.788583in}{3.601619in}}%
\pgfpathlineto{\pgfqpoint{5.789109in}{3.634682in}}%
\pgfpathlineto{\pgfqpoint{5.789123in}{3.630899in}}%
\pgfpathlineto{\pgfqpoint{5.789498in}{3.597261in}}%
\pgfpathlineto{\pgfqpoint{5.790399in}{3.625212in}}%
\pgfpathlineto{\pgfqpoint{5.790467in}{3.629769in}}%
\pgfpathlineto{\pgfqpoint{5.791207in}{3.602530in}}%
\pgfpathlineto{\pgfqpoint{5.791446in}{3.604716in}}%
\pgfpathlineto{\pgfqpoint{5.791884in}{3.589226in}}%
\pgfpathlineto{\pgfqpoint{5.792259in}{3.617985in}}%
\pgfpathlineto{\pgfqpoint{5.792536in}{3.609639in}}%
\pgfpathlineto{\pgfqpoint{5.793003in}{3.641328in}}%
\pgfpathlineto{\pgfqpoint{5.792599in}{3.600379in}}%
\pgfpathlineto{\pgfqpoint{5.793636in}{3.608119in}}%
\pgfpathlineto{\pgfqpoint{5.794109in}{3.591543in}}%
\pgfpathlineto{\pgfqpoint{5.794581in}{3.631594in}}%
\pgfpathlineto{\pgfqpoint{5.794668in}{3.620504in}}%
\pgfpathlineto{\pgfqpoint{5.795413in}{3.665539in}}%
\pgfpathlineto{\pgfqpoint{5.795769in}{3.615701in}}%
\pgfpathlineto{\pgfqpoint{5.795778in}{3.622069in}}%
\pgfpathlineto{\pgfqpoint{5.795959in}{3.609166in}}%
\pgfpathlineto{\pgfqpoint{5.796416in}{3.650852in}}%
\pgfpathlineto{\pgfqpoint{5.796854in}{3.627314in}}%
\pgfpathlineto{\pgfqpoint{5.797789in}{3.635282in}}%
\pgfpathlineto{\pgfqpoint{5.797161in}{3.594282in}}%
\pgfpathlineto{\pgfqpoint{5.797940in}{3.618630in}}%
\pgfpathlineto{\pgfqpoint{5.799425in}{3.590036in}}%
\pgfpathlineto{\pgfqpoint{5.798149in}{3.626094in}}%
\pgfpathlineto{\pgfqpoint{5.799449in}{3.591411in}}%
\pgfpathlineto{\pgfqpoint{5.800399in}{3.622464in}}%
\pgfpathlineto{\pgfqpoint{5.799532in}{3.587897in}}%
\pgfpathlineto{\pgfqpoint{5.800623in}{3.615421in}}%
\pgfpathlineto{\pgfqpoint{5.801742in}{3.581834in}}%
\pgfpathlineto{\pgfqpoint{5.801324in}{3.623720in}}%
\pgfpathlineto{\pgfqpoint{5.801791in}{3.583479in}}%
\pgfpathlineto{\pgfqpoint{5.802716in}{3.573642in}}%
\pgfpathlineto{\pgfqpoint{5.802356in}{3.605329in}}%
\pgfpathlineto{\pgfqpoint{5.802886in}{3.585533in}}%
\pgfpathlineto{\pgfqpoint{5.803106in}{3.602664in}}%
\pgfpathlineto{\pgfqpoint{5.803607in}{3.577993in}}%
\pgfpathlineto{\pgfqpoint{5.803651in}{3.569361in}}%
\pgfpathlineto{\pgfqpoint{5.804070in}{3.604775in}}%
\pgfpathlineto{\pgfqpoint{5.804693in}{3.597416in}}%
\pgfpathlineto{\pgfqpoint{5.805326in}{3.575507in}}%
\pgfpathlineto{\pgfqpoint{5.805788in}{3.600526in}}%
\pgfpathlineto{\pgfqpoint{5.805793in}{3.604238in}}%
\pgfpathlineto{\pgfqpoint{5.806270in}{3.577037in}}%
\pgfpathlineto{\pgfqpoint{5.806879in}{3.588990in}}%
\pgfpathlineto{\pgfqpoint{5.807886in}{3.551276in}}%
\pgfpathlineto{\pgfqpoint{5.806932in}{3.590696in}}%
\pgfpathlineto{\pgfqpoint{5.808047in}{3.558762in}}%
\pgfpathlineto{\pgfqpoint{5.808130in}{3.578327in}}%
\pgfpathlineto{\pgfqpoint{5.809031in}{3.549749in}}%
\pgfpathlineto{\pgfqpoint{5.809152in}{3.556985in}}%
\pgfpathlineto{\pgfqpoint{5.809181in}{3.554348in}}%
\pgfpathlineto{\pgfqpoint{5.810063in}{3.600468in}}%
\pgfpathlineto{\pgfqpoint{5.810126in}{3.591642in}}%
\pgfpathlineto{\pgfqpoint{5.811134in}{3.648229in}}%
\pgfpathlineto{\pgfqpoint{5.810160in}{3.589679in}}%
\pgfpathlineto{\pgfqpoint{5.811290in}{3.639037in}}%
\pgfpathlineto{\pgfqpoint{5.812351in}{3.654662in}}%
\pgfpathlineto{\pgfqpoint{5.811562in}{3.625126in}}%
\pgfpathlineto{\pgfqpoint{5.812438in}{3.647821in}}%
\pgfpathlineto{\pgfqpoint{5.813451in}{3.596520in}}%
\pgfpathlineto{\pgfqpoint{5.812468in}{3.649998in}}%
\pgfpathlineto{\pgfqpoint{5.813626in}{3.617397in}}%
\pgfpathlineto{\pgfqpoint{5.814741in}{3.647735in}}%
\pgfpathlineto{\pgfqpoint{5.813787in}{3.615305in}}%
\pgfpathlineto{\pgfqpoint{5.814766in}{3.644611in}}%
\pgfpathlineto{\pgfqpoint{5.814951in}{3.631300in}}%
\pgfpathlineto{\pgfqpoint{5.815457in}{3.681253in}}%
\pgfpathlineto{\pgfqpoint{5.815798in}{3.669156in}}%
\pgfpathlineto{\pgfqpoint{5.815944in}{3.691469in}}%
\pgfpathlineto{\pgfqpoint{5.816640in}{3.661481in}}%
\pgfpathlineto{\pgfqpoint{5.816883in}{3.664098in}}%
\pgfpathlineto{\pgfqpoint{5.817760in}{3.632510in}}%
\pgfpathlineto{\pgfqpoint{5.817526in}{3.670964in}}%
\pgfpathlineto{\pgfqpoint{5.818032in}{3.645918in}}%
\pgfpathlineto{\pgfqpoint{5.818933in}{3.707113in}}%
\pgfpathlineto{\pgfqpoint{5.819288in}{3.695088in}}%
\pgfpathlineto{\pgfqpoint{5.819298in}{3.696635in}}%
\pgfpathlineto{\pgfqpoint{5.819649in}{3.665967in}}%
\pgfpathlineto{\pgfqpoint{5.820277in}{3.652251in}}%
\pgfpathlineto{\pgfqpoint{5.820676in}{3.676386in}}%
\pgfpathlineto{\pgfqpoint{5.820764in}{3.662789in}}%
\pgfpathlineto{\pgfqpoint{5.820807in}{3.662339in}}%
\pgfpathlineto{\pgfqpoint{5.821085in}{3.682014in}}%
\pgfpathlineto{\pgfqpoint{5.821158in}{3.680251in}}%
\pgfpathlineto{\pgfqpoint{5.822200in}{3.706236in}}%
\pgfpathlineto{\pgfqpoint{5.821547in}{3.666762in}}%
\pgfpathlineto{\pgfqpoint{5.822302in}{3.696510in}}%
\pgfpathlineto{\pgfqpoint{5.823519in}{3.671084in}}%
\pgfpathlineto{\pgfqpoint{5.822750in}{3.703292in}}%
\pgfpathlineto{\pgfqpoint{5.823534in}{3.676011in}}%
\pgfpathlineto{\pgfqpoint{5.824863in}{3.725119in}}%
\pgfpathlineto{\pgfqpoint{5.823553in}{3.672185in}}%
\pgfpathlineto{\pgfqpoint{5.824887in}{3.718325in}}%
\pgfpathlineto{\pgfqpoint{5.825924in}{3.684748in}}%
\pgfpathlineto{\pgfqpoint{5.824989in}{3.729312in}}%
\pgfpathlineto{\pgfqpoint{5.826031in}{3.700526in}}%
\pgfpathlineto{\pgfqpoint{5.826070in}{3.689515in}}%
\pgfpathlineto{\pgfqpoint{5.826431in}{3.752938in}}%
\pgfpathlineto{\pgfqpoint{5.826981in}{3.739805in}}%
\pgfpathlineto{\pgfqpoint{5.827594in}{3.744901in}}%
\pgfpathlineto{\pgfqpoint{5.827054in}{3.726315in}}%
\pgfpathlineto{\pgfqpoint{5.827945in}{3.728680in}}%
\pgfpathlineto{\pgfqpoint{5.827984in}{3.728837in}}%
\pgfpathlineto{\pgfqpoint{5.829210in}{3.668021in}}%
\pgfpathlineto{\pgfqpoint{5.829235in}{3.666265in}}%
\pgfpathlineto{\pgfqpoint{5.829975in}{3.700368in}}%
\pgfpathlineto{\pgfqpoint{5.829999in}{3.699000in}}%
\pgfpathlineto{\pgfqpoint{5.831226in}{3.729516in}}%
\pgfpathlineto{\pgfqpoint{5.830408in}{3.695548in}}%
\pgfpathlineto{\pgfqpoint{5.831231in}{3.728002in}}%
\pgfpathlineto{\pgfqpoint{5.831241in}{3.722044in}}%
\pgfpathlineto{\pgfqpoint{5.832302in}{3.760128in}}%
\pgfpathlineto{\pgfqpoint{5.832307in}{3.758370in}}%
\pgfpathlineto{\pgfqpoint{5.832536in}{3.746347in}}%
\pgfpathlineto{\pgfqpoint{5.833558in}{3.833491in}}%
\pgfpathlineto{\pgfqpoint{5.834420in}{3.811374in}}%
\pgfpathlineto{\pgfqpoint{5.834001in}{3.848798in}}%
\pgfpathlineto{\pgfqpoint{5.834663in}{3.833731in}}%
\pgfpathlineto{\pgfqpoint{5.834829in}{3.856632in}}%
\pgfpathlineto{\pgfqpoint{5.835369in}{3.804763in}}%
\pgfpathlineto{\pgfqpoint{5.835763in}{3.825938in}}%
\pgfpathlineto{\pgfqpoint{5.836216in}{3.815186in}}%
\pgfpathlineto{\pgfqpoint{5.836698in}{3.853079in}}%
\pgfpathlineto{\pgfqpoint{5.836752in}{3.848190in}}%
\pgfpathlineto{\pgfqpoint{5.837998in}{3.883368in}}%
\pgfpathlineto{\pgfqpoint{5.836766in}{3.845117in}}%
\pgfpathlineto{\pgfqpoint{5.838310in}{3.870888in}}%
\pgfpathlineto{\pgfqpoint{5.838548in}{3.855072in}}%
\pgfpathlineto{\pgfqpoint{5.839191in}{3.900387in}}%
\pgfpathlineto{\pgfqpoint{5.839337in}{3.893115in}}%
\pgfpathlineto{\pgfqpoint{5.839687in}{3.901799in}}%
\pgfpathlineto{\pgfqpoint{5.839517in}{3.887580in}}%
\pgfpathlineto{\pgfqpoint{5.839692in}{3.900897in}}%
\pgfpathlineto{\pgfqpoint{5.840695in}{3.946797in}}%
\pgfpathlineto{\pgfqpoint{5.840004in}{3.885117in}}%
\pgfpathlineto{\pgfqpoint{5.840939in}{3.944588in}}%
\pgfpathlineto{\pgfqpoint{5.841635in}{3.918235in}}%
\pgfpathlineto{\pgfqpoint{5.841250in}{3.948041in}}%
\pgfpathlineto{\pgfqpoint{5.842063in}{3.942499in}}%
\pgfpathlineto{\pgfqpoint{5.842068in}{3.942636in}}%
\pgfpathlineto{\pgfqpoint{5.842346in}{3.923341in}}%
\pgfpathlineto{\pgfqpoint{5.843193in}{3.875162in}}%
\pgfpathlineto{\pgfqpoint{5.843490in}{3.890743in}}%
\pgfpathlineto{\pgfqpoint{5.843500in}{3.893409in}}%
\pgfpathlineto{\pgfqpoint{5.844162in}{3.869418in}}%
\pgfpathlineto{\pgfqpoint{5.845174in}{3.860922in}}%
\pgfpathlineto{\pgfqpoint{5.844580in}{3.889868in}}%
\pgfpathlineto{\pgfqpoint{5.845277in}{3.866707in}}%
\pgfpathlineto{\pgfqpoint{5.846279in}{3.889133in}}%
\pgfpathlineto{\pgfqpoint{5.845802in}{3.846429in}}%
\pgfpathlineto{\pgfqpoint{5.846396in}{3.876340in}}%
\pgfpathlineto{\pgfqpoint{5.846503in}{3.870125in}}%
\pgfpathlineto{\pgfqpoint{5.847346in}{3.914212in}}%
\pgfpathlineto{\pgfqpoint{5.847448in}{3.904646in}}%
\pgfpathlineto{\pgfqpoint{5.848134in}{3.929416in}}%
\pgfpathlineto{\pgfqpoint{5.847492in}{3.902644in}}%
\pgfpathlineto{\pgfqpoint{5.848548in}{3.904515in}}%
\pgfpathlineto{\pgfqpoint{5.849522in}{3.878810in}}%
\pgfpathlineto{\pgfqpoint{5.848582in}{3.906232in}}%
\pgfpathlineto{\pgfqpoint{5.849692in}{3.882715in}}%
\pgfpathlineto{\pgfqpoint{5.849731in}{3.885334in}}%
\pgfpathlineto{\pgfqpoint{5.850218in}{3.862309in}}%
\pgfpathlineto{\pgfqpoint{5.850666in}{3.878088in}}%
\pgfpathlineto{\pgfqpoint{5.850744in}{3.870538in}}%
\pgfpathlineto{\pgfqpoint{5.851654in}{3.922390in}}%
\pgfpathlineto{\pgfqpoint{5.851742in}{3.914391in}}%
\pgfpathlineto{\pgfqpoint{5.853256in}{3.996429in}}%
\pgfpathlineto{\pgfqpoint{5.854561in}{3.919548in}}%
\pgfpathlineto{\pgfqpoint{5.853271in}{3.997276in}}%
\pgfpathlineto{\pgfqpoint{5.854950in}{3.953957in}}%
\pgfpathlineto{\pgfqpoint{5.855554in}{3.963408in}}%
\pgfpathlineto{\pgfqpoint{5.855880in}{3.938093in}}%
\pgfpathlineto{\pgfqpoint{5.856002in}{3.941202in}}%
\pgfpathlineto{\pgfqpoint{5.857102in}{4.001706in}}%
\pgfpathlineto{\pgfqpoint{5.856021in}{3.934110in}}%
\pgfpathlineto{\pgfqpoint{5.858061in}{3.996142in}}%
\pgfpathlineto{\pgfqpoint{5.858295in}{3.981928in}}%
\pgfpathlineto{\pgfqpoint{5.858665in}{4.014279in}}%
\pgfpathlineto{\pgfqpoint{5.859162in}{4.003721in}}%
\pgfpathlineto{\pgfqpoint{5.860505in}{4.042135in}}%
\pgfpathlineto{\pgfqpoint{5.860198in}{3.999792in}}%
\pgfpathlineto{\pgfqpoint{5.860515in}{4.040018in}}%
\pgfpathlineto{\pgfqpoint{5.861143in}{4.021137in}}%
\pgfpathlineto{\pgfqpoint{5.860646in}{4.051948in}}%
\pgfpathlineto{\pgfqpoint{5.861649in}{4.033707in}}%
\pgfpathlineto{\pgfqpoint{5.862000in}{4.087655in}}%
\pgfpathlineto{\pgfqpoint{5.861683in}{4.033332in}}%
\pgfpathlineto{\pgfqpoint{5.862871in}{4.080918in}}%
\pgfpathlineto{\pgfqpoint{5.862886in}{4.078445in}}%
\pgfpathlineto{\pgfqpoint{5.863714in}{4.121286in}}%
\pgfpathlineto{\pgfqpoint{5.863864in}{4.107725in}}%
\pgfpathlineto{\pgfqpoint{5.864278in}{4.131730in}}%
\pgfpathlineto{\pgfqpoint{5.864872in}{4.077432in}}%
\pgfpathlineto{\pgfqpoint{5.864940in}{4.084069in}}%
\pgfpathlineto{\pgfqpoint{5.866226in}{4.032402in}}%
\pgfpathlineto{\pgfqpoint{5.865057in}{4.087063in}}%
\pgfpathlineto{\pgfqpoint{5.866235in}{4.032781in}}%
\pgfpathlineto{\pgfqpoint{5.866245in}{4.033913in}}%
\pgfpathlineto{\pgfqpoint{5.867107in}{3.989009in}}%
\pgfpathlineto{\pgfqpoint{5.868105in}{3.976340in}}%
\pgfpathlineto{\pgfqpoint{5.867910in}{4.000337in}}%
\pgfpathlineto{\pgfqpoint{5.868241in}{3.981887in}}%
\pgfpathlineto{\pgfqpoint{5.868382in}{3.991095in}}%
\pgfpathlineto{\pgfqpoint{5.868986in}{3.964807in}}%
\pgfpathlineto{\pgfqpoint{5.869142in}{3.968645in}}%
\pgfpathlineto{\pgfqpoint{5.869161in}{3.961886in}}%
\pgfpathlineto{\pgfqpoint{5.869984in}{4.034037in}}%
\pgfpathlineto{\pgfqpoint{5.870184in}{4.026155in}}%
\pgfpathlineto{\pgfqpoint{5.870958in}{4.050813in}}%
\pgfpathlineto{\pgfqpoint{5.870700in}{4.020553in}}%
\pgfpathlineto{\pgfqpoint{5.871269in}{4.024599in}}%
\pgfpathlineto{\pgfqpoint{5.871377in}{4.017209in}}%
\pgfpathlineto{\pgfqpoint{5.871703in}{4.041313in}}%
\pgfpathlineto{\pgfqpoint{5.872360in}{4.036254in}}%
\pgfpathlineto{\pgfqpoint{5.872550in}{4.042334in}}%
\pgfpathlineto{\pgfqpoint{5.872788in}{4.020332in}}%
\pgfpathlineto{\pgfqpoint{5.872793in}{4.020617in}}%
\pgfpathlineto{\pgfqpoint{5.873894in}{3.960348in}}%
\pgfpathlineto{\pgfqpoint{5.872837in}{4.022061in}}%
\pgfpathlineto{\pgfqpoint{5.874293in}{3.985124in}}%
\pgfpathlineto{\pgfqpoint{5.875588in}{4.023274in}}%
\pgfpathlineto{\pgfqpoint{5.874395in}{3.978027in}}%
\pgfpathlineto{\pgfqpoint{5.875593in}{4.021678in}}%
\pgfpathlineto{\pgfqpoint{5.876133in}{3.995273in}}%
\pgfpathlineto{\pgfqpoint{5.875622in}{4.026620in}}%
\pgfpathlineto{\pgfqpoint{5.876712in}{4.011076in}}%
\pgfpathlineto{\pgfqpoint{5.877827in}{4.050996in}}%
\pgfpathlineto{\pgfqpoint{5.877881in}{4.045656in}}%
\pgfpathlineto{\pgfqpoint{5.877983in}{4.039461in}}%
\pgfpathlineto{\pgfqpoint{5.878115in}{4.052927in}}%
\pgfpathlineto{\pgfqpoint{5.878261in}{4.051511in}}%
\pgfpathlineto{\pgfqpoint{5.879468in}{4.084179in}}%
\pgfpathlineto{\pgfqpoint{5.878368in}{4.042037in}}%
\pgfpathlineto{\pgfqpoint{5.879473in}{4.083491in}}%
\pgfpathlineto{\pgfqpoint{5.879814in}{4.058209in}}%
\pgfpathlineto{\pgfqpoint{5.879570in}{4.087330in}}%
\pgfpathlineto{\pgfqpoint{5.880607in}{4.064597in}}%
\pgfpathlineto{\pgfqpoint{5.880636in}{4.069214in}}%
\pgfpathlineto{\pgfqpoint{5.881109in}{4.023942in}}%
\pgfpathlineto{\pgfqpoint{5.881518in}{4.036723in}}%
\pgfpathlineto{\pgfqpoint{5.881737in}{4.020071in}}%
\pgfpathlineto{\pgfqpoint{5.882545in}{4.056785in}}%
\pgfpathlineto{\pgfqpoint{5.882594in}{4.052763in}}%
\pgfpathlineto{\pgfqpoint{5.882603in}{4.054688in}}%
\pgfpathlineto{\pgfqpoint{5.883402in}{4.031029in}}%
\pgfpathlineto{\pgfqpoint{5.883450in}{4.034041in}}%
\pgfpathlineto{\pgfqpoint{5.883592in}{4.031753in}}%
\pgfpathlineto{\pgfqpoint{5.883952in}{4.056294in}}%
\pgfpathlineto{\pgfqpoint{5.884405in}{4.045770in}}%
\pgfpathlineto{\pgfqpoint{5.885554in}{4.088341in}}%
\pgfpathlineto{\pgfqpoint{5.885622in}{4.077897in}}%
\pgfpathlineto{\pgfqpoint{5.886260in}{4.116842in}}%
\pgfpathlineto{\pgfqpoint{5.886274in}{4.114423in}}%
\pgfpathlineto{\pgfqpoint{5.886513in}{4.122131in}}%
\pgfpathlineto{\pgfqpoint{5.887209in}{4.082549in}}%
\pgfpathlineto{\pgfqpoint{5.887345in}{4.096199in}}%
\pgfpathlineto{\pgfqpoint{5.887413in}{4.086182in}}%
\pgfpathlineto{\pgfqpoint{5.888212in}{4.131534in}}%
\pgfpathlineto{\pgfqpoint{5.888226in}{4.130497in}}%
\pgfpathlineto{\pgfqpoint{5.888285in}{4.133162in}}%
\pgfpathlineto{\pgfqpoint{5.888796in}{4.101470in}}%
\pgfpathlineto{\pgfqpoint{5.889312in}{4.122808in}}%
\pgfpathlineto{\pgfqpoint{5.890524in}{4.054261in}}%
\pgfpathlineto{\pgfqpoint{5.890666in}{4.068934in}}%
\pgfpathlineto{\pgfqpoint{5.891595in}{4.098597in}}%
\pgfpathlineto{\pgfqpoint{5.891849in}{4.083728in}}%
\pgfpathlineto{\pgfqpoint{5.892618in}{4.065681in}}%
\pgfpathlineto{\pgfqpoint{5.892428in}{4.098849in}}%
\pgfpathlineto{\pgfqpoint{5.892964in}{4.078436in}}%
\pgfpathlineto{\pgfqpoint{5.893431in}{4.092754in}}%
\pgfpathlineto{\pgfqpoint{5.893889in}{4.059550in}}%
\pgfpathlineto{\pgfqpoint{5.893962in}{4.059851in}}%
\pgfpathlineto{\pgfqpoint{5.894916in}{4.035914in}}%
\pgfpathlineto{\pgfqpoint{5.893981in}{4.062639in}}%
\pgfpathlineto{\pgfqpoint{5.895101in}{4.050970in}}%
\pgfpathlineto{\pgfqpoint{5.896040in}{4.055662in}}%
\pgfpathlineto{\pgfqpoint{5.895505in}{4.029154in}}%
\pgfpathlineto{\pgfqpoint{5.896177in}{4.050701in}}%
\pgfpathlineto{\pgfqpoint{5.896250in}{4.043543in}}%
\pgfpathlineto{\pgfqpoint{5.896527in}{4.083892in}}%
\pgfpathlineto{\pgfqpoint{5.897262in}{4.064847in}}%
\pgfpathlineto{\pgfqpoint{5.898606in}{4.001628in}}%
\pgfpathlineto{\pgfqpoint{5.897574in}{4.077196in}}%
\pgfpathlineto{\pgfqpoint{5.898669in}{4.008001in}}%
\pgfpathlineto{\pgfqpoint{5.899473in}{4.029555in}}%
\pgfpathlineto{\pgfqpoint{5.899765in}{4.006969in}}%
\pgfpathlineto{\pgfqpoint{5.899779in}{4.008204in}}%
\pgfpathlineto{\pgfqpoint{5.900651in}{4.025798in}}%
\pgfpathlineto{\pgfqpoint{5.900315in}{3.989788in}}%
\pgfpathlineto{\pgfqpoint{5.900812in}{3.999199in}}%
\pgfpathlineto{\pgfqpoint{5.901552in}{3.977619in}}%
\pgfpathlineto{\pgfqpoint{5.901805in}{4.020586in}}%
\pgfpathlineto{\pgfqpoint{5.901907in}{4.009254in}}%
\pgfpathlineto{\pgfqpoint{5.902681in}{4.037982in}}%
\pgfpathlineto{\pgfqpoint{5.902257in}{4.002501in}}%
\pgfpathlineto{\pgfqpoint{5.903109in}{4.030619in}}%
\pgfpathlineto{\pgfqpoint{5.903139in}{4.022826in}}%
\pgfpathlineto{\pgfqpoint{5.904132in}{4.056762in}}%
\pgfpathlineto{\pgfqpoint{5.904190in}{4.053740in}}%
\pgfpathlineto{\pgfqpoint{5.904224in}{4.057485in}}%
\pgfpathlineto{\pgfqpoint{5.904945in}{4.034710in}}%
\pgfpathlineto{\pgfqpoint{5.905256in}{4.045745in}}%
\pgfpathlineto{\pgfqpoint{5.905388in}{4.040236in}}%
\pgfpathlineto{\pgfqpoint{5.905665in}{4.061972in}}%
\pgfpathlineto{\pgfqpoint{5.905680in}{4.064307in}}%
\pgfpathlineto{\pgfqpoint{5.906532in}{4.020238in}}%
\pgfpathlineto{\pgfqpoint{5.906702in}{4.039633in}}%
\pgfpathlineto{\pgfqpoint{5.906741in}{4.036207in}}%
\pgfpathlineto{\pgfqpoint{5.907647in}{4.073795in}}%
\pgfpathlineto{\pgfqpoint{5.907700in}{4.064106in}}%
\pgfpathlineto{\pgfqpoint{5.907739in}{4.067114in}}%
\pgfpathlineto{\pgfqpoint{5.908621in}{4.023564in}}%
\pgfpathlineto{\pgfqpoint{5.908718in}{4.033337in}}%
\pgfpathlineto{\pgfqpoint{5.909361in}{4.014907in}}%
\pgfpathlineto{\pgfqpoint{5.908903in}{4.041448in}}%
\pgfpathlineto{\pgfqpoint{5.909828in}{4.032597in}}%
\pgfpathlineto{\pgfqpoint{5.909916in}{4.022419in}}%
\pgfpathlineto{\pgfqpoint{5.910661in}{4.063934in}}%
\pgfpathlineto{\pgfqpoint{5.910709in}{4.061124in}}%
\pgfpathlineto{\pgfqpoint{5.910719in}{4.057391in}}%
\pgfpathlineto{\pgfqpoint{5.910962in}{4.050957in}}%
\pgfpathlineto{\pgfqpoint{5.911172in}{4.081540in}}%
\pgfpathlineto{\pgfqpoint{5.911790in}{4.071552in}}%
\pgfpathlineto{\pgfqpoint{5.912769in}{4.080018in}}%
\pgfpathlineto{\pgfqpoint{5.912199in}{4.049899in}}%
\pgfpathlineto{\pgfqpoint{5.912910in}{4.075761in}}%
\pgfpathlineto{\pgfqpoint{5.913874in}{4.066610in}}%
\pgfpathlineto{\pgfqpoint{5.913635in}{4.094266in}}%
\pgfpathlineto{\pgfqpoint{5.914025in}{4.074041in}}%
\pgfpathlineto{\pgfqpoint{5.915086in}{4.088335in}}%
\pgfpathlineto{\pgfqpoint{5.914438in}{4.049748in}}%
\pgfpathlineto{\pgfqpoint{5.915198in}{4.088087in}}%
\pgfpathlineto{\pgfqpoint{5.916868in}{4.012937in}}%
\pgfpathlineto{\pgfqpoint{5.915266in}{4.098167in}}%
\pgfpathlineto{\pgfqpoint{5.917067in}{4.022306in}}%
\pgfpathlineto{\pgfqpoint{5.917696in}{4.047366in}}%
\pgfpathlineto{\pgfqpoint{5.917150in}{4.013425in}}%
\pgfpathlineto{\pgfqpoint{5.918192in}{4.031230in}}%
\pgfpathlineto{\pgfqpoint{5.919214in}{4.041191in}}%
\pgfpathlineto{\pgfqpoint{5.918854in}{4.011583in}}%
\pgfpathlineto{\pgfqpoint{5.919278in}{4.028509in}}%
\pgfpathlineto{\pgfqpoint{5.920101in}{4.003843in}}%
\pgfpathlineto{\pgfqpoint{5.919857in}{4.037374in}}%
\pgfpathlineto{\pgfqpoint{5.920456in}{4.015782in}}%
\pgfpathlineto{\pgfqpoint{5.920495in}{4.028557in}}%
\pgfpathlineto{\pgfqpoint{5.920967in}{4.000204in}}%
\pgfpathlineto{\pgfqpoint{5.921537in}{4.008000in}}%
\pgfpathlineto{\pgfqpoint{5.921542in}{4.008191in}}%
\pgfpathlineto{\pgfqpoint{5.922033in}{3.984545in}}%
\pgfpathlineto{\pgfqpoint{5.922832in}{3.966023in}}%
\pgfpathlineto{\pgfqpoint{5.922179in}{3.998038in}}%
\pgfpathlineto{\pgfqpoint{5.923148in}{3.980562in}}%
\pgfpathlineto{\pgfqpoint{5.923470in}{4.008363in}}%
\pgfpathlineto{\pgfqpoint{5.924370in}{3.997676in}}%
\pgfpathlineto{\pgfqpoint{5.924614in}{3.977071in}}%
\pgfpathlineto{\pgfqpoint{5.924458in}{4.004477in}}%
\pgfpathlineto{\pgfqpoint{5.925490in}{3.989184in}}%
\pgfpathlineto{\pgfqpoint{5.926030in}{4.003865in}}%
\pgfpathlineto{\pgfqpoint{5.925665in}{3.981457in}}%
\pgfpathlineto{\pgfqpoint{5.926615in}{3.997459in}}%
\pgfpathlineto{\pgfqpoint{5.927768in}{4.044341in}}%
\pgfpathlineto{\pgfqpoint{5.926688in}{3.992660in}}%
\pgfpathlineto{\pgfqpoint{5.927783in}{4.043718in}}%
\pgfpathlineto{\pgfqpoint{5.928883in}{4.082399in}}%
\pgfpathlineto{\pgfqpoint{5.928182in}{4.043043in}}%
\pgfpathlineto{\pgfqpoint{5.928951in}{4.068366in}}%
\pgfpathlineto{\pgfqpoint{5.929877in}{4.049010in}}%
\pgfpathlineto{\pgfqpoint{5.929414in}{4.076013in}}%
\pgfpathlineto{\pgfqpoint{5.930071in}{4.057350in}}%
\pgfpathlineto{\pgfqpoint{5.931001in}{4.038358in}}%
\pgfpathlineto{\pgfqpoint{5.930344in}{4.068977in}}%
\pgfpathlineto{\pgfqpoint{5.931259in}{4.045247in}}%
\pgfpathlineto{\pgfqpoint{5.932174in}{4.081989in}}%
\pgfpathlineto{\pgfqpoint{5.931288in}{4.041298in}}%
\pgfpathlineto{\pgfqpoint{5.932413in}{4.077111in}}%
\pgfpathlineto{\pgfqpoint{5.933070in}{4.026647in}}%
\pgfpathlineto{\pgfqpoint{5.933611in}{4.059912in}}%
\pgfpathlineto{\pgfqpoint{5.934234in}{4.097469in}}%
\pgfpathlineto{\pgfqpoint{5.934774in}{4.087877in}}%
\pgfpathlineto{\pgfqpoint{5.935777in}{4.056454in}}%
\pgfpathlineto{\pgfqpoint{5.935057in}{4.093446in}}%
\pgfpathlineto{\pgfqpoint{5.935836in}{4.057217in}}%
\pgfpathlineto{\pgfqpoint{5.936176in}{4.044356in}}%
\pgfpathlineto{\pgfqpoint{5.936600in}{4.111438in}}%
\pgfpathlineto{\pgfqpoint{5.936892in}{4.106071in}}%
\pgfpathlineto{\pgfqpoint{5.937471in}{4.088338in}}%
\pgfpathlineto{\pgfqpoint{5.937170in}{4.119234in}}%
\pgfpathlineto{\pgfqpoint{5.937749in}{4.116769in}}%
\pgfpathlineto{\pgfqpoint{5.938752in}{4.188156in}}%
\pgfpathlineto{\pgfqpoint{5.937895in}{4.108683in}}%
\pgfpathlineto{\pgfqpoint{5.938951in}{4.166495in}}%
\pgfpathlineto{\pgfqpoint{5.939107in}{4.148470in}}%
\pgfpathlineto{\pgfqpoint{5.939959in}{4.175296in}}%
\pgfpathlineto{\pgfqpoint{5.940042in}{4.172511in}}%
\pgfpathlineto{\pgfqpoint{5.940821in}{4.200587in}}%
\pgfpathlineto{\pgfqpoint{5.941220in}{4.179746in}}%
\pgfpathlineto{\pgfqpoint{5.941653in}{4.153374in}}%
\pgfpathlineto{\pgfqpoint{5.942208in}{4.185600in}}%
\pgfpathlineto{\pgfqpoint{5.942301in}{4.184598in}}%
\pgfpathlineto{\pgfqpoint{5.942374in}{4.190835in}}%
\pgfpathlineto{\pgfqpoint{5.943158in}{4.141495in}}%
\pgfpathlineto{\pgfqpoint{5.943304in}{4.159086in}}%
\pgfpathlineto{\pgfqpoint{5.944472in}{4.117647in}}%
\pgfpathlineto{\pgfqpoint{5.944599in}{4.129185in}}%
\pgfpathlineto{\pgfqpoint{5.945509in}{4.093534in}}%
\pgfpathlineto{\pgfqpoint{5.945553in}{4.097235in}}%
\pgfpathlineto{\pgfqpoint{5.945865in}{4.076675in}}%
\pgfpathlineto{\pgfqpoint{5.946595in}{4.120283in}}%
\pgfpathlineto{\pgfqpoint{5.947715in}{4.143595in}}%
\pgfpathlineto{\pgfqpoint{5.947028in}{4.104986in}}%
\pgfpathlineto{\pgfqpoint{5.947754in}{4.143579in}}%
\pgfpathlineto{\pgfqpoint{5.948041in}{4.127352in}}%
\pgfpathlineto{\pgfqpoint{5.948649in}{4.170841in}}%
\pgfpathlineto{\pgfqpoint{5.948732in}{4.166934in}}%
\pgfpathlineto{\pgfqpoint{5.949711in}{4.195237in}}%
\pgfpathlineto{\pgfqpoint{5.949219in}{4.161216in}}%
\pgfpathlineto{\pgfqpoint{5.949862in}{4.176821in}}%
\pgfpathlineto{\pgfqpoint{5.950436in}{4.157202in}}%
\pgfpathlineto{\pgfqpoint{5.950013in}{4.185633in}}%
\pgfpathlineto{\pgfqpoint{5.951059in}{4.166524in}}%
\pgfpathlineto{\pgfqpoint{5.951196in}{4.182457in}}%
\pgfpathlineto{\pgfqpoint{5.951595in}{4.147806in}}%
\pgfpathlineto{\pgfqpoint{5.952130in}{4.154588in}}%
\pgfpathlineto{\pgfqpoint{5.952311in}{4.162135in}}%
\pgfpathlineto{\pgfqpoint{5.952418in}{4.138963in}}%
\pgfpathlineto{\pgfqpoint{5.952608in}{4.143593in}}%
\pgfpathlineto{\pgfqpoint{5.952953in}{4.103284in}}%
\pgfpathlineto{\pgfqpoint{5.953615in}{4.146020in}}%
\pgfpathlineto{\pgfqpoint{5.953722in}{4.140779in}}%
\pgfpathlineto{\pgfqpoint{5.953776in}{4.151082in}}%
\pgfpathlineto{\pgfqpoint{5.954350in}{4.104168in}}%
\pgfpathlineto{\pgfqpoint{5.954686in}{4.120954in}}%
\pgfpathlineto{\pgfqpoint{5.954691in}{4.120264in}}%
\pgfpathlineto{\pgfqpoint{5.955495in}{4.151838in}}%
\pgfpathlineto{\pgfqpoint{5.955753in}{4.126369in}}%
\pgfpathlineto{\pgfqpoint{5.956351in}{4.145449in}}%
\pgfpathlineto{\pgfqpoint{5.955845in}{4.115507in}}%
\pgfpathlineto{\pgfqpoint{5.956833in}{4.118111in}}%
\pgfpathlineto{\pgfqpoint{5.958294in}{4.166671in}}%
\pgfpathlineto{\pgfqpoint{5.957087in}{4.105103in}}%
\pgfpathlineto{\pgfqpoint{5.958396in}{4.161806in}}%
\pgfpathlineto{\pgfqpoint{5.959331in}{4.145998in}}%
\pgfpathlineto{\pgfqpoint{5.959496in}{4.167039in}}%
\pgfpathlineto{\pgfqpoint{5.959501in}{4.166296in}}%
\pgfpathlineto{\pgfqpoint{5.959842in}{4.215795in}}%
\pgfpathlineto{\pgfqpoint{5.960519in}{4.179522in}}%
\pgfpathlineto{\pgfqpoint{5.961497in}{4.221090in}}%
\pgfpathlineto{\pgfqpoint{5.960572in}{4.165772in}}%
\pgfpathlineto{\pgfqpoint{5.961721in}{4.199832in}}%
\pgfpathlineto{\pgfqpoint{5.962160in}{4.161718in}}%
\pgfpathlineto{\pgfqpoint{5.962841in}{4.186286in}}%
\pgfpathlineto{\pgfqpoint{5.962939in}{4.190307in}}%
\pgfpathlineto{\pgfqpoint{5.963416in}{4.158321in}}%
\pgfpathlineto{\pgfqpoint{5.963859in}{4.176652in}}%
\pgfpathlineto{\pgfqpoint{5.964754in}{4.154190in}}%
\pgfpathlineto{\pgfqpoint{5.964341in}{4.187325in}}%
\pgfpathlineto{\pgfqpoint{5.964993in}{4.165143in}}%
\pgfpathlineto{\pgfqpoint{5.965081in}{4.175291in}}%
\pgfpathlineto{\pgfqpoint{5.966015in}{4.150666in}}%
\pgfpathlineto{\pgfqpoint{5.966093in}{4.156054in}}%
\pgfpathlineto{\pgfqpoint{5.966799in}{4.167560in}}%
\pgfpathlineto{\pgfqpoint{5.966556in}{4.138223in}}%
\pgfpathlineto{\pgfqpoint{5.967194in}{4.154313in}}%
\pgfpathlineto{\pgfqpoint{5.967383in}{4.145048in}}%
\pgfpathlineto{\pgfqpoint{5.968211in}{4.197025in}}%
\pgfpathlineto{\pgfqpoint{5.968737in}{4.225427in}}%
\pgfpathlineto{\pgfqpoint{5.968469in}{4.189427in}}%
\pgfpathlineto{\pgfqpoint{5.969350in}{4.212466in}}%
\pgfpathlineto{\pgfqpoint{5.970022in}{4.207344in}}%
\pgfpathlineto{\pgfqpoint{5.969516in}{4.231317in}}%
\pgfpathlineto{\pgfqpoint{5.970446in}{4.219609in}}%
\pgfpathlineto{\pgfqpoint{5.970972in}{4.206443in}}%
\pgfpathlineto{\pgfqpoint{5.970602in}{4.238940in}}%
\pgfpathlineto{\pgfqpoint{5.971317in}{4.232358in}}%
\pgfpathlineto{\pgfqpoint{5.972461in}{4.284940in}}%
\pgfpathlineto{\pgfqpoint{5.971463in}{4.230525in}}%
\pgfpathlineto{\pgfqpoint{5.972529in}{4.282712in}}%
\pgfpathlineto{\pgfqpoint{5.973007in}{4.287667in}}%
\pgfpathlineto{\pgfqpoint{5.972885in}{4.263859in}}%
\pgfpathlineto{\pgfqpoint{5.973542in}{4.265732in}}%
\pgfpathlineto{\pgfqpoint{5.973878in}{4.270808in}}%
\pgfpathlineto{\pgfqpoint{5.974715in}{4.243805in}}%
\pgfpathlineto{\pgfqpoint{5.975436in}{4.251011in}}%
\pgfpathlineto{\pgfqpoint{5.975733in}{4.223744in}}%
\pgfpathlineto{\pgfqpoint{5.976785in}{4.204870in}}%
\pgfpathlineto{\pgfqpoint{5.975821in}{4.235350in}}%
\pgfpathlineto{\pgfqpoint{5.976872in}{4.220129in}}%
\pgfpathlineto{\pgfqpoint{5.977198in}{4.253506in}}%
\pgfpathlineto{\pgfqpoint{5.978031in}{4.227926in}}%
\pgfpathlineto{\pgfqpoint{5.978805in}{4.222769in}}%
\pgfpathlineto{\pgfqpoint{5.978469in}{4.268793in}}%
\pgfpathlineto{\pgfqpoint{5.979112in}{4.240004in}}%
\pgfpathlineto{\pgfqpoint{5.979326in}{4.261357in}}%
\pgfpathlineto{\pgfqpoint{5.979131in}{4.238070in}}%
\pgfpathlineto{\pgfqpoint{5.979331in}{4.259538in}}%
\pgfpathlineto{\pgfqpoint{5.979375in}{4.267290in}}%
\pgfpathlineto{\pgfqpoint{5.979842in}{4.239294in}}%
\pgfpathlineto{\pgfqpoint{5.980441in}{4.259693in}}%
\pgfpathlineto{\pgfqpoint{5.980465in}{4.264493in}}%
\pgfpathlineto{\pgfqpoint{5.980932in}{4.238346in}}%
\pgfpathlineto{\pgfqpoint{5.981259in}{4.238466in}}%
\pgfpathlineto{\pgfqpoint{5.981755in}{4.218518in}}%
\pgfpathlineto{\pgfqpoint{5.982325in}{4.244784in}}%
\pgfpathlineto{\pgfqpoint{5.982374in}{4.234544in}}%
\pgfpathlineto{\pgfqpoint{5.982748in}{4.238943in}}%
\pgfpathlineto{\pgfqpoint{5.983401in}{4.207457in}}%
\pgfpathlineto{\pgfqpoint{5.983406in}{4.208479in}}%
\pgfpathlineto{\pgfqpoint{5.983873in}{4.194988in}}%
\pgfpathlineto{\pgfqpoint{5.983581in}{4.222456in}}%
\pgfpathlineto{\pgfqpoint{5.984506in}{4.209212in}}%
\pgfpathlineto{\pgfqpoint{5.985884in}{4.263615in}}%
\pgfpathlineto{\pgfqpoint{5.984788in}{4.200919in}}%
\pgfpathlineto{\pgfqpoint{5.985913in}{4.258130in}}%
\pgfpathlineto{\pgfqpoint{5.986901in}{4.229222in}}%
\pgfpathlineto{\pgfqpoint{5.986015in}{4.258650in}}%
\pgfpathlineto{\pgfqpoint{5.987062in}{4.241109in}}%
\pgfpathlineto{\pgfqpoint{5.987189in}{4.253558in}}%
\pgfpathlineto{\pgfqpoint{5.987607in}{4.214286in}}%
\pgfpathlineto{\pgfqpoint{5.988138in}{4.230326in}}%
\pgfpathlineto{\pgfqpoint{5.988187in}{4.222315in}}%
\pgfpathlineto{\pgfqpoint{5.989160in}{4.245746in}}%
\pgfpathlineto{\pgfqpoint{5.989233in}{4.235256in}}%
\pgfpathlineto{\pgfqpoint{5.989301in}{4.231117in}}%
\pgfpathlineto{\pgfqpoint{5.989881in}{4.265686in}}%
\pgfpathlineto{\pgfqpoint{5.990007in}{4.258435in}}%
\pgfpathlineto{\pgfqpoint{5.990353in}{4.277000in}}%
\pgfpathlineto{\pgfqpoint{5.990650in}{4.251710in}}%
\pgfpathlineto{\pgfqpoint{5.991113in}{4.255652in}}%
\pgfpathlineto{\pgfqpoint{5.991117in}{4.254800in}}%
\pgfpathlineto{\pgfqpoint{5.991702in}{4.304548in}}%
\pgfpathlineto{\pgfqpoint{5.991965in}{4.300290in}}%
\pgfpathlineto{\pgfqpoint{5.992573in}{4.343261in}}%
\pgfpathlineto{\pgfqpoint{5.992008in}{4.299104in}}%
\pgfpathlineto{\pgfqpoint{5.993065in}{4.300733in}}%
\pgfpathlineto{\pgfqpoint{5.993527in}{4.315584in}}%
\pgfpathlineto{\pgfqpoint{5.994189in}{4.286768in}}%
\pgfpathlineto{\pgfqpoint{5.994608in}{4.312274in}}%
\pgfpathlineto{\pgfqpoint{5.994871in}{4.279787in}}%
\pgfpathlineto{\pgfqpoint{5.995304in}{4.293141in}}%
\pgfpathlineto{\pgfqpoint{5.995689in}{4.286075in}}%
\pgfpathlineto{\pgfqpoint{5.996244in}{4.320456in}}%
\pgfpathlineto{\pgfqpoint{5.996385in}{4.302118in}}%
\pgfpathlineto{\pgfqpoint{5.997242in}{4.329439in}}%
\pgfpathlineto{\pgfqpoint{5.996463in}{4.298817in}}%
\pgfpathlineto{\pgfqpoint{5.997534in}{4.309423in}}%
\pgfpathlineto{\pgfqpoint{5.998595in}{4.286452in}}%
\pgfpathlineto{\pgfqpoint{5.997656in}{4.319240in}}%
\pgfpathlineto{\pgfqpoint{5.998790in}{4.292569in}}%
\pgfpathlineto{\pgfqpoint{5.999340in}{4.318493in}}%
\pgfpathlineto{\pgfqpoint{5.999696in}{4.272233in}}%
\pgfpathlineto{\pgfqpoint{5.999876in}{4.281808in}}%
\pgfpathlineto{\pgfqpoint{6.001005in}{4.313747in}}%
\pgfpathlineto{\pgfqpoint{6.000114in}{4.262973in}}%
\pgfpathlineto{\pgfqpoint{6.001054in}{4.304982in}}%
\pgfpathlineto{\pgfqpoint{6.002768in}{4.259089in}}%
\pgfpathlineto{\pgfqpoint{6.002792in}{4.269682in}}%
\pgfpathlineto{\pgfqpoint{6.002821in}{4.276024in}}%
\pgfpathlineto{\pgfqpoint{6.003819in}{4.243720in}}%
\pgfpathlineto{\pgfqpoint{6.003878in}{4.254329in}}%
\pgfpathlineto{\pgfqpoint{6.004871in}{4.210882in}}%
\pgfpathlineto{\pgfqpoint{6.003951in}{4.260605in}}%
\pgfpathlineto{\pgfqpoint{6.005022in}{4.222976in}}%
\pgfpathlineto{\pgfqpoint{6.006312in}{4.268293in}}%
\pgfpathlineto{\pgfqpoint{6.005090in}{4.219622in}}%
\pgfpathlineto{\pgfqpoint{6.006575in}{4.251566in}}%
\pgfusepath{stroke}%
\end{pgfscope}%
\begin{pgfscope}%
\pgfsetrectcap%
\pgfsetmiterjoin%
\pgfsetlinewidth{0.803000pt}%
\definecolor{currentstroke}{rgb}{0.000000,0.000000,0.000000}%
\pgfsetstrokecolor{currentstroke}%
\pgfsetdash{}{0pt}%
\pgfpathmoveto{\pgfqpoint{0.894648in}{0.771190in}}%
\pgfpathlineto{\pgfqpoint{0.894648in}{4.601775in}}%
\pgfusepath{stroke}%
\end{pgfscope}%
\begin{pgfscope}%
\pgfsetrectcap%
\pgfsetmiterjoin%
\pgfsetlinewidth{0.803000pt}%
\definecolor{currentstroke}{rgb}{0.000000,0.000000,0.000000}%
\pgfsetstrokecolor{currentstroke}%
\pgfsetdash{}{0pt}%
\pgfpathmoveto{\pgfqpoint{6.250000in}{0.771190in}}%
\pgfpathlineto{\pgfqpoint{6.250000in}{4.601775in}}%
\pgfusepath{stroke}%
\end{pgfscope}%
\begin{pgfscope}%
\pgfsetrectcap%
\pgfsetmiterjoin%
\pgfsetlinewidth{0.803000pt}%
\definecolor{currentstroke}{rgb}{0.000000,0.000000,0.000000}%
\pgfsetstrokecolor{currentstroke}%
\pgfsetdash{}{0pt}%
\pgfpathmoveto{\pgfqpoint{0.894648in}{0.771190in}}%
\pgfpathlineto{\pgfqpoint{6.250000in}{0.771190in}}%
\pgfusepath{stroke}%
\end{pgfscope}%
\begin{pgfscope}%
\pgfsetrectcap%
\pgfsetmiterjoin%
\pgfsetlinewidth{0.803000pt}%
\definecolor{currentstroke}{rgb}{0.000000,0.000000,0.000000}%
\pgfsetstrokecolor{currentstroke}%
\pgfsetdash{}{0pt}%
\pgfpathmoveto{\pgfqpoint{0.894648in}{4.601775in}}%
\pgfpathlineto{\pgfqpoint{6.250000in}{4.601775in}}%
\pgfusepath{stroke}%
\end{pgfscope}%
\end{pgfpicture}%
\makeatother%
\endgroup%
}
    \caption{Plotted shape of potential and force $\Delta U = 0.1 k_B T$}
    \label{fig:No_flashing}
\end{figure}

\begin{figure}[H]
    \centering
    \resizebox{\columnwidth}{!}{%% Creator: Matplotlib, PGF backend
%%
%% To include the figure in your LaTeX document, write
%%   \input{<filename>.pgf}
%%
%% Make sure the required packages are loaded in your preamble
%%   \usepackage{pgf}
%%
%% Also ensure that all the required font packages are loaded; for instance,
%% the lmodern package is sometimes necessary when using math font.
%%   \usepackage{lmodern}
%%
%% Figures using additional raster images can only be included by \input if
%% they are in the same directory as the main LaTeX file. For loading figures
%% from other directories you can use the `import` package
%%   \usepackage{import}
%%
%% and then include the figures with
%%   \import{<path to file>}{<filename>.pgf}
%%
%% Matplotlib used the following preamble
%%   \def\mathdefault#1{#1}
%%   \everymath=\expandafter{\the\everymath\displaystyle}
%%   \IfFileExists{scrextend.sty}{
%%     \usepackage[fontsize=10.000000pt]{scrextend}
%%   }{
%%     \renewcommand{\normalsize}{\fontsize{10.000000}{12.000000}\selectfont}
%%     \normalsize
%%   }
%%   
%%   \makeatletter\@ifpackageloaded{underscore}{}{\usepackage[strings]{underscore}}\makeatother
%%
\begingroup%
\makeatletter%
\begin{pgfpicture}%
\pgfpathrectangle{\pgfpointorigin}{\pgfqpoint{6.400000in}{4.800000in}}%
\pgfusepath{use as bounding box, clip}%
\begin{pgfscope}%
\pgfsetbuttcap%
\pgfsetmiterjoin%
\definecolor{currentfill}{rgb}{1.000000,1.000000,1.000000}%
\pgfsetfillcolor{currentfill}%
\pgfsetlinewidth{0.000000pt}%
\definecolor{currentstroke}{rgb}{1.000000,1.000000,1.000000}%
\pgfsetstrokecolor{currentstroke}%
\pgfsetdash{}{0pt}%
\pgfpathmoveto{\pgfqpoint{0.000000in}{0.000000in}}%
\pgfpathlineto{\pgfqpoint{6.400000in}{0.000000in}}%
\pgfpathlineto{\pgfqpoint{6.400000in}{4.800000in}}%
\pgfpathlineto{\pgfqpoint{0.000000in}{4.800000in}}%
\pgfpathlineto{\pgfqpoint{0.000000in}{0.000000in}}%
\pgfpathclose%
\pgfusepath{fill}%
\end{pgfscope}%
\begin{pgfscope}%
\pgfsetbuttcap%
\pgfsetmiterjoin%
\definecolor{currentfill}{rgb}{1.000000,1.000000,1.000000}%
\pgfsetfillcolor{currentfill}%
\pgfsetlinewidth{0.000000pt}%
\definecolor{currentstroke}{rgb}{0.000000,0.000000,0.000000}%
\pgfsetstrokecolor{currentstroke}%
\pgfsetstrokeopacity{0.000000}%
\pgfsetdash{}{0pt}%
\pgfpathmoveto{\pgfqpoint{0.894648in}{0.771190in}}%
\pgfpathlineto{\pgfqpoint{6.250000in}{0.771190in}}%
\pgfpathlineto{\pgfqpoint{6.250000in}{4.601775in}}%
\pgfpathlineto{\pgfqpoint{0.894648in}{4.601775in}}%
\pgfpathlineto{\pgfqpoint{0.894648in}{0.771190in}}%
\pgfpathclose%
\pgfusepath{fill}%
\end{pgfscope}%
\begin{pgfscope}%
\pgfsetbuttcap%
\pgfsetroundjoin%
\definecolor{currentfill}{rgb}{0.000000,0.000000,0.000000}%
\pgfsetfillcolor{currentfill}%
\pgfsetlinewidth{0.803000pt}%
\definecolor{currentstroke}{rgb}{0.000000,0.000000,0.000000}%
\pgfsetstrokecolor{currentstroke}%
\pgfsetdash{}{0pt}%
\pgfsys@defobject{currentmarker}{\pgfqpoint{0.000000in}{-0.048611in}}{\pgfqpoint{0.000000in}{0.000000in}}{%
\pgfpathmoveto{\pgfqpoint{0.000000in}{0.000000in}}%
\pgfpathlineto{\pgfqpoint{0.000000in}{-0.048611in}}%
\pgfusepath{stroke,fill}%
}%
\begin{pgfscope}%
\pgfsys@transformshift{1.138073in}{0.771190in}%
\pgfsys@useobject{currentmarker}{}%
\end{pgfscope}%
\end{pgfscope}%
\begin{pgfscope}%
\definecolor{textcolor}{rgb}{0.000000,0.000000,0.000000}%
\pgfsetstrokecolor{textcolor}%
\pgfsetfillcolor{textcolor}%
\pgftext[x=1.138073in,y=0.673968in,,top]{\color{textcolor}{\rmfamily\fontsize{10.000000}{12.000000}\selectfont\catcode`\^=\active\def^{\ifmmode\sp\else\^{}\fi}\catcode`\%=\active\def%{\%}$\mathdefault{0}$}}%
\end{pgfscope}%
\begin{pgfscope}%
\pgfsetbuttcap%
\pgfsetroundjoin%
\definecolor{currentfill}{rgb}{0.000000,0.000000,0.000000}%
\pgfsetfillcolor{currentfill}%
\pgfsetlinewidth{0.803000pt}%
\definecolor{currentstroke}{rgb}{0.000000,0.000000,0.000000}%
\pgfsetstrokecolor{currentstroke}%
\pgfsetdash{}{0pt}%
\pgfsys@defobject{currentmarker}{\pgfqpoint{0.000000in}{-0.048611in}}{\pgfqpoint{0.000000in}{0.000000in}}{%
\pgfpathmoveto{\pgfqpoint{0.000000in}{0.000000in}}%
\pgfpathlineto{\pgfqpoint{0.000000in}{-0.048611in}}%
\pgfusepath{stroke,fill}%
}%
\begin{pgfscope}%
\pgfsys@transformshift{2.111773in}{0.771190in}%
\pgfsys@useobject{currentmarker}{}%
\end{pgfscope}%
\end{pgfscope}%
\begin{pgfscope}%
\definecolor{textcolor}{rgb}{0.000000,0.000000,0.000000}%
\pgfsetstrokecolor{textcolor}%
\pgfsetfillcolor{textcolor}%
\pgftext[x=2.111773in,y=0.673968in,,top]{\color{textcolor}{\rmfamily\fontsize{10.000000}{12.000000}\selectfont\catcode`\^=\active\def^{\ifmmode\sp\else\^{}\fi}\catcode`\%=\active\def%{\%}$\mathdefault{20}$}}%
\end{pgfscope}%
\begin{pgfscope}%
\pgfsetbuttcap%
\pgfsetroundjoin%
\definecolor{currentfill}{rgb}{0.000000,0.000000,0.000000}%
\pgfsetfillcolor{currentfill}%
\pgfsetlinewidth{0.803000pt}%
\definecolor{currentstroke}{rgb}{0.000000,0.000000,0.000000}%
\pgfsetstrokecolor{currentstroke}%
\pgfsetdash{}{0pt}%
\pgfsys@defobject{currentmarker}{\pgfqpoint{0.000000in}{-0.048611in}}{\pgfqpoint{0.000000in}{0.000000in}}{%
\pgfpathmoveto{\pgfqpoint{0.000000in}{0.000000in}}%
\pgfpathlineto{\pgfqpoint{0.000000in}{-0.048611in}}%
\pgfusepath{stroke,fill}%
}%
\begin{pgfscope}%
\pgfsys@transformshift{3.085474in}{0.771190in}%
\pgfsys@useobject{currentmarker}{}%
\end{pgfscope}%
\end{pgfscope}%
\begin{pgfscope}%
\definecolor{textcolor}{rgb}{0.000000,0.000000,0.000000}%
\pgfsetstrokecolor{textcolor}%
\pgfsetfillcolor{textcolor}%
\pgftext[x=3.085474in,y=0.673968in,,top]{\color{textcolor}{\rmfamily\fontsize{10.000000}{12.000000}\selectfont\catcode`\^=\active\def^{\ifmmode\sp\else\^{}\fi}\catcode`\%=\active\def%{\%}$\mathdefault{40}$}}%
\end{pgfscope}%
\begin{pgfscope}%
\pgfsetbuttcap%
\pgfsetroundjoin%
\definecolor{currentfill}{rgb}{0.000000,0.000000,0.000000}%
\pgfsetfillcolor{currentfill}%
\pgfsetlinewidth{0.803000pt}%
\definecolor{currentstroke}{rgb}{0.000000,0.000000,0.000000}%
\pgfsetstrokecolor{currentstroke}%
\pgfsetdash{}{0pt}%
\pgfsys@defobject{currentmarker}{\pgfqpoint{0.000000in}{-0.048611in}}{\pgfqpoint{0.000000in}{0.000000in}}{%
\pgfpathmoveto{\pgfqpoint{0.000000in}{0.000000in}}%
\pgfpathlineto{\pgfqpoint{0.000000in}{-0.048611in}}%
\pgfusepath{stroke,fill}%
}%
\begin{pgfscope}%
\pgfsys@transformshift{4.059174in}{0.771190in}%
\pgfsys@useobject{currentmarker}{}%
\end{pgfscope}%
\end{pgfscope}%
\begin{pgfscope}%
\definecolor{textcolor}{rgb}{0.000000,0.000000,0.000000}%
\pgfsetstrokecolor{textcolor}%
\pgfsetfillcolor{textcolor}%
\pgftext[x=4.059174in,y=0.673968in,,top]{\color{textcolor}{\rmfamily\fontsize{10.000000}{12.000000}\selectfont\catcode`\^=\active\def^{\ifmmode\sp\else\^{}\fi}\catcode`\%=\active\def%{\%}$\mathdefault{60}$}}%
\end{pgfscope}%
\begin{pgfscope}%
\pgfsetbuttcap%
\pgfsetroundjoin%
\definecolor{currentfill}{rgb}{0.000000,0.000000,0.000000}%
\pgfsetfillcolor{currentfill}%
\pgfsetlinewidth{0.803000pt}%
\definecolor{currentstroke}{rgb}{0.000000,0.000000,0.000000}%
\pgfsetstrokecolor{currentstroke}%
\pgfsetdash{}{0pt}%
\pgfsys@defobject{currentmarker}{\pgfqpoint{0.000000in}{-0.048611in}}{\pgfqpoint{0.000000in}{0.000000in}}{%
\pgfpathmoveto{\pgfqpoint{0.000000in}{0.000000in}}%
\pgfpathlineto{\pgfqpoint{0.000000in}{-0.048611in}}%
\pgfusepath{stroke,fill}%
}%
\begin{pgfscope}%
\pgfsys@transformshift{5.032874in}{0.771190in}%
\pgfsys@useobject{currentmarker}{}%
\end{pgfscope}%
\end{pgfscope}%
\begin{pgfscope}%
\definecolor{textcolor}{rgb}{0.000000,0.000000,0.000000}%
\pgfsetstrokecolor{textcolor}%
\pgfsetfillcolor{textcolor}%
\pgftext[x=5.032874in,y=0.673968in,,top]{\color{textcolor}{\rmfamily\fontsize{10.000000}{12.000000}\selectfont\catcode`\^=\active\def^{\ifmmode\sp\else\^{}\fi}\catcode`\%=\active\def%{\%}$\mathdefault{80}$}}%
\end{pgfscope}%
\begin{pgfscope}%
\pgfsetbuttcap%
\pgfsetroundjoin%
\definecolor{currentfill}{rgb}{0.000000,0.000000,0.000000}%
\pgfsetfillcolor{currentfill}%
\pgfsetlinewidth{0.803000pt}%
\definecolor{currentstroke}{rgb}{0.000000,0.000000,0.000000}%
\pgfsetstrokecolor{currentstroke}%
\pgfsetdash{}{0pt}%
\pgfsys@defobject{currentmarker}{\pgfqpoint{0.000000in}{-0.048611in}}{\pgfqpoint{0.000000in}{0.000000in}}{%
\pgfpathmoveto{\pgfqpoint{0.000000in}{0.000000in}}%
\pgfpathlineto{\pgfqpoint{0.000000in}{-0.048611in}}%
\pgfusepath{stroke,fill}%
}%
\begin{pgfscope}%
\pgfsys@transformshift{6.006575in}{0.771190in}%
\pgfsys@useobject{currentmarker}{}%
\end{pgfscope}%
\end{pgfscope}%
\begin{pgfscope}%
\definecolor{textcolor}{rgb}{0.000000,0.000000,0.000000}%
\pgfsetstrokecolor{textcolor}%
\pgfsetfillcolor{textcolor}%
\pgftext[x=6.006575in,y=0.673968in,,top]{\color{textcolor}{\rmfamily\fontsize{10.000000}{12.000000}\selectfont\catcode`\^=\active\def^{\ifmmode\sp\else\^{}\fi}\catcode`\%=\active\def%{\%}$\mathdefault{100}$}}%
\end{pgfscope}%
\begin{pgfscope}%
\definecolor{textcolor}{rgb}{0.000000,0.000000,0.000000}%
\pgfsetstrokecolor{textcolor}%
\pgfsetfillcolor{textcolor}%
\pgftext[x=3.572324in,y=0.494955in,,top]{\color{textcolor}{\rmfamily\fontsize{24.000000}{28.800000}\selectfont\catcode`\^=\active\def^{\ifmmode\sp\else\^{}\fi}\catcode`\%=\active\def%{\%}$t (s)$}}%
\end{pgfscope}%
\begin{pgfscope}%
\pgfsetbuttcap%
\pgfsetroundjoin%
\definecolor{currentfill}{rgb}{0.000000,0.000000,0.000000}%
\pgfsetfillcolor{currentfill}%
\pgfsetlinewidth{0.803000pt}%
\definecolor{currentstroke}{rgb}{0.000000,0.000000,0.000000}%
\pgfsetstrokecolor{currentstroke}%
\pgfsetdash{}{0pt}%
\pgfsys@defobject{currentmarker}{\pgfqpoint{-0.048611in}{0.000000in}}{\pgfqpoint{-0.000000in}{0.000000in}}{%
\pgfpathmoveto{\pgfqpoint{-0.000000in}{0.000000in}}%
\pgfpathlineto{\pgfqpoint{-0.048611in}{0.000000in}}%
\pgfusepath{stroke,fill}%
}%
\begin{pgfscope}%
\pgfsys@transformshift{0.894648in}{0.771190in}%
\pgfsys@useobject{currentmarker}{}%
\end{pgfscope}%
\end{pgfscope}%
\begin{pgfscope}%
\definecolor{textcolor}{rgb}{0.000000,0.000000,0.000000}%
\pgfsetstrokecolor{textcolor}%
\pgfsetfillcolor{textcolor}%
\pgftext[x=0.550511in, y=0.722965in, left, base]{\color{textcolor}{\rmfamily\fontsize{10.000000}{12.000000}\selectfont\catcode`\^=\active\def^{\ifmmode\sp\else\^{}\fi}\catcode`\%=\active\def%{\%}$\mathdefault{\ensuremath{-}40}$}}%
\end{pgfscope}%
\begin{pgfscope}%
\pgfsetbuttcap%
\pgfsetroundjoin%
\definecolor{currentfill}{rgb}{0.000000,0.000000,0.000000}%
\pgfsetfillcolor{currentfill}%
\pgfsetlinewidth{0.803000pt}%
\definecolor{currentstroke}{rgb}{0.000000,0.000000,0.000000}%
\pgfsetstrokecolor{currentstroke}%
\pgfsetdash{}{0pt}%
\pgfsys@defobject{currentmarker}{\pgfqpoint{-0.048611in}{0.000000in}}{\pgfqpoint{-0.000000in}{0.000000in}}{%
\pgfpathmoveto{\pgfqpoint{-0.000000in}{0.000000in}}%
\pgfpathlineto{\pgfqpoint{-0.048611in}{0.000000in}}%
\pgfusepath{stroke,fill}%
}%
\begin{pgfscope}%
\pgfsys@transformshift{0.894648in}{1.537307in}%
\pgfsys@useobject{currentmarker}{}%
\end{pgfscope}%
\end{pgfscope}%
\begin{pgfscope}%
\definecolor{textcolor}{rgb}{0.000000,0.000000,0.000000}%
\pgfsetstrokecolor{textcolor}%
\pgfsetfillcolor{textcolor}%
\pgftext[x=0.550511in, y=1.489082in, left, base]{\color{textcolor}{\rmfamily\fontsize{10.000000}{12.000000}\selectfont\catcode`\^=\active\def^{\ifmmode\sp\else\^{}\fi}\catcode`\%=\active\def%{\%}$\mathdefault{\ensuremath{-}20}$}}%
\end{pgfscope}%
\begin{pgfscope}%
\pgfsetbuttcap%
\pgfsetroundjoin%
\definecolor{currentfill}{rgb}{0.000000,0.000000,0.000000}%
\pgfsetfillcolor{currentfill}%
\pgfsetlinewidth{0.803000pt}%
\definecolor{currentstroke}{rgb}{0.000000,0.000000,0.000000}%
\pgfsetstrokecolor{currentstroke}%
\pgfsetdash{}{0pt}%
\pgfsys@defobject{currentmarker}{\pgfqpoint{-0.048611in}{0.000000in}}{\pgfqpoint{-0.000000in}{0.000000in}}{%
\pgfpathmoveto{\pgfqpoint{-0.000000in}{0.000000in}}%
\pgfpathlineto{\pgfqpoint{-0.048611in}{0.000000in}}%
\pgfusepath{stroke,fill}%
}%
\begin{pgfscope}%
\pgfsys@transformshift{0.894648in}{2.303424in}%
\pgfsys@useobject{currentmarker}{}%
\end{pgfscope}%
\end{pgfscope}%
\begin{pgfscope}%
\definecolor{textcolor}{rgb}{0.000000,0.000000,0.000000}%
\pgfsetstrokecolor{textcolor}%
\pgfsetfillcolor{textcolor}%
\pgftext[x=0.727981in, y=2.255199in, left, base]{\color{textcolor}{\rmfamily\fontsize{10.000000}{12.000000}\selectfont\catcode`\^=\active\def^{\ifmmode\sp\else\^{}\fi}\catcode`\%=\active\def%{\%}$\mathdefault{0}$}}%
\end{pgfscope}%
\begin{pgfscope}%
\pgfsetbuttcap%
\pgfsetroundjoin%
\definecolor{currentfill}{rgb}{0.000000,0.000000,0.000000}%
\pgfsetfillcolor{currentfill}%
\pgfsetlinewidth{0.803000pt}%
\definecolor{currentstroke}{rgb}{0.000000,0.000000,0.000000}%
\pgfsetstrokecolor{currentstroke}%
\pgfsetdash{}{0pt}%
\pgfsys@defobject{currentmarker}{\pgfqpoint{-0.048611in}{0.000000in}}{\pgfqpoint{-0.000000in}{0.000000in}}{%
\pgfpathmoveto{\pgfqpoint{-0.000000in}{0.000000in}}%
\pgfpathlineto{\pgfqpoint{-0.048611in}{0.000000in}}%
\pgfusepath{stroke,fill}%
}%
\begin{pgfscope}%
\pgfsys@transformshift{0.894648in}{3.069541in}%
\pgfsys@useobject{currentmarker}{}%
\end{pgfscope}%
\end{pgfscope}%
\begin{pgfscope}%
\definecolor{textcolor}{rgb}{0.000000,0.000000,0.000000}%
\pgfsetstrokecolor{textcolor}%
\pgfsetfillcolor{textcolor}%
\pgftext[x=0.658536in, y=3.021316in, left, base]{\color{textcolor}{\rmfamily\fontsize{10.000000}{12.000000}\selectfont\catcode`\^=\active\def^{\ifmmode\sp\else\^{}\fi}\catcode`\%=\active\def%{\%}$\mathdefault{20}$}}%
\end{pgfscope}%
\begin{pgfscope}%
\pgfsetbuttcap%
\pgfsetroundjoin%
\definecolor{currentfill}{rgb}{0.000000,0.000000,0.000000}%
\pgfsetfillcolor{currentfill}%
\pgfsetlinewidth{0.803000pt}%
\definecolor{currentstroke}{rgb}{0.000000,0.000000,0.000000}%
\pgfsetstrokecolor{currentstroke}%
\pgfsetdash{}{0pt}%
\pgfsys@defobject{currentmarker}{\pgfqpoint{-0.048611in}{0.000000in}}{\pgfqpoint{-0.000000in}{0.000000in}}{%
\pgfpathmoveto{\pgfqpoint{-0.000000in}{0.000000in}}%
\pgfpathlineto{\pgfqpoint{-0.048611in}{0.000000in}}%
\pgfusepath{stroke,fill}%
}%
\begin{pgfscope}%
\pgfsys@transformshift{0.894648in}{3.835658in}%
\pgfsys@useobject{currentmarker}{}%
\end{pgfscope}%
\end{pgfscope}%
\begin{pgfscope}%
\definecolor{textcolor}{rgb}{0.000000,0.000000,0.000000}%
\pgfsetstrokecolor{textcolor}%
\pgfsetfillcolor{textcolor}%
\pgftext[x=0.658536in, y=3.787432in, left, base]{\color{textcolor}{\rmfamily\fontsize{10.000000}{12.000000}\selectfont\catcode`\^=\active\def^{\ifmmode\sp\else\^{}\fi}\catcode`\%=\active\def%{\%}$\mathdefault{40}$}}%
\end{pgfscope}%
\begin{pgfscope}%
\pgfsetbuttcap%
\pgfsetroundjoin%
\definecolor{currentfill}{rgb}{0.000000,0.000000,0.000000}%
\pgfsetfillcolor{currentfill}%
\pgfsetlinewidth{0.803000pt}%
\definecolor{currentstroke}{rgb}{0.000000,0.000000,0.000000}%
\pgfsetstrokecolor{currentstroke}%
\pgfsetdash{}{0pt}%
\pgfsys@defobject{currentmarker}{\pgfqpoint{-0.048611in}{0.000000in}}{\pgfqpoint{-0.000000in}{0.000000in}}{%
\pgfpathmoveto{\pgfqpoint{-0.000000in}{0.000000in}}%
\pgfpathlineto{\pgfqpoint{-0.048611in}{0.000000in}}%
\pgfusepath{stroke,fill}%
}%
\begin{pgfscope}%
\pgfsys@transformshift{0.894648in}{4.601775in}%
\pgfsys@useobject{currentmarker}{}%
\end{pgfscope}%
\end{pgfscope}%
\begin{pgfscope}%
\definecolor{textcolor}{rgb}{0.000000,0.000000,0.000000}%
\pgfsetstrokecolor{textcolor}%
\pgfsetfillcolor{textcolor}%
\pgftext[x=0.658536in, y=4.553549in, left, base]{\color{textcolor}{\rmfamily\fontsize{10.000000}{12.000000}\selectfont\catcode`\^=\active\def^{\ifmmode\sp\else\^{}\fi}\catcode`\%=\active\def%{\%}$\mathdefault{60}$}}%
\end{pgfscope}%
\begin{pgfscope}%
\definecolor{textcolor}{rgb}{0.000000,0.000000,0.000000}%
\pgfsetstrokecolor{textcolor}%
\pgfsetfillcolor{textcolor}%
\pgftext[x=0.494955in,y=2.686482in,,bottom,rotate=90.000000]{\color{textcolor}{\rmfamily\fontsize{24.000000}{28.800000}\selectfont\catcode`\^=\active\def^{\ifmmode\sp\else\^{}\fi}\catcode`\%=\active\def%{\%}$x (\mu m)$}}%
\end{pgfscope}%
\begin{pgfscope}%
\pgfpathrectangle{\pgfqpoint{0.894648in}{0.771190in}}{\pgfqpoint{5.355353in}{3.830585in}}%
\pgfusepath{clip}%
\pgfsetrectcap%
\pgfsetroundjoin%
\pgfsetlinewidth{1.505625pt}%
\definecolor{currentstroke}{rgb}{0.121569,0.466667,0.705882}%
\pgfsetstrokecolor{currentstroke}%
\pgfsetdash{}{0pt}%
\pgfpathmoveto{\pgfqpoint{1.138073in}{2.303424in}}%
\pgfpathlineto{\pgfqpoint{1.138710in}{2.304443in}}%
\pgfpathlineto{\pgfqpoint{1.138725in}{2.305048in}}%
\pgfpathlineto{\pgfqpoint{1.139251in}{2.301076in}}%
\pgfpathlineto{\pgfqpoint{1.139801in}{2.303419in}}%
\pgfpathlineto{\pgfqpoint{1.140200in}{2.301914in}}%
\pgfpathlineto{\pgfqpoint{1.139913in}{2.304916in}}%
\pgfpathlineto{\pgfqpoint{1.140838in}{2.303692in}}%
\pgfpathlineto{\pgfqpoint{1.140965in}{2.304360in}}%
\pgfpathlineto{\pgfqpoint{1.141033in}{2.302302in}}%
\pgfpathlineto{\pgfqpoint{1.141933in}{2.303095in}}%
\pgfpathlineto{\pgfqpoint{1.142479in}{2.304413in}}%
\pgfpathlineto{\pgfqpoint{1.142834in}{2.301277in}}%
\pgfpathlineto{\pgfqpoint{1.142946in}{2.301947in}}%
\pgfpathlineto{\pgfqpoint{1.143676in}{2.301487in}}%
\pgfpathlineto{\pgfqpoint{1.143262in}{2.304284in}}%
\pgfpathlineto{\pgfqpoint{1.144017in}{2.303391in}}%
\pgfpathlineto{\pgfqpoint{1.144280in}{2.304357in}}%
\pgfpathlineto{\pgfqpoint{1.145064in}{2.301755in}}%
\pgfpathlineto{\pgfqpoint{1.145074in}{2.301855in}}%
\pgfpathlineto{\pgfqpoint{1.145799in}{2.299363in}}%
\pgfpathlineto{\pgfqpoint{1.145137in}{2.302474in}}%
\pgfpathlineto{\pgfqpoint{1.146257in}{2.299926in}}%
\pgfpathlineto{\pgfqpoint{1.146816in}{2.305348in}}%
\pgfpathlineto{\pgfqpoint{1.146315in}{2.299298in}}%
\pgfpathlineto{\pgfqpoint{1.147712in}{2.303926in}}%
\pgfpathlineto{\pgfqpoint{1.148024in}{2.301873in}}%
\pgfpathlineto{\pgfqpoint{1.148408in}{2.304296in}}%
\pgfpathlineto{\pgfqpoint{1.148817in}{2.303785in}}%
\pgfpathlineto{\pgfqpoint{1.149129in}{2.301031in}}%
\pgfpathlineto{\pgfqpoint{1.149937in}{2.304337in}}%
\pgfpathlineto{\pgfqpoint{1.150000in}{2.304997in}}%
\pgfpathlineto{\pgfqpoint{1.150648in}{2.300657in}}%
\pgfpathlineto{\pgfqpoint{1.150672in}{2.300886in}}%
\pgfpathlineto{\pgfqpoint{1.151393in}{2.300220in}}%
\pgfpathlineto{\pgfqpoint{1.151592in}{2.303160in}}%
\pgfpathlineto{\pgfqpoint{1.151773in}{2.301568in}}%
\pgfpathlineto{\pgfqpoint{1.152800in}{2.303861in}}%
\pgfpathlineto{\pgfqpoint{1.152079in}{2.300457in}}%
\pgfpathlineto{\pgfqpoint{1.153024in}{2.302907in}}%
\pgfpathlineto{\pgfqpoint{1.153861in}{2.301687in}}%
\pgfpathlineto{\pgfqpoint{1.153257in}{2.304098in}}%
\pgfpathlineto{\pgfqpoint{1.154134in}{2.302954in}}%
\pgfpathlineto{\pgfqpoint{1.155268in}{2.300326in}}%
\pgfpathlineto{\pgfqpoint{1.154694in}{2.302993in}}%
\pgfpathlineto{\pgfqpoint{1.155302in}{2.300399in}}%
\pgfpathlineto{\pgfqpoint{1.155424in}{2.299638in}}%
\pgfpathlineto{\pgfqpoint{1.156261in}{2.304386in}}%
\pgfpathlineto{\pgfqpoint{1.156291in}{2.304007in}}%
\pgfpathlineto{\pgfqpoint{1.156305in}{2.304285in}}%
\pgfpathlineto{\pgfqpoint{1.156962in}{2.300818in}}%
\pgfpathlineto{\pgfqpoint{1.157045in}{2.301156in}}%
\pgfpathlineto{\pgfqpoint{1.157537in}{2.296271in}}%
\pgfpathlineto{\pgfqpoint{1.158072in}{2.301292in}}%
\pgfpathlineto{\pgfqpoint{1.158160in}{2.300850in}}%
\pgfpathlineto{\pgfqpoint{1.158292in}{2.299844in}}%
\pgfpathlineto{\pgfqpoint{1.159036in}{2.303851in}}%
\pgfpathlineto{\pgfqpoint{1.159095in}{2.302932in}}%
\pgfpathlineto{\pgfqpoint{1.159864in}{2.303840in}}%
\pgfpathlineto{\pgfqpoint{1.160078in}{2.301470in}}%
\pgfpathlineto{\pgfqpoint{1.160185in}{2.301776in}}%
\pgfpathlineto{\pgfqpoint{1.160302in}{2.303175in}}%
\pgfpathlineto{\pgfqpoint{1.160838in}{2.300323in}}%
\pgfpathlineto{\pgfqpoint{1.161841in}{2.299649in}}%
\pgfpathlineto{\pgfqpoint{1.161140in}{2.302189in}}%
\pgfpathlineto{\pgfqpoint{1.161943in}{2.300475in}}%
\pgfpathlineto{\pgfqpoint{1.162542in}{2.302867in}}%
\pgfpathlineto{\pgfqpoint{1.162235in}{2.299168in}}%
\pgfpathlineto{\pgfqpoint{1.163063in}{2.300712in}}%
\pgfpathlineto{\pgfqpoint{1.163467in}{2.300024in}}%
\pgfpathlineto{\pgfqpoint{1.164070in}{2.301879in}}%
\pgfpathlineto{\pgfqpoint{1.164158in}{2.300872in}}%
\pgfpathlineto{\pgfqpoint{1.165273in}{2.306453in}}%
\pgfpathlineto{\pgfqpoint{1.165331in}{2.304989in}}%
\pgfpathlineto{\pgfqpoint{1.165448in}{2.304503in}}%
\pgfpathlineto{\pgfqpoint{1.165638in}{2.307683in}}%
\pgfpathlineto{\pgfqpoint{1.166334in}{2.306829in}}%
\pgfpathlineto{\pgfqpoint{1.166339in}{2.306899in}}%
\pgfpathlineto{\pgfqpoint{1.167103in}{2.302084in}}%
\pgfpathlineto{\pgfqpoint{1.167605in}{2.299575in}}%
\pgfpathlineto{\pgfqpoint{1.167250in}{2.303889in}}%
\pgfpathlineto{\pgfqpoint{1.168243in}{2.301260in}}%
\pgfpathlineto{\pgfqpoint{1.168282in}{2.300887in}}%
\pgfpathlineto{\pgfqpoint{1.168554in}{2.303609in}}%
\pgfpathlineto{\pgfqpoint{1.168559in}{2.303574in}}%
\pgfpathlineto{\pgfqpoint{1.169547in}{2.304373in}}%
\pgfpathlineto{\pgfqpoint{1.169095in}{2.300886in}}%
\pgfpathlineto{\pgfqpoint{1.169659in}{2.303227in}}%
\pgfpathlineto{\pgfqpoint{1.170789in}{2.305314in}}%
\pgfpathlineto{\pgfqpoint{1.170122in}{2.300353in}}%
\pgfpathlineto{\pgfqpoint{1.170925in}{2.304604in}}%
\pgfpathlineto{\pgfqpoint{1.171583in}{2.300386in}}%
\pgfpathlineto{\pgfqpoint{1.170959in}{2.304657in}}%
\pgfpathlineto{\pgfqpoint{1.172074in}{2.302990in}}%
\pgfpathlineto{\pgfqpoint{1.172766in}{2.303768in}}%
\pgfpathlineto{\pgfqpoint{1.172269in}{2.301659in}}%
\pgfpathlineto{\pgfqpoint{1.173184in}{2.302973in}}%
\pgfpathlineto{\pgfqpoint{1.173325in}{2.304750in}}%
\pgfpathlineto{\pgfqpoint{1.173983in}{2.301863in}}%
\pgfpathlineto{\pgfqpoint{1.174328in}{2.304421in}}%
\pgfpathlineto{\pgfqpoint{1.174752in}{2.300509in}}%
\pgfpathlineto{\pgfqpoint{1.174358in}{2.304830in}}%
\pgfpathlineto{\pgfqpoint{1.175458in}{2.302890in}}%
\pgfpathlineto{\pgfqpoint{1.175531in}{2.304106in}}%
\pgfpathlineto{\pgfqpoint{1.176285in}{2.299895in}}%
\pgfpathlineto{\pgfqpoint{1.176436in}{2.301225in}}%
\pgfpathlineto{\pgfqpoint{1.177478in}{2.299871in}}%
\pgfpathlineto{\pgfqpoint{1.177040in}{2.302246in}}%
\pgfpathlineto{\pgfqpoint{1.177571in}{2.300881in}}%
\pgfpathlineto{\pgfqpoint{1.178652in}{2.302452in}}%
\pgfpathlineto{\pgfqpoint{1.177985in}{2.299054in}}%
\pgfpathlineto{\pgfqpoint{1.178705in}{2.302450in}}%
\pgfpathlineto{\pgfqpoint{1.178739in}{2.301551in}}%
\pgfpathlineto{\pgfqpoint{1.179177in}{2.304245in}}%
\pgfpathlineto{\pgfqpoint{1.179796in}{2.303076in}}%
\pgfpathlineto{\pgfqpoint{1.180502in}{2.303939in}}%
\pgfpathlineto{\pgfqpoint{1.180648in}{2.302055in}}%
\pgfpathlineto{\pgfqpoint{1.180906in}{2.303242in}}%
\pgfpathlineto{\pgfqpoint{1.181529in}{2.301640in}}%
\pgfpathlineto{\pgfqpoint{1.180935in}{2.303738in}}%
\pgfpathlineto{\pgfqpoint{1.181543in}{2.301729in}}%
\pgfpathlineto{\pgfqpoint{1.182103in}{2.298085in}}%
\pgfpathlineto{\pgfqpoint{1.181573in}{2.301952in}}%
\pgfpathlineto{\pgfqpoint{1.182673in}{2.300439in}}%
\pgfpathlineto{\pgfqpoint{1.183530in}{2.303212in}}%
\pgfpathlineto{\pgfqpoint{1.183812in}{2.301282in}}%
\pgfpathlineto{\pgfqpoint{1.185502in}{2.306901in}}%
\pgfpathlineto{\pgfqpoint{1.184022in}{2.298812in}}%
\pgfpathlineto{\pgfqpoint{1.185872in}{2.303560in}}%
\pgfpathlineto{\pgfqpoint{1.186534in}{2.302235in}}%
\pgfpathlineto{\pgfqpoint{1.186076in}{2.304212in}}%
\pgfpathlineto{\pgfqpoint{1.186982in}{2.303427in}}%
\pgfpathlineto{\pgfqpoint{1.187381in}{2.303477in}}%
\pgfpathlineto{\pgfqpoint{1.187094in}{2.302087in}}%
\pgfpathlineto{\pgfqpoint{1.187459in}{2.302326in}}%
\pgfpathlineto{\pgfqpoint{1.188593in}{2.299357in}}%
\pgfpathlineto{\pgfqpoint{1.188014in}{2.304013in}}%
\pgfpathlineto{\pgfqpoint{1.188617in}{2.299809in}}%
\pgfpathlineto{\pgfqpoint{1.189270in}{2.303472in}}%
\pgfpathlineto{\pgfqpoint{1.188627in}{2.299782in}}%
\pgfpathlineto{\pgfqpoint{1.190073in}{2.301947in}}%
\pgfpathlineto{\pgfqpoint{1.190467in}{2.299410in}}%
\pgfpathlineto{\pgfqpoint{1.191154in}{2.303071in}}%
\pgfpathlineto{\pgfqpoint{1.191168in}{2.302835in}}%
\pgfpathlineto{\pgfqpoint{1.191207in}{2.302341in}}%
\pgfpathlineto{\pgfqpoint{1.191568in}{2.305515in}}%
\pgfpathlineto{\pgfqpoint{1.192084in}{2.304781in}}%
\pgfpathlineto{\pgfqpoint{1.192171in}{2.305692in}}%
\pgfpathlineto{\pgfqpoint{1.192863in}{2.302508in}}%
\pgfpathlineto{\pgfqpoint{1.193165in}{2.303108in}}%
\pgfpathlineto{\pgfqpoint{1.193715in}{2.301610in}}%
\pgfpathlineto{\pgfqpoint{1.194060in}{2.304517in}}%
\pgfpathlineto{\pgfqpoint{1.194070in}{2.304493in}}%
\pgfpathlineto{\pgfqpoint{1.194464in}{2.307070in}}%
\pgfpathlineto{\pgfqpoint{1.194995in}{2.302612in}}%
\pgfpathlineto{\pgfqpoint{1.195136in}{2.303337in}}%
\pgfpathlineto{\pgfqpoint{1.196251in}{2.301408in}}%
\pgfpathlineto{\pgfqpoint{1.195876in}{2.304839in}}%
\pgfpathlineto{\pgfqpoint{1.196305in}{2.301803in}}%
\pgfpathlineto{\pgfqpoint{1.197225in}{2.303624in}}%
\pgfpathlineto{\pgfqpoint{1.196646in}{2.299650in}}%
\pgfpathlineto{\pgfqpoint{1.197425in}{2.302462in}}%
\pgfpathlineto{\pgfqpoint{1.197673in}{2.301852in}}%
\pgfpathlineto{\pgfqpoint{1.197809in}{2.304442in}}%
\pgfpathlineto{\pgfqpoint{1.198384in}{2.303212in}}%
\pgfpathlineto{\pgfqpoint{1.198486in}{2.304420in}}%
\pgfpathlineto{\pgfqpoint{1.199187in}{2.299596in}}%
\pgfpathlineto{\pgfqpoint{1.199479in}{2.302601in}}%
\pgfpathlineto{\pgfqpoint{1.200745in}{2.304555in}}%
\pgfpathlineto{\pgfqpoint{1.199664in}{2.301829in}}%
\pgfpathlineto{\pgfqpoint{1.200896in}{2.303496in}}%
\pgfpathlineto{\pgfqpoint{1.201899in}{2.302566in}}%
\pgfpathlineto{\pgfqpoint{1.201081in}{2.304064in}}%
\pgfpathlineto{\pgfqpoint{1.202001in}{2.303531in}}%
\pgfpathlineto{\pgfqpoint{1.202751in}{2.303779in}}%
\pgfpathlineto{\pgfqpoint{1.202551in}{2.300753in}}%
\pgfpathlineto{\pgfqpoint{1.202994in}{2.302090in}}%
\pgfpathlineto{\pgfqpoint{1.203442in}{2.298440in}}%
\pgfpathlineto{\pgfqpoint{1.204051in}{2.303041in}}%
\pgfpathlineto{\pgfqpoint{1.204055in}{2.302990in}}%
\pgfpathlineto{\pgfqpoint{1.204971in}{2.308004in}}%
\pgfpathlineto{\pgfqpoint{1.205190in}{2.304253in}}%
\pgfpathlineto{\pgfqpoint{1.205594in}{2.302332in}}%
\pgfpathlineto{\pgfqpoint{1.206246in}{2.304773in}}%
\pgfpathlineto{\pgfqpoint{1.206251in}{2.305167in}}%
\pgfpathlineto{\pgfqpoint{1.206938in}{2.301276in}}%
\pgfpathlineto{\pgfqpoint{1.207312in}{2.301645in}}%
\pgfpathlineto{\pgfqpoint{1.207327in}{2.301155in}}%
\pgfpathlineto{\pgfqpoint{1.207629in}{2.305000in}}%
\pgfpathlineto{\pgfqpoint{1.208422in}{2.301630in}}%
\pgfpathlineto{\pgfqpoint{1.209435in}{2.299489in}}%
\pgfpathlineto{\pgfqpoint{1.208539in}{2.303589in}}%
\pgfpathlineto{\pgfqpoint{1.209567in}{2.300701in}}%
\pgfpathlineto{\pgfqpoint{1.209576in}{2.301093in}}%
\pgfpathlineto{\pgfqpoint{1.210404in}{2.297796in}}%
\pgfpathlineto{\pgfqpoint{1.210419in}{2.297613in}}%
\pgfpathlineto{\pgfqpoint{1.211202in}{2.304346in}}%
\pgfpathlineto{\pgfqpoint{1.211295in}{2.303852in}}%
\pgfpathlineto{\pgfqpoint{1.211300in}{2.304186in}}%
\pgfpathlineto{\pgfqpoint{1.211675in}{2.301271in}}%
\pgfpathlineto{\pgfqpoint{1.212390in}{2.302920in}}%
\pgfpathlineto{\pgfqpoint{1.213398in}{2.300273in}}%
\pgfpathlineto{\pgfqpoint{1.212410in}{2.303407in}}%
\pgfpathlineto{\pgfqpoint{1.213549in}{2.300761in}}%
\pgfpathlineto{\pgfqpoint{1.214742in}{2.304703in}}%
\pgfpathlineto{\pgfqpoint{1.213568in}{2.300452in}}%
\pgfpathlineto{\pgfqpoint{1.214795in}{2.304191in}}%
\pgfpathlineto{\pgfqpoint{1.215604in}{2.306037in}}%
\pgfpathlineto{\pgfqpoint{1.215014in}{2.303171in}}%
\pgfpathlineto{\pgfqpoint{1.215896in}{2.304104in}}%
\pgfpathlineto{\pgfqpoint{1.216855in}{2.299787in}}%
\pgfpathlineto{\pgfqpoint{1.216217in}{2.304378in}}%
\pgfpathlineto{\pgfqpoint{1.217132in}{2.300389in}}%
\pgfpathlineto{\pgfqpoint{1.217166in}{2.301140in}}%
\pgfpathlineto{\pgfqpoint{1.217677in}{2.299016in}}%
\pgfpathlineto{\pgfqpoint{1.218242in}{2.300614in}}%
\pgfpathlineto{\pgfqpoint{1.219211in}{2.303540in}}%
\pgfpathlineto{\pgfqpoint{1.218296in}{2.299804in}}%
\pgfpathlineto{\pgfqpoint{1.219790in}{2.301231in}}%
\pgfpathlineto{\pgfqpoint{1.219995in}{2.299645in}}%
\pgfpathlineto{\pgfqpoint{1.220730in}{2.304084in}}%
\pgfpathlineto{\pgfqpoint{1.220876in}{2.303488in}}%
\pgfpathlineto{\pgfqpoint{1.221587in}{2.298332in}}%
\pgfpathlineto{\pgfqpoint{1.222507in}{2.300872in}}%
\pgfpathlineto{\pgfqpoint{1.223232in}{2.302422in}}%
\pgfpathlineto{\pgfqpoint{1.223602in}{2.299686in}}%
\pgfpathlineto{\pgfqpoint{1.224192in}{2.297841in}}%
\pgfpathlineto{\pgfqpoint{1.223710in}{2.300208in}}%
\pgfpathlineto{\pgfqpoint{1.224576in}{2.300074in}}%
\pgfpathlineto{\pgfqpoint{1.224732in}{2.301675in}}%
\pgfpathlineto{\pgfqpoint{1.225623in}{2.297068in}}%
\pgfpathlineto{\pgfqpoint{1.225808in}{2.297975in}}%
\pgfpathlineto{\pgfqpoint{1.226416in}{2.293636in}}%
\pgfpathlineto{\pgfqpoint{1.226631in}{2.294835in}}%
\pgfpathlineto{\pgfqpoint{1.226640in}{2.294671in}}%
\pgfpathlineto{\pgfqpoint{1.227439in}{2.300380in}}%
\pgfpathlineto{\pgfqpoint{1.227561in}{2.300146in}}%
\pgfpathlineto{\pgfqpoint{1.228632in}{2.303743in}}%
\pgfpathlineto{\pgfqpoint{1.228184in}{2.299742in}}%
\pgfpathlineto{\pgfqpoint{1.228734in}{2.303016in}}%
\pgfpathlineto{\pgfqpoint{1.228758in}{2.302086in}}%
\pgfpathlineto{\pgfqpoint{1.229815in}{2.305398in}}%
\pgfpathlineto{\pgfqpoint{1.230370in}{2.307695in}}%
\pgfpathlineto{\pgfqpoint{1.230720in}{2.302733in}}%
\pgfpathlineto{\pgfqpoint{1.230900in}{2.304386in}}%
\pgfpathlineto{\pgfqpoint{1.231071in}{2.302256in}}%
\pgfpathlineto{\pgfqpoint{1.231918in}{2.304767in}}%
\pgfpathlineto{\pgfqpoint{1.232049in}{2.302918in}}%
\pgfpathlineto{\pgfqpoint{1.232410in}{2.300874in}}%
\pgfpathlineto{\pgfqpoint{1.232629in}{2.304232in}}%
\pgfpathlineto{\pgfqpoint{1.233140in}{2.303089in}}%
\pgfpathlineto{\pgfqpoint{1.233914in}{2.305015in}}%
\pgfpathlineto{\pgfqpoint{1.234235in}{2.302564in}}%
\pgfpathlineto{\pgfqpoint{1.234240in}{2.302694in}}%
\pgfpathlineto{\pgfqpoint{1.234756in}{2.301369in}}%
\pgfpathlineto{\pgfqpoint{1.234474in}{2.304460in}}%
\pgfpathlineto{\pgfqpoint{1.235272in}{2.303192in}}%
\pgfpathlineto{\pgfqpoint{1.235560in}{2.305346in}}%
\pgfpathlineto{\pgfqpoint{1.236319in}{2.302493in}}%
\pgfpathlineto{\pgfqpoint{1.236402in}{2.304149in}}%
\pgfpathlineto{\pgfqpoint{1.237127in}{2.301190in}}%
\pgfpathlineto{\pgfqpoint{1.237526in}{2.303256in}}%
\pgfpathlineto{\pgfqpoint{1.237789in}{2.303897in}}%
\pgfpathlineto{\pgfqpoint{1.238140in}{2.300971in}}%
\pgfpathlineto{\pgfqpoint{1.238150in}{2.300467in}}%
\pgfpathlineto{\pgfqpoint{1.238943in}{2.303832in}}%
\pgfpathlineto{\pgfqpoint{1.239196in}{2.303398in}}%
\pgfpathlineto{\pgfqpoint{1.239756in}{2.305058in}}%
\pgfpathlineto{\pgfqpoint{1.239337in}{2.302316in}}%
\pgfpathlineto{\pgfqpoint{1.240292in}{2.303082in}}%
\pgfpathlineto{\pgfqpoint{1.241110in}{2.298062in}}%
\pgfpathlineto{\pgfqpoint{1.241431in}{2.299888in}}%
\pgfpathlineto{\pgfqpoint{1.241820in}{2.298481in}}%
\pgfpathlineto{\pgfqpoint{1.242205in}{2.303064in}}%
\pgfpathlineto{\pgfqpoint{1.242507in}{2.301388in}}%
\pgfpathlineto{\pgfqpoint{1.242765in}{2.303530in}}%
\pgfpathlineto{\pgfqpoint{1.243339in}{2.301158in}}%
\pgfpathlineto{\pgfqpoint{1.243631in}{2.302377in}}%
\pgfpathlineto{\pgfqpoint{1.244084in}{2.301300in}}%
\pgfpathlineto{\pgfqpoint{1.244469in}{2.304659in}}%
\pgfpathlineto{\pgfqpoint{1.244605in}{2.304461in}}%
\pgfpathlineto{\pgfqpoint{1.244610in}{2.304729in}}%
\pgfpathlineto{\pgfqpoint{1.244946in}{2.300739in}}%
\pgfpathlineto{\pgfqpoint{1.245667in}{2.301994in}}%
\pgfpathlineto{\pgfqpoint{1.245963in}{2.300796in}}%
\pgfpathlineto{\pgfqpoint{1.246519in}{2.304074in}}%
\pgfpathlineto{\pgfqpoint{1.246923in}{2.306434in}}%
\pgfpathlineto{\pgfqpoint{1.247487in}{2.301999in}}%
\pgfpathlineto{\pgfqpoint{1.247570in}{2.302516in}}%
\pgfpathlineto{\pgfqpoint{1.247629in}{2.301573in}}%
\pgfpathlineto{\pgfqpoint{1.247833in}{2.304253in}}%
\pgfpathlineto{\pgfqpoint{1.248247in}{2.303650in}}%
\pgfpathlineto{\pgfqpoint{1.249259in}{2.305729in}}%
\pgfpathlineto{\pgfqpoint{1.248928in}{2.302695in}}%
\pgfpathlineto{\pgfqpoint{1.249362in}{2.303964in}}%
\pgfpathlineto{\pgfqpoint{1.249522in}{2.303466in}}%
\pgfpathlineto{\pgfqpoint{1.250272in}{2.306979in}}%
\pgfpathlineto{\pgfqpoint{1.250428in}{2.306714in}}%
\pgfpathlineto{\pgfqpoint{1.250886in}{2.301447in}}%
\pgfpathlineto{\pgfqpoint{1.250511in}{2.307112in}}%
\pgfpathlineto{\pgfqpoint{1.250915in}{2.301634in}}%
\pgfpathlineto{\pgfqpoint{1.251669in}{2.299773in}}%
\pgfpathlineto{\pgfqpoint{1.251426in}{2.302625in}}%
\pgfpathlineto{\pgfqpoint{1.252025in}{2.301698in}}%
\pgfpathlineto{\pgfqpoint{1.253179in}{2.304258in}}%
\pgfpathlineto{\pgfqpoint{1.252667in}{2.300836in}}%
\pgfpathlineto{\pgfqpoint{1.253183in}{2.304030in}}%
\pgfpathlineto{\pgfqpoint{1.254089in}{2.302148in}}%
\pgfpathlineto{\pgfqpoint{1.253617in}{2.305836in}}%
\pgfpathlineto{\pgfqpoint{1.254303in}{2.303396in}}%
\pgfpathlineto{\pgfqpoint{1.255399in}{2.304797in}}%
\pgfpathlineto{\pgfqpoint{1.255107in}{2.302462in}}%
\pgfpathlineto{\pgfqpoint{1.255433in}{2.304453in}}%
\pgfpathlineto{\pgfqpoint{1.256489in}{2.298271in}}%
\pgfpathlineto{\pgfqpoint{1.256635in}{2.298479in}}%
\pgfpathlineto{\pgfqpoint{1.257370in}{2.301179in}}%
\pgfpathlineto{\pgfqpoint{1.257658in}{2.296871in}}%
\pgfpathlineto{\pgfqpoint{1.257721in}{2.297497in}}%
\pgfpathlineto{\pgfqpoint{1.258709in}{2.294983in}}%
\pgfpathlineto{\pgfqpoint{1.258101in}{2.299244in}}%
\pgfpathlineto{\pgfqpoint{1.258826in}{2.297656in}}%
\pgfpathlineto{\pgfqpoint{1.258850in}{2.297201in}}%
\pgfpathlineto{\pgfqpoint{1.259245in}{2.302194in}}%
\pgfpathlineto{\pgfqpoint{1.259556in}{2.301912in}}%
\pgfpathlineto{\pgfqpoint{1.259975in}{2.305049in}}%
\pgfpathlineto{\pgfqpoint{1.260486in}{2.301633in}}%
\pgfpathlineto{\pgfqpoint{1.260696in}{2.302818in}}%
\pgfpathlineto{\pgfqpoint{1.261786in}{2.301300in}}%
\pgfpathlineto{\pgfqpoint{1.261572in}{2.304532in}}%
\pgfpathlineto{\pgfqpoint{1.261830in}{2.301624in}}%
\pgfpathlineto{\pgfqpoint{1.262697in}{2.304294in}}%
\pgfpathlineto{\pgfqpoint{1.262283in}{2.301074in}}%
\pgfpathlineto{\pgfqpoint{1.262964in}{2.303230in}}%
\pgfpathlineto{\pgfqpoint{1.263412in}{2.302102in}}%
\pgfpathlineto{\pgfqpoint{1.264045in}{2.304668in}}%
\pgfpathlineto{\pgfqpoint{1.264055in}{2.304374in}}%
\pgfpathlineto{\pgfqpoint{1.264746in}{2.301404in}}%
\pgfpathlineto{\pgfqpoint{1.264065in}{2.304377in}}%
\pgfpathlineto{\pgfqpoint{1.265238in}{2.302954in}}%
\pgfpathlineto{\pgfqpoint{1.265350in}{2.304081in}}%
\pgfpathlineto{\pgfqpoint{1.265647in}{2.300793in}}%
\pgfpathlineto{\pgfqpoint{1.266353in}{2.303697in}}%
\pgfpathlineto{\pgfqpoint{1.267507in}{2.300192in}}%
\pgfpathlineto{\pgfqpoint{1.266411in}{2.304086in}}%
\pgfpathlineto{\pgfqpoint{1.267526in}{2.300518in}}%
\pgfpathlineto{\pgfqpoint{1.268247in}{2.298033in}}%
\pgfpathlineto{\pgfqpoint{1.267667in}{2.301280in}}%
\pgfpathlineto{\pgfqpoint{1.268768in}{2.300688in}}%
\pgfpathlineto{\pgfqpoint{1.269478in}{2.303611in}}%
\pgfpathlineto{\pgfqpoint{1.268845in}{2.300386in}}%
\pgfpathlineto{\pgfqpoint{1.269887in}{2.301297in}}%
\pgfpathlineto{\pgfqpoint{1.270915in}{2.299404in}}%
\pgfpathlineto{\pgfqpoint{1.269897in}{2.301488in}}%
\pgfpathlineto{\pgfqpoint{1.271027in}{2.300480in}}%
\pgfpathlineto{\pgfqpoint{1.271173in}{2.302823in}}%
\pgfpathlineto{\pgfqpoint{1.272117in}{2.298901in}}%
\pgfpathlineto{\pgfqpoint{1.272575in}{2.296362in}}%
\pgfpathlineto{\pgfqpoint{1.273149in}{2.302234in}}%
\pgfpathlineto{\pgfqpoint{1.273188in}{2.302154in}}%
\pgfpathlineto{\pgfqpoint{1.273943in}{2.304334in}}%
\pgfpathlineto{\pgfqpoint{1.273407in}{2.301816in}}%
\pgfpathlineto{\pgfqpoint{1.274342in}{2.303177in}}%
\pgfpathlineto{\pgfqpoint{1.275476in}{2.301485in}}%
\pgfpathlineto{\pgfqpoint{1.274736in}{2.304537in}}%
\pgfpathlineto{\pgfqpoint{1.275481in}{2.301574in}}%
\pgfpathlineto{\pgfqpoint{1.275486in}{2.302000in}}%
\pgfpathlineto{\pgfqpoint{1.276304in}{2.296006in}}%
\pgfpathlineto{\pgfqpoint{1.276538in}{2.297357in}}%
\pgfpathlineto{\pgfqpoint{1.276825in}{2.295559in}}%
\pgfpathlineto{\pgfqpoint{1.277326in}{2.298270in}}%
\pgfpathlineto{\pgfqpoint{1.277643in}{2.297937in}}%
\pgfpathlineto{\pgfqpoint{1.278198in}{2.299021in}}%
\pgfpathlineto{\pgfqpoint{1.278393in}{2.296612in}}%
\pgfpathlineto{\pgfqpoint{1.278738in}{2.297399in}}%
\pgfpathlineto{\pgfqpoint{1.279809in}{2.294612in}}%
\pgfpathlineto{\pgfqpoint{1.279220in}{2.297553in}}%
\pgfpathlineto{\pgfqpoint{1.279873in}{2.295114in}}%
\pgfpathlineto{\pgfqpoint{1.281138in}{2.290200in}}%
\pgfpathlineto{\pgfqpoint{1.279936in}{2.295267in}}%
\pgfpathlineto{\pgfqpoint{1.281143in}{2.290429in}}%
\pgfpathlineto{\pgfqpoint{1.282838in}{2.294851in}}%
\pgfpathlineto{\pgfqpoint{1.281177in}{2.289913in}}%
\pgfpathlineto{\pgfqpoint{1.282896in}{2.294115in}}%
\pgfpathlineto{\pgfqpoint{1.283996in}{2.292281in}}%
\pgfpathlineto{\pgfqpoint{1.283743in}{2.295649in}}%
\pgfpathlineto{\pgfqpoint{1.284035in}{2.292584in}}%
\pgfpathlineto{\pgfqpoint{1.285097in}{2.297254in}}%
\pgfpathlineto{\pgfqpoint{1.284629in}{2.292095in}}%
\pgfpathlineto{\pgfqpoint{1.285174in}{2.296174in}}%
\pgfpathlineto{\pgfqpoint{1.285194in}{2.295839in}}%
\pgfpathlineto{\pgfqpoint{1.285773in}{2.298067in}}%
\pgfpathlineto{\pgfqpoint{1.286270in}{2.296864in}}%
\pgfpathlineto{\pgfqpoint{1.286314in}{2.296343in}}%
\pgfpathlineto{\pgfqpoint{1.286942in}{2.299762in}}%
\pgfpathlineto{\pgfqpoint{1.287195in}{2.299175in}}%
\pgfpathlineto{\pgfqpoint{1.288285in}{2.300089in}}%
\pgfpathlineto{\pgfqpoint{1.287560in}{2.296355in}}%
\pgfpathlineto{\pgfqpoint{1.288310in}{2.299780in}}%
\pgfpathlineto{\pgfqpoint{1.289220in}{2.297051in}}%
\pgfpathlineto{\pgfqpoint{1.288514in}{2.300505in}}%
\pgfpathlineto{\pgfqpoint{1.289434in}{2.298429in}}%
\pgfpathlineto{\pgfqpoint{1.290218in}{2.299061in}}%
\pgfpathlineto{\pgfqpoint{1.290437in}{2.296478in}}%
\pgfpathlineto{\pgfqpoint{1.290476in}{2.296801in}}%
\pgfpathlineto{\pgfqpoint{1.291542in}{2.292086in}}%
\pgfpathlineto{\pgfqpoint{1.291645in}{2.293508in}}%
\pgfpathlineto{\pgfqpoint{1.292209in}{2.296237in}}%
\pgfpathlineto{\pgfqpoint{1.292716in}{2.292249in}}%
\pgfpathlineto{\pgfqpoint{1.292721in}{2.292297in}}%
\pgfpathlineto{\pgfqpoint{1.293480in}{2.290367in}}%
\pgfpathlineto{\pgfqpoint{1.292969in}{2.293447in}}%
\pgfpathlineto{\pgfqpoint{1.293797in}{2.292297in}}%
\pgfpathlineto{\pgfqpoint{1.294950in}{2.299153in}}%
\pgfpathlineto{\pgfqpoint{1.294994in}{2.298502in}}%
\pgfpathlineto{\pgfqpoint{1.296265in}{2.293419in}}%
\pgfpathlineto{\pgfqpoint{1.295028in}{2.298513in}}%
\pgfpathlineto{\pgfqpoint{1.296328in}{2.294351in}}%
\pgfpathlineto{\pgfqpoint{1.297204in}{2.297664in}}%
\pgfpathlineto{\pgfqpoint{1.296761in}{2.292760in}}%
\pgfpathlineto{\pgfqpoint{1.297501in}{2.296613in}}%
\pgfpathlineto{\pgfqpoint{1.297545in}{2.295513in}}%
\pgfpathlineto{\pgfqpoint{1.298519in}{2.299322in}}%
\pgfpathlineto{\pgfqpoint{1.298524in}{2.299216in}}%
\pgfpathlineto{\pgfqpoint{1.298952in}{2.302570in}}%
\pgfpathlineto{\pgfqpoint{1.298563in}{2.299134in}}%
\pgfpathlineto{\pgfqpoint{1.299692in}{2.302467in}}%
\pgfpathlineto{\pgfqpoint{1.299814in}{2.300809in}}%
\pgfpathlineto{\pgfqpoint{1.300374in}{2.304227in}}%
\pgfpathlineto{\pgfqpoint{1.300754in}{2.302982in}}%
\pgfpathlineto{\pgfqpoint{1.301187in}{2.304727in}}%
\pgfpathlineto{\pgfqpoint{1.301771in}{2.301131in}}%
\pgfpathlineto{\pgfqpoint{1.301830in}{2.301232in}}%
\pgfpathlineto{\pgfqpoint{1.302983in}{2.299548in}}%
\pgfpathlineto{\pgfqpoint{1.302263in}{2.304850in}}%
\pgfpathlineto{\pgfqpoint{1.303008in}{2.299721in}}%
\pgfpathlineto{\pgfqpoint{1.303013in}{2.300162in}}%
\pgfpathlineto{\pgfqpoint{1.303894in}{2.293882in}}%
\pgfpathlineto{\pgfqpoint{1.304064in}{2.294896in}}%
\pgfpathlineto{\pgfqpoint{1.304517in}{2.294580in}}%
\pgfpathlineto{\pgfqpoint{1.304712in}{2.297036in}}%
\pgfpathlineto{\pgfqpoint{1.305111in}{2.296166in}}%
\pgfpathlineto{\pgfqpoint{1.307195in}{2.303864in}}%
\pgfpathlineto{\pgfqpoint{1.305179in}{2.295474in}}%
\pgfpathlineto{\pgfqpoint{1.307458in}{2.303487in}}%
\pgfpathlineto{\pgfqpoint{1.307526in}{2.302383in}}%
\pgfpathlineto{\pgfqpoint{1.308027in}{2.305203in}}%
\pgfpathlineto{\pgfqpoint{1.308538in}{2.304836in}}%
\pgfpathlineto{\pgfqpoint{1.309819in}{2.301265in}}%
\pgfpathlineto{\pgfqpoint{1.308645in}{2.305662in}}%
\pgfpathlineto{\pgfqpoint{1.309833in}{2.301485in}}%
\pgfpathlineto{\pgfqpoint{1.309838in}{2.301704in}}%
\pgfpathlineto{\pgfqpoint{1.310656in}{2.298996in}}%
\pgfpathlineto{\pgfqpoint{1.310900in}{2.299535in}}%
\pgfpathlineto{\pgfqpoint{1.312185in}{2.294423in}}%
\pgfpathlineto{\pgfqpoint{1.312199in}{2.294991in}}%
\pgfpathlineto{\pgfqpoint{1.313378in}{2.297872in}}%
\pgfpathlineto{\pgfqpoint{1.312467in}{2.293663in}}%
\pgfpathlineto{\pgfqpoint{1.313451in}{2.296470in}}%
\pgfpathlineto{\pgfqpoint{1.313456in}{2.296131in}}%
\pgfpathlineto{\pgfqpoint{1.314390in}{2.301068in}}%
\pgfpathlineto{\pgfqpoint{1.314517in}{2.299470in}}%
\pgfpathlineto{\pgfqpoint{1.314527in}{2.299181in}}%
\pgfpathlineto{\pgfqpoint{1.314955in}{2.304290in}}%
\pgfpathlineto{\pgfqpoint{1.315559in}{2.301732in}}%
\pgfpathlineto{\pgfqpoint{1.316820in}{2.306798in}}%
\pgfpathlineto{\pgfqpoint{1.315744in}{2.301192in}}%
\pgfpathlineto{\pgfqpoint{1.317000in}{2.304383in}}%
\pgfpathlineto{\pgfqpoint{1.317034in}{2.304872in}}%
\pgfpathlineto{\pgfqpoint{1.318241in}{2.301933in}}%
\pgfpathlineto{\pgfqpoint{1.319376in}{2.303746in}}%
\pgfpathlineto{\pgfqpoint{1.318879in}{2.299981in}}%
\pgfpathlineto{\pgfqpoint{1.319390in}{2.303430in}}%
\pgfpathlineto{\pgfqpoint{1.319634in}{2.303129in}}%
\pgfpathlineto{\pgfqpoint{1.320023in}{2.306041in}}%
\pgfpathlineto{\pgfqpoint{1.320276in}{2.305800in}}%
\pgfpathlineto{\pgfqpoint{1.320617in}{2.309065in}}%
\pgfpathlineto{\pgfqpoint{1.321216in}{2.303212in}}%
\pgfpathlineto{\pgfqpoint{1.321221in}{2.303311in}}%
\pgfpathlineto{\pgfqpoint{1.321786in}{2.300871in}}%
\pgfpathlineto{\pgfqpoint{1.321250in}{2.303483in}}%
\pgfpathlineto{\pgfqpoint{1.322409in}{2.301851in}}%
\pgfpathlineto{\pgfqpoint{1.322686in}{2.304121in}}%
\pgfpathlineto{\pgfqpoint{1.322477in}{2.301220in}}%
\pgfpathlineto{\pgfqpoint{1.323631in}{2.303521in}}%
\pgfpathlineto{\pgfqpoint{1.324088in}{2.302230in}}%
\pgfpathlineto{\pgfqpoint{1.324453in}{2.305026in}}%
\pgfpathlineto{\pgfqpoint{1.324736in}{2.303513in}}%
\pgfpathlineto{\pgfqpoint{1.324775in}{2.304236in}}%
\pgfpathlineto{\pgfqpoint{1.325724in}{2.300201in}}%
\pgfpathlineto{\pgfqpoint{1.325778in}{2.300968in}}%
\pgfpathlineto{\pgfqpoint{1.326776in}{2.298906in}}%
\pgfpathlineto{\pgfqpoint{1.326308in}{2.302338in}}%
\pgfpathlineto{\pgfqpoint{1.326888in}{2.300806in}}%
\pgfpathlineto{\pgfqpoint{1.328056in}{2.305009in}}%
\pgfpathlineto{\pgfqpoint{1.326951in}{2.300111in}}%
\pgfpathlineto{\pgfqpoint{1.328095in}{2.303975in}}%
\pgfpathlineto{\pgfqpoint{1.328796in}{2.305525in}}%
\pgfpathlineto{\pgfqpoint{1.328416in}{2.302102in}}%
\pgfpathlineto{\pgfqpoint{1.329103in}{2.303632in}}%
\pgfpathlineto{\pgfqpoint{1.330096in}{2.302690in}}%
\pgfpathlineto{\pgfqpoint{1.329458in}{2.305620in}}%
\pgfpathlineto{\pgfqpoint{1.330203in}{2.304449in}}%
\pgfpathlineto{\pgfqpoint{1.330335in}{2.305009in}}%
\pgfpathlineto{\pgfqpoint{1.330856in}{2.302113in}}%
\pgfpathlineto{\pgfqpoint{1.331216in}{2.303806in}}%
\pgfpathlineto{\pgfqpoint{1.332394in}{2.300201in}}%
\pgfpathlineto{\pgfqpoint{1.331430in}{2.304923in}}%
\pgfpathlineto{\pgfqpoint{1.332433in}{2.300731in}}%
\pgfpathlineto{\pgfqpoint{1.332462in}{2.300455in}}%
\pgfpathlineto{\pgfqpoint{1.332861in}{2.304158in}}%
\pgfpathlineto{\pgfqpoint{1.333061in}{2.303698in}}%
\pgfpathlineto{\pgfqpoint{1.333694in}{2.305853in}}%
\pgfpathlineto{\pgfqpoint{1.334088in}{2.303214in}}%
\pgfpathlineto{\pgfqpoint{1.334171in}{2.303688in}}%
\pgfpathlineto{\pgfqpoint{1.334200in}{2.304457in}}%
\pgfpathlineto{\pgfqpoint{1.335198in}{2.300824in}}%
\pgfpathlineto{\pgfqpoint{1.335213in}{2.301004in}}%
\pgfpathlineto{\pgfqpoint{1.335422in}{2.299793in}}%
\pgfpathlineto{\pgfqpoint{1.336133in}{2.304929in}}%
\pgfpathlineto{\pgfqpoint{1.336177in}{2.304573in}}%
\pgfpathlineto{\pgfqpoint{1.337126in}{2.308157in}}%
\pgfpathlineto{\pgfqpoint{1.336523in}{2.303505in}}%
\pgfpathlineto{\pgfqpoint{1.337311in}{2.305575in}}%
\pgfpathlineto{\pgfqpoint{1.338441in}{2.302032in}}%
\pgfpathlineto{\pgfqpoint{1.338543in}{2.303414in}}%
\pgfpathlineto{\pgfqpoint{1.338587in}{2.304017in}}%
\pgfpathlineto{\pgfqpoint{1.339010in}{2.300965in}}%
\pgfpathlineto{\pgfqpoint{1.339643in}{2.303429in}}%
\pgfpathlineto{\pgfqpoint{1.340412in}{2.302325in}}%
\pgfpathlineto{\pgfqpoint{1.340062in}{2.304630in}}%
\pgfpathlineto{\pgfqpoint{1.340768in}{2.302758in}}%
\pgfpathlineto{\pgfqpoint{1.341021in}{2.304440in}}%
\pgfpathlineto{\pgfqpoint{1.341177in}{2.302074in}}%
\pgfpathlineto{\pgfqpoint{1.341810in}{2.302486in}}%
\pgfpathlineto{\pgfqpoint{1.342228in}{2.300724in}}%
\pgfpathlineto{\pgfqpoint{1.342535in}{2.304960in}}%
\pgfpathlineto{\pgfqpoint{1.342876in}{2.304500in}}%
\pgfpathlineto{\pgfqpoint{1.342920in}{2.305193in}}%
\pgfpathlineto{\pgfqpoint{1.343392in}{2.302782in}}%
\pgfpathlineto{\pgfqpoint{1.343962in}{2.303867in}}%
\pgfpathlineto{\pgfqpoint{1.344609in}{2.301545in}}%
\pgfpathlineto{\pgfqpoint{1.344059in}{2.305817in}}%
\pgfpathlineto{\pgfqpoint{1.345169in}{2.303261in}}%
\pgfpathlineto{\pgfqpoint{1.345933in}{2.305884in}}%
\pgfpathlineto{\pgfqpoint{1.345778in}{2.301990in}}%
\pgfpathlineto{\pgfqpoint{1.346279in}{2.303420in}}%
\pgfpathlineto{\pgfqpoint{1.346367in}{2.302675in}}%
\pgfpathlineto{\pgfqpoint{1.347063in}{2.306768in}}%
\pgfpathlineto{\pgfqpoint{1.347360in}{2.304290in}}%
\pgfpathlineto{\pgfqpoint{1.348416in}{2.302638in}}%
\pgfpathlineto{\pgfqpoint{1.347399in}{2.304635in}}%
\pgfpathlineto{\pgfqpoint{1.348528in}{2.303688in}}%
\pgfpathlineto{\pgfqpoint{1.348587in}{2.304850in}}%
\pgfpathlineto{\pgfqpoint{1.348601in}{2.305090in}}%
\pgfpathlineto{\pgfqpoint{1.348825in}{2.301966in}}%
\pgfpathlineto{\pgfqpoint{1.349658in}{2.303674in}}%
\pgfpathlineto{\pgfqpoint{1.350607in}{2.300707in}}%
\pgfpathlineto{\pgfqpoint{1.349731in}{2.304130in}}%
\pgfpathlineto{\pgfqpoint{1.350802in}{2.301599in}}%
\pgfpathlineto{\pgfqpoint{1.351347in}{2.306511in}}%
\pgfpathlineto{\pgfqpoint{1.350812in}{2.301508in}}%
\pgfpathlineto{\pgfqpoint{1.352068in}{2.303484in}}%
\pgfpathlineto{\pgfqpoint{1.352783in}{2.302605in}}%
\pgfpathlineto{\pgfqpoint{1.352545in}{2.304575in}}%
\pgfpathlineto{\pgfqpoint{1.353178in}{2.303155in}}%
\pgfpathlineto{\pgfqpoint{1.353981in}{2.305118in}}%
\pgfpathlineto{\pgfqpoint{1.353630in}{2.302164in}}%
\pgfpathlineto{\pgfqpoint{1.354283in}{2.302957in}}%
\pgfpathlineto{\pgfqpoint{1.354599in}{2.301324in}}%
\pgfpathlineto{\pgfqpoint{1.355295in}{2.303690in}}%
\pgfpathlineto{\pgfqpoint{1.355368in}{2.303639in}}%
\pgfpathlineto{\pgfqpoint{1.355534in}{2.304543in}}%
\pgfpathlineto{\pgfqpoint{1.355889in}{2.301729in}}%
\pgfpathlineto{\pgfqpoint{1.356449in}{2.302329in}}%
\pgfpathlineto{\pgfqpoint{1.356605in}{2.301187in}}%
\pgfpathlineto{\pgfqpoint{1.357126in}{2.305419in}}%
\pgfpathlineto{\pgfqpoint{1.357530in}{2.303301in}}%
\pgfpathlineto{\pgfqpoint{1.357759in}{2.305325in}}%
\pgfpathlineto{\pgfqpoint{1.358528in}{2.302645in}}%
\pgfpathlineto{\pgfqpoint{1.358640in}{2.303463in}}%
\pgfpathlineto{\pgfqpoint{1.359609in}{2.301130in}}%
\pgfpathlineto{\pgfqpoint{1.359200in}{2.304462in}}%
\pgfpathlineto{\pgfqpoint{1.359755in}{2.302829in}}%
\pgfpathlineto{\pgfqpoint{1.360417in}{2.304077in}}%
\pgfpathlineto{\pgfqpoint{1.360096in}{2.301212in}}%
\pgfpathlineto{\pgfqpoint{1.360850in}{2.301791in}}%
\pgfpathlineto{\pgfqpoint{1.361663in}{2.304233in}}%
\pgfpathlineto{\pgfqpoint{1.361011in}{2.301021in}}%
\pgfpathlineto{\pgfqpoint{1.362004in}{2.303540in}}%
\pgfpathlineto{\pgfqpoint{1.362890in}{2.299709in}}%
\pgfpathlineto{\pgfqpoint{1.362345in}{2.304799in}}%
\pgfpathlineto{\pgfqpoint{1.363134in}{2.301818in}}%
\pgfpathlineto{\pgfqpoint{1.364093in}{2.297995in}}%
\pgfpathlineto{\pgfqpoint{1.364424in}{2.300199in}}%
\pgfpathlineto{\pgfqpoint{1.365524in}{2.304305in}}%
\pgfpathlineto{\pgfqpoint{1.364716in}{2.299431in}}%
\pgfpathlineto{\pgfqpoint{1.365578in}{2.303503in}}%
\pgfpathlineto{\pgfqpoint{1.366615in}{2.301675in}}%
\pgfpathlineto{\pgfqpoint{1.365889in}{2.304548in}}%
\pgfpathlineto{\pgfqpoint{1.366688in}{2.303460in}}%
\pgfpathlineto{\pgfqpoint{1.366829in}{2.304044in}}%
\pgfpathlineto{\pgfqpoint{1.367199in}{2.301563in}}%
\pgfpathlineto{\pgfqpoint{1.367730in}{2.301736in}}%
\pgfpathlineto{\pgfqpoint{1.368280in}{2.298846in}}%
\pgfpathlineto{\pgfqpoint{1.368806in}{2.303019in}}%
\pgfpathlineto{\pgfqpoint{1.368810in}{2.302813in}}%
\pgfpathlineto{\pgfqpoint{1.368893in}{2.303628in}}%
\pgfpathlineto{\pgfqpoint{1.369351in}{2.300760in}}%
\pgfpathlineto{\pgfqpoint{1.369828in}{2.301678in}}%
\pgfpathlineto{\pgfqpoint{1.370310in}{2.299920in}}%
\pgfpathlineto{\pgfqpoint{1.370899in}{2.302477in}}%
\pgfpathlineto{\pgfqpoint{1.371108in}{2.303891in}}%
\pgfpathlineto{\pgfqpoint{1.371814in}{2.301312in}}%
\pgfpathlineto{\pgfqpoint{1.371960in}{2.301449in}}%
\pgfpathlineto{\pgfqpoint{1.371985in}{2.301124in}}%
\pgfpathlineto{\pgfqpoint{1.372131in}{2.303665in}}%
\pgfpathlineto{\pgfqpoint{1.372910in}{2.302787in}}%
\pgfpathlineto{\pgfqpoint{1.373805in}{2.304525in}}%
\pgfpathlineto{\pgfqpoint{1.373986in}{2.302136in}}%
\pgfpathlineto{\pgfqpoint{1.374000in}{2.302227in}}%
\pgfpathlineto{\pgfqpoint{1.374916in}{2.300654in}}%
\pgfpathlineto{\pgfqpoint{1.374443in}{2.303507in}}%
\pgfpathlineto{\pgfqpoint{1.375120in}{2.302046in}}%
\pgfpathlineto{\pgfqpoint{1.376181in}{2.304114in}}%
\pgfpathlineto{\pgfqpoint{1.375768in}{2.301418in}}%
\pgfpathlineto{\pgfqpoint{1.376318in}{2.303553in}}%
\pgfpathlineto{\pgfqpoint{1.377525in}{2.298602in}}%
\pgfpathlineto{\pgfqpoint{1.377642in}{2.298783in}}%
\pgfpathlineto{\pgfqpoint{1.377695in}{2.300401in}}%
\pgfpathlineto{\pgfqpoint{1.378323in}{2.297907in}}%
\pgfpathlineto{\pgfqpoint{1.378786in}{2.300085in}}%
\pgfpathlineto{\pgfqpoint{1.379273in}{2.299047in}}%
\pgfpathlineto{\pgfqpoint{1.379701in}{2.303311in}}%
\pgfpathlineto{\pgfqpoint{1.379765in}{2.302811in}}%
\pgfpathlineto{\pgfqpoint{1.379911in}{2.304091in}}%
\pgfpathlineto{\pgfqpoint{1.380782in}{2.299888in}}%
\pgfpathlineto{\pgfqpoint{1.380802in}{2.300462in}}%
\pgfpathlineto{\pgfqpoint{1.380991in}{2.299527in}}%
\pgfpathlineto{\pgfqpoint{1.381853in}{2.302419in}}%
\pgfpathlineto{\pgfqpoint{1.381868in}{2.302221in}}%
\pgfpathlineto{\pgfqpoint{1.382462in}{2.304522in}}%
\pgfpathlineto{\pgfqpoint{1.382953in}{2.301242in}}%
\pgfpathlineto{\pgfqpoint{1.382963in}{2.301117in}}%
\pgfpathlineto{\pgfqpoint{1.383392in}{2.305355in}}%
\pgfpathlineto{\pgfqpoint{1.383961in}{2.303091in}}%
\pgfpathlineto{\pgfqpoint{1.384317in}{2.304966in}}%
\pgfpathlineto{\pgfqpoint{1.384643in}{2.301285in}}%
\pgfpathlineto{\pgfqpoint{1.385057in}{2.302303in}}%
\pgfpathlineto{\pgfqpoint{1.386142in}{2.300769in}}%
\pgfpathlineto{\pgfqpoint{1.385909in}{2.303566in}}%
\pgfpathlineto{\pgfqpoint{1.386201in}{2.300996in}}%
\pgfpathlineto{\pgfqpoint{1.386585in}{2.303562in}}%
\pgfpathlineto{\pgfqpoint{1.387165in}{2.300300in}}%
\pgfpathlineto{\pgfqpoint{1.387316in}{2.301270in}}%
\pgfpathlineto{\pgfqpoint{1.387802in}{2.298984in}}%
\pgfpathlineto{\pgfqpoint{1.388250in}{2.303534in}}%
\pgfpathlineto{\pgfqpoint{1.388377in}{2.302871in}}%
\pgfpathlineto{\pgfqpoint{1.389482in}{2.303621in}}%
\pgfpathlineto{\pgfqpoint{1.388917in}{2.299647in}}%
\pgfpathlineto{\pgfqpoint{1.389492in}{2.303377in}}%
\pgfpathlineto{\pgfqpoint{1.389667in}{2.305452in}}%
\pgfpathlineto{\pgfqpoint{1.390105in}{2.302971in}}%
\pgfpathlineto{\pgfqpoint{1.390689in}{2.304638in}}%
\pgfpathlineto{\pgfqpoint{1.391483in}{2.302812in}}%
\pgfpathlineto{\pgfqpoint{1.391055in}{2.306360in}}%
\pgfpathlineto{\pgfqpoint{1.391824in}{2.303126in}}%
\pgfpathlineto{\pgfqpoint{1.392408in}{2.304054in}}%
\pgfpathlineto{\pgfqpoint{1.392763in}{2.301772in}}%
\pgfpathlineto{\pgfqpoint{1.392929in}{2.302942in}}%
\pgfpathlineto{\pgfqpoint{1.394438in}{2.305033in}}%
\pgfpathlineto{\pgfqpoint{1.392987in}{2.302401in}}%
\pgfpathlineto{\pgfqpoint{1.394536in}{2.303805in}}%
\pgfpathlineto{\pgfqpoint{1.394662in}{2.301721in}}%
\pgfpathlineto{\pgfqpoint{1.395315in}{2.305322in}}%
\pgfpathlineto{\pgfqpoint{1.395646in}{2.303398in}}%
\pgfpathlineto{\pgfqpoint{1.396605in}{2.303854in}}%
\pgfpathlineto{\pgfqpoint{1.396473in}{2.301679in}}%
\pgfpathlineto{\pgfqpoint{1.396707in}{2.302167in}}%
\pgfpathlineto{\pgfqpoint{1.397354in}{2.299835in}}%
\pgfpathlineto{\pgfqpoint{1.397544in}{2.303024in}}%
\pgfpathlineto{\pgfqpoint{1.397836in}{2.301956in}}%
\pgfpathlineto{\pgfqpoint{1.397866in}{2.302750in}}%
\pgfpathlineto{\pgfqpoint{1.398343in}{2.299352in}}%
\pgfpathlineto{\pgfqpoint{1.398912in}{2.301176in}}%
\pgfpathlineto{\pgfqpoint{1.398922in}{2.300987in}}%
\pgfpathlineto{\pgfqpoint{1.399638in}{2.304575in}}%
\pgfpathlineto{\pgfqpoint{1.399954in}{2.303266in}}%
\pgfpathlineto{\pgfqpoint{1.400198in}{2.305449in}}%
\pgfpathlineto{\pgfqpoint{1.400947in}{2.302180in}}%
\pgfpathlineto{\pgfqpoint{1.401045in}{2.303001in}}%
\pgfpathlineto{\pgfqpoint{1.401809in}{2.301375in}}%
\pgfpathlineto{\pgfqpoint{1.402067in}{2.304569in}}%
\pgfpathlineto{\pgfqpoint{1.402276in}{2.305373in}}%
\pgfpathlineto{\pgfqpoint{1.402773in}{2.301885in}}%
\pgfpathlineto{\pgfqpoint{1.403148in}{2.303492in}}%
\pgfpathlineto{\pgfqpoint{1.404053in}{2.300854in}}%
\pgfpathlineto{\pgfqpoint{1.403464in}{2.304291in}}%
\pgfpathlineto{\pgfqpoint{1.404273in}{2.302577in}}%
\pgfpathlineto{\pgfqpoint{1.404404in}{2.303717in}}%
\pgfpathlineto{\pgfqpoint{1.405300in}{2.300907in}}%
\pgfpathlineto{\pgfqpoint{1.405368in}{2.301712in}}%
\pgfpathlineto{\pgfqpoint{1.405407in}{2.301516in}}%
\pgfpathlineto{\pgfqpoint{1.406006in}{2.304689in}}%
\pgfpathlineto{\pgfqpoint{1.406376in}{2.304402in}}%
\pgfpathlineto{\pgfqpoint{1.407179in}{2.304554in}}%
\pgfpathlineto{\pgfqpoint{1.406984in}{2.301694in}}%
\pgfpathlineto{\pgfqpoint{1.407398in}{2.303120in}}%
\pgfpathlineto{\pgfqpoint{1.408255in}{2.301680in}}%
\pgfpathlineto{\pgfqpoint{1.407807in}{2.303899in}}%
\pgfpathlineto{\pgfqpoint{1.408503in}{2.303452in}}%
\pgfpathlineto{\pgfqpoint{1.409735in}{2.301003in}}%
\pgfpathlineto{\pgfqpoint{1.409755in}{2.301278in}}%
\pgfpathlineto{\pgfqpoint{1.410884in}{2.303867in}}%
\pgfpathlineto{\pgfqpoint{1.410183in}{2.299226in}}%
\pgfpathlineto{\pgfqpoint{1.410957in}{2.303467in}}%
\pgfpathlineto{\pgfqpoint{1.411965in}{2.296756in}}%
\pgfpathlineto{\pgfqpoint{1.411152in}{2.303672in}}%
\pgfpathlineto{\pgfqpoint{1.412194in}{2.298690in}}%
\pgfpathlineto{\pgfqpoint{1.412627in}{2.300594in}}%
\pgfpathlineto{\pgfqpoint{1.413060in}{2.297528in}}%
\pgfpathlineto{\pgfqpoint{1.413265in}{2.298325in}}%
\pgfpathlineto{\pgfqpoint{1.413270in}{2.298152in}}%
\pgfpathlineto{\pgfqpoint{1.413615in}{2.301292in}}%
\pgfpathlineto{\pgfqpoint{1.414341in}{2.299368in}}%
\pgfpathlineto{\pgfqpoint{1.415602in}{2.302614in}}%
\pgfpathlineto{\pgfqpoint{1.414827in}{2.297632in}}%
\pgfpathlineto{\pgfqpoint{1.415606in}{2.302480in}}%
\pgfpathlineto{\pgfqpoint{1.416205in}{2.298329in}}%
\pgfpathlineto{\pgfqpoint{1.416770in}{2.298738in}}%
\pgfpathlineto{\pgfqpoint{1.416926in}{2.297580in}}%
\pgfpathlineto{\pgfqpoint{1.417763in}{2.300270in}}%
\pgfpathlineto{\pgfqpoint{1.417846in}{2.299609in}}%
\pgfpathlineto{\pgfqpoint{1.417890in}{2.300292in}}%
\pgfpathlineto{\pgfqpoint{1.418420in}{2.296021in}}%
\pgfpathlineto{\pgfqpoint{1.418508in}{2.296480in}}%
\pgfpathlineto{\pgfqpoint{1.418513in}{2.296439in}}%
\pgfpathlineto{\pgfqpoint{1.419044in}{2.299012in}}%
\pgfpathlineto{\pgfqpoint{1.419326in}{2.298701in}}%
\pgfpathlineto{\pgfqpoint{1.419414in}{2.299993in}}%
\pgfpathlineto{\pgfqpoint{1.420090in}{2.297657in}}%
\pgfpathlineto{\pgfqpoint{1.420431in}{2.298439in}}%
\pgfpathlineto{\pgfqpoint{1.420636in}{2.297480in}}%
\pgfpathlineto{\pgfqpoint{1.421303in}{2.300094in}}%
\pgfpathlineto{\pgfqpoint{1.421449in}{2.299373in}}%
\pgfpathlineto{\pgfqpoint{1.422466in}{2.302086in}}%
\pgfpathlineto{\pgfqpoint{1.421687in}{2.298564in}}%
\pgfpathlineto{\pgfqpoint{1.422602in}{2.300793in}}%
\pgfpathlineto{\pgfqpoint{1.423474in}{2.303191in}}%
\pgfpathlineto{\pgfqpoint{1.422972in}{2.300095in}}%
\pgfpathlineto{\pgfqpoint{1.423698in}{2.300599in}}%
\pgfpathlineto{\pgfqpoint{1.424521in}{2.296256in}}%
\pgfpathlineto{\pgfqpoint{1.423902in}{2.300733in}}%
\pgfpathlineto{\pgfqpoint{1.424852in}{2.298911in}}%
\pgfpathlineto{\pgfqpoint{1.425441in}{2.295893in}}%
\pgfpathlineto{\pgfqpoint{1.424998in}{2.299612in}}%
\pgfpathlineto{\pgfqpoint{1.426045in}{2.298005in}}%
\pgfpathlineto{\pgfqpoint{1.426142in}{2.296185in}}%
\pgfpathlineto{\pgfqpoint{1.426697in}{2.299632in}}%
\pgfpathlineto{\pgfqpoint{1.427203in}{2.304535in}}%
\pgfpathlineto{\pgfqpoint{1.427851in}{2.303458in}}%
\pgfpathlineto{\pgfqpoint{1.428907in}{2.302061in}}%
\pgfpathlineto{\pgfqpoint{1.428162in}{2.304713in}}%
\pgfpathlineto{\pgfqpoint{1.428966in}{2.302990in}}%
\pgfpathlineto{\pgfqpoint{1.429194in}{2.303831in}}%
\pgfpathlineto{\pgfqpoint{1.429871in}{2.301116in}}%
\pgfpathlineto{\pgfqpoint{1.430037in}{2.302429in}}%
\pgfpathlineto{\pgfqpoint{1.430056in}{2.301933in}}%
\pgfpathlineto{\pgfqpoint{1.430412in}{2.304270in}}%
\pgfpathlineto{\pgfqpoint{1.431132in}{2.303263in}}%
\pgfpathlineto{\pgfqpoint{1.432320in}{2.299533in}}%
\pgfpathlineto{\pgfqpoint{1.431161in}{2.303404in}}%
\pgfpathlineto{\pgfqpoint{1.432354in}{2.299804in}}%
\pgfpathlineto{\pgfqpoint{1.433157in}{2.303151in}}%
\pgfpathlineto{\pgfqpoint{1.433523in}{2.301216in}}%
\pgfpathlineto{\pgfqpoint{1.434082in}{2.300185in}}%
\pgfpathlineto{\pgfqpoint{1.434467in}{2.302038in}}%
\pgfpathlineto{\pgfqpoint{1.434628in}{2.301321in}}%
\pgfpathlineto{\pgfqpoint{1.434905in}{2.303209in}}%
\pgfpathlineto{\pgfqpoint{1.435699in}{2.299990in}}%
\pgfpathlineto{\pgfqpoint{1.435718in}{2.299447in}}%
\pgfpathlineto{\pgfqpoint{1.436030in}{2.301906in}}%
\pgfpathlineto{\pgfqpoint{1.436040in}{2.301897in}}%
\pgfpathlineto{\pgfqpoint{1.436190in}{2.302947in}}%
\pgfpathlineto{\pgfqpoint{1.436814in}{2.299471in}}%
\pgfpathlineto{\pgfqpoint{1.437096in}{2.300102in}}%
\pgfpathlineto{\pgfqpoint{1.438187in}{2.296806in}}%
\pgfpathlineto{\pgfqpoint{1.437627in}{2.300971in}}%
\pgfpathlineto{\pgfqpoint{1.438323in}{2.298332in}}%
\pgfpathlineto{\pgfqpoint{1.438897in}{2.300053in}}%
\pgfpathlineto{\pgfqpoint{1.439331in}{2.297278in}}%
\pgfpathlineto{\pgfqpoint{1.439413in}{2.297928in}}%
\pgfpathlineto{\pgfqpoint{1.439433in}{2.297169in}}%
\pgfpathlineto{\pgfqpoint{1.440455in}{2.302854in}}%
\pgfpathlineto{\pgfqpoint{1.440460in}{2.302837in}}%
\pgfpathlineto{\pgfqpoint{1.441434in}{2.306172in}}%
\pgfpathlineto{\pgfqpoint{1.440562in}{2.302012in}}%
\pgfpathlineto{\pgfqpoint{1.441643in}{2.304467in}}%
\pgfpathlineto{\pgfqpoint{1.442554in}{2.301540in}}%
\pgfpathlineto{\pgfqpoint{1.441687in}{2.305210in}}%
\pgfpathlineto{\pgfqpoint{1.442787in}{2.303655in}}%
\pgfpathlineto{\pgfqpoint{1.442846in}{2.304547in}}%
\pgfpathlineto{\pgfqpoint{1.443810in}{2.300017in}}%
\pgfpathlineto{\pgfqpoint{1.443824in}{2.300434in}}%
\pgfpathlineto{\pgfqpoint{1.443912in}{2.301543in}}%
\pgfpathlineto{\pgfqpoint{1.444555in}{2.297108in}}%
\pgfpathlineto{\pgfqpoint{1.444564in}{2.297253in}}%
\pgfpathlineto{\pgfqpoint{1.445149in}{2.295608in}}%
\pgfpathlineto{\pgfqpoint{1.445577in}{2.300282in}}%
\pgfpathlineto{\pgfqpoint{1.446478in}{2.302048in}}%
\pgfpathlineto{\pgfqpoint{1.446176in}{2.299959in}}%
\pgfpathlineto{\pgfqpoint{1.446692in}{2.300506in}}%
\pgfpathlineto{\pgfqpoint{1.447072in}{2.296704in}}%
\pgfpathlineto{\pgfqpoint{1.447914in}{2.298431in}}%
\pgfpathlineto{\pgfqpoint{1.449019in}{2.303598in}}%
\pgfpathlineto{\pgfqpoint{1.449126in}{2.303532in}}%
\pgfpathlineto{\pgfqpoint{1.449209in}{2.304197in}}%
\pgfpathlineto{\pgfqpoint{1.449968in}{2.300857in}}%
\pgfpathlineto{\pgfqpoint{1.449978in}{2.300928in}}%
\pgfpathlineto{\pgfqpoint{1.450144in}{2.299223in}}%
\pgfpathlineto{\pgfqpoint{1.450879in}{2.304331in}}%
\pgfpathlineto{\pgfqpoint{1.451030in}{2.303001in}}%
\pgfpathlineto{\pgfqpoint{1.451638in}{2.305149in}}%
\pgfpathlineto{\pgfqpoint{1.451220in}{2.302288in}}%
\pgfpathlineto{\pgfqpoint{1.452145in}{2.303424in}}%
\pgfpathlineto{\pgfqpoint{1.453099in}{2.302607in}}%
\pgfpathlineto{\pgfqpoint{1.452515in}{2.305731in}}%
\pgfpathlineto{\pgfqpoint{1.453255in}{2.303465in}}%
\pgfpathlineto{\pgfqpoint{1.454384in}{2.301147in}}%
\pgfpathlineto{\pgfqpoint{1.453371in}{2.304847in}}%
\pgfpathlineto{\pgfqpoint{1.454467in}{2.301378in}}%
\pgfpathlineto{\pgfqpoint{1.454813in}{2.303995in}}%
\pgfpathlineto{\pgfqpoint{1.455611in}{2.302833in}}%
\pgfpathlineto{\pgfqpoint{1.455776in}{2.301205in}}%
\pgfpathlineto{\pgfqpoint{1.456697in}{2.303738in}}%
\pgfpathlineto{\pgfqpoint{1.456745in}{2.304277in}}%
\pgfpathlineto{\pgfqpoint{1.457432in}{2.301237in}}%
\pgfpathlineto{\pgfqpoint{1.457743in}{2.301406in}}%
\pgfpathlineto{\pgfqpoint{1.457812in}{2.301936in}}%
\pgfpathlineto{\pgfqpoint{1.457777in}{2.302066in}}%
\pgfpathlineto{\pgfqpoint{1.457821in}{2.302341in}}%
\pgfpathlineto{\pgfqpoint{1.458031in}{2.304004in}}%
\pgfpathlineto{\pgfqpoint{1.458186in}{2.302117in}}%
\pgfpathlineto{\pgfqpoint{1.458936in}{2.303052in}}%
\pgfpathlineto{\pgfqpoint{1.459350in}{2.305889in}}%
\pgfpathlineto{\pgfqpoint{1.459618in}{2.302461in}}%
\pgfpathlineto{\pgfqpoint{1.460144in}{2.303739in}}%
\pgfpathlineto{\pgfqpoint{1.460192in}{2.302620in}}%
\pgfpathlineto{\pgfqpoint{1.460713in}{2.304982in}}%
\pgfpathlineto{\pgfqpoint{1.461229in}{2.304520in}}%
\pgfpathlineto{\pgfqpoint{1.461341in}{2.305024in}}%
\pgfpathlineto{\pgfqpoint{1.461580in}{2.302396in}}%
\pgfpathlineto{\pgfqpoint{1.462549in}{2.298144in}}%
\pgfpathlineto{\pgfqpoint{1.462096in}{2.304210in}}%
\pgfpathlineto{\pgfqpoint{1.462875in}{2.299923in}}%
\pgfpathlineto{\pgfqpoint{1.463732in}{2.301076in}}%
\pgfpathlineto{\pgfqpoint{1.463313in}{2.297979in}}%
\pgfpathlineto{\pgfqpoint{1.463985in}{2.300188in}}%
\pgfpathlineto{\pgfqpoint{1.464345in}{2.297901in}}%
\pgfpathlineto{\pgfqpoint{1.465075in}{2.301425in}}%
\pgfpathlineto{\pgfqpoint{1.465080in}{2.301086in}}%
\pgfpathlineto{\pgfqpoint{1.465625in}{2.303620in}}%
\pgfpathlineto{\pgfqpoint{1.465255in}{2.299921in}}%
\pgfpathlineto{\pgfqpoint{1.466205in}{2.301828in}}%
\pgfpathlineto{\pgfqpoint{1.466550in}{2.303139in}}%
\pgfpathlineto{\pgfqpoint{1.467115in}{2.301046in}}%
\pgfpathlineto{\pgfqpoint{1.467266in}{2.301604in}}%
\pgfpathlineto{\pgfqpoint{1.467602in}{2.302503in}}%
\pgfpathlineto{\pgfqpoint{1.468483in}{2.297260in}}%
\pgfpathlineto{\pgfqpoint{1.468537in}{2.297931in}}%
\pgfpathlineto{\pgfqpoint{1.469452in}{2.293035in}}%
\pgfpathlineto{\pgfqpoint{1.469467in}{2.293250in}}%
\pgfpathlineto{\pgfqpoint{1.469754in}{2.290597in}}%
\pgfpathlineto{\pgfqpoint{1.470484in}{2.295386in}}%
\pgfpathlineto{\pgfqpoint{1.470513in}{2.295159in}}%
\pgfpathlineto{\pgfqpoint{1.471141in}{2.295926in}}%
\pgfpathlineto{\pgfqpoint{1.470893in}{2.293740in}}%
\pgfpathlineto{\pgfqpoint{1.471594in}{2.294910in}}%
\pgfpathlineto{\pgfqpoint{1.471609in}{2.294056in}}%
\pgfpathlineto{\pgfqpoint{1.472612in}{2.296762in}}%
\pgfpathlineto{\pgfqpoint{1.472690in}{2.295898in}}%
\pgfpathlineto{\pgfqpoint{1.473576in}{2.298720in}}%
\pgfpathlineto{\pgfqpoint{1.472821in}{2.295773in}}%
\pgfpathlineto{\pgfqpoint{1.473882in}{2.297218in}}%
\pgfpathlineto{\pgfqpoint{1.474321in}{2.295990in}}%
\pgfpathlineto{\pgfqpoint{1.474661in}{2.300145in}}%
\pgfpathlineto{\pgfqpoint{1.474963in}{2.298074in}}%
\pgfpathlineto{\pgfqpoint{1.475085in}{2.298108in}}%
\pgfpathlineto{\pgfqpoint{1.475148in}{2.299553in}}%
\pgfpathlineto{\pgfqpoint{1.475659in}{2.301147in}}%
\pgfpathlineto{\pgfqpoint{1.475888in}{2.297213in}}%
\pgfpathlineto{\pgfqpoint{1.476268in}{2.299943in}}%
\pgfpathlineto{\pgfqpoint{1.476429in}{2.298405in}}%
\pgfpathlineto{\pgfqpoint{1.477271in}{2.301803in}}%
\pgfpathlineto{\pgfqpoint{1.477938in}{2.302498in}}%
\pgfpathlineto{\pgfqpoint{1.477597in}{2.300358in}}%
\pgfpathlineto{\pgfqpoint{1.478381in}{2.301794in}}%
\pgfpathlineto{\pgfqpoint{1.478503in}{2.302122in}}%
\pgfpathlineto{\pgfqpoint{1.478863in}{2.298370in}}%
\pgfpathlineto{\pgfqpoint{1.479053in}{2.298600in}}%
\pgfpathlineto{\pgfqpoint{1.480104in}{2.293661in}}%
\pgfpathlineto{\pgfqpoint{1.480280in}{2.293761in}}%
\pgfpathlineto{\pgfqpoint{1.480460in}{2.294260in}}%
\pgfpathlineto{\pgfqpoint{1.481044in}{2.290191in}}%
\pgfpathlineto{\pgfqpoint{1.481331in}{2.292194in}}%
\pgfpathlineto{\pgfqpoint{1.482543in}{2.287505in}}%
\pgfpathlineto{\pgfqpoint{1.481662in}{2.293774in}}%
\pgfpathlineto{\pgfqpoint{1.482690in}{2.288476in}}%
\pgfpathlineto{\pgfqpoint{1.482719in}{2.288072in}}%
\pgfpathlineto{\pgfqpoint{1.483347in}{2.291169in}}%
\pgfpathlineto{\pgfqpoint{1.483512in}{2.292484in}}%
\pgfpathlineto{\pgfqpoint{1.484374in}{2.288132in}}%
\pgfpathlineto{\pgfqpoint{1.485426in}{2.284951in}}%
\pgfpathlineto{\pgfqpoint{1.484554in}{2.288502in}}%
\pgfpathlineto{\pgfqpoint{1.485513in}{2.285836in}}%
\pgfpathlineto{\pgfqpoint{1.485528in}{2.286236in}}%
\pgfpathlineto{\pgfqpoint{1.486283in}{2.281775in}}%
\pgfpathlineto{\pgfqpoint{1.486541in}{2.283756in}}%
\pgfpathlineto{\pgfqpoint{1.487675in}{2.281215in}}%
\pgfpathlineto{\pgfqpoint{1.486988in}{2.285771in}}%
\pgfpathlineto{\pgfqpoint{1.487748in}{2.282248in}}%
\pgfpathlineto{\pgfqpoint{1.487870in}{2.283499in}}%
\pgfpathlineto{\pgfqpoint{1.488425in}{2.279524in}}%
\pgfpathlineto{\pgfqpoint{1.488868in}{2.282906in}}%
\pgfpathlineto{\pgfqpoint{1.489652in}{2.281432in}}%
\pgfpathlineto{\pgfqpoint{1.489549in}{2.283721in}}%
\pgfpathlineto{\pgfqpoint{1.489997in}{2.282636in}}%
\pgfpathlineto{\pgfqpoint{1.491010in}{2.285584in}}%
\pgfpathlineto{\pgfqpoint{1.490289in}{2.281210in}}%
\pgfpathlineto{\pgfqpoint{1.491146in}{2.284536in}}%
\pgfpathlineto{\pgfqpoint{1.491273in}{2.282975in}}%
\pgfpathlineto{\pgfqpoint{1.491701in}{2.286000in}}%
\pgfpathlineto{\pgfqpoint{1.492251in}{2.284848in}}%
\pgfpathlineto{\pgfqpoint{1.493400in}{2.288690in}}%
\pgfpathlineto{\pgfqpoint{1.492592in}{2.283572in}}%
\pgfpathlineto{\pgfqpoint{1.493468in}{2.287887in}}%
\pgfpathlineto{\pgfqpoint{1.493673in}{2.285808in}}%
\pgfpathlineto{\pgfqpoint{1.493980in}{2.288183in}}%
\pgfpathlineto{\pgfqpoint{1.494598in}{2.286680in}}%
\pgfpathlineto{\pgfqpoint{1.495328in}{2.290165in}}%
\pgfpathlineto{\pgfqpoint{1.494734in}{2.285302in}}%
\pgfpathlineto{\pgfqpoint{1.495776in}{2.289578in}}%
\pgfpathlineto{\pgfqpoint{1.495815in}{2.290016in}}%
\pgfpathlineto{\pgfqpoint{1.496107in}{2.286651in}}%
\pgfpathlineto{\pgfqpoint{1.496336in}{2.286807in}}%
\pgfpathlineto{\pgfqpoint{1.497242in}{2.284140in}}%
\pgfpathlineto{\pgfqpoint{1.496760in}{2.288061in}}%
\pgfpathlineto{\pgfqpoint{1.497470in}{2.285955in}}%
\pgfpathlineto{\pgfqpoint{1.498707in}{2.289426in}}%
\pgfpathlineto{\pgfqpoint{1.498016in}{2.285479in}}%
\pgfpathlineto{\pgfqpoint{1.498877in}{2.289026in}}%
\pgfpathlineto{\pgfqpoint{1.499004in}{2.288389in}}%
\pgfpathlineto{\pgfqpoint{1.499866in}{2.293595in}}%
\pgfpathlineto{\pgfqpoint{1.499934in}{2.291705in}}%
\pgfpathlineto{\pgfqpoint{1.500352in}{2.294508in}}%
\pgfpathlineto{\pgfqpoint{1.500883in}{2.294093in}}%
\pgfpathlineto{\pgfqpoint{1.502052in}{2.299501in}}%
\pgfpathlineto{\pgfqpoint{1.500927in}{2.294017in}}%
\pgfpathlineto{\pgfqpoint{1.502100in}{2.298844in}}%
\pgfpathlineto{\pgfqpoint{1.502534in}{2.297444in}}%
\pgfpathlineto{\pgfqpoint{1.502300in}{2.300543in}}%
\pgfpathlineto{\pgfqpoint{1.503157in}{2.300259in}}%
\pgfpathlineto{\pgfqpoint{1.504096in}{2.302115in}}%
\pgfpathlineto{\pgfqpoint{1.503697in}{2.298362in}}%
\pgfpathlineto{\pgfqpoint{1.504286in}{2.300947in}}%
\pgfpathlineto{\pgfqpoint{1.504379in}{2.300225in}}%
\pgfpathlineto{\pgfqpoint{1.504612in}{2.301992in}}%
\pgfpathlineto{\pgfqpoint{1.504627in}{2.301943in}}%
\pgfpathlineto{\pgfqpoint{1.504851in}{2.302534in}}%
\pgfpathlineto{\pgfqpoint{1.505309in}{2.300128in}}%
\pgfpathlineto{\pgfqpoint{1.505718in}{2.300920in}}%
\pgfpathlineto{\pgfqpoint{1.505766in}{2.300178in}}%
\pgfpathlineto{\pgfqpoint{1.506088in}{2.302362in}}%
\pgfpathlineto{\pgfqpoint{1.506808in}{2.301551in}}%
\pgfpathlineto{\pgfqpoint{1.507129in}{2.301845in}}%
\pgfpathlineto{\pgfqpoint{1.507821in}{2.298510in}}%
\pgfpathlineto{\pgfqpoint{1.507860in}{2.298994in}}%
\pgfpathlineto{\pgfqpoint{1.508415in}{2.296860in}}%
\pgfpathlineto{\pgfqpoint{1.508050in}{2.300081in}}%
\pgfpathlineto{\pgfqpoint{1.508994in}{2.297975in}}%
\pgfpathlineto{\pgfqpoint{1.510211in}{2.300626in}}%
\pgfpathlineto{\pgfqpoint{1.509072in}{2.297591in}}%
\pgfpathlineto{\pgfqpoint{1.510236in}{2.300211in}}%
\pgfpathlineto{\pgfqpoint{1.510240in}{2.299978in}}%
\pgfpathlineto{\pgfqpoint{1.510995in}{2.304032in}}%
\pgfpathlineto{\pgfqpoint{1.511287in}{2.303014in}}%
\pgfpathlineto{\pgfqpoint{1.511448in}{2.304038in}}%
\pgfpathlineto{\pgfqpoint{1.512275in}{2.299492in}}%
\pgfpathlineto{\pgfqpoint{1.512324in}{2.300160in}}%
\pgfpathlineto{\pgfqpoint{1.512826in}{2.298465in}}%
\pgfpathlineto{\pgfqpoint{1.513337in}{2.302961in}}%
\pgfpathlineto{\pgfqpoint{1.513361in}{2.302812in}}%
\pgfpathlineto{\pgfqpoint{1.514447in}{2.303829in}}%
\pgfpathlineto{\pgfqpoint{1.513751in}{2.300800in}}%
\pgfpathlineto{\pgfqpoint{1.514476in}{2.302962in}}%
\pgfpathlineto{\pgfqpoint{1.515581in}{2.301316in}}%
\pgfpathlineto{\pgfqpoint{1.514661in}{2.304691in}}%
\pgfpathlineto{\pgfqpoint{1.515601in}{2.302142in}}%
\pgfpathlineto{\pgfqpoint{1.515820in}{2.299776in}}%
\pgfpathlineto{\pgfqpoint{1.516331in}{2.302600in}}%
\pgfpathlineto{\pgfqpoint{1.516706in}{2.301805in}}%
\pgfpathlineto{\pgfqpoint{1.516774in}{2.303361in}}%
\pgfpathlineto{\pgfqpoint{1.517543in}{2.299841in}}%
\pgfpathlineto{\pgfqpoint{1.517806in}{2.300977in}}%
\pgfpathlineto{\pgfqpoint{1.517845in}{2.302033in}}%
\pgfpathlineto{\pgfqpoint{1.517850in}{2.301954in}}%
\pgfpathlineto{\pgfqpoint{1.518848in}{2.305282in}}%
\pgfpathlineto{\pgfqpoint{1.518059in}{2.301272in}}%
\pgfpathlineto{\pgfqpoint{1.519004in}{2.304596in}}%
\pgfpathlineto{\pgfqpoint{1.519013in}{2.304686in}}%
\pgfpathlineto{\pgfqpoint{1.519680in}{2.300300in}}%
\pgfpathlineto{\pgfqpoint{1.519934in}{2.301929in}}%
\pgfpathlineto{\pgfqpoint{1.520435in}{2.301371in}}%
\pgfpathlineto{\pgfqpoint{1.520109in}{2.303688in}}%
\pgfpathlineto{\pgfqpoint{1.521000in}{2.302862in}}%
\pgfpathlineto{\pgfqpoint{1.522105in}{2.304553in}}%
\pgfpathlineto{\pgfqpoint{1.521535in}{2.301941in}}%
\pgfpathlineto{\pgfqpoint{1.522124in}{2.304245in}}%
\pgfpathlineto{\pgfqpoint{1.523181in}{2.301561in}}%
\pgfpathlineto{\pgfqpoint{1.522290in}{2.304292in}}%
\pgfpathlineto{\pgfqpoint{1.523244in}{2.303845in}}%
\pgfpathlineto{\pgfqpoint{1.524121in}{2.302178in}}%
\pgfpathlineto{\pgfqpoint{1.523965in}{2.304701in}}%
\pgfpathlineto{\pgfqpoint{1.524379in}{2.303413in}}%
\pgfpathlineto{\pgfqpoint{1.524432in}{2.303844in}}%
\pgfpathlineto{\pgfqpoint{1.525396in}{2.299772in}}%
\pgfpathlineto{\pgfqpoint{1.525401in}{2.299789in}}%
\pgfpathlineto{\pgfqpoint{1.525961in}{2.297707in}}%
\pgfpathlineto{\pgfqpoint{1.525450in}{2.300293in}}%
\pgfpathlineto{\pgfqpoint{1.526530in}{2.298864in}}%
\pgfpathlineto{\pgfqpoint{1.527757in}{2.305503in}}%
\pgfpathlineto{\pgfqpoint{1.526545in}{2.298662in}}%
\pgfpathlineto{\pgfqpoint{1.528867in}{2.302142in}}%
\pgfpathlineto{\pgfqpoint{1.528921in}{2.301497in}}%
\pgfpathlineto{\pgfqpoint{1.529413in}{2.304461in}}%
\pgfpathlineto{\pgfqpoint{1.529914in}{2.303413in}}%
\pgfpathlineto{\pgfqpoint{1.529938in}{2.303619in}}%
\pgfpathlineto{\pgfqpoint{1.530109in}{2.300574in}}%
\pgfpathlineto{\pgfqpoint{1.530966in}{2.302026in}}%
\pgfpathlineto{\pgfqpoint{1.531102in}{2.301496in}}%
\pgfpathlineto{\pgfqpoint{1.531959in}{2.304109in}}%
\pgfpathlineto{\pgfqpoint{1.533088in}{2.306407in}}%
\pgfpathlineto{\pgfqpoint{1.532339in}{2.302613in}}%
\pgfpathlineto{\pgfqpoint{1.533108in}{2.305943in}}%
\pgfpathlineto{\pgfqpoint{1.534208in}{2.301014in}}%
\pgfpathlineto{\pgfqpoint{1.533171in}{2.306936in}}%
\pgfpathlineto{\pgfqpoint{1.534792in}{2.303177in}}%
\pgfpathlineto{\pgfqpoint{1.535440in}{2.304234in}}%
\pgfpathlineto{\pgfqpoint{1.535834in}{2.300077in}}%
\pgfpathlineto{\pgfqpoint{1.535868in}{2.300520in}}%
\pgfpathlineto{\pgfqpoint{1.536340in}{2.299275in}}%
\pgfpathlineto{\pgfqpoint{1.536628in}{2.302123in}}%
\pgfpathlineto{\pgfqpoint{1.537003in}{2.299884in}}%
\pgfpathlineto{\pgfqpoint{1.537631in}{2.302467in}}%
\pgfpathlineto{\pgfqpoint{1.537032in}{2.299511in}}%
\pgfpathlineto{\pgfqpoint{1.538117in}{2.300239in}}%
\pgfpathlineto{\pgfqpoint{1.538512in}{2.299497in}}%
\pgfpathlineto{\pgfqpoint{1.539159in}{2.302302in}}%
\pgfpathlineto{\pgfqpoint{1.539174in}{2.302023in}}%
\pgfpathlineto{\pgfqpoint{1.541389in}{2.305414in}}%
\pgfpathlineto{\pgfqpoint{1.539427in}{2.301582in}}%
\pgfpathlineto{\pgfqpoint{1.541589in}{2.304100in}}%
\pgfpathlineto{\pgfqpoint{1.542543in}{2.300428in}}%
\pgfpathlineto{\pgfqpoint{1.542066in}{2.304286in}}%
\pgfpathlineto{\pgfqpoint{1.542728in}{2.301927in}}%
\pgfpathlineto{\pgfqpoint{1.542966in}{2.304404in}}%
\pgfpathlineto{\pgfqpoint{1.543872in}{2.303847in}}%
\pgfpathlineto{\pgfqpoint{1.544807in}{2.305946in}}%
\pgfpathlineto{\pgfqpoint{1.544266in}{2.303306in}}%
\pgfpathlineto{\pgfqpoint{1.544967in}{2.303582in}}%
\pgfpathlineto{\pgfqpoint{1.545605in}{2.301543in}}%
\pgfpathlineto{\pgfqpoint{1.545294in}{2.303957in}}%
\pgfpathlineto{\pgfqpoint{1.546165in}{2.302860in}}%
\pgfpathlineto{\pgfqpoint{1.547246in}{2.305257in}}%
\pgfpathlineto{\pgfqpoint{1.546413in}{2.301748in}}%
\pgfpathlineto{\pgfqpoint{1.547319in}{2.304667in}}%
\pgfpathlineto{\pgfqpoint{1.548439in}{2.300296in}}%
\pgfpathlineto{\pgfqpoint{1.547898in}{2.305482in}}%
\pgfpathlineto{\pgfqpoint{1.548526in}{2.301120in}}%
\pgfpathlineto{\pgfqpoint{1.549656in}{2.306394in}}%
\pgfpathlineto{\pgfqpoint{1.548823in}{2.299998in}}%
\pgfpathlineto{\pgfqpoint{1.549768in}{2.304774in}}%
\pgfpathlineto{\pgfqpoint{1.550766in}{2.301448in}}%
\pgfpathlineto{\pgfqpoint{1.550975in}{2.302649in}}%
\pgfpathlineto{\pgfqpoint{1.551379in}{2.301882in}}%
\pgfpathlineto{\pgfqpoint{1.551998in}{2.303556in}}%
\pgfpathlineto{\pgfqpoint{1.552041in}{2.303094in}}%
\pgfpathlineto{\pgfqpoint{1.552295in}{2.305623in}}%
\pgfpathlineto{\pgfqpoint{1.553010in}{2.299508in}}%
\pgfpathlineto{\pgfqpoint{1.553122in}{2.300700in}}%
\pgfpathlineto{\pgfqpoint{1.554359in}{2.305672in}}%
\pgfpathlineto{\pgfqpoint{1.554364in}{2.305310in}}%
\pgfpathlineto{\pgfqpoint{1.555206in}{2.301780in}}%
\pgfpathlineto{\pgfqpoint{1.555493in}{2.303903in}}%
\pgfpathlineto{\pgfqpoint{1.556545in}{2.300252in}}%
\pgfpathlineto{\pgfqpoint{1.555615in}{2.304382in}}%
\pgfpathlineto{\pgfqpoint{1.557290in}{2.302999in}}%
\pgfpathlineto{\pgfqpoint{1.557664in}{2.303966in}}%
\pgfpathlineto{\pgfqpoint{1.558249in}{2.301334in}}%
\pgfpathlineto{\pgfqpoint{1.558400in}{2.302989in}}%
\pgfpathlineto{\pgfqpoint{1.558731in}{2.300709in}}%
\pgfpathlineto{\pgfqpoint{1.559417in}{2.303269in}}%
\pgfpathlineto{\pgfqpoint{1.559529in}{2.302369in}}%
\pgfpathlineto{\pgfqpoint{1.559914in}{2.304691in}}%
\pgfpathlineto{\pgfqpoint{1.559607in}{2.302027in}}%
\pgfpathlineto{\pgfqpoint{1.560615in}{2.302495in}}%
\pgfpathlineto{\pgfqpoint{1.561627in}{2.301360in}}%
\pgfpathlineto{\pgfqpoint{1.561024in}{2.304216in}}%
\pgfpathlineto{\pgfqpoint{1.561725in}{2.302390in}}%
\pgfpathlineto{\pgfqpoint{1.562246in}{2.299732in}}%
\pgfpathlineto{\pgfqpoint{1.561924in}{2.302552in}}%
\pgfpathlineto{\pgfqpoint{1.562859in}{2.301633in}}%
\pgfpathlineto{\pgfqpoint{1.563964in}{2.302920in}}%
\pgfpathlineto{\pgfqpoint{1.563531in}{2.300319in}}%
\pgfpathlineto{\pgfqpoint{1.563984in}{2.302550in}}%
\pgfpathlineto{\pgfqpoint{1.564081in}{2.302267in}}%
\pgfpathlineto{\pgfqpoint{1.564826in}{2.305565in}}%
\pgfpathlineto{\pgfqpoint{1.565001in}{2.305301in}}%
\pgfpathlineto{\pgfqpoint{1.565693in}{2.305616in}}%
\pgfpathlineto{\pgfqpoint{1.565391in}{2.302826in}}%
\pgfpathlineto{\pgfqpoint{1.566111in}{2.305289in}}%
\pgfpathlineto{\pgfqpoint{1.566496in}{2.301447in}}%
\pgfpathlineto{\pgfqpoint{1.566199in}{2.305381in}}%
\pgfpathlineto{\pgfqpoint{1.567275in}{2.302371in}}%
\pgfpathlineto{\pgfqpoint{1.567703in}{2.302771in}}%
\pgfpathlineto{\pgfqpoint{1.567489in}{2.300739in}}%
\pgfpathlineto{\pgfqpoint{1.568137in}{2.301534in}}%
\pgfpathlineto{\pgfqpoint{1.568365in}{2.300126in}}%
\pgfpathlineto{\pgfqpoint{1.569140in}{2.304159in}}%
\pgfpathlineto{\pgfqpoint{1.569534in}{2.305151in}}%
\pgfpathlineto{\pgfqpoint{1.570186in}{2.302195in}}%
\pgfpathlineto{\pgfqpoint{1.570225in}{2.303626in}}%
\pgfpathlineto{\pgfqpoint{1.571189in}{2.301219in}}%
\pgfpathlineto{\pgfqpoint{1.570430in}{2.304636in}}%
\pgfpathlineto{\pgfqpoint{1.571355in}{2.302621in}}%
\pgfpathlineto{\pgfqpoint{1.571866in}{2.304213in}}%
\pgfpathlineto{\pgfqpoint{1.572114in}{2.302285in}}%
\pgfpathlineto{\pgfqpoint{1.572479in}{2.303796in}}%
\pgfpathlineto{\pgfqpoint{1.572552in}{2.304942in}}%
\pgfpathlineto{\pgfqpoint{1.573555in}{2.302630in}}%
\pgfpathlineto{\pgfqpoint{1.573589in}{2.303815in}}%
\pgfpathlineto{\pgfqpoint{1.573628in}{2.304353in}}%
\pgfpathlineto{\pgfqpoint{1.574300in}{2.298692in}}%
\pgfpathlineto{\pgfqpoint{1.574617in}{2.300739in}}%
\pgfpathlineto{\pgfqpoint{1.575118in}{2.296907in}}%
\pgfpathlineto{\pgfqpoint{1.574660in}{2.301211in}}%
\pgfpathlineto{\pgfqpoint{1.575941in}{2.298661in}}%
\pgfpathlineto{\pgfqpoint{1.577216in}{2.301482in}}%
\pgfpathlineto{\pgfqpoint{1.576676in}{2.297524in}}%
\pgfpathlineto{\pgfqpoint{1.577333in}{2.300976in}}%
\pgfpathlineto{\pgfqpoint{1.577552in}{2.299101in}}%
\pgfpathlineto{\pgfqpoint{1.578175in}{2.303763in}}%
\pgfpathlineto{\pgfqpoint{1.578356in}{2.303134in}}%
\pgfpathlineto{\pgfqpoint{1.578365in}{2.303317in}}%
\pgfpathlineto{\pgfqpoint{1.579135in}{2.297235in}}%
\pgfpathlineto{\pgfqpoint{1.579359in}{2.299510in}}%
\pgfpathlineto{\pgfqpoint{1.579914in}{2.297616in}}%
\pgfpathlineto{\pgfqpoint{1.579767in}{2.299773in}}%
\pgfpathlineto{\pgfqpoint{1.580576in}{2.297837in}}%
\pgfpathlineto{\pgfqpoint{1.580590in}{2.298558in}}%
\pgfpathlineto{\pgfqpoint{1.581335in}{2.294761in}}%
\pgfpathlineto{\pgfqpoint{1.581671in}{2.296968in}}%
\pgfpathlineto{\pgfqpoint{1.581900in}{2.296680in}}%
\pgfpathlineto{\pgfqpoint{1.582543in}{2.299607in}}%
\pgfpathlineto{\pgfqpoint{1.583005in}{2.301619in}}%
\pgfpathlineto{\pgfqpoint{1.582650in}{2.298936in}}%
\pgfpathlineto{\pgfqpoint{1.583638in}{2.299121in}}%
\pgfpathlineto{\pgfqpoint{1.584378in}{2.295633in}}%
\pgfpathlineto{\pgfqpoint{1.583682in}{2.299565in}}%
\pgfpathlineto{\pgfqpoint{1.584763in}{2.298402in}}%
\pgfpathlineto{\pgfqpoint{1.585167in}{2.300066in}}%
\pgfpathlineto{\pgfqpoint{1.585580in}{2.297279in}}%
\pgfpathlineto{\pgfqpoint{1.585843in}{2.297756in}}%
\pgfpathlineto{\pgfqpoint{1.585853in}{2.297598in}}%
\pgfpathlineto{\pgfqpoint{1.586749in}{2.300545in}}%
\pgfpathlineto{\pgfqpoint{1.586788in}{2.300437in}}%
\pgfpathlineto{\pgfqpoint{1.586900in}{2.301573in}}%
\pgfpathlineto{\pgfqpoint{1.587820in}{2.297165in}}%
\pgfpathlineto{\pgfqpoint{1.588594in}{2.296502in}}%
\pgfpathlineto{\pgfqpoint{1.588312in}{2.298969in}}%
\pgfpathlineto{\pgfqpoint{1.588925in}{2.297284in}}%
\pgfpathlineto{\pgfqpoint{1.589057in}{2.299076in}}%
\pgfpathlineto{\pgfqpoint{1.589821in}{2.295357in}}%
\pgfpathlineto{\pgfqpoint{1.590001in}{2.296030in}}%
\pgfpathlineto{\pgfqpoint{1.590415in}{2.293734in}}%
\pgfpathlineto{\pgfqpoint{1.591087in}{2.297157in}}%
\pgfpathlineto{\pgfqpoint{1.591092in}{2.297393in}}%
\pgfpathlineto{\pgfqpoint{1.591647in}{2.293375in}}%
\pgfpathlineto{\pgfqpoint{1.592163in}{2.295647in}}%
\pgfpathlineto{\pgfqpoint{1.592197in}{2.295181in}}%
\pgfpathlineto{\pgfqpoint{1.593044in}{2.297888in}}%
\pgfpathlineto{\pgfqpoint{1.593219in}{2.296936in}}%
\pgfpathlineto{\pgfqpoint{1.594105in}{2.304259in}}%
\pgfpathlineto{\pgfqpoint{1.594431in}{2.301572in}}%
\pgfpathlineto{\pgfqpoint{1.594436in}{2.301566in}}%
\pgfpathlineto{\pgfqpoint{1.594539in}{2.302819in}}%
\pgfpathlineto{\pgfqpoint{1.594577in}{2.302793in}}%
\pgfpathlineto{\pgfqpoint{1.594587in}{2.303331in}}%
\pgfpathlineto{\pgfqpoint{1.595225in}{2.298737in}}%
\pgfpathlineto{\pgfqpoint{1.595673in}{2.301260in}}%
\pgfpathlineto{\pgfqpoint{1.595799in}{2.300649in}}%
\pgfpathlineto{\pgfqpoint{1.596861in}{2.304120in}}%
\pgfpathlineto{\pgfqpoint{1.597182in}{2.305230in}}%
\pgfpathlineto{\pgfqpoint{1.598073in}{2.300189in}}%
\pgfpathlineto{\pgfqpoint{1.598721in}{2.295823in}}%
\pgfpathlineto{\pgfqpoint{1.598102in}{2.300214in}}%
\pgfpathlineto{\pgfqpoint{1.599261in}{2.298935in}}%
\pgfpathlineto{\pgfqpoint{1.599870in}{2.301158in}}%
\pgfpathlineto{\pgfqpoint{1.599621in}{2.298221in}}%
\pgfpathlineto{\pgfqpoint{1.600386in}{2.299882in}}%
\pgfpathlineto{\pgfqpoint{1.601281in}{2.301887in}}%
\pgfpathlineto{\pgfqpoint{1.601510in}{2.299032in}}%
\pgfpathlineto{\pgfqpoint{1.601569in}{2.300298in}}%
\pgfpathlineto{\pgfqpoint{1.601773in}{2.297035in}}%
\pgfpathlineto{\pgfqpoint{1.602625in}{2.299399in}}%
\pgfpathlineto{\pgfqpoint{1.602718in}{2.298543in}}%
\pgfpathlineto{\pgfqpoint{1.603467in}{2.303864in}}%
\pgfpathlineto{\pgfqpoint{1.603691in}{2.302866in}}%
\pgfpathlineto{\pgfqpoint{1.604485in}{2.304510in}}%
\pgfpathlineto{\pgfqpoint{1.604013in}{2.301882in}}%
\pgfpathlineto{\pgfqpoint{1.604899in}{2.303038in}}%
\pgfpathlineto{\pgfqpoint{1.605230in}{2.300996in}}%
\pgfpathlineto{\pgfqpoint{1.605824in}{2.304236in}}%
\pgfpathlineto{\pgfqpoint{1.606014in}{2.302605in}}%
\pgfpathlineto{\pgfqpoint{1.606729in}{2.299663in}}%
\pgfpathlineto{\pgfqpoint{1.607153in}{2.301551in}}%
\pgfpathlineto{\pgfqpoint{1.607323in}{2.302772in}}%
\pgfpathlineto{\pgfqpoint{1.607878in}{2.297899in}}%
\pgfpathlineto{\pgfqpoint{1.608253in}{2.300529in}}%
\pgfpathlineto{\pgfqpoint{1.609412in}{2.297612in}}%
\pgfpathlineto{\pgfqpoint{1.608555in}{2.300987in}}%
\pgfpathlineto{\pgfqpoint{1.609422in}{2.297643in}}%
\pgfpathlineto{\pgfqpoint{1.609426in}{2.297802in}}%
\pgfpathlineto{\pgfqpoint{1.609626in}{2.294013in}}%
\pgfpathlineto{\pgfqpoint{1.610468in}{2.295637in}}%
\pgfpathlineto{\pgfqpoint{1.610532in}{2.295257in}}%
\pgfpathlineto{\pgfqpoint{1.610634in}{2.296463in}}%
\pgfpathlineto{\pgfqpoint{1.611442in}{2.301181in}}%
\pgfpathlineto{\pgfqpoint{1.611754in}{2.296986in}}%
\pgfpathlineto{\pgfqpoint{1.612859in}{2.292815in}}%
\pgfpathlineto{\pgfqpoint{1.611953in}{2.298010in}}%
\pgfpathlineto{\pgfqpoint{1.612980in}{2.293363in}}%
\pgfpathlineto{\pgfqpoint{1.614344in}{2.297054in}}%
\pgfpathlineto{\pgfqpoint{1.613156in}{2.291374in}}%
\pgfpathlineto{\pgfqpoint{1.614417in}{2.296303in}}%
\pgfpathlineto{\pgfqpoint{1.614436in}{2.296291in}}%
\pgfpathlineto{\pgfqpoint{1.614568in}{2.297617in}}%
\pgfpathlineto{\pgfqpoint{1.615288in}{2.300503in}}%
\pgfpathlineto{\pgfqpoint{1.614611in}{2.297289in}}%
\pgfpathlineto{\pgfqpoint{1.615780in}{2.300064in}}%
\pgfpathlineto{\pgfqpoint{1.616223in}{2.295801in}}%
\pgfpathlineto{\pgfqpoint{1.615819in}{2.300382in}}%
\pgfpathlineto{\pgfqpoint{1.616909in}{2.298558in}}%
\pgfpathlineto{\pgfqpoint{1.616919in}{2.299018in}}%
\pgfpathlineto{\pgfqpoint{1.617416in}{2.294700in}}%
\pgfpathlineto{\pgfqpoint{1.617971in}{2.295714in}}%
\pgfpathlineto{\pgfqpoint{1.618350in}{2.293826in}}%
\pgfpathlineto{\pgfqpoint{1.618526in}{2.296598in}}%
\pgfpathlineto{\pgfqpoint{1.618959in}{2.296465in}}%
\pgfpathlineto{\pgfqpoint{1.619967in}{2.299330in}}%
\pgfpathlineto{\pgfqpoint{1.619407in}{2.295764in}}%
\pgfpathlineto{\pgfqpoint{1.620108in}{2.297576in}}%
\pgfpathlineto{\pgfqpoint{1.621228in}{2.294988in}}%
\pgfpathlineto{\pgfqpoint{1.620716in}{2.297663in}}%
\pgfpathlineto{\pgfqpoint{1.621252in}{2.295705in}}%
\pgfpathlineto{\pgfqpoint{1.622294in}{2.298108in}}%
\pgfpathlineto{\pgfqpoint{1.621666in}{2.294211in}}%
\pgfpathlineto{\pgfqpoint{1.622377in}{2.297081in}}%
\pgfpathlineto{\pgfqpoint{1.622411in}{2.295544in}}%
\pgfpathlineto{\pgfqpoint{1.623350in}{2.299599in}}%
\pgfpathlineto{\pgfqpoint{1.623462in}{2.298532in}}%
\pgfpathlineto{\pgfqpoint{1.623832in}{2.295447in}}%
\pgfpathlineto{\pgfqpoint{1.624494in}{2.299231in}}%
\pgfpathlineto{\pgfqpoint{1.624567in}{2.298651in}}%
\pgfpathlineto{\pgfqpoint{1.624675in}{2.299448in}}%
\pgfpathlineto{\pgfqpoint{1.625395in}{2.295391in}}%
\pgfpathlineto{\pgfqpoint{1.625473in}{2.294105in}}%
\pgfpathlineto{\pgfqpoint{1.625673in}{2.298033in}}%
\pgfpathlineto{\pgfqpoint{1.626505in}{2.295409in}}%
\pgfpathlineto{\pgfqpoint{1.626554in}{2.295944in}}%
\pgfpathlineto{\pgfqpoint{1.627187in}{2.292359in}}%
\pgfpathlineto{\pgfqpoint{1.627771in}{2.290305in}}%
\pgfpathlineto{\pgfqpoint{1.628102in}{2.292720in}}%
\pgfpathlineto{\pgfqpoint{1.628277in}{2.292547in}}%
\pgfpathlineto{\pgfqpoint{1.628282in}{2.292806in}}%
\pgfpathlineto{\pgfqpoint{1.628784in}{2.290245in}}%
\pgfpathlineto{\pgfqpoint{1.629368in}{2.291812in}}%
\pgfpathlineto{\pgfqpoint{1.630478in}{2.289104in}}%
\pgfpathlineto{\pgfqpoint{1.629519in}{2.292545in}}%
\pgfpathlineto{\pgfqpoint{1.630507in}{2.289529in}}%
\pgfpathlineto{\pgfqpoint{1.631510in}{2.291511in}}%
\pgfpathlineto{\pgfqpoint{1.630570in}{2.288347in}}%
\pgfpathlineto{\pgfqpoint{1.631661in}{2.291362in}}%
\pgfpathlineto{\pgfqpoint{1.632537in}{2.287517in}}%
\pgfpathlineto{\pgfqpoint{1.631909in}{2.291681in}}%
\pgfpathlineto{\pgfqpoint{1.632795in}{2.290194in}}%
\pgfpathlineto{\pgfqpoint{1.632815in}{2.290932in}}%
\pgfpathlineto{\pgfqpoint{1.633623in}{2.286204in}}%
\pgfpathlineto{\pgfqpoint{1.633842in}{2.288239in}}%
\pgfpathlineto{\pgfqpoint{1.634154in}{2.286502in}}%
\pgfpathlineto{\pgfqpoint{1.634791in}{2.290941in}}%
\pgfpathlineto{\pgfqpoint{1.634937in}{2.289137in}}%
\pgfpathlineto{\pgfqpoint{1.638321in}{2.302410in}}%
\pgfpathlineto{\pgfqpoint{1.634967in}{2.289123in}}%
\pgfpathlineto{\pgfqpoint{1.639105in}{2.299940in}}%
\pgfpathlineto{\pgfqpoint{1.639188in}{2.299527in}}%
\pgfpathlineto{\pgfqpoint{1.640049in}{2.303557in}}%
\pgfpathlineto{\pgfqpoint{1.640877in}{2.304559in}}%
\pgfpathlineto{\pgfqpoint{1.640619in}{2.302076in}}%
\pgfpathlineto{\pgfqpoint{1.641159in}{2.303723in}}%
\pgfpathlineto{\pgfqpoint{1.641573in}{2.300821in}}%
\pgfpathlineto{\pgfqpoint{1.641909in}{2.304788in}}%
\pgfpathlineto{\pgfqpoint{1.642250in}{2.304069in}}%
\pgfpathlineto{\pgfqpoint{1.642274in}{2.304708in}}%
\pgfpathlineto{\pgfqpoint{1.642873in}{2.300850in}}%
\pgfpathlineto{\pgfqpoint{1.643350in}{2.303512in}}%
\pgfpathlineto{\pgfqpoint{1.643730in}{2.301918in}}%
\pgfpathlineto{\pgfqpoint{1.643477in}{2.304828in}}%
\pgfpathlineto{\pgfqpoint{1.644460in}{2.303544in}}%
\pgfpathlineto{\pgfqpoint{1.644928in}{2.304329in}}%
\pgfpathlineto{\pgfqpoint{1.645478in}{2.301314in}}%
\pgfpathlineto{\pgfqpoint{1.645560in}{2.303242in}}%
\pgfpathlineto{\pgfqpoint{1.646563in}{2.299522in}}%
\pgfpathlineto{\pgfqpoint{1.645638in}{2.304277in}}%
\pgfpathlineto{\pgfqpoint{1.646773in}{2.300418in}}%
\pgfpathlineto{\pgfqpoint{1.647824in}{2.304354in}}%
\pgfpathlineto{\pgfqpoint{1.647031in}{2.299524in}}%
\pgfpathlineto{\pgfqpoint{1.647927in}{2.303464in}}%
\pgfpathlineto{\pgfqpoint{1.648939in}{2.299223in}}%
\pgfpathlineto{\pgfqpoint{1.649669in}{2.301350in}}%
\pgfpathlineto{\pgfqpoint{1.650473in}{2.306286in}}%
\pgfpathlineto{\pgfqpoint{1.649728in}{2.301115in}}%
\pgfpathlineto{\pgfqpoint{1.650814in}{2.304288in}}%
\pgfpathlineto{\pgfqpoint{1.652873in}{2.298872in}}%
\pgfpathlineto{\pgfqpoint{1.650843in}{2.305178in}}%
\pgfpathlineto{\pgfqpoint{1.653228in}{2.301816in}}%
\pgfpathlineto{\pgfqpoint{1.654129in}{2.303661in}}%
\pgfpathlineto{\pgfqpoint{1.653326in}{2.300644in}}%
\pgfpathlineto{\pgfqpoint{1.654382in}{2.303569in}}%
\pgfpathlineto{\pgfqpoint{1.654665in}{2.302261in}}%
\pgfpathlineto{\pgfqpoint{1.655380in}{2.305796in}}%
\pgfpathlineto{\pgfqpoint{1.655473in}{2.305153in}}%
\pgfpathlineto{\pgfqpoint{1.655809in}{2.303451in}}%
\pgfpathlineto{\pgfqpoint{1.656057in}{2.306652in}}%
\pgfpathlineto{\pgfqpoint{1.656125in}{2.306338in}}%
\pgfpathlineto{\pgfqpoint{1.656471in}{2.307500in}}%
\pgfpathlineto{\pgfqpoint{1.656928in}{2.304188in}}%
\pgfpathlineto{\pgfqpoint{1.657211in}{2.304949in}}%
\pgfpathlineto{\pgfqpoint{1.657216in}{2.305020in}}%
\pgfpathlineto{\pgfqpoint{1.657532in}{2.302904in}}%
\pgfpathlineto{\pgfqpoint{1.658214in}{2.303439in}}%
\pgfpathlineto{\pgfqpoint{1.658545in}{2.302434in}}%
\pgfpathlineto{\pgfqpoint{1.658725in}{2.304565in}}%
\pgfpathlineto{\pgfqpoint{1.659324in}{2.303361in}}%
\pgfpathlineto{\pgfqpoint{1.659587in}{2.304452in}}%
\pgfpathlineto{\pgfqpoint{1.659966in}{2.300868in}}%
\pgfpathlineto{\pgfqpoint{1.660429in}{2.303273in}}%
\pgfpathlineto{\pgfqpoint{1.661247in}{2.301990in}}%
\pgfpathlineto{\pgfqpoint{1.660994in}{2.304774in}}%
\pgfpathlineto{\pgfqpoint{1.661407in}{2.303987in}}%
\pgfpathlineto{\pgfqpoint{1.661446in}{2.304583in}}%
\pgfpathlineto{\pgfqpoint{1.662352in}{2.302385in}}%
\pgfpathlineto{\pgfqpoint{1.662498in}{2.302791in}}%
\pgfpathlineto{\pgfqpoint{1.663472in}{2.301152in}}%
\pgfpathlineto{\pgfqpoint{1.662712in}{2.305506in}}%
\pgfpathlineto{\pgfqpoint{1.663618in}{2.302460in}}%
\pgfpathlineto{\pgfqpoint{1.664255in}{2.304365in}}%
\pgfpathlineto{\pgfqpoint{1.664392in}{2.302390in}}%
\pgfpathlineto{\pgfqpoint{1.664767in}{2.303981in}}%
\pgfpathlineto{\pgfqpoint{1.665400in}{2.300874in}}%
\pgfpathlineto{\pgfqpoint{1.664854in}{2.304539in}}%
\pgfpathlineto{\pgfqpoint{1.665940in}{2.303000in}}%
\pgfpathlineto{\pgfqpoint{1.667069in}{2.301446in}}%
\pgfpathlineto{\pgfqpoint{1.666782in}{2.304545in}}%
\pgfpathlineto{\pgfqpoint{1.667108in}{2.302444in}}%
\pgfpathlineto{\pgfqpoint{1.668092in}{2.305634in}}%
\pgfpathlineto{\pgfqpoint{1.667605in}{2.302064in}}%
\pgfpathlineto{\pgfqpoint{1.668253in}{2.303683in}}%
\pgfpathlineto{\pgfqpoint{1.669104in}{2.301412in}}%
\pgfpathlineto{\pgfqpoint{1.668540in}{2.304479in}}%
\pgfpathlineto{\pgfqpoint{1.669377in}{2.302666in}}%
\pgfpathlineto{\pgfqpoint{1.670507in}{2.303636in}}%
\pgfpathlineto{\pgfqpoint{1.669552in}{2.301474in}}%
\pgfpathlineto{\pgfqpoint{1.670521in}{2.303382in}}%
\pgfpathlineto{\pgfqpoint{1.671271in}{2.299804in}}%
\pgfpathlineto{\pgfqpoint{1.670726in}{2.304441in}}%
\pgfpathlineto{\pgfqpoint{1.671743in}{2.302290in}}%
\pgfpathlineto{\pgfqpoint{1.672220in}{2.303546in}}%
\pgfpathlineto{\pgfqpoint{1.672727in}{2.301331in}}%
\pgfpathlineto{\pgfqpoint{1.672732in}{2.301559in}}%
\pgfpathlineto{\pgfqpoint{1.672917in}{2.298666in}}%
\pgfpathlineto{\pgfqpoint{1.673783in}{2.303941in}}%
\pgfpathlineto{\pgfqpoint{1.673944in}{2.305155in}}%
\pgfpathlineto{\pgfqpoint{1.674586in}{2.302421in}}%
\pgfpathlineto{\pgfqpoint{1.674893in}{2.304136in}}%
\pgfpathlineto{\pgfqpoint{1.675404in}{2.301810in}}%
\pgfpathlineto{\pgfqpoint{1.675838in}{2.305136in}}%
\pgfpathlineto{\pgfqpoint{1.676028in}{2.303146in}}%
\pgfpathlineto{\pgfqpoint{1.676675in}{2.304206in}}%
\pgfpathlineto{\pgfqpoint{1.676957in}{2.301731in}}%
\pgfpathlineto{\pgfqpoint{1.676987in}{2.301845in}}%
\pgfpathlineto{\pgfqpoint{1.677045in}{2.301259in}}%
\pgfpathlineto{\pgfqpoint{1.677770in}{2.304415in}}%
\pgfpathlineto{\pgfqpoint{1.678058in}{2.303715in}}%
\pgfpathlineto{\pgfqpoint{1.678136in}{2.304251in}}%
\pgfpathlineto{\pgfqpoint{1.679061in}{2.301036in}}%
\pgfpathlineto{\pgfqpoint{1.679070in}{2.301472in}}%
\pgfpathlineto{\pgfqpoint{1.679343in}{2.299927in}}%
\pgfpathlineto{\pgfqpoint{1.679888in}{2.303074in}}%
\pgfpathlineto{\pgfqpoint{1.680166in}{2.301798in}}%
\pgfpathlineto{\pgfqpoint{1.681261in}{2.303975in}}%
\pgfpathlineto{\pgfqpoint{1.680244in}{2.301174in}}%
\pgfpathlineto{\pgfqpoint{1.681315in}{2.303658in}}%
\pgfpathlineto{\pgfqpoint{1.681918in}{2.301754in}}%
\pgfpathlineto{\pgfqpoint{1.681592in}{2.304095in}}%
\pgfpathlineto{\pgfqpoint{1.682444in}{2.303082in}}%
\pgfpathlineto{\pgfqpoint{1.682936in}{2.304542in}}%
\pgfpathlineto{\pgfqpoint{1.683428in}{2.302543in}}%
\pgfpathlineto{\pgfqpoint{1.683564in}{2.303572in}}%
\pgfpathlineto{\pgfqpoint{1.683705in}{2.302225in}}%
\pgfpathlineto{\pgfqpoint{1.684304in}{2.304835in}}%
\pgfpathlineto{\pgfqpoint{1.684669in}{2.303832in}}%
\pgfpathlineto{\pgfqpoint{1.685803in}{2.301037in}}%
\pgfpathlineto{\pgfqpoint{1.685356in}{2.304234in}}%
\pgfpathlineto{\pgfqpoint{1.685857in}{2.301450in}}%
\pgfpathlineto{\pgfqpoint{1.685945in}{2.302860in}}%
\pgfpathlineto{\pgfqpoint{1.686573in}{2.300664in}}%
\pgfpathlineto{\pgfqpoint{1.686991in}{2.302271in}}%
\pgfpathlineto{\pgfqpoint{1.687206in}{2.299981in}}%
\pgfpathlineto{\pgfqpoint{1.687001in}{2.302289in}}%
\pgfpathlineto{\pgfqpoint{1.688116in}{2.301763in}}%
\pgfpathlineto{\pgfqpoint{1.689012in}{2.303758in}}%
\pgfpathlineto{\pgfqpoint{1.688150in}{2.301039in}}%
\pgfpathlineto{\pgfqpoint{1.689241in}{2.302628in}}%
\pgfpathlineto{\pgfqpoint{1.690292in}{2.299867in}}%
\pgfpathlineto{\pgfqpoint{1.689479in}{2.303026in}}%
\pgfpathlineto{\pgfqpoint{1.690380in}{2.300635in}}%
\pgfpathlineto{\pgfqpoint{1.690618in}{2.299404in}}%
\pgfpathlineto{\pgfqpoint{1.691324in}{2.303143in}}%
\pgfpathlineto{\pgfqpoint{1.691329in}{2.303024in}}%
\pgfpathlineto{\pgfqpoint{1.692210in}{2.301271in}}%
\pgfpathlineto{\pgfqpoint{1.692459in}{2.304097in}}%
\pgfpathlineto{\pgfqpoint{1.693355in}{2.299903in}}%
\pgfpathlineto{\pgfqpoint{1.692595in}{2.304690in}}%
\pgfpathlineto{\pgfqpoint{1.693613in}{2.302360in}}%
\pgfpathlineto{\pgfqpoint{1.694713in}{2.305336in}}%
\pgfpathlineto{\pgfqpoint{1.693754in}{2.301884in}}%
\pgfpathlineto{\pgfqpoint{1.694810in}{2.304819in}}%
\pgfpathlineto{\pgfqpoint{1.695054in}{2.301589in}}%
\pgfpathlineto{\pgfqpoint{1.695940in}{2.303841in}}%
\pgfpathlineto{\pgfqpoint{1.695964in}{2.303997in}}%
\pgfpathlineto{\pgfqpoint{1.696621in}{2.299674in}}%
\pgfpathlineto{\pgfqpoint{1.696714in}{2.300097in}}%
\pgfpathlineto{\pgfqpoint{1.697395in}{2.296803in}}%
\pgfpathlineto{\pgfqpoint{1.696977in}{2.301483in}}%
\pgfpathlineto{\pgfqpoint{1.697838in}{2.299418in}}%
\pgfpathlineto{\pgfqpoint{1.699143in}{2.306194in}}%
\pgfpathlineto{\pgfqpoint{1.697946in}{2.298916in}}%
\pgfpathlineto{\pgfqpoint{1.699153in}{2.305669in}}%
\pgfpathlineto{\pgfqpoint{1.700292in}{2.302489in}}%
\pgfpathlineto{\pgfqpoint{1.700833in}{2.304557in}}%
\pgfpathlineto{\pgfqpoint{1.701397in}{2.302253in}}%
\pgfpathlineto{\pgfqpoint{1.702040in}{2.304460in}}%
\pgfpathlineto{\pgfqpoint{1.702459in}{2.301375in}}%
\pgfpathlineto{\pgfqpoint{1.702463in}{2.301401in}}%
\pgfpathlineto{\pgfqpoint{1.702804in}{2.299068in}}%
\pgfpathlineto{\pgfqpoint{1.703539in}{2.301991in}}%
\pgfpathlineto{\pgfqpoint{1.703544in}{2.301958in}}%
\pgfpathlineto{\pgfqpoint{1.703885in}{2.303622in}}%
\pgfpathlineto{\pgfqpoint{1.704567in}{2.299677in}}%
\pgfpathlineto{\pgfqpoint{1.704591in}{2.299705in}}%
\pgfpathlineto{\pgfqpoint{1.704937in}{2.296120in}}%
\pgfpathlineto{\pgfqpoint{1.705613in}{2.300357in}}%
\pgfpathlineto{\pgfqpoint{1.705696in}{2.299828in}}%
\pgfpathlineto{\pgfqpoint{1.706782in}{2.301926in}}%
\pgfpathlineto{\pgfqpoint{1.706159in}{2.297891in}}%
\pgfpathlineto{\pgfqpoint{1.706840in}{2.301298in}}%
\pgfpathlineto{\pgfqpoint{1.706884in}{2.300330in}}%
\pgfpathlineto{\pgfqpoint{1.707814in}{2.305042in}}%
\pgfpathlineto{\pgfqpoint{1.707921in}{2.303366in}}%
\pgfpathlineto{\pgfqpoint{1.708705in}{2.300503in}}%
\pgfpathlineto{\pgfqpoint{1.707975in}{2.303778in}}%
\pgfpathlineto{\pgfqpoint{1.709075in}{2.303296in}}%
\pgfpathlineto{\pgfqpoint{1.709484in}{2.307985in}}%
\pgfpathlineto{\pgfqpoint{1.710156in}{2.302242in}}%
\pgfpathlineto{\pgfqpoint{1.710161in}{2.302301in}}%
\pgfpathlineto{\pgfqpoint{1.710195in}{2.302123in}}%
\pgfpathlineto{\pgfqpoint{1.710433in}{2.304938in}}%
\pgfpathlineto{\pgfqpoint{1.711261in}{2.302631in}}%
\pgfpathlineto{\pgfqpoint{1.712244in}{2.304420in}}%
\pgfpathlineto{\pgfqpoint{1.711631in}{2.300387in}}%
\pgfpathlineto{\pgfqpoint{1.712420in}{2.304160in}}%
\pgfpathlineto{\pgfqpoint{1.712658in}{2.303580in}}%
\pgfpathlineto{\pgfqpoint{1.713140in}{2.306955in}}%
\pgfpathlineto{\pgfqpoint{1.713530in}{2.304082in}}%
\pgfpathlineto{\pgfqpoint{1.714669in}{2.300100in}}%
\pgfpathlineto{\pgfqpoint{1.714055in}{2.305262in}}%
\pgfpathlineto{\pgfqpoint{1.714786in}{2.301273in}}%
\pgfpathlineto{\pgfqpoint{1.715589in}{2.304313in}}%
\pgfpathlineto{\pgfqpoint{1.714980in}{2.300142in}}%
\pgfpathlineto{\pgfqpoint{1.715939in}{2.303505in}}%
\pgfpathlineto{\pgfqpoint{1.716894in}{2.301906in}}%
\pgfpathlineto{\pgfqpoint{1.716558in}{2.304619in}}%
\pgfpathlineto{\pgfqpoint{1.717054in}{2.303034in}}%
\pgfpathlineto{\pgfqpoint{1.717410in}{2.304879in}}%
\pgfpathlineto{\pgfqpoint{1.717225in}{2.301879in}}%
\pgfpathlineto{\pgfqpoint{1.718169in}{2.303574in}}%
\pgfpathlineto{\pgfqpoint{1.718335in}{2.303886in}}%
\pgfpathlineto{\pgfqpoint{1.719041in}{2.300105in}}%
\pgfpathlineto{\pgfqpoint{1.719075in}{2.300458in}}%
\pgfpathlineto{\pgfqpoint{1.719270in}{2.298317in}}%
\pgfpathlineto{\pgfqpoint{1.719703in}{2.301128in}}%
\pgfpathlineto{\pgfqpoint{1.720185in}{2.300297in}}%
\pgfpathlineto{\pgfqpoint{1.720896in}{2.302772in}}%
\pgfpathlineto{\pgfqpoint{1.720550in}{2.299708in}}%
\pgfpathlineto{\pgfqpoint{1.721339in}{2.302081in}}%
\pgfpathlineto{\pgfqpoint{1.722283in}{2.300033in}}%
\pgfpathlineto{\pgfqpoint{1.721543in}{2.303697in}}%
\pgfpathlineto{\pgfqpoint{1.722429in}{2.302877in}}%
\pgfpathlineto{\pgfqpoint{1.723091in}{2.305906in}}%
\pgfpathlineto{\pgfqpoint{1.722648in}{2.301877in}}%
\pgfpathlineto{\pgfqpoint{1.723544in}{2.303016in}}%
\pgfpathlineto{\pgfqpoint{1.724026in}{2.301784in}}%
\pgfpathlineto{\pgfqpoint{1.724304in}{2.305109in}}%
\pgfpathlineto{\pgfqpoint{1.724669in}{2.302313in}}%
\pgfpathlineto{\pgfqpoint{1.724678in}{2.302516in}}%
\pgfpathlineto{\pgfqpoint{1.724844in}{2.300650in}}%
\pgfpathlineto{\pgfqpoint{1.725248in}{2.300597in}}%
\pgfpathlineto{\pgfqpoint{1.726061in}{2.299931in}}%
\pgfpathlineto{\pgfqpoint{1.725652in}{2.302573in}}%
\pgfpathlineto{\pgfqpoint{1.726339in}{2.300803in}}%
\pgfpathlineto{\pgfqpoint{1.726787in}{2.302577in}}%
\pgfpathlineto{\pgfqpoint{1.726587in}{2.299923in}}%
\pgfpathlineto{\pgfqpoint{1.727458in}{2.301712in}}%
\pgfpathlineto{\pgfqpoint{1.727517in}{2.301186in}}%
\pgfpathlineto{\pgfqpoint{1.728150in}{2.304407in}}%
\pgfpathlineto{\pgfqpoint{1.728544in}{2.302244in}}%
\pgfpathlineto{\pgfqpoint{1.729162in}{2.304203in}}%
\pgfpathlineto{\pgfqpoint{1.728744in}{2.301625in}}%
\pgfpathlineto{\pgfqpoint{1.729756in}{2.303815in}}%
\pgfpathlineto{\pgfqpoint{1.730949in}{2.300332in}}%
\pgfpathlineto{\pgfqpoint{1.730190in}{2.305135in}}%
\pgfpathlineto{\pgfqpoint{1.730978in}{2.300870in}}%
\pgfpathlineto{\pgfqpoint{1.731003in}{2.301298in}}%
\pgfpathlineto{\pgfqpoint{1.731835in}{2.297960in}}%
\pgfpathlineto{\pgfqpoint{1.732044in}{2.299311in}}%
\pgfpathlineto{\pgfqpoint{1.733159in}{2.295693in}}%
\pgfpathlineto{\pgfqpoint{1.733208in}{2.296279in}}%
\pgfpathlineto{\pgfqpoint{1.734484in}{2.302306in}}%
\pgfpathlineto{\pgfqpoint{1.735433in}{2.300491in}}%
\pgfpathlineto{\pgfqpoint{1.735024in}{2.303499in}}%
\pgfpathlineto{\pgfqpoint{1.735642in}{2.301284in}}%
\pgfpathlineto{\pgfqpoint{1.736489in}{2.306681in}}%
\pgfpathlineto{\pgfqpoint{1.735662in}{2.301244in}}%
\pgfpathlineto{\pgfqpoint{1.736820in}{2.304221in}}%
\pgfpathlineto{\pgfqpoint{1.737152in}{2.300610in}}%
\pgfpathlineto{\pgfqpoint{1.736879in}{2.304375in}}%
\pgfpathlineto{\pgfqpoint{1.737979in}{2.302037in}}%
\pgfpathlineto{\pgfqpoint{1.738875in}{2.304580in}}%
\pgfpathlineto{\pgfqpoint{1.738218in}{2.301318in}}%
\pgfpathlineto{\pgfqpoint{1.739123in}{2.303430in}}%
\pgfpathlineto{\pgfqpoint{1.739308in}{2.303879in}}%
\pgfpathlineto{\pgfqpoint{1.740112in}{2.300433in}}%
\pgfpathlineto{\pgfqpoint{1.740151in}{2.300879in}}%
\pgfpathlineto{\pgfqpoint{1.740190in}{2.299626in}}%
\pgfpathlineto{\pgfqpoint{1.740764in}{2.303869in}}%
\pgfpathlineto{\pgfqpoint{1.741261in}{2.300624in}}%
\pgfpathlineto{\pgfqpoint{1.741338in}{2.301229in}}%
\pgfpathlineto{\pgfqpoint{1.741874in}{2.297855in}}%
\pgfpathlineto{\pgfqpoint{1.742322in}{2.299187in}}%
\pgfpathlineto{\pgfqpoint{1.742351in}{2.299397in}}%
\pgfpathlineto{\pgfqpoint{1.742911in}{2.295939in}}%
\pgfpathlineto{\pgfqpoint{1.743778in}{2.292799in}}%
\pgfpathlineto{\pgfqpoint{1.742965in}{2.296188in}}%
\pgfpathlineto{\pgfqpoint{1.744021in}{2.295714in}}%
\pgfpathlineto{\pgfqpoint{1.744109in}{2.297447in}}%
\pgfpathlineto{\pgfqpoint{1.744737in}{2.293782in}}%
\pgfpathlineto{\pgfqpoint{1.745116in}{2.294869in}}%
\pgfpathlineto{\pgfqpoint{1.745681in}{2.293190in}}%
\pgfpathlineto{\pgfqpoint{1.746080in}{2.296557in}}%
\pgfpathlineto{\pgfqpoint{1.746139in}{2.295765in}}%
\pgfpathlineto{\pgfqpoint{1.747288in}{2.305175in}}%
\pgfpathlineto{\pgfqpoint{1.747400in}{2.304036in}}%
\pgfpathlineto{\pgfqpoint{1.747828in}{2.301991in}}%
\pgfpathlineto{\pgfqpoint{1.747526in}{2.305509in}}%
\pgfpathlineto{\pgfqpoint{1.748646in}{2.302734in}}%
\pgfpathlineto{\pgfqpoint{1.749396in}{2.303335in}}%
\pgfpathlineto{\pgfqpoint{1.748938in}{2.301063in}}%
\pgfpathlineto{\pgfqpoint{1.749761in}{2.302846in}}%
\pgfpathlineto{\pgfqpoint{1.750788in}{2.301623in}}%
\pgfpathlineto{\pgfqpoint{1.750564in}{2.304665in}}%
\pgfpathlineto{\pgfqpoint{1.750871in}{2.302456in}}%
\pgfpathlineto{\pgfqpoint{1.751874in}{2.303474in}}%
\pgfpathlineto{\pgfqpoint{1.751704in}{2.300832in}}%
\pgfpathlineto{\pgfqpoint{1.751981in}{2.302904in}}%
\pgfpathlineto{\pgfqpoint{1.752629in}{2.302339in}}%
\pgfpathlineto{\pgfqpoint{1.752380in}{2.303966in}}%
\pgfpathlineto{\pgfqpoint{1.753081in}{2.303512in}}%
\pgfpathlineto{\pgfqpoint{1.753378in}{2.303067in}}%
\pgfpathlineto{\pgfqpoint{1.753739in}{2.305605in}}%
\pgfpathlineto{\pgfqpoint{1.753782in}{2.305314in}}%
\pgfpathlineto{\pgfqpoint{1.753807in}{2.305771in}}%
\pgfpathlineto{\pgfqpoint{1.754615in}{2.300307in}}%
\pgfpathlineto{\pgfqpoint{1.754805in}{2.301022in}}%
\pgfpathlineto{\pgfqpoint{1.755326in}{2.298495in}}%
\pgfpathlineto{\pgfqpoint{1.755866in}{2.302292in}}%
\pgfpathlineto{\pgfqpoint{1.755900in}{2.301755in}}%
\pgfpathlineto{\pgfqpoint{1.757161in}{2.303827in}}%
\pgfpathlineto{\pgfqpoint{1.756168in}{2.299953in}}%
\pgfpathlineto{\pgfqpoint{1.757185in}{2.303770in}}%
\pgfpathlineto{\pgfqpoint{1.757794in}{2.301846in}}%
\pgfpathlineto{\pgfqpoint{1.757901in}{2.303907in}}%
\pgfpathlineto{\pgfqpoint{1.758315in}{2.302880in}}%
\pgfpathlineto{\pgfqpoint{1.758724in}{2.304628in}}%
\pgfpathlineto{\pgfqpoint{1.759289in}{2.302460in}}%
\pgfpathlineto{\pgfqpoint{1.759430in}{2.302897in}}%
\pgfpathlineto{\pgfqpoint{1.760184in}{2.301648in}}%
\pgfpathlineto{\pgfqpoint{1.759688in}{2.305189in}}%
\pgfpathlineto{\pgfqpoint{1.760447in}{2.303317in}}%
\pgfpathlineto{\pgfqpoint{1.760632in}{2.305121in}}%
\pgfpathlineto{\pgfqpoint{1.761275in}{2.301305in}}%
\pgfpathlineto{\pgfqpoint{1.761533in}{2.302414in}}%
\pgfpathlineto{\pgfqpoint{1.762195in}{2.301717in}}%
\pgfpathlineto{\pgfqpoint{1.762492in}{2.304356in}}%
\pgfpathlineto{\pgfqpoint{1.763300in}{2.304851in}}%
\pgfpathlineto{\pgfqpoint{1.763519in}{2.302691in}}%
\pgfpathlineto{\pgfqpoint{1.763524in}{2.302717in}}%
\pgfpathlineto{\pgfqpoint{1.763958in}{2.301382in}}%
\pgfpathlineto{\pgfqpoint{1.764357in}{2.304133in}}%
\pgfpathlineto{\pgfqpoint{1.764625in}{2.302976in}}%
\pgfpathlineto{\pgfqpoint{1.764688in}{2.304634in}}%
\pgfpathlineto{\pgfqpoint{1.765306in}{2.301374in}}%
\pgfpathlineto{\pgfqpoint{1.765739in}{2.303248in}}%
\pgfpathlineto{\pgfqpoint{1.766348in}{2.302242in}}%
\pgfpathlineto{\pgfqpoint{1.766660in}{2.304396in}}%
\pgfpathlineto{\pgfqpoint{1.766845in}{2.303333in}}%
\pgfpathlineto{\pgfqpoint{1.767882in}{2.304411in}}%
\pgfpathlineto{\pgfqpoint{1.766986in}{2.301278in}}%
\pgfpathlineto{\pgfqpoint{1.767945in}{2.303213in}}%
\pgfpathlineto{\pgfqpoint{1.768373in}{2.299011in}}%
\pgfpathlineto{\pgfqpoint{1.768909in}{2.304388in}}%
\pgfpathlineto{\pgfqpoint{1.769094in}{2.301107in}}%
\pgfpathlineto{\pgfqpoint{1.769444in}{2.303328in}}%
\pgfpathlineto{\pgfqpoint{1.770082in}{2.300286in}}%
\pgfpathlineto{\pgfqpoint{1.770184in}{2.300613in}}%
\pgfpathlineto{\pgfqpoint{1.770671in}{2.299021in}}%
\pgfpathlineto{\pgfqpoint{1.770973in}{2.302891in}}%
\pgfpathlineto{\pgfqpoint{1.771124in}{2.302585in}}%
\pgfpathlineto{\pgfqpoint{1.771416in}{2.303833in}}%
\pgfpathlineto{\pgfqpoint{1.771888in}{2.301546in}}%
\pgfpathlineto{\pgfqpoint{1.772229in}{2.302398in}}%
\pgfpathlineto{\pgfqpoint{1.772443in}{2.303633in}}%
\pgfpathlineto{\pgfqpoint{1.773251in}{2.301402in}}%
\pgfpathlineto{\pgfqpoint{1.773295in}{2.301630in}}%
\pgfpathlineto{\pgfqpoint{1.774420in}{2.300300in}}%
\pgfpathlineto{\pgfqpoint{1.773918in}{2.303586in}}%
\pgfpathlineto{\pgfqpoint{1.774464in}{2.300767in}}%
\pgfpathlineto{\pgfqpoint{1.775608in}{2.304278in}}%
\pgfpathlineto{\pgfqpoint{1.774571in}{2.300244in}}%
\pgfpathlineto{\pgfqpoint{1.775730in}{2.303818in}}%
\pgfpathlineto{\pgfqpoint{1.776250in}{2.302573in}}%
\pgfpathlineto{\pgfqpoint{1.776528in}{2.304969in}}%
\pgfpathlineto{\pgfqpoint{1.776752in}{2.304308in}}%
\pgfpathlineto{\pgfqpoint{1.777507in}{2.305382in}}%
\pgfpathlineto{\pgfqpoint{1.777020in}{2.302599in}}%
\pgfpathlineto{\pgfqpoint{1.777862in}{2.304383in}}%
\pgfpathlineto{\pgfqpoint{1.777979in}{2.305220in}}%
\pgfpathlineto{\pgfqpoint{1.778646in}{2.302681in}}%
\pgfpathlineto{\pgfqpoint{1.778899in}{2.302884in}}%
\pgfpathlineto{\pgfqpoint{1.779800in}{2.302073in}}%
\pgfpathlineto{\pgfqpoint{1.779897in}{2.304251in}}%
\pgfpathlineto{\pgfqpoint{1.779990in}{2.303628in}}%
\pgfpathlineto{\pgfqpoint{1.780140in}{2.305058in}}%
\pgfpathlineto{\pgfqpoint{1.780613in}{2.302276in}}%
\pgfpathlineto{\pgfqpoint{1.781109in}{2.304467in}}%
\pgfpathlineto{\pgfqpoint{1.781523in}{2.302256in}}%
\pgfpathlineto{\pgfqpoint{1.781635in}{2.304538in}}%
\pgfpathlineto{\pgfqpoint{1.782307in}{2.302754in}}%
\pgfpathlineto{\pgfqpoint{1.782701in}{2.304199in}}%
\pgfpathlineto{\pgfqpoint{1.783344in}{2.301829in}}%
\pgfpathlineto{\pgfqpoint{1.783397in}{2.302424in}}%
\pgfpathlineto{\pgfqpoint{1.783456in}{2.301757in}}%
\pgfpathlineto{\pgfqpoint{1.783845in}{2.304087in}}%
\pgfpathlineto{\pgfqpoint{1.784483in}{2.303657in}}%
\pgfpathlineto{\pgfqpoint{1.784507in}{2.304653in}}%
\pgfpathlineto{\pgfqpoint{1.785374in}{2.297266in}}%
\pgfpathlineto{\pgfqpoint{1.785389in}{2.297393in}}%
\pgfpathlineto{\pgfqpoint{1.785394in}{2.297214in}}%
\pgfpathlineto{\pgfqpoint{1.786319in}{2.303696in}}%
\pgfpathlineto{\pgfqpoint{1.786328in}{2.303387in}}%
\pgfpathlineto{\pgfqpoint{1.786942in}{2.304485in}}%
\pgfpathlineto{\pgfqpoint{1.786445in}{2.302498in}}%
\pgfpathlineto{\pgfqpoint{1.787429in}{2.302995in}}%
\pgfpathlineto{\pgfqpoint{1.787541in}{2.303726in}}%
\pgfpathlineto{\pgfqpoint{1.787789in}{2.300970in}}%
\pgfpathlineto{\pgfqpoint{1.788412in}{2.301754in}}%
\pgfpathlineto{\pgfqpoint{1.789191in}{2.299239in}}%
\pgfpathlineto{\pgfqpoint{1.788548in}{2.303028in}}%
\pgfpathlineto{\pgfqpoint{1.789527in}{2.301659in}}%
\pgfpathlineto{\pgfqpoint{1.790131in}{2.303104in}}%
\pgfpathlineto{\pgfqpoint{1.790515in}{2.300618in}}%
\pgfpathlineto{\pgfqpoint{1.790647in}{2.302273in}}%
\pgfpathlineto{\pgfqpoint{1.791328in}{2.303122in}}%
\pgfpathlineto{\pgfqpoint{1.791640in}{2.299429in}}%
\pgfpathlineto{\pgfqpoint{1.791650in}{2.299779in}}%
\pgfpathlineto{\pgfqpoint{1.791908in}{2.298883in}}%
\pgfpathlineto{\pgfqpoint{1.792643in}{2.303564in}}%
\pgfpathlineto{\pgfqpoint{1.792696in}{2.303288in}}%
\pgfpathlineto{\pgfqpoint{1.793173in}{2.303907in}}%
\pgfpathlineto{\pgfqpoint{1.793962in}{2.300116in}}%
\pgfpathlineto{\pgfqpoint{1.794561in}{2.302196in}}%
\pgfpathlineto{\pgfqpoint{1.795009in}{2.299382in}}%
\pgfpathlineto{\pgfqpoint{1.795072in}{2.300277in}}%
\pgfpathlineto{\pgfqpoint{1.796075in}{2.298903in}}%
\pgfpathlineto{\pgfqpoint{1.795871in}{2.301702in}}%
\pgfpathlineto{\pgfqpoint{1.796202in}{2.299326in}}%
\pgfpathlineto{\pgfqpoint{1.796942in}{2.304149in}}%
\pgfpathlineto{\pgfqpoint{1.797370in}{2.303252in}}%
\pgfpathlineto{\pgfqpoint{1.797842in}{2.301389in}}%
\pgfpathlineto{\pgfqpoint{1.798300in}{2.304053in}}%
\pgfpathlineto{\pgfqpoint{1.798509in}{2.302739in}}%
\pgfpathlineto{\pgfqpoint{1.798563in}{2.303444in}}%
\pgfpathlineto{\pgfqpoint{1.799464in}{2.300796in}}%
\pgfpathlineto{\pgfqpoint{1.799517in}{2.300823in}}%
\pgfpathlineto{\pgfqpoint{1.799522in}{2.300808in}}%
\pgfpathlineto{\pgfqpoint{1.799760in}{2.303439in}}%
\pgfpathlineto{\pgfqpoint{1.799902in}{2.304026in}}%
\pgfpathlineto{\pgfqpoint{1.800189in}{2.301229in}}%
\pgfpathlineto{\pgfqpoint{1.800861in}{2.303258in}}%
\pgfpathlineto{\pgfqpoint{1.800890in}{2.302233in}}%
\pgfpathlineto{\pgfqpoint{1.801153in}{2.304799in}}%
\pgfpathlineto{\pgfqpoint{1.801966in}{2.303522in}}%
\pgfpathlineto{\pgfqpoint{1.802988in}{2.302141in}}%
\pgfpathlineto{\pgfqpoint{1.802258in}{2.305311in}}%
\pgfpathlineto{\pgfqpoint{1.803071in}{2.303693in}}%
\pgfpathlineto{\pgfqpoint{1.803874in}{2.300791in}}%
\pgfpathlineto{\pgfqpoint{1.803081in}{2.303913in}}%
\pgfpathlineto{\pgfqpoint{1.804186in}{2.303467in}}%
\pgfpathlineto{\pgfqpoint{1.804668in}{2.304000in}}%
\pgfpathlineto{\pgfqpoint{1.805038in}{2.300550in}}%
\pgfpathlineto{\pgfqpoint{1.805272in}{2.301626in}}%
\pgfpathlineto{\pgfqpoint{1.805817in}{2.304213in}}%
\pgfpathlineto{\pgfqpoint{1.806343in}{2.300139in}}%
\pgfpathlineto{\pgfqpoint{1.806357in}{2.300738in}}%
\pgfpathlineto{\pgfqpoint{1.807185in}{2.298954in}}%
\pgfpathlineto{\pgfqpoint{1.806922in}{2.301508in}}%
\pgfpathlineto{\pgfqpoint{1.807477in}{2.299861in}}%
\pgfpathlineto{\pgfqpoint{1.807482in}{2.300042in}}%
\pgfpathlineto{\pgfqpoint{1.808446in}{2.296019in}}%
\pgfpathlineto{\pgfqpoint{1.808499in}{2.296680in}}%
\pgfpathlineto{\pgfqpoint{1.809434in}{2.292835in}}%
\pgfpathlineto{\pgfqpoint{1.809683in}{2.293244in}}%
\pgfpathlineto{\pgfqpoint{1.810768in}{2.298156in}}%
\pgfpathlineto{\pgfqpoint{1.809945in}{2.292750in}}%
\pgfpathlineto{\pgfqpoint{1.810836in}{2.297915in}}%
\pgfpathlineto{\pgfqpoint{1.811479in}{2.296938in}}%
\pgfpathlineto{\pgfqpoint{1.811260in}{2.298552in}}%
\pgfpathlineto{\pgfqpoint{1.811956in}{2.297682in}}%
\pgfpathlineto{\pgfqpoint{1.813105in}{2.299927in}}%
\pgfpathlineto{\pgfqpoint{1.812044in}{2.297083in}}%
\pgfpathlineto{\pgfqpoint{1.813115in}{2.299898in}}%
\pgfpathlineto{\pgfqpoint{1.813519in}{2.297844in}}%
\pgfpathlineto{\pgfqpoint{1.814254in}{2.298098in}}%
\pgfpathlineto{\pgfqpoint{1.815057in}{2.295991in}}%
\pgfpathlineto{\pgfqpoint{1.814911in}{2.298225in}}%
\pgfpathlineto{\pgfqpoint{1.815403in}{2.296947in}}%
\pgfpathlineto{\pgfqpoint{1.817292in}{2.304997in}}%
\pgfpathlineto{\pgfqpoint{1.815452in}{2.296789in}}%
\pgfpathlineto{\pgfqpoint{1.817443in}{2.304229in}}%
\pgfpathlineto{\pgfqpoint{1.818178in}{2.302340in}}%
\pgfpathlineto{\pgfqpoint{1.817458in}{2.304397in}}%
\pgfpathlineto{\pgfqpoint{1.818543in}{2.304334in}}%
\pgfpathlineto{\pgfqpoint{1.819726in}{2.306105in}}%
\pgfpathlineto{\pgfqpoint{1.818855in}{2.302641in}}%
\pgfpathlineto{\pgfqpoint{1.819731in}{2.305967in}}%
\pgfpathlineto{\pgfqpoint{1.819999in}{2.308774in}}%
\pgfpathlineto{\pgfqpoint{1.820705in}{2.305365in}}%
\pgfpathlineto{\pgfqpoint{1.821766in}{2.301462in}}%
\pgfpathlineto{\pgfqpoint{1.821065in}{2.306519in}}%
\pgfpathlineto{\pgfqpoint{1.821917in}{2.303554in}}%
\pgfpathlineto{\pgfqpoint{1.822107in}{2.304487in}}%
\pgfpathlineto{\pgfqpoint{1.822925in}{2.298888in}}%
\pgfpathlineto{\pgfqpoint{1.822939in}{2.299018in}}%
\pgfpathlineto{\pgfqpoint{1.822944in}{2.298720in}}%
\pgfpathlineto{\pgfqpoint{1.823908in}{2.301502in}}%
\pgfpathlineto{\pgfqpoint{1.824030in}{2.299955in}}%
\pgfpathlineto{\pgfqpoint{1.825116in}{2.301055in}}%
\pgfpathlineto{\pgfqpoint{1.825047in}{2.299294in}}%
\pgfpathlineto{\pgfqpoint{1.825174in}{2.300833in}}%
\pgfpathlineto{\pgfqpoint{1.825271in}{2.300554in}}%
\pgfpathlineto{\pgfqpoint{1.825661in}{2.304564in}}%
\pgfpathlineto{\pgfqpoint{1.826245in}{2.302646in}}%
\pgfpathlineto{\pgfqpoint{1.826528in}{2.303992in}}%
\pgfpathlineto{\pgfqpoint{1.827136in}{2.299868in}}%
\pgfpathlineto{\pgfqpoint{1.827336in}{2.302044in}}%
\pgfpathlineto{\pgfqpoint{1.827394in}{2.300637in}}%
\pgfpathlineto{\pgfqpoint{1.828231in}{2.304252in}}%
\pgfpathlineto{\pgfqpoint{1.828421in}{2.303657in}}%
\pgfpathlineto{\pgfqpoint{1.828830in}{2.304918in}}%
\pgfpathlineto{\pgfqpoint{1.829512in}{2.301951in}}%
\pgfpathlineto{\pgfqpoint{1.829746in}{2.302063in}}%
\pgfpathlineto{\pgfqpoint{1.830072in}{2.298020in}}%
\pgfpathlineto{\pgfqpoint{1.830120in}{2.298200in}}%
\pgfpathlineto{\pgfqpoint{1.830593in}{2.294478in}}%
\pgfpathlineto{\pgfqpoint{1.830237in}{2.298275in}}%
\pgfpathlineto{\pgfqpoint{1.831240in}{2.297408in}}%
\pgfpathlineto{\pgfqpoint{1.832370in}{2.299652in}}%
\pgfpathlineto{\pgfqpoint{1.831605in}{2.295852in}}%
\pgfpathlineto{\pgfqpoint{1.832375in}{2.299622in}}%
\pgfpathlineto{\pgfqpoint{1.832637in}{2.301120in}}%
\pgfpathlineto{\pgfqpoint{1.832900in}{2.297874in}}%
\pgfpathlineto{\pgfqpoint{1.832910in}{2.297975in}}%
\pgfpathlineto{\pgfqpoint{1.832983in}{2.297772in}}%
\pgfpathlineto{\pgfqpoint{1.833791in}{2.301323in}}%
\pgfpathlineto{\pgfqpoint{1.833962in}{2.300023in}}%
\pgfpathlineto{\pgfqpoint{1.834926in}{2.301707in}}%
\pgfpathlineto{\pgfqpoint{1.834200in}{2.299125in}}%
\pgfpathlineto{\pgfqpoint{1.835043in}{2.299956in}}%
\pgfpathlineto{\pgfqpoint{1.835067in}{2.298684in}}%
\pgfpathlineto{\pgfqpoint{1.836060in}{2.303898in}}%
\pgfpathlineto{\pgfqpoint{1.836269in}{2.305024in}}%
\pgfpathlineto{\pgfqpoint{1.837029in}{2.300797in}}%
\pgfpathlineto{\pgfqpoint{1.837749in}{2.299177in}}%
\pgfpathlineto{\pgfqpoint{1.837311in}{2.302148in}}%
\pgfpathlineto{\pgfqpoint{1.838139in}{2.300841in}}%
\pgfpathlineto{\pgfqpoint{1.838723in}{2.304675in}}%
\pgfpathlineto{\pgfqpoint{1.838304in}{2.300144in}}%
\pgfpathlineto{\pgfqpoint{1.839283in}{2.302309in}}%
\pgfpathlineto{\pgfqpoint{1.839682in}{2.304467in}}%
\pgfpathlineto{\pgfqpoint{1.840189in}{2.301376in}}%
\pgfpathlineto{\pgfqpoint{1.840208in}{2.301490in}}%
\pgfpathlineto{\pgfqpoint{1.840218in}{2.301039in}}%
\pgfpathlineto{\pgfqpoint{1.841021in}{2.303897in}}%
\pgfpathlineto{\pgfqpoint{1.841279in}{2.302265in}}%
\pgfpathlineto{\pgfqpoint{1.841854in}{2.304799in}}%
\pgfpathlineto{\pgfqpoint{1.841333in}{2.301690in}}%
\pgfpathlineto{\pgfqpoint{1.842404in}{2.302897in}}%
\pgfpathlineto{\pgfqpoint{1.843314in}{2.301689in}}%
\pgfpathlineto{\pgfqpoint{1.842886in}{2.303825in}}%
\pgfpathlineto{\pgfqpoint{1.843480in}{2.303025in}}%
\pgfpathlineto{\pgfqpoint{1.843699in}{2.304093in}}%
\pgfpathlineto{\pgfqpoint{1.844385in}{2.300503in}}%
\pgfpathlineto{\pgfqpoint{1.844570in}{2.301882in}}%
\pgfpathlineto{\pgfqpoint{1.844950in}{2.301246in}}%
\pgfpathlineto{\pgfqpoint{1.845198in}{2.303417in}}%
\pgfpathlineto{\pgfqpoint{1.845666in}{2.302273in}}%
\pgfpathlineto{\pgfqpoint{1.845846in}{2.304318in}}%
\pgfpathlineto{\pgfqpoint{1.845690in}{2.302107in}}%
\pgfpathlineto{\pgfqpoint{1.846795in}{2.303575in}}%
\pgfpathlineto{\pgfqpoint{1.847598in}{2.299106in}}%
\pgfpathlineto{\pgfqpoint{1.846819in}{2.303771in}}%
\pgfpathlineto{\pgfqpoint{1.847983in}{2.300831in}}%
\pgfpathlineto{\pgfqpoint{1.848913in}{2.303920in}}%
\pgfpathlineto{\pgfqpoint{1.848256in}{2.299341in}}%
\pgfpathlineto{\pgfqpoint{1.849190in}{2.303061in}}%
\pgfpathlineto{\pgfqpoint{1.849862in}{2.300611in}}%
\pgfpathlineto{\pgfqpoint{1.849366in}{2.303832in}}%
\pgfpathlineto{\pgfqpoint{1.850325in}{2.301475in}}%
\pgfpathlineto{\pgfqpoint{1.851323in}{2.304664in}}%
\pgfpathlineto{\pgfqpoint{1.850432in}{2.301408in}}%
\pgfpathlineto{\pgfqpoint{1.851479in}{2.303729in}}%
\pgfpathlineto{\pgfqpoint{1.852345in}{2.302596in}}%
\pgfpathlineto{\pgfqpoint{1.851668in}{2.305228in}}%
\pgfpathlineto{\pgfqpoint{1.852555in}{2.303910in}}%
\pgfpathlineto{\pgfqpoint{1.852608in}{2.304948in}}%
\pgfpathlineto{\pgfqpoint{1.853533in}{2.299895in}}%
\pgfpathlineto{\pgfqpoint{1.853543in}{2.300379in}}%
\pgfpathlineto{\pgfqpoint{1.854629in}{2.297349in}}%
\pgfpathlineto{\pgfqpoint{1.853572in}{2.300511in}}%
\pgfpathlineto{\pgfqpoint{1.854731in}{2.297686in}}%
\pgfpathlineto{\pgfqpoint{1.855461in}{2.296248in}}%
\pgfpathlineto{\pgfqpoint{1.855125in}{2.300196in}}%
\pgfpathlineto{\pgfqpoint{1.855826in}{2.298030in}}%
\pgfpathlineto{\pgfqpoint{1.857087in}{2.301826in}}%
\pgfpathlineto{\pgfqpoint{1.856469in}{2.297695in}}%
\pgfpathlineto{\pgfqpoint{1.857223in}{2.299840in}}%
\pgfpathlineto{\pgfqpoint{1.857486in}{2.297302in}}%
\pgfpathlineto{\pgfqpoint{1.858212in}{2.300853in}}%
\pgfpathlineto{\pgfqpoint{1.858333in}{2.299686in}}%
\pgfpathlineto{\pgfqpoint{1.858572in}{2.300855in}}%
\pgfpathlineto{\pgfqpoint{1.859181in}{2.297856in}}%
\pgfpathlineto{\pgfqpoint{1.859439in}{2.299775in}}%
\pgfpathlineto{\pgfqpoint{1.859546in}{2.297021in}}%
\pgfpathlineto{\pgfqpoint{1.860451in}{2.302989in}}%
\pgfpathlineto{\pgfqpoint{1.860490in}{2.302806in}}%
\pgfpathlineto{\pgfqpoint{1.861600in}{2.304470in}}%
\pgfpathlineto{\pgfqpoint{1.861216in}{2.300684in}}%
\pgfpathlineto{\pgfqpoint{1.861698in}{2.303570in}}%
\pgfpathlineto{\pgfqpoint{1.862632in}{2.301305in}}%
\pgfpathlineto{\pgfqpoint{1.862141in}{2.305177in}}%
\pgfpathlineto{\pgfqpoint{1.862842in}{2.302687in}}%
\pgfpathlineto{\pgfqpoint{1.864336in}{2.298082in}}%
\pgfpathlineto{\pgfqpoint{1.862856in}{2.302758in}}%
\pgfpathlineto{\pgfqpoint{1.864390in}{2.298715in}}%
\pgfpathlineto{\pgfqpoint{1.864701in}{2.298701in}}%
\pgfpathlineto{\pgfqpoint{1.865553in}{2.300617in}}%
\pgfpathlineto{\pgfqpoint{1.865889in}{2.298947in}}%
\pgfpathlineto{\pgfqpoint{1.866503in}{2.302375in}}%
\pgfpathlineto{\pgfqpoint{1.866668in}{2.300091in}}%
\pgfpathlineto{\pgfqpoint{1.866848in}{2.301255in}}%
\pgfpathlineto{\pgfqpoint{1.867097in}{2.298836in}}%
\pgfpathlineto{\pgfqpoint{1.867759in}{2.299750in}}%
\pgfpathlineto{\pgfqpoint{1.868796in}{2.293617in}}%
\pgfpathlineto{\pgfqpoint{1.867803in}{2.300602in}}%
\pgfpathlineto{\pgfqpoint{1.869068in}{2.294504in}}%
\pgfpathlineto{\pgfqpoint{1.869273in}{2.297206in}}%
\pgfpathlineto{\pgfqpoint{1.870519in}{2.303551in}}%
\pgfpathlineto{\pgfqpoint{1.870758in}{2.303427in}}%
\pgfpathlineto{\pgfqpoint{1.872019in}{2.297902in}}%
\pgfpathlineto{\pgfqpoint{1.871332in}{2.303738in}}%
\pgfpathlineto{\pgfqpoint{1.872033in}{2.298042in}}%
\pgfpathlineto{\pgfqpoint{1.872598in}{2.299928in}}%
\pgfpathlineto{\pgfqpoint{1.872978in}{2.296936in}}%
\pgfpathlineto{\pgfqpoint{1.873134in}{2.297570in}}%
\pgfpathlineto{\pgfqpoint{1.874249in}{2.295560in}}%
\pgfpathlineto{\pgfqpoint{1.873397in}{2.298140in}}%
\pgfpathlineto{\pgfqpoint{1.874302in}{2.296351in}}%
\pgfpathlineto{\pgfqpoint{1.875471in}{2.303318in}}%
\pgfpathlineto{\pgfqpoint{1.874361in}{2.296321in}}%
\pgfpathlineto{\pgfqpoint{1.875602in}{2.303096in}}%
\pgfpathlineto{\pgfqpoint{1.876211in}{2.300647in}}%
\pgfpathlineto{\pgfqpoint{1.876595in}{2.304055in}}%
\pgfpathlineto{\pgfqpoint{1.876707in}{2.303232in}}%
\pgfpathlineto{\pgfqpoint{1.877155in}{2.303782in}}%
\pgfpathlineto{\pgfqpoint{1.877540in}{2.300657in}}%
\pgfpathlineto{\pgfqpoint{1.877788in}{2.302186in}}%
\pgfpathlineto{\pgfqpoint{1.877958in}{2.301045in}}%
\pgfpathlineto{\pgfqpoint{1.878601in}{2.304085in}}%
\pgfpathlineto{\pgfqpoint{1.878903in}{2.301781in}}%
\pgfpathlineto{\pgfqpoint{1.881064in}{2.305405in}}%
\pgfpathlineto{\pgfqpoint{1.878913in}{2.301642in}}%
\pgfpathlineto{\pgfqpoint{1.881269in}{2.303429in}}%
\pgfpathlineto{\pgfqpoint{1.881634in}{2.301230in}}%
\pgfpathlineto{\pgfqpoint{1.882165in}{2.306924in}}%
\pgfpathlineto{\pgfqpoint{1.882301in}{2.305781in}}%
\pgfpathlineto{\pgfqpoint{1.882306in}{2.305928in}}%
\pgfpathlineto{\pgfqpoint{1.883211in}{2.301411in}}%
\pgfpathlineto{\pgfqpoint{1.883294in}{2.302218in}}%
\pgfpathlineto{\pgfqpoint{1.883606in}{2.298863in}}%
\pgfpathlineto{\pgfqpoint{1.884458in}{2.299717in}}%
\pgfpathlineto{\pgfqpoint{1.884555in}{2.300131in}}%
\pgfpathlineto{\pgfqpoint{1.884969in}{2.295680in}}%
\pgfpathlineto{\pgfqpoint{1.885544in}{2.298792in}}%
\pgfpathlineto{\pgfqpoint{1.886619in}{2.299888in}}%
\pgfpathlineto{\pgfqpoint{1.886142in}{2.296894in}}%
\pgfpathlineto{\pgfqpoint{1.886658in}{2.298859in}}%
\pgfpathlineto{\pgfqpoint{1.887398in}{2.297923in}}%
\pgfpathlineto{\pgfqpoint{1.887661in}{2.301931in}}%
\pgfpathlineto{\pgfqpoint{1.887720in}{2.301730in}}%
\pgfpathlineto{\pgfqpoint{1.888479in}{2.304584in}}%
\pgfpathlineto{\pgfqpoint{1.887973in}{2.299661in}}%
\pgfpathlineto{\pgfqpoint{1.888878in}{2.303332in}}%
\pgfpathlineto{\pgfqpoint{1.890831in}{2.295561in}}%
\pgfpathlineto{\pgfqpoint{1.889020in}{2.304517in}}%
\pgfpathlineto{\pgfqpoint{1.891001in}{2.297302in}}%
\pgfpathlineto{\pgfqpoint{1.892194in}{2.302887in}}%
\pgfpathlineto{\pgfqpoint{1.891235in}{2.295935in}}%
\pgfpathlineto{\pgfqpoint{1.892208in}{2.302885in}}%
\pgfpathlineto{\pgfqpoint{1.892890in}{2.300608in}}%
\pgfpathlineto{\pgfqpoint{1.892554in}{2.303757in}}%
\pgfpathlineto{\pgfqpoint{1.893343in}{2.301497in}}%
\pgfpathlineto{\pgfqpoint{1.894463in}{2.305934in}}%
\pgfpathlineto{\pgfqpoint{1.894560in}{2.305513in}}%
\pgfpathlineto{\pgfqpoint{1.895407in}{2.298236in}}%
\pgfpathlineto{\pgfqpoint{1.895821in}{2.299594in}}%
\pgfpathlineto{\pgfqpoint{1.895879in}{2.300162in}}%
\pgfpathlineto{\pgfqpoint{1.896532in}{2.295056in}}%
\pgfpathlineto{\pgfqpoint{1.896848in}{2.295831in}}%
\pgfpathlineto{\pgfqpoint{1.896853in}{2.295669in}}%
\pgfpathlineto{\pgfqpoint{1.897578in}{2.300629in}}%
\pgfpathlineto{\pgfqpoint{1.897895in}{2.297834in}}%
\pgfpathlineto{\pgfqpoint{1.898143in}{2.298192in}}%
\pgfpathlineto{\pgfqpoint{1.898031in}{2.296778in}}%
\pgfpathlineto{\pgfqpoint{1.898206in}{2.296745in}}%
\pgfpathlineto{\pgfqpoint{1.898640in}{2.295986in}}%
\pgfpathlineto{\pgfqpoint{1.899248in}{2.299532in}}%
\pgfpathlineto{\pgfqpoint{1.899253in}{2.299353in}}%
\pgfpathlineto{\pgfqpoint{1.899721in}{2.300115in}}%
\pgfpathlineto{\pgfqpoint{1.900290in}{2.298047in}}%
\pgfpathlineto{\pgfqpoint{1.900358in}{2.299094in}}%
\pgfpathlineto{\pgfqpoint{1.900504in}{2.301716in}}%
\pgfpathlineto{\pgfqpoint{1.901162in}{2.298766in}}%
\pgfpathlineto{\pgfqpoint{1.901473in}{2.299335in}}%
\pgfpathlineto{\pgfqpoint{1.902574in}{2.303169in}}%
\pgfpathlineto{\pgfqpoint{1.901678in}{2.297993in}}%
\pgfpathlineto{\pgfqpoint{1.903036in}{2.300945in}}%
\pgfpathlineto{\pgfqpoint{1.903211in}{2.300059in}}%
\pgfpathlineto{\pgfqpoint{1.904068in}{2.303023in}}%
\pgfpathlineto{\pgfqpoint{1.904092in}{2.302304in}}%
\pgfpathlineto{\pgfqpoint{1.904511in}{2.302800in}}%
\pgfpathlineto{\pgfqpoint{1.904433in}{2.300922in}}%
\pgfpathlineto{\pgfqpoint{1.905164in}{2.301595in}}%
\pgfpathlineto{\pgfqpoint{1.905271in}{2.300147in}}%
\pgfpathlineto{\pgfqpoint{1.906186in}{2.303054in}}%
\pgfpathlineto{\pgfqpoint{1.906269in}{2.301618in}}%
\pgfpathlineto{\pgfqpoint{1.906566in}{2.304283in}}%
\pgfpathlineto{\pgfqpoint{1.906308in}{2.300986in}}%
\pgfpathlineto{\pgfqpoint{1.907403in}{2.303010in}}%
\pgfpathlineto{\pgfqpoint{1.908085in}{2.300530in}}%
\pgfpathlineto{\pgfqpoint{1.907656in}{2.304177in}}%
\pgfpathlineto{\pgfqpoint{1.908542in}{2.302928in}}%
\pgfpathlineto{\pgfqpoint{1.909574in}{2.304917in}}%
\pgfpathlineto{\pgfqpoint{1.908893in}{2.302281in}}%
\pgfpathlineto{\pgfqpoint{1.909657in}{2.303394in}}%
\pgfpathlineto{\pgfqpoint{1.910460in}{2.301904in}}%
\pgfpathlineto{\pgfqpoint{1.910344in}{2.304184in}}%
\pgfpathlineto{\pgfqpoint{1.910767in}{2.303155in}}%
\pgfpathlineto{\pgfqpoint{1.910792in}{2.304205in}}%
\pgfpathlineto{\pgfqpoint{1.911853in}{2.301664in}}%
\pgfpathlineto{\pgfqpoint{1.912544in}{2.300135in}}%
\pgfpathlineto{\pgfqpoint{1.912223in}{2.303396in}}%
\pgfpathlineto{\pgfqpoint{1.912953in}{2.301939in}}%
\pgfpathlineto{\pgfqpoint{1.914039in}{2.304367in}}%
\pgfpathlineto{\pgfqpoint{1.913143in}{2.300673in}}%
\pgfpathlineto{\pgfqpoint{1.914560in}{2.303136in}}%
\pgfpathlineto{\pgfqpoint{1.914905in}{2.300361in}}%
\pgfpathlineto{\pgfqpoint{1.915509in}{2.304383in}}%
\pgfpathlineto{\pgfqpoint{1.915665in}{2.303484in}}%
\pgfpathlineto{\pgfqpoint{1.916181in}{2.302735in}}%
\pgfpathlineto{\pgfqpoint{1.915972in}{2.304771in}}%
\pgfpathlineto{\pgfqpoint{1.916770in}{2.303513in}}%
\pgfpathlineto{\pgfqpoint{1.917009in}{2.304841in}}%
\pgfpathlineto{\pgfqpoint{1.917758in}{2.301575in}}%
\pgfpathlineto{\pgfqpoint{1.917856in}{2.301787in}}%
\pgfpathlineto{\pgfqpoint{1.918805in}{2.303809in}}%
\pgfpathlineto{\pgfqpoint{1.918352in}{2.299711in}}%
\pgfpathlineto{\pgfqpoint{1.918995in}{2.302826in}}%
\pgfpathlineto{\pgfqpoint{1.919920in}{2.299626in}}%
\pgfpathlineto{\pgfqpoint{1.919151in}{2.305100in}}%
\pgfpathlineto{\pgfqpoint{1.920193in}{2.300856in}}%
\pgfpathlineto{\pgfqpoint{1.921458in}{2.298473in}}%
\pgfpathlineto{\pgfqpoint{1.920976in}{2.303083in}}%
\pgfpathlineto{\pgfqpoint{1.921478in}{2.298760in}}%
\pgfpathlineto{\pgfqpoint{1.921502in}{2.299550in}}%
\pgfpathlineto{\pgfqpoint{1.922549in}{2.295732in}}%
\pgfpathlineto{\pgfqpoint{1.922559in}{2.296079in}}%
\pgfpathlineto{\pgfqpoint{1.922593in}{2.296217in}}%
\pgfpathlineto{\pgfqpoint{1.923732in}{2.293535in}}%
\pgfpathlineto{\pgfqpoint{1.924005in}{2.296709in}}%
\pgfpathlineto{\pgfqpoint{1.924793in}{2.293266in}}%
\pgfpathlineto{\pgfqpoint{1.924862in}{2.293785in}}%
\pgfpathlineto{\pgfqpoint{1.925042in}{2.291992in}}%
\pgfpathlineto{\pgfqpoint{1.925874in}{2.295955in}}%
\pgfpathlineto{\pgfqpoint{1.925928in}{2.295333in}}%
\pgfpathlineto{\pgfqpoint{1.926979in}{2.298565in}}%
\pgfpathlineto{\pgfqpoint{1.926093in}{2.294477in}}%
\pgfpathlineto{\pgfqpoint{1.927072in}{2.297529in}}%
\pgfpathlineto{\pgfqpoint{1.928060in}{2.293013in}}%
\pgfpathlineto{\pgfqpoint{1.927198in}{2.298552in}}%
\pgfpathlineto{\pgfqpoint{1.928323in}{2.294039in}}%
\pgfpathlineto{\pgfqpoint{1.929375in}{2.297521in}}%
\pgfpathlineto{\pgfqpoint{1.928635in}{2.292538in}}%
\pgfpathlineto{\pgfqpoint{1.929467in}{2.296365in}}%
\pgfpathlineto{\pgfqpoint{1.930455in}{2.299556in}}%
\pgfpathlineto{\pgfqpoint{1.929803in}{2.295791in}}%
\pgfpathlineto{\pgfqpoint{1.930718in}{2.298503in}}%
\pgfpathlineto{\pgfqpoint{1.930728in}{2.298140in}}%
\pgfpathlineto{\pgfqpoint{1.931677in}{2.303347in}}%
\pgfpathlineto{\pgfqpoint{1.931702in}{2.303310in}}%
\pgfpathlineto{\pgfqpoint{1.932822in}{2.305297in}}%
\pgfpathlineto{\pgfqpoint{1.932091in}{2.301883in}}%
\pgfpathlineto{\pgfqpoint{1.932846in}{2.304854in}}%
\pgfpathlineto{\pgfqpoint{1.933386in}{2.299364in}}%
\pgfpathlineto{\pgfqpoint{1.934024in}{2.301186in}}%
\pgfpathlineto{\pgfqpoint{1.934097in}{2.301888in}}%
\pgfpathlineto{\pgfqpoint{1.934267in}{2.298908in}}%
\pgfpathlineto{\pgfqpoint{1.934389in}{2.299102in}}%
\pgfpathlineto{\pgfqpoint{1.934414in}{2.298413in}}%
\pgfpathlineto{\pgfqpoint{1.935270in}{2.303700in}}%
\pgfpathlineto{\pgfqpoint{1.935426in}{2.303333in}}%
\pgfpathlineto{\pgfqpoint{1.935640in}{2.301128in}}%
\pgfpathlineto{\pgfqpoint{1.936395in}{2.304609in}}%
\pgfpathlineto{\pgfqpoint{1.936522in}{2.303499in}}%
\pgfpathlineto{\pgfqpoint{1.936638in}{2.304419in}}%
\pgfpathlineto{\pgfqpoint{1.936940in}{2.302476in}}%
\pgfpathlineto{\pgfqpoint{1.937578in}{2.302730in}}%
\pgfpathlineto{\pgfqpoint{1.938191in}{2.301934in}}%
\pgfpathlineto{\pgfqpoint{1.937977in}{2.304078in}}%
\pgfpathlineto{\pgfqpoint{1.938635in}{2.303262in}}%
\pgfpathlineto{\pgfqpoint{1.938927in}{2.304399in}}%
\pgfpathlineto{\pgfqpoint{1.939657in}{2.302041in}}%
\pgfpathlineto{\pgfqpoint{1.939730in}{2.302814in}}%
\pgfpathlineto{\pgfqpoint{1.939944in}{2.303840in}}%
\pgfpathlineto{\pgfqpoint{1.940499in}{2.301342in}}%
\pgfpathlineto{\pgfqpoint{1.940786in}{2.302493in}}%
\pgfpathlineto{\pgfqpoint{1.941770in}{2.302032in}}%
\pgfpathlineto{\pgfqpoint{1.941332in}{2.304455in}}%
\pgfpathlineto{\pgfqpoint{1.941892in}{2.302831in}}%
\pgfpathlineto{\pgfqpoint{1.942374in}{2.305678in}}%
\pgfpathlineto{\pgfqpoint{1.942821in}{2.302215in}}%
\pgfpathlineto{\pgfqpoint{1.943021in}{2.303447in}}%
\pgfpathlineto{\pgfqpoint{1.943810in}{2.299528in}}%
\pgfpathlineto{\pgfqpoint{1.944209in}{2.301107in}}%
\pgfpathlineto{\pgfqpoint{1.944925in}{2.304476in}}%
\pgfpathlineto{\pgfqpoint{1.944486in}{2.300976in}}%
\pgfpathlineto{\pgfqpoint{1.945373in}{2.302332in}}%
\pgfpathlineto{\pgfqpoint{1.945377in}{2.302088in}}%
\pgfpathlineto{\pgfqpoint{1.946078in}{2.305297in}}%
\pgfpathlineto{\pgfqpoint{1.946458in}{2.303506in}}%
\pgfpathlineto{\pgfqpoint{1.946731in}{2.305015in}}%
\pgfpathlineto{\pgfqpoint{1.947544in}{2.302105in}}%
\pgfpathlineto{\pgfqpoint{1.948128in}{2.299397in}}%
\pgfpathlineto{\pgfqpoint{1.947700in}{2.303305in}}%
\pgfpathlineto{\pgfqpoint{1.948683in}{2.300851in}}%
\pgfpathlineto{\pgfqpoint{1.949871in}{2.303976in}}%
\pgfpathlineto{\pgfqpoint{1.948800in}{2.299228in}}%
\pgfpathlineto{\pgfqpoint{1.949939in}{2.303240in}}%
\pgfpathlineto{\pgfqpoint{1.950830in}{2.300567in}}%
\pgfpathlineto{\pgfqpoint{1.950319in}{2.304889in}}%
\pgfpathlineto{\pgfqpoint{1.951122in}{2.302061in}}%
\pgfpathlineto{\pgfqpoint{1.951531in}{2.305551in}}%
\pgfpathlineto{\pgfqpoint{1.951137in}{2.301925in}}%
\pgfpathlineto{\pgfqpoint{1.952403in}{2.304157in}}%
\pgfpathlineto{\pgfqpoint{1.953547in}{2.300751in}}%
\pgfpathlineto{\pgfqpoint{1.952768in}{2.305363in}}%
\pgfpathlineto{\pgfqpoint{1.953556in}{2.301112in}}%
\pgfpathlineto{\pgfqpoint{1.953815in}{2.303693in}}%
\pgfpathlineto{\pgfqpoint{1.954579in}{2.299950in}}%
\pgfpathlineto{\pgfqpoint{1.954652in}{2.300546in}}%
\pgfpathlineto{\pgfqpoint{1.954676in}{2.300307in}}%
\pgfpathlineto{\pgfqpoint{1.955514in}{2.303598in}}%
\pgfpathlineto{\pgfqpoint{1.955616in}{2.303135in}}%
\pgfpathlineto{\pgfqpoint{1.955621in}{2.303280in}}%
\pgfpathlineto{\pgfqpoint{1.956551in}{2.298591in}}%
\pgfpathlineto{\pgfqpoint{1.957349in}{2.295685in}}%
\pgfpathlineto{\pgfqpoint{1.957670in}{2.296980in}}%
\pgfpathlineto{\pgfqpoint{1.958104in}{2.300054in}}%
\pgfpathlineto{\pgfqpoint{1.958805in}{2.298860in}}%
\pgfpathlineto{\pgfqpoint{1.959701in}{2.302786in}}%
\pgfpathlineto{\pgfqpoint{1.959009in}{2.297844in}}%
\pgfpathlineto{\pgfqpoint{1.959959in}{2.301163in}}%
\pgfpathlineto{\pgfqpoint{1.960533in}{2.300641in}}%
\pgfpathlineto{\pgfqpoint{1.961000in}{2.303648in}}%
\pgfpathlineto{\pgfqpoint{1.961025in}{2.303422in}}%
\pgfpathlineto{\pgfqpoint{1.961336in}{2.300858in}}%
\pgfpathlineto{\pgfqpoint{1.961108in}{2.303935in}}%
\pgfpathlineto{\pgfqpoint{1.962378in}{2.302268in}}%
\pgfpathlineto{\pgfqpoint{1.962602in}{2.303921in}}%
\pgfpathlineto{\pgfqpoint{1.963313in}{2.300127in}}%
\pgfpathlineto{\pgfqpoint{1.963449in}{2.300380in}}%
\pgfpathlineto{\pgfqpoint{1.963951in}{2.303506in}}%
\pgfpathlineto{\pgfqpoint{1.963464in}{2.300322in}}%
\pgfpathlineto{\pgfqpoint{1.965002in}{2.301750in}}%
\pgfpathlineto{\pgfqpoint{1.966039in}{2.300575in}}%
\pgfpathlineto{\pgfqpoint{1.965893in}{2.302316in}}%
\pgfpathlineto{\pgfqpoint{1.966132in}{2.301099in}}%
\pgfpathlineto{\pgfqpoint{1.967266in}{2.304779in}}%
\pgfpathlineto{\pgfqpoint{1.966229in}{2.301076in}}%
\pgfpathlineto{\pgfqpoint{1.967281in}{2.304363in}}%
\pgfpathlineto{\pgfqpoint{1.967962in}{2.301825in}}%
\pgfpathlineto{\pgfqpoint{1.967476in}{2.305145in}}%
\pgfpathlineto{\pgfqpoint{1.968459in}{2.303319in}}%
\pgfpathlineto{\pgfqpoint{1.968761in}{2.305577in}}%
\pgfpathlineto{\pgfqpoint{1.968995in}{2.301476in}}%
\pgfpathlineto{\pgfqpoint{1.969564in}{2.302977in}}%
\pgfpathlineto{\pgfqpoint{1.970158in}{2.299386in}}%
\pgfpathlineto{\pgfqpoint{1.969574in}{2.303018in}}%
\pgfpathlineto{\pgfqpoint{1.970728in}{2.302173in}}%
\pgfpathlineto{\pgfqpoint{1.971882in}{2.305863in}}%
\pgfpathlineto{\pgfqpoint{1.971438in}{2.301235in}}%
\pgfpathlineto{\pgfqpoint{1.972003in}{2.305278in}}%
\pgfpathlineto{\pgfqpoint{1.973099in}{2.302641in}}%
\pgfpathlineto{\pgfqpoint{1.972028in}{2.305907in}}%
\pgfpathlineto{\pgfqpoint{1.973293in}{2.304221in}}%
\pgfpathlineto{\pgfqpoint{1.973298in}{2.304333in}}%
\pgfpathlineto{\pgfqpoint{1.974102in}{2.301054in}}%
\pgfpathlineto{\pgfqpoint{1.974316in}{2.302614in}}%
\pgfpathlineto{\pgfqpoint{1.974540in}{2.301045in}}%
\pgfpathlineto{\pgfqpoint{1.975207in}{2.304152in}}%
\pgfpathlineto{\pgfqpoint{1.975416in}{2.303325in}}%
\pgfpathlineto{\pgfqpoint{1.975742in}{2.303730in}}%
\pgfpathlineto{\pgfqpoint{1.976297in}{2.300070in}}%
\pgfpathlineto{\pgfqpoint{1.976511in}{2.302620in}}%
\pgfpathlineto{\pgfqpoint{1.976589in}{2.303529in}}%
\pgfpathlineto{\pgfqpoint{1.977037in}{2.301224in}}%
\pgfpathlineto{\pgfqpoint{1.977188in}{2.301658in}}%
\pgfpathlineto{\pgfqpoint{1.977510in}{2.299498in}}%
\pgfpathlineto{\pgfqpoint{1.977787in}{2.302850in}}%
\pgfpathlineto{\pgfqpoint{1.978274in}{2.302397in}}%
\pgfpathlineto{\pgfqpoint{1.978512in}{2.304200in}}%
\pgfpathlineto{\pgfqpoint{1.979131in}{2.300360in}}%
\pgfpathlineto{\pgfqpoint{1.979218in}{2.300370in}}%
\pgfpathlineto{\pgfqpoint{1.980148in}{2.297694in}}%
\pgfpathlineto{\pgfqpoint{1.979413in}{2.300558in}}%
\pgfpathlineto{\pgfqpoint{1.980343in}{2.299687in}}%
\pgfpathlineto{\pgfqpoint{1.981249in}{2.303446in}}%
\pgfpathlineto{\pgfqpoint{1.980703in}{2.298388in}}%
\pgfpathlineto{\pgfqpoint{1.981502in}{2.303432in}}%
\pgfpathlineto{\pgfqpoint{1.981570in}{2.305026in}}%
\pgfpathlineto{\pgfqpoint{1.982393in}{2.300751in}}%
\pgfpathlineto{\pgfqpoint{1.982466in}{2.301012in}}%
\pgfpathlineto{\pgfqpoint{1.983590in}{2.297308in}}%
\pgfpathlineto{\pgfqpoint{1.982621in}{2.301821in}}%
\pgfpathlineto{\pgfqpoint{1.983795in}{2.298317in}}%
\pgfpathlineto{\pgfqpoint{1.984379in}{2.300185in}}%
\pgfpathlineto{\pgfqpoint{1.984856in}{2.296850in}}%
\pgfpathlineto{\pgfqpoint{1.985737in}{2.295161in}}%
\pgfpathlineto{\pgfqpoint{1.985119in}{2.298210in}}%
\pgfpathlineto{\pgfqpoint{1.985990in}{2.295969in}}%
\pgfpathlineto{\pgfqpoint{1.986049in}{2.297139in}}%
\pgfpathlineto{\pgfqpoint{1.986998in}{2.293436in}}%
\pgfpathlineto{\pgfqpoint{1.987062in}{2.294142in}}%
\pgfpathlineto{\pgfqpoint{1.987894in}{2.292045in}}%
\pgfpathlineto{\pgfqpoint{1.987553in}{2.295163in}}%
\pgfpathlineto{\pgfqpoint{1.988162in}{2.294111in}}%
\pgfpathlineto{\pgfqpoint{1.988941in}{2.296734in}}%
\pgfpathlineto{\pgfqpoint{1.989179in}{2.294073in}}%
\pgfpathlineto{\pgfqpoint{1.989286in}{2.294935in}}%
\pgfpathlineto{\pgfqpoint{1.989739in}{2.293620in}}%
\pgfpathlineto{\pgfqpoint{1.989525in}{2.296576in}}%
\pgfpathlineto{\pgfqpoint{1.990353in}{2.295601in}}%
\pgfpathlineto{\pgfqpoint{1.991516in}{2.298396in}}%
\pgfpathlineto{\pgfqpoint{1.990844in}{2.295214in}}%
\pgfpathlineto{\pgfqpoint{1.991521in}{2.298330in}}%
\pgfpathlineto{\pgfqpoint{1.992451in}{2.293711in}}%
\pgfpathlineto{\pgfqpoint{1.991623in}{2.299499in}}%
\pgfpathlineto{\pgfqpoint{1.992665in}{2.296014in}}%
\pgfpathlineto{\pgfqpoint{1.992748in}{2.297061in}}%
\pgfpathlineto{\pgfqpoint{1.993128in}{2.295265in}}%
\pgfpathlineto{\pgfqpoint{1.993790in}{2.296240in}}%
\pgfpathlineto{\pgfqpoint{1.993800in}{2.295963in}}%
\pgfpathlineto{\pgfqpoint{1.994632in}{2.299536in}}%
\pgfpathlineto{\pgfqpoint{1.994841in}{2.299269in}}%
\pgfpathlineto{\pgfqpoint{1.994856in}{2.299571in}}%
\pgfpathlineto{\pgfqpoint{1.995007in}{2.297963in}}%
\pgfpathlineto{\pgfqpoint{1.995017in}{2.298152in}}%
\pgfpathlineto{\pgfqpoint{1.995401in}{2.296811in}}%
\pgfpathlineto{\pgfqpoint{1.995971in}{2.303264in}}%
\pgfpathlineto{\pgfqpoint{1.996054in}{2.302871in}}%
\pgfpathlineto{\pgfqpoint{1.996136in}{2.302193in}}%
\pgfpathlineto{\pgfqpoint{1.996614in}{2.304997in}}%
\pgfpathlineto{\pgfqpoint{1.997139in}{2.303637in}}%
\pgfpathlineto{\pgfqpoint{1.997422in}{2.304618in}}%
\pgfpathlineto{\pgfqpoint{1.997733in}{2.301639in}}%
\pgfpathlineto{\pgfqpoint{1.998235in}{2.303323in}}%
\pgfpathlineto{\pgfqpoint{1.998444in}{2.301275in}}%
\pgfpathlineto{\pgfqpoint{1.999072in}{2.304626in}}%
\pgfpathlineto{\pgfqpoint{1.999121in}{2.304507in}}%
\pgfpathlineto{\pgfqpoint{1.999272in}{2.305260in}}%
\pgfpathlineto{\pgfqpoint{1.999744in}{2.302608in}}%
\pgfpathlineto{\pgfqpoint{2.000221in}{2.304287in}}%
\pgfpathlineto{\pgfqpoint{2.001093in}{2.300863in}}%
\pgfpathlineto{\pgfqpoint{2.000275in}{2.304584in}}%
\pgfpathlineto{\pgfqpoint{2.001394in}{2.302406in}}%
\pgfpathlineto{\pgfqpoint{2.001579in}{2.302978in}}%
\pgfpathlineto{\pgfqpoint{2.002027in}{2.300654in}}%
\pgfpathlineto{\pgfqpoint{2.002509in}{2.302698in}}%
\pgfpathlineto{\pgfqpoint{2.003293in}{2.306464in}}%
\pgfpathlineto{\pgfqpoint{2.002582in}{2.302124in}}%
\pgfpathlineto{\pgfqpoint{2.003629in}{2.303078in}}%
\pgfpathlineto{\pgfqpoint{2.003863in}{2.302160in}}%
\pgfpathlineto{\pgfqpoint{2.003989in}{2.304737in}}%
\pgfpathlineto{\pgfqpoint{2.004739in}{2.303135in}}%
\pgfpathlineto{\pgfqpoint{2.005796in}{2.304058in}}%
\pgfpathlineto{\pgfqpoint{2.004992in}{2.301989in}}%
\pgfpathlineto{\pgfqpoint{2.005834in}{2.302954in}}%
\pgfpathlineto{\pgfqpoint{2.006531in}{2.302049in}}%
\pgfpathlineto{\pgfqpoint{2.006331in}{2.305210in}}%
\pgfpathlineto{\pgfqpoint{2.006920in}{2.303959in}}%
\pgfpathlineto{\pgfqpoint{2.006944in}{2.304671in}}%
\pgfpathlineto{\pgfqpoint{2.007188in}{2.302192in}}%
\pgfpathlineto{\pgfqpoint{2.008001in}{2.303455in}}%
\pgfpathlineto{\pgfqpoint{2.008493in}{2.302393in}}%
\pgfpathlineto{\pgfqpoint{2.008790in}{2.304872in}}%
\pgfpathlineto{\pgfqpoint{2.009048in}{2.304439in}}%
\pgfpathlineto{\pgfqpoint{2.009053in}{2.304448in}}%
\pgfpathlineto{\pgfqpoint{2.009247in}{2.302648in}}%
\pgfpathlineto{\pgfqpoint{2.009403in}{2.302905in}}%
\pgfpathlineto{\pgfqpoint{2.010221in}{2.301132in}}%
\pgfpathlineto{\pgfqpoint{2.009632in}{2.305091in}}%
\pgfpathlineto{\pgfqpoint{2.010518in}{2.302623in}}%
\pgfpathlineto{\pgfqpoint{2.010820in}{2.304594in}}%
\pgfpathlineto{\pgfqpoint{2.011516in}{2.300307in}}%
\pgfpathlineto{\pgfqpoint{2.011550in}{2.300528in}}%
\pgfpathlineto{\pgfqpoint{2.012353in}{2.302796in}}%
\pgfpathlineto{\pgfqpoint{2.011609in}{2.300189in}}%
\pgfpathlineto{\pgfqpoint{2.012850in}{2.301785in}}%
\pgfpathlineto{\pgfqpoint{2.012855in}{2.301462in}}%
\pgfpathlineto{\pgfqpoint{2.013497in}{2.306145in}}%
\pgfpathlineto{\pgfqpoint{2.013945in}{2.302419in}}%
\pgfpathlineto{\pgfqpoint{2.014768in}{2.306505in}}%
\pgfpathlineto{\pgfqpoint{2.014213in}{2.301979in}}%
\pgfpathlineto{\pgfqpoint{2.015075in}{2.303478in}}%
\pgfpathlineto{\pgfqpoint{2.015411in}{2.304809in}}%
\pgfpathlineto{\pgfqpoint{2.015158in}{2.302414in}}%
\pgfpathlineto{\pgfqpoint{2.015752in}{2.302355in}}%
\pgfpathlineto{\pgfqpoint{2.016063in}{2.303189in}}%
\pgfpathlineto{\pgfqpoint{2.016954in}{2.300134in}}%
\pgfpathlineto{\pgfqpoint{2.017086in}{2.301069in}}%
\pgfpathlineto{\pgfqpoint{2.017723in}{2.297266in}}%
\pgfpathlineto{\pgfqpoint{2.018025in}{2.298008in}}%
\pgfpathlineto{\pgfqpoint{2.018054in}{2.296660in}}%
\pgfpathlineto{\pgfqpoint{2.018707in}{2.300611in}}%
\pgfpathlineto{\pgfqpoint{2.019111in}{2.299611in}}%
\pgfpathlineto{\pgfqpoint{2.019885in}{2.304203in}}%
\pgfpathlineto{\pgfqpoint{2.019160in}{2.299097in}}%
\pgfpathlineto{\pgfqpoint{2.020294in}{2.302701in}}%
\pgfpathlineto{\pgfqpoint{2.021433in}{2.301858in}}%
\pgfpathlineto{\pgfqpoint{2.020411in}{2.304052in}}%
\pgfpathlineto{\pgfqpoint{2.021457in}{2.302190in}}%
\pgfpathlineto{\pgfqpoint{2.021594in}{2.303544in}}%
\pgfpathlineto{\pgfqpoint{2.021915in}{2.301239in}}%
\pgfpathlineto{\pgfqpoint{2.022558in}{2.301946in}}%
\pgfpathlineto{\pgfqpoint{2.023609in}{2.299979in}}%
\pgfpathlineto{\pgfqpoint{2.023152in}{2.304420in}}%
\pgfpathlineto{\pgfqpoint{2.023721in}{2.301305in}}%
\pgfpathlineto{\pgfqpoint{2.023751in}{2.301427in}}%
\pgfpathlineto{\pgfqpoint{2.024179in}{2.299166in}}%
\pgfpathlineto{\pgfqpoint{2.024783in}{2.300079in}}%
\pgfpathlineto{\pgfqpoint{2.025089in}{2.298748in}}%
\pgfpathlineto{\pgfqpoint{2.025766in}{2.301626in}}%
\pgfpathlineto{\pgfqpoint{2.025873in}{2.301026in}}%
\pgfpathlineto{\pgfqpoint{2.026832in}{2.297286in}}%
\pgfpathlineto{\pgfqpoint{2.025961in}{2.302791in}}%
\pgfpathlineto{\pgfqpoint{2.027227in}{2.300143in}}%
\pgfpathlineto{\pgfqpoint{2.027509in}{2.303758in}}%
\pgfpathlineto{\pgfqpoint{2.027236in}{2.300073in}}%
\pgfpathlineto{\pgfqpoint{2.028385in}{2.303043in}}%
\pgfpathlineto{\pgfqpoint{2.028405in}{2.303301in}}%
\pgfpathlineto{\pgfqpoint{2.028648in}{2.301159in}}%
\pgfpathlineto{\pgfqpoint{2.029388in}{2.301980in}}%
\pgfpathlineto{\pgfqpoint{2.029612in}{2.300562in}}%
\pgfpathlineto{\pgfqpoint{2.029972in}{2.304642in}}%
\pgfpathlineto{\pgfqpoint{2.030489in}{2.302433in}}%
\pgfpathlineto{\pgfqpoint{2.030688in}{2.301487in}}%
\pgfpathlineto{\pgfqpoint{2.031263in}{2.304969in}}%
\pgfpathlineto{\pgfqpoint{2.031613in}{2.302078in}}%
\pgfpathlineto{\pgfqpoint{2.032095in}{2.300669in}}%
\pgfpathlineto{\pgfqpoint{2.032801in}{2.304734in}}%
\pgfpathlineto{\pgfqpoint{2.032952in}{2.305098in}}%
\pgfpathlineto{\pgfqpoint{2.033181in}{2.303303in}}%
\pgfpathlineto{\pgfqpoint{2.033283in}{2.303484in}}%
\pgfpathlineto{\pgfqpoint{2.033463in}{2.302688in}}%
\pgfpathlineto{\pgfqpoint{2.034164in}{2.305951in}}%
\pgfpathlineto{\pgfqpoint{2.034340in}{2.304878in}}%
\pgfpathlineto{\pgfqpoint{2.035537in}{2.301518in}}%
\pgfpathlineto{\pgfqpoint{2.034359in}{2.305233in}}%
\pgfpathlineto{\pgfqpoint{2.035610in}{2.303425in}}%
\pgfpathlineto{\pgfqpoint{2.036452in}{2.303848in}}%
\pgfpathlineto{\pgfqpoint{2.036199in}{2.302227in}}%
\pgfpathlineto{\pgfqpoint{2.036662in}{2.303058in}}%
\pgfpathlineto{\pgfqpoint{2.037105in}{2.300746in}}%
\pgfpathlineto{\pgfqpoint{2.037446in}{2.304008in}}%
\pgfpathlineto{\pgfqpoint{2.037806in}{2.301381in}}%
\pgfpathlineto{\pgfqpoint{2.038794in}{2.303783in}}%
\pgfpathlineto{\pgfqpoint{2.038234in}{2.300889in}}%
\pgfpathlineto{\pgfqpoint{2.038965in}{2.302775in}}%
\pgfpathlineto{\pgfqpoint{2.040172in}{2.294111in}}%
\pgfpathlineto{\pgfqpoint{2.040187in}{2.294218in}}%
\pgfpathlineto{\pgfqpoint{2.041058in}{2.299109in}}%
\pgfpathlineto{\pgfqpoint{2.040196in}{2.294102in}}%
\pgfpathlineto{\pgfqpoint{2.041389in}{2.298002in}}%
\pgfpathlineto{\pgfqpoint{2.043648in}{2.304385in}}%
\pgfpathlineto{\pgfqpoint{2.044184in}{2.302776in}}%
\pgfpathlineto{\pgfqpoint{2.044281in}{2.302513in}}%
\pgfpathlineto{\pgfqpoint{2.044320in}{2.303547in}}%
\pgfpathlineto{\pgfqpoint{2.044524in}{2.303647in}}%
\pgfpathlineto{\pgfqpoint{2.045011in}{2.304925in}}%
\pgfpathlineto{\pgfqpoint{2.045469in}{2.301773in}}%
\pgfpathlineto{\pgfqpoint{2.045625in}{2.302798in}}%
\pgfpathlineto{\pgfqpoint{2.045800in}{2.301749in}}%
\pgfpathlineto{\pgfqpoint{2.046404in}{2.303592in}}%
\pgfpathlineto{\pgfqpoint{2.046744in}{2.302175in}}%
\pgfpathlineto{\pgfqpoint{2.046944in}{2.301894in}}%
\pgfpathlineto{\pgfqpoint{2.047251in}{2.304325in}}%
\pgfpathlineto{\pgfqpoint{2.047796in}{2.303060in}}%
\pgfpathlineto{\pgfqpoint{2.047840in}{2.304083in}}%
\pgfpathlineto{\pgfqpoint{2.048536in}{2.299971in}}%
\pgfpathlineto{\pgfqpoint{2.048711in}{2.300108in}}%
\pgfpathlineto{\pgfqpoint{2.049442in}{2.297412in}}%
\pgfpathlineto{\pgfqpoint{2.049865in}{2.298326in}}%
\pgfpathlineto{\pgfqpoint{2.050698in}{2.301218in}}%
\pgfpathlineto{\pgfqpoint{2.050152in}{2.298114in}}%
\pgfpathlineto{\pgfqpoint{2.051004in}{2.299264in}}%
\pgfpathlineto{\pgfqpoint{2.051019in}{2.298715in}}%
\pgfpathlineto{\pgfqpoint{2.052061in}{2.301834in}}%
\pgfpathlineto{\pgfqpoint{2.052071in}{2.301467in}}%
\pgfpathlineto{\pgfqpoint{2.052777in}{2.304003in}}%
\pgfpathlineto{\pgfqpoint{2.052085in}{2.301159in}}%
\pgfpathlineto{\pgfqpoint{2.053200in}{2.301765in}}%
\pgfpathlineto{\pgfqpoint{2.054446in}{2.296427in}}%
\pgfpathlineto{\pgfqpoint{2.053371in}{2.301835in}}%
\pgfpathlineto{\pgfqpoint{2.054568in}{2.297482in}}%
\pgfpathlineto{\pgfqpoint{2.055780in}{2.301420in}}%
\pgfpathlineto{\pgfqpoint{2.054743in}{2.296487in}}%
\pgfpathlineto{\pgfqpoint{2.055829in}{2.300795in}}%
\pgfpathlineto{\pgfqpoint{2.056511in}{2.296326in}}%
\pgfpathlineto{\pgfqpoint{2.056009in}{2.301800in}}%
\pgfpathlineto{\pgfqpoint{2.056959in}{2.300013in}}%
\pgfpathlineto{\pgfqpoint{2.058108in}{2.301868in}}%
\pgfpathlineto{\pgfqpoint{2.057411in}{2.298563in}}%
\pgfpathlineto{\pgfqpoint{2.058117in}{2.301586in}}%
\pgfpathlineto{\pgfqpoint{2.058755in}{2.299064in}}%
\pgfpathlineto{\pgfqpoint{2.058565in}{2.302521in}}%
\pgfpathlineto{\pgfqpoint{2.059213in}{2.301920in}}%
\pgfpathlineto{\pgfqpoint{2.059437in}{2.302160in}}%
\pgfpathlineto{\pgfqpoint{2.059846in}{2.299670in}}%
\pgfpathlineto{\pgfqpoint{2.059875in}{2.299679in}}%
\pgfpathlineto{\pgfqpoint{2.061014in}{2.297352in}}%
\pgfpathlineto{\pgfqpoint{2.060172in}{2.300347in}}%
\pgfpathlineto{\pgfqpoint{2.061038in}{2.297446in}}%
\pgfpathlineto{\pgfqpoint{2.061345in}{2.295244in}}%
\pgfpathlineto{\pgfqpoint{2.061647in}{2.298479in}}%
\pgfpathlineto{\pgfqpoint{2.061910in}{2.298286in}}%
\pgfpathlineto{\pgfqpoint{2.063025in}{2.301475in}}%
\pgfpathlineto{\pgfqpoint{2.061973in}{2.298126in}}%
\pgfpathlineto{\pgfqpoint{2.063083in}{2.301009in}}%
\pgfpathlineto{\pgfqpoint{2.063288in}{2.300636in}}%
\pgfpathlineto{\pgfqpoint{2.063687in}{2.303566in}}%
\pgfpathlineto{\pgfqpoint{2.064183in}{2.301194in}}%
\pgfpathlineto{\pgfqpoint{2.064305in}{2.301813in}}%
\pgfpathlineto{\pgfqpoint{2.065211in}{2.299134in}}%
\pgfpathlineto{\pgfqpoint{2.065245in}{2.299547in}}%
\pgfpathlineto{\pgfqpoint{2.065396in}{2.297781in}}%
\pgfpathlineto{\pgfqpoint{2.066004in}{2.301789in}}%
\pgfpathlineto{\pgfqpoint{2.066068in}{2.301749in}}%
\pgfpathlineto{\pgfqpoint{2.066175in}{2.302298in}}%
\pgfpathlineto{\pgfqpoint{2.066754in}{2.298644in}}%
\pgfpathlineto{\pgfqpoint{2.066988in}{2.299550in}}%
\pgfpathlineto{\pgfqpoint{2.067105in}{2.298287in}}%
\pgfpathlineto{\pgfqpoint{2.067402in}{2.301749in}}%
\pgfpathlineto{\pgfqpoint{2.067479in}{2.301700in}}%
\pgfpathlineto{\pgfqpoint{2.068492in}{2.304014in}}%
\pgfpathlineto{\pgfqpoint{2.068628in}{2.303254in}}%
\pgfpathlineto{\pgfqpoint{2.069169in}{2.304300in}}%
\pgfpathlineto{\pgfqpoint{2.069295in}{2.302282in}}%
\pgfpathlineto{\pgfqpoint{2.069768in}{2.303718in}}%
\pgfpathlineto{\pgfqpoint{2.070931in}{2.301515in}}%
\pgfpathlineto{\pgfqpoint{2.070206in}{2.305265in}}%
\pgfpathlineto{\pgfqpoint{2.070951in}{2.301601in}}%
\pgfpathlineto{\pgfqpoint{2.072338in}{2.295701in}}%
\pgfpathlineto{\pgfqpoint{2.070965in}{2.301755in}}%
\pgfpathlineto{\pgfqpoint{2.072679in}{2.299188in}}%
\pgfpathlineto{\pgfqpoint{2.072908in}{2.300967in}}%
\pgfpathlineto{\pgfqpoint{2.073731in}{2.296525in}}%
\pgfpathlineto{\pgfqpoint{2.073755in}{2.296795in}}%
\pgfpathlineto{\pgfqpoint{2.073891in}{2.295561in}}%
\pgfpathlineto{\pgfqpoint{2.074271in}{2.300510in}}%
\pgfpathlineto{\pgfqpoint{2.074772in}{2.298281in}}%
\pgfpathlineto{\pgfqpoint{2.075639in}{2.300894in}}%
\pgfpathlineto{\pgfqpoint{2.075254in}{2.296978in}}%
\pgfpathlineto{\pgfqpoint{2.075921in}{2.300474in}}%
\pgfpathlineto{\pgfqpoint{2.077066in}{2.299148in}}%
\pgfpathlineto{\pgfqpoint{2.076472in}{2.303192in}}%
\pgfpathlineto{\pgfqpoint{2.077070in}{2.299357in}}%
\pgfpathlineto{\pgfqpoint{2.077104in}{2.300174in}}%
\pgfpathlineto{\pgfqpoint{2.077942in}{2.297659in}}%
\pgfpathlineto{\pgfqpoint{2.078161in}{2.298616in}}%
\pgfpathlineto{\pgfqpoint{2.078326in}{2.297275in}}%
\pgfpathlineto{\pgfqpoint{2.079081in}{2.302968in}}%
\pgfpathlineto{\pgfqpoint{2.079203in}{2.302433in}}%
\pgfpathlineto{\pgfqpoint{2.079592in}{2.303533in}}%
\pgfpathlineto{\pgfqpoint{2.079232in}{2.301525in}}%
\pgfpathlineto{\pgfqpoint{2.080118in}{2.301631in}}%
\pgfpathlineto{\pgfqpoint{2.080172in}{2.300729in}}%
\pgfpathlineto{\pgfqpoint{2.081058in}{2.303815in}}%
\pgfpathlineto{\pgfqpoint{2.081092in}{2.303570in}}%
\pgfpathlineto{\pgfqpoint{2.081097in}{2.303596in}}%
\pgfpathlineto{\pgfqpoint{2.081637in}{2.300654in}}%
\pgfpathlineto{\pgfqpoint{2.082002in}{2.297956in}}%
\pgfpathlineto{\pgfqpoint{2.082771in}{2.298571in}}%
\pgfpathlineto{\pgfqpoint{2.083297in}{2.296387in}}%
\pgfpathlineto{\pgfqpoint{2.082854in}{2.299039in}}%
\pgfpathlineto{\pgfqpoint{2.083862in}{2.298639in}}%
\pgfpathlineto{\pgfqpoint{2.084183in}{2.301894in}}%
\pgfpathlineto{\pgfqpoint{2.084991in}{2.300415in}}%
\pgfpathlineto{\pgfqpoint{2.086189in}{2.304506in}}%
\pgfpathlineto{\pgfqpoint{2.085035in}{2.299901in}}%
\pgfpathlineto{\pgfqpoint{2.086248in}{2.303898in}}%
\pgfpathlineto{\pgfqpoint{2.088351in}{2.295515in}}%
\pgfpathlineto{\pgfqpoint{2.086355in}{2.304730in}}%
\pgfpathlineto{\pgfqpoint{2.088492in}{2.296984in}}%
\pgfpathlineto{\pgfqpoint{2.089909in}{2.304292in}}%
\pgfpathlineto{\pgfqpoint{2.090045in}{2.303464in}}%
\pgfpathlineto{\pgfqpoint{2.090454in}{2.301243in}}%
\pgfpathlineto{\pgfqpoint{2.091024in}{2.303938in}}%
\pgfpathlineto{\pgfqpoint{2.091174in}{2.301708in}}%
\pgfpathlineto{\pgfqpoint{2.091871in}{2.305253in}}%
\pgfpathlineto{\pgfqpoint{2.091184in}{2.301555in}}%
\pgfpathlineto{\pgfqpoint{2.092333in}{2.302941in}}%
\pgfpathlineto{\pgfqpoint{2.093029in}{2.300883in}}%
\pgfpathlineto{\pgfqpoint{2.092776in}{2.303554in}}%
\pgfpathlineto{\pgfqpoint{2.093472in}{2.301426in}}%
\pgfpathlineto{\pgfqpoint{2.093701in}{2.299843in}}%
\pgfpathlineto{\pgfqpoint{2.094276in}{2.303151in}}%
\pgfpathlineto{\pgfqpoint{2.094553in}{2.301704in}}%
\pgfpathlineto{\pgfqpoint{2.095697in}{2.304338in}}%
\pgfpathlineto{\pgfqpoint{2.095240in}{2.300291in}}%
\pgfpathlineto{\pgfqpoint{2.095741in}{2.304221in}}%
\pgfpathlineto{\pgfqpoint{2.096574in}{2.301378in}}%
\pgfpathlineto{\pgfqpoint{2.096160in}{2.305492in}}%
\pgfpathlineto{\pgfqpoint{2.096885in}{2.303934in}}%
\pgfpathlineto{\pgfqpoint{2.096895in}{2.304248in}}%
\pgfpathlineto{\pgfqpoint{2.097455in}{2.301059in}}%
\pgfpathlineto{\pgfqpoint{2.097981in}{2.303281in}}%
\pgfpathlineto{\pgfqpoint{2.098467in}{2.301362in}}%
\pgfpathlineto{\pgfqpoint{2.099061in}{2.303805in}}%
\pgfpathlineto{\pgfqpoint{2.099071in}{2.303758in}}%
\pgfpathlineto{\pgfqpoint{2.099354in}{2.305225in}}%
\pgfpathlineto{\pgfqpoint{2.100089in}{2.302472in}}%
\pgfpathlineto{\pgfqpoint{2.100152in}{2.302519in}}%
\pgfpathlineto{\pgfqpoint{2.100259in}{2.301173in}}%
\pgfpathlineto{\pgfqpoint{2.101174in}{2.304323in}}%
\pgfpathlineto{\pgfqpoint{2.101247in}{2.303442in}}%
\pgfpathlineto{\pgfqpoint{2.101627in}{2.304206in}}%
\pgfpathlineto{\pgfqpoint{2.102056in}{2.300060in}}%
\pgfpathlineto{\pgfqpoint{2.102279in}{2.301739in}}%
\pgfpathlineto{\pgfqpoint{2.102732in}{2.300402in}}%
\pgfpathlineto{\pgfqpoint{2.103058in}{2.304059in}}%
\pgfpathlineto{\pgfqpoint{2.103346in}{2.303798in}}%
\pgfpathlineto{\pgfqpoint{2.103370in}{2.304007in}}%
\pgfpathlineto{\pgfqpoint{2.103599in}{2.301219in}}%
\pgfpathlineto{\pgfqpoint{2.104431in}{2.302930in}}%
\pgfpathlineto{\pgfqpoint{2.104860in}{2.301165in}}%
\pgfpathlineto{\pgfqpoint{2.105444in}{2.303656in}}%
\pgfpathlineto{\pgfqpoint{2.105517in}{2.303340in}}%
\pgfpathlineto{\pgfqpoint{2.105795in}{2.301914in}}%
\pgfpathlineto{\pgfqpoint{2.105999in}{2.304076in}}%
\pgfpathlineto{\pgfqpoint{2.106389in}{2.304179in}}%
\pgfpathlineto{\pgfqpoint{2.106437in}{2.305012in}}%
\pgfpathlineto{\pgfqpoint{2.107051in}{2.300493in}}%
\pgfpathlineto{\pgfqpoint{2.107469in}{2.303160in}}%
\pgfpathlineto{\pgfqpoint{2.107606in}{2.304255in}}%
\pgfpathlineto{\pgfqpoint{2.108297in}{2.301295in}}%
\pgfpathlineto{\pgfqpoint{2.108302in}{2.301512in}}%
\pgfpathlineto{\pgfqpoint{2.108389in}{2.300592in}}%
\pgfpathlineto{\pgfqpoint{2.109222in}{2.304601in}}%
\pgfpathlineto{\pgfqpoint{2.109397in}{2.302935in}}%
\pgfpathlineto{\pgfqpoint{2.109402in}{2.303068in}}%
\pgfpathlineto{\pgfqpoint{2.110313in}{2.300354in}}%
\pgfpathlineto{\pgfqpoint{2.110459in}{2.301666in}}%
\pgfpathlineto{\pgfqpoint{2.110639in}{2.300974in}}%
\pgfpathlineto{\pgfqpoint{2.111359in}{2.304581in}}%
\pgfpathlineto{\pgfqpoint{2.111393in}{2.304175in}}%
\pgfpathlineto{\pgfqpoint{2.111617in}{2.306340in}}%
\pgfpathlineto{\pgfqpoint{2.112425in}{2.301609in}}%
\pgfpathlineto{\pgfqpoint{2.112484in}{2.302439in}}%
\pgfpathlineto{\pgfqpoint{2.112810in}{2.304498in}}%
\pgfpathlineto{\pgfqpoint{2.113302in}{2.300905in}}%
\pgfpathlineto{\pgfqpoint{2.113584in}{2.301990in}}%
\pgfpathlineto{\pgfqpoint{2.113974in}{2.299787in}}%
\pgfpathlineto{\pgfqpoint{2.114246in}{2.302731in}}%
\pgfpathlineto{\pgfqpoint{2.114723in}{2.301024in}}%
\pgfpathlineto{\pgfqpoint{2.115751in}{2.297362in}}%
\pgfpathlineto{\pgfqpoint{2.115916in}{2.299978in}}%
\pgfpathlineto{\pgfqpoint{2.116749in}{2.300818in}}%
\pgfpathlineto{\pgfqpoint{2.116164in}{2.298315in}}%
\pgfpathlineto{\pgfqpoint{2.117016in}{2.299644in}}%
\pgfpathlineto{\pgfqpoint{2.118136in}{2.297022in}}%
\pgfpathlineto{\pgfqpoint{2.117172in}{2.300424in}}%
\pgfpathlineto{\pgfqpoint{2.118190in}{2.297718in}}%
\pgfpathlineto{\pgfqpoint{2.118350in}{2.297270in}}%
\pgfpathlineto{\pgfqpoint{2.119008in}{2.301442in}}%
\pgfpathlineto{\pgfqpoint{2.119090in}{2.301344in}}%
\pgfpathlineto{\pgfqpoint{2.119611in}{2.304317in}}%
\pgfpathlineto{\pgfqpoint{2.120220in}{2.302729in}}%
\pgfpathlineto{\pgfqpoint{2.120936in}{2.303998in}}%
\pgfpathlineto{\pgfqpoint{2.120687in}{2.300154in}}%
\pgfpathlineto{\pgfqpoint{2.121359in}{2.302968in}}%
\pgfpathlineto{\pgfqpoint{2.121744in}{2.301207in}}%
\pgfpathlineto{\pgfqpoint{2.121549in}{2.304907in}}%
\pgfpathlineto{\pgfqpoint{2.122464in}{2.303056in}}%
\pgfpathlineto{\pgfqpoint{2.122630in}{2.303838in}}%
\pgfpathlineto{\pgfqpoint{2.122859in}{2.300972in}}%
\pgfpathlineto{\pgfqpoint{2.122868in}{2.301185in}}%
\pgfpathlineto{\pgfqpoint{2.124003in}{2.297964in}}%
\pgfpathlineto{\pgfqpoint{2.122980in}{2.302438in}}%
\pgfpathlineto{\pgfqpoint{2.124022in}{2.298265in}}%
\pgfpathlineto{\pgfqpoint{2.125497in}{2.290253in}}%
\pgfpathlineto{\pgfqpoint{2.125789in}{2.292132in}}%
\pgfpathlineto{\pgfqpoint{2.125819in}{2.292790in}}%
\pgfpathlineto{\pgfqpoint{2.126310in}{2.286651in}}%
\pgfpathlineto{\pgfqpoint{2.126885in}{2.291141in}}%
\pgfpathlineto{\pgfqpoint{2.127255in}{2.290105in}}%
\pgfpathlineto{\pgfqpoint{2.127844in}{2.296987in}}%
\pgfpathlineto{\pgfqpoint{2.127883in}{2.296740in}}%
\pgfpathlineto{\pgfqpoint{2.129163in}{2.301344in}}%
\pgfpathlineto{\pgfqpoint{2.129168in}{2.301024in}}%
\pgfpathlineto{\pgfqpoint{2.129509in}{2.298972in}}%
\pgfpathlineto{\pgfqpoint{2.129937in}{2.301696in}}%
\pgfpathlineto{\pgfqpoint{2.130307in}{2.300217in}}%
\pgfpathlineto{\pgfqpoint{2.131106in}{2.303881in}}%
\pgfpathlineto{\pgfqpoint{2.131461in}{2.302925in}}%
\pgfpathlineto{\pgfqpoint{2.132123in}{2.304337in}}%
\pgfpathlineto{\pgfqpoint{2.131700in}{2.301486in}}%
\pgfpathlineto{\pgfqpoint{2.132552in}{2.302680in}}%
\pgfpathlineto{\pgfqpoint{2.133248in}{2.300501in}}%
\pgfpathlineto{\pgfqpoint{2.132625in}{2.303400in}}%
\pgfpathlineto{\pgfqpoint{2.133657in}{2.302845in}}%
\pgfpathlineto{\pgfqpoint{2.134796in}{2.304717in}}%
\pgfpathlineto{\pgfqpoint{2.134451in}{2.302233in}}%
\pgfpathlineto{\pgfqpoint{2.134816in}{2.304664in}}%
\pgfpathlineto{\pgfqpoint{2.135020in}{2.302905in}}%
\pgfpathlineto{\pgfqpoint{2.134840in}{2.305047in}}%
\pgfpathlineto{\pgfqpoint{2.135945in}{2.303804in}}%
\pgfpathlineto{\pgfqpoint{2.135965in}{2.304358in}}%
\pgfpathlineto{\pgfqpoint{2.136758in}{2.301136in}}%
\pgfpathlineto{\pgfqpoint{2.136826in}{2.301478in}}%
\pgfpathlineto{\pgfqpoint{2.137216in}{2.299244in}}%
\pgfpathlineto{\pgfqpoint{2.137644in}{2.303002in}}%
\pgfpathlineto{\pgfqpoint{2.137970in}{2.299858in}}%
\pgfpathlineto{\pgfqpoint{2.138000in}{2.299722in}}%
\pgfpathlineto{\pgfqpoint{2.138555in}{2.303020in}}%
\pgfpathlineto{\pgfqpoint{2.139012in}{2.301625in}}%
\pgfpathlineto{\pgfqpoint{2.139022in}{2.301925in}}%
\pgfpathlineto{\pgfqpoint{2.139674in}{2.298166in}}%
\pgfpathlineto{\pgfqpoint{2.139718in}{2.298516in}}%
\pgfpathlineto{\pgfqpoint{2.140570in}{2.295147in}}%
\pgfpathlineto{\pgfqpoint{2.139777in}{2.299124in}}%
\pgfpathlineto{\pgfqpoint{2.140974in}{2.296091in}}%
\pgfpathlineto{\pgfqpoint{2.142308in}{2.301529in}}%
\pgfpathlineto{\pgfqpoint{2.141052in}{2.295552in}}%
\pgfpathlineto{\pgfqpoint{2.142323in}{2.301431in}}%
\pgfpathlineto{\pgfqpoint{2.143375in}{2.300411in}}%
\pgfpathlineto{\pgfqpoint{2.143039in}{2.303528in}}%
\pgfpathlineto{\pgfqpoint{2.143443in}{2.300997in}}%
\pgfpathlineto{\pgfqpoint{2.143745in}{2.302157in}}%
\pgfpathlineto{\pgfqpoint{2.143842in}{2.300097in}}%
\pgfpathlineto{\pgfqpoint{2.144523in}{2.300730in}}%
\pgfpathlineto{\pgfqpoint{2.144558in}{2.299959in}}%
\pgfpathlineto{\pgfqpoint{2.145492in}{2.304582in}}%
\pgfpathlineto{\pgfqpoint{2.145570in}{2.303316in}}%
\pgfpathlineto{\pgfqpoint{2.146559in}{2.304121in}}%
\pgfpathlineto{\pgfqpoint{2.145823in}{2.301065in}}%
\pgfpathlineto{\pgfqpoint{2.146685in}{2.303790in}}%
\pgfpathlineto{\pgfqpoint{2.147673in}{2.301954in}}%
\pgfpathlineto{\pgfqpoint{2.146904in}{2.304324in}}%
\pgfpathlineto{\pgfqpoint{2.147810in}{2.303217in}}%
\pgfpathlineto{\pgfqpoint{2.147917in}{2.303838in}}%
\pgfpathlineto{\pgfqpoint{2.148525in}{2.298571in}}%
\pgfpathlineto{\pgfqpoint{2.148822in}{2.299967in}}%
\pgfpathlineto{\pgfqpoint{2.149723in}{2.296885in}}%
\pgfpathlineto{\pgfqpoint{2.149947in}{2.298594in}}%
\pgfpathlineto{\pgfqpoint{2.149996in}{2.297843in}}%
\pgfpathlineto{\pgfqpoint{2.150273in}{2.301157in}}%
\pgfpathlineto{\pgfqpoint{2.150371in}{2.300824in}}%
\pgfpathlineto{\pgfqpoint{2.150419in}{2.301997in}}%
\pgfpathlineto{\pgfqpoint{2.151403in}{2.298316in}}%
\pgfpathlineto{\pgfqpoint{2.151446in}{2.298474in}}%
\pgfpathlineto{\pgfqpoint{2.151461in}{2.298274in}}%
\pgfpathlineto{\pgfqpoint{2.152386in}{2.303665in}}%
\pgfpathlineto{\pgfqpoint{2.152401in}{2.303427in}}%
\pgfpathlineto{\pgfqpoint{2.152902in}{2.304536in}}%
\pgfpathlineto{\pgfqpoint{2.152513in}{2.301439in}}%
\pgfpathlineto{\pgfqpoint{2.153496in}{2.303210in}}%
\pgfpathlineto{\pgfqpoint{2.153973in}{2.301738in}}%
\pgfpathlineto{\pgfqpoint{2.154431in}{2.303773in}}%
\pgfpathlineto{\pgfqpoint{2.154616in}{2.302894in}}%
\pgfpathlineto{\pgfqpoint{2.155385in}{2.303718in}}%
\pgfpathlineto{\pgfqpoint{2.155619in}{2.300977in}}%
\pgfpathlineto{\pgfqpoint{2.155672in}{2.301106in}}%
\pgfpathlineto{\pgfqpoint{2.157347in}{2.305235in}}%
\pgfpathlineto{\pgfqpoint{2.155877in}{2.300736in}}%
\pgfpathlineto{\pgfqpoint{2.157410in}{2.303959in}}%
\pgfpathlineto{\pgfqpoint{2.158272in}{2.302527in}}%
\pgfpathlineto{\pgfqpoint{2.158014in}{2.304629in}}%
\pgfpathlineto{\pgfqpoint{2.158535in}{2.303118in}}%
\pgfpathlineto{\pgfqpoint{2.159377in}{2.305314in}}%
\pgfpathlineto{\pgfqpoint{2.158618in}{2.302567in}}%
\pgfpathlineto{\pgfqpoint{2.159674in}{2.305116in}}%
\pgfpathlineto{\pgfqpoint{2.159699in}{2.305514in}}%
\pgfpathlineto{\pgfqpoint{2.160687in}{2.301811in}}%
\pgfpathlineto{\pgfqpoint{2.160706in}{2.302455in}}%
\pgfpathlineto{\pgfqpoint{2.161018in}{2.302036in}}%
\pgfpathlineto{\pgfqpoint{2.161373in}{2.303664in}}%
\pgfpathlineto{\pgfqpoint{2.161714in}{2.303455in}}%
\pgfpathlineto{\pgfqpoint{2.162245in}{2.305390in}}%
\pgfpathlineto{\pgfqpoint{2.162775in}{2.302556in}}%
\pgfpathlineto{\pgfqpoint{2.162800in}{2.302936in}}%
\pgfpathlineto{\pgfqpoint{2.163666in}{2.301986in}}%
\pgfpathlineto{\pgfqpoint{2.163228in}{2.305063in}}%
\pgfpathlineto{\pgfqpoint{2.163900in}{2.303618in}}%
\pgfpathlineto{\pgfqpoint{2.165526in}{2.307926in}}%
\pgfpathlineto{\pgfqpoint{2.164470in}{2.302441in}}%
\pgfpathlineto{\pgfqpoint{2.165604in}{2.307109in}}%
\pgfpathlineto{\pgfqpoint{2.166018in}{2.311231in}}%
\pgfpathlineto{\pgfqpoint{2.166729in}{2.306436in}}%
\pgfpathlineto{\pgfqpoint{2.166826in}{2.307719in}}%
\pgfpathlineto{\pgfqpoint{2.167605in}{2.302009in}}%
\pgfpathlineto{\pgfqpoint{2.167771in}{2.302718in}}%
\pgfpathlineto{\pgfqpoint{2.168442in}{2.298747in}}%
\pgfpathlineto{\pgfqpoint{2.167844in}{2.302917in}}%
\pgfpathlineto{\pgfqpoint{2.168929in}{2.301254in}}%
\pgfpathlineto{\pgfqpoint{2.169153in}{2.302215in}}%
\pgfpathlineto{\pgfqpoint{2.169338in}{2.299524in}}%
\pgfpathlineto{\pgfqpoint{2.170030in}{2.300888in}}%
\pgfpathlineto{\pgfqpoint{2.170156in}{2.300265in}}%
\pgfpathlineto{\pgfqpoint{2.170731in}{2.303835in}}%
\pgfpathlineto{\pgfqpoint{2.171125in}{2.300887in}}%
\pgfpathlineto{\pgfqpoint{2.171169in}{2.302150in}}%
\pgfpathlineto{\pgfqpoint{2.172201in}{2.298970in}}%
\pgfpathlineto{\pgfqpoint{2.172216in}{2.299420in}}%
\pgfpathlineto{\pgfqpoint{2.172629in}{2.297245in}}%
\pgfpathlineto{\pgfqpoint{2.172994in}{2.299809in}}%
\pgfpathlineto{\pgfqpoint{2.173335in}{2.298496in}}%
\pgfpathlineto{\pgfqpoint{2.173379in}{2.298988in}}%
\pgfpathlineto{\pgfqpoint{2.174231in}{2.293187in}}%
\pgfpathlineto{\pgfqpoint{2.174285in}{2.293961in}}%
\pgfpathlineto{\pgfqpoint{2.175414in}{2.292845in}}%
\pgfpathlineto{\pgfqpoint{2.174840in}{2.294872in}}%
\pgfpathlineto{\pgfqpoint{2.175434in}{2.293118in}}%
\pgfpathlineto{\pgfqpoint{2.176753in}{2.295725in}}%
\pgfpathlineto{\pgfqpoint{2.175706in}{2.291416in}}%
\pgfpathlineto{\pgfqpoint{2.176821in}{2.295051in}}%
\pgfpathlineto{\pgfqpoint{2.177240in}{2.291552in}}%
\pgfpathlineto{\pgfqpoint{2.176865in}{2.295324in}}%
\pgfpathlineto{\pgfqpoint{2.177936in}{2.294323in}}%
\pgfpathlineto{\pgfqpoint{2.178871in}{2.298914in}}%
\pgfpathlineto{\pgfqpoint{2.178394in}{2.293941in}}%
\pgfpathlineto{\pgfqpoint{2.179104in}{2.298370in}}%
\pgfpathlineto{\pgfqpoint{2.179771in}{2.295315in}}%
\pgfpathlineto{\pgfqpoint{2.179163in}{2.298590in}}%
\pgfpathlineto{\pgfqpoint{2.180297in}{2.296250in}}%
\pgfpathlineto{\pgfqpoint{2.180701in}{2.300149in}}%
\pgfpathlineto{\pgfqpoint{2.180443in}{2.295767in}}%
\pgfpathlineto{\pgfqpoint{2.181475in}{2.299213in}}%
\pgfpathlineto{\pgfqpoint{2.181617in}{2.298776in}}%
\pgfpathlineto{\pgfqpoint{2.182079in}{2.301237in}}%
\pgfpathlineto{\pgfqpoint{2.182585in}{2.299077in}}%
\pgfpathlineto{\pgfqpoint{2.182790in}{2.299774in}}%
\pgfpathlineto{\pgfqpoint{2.183325in}{2.297484in}}%
\pgfpathlineto{\pgfqpoint{2.183676in}{2.298678in}}%
\pgfpathlineto{\pgfqpoint{2.183681in}{2.298460in}}%
\pgfpathlineto{\pgfqpoint{2.184518in}{2.304041in}}%
\pgfpathlineto{\pgfqpoint{2.184713in}{2.302000in}}%
\pgfpathlineto{\pgfqpoint{2.184718in}{2.301957in}}%
\pgfpathlineto{\pgfqpoint{2.185356in}{2.305330in}}%
\pgfpathlineto{\pgfqpoint{2.185682in}{2.303205in}}%
\pgfpathlineto{\pgfqpoint{2.186018in}{2.303855in}}%
\pgfpathlineto{\pgfqpoint{2.186646in}{2.300749in}}%
\pgfpathlineto{\pgfqpoint{2.186767in}{2.300223in}}%
\pgfpathlineto{\pgfqpoint{2.187322in}{2.303358in}}%
\pgfpathlineto{\pgfqpoint{2.187649in}{2.303145in}}%
\pgfpathlineto{\pgfqpoint{2.188160in}{2.304310in}}%
\pgfpathlineto{\pgfqpoint{2.188330in}{2.301364in}}%
\pgfpathlineto{\pgfqpoint{2.188744in}{2.302725in}}%
\pgfpathlineto{\pgfqpoint{2.189689in}{2.303896in}}%
\pgfpathlineto{\pgfqpoint{2.189095in}{2.302387in}}%
\pgfpathlineto{\pgfqpoint{2.189820in}{2.302529in}}%
\pgfpathlineto{\pgfqpoint{2.190117in}{2.300948in}}%
\pgfpathlineto{\pgfqpoint{2.190813in}{2.304678in}}%
\pgfpathlineto{\pgfqpoint{2.191076in}{2.305545in}}%
\pgfpathlineto{\pgfqpoint{2.191738in}{2.300644in}}%
\pgfpathlineto{\pgfqpoint{2.191840in}{2.300923in}}%
\pgfpathlineto{\pgfqpoint{2.192892in}{2.299063in}}%
\pgfpathlineto{\pgfqpoint{2.192021in}{2.301282in}}%
\pgfpathlineto{\pgfqpoint{2.192975in}{2.299257in}}%
\pgfpathlineto{\pgfqpoint{2.193958in}{2.303555in}}%
\pgfpathlineto{\pgfqpoint{2.193077in}{2.298440in}}%
\pgfpathlineto{\pgfqpoint{2.194114in}{2.301863in}}%
\pgfpathlineto{\pgfqpoint{2.194343in}{2.300699in}}%
\pgfpathlineto{\pgfqpoint{2.194947in}{2.303898in}}%
\pgfpathlineto{\pgfqpoint{2.195171in}{2.302895in}}%
\pgfpathlineto{\pgfqpoint{2.195565in}{2.303823in}}%
\pgfpathlineto{\pgfqpoint{2.196261in}{2.300757in}}%
\pgfpathlineto{\pgfqpoint{2.197449in}{2.304075in}}%
\pgfpathlineto{\pgfqpoint{2.196417in}{2.299436in}}%
\pgfpathlineto{\pgfqpoint{2.197488in}{2.303577in}}%
\pgfpathlineto{\pgfqpoint{2.198355in}{2.302436in}}%
\pgfpathlineto{\pgfqpoint{2.198247in}{2.304710in}}%
\pgfpathlineto{\pgfqpoint{2.198617in}{2.302567in}}%
\pgfpathlineto{\pgfqpoint{2.199757in}{2.304353in}}%
\pgfpathlineto{\pgfqpoint{2.199314in}{2.300440in}}%
\pgfpathlineto{\pgfqpoint{2.199796in}{2.303514in}}%
\pgfpathlineto{\pgfqpoint{2.200930in}{2.302146in}}%
\pgfpathlineto{\pgfqpoint{2.200570in}{2.305456in}}%
\pgfpathlineto{\pgfqpoint{2.200940in}{2.302210in}}%
\pgfpathlineto{\pgfqpoint{2.201451in}{2.304142in}}%
\pgfpathlineto{\pgfqpoint{2.201159in}{2.301081in}}%
\pgfpathlineto{\pgfqpoint{2.202098in}{2.303526in}}%
\pgfpathlineto{\pgfqpoint{2.202507in}{2.300334in}}%
\pgfpathlineto{\pgfqpoint{2.202269in}{2.304549in}}%
\pgfpathlineto{\pgfqpoint{2.203257in}{2.302557in}}%
\pgfpathlineto{\pgfqpoint{2.204002in}{2.304751in}}%
\pgfpathlineto{\pgfqpoint{2.203281in}{2.302196in}}%
\pgfpathlineto{\pgfqpoint{2.204401in}{2.302894in}}%
\pgfpathlineto{\pgfqpoint{2.205433in}{2.299376in}}%
\pgfpathlineto{\pgfqpoint{2.204416in}{2.302935in}}%
\pgfpathlineto{\pgfqpoint{2.205555in}{2.300124in}}%
\pgfpathlineto{\pgfqpoint{2.206597in}{2.303777in}}%
\pgfpathlineto{\pgfqpoint{2.206081in}{2.298904in}}%
\pgfpathlineto{\pgfqpoint{2.206704in}{2.303613in}}%
\pgfpathlineto{\pgfqpoint{2.207731in}{2.302472in}}%
\pgfpathlineto{\pgfqpoint{2.206835in}{2.304822in}}%
\pgfpathlineto{\pgfqpoint{2.207809in}{2.303652in}}%
\pgfpathlineto{\pgfqpoint{2.208632in}{2.306584in}}%
\pgfpathlineto{\pgfqpoint{2.208218in}{2.302144in}}%
\pgfpathlineto{\pgfqpoint{2.208924in}{2.304337in}}%
\pgfpathlineto{\pgfqpoint{2.209688in}{2.305709in}}%
\pgfpathlineto{\pgfqpoint{2.209211in}{2.301430in}}%
\pgfpathlineto{\pgfqpoint{2.210029in}{2.304327in}}%
\pgfpathlineto{\pgfqpoint{2.210443in}{2.300212in}}%
\pgfpathlineto{\pgfqpoint{2.211154in}{2.303615in}}%
\pgfpathlineto{\pgfqpoint{2.212366in}{2.300424in}}%
\pgfpathlineto{\pgfqpoint{2.211329in}{2.303660in}}%
\pgfpathlineto{\pgfqpoint{2.212381in}{2.300466in}}%
\pgfpathlineto{\pgfqpoint{2.212868in}{2.296969in}}%
\pgfpathlineto{\pgfqpoint{2.213442in}{2.292974in}}%
\pgfpathlineto{\pgfqpoint{2.214026in}{2.294859in}}%
\pgfpathlineto{\pgfqpoint{2.214518in}{2.298410in}}%
\pgfpathlineto{\pgfqpoint{2.214070in}{2.294808in}}%
\pgfpathlineto{\pgfqpoint{2.215146in}{2.294978in}}%
\pgfpathlineto{\pgfqpoint{2.215190in}{2.294513in}}%
\pgfpathlineto{\pgfqpoint{2.215565in}{2.297786in}}%
\pgfpathlineto{\pgfqpoint{2.216198in}{2.296898in}}%
\pgfpathlineto{\pgfqpoint{2.217322in}{2.300788in}}%
\pgfpathlineto{\pgfqpoint{2.217371in}{2.300028in}}%
\pgfpathlineto{\pgfqpoint{2.218296in}{2.295853in}}%
\pgfpathlineto{\pgfqpoint{2.217434in}{2.300746in}}%
\pgfpathlineto{\pgfqpoint{2.218578in}{2.298051in}}%
\pgfpathlineto{\pgfqpoint{2.219625in}{2.302000in}}%
\pgfpathlineto{\pgfqpoint{2.218588in}{2.298007in}}%
\pgfpathlineto{\pgfqpoint{2.220229in}{2.300163in}}%
\pgfpathlineto{\pgfqpoint{2.221105in}{2.299713in}}%
\pgfpathlineto{\pgfqpoint{2.220828in}{2.301767in}}%
\pgfpathlineto{\pgfqpoint{2.221290in}{2.301705in}}%
\pgfpathlineto{\pgfqpoint{2.221850in}{2.304142in}}%
\pgfpathlineto{\pgfqpoint{2.221543in}{2.301040in}}%
\pgfpathlineto{\pgfqpoint{2.222444in}{2.303384in}}%
\pgfpathlineto{\pgfqpoint{2.223583in}{2.305146in}}%
\pgfpathlineto{\pgfqpoint{2.222731in}{2.302506in}}%
\pgfpathlineto{\pgfqpoint{2.223705in}{2.303593in}}%
\pgfpathlineto{\pgfqpoint{2.223817in}{2.302569in}}%
\pgfpathlineto{\pgfqpoint{2.224343in}{2.305086in}}%
\pgfpathlineto{\pgfqpoint{2.224834in}{2.303246in}}%
\pgfpathlineto{\pgfqpoint{2.225399in}{2.305446in}}%
\pgfpathlineto{\pgfqpoint{2.225740in}{2.301255in}}%
\pgfpathlineto{\pgfqpoint{2.225891in}{2.302552in}}%
\pgfpathlineto{\pgfqpoint{2.225896in}{2.302359in}}%
\pgfpathlineto{\pgfqpoint{2.226728in}{2.308331in}}%
\pgfpathlineto{\pgfqpoint{2.226879in}{2.307398in}}%
\pgfpathlineto{\pgfqpoint{2.226913in}{2.307977in}}%
\pgfpathlineto{\pgfqpoint{2.227872in}{2.304164in}}%
\pgfpathlineto{\pgfqpoint{2.227882in}{2.304439in}}%
\pgfpathlineto{\pgfqpoint{2.228096in}{2.302963in}}%
\pgfpathlineto{\pgfqpoint{2.228515in}{2.307629in}}%
\pgfpathlineto{\pgfqpoint{2.228987in}{2.304496in}}%
\pgfpathlineto{\pgfqpoint{2.229075in}{2.305558in}}%
\pgfpathlineto{\pgfqpoint{2.229484in}{2.301937in}}%
\pgfpathlineto{\pgfqpoint{2.229956in}{2.303493in}}%
\pgfpathlineto{\pgfqpoint{2.230939in}{2.301057in}}%
\pgfpathlineto{\pgfqpoint{2.230087in}{2.305062in}}%
\pgfpathlineto{\pgfqpoint{2.231100in}{2.302534in}}%
\pgfpathlineto{\pgfqpoint{2.231412in}{2.303694in}}%
\pgfpathlineto{\pgfqpoint{2.232083in}{2.301255in}}%
\pgfpathlineto{\pgfqpoint{2.232200in}{2.302087in}}%
\pgfpathlineto{\pgfqpoint{2.232395in}{2.300972in}}%
\pgfpathlineto{\pgfqpoint{2.232911in}{2.304703in}}%
\pgfpathlineto{\pgfqpoint{2.233286in}{2.302101in}}%
\pgfpathlineto{\pgfqpoint{2.234046in}{2.305489in}}%
\pgfpathlineto{\pgfqpoint{2.233612in}{2.301811in}}%
\pgfpathlineto{\pgfqpoint{2.234425in}{2.303869in}}%
\pgfpathlineto{\pgfqpoint{2.234435in}{2.304223in}}%
\pgfpathlineto{\pgfqpoint{2.235394in}{2.301353in}}%
\pgfpathlineto{\pgfqpoint{2.235409in}{2.301454in}}%
\pgfpathlineto{\pgfqpoint{2.236416in}{2.297795in}}%
\pgfpathlineto{\pgfqpoint{2.235482in}{2.302219in}}%
\pgfpathlineto{\pgfqpoint{2.236592in}{2.299018in}}%
\pgfpathlineto{\pgfqpoint{2.236597in}{2.298922in}}%
\pgfpathlineto{\pgfqpoint{2.237049in}{2.302909in}}%
\pgfpathlineto{\pgfqpoint{2.237614in}{2.300554in}}%
\pgfpathlineto{\pgfqpoint{2.238831in}{2.305005in}}%
\pgfpathlineto{\pgfqpoint{2.237872in}{2.299587in}}%
\pgfpathlineto{\pgfqpoint{2.238865in}{2.304867in}}%
\pgfpathlineto{\pgfqpoint{2.239863in}{2.300942in}}%
\pgfpathlineto{\pgfqpoint{2.239216in}{2.305670in}}%
\pgfpathlineto{\pgfqpoint{2.240107in}{2.302320in}}%
\pgfpathlineto{\pgfqpoint{2.240589in}{2.304981in}}%
\pgfpathlineto{\pgfqpoint{2.241037in}{2.302197in}}%
\pgfpathlineto{\pgfqpoint{2.241217in}{2.302434in}}%
\pgfpathlineto{\pgfqpoint{2.241436in}{2.304437in}}%
\pgfpathlineto{\pgfqpoint{2.242142in}{2.300788in}}%
\pgfpathlineto{\pgfqpoint{2.242229in}{2.300968in}}%
\pgfpathlineto{\pgfqpoint{2.242293in}{2.299921in}}%
\pgfpathlineto{\pgfqpoint{2.243057in}{2.306765in}}%
\pgfpathlineto{\pgfqpoint{2.243291in}{2.305303in}}%
\pgfpathlineto{\pgfqpoint{2.243364in}{2.306565in}}%
\pgfpathlineto{\pgfqpoint{2.243982in}{2.302736in}}%
\pgfpathlineto{\pgfqpoint{2.244328in}{2.304139in}}%
\pgfpathlineto{\pgfqpoint{2.245340in}{2.300294in}}%
\pgfpathlineto{\pgfqpoint{2.244381in}{2.304507in}}%
\pgfpathlineto{\pgfqpoint{2.245482in}{2.301588in}}%
\pgfpathlineto{\pgfqpoint{2.246519in}{2.303715in}}%
\pgfpathlineto{\pgfqpoint{2.245657in}{2.300521in}}%
\pgfpathlineto{\pgfqpoint{2.246621in}{2.303130in}}%
\pgfpathlineto{\pgfqpoint{2.247171in}{2.301798in}}%
\pgfpathlineto{\pgfqpoint{2.246718in}{2.304445in}}%
\pgfpathlineto{\pgfqpoint{2.247400in}{2.303799in}}%
\pgfpathlineto{\pgfqpoint{2.247911in}{2.304213in}}%
\pgfpathlineto{\pgfqpoint{2.248383in}{2.301140in}}%
\pgfpathlineto{\pgfqpoint{2.248442in}{2.301563in}}%
\pgfpathlineto{\pgfqpoint{2.248451in}{2.300983in}}%
\pgfpathlineto{\pgfqpoint{2.249026in}{2.304056in}}%
\pgfpathlineto{\pgfqpoint{2.249493in}{2.303115in}}%
\pgfpathlineto{\pgfqpoint{2.249931in}{2.305249in}}%
\pgfpathlineto{\pgfqpoint{2.250370in}{2.301523in}}%
\pgfpathlineto{\pgfqpoint{2.250603in}{2.303071in}}%
\pgfpathlineto{\pgfqpoint{2.251859in}{2.300342in}}%
\pgfpathlineto{\pgfqpoint{2.250667in}{2.303962in}}%
\pgfpathlineto{\pgfqpoint{2.251879in}{2.300474in}}%
\pgfpathlineto{\pgfqpoint{2.251889in}{2.300701in}}%
\pgfpathlineto{\pgfqpoint{2.252405in}{2.297674in}}%
\pgfpathlineto{\pgfqpoint{2.252935in}{2.299459in}}%
\pgfpathlineto{\pgfqpoint{2.253729in}{2.298539in}}%
\pgfpathlineto{\pgfqpoint{2.253992in}{2.301438in}}%
\pgfpathlineto{\pgfqpoint{2.254021in}{2.300826in}}%
\pgfpathlineto{\pgfqpoint{2.254410in}{2.302997in}}%
\pgfpathlineto{\pgfqpoint{2.254853in}{2.300002in}}%
\pgfpathlineto{\pgfqpoint{2.255141in}{2.301061in}}%
\pgfpathlineto{\pgfqpoint{2.255686in}{2.297930in}}%
\pgfpathlineto{\pgfqpoint{2.255389in}{2.301302in}}%
\pgfpathlineto{\pgfqpoint{2.256260in}{2.299883in}}%
\pgfpathlineto{\pgfqpoint{2.257239in}{2.303494in}}%
\pgfpathlineto{\pgfqpoint{2.256314in}{2.299411in}}%
\pgfpathlineto{\pgfqpoint{2.257429in}{2.301988in}}%
\pgfpathlineto{\pgfqpoint{2.257473in}{2.301458in}}%
\pgfpathlineto{\pgfqpoint{2.257828in}{2.304359in}}%
\pgfpathlineto{\pgfqpoint{2.257901in}{2.304563in}}%
\pgfpathlineto{\pgfqpoint{2.258812in}{2.301691in}}%
\pgfpathlineto{\pgfqpoint{2.258836in}{2.302217in}}%
\pgfpathlineto{\pgfqpoint{2.259274in}{2.301073in}}%
\pgfpathlineto{\pgfqpoint{2.259722in}{2.304630in}}%
\pgfpathlineto{\pgfqpoint{2.259892in}{2.303758in}}%
\pgfpathlineto{\pgfqpoint{2.261192in}{2.300760in}}%
\pgfpathlineto{\pgfqpoint{2.259902in}{2.303790in}}%
\pgfpathlineto{\pgfqpoint{2.261280in}{2.302057in}}%
\pgfpathlineto{\pgfqpoint{2.261343in}{2.302962in}}%
\pgfpathlineto{\pgfqpoint{2.262117in}{2.298523in}}%
\pgfpathlineto{\pgfqpoint{2.262254in}{2.300100in}}%
\pgfpathlineto{\pgfqpoint{2.262882in}{2.298636in}}%
\pgfpathlineto{\pgfqpoint{2.262405in}{2.301094in}}%
\pgfpathlineto{\pgfqpoint{2.263378in}{2.299534in}}%
\pgfpathlineto{\pgfqpoint{2.264503in}{2.304697in}}%
\pgfpathlineto{\pgfqpoint{2.263422in}{2.298932in}}%
\pgfpathlineto{\pgfqpoint{2.264819in}{2.302924in}}%
\pgfpathlineto{\pgfqpoint{2.265934in}{2.299032in}}%
\pgfpathlineto{\pgfqpoint{2.265019in}{2.303789in}}%
\pgfpathlineto{\pgfqpoint{2.265978in}{2.299337in}}%
\pgfpathlineto{\pgfqpoint{2.266961in}{2.301173in}}%
\pgfpathlineto{\pgfqpoint{2.266504in}{2.298450in}}%
\pgfpathlineto{\pgfqpoint{2.267112in}{2.300272in}}%
\pgfpathlineto{\pgfqpoint{2.268252in}{2.298524in}}%
\pgfpathlineto{\pgfqpoint{2.267322in}{2.300694in}}%
\pgfpathlineto{\pgfqpoint{2.268310in}{2.299099in}}%
\pgfpathlineto{\pgfqpoint{2.268451in}{2.300472in}}%
\pgfpathlineto{\pgfqpoint{2.269157in}{2.294858in}}%
\pgfpathlineto{\pgfqpoint{2.269352in}{2.297369in}}%
\pgfpathlineto{\pgfqpoint{2.269639in}{2.294872in}}%
\pgfpathlineto{\pgfqpoint{2.270321in}{2.298280in}}%
\pgfpathlineto{\pgfqpoint{2.270472in}{2.296454in}}%
\pgfpathlineto{\pgfqpoint{2.271309in}{2.297543in}}%
\pgfpathlineto{\pgfqpoint{2.271012in}{2.294307in}}%
\pgfpathlineto{\pgfqpoint{2.271587in}{2.296451in}}%
\pgfpathlineto{\pgfqpoint{2.273758in}{2.302318in}}%
\pgfpathlineto{\pgfqpoint{2.271616in}{2.296272in}}%
\pgfpathlineto{\pgfqpoint{2.274216in}{2.301398in}}%
\pgfpathlineto{\pgfqpoint{2.274454in}{2.299528in}}%
\pgfpathlineto{\pgfqpoint{2.275311in}{2.302181in}}%
\pgfpathlineto{\pgfqpoint{2.275924in}{2.303366in}}%
\pgfpathlineto{\pgfqpoint{2.276392in}{2.300092in}}%
\pgfpathlineto{\pgfqpoint{2.277317in}{2.301940in}}%
\pgfpathlineto{\pgfqpoint{2.276557in}{2.298962in}}%
\pgfpathlineto{\pgfqpoint{2.277555in}{2.301161in}}%
\pgfpathlineto{\pgfqpoint{2.277906in}{2.297978in}}%
\pgfpathlineto{\pgfqpoint{2.278636in}{2.302156in}}%
\pgfpathlineto{\pgfqpoint{2.279030in}{2.301187in}}%
\pgfpathlineto{\pgfqpoint{2.279196in}{2.303702in}}%
\pgfpathlineto{\pgfqpoint{2.279386in}{2.303229in}}%
\pgfpathlineto{\pgfqpoint{2.279561in}{2.304630in}}%
\pgfpathlineto{\pgfqpoint{2.279868in}{2.301782in}}%
\pgfpathlineto{\pgfqpoint{2.279999in}{2.301972in}}%
\pgfpathlineto{\pgfqpoint{2.280360in}{2.299923in}}%
\pgfpathlineto{\pgfqpoint{2.280822in}{2.302679in}}%
\pgfpathlineto{\pgfqpoint{2.281114in}{2.301636in}}%
\pgfpathlineto{\pgfqpoint{2.282049in}{2.304777in}}%
\pgfpathlineto{\pgfqpoint{2.281309in}{2.300213in}}%
\pgfpathlineto{\pgfqpoint{2.282263in}{2.304264in}}%
\pgfpathlineto{\pgfqpoint{2.283359in}{2.301144in}}%
\pgfpathlineto{\pgfqpoint{2.282862in}{2.305087in}}%
\pgfpathlineto{\pgfqpoint{2.283412in}{2.301234in}}%
\pgfpathlineto{\pgfqpoint{2.284527in}{2.299962in}}%
\pgfpathlineto{\pgfqpoint{2.284006in}{2.302046in}}%
\pgfpathlineto{\pgfqpoint{2.284571in}{2.300668in}}%
\pgfpathlineto{\pgfqpoint{2.285199in}{2.303922in}}%
\pgfpathlineto{\pgfqpoint{2.284663in}{2.300198in}}%
\pgfpathlineto{\pgfqpoint{2.285730in}{2.302770in}}%
\pgfpathlineto{\pgfqpoint{2.286567in}{2.302097in}}%
\pgfpathlineto{\pgfqpoint{2.286796in}{2.304281in}}%
\pgfpathlineto{\pgfqpoint{2.287127in}{2.304453in}}%
\pgfpathlineto{\pgfqpoint{2.286903in}{2.302262in}}%
\pgfpathlineto{\pgfqpoint{2.287872in}{2.303522in}}%
\pgfpathlineto{\pgfqpoint{2.288505in}{2.302823in}}%
\pgfpathlineto{\pgfqpoint{2.288845in}{2.305812in}}%
\pgfpathlineto{\pgfqpoint{2.288957in}{2.305118in}}%
\pgfpathlineto{\pgfqpoint{2.289960in}{2.300952in}}%
\pgfpathlineto{\pgfqpoint{2.289040in}{2.305381in}}%
\pgfpathlineto{\pgfqpoint{2.290160in}{2.302116in}}%
\pgfpathlineto{\pgfqpoint{2.290174in}{2.301750in}}%
\pgfpathlineto{\pgfqpoint{2.290734in}{2.304746in}}%
\pgfpathlineto{\pgfqpoint{2.291221in}{2.302745in}}%
\pgfpathlineto{\pgfqpoint{2.291966in}{2.305201in}}%
\pgfpathlineto{\pgfqpoint{2.292132in}{2.302578in}}%
\pgfpathlineto{\pgfqpoint{2.292356in}{2.303915in}}%
\pgfpathlineto{\pgfqpoint{2.293582in}{2.297736in}}%
\pgfpathlineto{\pgfqpoint{2.292370in}{2.304079in}}%
\pgfpathlineto{\pgfqpoint{2.294069in}{2.298415in}}%
\pgfpathlineto{\pgfqpoint{2.294980in}{2.303125in}}%
\pgfpathlineto{\pgfqpoint{2.295228in}{2.301503in}}%
\pgfpathlineto{\pgfqpoint{2.296430in}{2.297511in}}%
\pgfpathlineto{\pgfqpoint{2.296484in}{2.296894in}}%
\pgfpathlineto{\pgfqpoint{2.296669in}{2.298905in}}%
\pgfpathlineto{\pgfqpoint{2.296805in}{2.300672in}}%
\pgfpathlineto{\pgfqpoint{2.297458in}{2.296963in}}%
\pgfpathlineto{\pgfqpoint{2.297808in}{2.300551in}}%
\pgfpathlineto{\pgfqpoint{2.297872in}{2.299992in}}%
\pgfpathlineto{\pgfqpoint{2.298251in}{2.302822in}}%
\pgfpathlineto{\pgfqpoint{2.298899in}{2.301725in}}%
\pgfpathlineto{\pgfqpoint{2.299814in}{2.300607in}}%
\pgfpathlineto{\pgfqpoint{2.299624in}{2.303462in}}%
\pgfpathlineto{\pgfqpoint{2.300019in}{2.301330in}}%
\pgfpathlineto{\pgfqpoint{2.300593in}{2.304278in}}%
\pgfpathlineto{\pgfqpoint{2.300062in}{2.301082in}}%
\pgfpathlineto{\pgfqpoint{2.301090in}{2.301237in}}%
\pgfpathlineto{\pgfqpoint{2.301927in}{2.300208in}}%
\pgfpathlineto{\pgfqpoint{2.302136in}{2.303267in}}%
\pgfpathlineto{\pgfqpoint{2.302190in}{2.301766in}}%
\pgfpathlineto{\pgfqpoint{2.302618in}{2.305020in}}%
\pgfpathlineto{\pgfqpoint{2.303266in}{2.300555in}}%
\pgfpathlineto{\pgfqpoint{2.303485in}{2.298643in}}%
\pgfpathlineto{\pgfqpoint{2.304342in}{2.302214in}}%
\pgfpathlineto{\pgfqpoint{2.304352in}{2.301975in}}%
\pgfpathlineto{\pgfqpoint{2.305608in}{2.305917in}}%
\pgfpathlineto{\pgfqpoint{2.304420in}{2.301525in}}%
\pgfpathlineto{\pgfqpoint{2.305875in}{2.302584in}}%
\pgfpathlineto{\pgfqpoint{2.307068in}{2.300624in}}%
\pgfpathlineto{\pgfqpoint{2.306226in}{2.305317in}}%
\pgfpathlineto{\pgfqpoint{2.307078in}{2.300913in}}%
\pgfpathlineto{\pgfqpoint{2.308159in}{2.304004in}}%
\pgfpathlineto{\pgfqpoint{2.308251in}{2.303717in}}%
\pgfpathlineto{\pgfqpoint{2.308631in}{2.301077in}}%
\pgfpathlineto{\pgfqpoint{2.309347in}{2.303891in}}%
\pgfpathlineto{\pgfqpoint{2.309361in}{2.303678in}}%
\pgfpathlineto{\pgfqpoint{2.310135in}{2.306111in}}%
\pgfpathlineto{\pgfqpoint{2.309600in}{2.300973in}}%
\pgfpathlineto{\pgfqpoint{2.310471in}{2.303902in}}%
\pgfpathlineto{\pgfqpoint{2.311484in}{2.300896in}}%
\pgfpathlineto{\pgfqpoint{2.310617in}{2.304200in}}%
\pgfpathlineto{\pgfqpoint{2.311615in}{2.301861in}}%
\pgfpathlineto{\pgfqpoint{2.311800in}{2.303642in}}%
\pgfpathlineto{\pgfqpoint{2.312613in}{2.299635in}}%
\pgfpathlineto{\pgfqpoint{2.312691in}{2.300318in}}%
\pgfpathlineto{\pgfqpoint{2.313845in}{2.296479in}}%
\pgfpathlineto{\pgfqpoint{2.312906in}{2.301157in}}%
\pgfpathlineto{\pgfqpoint{2.313850in}{2.296567in}}%
\pgfpathlineto{\pgfqpoint{2.313889in}{2.297597in}}%
\pgfpathlineto{\pgfqpoint{2.314556in}{2.292705in}}%
\pgfpathlineto{\pgfqpoint{2.314902in}{2.294247in}}%
\pgfpathlineto{\pgfqpoint{2.315286in}{2.297068in}}%
\pgfpathlineto{\pgfqpoint{2.314911in}{2.294188in}}%
\pgfpathlineto{\pgfqpoint{2.316060in}{2.295782in}}%
\pgfpathlineto{\pgfqpoint{2.316070in}{2.295687in}}%
\pgfpathlineto{\pgfqpoint{2.317365in}{2.300535in}}%
\pgfpathlineto{\pgfqpoint{2.317682in}{2.299878in}}%
\pgfpathlineto{\pgfqpoint{2.318319in}{2.302659in}}%
\pgfpathlineto{\pgfqpoint{2.318397in}{2.301993in}}%
\pgfpathlineto{\pgfqpoint{2.319624in}{2.304569in}}%
\pgfpathlineto{\pgfqpoint{2.318991in}{2.301483in}}%
\pgfpathlineto{\pgfqpoint{2.319629in}{2.304235in}}%
\pgfpathlineto{\pgfqpoint{2.319897in}{2.302005in}}%
\pgfpathlineto{\pgfqpoint{2.320617in}{2.304574in}}%
\pgfpathlineto{\pgfqpoint{2.320866in}{2.303096in}}%
\pgfpathlineto{\pgfqpoint{2.320900in}{2.304214in}}%
\pgfpathlineto{\pgfqpoint{2.321591in}{2.300881in}}%
\pgfpathlineto{\pgfqpoint{2.321971in}{2.302980in}}%
\pgfpathlineto{\pgfqpoint{2.322681in}{2.300924in}}%
\pgfpathlineto{\pgfqpoint{2.323047in}{2.303927in}}%
\pgfpathlineto{\pgfqpoint{2.323076in}{2.303005in}}%
\pgfpathlineto{\pgfqpoint{2.324015in}{2.304422in}}%
\pgfpathlineto{\pgfqpoint{2.323519in}{2.300982in}}%
\pgfpathlineto{\pgfqpoint{2.324205in}{2.303582in}}%
\pgfpathlineto{\pgfqpoint{2.324600in}{2.301616in}}%
\pgfpathlineto{\pgfqpoint{2.324755in}{2.303916in}}%
\pgfpathlineto{\pgfqpoint{2.325335in}{2.302147in}}%
\pgfpathlineto{\pgfqpoint{2.326187in}{2.300522in}}%
\pgfpathlineto{\pgfqpoint{2.326391in}{2.302828in}}%
\pgfpathlineto{\pgfqpoint{2.327414in}{2.301615in}}%
\pgfpathlineto{\pgfqpoint{2.327574in}{2.304606in}}%
\pgfpathlineto{\pgfqpoint{2.328426in}{2.302382in}}%
\pgfpathlineto{\pgfqpoint{2.328709in}{2.303413in}}%
\pgfpathlineto{\pgfqpoint{2.328937in}{2.304020in}}%
\pgfpathlineto{\pgfqpoint{2.329454in}{2.300863in}}%
\pgfpathlineto{\pgfqpoint{2.329799in}{2.303043in}}%
\pgfpathlineto{\pgfqpoint{2.329989in}{2.302294in}}%
\pgfpathlineto{\pgfqpoint{2.330077in}{2.304866in}}%
\pgfpathlineto{\pgfqpoint{2.330237in}{2.305944in}}%
\pgfpathlineto{\pgfqpoint{2.331046in}{2.299847in}}%
\pgfpathlineto{\pgfqpoint{2.331187in}{2.300831in}}%
\pgfpathlineto{\pgfqpoint{2.331411in}{2.298678in}}%
\pgfpathlineto{\pgfqpoint{2.331425in}{2.298673in}}%
\pgfpathlineto{\pgfqpoint{2.331742in}{2.296614in}}%
\pgfpathlineto{\pgfqpoint{2.332341in}{2.301561in}}%
\pgfpathlineto{\pgfqpoint{2.332380in}{2.300850in}}%
\pgfpathlineto{\pgfqpoint{2.333402in}{2.303625in}}%
\pgfpathlineto{\pgfqpoint{2.332414in}{2.300679in}}%
\pgfpathlineto{\pgfqpoint{2.333524in}{2.302108in}}%
\pgfpathlineto{\pgfqpoint{2.333621in}{2.301418in}}%
\pgfpathlineto{\pgfqpoint{2.333723in}{2.303243in}}%
\pgfpathlineto{\pgfqpoint{2.333967in}{2.303226in}}%
\pgfpathlineto{\pgfqpoint{2.334994in}{2.308253in}}%
\pgfpathlineto{\pgfqpoint{2.334171in}{2.303032in}}%
\pgfpathlineto{\pgfqpoint{2.335159in}{2.307059in}}%
\pgfpathlineto{\pgfqpoint{2.336313in}{2.304159in}}%
\pgfpathlineto{\pgfqpoint{2.335787in}{2.310018in}}%
\pgfpathlineto{\pgfqpoint{2.336328in}{2.304787in}}%
\pgfpathlineto{\pgfqpoint{2.336377in}{2.305273in}}%
\pgfpathlineto{\pgfqpoint{2.336693in}{2.302656in}}%
\pgfpathlineto{\pgfqpoint{2.337954in}{2.300577in}}%
\pgfpathlineto{\pgfqpoint{2.336936in}{2.303748in}}%
\pgfpathlineto{\pgfqpoint{2.337973in}{2.301255in}}%
\pgfpathlineto{\pgfqpoint{2.338163in}{2.302552in}}%
\pgfpathlineto{\pgfqpoint{2.338538in}{2.299239in}}%
\pgfpathlineto{\pgfqpoint{2.339079in}{2.301142in}}%
\pgfpathlineto{\pgfqpoint{2.339736in}{2.303290in}}%
\pgfpathlineto{\pgfqpoint{2.339298in}{2.299931in}}%
\pgfpathlineto{\pgfqpoint{2.340252in}{2.301007in}}%
\pgfpathlineto{\pgfqpoint{2.341284in}{2.298779in}}%
\pgfpathlineto{\pgfqpoint{2.340495in}{2.301511in}}%
\pgfpathlineto{\pgfqpoint{2.341396in}{2.299010in}}%
\pgfpathlineto{\pgfqpoint{2.341854in}{2.301757in}}%
\pgfpathlineto{\pgfqpoint{2.341566in}{2.298668in}}%
\pgfpathlineto{\pgfqpoint{2.342501in}{2.298963in}}%
\pgfpathlineto{\pgfqpoint{2.343738in}{2.297311in}}%
\pgfpathlineto{\pgfqpoint{2.342891in}{2.301031in}}%
\pgfpathlineto{\pgfqpoint{2.343747in}{2.297564in}}%
\pgfpathlineto{\pgfqpoint{2.344760in}{2.304066in}}%
\pgfpathlineto{\pgfqpoint{2.343913in}{2.297373in}}%
\pgfpathlineto{\pgfqpoint{2.345247in}{2.303545in}}%
\pgfpathlineto{\pgfqpoint{2.345505in}{2.302630in}}%
\pgfpathlineto{\pgfqpoint{2.346240in}{2.305358in}}%
\pgfpathlineto{\pgfqpoint{2.346342in}{2.304127in}}%
\pgfpathlineto{\pgfqpoint{2.346673in}{2.302242in}}%
\pgfpathlineto{\pgfqpoint{2.346897in}{2.304510in}}%
\pgfpathlineto{\pgfqpoint{2.347472in}{2.303669in}}%
\pgfpathlineto{\pgfqpoint{2.347676in}{2.304301in}}%
\pgfpathlineto{\pgfqpoint{2.348519in}{2.301681in}}%
\pgfpathlineto{\pgfqpoint{2.348533in}{2.301826in}}%
\pgfpathlineto{\pgfqpoint{2.348558in}{2.301512in}}%
\pgfpathlineto{\pgfqpoint{2.348908in}{2.303874in}}%
\pgfpathlineto{\pgfqpoint{2.349609in}{2.302674in}}%
\pgfpathlineto{\pgfqpoint{2.350768in}{2.301254in}}%
\pgfpathlineto{\pgfqpoint{2.349979in}{2.304432in}}%
\pgfpathlineto{\pgfqpoint{2.350948in}{2.301504in}}%
\pgfpathlineto{\pgfqpoint{2.352048in}{2.305428in}}%
\pgfpathlineto{\pgfqpoint{2.351425in}{2.300600in}}%
\pgfpathlineto{\pgfqpoint{2.352112in}{2.305226in}}%
\pgfpathlineto{\pgfqpoint{2.353197in}{2.302716in}}%
\pgfpathlineto{\pgfqpoint{2.352861in}{2.305500in}}%
\pgfpathlineto{\pgfqpoint{2.353363in}{2.303717in}}%
\pgfpathlineto{\pgfqpoint{2.353859in}{2.304269in}}%
\pgfpathlineto{\pgfqpoint{2.353704in}{2.301897in}}%
\pgfpathlineto{\pgfqpoint{2.354478in}{2.304012in}}%
\pgfpathlineto{\pgfqpoint{2.354512in}{2.304434in}}%
\pgfpathlineto{\pgfqpoint{2.355437in}{2.300158in}}%
\pgfpathlineto{\pgfqpoint{2.355466in}{2.300564in}}%
\pgfpathlineto{\pgfqpoint{2.355485in}{2.299996in}}%
\pgfpathlineto{\pgfqpoint{2.356289in}{2.303761in}}%
\pgfpathlineto{\pgfqpoint{2.356547in}{2.302555in}}%
\pgfpathlineto{\pgfqpoint{2.356654in}{2.303728in}}%
\pgfpathlineto{\pgfqpoint{2.357306in}{2.300695in}}%
\pgfpathlineto{\pgfqpoint{2.357637in}{2.301624in}}%
\pgfpathlineto{\pgfqpoint{2.358616in}{2.301503in}}%
\pgfpathlineto{\pgfqpoint{2.358840in}{2.304073in}}%
\pgfpathlineto{\pgfqpoint{2.359862in}{2.292922in}}%
\pgfpathlineto{\pgfqpoint{2.358859in}{2.304152in}}%
\pgfpathlineto{\pgfqpoint{2.360456in}{2.294785in}}%
\pgfpathlineto{\pgfqpoint{2.360471in}{2.295013in}}%
\pgfpathlineto{\pgfqpoint{2.361079in}{2.291408in}}%
\pgfpathlineto{\pgfqpoint{2.361274in}{2.292554in}}%
\pgfpathlineto{\pgfqpoint{2.361926in}{2.289147in}}%
\pgfpathlineto{\pgfqpoint{2.361381in}{2.292628in}}%
\pgfpathlineto{\pgfqpoint{2.362413in}{2.291342in}}%
\pgfpathlineto{\pgfqpoint{2.362467in}{2.292029in}}%
\pgfpathlineto{\pgfqpoint{2.363226in}{2.285687in}}%
\pgfpathlineto{\pgfqpoint{2.363319in}{2.286565in}}%
\pgfpathlineto{\pgfqpoint{2.363353in}{2.286447in}}%
\pgfpathlineto{\pgfqpoint{2.364049in}{2.290512in}}%
\pgfpathlineto{\pgfqpoint{2.364083in}{2.290330in}}%
\pgfpathlineto{\pgfqpoint{2.365237in}{2.293581in}}%
\pgfpathlineto{\pgfqpoint{2.364288in}{2.288124in}}%
\pgfpathlineto{\pgfqpoint{2.365276in}{2.292904in}}%
\pgfpathlineto{\pgfqpoint{2.366401in}{2.288279in}}%
\pgfpathlineto{\pgfqpoint{2.365378in}{2.293640in}}%
\pgfpathlineto{\pgfqpoint{2.366444in}{2.289800in}}%
\pgfpathlineto{\pgfqpoint{2.367442in}{2.293636in}}%
\pgfpathlineto{\pgfqpoint{2.366605in}{2.288633in}}%
\pgfpathlineto{\pgfqpoint{2.367666in}{2.293156in}}%
\pgfpathlineto{\pgfqpoint{2.367837in}{2.292257in}}%
\pgfpathlineto{\pgfqpoint{2.368650in}{2.295120in}}%
\pgfpathlineto{\pgfqpoint{2.368772in}{2.293388in}}%
\pgfpathlineto{\pgfqpoint{2.369984in}{2.299490in}}%
\pgfpathlineto{\pgfqpoint{2.369127in}{2.292460in}}%
\pgfpathlineto{\pgfqpoint{2.370086in}{2.298967in}}%
\pgfpathlineto{\pgfqpoint{2.370695in}{2.299386in}}%
\pgfpathlineto{\pgfqpoint{2.371040in}{2.297059in}}%
\pgfpathlineto{\pgfqpoint{2.371795in}{2.294703in}}%
\pgfpathlineto{\pgfqpoint{2.371264in}{2.297370in}}%
\pgfpathlineto{\pgfqpoint{2.372150in}{2.297121in}}%
\pgfpathlineto{\pgfqpoint{2.373523in}{2.301856in}}%
\pgfpathlineto{\pgfqpoint{2.372379in}{2.295313in}}%
\pgfpathlineto{\pgfqpoint{2.373698in}{2.299421in}}%
\pgfpathlineto{\pgfqpoint{2.373733in}{2.298669in}}%
\pgfpathlineto{\pgfqpoint{2.374229in}{2.300976in}}%
\pgfpathlineto{\pgfqpoint{2.374789in}{2.300874in}}%
\pgfpathlineto{\pgfqpoint{2.375529in}{2.295845in}}%
\pgfpathlineto{\pgfqpoint{2.375987in}{2.296569in}}%
\pgfpathlineto{\pgfqpoint{2.376108in}{2.297438in}}%
\pgfpathlineto{\pgfqpoint{2.377033in}{2.292176in}}%
\pgfpathlineto{\pgfqpoint{2.377038in}{2.292033in}}%
\pgfpathlineto{\pgfqpoint{2.377316in}{2.295907in}}%
\pgfpathlineto{\pgfqpoint{2.378066in}{2.294322in}}%
\pgfpathlineto{\pgfqpoint{2.379176in}{2.298979in}}%
\pgfpathlineto{\pgfqpoint{2.378484in}{2.293554in}}%
\pgfpathlineto{\pgfqpoint{2.379434in}{2.298049in}}%
\pgfpathlineto{\pgfqpoint{2.379463in}{2.297613in}}%
\pgfpathlineto{\pgfqpoint{2.380402in}{2.300997in}}%
\pgfpathlineto{\pgfqpoint{2.380519in}{2.299336in}}%
\pgfpathlineto{\pgfqpoint{2.381279in}{2.302366in}}%
\pgfpathlineto{\pgfqpoint{2.381658in}{2.301525in}}%
\pgfpathlineto{\pgfqpoint{2.382413in}{2.301294in}}%
\pgfpathlineto{\pgfqpoint{2.382136in}{2.303256in}}%
\pgfpathlineto{\pgfqpoint{2.382759in}{2.301761in}}%
\pgfpathlineto{\pgfqpoint{2.382910in}{2.303786in}}%
\pgfpathlineto{\pgfqpoint{2.383830in}{2.300763in}}%
\pgfpathlineto{\pgfqpoint{2.383835in}{2.300894in}}%
\pgfpathlineto{\pgfqpoint{2.383942in}{2.300256in}}%
\pgfpathlineto{\pgfqpoint{2.384570in}{2.304048in}}%
\pgfpathlineto{\pgfqpoint{2.384906in}{2.302984in}}%
\pgfpathlineto{\pgfqpoint{2.385373in}{2.306063in}}%
\pgfpathlineto{\pgfqpoint{2.385987in}{2.302187in}}%
\pgfpathlineto{\pgfqpoint{2.386001in}{2.302678in}}%
\pgfpathlineto{\pgfqpoint{2.386517in}{2.301471in}}%
\pgfpathlineto{\pgfqpoint{2.386902in}{2.304624in}}%
\pgfpathlineto{\pgfqpoint{2.387053in}{2.303625in}}%
\pgfpathlineto{\pgfqpoint{2.387924in}{2.304723in}}%
\pgfpathlineto{\pgfqpoint{2.388065in}{2.301603in}}%
\pgfpathlineto{\pgfqpoint{2.388095in}{2.302198in}}%
\pgfpathlineto{\pgfqpoint{2.389190in}{2.300617in}}%
\pgfpathlineto{\pgfqpoint{2.388577in}{2.303360in}}%
\pgfpathlineto{\pgfqpoint{2.389248in}{2.300827in}}%
\pgfpathlineto{\pgfqpoint{2.389862in}{2.306516in}}%
\pgfpathlineto{\pgfqpoint{2.389302in}{2.300357in}}%
\pgfpathlineto{\pgfqpoint{2.390436in}{2.302726in}}%
\pgfpathlineto{\pgfqpoint{2.390568in}{2.302268in}}%
\pgfpathlineto{\pgfqpoint{2.390836in}{2.305487in}}%
\pgfpathlineto{\pgfqpoint{2.391512in}{2.303588in}}%
\pgfpathlineto{\pgfqpoint{2.392642in}{2.304704in}}%
\pgfpathlineto{\pgfqpoint{2.392350in}{2.302140in}}%
\pgfpathlineto{\pgfqpoint{2.392647in}{2.304571in}}%
\pgfpathlineto{\pgfqpoint{2.393927in}{2.301398in}}%
\pgfpathlineto{\pgfqpoint{2.392861in}{2.305642in}}%
\pgfpathlineto{\pgfqpoint{2.394088in}{2.303165in}}%
\pgfpathlineto{\pgfqpoint{2.394906in}{2.304083in}}%
\pgfpathlineto{\pgfqpoint{2.394307in}{2.302065in}}%
\pgfpathlineto{\pgfqpoint{2.395193in}{2.302830in}}%
\pgfpathlineto{\pgfqpoint{2.395485in}{2.303834in}}%
\pgfpathlineto{\pgfqpoint{2.395743in}{2.305900in}}%
\pgfpathlineto{\pgfqpoint{2.396298in}{2.301579in}}%
\pgfpathlineto{\pgfqpoint{2.396556in}{2.302337in}}%
\pgfpathlineto{\pgfqpoint{2.396785in}{2.301371in}}%
\pgfpathlineto{\pgfqpoint{2.397525in}{2.304102in}}%
\pgfpathlineto{\pgfqpoint{2.397627in}{2.303439in}}%
\pgfpathlineto{\pgfqpoint{2.398362in}{2.302156in}}%
\pgfpathlineto{\pgfqpoint{2.398820in}{2.305058in}}%
\pgfpathlineto{\pgfqpoint{2.398990in}{2.305680in}}%
\pgfpathlineto{\pgfqpoint{2.399565in}{2.303060in}}%
\pgfpathlineto{\pgfqpoint{2.400704in}{2.296893in}}%
\pgfpathlineto{\pgfqpoint{2.399584in}{2.303112in}}%
\pgfpathlineto{\pgfqpoint{2.400787in}{2.297595in}}%
\pgfpathlineto{\pgfqpoint{2.401050in}{2.299764in}}%
\pgfpathlineto{\pgfqpoint{2.401420in}{2.297223in}}%
\pgfpathlineto{\pgfqpoint{2.401946in}{2.298674in}}%
\pgfpathlineto{\pgfqpoint{2.401950in}{2.298535in}}%
\pgfpathlineto{\pgfqpoint{2.402763in}{2.302449in}}%
\pgfpathlineto{\pgfqpoint{2.402948in}{2.301509in}}%
\pgfpathlineto{\pgfqpoint{2.402963in}{2.302190in}}%
\pgfpathlineto{\pgfqpoint{2.403547in}{2.299207in}}%
\pgfpathlineto{\pgfqpoint{2.404039in}{2.300364in}}%
\pgfpathlineto{\pgfqpoint{2.404054in}{2.300044in}}%
\pgfpathlineto{\pgfqpoint{2.404234in}{2.301368in}}%
\pgfpathlineto{\pgfqpoint{2.404511in}{2.303740in}}%
\pgfpathlineto{\pgfqpoint{2.405042in}{2.298884in}}%
\pgfpathlineto{\pgfqpoint{2.405334in}{2.300642in}}%
\pgfpathlineto{\pgfqpoint{2.406449in}{2.301421in}}%
\pgfpathlineto{\pgfqpoint{2.406055in}{2.297977in}}%
\pgfpathlineto{\pgfqpoint{2.406478in}{2.300998in}}%
\pgfpathlineto{\pgfqpoint{2.407291in}{2.299296in}}%
\pgfpathlineto{\pgfqpoint{2.406980in}{2.301933in}}%
\pgfpathlineto{\pgfqpoint{2.407622in}{2.299600in}}%
\pgfpathlineto{\pgfqpoint{2.407861in}{2.302228in}}%
\pgfpathlineto{\pgfqpoint{2.408494in}{2.298704in}}%
\pgfpathlineto{\pgfqpoint{2.408674in}{2.299089in}}%
\pgfpathlineto{\pgfqpoint{2.408688in}{2.298825in}}%
\pgfpathlineto{\pgfqpoint{2.409492in}{2.303738in}}%
\pgfpathlineto{\pgfqpoint{2.409648in}{2.303176in}}%
\pgfpathlineto{\pgfqpoint{2.410091in}{2.304155in}}%
\pgfpathlineto{\pgfqpoint{2.410266in}{2.302159in}}%
\pgfpathlineto{\pgfqpoint{2.410290in}{2.302154in}}%
\pgfpathlineto{\pgfqpoint{2.410855in}{2.298826in}}%
\pgfpathlineto{\pgfqpoint{2.411425in}{2.300418in}}%
\pgfpathlineto{\pgfqpoint{2.412165in}{2.304583in}}%
\pgfpathlineto{\pgfqpoint{2.411463in}{2.300276in}}%
\pgfpathlineto{\pgfqpoint{2.412583in}{2.302285in}}%
\pgfpathlineto{\pgfqpoint{2.413654in}{2.300704in}}%
\pgfpathlineto{\pgfqpoint{2.412895in}{2.305126in}}%
\pgfpathlineto{\pgfqpoint{2.413713in}{2.301384in}}%
\pgfpathlineto{\pgfqpoint{2.413844in}{2.302688in}}%
\pgfpathlineto{\pgfqpoint{2.414720in}{2.297507in}}%
\pgfpathlineto{\pgfqpoint{2.414794in}{2.299400in}}%
\pgfpathlineto{\pgfqpoint{2.414993in}{2.298787in}}%
\pgfpathlineto{\pgfqpoint{2.415665in}{2.304003in}}%
\pgfpathlineto{\pgfqpoint{2.415845in}{2.306206in}}%
\pgfpathlineto{\pgfqpoint{2.416507in}{2.301787in}}%
\pgfpathlineto{\pgfqpoint{2.416756in}{2.302780in}}%
\pgfpathlineto{\pgfqpoint{2.417126in}{2.299816in}}%
\pgfpathlineto{\pgfqpoint{2.416790in}{2.302957in}}%
\pgfpathlineto{\pgfqpoint{2.417904in}{2.302139in}}%
\pgfpathlineto{\pgfqpoint{2.418844in}{2.304357in}}%
\pgfpathlineto{\pgfqpoint{2.418279in}{2.300726in}}%
\pgfpathlineto{\pgfqpoint{2.419044in}{2.303459in}}%
\pgfpathlineto{\pgfqpoint{2.420368in}{2.306498in}}%
\pgfpathlineto{\pgfqpoint{2.419214in}{2.301996in}}%
\pgfpathlineto{\pgfqpoint{2.420631in}{2.305195in}}%
\pgfpathlineto{\pgfqpoint{2.421746in}{2.302625in}}%
\pgfpathlineto{\pgfqpoint{2.420952in}{2.305612in}}%
\pgfpathlineto{\pgfqpoint{2.421770in}{2.303416in}}%
\pgfpathlineto{\pgfqpoint{2.422169in}{2.304310in}}%
\pgfpathlineto{\pgfqpoint{2.422715in}{2.301366in}}%
\pgfpathlineto{\pgfqpoint{2.422841in}{2.302325in}}%
\pgfpathlineto{\pgfqpoint{2.423274in}{2.300332in}}%
\pgfpathlineto{\pgfqpoint{2.422968in}{2.303570in}}%
\pgfpathlineto{\pgfqpoint{2.423810in}{2.302443in}}%
\pgfpathlineto{\pgfqpoint{2.423980in}{2.303521in}}%
\pgfpathlineto{\pgfqpoint{2.424297in}{2.301352in}}%
\pgfpathlineto{\pgfqpoint{2.424910in}{2.301920in}}%
\pgfpathlineto{\pgfqpoint{2.425392in}{2.299290in}}%
\pgfpathlineto{\pgfqpoint{2.425976in}{2.304085in}}%
\pgfpathlineto{\pgfqpoint{2.426390in}{2.300339in}}%
\pgfpathlineto{\pgfqpoint{2.427101in}{2.303527in}}%
\pgfpathlineto{\pgfqpoint{2.427919in}{2.304566in}}%
\pgfpathlineto{\pgfqpoint{2.428192in}{2.302837in}}%
\pgfpathlineto{\pgfqpoint{2.428197in}{2.302857in}}%
\pgfpathlineto{\pgfqpoint{2.428532in}{2.299735in}}%
\pgfpathlineto{\pgfqpoint{2.429292in}{2.303601in}}%
\pgfpathlineto{\pgfqpoint{2.429297in}{2.303532in}}%
\pgfpathlineto{\pgfqpoint{2.429516in}{2.301182in}}%
\pgfpathlineto{\pgfqpoint{2.430300in}{2.305566in}}%
\pgfpathlineto{\pgfqpoint{2.430334in}{2.305399in}}%
\pgfpathlineto{\pgfqpoint{2.431317in}{2.308114in}}%
\pgfpathlineto{\pgfqpoint{2.431419in}{2.304898in}}%
\pgfpathlineto{\pgfqpoint{2.431600in}{2.302828in}}%
\pgfpathlineto{\pgfqpoint{2.431488in}{2.305101in}}%
\pgfpathlineto{\pgfqpoint{2.432559in}{2.303022in}}%
\pgfpathlineto{\pgfqpoint{2.432783in}{2.304719in}}%
\pgfpathlineto{\pgfqpoint{2.433133in}{2.300122in}}%
\pgfpathlineto{\pgfqpoint{2.433547in}{2.300569in}}%
\pgfpathlineto{\pgfqpoint{2.433557in}{2.300666in}}%
\pgfpathlineto{\pgfqpoint{2.433625in}{2.299207in}}%
\pgfpathlineto{\pgfqpoint{2.433630in}{2.299073in}}%
\pgfpathlineto{\pgfqpoint{2.434331in}{2.302538in}}%
\pgfpathlineto{\pgfqpoint{2.434633in}{2.301856in}}%
\pgfpathlineto{\pgfqpoint{2.435202in}{2.303346in}}%
\pgfpathlineto{\pgfqpoint{2.435650in}{2.299144in}}%
\pgfpathlineto{\pgfqpoint{2.435684in}{2.299613in}}%
\pgfpathlineto{\pgfqpoint{2.436911in}{2.303666in}}%
\pgfpathlineto{\pgfqpoint{2.435869in}{2.297420in}}%
\pgfpathlineto{\pgfqpoint{2.437223in}{2.301318in}}%
\pgfpathlineto{\pgfqpoint{2.438367in}{2.299657in}}%
\pgfpathlineto{\pgfqpoint{2.437710in}{2.303875in}}%
\pgfpathlineto{\pgfqpoint{2.438377in}{2.299822in}}%
\pgfpathlineto{\pgfqpoint{2.438415in}{2.299448in}}%
\pgfpathlineto{\pgfqpoint{2.438834in}{2.302389in}}%
\pgfpathlineto{\pgfqpoint{2.438902in}{2.302347in}}%
\pgfpathlineto{\pgfqpoint{2.439969in}{2.304162in}}%
\pgfpathlineto{\pgfqpoint{2.439194in}{2.300809in}}%
\pgfpathlineto{\pgfqpoint{2.440037in}{2.303607in}}%
\pgfpathlineto{\pgfqpoint{2.440548in}{2.301674in}}%
\pgfpathlineto{\pgfqpoint{2.441127in}{2.304272in}}%
\pgfpathlineto{\pgfqpoint{2.441142in}{2.303960in}}%
\pgfpathlineto{\pgfqpoint{2.441254in}{2.304774in}}%
\pgfpathlineto{\pgfqpoint{2.441619in}{2.302347in}}%
\pgfpathlineto{\pgfqpoint{2.442262in}{2.304611in}}%
\pgfpathlineto{\pgfqpoint{2.443440in}{2.302087in}}%
\pgfpathlineto{\pgfqpoint{2.442525in}{2.307711in}}%
\pgfpathlineto{\pgfqpoint{2.443445in}{2.302395in}}%
\pgfpathlineto{\pgfqpoint{2.443459in}{2.302548in}}%
\pgfpathlineto{\pgfqpoint{2.443995in}{2.300135in}}%
\pgfpathlineto{\pgfqpoint{2.444506in}{2.301334in}}%
\pgfpathlineto{\pgfqpoint{2.444822in}{2.300157in}}%
\pgfpathlineto{\pgfqpoint{2.445558in}{2.301959in}}%
\pgfpathlineto{\pgfqpoint{2.445616in}{2.301141in}}%
\pgfpathlineto{\pgfqpoint{2.445635in}{2.301588in}}%
\pgfpathlineto{\pgfqpoint{2.446483in}{2.298512in}}%
\pgfpathlineto{\pgfqpoint{2.446672in}{2.299554in}}%
\pgfpathlineto{\pgfqpoint{2.447578in}{2.298392in}}%
\pgfpathlineto{\pgfqpoint{2.447271in}{2.301531in}}%
\pgfpathlineto{\pgfqpoint{2.447782in}{2.299458in}}%
\pgfpathlineto{\pgfqpoint{2.448858in}{2.304682in}}%
\pgfpathlineto{\pgfqpoint{2.448036in}{2.298486in}}%
\pgfpathlineto{\pgfqpoint{2.449126in}{2.303911in}}%
\pgfpathlineto{\pgfqpoint{2.449871in}{2.301032in}}%
\pgfpathlineto{\pgfqpoint{2.450251in}{2.303040in}}%
\pgfpathlineto{\pgfqpoint{2.450577in}{2.304448in}}%
\pgfpathlineto{\pgfqpoint{2.450441in}{2.302172in}}%
\pgfpathlineto{\pgfqpoint{2.451366in}{2.303585in}}%
\pgfpathlineto{\pgfqpoint{2.451857in}{2.302285in}}%
\pgfpathlineto{\pgfqpoint{2.452325in}{2.304532in}}%
\pgfpathlineto{\pgfqpoint{2.452447in}{2.304041in}}%
\pgfpathlineto{\pgfqpoint{2.452753in}{2.305718in}}%
\pgfpathlineto{\pgfqpoint{2.453099in}{2.302711in}}%
\pgfpathlineto{\pgfqpoint{2.453547in}{2.304000in}}%
\pgfpathlineto{\pgfqpoint{2.454428in}{2.301601in}}%
\pgfpathlineto{\pgfqpoint{2.454102in}{2.304727in}}%
\pgfpathlineto{\pgfqpoint{2.454744in}{2.302130in}}%
\pgfpathlineto{\pgfqpoint{2.455421in}{2.304564in}}%
\pgfpathlineto{\pgfqpoint{2.454808in}{2.301465in}}%
\pgfpathlineto{\pgfqpoint{2.455845in}{2.301867in}}%
\pgfpathlineto{\pgfqpoint{2.455859in}{2.301488in}}%
\pgfpathlineto{\pgfqpoint{2.456312in}{2.304535in}}%
\pgfpathlineto{\pgfqpoint{2.456687in}{2.303805in}}%
\pgfpathlineto{\pgfqpoint{2.456789in}{2.304926in}}%
\pgfpathlineto{\pgfqpoint{2.457505in}{2.300866in}}%
\pgfpathlineto{\pgfqpoint{2.457763in}{2.302536in}}%
\pgfpathlineto{\pgfqpoint{2.458561in}{2.302094in}}%
\pgfpathlineto{\pgfqpoint{2.458225in}{2.304676in}}%
\pgfpathlineto{\pgfqpoint{2.458844in}{2.302977in}}%
\pgfpathlineto{\pgfqpoint{2.459408in}{2.304699in}}%
\pgfpathlineto{\pgfqpoint{2.459092in}{2.300514in}}%
\pgfpathlineto{\pgfqpoint{2.459968in}{2.303338in}}%
\pgfpathlineto{\pgfqpoint{2.459988in}{2.302620in}}%
\pgfpathlineto{\pgfqpoint{2.460538in}{2.306482in}}%
\pgfpathlineto{\pgfqpoint{2.461074in}{2.303359in}}%
\pgfpathlineto{\pgfqpoint{2.461478in}{2.304061in}}%
\pgfpathlineto{\pgfqpoint{2.461736in}{2.300873in}}%
\pgfpathlineto{\pgfqpoint{2.462203in}{2.303914in}}%
\pgfpathlineto{\pgfqpoint{2.463104in}{2.300819in}}%
\pgfpathlineto{\pgfqpoint{2.462300in}{2.304211in}}%
\pgfpathlineto{\pgfqpoint{2.463362in}{2.301696in}}%
\pgfpathlineto{\pgfqpoint{2.463664in}{2.303390in}}%
\pgfpathlineto{\pgfqpoint{2.464272in}{2.299456in}}%
\pgfpathlineto{\pgfqpoint{2.464326in}{2.299476in}}%
\pgfpathlineto{\pgfqpoint{2.465148in}{2.296368in}}%
\pgfpathlineto{\pgfqpoint{2.465450in}{2.298073in}}%
\pgfpathlineto{\pgfqpoint{2.465733in}{2.299515in}}%
\pgfpathlineto{\pgfqpoint{2.466361in}{2.295939in}}%
\pgfpathlineto{\pgfqpoint{2.466400in}{2.296321in}}%
\pgfpathlineto{\pgfqpoint{2.466740in}{2.294256in}}%
\pgfpathlineto{\pgfqpoint{2.466828in}{2.293620in}}%
\pgfpathlineto{\pgfqpoint{2.467743in}{2.299447in}}%
\pgfpathlineto{\pgfqpoint{2.468293in}{2.301425in}}%
\pgfpathlineto{\pgfqpoint{2.468493in}{2.298546in}}%
\pgfpathlineto{\pgfqpoint{2.468863in}{2.299862in}}%
\pgfpathlineto{\pgfqpoint{2.469968in}{2.296353in}}%
\pgfpathlineto{\pgfqpoint{2.470012in}{2.297198in}}%
\pgfpathlineto{\pgfqpoint{2.470893in}{2.296030in}}%
\pgfpathlineto{\pgfqpoint{2.470577in}{2.298320in}}%
\pgfpathlineto{\pgfqpoint{2.471132in}{2.297108in}}%
\pgfpathlineto{\pgfqpoint{2.471502in}{2.298553in}}%
\pgfpathlineto{\pgfqpoint{2.472193in}{2.295538in}}%
\pgfpathlineto{\pgfqpoint{2.472675in}{2.290868in}}%
\pgfpathlineto{\pgfqpoint{2.473342in}{2.293926in}}%
\pgfpathlineto{\pgfqpoint{2.474131in}{2.298323in}}%
\pgfpathlineto{\pgfqpoint{2.473464in}{2.293055in}}%
\pgfpathlineto{\pgfqpoint{2.474525in}{2.295765in}}%
\pgfpathlineto{\pgfqpoint{2.475820in}{2.299086in}}%
\pgfpathlineto{\pgfqpoint{2.474842in}{2.292783in}}%
\pgfpathlineto{\pgfqpoint{2.475942in}{2.298088in}}%
\pgfpathlineto{\pgfqpoint{2.476000in}{2.296851in}}%
\pgfpathlineto{\pgfqpoint{2.476361in}{2.299980in}}%
\pgfpathlineto{\pgfqpoint{2.477047in}{2.298440in}}%
\pgfpathlineto{\pgfqpoint{2.477422in}{2.301638in}}%
\pgfpathlineto{\pgfqpoint{2.477135in}{2.298249in}}%
\pgfpathlineto{\pgfqpoint{2.478196in}{2.300404in}}%
\pgfpathlineto{\pgfqpoint{2.478444in}{2.298841in}}%
\pgfpathlineto{\pgfqpoint{2.478727in}{2.303100in}}%
\pgfpathlineto{\pgfqpoint{2.479287in}{2.301145in}}%
\pgfpathlineto{\pgfqpoint{2.479803in}{2.306344in}}%
\pgfpathlineto{\pgfqpoint{2.479321in}{2.300581in}}%
\pgfpathlineto{\pgfqpoint{2.480440in}{2.302228in}}%
\pgfpathlineto{\pgfqpoint{2.481258in}{2.300316in}}%
\pgfpathlineto{\pgfqpoint{2.480796in}{2.304033in}}%
\pgfpathlineto{\pgfqpoint{2.481599in}{2.301410in}}%
\pgfpathlineto{\pgfqpoint{2.482407in}{2.301919in}}%
\pgfpathlineto{\pgfqpoint{2.482149in}{2.299350in}}%
\pgfpathlineto{\pgfqpoint{2.482621in}{2.300436in}}%
\pgfpathlineto{\pgfqpoint{2.483673in}{2.298592in}}%
\pgfpathlineto{\pgfqpoint{2.483357in}{2.302262in}}%
\pgfpathlineto{\pgfqpoint{2.483746in}{2.299605in}}%
\pgfpathlineto{\pgfqpoint{2.484798in}{2.303299in}}%
\pgfpathlineto{\pgfqpoint{2.484525in}{2.298101in}}%
\pgfpathlineto{\pgfqpoint{2.484954in}{2.302628in}}%
\pgfpathlineto{\pgfqpoint{2.485090in}{2.300787in}}%
\pgfpathlineto{\pgfqpoint{2.485776in}{2.303564in}}%
\pgfpathlineto{\pgfqpoint{2.486302in}{2.304084in}}%
\pgfpathlineto{\pgfqpoint{2.486414in}{2.302115in}}%
\pgfpathlineto{\pgfqpoint{2.486511in}{2.302240in}}%
\pgfpathlineto{\pgfqpoint{2.486867in}{2.300571in}}%
\pgfpathlineto{\pgfqpoint{2.487397in}{2.303874in}}%
\pgfpathlineto{\pgfqpoint{2.487587in}{2.303274in}}%
\pgfpathlineto{\pgfqpoint{2.488084in}{2.305253in}}%
\pgfpathlineto{\pgfqpoint{2.487860in}{2.301794in}}%
\pgfpathlineto{\pgfqpoint{2.488736in}{2.304382in}}%
\pgfpathlineto{\pgfqpoint{2.489871in}{2.301550in}}%
\pgfpathlineto{\pgfqpoint{2.489520in}{2.305169in}}%
\pgfpathlineto{\pgfqpoint{2.489890in}{2.302062in}}%
\pgfpathlineto{\pgfqpoint{2.489895in}{2.302174in}}%
\pgfpathlineto{\pgfqpoint{2.490421in}{2.299519in}}%
\pgfpathlineto{\pgfqpoint{2.490913in}{2.300332in}}%
\pgfpathlineto{\pgfqpoint{2.491331in}{2.299077in}}%
\pgfpathlineto{\pgfqpoint{2.491735in}{2.302291in}}%
\pgfpathlineto{\pgfqpoint{2.492023in}{2.300102in}}%
\pgfpathlineto{\pgfqpoint{2.492904in}{2.301647in}}%
\pgfpathlineto{\pgfqpoint{2.492212in}{2.299142in}}%
\pgfpathlineto{\pgfqpoint{2.493133in}{2.300122in}}%
\pgfpathlineto{\pgfqpoint{2.493619in}{2.301596in}}%
\pgfpathlineto{\pgfqpoint{2.493790in}{2.298682in}}%
\pgfpathlineto{\pgfqpoint{2.494286in}{2.297777in}}%
\pgfpathlineto{\pgfqpoint{2.494734in}{2.301141in}}%
\pgfpathlineto{\pgfqpoint{2.494866in}{2.300306in}}%
\pgfpathlineto{\pgfqpoint{2.494953in}{2.301905in}}%
\pgfpathlineto{\pgfqpoint{2.495771in}{2.297928in}}%
\pgfpathlineto{\pgfqpoint{2.495971in}{2.299877in}}%
\pgfpathlineto{\pgfqpoint{2.496156in}{2.299187in}}%
\pgfpathlineto{\pgfqpoint{2.495990in}{2.300180in}}%
\pgfpathlineto{\pgfqpoint{2.496385in}{2.300767in}}%
\pgfpathlineto{\pgfqpoint{2.497183in}{2.303798in}}%
\pgfpathlineto{\pgfqpoint{2.497514in}{2.301997in}}%
\pgfpathlineto{\pgfqpoint{2.497563in}{2.303096in}}%
\pgfpathlineto{\pgfqpoint{2.497909in}{2.300000in}}%
\pgfpathlineto{\pgfqpoint{2.498576in}{2.301395in}}%
\pgfpathlineto{\pgfqpoint{2.499788in}{2.298553in}}%
\pgfpathlineto{\pgfqpoint{2.498800in}{2.303667in}}%
\pgfpathlineto{\pgfqpoint{2.499802in}{2.298577in}}%
\pgfpathlineto{\pgfqpoint{2.500684in}{2.300492in}}%
\pgfpathlineto{\pgfqpoint{2.500289in}{2.296656in}}%
\pgfpathlineto{\pgfqpoint{2.501000in}{2.299070in}}%
\pgfpathlineto{\pgfqpoint{2.501059in}{2.298494in}}%
\pgfpathlineto{\pgfqpoint{2.501910in}{2.305058in}}%
\pgfpathlineto{\pgfqpoint{2.501935in}{2.304777in}}%
\pgfpathlineto{\pgfqpoint{2.502767in}{2.301343in}}%
\pgfpathlineto{\pgfqpoint{2.501979in}{2.304891in}}%
\pgfpathlineto{\pgfqpoint{2.503337in}{2.304192in}}%
\pgfpathlineto{\pgfqpoint{2.503751in}{2.306369in}}%
\pgfpathlineto{\pgfqpoint{2.504389in}{2.302571in}}%
\pgfpathlineto{\pgfqpoint{2.504423in}{2.302785in}}%
\pgfpathlineto{\pgfqpoint{2.505026in}{2.304592in}}%
\pgfpathlineto{\pgfqpoint{2.505416in}{2.301292in}}%
\pgfpathlineto{\pgfqpoint{2.505445in}{2.301545in}}%
\pgfpathlineto{\pgfqpoint{2.506550in}{2.299150in}}%
\pgfpathlineto{\pgfqpoint{2.505830in}{2.302118in}}%
\pgfpathlineto{\pgfqpoint{2.506604in}{2.299611in}}%
\pgfpathlineto{\pgfqpoint{2.506609in}{2.299613in}}%
\pgfpathlineto{\pgfqpoint{2.506657in}{2.298462in}}%
\pgfpathlineto{\pgfqpoint{2.507621in}{2.295807in}}%
\pgfpathlineto{\pgfqpoint{2.506823in}{2.299562in}}%
\pgfpathlineto{\pgfqpoint{2.507816in}{2.296355in}}%
\pgfpathlineto{\pgfqpoint{2.507884in}{2.296031in}}%
\pgfpathlineto{\pgfqpoint{2.509135in}{2.303026in}}%
\pgfpathlineto{\pgfqpoint{2.509793in}{2.302274in}}%
\pgfpathlineto{\pgfqpoint{2.509379in}{2.305233in}}%
\pgfpathlineto{\pgfqpoint{2.510202in}{2.304848in}}%
\pgfpathlineto{\pgfqpoint{2.510995in}{2.302143in}}%
\pgfpathlineto{\pgfqpoint{2.510211in}{2.304965in}}%
\pgfpathlineto{\pgfqpoint{2.511657in}{2.302838in}}%
\pgfpathlineto{\pgfqpoint{2.512680in}{2.304160in}}%
\pgfpathlineto{\pgfqpoint{2.512241in}{2.299829in}}%
\pgfpathlineto{\pgfqpoint{2.512772in}{2.303127in}}%
\pgfpathlineto{\pgfqpoint{2.513527in}{2.301694in}}%
\pgfpathlineto{\pgfqpoint{2.512933in}{2.305296in}}%
\pgfpathlineto{\pgfqpoint{2.513872in}{2.303188in}}%
\pgfpathlineto{\pgfqpoint{2.514724in}{2.305019in}}%
\pgfpathlineto{\pgfqpoint{2.514053in}{2.302710in}}%
\pgfpathlineto{\pgfqpoint{2.514987in}{2.303289in}}%
\pgfpathlineto{\pgfqpoint{2.515995in}{2.305767in}}%
\pgfpathlineto{\pgfqpoint{2.515299in}{2.301885in}}%
\pgfpathlineto{\pgfqpoint{2.516117in}{2.303943in}}%
\pgfpathlineto{\pgfqpoint{2.517232in}{2.302689in}}%
\pgfpathlineto{\pgfqpoint{2.516662in}{2.308249in}}%
\pgfpathlineto{\pgfqpoint{2.517246in}{2.303175in}}%
\pgfpathlineto{\pgfqpoint{2.517295in}{2.303652in}}%
\pgfpathlineto{\pgfqpoint{2.517923in}{2.299161in}}%
\pgfpathlineto{\pgfqpoint{2.518312in}{2.301257in}}%
\pgfpathlineto{\pgfqpoint{2.518347in}{2.300645in}}%
\pgfpathlineto{\pgfqpoint{2.519354in}{2.302971in}}%
\pgfpathlineto{\pgfqpoint{2.519359in}{2.302917in}}%
\pgfpathlineto{\pgfqpoint{2.520333in}{2.304282in}}%
\pgfpathlineto{\pgfqpoint{2.519530in}{2.302134in}}%
\pgfpathlineto{\pgfqpoint{2.520484in}{2.303634in}}%
\pgfpathlineto{\pgfqpoint{2.521594in}{2.300647in}}%
\pgfpathlineto{\pgfqpoint{2.520946in}{2.305845in}}%
\pgfpathlineto{\pgfqpoint{2.521657in}{2.301036in}}%
\pgfpathlineto{\pgfqpoint{2.521764in}{2.302243in}}%
\pgfpathlineto{\pgfqpoint{2.522689in}{2.297476in}}%
\pgfpathlineto{\pgfqpoint{2.522728in}{2.297690in}}%
\pgfpathlineto{\pgfqpoint{2.523780in}{2.296910in}}%
\pgfpathlineto{\pgfqpoint{2.523410in}{2.300541in}}%
\pgfpathlineto{\pgfqpoint{2.523838in}{2.297359in}}%
\pgfpathlineto{\pgfqpoint{2.524729in}{2.301228in}}%
\pgfpathlineto{\pgfqpoint{2.523892in}{2.296991in}}%
\pgfpathlineto{\pgfqpoint{2.525012in}{2.299787in}}%
\pgfpathlineto{\pgfqpoint{2.525601in}{2.297993in}}%
\pgfpathlineto{\pgfqpoint{2.525941in}{2.300458in}}%
\pgfpathlineto{\pgfqpoint{2.526141in}{2.298984in}}%
\pgfpathlineto{\pgfqpoint{2.526380in}{2.301518in}}%
\pgfpathlineto{\pgfqpoint{2.527295in}{2.300945in}}%
\pgfpathlineto{\pgfqpoint{2.527543in}{2.300473in}}%
\pgfpathlineto{\pgfqpoint{2.528001in}{2.304372in}}%
\pgfpathlineto{\pgfqpoint{2.528390in}{2.301484in}}%
\pgfpathlineto{\pgfqpoint{2.528770in}{2.299694in}}%
\pgfpathlineto{\pgfqpoint{2.529091in}{2.303476in}}%
\pgfpathlineto{\pgfqpoint{2.529442in}{2.302197in}}%
\pgfpathlineto{\pgfqpoint{2.529641in}{2.304353in}}%
\pgfpathlineto{\pgfqpoint{2.530401in}{2.297964in}}%
\pgfpathlineto{\pgfqpoint{2.530503in}{2.299316in}}%
\pgfpathlineto{\pgfqpoint{2.531681in}{2.304589in}}%
\pgfpathlineto{\pgfqpoint{2.531764in}{2.303999in}}%
\pgfpathlineto{\pgfqpoint{2.532149in}{2.301493in}}%
\pgfpathlineto{\pgfqpoint{2.532733in}{2.304520in}}%
\pgfpathlineto{\pgfqpoint{2.532825in}{2.304198in}}%
\pgfpathlineto{\pgfqpoint{2.532845in}{2.304414in}}%
\pgfpathlineto{\pgfqpoint{2.533760in}{2.297925in}}%
\pgfpathlineto{\pgfqpoint{2.533785in}{2.298426in}}%
\pgfpathlineto{\pgfqpoint{2.535440in}{2.293596in}}%
\pgfpathlineto{\pgfqpoint{2.533877in}{2.299874in}}%
\pgfpathlineto{\pgfqpoint{2.535615in}{2.294461in}}%
\pgfpathlineto{\pgfqpoint{2.536711in}{2.300000in}}%
\pgfpathlineto{\pgfqpoint{2.536871in}{2.299109in}}%
\pgfpathlineto{\pgfqpoint{2.537543in}{2.298327in}}%
\pgfpathlineto{\pgfqpoint{2.537835in}{2.301274in}}%
\pgfpathlineto{\pgfqpoint{2.537967in}{2.299863in}}%
\pgfpathlineto{\pgfqpoint{2.538750in}{2.305794in}}%
\pgfpathlineto{\pgfqpoint{2.539164in}{2.304019in}}%
\pgfpathlineto{\pgfqpoint{2.540255in}{2.301750in}}%
\pgfpathlineto{\pgfqpoint{2.539612in}{2.304475in}}%
\pgfpathlineto{\pgfqpoint{2.540304in}{2.302793in}}%
\pgfpathlineto{\pgfqpoint{2.540678in}{2.303893in}}%
\pgfpathlineto{\pgfqpoint{2.541277in}{2.299192in}}%
\pgfpathlineto{\pgfqpoint{2.542324in}{2.297594in}}%
\pgfpathlineto{\pgfqpoint{2.541896in}{2.300557in}}%
\pgfpathlineto{\pgfqpoint{2.542407in}{2.298104in}}%
\pgfpathlineto{\pgfqpoint{2.543127in}{2.299170in}}%
\pgfpathlineto{\pgfqpoint{2.542455in}{2.297042in}}%
\pgfpathlineto{\pgfqpoint{2.543268in}{2.297216in}}%
\pgfpathlineto{\pgfqpoint{2.543896in}{2.293802in}}%
\pgfpathlineto{\pgfqpoint{2.544447in}{2.295117in}}%
\pgfpathlineto{\pgfqpoint{2.544466in}{2.295454in}}%
\pgfpathlineto{\pgfqpoint{2.544875in}{2.293061in}}%
\pgfpathlineto{\pgfqpoint{2.545537in}{2.294794in}}%
\pgfpathlineto{\pgfqpoint{2.545566in}{2.294326in}}%
\pgfpathlineto{\pgfqpoint{2.546584in}{2.298622in}}%
\pgfpathlineto{\pgfqpoint{2.546608in}{2.298100in}}%
\pgfpathlineto{\pgfqpoint{2.547353in}{2.295831in}}%
\pgfpathlineto{\pgfqpoint{2.547115in}{2.299087in}}%
\pgfpathlineto{\pgfqpoint{2.547728in}{2.298017in}}%
\pgfpathlineto{\pgfqpoint{2.547747in}{2.298573in}}%
\pgfpathlineto{\pgfqpoint{2.548463in}{2.294850in}}%
\pgfpathlineto{\pgfqpoint{2.548809in}{2.295511in}}%
\pgfpathlineto{\pgfqpoint{2.549130in}{2.295009in}}%
\pgfpathlineto{\pgfqpoint{2.549651in}{2.300010in}}%
\pgfpathlineto{\pgfqpoint{2.549661in}{2.299945in}}%
\pgfpathlineto{\pgfqpoint{2.550664in}{2.303751in}}%
\pgfpathlineto{\pgfqpoint{2.550805in}{2.302875in}}%
\pgfpathlineto{\pgfqpoint{2.551413in}{2.304678in}}%
\pgfpathlineto{\pgfqpoint{2.550951in}{2.301945in}}%
\pgfpathlineto{\pgfqpoint{2.551939in}{2.303609in}}%
\pgfpathlineto{\pgfqpoint{2.552037in}{2.302481in}}%
\pgfpathlineto{\pgfqpoint{2.552548in}{2.305348in}}%
\pgfpathlineto{\pgfqpoint{2.553069in}{2.302904in}}%
\pgfpathlineto{\pgfqpoint{2.553074in}{2.302997in}}%
\pgfpathlineto{\pgfqpoint{2.553867in}{2.298656in}}%
\pgfpathlineto{\pgfqpoint{2.553989in}{2.299481in}}%
\pgfpathlineto{\pgfqpoint{2.554003in}{2.299069in}}%
\pgfpathlineto{\pgfqpoint{2.554705in}{2.303396in}}%
\pgfpathlineto{\pgfqpoint{2.555055in}{2.301689in}}%
\pgfpathlineto{\pgfqpoint{2.555975in}{2.304794in}}%
\pgfpathlineto{\pgfqpoint{2.555668in}{2.301091in}}%
\pgfpathlineto{\pgfqpoint{2.556301in}{2.303795in}}%
\pgfpathlineto{\pgfqpoint{2.556808in}{2.301253in}}%
\pgfpathlineto{\pgfqpoint{2.556511in}{2.304265in}}%
\pgfpathlineto{\pgfqpoint{2.557411in}{2.303695in}}%
\pgfpathlineto{\pgfqpoint{2.558322in}{2.300220in}}%
\pgfpathlineto{\pgfqpoint{2.557557in}{2.304066in}}%
\pgfpathlineto{\pgfqpoint{2.558580in}{2.301148in}}%
\pgfpathlineto{\pgfqpoint{2.558633in}{2.300113in}}%
\pgfpathlineto{\pgfqpoint{2.559802in}{2.304321in}}%
\pgfpathlineto{\pgfqpoint{2.560508in}{2.301247in}}%
\pgfpathlineto{\pgfqpoint{2.559997in}{2.305188in}}%
\pgfpathlineto{\pgfqpoint{2.560980in}{2.302555in}}%
\pgfpathlineto{\pgfqpoint{2.561423in}{2.304561in}}%
\pgfpathlineto{\pgfqpoint{2.561063in}{2.301726in}}%
\pgfpathlineto{\pgfqpoint{2.562110in}{2.303587in}}%
\pgfpathlineto{\pgfqpoint{2.562665in}{2.305551in}}%
\pgfpathlineto{\pgfqpoint{2.563093in}{2.301739in}}%
\pgfpathlineto{\pgfqpoint{2.563137in}{2.301926in}}%
\pgfpathlineto{\pgfqpoint{2.563857in}{2.298445in}}%
\pgfpathlineto{\pgfqpoint{2.564222in}{2.301941in}}%
\pgfpathlineto{\pgfqpoint{2.564247in}{2.301906in}}%
\pgfpathlineto{\pgfqpoint{2.564257in}{2.302091in}}%
\pgfpathlineto{\pgfqpoint{2.564461in}{2.299517in}}%
\pgfpathlineto{\pgfqpoint{2.565308in}{2.301264in}}%
\pgfpathlineto{\pgfqpoint{2.565332in}{2.300811in}}%
\pgfpathlineto{\pgfqpoint{2.566087in}{2.304570in}}%
\pgfpathlineto{\pgfqpoint{2.566306in}{2.304348in}}%
\pgfpathlineto{\pgfqpoint{2.566321in}{2.304466in}}%
\pgfpathlineto{\pgfqpoint{2.567056in}{2.300369in}}%
\pgfpathlineto{\pgfqpoint{2.567183in}{2.301235in}}%
\pgfpathlineto{\pgfqpoint{2.567265in}{2.300793in}}%
\pgfpathlineto{\pgfqpoint{2.568244in}{2.303965in}}%
\pgfpathlineto{\pgfqpoint{2.568682in}{2.304887in}}%
\pgfpathlineto{\pgfqpoint{2.568414in}{2.301887in}}%
\pgfpathlineto{\pgfqpoint{2.569193in}{2.301989in}}%
\pgfpathlineto{\pgfqpoint{2.570079in}{2.300108in}}%
\pgfpathlineto{\pgfqpoint{2.569748in}{2.303745in}}%
\pgfpathlineto{\pgfqpoint{2.570279in}{2.302522in}}%
\pgfpathlineto{\pgfqpoint{2.570498in}{2.304863in}}%
\pgfpathlineto{\pgfqpoint{2.571506in}{2.303737in}}%
\pgfpathlineto{\pgfqpoint{2.572265in}{2.302575in}}%
\pgfpathlineto{\pgfqpoint{2.571871in}{2.304511in}}%
\pgfpathlineto{\pgfqpoint{2.572616in}{2.303331in}}%
\pgfpathlineto{\pgfqpoint{2.573030in}{2.305160in}}%
\pgfpathlineto{\pgfqpoint{2.573161in}{2.301976in}}%
\pgfpathlineto{\pgfqpoint{2.573726in}{2.303457in}}%
\pgfpathlineto{\pgfqpoint{2.574763in}{2.304158in}}%
\pgfpathlineto{\pgfqpoint{2.574388in}{2.301815in}}%
\pgfpathlineto{\pgfqpoint{2.574870in}{2.303860in}}%
\pgfpathlineto{\pgfqpoint{2.575527in}{2.301368in}}%
\pgfpathlineto{\pgfqpoint{2.574953in}{2.304031in}}%
\pgfpathlineto{\pgfqpoint{2.576033in}{2.301929in}}%
\pgfpathlineto{\pgfqpoint{2.576924in}{2.303946in}}%
\pgfpathlineto{\pgfqpoint{2.576102in}{2.300998in}}%
\pgfpathlineto{\pgfqpoint{2.577173in}{2.302969in}}%
\pgfpathlineto{\pgfqpoint{2.577586in}{2.300424in}}%
\pgfpathlineto{\pgfqpoint{2.577236in}{2.303385in}}%
\pgfpathlineto{\pgfqpoint{2.578278in}{2.302840in}}%
\pgfpathlineto{\pgfqpoint{2.578322in}{2.304691in}}%
\pgfpathlineto{\pgfqpoint{2.579349in}{2.302082in}}%
\pgfpathlineto{\pgfqpoint{2.579388in}{2.303099in}}%
\pgfpathlineto{\pgfqpoint{2.579865in}{2.302478in}}%
\pgfpathlineto{\pgfqpoint{2.580108in}{2.304872in}}%
\pgfpathlineto{\pgfqpoint{2.580503in}{2.302923in}}%
\pgfpathlineto{\pgfqpoint{2.581632in}{2.300715in}}%
\pgfpathlineto{\pgfqpoint{2.580722in}{2.304498in}}%
\pgfpathlineto{\pgfqpoint{2.581734in}{2.302281in}}%
\pgfpathlineto{\pgfqpoint{2.582173in}{2.303404in}}%
\pgfpathlineto{\pgfqpoint{2.582752in}{2.298746in}}%
\pgfpathlineto{\pgfqpoint{2.583653in}{2.297294in}}%
\pgfpathlineto{\pgfqpoint{2.583161in}{2.299780in}}%
\pgfpathlineto{\pgfqpoint{2.583872in}{2.298838in}}%
\pgfpathlineto{\pgfqpoint{2.585259in}{2.305006in}}%
\pgfpathlineto{\pgfqpoint{2.583891in}{2.298354in}}%
\pgfpathlineto{\pgfqpoint{2.585376in}{2.302813in}}%
\pgfpathlineto{\pgfqpoint{2.586184in}{2.301001in}}%
\pgfpathlineto{\pgfqpoint{2.585907in}{2.304393in}}%
\pgfpathlineto{\pgfqpoint{2.586437in}{2.303724in}}%
\pgfpathlineto{\pgfqpoint{2.587698in}{2.306381in}}%
\pgfpathlineto{\pgfqpoint{2.586598in}{2.302515in}}%
\pgfpathlineto{\pgfqpoint{2.587742in}{2.305425in}}%
\pgfpathlineto{\pgfqpoint{2.588356in}{2.302249in}}%
\pgfpathlineto{\pgfqpoint{2.588916in}{2.303126in}}%
\pgfpathlineto{\pgfqpoint{2.589948in}{2.305607in}}%
\pgfpathlineto{\pgfqpoint{2.589461in}{2.302409in}}%
\pgfpathlineto{\pgfqpoint{2.590040in}{2.303844in}}%
\pgfpathlineto{\pgfqpoint{2.590999in}{2.304478in}}%
\pgfpathlineto{\pgfqpoint{2.590522in}{2.301325in}}%
\pgfpathlineto{\pgfqpoint{2.591155in}{2.304221in}}%
\pgfpathlineto{\pgfqpoint{2.591613in}{2.301882in}}%
\pgfpathlineto{\pgfqpoint{2.592061in}{2.305207in}}%
\pgfpathlineto{\pgfqpoint{2.592289in}{2.303411in}}%
\pgfpathlineto{\pgfqpoint{2.592854in}{2.306069in}}%
\pgfpathlineto{\pgfqpoint{2.593219in}{2.300838in}}%
\pgfpathlineto{\pgfqpoint{2.593321in}{2.302312in}}%
\pgfpathlineto{\pgfqpoint{2.593545in}{2.302651in}}%
\pgfpathlineto{\pgfqpoint{2.594505in}{2.298264in}}%
\pgfpathlineto{\pgfqpoint{2.594699in}{2.300367in}}%
\pgfpathlineto{\pgfqpoint{2.595473in}{2.293100in}}%
\pgfpathlineto{\pgfqpoint{2.595507in}{2.293955in}}%
\pgfpathlineto{\pgfqpoint{2.595819in}{2.295736in}}%
\pgfpathlineto{\pgfqpoint{2.596199in}{2.292950in}}%
\pgfpathlineto{\pgfqpoint{2.596238in}{2.292941in}}%
\pgfpathlineto{\pgfqpoint{2.596452in}{2.292016in}}%
\pgfpathlineto{\pgfqpoint{2.597095in}{2.297871in}}%
\pgfpathlineto{\pgfqpoint{2.598166in}{2.300838in}}%
\pgfpathlineto{\pgfqpoint{2.598273in}{2.300259in}}%
\pgfpathlineto{\pgfqpoint{2.599329in}{2.302808in}}%
\pgfpathlineto{\pgfqpoint{2.599466in}{2.301606in}}%
\pgfpathlineto{\pgfqpoint{2.599578in}{2.298855in}}%
\pgfpathlineto{\pgfqpoint{2.599889in}{2.302898in}}%
\pgfpathlineto{\pgfqpoint{2.600595in}{2.301112in}}%
\pgfpathlineto{\pgfqpoint{2.600619in}{2.301524in}}%
\pgfpathlineto{\pgfqpoint{2.601311in}{2.295571in}}%
\pgfpathlineto{\pgfqpoint{2.601393in}{2.294965in}}%
\pgfpathlineto{\pgfqpoint{2.602007in}{2.297852in}}%
\pgfpathlineto{\pgfqpoint{2.602382in}{2.296157in}}%
\pgfpathlineto{\pgfqpoint{2.602387in}{2.296358in}}%
\pgfpathlineto{\pgfqpoint{2.602951in}{2.292183in}}%
\pgfpathlineto{\pgfqpoint{2.603414in}{2.293190in}}%
\pgfpathlineto{\pgfqpoint{2.604558in}{2.288721in}}%
\pgfpathlineto{\pgfqpoint{2.603487in}{2.293193in}}%
\pgfpathlineto{\pgfqpoint{2.604607in}{2.289657in}}%
\pgfpathlineto{\pgfqpoint{2.605332in}{2.292157in}}%
\pgfpathlineto{\pgfqpoint{2.605683in}{2.287618in}}%
\pgfpathlineto{\pgfqpoint{2.606126in}{2.289528in}}%
\pgfpathlineto{\pgfqpoint{2.606656in}{2.286750in}}%
\pgfpathlineto{\pgfqpoint{2.606841in}{2.288577in}}%
\pgfpathlineto{\pgfqpoint{2.606934in}{2.287310in}}%
\pgfpathlineto{\pgfqpoint{2.607878in}{2.291350in}}%
\pgfpathlineto{\pgfqpoint{2.609076in}{2.289614in}}%
\pgfpathlineto{\pgfqpoint{2.608102in}{2.293886in}}%
\pgfpathlineto{\pgfqpoint{2.609105in}{2.290316in}}%
\pgfpathlineto{\pgfqpoint{2.610162in}{2.294425in}}%
\pgfpathlineto{\pgfqpoint{2.610293in}{2.293542in}}%
\pgfpathlineto{\pgfqpoint{2.610420in}{2.293107in}}%
\pgfpathlineto{\pgfqpoint{2.611291in}{2.295313in}}%
\pgfpathlineto{\pgfqpoint{2.611354in}{2.294926in}}%
\pgfpathlineto{\pgfqpoint{2.611807in}{2.297413in}}%
\pgfpathlineto{\pgfqpoint{2.612216in}{2.292971in}}%
\pgfpathlineto{\pgfqpoint{2.612425in}{2.293013in}}%
\pgfpathlineto{\pgfqpoint{2.613370in}{2.291565in}}%
\pgfpathlineto{\pgfqpoint{2.612727in}{2.293978in}}%
\pgfpathlineto{\pgfqpoint{2.613526in}{2.293090in}}%
\pgfpathlineto{\pgfqpoint{2.613847in}{2.295558in}}%
\pgfpathlineto{\pgfqpoint{2.614611in}{2.292746in}}%
\pgfpathlineto{\pgfqpoint{2.614631in}{2.292963in}}%
\pgfpathlineto{\pgfqpoint{2.614918in}{2.291013in}}%
\pgfpathlineto{\pgfqpoint{2.615571in}{2.296714in}}%
\pgfpathlineto{\pgfqpoint{2.615726in}{2.294042in}}%
\pgfpathlineto{\pgfqpoint{2.616734in}{2.295002in}}%
\pgfpathlineto{\pgfqpoint{2.615994in}{2.291501in}}%
\pgfpathlineto{\pgfqpoint{2.616802in}{2.293655in}}%
\pgfpathlineto{\pgfqpoint{2.616963in}{2.291419in}}%
\pgfpathlineto{\pgfqpoint{2.617854in}{2.294931in}}%
\pgfpathlineto{\pgfqpoint{2.617883in}{2.294432in}}%
\pgfpathlineto{\pgfqpoint{2.619349in}{2.300723in}}%
\pgfpathlineto{\pgfqpoint{2.619465in}{2.300443in}}%
\pgfpathlineto{\pgfqpoint{2.619616in}{2.298269in}}%
\pgfpathlineto{\pgfqpoint{2.620332in}{2.301775in}}%
\pgfpathlineto{\pgfqpoint{2.620571in}{2.300817in}}%
\pgfpathlineto{\pgfqpoint{2.621486in}{2.303682in}}%
\pgfpathlineto{\pgfqpoint{2.620663in}{2.300520in}}%
\pgfpathlineto{\pgfqpoint{2.621749in}{2.301118in}}%
\pgfpathlineto{\pgfqpoint{2.621768in}{2.300696in}}%
\pgfpathlineto{\pgfqpoint{2.622567in}{2.304944in}}%
\pgfpathlineto{\pgfqpoint{2.622829in}{2.302651in}}%
\pgfpathlineto{\pgfqpoint{2.623623in}{2.303908in}}%
\pgfpathlineto{\pgfqpoint{2.623896in}{2.300884in}}%
\pgfpathlineto{\pgfqpoint{2.623905in}{2.300933in}}%
\pgfpathlineto{\pgfqpoint{2.624047in}{2.299662in}}%
\pgfpathlineto{\pgfqpoint{2.624894in}{2.306005in}}%
\pgfpathlineto{\pgfqpoint{2.624981in}{2.304929in}}%
\pgfpathlineto{\pgfqpoint{2.625298in}{2.305846in}}%
\pgfpathlineto{\pgfqpoint{2.625468in}{2.303883in}}%
\pgfpathlineto{\pgfqpoint{2.625478in}{2.303629in}}%
\pgfpathlineto{\pgfqpoint{2.626310in}{2.306623in}}%
\pgfpathlineto{\pgfqpoint{2.626544in}{2.304845in}}%
\pgfpathlineto{\pgfqpoint{2.626607in}{2.306936in}}%
\pgfpathlineto{\pgfqpoint{2.627518in}{2.302205in}}%
\pgfpathlineto{\pgfqpoint{2.627620in}{2.303033in}}%
\pgfpathlineto{\pgfqpoint{2.628438in}{2.301031in}}%
\pgfpathlineto{\pgfqpoint{2.628195in}{2.304045in}}%
\pgfpathlineto{\pgfqpoint{2.628735in}{2.302656in}}%
\pgfpathlineto{\pgfqpoint{2.628784in}{2.303440in}}%
\pgfpathlineto{\pgfqpoint{2.629159in}{2.300241in}}%
\pgfpathlineto{\pgfqpoint{2.629733in}{2.301500in}}%
\pgfpathlineto{\pgfqpoint{2.629894in}{2.300035in}}%
\pgfpathlineto{\pgfqpoint{2.630522in}{2.304139in}}%
\pgfpathlineto{\pgfqpoint{2.630536in}{2.304048in}}%
\pgfpathlineto{\pgfqpoint{2.631320in}{2.304817in}}%
\pgfpathlineto{\pgfqpoint{2.631583in}{2.301235in}}%
\pgfpathlineto{\pgfqpoint{2.631612in}{2.302045in}}%
\pgfpathlineto{\pgfqpoint{2.631763in}{2.300688in}}%
\pgfpathlineto{\pgfqpoint{2.632381in}{2.297512in}}%
\pgfpathlineto{\pgfqpoint{2.631934in}{2.301079in}}%
\pgfpathlineto{\pgfqpoint{2.632893in}{2.299239in}}%
\pgfpathlineto{\pgfqpoint{2.634246in}{2.304536in}}%
\pgfpathlineto{\pgfqpoint{2.632902in}{2.299007in}}%
\pgfpathlineto{\pgfqpoint{2.634329in}{2.303347in}}%
\pgfpathlineto{\pgfqpoint{2.635074in}{2.300644in}}%
\pgfpathlineto{\pgfqpoint{2.634616in}{2.304673in}}%
\pgfpathlineto{\pgfqpoint{2.635463in}{2.302696in}}%
\pgfpathlineto{\pgfqpoint{2.635473in}{2.302899in}}%
\pgfpathlineto{\pgfqpoint{2.636174in}{2.298477in}}%
\pgfpathlineto{\pgfqpoint{2.636568in}{2.302666in}}%
\pgfpathlineto{\pgfqpoint{2.636778in}{2.301679in}}%
\pgfpathlineto{\pgfqpoint{2.637372in}{2.304543in}}%
\pgfpathlineto{\pgfqpoint{2.637635in}{2.303549in}}%
\pgfpathlineto{\pgfqpoint{2.638915in}{2.300660in}}%
\pgfpathlineto{\pgfqpoint{2.638121in}{2.305802in}}%
\pgfpathlineto{\pgfqpoint{2.639119in}{2.301978in}}%
\pgfpathlineto{\pgfqpoint{2.639402in}{2.303896in}}%
\pgfpathlineto{\pgfqpoint{2.640176in}{2.299258in}}%
\pgfpathlineto{\pgfqpoint{2.640181in}{2.299265in}}%
\pgfpathlineto{\pgfqpoint{2.640186in}{2.299191in}}%
\pgfpathlineto{\pgfqpoint{2.641013in}{2.302817in}}%
\pgfpathlineto{\pgfqpoint{2.641091in}{2.302144in}}%
\pgfpathlineto{\pgfqpoint{2.641841in}{2.304068in}}%
\pgfpathlineto{\pgfqpoint{2.642187in}{2.301231in}}%
\pgfpathlineto{\pgfqpoint{2.642333in}{2.298814in}}%
\pgfpathlineto{\pgfqpoint{2.642922in}{2.302454in}}%
\pgfpathlineto{\pgfqpoint{2.643287in}{2.301277in}}%
\pgfpathlineto{\pgfqpoint{2.644115in}{2.303990in}}%
\pgfpathlineto{\pgfqpoint{2.643443in}{2.300880in}}%
\pgfpathlineto{\pgfqpoint{2.644421in}{2.302263in}}%
\pgfpathlineto{\pgfqpoint{2.645522in}{2.301334in}}%
\pgfpathlineto{\pgfqpoint{2.644840in}{2.304056in}}%
\pgfpathlineto{\pgfqpoint{2.645546in}{2.301414in}}%
\pgfpathlineto{\pgfqpoint{2.645965in}{2.303899in}}%
\pgfpathlineto{\pgfqpoint{2.645585in}{2.300960in}}%
\pgfpathlineto{\pgfqpoint{2.646705in}{2.301990in}}%
\pgfpathlineto{\pgfqpoint{2.646802in}{2.301004in}}%
\pgfpathlineto{\pgfqpoint{2.647381in}{2.303763in}}%
\pgfpathlineto{\pgfqpoint{2.647810in}{2.302051in}}%
\pgfpathlineto{\pgfqpoint{2.649192in}{2.304542in}}%
\pgfpathlineto{\pgfqpoint{2.648720in}{2.301465in}}%
\pgfpathlineto{\pgfqpoint{2.649377in}{2.302709in}}%
\pgfpathlineto{\pgfqpoint{2.650200in}{2.301327in}}%
\pgfpathlineto{\pgfqpoint{2.649830in}{2.304646in}}%
\pgfpathlineto{\pgfqpoint{2.650458in}{2.304001in}}%
\pgfpathlineto{\pgfqpoint{2.650468in}{2.304373in}}%
\pgfpathlineto{\pgfqpoint{2.651451in}{2.301006in}}%
\pgfpathlineto{\pgfqpoint{2.651461in}{2.301179in}}%
\pgfpathlineto{\pgfqpoint{2.651641in}{2.300053in}}%
\pgfpathlineto{\pgfqpoint{2.651787in}{2.302663in}}%
\pgfpathlineto{\pgfqpoint{2.651841in}{2.303580in}}%
\pgfpathlineto{\pgfqpoint{2.652532in}{2.301129in}}%
\pgfpathlineto{\pgfqpoint{2.652863in}{2.301637in}}%
\pgfpathlineto{\pgfqpoint{2.653452in}{2.299626in}}%
\pgfpathlineto{\pgfqpoint{2.652888in}{2.301926in}}%
\pgfpathlineto{\pgfqpoint{2.653983in}{2.301496in}}%
\pgfpathlineto{\pgfqpoint{2.654105in}{2.303006in}}%
\pgfpathlineto{\pgfqpoint{2.654665in}{2.300569in}}%
\pgfpathlineto{\pgfqpoint{2.655103in}{2.301967in}}%
\pgfpathlineto{\pgfqpoint{2.656008in}{2.296332in}}%
\pgfpathlineto{\pgfqpoint{2.656320in}{2.297384in}}%
\pgfpathlineto{\pgfqpoint{2.656344in}{2.296997in}}%
\pgfpathlineto{\pgfqpoint{2.656427in}{2.298753in}}%
\pgfpathlineto{\pgfqpoint{2.656578in}{2.298602in}}%
\pgfpathlineto{\pgfqpoint{2.656670in}{2.299497in}}%
\pgfpathlineto{\pgfqpoint{2.657557in}{2.295692in}}%
\pgfpathlineto{\pgfqpoint{2.657595in}{2.295937in}}%
\pgfpathlineto{\pgfqpoint{2.657600in}{2.295908in}}%
\pgfpathlineto{\pgfqpoint{2.658155in}{2.298987in}}%
\pgfpathlineto{\pgfqpoint{2.658209in}{2.298990in}}%
\pgfpathlineto{\pgfqpoint{2.659041in}{2.302405in}}%
\pgfpathlineto{\pgfqpoint{2.659358in}{2.301343in}}%
\pgfpathlineto{\pgfqpoint{2.659791in}{2.303712in}}%
\pgfpathlineto{\pgfqpoint{2.660156in}{2.300798in}}%
\pgfpathlineto{\pgfqpoint{2.660468in}{2.301338in}}%
\pgfpathlineto{\pgfqpoint{2.660536in}{2.300547in}}%
\pgfpathlineto{\pgfqpoint{2.661519in}{2.304567in}}%
\pgfpathlineto{\pgfqpoint{2.663165in}{2.300893in}}%
\pgfpathlineto{\pgfqpoint{2.663262in}{2.302325in}}%
\pgfpathlineto{\pgfqpoint{2.663501in}{2.305108in}}%
\pgfpathlineto{\pgfqpoint{2.664392in}{2.304286in}}%
\pgfpathlineto{\pgfqpoint{2.665687in}{2.299391in}}%
\pgfpathlineto{\pgfqpoint{2.664402in}{2.304432in}}%
\pgfpathlineto{\pgfqpoint{2.665702in}{2.299716in}}%
\pgfpathlineto{\pgfqpoint{2.666169in}{2.302213in}}%
\pgfpathlineto{\pgfqpoint{2.666709in}{2.298394in}}%
\pgfpathlineto{\pgfqpoint{2.666777in}{2.299054in}}%
\pgfpathlineto{\pgfqpoint{2.666782in}{2.298649in}}%
\pgfpathlineto{\pgfqpoint{2.667824in}{2.304291in}}%
\pgfpathlineto{\pgfqpoint{2.667829in}{2.303993in}}%
\pgfpathlineto{\pgfqpoint{2.668774in}{2.304771in}}%
\pgfpathlineto{\pgfqpoint{2.668316in}{2.301812in}}%
\pgfpathlineto{\pgfqpoint{2.668934in}{2.303845in}}%
\pgfpathlineto{\pgfqpoint{2.669382in}{2.301484in}}%
\pgfpathlineto{\pgfqpoint{2.669139in}{2.304714in}}%
\pgfpathlineto{\pgfqpoint{2.670073in}{2.302900in}}%
\pgfpathlineto{\pgfqpoint{2.671096in}{2.306517in}}%
\pgfpathlineto{\pgfqpoint{2.670307in}{2.301979in}}%
\pgfpathlineto{\pgfqpoint{2.671203in}{2.305641in}}%
\pgfpathlineto{\pgfqpoint{2.672308in}{2.298962in}}%
\pgfpathlineto{\pgfqpoint{2.672410in}{2.300520in}}%
\pgfpathlineto{\pgfqpoint{2.673569in}{2.305570in}}%
\pgfpathlineto{\pgfqpoint{2.673009in}{2.299592in}}%
\pgfpathlineto{\pgfqpoint{2.673642in}{2.305066in}}%
\pgfpathlineto{\pgfqpoint{2.674611in}{2.302256in}}%
\pgfpathlineto{\pgfqpoint{2.673676in}{2.305378in}}%
\pgfpathlineto{\pgfqpoint{2.674786in}{2.303429in}}%
\pgfpathlineto{\pgfqpoint{2.674791in}{2.303617in}}%
\pgfpathlineto{\pgfqpoint{2.675789in}{2.299416in}}%
\pgfpathlineto{\pgfqpoint{2.676661in}{2.297337in}}%
\pgfpathlineto{\pgfqpoint{2.675984in}{2.299861in}}%
\pgfpathlineto{\pgfqpoint{2.677045in}{2.298420in}}%
\pgfpathlineto{\pgfqpoint{2.678199in}{2.301385in}}%
\pgfpathlineto{\pgfqpoint{2.677060in}{2.298397in}}%
\pgfpathlineto{\pgfqpoint{2.678243in}{2.300913in}}%
\pgfpathlineto{\pgfqpoint{2.678277in}{2.300207in}}%
\pgfpathlineto{\pgfqpoint{2.678803in}{2.303295in}}%
\pgfpathlineto{\pgfqpoint{2.679333in}{2.301559in}}%
\pgfpathlineto{\pgfqpoint{2.679353in}{2.302199in}}%
\pgfpathlineto{\pgfqpoint{2.679864in}{2.298147in}}%
\pgfpathlineto{\pgfqpoint{2.680424in}{2.300253in}}%
\pgfpathlineto{\pgfqpoint{2.680429in}{2.300322in}}%
\pgfpathlineto{\pgfqpoint{2.681208in}{2.297392in}}%
\pgfpathlineto{\pgfqpoint{2.681378in}{2.298492in}}%
\pgfpathlineto{\pgfqpoint{2.681860in}{2.296522in}}%
\pgfpathlineto{\pgfqpoint{2.682439in}{2.300781in}}%
\pgfpathlineto{\pgfqpoint{2.682469in}{2.300288in}}%
\pgfpathlineto{\pgfqpoint{2.682834in}{2.302594in}}%
\pgfpathlineto{\pgfqpoint{2.683111in}{2.298251in}}%
\pgfpathlineto{\pgfqpoint{2.683554in}{2.299450in}}%
\pgfpathlineto{\pgfqpoint{2.684367in}{2.297892in}}%
\pgfpathlineto{\pgfqpoint{2.684689in}{2.299115in}}%
\pgfpathlineto{\pgfqpoint{2.685677in}{2.302849in}}%
\pgfpathlineto{\pgfqpoint{2.684932in}{2.298294in}}%
\pgfpathlineto{\pgfqpoint{2.685847in}{2.301065in}}%
\pgfpathlineto{\pgfqpoint{2.686071in}{2.298736in}}%
\pgfpathlineto{\pgfqpoint{2.686933in}{2.302345in}}%
\pgfpathlineto{\pgfqpoint{2.686938in}{2.302224in}}%
\pgfpathlineto{\pgfqpoint{2.687269in}{2.305256in}}%
\pgfpathlineto{\pgfqpoint{2.687702in}{2.299984in}}%
\pgfpathlineto{\pgfqpoint{2.687756in}{2.300118in}}%
\pgfpathlineto{\pgfqpoint{2.689080in}{2.293767in}}%
\pgfpathlineto{\pgfqpoint{2.687800in}{2.300692in}}%
\pgfpathlineto{\pgfqpoint{2.689187in}{2.294095in}}%
\pgfpathlineto{\pgfqpoint{2.689655in}{2.296726in}}%
\pgfpathlineto{\pgfqpoint{2.690273in}{2.291583in}}%
\pgfpathlineto{\pgfqpoint{2.691032in}{2.294219in}}%
\pgfpathlineto{\pgfqpoint{2.690370in}{2.290883in}}%
\pgfpathlineto{\pgfqpoint{2.691417in}{2.292348in}}%
\pgfpathlineto{\pgfqpoint{2.691646in}{2.290548in}}%
\pgfpathlineto{\pgfqpoint{2.692444in}{2.294657in}}%
\pgfpathlineto{\pgfqpoint{2.692493in}{2.294394in}}%
\pgfpathlineto{\pgfqpoint{2.692551in}{2.293714in}}%
\pgfpathlineto{\pgfqpoint{2.693267in}{2.298860in}}%
\pgfpathlineto{\pgfqpoint{2.693710in}{2.303754in}}%
\pgfpathlineto{\pgfqpoint{2.694421in}{2.302056in}}%
\pgfpathlineto{\pgfqpoint{2.695536in}{2.297491in}}%
\pgfpathlineto{\pgfqpoint{2.695589in}{2.297948in}}%
\pgfpathlineto{\pgfqpoint{2.696651in}{2.294872in}}%
\pgfpathlineto{\pgfqpoint{2.695618in}{2.298042in}}%
\pgfpathlineto{\pgfqpoint{2.696816in}{2.296370in}}%
\pgfpathlineto{\pgfqpoint{2.697775in}{2.301136in}}%
\pgfpathlineto{\pgfqpoint{2.697975in}{2.300057in}}%
\pgfpathlineto{\pgfqpoint{2.698924in}{2.299812in}}%
\pgfpathlineto{\pgfqpoint{2.698481in}{2.302067in}}%
\pgfpathlineto{\pgfqpoint{2.699026in}{2.301241in}}%
\pgfpathlineto{\pgfqpoint{2.699046in}{2.302229in}}%
\pgfpathlineto{\pgfqpoint{2.699932in}{2.299104in}}%
\pgfpathlineto{\pgfqpoint{2.700784in}{2.295573in}}%
\pgfpathlineto{\pgfqpoint{2.700166in}{2.299941in}}%
\pgfpathlineto{\pgfqpoint{2.701120in}{2.297286in}}%
\pgfpathlineto{\pgfqpoint{2.701208in}{2.298790in}}%
\pgfpathlineto{\pgfqpoint{2.701680in}{2.295869in}}%
\pgfpathlineto{\pgfqpoint{2.702225in}{2.296825in}}%
\pgfpathlineto{\pgfqpoint{2.702585in}{2.294541in}}%
\pgfpathlineto{\pgfqpoint{2.703247in}{2.299000in}}%
\pgfpathlineto{\pgfqpoint{2.703281in}{2.298783in}}%
\pgfpathlineto{\pgfqpoint{2.703296in}{2.298998in}}%
\pgfpathlineto{\pgfqpoint{2.703919in}{2.296385in}}%
\pgfpathlineto{\pgfqpoint{2.704275in}{2.296971in}}%
\pgfpathlineto{\pgfqpoint{2.704528in}{2.294120in}}%
\pgfpathlineto{\pgfqpoint{2.705317in}{2.299733in}}%
\pgfpathlineto{\pgfqpoint{2.705321in}{2.299519in}}%
\pgfpathlineto{\pgfqpoint{2.706465in}{2.304668in}}%
\pgfpathlineto{\pgfqpoint{2.705448in}{2.298924in}}%
\pgfpathlineto{\pgfqpoint{2.706597in}{2.303274in}}%
\pgfpathlineto{\pgfqpoint{2.706641in}{2.302513in}}%
\pgfpathlineto{\pgfqpoint{2.707522in}{2.306238in}}%
\pgfpathlineto{\pgfqpoint{2.707673in}{2.305410in}}%
\pgfpathlineto{\pgfqpoint{2.708092in}{2.308884in}}%
\pgfpathlineto{\pgfqpoint{2.707692in}{2.305206in}}%
\pgfpathlineto{\pgfqpoint{2.708924in}{2.307992in}}%
\pgfpathlineto{\pgfqpoint{2.708929in}{2.308057in}}%
\pgfpathlineto{\pgfqpoint{2.709460in}{2.304194in}}%
\pgfpathlineto{\pgfqpoint{2.709757in}{2.304351in}}%
\pgfpathlineto{\pgfqpoint{2.710910in}{2.301946in}}%
\pgfpathlineto{\pgfqpoint{2.710355in}{2.305976in}}%
\pgfpathlineto{\pgfqpoint{2.710930in}{2.302035in}}%
\pgfpathlineto{\pgfqpoint{2.711577in}{2.303390in}}%
\pgfpathlineto{\pgfqpoint{2.711339in}{2.300628in}}%
\pgfpathlineto{\pgfqpoint{2.712001in}{2.301079in}}%
\pgfpathlineto{\pgfqpoint{2.712016in}{2.301257in}}%
\pgfpathlineto{\pgfqpoint{2.713169in}{2.297837in}}%
\pgfpathlineto{\pgfqpoint{2.713398in}{2.296628in}}%
\pgfpathlineto{\pgfqpoint{2.713676in}{2.300724in}}%
\pgfpathlineto{\pgfqpoint{2.713827in}{2.300726in}}%
\pgfpathlineto{\pgfqpoint{2.714825in}{2.302476in}}%
\pgfpathlineto{\pgfqpoint{2.714065in}{2.298262in}}%
\pgfpathlineto{\pgfqpoint{2.714971in}{2.301997in}}%
\pgfpathlineto{\pgfqpoint{2.716504in}{2.306129in}}%
\pgfpathlineto{\pgfqpoint{2.714995in}{2.301778in}}%
\pgfpathlineto{\pgfqpoint{2.716568in}{2.304397in}}%
\pgfpathlineto{\pgfqpoint{2.717327in}{2.301657in}}%
\pgfpathlineto{\pgfqpoint{2.717020in}{2.304709in}}%
\pgfpathlineto{\pgfqpoint{2.717712in}{2.303041in}}%
\pgfpathlineto{\pgfqpoint{2.717829in}{2.305375in}}%
\pgfpathlineto{\pgfqpoint{2.718184in}{2.302899in}}%
\pgfpathlineto{\pgfqpoint{2.718822in}{2.303421in}}%
\pgfpathlineto{\pgfqpoint{2.719221in}{2.302778in}}%
\pgfpathlineto{\pgfqpoint{2.719533in}{2.305235in}}%
\pgfpathlineto{\pgfqpoint{2.719927in}{2.303532in}}%
\pgfpathlineto{\pgfqpoint{2.720034in}{2.304329in}}%
\pgfpathlineto{\pgfqpoint{2.720492in}{2.301411in}}%
\pgfpathlineto{\pgfqpoint{2.721037in}{2.303516in}}%
\pgfpathlineto{\pgfqpoint{2.721178in}{2.302874in}}%
\pgfpathlineto{\pgfqpoint{2.722006in}{2.305266in}}%
\pgfpathlineto{\pgfqpoint{2.722147in}{2.303255in}}%
\pgfpathlineto{\pgfqpoint{2.723257in}{2.304176in}}%
\pgfpathlineto{\pgfqpoint{2.722984in}{2.301209in}}%
\pgfpathlineto{\pgfqpoint{2.723267in}{2.304023in}}%
\pgfpathlineto{\pgfqpoint{2.723890in}{2.302099in}}%
\pgfpathlineto{\pgfqpoint{2.724265in}{2.304064in}}%
\pgfpathlineto{\pgfqpoint{2.724386in}{2.303550in}}%
\pgfpathlineto{\pgfqpoint{2.724498in}{2.304957in}}%
\pgfpathlineto{\pgfqpoint{2.725428in}{2.300183in}}%
\pgfpathlineto{\pgfqpoint{2.725453in}{2.300407in}}%
\pgfpathlineto{\pgfqpoint{2.726051in}{2.302800in}}%
\pgfpathlineto{\pgfqpoint{2.725521in}{2.299824in}}%
\pgfpathlineto{\pgfqpoint{2.726538in}{2.300098in}}%
\pgfpathlineto{\pgfqpoint{2.727527in}{2.299427in}}%
\pgfpathlineto{\pgfqpoint{2.726903in}{2.302554in}}%
\pgfpathlineto{\pgfqpoint{2.727639in}{2.300764in}}%
\pgfpathlineto{\pgfqpoint{2.728510in}{2.303771in}}%
\pgfpathlineto{\pgfqpoint{2.727936in}{2.299054in}}%
\pgfpathlineto{\pgfqpoint{2.728778in}{2.302773in}}%
\pgfpathlineto{\pgfqpoint{2.729902in}{2.297494in}}%
\pgfpathlineto{\pgfqpoint{2.729055in}{2.303262in}}%
\pgfpathlineto{\pgfqpoint{2.729975in}{2.298467in}}%
\pgfpathlineto{\pgfqpoint{2.730457in}{2.295793in}}%
\pgfpathlineto{\pgfqpoint{2.730905in}{2.299358in}}%
\pgfpathlineto{\pgfqpoint{2.730910in}{2.299249in}}%
\pgfpathlineto{\pgfqpoint{2.731558in}{2.303688in}}%
\pgfpathlineto{\pgfqpoint{2.732059in}{2.301913in}}%
\pgfpathlineto{\pgfqpoint{2.732668in}{2.304409in}}%
\pgfpathlineto{\pgfqpoint{2.732965in}{2.301003in}}%
\pgfpathlineto{\pgfqpoint{2.733125in}{2.301574in}}%
\pgfpathlineto{\pgfqpoint{2.733627in}{2.305270in}}%
\pgfpathlineto{\pgfqpoint{2.734260in}{2.300156in}}%
\pgfpathlineto{\pgfqpoint{2.734420in}{2.299531in}}%
\pgfpathlineto{\pgfqpoint{2.734717in}{2.303576in}}%
\pgfpathlineto{\pgfqpoint{2.735185in}{2.302524in}}%
\pgfpathlineto{\pgfqpoint{2.735195in}{2.303140in}}%
\pgfpathlineto{\pgfqpoint{2.735988in}{2.300797in}}%
\pgfpathlineto{\pgfqpoint{2.736275in}{2.301430in}}%
\pgfpathlineto{\pgfqpoint{2.737419in}{2.298603in}}%
\pgfpathlineto{\pgfqpoint{2.736334in}{2.302024in}}%
\pgfpathlineto{\pgfqpoint{2.737444in}{2.298809in}}%
\pgfpathlineto{\pgfqpoint{2.738714in}{2.304307in}}%
\pgfpathlineto{\pgfqpoint{2.737570in}{2.298029in}}%
\pgfpathlineto{\pgfqpoint{2.738875in}{2.303864in}}%
\pgfpathlineto{\pgfqpoint{2.739566in}{2.300000in}}%
\pgfpathlineto{\pgfqpoint{2.738885in}{2.303901in}}%
\pgfpathlineto{\pgfqpoint{2.740039in}{2.302244in}}%
\pgfpathlineto{\pgfqpoint{2.740141in}{2.302940in}}%
\pgfpathlineto{\pgfqpoint{2.740199in}{2.303646in}}%
\pgfpathlineto{\pgfqpoint{2.741168in}{2.298747in}}%
\pgfpathlineto{\pgfqpoint{2.741222in}{2.300287in}}%
\pgfpathlineto{\pgfqpoint{2.741582in}{2.297035in}}%
\pgfpathlineto{\pgfqpoint{2.742249in}{2.300762in}}%
\pgfpathlineto{\pgfqpoint{2.742332in}{2.300225in}}%
\pgfpathlineto{\pgfqpoint{2.742351in}{2.299879in}}%
\pgfpathlineto{\pgfqpoint{2.742458in}{2.301662in}}%
\pgfpathlineto{\pgfqpoint{2.742497in}{2.301616in}}%
\pgfpathlineto{\pgfqpoint{2.743150in}{2.303855in}}%
\pgfpathlineto{\pgfqpoint{2.743515in}{2.300776in}}%
\pgfpathlineto{\pgfqpoint{2.743593in}{2.300969in}}%
\pgfpathlineto{\pgfqpoint{2.744751in}{2.297954in}}%
\pgfpathlineto{\pgfqpoint{2.743700in}{2.302068in}}%
\pgfpathlineto{\pgfqpoint{2.744756in}{2.298198in}}%
\pgfpathlineto{\pgfqpoint{2.744776in}{2.298726in}}%
\pgfpathlineto{\pgfqpoint{2.745574in}{2.295071in}}%
\pgfpathlineto{\pgfqpoint{2.745784in}{2.295706in}}%
\pgfpathlineto{\pgfqpoint{2.746256in}{2.296765in}}%
\pgfpathlineto{\pgfqpoint{2.746562in}{2.293881in}}%
\pgfpathlineto{\pgfqpoint{2.746786in}{2.294697in}}%
\pgfpathlineto{\pgfqpoint{2.747497in}{2.296154in}}%
\pgfpathlineto{\pgfqpoint{2.747098in}{2.293882in}}%
\pgfpathlineto{\pgfqpoint{2.747901in}{2.294675in}}%
\pgfpathlineto{\pgfqpoint{2.747911in}{2.294428in}}%
\pgfpathlineto{\pgfqpoint{2.748427in}{2.297726in}}%
\pgfpathlineto{\pgfqpoint{2.748958in}{2.296913in}}%
\pgfpathlineto{\pgfqpoint{2.749922in}{2.302266in}}%
\pgfpathlineto{\pgfqpoint{2.750375in}{2.300790in}}%
\pgfpathlineto{\pgfqpoint{2.750404in}{2.300293in}}%
\pgfpathlineto{\pgfqpoint{2.751275in}{2.303678in}}%
\pgfpathlineto{\pgfqpoint{2.751450in}{2.301670in}}%
\pgfpathlineto{\pgfqpoint{2.752526in}{2.303862in}}%
\pgfpathlineto{\pgfqpoint{2.751674in}{2.301145in}}%
\pgfpathlineto{\pgfqpoint{2.752638in}{2.303703in}}%
\pgfpathlineto{\pgfqpoint{2.753573in}{2.298916in}}%
\pgfpathlineto{\pgfqpoint{2.753802in}{2.299842in}}%
\pgfpathlineto{\pgfqpoint{2.754036in}{2.298759in}}%
\pgfpathlineto{\pgfqpoint{2.754966in}{2.303643in}}%
\pgfpathlineto{\pgfqpoint{2.756056in}{2.302271in}}%
\pgfpathlineto{\pgfqpoint{2.755482in}{2.304598in}}%
\pgfpathlineto{\pgfqpoint{2.756100in}{2.303209in}}%
\pgfpathlineto{\pgfqpoint{2.756192in}{2.304570in}}%
\pgfpathlineto{\pgfqpoint{2.756772in}{2.301457in}}%
\pgfpathlineto{\pgfqpoint{2.757239in}{2.303792in}}%
\pgfpathlineto{\pgfqpoint{2.757668in}{2.304716in}}%
\pgfpathlineto{\pgfqpoint{2.758417in}{2.301129in}}%
\pgfpathlineto{\pgfqpoint{2.759595in}{2.303053in}}%
\pgfpathlineto{\pgfqpoint{2.759006in}{2.298983in}}%
\pgfpathlineto{\pgfqpoint{2.759605in}{2.302912in}}%
\pgfpathlineto{\pgfqpoint{2.759815in}{2.301452in}}%
\pgfpathlineto{\pgfqpoint{2.760350in}{2.304663in}}%
\pgfpathlineto{\pgfqpoint{2.760701in}{2.303512in}}%
\pgfpathlineto{\pgfqpoint{2.760735in}{2.303346in}}%
\pgfpathlineto{\pgfqpoint{2.760861in}{2.305767in}}%
\pgfpathlineto{\pgfqpoint{2.761222in}{2.305244in}}%
\pgfpathlineto{\pgfqpoint{2.761557in}{2.306441in}}%
\pgfpathlineto{\pgfqpoint{2.762035in}{2.303307in}}%
\pgfpathlineto{\pgfqpoint{2.762336in}{2.305666in}}%
\pgfpathlineto{\pgfqpoint{2.762531in}{2.307206in}}%
\pgfpathlineto{\pgfqpoint{2.763101in}{2.303128in}}%
\pgfpathlineto{\pgfqpoint{2.763208in}{2.303323in}}%
\pgfpathlineto{\pgfqpoint{2.764084in}{2.302442in}}%
\pgfpathlineto{\pgfqpoint{2.763461in}{2.304248in}}%
\pgfpathlineto{\pgfqpoint{2.764318in}{2.303140in}}%
\pgfpathlineto{\pgfqpoint{2.764547in}{2.305199in}}%
\pgfpathlineto{\pgfqpoint{2.764917in}{2.302030in}}%
\pgfpathlineto{\pgfqpoint{2.765438in}{2.303707in}}%
\pgfpathlineto{\pgfqpoint{2.766557in}{2.302906in}}%
\pgfpathlineto{\pgfqpoint{2.766236in}{2.304520in}}%
\pgfpathlineto{\pgfqpoint{2.766562in}{2.303074in}}%
\pgfpathlineto{\pgfqpoint{2.767351in}{2.304202in}}%
\pgfpathlineto{\pgfqpoint{2.767010in}{2.301985in}}%
\pgfpathlineto{\pgfqpoint{2.767658in}{2.302509in}}%
\pgfpathlineto{\pgfqpoint{2.767818in}{2.303439in}}%
\pgfpathlineto{\pgfqpoint{2.767911in}{2.301501in}}%
\pgfpathlineto{\pgfqpoint{2.768997in}{2.297441in}}%
\pgfpathlineto{\pgfqpoint{2.767940in}{2.302032in}}%
\pgfpathlineto{\pgfqpoint{2.769074in}{2.298304in}}%
\pgfpathlineto{\pgfqpoint{2.769425in}{2.297948in}}%
\pgfpathlineto{\pgfqpoint{2.769907in}{2.300654in}}%
\pgfpathlineto{\pgfqpoint{2.770725in}{2.301099in}}%
\pgfpathlineto{\pgfqpoint{2.770301in}{2.297909in}}%
\pgfpathlineto{\pgfqpoint{2.771007in}{2.300189in}}%
\pgfpathlineto{\pgfqpoint{2.771241in}{2.299232in}}%
\pgfpathlineto{\pgfqpoint{2.771971in}{2.302262in}}%
\pgfpathlineto{\pgfqpoint{2.772098in}{2.301368in}}%
\pgfpathlineto{\pgfqpoint{2.772176in}{2.302016in}}%
\pgfpathlineto{\pgfqpoint{2.772930in}{2.299387in}}%
\pgfpathlineto{\pgfqpoint{2.772998in}{2.299511in}}%
\pgfpathlineto{\pgfqpoint{2.773529in}{2.296167in}}%
\pgfpathlineto{\pgfqpoint{2.774113in}{2.299165in}}%
\pgfpathlineto{\pgfqpoint{2.775530in}{2.304238in}}%
\pgfpathlineto{\pgfqpoint{2.774289in}{2.297981in}}%
\pgfpathlineto{\pgfqpoint{2.775778in}{2.303549in}}%
\pgfpathlineto{\pgfqpoint{2.776134in}{2.301209in}}%
\pgfpathlineto{\pgfqpoint{2.775793in}{2.303716in}}%
\pgfpathlineto{\pgfqpoint{2.776888in}{2.303436in}}%
\pgfpathlineto{\pgfqpoint{2.777697in}{2.304495in}}%
\pgfpathlineto{\pgfqpoint{2.777142in}{2.301425in}}%
\pgfpathlineto{\pgfqpoint{2.777989in}{2.302974in}}%
\pgfpathlineto{\pgfqpoint{2.778948in}{2.299642in}}%
\pgfpathlineto{\pgfqpoint{2.778144in}{2.304650in}}%
\pgfpathlineto{\pgfqpoint{2.779303in}{2.300418in}}%
\pgfpathlineto{\pgfqpoint{2.780481in}{2.304299in}}%
\pgfpathlineto{\pgfqpoint{2.780262in}{2.299731in}}%
\pgfpathlineto{\pgfqpoint{2.780501in}{2.303870in}}%
\pgfpathlineto{\pgfqpoint{2.780715in}{2.302243in}}%
\pgfpathlineto{\pgfqpoint{2.781440in}{2.306515in}}%
\pgfpathlineto{\pgfqpoint{2.781582in}{2.305436in}}%
\pgfpathlineto{\pgfqpoint{2.781640in}{2.306045in}}%
\pgfpathlineto{\pgfqpoint{2.782375in}{2.301555in}}%
\pgfpathlineto{\pgfqpoint{2.782585in}{2.301751in}}%
\pgfpathlineto{\pgfqpoint{2.782599in}{2.300906in}}%
\pgfpathlineto{\pgfqpoint{2.783086in}{2.304115in}}%
\pgfpathlineto{\pgfqpoint{2.783680in}{2.302902in}}%
\pgfpathlineto{\pgfqpoint{2.784556in}{2.303655in}}%
\pgfpathlineto{\pgfqpoint{2.783982in}{2.301425in}}%
\pgfpathlineto{\pgfqpoint{2.784770in}{2.302521in}}%
\pgfpathlineto{\pgfqpoint{2.785695in}{2.304797in}}%
\pgfpathlineto{\pgfqpoint{2.784902in}{2.302249in}}%
\pgfpathlineto{\pgfqpoint{2.786041in}{2.303089in}}%
\pgfpathlineto{\pgfqpoint{2.786883in}{2.298597in}}%
\pgfpathlineto{\pgfqpoint{2.787205in}{2.300693in}}%
\pgfpathlineto{\pgfqpoint{2.788174in}{2.303846in}}%
\pgfpathlineto{\pgfqpoint{2.787346in}{2.300119in}}%
\pgfpathlineto{\pgfqpoint{2.788422in}{2.303557in}}%
\pgfpathlineto{\pgfqpoint{2.788938in}{2.302309in}}%
\pgfpathlineto{\pgfqpoint{2.789503in}{2.305010in}}%
\pgfpathlineto{\pgfqpoint{2.789517in}{2.304597in}}%
\pgfpathlineto{\pgfqpoint{2.790968in}{2.298567in}}%
\pgfpathlineto{\pgfqpoint{2.790194in}{2.305338in}}%
\pgfpathlineto{\pgfqpoint{2.791002in}{2.298774in}}%
\pgfpathlineto{\pgfqpoint{2.792156in}{2.303613in}}%
\pgfpathlineto{\pgfqpoint{2.791319in}{2.298267in}}%
\pgfpathlineto{\pgfqpoint{2.792287in}{2.302338in}}%
\pgfpathlineto{\pgfqpoint{2.792292in}{2.302056in}}%
\pgfpathlineto{\pgfqpoint{2.793339in}{2.304759in}}%
\pgfpathlineto{\pgfqpoint{2.793363in}{2.303934in}}%
\pgfpathlineto{\pgfqpoint{2.793670in}{2.304751in}}%
\pgfpathlineto{\pgfqpoint{2.794425in}{2.302582in}}%
\pgfpathlineto{\pgfqpoint{2.794454in}{2.303149in}}%
\pgfpathlineto{\pgfqpoint{2.794464in}{2.302886in}}%
\pgfpathlineto{\pgfqpoint{2.794824in}{2.305927in}}%
\pgfpathlineto{\pgfqpoint{2.795520in}{2.303630in}}%
\pgfpathlineto{\pgfqpoint{2.795564in}{2.304044in}}%
\pgfpathlineto{\pgfqpoint{2.796304in}{2.299635in}}%
\pgfpathlineto{\pgfqpoint{2.796630in}{2.303679in}}%
\pgfpathlineto{\pgfqpoint{2.797404in}{2.302577in}}%
\pgfpathlineto{\pgfqpoint{2.797156in}{2.305469in}}%
\pgfpathlineto{\pgfqpoint{2.797745in}{2.303415in}}%
\pgfpathlineto{\pgfqpoint{2.797813in}{2.303926in}}%
\pgfpathlineto{\pgfqpoint{2.798617in}{2.300388in}}%
\pgfpathlineto{\pgfqpoint{2.798792in}{2.301667in}}%
\pgfpathlineto{\pgfqpoint{2.798957in}{2.299976in}}%
\pgfpathlineto{\pgfqpoint{2.799843in}{2.304427in}}%
\pgfpathlineto{\pgfqpoint{2.799848in}{2.304267in}}%
\pgfpathlineto{\pgfqpoint{2.800350in}{2.301789in}}%
\pgfpathlineto{\pgfqpoint{2.800637in}{2.304651in}}%
\pgfpathlineto{\pgfqpoint{2.801085in}{2.303218in}}%
\pgfpathlineto{\pgfqpoint{2.801289in}{2.304468in}}%
\pgfpathlineto{\pgfqpoint{2.802102in}{2.299294in}}%
\pgfpathlineto{\pgfqpoint{2.802385in}{2.298075in}}%
\pgfpathlineto{\pgfqpoint{2.802599in}{2.300259in}}%
\pgfpathlineto{\pgfqpoint{2.802891in}{2.300077in}}%
\pgfpathlineto{\pgfqpoint{2.803626in}{2.302779in}}%
\pgfpathlineto{\pgfqpoint{2.802940in}{2.299558in}}%
\pgfpathlineto{\pgfqpoint{2.804025in}{2.302214in}}%
\pgfpathlineto{\pgfqpoint{2.804605in}{2.302036in}}%
\pgfpathlineto{\pgfqpoint{2.804459in}{2.304667in}}%
\pgfpathlineto{\pgfqpoint{2.804804in}{2.303905in}}%
\pgfpathlineto{\pgfqpoint{2.804858in}{2.304467in}}%
\pgfpathlineto{\pgfqpoint{2.805588in}{2.302568in}}%
\pgfpathlineto{\pgfqpoint{2.805900in}{2.302936in}}%
\pgfpathlineto{\pgfqpoint{2.805963in}{2.302590in}}%
\pgfpathlineto{\pgfqpoint{2.806318in}{2.304852in}}%
\pgfpathlineto{\pgfqpoint{2.806976in}{2.304055in}}%
\pgfpathlineto{\pgfqpoint{2.807209in}{2.306775in}}%
\pgfpathlineto{\pgfqpoint{2.807857in}{2.302571in}}%
\pgfpathlineto{\pgfqpoint{2.808061in}{2.303660in}}%
\pgfpathlineto{\pgfqpoint{2.808383in}{2.303189in}}%
\pgfpathlineto{\pgfqpoint{2.808835in}{2.306187in}}%
\pgfpathlineto{\pgfqpoint{2.809128in}{2.304918in}}%
\pgfpathlineto{\pgfqpoint{2.809756in}{2.302302in}}%
\pgfpathlineto{\pgfqpoint{2.809288in}{2.305976in}}%
\pgfpathlineto{\pgfqpoint{2.810325in}{2.304119in}}%
\pgfpathlineto{\pgfqpoint{2.810330in}{2.304329in}}%
\pgfpathlineto{\pgfqpoint{2.811085in}{2.302346in}}%
\pgfpathlineto{\pgfqpoint{2.811401in}{2.303106in}}%
\pgfpathlineto{\pgfqpoint{2.812579in}{2.298636in}}%
\pgfpathlineto{\pgfqpoint{2.812988in}{2.303160in}}%
\pgfpathlineto{\pgfqpoint{2.812609in}{2.298481in}}%
\pgfpathlineto{\pgfqpoint{2.813826in}{2.299856in}}%
\pgfpathlineto{\pgfqpoint{2.813840in}{2.299550in}}%
\pgfpathlineto{\pgfqpoint{2.814736in}{2.304187in}}%
\pgfpathlineto{\pgfqpoint{2.814756in}{2.304169in}}%
\pgfpathlineto{\pgfqpoint{2.814785in}{2.304780in}}%
\pgfpathlineto{\pgfqpoint{2.815724in}{2.300453in}}%
\pgfpathlineto{\pgfqpoint{2.815744in}{2.300760in}}%
\pgfpathlineto{\pgfqpoint{2.815841in}{2.300078in}}%
\pgfpathlineto{\pgfqpoint{2.816343in}{2.304447in}}%
\pgfpathlineto{\pgfqpoint{2.816635in}{2.306033in}}%
\pgfpathlineto{\pgfqpoint{2.817131in}{2.302304in}}%
\pgfpathlineto{\pgfqpoint{2.818358in}{2.296379in}}%
\pgfpathlineto{\pgfqpoint{2.817331in}{2.302378in}}%
\pgfpathlineto{\pgfqpoint{2.818363in}{2.296768in}}%
\pgfpathlineto{\pgfqpoint{2.819196in}{2.293060in}}%
\pgfpathlineto{\pgfqpoint{2.818441in}{2.297463in}}%
\pgfpathlineto{\pgfqpoint{2.819566in}{2.296291in}}%
\pgfpathlineto{\pgfqpoint{2.819644in}{2.297556in}}%
\pgfpathlineto{\pgfqpoint{2.820247in}{2.292743in}}%
\pgfpathlineto{\pgfqpoint{2.820637in}{2.294839in}}%
\pgfpathlineto{\pgfqpoint{2.820919in}{2.293008in}}%
\pgfpathlineto{\pgfqpoint{2.821494in}{2.295934in}}%
\pgfpathlineto{\pgfqpoint{2.821771in}{2.293817in}}%
\pgfpathlineto{\pgfqpoint{2.821883in}{2.294930in}}%
\pgfpathlineto{\pgfqpoint{2.822798in}{2.289583in}}%
\pgfpathlineto{\pgfqpoint{2.823952in}{2.286354in}}%
\pgfpathlineto{\pgfqpoint{2.822891in}{2.290752in}}%
\pgfpathlineto{\pgfqpoint{2.824025in}{2.287600in}}%
\pgfpathlineto{\pgfqpoint{2.824624in}{2.289885in}}%
\pgfpathlineto{\pgfqpoint{2.825048in}{2.285894in}}%
\pgfpathlineto{\pgfqpoint{2.825106in}{2.286464in}}%
\pgfpathlineto{\pgfqpoint{2.825242in}{2.285837in}}%
\pgfpathlineto{\pgfqpoint{2.825700in}{2.288659in}}%
\pgfpathlineto{\pgfqpoint{2.826197in}{2.287574in}}%
\pgfpathlineto{\pgfqpoint{2.826469in}{2.288247in}}%
\pgfpathlineto{\pgfqpoint{2.827199in}{2.286504in}}%
\pgfpathlineto{\pgfqpoint{2.827316in}{2.287843in}}%
\pgfpathlineto{\pgfqpoint{2.827341in}{2.287448in}}%
\pgfpathlineto{\pgfqpoint{2.828120in}{2.289798in}}%
\pgfpathlineto{\pgfqpoint{2.828426in}{2.287807in}}%
\pgfpathlineto{\pgfqpoint{2.829770in}{2.281015in}}%
\pgfpathlineto{\pgfqpoint{2.828456in}{2.288337in}}%
\pgfpathlineto{\pgfqpoint{2.830096in}{2.281920in}}%
\pgfpathlineto{\pgfqpoint{2.830340in}{2.282854in}}%
\pgfpathlineto{\pgfqpoint{2.830904in}{2.278686in}}%
\pgfpathlineto{\pgfqpoint{2.831123in}{2.280545in}}%
\pgfpathlineto{\pgfqpoint{2.831216in}{2.278837in}}%
\pgfpathlineto{\pgfqpoint{2.831956in}{2.284016in}}%
\pgfpathlineto{\pgfqpoint{2.832214in}{2.282100in}}%
\pgfpathlineto{\pgfqpoint{2.833382in}{2.280196in}}%
\pgfpathlineto{\pgfqpoint{2.832944in}{2.284128in}}%
\pgfpathlineto{\pgfqpoint{2.833397in}{2.280692in}}%
\pgfpathlineto{\pgfqpoint{2.834176in}{2.286230in}}%
\pgfpathlineto{\pgfqpoint{2.833421in}{2.280399in}}%
\pgfpathlineto{\pgfqpoint{2.834551in}{2.283080in}}%
\pgfpathlineto{\pgfqpoint{2.834999in}{2.281551in}}%
\pgfpathlineto{\pgfqpoint{2.834629in}{2.284085in}}%
\pgfpathlineto{\pgfqpoint{2.835666in}{2.282945in}}%
\pgfpathlineto{\pgfqpoint{2.836717in}{2.288405in}}%
\pgfpathlineto{\pgfqpoint{2.835768in}{2.282066in}}%
\pgfpathlineto{\pgfqpoint{2.836873in}{2.287690in}}%
\pgfpathlineto{\pgfqpoint{2.836941in}{2.287225in}}%
\pgfpathlineto{\pgfqpoint{2.837370in}{2.290978in}}%
\pgfpathlineto{\pgfqpoint{2.837959in}{2.289361in}}%
\pgfpathlineto{\pgfqpoint{2.838957in}{2.293610in}}%
\pgfpathlineto{\pgfqpoint{2.838110in}{2.288843in}}%
\pgfpathlineto{\pgfqpoint{2.839230in}{2.292977in}}%
\pgfpathlineto{\pgfqpoint{2.839322in}{2.291813in}}%
\pgfpathlineto{\pgfqpoint{2.839424in}{2.294825in}}%
\pgfpathlineto{\pgfqpoint{2.840296in}{2.294370in}}%
\pgfpathlineto{\pgfqpoint{2.840422in}{2.295474in}}%
\pgfpathlineto{\pgfqpoint{2.841279in}{2.292815in}}%
\pgfpathlineto{\pgfqpoint{2.841381in}{2.291607in}}%
\pgfpathlineto{\pgfqpoint{2.842180in}{2.297852in}}%
\pgfpathlineto{\pgfqpoint{2.842229in}{2.297720in}}%
\pgfpathlineto{\pgfqpoint{2.843149in}{2.298851in}}%
\pgfpathlineto{\pgfqpoint{2.842900in}{2.296717in}}%
\pgfpathlineto{\pgfqpoint{2.843358in}{2.298214in}}%
\pgfpathlineto{\pgfqpoint{2.843392in}{2.296408in}}%
\pgfpathlineto{\pgfqpoint{2.844307in}{2.299532in}}%
\pgfpathlineto{\pgfqpoint{2.844380in}{2.299148in}}%
\pgfpathlineto{\pgfqpoint{2.846318in}{2.304539in}}%
\pgfpathlineto{\pgfqpoint{2.844512in}{2.298582in}}%
\pgfpathlineto{\pgfqpoint{2.846396in}{2.303655in}}%
\pgfpathlineto{\pgfqpoint{2.847019in}{2.301264in}}%
\pgfpathlineto{\pgfqpoint{2.846747in}{2.304666in}}%
\pgfpathlineto{\pgfqpoint{2.847525in}{2.302712in}}%
\pgfpathlineto{\pgfqpoint{2.848353in}{2.301987in}}%
\pgfpathlineto{\pgfqpoint{2.847774in}{2.304001in}}%
\pgfpathlineto{\pgfqpoint{2.848572in}{2.303710in}}%
\pgfpathlineto{\pgfqpoint{2.849010in}{2.304039in}}%
\pgfpathlineto{\pgfqpoint{2.849453in}{2.300784in}}%
\pgfpathlineto{\pgfqpoint{2.849585in}{2.301349in}}%
\pgfpathlineto{\pgfqpoint{2.849619in}{2.300421in}}%
\pgfpathlineto{\pgfqpoint{2.850067in}{2.303845in}}%
\pgfpathlineto{\pgfqpoint{2.850705in}{2.300490in}}%
\pgfpathlineto{\pgfqpoint{2.851274in}{2.303434in}}%
\pgfpathlineto{\pgfqpoint{2.851756in}{2.298633in}}%
\pgfpathlineto{\pgfqpoint{2.851781in}{2.298787in}}%
\pgfpathlineto{\pgfqpoint{2.851795in}{2.298729in}}%
\pgfpathlineto{\pgfqpoint{2.852472in}{2.302034in}}%
\pgfpathlineto{\pgfqpoint{2.852477in}{2.302010in}}%
\pgfpathlineto{\pgfqpoint{2.852681in}{2.303851in}}%
\pgfpathlineto{\pgfqpoint{2.853470in}{2.301046in}}%
\pgfpathlineto{\pgfqpoint{2.853567in}{2.301454in}}%
\pgfpathlineto{\pgfqpoint{2.854001in}{2.300559in}}%
\pgfpathlineto{\pgfqpoint{2.854434in}{2.304094in}}%
\pgfpathlineto{\pgfqpoint{2.854658in}{2.303279in}}%
\pgfpathlineto{\pgfqpoint{2.855364in}{2.304055in}}%
\pgfpathlineto{\pgfqpoint{2.855008in}{2.301254in}}%
\pgfpathlineto{\pgfqpoint{2.855724in}{2.302379in}}%
\pgfpathlineto{\pgfqpoint{2.855739in}{2.302043in}}%
\pgfpathlineto{\pgfqpoint{2.856688in}{2.304147in}}%
\pgfpathlineto{\pgfqpoint{2.856805in}{2.303269in}}%
\pgfpathlineto{\pgfqpoint{2.857253in}{2.303857in}}%
\pgfpathlineto{\pgfqpoint{2.857058in}{2.301845in}}%
\pgfpathlineto{\pgfqpoint{2.857910in}{2.302979in}}%
\pgfpathlineto{\pgfqpoint{2.858168in}{2.304496in}}%
\pgfpathlineto{\pgfqpoint{2.858811in}{2.300433in}}%
\pgfpathlineto{\pgfqpoint{2.859035in}{2.303984in}}%
\pgfpathlineto{\pgfqpoint{2.859302in}{2.302050in}}%
\pgfpathlineto{\pgfqpoint{2.859507in}{2.305887in}}%
\pgfpathlineto{\pgfqpoint{2.860154in}{2.303356in}}%
\pgfpathlineto{\pgfqpoint{2.860373in}{2.302223in}}%
\pgfpathlineto{\pgfqpoint{2.860963in}{2.305091in}}%
\pgfpathlineto{\pgfqpoint{2.861274in}{2.303130in}}%
\pgfpathlineto{\pgfqpoint{2.862277in}{2.304002in}}%
\pgfpathlineto{\pgfqpoint{2.861415in}{2.301709in}}%
\pgfpathlineto{\pgfqpoint{2.862350in}{2.303008in}}%
\pgfpathlineto{\pgfqpoint{2.862355in}{2.302756in}}%
\pgfpathlineto{\pgfqpoint{2.863134in}{2.305308in}}%
\pgfpathlineto{\pgfqpoint{2.863450in}{2.303497in}}%
\pgfpathlineto{\pgfqpoint{2.864195in}{2.301904in}}%
\pgfpathlineto{\pgfqpoint{2.864473in}{2.304402in}}%
\pgfpathlineto{\pgfqpoint{2.864478in}{2.304366in}}%
\pgfpathlineto{\pgfqpoint{2.865232in}{2.305250in}}%
\pgfpathlineto{\pgfqpoint{2.864814in}{2.302945in}}%
\pgfpathlineto{\pgfqpoint{2.865578in}{2.304253in}}%
\pgfpathlineto{\pgfqpoint{2.866060in}{2.302516in}}%
\pgfpathlineto{\pgfqpoint{2.865617in}{2.304521in}}%
\pgfpathlineto{\pgfqpoint{2.866688in}{2.303999in}}%
\pgfpathlineto{\pgfqpoint{2.867632in}{2.305086in}}%
\pgfpathlineto{\pgfqpoint{2.867223in}{2.301918in}}%
\pgfpathlineto{\pgfqpoint{2.867778in}{2.303667in}}%
\pgfpathlineto{\pgfqpoint{2.868251in}{2.301011in}}%
\pgfpathlineto{\pgfqpoint{2.867798in}{2.303763in}}%
\pgfpathlineto{\pgfqpoint{2.868898in}{2.303344in}}%
\pgfpathlineto{\pgfqpoint{2.869098in}{2.304612in}}%
\pgfpathlineto{\pgfqpoint{2.869950in}{2.302464in}}%
\pgfpathlineto{\pgfqpoint{2.869994in}{2.302983in}}%
\pgfpathlineto{\pgfqpoint{2.870899in}{2.297547in}}%
\pgfpathlineto{\pgfqpoint{2.871225in}{2.299632in}}%
\pgfpathlineto{\pgfqpoint{2.871474in}{2.299906in}}%
\pgfpathlineto{\pgfqpoint{2.872194in}{2.296241in}}%
\pgfpathlineto{\pgfqpoint{2.872262in}{2.296590in}}%
\pgfpathlineto{\pgfqpoint{2.873148in}{2.302169in}}%
\pgfpathlineto{\pgfqpoint{2.872369in}{2.295667in}}%
\pgfpathlineto{\pgfqpoint{2.873528in}{2.298935in}}%
\pgfpathlineto{\pgfqpoint{2.873884in}{2.298089in}}%
\pgfpathlineto{\pgfqpoint{2.874516in}{2.301354in}}%
\pgfpathlineto{\pgfqpoint{2.874609in}{2.300596in}}%
\pgfpathlineto{\pgfqpoint{2.874750in}{2.300253in}}%
\pgfpathlineto{\pgfqpoint{2.875534in}{2.304317in}}%
\pgfpathlineto{\pgfqpoint{2.876668in}{2.302281in}}%
\pgfpathlineto{\pgfqpoint{2.875777in}{2.304912in}}%
\pgfpathlineto{\pgfqpoint{2.876756in}{2.303144in}}%
\pgfpathlineto{\pgfqpoint{2.877024in}{2.304433in}}%
\pgfpathlineto{\pgfqpoint{2.877593in}{2.301957in}}%
\pgfpathlineto{\pgfqpoint{2.877866in}{2.303099in}}%
\pgfpathlineto{\pgfqpoint{2.877885in}{2.303495in}}%
\pgfpathlineto{\pgfqpoint{2.878791in}{2.298356in}}%
\pgfpathlineto{\pgfqpoint{2.878888in}{2.298754in}}%
\pgfpathlineto{\pgfqpoint{2.879945in}{2.297445in}}%
\pgfpathlineto{\pgfqpoint{2.879658in}{2.301098in}}%
\pgfpathlineto{\pgfqpoint{2.880013in}{2.298126in}}%
\pgfpathlineto{\pgfqpoint{2.880748in}{2.300128in}}%
\pgfpathlineto{\pgfqpoint{2.880140in}{2.296056in}}%
\pgfpathlineto{\pgfqpoint{2.881191in}{2.299069in}}%
\pgfpathlineto{\pgfqpoint{2.881678in}{2.296251in}}%
\pgfpathlineto{\pgfqpoint{2.882287in}{2.299647in}}%
\pgfpathlineto{\pgfqpoint{2.882296in}{2.299389in}}%
\pgfpathlineto{\pgfqpoint{2.882910in}{2.302864in}}%
\pgfpathlineto{\pgfqpoint{2.883460in}{2.301863in}}%
\pgfpathlineto{\pgfqpoint{2.883703in}{2.300749in}}%
\pgfpathlineto{\pgfqpoint{2.884414in}{2.304349in}}%
\pgfpathlineto{\pgfqpoint{2.884507in}{2.303358in}}%
\pgfpathlineto{\pgfqpoint{2.885315in}{2.304085in}}%
\pgfpathlineto{\pgfqpoint{2.884789in}{2.300938in}}%
\pgfpathlineto{\pgfqpoint{2.885568in}{2.301671in}}%
\pgfpathlineto{\pgfqpoint{2.886279in}{2.302811in}}%
\pgfpathlineto{\pgfqpoint{2.885782in}{2.299742in}}%
\pgfpathlineto{\pgfqpoint{2.886488in}{2.300680in}}%
\pgfpathlineto{\pgfqpoint{2.886732in}{2.298764in}}%
\pgfpathlineto{\pgfqpoint{2.887184in}{2.302908in}}%
\pgfpathlineto{\pgfqpoint{2.887588in}{2.301499in}}%
\pgfpathlineto{\pgfqpoint{2.887691in}{2.301150in}}%
\pgfpathlineto{\pgfqpoint{2.888221in}{2.304785in}}%
\pgfpathlineto{\pgfqpoint{2.888504in}{2.303539in}}%
\pgfpathlineto{\pgfqpoint{2.889185in}{2.304160in}}%
\pgfpathlineto{\pgfqpoint{2.889521in}{2.302292in}}%
\pgfpathlineto{\pgfqpoint{2.889745in}{2.304180in}}%
\pgfpathlineto{\pgfqpoint{2.890768in}{2.295352in}}%
\pgfpathlineto{\pgfqpoint{2.891118in}{2.295756in}}%
\pgfpathlineto{\pgfqpoint{2.890933in}{2.293271in}}%
\pgfpathlineto{\pgfqpoint{2.891512in}{2.294151in}}%
\pgfpathlineto{\pgfqpoint{2.891537in}{2.293792in}}%
\pgfpathlineto{\pgfqpoint{2.892282in}{2.299253in}}%
\pgfpathlineto{\pgfqpoint{2.892540in}{2.298778in}}%
\pgfpathlineto{\pgfqpoint{2.893942in}{2.302100in}}%
\pgfpathlineto{\pgfqpoint{2.894195in}{2.300778in}}%
\pgfpathlineto{\pgfqpoint{2.894297in}{2.299783in}}%
\pgfpathlineto{\pgfqpoint{2.895222in}{2.303153in}}%
\pgfpathlineto{\pgfqpoint{2.895247in}{2.302995in}}%
\pgfpathlineto{\pgfqpoint{2.896025in}{2.304080in}}%
\pgfpathlineto{\pgfqpoint{2.895767in}{2.301566in}}%
\pgfpathlineto{\pgfqpoint{2.896337in}{2.302686in}}%
\pgfpathlineto{\pgfqpoint{2.896989in}{2.298800in}}%
\pgfpathlineto{\pgfqpoint{2.896430in}{2.302999in}}%
\pgfpathlineto{\pgfqpoint{2.897457in}{2.302147in}}%
\pgfpathlineto{\pgfqpoint{2.898031in}{2.301810in}}%
\pgfpathlineto{\pgfqpoint{2.898129in}{2.303566in}}%
\pgfpathlineto{\pgfqpoint{2.898134in}{2.303501in}}%
\pgfpathlineto{\pgfqpoint{2.898158in}{2.304152in}}%
\pgfpathlineto{\pgfqpoint{2.899136in}{2.302058in}}%
\pgfpathlineto{\pgfqpoint{2.899244in}{2.303502in}}%
\pgfpathlineto{\pgfqpoint{2.899721in}{2.301918in}}%
\pgfpathlineto{\pgfqpoint{2.899429in}{2.304167in}}%
\pgfpathlineto{\pgfqpoint{2.900378in}{2.302509in}}%
\pgfpathlineto{\pgfqpoint{2.901283in}{2.305374in}}%
\pgfpathlineto{\pgfqpoint{2.900626in}{2.301915in}}%
\pgfpathlineto{\pgfqpoint{2.901542in}{2.304739in}}%
\pgfpathlineto{\pgfqpoint{2.901829in}{2.303653in}}%
\pgfpathlineto{\pgfqpoint{2.903002in}{2.299278in}}%
\pgfpathlineto{\pgfqpoint{2.901985in}{2.304374in}}%
\pgfpathlineto{\pgfqpoint{2.903056in}{2.299552in}}%
\pgfpathlineto{\pgfqpoint{2.904370in}{2.296385in}}%
\pgfpathlineto{\pgfqpoint{2.903450in}{2.301236in}}%
\pgfpathlineto{\pgfqpoint{2.904380in}{2.297012in}}%
\pgfpathlineto{\pgfqpoint{2.905504in}{2.299616in}}%
\pgfpathlineto{\pgfqpoint{2.905242in}{2.295196in}}%
\pgfpathlineto{\pgfqpoint{2.905539in}{2.299411in}}%
\pgfpathlineto{\pgfqpoint{2.905582in}{2.298296in}}%
\pgfpathlineto{\pgfqpoint{2.906391in}{2.300783in}}%
\pgfpathlineto{\pgfqpoint{2.906614in}{2.299491in}}%
\pgfpathlineto{\pgfqpoint{2.907286in}{2.304568in}}%
\pgfpathlineto{\pgfqpoint{2.906877in}{2.298965in}}%
\pgfpathlineto{\pgfqpoint{2.907768in}{2.303461in}}%
\pgfpathlineto{\pgfqpoint{2.908036in}{2.302441in}}%
\pgfpathlineto{\pgfqpoint{2.908236in}{2.305278in}}%
\pgfpathlineto{\pgfqpoint{2.908864in}{2.304270in}}%
\pgfpathlineto{\pgfqpoint{2.909940in}{2.302261in}}%
\pgfpathlineto{\pgfqpoint{2.909463in}{2.304488in}}%
\pgfpathlineto{\pgfqpoint{2.910018in}{2.303685in}}%
\pgfpathlineto{\pgfqpoint{2.910193in}{2.303987in}}%
\pgfpathlineto{\pgfqpoint{2.910475in}{2.301958in}}%
\pgfpathlineto{\pgfqpoint{2.911045in}{2.299669in}}%
\pgfpathlineto{\pgfqpoint{2.910704in}{2.303623in}}%
\pgfpathlineto{\pgfqpoint{2.911590in}{2.301795in}}%
\pgfpathlineto{\pgfqpoint{2.912656in}{2.304643in}}%
\pgfpathlineto{\pgfqpoint{2.911882in}{2.300748in}}%
\pgfpathlineto{\pgfqpoint{2.912739in}{2.303212in}}%
\pgfpathlineto{\pgfqpoint{2.912817in}{2.301670in}}%
\pgfpathlineto{\pgfqpoint{2.913060in}{2.304287in}}%
\pgfpathlineto{\pgfqpoint{2.913869in}{2.302430in}}%
\pgfpathlineto{\pgfqpoint{2.914019in}{2.303229in}}%
\pgfpathlineto{\pgfqpoint{2.914609in}{2.298589in}}%
\pgfpathlineto{\pgfqpoint{2.914613in}{2.298723in}}%
\pgfpathlineto{\pgfqpoint{2.914988in}{2.296168in}}%
\pgfpathlineto{\pgfqpoint{2.915733in}{2.298512in}}%
\pgfpathlineto{\pgfqpoint{2.916814in}{2.299341in}}%
\pgfpathlineto{\pgfqpoint{2.916347in}{2.295687in}}%
\pgfpathlineto{\pgfqpoint{2.916838in}{2.298840in}}%
\pgfpathlineto{\pgfqpoint{2.917345in}{2.296727in}}%
\pgfpathlineto{\pgfqpoint{2.917710in}{2.299505in}}%
\pgfpathlineto{\pgfqpoint{2.917948in}{2.298518in}}%
\pgfpathlineto{\pgfqpoint{2.918236in}{2.298826in}}%
\pgfpathlineto{\pgfqpoint{2.918664in}{2.296766in}}%
\pgfpathlineto{\pgfqpoint{2.918980in}{2.297820in}}%
\pgfpathlineto{\pgfqpoint{2.919531in}{2.295380in}}%
\pgfpathlineto{\pgfqpoint{2.919102in}{2.299145in}}%
\pgfpathlineto{\pgfqpoint{2.920086in}{2.297902in}}%
\pgfpathlineto{\pgfqpoint{2.920144in}{2.297486in}}%
\pgfpathlineto{\pgfqpoint{2.920485in}{2.300466in}}%
\pgfpathlineto{\pgfqpoint{2.920577in}{2.301300in}}%
\pgfpathlineto{\pgfqpoint{2.921459in}{2.296737in}}%
\pgfpathlineto{\pgfqpoint{2.921536in}{2.297415in}}%
\pgfpathlineto{\pgfqpoint{2.922126in}{2.295313in}}%
\pgfpathlineto{\pgfqpoint{2.921804in}{2.297784in}}%
\pgfpathlineto{\pgfqpoint{2.922695in}{2.295497in}}%
\pgfpathlineto{\pgfqpoint{2.923786in}{2.298662in}}%
\pgfpathlineto{\pgfqpoint{2.923255in}{2.295485in}}%
\pgfpathlineto{\pgfqpoint{2.923859in}{2.298468in}}%
\pgfpathlineto{\pgfqpoint{2.923873in}{2.297717in}}%
\pgfpathlineto{\pgfqpoint{2.924282in}{2.300951in}}%
\pgfpathlineto{\pgfqpoint{2.924969in}{2.298507in}}%
\pgfpathlineto{\pgfqpoint{2.924974in}{2.298404in}}%
\pgfpathlineto{\pgfqpoint{2.925538in}{2.305462in}}%
\pgfpathlineto{\pgfqpoint{2.925845in}{2.303390in}}%
\pgfpathlineto{\pgfqpoint{2.926960in}{2.304295in}}%
\pgfpathlineto{\pgfqpoint{2.926814in}{2.302545in}}%
\pgfpathlineto{\pgfqpoint{2.926984in}{2.303835in}}%
\pgfpathlineto{\pgfqpoint{2.928080in}{2.301166in}}%
\pgfpathlineto{\pgfqpoint{2.927418in}{2.305621in}}%
\pgfpathlineto{\pgfqpoint{2.928128in}{2.301768in}}%
\pgfpathlineto{\pgfqpoint{2.928810in}{2.303671in}}%
\pgfpathlineto{\pgfqpoint{2.929180in}{2.299425in}}%
\pgfpathlineto{\pgfqpoint{2.929199in}{2.299712in}}%
\pgfpathlineto{\pgfqpoint{2.930047in}{2.302110in}}%
\pgfpathlineto{\pgfqpoint{2.929701in}{2.299428in}}%
\pgfpathlineto{\pgfqpoint{2.930334in}{2.300276in}}%
\pgfpathlineto{\pgfqpoint{2.930533in}{2.298302in}}%
\pgfpathlineto{\pgfqpoint{2.931390in}{2.302056in}}%
\pgfpathlineto{\pgfqpoint{2.931429in}{2.301207in}}%
\pgfpathlineto{\pgfqpoint{2.931439in}{2.301342in}}%
\pgfpathlineto{\pgfqpoint{2.931605in}{2.299862in}}%
\pgfpathlineto{\pgfqpoint{2.931663in}{2.299344in}}%
\pgfpathlineto{\pgfqpoint{2.932228in}{2.304135in}}%
\pgfpathlineto{\pgfqpoint{2.932622in}{2.302091in}}%
\pgfpathlineto{\pgfqpoint{2.932997in}{2.302735in}}%
\pgfpathlineto{\pgfqpoint{2.933372in}{2.300048in}}%
\pgfpathlineto{\pgfqpoint{2.933664in}{2.300274in}}%
\pgfpathlineto{\pgfqpoint{2.933883in}{2.298611in}}%
\pgfpathlineto{\pgfqpoint{2.934438in}{2.301808in}}%
\pgfpathlineto{\pgfqpoint{2.934769in}{2.300683in}}%
\pgfpathlineto{\pgfqpoint{2.934774in}{2.300814in}}%
\pgfpathlineto{\pgfqpoint{2.935699in}{2.296564in}}%
\pgfpathlineto{\pgfqpoint{2.935757in}{2.297381in}}%
\pgfpathlineto{\pgfqpoint{2.935947in}{2.295704in}}%
\pgfpathlineto{\pgfqpoint{2.936390in}{2.301433in}}%
\pgfpathlineto{\pgfqpoint{2.936663in}{2.300785in}}%
\pgfpathlineto{\pgfqpoint{2.936692in}{2.301252in}}%
\pgfpathlineto{\pgfqpoint{2.936799in}{2.299726in}}%
\pgfpathlineto{\pgfqpoint{2.936892in}{2.299684in}}%
\pgfpathlineto{\pgfqpoint{2.937335in}{2.298937in}}%
\pgfpathlineto{\pgfqpoint{2.937817in}{2.301545in}}%
\pgfpathlineto{\pgfqpoint{2.937963in}{2.300949in}}%
\pgfpathlineto{\pgfqpoint{2.938601in}{2.304335in}}%
\pgfpathlineto{\pgfqpoint{2.937987in}{2.300719in}}%
\pgfpathlineto{\pgfqpoint{2.939107in}{2.303447in}}%
\pgfpathlineto{\pgfqpoint{2.939204in}{2.304206in}}%
\pgfpathlineto{\pgfqpoint{2.940066in}{2.301513in}}%
\pgfpathlineto{\pgfqpoint{2.940081in}{2.301987in}}%
\pgfpathlineto{\pgfqpoint{2.940718in}{2.301279in}}%
\pgfpathlineto{\pgfqpoint{2.941152in}{2.304116in}}%
\pgfpathlineto{\pgfqpoint{2.941161in}{2.303981in}}%
\pgfpathlineto{\pgfqpoint{2.941527in}{2.302107in}}%
\pgfpathlineto{\pgfqpoint{2.941239in}{2.304646in}}%
\pgfpathlineto{\pgfqpoint{2.942228in}{2.304372in}}%
\pgfpathlineto{\pgfqpoint{2.942310in}{2.305603in}}%
\pgfpathlineto{\pgfqpoint{2.942651in}{2.302567in}}%
\pgfpathlineto{\pgfqpoint{2.943333in}{2.304209in}}%
\pgfpathlineto{\pgfqpoint{2.944418in}{2.302842in}}%
\pgfpathlineto{\pgfqpoint{2.944141in}{2.305646in}}%
\pgfpathlineto{\pgfqpoint{2.944472in}{2.303698in}}%
\pgfpathlineto{\pgfqpoint{2.944959in}{2.299361in}}%
\pgfpathlineto{\pgfqpoint{2.944501in}{2.304170in}}%
\pgfpathlineto{\pgfqpoint{2.945543in}{2.303714in}}%
\pgfpathlineto{\pgfqpoint{2.945723in}{2.304128in}}%
\pgfpathlineto{\pgfqpoint{2.946619in}{2.301645in}}%
\pgfpathlineto{\pgfqpoint{2.947013in}{2.301238in}}%
\pgfpathlineto{\pgfqpoint{2.947432in}{2.304418in}}%
\pgfpathlineto{\pgfqpoint{2.947690in}{2.302966in}}%
\pgfpathlineto{\pgfqpoint{2.947841in}{2.303888in}}%
\pgfpathlineto{\pgfqpoint{2.948036in}{2.301075in}}%
\pgfpathlineto{\pgfqpoint{2.948162in}{2.301175in}}%
\pgfpathlineto{\pgfqpoint{2.948644in}{2.300106in}}%
\pgfpathlineto{\pgfqpoint{2.949014in}{2.302875in}}%
\pgfpathlineto{\pgfqpoint{2.949297in}{2.300394in}}%
\pgfpathlineto{\pgfqpoint{2.950119in}{2.302709in}}%
\pgfpathlineto{\pgfqpoint{2.949511in}{2.299427in}}%
\pgfpathlineto{\pgfqpoint{2.950426in}{2.301162in}}%
\pgfpathlineto{\pgfqpoint{2.951424in}{2.299471in}}%
\pgfpathlineto{\pgfqpoint{2.951025in}{2.302613in}}%
\pgfpathlineto{\pgfqpoint{2.951565in}{2.300918in}}%
\pgfpathlineto{\pgfqpoint{2.952330in}{2.304032in}}%
\pgfpathlineto{\pgfqpoint{2.951804in}{2.300272in}}%
\pgfpathlineto{\pgfqpoint{2.952714in}{2.302077in}}%
\pgfpathlineto{\pgfqpoint{2.953016in}{2.305331in}}%
\pgfpathlineto{\pgfqpoint{2.953712in}{2.300993in}}%
\pgfpathlineto{\pgfqpoint{2.954063in}{2.298103in}}%
\pgfpathlineto{\pgfqpoint{2.954477in}{2.301727in}}%
\pgfpathlineto{\pgfqpoint{2.954818in}{2.300896in}}%
\pgfpathlineto{\pgfqpoint{2.954861in}{2.301541in}}%
\pgfpathlineto{\pgfqpoint{2.955855in}{2.297131in}}%
\pgfpathlineto{\pgfqpoint{2.955869in}{2.297411in}}%
\pgfpathlineto{\pgfqpoint{2.956736in}{2.293806in}}%
\pgfpathlineto{\pgfqpoint{2.956741in}{2.293375in}}%
\pgfpathlineto{\pgfqpoint{2.957675in}{2.299173in}}%
\pgfpathlineto{\pgfqpoint{2.957787in}{2.298569in}}%
\pgfpathlineto{\pgfqpoint{2.958892in}{2.300823in}}%
\pgfpathlineto{\pgfqpoint{2.958255in}{2.296058in}}%
\pgfpathlineto{\pgfqpoint{2.958946in}{2.300572in}}%
\pgfpathlineto{\pgfqpoint{2.959014in}{2.299542in}}%
\pgfpathlineto{\pgfqpoint{2.959706in}{2.305357in}}%
\pgfpathlineto{\pgfqpoint{2.959891in}{2.304077in}}%
\pgfpathlineto{\pgfqpoint{2.959905in}{2.304398in}}%
\pgfpathlineto{\pgfqpoint{2.960261in}{2.300375in}}%
\pgfpathlineto{\pgfqpoint{2.960879in}{2.301834in}}%
\pgfpathlineto{\pgfqpoint{2.961254in}{2.301112in}}%
\pgfpathlineto{\pgfqpoint{2.961736in}{2.303509in}}%
\pgfpathlineto{\pgfqpoint{2.961921in}{2.302896in}}%
\pgfpathlineto{\pgfqpoint{2.962670in}{2.300201in}}%
\pgfpathlineto{\pgfqpoint{2.961974in}{2.303791in}}%
\pgfpathlineto{\pgfqpoint{2.963036in}{2.303068in}}%
\pgfpathlineto{\pgfqpoint{2.963284in}{2.305257in}}%
\pgfpathlineto{\pgfqpoint{2.963873in}{2.299837in}}%
\pgfpathlineto{\pgfqpoint{2.963946in}{2.300333in}}%
\pgfpathlineto{\pgfqpoint{2.964136in}{2.299034in}}%
\pgfpathlineto{\pgfqpoint{2.964482in}{2.303081in}}%
\pgfpathlineto{\pgfqpoint{2.964550in}{2.302826in}}%
\pgfpathlineto{\pgfqpoint{2.964759in}{2.304954in}}%
\pgfpathlineto{\pgfqpoint{2.965285in}{2.302048in}}%
\pgfpathlineto{\pgfqpoint{2.965665in}{2.303348in}}%
\pgfpathlineto{\pgfqpoint{2.965674in}{2.302896in}}%
\pgfpathlineto{\pgfqpoint{2.965815in}{2.305767in}}%
\pgfpathlineto{\pgfqpoint{2.966760in}{2.303997in}}%
\pgfpathlineto{\pgfqpoint{2.967690in}{2.300449in}}%
\pgfpathlineto{\pgfqpoint{2.967977in}{2.301380in}}%
\pgfpathlineto{\pgfqpoint{2.968099in}{2.301624in}}%
\pgfpathlineto{\pgfqpoint{2.968654in}{2.298481in}}%
\pgfpathlineto{\pgfqpoint{2.969058in}{2.300627in}}%
\pgfpathlineto{\pgfqpoint{2.969116in}{2.300054in}}%
\pgfpathlineto{\pgfqpoint{2.969788in}{2.304078in}}%
\pgfpathlineto{\pgfqpoint{2.970080in}{2.303512in}}%
\pgfpathlineto{\pgfqpoint{2.970265in}{2.303790in}}%
\pgfpathlineto{\pgfqpoint{2.970470in}{2.300731in}}%
\pgfpathlineto{\pgfqpoint{2.970845in}{2.301292in}}%
\pgfpathlineto{\pgfqpoint{2.971166in}{2.299933in}}%
\pgfpathlineto{\pgfqpoint{2.971750in}{2.303001in}}%
\pgfpathlineto{\pgfqpoint{2.971950in}{2.301659in}}%
\pgfpathlineto{\pgfqpoint{2.973118in}{2.298983in}}%
\pgfpathlineto{\pgfqpoint{2.972602in}{2.302404in}}%
\pgfpathlineto{\pgfqpoint{2.973172in}{2.299337in}}%
\pgfpathlineto{\pgfqpoint{2.973610in}{2.296675in}}%
\pgfpathlineto{\pgfqpoint{2.974165in}{2.300014in}}%
\pgfpathlineto{\pgfqpoint{2.974175in}{2.299860in}}%
\pgfpathlineto{\pgfqpoint{2.975182in}{2.303224in}}%
\pgfpathlineto{\pgfqpoint{2.974209in}{2.299450in}}%
\pgfpathlineto{\pgfqpoint{2.975324in}{2.302147in}}%
\pgfpathlineto{\pgfqpoint{2.975957in}{2.298786in}}%
\pgfpathlineto{\pgfqpoint{2.976292in}{2.302682in}}%
\pgfpathlineto{\pgfqpoint{2.976434in}{2.302033in}}%
\pgfpathlineto{\pgfqpoint{2.976531in}{2.304210in}}%
\pgfpathlineto{\pgfqpoint{2.977179in}{2.301420in}}%
\pgfpathlineto{\pgfqpoint{2.977524in}{2.301731in}}%
\pgfpathlineto{\pgfqpoint{2.977602in}{2.300699in}}%
\pgfpathlineto{\pgfqpoint{2.978357in}{2.304162in}}%
\pgfpathlineto{\pgfqpoint{2.978595in}{2.303903in}}%
\pgfpathlineto{\pgfqpoint{2.979442in}{2.305167in}}%
\pgfpathlineto{\pgfqpoint{2.978834in}{2.301622in}}%
\pgfpathlineto{\pgfqpoint{2.979598in}{2.303174in}}%
\pgfpathlineto{\pgfqpoint{2.979905in}{2.304725in}}%
\pgfpathlineto{\pgfqpoint{2.980294in}{2.302547in}}%
\pgfpathlineto{\pgfqpoint{2.980606in}{2.302830in}}%
\pgfpathlineto{\pgfqpoint{2.980879in}{2.300374in}}%
\pgfpathlineto{\pgfqpoint{2.981658in}{2.303785in}}%
\pgfpathlineto{\pgfqpoint{2.981701in}{2.303229in}}%
\pgfpathlineto{\pgfqpoint{2.982729in}{2.305249in}}%
\pgfpathlineto{\pgfqpoint{2.981964in}{2.299693in}}%
\pgfpathlineto{\pgfqpoint{2.982836in}{2.304524in}}%
\pgfpathlineto{\pgfqpoint{2.983624in}{2.302719in}}%
\pgfpathlineto{\pgfqpoint{2.982928in}{2.305220in}}%
\pgfpathlineto{\pgfqpoint{2.983960in}{2.303581in}}%
\pgfpathlineto{\pgfqpoint{2.984179in}{2.303938in}}%
\pgfpathlineto{\pgfqpoint{2.984515in}{2.301217in}}%
\pgfpathlineto{\pgfqpoint{2.984939in}{2.301261in}}%
\pgfpathlineto{\pgfqpoint{2.985397in}{2.300202in}}%
\pgfpathlineto{\pgfqpoint{2.985202in}{2.302657in}}%
\pgfpathlineto{\pgfqpoint{2.985942in}{2.301844in}}%
\pgfpathlineto{\pgfqpoint{2.986930in}{2.303273in}}%
\pgfpathlineto{\pgfqpoint{2.986477in}{2.300178in}}%
\pgfpathlineto{\pgfqpoint{2.987052in}{2.302067in}}%
\pgfpathlineto{\pgfqpoint{2.988147in}{2.296956in}}%
\pgfpathlineto{\pgfqpoint{2.988279in}{2.297241in}}%
\pgfpathlineto{\pgfqpoint{2.989145in}{2.299551in}}%
\pgfpathlineto{\pgfqpoint{2.988288in}{2.296999in}}%
\pgfpathlineto{\pgfqpoint{2.989476in}{2.298006in}}%
\pgfpathlineto{\pgfqpoint{2.989496in}{2.297287in}}%
\pgfpathlineto{\pgfqpoint{2.990533in}{2.302534in}}%
\pgfpathlineto{\pgfqpoint{2.990538in}{2.302471in}}%
\pgfpathlineto{\pgfqpoint{2.990903in}{2.304066in}}%
\pgfpathlineto{\pgfqpoint{2.991560in}{2.301804in}}%
\pgfpathlineto{\pgfqpoint{2.991784in}{2.300391in}}%
\pgfpathlineto{\pgfqpoint{2.992081in}{2.304346in}}%
\pgfpathlineto{\pgfqpoint{2.992656in}{2.302392in}}%
\pgfpathlineto{\pgfqpoint{2.993323in}{2.302306in}}%
\pgfpathlineto{\pgfqpoint{2.993814in}{2.304664in}}%
\pgfpathlineto{\pgfqpoint{2.994900in}{2.301669in}}%
\pgfpathlineto{\pgfqpoint{2.994622in}{2.305014in}}%
\pgfpathlineto{\pgfqpoint{2.994988in}{2.301971in}}%
\pgfpathlineto{\pgfqpoint{2.996102in}{2.305360in}}%
\pgfpathlineto{\pgfqpoint{2.995051in}{2.301880in}}%
\pgfpathlineto{\pgfqpoint{2.996210in}{2.303680in}}%
\pgfpathlineto{\pgfqpoint{2.997105in}{2.299541in}}%
\pgfpathlineto{\pgfqpoint{2.996365in}{2.303775in}}%
\pgfpathlineto{\pgfqpoint{2.997359in}{2.300840in}}%
\pgfpathlineto{\pgfqpoint{2.997777in}{2.304223in}}%
\pgfpathlineto{\pgfqpoint{2.998537in}{2.303342in}}%
\pgfpathlineto{\pgfqpoint{2.999822in}{2.300238in}}%
\pgfpathlineto{\pgfqpoint{2.999199in}{2.304232in}}%
\pgfpathlineto{\pgfqpoint{2.999841in}{2.300506in}}%
\pgfpathlineto{\pgfqpoint{3.000903in}{2.301672in}}%
\pgfpathlineto{\pgfqpoint{3.000659in}{2.299044in}}%
\pgfpathlineto{\pgfqpoint{3.000956in}{2.301267in}}%
\pgfpathlineto{\pgfqpoint{3.001117in}{2.299205in}}%
\pgfpathlineto{\pgfqpoint{3.001794in}{2.302987in}}%
\pgfpathlineto{\pgfqpoint{3.002032in}{2.302657in}}%
\pgfpathlineto{\pgfqpoint{3.002612in}{2.304275in}}%
\pgfpathlineto{\pgfqpoint{3.002996in}{2.301317in}}%
\pgfpathlineto{\pgfqpoint{3.003113in}{2.302203in}}%
\pgfpathlineto{\pgfqpoint{3.003449in}{2.299580in}}%
\pgfpathlineto{\pgfqpoint{3.004208in}{2.302723in}}%
\pgfpathlineto{\pgfqpoint{3.004218in}{2.302643in}}%
\pgfpathlineto{\pgfqpoint{3.004515in}{2.301990in}}%
\pgfpathlineto{\pgfqpoint{3.005056in}{2.304379in}}%
\pgfpathlineto{\pgfqpoint{3.005065in}{2.304219in}}%
\pgfpathlineto{\pgfqpoint{3.005430in}{2.305296in}}%
\pgfpathlineto{\pgfqpoint{3.005947in}{2.302046in}}%
\pgfpathlineto{\pgfqpoint{3.006141in}{2.303104in}}%
\pgfpathlineto{\pgfqpoint{3.007057in}{2.300028in}}%
\pgfpathlineto{\pgfqpoint{3.006467in}{2.304227in}}%
\pgfpathlineto{\pgfqpoint{3.007276in}{2.302578in}}%
\pgfpathlineto{\pgfqpoint{3.008322in}{2.303667in}}%
\pgfpathlineto{\pgfqpoint{3.008132in}{2.300828in}}%
\pgfpathlineto{\pgfqpoint{3.008405in}{2.303245in}}%
\pgfpathlineto{\pgfqpoint{3.009535in}{2.300920in}}%
\pgfpathlineto{\pgfqpoint{3.008473in}{2.304328in}}%
\pgfpathlineto{\pgfqpoint{3.009690in}{2.302923in}}%
\pgfpathlineto{\pgfqpoint{3.010289in}{2.303453in}}%
\pgfpathlineto{\pgfqpoint{3.009793in}{2.300702in}}%
\pgfpathlineto{\pgfqpoint{3.010805in}{2.303013in}}%
\pgfpathlineto{\pgfqpoint{3.011180in}{2.301797in}}%
\pgfpathlineto{\pgfqpoint{3.011380in}{2.304732in}}%
\pgfpathlineto{\pgfqpoint{3.011774in}{2.304070in}}%
\pgfpathlineto{\pgfqpoint{3.012884in}{2.304694in}}%
\pgfpathlineto{\pgfqpoint{3.012237in}{2.301990in}}%
\pgfpathlineto{\pgfqpoint{3.012889in}{2.304465in}}%
\pgfpathlineto{\pgfqpoint{3.013989in}{2.302478in}}%
\pgfpathlineto{\pgfqpoint{3.013551in}{2.304900in}}%
\pgfpathlineto{\pgfqpoint{3.014033in}{2.302949in}}%
\pgfpathlineto{\pgfqpoint{3.014481in}{2.304947in}}%
\pgfpathlineto{\pgfqpoint{3.014199in}{2.302261in}}%
\pgfpathlineto{\pgfqpoint{3.015270in}{2.303899in}}%
\pgfpathlineto{\pgfqpoint{3.016146in}{2.302578in}}%
\pgfpathlineto{\pgfqpoint{3.015742in}{2.304925in}}%
\pgfpathlineto{\pgfqpoint{3.016360in}{2.304532in}}%
\pgfpathlineto{\pgfqpoint{3.017329in}{2.305251in}}%
\pgfpathlineto{\pgfqpoint{3.016833in}{2.302344in}}%
\pgfpathlineto{\pgfqpoint{3.017465in}{2.304385in}}%
\pgfpathlineto{\pgfqpoint{3.018444in}{2.302721in}}%
\pgfpathlineto{\pgfqpoint{3.018093in}{2.304827in}}%
\pgfpathlineto{\pgfqpoint{3.018595in}{2.303721in}}%
\pgfpathlineto{\pgfqpoint{3.019447in}{2.304982in}}%
\pgfpathlineto{\pgfqpoint{3.019053in}{2.302827in}}%
\pgfpathlineto{\pgfqpoint{3.019715in}{2.304332in}}%
\pgfpathlineto{\pgfqpoint{3.019720in}{2.304375in}}%
\pgfpathlineto{\pgfqpoint{3.020094in}{2.301294in}}%
\pgfpathlineto{\pgfqpoint{3.020460in}{2.301445in}}%
\pgfpathlineto{\pgfqpoint{3.020679in}{2.299966in}}%
\pgfpathlineto{\pgfqpoint{3.021073in}{2.303898in}}%
\pgfpathlineto{\pgfqpoint{3.021560in}{2.302343in}}%
\pgfpathlineto{\pgfqpoint{3.021594in}{2.301734in}}%
\pgfpathlineto{\pgfqpoint{3.022329in}{2.304302in}}%
\pgfpathlineto{\pgfqpoint{3.022665in}{2.302532in}}%
\pgfpathlineto{\pgfqpoint{3.022918in}{2.305185in}}%
\pgfpathlineto{\pgfqpoint{3.023571in}{2.301378in}}%
\pgfpathlineto{\pgfqpoint{3.023804in}{2.303920in}}%
\pgfpathlineto{\pgfqpoint{3.024895in}{2.302679in}}%
\pgfpathlineto{\pgfqpoint{3.024130in}{2.304442in}}%
\pgfpathlineto{\pgfqpoint{3.024934in}{2.303392in}}%
\pgfpathlineto{\pgfqpoint{3.025075in}{2.305000in}}%
\pgfpathlineto{\pgfqpoint{3.025849in}{2.300603in}}%
\pgfpathlineto{\pgfqpoint{3.025922in}{2.301227in}}%
\pgfpathlineto{\pgfqpoint{3.026828in}{2.297493in}}%
\pgfpathlineto{\pgfqpoint{3.025995in}{2.301259in}}%
\pgfpathlineto{\pgfqpoint{3.027066in}{2.298987in}}%
\pgfpathlineto{\pgfqpoint{3.027368in}{2.300607in}}%
\pgfpathlineto{\pgfqpoint{3.027981in}{2.298669in}}%
\pgfpathlineto{\pgfqpoint{3.028186in}{2.299353in}}%
\pgfpathlineto{\pgfqpoint{3.028356in}{2.296399in}}%
\pgfpathlineto{\pgfqpoint{3.029199in}{2.304229in}}%
\pgfpathlineto{\pgfqpoint{3.029228in}{2.303279in}}%
\pgfpathlineto{\pgfqpoint{3.030362in}{2.304790in}}%
\pgfpathlineto{\pgfqpoint{3.030012in}{2.300837in}}%
\pgfpathlineto{\pgfqpoint{3.030367in}{2.304362in}}%
\pgfpathlineto{\pgfqpoint{3.030902in}{2.301535in}}%
\pgfpathlineto{\pgfqpoint{3.030377in}{2.304482in}}%
\pgfpathlineto{\pgfqpoint{3.031487in}{2.304399in}}%
\pgfpathlineto{\pgfqpoint{3.031832in}{2.305575in}}%
\pgfpathlineto{\pgfqpoint{3.032266in}{2.302191in}}%
\pgfpathlineto{\pgfqpoint{3.032543in}{2.303886in}}%
\pgfpathlineto{\pgfqpoint{3.033079in}{2.301156in}}%
\pgfpathlineto{\pgfqpoint{3.033254in}{2.304153in}}%
\pgfpathlineto{\pgfqpoint{3.033658in}{2.303557in}}%
\pgfpathlineto{\pgfqpoint{3.033663in}{2.303647in}}%
\pgfpathlineto{\pgfqpoint{3.034369in}{2.299662in}}%
\pgfpathlineto{\pgfqpoint{3.034729in}{2.302887in}}%
\pgfpathlineto{\pgfqpoint{3.035786in}{2.300121in}}%
\pgfpathlineto{\pgfqpoint{3.034958in}{2.304701in}}%
\pgfpathlineto{\pgfqpoint{3.035912in}{2.301354in}}%
\pgfpathlineto{\pgfqpoint{3.037076in}{2.302798in}}%
\pgfpathlineto{\pgfqpoint{3.036589in}{2.300046in}}%
\pgfpathlineto{\pgfqpoint{3.037081in}{2.302768in}}%
\pgfpathlineto{\pgfqpoint{3.037358in}{2.300375in}}%
\pgfpathlineto{\pgfqpoint{3.038157in}{2.305167in}}%
\pgfpathlineto{\pgfqpoint{3.038186in}{2.305831in}}%
\pgfpathlineto{\pgfqpoint{3.038945in}{2.302253in}}%
\pgfpathlineto{\pgfqpoint{3.039145in}{2.303014in}}%
\pgfpathlineto{\pgfqpoint{3.040299in}{2.296955in}}%
\pgfpathlineto{\pgfqpoint{3.039271in}{2.304131in}}%
\pgfpathlineto{\pgfqpoint{3.040508in}{2.297776in}}%
\pgfpathlineto{\pgfqpoint{3.040547in}{2.296766in}}%
\pgfpathlineto{\pgfqpoint{3.041014in}{2.299131in}}%
\pgfpathlineto{\pgfqpoint{3.041044in}{2.299626in}}%
\pgfpathlineto{\pgfqpoint{3.041993in}{2.296496in}}%
\pgfpathlineto{\pgfqpoint{3.042051in}{2.296890in}}%
\pgfpathlineto{\pgfqpoint{3.043108in}{2.294817in}}%
\pgfpathlineto{\pgfqpoint{3.042256in}{2.297443in}}%
\pgfpathlineto{\pgfqpoint{3.043191in}{2.295780in}}%
\pgfpathlineto{\pgfqpoint{3.043195in}{2.295695in}}%
\pgfpathlineto{\pgfqpoint{3.043965in}{2.298808in}}%
\pgfpathlineto{\pgfqpoint{3.044179in}{2.297637in}}%
\pgfpathlineto{\pgfqpoint{3.045138in}{2.300378in}}%
\pgfpathlineto{\pgfqpoint{3.044296in}{2.296993in}}%
\pgfpathlineto{\pgfqpoint{3.045304in}{2.298513in}}%
\pgfpathlineto{\pgfqpoint{3.046258in}{2.299367in}}%
\pgfpathlineto{\pgfqpoint{3.045552in}{2.296411in}}%
\pgfpathlineto{\pgfqpoint{3.046423in}{2.298672in}}%
\pgfpathlineto{\pgfqpoint{3.047329in}{2.295844in}}%
\pgfpathlineto{\pgfqpoint{3.046740in}{2.299488in}}%
\pgfpathlineto{\pgfqpoint{3.047597in}{2.296865in}}%
\pgfpathlineto{\pgfqpoint{3.048366in}{2.299879in}}%
\pgfpathlineto{\pgfqpoint{3.048692in}{2.296826in}}%
\pgfpathlineto{\pgfqpoint{3.048711in}{2.297274in}}%
\pgfpathlineto{\pgfqpoint{3.048716in}{2.297080in}}%
\pgfpathlineto{\pgfqpoint{3.049354in}{2.299727in}}%
\pgfpathlineto{\pgfqpoint{3.049778in}{2.299047in}}%
\pgfpathlineto{\pgfqpoint{3.049919in}{2.300339in}}%
\pgfpathlineto{\pgfqpoint{3.050761in}{2.296705in}}%
\pgfpathlineto{\pgfqpoint{3.050858in}{2.297807in}}%
\pgfpathlineto{\pgfqpoint{3.050863in}{2.297598in}}%
\pgfpathlineto{\pgfqpoint{3.051827in}{2.303794in}}%
\pgfpathlineto{\pgfqpoint{3.051832in}{2.303655in}}%
\pgfpathlineto{\pgfqpoint{3.052455in}{2.305348in}}%
\pgfpathlineto{\pgfqpoint{3.052027in}{2.302203in}}%
\pgfpathlineto{\pgfqpoint{3.052932in}{2.303270in}}%
\pgfpathlineto{\pgfqpoint{3.053682in}{2.302429in}}%
\pgfpathlineto{\pgfqpoint{3.053453in}{2.305022in}}%
\pgfpathlineto{\pgfqpoint{3.054042in}{2.303322in}}%
\pgfpathlineto{\pgfqpoint{3.054062in}{2.302769in}}%
\pgfpathlineto{\pgfqpoint{3.054213in}{2.304911in}}%
\pgfpathlineto{\pgfqpoint{3.054709in}{2.304922in}}%
\pgfpathlineto{\pgfqpoint{3.054719in}{2.305368in}}%
\pgfpathlineto{\pgfqpoint{3.055104in}{2.301694in}}%
\pgfpathlineto{\pgfqpoint{3.055800in}{2.303739in}}%
\pgfpathlineto{\pgfqpoint{3.056978in}{2.301603in}}%
\pgfpathlineto{\pgfqpoint{3.056214in}{2.304661in}}%
\pgfpathlineto{\pgfqpoint{3.057017in}{2.302064in}}%
\pgfpathlineto{\pgfqpoint{3.057421in}{2.304219in}}%
\pgfpathlineto{\pgfqpoint{3.057986in}{2.300724in}}%
\pgfpathlineto{\pgfqpoint{3.058117in}{2.301478in}}%
\pgfpathlineto{\pgfqpoint{3.058190in}{2.299889in}}%
\pgfpathlineto{\pgfqpoint{3.058818in}{2.304065in}}%
\pgfpathlineto{\pgfqpoint{3.059193in}{2.302711in}}%
\pgfpathlineto{\pgfqpoint{3.060152in}{2.303720in}}%
\pgfpathlineto{\pgfqpoint{3.059802in}{2.300530in}}%
\pgfpathlineto{\pgfqpoint{3.060298in}{2.302450in}}%
\pgfpathlineto{\pgfqpoint{3.060576in}{2.300762in}}%
\pgfpathlineto{\pgfqpoint{3.061253in}{2.303417in}}%
\pgfpathlineto{\pgfqpoint{3.061404in}{2.302508in}}%
\pgfpathlineto{\pgfqpoint{3.061842in}{2.304794in}}%
\pgfpathlineto{\pgfqpoint{3.062304in}{2.301413in}}%
\pgfpathlineto{\pgfqpoint{3.062538in}{2.304079in}}%
\pgfpathlineto{\pgfqpoint{3.062596in}{2.304469in}}%
\pgfpathlineto{\pgfqpoint{3.062927in}{2.301297in}}%
\pgfpathlineto{\pgfqpoint{3.063507in}{2.299961in}}%
\pgfpathlineto{\pgfqpoint{3.063921in}{2.303805in}}%
\pgfpathlineto{\pgfqpoint{3.064432in}{2.301527in}}%
\pgfpathlineto{\pgfqpoint{3.064851in}{2.304705in}}%
\pgfpathlineto{\pgfqpoint{3.064865in}{2.304652in}}%
\pgfpathlineto{\pgfqpoint{3.064894in}{2.305058in}}%
\pgfpathlineto{\pgfqpoint{3.065391in}{2.302003in}}%
\pgfpathlineto{\pgfqpoint{3.065547in}{2.302685in}}%
\pgfpathlineto{\pgfqpoint{3.066292in}{2.299197in}}%
\pgfpathlineto{\pgfqpoint{3.066686in}{2.299840in}}%
\pgfpathlineto{\pgfqpoint{3.067066in}{2.300433in}}%
\pgfpathlineto{\pgfqpoint{3.066735in}{2.298764in}}%
\pgfpathlineto{\pgfqpoint{3.067173in}{2.298641in}}%
\pgfpathlineto{\pgfqpoint{3.067309in}{2.298071in}}%
\pgfpathlineto{\pgfqpoint{3.067718in}{2.302921in}}%
\pgfpathlineto{\pgfqpoint{3.068258in}{2.300665in}}%
\pgfpathlineto{\pgfqpoint{3.069222in}{2.301720in}}%
\pgfpathlineto{\pgfqpoint{3.068575in}{2.298729in}}%
\pgfpathlineto{\pgfqpoint{3.069373in}{2.301073in}}%
\pgfpathlineto{\pgfqpoint{3.070104in}{2.298877in}}%
\pgfpathlineto{\pgfqpoint{3.069393in}{2.301371in}}%
\pgfpathlineto{\pgfqpoint{3.070527in}{2.300149in}}%
\pgfpathlineto{\pgfqpoint{3.071623in}{2.303259in}}%
\pgfpathlineto{\pgfqpoint{3.070785in}{2.298976in}}%
\pgfpathlineto{\pgfqpoint{3.071701in}{2.302751in}}%
\pgfpathlineto{\pgfqpoint{3.071871in}{2.300438in}}%
\pgfpathlineto{\pgfqpoint{3.072178in}{2.304433in}}%
\pgfpathlineto{\pgfqpoint{3.072840in}{2.301014in}}%
\pgfpathlineto{\pgfqpoint{3.073594in}{2.298637in}}%
\pgfpathlineto{\pgfqpoint{3.072923in}{2.302308in}}%
\pgfpathlineto{\pgfqpoint{3.074047in}{2.300767in}}%
\pgfpathlineto{\pgfqpoint{3.074495in}{2.304213in}}%
\pgfpathlineto{\pgfqpoint{3.075191in}{2.303731in}}%
\pgfpathlineto{\pgfqpoint{3.075303in}{2.304581in}}%
\pgfpathlineto{\pgfqpoint{3.075907in}{2.300369in}}%
\pgfpathlineto{\pgfqpoint{3.076145in}{2.301060in}}%
\pgfpathlineto{\pgfqpoint{3.077324in}{2.297267in}}%
\pgfpathlineto{\pgfqpoint{3.076199in}{2.301263in}}%
\pgfpathlineto{\pgfqpoint{3.077358in}{2.297717in}}%
\pgfpathlineto{\pgfqpoint{3.077811in}{2.297395in}}%
\pgfpathlineto{\pgfqpoint{3.078224in}{2.299695in}}%
\pgfpathlineto{\pgfqpoint{3.079135in}{2.303528in}}%
\pgfpathlineto{\pgfqpoint{3.079354in}{2.301319in}}%
\pgfpathlineto{\pgfqpoint{3.080425in}{2.299176in}}%
\pgfpathlineto{\pgfqpoint{3.079578in}{2.303272in}}%
\pgfpathlineto{\pgfqpoint{3.080522in}{2.299971in}}%
\pgfpathlineto{\pgfqpoint{3.080956in}{2.301588in}}%
\pgfpathlineto{\pgfqpoint{3.081126in}{2.298781in}}%
\pgfpathlineto{\pgfqpoint{3.081705in}{2.300509in}}%
\pgfpathlineto{\pgfqpoint{3.081895in}{2.299011in}}%
\pgfpathlineto{\pgfqpoint{3.082426in}{2.303711in}}%
\pgfpathlineto{\pgfqpoint{3.082781in}{2.303304in}}%
\pgfpathlineto{\pgfqpoint{3.083142in}{2.301734in}}%
\pgfpathlineto{\pgfqpoint{3.083677in}{2.304429in}}%
\pgfpathlineto{\pgfqpoint{3.083862in}{2.303756in}}%
\pgfpathlineto{\pgfqpoint{3.084198in}{2.306020in}}%
\pgfpathlineto{\pgfqpoint{3.084865in}{2.301356in}}%
\pgfpathlineto{\pgfqpoint{3.084884in}{2.301413in}}%
\pgfpathlineto{\pgfqpoint{3.085863in}{2.298695in}}%
\pgfpathlineto{\pgfqpoint{3.085035in}{2.303818in}}%
\pgfpathlineto{\pgfqpoint{3.086053in}{2.299829in}}%
\pgfpathlineto{\pgfqpoint{3.087460in}{2.304712in}}%
\pgfpathlineto{\pgfqpoint{3.086350in}{2.298690in}}%
\pgfpathlineto{\pgfqpoint{3.087723in}{2.302904in}}%
\pgfpathlineto{\pgfqpoint{3.087849in}{2.301614in}}%
\pgfpathlineto{\pgfqpoint{3.088020in}{2.303979in}}%
\pgfpathlineto{\pgfqpoint{3.088843in}{2.301780in}}%
\pgfpathlineto{\pgfqpoint{3.089880in}{2.303067in}}%
\pgfpathlineto{\pgfqpoint{3.089081in}{2.300296in}}%
\pgfpathlineto{\pgfqpoint{3.089967in}{2.302017in}}%
\pgfpathlineto{\pgfqpoint{3.090435in}{2.299354in}}%
\pgfpathlineto{\pgfqpoint{3.090040in}{2.302176in}}%
\pgfpathlineto{\pgfqpoint{3.091106in}{2.300368in}}%
\pgfpathlineto{\pgfqpoint{3.092289in}{2.304728in}}%
\pgfpathlineto{\pgfqpoint{3.091369in}{2.299762in}}%
\pgfpathlineto{\pgfqpoint{3.092343in}{2.302979in}}%
\pgfpathlineto{\pgfqpoint{3.093492in}{2.300108in}}%
\pgfpathlineto{\pgfqpoint{3.092976in}{2.304114in}}%
\pgfpathlineto{\pgfqpoint{3.093536in}{2.300212in}}%
\pgfpathlineto{\pgfqpoint{3.094310in}{2.303277in}}%
\pgfpathlineto{\pgfqpoint{3.093555in}{2.299939in}}%
\pgfpathlineto{\pgfqpoint{3.094904in}{2.301166in}}%
\pgfpathlineto{\pgfqpoint{3.095201in}{2.299880in}}%
\pgfpathlineto{\pgfqpoint{3.095556in}{2.303170in}}%
\pgfpathlineto{\pgfqpoint{3.095970in}{2.303140in}}%
\pgfpathlineto{\pgfqpoint{3.095985in}{2.303776in}}%
\pgfpathlineto{\pgfqpoint{3.096184in}{2.301764in}}%
\pgfpathlineto{\pgfqpoint{3.096204in}{2.301384in}}%
\pgfpathlineto{\pgfqpoint{3.097061in}{2.304979in}}%
\pgfpathlineto{\pgfqpoint{3.097285in}{2.302149in}}%
\pgfpathlineto{\pgfqpoint{3.097338in}{2.301740in}}%
\pgfpathlineto{\pgfqpoint{3.097840in}{2.303199in}}%
\pgfpathlineto{\pgfqpoint{3.098911in}{2.305606in}}%
\pgfpathlineto{\pgfqpoint{3.098015in}{2.302589in}}%
\pgfpathlineto{\pgfqpoint{3.098959in}{2.303833in}}%
\pgfpathlineto{\pgfqpoint{3.100235in}{2.298183in}}%
\pgfpathlineto{\pgfqpoint{3.099018in}{2.304406in}}%
\pgfpathlineto{\pgfqpoint{3.100254in}{2.298273in}}%
\pgfpathlineto{\pgfqpoint{3.101432in}{2.304481in}}%
\pgfpathlineto{\pgfqpoint{3.101437in}{2.304389in}}%
\pgfpathlineto{\pgfqpoint{3.101457in}{2.305170in}}%
\pgfpathlineto{\pgfqpoint{3.102275in}{2.300413in}}%
\pgfpathlineto{\pgfqpoint{3.102333in}{2.300724in}}%
\pgfpathlineto{\pgfqpoint{3.102547in}{2.300022in}}%
\pgfpathlineto{\pgfqpoint{3.103200in}{2.303802in}}%
\pgfpathlineto{\pgfqpoint{3.103399in}{2.302553in}}%
\pgfpathlineto{\pgfqpoint{3.103696in}{2.305025in}}%
\pgfpathlineto{\pgfqpoint{3.104125in}{2.302207in}}%
\pgfpathlineto{\pgfqpoint{3.104578in}{2.303715in}}%
\pgfpathlineto{\pgfqpoint{3.104860in}{2.302658in}}%
\pgfpathlineto{\pgfqpoint{3.105293in}{2.305118in}}%
\pgfpathlineto{\pgfqpoint{3.105697in}{2.303074in}}%
\pgfpathlineto{\pgfqpoint{3.106145in}{2.304689in}}%
\pgfpathlineto{\pgfqpoint{3.106720in}{2.300908in}}%
\pgfpathlineto{\pgfqpoint{3.106725in}{2.300930in}}%
\pgfpathlineto{\pgfqpoint{3.107158in}{2.300339in}}%
\pgfpathlineto{\pgfqpoint{3.107742in}{2.303360in}}%
\pgfpathlineto{\pgfqpoint{3.107786in}{2.302757in}}%
\pgfpathlineto{\pgfqpoint{3.108166in}{2.304777in}}%
\pgfpathlineto{\pgfqpoint{3.108876in}{2.302013in}}%
\pgfpathlineto{\pgfqpoint{3.108886in}{2.302327in}}%
\pgfpathlineto{\pgfqpoint{3.109422in}{2.304038in}}%
\pgfpathlineto{\pgfqpoint{3.109719in}{2.301610in}}%
\pgfpathlineto{\pgfqpoint{3.110050in}{2.303163in}}%
\pgfpathlineto{\pgfqpoint{3.110483in}{2.302290in}}%
\pgfpathlineto{\pgfqpoint{3.111135in}{2.304470in}}%
\pgfpathlineto{\pgfqpoint{3.111140in}{2.304236in}}%
\pgfpathlineto{\pgfqpoint{3.111145in}{2.304316in}}%
\pgfpathlineto{\pgfqpoint{3.111744in}{2.299658in}}%
\pgfpathlineto{\pgfqpoint{3.112094in}{2.301672in}}%
\pgfpathlineto{\pgfqpoint{3.113039in}{2.298012in}}%
\pgfpathlineto{\pgfqpoint{3.112328in}{2.301858in}}%
\pgfpathlineto{\pgfqpoint{3.113258in}{2.299809in}}%
\pgfpathlineto{\pgfqpoint{3.114232in}{2.302289in}}%
\pgfpathlineto{\pgfqpoint{3.113443in}{2.298917in}}%
\pgfpathlineto{\pgfqpoint{3.114412in}{2.301179in}}%
\pgfpathlineto{\pgfqpoint{3.114704in}{2.302479in}}%
\pgfpathlineto{\pgfqpoint{3.115171in}{2.299118in}}%
\pgfpathlineto{\pgfqpoint{3.115347in}{2.299982in}}%
\pgfpathlineto{\pgfqpoint{3.116432in}{2.299197in}}%
\pgfpathlineto{\pgfqpoint{3.116233in}{2.301378in}}%
\pgfpathlineto{\pgfqpoint{3.116466in}{2.299514in}}%
\pgfpathlineto{\pgfqpoint{3.117552in}{2.301722in}}%
\pgfpathlineto{\pgfqpoint{3.116944in}{2.298444in}}%
\pgfpathlineto{\pgfqpoint{3.117635in}{2.300198in}}%
\pgfpathlineto{\pgfqpoint{3.118068in}{2.298053in}}%
\pgfpathlineto{\pgfqpoint{3.118633in}{2.301036in}}%
\pgfpathlineto{\pgfqpoint{3.118643in}{2.300938in}}%
\pgfpathlineto{\pgfqpoint{3.119694in}{2.304642in}}%
\pgfpathlineto{\pgfqpoint{3.118730in}{2.300156in}}%
\pgfpathlineto{\pgfqpoint{3.119845in}{2.302463in}}%
\pgfpathlineto{\pgfqpoint{3.119938in}{2.302161in}}%
\pgfpathlineto{\pgfqpoint{3.120069in}{2.304073in}}%
\pgfpathlineto{\pgfqpoint{3.120931in}{2.303737in}}%
\pgfpathlineto{\pgfqpoint{3.121505in}{2.305358in}}%
\pgfpathlineto{\pgfqpoint{3.121189in}{2.301819in}}%
\pgfpathlineto{\pgfqpoint{3.122007in}{2.302112in}}%
\pgfpathlineto{\pgfqpoint{3.122216in}{2.301244in}}%
\pgfpathlineto{\pgfqpoint{3.123015in}{2.304676in}}%
\pgfpathlineto{\pgfqpoint{3.123019in}{2.304557in}}%
\pgfpathlineto{\pgfqpoint{3.124227in}{2.300657in}}%
\pgfpathlineto{\pgfqpoint{3.123414in}{2.305745in}}%
\pgfpathlineto{\pgfqpoint{3.124300in}{2.301474in}}%
\pgfpathlineto{\pgfqpoint{3.125405in}{2.302719in}}%
\pgfpathlineto{\pgfqpoint{3.124602in}{2.299257in}}%
\pgfpathlineto{\pgfqpoint{3.125444in}{2.302197in}}%
\pgfpathlineto{\pgfqpoint{3.125936in}{2.300510in}}%
\pgfpathlineto{\pgfqpoint{3.126135in}{2.304044in}}%
\pgfpathlineto{\pgfqpoint{3.126549in}{2.302352in}}%
\pgfpathlineto{\pgfqpoint{3.127494in}{2.304650in}}%
\pgfpathlineto{\pgfqpoint{3.126617in}{2.301557in}}%
\pgfpathlineto{\pgfqpoint{3.127703in}{2.303218in}}%
\pgfpathlineto{\pgfqpoint{3.128560in}{2.301428in}}%
\pgfpathlineto{\pgfqpoint{3.127903in}{2.304474in}}%
\pgfpathlineto{\pgfqpoint{3.128789in}{2.304215in}}%
\pgfpathlineto{\pgfqpoint{3.129884in}{2.302182in}}%
\pgfpathlineto{\pgfqpoint{3.128905in}{2.305489in}}%
\pgfpathlineto{\pgfqpoint{3.129981in}{2.303544in}}%
\pgfpathlineto{\pgfqpoint{3.130420in}{2.306634in}}%
\pgfpathlineto{\pgfqpoint{3.130921in}{2.302112in}}%
\pgfpathlineto{\pgfqpoint{3.131082in}{2.303547in}}%
\pgfpathlineto{\pgfqpoint{3.131403in}{2.301440in}}%
\pgfpathlineto{\pgfqpoint{3.132031in}{2.305324in}}%
\pgfpathlineto{\pgfqpoint{3.132211in}{2.302570in}}%
\pgfpathlineto{\pgfqpoint{3.133487in}{2.304942in}}%
\pgfpathlineto{\pgfqpoint{3.132279in}{2.301560in}}%
\pgfpathlineto{\pgfqpoint{3.133531in}{2.304219in}}%
\pgfpathlineto{\pgfqpoint{3.133925in}{2.298752in}}%
\pgfpathlineto{\pgfqpoint{3.134412in}{2.304268in}}%
\pgfpathlineto{\pgfqpoint{3.134718in}{2.301802in}}%
\pgfpathlineto{\pgfqpoint{3.135916in}{2.303984in}}%
\pgfpathlineto{\pgfqpoint{3.134923in}{2.301148in}}%
\pgfpathlineto{\pgfqpoint{3.135960in}{2.303723in}}%
\pgfpathlineto{\pgfqpoint{3.137007in}{2.299911in}}%
\pgfpathlineto{\pgfqpoint{3.136233in}{2.305803in}}%
\pgfpathlineto{\pgfqpoint{3.137270in}{2.302004in}}%
\pgfpathlineto{\pgfqpoint{3.137484in}{2.303332in}}%
\pgfpathlineto{\pgfqpoint{3.138190in}{2.298657in}}%
\pgfpathlineto{\pgfqpoint{3.138195in}{2.298746in}}%
\pgfpathlineto{\pgfqpoint{3.139348in}{2.297396in}}%
\pgfpathlineto{\pgfqpoint{3.138759in}{2.300813in}}%
\pgfpathlineto{\pgfqpoint{3.139368in}{2.297645in}}%
\pgfpathlineto{\pgfqpoint{3.140419in}{2.299399in}}%
\pgfpathlineto{\pgfqpoint{3.139850in}{2.296493in}}%
\pgfpathlineto{\pgfqpoint{3.140502in}{2.298675in}}%
\pgfpathlineto{\pgfqpoint{3.141062in}{2.296054in}}%
\pgfpathlineto{\pgfqpoint{3.141627in}{2.297859in}}%
\pgfpathlineto{\pgfqpoint{3.142727in}{2.299358in}}%
\pgfpathlineto{\pgfqpoint{3.141826in}{2.297555in}}%
\pgfpathlineto{\pgfqpoint{3.142766in}{2.298500in}}%
\pgfpathlineto{\pgfqpoint{3.143034in}{2.295779in}}%
\pgfpathlineto{\pgfqpoint{3.143711in}{2.298628in}}%
\pgfpathlineto{\pgfqpoint{3.143925in}{2.296555in}}%
\pgfpathlineto{\pgfqpoint{3.144908in}{2.300828in}}%
\pgfpathlineto{\pgfqpoint{3.144066in}{2.295101in}}%
\pgfpathlineto{\pgfqpoint{3.145191in}{2.299802in}}%
\pgfpathlineto{\pgfqpoint{3.145478in}{2.298621in}}%
\pgfpathlineto{\pgfqpoint{3.145926in}{2.301157in}}%
\pgfpathlineto{\pgfqpoint{3.146082in}{2.300964in}}%
\pgfpathlineto{\pgfqpoint{3.147153in}{2.304384in}}%
\pgfpathlineto{\pgfqpoint{3.146174in}{2.300833in}}%
\pgfpathlineto{\pgfqpoint{3.147240in}{2.303814in}}%
\pgfpathlineto{\pgfqpoint{3.148384in}{2.305504in}}%
\pgfpathlineto{\pgfqpoint{3.147586in}{2.301753in}}%
\pgfpathlineto{\pgfqpoint{3.148438in}{2.304880in}}%
\pgfpathlineto{\pgfqpoint{3.148710in}{2.302264in}}%
\pgfpathlineto{\pgfqpoint{3.148501in}{2.304892in}}%
\pgfpathlineto{\pgfqpoint{3.149562in}{2.303972in}}%
\pgfpathlineto{\pgfqpoint{3.149874in}{2.305171in}}%
\pgfpathlineto{\pgfqpoint{3.150312in}{2.301909in}}%
\pgfpathlineto{\pgfqpoint{3.150614in}{2.302792in}}%
\pgfpathlineto{\pgfqpoint{3.151646in}{2.301944in}}%
\pgfpathlineto{\pgfqpoint{3.151154in}{2.304059in}}%
\pgfpathlineto{\pgfqpoint{3.151724in}{2.302201in}}%
\pgfpathlineto{\pgfqpoint{3.152274in}{2.304946in}}%
\pgfpathlineto{\pgfqpoint{3.152844in}{2.303447in}}%
\pgfpathlineto{\pgfqpoint{3.153151in}{2.301377in}}%
\pgfpathlineto{\pgfqpoint{3.153438in}{2.304840in}}%
\pgfpathlineto{\pgfqpoint{3.154017in}{2.302722in}}%
\pgfpathlineto{\pgfqpoint{3.154937in}{2.304250in}}%
\pgfpathlineto{\pgfqpoint{3.154280in}{2.301504in}}%
\pgfpathlineto{\pgfqpoint{3.155122in}{2.302758in}}%
\pgfpathlineto{\pgfqpoint{3.155293in}{2.300988in}}%
\pgfpathlineto{\pgfqpoint{3.156130in}{2.304428in}}%
\pgfpathlineto{\pgfqpoint{3.156193in}{2.303872in}}%
\pgfpathlineto{\pgfqpoint{3.156300in}{2.304165in}}%
\pgfpathlineto{\pgfqpoint{3.156729in}{2.302358in}}%
\pgfpathlineto{\pgfqpoint{3.157289in}{2.303322in}}%
\pgfpathlineto{\pgfqpoint{3.157566in}{2.302059in}}%
\pgfpathlineto{\pgfqpoint{3.157732in}{2.304209in}}%
\pgfpathlineto{\pgfqpoint{3.158379in}{2.303255in}}%
\pgfpathlineto{\pgfqpoint{3.158443in}{2.304243in}}%
\pgfpathlineto{\pgfqpoint{3.159426in}{2.298946in}}%
\pgfpathlineto{\pgfqpoint{3.159431in}{2.299030in}}%
\pgfpathlineto{\pgfqpoint{3.160249in}{2.297499in}}%
\pgfpathlineto{\pgfqpoint{3.160005in}{2.300974in}}%
\pgfpathlineto{\pgfqpoint{3.160575in}{2.298274in}}%
\pgfpathlineto{\pgfqpoint{3.161641in}{2.303657in}}%
\pgfpathlineto{\pgfqpoint{3.160872in}{2.297484in}}%
\pgfpathlineto{\pgfqpoint{3.161904in}{2.302115in}}%
\pgfpathlineto{\pgfqpoint{3.161909in}{2.301977in}}%
\pgfpathlineto{\pgfqpoint{3.162070in}{2.305067in}}%
\pgfpathlineto{\pgfqpoint{3.162965in}{2.303445in}}%
\pgfpathlineto{\pgfqpoint{3.163428in}{2.302204in}}%
\pgfpathlineto{\pgfqpoint{3.163720in}{2.304970in}}%
\pgfpathlineto{\pgfqpoint{3.163744in}{2.304964in}}%
\pgfpathlineto{\pgfqpoint{3.163759in}{2.305052in}}%
\pgfpathlineto{\pgfqpoint{3.164314in}{2.300978in}}%
\pgfpathlineto{\pgfqpoint{3.164640in}{2.302379in}}%
\pgfpathlineto{\pgfqpoint{3.165560in}{2.299112in}}%
\pgfpathlineto{\pgfqpoint{3.165176in}{2.303498in}}%
\pgfpathlineto{\pgfqpoint{3.165760in}{2.301502in}}%
\pgfpathlineto{\pgfqpoint{3.165857in}{2.302828in}}%
\pgfpathlineto{\pgfqpoint{3.166602in}{2.298838in}}%
\pgfpathlineto{\pgfqpoint{3.166787in}{2.299838in}}%
\pgfpathlineto{\pgfqpoint{3.167279in}{2.297099in}}%
\pgfpathlineto{\pgfqpoint{3.167844in}{2.300681in}}%
\pgfpathlineto{\pgfqpoint{3.167897in}{2.299811in}}%
\pgfpathlineto{\pgfqpoint{3.167931in}{2.300127in}}%
\pgfpathlineto{\pgfqpoint{3.168481in}{2.294733in}}%
\pgfpathlineto{\pgfqpoint{3.168959in}{2.297133in}}%
\pgfpathlineto{\pgfqpoint{3.169080in}{2.297063in}}%
\pgfpathlineto{\pgfqpoint{3.169368in}{2.299246in}}%
\pgfpathlineto{\pgfqpoint{3.169377in}{2.299122in}}%
\pgfpathlineto{\pgfqpoint{3.169504in}{2.299315in}}%
\pgfpathlineto{\pgfqpoint{3.169592in}{2.296592in}}%
\pgfpathlineto{\pgfqpoint{3.170429in}{2.297419in}}%
\pgfpathlineto{\pgfqpoint{3.171203in}{2.296710in}}%
\pgfpathlineto{\pgfqpoint{3.171446in}{2.298859in}}%
\pgfpathlineto{\pgfqpoint{3.171510in}{2.298514in}}%
\pgfpathlineto{\pgfqpoint{3.172488in}{2.300593in}}%
\pgfpathlineto{\pgfqpoint{3.172001in}{2.297372in}}%
\pgfpathlineto{\pgfqpoint{3.172702in}{2.299682in}}%
\pgfpathlineto{\pgfqpoint{3.173856in}{2.296345in}}%
\pgfpathlineto{\pgfqpoint{3.173146in}{2.300596in}}%
\pgfpathlineto{\pgfqpoint{3.173905in}{2.297758in}}%
\pgfpathlineto{\pgfqpoint{3.174036in}{2.299116in}}%
\pgfpathlineto{\pgfqpoint{3.174616in}{2.294464in}}%
\pgfpathlineto{\pgfqpoint{3.175000in}{2.297037in}}%
\pgfpathlineto{\pgfqpoint{3.175789in}{2.294256in}}%
\pgfpathlineto{\pgfqpoint{3.175093in}{2.298196in}}%
\pgfpathlineto{\pgfqpoint{3.176179in}{2.295885in}}%
\pgfpathlineto{\pgfqpoint{3.177274in}{2.298112in}}%
\pgfpathlineto{\pgfqpoint{3.176777in}{2.294967in}}%
\pgfpathlineto{\pgfqpoint{3.177347in}{2.297581in}}%
\pgfpathlineto{\pgfqpoint{3.177878in}{2.296938in}}%
\pgfpathlineto{\pgfqpoint{3.178379in}{2.300095in}}%
\pgfpathlineto{\pgfqpoint{3.178608in}{2.302723in}}%
\pgfpathlineto{\pgfqpoint{3.179353in}{2.297858in}}%
\pgfpathlineto{\pgfqpoint{3.179460in}{2.299118in}}%
\pgfpathlineto{\pgfqpoint{3.179479in}{2.298892in}}%
\pgfpathlineto{\pgfqpoint{3.179689in}{2.301406in}}%
\pgfpathlineto{\pgfqpoint{3.180555in}{2.299636in}}%
\pgfpathlineto{\pgfqpoint{3.180930in}{2.300028in}}%
\pgfpathlineto{\pgfqpoint{3.181573in}{2.297587in}}%
\pgfpathlineto{\pgfqpoint{3.181607in}{2.298252in}}%
\pgfpathlineto{\pgfqpoint{3.182133in}{2.295124in}}%
\pgfpathlineto{\pgfqpoint{3.182683in}{2.299024in}}%
\pgfpathlineto{\pgfqpoint{3.182707in}{2.298652in}}%
\pgfpathlineto{\pgfqpoint{3.183642in}{2.299681in}}%
\pgfpathlineto{\pgfqpoint{3.183204in}{2.295639in}}%
\pgfpathlineto{\pgfqpoint{3.183735in}{2.297991in}}%
\pgfpathlineto{\pgfqpoint{3.184937in}{2.296318in}}%
\pgfpathlineto{\pgfqpoint{3.184114in}{2.299073in}}%
\pgfpathlineto{\pgfqpoint{3.184961in}{2.297228in}}%
\pgfpathlineto{\pgfqpoint{3.185093in}{2.298087in}}%
\pgfpathlineto{\pgfqpoint{3.185745in}{2.294497in}}%
\pgfpathlineto{\pgfqpoint{3.186023in}{2.296551in}}%
\pgfpathlineto{\pgfqpoint{3.186042in}{2.296194in}}%
\pgfpathlineto{\pgfqpoint{3.186534in}{2.298465in}}%
\pgfpathlineto{\pgfqpoint{3.187055in}{2.298213in}}%
\pgfpathlineto{\pgfqpoint{3.187327in}{2.300668in}}%
\pgfpathlineto{\pgfqpoint{3.188199in}{2.299247in}}%
\pgfpathlineto{\pgfqpoint{3.188204in}{2.298965in}}%
\pgfpathlineto{\pgfqpoint{3.189007in}{2.303776in}}%
\pgfpathlineto{\pgfqpoint{3.189246in}{2.302745in}}%
\pgfpathlineto{\pgfqpoint{3.190146in}{2.304081in}}%
\pgfpathlineto{\pgfqpoint{3.189358in}{2.301946in}}%
\pgfpathlineto{\pgfqpoint{3.190385in}{2.303751in}}%
\pgfpathlineto{\pgfqpoint{3.191388in}{2.299129in}}%
\pgfpathlineto{\pgfqpoint{3.190852in}{2.303908in}}%
\pgfpathlineto{\pgfqpoint{3.191519in}{2.301987in}}%
\pgfpathlineto{\pgfqpoint{3.192396in}{2.304610in}}%
\pgfpathlineto{\pgfqpoint{3.192211in}{2.301017in}}%
\pgfpathlineto{\pgfqpoint{3.192644in}{2.302585in}}%
\pgfpathlineto{\pgfqpoint{3.193428in}{2.298711in}}%
\pgfpathlineto{\pgfqpoint{3.193842in}{2.299655in}}%
\pgfpathlineto{\pgfqpoint{3.194625in}{2.303772in}}%
\pgfpathlineto{\pgfqpoint{3.193919in}{2.299357in}}%
\pgfpathlineto{\pgfqpoint{3.194630in}{2.303767in}}%
\pgfpathlineto{\pgfqpoint{3.194723in}{2.304140in}}%
\pgfpathlineto{\pgfqpoint{3.195560in}{2.298057in}}%
\pgfpathlineto{\pgfqpoint{3.195575in}{2.298069in}}%
\pgfpathlineto{\pgfqpoint{3.197069in}{2.303881in}}%
\pgfpathlineto{\pgfqpoint{3.195935in}{2.296344in}}%
\pgfpathlineto{\pgfqpoint{3.197327in}{2.302809in}}%
\pgfpathlineto{\pgfqpoint{3.197751in}{2.301837in}}%
\pgfpathlineto{\pgfqpoint{3.198101in}{2.305641in}}%
\pgfpathlineto{\pgfqpoint{3.198374in}{2.304578in}}%
\pgfpathlineto{\pgfqpoint{3.199099in}{2.305846in}}%
\pgfpathlineto{\pgfqpoint{3.199416in}{2.303199in}}%
\pgfpathlineto{\pgfqpoint{3.199489in}{2.302654in}}%
\pgfpathlineto{\pgfqpoint{3.199932in}{2.305115in}}%
\pgfpathlineto{\pgfqpoint{3.200521in}{2.303441in}}%
\pgfpathlineto{\pgfqpoint{3.200701in}{2.303743in}}%
\pgfpathlineto{\pgfqpoint{3.201081in}{2.300142in}}%
\pgfpathlineto{\pgfqpoint{3.201295in}{2.300430in}}%
\pgfpathlineto{\pgfqpoint{3.201315in}{2.299894in}}%
\pgfpathlineto{\pgfqpoint{3.202055in}{2.302507in}}%
\pgfpathlineto{\pgfqpoint{3.202391in}{2.301705in}}%
\pgfpathlineto{\pgfqpoint{3.203067in}{2.303065in}}%
\pgfpathlineto{\pgfqpoint{3.202707in}{2.300565in}}%
\pgfpathlineto{\pgfqpoint{3.203496in}{2.301594in}}%
\pgfpathlineto{\pgfqpoint{3.203768in}{2.302987in}}%
\pgfpathlineto{\pgfqpoint{3.204070in}{2.298609in}}%
\pgfpathlineto{\pgfqpoint{3.204099in}{2.298661in}}%
\pgfpathlineto{\pgfqpoint{3.204158in}{2.298095in}}%
\pgfpathlineto{\pgfqpoint{3.205088in}{2.303189in}}%
\pgfpathlineto{\pgfqpoint{3.205219in}{2.304289in}}%
\pgfpathlineto{\pgfqpoint{3.205852in}{2.300449in}}%
\pgfpathlineto{\pgfqpoint{3.206169in}{2.302009in}}%
\pgfpathlineto{\pgfqpoint{3.206456in}{2.300895in}}%
\pgfpathlineto{\pgfqpoint{3.207084in}{2.305145in}}%
\pgfpathlineto{\pgfqpoint{3.207254in}{2.302911in}}%
\pgfpathlineto{\pgfqpoint{3.208096in}{2.304637in}}%
\pgfpathlineto{\pgfqpoint{3.207731in}{2.301845in}}%
\pgfpathlineto{\pgfqpoint{3.208374in}{2.303535in}}%
\pgfpathlineto{\pgfqpoint{3.209450in}{2.301191in}}%
\pgfpathlineto{\pgfqpoint{3.208603in}{2.304143in}}%
\pgfpathlineto{\pgfqpoint{3.209650in}{2.302383in}}%
\pgfpathlineto{\pgfqpoint{3.210097in}{2.304779in}}%
\pgfpathlineto{\pgfqpoint{3.210609in}{2.301113in}}%
\pgfpathlineto{\pgfqpoint{3.210672in}{2.301377in}}%
\pgfpathlineto{\pgfqpoint{3.211957in}{2.298172in}}%
\pgfpathlineto{\pgfqpoint{3.211402in}{2.301904in}}%
\pgfpathlineto{\pgfqpoint{3.212045in}{2.299210in}}%
\pgfpathlineto{\pgfqpoint{3.212872in}{2.304403in}}%
\pgfpathlineto{\pgfqpoint{3.213238in}{2.303022in}}%
\pgfpathlineto{\pgfqpoint{3.214586in}{2.297381in}}%
\pgfpathlineto{\pgfqpoint{3.213447in}{2.303953in}}%
\pgfpathlineto{\pgfqpoint{3.214635in}{2.298175in}}%
\pgfpathlineto{\pgfqpoint{3.215876in}{2.302872in}}%
\pgfpathlineto{\pgfqpoint{3.214786in}{2.297995in}}%
\pgfpathlineto{\pgfqpoint{3.215930in}{2.302774in}}%
\pgfpathlineto{\pgfqpoint{3.216305in}{2.301145in}}%
\pgfpathlineto{\pgfqpoint{3.216066in}{2.303629in}}%
\pgfpathlineto{\pgfqpoint{3.217030in}{2.303381in}}%
\pgfpathlineto{\pgfqpoint{3.217220in}{2.303916in}}%
\pgfpathlineto{\pgfqpoint{3.217551in}{2.301595in}}%
\pgfpathlineto{\pgfqpoint{3.218121in}{2.303056in}}%
\pgfpathlineto{\pgfqpoint{3.218671in}{2.300986in}}%
\pgfpathlineto{\pgfqpoint{3.219021in}{2.305106in}}%
\pgfpathlineto{\pgfqpoint{3.219226in}{2.303152in}}%
\pgfpathlineto{\pgfqpoint{3.220375in}{2.304238in}}%
\pgfpathlineto{\pgfqpoint{3.219357in}{2.301650in}}%
\pgfpathlineto{\pgfqpoint{3.220399in}{2.303900in}}%
\pgfpathlineto{\pgfqpoint{3.221509in}{2.301778in}}%
\pgfpathlineto{\pgfqpoint{3.221276in}{2.304785in}}%
\pgfpathlineto{\pgfqpoint{3.221548in}{2.302122in}}%
\pgfpathlineto{\pgfqpoint{3.222629in}{2.305496in}}%
\pgfpathlineto{\pgfqpoint{3.221680in}{2.301310in}}%
\pgfpathlineto{\pgfqpoint{3.222756in}{2.305208in}}%
\pgfpathlineto{\pgfqpoint{3.222931in}{2.306309in}}%
\pgfpathlineto{\pgfqpoint{3.223350in}{2.302164in}}%
\pgfpathlineto{\pgfqpoint{3.223447in}{2.302241in}}%
\pgfpathlineto{\pgfqpoint{3.224031in}{2.301605in}}%
\pgfpathlineto{\pgfqpoint{3.224508in}{2.304098in}}%
\pgfpathlineto{\pgfqpoint{3.224547in}{2.302708in}}%
\pgfpathlineto{\pgfqpoint{3.224572in}{2.303521in}}%
\pgfpathlineto{\pgfqpoint{3.225482in}{2.300309in}}%
\pgfpathlineto{\pgfqpoint{3.225540in}{2.300462in}}%
\pgfpathlineto{\pgfqpoint{3.225555in}{2.299944in}}%
\pgfpathlineto{\pgfqpoint{3.226368in}{2.304018in}}%
\pgfpathlineto{\pgfqpoint{3.226568in}{2.303276in}}%
\pgfpathlineto{\pgfqpoint{3.226616in}{2.304254in}}%
\pgfpathlineto{\pgfqpoint{3.227118in}{2.301826in}}%
\pgfpathlineto{\pgfqpoint{3.227673in}{2.303280in}}%
\pgfpathlineto{\pgfqpoint{3.228812in}{2.299522in}}%
\pgfpathlineto{\pgfqpoint{3.228052in}{2.304141in}}%
\pgfpathlineto{\pgfqpoint{3.228943in}{2.299598in}}%
\pgfpathlineto{\pgfqpoint{3.229980in}{2.303271in}}%
\pgfpathlineto{\pgfqpoint{3.230117in}{2.303140in}}%
\pgfpathlineto{\pgfqpoint{3.230496in}{2.306405in}}%
\pgfpathlineto{\pgfqpoint{3.231086in}{2.302666in}}%
\pgfpathlineto{\pgfqpoint{3.231120in}{2.302687in}}%
\pgfpathlineto{\pgfqpoint{3.232269in}{2.300241in}}%
\pgfpathlineto{\pgfqpoint{3.231641in}{2.303839in}}%
\pgfpathlineto{\pgfqpoint{3.232303in}{2.301302in}}%
\pgfpathlineto{\pgfqpoint{3.233247in}{2.303996in}}%
\pgfpathlineto{\pgfqpoint{3.232629in}{2.301007in}}%
\pgfpathlineto{\pgfqpoint{3.233486in}{2.303626in}}%
\pgfpathlineto{\pgfqpoint{3.233666in}{2.302927in}}%
\pgfpathlineto{\pgfqpoint{3.234265in}{2.306857in}}%
\pgfpathlineto{\pgfqpoint{3.234508in}{2.305936in}}%
\pgfpathlineto{\pgfqpoint{3.234513in}{2.305994in}}%
\pgfpathlineto{\pgfqpoint{3.235160in}{2.300556in}}%
\pgfpathlineto{\pgfqpoint{3.235238in}{2.300022in}}%
\pgfpathlineto{\pgfqpoint{3.236115in}{2.304099in}}%
\pgfpathlineto{\pgfqpoint{3.236246in}{2.302382in}}%
\pgfpathlineto{\pgfqpoint{3.236305in}{2.301527in}}%
\pgfpathlineto{\pgfqpoint{3.237049in}{2.305824in}}%
\pgfpathlineto{\pgfqpoint{3.237332in}{2.302950in}}%
\pgfpathlineto{\pgfqpoint{3.237799in}{2.303823in}}%
\pgfpathlineto{\pgfqpoint{3.238106in}{2.301132in}}%
\pgfpathlineto{\pgfqpoint{3.238437in}{2.302597in}}%
\pgfpathlineto{\pgfqpoint{3.239021in}{2.300701in}}%
\pgfpathlineto{\pgfqpoint{3.239386in}{2.303598in}}%
\pgfpathlineto{\pgfqpoint{3.239542in}{2.302734in}}%
\pgfpathlineto{\pgfqpoint{3.240560in}{2.304728in}}%
\pgfpathlineto{\pgfqpoint{3.240010in}{2.302228in}}%
\pgfpathlineto{\pgfqpoint{3.240691in}{2.303451in}}%
\pgfpathlineto{\pgfqpoint{3.241139in}{2.304474in}}%
\pgfpathlineto{\pgfqpoint{3.241845in}{2.301799in}}%
\pgfpathlineto{\pgfqpoint{3.242166in}{2.304130in}}%
\pgfpathlineto{\pgfqpoint{3.242935in}{2.301424in}}%
\pgfpathlineto{\pgfqpoint{3.242940in}{2.301460in}}%
\pgfpathlineto{\pgfqpoint{3.243607in}{2.298548in}}%
\pgfpathlineto{\pgfqpoint{3.243968in}{2.303105in}}%
\pgfpathlineto{\pgfqpoint{3.244143in}{2.301384in}}%
\pgfpathlineto{\pgfqpoint{3.244484in}{2.304357in}}%
\pgfpathlineto{\pgfqpoint{3.245175in}{2.301739in}}%
\pgfpathlineto{\pgfqpoint{3.246178in}{2.305523in}}%
\pgfpathlineto{\pgfqpoint{3.245306in}{2.301275in}}%
\pgfpathlineto{\pgfqpoint{3.246304in}{2.303103in}}%
\pgfpathlineto{\pgfqpoint{3.246314in}{2.303462in}}%
\pgfpathlineto{\pgfqpoint{3.246543in}{2.301392in}}%
\pgfpathlineto{\pgfqpoint{3.246577in}{2.300815in}}%
\pgfpathlineto{\pgfqpoint{3.247536in}{2.303903in}}%
\pgfpathlineto{\pgfqpoint{3.247643in}{2.302109in}}%
\pgfpathlineto{\pgfqpoint{3.248860in}{2.304255in}}%
\pgfpathlineto{\pgfqpoint{3.247955in}{2.300564in}}%
\pgfpathlineto{\pgfqpoint{3.248914in}{2.303340in}}%
\pgfpathlineto{\pgfqpoint{3.250044in}{2.300154in}}%
\pgfpathlineto{\pgfqpoint{3.249143in}{2.304528in}}%
\pgfpathlineto{\pgfqpoint{3.250058in}{2.300306in}}%
\pgfpathlineto{\pgfqpoint{3.250112in}{2.299832in}}%
\pgfpathlineto{\pgfqpoint{3.250735in}{2.303565in}}%
\pgfpathlineto{\pgfqpoint{3.251061in}{2.302770in}}%
\pgfpathlineto{\pgfqpoint{3.251932in}{2.304895in}}%
\pgfpathlineto{\pgfqpoint{3.251119in}{2.302265in}}%
\pgfpathlineto{\pgfqpoint{3.252186in}{2.303475in}}%
\pgfpathlineto{\pgfqpoint{3.252838in}{2.302297in}}%
\pgfpathlineto{\pgfqpoint{3.252468in}{2.304811in}}%
\pgfpathlineto{\pgfqpoint{3.253305in}{2.302369in}}%
\pgfpathlineto{\pgfqpoint{3.254065in}{2.300795in}}%
\pgfpathlineto{\pgfqpoint{3.253758in}{2.303515in}}%
\pgfpathlineto{\pgfqpoint{3.254376in}{2.303273in}}%
\pgfpathlineto{\pgfqpoint{3.255287in}{2.304170in}}%
\pgfpathlineto{\pgfqpoint{3.255462in}{2.302471in}}%
\pgfpathlineto{\pgfqpoint{3.255477in}{2.302928in}}%
\pgfpathlineto{\pgfqpoint{3.256499in}{2.303860in}}%
\pgfpathlineto{\pgfqpoint{3.256051in}{2.300776in}}%
\pgfpathlineto{\pgfqpoint{3.256601in}{2.303042in}}%
\pgfpathlineto{\pgfqpoint{3.257560in}{2.300491in}}%
\pgfpathlineto{\pgfqpoint{3.256631in}{2.303464in}}%
\pgfpathlineto{\pgfqpoint{3.257828in}{2.301351in}}%
\pgfpathlineto{\pgfqpoint{3.258987in}{2.304252in}}%
\pgfpathlineto{\pgfqpoint{3.258013in}{2.301154in}}%
\pgfpathlineto{\pgfqpoint{3.259021in}{2.304112in}}%
\pgfpathlineto{\pgfqpoint{3.259732in}{2.300990in}}%
\pgfpathlineto{\pgfqpoint{3.259196in}{2.304228in}}%
\pgfpathlineto{\pgfqpoint{3.260160in}{2.302546in}}%
\pgfpathlineto{\pgfqpoint{3.260316in}{2.304645in}}%
\pgfpathlineto{\pgfqpoint{3.261071in}{2.300855in}}%
\pgfpathlineto{\pgfqpoint{3.261217in}{2.301942in}}%
\pgfpathlineto{\pgfqpoint{3.262210in}{2.298763in}}%
\pgfpathlineto{\pgfqpoint{3.261377in}{2.302191in}}%
\pgfpathlineto{\pgfqpoint{3.262395in}{2.299256in}}%
\pgfpathlineto{\pgfqpoint{3.263442in}{2.304111in}}%
\pgfpathlineto{\pgfqpoint{3.263573in}{2.303426in}}%
\pgfpathlineto{\pgfqpoint{3.264221in}{2.304563in}}%
\pgfpathlineto{\pgfqpoint{3.263967in}{2.301693in}}%
\pgfpathlineto{\pgfqpoint{3.264625in}{2.303053in}}%
\pgfpathlineto{\pgfqpoint{3.265715in}{2.300835in}}%
\pgfpathlineto{\pgfqpoint{3.264698in}{2.304113in}}%
\pgfpathlineto{\pgfqpoint{3.265749in}{2.301546in}}%
\pgfpathlineto{\pgfqpoint{3.266368in}{2.304367in}}%
\pgfpathlineto{\pgfqpoint{3.266889in}{2.302876in}}%
\pgfpathlineto{\pgfqpoint{3.268237in}{2.299825in}}%
\pgfpathlineto{\pgfqpoint{3.267366in}{2.304473in}}%
\pgfpathlineto{\pgfqpoint{3.268554in}{2.300516in}}%
\pgfpathlineto{\pgfqpoint{3.269722in}{2.304582in}}%
\pgfpathlineto{\pgfqpoint{3.268778in}{2.298242in}}%
\pgfpathlineto{\pgfqpoint{3.269727in}{2.304568in}}%
\pgfpathlineto{\pgfqpoint{3.270701in}{2.299507in}}%
\pgfpathlineto{\pgfqpoint{3.270973in}{2.300362in}}%
\pgfpathlineto{\pgfqpoint{3.270988in}{2.300148in}}%
\pgfpathlineto{\pgfqpoint{3.271173in}{2.302312in}}%
\pgfpathlineto{\pgfqpoint{3.272142in}{2.304102in}}%
\pgfpathlineto{\pgfqpoint{3.271557in}{2.302293in}}%
\pgfpathlineto{\pgfqpoint{3.272307in}{2.303591in}}%
\pgfpathlineto{\pgfqpoint{3.272385in}{2.302069in}}%
\pgfpathlineto{\pgfqpoint{3.273237in}{2.305450in}}%
\pgfpathlineto{\pgfqpoint{3.273412in}{2.303708in}}%
\pgfpathlineto{\pgfqpoint{3.273797in}{2.305567in}}%
\pgfpathlineto{\pgfqpoint{3.274493in}{2.301677in}}%
\pgfpathlineto{\pgfqpoint{3.275686in}{2.303927in}}%
\pgfpathlineto{\pgfqpoint{3.275019in}{2.300327in}}%
\pgfpathlineto{\pgfqpoint{3.275696in}{2.303713in}}%
\pgfpathlineto{\pgfqpoint{3.276757in}{2.297351in}}%
\pgfpathlineto{\pgfqpoint{3.276903in}{2.298642in}}%
\pgfpathlineto{\pgfqpoint{3.278179in}{2.305127in}}%
\pgfpathlineto{\pgfqpoint{3.278193in}{2.305078in}}%
\pgfpathlineto{\pgfqpoint{3.278563in}{2.306439in}}%
\pgfpathlineto{\pgfqpoint{3.279186in}{2.302613in}}%
\pgfpathlineto{\pgfqpoint{3.280282in}{2.300586in}}%
\pgfpathlineto{\pgfqpoint{3.279478in}{2.303798in}}%
\pgfpathlineto{\pgfqpoint{3.280316in}{2.301412in}}%
\pgfpathlineto{\pgfqpoint{3.281733in}{2.307183in}}%
\pgfpathlineto{\pgfqpoint{3.280890in}{2.300402in}}%
\pgfpathlineto{\pgfqpoint{3.281937in}{2.306041in}}%
\pgfpathlineto{\pgfqpoint{3.282278in}{2.302305in}}%
\pgfpathlineto{\pgfqpoint{3.282005in}{2.306799in}}%
\pgfpathlineto{\pgfqpoint{3.283101in}{2.303076in}}%
\pgfpathlineto{\pgfqpoint{3.283388in}{2.305508in}}%
\pgfpathlineto{\pgfqpoint{3.283144in}{2.302157in}}%
\pgfpathlineto{\pgfqpoint{3.284245in}{2.303998in}}%
\pgfpathlineto{\pgfqpoint{3.284668in}{2.302476in}}%
\pgfpathlineto{\pgfqpoint{3.285301in}{2.306872in}}%
\pgfpathlineto{\pgfqpoint{3.285472in}{2.307670in}}%
\pgfpathlineto{\pgfqpoint{3.285720in}{2.305729in}}%
\pgfpathlineto{\pgfqpoint{3.286679in}{2.302248in}}%
\pgfpathlineto{\pgfqpoint{3.285954in}{2.306293in}}%
\pgfpathlineto{\pgfqpoint{3.286849in}{2.304966in}}%
\pgfpathlineto{\pgfqpoint{3.287112in}{2.306543in}}%
\pgfpathlineto{\pgfqpoint{3.287438in}{2.302916in}}%
\pgfpathlineto{\pgfqpoint{3.287847in}{2.303807in}}%
\pgfpathlineto{\pgfqpoint{3.288232in}{2.302025in}}%
\pgfpathlineto{\pgfqpoint{3.287989in}{2.304848in}}%
\pgfpathlineto{\pgfqpoint{3.288962in}{2.303389in}}%
\pgfpathlineto{\pgfqpoint{3.289172in}{2.304201in}}%
\pgfpathlineto{\pgfqpoint{3.289829in}{2.301000in}}%
\pgfpathlineto{\pgfqpoint{3.289994in}{2.302508in}}%
\pgfpathlineto{\pgfqpoint{3.290364in}{2.303375in}}%
\pgfpathlineto{\pgfqpoint{3.291168in}{2.299688in}}%
\pgfpathlineto{\pgfqpoint{3.292190in}{2.301559in}}%
\pgfpathlineto{\pgfqpoint{3.291280in}{2.298764in}}%
\pgfpathlineto{\pgfqpoint{3.292297in}{2.300528in}}%
\pgfpathlineto{\pgfqpoint{3.292726in}{2.297765in}}%
\pgfpathlineto{\pgfqpoint{3.292341in}{2.301373in}}%
\pgfpathlineto{\pgfqpoint{3.293398in}{2.300669in}}%
\pgfpathlineto{\pgfqpoint{3.294337in}{2.303589in}}%
\pgfpathlineto{\pgfqpoint{3.294532in}{2.301383in}}%
\pgfpathlineto{\pgfqpoint{3.294775in}{2.300768in}}%
\pgfpathlineto{\pgfqpoint{3.295564in}{2.303484in}}%
\pgfpathlineto{\pgfqpoint{3.295569in}{2.303393in}}%
\pgfpathlineto{\pgfqpoint{3.295817in}{2.306164in}}%
\pgfpathlineto{\pgfqpoint{3.296655in}{2.302401in}}%
\pgfpathlineto{\pgfqpoint{3.297769in}{2.300784in}}%
\pgfpathlineto{\pgfqpoint{3.296942in}{2.303916in}}%
\pgfpathlineto{\pgfqpoint{3.297799in}{2.301404in}}%
\pgfpathlineto{\pgfqpoint{3.297823in}{2.301768in}}%
\pgfpathlineto{\pgfqpoint{3.298841in}{2.297968in}}%
\pgfpathlineto{\pgfqpoint{3.298850in}{2.298354in}}%
\pgfpathlineto{\pgfqpoint{3.298884in}{2.298700in}}%
\pgfpathlineto{\pgfqpoint{3.299546in}{2.295735in}}%
\pgfpathlineto{\pgfqpoint{3.299955in}{2.298346in}}%
\pgfpathlineto{\pgfqpoint{3.300121in}{2.296866in}}%
\pgfpathlineto{\pgfqpoint{3.300749in}{2.300844in}}%
\pgfpathlineto{\pgfqpoint{3.301065in}{2.298317in}}%
\pgfpathlineto{\pgfqpoint{3.301640in}{2.296911in}}%
\pgfpathlineto{\pgfqpoint{3.301377in}{2.299887in}}%
\pgfpathlineto{\pgfqpoint{3.302180in}{2.298249in}}%
\pgfpathlineto{\pgfqpoint{3.302769in}{2.300707in}}%
\pgfpathlineto{\pgfqpoint{3.302365in}{2.297077in}}%
\pgfpathlineto{\pgfqpoint{3.303320in}{2.299506in}}%
\pgfpathlineto{\pgfqpoint{3.304454in}{2.297354in}}%
\pgfpathlineto{\pgfqpoint{3.303412in}{2.300770in}}%
\pgfpathlineto{\pgfqpoint{3.304459in}{2.297495in}}%
\pgfpathlineto{\pgfqpoint{3.305350in}{2.298518in}}%
\pgfpathlineto{\pgfqpoint{3.305111in}{2.296351in}}%
\pgfpathlineto{\pgfqpoint{3.305598in}{2.297980in}}%
\pgfpathlineto{\pgfqpoint{3.305608in}{2.297886in}}%
\pgfpathlineto{\pgfqpoint{3.305798in}{2.301461in}}%
\pgfpathlineto{\pgfqpoint{3.306513in}{2.300615in}}%
\pgfpathlineto{\pgfqpoint{3.307414in}{2.301353in}}%
\pgfpathlineto{\pgfqpoint{3.306922in}{2.298116in}}%
\pgfpathlineto{\pgfqpoint{3.307623in}{2.300939in}}%
\pgfpathlineto{\pgfqpoint{3.308363in}{2.299577in}}%
\pgfpathlineto{\pgfqpoint{3.307974in}{2.302560in}}%
\pgfpathlineto{\pgfqpoint{3.308738in}{2.300327in}}%
\pgfpathlineto{\pgfqpoint{3.309098in}{2.297440in}}%
\pgfpathlineto{\pgfqpoint{3.309746in}{2.301951in}}%
\pgfpathlineto{\pgfqpoint{3.309824in}{2.301720in}}%
\pgfpathlineto{\pgfqpoint{3.310564in}{2.300506in}}%
\pgfpathlineto{\pgfqpoint{3.310875in}{2.303677in}}%
\pgfpathlineto{\pgfqpoint{3.310890in}{2.303402in}}%
\pgfpathlineto{\pgfqpoint{3.311528in}{2.305323in}}%
\pgfpathlineto{\pgfqpoint{3.311878in}{2.301298in}}%
\pgfpathlineto{\pgfqpoint{3.312054in}{2.304551in}}%
\pgfpathlineto{\pgfqpoint{3.313207in}{2.301687in}}%
\pgfpathlineto{\pgfqpoint{3.312102in}{2.305247in}}%
\pgfpathlineto{\pgfqpoint{3.313212in}{2.301795in}}%
\pgfpathlineto{\pgfqpoint{3.313748in}{2.302877in}}%
\pgfpathlineto{\pgfqpoint{3.314167in}{2.298897in}}%
\pgfpathlineto{\pgfqpoint{3.314215in}{2.299544in}}%
\pgfpathlineto{\pgfqpoint{3.315359in}{2.297408in}}%
\pgfpathlineto{\pgfqpoint{3.314751in}{2.301317in}}%
\pgfpathlineto{\pgfqpoint{3.315369in}{2.297996in}}%
\pgfpathlineto{\pgfqpoint{3.315393in}{2.298660in}}%
\pgfpathlineto{\pgfqpoint{3.316051in}{2.294645in}}%
\pgfpathlineto{\pgfqpoint{3.316460in}{2.297106in}}%
\pgfpathlineto{\pgfqpoint{3.317146in}{2.298956in}}%
\pgfpathlineto{\pgfqpoint{3.316864in}{2.296849in}}%
\pgfpathlineto{\pgfqpoint{3.317643in}{2.297612in}}%
\pgfpathlineto{\pgfqpoint{3.318699in}{2.295248in}}%
\pgfpathlineto{\pgfqpoint{3.318154in}{2.298617in}}%
\pgfpathlineto{\pgfqpoint{3.318792in}{2.296090in}}%
\pgfpathlineto{\pgfqpoint{3.318821in}{2.296481in}}%
\pgfpathlineto{\pgfqpoint{3.319819in}{2.292908in}}%
\pgfpathlineto{\pgfqpoint{3.319853in}{2.292958in}}%
\pgfpathlineto{\pgfqpoint{3.319921in}{2.292167in}}%
\pgfpathlineto{\pgfqpoint{3.320014in}{2.294141in}}%
\pgfpathlineto{\pgfqpoint{3.321167in}{2.300865in}}%
\pgfpathlineto{\pgfqpoint{3.321192in}{2.300852in}}%
\pgfpathlineto{\pgfqpoint{3.322774in}{2.303906in}}%
\pgfpathlineto{\pgfqpoint{3.321279in}{2.300372in}}%
\pgfpathlineto{\pgfqpoint{3.322852in}{2.303033in}}%
\pgfpathlineto{\pgfqpoint{3.323266in}{2.301040in}}%
\pgfpathlineto{\pgfqpoint{3.323865in}{2.304079in}}%
\pgfpathlineto{\pgfqpoint{3.323952in}{2.303344in}}%
\pgfpathlineto{\pgfqpoint{3.325335in}{2.297923in}}%
\pgfpathlineto{\pgfqpoint{3.324424in}{2.304363in}}%
\pgfpathlineto{\pgfqpoint{3.325452in}{2.298556in}}%
\pgfpathlineto{\pgfqpoint{3.325471in}{2.299409in}}%
\pgfpathlineto{\pgfqpoint{3.326143in}{2.293625in}}%
\pgfpathlineto{\pgfqpoint{3.326533in}{2.296782in}}%
\pgfpathlineto{\pgfqpoint{3.328665in}{2.292616in}}%
\pgfpathlineto{\pgfqpoint{3.328714in}{2.293280in}}%
\pgfpathlineto{\pgfqpoint{3.329970in}{2.299723in}}%
\pgfpathlineto{\pgfqpoint{3.329059in}{2.293126in}}%
\pgfpathlineto{\pgfqpoint{3.330043in}{2.298562in}}%
\pgfpathlineto{\pgfqpoint{3.331479in}{2.301472in}}%
\pgfpathlineto{\pgfqpoint{3.330155in}{2.297024in}}%
\pgfpathlineto{\pgfqpoint{3.331669in}{2.300677in}}%
\pgfpathlineto{\pgfqpoint{3.332657in}{2.299996in}}%
\pgfpathlineto{\pgfqpoint{3.332073in}{2.302593in}}%
\pgfpathlineto{\pgfqpoint{3.332740in}{2.300780in}}%
\pgfpathlineto{\pgfqpoint{3.333426in}{2.302074in}}%
\pgfpathlineto{\pgfqpoint{3.333694in}{2.299397in}}%
\pgfpathlineto{\pgfqpoint{3.333850in}{2.300894in}}%
\pgfpathlineto{\pgfqpoint{3.334127in}{2.301722in}}%
\pgfpathlineto{\pgfqpoint{3.334726in}{2.298918in}}%
\pgfpathlineto{\pgfqpoint{3.334877in}{2.299752in}}%
\pgfpathlineto{\pgfqpoint{3.334931in}{2.298913in}}%
\pgfpathlineto{\pgfqpoint{3.335744in}{2.305405in}}%
\pgfpathlineto{\pgfqpoint{3.335846in}{2.304814in}}%
\pgfpathlineto{\pgfqpoint{3.335856in}{2.305539in}}%
\pgfpathlineto{\pgfqpoint{3.336873in}{2.301581in}}%
\pgfpathlineto{\pgfqpoint{3.336922in}{2.301903in}}%
\pgfpathlineto{\pgfqpoint{3.338090in}{2.298748in}}%
\pgfpathlineto{\pgfqpoint{3.337681in}{2.302473in}}%
\pgfpathlineto{\pgfqpoint{3.338134in}{2.299201in}}%
\pgfpathlineto{\pgfqpoint{3.338889in}{2.302125in}}%
\pgfpathlineto{\pgfqpoint{3.338485in}{2.298435in}}%
\pgfpathlineto{\pgfqpoint{3.339458in}{2.300567in}}%
\pgfpathlineto{\pgfqpoint{3.339639in}{2.299291in}}%
\pgfpathlineto{\pgfqpoint{3.340398in}{2.303906in}}%
\pgfpathlineto{\pgfqpoint{3.340525in}{2.302972in}}%
\pgfpathlineto{\pgfqpoint{3.340646in}{2.304983in}}%
\pgfpathlineto{\pgfqpoint{3.341089in}{2.301562in}}%
\pgfpathlineto{\pgfqpoint{3.341640in}{2.303210in}}%
\pgfpathlineto{\pgfqpoint{3.342258in}{2.304148in}}%
\pgfpathlineto{\pgfqpoint{3.341951in}{2.300249in}}%
\pgfpathlineto{\pgfqpoint{3.342691in}{2.302780in}}%
\pgfpathlineto{\pgfqpoint{3.343665in}{2.296133in}}%
\pgfpathlineto{\pgfqpoint{3.343967in}{2.296487in}}%
\pgfpathlineto{\pgfqpoint{3.344400in}{2.297384in}}%
\pgfpathlineto{\pgfqpoint{3.344697in}{2.294597in}}%
\pgfpathlineto{\pgfqpoint{3.345052in}{2.296220in}}%
\pgfpathlineto{\pgfqpoint{3.346182in}{2.292493in}}%
\pgfpathlineto{\pgfqpoint{3.346231in}{2.293254in}}%
\pgfpathlineto{\pgfqpoint{3.346917in}{2.295014in}}%
\pgfpathlineto{\pgfqpoint{3.346435in}{2.292381in}}%
\pgfpathlineto{\pgfqpoint{3.347365in}{2.293659in}}%
\pgfpathlineto{\pgfqpoint{3.347793in}{2.289952in}}%
\pgfpathlineto{\pgfqpoint{3.348134in}{2.296249in}}%
\pgfpathlineto{\pgfqpoint{3.348382in}{2.295858in}}%
\pgfpathlineto{\pgfqpoint{3.348572in}{2.296519in}}%
\pgfpathlineto{\pgfqpoint{3.349449in}{2.293881in}}%
\pgfpathlineto{\pgfqpoint{3.350851in}{2.301069in}}%
\pgfpathlineto{\pgfqpoint{3.349638in}{2.293154in}}%
\pgfpathlineto{\pgfqpoint{3.351118in}{2.299542in}}%
\pgfpathlineto{\pgfqpoint{3.352141in}{2.295389in}}%
\pgfpathlineto{\pgfqpoint{3.351226in}{2.300813in}}%
\pgfpathlineto{\pgfqpoint{3.352282in}{2.296909in}}%
\pgfpathlineto{\pgfqpoint{3.352579in}{2.295578in}}%
\pgfpathlineto{\pgfqpoint{3.353251in}{2.299274in}}%
\pgfpathlineto{\pgfqpoint{3.353295in}{2.298933in}}%
\pgfpathlineto{\pgfqpoint{3.354410in}{2.300183in}}%
\pgfpathlineto{\pgfqpoint{3.353894in}{2.295136in}}%
\pgfpathlineto{\pgfqpoint{3.354449in}{2.299781in}}%
\pgfpathlineto{\pgfqpoint{3.354468in}{2.299180in}}%
\pgfpathlineto{\pgfqpoint{3.354819in}{2.302241in}}%
\pgfpathlineto{\pgfqpoint{3.355330in}{2.302143in}}%
\pgfpathlineto{\pgfqpoint{3.356065in}{2.300069in}}%
\pgfpathlineto{\pgfqpoint{3.356498in}{2.304656in}}%
\pgfpathlineto{\pgfqpoint{3.356523in}{2.304112in}}%
\pgfpathlineto{\pgfqpoint{3.356664in}{2.303798in}}%
\pgfpathlineto{\pgfqpoint{3.357443in}{2.300189in}}%
\pgfpathlineto{\pgfqpoint{3.356824in}{2.305315in}}%
\pgfpathlineto{\pgfqpoint{3.357813in}{2.301102in}}%
\pgfpathlineto{\pgfqpoint{3.358105in}{2.302942in}}%
\pgfpathlineto{\pgfqpoint{3.358533in}{2.298762in}}%
\pgfpathlineto{\pgfqpoint{3.358874in}{2.300052in}}%
\pgfpathlineto{\pgfqpoint{3.359560in}{2.299209in}}%
\pgfpathlineto{\pgfqpoint{3.359190in}{2.302735in}}%
\pgfpathlineto{\pgfqpoint{3.359843in}{2.301318in}}%
\pgfpathlineto{\pgfqpoint{3.361055in}{2.305460in}}%
\pgfpathlineto{\pgfqpoint{3.361065in}{2.305126in}}%
\pgfpathlineto{\pgfqpoint{3.362384in}{2.298647in}}%
\pgfpathlineto{\pgfqpoint{3.362467in}{2.299077in}}%
\pgfpathlineto{\pgfqpoint{3.363743in}{2.303741in}}%
\pgfpathlineto{\pgfqpoint{3.362793in}{2.297398in}}%
\pgfpathlineto{\pgfqpoint{3.363791in}{2.302946in}}%
\pgfpathlineto{\pgfqpoint{3.364181in}{2.299255in}}%
\pgfpathlineto{\pgfqpoint{3.364677in}{2.303168in}}%
\pgfpathlineto{\pgfqpoint{3.364921in}{2.302510in}}%
\pgfpathlineto{\pgfqpoint{3.365466in}{2.303993in}}%
\pgfpathlineto{\pgfqpoint{3.365213in}{2.300513in}}%
\pgfpathlineto{\pgfqpoint{3.366016in}{2.302356in}}%
\pgfpathlineto{\pgfqpoint{3.366980in}{2.298358in}}%
\pgfpathlineto{\pgfqpoint{3.367180in}{2.299521in}}%
\pgfpathlineto{\pgfqpoint{3.367321in}{2.299907in}}%
\pgfpathlineto{\pgfqpoint{3.367540in}{2.297484in}}%
\pgfpathlineto{\pgfqpoint{3.368207in}{2.298239in}}%
\pgfpathlineto{\pgfqpoint{3.369288in}{2.292162in}}%
\pgfpathlineto{\pgfqpoint{3.368217in}{2.298269in}}%
\pgfpathlineto{\pgfqpoint{3.369458in}{2.293224in}}%
\pgfpathlineto{\pgfqpoint{3.370573in}{2.292231in}}%
\pgfpathlineto{\pgfqpoint{3.370208in}{2.295079in}}%
\pgfpathlineto{\pgfqpoint{3.370602in}{2.292429in}}%
\pgfpathlineto{\pgfqpoint{3.371279in}{2.290410in}}%
\pgfpathlineto{\pgfqpoint{3.370997in}{2.293476in}}%
\pgfpathlineto{\pgfqpoint{3.371625in}{2.293114in}}%
\pgfpathlineto{\pgfqpoint{3.372735in}{2.293475in}}%
\pgfpathlineto{\pgfqpoint{3.371858in}{2.291580in}}%
\pgfpathlineto{\pgfqpoint{3.372739in}{2.293389in}}%
\pgfpathlineto{\pgfqpoint{3.373022in}{2.292094in}}%
\pgfpathlineto{\pgfqpoint{3.373577in}{2.298064in}}%
\pgfpathlineto{\pgfqpoint{3.373674in}{2.297266in}}%
\pgfpathlineto{\pgfqpoint{3.373854in}{2.298961in}}%
\pgfpathlineto{\pgfqpoint{3.374127in}{2.296231in}}%
\pgfpathlineto{\pgfqpoint{3.374784in}{2.297695in}}%
\pgfpathlineto{\pgfqpoint{3.375724in}{2.295512in}}%
\pgfpathlineto{\pgfqpoint{3.375305in}{2.297808in}}%
\pgfpathlineto{\pgfqpoint{3.375919in}{2.296859in}}%
\pgfpathlineto{\pgfqpoint{3.376143in}{2.299506in}}%
\pgfpathlineto{\pgfqpoint{3.376556in}{2.295976in}}%
\pgfpathlineto{\pgfqpoint{3.377043in}{2.297868in}}%
\pgfpathlineto{\pgfqpoint{3.378046in}{2.297106in}}%
\pgfpathlineto{\pgfqpoint{3.377472in}{2.301534in}}%
\pgfpathlineto{\pgfqpoint{3.378144in}{2.297755in}}%
\pgfpathlineto{\pgfqpoint{3.378776in}{2.301329in}}%
\pgfpathlineto{\pgfqpoint{3.379273in}{2.300017in}}%
\pgfpathlineto{\pgfqpoint{3.380208in}{2.298488in}}%
\pgfpathlineto{\pgfqpoint{3.380023in}{2.300355in}}%
\pgfpathlineto{\pgfqpoint{3.380383in}{2.299808in}}%
\pgfpathlineto{\pgfqpoint{3.381313in}{2.301730in}}%
\pgfpathlineto{\pgfqpoint{3.380797in}{2.298827in}}%
\pgfpathlineto{\pgfqpoint{3.381513in}{2.300981in}}%
\pgfpathlineto{\pgfqpoint{3.381600in}{2.300417in}}%
\pgfpathlineto{\pgfqpoint{3.382019in}{2.304005in}}%
\pgfpathlineto{\pgfqpoint{3.382569in}{2.303288in}}%
\pgfpathlineto{\pgfqpoint{3.383815in}{2.300557in}}%
\pgfpathlineto{\pgfqpoint{3.382876in}{2.304168in}}%
\pgfpathlineto{\pgfqpoint{3.383820in}{2.300720in}}%
\pgfpathlineto{\pgfqpoint{3.384843in}{2.303527in}}%
\pgfpathlineto{\pgfqpoint{3.385018in}{2.302259in}}%
\pgfpathlineto{\pgfqpoint{3.385763in}{2.300404in}}%
\pgfpathlineto{\pgfqpoint{3.385232in}{2.303233in}}%
\pgfpathlineto{\pgfqpoint{3.386172in}{2.300589in}}%
\pgfpathlineto{\pgfqpoint{3.387223in}{2.301790in}}%
\pgfpathlineto{\pgfqpoint{3.386693in}{2.299768in}}%
\pgfpathlineto{\pgfqpoint{3.387282in}{2.300657in}}%
\pgfpathlineto{\pgfqpoint{3.387321in}{2.300501in}}%
\pgfpathlineto{\pgfqpoint{3.387939in}{2.304908in}}%
\pgfpathlineto{\pgfqpoint{3.388309in}{2.301802in}}%
\pgfpathlineto{\pgfqpoint{3.388854in}{2.300279in}}%
\pgfpathlineto{\pgfqpoint{3.389443in}{2.303611in}}%
\pgfpathlineto{\pgfqpoint{3.389711in}{2.302731in}}%
\pgfpathlineto{\pgfqpoint{3.389833in}{2.304553in}}%
\pgfpathlineto{\pgfqpoint{3.390558in}{2.303253in}}%
\pgfpathlineto{\pgfqpoint{3.391279in}{2.302820in}}%
\pgfpathlineto{\pgfqpoint{3.390777in}{2.305448in}}%
\pgfpathlineto{\pgfqpoint{3.391512in}{2.304276in}}%
\pgfpathlineto{\pgfqpoint{3.391863in}{2.304723in}}%
\pgfpathlineto{\pgfqpoint{3.392170in}{2.301612in}}%
\pgfpathlineto{\pgfqpoint{3.392433in}{2.302645in}}%
\pgfpathlineto{\pgfqpoint{3.393002in}{2.304872in}}%
\pgfpathlineto{\pgfqpoint{3.393007in}{2.304793in}}%
\pgfpathlineto{\pgfqpoint{3.393065in}{2.306279in}}%
\pgfpathlineto{\pgfqpoint{3.393742in}{2.301526in}}%
\pgfpathlineto{\pgfqpoint{3.394083in}{2.303221in}}%
\pgfpathlineto{\pgfqpoint{3.394346in}{2.304078in}}%
\pgfpathlineto{\pgfqpoint{3.394536in}{2.301183in}}%
\pgfpathlineto{\pgfqpoint{3.394555in}{2.301220in}}%
\pgfpathlineto{\pgfqpoint{3.394560in}{2.300887in}}%
\pgfpathlineto{\pgfqpoint{3.395466in}{2.304496in}}%
\pgfpathlineto{\pgfqpoint{3.395626in}{2.303226in}}%
\pgfpathlineto{\pgfqpoint{3.395694in}{2.304362in}}%
\pgfpathlineto{\pgfqpoint{3.395855in}{2.301041in}}%
\pgfpathlineto{\pgfqpoint{3.396722in}{2.302182in}}%
\pgfpathlineto{\pgfqpoint{3.397895in}{2.297376in}}%
\pgfpathlineto{\pgfqpoint{3.397910in}{2.297565in}}%
\pgfpathlineto{\pgfqpoint{3.398177in}{2.296092in}}%
\pgfpathlineto{\pgfqpoint{3.398572in}{2.298783in}}%
\pgfpathlineto{\pgfqpoint{3.398601in}{2.298700in}}%
\pgfpathlineto{\pgfqpoint{3.398606in}{2.298771in}}%
\pgfpathlineto{\pgfqpoint{3.399395in}{2.294802in}}%
\pgfpathlineto{\pgfqpoint{3.399492in}{2.295698in}}%
\pgfpathlineto{\pgfqpoint{3.399726in}{2.293916in}}%
\pgfpathlineto{\pgfqpoint{3.400568in}{2.296748in}}%
\pgfpathlineto{\pgfqpoint{3.400587in}{2.296293in}}%
\pgfpathlineto{\pgfqpoint{3.400753in}{2.297462in}}%
\pgfpathlineto{\pgfqpoint{3.401464in}{2.294298in}}%
\pgfpathlineto{\pgfqpoint{3.401493in}{2.294507in}}%
\pgfpathlineto{\pgfqpoint{3.401561in}{2.293496in}}%
\pgfpathlineto{\pgfqpoint{3.402418in}{2.295591in}}%
\pgfpathlineto{\pgfqpoint{3.402627in}{2.293541in}}%
\pgfpathlineto{\pgfqpoint{3.403508in}{2.294549in}}%
\pgfpathlineto{\pgfqpoint{3.403109in}{2.292351in}}%
\pgfpathlineto{\pgfqpoint{3.403693in}{2.292485in}}%
\pgfpathlineto{\pgfqpoint{3.404185in}{2.288818in}}%
\pgfpathlineto{\pgfqpoint{3.404818in}{2.291328in}}%
\pgfpathlineto{\pgfqpoint{3.405772in}{2.294299in}}%
\pgfpathlineto{\pgfqpoint{3.406094in}{2.293400in}}%
\pgfpathlineto{\pgfqpoint{3.406123in}{2.292632in}}%
\pgfpathlineto{\pgfqpoint{3.406595in}{2.295745in}}%
\pgfpathlineto{\pgfqpoint{3.407184in}{2.294176in}}%
\pgfpathlineto{\pgfqpoint{3.407632in}{2.292185in}}%
\pgfpathlineto{\pgfqpoint{3.407247in}{2.294946in}}%
\pgfpathlineto{\pgfqpoint{3.408143in}{2.294799in}}%
\pgfpathlineto{\pgfqpoint{3.408465in}{2.294596in}}%
\pgfpathlineto{\pgfqpoint{3.409278in}{2.296154in}}%
\pgfpathlineto{\pgfqpoint{3.410422in}{2.294371in}}%
\pgfpathlineto{\pgfqpoint{3.409594in}{2.296560in}}%
\pgfpathlineto{\pgfqpoint{3.410456in}{2.295090in}}%
\pgfpathlineto{\pgfqpoint{3.410685in}{2.296021in}}%
\pgfpathlineto{\pgfqpoint{3.411196in}{2.293292in}}%
\pgfpathlineto{\pgfqpoint{3.411561in}{2.294942in}}%
\pgfpathlineto{\pgfqpoint{3.411712in}{2.294471in}}%
\pgfpathlineto{\pgfqpoint{3.412535in}{2.298346in}}%
\pgfpathlineto{\pgfqpoint{3.412544in}{2.298275in}}%
\pgfpathlineto{\pgfqpoint{3.412978in}{2.300579in}}%
\pgfpathlineto{\pgfqpoint{3.412690in}{2.297767in}}%
\pgfpathlineto{\pgfqpoint{3.413664in}{2.299378in}}%
\pgfpathlineto{\pgfqpoint{3.414326in}{2.298651in}}%
\pgfpathlineto{\pgfqpoint{3.413995in}{2.302300in}}%
\pgfpathlineto{\pgfqpoint{3.414638in}{2.301727in}}%
\pgfpathlineto{\pgfqpoint{3.415787in}{2.305260in}}%
\pgfpathlineto{\pgfqpoint{3.414813in}{2.301140in}}%
\pgfpathlineto{\pgfqpoint{3.415874in}{2.304182in}}%
\pgfpathlineto{\pgfqpoint{3.416634in}{2.304579in}}%
\pgfpathlineto{\pgfqpoint{3.417028in}{2.301406in}}%
\pgfpathlineto{\pgfqpoint{3.417505in}{2.300788in}}%
\pgfpathlineto{\pgfqpoint{3.417676in}{2.303564in}}%
\pgfpathlineto{\pgfqpoint{3.418056in}{2.302160in}}%
\pgfpathlineto{\pgfqpoint{3.418771in}{2.304575in}}%
\pgfpathlineto{\pgfqpoint{3.418211in}{2.301149in}}%
\pgfpathlineto{\pgfqpoint{3.419185in}{2.302974in}}%
\pgfpathlineto{\pgfqpoint{3.419920in}{2.300661in}}%
\pgfpathlineto{\pgfqpoint{3.419419in}{2.303293in}}%
\pgfpathlineto{\pgfqpoint{3.420339in}{2.301092in}}%
\pgfpathlineto{\pgfqpoint{3.420904in}{2.303556in}}%
\pgfpathlineto{\pgfqpoint{3.420373in}{2.300427in}}%
\pgfpathlineto{\pgfqpoint{3.421463in}{2.301587in}}%
\pgfpathlineto{\pgfqpoint{3.421483in}{2.301144in}}%
\pgfpathlineto{\pgfqpoint{3.421843in}{2.304872in}}%
\pgfpathlineto{\pgfqpoint{3.422559in}{2.302848in}}%
\pgfpathlineto{\pgfqpoint{3.422598in}{2.302480in}}%
\pgfpathlineto{\pgfqpoint{3.422807in}{2.305453in}}%
\pgfpathlineto{\pgfqpoint{3.423601in}{2.303721in}}%
\pgfpathlineto{\pgfqpoint{3.424248in}{2.305397in}}%
\pgfpathlineto{\pgfqpoint{3.424531in}{2.301810in}}%
\pgfpathlineto{\pgfqpoint{3.424682in}{2.302917in}}%
\pgfpathlineto{\pgfqpoint{3.425334in}{2.302402in}}%
\pgfpathlineto{\pgfqpoint{3.424891in}{2.305631in}}%
\pgfpathlineto{\pgfqpoint{3.425748in}{2.303558in}}%
\pgfpathlineto{\pgfqpoint{3.425908in}{2.304556in}}%
\pgfpathlineto{\pgfqpoint{3.426580in}{2.300430in}}%
\pgfpathlineto{\pgfqpoint{3.426644in}{2.300611in}}%
\pgfpathlineto{\pgfqpoint{3.427155in}{2.295309in}}%
\pgfpathlineto{\pgfqpoint{3.427861in}{2.296109in}}%
\pgfpathlineto{\pgfqpoint{3.428722in}{2.303319in}}%
\pgfpathlineto{\pgfqpoint{3.428041in}{2.295450in}}%
\pgfpathlineto{\pgfqpoint{3.429136in}{2.302651in}}%
\pgfpathlineto{\pgfqpoint{3.430086in}{2.301443in}}%
\pgfpathlineto{\pgfqpoint{3.429730in}{2.304036in}}%
\pgfpathlineto{\pgfqpoint{3.430232in}{2.302893in}}%
\pgfpathlineto{\pgfqpoint{3.431137in}{2.304186in}}%
\pgfpathlineto{\pgfqpoint{3.430796in}{2.301994in}}%
\pgfpathlineto{\pgfqpoint{3.431337in}{2.302643in}}%
\pgfpathlineto{\pgfqpoint{3.431395in}{2.300916in}}%
\pgfpathlineto{\pgfqpoint{3.432398in}{2.303966in}}%
\pgfpathlineto{\pgfqpoint{3.432929in}{2.304924in}}%
\pgfpathlineto{\pgfqpoint{3.433265in}{2.301504in}}%
\pgfpathlineto{\pgfqpoint{3.433484in}{2.302947in}}%
\pgfpathlineto{\pgfqpoint{3.433688in}{2.301580in}}%
\pgfpathlineto{\pgfqpoint{3.434409in}{2.305748in}}%
\pgfpathlineto{\pgfqpoint{3.434574in}{2.304109in}}%
\pgfpathlineto{\pgfqpoint{3.434939in}{2.304907in}}%
\pgfpathlineto{\pgfqpoint{3.435568in}{2.301730in}}%
\pgfpathlineto{\pgfqpoint{3.435626in}{2.301983in}}%
\pgfpathlineto{\pgfqpoint{3.436181in}{2.301056in}}%
\pgfpathlineto{\pgfqpoint{3.436585in}{2.304627in}}%
\pgfpathlineto{\pgfqpoint{3.436614in}{2.303980in}}%
\pgfpathlineto{\pgfqpoint{3.437593in}{2.305116in}}%
\pgfpathlineto{\pgfqpoint{3.436833in}{2.301656in}}%
\pgfpathlineto{\pgfqpoint{3.437719in}{2.303760in}}%
\pgfpathlineto{\pgfqpoint{3.437875in}{2.305599in}}%
\pgfpathlineto{\pgfqpoint{3.438742in}{2.302353in}}%
\pgfpathlineto{\pgfqpoint{3.438761in}{2.302464in}}%
\pgfpathlineto{\pgfqpoint{3.439492in}{2.300315in}}%
\pgfpathlineto{\pgfqpoint{3.438980in}{2.303739in}}%
\pgfpathlineto{\pgfqpoint{3.439847in}{2.302855in}}%
\pgfpathlineto{\pgfqpoint{3.440996in}{2.304179in}}%
\pgfpathlineto{\pgfqpoint{3.440509in}{2.300375in}}%
\pgfpathlineto{\pgfqpoint{3.441001in}{2.304052in}}%
\pgfpathlineto{\pgfqpoint{3.441677in}{2.298099in}}%
\pgfpathlineto{\pgfqpoint{3.442427in}{2.301662in}}%
\pgfpathlineto{\pgfqpoint{3.442788in}{2.303815in}}%
\pgfpathlineto{\pgfqpoint{3.443430in}{2.299425in}}%
\pgfpathlineto{\pgfqpoint{3.443464in}{2.299847in}}%
\pgfpathlineto{\pgfqpoint{3.444788in}{2.295444in}}%
\pgfpathlineto{\pgfqpoint{3.443489in}{2.300247in}}%
\pgfpathlineto{\pgfqpoint{3.444925in}{2.296023in}}%
\pgfpathlineto{\pgfqpoint{3.445553in}{2.297669in}}%
\pgfpathlineto{\pgfqpoint{3.445860in}{2.295244in}}%
\pgfpathlineto{\pgfqpoint{3.446006in}{2.295778in}}%
\pgfpathlineto{\pgfqpoint{3.446629in}{2.294103in}}%
\pgfpathlineto{\pgfqpoint{3.446424in}{2.296815in}}%
\pgfpathlineto{\pgfqpoint{3.447116in}{2.295563in}}%
\pgfpathlineto{\pgfqpoint{3.447349in}{2.297320in}}%
\pgfpathlineto{\pgfqpoint{3.448143in}{2.294663in}}%
\pgfpathlineto{\pgfqpoint{3.448216in}{2.294834in}}%
\pgfpathlineto{\pgfqpoint{3.448401in}{2.293361in}}%
\pgfpathlineto{\pgfqpoint{3.449272in}{2.296161in}}%
\pgfpathlineto{\pgfqpoint{3.449306in}{2.295609in}}%
\pgfpathlineto{\pgfqpoint{3.449336in}{2.295483in}}%
\pgfpathlineto{\pgfqpoint{3.449448in}{2.297193in}}%
\pgfpathlineto{\pgfqpoint{3.450470in}{2.299936in}}%
\pgfpathlineto{\pgfqpoint{3.450236in}{2.295627in}}%
\pgfpathlineto{\pgfqpoint{3.450679in}{2.299487in}}%
\pgfpathlineto{\pgfqpoint{3.451926in}{2.296191in}}%
\pgfpathlineto{\pgfqpoint{3.450718in}{2.300536in}}%
\pgfpathlineto{\pgfqpoint{3.452023in}{2.297450in}}%
\pgfpathlineto{\pgfqpoint{3.452028in}{2.297876in}}%
\pgfpathlineto{\pgfqpoint{3.453050in}{2.293766in}}%
\pgfpathlineto{\pgfqpoint{3.453104in}{2.294962in}}%
\pgfpathlineto{\pgfqpoint{3.453279in}{2.293758in}}%
\pgfpathlineto{\pgfqpoint{3.453966in}{2.298675in}}%
\pgfpathlineto{\pgfqpoint{3.454175in}{2.297063in}}%
\pgfpathlineto{\pgfqpoint{3.454934in}{2.300247in}}%
\pgfpathlineto{\pgfqpoint{3.454477in}{2.296193in}}%
\pgfpathlineto{\pgfqpoint{3.455290in}{2.297643in}}%
\pgfpathlineto{\pgfqpoint{3.456361in}{2.295178in}}%
\pgfpathlineto{\pgfqpoint{3.455820in}{2.298078in}}%
\pgfpathlineto{\pgfqpoint{3.456526in}{2.296400in}}%
\pgfpathlineto{\pgfqpoint{3.457554in}{2.301177in}}%
\pgfpathlineto{\pgfqpoint{3.456580in}{2.296144in}}%
\pgfpathlineto{\pgfqpoint{3.457695in}{2.300249in}}%
\pgfpathlineto{\pgfqpoint{3.458873in}{2.297201in}}%
\pgfpathlineto{\pgfqpoint{3.459102in}{2.298401in}}%
\pgfpathlineto{\pgfqpoint{3.459676in}{2.294072in}}%
\pgfpathlineto{\pgfqpoint{3.459681in}{2.294124in}}%
\pgfpathlineto{\pgfqpoint{3.460032in}{2.291911in}}%
\pgfpathlineto{\pgfqpoint{3.460670in}{2.295867in}}%
\pgfpathlineto{\pgfqpoint{3.460757in}{2.295744in}}%
\pgfpathlineto{\pgfqpoint{3.461857in}{2.296879in}}%
\pgfpathlineto{\pgfqpoint{3.461448in}{2.292931in}}%
\pgfpathlineto{\pgfqpoint{3.461901in}{2.296864in}}%
\pgfpathlineto{\pgfqpoint{3.461930in}{2.296925in}}%
\pgfpathlineto{\pgfqpoint{3.462558in}{2.292274in}}%
\pgfpathlineto{\pgfqpoint{3.462563in}{2.292421in}}%
\pgfpathlineto{\pgfqpoint{3.462894in}{2.290823in}}%
\pgfpathlineto{\pgfqpoint{3.463362in}{2.295063in}}%
\pgfpathlineto{\pgfqpoint{3.463610in}{2.294689in}}%
\pgfpathlineto{\pgfqpoint{3.464783in}{2.300092in}}%
\pgfpathlineto{\pgfqpoint{3.463917in}{2.294086in}}%
\pgfpathlineto{\pgfqpoint{3.464793in}{2.299935in}}%
\pgfpathlineto{\pgfqpoint{3.464808in}{2.299349in}}%
\pgfpathlineto{\pgfqpoint{3.465577in}{2.302299in}}%
\pgfpathlineto{\pgfqpoint{3.465884in}{2.300544in}}%
\pgfpathlineto{\pgfqpoint{3.466736in}{2.301758in}}%
\pgfpathlineto{\pgfqpoint{3.466448in}{2.299935in}}%
\pgfpathlineto{\pgfqpoint{3.466930in}{2.299906in}}%
\pgfpathlineto{\pgfqpoint{3.468104in}{2.294167in}}%
\pgfpathlineto{\pgfqpoint{3.466984in}{2.300120in}}%
\pgfpathlineto{\pgfqpoint{3.468157in}{2.295213in}}%
\pgfpathlineto{\pgfqpoint{3.469311in}{2.297170in}}%
\pgfpathlineto{\pgfqpoint{3.469024in}{2.294091in}}%
\pgfpathlineto{\pgfqpoint{3.469316in}{2.296963in}}%
\pgfpathlineto{\pgfqpoint{3.470372in}{2.294810in}}%
\pgfpathlineto{\pgfqpoint{3.469511in}{2.298342in}}%
\pgfpathlineto{\pgfqpoint{3.470480in}{2.295867in}}%
\pgfpathlineto{\pgfqpoint{3.470518in}{2.296880in}}%
\pgfpathlineto{\pgfqpoint{3.471478in}{2.291444in}}%
\pgfpathlineto{\pgfqpoint{3.471497in}{2.291620in}}%
\pgfpathlineto{\pgfqpoint{3.471541in}{2.290711in}}%
\pgfpathlineto{\pgfqpoint{3.472320in}{2.295091in}}%
\pgfpathlineto{\pgfqpoint{3.472558in}{2.293532in}}%
\pgfpathlineto{\pgfqpoint{3.473259in}{2.292674in}}%
\pgfpathlineto{\pgfqpoint{3.473688in}{2.295491in}}%
\pgfpathlineto{\pgfqpoint{3.473756in}{2.294037in}}%
\pgfpathlineto{\pgfqpoint{3.474905in}{2.291929in}}%
\pgfpathlineto{\pgfqpoint{3.473810in}{2.295000in}}%
\pgfpathlineto{\pgfqpoint{3.474929in}{2.292158in}}%
\pgfpathlineto{\pgfqpoint{3.476638in}{2.297410in}}%
\pgfpathlineto{\pgfqpoint{3.474968in}{2.292024in}}%
\pgfpathlineto{\pgfqpoint{3.476677in}{2.296837in}}%
\pgfpathlineto{\pgfqpoint{3.477330in}{2.295105in}}%
\pgfpathlineto{\pgfqpoint{3.476852in}{2.297426in}}%
\pgfpathlineto{\pgfqpoint{3.477812in}{2.295938in}}%
\pgfpathlineto{\pgfqpoint{3.477889in}{2.296168in}}%
\pgfpathlineto{\pgfqpoint{3.478133in}{2.292854in}}%
\pgfpathlineto{\pgfqpoint{3.478863in}{2.294851in}}%
\pgfpathlineto{\pgfqpoint{3.479851in}{2.293269in}}%
\pgfpathlineto{\pgfqpoint{3.479043in}{2.295546in}}%
\pgfpathlineto{\pgfqpoint{3.480007in}{2.294907in}}%
\pgfpathlineto{\pgfqpoint{3.480996in}{2.302887in}}%
\pgfpathlineto{\pgfqpoint{3.480027in}{2.294776in}}%
\pgfpathlineto{\pgfqpoint{3.481210in}{2.302749in}}%
\pgfpathlineto{\pgfqpoint{3.482412in}{2.298513in}}%
\pgfpathlineto{\pgfqpoint{3.481585in}{2.303713in}}%
\pgfpathlineto{\pgfqpoint{3.482441in}{2.298741in}}%
\pgfpathlineto{\pgfqpoint{3.483561in}{2.300192in}}%
\pgfpathlineto{\pgfqpoint{3.483298in}{2.297263in}}%
\pgfpathlineto{\pgfqpoint{3.483566in}{2.300171in}}%
\pgfpathlineto{\pgfqpoint{3.483625in}{2.299047in}}%
\pgfpathlineto{\pgfqpoint{3.484408in}{2.303595in}}%
\pgfpathlineto{\pgfqpoint{3.484593in}{2.302718in}}%
\pgfpathlineto{\pgfqpoint{3.485630in}{2.303884in}}%
\pgfpathlineto{\pgfqpoint{3.485075in}{2.301399in}}%
\pgfpathlineto{\pgfqpoint{3.485737in}{2.303304in}}%
\pgfpathlineto{\pgfqpoint{3.486604in}{2.302440in}}%
\pgfpathlineto{\pgfqpoint{3.485981in}{2.304129in}}%
\pgfpathlineto{\pgfqpoint{3.486838in}{2.303270in}}%
\pgfpathlineto{\pgfqpoint{3.486911in}{2.304190in}}%
\pgfpathlineto{\pgfqpoint{3.487714in}{2.301112in}}%
\pgfpathlineto{\pgfqpoint{3.487948in}{2.303419in}}%
\pgfpathlineto{\pgfqpoint{3.488240in}{2.300547in}}%
\pgfpathlineto{\pgfqpoint{3.488766in}{2.305150in}}%
\pgfpathlineto{\pgfqpoint{3.489058in}{2.303443in}}%
\pgfpathlineto{\pgfqpoint{3.489428in}{2.302356in}}%
\pgfpathlineto{\pgfqpoint{3.489501in}{2.302102in}}%
\pgfpathlineto{\pgfqpoint{3.490231in}{2.306046in}}%
\pgfpathlineto{\pgfqpoint{3.490455in}{2.303463in}}%
\pgfpathlineto{\pgfqpoint{3.491570in}{2.304263in}}%
\pgfpathlineto{\pgfqpoint{3.491049in}{2.301278in}}%
\pgfpathlineto{\pgfqpoint{3.491575in}{2.304155in}}%
\pgfpathlineto{\pgfqpoint{3.492436in}{2.300742in}}%
\pgfpathlineto{\pgfqpoint{3.492729in}{2.301579in}}%
\pgfpathlineto{\pgfqpoint{3.493235in}{2.304265in}}%
\pgfpathlineto{\pgfqpoint{3.492738in}{2.301528in}}%
\pgfpathlineto{\pgfqpoint{3.493843in}{2.301756in}}%
\pgfpathlineto{\pgfqpoint{3.495031in}{2.298401in}}%
\pgfpathlineto{\pgfqpoint{3.494389in}{2.302495in}}%
\pgfpathlineto{\pgfqpoint{3.495051in}{2.298650in}}%
\pgfpathlineto{\pgfqpoint{3.495952in}{2.303825in}}%
\pgfpathlineto{\pgfqpoint{3.496331in}{2.303680in}}%
\pgfpathlineto{\pgfqpoint{3.497363in}{2.302482in}}%
\pgfpathlineto{\pgfqpoint{3.497115in}{2.306038in}}%
\pgfpathlineto{\pgfqpoint{3.497432in}{2.303573in}}%
\pgfpathlineto{\pgfqpoint{3.498040in}{2.305948in}}%
\pgfpathlineto{\pgfqpoint{3.497573in}{2.302999in}}%
\pgfpathlineto{\pgfqpoint{3.498527in}{2.303058in}}%
\pgfpathlineto{\pgfqpoint{3.499613in}{2.301947in}}%
\pgfpathlineto{\pgfqpoint{3.499028in}{2.303715in}}%
\pgfpathlineto{\pgfqpoint{3.499681in}{2.302146in}}%
\pgfpathlineto{\pgfqpoint{3.499885in}{2.303874in}}%
\pgfpathlineto{\pgfqpoint{3.499705in}{2.301777in}}%
\pgfpathlineto{\pgfqpoint{3.500732in}{2.301857in}}%
\pgfpathlineto{\pgfqpoint{3.500976in}{2.299053in}}%
\pgfpathlineto{\pgfqpoint{3.501808in}{2.303629in}}%
\pgfpathlineto{\pgfqpoint{3.501818in}{2.303479in}}%
\pgfpathlineto{\pgfqpoint{3.502042in}{2.304915in}}%
\pgfpathlineto{\pgfqpoint{3.502836in}{2.301411in}}%
\pgfpathlineto{\pgfqpoint{3.502889in}{2.301720in}}%
\pgfpathlineto{\pgfqpoint{3.504160in}{2.299084in}}%
\pgfpathlineto{\pgfqpoint{3.503245in}{2.303075in}}%
\pgfpathlineto{\pgfqpoint{3.504189in}{2.299563in}}%
\pgfpathlineto{\pgfqpoint{3.505839in}{2.305901in}}%
\pgfpathlineto{\pgfqpoint{3.504199in}{2.299492in}}%
\pgfpathlineto{\pgfqpoint{3.506107in}{2.304708in}}%
\pgfpathlineto{\pgfqpoint{3.506774in}{2.301256in}}%
\pgfpathlineto{\pgfqpoint{3.506180in}{2.305862in}}%
\pgfpathlineto{\pgfqpoint{3.507271in}{2.303600in}}%
\pgfpathlineto{\pgfqpoint{3.508313in}{2.305376in}}%
\pgfpathlineto{\pgfqpoint{3.507806in}{2.301525in}}%
\pgfpathlineto{\pgfqpoint{3.508400in}{2.305199in}}%
\pgfpathlineto{\pgfqpoint{3.509335in}{2.301400in}}%
\pgfpathlineto{\pgfqpoint{3.508439in}{2.305830in}}%
\pgfpathlineto{\pgfqpoint{3.509695in}{2.303304in}}%
\pgfpathlineto{\pgfqpoint{3.510392in}{2.304052in}}%
\pgfpathlineto{\pgfqpoint{3.510236in}{2.300457in}}%
\pgfpathlineto{\pgfqpoint{3.510781in}{2.302915in}}%
\pgfpathlineto{\pgfqpoint{3.510976in}{2.301006in}}%
\pgfpathlineto{\pgfqpoint{3.511409in}{2.303667in}}%
\pgfpathlineto{\pgfqpoint{3.511925in}{2.301830in}}%
\pgfpathlineto{\pgfqpoint{3.512436in}{2.303907in}}%
\pgfpathlineto{\pgfqpoint{3.512913in}{2.298667in}}%
\pgfpathlineto{\pgfqpoint{3.513006in}{2.300326in}}%
\pgfpathlineto{\pgfqpoint{3.513016in}{2.300165in}}%
\pgfpathlineto{\pgfqpoint{3.513731in}{2.303793in}}%
\pgfpathlineto{\pgfqpoint{3.513907in}{2.302749in}}%
\pgfpathlineto{\pgfqpoint{3.514199in}{2.303998in}}%
\pgfpathlineto{\pgfqpoint{3.514642in}{2.300155in}}%
\pgfpathlineto{\pgfqpoint{3.514973in}{2.301413in}}%
\pgfpathlineto{\pgfqpoint{3.515742in}{2.301324in}}%
\pgfpathlineto{\pgfqpoint{3.515908in}{2.303721in}}%
\pgfpathlineto{\pgfqpoint{3.515981in}{2.303371in}}%
\pgfpathlineto{\pgfqpoint{3.516409in}{2.303796in}}%
\pgfpathlineto{\pgfqpoint{3.516696in}{2.300582in}}%
\pgfpathlineto{\pgfqpoint{3.517018in}{2.302084in}}%
\pgfpathlineto{\pgfqpoint{3.517100in}{2.301520in}}%
\pgfpathlineto{\pgfqpoint{3.517933in}{2.305189in}}%
\pgfpathlineto{\pgfqpoint{3.518103in}{2.303065in}}%
\pgfpathlineto{\pgfqpoint{3.519009in}{2.305049in}}%
\pgfpathlineto{\pgfqpoint{3.518303in}{2.302864in}}%
\pgfpathlineto{\pgfqpoint{3.519204in}{2.303079in}}%
\pgfpathlineto{\pgfqpoint{3.520012in}{2.303526in}}%
\pgfpathlineto{\pgfqpoint{3.520357in}{2.300144in}}%
\pgfpathlineto{\pgfqpoint{3.520722in}{2.298499in}}%
\pgfpathlineto{\pgfqpoint{3.521287in}{2.301299in}}%
\pgfpathlineto{\pgfqpoint{3.521438in}{2.300865in}}%
\pgfpathlineto{\pgfqpoint{3.522241in}{2.303941in}}%
\pgfpathlineto{\pgfqpoint{3.521463in}{2.300779in}}%
\pgfpathlineto{\pgfqpoint{3.522597in}{2.302458in}}%
\pgfpathlineto{\pgfqpoint{3.522709in}{2.301914in}}%
\pgfpathlineto{\pgfqpoint{3.523390in}{2.305807in}}%
\pgfpathlineto{\pgfqpoint{3.523687in}{2.303029in}}%
\pgfpathlineto{\pgfqpoint{3.523775in}{2.304407in}}%
\pgfpathlineto{\pgfqpoint{3.524457in}{2.300349in}}%
\pgfpathlineto{\pgfqpoint{3.524797in}{2.302995in}}%
\pgfpathlineto{\pgfqpoint{3.524802in}{2.302990in}}%
\pgfpathlineto{\pgfqpoint{3.524870in}{2.304402in}}%
\pgfpathlineto{\pgfqpoint{3.524890in}{2.304956in}}%
\pgfpathlineto{\pgfqpoint{3.525766in}{2.302901in}}%
\pgfpathlineto{\pgfqpoint{3.525976in}{2.304133in}}%
\pgfpathlineto{\pgfqpoint{3.526623in}{2.299835in}}%
\pgfpathlineto{\pgfqpoint{3.527105in}{2.302038in}}%
\pgfpathlineto{\pgfqpoint{3.527806in}{2.300836in}}%
\pgfpathlineto{\pgfqpoint{3.527446in}{2.303545in}}%
\pgfpathlineto{\pgfqpoint{3.528201in}{2.303026in}}%
\pgfpathlineto{\pgfqpoint{3.528278in}{2.303425in}}%
\pgfpathlineto{\pgfqpoint{3.528848in}{2.298085in}}%
\pgfpathlineto{\pgfqpoint{3.529164in}{2.298977in}}%
\pgfpathlineto{\pgfqpoint{3.529330in}{2.299610in}}%
\pgfpathlineto{\pgfqpoint{3.529783in}{2.297217in}}%
\pgfpathlineto{\pgfqpoint{3.530109in}{2.297971in}}%
\pgfpathlineto{\pgfqpoint{3.530766in}{2.295156in}}%
\pgfpathlineto{\pgfqpoint{3.531229in}{2.297379in}}%
\pgfpathlineto{\pgfqpoint{3.532412in}{2.302399in}}%
\pgfpathlineto{\pgfqpoint{3.531287in}{2.297076in}}%
\pgfpathlineto{\pgfqpoint{3.532446in}{2.301954in}}%
\pgfpathlineto{\pgfqpoint{3.532660in}{2.300332in}}%
\pgfpathlineto{\pgfqpoint{3.532986in}{2.303328in}}%
\pgfpathlineto{\pgfqpoint{3.533600in}{2.301394in}}%
\pgfpathlineto{\pgfqpoint{3.533605in}{2.301554in}}%
\pgfpathlineto{\pgfqpoint{3.534369in}{2.294695in}}%
\pgfpathlineto{\pgfqpoint{3.534515in}{2.295006in}}%
\pgfpathlineto{\pgfqpoint{3.534525in}{2.294732in}}%
\pgfpathlineto{\pgfqpoint{3.535284in}{2.298594in}}%
\pgfpathlineto{\pgfqpoint{3.535450in}{2.298250in}}%
\pgfpathlineto{\pgfqpoint{3.535995in}{2.300778in}}%
\pgfpathlineto{\pgfqpoint{3.535562in}{2.297821in}}%
\pgfpathlineto{\pgfqpoint{3.536579in}{2.298713in}}%
\pgfpathlineto{\pgfqpoint{3.536589in}{2.298335in}}%
\pgfpathlineto{\pgfqpoint{3.537266in}{2.303382in}}%
\pgfpathlineto{\pgfqpoint{3.537675in}{2.299672in}}%
\pgfpathlineto{\pgfqpoint{3.537679in}{2.299497in}}%
\pgfpathlineto{\pgfqpoint{3.538527in}{2.303067in}}%
\pgfpathlineto{\pgfqpoint{3.538726in}{2.301710in}}%
\pgfpathlineto{\pgfqpoint{3.538872in}{2.303685in}}%
\pgfpathlineto{\pgfqpoint{3.539057in}{2.300866in}}%
\pgfpathlineto{\pgfqpoint{3.539890in}{2.302936in}}%
\pgfpathlineto{\pgfqpoint{3.540678in}{2.301783in}}%
\pgfpathlineto{\pgfqpoint{3.540177in}{2.304343in}}%
\pgfpathlineto{\pgfqpoint{3.541010in}{2.302478in}}%
\pgfpathlineto{\pgfqpoint{3.542129in}{2.304623in}}%
\pgfpathlineto{\pgfqpoint{3.541584in}{2.300824in}}%
\pgfpathlineto{\pgfqpoint{3.542149in}{2.304587in}}%
\pgfpathlineto{\pgfqpoint{3.542903in}{2.302217in}}%
\pgfpathlineto{\pgfqpoint{3.542485in}{2.305094in}}%
\pgfpathlineto{\pgfqpoint{3.543293in}{2.303652in}}%
\pgfpathlineto{\pgfqpoint{3.543536in}{2.304304in}}%
\pgfpathlineto{\pgfqpoint{3.543731in}{2.302672in}}%
\pgfpathlineto{\pgfqpoint{3.544408in}{2.303914in}}%
\pgfpathlineto{\pgfqpoint{3.544627in}{2.303284in}}%
\pgfpathlineto{\pgfqpoint{3.545080in}{2.306010in}}%
\pgfpathlineto{\pgfqpoint{3.545493in}{2.304011in}}%
\pgfpathlineto{\pgfqpoint{3.545586in}{2.304871in}}%
\pgfpathlineto{\pgfqpoint{3.545844in}{2.301036in}}%
\pgfpathlineto{\pgfqpoint{3.546565in}{2.301777in}}%
\pgfpathlineto{\pgfqpoint{3.547490in}{2.306205in}}%
\pgfpathlineto{\pgfqpoint{3.547806in}{2.303175in}}%
\pgfpathlineto{\pgfqpoint{3.548872in}{2.301726in}}%
\pgfpathlineto{\pgfqpoint{3.548342in}{2.304273in}}%
\pgfpathlineto{\pgfqpoint{3.548926in}{2.302426in}}%
\pgfpathlineto{\pgfqpoint{3.549724in}{2.304843in}}%
\pgfpathlineto{\pgfqpoint{3.549271in}{2.301729in}}%
\pgfpathlineto{\pgfqpoint{3.550045in}{2.303788in}}%
\pgfpathlineto{\pgfqpoint{3.550576in}{2.299974in}}%
\pgfpathlineto{\pgfqpoint{3.550201in}{2.305641in}}%
\pgfpathlineto{\pgfqpoint{3.551219in}{2.301245in}}%
\pgfpathlineto{\pgfqpoint{3.551735in}{2.303873in}}%
\pgfpathlineto{\pgfqpoint{3.552275in}{2.301180in}}%
\pgfpathlineto{\pgfqpoint{3.552329in}{2.301281in}}%
\pgfpathlineto{\pgfqpoint{3.553332in}{2.304366in}}%
\pgfpathlineto{\pgfqpoint{3.552377in}{2.300933in}}%
\pgfpathlineto{\pgfqpoint{3.553449in}{2.301851in}}%
\pgfpathlineto{\pgfqpoint{3.553453in}{2.301840in}}%
\pgfpathlineto{\pgfqpoint{3.553731in}{2.303378in}}%
\pgfpathlineto{\pgfqpoint{3.553736in}{2.303344in}}%
\pgfpathlineto{\pgfqpoint{3.554018in}{2.305722in}}%
\pgfpathlineto{\pgfqpoint{3.554700in}{2.301405in}}%
\pgfpathlineto{\pgfqpoint{3.554841in}{2.303145in}}%
\pgfpathlineto{\pgfqpoint{3.555766in}{2.300746in}}%
\pgfpathlineto{\pgfqpoint{3.555323in}{2.304674in}}%
\pgfpathlineto{\pgfqpoint{3.556058in}{2.302756in}}%
\pgfpathlineto{\pgfqpoint{3.556613in}{2.303674in}}%
\pgfpathlineto{\pgfqpoint{3.557124in}{2.300883in}}%
\pgfpathlineto{\pgfqpoint{3.557134in}{2.300935in}}%
\pgfpathlineto{\pgfqpoint{3.558142in}{2.296807in}}%
\pgfpathlineto{\pgfqpoint{3.557231in}{2.301049in}}%
\pgfpathlineto{\pgfqpoint{3.558385in}{2.297435in}}%
\pgfpathlineto{\pgfqpoint{3.558453in}{2.296870in}}%
\pgfpathlineto{\pgfqpoint{3.558857in}{2.299817in}}%
\pgfpathlineto{\pgfqpoint{3.559417in}{2.298905in}}%
\pgfpathlineto{\pgfqpoint{3.560561in}{2.296588in}}%
\pgfpathlineto{\pgfqpoint{3.559773in}{2.300433in}}%
\pgfpathlineto{\pgfqpoint{3.560683in}{2.297415in}}%
\pgfpathlineto{\pgfqpoint{3.561691in}{2.299135in}}%
\pgfpathlineto{\pgfqpoint{3.560931in}{2.296193in}}%
\pgfpathlineto{\pgfqpoint{3.561793in}{2.297445in}}%
\pgfpathlineto{\pgfqpoint{3.562358in}{2.299535in}}%
\pgfpathlineto{\pgfqpoint{3.562509in}{2.297033in}}%
\pgfpathlineto{\pgfqpoint{3.562913in}{2.297783in}}%
\pgfpathlineto{\pgfqpoint{3.563127in}{2.294380in}}%
\pgfpathlineto{\pgfqpoint{3.562937in}{2.298137in}}%
\pgfpathlineto{\pgfqpoint{3.564062in}{2.297671in}}%
\pgfpathlineto{\pgfqpoint{3.564855in}{2.300786in}}%
\pgfpathlineto{\pgfqpoint{3.564481in}{2.297071in}}%
\pgfpathlineto{\pgfqpoint{3.565167in}{2.297392in}}%
\pgfpathlineto{\pgfqpoint{3.565985in}{2.299156in}}%
\pgfpathlineto{\pgfqpoint{3.565435in}{2.296954in}}%
\pgfpathlineto{\pgfqpoint{3.566321in}{2.298285in}}%
\pgfpathlineto{\pgfqpoint{3.567212in}{2.297069in}}%
\pgfpathlineto{\pgfqpoint{3.566589in}{2.300863in}}%
\pgfpathlineto{\pgfqpoint{3.567436in}{2.298157in}}%
\pgfpathlineto{\pgfqpoint{3.567465in}{2.298435in}}%
\pgfpathlineto{\pgfqpoint{3.568205in}{2.294887in}}%
\pgfpathlineto{\pgfqpoint{3.568502in}{2.296169in}}%
\pgfpathlineto{\pgfqpoint{3.568507in}{2.295850in}}%
\pgfpathlineto{\pgfqpoint{3.568828in}{2.298207in}}%
\pgfpathlineto{\pgfqpoint{3.569612in}{2.296065in}}%
\pgfpathlineto{\pgfqpoint{3.570065in}{2.296981in}}%
\pgfpathlineto{\pgfqpoint{3.570415in}{2.293825in}}%
\pgfpathlineto{\pgfqpoint{3.570683in}{2.294912in}}%
\pgfpathlineto{\pgfqpoint{3.571462in}{2.296303in}}%
\pgfpathlineto{\pgfqpoint{3.570800in}{2.294144in}}%
\pgfpathlineto{\pgfqpoint{3.571589in}{2.294368in}}%
\pgfpathlineto{\pgfqpoint{3.572626in}{2.291222in}}%
\pgfpathlineto{\pgfqpoint{3.571696in}{2.294973in}}%
\pgfpathlineto{\pgfqpoint{3.572738in}{2.292519in}}%
\pgfpathlineto{\pgfqpoint{3.572864in}{2.294010in}}%
\pgfpathlineto{\pgfqpoint{3.573512in}{2.290132in}}%
\pgfpathlineto{\pgfqpoint{3.573857in}{2.293054in}}%
\pgfpathlineto{\pgfqpoint{3.573950in}{2.293584in}}%
\pgfpathlineto{\pgfqpoint{3.574072in}{2.291642in}}%
\pgfpathlineto{\pgfqpoint{3.574076in}{2.291654in}}%
\pgfpathlineto{\pgfqpoint{3.574096in}{2.291179in}}%
\pgfpathlineto{\pgfqpoint{3.574919in}{2.295793in}}%
\pgfpathlineto{\pgfqpoint{3.575079in}{2.294818in}}%
\pgfpathlineto{\pgfqpoint{3.575814in}{2.299209in}}%
\pgfpathlineto{\pgfqpoint{3.575313in}{2.294549in}}%
\pgfpathlineto{\pgfqpoint{3.576296in}{2.297649in}}%
\pgfpathlineto{\pgfqpoint{3.577329in}{2.296186in}}%
\pgfpathlineto{\pgfqpoint{3.576710in}{2.299105in}}%
\pgfpathlineto{\pgfqpoint{3.577441in}{2.296835in}}%
\pgfpathlineto{\pgfqpoint{3.578030in}{2.300100in}}%
\pgfpathlineto{\pgfqpoint{3.578565in}{2.298687in}}%
\pgfpathlineto{\pgfqpoint{3.579773in}{2.296740in}}%
\pgfpathlineto{\pgfqpoint{3.578804in}{2.301027in}}%
\pgfpathlineto{\pgfqpoint{3.579782in}{2.297130in}}%
\pgfpathlineto{\pgfqpoint{3.580147in}{2.299578in}}%
\pgfpathlineto{\pgfqpoint{3.580566in}{2.296868in}}%
\pgfpathlineto{\pgfqpoint{3.580980in}{2.299389in}}%
\pgfpathlineto{\pgfqpoint{3.581428in}{2.299982in}}%
\pgfpathlineto{\pgfqpoint{3.582158in}{2.297185in}}%
\pgfpathlineto{\pgfqpoint{3.582859in}{2.303143in}}%
\pgfpathlineto{\pgfqpoint{3.582197in}{2.297128in}}%
\pgfpathlineto{\pgfqpoint{3.583370in}{2.301433in}}%
\pgfpathlineto{\pgfqpoint{3.584505in}{2.304781in}}%
\pgfpathlineto{\pgfqpoint{3.584023in}{2.300381in}}%
\pgfpathlineto{\pgfqpoint{3.584573in}{2.304072in}}%
\pgfpathlineto{\pgfqpoint{3.585654in}{2.301056in}}%
\pgfpathlineto{\pgfqpoint{3.584665in}{2.304777in}}%
\pgfpathlineto{\pgfqpoint{3.585707in}{2.301393in}}%
\pgfpathlineto{\pgfqpoint{3.586749in}{2.302564in}}%
\pgfpathlineto{\pgfqpoint{3.586024in}{2.299973in}}%
\pgfpathlineto{\pgfqpoint{3.586822in}{2.301481in}}%
\pgfpathlineto{\pgfqpoint{3.586881in}{2.300144in}}%
\pgfpathlineto{\pgfqpoint{3.587893in}{2.304199in}}%
\pgfpathlineto{\pgfqpoint{3.587903in}{2.304060in}}%
\pgfpathlineto{\pgfqpoint{3.589101in}{2.301045in}}%
\pgfpathlineto{\pgfqpoint{3.588176in}{2.304951in}}%
\pgfpathlineto{\pgfqpoint{3.589105in}{2.301140in}}%
\pgfpathlineto{\pgfqpoint{3.590191in}{2.303948in}}%
\pgfpathlineto{\pgfqpoint{3.589203in}{2.300431in}}%
\pgfpathlineto{\pgfqpoint{3.590264in}{2.303281in}}%
\pgfpathlineto{\pgfqpoint{3.590936in}{2.300440in}}%
\pgfpathlineto{\pgfqpoint{3.591408in}{2.301274in}}%
\pgfpathlineto{\pgfqpoint{3.592363in}{2.304440in}}%
\pgfpathlineto{\pgfqpoint{3.592543in}{2.303330in}}%
\pgfpathlineto{\pgfqpoint{3.593487in}{2.302224in}}%
\pgfpathlineto{\pgfqpoint{3.593025in}{2.305229in}}%
\pgfpathlineto{\pgfqpoint{3.593662in}{2.302746in}}%
\pgfpathlineto{\pgfqpoint{3.594222in}{2.303992in}}%
\pgfpathlineto{\pgfqpoint{3.594558in}{2.300409in}}%
\pgfpathlineto{\pgfqpoint{3.594733in}{2.300710in}}%
\pgfpathlineto{\pgfqpoint{3.595313in}{2.299757in}}%
\pgfpathlineto{\pgfqpoint{3.595770in}{2.302832in}}%
\pgfpathlineto{\pgfqpoint{3.595785in}{2.302572in}}%
\pgfpathlineto{\pgfqpoint{3.596778in}{2.304152in}}%
\pgfpathlineto{\pgfqpoint{3.596228in}{2.301169in}}%
\pgfpathlineto{\pgfqpoint{3.596880in}{2.302376in}}%
\pgfpathlineto{\pgfqpoint{3.597572in}{2.299159in}}%
\pgfpathlineto{\pgfqpoint{3.596890in}{2.302576in}}%
\pgfpathlineto{\pgfqpoint{3.598025in}{2.301621in}}%
\pgfpathlineto{\pgfqpoint{3.598570in}{2.305077in}}%
\pgfpathlineto{\pgfqpoint{3.599159in}{2.302476in}}%
\pgfpathlineto{\pgfqpoint{3.600235in}{2.301463in}}%
\pgfpathlineto{\pgfqpoint{3.599787in}{2.304688in}}%
\pgfpathlineto{\pgfqpoint{3.600274in}{2.301916in}}%
\pgfpathlineto{\pgfqpoint{3.601228in}{2.302605in}}%
\pgfpathlineto{\pgfqpoint{3.600702in}{2.299785in}}%
\pgfpathlineto{\pgfqpoint{3.601340in}{2.301464in}}%
\pgfpathlineto{\pgfqpoint{3.601739in}{2.299544in}}%
\pgfpathlineto{\pgfqpoint{3.602353in}{2.302676in}}%
\pgfpathlineto{\pgfqpoint{3.602406in}{2.302416in}}%
\pgfpathlineto{\pgfqpoint{3.602557in}{2.303894in}}%
\pgfpathlineto{\pgfqpoint{3.603424in}{2.300518in}}%
\pgfpathlineto{\pgfqpoint{3.603487in}{2.301231in}}%
\pgfpathlineto{\pgfqpoint{3.604534in}{2.301643in}}%
\pgfpathlineto{\pgfqpoint{3.604071in}{2.299194in}}%
\pgfpathlineto{\pgfqpoint{3.604573in}{2.300847in}}%
\pgfpathlineto{\pgfqpoint{3.605454in}{2.297627in}}%
\pgfpathlineto{\pgfqpoint{3.604792in}{2.301468in}}%
\pgfpathlineto{\pgfqpoint{3.605863in}{2.299514in}}%
\pgfpathlineto{\pgfqpoint{3.606715in}{2.304709in}}%
\pgfpathlineto{\pgfqpoint{3.607182in}{2.304346in}}%
\pgfpathlineto{\pgfqpoint{3.608180in}{2.302172in}}%
\pgfpathlineto{\pgfqpoint{3.607294in}{2.306947in}}%
\pgfpathlineto{\pgfqpoint{3.608360in}{2.303465in}}%
\pgfpathlineto{\pgfqpoint{3.608818in}{2.301000in}}%
\pgfpathlineto{\pgfqpoint{3.608541in}{2.304291in}}%
\pgfpathlineto{\pgfqpoint{3.609636in}{2.302290in}}%
\pgfpathlineto{\pgfqpoint{3.610551in}{2.304519in}}%
\pgfpathlineto{\pgfqpoint{3.609870in}{2.301029in}}%
\pgfpathlineto{\pgfqpoint{3.610731in}{2.302147in}}%
\pgfpathlineto{\pgfqpoint{3.611656in}{2.300769in}}%
\pgfpathlineto{\pgfqpoint{3.611131in}{2.304176in}}%
\pgfpathlineto{\pgfqpoint{3.611778in}{2.303278in}}%
\pgfpathlineto{\pgfqpoint{3.611798in}{2.303821in}}%
\pgfpathlineto{\pgfqpoint{3.612479in}{2.301726in}}%
\pgfpathlineto{\pgfqpoint{3.612796in}{2.302166in}}%
\pgfpathlineto{\pgfqpoint{3.612805in}{2.302033in}}%
\pgfpathlineto{\pgfqpoint{3.613278in}{2.304314in}}%
\pgfpathlineto{\pgfqpoint{3.613852in}{2.303317in}}%
\pgfpathlineto{\pgfqpoint{3.614826in}{2.305373in}}%
\pgfpathlineto{\pgfqpoint{3.614198in}{2.302502in}}%
\pgfpathlineto{\pgfqpoint{3.614962in}{2.303372in}}%
\pgfpathlineto{\pgfqpoint{3.616087in}{2.301979in}}%
\pgfpathlineto{\pgfqpoint{3.615439in}{2.305661in}}%
\pgfpathlineto{\pgfqpoint{3.616096in}{2.302809in}}%
\pgfpathlineto{\pgfqpoint{3.616194in}{2.303672in}}%
\pgfpathlineto{\pgfqpoint{3.616963in}{2.302362in}}%
\pgfpathlineto{\pgfqpoint{3.617168in}{2.302510in}}%
\pgfpathlineto{\pgfqpoint{3.617547in}{2.301622in}}%
\pgfpathlineto{\pgfqpoint{3.617937in}{2.303520in}}%
\pgfpathlineto{\pgfqpoint{3.618316in}{2.302016in}}%
\pgfpathlineto{\pgfqpoint{3.618355in}{2.302313in}}%
\pgfpathlineto{\pgfqpoint{3.618385in}{2.300941in}}%
\pgfpathlineto{\pgfqpoint{3.618424in}{2.300282in}}%
\pgfpathlineto{\pgfqpoint{3.619027in}{2.304136in}}%
\pgfpathlineto{\pgfqpoint{3.619407in}{2.302881in}}%
\pgfpathlineto{\pgfqpoint{3.619422in}{2.303268in}}%
\pgfpathlineto{\pgfqpoint{3.620356in}{2.299035in}}%
\pgfpathlineto{\pgfqpoint{3.620478in}{2.300842in}}%
\pgfpathlineto{\pgfqpoint{3.620512in}{2.299880in}}%
\pgfpathlineto{\pgfqpoint{3.620999in}{2.304330in}}%
\pgfpathlineto{\pgfqpoint{3.621462in}{2.303301in}}%
\pgfpathlineto{\pgfqpoint{3.622279in}{2.304851in}}%
\pgfpathlineto{\pgfqpoint{3.622547in}{2.302280in}}%
\pgfpathlineto{\pgfqpoint{3.623458in}{2.304488in}}%
\pgfpathlineto{\pgfqpoint{3.622732in}{2.301668in}}%
\pgfpathlineto{\pgfqpoint{3.623706in}{2.303437in}}%
\pgfpathlineto{\pgfqpoint{3.624178in}{2.301817in}}%
\pgfpathlineto{\pgfqpoint{3.624767in}{2.304586in}}%
\pgfpathlineto{\pgfqpoint{3.624796in}{2.304179in}}%
\pgfpathlineto{\pgfqpoint{3.624821in}{2.304765in}}%
\pgfpathlineto{\pgfqpoint{3.625639in}{2.302576in}}%
\pgfpathlineto{\pgfqpoint{3.625863in}{2.303541in}}%
\pgfpathlineto{\pgfqpoint{3.626072in}{2.302851in}}%
\pgfpathlineto{\pgfqpoint{3.626569in}{2.306884in}}%
\pgfpathlineto{\pgfqpoint{3.626953in}{2.303780in}}%
\pgfpathlineto{\pgfqpoint{3.627021in}{2.304699in}}%
\pgfpathlineto{\pgfqpoint{3.627737in}{2.299772in}}%
\pgfpathlineto{\pgfqpoint{3.628034in}{2.301876in}}%
\pgfpathlineto{\pgfqpoint{3.628964in}{2.304964in}}%
\pgfpathlineto{\pgfqpoint{3.628044in}{2.301843in}}%
\pgfpathlineto{\pgfqpoint{3.629222in}{2.303797in}}%
\pgfpathlineto{\pgfqpoint{3.630118in}{2.301079in}}%
\pgfpathlineto{\pgfqpoint{3.629324in}{2.304320in}}%
\pgfpathlineto{\pgfqpoint{3.630332in}{2.303374in}}%
\pgfpathlineto{\pgfqpoint{3.630381in}{2.303987in}}%
\pgfpathlineto{\pgfqpoint{3.630629in}{2.302126in}}%
\pgfpathlineto{\pgfqpoint{3.631427in}{2.302558in}}%
\pgfpathlineto{\pgfqpoint{3.632586in}{2.299701in}}%
\pgfpathlineto{\pgfqpoint{3.631676in}{2.304423in}}%
\pgfpathlineto{\pgfqpoint{3.632708in}{2.301714in}}%
\pgfpathlineto{\pgfqpoint{3.633053in}{2.303983in}}%
\pgfpathlineto{\pgfqpoint{3.633545in}{2.299964in}}%
\pgfpathlineto{\pgfqpoint{3.633813in}{2.301310in}}%
\pgfpathlineto{\pgfqpoint{3.634008in}{2.300770in}}%
\pgfpathlineto{\pgfqpoint{3.634460in}{2.304214in}}%
\pgfpathlineto{\pgfqpoint{3.634718in}{2.303385in}}%
\pgfpathlineto{\pgfqpoint{3.635463in}{2.304071in}}%
\pgfpathlineto{\pgfqpoint{3.634894in}{2.300948in}}%
\pgfpathlineto{\pgfqpoint{3.635794in}{2.301287in}}%
\pgfpathlineto{\pgfqpoint{3.635916in}{2.300559in}}%
\pgfpathlineto{\pgfqpoint{3.636087in}{2.302563in}}%
\pgfpathlineto{\pgfqpoint{3.636914in}{2.301136in}}%
\pgfpathlineto{\pgfqpoint{3.637455in}{2.304439in}}%
\pgfpathlineto{\pgfqpoint{3.636943in}{2.300659in}}%
\pgfpathlineto{\pgfqpoint{3.638049in}{2.302282in}}%
\pgfpathlineto{\pgfqpoint{3.638876in}{2.298165in}}%
\pgfpathlineto{\pgfqpoint{3.638263in}{2.303821in}}%
\pgfpathlineto{\pgfqpoint{3.639193in}{2.300809in}}%
\pgfpathlineto{\pgfqpoint{3.639197in}{2.300729in}}%
\pgfpathlineto{\pgfqpoint{3.639582in}{2.304746in}}%
\pgfpathlineto{\pgfqpoint{3.640191in}{2.302244in}}%
\pgfpathlineto{\pgfqpoint{3.640493in}{2.305246in}}%
\pgfpathlineto{\pgfqpoint{3.641233in}{2.301344in}}%
\pgfpathlineto{\pgfqpoint{3.641296in}{2.301909in}}%
\pgfpathlineto{\pgfqpoint{3.641301in}{2.301803in}}%
\pgfpathlineto{\pgfqpoint{3.642250in}{2.305242in}}%
\pgfpathlineto{\pgfqpoint{3.642260in}{2.305047in}}%
\pgfpathlineto{\pgfqpoint{3.642265in}{2.305107in}}%
\pgfpathlineto{\pgfqpoint{3.642518in}{2.302686in}}%
\pgfpathlineto{\pgfqpoint{3.643238in}{2.303346in}}%
\pgfpathlineto{\pgfqpoint{3.643345in}{2.304109in}}%
\pgfpathlineto{\pgfqpoint{3.643623in}{2.301645in}}%
\pgfpathlineto{\pgfqpoint{3.644801in}{2.297785in}}%
\pgfpathlineto{\pgfqpoint{3.644008in}{2.302674in}}%
\pgfpathlineto{\pgfqpoint{3.644811in}{2.298083in}}%
\pgfpathlineto{\pgfqpoint{3.644850in}{2.297970in}}%
\pgfpathlineto{\pgfqpoint{3.645799in}{2.300158in}}%
\pgfpathlineto{\pgfqpoint{3.645814in}{2.299964in}}%
\pgfpathlineto{\pgfqpoint{3.646476in}{2.303444in}}%
\pgfpathlineto{\pgfqpoint{3.645828in}{2.299891in}}%
\pgfpathlineto{\pgfqpoint{3.647119in}{2.302262in}}%
\pgfpathlineto{\pgfqpoint{3.648019in}{2.299902in}}%
\pgfpathlineto{\pgfqpoint{3.647265in}{2.303704in}}%
\pgfpathlineto{\pgfqpoint{3.648423in}{2.300281in}}%
\pgfpathlineto{\pgfqpoint{3.649543in}{2.304658in}}%
\pgfpathlineto{\pgfqpoint{3.648477in}{2.299900in}}%
\pgfpathlineto{\pgfqpoint{3.649582in}{2.304513in}}%
\pgfpathlineto{\pgfqpoint{3.650897in}{2.297446in}}%
\pgfpathlineto{\pgfqpoint{3.649601in}{2.304942in}}%
\pgfpathlineto{\pgfqpoint{3.650911in}{2.298097in}}%
\pgfpathlineto{\pgfqpoint{3.651924in}{2.302307in}}%
\pgfpathlineto{\pgfqpoint{3.650921in}{2.298079in}}%
\pgfpathlineto{\pgfqpoint{3.652080in}{2.300797in}}%
\pgfpathlineto{\pgfqpoint{3.652571in}{2.299396in}}%
\pgfpathlineto{\pgfqpoint{3.653082in}{2.303562in}}%
\pgfpathlineto{\pgfqpoint{3.653136in}{2.303081in}}%
\pgfpathlineto{\pgfqpoint{3.653827in}{2.301859in}}%
\pgfpathlineto{\pgfqpoint{3.653375in}{2.303920in}}%
\pgfpathlineto{\pgfqpoint{3.654212in}{2.303359in}}%
\pgfpathlineto{\pgfqpoint{3.654373in}{2.303623in}}%
\pgfpathlineto{\pgfqpoint{3.654723in}{2.300589in}}%
\pgfpathlineto{\pgfqpoint{3.655200in}{2.302228in}}%
\pgfpathlineto{\pgfqpoint{3.655629in}{2.300294in}}%
\pgfpathlineto{\pgfqpoint{3.655994in}{2.304380in}}%
\pgfpathlineto{\pgfqpoint{3.656296in}{2.303367in}}%
\pgfpathlineto{\pgfqpoint{3.656885in}{2.301599in}}%
\pgfpathlineto{\pgfqpoint{3.657313in}{2.304819in}}%
\pgfpathlineto{\pgfqpoint{3.657338in}{2.304380in}}%
\pgfpathlineto{\pgfqpoint{3.657352in}{2.304606in}}%
\pgfpathlineto{\pgfqpoint{3.657824in}{2.300979in}}%
\pgfpathlineto{\pgfqpoint{3.658058in}{2.301728in}}%
\pgfpathlineto{\pgfqpoint{3.659139in}{2.299117in}}%
\pgfpathlineto{\pgfqpoint{3.658078in}{2.301911in}}%
\pgfpathlineto{\pgfqpoint{3.659188in}{2.300270in}}%
\pgfpathlineto{\pgfqpoint{3.659231in}{2.299690in}}%
\pgfpathlineto{\pgfqpoint{3.659334in}{2.301719in}}%
\pgfpathlineto{\pgfqpoint{3.659460in}{2.301541in}}%
\pgfpathlineto{\pgfqpoint{3.660059in}{2.302895in}}%
\pgfpathlineto{\pgfqpoint{3.660293in}{2.299219in}}%
\pgfpathlineto{\pgfqpoint{3.660546in}{2.300602in}}%
\pgfpathlineto{\pgfqpoint{3.661481in}{2.296694in}}%
\pgfpathlineto{\pgfqpoint{3.660736in}{2.301159in}}%
\pgfpathlineto{\pgfqpoint{3.661709in}{2.298427in}}%
\pgfpathlineto{\pgfqpoint{3.661977in}{2.300396in}}%
\pgfpathlineto{\pgfqpoint{3.662746in}{2.296970in}}%
\pgfpathlineto{\pgfqpoint{3.662800in}{2.297199in}}%
\pgfpathlineto{\pgfqpoint{3.663662in}{2.293718in}}%
\pgfpathlineto{\pgfqpoint{3.662863in}{2.297526in}}%
\pgfpathlineto{\pgfqpoint{3.664105in}{2.294075in}}%
\pgfpathlineto{\pgfqpoint{3.665288in}{2.299118in}}%
\pgfpathlineto{\pgfqpoint{3.664582in}{2.294012in}}%
\pgfpathlineto{\pgfqpoint{3.665298in}{2.299110in}}%
\pgfpathlineto{\pgfqpoint{3.665414in}{2.298708in}}%
\pgfpathlineto{\pgfqpoint{3.665999in}{2.301618in}}%
\pgfpathlineto{\pgfqpoint{3.666320in}{2.301575in}}%
\pgfpathlineto{\pgfqpoint{3.667449in}{2.303622in}}%
\pgfpathlineto{\pgfqpoint{3.666865in}{2.300932in}}%
\pgfpathlineto{\pgfqpoint{3.667576in}{2.302831in}}%
\pgfpathlineto{\pgfqpoint{3.667615in}{2.302303in}}%
\pgfpathlineto{\pgfqpoint{3.668365in}{2.307362in}}%
\pgfpathlineto{\pgfqpoint{3.668662in}{2.304554in}}%
\pgfpathlineto{\pgfqpoint{3.669781in}{2.307687in}}%
\pgfpathlineto{\pgfqpoint{3.669859in}{2.306960in}}%
\pgfpathlineto{\pgfqpoint{3.671076in}{2.300531in}}%
\pgfpathlineto{\pgfqpoint{3.670108in}{2.307333in}}%
\pgfpathlineto{\pgfqpoint{3.671081in}{2.300806in}}%
\pgfpathlineto{\pgfqpoint{3.671339in}{2.301252in}}%
\pgfpathlineto{\pgfqpoint{3.671646in}{2.298191in}}%
\pgfpathlineto{\pgfqpoint{3.672143in}{2.299722in}}%
\pgfpathlineto{\pgfqpoint{3.672673in}{2.298503in}}%
\pgfpathlineto{\pgfqpoint{3.672444in}{2.300902in}}%
\pgfpathlineto{\pgfqpoint{3.673233in}{2.300462in}}%
\pgfpathlineto{\pgfqpoint{3.673983in}{2.298290in}}%
\pgfpathlineto{\pgfqpoint{3.673316in}{2.301008in}}%
\pgfpathlineto{\pgfqpoint{3.674363in}{2.300061in}}%
\pgfpathlineto{\pgfqpoint{3.674986in}{2.303072in}}%
\pgfpathlineto{\pgfqpoint{3.674436in}{2.299681in}}%
\pgfpathlineto{\pgfqpoint{3.675482in}{2.300391in}}%
\pgfpathlineto{\pgfqpoint{3.676480in}{2.298275in}}%
\pgfpathlineto{\pgfqpoint{3.676247in}{2.301095in}}%
\pgfpathlineto{\pgfqpoint{3.676651in}{2.298624in}}%
\pgfpathlineto{\pgfqpoint{3.677001in}{2.297017in}}%
\pgfpathlineto{\pgfqpoint{3.677513in}{2.300408in}}%
\pgfpathlineto{\pgfqpoint{3.677707in}{2.299456in}}%
\pgfpathlineto{\pgfqpoint{3.678949in}{2.304304in}}%
\pgfpathlineto{\pgfqpoint{3.678968in}{2.304010in}}%
\pgfpathlineto{\pgfqpoint{3.679996in}{2.302269in}}%
\pgfpathlineto{\pgfqpoint{3.679329in}{2.306695in}}%
\pgfpathlineto{\pgfqpoint{3.680117in}{2.302954in}}%
\pgfpathlineto{\pgfqpoint{3.681252in}{2.305146in}}%
\pgfpathlineto{\pgfqpoint{3.680891in}{2.302249in}}%
\pgfpathlineto{\pgfqpoint{3.681271in}{2.304767in}}%
\pgfpathlineto{\pgfqpoint{3.681325in}{2.305552in}}%
\pgfpathlineto{\pgfqpoint{3.681811in}{2.301653in}}%
\pgfpathlineto{\pgfqpoint{3.682094in}{2.302022in}}%
\pgfpathlineto{\pgfqpoint{3.682186in}{2.301413in}}%
\pgfpathlineto{\pgfqpoint{3.683068in}{2.304310in}}%
\pgfpathlineto{\pgfqpoint{3.683184in}{2.302623in}}%
\pgfpathlineto{\pgfqpoint{3.683647in}{2.304211in}}%
\pgfpathlineto{\pgfqpoint{3.683311in}{2.301812in}}%
\pgfpathlineto{\pgfqpoint{3.684304in}{2.303278in}}%
\pgfpathlineto{\pgfqpoint{3.684854in}{2.303375in}}%
\pgfpathlineto{\pgfqpoint{3.684484in}{2.300641in}}%
\pgfpathlineto{\pgfqpoint{3.685258in}{2.301228in}}%
\pgfpathlineto{\pgfqpoint{3.686480in}{2.295920in}}%
\pgfpathlineto{\pgfqpoint{3.685297in}{2.301884in}}%
\pgfpathlineto{\pgfqpoint{3.686592in}{2.296538in}}%
\pgfpathlineto{\pgfqpoint{3.686860in}{2.297410in}}%
\pgfpathlineto{\pgfqpoint{3.687498in}{2.294673in}}%
\pgfpathlineto{\pgfqpoint{3.687639in}{2.295526in}}%
\pgfpathlineto{\pgfqpoint{3.687741in}{2.295283in}}%
\pgfpathlineto{\pgfqpoint{3.688296in}{2.298821in}}%
\pgfpathlineto{\pgfqpoint{3.688316in}{2.298607in}}%
\pgfpathlineto{\pgfqpoint{3.689119in}{2.303799in}}%
\pgfpathlineto{\pgfqpoint{3.689475in}{2.301058in}}%
\pgfpathlineto{\pgfqpoint{3.689854in}{2.300127in}}%
\pgfpathlineto{\pgfqpoint{3.689582in}{2.302756in}}%
\pgfpathlineto{\pgfqpoint{3.690585in}{2.300981in}}%
\pgfpathlineto{\pgfqpoint{3.691568in}{2.304177in}}%
\pgfpathlineto{\pgfqpoint{3.690648in}{2.300611in}}%
\pgfpathlineto{\pgfqpoint{3.691734in}{2.302834in}}%
\pgfpathlineto{\pgfqpoint{3.692673in}{2.301621in}}%
\pgfpathlineto{\pgfqpoint{3.692318in}{2.304179in}}%
\pgfpathlineto{\pgfqpoint{3.692873in}{2.301886in}}%
\pgfpathlineto{\pgfqpoint{3.692990in}{2.302922in}}%
\pgfpathlineto{\pgfqpoint{3.693520in}{2.299531in}}%
\pgfpathlineto{\pgfqpoint{3.693963in}{2.300843in}}%
\pgfpathlineto{\pgfqpoint{3.695283in}{2.303983in}}%
\pgfpathlineto{\pgfqpoint{3.694387in}{2.299991in}}%
\pgfpathlineto{\pgfqpoint{3.695312in}{2.303876in}}%
\pgfpathlineto{\pgfqpoint{3.695409in}{2.302486in}}%
\pgfpathlineto{\pgfqpoint{3.695716in}{2.305158in}}%
\pgfpathlineto{\pgfqpoint{3.696339in}{2.304200in}}%
\pgfpathlineto{\pgfqpoint{3.696427in}{2.306057in}}%
\pgfpathlineto{\pgfqpoint{3.697420in}{2.303104in}}%
\pgfpathlineto{\pgfqpoint{3.698194in}{2.301708in}}%
\pgfpathlineto{\pgfqpoint{3.697858in}{2.304499in}}%
\pgfpathlineto{\pgfqpoint{3.698535in}{2.302695in}}%
\pgfpathlineto{\pgfqpoint{3.698983in}{2.304733in}}%
\pgfpathlineto{\pgfqpoint{3.699606in}{2.302251in}}%
\pgfpathlineto{\pgfqpoint{3.699640in}{2.302473in}}%
\pgfpathlineto{\pgfqpoint{3.700322in}{2.301574in}}%
\pgfpathlineto{\pgfqpoint{3.700536in}{2.303692in}}%
\pgfpathlineto{\pgfqpoint{3.700774in}{2.301641in}}%
\pgfpathlineto{\pgfqpoint{3.701388in}{2.301052in}}%
\pgfpathlineto{\pgfqpoint{3.701198in}{2.303832in}}%
\pgfpathlineto{\pgfqpoint{3.701826in}{2.302977in}}%
\pgfpathlineto{\pgfqpoint{3.702449in}{2.301211in}}%
\pgfpathlineto{\pgfqpoint{3.702103in}{2.303676in}}%
\pgfpathlineto{\pgfqpoint{3.702907in}{2.303490in}}%
\pgfpathlineto{\pgfqpoint{3.703223in}{2.305045in}}%
\pgfpathlineto{\pgfqpoint{3.703739in}{2.300650in}}%
\pgfpathlineto{\pgfqpoint{3.703788in}{2.300937in}}%
\pgfpathlineto{\pgfqpoint{3.705599in}{2.305921in}}%
\pgfpathlineto{\pgfqpoint{3.703841in}{2.300586in}}%
\pgfpathlineto{\pgfqpoint{3.705716in}{2.304064in}}%
\pgfpathlineto{\pgfqpoint{3.706845in}{2.301827in}}%
\pgfpathlineto{\pgfqpoint{3.705755in}{2.305048in}}%
\pgfpathlineto{\pgfqpoint{3.706855in}{2.301845in}}%
\pgfpathlineto{\pgfqpoint{3.707863in}{2.303534in}}%
\pgfpathlineto{\pgfqpoint{3.706928in}{2.300956in}}%
\pgfpathlineto{\pgfqpoint{3.707985in}{2.301898in}}%
\pgfpathlineto{\pgfqpoint{3.708554in}{2.300976in}}%
\pgfpathlineto{\pgfqpoint{3.708277in}{2.303451in}}%
\pgfpathlineto{\pgfqpoint{3.708900in}{2.303019in}}%
\pgfpathlineto{\pgfqpoint{3.708992in}{2.303710in}}%
\pgfpathlineto{\pgfqpoint{3.709367in}{2.299424in}}%
\pgfpathlineto{\pgfqpoint{3.710010in}{2.303006in}}%
\pgfpathlineto{\pgfqpoint{3.710263in}{2.302066in}}%
\pgfpathlineto{\pgfqpoint{3.710755in}{2.304711in}}%
\pgfpathlineto{\pgfqpoint{3.711052in}{2.303583in}}%
\pgfpathlineto{\pgfqpoint{3.711728in}{2.304809in}}%
\pgfpathlineto{\pgfqpoint{3.711324in}{2.302188in}}%
\pgfpathlineto{\pgfqpoint{3.712167in}{2.303761in}}%
\pgfpathlineto{\pgfqpoint{3.712600in}{2.302153in}}%
\pgfpathlineto{\pgfqpoint{3.712936in}{2.305033in}}%
\pgfpathlineto{\pgfqpoint{3.713301in}{2.303245in}}%
\pgfpathlineto{\pgfqpoint{3.713384in}{2.304410in}}%
\pgfpathlineto{\pgfqpoint{3.713525in}{2.302369in}}%
\pgfpathlineto{\pgfqpoint{3.714348in}{2.302341in}}%
\pgfpathlineto{\pgfqpoint{3.715424in}{2.301695in}}%
\pgfpathlineto{\pgfqpoint{3.714888in}{2.304547in}}%
\pgfpathlineto{\pgfqpoint{3.715472in}{2.301884in}}%
\pgfpathlineto{\pgfqpoint{3.716612in}{2.303383in}}%
\pgfpathlineto{\pgfqpoint{3.716188in}{2.300035in}}%
\pgfpathlineto{\pgfqpoint{3.716616in}{2.302986in}}%
\pgfpathlineto{\pgfqpoint{3.717303in}{2.303444in}}%
\pgfpathlineto{\pgfqpoint{3.717006in}{2.301549in}}%
\pgfpathlineto{\pgfqpoint{3.717580in}{2.302476in}}%
\pgfpathlineto{\pgfqpoint{3.718642in}{2.299710in}}%
\pgfpathlineto{\pgfqpoint{3.718793in}{2.301001in}}%
\pgfpathlineto{\pgfqpoint{3.719761in}{2.305248in}}%
\pgfpathlineto{\pgfqpoint{3.718856in}{2.300641in}}%
\pgfpathlineto{\pgfqpoint{3.720097in}{2.303221in}}%
\pgfpathlineto{\pgfqpoint{3.720292in}{2.301815in}}%
\pgfpathlineto{\pgfqpoint{3.721027in}{2.305022in}}%
\pgfpathlineto{\pgfqpoint{3.721212in}{2.302918in}}%
\pgfpathlineto{\pgfqpoint{3.722356in}{2.304291in}}%
\pgfpathlineto{\pgfqpoint{3.721714in}{2.301296in}}%
\pgfpathlineto{\pgfqpoint{3.722459in}{2.303522in}}%
\pgfpathlineto{\pgfqpoint{3.723452in}{2.302288in}}%
\pgfpathlineto{\pgfqpoint{3.722872in}{2.304726in}}%
\pgfpathlineto{\pgfqpoint{3.723564in}{2.303524in}}%
\pgfpathlineto{\pgfqpoint{3.724645in}{2.304580in}}%
\pgfpathlineto{\pgfqpoint{3.724041in}{2.302038in}}%
\pgfpathlineto{\pgfqpoint{3.724684in}{2.304020in}}%
\pgfpathlineto{\pgfqpoint{3.725682in}{2.304779in}}%
\pgfpathlineto{\pgfqpoint{3.725170in}{2.301980in}}%
\pgfpathlineto{\pgfqpoint{3.725798in}{2.304184in}}%
\pgfpathlineto{\pgfqpoint{3.726631in}{2.301369in}}%
\pgfpathlineto{\pgfqpoint{3.726933in}{2.302343in}}%
\pgfpathlineto{\pgfqpoint{3.727347in}{2.304515in}}%
\pgfpathlineto{\pgfqpoint{3.726952in}{2.302036in}}%
\pgfpathlineto{\pgfqpoint{3.728091in}{2.303542in}}%
\pgfpathlineto{\pgfqpoint{3.728500in}{2.300425in}}%
\pgfpathlineto{\pgfqpoint{3.729216in}{2.302439in}}%
\pgfpathlineto{\pgfqpoint{3.729270in}{2.302864in}}%
\pgfpathlineto{\pgfqpoint{3.730127in}{2.297199in}}%
\pgfpathlineto{\pgfqpoint{3.731168in}{2.291794in}}%
\pgfpathlineto{\pgfqpoint{3.730394in}{2.298260in}}%
\pgfpathlineto{\pgfqpoint{3.731295in}{2.291835in}}%
\pgfpathlineto{\pgfqpoint{3.733393in}{2.287511in}}%
\pgfpathlineto{\pgfqpoint{3.731631in}{2.292365in}}%
\pgfpathlineto{\pgfqpoint{3.733642in}{2.289445in}}%
\pgfpathlineto{\pgfqpoint{3.734747in}{2.292486in}}%
\pgfpathlineto{\pgfqpoint{3.734124in}{2.288262in}}%
\pgfpathlineto{\pgfqpoint{3.734810in}{2.291215in}}%
\pgfpathlineto{\pgfqpoint{3.735370in}{2.296634in}}%
\pgfpathlineto{\pgfqpoint{3.734829in}{2.291124in}}%
\pgfpathlineto{\pgfqpoint{3.736412in}{2.293462in}}%
\pgfpathlineto{\pgfqpoint{3.736582in}{2.292730in}}%
\pgfpathlineto{\pgfqpoint{3.737235in}{2.296327in}}%
\pgfpathlineto{\pgfqpoint{3.737488in}{2.295742in}}%
\pgfpathlineto{\pgfqpoint{3.737960in}{2.293842in}}%
\pgfpathlineto{\pgfqpoint{3.738062in}{2.296367in}}%
\pgfpathlineto{\pgfqpoint{3.738505in}{2.296041in}}%
\pgfpathlineto{\pgfqpoint{3.738802in}{2.296668in}}%
\pgfpathlineto{\pgfqpoint{3.739552in}{2.293808in}}%
\pgfpathlineto{\pgfqpoint{3.739567in}{2.294029in}}%
\pgfpathlineto{\pgfqpoint{3.739756in}{2.292562in}}%
\pgfpathlineto{\pgfqpoint{3.740350in}{2.298714in}}%
\pgfpathlineto{\pgfqpoint{3.740375in}{2.298496in}}%
\pgfpathlineto{\pgfqpoint{3.740935in}{2.299129in}}%
\pgfpathlineto{\pgfqpoint{3.741344in}{2.296013in}}%
\pgfpathlineto{\pgfqpoint{3.741451in}{2.296583in}}%
\pgfpathlineto{\pgfqpoint{3.742385in}{2.290636in}}%
\pgfpathlineto{\pgfqpoint{3.741465in}{2.296598in}}%
\pgfpathlineto{\pgfqpoint{3.742833in}{2.291409in}}%
\pgfpathlineto{\pgfqpoint{3.744016in}{2.297521in}}%
\pgfpathlineto{\pgfqpoint{3.742872in}{2.291232in}}%
\pgfpathlineto{\pgfqpoint{3.744026in}{2.297501in}}%
\pgfpathlineto{\pgfqpoint{3.744820in}{2.296510in}}%
\pgfpathlineto{\pgfqpoint{3.744406in}{2.299799in}}%
\pgfpathlineto{\pgfqpoint{3.745122in}{2.298051in}}%
\pgfpathlineto{\pgfqpoint{3.746144in}{2.298544in}}%
\pgfpathlineto{\pgfqpoint{3.745370in}{2.296221in}}%
\pgfpathlineto{\pgfqpoint{3.746227in}{2.297952in}}%
\pgfpathlineto{\pgfqpoint{3.746743in}{2.299211in}}%
\pgfpathlineto{\pgfqpoint{3.746300in}{2.297172in}}%
\pgfpathlineto{\pgfqpoint{3.746894in}{2.297403in}}%
\pgfpathlineto{\pgfqpoint{3.747307in}{2.296178in}}%
\pgfpathlineto{\pgfqpoint{3.747799in}{2.299348in}}%
\pgfpathlineto{\pgfqpoint{3.747984in}{2.298740in}}%
\pgfpathlineto{\pgfqpoint{3.749060in}{2.302915in}}%
\pgfpathlineto{\pgfqpoint{3.749187in}{2.301808in}}%
\pgfpathlineto{\pgfqpoint{3.749367in}{2.300992in}}%
\pgfpathlineto{\pgfqpoint{3.750039in}{2.304634in}}%
\pgfpathlineto{\pgfqpoint{3.750199in}{2.303209in}}%
\pgfpathlineto{\pgfqpoint{3.751090in}{2.304988in}}%
\pgfpathlineto{\pgfqpoint{3.750579in}{2.301312in}}%
\pgfpathlineto{\pgfqpoint{3.751314in}{2.303434in}}%
\pgfpathlineto{\pgfqpoint{3.751465in}{2.301926in}}%
\pgfpathlineto{\pgfqpoint{3.752293in}{2.305407in}}%
\pgfpathlineto{\pgfqpoint{3.752400in}{2.304789in}}%
\pgfpathlineto{\pgfqpoint{3.753841in}{2.300492in}}%
\pgfpathlineto{\pgfqpoint{3.752439in}{2.304855in}}%
\pgfpathlineto{\pgfqpoint{3.753972in}{2.301709in}}%
\pgfpathlineto{\pgfqpoint{3.754084in}{2.303352in}}%
\pgfpathlineto{\pgfqpoint{3.754785in}{2.299308in}}%
\pgfpathlineto{\pgfqpoint{3.755078in}{2.301521in}}%
\pgfpathlineto{\pgfqpoint{3.755165in}{2.302611in}}%
\pgfpathlineto{\pgfqpoint{3.755418in}{2.300292in}}%
\pgfpathlineto{\pgfqpoint{3.755973in}{2.300756in}}%
\pgfpathlineto{\pgfqpoint{3.756845in}{2.298291in}}%
\pgfpathlineto{\pgfqpoint{3.756562in}{2.301401in}}%
\pgfpathlineto{\pgfqpoint{3.757093in}{2.300389in}}%
\pgfpathlineto{\pgfqpoint{3.757439in}{2.301176in}}%
\pgfpathlineto{\pgfqpoint{3.757672in}{2.299265in}}%
\pgfpathlineto{\pgfqpoint{3.758208in}{2.300597in}}%
\pgfpathlineto{\pgfqpoint{3.758783in}{2.299816in}}%
\pgfpathlineto{\pgfqpoint{3.758943in}{2.302510in}}%
\pgfpathlineto{\pgfqpoint{3.759269in}{2.301639in}}%
\pgfpathlineto{\pgfqpoint{3.760000in}{2.304177in}}%
\pgfpathlineto{\pgfqpoint{3.759401in}{2.301208in}}%
\pgfpathlineto{\pgfqpoint{3.760316in}{2.301211in}}%
\pgfpathlineto{\pgfqpoint{3.761032in}{2.298540in}}%
\pgfpathlineto{\pgfqpoint{3.761431in}{2.300652in}}%
\pgfpathlineto{\pgfqpoint{3.761889in}{2.298736in}}%
\pgfpathlineto{\pgfqpoint{3.761747in}{2.302040in}}%
\pgfpathlineto{\pgfqpoint{3.762507in}{2.301343in}}%
\pgfpathlineto{\pgfqpoint{3.762653in}{2.302511in}}%
\pgfpathlineto{\pgfqpoint{3.763437in}{2.300336in}}%
\pgfpathlineto{\pgfqpoint{3.764415in}{2.299225in}}%
\pgfpathlineto{\pgfqpoint{3.763982in}{2.302501in}}%
\pgfpathlineto{\pgfqpoint{3.764547in}{2.300162in}}%
\pgfpathlineto{\pgfqpoint{3.765521in}{2.303970in}}%
\pgfpathlineto{\pgfqpoint{3.764600in}{2.300077in}}%
\pgfpathlineto{\pgfqpoint{3.765715in}{2.303322in}}%
\pgfpathlineto{\pgfqpoint{3.765720in}{2.303045in}}%
\pgfpathlineto{\pgfqpoint{3.766528in}{2.305908in}}%
\pgfpathlineto{\pgfqpoint{3.766791in}{2.305140in}}%
\pgfpathlineto{\pgfqpoint{3.766820in}{2.305359in}}%
\pgfpathlineto{\pgfqpoint{3.766976in}{2.302682in}}%
\pgfpathlineto{\pgfqpoint{3.767755in}{2.303985in}}%
\pgfpathlineto{\pgfqpoint{3.768666in}{2.303151in}}%
\pgfpathlineto{\pgfqpoint{3.768320in}{2.305169in}}%
\pgfpathlineto{\pgfqpoint{3.768865in}{2.303901in}}%
\pgfpathlineto{\pgfqpoint{3.769990in}{2.305264in}}%
\pgfpathlineto{\pgfqpoint{3.769440in}{2.301930in}}%
\pgfpathlineto{\pgfqpoint{3.769995in}{2.305141in}}%
\pgfpathlineto{\pgfqpoint{3.770423in}{2.303452in}}%
\pgfpathlineto{\pgfqpoint{3.770112in}{2.305641in}}%
\pgfpathlineto{\pgfqpoint{3.771110in}{2.304593in}}%
\pgfpathlineto{\pgfqpoint{3.771144in}{2.304758in}}%
\pgfpathlineto{\pgfqpoint{3.771358in}{2.302993in}}%
\pgfpathlineto{\pgfqpoint{3.771903in}{2.303348in}}%
\pgfpathlineto{\pgfqpoint{3.771942in}{2.301934in}}%
\pgfpathlineto{\pgfqpoint{3.772663in}{2.304902in}}%
\pgfpathlineto{\pgfqpoint{3.773008in}{2.303914in}}%
\pgfpathlineto{\pgfqpoint{3.773315in}{2.305285in}}%
\pgfpathlineto{\pgfqpoint{3.773646in}{2.301857in}}%
\pgfpathlineto{\pgfqpoint{3.774040in}{2.302726in}}%
\pgfpathlineto{\pgfqpoint{3.774936in}{2.299036in}}%
\pgfpathlineto{\pgfqpoint{3.774167in}{2.303969in}}%
\pgfpathlineto{\pgfqpoint{3.775155in}{2.302233in}}%
\pgfpathlineto{\pgfqpoint{3.775345in}{2.300839in}}%
\pgfpathlineto{\pgfqpoint{3.775978in}{2.304132in}}%
\pgfpathlineto{\pgfqpoint{3.776207in}{2.303558in}}%
\pgfpathlineto{\pgfqpoint{3.776226in}{2.303967in}}%
\pgfpathlineto{\pgfqpoint{3.776708in}{2.301697in}}%
\pgfpathlineto{\pgfqpoint{3.777205in}{2.301808in}}%
\pgfpathlineto{\pgfqpoint{3.778354in}{2.296615in}}%
\pgfpathlineto{\pgfqpoint{3.777312in}{2.301988in}}%
\pgfpathlineto{\pgfqpoint{3.778442in}{2.298332in}}%
\pgfpathlineto{\pgfqpoint{3.779415in}{2.303088in}}%
\pgfpathlineto{\pgfqpoint{3.778515in}{2.297582in}}%
\pgfpathlineto{\pgfqpoint{3.779600in}{2.299402in}}%
\pgfpathlineto{\pgfqpoint{3.780009in}{2.297459in}}%
\pgfpathlineto{\pgfqpoint{3.780472in}{2.300082in}}%
\pgfpathlineto{\pgfqpoint{3.780701in}{2.299846in}}%
\pgfpathlineto{\pgfqpoint{3.781450in}{2.302954in}}%
\pgfpathlineto{\pgfqpoint{3.780720in}{2.299668in}}%
\pgfpathlineto{\pgfqpoint{3.781849in}{2.302315in}}%
\pgfpathlineto{\pgfqpoint{3.781864in}{2.302240in}}%
\pgfpathlineto{\pgfqpoint{3.781971in}{2.303981in}}%
\pgfpathlineto{\pgfqpoint{3.781981in}{2.303947in}}%
\pgfpathlineto{\pgfqpoint{3.782005in}{2.304773in}}%
\pgfpathlineto{\pgfqpoint{3.782682in}{2.301340in}}%
\pgfpathlineto{\pgfqpoint{3.783057in}{2.302753in}}%
\pgfpathlineto{\pgfqpoint{3.783203in}{2.302148in}}%
\pgfpathlineto{\pgfqpoint{3.783466in}{2.303882in}}%
\pgfpathlineto{\pgfqpoint{3.784099in}{2.302999in}}%
\pgfpathlineto{\pgfqpoint{3.784717in}{2.303792in}}%
\pgfpathlineto{\pgfqpoint{3.785155in}{2.300272in}}%
\pgfpathlineto{\pgfqpoint{3.785403in}{2.299206in}}%
\pgfpathlineto{\pgfqpoint{3.786002in}{2.302812in}}%
\pgfpathlineto{\pgfqpoint{3.786221in}{2.302039in}}%
\pgfpathlineto{\pgfqpoint{3.786363in}{2.302764in}}%
\pgfpathlineto{\pgfqpoint{3.787107in}{2.297756in}}%
\pgfpathlineto{\pgfqpoint{3.787219in}{2.298066in}}%
\pgfpathlineto{\pgfqpoint{3.787487in}{2.298735in}}%
\pgfpathlineto{\pgfqpoint{3.787560in}{2.296825in}}%
\pgfpathlineto{\pgfqpoint{3.788622in}{2.293781in}}%
\pgfpathlineto{\pgfqpoint{3.788724in}{2.294741in}}%
\pgfpathlineto{\pgfqpoint{3.789566in}{2.296858in}}%
\pgfpathlineto{\pgfqpoint{3.789220in}{2.293518in}}%
\pgfpathlineto{\pgfqpoint{3.789839in}{2.294921in}}%
\pgfpathlineto{\pgfqpoint{3.789878in}{2.295055in}}%
\pgfpathlineto{\pgfqpoint{3.790808in}{2.292578in}}%
\pgfpathlineto{\pgfqpoint{3.790812in}{2.292641in}}%
\pgfpathlineto{\pgfqpoint{3.791436in}{2.289295in}}%
\pgfpathlineto{\pgfqpoint{3.791621in}{2.292945in}}%
\pgfpathlineto{\pgfqpoint{3.791942in}{2.292372in}}%
\pgfpathlineto{\pgfqpoint{3.792424in}{2.294577in}}%
\pgfpathlineto{\pgfqpoint{3.792677in}{2.290756in}}%
\pgfpathlineto{\pgfqpoint{3.792823in}{2.291289in}}%
\pgfpathlineto{\pgfqpoint{3.793772in}{2.290037in}}%
\pgfpathlineto{\pgfqpoint{3.793203in}{2.293199in}}%
\pgfpathlineto{\pgfqpoint{3.793923in}{2.292064in}}%
\pgfpathlineto{\pgfqpoint{3.794576in}{2.290552in}}%
\pgfpathlineto{\pgfqpoint{3.794177in}{2.293210in}}%
\pgfpathlineto{\pgfqpoint{3.795043in}{2.291658in}}%
\pgfpathlineto{\pgfqpoint{3.796051in}{2.287873in}}%
\pgfpathlineto{\pgfqpoint{3.795822in}{2.291845in}}%
\pgfpathlineto{\pgfqpoint{3.796304in}{2.288603in}}%
\pgfpathlineto{\pgfqpoint{3.797307in}{2.290940in}}%
\pgfpathlineto{\pgfqpoint{3.796801in}{2.286991in}}%
\pgfpathlineto{\pgfqpoint{3.797419in}{2.289178in}}%
\pgfpathlineto{\pgfqpoint{3.798436in}{2.287703in}}%
\pgfpathlineto{\pgfqpoint{3.797726in}{2.290076in}}%
\pgfpathlineto{\pgfqpoint{3.798529in}{2.288850in}}%
\pgfpathlineto{\pgfqpoint{3.799668in}{2.292492in}}%
\pgfpathlineto{\pgfqpoint{3.798884in}{2.287809in}}%
\pgfpathlineto{\pgfqpoint{3.799693in}{2.291839in}}%
\pgfpathlineto{\pgfqpoint{3.800316in}{2.295106in}}%
\pgfpathlineto{\pgfqpoint{3.799843in}{2.290560in}}%
\pgfpathlineto{\pgfqpoint{3.800846in}{2.292580in}}%
\pgfpathlineto{\pgfqpoint{3.800851in}{2.292254in}}%
\pgfpathlineto{\pgfqpoint{3.801445in}{2.296108in}}%
\pgfpathlineto{\pgfqpoint{3.801913in}{2.295222in}}%
\pgfpathlineto{\pgfqpoint{3.803037in}{2.297372in}}%
\pgfpathlineto{\pgfqpoint{3.802521in}{2.293583in}}%
\pgfpathlineto{\pgfqpoint{3.803120in}{2.296990in}}%
\pgfpathlineto{\pgfqpoint{3.803135in}{2.296719in}}%
\pgfpathlineto{\pgfqpoint{3.803753in}{2.300194in}}%
\pgfpathlineto{\pgfqpoint{3.804147in}{2.300837in}}%
\pgfpathlineto{\pgfqpoint{3.804551in}{2.296744in}}%
\pgfpathlineto{\pgfqpoint{3.804814in}{2.297720in}}%
\pgfpathlineto{\pgfqpoint{3.805004in}{2.295107in}}%
\pgfpathlineto{\pgfqpoint{3.805452in}{2.298622in}}%
\pgfpathlineto{\pgfqpoint{3.805914in}{2.298133in}}%
\pgfpathlineto{\pgfqpoint{3.805958in}{2.298613in}}%
\pgfpathlineto{\pgfqpoint{3.806182in}{2.296600in}}%
\pgfpathlineto{\pgfqpoint{3.806348in}{2.294927in}}%
\pgfpathlineto{\pgfqpoint{3.807253in}{2.298979in}}%
\pgfpathlineto{\pgfqpoint{3.807779in}{2.301906in}}%
\pgfpathlineto{\pgfqpoint{3.807312in}{2.298240in}}%
\pgfpathlineto{\pgfqpoint{3.808553in}{2.299996in}}%
\pgfpathlineto{\pgfqpoint{3.808563in}{2.299607in}}%
\pgfpathlineto{\pgfqpoint{3.809580in}{2.302767in}}%
\pgfpathlineto{\pgfqpoint{3.809639in}{2.301743in}}%
\pgfpathlineto{\pgfqpoint{3.810131in}{2.299594in}}%
\pgfpathlineto{\pgfqpoint{3.810690in}{2.302666in}}%
\pgfpathlineto{\pgfqpoint{3.811314in}{2.304935in}}%
\pgfpathlineto{\pgfqpoint{3.811825in}{2.303462in}}%
\pgfpathlineto{\pgfqpoint{3.812803in}{2.302547in}}%
\pgfpathlineto{\pgfqpoint{3.812448in}{2.305263in}}%
\pgfpathlineto{\pgfqpoint{3.812935in}{2.303527in}}%
\pgfpathlineto{\pgfqpoint{3.813874in}{2.299124in}}%
\pgfpathlineto{\pgfqpoint{3.813076in}{2.304150in}}%
\pgfpathlineto{\pgfqpoint{3.814142in}{2.299413in}}%
\pgfpathlineto{\pgfqpoint{3.814941in}{2.301512in}}%
\pgfpathlineto{\pgfqpoint{3.815257in}{2.299819in}}%
\pgfpathlineto{\pgfqpoint{3.816333in}{2.293639in}}%
\pgfpathlineto{\pgfqpoint{3.816430in}{2.294033in}}%
\pgfpathlineto{\pgfqpoint{3.817058in}{2.296256in}}%
\pgfpathlineto{\pgfqpoint{3.817365in}{2.293534in}}%
\pgfpathlineto{\pgfqpoint{3.817540in}{2.294068in}}%
\pgfpathlineto{\pgfqpoint{3.818203in}{2.291849in}}%
\pgfpathlineto{\pgfqpoint{3.818490in}{2.294943in}}%
\pgfpathlineto{\pgfqpoint{3.818670in}{2.292764in}}%
\pgfpathlineto{\pgfqpoint{3.818719in}{2.292993in}}%
\pgfpathlineto{\pgfqpoint{3.819327in}{2.289689in}}%
\pgfpathlineto{\pgfqpoint{3.819332in}{2.289774in}}%
\pgfpathlineto{\pgfqpoint{3.819337in}{2.289580in}}%
\pgfpathlineto{\pgfqpoint{3.819668in}{2.292125in}}%
\pgfpathlineto{\pgfqpoint{3.820398in}{2.291556in}}%
\pgfpathlineto{\pgfqpoint{3.820812in}{2.290274in}}%
\pgfpathlineto{\pgfqpoint{3.820530in}{2.292691in}}%
\pgfpathlineto{\pgfqpoint{3.821318in}{2.292159in}}%
\pgfpathlineto{\pgfqpoint{3.822526in}{2.296940in}}%
\pgfpathlineto{\pgfqpoint{3.821328in}{2.292027in}}%
\pgfpathlineto{\pgfqpoint{3.822575in}{2.296057in}}%
\pgfpathlineto{\pgfqpoint{3.823256in}{2.293505in}}%
\pgfpathlineto{\pgfqpoint{3.823665in}{2.296687in}}%
\pgfpathlineto{\pgfqpoint{3.824624in}{2.301288in}}%
\pgfpathlineto{\pgfqpoint{3.823767in}{2.295339in}}%
\pgfpathlineto{\pgfqpoint{3.824863in}{2.300837in}}%
\pgfpathlineto{\pgfqpoint{3.825018in}{2.298987in}}%
\pgfpathlineto{\pgfqpoint{3.825802in}{2.303336in}}%
\pgfpathlineto{\pgfqpoint{3.825963in}{2.301645in}}%
\pgfpathlineto{\pgfqpoint{3.827112in}{2.303524in}}%
\pgfpathlineto{\pgfqpoint{3.826737in}{2.300037in}}%
\pgfpathlineto{\pgfqpoint{3.827146in}{2.303062in}}%
\pgfpathlineto{\pgfqpoint{3.827789in}{2.301963in}}%
\pgfpathlineto{\pgfqpoint{3.827555in}{2.306005in}}%
\pgfpathlineto{\pgfqpoint{3.828246in}{2.303552in}}%
\pgfpathlineto{\pgfqpoint{3.828831in}{2.305257in}}%
\pgfpathlineto{\pgfqpoint{3.829239in}{2.303022in}}%
\pgfpathlineto{\pgfqpoint{3.829376in}{2.303754in}}%
\pgfpathlineto{\pgfqpoint{3.830087in}{2.300843in}}%
\pgfpathlineto{\pgfqpoint{3.829644in}{2.304138in}}%
\pgfpathlineto{\pgfqpoint{3.830525in}{2.301372in}}%
\pgfpathlineto{\pgfqpoint{3.831382in}{2.304625in}}%
\pgfpathlineto{\pgfqpoint{3.830627in}{2.301004in}}%
\pgfpathlineto{\pgfqpoint{3.831674in}{2.302830in}}%
\pgfpathlineto{\pgfqpoint{3.831693in}{2.302097in}}%
\pgfpathlineto{\pgfqpoint{3.832102in}{2.305060in}}%
\pgfpathlineto{\pgfqpoint{3.832725in}{2.303215in}}%
\pgfpathlineto{\pgfqpoint{3.833202in}{2.304948in}}%
\pgfpathlineto{\pgfqpoint{3.833670in}{2.302422in}}%
\pgfpathlineto{\pgfqpoint{3.833845in}{2.304036in}}%
\pgfpathlineto{\pgfqpoint{3.834984in}{2.302169in}}%
\pgfpathlineto{\pgfqpoint{3.834405in}{2.306483in}}%
\pgfpathlineto{\pgfqpoint{3.835014in}{2.302796in}}%
\pgfpathlineto{\pgfqpoint{3.835033in}{2.303591in}}%
\pgfpathlineto{\pgfqpoint{3.836065in}{2.301123in}}%
\pgfpathlineto{\pgfqpoint{3.836094in}{2.301690in}}%
\pgfpathlineto{\pgfqpoint{3.836683in}{2.300785in}}%
\pgfpathlineto{\pgfqpoint{3.837014in}{2.302911in}}%
\pgfpathlineto{\pgfqpoint{3.837156in}{2.302674in}}%
\pgfpathlineto{\pgfqpoint{3.837287in}{2.303129in}}%
\pgfpathlineto{\pgfqpoint{3.837993in}{2.299356in}}%
\pgfpathlineto{\pgfqpoint{3.838013in}{2.299015in}}%
\pgfpathlineto{\pgfqpoint{3.838314in}{2.302671in}}%
\pgfpathlineto{\pgfqpoint{3.838869in}{2.301997in}}%
\pgfpathlineto{\pgfqpoint{3.839045in}{2.303947in}}%
\pgfpathlineto{\pgfqpoint{3.839882in}{2.299155in}}%
\pgfpathlineto{\pgfqpoint{3.839945in}{2.299909in}}%
\pgfpathlineto{\pgfqpoint{3.840121in}{2.298742in}}%
\pgfpathlineto{\pgfqpoint{3.842214in}{2.303987in}}%
\pgfpathlineto{\pgfqpoint{3.842677in}{2.304376in}}%
\pgfpathlineto{\pgfqpoint{3.843027in}{2.302179in}}%
\pgfpathlineto{\pgfqpoint{3.843319in}{2.303790in}}%
\pgfpathlineto{\pgfqpoint{3.843451in}{2.304897in}}%
\pgfpathlineto{\pgfqpoint{3.844288in}{2.302131in}}%
\pgfpathlineto{\pgfqpoint{3.844327in}{2.301335in}}%
\pgfpathlineto{\pgfqpoint{3.844658in}{2.303566in}}%
\pgfpathlineto{\pgfqpoint{3.845315in}{2.302981in}}%
\pgfpathlineto{\pgfqpoint{3.845997in}{2.303982in}}%
\pgfpathlineto{\pgfqpoint{3.846377in}{2.300581in}}%
\pgfpathlineto{\pgfqpoint{3.847083in}{2.297026in}}%
\pgfpathlineto{\pgfqpoint{3.846552in}{2.301448in}}%
\pgfpathlineto{\pgfqpoint{3.847540in}{2.298961in}}%
\pgfpathlineto{\pgfqpoint{3.847973in}{2.301657in}}%
\pgfpathlineto{\pgfqpoint{3.847560in}{2.298741in}}%
\pgfpathlineto{\pgfqpoint{3.848679in}{2.299684in}}%
\pgfpathlineto{\pgfqpoint{3.848806in}{2.298424in}}%
\pgfpathlineto{\pgfqpoint{3.849653in}{2.303149in}}%
\pgfpathlineto{\pgfqpoint{3.849716in}{2.304477in}}%
\pgfpathlineto{\pgfqpoint{3.850471in}{2.298765in}}%
\pgfpathlineto{\pgfqpoint{3.850651in}{2.299947in}}%
\pgfpathlineto{\pgfqpoint{3.851781in}{2.298344in}}%
\pgfpathlineto{\pgfqpoint{3.851260in}{2.300734in}}%
\pgfpathlineto{\pgfqpoint{3.851790in}{2.298530in}}%
\pgfpathlineto{\pgfqpoint{3.851815in}{2.298188in}}%
\pgfpathlineto{\pgfqpoint{3.852691in}{2.302606in}}%
\pgfpathlineto{\pgfqpoint{3.853027in}{2.304051in}}%
\pgfpathlineto{\pgfqpoint{3.853485in}{2.300757in}}%
\pgfpathlineto{\pgfqpoint{3.853757in}{2.301847in}}%
\pgfpathlineto{\pgfqpoint{3.854848in}{2.300077in}}%
\pgfpathlineto{\pgfqpoint{3.854098in}{2.303835in}}%
\pgfpathlineto{\pgfqpoint{3.854877in}{2.301012in}}%
\pgfpathlineto{\pgfqpoint{3.855773in}{2.302888in}}%
\pgfpathlineto{\pgfqpoint{3.854999in}{2.299643in}}%
\pgfpathlineto{\pgfqpoint{3.856011in}{2.301959in}}%
\pgfpathlineto{\pgfqpoint{3.856060in}{2.301638in}}%
\pgfpathlineto{\pgfqpoint{3.856644in}{2.305278in}}%
\pgfpathlineto{\pgfqpoint{3.857009in}{2.304037in}}%
\pgfpathlineto{\pgfqpoint{3.857462in}{2.305142in}}%
\pgfpathlineto{\pgfqpoint{3.858066in}{2.302900in}}%
\pgfpathlineto{\pgfqpoint{3.858743in}{2.300562in}}%
\pgfpathlineto{\pgfqpoint{3.858937in}{2.303936in}}%
\pgfpathlineto{\pgfqpoint{3.859176in}{2.302788in}}%
\pgfpathlineto{\pgfqpoint{3.859711in}{2.302012in}}%
\pgfpathlineto{\pgfqpoint{3.859453in}{2.304519in}}%
\pgfpathlineto{\pgfqpoint{3.859877in}{2.304127in}}%
\pgfpathlineto{\pgfqpoint{3.859935in}{2.304465in}}%
\pgfpathlineto{\pgfqpoint{3.860933in}{2.300385in}}%
\pgfpathlineto{\pgfqpoint{3.861936in}{2.304133in}}%
\pgfpathlineto{\pgfqpoint{3.861216in}{2.298790in}}%
\pgfpathlineto{\pgfqpoint{3.862136in}{2.301875in}}%
\pgfpathlineto{\pgfqpoint{3.862185in}{2.302657in}}%
\pgfpathlineto{\pgfqpoint{3.862292in}{2.300549in}}%
\pgfpathlineto{\pgfqpoint{3.862336in}{2.300771in}}%
\pgfpathlineto{\pgfqpoint{3.863119in}{2.296109in}}%
\pgfpathlineto{\pgfqpoint{3.863475in}{2.296917in}}%
\pgfpathlineto{\pgfqpoint{3.864619in}{2.299160in}}%
\pgfpathlineto{\pgfqpoint{3.863859in}{2.294870in}}%
\pgfpathlineto{\pgfqpoint{3.864638in}{2.299129in}}%
\pgfpathlineto{\pgfqpoint{3.864760in}{2.297976in}}%
\pgfpathlineto{\pgfqpoint{3.865417in}{2.304263in}}%
\pgfpathlineto{\pgfqpoint{3.865617in}{2.302864in}}%
\pgfpathlineto{\pgfqpoint{3.866732in}{2.297404in}}%
\pgfpathlineto{\pgfqpoint{3.865724in}{2.303355in}}%
\pgfpathlineto{\pgfqpoint{3.866946in}{2.299007in}}%
\pgfpathlineto{\pgfqpoint{3.867345in}{2.300900in}}%
\pgfpathlineto{\pgfqpoint{3.868027in}{2.297772in}}%
\pgfpathlineto{\pgfqpoint{3.868032in}{2.297767in}}%
\pgfpathlineto{\pgfqpoint{3.868158in}{2.299036in}}%
\pgfpathlineto{\pgfqpoint{3.868280in}{2.299532in}}%
\pgfpathlineto{\pgfqpoint{3.869020in}{2.296486in}}%
\pgfpathlineto{\pgfqpoint{3.869259in}{2.298461in}}%
\pgfpathlineto{\pgfqpoint{3.870578in}{2.293034in}}%
\pgfpathlineto{\pgfqpoint{3.869468in}{2.299182in}}%
\pgfpathlineto{\pgfqpoint{3.870690in}{2.294084in}}%
\pgfpathlineto{\pgfqpoint{3.872136in}{2.302782in}}%
\pgfpathlineto{\pgfqpoint{3.870967in}{2.293755in}}%
\pgfpathlineto{\pgfqpoint{3.872185in}{2.302265in}}%
\pgfpathlineto{\pgfqpoint{3.872808in}{2.300356in}}%
\pgfpathlineto{\pgfqpoint{3.873178in}{2.303353in}}%
\pgfpathlineto{\pgfqpoint{3.873309in}{2.301916in}}%
\pgfpathlineto{\pgfqpoint{3.874297in}{2.304199in}}%
\pgfpathlineto{\pgfqpoint{3.873543in}{2.301186in}}%
\pgfpathlineto{\pgfqpoint{3.874638in}{2.303331in}}%
\pgfpathlineto{\pgfqpoint{3.874682in}{2.303080in}}%
\pgfpathlineto{\pgfqpoint{3.875578in}{2.305813in}}%
\pgfpathlineto{\pgfqpoint{3.875695in}{2.305717in}}%
\pgfpathlineto{\pgfqpoint{3.875719in}{2.306008in}}%
\pgfpathlineto{\pgfqpoint{3.876464in}{2.301399in}}%
\pgfpathlineto{\pgfqpoint{3.876503in}{2.301536in}}%
\pgfpathlineto{\pgfqpoint{3.876508in}{2.301350in}}%
\pgfpathlineto{\pgfqpoint{3.877150in}{2.306060in}}%
\pgfpathlineto{\pgfqpoint{3.877550in}{2.303970in}}%
\pgfpathlineto{\pgfqpoint{3.878465in}{2.302264in}}%
\pgfpathlineto{\pgfqpoint{3.877934in}{2.304280in}}%
\pgfpathlineto{\pgfqpoint{3.878660in}{2.303845in}}%
\pgfpathlineto{\pgfqpoint{3.879020in}{2.304286in}}%
\pgfpathlineto{\pgfqpoint{3.879560in}{2.301012in}}%
\pgfpathlineto{\pgfqpoint{3.879711in}{2.301070in}}%
\pgfpathlineto{\pgfqpoint{3.880958in}{2.295640in}}%
\pgfpathlineto{\pgfqpoint{3.881396in}{2.297870in}}%
\pgfpathlineto{\pgfqpoint{3.882072in}{2.300887in}}%
\pgfpathlineto{\pgfqpoint{3.882442in}{2.297581in}}%
\pgfpathlineto{\pgfqpoint{3.882540in}{2.298446in}}%
\pgfpathlineto{\pgfqpoint{3.882988in}{2.297344in}}%
\pgfpathlineto{\pgfqpoint{3.883139in}{2.300197in}}%
\pgfpathlineto{\pgfqpoint{3.883640in}{2.298853in}}%
\pgfpathlineto{\pgfqpoint{3.884663in}{2.300430in}}%
\pgfpathlineto{\pgfqpoint{3.884273in}{2.297029in}}%
\pgfpathlineto{\pgfqpoint{3.884779in}{2.299693in}}%
\pgfpathlineto{\pgfqpoint{3.886138in}{2.294071in}}%
\pgfpathlineto{\pgfqpoint{3.885169in}{2.300415in}}%
\pgfpathlineto{\pgfqpoint{3.886152in}{2.294342in}}%
\pgfpathlineto{\pgfqpoint{3.887082in}{2.298197in}}%
\pgfpathlineto{\pgfqpoint{3.887301in}{2.296379in}}%
\pgfpathlineto{\pgfqpoint{3.887399in}{2.296167in}}%
\pgfpathlineto{\pgfqpoint{3.888134in}{2.299061in}}%
\pgfpathlineto{\pgfqpoint{3.888168in}{2.298832in}}%
\pgfpathlineto{\pgfqpoint{3.888596in}{2.300794in}}%
\pgfpathlineto{\pgfqpoint{3.889059in}{2.297636in}}%
\pgfpathlineto{\pgfqpoint{3.889258in}{2.298782in}}%
\pgfpathlineto{\pgfqpoint{3.889979in}{2.298235in}}%
\pgfpathlineto{\pgfqpoint{3.890217in}{2.301289in}}%
\pgfpathlineto{\pgfqpoint{3.890339in}{2.300813in}}%
\pgfpathlineto{\pgfqpoint{3.890344in}{2.300771in}}%
\pgfpathlineto{\pgfqpoint{3.890914in}{2.304827in}}%
\pgfpathlineto{\pgfqpoint{3.890992in}{2.304350in}}%
\pgfpathlineto{\pgfqpoint{3.891021in}{2.305195in}}%
\pgfpathlineto{\pgfqpoint{3.891839in}{2.301906in}}%
\pgfpathlineto{\pgfqpoint{3.892087in}{2.303374in}}%
\pgfpathlineto{\pgfqpoint{3.893197in}{2.301491in}}%
\pgfpathlineto{\pgfqpoint{3.892540in}{2.306117in}}%
\pgfpathlineto{\pgfqpoint{3.893251in}{2.302212in}}%
\pgfpathlineto{\pgfqpoint{3.894288in}{2.303867in}}%
\pgfpathlineto{\pgfqpoint{3.893489in}{2.301922in}}%
\pgfpathlineto{\pgfqpoint{3.894375in}{2.303073in}}%
\pgfpathlineto{\pgfqpoint{3.895422in}{2.301313in}}%
\pgfpathlineto{\pgfqpoint{3.894619in}{2.304079in}}%
\pgfpathlineto{\pgfqpoint{3.895553in}{2.302025in}}%
\pgfpathlineto{\pgfqpoint{3.896269in}{2.304353in}}%
\pgfpathlineto{\pgfqpoint{3.896683in}{2.303362in}}%
\pgfpathlineto{\pgfqpoint{3.897954in}{2.299766in}}%
\pgfpathlineto{\pgfqpoint{3.896946in}{2.304566in}}%
\pgfpathlineto{\pgfqpoint{3.897958in}{2.299937in}}%
\pgfpathlineto{\pgfqpoint{3.898109in}{2.299079in}}%
\pgfpathlineto{\pgfqpoint{3.898659in}{2.302356in}}%
\pgfpathlineto{\pgfqpoint{3.898835in}{2.301979in}}%
\pgfpathlineto{\pgfqpoint{3.898966in}{2.303419in}}%
\pgfpathlineto{\pgfqpoint{3.899331in}{2.299494in}}%
\pgfpathlineto{\pgfqpoint{3.899920in}{2.301305in}}%
\pgfpathlineto{\pgfqpoint{3.899925in}{2.301052in}}%
\pgfpathlineto{\pgfqpoint{3.900208in}{2.304178in}}%
\pgfpathlineto{\pgfqpoint{3.900982in}{2.303895in}}%
\pgfpathlineto{\pgfqpoint{3.902067in}{2.305237in}}%
\pgfpathlineto{\pgfqpoint{3.901147in}{2.302518in}}%
\pgfpathlineto{\pgfqpoint{3.902160in}{2.304269in}}%
\pgfpathlineto{\pgfqpoint{3.902175in}{2.304685in}}%
\pgfpathlineto{\pgfqpoint{3.902306in}{2.302924in}}%
\pgfpathlineto{\pgfqpoint{3.902880in}{2.303066in}}%
\pgfpathlineto{\pgfqpoint{3.904015in}{2.301845in}}%
\pgfpathlineto{\pgfqpoint{3.903241in}{2.305076in}}%
\pgfpathlineto{\pgfqpoint{3.904137in}{2.302262in}}%
\pgfpathlineto{\pgfqpoint{3.904915in}{2.303242in}}%
\pgfpathlineto{\pgfqpoint{3.904443in}{2.300496in}}%
\pgfpathlineto{\pgfqpoint{3.905247in}{2.302215in}}%
\pgfpathlineto{\pgfqpoint{3.905514in}{2.303088in}}%
\pgfpathlineto{\pgfqpoint{3.905870in}{2.300274in}}%
\pgfpathlineto{\pgfqpoint{3.905977in}{2.300400in}}%
\pgfpathlineto{\pgfqpoint{3.906829in}{2.297118in}}%
\pgfpathlineto{\pgfqpoint{3.906371in}{2.301567in}}%
\pgfpathlineto{\pgfqpoint{3.907121in}{2.298715in}}%
\pgfpathlineto{\pgfqpoint{3.907204in}{2.298540in}}%
\pgfpathlineto{\pgfqpoint{3.908343in}{2.304059in}}%
\pgfpathlineto{\pgfqpoint{3.909443in}{2.301623in}}%
\pgfpathlineto{\pgfqpoint{3.908815in}{2.306489in}}%
\pgfpathlineto{\pgfqpoint{3.909492in}{2.301864in}}%
\pgfpathlineto{\pgfqpoint{3.910261in}{2.300906in}}%
\pgfpathlineto{\pgfqpoint{3.909769in}{2.303351in}}%
\pgfpathlineto{\pgfqpoint{3.910597in}{2.301861in}}%
\pgfpathlineto{\pgfqpoint{3.911371in}{2.303775in}}%
\pgfpathlineto{\pgfqpoint{3.910806in}{2.301296in}}%
\pgfpathlineto{\pgfqpoint{3.911727in}{2.302203in}}%
\pgfpathlineto{\pgfqpoint{3.912471in}{2.299779in}}%
\pgfpathlineto{\pgfqpoint{3.912866in}{2.301127in}}%
\pgfpathlineto{\pgfqpoint{3.914015in}{2.305515in}}%
\pgfpathlineto{\pgfqpoint{3.914024in}{2.305334in}}%
\pgfpathlineto{\pgfqpoint{3.915042in}{2.302110in}}%
\pgfpathlineto{\pgfqpoint{3.914049in}{2.305394in}}%
\pgfpathlineto{\pgfqpoint{3.915237in}{2.303393in}}%
\pgfpathlineto{\pgfqpoint{3.915490in}{2.304777in}}%
\pgfpathlineto{\pgfqpoint{3.915840in}{2.302009in}}%
\pgfpathlineto{\pgfqpoint{3.916332in}{2.302861in}}%
\pgfpathlineto{\pgfqpoint{3.916556in}{2.302367in}}%
\pgfpathlineto{\pgfqpoint{3.917243in}{2.304909in}}%
\pgfpathlineto{\pgfqpoint{3.917403in}{2.304310in}}%
\pgfpathlineto{\pgfqpoint{3.918090in}{2.305699in}}%
\pgfpathlineto{\pgfqpoint{3.917895in}{2.302404in}}%
\pgfpathlineto{\pgfqpoint{3.918474in}{2.302394in}}%
\pgfpathlineto{\pgfqpoint{3.918630in}{2.301667in}}%
\pgfpathlineto{\pgfqpoint{3.919136in}{2.304242in}}%
\pgfpathlineto{\pgfqpoint{3.919497in}{2.303540in}}%
\pgfpathlineto{\pgfqpoint{3.920003in}{2.304455in}}%
\pgfpathlineto{\pgfqpoint{3.920300in}{2.301893in}}%
\pgfpathlineto{\pgfqpoint{3.920597in}{2.303318in}}%
\pgfpathlineto{\pgfqpoint{3.920621in}{2.302464in}}%
\pgfpathlineto{\pgfqpoint{3.921580in}{2.305262in}}%
\pgfpathlineto{\pgfqpoint{3.921692in}{2.303895in}}%
\pgfpathlineto{\pgfqpoint{3.922754in}{2.305607in}}%
\pgfpathlineto{\pgfqpoint{3.922359in}{2.302391in}}%
\pgfpathlineto{\pgfqpoint{3.922817in}{2.304505in}}%
\pgfpathlineto{\pgfqpoint{3.923421in}{2.303192in}}%
\pgfpathlineto{\pgfqpoint{3.923761in}{2.306544in}}%
\pgfpathlineto{\pgfqpoint{3.923912in}{2.305118in}}%
\pgfpathlineto{\pgfqpoint{3.924594in}{2.301264in}}%
\pgfpathlineto{\pgfqpoint{3.924010in}{2.305556in}}%
\pgfpathlineto{\pgfqpoint{3.925178in}{2.302635in}}%
\pgfpathlineto{\pgfqpoint{3.925295in}{2.304379in}}%
\pgfpathlineto{\pgfqpoint{3.925762in}{2.300289in}}%
\pgfpathlineto{\pgfqpoint{3.926269in}{2.301841in}}%
\pgfpathlineto{\pgfqpoint{3.926332in}{2.300958in}}%
\pgfpathlineto{\pgfqpoint{3.927194in}{2.303990in}}%
\pgfpathlineto{\pgfqpoint{3.927320in}{2.303374in}}%
\pgfpathlineto{\pgfqpoint{3.927973in}{2.304587in}}%
\pgfpathlineto{\pgfqpoint{3.927822in}{2.302422in}}%
\pgfpathlineto{\pgfqpoint{3.928367in}{2.303329in}}%
\pgfpathlineto{\pgfqpoint{3.929107in}{2.298722in}}%
\pgfpathlineto{\pgfqpoint{3.929482in}{2.302742in}}%
\pgfpathlineto{\pgfqpoint{3.930227in}{2.305898in}}%
\pgfpathlineto{\pgfqpoint{3.929540in}{2.302326in}}%
\pgfpathlineto{\pgfqpoint{3.930753in}{2.305027in}}%
\pgfpathlineto{\pgfqpoint{3.931376in}{2.300438in}}%
\pgfpathlineto{\pgfqpoint{3.931902in}{2.300738in}}%
\pgfpathlineto{\pgfqpoint{3.932846in}{2.299541in}}%
\pgfpathlineto{\pgfqpoint{3.932418in}{2.303140in}}%
\pgfpathlineto{\pgfqpoint{3.933002in}{2.301315in}}%
\pgfpathlineto{\pgfqpoint{3.933571in}{2.303597in}}%
\pgfpathlineto{\pgfqpoint{3.933313in}{2.300301in}}%
\pgfpathlineto{\pgfqpoint{3.934131in}{2.301648in}}%
\pgfpathlineto{\pgfqpoint{3.934214in}{2.300549in}}%
\pgfpathlineto{\pgfqpoint{3.935008in}{2.303296in}}%
\pgfpathlineto{\pgfqpoint{3.935198in}{2.302320in}}%
\pgfpathlineto{\pgfqpoint{3.935222in}{2.302989in}}%
\pgfpathlineto{\pgfqpoint{3.936249in}{2.298733in}}%
\pgfpathlineto{\pgfqpoint{3.936254in}{2.299007in}}%
\pgfpathlineto{\pgfqpoint{3.937150in}{2.302416in}}%
\pgfpathlineto{\pgfqpoint{3.936332in}{2.298901in}}%
\pgfpathlineto{\pgfqpoint{3.937408in}{2.300779in}}%
\pgfpathlineto{\pgfqpoint{3.937997in}{2.299154in}}%
\pgfpathlineto{\pgfqpoint{3.938416in}{2.302595in}}%
\pgfpathlineto{\pgfqpoint{3.938503in}{2.301663in}}%
\pgfpathlineto{\pgfqpoint{3.939545in}{2.306854in}}%
\pgfpathlineto{\pgfqpoint{3.938693in}{2.301121in}}%
\pgfpathlineto{\pgfqpoint{3.939793in}{2.306029in}}%
\pgfpathlineto{\pgfqpoint{3.940923in}{2.303482in}}%
\pgfpathlineto{\pgfqpoint{3.940548in}{2.308788in}}%
\pgfpathlineto{\pgfqpoint{3.940938in}{2.303639in}}%
\pgfpathlineto{\pgfqpoint{3.941230in}{2.305062in}}%
\pgfpathlineto{\pgfqpoint{3.941809in}{2.301895in}}%
\pgfpathlineto{\pgfqpoint{3.942057in}{2.304046in}}%
\pgfpathlineto{\pgfqpoint{3.943046in}{2.301697in}}%
\pgfpathlineto{\pgfqpoint{3.942457in}{2.304131in}}%
\pgfpathlineto{\pgfqpoint{3.943255in}{2.302240in}}%
\pgfpathlineto{\pgfqpoint{3.943586in}{2.305260in}}%
\pgfpathlineto{\pgfqpoint{3.944282in}{2.302094in}}%
\pgfpathlineto{\pgfqpoint{3.944384in}{2.303045in}}%
\pgfpathlineto{\pgfqpoint{3.945178in}{2.301646in}}%
\pgfpathlineto{\pgfqpoint{3.944759in}{2.304192in}}%
\pgfpathlineto{\pgfqpoint{3.945543in}{2.302962in}}%
\pgfpathlineto{\pgfqpoint{3.946176in}{2.305156in}}%
\pgfpathlineto{\pgfqpoint{3.945767in}{2.301135in}}%
\pgfpathlineto{\pgfqpoint{3.946653in}{2.303097in}}%
\pgfpathlineto{\pgfqpoint{3.946911in}{2.301570in}}%
\pgfpathlineto{\pgfqpoint{3.947729in}{2.303947in}}%
\pgfpathlineto{\pgfqpoint{3.947734in}{2.304158in}}%
\pgfpathlineto{\pgfqpoint{3.948669in}{2.301466in}}%
\pgfpathlineto{\pgfqpoint{3.948829in}{2.303512in}}%
\pgfpathlineto{\pgfqpoint{3.949078in}{2.302282in}}%
\pgfpathlineto{\pgfqpoint{3.949672in}{2.305195in}}%
\pgfpathlineto{\pgfqpoint{3.949915in}{2.304723in}}%
\pgfpathlineto{\pgfqpoint{3.950796in}{2.305597in}}%
\pgfpathlineto{\pgfqpoint{3.950353in}{2.304020in}}%
\pgfpathlineto{\pgfqpoint{3.950894in}{2.304297in}}%
\pgfpathlineto{\pgfqpoint{3.951867in}{2.302260in}}%
\pgfpathlineto{\pgfqpoint{3.951157in}{2.305428in}}%
\pgfpathlineto{\pgfqpoint{3.952043in}{2.303039in}}%
\pgfpathlineto{\pgfqpoint{3.952807in}{2.305063in}}%
\pgfpathlineto{\pgfqpoint{3.952262in}{2.301278in}}%
\pgfpathlineto{\pgfqpoint{3.953123in}{2.302595in}}%
\pgfpathlineto{\pgfqpoint{3.953226in}{2.301878in}}%
\pgfpathlineto{\pgfqpoint{3.953479in}{2.304297in}}%
\pgfpathlineto{\pgfqpoint{3.954243in}{2.302165in}}%
\pgfpathlineto{\pgfqpoint{3.955144in}{2.305017in}}%
\pgfpathlineto{\pgfqpoint{3.955368in}{2.303537in}}%
\pgfpathlineto{\pgfqpoint{3.956332in}{2.301968in}}%
\pgfpathlineto{\pgfqpoint{3.955558in}{2.303698in}}%
\pgfpathlineto{\pgfqpoint{3.956478in}{2.303314in}}%
\pgfpathlineto{\pgfqpoint{3.956989in}{2.303881in}}%
\pgfpathlineto{\pgfqpoint{3.957369in}{2.301460in}}%
\pgfpathlineto{\pgfqpoint{3.957588in}{2.303380in}}%
\pgfpathlineto{\pgfqpoint{3.957875in}{2.302640in}}%
\pgfpathlineto{\pgfqpoint{3.958235in}{2.305084in}}%
\pgfpathlineto{\pgfqpoint{3.958260in}{2.305453in}}%
\pgfpathlineto{\pgfqpoint{3.958683in}{2.300633in}}%
\pgfpathlineto{\pgfqpoint{3.959277in}{2.303120in}}%
\pgfpathlineto{\pgfqpoint{3.959642in}{2.300666in}}%
\pgfpathlineto{\pgfqpoint{3.959438in}{2.303462in}}%
\pgfpathlineto{\pgfqpoint{3.960533in}{2.302410in}}%
\pgfpathlineto{\pgfqpoint{3.961400in}{2.305338in}}%
\pgfpathlineto{\pgfqpoint{3.961634in}{2.302186in}}%
\pgfpathlineto{\pgfqpoint{3.961643in}{2.302513in}}%
\pgfpathlineto{\pgfqpoint{3.963162in}{2.298590in}}%
\pgfpathlineto{\pgfqpoint{3.961838in}{2.302628in}}%
\pgfpathlineto{\pgfqpoint{3.963216in}{2.299563in}}%
\pgfpathlineto{\pgfqpoint{3.963873in}{2.301832in}}%
\pgfpathlineto{\pgfqpoint{3.963459in}{2.297739in}}%
\pgfpathlineto{\pgfqpoint{3.964389in}{2.300049in}}%
\pgfpathlineto{\pgfqpoint{3.964413in}{2.299804in}}%
\pgfpathlineto{\pgfqpoint{3.964900in}{2.303575in}}%
\pgfpathlineto{\pgfqpoint{3.965421in}{2.301979in}}%
\pgfpathlineto{\pgfqpoint{3.965942in}{2.304540in}}%
\pgfpathlineto{\pgfqpoint{3.966229in}{2.300159in}}%
\pgfpathlineto{\pgfqpoint{3.966531in}{2.302045in}}%
\pgfpathlineto{\pgfqpoint{3.967252in}{2.299026in}}%
\pgfpathlineto{\pgfqpoint{3.966614in}{2.303074in}}%
\pgfpathlineto{\pgfqpoint{3.967675in}{2.300714in}}%
\pgfpathlineto{\pgfqpoint{3.968050in}{2.302993in}}%
\pgfpathlineto{\pgfqpoint{3.968722in}{2.299662in}}%
\pgfpathlineto{\pgfqpoint{3.968771in}{2.300016in}}%
\pgfpathlineto{\pgfqpoint{3.969423in}{2.303835in}}%
\pgfpathlineto{\pgfqpoint{3.968815in}{2.299618in}}%
\pgfpathlineto{\pgfqpoint{3.969891in}{2.300504in}}%
\pgfpathlineto{\pgfqpoint{3.970957in}{2.297117in}}%
\pgfpathlineto{\pgfqpoint{3.970421in}{2.302695in}}%
\pgfpathlineto{\pgfqpoint{3.971035in}{2.297633in}}%
\pgfpathlineto{\pgfqpoint{3.972242in}{2.301881in}}%
\pgfpathlineto{\pgfqpoint{3.971453in}{2.297500in}}%
\pgfpathlineto{\pgfqpoint{3.972271in}{2.301732in}}%
\pgfpathlineto{\pgfqpoint{3.973333in}{2.298035in}}%
\pgfpathlineto{\pgfqpoint{3.972700in}{2.303896in}}%
\pgfpathlineto{\pgfqpoint{3.973464in}{2.299792in}}%
\pgfpathlineto{\pgfqpoint{3.974808in}{2.302561in}}%
\pgfpathlineto{\pgfqpoint{3.973634in}{2.298300in}}%
\pgfpathlineto{\pgfqpoint{3.974822in}{2.302211in}}%
\pgfpathlineto{\pgfqpoint{3.975528in}{2.300086in}}%
\pgfpathlineto{\pgfqpoint{3.975134in}{2.302747in}}%
\pgfpathlineto{\pgfqpoint{3.975927in}{2.302310in}}%
\pgfpathlineto{\pgfqpoint{3.976823in}{2.300759in}}%
\pgfpathlineto{\pgfqpoint{3.975966in}{2.302644in}}%
\pgfpathlineto{\pgfqpoint{3.976960in}{2.302904in}}%
\pgfpathlineto{\pgfqpoint{3.978060in}{2.304326in}}%
\pgfpathlineto{\pgfqpoint{3.977276in}{2.301039in}}%
\pgfpathlineto{\pgfqpoint{3.978084in}{2.304212in}}%
\pgfpathlineto{\pgfqpoint{3.978858in}{2.301856in}}%
\pgfpathlineto{\pgfqpoint{3.979092in}{2.304570in}}%
\pgfpathlineto{\pgfqpoint{3.979194in}{2.303741in}}%
\pgfpathlineto{\pgfqpoint{3.979204in}{2.304289in}}%
\pgfpathlineto{\pgfqpoint{3.979618in}{2.300554in}}%
\pgfpathlineto{\pgfqpoint{3.980295in}{2.302709in}}%
\pgfpathlineto{\pgfqpoint{3.980480in}{2.303882in}}%
\pgfpathlineto{\pgfqpoint{3.981078in}{2.300083in}}%
\pgfpathlineto{\pgfqpoint{3.981088in}{2.300220in}}%
\pgfpathlineto{\pgfqpoint{3.981103in}{2.299897in}}%
\pgfpathlineto{\pgfqpoint{3.981955in}{2.304176in}}%
\pgfpathlineto{\pgfqpoint{3.982159in}{2.302756in}}%
\pgfpathlineto{\pgfqpoint{3.983405in}{2.304328in}}%
\pgfpathlineto{\pgfqpoint{3.982583in}{2.301909in}}%
\pgfpathlineto{\pgfqpoint{3.983615in}{2.303614in}}%
\pgfpathlineto{\pgfqpoint{3.984068in}{2.301945in}}%
\pgfpathlineto{\pgfqpoint{3.984535in}{2.306166in}}%
\pgfpathlineto{\pgfqpoint{3.984700in}{2.304900in}}%
\pgfpathlineto{\pgfqpoint{3.985027in}{2.306297in}}%
\pgfpathlineto{\pgfqpoint{3.985514in}{2.302213in}}%
\pgfpathlineto{\pgfqpoint{3.985684in}{2.302802in}}%
\pgfpathlineto{\pgfqpoint{3.985728in}{2.302465in}}%
\pgfpathlineto{\pgfqpoint{3.986103in}{2.305155in}}%
\pgfpathlineto{\pgfqpoint{3.986755in}{2.303515in}}%
\pgfpathlineto{\pgfqpoint{3.987865in}{2.304397in}}%
\pgfpathlineto{\pgfqpoint{3.986896in}{2.302192in}}%
\pgfpathlineto{\pgfqpoint{3.987870in}{2.304353in}}%
\pgfpathlineto{\pgfqpoint{3.989063in}{2.301240in}}%
\pgfpathlineto{\pgfqpoint{3.987933in}{2.304645in}}%
\pgfpathlineto{\pgfqpoint{3.989072in}{2.301439in}}%
\pgfpathlineto{\pgfqpoint{3.989287in}{2.299414in}}%
\pgfpathlineto{\pgfqpoint{3.990080in}{2.302295in}}%
\pgfpathlineto{\pgfqpoint{3.990163in}{2.301673in}}%
\pgfpathlineto{\pgfqpoint{3.990484in}{2.303522in}}%
\pgfpathlineto{\pgfqpoint{3.990781in}{2.300198in}}%
\pgfpathlineto{\pgfqpoint{3.991258in}{2.301468in}}%
\pgfpathlineto{\pgfqpoint{3.991721in}{2.298042in}}%
\pgfpathlineto{\pgfqpoint{3.992373in}{2.300717in}}%
\pgfpathlineto{\pgfqpoint{3.992412in}{2.301426in}}%
\pgfpathlineto{\pgfqpoint{3.992957in}{2.297628in}}%
\pgfpathlineto{\pgfqpoint{3.993464in}{2.299525in}}%
\pgfpathlineto{\pgfqpoint{3.994447in}{2.303733in}}%
\pgfpathlineto{\pgfqpoint{3.993775in}{2.297619in}}%
\pgfpathlineto{\pgfqpoint{3.994691in}{2.302494in}}%
\pgfpathlineto{\pgfqpoint{3.994808in}{2.302191in}}%
\pgfpathlineto{\pgfqpoint{3.995153in}{2.304166in}}%
\pgfpathlineto{\pgfqpoint{3.995251in}{2.304124in}}%
\pgfpathlineto{\pgfqpoint{3.995255in}{2.304354in}}%
\pgfpathlineto{\pgfqpoint{3.995679in}{2.302531in}}%
\pgfpathlineto{\pgfqpoint{3.996346in}{2.303414in}}%
\pgfpathlineto{\pgfqpoint{3.996945in}{2.303994in}}%
\pgfpathlineto{\pgfqpoint{3.996750in}{2.301226in}}%
\pgfpathlineto{\pgfqpoint{3.997232in}{2.302497in}}%
\pgfpathlineto{\pgfqpoint{3.998342in}{2.301069in}}%
\pgfpathlineto{\pgfqpoint{3.997855in}{2.303041in}}%
\pgfpathlineto{\pgfqpoint{3.998362in}{2.301169in}}%
\pgfpathlineto{\pgfqpoint{3.999345in}{2.304940in}}%
\pgfpathlineto{\pgfqpoint{3.998415in}{2.300614in}}%
\pgfpathlineto{\pgfqpoint{3.999759in}{2.303529in}}%
\pgfpathlineto{\pgfqpoint{4.000090in}{2.301205in}}%
\pgfpathlineto{\pgfqpoint{4.000582in}{2.305134in}}%
\pgfpathlineto{\pgfqpoint{4.000854in}{2.303981in}}%
\pgfpathlineto{\pgfqpoint{4.001112in}{2.305465in}}%
\pgfpathlineto{\pgfqpoint{4.001701in}{2.303153in}}%
\pgfpathlineto{\pgfqpoint{4.001959in}{2.304149in}}%
\pgfpathlineto{\pgfqpoint{4.002159in}{2.302669in}}%
\pgfpathlineto{\pgfqpoint{4.002514in}{2.305843in}}%
\pgfpathlineto{\pgfqpoint{4.003079in}{2.303647in}}%
\pgfpathlineto{\pgfqpoint{4.003254in}{2.305372in}}%
\pgfpathlineto{\pgfqpoint{4.003780in}{2.302468in}}%
\pgfpathlineto{\pgfqpoint{4.004170in}{2.300510in}}%
\pgfpathlineto{\pgfqpoint{4.004510in}{2.302699in}}%
\pgfpathlineto{\pgfqpoint{4.004895in}{2.301590in}}%
\pgfpathlineto{\pgfqpoint{4.005226in}{2.302144in}}%
\pgfpathlineto{\pgfqpoint{4.005523in}{2.299851in}}%
\pgfpathlineto{\pgfqpoint{4.005995in}{2.301210in}}%
\pgfpathlineto{\pgfqpoint{4.006107in}{2.300316in}}%
\pgfpathlineto{\pgfqpoint{4.006989in}{2.304991in}}%
\pgfpathlineto{\pgfqpoint{4.007646in}{2.305797in}}%
\pgfpathlineto{\pgfqpoint{4.007991in}{2.302008in}}%
\pgfpathlineto{\pgfqpoint{4.008030in}{2.302557in}}%
\pgfpathlineto{\pgfqpoint{4.008839in}{2.300098in}}%
\pgfpathlineto{\pgfqpoint{4.008147in}{2.303136in}}%
\pgfpathlineto{\pgfqpoint{4.009179in}{2.301447in}}%
\pgfpathlineto{\pgfqpoint{4.010090in}{2.306202in}}%
\pgfpathlineto{\pgfqpoint{4.009267in}{2.300916in}}%
\pgfpathlineto{\pgfqpoint{4.010435in}{2.304684in}}%
\pgfpathlineto{\pgfqpoint{4.011219in}{2.301050in}}%
\pgfpathlineto{\pgfqpoint{4.010465in}{2.305043in}}%
\pgfpathlineto{\pgfqpoint{4.011716in}{2.302172in}}%
\pgfpathlineto{\pgfqpoint{4.012850in}{2.304651in}}%
\pgfpathlineto{\pgfqpoint{4.011857in}{2.300968in}}%
\pgfpathlineto{\pgfqpoint{4.012889in}{2.304321in}}%
\pgfpathlineto{\pgfqpoint{4.013147in}{2.301499in}}%
\pgfpathlineto{\pgfqpoint{4.012952in}{2.304536in}}%
\pgfpathlineto{\pgfqpoint{4.014033in}{2.303165in}}%
\pgfpathlineto{\pgfqpoint{4.014559in}{2.304110in}}%
\pgfpathlineto{\pgfqpoint{4.014471in}{2.302739in}}%
\pgfpathlineto{\pgfqpoint{4.015143in}{2.303346in}}%
\pgfpathlineto{\pgfqpoint{4.016253in}{2.302709in}}%
\pgfpathlineto{\pgfqpoint{4.016024in}{2.304854in}}%
\pgfpathlineto{\pgfqpoint{4.016263in}{2.303079in}}%
\pgfpathlineto{\pgfqpoint{4.016321in}{2.303935in}}%
\pgfpathlineto{\pgfqpoint{4.017305in}{2.299510in}}%
\pgfpathlineto{\pgfqpoint{4.017310in}{2.299612in}}%
\pgfpathlineto{\pgfqpoint{4.017319in}{2.299228in}}%
\pgfpathlineto{\pgfqpoint{4.018279in}{2.302302in}}%
\pgfpathlineto{\pgfqpoint{4.018395in}{2.300837in}}%
\pgfpathlineto{\pgfqpoint{4.019389in}{2.303946in}}%
\pgfpathlineto{\pgfqpoint{4.018444in}{2.300337in}}%
\pgfpathlineto{\pgfqpoint{4.019559in}{2.302773in}}%
\pgfpathlineto{\pgfqpoint{4.020187in}{2.301944in}}%
\pgfpathlineto{\pgfqpoint{4.019632in}{2.304178in}}%
\pgfpathlineto{\pgfqpoint{4.020650in}{2.303822in}}%
\pgfpathlineto{\pgfqpoint{4.021725in}{2.301414in}}%
\pgfpathlineto{\pgfqpoint{4.021000in}{2.306792in}}%
\pgfpathlineto{\pgfqpoint{4.021823in}{2.301814in}}%
\pgfpathlineto{\pgfqpoint{4.022845in}{2.305398in}}%
\pgfpathlineto{\pgfqpoint{4.022095in}{2.301422in}}%
\pgfpathlineto{\pgfqpoint{4.022957in}{2.304470in}}%
\pgfpathlineto{\pgfqpoint{4.023838in}{2.303110in}}%
\pgfpathlineto{\pgfqpoint{4.023434in}{2.307129in}}%
\pgfpathlineto{\pgfqpoint{4.024062in}{2.304440in}}%
\pgfpathlineto{\pgfqpoint{4.024082in}{2.304786in}}%
\pgfpathlineto{\pgfqpoint{4.024622in}{2.302654in}}%
\pgfpathlineto{\pgfqpoint{4.025143in}{2.303762in}}%
\pgfpathlineto{\pgfqpoint{4.026093in}{2.301375in}}%
\pgfpathlineto{\pgfqpoint{4.025430in}{2.304521in}}%
\pgfpathlineto{\pgfqpoint{4.026287in}{2.301734in}}%
\pgfpathlineto{\pgfqpoint{4.026730in}{2.303359in}}%
\pgfpathlineto{\pgfqpoint{4.027271in}{2.299219in}}%
\pgfpathlineto{\pgfqpoint{4.027373in}{2.299704in}}%
\pgfpathlineto{\pgfqpoint{4.027446in}{2.299450in}}%
\pgfpathlineto{\pgfqpoint{4.028152in}{2.303743in}}%
\pgfpathlineto{\pgfqpoint{4.028429in}{2.302155in}}%
\pgfpathlineto{\pgfqpoint{4.030226in}{2.297302in}}%
\pgfpathlineto{\pgfqpoint{4.028795in}{2.302476in}}%
\pgfpathlineto{\pgfqpoint{4.030445in}{2.297873in}}%
\pgfpathlineto{\pgfqpoint{4.031117in}{2.301602in}}%
\pgfpathlineto{\pgfqpoint{4.031604in}{2.299518in}}%
\pgfpathlineto{\pgfqpoint{4.032183in}{2.301252in}}%
\pgfpathlineto{\pgfqpoint{4.032446in}{2.297341in}}%
\pgfpathlineto{\pgfqpoint{4.032626in}{2.298304in}}%
\pgfpathlineto{\pgfqpoint{4.033702in}{2.295789in}}%
\pgfpathlineto{\pgfqpoint{4.033225in}{2.299869in}}%
\pgfpathlineto{\pgfqpoint{4.033765in}{2.295845in}}%
\pgfpathlineto{\pgfqpoint{4.034836in}{2.297255in}}%
\pgfpathlineto{\pgfqpoint{4.034354in}{2.294496in}}%
\pgfpathlineto{\pgfqpoint{4.034890in}{2.296111in}}%
\pgfpathlineto{\pgfqpoint{4.034895in}{2.295947in}}%
\pgfpathlineto{\pgfqpoint{4.035898in}{2.298702in}}%
\pgfpathlineto{\pgfqpoint{4.035966in}{2.297022in}}%
\pgfpathlineto{\pgfqpoint{4.037056in}{2.298332in}}%
\pgfpathlineto{\pgfqpoint{4.036423in}{2.294459in}}%
\pgfpathlineto{\pgfqpoint{4.037076in}{2.297216in}}%
\pgfpathlineto{\pgfqpoint{4.038079in}{2.295457in}}%
\pgfpathlineto{\pgfqpoint{4.037470in}{2.298287in}}%
\pgfpathlineto{\pgfqpoint{4.038181in}{2.297112in}}%
\pgfpathlineto{\pgfqpoint{4.038196in}{2.297577in}}%
\pgfpathlineto{\pgfqpoint{4.038936in}{2.293121in}}%
\pgfpathlineto{\pgfqpoint{4.039242in}{2.293859in}}%
\pgfpathlineto{\pgfqpoint{4.040187in}{2.297182in}}%
\pgfpathlineto{\pgfqpoint{4.039379in}{2.292842in}}%
\pgfpathlineto{\pgfqpoint{4.040435in}{2.295608in}}%
\pgfpathlineto{\pgfqpoint{4.040679in}{2.295420in}}%
\pgfpathlineto{\pgfqpoint{4.041156in}{2.300658in}}%
\pgfpathlineto{\pgfqpoint{4.041404in}{2.298476in}}%
\pgfpathlineto{\pgfqpoint{4.042397in}{2.304658in}}%
\pgfpathlineto{\pgfqpoint{4.041423in}{2.298222in}}%
\pgfpathlineto{\pgfqpoint{4.042602in}{2.303757in}}%
\pgfpathlineto{\pgfqpoint{4.043639in}{2.299376in}}%
\pgfpathlineto{\pgfqpoint{4.042762in}{2.304026in}}%
\pgfpathlineto{\pgfqpoint{4.043804in}{2.300123in}}%
\pgfpathlineto{\pgfqpoint{4.043936in}{2.300920in}}%
\pgfpathlineto{\pgfqpoint{4.044598in}{2.296392in}}%
\pgfpathlineto{\pgfqpoint{4.044885in}{2.298275in}}%
\pgfpathlineto{\pgfqpoint{4.044895in}{2.298749in}}%
\pgfpathlineto{\pgfqpoint{4.045805in}{2.295526in}}%
\pgfpathlineto{\pgfqpoint{4.045829in}{2.295664in}}%
\pgfpathlineto{\pgfqpoint{4.045844in}{2.295552in}}%
\pgfpathlineto{\pgfqpoint{4.046448in}{2.301441in}}%
\pgfpathlineto{\pgfqpoint{4.046672in}{2.300288in}}%
\pgfpathlineto{\pgfqpoint{4.047183in}{2.301630in}}%
\pgfpathlineto{\pgfqpoint{4.046691in}{2.299966in}}%
\pgfpathlineto{\pgfqpoint{4.047188in}{2.301582in}}%
\pgfpathlineto{\pgfqpoint{4.047889in}{2.305547in}}%
\pgfpathlineto{\pgfqpoint{4.047319in}{2.300659in}}%
\pgfpathlineto{\pgfqpoint{4.048332in}{2.304030in}}%
\pgfpathlineto{\pgfqpoint{4.048658in}{2.301763in}}%
\pgfpathlineto{\pgfqpoint{4.048361in}{2.304545in}}%
\pgfpathlineto{\pgfqpoint{4.049505in}{2.303286in}}%
\pgfpathlineto{\pgfqpoint{4.049749in}{2.305318in}}%
\pgfpathlineto{\pgfqpoint{4.050532in}{2.301380in}}%
\pgfpathlineto{\pgfqpoint{4.050591in}{2.301609in}}%
\pgfpathlineto{\pgfqpoint{4.051277in}{2.299591in}}%
\pgfpathlineto{\pgfqpoint{4.051662in}{2.302635in}}%
\pgfpathlineto{\pgfqpoint{4.052791in}{2.304650in}}%
\pgfpathlineto{\pgfqpoint{4.051740in}{2.301849in}}%
\pgfpathlineto{\pgfqpoint{4.052801in}{2.304501in}}%
\pgfpathlineto{\pgfqpoint{4.053979in}{2.301491in}}%
\pgfpathlineto{\pgfqpoint{4.052864in}{2.304866in}}%
\pgfpathlineto{\pgfqpoint{4.053984in}{2.301611in}}%
\pgfpathlineto{\pgfqpoint{4.055060in}{2.306918in}}%
\pgfpathlineto{\pgfqpoint{4.054140in}{2.300299in}}%
\pgfpathlineto{\pgfqpoint{4.055153in}{2.305935in}}%
\pgfpathlineto{\pgfqpoint{4.056521in}{2.299027in}}%
\pgfpathlineto{\pgfqpoint{4.056618in}{2.300822in}}%
\pgfpathlineto{\pgfqpoint{4.057085in}{2.304810in}}%
\pgfpathlineto{\pgfqpoint{4.057767in}{2.303324in}}%
\pgfpathlineto{\pgfqpoint{4.058926in}{2.302081in}}%
\pgfpathlineto{\pgfqpoint{4.058132in}{2.305210in}}%
\pgfpathlineto{\pgfqpoint{4.058999in}{2.302333in}}%
\pgfpathlineto{\pgfqpoint{4.059286in}{2.303850in}}%
\pgfpathlineto{\pgfqpoint{4.059826in}{2.300572in}}%
\pgfpathlineto{\pgfqpoint{4.060070in}{2.301717in}}%
\pgfpathlineto{\pgfqpoint{4.060888in}{2.301028in}}%
\pgfpathlineto{\pgfqpoint{4.060639in}{2.304006in}}%
\pgfpathlineto{\pgfqpoint{4.061155in}{2.302875in}}%
\pgfpathlineto{\pgfqpoint{4.061462in}{2.304521in}}%
\pgfpathlineto{\pgfqpoint{4.062144in}{2.300849in}}%
\pgfpathlineto{\pgfqpoint{4.062178in}{2.301109in}}%
\pgfpathlineto{\pgfqpoint{4.063293in}{2.300036in}}%
\pgfpathlineto{\pgfqpoint{4.062699in}{2.304121in}}%
\pgfpathlineto{\pgfqpoint{4.063302in}{2.300346in}}%
\pgfpathlineto{\pgfqpoint{4.063531in}{2.299491in}}%
\pgfpathlineto{\pgfqpoint{4.063945in}{2.301648in}}%
\pgfpathlineto{\pgfqpoint{4.064408in}{2.300464in}}%
\pgfpathlineto{\pgfqpoint{4.064427in}{2.300725in}}%
\pgfpathlineto{\pgfqpoint{4.065318in}{2.295900in}}%
\pgfpathlineto{\pgfqpoint{4.065362in}{2.296545in}}%
\pgfpathlineto{\pgfqpoint{4.065625in}{2.295971in}}%
\pgfpathlineto{\pgfqpoint{4.066136in}{2.299286in}}%
\pgfpathlineto{\pgfqpoint{4.066448in}{2.296886in}}%
\pgfpathlineto{\pgfqpoint{4.066998in}{2.298533in}}%
\pgfpathlineto{\pgfqpoint{4.067402in}{2.295209in}}%
\pgfpathlineto{\pgfqpoint{4.067553in}{2.296473in}}%
\pgfpathlineto{\pgfqpoint{4.067592in}{2.297588in}}%
\pgfpathlineto{\pgfqpoint{4.068298in}{2.293995in}}%
\pgfpathlineto{\pgfqpoint{4.068541in}{2.294948in}}%
\pgfpathlineto{\pgfqpoint{4.069291in}{2.289795in}}%
\pgfpathlineto{\pgfqpoint{4.068575in}{2.295052in}}%
\pgfpathlineto{\pgfqpoint{4.069739in}{2.291897in}}%
\pgfpathlineto{\pgfqpoint{4.070912in}{2.297880in}}%
\pgfpathlineto{\pgfqpoint{4.072041in}{2.301396in}}%
\pgfpathlineto{\pgfqpoint{4.071564in}{2.297585in}}%
\pgfpathlineto{\pgfqpoint{4.072076in}{2.301033in}}%
\pgfpathlineto{\pgfqpoint{4.072105in}{2.301565in}}%
\pgfpathlineto{\pgfqpoint{4.073025in}{2.297767in}}%
\pgfpathlineto{\pgfqpoint{4.073035in}{2.297471in}}%
\pgfpathlineto{\pgfqpoint{4.073862in}{2.301671in}}%
\pgfpathlineto{\pgfqpoint{4.074072in}{2.300492in}}%
\pgfpathlineto{\pgfqpoint{4.075148in}{2.305051in}}%
\pgfpathlineto{\pgfqpoint{4.075357in}{2.304116in}}%
\pgfpathlineto{\pgfqpoint{4.075659in}{2.300831in}}%
\pgfpathlineto{\pgfqpoint{4.075668in}{2.301058in}}%
\pgfpathlineto{\pgfqpoint{4.075678in}{2.300647in}}%
\pgfpathlineto{\pgfqpoint{4.076632in}{2.305035in}}%
\pgfpathlineto{\pgfqpoint{4.076735in}{2.304027in}}%
\pgfpathlineto{\pgfqpoint{4.077304in}{2.305551in}}%
\pgfpathlineto{\pgfqpoint{4.076983in}{2.302387in}}%
\pgfpathlineto{\pgfqpoint{4.077407in}{2.303042in}}%
\pgfpathlineto{\pgfqpoint{4.078117in}{2.298545in}}%
\pgfpathlineto{\pgfqpoint{4.077465in}{2.303328in}}%
\pgfpathlineto{\pgfqpoint{4.078536in}{2.301225in}}%
\pgfpathlineto{\pgfqpoint{4.079140in}{2.298615in}}%
\pgfpathlineto{\pgfqpoint{4.078629in}{2.301589in}}%
\pgfpathlineto{\pgfqpoint{4.079627in}{2.301442in}}%
\pgfpathlineto{\pgfqpoint{4.080298in}{2.302374in}}%
\pgfpathlineto{\pgfqpoint{4.080644in}{2.299511in}}%
\pgfpathlineto{\pgfqpoint{4.080707in}{2.299584in}}%
\pgfpathlineto{\pgfqpoint{4.080990in}{2.296214in}}%
\pgfpathlineto{\pgfqpoint{4.081847in}{2.297422in}}%
\pgfpathlineto{\pgfqpoint{4.082874in}{2.299137in}}%
\pgfpathlineto{\pgfqpoint{4.082465in}{2.295215in}}%
\pgfpathlineto{\pgfqpoint{4.082976in}{2.298023in}}%
\pgfpathlineto{\pgfqpoint{4.083356in}{2.294877in}}%
\pgfpathlineto{\pgfqpoint{4.084140in}{2.294903in}}%
\pgfpathlineto{\pgfqpoint{4.085269in}{2.298521in}}%
\pgfpathlineto{\pgfqpoint{4.084310in}{2.294264in}}%
\pgfpathlineto{\pgfqpoint{4.085303in}{2.297827in}}%
\pgfpathlineto{\pgfqpoint{4.086433in}{2.294800in}}%
\pgfpathlineto{\pgfqpoint{4.085362in}{2.298375in}}%
\pgfpathlineto{\pgfqpoint{4.086467in}{2.295010in}}%
\pgfpathlineto{\pgfqpoint{4.087533in}{2.298775in}}%
\pgfpathlineto{\pgfqpoint{4.087611in}{2.297835in}}%
\pgfpathlineto{\pgfqpoint{4.088492in}{2.300552in}}%
\pgfpathlineto{\pgfqpoint{4.088774in}{2.298477in}}%
\pgfpathlineto{\pgfqpoint{4.089329in}{2.296911in}}%
\pgfpathlineto{\pgfqpoint{4.089101in}{2.299493in}}%
\pgfpathlineto{\pgfqpoint{4.089909in}{2.297536in}}%
\pgfpathlineto{\pgfqpoint{4.089919in}{2.297512in}}%
\pgfpathlineto{\pgfqpoint{4.090069in}{2.300222in}}%
\pgfpathlineto{\pgfqpoint{4.090079in}{2.300213in}}%
\pgfpathlineto{\pgfqpoint{4.090201in}{2.300642in}}%
\pgfpathlineto{\pgfqpoint{4.090771in}{2.297938in}}%
\pgfpathlineto{\pgfqpoint{4.091180in}{2.299478in}}%
\pgfpathlineto{\pgfqpoint{4.091277in}{2.298097in}}%
\pgfpathlineto{\pgfqpoint{4.091949in}{2.303727in}}%
\pgfpathlineto{\pgfqpoint{4.092168in}{2.301614in}}%
\pgfpathlineto{\pgfqpoint{4.092192in}{2.301976in}}%
\pgfpathlineto{\pgfqpoint{4.092947in}{2.299382in}}%
\pgfpathlineto{\pgfqpoint{4.093054in}{2.299732in}}%
\pgfpathlineto{\pgfqpoint{4.093064in}{2.299481in}}%
\pgfpathlineto{\pgfqpoint{4.093550in}{2.304913in}}%
\pgfpathlineto{\pgfqpoint{4.094096in}{2.303146in}}%
\pgfpathlineto{\pgfqpoint{4.095152in}{2.304275in}}%
\pgfpathlineto{\pgfqpoint{4.094777in}{2.302309in}}%
\pgfpathlineto{\pgfqpoint{4.095215in}{2.303897in}}%
\pgfpathlineto{\pgfqpoint{4.095605in}{2.301373in}}%
\pgfpathlineto{\pgfqpoint{4.095293in}{2.304465in}}%
\pgfpathlineto{\pgfqpoint{4.096364in}{2.302552in}}%
\pgfpathlineto{\pgfqpoint{4.097075in}{2.304413in}}%
\pgfpathlineto{\pgfqpoint{4.097338in}{2.302131in}}%
\pgfpathlineto{\pgfqpoint{4.097484in}{2.303225in}}%
\pgfpathlineto{\pgfqpoint{4.097601in}{2.304573in}}%
\pgfpathlineto{\pgfqpoint{4.097961in}{2.301726in}}%
\pgfpathlineto{\pgfqpoint{4.098497in}{2.301796in}}%
\pgfpathlineto{\pgfqpoint{4.098808in}{2.304228in}}%
\pgfpathlineto{\pgfqpoint{4.099276in}{2.300412in}}%
\pgfpathlineto{\pgfqpoint{4.099422in}{2.300641in}}%
\pgfpathlineto{\pgfqpoint{4.099743in}{2.298888in}}%
\pgfpathlineto{\pgfqpoint{4.100357in}{2.301886in}}%
\pgfpathlineto{\pgfqpoint{4.100488in}{2.301723in}}%
\pgfpathlineto{\pgfqpoint{4.100576in}{2.299679in}}%
\pgfpathlineto{\pgfqpoint{4.101058in}{2.303794in}}%
\pgfpathlineto{\pgfqpoint{4.101155in}{2.303620in}}%
\pgfpathlineto{\pgfqpoint{4.101476in}{2.305451in}}%
\pgfpathlineto{\pgfqpoint{4.102207in}{2.301161in}}%
\pgfpathlineto{\pgfqpoint{4.102212in}{2.301231in}}%
\pgfpathlineto{\pgfqpoint{4.102728in}{2.301629in}}%
\pgfpathlineto{\pgfqpoint{4.103361in}{2.298441in}}%
\pgfpathlineto{\pgfqpoint{4.103380in}{2.298907in}}%
\pgfpathlineto{\pgfqpoint{4.104329in}{2.295241in}}%
\pgfpathlineto{\pgfqpoint{4.104422in}{2.295897in}}%
\pgfpathlineto{\pgfqpoint{4.104870in}{2.299182in}}%
\pgfpathlineto{\pgfqpoint{4.104446in}{2.295515in}}%
\pgfpathlineto{\pgfqpoint{4.105585in}{2.296906in}}%
\pgfpathlineto{\pgfqpoint{4.106106in}{2.295553in}}%
\pgfpathlineto{\pgfqpoint{4.106311in}{2.299118in}}%
\pgfpathlineto{\pgfqpoint{4.106642in}{2.298134in}}%
\pgfpathlineto{\pgfqpoint{4.106666in}{2.298647in}}%
\pgfpathlineto{\pgfqpoint{4.107664in}{2.295349in}}%
\pgfpathlineto{\pgfqpoint{4.107708in}{2.296156in}}%
\pgfpathlineto{\pgfqpoint{4.107713in}{2.296097in}}%
\pgfpathlineto{\pgfqpoint{4.108404in}{2.300228in}}%
\pgfpathlineto{\pgfqpoint{4.108516in}{2.299658in}}%
\pgfpathlineto{\pgfqpoint{4.108915in}{2.302500in}}%
\pgfpathlineto{\pgfqpoint{4.108594in}{2.299199in}}%
\pgfpathlineto{\pgfqpoint{4.109670in}{2.301506in}}%
\pgfpathlineto{\pgfqpoint{4.110167in}{2.303823in}}%
\pgfpathlineto{\pgfqpoint{4.110172in}{2.303581in}}%
\pgfpathlineto{\pgfqpoint{4.111286in}{2.305835in}}%
\pgfpathlineto{\pgfqpoint{4.110946in}{2.302271in}}%
\pgfpathlineto{\pgfqpoint{4.111311in}{2.305259in}}%
\pgfpathlineto{\pgfqpoint{4.112109in}{2.302236in}}%
\pgfpathlineto{\pgfqpoint{4.111676in}{2.306733in}}%
\pgfpathlineto{\pgfqpoint{4.112455in}{2.303031in}}%
\pgfpathlineto{\pgfqpoint{4.112820in}{2.304375in}}%
\pgfpathlineto{\pgfqpoint{4.113375in}{2.302371in}}%
\pgfpathlineto{\pgfqpoint{4.113589in}{2.303536in}}%
\pgfpathlineto{\pgfqpoint{4.114704in}{2.300507in}}%
\pgfpathlineto{\pgfqpoint{4.113716in}{2.304138in}}%
\pgfpathlineto{\pgfqpoint{4.114899in}{2.302813in}}%
\pgfpathlineto{\pgfqpoint{4.115673in}{2.304051in}}%
\pgfpathlineto{\pgfqpoint{4.115376in}{2.302114in}}%
\pgfpathlineto{\pgfqpoint{4.116023in}{2.303557in}}%
\pgfpathlineto{\pgfqpoint{4.116948in}{2.302621in}}%
\pgfpathlineto{\pgfqpoint{4.116155in}{2.304998in}}%
\pgfpathlineto{\pgfqpoint{4.117138in}{2.303044in}}%
\pgfpathlineto{\pgfqpoint{4.117299in}{2.304511in}}%
\pgfpathlineto{\pgfqpoint{4.117805in}{2.301688in}}%
\pgfpathlineto{\pgfqpoint{4.118248in}{2.303043in}}%
\pgfpathlineto{\pgfqpoint{4.119198in}{2.301214in}}%
\pgfpathlineto{\pgfqpoint{4.118458in}{2.303778in}}%
\pgfpathlineto{\pgfqpoint{4.119373in}{2.301957in}}%
\pgfpathlineto{\pgfqpoint{4.120225in}{2.304683in}}%
\pgfpathlineto{\pgfqpoint{4.119451in}{2.301676in}}%
\pgfpathlineto{\pgfqpoint{4.120493in}{2.302706in}}%
\pgfpathlineto{\pgfqpoint{4.121092in}{2.303615in}}%
\pgfpathlineto{\pgfqpoint{4.121316in}{2.300196in}}%
\pgfpathlineto{\pgfqpoint{4.121549in}{2.301720in}}%
\pgfpathlineto{\pgfqpoint{4.122465in}{2.300520in}}%
\pgfpathlineto{\pgfqpoint{4.121992in}{2.302743in}}%
\pgfpathlineto{\pgfqpoint{4.122684in}{2.301253in}}%
\pgfpathlineto{\pgfqpoint{4.122869in}{2.303174in}}%
\pgfpathlineto{\pgfqpoint{4.123433in}{2.299495in}}%
\pgfpathlineto{\pgfqpoint{4.123784in}{2.300094in}}%
\pgfpathlineto{\pgfqpoint{4.123833in}{2.299487in}}%
\pgfpathlineto{\pgfqpoint{4.124178in}{2.302332in}}%
\pgfpathlineto{\pgfqpoint{4.124704in}{2.301985in}}%
\pgfpathlineto{\pgfqpoint{4.125843in}{2.303633in}}%
\pgfpathlineto{\pgfqpoint{4.125103in}{2.298425in}}%
\pgfpathlineto{\pgfqpoint{4.125853in}{2.303526in}}%
\pgfpathlineto{\pgfqpoint{4.126476in}{2.302472in}}%
\pgfpathlineto{\pgfqpoint{4.126072in}{2.304663in}}%
\pgfpathlineto{\pgfqpoint{4.126973in}{2.303090in}}%
\pgfpathlineto{\pgfqpoint{4.127576in}{2.305221in}}%
\pgfpathlineto{\pgfqpoint{4.127250in}{2.301958in}}%
\pgfpathlineto{\pgfqpoint{4.128102in}{2.303330in}}%
\pgfpathlineto{\pgfqpoint{4.128268in}{2.305621in}}%
\pgfpathlineto{\pgfqpoint{4.129241in}{2.301904in}}%
\pgfpathlineto{\pgfqpoint{4.129728in}{2.302772in}}%
\pgfpathlineto{\pgfqpoint{4.129879in}{2.300772in}}%
\pgfpathlineto{\pgfqpoint{4.130342in}{2.301662in}}%
\pgfpathlineto{\pgfqpoint{4.130697in}{2.300232in}}%
\pgfpathlineto{\pgfqpoint{4.130984in}{2.303587in}}%
\pgfpathlineto{\pgfqpoint{4.131384in}{2.302732in}}%
\pgfpathlineto{\pgfqpoint{4.131987in}{2.305952in}}%
\pgfpathlineto{\pgfqpoint{4.131515in}{2.301911in}}%
\pgfpathlineto{\pgfqpoint{4.132995in}{2.303681in}}%
\pgfpathlineto{\pgfqpoint{4.133555in}{2.302516in}}%
\pgfpathlineto{\pgfqpoint{4.134003in}{2.304438in}}%
\pgfpathlineto{\pgfqpoint{4.134095in}{2.304261in}}%
\pgfpathlineto{\pgfqpoint{4.135386in}{2.298738in}}%
\pgfpathlineto{\pgfqpoint{4.134261in}{2.304621in}}%
\pgfpathlineto{\pgfqpoint{4.135780in}{2.301255in}}%
\pgfpathlineto{\pgfqpoint{4.135926in}{2.302436in}}%
\pgfpathlineto{\pgfqpoint{4.136729in}{2.298596in}}%
\pgfpathlineto{\pgfqpoint{4.136861in}{2.300100in}}%
\pgfpathlineto{\pgfqpoint{4.136968in}{2.302132in}}%
\pgfpathlineto{\pgfqpoint{4.136880in}{2.299674in}}%
\pgfpathlineto{\pgfqpoint{4.136973in}{2.302073in}}%
\pgfpathlineto{\pgfqpoint{4.137377in}{2.302596in}}%
\pgfpathlineto{\pgfqpoint{4.137868in}{2.298987in}}%
\pgfpathlineto{\pgfqpoint{4.138058in}{2.301074in}}%
\pgfpathlineto{\pgfqpoint{4.139129in}{2.298438in}}%
\pgfpathlineto{\pgfqpoint{4.138501in}{2.302302in}}%
\pgfpathlineto{\pgfqpoint{4.139275in}{2.298979in}}%
\pgfpathlineto{\pgfqpoint{4.139748in}{2.299819in}}%
\pgfpathlineto{\pgfqpoint{4.140254in}{2.296870in}}%
\pgfpathlineto{\pgfqpoint{4.140293in}{2.297591in}}%
\pgfpathlineto{\pgfqpoint{4.141349in}{2.293670in}}%
\pgfpathlineto{\pgfqpoint{4.140347in}{2.297740in}}%
\pgfpathlineto{\pgfqpoint{4.141471in}{2.294133in}}%
\pgfpathlineto{\pgfqpoint{4.142581in}{2.300340in}}%
\pgfpathlineto{\pgfqpoint{4.141525in}{2.293602in}}%
\pgfpathlineto{\pgfqpoint{4.142708in}{2.300056in}}%
\pgfpathlineto{\pgfqpoint{4.143545in}{2.298096in}}%
\pgfpathlineto{\pgfqpoint{4.143087in}{2.300546in}}%
\pgfpathlineto{\pgfqpoint{4.143701in}{2.300512in}}%
\pgfpathlineto{\pgfqpoint{4.144509in}{2.303422in}}%
\pgfpathlineto{\pgfqpoint{4.143852in}{2.299877in}}%
\pgfpathlineto{\pgfqpoint{4.144821in}{2.300898in}}%
\pgfpathlineto{\pgfqpoint{4.145955in}{2.302311in}}%
\pgfpathlineto{\pgfqpoint{4.145108in}{2.299263in}}%
\pgfpathlineto{\pgfqpoint{4.146077in}{2.301487in}}%
\pgfpathlineto{\pgfqpoint{4.146364in}{2.300148in}}%
\pgfpathlineto{\pgfqpoint{4.147080in}{2.303468in}}%
\pgfpathlineto{\pgfqpoint{4.147162in}{2.302820in}}%
\pgfpathlineto{\pgfqpoint{4.147440in}{2.301373in}}%
\pgfpathlineto{\pgfqpoint{4.147771in}{2.304164in}}%
\pgfpathlineto{\pgfqpoint{4.148253in}{2.303244in}}%
\pgfpathlineto{\pgfqpoint{4.148579in}{2.305086in}}%
\pgfpathlineto{\pgfqpoint{4.149183in}{2.302119in}}%
\pgfpathlineto{\pgfqpoint{4.149358in}{2.303170in}}%
\pgfpathlineto{\pgfqpoint{4.149855in}{2.305304in}}%
\pgfpathlineto{\pgfqpoint{4.150015in}{2.301985in}}%
\pgfpathlineto{\pgfqpoint{4.150974in}{2.300800in}}%
\pgfpathlineto{\pgfqpoint{4.150346in}{2.305070in}}%
\pgfpathlineto{\pgfqpoint{4.151130in}{2.301752in}}%
\pgfpathlineto{\pgfqpoint{4.151500in}{2.304109in}}%
\pgfpathlineto{\pgfqpoint{4.152221in}{2.301347in}}%
\pgfpathlineto{\pgfqpoint{4.152231in}{2.301550in}}%
\pgfpathlineto{\pgfqpoint{4.152367in}{2.300703in}}%
\pgfpathlineto{\pgfqpoint{4.152883in}{2.304385in}}%
\pgfpathlineto{\pgfqpoint{4.153302in}{2.303157in}}%
\pgfpathlineto{\pgfqpoint{4.154056in}{2.304138in}}%
\pgfpathlineto{\pgfqpoint{4.153711in}{2.301145in}}%
\pgfpathlineto{\pgfqpoint{4.154412in}{2.303092in}}%
\pgfpathlineto{\pgfqpoint{4.155692in}{2.299429in}}%
\pgfpathlineto{\pgfqpoint{4.154455in}{2.303480in}}%
\pgfpathlineto{\pgfqpoint{4.156013in}{2.300366in}}%
\pgfpathlineto{\pgfqpoint{4.156023in}{2.300254in}}%
\pgfpathlineto{\pgfqpoint{4.156992in}{2.302017in}}%
\pgfpathlineto{\pgfqpoint{4.157912in}{2.302727in}}%
\pgfpathlineto{\pgfqpoint{4.157484in}{2.299950in}}%
\pgfpathlineto{\pgfqpoint{4.158034in}{2.300201in}}%
\pgfpathlineto{\pgfqpoint{4.158267in}{2.302232in}}%
\pgfpathlineto{\pgfqpoint{4.158068in}{2.299949in}}%
\pgfpathlineto{\pgfqpoint{4.158272in}{2.302174in}}%
\pgfpathlineto{\pgfqpoint{4.158550in}{2.305360in}}%
\pgfpathlineto{\pgfqpoint{4.159358in}{2.300709in}}%
\pgfpathlineto{\pgfqpoint{4.159363in}{2.300970in}}%
\pgfpathlineto{\pgfqpoint{4.159387in}{2.300725in}}%
\pgfpathlineto{\pgfqpoint{4.159713in}{2.303298in}}%
\pgfpathlineto{\pgfqpoint{4.160371in}{2.302498in}}%
\pgfpathlineto{\pgfqpoint{4.160395in}{2.303189in}}%
\pgfpathlineto{\pgfqpoint{4.160799in}{2.300066in}}%
\pgfpathlineto{\pgfqpoint{4.161461in}{2.301777in}}%
\pgfpathlineto{\pgfqpoint{4.162435in}{2.299744in}}%
\pgfpathlineto{\pgfqpoint{4.161744in}{2.303165in}}%
\pgfpathlineto{\pgfqpoint{4.162586in}{2.300869in}}%
\pgfpathlineto{\pgfqpoint{4.162776in}{2.303535in}}%
\pgfpathlineto{\pgfqpoint{4.163623in}{2.300202in}}%
\pgfpathlineto{\pgfqpoint{4.163662in}{2.300552in}}%
\pgfpathlineto{\pgfqpoint{4.163759in}{2.299702in}}%
\pgfpathlineto{\pgfqpoint{4.164723in}{2.303242in}}%
\pgfpathlineto{\pgfqpoint{4.165346in}{2.304063in}}%
\pgfpathlineto{\pgfqpoint{4.165648in}{2.299890in}}%
\pgfpathlineto{\pgfqpoint{4.165736in}{2.300458in}}%
\pgfpathlineto{\pgfqpoint{4.166305in}{2.297387in}}%
\pgfpathlineto{\pgfqpoint{4.166807in}{2.301784in}}%
\pgfpathlineto{\pgfqpoint{4.167527in}{2.305969in}}%
\pgfpathlineto{\pgfqpoint{4.166958in}{2.300724in}}%
\pgfpathlineto{\pgfqpoint{4.167970in}{2.303445in}}%
\pgfpathlineto{\pgfqpoint{4.168905in}{2.302241in}}%
\pgfpathlineto{\pgfqpoint{4.168569in}{2.305167in}}%
\pgfpathlineto{\pgfqpoint{4.169046in}{2.303509in}}%
\pgfpathlineto{\pgfqpoint{4.169718in}{2.306575in}}%
\pgfpathlineto{\pgfqpoint{4.170186in}{2.305054in}}%
\pgfpathlineto{\pgfqpoint{4.171203in}{2.302207in}}%
\pgfpathlineto{\pgfqpoint{4.170448in}{2.306175in}}%
\pgfpathlineto{\pgfqpoint{4.171485in}{2.304157in}}%
\pgfpathlineto{\pgfqpoint{4.171943in}{2.307015in}}%
\pgfpathlineto{\pgfqpoint{4.172372in}{2.303293in}}%
\pgfpathlineto{\pgfqpoint{4.172586in}{2.303574in}}%
\pgfpathlineto{\pgfqpoint{4.173803in}{2.299090in}}%
\pgfpathlineto{\pgfqpoint{4.173048in}{2.304357in}}%
\pgfpathlineto{\pgfqpoint{4.173852in}{2.299502in}}%
\pgfpathlineto{\pgfqpoint{4.174757in}{2.304868in}}%
\pgfpathlineto{\pgfqpoint{4.175088in}{2.303726in}}%
\pgfpathlineto{\pgfqpoint{4.176398in}{2.299171in}}%
\pgfpathlineto{\pgfqpoint{4.175755in}{2.304777in}}%
\pgfpathlineto{\pgfqpoint{4.176442in}{2.299413in}}%
\pgfpathlineto{\pgfqpoint{4.177410in}{2.302200in}}%
\pgfpathlineto{\pgfqpoint{4.176641in}{2.298045in}}%
\pgfpathlineto{\pgfqpoint{4.177591in}{2.301007in}}%
\pgfpathlineto{\pgfqpoint{4.177595in}{2.300828in}}%
\pgfpathlineto{\pgfqpoint{4.178248in}{2.303192in}}%
\pgfpathlineto{\pgfqpoint{4.178666in}{2.302255in}}%
\pgfpathlineto{\pgfqpoint{4.178832in}{2.303536in}}%
\pgfpathlineto{\pgfqpoint{4.178925in}{2.301401in}}%
\pgfpathlineto{\pgfqpoint{4.179723in}{2.301748in}}%
\pgfpathlineto{\pgfqpoint{4.180385in}{2.298114in}}%
\pgfpathlineto{\pgfqpoint{4.179845in}{2.302896in}}%
\pgfpathlineto{\pgfqpoint{4.180872in}{2.300803in}}%
\pgfpathlineto{\pgfqpoint{4.181276in}{2.301346in}}%
\pgfpathlineto{\pgfqpoint{4.181534in}{2.298887in}}%
\pgfpathlineto{\pgfqpoint{4.181977in}{2.300855in}}%
\pgfpathlineto{\pgfqpoint{4.182089in}{2.298460in}}%
\pgfpathlineto{\pgfqpoint{4.183004in}{2.301494in}}%
\pgfpathlineto{\pgfqpoint{4.183087in}{2.300387in}}%
\pgfpathlineto{\pgfqpoint{4.183360in}{2.302002in}}%
\pgfpathlineto{\pgfqpoint{4.184007in}{2.300001in}}%
\pgfpathlineto{\pgfqpoint{4.184202in}{2.300687in}}%
\pgfpathlineto{\pgfqpoint{4.184840in}{2.299059in}}%
\pgfpathlineto{\pgfqpoint{4.184553in}{2.301893in}}%
\pgfpathlineto{\pgfqpoint{4.185244in}{2.301551in}}%
\pgfpathlineto{\pgfqpoint{4.186524in}{2.303966in}}%
\pgfpathlineto{\pgfqpoint{4.185852in}{2.299912in}}%
\pgfpathlineto{\pgfqpoint{4.186612in}{2.302823in}}%
\pgfpathlineto{\pgfqpoint{4.187123in}{2.300088in}}%
\pgfpathlineto{\pgfqpoint{4.187693in}{2.303739in}}%
\pgfpathlineto{\pgfqpoint{4.187707in}{2.303456in}}%
\pgfpathlineto{\pgfqpoint{4.189036in}{2.300774in}}%
\pgfpathlineto{\pgfqpoint{4.188180in}{2.304439in}}%
\pgfpathlineto{\pgfqpoint{4.189046in}{2.300846in}}%
\pgfpathlineto{\pgfqpoint{4.190034in}{2.296632in}}%
\pgfpathlineto{\pgfqpoint{4.189105in}{2.301789in}}%
\pgfpathlineto{\pgfqpoint{4.190200in}{2.297764in}}%
\pgfpathlineto{\pgfqpoint{4.191037in}{2.303666in}}%
\pgfpathlineto{\pgfqpoint{4.190317in}{2.297201in}}%
\pgfpathlineto{\pgfqpoint{4.191466in}{2.303201in}}%
\pgfpathlineto{\pgfqpoint{4.191904in}{2.301541in}}%
\pgfpathlineto{\pgfqpoint{4.192366in}{2.305063in}}%
\pgfpathlineto{\pgfqpoint{4.192571in}{2.303595in}}%
\pgfpathlineto{\pgfqpoint{4.193330in}{2.302166in}}%
\pgfpathlineto{\pgfqpoint{4.192659in}{2.305016in}}%
\pgfpathlineto{\pgfqpoint{4.193715in}{2.303067in}}%
\pgfpathlineto{\pgfqpoint{4.194319in}{2.304142in}}%
\pgfpathlineto{\pgfqpoint{4.194470in}{2.302540in}}%
\pgfpathlineto{\pgfqpoint{4.194820in}{2.302982in}}%
\pgfpathlineto{\pgfqpoint{4.195312in}{2.301374in}}%
\pgfpathlineto{\pgfqpoint{4.195813in}{2.304257in}}%
\pgfpathlineto{\pgfqpoint{4.195916in}{2.303347in}}%
\pgfpathlineto{\pgfqpoint{4.196485in}{2.305187in}}%
\pgfpathlineto{\pgfqpoint{4.196948in}{2.302725in}}%
\pgfpathlineto{\pgfqpoint{4.197026in}{2.303549in}}%
\pgfpathlineto{\pgfqpoint{4.197493in}{2.302661in}}%
\pgfpathlineto{\pgfqpoint{4.198067in}{2.304849in}}%
\pgfpathlineto{\pgfqpoint{4.198111in}{2.304459in}}%
\pgfpathlineto{\pgfqpoint{4.198121in}{2.304780in}}%
\pgfpathlineto{\pgfqpoint{4.198924in}{2.299833in}}%
\pgfpathlineto{\pgfqpoint{4.199114in}{2.301171in}}%
\pgfpathlineto{\pgfqpoint{4.199153in}{2.300543in}}%
\pgfpathlineto{\pgfqpoint{4.199771in}{2.304935in}}%
\pgfpathlineto{\pgfqpoint{4.200141in}{2.304511in}}%
\pgfpathlineto{\pgfqpoint{4.200234in}{2.305163in}}%
\pgfpathlineto{\pgfqpoint{4.201115in}{2.298702in}}%
\pgfpathlineto{\pgfqpoint{4.201140in}{2.299154in}}%
\pgfpathlineto{\pgfqpoint{4.201169in}{2.300294in}}%
\pgfpathlineto{\pgfqpoint{4.202016in}{2.294246in}}%
\pgfpathlineto{\pgfqpoint{4.202040in}{2.294530in}}%
\pgfpathlineto{\pgfqpoint{4.202410in}{2.292767in}}%
\pgfpathlineto{\pgfqpoint{4.203092in}{2.296096in}}%
\pgfpathlineto{\pgfqpoint{4.203136in}{2.295691in}}%
\pgfpathlineto{\pgfqpoint{4.204382in}{2.303460in}}%
\pgfpathlineto{\pgfqpoint{4.204504in}{2.303142in}}%
\pgfpathlineto{\pgfqpoint{4.205657in}{2.300739in}}%
\pgfpathlineto{\pgfqpoint{4.204596in}{2.304268in}}%
\pgfpathlineto{\pgfqpoint{4.205662in}{2.300759in}}%
\pgfpathlineto{\pgfqpoint{4.206841in}{2.303210in}}%
\pgfpathlineto{\pgfqpoint{4.205872in}{2.300221in}}%
\pgfpathlineto{\pgfqpoint{4.206850in}{2.303006in}}%
\pgfpathlineto{\pgfqpoint{4.207361in}{2.299643in}}%
\pgfpathlineto{\pgfqpoint{4.207814in}{2.303328in}}%
\pgfpathlineto{\pgfqpoint{4.207994in}{2.302082in}}%
\pgfpathlineto{\pgfqpoint{4.207999in}{2.302365in}}%
\pgfpathlineto{\pgfqpoint{4.208973in}{2.297423in}}%
\pgfpathlineto{\pgfqpoint{4.209046in}{2.298681in}}%
\pgfpathlineto{\pgfqpoint{4.210214in}{2.301448in}}%
\pgfpathlineto{\pgfqpoint{4.209421in}{2.298200in}}%
\pgfpathlineto{\pgfqpoint{4.210297in}{2.300775in}}%
\pgfpathlineto{\pgfqpoint{4.210526in}{2.300084in}}%
\pgfpathlineto{\pgfqpoint{4.211144in}{2.304153in}}%
\pgfpathlineto{\pgfqpoint{4.211271in}{2.303856in}}%
\pgfpathlineto{\pgfqpoint{4.212064in}{2.304928in}}%
\pgfpathlineto{\pgfqpoint{4.211505in}{2.301699in}}%
\pgfpathlineto{\pgfqpoint{4.212303in}{2.302860in}}%
\pgfpathlineto{\pgfqpoint{4.213447in}{2.300138in}}%
\pgfpathlineto{\pgfqpoint{4.212985in}{2.303285in}}%
\pgfpathlineto{\pgfqpoint{4.213452in}{2.300437in}}%
\pgfpathlineto{\pgfqpoint{4.213802in}{2.302638in}}%
\pgfpathlineto{\pgfqpoint{4.214440in}{2.299251in}}%
\pgfpathlineto{\pgfqpoint{4.214445in}{2.299474in}}%
\pgfpathlineto{\pgfqpoint{4.214727in}{2.297210in}}%
\pgfpathlineto{\pgfqpoint{4.215487in}{2.301247in}}%
\pgfpathlineto{\pgfqpoint{4.215550in}{2.299883in}}%
\pgfpathlineto{\pgfqpoint{4.215979in}{2.302304in}}%
\pgfpathlineto{\pgfqpoint{4.216422in}{2.297806in}}%
\pgfpathlineto{\pgfqpoint{4.216427in}{2.297871in}}%
\pgfpathlineto{\pgfqpoint{4.216714in}{2.297102in}}%
\pgfpathlineto{\pgfqpoint{4.217186in}{2.300183in}}%
\pgfpathlineto{\pgfqpoint{4.217507in}{2.299396in}}%
\pgfpathlineto{\pgfqpoint{4.217600in}{2.298489in}}%
\pgfpathlineto{\pgfqpoint{4.217882in}{2.301192in}}%
\pgfpathlineto{\pgfqpoint{4.218247in}{2.301911in}}%
\pgfpathlineto{\pgfqpoint{4.218676in}{2.298875in}}%
\pgfpathlineto{\pgfqpoint{4.218895in}{2.300118in}}%
\pgfpathlineto{\pgfqpoint{4.219318in}{2.298136in}}%
\pgfpathlineto{\pgfqpoint{4.219557in}{2.301702in}}%
\pgfpathlineto{\pgfqpoint{4.219562in}{2.301701in}}%
\pgfpathlineto{\pgfqpoint{4.220750in}{2.304709in}}%
\pgfpathlineto{\pgfqpoint{4.219684in}{2.301122in}}%
\pgfpathlineto{\pgfqpoint{4.220774in}{2.304597in}}%
\pgfpathlineto{\pgfqpoint{4.221655in}{2.299854in}}%
\pgfpathlineto{\pgfqpoint{4.221962in}{2.300835in}}%
\pgfpathlineto{\pgfqpoint{4.222021in}{2.299818in}}%
\pgfpathlineto{\pgfqpoint{4.222025in}{2.299664in}}%
\pgfpathlineto{\pgfqpoint{4.222975in}{2.303054in}}%
\pgfpathlineto{\pgfqpoint{4.223043in}{2.302551in}}%
\pgfpathlineto{\pgfqpoint{4.223301in}{2.304412in}}%
\pgfpathlineto{\pgfqpoint{4.223053in}{2.302490in}}%
\pgfpathlineto{\pgfqpoint{4.224202in}{2.302956in}}%
\pgfpathlineto{\pgfqpoint{4.224888in}{2.300800in}}%
\pgfpathlineto{\pgfqpoint{4.224596in}{2.303913in}}%
\pgfpathlineto{\pgfqpoint{4.225326in}{2.301852in}}%
\pgfpathlineto{\pgfqpoint{4.225506in}{2.301379in}}%
\pgfpathlineto{\pgfqpoint{4.226514in}{2.304654in}}%
\pgfpathlineto{\pgfqpoint{4.227799in}{2.299033in}}%
\pgfpathlineto{\pgfqpoint{4.226738in}{2.305064in}}%
\pgfpathlineto{\pgfqpoint{4.227809in}{2.299478in}}%
\pgfpathlineto{\pgfqpoint{4.228982in}{2.301304in}}%
\pgfpathlineto{\pgfqpoint{4.228413in}{2.297017in}}%
\pgfpathlineto{\pgfqpoint{4.228992in}{2.301289in}}%
\pgfpathlineto{\pgfqpoint{4.229294in}{2.300177in}}%
\pgfpathlineto{\pgfqpoint{4.229976in}{2.303540in}}%
\pgfpathlineto{\pgfqpoint{4.230049in}{2.302906in}}%
\pgfpathlineto{\pgfqpoint{4.231276in}{2.306086in}}%
\pgfpathlineto{\pgfqpoint{4.230584in}{2.301003in}}%
\pgfpathlineto{\pgfqpoint{4.231329in}{2.304476in}}%
\pgfpathlineto{\pgfqpoint{4.232181in}{2.301465in}}%
\pgfpathlineto{\pgfqpoint{4.231723in}{2.305018in}}%
\pgfpathlineto{\pgfqpoint{4.232444in}{2.303713in}}%
\pgfpathlineto{\pgfqpoint{4.232887in}{2.304707in}}%
\pgfpathlineto{\pgfqpoint{4.233223in}{2.300884in}}%
\pgfpathlineto{\pgfqpoint{4.233544in}{2.303264in}}%
\pgfpathlineto{\pgfqpoint{4.234109in}{2.299479in}}%
\pgfpathlineto{\pgfqpoint{4.234766in}{2.300806in}}%
\pgfpathlineto{\pgfqpoint{4.236013in}{2.305596in}}%
\pgfpathlineto{\pgfqpoint{4.235068in}{2.299883in}}%
\pgfpathlineto{\pgfqpoint{4.236061in}{2.304777in}}%
\pgfpathlineto{\pgfqpoint{4.236149in}{2.304738in}}%
\pgfpathlineto{\pgfqpoint{4.236266in}{2.306211in}}%
\pgfpathlineto{\pgfqpoint{4.236285in}{2.307196in}}%
\pgfpathlineto{\pgfqpoint{4.237230in}{2.302670in}}%
\pgfpathlineto{\pgfqpoint{4.237361in}{2.304941in}}%
\pgfpathlineto{\pgfqpoint{4.238505in}{2.300935in}}%
\pgfpathlineto{\pgfqpoint{4.238525in}{2.301289in}}%
\pgfpathlineto{\pgfqpoint{4.238822in}{2.299380in}}%
\pgfpathlineto{\pgfqpoint{4.239659in}{2.302841in}}%
\pgfpathlineto{\pgfqpoint{4.240647in}{2.300695in}}%
\pgfpathlineto{\pgfqpoint{4.239883in}{2.304024in}}%
\pgfpathlineto{\pgfqpoint{4.240774in}{2.302251in}}%
\pgfpathlineto{\pgfqpoint{4.241095in}{2.305449in}}%
\pgfpathlineto{\pgfqpoint{4.241616in}{2.301664in}}%
\pgfpathlineto{\pgfqpoint{4.241908in}{2.303611in}}%
\pgfpathlineto{\pgfqpoint{4.242975in}{2.300587in}}%
\pgfpathlineto{\pgfqpoint{4.242532in}{2.304045in}}%
\pgfpathlineto{\pgfqpoint{4.243052in}{2.301638in}}%
\pgfpathlineto{\pgfqpoint{4.243087in}{2.302388in}}%
\pgfpathlineto{\pgfqpoint{4.243325in}{2.299170in}}%
\pgfpathlineto{\pgfqpoint{4.244104in}{2.301215in}}%
\pgfpathlineto{\pgfqpoint{4.245428in}{2.295332in}}%
\pgfpathlineto{\pgfqpoint{4.244221in}{2.302115in}}%
\pgfpathlineto{\pgfqpoint{4.245492in}{2.296104in}}%
\pgfpathlineto{\pgfqpoint{4.245949in}{2.295222in}}%
\pgfpathlineto{\pgfqpoint{4.245686in}{2.297013in}}%
\pgfpathlineto{\pgfqpoint{4.246003in}{2.296850in}}%
\pgfpathlineto{\pgfqpoint{4.246159in}{2.298171in}}%
\pgfpathlineto{\pgfqpoint{4.246796in}{2.295585in}}%
\pgfpathlineto{\pgfqpoint{4.247118in}{2.297139in}}%
\pgfpathlineto{\pgfqpoint{4.248257in}{2.301124in}}%
\pgfpathlineto{\pgfqpoint{4.247327in}{2.294404in}}%
\pgfpathlineto{\pgfqpoint{4.248354in}{2.299998in}}%
\pgfpathlineto{\pgfqpoint{4.249182in}{2.299269in}}%
\pgfpathlineto{\pgfqpoint{4.248768in}{2.301597in}}%
\pgfpathlineto{\pgfqpoint{4.249318in}{2.300625in}}%
\pgfpathlineto{\pgfqpoint{4.250253in}{2.307776in}}%
\pgfpathlineto{\pgfqpoint{4.249338in}{2.300185in}}%
\pgfpathlineto{\pgfqpoint{4.250603in}{2.305243in}}%
\pgfpathlineto{\pgfqpoint{4.251426in}{2.302432in}}%
\pgfpathlineto{\pgfqpoint{4.250677in}{2.306354in}}%
\pgfpathlineto{\pgfqpoint{4.251816in}{2.304388in}}%
\pgfpathlineto{\pgfqpoint{4.252629in}{2.305321in}}%
\pgfpathlineto{\pgfqpoint{4.252006in}{2.302846in}}%
\pgfpathlineto{\pgfqpoint{4.252935in}{2.304838in}}%
\pgfpathlineto{\pgfqpoint{4.253033in}{2.305175in}}%
\pgfpathlineto{\pgfqpoint{4.253491in}{2.299418in}}%
\pgfpathlineto{\pgfqpoint{4.253836in}{2.300773in}}%
\pgfpathlineto{\pgfqpoint{4.254878in}{2.299342in}}%
\pgfpathlineto{\pgfqpoint{4.254323in}{2.301403in}}%
\pgfpathlineto{\pgfqpoint{4.254971in}{2.299663in}}%
\pgfpathlineto{\pgfqpoint{4.254995in}{2.299808in}}%
\pgfpathlineto{\pgfqpoint{4.255268in}{2.297745in}}%
\pgfpathlineto{\pgfqpoint{4.256056in}{2.299030in}}%
\pgfpathlineto{\pgfqpoint{4.256144in}{2.298100in}}%
\pgfpathlineto{\pgfqpoint{4.256796in}{2.301192in}}%
\pgfpathlineto{\pgfqpoint{4.257132in}{2.300187in}}%
\pgfpathlineto{\pgfqpoint{4.257799in}{2.301279in}}%
\pgfpathlineto{\pgfqpoint{4.257673in}{2.299554in}}%
\pgfpathlineto{\pgfqpoint{4.258213in}{2.299998in}}%
\pgfpathlineto{\pgfqpoint{4.258330in}{2.298348in}}%
\pgfpathlineto{\pgfqpoint{4.259172in}{2.303976in}}%
\pgfpathlineto{\pgfqpoint{4.259284in}{2.303070in}}%
\pgfpathlineto{\pgfqpoint{4.259323in}{2.303481in}}%
\pgfpathlineto{\pgfqpoint{4.260014in}{2.299279in}}%
\pgfpathlineto{\pgfqpoint{4.260375in}{2.302994in}}%
\pgfpathlineto{\pgfqpoint{4.260599in}{2.301802in}}%
\pgfpathlineto{\pgfqpoint{4.260803in}{2.303361in}}%
\pgfpathlineto{\pgfqpoint{4.261480in}{2.303016in}}%
\pgfpathlineto{\pgfqpoint{4.262249in}{2.304869in}}%
\pgfpathlineto{\pgfqpoint{4.261996in}{2.301549in}}%
\pgfpathlineto{\pgfqpoint{4.262624in}{2.304173in}}%
\pgfpathlineto{\pgfqpoint{4.263646in}{2.302137in}}%
\pgfpathlineto{\pgfqpoint{4.263451in}{2.304797in}}%
\pgfpathlineto{\pgfqpoint{4.263744in}{2.303158in}}%
\pgfpathlineto{\pgfqpoint{4.264712in}{2.304096in}}%
\pgfpathlineto{\pgfqpoint{4.264396in}{2.301029in}}%
\pgfpathlineto{\pgfqpoint{4.264820in}{2.302524in}}%
\pgfpathlineto{\pgfqpoint{4.266261in}{2.297001in}}%
\pgfpathlineto{\pgfqpoint{4.264966in}{2.304053in}}%
\pgfpathlineto{\pgfqpoint{4.266436in}{2.298211in}}%
\pgfpathlineto{\pgfqpoint{4.267512in}{2.303458in}}%
\pgfpathlineto{\pgfqpoint{4.267668in}{2.303203in}}%
\pgfpathlineto{\pgfqpoint{4.267741in}{2.303591in}}%
\pgfpathlineto{\pgfqpoint{4.268227in}{2.300337in}}%
\pgfpathlineto{\pgfqpoint{4.268744in}{2.302301in}}%
\pgfpathlineto{\pgfqpoint{4.270082in}{2.304378in}}%
\pgfpathlineto{\pgfqpoint{4.268753in}{2.302015in}}%
\pgfpathlineto{\pgfqpoint{4.270131in}{2.303697in}}%
\pgfpathlineto{\pgfqpoint{4.270379in}{2.301804in}}%
\pgfpathlineto{\pgfqpoint{4.270214in}{2.304179in}}%
\pgfpathlineto{\pgfqpoint{4.271256in}{2.303435in}}%
\pgfpathlineto{\pgfqpoint{4.272327in}{2.305596in}}%
\pgfpathlineto{\pgfqpoint{4.271893in}{2.302457in}}%
\pgfpathlineto{\pgfqpoint{4.272444in}{2.305348in}}%
\pgfpathlineto{\pgfqpoint{4.273023in}{2.301232in}}%
\pgfpathlineto{\pgfqpoint{4.272483in}{2.305980in}}%
\pgfpathlineto{\pgfqpoint{4.273612in}{2.303239in}}%
\pgfpathlineto{\pgfqpoint{4.274790in}{2.300563in}}%
\pgfpathlineto{\pgfqpoint{4.273812in}{2.303406in}}%
\pgfpathlineto{\pgfqpoint{4.274853in}{2.301188in}}%
\pgfpathlineto{\pgfqpoint{4.275876in}{2.304018in}}%
\pgfpathlineto{\pgfqpoint{4.275097in}{2.300050in}}%
\pgfpathlineto{\pgfqpoint{4.275978in}{2.302490in}}%
\pgfpathlineto{\pgfqpoint{4.276538in}{2.302095in}}%
\pgfpathlineto{\pgfqpoint{4.276884in}{2.304331in}}%
\pgfpathlineto{\pgfqpoint{4.277074in}{2.302829in}}%
\pgfpathlineto{\pgfqpoint{4.277994in}{2.305575in}}%
\pgfpathlineto{\pgfqpoint{4.277473in}{2.301351in}}%
\pgfpathlineto{\pgfqpoint{4.278213in}{2.304947in}}%
\pgfpathlineto{\pgfqpoint{4.279269in}{2.301200in}}%
\pgfpathlineto{\pgfqpoint{4.278271in}{2.305497in}}%
\pgfpathlineto{\pgfqpoint{4.279522in}{2.303260in}}%
\pgfpathlineto{\pgfqpoint{4.279819in}{2.303996in}}%
\pgfpathlineto{\pgfqpoint{4.279746in}{2.302102in}}%
\pgfpathlineto{\pgfqpoint{4.280608in}{2.302309in}}%
\pgfpathlineto{\pgfqpoint{4.280783in}{2.301387in}}%
\pgfpathlineto{\pgfqpoint{4.281509in}{2.303770in}}%
\pgfpathlineto{\pgfqpoint{4.281694in}{2.302876in}}%
\pgfpathlineto{\pgfqpoint{4.281893in}{2.305847in}}%
\pgfpathlineto{\pgfqpoint{4.282336in}{2.302303in}}%
\pgfpathlineto{\pgfqpoint{4.282804in}{2.302867in}}%
\pgfpathlineto{\pgfqpoint{4.283218in}{2.301108in}}%
\pgfpathlineto{\pgfqpoint{4.283656in}{2.304866in}}%
\pgfpathlineto{\pgfqpoint{4.283889in}{2.303550in}}%
\pgfpathlineto{\pgfqpoint{4.283958in}{2.304027in}}%
\pgfpathlineto{\pgfqpoint{4.284420in}{2.301979in}}%
\pgfpathlineto{\pgfqpoint{4.284980in}{2.303005in}}%
\pgfpathlineto{\pgfqpoint{4.285968in}{2.299493in}}%
\pgfpathlineto{\pgfqpoint{4.285180in}{2.303341in}}%
\pgfpathlineto{\pgfqpoint{4.286187in}{2.300428in}}%
\pgfpathlineto{\pgfqpoint{4.286484in}{2.303460in}}%
\pgfpathlineto{\pgfqpoint{4.286918in}{2.299655in}}%
\pgfpathlineto{\pgfqpoint{4.287312in}{2.301420in}}%
\pgfpathlineto{\pgfqpoint{4.287448in}{2.299957in}}%
\pgfpathlineto{\pgfqpoint{4.287882in}{2.303256in}}%
\pgfpathlineto{\pgfqpoint{4.288427in}{2.301236in}}%
\pgfpathlineto{\pgfqpoint{4.288923in}{2.305387in}}%
\pgfpathlineto{\pgfqpoint{4.289478in}{2.300587in}}%
\pgfpathlineto{\pgfqpoint{4.289532in}{2.301159in}}%
\pgfpathlineto{\pgfqpoint{4.289537in}{2.301075in}}%
\pgfpathlineto{\pgfqpoint{4.290175in}{2.303931in}}%
\pgfpathlineto{\pgfqpoint{4.290545in}{2.302584in}}%
\pgfpathlineto{\pgfqpoint{4.291426in}{2.305898in}}%
\pgfpathlineto{\pgfqpoint{4.291684in}{2.304126in}}%
\pgfpathlineto{\pgfqpoint{4.292624in}{2.302055in}}%
\pgfpathlineto{\pgfqpoint{4.291825in}{2.305126in}}%
\pgfpathlineto{\pgfqpoint{4.292799in}{2.303827in}}%
\pgfpathlineto{\pgfqpoint{4.292847in}{2.304246in}}%
\pgfpathlineto{\pgfqpoint{4.293612in}{2.301544in}}%
\pgfpathlineto{\pgfqpoint{4.293884in}{2.303142in}}%
\pgfpathlineto{\pgfqpoint{4.294941in}{2.299490in}}%
\pgfpathlineto{\pgfqpoint{4.294298in}{2.303646in}}%
\pgfpathlineto{\pgfqpoint{4.295033in}{2.300255in}}%
\pgfpathlineto{\pgfqpoint{4.296007in}{2.302747in}}%
\pgfpathlineto{\pgfqpoint{4.295092in}{2.299479in}}%
\pgfpathlineto{\pgfqpoint{4.296275in}{2.301082in}}%
\pgfpathlineto{\pgfqpoint{4.296528in}{2.298978in}}%
\pgfpathlineto{\pgfqpoint{4.297253in}{2.303780in}}%
\pgfpathlineto{\pgfqpoint{4.297336in}{2.302330in}}%
\pgfpathlineto{\pgfqpoint{4.297959in}{2.303845in}}%
\pgfpathlineto{\pgfqpoint{4.297721in}{2.300763in}}%
\pgfpathlineto{\pgfqpoint{4.298446in}{2.302570in}}%
\pgfpathlineto{\pgfqpoint{4.299157in}{2.300410in}}%
\pgfpathlineto{\pgfqpoint{4.299318in}{2.303069in}}%
\pgfpathlineto{\pgfqpoint{4.299551in}{2.302799in}}%
\pgfpathlineto{\pgfqpoint{4.299663in}{2.302171in}}%
\pgfpathlineto{\pgfqpoint{4.300248in}{2.305274in}}%
\pgfpathlineto{\pgfqpoint{4.300598in}{2.303576in}}%
\pgfpathlineto{\pgfqpoint{4.300603in}{2.303834in}}%
\pgfpathlineto{\pgfqpoint{4.301538in}{2.298411in}}%
\pgfpathlineto{\pgfqpoint{4.301645in}{2.300152in}}%
\pgfpathlineto{\pgfqpoint{4.302779in}{2.304246in}}%
\pgfpathlineto{\pgfqpoint{4.301669in}{2.300145in}}%
\pgfpathlineto{\pgfqpoint{4.303247in}{2.303437in}}%
\pgfpathlineto{\pgfqpoint{4.304103in}{2.300207in}}%
\pgfpathlineto{\pgfqpoint{4.303923in}{2.304299in}}%
\pgfpathlineto{\pgfqpoint{4.304371in}{2.302224in}}%
\pgfpathlineto{\pgfqpoint{4.304454in}{2.301552in}}%
\pgfpathlineto{\pgfqpoint{4.305325in}{2.303775in}}%
\pgfpathlineto{\pgfqpoint{4.305345in}{2.303746in}}%
\pgfpathlineto{\pgfqpoint{4.305360in}{2.304119in}}%
\pgfpathlineto{\pgfqpoint{4.306309in}{2.300999in}}%
\pgfpathlineto{\pgfqpoint{4.306416in}{2.302100in}}%
\pgfpathlineto{\pgfqpoint{4.306620in}{2.300552in}}%
\pgfpathlineto{\pgfqpoint{4.306805in}{2.304155in}}%
\pgfpathlineto{\pgfqpoint{4.307487in}{2.303954in}}%
\pgfpathlineto{\pgfqpoint{4.307735in}{2.304462in}}%
\pgfpathlineto{\pgfqpoint{4.308485in}{2.300413in}}%
\pgfpathlineto{\pgfqpoint{4.308500in}{2.300558in}}%
\pgfpathlineto{\pgfqpoint{4.308923in}{2.299224in}}%
\pgfpathlineto{\pgfqpoint{4.309318in}{2.304150in}}%
\pgfpathlineto{\pgfqpoint{4.309610in}{2.300439in}}%
\pgfpathlineto{\pgfqpoint{4.310408in}{2.294670in}}%
\pgfpathlineto{\pgfqpoint{4.309829in}{2.300744in}}%
\pgfpathlineto{\pgfqpoint{4.310734in}{2.300037in}}%
\pgfpathlineto{\pgfqpoint{4.310992in}{2.301647in}}%
\pgfpathlineto{\pgfqpoint{4.311562in}{2.296054in}}%
\pgfpathlineto{\pgfqpoint{4.311796in}{2.297804in}}%
\pgfpathlineto{\pgfqpoint{4.312531in}{2.295680in}}%
\pgfpathlineto{\pgfqpoint{4.312799in}{2.298983in}}%
\pgfpathlineto{\pgfqpoint{4.312823in}{2.298679in}}%
\pgfpathlineto{\pgfqpoint{4.313393in}{2.303405in}}%
\pgfpathlineto{\pgfqpoint{4.312862in}{2.298123in}}%
\pgfpathlineto{\pgfqpoint{4.313991in}{2.300404in}}%
\pgfpathlineto{\pgfqpoint{4.314400in}{2.300203in}}%
\pgfpathlineto{\pgfqpoint{4.314955in}{2.304420in}}%
\pgfpathlineto{\pgfqpoint{4.315014in}{2.303585in}}%
\pgfpathlineto{\pgfqpoint{4.315033in}{2.304228in}}%
\pgfpathlineto{\pgfqpoint{4.315403in}{2.301656in}}%
\pgfpathlineto{\pgfqpoint{4.316022in}{2.300773in}}%
\pgfpathlineto{\pgfqpoint{4.316236in}{2.303764in}}%
\pgfpathlineto{\pgfqpoint{4.316484in}{2.302468in}}%
\pgfpathlineto{\pgfqpoint{4.316942in}{2.304331in}}%
\pgfpathlineto{\pgfqpoint{4.317516in}{2.300730in}}%
\pgfpathlineto{\pgfqpoint{4.317575in}{2.300848in}}%
\pgfpathlineto{\pgfqpoint{4.318222in}{2.296555in}}%
\pgfpathlineto{\pgfqpoint{4.318753in}{2.297497in}}%
\pgfpathlineto{\pgfqpoint{4.318923in}{2.295228in}}%
\pgfpathlineto{\pgfqpoint{4.319444in}{2.300567in}}%
\pgfpathlineto{\pgfqpoint{4.319751in}{2.299385in}}%
\pgfpathlineto{\pgfqpoint{4.319897in}{2.301031in}}%
\pgfpathlineto{\pgfqpoint{4.320286in}{2.299256in}}%
\pgfpathlineto{\pgfqpoint{4.320866in}{2.299572in}}%
\pgfpathlineto{\pgfqpoint{4.321007in}{2.301340in}}%
\pgfpathlineto{\pgfqpoint{4.321518in}{2.296158in}}%
\pgfpathlineto{\pgfqpoint{4.321577in}{2.296690in}}%
\pgfpathlineto{\pgfqpoint{4.322122in}{2.295642in}}%
\pgfpathlineto{\pgfqpoint{4.322623in}{2.299232in}}%
\pgfpathlineto{\pgfqpoint{4.322662in}{2.298608in}}%
\pgfpathlineto{\pgfqpoint{4.324094in}{2.304149in}}%
\pgfpathlineto{\pgfqpoint{4.322940in}{2.296810in}}%
\pgfpathlineto{\pgfqpoint{4.324157in}{2.303407in}}%
\pgfpathlineto{\pgfqpoint{4.324522in}{2.301625in}}%
\pgfpathlineto{\pgfqpoint{4.325150in}{2.306629in}}%
\pgfpathlineto{\pgfqpoint{4.325189in}{2.306273in}}%
\pgfpathlineto{\pgfqpoint{4.325194in}{2.306322in}}%
\pgfpathlineto{\pgfqpoint{4.325793in}{2.301315in}}%
\pgfpathlineto{\pgfqpoint{4.325836in}{2.301627in}}%
\pgfpathlineto{\pgfqpoint{4.325890in}{2.300796in}}%
\pgfpathlineto{\pgfqpoint{4.326538in}{2.304112in}}%
\pgfpathlineto{\pgfqpoint{4.326912in}{2.302808in}}%
\pgfpathlineto{\pgfqpoint{4.327190in}{2.301895in}}%
\pgfpathlineto{\pgfqpoint{4.327686in}{2.304981in}}%
\pgfpathlineto{\pgfqpoint{4.327964in}{2.302986in}}%
\pgfpathlineto{\pgfqpoint{4.328845in}{2.303952in}}%
\pgfpathlineto{\pgfqpoint{4.328081in}{2.300631in}}%
\pgfpathlineto{\pgfqpoint{4.329069in}{2.302546in}}%
\pgfpathlineto{\pgfqpoint{4.329322in}{2.303739in}}%
\pgfpathlineto{\pgfqpoint{4.329858in}{2.300748in}}%
\pgfpathlineto{\pgfqpoint{4.330218in}{2.303490in}}%
\pgfpathlineto{\pgfqpoint{4.330778in}{2.304012in}}%
\pgfpathlineto{\pgfqpoint{4.330861in}{2.301850in}}%
\pgfpathlineto{\pgfqpoint{4.331270in}{2.302509in}}%
\pgfpathlineto{\pgfqpoint{4.331484in}{2.304282in}}%
\pgfpathlineto{\pgfqpoint{4.332049in}{2.300783in}}%
\pgfpathlineto{\pgfqpoint{4.332054in}{2.300795in}}%
\pgfpathlineto{\pgfqpoint{4.332107in}{2.300089in}}%
\pgfpathlineto{\pgfqpoint{4.332803in}{2.303831in}}%
\pgfpathlineto{\pgfqpoint{4.333120in}{2.302127in}}%
\pgfpathlineto{\pgfqpoint{4.334196in}{2.305557in}}%
\pgfpathlineto{\pgfqpoint{4.333412in}{2.300894in}}%
\pgfpathlineto{\pgfqpoint{4.334298in}{2.305411in}}%
\pgfpathlineto{\pgfqpoint{4.335447in}{2.300743in}}%
\pgfpathlineto{\pgfqpoint{4.335471in}{2.301496in}}%
\pgfpathlineto{\pgfqpoint{4.336416in}{2.304099in}}%
\pgfpathlineto{\pgfqpoint{4.336649in}{2.303919in}}%
\pgfpathlineto{\pgfqpoint{4.337350in}{2.300531in}}%
\pgfpathlineto{\pgfqpoint{4.337667in}{2.304370in}}%
\pgfpathlineto{\pgfqpoint{4.337784in}{2.302628in}}%
\pgfpathlineto{\pgfqpoint{4.338650in}{2.303601in}}%
\pgfpathlineto{\pgfqpoint{4.338285in}{2.301560in}}%
\pgfpathlineto{\pgfqpoint{4.338826in}{2.302235in}}%
\pgfpathlineto{\pgfqpoint{4.340033in}{2.296434in}}%
\pgfpathlineto{\pgfqpoint{4.339210in}{2.302756in}}%
\pgfpathlineto{\pgfqpoint{4.340057in}{2.297543in}}%
\pgfpathlineto{\pgfqpoint{4.340291in}{2.296157in}}%
\pgfpathlineto{\pgfqpoint{4.340476in}{2.298450in}}%
\pgfpathlineto{\pgfqpoint{4.340783in}{2.298467in}}%
\pgfpathlineto{\pgfqpoint{4.341318in}{2.299702in}}%
\pgfpathlineto{\pgfqpoint{4.341854in}{2.295900in}}%
\pgfpathlineto{\pgfqpoint{4.342584in}{2.298031in}}%
\pgfpathlineto{\pgfqpoint{4.342881in}{2.294964in}}%
\pgfpathlineto{\pgfqpoint{4.343076in}{2.294060in}}%
\pgfpathlineto{\pgfqpoint{4.343543in}{2.299117in}}%
\pgfpathlineto{\pgfqpoint{4.343938in}{2.297452in}}%
\pgfpathlineto{\pgfqpoint{4.344911in}{2.299696in}}%
\pgfpathlineto{\pgfqpoint{4.344093in}{2.295880in}}%
\pgfpathlineto{\pgfqpoint{4.345062in}{2.297968in}}%
\pgfpathlineto{\pgfqpoint{4.345992in}{2.297652in}}%
\pgfpathlineto{\pgfqpoint{4.345301in}{2.301442in}}%
\pgfpathlineto{\pgfqpoint{4.346085in}{2.299106in}}%
\pgfpathlineto{\pgfqpoint{4.346805in}{2.300271in}}%
\pgfpathlineto{\pgfqpoint{4.346961in}{2.298080in}}%
\pgfpathlineto{\pgfqpoint{4.347204in}{2.299633in}}%
\pgfpathlineto{\pgfqpoint{4.348149in}{2.300995in}}%
\pgfpathlineto{\pgfqpoint{4.347745in}{2.298728in}}%
\pgfpathlineto{\pgfqpoint{4.348236in}{2.299231in}}%
\pgfpathlineto{\pgfqpoint{4.348241in}{2.299091in}}%
\pgfpathlineto{\pgfqpoint{4.348572in}{2.301982in}}%
\pgfpathlineto{\pgfqpoint{4.349264in}{2.301368in}}%
\pgfpathlineto{\pgfqpoint{4.350418in}{2.304849in}}%
\pgfpathlineto{\pgfqpoint{4.349376in}{2.300328in}}%
\pgfpathlineto{\pgfqpoint{4.350437in}{2.304783in}}%
\pgfpathlineto{\pgfqpoint{4.351430in}{2.299294in}}%
\pgfpathlineto{\pgfqpoint{4.350578in}{2.305632in}}%
\pgfpathlineto{\pgfqpoint{4.351576in}{2.303088in}}%
\pgfpathlineto{\pgfqpoint{4.351936in}{2.304119in}}%
\pgfpathlineto{\pgfqpoint{4.352238in}{2.300282in}}%
\pgfpathlineto{\pgfqpoint{4.352365in}{2.300864in}}%
\pgfpathlineto{\pgfqpoint{4.352925in}{2.297941in}}%
\pgfpathlineto{\pgfqpoint{4.353509in}{2.298234in}}%
\pgfpathlineto{\pgfqpoint{4.353976in}{2.296894in}}%
\pgfpathlineto{\pgfqpoint{4.354307in}{2.300149in}}%
\pgfpathlineto{\pgfqpoint{4.354371in}{2.299951in}}%
\pgfpathlineto{\pgfqpoint{4.354522in}{2.301328in}}%
\pgfpathlineto{\pgfqpoint{4.354950in}{2.298523in}}%
\pgfpathlineto{\pgfqpoint{4.355452in}{2.299066in}}%
\pgfpathlineto{\pgfqpoint{4.355865in}{2.295636in}}%
\pgfpathlineto{\pgfqpoint{4.356469in}{2.301180in}}%
\pgfpathlineto{\pgfqpoint{4.356532in}{2.300393in}}%
\pgfpathlineto{\pgfqpoint{4.357316in}{2.298584in}}%
\pgfpathlineto{\pgfqpoint{4.357058in}{2.301080in}}%
\pgfpathlineto{\pgfqpoint{4.357696in}{2.299478in}}%
\pgfpathlineto{\pgfqpoint{4.357735in}{2.300906in}}%
\pgfpathlineto{\pgfqpoint{4.358451in}{2.295779in}}%
\pgfpathlineto{\pgfqpoint{4.358757in}{2.297246in}}%
\pgfpathlineto{\pgfqpoint{4.358791in}{2.296707in}}%
\pgfpathlineto{\pgfqpoint{4.359045in}{2.298805in}}%
\pgfpathlineto{\pgfqpoint{4.359862in}{2.297239in}}%
\pgfpathlineto{\pgfqpoint{4.360500in}{2.300994in}}%
\pgfpathlineto{\pgfqpoint{4.360997in}{2.299386in}}%
\pgfpathlineto{\pgfqpoint{4.361469in}{2.295777in}}%
\pgfpathlineto{\pgfqpoint{4.361099in}{2.300131in}}%
\pgfpathlineto{\pgfqpoint{4.362121in}{2.298191in}}%
\pgfpathlineto{\pgfqpoint{4.362988in}{2.301678in}}%
\pgfpathlineto{\pgfqpoint{4.362190in}{2.298025in}}%
\pgfpathlineto{\pgfqpoint{4.363373in}{2.301054in}}%
\pgfpathlineto{\pgfqpoint{4.363796in}{2.298529in}}%
\pgfpathlineto{\pgfqpoint{4.364337in}{2.302387in}}%
\pgfpathlineto{\pgfqpoint{4.364512in}{2.300217in}}%
\pgfpathlineto{\pgfqpoint{4.364619in}{2.301091in}}%
\pgfpathlineto{\pgfqpoint{4.365135in}{2.296895in}}%
\pgfpathlineto{\pgfqpoint{4.365593in}{2.299221in}}%
\pgfpathlineto{\pgfqpoint{4.365700in}{2.297540in}}%
\pgfpathlineto{\pgfqpoint{4.366552in}{2.303287in}}%
\pgfpathlineto{\pgfqpoint{4.367087in}{2.304751in}}%
\pgfpathlineto{\pgfqpoint{4.366863in}{2.301915in}}%
\pgfpathlineto{\pgfqpoint{4.367662in}{2.303378in}}%
\pgfpathlineto{\pgfqpoint{4.368412in}{2.302203in}}%
\pgfpathlineto{\pgfqpoint{4.367783in}{2.304249in}}%
\pgfpathlineto{\pgfqpoint{4.368601in}{2.303746in}}%
\pgfpathlineto{\pgfqpoint{4.368704in}{2.304694in}}%
\pgfpathlineto{\pgfqpoint{4.369648in}{2.301987in}}%
\pgfpathlineto{\pgfqpoint{4.369682in}{2.302465in}}%
\pgfpathlineto{\pgfqpoint{4.370617in}{2.304751in}}%
\pgfpathlineto{\pgfqpoint{4.370198in}{2.301567in}}%
\pgfpathlineto{\pgfqpoint{4.370948in}{2.303256in}}%
\pgfpathlineto{\pgfqpoint{4.370958in}{2.303012in}}%
\pgfpathlineto{\pgfqpoint{4.371683in}{2.305724in}}%
\pgfpathlineto{\pgfqpoint{4.371990in}{2.305273in}}%
\pgfpathlineto{\pgfqpoint{4.372891in}{2.306322in}}%
\pgfpathlineto{\pgfqpoint{4.372233in}{2.303803in}}%
\pgfpathlineto{\pgfqpoint{4.373027in}{2.304887in}}%
\pgfpathlineto{\pgfqpoint{4.374122in}{2.300126in}}%
\pgfpathlineto{\pgfqpoint{4.374215in}{2.301606in}}%
\pgfpathlineto{\pgfqpoint{4.375042in}{2.305097in}}%
\pgfpathlineto{\pgfqpoint{4.374380in}{2.299796in}}%
\pgfpathlineto{\pgfqpoint{4.375359in}{2.304224in}}%
\pgfpathlineto{\pgfqpoint{4.376381in}{2.301659in}}%
\pgfpathlineto{\pgfqpoint{4.375763in}{2.305161in}}%
\pgfpathlineto{\pgfqpoint{4.376503in}{2.303452in}}%
\pgfpathlineto{\pgfqpoint{4.376727in}{2.304729in}}%
\pgfpathlineto{\pgfqpoint{4.377389in}{2.301774in}}%
\pgfpathlineto{\pgfqpoint{4.377608in}{2.303030in}}%
\pgfpathlineto{\pgfqpoint{4.378280in}{2.304719in}}%
\pgfpathlineto{\pgfqpoint{4.378630in}{2.302854in}}%
\pgfpathlineto{\pgfqpoint{4.378738in}{2.303594in}}%
\pgfpathlineto{\pgfqpoint{4.379658in}{2.302045in}}%
\pgfpathlineto{\pgfqpoint{4.378811in}{2.304151in}}%
\pgfpathlineto{\pgfqpoint{4.379906in}{2.302626in}}%
\pgfpathlineto{\pgfqpoint{4.380748in}{2.303658in}}%
\pgfpathlineto{\pgfqpoint{4.380442in}{2.301235in}}%
\pgfpathlineto{\pgfqpoint{4.381026in}{2.303160in}}%
\pgfpathlineto{\pgfqpoint{4.381040in}{2.303492in}}%
\pgfpathlineto{\pgfqpoint{4.381717in}{2.299901in}}%
\pgfpathlineto{\pgfqpoint{4.382087in}{2.301757in}}%
\pgfpathlineto{\pgfqpoint{4.382175in}{2.300233in}}%
\pgfpathlineto{\pgfqpoint{4.382774in}{2.303767in}}%
\pgfpathlineto{\pgfqpoint{4.382915in}{2.303053in}}%
\pgfpathlineto{\pgfqpoint{4.383742in}{2.305131in}}%
\pgfpathlineto{\pgfqpoint{4.383251in}{2.302167in}}%
\pgfpathlineto{\pgfqpoint{4.384025in}{2.303366in}}%
\pgfpathlineto{\pgfqpoint{4.384575in}{2.304899in}}%
\pgfpathlineto{\pgfqpoint{4.384147in}{2.302494in}}%
\pgfpathlineto{\pgfqpoint{4.385076in}{2.303065in}}%
\pgfpathlineto{\pgfqpoint{4.386191in}{2.300717in}}%
\pgfpathlineto{\pgfqpoint{4.385831in}{2.304459in}}%
\pgfpathlineto{\pgfqpoint{4.386230in}{2.301316in}}%
\pgfpathlineto{\pgfqpoint{4.387087in}{2.298273in}}%
\pgfpathlineto{\pgfqpoint{4.386269in}{2.301870in}}%
\pgfpathlineto{\pgfqpoint{4.387457in}{2.300079in}}%
\pgfpathlineto{\pgfqpoint{4.388241in}{2.303702in}}%
\pgfpathlineto{\pgfqpoint{4.388621in}{2.301354in}}%
\pgfpathlineto{\pgfqpoint{4.388713in}{2.300247in}}%
\pgfpathlineto{\pgfqpoint{4.388903in}{2.303261in}}%
\pgfpathlineto{\pgfqpoint{4.389658in}{2.302266in}}%
\pgfpathlineto{\pgfqpoint{4.390578in}{2.304274in}}%
\pgfpathlineto{\pgfqpoint{4.389959in}{2.300525in}}%
\pgfpathlineto{\pgfqpoint{4.390768in}{2.302289in}}%
\pgfpathlineto{\pgfqpoint{4.391975in}{2.305350in}}%
\pgfpathlineto{\pgfqpoint{4.391001in}{2.300837in}}%
\pgfpathlineto{\pgfqpoint{4.392131in}{2.303713in}}%
\pgfpathlineto{\pgfqpoint{4.393129in}{2.300999in}}%
\pgfpathlineto{\pgfqpoint{4.392725in}{2.304109in}}%
\pgfpathlineto{\pgfqpoint{4.393270in}{2.302277in}}%
\pgfpathlineto{\pgfqpoint{4.393416in}{2.302734in}}%
\pgfpathlineto{\pgfqpoint{4.394073in}{2.297908in}}%
\pgfpathlineto{\pgfqpoint{4.394088in}{2.298034in}}%
\pgfpathlineto{\pgfqpoint{4.394200in}{2.297038in}}%
\pgfpathlineto{\pgfqpoint{4.394911in}{2.300379in}}%
\pgfpathlineto{\pgfqpoint{4.395120in}{2.299757in}}%
\pgfpathlineto{\pgfqpoint{4.395276in}{2.302434in}}%
\pgfpathlineto{\pgfqpoint{4.395529in}{2.298865in}}%
\pgfpathlineto{\pgfqpoint{4.396279in}{2.301077in}}%
\pgfpathlineto{\pgfqpoint{4.396439in}{2.300456in}}%
\pgfpathlineto{\pgfqpoint{4.396698in}{2.302965in}}%
\pgfpathlineto{\pgfqpoint{4.396809in}{2.302829in}}%
\pgfpathlineto{\pgfqpoint{4.397476in}{2.304607in}}%
\pgfpathlineto{\pgfqpoint{4.397199in}{2.301609in}}%
\pgfpathlineto{\pgfqpoint{4.397934in}{2.303337in}}%
\pgfpathlineto{\pgfqpoint{4.398747in}{2.299282in}}%
\pgfpathlineto{\pgfqpoint{4.398207in}{2.304973in}}%
\pgfpathlineto{\pgfqpoint{4.399122in}{2.299829in}}%
\pgfpathlineto{\pgfqpoint{4.400067in}{2.304283in}}%
\pgfpathlineto{\pgfqpoint{4.399205in}{2.299381in}}%
\pgfpathlineto{\pgfqpoint{4.400271in}{2.302756in}}%
\pgfpathlineto{\pgfqpoint{4.400904in}{2.304032in}}%
\pgfpathlineto{\pgfqpoint{4.400475in}{2.300209in}}%
\pgfpathlineto{\pgfqpoint{4.401391in}{2.303172in}}%
\pgfpathlineto{\pgfqpoint{4.402462in}{2.300818in}}%
\pgfpathlineto{\pgfqpoint{4.401887in}{2.305655in}}%
\pgfpathlineto{\pgfqpoint{4.402535in}{2.301003in}}%
\pgfpathlineto{\pgfqpoint{4.402764in}{2.300748in}}%
\pgfpathlineto{\pgfqpoint{4.403708in}{2.304257in}}%
\pgfpathlineto{\pgfqpoint{4.403713in}{2.304277in}}%
\pgfpathlineto{\pgfqpoint{4.404141in}{2.301898in}}%
\pgfpathlineto{\pgfqpoint{4.404409in}{2.302598in}}%
\pgfpathlineto{\pgfqpoint{4.404672in}{2.302825in}}%
\pgfpathlineto{\pgfqpoint{4.405675in}{2.297947in}}%
\pgfpathlineto{\pgfqpoint{4.406770in}{2.303525in}}%
\pgfpathlineto{\pgfqpoint{4.406882in}{2.302944in}}%
\pgfpathlineto{\pgfqpoint{4.407092in}{2.303683in}}%
\pgfpathlineto{\pgfqpoint{4.407778in}{2.297963in}}%
\pgfpathlineto{\pgfqpoint{4.407919in}{2.300327in}}%
\pgfpathlineto{\pgfqpoint{4.408007in}{2.299219in}}%
\pgfpathlineto{\pgfqpoint{4.408937in}{2.302851in}}%
\pgfpathlineto{\pgfqpoint{4.408956in}{2.302700in}}%
\pgfpathlineto{\pgfqpoint{4.409404in}{2.303820in}}%
\pgfpathlineto{\pgfqpoint{4.409804in}{2.301087in}}%
\pgfpathlineto{\pgfqpoint{4.410066in}{2.302930in}}%
\pgfpathlineto{\pgfqpoint{4.410315in}{2.301902in}}%
\pgfpathlineto{\pgfqpoint{4.410651in}{2.305659in}}%
\pgfpathlineto{\pgfqpoint{4.411167in}{2.303236in}}%
\pgfpathlineto{\pgfqpoint{4.411975in}{2.304855in}}%
\pgfpathlineto{\pgfqpoint{4.411512in}{2.301514in}}%
\pgfpathlineto{\pgfqpoint{4.412223in}{2.302242in}}%
\pgfpathlineto{\pgfqpoint{4.413138in}{2.301569in}}%
\pgfpathlineto{\pgfqpoint{4.412476in}{2.305227in}}%
\pgfpathlineto{\pgfqpoint{4.413338in}{2.301700in}}%
\pgfpathlineto{\pgfqpoint{4.413353in}{2.302012in}}%
\pgfpathlineto{\pgfqpoint{4.413922in}{2.298986in}}%
\pgfpathlineto{\pgfqpoint{4.414331in}{2.299507in}}%
\pgfpathlineto{\pgfqpoint{4.414584in}{2.299245in}}%
\pgfpathlineto{\pgfqpoint{4.415242in}{2.301746in}}%
\pgfpathlineto{\pgfqpoint{4.415266in}{2.301334in}}%
\pgfpathlineto{\pgfqpoint{4.415787in}{2.303833in}}%
\pgfpathlineto{\pgfqpoint{4.415305in}{2.301184in}}%
\pgfpathlineto{\pgfqpoint{4.416415in}{2.302911in}}%
\pgfpathlineto{\pgfqpoint{4.416551in}{2.300854in}}%
\pgfpathlineto{\pgfqpoint{4.417437in}{2.304195in}}%
\pgfpathlineto{\pgfqpoint{4.417505in}{2.303599in}}%
\pgfpathlineto{\pgfqpoint{4.417725in}{2.304826in}}%
\pgfpathlineto{\pgfqpoint{4.418270in}{2.301807in}}%
\pgfpathlineto{\pgfqpoint{4.418533in}{2.302999in}}%
\pgfpathlineto{\pgfqpoint{4.419097in}{2.301475in}}%
\pgfpathlineto{\pgfqpoint{4.418747in}{2.303693in}}%
\pgfpathlineto{\pgfqpoint{4.419657in}{2.302247in}}%
\pgfpathlineto{\pgfqpoint{4.420066in}{2.305279in}}%
\pgfpathlineto{\pgfqpoint{4.420724in}{2.301314in}}%
\pgfpathlineto{\pgfqpoint{4.421220in}{2.298573in}}%
\pgfpathlineto{\pgfqpoint{4.420884in}{2.301437in}}%
\pgfpathlineto{\pgfqpoint{4.421887in}{2.300064in}}%
\pgfpathlineto{\pgfqpoint{4.422725in}{2.301313in}}%
\pgfpathlineto{\pgfqpoint{4.422452in}{2.299056in}}%
\pgfpathlineto{\pgfqpoint{4.423041in}{2.301170in}}%
\pgfpathlineto{\pgfqpoint{4.423796in}{2.299433in}}%
\pgfpathlineto{\pgfqpoint{4.423445in}{2.301965in}}%
\pgfpathlineto{\pgfqpoint{4.424175in}{2.300718in}}%
\pgfpathlineto{\pgfqpoint{4.424243in}{2.301154in}}%
\pgfpathlineto{\pgfqpoint{4.424721in}{2.298390in}}%
\pgfpathlineto{\pgfqpoint{4.425105in}{2.299206in}}%
\pgfpathlineto{\pgfqpoint{4.425553in}{2.297244in}}%
\pgfpathlineto{\pgfqpoint{4.425285in}{2.299995in}}%
\pgfpathlineto{\pgfqpoint{4.426225in}{2.298662in}}%
\pgfpathlineto{\pgfqpoint{4.426692in}{2.298063in}}%
\pgfpathlineto{\pgfqpoint{4.427140in}{2.300884in}}%
\pgfpathlineto{\pgfqpoint{4.427262in}{2.300707in}}%
\pgfpathlineto{\pgfqpoint{4.427272in}{2.300618in}}%
\pgfpathlineto{\pgfqpoint{4.427836in}{2.303556in}}%
\pgfpathlineto{\pgfqpoint{4.428143in}{2.302466in}}%
\pgfpathlineto{\pgfqpoint{4.428708in}{2.303286in}}%
\pgfpathlineto{\pgfqpoint{4.428289in}{2.300942in}}%
\pgfpathlineto{\pgfqpoint{4.429229in}{2.301305in}}%
\pgfpathlineto{\pgfqpoint{4.429628in}{2.298482in}}%
\pgfpathlineto{\pgfqpoint{4.430271in}{2.303178in}}%
\pgfpathlineto{\pgfqpoint{4.430305in}{2.302729in}}%
\pgfpathlineto{\pgfqpoint{4.430792in}{2.304235in}}%
\pgfpathlineto{\pgfqpoint{4.431337in}{2.301235in}}%
\pgfpathlineto{\pgfqpoint{4.431405in}{2.302387in}}%
\pgfpathlineto{\pgfqpoint{4.432583in}{2.297283in}}%
\pgfpathlineto{\pgfqpoint{4.431760in}{2.302418in}}%
\pgfpathlineto{\pgfqpoint{4.432588in}{2.297577in}}%
\pgfpathlineto{\pgfqpoint{4.432647in}{2.298087in}}%
\pgfpathlineto{\pgfqpoint{4.432710in}{2.295968in}}%
\pgfpathlineto{\pgfqpoint{4.432793in}{2.296309in}}%
\pgfpathlineto{\pgfqpoint{4.432827in}{2.295348in}}%
\pgfpathlineto{\pgfqpoint{4.433664in}{2.301805in}}%
\pgfpathlineto{\pgfqpoint{4.433834in}{2.301218in}}%
\pgfpathlineto{\pgfqpoint{4.434102in}{2.303310in}}%
\pgfpathlineto{\pgfqpoint{4.434560in}{2.298767in}}%
\pgfpathlineto{\pgfqpoint{4.434910in}{2.299855in}}%
\pgfpathlineto{\pgfqpoint{4.435373in}{2.299007in}}%
\pgfpathlineto{\pgfqpoint{4.435947in}{2.303116in}}%
\pgfpathlineto{\pgfqpoint{4.435957in}{2.302823in}}%
\pgfpathlineto{\pgfqpoint{4.435981in}{2.303341in}}%
\pgfpathlineto{\pgfqpoint{4.436843in}{2.299004in}}%
\pgfpathlineto{\pgfqpoint{4.436975in}{2.299371in}}%
\pgfpathlineto{\pgfqpoint{4.437827in}{2.297005in}}%
\pgfpathlineto{\pgfqpoint{4.438055in}{2.299996in}}%
\pgfpathlineto{\pgfqpoint{4.438070in}{2.299551in}}%
\pgfpathlineto{\pgfqpoint{4.438158in}{2.300189in}}%
\pgfpathlineto{\pgfqpoint{4.439088in}{2.296287in}}%
\pgfpathlineto{\pgfqpoint{4.439112in}{2.296647in}}%
\pgfpathlineto{\pgfqpoint{4.439394in}{2.298488in}}%
\pgfpathlineto{\pgfqpoint{4.439165in}{2.295186in}}%
\pgfpathlineto{\pgfqpoint{4.440207in}{2.296436in}}%
\pgfpathlineto{\pgfqpoint{4.441356in}{2.291477in}}%
\pgfpathlineto{\pgfqpoint{4.441395in}{2.291825in}}%
\pgfpathlineto{\pgfqpoint{4.441872in}{2.293476in}}%
\pgfpathlineto{\pgfqpoint{4.442218in}{2.288060in}}%
\pgfpathlineto{\pgfqpoint{4.442457in}{2.288502in}}%
\pgfpathlineto{\pgfqpoint{4.442642in}{2.286355in}}%
\pgfpathlineto{\pgfqpoint{4.442656in}{2.285814in}}%
\pgfpathlineto{\pgfqpoint{4.443542in}{2.292091in}}%
\pgfpathlineto{\pgfqpoint{4.443571in}{2.292025in}}%
\pgfpathlineto{\pgfqpoint{4.444380in}{2.293323in}}%
\pgfpathlineto{\pgfqpoint{4.444088in}{2.291060in}}%
\pgfpathlineto{\pgfqpoint{4.444691in}{2.292279in}}%
\pgfpathlineto{\pgfqpoint{4.445645in}{2.288673in}}%
\pgfpathlineto{\pgfqpoint{4.444886in}{2.293380in}}%
\pgfpathlineto{\pgfqpoint{4.445865in}{2.291384in}}%
\pgfpathlineto{\pgfqpoint{4.446463in}{2.289939in}}%
\pgfpathlineto{\pgfqpoint{4.447150in}{2.294212in}}%
\pgfpathlineto{\pgfqpoint{4.447895in}{2.292543in}}%
\pgfpathlineto{\pgfqpoint{4.448211in}{2.296367in}}%
\pgfpathlineto{\pgfqpoint{4.449112in}{2.294006in}}%
\pgfpathlineto{\pgfqpoint{4.449282in}{2.296809in}}%
\pgfpathlineto{\pgfqpoint{4.449311in}{2.296374in}}%
\pgfpathlineto{\pgfqpoint{4.450027in}{2.299512in}}%
\pgfpathlineto{\pgfqpoint{4.450305in}{2.295496in}}%
\pgfpathlineto{\pgfqpoint{4.450412in}{2.296029in}}%
\pgfpathlineto{\pgfqpoint{4.451497in}{2.295045in}}%
\pgfpathlineto{\pgfqpoint{4.450816in}{2.297005in}}%
\pgfpathlineto{\pgfqpoint{4.451541in}{2.295181in}}%
\pgfpathlineto{\pgfqpoint{4.452082in}{2.295992in}}%
\pgfpathlineto{\pgfqpoint{4.452242in}{2.293469in}}%
\pgfpathlineto{\pgfqpoint{4.452637in}{2.294574in}}%
\pgfpathlineto{\pgfqpoint{4.452734in}{2.295545in}}%
\pgfpathlineto{\pgfqpoint{4.453055in}{2.293427in}}%
\pgfpathlineto{\pgfqpoint{4.453255in}{2.292355in}}%
\pgfpathlineto{\pgfqpoint{4.453591in}{2.295396in}}%
\pgfpathlineto{\pgfqpoint{4.454156in}{2.294259in}}%
\pgfpathlineto{\pgfqpoint{4.454185in}{2.293639in}}%
\pgfpathlineto{\pgfqpoint{4.454506in}{2.297323in}}%
\pgfpathlineto{\pgfqpoint{4.455222in}{2.296061in}}%
\pgfpathlineto{\pgfqpoint{4.455402in}{2.295227in}}%
\pgfpathlineto{\pgfqpoint{4.456001in}{2.299557in}}%
\pgfpathlineto{\pgfqpoint{4.456147in}{2.298003in}}%
\pgfpathlineto{\pgfqpoint{4.457013in}{2.301531in}}%
\pgfpathlineto{\pgfqpoint{4.456157in}{2.297976in}}%
\pgfpathlineto{\pgfqpoint{4.457310in}{2.301323in}}%
\pgfpathlineto{\pgfqpoint{4.457383in}{2.300486in}}%
\pgfpathlineto{\pgfqpoint{4.458294in}{2.304967in}}%
\pgfpathlineto{\pgfqpoint{4.458308in}{2.304789in}}%
\pgfpathlineto{\pgfqpoint{4.459311in}{2.306441in}}%
\pgfpathlineto{\pgfqpoint{4.458912in}{2.302579in}}%
\pgfpathlineto{\pgfqpoint{4.459428in}{2.305371in}}%
\pgfpathlineto{\pgfqpoint{4.460266in}{2.301917in}}%
\pgfpathlineto{\pgfqpoint{4.460553in}{2.304403in}}%
\pgfpathlineto{\pgfqpoint{4.461298in}{2.301253in}}%
\pgfpathlineto{\pgfqpoint{4.461064in}{2.305125in}}%
\pgfpathlineto{\pgfqpoint{4.461770in}{2.303490in}}%
\pgfpathlineto{\pgfqpoint{4.461775in}{2.303891in}}%
\pgfpathlineto{\pgfqpoint{4.462062in}{2.301463in}}%
\pgfpathlineto{\pgfqpoint{4.462870in}{2.303052in}}%
\pgfpathlineto{\pgfqpoint{4.463143in}{2.301949in}}%
\pgfpathlineto{\pgfqpoint{4.463834in}{2.304629in}}%
\pgfpathlineto{\pgfqpoint{4.463995in}{2.302481in}}%
\pgfpathlineto{\pgfqpoint{4.464219in}{2.300330in}}%
\pgfpathlineto{\pgfqpoint{4.464808in}{2.303092in}}%
\pgfpathlineto{\pgfqpoint{4.464827in}{2.303064in}}%
\pgfpathlineto{\pgfqpoint{4.465947in}{2.304908in}}%
\pgfpathlineto{\pgfqpoint{4.465553in}{2.300408in}}%
\pgfpathlineto{\pgfqpoint{4.465952in}{2.304895in}}%
\pgfpathlineto{\pgfqpoint{4.466512in}{2.300925in}}%
\pgfpathlineto{\pgfqpoint{4.465967in}{2.304911in}}%
\pgfpathlineto{\pgfqpoint{4.467227in}{2.301760in}}%
\pgfpathlineto{\pgfqpoint{4.468800in}{2.303977in}}%
\pgfpathlineto{\pgfqpoint{4.467636in}{2.300751in}}%
\pgfpathlineto{\pgfqpoint{4.469087in}{2.302251in}}%
\pgfpathlineto{\pgfqpoint{4.470295in}{2.298000in}}%
\pgfpathlineto{\pgfqpoint{4.469238in}{2.303549in}}%
\pgfpathlineto{\pgfqpoint{4.470455in}{2.299510in}}%
\pgfpathlineto{\pgfqpoint{4.471741in}{2.304172in}}%
\pgfpathlineto{\pgfqpoint{4.470606in}{2.298572in}}%
\pgfpathlineto{\pgfqpoint{4.471745in}{2.303756in}}%
\pgfpathlineto{\pgfqpoint{4.472247in}{2.300832in}}%
\pgfpathlineto{\pgfqpoint{4.471857in}{2.304875in}}%
\pgfpathlineto{\pgfqpoint{4.472860in}{2.303415in}}%
\pgfpathlineto{\pgfqpoint{4.473289in}{2.301848in}}%
\pgfpathlineto{\pgfqpoint{4.473586in}{2.304976in}}%
\pgfpathlineto{\pgfqpoint{4.473990in}{2.303050in}}%
\pgfpathlineto{\pgfqpoint{4.474511in}{2.301829in}}%
\pgfpathlineto{\pgfqpoint{4.474160in}{2.303881in}}%
\pgfpathlineto{\pgfqpoint{4.475100in}{2.302985in}}%
\pgfpathlineto{\pgfqpoint{4.475173in}{2.303730in}}%
\pgfpathlineto{\pgfqpoint{4.476103in}{2.301548in}}%
\pgfpathlineto{\pgfqpoint{4.476195in}{2.302543in}}%
\pgfpathlineto{\pgfqpoint{4.477135in}{2.298747in}}%
\pgfpathlineto{\pgfqpoint{4.477330in}{2.301399in}}%
\pgfpathlineto{\pgfqpoint{4.477520in}{2.302654in}}%
\pgfpathlineto{\pgfqpoint{4.478060in}{2.300144in}}%
\pgfpathlineto{\pgfqpoint{4.478440in}{2.301610in}}%
\pgfpathlineto{\pgfqpoint{4.478654in}{2.300068in}}%
\pgfpathlineto{\pgfqpoint{4.479034in}{2.303633in}}%
\pgfpathlineto{\pgfqpoint{4.479535in}{2.302399in}}%
\pgfpathlineto{\pgfqpoint{4.480231in}{2.301105in}}%
\pgfpathlineto{\pgfqpoint{4.480450in}{2.304353in}}%
\pgfpathlineto{\pgfqpoint{4.480587in}{2.303997in}}%
\pgfpathlineto{\pgfqpoint{4.480966in}{2.303906in}}%
\pgfpathlineto{\pgfqpoint{4.480811in}{2.302456in}}%
\pgfpathlineto{\pgfqpoint{4.481161in}{2.302791in}}%
\pgfpathlineto{\pgfqpoint{4.481234in}{2.302578in}}%
\pgfpathlineto{\pgfqpoint{4.482052in}{2.304679in}}%
\pgfpathlineto{\pgfqpoint{4.482266in}{2.302713in}}%
\pgfpathlineto{\pgfqpoint{4.483162in}{2.304486in}}%
\pgfpathlineto{\pgfqpoint{4.482646in}{2.301166in}}%
\pgfpathlineto{\pgfqpoint{4.483391in}{2.303472in}}%
\pgfpathlineto{\pgfqpoint{4.483956in}{2.299865in}}%
\pgfpathlineto{\pgfqpoint{4.483629in}{2.303937in}}%
\pgfpathlineto{\pgfqpoint{4.484598in}{2.300384in}}%
\pgfpathlineto{\pgfqpoint{4.485888in}{2.304226in}}%
\pgfpathlineto{\pgfqpoint{4.484662in}{2.299254in}}%
\pgfpathlineto{\pgfqpoint{4.485918in}{2.303740in}}%
\pgfpathlineto{\pgfqpoint{4.486492in}{2.302650in}}%
\pgfpathlineto{\pgfqpoint{4.486341in}{2.304927in}}%
\pgfpathlineto{\pgfqpoint{4.486643in}{2.304908in}}%
\pgfpathlineto{\pgfqpoint{4.486682in}{2.305295in}}%
\pgfpathlineto{\pgfqpoint{4.487592in}{2.302302in}}%
\pgfpathlineto{\pgfqpoint{4.487651in}{2.302789in}}%
\pgfpathlineto{\pgfqpoint{4.488717in}{2.300781in}}%
\pgfpathlineto{\pgfqpoint{4.487933in}{2.303943in}}%
\pgfpathlineto{\pgfqpoint{4.488829in}{2.301506in}}%
\pgfpathlineto{\pgfqpoint{4.489812in}{2.303424in}}%
\pgfpathlineto{\pgfqpoint{4.489194in}{2.299950in}}%
\pgfpathlineto{\pgfqpoint{4.489959in}{2.302767in}}%
\pgfpathlineto{\pgfqpoint{4.490952in}{2.298185in}}%
\pgfpathlineto{\pgfqpoint{4.490173in}{2.304295in}}%
\pgfpathlineto{\pgfqpoint{4.491127in}{2.300367in}}%
\pgfpathlineto{\pgfqpoint{4.492398in}{2.302937in}}%
\pgfpathlineto{\pgfqpoint{4.491521in}{2.299606in}}%
\pgfpathlineto{\pgfqpoint{4.492412in}{2.302595in}}%
\pgfpathlineto{\pgfqpoint{4.492958in}{2.300489in}}%
\pgfpathlineto{\pgfqpoint{4.493113in}{2.304061in}}%
\pgfpathlineto{\pgfqpoint{4.493522in}{2.302623in}}%
\pgfpathlineto{\pgfqpoint{4.494063in}{2.301001in}}%
\pgfpathlineto{\pgfqpoint{4.493736in}{2.303576in}}%
\pgfpathlineto{\pgfqpoint{4.494681in}{2.301882in}}%
\pgfpathlineto{\pgfqpoint{4.495801in}{2.306013in}}%
\pgfpathlineto{\pgfqpoint{4.495845in}{2.305597in}}%
\pgfpathlineto{\pgfqpoint{4.496857in}{2.302502in}}%
\pgfpathlineto{\pgfqpoint{4.495864in}{2.305864in}}%
\pgfpathlineto{\pgfqpoint{4.496989in}{2.303776in}}%
\pgfpathlineto{\pgfqpoint{4.497466in}{2.304133in}}%
\pgfpathlineto{\pgfqpoint{4.497276in}{2.301915in}}%
\pgfpathlineto{\pgfqpoint{4.498089in}{2.303677in}}%
\pgfpathlineto{\pgfqpoint{4.498133in}{2.302702in}}%
\pgfpathlineto{\pgfqpoint{4.499102in}{2.305932in}}%
\pgfpathlineto{\pgfqpoint{4.499175in}{2.305495in}}%
\pgfpathlineto{\pgfqpoint{4.499214in}{2.306397in}}%
\pgfpathlineto{\pgfqpoint{4.500192in}{2.300834in}}%
\pgfpathlineto{\pgfqpoint{4.500226in}{2.300983in}}%
\pgfpathlineto{\pgfqpoint{4.501229in}{2.303674in}}%
\pgfpathlineto{\pgfqpoint{4.500513in}{2.299716in}}%
\pgfpathlineto{\pgfqpoint{4.501375in}{2.303021in}}%
\pgfpathlineto{\pgfqpoint{4.501497in}{2.301557in}}%
\pgfpathlineto{\pgfqpoint{4.502115in}{2.304464in}}%
\pgfpathlineto{\pgfqpoint{4.502500in}{2.302256in}}%
\pgfpathlineto{\pgfqpoint{4.503035in}{2.299721in}}%
\pgfpathlineto{\pgfqpoint{4.502709in}{2.302957in}}%
\pgfpathlineto{\pgfqpoint{4.503702in}{2.301319in}}%
\pgfpathlineto{\pgfqpoint{4.503975in}{2.302502in}}%
\pgfpathlineto{\pgfqpoint{4.504661in}{2.298353in}}%
\pgfpathlineto{\pgfqpoint{4.504793in}{2.299965in}}%
\pgfpathlineto{\pgfqpoint{4.505333in}{2.304949in}}%
\pgfpathlineto{\pgfqpoint{4.505952in}{2.301559in}}%
\pgfpathlineto{\pgfqpoint{4.505956in}{2.301185in}}%
\pgfpathlineto{\pgfqpoint{4.506886in}{2.304128in}}%
\pgfpathlineto{\pgfqpoint{4.507037in}{2.303400in}}%
\pgfpathlineto{\pgfqpoint{4.508089in}{2.302381in}}%
\pgfpathlineto{\pgfqpoint{4.508157in}{2.304410in}}%
\pgfpathlineto{\pgfqpoint{4.508882in}{2.302440in}}%
\pgfpathlineto{\pgfqpoint{4.509136in}{2.305222in}}%
\pgfpathlineto{\pgfqpoint{4.509345in}{2.302983in}}%
\pgfpathlineto{\pgfqpoint{4.509807in}{2.304761in}}%
\pgfpathlineto{\pgfqpoint{4.510323in}{2.301059in}}%
\pgfpathlineto{\pgfqpoint{4.510382in}{2.301797in}}%
\pgfpathlineto{\pgfqpoint{4.511409in}{2.304484in}}%
\pgfpathlineto{\pgfqpoint{4.511726in}{2.303470in}}%
\pgfpathlineto{\pgfqpoint{4.512743in}{2.301690in}}%
\pgfpathlineto{\pgfqpoint{4.512519in}{2.304333in}}%
\pgfpathlineto{\pgfqpoint{4.512865in}{2.301925in}}%
\pgfpathlineto{\pgfqpoint{4.513025in}{2.300856in}}%
\pgfpathlineto{\pgfqpoint{4.513332in}{2.304438in}}%
\pgfpathlineto{\pgfqpoint{4.513848in}{2.303083in}}%
\pgfpathlineto{\pgfqpoint{4.514841in}{2.304771in}}%
\pgfpathlineto{\pgfqpoint{4.514359in}{2.301670in}}%
\pgfpathlineto{\pgfqpoint{4.514968in}{2.304146in}}%
\pgfpathlineto{\pgfqpoint{4.515465in}{2.301999in}}%
\pgfpathlineto{\pgfqpoint{4.516044in}{2.304361in}}%
\pgfpathlineto{\pgfqpoint{4.516098in}{2.303161in}}%
\pgfpathlineto{\pgfqpoint{4.516317in}{2.304612in}}%
\pgfpathlineto{\pgfqpoint{4.516901in}{2.299644in}}%
\pgfpathlineto{\pgfqpoint{4.516979in}{2.299956in}}%
\pgfpathlineto{\pgfqpoint{4.518098in}{2.297220in}}%
\pgfpathlineto{\pgfqpoint{4.517621in}{2.300799in}}%
\pgfpathlineto{\pgfqpoint{4.518113in}{2.297519in}}%
\pgfpathlineto{\pgfqpoint{4.518337in}{2.295875in}}%
\pgfpathlineto{\pgfqpoint{4.518873in}{2.298848in}}%
\pgfpathlineto{\pgfqpoint{4.519286in}{2.296123in}}%
\pgfpathlineto{\pgfqpoint{4.519690in}{2.298637in}}%
\pgfpathlineto{\pgfqpoint{4.520343in}{2.294933in}}%
\pgfpathlineto{\pgfqpoint{4.520392in}{2.295963in}}%
\pgfpathlineto{\pgfqpoint{4.520650in}{2.295064in}}%
\pgfpathlineto{\pgfqpoint{4.521209in}{2.298790in}}%
\pgfpathlineto{\pgfqpoint{4.521467in}{2.297606in}}%
\pgfpathlineto{\pgfqpoint{4.521833in}{2.299871in}}%
\pgfpathlineto{\pgfqpoint{4.521599in}{2.296927in}}%
\pgfpathlineto{\pgfqpoint{4.522607in}{2.298835in}}%
\pgfpathlineto{\pgfqpoint{4.522626in}{2.298269in}}%
\pgfpathlineto{\pgfqpoint{4.523576in}{2.303286in}}%
\pgfpathlineto{\pgfqpoint{4.523614in}{2.303007in}}%
\pgfpathlineto{\pgfqpoint{4.524647in}{2.304620in}}%
\pgfpathlineto{\pgfqpoint{4.524277in}{2.301751in}}%
\pgfpathlineto{\pgfqpoint{4.524744in}{2.303672in}}%
\pgfpathlineto{\pgfqpoint{4.525523in}{2.299906in}}%
\pgfpathlineto{\pgfqpoint{4.524905in}{2.303880in}}%
\pgfpathlineto{\pgfqpoint{4.525917in}{2.300979in}}%
\pgfpathlineto{\pgfqpoint{4.527134in}{2.305155in}}%
\pgfpathlineto{\pgfqpoint{4.527159in}{2.304439in}}%
\pgfpathlineto{\pgfqpoint{4.527558in}{2.301320in}}%
\pgfpathlineto{\pgfqpoint{4.528332in}{2.302849in}}%
\pgfpathlineto{\pgfqpoint{4.528347in}{2.303484in}}%
\pgfpathlineto{\pgfqpoint{4.528644in}{2.300870in}}%
\pgfpathlineto{\pgfqpoint{4.529423in}{2.301726in}}%
\pgfpathlineto{\pgfqpoint{4.529637in}{2.300209in}}%
\pgfpathlineto{\pgfqpoint{4.530172in}{2.305325in}}%
\pgfpathlineto{\pgfqpoint{4.530469in}{2.303613in}}%
\pgfpathlineto{\pgfqpoint{4.531657in}{2.306126in}}%
\pgfpathlineto{\pgfqpoint{4.530674in}{2.302554in}}%
\pgfpathlineto{\pgfqpoint{4.531740in}{2.304925in}}%
\pgfpathlineto{\pgfqpoint{4.532237in}{2.302059in}}%
\pgfpathlineto{\pgfqpoint{4.532733in}{2.306648in}}%
\pgfpathlineto{\pgfqpoint{4.532816in}{2.306284in}}%
\pgfpathlineto{\pgfqpoint{4.532845in}{2.306852in}}%
\pgfpathlineto{\pgfqpoint{4.533717in}{2.302253in}}%
\pgfpathlineto{\pgfqpoint{4.533760in}{2.302711in}}%
\pgfpathlineto{\pgfqpoint{4.533965in}{2.302024in}}%
\pgfpathlineto{\pgfqpoint{4.534252in}{2.304604in}}%
\pgfpathlineto{\pgfqpoint{4.534783in}{2.303651in}}%
\pgfpathlineto{\pgfqpoint{4.535021in}{2.306555in}}%
\pgfpathlineto{\pgfqpoint{4.535430in}{2.302502in}}%
\pgfpathlineto{\pgfqpoint{4.535888in}{2.303552in}}%
\pgfpathlineto{\pgfqpoint{4.536774in}{2.301361in}}%
\pgfpathlineto{\pgfqpoint{4.535966in}{2.303863in}}%
\pgfpathlineto{\pgfqpoint{4.537032in}{2.302310in}}%
\pgfpathlineto{\pgfqpoint{4.537699in}{2.304265in}}%
\pgfpathlineto{\pgfqpoint{4.537164in}{2.302250in}}%
\pgfpathlineto{\pgfqpoint{4.538235in}{2.303632in}}%
\pgfpathlineto{\pgfqpoint{4.539335in}{2.301828in}}%
\pgfpathlineto{\pgfqpoint{4.538351in}{2.304201in}}%
\pgfpathlineto{\pgfqpoint{4.539384in}{2.302589in}}%
\pgfpathlineto{\pgfqpoint{4.539836in}{2.304174in}}%
\pgfpathlineto{\pgfqpoint{4.540328in}{2.301171in}}%
\pgfpathlineto{\pgfqpoint{4.540513in}{2.303469in}}%
\pgfpathlineto{\pgfqpoint{4.541015in}{2.302238in}}%
\pgfpathlineto{\pgfqpoint{4.540630in}{2.304714in}}%
\pgfpathlineto{\pgfqpoint{4.541599in}{2.304381in}}%
\pgfpathlineto{\pgfqpoint{4.542694in}{2.304535in}}%
\pgfpathlineto{\pgfqpoint{4.541940in}{2.300541in}}%
\pgfpathlineto{\pgfqpoint{4.542704in}{2.304465in}}%
\pgfpathlineto{\pgfqpoint{4.543814in}{2.299879in}}%
\pgfpathlineto{\pgfqpoint{4.542845in}{2.305196in}}%
\pgfpathlineto{\pgfqpoint{4.543916in}{2.301394in}}%
\pgfpathlineto{\pgfqpoint{4.544909in}{2.304103in}}%
\pgfpathlineto{\pgfqpoint{4.543931in}{2.301191in}}%
\pgfpathlineto{\pgfqpoint{4.545089in}{2.301979in}}%
\pgfpathlineto{\pgfqpoint{4.545333in}{2.301538in}}%
\pgfpathlineto{\pgfqpoint{4.545644in}{2.304746in}}%
\pgfpathlineto{\pgfqpoint{4.546190in}{2.302284in}}%
\pgfpathlineto{\pgfqpoint{4.546404in}{2.303259in}}%
\pgfpathlineto{\pgfqpoint{4.546876in}{2.300975in}}%
\pgfpathlineto{\pgfqpoint{4.547329in}{2.302541in}}%
\pgfpathlineto{\pgfqpoint{4.548478in}{2.298210in}}%
\pgfpathlineto{\pgfqpoint{4.547387in}{2.303135in}}%
\pgfpathlineto{\pgfqpoint{4.548483in}{2.298391in}}%
\pgfpathlineto{\pgfqpoint{4.549763in}{2.302178in}}%
\pgfpathlineto{\pgfqpoint{4.548751in}{2.296527in}}%
\pgfpathlineto{\pgfqpoint{4.549773in}{2.302006in}}%
\pgfpathlineto{\pgfqpoint{4.550547in}{2.303939in}}%
\pgfpathlineto{\pgfqpoint{4.550274in}{2.300298in}}%
\pgfpathlineto{\pgfqpoint{4.550854in}{2.301803in}}%
\pgfpathlineto{\pgfqpoint{4.550883in}{2.301562in}}%
\pgfpathlineto{\pgfqpoint{4.551842in}{2.306640in}}%
\pgfpathlineto{\pgfqpoint{4.552163in}{2.308369in}}%
\pgfpathlineto{\pgfqpoint{4.552806in}{2.302602in}}%
\pgfpathlineto{\pgfqpoint{4.552894in}{2.303850in}}%
\pgfpathlineto{\pgfqpoint{4.553682in}{2.301467in}}%
\pgfpathlineto{\pgfqpoint{4.552981in}{2.304324in}}%
\pgfpathlineto{\pgfqpoint{4.554077in}{2.303382in}}%
\pgfpathlineto{\pgfqpoint{4.554325in}{2.304017in}}%
\pgfpathlineto{\pgfqpoint{4.554826in}{2.300029in}}%
\pgfpathlineto{\pgfqpoint{4.554968in}{2.300948in}}%
\pgfpathlineto{\pgfqpoint{4.555971in}{2.295481in}}%
\pgfpathlineto{\pgfqpoint{4.555133in}{2.301323in}}%
\pgfpathlineto{\pgfqpoint{4.556136in}{2.297208in}}%
\pgfpathlineto{\pgfqpoint{4.556827in}{2.299100in}}%
\pgfpathlineto{\pgfqpoint{4.556238in}{2.296019in}}%
\pgfpathlineto{\pgfqpoint{4.557305in}{2.298025in}}%
\pgfpathlineto{\pgfqpoint{4.557314in}{2.297903in}}%
\pgfpathlineto{\pgfqpoint{4.558025in}{2.301993in}}%
\pgfpathlineto{\pgfqpoint{4.558166in}{2.300487in}}%
\pgfpathlineto{\pgfqpoint{4.559174in}{2.305043in}}%
\pgfpathlineto{\pgfqpoint{4.559330in}{2.304404in}}%
\pgfpathlineto{\pgfqpoint{4.560157in}{2.305421in}}%
\pgfpathlineto{\pgfqpoint{4.559807in}{2.302638in}}%
\pgfpathlineto{\pgfqpoint{4.560415in}{2.304166in}}%
\pgfpathlineto{\pgfqpoint{4.560907in}{2.302885in}}%
\pgfpathlineto{\pgfqpoint{4.561170in}{2.305545in}}%
\pgfpathlineto{\pgfqpoint{4.561521in}{2.304268in}}%
\pgfpathlineto{\pgfqpoint{4.561564in}{2.304469in}}%
\pgfpathlineto{\pgfqpoint{4.562339in}{2.300523in}}%
\pgfpathlineto{\pgfqpoint{4.562382in}{2.300711in}}%
\pgfpathlineto{\pgfqpoint{4.562489in}{2.300293in}}%
\pgfpathlineto{\pgfqpoint{4.562811in}{2.303320in}}%
\pgfpathlineto{\pgfqpoint{4.563483in}{2.300789in}}%
\pgfpathlineto{\pgfqpoint{4.563974in}{2.300161in}}%
\pgfpathlineto{\pgfqpoint{4.564632in}{2.302134in}}%
\pgfpathlineto{\pgfqpoint{4.564909in}{2.301096in}}%
\pgfpathlineto{\pgfqpoint{4.565600in}{2.305489in}}%
\pgfpathlineto{\pgfqpoint{4.565888in}{2.306422in}}%
\pgfpathlineto{\pgfqpoint{4.566448in}{2.300980in}}%
\pgfpathlineto{\pgfqpoint{4.566647in}{2.301700in}}%
\pgfpathlineto{\pgfqpoint{4.566764in}{2.300865in}}%
\pgfpathlineto{\pgfqpoint{4.567499in}{2.304478in}}%
\pgfpathlineto{\pgfqpoint{4.567718in}{2.303004in}}%
\pgfpathlineto{\pgfqpoint{4.568098in}{2.304549in}}%
\pgfpathlineto{\pgfqpoint{4.568634in}{2.301948in}}%
\pgfpathlineto{\pgfqpoint{4.568843in}{2.304209in}}%
\pgfpathlineto{\pgfqpoint{4.569266in}{2.305235in}}%
\pgfpathlineto{\pgfqpoint{4.569753in}{2.302504in}}%
\pgfpathlineto{\pgfqpoint{4.570561in}{2.300814in}}%
\pgfpathlineto{\pgfqpoint{4.570235in}{2.303483in}}%
\pgfpathlineto{\pgfqpoint{4.570878in}{2.302170in}}%
\pgfpathlineto{\pgfqpoint{4.570902in}{2.303242in}}%
\pgfpathlineto{\pgfqpoint{4.571584in}{2.300277in}}%
\pgfpathlineto{\pgfqpoint{4.571973in}{2.301843in}}%
\pgfpathlineto{\pgfqpoint{4.572694in}{2.300499in}}%
\pgfpathlineto{\pgfqpoint{4.572158in}{2.303439in}}%
\pgfpathlineto{\pgfqpoint{4.572908in}{2.303462in}}%
\pgfpathlineto{\pgfqpoint{4.572981in}{2.304253in}}%
\pgfpathlineto{\pgfqpoint{4.573741in}{2.300695in}}%
\pgfpathlineto{\pgfqpoint{4.574003in}{2.302902in}}%
\pgfpathlineto{\pgfqpoint{4.574349in}{2.302023in}}%
\pgfpathlineto{\pgfqpoint{4.574632in}{2.305252in}}%
\pgfpathlineto{\pgfqpoint{4.575031in}{2.304042in}}%
\pgfpathlineto{\pgfqpoint{4.575863in}{2.300435in}}%
\pgfpathlineto{\pgfqpoint{4.576155in}{2.304887in}}%
\pgfpathlineto{\pgfqpoint{4.576681in}{2.300718in}}%
\pgfpathlineto{\pgfqpoint{4.577187in}{2.304949in}}%
\pgfpathlineto{\pgfqpoint{4.577270in}{2.304388in}}%
\pgfpathlineto{\pgfqpoint{4.577825in}{2.306858in}}%
\pgfpathlineto{\pgfqpoint{4.578171in}{2.303071in}}%
\pgfpathlineto{\pgfqpoint{4.578336in}{2.303122in}}%
\pgfpathlineto{\pgfqpoint{4.579554in}{2.299464in}}%
\pgfpathlineto{\pgfqpoint{4.579232in}{2.303660in}}%
\pgfpathlineto{\pgfqpoint{4.579578in}{2.300114in}}%
\pgfpathlineto{\pgfqpoint{4.580138in}{2.303674in}}%
\pgfpathlineto{\pgfqpoint{4.579773in}{2.299260in}}%
\pgfpathlineto{\pgfqpoint{4.580766in}{2.302777in}}%
\pgfpathlineto{\pgfqpoint{4.581900in}{2.300522in}}%
\pgfpathlineto{\pgfqpoint{4.580810in}{2.303301in}}%
\pgfpathlineto{\pgfqpoint{4.581944in}{2.301436in}}%
\pgfpathlineto{\pgfqpoint{4.581954in}{2.301743in}}%
\pgfpathlineto{\pgfqpoint{4.582265in}{2.298907in}}%
\pgfpathlineto{\pgfqpoint{4.583039in}{2.301121in}}%
\pgfpathlineto{\pgfqpoint{4.583551in}{2.298677in}}%
\pgfpathlineto{\pgfqpoint{4.583088in}{2.302206in}}%
\pgfpathlineto{\pgfqpoint{4.584193in}{2.298886in}}%
\pgfpathlineto{\pgfqpoint{4.584412in}{2.301723in}}%
\pgfpathlineto{\pgfqpoint{4.585367in}{2.301204in}}%
\pgfpathlineto{\pgfqpoint{4.585381in}{2.301038in}}%
\pgfpathlineto{\pgfqpoint{4.585474in}{2.302751in}}%
\pgfpathlineto{\pgfqpoint{4.585790in}{2.303765in}}%
\pgfpathlineto{\pgfqpoint{4.586443in}{2.299599in}}%
\pgfpathlineto{\pgfqpoint{4.586550in}{2.300175in}}%
\pgfpathlineto{\pgfqpoint{4.586574in}{2.299764in}}%
\pgfpathlineto{\pgfqpoint{4.586847in}{2.302192in}}%
\pgfpathlineto{\pgfqpoint{4.587660in}{2.300087in}}%
\pgfpathlineto{\pgfqpoint{4.587937in}{2.300838in}}%
\pgfpathlineto{\pgfqpoint{4.588336in}{2.298529in}}%
\pgfpathlineto{\pgfqpoint{4.588492in}{2.298646in}}%
\pgfpathlineto{\pgfqpoint{4.588551in}{2.298070in}}%
\pgfpathlineto{\pgfqpoint{4.588896in}{2.300974in}}%
\pgfpathlineto{\pgfqpoint{4.589544in}{2.300408in}}%
\pgfpathlineto{\pgfqpoint{4.590264in}{2.303983in}}%
\pgfpathlineto{\pgfqpoint{4.589578in}{2.300286in}}%
\pgfpathlineto{\pgfqpoint{4.590683in}{2.301612in}}%
\pgfpathlineto{\pgfqpoint{4.591092in}{2.299200in}}%
\pgfpathlineto{\pgfqpoint{4.591759in}{2.302673in}}%
\pgfpathlineto{\pgfqpoint{4.591769in}{2.302647in}}%
\pgfpathlineto{\pgfqpoint{4.592494in}{2.300148in}}%
\pgfpathlineto{\pgfqpoint{4.591851in}{2.303390in}}%
\pgfpathlineto{\pgfqpoint{4.592874in}{2.302838in}}%
\pgfpathlineto{\pgfqpoint{4.593107in}{2.303748in}}%
\pgfpathlineto{\pgfqpoint{4.593487in}{2.300138in}}%
\pgfpathlineto{\pgfqpoint{4.593916in}{2.301186in}}%
\pgfpathlineto{\pgfqpoint{4.594037in}{2.300864in}}%
\pgfpathlineto{\pgfqpoint{4.594252in}{2.303910in}}%
\pgfpathlineto{\pgfqpoint{4.594256in}{2.303832in}}%
\pgfpathlineto{\pgfqpoint{4.594271in}{2.304081in}}%
\pgfpathlineto{\pgfqpoint{4.595196in}{2.299793in}}%
\pgfpathlineto{\pgfqpoint{4.595235in}{2.300219in}}%
\pgfpathlineto{\pgfqpoint{4.596379in}{2.298777in}}%
\pgfpathlineto{\pgfqpoint{4.595727in}{2.303090in}}%
\pgfpathlineto{\pgfqpoint{4.596384in}{2.298833in}}%
\pgfpathlineto{\pgfqpoint{4.596890in}{2.302409in}}%
\pgfpathlineto{\pgfqpoint{4.596481in}{2.298499in}}%
\pgfpathlineto{\pgfqpoint{4.597523in}{2.299570in}}%
\pgfpathlineto{\pgfqpoint{4.597942in}{2.298468in}}%
\pgfpathlineto{\pgfqpoint{4.598263in}{2.302222in}}%
\pgfpathlineto{\pgfqpoint{4.598594in}{2.301385in}}%
\pgfpathlineto{\pgfqpoint{4.599242in}{2.305303in}}%
\pgfpathlineto{\pgfqpoint{4.598614in}{2.301200in}}%
\pgfpathlineto{\pgfqpoint{4.599724in}{2.302387in}}%
\pgfpathlineto{\pgfqpoint{4.600848in}{2.301272in}}%
\pgfpathlineto{\pgfqpoint{4.600284in}{2.303292in}}%
\pgfpathlineto{\pgfqpoint{4.600873in}{2.301561in}}%
\pgfpathlineto{\pgfqpoint{4.601627in}{2.304036in}}%
\pgfpathlineto{\pgfqpoint{4.600897in}{2.301494in}}%
\pgfpathlineto{\pgfqpoint{4.602007in}{2.303267in}}%
\pgfpathlineto{\pgfqpoint{4.603078in}{2.300007in}}%
\pgfpathlineto{\pgfqpoint{4.602046in}{2.303271in}}%
\pgfpathlineto{\pgfqpoint{4.603156in}{2.300653in}}%
\pgfpathlineto{\pgfqpoint{4.603200in}{2.301086in}}%
\pgfpathlineto{\pgfqpoint{4.603794in}{2.298291in}}%
\pgfpathlineto{\pgfqpoint{4.604130in}{2.298948in}}%
\pgfpathlineto{\pgfqpoint{4.604880in}{2.297572in}}%
\pgfpathlineto{\pgfqpoint{4.605196in}{2.300853in}}%
\pgfpathlineto{\pgfqpoint{4.605215in}{2.300288in}}%
\pgfpathlineto{\pgfqpoint{4.605225in}{2.300480in}}%
\pgfpathlineto{\pgfqpoint{4.606048in}{2.297415in}}%
\pgfpathlineto{\pgfqpoint{4.606238in}{2.298278in}}%
\pgfpathlineto{\pgfqpoint{4.606632in}{2.296147in}}%
\pgfpathlineto{\pgfqpoint{4.607139in}{2.298867in}}%
\pgfpathlineto{\pgfqpoint{4.607338in}{2.298618in}}%
\pgfpathlineto{\pgfqpoint{4.607343in}{2.298602in}}%
\pgfpathlineto{\pgfqpoint{4.607650in}{2.300805in}}%
\pgfpathlineto{\pgfqpoint{4.607869in}{2.300582in}}%
\pgfpathlineto{\pgfqpoint{4.608497in}{2.304939in}}%
\pgfpathlineto{\pgfqpoint{4.609042in}{2.303092in}}%
\pgfpathlineto{\pgfqpoint{4.610396in}{2.298486in}}%
\pgfpathlineto{\pgfqpoint{4.609631in}{2.305430in}}%
\pgfpathlineto{\pgfqpoint{4.610473in}{2.299690in}}%
\pgfpathlineto{\pgfqpoint{4.611603in}{2.301214in}}%
\pgfpathlineto{\pgfqpoint{4.610576in}{2.299201in}}%
\pgfpathlineto{\pgfqpoint{4.611608in}{2.301058in}}%
\pgfpathlineto{\pgfqpoint{4.611613in}{2.300818in}}%
\pgfpathlineto{\pgfqpoint{4.612211in}{2.304063in}}%
\pgfpathlineto{\pgfqpoint{4.612684in}{2.302674in}}%
\pgfpathlineto{\pgfqpoint{4.612713in}{2.302766in}}%
\pgfpathlineto{\pgfqpoint{4.612913in}{2.300372in}}%
\pgfpathlineto{\pgfqpoint{4.613039in}{2.299157in}}%
\pgfpathlineto{\pgfqpoint{4.613545in}{2.304779in}}%
\pgfpathlineto{\pgfqpoint{4.613954in}{2.303468in}}%
\pgfpathlineto{\pgfqpoint{4.613969in}{2.303596in}}%
\pgfpathlineto{\pgfqpoint{4.614787in}{2.300470in}}%
\pgfpathlineto{\pgfqpoint{4.614952in}{2.300615in}}%
\pgfpathlineto{\pgfqpoint{4.616028in}{2.303233in}}%
\pgfpathlineto{\pgfqpoint{4.615021in}{2.299114in}}%
\pgfpathlineto{\pgfqpoint{4.616218in}{2.302070in}}%
\pgfpathlineto{\pgfqpoint{4.616437in}{2.301118in}}%
\pgfpathlineto{\pgfqpoint{4.617158in}{2.304464in}}%
\pgfpathlineto{\pgfqpoint{4.617270in}{2.303716in}}%
\pgfpathlineto{\pgfqpoint{4.618175in}{2.302748in}}%
\pgfpathlineto{\pgfqpoint{4.618404in}{2.304604in}}%
\pgfpathlineto{\pgfqpoint{4.619212in}{2.300369in}}%
\pgfpathlineto{\pgfqpoint{4.619534in}{2.303116in}}%
\pgfpathlineto{\pgfqpoint{4.619826in}{2.300884in}}%
\pgfpathlineto{\pgfqpoint{4.620171in}{2.303683in}}%
\pgfpathlineto{\pgfqpoint{4.620663in}{2.301961in}}%
\pgfpathlineto{\pgfqpoint{4.621525in}{2.304489in}}%
\pgfpathlineto{\pgfqpoint{4.620868in}{2.299998in}}%
\pgfpathlineto{\pgfqpoint{4.621778in}{2.302329in}}%
\pgfpathlineto{\pgfqpoint{4.621788in}{2.302143in}}%
\pgfpathlineto{\pgfqpoint{4.622226in}{2.304843in}}%
\pgfpathlineto{\pgfqpoint{4.622786in}{2.303577in}}%
\pgfpathlineto{\pgfqpoint{4.623375in}{2.304398in}}%
\pgfpathlineto{\pgfqpoint{4.622874in}{2.302716in}}%
\pgfpathlineto{\pgfqpoint{4.623906in}{2.304072in}}%
\pgfpathlineto{\pgfqpoint{4.624743in}{2.302452in}}%
\pgfpathlineto{\pgfqpoint{4.624008in}{2.305316in}}%
\pgfpathlineto{\pgfqpoint{4.625011in}{2.304168in}}%
\pgfpathlineto{\pgfqpoint{4.626374in}{2.299959in}}%
\pgfpathlineto{\pgfqpoint{4.626379in}{2.299966in}}%
\pgfpathlineto{\pgfqpoint{4.626408in}{2.299456in}}%
\pgfpathlineto{\pgfqpoint{4.627168in}{2.302053in}}%
\pgfpathlineto{\pgfqpoint{4.627465in}{2.301095in}}%
\pgfpathlineto{\pgfqpoint{4.628545in}{2.304127in}}%
\pgfpathlineto{\pgfqpoint{4.627640in}{2.300766in}}%
\pgfpathlineto{\pgfqpoint{4.628652in}{2.303381in}}%
\pgfpathlineto{\pgfqpoint{4.628847in}{2.301348in}}%
\pgfpathlineto{\pgfqpoint{4.629329in}{2.304008in}}%
\pgfpathlineto{\pgfqpoint{4.629787in}{2.301742in}}%
\pgfpathlineto{\pgfqpoint{4.629821in}{2.301903in}}%
\pgfpathlineto{\pgfqpoint{4.630147in}{2.299177in}}%
\pgfpathlineto{\pgfqpoint{4.630425in}{2.300328in}}%
\pgfpathlineto{\pgfqpoint{4.631204in}{2.296902in}}%
\pgfpathlineto{\pgfqpoint{4.630595in}{2.300367in}}%
\pgfpathlineto{\pgfqpoint{4.631559in}{2.297723in}}%
\pgfpathlineto{\pgfqpoint{4.631822in}{2.297221in}}%
\pgfpathlineto{\pgfqpoint{4.632722in}{2.301669in}}%
\pgfpathlineto{\pgfqpoint{4.633214in}{2.304302in}}%
\pgfpathlineto{\pgfqpoint{4.632747in}{2.301519in}}%
\pgfpathlineto{\pgfqpoint{4.633852in}{2.302788in}}%
\pgfpathlineto{\pgfqpoint{4.634436in}{2.299892in}}%
\pgfpathlineto{\pgfqpoint{4.634967in}{2.302392in}}%
\pgfpathlineto{\pgfqpoint{4.635064in}{2.303777in}}%
\pgfpathlineto{\pgfqpoint{4.635775in}{2.301629in}}%
\pgfpathlineto{\pgfqpoint{4.636087in}{2.302916in}}%
\pgfpathlineto{\pgfqpoint{4.636150in}{2.301626in}}%
\pgfpathlineto{\pgfqpoint{4.636666in}{2.305197in}}%
\pgfpathlineto{\pgfqpoint{4.637172in}{2.303862in}}%
\pgfpathlineto{\pgfqpoint{4.638097in}{2.302340in}}%
\pgfpathlineto{\pgfqpoint{4.637523in}{2.305391in}}%
\pgfpathlineto{\pgfqpoint{4.638321in}{2.303150in}}%
\pgfpathlineto{\pgfqpoint{4.638331in}{2.303869in}}%
\pgfpathlineto{\pgfqpoint{4.639061in}{2.300870in}}%
\pgfpathlineto{\pgfqpoint{4.639426in}{2.302958in}}%
\pgfpathlineto{\pgfqpoint{4.640278in}{2.304148in}}%
\pgfpathlineto{\pgfqpoint{4.640030in}{2.300988in}}%
\pgfpathlineto{\pgfqpoint{4.640449in}{2.301003in}}%
\pgfpathlineto{\pgfqpoint{4.640624in}{2.299539in}}%
\pgfpathlineto{\pgfqpoint{4.641311in}{2.304003in}}%
\pgfpathlineto{\pgfqpoint{4.641534in}{2.302392in}}%
\pgfpathlineto{\pgfqpoint{4.641831in}{2.305220in}}%
\pgfpathlineto{\pgfqpoint{4.641544in}{2.302260in}}%
\pgfpathlineto{\pgfqpoint{4.642718in}{2.303458in}}%
\pgfpathlineto{\pgfqpoint{4.643720in}{2.300520in}}%
\pgfpathlineto{\pgfqpoint{4.642815in}{2.303843in}}%
\pgfpathlineto{\pgfqpoint{4.643852in}{2.301744in}}%
\pgfpathlineto{\pgfqpoint{4.644938in}{2.303387in}}%
\pgfpathlineto{\pgfqpoint{4.644417in}{2.300429in}}%
\pgfpathlineto{\pgfqpoint{4.644981in}{2.302800in}}%
\pgfpathlineto{\pgfqpoint{4.645293in}{2.301348in}}%
\pgfpathlineto{\pgfqpoint{4.645478in}{2.304149in}}%
\pgfpathlineto{\pgfqpoint{4.646018in}{2.303969in}}%
\pgfpathlineto{\pgfqpoint{4.646544in}{2.304515in}}%
\pgfpathlineto{\pgfqpoint{4.646909in}{2.302363in}}%
\pgfpathlineto{\pgfqpoint{4.646934in}{2.302844in}}%
\pgfpathlineto{\pgfqpoint{4.647941in}{2.300403in}}%
\pgfpathlineto{\pgfqpoint{4.648053in}{2.301909in}}%
\pgfpathlineto{\pgfqpoint{4.649266in}{2.304141in}}%
\pgfpathlineto{\pgfqpoint{4.648346in}{2.301590in}}%
\pgfpathlineto{\pgfqpoint{4.649295in}{2.303925in}}%
\pgfpathlineto{\pgfqpoint{4.649986in}{2.299835in}}%
\pgfpathlineto{\pgfqpoint{4.649675in}{2.304982in}}%
\pgfpathlineto{\pgfqpoint{4.650434in}{2.302910in}}%
\pgfpathlineto{\pgfqpoint{4.651184in}{2.304194in}}%
\pgfpathlineto{\pgfqpoint{4.650814in}{2.301092in}}%
\pgfpathlineto{\pgfqpoint{4.651422in}{2.302144in}}%
\pgfpathlineto{\pgfqpoint{4.652649in}{2.296173in}}%
\pgfpathlineto{\pgfqpoint{4.651515in}{2.303575in}}%
\pgfpathlineto{\pgfqpoint{4.652708in}{2.296904in}}%
\pgfpathlineto{\pgfqpoint{4.653394in}{2.296119in}}%
\pgfpathlineto{\pgfqpoint{4.652795in}{2.298844in}}%
\pgfpathlineto{\pgfqpoint{4.653652in}{2.298815in}}%
\pgfpathlineto{\pgfqpoint{4.654407in}{2.304166in}}%
\pgfpathlineto{\pgfqpoint{4.653745in}{2.297457in}}%
\pgfpathlineto{\pgfqpoint{4.654830in}{2.302365in}}%
\pgfpathlineto{\pgfqpoint{4.655200in}{2.301754in}}%
\pgfpathlineto{\pgfqpoint{4.655746in}{2.304318in}}%
\pgfpathlineto{\pgfqpoint{4.655838in}{2.303700in}}%
\pgfpathlineto{\pgfqpoint{4.655955in}{2.305225in}}%
\pgfpathlineto{\pgfqpoint{4.656729in}{2.299875in}}%
\pgfpathlineto{\pgfqpoint{4.656914in}{2.301376in}}%
\pgfpathlineto{\pgfqpoint{4.658170in}{2.298636in}}%
\pgfpathlineto{\pgfqpoint{4.656982in}{2.302899in}}%
\pgfpathlineto{\pgfqpoint{4.658219in}{2.299045in}}%
\pgfpathlineto{\pgfqpoint{4.658886in}{2.300315in}}%
\pgfpathlineto{\pgfqpoint{4.659173in}{2.297284in}}%
\pgfpathlineto{\pgfqpoint{4.659309in}{2.298711in}}%
\pgfpathlineto{\pgfqpoint{4.659319in}{2.298470in}}%
\pgfpathlineto{\pgfqpoint{4.659825in}{2.301478in}}%
\pgfpathlineto{\pgfqpoint{4.660249in}{2.300377in}}%
\pgfpathlineto{\pgfqpoint{4.660561in}{2.304015in}}%
\pgfpathlineto{\pgfqpoint{4.661413in}{2.301796in}}%
\pgfpathlineto{\pgfqpoint{4.661778in}{2.303165in}}%
\pgfpathlineto{\pgfqpoint{4.662119in}{2.300638in}}%
\pgfpathlineto{\pgfqpoint{4.662362in}{2.301360in}}%
\pgfpathlineto{\pgfqpoint{4.662401in}{2.300886in}}%
\pgfpathlineto{\pgfqpoint{4.663375in}{2.304660in}}%
\pgfpathlineto{\pgfqpoint{4.663394in}{2.304360in}}%
\pgfpathlineto{\pgfqpoint{4.664674in}{2.300237in}}%
\pgfpathlineto{\pgfqpoint{4.663545in}{2.304782in}}%
\pgfpathlineto{\pgfqpoint{4.664723in}{2.300374in}}%
\pgfpathlineto{\pgfqpoint{4.664782in}{2.300728in}}%
\pgfpathlineto{\pgfqpoint{4.665371in}{2.297923in}}%
\pgfpathlineto{\pgfqpoint{4.665833in}{2.300362in}}%
\pgfpathlineto{\pgfqpoint{4.666968in}{2.298058in}}%
\pgfpathlineto{\pgfqpoint{4.666072in}{2.300795in}}%
\pgfpathlineto{\pgfqpoint{4.666992in}{2.298439in}}%
\pgfpathlineto{\pgfqpoint{4.668126in}{2.304578in}}%
\pgfpathlineto{\pgfqpoint{4.668175in}{2.303854in}}%
\pgfpathlineto{\pgfqpoint{4.668934in}{2.300530in}}%
\pgfpathlineto{\pgfqpoint{4.669304in}{2.302263in}}%
\pgfpathlineto{\pgfqpoint{4.669392in}{2.301096in}}%
\pgfpathlineto{\pgfqpoint{4.670298in}{2.304853in}}%
\pgfpathlineto{\pgfqpoint{4.670356in}{2.304090in}}%
\pgfpathlineto{\pgfqpoint{4.670366in}{2.304364in}}%
\pgfpathlineto{\pgfqpoint{4.670828in}{2.299956in}}%
\pgfpathlineto{\pgfqpoint{4.671344in}{2.300985in}}%
\pgfpathlineto{\pgfqpoint{4.671524in}{2.299886in}}%
\pgfpathlineto{\pgfqpoint{4.672328in}{2.302754in}}%
\pgfpathlineto{\pgfqpoint{4.672430in}{2.302135in}}%
\pgfpathlineto{\pgfqpoint{4.672683in}{2.304905in}}%
\pgfpathlineto{\pgfqpoint{4.673131in}{2.301420in}}%
\pgfpathlineto{\pgfqpoint{4.673730in}{2.304210in}}%
\pgfpathlineto{\pgfqpoint{4.674879in}{2.301694in}}%
\pgfpathlineto{\pgfqpoint{4.674942in}{2.300843in}}%
\pgfpathlineto{\pgfqpoint{4.676077in}{2.296210in}}%
\pgfpathlineto{\pgfqpoint{4.676111in}{2.296691in}}%
\pgfpathlineto{\pgfqpoint{4.676481in}{2.297627in}}%
\pgfpathlineto{\pgfqpoint{4.677113in}{2.294739in}}%
\pgfpathlineto{\pgfqpoint{4.677128in}{2.294779in}}%
\pgfpathlineto{\pgfqpoint{4.677780in}{2.293447in}}%
\pgfpathlineto{\pgfqpoint{4.678160in}{2.297479in}}%
\pgfpathlineto{\pgfqpoint{4.678185in}{2.296975in}}%
\pgfpathlineto{\pgfqpoint{4.679509in}{2.304047in}}%
\pgfpathlineto{\pgfqpoint{4.679626in}{2.302845in}}%
\pgfpathlineto{\pgfqpoint{4.680565in}{2.302092in}}%
\pgfpathlineto{\pgfqpoint{4.680239in}{2.305135in}}%
\pgfpathlineto{\pgfqpoint{4.680721in}{2.303617in}}%
\pgfpathlineto{\pgfqpoint{4.680891in}{2.305179in}}%
\pgfpathlineto{\pgfqpoint{4.681339in}{2.301691in}}%
\pgfpathlineto{\pgfqpoint{4.681690in}{2.302463in}}%
\pgfpathlineto{\pgfqpoint{4.682099in}{2.299879in}}%
\pgfpathlineto{\pgfqpoint{4.681724in}{2.302582in}}%
\pgfpathlineto{\pgfqpoint{4.682819in}{2.301670in}}%
\pgfpathlineto{\pgfqpoint{4.683569in}{2.303784in}}%
\pgfpathlineto{\pgfqpoint{4.683287in}{2.300524in}}%
\pgfpathlineto{\pgfqpoint{4.683934in}{2.302261in}}%
\pgfpathlineto{\pgfqpoint{4.684903in}{2.299339in}}%
\pgfpathlineto{\pgfqpoint{4.684075in}{2.303779in}}%
\pgfpathlineto{\pgfqpoint{4.685049in}{2.301979in}}%
\pgfpathlineto{\pgfqpoint{4.685239in}{2.301361in}}%
\pgfpathlineto{\pgfqpoint{4.686091in}{2.303816in}}%
\pgfpathlineto{\pgfqpoint{4.686110in}{2.303722in}}%
\pgfpathlineto{\pgfqpoint{4.686130in}{2.304154in}}%
\pgfpathlineto{\pgfqpoint{4.687070in}{2.299078in}}%
\pgfpathlineto{\pgfqpoint{4.687138in}{2.299315in}}%
\pgfpathlineto{\pgfqpoint{4.687147in}{2.299092in}}%
\pgfpathlineto{\pgfqpoint{4.687800in}{2.303545in}}%
\pgfpathlineto{\pgfqpoint{4.688043in}{2.301357in}}%
\pgfpathlineto{\pgfqpoint{4.689046in}{2.302111in}}%
\pgfpathlineto{\pgfqpoint{4.688389in}{2.298896in}}%
\pgfpathlineto{\pgfqpoint{4.689124in}{2.300894in}}%
\pgfpathlineto{\pgfqpoint{4.689523in}{2.299886in}}%
\pgfpathlineto{\pgfqpoint{4.689849in}{2.302436in}}%
\pgfpathlineto{\pgfqpoint{4.690093in}{2.302061in}}%
\pgfpathlineto{\pgfqpoint{4.690687in}{2.303068in}}%
\pgfpathlineto{\pgfqpoint{4.690307in}{2.300440in}}%
\pgfpathlineto{\pgfqpoint{4.691183in}{2.300742in}}%
\pgfpathlineto{\pgfqpoint{4.692289in}{2.306438in}}%
\pgfpathlineto{\pgfqpoint{4.692483in}{2.305019in}}%
\pgfpathlineto{\pgfqpoint{4.693228in}{2.300691in}}%
\pgfpathlineto{\pgfqpoint{4.693681in}{2.303413in}}%
\pgfpathlineto{\pgfqpoint{4.693861in}{2.305233in}}%
\pgfpathlineto{\pgfqpoint{4.694075in}{2.302137in}}%
\pgfpathlineto{\pgfqpoint{4.694801in}{2.304014in}}%
\pgfpathlineto{\pgfqpoint{4.695818in}{2.300815in}}%
\pgfpathlineto{\pgfqpoint{4.695964in}{2.301867in}}%
\pgfpathlineto{\pgfqpoint{4.696612in}{2.305505in}}%
\pgfpathlineto{\pgfqpoint{4.696914in}{2.300859in}}%
\pgfpathlineto{\pgfqpoint{4.697099in}{2.303197in}}%
\pgfpathlineto{\pgfqpoint{4.697274in}{2.304760in}}%
\pgfpathlineto{\pgfqpoint{4.697590in}{2.302161in}}%
\pgfpathlineto{\pgfqpoint{4.698209in}{2.303324in}}%
\pgfpathlineto{\pgfqpoint{4.698598in}{2.300631in}}%
\pgfpathlineto{\pgfqpoint{4.698997in}{2.304595in}}%
\pgfpathlineto{\pgfqpoint{4.699270in}{2.303963in}}%
\pgfpathlineto{\pgfqpoint{4.700073in}{2.305081in}}%
\pgfpathlineto{\pgfqpoint{4.699479in}{2.302651in}}%
\pgfpathlineto{\pgfqpoint{4.700244in}{2.303067in}}%
\pgfpathlineto{\pgfqpoint{4.700867in}{2.301070in}}%
\pgfpathlineto{\pgfqpoint{4.701159in}{2.304022in}}%
\pgfpathlineto{\pgfqpoint{4.701339in}{2.303247in}}%
\pgfpathlineto{\pgfqpoint{4.702069in}{2.303712in}}%
\pgfpathlineto{\pgfqpoint{4.701544in}{2.301561in}}%
\pgfpathlineto{\pgfqpoint{4.702444in}{2.302956in}}%
\pgfpathlineto{\pgfqpoint{4.703335in}{2.301809in}}%
\pgfpathlineto{\pgfqpoint{4.702897in}{2.304604in}}%
\pgfpathlineto{\pgfqpoint{4.703564in}{2.302466in}}%
\pgfpathlineto{\pgfqpoint{4.703579in}{2.302796in}}%
\pgfpathlineto{\pgfqpoint{4.704207in}{2.300268in}}%
\pgfpathlineto{\pgfqpoint{4.704212in}{2.300276in}}%
\pgfpathlineto{\pgfqpoint{4.704689in}{2.298092in}}%
\pgfpathlineto{\pgfqpoint{4.704270in}{2.300767in}}%
\pgfpathlineto{\pgfqpoint{4.705356in}{2.299637in}}%
\pgfpathlineto{\pgfqpoint{4.705546in}{2.298793in}}%
\pgfpathlineto{\pgfqpoint{4.706519in}{2.301486in}}%
\pgfpathlineto{\pgfqpoint{4.707293in}{2.300717in}}%
\pgfpathlineto{\pgfqpoint{4.707576in}{2.303473in}}%
\pgfpathlineto{\pgfqpoint{4.707595in}{2.303274in}}%
\pgfpathlineto{\pgfqpoint{4.707795in}{2.304510in}}%
\pgfpathlineto{\pgfqpoint{4.708457in}{2.302040in}}%
\pgfpathlineto{\pgfqpoint{4.709246in}{2.299624in}}%
\pgfpathlineto{\pgfqpoint{4.708720in}{2.302270in}}%
\pgfpathlineto{\pgfqpoint{4.709577in}{2.301645in}}%
\pgfpathlineto{\pgfqpoint{4.710648in}{2.303558in}}%
\pgfpathlineto{\pgfqpoint{4.709703in}{2.300199in}}%
\pgfpathlineto{\pgfqpoint{4.710799in}{2.303275in}}%
\pgfpathlineto{\pgfqpoint{4.711164in}{2.300790in}}%
\pgfpathlineto{\pgfqpoint{4.711646in}{2.304623in}}%
\pgfpathlineto{\pgfqpoint{4.711904in}{2.303270in}}%
\pgfpathlineto{\pgfqpoint{4.713087in}{2.304898in}}%
\pgfpathlineto{\pgfqpoint{4.712605in}{2.302534in}}%
\pgfpathlineto{\pgfqpoint{4.713092in}{2.304750in}}%
\pgfpathlineto{\pgfqpoint{4.713803in}{2.301375in}}%
\pgfpathlineto{\pgfqpoint{4.713174in}{2.304940in}}%
\pgfpathlineto{\pgfqpoint{4.714285in}{2.302268in}}%
\pgfpathlineto{\pgfqpoint{4.714835in}{2.300514in}}%
\pgfpathlineto{\pgfqpoint{4.714426in}{2.302877in}}%
\pgfpathlineto{\pgfqpoint{4.715438in}{2.301702in}}%
\pgfpathlineto{\pgfqpoint{4.715901in}{2.304541in}}%
\pgfpathlineto{\pgfqpoint{4.715502in}{2.300895in}}%
\pgfpathlineto{\pgfqpoint{4.716548in}{2.301794in}}%
\pgfpathlineto{\pgfqpoint{4.716665in}{2.302253in}}%
\pgfpathlineto{\pgfqpoint{4.717814in}{2.296371in}}%
\pgfpathlineto{\pgfqpoint{4.718647in}{2.298883in}}%
\pgfpathlineto{\pgfqpoint{4.718369in}{2.295706in}}%
\pgfpathlineto{\pgfqpoint{4.718919in}{2.296375in}}%
\pgfpathlineto{\pgfqpoint{4.718924in}{2.296034in}}%
\pgfpathlineto{\pgfqpoint{4.719966in}{2.298771in}}%
\pgfpathlineto{\pgfqpoint{4.720000in}{2.298305in}}%
\pgfpathlineto{\pgfqpoint{4.720881in}{2.300537in}}%
\pgfpathlineto{\pgfqpoint{4.720161in}{2.297367in}}%
\pgfpathlineto{\pgfqpoint{4.721144in}{2.299930in}}%
\pgfpathlineto{\pgfqpoint{4.721169in}{2.299592in}}%
\pgfpathlineto{\pgfqpoint{4.721436in}{2.301567in}}%
\pgfpathlineto{\pgfqpoint{4.722210in}{2.300228in}}%
\pgfpathlineto{\pgfqpoint{4.723316in}{2.302909in}}%
\pgfpathlineto{\pgfqpoint{4.722405in}{2.298308in}}%
\pgfpathlineto{\pgfqpoint{4.723374in}{2.302787in}}%
\pgfpathlineto{\pgfqpoint{4.724109in}{2.300761in}}%
\pgfpathlineto{\pgfqpoint{4.723457in}{2.303343in}}%
\pgfpathlineto{\pgfqpoint{4.724489in}{2.302517in}}%
\pgfpathlineto{\pgfqpoint{4.725088in}{2.304257in}}%
\pgfpathlineto{\pgfqpoint{4.724835in}{2.301716in}}%
\pgfpathlineto{\pgfqpoint{4.725667in}{2.303694in}}%
\pgfpathlineto{\pgfqpoint{4.726743in}{2.298693in}}%
\pgfpathlineto{\pgfqpoint{4.725764in}{2.304157in}}%
\pgfpathlineto{\pgfqpoint{4.726850in}{2.299421in}}%
\pgfpathlineto{\pgfqpoint{4.727585in}{2.296625in}}%
\pgfpathlineto{\pgfqpoint{4.727181in}{2.299751in}}%
\pgfpathlineto{\pgfqpoint{4.728082in}{2.298409in}}%
\pgfpathlineto{\pgfqpoint{4.728096in}{2.298748in}}%
\pgfpathlineto{\pgfqpoint{4.729046in}{2.295237in}}%
\pgfpathlineto{\pgfqpoint{4.729138in}{2.295486in}}%
\pgfpathlineto{\pgfqpoint{4.730302in}{2.297020in}}%
\pgfpathlineto{\pgfqpoint{4.729640in}{2.292911in}}%
\pgfpathlineto{\pgfqpoint{4.730321in}{2.297013in}}%
\pgfpathlineto{\pgfqpoint{4.730351in}{2.296284in}}%
\pgfpathlineto{\pgfqpoint{4.730526in}{2.298677in}}%
\pgfpathlineto{\pgfqpoint{4.731431in}{2.296927in}}%
\pgfpathlineto{\pgfqpoint{4.731685in}{2.296051in}}%
\pgfpathlineto{\pgfqpoint{4.732575in}{2.299278in}}%
\pgfpathlineto{\pgfqpoint{4.733822in}{2.296478in}}%
\pgfpathlineto{\pgfqpoint{4.732911in}{2.300073in}}%
\pgfpathlineto{\pgfqpoint{4.733909in}{2.296994in}}%
\pgfpathlineto{\pgfqpoint{4.734766in}{2.299226in}}%
\pgfpathlineto{\pgfqpoint{4.734124in}{2.296655in}}%
\pgfpathlineto{\pgfqpoint{4.735073in}{2.298166in}}%
\pgfpathlineto{\pgfqpoint{4.735818in}{2.295105in}}%
\pgfpathlineto{\pgfqpoint{4.736115in}{2.298282in}}%
\pgfpathlineto{\pgfqpoint{4.736183in}{2.297846in}}%
\pgfpathlineto{\pgfqpoint{4.736324in}{2.298741in}}%
\pgfpathlineto{\pgfqpoint{4.736699in}{2.294900in}}%
\pgfpathlineto{\pgfqpoint{4.737249in}{2.296062in}}%
\pgfpathlineto{\pgfqpoint{4.738252in}{2.294126in}}%
\pgfpathlineto{\pgfqpoint{4.737293in}{2.296700in}}%
\pgfpathlineto{\pgfqpoint{4.738364in}{2.295782in}}%
\pgfpathlineto{\pgfqpoint{4.738384in}{2.295579in}}%
\pgfpathlineto{\pgfqpoint{4.738500in}{2.297121in}}%
\pgfpathlineto{\pgfqpoint{4.739075in}{2.300508in}}%
\pgfpathlineto{\pgfqpoint{4.738510in}{2.296979in}}%
\pgfpathlineto{\pgfqpoint{4.739664in}{2.298844in}}%
\pgfpathlineto{\pgfqpoint{4.739980in}{2.298556in}}%
\pgfpathlineto{\pgfqpoint{4.740686in}{2.301689in}}%
\pgfpathlineto{\pgfqpoint{4.740711in}{2.301546in}}%
\pgfpathlineto{\pgfqpoint{4.740910in}{2.304864in}}%
\pgfpathlineto{\pgfqpoint{4.741344in}{2.303507in}}%
\pgfpathlineto{\pgfqpoint{4.742454in}{2.306442in}}%
\pgfpathlineto{\pgfqpoint{4.741378in}{2.303020in}}%
\pgfpathlineto{\pgfqpoint{4.742585in}{2.305723in}}%
\pgfpathlineto{\pgfqpoint{4.744148in}{2.300590in}}%
\pgfpathlineto{\pgfqpoint{4.744455in}{2.301961in}}%
\pgfpathlineto{\pgfqpoint{4.744528in}{2.303002in}}%
\pgfpathlineto{\pgfqpoint{4.745287in}{2.297173in}}%
\pgfpathlineto{\pgfqpoint{4.745511in}{2.300324in}}%
\pgfpathlineto{\pgfqpoint{4.745516in}{2.300221in}}%
\pgfpathlineto{\pgfqpoint{4.746178in}{2.303427in}}%
\pgfpathlineto{\pgfqpoint{4.746499in}{2.302694in}}%
\pgfpathlineto{\pgfqpoint{4.746519in}{2.303503in}}%
\pgfpathlineto{\pgfqpoint{4.747458in}{2.299480in}}%
\pgfpathlineto{\pgfqpoint{4.747570in}{2.300719in}}%
\pgfpathlineto{\pgfqpoint{4.748583in}{2.296726in}}%
\pgfpathlineto{\pgfqpoint{4.747872in}{2.302307in}}%
\pgfpathlineto{\pgfqpoint{4.748724in}{2.297575in}}%
\pgfpathlineto{\pgfqpoint{4.748749in}{2.297997in}}%
\pgfpathlineto{\pgfqpoint{4.749625in}{2.293808in}}%
\pgfpathlineto{\pgfqpoint{4.749810in}{2.296385in}}%
\pgfpathlineto{\pgfqpoint{4.750209in}{2.295622in}}%
\pgfpathlineto{\pgfqpoint{4.750560in}{2.297998in}}%
\pgfpathlineto{\pgfqpoint{4.750628in}{2.297228in}}%
\pgfpathlineto{\pgfqpoint{4.750740in}{2.298870in}}%
\pgfpathlineto{\pgfqpoint{4.751714in}{2.295543in}}%
\pgfpathlineto{\pgfqpoint{4.751728in}{2.296005in}}%
\pgfpathlineto{\pgfqpoint{4.752273in}{2.297558in}}%
\pgfpathlineto{\pgfqpoint{4.752716in}{2.293760in}}%
\pgfpathlineto{\pgfqpoint{4.752799in}{2.293996in}}%
\pgfpathlineto{\pgfqpoint{4.753320in}{2.295834in}}%
\pgfpathlineto{\pgfqpoint{4.753846in}{2.292819in}}%
\pgfpathlineto{\pgfqpoint{4.753880in}{2.293017in}}%
\pgfpathlineto{\pgfqpoint{4.755292in}{2.287314in}}%
\pgfpathlineto{\pgfqpoint{4.754435in}{2.293584in}}%
\pgfpathlineto{\pgfqpoint{4.755394in}{2.287863in}}%
\pgfpathlineto{\pgfqpoint{4.755418in}{2.288511in}}%
\pgfpathlineto{\pgfqpoint{4.756197in}{2.284333in}}%
\pgfpathlineto{\pgfqpoint{4.756480in}{2.286570in}}%
\pgfpathlineto{\pgfqpoint{4.756567in}{2.285890in}}%
\pgfpathlineto{\pgfqpoint{4.757181in}{2.293360in}}%
\pgfpathlineto{\pgfqpoint{4.757191in}{2.293227in}}%
\pgfpathlineto{\pgfqpoint{4.757195in}{2.293415in}}%
\pgfpathlineto{\pgfqpoint{4.758043in}{2.288378in}}%
\pgfpathlineto{\pgfqpoint{4.758218in}{2.290315in}}%
\pgfpathlineto{\pgfqpoint{4.758963in}{2.287525in}}%
\pgfpathlineto{\pgfqpoint{4.759260in}{2.290494in}}%
\pgfpathlineto{\pgfqpoint{4.759343in}{2.289800in}}%
\pgfpathlineto{\pgfqpoint{4.760345in}{2.292239in}}%
\pgfpathlineto{\pgfqpoint{4.759834in}{2.288795in}}%
\pgfpathlineto{\pgfqpoint{4.760487in}{2.290659in}}%
\pgfpathlineto{\pgfqpoint{4.760564in}{2.290211in}}%
\pgfpathlineto{\pgfqpoint{4.761295in}{2.292626in}}%
\pgfpathlineto{\pgfqpoint{4.761514in}{2.291336in}}%
\pgfpathlineto{\pgfqpoint{4.761533in}{2.291882in}}%
\pgfpathlineto{\pgfqpoint{4.762400in}{2.286578in}}%
\pgfpathlineto{\pgfqpoint{4.762872in}{2.283933in}}%
\pgfpathlineto{\pgfqpoint{4.762522in}{2.287057in}}%
\pgfpathlineto{\pgfqpoint{4.763515in}{2.285982in}}%
\pgfpathlineto{\pgfqpoint{4.764571in}{2.293692in}}%
\pgfpathlineto{\pgfqpoint{4.764737in}{2.291770in}}%
\pgfpathlineto{\pgfqpoint{4.766027in}{2.287877in}}%
\pgfpathlineto{\pgfqpoint{4.765238in}{2.291923in}}%
\pgfpathlineto{\pgfqpoint{4.766076in}{2.288676in}}%
\pgfpathlineto{\pgfqpoint{4.766645in}{2.291252in}}%
\pgfpathlineto{\pgfqpoint{4.767059in}{2.287716in}}%
\pgfpathlineto{\pgfqpoint{4.767074in}{2.287753in}}%
\pgfpathlineto{\pgfqpoint{4.767161in}{2.287100in}}%
\pgfpathlineto{\pgfqpoint{4.767604in}{2.290327in}}%
\pgfpathlineto{\pgfqpoint{4.768135in}{2.289322in}}%
\pgfpathlineto{\pgfqpoint{4.769211in}{2.285421in}}%
\pgfpathlineto{\pgfqpoint{4.768232in}{2.289738in}}%
\pgfpathlineto{\pgfqpoint{4.769372in}{2.287514in}}%
\pgfpathlineto{\pgfqpoint{4.770350in}{2.291607in}}%
\pgfpathlineto{\pgfqpoint{4.769420in}{2.287052in}}%
\pgfpathlineto{\pgfqpoint{4.770535in}{2.290924in}}%
\pgfpathlineto{\pgfqpoint{4.770550in}{2.290818in}}%
\pgfpathlineto{\pgfqpoint{4.771110in}{2.294629in}}%
\pgfpathlineto{\pgfqpoint{4.771197in}{2.294123in}}%
\pgfpathlineto{\pgfqpoint{4.773067in}{2.303704in}}%
\pgfpathlineto{\pgfqpoint{4.771231in}{2.293770in}}%
\pgfpathlineto{\pgfqpoint{4.773383in}{2.301084in}}%
\pgfpathlineto{\pgfqpoint{4.773451in}{2.299418in}}%
\pgfpathlineto{\pgfqpoint{4.774347in}{2.302315in}}%
\pgfpathlineto{\pgfqpoint{4.774498in}{2.300733in}}%
\pgfpathlineto{\pgfqpoint{4.775277in}{2.304409in}}%
\pgfpathlineto{\pgfqpoint{4.775652in}{2.301514in}}%
\pgfpathlineto{\pgfqpoint{4.776299in}{2.299879in}}%
\pgfpathlineto{\pgfqpoint{4.775827in}{2.302976in}}%
\pgfpathlineto{\pgfqpoint{4.776816in}{2.300885in}}%
\pgfpathlineto{\pgfqpoint{4.777648in}{2.303152in}}%
\pgfpathlineto{\pgfqpoint{4.777298in}{2.299397in}}%
\pgfpathlineto{\pgfqpoint{4.777945in}{2.301730in}}%
\pgfpathlineto{\pgfqpoint{4.778135in}{2.299411in}}%
\pgfpathlineto{\pgfqpoint{4.778544in}{2.302429in}}%
\pgfpathlineto{\pgfqpoint{4.779016in}{2.301991in}}%
\pgfpathlineto{\pgfqpoint{4.779333in}{2.304188in}}%
\pgfpathlineto{\pgfqpoint{4.779625in}{2.300850in}}%
\pgfpathlineto{\pgfqpoint{4.780112in}{2.302008in}}%
\pgfpathlineto{\pgfqpoint{4.781178in}{2.297785in}}%
\pgfpathlineto{\pgfqpoint{4.780496in}{2.302312in}}%
\pgfpathlineto{\pgfqpoint{4.781256in}{2.298420in}}%
\pgfpathlineto{\pgfqpoint{4.781314in}{2.299439in}}%
\pgfpathlineto{\pgfqpoint{4.782005in}{2.293519in}}%
\pgfpathlineto{\pgfqpoint{4.782307in}{2.296361in}}%
\pgfpathlineto{\pgfqpoint{4.782794in}{2.295152in}}%
\pgfpathlineto{\pgfqpoint{4.783218in}{2.298109in}}%
\pgfpathlineto{\pgfqpoint{4.783349in}{2.297728in}}%
\pgfpathlineto{\pgfqpoint{4.784084in}{2.299030in}}%
\pgfpathlineto{\pgfqpoint{4.784425in}{2.296428in}}%
\pgfpathlineto{\pgfqpoint{4.785657in}{2.294286in}}%
\pgfpathlineto{\pgfqpoint{4.784839in}{2.299164in}}%
\pgfpathlineto{\pgfqpoint{4.785686in}{2.295056in}}%
\pgfpathlineto{\pgfqpoint{4.786416in}{2.296825in}}%
\pgfpathlineto{\pgfqpoint{4.786270in}{2.293696in}}%
\pgfpathlineto{\pgfqpoint{4.786981in}{2.296297in}}%
\pgfpathlineto{\pgfqpoint{4.787463in}{2.293481in}}%
\pgfpathlineto{\pgfqpoint{4.788057in}{2.296512in}}%
\pgfpathlineto{\pgfqpoint{4.788086in}{2.296257in}}%
\pgfpathlineto{\pgfqpoint{4.789094in}{2.300261in}}%
\pgfpathlineto{\pgfqpoint{4.788281in}{2.296051in}}%
\pgfpathlineto{\pgfqpoint{4.789235in}{2.298408in}}%
\pgfpathlineto{\pgfqpoint{4.789654in}{2.297661in}}%
\pgfpathlineto{\pgfqpoint{4.790121in}{2.299650in}}%
\pgfpathlineto{\pgfqpoint{4.790355in}{2.298065in}}%
\pgfpathlineto{\pgfqpoint{4.790413in}{2.298753in}}%
\pgfpathlineto{\pgfqpoint{4.791294in}{2.293280in}}%
\pgfpathlineto{\pgfqpoint{4.791377in}{2.293931in}}%
\pgfpathlineto{\pgfqpoint{4.791382in}{2.293899in}}%
\pgfpathlineto{\pgfqpoint{4.791572in}{2.295634in}}%
\pgfpathlineto{\pgfqpoint{4.792361in}{2.294838in}}%
\pgfpathlineto{\pgfqpoint{4.793544in}{2.298897in}}%
\pgfpathlineto{\pgfqpoint{4.792463in}{2.294438in}}%
\pgfpathlineto{\pgfqpoint{4.793607in}{2.298584in}}%
\pgfpathlineto{\pgfqpoint{4.794600in}{2.293447in}}%
\pgfpathlineto{\pgfqpoint{4.794824in}{2.295130in}}%
\pgfpathlineto{\pgfqpoint{4.796017in}{2.298108in}}%
\pgfpathlineto{\pgfqpoint{4.795292in}{2.293629in}}%
\pgfpathlineto{\pgfqpoint{4.796022in}{2.297957in}}%
\pgfpathlineto{\pgfqpoint{4.796995in}{2.296869in}}%
\pgfpathlineto{\pgfqpoint{4.796392in}{2.299108in}}%
\pgfpathlineto{\pgfqpoint{4.797112in}{2.298779in}}%
\pgfpathlineto{\pgfqpoint{4.797940in}{2.304467in}}%
\pgfpathlineto{\pgfqpoint{4.798329in}{2.302027in}}%
\pgfpathlineto{\pgfqpoint{4.798675in}{2.300612in}}%
\pgfpathlineto{\pgfqpoint{4.798899in}{2.302907in}}%
\pgfpathlineto{\pgfqpoint{4.798904in}{2.303184in}}%
\pgfpathlineto{\pgfqpoint{4.799254in}{2.300593in}}%
\pgfpathlineto{\pgfqpoint{4.800004in}{2.302614in}}%
\pgfpathlineto{\pgfqpoint{4.801075in}{2.304163in}}%
\pgfpathlineto{\pgfqpoint{4.800837in}{2.301905in}}%
\pgfpathlineto{\pgfqpoint{4.801124in}{2.302794in}}%
\pgfpathlineto{\pgfqpoint{4.801543in}{2.299353in}}%
\pgfpathlineto{\pgfqpoint{4.801139in}{2.302947in}}%
\pgfpathlineto{\pgfqpoint{4.802234in}{2.302806in}}%
\pgfpathlineto{\pgfqpoint{4.802682in}{2.304481in}}%
\pgfpathlineto{\pgfqpoint{4.803174in}{2.300695in}}%
\pgfpathlineto{\pgfqpoint{4.803188in}{2.300723in}}%
\pgfpathlineto{\pgfqpoint{4.803573in}{2.299454in}}%
\pgfpathlineto{\pgfqpoint{4.804138in}{2.303943in}}%
\pgfpathlineto{\pgfqpoint{4.804201in}{2.304536in}}%
\pgfpathlineto{\pgfqpoint{4.805155in}{2.301341in}}%
\pgfpathlineto{\pgfqpoint{4.805228in}{2.302124in}}%
\pgfpathlineto{\pgfqpoint{4.805554in}{2.300944in}}%
\pgfpathlineto{\pgfqpoint{4.805871in}{2.304109in}}%
\pgfpathlineto{\pgfqpoint{4.806289in}{2.303590in}}%
\pgfpathlineto{\pgfqpoint{4.807341in}{2.305167in}}%
\pgfpathlineto{\pgfqpoint{4.807000in}{2.301778in}}%
\pgfpathlineto{\pgfqpoint{4.807414in}{2.304662in}}%
\pgfpathlineto{\pgfqpoint{4.807696in}{2.306103in}}%
\pgfpathlineto{\pgfqpoint{4.808247in}{2.299684in}}%
\pgfpathlineto{\pgfqpoint{4.809425in}{2.294204in}}%
\pgfpathlineto{\pgfqpoint{4.809551in}{2.295216in}}%
\pgfpathlineto{\pgfqpoint{4.809873in}{2.297084in}}%
\pgfpathlineto{\pgfqpoint{4.810233in}{2.294179in}}%
\pgfpathlineto{\pgfqpoint{4.810676in}{2.296923in}}%
\pgfpathlineto{\pgfqpoint{4.811601in}{2.295005in}}%
\pgfpathlineto{\pgfqpoint{4.811027in}{2.298616in}}%
\pgfpathlineto{\pgfqpoint{4.811879in}{2.295830in}}%
\pgfpathlineto{\pgfqpoint{4.812989in}{2.303109in}}%
\pgfpathlineto{\pgfqpoint{4.813412in}{2.302301in}}%
\pgfpathlineto{\pgfqpoint{4.814512in}{2.297167in}}%
\pgfpathlineto{\pgfqpoint{4.813441in}{2.302455in}}%
\pgfpathlineto{\pgfqpoint{4.814590in}{2.298293in}}%
\pgfpathlineto{\pgfqpoint{4.815744in}{2.301560in}}%
\pgfpathlineto{\pgfqpoint{4.814702in}{2.297433in}}%
\pgfpathlineto{\pgfqpoint{4.815749in}{2.301278in}}%
\pgfpathlineto{\pgfqpoint{4.815832in}{2.300261in}}%
\pgfpathlineto{\pgfqpoint{4.816684in}{2.303462in}}%
\pgfpathlineto{\pgfqpoint{4.816810in}{2.302491in}}%
\pgfpathlineto{\pgfqpoint{4.817297in}{2.305731in}}%
\pgfpathlineto{\pgfqpoint{4.817998in}{2.304006in}}%
\pgfpathlineto{\pgfqpoint{4.818782in}{2.298749in}}%
\pgfpathlineto{\pgfqpoint{4.819293in}{2.301040in}}%
\pgfpathlineto{\pgfqpoint{4.820496in}{2.306861in}}%
\pgfpathlineto{\pgfqpoint{4.819361in}{2.300935in}}%
\pgfpathlineto{\pgfqpoint{4.820540in}{2.305779in}}%
\pgfpathlineto{\pgfqpoint{4.821810in}{2.299982in}}%
\pgfpathlineto{\pgfqpoint{4.821830in}{2.300042in}}%
\pgfpathlineto{\pgfqpoint{4.822930in}{2.304408in}}%
\pgfpathlineto{\pgfqpoint{4.821839in}{2.299561in}}%
\pgfpathlineto{\pgfqpoint{4.823023in}{2.303936in}}%
\pgfpathlineto{\pgfqpoint{4.824016in}{2.301956in}}%
\pgfpathlineto{\pgfqpoint{4.823188in}{2.304281in}}%
\pgfpathlineto{\pgfqpoint{4.824147in}{2.303633in}}%
\pgfpathlineto{\pgfqpoint{4.825277in}{2.304129in}}%
\pgfpathlineto{\pgfqpoint{4.824702in}{2.301690in}}%
\pgfpathlineto{\pgfqpoint{4.825281in}{2.304036in}}%
\pgfpathlineto{\pgfqpoint{4.826182in}{2.302229in}}%
\pgfpathlineto{\pgfqpoint{4.825783in}{2.304523in}}%
\pgfpathlineto{\pgfqpoint{4.826406in}{2.303148in}}%
\pgfpathlineto{\pgfqpoint{4.827044in}{2.304679in}}%
\pgfpathlineto{\pgfqpoint{4.826776in}{2.301950in}}%
\pgfpathlineto{\pgfqpoint{4.827497in}{2.302321in}}%
\pgfpathlineto{\pgfqpoint{4.828587in}{2.300917in}}%
\pgfpathlineto{\pgfqpoint{4.827842in}{2.304022in}}%
\pgfpathlineto{\pgfqpoint{4.828626in}{2.301354in}}%
\pgfpathlineto{\pgfqpoint{4.829132in}{2.304020in}}%
\pgfpathlineto{\pgfqpoint{4.829517in}{2.299286in}}%
\pgfpathlineto{\pgfqpoint{4.829712in}{2.300060in}}%
\pgfpathlineto{\pgfqpoint{4.829722in}{2.299656in}}%
\pgfpathlineto{\pgfqpoint{4.829921in}{2.303298in}}%
\pgfpathlineto{\pgfqpoint{4.830739in}{2.302768in}}%
\pgfpathlineto{\pgfqpoint{4.831289in}{2.304284in}}%
\pgfpathlineto{\pgfqpoint{4.830866in}{2.301907in}}%
\pgfpathlineto{\pgfqpoint{4.831834in}{2.302108in}}%
\pgfpathlineto{\pgfqpoint{4.831849in}{2.301804in}}%
\pgfpathlineto{\pgfqpoint{4.832798in}{2.304488in}}%
\pgfpathlineto{\pgfqpoint{4.832906in}{2.303584in}}%
\pgfpathlineto{\pgfqpoint{4.832979in}{2.305275in}}%
\pgfpathlineto{\pgfqpoint{4.833300in}{2.302188in}}%
\pgfpathlineto{\pgfqpoint{4.834006in}{2.302879in}}%
\pgfpathlineto{\pgfqpoint{4.834278in}{2.304370in}}%
\pgfpathlineto{\pgfqpoint{4.835028in}{2.302440in}}%
\pgfpathlineto{\pgfqpoint{4.835900in}{2.300695in}}%
\pgfpathlineto{\pgfqpoint{4.835301in}{2.304547in}}%
\pgfpathlineto{\pgfqpoint{4.836129in}{2.303251in}}%
\pgfpathlineto{\pgfqpoint{4.837039in}{2.302398in}}%
\pgfpathlineto{\pgfqpoint{4.836348in}{2.304087in}}%
\pgfpathlineto{\pgfqpoint{4.837224in}{2.303832in}}%
\pgfpathlineto{\pgfqpoint{4.837906in}{2.305009in}}%
\pgfpathlineto{\pgfqpoint{4.837613in}{2.302381in}}%
\pgfpathlineto{\pgfqpoint{4.838324in}{2.303346in}}%
\pgfpathlineto{\pgfqpoint{4.838582in}{2.302338in}}%
\pgfpathlineto{\pgfqpoint{4.838928in}{2.305565in}}%
\pgfpathlineto{\pgfqpoint{4.839420in}{2.303260in}}%
\pgfpathlineto{\pgfqpoint{4.840345in}{2.304388in}}%
\pgfpathlineto{\pgfqpoint{4.840014in}{2.300526in}}%
\pgfpathlineto{\pgfqpoint{4.840535in}{2.304182in}}%
\pgfpathlineto{\pgfqpoint{4.841265in}{2.302704in}}%
\pgfpathlineto{\pgfqpoint{4.841182in}{2.305057in}}%
\pgfpathlineto{\pgfqpoint{4.841659in}{2.303400in}}%
\pgfpathlineto{\pgfqpoint{4.842122in}{2.305637in}}%
\pgfpathlineto{\pgfqpoint{4.842677in}{2.302201in}}%
\pgfpathlineto{\pgfqpoint{4.842706in}{2.302784in}}%
\pgfpathlineto{\pgfqpoint{4.842959in}{2.301241in}}%
\pgfpathlineto{\pgfqpoint{4.843572in}{2.304575in}}%
\pgfpathlineto{\pgfqpoint{4.843806in}{2.303069in}}%
\pgfpathlineto{\pgfqpoint{4.844342in}{2.306523in}}%
\pgfpathlineto{\pgfqpoint{4.844853in}{2.302177in}}%
\pgfpathlineto{\pgfqpoint{4.844941in}{2.303960in}}%
\pgfpathlineto{\pgfqpoint{4.845797in}{2.301519in}}%
\pgfpathlineto{\pgfqpoint{4.845203in}{2.304193in}}%
\pgfpathlineto{\pgfqpoint{4.846046in}{2.304048in}}%
\pgfpathlineto{\pgfqpoint{4.846372in}{2.304838in}}%
\pgfpathlineto{\pgfqpoint{4.846620in}{2.302702in}}%
\pgfpathlineto{\pgfqpoint{4.847088in}{2.303651in}}%
\pgfpathlineto{\pgfqpoint{4.847389in}{2.301701in}}%
\pgfpathlineto{\pgfqpoint{4.847998in}{2.304412in}}%
\pgfpathlineto{\pgfqpoint{4.848251in}{2.301842in}}%
\pgfpathlineto{\pgfqpoint{4.848295in}{2.301503in}}%
\pgfpathlineto{\pgfqpoint{4.848529in}{2.304355in}}%
\pgfpathlineto{\pgfqpoint{4.849001in}{2.303439in}}%
\pgfpathlineto{\pgfqpoint{4.849006in}{2.303536in}}%
\pgfpathlineto{\pgfqpoint{4.849902in}{2.299967in}}%
\pgfpathlineto{\pgfqpoint{4.849950in}{2.300496in}}%
\pgfpathlineto{\pgfqpoint{4.850072in}{2.298975in}}%
\pgfpathlineto{\pgfqpoint{4.850924in}{2.301875in}}%
\pgfpathlineto{\pgfqpoint{4.851036in}{2.301768in}}%
\pgfpathlineto{\pgfqpoint{4.852190in}{2.304451in}}%
\pgfpathlineto{\pgfqpoint{4.852199in}{2.304335in}}%
\pgfpathlineto{\pgfqpoint{4.854054in}{2.299022in}}%
\pgfpathlineto{\pgfqpoint{4.854161in}{2.301535in}}%
\pgfpathlineto{\pgfqpoint{4.855276in}{2.305251in}}%
\pgfpathlineto{\pgfqpoint{4.855330in}{2.303803in}}%
\pgfpathlineto{\pgfqpoint{4.855374in}{2.303995in}}%
\pgfpathlineto{\pgfqpoint{4.855909in}{2.300924in}}%
\pgfpathlineto{\pgfqpoint{4.856274in}{2.301655in}}%
\pgfpathlineto{\pgfqpoint{4.857316in}{2.304580in}}%
\pgfpathlineto{\pgfqpoint{4.856566in}{2.300146in}}%
\pgfpathlineto{\pgfqpoint{4.857681in}{2.303217in}}%
\pgfpathlineto{\pgfqpoint{4.858407in}{2.302145in}}%
\pgfpathlineto{\pgfqpoint{4.858684in}{2.305643in}}%
\pgfpathlineto{\pgfqpoint{4.858772in}{2.304600in}}%
\pgfpathlineto{\pgfqpoint{4.858777in}{2.304952in}}%
\pgfpathlineto{\pgfqpoint{4.859824in}{2.300608in}}%
\pgfpathlineto{\pgfqpoint{4.859974in}{2.297675in}}%
\pgfpathlineto{\pgfqpoint{4.860822in}{2.303072in}}%
\pgfpathlineto{\pgfqpoint{4.860875in}{2.302431in}}%
\pgfpathlineto{\pgfqpoint{4.861065in}{2.303938in}}%
\pgfpathlineto{\pgfqpoint{4.861654in}{2.301253in}}%
\pgfpathlineto{\pgfqpoint{4.861990in}{2.303214in}}%
\pgfpathlineto{\pgfqpoint{4.862248in}{2.301351in}}%
\pgfpathlineto{\pgfqpoint{4.862964in}{2.305910in}}%
\pgfpathlineto{\pgfqpoint{4.863081in}{2.304069in}}%
\pgfpathlineto{\pgfqpoint{4.863679in}{2.299912in}}%
\pgfpathlineto{\pgfqpoint{4.864327in}{2.302585in}}%
\pgfpathlineto{\pgfqpoint{4.865135in}{2.301081in}}%
\pgfpathlineto{\pgfqpoint{4.864575in}{2.303794in}}%
\pgfpathlineto{\pgfqpoint{4.865544in}{2.302134in}}%
\pgfpathlineto{\pgfqpoint{4.866596in}{2.303322in}}%
\pgfpathlineto{\pgfqpoint{4.866104in}{2.301513in}}%
\pgfpathlineto{\pgfqpoint{4.866678in}{2.303227in}}%
\pgfpathlineto{\pgfqpoint{4.867136in}{2.305276in}}%
\pgfpathlineto{\pgfqpoint{4.867686in}{2.302740in}}%
\pgfpathlineto{\pgfqpoint{4.867822in}{2.302265in}}%
\pgfpathlineto{\pgfqpoint{4.868178in}{2.304831in}}%
\pgfpathlineto{\pgfqpoint{4.868791in}{2.302862in}}%
\pgfpathlineto{\pgfqpoint{4.868816in}{2.303625in}}%
\pgfpathlineto{\pgfqpoint{4.869726in}{2.298780in}}%
\pgfpathlineto{\pgfqpoint{4.869867in}{2.299782in}}%
\pgfpathlineto{\pgfqpoint{4.870427in}{2.303704in}}%
\pgfpathlineto{\pgfqpoint{4.869892in}{2.299430in}}%
\pgfpathlineto{\pgfqpoint{4.870987in}{2.299894in}}%
\pgfpathlineto{\pgfqpoint{4.871060in}{2.301018in}}%
\pgfpathlineto{\pgfqpoint{4.871581in}{2.297691in}}%
\pgfpathlineto{\pgfqpoint{4.871635in}{2.297909in}}%
\pgfpathlineto{\pgfqpoint{4.872949in}{2.293156in}}%
\pgfpathlineto{\pgfqpoint{4.871975in}{2.299144in}}%
\pgfpathlineto{\pgfqpoint{4.873042in}{2.293725in}}%
\pgfpathlineto{\pgfqpoint{4.873241in}{2.294677in}}%
\pgfpathlineto{\pgfqpoint{4.873889in}{2.291317in}}%
\pgfpathlineto{\pgfqpoint{4.874132in}{2.293476in}}%
\pgfpathlineto{\pgfqpoint{4.874152in}{2.294062in}}%
\pgfpathlineto{\pgfqpoint{4.874361in}{2.291722in}}%
\pgfpathlineto{\pgfqpoint{4.874590in}{2.291853in}}%
\pgfpathlineto{\pgfqpoint{4.874872in}{2.290101in}}%
\pgfpathlineto{\pgfqpoint{4.875636in}{2.292796in}}%
\pgfpathlineto{\pgfqpoint{4.875641in}{2.292862in}}%
\pgfpathlineto{\pgfqpoint{4.876342in}{2.288515in}}%
\pgfpathlineto{\pgfqpoint{4.876430in}{2.288533in}}%
\pgfpathlineto{\pgfqpoint{4.877194in}{2.285495in}}%
\pgfpathlineto{\pgfqpoint{4.876440in}{2.288777in}}%
\pgfpathlineto{\pgfqpoint{4.877681in}{2.287431in}}%
\pgfpathlineto{\pgfqpoint{4.878426in}{2.289368in}}%
\pgfpathlineto{\pgfqpoint{4.877973in}{2.285562in}}%
\pgfpathlineto{\pgfqpoint{4.878786in}{2.287424in}}%
\pgfpathlineto{\pgfqpoint{4.879672in}{2.285650in}}%
\pgfpathlineto{\pgfqpoint{4.879220in}{2.288536in}}%
\pgfpathlineto{\pgfqpoint{4.879916in}{2.287018in}}%
\pgfpathlineto{\pgfqpoint{4.880817in}{2.288165in}}%
\pgfpathlineto{\pgfqpoint{4.880349in}{2.285698in}}%
\pgfpathlineto{\pgfqpoint{4.880963in}{2.286692in}}%
\pgfpathlineto{\pgfqpoint{4.882073in}{2.284906in}}%
\pgfpathlineto{\pgfqpoint{4.881878in}{2.287393in}}%
\pgfpathlineto{\pgfqpoint{4.882121in}{2.284976in}}%
\pgfpathlineto{\pgfqpoint{4.882881in}{2.283396in}}%
\pgfpathlineto{\pgfqpoint{4.882423in}{2.286225in}}%
\pgfpathlineto{\pgfqpoint{4.883236in}{2.284814in}}%
\pgfpathlineto{\pgfqpoint{4.883270in}{2.285651in}}%
\pgfpathlineto{\pgfqpoint{4.883669in}{2.283011in}}%
\pgfpathlineto{\pgfqpoint{4.884293in}{2.283727in}}%
\pgfpathlineto{\pgfqpoint{4.884327in}{2.282857in}}%
\pgfpathlineto{\pgfqpoint{4.885354in}{2.285332in}}%
\pgfpathlineto{\pgfqpoint{4.885369in}{2.285180in}}%
\pgfpathlineto{\pgfqpoint{4.885398in}{2.285445in}}%
\pgfpathlineto{\pgfqpoint{4.886060in}{2.283088in}}%
\pgfpathlineto{\pgfqpoint{4.886376in}{2.283241in}}%
\pgfpathlineto{\pgfqpoint{4.887389in}{2.281430in}}%
\pgfpathlineto{\pgfqpoint{4.886795in}{2.285415in}}%
\pgfpathlineto{\pgfqpoint{4.887525in}{2.282283in}}%
\pgfpathlineto{\pgfqpoint{4.887701in}{2.281135in}}%
\pgfpathlineto{\pgfqpoint{4.888777in}{2.288353in}}%
\pgfpathlineto{\pgfqpoint{4.888781in}{2.288053in}}%
\pgfpathlineto{\pgfqpoint{4.889122in}{2.291336in}}%
\pgfpathlineto{\pgfqpoint{4.889857in}{2.290174in}}%
\pgfpathlineto{\pgfqpoint{4.891006in}{2.294593in}}%
\pgfpathlineto{\pgfqpoint{4.890145in}{2.290108in}}%
\pgfpathlineto{\pgfqpoint{4.891094in}{2.293931in}}%
\pgfpathlineto{\pgfqpoint{4.892004in}{2.289762in}}%
\pgfpathlineto{\pgfqpoint{4.891313in}{2.294419in}}%
\pgfpathlineto{\pgfqpoint{4.892408in}{2.291743in}}%
\pgfpathlineto{\pgfqpoint{4.892472in}{2.293332in}}%
\pgfpathlineto{\pgfqpoint{4.893207in}{2.286940in}}%
\pgfpathlineto{\pgfqpoint{4.893475in}{2.289064in}}%
\pgfpathlineto{\pgfqpoint{4.894468in}{2.285501in}}%
\pgfpathlineto{\pgfqpoint{4.894643in}{2.287617in}}%
\pgfpathlineto{\pgfqpoint{4.895787in}{2.290238in}}%
\pgfpathlineto{\pgfqpoint{4.895359in}{2.286609in}}%
\pgfpathlineto{\pgfqpoint{4.895802in}{2.289479in}}%
\pgfpathlineto{\pgfqpoint{4.896006in}{2.286943in}}%
\pgfpathlineto{\pgfqpoint{4.896785in}{2.290938in}}%
\pgfpathlineto{\pgfqpoint{4.896897in}{2.290316in}}%
\pgfpathlineto{\pgfqpoint{4.897934in}{2.297747in}}%
\pgfpathlineto{\pgfqpoint{4.898182in}{2.296637in}}%
\pgfpathlineto{\pgfqpoint{4.898981in}{2.293733in}}%
\pgfpathlineto{\pgfqpoint{4.898363in}{2.297535in}}%
\pgfpathlineto{\pgfqpoint{4.899312in}{2.295371in}}%
\pgfpathlineto{\pgfqpoint{4.899395in}{2.294018in}}%
\pgfpathlineto{\pgfqpoint{4.899950in}{2.296977in}}%
\pgfpathlineto{\pgfqpoint{4.900407in}{2.296005in}}%
\pgfpathlineto{\pgfqpoint{4.900592in}{2.295346in}}%
\pgfpathlineto{\pgfqpoint{4.901079in}{2.301158in}}%
\pgfpathlineto{\pgfqpoint{4.901211in}{2.300741in}}%
\pgfpathlineto{\pgfqpoint{4.901712in}{2.303370in}}%
\pgfpathlineto{\pgfqpoint{4.902136in}{2.300630in}}%
\pgfpathlineto{\pgfqpoint{4.902340in}{2.301336in}}%
\pgfpathlineto{\pgfqpoint{4.903017in}{2.302736in}}%
\pgfpathlineto{\pgfqpoint{4.903363in}{2.300548in}}%
\pgfpathlineto{\pgfqpoint{4.903426in}{2.299633in}}%
\pgfpathlineto{\pgfqpoint{4.904244in}{2.305217in}}%
\pgfpathlineto{\pgfqpoint{4.904463in}{2.301254in}}%
\pgfpathlineto{\pgfqpoint{4.904585in}{2.298777in}}%
\pgfpathlineto{\pgfqpoint{4.905485in}{2.301840in}}%
\pgfpathlineto{\pgfqpoint{4.906342in}{2.304642in}}%
\pgfpathlineto{\pgfqpoint{4.906615in}{2.302267in}}%
\pgfpathlineto{\pgfqpoint{4.908251in}{2.293207in}}%
\pgfpathlineto{\pgfqpoint{4.906980in}{2.302591in}}%
\pgfpathlineto{\pgfqpoint{4.908582in}{2.296099in}}%
\pgfpathlineto{\pgfqpoint{4.909560in}{2.296969in}}%
\pgfpathlineto{\pgfqpoint{4.909278in}{2.294468in}}%
\pgfpathlineto{\pgfqpoint{4.909701in}{2.296796in}}%
\pgfpathlineto{\pgfqpoint{4.910178in}{2.295517in}}%
\pgfpathlineto{\pgfqpoint{4.909974in}{2.298343in}}%
\pgfpathlineto{\pgfqpoint{4.910699in}{2.297805in}}%
\pgfpathlineto{\pgfqpoint{4.911313in}{2.302036in}}%
\pgfpathlineto{\pgfqpoint{4.910753in}{2.297359in}}%
\pgfpathlineto{\pgfqpoint{4.911887in}{2.300039in}}%
\pgfpathlineto{\pgfqpoint{4.911921in}{2.299804in}}%
\pgfpathlineto{\pgfqpoint{4.912842in}{2.304055in}}%
\pgfpathlineto{\pgfqpoint{4.912949in}{2.301949in}}%
\pgfpathlineto{\pgfqpoint{4.913367in}{2.304485in}}%
\pgfpathlineto{\pgfqpoint{4.914078in}{2.302562in}}%
\pgfpathlineto{\pgfqpoint{4.914258in}{2.301863in}}%
\pgfpathlineto{\pgfqpoint{4.914526in}{2.304873in}}%
\pgfpathlineto{\pgfqpoint{4.914901in}{2.304192in}}%
\pgfpathlineto{\pgfqpoint{4.915806in}{2.304684in}}%
\pgfpathlineto{\pgfqpoint{4.915952in}{2.302512in}}%
\pgfpathlineto{\pgfqpoint{4.915967in}{2.302691in}}%
\pgfpathlineto{\pgfqpoint{4.916697in}{2.299310in}}%
\pgfpathlineto{\pgfqpoint{4.917111in}{2.301866in}}%
\pgfpathlineto{\pgfqpoint{4.917199in}{2.303151in}}%
\pgfpathlineto{\pgfqpoint{4.917983in}{2.300088in}}%
\pgfpathlineto{\pgfqpoint{4.918197in}{2.300981in}}%
\pgfpathlineto{\pgfqpoint{4.918372in}{2.300255in}}%
\pgfpathlineto{\pgfqpoint{4.918766in}{2.303321in}}%
\pgfpathlineto{\pgfqpoint{4.919312in}{2.300735in}}%
\pgfpathlineto{\pgfqpoint{4.920242in}{2.306756in}}%
\pgfpathlineto{\pgfqpoint{4.919331in}{2.300588in}}%
\pgfpathlineto{\pgfqpoint{4.920582in}{2.304466in}}%
\pgfpathlineto{\pgfqpoint{4.921688in}{2.297968in}}%
\pgfpathlineto{\pgfqpoint{4.920597in}{2.304622in}}%
\pgfpathlineto{\pgfqpoint{4.921795in}{2.298483in}}%
\pgfpathlineto{\pgfqpoint{4.923538in}{2.304883in}}%
\pgfpathlineto{\pgfqpoint{4.922101in}{2.296805in}}%
\pgfpathlineto{\pgfqpoint{4.923581in}{2.304295in}}%
\pgfpathlineto{\pgfqpoint{4.924463in}{2.300907in}}%
\pgfpathlineto{\pgfqpoint{4.923825in}{2.304429in}}%
\pgfpathlineto{\pgfqpoint{4.924730in}{2.302250in}}%
\pgfpathlineto{\pgfqpoint{4.925276in}{2.304564in}}%
\pgfpathlineto{\pgfqpoint{4.925738in}{2.301645in}}%
\pgfpathlineto{\pgfqpoint{4.925836in}{2.302158in}}%
\pgfpathlineto{\pgfqpoint{4.925904in}{2.301478in}}%
\pgfpathlineto{\pgfqpoint{4.926522in}{2.304466in}}%
\pgfpathlineto{\pgfqpoint{4.926941in}{2.302289in}}%
\pgfpathlineto{\pgfqpoint{4.927277in}{2.304181in}}%
\pgfpathlineto{\pgfqpoint{4.927715in}{2.301509in}}%
\pgfpathlineto{\pgfqpoint{4.928051in}{2.302393in}}%
\pgfpathlineto{\pgfqpoint{4.928659in}{2.299859in}}%
\pgfpathlineto{\pgfqpoint{4.928177in}{2.302820in}}%
\pgfpathlineto{\pgfqpoint{4.929151in}{2.302653in}}%
\pgfpathlineto{\pgfqpoint{4.929735in}{2.302940in}}%
\pgfpathlineto{\pgfqpoint{4.930144in}{2.300319in}}%
\pgfpathlineto{\pgfqpoint{4.930198in}{2.300956in}}%
\pgfpathlineto{\pgfqpoint{4.931761in}{2.295762in}}%
\pgfpathlineto{\pgfqpoint{4.930592in}{2.303225in}}%
\pgfpathlineto{\pgfqpoint{4.931819in}{2.295882in}}%
\pgfpathlineto{\pgfqpoint{4.932934in}{2.297344in}}%
\pgfpathlineto{\pgfqpoint{4.932354in}{2.294379in}}%
\pgfpathlineto{\pgfqpoint{4.932968in}{2.296824in}}%
\pgfpathlineto{\pgfqpoint{4.934039in}{2.291363in}}%
\pgfpathlineto{\pgfqpoint{4.934141in}{2.292405in}}%
\pgfpathlineto{\pgfqpoint{4.934175in}{2.291806in}}%
\pgfpathlineto{\pgfqpoint{4.934472in}{2.295965in}}%
\pgfpathlineto{\pgfqpoint{4.935095in}{2.295124in}}%
\pgfpathlineto{\pgfqpoint{4.935237in}{2.295070in}}%
\pgfpathlineto{\pgfqpoint{4.935650in}{2.298394in}}%
\pgfpathlineto{\pgfqpoint{4.935772in}{2.299284in}}%
\pgfpathlineto{\pgfqpoint{4.936605in}{2.295363in}}%
\pgfpathlineto{\pgfqpoint{4.936663in}{2.295977in}}%
\pgfpathlineto{\pgfqpoint{4.937583in}{2.292671in}}%
\pgfpathlineto{\pgfqpoint{4.937788in}{2.295018in}}%
\pgfpathlineto{\pgfqpoint{4.938474in}{2.297683in}}%
\pgfpathlineto{\pgfqpoint{4.937856in}{2.294868in}}%
\pgfpathlineto{\pgfqpoint{4.938951in}{2.297565in}}%
\pgfpathlineto{\pgfqpoint{4.939545in}{2.295903in}}%
\pgfpathlineto{\pgfqpoint{4.940008in}{2.299661in}}%
\pgfpathlineto{\pgfqpoint{4.940017in}{2.299427in}}%
\pgfpathlineto{\pgfqpoint{4.941220in}{2.302413in}}%
\pgfpathlineto{\pgfqpoint{4.940100in}{2.299015in}}%
\pgfpathlineto{\pgfqpoint{4.941225in}{2.302385in}}%
\pgfpathlineto{\pgfqpoint{4.941887in}{2.299485in}}%
\pgfpathlineto{\pgfqpoint{4.942306in}{2.302614in}}%
\pgfpathlineto{\pgfqpoint{4.942345in}{2.301722in}}%
\pgfpathlineto{\pgfqpoint{4.943406in}{2.299110in}}%
\pgfpathlineto{\pgfqpoint{4.943065in}{2.302887in}}%
\pgfpathlineto{\pgfqpoint{4.943494in}{2.300895in}}%
\pgfpathlineto{\pgfqpoint{4.944307in}{2.303873in}}%
\pgfpathlineto{\pgfqpoint{4.944662in}{2.303544in}}%
\pgfpathlineto{\pgfqpoint{4.944672in}{2.303790in}}%
\pgfpathlineto{\pgfqpoint{4.945607in}{2.299889in}}%
\pgfpathlineto{\pgfqpoint{4.946765in}{2.296041in}}%
\pgfpathlineto{\pgfqpoint{4.945884in}{2.301813in}}%
\pgfpathlineto{\pgfqpoint{4.946775in}{2.296239in}}%
\pgfpathlineto{\pgfqpoint{4.947359in}{2.299149in}}%
\pgfpathlineto{\pgfqpoint{4.947856in}{2.295273in}}%
\pgfpathlineto{\pgfqpoint{4.947987in}{2.294956in}}%
\pgfpathlineto{\pgfqpoint{4.948421in}{2.299109in}}%
\pgfpathlineto{\pgfqpoint{4.948883in}{2.298166in}}%
\pgfpathlineto{\pgfqpoint{4.949039in}{2.296862in}}%
\pgfpathlineto{\pgfqpoint{4.949384in}{2.300311in}}%
\pgfpathlineto{\pgfqpoint{4.949696in}{2.300044in}}%
\pgfpathlineto{\pgfqpoint{4.949740in}{2.300865in}}%
\pgfpathlineto{\pgfqpoint{4.950324in}{2.296857in}}%
\pgfpathlineto{\pgfqpoint{4.950777in}{2.299531in}}%
\pgfpathlineto{\pgfqpoint{4.951916in}{2.295273in}}%
\pgfpathlineto{\pgfqpoint{4.951283in}{2.300175in}}%
\pgfpathlineto{\pgfqpoint{4.952004in}{2.296358in}}%
\pgfpathlineto{\pgfqpoint{4.952778in}{2.300204in}}%
\pgfpathlineto{\pgfqpoint{4.952013in}{2.296337in}}%
\pgfpathlineto{\pgfqpoint{4.953182in}{2.299485in}}%
\pgfpathlineto{\pgfqpoint{4.953416in}{2.298531in}}%
\pgfpathlineto{\pgfqpoint{4.954190in}{2.305752in}}%
\pgfpathlineto{\pgfqpoint{4.954229in}{2.304286in}}%
\pgfpathlineto{\pgfqpoint{4.954238in}{2.305036in}}%
\pgfpathlineto{\pgfqpoint{4.954978in}{2.302087in}}%
\pgfpathlineto{\pgfqpoint{4.955324in}{2.303412in}}%
\pgfpathlineto{\pgfqpoint{4.955777in}{2.303595in}}%
\pgfpathlineto{\pgfqpoint{4.956249in}{2.300255in}}%
\pgfpathlineto{\pgfqpoint{4.956346in}{2.301389in}}%
\pgfpathlineto{\pgfqpoint{4.957018in}{2.302899in}}%
\pgfpathlineto{\pgfqpoint{4.956653in}{2.299970in}}%
\pgfpathlineto{\pgfqpoint{4.957383in}{2.300790in}}%
\pgfpathlineto{\pgfqpoint{4.957598in}{2.298979in}}%
\pgfpathlineto{\pgfqpoint{4.958391in}{2.304474in}}%
\pgfpathlineto{\pgfqpoint{4.958430in}{2.304408in}}%
\pgfpathlineto{\pgfqpoint{4.958435in}{2.304417in}}%
\pgfpathlineto{\pgfqpoint{4.958659in}{2.302620in}}%
\pgfpathlineto{\pgfqpoint{4.958669in}{2.302501in}}%
\pgfpathlineto{\pgfqpoint{4.959423in}{2.307073in}}%
\pgfpathlineto{\pgfqpoint{4.959681in}{2.304632in}}%
\pgfpathlineto{\pgfqpoint{4.960743in}{2.305394in}}%
\pgfpathlineto{\pgfqpoint{4.959920in}{2.302456in}}%
\pgfpathlineto{\pgfqpoint{4.960791in}{2.304656in}}%
\pgfpathlineto{\pgfqpoint{4.961750in}{2.300217in}}%
\pgfpathlineto{\pgfqpoint{4.960825in}{2.304793in}}%
\pgfpathlineto{\pgfqpoint{4.961989in}{2.300998in}}%
\pgfpathlineto{\pgfqpoint{4.962398in}{2.303583in}}%
\pgfpathlineto{\pgfqpoint{4.962062in}{2.299866in}}%
\pgfpathlineto{\pgfqpoint{4.963157in}{2.302516in}}%
\pgfpathlineto{\pgfqpoint{4.964248in}{2.301252in}}%
\pgfpathlineto{\pgfqpoint{4.963527in}{2.305432in}}%
\pgfpathlineto{\pgfqpoint{4.964292in}{2.301796in}}%
\pgfpathlineto{\pgfqpoint{4.964331in}{2.300805in}}%
\pgfpathlineto{\pgfqpoint{4.964959in}{2.304339in}}%
\pgfpathlineto{\pgfqpoint{4.965378in}{2.302144in}}%
\pgfpathlineto{\pgfqpoint{4.965392in}{2.302291in}}%
\pgfpathlineto{\pgfqpoint{4.966327in}{2.297952in}}%
\pgfpathlineto{\pgfqpoint{4.967544in}{2.300324in}}%
\pgfpathlineto{\pgfqpoint{4.966833in}{2.296634in}}%
\pgfpathlineto{\pgfqpoint{4.967549in}{2.300171in}}%
\pgfpathlineto{\pgfqpoint{4.967559in}{2.299697in}}%
\pgfpathlineto{\pgfqpoint{4.968299in}{2.303520in}}%
\pgfpathlineto{\pgfqpoint{4.968576in}{2.302816in}}%
\pgfpathlineto{\pgfqpoint{4.969345in}{2.304878in}}%
\pgfpathlineto{\pgfqpoint{4.969618in}{2.302085in}}%
\pgfpathlineto{\pgfqpoint{4.969672in}{2.302345in}}%
\pgfpathlineto{\pgfqpoint{4.970830in}{2.299900in}}%
\pgfpathlineto{\pgfqpoint{4.969993in}{2.303430in}}%
\pgfpathlineto{\pgfqpoint{4.970840in}{2.300038in}}%
\pgfpathlineto{\pgfqpoint{4.971989in}{2.303203in}}%
\pgfpathlineto{\pgfqpoint{4.971590in}{2.299865in}}%
\pgfpathlineto{\pgfqpoint{4.971994in}{2.303152in}}%
\pgfpathlineto{\pgfqpoint{4.972111in}{2.301698in}}%
\pgfpathlineto{\pgfqpoint{4.972768in}{2.305712in}}%
\pgfpathlineto{\pgfqpoint{4.973084in}{2.303809in}}%
\pgfpathlineto{\pgfqpoint{4.973240in}{2.304124in}}%
\pgfpathlineto{\pgfqpoint{4.974053in}{2.299876in}}%
\pgfpathlineto{\pgfqpoint{4.974121in}{2.300442in}}%
\pgfpathlineto{\pgfqpoint{4.974126in}{2.300543in}}%
\pgfpathlineto{\pgfqpoint{4.974428in}{2.297776in}}%
\pgfpathlineto{\pgfqpoint{4.975153in}{2.298812in}}%
\pgfpathlineto{\pgfqpoint{4.975158in}{2.298748in}}%
\pgfpathlineto{\pgfqpoint{4.975762in}{2.303989in}}%
\pgfpathlineto{\pgfqpoint{4.975918in}{2.303202in}}%
\pgfpathlineto{\pgfqpoint{4.975923in}{2.303380in}}%
\pgfpathlineto{\pgfqpoint{4.976872in}{2.299214in}}%
\pgfpathlineto{\pgfqpoint{4.976945in}{2.300449in}}%
\pgfpathlineto{\pgfqpoint{4.977252in}{2.298506in}}%
\pgfpathlineto{\pgfqpoint{4.978011in}{2.302342in}}%
\pgfpathlineto{\pgfqpoint{4.979180in}{2.306436in}}%
\pgfpathlineto{\pgfqpoint{4.978036in}{2.301998in}}%
\pgfpathlineto{\pgfqpoint{4.979301in}{2.304692in}}%
\pgfpathlineto{\pgfqpoint{4.980348in}{2.300218in}}%
\pgfpathlineto{\pgfqpoint{4.980470in}{2.300772in}}%
\pgfpathlineto{\pgfqpoint{4.981716in}{2.295487in}}%
\pgfpathlineto{\pgfqpoint{4.980558in}{2.301004in}}%
\pgfpathlineto{\pgfqpoint{4.982490in}{2.297489in}}%
\pgfpathlineto{\pgfqpoint{4.983469in}{2.302516in}}%
\pgfpathlineto{\pgfqpoint{4.982534in}{2.297219in}}%
\pgfpathlineto{\pgfqpoint{4.983790in}{2.302367in}}%
\pgfpathlineto{\pgfqpoint{4.984642in}{2.301341in}}%
\pgfpathlineto{\pgfqpoint{4.983970in}{2.303540in}}%
\pgfpathlineto{\pgfqpoint{4.984890in}{2.302558in}}%
\pgfpathlineto{\pgfqpoint{4.985402in}{2.304980in}}%
\pgfpathlineto{\pgfqpoint{4.985976in}{2.302339in}}%
\pgfpathlineto{\pgfqpoint{4.986005in}{2.303078in}}%
\pgfpathlineto{\pgfqpoint{4.986663in}{2.300632in}}%
\pgfpathlineto{\pgfqpoint{4.987062in}{2.304684in}}%
\pgfpathlineto{\pgfqpoint{4.987091in}{2.304265in}}%
\pgfpathlineto{\pgfqpoint{4.988362in}{2.297075in}}%
\pgfpathlineto{\pgfqpoint{4.988488in}{2.297765in}}%
\pgfpathlineto{\pgfqpoint{4.988693in}{2.299695in}}%
\pgfpathlineto{\pgfqpoint{4.989306in}{2.296497in}}%
\pgfpathlineto{\pgfqpoint{4.989613in}{2.298117in}}%
\pgfpathlineto{\pgfqpoint{4.989618in}{2.298106in}}%
\pgfpathlineto{\pgfqpoint{4.990036in}{2.299954in}}%
\pgfpathlineto{\pgfqpoint{4.991356in}{2.303903in}}%
\pgfpathlineto{\pgfqpoint{4.990217in}{2.299315in}}%
\pgfpathlineto{\pgfqpoint{4.991361in}{2.303566in}}%
\pgfpathlineto{\pgfqpoint{4.992442in}{2.299434in}}%
\pgfpathlineto{\pgfqpoint{4.991594in}{2.304110in}}%
\pgfpathlineto{\pgfqpoint{4.992812in}{2.301665in}}%
\pgfpathlineto{\pgfqpoint{4.993532in}{2.304026in}}%
\pgfpathlineto{\pgfqpoint{4.993756in}{2.299621in}}%
\pgfpathlineto{\pgfqpoint{4.993926in}{2.301862in}}%
\pgfpathlineto{\pgfqpoint{4.994068in}{2.300298in}}%
\pgfpathlineto{\pgfqpoint{4.994788in}{2.302970in}}%
\pgfpathlineto{\pgfqpoint{4.995036in}{2.301571in}}%
\pgfpathlineto{\pgfqpoint{4.995859in}{2.305565in}}%
\pgfpathlineto{\pgfqpoint{4.995192in}{2.301105in}}%
\pgfpathlineto{\pgfqpoint{4.996234in}{2.304958in}}%
\pgfpathlineto{\pgfqpoint{4.997320in}{2.303134in}}%
\pgfpathlineto{\pgfqpoint{4.996478in}{2.305894in}}%
\pgfpathlineto{\pgfqpoint{4.997368in}{2.303613in}}%
\pgfpathlineto{\pgfqpoint{4.998279in}{2.301892in}}%
\pgfpathlineto{\pgfqpoint{4.997398in}{2.303931in}}%
\pgfpathlineto{\pgfqpoint{4.998605in}{2.302663in}}%
\pgfpathlineto{\pgfqpoint{4.998625in}{2.302965in}}%
\pgfpathlineto{\pgfqpoint{4.999403in}{2.297849in}}%
\pgfpathlineto{\pgfqpoint{4.999530in}{2.299403in}}%
\pgfpathlineto{\pgfqpoint{4.999608in}{2.298884in}}%
\pgfpathlineto{\pgfqpoint{5.000100in}{2.301924in}}%
\pgfpathlineto{\pgfqpoint{5.000621in}{2.299759in}}%
\pgfpathlineto{\pgfqpoint{5.000796in}{2.301261in}}%
\pgfpathlineto{\pgfqpoint{5.001380in}{2.297628in}}%
\pgfpathlineto{\pgfqpoint{5.001711in}{2.298940in}}%
\pgfpathlineto{\pgfqpoint{5.002315in}{2.294124in}}%
\pgfpathlineto{\pgfqpoint{5.002807in}{2.298984in}}%
\pgfpathlineto{\pgfqpoint{5.002821in}{2.298593in}}%
\pgfpathlineto{\pgfqpoint{5.003848in}{2.303955in}}%
\pgfpathlineto{\pgfqpoint{5.002992in}{2.297069in}}%
\pgfpathlineto{\pgfqpoint{5.004705in}{2.301213in}}%
\pgfpathlineto{\pgfqpoint{5.004881in}{2.299739in}}%
\pgfpathlineto{\pgfqpoint{5.005762in}{2.303176in}}%
\pgfpathlineto{\pgfqpoint{5.005776in}{2.303147in}}%
\pgfpathlineto{\pgfqpoint{5.006453in}{2.304552in}}%
\pgfpathlineto{\pgfqpoint{5.006711in}{2.302447in}}%
\pgfpathlineto{\pgfqpoint{5.006886in}{2.303376in}}%
\pgfpathlineto{\pgfqpoint{5.007110in}{2.302759in}}%
\pgfpathlineto{\pgfqpoint{5.007446in}{2.305868in}}%
\pgfpathlineto{\pgfqpoint{5.008001in}{2.303191in}}%
\pgfpathlineto{\pgfqpoint{5.008483in}{2.303875in}}%
\pgfpathlineto{\pgfqpoint{5.008634in}{2.301816in}}%
\pgfpathlineto{\pgfqpoint{5.008994in}{2.302006in}}%
\pgfpathlineto{\pgfqpoint{5.010216in}{2.297507in}}%
\pgfpathlineto{\pgfqpoint{5.009165in}{2.303166in}}%
\pgfpathlineto{\pgfqpoint{5.010265in}{2.298062in}}%
\pgfpathlineto{\pgfqpoint{5.011404in}{2.301850in}}%
\pgfpathlineto{\pgfqpoint{5.011029in}{2.297772in}}%
\pgfpathlineto{\pgfqpoint{5.011419in}{2.301535in}}%
\pgfpathlineto{\pgfqpoint{5.011570in}{2.300868in}}%
\pgfpathlineto{\pgfqpoint{5.012402in}{2.304352in}}%
\pgfpathlineto{\pgfqpoint{5.012446in}{2.303435in}}%
\pgfpathlineto{\pgfqpoint{5.013542in}{2.304564in}}%
\pgfpathlineto{\pgfqpoint{5.012704in}{2.301882in}}%
\pgfpathlineto{\pgfqpoint{5.013576in}{2.304382in}}%
\pgfpathlineto{\pgfqpoint{5.014895in}{2.301305in}}%
\pgfpathlineto{\pgfqpoint{5.014330in}{2.305260in}}%
\pgfpathlineto{\pgfqpoint{5.014958in}{2.302193in}}%
\pgfpathlineto{\pgfqpoint{5.015908in}{2.303759in}}%
\pgfpathlineto{\pgfqpoint{5.015168in}{2.300595in}}%
\pgfpathlineto{\pgfqpoint{5.016068in}{2.302099in}}%
\pgfpathlineto{\pgfqpoint{5.016390in}{2.303897in}}%
\pgfpathlineto{\pgfqpoint{5.016842in}{2.301069in}}%
\pgfpathlineto{\pgfqpoint{5.017159in}{2.301478in}}%
\pgfpathlineto{\pgfqpoint{5.017164in}{2.301536in}}%
\pgfpathlineto{\pgfqpoint{5.017553in}{2.298797in}}%
\pgfpathlineto{\pgfqpoint{5.018196in}{2.300665in}}%
\pgfpathlineto{\pgfqpoint{5.018931in}{2.299331in}}%
\pgfpathlineto{\pgfqpoint{5.019277in}{2.302612in}}%
\pgfpathlineto{\pgfqpoint{5.019759in}{2.304336in}}%
\pgfpathlineto{\pgfqpoint{5.020294in}{2.299397in}}%
\pgfpathlineto{\pgfqpoint{5.020333in}{2.299784in}}%
\pgfpathlineto{\pgfqpoint{5.020431in}{2.298875in}}%
\pgfpathlineto{\pgfqpoint{5.020732in}{2.303574in}}%
\pgfpathlineto{\pgfqpoint{5.020737in}{2.303436in}}%
\pgfpathlineto{\pgfqpoint{5.020776in}{2.304056in}}%
\pgfpathlineto{\pgfqpoint{5.021146in}{2.300196in}}%
\pgfpathlineto{\pgfqpoint{5.021852in}{2.303890in}}%
\pgfpathlineto{\pgfqpoint{5.022212in}{2.302450in}}%
\pgfpathlineto{\pgfqpoint{5.021911in}{2.304961in}}%
\pgfpathlineto{\pgfqpoint{5.022913in}{2.304099in}}%
\pgfpathlineto{\pgfqpoint{5.022923in}{2.304775in}}%
\pgfpathlineto{\pgfqpoint{5.023834in}{2.301574in}}%
\pgfpathlineto{\pgfqpoint{5.023985in}{2.302145in}}%
\pgfpathlineto{\pgfqpoint{5.023999in}{2.302569in}}%
\pgfpathlineto{\pgfqpoint{5.024540in}{2.298432in}}%
\pgfpathlineto{\pgfqpoint{5.025036in}{2.301023in}}%
\pgfpathlineto{\pgfqpoint{5.025796in}{2.296873in}}%
\pgfpathlineto{\pgfqpoint{5.025197in}{2.301177in}}%
\pgfpathlineto{\pgfqpoint{5.026234in}{2.297680in}}%
\pgfpathlineto{\pgfqpoint{5.027086in}{2.301633in}}%
\pgfpathlineto{\pgfqpoint{5.027500in}{2.301348in}}%
\pgfpathlineto{\pgfqpoint{5.027660in}{2.300481in}}%
\pgfpathlineto{\pgfqpoint{5.028556in}{2.305100in}}%
\pgfpathlineto{\pgfqpoint{5.029763in}{2.301794in}}%
\pgfpathlineto{\pgfqpoint{5.028600in}{2.305642in}}%
\pgfpathlineto{\pgfqpoint{5.029773in}{2.302258in}}%
\pgfpathlineto{\pgfqpoint{5.030070in}{2.304217in}}%
\pgfpathlineto{\pgfqpoint{5.030679in}{2.300437in}}%
\pgfpathlineto{\pgfqpoint{5.030844in}{2.301094in}}%
\pgfpathlineto{\pgfqpoint{5.031886in}{2.300239in}}%
\pgfpathlineto{\pgfqpoint{5.031565in}{2.304951in}}%
\pgfpathlineto{\pgfqpoint{5.031964in}{2.300699in}}%
\pgfpathlineto{\pgfqpoint{5.033176in}{2.303830in}}%
\pgfpathlineto{\pgfqpoint{5.032052in}{2.300416in}}%
\pgfpathlineto{\pgfqpoint{5.033639in}{2.302217in}}%
\pgfpathlineto{\pgfqpoint{5.034023in}{2.299912in}}%
\pgfpathlineto{\pgfqpoint{5.034734in}{2.302340in}}%
\pgfpathlineto{\pgfqpoint{5.034754in}{2.302071in}}%
\pgfpathlineto{\pgfqpoint{5.036316in}{2.295832in}}%
\pgfpathlineto{\pgfqpoint{5.034827in}{2.302094in}}%
\pgfpathlineto{\pgfqpoint{5.036467in}{2.296277in}}%
\pgfpathlineto{\pgfqpoint{5.036579in}{2.297760in}}%
\pgfpathlineto{\pgfqpoint{5.037022in}{2.295561in}}%
\pgfpathlineto{\pgfqpoint{5.037582in}{2.296577in}}%
\pgfpathlineto{\pgfqpoint{5.037874in}{2.300318in}}%
\pgfpathlineto{\pgfqpoint{5.038795in}{2.298999in}}%
\pgfpathlineto{\pgfqpoint{5.038843in}{2.298383in}}%
\pgfpathlineto{\pgfqpoint{5.039608in}{2.303411in}}%
\pgfpathlineto{\pgfqpoint{5.039617in}{2.303280in}}%
\pgfpathlineto{\pgfqpoint{5.040518in}{2.304412in}}%
\pgfpathlineto{\pgfqpoint{5.040036in}{2.302220in}}%
\pgfpathlineto{\pgfqpoint{5.040732in}{2.303733in}}%
\pgfpathlineto{\pgfqpoint{5.041857in}{2.301429in}}%
\pgfpathlineto{\pgfqpoint{5.040815in}{2.304784in}}%
\pgfpathlineto{\pgfqpoint{5.041906in}{2.302001in}}%
\pgfpathlineto{\pgfqpoint{5.043011in}{2.303198in}}%
\pgfpathlineto{\pgfqpoint{5.042349in}{2.300538in}}%
\pgfpathlineto{\pgfqpoint{5.043030in}{2.302909in}}%
\pgfpathlineto{\pgfqpoint{5.043035in}{2.302907in}}%
\pgfpathlineto{\pgfqpoint{5.043108in}{2.303899in}}%
\pgfpathlineto{\pgfqpoint{5.043152in}{2.304175in}}%
\pgfpathlineto{\pgfqpoint{5.043819in}{2.298875in}}%
\pgfpathlineto{\pgfqpoint{5.044091in}{2.300195in}}%
\pgfpathlineto{\pgfqpoint{5.044476in}{2.299170in}}%
\pgfpathlineto{\pgfqpoint{5.044744in}{2.302115in}}%
\pgfpathlineto{\pgfqpoint{5.045007in}{2.301411in}}%
\pgfpathlineto{\pgfqpoint{5.046088in}{2.303788in}}%
\pgfpathlineto{\pgfqpoint{5.045080in}{2.300865in}}%
\pgfpathlineto{\pgfqpoint{5.046131in}{2.302780in}}%
\pgfpathlineto{\pgfqpoint{5.047261in}{2.298527in}}%
\pgfpathlineto{\pgfqpoint{5.046151in}{2.302981in}}%
\pgfpathlineto{\pgfqpoint{5.047592in}{2.301195in}}%
\pgfpathlineto{\pgfqpoint{5.047933in}{2.299855in}}%
\pgfpathlineto{\pgfqpoint{5.048770in}{2.303586in}}%
\pgfpathlineto{\pgfqpoint{5.049719in}{2.300470in}}%
\pgfpathlineto{\pgfqpoint{5.048809in}{2.303675in}}%
\pgfpathlineto{\pgfqpoint{5.049919in}{2.302275in}}%
\pgfpathlineto{\pgfqpoint{5.050055in}{2.304796in}}%
\pgfpathlineto{\pgfqpoint{5.051010in}{2.302054in}}%
\pgfpathlineto{\pgfqpoint{5.051029in}{2.302325in}}%
\pgfpathlineto{\pgfqpoint{5.051501in}{2.300794in}}%
\pgfpathlineto{\pgfqpoint{5.051156in}{2.303294in}}%
\pgfpathlineto{\pgfqpoint{5.052120in}{2.302844in}}%
\pgfpathlineto{\pgfqpoint{5.052392in}{2.303662in}}%
\pgfpathlineto{\pgfqpoint{5.052762in}{2.301036in}}%
\pgfpathlineto{\pgfqpoint{5.052933in}{2.301411in}}%
\pgfpathlineto{\pgfqpoint{5.053449in}{2.298572in}}%
\pgfpathlineto{\pgfqpoint{5.053838in}{2.302062in}}%
\pgfpathlineto{\pgfqpoint{5.054038in}{2.301436in}}%
\pgfpathlineto{\pgfqpoint{5.054461in}{2.303960in}}%
\pgfpathlineto{\pgfqpoint{5.054116in}{2.301188in}}%
\pgfpathlineto{\pgfqpoint{5.055182in}{2.303204in}}%
\pgfpathlineto{\pgfqpoint{5.055255in}{2.301742in}}%
\pgfpathlineto{\pgfqpoint{5.055781in}{2.305795in}}%
\pgfpathlineto{\pgfqpoint{5.056253in}{2.304323in}}%
\pgfpathlineto{\pgfqpoint{5.056599in}{2.302380in}}%
\pgfpathlineto{\pgfqpoint{5.057051in}{2.304521in}}%
\pgfpathlineto{\pgfqpoint{5.057378in}{2.303892in}}%
\pgfpathlineto{\pgfqpoint{5.058035in}{2.306813in}}%
\pgfpathlineto{\pgfqpoint{5.057587in}{2.303626in}}%
\pgfpathlineto{\pgfqpoint{5.058497in}{2.304370in}}%
\pgfpathlineto{\pgfqpoint{5.058936in}{2.301789in}}%
\pgfpathlineto{\pgfqpoint{5.058531in}{2.304757in}}%
\pgfpathlineto{\pgfqpoint{5.059627in}{2.303647in}}%
\pgfpathlineto{\pgfqpoint{5.059851in}{2.303888in}}%
\pgfpathlineto{\pgfqpoint{5.059656in}{2.303087in}}%
\pgfpathlineto{\pgfqpoint{5.059904in}{2.302640in}}%
\pgfpathlineto{\pgfqpoint{5.059943in}{2.301707in}}%
\pgfpathlineto{\pgfqpoint{5.060430in}{2.303582in}}%
\pgfpathlineto{\pgfqpoint{5.061014in}{2.302276in}}%
\pgfpathlineto{\pgfqpoint{5.062061in}{2.303463in}}%
\pgfpathlineto{\pgfqpoint{5.061448in}{2.300424in}}%
\pgfpathlineto{\pgfqpoint{5.062129in}{2.302573in}}%
\pgfpathlineto{\pgfqpoint{5.062280in}{2.303902in}}%
\pgfpathlineto{\pgfqpoint{5.062752in}{2.299594in}}%
\pgfpathlineto{\pgfqpoint{5.062772in}{2.299785in}}%
\pgfpathlineto{\pgfqpoint{5.063200in}{2.298446in}}%
\pgfpathlineto{\pgfqpoint{5.063785in}{2.301378in}}%
\pgfpathlineto{\pgfqpoint{5.063853in}{2.301171in}}%
\pgfpathlineto{\pgfqpoint{5.064271in}{2.303668in}}%
\pgfpathlineto{\pgfqpoint{5.064700in}{2.299427in}}%
\pgfpathlineto{\pgfqpoint{5.064987in}{2.302119in}}%
\pgfpathlineto{\pgfqpoint{5.066092in}{2.298811in}}%
\pgfpathlineto{\pgfqpoint{5.066194in}{2.299509in}}%
\pgfpathlineto{\pgfqpoint{5.067397in}{2.303896in}}%
\pgfpathlineto{\pgfqpoint{5.066219in}{2.299459in}}%
\pgfpathlineto{\pgfqpoint{5.067412in}{2.303649in}}%
\pgfpathlineto{\pgfqpoint{5.067616in}{2.304343in}}%
\pgfpathlineto{\pgfqpoint{5.068074in}{2.301942in}}%
\pgfpathlineto{\pgfqpoint{5.068156in}{2.302470in}}%
\pgfpathlineto{\pgfqpoint{5.068161in}{2.302213in}}%
\pgfpathlineto{\pgfqpoint{5.068522in}{2.304000in}}%
\pgfpathlineto{\pgfqpoint{5.069262in}{2.302702in}}%
\pgfpathlineto{\pgfqpoint{5.069266in}{2.303073in}}%
\pgfpathlineto{\pgfqpoint{5.069729in}{2.298413in}}%
\pgfpathlineto{\pgfqpoint{5.070342in}{2.300437in}}%
\pgfpathlineto{\pgfqpoint{5.071272in}{2.298605in}}%
\pgfpathlineto{\pgfqpoint{5.071151in}{2.301027in}}%
\pgfpathlineto{\pgfqpoint{5.071452in}{2.300216in}}%
\pgfpathlineto{\pgfqpoint{5.072635in}{2.297797in}}%
\pgfpathlineto{\pgfqpoint{5.072295in}{2.301927in}}%
\pgfpathlineto{\pgfqpoint{5.072767in}{2.298710in}}%
\pgfpathlineto{\pgfqpoint{5.073833in}{2.301540in}}%
\pgfpathlineto{\pgfqpoint{5.072806in}{2.298667in}}%
\pgfpathlineto{\pgfqpoint{5.073916in}{2.300678in}}%
\pgfpathlineto{\pgfqpoint{5.074894in}{2.296967in}}%
\pgfpathlineto{\pgfqpoint{5.075041in}{2.299636in}}%
\pgfpathlineto{\pgfqpoint{5.075089in}{2.300225in}}%
\pgfpathlineto{\pgfqpoint{5.075561in}{2.297068in}}%
\pgfpathlineto{\pgfqpoint{5.076141in}{2.299145in}}%
\pgfpathlineto{\pgfqpoint{5.076189in}{2.298360in}}%
\pgfpathlineto{\pgfqpoint{5.076788in}{2.302137in}}%
\pgfpathlineto{\pgfqpoint{5.077168in}{2.301234in}}%
\pgfpathlineto{\pgfqpoint{5.078040in}{2.303425in}}%
\pgfpathlineto{\pgfqpoint{5.077275in}{2.299406in}}%
\pgfpathlineto{\pgfqpoint{5.078288in}{2.301859in}}%
\pgfpathlineto{\pgfqpoint{5.078293in}{2.301505in}}%
\pgfpathlineto{\pgfqpoint{5.079242in}{2.304327in}}%
\pgfpathlineto{\pgfqpoint{5.079388in}{2.302545in}}%
\pgfpathlineto{\pgfqpoint{5.079963in}{2.304828in}}%
\pgfpathlineto{\pgfqpoint{5.080420in}{2.300045in}}%
\pgfpathlineto{\pgfqpoint{5.080756in}{2.298914in}}%
\pgfpathlineto{\pgfqpoint{5.081199in}{2.301325in}}%
\pgfpathlineto{\pgfqpoint{5.081511in}{2.300301in}}%
\pgfpathlineto{\pgfqpoint{5.082202in}{2.303365in}}%
\pgfpathlineto{\pgfqpoint{5.082665in}{2.302848in}}%
\pgfpathlineto{\pgfqpoint{5.083186in}{2.304386in}}%
\pgfpathlineto{\pgfqpoint{5.082704in}{2.302011in}}%
\pgfpathlineto{\pgfqpoint{5.083799in}{2.303639in}}%
\pgfpathlineto{\pgfqpoint{5.084461in}{2.298601in}}%
\pgfpathlineto{\pgfqpoint{5.083814in}{2.303890in}}%
\pgfpathlineto{\pgfqpoint{5.085011in}{2.301440in}}%
\pgfpathlineto{\pgfqpoint{5.085313in}{2.304166in}}%
\pgfpathlineto{\pgfqpoint{5.085031in}{2.301246in}}%
\pgfpathlineto{\pgfqpoint{5.086175in}{2.303005in}}%
\pgfpathlineto{\pgfqpoint{5.086944in}{2.299197in}}%
\pgfpathlineto{\pgfqpoint{5.086223in}{2.303357in}}%
\pgfpathlineto{\pgfqpoint{5.087290in}{2.302709in}}%
\pgfpathlineto{\pgfqpoint{5.088663in}{2.307586in}}%
\pgfpathlineto{\pgfqpoint{5.087733in}{2.302293in}}%
\pgfpathlineto{\pgfqpoint{5.089106in}{2.305113in}}%
\pgfpathlineto{\pgfqpoint{5.090264in}{2.300886in}}%
\pgfpathlineto{\pgfqpoint{5.089349in}{2.305153in}}%
\pgfpathlineto{\pgfqpoint{5.090279in}{2.301054in}}%
\pgfpathlineto{\pgfqpoint{5.091146in}{2.304204in}}%
\pgfpathlineto{\pgfqpoint{5.090522in}{2.299736in}}%
\pgfpathlineto{\pgfqpoint{5.091511in}{2.302384in}}%
\pgfpathlineto{\pgfqpoint{5.092587in}{2.300198in}}%
\pgfpathlineto{\pgfqpoint{5.092114in}{2.303704in}}%
\pgfpathlineto{\pgfqpoint{5.092665in}{2.300904in}}%
\pgfpathlineto{\pgfqpoint{5.093750in}{2.303901in}}%
\pgfpathlineto{\pgfqpoint{5.093288in}{2.300606in}}%
\pgfpathlineto{\pgfqpoint{5.093872in}{2.303124in}}%
\pgfpathlineto{\pgfqpoint{5.094164in}{2.301261in}}%
\pgfpathlineto{\pgfqpoint{5.094850in}{2.307189in}}%
\pgfpathlineto{\pgfqpoint{5.094885in}{2.306707in}}%
\pgfpathlineto{\pgfqpoint{5.094987in}{2.305781in}}%
\pgfpathlineto{\pgfqpoint{5.094992in}{2.305816in}}%
\pgfpathlineto{\pgfqpoint{5.096111in}{2.301193in}}%
\pgfpathlineto{\pgfqpoint{5.096150in}{2.301440in}}%
\pgfpathlineto{\pgfqpoint{5.096978in}{2.303975in}}%
\pgfpathlineto{\pgfqpoint{5.096477in}{2.299880in}}%
\pgfpathlineto{\pgfqpoint{5.097324in}{2.302657in}}%
\pgfpathlineto{\pgfqpoint{5.098439in}{2.301330in}}%
\pgfpathlineto{\pgfqpoint{5.097830in}{2.303993in}}%
\pgfpathlineto{\pgfqpoint{5.098453in}{2.301507in}}%
\pgfpathlineto{\pgfqpoint{5.099568in}{2.302120in}}%
\pgfpathlineto{\pgfqpoint{5.098750in}{2.299504in}}%
\pgfpathlineto{\pgfqpoint{5.099583in}{2.301784in}}%
\pgfpathlineto{\pgfqpoint{5.100259in}{2.298633in}}%
\pgfpathlineto{\pgfqpoint{5.099865in}{2.301975in}}%
\pgfpathlineto{\pgfqpoint{5.100736in}{2.299281in}}%
\pgfpathlineto{\pgfqpoint{5.101238in}{2.300197in}}%
\pgfpathlineto{\pgfqpoint{5.101749in}{2.297031in}}%
\pgfpathlineto{\pgfqpoint{5.101788in}{2.296724in}}%
\pgfpathlineto{\pgfqpoint{5.102114in}{2.299232in}}%
\pgfpathlineto{\pgfqpoint{5.103180in}{2.303295in}}%
\pgfpathlineto{\pgfqpoint{5.102518in}{2.298869in}}%
\pgfpathlineto{\pgfqpoint{5.103380in}{2.302404in}}%
\pgfpathlineto{\pgfqpoint{5.103589in}{2.300979in}}%
\pgfpathlineto{\pgfqpoint{5.104008in}{2.303709in}}%
\pgfpathlineto{\pgfqpoint{5.104441in}{2.303158in}}%
\pgfpathlineto{\pgfqpoint{5.104544in}{2.304958in}}%
\pgfpathlineto{\pgfqpoint{5.104777in}{2.301235in}}%
\pgfpathlineto{\pgfqpoint{5.105542in}{2.303064in}}%
\pgfpathlineto{\pgfqpoint{5.106038in}{2.301438in}}%
\pgfpathlineto{\pgfqpoint{5.106559in}{2.304448in}}%
\pgfpathlineto{\pgfqpoint{5.106627in}{2.303846in}}%
\pgfpathlineto{\pgfqpoint{5.106851in}{2.305378in}}%
\pgfpathlineto{\pgfqpoint{5.107479in}{2.302522in}}%
\pgfpathlineto{\pgfqpoint{5.107504in}{2.302915in}}%
\pgfpathlineto{\pgfqpoint{5.108662in}{2.295715in}}%
\pgfpathlineto{\pgfqpoint{5.107844in}{2.303838in}}%
\pgfpathlineto{\pgfqpoint{5.108696in}{2.296072in}}%
\pgfpathlineto{\pgfqpoint{5.109836in}{2.300348in}}%
\pgfpathlineto{\pgfqpoint{5.109471in}{2.294734in}}%
\pgfpathlineto{\pgfqpoint{5.109870in}{2.299665in}}%
\pgfpathlineto{\pgfqpoint{5.109918in}{2.298705in}}%
\pgfpathlineto{\pgfqpoint{5.110902in}{2.301684in}}%
\pgfpathlineto{\pgfqpoint{5.110907in}{2.301656in}}%
\pgfpathlineto{\pgfqpoint{5.112041in}{2.304072in}}%
\pgfpathlineto{\pgfqpoint{5.111267in}{2.301444in}}%
\pgfpathlineto{\pgfqpoint{5.112056in}{2.303452in}}%
\pgfpathlineto{\pgfqpoint{5.112840in}{2.299122in}}%
\pgfpathlineto{\pgfqpoint{5.112177in}{2.304064in}}%
\pgfpathlineto{\pgfqpoint{5.113180in}{2.301877in}}%
\pgfpathlineto{\pgfqpoint{5.114183in}{2.304286in}}%
\pgfpathlineto{\pgfqpoint{5.113633in}{2.300268in}}%
\pgfpathlineto{\pgfqpoint{5.114315in}{2.303535in}}%
\pgfpathlineto{\pgfqpoint{5.114850in}{2.301501in}}%
\pgfpathlineto{\pgfqpoint{5.114446in}{2.303959in}}%
\pgfpathlineto{\pgfqpoint{5.115464in}{2.302986in}}%
\pgfpathlineto{\pgfqpoint{5.116160in}{2.305188in}}%
\pgfpathlineto{\pgfqpoint{5.115702in}{2.301772in}}%
\pgfpathlineto{\pgfqpoint{5.116574in}{2.303064in}}%
\pgfpathlineto{\pgfqpoint{5.116705in}{2.302056in}}%
\pgfpathlineto{\pgfqpoint{5.117518in}{2.303986in}}%
\pgfpathlineto{\pgfqpoint{5.117689in}{2.302739in}}%
\pgfpathlineto{\pgfqpoint{5.118258in}{2.303932in}}%
\pgfpathlineto{\pgfqpoint{5.118472in}{2.300797in}}%
\pgfpathlineto{\pgfqpoint{5.118789in}{2.302581in}}%
\pgfpathlineto{\pgfqpoint{5.119451in}{2.299825in}}%
\pgfpathlineto{\pgfqpoint{5.119884in}{2.303325in}}%
\pgfpathlineto{\pgfqpoint{5.119894in}{2.303059in}}%
\pgfpathlineto{\pgfqpoint{5.120678in}{2.304673in}}%
\pgfpathlineto{\pgfqpoint{5.119957in}{2.302376in}}%
\pgfpathlineto{\pgfqpoint{5.121024in}{2.303448in}}%
\pgfpathlineto{\pgfqpoint{5.121812in}{2.300822in}}%
\pgfpathlineto{\pgfqpoint{5.121189in}{2.304046in}}%
\pgfpathlineto{\pgfqpoint{5.122138in}{2.303202in}}%
\pgfpathlineto{\pgfqpoint{5.122211in}{2.303839in}}%
\pgfpathlineto{\pgfqpoint{5.122396in}{2.300610in}}%
\pgfpathlineto{\pgfqpoint{5.123088in}{2.301530in}}%
\pgfpathlineto{\pgfqpoint{5.123117in}{2.301103in}}%
\pgfpathlineto{\pgfqpoint{5.123584in}{2.304953in}}%
\pgfpathlineto{\pgfqpoint{5.124188in}{2.301891in}}%
\pgfpathlineto{\pgfqpoint{5.124816in}{2.303691in}}%
\pgfpathlineto{\pgfqpoint{5.125176in}{2.299921in}}%
\pgfpathlineto{\pgfqpoint{5.125259in}{2.300001in}}%
\pgfpathlineto{\pgfqpoint{5.125668in}{2.298933in}}%
\pgfpathlineto{\pgfqpoint{5.126228in}{2.301689in}}%
\pgfpathlineto{\pgfqpoint{5.126350in}{2.300024in}}%
\pgfpathlineto{\pgfqpoint{5.127036in}{2.304365in}}%
\pgfpathlineto{\pgfqpoint{5.126510in}{2.298997in}}%
\pgfpathlineto{\pgfqpoint{5.127484in}{2.302577in}}%
\pgfpathlineto{\pgfqpoint{5.128545in}{2.298810in}}%
\pgfpathlineto{\pgfqpoint{5.127732in}{2.302843in}}%
\pgfpathlineto{\pgfqpoint{5.128667in}{2.300429in}}%
\pgfpathlineto{\pgfqpoint{5.129134in}{2.299142in}}%
\pgfpathlineto{\pgfqpoint{5.129441in}{2.302241in}}%
\pgfpathlineto{\pgfqpoint{5.129836in}{2.299508in}}%
\pgfpathlineto{\pgfqpoint{5.130980in}{2.303492in}}%
\pgfpathlineto{\pgfqpoint{5.130999in}{2.303314in}}%
\pgfpathlineto{\pgfqpoint{5.132036in}{2.301503in}}%
\pgfpathlineto{\pgfqpoint{5.131447in}{2.304040in}}%
\pgfpathlineto{\pgfqpoint{5.132163in}{2.302563in}}%
\pgfpathlineto{\pgfqpoint{5.133146in}{2.306193in}}%
\pgfpathlineto{\pgfqpoint{5.132206in}{2.302215in}}%
\pgfpathlineto{\pgfqpoint{5.133336in}{2.305136in}}%
\pgfpathlineto{\pgfqpoint{5.134251in}{2.300715in}}%
\pgfpathlineto{\pgfqpoint{5.134719in}{2.304320in}}%
\pgfpathlineto{\pgfqpoint{5.134782in}{2.304665in}}%
\pgfpathlineto{\pgfqpoint{5.135420in}{2.300725in}}%
\pgfpathlineto{\pgfqpoint{5.135785in}{2.303232in}}%
\pgfpathlineto{\pgfqpoint{5.136413in}{2.300904in}}%
\pgfpathlineto{\pgfqpoint{5.136895in}{2.302963in}}%
\pgfpathlineto{\pgfqpoint{5.136924in}{2.302295in}}%
\pgfpathlineto{\pgfqpoint{5.137732in}{2.304571in}}%
\pgfpathlineto{\pgfqpoint{5.137903in}{2.303571in}}%
\pgfpathlineto{\pgfqpoint{5.138526in}{2.304955in}}%
\pgfpathlineto{\pgfqpoint{5.138983in}{2.302517in}}%
\pgfpathlineto{\pgfqpoint{5.139280in}{2.300656in}}%
\pgfpathlineto{\pgfqpoint{5.139616in}{2.304417in}}%
\pgfpathlineto{\pgfqpoint{5.140079in}{2.303641in}}%
\pgfpathlineto{\pgfqpoint{5.140128in}{2.304360in}}%
\pgfpathlineto{\pgfqpoint{5.140872in}{2.300353in}}%
\pgfpathlineto{\pgfqpoint{5.141140in}{2.301937in}}%
\pgfpathlineto{\pgfqpoint{5.142041in}{2.300598in}}%
\pgfpathlineto{\pgfqpoint{5.141291in}{2.303973in}}%
\pgfpathlineto{\pgfqpoint{5.142240in}{2.301966in}}%
\pgfpathlineto{\pgfqpoint{5.142245in}{2.302297in}}%
\pgfpathlineto{\pgfqpoint{5.143015in}{2.297084in}}%
\pgfpathlineto{\pgfqpoint{5.143307in}{2.299120in}}%
\pgfpathlineto{\pgfqpoint{5.143759in}{2.298031in}}%
\pgfpathlineto{\pgfqpoint{5.144271in}{2.300340in}}%
\pgfpathlineto{\pgfqpoint{5.144412in}{2.299133in}}%
\pgfpathlineto{\pgfqpoint{5.144913in}{2.301038in}}%
\pgfpathlineto{\pgfqpoint{5.145371in}{2.297628in}}%
\pgfpathlineto{\pgfqpoint{5.145400in}{2.297087in}}%
\pgfpathlineto{\pgfqpoint{5.146233in}{2.300777in}}%
\pgfpathlineto{\pgfqpoint{5.146393in}{2.300008in}}%
\pgfpathlineto{\pgfqpoint{5.146700in}{2.301535in}}%
\pgfpathlineto{\pgfqpoint{5.147304in}{2.298682in}}%
\pgfpathlineto{\pgfqpoint{5.147333in}{2.298696in}}%
\pgfpathlineto{\pgfqpoint{5.147343in}{2.298381in}}%
\pgfpathlineto{\pgfqpoint{5.148389in}{2.300741in}}%
\pgfpathlineto{\pgfqpoint{5.148847in}{2.301777in}}%
\pgfpathlineto{\pgfqpoint{5.149017in}{2.299390in}}%
\pgfpathlineto{\pgfqpoint{5.149480in}{2.300491in}}%
\pgfpathlineto{\pgfqpoint{5.149499in}{2.300017in}}%
\pgfpathlineto{\pgfqpoint{5.149972in}{2.305201in}}%
\pgfpathlineto{\pgfqpoint{5.150429in}{2.302421in}}%
\pgfpathlineto{\pgfqpoint{5.150936in}{2.306091in}}%
\pgfpathlineto{\pgfqpoint{5.151593in}{2.304118in}}%
\pgfpathlineto{\pgfqpoint{5.151671in}{2.304859in}}%
\pgfpathlineto{\pgfqpoint{5.152484in}{2.300726in}}%
\pgfpathlineto{\pgfqpoint{5.152864in}{2.298875in}}%
\pgfpathlineto{\pgfqpoint{5.153433in}{2.303135in}}%
\pgfpathlineto{\pgfqpoint{5.153569in}{2.301827in}}%
\pgfpathlineto{\pgfqpoint{5.154782in}{2.303422in}}%
\pgfpathlineto{\pgfqpoint{5.154115in}{2.300554in}}%
\pgfpathlineto{\pgfqpoint{5.154811in}{2.303007in}}%
\pgfpathlineto{\pgfqpoint{5.155342in}{2.301532in}}%
\pgfpathlineto{\pgfqpoint{5.155838in}{2.304303in}}%
\pgfpathlineto{\pgfqpoint{5.155906in}{2.303822in}}%
\pgfpathlineto{\pgfqpoint{5.156525in}{2.304249in}}%
\pgfpathlineto{\pgfqpoint{5.156778in}{2.301268in}}%
\pgfpathlineto{\pgfqpoint{5.156895in}{2.301456in}}%
\pgfpathlineto{\pgfqpoint{5.156909in}{2.301272in}}%
\pgfpathlineto{\pgfqpoint{5.157133in}{2.303140in}}%
\pgfpathlineto{\pgfqpoint{5.157800in}{2.303079in}}%
\pgfpathlineto{\pgfqpoint{5.158370in}{2.304348in}}%
\pgfpathlineto{\pgfqpoint{5.158803in}{2.301400in}}%
\pgfpathlineto{\pgfqpoint{5.158823in}{2.301633in}}%
\pgfpathlineto{\pgfqpoint{5.159256in}{2.301222in}}%
\pgfpathlineto{\pgfqpoint{5.159723in}{2.303979in}}%
\pgfpathlineto{\pgfqpoint{5.159918in}{2.302303in}}%
\pgfpathlineto{\pgfqpoint{5.160867in}{2.304549in}}%
\pgfpathlineto{\pgfqpoint{5.160093in}{2.301290in}}%
\pgfpathlineto{\pgfqpoint{5.161018in}{2.302388in}}%
\pgfpathlineto{\pgfqpoint{5.161495in}{2.300415in}}%
\pgfpathlineto{\pgfqpoint{5.162060in}{2.303699in}}%
\pgfpathlineto{\pgfqpoint{5.162104in}{2.303490in}}%
\pgfpathlineto{\pgfqpoint{5.162255in}{2.301713in}}%
\pgfpathlineto{\pgfqpoint{5.162742in}{2.304904in}}%
\pgfpathlineto{\pgfqpoint{5.163219in}{2.303265in}}%
\pgfpathlineto{\pgfqpoint{5.164197in}{2.303613in}}%
\pgfpathlineto{\pgfqpoint{5.163788in}{2.300899in}}%
\pgfpathlineto{\pgfqpoint{5.164270in}{2.302725in}}%
\pgfpathlineto{\pgfqpoint{5.165010in}{2.300000in}}%
\pgfpathlineto{\pgfqpoint{5.164465in}{2.304006in}}%
\pgfpathlineto{\pgfqpoint{5.165434in}{2.300786in}}%
\pgfpathlineto{\pgfqpoint{5.165872in}{2.302653in}}%
\pgfpathlineto{\pgfqpoint{5.165468in}{2.300195in}}%
\pgfpathlineto{\pgfqpoint{5.166534in}{2.300588in}}%
\pgfpathlineto{\pgfqpoint{5.166598in}{2.300205in}}%
\pgfpathlineto{\pgfqpoint{5.167143in}{2.303346in}}%
\pgfpathlineto{\pgfqpoint{5.167615in}{2.302137in}}%
\pgfpathlineto{\pgfqpoint{5.167771in}{2.303452in}}%
\pgfpathlineto{\pgfqpoint{5.168287in}{2.300914in}}%
\pgfpathlineto{\pgfqpoint{5.168433in}{2.301144in}}%
\pgfpathlineto{\pgfqpoint{5.169465in}{2.296790in}}%
\pgfpathlineto{\pgfqpoint{5.169577in}{2.298564in}}%
\pgfpathlineto{\pgfqpoint{5.170629in}{2.304448in}}%
\pgfpathlineto{\pgfqpoint{5.170770in}{2.303032in}}%
\pgfpathlineto{\pgfqpoint{5.171081in}{2.299145in}}%
\pgfpathlineto{\pgfqpoint{5.171938in}{2.300570in}}%
\pgfpathlineto{\pgfqpoint{5.171987in}{2.299758in}}%
\pgfpathlineto{\pgfqpoint{5.172484in}{2.302256in}}%
\pgfpathlineto{\pgfqpoint{5.172493in}{2.302235in}}%
\pgfpathlineto{\pgfqpoint{5.173569in}{2.305647in}}%
\pgfpathlineto{\pgfqpoint{5.172571in}{2.301356in}}%
\pgfpathlineto{\pgfqpoint{5.173861in}{2.303418in}}%
\pgfpathlineto{\pgfqpoint{5.174553in}{2.299431in}}%
\pgfpathlineto{\pgfqpoint{5.173993in}{2.303682in}}%
\pgfpathlineto{\pgfqpoint{5.174991in}{2.301596in}}%
\pgfpathlineto{\pgfqpoint{5.175994in}{2.301120in}}%
\pgfpathlineto{\pgfqpoint{5.175205in}{2.303739in}}%
\pgfpathlineto{\pgfqpoint{5.176057in}{2.301886in}}%
\pgfpathlineto{\pgfqpoint{5.176330in}{2.303152in}}%
\pgfpathlineto{\pgfqpoint{5.176880in}{2.300433in}}%
\pgfpathlineto{\pgfqpoint{5.177157in}{2.301649in}}%
\pgfpathlineto{\pgfqpoint{5.177162in}{2.301287in}}%
\pgfpathlineto{\pgfqpoint{5.177883in}{2.305389in}}%
\pgfpathlineto{\pgfqpoint{5.178243in}{2.303520in}}%
\pgfpathlineto{\pgfqpoint{5.178881in}{2.301866in}}%
\pgfpathlineto{\pgfqpoint{5.178691in}{2.304622in}}%
\pgfpathlineto{\pgfqpoint{5.179309in}{2.303582in}}%
\pgfpathlineto{\pgfqpoint{5.179407in}{2.304270in}}%
\pgfpathlineto{\pgfqpoint{5.180332in}{2.302168in}}%
\pgfpathlineto{\pgfqpoint{5.180405in}{2.302804in}}%
\pgfpathlineto{\pgfqpoint{5.181408in}{2.301238in}}%
\pgfpathlineto{\pgfqpoint{5.180585in}{2.304008in}}%
\pgfpathlineto{\pgfqpoint{5.181505in}{2.303037in}}%
\pgfpathlineto{\pgfqpoint{5.182284in}{2.305533in}}%
\pgfpathlineto{\pgfqpoint{5.181885in}{2.301549in}}%
\pgfpathlineto{\pgfqpoint{5.182649in}{2.304024in}}%
\pgfpathlineto{\pgfqpoint{5.183749in}{2.302416in}}%
\pgfpathlineto{\pgfqpoint{5.183642in}{2.305052in}}%
\pgfpathlineto{\pgfqpoint{5.183783in}{2.302681in}}%
\pgfpathlineto{\pgfqpoint{5.184324in}{2.303669in}}%
\pgfpathlineto{\pgfqpoint{5.184850in}{2.300245in}}%
\pgfpathlineto{\pgfqpoint{5.184864in}{2.300482in}}%
\pgfpathlineto{\pgfqpoint{5.185809in}{2.304239in}}%
\pgfpathlineto{\pgfqpoint{5.185030in}{2.299857in}}%
\pgfpathlineto{\pgfqpoint{5.186038in}{2.302401in}}%
\pgfpathlineto{\pgfqpoint{5.186787in}{2.303972in}}%
\pgfpathlineto{\pgfqpoint{5.186529in}{2.300600in}}%
\pgfpathlineto{\pgfqpoint{5.186972in}{2.301571in}}%
\pgfpathlineto{\pgfqpoint{5.186977in}{2.301345in}}%
\pgfpathlineto{\pgfqpoint{5.187410in}{2.303756in}}%
\pgfpathlineto{\pgfqpoint{5.188053in}{2.302909in}}%
\pgfpathlineto{\pgfqpoint{5.189075in}{2.302072in}}%
\pgfpathlineto{\pgfqpoint{5.188754in}{2.305487in}}%
\pgfpathlineto{\pgfqpoint{5.189183in}{2.302473in}}%
\pgfpathlineto{\pgfqpoint{5.190039in}{2.304338in}}%
\pgfpathlineto{\pgfqpoint{5.189319in}{2.301987in}}%
\pgfpathlineto{\pgfqpoint{5.190312in}{2.303614in}}%
\pgfpathlineto{\pgfqpoint{5.190658in}{2.300826in}}%
\pgfpathlineto{\pgfqpoint{5.191286in}{2.304235in}}%
\pgfpathlineto{\pgfqpoint{5.191417in}{2.303738in}}%
\pgfpathlineto{\pgfqpoint{5.192157in}{2.301182in}}%
\pgfpathlineto{\pgfqpoint{5.191524in}{2.304577in}}%
\pgfpathlineto{\pgfqpoint{5.192566in}{2.302889in}}%
\pgfpathlineto{\pgfqpoint{5.192970in}{2.304429in}}%
\pgfpathlineto{\pgfqpoint{5.193452in}{2.301087in}}%
\pgfpathlineto{\pgfqpoint{5.193671in}{2.302396in}}%
\pgfpathlineto{\pgfqpoint{5.193822in}{2.300986in}}%
\pgfpathlineto{\pgfqpoint{5.193681in}{2.302465in}}%
\pgfpathlineto{\pgfqpoint{5.193856in}{2.301193in}}%
\pgfpathlineto{\pgfqpoint{5.194027in}{2.298433in}}%
\pgfpathlineto{\pgfqpoint{5.194854in}{2.302423in}}%
\pgfpathlineto{\pgfqpoint{5.194966in}{2.300928in}}%
\pgfpathlineto{\pgfqpoint{5.195263in}{2.302089in}}%
\pgfpathlineto{\pgfqpoint{5.195813in}{2.299168in}}%
\pgfpathlineto{\pgfqpoint{5.196028in}{2.300487in}}%
\pgfpathlineto{\pgfqpoint{5.196398in}{2.298320in}}%
\pgfpathlineto{\pgfqpoint{5.197084in}{2.301770in}}%
\pgfpathlineto{\pgfqpoint{5.197133in}{2.300861in}}%
\pgfpathlineto{\pgfqpoint{5.197186in}{2.300417in}}%
\pgfpathlineto{\pgfqpoint{5.197620in}{2.305165in}}%
\pgfpathlineto{\pgfqpoint{5.198180in}{2.302716in}}%
\pgfpathlineto{\pgfqpoint{5.199002in}{2.304986in}}%
\pgfpathlineto{\pgfqpoint{5.198472in}{2.300983in}}%
\pgfpathlineto{\pgfqpoint{5.199299in}{2.303163in}}%
\pgfpathlineto{\pgfqpoint{5.199674in}{2.300788in}}%
\pgfpathlineto{\pgfqpoint{5.199319in}{2.303217in}}%
\pgfpathlineto{\pgfqpoint{5.200409in}{2.303109in}}%
\pgfpathlineto{\pgfqpoint{5.200701in}{2.304786in}}%
\pgfpathlineto{\pgfqpoint{5.201183in}{2.302431in}}%
\pgfpathlineto{\pgfqpoint{5.201213in}{2.302414in}}%
\pgfpathlineto{\pgfqpoint{5.201982in}{2.301842in}}%
\pgfpathlineto{\pgfqpoint{5.201763in}{2.303823in}}%
\pgfpathlineto{\pgfqpoint{5.202303in}{2.303341in}}%
\pgfpathlineto{\pgfqpoint{5.203131in}{2.304836in}}%
\pgfpathlineto{\pgfqpoint{5.203355in}{2.302133in}}%
\pgfpathlineto{\pgfqpoint{5.203384in}{2.302471in}}%
\pgfpathlineto{\pgfqpoint{5.204479in}{2.301760in}}%
\pgfpathlineto{\pgfqpoint{5.203735in}{2.305672in}}%
\pgfpathlineto{\pgfqpoint{5.204509in}{2.302061in}}%
\pgfpathlineto{\pgfqpoint{5.204523in}{2.301907in}}%
\pgfpathlineto{\pgfqpoint{5.204543in}{2.302904in}}%
\pgfpathlineto{\pgfqpoint{5.204606in}{2.303161in}}%
\pgfpathlineto{\pgfqpoint{5.204845in}{2.304472in}}%
\pgfpathlineto{\pgfqpoint{5.205560in}{2.299655in}}%
\pgfpathlineto{\pgfqpoint{5.205570in}{2.299445in}}%
\pgfpathlineto{\pgfqpoint{5.206247in}{2.303004in}}%
\pgfpathlineto{\pgfqpoint{5.206636in}{2.300796in}}%
\pgfpathlineto{\pgfqpoint{5.206641in}{2.300675in}}%
\pgfpathlineto{\pgfqpoint{5.207191in}{2.303552in}}%
\pgfpathlineto{\pgfqpoint{5.207693in}{2.302119in}}%
\pgfpathlineto{\pgfqpoint{5.208316in}{2.303569in}}%
\pgfpathlineto{\pgfqpoint{5.208657in}{2.300132in}}%
\pgfpathlineto{\pgfqpoint{5.208793in}{2.301443in}}%
\pgfpathlineto{\pgfqpoint{5.208803in}{2.301108in}}%
\pgfpathlineto{\pgfqpoint{5.209241in}{2.304411in}}%
\pgfpathlineto{\pgfqpoint{5.209903in}{2.301444in}}%
\pgfpathlineto{\pgfqpoint{5.209908in}{2.301386in}}%
\pgfpathlineto{\pgfqpoint{5.210750in}{2.304398in}}%
\pgfpathlineto{\pgfqpoint{5.210779in}{2.304053in}}%
\pgfpathlineto{\pgfqpoint{5.210784in}{2.304412in}}%
\pgfpathlineto{\pgfqpoint{5.211724in}{2.301124in}}%
\pgfpathlineto{\pgfqpoint{5.211865in}{2.302188in}}%
\pgfpathlineto{\pgfqpoint{5.212181in}{2.304274in}}%
\pgfpathlineto{\pgfqpoint{5.212819in}{2.298704in}}%
\pgfpathlineto{\pgfqpoint{5.212868in}{2.299217in}}%
\pgfpathlineto{\pgfqpoint{5.213364in}{2.297836in}}%
\pgfpathlineto{\pgfqpoint{5.213705in}{2.300550in}}%
\pgfpathlineto{\pgfqpoint{5.213744in}{2.300306in}}%
\pgfpathlineto{\pgfqpoint{5.214294in}{2.301542in}}%
\pgfpathlineto{\pgfqpoint{5.214494in}{2.298221in}}%
\pgfpathlineto{\pgfqpoint{5.214810in}{2.299159in}}%
\pgfpathlineto{\pgfqpoint{5.215619in}{2.296449in}}%
\pgfpathlineto{\pgfqpoint{5.215102in}{2.299789in}}%
\pgfpathlineto{\pgfqpoint{5.215974in}{2.299075in}}%
\pgfpathlineto{\pgfqpoint{5.216602in}{2.303657in}}%
\pgfpathlineto{\pgfqpoint{5.216032in}{2.298900in}}%
\pgfpathlineto{\pgfqpoint{5.217211in}{2.302306in}}%
\pgfpathlineto{\pgfqpoint{5.218301in}{2.296927in}}%
\pgfpathlineto{\pgfqpoint{5.217391in}{2.303711in}}%
\pgfpathlineto{\pgfqpoint{5.218442in}{2.297549in}}%
\pgfpathlineto{\pgfqpoint{5.219572in}{2.298714in}}%
\pgfpathlineto{\pgfqpoint{5.219177in}{2.295429in}}%
\pgfpathlineto{\pgfqpoint{5.219616in}{2.298478in}}%
\pgfpathlineto{\pgfqpoint{5.219728in}{2.297118in}}%
\pgfpathlineto{\pgfqpoint{5.220365in}{2.302262in}}%
\pgfpathlineto{\pgfqpoint{5.221344in}{2.304224in}}%
\pgfpathlineto{\pgfqpoint{5.220618in}{2.300888in}}%
\pgfpathlineto{\pgfqpoint{5.221490in}{2.303582in}}%
\pgfpathlineto{\pgfqpoint{5.222512in}{2.301915in}}%
\pgfpathlineto{\pgfqpoint{5.221714in}{2.304297in}}%
\pgfpathlineto{\pgfqpoint{5.222619in}{2.302071in}}%
\pgfpathlineto{\pgfqpoint{5.224596in}{2.294632in}}%
\pgfpathlineto{\pgfqpoint{5.223043in}{2.303089in}}%
\pgfpathlineto{\pgfqpoint{5.225088in}{2.295479in}}%
\pgfpathlineto{\pgfqpoint{5.226183in}{2.300017in}}%
\pgfpathlineto{\pgfqpoint{5.225317in}{2.294891in}}%
\pgfpathlineto{\pgfqpoint{5.226251in}{2.298929in}}%
\pgfpathlineto{\pgfqpoint{5.226777in}{2.300912in}}%
\pgfpathlineto{\pgfqpoint{5.227235in}{2.297098in}}%
\pgfpathlineto{\pgfqpoint{5.227503in}{2.294648in}}%
\pgfpathlineto{\pgfqpoint{5.228072in}{2.297259in}}%
\pgfpathlineto{\pgfqpoint{5.228369in}{2.295627in}}%
\pgfpathlineto{\pgfqpoint{5.229187in}{2.294268in}}%
\pgfpathlineto{\pgfqpoint{5.228583in}{2.298253in}}%
\pgfpathlineto{\pgfqpoint{5.229469in}{2.295683in}}%
\pgfpathlineto{\pgfqpoint{5.229474in}{2.295795in}}%
\pgfpathlineto{\pgfqpoint{5.230224in}{2.291602in}}%
\pgfpathlineto{\pgfqpoint{5.230463in}{2.293172in}}%
\pgfpathlineto{\pgfqpoint{5.230823in}{2.292152in}}%
\pgfpathlineto{\pgfqpoint{5.231490in}{2.295166in}}%
\pgfpathlineto{\pgfqpoint{5.231646in}{2.295960in}}%
\pgfpathlineto{\pgfqpoint{5.232347in}{2.293103in}}%
\pgfpathlineto{\pgfqpoint{5.232527in}{2.293814in}}%
\pgfpathlineto{\pgfqpoint{5.233145in}{2.291973in}}%
\pgfpathlineto{\pgfqpoint{5.232746in}{2.294841in}}%
\pgfpathlineto{\pgfqpoint{5.233676in}{2.292157in}}%
\pgfpathlineto{\pgfqpoint{5.233978in}{2.293290in}}%
\pgfpathlineto{\pgfqpoint{5.234474in}{2.289177in}}%
\pgfpathlineto{\pgfqpoint{5.234737in}{2.290150in}}%
\pgfpathlineto{\pgfqpoint{5.235833in}{2.296036in}}%
\pgfpathlineto{\pgfqpoint{5.236621in}{2.293926in}}%
\pgfpathlineto{\pgfqpoint{5.236636in}{2.293682in}}%
\pgfpathlineto{\pgfqpoint{5.237507in}{2.297737in}}%
\pgfpathlineto{\pgfqpoint{5.237663in}{2.296468in}}%
\pgfpathlineto{\pgfqpoint{5.238048in}{2.296268in}}%
\pgfpathlineto{\pgfqpoint{5.238413in}{2.299085in}}%
\pgfpathlineto{\pgfqpoint{5.238418in}{2.298922in}}%
\pgfpathlineto{\pgfqpoint{5.239596in}{2.302406in}}%
\pgfpathlineto{\pgfqpoint{5.238457in}{2.298095in}}%
\pgfpathlineto{\pgfqpoint{5.239601in}{2.302393in}}%
\pgfpathlineto{\pgfqpoint{5.239791in}{2.301290in}}%
\pgfpathlineto{\pgfqpoint{5.240253in}{2.304168in}}%
\pgfpathlineto{\pgfqpoint{5.240672in}{2.302875in}}%
\pgfpathlineto{\pgfqpoint{5.241105in}{2.304763in}}%
\pgfpathlineto{\pgfqpoint{5.241762in}{2.302361in}}%
\pgfpathlineto{\pgfqpoint{5.241772in}{2.302552in}}%
\pgfpathlineto{\pgfqpoint{5.242220in}{2.301844in}}%
\pgfpathlineto{\pgfqpoint{5.242658in}{2.305060in}}%
\pgfpathlineto{\pgfqpoint{5.242819in}{2.304850in}}%
\pgfpathlineto{\pgfqpoint{5.243184in}{2.306637in}}%
\pgfpathlineto{\pgfqpoint{5.243612in}{2.302279in}}%
\pgfpathlineto{\pgfqpoint{5.243914in}{2.304317in}}%
\pgfpathlineto{\pgfqpoint{5.244645in}{2.302251in}}%
\pgfpathlineto{\pgfqpoint{5.243973in}{2.304840in}}%
\pgfpathlineto{\pgfqpoint{5.245073in}{2.302753in}}%
\pgfpathlineto{\pgfqpoint{5.245974in}{2.305692in}}%
\pgfpathlineto{\pgfqpoint{5.245136in}{2.301891in}}%
\pgfpathlineto{\pgfqpoint{5.246198in}{2.302895in}}%
\pgfpathlineto{\pgfqpoint{5.246680in}{2.301404in}}%
\pgfpathlineto{\pgfqpoint{5.247210in}{2.304072in}}%
\pgfpathlineto{\pgfqpoint{5.247298in}{2.303529in}}%
\pgfpathlineto{\pgfqpoint{5.247386in}{2.302544in}}%
\pgfpathlineto{\pgfqpoint{5.247532in}{2.304554in}}%
\pgfpathlineto{\pgfqpoint{5.247794in}{2.304501in}}%
\pgfpathlineto{\pgfqpoint{5.248325in}{2.304612in}}%
\pgfpathlineto{\pgfqpoint{5.247989in}{2.301978in}}%
\pgfpathlineto{\pgfqpoint{5.248846in}{2.303623in}}%
\pgfpathlineto{\pgfqpoint{5.248856in}{2.303823in}}%
\pgfpathlineto{\pgfqpoint{5.249153in}{2.301113in}}%
\pgfpathlineto{\pgfqpoint{5.249854in}{2.299404in}}%
\pgfpathlineto{\pgfqpoint{5.249309in}{2.302433in}}%
\pgfpathlineto{\pgfqpoint{5.250253in}{2.301648in}}%
\pgfpathlineto{\pgfqpoint{5.250613in}{2.302912in}}%
\pgfpathlineto{\pgfqpoint{5.250901in}{2.297862in}}%
\pgfpathlineto{\pgfqpoint{5.251159in}{2.298995in}}%
\pgfpathlineto{\pgfqpoint{5.251202in}{2.298270in}}%
\pgfpathlineto{\pgfqpoint{5.251826in}{2.301637in}}%
\pgfpathlineto{\pgfqpoint{5.251913in}{2.301211in}}%
\pgfpathlineto{\pgfqpoint{5.251972in}{2.301990in}}%
\pgfpathlineto{\pgfqpoint{5.252794in}{2.297082in}}%
\pgfpathlineto{\pgfqpoint{5.252989in}{2.298332in}}%
\pgfpathlineto{\pgfqpoint{5.253315in}{2.297106in}}%
\pgfpathlineto{\pgfqpoint{5.253997in}{2.300936in}}%
\pgfpathlineto{\pgfqpoint{5.254002in}{2.300769in}}%
\pgfpathlineto{\pgfqpoint{5.254060in}{2.299928in}}%
\pgfpathlineto{\pgfqpoint{5.255214in}{2.304451in}}%
\pgfpathlineto{\pgfqpoint{5.255686in}{2.300200in}}%
\pgfpathlineto{\pgfqpoint{5.256339in}{2.303136in}}%
\pgfpathlineto{\pgfqpoint{5.257215in}{2.305415in}}%
\pgfpathlineto{\pgfqpoint{5.256694in}{2.302154in}}%
\pgfpathlineto{\pgfqpoint{5.257488in}{2.304932in}}%
\pgfpathlineto{\pgfqpoint{5.258495in}{2.303022in}}%
\pgfpathlineto{\pgfqpoint{5.258091in}{2.306903in}}%
\pgfpathlineto{\pgfqpoint{5.258656in}{2.303426in}}%
\pgfpathlineto{\pgfqpoint{5.259747in}{2.300760in}}%
\pgfpathlineto{\pgfqpoint{5.258690in}{2.303684in}}%
\pgfpathlineto{\pgfqpoint{5.259893in}{2.301754in}}%
\pgfpathlineto{\pgfqpoint{5.260282in}{2.302956in}}%
\pgfpathlineto{\pgfqpoint{5.260944in}{2.298547in}}%
\pgfpathlineto{\pgfqpoint{5.261616in}{2.295132in}}%
\pgfpathlineto{\pgfqpoint{5.262103in}{2.297123in}}%
\pgfpathlineto{\pgfqpoint{5.262390in}{2.298312in}}%
\pgfpathlineto{\pgfqpoint{5.263125in}{2.295723in}}%
\pgfpathlineto{\pgfqpoint{5.263155in}{2.295803in}}%
\pgfpathlineto{\pgfqpoint{5.263938in}{2.292893in}}%
\pgfpathlineto{\pgfqpoint{5.263169in}{2.295811in}}%
\pgfpathlineto{\pgfqpoint{5.264299in}{2.294051in}}%
\pgfpathlineto{\pgfqpoint{5.265160in}{2.296180in}}%
\pgfpathlineto{\pgfqpoint{5.264450in}{2.292519in}}%
\pgfpathlineto{\pgfqpoint{5.265336in}{2.293305in}}%
\pgfpathlineto{\pgfqpoint{5.266684in}{2.289259in}}%
\pgfpathlineto{\pgfqpoint{5.265993in}{2.294792in}}%
\pgfpathlineto{\pgfqpoint{5.266699in}{2.289539in}}%
\pgfpathlineto{\pgfqpoint{5.267171in}{2.291043in}}%
\pgfpathlineto{\pgfqpoint{5.266796in}{2.288875in}}%
\pgfpathlineto{\pgfqpoint{5.267176in}{2.291031in}}%
\pgfpathlineto{\pgfqpoint{5.268213in}{2.292739in}}%
\pgfpathlineto{\pgfqpoint{5.267746in}{2.289069in}}%
\pgfpathlineto{\pgfqpoint{5.268305in}{2.292031in}}%
\pgfpathlineto{\pgfqpoint{5.269196in}{2.290270in}}%
\pgfpathlineto{\pgfqpoint{5.268787in}{2.292935in}}%
\pgfpathlineto{\pgfqpoint{5.269396in}{2.292619in}}%
\pgfpathlineto{\pgfqpoint{5.269552in}{2.293113in}}%
\pgfpathlineto{\pgfqpoint{5.270141in}{2.290444in}}%
\pgfpathlineto{\pgfqpoint{5.270477in}{2.291376in}}%
\pgfpathlineto{\pgfqpoint{5.270886in}{2.290453in}}%
\pgfpathlineto{\pgfqpoint{5.271460in}{2.294294in}}%
\pgfpathlineto{\pgfqpoint{5.271528in}{2.293604in}}%
\pgfpathlineto{\pgfqpoint{5.271845in}{2.289662in}}%
\pgfpathlineto{\pgfqpoint{5.272556in}{2.294230in}}%
\pgfpathlineto{\pgfqpoint{5.272609in}{2.293761in}}%
\pgfpathlineto{\pgfqpoint{5.273778in}{2.298167in}}%
\pgfpathlineto{\pgfqpoint{5.272760in}{2.293441in}}%
\pgfpathlineto{\pgfqpoint{5.273807in}{2.298070in}}%
\pgfpathlineto{\pgfqpoint{5.274552in}{2.295331in}}%
\pgfpathlineto{\pgfqpoint{5.274279in}{2.298785in}}%
\pgfpathlineto{\pgfqpoint{5.274966in}{2.295750in}}%
\pgfpathlineto{\pgfqpoint{5.275043in}{2.296333in}}%
\pgfpathlineto{\pgfqpoint{5.275248in}{2.294050in}}%
\pgfpathlineto{\pgfqpoint{5.275594in}{2.294335in}}%
\pgfpathlineto{\pgfqpoint{5.275725in}{2.293722in}}%
\pgfpathlineto{\pgfqpoint{5.276577in}{2.297515in}}%
\pgfpathlineto{\pgfqpoint{5.276650in}{2.296699in}}%
\pgfpathlineto{\pgfqpoint{5.277814in}{2.292471in}}%
\pgfpathlineto{\pgfqpoint{5.277867in}{2.293382in}}%
\pgfpathlineto{\pgfqpoint{5.277994in}{2.292435in}}%
\pgfpathlineto{\pgfqpoint{5.278057in}{2.291381in}}%
\pgfpathlineto{\pgfqpoint{5.279055in}{2.295249in}}%
\pgfpathlineto{\pgfqpoint{5.279513in}{2.298714in}}%
\pgfpathlineto{\pgfqpoint{5.279099in}{2.294643in}}%
\pgfpathlineto{\pgfqpoint{5.280185in}{2.295291in}}%
\pgfpathlineto{\pgfqpoint{5.281353in}{2.290384in}}%
\pgfpathlineto{\pgfqpoint{5.280535in}{2.295887in}}%
\pgfpathlineto{\pgfqpoint{5.281562in}{2.291816in}}%
\pgfpathlineto{\pgfqpoint{5.282570in}{2.295327in}}%
\pgfpathlineto{\pgfqpoint{5.281596in}{2.291253in}}%
\pgfpathlineto{\pgfqpoint{5.282736in}{2.293702in}}%
\pgfpathlineto{\pgfqpoint{5.283373in}{2.291054in}}%
\pgfpathlineto{\pgfqpoint{5.283924in}{2.292026in}}%
\pgfpathlineto{\pgfqpoint{5.284070in}{2.293062in}}%
\pgfpathlineto{\pgfqpoint{5.284907in}{2.285680in}}%
\pgfpathlineto{\pgfqpoint{5.284912in}{2.285643in}}%
\pgfpathlineto{\pgfqpoint{5.285569in}{2.288694in}}%
\pgfpathlineto{\pgfqpoint{5.285657in}{2.288328in}}%
\pgfpathlineto{\pgfqpoint{5.285988in}{2.289143in}}%
\pgfpathlineto{\pgfqpoint{5.286455in}{2.286211in}}%
\pgfpathlineto{\pgfqpoint{5.286767in}{2.288302in}}%
\pgfpathlineto{\pgfqpoint{5.288106in}{2.295510in}}%
\pgfpathlineto{\pgfqpoint{5.286845in}{2.287665in}}%
\pgfpathlineto{\pgfqpoint{5.288125in}{2.295291in}}%
\pgfpathlineto{\pgfqpoint{5.288325in}{2.294711in}}%
\pgfpathlineto{\pgfqpoint{5.289186in}{2.298062in}}%
\pgfpathlineto{\pgfqpoint{5.289961in}{2.298919in}}%
\pgfpathlineto{\pgfqpoint{5.289318in}{2.296603in}}%
\pgfpathlineto{\pgfqpoint{5.290228in}{2.297665in}}%
\pgfpathlineto{\pgfqpoint{5.290238in}{2.297113in}}%
\pgfpathlineto{\pgfqpoint{5.291265in}{2.299601in}}%
\pgfpathlineto{\pgfqpoint{5.291324in}{2.298574in}}%
\pgfpathlineto{\pgfqpoint{5.292380in}{2.299062in}}%
\pgfpathlineto{\pgfqpoint{5.292088in}{2.295740in}}%
\pgfpathlineto{\pgfqpoint{5.292434in}{2.298709in}}%
\pgfpathlineto{\pgfqpoint{5.293588in}{2.296540in}}%
\pgfpathlineto{\pgfqpoint{5.292969in}{2.299408in}}%
\pgfpathlineto{\pgfqpoint{5.293592in}{2.296635in}}%
\pgfpathlineto{\pgfqpoint{5.294917in}{2.302352in}}%
\pgfpathlineto{\pgfqpoint{5.293919in}{2.294930in}}%
\pgfpathlineto{\pgfqpoint{5.294931in}{2.302320in}}%
\pgfpathlineto{\pgfqpoint{5.294951in}{2.302083in}}%
\pgfpathlineto{\pgfqpoint{5.295535in}{2.304916in}}%
\pgfpathlineto{\pgfqpoint{5.295959in}{2.303774in}}%
\pgfpathlineto{\pgfqpoint{5.297069in}{2.304156in}}%
\pgfpathlineto{\pgfqpoint{5.296630in}{2.300157in}}%
\pgfpathlineto{\pgfqpoint{5.297073in}{2.303990in}}%
\pgfpathlineto{\pgfqpoint{5.298062in}{2.302543in}}%
\pgfpathlineto{\pgfqpoint{5.297828in}{2.305482in}}%
\pgfpathlineto{\pgfqpoint{5.298169in}{2.304296in}}%
\pgfpathlineto{\pgfqpoint{5.298330in}{2.304864in}}%
\pgfpathlineto{\pgfqpoint{5.299211in}{2.302343in}}%
\pgfpathlineto{\pgfqpoint{5.299493in}{2.301658in}}%
\pgfpathlineto{\pgfqpoint{5.299761in}{2.305393in}}%
\pgfpathlineto{\pgfqpoint{5.300306in}{2.303524in}}%
\pgfpathlineto{\pgfqpoint{5.300418in}{2.304051in}}%
\pgfpathlineto{\pgfqpoint{5.300881in}{2.301691in}}%
\pgfpathlineto{\pgfqpoint{5.301144in}{2.299596in}}%
\pgfpathlineto{\pgfqpoint{5.301937in}{2.303216in}}%
\pgfpathlineto{\pgfqpoint{5.301942in}{2.303388in}}%
\pgfpathlineto{\pgfqpoint{5.302575in}{2.299223in}}%
\pgfpathlineto{\pgfqpoint{5.302945in}{2.299688in}}%
\pgfpathlineto{\pgfqpoint{5.303023in}{2.298908in}}%
\pgfpathlineto{\pgfqpoint{5.303880in}{2.301574in}}%
\pgfpathlineto{\pgfqpoint{5.303962in}{2.300873in}}%
\pgfpathlineto{\pgfqpoint{5.304736in}{2.304164in}}%
\pgfpathlineto{\pgfqpoint{5.304439in}{2.300551in}}%
\pgfpathlineto{\pgfqpoint{5.305106in}{2.302440in}}%
\pgfpathlineto{\pgfqpoint{5.305520in}{2.298920in}}%
\pgfpathlineto{\pgfqpoint{5.305204in}{2.302854in}}%
\pgfpathlineto{\pgfqpoint{5.306275in}{2.299662in}}%
\pgfpathlineto{\pgfqpoint{5.306528in}{2.302047in}}%
\pgfpathlineto{\pgfqpoint{5.307395in}{2.300573in}}%
\pgfpathlineto{\pgfqpoint{5.307950in}{2.296214in}}%
\pgfpathlineto{\pgfqpoint{5.307414in}{2.300608in}}%
\pgfpathlineto{\pgfqpoint{5.308568in}{2.297566in}}%
\pgfpathlineto{\pgfqpoint{5.309566in}{2.302059in}}%
\pgfpathlineto{\pgfqpoint{5.309736in}{2.300504in}}%
\pgfpathlineto{\pgfqpoint{5.310637in}{2.299784in}}%
\pgfpathlineto{\pgfqpoint{5.310355in}{2.302344in}}%
\pgfpathlineto{\pgfqpoint{5.310837in}{2.301245in}}%
\pgfpathlineto{\pgfqpoint{5.311650in}{2.297417in}}%
\pgfpathlineto{\pgfqpoint{5.310876in}{2.301928in}}%
\pgfpathlineto{\pgfqpoint{5.311971in}{2.300125in}}%
\pgfpathlineto{\pgfqpoint{5.312312in}{2.303668in}}%
\pgfpathlineto{\pgfqpoint{5.313096in}{2.301018in}}%
\pgfpathlineto{\pgfqpoint{5.313982in}{2.304187in}}%
\pgfpathlineto{\pgfqpoint{5.313383in}{2.299564in}}%
\pgfpathlineto{\pgfqpoint{5.314381in}{2.302554in}}%
\pgfpathlineto{\pgfqpoint{5.314805in}{2.301716in}}%
\pgfpathlineto{\pgfqpoint{5.314595in}{2.304073in}}%
\pgfpathlineto{\pgfqpoint{5.315413in}{2.303328in}}%
\pgfpathlineto{\pgfqpoint{5.316474in}{2.303971in}}%
\pgfpathlineto{\pgfqpoint{5.316246in}{2.301716in}}%
\pgfpathlineto{\pgfqpoint{5.316557in}{2.303696in}}%
\pgfpathlineto{\pgfqpoint{5.316713in}{2.304035in}}%
\pgfpathlineto{\pgfqpoint{5.316990in}{2.302219in}}%
\pgfpathlineto{\pgfqpoint{5.317453in}{2.302752in}}%
\pgfpathlineto{\pgfqpoint{5.317463in}{2.302187in}}%
\pgfpathlineto{\pgfqpoint{5.317993in}{2.305489in}}%
\pgfpathlineto{\pgfqpoint{5.318563in}{2.302724in}}%
\pgfpathlineto{\pgfqpoint{5.319668in}{2.304724in}}%
\pgfpathlineto{\pgfqpoint{5.319060in}{2.300864in}}%
\pgfpathlineto{\pgfqpoint{5.319770in}{2.303310in}}%
\pgfpathlineto{\pgfqpoint{5.319892in}{2.301155in}}%
\pgfpathlineto{\pgfqpoint{5.320428in}{2.303555in}}%
\pgfpathlineto{\pgfqpoint{5.320890in}{2.302842in}}%
\pgfpathlineto{\pgfqpoint{5.320895in}{2.302855in}}%
\pgfpathlineto{\pgfqpoint{5.321304in}{2.301047in}}%
\pgfpathlineto{\pgfqpoint{5.321362in}{2.301410in}}%
\pgfpathlineto{\pgfqpoint{5.322283in}{2.297541in}}%
\pgfpathlineto{\pgfqpoint{5.322487in}{2.299792in}}%
\pgfpathlineto{\pgfqpoint{5.323446in}{2.303793in}}%
\pgfpathlineto{\pgfqpoint{5.322779in}{2.298996in}}%
\pgfpathlineto{\pgfqpoint{5.323826in}{2.301598in}}%
\pgfpathlineto{\pgfqpoint{5.324868in}{2.297720in}}%
\pgfpathlineto{\pgfqpoint{5.323860in}{2.301680in}}%
\pgfpathlineto{\pgfqpoint{5.324989in}{2.299505in}}%
\pgfpathlineto{\pgfqpoint{5.326031in}{2.306045in}}%
\pgfpathlineto{\pgfqpoint{5.325033in}{2.299117in}}%
\pgfpathlineto{\pgfqpoint{5.326616in}{2.303963in}}%
\pgfpathlineto{\pgfqpoint{5.327044in}{2.300588in}}%
\pgfpathlineto{\pgfqpoint{5.326762in}{2.304084in}}%
\pgfpathlineto{\pgfqpoint{5.327745in}{2.303486in}}%
\pgfpathlineto{\pgfqpoint{5.328826in}{2.304872in}}%
\pgfpathlineto{\pgfqpoint{5.328446in}{2.300395in}}%
\pgfpathlineto{\pgfqpoint{5.328860in}{2.304223in}}%
\pgfpathlineto{\pgfqpoint{5.329103in}{2.302209in}}%
\pgfpathlineto{\pgfqpoint{5.329941in}{2.306004in}}%
\pgfpathlineto{\pgfqpoint{5.330837in}{2.302009in}}%
\pgfpathlineto{\pgfqpoint{5.330072in}{2.306453in}}%
\pgfpathlineto{\pgfqpoint{5.331275in}{2.303830in}}%
\pgfpathlineto{\pgfqpoint{5.332195in}{2.299569in}}%
\pgfpathlineto{\pgfqpoint{5.331494in}{2.304164in}}%
\pgfpathlineto{\pgfqpoint{5.332604in}{2.303180in}}%
\pgfpathlineto{\pgfqpoint{5.332910in}{2.304613in}}%
\pgfpathlineto{\pgfqpoint{5.333354in}{2.300341in}}%
\pgfpathlineto{\pgfqpoint{5.333680in}{2.301161in}}%
\pgfpathlineto{\pgfqpoint{5.334269in}{2.303064in}}%
\pgfpathlineto{\pgfqpoint{5.333748in}{2.300812in}}%
\pgfpathlineto{\pgfqpoint{5.334775in}{2.301084in}}%
\pgfpathlineto{\pgfqpoint{5.335700in}{2.298655in}}%
\pgfpathlineto{\pgfqpoint{5.334989in}{2.302430in}}%
\pgfpathlineto{\pgfqpoint{5.335909in}{2.299687in}}%
\pgfpathlineto{\pgfqpoint{5.336981in}{2.304078in}}%
\pgfpathlineto{\pgfqpoint{5.336168in}{2.297470in}}%
\pgfpathlineto{\pgfqpoint{5.337161in}{2.303645in}}%
\pgfpathlineto{\pgfqpoint{5.337501in}{2.301871in}}%
\pgfpathlineto{\pgfqpoint{5.337229in}{2.304631in}}%
\pgfpathlineto{\pgfqpoint{5.338276in}{2.303047in}}%
\pgfpathlineto{\pgfqpoint{5.338923in}{2.303556in}}%
\pgfpathlineto{\pgfqpoint{5.339230in}{2.300334in}}%
\pgfpathlineto{\pgfqpoint{5.339361in}{2.301585in}}%
\pgfpathlineto{\pgfqpoint{5.340009in}{2.304340in}}%
\pgfpathlineto{\pgfqpoint{5.339473in}{2.301219in}}%
\pgfpathlineto{\pgfqpoint{5.340520in}{2.303738in}}%
\pgfpathlineto{\pgfqpoint{5.341260in}{2.305894in}}%
\pgfpathlineto{\pgfqpoint{5.341192in}{2.302997in}}%
\pgfpathlineto{\pgfqpoint{5.341567in}{2.303319in}}%
\pgfpathlineto{\pgfqpoint{5.342730in}{2.298360in}}%
\pgfpathlineto{\pgfqpoint{5.342784in}{2.298503in}}%
\pgfpathlineto{\pgfqpoint{5.343110in}{2.300675in}}%
\pgfpathlineto{\pgfqpoint{5.342798in}{2.298386in}}%
\pgfpathlineto{\pgfqpoint{5.343957in}{2.299528in}}%
\pgfpathlineto{\pgfqpoint{5.343967in}{2.299008in}}%
\pgfpathlineto{\pgfqpoint{5.344853in}{2.301326in}}%
\pgfpathlineto{\pgfqpoint{5.345057in}{2.300142in}}%
\pgfpathlineto{\pgfqpoint{5.345077in}{2.300785in}}%
\pgfpathlineto{\pgfqpoint{5.346012in}{2.296884in}}%
\pgfpathlineto{\pgfqpoint{5.346094in}{2.297487in}}%
\pgfpathlineto{\pgfqpoint{5.347273in}{2.293121in}}%
\pgfpathlineto{\pgfqpoint{5.346279in}{2.298621in}}%
\pgfpathlineto{\pgfqpoint{5.347277in}{2.293248in}}%
\pgfpathlineto{\pgfqpoint{5.348426in}{2.297427in}}%
\pgfpathlineto{\pgfqpoint{5.347375in}{2.292574in}}%
\pgfpathlineto{\pgfqpoint{5.348436in}{2.297028in}}%
\pgfpathlineto{\pgfqpoint{5.349035in}{2.296005in}}%
\pgfpathlineto{\pgfqpoint{5.348719in}{2.299237in}}%
\pgfpathlineto{\pgfqpoint{5.349546in}{2.296799in}}%
\pgfpathlineto{\pgfqpoint{5.350515in}{2.298392in}}%
\pgfpathlineto{\pgfqpoint{5.350111in}{2.294466in}}%
\pgfpathlineto{\pgfqpoint{5.350671in}{2.297384in}}%
\pgfpathlineto{\pgfqpoint{5.350778in}{2.295600in}}%
\pgfpathlineto{\pgfqpoint{5.351581in}{2.302443in}}%
\pgfpathlineto{\pgfqpoint{5.351625in}{2.302374in}}%
\pgfpathlineto{\pgfqpoint{5.351713in}{2.303112in}}%
\pgfpathlineto{\pgfqpoint{5.351976in}{2.300628in}}%
\pgfpathlineto{\pgfqpoint{5.352423in}{2.300572in}}%
\pgfpathlineto{\pgfqpoint{5.352808in}{2.297572in}}%
\pgfpathlineto{\pgfqpoint{5.353524in}{2.300889in}}%
\pgfpathlineto{\pgfqpoint{5.353903in}{2.300699in}}%
\pgfpathlineto{\pgfqpoint{5.353587in}{2.301509in}}%
\pgfpathlineto{\pgfqpoint{5.353967in}{2.301753in}}%
\pgfpathlineto{\pgfqpoint{5.354687in}{2.304149in}}%
\pgfpathlineto{\pgfqpoint{5.355009in}{2.299445in}}%
\pgfpathlineto{\pgfqpoint{5.355043in}{2.299684in}}%
\pgfpathlineto{\pgfqpoint{5.356240in}{2.304477in}}%
\pgfpathlineto{\pgfqpoint{5.355052in}{2.299581in}}%
\pgfpathlineto{\pgfqpoint{5.356309in}{2.302954in}}%
\pgfpathlineto{\pgfqpoint{5.356333in}{2.302539in}}%
\pgfpathlineto{\pgfqpoint{5.356640in}{2.304257in}}%
\pgfpathlineto{\pgfqpoint{5.357307in}{2.303590in}}%
\pgfpathlineto{\pgfqpoint{5.357803in}{2.305995in}}%
\pgfpathlineto{\pgfqpoint{5.358266in}{2.302127in}}%
\pgfpathlineto{\pgfqpoint{5.358363in}{2.302435in}}%
\pgfpathlineto{\pgfqpoint{5.359361in}{2.301006in}}%
\pgfpathlineto{\pgfqpoint{5.358733in}{2.305046in}}%
\pgfpathlineto{\pgfqpoint{5.359468in}{2.302673in}}%
\pgfpathlineto{\pgfqpoint{5.360471in}{2.304133in}}%
\pgfpathlineto{\pgfqpoint{5.359541in}{2.302384in}}%
\pgfpathlineto{\pgfqpoint{5.360593in}{2.303647in}}%
\pgfpathlineto{\pgfqpoint{5.361128in}{2.304612in}}%
\pgfpathlineto{\pgfqpoint{5.361766in}{2.299568in}}%
\pgfpathlineto{\pgfqpoint{5.362949in}{2.304497in}}%
\pgfpathlineto{\pgfqpoint{5.362993in}{2.304238in}}%
\pgfpathlineto{\pgfqpoint{5.364049in}{2.299558in}}%
\pgfpathlineto{\pgfqpoint{5.364191in}{2.300577in}}%
\pgfpathlineto{\pgfqpoint{5.364210in}{2.300321in}}%
\pgfpathlineto{\pgfqpoint{5.364278in}{2.301625in}}%
\pgfpathlineto{\pgfqpoint{5.365349in}{2.303308in}}%
\pgfpathlineto{\pgfqpoint{5.365023in}{2.300267in}}%
\pgfpathlineto{\pgfqpoint{5.365408in}{2.302485in}}%
\pgfpathlineto{\pgfqpoint{5.366211in}{2.300630in}}%
\pgfpathlineto{\pgfqpoint{5.365534in}{2.303024in}}%
\pgfpathlineto{\pgfqpoint{5.366523in}{2.302293in}}%
\pgfpathlineto{\pgfqpoint{5.367379in}{2.304639in}}%
\pgfpathlineto{\pgfqpoint{5.366591in}{2.302030in}}%
\pgfpathlineto{\pgfqpoint{5.367749in}{2.303423in}}%
\pgfpathlineto{\pgfqpoint{5.367905in}{2.302809in}}%
\pgfpathlineto{\pgfqpoint{5.368728in}{2.304813in}}%
\pgfpathlineto{\pgfqpoint{5.368840in}{2.304079in}}%
\pgfpathlineto{\pgfqpoint{5.369502in}{2.302578in}}%
\pgfpathlineto{\pgfqpoint{5.368894in}{2.304672in}}%
\pgfpathlineto{\pgfqpoint{5.369785in}{2.304683in}}%
\pgfpathlineto{\pgfqpoint{5.369799in}{2.305020in}}%
\pgfpathlineto{\pgfqpoint{5.370783in}{2.299614in}}%
\pgfpathlineto{\pgfqpoint{5.370787in}{2.299660in}}%
\pgfpathlineto{\pgfqpoint{5.371041in}{2.298415in}}%
\pgfpathlineto{\pgfqpoint{5.371537in}{2.301280in}}%
\pgfpathlineto{\pgfqpoint{5.371820in}{2.300046in}}%
\pgfpathlineto{\pgfqpoint{5.372891in}{2.301860in}}%
\pgfpathlineto{\pgfqpoint{5.372170in}{2.299618in}}%
\pgfpathlineto{\pgfqpoint{5.372949in}{2.301638in}}%
\pgfpathlineto{\pgfqpoint{5.373962in}{2.299248in}}%
\pgfpathlineto{\pgfqpoint{5.373499in}{2.302375in}}%
\pgfpathlineto{\pgfqpoint{5.374079in}{2.300876in}}%
\pgfpathlineto{\pgfqpoint{5.374361in}{2.301664in}}%
\pgfpathlineto{\pgfqpoint{5.374721in}{2.299039in}}%
\pgfpathlineto{\pgfqpoint{5.375159in}{2.300413in}}%
\pgfpathlineto{\pgfqpoint{5.375262in}{2.298822in}}%
\pgfpathlineto{\pgfqpoint{5.376123in}{2.303110in}}%
\pgfpathlineto{\pgfqpoint{5.376216in}{2.302704in}}%
\pgfpathlineto{\pgfqpoint{5.376347in}{2.304139in}}%
\pgfpathlineto{\pgfqpoint{5.377194in}{2.298613in}}%
\pgfpathlineto{\pgfqpoint{5.377204in}{2.298868in}}%
\pgfpathlineto{\pgfqpoint{5.377443in}{2.297750in}}%
\pgfpathlineto{\pgfqpoint{5.377686in}{2.300101in}}%
\pgfpathlineto{\pgfqpoint{5.378256in}{2.299599in}}%
\pgfpathlineto{\pgfqpoint{5.379556in}{2.297237in}}%
\pgfpathlineto{\pgfqpoint{5.378684in}{2.300007in}}%
\pgfpathlineto{\pgfqpoint{5.379595in}{2.298034in}}%
\pgfpathlineto{\pgfqpoint{5.380159in}{2.299113in}}%
\pgfpathlineto{\pgfqpoint{5.380456in}{2.295930in}}%
\pgfpathlineto{\pgfqpoint{5.380700in}{2.297800in}}%
\pgfpathlineto{\pgfqpoint{5.381211in}{2.300215in}}%
\pgfpathlineto{\pgfqpoint{5.381299in}{2.299329in}}%
\pgfpathlineto{\pgfqpoint{5.382448in}{2.304131in}}%
\pgfpathlineto{\pgfqpoint{5.383334in}{2.298772in}}%
\pgfpathlineto{\pgfqpoint{5.382764in}{2.304412in}}%
\pgfpathlineto{\pgfqpoint{5.383606in}{2.301306in}}%
\pgfpathlineto{\pgfqpoint{5.383908in}{2.303554in}}%
\pgfpathlineto{\pgfqpoint{5.384633in}{2.300332in}}%
\pgfpathlineto{\pgfqpoint{5.384706in}{2.300709in}}%
\pgfpathlineto{\pgfqpoint{5.385106in}{2.304116in}}%
\pgfpathlineto{\pgfqpoint{5.385408in}{2.299642in}}%
\pgfpathlineto{\pgfqpoint{5.385875in}{2.301850in}}%
\pgfpathlineto{\pgfqpoint{5.387004in}{2.299442in}}%
\pgfpathlineto{\pgfqpoint{5.386182in}{2.302942in}}%
\pgfpathlineto{\pgfqpoint{5.387009in}{2.299606in}}%
\pgfpathlineto{\pgfqpoint{5.387175in}{2.301063in}}%
\pgfpathlineto{\pgfqpoint{5.387584in}{2.295869in}}%
\pgfpathlineto{\pgfqpoint{5.387910in}{2.296622in}}%
\pgfpathlineto{\pgfqpoint{5.388514in}{2.295032in}}%
\pgfpathlineto{\pgfqpoint{5.388913in}{2.298808in}}%
\pgfpathlineto{\pgfqpoint{5.388991in}{2.297875in}}%
\pgfpathlineto{\pgfqpoint{5.389448in}{2.296198in}}%
\pgfpathlineto{\pgfqpoint{5.389147in}{2.299220in}}%
\pgfpathlineto{\pgfqpoint{5.390018in}{2.298412in}}%
\pgfpathlineto{\pgfqpoint{5.391240in}{2.303523in}}%
\pgfpathlineto{\pgfqpoint{5.390033in}{2.298370in}}%
\pgfpathlineto{\pgfqpoint{5.391289in}{2.303275in}}%
\pgfpathlineto{\pgfqpoint{5.391985in}{2.302516in}}%
\pgfpathlineto{\pgfqpoint{5.391464in}{2.304533in}}%
\pgfpathlineto{\pgfqpoint{5.392394in}{2.303412in}}%
\pgfpathlineto{\pgfqpoint{5.393509in}{2.300114in}}%
\pgfpathlineto{\pgfqpoint{5.392486in}{2.303774in}}%
\pgfpathlineto{\pgfqpoint{5.393630in}{2.301937in}}%
\pgfpathlineto{\pgfqpoint{5.394317in}{2.306052in}}%
\pgfpathlineto{\pgfqpoint{5.393791in}{2.301917in}}%
\pgfpathlineto{\pgfqpoint{5.394789in}{2.303546in}}%
\pgfpathlineto{\pgfqpoint{5.395919in}{2.301534in}}%
\pgfpathlineto{\pgfqpoint{5.394887in}{2.304167in}}%
\pgfpathlineto{\pgfqpoint{5.395967in}{2.301678in}}%
\pgfpathlineto{\pgfqpoint{5.397097in}{2.304278in}}%
\pgfpathlineto{\pgfqpoint{5.395987in}{2.301415in}}%
\pgfpathlineto{\pgfqpoint{5.397180in}{2.303491in}}%
\pgfpathlineto{\pgfqpoint{5.397725in}{2.299071in}}%
\pgfpathlineto{\pgfqpoint{5.397219in}{2.303973in}}%
\pgfpathlineto{\pgfqpoint{5.398314in}{2.302117in}}%
\pgfpathlineto{\pgfqpoint{5.398319in}{2.302472in}}%
\pgfpathlineto{\pgfqpoint{5.398952in}{2.299101in}}%
\pgfpathlineto{\pgfqpoint{5.399400in}{2.300319in}}%
\pgfpathlineto{\pgfqpoint{5.400378in}{2.293574in}}%
\pgfpathlineto{\pgfqpoint{5.400602in}{2.294530in}}%
\pgfpathlineto{\pgfqpoint{5.401722in}{2.296958in}}%
\pgfpathlineto{\pgfqpoint{5.401143in}{2.293167in}}%
\pgfpathlineto{\pgfqpoint{5.401751in}{2.296857in}}%
\pgfpathlineto{\pgfqpoint{5.402949in}{2.290943in}}%
\pgfpathlineto{\pgfqpoint{5.401951in}{2.297535in}}%
\pgfpathlineto{\pgfqpoint{5.402978in}{2.291343in}}%
\pgfpathlineto{\pgfqpoint{5.403173in}{2.293501in}}%
\pgfpathlineto{\pgfqpoint{5.403786in}{2.287892in}}%
\pgfpathlineto{\pgfqpoint{5.404030in}{2.287994in}}%
\pgfpathlineto{\pgfqpoint{5.404151in}{2.287455in}}%
\pgfpathlineto{\pgfqpoint{5.404706in}{2.293688in}}%
\pgfpathlineto{\pgfqpoint{5.404979in}{2.291679in}}%
\pgfpathlineto{\pgfqpoint{5.405471in}{2.296161in}}%
\pgfpathlineto{\pgfqpoint{5.406094in}{2.292161in}}%
\pgfpathlineto{\pgfqpoint{5.406274in}{2.291841in}}%
\pgfpathlineto{\pgfqpoint{5.406960in}{2.295670in}}%
\pgfpathlineto{\pgfqpoint{5.407063in}{2.294950in}}%
\pgfpathlineto{\pgfqpoint{5.407710in}{2.299026in}}%
\pgfpathlineto{\pgfqpoint{5.407102in}{2.294588in}}%
\pgfpathlineto{\pgfqpoint{5.408450in}{2.296147in}}%
\pgfpathlineto{\pgfqpoint{5.408455in}{2.295963in}}%
\pgfpathlineto{\pgfqpoint{5.408976in}{2.298746in}}%
\pgfpathlineto{\pgfqpoint{5.409531in}{2.297287in}}%
\pgfpathlineto{\pgfqpoint{5.410378in}{2.300456in}}%
\pgfpathlineto{\pgfqpoint{5.410057in}{2.296446in}}%
\pgfpathlineto{\pgfqpoint{5.410656in}{2.298265in}}%
\pgfpathlineto{\pgfqpoint{5.411415in}{2.294884in}}%
\pgfpathlineto{\pgfqpoint{5.410724in}{2.299125in}}%
\pgfpathlineto{\pgfqpoint{5.411790in}{2.296652in}}%
\pgfpathlineto{\pgfqpoint{5.413002in}{2.299110in}}%
\pgfpathlineto{\pgfqpoint{5.412345in}{2.293669in}}%
\pgfpathlineto{\pgfqpoint{5.413027in}{2.298528in}}%
\pgfpathlineto{\pgfqpoint{5.414151in}{2.294020in}}%
\pgfpathlineto{\pgfqpoint{5.413085in}{2.299608in}}%
\pgfpathlineto{\pgfqpoint{5.414244in}{2.294150in}}%
\pgfpathlineto{\pgfqpoint{5.415573in}{2.300948in}}%
\pgfpathlineto{\pgfqpoint{5.415583in}{2.300511in}}%
\pgfpathlineto{\pgfqpoint{5.415607in}{2.299754in}}%
\pgfpathlineto{\pgfqpoint{5.416284in}{2.304041in}}%
\pgfpathlineto{\pgfqpoint{5.416508in}{2.303099in}}%
\pgfpathlineto{\pgfqpoint{5.417145in}{2.306632in}}%
\pgfpathlineto{\pgfqpoint{5.417637in}{2.303435in}}%
\pgfpathlineto{\pgfqpoint{5.418095in}{2.302050in}}%
\pgfpathlineto{\pgfqpoint{5.418508in}{2.304809in}}%
\pgfpathlineto{\pgfqpoint{5.418732in}{2.303718in}}%
\pgfpathlineto{\pgfqpoint{5.419049in}{2.304820in}}%
\pgfpathlineto{\pgfqpoint{5.419643in}{2.301628in}}%
\pgfpathlineto{\pgfqpoint{5.419833in}{2.303235in}}%
\pgfpathlineto{\pgfqpoint{5.419852in}{2.302948in}}%
\pgfpathlineto{\pgfqpoint{5.419901in}{2.304360in}}%
\pgfpathlineto{\pgfqpoint{5.419930in}{2.304762in}}%
\pgfpathlineto{\pgfqpoint{5.420675in}{2.302031in}}%
\pgfpathlineto{\pgfqpoint{5.421001in}{2.303770in}}%
\pgfpathlineto{\pgfqpoint{5.421157in}{2.302418in}}%
\pgfpathlineto{\pgfqpoint{5.421994in}{2.304378in}}%
\pgfpathlineto{\pgfqpoint{5.422097in}{2.304046in}}%
\pgfpathlineto{\pgfqpoint{5.422428in}{2.305368in}}%
\pgfpathlineto{\pgfqpoint{5.422209in}{2.302849in}}%
\pgfpathlineto{\pgfqpoint{5.422851in}{2.303379in}}%
\pgfpathlineto{\pgfqpoint{5.423114in}{2.302401in}}%
\pgfpathlineto{\pgfqpoint{5.423338in}{2.304458in}}%
\pgfpathlineto{\pgfqpoint{5.423966in}{2.302928in}}%
\pgfpathlineto{\pgfqpoint{5.424136in}{2.303770in}}%
\pgfpathlineto{\pgfqpoint{5.424823in}{2.300652in}}%
\pgfpathlineto{\pgfqpoint{5.425066in}{2.302543in}}%
\pgfpathlineto{\pgfqpoint{5.425879in}{2.296146in}}%
\pgfpathlineto{\pgfqpoint{5.425110in}{2.302906in}}%
\pgfpathlineto{\pgfqpoint{5.426303in}{2.300103in}}%
\pgfpathlineto{\pgfqpoint{5.426313in}{2.300278in}}%
\pgfpathlineto{\pgfqpoint{5.427306in}{2.297334in}}%
\pgfpathlineto{\pgfqpoint{5.427340in}{2.297802in}}%
\pgfpathlineto{\pgfqpoint{5.428396in}{2.294617in}}%
\pgfpathlineto{\pgfqpoint{5.427559in}{2.298054in}}%
\pgfpathlineto{\pgfqpoint{5.428499in}{2.295321in}}%
\pgfpathlineto{\pgfqpoint{5.429726in}{2.304570in}}%
\pgfpathlineto{\pgfqpoint{5.428533in}{2.295186in}}%
\pgfpathlineto{\pgfqpoint{5.429760in}{2.303782in}}%
\pgfpathlineto{\pgfqpoint{5.429940in}{2.305075in}}%
\pgfpathlineto{\pgfqpoint{5.430626in}{2.301974in}}%
\pgfpathlineto{\pgfqpoint{5.430879in}{2.304336in}}%
\pgfpathlineto{\pgfqpoint{5.431722in}{2.302924in}}%
\pgfpathlineto{\pgfqpoint{5.431142in}{2.304763in}}%
\pgfpathlineto{\pgfqpoint{5.431994in}{2.303935in}}%
\pgfpathlineto{\pgfqpoint{5.432019in}{2.304371in}}%
\pgfpathlineto{\pgfqpoint{5.432987in}{2.300150in}}%
\pgfpathlineto{\pgfqpoint{5.433007in}{2.300550in}}%
\pgfpathlineto{\pgfqpoint{5.433903in}{2.305719in}}%
\pgfpathlineto{\pgfqpoint{5.433095in}{2.300109in}}%
\pgfpathlineto{\pgfqpoint{5.434380in}{2.303227in}}%
\pgfpathlineto{\pgfqpoint{5.434964in}{2.301635in}}%
\pgfpathlineto{\pgfqpoint{5.434667in}{2.304671in}}%
\pgfpathlineto{\pgfqpoint{5.435480in}{2.304128in}}%
\pgfpathlineto{\pgfqpoint{5.435655in}{2.303129in}}%
\pgfpathlineto{\pgfqpoint{5.436342in}{2.305840in}}%
\pgfpathlineto{\pgfqpoint{5.436571in}{2.304313in}}%
\pgfpathlineto{\pgfqpoint{5.436858in}{2.306852in}}%
\pgfpathlineto{\pgfqpoint{5.437223in}{2.302192in}}%
\pgfpathlineto{\pgfqpoint{5.437608in}{2.302543in}}%
\pgfpathlineto{\pgfqpoint{5.437968in}{2.300210in}}%
\pgfpathlineto{\pgfqpoint{5.438494in}{2.302781in}}%
\pgfpathlineto{\pgfqpoint{5.438747in}{2.300875in}}%
\pgfpathlineto{\pgfqpoint{5.439316in}{2.304282in}}%
\pgfpathlineto{\pgfqpoint{5.439039in}{2.300750in}}%
\pgfpathlineto{\pgfqpoint{5.439891in}{2.303738in}}%
\pgfpathlineto{\pgfqpoint{5.440782in}{2.302135in}}%
\pgfpathlineto{\pgfqpoint{5.440256in}{2.304129in}}%
\pgfpathlineto{\pgfqpoint{5.441025in}{2.302230in}}%
\pgfpathlineto{\pgfqpoint{5.441201in}{2.301047in}}%
\pgfpathlineto{\pgfqpoint{5.441605in}{2.303332in}}%
\pgfpathlineto{\pgfqpoint{5.441634in}{2.303022in}}%
\pgfpathlineto{\pgfqpoint{5.442744in}{2.304984in}}%
\pgfpathlineto{\pgfqpoint{5.442116in}{2.302072in}}%
\pgfpathlineto{\pgfqpoint{5.442768in}{2.304463in}}%
\pgfpathlineto{\pgfqpoint{5.442968in}{2.303376in}}%
\pgfpathlineto{\pgfqpoint{5.443226in}{2.305849in}}%
\pgfpathlineto{\pgfqpoint{5.443382in}{2.305707in}}%
\pgfpathlineto{\pgfqpoint{5.443445in}{2.306812in}}%
\pgfpathlineto{\pgfqpoint{5.444088in}{2.299732in}}%
\pgfpathlineto{\pgfqpoint{5.444229in}{2.300206in}}%
\pgfpathlineto{\pgfqpoint{5.444988in}{2.296384in}}%
\pgfpathlineto{\pgfqpoint{5.444482in}{2.300548in}}%
\pgfpathlineto{\pgfqpoint{5.445519in}{2.299505in}}%
\pgfpathlineto{\pgfqpoint{5.445913in}{2.301844in}}%
\pgfpathlineto{\pgfqpoint{5.446342in}{2.298259in}}%
\pgfpathlineto{\pgfqpoint{5.446629in}{2.299850in}}%
\pgfpathlineto{\pgfqpoint{5.446916in}{2.298514in}}%
\pgfpathlineto{\pgfqpoint{5.447500in}{2.301845in}}%
\pgfpathlineto{\pgfqpoint{5.447715in}{2.301153in}}%
\pgfpathlineto{\pgfqpoint{5.447817in}{2.300241in}}%
\pgfpathlineto{\pgfqpoint{5.448474in}{2.304133in}}%
\pgfpathlineto{\pgfqpoint{5.448659in}{2.303737in}}%
\pgfpathlineto{\pgfqpoint{5.448830in}{2.304753in}}%
\pgfpathlineto{\pgfqpoint{5.449346in}{2.302167in}}%
\pgfpathlineto{\pgfqpoint{5.449657in}{2.302690in}}%
\pgfpathlineto{\pgfqpoint{5.450436in}{2.301466in}}%
\pgfpathlineto{\pgfqpoint{5.450699in}{2.303662in}}%
\pgfpathlineto{\pgfqpoint{5.450704in}{2.303593in}}%
\pgfpathlineto{\pgfqpoint{5.450972in}{2.304641in}}%
\pgfpathlineto{\pgfqpoint{5.451376in}{2.302449in}}%
\pgfpathlineto{\pgfqpoint{5.451819in}{2.304278in}}%
\pgfpathlineto{\pgfqpoint{5.452101in}{2.302753in}}%
\pgfpathlineto{\pgfqpoint{5.452807in}{2.305447in}}%
\pgfpathlineto{\pgfqpoint{5.452812in}{2.305711in}}%
\pgfpathlineto{\pgfqpoint{5.453333in}{2.301540in}}%
\pgfpathlineto{\pgfqpoint{5.453903in}{2.304616in}}%
\pgfpathlineto{\pgfqpoint{5.453917in}{2.305118in}}%
\pgfpathlineto{\pgfqpoint{5.454677in}{2.299777in}}%
\pgfpathlineto{\pgfqpoint{5.454969in}{2.302996in}}%
\pgfpathlineto{\pgfqpoint{5.456030in}{2.300965in}}%
\pgfpathlineto{\pgfqpoint{5.455222in}{2.304223in}}%
\pgfpathlineto{\pgfqpoint{5.456088in}{2.302441in}}%
\pgfpathlineto{\pgfqpoint{5.457091in}{2.301186in}}%
\pgfpathlineto{\pgfqpoint{5.456609in}{2.305027in}}%
\pgfpathlineto{\pgfqpoint{5.457164in}{2.303158in}}%
\pgfpathlineto{\pgfqpoint{5.458138in}{2.306016in}}%
\pgfpathlineto{\pgfqpoint{5.457539in}{2.301690in}}%
\pgfpathlineto{\pgfqpoint{5.458357in}{2.305229in}}%
\pgfpathlineto{\pgfqpoint{5.459511in}{2.298654in}}%
\pgfpathlineto{\pgfqpoint{5.458445in}{2.305253in}}%
\pgfpathlineto{\pgfqpoint{5.459623in}{2.299481in}}%
\pgfpathlineto{\pgfqpoint{5.460344in}{2.297545in}}%
\pgfpathlineto{\pgfqpoint{5.459857in}{2.301284in}}%
\pgfpathlineto{\pgfqpoint{5.460757in}{2.299184in}}%
\pgfpathlineto{\pgfqpoint{5.462057in}{2.304235in}}%
\pgfpathlineto{\pgfqpoint{5.460947in}{2.297831in}}%
\pgfpathlineto{\pgfqpoint{5.462130in}{2.303991in}}%
\pgfpathlineto{\pgfqpoint{5.462529in}{2.302544in}}%
\pgfpathlineto{\pgfqpoint{5.462997in}{2.305741in}}%
\pgfpathlineto{\pgfqpoint{5.463182in}{2.305000in}}%
\pgfpathlineto{\pgfqpoint{5.463664in}{2.305515in}}%
\pgfpathlineto{\pgfqpoint{5.464039in}{2.300983in}}%
\pgfpathlineto{\pgfqpoint{5.464238in}{2.303837in}}%
\pgfpathlineto{\pgfqpoint{5.465339in}{2.302988in}}%
\pgfpathlineto{\pgfqpoint{5.465027in}{2.305424in}}%
\pgfpathlineto{\pgfqpoint{5.465353in}{2.303565in}}%
\pgfpathlineto{\pgfqpoint{5.466332in}{2.302013in}}%
\pgfpathlineto{\pgfqpoint{5.465752in}{2.304197in}}%
\pgfpathlineto{\pgfqpoint{5.466497in}{2.302784in}}%
\pgfpathlineto{\pgfqpoint{5.466765in}{2.306098in}}%
\pgfpathlineto{\pgfqpoint{5.466517in}{2.302373in}}%
\pgfpathlineto{\pgfqpoint{5.467607in}{2.303161in}}%
\pgfpathlineto{\pgfqpoint{5.468016in}{2.302479in}}%
\pgfpathlineto{\pgfqpoint{5.468605in}{2.304483in}}%
\pgfpathlineto{\pgfqpoint{5.468722in}{2.302798in}}%
\pgfpathlineto{\pgfqpoint{5.469326in}{2.304962in}}%
\pgfpathlineto{\pgfqpoint{5.468737in}{2.302771in}}%
\pgfpathlineto{\pgfqpoint{5.469852in}{2.303785in}}%
\pgfpathlineto{\pgfqpoint{5.471035in}{2.301922in}}%
\pgfpathlineto{\pgfqpoint{5.469944in}{2.304383in}}%
\pgfpathlineto{\pgfqpoint{5.471040in}{2.301980in}}%
\pgfpathlineto{\pgfqpoint{5.472028in}{2.303793in}}%
\pgfpathlineto{\pgfqpoint{5.471346in}{2.299996in}}%
\pgfpathlineto{\pgfqpoint{5.472184in}{2.303635in}}%
\pgfpathlineto{\pgfqpoint{5.472568in}{2.301988in}}%
\pgfpathlineto{\pgfqpoint{5.473002in}{2.303935in}}%
\pgfpathlineto{\pgfqpoint{5.473299in}{2.303151in}}%
\pgfpathlineto{\pgfqpoint{5.474267in}{2.303953in}}%
\pgfpathlineto{\pgfqpoint{5.473493in}{2.301228in}}%
\pgfpathlineto{\pgfqpoint{5.474418in}{2.303416in}}%
\pgfpathlineto{\pgfqpoint{5.475957in}{2.295225in}}%
\pgfpathlineto{\pgfqpoint{5.474686in}{2.304204in}}%
\pgfpathlineto{\pgfqpoint{5.476006in}{2.295745in}}%
\pgfpathlineto{\pgfqpoint{5.477154in}{2.298508in}}%
\pgfpathlineto{\pgfqpoint{5.476244in}{2.295419in}}%
\pgfpathlineto{\pgfqpoint{5.477174in}{2.298501in}}%
\pgfpathlineto{\pgfqpoint{5.478376in}{2.295967in}}%
\pgfpathlineto{\pgfqpoint{5.477335in}{2.300188in}}%
\pgfpathlineto{\pgfqpoint{5.478406in}{2.296765in}}%
\pgfpathlineto{\pgfqpoint{5.479287in}{2.301842in}}%
\pgfpathlineto{\pgfqpoint{5.479603in}{2.300964in}}%
\pgfpathlineto{\pgfqpoint{5.479681in}{2.299929in}}%
\pgfpathlineto{\pgfqpoint{5.480056in}{2.303094in}}%
\pgfpathlineto{\pgfqpoint{5.480708in}{2.301115in}}%
\pgfpathlineto{\pgfqpoint{5.481156in}{2.304729in}}%
\pgfpathlineto{\pgfqpoint{5.480796in}{2.300875in}}%
\pgfpathlineto{\pgfqpoint{5.481833in}{2.301834in}}%
\pgfpathlineto{\pgfqpoint{5.482057in}{2.300238in}}%
\pgfpathlineto{\pgfqpoint{5.482782in}{2.303179in}}%
\pgfpathlineto{\pgfqpoint{5.482938in}{2.302038in}}%
\pgfpathlineto{\pgfqpoint{5.484131in}{2.298445in}}%
\pgfpathlineto{\pgfqpoint{5.483182in}{2.302331in}}%
\pgfpathlineto{\pgfqpoint{5.484180in}{2.299345in}}%
\pgfpathlineto{\pgfqpoint{5.485173in}{2.302361in}}%
\pgfpathlineto{\pgfqpoint{5.484594in}{2.297823in}}%
\pgfpathlineto{\pgfqpoint{5.485368in}{2.302084in}}%
\pgfpathlineto{\pgfqpoint{5.485377in}{2.301836in}}%
\pgfpathlineto{\pgfqpoint{5.485932in}{2.304818in}}%
\pgfpathlineto{\pgfqpoint{5.486468in}{2.302493in}}%
\pgfpathlineto{\pgfqpoint{5.487938in}{2.296944in}}%
\pgfpathlineto{\pgfqpoint{5.486487in}{2.302885in}}%
\pgfpathlineto{\pgfqpoint{5.487948in}{2.297204in}}%
\pgfpathlineto{\pgfqpoint{5.488892in}{2.299623in}}%
\pgfpathlineto{\pgfqpoint{5.489107in}{2.297951in}}%
\pgfpathlineto{\pgfqpoint{5.489189in}{2.297184in}}%
\pgfpathlineto{\pgfqpoint{5.490071in}{2.300644in}}%
\pgfpathlineto{\pgfqpoint{5.490153in}{2.300557in}}%
\pgfpathlineto{\pgfqpoint{5.490976in}{2.302445in}}%
\pgfpathlineto{\pgfqpoint{5.490192in}{2.300420in}}%
\pgfpathlineto{\pgfqpoint{5.491307in}{2.300770in}}%
\pgfpathlineto{\pgfqpoint{5.491419in}{2.300017in}}%
\pgfpathlineto{\pgfqpoint{5.492179in}{2.304195in}}%
\pgfpathlineto{\pgfqpoint{5.492281in}{2.303399in}}%
\pgfpathlineto{\pgfqpoint{5.493026in}{2.304160in}}%
\pgfpathlineto{\pgfqpoint{5.493264in}{2.301169in}}%
\pgfpathlineto{\pgfqpoint{5.493357in}{2.302285in}}%
\pgfpathlineto{\pgfqpoint{5.494481in}{2.300829in}}%
\pgfpathlineto{\pgfqpoint{5.494141in}{2.303967in}}%
\pgfpathlineto{\pgfqpoint{5.494491in}{2.300907in}}%
\pgfpathlineto{\pgfqpoint{5.494496in}{2.301345in}}%
\pgfpathlineto{\pgfqpoint{5.495231in}{2.296948in}}%
\pgfpathlineto{\pgfqpoint{5.495582in}{2.299192in}}%
\pgfpathlineto{\pgfqpoint{5.496473in}{2.301444in}}%
\pgfpathlineto{\pgfqpoint{5.495786in}{2.297649in}}%
\pgfpathlineto{\pgfqpoint{5.496789in}{2.300252in}}%
\pgfpathlineto{\pgfqpoint{5.497388in}{2.296701in}}%
\pgfpathlineto{\pgfqpoint{5.496818in}{2.300339in}}%
\pgfpathlineto{\pgfqpoint{5.497958in}{2.299319in}}%
\pgfpathlineto{\pgfqpoint{5.497972in}{2.299539in}}%
\pgfpathlineto{\pgfqpoint{5.498259in}{2.296815in}}%
\pgfpathlineto{\pgfqpoint{5.499053in}{2.298726in}}%
\pgfpathlineto{\pgfqpoint{5.499111in}{2.298167in}}%
\pgfpathlineto{\pgfqpoint{5.499618in}{2.303828in}}%
\pgfpathlineto{\pgfqpoint{5.499681in}{2.303793in}}%
\pgfpathlineto{\pgfqpoint{5.500772in}{2.306016in}}%
\pgfpathlineto{\pgfqpoint{5.500499in}{2.302177in}}%
\pgfpathlineto{\pgfqpoint{5.500845in}{2.305630in}}%
\pgfpathlineto{\pgfqpoint{5.502914in}{2.298263in}}%
\pgfpathlineto{\pgfqpoint{5.503741in}{2.300470in}}%
\pgfpathlineto{\pgfqpoint{5.504525in}{2.302963in}}%
\pgfpathlineto{\pgfqpoint{5.504267in}{2.299409in}}%
\pgfpathlineto{\pgfqpoint{5.504851in}{2.300803in}}%
\pgfpathlineto{\pgfqpoint{5.505294in}{2.299275in}}%
\pgfpathlineto{\pgfqpoint{5.505840in}{2.302496in}}%
\pgfpathlineto{\pgfqpoint{5.506424in}{2.304570in}}%
\pgfpathlineto{\pgfqpoint{5.506901in}{2.301238in}}%
\pgfpathlineto{\pgfqpoint{5.506930in}{2.301886in}}%
\pgfpathlineto{\pgfqpoint{5.506935in}{2.301709in}}%
\pgfpathlineto{\pgfqpoint{5.507768in}{2.303839in}}%
\pgfpathlineto{\pgfqpoint{5.507992in}{2.303655in}}%
\pgfpathlineto{\pgfqpoint{5.508566in}{2.304201in}}%
\pgfpathlineto{\pgfqpoint{5.508868in}{2.300509in}}%
\pgfpathlineto{\pgfqpoint{5.509038in}{2.301460in}}%
\pgfpathlineto{\pgfqpoint{5.510153in}{2.300237in}}%
\pgfpathlineto{\pgfqpoint{5.509574in}{2.304022in}}%
\pgfpathlineto{\pgfqpoint{5.510158in}{2.300427in}}%
\pgfpathlineto{\pgfqpoint{5.510698in}{2.301573in}}%
\pgfpathlineto{\pgfqpoint{5.511283in}{2.300031in}}%
\pgfpathlineto{\pgfqpoint{5.512018in}{2.301945in}}%
\pgfpathlineto{\pgfqpoint{5.511375in}{2.299674in}}%
\pgfpathlineto{\pgfqpoint{5.512412in}{2.301068in}}%
\pgfpathlineto{\pgfqpoint{5.512865in}{2.300061in}}%
\pgfpathlineto{\pgfqpoint{5.513123in}{2.302646in}}%
\pgfpathlineto{\pgfqpoint{5.513400in}{2.302183in}}%
\pgfpathlineto{\pgfqpoint{5.513405in}{2.302322in}}%
\pgfpathlineto{\pgfqpoint{5.514262in}{2.297427in}}%
\pgfpathlineto{\pgfqpoint{5.514384in}{2.298641in}}%
\pgfpathlineto{\pgfqpoint{5.514827in}{2.296939in}}%
\pgfpathlineto{\pgfqpoint{5.514501in}{2.299514in}}%
\pgfpathlineto{\pgfqpoint{5.515397in}{2.299159in}}%
\pgfpathlineto{\pgfqpoint{5.515601in}{2.300618in}}%
\pgfpathlineto{\pgfqpoint{5.516443in}{2.296803in}}%
\pgfpathlineto{\pgfqpoint{5.516468in}{2.297050in}}%
\pgfpathlineto{\pgfqpoint{5.517617in}{2.293912in}}%
\pgfpathlineto{\pgfqpoint{5.516765in}{2.299406in}}%
\pgfpathlineto{\pgfqpoint{5.517621in}{2.294121in}}%
\pgfpathlineto{\pgfqpoint{5.518084in}{2.297477in}}%
\pgfpathlineto{\pgfqpoint{5.517636in}{2.293947in}}%
\pgfpathlineto{\pgfqpoint{5.518775in}{2.295735in}}%
\pgfpathlineto{\pgfqpoint{5.518863in}{2.295233in}}%
\pgfpathlineto{\pgfqpoint{5.518931in}{2.296845in}}%
\pgfpathlineto{\pgfqpoint{5.518936in}{2.296707in}}%
\pgfpathlineto{\pgfqpoint{5.520070in}{2.301508in}}%
\pgfpathlineto{\pgfqpoint{5.518955in}{2.296532in}}%
\pgfpathlineto{\pgfqpoint{5.520138in}{2.301441in}}%
\pgfpathlineto{\pgfqpoint{5.520825in}{2.302281in}}%
\pgfpathlineto{\pgfqpoint{5.520470in}{2.298344in}}%
\pgfpathlineto{\pgfqpoint{5.521141in}{2.299554in}}%
\pgfpathlineto{\pgfqpoint{5.522217in}{2.297564in}}%
\pgfpathlineto{\pgfqpoint{5.521828in}{2.300055in}}%
\pgfpathlineto{\pgfqpoint{5.522281in}{2.298276in}}%
\pgfpathlineto{\pgfqpoint{5.523542in}{2.306155in}}%
\pgfpathlineto{\pgfqpoint{5.523566in}{2.305850in}}%
\pgfpathlineto{\pgfqpoint{5.523571in}{2.305943in}}%
\pgfpathlineto{\pgfqpoint{5.524355in}{2.299855in}}%
\pgfpathlineto{\pgfqpoint{5.524364in}{2.299972in}}%
\pgfpathlineto{\pgfqpoint{5.525187in}{2.298838in}}%
\pgfpathlineto{\pgfqpoint{5.524856in}{2.303561in}}%
\pgfpathlineto{\pgfqpoint{5.525435in}{2.301498in}}%
\pgfpathlineto{\pgfqpoint{5.525839in}{2.302698in}}%
\pgfpathlineto{\pgfqpoint{5.526365in}{2.299208in}}%
\pgfpathlineto{\pgfqpoint{5.526458in}{2.300300in}}%
\pgfpathlineto{\pgfqpoint{5.526726in}{2.298180in}}%
\pgfpathlineto{\pgfqpoint{5.526541in}{2.301052in}}%
\pgfpathlineto{\pgfqpoint{5.527383in}{2.300950in}}%
\pgfpathlineto{\pgfqpoint{5.528021in}{2.302664in}}%
\pgfpathlineto{\pgfqpoint{5.527631in}{2.299257in}}%
\pgfpathlineto{\pgfqpoint{5.528478in}{2.300550in}}%
\pgfpathlineto{\pgfqpoint{5.528843in}{2.303064in}}%
\pgfpathlineto{\pgfqpoint{5.529432in}{2.297599in}}%
\pgfpathlineto{\pgfqpoint{5.529496in}{2.296865in}}%
\pgfpathlineto{\pgfqpoint{5.530401in}{2.301965in}}%
\pgfpathlineto{\pgfqpoint{5.530426in}{2.301883in}}%
\pgfpathlineto{\pgfqpoint{5.531628in}{2.303557in}}%
\pgfpathlineto{\pgfqpoint{5.531039in}{2.299554in}}%
\pgfpathlineto{\pgfqpoint{5.531638in}{2.303436in}}%
\pgfpathlineto{\pgfqpoint{5.532616in}{2.299357in}}%
\pgfpathlineto{\pgfqpoint{5.532777in}{2.301603in}}%
\pgfpathlineto{\pgfqpoint{5.533775in}{2.303073in}}%
\pgfpathlineto{\pgfqpoint{5.533128in}{2.301122in}}%
\pgfpathlineto{\pgfqpoint{5.533897in}{2.302337in}}%
\pgfpathlineto{\pgfqpoint{5.534096in}{2.300929in}}%
\pgfpathlineto{\pgfqpoint{5.534525in}{2.303862in}}%
\pgfpathlineto{\pgfqpoint{5.535012in}{2.301518in}}%
\pgfpathlineto{\pgfqpoint{5.535484in}{2.304128in}}%
\pgfpathlineto{\pgfqpoint{5.535864in}{2.300066in}}%
\pgfpathlineto{\pgfqpoint{5.536063in}{2.300346in}}%
\pgfpathlineto{\pgfqpoint{5.536097in}{2.299922in}}%
\pgfpathlineto{\pgfqpoint{5.536560in}{2.303047in}}%
\pgfpathlineto{\pgfqpoint{5.537154in}{2.301277in}}%
\pgfpathlineto{\pgfqpoint{5.538205in}{2.302438in}}%
\pgfpathlineto{\pgfqpoint{5.537534in}{2.299417in}}%
\pgfpathlineto{\pgfqpoint{5.538264in}{2.301293in}}%
\pgfpathlineto{\pgfqpoint{5.538478in}{2.304534in}}%
\pgfpathlineto{\pgfqpoint{5.539427in}{2.303612in}}%
\pgfpathlineto{\pgfqpoint{5.539700in}{2.302482in}}%
\pgfpathlineto{\pgfqpoint{5.539856in}{2.304325in}}%
\pgfpathlineto{\pgfqpoint{5.540528in}{2.303969in}}%
\pgfpathlineto{\pgfqpoint{5.540533in}{2.304106in}}%
\pgfpathlineto{\pgfqpoint{5.541185in}{2.299486in}}%
\pgfpathlineto{\pgfqpoint{5.541506in}{2.300473in}}%
\pgfpathlineto{\pgfqpoint{5.541511in}{2.300350in}}%
\pgfpathlineto{\pgfqpoint{5.542383in}{2.304312in}}%
\pgfpathlineto{\pgfqpoint{5.542490in}{2.303642in}}%
\pgfpathlineto{\pgfqpoint{5.543352in}{2.304354in}}%
\pgfpathlineto{\pgfqpoint{5.543123in}{2.302404in}}%
\pgfpathlineto{\pgfqpoint{5.543546in}{2.303360in}}%
\pgfpathlineto{\pgfqpoint{5.543931in}{2.299911in}}%
\pgfpathlineto{\pgfqpoint{5.544700in}{2.301056in}}%
\pgfpathlineto{\pgfqpoint{5.544739in}{2.301916in}}%
\pgfpathlineto{\pgfqpoint{5.545737in}{2.298575in}}%
\pgfpathlineto{\pgfqpoint{5.545776in}{2.298715in}}%
\pgfpathlineto{\pgfqpoint{5.545830in}{2.298970in}}%
\pgfpathlineto{\pgfqpoint{5.546200in}{2.296122in}}%
\pgfpathlineto{\pgfqpoint{5.546574in}{2.297023in}}%
\pgfpathlineto{\pgfqpoint{5.547680in}{2.294679in}}%
\pgfpathlineto{\pgfqpoint{5.546769in}{2.298594in}}%
\pgfpathlineto{\pgfqpoint{5.547733in}{2.295760in}}%
\pgfpathlineto{\pgfqpoint{5.548322in}{2.299282in}}%
\pgfpathlineto{\pgfqpoint{5.548882in}{2.298332in}}%
\pgfpathlineto{\pgfqpoint{5.548941in}{2.297466in}}%
\pgfpathlineto{\pgfqpoint{5.549846in}{2.301266in}}%
\pgfpathlineto{\pgfqpoint{5.549885in}{2.300679in}}%
\pgfpathlineto{\pgfqpoint{5.551049in}{2.299967in}}%
\pgfpathlineto{\pgfqpoint{5.550596in}{2.303630in}}%
\pgfpathlineto{\pgfqpoint{5.551053in}{2.300206in}}%
\pgfpathlineto{\pgfqpoint{5.552042in}{2.296432in}}%
\pgfpathlineto{\pgfqpoint{5.551282in}{2.301765in}}%
\pgfpathlineto{\pgfqpoint{5.552485in}{2.298620in}}%
\pgfpathlineto{\pgfqpoint{5.553751in}{2.300909in}}%
\pgfpathlineto{\pgfqpoint{5.552869in}{2.297263in}}%
\pgfpathlineto{\pgfqpoint{5.553897in}{2.299504in}}%
\pgfpathlineto{\pgfqpoint{5.553940in}{2.298783in}}%
\pgfpathlineto{\pgfqpoint{5.554627in}{2.304002in}}%
\pgfpathlineto{\pgfqpoint{5.554948in}{2.303040in}}%
\pgfpathlineto{\pgfqpoint{5.555664in}{2.305529in}}%
\pgfpathlineto{\pgfqpoint{5.555386in}{2.302582in}}%
\pgfpathlineto{\pgfqpoint{5.556199in}{2.304603in}}%
\pgfpathlineto{\pgfqpoint{5.556818in}{2.301951in}}%
\pgfpathlineto{\pgfqpoint{5.556263in}{2.305523in}}%
\pgfpathlineto{\pgfqpoint{5.557319in}{2.304254in}}%
\pgfpathlineto{\pgfqpoint{5.557490in}{2.304590in}}%
\pgfpathlineto{\pgfqpoint{5.557952in}{2.302135in}}%
\pgfpathlineto{\pgfqpoint{5.558317in}{2.302877in}}%
\pgfpathlineto{\pgfqpoint{5.559213in}{2.301439in}}%
\pgfpathlineto{\pgfqpoint{5.558668in}{2.304473in}}%
\pgfpathlineto{\pgfqpoint{5.559466in}{2.301959in}}%
\pgfpathlineto{\pgfqpoint{5.560479in}{2.303756in}}%
\pgfpathlineto{\pgfqpoint{5.559977in}{2.299748in}}%
\pgfpathlineto{\pgfqpoint{5.560630in}{2.303608in}}%
\pgfpathlineto{\pgfqpoint{5.561014in}{2.302326in}}%
\pgfpathlineto{\pgfqpoint{5.561652in}{2.304345in}}%
\pgfpathlineto{\pgfqpoint{5.561745in}{2.303116in}}%
\pgfpathlineto{\pgfqpoint{5.562290in}{2.304222in}}%
\pgfpathlineto{\pgfqpoint{5.562480in}{2.302129in}}%
\pgfpathlineto{\pgfqpoint{5.562840in}{2.302373in}}%
\pgfpathlineto{\pgfqpoint{5.563561in}{2.301157in}}%
\pgfpathlineto{\pgfqpoint{5.562976in}{2.304417in}}%
\pgfpathlineto{\pgfqpoint{5.563931in}{2.302678in}}%
\pgfpathlineto{\pgfqpoint{5.564082in}{2.302117in}}%
\pgfpathlineto{\pgfqpoint{5.565119in}{2.305383in}}%
\pgfpathlineto{\pgfqpoint{5.565493in}{2.302031in}}%
\pgfpathlineto{\pgfqpoint{5.565128in}{2.305480in}}%
\pgfpathlineto{\pgfqpoint{5.566336in}{2.303967in}}%
\pgfpathlineto{\pgfqpoint{5.569593in}{2.292671in}}%
\pgfpathlineto{\pgfqpoint{5.569671in}{2.293447in}}%
\pgfpathlineto{\pgfqpoint{5.571019in}{2.299537in}}%
\pgfpathlineto{\pgfqpoint{5.569714in}{2.293129in}}%
\pgfpathlineto{\pgfqpoint{5.571029in}{2.299067in}}%
\pgfpathlineto{\pgfqpoint{5.571039in}{2.298903in}}%
\pgfpathlineto{\pgfqpoint{5.571818in}{2.304412in}}%
\pgfpathlineto{\pgfqpoint{5.571823in}{2.304310in}}%
\pgfpathlineto{\pgfqpoint{5.572592in}{2.305212in}}%
\pgfpathlineto{\pgfqpoint{5.572051in}{2.302956in}}%
\pgfpathlineto{\pgfqpoint{5.572937in}{2.304514in}}%
\pgfpathlineto{\pgfqpoint{5.574106in}{2.301664in}}%
\pgfpathlineto{\pgfqpoint{5.574111in}{2.301794in}}%
\pgfpathlineto{\pgfqpoint{5.574890in}{2.305082in}}%
\pgfpathlineto{\pgfqpoint{5.574145in}{2.301778in}}%
\pgfpathlineto{\pgfqpoint{5.575240in}{2.303336in}}%
\pgfpathlineto{\pgfqpoint{5.575683in}{2.304709in}}%
\pgfpathlineto{\pgfqpoint{5.576457in}{2.298555in}}%
\pgfpathlineto{\pgfqpoint{5.576930in}{2.302749in}}%
\pgfpathlineto{\pgfqpoint{5.577772in}{2.299912in}}%
\pgfpathlineto{\pgfqpoint{5.577786in}{2.299147in}}%
\pgfpathlineto{\pgfqpoint{5.578424in}{2.302899in}}%
\pgfpathlineto{\pgfqpoint{5.578867in}{2.300883in}}%
\pgfpathlineto{\pgfqpoint{5.579787in}{2.305018in}}%
\pgfpathlineto{\pgfqpoint{5.580006in}{2.302259in}}%
\pgfpathlineto{\pgfqpoint{5.580172in}{2.299962in}}%
\pgfpathlineto{\pgfqpoint{5.581092in}{2.303408in}}%
\pgfpathlineto{\pgfqpoint{5.581097in}{2.303142in}}%
\pgfpathlineto{\pgfqpoint{5.582183in}{2.303848in}}%
\pgfpathlineto{\pgfqpoint{5.581749in}{2.301868in}}%
\pgfpathlineto{\pgfqpoint{5.582217in}{2.303438in}}%
\pgfpathlineto{\pgfqpoint{5.582835in}{2.300555in}}%
\pgfpathlineto{\pgfqpoint{5.582280in}{2.303672in}}%
\pgfpathlineto{\pgfqpoint{5.583337in}{2.303117in}}%
\pgfpathlineto{\pgfqpoint{5.583658in}{2.304577in}}%
\pgfpathlineto{\pgfqpoint{5.584159in}{2.300765in}}%
\pgfpathlineto{\pgfqpoint{5.584417in}{2.302596in}}%
\pgfpathlineto{\pgfqpoint{5.584529in}{2.299917in}}%
\pgfpathlineto{\pgfqpoint{5.585508in}{2.302796in}}%
\pgfpathlineto{\pgfqpoint{5.585527in}{2.302467in}}%
\pgfpathlineto{\pgfqpoint{5.587314in}{2.296438in}}%
\pgfpathlineto{\pgfqpoint{5.586097in}{2.303720in}}%
\pgfpathlineto{\pgfqpoint{5.587528in}{2.296773in}}%
\pgfpathlineto{\pgfqpoint{5.587592in}{2.296326in}}%
\pgfpathlineto{\pgfqpoint{5.587772in}{2.298165in}}%
\pgfpathlineto{\pgfqpoint{5.587830in}{2.298037in}}%
\pgfpathlineto{\pgfqpoint{5.588526in}{2.301679in}}%
\pgfpathlineto{\pgfqpoint{5.587884in}{2.296636in}}%
\pgfpathlineto{\pgfqpoint{5.588989in}{2.299880in}}%
\pgfpathlineto{\pgfqpoint{5.589291in}{2.297497in}}%
\pgfpathlineto{\pgfqpoint{5.589836in}{2.300706in}}%
\pgfpathlineto{\pgfqpoint{5.590113in}{2.299234in}}%
\pgfpathlineto{\pgfqpoint{5.590790in}{2.304469in}}%
\pgfpathlineto{\pgfqpoint{5.590191in}{2.298703in}}%
\pgfpathlineto{\pgfqpoint{5.591306in}{2.303998in}}%
\pgfpathlineto{\pgfqpoint{5.591701in}{2.300555in}}%
\pgfpathlineto{\pgfqpoint{5.592299in}{2.305575in}}%
\pgfpathlineto{\pgfqpoint{5.592368in}{2.304807in}}%
\pgfpathlineto{\pgfqpoint{5.592372in}{2.304884in}}%
\pgfpathlineto{\pgfqpoint{5.592567in}{2.301138in}}%
\pgfpathlineto{\pgfqpoint{5.593346in}{2.302907in}}%
\pgfpathlineto{\pgfqpoint{5.593444in}{2.301198in}}%
\pgfpathlineto{\pgfqpoint{5.594101in}{2.306284in}}%
\pgfpathlineto{\pgfqpoint{5.594412in}{2.304917in}}%
\pgfpathlineto{\pgfqpoint{5.594422in}{2.305144in}}%
\pgfpathlineto{\pgfqpoint{5.595186in}{2.301182in}}%
\pgfpathlineto{\pgfqpoint{5.595347in}{2.301564in}}%
\pgfpathlineto{\pgfqpoint{5.595440in}{2.300887in}}%
\pgfpathlineto{\pgfqpoint{5.596048in}{2.302608in}}%
\pgfpathlineto{\pgfqpoint{5.596457in}{2.301572in}}%
\pgfpathlineto{\pgfqpoint{5.597295in}{2.305276in}}%
\pgfpathlineto{\pgfqpoint{5.596988in}{2.301242in}}%
\pgfpathlineto{\pgfqpoint{5.597708in}{2.303354in}}%
\pgfpathlineto{\pgfqpoint{5.598624in}{2.296073in}}%
\pgfpathlineto{\pgfqpoint{5.598974in}{2.297003in}}%
\pgfpathlineto{\pgfqpoint{5.599700in}{2.299251in}}%
\pgfpathlineto{\pgfqpoint{5.599042in}{2.296317in}}%
\pgfpathlineto{\pgfqpoint{5.600147in}{2.297904in}}%
\pgfpathlineto{\pgfqpoint{5.600980in}{2.299517in}}%
\pgfpathlineto{\pgfqpoint{5.601267in}{2.296697in}}%
\pgfpathlineto{\pgfqpoint{5.601915in}{2.298116in}}%
\pgfpathlineto{\pgfqpoint{5.602217in}{2.294745in}}%
\pgfpathlineto{\pgfqpoint{5.602363in}{2.296138in}}%
\pgfpathlineto{\pgfqpoint{5.603244in}{2.294103in}}%
\pgfpathlineto{\pgfqpoint{5.602786in}{2.296906in}}%
\pgfpathlineto{\pgfqpoint{5.603482in}{2.295724in}}%
\pgfpathlineto{\pgfqpoint{5.603867in}{2.295187in}}%
\pgfpathlineto{\pgfqpoint{5.604130in}{2.298821in}}%
\pgfpathlineto{\pgfqpoint{5.604495in}{2.296926in}}%
\pgfpathlineto{\pgfqpoint{5.605308in}{2.298625in}}%
\pgfpathlineto{\pgfqpoint{5.604709in}{2.295184in}}%
\pgfpathlineto{\pgfqpoint{5.605629in}{2.298027in}}%
\pgfpathlineto{\pgfqpoint{5.605712in}{2.296961in}}%
\pgfpathlineto{\pgfqpoint{5.606384in}{2.300521in}}%
\pgfpathlineto{\pgfqpoint{5.606705in}{2.299060in}}%
\pgfpathlineto{\pgfqpoint{5.607708in}{2.301741in}}%
\pgfpathlineto{\pgfqpoint{5.606754in}{2.299036in}}%
\pgfpathlineto{\pgfqpoint{5.607854in}{2.300862in}}%
\pgfpathlineto{\pgfqpoint{5.608521in}{2.297182in}}%
\pgfpathlineto{\pgfqpoint{5.609003in}{2.297836in}}%
\pgfpathlineto{\pgfqpoint{5.609115in}{2.297019in}}%
\pgfpathlineto{\pgfqpoint{5.609558in}{2.300305in}}%
\pgfpathlineto{\pgfqpoint{5.609685in}{2.301203in}}%
\pgfpathlineto{\pgfqpoint{5.610123in}{2.298896in}}%
\pgfpathlineto{\pgfqpoint{5.610649in}{2.299980in}}%
\pgfpathlineto{\pgfqpoint{5.610790in}{2.298132in}}%
\pgfpathlineto{\pgfqpoint{5.611603in}{2.302363in}}%
\pgfpathlineto{\pgfqpoint{5.611730in}{2.300959in}}%
\pgfpathlineto{\pgfqpoint{5.611788in}{2.301546in}}%
\pgfpathlineto{\pgfqpoint{5.612392in}{2.298484in}}%
\pgfpathlineto{\pgfqpoint{5.612810in}{2.299502in}}%
\pgfpathlineto{\pgfqpoint{5.613346in}{2.301883in}}%
\pgfpathlineto{\pgfqpoint{5.613030in}{2.298467in}}%
\pgfpathlineto{\pgfqpoint{5.613969in}{2.300683in}}%
\pgfpathlineto{\pgfqpoint{5.614733in}{2.298086in}}%
\pgfpathlineto{\pgfqpoint{5.614076in}{2.301350in}}%
\pgfpathlineto{\pgfqpoint{5.615103in}{2.299041in}}%
\pgfpathlineto{\pgfqpoint{5.615113in}{2.299579in}}%
\pgfpathlineto{\pgfqpoint{5.616043in}{2.294306in}}%
\pgfpathlineto{\pgfqpoint{5.616175in}{2.295548in}}%
\pgfpathlineto{\pgfqpoint{5.616306in}{2.297039in}}%
\pgfpathlineto{\pgfqpoint{5.616968in}{2.294441in}}%
\pgfpathlineto{\pgfqpoint{5.617260in}{2.294787in}}%
\pgfpathlineto{\pgfqpoint{5.617664in}{2.293176in}}%
\pgfpathlineto{\pgfqpoint{5.617299in}{2.295166in}}%
\pgfpathlineto{\pgfqpoint{5.618385in}{2.294422in}}%
\pgfpathlineto{\pgfqpoint{5.618818in}{2.295483in}}%
\pgfpathlineto{\pgfqpoint{5.618604in}{2.293107in}}%
\pgfpathlineto{\pgfqpoint{5.619490in}{2.294192in}}%
\pgfpathlineto{\pgfqpoint{5.619514in}{2.294052in}}%
\pgfpathlineto{\pgfqpoint{5.619743in}{2.296005in}}%
\pgfpathlineto{\pgfqpoint{5.620814in}{2.300617in}}%
\pgfpathlineto{\pgfqpoint{5.620907in}{2.299016in}}%
\pgfpathlineto{\pgfqpoint{5.620912in}{2.298768in}}%
\pgfpathlineto{\pgfqpoint{5.621773in}{2.302778in}}%
\pgfpathlineto{\pgfqpoint{5.621992in}{2.300189in}}%
\pgfpathlineto{\pgfqpoint{5.623093in}{2.306133in}}%
\pgfpathlineto{\pgfqpoint{5.622022in}{2.299828in}}%
\pgfpathlineto{\pgfqpoint{5.623409in}{2.303882in}}%
\pgfpathlineto{\pgfqpoint{5.624115in}{2.302421in}}%
\pgfpathlineto{\pgfqpoint{5.623794in}{2.306109in}}%
\pgfpathlineto{\pgfqpoint{5.624514in}{2.304266in}}%
\pgfpathlineto{\pgfqpoint{5.624553in}{2.305545in}}%
\pgfpathlineto{\pgfqpoint{5.625381in}{2.302032in}}%
\pgfpathlineto{\pgfqpoint{5.625527in}{2.302593in}}%
\pgfpathlineto{\pgfqpoint{5.626715in}{2.300266in}}%
\pgfpathlineto{\pgfqpoint{5.625751in}{2.303881in}}%
\pgfpathlineto{\pgfqpoint{5.626734in}{2.300595in}}%
\pgfpathlineto{\pgfqpoint{5.626827in}{2.301463in}}%
\pgfpathlineto{\pgfqpoint{5.627411in}{2.298404in}}%
\pgfpathlineto{\pgfqpoint{5.627820in}{2.300172in}}%
\pgfpathlineto{\pgfqpoint{5.628185in}{2.298429in}}%
\pgfpathlineto{\pgfqpoint{5.628891in}{2.301471in}}%
\pgfpathlineto{\pgfqpoint{5.628901in}{2.301442in}}%
\pgfpathlineto{\pgfqpoint{5.629967in}{2.304349in}}%
\pgfpathlineto{\pgfqpoint{5.630035in}{2.303593in}}%
\pgfpathlineto{\pgfqpoint{5.630254in}{2.302522in}}%
\pgfpathlineto{\pgfqpoint{5.630663in}{2.304421in}}%
\pgfpathlineto{\pgfqpoint{5.631155in}{2.303188in}}%
\pgfpathlineto{\pgfqpoint{5.631423in}{2.305749in}}%
\pgfpathlineto{\pgfqpoint{5.632026in}{2.302057in}}%
\pgfpathlineto{\pgfqpoint{5.632221in}{2.302432in}}%
\pgfpathlineto{\pgfqpoint{5.632289in}{2.301086in}}%
\pgfpathlineto{\pgfqpoint{5.633263in}{2.305453in}}%
\pgfpathlineto{\pgfqpoint{5.633292in}{2.305404in}}%
\pgfpathlineto{\pgfqpoint{5.633911in}{2.302698in}}%
\pgfpathlineto{\pgfqpoint{5.633380in}{2.306160in}}%
\pgfpathlineto{\pgfqpoint{5.634597in}{2.303288in}}%
\pgfpathlineto{\pgfqpoint{5.635892in}{2.305422in}}%
\pgfpathlineto{\pgfqpoint{5.635483in}{2.302860in}}%
\pgfpathlineto{\pgfqpoint{5.635931in}{2.304335in}}%
\pgfpathlineto{\pgfqpoint{5.636252in}{2.301586in}}%
\pgfpathlineto{\pgfqpoint{5.636613in}{2.304641in}}%
\pgfpathlineto{\pgfqpoint{5.637065in}{2.302462in}}%
\pgfpathlineto{\pgfqpoint{5.637966in}{2.305661in}}%
\pgfpathlineto{\pgfqpoint{5.638258in}{2.304063in}}%
\pgfpathlineto{\pgfqpoint{5.639329in}{2.302968in}}%
\pgfpathlineto{\pgfqpoint{5.638789in}{2.305484in}}%
\pgfpathlineto{\pgfqpoint{5.639388in}{2.303455in}}%
\pgfpathlineto{\pgfqpoint{5.639451in}{2.304124in}}%
\pgfpathlineto{\pgfqpoint{5.640386in}{2.298935in}}%
\pgfpathlineto{\pgfqpoint{5.640395in}{2.298994in}}%
\pgfpathlineto{\pgfqpoint{5.640425in}{2.298154in}}%
\pgfpathlineto{\pgfqpoint{5.641413in}{2.300823in}}%
\pgfpathlineto{\pgfqpoint{5.641491in}{2.299931in}}%
\pgfpathlineto{\pgfqpoint{5.641924in}{2.303409in}}%
\pgfpathlineto{\pgfqpoint{5.642552in}{2.299705in}}%
\pgfpathlineto{\pgfqpoint{5.642615in}{2.298593in}}%
\pgfpathlineto{\pgfqpoint{5.643497in}{2.303747in}}%
\pgfpathlineto{\pgfqpoint{5.643623in}{2.302257in}}%
\pgfpathlineto{\pgfqpoint{5.644451in}{2.304162in}}%
\pgfpathlineto{\pgfqpoint{5.643750in}{2.301224in}}%
\pgfpathlineto{\pgfqpoint{5.644753in}{2.303161in}}%
\pgfpathlineto{\pgfqpoint{5.645089in}{2.302060in}}%
\pgfpathlineto{\pgfqpoint{5.645809in}{2.304290in}}%
\pgfpathlineto{\pgfqpoint{5.645853in}{2.303842in}}%
\pgfpathlineto{\pgfqpoint{5.646569in}{2.307334in}}%
\pgfpathlineto{\pgfqpoint{5.645911in}{2.303049in}}%
\pgfpathlineto{\pgfqpoint{5.646963in}{2.303879in}}%
\pgfpathlineto{\pgfqpoint{5.647391in}{2.302214in}}%
\pgfpathlineto{\pgfqpoint{5.647669in}{2.305439in}}%
\pgfpathlineto{\pgfqpoint{5.648097in}{2.302610in}}%
\pgfpathlineto{\pgfqpoint{5.648638in}{2.304996in}}%
\pgfpathlineto{\pgfqpoint{5.648175in}{2.301546in}}%
\pgfpathlineto{\pgfqpoint{5.649451in}{2.303238in}}%
\pgfpathlineto{\pgfqpoint{5.650381in}{2.302169in}}%
\pgfpathlineto{\pgfqpoint{5.649709in}{2.304246in}}%
\pgfpathlineto{\pgfqpoint{5.650566in}{2.303076in}}%
\pgfpathlineto{\pgfqpoint{5.650580in}{2.303556in}}%
\pgfpathlineto{\pgfqpoint{5.651354in}{2.299717in}}%
\pgfpathlineto{\pgfqpoint{5.651642in}{2.301496in}}%
\pgfpathlineto{\pgfqpoint{5.652460in}{2.304872in}}%
\pgfpathlineto{\pgfqpoint{5.651763in}{2.300968in}}%
\pgfpathlineto{\pgfqpoint{5.652932in}{2.303785in}}%
\pgfpathlineto{\pgfqpoint{5.653886in}{2.301438in}}%
\pgfpathlineto{\pgfqpoint{5.653404in}{2.304357in}}%
\pgfpathlineto{\pgfqpoint{5.654081in}{2.302203in}}%
\pgfpathlineto{\pgfqpoint{5.654835in}{2.304911in}}%
\pgfpathlineto{\pgfqpoint{5.654115in}{2.301497in}}%
\pgfpathlineto{\pgfqpoint{5.655235in}{2.303142in}}%
\pgfpathlineto{\pgfqpoint{5.655916in}{2.301571in}}%
\pgfpathlineto{\pgfqpoint{5.656242in}{2.304145in}}%
\pgfpathlineto{\pgfqpoint{5.656310in}{2.303348in}}%
\pgfpathlineto{\pgfqpoint{5.657255in}{2.305260in}}%
\pgfpathlineto{\pgfqpoint{5.656598in}{2.301866in}}%
\pgfpathlineto{\pgfqpoint{5.657421in}{2.303505in}}%
\pgfpathlineto{\pgfqpoint{5.658209in}{2.302610in}}%
\pgfpathlineto{\pgfqpoint{5.657606in}{2.304796in}}%
\pgfpathlineto{\pgfqpoint{5.658511in}{2.303809in}}%
\pgfpathlineto{\pgfqpoint{5.659392in}{2.305462in}}%
\pgfpathlineto{\pgfqpoint{5.659124in}{2.302712in}}%
\pgfpathlineto{\pgfqpoint{5.659621in}{2.304250in}}%
\pgfpathlineto{\pgfqpoint{5.660035in}{2.300865in}}%
\pgfpathlineto{\pgfqpoint{5.659704in}{2.304809in}}%
\pgfpathlineto{\pgfqpoint{5.660726in}{2.304503in}}%
\pgfpathlineto{\pgfqpoint{5.660755in}{2.305137in}}%
\pgfpathlineto{\pgfqpoint{5.661646in}{2.301810in}}%
\pgfpathlineto{\pgfqpoint{5.661656in}{2.301975in}}%
\pgfpathlineto{\pgfqpoint{5.661846in}{2.299823in}}%
\pgfpathlineto{\pgfqpoint{5.662717in}{2.302812in}}%
\pgfpathlineto{\pgfqpoint{5.662747in}{2.302344in}}%
\pgfpathlineto{\pgfqpoint{5.663438in}{2.304914in}}%
\pgfpathlineto{\pgfqpoint{5.662917in}{2.301429in}}%
\pgfpathlineto{\pgfqpoint{5.663852in}{2.302111in}}%
\pgfpathlineto{\pgfqpoint{5.664129in}{2.301657in}}%
\pgfpathlineto{\pgfqpoint{5.664626in}{2.306038in}}%
\pgfpathlineto{\pgfqpoint{5.664918in}{2.304510in}}%
\pgfpathlineto{\pgfqpoint{5.664991in}{2.305475in}}%
\pgfpathlineto{\pgfqpoint{5.665351in}{2.302051in}}%
\pgfpathlineto{\pgfqpoint{5.666520in}{2.299288in}}%
\pgfpathlineto{\pgfqpoint{5.665945in}{2.304582in}}%
\pgfpathlineto{\pgfqpoint{5.666525in}{2.299585in}}%
\pgfpathlineto{\pgfqpoint{5.667459in}{2.302506in}}%
\pgfpathlineto{\pgfqpoint{5.666836in}{2.298426in}}%
\pgfpathlineto{\pgfqpoint{5.667790in}{2.300906in}}%
\pgfpathlineto{\pgfqpoint{5.668102in}{2.298712in}}%
\pgfpathlineto{\pgfqpoint{5.668540in}{2.303574in}}%
\pgfpathlineto{\pgfqpoint{5.668866in}{2.302144in}}%
\pgfpathlineto{\pgfqpoint{5.669095in}{2.302184in}}%
\pgfpathlineto{\pgfqpoint{5.669149in}{2.303576in}}%
\pgfpathlineto{\pgfqpoint{5.669543in}{2.305963in}}%
\pgfpathlineto{\pgfqpoint{5.670225in}{2.301796in}}%
\pgfpathlineto{\pgfqpoint{5.670230in}{2.301878in}}%
\pgfpathlineto{\pgfqpoint{5.670244in}{2.301754in}}%
\pgfpathlineto{\pgfqpoint{5.670619in}{2.304443in}}%
\pgfpathlineto{\pgfqpoint{5.670624in}{2.304392in}}%
\pgfpathlineto{\pgfqpoint{5.670731in}{2.305128in}}%
\pgfpathlineto{\pgfqpoint{5.671612in}{2.301567in}}%
\pgfpathlineto{\pgfqpoint{5.672815in}{2.296413in}}%
\pgfpathlineto{\pgfqpoint{5.671763in}{2.302095in}}%
\pgfpathlineto{\pgfqpoint{5.672829in}{2.296836in}}%
\pgfpathlineto{\pgfqpoint{5.672951in}{2.295618in}}%
\pgfpathlineto{\pgfqpoint{5.673798in}{2.298125in}}%
\pgfpathlineto{\pgfqpoint{5.673920in}{2.297288in}}%
\pgfpathlineto{\pgfqpoint{5.674543in}{2.300304in}}%
\pgfpathlineto{\pgfqpoint{5.674261in}{2.296387in}}%
\pgfpathlineto{\pgfqpoint{5.675064in}{2.297754in}}%
\pgfpathlineto{\pgfqpoint{5.675069in}{2.297639in}}%
\pgfpathlineto{\pgfqpoint{5.675799in}{2.300571in}}%
\pgfpathlineto{\pgfqpoint{5.676096in}{2.299502in}}%
\pgfpathlineto{\pgfqpoint{5.676919in}{2.300401in}}%
\pgfpathlineto{\pgfqpoint{5.677065in}{2.297841in}}%
\pgfpathlineto{\pgfqpoint{5.677172in}{2.298252in}}%
\pgfpathlineto{\pgfqpoint{5.677177in}{2.298034in}}%
\pgfpathlineto{\pgfqpoint{5.677737in}{2.302060in}}%
\pgfpathlineto{\pgfqpoint{5.678209in}{2.301461in}}%
\pgfpathlineto{\pgfqpoint{5.679538in}{2.298493in}}%
\pgfpathlineto{\pgfqpoint{5.678525in}{2.302791in}}%
\pgfpathlineto{\pgfqpoint{5.679597in}{2.299678in}}%
\pgfpathlineto{\pgfqpoint{5.679811in}{2.301526in}}%
\pgfpathlineto{\pgfqpoint{5.680507in}{2.298059in}}%
\pgfpathlineto{\pgfqpoint{5.680702in}{2.299417in}}%
\pgfpathlineto{\pgfqpoint{5.681281in}{2.297424in}}%
\pgfpathlineto{\pgfqpoint{5.681768in}{2.300253in}}%
\pgfpathlineto{\pgfqpoint{5.681792in}{2.299685in}}%
\pgfpathlineto{\pgfqpoint{5.682922in}{2.302744in}}%
\pgfpathlineto{\pgfqpoint{5.681963in}{2.297552in}}%
\pgfpathlineto{\pgfqpoint{5.682936in}{2.302339in}}%
\pgfpathlineto{\pgfqpoint{5.683779in}{2.297899in}}%
\pgfpathlineto{\pgfqpoint{5.684080in}{2.299527in}}%
\pgfpathlineto{\pgfqpoint{5.684139in}{2.299424in}}%
\pgfpathlineto{\pgfqpoint{5.685244in}{2.302858in}}%
\pgfpathlineto{\pgfqpoint{5.686354in}{2.304443in}}%
\pgfpathlineto{\pgfqpoint{5.685843in}{2.301241in}}%
\pgfpathlineto{\pgfqpoint{5.686427in}{2.303130in}}%
\pgfpathlineto{\pgfqpoint{5.686787in}{2.302249in}}%
\pgfpathlineto{\pgfqpoint{5.687381in}{2.305094in}}%
\pgfpathlineto{\pgfqpoint{5.687532in}{2.303082in}}%
\pgfpathlineto{\pgfqpoint{5.688613in}{2.303915in}}%
\pgfpathlineto{\pgfqpoint{5.687965in}{2.300870in}}%
\pgfpathlineto{\pgfqpoint{5.688652in}{2.303372in}}%
\pgfpathlineto{\pgfqpoint{5.689538in}{2.298577in}}%
\pgfpathlineto{\pgfqpoint{5.689864in}{2.301994in}}%
\pgfpathlineto{\pgfqpoint{5.691067in}{2.304201in}}%
\pgfpathlineto{\pgfqpoint{5.689991in}{2.301344in}}%
\pgfpathlineto{\pgfqpoint{5.691081in}{2.303474in}}%
\pgfpathlineto{\pgfqpoint{5.692099in}{2.300760in}}%
\pgfpathlineto{\pgfqpoint{5.691252in}{2.304250in}}%
\pgfpathlineto{\pgfqpoint{5.692235in}{2.301504in}}%
\pgfpathlineto{\pgfqpoint{5.692289in}{2.302345in}}%
\pgfpathlineto{\pgfqpoint{5.693136in}{2.299023in}}%
\pgfpathlineto{\pgfqpoint{5.693326in}{2.300202in}}%
\pgfpathlineto{\pgfqpoint{5.694012in}{2.302323in}}%
\pgfpathlineto{\pgfqpoint{5.693350in}{2.299525in}}%
\pgfpathlineto{\pgfqpoint{5.694484in}{2.300667in}}%
\pgfpathlineto{\pgfqpoint{5.695570in}{2.297732in}}%
\pgfpathlineto{\pgfqpoint{5.695127in}{2.301357in}}%
\pgfpathlineto{\pgfqpoint{5.695619in}{2.298564in}}%
\pgfpathlineto{\pgfqpoint{5.695672in}{2.297557in}}%
\pgfpathlineto{\pgfqpoint{5.696120in}{2.301001in}}%
\pgfpathlineto{\pgfqpoint{5.696704in}{2.299715in}}%
\pgfpathlineto{\pgfqpoint{5.697488in}{2.303841in}}%
\pgfpathlineto{\pgfqpoint{5.696880in}{2.299554in}}%
\pgfpathlineto{\pgfqpoint{5.697897in}{2.303250in}}%
\pgfpathlineto{\pgfqpoint{5.699402in}{2.296267in}}%
\pgfpathlineto{\pgfqpoint{5.698316in}{2.303857in}}%
\pgfpathlineto{\pgfqpoint{5.699713in}{2.299357in}}%
\pgfpathlineto{\pgfqpoint{5.700181in}{2.304429in}}%
\pgfpathlineto{\pgfqpoint{5.700882in}{2.301589in}}%
\pgfpathlineto{\pgfqpoint{5.701295in}{2.298503in}}%
\pgfpathlineto{\pgfqpoint{5.700921in}{2.302292in}}%
\pgfpathlineto{\pgfqpoint{5.702026in}{2.299719in}}%
\pgfpathlineto{\pgfqpoint{5.703155in}{2.306790in}}%
\pgfpathlineto{\pgfqpoint{5.702094in}{2.299621in}}%
\pgfpathlineto{\pgfqpoint{5.703257in}{2.305855in}}%
\pgfpathlineto{\pgfqpoint{5.704358in}{2.303274in}}%
\pgfpathlineto{\pgfqpoint{5.703817in}{2.307816in}}%
\pgfpathlineto{\pgfqpoint{5.704411in}{2.303929in}}%
\pgfpathlineto{\pgfqpoint{5.705356in}{2.305254in}}%
\pgfpathlineto{\pgfqpoint{5.704703in}{2.301617in}}%
\pgfpathlineto{\pgfqpoint{5.705531in}{2.304407in}}%
\pgfpathlineto{\pgfqpoint{5.705589in}{2.305170in}}%
\pgfpathlineto{\pgfqpoint{5.706349in}{2.300715in}}%
\pgfpathlineto{\pgfqpoint{5.706383in}{2.300815in}}%
\pgfpathlineto{\pgfqpoint{5.706432in}{2.299946in}}%
\pgfpathlineto{\pgfqpoint{5.707410in}{2.302254in}}%
\pgfpathlineto{\pgfqpoint{5.707425in}{2.303188in}}%
\pgfpathlineto{\pgfqpoint{5.707999in}{2.300479in}}%
\pgfpathlineto{\pgfqpoint{5.708520in}{2.302471in}}%
\pgfpathlineto{\pgfqpoint{5.708681in}{2.301518in}}%
\pgfpathlineto{\pgfqpoint{5.709557in}{2.304553in}}%
\pgfpathlineto{\pgfqpoint{5.709567in}{2.304378in}}%
\pgfpathlineto{\pgfqpoint{5.710682in}{2.306291in}}%
\pgfpathlineto{\pgfqpoint{5.710015in}{2.303327in}}%
\pgfpathlineto{\pgfqpoint{5.710701in}{2.306077in}}%
\pgfpathlineto{\pgfqpoint{5.711831in}{2.304184in}}%
\pgfpathlineto{\pgfqpoint{5.711495in}{2.307222in}}%
\pgfpathlineto{\pgfqpoint{5.711855in}{2.304509in}}%
\pgfpathlineto{\pgfqpoint{5.711948in}{2.305575in}}%
\pgfpathlineto{\pgfqpoint{5.712269in}{2.301378in}}%
\pgfpathlineto{\pgfqpoint{5.712926in}{2.303348in}}%
\pgfpathlineto{\pgfqpoint{5.713014in}{2.301798in}}%
\pgfpathlineto{\pgfqpoint{5.713584in}{2.304551in}}%
\pgfpathlineto{\pgfqpoint{5.714022in}{2.304130in}}%
\pgfpathlineto{\pgfqpoint{5.714830in}{2.306787in}}%
\pgfpathlineto{\pgfqpoint{5.715093in}{2.302151in}}%
\pgfpathlineto{\pgfqpoint{5.716276in}{2.304085in}}%
\pgfpathlineto{\pgfqpoint{5.715351in}{2.300355in}}%
\pgfpathlineto{\pgfqpoint{5.716295in}{2.303840in}}%
\pgfpathlineto{\pgfqpoint{5.716875in}{2.300770in}}%
\pgfpathlineto{\pgfqpoint{5.717425in}{2.302785in}}%
\pgfpathlineto{\pgfqpoint{5.718067in}{2.307256in}}%
\pgfpathlineto{\pgfqpoint{5.717454in}{2.302505in}}%
\pgfpathlineto{\pgfqpoint{5.718622in}{2.303154in}}%
\pgfpathlineto{\pgfqpoint{5.719776in}{2.301059in}}%
\pgfpathlineto{\pgfqpoint{5.718949in}{2.305781in}}%
\pgfpathlineto{\pgfqpoint{5.719805in}{2.302097in}}%
\pgfpathlineto{\pgfqpoint{5.720253in}{2.304234in}}%
\pgfpathlineto{\pgfqpoint{5.720546in}{2.301167in}}%
\pgfpathlineto{\pgfqpoint{5.720935in}{2.302588in}}%
\pgfpathlineto{\pgfqpoint{5.722045in}{2.300395in}}%
\pgfpathlineto{\pgfqpoint{5.721139in}{2.304104in}}%
\pgfpathlineto{\pgfqpoint{5.722230in}{2.301384in}}%
\pgfpathlineto{\pgfqpoint{5.722766in}{2.302452in}}%
\pgfpathlineto{\pgfqpoint{5.723077in}{2.300148in}}%
\pgfpathlineto{\pgfqpoint{5.723345in}{2.301785in}}%
\pgfpathlineto{\pgfqpoint{5.723608in}{2.301244in}}%
\pgfpathlineto{\pgfqpoint{5.724065in}{2.304556in}}%
\pgfpathlineto{\pgfqpoint{5.724416in}{2.303565in}}%
\pgfpathlineto{\pgfqpoint{5.725594in}{2.300715in}}%
\pgfpathlineto{\pgfqpoint{5.724572in}{2.304634in}}%
\pgfpathlineto{\pgfqpoint{5.725721in}{2.301595in}}%
\pgfpathlineto{\pgfqpoint{5.726626in}{2.303896in}}%
\pgfpathlineto{\pgfqpoint{5.726105in}{2.299979in}}%
\pgfpathlineto{\pgfqpoint{5.726860in}{2.302488in}}%
\pgfpathlineto{\pgfqpoint{5.726991in}{2.303923in}}%
\pgfpathlineto{\pgfqpoint{5.727410in}{2.300723in}}%
\pgfpathlineto{\pgfqpoint{5.727907in}{2.301635in}}%
\pgfpathlineto{\pgfqpoint{5.728754in}{2.299181in}}%
\pgfpathlineto{\pgfqpoint{5.728209in}{2.302520in}}%
\pgfpathlineto{\pgfqpoint{5.729017in}{2.301623in}}%
\pgfpathlineto{\pgfqpoint{5.729382in}{2.299050in}}%
\pgfpathlineto{\pgfqpoint{5.729674in}{2.303129in}}%
\pgfpathlineto{\pgfqpoint{5.730000in}{2.302756in}}%
\pgfpathlineto{\pgfqpoint{5.731120in}{2.305367in}}%
\pgfpathlineto{\pgfqpoint{5.730468in}{2.301007in}}%
\pgfpathlineto{\pgfqpoint{5.731139in}{2.305237in}}%
\pgfpathlineto{\pgfqpoint{5.731597in}{2.301573in}}%
\pgfpathlineto{\pgfqpoint{5.731281in}{2.305757in}}%
\pgfpathlineto{\pgfqpoint{5.732327in}{2.302061in}}%
\pgfpathlineto{\pgfqpoint{5.733150in}{2.303874in}}%
\pgfpathlineto{\pgfqpoint{5.733350in}{2.301479in}}%
\pgfpathlineto{\pgfqpoint{5.734085in}{2.299575in}}%
\pgfpathlineto{\pgfqpoint{5.734440in}{2.302785in}}%
\pgfpathlineto{\pgfqpoint{5.734542in}{2.303640in}}%
\pgfpathlineto{\pgfqpoint{5.734713in}{2.301545in}}%
\pgfpathlineto{\pgfqpoint{5.734835in}{2.301624in}}%
\pgfpathlineto{\pgfqpoint{5.734903in}{2.300446in}}%
\pgfpathlineto{\pgfqpoint{5.735657in}{2.304643in}}%
\pgfpathlineto{\pgfqpoint{5.735789in}{2.303516in}}%
\pgfpathlineto{\pgfqpoint{5.736383in}{2.306400in}}%
\pgfpathlineto{\pgfqpoint{5.735813in}{2.302909in}}%
\pgfpathlineto{\pgfqpoint{5.736904in}{2.304014in}}%
\pgfpathlineto{\pgfqpoint{5.737094in}{2.302266in}}%
\pgfpathlineto{\pgfqpoint{5.737644in}{2.304189in}}%
\pgfpathlineto{\pgfqpoint{5.738033in}{2.303261in}}%
\pgfpathlineto{\pgfqpoint{5.738359in}{2.304958in}}%
\pgfpathlineto{\pgfqpoint{5.738754in}{2.301784in}}%
\pgfpathlineto{\pgfqpoint{5.739153in}{2.303643in}}%
\pgfpathlineto{\pgfqpoint{5.740136in}{2.304721in}}%
\pgfpathlineto{\pgfqpoint{5.739241in}{2.302072in}}%
\pgfpathlineto{\pgfqpoint{5.740273in}{2.303881in}}%
\pgfpathlineto{\pgfqpoint{5.740521in}{2.301638in}}%
\pgfpathlineto{\pgfqpoint{5.741081in}{2.305231in}}%
\pgfpathlineto{\pgfqpoint{5.741392in}{2.303745in}}%
\pgfpathlineto{\pgfqpoint{5.742376in}{2.298806in}}%
\pgfpathlineto{\pgfqpoint{5.741509in}{2.304495in}}%
\pgfpathlineto{\pgfqpoint{5.742785in}{2.300983in}}%
\pgfpathlineto{\pgfqpoint{5.743953in}{2.304520in}}%
\pgfpathlineto{\pgfqpoint{5.742911in}{2.300657in}}%
\pgfpathlineto{\pgfqpoint{5.743963in}{2.304272in}}%
\pgfpathlineto{\pgfqpoint{5.744547in}{2.302611in}}%
\pgfpathlineto{\pgfqpoint{5.744839in}{2.306100in}}%
\pgfpathlineto{\pgfqpoint{5.745073in}{2.304003in}}%
\pgfpathlineto{\pgfqpoint{5.745161in}{2.304669in}}%
\pgfpathlineto{\pgfqpoint{5.745974in}{2.301251in}}%
\pgfpathlineto{\pgfqpoint{5.746144in}{2.302514in}}%
\pgfpathlineto{\pgfqpoint{5.746699in}{2.303766in}}%
\pgfpathlineto{\pgfqpoint{5.746850in}{2.301247in}}%
\pgfpathlineto{\pgfqpoint{5.746889in}{2.301299in}}%
\pgfpathlineto{\pgfqpoint{5.747454in}{2.299079in}}%
\pgfpathlineto{\pgfqpoint{5.747283in}{2.301887in}}%
\pgfpathlineto{\pgfqpoint{5.748018in}{2.299633in}}%
\pgfpathlineto{\pgfqpoint{5.749187in}{2.303783in}}%
\pgfpathlineto{\pgfqpoint{5.748286in}{2.299266in}}%
\pgfpathlineto{\pgfqpoint{5.749450in}{2.302956in}}%
\pgfpathlineto{\pgfqpoint{5.750185in}{2.302093in}}%
\pgfpathlineto{\pgfqpoint{5.749635in}{2.304089in}}%
\pgfpathlineto{\pgfqpoint{5.750560in}{2.302539in}}%
\pgfpathlineto{\pgfqpoint{5.751495in}{2.303326in}}%
\pgfpathlineto{\pgfqpoint{5.750940in}{2.298393in}}%
\pgfpathlineto{\pgfqpoint{5.751694in}{2.303222in}}%
\pgfpathlineto{\pgfqpoint{5.753471in}{2.298233in}}%
\pgfpathlineto{\pgfqpoint{5.751709in}{2.303596in}}%
\pgfpathlineto{\pgfqpoint{5.753646in}{2.299569in}}%
\pgfpathlineto{\pgfqpoint{5.754211in}{2.302863in}}%
\pgfpathlineto{\pgfqpoint{5.753656in}{2.299475in}}%
\pgfpathlineto{\pgfqpoint{5.754844in}{2.302483in}}%
\pgfpathlineto{\pgfqpoint{5.755151in}{2.301579in}}%
\pgfpathlineto{\pgfqpoint{5.755331in}{2.304333in}}%
\pgfpathlineto{\pgfqpoint{5.755935in}{2.303271in}}%
\pgfpathlineto{\pgfqpoint{5.756889in}{2.304066in}}%
\pgfpathlineto{\pgfqpoint{5.756684in}{2.300163in}}%
\pgfpathlineto{\pgfqpoint{5.757069in}{2.303594in}}%
\pgfpathlineto{\pgfqpoint{5.757556in}{2.302364in}}%
\pgfpathlineto{\pgfqpoint{5.758106in}{2.304236in}}%
\pgfpathlineto{\pgfqpoint{5.758174in}{2.303802in}}%
\pgfpathlineto{\pgfqpoint{5.759289in}{2.302331in}}%
\pgfpathlineto{\pgfqpoint{5.758890in}{2.305230in}}%
\pgfpathlineto{\pgfqpoint{5.759309in}{2.303259in}}%
\pgfpathlineto{\pgfqpoint{5.759961in}{2.304128in}}%
\pgfpathlineto{\pgfqpoint{5.760384in}{2.301830in}}%
\pgfpathlineto{\pgfqpoint{5.760764in}{2.303367in}}%
\pgfpathlineto{\pgfqpoint{5.761183in}{2.299724in}}%
\pgfpathlineto{\pgfqpoint{5.761353in}{2.300640in}}%
\pgfpathlineto{\pgfqpoint{5.761373in}{2.300304in}}%
\pgfpathlineto{\pgfqpoint{5.762327in}{2.306104in}}%
\pgfpathlineto{\pgfqpoint{5.762371in}{2.305604in}}%
\pgfpathlineto{\pgfqpoint{5.762595in}{2.305954in}}%
\pgfpathlineto{\pgfqpoint{5.762814in}{2.303077in}}%
\pgfpathlineto{\pgfqpoint{5.763437in}{2.304193in}}%
\pgfpathlineto{\pgfqpoint{5.764075in}{2.300796in}}%
\pgfpathlineto{\pgfqpoint{5.763656in}{2.305366in}}%
\pgfpathlineto{\pgfqpoint{5.764596in}{2.303740in}}%
\pgfpathlineto{\pgfqpoint{5.765297in}{2.304743in}}%
\pgfpathlineto{\pgfqpoint{5.764868in}{2.302422in}}%
\pgfpathlineto{\pgfqpoint{5.765691in}{2.303461in}}%
\pgfpathlineto{\pgfqpoint{5.765750in}{2.302351in}}%
\pgfpathlineto{\pgfqpoint{5.766222in}{2.306458in}}%
\pgfpathlineto{\pgfqpoint{5.766791in}{2.303598in}}%
\pgfpathlineto{\pgfqpoint{5.767181in}{2.305626in}}%
\pgfpathlineto{\pgfqpoint{5.767770in}{2.301615in}}%
\pgfpathlineto{\pgfqpoint{5.767794in}{2.301742in}}%
\pgfpathlineto{\pgfqpoint{5.767799in}{2.301668in}}%
\pgfpathlineto{\pgfqpoint{5.768169in}{2.303936in}}%
\pgfpathlineto{\pgfqpoint{5.768880in}{2.302028in}}%
\pgfpathlineto{\pgfqpoint{5.769372in}{2.304358in}}%
\pgfpathlineto{\pgfqpoint{5.769927in}{2.299746in}}%
\pgfpathlineto{\pgfqpoint{5.769966in}{2.299974in}}%
\pgfpathlineto{\pgfqpoint{5.770350in}{2.297189in}}%
\pgfpathlineto{\pgfqpoint{5.770978in}{2.301695in}}%
\pgfpathlineto{\pgfqpoint{5.770993in}{2.301579in}}%
\pgfpathlineto{\pgfqpoint{5.771363in}{2.300673in}}%
\pgfpathlineto{\pgfqpoint{5.772142in}{2.303279in}}%
\pgfpathlineto{\pgfqpoint{5.772220in}{2.304475in}}%
\pgfpathlineto{\pgfqpoint{5.772994in}{2.299124in}}%
\pgfpathlineto{\pgfqpoint{5.772999in}{2.299208in}}%
\pgfpathlineto{\pgfqpoint{5.773729in}{2.297697in}}%
\pgfpathlineto{\pgfqpoint{5.773291in}{2.300713in}}%
\pgfpathlineto{\pgfqpoint{5.774109in}{2.299117in}}%
\pgfpathlineto{\pgfqpoint{5.775175in}{2.300432in}}%
\pgfpathlineto{\pgfqpoint{5.774839in}{2.297420in}}%
\pgfpathlineto{\pgfqpoint{5.775214in}{2.299039in}}%
\pgfpathlineto{\pgfqpoint{5.775404in}{2.298750in}}%
\pgfpathlineto{\pgfqpoint{5.775861in}{2.302053in}}%
\pgfpathlineto{\pgfqpoint{5.776090in}{2.300971in}}%
\pgfpathlineto{\pgfqpoint{5.776928in}{2.303629in}}%
\pgfpathlineto{\pgfqpoint{5.776149in}{2.300339in}}%
\pgfpathlineto{\pgfqpoint{5.777234in}{2.302377in}}%
\pgfpathlineto{\pgfqpoint{5.777239in}{2.302422in}}%
\pgfpathlineto{\pgfqpoint{5.777575in}{2.299721in}}%
\pgfpathlineto{\pgfqpoint{5.778276in}{2.301830in}}%
\pgfpathlineto{\pgfqpoint{5.779323in}{2.299155in}}%
\pgfpathlineto{\pgfqpoint{5.778529in}{2.302196in}}%
\pgfpathlineto{\pgfqpoint{5.779411in}{2.299905in}}%
\pgfpathlineto{\pgfqpoint{5.780521in}{2.301611in}}%
\pgfpathlineto{\pgfqpoint{5.779776in}{2.298452in}}%
\pgfpathlineto{\pgfqpoint{5.780545in}{2.301381in}}%
\pgfpathlineto{\pgfqpoint{5.780939in}{2.299851in}}%
\pgfpathlineto{\pgfqpoint{5.781363in}{2.302279in}}%
\pgfpathlineto{\pgfqpoint{5.781485in}{2.304204in}}%
\pgfpathlineto{\pgfqpoint{5.782264in}{2.299471in}}%
\pgfpathlineto{\pgfqpoint{5.782453in}{2.300775in}}%
\pgfpathlineto{\pgfqpoint{5.783193in}{2.296668in}}%
\pgfpathlineto{\pgfqpoint{5.782604in}{2.301915in}}%
\pgfpathlineto{\pgfqpoint{5.783632in}{2.299087in}}%
\pgfpathlineto{\pgfqpoint{5.784016in}{2.300578in}}%
\pgfpathlineto{\pgfqpoint{5.784274in}{2.298042in}}%
\pgfpathlineto{\pgfqpoint{5.784742in}{2.299254in}}%
\pgfpathlineto{\pgfqpoint{5.784970in}{2.297095in}}%
\pgfpathlineto{\pgfqpoint{5.785598in}{2.299480in}}%
\pgfpathlineto{\pgfqpoint{5.785861in}{2.298325in}}%
\pgfpathlineto{\pgfqpoint{5.786129in}{2.299075in}}%
\pgfpathlineto{\pgfqpoint{5.786519in}{2.293760in}}%
\pgfpathlineto{\pgfqpoint{5.786947in}{2.296999in}}%
\pgfpathlineto{\pgfqpoint{5.787302in}{2.299487in}}%
\pgfpathlineto{\pgfqpoint{5.787687in}{2.295720in}}%
\pgfpathlineto{\pgfqpoint{5.788164in}{2.293343in}}%
\pgfpathlineto{\pgfqpoint{5.788714in}{2.296196in}}%
\pgfpathlineto{\pgfqpoint{5.788782in}{2.296022in}}%
\pgfpathlineto{\pgfqpoint{5.789776in}{2.298344in}}%
\pgfpathlineto{\pgfqpoint{5.789255in}{2.295084in}}%
\pgfpathlineto{\pgfqpoint{5.789917in}{2.297402in}}%
\pgfpathlineto{\pgfqpoint{5.790501in}{2.296138in}}%
\pgfpathlineto{\pgfqpoint{5.790170in}{2.298674in}}%
\pgfpathlineto{\pgfqpoint{5.790842in}{2.297867in}}%
\pgfpathlineto{\pgfqpoint{5.791153in}{2.298662in}}%
\pgfpathlineto{\pgfqpoint{5.791816in}{2.292210in}}%
\pgfpathlineto{\pgfqpoint{5.791845in}{2.293183in}}%
\pgfpathlineto{\pgfqpoint{5.792945in}{2.292239in}}%
\pgfpathlineto{\pgfqpoint{5.792419in}{2.294308in}}%
\pgfpathlineto{\pgfqpoint{5.792974in}{2.292536in}}%
\pgfpathlineto{\pgfqpoint{5.794148in}{2.296464in}}%
\pgfpathlineto{\pgfqpoint{5.793213in}{2.292138in}}%
\pgfpathlineto{\pgfqpoint{5.794216in}{2.295339in}}%
\pgfpathlineto{\pgfqpoint{5.794897in}{2.291171in}}%
\pgfpathlineto{\pgfqpoint{5.795345in}{2.294042in}}%
\pgfpathlineto{\pgfqpoint{5.795408in}{2.293904in}}%
\pgfpathlineto{\pgfqpoint{5.795949in}{2.298108in}}%
\pgfpathlineto{\pgfqpoint{5.796348in}{2.296506in}}%
\pgfpathlineto{\pgfqpoint{5.797117in}{2.299301in}}%
\pgfpathlineto{\pgfqpoint{5.797497in}{2.298640in}}%
\pgfpathlineto{\pgfqpoint{5.799303in}{2.304583in}}%
\pgfpathlineto{\pgfqpoint{5.797555in}{2.298001in}}%
\pgfpathlineto{\pgfqpoint{5.799386in}{2.303438in}}%
\pgfpathlineto{\pgfqpoint{5.800739in}{2.299196in}}%
\pgfpathlineto{\pgfqpoint{5.799727in}{2.304576in}}%
\pgfpathlineto{\pgfqpoint{5.800744in}{2.299362in}}%
\pgfpathlineto{\pgfqpoint{5.801849in}{2.300676in}}%
\pgfpathlineto{\pgfqpoint{5.801236in}{2.296912in}}%
\pgfpathlineto{\pgfqpoint{5.801884in}{2.300242in}}%
\pgfpathlineto{\pgfqpoint{5.802692in}{2.299571in}}%
\pgfpathlineto{\pgfqpoint{5.802293in}{2.303181in}}%
\pgfpathlineto{\pgfqpoint{5.802901in}{2.302574in}}%
\pgfpathlineto{\pgfqpoint{5.803154in}{2.300375in}}%
\pgfpathlineto{\pgfqpoint{5.804133in}{2.306150in}}%
\pgfpathlineto{\pgfqpoint{5.805282in}{2.302035in}}%
\pgfpathlineto{\pgfqpoint{5.804177in}{2.306222in}}%
\pgfpathlineto{\pgfqpoint{5.805350in}{2.302777in}}%
\pgfpathlineto{\pgfqpoint{5.805895in}{2.305241in}}%
\pgfpathlineto{\pgfqpoint{5.805506in}{2.301483in}}%
\pgfpathlineto{\pgfqpoint{5.806509in}{2.304477in}}%
\pgfpathlineto{\pgfqpoint{5.807521in}{2.299661in}}%
\pgfpathlineto{\pgfqpoint{5.808232in}{2.302844in}}%
\pgfpathlineto{\pgfqpoint{5.808281in}{2.303757in}}%
\pgfpathlineto{\pgfqpoint{5.808933in}{2.301026in}}%
\pgfpathlineto{\pgfqpoint{5.809352in}{2.303331in}}%
\pgfpathlineto{\pgfqpoint{5.810262in}{2.302078in}}%
\pgfpathlineto{\pgfqpoint{5.809590in}{2.304591in}}%
\pgfpathlineto{\pgfqpoint{5.810457in}{2.303243in}}%
\pgfpathlineto{\pgfqpoint{5.810686in}{2.303939in}}%
\pgfpathlineto{\pgfqpoint{5.811450in}{2.301203in}}%
\pgfpathlineto{\pgfqpoint{5.811513in}{2.301340in}}%
\pgfpathlineto{\pgfqpoint{5.811601in}{2.300571in}}%
\pgfpathlineto{\pgfqpoint{5.812161in}{2.303963in}}%
\pgfpathlineto{\pgfqpoint{5.812507in}{2.302520in}}%
\pgfpathlineto{\pgfqpoint{5.813563in}{2.304018in}}%
\pgfpathlineto{\pgfqpoint{5.812989in}{2.299914in}}%
\pgfpathlineto{\pgfqpoint{5.813626in}{2.303380in}}%
\pgfpathlineto{\pgfqpoint{5.814396in}{2.301433in}}%
\pgfpathlineto{\pgfqpoint{5.813719in}{2.304303in}}%
\pgfpathlineto{\pgfqpoint{5.814717in}{2.303667in}}%
\pgfpathlineto{\pgfqpoint{5.815661in}{2.305663in}}%
\pgfpathlineto{\pgfqpoint{5.814965in}{2.302994in}}%
\pgfpathlineto{\pgfqpoint{5.815861in}{2.304361in}}%
\pgfpathlineto{\pgfqpoint{5.816913in}{2.302481in}}%
\pgfpathlineto{\pgfqpoint{5.815876in}{2.304711in}}%
\pgfpathlineto{\pgfqpoint{5.817015in}{2.303426in}}%
\pgfpathlineto{\pgfqpoint{5.817998in}{2.302147in}}%
\pgfpathlineto{\pgfqpoint{5.817229in}{2.304461in}}%
\pgfpathlineto{\pgfqpoint{5.818115in}{2.303593in}}%
\pgfpathlineto{\pgfqpoint{5.818373in}{2.304788in}}%
\pgfpathlineto{\pgfqpoint{5.818485in}{2.301934in}}%
\pgfpathlineto{\pgfqpoint{5.819225in}{2.303795in}}%
\pgfpathlineto{\pgfqpoint{5.820428in}{2.301054in}}%
\pgfpathlineto{\pgfqpoint{5.819361in}{2.305018in}}%
\pgfpathlineto{\pgfqpoint{5.820442in}{2.301369in}}%
\pgfpathlineto{\pgfqpoint{5.821109in}{2.304185in}}%
\pgfpathlineto{\pgfqpoint{5.821499in}{2.300457in}}%
\pgfpathlineto{\pgfqpoint{5.821956in}{2.298480in}}%
\pgfpathlineto{\pgfqpoint{5.821684in}{2.301992in}}%
\pgfpathlineto{\pgfqpoint{5.822614in}{2.300213in}}%
\pgfpathlineto{\pgfqpoint{5.822662in}{2.299976in}}%
\pgfpathlineto{\pgfqpoint{5.823027in}{2.303821in}}%
\pgfpathlineto{\pgfqpoint{5.823363in}{2.302825in}}%
\pgfpathlineto{\pgfqpoint{5.824288in}{2.304057in}}%
\pgfpathlineto{\pgfqpoint{5.823797in}{2.301325in}}%
\pgfpathlineto{\pgfqpoint{5.824454in}{2.301433in}}%
\pgfpathlineto{\pgfqpoint{5.825613in}{2.304714in}}%
\pgfpathlineto{\pgfqpoint{5.824542in}{2.299905in}}%
\pgfpathlineto{\pgfqpoint{5.825676in}{2.304032in}}%
\pgfpathlineto{\pgfqpoint{5.825997in}{2.305506in}}%
\pgfpathlineto{\pgfqpoint{5.825764in}{2.303477in}}%
\pgfpathlineto{\pgfqpoint{5.826260in}{2.303357in}}%
\pgfpathlineto{\pgfqpoint{5.826309in}{2.302659in}}%
\pgfpathlineto{\pgfqpoint{5.826611in}{2.305342in}}%
\pgfpathlineto{\pgfqpoint{5.827365in}{2.303894in}}%
\pgfpathlineto{\pgfqpoint{5.828183in}{2.298189in}}%
\pgfpathlineto{\pgfqpoint{5.827380in}{2.303967in}}%
\pgfpathlineto{\pgfqpoint{5.828529in}{2.301458in}}%
\pgfpathlineto{\pgfqpoint{5.829522in}{2.297441in}}%
\pgfpathlineto{\pgfqpoint{5.828548in}{2.302313in}}%
\pgfpathlineto{\pgfqpoint{5.829785in}{2.298026in}}%
\pgfpathlineto{\pgfqpoint{5.830087in}{2.298498in}}%
\pgfpathlineto{\pgfqpoint{5.830501in}{2.294815in}}%
\pgfpathlineto{\pgfqpoint{5.830837in}{2.297234in}}%
\pgfpathlineto{\pgfqpoint{5.830919in}{2.296244in}}%
\pgfpathlineto{\pgfqpoint{5.831849in}{2.302989in}}%
\pgfpathlineto{\pgfqpoint{5.831888in}{2.302346in}}%
\pgfpathlineto{\pgfqpoint{5.832219in}{2.303873in}}%
\pgfpathlineto{\pgfqpoint{5.832740in}{2.298923in}}%
\pgfpathlineto{\pgfqpoint{5.832940in}{2.299596in}}%
\pgfpathlineto{\pgfqpoint{5.833446in}{2.301944in}}%
\pgfpathlineto{\pgfqpoint{5.833957in}{2.298294in}}%
\pgfpathlineto{\pgfqpoint{5.834084in}{2.300851in}}%
\pgfpathlineto{\pgfqpoint{5.835111in}{2.296785in}}%
\pgfpathlineto{\pgfqpoint{5.834376in}{2.302412in}}%
\pgfpathlineto{\pgfqpoint{5.835238in}{2.298407in}}%
\pgfpathlineto{\pgfqpoint{5.836547in}{2.295025in}}%
\pgfpathlineto{\pgfqpoint{5.835282in}{2.299175in}}%
\pgfpathlineto{\pgfqpoint{5.836606in}{2.295883in}}%
\pgfpathlineto{\pgfqpoint{5.836888in}{2.298216in}}%
\pgfpathlineto{\pgfqpoint{5.836640in}{2.295838in}}%
\pgfpathlineto{\pgfqpoint{5.837842in}{2.297157in}}%
\pgfpathlineto{\pgfqpoint{5.838212in}{2.294640in}}%
\pgfpathlineto{\pgfqpoint{5.838762in}{2.298141in}}%
\pgfpathlineto{\pgfqpoint{5.838957in}{2.296494in}}%
\pgfpathlineto{\pgfqpoint{5.839162in}{2.298304in}}%
\pgfpathlineto{\pgfqpoint{5.839965in}{2.294950in}}%
\pgfpathlineto{\pgfqpoint{5.839970in}{2.294964in}}%
\pgfpathlineto{\pgfqpoint{5.840793in}{2.292918in}}%
\pgfpathlineto{\pgfqpoint{5.840247in}{2.297104in}}%
\pgfpathlineto{\pgfqpoint{5.841085in}{2.294653in}}%
\pgfpathlineto{\pgfqpoint{5.842097in}{2.297862in}}%
\pgfpathlineto{\pgfqpoint{5.841206in}{2.293987in}}%
\pgfpathlineto{\pgfqpoint{5.842239in}{2.296930in}}%
\pgfpathlineto{\pgfqpoint{5.843538in}{2.302915in}}%
\pgfpathlineto{\pgfqpoint{5.843728in}{2.302064in}}%
\pgfpathlineto{\pgfqpoint{5.844687in}{2.297795in}}%
\pgfpathlineto{\pgfqpoint{5.843787in}{2.302468in}}%
\pgfpathlineto{\pgfqpoint{5.844877in}{2.299406in}}%
\pgfpathlineto{\pgfqpoint{5.845004in}{2.298800in}}%
\pgfpathlineto{\pgfqpoint{5.845744in}{2.301196in}}%
\pgfpathlineto{\pgfqpoint{5.845944in}{2.300901in}}%
\pgfpathlineto{\pgfqpoint{5.846791in}{2.302375in}}%
\pgfpathlineto{\pgfqpoint{5.846421in}{2.299175in}}%
\pgfpathlineto{\pgfqpoint{5.847092in}{2.301930in}}%
\pgfpathlineto{\pgfqpoint{5.847531in}{2.300004in}}%
\pgfpathlineto{\pgfqpoint{5.848129in}{2.302953in}}%
\pgfpathlineto{\pgfqpoint{5.848134in}{2.302951in}}%
\pgfpathlineto{\pgfqpoint{5.848509in}{2.304817in}}%
\pgfpathlineto{\pgfqpoint{5.849025in}{2.299756in}}%
\pgfpathlineto{\pgfqpoint{5.849166in}{2.300282in}}%
\pgfpathlineto{\pgfqpoint{5.849425in}{2.299992in}}%
\pgfpathlineto{\pgfqpoint{5.849580in}{2.302244in}}%
\pgfpathlineto{\pgfqpoint{5.849795in}{2.302098in}}%
\pgfpathlineto{\pgfqpoint{5.850345in}{2.304743in}}%
\pgfpathlineto{\pgfqpoint{5.849833in}{2.301490in}}%
\pgfpathlineto{\pgfqpoint{5.850905in}{2.302145in}}%
\pgfpathlineto{\pgfqpoint{5.850968in}{2.300830in}}%
\pgfpathlineto{\pgfqpoint{5.851313in}{2.303942in}}%
\pgfpathlineto{\pgfqpoint{5.851601in}{2.303301in}}%
\pgfpathlineto{\pgfqpoint{5.852039in}{2.299805in}}%
\pgfpathlineto{\pgfqpoint{5.852720in}{2.304303in}}%
\pgfpathlineto{\pgfqpoint{5.853792in}{2.299343in}}%
\pgfpathlineto{\pgfqpoint{5.853884in}{2.300744in}}%
\pgfpathlineto{\pgfqpoint{5.854249in}{2.301409in}}%
\pgfpathlineto{\pgfqpoint{5.854804in}{2.296859in}}%
\pgfpathlineto{\pgfqpoint{5.854838in}{2.297522in}}%
\pgfpathlineto{\pgfqpoint{5.855379in}{2.295798in}}%
\pgfpathlineto{\pgfqpoint{5.855807in}{2.299556in}}%
\pgfpathlineto{\pgfqpoint{5.855890in}{2.298516in}}%
\pgfpathlineto{\pgfqpoint{5.856046in}{2.299315in}}%
\pgfpathlineto{\pgfqpoint{5.856859in}{2.296300in}}%
\pgfpathlineto{\pgfqpoint{5.856980in}{2.296655in}}%
\pgfpathlineto{\pgfqpoint{5.857716in}{2.301326in}}%
\pgfpathlineto{\pgfqpoint{5.857102in}{2.296388in}}%
\pgfpathlineto{\pgfqpoint{5.858188in}{2.298853in}}%
\pgfpathlineto{\pgfqpoint{5.859244in}{2.296576in}}%
\pgfpathlineto{\pgfqpoint{5.858957in}{2.299340in}}%
\pgfpathlineto{\pgfqpoint{5.859361in}{2.297646in}}%
\pgfpathlineto{\pgfqpoint{5.859600in}{2.299013in}}%
\pgfpathlineto{\pgfqpoint{5.860257in}{2.293464in}}%
\pgfpathlineto{\pgfqpoint{5.860393in}{2.293608in}}%
\pgfpathlineto{\pgfqpoint{5.860427in}{2.292781in}}%
\pgfpathlineto{\pgfqpoint{5.861289in}{2.300831in}}%
\pgfpathlineto{\pgfqpoint{5.861411in}{2.298937in}}%
\pgfpathlineto{\pgfqpoint{5.861878in}{2.302309in}}%
\pgfpathlineto{\pgfqpoint{5.862565in}{2.300442in}}%
\pgfpathlineto{\pgfqpoint{5.862832in}{2.301635in}}%
\pgfpathlineto{\pgfqpoint{5.863602in}{2.305257in}}%
\pgfpathlineto{\pgfqpoint{5.864040in}{2.303542in}}%
\pgfpathlineto{\pgfqpoint{5.864663in}{2.301707in}}%
\pgfpathlineto{\pgfqpoint{5.864488in}{2.303938in}}%
\pgfpathlineto{\pgfqpoint{5.865145in}{2.303459in}}%
\pgfpathlineto{\pgfqpoint{5.865491in}{2.305095in}}%
\pgfpathlineto{\pgfqpoint{5.865963in}{2.301689in}}%
\pgfpathlineto{\pgfqpoint{5.866245in}{2.302866in}}%
\pgfpathlineto{\pgfqpoint{5.866781in}{2.301738in}}%
\pgfpathlineto{\pgfqpoint{5.866537in}{2.304270in}}%
\pgfpathlineto{\pgfqpoint{5.867380in}{2.302250in}}%
\pgfpathlineto{\pgfqpoint{5.867662in}{2.304370in}}%
\pgfpathlineto{\pgfqpoint{5.867472in}{2.301793in}}%
\pgfpathlineto{\pgfqpoint{5.868514in}{2.303131in}}%
\pgfpathlineto{\pgfqpoint{5.868567in}{2.303068in}}%
\pgfpathlineto{\pgfqpoint{5.868592in}{2.302001in}}%
\pgfpathlineto{\pgfqpoint{5.868597in}{2.302097in}}%
\pgfpathlineto{\pgfqpoint{5.869522in}{2.300046in}}%
\pgfpathlineto{\pgfqpoint{5.869181in}{2.303400in}}%
\pgfpathlineto{\pgfqpoint{5.869721in}{2.301240in}}%
\pgfpathlineto{\pgfqpoint{5.869755in}{2.302323in}}%
\pgfpathlineto{\pgfqpoint{5.870232in}{2.299400in}}%
\pgfpathlineto{\pgfqpoint{5.870841in}{2.301724in}}%
\pgfpathlineto{\pgfqpoint{5.870875in}{2.302467in}}%
\pgfpathlineto{\pgfqpoint{5.871231in}{2.299275in}}%
\pgfpathlineto{\pgfqpoint{5.871902in}{2.300074in}}%
\pgfpathlineto{\pgfqpoint{5.872745in}{2.304709in}}%
\pgfpathlineto{\pgfqpoint{5.871946in}{2.299964in}}%
\pgfpathlineto{\pgfqpoint{5.873212in}{2.302437in}}%
\pgfpathlineto{\pgfqpoint{5.873246in}{2.301710in}}%
\pgfpathlineto{\pgfqpoint{5.874195in}{2.303844in}}%
\pgfpathlineto{\pgfqpoint{5.874283in}{2.303665in}}%
\pgfpathlineto{\pgfqpoint{5.874765in}{2.304835in}}%
\pgfpathlineto{\pgfqpoint{5.875193in}{2.302129in}}%
\pgfpathlineto{\pgfqpoint{5.875393in}{2.303792in}}%
\pgfpathlineto{\pgfqpoint{5.876162in}{2.299798in}}%
\pgfpathlineto{\pgfqpoint{5.876557in}{2.301150in}}%
\pgfpathlineto{\pgfqpoint{5.877321in}{2.298795in}}%
\pgfpathlineto{\pgfqpoint{5.876630in}{2.301652in}}%
\pgfpathlineto{\pgfqpoint{5.877657in}{2.301480in}}%
\pgfpathlineto{\pgfqpoint{5.878821in}{2.303440in}}%
\pgfpathlineto{\pgfqpoint{5.878037in}{2.300283in}}%
\pgfpathlineto{\pgfqpoint{5.878889in}{2.302536in}}%
\pgfpathlineto{\pgfqpoint{5.879979in}{2.300023in}}%
\pgfpathlineto{\pgfqpoint{5.879517in}{2.303946in}}%
\pgfpathlineto{\pgfqpoint{5.880028in}{2.300881in}}%
\pgfpathlineto{\pgfqpoint{5.880510in}{2.302087in}}%
\pgfpathlineto{\pgfqpoint{5.880860in}{2.299534in}}%
\pgfpathlineto{\pgfqpoint{5.881138in}{2.300922in}}%
\pgfpathlineto{\pgfqpoint{5.881484in}{2.297995in}}%
\pgfpathlineto{\pgfqpoint{5.882297in}{2.299426in}}%
\pgfpathlineto{\pgfqpoint{5.882306in}{2.299306in}}%
\pgfpathlineto{\pgfqpoint{5.882462in}{2.301432in}}%
\pgfpathlineto{\pgfqpoint{5.882618in}{2.301348in}}%
\pgfpathlineto{\pgfqpoint{5.883017in}{2.303498in}}%
\pgfpathlineto{\pgfqpoint{5.882735in}{2.301062in}}%
\pgfpathlineto{\pgfqpoint{5.883752in}{2.302683in}}%
\pgfpathlineto{\pgfqpoint{5.884239in}{2.301259in}}%
\pgfpathlineto{\pgfqpoint{5.884624in}{2.304163in}}%
\pgfpathlineto{\pgfqpoint{5.885632in}{2.306737in}}%
\pgfpathlineto{\pgfqpoint{5.885198in}{2.303045in}}%
\pgfpathlineto{\pgfqpoint{5.885753in}{2.305552in}}%
\pgfpathlineto{\pgfqpoint{5.886854in}{2.300182in}}%
\pgfpathlineto{\pgfqpoint{5.885792in}{2.305733in}}%
\pgfpathlineto{\pgfqpoint{5.887053in}{2.301538in}}%
\pgfpathlineto{\pgfqpoint{5.888144in}{2.304329in}}%
\pgfpathlineto{\pgfqpoint{5.888183in}{2.304078in}}%
\pgfpathlineto{\pgfqpoint{5.889132in}{2.300033in}}%
\pgfpathlineto{\pgfqpoint{5.888480in}{2.305211in}}%
\pgfpathlineto{\pgfqpoint{5.889312in}{2.301876in}}%
\pgfpathlineto{\pgfqpoint{5.890451in}{2.305621in}}%
\pgfpathlineto{\pgfqpoint{5.889896in}{2.301744in}}%
\pgfpathlineto{\pgfqpoint{5.890495in}{2.304691in}}%
\pgfpathlineto{\pgfqpoint{5.891221in}{2.301564in}}%
\pgfpathlineto{\pgfqpoint{5.890768in}{2.306253in}}%
\pgfpathlineto{\pgfqpoint{5.891693in}{2.302031in}}%
\pgfpathlineto{\pgfqpoint{5.892808in}{2.305373in}}%
\pgfpathlineto{\pgfqpoint{5.892019in}{2.301640in}}%
\pgfpathlineto{\pgfqpoint{5.892915in}{2.304569in}}%
\pgfpathlineto{\pgfqpoint{5.892920in}{2.304595in}}%
\pgfpathlineto{\pgfqpoint{5.893105in}{2.302652in}}%
\pgfpathlineto{\pgfqpoint{5.893806in}{2.303303in}}%
\pgfpathlineto{\pgfqpoint{5.894190in}{2.301972in}}%
\pgfpathlineto{\pgfqpoint{5.894765in}{2.304341in}}%
\pgfpathlineto{\pgfqpoint{5.894911in}{2.303599in}}%
\pgfpathlineto{\pgfqpoint{5.896259in}{2.301717in}}%
\pgfpathlineto{\pgfqpoint{5.895544in}{2.305004in}}%
\pgfpathlineto{\pgfqpoint{5.896274in}{2.302185in}}%
\pgfpathlineto{\pgfqpoint{5.896678in}{2.304967in}}%
\pgfpathlineto{\pgfqpoint{5.896323in}{2.301845in}}%
\pgfpathlineto{\pgfqpoint{5.897477in}{2.303801in}}%
\pgfpathlineto{\pgfqpoint{5.898436in}{2.301769in}}%
\pgfpathlineto{\pgfqpoint{5.897720in}{2.304986in}}%
\pgfpathlineto{\pgfqpoint{5.898601in}{2.302629in}}%
\pgfpathlineto{\pgfqpoint{5.899171in}{2.303310in}}%
\pgfpathlineto{\pgfqpoint{5.898811in}{2.301341in}}%
\pgfpathlineto{\pgfqpoint{5.899536in}{2.302242in}}%
\pgfpathlineto{\pgfqpoint{5.900485in}{2.299191in}}%
\pgfpathlineto{\pgfqpoint{5.899609in}{2.303032in}}%
\pgfpathlineto{\pgfqpoint{5.900661in}{2.300982in}}%
\pgfpathlineto{\pgfqpoint{5.900709in}{2.301718in}}%
\pgfpathlineto{\pgfqpoint{5.900841in}{2.299594in}}%
\pgfpathlineto{\pgfqpoint{5.901294in}{2.299724in}}%
\pgfpathlineto{\pgfqpoint{5.901980in}{2.297938in}}%
\pgfpathlineto{\pgfqpoint{5.901376in}{2.300728in}}%
\pgfpathlineto{\pgfqpoint{5.902389in}{2.299906in}}%
\pgfpathlineto{\pgfqpoint{5.902803in}{2.299022in}}%
\pgfpathlineto{\pgfqpoint{5.903567in}{2.305468in}}%
\pgfpathlineto{\pgfqpoint{5.904682in}{2.303596in}}%
\pgfpathlineto{\pgfqpoint{5.903849in}{2.306659in}}%
\pgfpathlineto{\pgfqpoint{5.904886in}{2.303652in}}%
\pgfpathlineto{\pgfqpoint{5.905456in}{2.306174in}}%
\pgfpathlineto{\pgfqpoint{5.905880in}{2.302594in}}%
\pgfpathlineto{\pgfqpoint{5.905992in}{2.303099in}}%
\pgfpathlineto{\pgfqpoint{5.906995in}{2.301660in}}%
\pgfpathlineto{\pgfqpoint{5.906113in}{2.304971in}}%
\pgfpathlineto{\pgfqpoint{5.907102in}{2.302923in}}%
\pgfpathlineto{\pgfqpoint{5.907218in}{2.303047in}}%
\pgfpathlineto{\pgfqpoint{5.907588in}{2.299529in}}%
\pgfpathlineto{\pgfqpoint{5.908173in}{2.302381in}}%
\pgfpathlineto{\pgfqpoint{5.908212in}{2.301134in}}%
\pgfpathlineto{\pgfqpoint{5.908884in}{2.304946in}}%
\pgfpathlineto{\pgfqpoint{5.909263in}{2.303512in}}%
\pgfpathlineto{\pgfqpoint{5.909711in}{2.304223in}}%
\pgfpathlineto{\pgfqpoint{5.910057in}{2.301114in}}%
\pgfpathlineto{\pgfqpoint{5.910281in}{2.301704in}}%
\pgfpathlineto{\pgfqpoint{5.911328in}{2.300552in}}%
\pgfpathlineto{\pgfqpoint{5.910826in}{2.304139in}}%
\pgfpathlineto{\pgfqpoint{5.911396in}{2.301416in}}%
\pgfpathlineto{\pgfqpoint{5.911610in}{2.302760in}}%
\pgfpathlineto{\pgfqpoint{5.912442in}{2.298799in}}%
\pgfpathlineto{\pgfqpoint{5.912525in}{2.297419in}}%
\pgfpathlineto{\pgfqpoint{5.913484in}{2.301086in}}%
\pgfpathlineto{\pgfqpoint{5.913489in}{2.300951in}}%
\pgfpathlineto{\pgfqpoint{5.914249in}{2.302290in}}%
\pgfpathlineto{\pgfqpoint{5.914419in}{2.299899in}}%
\pgfpathlineto{\pgfqpoint{5.914560in}{2.300187in}}%
\pgfpathlineto{\pgfqpoint{5.915290in}{2.298323in}}%
\pgfpathlineto{\pgfqpoint{5.914857in}{2.302267in}}%
\pgfpathlineto{\pgfqpoint{5.915665in}{2.300410in}}%
\pgfpathlineto{\pgfqpoint{5.916683in}{2.299301in}}%
\pgfpathlineto{\pgfqpoint{5.915889in}{2.302469in}}%
\pgfpathlineto{\pgfqpoint{5.916814in}{2.299748in}}%
\pgfpathlineto{\pgfqpoint{5.917939in}{2.301321in}}%
\pgfpathlineto{\pgfqpoint{5.917024in}{2.297897in}}%
\pgfpathlineto{\pgfqpoint{5.917973in}{2.300913in}}%
\pgfpathlineto{\pgfqpoint{5.917988in}{2.300616in}}%
\pgfpathlineto{\pgfqpoint{5.918440in}{2.304119in}}%
\pgfpathlineto{\pgfqpoint{5.919073in}{2.301244in}}%
\pgfpathlineto{\pgfqpoint{5.919331in}{2.303080in}}%
\pgfpathlineto{\pgfqpoint{5.919468in}{2.300407in}}%
\pgfpathlineto{\pgfqpoint{5.920198in}{2.301568in}}%
\pgfpathlineto{\pgfqpoint{5.920889in}{2.300422in}}%
\pgfpathlineto{\pgfqpoint{5.920573in}{2.303767in}}%
\pgfpathlineto{\pgfqpoint{5.921220in}{2.303042in}}%
\pgfpathlineto{\pgfqpoint{5.921254in}{2.304000in}}%
\pgfpathlineto{\pgfqpoint{5.922306in}{2.301459in}}%
\pgfpathlineto{\pgfqpoint{5.922613in}{2.300721in}}%
\pgfpathlineto{\pgfqpoint{5.923119in}{2.303798in}}%
\pgfpathlineto{\pgfqpoint{5.923397in}{2.302556in}}%
\pgfpathlineto{\pgfqpoint{5.924482in}{2.300500in}}%
\pgfpathlineto{\pgfqpoint{5.923655in}{2.303288in}}%
\pgfpathlineto{\pgfqpoint{5.924575in}{2.301687in}}%
\pgfpathlineto{\pgfqpoint{5.925397in}{2.304328in}}%
\pgfpathlineto{\pgfqpoint{5.925699in}{2.303278in}}%
\pgfpathlineto{\pgfqpoint{5.926717in}{2.299860in}}%
\pgfpathlineto{\pgfqpoint{5.926839in}{2.300920in}}%
\pgfpathlineto{\pgfqpoint{5.927676in}{2.303752in}}%
\pgfpathlineto{\pgfqpoint{5.927024in}{2.298631in}}%
\pgfpathlineto{\pgfqpoint{5.928031in}{2.301694in}}%
\pgfpathlineto{\pgfqpoint{5.928036in}{2.301467in}}%
\pgfpathlineto{\pgfqpoint{5.929073in}{2.304856in}}%
\pgfpathlineto{\pgfqpoint{5.929127in}{2.305454in}}%
\pgfpathlineto{\pgfqpoint{5.929847in}{2.301498in}}%
\pgfpathlineto{\pgfqpoint{5.929881in}{2.301872in}}%
\pgfpathlineto{\pgfqpoint{5.930913in}{2.297297in}}%
\pgfpathlineto{\pgfqpoint{5.929959in}{2.302288in}}%
\pgfpathlineto{\pgfqpoint{5.931172in}{2.298145in}}%
\pgfpathlineto{\pgfqpoint{5.931556in}{2.298836in}}%
\pgfpathlineto{\pgfqpoint{5.931980in}{2.295998in}}%
\pgfpathlineto{\pgfqpoint{5.932228in}{2.296265in}}%
\pgfpathlineto{\pgfqpoint{5.933104in}{2.301618in}}%
\pgfpathlineto{\pgfqpoint{5.932286in}{2.296217in}}%
\pgfpathlineto{\pgfqpoint{5.933791in}{2.299970in}}%
\pgfpathlineto{\pgfqpoint{5.934643in}{2.298633in}}%
\pgfpathlineto{\pgfqpoint{5.934287in}{2.301123in}}%
\pgfpathlineto{\pgfqpoint{5.934789in}{2.300764in}}%
\pgfpathlineto{\pgfqpoint{5.935816in}{2.303815in}}%
\pgfpathlineto{\pgfqpoint{5.935923in}{2.302107in}}%
\pgfpathlineto{\pgfqpoint{5.936848in}{2.304278in}}%
\pgfpathlineto{\pgfqpoint{5.936556in}{2.301143in}}%
\pgfpathlineto{\pgfqpoint{5.937150in}{2.303081in}}%
\pgfpathlineto{\pgfqpoint{5.937948in}{2.301928in}}%
\pgfpathlineto{\pgfqpoint{5.937374in}{2.304802in}}%
\pgfpathlineto{\pgfqpoint{5.938241in}{2.303315in}}%
\pgfpathlineto{\pgfqpoint{5.939039in}{2.304360in}}%
\pgfpathlineto{\pgfqpoint{5.938362in}{2.302117in}}%
\pgfpathlineto{\pgfqpoint{5.939302in}{2.302296in}}%
\pgfpathlineto{\pgfqpoint{5.940475in}{2.300932in}}%
\pgfpathlineto{\pgfqpoint{5.939716in}{2.305234in}}%
\pgfpathlineto{\pgfqpoint{5.940485in}{2.301074in}}%
\pgfpathlineto{\pgfqpoint{5.940874in}{2.302958in}}%
\pgfpathlineto{\pgfqpoint{5.940524in}{2.300224in}}%
\pgfpathlineto{\pgfqpoint{5.941580in}{2.300870in}}%
\pgfpathlineto{\pgfqpoint{5.942111in}{2.300554in}}%
\pgfpathlineto{\pgfqpoint{5.942471in}{2.302837in}}%
\pgfpathlineto{\pgfqpoint{5.942656in}{2.302627in}}%
\pgfpathlineto{\pgfqpoint{5.942880in}{2.304096in}}%
\pgfpathlineto{\pgfqpoint{5.943056in}{2.301396in}}%
\pgfpathlineto{\pgfqpoint{5.943761in}{2.302107in}}%
\pgfpathlineto{\pgfqpoint{5.944803in}{2.300197in}}%
\pgfpathlineto{\pgfqpoint{5.944287in}{2.304610in}}%
\pgfpathlineto{\pgfqpoint{5.944871in}{2.301712in}}%
\pgfpathlineto{\pgfqpoint{5.946001in}{2.303971in}}%
\pgfpathlineto{\pgfqpoint{5.945266in}{2.300834in}}%
\pgfpathlineto{\pgfqpoint{5.946016in}{2.303734in}}%
\pgfpathlineto{\pgfqpoint{5.946347in}{2.302323in}}%
\pgfpathlineto{\pgfqpoint{5.947023in}{2.304328in}}%
\pgfpathlineto{\pgfqpoint{5.947184in}{2.304938in}}%
\pgfpathlineto{\pgfqpoint{5.948094in}{2.302802in}}%
\pgfpathlineto{\pgfqpoint{5.948596in}{2.300662in}}%
\pgfpathlineto{\pgfqpoint{5.948309in}{2.303287in}}%
\pgfpathlineto{\pgfqpoint{5.949375in}{2.301211in}}%
\pgfpathlineto{\pgfqpoint{5.950222in}{2.302900in}}%
\pgfpathlineto{\pgfqpoint{5.949462in}{2.300274in}}%
\pgfpathlineto{\pgfqpoint{5.950495in}{2.301587in}}%
\pgfpathlineto{\pgfqpoint{5.950835in}{2.300180in}}%
\pgfpathlineto{\pgfqpoint{5.951478in}{2.304946in}}%
\pgfpathlineto{\pgfqpoint{5.951527in}{2.304545in}}%
\pgfpathlineto{\pgfqpoint{5.951551in}{2.304707in}}%
\pgfpathlineto{\pgfqpoint{5.952388in}{2.301038in}}%
\pgfpathlineto{\pgfqpoint{5.952486in}{2.302003in}}%
\pgfpathlineto{\pgfqpoint{5.953240in}{2.298307in}}%
\pgfpathlineto{\pgfqpoint{5.953606in}{2.300746in}}%
\pgfpathlineto{\pgfqpoint{5.953674in}{2.301468in}}%
\pgfpathlineto{\pgfqpoint{5.953937in}{2.299349in}}%
\pgfpathlineto{\pgfqpoint{5.954248in}{2.297788in}}%
\pgfpathlineto{\pgfqpoint{5.954974in}{2.302569in}}%
\pgfpathlineto{\pgfqpoint{5.954983in}{2.302559in}}%
\pgfpathlineto{\pgfqpoint{5.954993in}{2.302326in}}%
\pgfpathlineto{\pgfqpoint{5.955465in}{2.305620in}}%
\pgfpathlineto{\pgfqpoint{5.955991in}{2.304663in}}%
\pgfpathlineto{\pgfqpoint{5.956853in}{2.305166in}}%
\pgfpathlineto{\pgfqpoint{5.956191in}{2.301409in}}%
\pgfpathlineto{\pgfqpoint{5.957101in}{2.304965in}}%
\pgfpathlineto{\pgfqpoint{5.957637in}{2.307339in}}%
\pgfpathlineto{\pgfqpoint{5.958226in}{2.303661in}}%
\pgfpathlineto{\pgfqpoint{5.959253in}{2.304971in}}%
\pgfpathlineto{\pgfqpoint{5.958761in}{2.302651in}}%
\pgfpathlineto{\pgfqpoint{5.959336in}{2.303863in}}%
\pgfpathlineto{\pgfqpoint{5.961035in}{2.295888in}}%
\pgfpathlineto{\pgfqpoint{5.959399in}{2.305115in}}%
\pgfpathlineto{\pgfqpoint{5.961050in}{2.296753in}}%
\pgfpathlineto{\pgfqpoint{5.961731in}{2.299874in}}%
\pgfpathlineto{\pgfqpoint{5.961273in}{2.295568in}}%
\pgfpathlineto{\pgfqpoint{5.962160in}{2.296957in}}%
\pgfpathlineto{\pgfqpoint{5.962247in}{2.296323in}}%
\pgfpathlineto{\pgfqpoint{5.963041in}{2.301870in}}%
\pgfpathlineto{\pgfqpoint{5.963177in}{2.300829in}}%
\pgfpathlineto{\pgfqpoint{5.964136in}{2.297575in}}%
\pgfpathlineto{\pgfqpoint{5.963240in}{2.301443in}}%
\pgfpathlineto{\pgfqpoint{5.964691in}{2.299948in}}%
\pgfpathlineto{\pgfqpoint{5.965222in}{2.304858in}}%
\pgfpathlineto{\pgfqpoint{5.965816in}{2.300752in}}%
\pgfpathlineto{\pgfqpoint{5.965835in}{2.301325in}}%
\pgfpathlineto{\pgfqpoint{5.966760in}{2.297581in}}%
\pgfpathlineto{\pgfqpoint{5.966828in}{2.298922in}}%
\pgfpathlineto{\pgfqpoint{5.966833in}{2.298817in}}%
\pgfpathlineto{\pgfqpoint{5.967724in}{2.303241in}}%
\pgfpathlineto{\pgfqpoint{5.967788in}{2.302014in}}%
\pgfpathlineto{\pgfqpoint{5.967846in}{2.303785in}}%
\pgfpathlineto{\pgfqpoint{5.968503in}{2.299837in}}%
\pgfpathlineto{\pgfqpoint{5.968888in}{2.301561in}}%
\pgfpathlineto{\pgfqpoint{5.969058in}{2.300018in}}%
\pgfpathlineto{\pgfqpoint{5.969677in}{2.304482in}}%
\pgfpathlineto{\pgfqpoint{5.969701in}{2.304935in}}%
\pgfpathlineto{\pgfqpoint{5.970217in}{2.301414in}}%
\pgfpathlineto{\pgfqpoint{5.970694in}{2.303681in}}%
\pgfpathlineto{\pgfqpoint{5.971804in}{2.302259in}}%
\pgfpathlineto{\pgfqpoint{5.971249in}{2.304421in}}%
\pgfpathlineto{\pgfqpoint{5.971828in}{2.302737in}}%
\pgfpathlineto{\pgfqpoint{5.972919in}{2.304392in}}%
\pgfpathlineto{\pgfqpoint{5.971843in}{2.302637in}}%
\pgfpathlineto{\pgfqpoint{5.973036in}{2.303667in}}%
\pgfpathlineto{\pgfqpoint{5.974121in}{2.301325in}}%
\pgfpathlineto{\pgfqpoint{5.973255in}{2.305374in}}%
\pgfpathlineto{\pgfqpoint{5.974165in}{2.302267in}}%
\pgfpathlineto{\pgfqpoint{5.975270in}{2.301601in}}%
\pgfpathlineto{\pgfqpoint{5.974915in}{2.303985in}}%
\pgfpathlineto{\pgfqpoint{5.975290in}{2.302012in}}%
\pgfpathlineto{\pgfqpoint{5.975504in}{2.303668in}}%
\pgfpathlineto{\pgfqpoint{5.976264in}{2.301090in}}%
\pgfpathlineto{\pgfqpoint{5.976303in}{2.301101in}}%
\pgfpathlineto{\pgfqpoint{5.976337in}{2.300886in}}%
\pgfpathlineto{\pgfqpoint{5.977335in}{2.303154in}}%
\pgfpathlineto{\pgfqpoint{5.977340in}{2.302968in}}%
\pgfpathlineto{\pgfqpoint{5.977486in}{2.304603in}}%
\pgfpathlineto{\pgfqpoint{5.978153in}{2.301339in}}%
\pgfpathlineto{\pgfqpoint{5.978274in}{2.301581in}}%
\pgfpathlineto{\pgfqpoint{5.978303in}{2.301143in}}%
\pgfpathlineto{\pgfqpoint{5.979229in}{2.304440in}}%
\pgfpathlineto{\pgfqpoint{5.979340in}{2.302803in}}%
\pgfpathlineto{\pgfqpoint{5.979370in}{2.304056in}}%
\pgfpathlineto{\pgfqpoint{5.979550in}{2.304911in}}%
\pgfpathlineto{\pgfqpoint{5.980178in}{2.300950in}}%
\pgfpathlineto{\pgfqpoint{5.980451in}{2.302899in}}%
\pgfpathlineto{\pgfqpoint{5.981556in}{2.301220in}}%
\pgfpathlineto{\pgfqpoint{5.980704in}{2.304594in}}%
\pgfpathlineto{\pgfqpoint{5.981585in}{2.301772in}}%
\pgfpathlineto{\pgfqpoint{5.982276in}{2.305689in}}%
\pgfpathlineto{\pgfqpoint{5.982719in}{2.303634in}}%
\pgfpathlineto{\pgfqpoint{5.983187in}{2.301651in}}%
\pgfpathlineto{\pgfqpoint{5.982826in}{2.304026in}}%
\pgfpathlineto{\pgfqpoint{5.983834in}{2.302789in}}%
\pgfpathlineto{\pgfqpoint{5.983858in}{2.303485in}}%
\pgfpathlineto{\pgfqpoint{5.984589in}{2.299518in}}%
\pgfpathlineto{\pgfqpoint{5.984930in}{2.301728in}}%
\pgfpathlineto{\pgfqpoint{5.986078in}{2.305944in}}%
\pgfpathlineto{\pgfqpoint{5.985504in}{2.300888in}}%
\pgfpathlineto{\pgfqpoint{5.986127in}{2.304730in}}%
\pgfpathlineto{\pgfqpoint{5.986930in}{2.302202in}}%
\pgfpathlineto{\pgfqpoint{5.986205in}{2.306115in}}%
\pgfpathlineto{\pgfqpoint{5.987320in}{2.303328in}}%
\pgfpathlineto{\pgfqpoint{5.987524in}{2.305251in}}%
\pgfpathlineto{\pgfqpoint{5.988109in}{2.299435in}}%
\pgfpathlineto{\pgfqpoint{5.988401in}{2.301137in}}%
\pgfpathlineto{\pgfqpoint{5.989336in}{2.298496in}}%
\pgfpathlineto{\pgfqpoint{5.989521in}{2.300251in}}%
\pgfpathlineto{\pgfqpoint{5.990290in}{2.304442in}}%
\pgfpathlineto{\pgfqpoint{5.990655in}{2.302698in}}%
\pgfpathlineto{\pgfqpoint{5.991385in}{2.299169in}}%
\pgfpathlineto{\pgfqpoint{5.990767in}{2.303358in}}%
\pgfpathlineto{\pgfqpoint{5.991814in}{2.300602in}}%
\pgfpathlineto{\pgfqpoint{5.992549in}{2.304289in}}%
\pgfpathlineto{\pgfqpoint{5.992159in}{2.298315in}}%
\pgfpathlineto{\pgfqpoint{5.992982in}{2.302914in}}%
\pgfpathlineto{\pgfqpoint{5.994170in}{2.305439in}}%
\pgfpathlineto{\pgfqpoint{5.993303in}{2.301129in}}%
\pgfpathlineto{\pgfqpoint{5.994209in}{2.303994in}}%
\pgfpathlineto{\pgfqpoint{5.994710in}{2.301476in}}%
\pgfpathlineto{\pgfqpoint{5.995260in}{2.305435in}}%
\pgfpathlineto{\pgfqpoint{5.995314in}{2.304263in}}%
\pgfpathlineto{\pgfqpoint{5.996229in}{2.301543in}}%
\pgfpathlineto{\pgfqpoint{5.995738in}{2.305203in}}%
\pgfpathlineto{\pgfqpoint{5.996444in}{2.303583in}}%
\pgfpathlineto{\pgfqpoint{5.996502in}{2.303877in}}%
\pgfpathlineto{\pgfqpoint{5.997266in}{2.298886in}}%
\pgfpathlineto{\pgfqpoint{5.997485in}{2.300024in}}%
\pgfpathlineto{\pgfqpoint{5.997500in}{2.299567in}}%
\pgfpathlineto{\pgfqpoint{5.997948in}{2.303685in}}%
\pgfpathlineto{\pgfqpoint{5.998449in}{2.302850in}}%
\pgfpathlineto{\pgfqpoint{5.998776in}{2.304862in}}%
\pgfpathlineto{\pgfqpoint{5.998488in}{2.302141in}}%
\pgfpathlineto{\pgfqpoint{5.999559in}{2.302923in}}%
\pgfpathlineto{\pgfqpoint{5.999603in}{2.302249in}}%
\pgfpathlineto{\pgfqpoint{5.999954in}{2.304401in}}%
\pgfpathlineto{\pgfqpoint{6.000606in}{2.303666in}}%
\pgfpathlineto{\pgfqpoint{6.000611in}{2.303734in}}%
\pgfpathlineto{\pgfqpoint{6.000869in}{2.299555in}}%
\pgfpathlineto{\pgfqpoint{6.001512in}{2.300860in}}%
\pgfpathlineto{\pgfqpoint{6.002651in}{2.299671in}}%
\pgfpathlineto{\pgfqpoint{6.001960in}{2.303738in}}%
\pgfpathlineto{\pgfqpoint{6.002661in}{2.300070in}}%
\pgfpathlineto{\pgfqpoint{6.003698in}{2.301945in}}%
\pgfpathlineto{\pgfqpoint{6.003167in}{2.298928in}}%
\pgfpathlineto{\pgfqpoint{6.003795in}{2.301729in}}%
\pgfpathlineto{\pgfqpoint{6.004223in}{2.302697in}}%
\pgfpathlineto{\pgfqpoint{6.005075in}{2.298609in}}%
\pgfpathlineto{\pgfqpoint{6.005236in}{2.300822in}}%
\pgfpathlineto{\pgfqpoint{6.005655in}{2.297992in}}%
\pgfpathlineto{\pgfqpoint{6.006200in}{2.298945in}}%
\pgfpathlineto{\pgfqpoint{6.006268in}{2.298576in}}%
\pgfpathlineto{\pgfqpoint{6.006575in}{2.300752in}}%
\pgfpathlineto{\pgfqpoint{6.006575in}{2.300752in}}%
\pgfusepath{stroke}%
\end{pgfscope}%
\begin{pgfscope}%
\pgfsetrectcap%
\pgfsetmiterjoin%
\pgfsetlinewidth{0.803000pt}%
\definecolor{currentstroke}{rgb}{0.000000,0.000000,0.000000}%
\pgfsetstrokecolor{currentstroke}%
\pgfsetdash{}{0pt}%
\pgfpathmoveto{\pgfqpoint{0.894648in}{0.771190in}}%
\pgfpathlineto{\pgfqpoint{0.894648in}{4.601775in}}%
\pgfusepath{stroke}%
\end{pgfscope}%
\begin{pgfscope}%
\pgfsetrectcap%
\pgfsetmiterjoin%
\pgfsetlinewidth{0.803000pt}%
\definecolor{currentstroke}{rgb}{0.000000,0.000000,0.000000}%
\pgfsetstrokecolor{currentstroke}%
\pgfsetdash{}{0pt}%
\pgfpathmoveto{\pgfqpoint{6.250000in}{0.771190in}}%
\pgfpathlineto{\pgfqpoint{6.250000in}{4.601775in}}%
\pgfusepath{stroke}%
\end{pgfscope}%
\begin{pgfscope}%
\pgfsetrectcap%
\pgfsetmiterjoin%
\pgfsetlinewidth{0.803000pt}%
\definecolor{currentstroke}{rgb}{0.000000,0.000000,0.000000}%
\pgfsetstrokecolor{currentstroke}%
\pgfsetdash{}{0pt}%
\pgfpathmoveto{\pgfqpoint{0.894648in}{0.771190in}}%
\pgfpathlineto{\pgfqpoint{6.250000in}{0.771190in}}%
\pgfusepath{stroke}%
\end{pgfscope}%
\begin{pgfscope}%
\pgfsetrectcap%
\pgfsetmiterjoin%
\pgfsetlinewidth{0.803000pt}%
\definecolor{currentstroke}{rgb}{0.000000,0.000000,0.000000}%
\pgfsetstrokecolor{currentstroke}%
\pgfsetdash{}{0pt}%
\pgfpathmoveto{\pgfqpoint{0.894648in}{4.601775in}}%
\pgfpathlineto{\pgfqpoint{6.250000in}{4.601775in}}%
\pgfusepath{stroke}%
\end{pgfscope}%
\end{pgfpicture}%
\makeatother%
\endgroup%
}
    \caption{Plotted shape of potential and force $\Delta U = 10 k_B T$}
    \label{fig:No_flashing}
\end{figure}


\begin{figure}[H]
    \centering
    \resizebox{\columnwidth}{!}{%% Creator: Matplotlib, PGF backend
%%
%% To include the figure in your LaTeX document, write
%%   \input{<filename>.pgf}
%%
%% Make sure the required packages are loaded in your preamble
%%   \usepackage{pgf}
%%
%% Also ensure that all the required font packages are loaded; for instance,
%% the lmodern package is sometimes necessary when using math font.
%%   \usepackage{lmodern}
%%
%% Figures using additional raster images can only be included by \input if
%% they are in the same directory as the main LaTeX file. For loading figures
%% from other directories you can use the `import` package
%%   \usepackage{import}
%%
%% and then include the figures with
%%   \import{<path to file>}{<filename>.pgf}
%%
%% Matplotlib used the following preamble
%%   \def\mathdefault#1{#1}
%%   \everymath=\expandafter{\the\everymath\displaystyle}
%%   \IfFileExists{scrextend.sty}{
%%     \usepackage[fontsize=10.000000pt]{scrextend}
%%   }{
%%     \renewcommand{\normalsize}{\fontsize{10.000000}{12.000000}\selectfont}
%%     \normalsize
%%   }
%%   
%%   \makeatletter\@ifpackageloaded{underscore}{}{\usepackage[strings]{underscore}}\makeatother
%%
\begingroup%
\makeatletter%
\begin{pgfpicture}%
\pgfpathrectangle{\pgfpointorigin}{\pgfqpoint{6.400000in}{4.800000in}}%
\pgfusepath{use as bounding box, clip}%
\begin{pgfscope}%
\pgfsetbuttcap%
\pgfsetmiterjoin%
\definecolor{currentfill}{rgb}{1.000000,1.000000,1.000000}%
\pgfsetfillcolor{currentfill}%
\pgfsetlinewidth{0.000000pt}%
\definecolor{currentstroke}{rgb}{1.000000,1.000000,1.000000}%
\pgfsetstrokecolor{currentstroke}%
\pgfsetdash{}{0pt}%
\pgfpathmoveto{\pgfqpoint{0.000000in}{0.000000in}}%
\pgfpathlineto{\pgfqpoint{6.400000in}{0.000000in}}%
\pgfpathlineto{\pgfqpoint{6.400000in}{4.800000in}}%
\pgfpathlineto{\pgfqpoint{0.000000in}{4.800000in}}%
\pgfpathlineto{\pgfqpoint{0.000000in}{0.000000in}}%
\pgfpathclose%
\pgfusepath{fill}%
\end{pgfscope}%
\begin{pgfscope}%
\pgfsetbuttcap%
\pgfsetmiterjoin%
\definecolor{currentfill}{rgb}{1.000000,1.000000,1.000000}%
\pgfsetfillcolor{currentfill}%
\pgfsetlinewidth{0.000000pt}%
\definecolor{currentstroke}{rgb}{0.000000,0.000000,0.000000}%
\pgfsetstrokecolor{currentstroke}%
\pgfsetstrokeopacity{0.000000}%
\pgfsetdash{}{0pt}%
\pgfpathmoveto{\pgfqpoint{0.894648in}{0.771190in}}%
\pgfpathlineto{\pgfqpoint{6.250000in}{0.771190in}}%
\pgfpathlineto{\pgfqpoint{6.250000in}{4.601775in}}%
\pgfpathlineto{\pgfqpoint{0.894648in}{4.601775in}}%
\pgfpathlineto{\pgfqpoint{0.894648in}{0.771190in}}%
\pgfpathclose%
\pgfusepath{fill}%
\end{pgfscope}%
\begin{pgfscope}%
\pgfsetbuttcap%
\pgfsetroundjoin%
\definecolor{currentfill}{rgb}{0.000000,0.000000,0.000000}%
\pgfsetfillcolor{currentfill}%
\pgfsetlinewidth{0.803000pt}%
\definecolor{currentstroke}{rgb}{0.000000,0.000000,0.000000}%
\pgfsetstrokecolor{currentstroke}%
\pgfsetdash{}{0pt}%
\pgfsys@defobject{currentmarker}{\pgfqpoint{0.000000in}{-0.048611in}}{\pgfqpoint{0.000000in}{0.000000in}}{%
\pgfpathmoveto{\pgfqpoint{0.000000in}{0.000000in}}%
\pgfpathlineto{\pgfqpoint{0.000000in}{-0.048611in}}%
\pgfusepath{stroke,fill}%
}%
\begin{pgfscope}%
\pgfsys@transformshift{1.138073in}{0.771190in}%
\pgfsys@useobject{currentmarker}{}%
\end{pgfscope}%
\end{pgfscope}%
\begin{pgfscope}%
\definecolor{textcolor}{rgb}{0.000000,0.000000,0.000000}%
\pgfsetstrokecolor{textcolor}%
\pgfsetfillcolor{textcolor}%
\pgftext[x=1.138073in,y=0.673968in,,top]{\color{textcolor}{\rmfamily\fontsize{10.000000}{12.000000}\selectfont\catcode`\^=\active\def^{\ifmmode\sp\else\^{}\fi}\catcode`\%=\active\def%{\%}$\mathdefault{0}$}}%
\end{pgfscope}%
\begin{pgfscope}%
\pgfsetbuttcap%
\pgfsetroundjoin%
\definecolor{currentfill}{rgb}{0.000000,0.000000,0.000000}%
\pgfsetfillcolor{currentfill}%
\pgfsetlinewidth{0.803000pt}%
\definecolor{currentstroke}{rgb}{0.000000,0.000000,0.000000}%
\pgfsetstrokecolor{currentstroke}%
\pgfsetdash{}{0pt}%
\pgfsys@defobject{currentmarker}{\pgfqpoint{0.000000in}{-0.048611in}}{\pgfqpoint{0.000000in}{0.000000in}}{%
\pgfpathmoveto{\pgfqpoint{0.000000in}{0.000000in}}%
\pgfpathlineto{\pgfqpoint{0.000000in}{-0.048611in}}%
\pgfusepath{stroke,fill}%
}%
\begin{pgfscope}%
\pgfsys@transformshift{1.887073in}{0.771190in}%
\pgfsys@useobject{currentmarker}{}%
\end{pgfscope}%
\end{pgfscope}%
\begin{pgfscope}%
\definecolor{textcolor}{rgb}{0.000000,0.000000,0.000000}%
\pgfsetstrokecolor{textcolor}%
\pgfsetfillcolor{textcolor}%
\pgftext[x=1.887073in,y=0.673968in,,top]{\color{textcolor}{\rmfamily\fontsize{10.000000}{12.000000}\selectfont\catcode`\^=\active\def^{\ifmmode\sp\else\^{}\fi}\catcode`\%=\active\def%{\%}$\mathdefault{100}$}}%
\end{pgfscope}%
\begin{pgfscope}%
\pgfsetbuttcap%
\pgfsetroundjoin%
\definecolor{currentfill}{rgb}{0.000000,0.000000,0.000000}%
\pgfsetfillcolor{currentfill}%
\pgfsetlinewidth{0.803000pt}%
\definecolor{currentstroke}{rgb}{0.000000,0.000000,0.000000}%
\pgfsetstrokecolor{currentstroke}%
\pgfsetdash{}{0pt}%
\pgfsys@defobject{currentmarker}{\pgfqpoint{0.000000in}{-0.048611in}}{\pgfqpoint{0.000000in}{0.000000in}}{%
\pgfpathmoveto{\pgfqpoint{0.000000in}{0.000000in}}%
\pgfpathlineto{\pgfqpoint{0.000000in}{-0.048611in}}%
\pgfusepath{stroke,fill}%
}%
\begin{pgfscope}%
\pgfsys@transformshift{2.636073in}{0.771190in}%
\pgfsys@useobject{currentmarker}{}%
\end{pgfscope}%
\end{pgfscope}%
\begin{pgfscope}%
\definecolor{textcolor}{rgb}{0.000000,0.000000,0.000000}%
\pgfsetstrokecolor{textcolor}%
\pgfsetfillcolor{textcolor}%
\pgftext[x=2.636073in,y=0.673968in,,top]{\color{textcolor}{\rmfamily\fontsize{10.000000}{12.000000}\selectfont\catcode`\^=\active\def^{\ifmmode\sp\else\^{}\fi}\catcode`\%=\active\def%{\%}$\mathdefault{200}$}}%
\end{pgfscope}%
\begin{pgfscope}%
\pgfsetbuttcap%
\pgfsetroundjoin%
\definecolor{currentfill}{rgb}{0.000000,0.000000,0.000000}%
\pgfsetfillcolor{currentfill}%
\pgfsetlinewidth{0.803000pt}%
\definecolor{currentstroke}{rgb}{0.000000,0.000000,0.000000}%
\pgfsetstrokecolor{currentstroke}%
\pgfsetdash{}{0pt}%
\pgfsys@defobject{currentmarker}{\pgfqpoint{0.000000in}{-0.048611in}}{\pgfqpoint{0.000000in}{0.000000in}}{%
\pgfpathmoveto{\pgfqpoint{0.000000in}{0.000000in}}%
\pgfpathlineto{\pgfqpoint{0.000000in}{-0.048611in}}%
\pgfusepath{stroke,fill}%
}%
\begin{pgfscope}%
\pgfsys@transformshift{3.385074in}{0.771190in}%
\pgfsys@useobject{currentmarker}{}%
\end{pgfscope}%
\end{pgfscope}%
\begin{pgfscope}%
\definecolor{textcolor}{rgb}{0.000000,0.000000,0.000000}%
\pgfsetstrokecolor{textcolor}%
\pgfsetfillcolor{textcolor}%
\pgftext[x=3.385074in,y=0.673968in,,top]{\color{textcolor}{\rmfamily\fontsize{10.000000}{12.000000}\selectfont\catcode`\^=\active\def^{\ifmmode\sp\else\^{}\fi}\catcode`\%=\active\def%{\%}$\mathdefault{300}$}}%
\end{pgfscope}%
\begin{pgfscope}%
\pgfsetbuttcap%
\pgfsetroundjoin%
\definecolor{currentfill}{rgb}{0.000000,0.000000,0.000000}%
\pgfsetfillcolor{currentfill}%
\pgfsetlinewidth{0.803000pt}%
\definecolor{currentstroke}{rgb}{0.000000,0.000000,0.000000}%
\pgfsetstrokecolor{currentstroke}%
\pgfsetdash{}{0pt}%
\pgfsys@defobject{currentmarker}{\pgfqpoint{0.000000in}{-0.048611in}}{\pgfqpoint{0.000000in}{0.000000in}}{%
\pgfpathmoveto{\pgfqpoint{0.000000in}{0.000000in}}%
\pgfpathlineto{\pgfqpoint{0.000000in}{-0.048611in}}%
\pgfusepath{stroke,fill}%
}%
\begin{pgfscope}%
\pgfsys@transformshift{4.134074in}{0.771190in}%
\pgfsys@useobject{currentmarker}{}%
\end{pgfscope}%
\end{pgfscope}%
\begin{pgfscope}%
\definecolor{textcolor}{rgb}{0.000000,0.000000,0.000000}%
\pgfsetstrokecolor{textcolor}%
\pgfsetfillcolor{textcolor}%
\pgftext[x=4.134074in,y=0.673968in,,top]{\color{textcolor}{\rmfamily\fontsize{10.000000}{12.000000}\selectfont\catcode`\^=\active\def^{\ifmmode\sp\else\^{}\fi}\catcode`\%=\active\def%{\%}$\mathdefault{400}$}}%
\end{pgfscope}%
\begin{pgfscope}%
\pgfsetbuttcap%
\pgfsetroundjoin%
\definecolor{currentfill}{rgb}{0.000000,0.000000,0.000000}%
\pgfsetfillcolor{currentfill}%
\pgfsetlinewidth{0.803000pt}%
\definecolor{currentstroke}{rgb}{0.000000,0.000000,0.000000}%
\pgfsetstrokecolor{currentstroke}%
\pgfsetdash{}{0pt}%
\pgfsys@defobject{currentmarker}{\pgfqpoint{0.000000in}{-0.048611in}}{\pgfqpoint{0.000000in}{0.000000in}}{%
\pgfpathmoveto{\pgfqpoint{0.000000in}{0.000000in}}%
\pgfpathlineto{\pgfqpoint{0.000000in}{-0.048611in}}%
\pgfusepath{stroke,fill}%
}%
\begin{pgfscope}%
\pgfsys@transformshift{4.883074in}{0.771190in}%
\pgfsys@useobject{currentmarker}{}%
\end{pgfscope}%
\end{pgfscope}%
\begin{pgfscope}%
\definecolor{textcolor}{rgb}{0.000000,0.000000,0.000000}%
\pgfsetstrokecolor{textcolor}%
\pgfsetfillcolor{textcolor}%
\pgftext[x=4.883074in,y=0.673968in,,top]{\color{textcolor}{\rmfamily\fontsize{10.000000}{12.000000}\selectfont\catcode`\^=\active\def^{\ifmmode\sp\else\^{}\fi}\catcode`\%=\active\def%{\%}$\mathdefault{500}$}}%
\end{pgfscope}%
\begin{pgfscope}%
\pgfsetbuttcap%
\pgfsetroundjoin%
\definecolor{currentfill}{rgb}{0.000000,0.000000,0.000000}%
\pgfsetfillcolor{currentfill}%
\pgfsetlinewidth{0.803000pt}%
\definecolor{currentstroke}{rgb}{0.000000,0.000000,0.000000}%
\pgfsetstrokecolor{currentstroke}%
\pgfsetdash{}{0pt}%
\pgfsys@defobject{currentmarker}{\pgfqpoint{0.000000in}{-0.048611in}}{\pgfqpoint{0.000000in}{0.000000in}}{%
\pgfpathmoveto{\pgfqpoint{0.000000in}{0.000000in}}%
\pgfpathlineto{\pgfqpoint{0.000000in}{-0.048611in}}%
\pgfusepath{stroke,fill}%
}%
\begin{pgfscope}%
\pgfsys@transformshift{5.632075in}{0.771190in}%
\pgfsys@useobject{currentmarker}{}%
\end{pgfscope}%
\end{pgfscope}%
\begin{pgfscope}%
\definecolor{textcolor}{rgb}{0.000000,0.000000,0.000000}%
\pgfsetstrokecolor{textcolor}%
\pgfsetfillcolor{textcolor}%
\pgftext[x=5.632075in,y=0.673968in,,top]{\color{textcolor}{\rmfamily\fontsize{10.000000}{12.000000}\selectfont\catcode`\^=\active\def^{\ifmmode\sp\else\^{}\fi}\catcode`\%=\active\def%{\%}$\mathdefault{600}$}}%
\end{pgfscope}%
\begin{pgfscope}%
\definecolor{textcolor}{rgb}{0.000000,0.000000,0.000000}%
\pgfsetstrokecolor{textcolor}%
\pgfsetfillcolor{textcolor}%
\pgftext[x=3.572324in,y=0.494955in,,top]{\color{textcolor}{\rmfamily\fontsize{24.000000}{28.800000}\selectfont\catcode`\^=\active\def^{\ifmmode\sp\else\^{}\fi}\catcode`\%=\active\def%{\%}$t (s)$}}%
\end{pgfscope}%
\begin{pgfscope}%
\pgfsetbuttcap%
\pgfsetroundjoin%
\definecolor{currentfill}{rgb}{0.000000,0.000000,0.000000}%
\pgfsetfillcolor{currentfill}%
\pgfsetlinewidth{0.803000pt}%
\definecolor{currentstroke}{rgb}{0.000000,0.000000,0.000000}%
\pgfsetstrokecolor{currentstroke}%
\pgfsetdash{}{0pt}%
\pgfsys@defobject{currentmarker}{\pgfqpoint{-0.048611in}{0.000000in}}{\pgfqpoint{-0.000000in}{0.000000in}}{%
\pgfpathmoveto{\pgfqpoint{-0.000000in}{0.000000in}}%
\pgfpathlineto{\pgfqpoint{-0.048611in}{0.000000in}}%
\pgfusepath{stroke,fill}%
}%
\begin{pgfscope}%
\pgfsys@transformshift{0.894648in}{0.771190in}%
\pgfsys@useobject{currentmarker}{}%
\end{pgfscope}%
\end{pgfscope}%
\begin{pgfscope}%
\definecolor{textcolor}{rgb}{0.000000,0.000000,0.000000}%
\pgfsetstrokecolor{textcolor}%
\pgfsetfillcolor{textcolor}%
\pgftext[x=0.550511in, y=0.722965in, left, base]{\color{textcolor}{\rmfamily\fontsize{10.000000}{12.000000}\selectfont\catcode`\^=\active\def^{\ifmmode\sp\else\^{}\fi}\catcode`\%=\active\def%{\%}$\mathdefault{\ensuremath{-}40}$}}%
\end{pgfscope}%
\begin{pgfscope}%
\pgfsetbuttcap%
\pgfsetroundjoin%
\definecolor{currentfill}{rgb}{0.000000,0.000000,0.000000}%
\pgfsetfillcolor{currentfill}%
\pgfsetlinewidth{0.803000pt}%
\definecolor{currentstroke}{rgb}{0.000000,0.000000,0.000000}%
\pgfsetstrokecolor{currentstroke}%
\pgfsetdash{}{0pt}%
\pgfsys@defobject{currentmarker}{\pgfqpoint{-0.048611in}{0.000000in}}{\pgfqpoint{-0.000000in}{0.000000in}}{%
\pgfpathmoveto{\pgfqpoint{-0.000000in}{0.000000in}}%
\pgfpathlineto{\pgfqpoint{-0.048611in}{0.000000in}}%
\pgfusepath{stroke,fill}%
}%
\begin{pgfscope}%
\pgfsys@transformshift{0.894648in}{1.537307in}%
\pgfsys@useobject{currentmarker}{}%
\end{pgfscope}%
\end{pgfscope}%
\begin{pgfscope}%
\definecolor{textcolor}{rgb}{0.000000,0.000000,0.000000}%
\pgfsetstrokecolor{textcolor}%
\pgfsetfillcolor{textcolor}%
\pgftext[x=0.550511in, y=1.489082in, left, base]{\color{textcolor}{\rmfamily\fontsize{10.000000}{12.000000}\selectfont\catcode`\^=\active\def^{\ifmmode\sp\else\^{}\fi}\catcode`\%=\active\def%{\%}$\mathdefault{\ensuremath{-}20}$}}%
\end{pgfscope}%
\begin{pgfscope}%
\pgfsetbuttcap%
\pgfsetroundjoin%
\definecolor{currentfill}{rgb}{0.000000,0.000000,0.000000}%
\pgfsetfillcolor{currentfill}%
\pgfsetlinewidth{0.803000pt}%
\definecolor{currentstroke}{rgb}{0.000000,0.000000,0.000000}%
\pgfsetstrokecolor{currentstroke}%
\pgfsetdash{}{0pt}%
\pgfsys@defobject{currentmarker}{\pgfqpoint{-0.048611in}{0.000000in}}{\pgfqpoint{-0.000000in}{0.000000in}}{%
\pgfpathmoveto{\pgfqpoint{-0.000000in}{0.000000in}}%
\pgfpathlineto{\pgfqpoint{-0.048611in}{0.000000in}}%
\pgfusepath{stroke,fill}%
}%
\begin{pgfscope}%
\pgfsys@transformshift{0.894648in}{2.303424in}%
\pgfsys@useobject{currentmarker}{}%
\end{pgfscope}%
\end{pgfscope}%
\begin{pgfscope}%
\definecolor{textcolor}{rgb}{0.000000,0.000000,0.000000}%
\pgfsetstrokecolor{textcolor}%
\pgfsetfillcolor{textcolor}%
\pgftext[x=0.727981in, y=2.255199in, left, base]{\color{textcolor}{\rmfamily\fontsize{10.000000}{12.000000}\selectfont\catcode`\^=\active\def^{\ifmmode\sp\else\^{}\fi}\catcode`\%=\active\def%{\%}$\mathdefault{0}$}}%
\end{pgfscope}%
\begin{pgfscope}%
\pgfsetbuttcap%
\pgfsetroundjoin%
\definecolor{currentfill}{rgb}{0.000000,0.000000,0.000000}%
\pgfsetfillcolor{currentfill}%
\pgfsetlinewidth{0.803000pt}%
\definecolor{currentstroke}{rgb}{0.000000,0.000000,0.000000}%
\pgfsetstrokecolor{currentstroke}%
\pgfsetdash{}{0pt}%
\pgfsys@defobject{currentmarker}{\pgfqpoint{-0.048611in}{0.000000in}}{\pgfqpoint{-0.000000in}{0.000000in}}{%
\pgfpathmoveto{\pgfqpoint{-0.000000in}{0.000000in}}%
\pgfpathlineto{\pgfqpoint{-0.048611in}{0.000000in}}%
\pgfusepath{stroke,fill}%
}%
\begin{pgfscope}%
\pgfsys@transformshift{0.894648in}{3.069541in}%
\pgfsys@useobject{currentmarker}{}%
\end{pgfscope}%
\end{pgfscope}%
\begin{pgfscope}%
\definecolor{textcolor}{rgb}{0.000000,0.000000,0.000000}%
\pgfsetstrokecolor{textcolor}%
\pgfsetfillcolor{textcolor}%
\pgftext[x=0.658536in, y=3.021316in, left, base]{\color{textcolor}{\rmfamily\fontsize{10.000000}{12.000000}\selectfont\catcode`\^=\active\def^{\ifmmode\sp\else\^{}\fi}\catcode`\%=\active\def%{\%}$\mathdefault{20}$}}%
\end{pgfscope}%
\begin{pgfscope}%
\pgfsetbuttcap%
\pgfsetroundjoin%
\definecolor{currentfill}{rgb}{0.000000,0.000000,0.000000}%
\pgfsetfillcolor{currentfill}%
\pgfsetlinewidth{0.803000pt}%
\definecolor{currentstroke}{rgb}{0.000000,0.000000,0.000000}%
\pgfsetstrokecolor{currentstroke}%
\pgfsetdash{}{0pt}%
\pgfsys@defobject{currentmarker}{\pgfqpoint{-0.048611in}{0.000000in}}{\pgfqpoint{-0.000000in}{0.000000in}}{%
\pgfpathmoveto{\pgfqpoint{-0.000000in}{0.000000in}}%
\pgfpathlineto{\pgfqpoint{-0.048611in}{0.000000in}}%
\pgfusepath{stroke,fill}%
}%
\begin{pgfscope}%
\pgfsys@transformshift{0.894648in}{3.835658in}%
\pgfsys@useobject{currentmarker}{}%
\end{pgfscope}%
\end{pgfscope}%
\begin{pgfscope}%
\definecolor{textcolor}{rgb}{0.000000,0.000000,0.000000}%
\pgfsetstrokecolor{textcolor}%
\pgfsetfillcolor{textcolor}%
\pgftext[x=0.658536in, y=3.787432in, left, base]{\color{textcolor}{\rmfamily\fontsize{10.000000}{12.000000}\selectfont\catcode`\^=\active\def^{\ifmmode\sp\else\^{}\fi}\catcode`\%=\active\def%{\%}$\mathdefault{40}$}}%
\end{pgfscope}%
\begin{pgfscope}%
\pgfsetbuttcap%
\pgfsetroundjoin%
\definecolor{currentfill}{rgb}{0.000000,0.000000,0.000000}%
\pgfsetfillcolor{currentfill}%
\pgfsetlinewidth{0.803000pt}%
\definecolor{currentstroke}{rgb}{0.000000,0.000000,0.000000}%
\pgfsetstrokecolor{currentstroke}%
\pgfsetdash{}{0pt}%
\pgfsys@defobject{currentmarker}{\pgfqpoint{-0.048611in}{0.000000in}}{\pgfqpoint{-0.000000in}{0.000000in}}{%
\pgfpathmoveto{\pgfqpoint{-0.000000in}{0.000000in}}%
\pgfpathlineto{\pgfqpoint{-0.048611in}{0.000000in}}%
\pgfusepath{stroke,fill}%
}%
\begin{pgfscope}%
\pgfsys@transformshift{0.894648in}{4.601775in}%
\pgfsys@useobject{currentmarker}{}%
\end{pgfscope}%
\end{pgfscope}%
\begin{pgfscope}%
\definecolor{textcolor}{rgb}{0.000000,0.000000,0.000000}%
\pgfsetstrokecolor{textcolor}%
\pgfsetfillcolor{textcolor}%
\pgftext[x=0.658536in, y=4.553549in, left, base]{\color{textcolor}{\rmfamily\fontsize{10.000000}{12.000000}\selectfont\catcode`\^=\active\def^{\ifmmode\sp\else\^{}\fi}\catcode`\%=\active\def%{\%}$\mathdefault{60}$}}%
\end{pgfscope}%
\begin{pgfscope}%
\definecolor{textcolor}{rgb}{0.000000,0.000000,0.000000}%
\pgfsetstrokecolor{textcolor}%
\pgfsetfillcolor{textcolor}%
\pgftext[x=0.494955in,y=2.686482in,,bottom,rotate=90.000000]{\color{textcolor}{\rmfamily\fontsize{24.000000}{28.800000}\selectfont\catcode`\^=\active\def^{\ifmmode\sp\else\^{}\fi}\catcode`\%=\active\def%{\%}$x (\mu m)$}}%
\end{pgfscope}%
\begin{pgfscope}%
\pgfpathrectangle{\pgfqpoint{0.894648in}{0.771190in}}{\pgfqpoint{5.355353in}{3.830585in}}%
\pgfusepath{clip}%
\pgfsetrectcap%
\pgfsetroundjoin%
\pgfsetlinewidth{1.505625pt}%
\definecolor{currentstroke}{rgb}{0.121569,0.466667,0.705882}%
\pgfsetstrokecolor{currentstroke}%
\pgfsetdash{}{0pt}%
\pgfpathmoveto{\pgfqpoint{1.138073in}{2.303424in}}%
\pgfpathlineto{\pgfqpoint{1.138262in}{2.303574in}}%
\pgfpathlineto{\pgfqpoint{1.139024in}{2.303085in}}%
\pgfpathlineto{\pgfqpoint{1.139183in}{2.303290in}}%
\pgfpathlineto{\pgfqpoint{1.140299in}{2.302977in}}%
\pgfpathlineto{\pgfqpoint{1.139186in}{2.303311in}}%
\pgfpathlineto{\pgfqpoint{1.140421in}{2.303067in}}%
\pgfpathlineto{\pgfqpoint{1.141571in}{2.303248in}}%
\pgfpathlineto{\pgfqpoint{1.140571in}{2.302767in}}%
\pgfpathlineto{\pgfqpoint{1.141572in}{2.303236in}}%
\pgfpathlineto{\pgfqpoint{1.141648in}{2.303160in}}%
\pgfpathlineto{\pgfqpoint{1.142102in}{2.303631in}}%
\pgfpathlineto{\pgfqpoint{1.142656in}{2.303570in}}%
\pgfpathlineto{\pgfqpoint{1.143658in}{2.303632in}}%
\pgfpathlineto{\pgfqpoint{1.143161in}{2.303242in}}%
\pgfpathlineto{\pgfqpoint{1.143756in}{2.303434in}}%
\pgfpathlineto{\pgfqpoint{1.144367in}{2.302747in}}%
\pgfpathlineto{\pgfqpoint{1.144880in}{2.303061in}}%
\pgfpathlineto{\pgfqpoint{1.146137in}{2.302282in}}%
\pgfpathlineto{\pgfqpoint{1.144927in}{2.303175in}}%
\pgfpathlineto{\pgfqpoint{1.146142in}{2.302324in}}%
\pgfpathlineto{\pgfqpoint{1.146356in}{2.302539in}}%
\pgfpathlineto{\pgfqpoint{1.147100in}{2.301799in}}%
\pgfpathlineto{\pgfqpoint{1.147232in}{2.301875in}}%
\pgfpathlineto{\pgfqpoint{1.147589in}{2.301400in}}%
\pgfpathlineto{\pgfqpoint{1.147237in}{2.301878in}}%
\pgfpathlineto{\pgfqpoint{1.148365in}{2.301554in}}%
\pgfpathlineto{\pgfqpoint{1.148884in}{2.301628in}}%
\pgfpathlineto{\pgfqpoint{1.148767in}{2.301265in}}%
\pgfpathlineto{\pgfqpoint{1.149435in}{2.301244in}}%
\pgfpathlineto{\pgfqpoint{1.149610in}{2.300964in}}%
\pgfpathlineto{\pgfqpoint{1.150432in}{2.301904in}}%
\pgfpathlineto{\pgfqpoint{1.150444in}{2.301877in}}%
\pgfpathlineto{\pgfqpoint{1.151421in}{2.301911in}}%
\pgfpathlineto{\pgfqpoint{1.150771in}{2.301404in}}%
\pgfpathlineto{\pgfqpoint{1.151534in}{2.301618in}}%
\pgfpathlineto{\pgfqpoint{1.152474in}{2.301509in}}%
\pgfpathlineto{\pgfqpoint{1.151738in}{2.301836in}}%
\pgfpathlineto{\pgfqpoint{1.152644in}{2.301702in}}%
\pgfpathlineto{\pgfqpoint{1.153683in}{2.301880in}}%
\pgfpathlineto{\pgfqpoint{1.153023in}{2.301402in}}%
\pgfpathlineto{\pgfqpoint{1.153766in}{2.301813in}}%
\pgfpathlineto{\pgfqpoint{1.154366in}{2.301373in}}%
\pgfpathlineto{\pgfqpoint{1.153933in}{2.301886in}}%
\pgfpathlineto{\pgfqpoint{1.154908in}{2.301506in}}%
\pgfpathlineto{\pgfqpoint{1.155927in}{2.301833in}}%
\pgfpathlineto{\pgfqpoint{1.155124in}{2.301269in}}%
\pgfpathlineto{\pgfqpoint{1.156094in}{2.301718in}}%
\pgfpathlineto{\pgfqpoint{1.156187in}{2.301527in}}%
\pgfpathlineto{\pgfqpoint{1.156750in}{2.302177in}}%
\pgfpathlineto{\pgfqpoint{1.157188in}{2.301913in}}%
\pgfpathlineto{\pgfqpoint{1.157546in}{2.302139in}}%
\pgfpathlineto{\pgfqpoint{1.157830in}{2.301579in}}%
\pgfpathlineto{\pgfqpoint{1.158302in}{2.302014in}}%
\pgfpathlineto{\pgfqpoint{1.158576in}{2.302149in}}%
\pgfpathlineto{\pgfqpoint{1.159275in}{2.301397in}}%
\pgfpathlineto{\pgfqpoint{1.159399in}{2.301810in}}%
\pgfpathlineto{\pgfqpoint{1.160720in}{2.302393in}}%
\pgfpathlineto{\pgfqpoint{1.159474in}{2.301637in}}%
\pgfpathlineto{\pgfqpoint{1.160754in}{2.302324in}}%
\pgfpathlineto{\pgfqpoint{1.161726in}{2.302090in}}%
\pgfpathlineto{\pgfqpoint{1.160864in}{2.302518in}}%
\pgfpathlineto{\pgfqpoint{1.161878in}{2.302260in}}%
\pgfpathlineto{\pgfqpoint{1.161925in}{2.302404in}}%
\pgfpathlineto{\pgfqpoint{1.162610in}{2.301801in}}%
\pgfpathlineto{\pgfqpoint{1.162981in}{2.301945in}}%
\pgfpathlineto{\pgfqpoint{1.164050in}{2.302229in}}%
\pgfpathlineto{\pgfqpoint{1.163051in}{2.301907in}}%
\pgfpathlineto{\pgfqpoint{1.164212in}{2.302096in}}%
\pgfpathlineto{\pgfqpoint{1.164240in}{2.302048in}}%
\pgfpathlineto{\pgfqpoint{1.165013in}{2.302945in}}%
\pgfpathlineto{\pgfqpoint{1.165257in}{2.302786in}}%
\pgfpathlineto{\pgfqpoint{1.166359in}{2.303309in}}%
\pgfpathlineto{\pgfqpoint{1.165367in}{2.302654in}}%
\pgfpathlineto{\pgfqpoint{1.166389in}{2.303263in}}%
\pgfpathlineto{\pgfqpoint{1.167075in}{2.303157in}}%
\pgfpathlineto{\pgfqpoint{1.167477in}{2.303655in}}%
\pgfpathlineto{\pgfqpoint{1.167479in}{2.303652in}}%
\pgfpathlineto{\pgfqpoint{1.167514in}{2.303768in}}%
\pgfpathlineto{\pgfqpoint{1.168036in}{2.303268in}}%
\pgfpathlineto{\pgfqpoint{1.168579in}{2.303462in}}%
\pgfpathlineto{\pgfqpoint{1.169634in}{2.303253in}}%
\pgfpathlineto{\pgfqpoint{1.168765in}{2.303721in}}%
\pgfpathlineto{\pgfqpoint{1.169728in}{2.303414in}}%
\pgfpathlineto{\pgfqpoint{1.170063in}{2.303690in}}%
\pgfpathlineto{\pgfqpoint{1.169771in}{2.303261in}}%
\pgfpathlineto{\pgfqpoint{1.170848in}{2.303473in}}%
\pgfpathlineto{\pgfqpoint{1.171130in}{2.303270in}}%
\pgfpathlineto{\pgfqpoint{1.171914in}{2.303840in}}%
\pgfpathlineto{\pgfqpoint{1.171921in}{2.303816in}}%
\pgfpathlineto{\pgfqpoint{1.172752in}{2.303952in}}%
\pgfpathlineto{\pgfqpoint{1.171955in}{2.303678in}}%
\pgfpathlineto{\pgfqpoint{1.172990in}{2.303524in}}%
\pgfpathlineto{\pgfqpoint{1.172998in}{2.303465in}}%
\pgfpathlineto{\pgfqpoint{1.173784in}{2.303889in}}%
\pgfpathlineto{\pgfqpoint{1.174085in}{2.303865in}}%
\pgfpathlineto{\pgfqpoint{1.175451in}{2.304774in}}%
\pgfpathlineto{\pgfqpoint{1.174095in}{2.303826in}}%
\pgfpathlineto{\pgfqpoint{1.175561in}{2.304629in}}%
\pgfpathlineto{\pgfqpoint{1.176695in}{2.304298in}}%
\pgfpathlineto{\pgfqpoint{1.176227in}{2.304910in}}%
\pgfpathlineto{\pgfqpoint{1.176699in}{2.304305in}}%
\pgfpathlineto{\pgfqpoint{1.177059in}{2.304621in}}%
\pgfpathlineto{\pgfqpoint{1.177748in}{2.304134in}}%
\pgfpathlineto{\pgfqpoint{1.177805in}{2.304202in}}%
\pgfpathlineto{\pgfqpoint{1.178125in}{2.303855in}}%
\pgfpathlineto{\pgfqpoint{1.178585in}{2.304326in}}%
\pgfpathlineto{\pgfqpoint{1.178916in}{2.304218in}}%
\pgfpathlineto{\pgfqpoint{1.180072in}{2.304598in}}%
\pgfpathlineto{\pgfqpoint{1.179293in}{2.303964in}}%
\pgfpathlineto{\pgfqpoint{1.180110in}{2.304520in}}%
\pgfpathlineto{\pgfqpoint{1.181143in}{2.303913in}}%
\pgfpathlineto{\pgfqpoint{1.180309in}{2.304600in}}%
\pgfpathlineto{\pgfqpoint{1.181238in}{2.304021in}}%
\pgfpathlineto{\pgfqpoint{1.181429in}{2.304363in}}%
\pgfpathlineto{\pgfqpoint{1.181779in}{2.303920in}}%
\pgfpathlineto{\pgfqpoint{1.182354in}{2.304010in}}%
\pgfpathlineto{\pgfqpoint{1.183481in}{2.303205in}}%
\pgfpathlineto{\pgfqpoint{1.183535in}{2.303372in}}%
\pgfpathlineto{\pgfqpoint{1.183549in}{2.303408in}}%
\pgfpathlineto{\pgfqpoint{1.184019in}{2.302855in}}%
\pgfpathlineto{\pgfqpoint{1.184630in}{2.303014in}}%
\pgfpathlineto{\pgfqpoint{1.185607in}{2.302958in}}%
\pgfpathlineto{\pgfqpoint{1.185102in}{2.303304in}}%
\pgfpathlineto{\pgfqpoint{1.185730in}{2.303170in}}%
\pgfpathlineto{\pgfqpoint{1.186008in}{2.303403in}}%
\pgfpathlineto{\pgfqpoint{1.186406in}{2.303031in}}%
\pgfpathlineto{\pgfqpoint{1.186837in}{2.303082in}}%
\pgfpathlineto{\pgfqpoint{1.187835in}{2.302904in}}%
\pgfpathlineto{\pgfqpoint{1.187419in}{2.303300in}}%
\pgfpathlineto{\pgfqpoint{1.187949in}{2.303065in}}%
\pgfpathlineto{\pgfqpoint{1.188879in}{2.303664in}}%
\pgfpathlineto{\pgfqpoint{1.189089in}{2.303505in}}%
\pgfpathlineto{\pgfqpoint{1.190264in}{2.303275in}}%
\pgfpathlineto{\pgfqpoint{1.189180in}{2.303596in}}%
\pgfpathlineto{\pgfqpoint{1.190334in}{2.303410in}}%
\pgfpathlineto{\pgfqpoint{1.190870in}{2.303691in}}%
\pgfpathlineto{\pgfqpoint{1.191361in}{2.303326in}}%
\pgfpathlineto{\pgfqpoint{1.191450in}{2.303465in}}%
\pgfpathlineto{\pgfqpoint{1.192449in}{2.303101in}}%
\pgfpathlineto{\pgfqpoint{1.191539in}{2.303567in}}%
\pgfpathlineto{\pgfqpoint{1.192560in}{2.303511in}}%
\pgfpathlineto{\pgfqpoint{1.195508in}{2.303436in}}%
\pgfpathlineto{\pgfqpoint{1.195513in}{2.303390in}}%
\pgfpathlineto{\pgfqpoint{1.196659in}{2.303361in}}%
\pgfpathlineto{\pgfqpoint{1.195527in}{2.303450in}}%
\pgfpathlineto{\pgfqpoint{1.196664in}{2.303407in}}%
\pgfpathlineto{\pgfqpoint{1.197707in}{2.303451in}}%
\pgfpathlineto{\pgfqpoint{1.196776in}{2.303364in}}%
\pgfpathlineto{\pgfqpoint{1.197777in}{2.303432in}}%
\pgfpathlineto{\pgfqpoint{1.198882in}{2.303360in}}%
\pgfpathlineto{\pgfqpoint{1.198024in}{2.303458in}}%
\pgfpathlineto{\pgfqpoint{1.198888in}{2.303410in}}%
\pgfpathlineto{\pgfqpoint{1.199870in}{2.303451in}}%
\pgfpathlineto{\pgfqpoint{1.198903in}{2.303361in}}%
\pgfpathlineto{\pgfqpoint{1.200000in}{2.303427in}}%
\pgfpathlineto{\pgfqpoint{1.200716in}{2.303352in}}%
\pgfpathlineto{\pgfqpoint{1.200024in}{2.303448in}}%
\pgfpathlineto{\pgfqpoint{1.201112in}{2.303407in}}%
\pgfpathlineto{\pgfqpoint{1.202177in}{2.303440in}}%
\pgfpathlineto{\pgfqpoint{1.201337in}{2.303366in}}%
\pgfpathlineto{\pgfqpoint{1.202224in}{2.303400in}}%
\pgfpathlineto{\pgfqpoint{1.203265in}{2.303447in}}%
\pgfpathlineto{\pgfqpoint{1.202345in}{2.303367in}}%
\pgfpathlineto{\pgfqpoint{1.203346in}{2.303426in}}%
\pgfpathlineto{\pgfqpoint{1.204159in}{2.303352in}}%
\pgfpathlineto{\pgfqpoint{1.203378in}{2.303447in}}%
\pgfpathlineto{\pgfqpoint{1.204459in}{2.303413in}}%
\pgfpathlineto{\pgfqpoint{1.205532in}{2.303448in}}%
\pgfpathlineto{\pgfqpoint{1.204492in}{2.303356in}}%
\pgfpathlineto{\pgfqpoint{1.205573in}{2.303425in}}%
\pgfpathlineto{\pgfqpoint{1.206616in}{2.303376in}}%
\pgfpathlineto{\pgfqpoint{1.205608in}{2.303446in}}%
\pgfpathlineto{\pgfqpoint{1.206685in}{2.303424in}}%
\pgfpathlineto{\pgfqpoint{1.207739in}{2.303445in}}%
\pgfpathlineto{\pgfqpoint{1.206728in}{2.303371in}}%
\pgfpathlineto{\pgfqpoint{1.207799in}{2.303410in}}%
\pgfpathlineto{\pgfqpoint{1.208891in}{2.303380in}}%
\pgfpathlineto{\pgfqpoint{1.207820in}{2.303440in}}%
\pgfpathlineto{\pgfqpoint{1.208912in}{2.303408in}}%
\pgfpathlineto{\pgfqpoint{1.210015in}{2.303442in}}%
\pgfpathlineto{\pgfqpoint{1.209145in}{2.303341in}}%
\pgfpathlineto{\pgfqpoint{1.210025in}{2.303425in}}%
\pgfpathlineto{\pgfqpoint{1.210903in}{2.303354in}}%
\pgfpathlineto{\pgfqpoint{1.210084in}{2.303446in}}%
\pgfpathlineto{\pgfqpoint{1.211141in}{2.303389in}}%
\pgfpathlineto{\pgfqpoint{1.212218in}{2.303453in}}%
\pgfpathlineto{\pgfqpoint{1.211302in}{2.303312in}}%
\pgfpathlineto{\pgfqpoint{1.212254in}{2.303420in}}%
\pgfpathlineto{\pgfqpoint{1.213064in}{2.303352in}}%
\pgfpathlineto{\pgfqpoint{1.213221in}{2.303787in}}%
\pgfpathlineto{\pgfqpoint{1.213346in}{2.303677in}}%
\pgfpathlineto{\pgfqpoint{1.213911in}{2.304062in}}%
\pgfpathlineto{\pgfqpoint{1.213458in}{2.303484in}}%
\pgfpathlineto{\pgfqpoint{1.214441in}{2.303481in}}%
\pgfpathlineto{\pgfqpoint{1.215366in}{2.302647in}}%
\pgfpathlineto{\pgfqpoint{1.214461in}{2.303499in}}%
\pgfpathlineto{\pgfqpoint{1.215570in}{2.302939in}}%
\pgfpathlineto{\pgfqpoint{1.216143in}{2.302768in}}%
\pgfpathlineto{\pgfqpoint{1.215965in}{2.303104in}}%
\pgfpathlineto{\pgfqpoint{1.216688in}{2.302997in}}%
\pgfpathlineto{\pgfqpoint{1.217601in}{2.303288in}}%
\pgfpathlineto{\pgfqpoint{1.216701in}{2.302948in}}%
\pgfpathlineto{\pgfqpoint{1.217805in}{2.303096in}}%
\pgfpathlineto{\pgfqpoint{1.218965in}{2.301903in}}%
\pgfpathlineto{\pgfqpoint{1.217861in}{2.303262in}}%
\pgfpathlineto{\pgfqpoint{1.218971in}{2.301920in}}%
\pgfpathlineto{\pgfqpoint{1.219547in}{2.301923in}}%
\pgfpathlineto{\pgfqpoint{1.219850in}{2.302327in}}%
\pgfpathlineto{\pgfqpoint{1.220062in}{2.302218in}}%
\pgfpathlineto{\pgfqpoint{1.220627in}{2.302876in}}%
\pgfpathlineto{\pgfqpoint{1.220305in}{2.302093in}}%
\pgfpathlineto{\pgfqpoint{1.221195in}{2.302478in}}%
\pgfpathlineto{\pgfqpoint{1.221502in}{2.302687in}}%
\pgfpathlineto{\pgfqpoint{1.221943in}{2.302069in}}%
\pgfpathlineto{\pgfqpoint{1.222280in}{2.302207in}}%
\pgfpathlineto{\pgfqpoint{1.222397in}{2.302135in}}%
\pgfpathlineto{\pgfqpoint{1.223042in}{2.302542in}}%
\pgfpathlineto{\pgfqpoint{1.223378in}{2.302483in}}%
\pgfpathlineto{\pgfqpoint{1.224567in}{2.303071in}}%
\pgfpathlineto{\pgfqpoint{1.223392in}{2.302422in}}%
\pgfpathlineto{\pgfqpoint{1.224586in}{2.303016in}}%
\pgfpathlineto{\pgfqpoint{1.224605in}{2.302994in}}%
\pgfpathlineto{\pgfqpoint{1.225545in}{2.303653in}}%
\pgfpathlineto{\pgfqpoint{1.225569in}{2.303594in}}%
\pgfpathlineto{\pgfqpoint{1.226697in}{2.303881in}}%
\pgfpathlineto{\pgfqpoint{1.226297in}{2.303308in}}%
\pgfpathlineto{\pgfqpoint{1.226712in}{2.303842in}}%
\pgfpathlineto{\pgfqpoint{1.226893in}{2.303724in}}%
\pgfpathlineto{\pgfqpoint{1.227299in}{2.304185in}}%
\pgfpathlineto{\pgfqpoint{1.227822in}{2.303885in}}%
\pgfpathlineto{\pgfqpoint{1.227822in}{2.303905in}}%
\pgfpathlineto{\pgfqpoint{1.228750in}{2.303293in}}%
\pgfpathlineto{\pgfqpoint{1.228920in}{2.303539in}}%
\pgfpathlineto{\pgfqpoint{1.229507in}{2.303163in}}%
\pgfpathlineto{\pgfqpoint{1.229259in}{2.303665in}}%
\pgfpathlineto{\pgfqpoint{1.230048in}{2.303656in}}%
\pgfpathlineto{\pgfqpoint{1.231067in}{2.303938in}}%
\pgfpathlineto{\pgfqpoint{1.230260in}{2.303205in}}%
\pgfpathlineto{\pgfqpoint{1.231170in}{2.303819in}}%
\pgfpathlineto{\pgfqpoint{1.231810in}{2.303381in}}%
\pgfpathlineto{\pgfqpoint{1.231361in}{2.303949in}}%
\pgfpathlineto{\pgfqpoint{1.232293in}{2.303727in}}%
\pgfpathlineto{\pgfqpoint{1.232800in}{2.304429in}}%
\pgfpathlineto{\pgfqpoint{1.233477in}{2.304298in}}%
\pgfpathlineto{\pgfqpoint{1.234510in}{2.303941in}}%
\pgfpathlineto{\pgfqpoint{1.233823in}{2.304513in}}%
\pgfpathlineto{\pgfqpoint{1.234699in}{2.304118in}}%
\pgfpathlineto{\pgfqpoint{1.235790in}{2.303944in}}%
\pgfpathlineto{\pgfqpoint{1.235095in}{2.303656in}}%
\pgfpathlineto{\pgfqpoint{1.235792in}{2.303920in}}%
\pgfpathlineto{\pgfqpoint{1.237070in}{2.303227in}}%
\pgfpathlineto{\pgfqpoint{1.235915in}{2.304129in}}%
\pgfpathlineto{\pgfqpoint{1.237187in}{2.303373in}}%
\pgfpathlineto{\pgfqpoint{1.237479in}{2.303730in}}%
\pgfpathlineto{\pgfqpoint{1.238049in}{2.302821in}}%
\pgfpathlineto{\pgfqpoint{1.238270in}{2.302844in}}%
\pgfpathlineto{\pgfqpoint{1.239206in}{2.303244in}}%
\pgfpathlineto{\pgfqpoint{1.238752in}{2.302635in}}%
\pgfpathlineto{\pgfqpoint{1.239401in}{2.302900in}}%
\pgfpathlineto{\pgfqpoint{1.239934in}{2.302290in}}%
\pgfpathlineto{\pgfqpoint{1.240529in}{2.302506in}}%
\pgfpathlineto{\pgfqpoint{1.241530in}{2.302676in}}%
\pgfpathlineto{\pgfqpoint{1.240722in}{2.302237in}}%
\pgfpathlineto{\pgfqpoint{1.241753in}{2.302414in}}%
\pgfpathlineto{\pgfqpoint{1.242751in}{2.302152in}}%
\pgfpathlineto{\pgfqpoint{1.242297in}{2.302566in}}%
\pgfpathlineto{\pgfqpoint{1.242878in}{2.302249in}}%
\pgfpathlineto{\pgfqpoint{1.243867in}{2.302478in}}%
\pgfpathlineto{\pgfqpoint{1.243169in}{2.301957in}}%
\pgfpathlineto{\pgfqpoint{1.243992in}{2.302284in}}%
\pgfpathlineto{\pgfqpoint{1.245464in}{2.302170in}}%
\pgfpathlineto{\pgfqpoint{1.244092in}{2.302415in}}%
\pgfpathlineto{\pgfqpoint{1.245559in}{2.302362in}}%
\pgfpathlineto{\pgfqpoint{1.246644in}{2.302575in}}%
\pgfpathlineto{\pgfqpoint{1.246036in}{2.301998in}}%
\pgfpathlineto{\pgfqpoint{1.246678in}{2.302543in}}%
\pgfpathlineto{\pgfqpoint{1.247573in}{2.302602in}}%
\pgfpathlineto{\pgfqpoint{1.246862in}{2.302783in}}%
\pgfpathlineto{\pgfqpoint{1.247721in}{2.302943in}}%
\pgfpathlineto{\pgfqpoint{1.247956in}{2.303449in}}%
\pgfpathlineto{\pgfqpoint{1.248827in}{2.302803in}}%
\pgfpathlineto{\pgfqpoint{1.249825in}{2.302689in}}%
\pgfpathlineto{\pgfqpoint{1.249449in}{2.302949in}}%
\pgfpathlineto{\pgfqpoint{1.249941in}{2.302876in}}%
\pgfpathlineto{\pgfqpoint{1.250868in}{2.303382in}}%
\pgfpathlineto{\pgfqpoint{1.250047in}{2.302752in}}%
\pgfpathlineto{\pgfqpoint{1.251069in}{2.303268in}}%
\pgfpathlineto{\pgfqpoint{1.251361in}{2.303316in}}%
\pgfpathlineto{\pgfqpoint{1.252199in}{2.303039in}}%
\pgfpathlineto{\pgfqpoint{1.253355in}{2.303338in}}%
\pgfpathlineto{\pgfqpoint{1.252446in}{2.302829in}}%
\pgfpathlineto{\pgfqpoint{1.253357in}{2.303314in}}%
\pgfpathlineto{\pgfqpoint{1.254341in}{2.303064in}}%
\pgfpathlineto{\pgfqpoint{1.254012in}{2.303597in}}%
\pgfpathlineto{\pgfqpoint{1.254478in}{2.303093in}}%
\pgfpathlineto{\pgfqpoint{1.255555in}{2.303374in}}%
\pgfpathlineto{\pgfqpoint{1.254603in}{2.302944in}}%
\pgfpathlineto{\pgfqpoint{1.255599in}{2.303273in}}%
\pgfpathlineto{\pgfqpoint{1.256571in}{2.303709in}}%
\pgfpathlineto{\pgfqpoint{1.255681in}{2.303038in}}%
\pgfpathlineto{\pgfqpoint{1.256720in}{2.303399in}}%
\pgfpathlineto{\pgfqpoint{1.258039in}{2.303267in}}%
\pgfpathlineto{\pgfqpoint{1.256725in}{2.303417in}}%
\pgfpathlineto{\pgfqpoint{1.258041in}{2.303305in}}%
\pgfpathlineto{\pgfqpoint{1.259017in}{2.303830in}}%
\pgfpathlineto{\pgfqpoint{1.259170in}{2.303661in}}%
\pgfpathlineto{\pgfqpoint{1.259375in}{2.303336in}}%
\pgfpathlineto{\pgfqpoint{1.259553in}{2.303918in}}%
\pgfpathlineto{\pgfqpoint{1.260281in}{2.303769in}}%
\pgfpathlineto{\pgfqpoint{1.260965in}{2.303848in}}%
\pgfpathlineto{\pgfqpoint{1.260634in}{2.303407in}}%
\pgfpathlineto{\pgfqpoint{1.261372in}{2.303491in}}%
\pgfpathlineto{\pgfqpoint{1.262199in}{2.303007in}}%
\pgfpathlineto{\pgfqpoint{1.261655in}{2.303646in}}%
\pgfpathlineto{\pgfqpoint{1.262595in}{2.303249in}}%
\pgfpathlineto{\pgfqpoint{1.263499in}{2.303644in}}%
\pgfpathlineto{\pgfqpoint{1.262780in}{2.303080in}}%
\pgfpathlineto{\pgfqpoint{1.263721in}{2.303531in}}%
\pgfpathlineto{\pgfqpoint{1.264247in}{2.303098in}}%
\pgfpathlineto{\pgfqpoint{1.264787in}{2.303966in}}%
\pgfpathlineto{\pgfqpoint{1.264801in}{2.303925in}}%
\pgfpathlineto{\pgfqpoint{1.265248in}{2.304012in}}%
\pgfpathlineto{\pgfqpoint{1.265712in}{2.303101in}}%
\pgfpathlineto{\pgfqpoint{1.265814in}{2.303205in}}%
\pgfpathlineto{\pgfqpoint{1.265831in}{2.303166in}}%
\pgfpathlineto{\pgfqpoint{1.266682in}{2.303875in}}%
\pgfpathlineto{\pgfqpoint{1.266872in}{2.303780in}}%
\pgfpathlineto{\pgfqpoint{1.267444in}{2.304036in}}%
\pgfpathlineto{\pgfqpoint{1.266930in}{2.303609in}}%
\pgfpathlineto{\pgfqpoint{1.267978in}{2.303645in}}%
\pgfpathlineto{\pgfqpoint{1.269130in}{2.303352in}}%
\pgfpathlineto{\pgfqpoint{1.267990in}{2.303683in}}%
\pgfpathlineto{\pgfqpoint{1.269196in}{2.303414in}}%
\pgfpathlineto{\pgfqpoint{1.270288in}{2.303444in}}%
\pgfpathlineto{\pgfqpoint{1.269237in}{2.303369in}}%
\pgfpathlineto{\pgfqpoint{1.270309in}{2.303428in}}%
\pgfpathlineto{\pgfqpoint{1.271223in}{2.303368in}}%
\pgfpathlineto{\pgfqpoint{1.270355in}{2.303443in}}%
\pgfpathlineto{\pgfqpoint{1.271419in}{2.303413in}}%
\pgfpathlineto{\pgfqpoint{1.272395in}{2.303451in}}%
\pgfpathlineto{\pgfqpoint{1.271473in}{2.303358in}}%
\pgfpathlineto{\pgfqpoint{1.272531in}{2.303417in}}%
\pgfpathlineto{\pgfqpoint{1.272599in}{2.303372in}}%
\pgfpathlineto{\pgfqpoint{1.273646in}{2.303436in}}%
\pgfpathlineto{\pgfqpoint{1.274718in}{2.303355in}}%
\pgfpathlineto{\pgfqpoint{1.273784in}{2.303442in}}%
\pgfpathlineto{\pgfqpoint{1.274758in}{2.303391in}}%
\pgfpathlineto{\pgfqpoint{1.275851in}{2.303370in}}%
\pgfpathlineto{\pgfqpoint{1.274800in}{2.303446in}}%
\pgfpathlineto{\pgfqpoint{1.275875in}{2.303414in}}%
\pgfpathlineto{\pgfqpoint{1.277115in}{2.303432in}}%
\pgfpathlineto{\pgfqpoint{1.275902in}{2.303360in}}%
\pgfpathlineto{\pgfqpoint{1.277122in}{2.303388in}}%
\pgfpathlineto{\pgfqpoint{1.278300in}{2.303386in}}%
\pgfpathlineto{\pgfqpoint{1.277129in}{2.303429in}}%
\pgfpathlineto{\pgfqpoint{1.278317in}{2.303422in}}%
\pgfpathlineto{\pgfqpoint{1.279446in}{2.303439in}}%
\pgfpathlineto{\pgfqpoint{1.278321in}{2.303383in}}%
\pgfpathlineto{\pgfqpoint{1.279447in}{2.303425in}}%
\pgfpathlineto{\pgfqpoint{1.280379in}{2.303347in}}%
\pgfpathlineto{\pgfqpoint{1.279488in}{2.303444in}}%
\pgfpathlineto{\pgfqpoint{1.280560in}{2.303421in}}%
\pgfpathlineto{\pgfqpoint{1.281631in}{2.303441in}}%
\pgfpathlineto{\pgfqpoint{1.280572in}{2.303393in}}%
\pgfpathlineto{\pgfqpoint{1.281690in}{2.303393in}}%
\pgfpathlineto{\pgfqpoint{1.282739in}{2.303376in}}%
\pgfpathlineto{\pgfqpoint{1.281714in}{2.303453in}}%
\pgfpathlineto{\pgfqpoint{1.282809in}{2.303413in}}%
\pgfpathlineto{\pgfqpoint{1.284111in}{2.303441in}}%
\pgfpathlineto{\pgfqpoint{1.282815in}{2.303399in}}%
\pgfpathlineto{\pgfqpoint{1.284114in}{2.303416in}}%
\pgfpathlineto{\pgfqpoint{1.285134in}{2.303363in}}%
\pgfpathlineto{\pgfqpoint{1.284125in}{2.303439in}}%
\pgfpathlineto{\pgfqpoint{1.285229in}{2.303410in}}%
\pgfpathlineto{\pgfqpoint{1.286209in}{2.303359in}}%
\pgfpathlineto{\pgfqpoint{1.285278in}{2.303443in}}%
\pgfpathlineto{\pgfqpoint{1.286355in}{2.303401in}}%
\pgfpathlineto{\pgfqpoint{1.287383in}{2.303448in}}%
\pgfpathlineto{\pgfqpoint{1.286500in}{2.303339in}}%
\pgfpathlineto{\pgfqpoint{1.287467in}{2.303424in}}%
\pgfpathlineto{\pgfqpoint{1.288471in}{2.303296in}}%
\pgfpathlineto{\pgfqpoint{1.288085in}{2.303686in}}%
\pgfpathlineto{\pgfqpoint{1.288576in}{2.303491in}}%
\pgfpathlineto{\pgfqpoint{1.291452in}{2.304078in}}%
\pgfpathlineto{\pgfqpoint{1.288577in}{2.303505in}}%
\pgfpathlineto{\pgfqpoint{1.291467in}{2.304206in}}%
\pgfpathlineto{\pgfqpoint{1.291468in}{2.304217in}}%
\pgfpathlineto{\pgfqpoint{1.292034in}{2.303531in}}%
\pgfpathlineto{\pgfqpoint{1.292555in}{2.303855in}}%
\pgfpathlineto{\pgfqpoint{1.292559in}{2.303848in}}%
\pgfpathlineto{\pgfqpoint{1.292671in}{2.304160in}}%
\pgfpathlineto{\pgfqpoint{1.293652in}{2.304031in}}%
\pgfpathlineto{\pgfqpoint{1.293778in}{2.304162in}}%
\pgfpathlineto{\pgfqpoint{1.294246in}{2.303535in}}%
\pgfpathlineto{\pgfqpoint{1.294732in}{2.303637in}}%
\pgfpathlineto{\pgfqpoint{1.294758in}{2.303625in}}%
\pgfpathlineto{\pgfqpoint{1.295546in}{2.304122in}}%
\pgfpathlineto{\pgfqpoint{1.295819in}{2.303929in}}%
\pgfpathlineto{\pgfqpoint{1.296580in}{2.304055in}}%
\pgfpathlineto{\pgfqpoint{1.296208in}{2.303700in}}%
\pgfpathlineto{\pgfqpoint{1.296936in}{2.303960in}}%
\pgfpathlineto{\pgfqpoint{1.296998in}{2.303860in}}%
\pgfpathlineto{\pgfqpoint{1.297797in}{2.304489in}}%
\pgfpathlineto{\pgfqpoint{1.298016in}{2.304417in}}%
\pgfpathlineto{\pgfqpoint{1.299141in}{2.304919in}}%
\pgfpathlineto{\pgfqpoint{1.298228in}{2.304183in}}%
\pgfpathlineto{\pgfqpoint{1.299190in}{2.304787in}}%
\pgfpathlineto{\pgfqpoint{1.299281in}{2.304594in}}%
\pgfpathlineto{\pgfqpoint{1.300180in}{2.305050in}}%
\pgfpathlineto{\pgfqpoint{1.300293in}{2.304968in}}%
\pgfpathlineto{\pgfqpoint{1.301379in}{2.305357in}}%
\pgfpathlineto{\pgfqpoint{1.300530in}{2.304908in}}%
\pgfpathlineto{\pgfqpoint{1.301418in}{2.305280in}}%
\pgfpathlineto{\pgfqpoint{1.301618in}{2.305005in}}%
\pgfpathlineto{\pgfqpoint{1.301932in}{2.305454in}}%
\pgfpathlineto{\pgfqpoint{1.302528in}{2.305313in}}%
\pgfpathlineto{\pgfqpoint{1.302546in}{2.305412in}}%
\pgfpathlineto{\pgfqpoint{1.303338in}{2.304733in}}%
\pgfpathlineto{\pgfqpoint{1.303608in}{2.304783in}}%
\pgfpathlineto{\pgfqpoint{1.304601in}{2.304728in}}%
\pgfpathlineto{\pgfqpoint{1.303998in}{2.305267in}}%
\pgfpathlineto{\pgfqpoint{1.304710in}{2.304917in}}%
\pgfpathlineto{\pgfqpoint{1.305236in}{2.305000in}}%
\pgfpathlineto{\pgfqpoint{1.305691in}{2.304381in}}%
\pgfpathlineto{\pgfqpoint{1.305757in}{2.304423in}}%
\pgfpathlineto{\pgfqpoint{1.305978in}{2.304065in}}%
\pgfpathlineto{\pgfqpoint{1.306839in}{2.304687in}}%
\pgfpathlineto{\pgfqpoint{1.306861in}{2.304631in}}%
\pgfpathlineto{\pgfqpoint{1.308128in}{2.303559in}}%
\pgfpathlineto{\pgfqpoint{1.308254in}{2.303763in}}%
\pgfpathlineto{\pgfqpoint{1.309340in}{2.303973in}}%
\pgfpathlineto{\pgfqpoint{1.308762in}{2.303562in}}%
\pgfpathlineto{\pgfqpoint{1.309381in}{2.303940in}}%
\pgfpathlineto{\pgfqpoint{1.310492in}{2.303474in}}%
\pgfpathlineto{\pgfqpoint{1.309555in}{2.304093in}}%
\pgfpathlineto{\pgfqpoint{1.310507in}{2.303482in}}%
\pgfpathlineto{\pgfqpoint{1.310534in}{2.303586in}}%
\pgfpathlineto{\pgfqpoint{1.311476in}{2.302916in}}%
\pgfpathlineto{\pgfqpoint{1.311607in}{2.303190in}}%
\pgfpathlineto{\pgfqpoint{1.311649in}{2.303112in}}%
\pgfpathlineto{\pgfqpoint{1.312496in}{2.303596in}}%
\pgfpathlineto{\pgfqpoint{1.312716in}{2.303253in}}%
\pgfpathlineto{\pgfqpoint{1.312766in}{2.303286in}}%
\pgfpathlineto{\pgfqpoint{1.313794in}{2.302770in}}%
\pgfpathlineto{\pgfqpoint{1.313796in}{2.302782in}}%
\pgfpathlineto{\pgfqpoint{1.314891in}{2.302271in}}%
\pgfpathlineto{\pgfqpoint{1.313938in}{2.302896in}}%
\pgfpathlineto{\pgfqpoint{1.314925in}{2.302351in}}%
\pgfpathlineto{\pgfqpoint{1.315815in}{2.302317in}}%
\pgfpathlineto{\pgfqpoint{1.316024in}{2.302839in}}%
\pgfpathlineto{\pgfqpoint{1.316461in}{2.303394in}}%
\pgfpathlineto{\pgfqpoint{1.316047in}{2.302809in}}%
\pgfpathlineto{\pgfqpoint{1.317170in}{2.303160in}}%
\pgfpathlineto{\pgfqpoint{1.318305in}{2.302393in}}%
\pgfpathlineto{\pgfqpoint{1.317564in}{2.303492in}}%
\pgfpathlineto{\pgfqpoint{1.318339in}{2.302511in}}%
\pgfpathlineto{\pgfqpoint{1.320534in}{2.302632in}}%
\pgfpathlineto{\pgfqpoint{1.318347in}{2.302495in}}%
\pgfpathlineto{\pgfqpoint{1.320580in}{2.302452in}}%
\pgfpathlineto{\pgfqpoint{1.321268in}{2.302663in}}%
\pgfpathlineto{\pgfqpoint{1.321432in}{2.303086in}}%
\pgfpathlineto{\pgfqpoint{1.321580in}{2.302934in}}%
\pgfpathlineto{\pgfqpoint{1.322090in}{2.303116in}}%
\pgfpathlineto{\pgfqpoint{1.322423in}{2.302622in}}%
\pgfpathlineto{\pgfqpoint{1.322688in}{2.302864in}}%
\pgfpathlineto{\pgfqpoint{1.324542in}{2.302418in}}%
\pgfpathlineto{\pgfqpoint{1.322966in}{2.303295in}}%
\pgfpathlineto{\pgfqpoint{1.324577in}{2.302496in}}%
\pgfpathlineto{\pgfqpoint{1.325548in}{2.302649in}}%
\pgfpathlineto{\pgfqpoint{1.325836in}{2.302329in}}%
\pgfpathlineto{\pgfqpoint{1.326204in}{2.301903in}}%
\pgfpathlineto{\pgfqpoint{1.326785in}{2.302631in}}%
\pgfpathlineto{\pgfqpoint{1.326929in}{2.302584in}}%
\pgfpathlineto{\pgfqpoint{1.327965in}{2.302795in}}%
\pgfpathlineto{\pgfqpoint{1.327541in}{2.302199in}}%
\pgfpathlineto{\pgfqpoint{1.328045in}{2.302637in}}%
\pgfpathlineto{\pgfqpoint{1.329044in}{2.303096in}}%
\pgfpathlineto{\pgfqpoint{1.328236in}{2.302477in}}%
\pgfpathlineto{\pgfqpoint{1.329191in}{2.303048in}}%
\pgfpathlineto{\pgfqpoint{1.329372in}{2.302934in}}%
\pgfpathlineto{\pgfqpoint{1.330194in}{2.303671in}}%
\pgfpathlineto{\pgfqpoint{1.330268in}{2.303570in}}%
\pgfpathlineto{\pgfqpoint{1.331419in}{2.304145in}}%
\pgfpathlineto{\pgfqpoint{1.330305in}{2.303469in}}%
\pgfpathlineto{\pgfqpoint{1.331420in}{2.304134in}}%
\pgfpathlineto{\pgfqpoint{1.332554in}{2.303777in}}%
\pgfpathlineto{\pgfqpoint{1.331537in}{2.304253in}}%
\pgfpathlineto{\pgfqpoint{1.332556in}{2.303810in}}%
\pgfpathlineto{\pgfqpoint{1.333443in}{2.304215in}}%
\pgfpathlineto{\pgfqpoint{1.332563in}{2.303874in}}%
\pgfpathlineto{\pgfqpoint{1.333515in}{2.304414in}}%
\pgfpathlineto{\pgfqpoint{1.333515in}{2.304429in}}%
\pgfpathlineto{\pgfqpoint{1.333974in}{2.303993in}}%
\pgfpathlineto{\pgfqpoint{1.334610in}{2.304086in}}%
\pgfpathlineto{\pgfqpoint{1.335173in}{2.304124in}}%
\pgfpathlineto{\pgfqpoint{1.335543in}{2.303699in}}%
\pgfpathlineto{\pgfqpoint{1.335687in}{2.303693in}}%
\pgfpathlineto{\pgfqpoint{1.336473in}{2.303988in}}%
\pgfpathlineto{\pgfqpoint{1.336805in}{2.303677in}}%
\pgfpathlineto{\pgfqpoint{1.336806in}{2.303664in}}%
\pgfpathlineto{\pgfqpoint{1.337452in}{2.304216in}}%
\pgfpathlineto{\pgfqpoint{1.337901in}{2.303947in}}%
\pgfpathlineto{\pgfqpoint{1.338063in}{2.303955in}}%
\pgfpathlineto{\pgfqpoint{1.338609in}{2.303423in}}%
\pgfpathlineto{\pgfqpoint{1.338999in}{2.303785in}}%
\pgfpathlineto{\pgfqpoint{1.339316in}{2.303429in}}%
\pgfpathlineto{\pgfqpoint{1.339804in}{2.303826in}}%
\pgfpathlineto{\pgfqpoint{1.340115in}{2.303597in}}%
\pgfpathlineto{\pgfqpoint{1.341282in}{2.303917in}}%
\pgfpathlineto{\pgfqpoint{1.340136in}{2.303678in}}%
\pgfpathlineto{\pgfqpoint{1.341307in}{2.304051in}}%
\pgfpathlineto{\pgfqpoint{1.341887in}{2.304216in}}%
\pgfpathlineto{\pgfqpoint{1.342210in}{2.303627in}}%
\pgfpathlineto{\pgfqpoint{1.342395in}{2.303819in}}%
\pgfpathlineto{\pgfqpoint{1.343420in}{2.303948in}}%
\pgfpathlineto{\pgfqpoint{1.342675in}{2.304150in}}%
\pgfpathlineto{\pgfqpoint{1.343466in}{2.304128in}}%
\pgfpathlineto{\pgfqpoint{1.343773in}{2.304270in}}%
\pgfpathlineto{\pgfqpoint{1.344083in}{2.303349in}}%
\pgfpathlineto{\pgfqpoint{1.344469in}{2.303427in}}%
\pgfpathlineto{\pgfqpoint{1.345230in}{2.303328in}}%
\pgfpathlineto{\pgfqpoint{1.344800in}{2.303455in}}%
\pgfpathlineto{\pgfqpoint{1.345582in}{2.303385in}}%
\pgfpathlineto{\pgfqpoint{1.346693in}{2.303443in}}%
\pgfpathlineto{\pgfqpoint{1.345615in}{2.303378in}}%
\pgfpathlineto{\pgfqpoint{1.346698in}{2.303415in}}%
\pgfpathlineto{\pgfqpoint{1.347739in}{2.303374in}}%
\pgfpathlineto{\pgfqpoint{1.346710in}{2.303437in}}%
\pgfpathlineto{\pgfqpoint{1.347823in}{2.303406in}}%
\pgfpathlineto{\pgfqpoint{1.348853in}{2.303448in}}%
\pgfpathlineto{\pgfqpoint{1.347974in}{2.303365in}}%
\pgfpathlineto{\pgfqpoint{1.348936in}{2.303430in}}%
\pgfpathlineto{\pgfqpoint{1.349489in}{2.303314in}}%
\pgfpathlineto{\pgfqpoint{1.348948in}{2.303446in}}%
\pgfpathlineto{\pgfqpoint{1.350047in}{2.303417in}}%
\pgfpathlineto{\pgfqpoint{1.351191in}{2.303375in}}%
\pgfpathlineto{\pgfqpoint{1.350049in}{2.303436in}}%
\pgfpathlineto{\pgfqpoint{1.351216in}{2.303405in}}%
\pgfpathlineto{\pgfqpoint{1.352279in}{2.303444in}}%
\pgfpathlineto{\pgfqpoint{1.351523in}{2.303362in}}%
\pgfpathlineto{\pgfqpoint{1.352327in}{2.303418in}}%
\pgfpathlineto{\pgfqpoint{1.353329in}{2.303372in}}%
\pgfpathlineto{\pgfqpoint{1.352377in}{2.303452in}}%
\pgfpathlineto{\pgfqpoint{1.353449in}{2.303416in}}%
\pgfpathlineto{\pgfqpoint{1.354547in}{2.303444in}}%
\pgfpathlineto{\pgfqpoint{1.353612in}{2.303362in}}%
\pgfpathlineto{\pgfqpoint{1.354560in}{2.303434in}}%
\pgfpathlineto{\pgfqpoint{1.355507in}{2.303363in}}%
\pgfpathlineto{\pgfqpoint{1.354704in}{2.303452in}}%
\pgfpathlineto{\pgfqpoint{1.355672in}{2.303401in}}%
\pgfpathlineto{\pgfqpoint{1.356736in}{2.303448in}}%
\pgfpathlineto{\pgfqpoint{1.355705in}{2.303384in}}%
\pgfpathlineto{\pgfqpoint{1.356808in}{2.303398in}}%
\pgfpathlineto{\pgfqpoint{1.357975in}{2.303442in}}%
\pgfpathlineto{\pgfqpoint{1.356811in}{2.303386in}}%
\pgfpathlineto{\pgfqpoint{1.357984in}{2.303426in}}%
\pgfpathlineto{\pgfqpoint{1.359034in}{2.303375in}}%
\pgfpathlineto{\pgfqpoint{1.358251in}{2.303452in}}%
\pgfpathlineto{\pgfqpoint{1.359095in}{2.303417in}}%
\pgfpathlineto{\pgfqpoint{1.359912in}{2.303454in}}%
\pgfpathlineto{\pgfqpoint{1.359395in}{2.303354in}}%
\pgfpathlineto{\pgfqpoint{1.360206in}{2.303434in}}%
\pgfpathlineto{\pgfqpoint{1.361172in}{2.303365in}}%
\pgfpathlineto{\pgfqpoint{1.360272in}{2.303446in}}%
\pgfpathlineto{\pgfqpoint{1.361318in}{2.303407in}}%
\pgfpathlineto{\pgfqpoint{1.362914in}{2.303309in}}%
\pgfpathlineto{\pgfqpoint{1.361329in}{2.303386in}}%
\pgfpathlineto{\pgfqpoint{1.362918in}{2.303262in}}%
\pgfpathlineto{\pgfqpoint{1.364042in}{2.302536in}}%
\pgfpathlineto{\pgfqpoint{1.364062in}{2.302594in}}%
\pgfpathlineto{\pgfqpoint{1.364337in}{2.302843in}}%
\pgfpathlineto{\pgfqpoint{1.365011in}{2.302261in}}%
\pgfpathlineto{\pgfqpoint{1.365152in}{2.302329in}}%
\pgfpathlineto{\pgfqpoint{1.366243in}{2.302415in}}%
\pgfpathlineto{\pgfqpoint{1.365973in}{2.302783in}}%
\pgfpathlineto{\pgfqpoint{1.366259in}{2.302460in}}%
\pgfpathlineto{\pgfqpoint{1.366963in}{2.302632in}}%
\pgfpathlineto{\pgfqpoint{1.367236in}{2.302182in}}%
\pgfpathlineto{\pgfqpoint{1.367377in}{2.302381in}}%
\pgfpathlineto{\pgfqpoint{1.368718in}{2.301293in}}%
\pgfpathlineto{\pgfqpoint{1.368741in}{2.301441in}}%
\pgfpathlineto{\pgfqpoint{1.370127in}{2.302195in}}%
\pgfpathlineto{\pgfqpoint{1.368743in}{2.301428in}}%
\pgfpathlineto{\pgfqpoint{1.370192in}{2.302019in}}%
\pgfpathlineto{\pgfqpoint{1.370571in}{2.301712in}}%
\pgfpathlineto{\pgfqpoint{1.370275in}{2.302188in}}%
\pgfpathlineto{\pgfqpoint{1.371296in}{2.302142in}}%
\pgfpathlineto{\pgfqpoint{1.372519in}{2.302304in}}%
\pgfpathlineto{\pgfqpoint{1.371303in}{2.302120in}}%
\pgfpathlineto{\pgfqpoint{1.372538in}{2.302205in}}%
\pgfpathlineto{\pgfqpoint{1.372547in}{2.302119in}}%
\pgfpathlineto{\pgfqpoint{1.373631in}{2.302695in}}%
\pgfpathlineto{\pgfqpoint{1.374154in}{2.302733in}}%
\pgfpathlineto{\pgfqpoint{1.374466in}{2.302142in}}%
\pgfpathlineto{\pgfqpoint{1.374725in}{2.302372in}}%
\pgfpathlineto{\pgfqpoint{1.375780in}{2.303109in}}%
\pgfpathlineto{\pgfqpoint{1.374886in}{2.302299in}}%
\pgfpathlineto{\pgfqpoint{1.375912in}{2.302908in}}%
\pgfpathlineto{\pgfqpoint{1.376942in}{2.302537in}}%
\pgfpathlineto{\pgfqpoint{1.375956in}{2.302961in}}%
\pgfpathlineto{\pgfqpoint{1.377056in}{2.302710in}}%
\pgfpathlineto{\pgfqpoint{1.377152in}{2.302770in}}%
\pgfpathlineto{\pgfqpoint{1.378071in}{2.302170in}}%
\pgfpathlineto{\pgfqpoint{1.378080in}{2.302196in}}%
\pgfpathlineto{\pgfqpoint{1.379795in}{2.301446in}}%
\pgfpathlineto{\pgfqpoint{1.378087in}{2.302260in}}%
\pgfpathlineto{\pgfqpoint{1.379816in}{2.301521in}}%
\pgfpathlineto{\pgfqpoint{1.380917in}{2.301823in}}%
\pgfpathlineto{\pgfqpoint{1.380940in}{2.301771in}}%
\pgfpathlineto{\pgfqpoint{1.382808in}{2.302874in}}%
\pgfpathlineto{\pgfqpoint{1.380943in}{2.301760in}}%
\pgfpathlineto{\pgfqpoint{1.382836in}{2.302785in}}%
\pgfpathlineto{\pgfqpoint{1.383580in}{2.302520in}}%
\pgfpathlineto{\pgfqpoint{1.383178in}{2.302871in}}%
\pgfpathlineto{\pgfqpoint{1.383951in}{2.302696in}}%
\pgfpathlineto{\pgfqpoint{1.385655in}{2.301933in}}%
\pgfpathlineto{\pgfqpoint{1.383968in}{2.302780in}}%
\pgfpathlineto{\pgfqpoint{1.385823in}{2.302217in}}%
\pgfpathlineto{\pgfqpoint{1.386169in}{2.302426in}}%
\pgfpathlineto{\pgfqpoint{1.386868in}{2.301947in}}%
\pgfpathlineto{\pgfqpoint{1.386934in}{2.302199in}}%
\pgfpathlineto{\pgfqpoint{1.388363in}{2.301436in}}%
\pgfpathlineto{\pgfqpoint{1.386978in}{2.302355in}}%
\pgfpathlineto{\pgfqpoint{1.388502in}{2.301581in}}%
\pgfpathlineto{\pgfqpoint{1.389219in}{2.301751in}}%
\pgfpathlineto{\pgfqpoint{1.389415in}{2.301304in}}%
\pgfpathlineto{\pgfqpoint{1.389606in}{2.301432in}}%
\pgfpathlineto{\pgfqpoint{1.389606in}{2.301425in}}%
\pgfpathlineto{\pgfqpoint{1.390624in}{2.302106in}}%
\pgfpathlineto{\pgfqpoint{1.390663in}{2.301972in}}%
\pgfpathlineto{\pgfqpoint{1.390665in}{2.302000in}}%
\pgfpathlineto{\pgfqpoint{1.391647in}{2.301328in}}%
\pgfpathlineto{\pgfqpoint{1.391745in}{2.301359in}}%
\pgfpathlineto{\pgfqpoint{1.392625in}{2.301519in}}%
\pgfpathlineto{\pgfqpoint{1.392336in}{2.300899in}}%
\pgfpathlineto{\pgfqpoint{1.392867in}{2.301385in}}%
\pgfpathlineto{\pgfqpoint{1.393984in}{2.300865in}}%
\pgfpathlineto{\pgfqpoint{1.392938in}{2.301400in}}%
\pgfpathlineto{\pgfqpoint{1.393993in}{2.300875in}}%
\pgfpathlineto{\pgfqpoint{1.393994in}{2.300879in}}%
\pgfpathlineto{\pgfqpoint{1.394480in}{2.300189in}}%
\pgfpathlineto{\pgfqpoint{1.395026in}{2.300379in}}%
\pgfpathlineto{\pgfqpoint{1.396004in}{2.300132in}}%
\pgfpathlineto{\pgfqpoint{1.395193in}{2.300451in}}%
\pgfpathlineto{\pgfqpoint{1.396140in}{2.300391in}}%
\pgfpathlineto{\pgfqpoint{1.396711in}{2.300499in}}%
\pgfpathlineto{\pgfqpoint{1.397175in}{2.299825in}}%
\pgfpathlineto{\pgfqpoint{1.397208in}{2.299891in}}%
\pgfpathlineto{\pgfqpoint{1.398185in}{2.299472in}}%
\pgfpathlineto{\pgfqpoint{1.397375in}{2.299950in}}%
\pgfpathlineto{\pgfqpoint{1.398340in}{2.299488in}}%
\pgfpathlineto{\pgfqpoint{1.398437in}{2.299733in}}%
\pgfpathlineto{\pgfqpoint{1.398866in}{2.299248in}}%
\pgfpathlineto{\pgfqpoint{1.399442in}{2.299310in}}%
\pgfpathlineto{\pgfqpoint{1.399555in}{2.299262in}}%
\pgfpathlineto{\pgfqpoint{1.400195in}{2.299889in}}%
\pgfpathlineto{\pgfqpoint{1.400486in}{2.299845in}}%
\pgfpathlineto{\pgfqpoint{1.401578in}{2.300004in}}%
\pgfpathlineto{\pgfqpoint{1.400728in}{2.299724in}}%
\pgfpathlineto{\pgfqpoint{1.401609in}{2.299900in}}%
\pgfpathlineto{\pgfqpoint{1.401644in}{2.299830in}}%
\pgfpathlineto{\pgfqpoint{1.402524in}{2.300718in}}%
\pgfpathlineto{\pgfqpoint{1.402606in}{2.300646in}}%
\pgfpathlineto{\pgfqpoint{1.403521in}{2.301292in}}%
\pgfpathlineto{\pgfqpoint{1.402631in}{2.300619in}}%
\pgfpathlineto{\pgfqpoint{1.403748in}{2.301167in}}%
\pgfpathlineto{\pgfqpoint{1.403921in}{2.301019in}}%
\pgfpathlineto{\pgfqpoint{1.404820in}{2.301849in}}%
\pgfpathlineto{\pgfqpoint{1.404824in}{2.301843in}}%
\pgfpathlineto{\pgfqpoint{1.406335in}{2.303393in}}%
\pgfpathlineto{\pgfqpoint{1.406425in}{2.303155in}}%
\pgfpathlineto{\pgfqpoint{1.407598in}{2.302747in}}%
\pgfpathlineto{\pgfqpoint{1.406508in}{2.303360in}}%
\pgfpathlineto{\pgfqpoint{1.407599in}{2.302755in}}%
\pgfpathlineto{\pgfqpoint{1.407848in}{2.302946in}}%
\pgfpathlineto{\pgfqpoint{1.408683in}{2.302219in}}%
\pgfpathlineto{\pgfqpoint{1.408696in}{2.302297in}}%
\pgfpathlineto{\pgfqpoint{1.410619in}{2.302127in}}%
\pgfpathlineto{\pgfqpoint{1.408813in}{2.302628in}}%
\pgfpathlineto{\pgfqpoint{1.410646in}{2.302244in}}%
\pgfpathlineto{\pgfqpoint{1.411599in}{2.302988in}}%
\pgfpathlineto{\pgfqpoint{1.410650in}{2.302241in}}%
\pgfpathlineto{\pgfqpoint{1.411826in}{2.302847in}}%
\pgfpathlineto{\pgfqpoint{1.412879in}{2.302389in}}%
\pgfpathlineto{\pgfqpoint{1.412285in}{2.302978in}}%
\pgfpathlineto{\pgfqpoint{1.412954in}{2.302647in}}%
\pgfpathlineto{\pgfqpoint{1.413737in}{2.303128in}}%
\pgfpathlineto{\pgfqpoint{1.414076in}{2.302930in}}%
\pgfpathlineto{\pgfqpoint{1.414097in}{2.302769in}}%
\pgfpathlineto{\pgfqpoint{1.414735in}{2.303229in}}%
\pgfpathlineto{\pgfqpoint{1.415184in}{2.303020in}}%
\pgfpathlineto{\pgfqpoint{1.415380in}{2.303260in}}%
\pgfpathlineto{\pgfqpoint{1.416155in}{2.302333in}}%
\pgfpathlineto{\pgfqpoint{1.416276in}{2.302382in}}%
\pgfpathlineto{\pgfqpoint{1.417391in}{2.302389in}}%
\pgfpathlineto{\pgfqpoint{1.417080in}{2.301855in}}%
\pgfpathlineto{\pgfqpoint{1.417394in}{2.302363in}}%
\pgfpathlineto{\pgfqpoint{1.418014in}{2.302081in}}%
\pgfpathlineto{\pgfqpoint{1.417429in}{2.302429in}}%
\pgfpathlineto{\pgfqpoint{1.418456in}{2.302682in}}%
\pgfpathlineto{\pgfqpoint{1.419550in}{2.303445in}}%
\pgfpathlineto{\pgfqpoint{1.418731in}{2.302160in}}%
\pgfpathlineto{\pgfqpoint{1.419594in}{2.303401in}}%
\pgfpathlineto{\pgfqpoint{1.420536in}{2.303348in}}%
\pgfpathlineto{\pgfqpoint{1.419637in}{2.303444in}}%
\pgfpathlineto{\pgfqpoint{1.420706in}{2.303406in}}%
\pgfpathlineto{\pgfqpoint{1.421669in}{2.303451in}}%
\pgfpathlineto{\pgfqpoint{1.420816in}{2.303354in}}%
\pgfpathlineto{\pgfqpoint{1.421820in}{2.303435in}}%
\pgfpathlineto{\pgfqpoint{1.422604in}{2.303321in}}%
\pgfpathlineto{\pgfqpoint{1.421861in}{2.303451in}}%
\pgfpathlineto{\pgfqpoint{1.422933in}{2.303411in}}%
\pgfpathlineto{\pgfqpoint{1.425234in}{2.303343in}}%
\pgfpathlineto{\pgfqpoint{1.422935in}{2.303444in}}%
\pgfpathlineto{\pgfqpoint{1.425245in}{2.303417in}}%
\pgfpathlineto{\pgfqpoint{1.426314in}{2.303446in}}%
\pgfpathlineto{\pgfqpoint{1.425590in}{2.303365in}}%
\pgfpathlineto{\pgfqpoint{1.426355in}{2.303426in}}%
\pgfpathlineto{\pgfqpoint{1.427274in}{2.303363in}}%
\pgfpathlineto{\pgfqpoint{1.426394in}{2.303454in}}%
\pgfpathlineto{\pgfqpoint{1.427468in}{2.303403in}}%
\pgfpathlineto{\pgfqpoint{1.428644in}{2.303446in}}%
\pgfpathlineto{\pgfqpoint{1.427496in}{2.303389in}}%
\pgfpathlineto{\pgfqpoint{1.428682in}{2.303416in}}%
\pgfpathlineto{\pgfqpoint{1.429677in}{2.303380in}}%
\pgfpathlineto{\pgfqpoint{1.428823in}{2.303450in}}%
\pgfpathlineto{\pgfqpoint{1.429796in}{2.303405in}}%
\pgfpathlineto{\pgfqpoint{1.430849in}{2.303455in}}%
\pgfpathlineto{\pgfqpoint{1.429808in}{2.303382in}}%
\pgfpathlineto{\pgfqpoint{1.430908in}{2.303431in}}%
\pgfpathlineto{\pgfqpoint{1.431968in}{2.303349in}}%
\pgfpathlineto{\pgfqpoint{1.431080in}{2.303446in}}%
\pgfpathlineto{\pgfqpoint{1.432020in}{2.303379in}}%
\pgfpathlineto{\pgfqpoint{1.433114in}{2.303378in}}%
\pgfpathlineto{\pgfqpoint{1.432034in}{2.303442in}}%
\pgfpathlineto{\pgfqpoint{1.433134in}{2.303401in}}%
\pgfpathlineto{\pgfqpoint{1.434254in}{2.303374in}}%
\pgfpathlineto{\pgfqpoint{1.433174in}{2.303446in}}%
\pgfpathlineto{\pgfqpoint{1.434266in}{2.303405in}}%
\pgfpathlineto{\pgfqpoint{1.435326in}{2.303438in}}%
\pgfpathlineto{\pgfqpoint{1.434298in}{2.303365in}}%
\pgfpathlineto{\pgfqpoint{1.435378in}{2.303411in}}%
\pgfpathlineto{\pgfqpoint{1.436580in}{2.303384in}}%
\pgfpathlineto{\pgfqpoint{1.435396in}{2.303443in}}%
\pgfpathlineto{\pgfqpoint{1.436593in}{2.303415in}}%
\pgfpathlineto{\pgfqpoint{1.437286in}{2.303456in}}%
\pgfpathlineto{\pgfqpoint{1.437686in}{2.303315in}}%
\pgfpathlineto{\pgfqpoint{1.437704in}{2.303407in}}%
\pgfpathlineto{\pgfqpoint{1.440048in}{2.301875in}}%
\pgfpathlineto{\pgfqpoint{1.437707in}{2.303441in}}%
\pgfpathlineto{\pgfqpoint{1.440082in}{2.302003in}}%
\pgfpathlineto{\pgfqpoint{1.441177in}{2.302482in}}%
\pgfpathlineto{\pgfqpoint{1.440113in}{2.301882in}}%
\pgfpathlineto{\pgfqpoint{1.441211in}{2.302424in}}%
\pgfpathlineto{\pgfqpoint{1.442211in}{2.301974in}}%
\pgfpathlineto{\pgfqpoint{1.441323in}{2.302639in}}%
\pgfpathlineto{\pgfqpoint{1.442345in}{2.302169in}}%
\pgfpathlineto{\pgfqpoint{1.442869in}{2.302560in}}%
\pgfpathlineto{\pgfqpoint{1.442448in}{2.302106in}}%
\pgfpathlineto{\pgfqpoint{1.443459in}{2.302181in}}%
\pgfpathlineto{\pgfqpoint{1.444366in}{2.302040in}}%
\pgfpathlineto{\pgfqpoint{1.443822in}{2.302449in}}%
\pgfpathlineto{\pgfqpoint{1.444586in}{2.302150in}}%
\pgfpathlineto{\pgfqpoint{1.445753in}{2.302632in}}%
\pgfpathlineto{\pgfqpoint{1.444912in}{2.301949in}}%
\pgfpathlineto{\pgfqpoint{1.445763in}{2.302559in}}%
\pgfpathlineto{\pgfqpoint{1.446889in}{2.302425in}}%
\pgfpathlineto{\pgfqpoint{1.446606in}{2.302785in}}%
\pgfpathlineto{\pgfqpoint{1.446897in}{2.302444in}}%
\pgfpathlineto{\pgfqpoint{1.448026in}{2.302788in}}%
\pgfpathlineto{\pgfqpoint{1.447646in}{2.302061in}}%
\pgfpathlineto{\pgfqpoint{1.448036in}{2.302770in}}%
\pgfpathlineto{\pgfqpoint{1.448039in}{2.302759in}}%
\pgfpathlineto{\pgfqpoint{1.449034in}{2.303783in}}%
\pgfpathlineto{\pgfqpoint{1.449042in}{2.303746in}}%
\pgfpathlineto{\pgfqpoint{1.449227in}{2.304200in}}%
\pgfpathlineto{\pgfqpoint{1.450126in}{2.303675in}}%
\pgfpathlineto{\pgfqpoint{1.450152in}{2.303703in}}%
\pgfpathlineto{\pgfqpoint{1.451227in}{2.303376in}}%
\pgfpathlineto{\pgfqpoint{1.450490in}{2.304015in}}%
\pgfpathlineto{\pgfqpoint{1.451288in}{2.303498in}}%
\pgfpathlineto{\pgfqpoint{1.451777in}{2.303634in}}%
\pgfpathlineto{\pgfqpoint{1.452284in}{2.303203in}}%
\pgfpathlineto{\pgfqpoint{1.452397in}{2.303441in}}%
\pgfpathlineto{\pgfqpoint{1.452833in}{2.303422in}}%
\pgfpathlineto{\pgfqpoint{1.453489in}{2.303973in}}%
\pgfpathlineto{\pgfqpoint{1.453970in}{2.304353in}}%
\pgfpathlineto{\pgfqpoint{1.453504in}{2.303913in}}%
\pgfpathlineto{\pgfqpoint{1.454683in}{2.303951in}}%
\pgfpathlineto{\pgfqpoint{1.455700in}{2.303340in}}%
\pgfpathlineto{\pgfqpoint{1.455814in}{2.303403in}}%
\pgfpathlineto{\pgfqpoint{1.455885in}{2.303557in}}%
\pgfpathlineto{\pgfqpoint{1.456361in}{2.302943in}}%
\pgfpathlineto{\pgfqpoint{1.456900in}{2.302957in}}%
\pgfpathlineto{\pgfqpoint{1.457462in}{2.302541in}}%
\pgfpathlineto{\pgfqpoint{1.456918in}{2.302990in}}%
\pgfpathlineto{\pgfqpoint{1.458028in}{2.302905in}}%
\pgfpathlineto{\pgfqpoint{1.459797in}{2.303332in}}%
\pgfpathlineto{\pgfqpoint{1.458039in}{2.302888in}}%
\pgfpathlineto{\pgfqpoint{1.459839in}{2.303218in}}%
\pgfpathlineto{\pgfqpoint{1.460558in}{2.303096in}}%
\pgfpathlineto{\pgfqpoint{1.460213in}{2.303482in}}%
\pgfpathlineto{\pgfqpoint{1.460946in}{2.303310in}}%
\pgfpathlineto{\pgfqpoint{1.462059in}{2.303485in}}%
\pgfpathlineto{\pgfqpoint{1.461418in}{2.302945in}}%
\pgfpathlineto{\pgfqpoint{1.462062in}{2.303474in}}%
\pgfpathlineto{\pgfqpoint{1.462173in}{2.303589in}}%
\pgfpathlineto{\pgfqpoint{1.463125in}{2.302713in}}%
\pgfpathlineto{\pgfqpoint{1.463143in}{2.302786in}}%
\pgfpathlineto{\pgfqpoint{1.464308in}{2.302442in}}%
\pgfpathlineto{\pgfqpoint{1.463166in}{2.302832in}}%
\pgfpathlineto{\pgfqpoint{1.464312in}{2.302482in}}%
\pgfpathlineto{\pgfqpoint{1.464918in}{2.302249in}}%
\pgfpathlineto{\pgfqpoint{1.465454in}{2.302827in}}%
\pgfpathlineto{\pgfqpoint{1.465479in}{2.302763in}}%
\pgfpathlineto{\pgfqpoint{1.466145in}{2.303328in}}%
\pgfpathlineto{\pgfqpoint{1.466561in}{2.302982in}}%
\pgfpathlineto{\pgfqpoint{1.467493in}{2.302829in}}%
\pgfpathlineto{\pgfqpoint{1.467207in}{2.303359in}}%
\pgfpathlineto{\pgfqpoint{1.467669in}{2.303066in}}%
\pgfpathlineto{\pgfqpoint{1.467783in}{2.303231in}}%
\pgfpathlineto{\pgfqpoint{1.468289in}{2.302408in}}%
\pgfpathlineto{\pgfqpoint{1.468765in}{2.302847in}}%
\pgfpathlineto{\pgfqpoint{1.468795in}{2.302681in}}%
\pgfpathlineto{\pgfqpoint{1.469565in}{2.303165in}}%
\pgfpathlineto{\pgfqpoint{1.469872in}{2.302977in}}%
\pgfpathlineto{\pgfqpoint{1.471113in}{2.303780in}}%
\pgfpathlineto{\pgfqpoint{1.470075in}{2.302843in}}%
\pgfpathlineto{\pgfqpoint{1.471139in}{2.303664in}}%
\pgfpathlineto{\pgfqpoint{1.472155in}{2.303094in}}%
\pgfpathlineto{\pgfqpoint{1.471208in}{2.303745in}}%
\pgfpathlineto{\pgfqpoint{1.472262in}{2.303167in}}%
\pgfpathlineto{\pgfqpoint{1.473530in}{2.303331in}}%
\pgfpathlineto{\pgfqpoint{1.472264in}{2.303144in}}%
\pgfpathlineto{\pgfqpoint{1.473581in}{2.303179in}}%
\pgfpathlineto{\pgfqpoint{1.474251in}{2.302886in}}%
\pgfpathlineto{\pgfqpoint{1.474509in}{2.303454in}}%
\pgfpathlineto{\pgfqpoint{1.474691in}{2.303229in}}%
\pgfpathlineto{\pgfqpoint{1.475412in}{2.303501in}}%
\pgfpathlineto{\pgfqpoint{1.474796in}{2.303150in}}%
\pgfpathlineto{\pgfqpoint{1.475842in}{2.303312in}}%
\pgfpathlineto{\pgfqpoint{1.477476in}{2.302624in}}%
\pgfpathlineto{\pgfqpoint{1.475852in}{2.303335in}}%
\pgfpathlineto{\pgfqpoint{1.477517in}{2.302755in}}%
\pgfpathlineto{\pgfqpoint{1.478635in}{2.303131in}}%
\pgfpathlineto{\pgfqpoint{1.477612in}{2.302739in}}%
\pgfpathlineto{\pgfqpoint{1.478657in}{2.303074in}}%
\pgfpathlineto{\pgfqpoint{1.479418in}{2.302739in}}%
\pgfpathlineto{\pgfqpoint{1.478814in}{2.303160in}}%
\pgfpathlineto{\pgfqpoint{1.479767in}{2.303123in}}%
\pgfpathlineto{\pgfqpoint{1.480758in}{2.303853in}}%
\pgfpathlineto{\pgfqpoint{1.479787in}{2.303085in}}%
\pgfpathlineto{\pgfqpoint{1.481108in}{2.303586in}}%
\pgfpathlineto{\pgfqpoint{1.481290in}{2.303402in}}%
\pgfpathlineto{\pgfqpoint{1.481644in}{2.303751in}}%
\pgfpathlineto{\pgfqpoint{1.482219in}{2.303612in}}%
\pgfpathlineto{\pgfqpoint{1.482682in}{2.303793in}}%
\pgfpathlineto{\pgfqpoint{1.483309in}{2.303232in}}%
\pgfpathlineto{\pgfqpoint{1.484425in}{2.303352in}}%
\pgfpathlineto{\pgfqpoint{1.483747in}{2.302871in}}%
\pgfpathlineto{\pgfqpoint{1.484514in}{2.303220in}}%
\pgfpathlineto{\pgfqpoint{1.484831in}{2.303146in}}%
\pgfpathlineto{\pgfqpoint{1.485473in}{2.303736in}}%
\pgfpathlineto{\pgfqpoint{1.485605in}{2.303432in}}%
\pgfpathlineto{\pgfqpoint{1.486713in}{2.303905in}}%
\pgfpathlineto{\pgfqpoint{1.485994in}{2.303408in}}%
\pgfpathlineto{\pgfqpoint{1.486728in}{2.303827in}}%
\pgfpathlineto{\pgfqpoint{1.487753in}{2.303589in}}%
\pgfpathlineto{\pgfqpoint{1.486779in}{2.303997in}}%
\pgfpathlineto{\pgfqpoint{1.487842in}{2.303764in}}%
\pgfpathlineto{\pgfqpoint{1.488755in}{2.304222in}}%
\pgfpathlineto{\pgfqpoint{1.488971in}{2.303924in}}%
\pgfpathlineto{\pgfqpoint{1.490042in}{2.303287in}}%
\pgfpathlineto{\pgfqpoint{1.489569in}{2.304390in}}%
\pgfpathlineto{\pgfqpoint{1.490110in}{2.303485in}}%
\pgfpathlineto{\pgfqpoint{1.490721in}{2.303453in}}%
\pgfpathlineto{\pgfqpoint{1.490348in}{2.303131in}}%
\pgfpathlineto{\pgfqpoint{1.491187in}{2.303223in}}%
\pgfpathlineto{\pgfqpoint{1.492083in}{2.302894in}}%
\pgfpathlineto{\pgfqpoint{1.491482in}{2.303460in}}%
\pgfpathlineto{\pgfqpoint{1.492304in}{2.303125in}}%
\pgfpathlineto{\pgfqpoint{1.492401in}{2.302979in}}%
\pgfpathlineto{\pgfqpoint{1.492836in}{2.303469in}}%
\pgfpathlineto{\pgfqpoint{1.493397in}{2.303312in}}%
\pgfpathlineto{\pgfqpoint{1.494232in}{2.303458in}}%
\pgfpathlineto{\pgfqpoint{1.493638in}{2.303030in}}%
\pgfpathlineto{\pgfqpoint{1.494511in}{2.303402in}}%
\pgfpathlineto{\pgfqpoint{1.495605in}{2.303444in}}%
\pgfpathlineto{\pgfqpoint{1.494595in}{2.303346in}}%
\pgfpathlineto{\pgfqpoint{1.495639in}{2.303391in}}%
\pgfpathlineto{\pgfqpoint{1.496676in}{2.303367in}}%
\pgfpathlineto{\pgfqpoint{1.495667in}{2.303457in}}%
\pgfpathlineto{\pgfqpoint{1.496756in}{2.303391in}}%
\pgfpathlineto{\pgfqpoint{1.497791in}{2.303439in}}%
\pgfpathlineto{\pgfqpoint{1.496764in}{2.303377in}}%
\pgfpathlineto{\pgfqpoint{1.497871in}{2.303435in}}%
\pgfpathlineto{\pgfqpoint{1.498960in}{2.303365in}}%
\pgfpathlineto{\pgfqpoint{1.497921in}{2.303445in}}%
\pgfpathlineto{\pgfqpoint{1.498983in}{2.303398in}}%
\pgfpathlineto{\pgfqpoint{1.500058in}{2.303370in}}%
\pgfpathlineto{\pgfqpoint{1.499005in}{2.303436in}}%
\pgfpathlineto{\pgfqpoint{1.500114in}{2.303383in}}%
\pgfpathlineto{\pgfqpoint{1.501204in}{2.303452in}}%
\pgfpathlineto{\pgfqpoint{1.500429in}{2.303335in}}%
\pgfpathlineto{\pgfqpoint{1.501225in}{2.303418in}}%
\pgfpathlineto{\pgfqpoint{1.502352in}{2.303440in}}%
\pgfpathlineto{\pgfqpoint{1.501238in}{2.303378in}}%
\pgfpathlineto{\pgfqpoint{1.502355in}{2.303425in}}%
\pgfpathlineto{\pgfqpoint{1.503281in}{2.303362in}}%
\pgfpathlineto{\pgfqpoint{1.502450in}{2.303444in}}%
\pgfpathlineto{\pgfqpoint{1.503467in}{2.303400in}}%
\pgfpathlineto{\pgfqpoint{1.504621in}{2.303434in}}%
\pgfpathlineto{\pgfqpoint{1.503508in}{2.303370in}}%
\pgfpathlineto{\pgfqpoint{1.504631in}{2.303393in}}%
\pgfpathlineto{\pgfqpoint{1.505707in}{2.303361in}}%
\pgfpathlineto{\pgfqpoint{1.504682in}{2.303444in}}%
\pgfpathlineto{\pgfqpoint{1.505741in}{2.303415in}}%
\pgfpathlineto{\pgfqpoint{1.506353in}{2.303461in}}%
\pgfpathlineto{\pgfqpoint{1.505794in}{2.303369in}}%
\pgfpathlineto{\pgfqpoint{1.506852in}{2.303397in}}%
\pgfpathlineto{\pgfqpoint{1.507958in}{2.303446in}}%
\pgfpathlineto{\pgfqpoint{1.506954in}{2.303367in}}%
\pgfpathlineto{\pgfqpoint{1.507966in}{2.303433in}}%
\pgfpathlineto{\pgfqpoint{1.508978in}{2.303372in}}%
\pgfpathlineto{\pgfqpoint{1.508232in}{2.303453in}}%
\pgfpathlineto{\pgfqpoint{1.509079in}{2.303410in}}%
\pgfpathlineto{\pgfqpoint{1.510151in}{2.303443in}}%
\pgfpathlineto{\pgfqpoint{1.509309in}{2.303327in}}%
\pgfpathlineto{\pgfqpoint{1.510194in}{2.303425in}}%
\pgfpathlineto{\pgfqpoint{1.511150in}{2.303360in}}%
\pgfpathlineto{\pgfqpoint{1.510236in}{2.303442in}}%
\pgfpathlineto{\pgfqpoint{1.511307in}{2.303387in}}%
\pgfpathlineto{\pgfqpoint{1.512411in}{2.303452in}}%
\pgfpathlineto{\pgfqpoint{1.511847in}{2.303354in}}%
\pgfpathlineto{\pgfqpoint{1.512418in}{2.303391in}}%
\pgfpathlineto{\pgfqpoint{1.513156in}{2.303240in}}%
\pgfpathlineto{\pgfqpoint{1.512678in}{2.303530in}}%
\pgfpathlineto{\pgfqpoint{1.513523in}{2.303517in}}%
\pgfpathlineto{\pgfqpoint{1.513528in}{2.303576in}}%
\pgfpathlineto{\pgfqpoint{1.514538in}{2.302793in}}%
\pgfpathlineto{\pgfqpoint{1.514582in}{2.302836in}}%
\pgfpathlineto{\pgfqpoint{1.515429in}{2.302230in}}%
\pgfpathlineto{\pgfqpoint{1.514624in}{2.302878in}}%
\pgfpathlineto{\pgfqpoint{1.515716in}{2.302307in}}%
\pgfpathlineto{\pgfqpoint{1.517115in}{2.302080in}}%
\pgfpathlineto{\pgfqpoint{1.515796in}{2.302118in}}%
\pgfpathlineto{\pgfqpoint{1.517128in}{2.301971in}}%
\pgfpathlineto{\pgfqpoint{1.518140in}{2.301273in}}%
\pgfpathlineto{\pgfqpoint{1.517397in}{2.302076in}}%
\pgfpathlineto{\pgfqpoint{1.518273in}{2.301434in}}%
\pgfpathlineto{\pgfqpoint{1.518342in}{2.301389in}}%
\pgfpathlineto{\pgfqpoint{1.519178in}{2.302084in}}%
\pgfpathlineto{\pgfqpoint{1.519319in}{2.301937in}}%
\pgfpathlineto{\pgfqpoint{1.520110in}{2.302033in}}%
\pgfpathlineto{\pgfqpoint{1.519484in}{2.301640in}}%
\pgfpathlineto{\pgfqpoint{1.520434in}{2.301956in}}%
\pgfpathlineto{\pgfqpoint{1.521668in}{2.302054in}}%
\pgfpathlineto{\pgfqpoint{1.520764in}{2.301493in}}%
\pgfpathlineto{\pgfqpoint{1.521690in}{2.301996in}}%
\pgfpathlineto{\pgfqpoint{1.521720in}{2.301848in}}%
\pgfpathlineto{\pgfqpoint{1.522180in}{2.302365in}}%
\pgfpathlineto{\pgfqpoint{1.522787in}{2.302179in}}%
\pgfpathlineto{\pgfqpoint{1.523454in}{2.302356in}}%
\pgfpathlineto{\pgfqpoint{1.523117in}{2.301987in}}%
\pgfpathlineto{\pgfqpoint{1.523906in}{2.302256in}}%
\pgfpathlineto{\pgfqpoint{1.524545in}{2.302815in}}%
\pgfpathlineto{\pgfqpoint{1.523910in}{2.302228in}}%
\pgfpathlineto{\pgfqpoint{1.525105in}{2.302430in}}%
\pgfpathlineto{\pgfqpoint{1.526418in}{2.302053in}}%
\pgfpathlineto{\pgfqpoint{1.525126in}{2.302508in}}%
\pgfpathlineto{\pgfqpoint{1.526425in}{2.302089in}}%
\pgfpathlineto{\pgfqpoint{1.527064in}{2.302767in}}%
\pgfpathlineto{\pgfqpoint{1.526468in}{2.302009in}}%
\pgfpathlineto{\pgfqpoint{1.527596in}{2.302420in}}%
\pgfpathlineto{\pgfqpoint{1.528503in}{2.302043in}}%
\pgfpathlineto{\pgfqpoint{1.527603in}{2.302438in}}%
\pgfpathlineto{\pgfqpoint{1.528732in}{2.302151in}}%
\pgfpathlineto{\pgfqpoint{1.528961in}{2.302427in}}%
\pgfpathlineto{\pgfqpoint{1.529347in}{2.302069in}}%
\pgfpathlineto{\pgfqpoint{1.529845in}{2.302210in}}%
\pgfpathlineto{\pgfqpoint{1.530615in}{2.302198in}}%
\pgfpathlineto{\pgfqpoint{1.530218in}{2.302603in}}%
\pgfpathlineto{\pgfqpoint{1.530922in}{2.302490in}}%
\pgfpathlineto{\pgfqpoint{1.531525in}{2.302654in}}%
\pgfpathlineto{\pgfqpoint{1.531992in}{2.302085in}}%
\pgfpathlineto{\pgfqpoint{1.531995in}{2.302093in}}%
\pgfpathlineto{\pgfqpoint{1.532081in}{2.302028in}}%
\pgfpathlineto{\pgfqpoint{1.532947in}{2.302703in}}%
\pgfpathlineto{\pgfqpoint{1.533083in}{2.302524in}}%
\pgfpathlineto{\pgfqpoint{1.533247in}{2.302456in}}%
\pgfpathlineto{\pgfqpoint{1.533710in}{2.301755in}}%
\pgfpathlineto{\pgfqpoint{1.533915in}{2.301795in}}%
\pgfpathlineto{\pgfqpoint{1.535007in}{2.301217in}}%
\pgfpathlineto{\pgfqpoint{1.533996in}{2.301942in}}%
\pgfpathlineto{\pgfqpoint{1.535048in}{2.301317in}}%
\pgfpathlineto{\pgfqpoint{1.535885in}{2.301607in}}%
\pgfpathlineto{\pgfqpoint{1.535387in}{2.301210in}}%
\pgfpathlineto{\pgfqpoint{1.536176in}{2.301511in}}%
\pgfpathlineto{\pgfqpoint{1.537112in}{2.300713in}}%
\pgfpathlineto{\pgfqpoint{1.536268in}{2.301535in}}%
\pgfpathlineto{\pgfqpoint{1.537315in}{2.300792in}}%
\pgfpathlineto{\pgfqpoint{1.538075in}{2.300654in}}%
\pgfpathlineto{\pgfqpoint{1.538398in}{2.301012in}}%
\pgfpathlineto{\pgfqpoint{1.538414in}{2.300960in}}%
\pgfpathlineto{\pgfqpoint{1.539251in}{2.301067in}}%
\pgfpathlineto{\pgfqpoint{1.538731in}{2.300713in}}%
\pgfpathlineto{\pgfqpoint{1.539528in}{2.300878in}}%
\pgfpathlineto{\pgfqpoint{1.540335in}{2.300832in}}%
\pgfpathlineto{\pgfqpoint{1.539921in}{2.301383in}}%
\pgfpathlineto{\pgfqpoint{1.540624in}{2.301117in}}%
\pgfpathlineto{\pgfqpoint{1.540669in}{2.300905in}}%
\pgfpathlineto{\pgfqpoint{1.541227in}{2.301846in}}%
\pgfpathlineto{\pgfqpoint{1.541704in}{2.301418in}}%
\pgfpathlineto{\pgfqpoint{1.542349in}{2.301583in}}%
\pgfpathlineto{\pgfqpoint{1.542806in}{2.301243in}}%
\pgfpathlineto{\pgfqpoint{1.542809in}{2.301260in}}%
\pgfpathlineto{\pgfqpoint{1.543328in}{2.300960in}}%
\pgfpathlineto{\pgfqpoint{1.543028in}{2.301440in}}%
\pgfpathlineto{\pgfqpoint{1.543924in}{2.301214in}}%
\pgfpathlineto{\pgfqpoint{1.543975in}{2.301369in}}%
\pgfpathlineto{\pgfqpoint{1.544993in}{2.300830in}}%
\pgfpathlineto{\pgfqpoint{1.545026in}{2.300922in}}%
\pgfpathlineto{\pgfqpoint{1.546100in}{2.300415in}}%
\pgfpathlineto{\pgfqpoint{1.545203in}{2.300999in}}%
\pgfpathlineto{\pgfqpoint{1.546178in}{2.300527in}}%
\pgfpathlineto{\pgfqpoint{1.547025in}{2.301010in}}%
\pgfpathlineto{\pgfqpoint{1.546308in}{2.300452in}}%
\pgfpathlineto{\pgfqpoint{1.547297in}{2.300530in}}%
\pgfpathlineto{\pgfqpoint{1.547362in}{2.300427in}}%
\pgfpathlineto{\pgfqpoint{1.548273in}{2.301028in}}%
\pgfpathlineto{\pgfqpoint{1.548386in}{2.300934in}}%
\pgfpathlineto{\pgfqpoint{1.548454in}{2.301130in}}%
\pgfpathlineto{\pgfqpoint{1.549585in}{2.300469in}}%
\pgfpathlineto{\pgfqpoint{1.549703in}{2.300568in}}%
\pgfpathlineto{\pgfqpoint{1.550285in}{2.299708in}}%
\pgfpathlineto{\pgfqpoint{1.550668in}{2.300000in}}%
\pgfpathlineto{\pgfqpoint{1.550720in}{2.300174in}}%
\pgfpathlineto{\pgfqpoint{1.551626in}{2.299298in}}%
\pgfpathlineto{\pgfqpoint{1.551713in}{2.299394in}}%
\pgfpathlineto{\pgfqpoint{1.552945in}{2.298862in}}%
\pgfpathlineto{\pgfqpoint{1.552975in}{2.298970in}}%
\pgfpathlineto{\pgfqpoint{1.554131in}{2.299519in}}%
\pgfpathlineto{\pgfqpoint{1.553370in}{2.298648in}}%
\pgfpathlineto{\pgfqpoint{1.554136in}{2.299502in}}%
\pgfpathlineto{\pgfqpoint{1.554295in}{2.299624in}}%
\pgfpathlineto{\pgfqpoint{1.555272in}{2.298890in}}%
\pgfpathlineto{\pgfqpoint{1.556434in}{2.299310in}}%
\pgfpathlineto{\pgfqpoint{1.555454in}{2.298718in}}%
\pgfpathlineto{\pgfqpoint{1.556469in}{2.299288in}}%
\pgfpathlineto{\pgfqpoint{1.557589in}{2.299123in}}%
\pgfpathlineto{\pgfqpoint{1.557212in}{2.299592in}}%
\pgfpathlineto{\pgfqpoint{1.557591in}{2.299129in}}%
\pgfpathlineto{\pgfqpoint{1.559068in}{2.299641in}}%
\pgfpathlineto{\pgfqpoint{1.557813in}{2.298900in}}%
\pgfpathlineto{\pgfqpoint{1.559075in}{2.299583in}}%
\pgfpathlineto{\pgfqpoint{1.559110in}{2.299440in}}%
\pgfpathlineto{\pgfqpoint{1.559570in}{2.299955in}}%
\pgfpathlineto{\pgfqpoint{1.560180in}{2.299715in}}%
\pgfpathlineto{\pgfqpoint{1.560187in}{2.299739in}}%
\pgfpathlineto{\pgfqpoint{1.560885in}{2.298989in}}%
\pgfpathlineto{\pgfqpoint{1.561266in}{2.299353in}}%
\pgfpathlineto{\pgfqpoint{1.562369in}{2.298923in}}%
\pgfpathlineto{\pgfqpoint{1.561538in}{2.299552in}}%
\pgfpathlineto{\pgfqpoint{1.562447in}{2.299221in}}%
\pgfpathlineto{\pgfqpoint{1.563593in}{2.299592in}}%
\pgfpathlineto{\pgfqpoint{1.562494in}{2.299171in}}%
\pgfpathlineto{\pgfqpoint{1.563598in}{2.299567in}}%
\pgfpathlineto{\pgfqpoint{1.564386in}{2.298930in}}%
\pgfpathlineto{\pgfqpoint{1.563612in}{2.299592in}}%
\pgfpathlineto{\pgfqpoint{1.564771in}{2.299077in}}%
\pgfpathlineto{\pgfqpoint{1.564780in}{2.299157in}}%
\pgfpathlineto{\pgfqpoint{1.565464in}{2.298457in}}%
\pgfpathlineto{\pgfqpoint{1.565872in}{2.298745in}}%
\pgfpathlineto{\pgfqpoint{1.566029in}{2.298431in}}%
\pgfpathlineto{\pgfqpoint{1.566477in}{2.299019in}}%
\pgfpathlineto{\pgfqpoint{1.566982in}{2.298763in}}%
\pgfpathlineto{\pgfqpoint{1.567035in}{2.298920in}}%
\pgfpathlineto{\pgfqpoint{1.567748in}{2.298500in}}%
\pgfpathlineto{\pgfqpoint{1.568089in}{2.298637in}}%
\pgfpathlineto{\pgfqpoint{1.569658in}{2.303443in}}%
\pgfpathlineto{\pgfqpoint{1.568134in}{2.298578in}}%
\pgfpathlineto{\pgfqpoint{1.569704in}{2.303405in}}%
\pgfpathlineto{\pgfqpoint{1.570600in}{2.303349in}}%
\pgfpathlineto{\pgfqpoint{1.569937in}{2.303456in}}%
\pgfpathlineto{\pgfqpoint{1.570816in}{2.303418in}}%
\pgfpathlineto{\pgfqpoint{1.571907in}{2.303443in}}%
\pgfpathlineto{\pgfqpoint{1.570918in}{2.303369in}}%
\pgfpathlineto{\pgfqpoint{1.571929in}{2.303425in}}%
\pgfpathlineto{\pgfqpoint{1.572845in}{2.303347in}}%
\pgfpathlineto{\pgfqpoint{1.572022in}{2.303443in}}%
\pgfpathlineto{\pgfqpoint{1.573042in}{2.303402in}}%
\pgfpathlineto{\pgfqpoint{1.574147in}{2.303447in}}%
\pgfpathlineto{\pgfqpoint{1.573078in}{2.303358in}}%
\pgfpathlineto{\pgfqpoint{1.574171in}{2.303415in}}%
\pgfpathlineto{\pgfqpoint{1.575451in}{2.303364in}}%
\pgfpathlineto{\pgfqpoint{1.574172in}{2.303433in}}%
\pgfpathlineto{\pgfqpoint{1.575503in}{2.303430in}}%
\pgfpathlineto{\pgfqpoint{1.577981in}{2.303435in}}%
\pgfpathlineto{\pgfqpoint{1.575505in}{2.303398in}}%
\pgfpathlineto{\pgfqpoint{1.577984in}{2.303395in}}%
\pgfpathlineto{\pgfqpoint{1.578987in}{2.303366in}}%
\pgfpathlineto{\pgfqpoint{1.578040in}{2.303444in}}%
\pgfpathlineto{\pgfqpoint{1.579095in}{2.303420in}}%
\pgfpathlineto{\pgfqpoint{1.580181in}{2.303366in}}%
\pgfpathlineto{\pgfqpoint{1.579097in}{2.303433in}}%
\pgfpathlineto{\pgfqpoint{1.580221in}{2.303413in}}%
\pgfpathlineto{\pgfqpoint{1.581250in}{2.303447in}}%
\pgfpathlineto{\pgfqpoint{1.580536in}{2.303361in}}%
\pgfpathlineto{\pgfqpoint{1.581336in}{2.303427in}}%
\pgfpathlineto{\pgfqpoint{1.582429in}{2.303363in}}%
\pgfpathlineto{\pgfqpoint{1.581338in}{2.303453in}}%
\pgfpathlineto{\pgfqpoint{1.582450in}{2.303419in}}%
\pgfpathlineto{\pgfqpoint{1.583494in}{2.303442in}}%
\pgfpathlineto{\pgfqpoint{1.582717in}{2.303362in}}%
\pgfpathlineto{\pgfqpoint{1.583561in}{2.303389in}}%
\pgfpathlineto{\pgfqpoint{1.584329in}{2.303345in}}%
\pgfpathlineto{\pgfqpoint{1.583656in}{2.303455in}}%
\pgfpathlineto{\pgfqpoint{1.584675in}{2.303407in}}%
\pgfpathlineto{\pgfqpoint{1.584733in}{2.303380in}}%
\pgfpathlineto{\pgfqpoint{1.585794in}{2.303437in}}%
\pgfpathlineto{\pgfqpoint{1.586856in}{2.303349in}}%
\pgfpathlineto{\pgfqpoint{1.585976in}{2.303454in}}%
\pgfpathlineto{\pgfqpoint{1.586906in}{2.303393in}}%
\pgfpathlineto{\pgfqpoint{1.588022in}{2.304098in}}%
\pgfpathlineto{\pgfqpoint{1.587063in}{2.303353in}}%
\pgfpathlineto{\pgfqpoint{1.588050in}{2.304069in}}%
\pgfpathlineto{\pgfqpoint{1.588210in}{2.303877in}}%
\pgfpathlineto{\pgfqpoint{1.589062in}{2.304826in}}%
\pgfpathlineto{\pgfqpoint{1.589137in}{2.304529in}}%
\pgfpathlineto{\pgfqpoint{1.590240in}{2.305233in}}%
\pgfpathlineto{\pgfqpoint{1.589216in}{2.304301in}}%
\pgfpathlineto{\pgfqpoint{1.590401in}{2.305133in}}%
\pgfpathlineto{\pgfqpoint{1.590544in}{2.305097in}}%
\pgfpathlineto{\pgfqpoint{1.590965in}{2.305620in}}%
\pgfpathlineto{\pgfqpoint{1.591502in}{2.305392in}}%
\pgfpathlineto{\pgfqpoint{1.592564in}{2.305483in}}%
\pgfpathlineto{\pgfqpoint{1.591647in}{2.305708in}}%
\pgfpathlineto{\pgfqpoint{1.592602in}{2.305579in}}%
\pgfpathlineto{\pgfqpoint{1.592606in}{2.305643in}}%
\pgfpathlineto{\pgfqpoint{1.593696in}{2.305020in}}%
\pgfpathlineto{\pgfqpoint{1.593697in}{2.305023in}}%
\pgfpathlineto{\pgfqpoint{1.600289in}{2.304254in}}%
\pgfpathlineto{\pgfqpoint{1.593707in}{2.305077in}}%
\pgfpathlineto{\pgfqpoint{1.600308in}{2.304410in}}%
\pgfpathlineto{\pgfqpoint{1.600800in}{2.304793in}}%
\pgfpathlineto{\pgfqpoint{1.600323in}{2.304396in}}%
\pgfpathlineto{\pgfqpoint{1.601422in}{2.304350in}}%
\pgfpathlineto{\pgfqpoint{1.602207in}{2.304255in}}%
\pgfpathlineto{\pgfqpoint{1.601964in}{2.304640in}}%
\pgfpathlineto{\pgfqpoint{1.602534in}{2.304369in}}%
\pgfpathlineto{\pgfqpoint{1.603653in}{2.304613in}}%
\pgfpathlineto{\pgfqpoint{1.602949in}{2.304070in}}%
\pgfpathlineto{\pgfqpoint{1.603660in}{2.304603in}}%
\pgfpathlineto{\pgfqpoint{1.604746in}{2.304143in}}%
\pgfpathlineto{\pgfqpoint{1.603725in}{2.304693in}}%
\pgfpathlineto{\pgfqpoint{1.604851in}{2.304348in}}%
\pgfpathlineto{\pgfqpoint{1.605701in}{2.304665in}}%
\pgfpathlineto{\pgfqpoint{1.605101in}{2.304139in}}%
\pgfpathlineto{\pgfqpoint{1.605990in}{2.304539in}}%
\pgfpathlineto{\pgfqpoint{1.606509in}{2.304291in}}%
\pgfpathlineto{\pgfqpoint{1.607032in}{2.304765in}}%
\pgfpathlineto{\pgfqpoint{1.607093in}{2.304693in}}%
\pgfpathlineto{\pgfqpoint{1.607312in}{2.305107in}}%
\pgfpathlineto{\pgfqpoint{1.607682in}{2.304586in}}%
\pgfpathlineto{\pgfqpoint{1.608209in}{2.304801in}}%
\pgfpathlineto{\pgfqpoint{1.609165in}{2.304336in}}%
\pgfpathlineto{\pgfqpoint{1.608443in}{2.304984in}}%
\pgfpathlineto{\pgfqpoint{1.609340in}{2.304498in}}%
\pgfpathlineto{\pgfqpoint{1.609665in}{2.304494in}}%
\pgfpathlineto{\pgfqpoint{1.610174in}{2.304833in}}%
\pgfpathlineto{\pgfqpoint{1.610424in}{2.304739in}}%
\pgfpathlineto{\pgfqpoint{1.610646in}{2.305079in}}%
\pgfpathlineto{\pgfqpoint{1.611508in}{2.304411in}}%
\pgfpathlineto{\pgfqpoint{1.611520in}{2.304473in}}%
\pgfpathlineto{\pgfqpoint{1.612651in}{2.304107in}}%
\pgfpathlineto{\pgfqpoint{1.611612in}{2.304532in}}%
\pgfpathlineto{\pgfqpoint{1.612656in}{2.304152in}}%
\pgfpathlineto{\pgfqpoint{1.613795in}{2.304455in}}%
\pgfpathlineto{\pgfqpoint{1.613236in}{2.303858in}}%
\pgfpathlineto{\pgfqpoint{1.613818in}{2.304444in}}%
\pgfpathlineto{\pgfqpoint{1.614503in}{2.304513in}}%
\pgfpathlineto{\pgfqpoint{1.614240in}{2.304784in}}%
\pgfpathlineto{\pgfqpoint{1.614869in}{2.304862in}}%
\pgfpathlineto{\pgfqpoint{1.614873in}{2.304898in}}%
\pgfpathlineto{\pgfqpoint{1.615371in}{2.304317in}}%
\pgfpathlineto{\pgfqpoint{1.615972in}{2.304704in}}%
\pgfpathlineto{\pgfqpoint{1.616668in}{2.304519in}}%
\pgfpathlineto{\pgfqpoint{1.615998in}{2.304801in}}%
\pgfpathlineto{\pgfqpoint{1.617079in}{2.304813in}}%
\pgfpathlineto{\pgfqpoint{1.617785in}{2.304941in}}%
\pgfpathlineto{\pgfqpoint{1.617425in}{2.304681in}}%
\pgfpathlineto{\pgfqpoint{1.618186in}{2.304729in}}%
\pgfpathlineto{\pgfqpoint{1.619326in}{2.304025in}}%
\pgfpathlineto{\pgfqpoint{1.618213in}{2.304792in}}%
\pgfpathlineto{\pgfqpoint{1.619328in}{2.304040in}}%
\pgfpathlineto{\pgfqpoint{1.619331in}{2.304103in}}%
\pgfpathlineto{\pgfqpoint{1.620388in}{2.303390in}}%
\pgfpathlineto{\pgfqpoint{1.620420in}{2.303461in}}%
\pgfpathlineto{\pgfqpoint{1.623691in}{2.302594in}}%
\pgfpathlineto{\pgfqpoint{1.620430in}{2.303536in}}%
\pgfpathlineto{\pgfqpoint{1.623743in}{2.302826in}}%
\pgfpathlineto{\pgfqpoint{1.624287in}{2.302887in}}%
\pgfpathlineto{\pgfqpoint{1.624787in}{2.302355in}}%
\pgfpathlineto{\pgfqpoint{1.624801in}{2.302371in}}%
\pgfpathlineto{\pgfqpoint{1.625207in}{2.302213in}}%
\pgfpathlineto{\pgfqpoint{1.625891in}{2.302706in}}%
\pgfpathlineto{\pgfqpoint{1.626863in}{2.302436in}}%
\pgfpathlineto{\pgfqpoint{1.626274in}{2.302999in}}%
\pgfpathlineto{\pgfqpoint{1.627150in}{2.302675in}}%
\pgfpathlineto{\pgfqpoint{1.628126in}{2.302969in}}%
\pgfpathlineto{\pgfqpoint{1.627280in}{2.302477in}}%
\pgfpathlineto{\pgfqpoint{1.628275in}{2.302876in}}%
\pgfpathlineto{\pgfqpoint{1.628311in}{2.302742in}}%
\pgfpathlineto{\pgfqpoint{1.629021in}{2.303728in}}%
\pgfpathlineto{\pgfqpoint{1.629325in}{2.303519in}}%
\pgfpathlineto{\pgfqpoint{1.630316in}{2.304265in}}%
\pgfpathlineto{\pgfqpoint{1.629340in}{2.303481in}}%
\pgfpathlineto{\pgfqpoint{1.630480in}{2.303983in}}%
\pgfpathlineto{\pgfqpoint{1.631479in}{2.304232in}}%
\pgfpathlineto{\pgfqpoint{1.630858in}{2.304417in}}%
\pgfpathlineto{\pgfqpoint{1.631525in}{2.304365in}}%
\pgfpathlineto{\pgfqpoint{1.631528in}{2.304385in}}%
\pgfpathlineto{\pgfqpoint{1.632070in}{2.303958in}}%
\pgfpathlineto{\pgfqpoint{1.632620in}{2.304156in}}%
\pgfpathlineto{\pgfqpoint{1.632932in}{2.304161in}}%
\pgfpathlineto{\pgfqpoint{1.633577in}{2.304842in}}%
\pgfpathlineto{\pgfqpoint{1.633605in}{2.304778in}}%
\pgfpathlineto{\pgfqpoint{1.633758in}{2.304971in}}%
\pgfpathlineto{\pgfqpoint{1.634700in}{2.304589in}}%
\pgfpathlineto{\pgfqpoint{1.634712in}{2.304688in}}%
\pgfpathlineto{\pgfqpoint{1.634955in}{2.304630in}}%
\pgfpathlineto{\pgfqpoint{1.635239in}{2.305208in}}%
\pgfpathlineto{\pgfqpoint{1.635777in}{2.305210in}}%
\pgfpathlineto{\pgfqpoint{1.636248in}{2.305463in}}%
\pgfpathlineto{\pgfqpoint{1.636774in}{2.304984in}}%
\pgfpathlineto{\pgfqpoint{1.636883in}{2.305100in}}%
\pgfpathlineto{\pgfqpoint{1.637680in}{2.305016in}}%
\pgfpathlineto{\pgfqpoint{1.637284in}{2.305533in}}%
\pgfpathlineto{\pgfqpoint{1.637995in}{2.305130in}}%
\pgfpathlineto{\pgfqpoint{1.639232in}{2.305457in}}%
\pgfpathlineto{\pgfqpoint{1.638033in}{2.305069in}}%
\pgfpathlineto{\pgfqpoint{1.639244in}{2.305347in}}%
\pgfpathlineto{\pgfqpoint{1.640353in}{2.304815in}}%
\pgfpathlineto{\pgfqpoint{1.639247in}{2.305350in}}%
\pgfpathlineto{\pgfqpoint{1.640379in}{2.304859in}}%
\pgfpathlineto{\pgfqpoint{1.640692in}{2.305284in}}%
\pgfpathlineto{\pgfqpoint{1.640391in}{2.304842in}}%
\pgfpathlineto{\pgfqpoint{1.641509in}{2.304929in}}%
\pgfpathlineto{\pgfqpoint{1.643118in}{2.304727in}}%
\pgfpathlineto{\pgfqpoint{1.643162in}{2.304836in}}%
\pgfpathlineto{\pgfqpoint{1.643321in}{2.305152in}}%
\pgfpathlineto{\pgfqpoint{1.643873in}{2.303368in}}%
\pgfpathlineto{\pgfqpoint{1.644204in}{2.303440in}}%
\pgfpathlineto{\pgfqpoint{1.645221in}{2.303351in}}%
\pgfpathlineto{\pgfqpoint{1.644437in}{2.303455in}}%
\pgfpathlineto{\pgfqpoint{1.645316in}{2.303422in}}%
\pgfpathlineto{\pgfqpoint{1.646344in}{2.303448in}}%
\pgfpathlineto{\pgfqpoint{1.645341in}{2.303357in}}%
\pgfpathlineto{\pgfqpoint{1.646427in}{2.303411in}}%
\pgfpathlineto{\pgfqpoint{1.648924in}{2.303429in}}%
\pgfpathlineto{\pgfqpoint{1.646429in}{2.303392in}}%
\pgfpathlineto{\pgfqpoint{1.648925in}{2.303392in}}%
\pgfpathlineto{\pgfqpoint{1.649761in}{2.303375in}}%
\pgfpathlineto{\pgfqpoint{1.648949in}{2.303440in}}%
\pgfpathlineto{\pgfqpoint{1.650036in}{2.303415in}}%
\pgfpathlineto{\pgfqpoint{1.651118in}{2.303353in}}%
\pgfpathlineto{\pgfqpoint{1.650107in}{2.303447in}}%
\pgfpathlineto{\pgfqpoint{1.651153in}{2.303415in}}%
\pgfpathlineto{\pgfqpoint{1.652156in}{2.303448in}}%
\pgfpathlineto{\pgfqpoint{1.651172in}{2.303357in}}%
\pgfpathlineto{\pgfqpoint{1.652265in}{2.303432in}}%
\pgfpathlineto{\pgfqpoint{1.653367in}{2.303373in}}%
\pgfpathlineto{\pgfqpoint{1.652331in}{2.303459in}}%
\pgfpathlineto{\pgfqpoint{1.653380in}{2.303420in}}%
\pgfpathlineto{\pgfqpoint{1.654355in}{2.303443in}}%
\pgfpathlineto{\pgfqpoint{1.653826in}{2.303342in}}%
\pgfpathlineto{\pgfqpoint{1.654490in}{2.303370in}}%
\pgfpathlineto{\pgfqpoint{1.655564in}{2.303441in}}%
\pgfpathlineto{\pgfqpoint{1.655604in}{2.303405in}}%
\pgfpathlineto{\pgfqpoint{1.656449in}{2.303342in}}%
\pgfpathlineto{\pgfqpoint{1.655623in}{2.303442in}}%
\pgfpathlineto{\pgfqpoint{1.656719in}{2.303418in}}%
\pgfpathlineto{\pgfqpoint{1.660220in}{2.303437in}}%
\pgfpathlineto{\pgfqpoint{1.656720in}{2.303411in}}%
\pgfpathlineto{\pgfqpoint{1.660225in}{2.303405in}}%
\pgfpathlineto{\pgfqpoint{1.661204in}{2.303359in}}%
\pgfpathlineto{\pgfqpoint{1.660237in}{2.303455in}}%
\pgfpathlineto{\pgfqpoint{1.661338in}{2.303417in}}%
\pgfpathlineto{\pgfqpoint{1.662437in}{2.303480in}}%
\pgfpathlineto{\pgfqpoint{1.661367in}{2.303355in}}%
\pgfpathlineto{\pgfqpoint{1.662449in}{2.303407in}}%
\pgfpathlineto{\pgfqpoint{1.662630in}{2.303296in}}%
\pgfpathlineto{\pgfqpoint{1.663472in}{2.304216in}}%
\pgfpathlineto{\pgfqpoint{1.663509in}{2.304156in}}%
\pgfpathlineto{\pgfqpoint{1.663537in}{2.304230in}}%
\pgfpathlineto{\pgfqpoint{1.664293in}{2.303635in}}%
\pgfpathlineto{\pgfqpoint{1.664599in}{2.303812in}}%
\pgfpathlineto{\pgfqpoint{1.665577in}{2.303709in}}%
\pgfpathlineto{\pgfqpoint{1.665052in}{2.304001in}}%
\pgfpathlineto{\pgfqpoint{1.665707in}{2.303878in}}%
\pgfpathlineto{\pgfqpoint{1.665891in}{2.303966in}}%
\pgfpathlineto{\pgfqpoint{1.666515in}{2.303394in}}%
\pgfpathlineto{\pgfqpoint{1.666746in}{2.303371in}}%
\pgfpathlineto{\pgfqpoint{1.666831in}{2.303156in}}%
\pgfpathlineto{\pgfqpoint{1.667562in}{2.303868in}}%
\pgfpathlineto{\pgfqpoint{1.667844in}{2.303797in}}%
\pgfpathlineto{\pgfqpoint{1.669022in}{2.303900in}}%
\pgfpathlineto{\pgfqpoint{1.667856in}{2.303772in}}%
\pgfpathlineto{\pgfqpoint{1.669085in}{2.303792in}}%
\pgfpathlineto{\pgfqpoint{1.671107in}{2.303360in}}%
\pgfpathlineto{\pgfqpoint{1.669088in}{2.303773in}}%
\pgfpathlineto{\pgfqpoint{1.671108in}{2.303356in}}%
\pgfpathlineto{\pgfqpoint{1.671456in}{2.303129in}}%
\pgfpathlineto{\pgfqpoint{1.672188in}{2.303596in}}%
\pgfpathlineto{\pgfqpoint{1.672215in}{2.303501in}}%
\pgfpathlineto{\pgfqpoint{1.673947in}{2.303719in}}%
\pgfpathlineto{\pgfqpoint{1.673961in}{2.303637in}}%
\pgfpathlineto{\pgfqpoint{1.675098in}{2.303310in}}%
\pgfpathlineto{\pgfqpoint{1.674058in}{2.303759in}}%
\pgfpathlineto{\pgfqpoint{1.675272in}{2.303530in}}%
\pgfpathlineto{\pgfqpoint{1.676241in}{2.303660in}}%
\pgfpathlineto{\pgfqpoint{1.675896in}{2.303021in}}%
\pgfpathlineto{\pgfqpoint{1.676387in}{2.303451in}}%
\pgfpathlineto{\pgfqpoint{1.676493in}{2.303394in}}%
\pgfpathlineto{\pgfqpoint{1.677467in}{2.303926in}}%
\pgfpathlineto{\pgfqpoint{1.677473in}{2.303906in}}%
\pgfpathlineto{\pgfqpoint{1.678979in}{2.303597in}}%
\pgfpathlineto{\pgfqpoint{1.677506in}{2.304017in}}%
\pgfpathlineto{\pgfqpoint{1.679113in}{2.303894in}}%
\pgfpathlineto{\pgfqpoint{1.679365in}{2.304085in}}%
\pgfpathlineto{\pgfqpoint{1.680096in}{2.303561in}}%
\pgfpathlineto{\pgfqpoint{1.680215in}{2.303758in}}%
\pgfpathlineto{\pgfqpoint{1.681376in}{2.302961in}}%
\pgfpathlineto{\pgfqpoint{1.680399in}{2.303914in}}%
\pgfpathlineto{\pgfqpoint{1.681383in}{2.302981in}}%
\pgfpathlineto{\pgfqpoint{1.681875in}{2.303123in}}%
\pgfpathlineto{\pgfqpoint{1.682224in}{2.302693in}}%
\pgfpathlineto{\pgfqpoint{1.682488in}{2.302809in}}%
\pgfpathlineto{\pgfqpoint{1.683501in}{2.302733in}}%
\pgfpathlineto{\pgfqpoint{1.682892in}{2.303408in}}%
\pgfpathlineto{\pgfqpoint{1.683595in}{2.302913in}}%
\pgfpathlineto{\pgfqpoint{1.683958in}{2.303272in}}%
\pgfpathlineto{\pgfqpoint{1.683618in}{2.302898in}}%
\pgfpathlineto{\pgfqpoint{1.684707in}{2.302881in}}%
\pgfpathlineto{\pgfqpoint{1.685423in}{2.302768in}}%
\pgfpathlineto{\pgfqpoint{1.684996in}{2.303177in}}%
\pgfpathlineto{\pgfqpoint{1.685811in}{2.303047in}}%
\pgfpathlineto{\pgfqpoint{1.686939in}{2.302897in}}%
\pgfpathlineto{\pgfqpoint{1.686496in}{2.303280in}}%
\pgfpathlineto{\pgfqpoint{1.686939in}{2.302910in}}%
\pgfpathlineto{\pgfqpoint{1.687912in}{2.302839in}}%
\pgfpathlineto{\pgfqpoint{1.687041in}{2.302647in}}%
\pgfpathlineto{\pgfqpoint{1.688031in}{2.302681in}}%
\pgfpathlineto{\pgfqpoint{1.688671in}{2.302179in}}%
\pgfpathlineto{\pgfqpoint{1.688308in}{2.302794in}}%
\pgfpathlineto{\pgfqpoint{1.689164in}{2.302216in}}%
\pgfpathlineto{\pgfqpoint{1.690064in}{2.302470in}}%
\pgfpathlineto{\pgfqpoint{1.689506in}{2.301938in}}%
\pgfpathlineto{\pgfqpoint{1.690292in}{2.302450in}}%
\pgfpathlineto{\pgfqpoint{1.690458in}{2.302255in}}%
\pgfpathlineto{\pgfqpoint{1.690857in}{2.302843in}}%
\pgfpathlineto{\pgfqpoint{1.691361in}{2.302831in}}%
\pgfpathlineto{\pgfqpoint{1.692407in}{2.303054in}}%
\pgfpathlineto{\pgfqpoint{1.691493in}{2.302661in}}%
\pgfpathlineto{\pgfqpoint{1.692486in}{2.303000in}}%
\pgfpathlineto{\pgfqpoint{1.692703in}{2.302863in}}%
\pgfpathlineto{\pgfqpoint{1.693560in}{2.303515in}}%
\pgfpathlineto{\pgfqpoint{1.694745in}{2.302849in}}%
\pgfpathlineto{\pgfqpoint{1.693918in}{2.303820in}}%
\pgfpathlineto{\pgfqpoint{1.694762in}{2.302875in}}%
\pgfpathlineto{\pgfqpoint{1.695895in}{2.303532in}}%
\pgfpathlineto{\pgfqpoint{1.694950in}{2.302725in}}%
\pgfpathlineto{\pgfqpoint{1.696031in}{2.303386in}}%
\pgfpathlineto{\pgfqpoint{1.697433in}{2.303597in}}%
\pgfpathlineto{\pgfqpoint{1.696070in}{2.303274in}}%
\pgfpathlineto{\pgfqpoint{1.697547in}{2.303302in}}%
\pgfpathlineto{\pgfqpoint{1.698613in}{2.302955in}}%
\pgfpathlineto{\pgfqpoint{1.698034in}{2.303533in}}%
\pgfpathlineto{\pgfqpoint{1.698672in}{2.302992in}}%
\pgfpathlineto{\pgfqpoint{1.699549in}{2.303511in}}%
\pgfpathlineto{\pgfqpoint{1.699804in}{2.303397in}}%
\pgfpathlineto{\pgfqpoint{1.699843in}{2.303350in}}%
\pgfpathlineto{\pgfqpoint{1.700442in}{2.304236in}}%
\pgfpathlineto{\pgfqpoint{1.700797in}{2.303944in}}%
\pgfpathlineto{\pgfqpoint{1.701858in}{2.304239in}}%
\pgfpathlineto{\pgfqpoint{1.701538in}{2.303626in}}%
\pgfpathlineto{\pgfqpoint{1.701924in}{2.304152in}}%
\pgfpathlineto{\pgfqpoint{1.702333in}{2.304010in}}%
\pgfpathlineto{\pgfqpoint{1.702813in}{2.304387in}}%
\pgfpathlineto{\pgfqpoint{1.703041in}{2.304105in}}%
\pgfpathlineto{\pgfqpoint{1.703335in}{2.304339in}}%
\pgfpathlineto{\pgfqpoint{1.704071in}{2.303659in}}%
\pgfpathlineto{\pgfqpoint{1.704136in}{2.303837in}}%
\pgfpathlineto{\pgfqpoint{1.706267in}{2.303140in}}%
\pgfpathlineto{\pgfqpoint{1.704149in}{2.303883in}}%
\pgfpathlineto{\pgfqpoint{1.706311in}{2.303242in}}%
\pgfpathlineto{\pgfqpoint{1.706983in}{2.303621in}}%
\pgfpathlineto{\pgfqpoint{1.706543in}{2.303129in}}%
\pgfpathlineto{\pgfqpoint{1.707430in}{2.303438in}}%
\pgfpathlineto{\pgfqpoint{1.709236in}{2.304732in}}%
\pgfpathlineto{\pgfqpoint{1.710351in}{2.305020in}}%
\pgfpathlineto{\pgfqpoint{1.709529in}{2.304577in}}%
\pgfpathlineto{\pgfqpoint{1.710365in}{2.304939in}}%
\pgfpathlineto{\pgfqpoint{1.710746in}{2.305117in}}%
\pgfpathlineto{\pgfqpoint{1.711411in}{2.304607in}}%
\pgfpathlineto{\pgfqpoint{1.711449in}{2.304645in}}%
\pgfpathlineto{\pgfqpoint{1.712104in}{2.304558in}}%
\pgfpathlineto{\pgfqpoint{1.712199in}{2.304959in}}%
\pgfpathlineto{\pgfqpoint{1.712557in}{2.304734in}}%
\pgfpathlineto{\pgfqpoint{1.713677in}{2.305278in}}%
\pgfpathlineto{\pgfqpoint{1.712622in}{2.304557in}}%
\pgfpathlineto{\pgfqpoint{1.713801in}{2.305118in}}%
\pgfpathlineto{\pgfqpoint{1.714881in}{2.304812in}}%
\pgfpathlineto{\pgfqpoint{1.713998in}{2.305240in}}%
\pgfpathlineto{\pgfqpoint{1.714938in}{2.304968in}}%
\pgfpathlineto{\pgfqpoint{1.716045in}{2.304894in}}%
\pgfpathlineto{\pgfqpoint{1.715550in}{2.304690in}}%
\pgfpathlineto{\pgfqpoint{1.716054in}{2.304856in}}%
\pgfpathlineto{\pgfqpoint{1.717043in}{2.304697in}}%
\pgfpathlineto{\pgfqpoint{1.716843in}{2.305045in}}%
\pgfpathlineto{\pgfqpoint{1.717161in}{2.304925in}}%
\pgfpathlineto{\pgfqpoint{1.717415in}{2.305217in}}%
\pgfpathlineto{\pgfqpoint{1.717851in}{2.304754in}}%
\pgfpathlineto{\pgfqpoint{1.718280in}{2.304976in}}%
\pgfpathlineto{\pgfqpoint{1.719359in}{2.303365in}}%
\pgfpathlineto{\pgfqpoint{1.719470in}{2.303421in}}%
\pgfpathlineto{\pgfqpoint{1.719508in}{2.303352in}}%
\pgfpathlineto{\pgfqpoint{1.720595in}{2.303441in}}%
\pgfpathlineto{\pgfqpoint{1.721700in}{2.303374in}}%
\pgfpathlineto{\pgfqpoint{1.720615in}{2.303450in}}%
\pgfpathlineto{\pgfqpoint{1.721707in}{2.303413in}}%
\pgfpathlineto{\pgfqpoint{1.722815in}{2.303381in}}%
\pgfpathlineto{\pgfqpoint{1.721724in}{2.303435in}}%
\pgfpathlineto{\pgfqpoint{1.722828in}{2.303404in}}%
\pgfpathlineto{\pgfqpoint{1.723800in}{2.303449in}}%
\pgfpathlineto{\pgfqpoint{1.723033in}{2.303351in}}%
\pgfpathlineto{\pgfqpoint{1.723941in}{2.303419in}}%
\pgfpathlineto{\pgfqpoint{1.724995in}{2.303377in}}%
\pgfpathlineto{\pgfqpoint{1.723962in}{2.303450in}}%
\pgfpathlineto{\pgfqpoint{1.725078in}{2.303403in}}%
\pgfpathlineto{\pgfqpoint{1.726120in}{2.303449in}}%
\pgfpathlineto{\pgfqpoint{1.725114in}{2.303337in}}%
\pgfpathlineto{\pgfqpoint{1.726189in}{2.303401in}}%
\pgfpathlineto{\pgfqpoint{1.727483in}{2.303439in}}%
\pgfpathlineto{\pgfqpoint{1.727486in}{2.303391in}}%
\pgfpathlineto{\pgfqpoint{1.728584in}{2.303370in}}%
\pgfpathlineto{\pgfqpoint{1.727522in}{2.303441in}}%
\pgfpathlineto{\pgfqpoint{1.728599in}{2.303416in}}%
\pgfpathlineto{\pgfqpoint{1.729686in}{2.303444in}}%
\pgfpathlineto{\pgfqpoint{1.728658in}{2.303353in}}%
\pgfpathlineto{\pgfqpoint{1.729712in}{2.303431in}}%
\pgfpathlineto{\pgfqpoint{1.730721in}{2.303350in}}%
\pgfpathlineto{\pgfqpoint{1.729862in}{2.303466in}}%
\pgfpathlineto{\pgfqpoint{1.730825in}{2.303370in}}%
\pgfpathlineto{\pgfqpoint{1.731948in}{2.303448in}}%
\pgfpathlineto{\pgfqpoint{1.730828in}{2.303367in}}%
\pgfpathlineto{\pgfqpoint{1.731949in}{2.303427in}}%
\pgfpathlineto{\pgfqpoint{1.732932in}{2.303324in}}%
\pgfpathlineto{\pgfqpoint{1.732089in}{2.303450in}}%
\pgfpathlineto{\pgfqpoint{1.733064in}{2.303339in}}%
\pgfpathlineto{\pgfqpoint{1.734159in}{2.303444in}}%
\pgfpathlineto{\pgfqpoint{1.733067in}{2.303334in}}%
\pgfpathlineto{\pgfqpoint{1.734185in}{2.303409in}}%
\pgfpathlineto{\pgfqpoint{1.735256in}{2.303363in}}%
\pgfpathlineto{\pgfqpoint{1.734221in}{2.303449in}}%
\pgfpathlineto{\pgfqpoint{1.735298in}{2.303410in}}%
\pgfpathlineto{\pgfqpoint{1.736342in}{2.303449in}}%
\pgfpathlineto{\pgfqpoint{1.735335in}{2.303360in}}%
\pgfpathlineto{\pgfqpoint{1.736409in}{2.303412in}}%
\pgfpathlineto{\pgfqpoint{1.737396in}{2.303318in}}%
\pgfpathlineto{\pgfqpoint{1.736412in}{2.303443in}}%
\pgfpathlineto{\pgfqpoint{1.737517in}{2.303516in}}%
\pgfpathlineto{\pgfqpoint{1.738321in}{2.303795in}}%
\pgfpathlineto{\pgfqpoint{1.737691in}{2.303361in}}%
\pgfpathlineto{\pgfqpoint{1.738632in}{2.303543in}}%
\pgfpathlineto{\pgfqpoint{1.738668in}{2.303403in}}%
\pgfpathlineto{\pgfqpoint{1.739587in}{2.304562in}}%
\pgfpathlineto{\pgfqpoint{1.739718in}{2.304454in}}%
\pgfpathlineto{\pgfqpoint{1.740254in}{2.304370in}}%
\pgfpathlineto{\pgfqpoint{1.740858in}{2.305000in}}%
\pgfpathlineto{\pgfqpoint{1.741933in}{2.304771in}}%
\pgfpathlineto{\pgfqpoint{1.740969in}{2.305225in}}%
\pgfpathlineto{\pgfqpoint{1.742007in}{2.304842in}}%
\pgfpathlineto{\pgfqpoint{1.742939in}{2.305290in}}%
\pgfpathlineto{\pgfqpoint{1.742061in}{2.304632in}}%
\pgfpathlineto{\pgfqpoint{1.743130in}{2.305055in}}%
\pgfpathlineto{\pgfqpoint{1.744269in}{2.305575in}}%
\pgfpathlineto{\pgfqpoint{1.743148in}{2.305014in}}%
\pgfpathlineto{\pgfqpoint{1.744298in}{2.305444in}}%
\pgfpathlineto{\pgfqpoint{1.744761in}{2.305127in}}%
\pgfpathlineto{\pgfqpoint{1.745245in}{2.305649in}}%
\pgfpathlineto{\pgfqpoint{1.745404in}{2.305548in}}%
\pgfpathlineto{\pgfqpoint{1.746406in}{2.306218in}}%
\pgfpathlineto{\pgfqpoint{1.746549in}{2.305983in}}%
\pgfpathlineto{\pgfqpoint{1.747148in}{2.305762in}}%
\pgfpathlineto{\pgfqpoint{1.747542in}{2.306202in}}%
\pgfpathlineto{\pgfqpoint{1.747656in}{2.306084in}}%
\pgfpathlineto{\pgfqpoint{1.748879in}{2.306666in}}%
\pgfpathlineto{\pgfqpoint{1.747975in}{2.305640in}}%
\pgfpathlineto{\pgfqpoint{1.748882in}{2.306653in}}%
\pgfpathlineto{\pgfqpoint{1.749939in}{2.306245in}}%
\pgfpathlineto{\pgfqpoint{1.749229in}{2.306876in}}%
\pgfpathlineto{\pgfqpoint{1.750069in}{2.306311in}}%
\pgfpathlineto{\pgfqpoint{1.750461in}{2.306488in}}%
\pgfpathlineto{\pgfqpoint{1.750929in}{2.305809in}}%
\pgfpathlineto{\pgfqpoint{1.751070in}{2.305840in}}%
\pgfpathlineto{\pgfqpoint{1.752129in}{2.304953in}}%
\pgfpathlineto{\pgfqpoint{1.751112in}{2.305938in}}%
\pgfpathlineto{\pgfqpoint{1.752229in}{2.305103in}}%
\pgfpathlineto{\pgfqpoint{1.753024in}{2.305175in}}%
\pgfpathlineto{\pgfqpoint{1.752572in}{2.304765in}}%
\pgfpathlineto{\pgfqpoint{1.753272in}{2.304717in}}%
\pgfpathlineto{\pgfqpoint{1.753275in}{2.304690in}}%
\pgfpathlineto{\pgfqpoint{1.754214in}{2.305433in}}%
\pgfpathlineto{\pgfqpoint{1.754355in}{2.305129in}}%
\pgfpathlineto{\pgfqpoint{1.755557in}{2.304730in}}%
\pgfpathlineto{\pgfqpoint{1.754825in}{2.305383in}}%
\pgfpathlineto{\pgfqpoint{1.755622in}{2.304851in}}%
\pgfpathlineto{\pgfqpoint{1.756559in}{2.304980in}}%
\pgfpathlineto{\pgfqpoint{1.756009in}{2.304623in}}%
\pgfpathlineto{\pgfqpoint{1.756721in}{2.304667in}}%
\pgfpathlineto{\pgfqpoint{1.757630in}{2.304047in}}%
\pgfpathlineto{\pgfqpoint{1.756727in}{2.304674in}}%
\pgfpathlineto{\pgfqpoint{1.757856in}{2.304221in}}%
\pgfpathlineto{\pgfqpoint{1.758416in}{2.304733in}}%
\pgfpathlineto{\pgfqpoint{1.758954in}{2.304109in}}%
\pgfpathlineto{\pgfqpoint{1.758965in}{2.304157in}}%
\pgfpathlineto{\pgfqpoint{1.759774in}{2.303795in}}%
\pgfpathlineto{\pgfqpoint{1.759183in}{2.304350in}}%
\pgfpathlineto{\pgfqpoint{1.760087in}{2.304039in}}%
\pgfpathlineto{\pgfqpoint{1.760454in}{2.303772in}}%
\pgfpathlineto{\pgfqpoint{1.760976in}{2.304580in}}%
\pgfpathlineto{\pgfqpoint{1.761078in}{2.304548in}}%
\pgfpathlineto{\pgfqpoint{1.761301in}{2.304579in}}%
\pgfpathlineto{\pgfqpoint{1.761802in}{2.304083in}}%
\pgfpathlineto{\pgfqpoint{1.762144in}{2.304083in}}%
\pgfpathlineto{\pgfqpoint{1.762729in}{2.303893in}}%
\pgfpathlineto{\pgfqpoint{1.763233in}{2.304533in}}%
\pgfpathlineto{\pgfqpoint{1.763277in}{2.304568in}}%
\pgfpathlineto{\pgfqpoint{1.764202in}{2.303941in}}%
\pgfpathlineto{\pgfqpoint{1.764304in}{2.304147in}}%
\pgfpathlineto{\pgfqpoint{1.765030in}{2.303741in}}%
\pgfpathlineto{\pgfqpoint{1.764600in}{2.304355in}}%
\pgfpathlineto{\pgfqpoint{1.765437in}{2.304140in}}%
\pgfpathlineto{\pgfqpoint{1.765769in}{2.304269in}}%
\pgfpathlineto{\pgfqpoint{1.766286in}{2.303756in}}%
\pgfpathlineto{\pgfqpoint{1.766539in}{2.303896in}}%
\pgfpathlineto{\pgfqpoint{1.766703in}{2.303729in}}%
\pgfpathlineto{\pgfqpoint{1.767527in}{2.304305in}}%
\pgfpathlineto{\pgfqpoint{1.767649in}{2.303952in}}%
\pgfpathlineto{\pgfqpoint{1.768763in}{2.304559in}}%
\pgfpathlineto{\pgfqpoint{1.767768in}{2.303747in}}%
\pgfpathlineto{\pgfqpoint{1.768805in}{2.304419in}}%
\pgfpathlineto{\pgfqpoint{1.769422in}{2.303904in}}%
\pgfpathlineto{\pgfqpoint{1.768906in}{2.304521in}}%
\pgfpathlineto{\pgfqpoint{1.769936in}{2.304076in}}%
\pgfpathlineto{\pgfqpoint{1.769947in}{2.304109in}}%
\pgfpathlineto{\pgfqpoint{1.770875in}{2.303314in}}%
\pgfpathlineto{\pgfqpoint{1.770996in}{2.303533in}}%
\pgfpathlineto{\pgfqpoint{1.772156in}{2.303220in}}%
\pgfpathlineto{\pgfqpoint{1.771045in}{2.303633in}}%
\pgfpathlineto{\pgfqpoint{1.772193in}{2.303328in}}%
\pgfpathlineto{\pgfqpoint{1.772989in}{2.303661in}}%
\pgfpathlineto{\pgfqpoint{1.773289in}{2.303122in}}%
\pgfpathlineto{\pgfqpoint{1.773292in}{2.303154in}}%
\pgfpathlineto{\pgfqpoint{1.774260in}{2.303111in}}%
\pgfpathlineto{\pgfqpoint{1.774020in}{2.303608in}}%
\pgfpathlineto{\pgfqpoint{1.774396in}{2.303483in}}%
\pgfpathlineto{\pgfqpoint{1.775273in}{2.303332in}}%
\pgfpathlineto{\pgfqpoint{1.775458in}{2.303852in}}%
\pgfpathlineto{\pgfqpoint{1.775473in}{2.303822in}}%
\pgfpathlineto{\pgfqpoint{1.776463in}{2.304188in}}%
\pgfpathlineto{\pgfqpoint{1.775480in}{2.303812in}}%
\pgfpathlineto{\pgfqpoint{1.776612in}{2.304160in}}%
\pgfpathlineto{\pgfqpoint{1.777663in}{2.303946in}}%
\pgfpathlineto{\pgfqpoint{1.776700in}{2.304408in}}%
\pgfpathlineto{\pgfqpoint{1.777725in}{2.304127in}}%
\pgfpathlineto{\pgfqpoint{1.778290in}{2.304457in}}%
\pgfpathlineto{\pgfqpoint{1.778639in}{2.304030in}}%
\pgfpathlineto{\pgfqpoint{1.778838in}{2.304233in}}%
\pgfpathlineto{\pgfqpoint{1.779991in}{2.303626in}}%
\pgfpathlineto{\pgfqpoint{1.779518in}{2.304433in}}%
\pgfpathlineto{\pgfqpoint{1.779991in}{2.303636in}}%
\pgfpathlineto{\pgfqpoint{1.781041in}{2.303902in}}%
\pgfpathlineto{\pgfqpoint{1.780318in}{2.303390in}}%
\pgfpathlineto{\pgfqpoint{1.781116in}{2.303724in}}%
\pgfpathlineto{\pgfqpoint{1.781217in}{2.303658in}}%
\pgfpathlineto{\pgfqpoint{1.781673in}{2.304167in}}%
\pgfpathlineto{\pgfqpoint{1.782221in}{2.303899in}}%
\pgfpathlineto{\pgfqpoint{1.782523in}{2.303603in}}%
\pgfpathlineto{\pgfqpoint{1.782769in}{2.304070in}}%
\pgfpathlineto{\pgfqpoint{1.783347in}{2.303800in}}%
\pgfpathlineto{\pgfqpoint{1.784646in}{2.304056in}}%
\pgfpathlineto{\pgfqpoint{1.783413in}{2.303700in}}%
\pgfpathlineto{\pgfqpoint{1.784676in}{2.303980in}}%
\pgfpathlineto{\pgfqpoint{1.784679in}{2.303945in}}%
\pgfpathlineto{\pgfqpoint{1.785742in}{2.304762in}}%
\pgfpathlineto{\pgfqpoint{1.785760in}{2.304697in}}%
\pgfpathlineto{\pgfqpoint{1.786928in}{2.303597in}}%
\pgfpathlineto{\pgfqpoint{1.785991in}{2.304782in}}%
\pgfpathlineto{\pgfqpoint{1.786937in}{2.303600in}}%
\pgfpathlineto{\pgfqpoint{1.788499in}{2.303927in}}%
\pgfpathlineto{\pgfqpoint{1.786938in}{2.303592in}}%
\pgfpathlineto{\pgfqpoint{1.788549in}{2.303774in}}%
\pgfpathlineto{\pgfqpoint{1.789601in}{2.303021in}}%
\pgfpathlineto{\pgfqpoint{1.789680in}{2.303161in}}%
\pgfpathlineto{\pgfqpoint{1.790835in}{2.303819in}}%
\pgfpathlineto{\pgfqpoint{1.789687in}{2.303128in}}%
\pgfpathlineto{\pgfqpoint{1.790839in}{2.303774in}}%
\pgfpathlineto{\pgfqpoint{1.792091in}{2.302848in}}%
\pgfpathlineto{\pgfqpoint{1.790927in}{2.303850in}}%
\pgfpathlineto{\pgfqpoint{1.792111in}{2.302934in}}%
\pgfpathlineto{\pgfqpoint{1.792636in}{2.303009in}}%
\pgfpathlineto{\pgfqpoint{1.793106in}{2.302433in}}%
\pgfpathlineto{\pgfqpoint{1.793206in}{2.302644in}}%
\pgfpathlineto{\pgfqpoint{1.793365in}{2.302619in}}%
\pgfpathlineto{\pgfqpoint{1.793613in}{2.303450in}}%
\pgfpathlineto{\pgfqpoint{1.794017in}{2.303412in}}%
\pgfpathlineto{\pgfqpoint{1.795016in}{2.303442in}}%
\pgfpathlineto{\pgfqpoint{1.794039in}{2.303356in}}%
\pgfpathlineto{\pgfqpoint{1.795129in}{2.303412in}}%
\pgfpathlineto{\pgfqpoint{1.796162in}{2.303451in}}%
\pgfpathlineto{\pgfqpoint{1.795135in}{2.303377in}}%
\pgfpathlineto{\pgfqpoint{1.796253in}{2.303407in}}%
\pgfpathlineto{\pgfqpoint{1.797069in}{2.303350in}}%
\pgfpathlineto{\pgfqpoint{1.796338in}{2.303446in}}%
\pgfpathlineto{\pgfqpoint{1.797367in}{2.303421in}}%
\pgfpathlineto{\pgfqpoint{1.798468in}{2.303444in}}%
\pgfpathlineto{\pgfqpoint{1.797621in}{2.303364in}}%
\pgfpathlineto{\pgfqpoint{1.798478in}{2.303421in}}%
\pgfpathlineto{\pgfqpoint{1.799587in}{2.303376in}}%
\pgfpathlineto{\pgfqpoint{1.798576in}{2.303445in}}%
\pgfpathlineto{\pgfqpoint{1.799592in}{2.303401in}}%
\pgfpathlineto{\pgfqpoint{1.800646in}{2.303459in}}%
\pgfpathlineto{\pgfqpoint{1.799841in}{2.303351in}}%
\pgfpathlineto{\pgfqpoint{1.800704in}{2.303423in}}%
\pgfpathlineto{\pgfqpoint{1.801804in}{2.303442in}}%
\pgfpathlineto{\pgfqpoint{1.800824in}{2.303364in}}%
\pgfpathlineto{\pgfqpoint{1.801824in}{2.303416in}}%
\pgfpathlineto{\pgfqpoint{1.802822in}{2.303358in}}%
\pgfpathlineto{\pgfqpoint{1.801838in}{2.303434in}}%
\pgfpathlineto{\pgfqpoint{1.802944in}{2.303417in}}%
\pgfpathlineto{\pgfqpoint{1.803064in}{2.303363in}}%
\pgfpathlineto{\pgfqpoint{1.804057in}{2.303443in}}%
\pgfpathlineto{\pgfqpoint{1.804878in}{2.303365in}}%
\pgfpathlineto{\pgfqpoint{1.804203in}{2.303461in}}%
\pgfpathlineto{\pgfqpoint{1.805170in}{2.303412in}}%
\pgfpathlineto{\pgfqpoint{1.805233in}{2.303368in}}%
\pgfpathlineto{\pgfqpoint{1.806282in}{2.303444in}}%
\pgfpathlineto{\pgfqpoint{1.807345in}{2.303345in}}%
\pgfpathlineto{\pgfqpoint{1.806350in}{2.303450in}}%
\pgfpathlineto{\pgfqpoint{1.807397in}{2.303417in}}%
\pgfpathlineto{\pgfqpoint{1.808498in}{2.303436in}}%
\pgfpathlineto{\pgfqpoint{1.807426in}{2.303380in}}%
\pgfpathlineto{\pgfqpoint{1.808524in}{2.303425in}}%
\pgfpathlineto{\pgfqpoint{1.809250in}{2.303332in}}%
\pgfpathlineto{\pgfqpoint{1.808622in}{2.303451in}}%
\pgfpathlineto{\pgfqpoint{1.809635in}{2.303412in}}%
\pgfpathlineto{\pgfqpoint{1.810617in}{2.303370in}}%
\pgfpathlineto{\pgfqpoint{1.809642in}{2.303437in}}%
\pgfpathlineto{\pgfqpoint{1.810752in}{2.303421in}}%
\pgfpathlineto{\pgfqpoint{1.811810in}{2.303451in}}%
\pgfpathlineto{\pgfqpoint{1.811125in}{2.303349in}}%
\pgfpathlineto{\pgfqpoint{1.811864in}{2.303426in}}%
\pgfpathlineto{\pgfqpoint{1.812649in}{2.303268in}}%
\pgfpathlineto{\pgfqpoint{1.812947in}{2.303767in}}%
\pgfpathlineto{\pgfqpoint{1.812970in}{2.303694in}}%
\pgfpathlineto{\pgfqpoint{1.813621in}{2.303998in}}%
\pgfpathlineto{\pgfqpoint{1.813986in}{2.303444in}}%
\pgfpathlineto{\pgfqpoint{1.814089in}{2.303782in}}%
\pgfpathlineto{\pgfqpoint{1.815196in}{2.303621in}}%
\pgfpathlineto{\pgfqpoint{1.814203in}{2.303965in}}%
\pgfpathlineto{\pgfqpoint{1.815225in}{2.303713in}}%
\pgfpathlineto{\pgfqpoint{1.816573in}{2.304696in}}%
\pgfpathlineto{\pgfqpoint{1.815330in}{2.303529in}}%
\pgfpathlineto{\pgfqpoint{1.816601in}{2.304586in}}%
\pgfpathlineto{\pgfqpoint{1.816750in}{2.304464in}}%
\pgfpathlineto{\pgfqpoint{1.817674in}{2.305018in}}%
\pgfpathlineto{\pgfqpoint{1.817687in}{2.304956in}}%
\pgfpathlineto{\pgfqpoint{1.818694in}{2.304984in}}%
\pgfpathlineto{\pgfqpoint{1.817942in}{2.304674in}}%
\pgfpathlineto{\pgfqpoint{1.818800in}{2.304875in}}%
\pgfpathlineto{\pgfqpoint{1.819591in}{2.304415in}}%
\pgfpathlineto{\pgfqpoint{1.818848in}{2.304939in}}%
\pgfpathlineto{\pgfqpoint{1.819922in}{2.304721in}}%
\pgfpathlineto{\pgfqpoint{1.820244in}{2.305039in}}%
\pgfpathlineto{\pgfqpoint{1.820493in}{2.304472in}}%
\pgfpathlineto{\pgfqpoint{1.821037in}{2.304767in}}%
\pgfpathlineto{\pgfqpoint{1.821431in}{2.304568in}}%
\pgfpathlineto{\pgfqpoint{1.821644in}{2.305047in}}%
\pgfpathlineto{\pgfqpoint{1.822142in}{2.304872in}}%
\pgfpathlineto{\pgfqpoint{1.823272in}{2.305761in}}%
\pgfpathlineto{\pgfqpoint{1.822287in}{2.304750in}}%
\pgfpathlineto{\pgfqpoint{1.823286in}{2.305734in}}%
\pgfpathlineto{\pgfqpoint{1.823913in}{2.305451in}}%
\pgfpathlineto{\pgfqpoint{1.823409in}{2.305805in}}%
\pgfpathlineto{\pgfqpoint{1.824407in}{2.305549in}}%
\pgfpathlineto{\pgfqpoint{1.825555in}{2.305352in}}%
\pgfpathlineto{\pgfqpoint{1.824918in}{2.305884in}}%
\pgfpathlineto{\pgfqpoint{1.825561in}{2.305390in}}%
\pgfpathlineto{\pgfqpoint{1.826635in}{2.305551in}}%
\pgfpathlineto{\pgfqpoint{1.825621in}{2.305312in}}%
\pgfpathlineto{\pgfqpoint{1.826677in}{2.305409in}}%
\pgfpathlineto{\pgfqpoint{1.827799in}{2.304737in}}%
\pgfpathlineto{\pgfqpoint{1.828052in}{2.304943in}}%
\pgfpathlineto{\pgfqpoint{1.828231in}{2.305156in}}%
\pgfpathlineto{\pgfqpoint{1.828638in}{2.304718in}}%
\pgfpathlineto{\pgfqpoint{1.829167in}{2.304910in}}%
\pgfpathlineto{\pgfqpoint{1.830225in}{2.304575in}}%
\pgfpathlineto{\pgfqpoint{1.830298in}{2.304692in}}%
\pgfpathlineto{\pgfqpoint{1.831332in}{2.304890in}}%
\pgfpathlineto{\pgfqpoint{1.830884in}{2.305244in}}%
\pgfpathlineto{\pgfqpoint{1.831412in}{2.305020in}}%
\pgfpathlineto{\pgfqpoint{1.831459in}{2.305165in}}%
\pgfpathlineto{\pgfqpoint{1.832176in}{2.304461in}}%
\pgfpathlineto{\pgfqpoint{1.832459in}{2.304441in}}%
\pgfpathlineto{\pgfqpoint{1.836969in}{2.304523in}}%
\pgfpathlineto{\pgfqpoint{1.832461in}{2.304455in}}%
\pgfpathlineto{\pgfqpoint{1.836970in}{2.304545in}}%
\pgfpathlineto{\pgfqpoint{1.837939in}{2.304672in}}%
\pgfpathlineto{\pgfqpoint{1.837109in}{2.304276in}}%
\pgfpathlineto{\pgfqpoint{1.838070in}{2.304393in}}%
\pgfpathlineto{\pgfqpoint{1.838386in}{2.303880in}}%
\pgfpathlineto{\pgfqpoint{1.838998in}{2.304597in}}%
\pgfpathlineto{\pgfqpoint{1.839180in}{2.304424in}}%
\pgfpathlineto{\pgfqpoint{1.839917in}{2.305055in}}%
\pgfpathlineto{\pgfqpoint{1.839302in}{2.304182in}}%
\pgfpathlineto{\pgfqpoint{1.840304in}{2.304658in}}%
\pgfpathlineto{\pgfqpoint{1.841369in}{2.304273in}}%
\pgfpathlineto{\pgfqpoint{1.840408in}{2.304743in}}%
\pgfpathlineto{\pgfqpoint{1.841515in}{2.304354in}}%
\pgfpathlineto{\pgfqpoint{1.842591in}{2.304886in}}%
\pgfpathlineto{\pgfqpoint{1.841856in}{2.304051in}}%
\pgfpathlineto{\pgfqpoint{1.842693in}{2.304615in}}%
\pgfpathlineto{\pgfqpoint{1.842988in}{2.304496in}}%
\pgfpathlineto{\pgfqpoint{1.843754in}{2.305019in}}%
\pgfpathlineto{\pgfqpoint{1.843761in}{2.304987in}}%
\pgfpathlineto{\pgfqpoint{1.846388in}{2.304291in}}%
\pgfpathlineto{\pgfqpoint{1.843767in}{2.304941in}}%
\pgfpathlineto{\pgfqpoint{1.846389in}{2.304285in}}%
\pgfpathlineto{\pgfqpoint{1.846390in}{2.304272in}}%
\pgfpathlineto{\pgfqpoint{1.846951in}{2.304834in}}%
\pgfpathlineto{\pgfqpoint{1.847488in}{2.304508in}}%
\pgfpathlineto{\pgfqpoint{1.847861in}{2.304737in}}%
\pgfpathlineto{\pgfqpoint{1.847619in}{2.304370in}}%
\pgfpathlineto{\pgfqpoint{1.848593in}{2.304399in}}%
\pgfpathlineto{\pgfqpoint{1.849686in}{2.304191in}}%
\pgfpathlineto{\pgfqpoint{1.848964in}{2.304832in}}%
\pgfpathlineto{\pgfqpoint{1.849714in}{2.304275in}}%
\pgfpathlineto{\pgfqpoint{1.849892in}{2.304574in}}%
\pgfpathlineto{\pgfqpoint{1.849719in}{2.304247in}}%
\pgfpathlineto{\pgfqpoint{1.850834in}{2.304415in}}%
\pgfpathlineto{\pgfqpoint{1.851728in}{2.304354in}}%
\pgfpathlineto{\pgfqpoint{1.851214in}{2.304678in}}%
\pgfpathlineto{\pgfqpoint{1.851946in}{2.304411in}}%
\pgfpathlineto{\pgfqpoint{1.852940in}{2.303947in}}%
\pgfpathlineto{\pgfqpoint{1.853118in}{2.304176in}}%
\pgfpathlineto{\pgfqpoint{1.853182in}{2.304302in}}%
\pgfpathlineto{\pgfqpoint{1.853739in}{2.303402in}}%
\pgfpathlineto{\pgfqpoint{1.854076in}{2.303615in}}%
\pgfpathlineto{\pgfqpoint{1.855296in}{2.303347in}}%
\pgfpathlineto{\pgfqpoint{1.854229in}{2.304014in}}%
\pgfpathlineto{\pgfqpoint{1.855313in}{2.303365in}}%
\pgfpathlineto{\pgfqpoint{1.855742in}{2.303578in}}%
\pgfpathlineto{\pgfqpoint{1.856192in}{2.303163in}}%
\pgfpathlineto{\pgfqpoint{1.856420in}{2.303263in}}%
\pgfpathlineto{\pgfqpoint{1.856771in}{2.302575in}}%
\pgfpathlineto{\pgfqpoint{1.856457in}{2.303332in}}%
\pgfpathlineto{\pgfqpoint{1.857553in}{2.302931in}}%
\pgfpathlineto{\pgfqpoint{1.857589in}{2.303037in}}%
\pgfpathlineto{\pgfqpoint{1.858639in}{2.302567in}}%
\pgfpathlineto{\pgfqpoint{1.858648in}{2.302592in}}%
\pgfpathlineto{\pgfqpoint{1.859498in}{2.302203in}}%
\pgfpathlineto{\pgfqpoint{1.858767in}{2.302808in}}%
\pgfpathlineto{\pgfqpoint{1.859807in}{2.302374in}}%
\pgfpathlineto{\pgfqpoint{1.860099in}{2.302560in}}%
\pgfpathlineto{\pgfqpoint{1.860544in}{2.301931in}}%
\pgfpathlineto{\pgfqpoint{1.860909in}{2.302188in}}%
\pgfpathlineto{\pgfqpoint{1.861934in}{2.302055in}}%
\pgfpathlineto{\pgfqpoint{1.860920in}{2.302117in}}%
\pgfpathlineto{\pgfqpoint{1.862093in}{2.301771in}}%
\pgfpathlineto{\pgfqpoint{1.862104in}{2.301747in}}%
\pgfpathlineto{\pgfqpoint{1.863111in}{2.302492in}}%
\pgfpathlineto{\pgfqpoint{1.863172in}{2.302387in}}%
\pgfpathlineto{\pgfqpoint{1.871269in}{2.303414in}}%
\pgfpathlineto{\pgfqpoint{1.872357in}{2.303446in}}%
\pgfpathlineto{\pgfqpoint{1.871629in}{2.303362in}}%
\pgfpathlineto{\pgfqpoint{1.872381in}{2.303427in}}%
\pgfpathlineto{\pgfqpoint{1.873384in}{2.303345in}}%
\pgfpathlineto{\pgfqpoint{1.872500in}{2.303451in}}%
\pgfpathlineto{\pgfqpoint{1.873497in}{2.303422in}}%
\pgfpathlineto{\pgfqpoint{1.874512in}{2.303443in}}%
\pgfpathlineto{\pgfqpoint{1.874082in}{2.303346in}}%
\pgfpathlineto{\pgfqpoint{1.874607in}{2.303417in}}%
\pgfpathlineto{\pgfqpoint{1.875551in}{2.303317in}}%
\pgfpathlineto{\pgfqpoint{1.874707in}{2.303455in}}%
\pgfpathlineto{\pgfqpoint{1.875723in}{2.303420in}}%
\pgfpathlineto{\pgfqpoint{1.876840in}{2.303448in}}%
\pgfpathlineto{\pgfqpoint{1.875738in}{2.303375in}}%
\pgfpathlineto{\pgfqpoint{1.876841in}{2.303440in}}%
\pgfpathlineto{\pgfqpoint{1.877951in}{2.303365in}}%
\pgfpathlineto{\pgfqpoint{1.876898in}{2.303448in}}%
\pgfpathlineto{\pgfqpoint{1.877954in}{2.303393in}}%
\pgfpathlineto{\pgfqpoint{1.879060in}{2.303437in}}%
\pgfpathlineto{\pgfqpoint{1.878342in}{2.303352in}}%
\pgfpathlineto{\pgfqpoint{1.879071in}{2.303435in}}%
\pgfpathlineto{\pgfqpoint{1.879281in}{2.303333in}}%
\pgfpathlineto{\pgfqpoint{1.879086in}{2.303449in}}%
\pgfpathlineto{\pgfqpoint{1.880182in}{2.303429in}}%
\pgfpathlineto{\pgfqpoint{1.881244in}{2.303371in}}%
\pgfpathlineto{\pgfqpoint{1.880250in}{2.303453in}}%
\pgfpathlineto{\pgfqpoint{1.881296in}{2.303415in}}%
\pgfpathlineto{\pgfqpoint{1.882404in}{2.303451in}}%
\pgfpathlineto{\pgfqpoint{1.881404in}{2.303364in}}%
\pgfpathlineto{\pgfqpoint{1.882408in}{2.303418in}}%
\pgfpathlineto{\pgfqpoint{1.883519in}{2.303372in}}%
\pgfpathlineto{\pgfqpoint{1.882425in}{2.303440in}}%
\pgfpathlineto{\pgfqpoint{1.883524in}{2.303412in}}%
\pgfpathlineto{\pgfqpoint{1.884624in}{2.303445in}}%
\pgfpathlineto{\pgfqpoint{1.883883in}{2.303361in}}%
\pgfpathlineto{\pgfqpoint{1.884635in}{2.303396in}}%
\pgfpathlineto{\pgfqpoint{1.885740in}{2.303369in}}%
\pgfpathlineto{\pgfqpoint{1.884750in}{2.303451in}}%
\pgfpathlineto{\pgfqpoint{1.885750in}{2.303406in}}%
\pgfpathlineto{\pgfqpoint{1.886809in}{2.303443in}}%
\pgfpathlineto{\pgfqpoint{1.885880in}{2.303362in}}%
\pgfpathlineto{\pgfqpoint{1.886871in}{2.303429in}}%
\pgfpathlineto{\pgfqpoint{1.887350in}{2.303139in}}%
\pgfpathlineto{\pgfqpoint{1.887973in}{2.303676in}}%
\pgfpathlineto{\pgfqpoint{1.887976in}{2.303663in}}%
\pgfpathlineto{\pgfqpoint{1.889108in}{2.304464in}}%
\pgfpathlineto{\pgfqpoint{1.888079in}{2.303429in}}%
\pgfpathlineto{\pgfqpoint{1.889125in}{2.304386in}}%
\pgfpathlineto{\pgfqpoint{1.889963in}{2.304086in}}%
\pgfpathlineto{\pgfqpoint{1.889444in}{2.304438in}}%
\pgfpathlineto{\pgfqpoint{1.890252in}{2.304322in}}%
\pgfpathlineto{\pgfqpoint{1.891183in}{2.304806in}}%
\pgfpathlineto{\pgfqpoint{1.890336in}{2.304149in}}%
\pgfpathlineto{\pgfqpoint{1.891382in}{2.304737in}}%
\pgfpathlineto{\pgfqpoint{1.891990in}{2.304518in}}%
\pgfpathlineto{\pgfqpoint{1.892413in}{2.305208in}}%
\pgfpathlineto{\pgfqpoint{1.892463in}{2.305067in}}%
\pgfpathlineto{\pgfqpoint{1.893604in}{2.306083in}}%
\pgfpathlineto{\pgfqpoint{1.892484in}{2.305007in}}%
\pgfpathlineto{\pgfqpoint{1.893605in}{2.306073in}}%
\pgfpathlineto{\pgfqpoint{1.894723in}{2.306709in}}%
\pgfpathlineto{\pgfqpoint{1.893657in}{2.305973in}}%
\pgfpathlineto{\pgfqpoint{1.894876in}{2.306486in}}%
\pgfpathlineto{\pgfqpoint{1.895943in}{2.305810in}}%
\pgfpathlineto{\pgfqpoint{1.895229in}{2.306514in}}%
\pgfpathlineto{\pgfqpoint{1.896006in}{2.305925in}}%
\pgfpathlineto{\pgfqpoint{1.896007in}{2.305942in}}%
\pgfpathlineto{\pgfqpoint{1.897005in}{2.305147in}}%
\pgfpathlineto{\pgfqpoint{1.897065in}{2.305194in}}%
\pgfpathlineto{\pgfqpoint{1.897393in}{2.304954in}}%
\pgfpathlineto{\pgfqpoint{1.897811in}{2.305459in}}%
\pgfpathlineto{\pgfqpoint{1.898154in}{2.305433in}}%
\pgfpathlineto{\pgfqpoint{1.899501in}{2.306928in}}%
\pgfpathlineto{\pgfqpoint{1.898155in}{2.305412in}}%
\pgfpathlineto{\pgfqpoint{1.899566in}{2.306805in}}%
\pgfpathlineto{\pgfqpoint{1.900720in}{2.306494in}}%
\pgfpathlineto{\pgfqpoint{1.899588in}{2.306826in}}%
\pgfpathlineto{\pgfqpoint{1.900720in}{2.306509in}}%
\pgfpathlineto{\pgfqpoint{1.900753in}{2.306559in}}%
\pgfpathlineto{\pgfqpoint{1.901430in}{2.305696in}}%
\pgfpathlineto{\pgfqpoint{1.901816in}{2.306003in}}%
\pgfpathlineto{\pgfqpoint{1.902762in}{2.306663in}}%
\pgfpathlineto{\pgfqpoint{1.902097in}{2.305929in}}%
\pgfpathlineto{\pgfqpoint{1.902976in}{2.306631in}}%
\pgfpathlineto{\pgfqpoint{1.903695in}{2.306054in}}%
\pgfpathlineto{\pgfqpoint{1.904107in}{2.306335in}}%
\pgfpathlineto{\pgfqpoint{1.905053in}{2.305982in}}%
\pgfpathlineto{\pgfqpoint{1.904453in}{2.306553in}}%
\pgfpathlineto{\pgfqpoint{1.905247in}{2.306299in}}%
\pgfpathlineto{\pgfqpoint{1.906068in}{2.306696in}}%
\pgfpathlineto{\pgfqpoint{1.905262in}{2.306279in}}%
\pgfpathlineto{\pgfqpoint{1.906368in}{2.306588in}}%
\pgfpathlineto{\pgfqpoint{1.908888in}{2.306380in}}%
\pgfpathlineto{\pgfqpoint{1.906421in}{2.306705in}}%
\pgfpathlineto{\pgfqpoint{1.908921in}{2.306502in}}%
\pgfpathlineto{\pgfqpoint{1.908922in}{2.306504in}}%
\pgfpathlineto{\pgfqpoint{1.909366in}{2.305790in}}%
\pgfpathlineto{\pgfqpoint{1.909926in}{2.306011in}}%
\pgfpathlineto{\pgfqpoint{1.910120in}{2.305752in}}%
\pgfpathlineto{\pgfqpoint{1.910764in}{2.306289in}}%
\pgfpathlineto{\pgfqpoint{1.911026in}{2.306182in}}%
\pgfpathlineto{\pgfqpoint{1.911693in}{2.306369in}}%
\pgfpathlineto{\pgfqpoint{1.911266in}{2.305929in}}%
\pgfpathlineto{\pgfqpoint{1.912131in}{2.306010in}}%
\pgfpathlineto{\pgfqpoint{1.912143in}{2.305977in}}%
\pgfpathlineto{\pgfqpoint{1.913155in}{2.306884in}}%
\pgfpathlineto{\pgfqpoint{1.913158in}{2.306874in}}%
\pgfpathlineto{\pgfqpoint{1.913159in}{2.306888in}}%
\pgfpathlineto{\pgfqpoint{1.914235in}{2.306219in}}%
\pgfpathlineto{\pgfqpoint{1.914238in}{2.306241in}}%
\pgfpathlineto{\pgfqpoint{1.915199in}{2.306033in}}%
\pgfpathlineto{\pgfqpoint{1.914483in}{2.306469in}}%
\pgfpathlineto{\pgfqpoint{1.915360in}{2.306080in}}%
\pgfpathlineto{\pgfqpoint{1.915361in}{2.306097in}}%
\pgfpathlineto{\pgfqpoint{1.916399in}{2.305541in}}%
\pgfpathlineto{\pgfqpoint{1.916443in}{2.305590in}}%
\pgfpathlineto{\pgfqpoint{1.917726in}{2.304399in}}%
\pgfpathlineto{\pgfqpoint{1.916513in}{2.305678in}}%
\pgfpathlineto{\pgfqpoint{1.917740in}{2.304456in}}%
\pgfpathlineto{\pgfqpoint{1.918359in}{2.304785in}}%
\pgfpathlineto{\pgfqpoint{1.918818in}{2.304262in}}%
\pgfpathlineto{\pgfqpoint{1.918846in}{2.304324in}}%
\pgfpathlineto{\pgfqpoint{1.919999in}{2.304649in}}%
\pgfpathlineto{\pgfqpoint{1.919136in}{2.303881in}}%
\pgfpathlineto{\pgfqpoint{1.920074in}{2.304507in}}%
\pgfpathlineto{\pgfqpoint{1.920583in}{2.304331in}}%
\pgfpathlineto{\pgfqpoint{1.921003in}{2.304858in}}%
\pgfpathlineto{\pgfqpoint{1.921182in}{2.304594in}}%
\pgfpathlineto{\pgfqpoint{1.922080in}{2.304869in}}%
\pgfpathlineto{\pgfqpoint{1.921916in}{2.304510in}}%
\pgfpathlineto{\pgfqpoint{1.922297in}{2.304729in}}%
\pgfpathlineto{\pgfqpoint{1.923329in}{2.304372in}}%
\pgfpathlineto{\pgfqpoint{1.922614in}{2.304899in}}%
\pgfpathlineto{\pgfqpoint{1.923430in}{2.304672in}}%
\pgfpathlineto{\pgfqpoint{1.924591in}{2.305208in}}%
\pgfpathlineto{\pgfqpoint{1.923591in}{2.304476in}}%
\pgfpathlineto{\pgfqpoint{1.924594in}{2.305185in}}%
\pgfpathlineto{\pgfqpoint{1.925126in}{2.305093in}}%
\pgfpathlineto{\pgfqpoint{1.924906in}{2.305475in}}%
\pgfpathlineto{\pgfqpoint{1.925697in}{2.305379in}}%
\pgfpathlineto{\pgfqpoint{1.939307in}{2.305259in}}%
\pgfpathlineto{\pgfqpoint{1.939309in}{2.305312in}}%
\pgfpathlineto{\pgfqpoint{1.940146in}{2.305625in}}%
\pgfpathlineto{\pgfqpoint{1.939356in}{2.305219in}}%
\pgfpathlineto{\pgfqpoint{1.940435in}{2.305423in}}%
\pgfpathlineto{\pgfqpoint{1.941492in}{2.306181in}}%
\pgfpathlineto{\pgfqpoint{1.940618in}{2.305214in}}%
\pgfpathlineto{\pgfqpoint{1.941687in}{2.305929in}}%
\pgfpathlineto{\pgfqpoint{1.942270in}{2.305681in}}%
\pgfpathlineto{\pgfqpoint{1.942736in}{2.306613in}}%
\pgfpathlineto{\pgfqpoint{1.942775in}{2.306531in}}%
\pgfpathlineto{\pgfqpoint{1.944581in}{2.303381in}}%
\pgfpathlineto{\pgfqpoint{1.943246in}{2.306966in}}%
\pgfpathlineto{\pgfqpoint{1.944621in}{2.303417in}}%
\pgfpathlineto{\pgfqpoint{1.945692in}{2.303445in}}%
\pgfpathlineto{\pgfqpoint{1.944743in}{2.303365in}}%
\pgfpathlineto{\pgfqpoint{1.945732in}{2.303399in}}%
\pgfpathlineto{\pgfqpoint{1.946810in}{2.303443in}}%
\pgfpathlineto{\pgfqpoint{1.945762in}{2.303384in}}%
\pgfpathlineto{\pgfqpoint{1.946848in}{2.303423in}}%
\pgfpathlineto{\pgfqpoint{1.948126in}{2.303377in}}%
\pgfpathlineto{\pgfqpoint{1.946849in}{2.303426in}}%
\pgfpathlineto{\pgfqpoint{1.948128in}{2.303417in}}%
\pgfpathlineto{\pgfqpoint{1.949160in}{2.303453in}}%
\pgfpathlineto{\pgfqpoint{1.948310in}{2.303358in}}%
\pgfpathlineto{\pgfqpoint{1.949240in}{2.303413in}}%
\pgfpathlineto{\pgfqpoint{1.950196in}{2.303356in}}%
\pgfpathlineto{\pgfqpoint{1.949295in}{2.303448in}}%
\pgfpathlineto{\pgfqpoint{1.950358in}{2.303394in}}%
\pgfpathlineto{\pgfqpoint{1.951427in}{2.303446in}}%
\pgfpathlineto{\pgfqpoint{1.950454in}{2.303364in}}%
\pgfpathlineto{\pgfqpoint{1.951491in}{2.303433in}}%
\pgfpathlineto{\pgfqpoint{1.952605in}{2.303373in}}%
\pgfpathlineto{\pgfqpoint{1.951536in}{2.303445in}}%
\pgfpathlineto{\pgfqpoint{1.952607in}{2.303376in}}%
\pgfpathlineto{\pgfqpoint{1.953626in}{2.303362in}}%
\pgfpathlineto{\pgfqpoint{1.952612in}{2.303441in}}%
\pgfpathlineto{\pgfqpoint{1.953718in}{2.303422in}}%
\pgfpathlineto{\pgfqpoint{1.954822in}{2.303444in}}%
\pgfpathlineto{\pgfqpoint{1.953834in}{2.303349in}}%
\pgfpathlineto{\pgfqpoint{1.954831in}{2.303430in}}%
\pgfpathlineto{\pgfqpoint{1.955836in}{2.303368in}}%
\pgfpathlineto{\pgfqpoint{1.954932in}{2.303449in}}%
\pgfpathlineto{\pgfqpoint{1.955944in}{2.303406in}}%
\pgfpathlineto{\pgfqpoint{1.957024in}{2.303449in}}%
\pgfpathlineto{\pgfqpoint{1.956051in}{2.303375in}}%
\pgfpathlineto{\pgfqpoint{1.957056in}{2.303407in}}%
\pgfpathlineto{\pgfqpoint{1.957961in}{2.303342in}}%
\pgfpathlineto{\pgfqpoint{1.957091in}{2.303444in}}%
\pgfpathlineto{\pgfqpoint{1.958172in}{2.303417in}}%
\pgfpathlineto{\pgfqpoint{1.959278in}{2.303435in}}%
\pgfpathlineto{\pgfqpoint{1.958237in}{2.303376in}}%
\pgfpathlineto{\pgfqpoint{1.959284in}{2.303422in}}%
\pgfpathlineto{\pgfqpoint{1.960374in}{2.303440in}}%
\pgfpathlineto{\pgfqpoint{1.959298in}{2.303378in}}%
\pgfpathlineto{\pgfqpoint{1.960407in}{2.303424in}}%
\pgfpathlineto{\pgfqpoint{1.960421in}{2.303297in}}%
\pgfpathlineto{\pgfqpoint{1.960478in}{2.303442in}}%
\pgfpathlineto{\pgfqpoint{1.961521in}{2.303398in}}%
\pgfpathlineto{\pgfqpoint{1.962350in}{2.303845in}}%
\pgfpathlineto{\pgfqpoint{1.961992in}{2.303347in}}%
\pgfpathlineto{\pgfqpoint{1.962637in}{2.303506in}}%
\pgfpathlineto{\pgfqpoint{1.963067in}{2.303435in}}%
\pgfpathlineto{\pgfqpoint{1.963714in}{2.304089in}}%
\pgfpathlineto{\pgfqpoint{1.963714in}{2.304073in}}%
\pgfpathlineto{\pgfqpoint{1.964147in}{2.304322in}}%
\pgfpathlineto{\pgfqpoint{1.964726in}{2.303958in}}%
\pgfpathlineto{\pgfqpoint{1.964823in}{2.303987in}}%
\pgfpathlineto{\pgfqpoint{1.965416in}{2.304244in}}%
\pgfpathlineto{\pgfqpoint{1.965002in}{2.303854in}}%
\pgfpathlineto{\pgfqpoint{1.965935in}{2.303965in}}%
\pgfpathlineto{\pgfqpoint{1.966719in}{2.303708in}}%
\pgfpathlineto{\pgfqpoint{1.966047in}{2.304118in}}%
\pgfpathlineto{\pgfqpoint{1.967064in}{2.303892in}}%
\pgfpathlineto{\pgfqpoint{1.967169in}{2.303994in}}%
\pgfpathlineto{\pgfqpoint{1.967878in}{2.303538in}}%
\pgfpathlineto{\pgfqpoint{1.968163in}{2.303663in}}%
\pgfpathlineto{\pgfqpoint{1.969211in}{2.303744in}}%
\pgfpathlineto{\pgfqpoint{1.968781in}{2.304207in}}%
\pgfpathlineto{\pgfqpoint{1.969253in}{2.303884in}}%
\pgfpathlineto{\pgfqpoint{1.970226in}{2.304042in}}%
\pgfpathlineto{\pgfqpoint{1.969288in}{2.303842in}}%
\pgfpathlineto{\pgfqpoint{1.970392in}{2.303869in}}%
\pgfpathlineto{\pgfqpoint{1.971063in}{2.303368in}}%
\pgfpathlineto{\pgfqpoint{1.970398in}{2.303874in}}%
\pgfpathlineto{\pgfqpoint{1.971510in}{2.303798in}}%
\pgfpathlineto{\pgfqpoint{1.972247in}{2.303993in}}%
\pgfpathlineto{\pgfqpoint{1.971989in}{2.303514in}}%
\pgfpathlineto{\pgfqpoint{1.972619in}{2.303738in}}%
\pgfpathlineto{\pgfqpoint{1.973733in}{2.303321in}}%
\pgfpathlineto{\pgfqpoint{1.973009in}{2.303885in}}%
\pgfpathlineto{\pgfqpoint{1.973755in}{2.303351in}}%
\pgfpathlineto{\pgfqpoint{1.974626in}{2.302805in}}%
\pgfpathlineto{\pgfqpoint{1.974473in}{2.302586in}}%
\pgfpathlineto{\pgfqpoint{1.974661in}{2.302717in}}%
\pgfpathlineto{\pgfqpoint{1.975484in}{2.302629in}}%
\pgfpathlineto{\pgfqpoint{1.974955in}{2.302985in}}%
\pgfpathlineto{\pgfqpoint{1.975769in}{2.302788in}}%
\pgfpathlineto{\pgfqpoint{1.976549in}{2.303241in}}%
\pgfpathlineto{\pgfqpoint{1.975981in}{2.302535in}}%
\pgfpathlineto{\pgfqpoint{1.976926in}{2.303039in}}%
\pgfpathlineto{\pgfqpoint{1.977232in}{2.302925in}}%
\pgfpathlineto{\pgfqpoint{1.977985in}{2.303465in}}%
\pgfpathlineto{\pgfqpoint{1.978020in}{2.303420in}}%
\pgfpathlineto{\pgfqpoint{1.984516in}{2.301958in}}%
\pgfpathlineto{\pgfqpoint{1.984541in}{2.302075in}}%
\pgfpathlineto{\pgfqpoint{1.985704in}{2.302388in}}%
\pgfpathlineto{\pgfqpoint{1.985449in}{2.301785in}}%
\pgfpathlineto{\pgfqpoint{1.985716in}{2.302377in}}%
\pgfpathlineto{\pgfqpoint{1.986504in}{2.302322in}}%
\pgfpathlineto{\pgfqpoint{1.986099in}{2.302594in}}%
\pgfpathlineto{\pgfqpoint{1.986813in}{2.302558in}}%
\pgfpathlineto{\pgfqpoint{1.987837in}{2.302925in}}%
\pgfpathlineto{\pgfqpoint{1.986966in}{2.302318in}}%
\pgfpathlineto{\pgfqpoint{1.987945in}{2.302775in}}%
\pgfpathlineto{\pgfqpoint{1.988732in}{2.302517in}}%
\pgfpathlineto{\pgfqpoint{1.988369in}{2.302940in}}%
\pgfpathlineto{\pgfqpoint{1.989055in}{2.302836in}}%
\pgfpathlineto{\pgfqpoint{1.990161in}{2.303103in}}%
\pgfpathlineto{\pgfqpoint{1.989428in}{2.302628in}}%
\pgfpathlineto{\pgfqpoint{1.990183in}{2.303057in}}%
\pgfpathlineto{\pgfqpoint{1.991426in}{2.301808in}}%
\pgfpathlineto{\pgfqpoint{1.990189in}{2.303073in}}%
\pgfpathlineto{\pgfqpoint{1.991427in}{2.301825in}}%
\pgfpathlineto{\pgfqpoint{1.993082in}{2.302318in}}%
\pgfpathlineto{\pgfqpoint{1.991468in}{2.301727in}}%
\pgfpathlineto{\pgfqpoint{1.993096in}{2.302261in}}%
\pgfpathlineto{\pgfqpoint{1.993760in}{2.302321in}}%
\pgfpathlineto{\pgfqpoint{1.994232in}{2.302019in}}%
\pgfpathlineto{\pgfqpoint{1.994686in}{2.302299in}}%
\pgfpathlineto{\pgfqpoint{1.995331in}{2.301804in}}%
\pgfpathlineto{\pgfqpoint{1.995333in}{2.301804in}}%
\pgfpathlineto{\pgfqpoint{1.999548in}{2.301441in}}%
\pgfpathlineto{\pgfqpoint{1.999556in}{2.301357in}}%
\pgfpathlineto{\pgfqpoint{1.999993in}{2.301315in}}%
\pgfpathlineto{\pgfqpoint{2.000397in}{2.301811in}}%
\pgfpathlineto{\pgfqpoint{2.000646in}{2.301793in}}%
\pgfpathlineto{\pgfqpoint{2.001610in}{2.302083in}}%
\pgfpathlineto{\pgfqpoint{2.001000in}{2.301770in}}%
\pgfpathlineto{\pgfqpoint{2.001764in}{2.301928in}}%
\pgfpathlineto{\pgfqpoint{2.002472in}{2.301709in}}%
\pgfpathlineto{\pgfqpoint{2.001883in}{2.302172in}}%
\pgfpathlineto{\pgfqpoint{2.002871in}{2.302017in}}%
\pgfpathlineto{\pgfqpoint{2.003719in}{2.302003in}}%
\pgfpathlineto{\pgfqpoint{2.003235in}{2.301529in}}%
\pgfpathlineto{\pgfqpoint{2.003973in}{2.301886in}}%
\pgfpathlineto{\pgfqpoint{2.005053in}{2.301436in}}%
\pgfpathlineto{\pgfqpoint{2.004501in}{2.301925in}}%
\pgfpathlineto{\pgfqpoint{2.005104in}{2.301495in}}%
\pgfpathlineto{\pgfqpoint{2.005989in}{2.301727in}}%
\pgfpathlineto{\pgfqpoint{2.006181in}{2.301329in}}%
\pgfpathlineto{\pgfqpoint{2.006212in}{2.301408in}}%
\pgfpathlineto{\pgfqpoint{2.007344in}{2.300955in}}%
\pgfpathlineto{\pgfqpoint{2.006222in}{2.301414in}}%
\pgfpathlineto{\pgfqpoint{2.007356in}{2.300996in}}%
\pgfpathlineto{\pgfqpoint{2.008236in}{2.301094in}}%
\pgfpathlineto{\pgfqpoint{2.007637in}{2.300605in}}%
\pgfpathlineto{\pgfqpoint{2.008454in}{2.300813in}}%
\pgfpathlineto{\pgfqpoint{2.009630in}{2.300700in}}%
\pgfpathlineto{\pgfqpoint{2.008633in}{2.301101in}}%
\pgfpathlineto{\pgfqpoint{2.009654in}{2.300806in}}%
\pgfpathlineto{\pgfqpoint{2.009700in}{2.300862in}}%
\pgfpathlineto{\pgfqpoint{2.010488in}{2.300356in}}%
\pgfpathlineto{\pgfqpoint{2.010744in}{2.300437in}}%
\pgfpathlineto{\pgfqpoint{2.011735in}{2.300486in}}%
\pgfpathlineto{\pgfqpoint{2.011111in}{2.300748in}}%
\pgfpathlineto{\pgfqpoint{2.011839in}{2.300611in}}%
\pgfpathlineto{\pgfqpoint{2.012039in}{2.300951in}}%
\pgfpathlineto{\pgfqpoint{2.012901in}{2.300380in}}%
\pgfpathlineto{\pgfqpoint{2.012945in}{2.300404in}}%
\pgfpathlineto{\pgfqpoint{2.012951in}{2.300369in}}%
\pgfpathlineto{\pgfqpoint{2.013387in}{2.300908in}}%
\pgfpathlineto{\pgfqpoint{2.014023in}{2.300745in}}%
\pgfpathlineto{\pgfqpoint{2.014811in}{2.301311in}}%
\pgfpathlineto{\pgfqpoint{2.014137in}{2.300425in}}%
\pgfpathlineto{\pgfqpoint{2.015136in}{2.300792in}}%
\pgfpathlineto{\pgfqpoint{2.015297in}{2.300556in}}%
\pgfpathlineto{\pgfqpoint{2.015961in}{2.301381in}}%
\pgfpathlineto{\pgfqpoint{2.016234in}{2.301210in}}%
\pgfpathlineto{\pgfqpoint{2.017130in}{2.301275in}}%
\pgfpathlineto{\pgfqpoint{2.016785in}{2.301548in}}%
\pgfpathlineto{\pgfqpoint{2.017328in}{2.301443in}}%
\pgfpathlineto{\pgfqpoint{2.018697in}{2.303441in}}%
\pgfpathlineto{\pgfqpoint{2.017818in}{2.301304in}}%
\pgfpathlineto{\pgfqpoint{2.018763in}{2.303408in}}%
\pgfpathlineto{\pgfqpoint{2.019790in}{2.303361in}}%
\pgfpathlineto{\pgfqpoint{2.018786in}{2.303444in}}%
\pgfpathlineto{\pgfqpoint{2.019874in}{2.303407in}}%
\pgfpathlineto{\pgfqpoint{2.020981in}{2.303450in}}%
\pgfpathlineto{\pgfqpoint{2.020036in}{2.303356in}}%
\pgfpathlineto{\pgfqpoint{2.020987in}{2.303439in}}%
\pgfpathlineto{\pgfqpoint{2.021957in}{2.303368in}}%
\pgfpathlineto{\pgfqpoint{2.021089in}{2.303449in}}%
\pgfpathlineto{\pgfqpoint{2.022100in}{2.303413in}}%
\pgfpathlineto{\pgfqpoint{2.023047in}{2.303360in}}%
\pgfpathlineto{\pgfqpoint{2.022124in}{2.303441in}}%
\pgfpathlineto{\pgfqpoint{2.023213in}{2.303418in}}%
\pgfpathlineto{\pgfqpoint{2.024315in}{2.303450in}}%
\pgfpathlineto{\pgfqpoint{2.023587in}{2.303344in}}%
\pgfpathlineto{\pgfqpoint{2.024324in}{2.303424in}}%
\pgfpathlineto{\pgfqpoint{2.025450in}{2.303383in}}%
\pgfpathlineto{\pgfqpoint{2.024328in}{2.303439in}}%
\pgfpathlineto{\pgfqpoint{2.025463in}{2.303423in}}%
\pgfpathlineto{\pgfqpoint{2.026557in}{2.303450in}}%
\pgfpathlineto{\pgfqpoint{2.025466in}{2.303351in}}%
\pgfpathlineto{\pgfqpoint{2.026594in}{2.303405in}}%
\pgfpathlineto{\pgfqpoint{2.027574in}{2.303359in}}%
\pgfpathlineto{\pgfqpoint{2.026660in}{2.303441in}}%
\pgfpathlineto{\pgfqpoint{2.027706in}{2.303422in}}%
\pgfpathlineto{\pgfqpoint{2.028794in}{2.303439in}}%
\pgfpathlineto{\pgfqpoint{2.027860in}{2.303333in}}%
\pgfpathlineto{\pgfqpoint{2.028830in}{2.303404in}}%
\pgfpathlineto{\pgfqpoint{2.029811in}{2.303371in}}%
\pgfpathlineto{\pgfqpoint{2.028865in}{2.303452in}}%
\pgfpathlineto{\pgfqpoint{2.029941in}{2.303424in}}%
\pgfpathlineto{\pgfqpoint{2.030745in}{2.303347in}}%
\pgfpathlineto{\pgfqpoint{2.031053in}{2.303449in}}%
\pgfpathlineto{\pgfqpoint{2.032055in}{2.303300in}}%
\pgfpathlineto{\pgfqpoint{2.032165in}{2.303407in}}%
\pgfpathlineto{\pgfqpoint{2.032225in}{2.303376in}}%
\pgfpathlineto{\pgfqpoint{2.033282in}{2.303436in}}%
\pgfpathlineto{\pgfqpoint{2.034236in}{2.303306in}}%
\pgfpathlineto{\pgfqpoint{2.034399in}{2.303420in}}%
\pgfpathlineto{\pgfqpoint{2.035487in}{2.303444in}}%
\pgfpathlineto{\pgfqpoint{2.034842in}{2.303346in}}%
\pgfpathlineto{\pgfqpoint{2.035510in}{2.303405in}}%
\pgfpathlineto{\pgfqpoint{2.036623in}{2.303379in}}%
\pgfpathlineto{\pgfqpoint{2.035547in}{2.303447in}}%
\pgfpathlineto{\pgfqpoint{2.036637in}{2.303400in}}%
\pgfpathlineto{\pgfqpoint{2.037578in}{2.303836in}}%
\pgfpathlineto{\pgfqpoint{2.037047in}{2.303311in}}%
\pgfpathlineto{\pgfqpoint{2.037768in}{2.303678in}}%
\pgfpathlineto{\pgfqpoint{2.038793in}{2.302919in}}%
\pgfpathlineto{\pgfqpoint{2.037798in}{2.303714in}}%
\pgfpathlineto{\pgfqpoint{2.038895in}{2.303208in}}%
\pgfpathlineto{\pgfqpoint{2.040019in}{2.303318in}}%
\pgfpathlineto{\pgfqpoint{2.039092in}{2.302786in}}%
\pgfpathlineto{\pgfqpoint{2.040046in}{2.303164in}}%
\pgfpathlineto{\pgfqpoint{2.040670in}{2.303541in}}%
\pgfpathlineto{\pgfqpoint{2.041159in}{2.303060in}}%
\pgfpathlineto{\pgfqpoint{2.041607in}{2.303205in}}%
\pgfpathlineto{\pgfqpoint{2.041905in}{2.302556in}}%
\pgfpathlineto{\pgfqpoint{2.042254in}{2.302841in}}%
\pgfpathlineto{\pgfqpoint{2.042437in}{2.302645in}}%
\pgfpathlineto{\pgfqpoint{2.043275in}{2.303256in}}%
\pgfpathlineto{\pgfqpoint{2.043345in}{2.303163in}}%
\pgfpathlineto{\pgfqpoint{2.044321in}{2.303505in}}%
\pgfpathlineto{\pgfqpoint{2.044039in}{2.302939in}}%
\pgfpathlineto{\pgfqpoint{2.044477in}{2.303396in}}%
\pgfpathlineto{\pgfqpoint{2.045624in}{2.303073in}}%
\pgfpathlineto{\pgfqpoint{2.045274in}{2.303656in}}%
\pgfpathlineto{\pgfqpoint{2.045630in}{2.303093in}}%
\pgfpathlineto{\pgfqpoint{2.045751in}{2.303153in}}%
\pgfpathlineto{\pgfqpoint{2.046343in}{2.302533in}}%
\pgfpathlineto{\pgfqpoint{2.046648in}{2.302515in}}%
\pgfpathlineto{\pgfqpoint{2.047776in}{2.302536in}}%
\pgfpathlineto{\pgfqpoint{2.046806in}{2.302833in}}%
\pgfpathlineto{\pgfqpoint{2.047790in}{2.302563in}}%
\pgfpathlineto{\pgfqpoint{2.048842in}{2.303090in}}%
\pgfpathlineto{\pgfqpoint{2.048921in}{2.302913in}}%
\pgfpathlineto{\pgfqpoint{2.049556in}{2.302744in}}%
\pgfpathlineto{\pgfqpoint{2.049880in}{2.303133in}}%
\pgfpathlineto{\pgfqpoint{2.050031in}{2.302968in}}%
\pgfpathlineto{\pgfqpoint{2.050569in}{2.302962in}}%
\pgfpathlineto{\pgfqpoint{2.050772in}{2.302578in}}%
\pgfpathlineto{\pgfqpoint{2.051103in}{2.302694in}}%
\pgfpathlineto{\pgfqpoint{2.051127in}{2.302594in}}%
\pgfpathlineto{\pgfqpoint{2.052114in}{2.303198in}}%
\pgfpathlineto{\pgfqpoint{2.052198in}{2.303152in}}%
\pgfpathlineto{\pgfqpoint{2.053300in}{2.303291in}}%
\pgfpathlineto{\pgfqpoint{2.052479in}{2.302882in}}%
\pgfpathlineto{\pgfqpoint{2.053406in}{2.303015in}}%
\pgfpathlineto{\pgfqpoint{2.054149in}{2.302932in}}%
\pgfpathlineto{\pgfqpoint{2.053830in}{2.303288in}}%
\pgfpathlineto{\pgfqpoint{2.054525in}{2.303030in}}%
\pgfpathlineto{\pgfqpoint{2.054768in}{2.303152in}}%
\pgfpathlineto{\pgfqpoint{2.055615in}{2.302470in}}%
\pgfpathlineto{\pgfqpoint{2.058625in}{2.303011in}}%
\pgfpathlineto{\pgfqpoint{2.055623in}{2.302520in}}%
\pgfpathlineto{\pgfqpoint{2.058630in}{2.303049in}}%
\pgfpathlineto{\pgfqpoint{2.058936in}{2.303453in}}%
\pgfpathlineto{\pgfqpoint{2.059307in}{2.302976in}}%
\pgfpathlineto{\pgfqpoint{2.059752in}{2.303228in}}%
\pgfpathlineto{\pgfqpoint{2.060563in}{2.302797in}}%
\pgfpathlineto{\pgfqpoint{2.060891in}{2.302929in}}%
\pgfpathlineto{\pgfqpoint{2.062127in}{2.303583in}}%
\pgfpathlineto{\pgfqpoint{2.062215in}{2.303383in}}%
\pgfpathlineto{\pgfqpoint{2.063241in}{2.302706in}}%
\pgfpathlineto{\pgfqpoint{2.062256in}{2.303417in}}%
\pgfpathlineto{\pgfqpoint{2.063340in}{2.302836in}}%
\pgfpathlineto{\pgfqpoint{2.064542in}{2.304077in}}%
\pgfpathlineto{\pgfqpoint{2.063346in}{2.302779in}}%
\pgfpathlineto{\pgfqpoint{2.064675in}{2.303978in}}%
\pgfpathlineto{\pgfqpoint{2.065791in}{2.303871in}}%
\pgfpathlineto{\pgfqpoint{2.064870in}{2.304202in}}%
\pgfpathlineto{\pgfqpoint{2.065801in}{2.303906in}}%
\pgfpathlineto{\pgfqpoint{2.066811in}{2.304435in}}%
\pgfpathlineto{\pgfqpoint{2.066053in}{2.303602in}}%
\pgfpathlineto{\pgfqpoint{2.066922in}{2.304173in}}%
\pgfpathlineto{\pgfqpoint{2.068131in}{2.303505in}}%
\pgfpathlineto{\pgfqpoint{2.066965in}{2.304250in}}%
\pgfpathlineto{\pgfqpoint{2.068136in}{2.303562in}}%
\pgfpathlineto{\pgfqpoint{2.068337in}{2.303603in}}%
\pgfpathlineto{\pgfqpoint{2.069193in}{2.303005in}}%
\pgfpathlineto{\pgfqpoint{2.069222in}{2.303096in}}%
\pgfpathlineto{\pgfqpoint{2.069980in}{2.302477in}}%
\pgfpathlineto{\pgfqpoint{2.069240in}{2.303110in}}%
\pgfpathlineto{\pgfqpoint{2.070335in}{2.303135in}}%
\pgfpathlineto{\pgfqpoint{2.070550in}{2.303490in}}%
\pgfpathlineto{\pgfqpoint{2.071418in}{2.302864in}}%
\pgfpathlineto{\pgfqpoint{2.071427in}{2.302880in}}%
\pgfpathlineto{\pgfqpoint{2.072221in}{2.302337in}}%
\pgfpathlineto{\pgfqpoint{2.072565in}{2.302628in}}%
\pgfpathlineto{\pgfqpoint{2.073450in}{2.303218in}}%
\pgfpathlineto{\pgfqpoint{2.072608in}{2.302772in}}%
\pgfpathlineto{\pgfqpoint{2.073477in}{2.303314in}}%
\pgfpathlineto{\pgfqpoint{2.074629in}{2.304234in}}%
\pgfpathlineto{\pgfqpoint{2.073484in}{2.303289in}}%
\pgfpathlineto{\pgfqpoint{2.074643in}{2.304182in}}%
\pgfpathlineto{\pgfqpoint{2.074883in}{2.304440in}}%
\pgfpathlineto{\pgfqpoint{2.075588in}{2.303502in}}%
\pgfpathlineto{\pgfqpoint{2.075590in}{2.303509in}}%
\pgfpathlineto{\pgfqpoint{2.075724in}{2.303441in}}%
\pgfpathlineto{\pgfqpoint{2.076220in}{2.303834in}}%
\pgfpathlineto{\pgfqpoint{2.076701in}{2.303532in}}%
\pgfpathlineto{\pgfqpoint{2.077294in}{2.304090in}}%
\pgfpathlineto{\pgfqpoint{2.076709in}{2.303473in}}%
\pgfpathlineto{\pgfqpoint{2.077826in}{2.303612in}}%
\pgfpathlineto{\pgfqpoint{2.078272in}{2.303241in}}%
\pgfpathlineto{\pgfqpoint{2.078802in}{2.304012in}}%
\pgfpathlineto{\pgfqpoint{2.078910in}{2.303853in}}%
\pgfpathlineto{\pgfqpoint{2.078918in}{2.303934in}}%
\pgfpathlineto{\pgfqpoint{2.079684in}{2.303252in}}%
\pgfpathlineto{\pgfqpoint{2.079997in}{2.303485in}}%
\pgfpathlineto{\pgfqpoint{2.081201in}{2.303428in}}%
\pgfpathlineto{\pgfqpoint{2.080214in}{2.303650in}}%
\pgfpathlineto{\pgfqpoint{2.081242in}{2.303513in}}%
\pgfpathlineto{\pgfqpoint{2.081248in}{2.303544in}}%
\pgfpathlineto{\pgfqpoint{2.082298in}{2.302819in}}%
\pgfpathlineto{\pgfqpoint{2.082323in}{2.302891in}}%
\pgfpathlineto{\pgfqpoint{2.083328in}{2.302744in}}%
\pgfpathlineto{\pgfqpoint{2.082524in}{2.302478in}}%
\pgfpathlineto{\pgfqpoint{2.083396in}{2.302597in}}%
\pgfpathlineto{\pgfqpoint{2.084237in}{2.302096in}}%
\pgfpathlineto{\pgfqpoint{2.084519in}{2.302280in}}%
\pgfpathlineto{\pgfqpoint{2.086053in}{2.302527in}}%
\pgfpathlineto{\pgfqpoint{2.084524in}{2.302331in}}%
\pgfpathlineto{\pgfqpoint{2.086058in}{2.302577in}}%
\pgfpathlineto{\pgfqpoint{2.086986in}{2.303021in}}%
\pgfpathlineto{\pgfqpoint{2.086315in}{2.302488in}}%
\pgfpathlineto{\pgfqpoint{2.087195in}{2.302799in}}%
\pgfpathlineto{\pgfqpoint{2.088421in}{2.302807in}}%
\pgfpathlineto{\pgfqpoint{2.087335in}{2.302930in}}%
\pgfpathlineto{\pgfqpoint{2.088444in}{2.302848in}}%
\pgfpathlineto{\pgfqpoint{2.088447in}{2.302861in}}%
\pgfpathlineto{\pgfqpoint{2.089000in}{2.302483in}}%
\pgfpathlineto{\pgfqpoint{2.089543in}{2.302645in}}%
\pgfpathlineto{\pgfqpoint{2.090662in}{2.302265in}}%
\pgfpathlineto{\pgfqpoint{2.089573in}{2.302696in}}%
\pgfpathlineto{\pgfqpoint{2.090701in}{2.302348in}}%
\pgfpathlineto{\pgfqpoint{2.090710in}{2.302368in}}%
\pgfpathlineto{\pgfqpoint{2.091599in}{2.301697in}}%
\pgfpathlineto{\pgfqpoint{2.091781in}{2.301914in}}%
\pgfpathlineto{\pgfqpoint{2.092904in}{2.301488in}}%
\pgfpathlineto{\pgfqpoint{2.091802in}{2.301977in}}%
\pgfpathlineto{\pgfqpoint{2.092946in}{2.301588in}}%
\pgfpathlineto{\pgfqpoint{2.094439in}{2.303438in}}%
\pgfpathlineto{\pgfqpoint{2.092987in}{2.301443in}}%
\pgfpathlineto{\pgfqpoint{2.094458in}{2.303414in}}%
\pgfpathlineto{\pgfqpoint{2.095539in}{2.303375in}}%
\pgfpathlineto{\pgfqpoint{2.094506in}{2.303443in}}%
\pgfpathlineto{\pgfqpoint{2.095570in}{2.303411in}}%
\pgfpathlineto{\pgfqpoint{2.096668in}{2.303446in}}%
\pgfpathlineto{\pgfqpoint{2.095877in}{2.303349in}}%
\pgfpathlineto{\pgfqpoint{2.096681in}{2.303414in}}%
\pgfpathlineto{\pgfqpoint{2.097814in}{2.303390in}}%
\pgfpathlineto{\pgfqpoint{2.096684in}{2.303444in}}%
\pgfpathlineto{\pgfqpoint{2.097815in}{2.303391in}}%
\pgfpathlineto{\pgfqpoint{2.098881in}{2.303449in}}%
\pgfpathlineto{\pgfqpoint{2.098361in}{2.303341in}}%
\pgfpathlineto{\pgfqpoint{2.098928in}{2.303394in}}%
\pgfpathlineto{\pgfqpoint{2.100042in}{2.303442in}}%
\pgfpathlineto{\pgfqpoint{2.099055in}{2.303361in}}%
\pgfpathlineto{\pgfqpoint{2.100044in}{2.303428in}}%
\pgfpathlineto{\pgfqpoint{2.100788in}{2.303356in}}%
\pgfpathlineto{\pgfqpoint{2.100073in}{2.303448in}}%
\pgfpathlineto{\pgfqpoint{2.101156in}{2.303387in}}%
\pgfpathlineto{\pgfqpoint{2.102271in}{2.303441in}}%
\pgfpathlineto{\pgfqpoint{2.102282in}{2.303408in}}%
\pgfpathlineto{\pgfqpoint{2.102975in}{2.303365in}}%
\pgfpathlineto{\pgfqpoint{2.102285in}{2.303447in}}%
\pgfpathlineto{\pgfqpoint{2.103393in}{2.303407in}}%
\pgfpathlineto{\pgfqpoint{2.104327in}{2.303448in}}%
\pgfpathlineto{\pgfqpoint{2.103624in}{2.303358in}}%
\pgfpathlineto{\pgfqpoint{2.104506in}{2.303428in}}%
\pgfpathlineto{\pgfqpoint{2.105138in}{2.303356in}}%
\pgfpathlineto{\pgfqpoint{2.104556in}{2.303449in}}%
\pgfpathlineto{\pgfqpoint{2.105618in}{2.303409in}}%
\pgfpathlineto{\pgfqpoint{2.106705in}{2.303455in}}%
\pgfpathlineto{\pgfqpoint{2.105623in}{2.303382in}}%
\pgfpathlineto{\pgfqpoint{2.106732in}{2.303426in}}%
\pgfpathlineto{\pgfqpoint{2.107833in}{2.303374in}}%
\pgfpathlineto{\pgfqpoint{2.106748in}{2.303442in}}%
\pgfpathlineto{\pgfqpoint{2.107844in}{2.303410in}}%
\pgfpathlineto{\pgfqpoint{2.109177in}{2.303378in}}%
\pgfpathlineto{\pgfqpoint{2.107848in}{2.303447in}}%
\pgfpathlineto{\pgfqpoint{2.109179in}{2.303421in}}%
\pgfpathlineto{\pgfqpoint{2.110179in}{2.303453in}}%
\pgfpathlineto{\pgfqpoint{2.109310in}{2.303368in}}%
\pgfpathlineto{\pgfqpoint{2.110290in}{2.303397in}}%
\pgfpathlineto{\pgfqpoint{2.111544in}{2.303399in}}%
\pgfpathlineto{\pgfqpoint{2.110294in}{2.303438in}}%
\pgfpathlineto{\pgfqpoint{2.111545in}{2.303415in}}%
\pgfpathlineto{\pgfqpoint{2.112677in}{2.303970in}}%
\pgfpathlineto{\pgfqpoint{2.111721in}{2.303367in}}%
\pgfpathlineto{\pgfqpoint{2.112679in}{2.303955in}}%
\pgfpathlineto{\pgfqpoint{2.113154in}{2.303576in}}%
\pgfpathlineto{\pgfqpoint{2.113768in}{2.304186in}}%
\pgfpathlineto{\pgfqpoint{2.113782in}{2.304081in}}%
\pgfpathlineto{\pgfqpoint{2.114816in}{2.305019in}}%
\pgfpathlineto{\pgfqpoint{2.113810in}{2.304064in}}%
\pgfpathlineto{\pgfqpoint{2.114964in}{2.304834in}}%
\pgfpathlineto{\pgfqpoint{2.116082in}{2.305034in}}%
\pgfpathlineto{\pgfqpoint{2.115245in}{2.304525in}}%
\pgfpathlineto{\pgfqpoint{2.116146in}{2.304823in}}%
\pgfpathlineto{\pgfqpoint{2.116148in}{2.304813in}}%
\pgfpathlineto{\pgfqpoint{2.116930in}{2.305406in}}%
\pgfpathlineto{\pgfqpoint{2.117209in}{2.305159in}}%
\pgfpathlineto{\pgfqpoint{2.118011in}{2.305240in}}%
\pgfpathlineto{\pgfqpoint{2.117811in}{2.304990in}}%
\pgfpathlineto{\pgfqpoint{2.118316in}{2.305046in}}%
\pgfpathlineto{\pgfqpoint{2.119029in}{2.304756in}}%
\pgfpathlineto{\pgfqpoint{2.118565in}{2.305101in}}%
\pgfpathlineto{\pgfqpoint{2.119436in}{2.304977in}}%
\pgfpathlineto{\pgfqpoint{2.120483in}{2.305633in}}%
\pgfpathlineto{\pgfqpoint{2.119674in}{2.304740in}}%
\pgfpathlineto{\pgfqpoint{2.120584in}{2.305411in}}%
\pgfpathlineto{\pgfqpoint{2.121701in}{2.305246in}}%
\pgfpathlineto{\pgfqpoint{2.120864in}{2.305710in}}%
\pgfpathlineto{\pgfqpoint{2.121704in}{2.305270in}}%
\pgfpathlineto{\pgfqpoint{2.121918in}{2.305354in}}%
\pgfpathlineto{\pgfqpoint{2.122747in}{2.304783in}}%
\pgfpathlineto{\pgfqpoint{2.123766in}{2.304548in}}%
\pgfpathlineto{\pgfqpoint{2.122979in}{2.304943in}}%
\pgfpathlineto{\pgfqpoint{2.123863in}{2.304762in}}%
\pgfpathlineto{\pgfqpoint{2.124963in}{2.304837in}}%
\pgfpathlineto{\pgfqpoint{2.124782in}{2.304383in}}%
\pgfpathlineto{\pgfqpoint{2.124978in}{2.304753in}}%
\pgfpathlineto{\pgfqpoint{2.125743in}{2.303994in}}%
\pgfpathlineto{\pgfqpoint{2.126116in}{2.304305in}}%
\pgfpathlineto{\pgfqpoint{2.126457in}{2.304205in}}%
\pgfpathlineto{\pgfqpoint{2.126890in}{2.304704in}}%
\pgfpathlineto{\pgfqpoint{2.127211in}{2.304579in}}%
\pgfpathlineto{\pgfqpoint{2.128253in}{2.305340in}}%
\pgfpathlineto{\pgfqpoint{2.127368in}{2.304486in}}%
\pgfpathlineto{\pgfqpoint{2.128399in}{2.305112in}}%
\pgfpathlineto{\pgfqpoint{2.128587in}{2.304923in}}%
\pgfpathlineto{\pgfqpoint{2.129118in}{2.305506in}}%
\pgfpathlineto{\pgfqpoint{2.129504in}{2.305228in}}%
\pgfpathlineto{\pgfqpoint{2.130910in}{2.305005in}}%
\pgfpathlineto{\pgfqpoint{2.129518in}{2.305285in}}%
\pgfpathlineto{\pgfqpoint{2.130983in}{2.305097in}}%
\pgfpathlineto{\pgfqpoint{2.132048in}{2.305800in}}%
\pgfpathlineto{\pgfqpoint{2.132111in}{2.305721in}}%
\pgfpathlineto{\pgfqpoint{2.132269in}{2.305951in}}%
\pgfpathlineto{\pgfqpoint{2.133126in}{2.304782in}}%
\pgfpathlineto{\pgfqpoint{2.133156in}{2.304874in}}%
\pgfpathlineto{\pgfqpoint{2.134289in}{2.304540in}}%
\pgfpathlineto{\pgfqpoint{2.133197in}{2.304984in}}%
\pgfpathlineto{\pgfqpoint{2.134290in}{2.304549in}}%
\pgfpathlineto{\pgfqpoint{2.135742in}{2.303930in}}%
\pgfpathlineto{\pgfqpoint{2.134312in}{2.304625in}}%
\pgfpathlineto{\pgfqpoint{2.135913in}{2.304129in}}%
\pgfpathlineto{\pgfqpoint{2.137140in}{2.305103in}}%
\pgfpathlineto{\pgfqpoint{2.135916in}{2.304120in}}%
\pgfpathlineto{\pgfqpoint{2.137226in}{2.305036in}}%
\pgfpathlineto{\pgfqpoint{2.137805in}{2.304817in}}%
\pgfpathlineto{\pgfqpoint{2.138321in}{2.305417in}}%
\pgfpathlineto{\pgfqpoint{2.139025in}{2.305681in}}%
\pgfpathlineto{\pgfqpoint{2.138909in}{2.305249in}}%
\pgfpathlineto{\pgfqpoint{2.139432in}{2.305391in}}%
\pgfpathlineto{\pgfqpoint{2.139953in}{2.305174in}}%
\pgfpathlineto{\pgfqpoint{2.140241in}{2.305640in}}%
\pgfpathlineto{\pgfqpoint{2.140526in}{2.305585in}}%
\pgfpathlineto{\pgfqpoint{2.142376in}{2.305310in}}%
\pgfpathlineto{\pgfqpoint{2.142407in}{2.305199in}}%
\pgfpathlineto{\pgfqpoint{2.143534in}{2.304373in}}%
\pgfpathlineto{\pgfqpoint{2.142457in}{2.305343in}}%
\pgfpathlineto{\pgfqpoint{2.143552in}{2.304436in}}%
\pgfpathlineto{\pgfqpoint{2.143645in}{2.304378in}}%
\pgfpathlineto{\pgfqpoint{2.144599in}{2.305194in}}%
\pgfpathlineto{\pgfqpoint{2.145254in}{2.305895in}}%
\pgfpathlineto{\pgfqpoint{2.144618in}{2.305144in}}%
\pgfpathlineto{\pgfqpoint{2.145722in}{2.305387in}}%
\pgfpathlineto{\pgfqpoint{2.145725in}{2.305372in}}%
\pgfpathlineto{\pgfqpoint{2.146022in}{2.305941in}}%
\pgfpathlineto{\pgfqpoint{2.146778in}{2.305925in}}%
\pgfpathlineto{\pgfqpoint{2.148189in}{2.306432in}}%
\pgfpathlineto{\pgfqpoint{2.147319in}{2.305595in}}%
\pgfpathlineto{\pgfqpoint{2.148207in}{2.306374in}}%
\pgfpathlineto{\pgfqpoint{2.149126in}{2.305864in}}%
\pgfpathlineto{\pgfqpoint{2.148291in}{2.306442in}}%
\pgfpathlineto{\pgfqpoint{2.149360in}{2.306118in}}%
\pgfpathlineto{\pgfqpoint{2.149619in}{2.306606in}}%
\pgfpathlineto{\pgfqpoint{2.149377in}{2.306109in}}%
\pgfpathlineto{\pgfqpoint{2.150471in}{2.306090in}}%
\pgfpathlineto{\pgfqpoint{2.150505in}{2.306006in}}%
\pgfpathlineto{\pgfqpoint{2.151121in}{2.306601in}}%
\pgfpathlineto{\pgfqpoint{2.151552in}{2.306521in}}%
\pgfpathlineto{\pgfqpoint{2.152641in}{2.306586in}}%
\pgfpathlineto{\pgfqpoint{2.151904in}{2.306095in}}%
\pgfpathlineto{\pgfqpoint{2.152665in}{2.306524in}}%
\pgfpathlineto{\pgfqpoint{2.153322in}{2.306550in}}%
\pgfpathlineto{\pgfqpoint{2.153553in}{2.306896in}}%
\pgfpathlineto{\pgfqpoint{2.153761in}{2.306713in}}%
\pgfpathlineto{\pgfqpoint{2.153999in}{2.306948in}}%
\pgfpathlineto{\pgfqpoint{2.154592in}{2.306356in}}%
\pgfpathlineto{\pgfqpoint{2.154865in}{2.306496in}}%
\pgfpathlineto{\pgfqpoint{2.154928in}{2.306538in}}%
\pgfpathlineto{\pgfqpoint{2.155637in}{2.305951in}}%
\pgfpathlineto{\pgfqpoint{2.155918in}{2.306024in}}%
\pgfpathlineto{\pgfqpoint{2.156915in}{2.305881in}}%
\pgfpathlineto{\pgfqpoint{2.156311in}{2.306452in}}%
\pgfpathlineto{\pgfqpoint{2.157028in}{2.306075in}}%
\pgfpathlineto{\pgfqpoint{2.157741in}{2.306093in}}%
\pgfpathlineto{\pgfqpoint{2.158019in}{2.306765in}}%
\pgfpathlineto{\pgfqpoint{2.158034in}{2.306714in}}%
\pgfpathlineto{\pgfqpoint{2.158052in}{2.306751in}}%
\pgfpathlineto{\pgfqpoint{2.158568in}{2.306323in}}%
\pgfpathlineto{\pgfqpoint{2.159128in}{2.306401in}}%
\pgfpathlineto{\pgfqpoint{2.159178in}{2.306335in}}%
\pgfpathlineto{\pgfqpoint{2.159902in}{2.306821in}}%
\pgfpathlineto{\pgfqpoint{2.160237in}{2.306451in}}%
\pgfpathlineto{\pgfqpoint{2.161272in}{2.305976in}}%
\pgfpathlineto{\pgfqpoint{2.161585in}{2.306676in}}%
\pgfpathlineto{\pgfqpoint{2.162181in}{2.307157in}}%
\pgfpathlineto{\pgfqpoint{2.161775in}{2.306620in}}%
\pgfpathlineto{\pgfqpoint{2.162708in}{2.306850in}}%
\pgfpathlineto{\pgfqpoint{2.163892in}{2.306697in}}%
\pgfpathlineto{\pgfqpoint{2.163101in}{2.307026in}}%
\pgfpathlineto{\pgfqpoint{2.163896in}{2.306700in}}%
\pgfpathlineto{\pgfqpoint{2.165014in}{2.307120in}}%
\pgfpathlineto{\pgfqpoint{2.164011in}{2.306590in}}%
\pgfpathlineto{\pgfqpoint{2.165023in}{2.307086in}}%
\pgfpathlineto{\pgfqpoint{2.166078in}{2.306731in}}%
\pgfpathlineto{\pgfqpoint{2.165303in}{2.307321in}}%
\pgfpathlineto{\pgfqpoint{2.166172in}{2.306835in}}%
\pgfpathlineto{\pgfqpoint{2.167672in}{2.306337in}}%
\pgfpathlineto{\pgfqpoint{2.166207in}{2.306680in}}%
\pgfpathlineto{\pgfqpoint{2.167679in}{2.306291in}}%
\pgfpathlineto{\pgfqpoint{2.168836in}{2.303365in}}%
\pgfpathlineto{\pgfqpoint{2.167932in}{2.306472in}}%
\pgfpathlineto{\pgfqpoint{2.169054in}{2.303404in}}%
\pgfpathlineto{\pgfqpoint{2.170051in}{2.303458in}}%
\pgfpathlineto{\pgfqpoint{2.169624in}{2.303343in}}%
\pgfpathlineto{\pgfqpoint{2.170166in}{2.303395in}}%
\pgfpathlineto{\pgfqpoint{2.171375in}{2.303444in}}%
\pgfpathlineto{\pgfqpoint{2.170176in}{2.303373in}}%
\pgfpathlineto{\pgfqpoint{2.171389in}{2.303412in}}%
\pgfpathlineto{\pgfqpoint{2.172480in}{2.303381in}}%
\pgfpathlineto{\pgfqpoint{2.171448in}{2.303448in}}%
\pgfpathlineto{\pgfqpoint{2.172502in}{2.303415in}}%
\pgfpathlineto{\pgfqpoint{2.173494in}{2.303457in}}%
\pgfpathlineto{\pgfqpoint{2.172647in}{2.303351in}}%
\pgfpathlineto{\pgfqpoint{2.173614in}{2.303426in}}%
\pgfpathlineto{\pgfqpoint{2.174589in}{2.303349in}}%
\pgfpathlineto{\pgfqpoint{2.173708in}{2.303445in}}%
\pgfpathlineto{\pgfqpoint{2.174725in}{2.303404in}}%
\pgfpathlineto{\pgfqpoint{2.175820in}{2.303369in}}%
\pgfpathlineto{\pgfqpoint{2.174748in}{2.303438in}}%
\pgfpathlineto{\pgfqpoint{2.175846in}{2.303420in}}%
\pgfpathlineto{\pgfqpoint{2.176919in}{2.303447in}}%
\pgfpathlineto{\pgfqpoint{2.176272in}{2.303361in}}%
\pgfpathlineto{\pgfqpoint{2.176958in}{2.303441in}}%
\pgfpathlineto{\pgfqpoint{2.178016in}{2.303359in}}%
\pgfpathlineto{\pgfqpoint{2.178071in}{2.303397in}}%
\pgfpathlineto{\pgfqpoint{2.179127in}{2.303450in}}%
\pgfpathlineto{\pgfqpoint{2.178153in}{2.303344in}}%
\pgfpathlineto{\pgfqpoint{2.179184in}{2.303407in}}%
\pgfpathlineto{\pgfqpoint{2.180038in}{2.303362in}}%
\pgfpathlineto{\pgfqpoint{2.179296in}{2.303457in}}%
\pgfpathlineto{\pgfqpoint{2.180295in}{2.303417in}}%
\pgfpathlineto{\pgfqpoint{2.180301in}{2.303400in}}%
\pgfpathlineto{\pgfqpoint{2.181732in}{2.303426in}}%
\pgfpathlineto{\pgfqpoint{2.182652in}{2.303302in}}%
\pgfpathlineto{\pgfqpoint{2.181763in}{2.303440in}}%
\pgfpathlineto{\pgfqpoint{2.182843in}{2.303424in}}%
\pgfpathlineto{\pgfqpoint{2.183151in}{2.303335in}}%
\pgfpathlineto{\pgfqpoint{2.182927in}{2.303447in}}%
\pgfpathlineto{\pgfqpoint{2.183955in}{2.303401in}}%
\pgfpathlineto{\pgfqpoint{2.184996in}{2.303449in}}%
\pgfpathlineto{\pgfqpoint{2.184027in}{2.303373in}}%
\pgfpathlineto{\pgfqpoint{2.185066in}{2.303426in}}%
\pgfpathlineto{\pgfqpoint{2.185777in}{2.303360in}}%
\pgfpathlineto{\pgfqpoint{2.185284in}{2.303453in}}%
\pgfpathlineto{\pgfqpoint{2.186179in}{2.303419in}}%
\pgfpathlineto{\pgfqpoint{2.187376in}{2.303095in}}%
\pgfpathlineto{\pgfqpoint{2.186197in}{2.303444in}}%
\pgfpathlineto{\pgfqpoint{2.187394in}{2.303140in}}%
\pgfpathlineto{\pgfqpoint{2.188518in}{2.303362in}}%
\pgfpathlineto{\pgfqpoint{2.187964in}{2.302754in}}%
\pgfpathlineto{\pgfqpoint{2.188523in}{2.303337in}}%
\pgfpathlineto{\pgfqpoint{2.188524in}{2.303340in}}%
\pgfpathlineto{\pgfqpoint{2.188895in}{2.302786in}}%
\pgfpathlineto{\pgfqpoint{2.189548in}{2.302888in}}%
\pgfpathlineto{\pgfqpoint{2.190578in}{2.302737in}}%
\pgfpathlineto{\pgfqpoint{2.189857in}{2.303058in}}%
\pgfpathlineto{\pgfqpoint{2.190654in}{2.302974in}}%
\pgfpathlineto{\pgfqpoint{2.190757in}{2.303080in}}%
\pgfpathlineto{\pgfqpoint{2.191502in}{2.302488in}}%
\pgfpathlineto{\pgfqpoint{2.191745in}{2.302569in}}%
\pgfpathlineto{\pgfqpoint{2.192312in}{2.302311in}}%
\pgfpathlineto{\pgfqpoint{2.191956in}{2.302719in}}%
\pgfpathlineto{\pgfqpoint{2.192873in}{2.302591in}}%
\pgfpathlineto{\pgfqpoint{2.192879in}{2.302625in}}%
\pgfpathlineto{\pgfqpoint{2.193174in}{2.302158in}}%
\pgfpathlineto{\pgfqpoint{2.193972in}{2.302359in}}%
\pgfpathlineto{\pgfqpoint{2.194625in}{2.302079in}}%
\pgfpathlineto{\pgfqpoint{2.194910in}{2.302653in}}%
\pgfpathlineto{\pgfqpoint{2.195082in}{2.302409in}}%
\pgfpathlineto{\pgfqpoint{2.195138in}{2.302552in}}%
\pgfpathlineto{\pgfqpoint{2.196105in}{2.301883in}}%
\pgfpathlineto{\pgfqpoint{2.196157in}{2.301964in}}%
\pgfpathlineto{\pgfqpoint{2.196835in}{2.301713in}}%
\pgfpathlineto{\pgfqpoint{2.196355in}{2.302202in}}%
\pgfpathlineto{\pgfqpoint{2.197267in}{2.302008in}}%
\pgfpathlineto{\pgfqpoint{2.198393in}{2.302270in}}%
\pgfpathlineto{\pgfqpoint{2.197388in}{2.301834in}}%
\pgfpathlineto{\pgfqpoint{2.198394in}{2.302266in}}%
\pgfpathlineto{\pgfqpoint{2.199481in}{2.301749in}}%
\pgfpathlineto{\pgfqpoint{2.198611in}{2.302449in}}%
\pgfpathlineto{\pgfqpoint{2.199627in}{2.301896in}}%
\pgfpathlineto{\pgfqpoint{2.200754in}{2.302116in}}%
\pgfpathlineto{\pgfqpoint{2.199745in}{2.301700in}}%
\pgfpathlineto{\pgfqpoint{2.200755in}{2.302100in}}%
\pgfpathlineto{\pgfqpoint{2.201938in}{2.302524in}}%
\pgfpathlineto{\pgfqpoint{2.200757in}{2.302099in}}%
\pgfpathlineto{\pgfqpoint{2.202001in}{2.302312in}}%
\pgfpathlineto{\pgfqpoint{2.202110in}{2.302224in}}%
\pgfpathlineto{\pgfqpoint{2.202846in}{2.302746in}}%
\pgfpathlineto{\pgfqpoint{2.203100in}{2.302590in}}%
\pgfpathlineto{\pgfqpoint{2.204845in}{2.302501in}}%
\pgfpathlineto{\pgfqpoint{2.203334in}{2.303336in}}%
\pgfpathlineto{\pgfqpoint{2.204866in}{2.302609in}}%
\pgfpathlineto{\pgfqpoint{2.205956in}{2.302881in}}%
\pgfpathlineto{\pgfqpoint{2.205107in}{2.302425in}}%
\pgfpathlineto{\pgfqpoint{2.205985in}{2.302800in}}%
\pgfpathlineto{\pgfqpoint{2.207051in}{2.302784in}}%
\pgfpathlineto{\pgfqpoint{2.206808in}{2.303150in}}%
\pgfpathlineto{\pgfqpoint{2.207089in}{2.302948in}}%
\pgfpathlineto{\pgfqpoint{2.208158in}{2.303252in}}%
\pgfpathlineto{\pgfqpoint{2.207183in}{2.302856in}}%
\pgfpathlineto{\pgfqpoint{2.208220in}{2.303143in}}%
\pgfpathlineto{\pgfqpoint{2.208679in}{2.302959in}}%
\pgfpathlineto{\pgfqpoint{2.209278in}{2.303580in}}%
\pgfpathlineto{\pgfqpoint{2.209290in}{2.303533in}}%
\pgfpathlineto{\pgfqpoint{2.209730in}{2.303877in}}%
\pgfpathlineto{\pgfqpoint{2.210031in}{2.303172in}}%
\pgfpathlineto{\pgfqpoint{2.210372in}{2.303236in}}%
\pgfpathlineto{\pgfqpoint{2.211526in}{2.302872in}}%
\pgfpathlineto{\pgfqpoint{2.210403in}{2.303293in}}%
\pgfpathlineto{\pgfqpoint{2.211541in}{2.302935in}}%
\pgfpathlineto{\pgfqpoint{2.212464in}{2.303761in}}%
\pgfpathlineto{\pgfqpoint{2.211790in}{2.302715in}}%
\pgfpathlineto{\pgfqpoint{2.212711in}{2.303540in}}%
\pgfpathlineto{\pgfqpoint{2.213375in}{2.303301in}}%
\pgfpathlineto{\pgfqpoint{2.212787in}{2.303704in}}%
\pgfpathlineto{\pgfqpoint{2.213824in}{2.303577in}}%
\pgfpathlineto{\pgfqpoint{2.213919in}{2.303659in}}%
\pgfpathlineto{\pgfqpoint{2.214820in}{2.302886in}}%
\pgfpathlineto{\pgfqpoint{2.214919in}{2.303049in}}%
\pgfpathlineto{\pgfqpoint{2.214965in}{2.303015in}}%
\pgfpathlineto{\pgfqpoint{2.215850in}{2.303827in}}%
\pgfpathlineto{\pgfqpoint{2.215922in}{2.303702in}}%
\pgfpathlineto{\pgfqpoint{2.216751in}{2.304111in}}%
\pgfpathlineto{\pgfqpoint{2.217020in}{2.303499in}}%
\pgfpathlineto{\pgfqpoint{2.217031in}{2.303653in}}%
\pgfpathlineto{\pgfqpoint{2.217044in}{2.303563in}}%
\pgfpathlineto{\pgfqpoint{2.217512in}{2.304107in}}%
\pgfpathlineto{\pgfqpoint{2.218139in}{2.303764in}}%
\pgfpathlineto{\pgfqpoint{2.218449in}{2.304268in}}%
\pgfpathlineto{\pgfqpoint{2.218938in}{2.303482in}}%
\pgfpathlineto{\pgfqpoint{2.219240in}{2.303603in}}%
\pgfpathlineto{\pgfqpoint{2.219310in}{2.303568in}}%
\pgfpathlineto{\pgfqpoint{2.220229in}{2.304589in}}%
\pgfpathlineto{\pgfqpoint{2.220281in}{2.304475in}}%
\pgfpathlineto{\pgfqpoint{2.220840in}{2.304290in}}%
\pgfpathlineto{\pgfqpoint{2.221321in}{2.304926in}}%
\pgfpathlineto{\pgfqpoint{2.221338in}{2.304878in}}%
\pgfpathlineto{\pgfqpoint{2.221395in}{2.304988in}}%
\pgfpathlineto{\pgfqpoint{2.221669in}{2.304532in}}%
\pgfpathlineto{\pgfqpoint{2.222448in}{2.304845in}}%
\pgfpathlineto{\pgfqpoint{2.223489in}{2.304124in}}%
\pgfpathlineto{\pgfqpoint{2.223584in}{2.304337in}}%
\pgfpathlineto{\pgfqpoint{2.223845in}{2.304608in}}%
\pgfpathlineto{\pgfqpoint{2.224149in}{2.303943in}}%
\pgfpathlineto{\pgfqpoint{2.224692in}{2.304253in}}%
\pgfpathlineto{\pgfqpoint{2.224804in}{2.304470in}}%
\pgfpathlineto{\pgfqpoint{2.225107in}{2.304022in}}%
\pgfpathlineto{\pgfqpoint{2.225785in}{2.304055in}}%
\pgfpathlineto{\pgfqpoint{2.226609in}{2.303623in}}%
\pgfpathlineto{\pgfqpoint{2.225787in}{2.304065in}}%
\pgfpathlineto{\pgfqpoint{2.226923in}{2.303808in}}%
\pgfpathlineto{\pgfqpoint{2.227034in}{2.303844in}}%
\pgfpathlineto{\pgfqpoint{2.227189in}{2.303502in}}%
\pgfpathlineto{\pgfqpoint{2.228011in}{2.303529in}}%
\pgfpathlineto{\pgfqpoint{2.229145in}{2.303126in}}%
\pgfpathlineto{\pgfqpoint{2.228607in}{2.303699in}}%
\pgfpathlineto{\pgfqpoint{2.229147in}{2.303133in}}%
\pgfpathlineto{\pgfqpoint{2.229325in}{2.302902in}}%
\pgfpathlineto{\pgfqpoint{2.229756in}{2.303416in}}%
\pgfpathlineto{\pgfqpoint{2.230215in}{2.303442in}}%
\pgfpathlineto{\pgfqpoint{2.231296in}{2.303441in}}%
\pgfpathlineto{\pgfqpoint{2.231051in}{2.302906in}}%
\pgfpathlineto{\pgfqpoint{2.231317in}{2.303296in}}%
\pgfpathlineto{\pgfqpoint{2.232470in}{2.303185in}}%
\pgfpathlineto{\pgfqpoint{2.231934in}{2.303663in}}%
\pgfpathlineto{\pgfqpoint{2.232498in}{2.303297in}}%
\pgfpathlineto{\pgfqpoint{2.233488in}{2.303450in}}%
\pgfpathlineto{\pgfqpoint{2.233116in}{2.302913in}}%
\pgfpathlineto{\pgfqpoint{2.233609in}{2.303341in}}%
\pgfpathlineto{\pgfqpoint{2.233668in}{2.303196in}}%
\pgfpathlineto{\pgfqpoint{2.234587in}{2.303714in}}%
\pgfpathlineto{\pgfqpoint{2.234710in}{2.303590in}}%
\pgfpathlineto{\pgfqpoint{2.234760in}{2.303734in}}%
\pgfpathlineto{\pgfqpoint{2.235172in}{2.303225in}}%
\pgfpathlineto{\pgfqpoint{2.235821in}{2.303570in}}%
\pgfpathlineto{\pgfqpoint{2.236458in}{2.303306in}}%
\pgfpathlineto{\pgfqpoint{2.235882in}{2.303663in}}%
\pgfpathlineto{\pgfqpoint{2.236939in}{2.303488in}}%
\pgfpathlineto{\pgfqpoint{2.237969in}{2.304119in}}%
\pgfpathlineto{\pgfqpoint{2.237016in}{2.303283in}}%
\pgfpathlineto{\pgfqpoint{2.238066in}{2.303953in}}%
\pgfpathlineto{\pgfqpoint{2.239073in}{2.303867in}}%
\pgfpathlineto{\pgfqpoint{2.238386in}{2.304225in}}%
\pgfpathlineto{\pgfqpoint{2.239189in}{2.304038in}}%
\pgfpathlineto{\pgfqpoint{2.239347in}{2.304082in}}%
\pgfpathlineto{\pgfqpoint{2.239657in}{2.303669in}}%
\pgfpathlineto{\pgfqpoint{2.240287in}{2.303851in}}%
\pgfpathlineto{\pgfqpoint{2.241471in}{2.304244in}}%
\pgfpathlineto{\pgfqpoint{2.240427in}{2.303712in}}%
\pgfpathlineto{\pgfqpoint{2.241481in}{2.304206in}}%
\pgfpathlineto{\pgfqpoint{2.242544in}{2.304067in}}%
\pgfpathlineto{\pgfqpoint{2.242256in}{2.304482in}}%
\pgfpathlineto{\pgfqpoint{2.242597in}{2.304155in}}%
\pgfpathlineto{\pgfqpoint{2.242745in}{2.304363in}}%
\pgfpathlineto{\pgfqpoint{2.242992in}{2.303360in}}%
\pgfpathlineto{\pgfqpoint{2.243663in}{2.303426in}}%
\pgfpathlineto{\pgfqpoint{2.244699in}{2.303374in}}%
\pgfpathlineto{\pgfqpoint{2.243712in}{2.303451in}}%
\pgfpathlineto{\pgfqpoint{2.244774in}{2.303419in}}%
\pgfpathlineto{\pgfqpoint{2.245738in}{2.303452in}}%
\pgfpathlineto{\pgfqpoint{2.244910in}{2.303367in}}%
\pgfpathlineto{\pgfqpoint{2.245886in}{2.303424in}}%
\pgfpathlineto{\pgfqpoint{2.246944in}{2.303375in}}%
\pgfpathlineto{\pgfqpoint{2.245924in}{2.303445in}}%
\pgfpathlineto{\pgfqpoint{2.246998in}{2.303398in}}%
\pgfpathlineto{\pgfqpoint{2.248069in}{2.303441in}}%
\pgfpathlineto{\pgfqpoint{2.247181in}{2.303332in}}%
\pgfpathlineto{\pgfqpoint{2.248111in}{2.303405in}}%
\pgfpathlineto{\pgfqpoint{2.249300in}{2.303379in}}%
\pgfpathlineto{\pgfqpoint{2.248120in}{2.303439in}}%
\pgfpathlineto{\pgfqpoint{2.249316in}{2.303422in}}%
\pgfpathlineto{\pgfqpoint{2.250424in}{2.303441in}}%
\pgfpathlineto{\pgfqpoint{2.249526in}{2.303340in}}%
\pgfpathlineto{\pgfqpoint{2.250428in}{2.303424in}}%
\pgfpathlineto{\pgfqpoint{2.251277in}{2.303337in}}%
\pgfpathlineto{\pgfqpoint{2.250544in}{2.303448in}}%
\pgfpathlineto{\pgfqpoint{2.251540in}{2.303415in}}%
\pgfpathlineto{\pgfqpoint{2.252746in}{2.303353in}}%
\pgfpathlineto{\pgfqpoint{2.251562in}{2.303438in}}%
\pgfpathlineto{\pgfqpoint{2.252760in}{2.303410in}}%
\pgfpathlineto{\pgfqpoint{2.253820in}{2.303451in}}%
\pgfpathlineto{\pgfqpoint{2.252888in}{2.303363in}}%
\pgfpathlineto{\pgfqpoint{2.253870in}{2.303425in}}%
\pgfpathlineto{\pgfqpoint{2.254822in}{2.303363in}}%
\pgfpathlineto{\pgfqpoint{2.254083in}{2.303451in}}%
\pgfpathlineto{\pgfqpoint{2.254981in}{2.303419in}}%
\pgfpathlineto{\pgfqpoint{2.256065in}{2.303440in}}%
\pgfpathlineto{\pgfqpoint{2.255198in}{2.303357in}}%
\pgfpathlineto{\pgfqpoint{2.256093in}{2.303435in}}%
\pgfpathlineto{\pgfqpoint{2.256887in}{2.303354in}}%
\pgfpathlineto{\pgfqpoint{2.256151in}{2.303446in}}%
\pgfpathlineto{\pgfqpoint{2.257205in}{2.303407in}}%
\pgfpathlineto{\pgfqpoint{2.258385in}{2.303434in}}%
\pgfpathlineto{\pgfqpoint{2.257220in}{2.303359in}}%
\pgfpathlineto{\pgfqpoint{2.258388in}{2.303423in}}%
\pgfpathlineto{\pgfqpoint{2.259165in}{2.303355in}}%
\pgfpathlineto{\pgfqpoint{2.258464in}{2.303446in}}%
\pgfpathlineto{\pgfqpoint{2.259500in}{2.303419in}}%
\pgfpathlineto{\pgfqpoint{2.260619in}{2.303369in}}%
\pgfpathlineto{\pgfqpoint{2.259545in}{2.303446in}}%
\pgfpathlineto{\pgfqpoint{2.260623in}{2.303386in}}%
\pgfpathlineto{\pgfqpoint{2.261731in}{2.303541in}}%
\pgfpathlineto{\pgfqpoint{2.261645in}{2.303292in}}%
\pgfpathlineto{\pgfqpoint{2.261740in}{2.303506in}}%
\pgfpathlineto{\pgfqpoint{2.264555in}{2.302573in}}%
\pgfpathlineto{\pgfqpoint{2.265666in}{2.302314in}}%
\pgfpathlineto{\pgfqpoint{2.265171in}{2.302942in}}%
\pgfpathlineto{\pgfqpoint{2.265673in}{2.302344in}}%
\pgfpathlineto{\pgfqpoint{2.265848in}{2.302441in}}%
\pgfpathlineto{\pgfqpoint{2.266740in}{2.301634in}}%
\pgfpathlineto{\pgfqpoint{2.267851in}{2.300961in}}%
\pgfpathlineto{\pgfqpoint{2.267140in}{2.301844in}}%
\pgfpathlineto{\pgfqpoint{2.267883in}{2.301035in}}%
\pgfpathlineto{\pgfqpoint{2.267931in}{2.301058in}}%
\pgfpathlineto{\pgfqpoint{2.268608in}{2.300629in}}%
\pgfpathlineto{\pgfqpoint{2.268979in}{2.300850in}}%
\pgfpathlineto{\pgfqpoint{2.269708in}{2.300431in}}%
\pgfpathlineto{\pgfqpoint{2.268981in}{2.300853in}}%
\pgfpathlineto{\pgfqpoint{2.270115in}{2.300492in}}%
\pgfpathlineto{\pgfqpoint{2.270260in}{2.300671in}}%
\pgfpathlineto{\pgfqpoint{2.270483in}{2.300229in}}%
\pgfpathlineto{\pgfqpoint{2.271223in}{2.300406in}}%
\pgfpathlineto{\pgfqpoint{2.272087in}{2.300677in}}%
\pgfpathlineto{\pgfqpoint{2.271358in}{2.300224in}}%
\pgfpathlineto{\pgfqpoint{2.272352in}{2.300436in}}%
\pgfpathlineto{\pgfqpoint{2.273282in}{2.300177in}}%
\pgfpathlineto{\pgfqpoint{2.272906in}{2.300608in}}%
\pgfpathlineto{\pgfqpoint{2.273464in}{2.300416in}}%
\pgfpathlineto{\pgfqpoint{2.274604in}{2.301031in}}%
\pgfpathlineto{\pgfqpoint{2.274611in}{2.301011in}}%
\pgfpathlineto{\pgfqpoint{2.275201in}{2.301184in}}%
\pgfpathlineto{\pgfqpoint{2.274679in}{2.300841in}}%
\pgfpathlineto{\pgfqpoint{2.275716in}{2.300730in}}%
\pgfpathlineto{\pgfqpoint{2.276261in}{2.300209in}}%
\pgfpathlineto{\pgfqpoint{2.275759in}{2.300757in}}%
\pgfpathlineto{\pgfqpoint{2.276845in}{2.300456in}}%
\pgfpathlineto{\pgfqpoint{2.277065in}{2.300363in}}%
\pgfpathlineto{\pgfqpoint{2.277602in}{2.300874in}}%
\pgfpathlineto{\pgfqpoint{2.277960in}{2.300474in}}%
\pgfpathlineto{\pgfqpoint{2.278415in}{2.300638in}}%
\pgfpathlineto{\pgfqpoint{2.278742in}{2.300118in}}%
\pgfpathlineto{\pgfqpoint{2.279041in}{2.300179in}}%
\pgfpathlineto{\pgfqpoint{2.279284in}{2.299873in}}%
\pgfpathlineto{\pgfqpoint{2.279847in}{2.300536in}}%
\pgfpathlineto{\pgfqpoint{2.280149in}{2.300242in}}%
\pgfpathlineto{\pgfqpoint{2.281252in}{2.300341in}}%
\pgfpathlineto{\pgfqpoint{2.281000in}{2.299741in}}%
\pgfpathlineto{\pgfqpoint{2.281266in}{2.300326in}}%
\pgfpathlineto{\pgfqpoint{2.282288in}{2.299673in}}%
\pgfpathlineto{\pgfqpoint{2.281321in}{2.300387in}}%
\pgfpathlineto{\pgfqpoint{2.282443in}{2.299868in}}%
\pgfpathlineto{\pgfqpoint{2.283440in}{2.299823in}}%
\pgfpathlineto{\pgfqpoint{2.282805in}{2.299611in}}%
\pgfpathlineto{\pgfqpoint{2.283513in}{2.299565in}}%
\pgfpathlineto{\pgfqpoint{2.284225in}{2.299711in}}%
\pgfpathlineto{\pgfqpoint{2.284643in}{2.299112in}}%
\pgfpathlineto{\pgfqpoint{2.285780in}{2.299724in}}%
\pgfpathlineto{\pgfqpoint{2.284682in}{2.299061in}}%
\pgfpathlineto{\pgfqpoint{2.285852in}{2.299675in}}%
\pgfpathlineto{\pgfqpoint{2.286254in}{2.299522in}}%
\pgfpathlineto{\pgfqpoint{2.286840in}{2.300174in}}%
\pgfpathlineto{\pgfqpoint{2.286926in}{2.300117in}}%
\pgfpathlineto{\pgfqpoint{2.287515in}{2.300467in}}%
\pgfpathlineto{\pgfqpoint{2.286932in}{2.300084in}}%
\pgfpathlineto{\pgfqpoint{2.288062in}{2.300296in}}%
\pgfpathlineto{\pgfqpoint{2.288663in}{2.299839in}}%
\pgfpathlineto{\pgfqpoint{2.288231in}{2.300423in}}%
\pgfpathlineto{\pgfqpoint{2.289211in}{2.300205in}}%
\pgfpathlineto{\pgfqpoint{2.290321in}{2.300727in}}%
\pgfpathlineto{\pgfqpoint{2.290333in}{2.300662in}}%
\pgfpathlineto{\pgfqpoint{2.291340in}{2.300339in}}%
\pgfpathlineto{\pgfqpoint{2.291458in}{2.300610in}}%
\pgfpathlineto{\pgfqpoint{2.291627in}{2.300846in}}%
\pgfpathlineto{\pgfqpoint{2.291926in}{2.300243in}}%
\pgfpathlineto{\pgfqpoint{2.292546in}{2.300241in}}%
\pgfpathlineto{\pgfqpoint{2.292848in}{2.300189in}}%
\pgfpathlineto{\pgfqpoint{2.293555in}{2.300828in}}%
\pgfpathlineto{\pgfqpoint{2.293627in}{2.300745in}}%
\pgfpathlineto{\pgfqpoint{2.294448in}{2.300949in}}%
\pgfpathlineto{\pgfqpoint{2.293869in}{2.300546in}}%
\pgfpathlineto{\pgfqpoint{2.294713in}{2.300509in}}%
\pgfpathlineto{\pgfqpoint{2.295287in}{2.300137in}}%
\pgfpathlineto{\pgfqpoint{2.294770in}{2.300676in}}%
\pgfpathlineto{\pgfqpoint{2.295830in}{2.300414in}}%
\pgfpathlineto{\pgfqpoint{2.296960in}{2.300636in}}%
\pgfpathlineto{\pgfqpoint{2.295896in}{2.300327in}}%
\pgfpathlineto{\pgfqpoint{2.296961in}{2.300632in}}%
\pgfpathlineto{\pgfqpoint{2.296969in}{2.300579in}}%
\pgfpathlineto{\pgfqpoint{2.297865in}{2.301444in}}%
\pgfpathlineto{\pgfqpoint{2.297995in}{2.301268in}}%
\pgfpathlineto{\pgfqpoint{2.299088in}{2.301476in}}%
\pgfpathlineto{\pgfqpoint{2.298290in}{2.301020in}}%
\pgfpathlineto{\pgfqpoint{2.299118in}{2.301439in}}%
\pgfpathlineto{\pgfqpoint{2.300079in}{2.301119in}}%
\pgfpathlineto{\pgfqpoint{2.299694in}{2.301652in}}%
\pgfpathlineto{\pgfqpoint{2.300240in}{2.301310in}}%
\pgfpathlineto{\pgfqpoint{2.300352in}{2.301451in}}%
\pgfpathlineto{\pgfqpoint{2.300695in}{2.300871in}}%
\pgfpathlineto{\pgfqpoint{2.301304in}{2.300929in}}%
\pgfpathlineto{\pgfqpoint{2.302140in}{2.300648in}}%
\pgfpathlineto{\pgfqpoint{2.301608in}{2.301321in}}%
\pgfpathlineto{\pgfqpoint{2.302424in}{2.300879in}}%
\pgfpathlineto{\pgfqpoint{2.302920in}{2.300429in}}%
\pgfpathlineto{\pgfqpoint{2.302498in}{2.301011in}}%
\pgfpathlineto{\pgfqpoint{2.303525in}{2.301046in}}%
\pgfpathlineto{\pgfqpoint{2.303733in}{2.301346in}}%
\pgfpathlineto{\pgfqpoint{2.304586in}{2.300898in}}%
\pgfpathlineto{\pgfqpoint{2.304631in}{2.300922in}}%
\pgfpathlineto{\pgfqpoint{2.305043in}{2.300501in}}%
\pgfpathlineto{\pgfqpoint{2.305751in}{2.300806in}}%
\pgfpathlineto{\pgfqpoint{2.306021in}{2.300412in}}%
\pgfpathlineto{\pgfqpoint{2.305754in}{2.300825in}}%
\pgfpathlineto{\pgfqpoint{2.306815in}{2.301114in}}%
\pgfpathlineto{\pgfqpoint{2.307602in}{2.301408in}}%
\pgfpathlineto{\pgfqpoint{2.307166in}{2.300837in}}%
\pgfpathlineto{\pgfqpoint{2.307930in}{2.301198in}}%
\pgfpathlineto{\pgfqpoint{2.307930in}{2.301187in}}%
\pgfpathlineto{\pgfqpoint{2.308864in}{2.302043in}}%
\pgfpathlineto{\pgfqpoint{2.308998in}{2.301844in}}%
\pgfpathlineto{\pgfqpoint{2.310105in}{2.302337in}}%
\pgfpathlineto{\pgfqpoint{2.309007in}{2.301824in}}%
\pgfpathlineto{\pgfqpoint{2.310197in}{2.302266in}}%
\pgfpathlineto{\pgfqpoint{2.310917in}{2.301962in}}%
\pgfpathlineto{\pgfqpoint{2.311234in}{2.302443in}}%
\pgfpathlineto{\pgfqpoint{2.311308in}{2.302266in}}%
\pgfpathlineto{\pgfqpoint{2.312175in}{2.301970in}}%
\pgfpathlineto{\pgfqpoint{2.311531in}{2.302581in}}%
\pgfpathlineto{\pgfqpoint{2.312448in}{2.302166in}}%
\pgfpathlineto{\pgfqpoint{2.313507in}{2.302691in}}%
\pgfpathlineto{\pgfqpoint{2.312529in}{2.302087in}}%
\pgfpathlineto{\pgfqpoint{2.313610in}{2.302610in}}%
\pgfpathlineto{\pgfqpoint{2.314686in}{2.302461in}}%
\pgfpathlineto{\pgfqpoint{2.313955in}{2.303038in}}%
\pgfpathlineto{\pgfqpoint{2.314732in}{2.302558in}}%
\pgfpathlineto{\pgfqpoint{2.316294in}{2.303070in}}%
\pgfpathlineto{\pgfqpoint{2.314827in}{2.302446in}}%
\pgfpathlineto{\pgfqpoint{2.316330in}{2.303007in}}%
\pgfpathlineto{\pgfqpoint{2.316691in}{2.302688in}}%
\pgfpathlineto{\pgfqpoint{2.317266in}{2.303436in}}%
\pgfpathlineto{\pgfqpoint{2.317419in}{2.303355in}}%
\pgfpathlineto{\pgfqpoint{2.317704in}{2.303246in}}%
\pgfpathlineto{\pgfqpoint{2.317644in}{2.303439in}}%
\pgfpathlineto{\pgfqpoint{2.318530in}{2.303424in}}%
\pgfpathlineto{\pgfqpoint{2.319682in}{2.303399in}}%
\pgfpathlineto{\pgfqpoint{2.318534in}{2.303440in}}%
\pgfpathlineto{\pgfqpoint{2.319684in}{2.303418in}}%
\pgfpathlineto{\pgfqpoint{2.320818in}{2.303437in}}%
\pgfpathlineto{\pgfqpoint{2.319696in}{2.303388in}}%
\pgfpathlineto{\pgfqpoint{2.320821in}{2.303423in}}%
\pgfpathlineto{\pgfqpoint{2.321853in}{2.303376in}}%
\pgfpathlineto{\pgfqpoint{2.320875in}{2.303444in}}%
\pgfpathlineto{\pgfqpoint{2.321939in}{2.303414in}}%
\pgfpathlineto{\pgfqpoint{2.322990in}{2.303451in}}%
\pgfpathlineto{\pgfqpoint{2.322045in}{2.303364in}}%
\pgfpathlineto{\pgfqpoint{2.323052in}{2.303418in}}%
\pgfpathlineto{\pgfqpoint{2.324207in}{2.303440in}}%
\pgfpathlineto{\pgfqpoint{2.323108in}{2.303374in}}%
\pgfpathlineto{\pgfqpoint{2.324217in}{2.303380in}}%
\pgfpathlineto{\pgfqpoint{2.325292in}{2.303378in}}%
\pgfpathlineto{\pgfqpoint{2.324280in}{2.303450in}}%
\pgfpathlineto{\pgfqpoint{2.325335in}{2.303416in}}%
\pgfpathlineto{\pgfqpoint{2.326389in}{2.303443in}}%
\pgfpathlineto{\pgfqpoint{2.325801in}{2.303327in}}%
\pgfpathlineto{\pgfqpoint{2.326448in}{2.303428in}}%
\pgfpathlineto{\pgfqpoint{2.327460in}{2.303337in}}%
\pgfpathlineto{\pgfqpoint{2.326662in}{2.303457in}}%
\pgfpathlineto{\pgfqpoint{2.327559in}{2.303430in}}%
\pgfpathlineto{\pgfqpoint{2.328430in}{2.303362in}}%
\pgfpathlineto{\pgfqpoint{2.327633in}{2.303453in}}%
\pgfpathlineto{\pgfqpoint{2.328672in}{2.303419in}}%
\pgfpathlineto{\pgfqpoint{2.329729in}{2.303443in}}%
\pgfpathlineto{\pgfqpoint{2.328726in}{2.303363in}}%
\pgfpathlineto{\pgfqpoint{2.329784in}{2.303430in}}%
\pgfpathlineto{\pgfqpoint{2.330803in}{2.303380in}}%
\pgfpathlineto{\pgfqpoint{2.329836in}{2.303449in}}%
\pgfpathlineto{\pgfqpoint{2.330895in}{2.303416in}}%
\pgfpathlineto{\pgfqpoint{2.332039in}{2.303365in}}%
\pgfpathlineto{\pgfqpoint{2.330947in}{2.303450in}}%
\pgfpathlineto{\pgfqpoint{2.332042in}{2.303397in}}%
\pgfpathlineto{\pgfqpoint{2.333014in}{2.303449in}}%
\pgfpathlineto{\pgfqpoint{2.332230in}{2.303371in}}%
\pgfpathlineto{\pgfqpoint{2.333156in}{2.303434in}}%
\pgfpathlineto{\pgfqpoint{2.333676in}{2.303354in}}%
\pgfpathlineto{\pgfqpoint{2.333349in}{2.303452in}}%
\pgfpathlineto{\pgfqpoint{2.334268in}{2.303412in}}%
\pgfpathlineto{\pgfqpoint{2.335367in}{2.303443in}}%
\pgfpathlineto{\pgfqpoint{2.335231in}{2.303322in}}%
\pgfpathlineto{\pgfqpoint{2.335380in}{2.303428in}}%
\pgfpathlineto{\pgfqpoint{2.336487in}{2.303303in}}%
\pgfpathlineto{\pgfqpoint{2.335517in}{2.303446in}}%
\pgfpathlineto{\pgfqpoint{2.336494in}{2.303332in}}%
\pgfpathlineto{\pgfqpoint{2.337074in}{2.302822in}}%
\pgfpathlineto{\pgfqpoint{2.336511in}{2.303425in}}%
\pgfpathlineto{\pgfqpoint{2.337623in}{2.303175in}}%
\pgfpathlineto{\pgfqpoint{2.338549in}{2.303422in}}%
\pgfpathlineto{\pgfqpoint{2.337767in}{2.303035in}}%
\pgfpathlineto{\pgfqpoint{2.338734in}{2.303197in}}%
\pgfpathlineto{\pgfqpoint{2.339705in}{2.303314in}}%
\pgfpathlineto{\pgfqpoint{2.339084in}{2.302922in}}%
\pgfpathlineto{\pgfqpoint{2.339868in}{2.303228in}}%
\pgfpathlineto{\pgfqpoint{2.339940in}{2.303123in}}%
\pgfpathlineto{\pgfqpoint{2.340763in}{2.303842in}}%
\pgfpathlineto{\pgfqpoint{2.340956in}{2.303540in}}%
\pgfpathlineto{\pgfqpoint{2.342018in}{2.304141in}}%
\pgfpathlineto{\pgfqpoint{2.340975in}{2.303520in}}%
\pgfpathlineto{\pgfqpoint{2.342130in}{2.303956in}}%
\pgfpathlineto{\pgfqpoint{2.343252in}{2.303714in}}%
\pgfpathlineto{\pgfqpoint{2.342757in}{2.304196in}}%
\pgfpathlineto{\pgfqpoint{2.343263in}{2.303759in}}%
\pgfpathlineto{\pgfqpoint{2.344370in}{2.303598in}}%
\pgfpathlineto{\pgfqpoint{2.343817in}{2.304277in}}%
\pgfpathlineto{\pgfqpoint{2.344398in}{2.303703in}}%
\pgfpathlineto{\pgfqpoint{2.345272in}{2.304144in}}%
\pgfpathlineto{\pgfqpoint{2.344586in}{2.303415in}}%
\pgfpathlineto{\pgfqpoint{2.345526in}{2.303858in}}%
\pgfpathlineto{\pgfqpoint{2.345527in}{2.303854in}}%
\pgfpathlineto{\pgfqpoint{2.346405in}{2.304374in}}%
\pgfpathlineto{\pgfqpoint{2.346587in}{2.304244in}}%
\pgfpathlineto{\pgfqpoint{2.346669in}{2.304268in}}%
\pgfpathlineto{\pgfqpoint{2.347087in}{2.303746in}}%
\pgfpathlineto{\pgfqpoint{2.347684in}{2.304000in}}%
\pgfpathlineto{\pgfqpoint{2.348889in}{2.303783in}}%
\pgfpathlineto{\pgfqpoint{2.348070in}{2.304424in}}%
\pgfpathlineto{\pgfqpoint{2.348986in}{2.303976in}}%
\pgfpathlineto{\pgfqpoint{2.349809in}{2.304162in}}%
\pgfpathlineto{\pgfqpoint{2.348999in}{2.303970in}}%
\pgfpathlineto{\pgfqpoint{2.350077in}{2.303780in}}%
\pgfpathlineto{\pgfqpoint{2.351114in}{2.303511in}}%
\pgfpathlineto{\pgfqpoint{2.350669in}{2.304041in}}%
\pgfpathlineto{\pgfqpoint{2.351193in}{2.303687in}}%
\pgfpathlineto{\pgfqpoint{2.351313in}{2.303759in}}%
\pgfpathlineto{\pgfqpoint{2.352043in}{2.302883in}}%
\pgfpathlineto{\pgfqpoint{2.352100in}{2.302927in}}%
\pgfpathlineto{\pgfqpoint{2.353016in}{2.302490in}}%
\pgfpathlineto{\pgfqpoint{2.352156in}{2.303030in}}%
\pgfpathlineto{\pgfqpoint{2.353217in}{2.302814in}}%
\pgfpathlineto{\pgfqpoint{2.354538in}{2.303312in}}%
\pgfpathlineto{\pgfqpoint{2.353276in}{2.302722in}}%
\pgfpathlineto{\pgfqpoint{2.354606in}{2.303219in}}%
\pgfpathlineto{\pgfqpoint{2.355226in}{2.303081in}}%
\pgfpathlineto{\pgfqpoint{2.354857in}{2.303375in}}%
\pgfpathlineto{\pgfqpoint{2.355716in}{2.303276in}}%
\pgfpathlineto{\pgfqpoint{2.356755in}{2.303185in}}%
\pgfpathlineto{\pgfqpoint{2.356070in}{2.303472in}}%
\pgfpathlineto{\pgfqpoint{2.356859in}{2.303281in}}%
\pgfpathlineto{\pgfqpoint{2.356922in}{2.303366in}}%
\pgfpathlineto{\pgfqpoint{2.357893in}{2.302555in}}%
\pgfpathlineto{\pgfqpoint{2.357898in}{2.302575in}}%
\pgfpathlineto{\pgfqpoint{2.358253in}{2.302276in}}%
\pgfpathlineto{\pgfqpoint{2.358958in}{2.302684in}}%
\pgfpathlineto{\pgfqpoint{2.359011in}{2.302548in}}%
\pgfpathlineto{\pgfqpoint{2.360115in}{2.302807in}}%
\pgfpathlineto{\pgfqpoint{2.359599in}{2.302319in}}%
\pgfpathlineto{\pgfqpoint{2.360143in}{2.302755in}}%
\pgfpathlineto{\pgfqpoint{2.360238in}{2.302639in}}%
\pgfpathlineto{\pgfqpoint{2.360640in}{2.303316in}}%
\pgfpathlineto{\pgfqpoint{2.361240in}{2.303219in}}%
\pgfpathlineto{\pgfqpoint{2.362667in}{2.303334in}}%
\pgfpathlineto{\pgfqpoint{2.362753in}{2.303051in}}%
\pgfpathlineto{\pgfqpoint{2.363890in}{2.302426in}}%
\pgfpathlineto{\pgfqpoint{2.362852in}{2.303197in}}%
\pgfpathlineto{\pgfqpoint{2.363918in}{2.302479in}}%
\pgfpathlineto{\pgfqpoint{2.376151in}{2.301315in}}%
\pgfpathlineto{\pgfqpoint{2.363923in}{2.302530in}}%
\pgfpathlineto{\pgfqpoint{2.376153in}{2.301363in}}%
\pgfpathlineto{\pgfqpoint{2.376745in}{2.301773in}}%
\pgfpathlineto{\pgfqpoint{2.376257in}{2.301162in}}%
\pgfpathlineto{\pgfqpoint{2.377279in}{2.301398in}}%
\pgfpathlineto{\pgfqpoint{2.378462in}{2.300848in}}%
\pgfpathlineto{\pgfqpoint{2.378598in}{2.301220in}}%
\pgfpathlineto{\pgfqpoint{2.379575in}{2.301267in}}%
\pgfpathlineto{\pgfqpoint{2.379216in}{2.300869in}}%
\pgfpathlineto{\pgfqpoint{2.379702in}{2.301017in}}%
\pgfpathlineto{\pgfqpoint{2.379876in}{2.301034in}}%
\pgfpathlineto{\pgfqpoint{2.380069in}{2.300645in}}%
\pgfpathlineto{\pgfqpoint{2.380807in}{2.300896in}}%
\pgfpathlineto{\pgfqpoint{2.381619in}{2.301023in}}%
\pgfpathlineto{\pgfqpoint{2.381993in}{2.300657in}}%
\pgfpathlineto{\pgfqpoint{2.383104in}{2.300574in}}%
\pgfpathlineto{\pgfqpoint{2.382711in}{2.301060in}}%
\pgfpathlineto{\pgfqpoint{2.383108in}{2.300602in}}%
\pgfpathlineto{\pgfqpoint{2.383280in}{2.300210in}}%
\pgfpathlineto{\pgfqpoint{2.383982in}{2.300822in}}%
\pgfpathlineto{\pgfqpoint{2.384210in}{2.300747in}}%
\pgfpathlineto{\pgfqpoint{2.385354in}{2.301283in}}%
\pgfpathlineto{\pgfqpoint{2.384264in}{2.300636in}}%
\pgfpathlineto{\pgfqpoint{2.385378in}{2.301212in}}%
\pgfpathlineto{\pgfqpoint{2.386668in}{2.301803in}}%
\pgfpathlineto{\pgfqpoint{2.385754in}{2.300933in}}%
\pgfpathlineto{\pgfqpoint{2.386678in}{2.301764in}}%
\pgfpathlineto{\pgfqpoint{2.386815in}{2.301943in}}%
\pgfpathlineto{\pgfqpoint{2.387807in}{2.301419in}}%
\pgfpathlineto{\pgfqpoint{2.387892in}{2.301487in}}%
\pgfpathlineto{\pgfqpoint{2.388550in}{2.300725in}}%
\pgfpathlineto{\pgfqpoint{2.388868in}{2.300929in}}%
\pgfpathlineto{\pgfqpoint{2.388874in}{2.300850in}}%
\pgfpathlineto{\pgfqpoint{2.389956in}{2.301432in}}%
\pgfpathlineto{\pgfqpoint{2.390794in}{2.301040in}}%
\pgfpathlineto{\pgfqpoint{2.389991in}{2.301542in}}%
\pgfpathlineto{\pgfqpoint{2.391111in}{2.301176in}}%
\pgfpathlineto{\pgfqpoint{2.391942in}{2.301521in}}%
\pgfpathlineto{\pgfqpoint{2.391446in}{2.300870in}}%
\pgfpathlineto{\pgfqpoint{2.392233in}{2.301454in}}%
\pgfpathlineto{\pgfqpoint{2.392354in}{2.301320in}}%
\pgfpathlineto{\pgfqpoint{2.393001in}{2.303445in}}%
\pgfpathlineto{\pgfqpoint{2.393181in}{2.303396in}}%
\pgfpathlineto{\pgfqpoint{2.394251in}{2.303447in}}%
\pgfpathlineto{\pgfqpoint{2.393283in}{2.303370in}}%
\pgfpathlineto{\pgfqpoint{2.394292in}{2.303398in}}%
\pgfpathlineto{\pgfqpoint{2.395429in}{2.303394in}}%
\pgfpathlineto{\pgfqpoint{2.394341in}{2.303442in}}%
\pgfpathlineto{\pgfqpoint{2.395436in}{2.303411in}}%
\pgfpathlineto{\pgfqpoint{2.396477in}{2.303442in}}%
\pgfpathlineto{\pgfqpoint{2.395448in}{2.303368in}}%
\pgfpathlineto{\pgfqpoint{2.396550in}{2.303431in}}%
\pgfpathlineto{\pgfqpoint{2.397735in}{2.303400in}}%
\pgfpathlineto{\pgfqpoint{2.397741in}{2.303416in}}%
\pgfpathlineto{\pgfqpoint{2.398850in}{2.303448in}}%
\pgfpathlineto{\pgfqpoint{2.397769in}{2.303363in}}%
\pgfpathlineto{\pgfqpoint{2.398854in}{2.303439in}}%
\pgfpathlineto{\pgfqpoint{2.399699in}{2.303357in}}%
\pgfpathlineto{\pgfqpoint{2.398866in}{2.303454in}}%
\pgfpathlineto{\pgfqpoint{2.399966in}{2.303419in}}%
\pgfpathlineto{\pgfqpoint{2.402207in}{2.303395in}}%
\pgfpathlineto{\pgfqpoint{2.399968in}{2.303424in}}%
\pgfpathlineto{\pgfqpoint{2.402213in}{2.303423in}}%
\pgfpathlineto{\pgfqpoint{2.403217in}{2.303456in}}%
\pgfpathlineto{\pgfqpoint{2.402438in}{2.303373in}}%
\pgfpathlineto{\pgfqpoint{2.403324in}{2.303412in}}%
\pgfpathlineto{\pgfqpoint{2.404383in}{2.303341in}}%
\pgfpathlineto{\pgfqpoint{2.403328in}{2.303448in}}%
\pgfpathlineto{\pgfqpoint{2.404438in}{2.303393in}}%
\pgfpathlineto{\pgfqpoint{2.405228in}{2.303329in}}%
\pgfpathlineto{\pgfqpoint{2.405554in}{2.303438in}}%
\pgfpathlineto{\pgfqpoint{2.406526in}{2.303341in}}%
\pgfpathlineto{\pgfqpoint{2.405679in}{2.303447in}}%
\pgfpathlineto{\pgfqpoint{2.406666in}{2.303421in}}%
\pgfpathlineto{\pgfqpoint{2.407785in}{2.303450in}}%
\pgfpathlineto{\pgfqpoint{2.406689in}{2.303368in}}%
\pgfpathlineto{\pgfqpoint{2.407785in}{2.303432in}}%
\pgfpathlineto{\pgfqpoint{2.408854in}{2.303340in}}%
\pgfpathlineto{\pgfqpoint{2.407907in}{2.303446in}}%
\pgfpathlineto{\pgfqpoint{2.408899in}{2.303412in}}%
\pgfpathlineto{\pgfqpoint{2.409842in}{2.303451in}}%
\pgfpathlineto{\pgfqpoint{2.408921in}{2.303361in}}%
\pgfpathlineto{\pgfqpoint{2.410012in}{2.303428in}}%
\pgfpathlineto{\pgfqpoint{2.411114in}{2.303360in}}%
\pgfpathlineto{\pgfqpoint{2.410110in}{2.303447in}}%
\pgfpathlineto{\pgfqpoint{2.411124in}{2.303411in}}%
\pgfpathlineto{\pgfqpoint{2.411721in}{2.303739in}}%
\pgfpathlineto{\pgfqpoint{2.412145in}{2.303230in}}%
\pgfpathlineto{\pgfqpoint{2.412232in}{2.303322in}}%
\pgfpathlineto{\pgfqpoint{2.413848in}{2.303300in}}%
\pgfpathlineto{\pgfqpoint{2.412250in}{2.303258in}}%
\pgfpathlineto{\pgfqpoint{2.413904in}{2.303173in}}%
\pgfpathlineto{\pgfqpoint{2.414449in}{2.302887in}}%
\pgfpathlineto{\pgfqpoint{2.413918in}{2.303279in}}%
\pgfpathlineto{\pgfqpoint{2.415022in}{2.302917in}}%
\pgfpathlineto{\pgfqpoint{2.415457in}{2.302608in}}%
\pgfpathlineto{\pgfqpoint{2.415790in}{2.303547in}}%
\pgfpathlineto{\pgfqpoint{2.415971in}{2.303530in}}%
\pgfpathlineto{\pgfqpoint{2.416858in}{2.303965in}}%
\pgfpathlineto{\pgfqpoint{2.416238in}{2.303391in}}%
\pgfpathlineto{\pgfqpoint{2.417096in}{2.303687in}}%
\pgfpathlineto{\pgfqpoint{2.421105in}{2.303153in}}%
\pgfpathlineto{\pgfqpoint{2.421106in}{2.303129in}}%
\pgfpathlineto{\pgfqpoint{2.422207in}{2.302905in}}%
\pgfpathlineto{\pgfqpoint{2.421764in}{2.303480in}}%
\pgfpathlineto{\pgfqpoint{2.422225in}{2.302928in}}%
\pgfpathlineto{\pgfqpoint{2.422826in}{2.302531in}}%
\pgfpathlineto{\pgfqpoint{2.422491in}{2.303217in}}%
\pgfpathlineto{\pgfqpoint{2.423362in}{2.302783in}}%
\pgfpathlineto{\pgfqpoint{2.423762in}{2.302983in}}%
\pgfpathlineto{\pgfqpoint{2.424031in}{2.302591in}}%
\pgfpathlineto{\pgfqpoint{2.424477in}{2.302844in}}%
\pgfpathlineto{\pgfqpoint{2.424505in}{2.302778in}}%
\pgfpathlineto{\pgfqpoint{2.425123in}{2.303290in}}%
\pgfpathlineto{\pgfqpoint{2.425560in}{2.303123in}}%
\pgfpathlineto{\pgfqpoint{2.426605in}{2.303336in}}%
\pgfpathlineto{\pgfqpoint{2.426183in}{2.302786in}}%
\pgfpathlineto{\pgfqpoint{2.426675in}{2.303241in}}%
\pgfpathlineto{\pgfqpoint{2.427770in}{2.303107in}}%
\pgfpathlineto{\pgfqpoint{2.427125in}{2.303702in}}%
\pgfpathlineto{\pgfqpoint{2.427795in}{2.303138in}}%
\pgfpathlineto{\pgfqpoint{2.428117in}{2.302835in}}%
\pgfpathlineto{\pgfqpoint{2.428916in}{2.303345in}}%
\pgfpathlineto{\pgfqpoint{2.430068in}{2.303061in}}%
\pgfpathlineto{\pgfqpoint{2.429045in}{2.303531in}}%
\pgfpathlineto{\pgfqpoint{2.430157in}{2.303343in}}%
\pgfpathlineto{\pgfqpoint{2.430790in}{2.303502in}}%
\pgfpathlineto{\pgfqpoint{2.430474in}{2.303057in}}%
\pgfpathlineto{\pgfqpoint{2.431256in}{2.303141in}}%
\pgfpathlineto{\pgfqpoint{2.432375in}{2.302875in}}%
\pgfpathlineto{\pgfqpoint{2.431383in}{2.303195in}}%
\pgfpathlineto{\pgfqpoint{2.432377in}{2.302908in}}%
\pgfpathlineto{\pgfqpoint{2.432720in}{2.303424in}}%
\pgfpathlineto{\pgfqpoint{2.433491in}{2.303023in}}%
\pgfpathlineto{\pgfqpoint{2.434786in}{2.301937in}}%
\pgfpathlineto{\pgfqpoint{2.433518in}{2.303041in}}%
\pgfpathlineto{\pgfqpoint{2.435042in}{2.302147in}}%
\pgfpathlineto{\pgfqpoint{2.435163in}{2.302281in}}%
\pgfpathlineto{\pgfqpoint{2.435801in}{2.301871in}}%
\pgfpathlineto{\pgfqpoint{2.436155in}{2.302166in}}%
\pgfpathlineto{\pgfqpoint{2.436212in}{2.302068in}}%
\pgfpathlineto{\pgfqpoint{2.437147in}{2.302780in}}%
\pgfpathlineto{\pgfqpoint{2.437223in}{2.302758in}}%
\pgfpathlineto{\pgfqpoint{2.437957in}{2.303102in}}%
\pgfpathlineto{\pgfqpoint{2.438230in}{2.302606in}}%
\pgfpathlineto{\pgfqpoint{2.438332in}{2.302680in}}%
\pgfpathlineto{\pgfqpoint{2.439628in}{2.302159in}}%
\pgfpathlineto{\pgfqpoint{2.438399in}{2.302797in}}%
\pgfpathlineto{\pgfqpoint{2.439637in}{2.302201in}}%
\pgfpathlineto{\pgfqpoint{2.440628in}{2.302539in}}%
\pgfpathlineto{\pgfqpoint{2.439751in}{2.302065in}}%
\pgfpathlineto{\pgfqpoint{2.440784in}{2.302429in}}%
\pgfpathlineto{\pgfqpoint{2.441867in}{2.301789in}}%
\pgfpathlineto{\pgfqpoint{2.441157in}{2.302838in}}%
\pgfpathlineto{\pgfqpoint{2.441921in}{2.301943in}}%
\pgfpathlineto{\pgfqpoint{2.443003in}{2.303082in}}%
\pgfpathlineto{\pgfqpoint{2.441929in}{2.301891in}}%
\pgfpathlineto{\pgfqpoint{2.443172in}{2.302831in}}%
\pgfpathlineto{\pgfqpoint{2.443812in}{2.302329in}}%
\pgfpathlineto{\pgfqpoint{2.444291in}{2.302469in}}%
\pgfpathlineto{\pgfqpoint{2.445266in}{2.302344in}}%
\pgfpathlineto{\pgfqpoint{2.444742in}{2.302705in}}%
\pgfpathlineto{\pgfqpoint{2.445403in}{2.302448in}}%
\pgfpathlineto{\pgfqpoint{2.445877in}{2.302921in}}%
\pgfpathlineto{\pgfqpoint{2.445685in}{2.302309in}}%
\pgfpathlineto{\pgfqpoint{2.446532in}{2.302665in}}%
\pgfpathlineto{\pgfqpoint{2.446601in}{2.302462in}}%
\pgfpathlineto{\pgfqpoint{2.447414in}{2.303363in}}%
\pgfpathlineto{\pgfqpoint{2.447622in}{2.303188in}}%
\pgfpathlineto{\pgfqpoint{2.447728in}{2.303314in}}%
\pgfpathlineto{\pgfqpoint{2.448401in}{2.302764in}}%
\pgfpathlineto{\pgfqpoint{2.448727in}{2.303009in}}%
\pgfpathlineto{\pgfqpoint{2.449600in}{2.302914in}}%
\pgfpathlineto{\pgfqpoint{2.449041in}{2.303531in}}%
\pgfpathlineto{\pgfqpoint{2.449815in}{2.303260in}}%
\pgfpathlineto{\pgfqpoint{2.450813in}{2.303138in}}%
\pgfpathlineto{\pgfqpoint{2.450332in}{2.302977in}}%
\pgfpathlineto{\pgfqpoint{2.450900in}{2.302988in}}%
\pgfpathlineto{\pgfqpoint{2.451869in}{2.302764in}}%
\pgfpathlineto{\pgfqpoint{2.451061in}{2.303142in}}%
\pgfpathlineto{\pgfqpoint{2.452014in}{2.302896in}}%
\pgfpathlineto{\pgfqpoint{2.452903in}{2.303556in}}%
\pgfpathlineto{\pgfqpoint{2.452016in}{2.302887in}}%
\pgfpathlineto{\pgfqpoint{2.453161in}{2.303448in}}%
\pgfpathlineto{\pgfqpoint{2.453879in}{2.303080in}}%
\pgfpathlineto{\pgfqpoint{2.453183in}{2.303540in}}%
\pgfpathlineto{\pgfqpoint{2.454285in}{2.303336in}}%
\pgfpathlineto{\pgfqpoint{2.454570in}{2.303432in}}%
\pgfpathlineto{\pgfqpoint{2.455280in}{2.302680in}}%
\pgfpathlineto{\pgfqpoint{2.455303in}{2.302725in}}%
\pgfpathlineto{\pgfqpoint{2.456411in}{2.302729in}}%
\pgfpathlineto{\pgfqpoint{2.456008in}{2.303335in}}%
\pgfpathlineto{\pgfqpoint{2.456419in}{2.302761in}}%
\pgfpathlineto{\pgfqpoint{2.456776in}{2.302821in}}%
\pgfpathlineto{\pgfqpoint{2.456943in}{2.302505in}}%
\pgfpathlineto{\pgfqpoint{2.457518in}{2.302598in}}%
\pgfpathlineto{\pgfqpoint{2.457599in}{2.302444in}}%
\pgfpathlineto{\pgfqpoint{2.458351in}{2.303330in}}%
\pgfpathlineto{\pgfqpoint{2.458608in}{2.303055in}}%
\pgfpathlineto{\pgfqpoint{2.458744in}{2.303204in}}%
\pgfpathlineto{\pgfqpoint{2.459436in}{2.302404in}}%
\pgfpathlineto{\pgfqpoint{2.459696in}{2.302439in}}%
\pgfpathlineto{\pgfqpoint{2.459712in}{2.302415in}}%
\pgfpathlineto{\pgfqpoint{2.460708in}{2.303037in}}%
\pgfpathlineto{\pgfqpoint{2.460709in}{2.303034in}}%
\pgfpathlineto{\pgfqpoint{2.461782in}{2.303714in}}%
\pgfpathlineto{\pgfqpoint{2.460833in}{2.302848in}}%
\pgfpathlineto{\pgfqpoint{2.461893in}{2.303559in}}%
\pgfpathlineto{\pgfqpoint{2.467544in}{2.301814in}}%
\pgfpathlineto{\pgfqpoint{2.461902in}{2.303619in}}%
\pgfpathlineto{\pgfqpoint{2.467721in}{2.303201in}}%
\pgfpathlineto{\pgfqpoint{2.468798in}{2.303445in}}%
\pgfpathlineto{\pgfqpoint{2.467724in}{2.303197in}}%
\pgfpathlineto{\pgfqpoint{2.468843in}{2.303422in}}%
\pgfpathlineto{\pgfqpoint{2.470122in}{2.303448in}}%
\pgfpathlineto{\pgfqpoint{2.468851in}{2.303383in}}%
\pgfpathlineto{\pgfqpoint{2.470131in}{2.303428in}}%
\pgfpathlineto{\pgfqpoint{2.471198in}{2.303350in}}%
\pgfpathlineto{\pgfqpoint{2.470260in}{2.303449in}}%
\pgfpathlineto{\pgfqpoint{2.471242in}{2.303406in}}%
\pgfpathlineto{\pgfqpoint{2.476739in}{2.303437in}}%
\pgfpathlineto{\pgfqpoint{2.471244in}{2.303392in}}%
\pgfpathlineto{\pgfqpoint{2.476741in}{2.303386in}}%
\pgfpathlineto{\pgfqpoint{2.477523in}{2.303331in}}%
\pgfpathlineto{\pgfqpoint{2.476780in}{2.303450in}}%
\pgfpathlineto{\pgfqpoint{2.477853in}{2.303422in}}%
\pgfpathlineto{\pgfqpoint{2.478911in}{2.303439in}}%
\pgfpathlineto{\pgfqpoint{2.478192in}{2.303349in}}%
\pgfpathlineto{\pgfqpoint{2.478967in}{2.303426in}}%
\pgfpathlineto{\pgfqpoint{2.479982in}{2.303369in}}%
\pgfpathlineto{\pgfqpoint{2.479003in}{2.303446in}}%
\pgfpathlineto{\pgfqpoint{2.480079in}{2.303396in}}%
\pgfpathlineto{\pgfqpoint{2.481192in}{2.303444in}}%
\pgfpathlineto{\pgfqpoint{2.480192in}{2.303371in}}%
\pgfpathlineto{\pgfqpoint{2.481194in}{2.303433in}}%
\pgfpathlineto{\pgfqpoint{2.482213in}{2.303367in}}%
\pgfpathlineto{\pgfqpoint{2.481199in}{2.303436in}}%
\pgfpathlineto{\pgfqpoint{2.482313in}{2.303423in}}%
\pgfpathlineto{\pgfqpoint{2.483477in}{2.303446in}}%
\pgfpathlineto{\pgfqpoint{2.482325in}{2.303370in}}%
\pgfpathlineto{\pgfqpoint{2.483480in}{2.303391in}}%
\pgfpathlineto{\pgfqpoint{2.484500in}{2.303367in}}%
\pgfpathlineto{\pgfqpoint{2.483540in}{2.303443in}}%
\pgfpathlineto{\pgfqpoint{2.484590in}{2.303419in}}%
\pgfpathlineto{\pgfqpoint{2.485667in}{2.303448in}}%
\pgfpathlineto{\pgfqpoint{2.484706in}{2.303364in}}%
\pgfpathlineto{\pgfqpoint{2.485701in}{2.303427in}}%
\pgfpathlineto{\pgfqpoint{2.486580in}{2.303268in}}%
\pgfpathlineto{\pgfqpoint{2.486366in}{2.303507in}}%
\pgfpathlineto{\pgfqpoint{2.486811in}{2.303467in}}%
\pgfpathlineto{\pgfqpoint{2.486881in}{2.303340in}}%
\pgfpathlineto{\pgfqpoint{2.487452in}{2.303747in}}%
\pgfpathlineto{\pgfqpoint{2.487902in}{2.303720in}}%
\pgfpathlineto{\pgfqpoint{2.489016in}{2.304286in}}%
\pgfpathlineto{\pgfqpoint{2.487923in}{2.303649in}}%
\pgfpathlineto{\pgfqpoint{2.489038in}{2.304241in}}%
\pgfpathlineto{\pgfqpoint{2.489312in}{2.303964in}}%
\pgfpathlineto{\pgfqpoint{2.489961in}{2.304848in}}%
\pgfpathlineto{\pgfqpoint{2.490071in}{2.304658in}}%
\pgfpathlineto{\pgfqpoint{2.490951in}{2.304854in}}%
\pgfpathlineto{\pgfqpoint{2.490605in}{2.304480in}}%
\pgfpathlineto{\pgfqpoint{2.491183in}{2.304652in}}%
\pgfpathlineto{\pgfqpoint{2.492343in}{2.305443in}}%
\pgfpathlineto{\pgfqpoint{2.493569in}{2.304165in}}%
\pgfpathlineto{\pgfqpoint{2.492373in}{2.305504in}}%
\pgfpathlineto{\pgfqpoint{2.493569in}{2.304178in}}%
\pgfpathlineto{\pgfqpoint{2.493599in}{2.304228in}}%
\pgfpathlineto{\pgfqpoint{2.494589in}{2.303379in}}%
\pgfpathlineto{\pgfqpoint{2.494655in}{2.303459in}}%
\pgfpathlineto{\pgfqpoint{2.495015in}{2.303206in}}%
\pgfpathlineto{\pgfqpoint{2.494715in}{2.303599in}}%
\pgfpathlineto{\pgfqpoint{2.495774in}{2.303413in}}%
\pgfpathlineto{\pgfqpoint{2.497188in}{2.302571in}}%
\pgfpathlineto{\pgfqpoint{2.495782in}{2.303417in}}%
\pgfpathlineto{\pgfqpoint{2.497230in}{2.302684in}}%
\pgfpathlineto{\pgfqpoint{2.498337in}{2.303107in}}%
\pgfpathlineto{\pgfqpoint{2.497857in}{2.302661in}}%
\pgfpathlineto{\pgfqpoint{2.498352in}{2.303013in}}%
\pgfpathlineto{\pgfqpoint{2.499178in}{2.302751in}}%
\pgfpathlineto{\pgfqpoint{2.498510in}{2.303342in}}%
\pgfpathlineto{\pgfqpoint{2.499503in}{2.302986in}}%
\pgfpathlineto{\pgfqpoint{2.499682in}{2.303043in}}%
\pgfpathlineto{\pgfqpoint{2.500564in}{2.302526in}}%
\pgfpathlineto{\pgfqpoint{2.500592in}{2.302570in}}%
\pgfpathlineto{\pgfqpoint{2.501176in}{2.303317in}}%
\pgfpathlineto{\pgfqpoint{2.500602in}{2.302537in}}%
\pgfpathlineto{\pgfqpoint{2.501738in}{2.302877in}}%
\pgfpathlineto{\pgfqpoint{2.501740in}{2.302846in}}%
\pgfpathlineto{\pgfqpoint{2.502786in}{2.303556in}}%
\pgfpathlineto{\pgfqpoint{2.502792in}{2.303537in}}%
\pgfpathlineto{\pgfqpoint{2.503643in}{2.303804in}}%
\pgfpathlineto{\pgfqpoint{2.502939in}{2.303421in}}%
\pgfpathlineto{\pgfqpoint{2.503906in}{2.303520in}}%
\pgfpathlineto{\pgfqpoint{2.504822in}{2.303485in}}%
\pgfpathlineto{\pgfqpoint{2.504654in}{2.303782in}}%
\pgfpathlineto{\pgfqpoint{2.504998in}{2.303785in}}%
\pgfpathlineto{\pgfqpoint{2.505538in}{2.304033in}}%
\pgfpathlineto{\pgfqpoint{2.505846in}{2.303406in}}%
\pgfpathlineto{\pgfqpoint{2.506058in}{2.303450in}}%
\pgfpathlineto{\pgfqpoint{2.507053in}{2.303244in}}%
\pgfpathlineto{\pgfqpoint{2.506378in}{2.303558in}}%
\pgfpathlineto{\pgfqpoint{2.507167in}{2.303507in}}%
\pgfpathlineto{\pgfqpoint{2.508253in}{2.303858in}}%
\pgfpathlineto{\pgfqpoint{2.507390in}{2.303397in}}%
\pgfpathlineto{\pgfqpoint{2.508290in}{2.303829in}}%
\pgfpathlineto{\pgfqpoint{2.509409in}{2.303397in}}%
\pgfpathlineto{\pgfqpoint{2.508769in}{2.304222in}}%
\pgfpathlineto{\pgfqpoint{2.509431in}{2.303415in}}%
\pgfpathlineto{\pgfqpoint{2.510423in}{2.304363in}}%
\pgfpathlineto{\pgfqpoint{2.510574in}{2.304276in}}%
\pgfpathlineto{\pgfqpoint{2.511789in}{2.304446in}}%
\pgfpathlineto{\pgfqpoint{2.510575in}{2.304280in}}%
\pgfpathlineto{\pgfqpoint{2.511888in}{2.304834in}}%
\pgfpathlineto{\pgfqpoint{2.511972in}{2.304915in}}%
\pgfpathlineto{\pgfqpoint{2.512604in}{2.304394in}}%
\pgfpathlineto{\pgfqpoint{2.512974in}{2.304385in}}%
\pgfpathlineto{\pgfqpoint{2.512983in}{2.304361in}}%
\pgfpathlineto{\pgfqpoint{2.513959in}{2.305001in}}%
\pgfpathlineto{\pgfqpoint{2.513982in}{2.304948in}}%
\pgfpathlineto{\pgfqpoint{2.514129in}{2.305268in}}%
\pgfpathlineto{\pgfqpoint{2.514529in}{2.304669in}}%
\pgfpathlineto{\pgfqpoint{2.515096in}{2.304950in}}%
\pgfpathlineto{\pgfqpoint{2.516172in}{2.304128in}}%
\pgfpathlineto{\pgfqpoint{2.515224in}{2.305112in}}%
\pgfpathlineto{\pgfqpoint{2.516325in}{2.304383in}}%
\pgfpathlineto{\pgfqpoint{2.517236in}{2.304609in}}%
\pgfpathlineto{\pgfqpoint{2.517088in}{2.304184in}}%
\pgfpathlineto{\pgfqpoint{2.517429in}{2.304226in}}%
\pgfpathlineto{\pgfqpoint{2.518536in}{2.304137in}}%
\pgfpathlineto{\pgfqpoint{2.517716in}{2.304523in}}%
\pgfpathlineto{\pgfqpoint{2.518544in}{2.304178in}}%
\pgfpathlineto{\pgfqpoint{2.519596in}{2.303986in}}%
\pgfpathlineto{\pgfqpoint{2.518669in}{2.304291in}}%
\pgfpathlineto{\pgfqpoint{2.519680in}{2.304121in}}%
\pgfpathlineto{\pgfqpoint{2.520846in}{2.304877in}}%
\pgfpathlineto{\pgfqpoint{2.520870in}{2.304812in}}%
\pgfpathlineto{\pgfqpoint{2.522054in}{2.304212in}}%
\pgfpathlineto{\pgfqpoint{2.520877in}{2.304843in}}%
\pgfpathlineto{\pgfqpoint{2.522104in}{2.304293in}}%
\pgfpathlineto{\pgfqpoint{2.522737in}{2.304987in}}%
\pgfpathlineto{\pgfqpoint{2.522219in}{2.304262in}}%
\pgfpathlineto{\pgfqpoint{2.523246in}{2.304593in}}%
\pgfpathlineto{\pgfqpoint{2.523938in}{2.304252in}}%
\pgfpathlineto{\pgfqpoint{2.524318in}{2.304929in}}%
\pgfpathlineto{\pgfqpoint{2.524626in}{2.305170in}}%
\pgfpathlineto{\pgfqpoint{2.525000in}{2.304627in}}%
\pgfpathlineto{\pgfqpoint{2.525435in}{2.305087in}}%
\pgfpathlineto{\pgfqpoint{2.527632in}{2.305180in}}%
\pgfpathlineto{\pgfqpoint{2.527667in}{2.305045in}}%
\pgfpathlineto{\pgfqpoint{2.528189in}{2.305048in}}%
\pgfpathlineto{\pgfqpoint{2.527948in}{2.305413in}}%
\pgfpathlineto{\pgfqpoint{2.528763in}{2.305267in}}%
\pgfpathlineto{\pgfqpoint{2.529494in}{2.305591in}}%
\pgfpathlineto{\pgfqpoint{2.528881in}{2.305046in}}%
\pgfpathlineto{\pgfqpoint{2.529861in}{2.305046in}}%
\pgfpathlineto{\pgfqpoint{2.530747in}{2.304905in}}%
\pgfpathlineto{\pgfqpoint{2.530262in}{2.305483in}}%
\pgfpathlineto{\pgfqpoint{2.530980in}{2.305150in}}%
\pgfpathlineto{\pgfqpoint{2.531948in}{2.305387in}}%
\pgfpathlineto{\pgfqpoint{2.531109in}{2.304970in}}%
\pgfpathlineto{\pgfqpoint{2.532095in}{2.305121in}}%
\pgfpathlineto{\pgfqpoint{2.532448in}{2.304851in}}%
\pgfpathlineto{\pgfqpoint{2.533211in}{2.305076in}}%
\pgfpathlineto{\pgfqpoint{2.534201in}{2.305400in}}%
\pgfpathlineto{\pgfqpoint{2.533272in}{2.304947in}}%
\pgfpathlineto{\pgfqpoint{2.534327in}{2.305127in}}%
\pgfpathlineto{\pgfqpoint{2.535268in}{2.304291in}}%
\pgfpathlineto{\pgfqpoint{2.534349in}{2.305172in}}%
\pgfpathlineto{\pgfqpoint{2.535479in}{2.304485in}}%
\pgfpathlineto{\pgfqpoint{2.536047in}{2.304788in}}%
\pgfpathlineto{\pgfqpoint{2.535506in}{2.304426in}}%
\pgfpathlineto{\pgfqpoint{2.536614in}{2.304666in}}%
\pgfpathlineto{\pgfqpoint{2.537046in}{2.303946in}}%
\pgfpathlineto{\pgfqpoint{2.536617in}{2.304678in}}%
\pgfpathlineto{\pgfqpoint{2.537757in}{2.304428in}}%
\pgfpathlineto{\pgfqpoint{2.537979in}{2.304883in}}%
\pgfpathlineto{\pgfqpoint{2.537762in}{2.304427in}}%
\pgfpathlineto{\pgfqpoint{2.538879in}{2.304564in}}%
\pgfpathlineto{\pgfqpoint{2.539305in}{2.304220in}}%
\pgfpathlineto{\pgfqpoint{2.539839in}{2.304977in}}%
\pgfpathlineto{\pgfqpoint{2.539980in}{2.304869in}}%
\pgfpathlineto{\pgfqpoint{2.541168in}{2.304380in}}%
\pgfpathlineto{\pgfqpoint{2.540030in}{2.304965in}}%
\pgfpathlineto{\pgfqpoint{2.541197in}{2.304495in}}%
\pgfpathlineto{\pgfqpoint{2.541513in}{2.304766in}}%
\pgfpathlineto{\pgfqpoint{2.541761in}{2.304110in}}%
\pgfpathlineto{\pgfqpoint{2.542318in}{2.304628in}}%
\pgfpathlineto{\pgfqpoint{2.546539in}{2.303392in}}%
\pgfpathlineto{\pgfqpoint{2.542334in}{2.304684in}}%
\pgfpathlineto{\pgfqpoint{2.546543in}{2.303422in}}%
\pgfpathlineto{\pgfqpoint{2.547646in}{2.303447in}}%
\pgfpathlineto{\pgfqpoint{2.546834in}{2.303356in}}%
\pgfpathlineto{\pgfqpoint{2.547655in}{2.303435in}}%
\pgfpathlineto{\pgfqpoint{2.548594in}{2.303371in}}%
\pgfpathlineto{\pgfqpoint{2.547848in}{2.303455in}}%
\pgfpathlineto{\pgfqpoint{2.548766in}{2.303424in}}%
\pgfpathlineto{\pgfqpoint{2.549882in}{2.303445in}}%
\pgfpathlineto{\pgfqpoint{2.548774in}{2.303384in}}%
\pgfpathlineto{\pgfqpoint{2.549901in}{2.303428in}}%
\pgfpathlineto{\pgfqpoint{2.550660in}{2.303351in}}%
\pgfpathlineto{\pgfqpoint{2.549920in}{2.303447in}}%
\pgfpathlineto{\pgfqpoint{2.551013in}{2.303401in}}%
\pgfpathlineto{\pgfqpoint{2.552177in}{2.303344in}}%
\pgfpathlineto{\pgfqpoint{2.551018in}{2.303430in}}%
\pgfpathlineto{\pgfqpoint{2.552189in}{2.303415in}}%
\pgfpathlineto{\pgfqpoint{2.553259in}{2.303445in}}%
\pgfpathlineto{\pgfqpoint{2.552222in}{2.303369in}}%
\pgfpathlineto{\pgfqpoint{2.553304in}{2.303434in}}%
\pgfpathlineto{\pgfqpoint{2.553615in}{2.303355in}}%
\pgfpathlineto{\pgfqpoint{2.553450in}{2.303456in}}%
\pgfpathlineto{\pgfqpoint{2.554416in}{2.303401in}}%
\pgfpathlineto{\pgfqpoint{2.555529in}{2.303440in}}%
\pgfpathlineto{\pgfqpoint{2.554458in}{2.303366in}}%
\pgfpathlineto{\pgfqpoint{2.555531in}{2.303409in}}%
\pgfpathlineto{\pgfqpoint{2.556667in}{2.303444in}}%
\pgfpathlineto{\pgfqpoint{2.555538in}{2.303391in}}%
\pgfpathlineto{\pgfqpoint{2.556671in}{2.303404in}}%
\pgfpathlineto{\pgfqpoint{2.557749in}{2.303372in}}%
\pgfpathlineto{\pgfqpoint{2.556686in}{2.303438in}}%
\pgfpathlineto{\pgfqpoint{2.557792in}{2.303407in}}%
\pgfpathlineto{\pgfqpoint{2.557800in}{2.303367in}}%
\pgfpathlineto{\pgfqpoint{2.558904in}{2.303447in}}%
\pgfpathlineto{\pgfqpoint{2.559163in}{2.303322in}}%
\pgfpathlineto{\pgfqpoint{2.558907in}{2.303448in}}%
\pgfpathlineto{\pgfqpoint{2.560017in}{2.303408in}}%
\pgfpathlineto{\pgfqpoint{2.561114in}{2.303443in}}%
\pgfpathlineto{\pgfqpoint{2.560045in}{2.303366in}}%
\pgfpathlineto{\pgfqpoint{2.561129in}{2.303393in}}%
\pgfpathlineto{\pgfqpoint{2.561699in}{2.303218in}}%
\pgfpathlineto{\pgfqpoint{2.561367in}{2.303689in}}%
\pgfpathlineto{\pgfqpoint{2.562243in}{2.303621in}}%
\pgfpathlineto{\pgfqpoint{2.562503in}{2.303985in}}%
\pgfpathlineto{\pgfqpoint{2.563293in}{2.303452in}}%
\pgfpathlineto{\pgfqpoint{2.563345in}{2.303465in}}%
\pgfpathlineto{\pgfqpoint{2.564112in}{2.303458in}}%
\pgfpathlineto{\pgfqpoint{2.564424in}{2.303891in}}%
\pgfpathlineto{\pgfqpoint{2.565557in}{2.304816in}}%
\pgfpathlineto{\pgfqpoint{2.565576in}{2.304764in}}%
\pgfpathlineto{\pgfqpoint{2.566648in}{2.305257in}}%
\pgfpathlineto{\pgfqpoint{2.565874in}{2.304669in}}%
\pgfpathlineto{\pgfqpoint{2.566723in}{2.305098in}}%
\pgfpathlineto{\pgfqpoint{2.567869in}{2.304703in}}%
\pgfpathlineto{\pgfqpoint{2.566769in}{2.305109in}}%
\pgfpathlineto{\pgfqpoint{2.567882in}{2.304735in}}%
\pgfpathlineto{\pgfqpoint{2.568790in}{2.305182in}}%
\pgfpathlineto{\pgfqpoint{2.567884in}{2.304734in}}%
\pgfpathlineto{\pgfqpoint{2.569015in}{2.305025in}}%
\pgfpathlineto{\pgfqpoint{2.569979in}{2.304739in}}%
\pgfpathlineto{\pgfqpoint{2.569577in}{2.305175in}}%
\pgfpathlineto{\pgfqpoint{2.570150in}{2.304979in}}%
\pgfpathlineto{\pgfqpoint{2.571214in}{2.304765in}}%
\pgfpathlineto{\pgfqpoint{2.570152in}{2.304973in}}%
\pgfpathlineto{\pgfqpoint{2.571282in}{2.304436in}}%
\pgfpathlineto{\pgfqpoint{2.572607in}{2.303859in}}%
\pgfpathlineto{\pgfqpoint{2.571389in}{2.304678in}}%
\pgfpathlineto{\pgfqpoint{2.572618in}{2.303870in}}%
\pgfpathlineto{\pgfqpoint{2.573691in}{2.303960in}}%
\pgfpathlineto{\pgfqpoint{2.573360in}{2.303558in}}%
\pgfpathlineto{\pgfqpoint{2.573743in}{2.303846in}}%
\pgfpathlineto{\pgfqpoint{2.574551in}{2.303459in}}%
\pgfpathlineto{\pgfqpoint{2.573833in}{2.303907in}}%
\pgfpathlineto{\pgfqpoint{2.574881in}{2.303571in}}%
\pgfpathlineto{\pgfqpoint{2.575765in}{2.303653in}}%
\pgfpathlineto{\pgfqpoint{2.575508in}{2.303210in}}%
\pgfpathlineto{\pgfqpoint{2.575983in}{2.303436in}}%
\pgfpathlineto{\pgfqpoint{2.576086in}{2.303276in}}%
\pgfpathlineto{\pgfqpoint{2.576272in}{2.303694in}}%
\pgfpathlineto{\pgfqpoint{2.577099in}{2.303386in}}%
\pgfpathlineto{\pgfqpoint{2.577504in}{2.303677in}}%
\pgfpathlineto{\pgfqpoint{2.578201in}{2.303133in}}%
\pgfpathlineto{\pgfqpoint{2.579238in}{2.302872in}}%
\pgfpathlineto{\pgfqpoint{2.578575in}{2.303650in}}%
\pgfpathlineto{\pgfqpoint{2.579329in}{2.303124in}}%
\pgfpathlineto{\pgfqpoint{2.580133in}{2.303065in}}%
\pgfpathlineto{\pgfqpoint{2.579850in}{2.302669in}}%
\pgfpathlineto{\pgfqpoint{2.580423in}{2.302932in}}%
\pgfpathlineto{\pgfqpoint{2.580631in}{2.302753in}}%
\pgfpathlineto{\pgfqpoint{2.581280in}{2.303236in}}%
\pgfpathlineto{\pgfqpoint{2.581526in}{2.303114in}}%
\pgfpathlineto{\pgfqpoint{2.582243in}{2.303567in}}%
\pgfpathlineto{\pgfqpoint{2.581542in}{2.303094in}}%
\pgfpathlineto{\pgfqpoint{2.582643in}{2.303249in}}%
\pgfpathlineto{\pgfqpoint{2.582893in}{2.303105in}}%
\pgfpathlineto{\pgfqpoint{2.583708in}{2.303754in}}%
\pgfpathlineto{\pgfqpoint{2.583745in}{2.303648in}}%
\pgfpathlineto{\pgfqpoint{2.585310in}{2.303246in}}%
\pgfpathlineto{\pgfqpoint{2.585318in}{2.303188in}}%
\pgfpathlineto{\pgfqpoint{2.585384in}{2.303061in}}%
\pgfpathlineto{\pgfqpoint{2.586400in}{2.303596in}}%
\pgfpathlineto{\pgfqpoint{2.587642in}{2.303757in}}%
\pgfpathlineto{\pgfqpoint{2.586406in}{2.303562in}}%
\pgfpathlineto{\pgfqpoint{2.587773in}{2.303508in}}%
\pgfpathlineto{\pgfqpoint{2.588842in}{2.303201in}}%
\pgfpathlineto{\pgfqpoint{2.587955in}{2.303736in}}%
\pgfpathlineto{\pgfqpoint{2.588898in}{2.303294in}}%
\pgfpathlineto{\pgfqpoint{2.597668in}{2.303944in}}%
\pgfpathlineto{\pgfqpoint{2.597668in}{2.303973in}}%
\pgfpathlineto{\pgfqpoint{2.598724in}{2.304310in}}%
\pgfpathlineto{\pgfqpoint{2.598477in}{2.303873in}}%
\pgfpathlineto{\pgfqpoint{2.598791in}{2.304248in}}%
\pgfpathlineto{\pgfqpoint{2.599327in}{2.303774in}}%
\pgfpathlineto{\pgfqpoint{2.598824in}{2.304304in}}%
\pgfpathlineto{\pgfqpoint{2.599954in}{2.304137in}}%
\pgfpathlineto{\pgfqpoint{2.601010in}{2.304560in}}%
\pgfpathlineto{\pgfqpoint{2.601093in}{2.304487in}}%
\pgfpathlineto{\pgfqpoint{2.601217in}{2.304311in}}%
\pgfpathlineto{\pgfqpoint{2.602030in}{2.304745in}}%
\pgfpathlineto{\pgfqpoint{2.602214in}{2.304476in}}%
\pgfpathlineto{\pgfqpoint{2.602947in}{2.304866in}}%
\pgfpathlineto{\pgfqpoint{2.602678in}{2.304405in}}%
\pgfpathlineto{\pgfqpoint{2.603325in}{2.304467in}}%
\pgfpathlineto{\pgfqpoint{2.603743in}{2.303912in}}%
\pgfpathlineto{\pgfqpoint{2.604243in}{2.304664in}}%
\pgfpathlineto{\pgfqpoint{2.604467in}{2.304408in}}%
\pgfpathlineto{\pgfqpoint{2.605298in}{2.304577in}}%
\pgfpathlineto{\pgfqpoint{2.604471in}{2.304396in}}%
\pgfpathlineto{\pgfqpoint{2.605596in}{2.304451in}}%
\pgfpathlineto{\pgfqpoint{2.606529in}{2.304149in}}%
\pgfpathlineto{\pgfqpoint{2.606022in}{2.304702in}}%
\pgfpathlineto{\pgfqpoint{2.606714in}{2.304261in}}%
\pgfpathlineto{\pgfqpoint{2.607232in}{2.304513in}}%
\pgfpathlineto{\pgfqpoint{2.607661in}{2.303749in}}%
\pgfpathlineto{\pgfqpoint{2.607806in}{2.303875in}}%
\pgfpathlineto{\pgfqpoint{2.610085in}{2.305157in}}%
\pgfpathlineto{\pgfqpoint{2.610209in}{2.305235in}}%
\pgfpathlineto{\pgfqpoint{2.610768in}{2.304702in}}%
\pgfpathlineto{\pgfqpoint{2.611189in}{2.304880in}}%
\pgfpathlineto{\pgfqpoint{2.611489in}{2.305221in}}%
\pgfpathlineto{\pgfqpoint{2.612138in}{2.304734in}}%
\pgfpathlineto{\pgfqpoint{2.612306in}{2.304936in}}%
\pgfpathlineto{\pgfqpoint{2.612391in}{2.304846in}}%
\pgfpathlineto{\pgfqpoint{2.612942in}{2.305431in}}%
\pgfpathlineto{\pgfqpoint{2.613381in}{2.305320in}}%
\pgfpathlineto{\pgfqpoint{2.614595in}{2.305406in}}%
\pgfpathlineto{\pgfqpoint{2.613427in}{2.305452in}}%
\pgfpathlineto{\pgfqpoint{2.614632in}{2.305465in}}%
\pgfpathlineto{\pgfqpoint{2.615509in}{2.305642in}}%
\pgfpathlineto{\pgfqpoint{2.615172in}{2.305252in}}%
\pgfpathlineto{\pgfqpoint{2.615739in}{2.305364in}}%
\pgfpathlineto{\pgfqpoint{2.616586in}{2.304906in}}%
\pgfpathlineto{\pgfqpoint{2.615803in}{2.305455in}}%
\pgfpathlineto{\pgfqpoint{2.616844in}{2.305488in}}%
\pgfpathlineto{\pgfqpoint{2.616845in}{2.305500in}}%
\pgfpathlineto{\pgfqpoint{2.617422in}{2.303372in}}%
\pgfpathlineto{\pgfqpoint{2.617828in}{2.303387in}}%
\pgfpathlineto{\pgfqpoint{2.618765in}{2.303453in}}%
\pgfpathlineto{\pgfqpoint{2.617884in}{2.303342in}}%
\pgfpathlineto{\pgfqpoint{2.618944in}{2.303406in}}%
\pgfpathlineto{\pgfqpoint{2.619857in}{2.303456in}}%
\pgfpathlineto{\pgfqpoint{2.619071in}{2.303364in}}%
\pgfpathlineto{\pgfqpoint{2.620060in}{2.303400in}}%
\pgfpathlineto{\pgfqpoint{2.621078in}{2.303365in}}%
\pgfpathlineto{\pgfqpoint{2.620193in}{2.303451in}}%
\pgfpathlineto{\pgfqpoint{2.621172in}{2.303413in}}%
\pgfpathlineto{\pgfqpoint{2.622274in}{2.303450in}}%
\pgfpathlineto{\pgfqpoint{2.621255in}{2.303359in}}%
\pgfpathlineto{\pgfqpoint{2.622286in}{2.303424in}}%
\pgfpathlineto{\pgfqpoint{2.622774in}{2.303345in}}%
\pgfpathlineto{\pgfqpoint{2.622328in}{2.303454in}}%
\pgfpathlineto{\pgfqpoint{2.623397in}{2.303408in}}%
\pgfpathlineto{\pgfqpoint{2.624507in}{2.303443in}}%
\pgfpathlineto{\pgfqpoint{2.623539in}{2.303353in}}%
\pgfpathlineto{\pgfqpoint{2.624510in}{2.303434in}}%
\pgfpathlineto{\pgfqpoint{2.625513in}{2.303360in}}%
\pgfpathlineto{\pgfqpoint{2.624775in}{2.303451in}}%
\pgfpathlineto{\pgfqpoint{2.625622in}{2.303408in}}%
\pgfpathlineto{\pgfqpoint{2.626521in}{2.303455in}}%
\pgfpathlineto{\pgfqpoint{2.625643in}{2.303385in}}%
\pgfpathlineto{\pgfqpoint{2.626734in}{2.303398in}}%
\pgfpathlineto{\pgfqpoint{2.627811in}{2.303448in}}%
\pgfpathlineto{\pgfqpoint{2.626788in}{2.303349in}}%
\pgfpathlineto{\pgfqpoint{2.627855in}{2.303429in}}%
\pgfpathlineto{\pgfqpoint{2.628630in}{2.303358in}}%
\pgfpathlineto{\pgfqpoint{2.628007in}{2.303446in}}%
\pgfpathlineto{\pgfqpoint{2.628967in}{2.303404in}}%
\pgfpathlineto{\pgfqpoint{2.629994in}{2.303444in}}%
\pgfpathlineto{\pgfqpoint{2.629015in}{2.303378in}}%
\pgfpathlineto{\pgfqpoint{2.630079in}{2.303426in}}%
\pgfpathlineto{\pgfqpoint{2.631056in}{2.303352in}}%
\pgfpathlineto{\pgfqpoint{2.630090in}{2.303438in}}%
\pgfpathlineto{\pgfqpoint{2.631190in}{2.303427in}}%
\pgfpathlineto{\pgfqpoint{2.632297in}{2.303439in}}%
\pgfpathlineto{\pgfqpoint{2.631221in}{2.303353in}}%
\pgfpathlineto{\pgfqpoint{2.632313in}{2.303427in}}%
\pgfpathlineto{\pgfqpoint{2.632503in}{2.303335in}}%
\pgfpathlineto{\pgfqpoint{2.633075in}{2.303463in}}%
\pgfpathlineto{\pgfqpoint{2.633424in}{2.303394in}}%
\pgfpathlineto{\pgfqpoint{2.634161in}{2.303457in}}%
\pgfpathlineto{\pgfqpoint{2.633447in}{2.303377in}}%
\pgfpathlineto{\pgfqpoint{2.634536in}{2.303414in}}%
\pgfpathlineto{\pgfqpoint{2.635631in}{2.303372in}}%
\pgfpathlineto{\pgfqpoint{2.634624in}{2.303443in}}%
\pgfpathlineto{\pgfqpoint{2.635648in}{2.303419in}}%
\pgfpathlineto{\pgfqpoint{2.636376in}{2.303566in}}%
\pgfpathlineto{\pgfqpoint{2.636158in}{2.303085in}}%
\pgfpathlineto{\pgfqpoint{2.636756in}{2.303349in}}%
\pgfpathlineto{\pgfqpoint{2.637015in}{2.302984in}}%
\pgfpathlineto{\pgfqpoint{2.637806in}{2.303727in}}%
\pgfpathlineto{\pgfqpoint{2.637824in}{2.303709in}}%
\pgfpathlineto{\pgfqpoint{2.637860in}{2.303831in}}%
\pgfpathlineto{\pgfqpoint{2.638594in}{2.303075in}}%
\pgfpathlineto{\pgfqpoint{2.638911in}{2.303374in}}%
\pgfpathlineto{\pgfqpoint{2.640716in}{2.303770in}}%
\pgfpathlineto{\pgfqpoint{2.639051in}{2.303047in}}%
\pgfpathlineto{\pgfqpoint{2.640867in}{2.303458in}}%
\pgfpathlineto{\pgfqpoint{2.641767in}{2.303396in}}%
\pgfpathlineto{\pgfqpoint{2.640978in}{2.303766in}}%
\pgfpathlineto{\pgfqpoint{2.641976in}{2.303521in}}%
\pgfpathlineto{\pgfqpoint{2.642564in}{2.302874in}}%
\pgfpathlineto{\pgfqpoint{2.641980in}{2.303544in}}%
\pgfpathlineto{\pgfqpoint{2.643242in}{2.303460in}}%
\pgfpathlineto{\pgfqpoint{2.644157in}{2.303738in}}%
\pgfpathlineto{\pgfqpoint{2.643724in}{2.303220in}}%
\pgfpathlineto{\pgfqpoint{2.644354in}{2.303431in}}%
\pgfpathlineto{\pgfqpoint{2.645063in}{2.303326in}}%
\pgfpathlineto{\pgfqpoint{2.645332in}{2.303690in}}%
\pgfpathlineto{\pgfqpoint{2.645464in}{2.303477in}}%
\pgfpathlineto{\pgfqpoint{2.646584in}{2.303117in}}%
\pgfpathlineto{\pgfqpoint{2.645471in}{2.303539in}}%
\pgfpathlineto{\pgfqpoint{2.646665in}{2.303228in}}%
\pgfpathlineto{\pgfqpoint{2.646783in}{2.303356in}}%
\pgfpathlineto{\pgfqpoint{2.647418in}{2.302817in}}%
\pgfpathlineto{\pgfqpoint{2.647767in}{2.303038in}}%
\pgfpathlineto{\pgfqpoint{2.648889in}{2.302679in}}%
\pgfpathlineto{\pgfqpoint{2.648073in}{2.303371in}}%
\pgfpathlineto{\pgfqpoint{2.648889in}{2.302687in}}%
\pgfpathlineto{\pgfqpoint{2.650020in}{2.303295in}}%
\pgfpathlineto{\pgfqpoint{2.649043in}{2.302610in}}%
\pgfpathlineto{\pgfqpoint{2.650055in}{2.303241in}}%
\pgfpathlineto{\pgfqpoint{2.650817in}{2.302780in}}%
\pgfpathlineto{\pgfqpoint{2.650530in}{2.303366in}}%
\pgfpathlineto{\pgfqpoint{2.651161in}{2.303379in}}%
\pgfpathlineto{\pgfqpoint{2.652283in}{2.303501in}}%
\pgfpathlineto{\pgfqpoint{2.651342in}{2.303052in}}%
\pgfpathlineto{\pgfqpoint{2.652300in}{2.303401in}}%
\pgfpathlineto{\pgfqpoint{2.652885in}{2.302832in}}%
\pgfpathlineto{\pgfqpoint{2.652303in}{2.303431in}}%
\pgfpathlineto{\pgfqpoint{2.653441in}{2.303179in}}%
\pgfpathlineto{\pgfqpoint{2.655058in}{2.304383in}}%
\pgfpathlineto{\pgfqpoint{2.653664in}{2.302985in}}%
\pgfpathlineto{\pgfqpoint{2.655074in}{2.304297in}}%
\pgfpathlineto{\pgfqpoint{2.655330in}{2.304051in}}%
\pgfpathlineto{\pgfqpoint{2.656096in}{2.304608in}}%
\pgfpathlineto{\pgfqpoint{2.656179in}{2.304412in}}%
\pgfpathlineto{\pgfqpoint{2.656561in}{2.304701in}}%
\pgfpathlineto{\pgfqpoint{2.656729in}{2.304253in}}%
\pgfpathlineto{\pgfqpoint{2.657300in}{2.304511in}}%
\pgfpathlineto{\pgfqpoint{2.657606in}{2.304842in}}%
\pgfpathlineto{\pgfqpoint{2.658304in}{2.303968in}}%
\pgfpathlineto{\pgfqpoint{2.659301in}{2.303644in}}%
\pgfpathlineto{\pgfqpoint{2.658447in}{2.304132in}}%
\pgfpathlineto{\pgfqpoint{2.659429in}{2.303693in}}%
\pgfpathlineto{\pgfqpoint{2.660620in}{2.304356in}}%
\pgfpathlineto{\pgfqpoint{2.659677in}{2.303431in}}%
\pgfpathlineto{\pgfqpoint{2.660661in}{2.304259in}}%
\pgfpathlineto{\pgfqpoint{2.660736in}{2.304031in}}%
\pgfpathlineto{\pgfqpoint{2.661398in}{2.304684in}}%
\pgfpathlineto{\pgfqpoint{2.661765in}{2.304494in}}%
\pgfpathlineto{\pgfqpoint{2.662838in}{2.304845in}}%
\pgfpathlineto{\pgfqpoint{2.662284in}{2.304251in}}%
\pgfpathlineto{\pgfqpoint{2.662906in}{2.304751in}}%
\pgfpathlineto{\pgfqpoint{2.663530in}{2.304246in}}%
\pgfpathlineto{\pgfqpoint{2.664025in}{2.304535in}}%
\pgfpathlineto{\pgfqpoint{2.664589in}{2.304968in}}%
\pgfpathlineto{\pgfqpoint{2.664116in}{2.304490in}}%
\pgfpathlineto{\pgfqpoint{2.665144in}{2.304715in}}%
\pgfpathlineto{\pgfqpoint{2.666210in}{2.304340in}}%
\pgfpathlineto{\pgfqpoint{2.665501in}{2.304818in}}%
\pgfpathlineto{\pgfqpoint{2.666278in}{2.304519in}}%
\pgfpathlineto{\pgfqpoint{2.667350in}{2.304405in}}%
\pgfpathlineto{\pgfqpoint{2.667111in}{2.304866in}}%
\pgfpathlineto{\pgfqpoint{2.667409in}{2.304502in}}%
\pgfpathlineto{\pgfqpoint{2.668451in}{2.304688in}}%
\pgfpathlineto{\pgfqpoint{2.667528in}{2.304313in}}%
\pgfpathlineto{\pgfqpoint{2.668525in}{2.304528in}}%
\pgfpathlineto{\pgfqpoint{2.669159in}{2.304214in}}%
\pgfpathlineto{\pgfqpoint{2.668601in}{2.304662in}}%
\pgfpathlineto{\pgfqpoint{2.669623in}{2.304698in}}%
\pgfpathlineto{\pgfqpoint{2.669964in}{2.304985in}}%
\pgfpathlineto{\pgfqpoint{2.670588in}{2.304581in}}%
\pgfpathlineto{\pgfqpoint{2.670738in}{2.304742in}}%
\pgfpathlineto{\pgfqpoint{2.671234in}{2.304961in}}%
\pgfpathlineto{\pgfqpoint{2.671807in}{2.304410in}}%
\pgfpathlineto{\pgfqpoint{2.671828in}{2.304468in}}%
\pgfpathlineto{\pgfqpoint{2.672857in}{2.304403in}}%
\pgfpathlineto{\pgfqpoint{2.671981in}{2.304631in}}%
\pgfpathlineto{\pgfqpoint{2.672941in}{2.304522in}}%
\pgfpathlineto{\pgfqpoint{2.674079in}{2.304979in}}%
\pgfpathlineto{\pgfqpoint{2.673107in}{2.304337in}}%
\pgfpathlineto{\pgfqpoint{2.674082in}{2.304959in}}%
\pgfpathlineto{\pgfqpoint{2.674998in}{2.304717in}}%
\pgfpathlineto{\pgfqpoint{2.675155in}{2.305271in}}%
\pgfpathlineto{\pgfqpoint{2.675179in}{2.305249in}}%
\pgfpathlineto{\pgfqpoint{2.675248in}{2.305160in}}%
\pgfpathlineto{\pgfqpoint{2.676418in}{2.306074in}}%
\pgfpathlineto{\pgfqpoint{2.683012in}{2.307303in}}%
\pgfpathlineto{\pgfqpoint{2.683040in}{2.307181in}}%
\pgfpathlineto{\pgfqpoint{2.684140in}{2.306672in}}%
\pgfpathlineto{\pgfqpoint{2.683478in}{2.307404in}}%
\pgfpathlineto{\pgfqpoint{2.684182in}{2.306746in}}%
\pgfpathlineto{\pgfqpoint{2.685193in}{2.306802in}}%
\pgfpathlineto{\pgfqpoint{2.684368in}{2.306472in}}%
\pgfpathlineto{\pgfqpoint{2.685281in}{2.306576in}}%
\pgfpathlineto{\pgfqpoint{2.685590in}{2.306393in}}%
\pgfpathlineto{\pgfqpoint{2.686292in}{2.306765in}}%
\pgfpathlineto{\pgfqpoint{2.686385in}{2.306700in}}%
\pgfpathlineto{\pgfqpoint{2.687383in}{2.306840in}}%
\pgfpathlineto{\pgfqpoint{2.686942in}{2.306479in}}%
\pgfpathlineto{\pgfqpoint{2.687491in}{2.306554in}}%
\pgfpathlineto{\pgfqpoint{2.688512in}{2.305710in}}%
\pgfpathlineto{\pgfqpoint{2.687498in}{2.306513in}}%
\pgfpathlineto{\pgfqpoint{2.688515in}{2.305689in}}%
\pgfpathlineto{\pgfqpoint{2.689550in}{2.305447in}}%
\pgfpathlineto{\pgfqpoint{2.689081in}{2.306037in}}%
\pgfpathlineto{\pgfqpoint{2.689628in}{2.305660in}}%
\pgfpathlineto{\pgfqpoint{2.690793in}{2.304507in}}%
\pgfpathlineto{\pgfqpoint{2.690980in}{2.304767in}}%
\pgfpathlineto{\pgfqpoint{2.691176in}{2.304838in}}%
\pgfpathlineto{\pgfqpoint{2.691695in}{2.304287in}}%
\pgfpathlineto{\pgfqpoint{2.692072in}{2.304381in}}%
\pgfpathlineto{\pgfqpoint{2.692240in}{2.304395in}}%
\pgfpathlineto{\pgfqpoint{2.692304in}{2.303382in}}%
\pgfpathlineto{\pgfqpoint{2.692973in}{2.303411in}}%
\pgfpathlineto{\pgfqpoint{2.693887in}{2.303349in}}%
\pgfpathlineto{\pgfqpoint{2.693011in}{2.303456in}}%
\pgfpathlineto{\pgfqpoint{2.694092in}{2.303414in}}%
\pgfpathlineto{\pgfqpoint{2.695150in}{2.303451in}}%
\pgfpathlineto{\pgfqpoint{2.694187in}{2.303359in}}%
\pgfpathlineto{\pgfqpoint{2.695202in}{2.303386in}}%
\pgfpathlineto{\pgfqpoint{2.696294in}{2.303447in}}%
\pgfpathlineto{\pgfqpoint{2.695256in}{2.303364in}}%
\pgfpathlineto{\pgfqpoint{2.696328in}{2.303394in}}%
\pgfpathlineto{\pgfqpoint{2.697133in}{2.303356in}}%
\pgfpathlineto{\pgfqpoint{2.696344in}{2.303439in}}%
\pgfpathlineto{\pgfqpoint{2.697443in}{2.303413in}}%
\pgfpathlineto{\pgfqpoint{2.698499in}{2.303449in}}%
\pgfpathlineto{\pgfqpoint{2.697558in}{2.303361in}}%
\pgfpathlineto{\pgfqpoint{2.698555in}{2.303413in}}%
\pgfpathlineto{\pgfqpoint{2.699639in}{2.303378in}}%
\pgfpathlineto{\pgfqpoint{2.698625in}{2.303445in}}%
\pgfpathlineto{\pgfqpoint{2.699671in}{2.303411in}}%
\pgfpathlineto{\pgfqpoint{2.700778in}{2.303456in}}%
\pgfpathlineto{\pgfqpoint{2.699780in}{2.303366in}}%
\pgfpathlineto{\pgfqpoint{2.700801in}{2.303428in}}%
\pgfpathlineto{\pgfqpoint{2.701788in}{2.303356in}}%
\pgfpathlineto{\pgfqpoint{2.700819in}{2.303446in}}%
\pgfpathlineto{\pgfqpoint{2.701913in}{2.303402in}}%
\pgfpathlineto{\pgfqpoint{2.703474in}{2.303434in}}%
\pgfpathlineto{\pgfqpoint{2.703490in}{2.303377in}}%
\pgfpathlineto{\pgfqpoint{2.704517in}{2.303375in}}%
\pgfpathlineto{\pgfqpoint{2.703545in}{2.303448in}}%
\pgfpathlineto{\pgfqpoint{2.704602in}{2.303385in}}%
\pgfpathlineto{\pgfqpoint{2.705732in}{2.303449in}}%
\pgfpathlineto{\pgfqpoint{2.704622in}{2.303378in}}%
\pgfpathlineto{\pgfqpoint{2.705737in}{2.303406in}}%
\pgfpathlineto{\pgfqpoint{2.706838in}{2.303351in}}%
\pgfpathlineto{\pgfqpoint{2.705837in}{2.303443in}}%
\pgfpathlineto{\pgfqpoint{2.706849in}{2.303408in}}%
\pgfpathlineto{\pgfqpoint{2.707935in}{2.303440in}}%
\pgfpathlineto{\pgfqpoint{2.706986in}{2.303336in}}%
\pgfpathlineto{\pgfqpoint{2.707960in}{2.303374in}}%
\pgfpathlineto{\pgfqpoint{2.709153in}{2.303440in}}%
\pgfpathlineto{\pgfqpoint{2.709154in}{2.303432in}}%
\pgfpathlineto{\pgfqpoint{2.710017in}{2.303359in}}%
\pgfpathlineto{\pgfqpoint{2.709225in}{2.303453in}}%
\pgfpathlineto{\pgfqpoint{2.710266in}{2.303401in}}%
\pgfpathlineto{\pgfqpoint{2.711393in}{2.304167in}}%
\pgfpathlineto{\pgfqpoint{2.710537in}{2.303325in}}%
\pgfpathlineto{\pgfqpoint{2.711398in}{2.304135in}}%
\pgfpathlineto{\pgfqpoint{2.712575in}{2.304324in}}%
\pgfpathlineto{\pgfqpoint{2.711494in}{2.304023in}}%
\pgfpathlineto{\pgfqpoint{2.712581in}{2.304312in}}%
\pgfpathlineto{\pgfqpoint{2.712629in}{2.304231in}}%
\pgfpathlineto{\pgfqpoint{2.713657in}{2.305154in}}%
\pgfpathlineto{\pgfqpoint{2.714636in}{2.305251in}}%
\pgfpathlineto{\pgfqpoint{2.714798in}{2.305388in}}%
\pgfpathlineto{\pgfqpoint{2.715926in}{2.305669in}}%
\pgfpathlineto{\pgfqpoint{2.714985in}{2.305162in}}%
\pgfpathlineto{\pgfqpoint{2.715929in}{2.305654in}}%
\pgfpathlineto{\pgfqpoint{2.716983in}{2.305530in}}%
\pgfpathlineto{\pgfqpoint{2.716114in}{2.305978in}}%
\pgfpathlineto{\pgfqpoint{2.717105in}{2.305718in}}%
\pgfpathlineto{\pgfqpoint{2.717143in}{2.305539in}}%
\pgfpathlineto{\pgfqpoint{2.718302in}{2.306378in}}%
\pgfpathlineto{\pgfqpoint{2.719397in}{2.305941in}}%
\pgfpathlineto{\pgfqpoint{2.719113in}{2.306513in}}%
\pgfpathlineto{\pgfqpoint{2.719476in}{2.305972in}}%
\pgfpathlineto{\pgfqpoint{2.720617in}{2.306579in}}%
\pgfpathlineto{\pgfqpoint{2.719572in}{2.305899in}}%
\pgfpathlineto{\pgfqpoint{2.720666in}{2.306514in}}%
\pgfpathlineto{\pgfqpoint{2.721458in}{2.306101in}}%
\pgfpathlineto{\pgfqpoint{2.721785in}{2.306365in}}%
\pgfpathlineto{\pgfqpoint{2.723470in}{2.305981in}}%
\pgfpathlineto{\pgfqpoint{2.723507in}{2.306077in}}%
\pgfpathlineto{\pgfqpoint{2.723840in}{2.306562in}}%
\pgfpathlineto{\pgfqpoint{2.723530in}{2.306043in}}%
\pgfpathlineto{\pgfqpoint{2.724631in}{2.306453in}}%
\pgfpathlineto{\pgfqpoint{2.725117in}{2.306758in}}%
\pgfpathlineto{\pgfqpoint{2.724759in}{2.306288in}}%
\pgfpathlineto{\pgfqpoint{2.725740in}{2.306134in}}%
\pgfpathlineto{\pgfqpoint{2.726829in}{2.305603in}}%
\pgfpathlineto{\pgfqpoint{2.726871in}{2.305699in}}%
\pgfpathlineto{\pgfqpoint{2.726889in}{2.305758in}}%
\pgfpathlineto{\pgfqpoint{2.727471in}{2.305174in}}%
\pgfpathlineto{\pgfqpoint{2.727964in}{2.305390in}}%
\pgfpathlineto{\pgfqpoint{2.729035in}{2.305338in}}%
\pgfpathlineto{\pgfqpoint{2.728489in}{2.305862in}}%
\pgfpathlineto{\pgfqpoint{2.729077in}{2.305451in}}%
\pgfpathlineto{\pgfqpoint{2.729194in}{2.305649in}}%
\pgfpathlineto{\pgfqpoint{2.730013in}{2.305248in}}%
\pgfpathlineto{\pgfqpoint{2.730190in}{2.305428in}}%
\pgfpathlineto{\pgfqpoint{2.731462in}{2.306349in}}%
\pgfpathlineto{\pgfqpoint{2.731463in}{2.306360in}}%
\pgfpathlineto{\pgfqpoint{2.732205in}{2.306536in}}%
\pgfpathlineto{\pgfqpoint{2.731769in}{2.306212in}}%
\pgfpathlineto{\pgfqpoint{2.732569in}{2.306238in}}%
\pgfpathlineto{\pgfqpoint{2.733681in}{2.305787in}}%
\pgfpathlineto{\pgfqpoint{2.732658in}{2.306265in}}%
\pgfpathlineto{\pgfqpoint{2.733709in}{2.305817in}}%
\pgfpathlineto{\pgfqpoint{2.734841in}{2.306258in}}%
\pgfpathlineto{\pgfqpoint{2.733902in}{2.305644in}}%
\pgfpathlineto{\pgfqpoint{2.734842in}{2.306235in}}%
\pgfpathlineto{\pgfqpoint{2.734842in}{2.306222in}}%
\pgfpathlineto{\pgfqpoint{2.735164in}{2.306694in}}%
\pgfpathlineto{\pgfqpoint{2.735936in}{2.306557in}}%
\pgfpathlineto{\pgfqpoint{2.736155in}{2.306601in}}%
\pgfpathlineto{\pgfqpoint{2.736705in}{2.305892in}}%
\pgfpathlineto{\pgfqpoint{2.737020in}{2.306149in}}%
\pgfpathlineto{\pgfqpoint{2.737531in}{2.305921in}}%
\pgfpathlineto{\pgfqpoint{2.737225in}{2.306406in}}%
\pgfpathlineto{\pgfqpoint{2.738128in}{2.306224in}}%
\pgfpathlineto{\pgfqpoint{2.739142in}{2.306524in}}%
\pgfpathlineto{\pgfqpoint{2.738156in}{2.306221in}}%
\pgfpathlineto{\pgfqpoint{2.739242in}{2.306320in}}%
\pgfpathlineto{\pgfqpoint{2.740505in}{2.305699in}}%
\pgfpathlineto{\pgfqpoint{2.739244in}{2.306299in}}%
\pgfpathlineto{\pgfqpoint{2.740513in}{2.305667in}}%
\pgfpathlineto{\pgfqpoint{2.741392in}{2.304949in}}%
\pgfpathlineto{\pgfqpoint{2.740635in}{2.305678in}}%
\pgfpathlineto{\pgfqpoint{2.741633in}{2.305387in}}%
\pgfpathlineto{\pgfqpoint{2.742759in}{2.305430in}}%
\pgfpathlineto{\pgfqpoint{2.741998in}{2.305059in}}%
\pgfpathlineto{\pgfqpoint{2.742760in}{2.305430in}}%
\pgfpathlineto{\pgfqpoint{2.743356in}{2.305203in}}%
\pgfpathlineto{\pgfqpoint{2.743058in}{2.305578in}}%
\pgfpathlineto{\pgfqpoint{2.743863in}{2.305568in}}%
\pgfpathlineto{\pgfqpoint{2.743866in}{2.305581in}}%
\pgfpathlineto{\pgfqpoint{2.744836in}{2.304994in}}%
\pgfpathlineto{\pgfqpoint{2.744944in}{2.305086in}}%
\pgfpathlineto{\pgfqpoint{2.745176in}{2.304958in}}%
\pgfpathlineto{\pgfqpoint{2.745870in}{2.305534in}}%
\pgfpathlineto{\pgfqpoint{2.746038in}{2.305323in}}%
\pgfpathlineto{\pgfqpoint{2.746117in}{2.305511in}}%
\pgfpathlineto{\pgfqpoint{2.746973in}{2.305079in}}%
\pgfpathlineto{\pgfqpoint{2.747151in}{2.305307in}}%
\pgfpathlineto{\pgfqpoint{2.748342in}{2.304971in}}%
\pgfpathlineto{\pgfqpoint{2.747327in}{2.305395in}}%
\pgfpathlineto{\pgfqpoint{2.748354in}{2.305014in}}%
\pgfpathlineto{\pgfqpoint{2.748680in}{2.305325in}}%
\pgfpathlineto{\pgfqpoint{2.749313in}{2.304876in}}%
\pgfpathlineto{\pgfqpoint{2.749466in}{2.305054in}}%
\pgfpathlineto{\pgfqpoint{2.750563in}{2.304935in}}%
\pgfpathlineto{\pgfqpoint{2.750112in}{2.305433in}}%
\pgfpathlineto{\pgfqpoint{2.750582in}{2.304952in}}%
\pgfpathlineto{\pgfqpoint{2.751278in}{2.304973in}}%
\pgfpathlineto{\pgfqpoint{2.751534in}{2.304559in}}%
\pgfpathlineto{\pgfqpoint{2.751664in}{2.304562in}}%
\pgfpathlineto{\pgfqpoint{2.751715in}{2.304411in}}%
\pgfpathlineto{\pgfqpoint{2.752663in}{2.305053in}}%
\pgfpathlineto{\pgfqpoint{2.752767in}{2.304831in}}%
\pgfpathlineto{\pgfqpoint{2.756752in}{2.303328in}}%
\pgfpathlineto{\pgfqpoint{2.756790in}{2.303481in}}%
\pgfpathlineto{\pgfqpoint{2.757731in}{2.303665in}}%
\pgfpathlineto{\pgfqpoint{2.757887in}{2.303260in}}%
\pgfpathlineto{\pgfqpoint{2.758858in}{2.302989in}}%
\pgfpathlineto{\pgfqpoint{2.758241in}{2.303587in}}%
\pgfpathlineto{\pgfqpoint{2.759028in}{2.303199in}}%
\pgfpathlineto{\pgfqpoint{2.760071in}{2.303624in}}%
\pgfpathlineto{\pgfqpoint{2.759246in}{2.303011in}}%
\pgfpathlineto{\pgfqpoint{2.760170in}{2.303592in}}%
\pgfpathlineto{\pgfqpoint{2.761439in}{2.304764in}}%
\pgfpathlineto{\pgfqpoint{2.760196in}{2.303585in}}%
\pgfpathlineto{\pgfqpoint{2.761469in}{2.304709in}}%
\pgfpathlineto{\pgfqpoint{2.762567in}{2.304283in}}%
\pgfpathlineto{\pgfqpoint{2.761598in}{2.304905in}}%
\pgfpathlineto{\pgfqpoint{2.762595in}{2.304410in}}%
\pgfpathlineto{\pgfqpoint{2.762771in}{2.304471in}}%
\pgfpathlineto{\pgfqpoint{2.763513in}{2.303729in}}%
\pgfpathlineto{\pgfqpoint{2.763660in}{2.303994in}}%
\pgfpathlineto{\pgfqpoint{2.763841in}{2.303857in}}%
\pgfpathlineto{\pgfqpoint{2.764704in}{2.304298in}}%
\pgfpathlineto{\pgfqpoint{2.764759in}{2.304225in}}%
\pgfpathlineto{\pgfqpoint{2.765877in}{2.303798in}}%
\pgfpathlineto{\pgfqpoint{2.764785in}{2.304314in}}%
\pgfpathlineto{\pgfqpoint{2.765898in}{2.303861in}}%
\pgfpathlineto{\pgfqpoint{2.767050in}{2.303548in}}%
\pgfpathlineto{\pgfqpoint{2.766393in}{2.304047in}}%
\pgfpathlineto{\pgfqpoint{2.767053in}{2.303560in}}%
\pgfpathlineto{\pgfqpoint{2.767141in}{2.303858in}}%
\pgfpathlineto{\pgfqpoint{2.767166in}{2.303382in}}%
\pgfpathlineto{\pgfqpoint{2.768152in}{2.303411in}}%
\pgfpathlineto{\pgfqpoint{2.769178in}{2.303377in}}%
\pgfpathlineto{\pgfqpoint{2.768241in}{2.303448in}}%
\pgfpathlineto{\pgfqpoint{2.769266in}{2.303414in}}%
\pgfpathlineto{\pgfqpoint{2.770308in}{2.303451in}}%
\pgfpathlineto{\pgfqpoint{2.769305in}{2.303384in}}%
\pgfpathlineto{\pgfqpoint{2.770377in}{2.303395in}}%
\pgfpathlineto{\pgfqpoint{2.771510in}{2.303443in}}%
\pgfpathlineto{\pgfqpoint{2.770439in}{2.303341in}}%
\pgfpathlineto{\pgfqpoint{2.771544in}{2.303396in}}%
\pgfpathlineto{\pgfqpoint{2.772099in}{2.303323in}}%
\pgfpathlineto{\pgfqpoint{2.771610in}{2.303443in}}%
\pgfpathlineto{\pgfqpoint{2.772656in}{2.303395in}}%
\pgfpathlineto{\pgfqpoint{2.774056in}{2.303378in}}%
\pgfpathlineto{\pgfqpoint{2.772663in}{2.303434in}}%
\pgfpathlineto{\pgfqpoint{2.774062in}{2.303422in}}%
\pgfpathlineto{\pgfqpoint{2.775134in}{2.303438in}}%
\pgfpathlineto{\pgfqpoint{2.774211in}{2.303368in}}%
\pgfpathlineto{\pgfqpoint{2.775174in}{2.303412in}}%
\pgfpathlineto{\pgfqpoint{2.776269in}{2.303441in}}%
\pgfpathlineto{\pgfqpoint{2.775220in}{2.303374in}}%
\pgfpathlineto{\pgfqpoint{2.776306in}{2.303428in}}%
\pgfpathlineto{\pgfqpoint{2.776790in}{2.303328in}}%
\pgfpathlineto{\pgfqpoint{2.776339in}{2.303445in}}%
\pgfpathlineto{\pgfqpoint{2.777416in}{2.303411in}}%
\pgfpathlineto{\pgfqpoint{2.778481in}{2.303448in}}%
\pgfpathlineto{\pgfqpoint{2.777735in}{2.303359in}}%
\pgfpathlineto{\pgfqpoint{2.778528in}{2.303387in}}%
\pgfpathlineto{\pgfqpoint{2.779644in}{2.303438in}}%
\pgfpathlineto{\pgfqpoint{2.778536in}{2.303370in}}%
\pgfpathlineto{\pgfqpoint{2.779646in}{2.303438in}}%
\pgfpathlineto{\pgfqpoint{2.780521in}{2.303365in}}%
\pgfpathlineto{\pgfqpoint{2.779892in}{2.303457in}}%
\pgfpathlineto{\pgfqpoint{2.780758in}{2.303411in}}%
\pgfpathlineto{\pgfqpoint{2.781838in}{2.303447in}}%
\pgfpathlineto{\pgfqpoint{2.780801in}{2.303365in}}%
\pgfpathlineto{\pgfqpoint{2.781881in}{2.303415in}}%
\pgfpathlineto{\pgfqpoint{2.782928in}{2.303347in}}%
\pgfpathlineto{\pgfqpoint{2.781888in}{2.303443in}}%
\pgfpathlineto{\pgfqpoint{2.783037in}{2.303416in}}%
\pgfpathlineto{\pgfqpoint{2.784110in}{2.303447in}}%
\pgfpathlineto{\pgfqpoint{2.783189in}{2.303368in}}%
\pgfpathlineto{\pgfqpoint{2.784148in}{2.303425in}}%
\pgfpathlineto{\pgfqpoint{2.785254in}{2.303359in}}%
\pgfpathlineto{\pgfqpoint{2.784272in}{2.303451in}}%
\pgfpathlineto{\pgfqpoint{2.785260in}{2.303382in}}%
\pgfpathlineto{\pgfqpoint{2.786378in}{2.303437in}}%
\pgfpathlineto{\pgfqpoint{2.786001in}{2.303236in}}%
\pgfpathlineto{\pgfqpoint{2.786380in}{2.303414in}}%
\pgfpathlineto{\pgfqpoint{2.786732in}{2.303006in}}%
\pgfpathlineto{\pgfqpoint{2.787498in}{2.303283in}}%
\pgfpathlineto{\pgfqpoint{2.787980in}{2.303347in}}%
\pgfpathlineto{\pgfqpoint{2.788510in}{2.302779in}}%
\pgfpathlineto{\pgfqpoint{2.788547in}{2.302841in}}%
\pgfpathlineto{\pgfqpoint{2.788674in}{2.302407in}}%
\pgfpathlineto{\pgfqpoint{2.789662in}{2.302770in}}%
\pgfpathlineto{\pgfqpoint{2.790827in}{2.302969in}}%
\pgfpathlineto{\pgfqpoint{2.790018in}{2.302446in}}%
\pgfpathlineto{\pgfqpoint{2.790836in}{2.302928in}}%
\pgfpathlineto{\pgfqpoint{2.795073in}{2.303745in}}%
\pgfpathlineto{\pgfqpoint{2.796129in}{2.304056in}}%
\pgfpathlineto{\pgfqpoint{2.795451in}{2.303692in}}%
\pgfpathlineto{\pgfqpoint{2.796191in}{2.304003in}}%
\pgfpathlineto{\pgfqpoint{2.796667in}{2.303840in}}%
\pgfpathlineto{\pgfqpoint{2.797186in}{2.304231in}}%
\pgfpathlineto{\pgfqpoint{2.797298in}{2.304124in}}%
\pgfpathlineto{\pgfqpoint{2.802289in}{2.304949in}}%
\pgfpathlineto{\pgfqpoint{2.797302in}{2.304110in}}%
\pgfpathlineto{\pgfqpoint{2.802299in}{2.304865in}}%
\pgfpathlineto{\pgfqpoint{2.803407in}{2.304732in}}%
\pgfpathlineto{\pgfqpoint{2.802804in}{2.305341in}}%
\pgfpathlineto{\pgfqpoint{2.803429in}{2.304751in}}%
\pgfpathlineto{\pgfqpoint{2.804526in}{2.305026in}}%
\pgfpathlineto{\pgfqpoint{2.803754in}{2.304469in}}%
\pgfpathlineto{\pgfqpoint{2.804551in}{2.304965in}}%
\pgfpathlineto{\pgfqpoint{2.805731in}{2.304331in}}%
\pgfpathlineto{\pgfqpoint{2.804696in}{2.305145in}}%
\pgfpathlineto{\pgfqpoint{2.805811in}{2.304456in}}%
\pgfpathlineto{\pgfqpoint{2.806221in}{2.304326in}}%
\pgfpathlineto{\pgfqpoint{2.806949in}{2.304987in}}%
\pgfpathlineto{\pgfqpoint{2.808060in}{2.304743in}}%
\pgfpathlineto{\pgfqpoint{2.807002in}{2.305066in}}%
\pgfpathlineto{\pgfqpoint{2.808069in}{2.304815in}}%
\pgfpathlineto{\pgfqpoint{2.809397in}{2.304503in}}%
\pgfpathlineto{\pgfqpoint{2.808084in}{2.304858in}}%
\pgfpathlineto{\pgfqpoint{2.809402in}{2.304533in}}%
\pgfpathlineto{\pgfqpoint{2.810411in}{2.304851in}}%
\pgfpathlineto{\pgfqpoint{2.809841in}{2.304332in}}%
\pgfpathlineto{\pgfqpoint{2.810518in}{2.304603in}}%
\pgfpathlineto{\pgfqpoint{2.811654in}{2.305433in}}%
\pgfpathlineto{\pgfqpoint{2.810605in}{2.304480in}}%
\pgfpathlineto{\pgfqpoint{2.811799in}{2.305157in}}%
\pgfpathlineto{\pgfqpoint{2.812754in}{2.304447in}}%
\pgfpathlineto{\pgfqpoint{2.812975in}{2.304592in}}%
\pgfpathlineto{\pgfqpoint{2.813803in}{2.304724in}}%
\pgfpathlineto{\pgfqpoint{2.813353in}{2.304385in}}%
\pgfpathlineto{\pgfqpoint{2.814123in}{2.304598in}}%
\pgfpathlineto{\pgfqpoint{2.814123in}{2.304583in}}%
\pgfpathlineto{\pgfqpoint{2.815189in}{2.305260in}}%
\pgfpathlineto{\pgfqpoint{2.815203in}{2.305202in}}%
\pgfpathlineto{\pgfqpoint{2.816612in}{2.305515in}}%
\pgfpathlineto{\pgfqpoint{2.815205in}{2.305191in}}%
\pgfpathlineto{\pgfqpoint{2.816669in}{2.305374in}}%
\pgfpathlineto{\pgfqpoint{2.817397in}{2.304690in}}%
\pgfpathlineto{\pgfqpoint{2.816925in}{2.305572in}}%
\pgfpathlineto{\pgfqpoint{2.817800in}{2.304902in}}%
\pgfpathlineto{\pgfqpoint{2.818268in}{2.304728in}}%
\pgfpathlineto{\pgfqpoint{2.818582in}{2.305156in}}%
\pgfpathlineto{\pgfqpoint{2.818903in}{2.305031in}}%
\pgfpathlineto{\pgfqpoint{2.819508in}{2.305161in}}%
\pgfpathlineto{\pgfqpoint{2.819085in}{2.304815in}}%
\pgfpathlineto{\pgfqpoint{2.820009in}{2.304931in}}%
\pgfpathlineto{\pgfqpoint{2.821227in}{2.304119in}}%
\pgfpathlineto{\pgfqpoint{2.820012in}{2.304953in}}%
\pgfpathlineto{\pgfqpoint{2.821228in}{2.304130in}}%
\pgfpathlineto{\pgfqpoint{2.821650in}{2.303772in}}%
\pgfpathlineto{\pgfqpoint{2.822383in}{2.304214in}}%
\pgfpathlineto{\pgfqpoint{2.822855in}{2.304851in}}%
\pgfpathlineto{\pgfqpoint{2.823522in}{2.304590in}}%
\pgfpathlineto{\pgfqpoint{2.824649in}{2.304184in}}%
\pgfpathlineto{\pgfqpoint{2.823816in}{2.304843in}}%
\pgfpathlineto{\pgfqpoint{2.824659in}{2.304190in}}%
\pgfpathlineto{\pgfqpoint{2.825611in}{2.304011in}}%
\pgfpathlineto{\pgfqpoint{2.824744in}{2.304438in}}%
\pgfpathlineto{\pgfqpoint{2.825777in}{2.304166in}}%
\pgfpathlineto{\pgfqpoint{2.827006in}{2.303273in}}%
\pgfpathlineto{\pgfqpoint{2.825810in}{2.304224in}}%
\pgfpathlineto{\pgfqpoint{2.827138in}{2.303454in}}%
\pgfpathlineto{\pgfqpoint{2.827183in}{2.303581in}}%
\pgfpathlineto{\pgfqpoint{2.827835in}{2.303060in}}%
\pgfpathlineto{\pgfqpoint{2.828238in}{2.303056in}}%
\pgfpathlineto{\pgfqpoint{2.829378in}{2.303151in}}%
\pgfpathlineto{\pgfqpoint{2.828326in}{2.302875in}}%
\pgfpathlineto{\pgfqpoint{2.829384in}{2.303109in}}%
\pgfpathlineto{\pgfqpoint{2.829510in}{2.302992in}}%
\pgfpathlineto{\pgfqpoint{2.830285in}{2.303451in}}%
\pgfpathlineto{\pgfqpoint{2.830494in}{2.303130in}}%
\pgfpathlineto{\pgfqpoint{2.831552in}{2.303935in}}%
\pgfpathlineto{\pgfqpoint{2.830721in}{2.302793in}}%
\pgfpathlineto{\pgfqpoint{2.831643in}{2.303849in}}%
\pgfpathlineto{\pgfqpoint{2.831714in}{2.303820in}}%
\pgfpathlineto{\pgfqpoint{2.832531in}{2.304648in}}%
\pgfpathlineto{\pgfqpoint{2.832724in}{2.304442in}}%
\pgfpathlineto{\pgfqpoint{2.832838in}{2.304369in}}%
\pgfpathlineto{\pgfqpoint{2.833736in}{2.304861in}}%
\pgfpathlineto{\pgfqpoint{2.833815in}{2.304662in}}%
\pgfpathlineto{\pgfqpoint{2.834604in}{2.304887in}}%
\pgfpathlineto{\pgfqpoint{2.834428in}{2.304454in}}%
\pgfpathlineto{\pgfqpoint{2.834927in}{2.304641in}}%
\pgfpathlineto{\pgfqpoint{2.834970in}{2.304484in}}%
\pgfpathlineto{\pgfqpoint{2.835687in}{2.305099in}}%
\pgfpathlineto{\pgfqpoint{2.836036in}{2.304702in}}%
\pgfpathlineto{\pgfqpoint{2.836832in}{2.304045in}}%
\pgfpathlineto{\pgfqpoint{2.836055in}{2.304766in}}%
\pgfpathlineto{\pgfqpoint{2.837173in}{2.304433in}}%
\pgfpathlineto{\pgfqpoint{2.837238in}{2.304450in}}%
\pgfpathlineto{\pgfqpoint{2.838154in}{2.303203in}}%
\pgfpathlineto{\pgfqpoint{2.838157in}{2.303210in}}%
\pgfpathlineto{\pgfqpoint{2.839218in}{2.302587in}}%
\pgfpathlineto{\pgfqpoint{2.838311in}{2.303228in}}%
\pgfpathlineto{\pgfqpoint{2.839285in}{2.302719in}}%
\pgfpathlineto{\pgfqpoint{2.840263in}{2.302336in}}%
\pgfpathlineto{\pgfqpoint{2.839464in}{2.302920in}}%
\pgfpathlineto{\pgfqpoint{2.840414in}{2.302432in}}%
\pgfpathlineto{\pgfqpoint{2.840949in}{2.302653in}}%
\pgfpathlineto{\pgfqpoint{2.841507in}{2.302139in}}%
\pgfpathlineto{\pgfqpoint{2.841515in}{2.302175in}}%
\pgfpathlineto{\pgfqpoint{2.841520in}{2.302149in}}%
\pgfpathlineto{\pgfqpoint{2.842267in}{2.303452in}}%
\pgfpathlineto{\pgfqpoint{2.842427in}{2.303419in}}%
\pgfpathlineto{\pgfqpoint{2.843534in}{2.303448in}}%
\pgfpathlineto{\pgfqpoint{2.842544in}{2.303369in}}%
\pgfpathlineto{\pgfqpoint{2.843540in}{2.303387in}}%
\pgfpathlineto{\pgfqpoint{2.844523in}{2.303358in}}%
\pgfpathlineto{\pgfqpoint{2.843545in}{2.303439in}}%
\pgfpathlineto{\pgfqpoint{2.844656in}{2.303420in}}%
\pgfpathlineto{\pgfqpoint{2.846155in}{2.303429in}}%
\pgfpathlineto{\pgfqpoint{2.844673in}{2.303387in}}%
\pgfpathlineto{\pgfqpoint{2.846158in}{2.303378in}}%
\pgfpathlineto{\pgfqpoint{2.847240in}{2.303367in}}%
\pgfpathlineto{\pgfqpoint{2.846178in}{2.303444in}}%
\pgfpathlineto{\pgfqpoint{2.847270in}{2.303410in}}%
\pgfpathlineto{\pgfqpoint{2.848301in}{2.303459in}}%
\pgfpathlineto{\pgfqpoint{2.847293in}{2.303388in}}%
\pgfpathlineto{\pgfqpoint{2.848395in}{2.303407in}}%
\pgfpathlineto{\pgfqpoint{2.849885in}{2.303393in}}%
\pgfpathlineto{\pgfqpoint{2.848398in}{2.303434in}}%
\pgfpathlineto{\pgfqpoint{2.849891in}{2.303412in}}%
\pgfpathlineto{\pgfqpoint{2.850863in}{2.303447in}}%
\pgfpathlineto{\pgfqpoint{2.850072in}{2.303374in}}%
\pgfpathlineto{\pgfqpoint{2.851003in}{2.303431in}}%
\pgfpathlineto{\pgfqpoint{2.855599in}{2.303408in}}%
\pgfpathlineto{\pgfqpoint{2.855601in}{2.303422in}}%
\pgfpathlineto{\pgfqpoint{2.856615in}{2.303446in}}%
\pgfpathlineto{\pgfqpoint{2.856384in}{2.303328in}}%
\pgfpathlineto{\pgfqpoint{2.856712in}{2.303421in}}%
\pgfpathlineto{\pgfqpoint{2.857789in}{2.303448in}}%
\pgfpathlineto{\pgfqpoint{2.856739in}{2.303379in}}%
\pgfpathlineto{\pgfqpoint{2.857828in}{2.303426in}}%
\pgfpathlineto{\pgfqpoint{2.858920in}{2.303381in}}%
\pgfpathlineto{\pgfqpoint{2.857837in}{2.303445in}}%
\pgfpathlineto{\pgfqpoint{2.858944in}{2.303422in}}%
\pgfpathlineto{\pgfqpoint{2.859890in}{2.303451in}}%
\pgfpathlineto{\pgfqpoint{2.859082in}{2.303352in}}%
\pgfpathlineto{\pgfqpoint{2.860055in}{2.303408in}}%
\pgfpathlineto{\pgfqpoint{2.860822in}{2.303362in}}%
\pgfpathlineto{\pgfqpoint{2.861074in}{2.303625in}}%
\pgfpathlineto{\pgfqpoint{2.861163in}{2.303490in}}%
\pgfpathlineto{\pgfqpoint{2.862364in}{2.304471in}}%
\pgfpathlineto{\pgfqpoint{2.861167in}{2.303476in}}%
\pgfpathlineto{\pgfqpoint{2.862408in}{2.304433in}}%
\pgfpathlineto{\pgfqpoint{2.863204in}{2.303753in}}%
\pgfpathlineto{\pgfqpoint{2.862415in}{2.304447in}}%
\pgfpathlineto{\pgfqpoint{2.863533in}{2.304144in}}%
\pgfpathlineto{\pgfqpoint{2.864908in}{2.304765in}}%
\pgfpathlineto{\pgfqpoint{2.863966in}{2.303894in}}%
\pgfpathlineto{\pgfqpoint{2.865065in}{2.304638in}}%
\pgfpathlineto{\pgfqpoint{2.865122in}{2.304550in}}%
\pgfpathlineto{\pgfqpoint{2.865823in}{2.305291in}}%
\pgfpathlineto{\pgfqpoint{2.866156in}{2.305180in}}%
\pgfpathlineto{\pgfqpoint{2.866864in}{2.305806in}}%
\pgfpathlineto{\pgfqpoint{2.866210in}{2.305138in}}%
\pgfpathlineto{\pgfqpoint{2.867335in}{2.305488in}}%
\pgfpathlineto{\pgfqpoint{2.868293in}{2.305283in}}%
\pgfpathlineto{\pgfqpoint{2.867842in}{2.305760in}}%
\pgfpathlineto{\pgfqpoint{2.868454in}{2.305354in}}%
\pgfpathlineto{\pgfqpoint{2.869729in}{2.304839in}}%
\pgfpathlineto{\pgfqpoint{2.869734in}{2.304909in}}%
\pgfpathlineto{\pgfqpoint{2.870709in}{2.305133in}}%
\pgfpathlineto{\pgfqpoint{2.870178in}{2.304757in}}%
\pgfpathlineto{\pgfqpoint{2.870852in}{2.304999in}}%
\pgfpathlineto{\pgfqpoint{2.871832in}{2.305463in}}%
\pgfpathlineto{\pgfqpoint{2.872003in}{2.305249in}}%
\pgfpathlineto{\pgfqpoint{2.872561in}{2.304939in}}%
\pgfpathlineto{\pgfqpoint{2.872946in}{2.305469in}}%
\pgfpathlineto{\pgfqpoint{2.873120in}{2.305111in}}%
\pgfpathlineto{\pgfqpoint{2.873121in}{2.305127in}}%
\pgfpathlineto{\pgfqpoint{2.873876in}{2.304530in}}%
\pgfpathlineto{\pgfqpoint{2.874222in}{2.304891in}}%
\pgfpathlineto{\pgfqpoint{2.877543in}{2.304878in}}%
\pgfpathlineto{\pgfqpoint{2.874224in}{2.304878in}}%
\pgfpathlineto{\pgfqpoint{2.877556in}{2.304795in}}%
\pgfpathlineto{\pgfqpoint{2.877672in}{2.304697in}}%
\pgfpathlineto{\pgfqpoint{2.878225in}{2.305467in}}%
\pgfpathlineto{\pgfqpoint{2.878639in}{2.305271in}}%
\pgfpathlineto{\pgfqpoint{2.879102in}{2.305643in}}%
\pgfpathlineto{\pgfqpoint{2.878657in}{2.305246in}}%
\pgfpathlineto{\pgfqpoint{2.879761in}{2.305389in}}%
\pgfpathlineto{\pgfqpoint{2.880878in}{2.305216in}}%
\pgfpathlineto{\pgfqpoint{2.879770in}{2.305422in}}%
\pgfpathlineto{\pgfqpoint{2.880950in}{2.305406in}}%
\pgfpathlineto{\pgfqpoint{2.881576in}{2.305481in}}%
\pgfpathlineto{\pgfqpoint{2.881744in}{2.305143in}}%
\pgfpathlineto{\pgfqpoint{2.882057in}{2.305314in}}%
\pgfpathlineto{\pgfqpoint{2.883160in}{2.304880in}}%
\pgfpathlineto{\pgfqpoint{2.882500in}{2.305414in}}%
\pgfpathlineto{\pgfqpoint{2.883184in}{2.305021in}}%
\pgfpathlineto{\pgfqpoint{2.883988in}{2.305280in}}%
\pgfpathlineto{\pgfqpoint{2.883309in}{2.304721in}}%
\pgfpathlineto{\pgfqpoint{2.884285in}{2.304840in}}%
\pgfpathlineto{\pgfqpoint{2.884926in}{2.304505in}}%
\pgfpathlineto{\pgfqpoint{2.884339in}{2.304861in}}%
\pgfpathlineto{\pgfqpoint{2.885402in}{2.304778in}}%
\pgfpathlineto{\pgfqpoint{2.886602in}{2.304848in}}%
\pgfpathlineto{\pgfqpoint{2.885739in}{2.305190in}}%
\pgfpathlineto{\pgfqpoint{2.886608in}{2.304894in}}%
\pgfpathlineto{\pgfqpoint{2.887629in}{2.305401in}}%
\pgfpathlineto{\pgfqpoint{2.887087in}{2.304714in}}%
\pgfpathlineto{\pgfqpoint{2.887797in}{2.305175in}}%
\pgfpathlineto{\pgfqpoint{2.888868in}{2.305412in}}%
\pgfpathlineto{\pgfqpoint{2.888129in}{2.305706in}}%
\pgfpathlineto{\pgfqpoint{2.888872in}{2.305450in}}%
\pgfpathlineto{\pgfqpoint{2.888973in}{2.305719in}}%
\pgfpathlineto{\pgfqpoint{2.889286in}{2.305189in}}%
\pgfpathlineto{\pgfqpoint{2.889981in}{2.305404in}}%
\pgfpathlineto{\pgfqpoint{2.890615in}{2.305315in}}%
\pgfpathlineto{\pgfqpoint{2.891082in}{2.305817in}}%
\pgfpathlineto{\pgfqpoint{2.891085in}{2.305799in}}%
\pgfpathlineto{\pgfqpoint{2.892251in}{2.305704in}}%
\pgfpathlineto{\pgfqpoint{2.891090in}{2.305821in}}%
\pgfpathlineto{\pgfqpoint{2.892264in}{2.305750in}}%
\pgfpathlineto{\pgfqpoint{2.893187in}{2.306772in}}%
\pgfpathlineto{\pgfqpoint{2.893434in}{2.306453in}}%
\pgfpathlineto{\pgfqpoint{2.894153in}{2.306543in}}%
\pgfpathlineto{\pgfqpoint{2.893734in}{2.306198in}}%
\pgfpathlineto{\pgfqpoint{2.894540in}{2.306305in}}%
\pgfpathlineto{\pgfqpoint{2.895177in}{2.305730in}}%
\pgfpathlineto{\pgfqpoint{2.894551in}{2.306313in}}%
\pgfpathlineto{\pgfqpoint{2.895672in}{2.306080in}}%
\pgfpathlineto{\pgfqpoint{2.896639in}{2.306540in}}%
\pgfpathlineto{\pgfqpoint{2.896792in}{2.306262in}}%
\pgfpathlineto{\pgfqpoint{2.897113in}{2.306045in}}%
\pgfpathlineto{\pgfqpoint{2.897839in}{2.306694in}}%
\pgfpathlineto{\pgfqpoint{2.897888in}{2.306656in}}%
\pgfpathlineto{\pgfqpoint{2.899184in}{2.306560in}}%
\pgfpathlineto{\pgfqpoint{2.897893in}{2.306641in}}%
\pgfpathlineto{\pgfqpoint{2.899189in}{2.306516in}}%
\pgfpathlineto{\pgfqpoint{2.899743in}{2.306208in}}%
\pgfpathlineto{\pgfqpoint{2.899279in}{2.306788in}}%
\pgfpathlineto{\pgfqpoint{2.900297in}{2.306609in}}%
\pgfpathlineto{\pgfqpoint{2.900302in}{2.306639in}}%
\pgfpathlineto{\pgfqpoint{2.901084in}{2.306123in}}%
\pgfpathlineto{\pgfqpoint{2.901397in}{2.306396in}}%
\pgfpathlineto{\pgfqpoint{2.902461in}{2.305932in}}%
\pgfpathlineto{\pgfqpoint{2.901440in}{2.306612in}}%
\pgfpathlineto{\pgfqpoint{2.902560in}{2.306110in}}%
\pgfpathlineto{\pgfqpoint{2.903677in}{2.306583in}}%
\pgfpathlineto{\pgfqpoint{2.902759in}{2.305918in}}%
\pgfpathlineto{\pgfqpoint{2.903714in}{2.306565in}}%
\pgfpathlineto{\pgfqpoint{2.904574in}{2.306170in}}%
\pgfpathlineto{\pgfqpoint{2.904359in}{2.305943in}}%
\pgfpathlineto{\pgfqpoint{2.904629in}{2.305924in}}%
\pgfpathlineto{\pgfqpoint{2.905238in}{2.305495in}}%
\pgfpathlineto{\pgfqpoint{2.904912in}{2.306358in}}%
\pgfpathlineto{\pgfqpoint{2.905721in}{2.306120in}}%
\pgfpathlineto{\pgfqpoint{2.906337in}{2.306162in}}%
\pgfpathlineto{\pgfqpoint{2.906056in}{2.305703in}}%
\pgfpathlineto{\pgfqpoint{2.906828in}{2.306017in}}%
\pgfpathlineto{\pgfqpoint{2.907929in}{2.306113in}}%
\pgfpathlineto{\pgfqpoint{2.907396in}{2.306478in}}%
\pgfpathlineto{\pgfqpoint{2.907934in}{2.306135in}}%
\pgfpathlineto{\pgfqpoint{2.908274in}{2.306204in}}%
\pgfpathlineto{\pgfqpoint{2.908842in}{2.305596in}}%
\pgfpathlineto{\pgfqpoint{2.908996in}{2.305653in}}%
\pgfpathlineto{\pgfqpoint{2.909613in}{2.305340in}}%
\pgfpathlineto{\pgfqpoint{2.909151in}{2.305702in}}%
\pgfpathlineto{\pgfqpoint{2.910109in}{2.305660in}}%
\pgfpathlineto{\pgfqpoint{2.910973in}{2.305913in}}%
\pgfpathlineto{\pgfqpoint{2.910580in}{2.305377in}}%
\pgfpathlineto{\pgfqpoint{2.911234in}{2.305712in}}%
\pgfpathlineto{\pgfqpoint{2.911962in}{2.305693in}}%
\pgfpathlineto{\pgfqpoint{2.911541in}{2.305972in}}%
\pgfpathlineto{\pgfqpoint{2.912303in}{2.306056in}}%
\pgfpathlineto{\pgfqpoint{2.912871in}{2.306208in}}%
\pgfpathlineto{\pgfqpoint{2.912533in}{2.305801in}}%
\pgfpathlineto{\pgfqpoint{2.913412in}{2.305948in}}%
\pgfpathlineto{\pgfqpoint{2.913448in}{2.305868in}}%
\pgfpathlineto{\pgfqpoint{2.913944in}{2.306546in}}%
\pgfpathlineto{\pgfqpoint{2.914477in}{2.306523in}}%
\pgfpathlineto{\pgfqpoint{2.915575in}{2.306505in}}%
\pgfpathlineto{\pgfqpoint{2.914651in}{2.306380in}}%
\pgfpathlineto{\pgfqpoint{2.915641in}{2.306300in}}%
\pgfpathlineto{\pgfqpoint{2.916533in}{2.306013in}}%
\pgfpathlineto{\pgfqpoint{2.915655in}{2.306316in}}%
\pgfpathlineto{\pgfqpoint{2.916802in}{2.306289in}}%
\pgfpathlineto{\pgfqpoint{2.916944in}{2.306345in}}%
\pgfpathlineto{\pgfqpoint{2.917077in}{2.303377in}}%
\pgfpathlineto{\pgfqpoint{2.917525in}{2.303420in}}%
\pgfpathlineto{\pgfqpoint{2.918671in}{2.303396in}}%
\pgfpathlineto{\pgfqpoint{2.917530in}{2.303441in}}%
\pgfpathlineto{\pgfqpoint{2.918673in}{2.303424in}}%
\pgfpathlineto{\pgfqpoint{2.919757in}{2.303455in}}%
\pgfpathlineto{\pgfqpoint{2.919047in}{2.303344in}}%
\pgfpathlineto{\pgfqpoint{2.919784in}{2.303425in}}%
\pgfpathlineto{\pgfqpoint{2.920851in}{2.303351in}}%
\pgfpathlineto{\pgfqpoint{2.919859in}{2.303449in}}%
\pgfpathlineto{\pgfqpoint{2.920896in}{2.303409in}}%
\pgfpathlineto{\pgfqpoint{2.921927in}{2.303446in}}%
\pgfpathlineto{\pgfqpoint{2.921303in}{2.303328in}}%
\pgfpathlineto{\pgfqpoint{2.922006in}{2.303399in}}%
\pgfpathlineto{\pgfqpoint{2.923196in}{2.303398in}}%
\pgfpathlineto{\pgfqpoint{2.922009in}{2.303450in}}%
\pgfpathlineto{\pgfqpoint{2.923203in}{2.303423in}}%
\pgfpathlineto{\pgfqpoint{2.924277in}{2.303443in}}%
\pgfpathlineto{\pgfqpoint{2.923277in}{2.303374in}}%
\pgfpathlineto{\pgfqpoint{2.924318in}{2.303430in}}%
\pgfpathlineto{\pgfqpoint{2.925404in}{2.303326in}}%
\pgfpathlineto{\pgfqpoint{2.924586in}{2.303457in}}%
\pgfpathlineto{\pgfqpoint{2.925430in}{2.303398in}}%
\pgfpathlineto{\pgfqpoint{2.926464in}{2.303446in}}%
\pgfpathlineto{\pgfqpoint{2.925462in}{2.303373in}}%
\pgfpathlineto{\pgfqpoint{2.926550in}{2.303427in}}%
\pgfpathlineto{\pgfqpoint{2.926790in}{2.303335in}}%
\pgfpathlineto{\pgfqpoint{2.926576in}{2.303444in}}%
\pgfpathlineto{\pgfqpoint{2.927661in}{2.303410in}}%
\pgfpathlineto{\pgfqpoint{2.928655in}{2.303449in}}%
\pgfpathlineto{\pgfqpoint{2.927680in}{2.303379in}}%
\pgfpathlineto{\pgfqpoint{2.928773in}{2.303427in}}%
\pgfpathlineto{\pgfqpoint{2.929692in}{2.303368in}}%
\pgfpathlineto{\pgfqpoint{2.928882in}{2.303453in}}%
\pgfpathlineto{\pgfqpoint{2.929886in}{2.303404in}}%
\pgfpathlineto{\pgfqpoint{2.930984in}{2.303437in}}%
\pgfpathlineto{\pgfqpoint{2.930016in}{2.303333in}}%
\pgfpathlineto{\pgfqpoint{2.931003in}{2.303394in}}%
\pgfpathlineto{\pgfqpoint{2.931779in}{2.303321in}}%
\pgfpathlineto{\pgfqpoint{2.931057in}{2.303447in}}%
\pgfpathlineto{\pgfqpoint{2.932114in}{2.303408in}}%
\pgfpathlineto{\pgfqpoint{2.933194in}{2.303449in}}%
\pgfpathlineto{\pgfqpoint{2.932203in}{2.303366in}}%
\pgfpathlineto{\pgfqpoint{2.933226in}{2.303402in}}%
\pgfpathlineto{\pgfqpoint{2.934333in}{2.303437in}}%
\pgfpathlineto{\pgfqpoint{2.933280in}{2.303374in}}%
\pgfpathlineto{\pgfqpoint{2.934344in}{2.303426in}}%
\pgfpathlineto{\pgfqpoint{2.934830in}{2.303333in}}%
\pgfpathlineto{\pgfqpoint{2.934565in}{2.303449in}}%
\pgfpathlineto{\pgfqpoint{2.935455in}{2.303419in}}%
\pgfpathlineto{\pgfqpoint{2.936131in}{2.303757in}}%
\pgfpathlineto{\pgfqpoint{2.935461in}{2.303385in}}%
\pgfpathlineto{\pgfqpoint{2.936568in}{2.303414in}}%
\pgfpathlineto{\pgfqpoint{2.937621in}{2.303231in}}%
\pgfpathlineto{\pgfqpoint{2.936575in}{2.303446in}}%
\pgfpathlineto{\pgfqpoint{2.937738in}{2.303462in}}%
\pgfpathlineto{\pgfqpoint{2.938613in}{2.304128in}}%
\pgfpathlineto{\pgfqpoint{2.937742in}{2.303445in}}%
\pgfpathlineto{\pgfqpoint{2.938917in}{2.303906in}}%
\pgfpathlineto{\pgfqpoint{2.939817in}{2.303183in}}%
\pgfpathlineto{\pgfqpoint{2.939001in}{2.304031in}}%
\pgfpathlineto{\pgfqpoint{2.940090in}{2.303345in}}%
\pgfpathlineto{\pgfqpoint{2.941047in}{2.304024in}}%
\pgfpathlineto{\pgfqpoint{2.940129in}{2.303309in}}%
\pgfpathlineto{\pgfqpoint{2.941230in}{2.303871in}}%
\pgfpathlineto{\pgfqpoint{2.942359in}{2.303299in}}%
\pgfpathlineto{\pgfqpoint{2.941275in}{2.303941in}}%
\pgfpathlineto{\pgfqpoint{2.942364in}{2.303317in}}%
\pgfpathlineto{\pgfqpoint{2.942828in}{2.303521in}}%
\pgfpathlineto{\pgfqpoint{2.943035in}{2.303084in}}%
\pgfpathlineto{\pgfqpoint{2.943483in}{2.303366in}}%
\pgfpathlineto{\pgfqpoint{2.943608in}{2.303289in}}%
\pgfpathlineto{\pgfqpoint{2.944244in}{2.304065in}}%
\pgfpathlineto{\pgfqpoint{2.944550in}{2.303903in}}%
\pgfpathlineto{\pgfqpoint{2.945702in}{2.303820in}}%
\pgfpathlineto{\pgfqpoint{2.944771in}{2.303654in}}%
\pgfpathlineto{\pgfqpoint{2.945722in}{2.303736in}}%
\pgfpathlineto{\pgfqpoint{2.946787in}{2.303263in}}%
\pgfpathlineto{\pgfqpoint{2.945776in}{2.303873in}}%
\pgfpathlineto{\pgfqpoint{2.946850in}{2.303420in}}%
\pgfpathlineto{\pgfqpoint{2.946891in}{2.303456in}}%
\pgfpathlineto{\pgfqpoint{2.947827in}{2.302840in}}%
\pgfpathlineto{\pgfqpoint{2.947944in}{2.303017in}}%
\pgfpathlineto{\pgfqpoint{2.949015in}{2.303264in}}%
\pgfpathlineto{\pgfqpoint{2.948553in}{2.302921in}}%
\pgfpathlineto{\pgfqpoint{2.949070in}{2.303166in}}%
\pgfpathlineto{\pgfqpoint{2.949102in}{2.303057in}}%
\pgfpathlineto{\pgfqpoint{2.949842in}{2.303789in}}%
\pgfpathlineto{\pgfqpoint{2.950156in}{2.303669in}}%
\pgfpathlineto{\pgfqpoint{2.950807in}{2.303789in}}%
\pgfpathlineto{\pgfqpoint{2.950260in}{2.303456in}}%
\pgfpathlineto{\pgfqpoint{2.951242in}{2.303392in}}%
\pgfpathlineto{\pgfqpoint{2.953443in}{2.303359in}}%
\pgfpathlineto{\pgfqpoint{2.951245in}{2.303411in}}%
\pgfpathlineto{\pgfqpoint{2.953460in}{2.303455in}}%
\pgfpathlineto{\pgfqpoint{2.953500in}{2.303566in}}%
\pgfpathlineto{\pgfqpoint{2.954439in}{2.303020in}}%
\pgfpathlineto{\pgfqpoint{2.954554in}{2.303087in}}%
\pgfpathlineto{\pgfqpoint{2.955652in}{2.302793in}}%
\pgfpathlineto{\pgfqpoint{2.955192in}{2.303184in}}%
\pgfpathlineto{\pgfqpoint{2.955673in}{2.302897in}}%
\pgfpathlineto{\pgfqpoint{2.956989in}{2.302907in}}%
\pgfpathlineto{\pgfqpoint{2.955832in}{2.302543in}}%
\pgfpathlineto{\pgfqpoint{2.957044in}{2.302702in}}%
\pgfpathlineto{\pgfqpoint{2.957727in}{2.302504in}}%
\pgfpathlineto{\pgfqpoint{2.957284in}{2.302994in}}%
\pgfpathlineto{\pgfqpoint{2.958162in}{2.302603in}}%
\pgfpathlineto{\pgfqpoint{2.959340in}{2.303351in}}%
\pgfpathlineto{\pgfqpoint{2.958182in}{2.302538in}}%
\pgfpathlineto{\pgfqpoint{2.959369in}{2.303324in}}%
\pgfpathlineto{\pgfqpoint{2.959738in}{2.303036in}}%
\pgfpathlineto{\pgfqpoint{2.960368in}{2.303542in}}%
\pgfpathlineto{\pgfqpoint{2.960474in}{2.303484in}}%
\pgfpathlineto{\pgfqpoint{2.961586in}{2.303794in}}%
\pgfpathlineto{\pgfqpoint{2.960735in}{2.303319in}}%
\pgfpathlineto{\pgfqpoint{2.961594in}{2.303762in}}%
\pgfpathlineto{\pgfqpoint{2.962112in}{2.302831in}}%
\pgfpathlineto{\pgfqpoint{2.961596in}{2.303767in}}%
\pgfpathlineto{\pgfqpoint{2.962852in}{2.303363in}}%
\pgfpathlineto{\pgfqpoint{2.964276in}{2.303662in}}%
\pgfpathlineto{\pgfqpoint{2.964342in}{2.303528in}}%
\pgfpathlineto{\pgfqpoint{2.965034in}{2.303068in}}%
\pgfpathlineto{\pgfqpoint{2.964368in}{2.303596in}}%
\pgfpathlineto{\pgfqpoint{2.965463in}{2.303463in}}%
\pgfpathlineto{\pgfqpoint{2.966353in}{2.303666in}}%
\pgfpathlineto{\pgfqpoint{2.966100in}{2.303294in}}%
\pgfpathlineto{\pgfqpoint{2.966581in}{2.303487in}}%
\pgfpathlineto{\pgfqpoint{2.967389in}{2.303201in}}%
\pgfpathlineto{\pgfqpoint{2.966585in}{2.303490in}}%
\pgfpathlineto{\pgfqpoint{2.967677in}{2.303737in}}%
\pgfpathlineto{\pgfqpoint{2.967742in}{2.303974in}}%
\pgfpathlineto{\pgfqpoint{2.968310in}{2.303232in}}%
\pgfpathlineto{\pgfqpoint{2.968792in}{2.303853in}}%
\pgfpathlineto{\pgfqpoint{2.969811in}{2.303903in}}%
\pgfpathlineto{\pgfqpoint{2.969899in}{2.303584in}}%
\pgfpathlineto{\pgfqpoint{2.970528in}{2.303311in}}%
\pgfpathlineto{\pgfqpoint{2.970033in}{2.303831in}}%
\pgfpathlineto{\pgfqpoint{2.971020in}{2.303480in}}%
\pgfpathlineto{\pgfqpoint{2.971024in}{2.303505in}}%
\pgfpathlineto{\pgfqpoint{2.971921in}{2.302627in}}%
\pgfpathlineto{\pgfqpoint{2.972084in}{2.302801in}}%
\pgfpathlineto{\pgfqpoint{2.972749in}{2.302535in}}%
\pgfpathlineto{\pgfqpoint{2.972110in}{2.302853in}}%
\pgfpathlineto{\pgfqpoint{2.973189in}{2.302943in}}%
\pgfpathlineto{\pgfqpoint{2.974321in}{2.303157in}}%
\pgfpathlineto{\pgfqpoint{2.973199in}{2.302973in}}%
\pgfpathlineto{\pgfqpoint{2.974350in}{2.303325in}}%
\pgfpathlineto{\pgfqpoint{2.975048in}{2.303432in}}%
\pgfpathlineto{\pgfqpoint{2.974603in}{2.302964in}}%
\pgfpathlineto{\pgfqpoint{2.975454in}{2.303112in}}%
\pgfpathlineto{\pgfqpoint{2.976782in}{2.302520in}}%
\pgfpathlineto{\pgfqpoint{2.975462in}{2.303138in}}%
\pgfpathlineto{\pgfqpoint{2.976797in}{2.302565in}}%
\pgfpathlineto{\pgfqpoint{2.977406in}{2.302804in}}%
\pgfpathlineto{\pgfqpoint{2.977123in}{2.302379in}}%
\pgfpathlineto{\pgfqpoint{2.977905in}{2.302485in}}%
\pgfpathlineto{\pgfqpoint{2.978810in}{2.302948in}}%
\pgfpathlineto{\pgfqpoint{2.978271in}{2.302330in}}%
\pgfpathlineto{\pgfqpoint{2.979044in}{2.302624in}}%
\pgfpathlineto{\pgfqpoint{2.986907in}{2.304091in}}%
\pgfpathlineto{\pgfqpoint{2.979047in}{2.302606in}}%
\pgfpathlineto{\pgfqpoint{2.986913in}{2.304036in}}%
\pgfpathlineto{\pgfqpoint{2.987988in}{2.303221in}}%
\pgfpathlineto{\pgfqpoint{2.987116in}{2.304300in}}%
\pgfpathlineto{\pgfqpoint{2.988074in}{2.303313in}}%
\pgfpathlineto{\pgfqpoint{2.988405in}{2.303381in}}%
\pgfpathlineto{\pgfqpoint{2.988915in}{2.302994in}}%
\pgfpathlineto{\pgfqpoint{2.989185in}{2.303298in}}%
\pgfpathlineto{\pgfqpoint{2.989926in}{2.303631in}}%
\pgfpathlineto{\pgfqpoint{2.989469in}{2.303043in}}%
\pgfpathlineto{\pgfqpoint{2.990315in}{2.303564in}}%
\pgfpathlineto{\pgfqpoint{2.990653in}{2.303543in}}%
\pgfpathlineto{\pgfqpoint{2.990987in}{2.303960in}}%
\pgfpathlineto{\pgfqpoint{2.991413in}{2.303792in}}%
\pgfpathlineto{\pgfqpoint{2.991678in}{2.304031in}}%
\pgfpathlineto{\pgfqpoint{2.991871in}{2.303384in}}%
\pgfpathlineto{\pgfqpoint{2.992448in}{2.303392in}}%
\pgfpathlineto{\pgfqpoint{2.993496in}{2.303367in}}%
\pgfpathlineto{\pgfqpoint{2.992455in}{2.303444in}}%
\pgfpathlineto{\pgfqpoint{2.993561in}{2.303395in}}%
\pgfpathlineto{\pgfqpoint{2.994598in}{2.303443in}}%
\pgfpathlineto{\pgfqpoint{2.993660in}{2.303330in}}%
\pgfpathlineto{\pgfqpoint{2.994673in}{2.303415in}}%
\pgfpathlineto{\pgfqpoint{2.995744in}{2.303445in}}%
\pgfpathlineto{\pgfqpoint{2.994758in}{2.303351in}}%
\pgfpathlineto{\pgfqpoint{2.995787in}{2.303404in}}%
\pgfpathlineto{\pgfqpoint{2.996824in}{2.303451in}}%
\pgfpathlineto{\pgfqpoint{2.995837in}{2.303366in}}%
\pgfpathlineto{\pgfqpoint{2.996903in}{2.303431in}}%
\pgfpathlineto{\pgfqpoint{2.997966in}{2.303367in}}%
\pgfpathlineto{\pgfqpoint{2.996969in}{2.303448in}}%
\pgfpathlineto{\pgfqpoint{2.998016in}{2.303390in}}%
\pgfpathlineto{\pgfqpoint{2.999138in}{2.303434in}}%
\pgfpathlineto{\pgfqpoint{2.998034in}{2.303382in}}%
\pgfpathlineto{\pgfqpoint{2.999141in}{2.303412in}}%
\pgfpathlineto{\pgfqpoint{3.000266in}{2.303363in}}%
\pgfpathlineto{\pgfqpoint{2.999256in}{2.303458in}}%
\pgfpathlineto{\pgfqpoint{3.000266in}{2.303372in}}%
\pgfpathlineto{\pgfqpoint{3.001353in}{2.303442in}}%
\pgfpathlineto{\pgfqpoint{3.000287in}{2.303369in}}%
\pgfpathlineto{\pgfqpoint{3.001382in}{2.303411in}}%
\pgfpathlineto{\pgfqpoint{3.002508in}{2.303390in}}%
\pgfpathlineto{\pgfqpoint{3.001412in}{2.303444in}}%
\pgfpathlineto{\pgfqpoint{3.002531in}{2.303406in}}%
\pgfpathlineto{\pgfqpoint{3.003641in}{2.303446in}}%
\pgfpathlineto{\pgfqpoint{3.003042in}{2.303351in}}%
\pgfpathlineto{\pgfqpoint{3.003642in}{2.303426in}}%
\pgfpathlineto{\pgfqpoint{3.004730in}{2.303360in}}%
\pgfpathlineto{\pgfqpoint{3.003700in}{2.303450in}}%
\pgfpathlineto{\pgfqpoint{3.004754in}{2.303412in}}%
\pgfpathlineto{\pgfqpoint{3.005806in}{2.303447in}}%
\pgfpathlineto{\pgfqpoint{3.004866in}{2.303378in}}%
\pgfpathlineto{\pgfqpoint{3.005875in}{2.303408in}}%
\pgfpathlineto{\pgfqpoint{3.006736in}{2.303338in}}%
\pgfpathlineto{\pgfqpoint{3.005880in}{2.303448in}}%
\pgfpathlineto{\pgfqpoint{3.006988in}{2.303416in}}%
\pgfpathlineto{\pgfqpoint{3.008659in}{2.303384in}}%
\pgfpathlineto{\pgfqpoint{3.006992in}{2.303434in}}%
\pgfpathlineto{\pgfqpoint{3.008664in}{2.303423in}}%
\pgfpathlineto{\pgfqpoint{3.009801in}{2.303443in}}%
\pgfpathlineto{\pgfqpoint{3.008679in}{2.303372in}}%
\pgfpathlineto{\pgfqpoint{3.009802in}{2.303405in}}%
\pgfpathlineto{\pgfqpoint{3.010966in}{2.303169in}}%
\pgfpathlineto{\pgfqpoint{3.009872in}{2.303444in}}%
\pgfpathlineto{\pgfqpoint{3.010968in}{2.303174in}}%
\pgfpathlineto{\pgfqpoint{3.010969in}{2.303184in}}%
\pgfpathlineto{\pgfqpoint{3.011545in}{2.302383in}}%
\pgfpathlineto{\pgfqpoint{3.012033in}{2.302517in}}%
\pgfpathlineto{\pgfqpoint{3.012043in}{2.302472in}}%
\pgfpathlineto{\pgfqpoint{3.013127in}{2.303063in}}%
\pgfpathlineto{\pgfqpoint{3.014253in}{2.303319in}}%
\pgfpathlineto{\pgfqpoint{3.013331in}{2.303008in}}%
\pgfpathlineto{\pgfqpoint{3.014256in}{2.303276in}}%
\pgfpathlineto{\pgfqpoint{3.015396in}{2.302305in}}%
\pgfpathlineto{\pgfqpoint{3.014292in}{2.303398in}}%
\pgfpathlineto{\pgfqpoint{3.015433in}{2.302336in}}%
\pgfpathlineto{\pgfqpoint{3.016434in}{2.302529in}}%
\pgfpathlineto{\pgfqpoint{3.015929in}{2.302176in}}%
\pgfpathlineto{\pgfqpoint{3.016546in}{2.302358in}}%
\pgfpathlineto{\pgfqpoint{3.016547in}{2.302356in}}%
\pgfpathlineto{\pgfqpoint{3.017278in}{2.303329in}}%
\pgfpathlineto{\pgfqpoint{3.017353in}{2.303278in}}%
\pgfpathlineto{\pgfqpoint{3.018363in}{2.303509in}}%
\pgfpathlineto{\pgfqpoint{3.017584in}{2.302953in}}%
\pgfpathlineto{\pgfqpoint{3.018471in}{2.303374in}}%
\pgfpathlineto{\pgfqpoint{3.021697in}{2.304110in}}%
\pgfpathlineto{\pgfqpoint{3.021709in}{2.304064in}}%
\pgfpathlineto{\pgfqpoint{3.021850in}{2.303774in}}%
\pgfpathlineto{\pgfqpoint{3.022566in}{2.304258in}}%
\pgfpathlineto{\pgfqpoint{3.022822in}{2.304027in}}%
\pgfpathlineto{\pgfqpoint{3.022844in}{2.303993in}}%
\pgfpathlineto{\pgfqpoint{3.023512in}{2.305026in}}%
\pgfpathlineto{\pgfqpoint{3.023864in}{2.304798in}}%
\pgfpathlineto{\pgfqpoint{3.024453in}{2.305154in}}%
\pgfpathlineto{\pgfqpoint{3.024925in}{2.304470in}}%
\pgfpathlineto{\pgfqpoint{3.024950in}{2.304541in}}%
\pgfpathlineto{\pgfqpoint{3.025009in}{2.304383in}}%
\pgfpathlineto{\pgfqpoint{3.025759in}{2.305136in}}%
\pgfpathlineto{\pgfqpoint{3.026037in}{2.304804in}}%
\pgfpathlineto{\pgfqpoint{3.027101in}{2.304951in}}%
\pgfpathlineto{\pgfqpoint{3.026855in}{2.304540in}}%
\pgfpathlineto{\pgfqpoint{3.027152in}{2.304907in}}%
\pgfpathlineto{\pgfqpoint{3.027498in}{2.304665in}}%
\pgfpathlineto{\pgfqpoint{3.028158in}{2.305392in}}%
\pgfpathlineto{\pgfqpoint{3.028245in}{2.305285in}}%
\pgfpathlineto{\pgfqpoint{3.029537in}{2.305934in}}%
\pgfpathlineto{\pgfqpoint{3.028400in}{2.305146in}}%
\pgfpathlineto{\pgfqpoint{3.029558in}{2.305842in}}%
\pgfpathlineto{\pgfqpoint{3.030584in}{2.305657in}}%
\pgfpathlineto{\pgfqpoint{3.029778in}{2.305931in}}%
\pgfpathlineto{\pgfqpoint{3.030665in}{2.305937in}}%
\pgfpathlineto{\pgfqpoint{3.031420in}{2.306193in}}%
\pgfpathlineto{\pgfqpoint{3.031762in}{2.305673in}}%
\pgfpathlineto{\pgfqpoint{3.031834in}{2.305797in}}%
\pgfpathlineto{\pgfqpoint{3.032238in}{2.305192in}}%
\pgfpathlineto{\pgfqpoint{3.032844in}{2.305401in}}%
\pgfpathlineto{\pgfqpoint{3.033113in}{2.305226in}}%
\pgfpathlineto{\pgfqpoint{3.033245in}{2.305592in}}%
\pgfpathlineto{\pgfqpoint{3.033949in}{2.305537in}}%
\pgfpathlineto{\pgfqpoint{3.034752in}{2.305376in}}%
\pgfpathlineto{\pgfqpoint{3.033977in}{2.305640in}}%
\pgfpathlineto{\pgfqpoint{3.035080in}{2.305543in}}%
\pgfpathlineto{\pgfqpoint{3.036195in}{2.305587in}}%
\pgfpathlineto{\pgfqpoint{3.035696in}{2.305172in}}%
\pgfpathlineto{\pgfqpoint{3.036199in}{2.305586in}}%
\pgfpathlineto{\pgfqpoint{3.037323in}{2.305207in}}%
\pgfpathlineto{\pgfqpoint{3.037327in}{2.305221in}}%
\pgfpathlineto{\pgfqpoint{3.037439in}{2.305095in}}%
\pgfpathlineto{\pgfqpoint{3.037997in}{2.305714in}}%
\pgfpathlineto{\pgfqpoint{3.038340in}{2.305684in}}%
\pgfpathlineto{\pgfqpoint{3.038388in}{2.305799in}}%
\pgfpathlineto{\pgfqpoint{3.039414in}{2.305111in}}%
\pgfpathlineto{\pgfqpoint{3.039435in}{2.305155in}}%
\pgfpathlineto{\pgfqpoint{3.040653in}{2.304582in}}%
\pgfpathlineto{\pgfqpoint{3.039511in}{2.305206in}}%
\pgfpathlineto{\pgfqpoint{3.040715in}{2.304752in}}%
\pgfpathlineto{\pgfqpoint{3.041809in}{2.305122in}}%
\pgfpathlineto{\pgfqpoint{3.041152in}{2.304524in}}%
\pgfpathlineto{\pgfqpoint{3.041903in}{2.304973in}}%
\pgfpathlineto{\pgfqpoint{3.042011in}{2.304910in}}%
\pgfpathlineto{\pgfqpoint{3.042616in}{2.305460in}}%
\pgfpathlineto{\pgfqpoint{3.043012in}{2.305055in}}%
\pgfpathlineto{\pgfqpoint{3.043485in}{2.305375in}}%
\pgfpathlineto{\pgfqpoint{3.043097in}{2.305024in}}%
\pgfpathlineto{\pgfqpoint{3.044124in}{2.305073in}}%
\pgfpathlineto{\pgfqpoint{3.045024in}{2.304734in}}%
\pgfpathlineto{\pgfqpoint{3.044321in}{2.305330in}}%
\pgfpathlineto{\pgfqpoint{3.045280in}{2.305007in}}%
\pgfpathlineto{\pgfqpoint{3.045291in}{2.305061in}}%
\pgfpathlineto{\pgfqpoint{3.045756in}{2.304178in}}%
\pgfpathlineto{\pgfqpoint{3.046334in}{2.304217in}}%
\pgfpathlineto{\pgfqpoint{3.046486in}{2.304029in}}%
\pgfpathlineto{\pgfqpoint{3.047127in}{2.304648in}}%
\pgfpathlineto{\pgfqpoint{3.047440in}{2.304345in}}%
\pgfpathlineto{\pgfqpoint{3.048713in}{2.304703in}}%
\pgfpathlineto{\pgfqpoint{3.047442in}{2.304340in}}%
\pgfpathlineto{\pgfqpoint{3.048718in}{2.304669in}}%
\pgfpathlineto{\pgfqpoint{3.048808in}{2.304625in}}%
\pgfpathlineto{\pgfqpoint{3.049301in}{2.305464in}}%
\pgfpathlineto{\pgfqpoint{3.049808in}{2.305204in}}%
\pgfpathlineto{\pgfqpoint{3.050688in}{2.304640in}}%
\pgfpathlineto{\pgfqpoint{3.050942in}{2.304995in}}%
\pgfpathlineto{\pgfqpoint{3.051807in}{2.305765in}}%
\pgfpathlineto{\pgfqpoint{3.050954in}{2.304963in}}%
\pgfpathlineto{\pgfqpoint{3.052239in}{2.305278in}}%
\pgfpathlineto{\pgfqpoint{3.053613in}{2.305242in}}%
\pgfpathlineto{\pgfqpoint{3.052246in}{2.305325in}}%
\pgfpathlineto{\pgfqpoint{3.053651in}{2.305346in}}%
\pgfpathlineto{\pgfqpoint{3.053652in}{2.305374in}}%
\pgfpathlineto{\pgfqpoint{3.054396in}{2.304589in}}%
\pgfpathlineto{\pgfqpoint{3.054737in}{2.304721in}}%
\pgfpathlineto{\pgfqpoint{3.055808in}{2.304149in}}%
\pgfpathlineto{\pgfqpoint{3.055879in}{2.304280in}}%
\pgfpathlineto{\pgfqpoint{3.056725in}{2.304827in}}%
\pgfpathlineto{\pgfqpoint{3.055882in}{2.304240in}}%
\pgfpathlineto{\pgfqpoint{3.057074in}{2.304601in}}%
\pgfpathlineto{\pgfqpoint{3.058136in}{2.304478in}}%
\pgfpathlineto{\pgfqpoint{3.057739in}{2.304926in}}%
\pgfpathlineto{\pgfqpoint{3.058184in}{2.304645in}}%
\pgfpathlineto{\pgfqpoint{3.059303in}{2.305178in}}%
\pgfpathlineto{\pgfqpoint{3.058315in}{2.304515in}}%
\pgfpathlineto{\pgfqpoint{3.059313in}{2.305141in}}%
\pgfpathlineto{\pgfqpoint{3.060234in}{2.304883in}}%
\pgfpathlineto{\pgfqpoint{3.059614in}{2.305447in}}%
\pgfpathlineto{\pgfqpoint{3.060427in}{2.305145in}}%
\pgfpathlineto{\pgfqpoint{3.061280in}{2.305656in}}%
\pgfpathlineto{\pgfqpoint{3.060706in}{2.304913in}}%
\pgfpathlineto{\pgfqpoint{3.061557in}{2.305530in}}%
\pgfpathlineto{\pgfqpoint{3.062464in}{2.304997in}}%
\pgfpathlineto{\pgfqpoint{3.061567in}{2.305595in}}%
\pgfpathlineto{\pgfqpoint{3.062706in}{2.305138in}}%
\pgfpathlineto{\pgfqpoint{3.063444in}{2.305662in}}%
\pgfpathlineto{\pgfqpoint{3.063825in}{2.305327in}}%
\pgfpathlineto{\pgfqpoint{3.064620in}{2.304673in}}%
\pgfpathlineto{\pgfqpoint{3.063829in}{2.305329in}}%
\pgfpathlineto{\pgfqpoint{3.064983in}{2.304755in}}%
\pgfpathlineto{\pgfqpoint{3.065519in}{2.304953in}}%
\pgfpathlineto{\pgfqpoint{3.066051in}{2.304345in}}%
\pgfpathlineto{\pgfqpoint{3.066074in}{2.304439in}}%
\pgfpathlineto{\pgfqpoint{3.067226in}{2.303372in}}%
\pgfpathlineto{\pgfqpoint{3.066695in}{2.305104in}}%
\pgfpathlineto{\pgfqpoint{3.067255in}{2.303413in}}%
\pgfpathlineto{\pgfqpoint{3.068303in}{2.303443in}}%
\pgfpathlineto{\pgfqpoint{3.067345in}{2.303375in}}%
\pgfpathlineto{\pgfqpoint{3.068366in}{2.303425in}}%
\pgfpathlineto{\pgfqpoint{3.069426in}{2.303369in}}%
\pgfpathlineto{\pgfqpoint{3.068470in}{2.303454in}}%
\pgfpathlineto{\pgfqpoint{3.069479in}{2.303418in}}%
\pgfpathlineto{\pgfqpoint{3.070598in}{2.303384in}}%
\pgfpathlineto{\pgfqpoint{3.069574in}{2.303452in}}%
\pgfpathlineto{\pgfqpoint{3.070600in}{2.303408in}}%
\pgfpathlineto{\pgfqpoint{3.072071in}{2.303425in}}%
\pgfpathlineto{\pgfqpoint{3.070604in}{2.303390in}}%
\pgfpathlineto{\pgfqpoint{3.072072in}{2.303399in}}%
\pgfpathlineto{\pgfqpoint{3.072888in}{2.303348in}}%
\pgfpathlineto{\pgfqpoint{3.072197in}{2.303457in}}%
\pgfpathlineto{\pgfqpoint{3.073183in}{2.303413in}}%
\pgfpathlineto{\pgfqpoint{3.074248in}{2.303446in}}%
\pgfpathlineto{\pgfqpoint{3.073335in}{2.303344in}}%
\pgfpathlineto{\pgfqpoint{3.074302in}{2.303414in}}%
\pgfpathlineto{\pgfqpoint{3.075268in}{2.303369in}}%
\pgfpathlineto{\pgfqpoint{3.074498in}{2.303448in}}%
\pgfpathlineto{\pgfqpoint{3.075413in}{2.303422in}}%
\pgfpathlineto{\pgfqpoint{3.076500in}{2.303445in}}%
\pgfpathlineto{\pgfqpoint{3.075675in}{2.303356in}}%
\pgfpathlineto{\pgfqpoint{3.076526in}{2.303432in}}%
\pgfpathlineto{\pgfqpoint{3.077470in}{2.303345in}}%
\pgfpathlineto{\pgfqpoint{3.076617in}{2.303443in}}%
\pgfpathlineto{\pgfqpoint{3.077642in}{2.303371in}}%
\pgfpathlineto{\pgfqpoint{3.078349in}{2.303462in}}%
\pgfpathlineto{\pgfqpoint{3.078754in}{2.303398in}}%
\pgfpathlineto{\pgfqpoint{3.079798in}{2.303443in}}%
\pgfpathlineto{\pgfqpoint{3.078789in}{2.303361in}}%
\pgfpathlineto{\pgfqpoint{3.079877in}{2.303432in}}%
\pgfpathlineto{\pgfqpoint{3.080673in}{2.303357in}}%
\pgfpathlineto{\pgfqpoint{3.080001in}{2.303450in}}%
\pgfpathlineto{\pgfqpoint{3.080989in}{2.303401in}}%
\pgfpathlineto{\pgfqpoint{3.082121in}{2.303389in}}%
\pgfpathlineto{\pgfqpoint{3.080995in}{2.303436in}}%
\pgfpathlineto{\pgfqpoint{3.082142in}{2.303413in}}%
\pgfpathlineto{\pgfqpoint{3.083209in}{2.303448in}}%
\pgfpathlineto{\pgfqpoint{3.082405in}{2.303360in}}%
\pgfpathlineto{\pgfqpoint{3.083254in}{2.303427in}}%
\pgfpathlineto{\pgfqpoint{3.084321in}{2.303364in}}%
\pgfpathlineto{\pgfqpoint{3.083404in}{2.303449in}}%
\pgfpathlineto{\pgfqpoint{3.084366in}{2.303409in}}%
\pgfpathlineto{\pgfqpoint{3.085435in}{2.303448in}}%
\pgfpathlineto{\pgfqpoint{3.084680in}{2.303339in}}%
\pgfpathlineto{\pgfqpoint{3.085482in}{2.303427in}}%
\pgfpathlineto{\pgfqpoint{3.086571in}{2.303402in}}%
\pgfpathlineto{\pgfqpoint{3.085733in}{2.303601in}}%
\pgfpathlineto{\pgfqpoint{3.086603in}{2.303473in}}%
\pgfpathlineto{\pgfqpoint{3.087497in}{2.303772in}}%
\pgfpathlineto{\pgfqpoint{3.086819in}{2.303322in}}%
\pgfpathlineto{\pgfqpoint{3.087721in}{2.303587in}}%
\pgfpathlineto{\pgfqpoint{3.087857in}{2.303553in}}%
\pgfpathlineto{\pgfqpoint{3.088723in}{2.304299in}}%
\pgfpathlineto{\pgfqpoint{3.088792in}{2.304182in}}%
\pgfpathlineto{\pgfqpoint{3.089467in}{2.303847in}}%
\pgfpathlineto{\pgfqpoint{3.089861in}{2.304469in}}%
\pgfpathlineto{\pgfqpoint{3.089885in}{2.304387in}}%
\pgfpathlineto{\pgfqpoint{3.091213in}{2.304394in}}%
\pgfpathlineto{\pgfqpoint{3.089976in}{2.304150in}}%
\pgfpathlineto{\pgfqpoint{3.091257in}{2.304189in}}%
\pgfpathlineto{\pgfqpoint{3.091889in}{2.303735in}}%
\pgfpathlineto{\pgfqpoint{3.091325in}{2.304278in}}%
\pgfpathlineto{\pgfqpoint{3.092361in}{2.304316in}}%
\pgfpathlineto{\pgfqpoint{3.092925in}{2.304544in}}%
\pgfpathlineto{\pgfqpoint{3.092497in}{2.304152in}}%
\pgfpathlineto{\pgfqpoint{3.093486in}{2.304411in}}%
\pgfpathlineto{\pgfqpoint{3.094357in}{2.304069in}}%
\pgfpathlineto{\pgfqpoint{3.093725in}{2.304529in}}%
\pgfpathlineto{\pgfqpoint{3.094608in}{2.304349in}}%
\pgfpathlineto{\pgfqpoint{3.096160in}{2.304595in}}%
\pgfpathlineto{\pgfqpoint{3.094992in}{2.303955in}}%
\pgfpathlineto{\pgfqpoint{3.096204in}{2.304435in}}%
\pgfpathlineto{\pgfqpoint{3.096467in}{2.304175in}}%
\pgfpathlineto{\pgfqpoint{3.097299in}{2.304781in}}%
\pgfpathlineto{\pgfqpoint{3.097299in}{2.304778in}}%
\pgfpathlineto{\pgfqpoint{3.098856in}{2.305217in}}%
\pgfpathlineto{\pgfqpoint{3.097344in}{2.304658in}}%
\pgfpathlineto{\pgfqpoint{3.098918in}{2.305128in}}%
\pgfpathlineto{\pgfqpoint{3.099295in}{2.304952in}}%
\pgfpathlineto{\pgfqpoint{3.099940in}{2.305311in}}%
\pgfpathlineto{\pgfqpoint{3.100033in}{2.305038in}}%
\pgfpathlineto{\pgfqpoint{3.102578in}{2.306134in}}%
\pgfpathlineto{\pgfqpoint{3.100038in}{2.305014in}}%
\pgfpathlineto{\pgfqpoint{3.102598in}{2.306019in}}%
\pgfpathlineto{\pgfqpoint{3.103717in}{2.305431in}}%
\pgfpathlineto{\pgfqpoint{3.102691in}{2.306077in}}%
\pgfpathlineto{\pgfqpoint{3.103751in}{2.305542in}}%
\pgfpathlineto{\pgfqpoint{3.104549in}{2.305881in}}%
\pgfpathlineto{\pgfqpoint{3.103851in}{2.305443in}}%
\pgfpathlineto{\pgfqpoint{3.104884in}{2.305737in}}%
\pgfpathlineto{\pgfqpoint{3.106842in}{2.304715in}}%
\pgfpathlineto{\pgfqpoint{3.104954in}{2.305836in}}%
\pgfpathlineto{\pgfqpoint{3.106944in}{2.304878in}}%
\pgfpathlineto{\pgfqpoint{3.107716in}{2.305329in}}%
\pgfpathlineto{\pgfqpoint{3.107013in}{2.304848in}}%
\pgfpathlineto{\pgfqpoint{3.108055in}{2.304884in}}%
\pgfpathlineto{\pgfqpoint{3.108966in}{2.304288in}}%
\pgfpathlineto{\pgfqpoint{3.108245in}{2.304960in}}%
\pgfpathlineto{\pgfqpoint{3.109191in}{2.304589in}}%
\pgfpathlineto{\pgfqpoint{3.109251in}{2.304814in}}%
\pgfpathlineto{\pgfqpoint{3.109994in}{2.304236in}}%
\pgfpathlineto{\pgfqpoint{3.110288in}{2.304401in}}%
\pgfpathlineto{\pgfqpoint{3.111372in}{2.303873in}}%
\pgfpathlineto{\pgfqpoint{3.110516in}{2.304506in}}%
\pgfpathlineto{\pgfqpoint{3.111528in}{2.304094in}}%
\pgfpathlineto{\pgfqpoint{3.112604in}{2.303728in}}%
\pgfpathlineto{\pgfqpoint{3.111654in}{2.303862in}}%
\pgfpathlineto{\pgfqpoint{3.112608in}{2.303697in}}%
\pgfpathlineto{\pgfqpoint{3.112623in}{2.303673in}}%
\pgfpathlineto{\pgfqpoint{3.113456in}{2.304494in}}%
\pgfpathlineto{\pgfqpoint{3.113596in}{2.304297in}}%
\pgfpathlineto{\pgfqpoint{3.113890in}{2.304512in}}%
\pgfpathlineto{\pgfqpoint{3.113646in}{2.304107in}}%
\pgfpathlineto{\pgfqpoint{3.114651in}{2.303937in}}%
\pgfpathlineto{\pgfqpoint{3.115629in}{2.303301in}}%
\pgfpathlineto{\pgfqpoint{3.115011in}{2.304109in}}%
\pgfpathlineto{\pgfqpoint{3.115786in}{2.303363in}}%
\pgfpathlineto{\pgfqpoint{3.117000in}{2.304228in}}%
\pgfpathlineto{\pgfqpoint{3.117002in}{2.304213in}}%
\pgfpathlineto{\pgfqpoint{3.117705in}{2.303822in}}%
\pgfpathlineto{\pgfqpoint{3.117283in}{2.304423in}}%
\pgfpathlineto{\pgfqpoint{3.118120in}{2.304131in}}%
\pgfpathlineto{\pgfqpoint{3.118274in}{2.304257in}}%
\pgfpathlineto{\pgfqpoint{3.119184in}{2.303846in}}%
\pgfpathlineto{\pgfqpoint{3.119219in}{2.303913in}}%
\pgfpathlineto{\pgfqpoint{3.133670in}{2.302327in}}%
\pgfpathlineto{\pgfqpoint{3.134681in}{2.301518in}}%
\pgfpathlineto{\pgfqpoint{3.134857in}{2.301584in}}%
\pgfpathlineto{\pgfqpoint{3.136225in}{2.301555in}}%
\pgfpathlineto{\pgfqpoint{3.135003in}{2.301715in}}%
\pgfpathlineto{\pgfqpoint{3.136236in}{2.301593in}}%
\pgfpathlineto{\pgfqpoint{3.136523in}{2.302023in}}%
\pgfpathlineto{\pgfqpoint{3.137113in}{2.301292in}}%
\pgfpathlineto{\pgfqpoint{3.137342in}{2.301457in}}%
\pgfpathlineto{\pgfqpoint{3.138243in}{2.301719in}}%
\pgfpathlineto{\pgfqpoint{3.137441in}{2.301339in}}%
\pgfpathlineto{\pgfqpoint{3.138434in}{2.301266in}}%
\pgfpathlineto{\pgfqpoint{3.139007in}{2.301135in}}%
\pgfpathlineto{\pgfqpoint{3.138855in}{2.301608in}}%
\pgfpathlineto{\pgfqpoint{3.139539in}{2.301442in}}%
\pgfpathlineto{\pgfqpoint{3.140690in}{2.302164in}}%
\pgfpathlineto{\pgfqpoint{3.139656in}{2.301398in}}%
\pgfpathlineto{\pgfqpoint{3.140708in}{2.302130in}}%
\pgfpathlineto{\pgfqpoint{3.141976in}{2.303449in}}%
\pgfpathlineto{\pgfqpoint{3.140766in}{2.302077in}}%
\pgfpathlineto{\pgfqpoint{3.142016in}{2.303412in}}%
\pgfpathlineto{\pgfqpoint{3.143094in}{2.303378in}}%
\pgfpathlineto{\pgfqpoint{3.142046in}{2.303441in}}%
\pgfpathlineto{\pgfqpoint{3.143132in}{2.303407in}}%
\pgfpathlineto{\pgfqpoint{3.144240in}{2.303457in}}%
\pgfpathlineto{\pgfqpoint{3.143278in}{2.303335in}}%
\pgfpathlineto{\pgfqpoint{3.144244in}{2.303406in}}%
\pgfpathlineto{\pgfqpoint{3.145377in}{2.303440in}}%
\pgfpathlineto{\pgfqpoint{3.144253in}{2.303387in}}%
\pgfpathlineto{\pgfqpoint{3.145403in}{2.303430in}}%
\pgfpathlineto{\pgfqpoint{3.146263in}{2.303363in}}%
\pgfpathlineto{\pgfqpoint{3.145767in}{2.303460in}}%
\pgfpathlineto{\pgfqpoint{3.146515in}{2.303406in}}%
\pgfpathlineto{\pgfqpoint{3.151129in}{2.303444in}}%
\pgfpathlineto{\pgfqpoint{3.151132in}{2.303388in}}%
\pgfpathlineto{\pgfqpoint{3.151929in}{2.303338in}}%
\pgfpathlineto{\pgfqpoint{3.151189in}{2.303444in}}%
\pgfpathlineto{\pgfqpoint{3.152242in}{2.303405in}}%
\pgfpathlineto{\pgfqpoint{3.153293in}{2.303346in}}%
\pgfpathlineto{\pgfqpoint{3.152257in}{2.303437in}}%
\pgfpathlineto{\pgfqpoint{3.153361in}{2.303418in}}%
\pgfpathlineto{\pgfqpoint{3.154492in}{2.303434in}}%
\pgfpathlineto{\pgfqpoint{3.153376in}{2.303388in}}%
\pgfpathlineto{\pgfqpoint{3.154513in}{2.303398in}}%
\pgfpathlineto{\pgfqpoint{3.155636in}{2.303386in}}%
\pgfpathlineto{\pgfqpoint{3.154525in}{2.303435in}}%
\pgfpathlineto{\pgfqpoint{3.155638in}{2.303403in}}%
\pgfpathlineto{\pgfqpoint{3.156717in}{2.303442in}}%
\pgfpathlineto{\pgfqpoint{3.155648in}{2.303360in}}%
\pgfpathlineto{\pgfqpoint{3.156750in}{2.303410in}}%
\pgfpathlineto{\pgfqpoint{3.157892in}{2.303452in}}%
\pgfpathlineto{\pgfqpoint{3.156775in}{2.303376in}}%
\pgfpathlineto{\pgfqpoint{3.157911in}{2.303415in}}%
\pgfpathlineto{\pgfqpoint{3.159015in}{2.303364in}}%
\pgfpathlineto{\pgfqpoint{3.157920in}{2.303440in}}%
\pgfpathlineto{\pgfqpoint{3.159042in}{2.303419in}}%
\pgfpathlineto{\pgfqpoint{3.160147in}{2.303446in}}%
\pgfpathlineto{\pgfqpoint{3.159138in}{2.303362in}}%
\pgfpathlineto{\pgfqpoint{3.160164in}{2.303425in}}%
\pgfpathlineto{\pgfqpoint{3.160342in}{2.303362in}}%
\pgfpathlineto{\pgfqpoint{3.161242in}{2.303879in}}%
\pgfpathlineto{\pgfqpoint{3.161263in}{2.303810in}}%
\pgfpathlineto{\pgfqpoint{3.161737in}{2.303503in}}%
\pgfpathlineto{\pgfqpoint{3.161322in}{2.303866in}}%
\pgfpathlineto{\pgfqpoint{3.162364in}{2.303955in}}%
\pgfpathlineto{\pgfqpoint{3.163098in}{2.304111in}}%
\pgfpathlineto{\pgfqpoint{3.162934in}{2.303729in}}%
\pgfpathlineto{\pgfqpoint{3.163473in}{2.303888in}}%
\pgfpathlineto{\pgfqpoint{3.163483in}{2.303822in}}%
\pgfpathlineto{\pgfqpoint{3.163971in}{2.304412in}}%
\pgfpathlineto{\pgfqpoint{3.164543in}{2.304265in}}%
\pgfpathlineto{\pgfqpoint{3.165524in}{2.304786in}}%
\pgfpathlineto{\pgfqpoint{3.165664in}{2.304582in}}%
\pgfpathlineto{\pgfqpoint{3.165696in}{2.304550in}}%
\pgfpathlineto{\pgfqpoint{3.166304in}{2.305181in}}%
\pgfpathlineto{\pgfqpoint{3.166707in}{2.305106in}}%
\pgfpathlineto{\pgfqpoint{3.167083in}{2.305234in}}%
\pgfpathlineto{\pgfqpoint{3.167640in}{2.304596in}}%
\pgfpathlineto{\pgfqpoint{3.167766in}{2.304737in}}%
\pgfpathlineto{\pgfqpoint{3.168547in}{2.304658in}}%
\pgfpathlineto{\pgfqpoint{3.168828in}{2.305326in}}%
\pgfpathlineto{\pgfqpoint{3.168845in}{2.305299in}}%
\pgfpathlineto{\pgfqpoint{3.168860in}{2.305380in}}%
\pgfpathlineto{\pgfqpoint{3.169881in}{2.304911in}}%
\pgfpathlineto{\pgfqpoint{3.169943in}{2.305051in}}%
\pgfpathlineto{\pgfqpoint{3.170724in}{2.305545in}}%
\pgfpathlineto{\pgfqpoint{3.169953in}{2.305015in}}%
\pgfpathlineto{\pgfqpoint{3.171070in}{2.305090in}}%
\pgfpathlineto{\pgfqpoint{3.171565in}{2.305081in}}%
\pgfpathlineto{\pgfqpoint{3.172132in}{2.305551in}}%
\pgfpathlineto{\pgfqpoint{3.172137in}{2.305518in}}%
\pgfpathlineto{\pgfqpoint{3.172934in}{2.306190in}}%
\pgfpathlineto{\pgfqpoint{3.172169in}{2.305455in}}%
\pgfpathlineto{\pgfqpoint{3.173258in}{2.305863in}}%
\pgfpathlineto{\pgfqpoint{3.173513in}{2.305931in}}%
\pgfpathlineto{\pgfqpoint{3.173796in}{2.305301in}}%
\pgfpathlineto{\pgfqpoint{3.174333in}{2.305568in}}%
\pgfpathlineto{\pgfqpoint{3.175469in}{2.304990in}}%
\pgfpathlineto{\pgfqpoint{3.175470in}{2.305012in}}%
\pgfpathlineto{\pgfqpoint{3.188233in}{2.306918in}}%
\pgfpathlineto{\pgfqpoint{3.175472in}{2.305001in}}%
\pgfpathlineto{\pgfqpoint{3.188236in}{2.306881in}}%
\pgfpathlineto{\pgfqpoint{3.189427in}{2.306561in}}%
\pgfpathlineto{\pgfqpoint{3.188349in}{2.307059in}}%
\pgfpathlineto{\pgfqpoint{3.189440in}{2.306603in}}%
\pgfpathlineto{\pgfqpoint{3.190519in}{2.307161in}}%
\pgfpathlineto{\pgfqpoint{3.189451in}{2.306547in}}%
\pgfpathlineto{\pgfqpoint{3.190787in}{2.306781in}}%
\pgfpathlineto{\pgfqpoint{3.191812in}{2.306042in}}%
\pgfpathlineto{\pgfqpoint{3.190876in}{2.306894in}}%
\pgfpathlineto{\pgfqpoint{3.191938in}{2.306108in}}%
\pgfpathlineto{\pgfqpoint{3.192127in}{2.306006in}}%
\pgfpathlineto{\pgfqpoint{3.192537in}{2.306715in}}%
\pgfpathlineto{\pgfqpoint{3.193017in}{2.306451in}}%
\pgfpathlineto{\pgfqpoint{3.193402in}{2.306171in}}%
\pgfpathlineto{\pgfqpoint{3.194149in}{2.306782in}}%
\pgfpathlineto{\pgfqpoint{3.194954in}{2.306569in}}%
\pgfpathlineto{\pgfqpoint{3.194325in}{2.306895in}}%
\pgfpathlineto{\pgfqpoint{3.195259in}{2.306818in}}%
\pgfpathlineto{\pgfqpoint{3.196391in}{2.307263in}}%
\pgfpathlineto{\pgfqpoint{3.195874in}{2.306592in}}%
\pgfpathlineto{\pgfqpoint{3.196393in}{2.307251in}}%
\pgfpathlineto{\pgfqpoint{3.197281in}{2.307405in}}%
\pgfpathlineto{\pgfqpoint{3.196606in}{2.307150in}}%
\pgfpathlineto{\pgfqpoint{3.197478in}{2.307015in}}%
\pgfpathlineto{\pgfqpoint{3.198948in}{2.306180in}}%
\pgfpathlineto{\pgfqpoint{3.197513in}{2.307077in}}%
\pgfpathlineto{\pgfqpoint{3.198958in}{2.306202in}}%
\pgfpathlineto{\pgfqpoint{3.199276in}{2.306615in}}%
\pgfpathlineto{\pgfqpoint{3.198975in}{2.306099in}}%
\pgfpathlineto{\pgfqpoint{3.200080in}{2.306350in}}%
\pgfpathlineto{\pgfqpoint{3.201023in}{2.305998in}}%
\pgfpathlineto{\pgfqpoint{3.200093in}{2.306353in}}%
\pgfpathlineto{\pgfqpoint{3.201208in}{2.306093in}}%
\pgfpathlineto{\pgfqpoint{3.202219in}{2.306221in}}%
\pgfpathlineto{\pgfqpoint{3.201480in}{2.305901in}}%
\pgfpathlineto{\pgfqpoint{3.202451in}{2.306030in}}%
\pgfpathlineto{\pgfqpoint{3.203295in}{2.305232in}}%
\pgfpathlineto{\pgfqpoint{3.202461in}{2.306055in}}%
\pgfpathlineto{\pgfqpoint{3.203586in}{2.305354in}}%
\pgfpathlineto{\pgfqpoint{3.203786in}{2.305460in}}%
\pgfpathlineto{\pgfqpoint{3.204675in}{2.304971in}}%
\pgfpathlineto{\pgfqpoint{3.205335in}{2.304465in}}%
\pgfpathlineto{\pgfqpoint{3.204697in}{2.305005in}}%
\pgfpathlineto{\pgfqpoint{3.205808in}{2.304801in}}%
\pgfpathlineto{\pgfqpoint{3.205952in}{2.305016in}}%
\pgfpathlineto{\pgfqpoint{3.206907in}{2.304440in}}%
\pgfpathlineto{\pgfqpoint{3.207255in}{2.304079in}}%
\pgfpathlineto{\pgfqpoint{3.206928in}{2.304481in}}%
\pgfpathlineto{\pgfqpoint{3.208050in}{2.304276in}}%
\pgfpathlineto{\pgfqpoint{3.208559in}{2.304403in}}%
\pgfpathlineto{\pgfqpoint{3.208892in}{2.304018in}}%
\pgfpathlineto{\pgfqpoint{3.209148in}{2.304030in}}%
\pgfpathlineto{\pgfqpoint{3.210279in}{2.303642in}}%
\pgfpathlineto{\pgfqpoint{3.209216in}{2.304106in}}%
\pgfpathlineto{\pgfqpoint{3.210282in}{2.303667in}}%
\pgfpathlineto{\pgfqpoint{3.210831in}{2.303954in}}%
\pgfpathlineto{\pgfqpoint{3.210580in}{2.303486in}}%
\pgfpathlineto{\pgfqpoint{3.211393in}{2.303638in}}%
\pgfpathlineto{\pgfqpoint{3.211774in}{2.303485in}}%
\pgfpathlineto{\pgfqpoint{3.211968in}{2.303841in}}%
\pgfpathlineto{\pgfqpoint{3.212512in}{2.303584in}}%
\pgfpathlineto{\pgfqpoint{3.213394in}{2.303852in}}%
\pgfpathlineto{\pgfqpoint{3.212854in}{2.303339in}}%
\pgfpathlineto{\pgfqpoint{3.213642in}{2.303705in}}%
\pgfpathlineto{\pgfqpoint{3.213771in}{2.303435in}}%
\pgfpathlineto{\pgfqpoint{3.214330in}{2.304007in}}%
\pgfpathlineto{\pgfqpoint{3.214740in}{2.303890in}}%
\pgfpathlineto{\pgfqpoint{3.215365in}{2.304153in}}%
\pgfpathlineto{\pgfqpoint{3.214744in}{2.303880in}}%
\pgfpathlineto{\pgfqpoint{3.215872in}{2.304077in}}%
\pgfpathlineto{\pgfqpoint{3.217011in}{2.303382in}}%
\pgfpathlineto{\pgfqpoint{3.216064in}{2.304273in}}%
\pgfpathlineto{\pgfqpoint{3.217068in}{2.303402in}}%
\pgfpathlineto{\pgfqpoint{3.218169in}{2.303443in}}%
\pgfpathlineto{\pgfqpoint{3.217133in}{2.303366in}}%
\pgfpathlineto{\pgfqpoint{3.218181in}{2.303416in}}%
\pgfpathlineto{\pgfqpoint{3.229762in}{2.303419in}}%
\pgfpathlineto{\pgfqpoint{3.230536in}{2.303454in}}%
\pgfpathlineto{\pgfqpoint{3.230038in}{2.303366in}}%
\pgfpathlineto{\pgfqpoint{3.230872in}{2.303393in}}%
\pgfpathlineto{\pgfqpoint{3.231992in}{2.303444in}}%
\pgfpathlineto{\pgfqpoint{3.230975in}{2.303323in}}%
\pgfpathlineto{\pgfqpoint{3.231994in}{2.303402in}}%
\pgfpathlineto{\pgfqpoint{3.233004in}{2.303364in}}%
\pgfpathlineto{\pgfqpoint{3.232009in}{2.303448in}}%
\pgfpathlineto{\pgfqpoint{3.233106in}{2.303374in}}%
\pgfpathlineto{\pgfqpoint{3.234131in}{2.303444in}}%
\pgfpathlineto{\pgfqpoint{3.233134in}{2.303358in}}%
\pgfpathlineto{\pgfqpoint{3.234221in}{2.303433in}}%
\pgfpathlineto{\pgfqpoint{3.235120in}{2.303354in}}%
\pgfpathlineto{\pgfqpoint{3.235329in}{2.303527in}}%
\pgfpathlineto{\pgfqpoint{3.235330in}{2.303509in}}%
\pgfpathlineto{\pgfqpoint{3.236450in}{2.304200in}}%
\pgfpathlineto{\pgfqpoint{3.235803in}{2.303336in}}%
\pgfpathlineto{\pgfqpoint{3.236474in}{2.304149in}}%
\pgfpathlineto{\pgfqpoint{3.237101in}{2.303949in}}%
\pgfpathlineto{\pgfqpoint{3.237458in}{2.304398in}}%
\pgfpathlineto{\pgfqpoint{3.237580in}{2.304263in}}%
\pgfpathlineto{\pgfqpoint{3.238736in}{2.304537in}}%
\pgfpathlineto{\pgfqpoint{3.237706in}{2.304031in}}%
\pgfpathlineto{\pgfqpoint{3.238795in}{2.304468in}}%
\pgfpathlineto{\pgfqpoint{3.239722in}{2.303733in}}%
\pgfpathlineto{\pgfqpoint{3.238982in}{2.304541in}}%
\pgfpathlineto{\pgfqpoint{3.239924in}{2.303892in}}%
\pgfpathlineto{\pgfqpoint{3.241126in}{2.304311in}}%
\pgfpathlineto{\pgfqpoint{3.240108in}{2.303774in}}%
\pgfpathlineto{\pgfqpoint{3.241129in}{2.304296in}}%
\pgfpathlineto{\pgfqpoint{3.241517in}{2.304669in}}%
\pgfpathlineto{\pgfqpoint{3.241145in}{2.304248in}}%
\pgfpathlineto{\pgfqpoint{3.242185in}{2.303871in}}%
\pgfpathlineto{\pgfqpoint{3.243185in}{2.303708in}}%
\pgfpathlineto{\pgfqpoint{3.242618in}{2.304093in}}%
\pgfpathlineto{\pgfqpoint{3.243293in}{2.303949in}}%
\pgfpathlineto{\pgfqpoint{3.243392in}{2.304105in}}%
\pgfpathlineto{\pgfqpoint{3.244211in}{2.303249in}}%
\pgfpathlineto{\pgfqpoint{3.244381in}{2.303321in}}%
\pgfpathlineto{\pgfqpoint{3.244410in}{2.303236in}}%
\pgfpathlineto{\pgfqpoint{3.245436in}{2.303956in}}%
\pgfpathlineto{\pgfqpoint{3.245451in}{2.303900in}}%
\pgfpathlineto{\pgfqpoint{3.245585in}{2.304049in}}%
\pgfpathlineto{\pgfqpoint{3.246244in}{2.302904in}}%
\pgfpathlineto{\pgfqpoint{3.246524in}{2.303367in}}%
\pgfpathlineto{\pgfqpoint{3.246535in}{2.303334in}}%
\pgfpathlineto{\pgfqpoint{3.247497in}{2.304039in}}%
\pgfpathlineto{\pgfqpoint{3.247572in}{2.303963in}}%
\pgfpathlineto{\pgfqpoint{3.247698in}{2.304035in}}%
\pgfpathlineto{\pgfqpoint{3.248561in}{2.303451in}}%
\pgfpathlineto{\pgfqpoint{3.248660in}{2.303660in}}%
\pgfpathlineto{\pgfqpoint{3.249839in}{2.303105in}}%
\pgfpathlineto{\pgfqpoint{3.248776in}{2.303817in}}%
\pgfpathlineto{\pgfqpoint{3.249861in}{2.303168in}}%
\pgfpathlineto{\pgfqpoint{3.250345in}{2.303465in}}%
\pgfpathlineto{\pgfqpoint{3.250788in}{2.302972in}}%
\pgfpathlineto{\pgfqpoint{3.250966in}{2.303056in}}%
\pgfpathlineto{\pgfqpoint{3.250997in}{2.302996in}}%
\pgfpathlineto{\pgfqpoint{3.251903in}{2.303527in}}%
\pgfpathlineto{\pgfqpoint{3.252061in}{2.303332in}}%
\pgfpathlineto{\pgfqpoint{3.253012in}{2.303492in}}%
\pgfpathlineto{\pgfqpoint{3.252724in}{2.303012in}}%
\pgfpathlineto{\pgfqpoint{3.253184in}{2.303338in}}%
\pgfpathlineto{\pgfqpoint{3.253668in}{2.302781in}}%
\pgfpathlineto{\pgfqpoint{3.253265in}{2.303382in}}%
\pgfpathlineto{\pgfqpoint{3.254306in}{2.302997in}}%
\pgfpathlineto{\pgfqpoint{3.255521in}{2.302707in}}%
\pgfpathlineto{\pgfqpoint{3.254316in}{2.303054in}}%
\pgfpathlineto{\pgfqpoint{3.255532in}{2.302775in}}%
\pgfpathlineto{\pgfqpoint{3.256375in}{2.303051in}}%
\pgfpathlineto{\pgfqpoint{3.256598in}{2.302587in}}%
\pgfpathlineto{\pgfqpoint{3.256642in}{2.302725in}}%
\pgfpathlineto{\pgfqpoint{3.257250in}{2.301996in}}%
\pgfpathlineto{\pgfqpoint{3.256783in}{2.302818in}}%
\pgfpathlineto{\pgfqpoint{3.257768in}{2.302481in}}%
\pgfpathlineto{\pgfqpoint{3.258930in}{2.303524in}}%
\pgfpathlineto{\pgfqpoint{3.257776in}{2.302456in}}%
\pgfpathlineto{\pgfqpoint{3.258967in}{2.303465in}}%
\pgfpathlineto{\pgfqpoint{3.259237in}{2.303296in}}%
\pgfpathlineto{\pgfqpoint{3.259703in}{2.303611in}}%
\pgfpathlineto{\pgfqpoint{3.260084in}{2.303431in}}%
\pgfpathlineto{\pgfqpoint{3.261054in}{2.303501in}}%
\pgfpathlineto{\pgfqpoint{3.260138in}{2.303308in}}%
\pgfpathlineto{\pgfqpoint{3.261187in}{2.303313in}}%
\pgfpathlineto{\pgfqpoint{3.261927in}{2.302649in}}%
\pgfpathlineto{\pgfqpoint{3.262326in}{2.302844in}}%
\pgfpathlineto{\pgfqpoint{3.262691in}{2.303033in}}%
\pgfpathlineto{\pgfqpoint{3.263333in}{2.302451in}}%
\pgfpathlineto{\pgfqpoint{3.263399in}{2.302474in}}%
\pgfpathlineto{\pgfqpoint{3.264972in}{2.301910in}}%
\pgfpathlineto{\pgfqpoint{3.263581in}{2.302700in}}%
\pgfpathlineto{\pgfqpoint{3.265012in}{2.302064in}}%
\pgfpathlineto{\pgfqpoint{3.265049in}{2.302167in}}%
\pgfpathlineto{\pgfqpoint{3.265679in}{2.301684in}}%
\pgfpathlineto{\pgfqpoint{3.266105in}{2.301769in}}%
\pgfpathlineto{\pgfqpoint{3.267242in}{2.301308in}}%
\pgfpathlineto{\pgfqpoint{3.266677in}{2.301907in}}%
\pgfpathlineto{\pgfqpoint{3.267243in}{2.301310in}}%
\pgfpathlineto{\pgfqpoint{3.268379in}{2.301289in}}%
\pgfpathlineto{\pgfqpoint{3.267362in}{2.301595in}}%
\pgfpathlineto{\pgfqpoint{3.268470in}{2.301531in}}%
\pgfpathlineto{\pgfqpoint{3.269588in}{2.301615in}}%
\pgfpathlineto{\pgfqpoint{3.269209in}{2.301145in}}%
\pgfpathlineto{\pgfqpoint{3.269588in}{2.301610in}}%
\pgfpathlineto{\pgfqpoint{3.272168in}{2.302249in}}%
\pgfpathlineto{\pgfqpoint{3.269607in}{2.301488in}}%
\pgfpathlineto{\pgfqpoint{3.272241in}{2.301970in}}%
\pgfpathlineto{\pgfqpoint{3.272342in}{2.301737in}}%
\pgfpathlineto{\pgfqpoint{3.273192in}{2.302549in}}%
\pgfpathlineto{\pgfqpoint{3.273321in}{2.302414in}}%
\pgfpathlineto{\pgfqpoint{3.273407in}{2.302505in}}%
\pgfpathlineto{\pgfqpoint{3.273816in}{2.301502in}}%
\pgfpathlineto{\pgfqpoint{3.274213in}{2.301754in}}%
\pgfpathlineto{\pgfqpoint{3.275195in}{2.301525in}}%
\pgfpathlineto{\pgfqpoint{3.274824in}{2.302030in}}%
\pgfpathlineto{\pgfqpoint{3.275324in}{2.301774in}}%
\pgfpathlineto{\pgfqpoint{3.275568in}{2.302225in}}%
\pgfpathlineto{\pgfqpoint{3.276407in}{2.301482in}}%
\pgfpathlineto{\pgfqpoint{3.276408in}{2.301485in}}%
\pgfpathlineto{\pgfqpoint{3.276633in}{2.301341in}}%
\pgfpathlineto{\pgfqpoint{3.276876in}{2.301727in}}%
\pgfpathlineto{\pgfqpoint{3.277521in}{2.301418in}}%
\pgfpathlineto{\pgfqpoint{3.278440in}{2.301956in}}%
\pgfpathlineto{\pgfqpoint{3.277536in}{2.301338in}}%
\pgfpathlineto{\pgfqpoint{3.278673in}{2.301741in}}%
\pgfpathlineto{\pgfqpoint{3.286075in}{2.303376in}}%
\pgfpathlineto{\pgfqpoint{3.286575in}{2.303767in}}%
\pgfpathlineto{\pgfqpoint{3.286153in}{2.303315in}}%
\pgfpathlineto{\pgfqpoint{3.287193in}{2.303519in}}%
\pgfpathlineto{\pgfqpoint{3.288074in}{2.303302in}}%
\pgfpathlineto{\pgfqpoint{3.287330in}{2.303630in}}%
\pgfpathlineto{\pgfqpoint{3.288277in}{2.303760in}}%
\pgfpathlineto{\pgfqpoint{3.289239in}{2.304091in}}%
\pgfpathlineto{\pgfqpoint{3.288650in}{2.303444in}}%
\pgfpathlineto{\pgfqpoint{3.289393in}{2.303898in}}%
\pgfpathlineto{\pgfqpoint{3.289500in}{2.304025in}}%
\pgfpathlineto{\pgfqpoint{3.290309in}{2.303503in}}%
\pgfpathlineto{\pgfqpoint{3.290475in}{2.303536in}}%
\pgfpathlineto{\pgfqpoint{3.291326in}{2.303218in}}%
\pgfpathlineto{\pgfqpoint{3.290923in}{2.303726in}}%
\pgfpathlineto{\pgfqpoint{3.291591in}{2.303413in}}%
\pgfpathlineto{\pgfqpoint{3.292678in}{2.303380in}}%
\pgfpathlineto{\pgfqpoint{3.291607in}{2.303444in}}%
\pgfpathlineto{\pgfqpoint{3.292705in}{2.303419in}}%
\pgfpathlineto{\pgfqpoint{3.293664in}{2.303451in}}%
\pgfpathlineto{\pgfqpoint{3.293020in}{2.303360in}}%
\pgfpathlineto{\pgfqpoint{3.293816in}{2.303404in}}%
\pgfpathlineto{\pgfqpoint{3.295126in}{2.303433in}}%
\pgfpathlineto{\pgfqpoint{3.293841in}{2.303366in}}%
\pgfpathlineto{\pgfqpoint{3.295128in}{2.303409in}}%
\pgfpathlineto{\pgfqpoint{3.296170in}{2.303376in}}%
\pgfpathlineto{\pgfqpoint{3.295198in}{2.303441in}}%
\pgfpathlineto{\pgfqpoint{3.296239in}{2.303421in}}%
\pgfpathlineto{\pgfqpoint{3.297393in}{2.303443in}}%
\pgfpathlineto{\pgfqpoint{3.296293in}{2.303370in}}%
\pgfpathlineto{\pgfqpoint{3.297396in}{2.303434in}}%
\pgfpathlineto{\pgfqpoint{3.298356in}{2.303362in}}%
\pgfpathlineto{\pgfqpoint{3.297766in}{2.303452in}}%
\pgfpathlineto{\pgfqpoint{3.298508in}{2.303400in}}%
\pgfpathlineto{\pgfqpoint{3.299619in}{2.303373in}}%
\pgfpathlineto{\pgfqpoint{3.298516in}{2.303440in}}%
\pgfpathlineto{\pgfqpoint{3.299634in}{2.303408in}}%
\pgfpathlineto{\pgfqpoint{3.300574in}{2.303446in}}%
\pgfpathlineto{\pgfqpoint{3.300192in}{2.303359in}}%
\pgfpathlineto{\pgfqpoint{3.300745in}{2.303419in}}%
\pgfpathlineto{\pgfqpoint{3.301807in}{2.303378in}}%
\pgfpathlineto{\pgfqpoint{3.300804in}{2.303451in}}%
\pgfpathlineto{\pgfqpoint{3.301860in}{2.303397in}}%
\pgfpathlineto{\pgfqpoint{3.302956in}{2.303449in}}%
\pgfpathlineto{\pgfqpoint{3.302289in}{2.303349in}}%
\pgfpathlineto{\pgfqpoint{3.302973in}{2.303422in}}%
\pgfpathlineto{\pgfqpoint{3.304078in}{2.303444in}}%
\pgfpathlineto{\pgfqpoint{3.302987in}{2.303371in}}%
\pgfpathlineto{\pgfqpoint{3.304108in}{2.303405in}}%
\pgfpathlineto{\pgfqpoint{3.305206in}{2.303380in}}%
\pgfpathlineto{\pgfqpoint{3.304237in}{2.303445in}}%
\pgfpathlineto{\pgfqpoint{3.305220in}{2.303422in}}%
\pgfpathlineto{\pgfqpoint{3.306322in}{2.303440in}}%
\pgfpathlineto{\pgfqpoint{3.305379in}{2.303359in}}%
\pgfpathlineto{\pgfqpoint{3.306330in}{2.303374in}}%
\pgfpathlineto{\pgfqpoint{3.306450in}{2.303331in}}%
\pgfpathlineto{\pgfqpoint{3.306393in}{2.303448in}}%
\pgfpathlineto{\pgfqpoint{3.307440in}{2.303414in}}%
\pgfpathlineto{\pgfqpoint{3.310174in}{2.303430in}}%
\pgfpathlineto{\pgfqpoint{3.307443in}{2.303388in}}%
\pgfpathlineto{\pgfqpoint{3.310186in}{2.303380in}}%
\pgfpathlineto{\pgfqpoint{3.310226in}{2.303289in}}%
\pgfpathlineto{\pgfqpoint{3.311277in}{2.303990in}}%
\pgfpathlineto{\pgfqpoint{3.311947in}{2.304038in}}%
\pgfpathlineto{\pgfqpoint{3.312257in}{2.303444in}}%
\pgfpathlineto{\pgfqpoint{3.312257in}{2.303445in}}%
\pgfpathlineto{\pgfqpoint{3.312340in}{2.303296in}}%
\pgfpathlineto{\pgfqpoint{3.313339in}{2.304078in}}%
\pgfpathlineto{\pgfqpoint{3.313394in}{2.304148in}}%
\pgfpathlineto{\pgfqpoint{3.314051in}{2.303681in}}%
\pgfpathlineto{\pgfqpoint{3.314412in}{2.303718in}}%
\pgfpathlineto{\pgfqpoint{3.315359in}{2.302968in}}%
\pgfpathlineto{\pgfqpoint{3.314422in}{2.303806in}}%
\pgfpathlineto{\pgfqpoint{3.315574in}{2.303036in}}%
\pgfpathlineto{\pgfqpoint{3.315894in}{2.303288in}}%
\pgfpathlineto{\pgfqpoint{3.315721in}{2.302895in}}%
\pgfpathlineto{\pgfqpoint{3.316686in}{2.303065in}}%
\pgfpathlineto{\pgfqpoint{3.317437in}{2.302665in}}%
\pgfpathlineto{\pgfqpoint{3.316904in}{2.303229in}}%
\pgfpathlineto{\pgfqpoint{3.317811in}{2.302985in}}%
\pgfpathlineto{\pgfqpoint{3.318707in}{2.303520in}}%
\pgfpathlineto{\pgfqpoint{3.318000in}{2.302868in}}%
\pgfpathlineto{\pgfqpoint{3.318934in}{2.303284in}}%
\pgfpathlineto{\pgfqpoint{3.318959in}{2.303241in}}%
\pgfpathlineto{\pgfqpoint{3.319413in}{2.303692in}}%
\pgfpathlineto{\pgfqpoint{3.319986in}{2.303714in}}%
\pgfpathlineto{\pgfqpoint{3.320957in}{2.304180in}}%
\pgfpathlineto{\pgfqpoint{3.320519in}{2.303542in}}%
\pgfpathlineto{\pgfqpoint{3.321133in}{2.303979in}}%
\pgfpathlineto{\pgfqpoint{3.321283in}{2.303844in}}%
\pgfpathlineto{\pgfqpoint{3.322152in}{2.304530in}}%
\pgfpathlineto{\pgfqpoint{3.322190in}{2.304383in}}%
\pgfpathlineto{\pgfqpoint{3.322479in}{2.304810in}}%
\pgfpathlineto{\pgfqpoint{3.323310in}{2.304601in}}%
\pgfpathlineto{\pgfqpoint{3.323765in}{2.304279in}}%
\pgfpathlineto{\pgfqpoint{3.324263in}{2.304883in}}%
\pgfpathlineto{\pgfqpoint{3.324401in}{2.304845in}}%
\pgfpathlineto{\pgfqpoint{3.325432in}{2.304995in}}%
\pgfpathlineto{\pgfqpoint{3.324710in}{2.304520in}}%
\pgfpathlineto{\pgfqpoint{3.325517in}{2.304910in}}%
\pgfpathlineto{\pgfqpoint{3.327050in}{2.304513in}}%
\pgfpathlineto{\pgfqpoint{3.325525in}{2.304874in}}%
\pgfpathlineto{\pgfqpoint{3.327064in}{2.304439in}}%
\pgfpathlineto{\pgfqpoint{3.328051in}{2.304128in}}%
\pgfpathlineto{\pgfqpoint{3.327506in}{2.304717in}}%
\pgfpathlineto{\pgfqpoint{3.328180in}{2.304288in}}%
\pgfpathlineto{\pgfqpoint{3.329014in}{2.304302in}}%
\pgfpathlineto{\pgfqpoint{3.328362in}{2.304490in}}%
\pgfpathlineto{\pgfqpoint{3.329252in}{2.304581in}}%
\pgfpathlineto{\pgfqpoint{3.329271in}{2.304611in}}%
\pgfpathlineto{\pgfqpoint{3.329938in}{2.303938in}}%
\pgfpathlineto{\pgfqpoint{3.330332in}{2.304224in}}%
\pgfpathlineto{\pgfqpoint{3.330487in}{2.304134in}}%
\pgfpathlineto{\pgfqpoint{3.331395in}{2.304906in}}%
\pgfpathlineto{\pgfqpoint{3.331400in}{2.304858in}}%
\pgfpathlineto{\pgfqpoint{3.332481in}{2.305518in}}%
\pgfpathlineto{\pgfqpoint{3.332528in}{2.305461in}}%
\pgfpathlineto{\pgfqpoint{3.334951in}{2.305837in}}%
\pgfpathlineto{\pgfqpoint{3.332534in}{2.305505in}}%
\pgfpathlineto{\pgfqpoint{3.334975in}{2.305968in}}%
\pgfpathlineto{\pgfqpoint{3.335649in}{2.306498in}}%
\pgfpathlineto{\pgfqpoint{3.335089in}{2.305895in}}%
\pgfpathlineto{\pgfqpoint{3.336091in}{2.306029in}}%
\pgfpathlineto{\pgfqpoint{3.336452in}{2.305822in}}%
\pgfpathlineto{\pgfqpoint{3.336993in}{2.306669in}}%
\pgfpathlineto{\pgfqpoint{3.337017in}{2.306760in}}%
\pgfpathlineto{\pgfqpoint{3.337947in}{2.306231in}}%
\pgfpathlineto{\pgfqpoint{3.338086in}{2.306326in}}%
\pgfpathlineto{\pgfqpoint{3.339227in}{2.306990in}}%
\pgfpathlineto{\pgfqpoint{3.339269in}{2.306944in}}%
\pgfpathlineto{\pgfqpoint{3.340081in}{2.306645in}}%
\pgfpathlineto{\pgfqpoint{3.339737in}{2.307209in}}%
\pgfpathlineto{\pgfqpoint{3.340378in}{2.306998in}}%
\pgfpathlineto{\pgfqpoint{3.340755in}{2.307223in}}%
\pgfpathlineto{\pgfqpoint{3.341438in}{2.306423in}}%
\pgfpathlineto{\pgfqpoint{3.341439in}{2.306429in}}%
\pgfpathlineto{\pgfqpoint{3.342447in}{2.305447in}}%
\pgfpathlineto{\pgfqpoint{3.341459in}{2.306434in}}%
\pgfpathlineto{\pgfqpoint{3.342728in}{2.305562in}}%
\pgfpathlineto{\pgfqpoint{3.343061in}{2.305672in}}%
\pgfpathlineto{\pgfqpoint{3.343694in}{2.305037in}}%
\pgfpathlineto{\pgfqpoint{3.343807in}{2.305094in}}%
\pgfpathlineto{\pgfqpoint{3.344979in}{2.305489in}}%
\pgfpathlineto{\pgfqpoint{3.344109in}{2.304935in}}%
\pgfpathlineto{\pgfqpoint{3.344992in}{2.305464in}}%
\pgfpathlineto{\pgfqpoint{3.346950in}{2.304966in}}%
\pgfpathlineto{\pgfqpoint{3.345025in}{2.305540in}}%
\pgfpathlineto{\pgfqpoint{3.346993in}{2.305109in}}%
\pgfpathlineto{\pgfqpoint{3.347517in}{2.305248in}}%
\pgfpathlineto{\pgfqpoint{3.348018in}{2.304766in}}%
\pgfpathlineto{\pgfqpoint{3.348093in}{2.304945in}}%
\pgfpathlineto{\pgfqpoint{3.349092in}{2.304910in}}%
\pgfpathlineto{\pgfqpoint{3.348551in}{2.305335in}}%
\pgfpathlineto{\pgfqpoint{3.349195in}{2.305086in}}%
\pgfpathlineto{\pgfqpoint{3.349450in}{2.305273in}}%
\pgfpathlineto{\pgfqpoint{3.350147in}{2.304093in}}%
\pgfpathlineto{\pgfqpoint{3.350255in}{2.304311in}}%
\pgfpathlineto{\pgfqpoint{3.351342in}{2.304223in}}%
\pgfpathlineto{\pgfqpoint{3.351031in}{2.303820in}}%
\pgfpathlineto{\pgfqpoint{3.351362in}{2.304133in}}%
\pgfpathlineto{\pgfqpoint{3.352052in}{2.303597in}}%
\pgfpathlineto{\pgfqpoint{3.351490in}{2.304141in}}%
\pgfpathlineto{\pgfqpoint{3.352477in}{2.304083in}}%
\pgfpathlineto{\pgfqpoint{3.353023in}{2.304679in}}%
\pgfpathlineto{\pgfqpoint{3.352613in}{2.303996in}}%
\pgfpathlineto{\pgfqpoint{3.353630in}{2.304201in}}%
\pgfpathlineto{\pgfqpoint{3.353807in}{2.303962in}}%
\pgfpathlineto{\pgfqpoint{3.354192in}{2.304596in}}%
\pgfpathlineto{\pgfqpoint{3.354750in}{2.304170in}}%
\pgfpathlineto{\pgfqpoint{3.355754in}{2.304274in}}%
\pgfpathlineto{\pgfqpoint{3.355141in}{2.303899in}}%
\pgfpathlineto{\pgfqpoint{3.355858in}{2.304092in}}%
\pgfpathlineto{\pgfqpoint{3.357209in}{2.304566in}}%
\pgfpathlineto{\pgfqpoint{3.355860in}{2.304080in}}%
\pgfpathlineto{\pgfqpoint{3.357337in}{2.304390in}}%
\pgfpathlineto{\pgfqpoint{3.357414in}{2.304289in}}%
\pgfpathlineto{\pgfqpoint{3.358244in}{2.304666in}}%
\pgfpathlineto{\pgfqpoint{3.358448in}{2.304390in}}%
\pgfpathlineto{\pgfqpoint{3.359064in}{2.304763in}}%
\pgfpathlineto{\pgfqpoint{3.358785in}{2.304263in}}%
\pgfpathlineto{\pgfqpoint{3.359565in}{2.304356in}}%
\pgfpathlineto{\pgfqpoint{3.359582in}{2.304271in}}%
\pgfpathlineto{\pgfqpoint{3.360080in}{2.304669in}}%
\pgfpathlineto{\pgfqpoint{3.360664in}{2.304657in}}%
\pgfpathlineto{\pgfqpoint{3.361747in}{2.305330in}}%
\pgfpathlineto{\pgfqpoint{3.360667in}{2.304635in}}%
\pgfpathlineto{\pgfqpoint{3.361825in}{2.305171in}}%
\pgfpathlineto{\pgfqpoint{3.363322in}{2.305191in}}%
\pgfpathlineto{\pgfqpoint{3.363341in}{2.305263in}}%
\pgfpathlineto{\pgfqpoint{3.363830in}{2.305439in}}%
\pgfpathlineto{\pgfqpoint{3.363424in}{2.305061in}}%
\pgfpathlineto{\pgfqpoint{3.364456in}{2.305317in}}%
\pgfpathlineto{\pgfqpoint{3.365469in}{2.305179in}}%
\pgfpathlineto{\pgfqpoint{3.364789in}{2.305692in}}%
\pgfpathlineto{\pgfqpoint{3.365631in}{2.305386in}}%
\pgfpathlineto{\pgfqpoint{3.366251in}{2.305667in}}%
\pgfpathlineto{\pgfqpoint{3.366562in}{2.303378in}}%
\pgfpathlineto{\pgfqpoint{3.366694in}{2.303405in}}%
\pgfpathlineto{\pgfqpoint{3.367725in}{2.303381in}}%
\pgfpathlineto{\pgfqpoint{3.366780in}{2.303450in}}%
\pgfpathlineto{\pgfqpoint{3.367809in}{2.303407in}}%
\pgfpathlineto{\pgfqpoint{3.368878in}{2.303441in}}%
\pgfpathlineto{\pgfqpoint{3.367904in}{2.303351in}}%
\pgfpathlineto{\pgfqpoint{3.368922in}{2.303414in}}%
\pgfpathlineto{\pgfqpoint{3.369836in}{2.303370in}}%
\pgfpathlineto{\pgfqpoint{3.368970in}{2.303445in}}%
\pgfpathlineto{\pgfqpoint{3.370037in}{2.303407in}}%
\pgfpathlineto{\pgfqpoint{3.371079in}{2.303450in}}%
\pgfpathlineto{\pgfqpoint{3.370355in}{2.303361in}}%
\pgfpathlineto{\pgfqpoint{3.371148in}{2.303431in}}%
\pgfpathlineto{\pgfqpoint{3.372152in}{2.303373in}}%
\pgfpathlineto{\pgfqpoint{3.371207in}{2.303452in}}%
\pgfpathlineto{\pgfqpoint{3.372260in}{2.303394in}}%
\pgfpathlineto{\pgfqpoint{3.373347in}{2.303450in}}%
\pgfpathlineto{\pgfqpoint{3.372911in}{2.303334in}}%
\pgfpathlineto{\pgfqpoint{3.373371in}{2.303400in}}%
\pgfpathlineto{\pgfqpoint{3.374399in}{2.303451in}}%
\pgfpathlineto{\pgfqpoint{3.373380in}{2.303387in}}%
\pgfpathlineto{\pgfqpoint{3.374484in}{2.303401in}}%
\pgfpathlineto{\pgfqpoint{3.375468in}{2.303357in}}%
\pgfpathlineto{\pgfqpoint{3.374500in}{2.303441in}}%
\pgfpathlineto{\pgfqpoint{3.375595in}{2.303402in}}%
\pgfpathlineto{\pgfqpoint{3.376719in}{2.303450in}}%
\pgfpathlineto{\pgfqpoint{3.375673in}{2.303375in}}%
\pgfpathlineto{\pgfqpoint{3.376734in}{2.303377in}}%
\pgfpathlineto{\pgfqpoint{3.378077in}{2.303392in}}%
\pgfpathlineto{\pgfqpoint{3.376739in}{2.303425in}}%
\pgfpathlineto{\pgfqpoint{3.378080in}{2.303420in}}%
\pgfpathlineto{\pgfqpoint{3.379038in}{2.303444in}}%
\pgfpathlineto{\pgfqpoint{3.378491in}{2.303345in}}%
\pgfpathlineto{\pgfqpoint{3.379193in}{2.303432in}}%
\pgfpathlineto{\pgfqpoint{3.380127in}{2.303358in}}%
\pgfpathlineto{\pgfqpoint{3.379218in}{2.303447in}}%
\pgfpathlineto{\pgfqpoint{3.380306in}{2.303396in}}%
\pgfpathlineto{\pgfqpoint{3.381222in}{2.303450in}}%
\pgfpathlineto{\pgfqpoint{3.380601in}{2.303336in}}%
\pgfpathlineto{\pgfqpoint{3.381419in}{2.303421in}}%
\pgfpathlineto{\pgfqpoint{3.382334in}{2.303460in}}%
\pgfpathlineto{\pgfqpoint{3.381499in}{2.303369in}}%
\pgfpathlineto{\pgfqpoint{3.382530in}{2.303414in}}%
\pgfpathlineto{\pgfqpoint{3.383396in}{2.303364in}}%
\pgfpathlineto{\pgfqpoint{3.382621in}{2.303443in}}%
\pgfpathlineto{\pgfqpoint{3.383643in}{2.303417in}}%
\pgfpathlineto{\pgfqpoint{3.384728in}{2.303445in}}%
\pgfpathlineto{\pgfqpoint{3.383764in}{2.303351in}}%
\pgfpathlineto{\pgfqpoint{3.384758in}{2.303411in}}%
\pgfpathlineto{\pgfqpoint{3.385613in}{2.302531in}}%
\pgfpathlineto{\pgfqpoint{3.385103in}{2.303497in}}%
\pgfpathlineto{\pgfqpoint{3.385907in}{2.302832in}}%
\pgfpathlineto{\pgfqpoint{3.386720in}{2.303079in}}%
\pgfpathlineto{\pgfqpoint{3.386321in}{2.302525in}}%
\pgfpathlineto{\pgfqpoint{3.387017in}{2.302770in}}%
\pgfpathlineto{\pgfqpoint{3.387901in}{2.302415in}}%
\pgfpathlineto{\pgfqpoint{3.387087in}{2.302840in}}%
\pgfpathlineto{\pgfqpoint{3.388133in}{2.302649in}}%
\pgfpathlineto{\pgfqpoint{3.389198in}{2.302697in}}%
\pgfpathlineto{\pgfqpoint{3.388742in}{2.302210in}}%
\pgfpathlineto{\pgfqpoint{3.389250in}{2.302630in}}%
\pgfpathlineto{\pgfqpoint{3.390343in}{2.302502in}}%
\pgfpathlineto{\pgfqpoint{3.389485in}{2.303012in}}%
\pgfpathlineto{\pgfqpoint{3.390392in}{2.302577in}}%
\pgfpathlineto{\pgfqpoint{3.391466in}{2.302621in}}%
\pgfpathlineto{\pgfqpoint{3.390780in}{2.302142in}}%
\pgfpathlineto{\pgfqpoint{3.391503in}{2.302563in}}%
\pgfpathlineto{\pgfqpoint{3.391520in}{2.302526in}}%
\pgfpathlineto{\pgfqpoint{3.392452in}{2.303263in}}%
\pgfpathlineto{\pgfqpoint{3.392591in}{2.303016in}}%
\pgfpathlineto{\pgfqpoint{3.393110in}{2.303187in}}%
\pgfpathlineto{\pgfqpoint{3.393608in}{2.302710in}}%
\pgfpathlineto{\pgfqpoint{3.393672in}{2.302727in}}%
\pgfpathlineto{\pgfqpoint{3.394722in}{2.302409in}}%
\pgfpathlineto{\pgfqpoint{3.393970in}{2.303000in}}%
\pgfpathlineto{\pgfqpoint{3.394790in}{2.302610in}}%
\pgfpathlineto{\pgfqpoint{3.395829in}{2.303249in}}%
\pgfpathlineto{\pgfqpoint{3.394825in}{2.302515in}}%
\pgfpathlineto{\pgfqpoint{3.395916in}{2.303029in}}%
\pgfpathlineto{\pgfqpoint{3.396897in}{2.303156in}}%
\pgfpathlineto{\pgfqpoint{3.396640in}{2.302760in}}%
\pgfpathlineto{\pgfqpoint{3.397030in}{2.303055in}}%
\pgfpathlineto{\pgfqpoint{3.397301in}{2.302806in}}%
\pgfpathlineto{\pgfqpoint{3.397609in}{2.303210in}}%
\pgfpathlineto{\pgfqpoint{3.398135in}{2.303176in}}%
\pgfpathlineto{\pgfqpoint{3.399272in}{2.303501in}}%
\pgfpathlineto{\pgfqpoint{3.398408in}{2.303102in}}%
\pgfpathlineto{\pgfqpoint{3.399285in}{2.303423in}}%
\pgfpathlineto{\pgfqpoint{3.400425in}{2.303142in}}%
\pgfpathlineto{\pgfqpoint{3.399758in}{2.303597in}}%
\pgfpathlineto{\pgfqpoint{3.400428in}{2.303157in}}%
\pgfpathlineto{\pgfqpoint{3.401214in}{2.303332in}}%
\pgfpathlineto{\pgfqpoint{3.400664in}{2.302830in}}%
\pgfpathlineto{\pgfqpoint{3.401544in}{2.303227in}}%
\pgfpathlineto{\pgfqpoint{3.402600in}{2.302816in}}%
\pgfpathlineto{\pgfqpoint{3.401715in}{2.303388in}}%
\pgfpathlineto{\pgfqpoint{3.402674in}{2.302984in}}%
\pgfpathlineto{\pgfqpoint{3.405251in}{2.304657in}}%
\pgfpathlineto{\pgfqpoint{3.405258in}{2.304609in}}%
\pgfpathlineto{\pgfqpoint{3.405942in}{2.304215in}}%
\pgfpathlineto{\pgfqpoint{3.405577in}{2.305006in}}%
\pgfpathlineto{\pgfqpoint{3.406390in}{2.304378in}}%
\pgfpathlineto{\pgfqpoint{3.407460in}{2.305140in}}%
\pgfpathlineto{\pgfqpoint{3.406422in}{2.304317in}}%
\pgfpathlineto{\pgfqpoint{3.407541in}{2.305016in}}%
\pgfpathlineto{\pgfqpoint{3.408536in}{2.304147in}}%
\pgfpathlineto{\pgfqpoint{3.407544in}{2.305040in}}%
\pgfpathlineto{\pgfqpoint{3.408812in}{2.304511in}}%
\pgfpathlineto{\pgfqpoint{3.410648in}{2.304902in}}%
\pgfpathlineto{\pgfqpoint{3.408816in}{2.304493in}}%
\pgfpathlineto{\pgfqpoint{3.410684in}{2.304780in}}%
\pgfpathlineto{\pgfqpoint{3.411573in}{2.304450in}}%
\pgfpathlineto{\pgfqpoint{3.410976in}{2.305003in}}%
\pgfpathlineto{\pgfqpoint{3.411799in}{2.304837in}}%
\pgfpathlineto{\pgfqpoint{3.412853in}{2.305630in}}%
\pgfpathlineto{\pgfqpoint{3.411802in}{2.304823in}}%
\pgfpathlineto{\pgfqpoint{3.412951in}{2.305603in}}%
\pgfpathlineto{\pgfqpoint{3.414142in}{2.306482in}}%
\pgfpathlineto{\pgfqpoint{3.412955in}{2.305600in}}%
\pgfpathlineto{\pgfqpoint{3.414184in}{2.306402in}}%
\pgfpathlineto{\pgfqpoint{3.415372in}{2.305545in}}%
\pgfpathlineto{\pgfqpoint{3.414229in}{2.306459in}}%
\pgfpathlineto{\pgfqpoint{3.415377in}{2.305558in}}%
\pgfpathlineto{\pgfqpoint{3.416288in}{2.305221in}}%
\pgfpathlineto{\pgfqpoint{3.415547in}{2.305711in}}%
\pgfpathlineto{\pgfqpoint{3.416587in}{2.305476in}}%
\pgfpathlineto{\pgfqpoint{3.417376in}{2.305589in}}%
\pgfpathlineto{\pgfqpoint{3.416764in}{2.305352in}}%
\pgfpathlineto{\pgfqpoint{3.417691in}{2.305328in}}%
\pgfpathlineto{\pgfqpoint{3.423380in}{2.304725in}}%
\pgfpathlineto{\pgfqpoint{3.417704in}{2.305402in}}%
\pgfpathlineto{\pgfqpoint{3.423394in}{2.304792in}}%
\pgfpathlineto{\pgfqpoint{3.423724in}{2.304968in}}%
\pgfpathlineto{\pgfqpoint{3.424353in}{2.304336in}}%
\pgfpathlineto{\pgfqpoint{3.424479in}{2.304546in}}%
\pgfpathlineto{\pgfqpoint{3.427920in}{2.305229in}}%
\pgfpathlineto{\pgfqpoint{3.424557in}{2.304276in}}%
\pgfpathlineto{\pgfqpoint{3.427939in}{2.305111in}}%
\pgfpathlineto{\pgfqpoint{3.427997in}{2.305022in}}%
\pgfpathlineto{\pgfqpoint{3.428945in}{2.305768in}}%
\pgfpathlineto{\pgfqpoint{3.429021in}{2.305627in}}%
\pgfpathlineto{\pgfqpoint{3.430265in}{2.305335in}}%
\pgfpathlineto{\pgfqpoint{3.429178in}{2.305961in}}%
\pgfpathlineto{\pgfqpoint{3.430342in}{2.305504in}}%
\pgfpathlineto{\pgfqpoint{3.430348in}{2.305536in}}%
\pgfpathlineto{\pgfqpoint{3.431330in}{2.304821in}}%
\pgfpathlineto{\pgfqpoint{3.431435in}{2.304990in}}%
\pgfpathlineto{\pgfqpoint{3.432485in}{2.304885in}}%
\pgfpathlineto{\pgfqpoint{3.431790in}{2.305461in}}%
\pgfpathlineto{\pgfqpoint{3.432561in}{2.304950in}}%
\pgfpathlineto{\pgfqpoint{3.433529in}{2.305416in}}%
\pgfpathlineto{\pgfqpoint{3.432795in}{2.304827in}}%
\pgfpathlineto{\pgfqpoint{3.433688in}{2.305143in}}%
\pgfpathlineto{\pgfqpoint{3.434815in}{2.304635in}}%
\pgfpathlineto{\pgfqpoint{3.434340in}{2.305309in}}%
\pgfpathlineto{\pgfqpoint{3.434816in}{2.304635in}}%
\pgfpathlineto{\pgfqpoint{3.435774in}{2.304609in}}%
\pgfpathlineto{\pgfqpoint{3.435133in}{2.304370in}}%
\pgfpathlineto{\pgfqpoint{3.435901in}{2.304388in}}%
\pgfpathlineto{\pgfqpoint{3.436730in}{2.304182in}}%
\pgfpathlineto{\pgfqpoint{3.436085in}{2.304561in}}%
\pgfpathlineto{\pgfqpoint{3.437013in}{2.304329in}}%
\pgfpathlineto{\pgfqpoint{3.438123in}{2.304379in}}%
\pgfpathlineto{\pgfqpoint{3.437451in}{2.303845in}}%
\pgfpathlineto{\pgfqpoint{3.438127in}{2.304346in}}%
\pgfpathlineto{\pgfqpoint{3.439250in}{2.304050in}}%
\pgfpathlineto{\pgfqpoint{3.438139in}{2.304282in}}%
\pgfpathlineto{\pgfqpoint{3.439256in}{2.304003in}}%
\pgfpathlineto{\pgfqpoint{3.439811in}{2.303494in}}%
\pgfpathlineto{\pgfqpoint{3.439542in}{2.304048in}}%
\pgfpathlineto{\pgfqpoint{3.440386in}{2.303694in}}%
\pgfpathlineto{\pgfqpoint{3.441236in}{2.304189in}}%
\pgfpathlineto{\pgfqpoint{3.441348in}{2.303350in}}%
\pgfpathlineto{\pgfqpoint{3.441462in}{2.303430in}}%
\pgfpathlineto{\pgfqpoint{3.442353in}{2.303349in}}%
\pgfpathlineto{\pgfqpoint{3.441565in}{2.303451in}}%
\pgfpathlineto{\pgfqpoint{3.442575in}{2.303407in}}%
\pgfpathlineto{\pgfqpoint{3.443654in}{2.303445in}}%
\pgfpathlineto{\pgfqpoint{3.442684in}{2.303376in}}%
\pgfpathlineto{\pgfqpoint{3.443690in}{2.303403in}}%
\pgfpathlineto{\pgfqpoint{3.444827in}{2.303376in}}%
\pgfpathlineto{\pgfqpoint{3.443703in}{2.303444in}}%
\pgfpathlineto{\pgfqpoint{3.444848in}{2.303410in}}%
\pgfpathlineto{\pgfqpoint{3.445958in}{2.303445in}}%
\pgfpathlineto{\pgfqpoint{3.444880in}{2.303368in}}%
\pgfpathlineto{\pgfqpoint{3.445967in}{2.303424in}}%
\pgfpathlineto{\pgfqpoint{3.447007in}{2.303375in}}%
\pgfpathlineto{\pgfqpoint{3.446156in}{2.303450in}}%
\pgfpathlineto{\pgfqpoint{3.447080in}{2.303412in}}%
\pgfpathlineto{\pgfqpoint{3.448165in}{2.303448in}}%
\pgfpathlineto{\pgfqpoint{3.447159in}{2.303360in}}%
\pgfpathlineto{\pgfqpoint{3.448196in}{2.303427in}}%
\pgfpathlineto{\pgfqpoint{3.449030in}{2.303354in}}%
\pgfpathlineto{\pgfqpoint{3.448384in}{2.303450in}}%
\pgfpathlineto{\pgfqpoint{3.449308in}{2.303411in}}%
\pgfpathlineto{\pgfqpoint{3.450403in}{2.303455in}}%
\pgfpathlineto{\pgfqpoint{3.449455in}{2.303328in}}%
\pgfpathlineto{\pgfqpoint{3.450419in}{2.303408in}}%
\pgfpathlineto{\pgfqpoint{3.451371in}{2.303357in}}%
\pgfpathlineto{\pgfqpoint{3.450670in}{2.303452in}}%
\pgfpathlineto{\pgfqpoint{3.451529in}{2.303416in}}%
\pgfpathlineto{\pgfqpoint{3.452562in}{2.303446in}}%
\pgfpathlineto{\pgfqpoint{3.451804in}{2.303356in}}%
\pgfpathlineto{\pgfqpoint{3.452640in}{2.303409in}}%
\pgfpathlineto{\pgfqpoint{3.453723in}{2.303443in}}%
\pgfpathlineto{\pgfqpoint{3.452648in}{2.303376in}}%
\pgfpathlineto{\pgfqpoint{3.453758in}{2.303415in}}%
\pgfpathlineto{\pgfqpoint{3.454687in}{2.303334in}}%
\pgfpathlineto{\pgfqpoint{3.453762in}{2.303442in}}%
\pgfpathlineto{\pgfqpoint{3.454873in}{2.303416in}}%
\pgfpathlineto{\pgfqpoint{3.455886in}{2.303448in}}%
\pgfpathlineto{\pgfqpoint{3.454887in}{2.303382in}}%
\pgfpathlineto{\pgfqpoint{3.455985in}{2.303400in}}%
\pgfpathlineto{\pgfqpoint{3.457020in}{2.303346in}}%
\pgfpathlineto{\pgfqpoint{3.456010in}{2.303440in}}%
\pgfpathlineto{\pgfqpoint{3.457127in}{2.303415in}}%
\pgfpathlineto{\pgfqpoint{3.458160in}{2.303452in}}%
\pgfpathlineto{\pgfqpoint{3.457216in}{2.303318in}}%
\pgfpathlineto{\pgfqpoint{3.458239in}{2.303401in}}%
\pgfpathlineto{\pgfqpoint{3.459399in}{2.303436in}}%
\pgfpathlineto{\pgfqpoint{3.458241in}{2.303383in}}%
\pgfpathlineto{\pgfqpoint{3.459423in}{2.303404in}}%
\pgfpathlineto{\pgfqpoint{3.460551in}{2.302999in}}%
\pgfpathlineto{\pgfqpoint{3.459436in}{2.303443in}}%
\pgfpathlineto{\pgfqpoint{3.460610in}{2.303134in}}%
\pgfpathlineto{\pgfqpoint{3.461314in}{2.303392in}}%
\pgfpathlineto{\pgfqpoint{3.460873in}{2.302815in}}%
\pgfpathlineto{\pgfqpoint{3.461722in}{2.303158in}}%
\pgfpathlineto{\pgfqpoint{3.462430in}{2.302335in}}%
\pgfpathlineto{\pgfqpoint{3.461817in}{2.303195in}}%
\pgfpathlineto{\pgfqpoint{3.462845in}{2.302861in}}%
\pgfpathlineto{\pgfqpoint{3.463973in}{2.303454in}}%
\pgfpathlineto{\pgfqpoint{3.463990in}{2.303362in}}%
\pgfpathlineto{\pgfqpoint{3.465042in}{2.302649in}}%
\pgfpathlineto{\pgfqpoint{3.464256in}{2.303447in}}%
\pgfpathlineto{\pgfqpoint{3.465151in}{2.302746in}}%
\pgfpathlineto{\pgfqpoint{3.465686in}{2.303202in}}%
\pgfpathlineto{\pgfqpoint{3.466274in}{2.302854in}}%
\pgfpathlineto{\pgfqpoint{3.467376in}{2.301940in}}%
\pgfpathlineto{\pgfqpoint{3.466336in}{2.302992in}}%
\pgfpathlineto{\pgfqpoint{3.467926in}{2.302502in}}%
\pgfpathlineto{\pgfqpoint{3.468474in}{2.302696in}}%
\pgfpathlineto{\pgfqpoint{3.469021in}{2.302181in}}%
\pgfpathlineto{\pgfqpoint{3.469231in}{2.301776in}}%
\pgfpathlineto{\pgfqpoint{3.470145in}{2.302008in}}%
\pgfpathlineto{\pgfqpoint{3.471997in}{2.303109in}}%
\pgfpathlineto{\pgfqpoint{3.470151in}{2.301994in}}%
\pgfpathlineto{\pgfqpoint{3.472253in}{2.302715in}}%
\pgfpathlineto{\pgfqpoint{3.473151in}{2.302088in}}%
\pgfpathlineto{\pgfqpoint{3.472265in}{2.302721in}}%
\pgfpathlineto{\pgfqpoint{3.473442in}{2.302144in}}%
\pgfpathlineto{\pgfqpoint{3.473653in}{2.302450in}}%
\pgfpathlineto{\pgfqpoint{3.474535in}{2.301873in}}%
\pgfpathlineto{\pgfqpoint{3.474538in}{2.301902in}}%
\pgfpathlineto{\pgfqpoint{3.475669in}{2.301201in}}%
\pgfpathlineto{\pgfqpoint{3.474888in}{2.301988in}}%
\pgfpathlineto{\pgfqpoint{3.475690in}{2.301295in}}%
\pgfpathlineto{\pgfqpoint{3.475711in}{2.301384in}}%
\pgfpathlineto{\pgfqpoint{3.476497in}{2.300665in}}%
\pgfpathlineto{\pgfqpoint{3.476785in}{2.300826in}}%
\pgfpathlineto{\pgfqpoint{3.478606in}{2.300591in}}%
\pgfpathlineto{\pgfqpoint{3.476787in}{2.300816in}}%
\pgfpathlineto{\pgfqpoint{3.478617in}{2.300521in}}%
\pgfpathlineto{\pgfqpoint{3.478921in}{2.300294in}}%
\pgfpathlineto{\pgfqpoint{3.479266in}{2.300735in}}%
\pgfpathlineto{\pgfqpoint{3.479717in}{2.300694in}}%
\pgfpathlineto{\pgfqpoint{3.480881in}{2.300680in}}%
\pgfpathlineto{\pgfqpoint{3.479734in}{2.300719in}}%
\pgfpathlineto{\pgfqpoint{3.480905in}{2.300727in}}%
\pgfpathlineto{\pgfqpoint{3.482017in}{2.300978in}}%
\pgfpathlineto{\pgfqpoint{3.481274in}{2.300546in}}%
\pgfpathlineto{\pgfqpoint{3.482029in}{2.300921in}}%
\pgfpathlineto{\pgfqpoint{3.482930in}{2.300545in}}%
\pgfpathlineto{\pgfqpoint{3.482143in}{2.301086in}}%
\pgfpathlineto{\pgfqpoint{3.483144in}{2.300881in}}%
\pgfpathlineto{\pgfqpoint{3.484439in}{2.300308in}}%
\pgfpathlineto{\pgfqpoint{3.484447in}{2.300376in}}%
\pgfpathlineto{\pgfqpoint{3.485583in}{2.300714in}}%
\pgfpathlineto{\pgfqpoint{3.484985in}{2.300158in}}%
\pgfpathlineto{\pgfqpoint{3.485598in}{2.300684in}}%
\pgfpathlineto{\pgfqpoint{3.486185in}{2.300596in}}%
\pgfpathlineto{\pgfqpoint{3.485705in}{2.300850in}}%
\pgfpathlineto{\pgfqpoint{3.486719in}{2.300809in}}%
\pgfpathlineto{\pgfqpoint{3.487616in}{2.301017in}}%
\pgfpathlineto{\pgfqpoint{3.487171in}{2.300597in}}%
\pgfpathlineto{\pgfqpoint{3.487840in}{2.300819in}}%
\pgfpathlineto{\pgfqpoint{3.488560in}{2.300301in}}%
\pgfpathlineto{\pgfqpoint{3.487884in}{2.300947in}}%
\pgfpathlineto{\pgfqpoint{3.488962in}{2.300667in}}%
\pgfpathlineto{\pgfqpoint{3.489121in}{2.300769in}}%
\pgfpathlineto{\pgfqpoint{3.490033in}{2.300010in}}%
\pgfpathlineto{\pgfqpoint{3.490046in}{2.300036in}}%
\pgfpathlineto{\pgfqpoint{3.490627in}{2.300150in}}%
\pgfpathlineto{\pgfqpoint{3.490225in}{2.299682in}}%
\pgfpathlineto{\pgfqpoint{3.491143in}{2.299787in}}%
\pgfpathlineto{\pgfqpoint{3.491943in}{2.299579in}}%
\pgfpathlineto{\pgfqpoint{3.491274in}{2.299608in}}%
\pgfpathlineto{\pgfqpoint{3.492105in}{2.299230in}}%
\pgfpathlineto{\pgfqpoint{3.492212in}{2.299193in}}%
\pgfpathlineto{\pgfqpoint{3.493065in}{2.300383in}}%
\pgfpathlineto{\pgfqpoint{3.493165in}{2.300304in}}%
\pgfpathlineto{\pgfqpoint{3.493744in}{2.300036in}}%
\pgfpathlineto{\pgfqpoint{3.494208in}{2.300587in}}%
\pgfpathlineto{\pgfqpoint{3.494255in}{2.300538in}}%
\pgfpathlineto{\pgfqpoint{3.495312in}{2.300889in}}%
\pgfpathlineto{\pgfqpoint{3.494592in}{2.300084in}}%
\pgfpathlineto{\pgfqpoint{3.495395in}{2.300810in}}%
\pgfpathlineto{\pgfqpoint{3.495613in}{2.300896in}}%
\pgfpathlineto{\pgfqpoint{3.496151in}{2.301651in}}%
\pgfpathlineto{\pgfqpoint{3.496166in}{2.301655in}}%
\pgfpathlineto{\pgfqpoint{3.496222in}{2.301761in}}%
\pgfpathlineto{\pgfqpoint{3.497203in}{2.301266in}}%
\pgfpathlineto{\pgfqpoint{3.497266in}{2.301313in}}%
\pgfpathlineto{\pgfqpoint{3.497365in}{2.301260in}}%
\pgfpathlineto{\pgfqpoint{3.497982in}{2.301810in}}%
\pgfpathlineto{\pgfqpoint{3.498365in}{2.301484in}}%
\pgfpathlineto{\pgfqpoint{3.499505in}{2.301705in}}%
\pgfpathlineto{\pgfqpoint{3.498437in}{2.301387in}}%
\pgfpathlineto{\pgfqpoint{3.499510in}{2.301660in}}%
\pgfpathlineto{\pgfqpoint{3.500587in}{2.301644in}}%
\pgfpathlineto{\pgfqpoint{3.500420in}{2.301941in}}%
\pgfpathlineto{\pgfqpoint{3.500620in}{2.301725in}}%
\pgfpathlineto{\pgfqpoint{3.501452in}{2.301471in}}%
\pgfpathlineto{\pgfqpoint{3.501748in}{2.301874in}}%
\pgfpathlineto{\pgfqpoint{3.502032in}{2.301704in}}%
\pgfpathlineto{\pgfqpoint{3.502790in}{2.302369in}}%
\pgfpathlineto{\pgfqpoint{3.502830in}{2.302222in}}%
\pgfpathlineto{\pgfqpoint{3.502885in}{2.302302in}}%
\pgfpathlineto{\pgfqpoint{3.503821in}{2.301489in}}%
\pgfpathlineto{\pgfqpoint{3.503867in}{2.301547in}}%
\pgfpathlineto{\pgfqpoint{3.504710in}{2.301524in}}%
\pgfpathlineto{\pgfqpoint{3.504318in}{2.301808in}}%
\pgfpathlineto{\pgfqpoint{3.504953in}{2.301854in}}%
\pgfpathlineto{\pgfqpoint{3.505023in}{2.301947in}}%
\pgfpathlineto{\pgfqpoint{3.505854in}{2.301342in}}%
\pgfpathlineto{\pgfqpoint{3.506049in}{2.301416in}}%
\pgfpathlineto{\pgfqpoint{3.506189in}{2.301211in}}%
\pgfpathlineto{\pgfqpoint{3.507121in}{2.301983in}}%
\pgfpathlineto{\pgfqpoint{3.507140in}{2.301887in}}%
\pgfpathlineto{\pgfqpoint{3.507227in}{2.301798in}}%
\pgfpathlineto{\pgfqpoint{3.507411in}{2.302221in}}%
\pgfpathlineto{\pgfqpoint{3.508250in}{2.301946in}}%
\pgfpathlineto{\pgfqpoint{3.509993in}{2.302254in}}%
\pgfpathlineto{\pgfqpoint{3.508455in}{2.301600in}}%
\pgfpathlineto{\pgfqpoint{3.510006in}{2.302211in}}%
\pgfpathlineto{\pgfqpoint{3.510421in}{2.301804in}}%
\pgfpathlineto{\pgfqpoint{3.510059in}{2.302316in}}%
\pgfpathlineto{\pgfqpoint{3.511118in}{2.302243in}}%
\pgfpathlineto{\pgfqpoint{3.512272in}{2.302924in}}%
\pgfpathlineto{\pgfqpoint{3.511248in}{2.302225in}}%
\pgfpathlineto{\pgfqpoint{3.512276in}{2.302922in}}%
\pgfpathlineto{\pgfqpoint{3.513315in}{2.302383in}}%
\pgfpathlineto{\pgfqpoint{3.512328in}{2.302979in}}%
\pgfpathlineto{\pgfqpoint{3.513452in}{2.302433in}}%
\pgfpathlineto{\pgfqpoint{3.514450in}{2.302492in}}%
\pgfpathlineto{\pgfqpoint{3.514135in}{2.302054in}}%
\pgfpathlineto{\pgfqpoint{3.514579in}{2.302371in}}%
\pgfpathlineto{\pgfqpoint{3.515174in}{2.302195in}}%
\pgfpathlineto{\pgfqpoint{3.514699in}{2.302632in}}%
\pgfpathlineto{\pgfqpoint{3.515693in}{2.302343in}}%
\pgfpathlineto{\pgfqpoint{3.516871in}{2.303444in}}%
\pgfpathlineto{\pgfqpoint{3.515888in}{2.302285in}}%
\pgfpathlineto{\pgfqpoint{3.516878in}{2.303432in}}%
\pgfpathlineto{\pgfqpoint{3.517752in}{2.303356in}}%
\pgfpathlineto{\pgfqpoint{3.517283in}{2.303455in}}%
\pgfpathlineto{\pgfqpoint{3.517991in}{2.303402in}}%
\pgfpathlineto{\pgfqpoint{3.519047in}{2.303448in}}%
\pgfpathlineto{\pgfqpoint{3.517999in}{2.303383in}}%
\pgfpathlineto{\pgfqpoint{3.519103in}{2.303421in}}%
\pgfpathlineto{\pgfqpoint{3.520105in}{2.303374in}}%
\pgfpathlineto{\pgfqpoint{3.519117in}{2.303448in}}%
\pgfpathlineto{\pgfqpoint{3.520223in}{2.303407in}}%
\pgfpathlineto{\pgfqpoint{3.521339in}{2.303441in}}%
\pgfpathlineto{\pgfqpoint{3.520236in}{2.303367in}}%
\pgfpathlineto{\pgfqpoint{3.521345in}{2.303424in}}%
\pgfpathlineto{\pgfqpoint{3.521893in}{2.303323in}}%
\pgfpathlineto{\pgfqpoint{3.521356in}{2.303442in}}%
\pgfpathlineto{\pgfqpoint{3.522458in}{2.303407in}}%
\pgfpathlineto{\pgfqpoint{3.523526in}{2.303377in}}%
\pgfpathlineto{\pgfqpoint{3.522515in}{2.303457in}}%
\pgfpathlineto{\pgfqpoint{3.523570in}{2.303424in}}%
\pgfpathlineto{\pgfqpoint{3.524673in}{2.303443in}}%
\pgfpathlineto{\pgfqpoint{3.523631in}{2.303360in}}%
\pgfpathlineto{\pgfqpoint{3.524682in}{2.303391in}}%
\pgfpathlineto{\pgfqpoint{3.525783in}{2.303386in}}%
\pgfpathlineto{\pgfqpoint{3.524705in}{2.303439in}}%
\pgfpathlineto{\pgfqpoint{3.525812in}{2.303426in}}%
\pgfpathlineto{\pgfqpoint{3.526916in}{2.303437in}}%
\pgfpathlineto{\pgfqpoint{3.525929in}{2.303353in}}%
\pgfpathlineto{\pgfqpoint{3.526924in}{2.303427in}}%
\pgfpathlineto{\pgfqpoint{3.527649in}{2.303352in}}%
\pgfpathlineto{\pgfqpoint{3.526994in}{2.303445in}}%
\pgfpathlineto{\pgfqpoint{3.528036in}{2.303398in}}%
\pgfpathlineto{\pgfqpoint{3.529142in}{2.303445in}}%
\pgfpathlineto{\pgfqpoint{3.528068in}{2.303338in}}%
\pgfpathlineto{\pgfqpoint{3.529177in}{2.303437in}}%
\pgfpathlineto{\pgfqpoint{3.530207in}{2.303352in}}%
\pgfpathlineto{\pgfqpoint{3.529190in}{2.303450in}}%
\pgfpathlineto{\pgfqpoint{3.530288in}{2.303417in}}%
\pgfpathlineto{\pgfqpoint{3.531390in}{2.303440in}}%
\pgfpathlineto{\pgfqpoint{3.530376in}{2.303350in}}%
\pgfpathlineto{\pgfqpoint{3.531404in}{2.303438in}}%
\pgfpathlineto{\pgfqpoint{3.532401in}{2.303330in}}%
\pgfpathlineto{\pgfqpoint{3.531437in}{2.303446in}}%
\pgfpathlineto{\pgfqpoint{3.532515in}{2.303405in}}%
\pgfpathlineto{\pgfqpoint{3.533690in}{2.303432in}}%
\pgfpathlineto{\pgfqpoint{3.532563in}{2.303363in}}%
\pgfpathlineto{\pgfqpoint{3.533694in}{2.303426in}}%
\pgfpathlineto{\pgfqpoint{3.534664in}{2.303362in}}%
\pgfpathlineto{\pgfqpoint{3.534015in}{2.303459in}}%
\pgfpathlineto{\pgfqpoint{3.534806in}{2.303400in}}%
\pgfpathlineto{\pgfqpoint{3.536111in}{2.304360in}}%
\pgfpathlineto{\pgfqpoint{3.534888in}{2.303363in}}%
\pgfpathlineto{\pgfqpoint{3.536147in}{2.304292in}}%
\pgfpathlineto{\pgfqpoint{3.537293in}{2.303891in}}%
\pgfpathlineto{\pgfqpoint{3.536774in}{2.304504in}}%
\pgfpathlineto{\pgfqpoint{3.537297in}{2.303897in}}%
\pgfpathlineto{\pgfqpoint{3.537978in}{2.303651in}}%
\pgfpathlineto{\pgfqpoint{3.537345in}{2.303971in}}%
\pgfpathlineto{\pgfqpoint{3.538423in}{2.303938in}}%
\pgfpathlineto{\pgfqpoint{3.538961in}{2.304200in}}%
\pgfpathlineto{\pgfqpoint{3.538461in}{2.303864in}}%
\pgfpathlineto{\pgfqpoint{3.539535in}{2.303899in}}%
\pgfpathlineto{\pgfqpoint{3.540427in}{2.303641in}}%
\pgfpathlineto{\pgfqpoint{3.540631in}{2.304119in}}%
\pgfpathlineto{\pgfqpoint{3.541741in}{2.304827in}}%
\pgfpathlineto{\pgfqpoint{3.540691in}{2.304094in}}%
\pgfpathlineto{\pgfqpoint{3.541787in}{2.304731in}}%
\pgfpathlineto{\pgfqpoint{3.542904in}{2.304444in}}%
\pgfpathlineto{\pgfqpoint{3.542242in}{2.304917in}}%
\pgfpathlineto{\pgfqpoint{3.542916in}{2.304453in}}%
\pgfpathlineto{\pgfqpoint{3.543261in}{2.304623in}}%
\pgfpathlineto{\pgfqpoint{3.543905in}{2.304065in}}%
\pgfpathlineto{\pgfqpoint{3.544022in}{2.304306in}}%
\pgfpathlineto{\pgfqpoint{3.544863in}{2.303787in}}%
\pgfpathlineto{\pgfqpoint{3.545147in}{2.304210in}}%
\pgfpathlineto{\pgfqpoint{3.545309in}{2.304600in}}%
\pgfpathlineto{\pgfqpoint{3.546269in}{2.304318in}}%
\pgfpathlineto{\pgfqpoint{3.547215in}{2.304348in}}%
\pgfpathlineto{\pgfqpoint{3.547101in}{2.304068in}}%
\pgfpathlineto{\pgfqpoint{3.547354in}{2.304082in}}%
\pgfpathlineto{\pgfqpoint{3.547362in}{2.304070in}}%
\pgfpathlineto{\pgfqpoint{3.547785in}{2.304678in}}%
\pgfpathlineto{\pgfqpoint{3.548416in}{2.304510in}}%
\pgfpathlineto{\pgfqpoint{3.549398in}{2.305088in}}%
\pgfpathlineto{\pgfqpoint{3.548701in}{2.304461in}}%
\pgfpathlineto{\pgfqpoint{3.549544in}{2.304971in}}%
\pgfpathlineto{\pgfqpoint{3.550612in}{2.304793in}}%
\pgfpathlineto{\pgfqpoint{3.550153in}{2.305160in}}%
\pgfpathlineto{\pgfqpoint{3.550661in}{2.304889in}}%
\pgfpathlineto{\pgfqpoint{3.551596in}{2.304922in}}%
\pgfpathlineto{\pgfqpoint{3.551331in}{2.305269in}}%
\pgfpathlineto{\pgfqpoint{3.551762in}{2.305039in}}%
\pgfpathlineto{\pgfqpoint{3.553849in}{2.305211in}}%
\pgfpathlineto{\pgfqpoint{3.551775in}{2.304959in}}%
\pgfpathlineto{\pgfqpoint{3.553872in}{2.305107in}}%
\pgfpathlineto{\pgfqpoint{3.554193in}{2.304904in}}%
\pgfpathlineto{\pgfqpoint{3.554784in}{2.305289in}}%
\pgfpathlineto{\pgfqpoint{3.554980in}{2.305187in}}%
\pgfpathlineto{\pgfqpoint{3.556372in}{2.305413in}}%
\pgfpathlineto{\pgfqpoint{3.554981in}{2.305183in}}%
\pgfpathlineto{\pgfqpoint{3.556545in}{2.305185in}}%
\pgfpathlineto{\pgfqpoint{3.556605in}{2.305098in}}%
\pgfpathlineto{\pgfqpoint{3.557551in}{2.305902in}}%
\pgfpathlineto{\pgfqpoint{3.557634in}{2.305858in}}%
\pgfpathlineto{\pgfqpoint{3.558775in}{2.305345in}}%
\pgfpathlineto{\pgfqpoint{3.557889in}{2.306165in}}%
\pgfpathlineto{\pgfqpoint{3.558781in}{2.305364in}}%
\pgfpathlineto{\pgfqpoint{3.558783in}{2.305380in}}%
\pgfpathlineto{\pgfqpoint{3.559730in}{2.304622in}}%
\pgfpathlineto{\pgfqpoint{3.559857in}{2.304739in}}%
\pgfpathlineto{\pgfqpoint{3.560126in}{2.304628in}}%
\pgfpathlineto{\pgfqpoint{3.560521in}{2.305203in}}%
\pgfpathlineto{\pgfqpoint{3.560931in}{2.305041in}}%
\pgfpathlineto{\pgfqpoint{3.564288in}{2.303822in}}%
\pgfpathlineto{\pgfqpoint{3.565247in}{2.303661in}}%
\pgfpathlineto{\pgfqpoint{3.565390in}{2.304029in}}%
\pgfpathlineto{\pgfqpoint{3.565393in}{2.303999in}}%
\pgfpathlineto{\pgfqpoint{3.565471in}{2.303797in}}%
\pgfpathlineto{\pgfqpoint{3.566452in}{2.304484in}}%
\pgfpathlineto{\pgfqpoint{3.566452in}{2.304481in}}%
\pgfpathlineto{\pgfqpoint{3.566502in}{2.304608in}}%
\pgfpathlineto{\pgfqpoint{3.567440in}{2.304271in}}%
\pgfpathlineto{\pgfqpoint{3.567556in}{2.304269in}}%
\pgfpathlineto{\pgfqpoint{3.568782in}{2.303966in}}%
\pgfpathlineto{\pgfqpoint{3.567812in}{2.304462in}}%
\pgfpathlineto{\pgfqpoint{3.568788in}{2.303995in}}%
\pgfpathlineto{\pgfqpoint{3.569631in}{2.304440in}}%
\pgfpathlineto{\pgfqpoint{3.568939in}{2.303928in}}%
\pgfpathlineto{\pgfqpoint{3.569912in}{2.304282in}}%
\pgfpathlineto{\pgfqpoint{3.570948in}{2.304038in}}%
\pgfpathlineto{\pgfqpoint{3.570261in}{2.304674in}}%
\pgfpathlineto{\pgfqpoint{3.571036in}{2.304095in}}%
\pgfpathlineto{\pgfqpoint{3.571942in}{2.304498in}}%
\pgfpathlineto{\pgfqpoint{3.571066in}{2.304074in}}%
\pgfpathlineto{\pgfqpoint{3.572155in}{2.304260in}}%
\pgfpathlineto{\pgfqpoint{3.572595in}{2.303993in}}%
\pgfpathlineto{\pgfqpoint{3.573090in}{2.304655in}}%
\pgfpathlineto{\pgfqpoint{3.573241in}{2.304553in}}%
\pgfpathlineto{\pgfqpoint{3.573724in}{2.304807in}}%
\pgfpathlineto{\pgfqpoint{3.574336in}{2.304091in}}%
\pgfpathlineto{\pgfqpoint{3.574493in}{2.304064in}}%
\pgfpathlineto{\pgfqpoint{3.574931in}{2.303448in}}%
\pgfpathlineto{\pgfqpoint{3.575283in}{2.303525in}}%
\pgfpathlineto{\pgfqpoint{3.575944in}{2.303249in}}%
\pgfpathlineto{\pgfqpoint{3.575297in}{2.303594in}}%
\pgfpathlineto{\pgfqpoint{3.576404in}{2.303348in}}%
\pgfpathlineto{\pgfqpoint{3.577523in}{2.303183in}}%
\pgfpathlineto{\pgfqpoint{3.577043in}{2.303754in}}%
\pgfpathlineto{\pgfqpoint{3.577559in}{2.303249in}}%
\pgfpathlineto{\pgfqpoint{3.578682in}{2.303630in}}%
\pgfpathlineto{\pgfqpoint{3.577935in}{2.303137in}}%
\pgfpathlineto{\pgfqpoint{3.578695in}{2.303542in}}%
\pgfpathlineto{\pgfqpoint{3.579042in}{2.303722in}}%
\pgfpathlineto{\pgfqpoint{3.579258in}{2.303257in}}%
\pgfpathlineto{\pgfqpoint{3.579805in}{2.303470in}}%
\pgfpathlineto{\pgfqpoint{3.581156in}{2.303352in}}%
\pgfpathlineto{\pgfqpoint{3.579829in}{2.303517in}}%
\pgfpathlineto{\pgfqpoint{3.581176in}{2.303432in}}%
\pgfpathlineto{\pgfqpoint{3.581357in}{2.303547in}}%
\pgfpathlineto{\pgfqpoint{3.582214in}{2.302885in}}%
\pgfpathlineto{\pgfqpoint{3.582258in}{2.303023in}}%
\pgfpathlineto{\pgfqpoint{3.583155in}{2.302386in}}%
\pgfpathlineto{\pgfqpoint{3.582261in}{2.303023in}}%
\pgfpathlineto{\pgfqpoint{3.583383in}{2.302613in}}%
\pgfpathlineto{\pgfqpoint{3.583583in}{2.302555in}}%
\pgfpathlineto{\pgfqpoint{3.584219in}{2.303091in}}%
\pgfpathlineto{\pgfqpoint{3.584449in}{2.302927in}}%
\pgfpathlineto{\pgfqpoint{3.585806in}{2.302274in}}%
\pgfpathlineto{\pgfqpoint{3.584459in}{2.303008in}}%
\pgfpathlineto{\pgfqpoint{3.586013in}{2.302663in}}%
\pgfpathlineto{\pgfqpoint{3.586061in}{2.302766in}}%
\pgfpathlineto{\pgfqpoint{3.586347in}{2.302394in}}%
\pgfpathlineto{\pgfqpoint{3.587108in}{2.302379in}}%
\pgfpathlineto{\pgfqpoint{3.588074in}{2.302140in}}%
\pgfpathlineto{\pgfqpoint{3.587498in}{2.302515in}}%
\pgfpathlineto{\pgfqpoint{3.588230in}{2.302390in}}%
\pgfpathlineto{\pgfqpoint{3.589248in}{2.302911in}}%
\pgfpathlineto{\pgfqpoint{3.588249in}{2.302310in}}%
\pgfpathlineto{\pgfqpoint{3.589375in}{2.302900in}}%
\pgfpathlineto{\pgfqpoint{3.589383in}{2.302851in}}%
\pgfpathlineto{\pgfqpoint{3.590041in}{2.303563in}}%
\pgfpathlineto{\pgfqpoint{3.590437in}{2.303525in}}%
\pgfpathlineto{\pgfqpoint{3.590901in}{2.303727in}}%
\pgfpathlineto{\pgfqpoint{3.591103in}{2.303370in}}%
\pgfpathlineto{\pgfqpoint{3.591542in}{2.303412in}}%
\pgfpathlineto{\pgfqpoint{3.592626in}{2.303337in}}%
\pgfpathlineto{\pgfqpoint{3.591624in}{2.303448in}}%
\pgfpathlineto{\pgfqpoint{3.592653in}{2.303423in}}%
\pgfpathlineto{\pgfqpoint{3.593732in}{2.303453in}}%
\pgfpathlineto{\pgfqpoint{3.592670in}{2.303363in}}%
\pgfpathlineto{\pgfqpoint{3.593768in}{2.303428in}}%
\pgfpathlineto{\pgfqpoint{3.594375in}{2.303341in}}%
\pgfpathlineto{\pgfqpoint{3.593843in}{2.303445in}}%
\pgfpathlineto{\pgfqpoint{3.594881in}{2.303416in}}%
\pgfpathlineto{\pgfqpoint{3.595989in}{2.303446in}}%
\pgfpathlineto{\pgfqpoint{3.594907in}{2.303379in}}%
\pgfpathlineto{\pgfqpoint{3.596000in}{2.303399in}}%
\pgfpathlineto{\pgfqpoint{3.597031in}{2.303351in}}%
\pgfpathlineto{\pgfqpoint{3.596007in}{2.303455in}}%
\pgfpathlineto{\pgfqpoint{3.597123in}{2.303411in}}%
\pgfpathlineto{\pgfqpoint{3.597949in}{2.303452in}}%
\pgfpathlineto{\pgfqpoint{3.597907in}{2.303357in}}%
\pgfpathlineto{\pgfqpoint{3.598235in}{2.303423in}}%
\pgfpathlineto{\pgfqpoint{3.599339in}{2.303442in}}%
\pgfpathlineto{\pgfqpoint{3.598357in}{2.303342in}}%
\pgfpathlineto{\pgfqpoint{3.599346in}{2.303396in}}%
\pgfpathlineto{\pgfqpoint{3.600446in}{2.303376in}}%
\pgfpathlineto{\pgfqpoint{3.599374in}{2.303450in}}%
\pgfpathlineto{\pgfqpoint{3.600460in}{2.303392in}}%
\pgfpathlineto{\pgfqpoint{3.605486in}{2.303417in}}%
\pgfpathlineto{\pgfqpoint{3.600467in}{2.303433in}}%
\pgfpathlineto{\pgfqpoint{3.605486in}{2.303422in}}%
\pgfpathlineto{\pgfqpoint{3.606543in}{2.303448in}}%
\pgfpathlineto{\pgfqpoint{3.605529in}{2.303374in}}%
\pgfpathlineto{\pgfqpoint{3.606600in}{2.303413in}}%
\pgfpathlineto{\pgfqpoint{3.607624in}{2.303362in}}%
\pgfpathlineto{\pgfqpoint{3.606647in}{2.303442in}}%
\pgfpathlineto{\pgfqpoint{3.607719in}{2.303407in}}%
\pgfpathlineto{\pgfqpoint{3.608772in}{2.303457in}}%
\pgfpathlineto{\pgfqpoint{3.607949in}{2.303328in}}%
\pgfpathlineto{\pgfqpoint{3.608830in}{2.303406in}}%
\pgfpathlineto{\pgfqpoint{3.609461in}{2.303324in}}%
\pgfpathlineto{\pgfqpoint{3.609914in}{2.303744in}}%
\pgfpathlineto{\pgfqpoint{3.609915in}{2.303727in}}%
\pgfpathlineto{\pgfqpoint{3.610568in}{2.303861in}}%
\pgfpathlineto{\pgfqpoint{3.610028in}{2.303396in}}%
\pgfpathlineto{\pgfqpoint{3.611037in}{2.303772in}}%
\pgfpathlineto{\pgfqpoint{3.611465in}{2.303482in}}%
\pgfpathlineto{\pgfqpoint{3.612016in}{2.304108in}}%
\pgfpathlineto{\pgfqpoint{3.612147in}{2.303804in}}%
\pgfpathlineto{\pgfqpoint{3.613105in}{2.304296in}}%
\pgfpathlineto{\pgfqpoint{3.612222in}{2.303687in}}%
\pgfpathlineto{\pgfqpoint{3.613277in}{2.303960in}}%
\pgfpathlineto{\pgfqpoint{3.614206in}{2.303933in}}%
\pgfpathlineto{\pgfqpoint{3.613809in}{2.304360in}}%
\pgfpathlineto{\pgfqpoint{3.614370in}{2.304169in}}%
\pgfpathlineto{\pgfqpoint{3.615258in}{2.304323in}}%
\pgfpathlineto{\pgfqpoint{3.614806in}{2.304083in}}%
\pgfpathlineto{\pgfqpoint{3.615483in}{2.304221in}}%
\pgfpathlineto{\pgfqpoint{3.616887in}{2.303815in}}%
\pgfpathlineto{\pgfqpoint{3.615514in}{2.304243in}}%
\pgfpathlineto{\pgfqpoint{3.616896in}{2.303885in}}%
\pgfpathlineto{\pgfqpoint{3.617040in}{2.304099in}}%
\pgfpathlineto{\pgfqpoint{3.617841in}{2.303528in}}%
\pgfpathlineto{\pgfqpoint{3.617998in}{2.303656in}}%
\pgfpathlineto{\pgfqpoint{3.618682in}{2.303446in}}%
\pgfpathlineto{\pgfqpoint{3.619097in}{2.303941in}}%
\pgfpathlineto{\pgfqpoint{3.619099in}{2.303916in}}%
\pgfpathlineto{\pgfqpoint{3.619852in}{2.304467in}}%
\pgfpathlineto{\pgfqpoint{3.620272in}{2.304276in}}%
\pgfpathlineto{\pgfqpoint{3.621105in}{2.303941in}}%
\pgfpathlineto{\pgfqpoint{3.620274in}{2.304289in}}%
\pgfpathlineto{\pgfqpoint{3.621390in}{2.304366in}}%
\pgfpathlineto{\pgfqpoint{3.621512in}{2.304645in}}%
\pgfpathlineto{\pgfqpoint{3.622137in}{2.304093in}}%
\pgfpathlineto{\pgfqpoint{3.622502in}{2.304367in}}%
\pgfpathlineto{\pgfqpoint{3.623460in}{2.303847in}}%
\pgfpathlineto{\pgfqpoint{3.622611in}{2.304492in}}%
\pgfpathlineto{\pgfqpoint{3.623668in}{2.304191in}}%
\pgfpathlineto{\pgfqpoint{3.624523in}{2.304527in}}%
\pgfpathlineto{\pgfqpoint{3.623985in}{2.304128in}}%
\pgfpathlineto{\pgfqpoint{3.624799in}{2.304303in}}%
\pgfpathlineto{\pgfqpoint{3.624952in}{2.304176in}}%
\pgfpathlineto{\pgfqpoint{3.625703in}{2.305062in}}%
\pgfpathlineto{\pgfqpoint{3.625874in}{2.305031in}}%
\pgfpathlineto{\pgfqpoint{3.625892in}{2.305113in}}%
\pgfpathlineto{\pgfqpoint{3.626640in}{2.304227in}}%
\pgfpathlineto{\pgfqpoint{3.626945in}{2.304375in}}%
\pgfpathlineto{\pgfqpoint{3.628087in}{2.304130in}}%
\pgfpathlineto{\pgfqpoint{3.627303in}{2.304520in}}%
\pgfpathlineto{\pgfqpoint{3.628125in}{2.304207in}}%
\pgfpathlineto{\pgfqpoint{3.628404in}{2.304408in}}%
\pgfpathlineto{\pgfqpoint{3.629159in}{2.303984in}}%
\pgfpathlineto{\pgfqpoint{3.629224in}{2.303995in}}%
\pgfpathlineto{\pgfqpoint{3.630269in}{2.303467in}}%
\pgfpathlineto{\pgfqpoint{3.629273in}{2.304110in}}%
\pgfpathlineto{\pgfqpoint{3.630358in}{2.303568in}}%
\pgfpathlineto{\pgfqpoint{3.631508in}{2.304298in}}%
\pgfpathlineto{\pgfqpoint{3.631523in}{2.304278in}}%
\pgfpathlineto{\pgfqpoint{3.631795in}{2.304105in}}%
\pgfpathlineto{\pgfqpoint{3.632114in}{2.304528in}}%
\pgfpathlineto{\pgfqpoint{3.632621in}{2.304436in}}%
\pgfpathlineto{\pgfqpoint{3.633310in}{2.304523in}}%
\pgfpathlineto{\pgfqpoint{3.632704in}{2.304187in}}%
\pgfpathlineto{\pgfqpoint{3.633721in}{2.304264in}}%
\pgfpathlineto{\pgfqpoint{3.633961in}{2.304089in}}%
\pgfpathlineto{\pgfqpoint{3.634679in}{2.304668in}}%
\pgfpathlineto{\pgfqpoint{3.634829in}{2.304334in}}%
\pgfpathlineto{\pgfqpoint{3.634833in}{2.304390in}}%
\pgfpathlineto{\pgfqpoint{3.635177in}{2.303949in}}%
\pgfpathlineto{\pgfqpoint{3.635936in}{2.304128in}}%
\pgfpathlineto{\pgfqpoint{3.636467in}{2.304237in}}%
\pgfpathlineto{\pgfqpoint{3.636740in}{2.303891in}}%
\pgfpathlineto{\pgfqpoint{3.637033in}{2.303931in}}%
\pgfpathlineto{\pgfqpoint{3.637914in}{2.303753in}}%
\pgfpathlineto{\pgfqpoint{3.637497in}{2.304334in}}%
\pgfpathlineto{\pgfqpoint{3.638149in}{2.303831in}}%
\pgfpathlineto{\pgfqpoint{3.639436in}{2.303290in}}%
\pgfpathlineto{\pgfqpoint{3.638194in}{2.304025in}}%
\pgfpathlineto{\pgfqpoint{3.639466in}{2.303457in}}%
\pgfpathlineto{\pgfqpoint{3.640637in}{2.304225in}}%
\pgfpathlineto{\pgfqpoint{3.639476in}{2.303412in}}%
\pgfpathlineto{\pgfqpoint{3.640638in}{2.304217in}}%
\pgfpathlineto{\pgfqpoint{3.641826in}{2.303341in}}%
\pgfpathlineto{\pgfqpoint{3.641057in}{2.304400in}}%
\pgfpathlineto{\pgfqpoint{3.641833in}{2.303367in}}%
\pgfpathlineto{\pgfqpoint{3.641848in}{2.303270in}}%
\pgfpathlineto{\pgfqpoint{3.642485in}{2.303916in}}%
\pgfpathlineto{\pgfqpoint{3.642855in}{2.303813in}}%
\pgfpathlineto{\pgfqpoint{3.643853in}{2.304084in}}%
\pgfpathlineto{\pgfqpoint{3.642910in}{2.303623in}}%
\pgfpathlineto{\pgfqpoint{3.643966in}{2.303798in}}%
\pgfpathlineto{\pgfqpoint{3.643972in}{2.303755in}}%
\pgfpathlineto{\pgfqpoint{3.644196in}{2.304182in}}%
\pgfpathlineto{\pgfqpoint{3.645055in}{2.304067in}}%
\pgfpathlineto{\pgfqpoint{3.645091in}{2.304153in}}%
\pgfpathlineto{\pgfqpoint{3.645837in}{2.303163in}}%
\pgfpathlineto{\pgfqpoint{3.646031in}{2.303275in}}%
\pgfpathlineto{\pgfqpoint{3.646454in}{2.302810in}}%
\pgfpathlineto{\pgfqpoint{3.646069in}{2.303374in}}%
\pgfpathlineto{\pgfqpoint{3.647153in}{2.303307in}}%
\pgfpathlineto{\pgfqpoint{3.648069in}{2.303513in}}%
\pgfpathlineto{\pgfqpoint{3.647815in}{2.303178in}}%
\pgfpathlineto{\pgfqpoint{3.648260in}{2.303243in}}%
\pgfpathlineto{\pgfqpoint{3.648508in}{2.303011in}}%
\pgfpathlineto{\pgfqpoint{3.649349in}{2.303498in}}%
\pgfpathlineto{\pgfqpoint{3.649366in}{2.303490in}}%
\pgfpathlineto{\pgfqpoint{3.650310in}{2.303668in}}%
\pgfpathlineto{\pgfqpoint{3.649573in}{2.303205in}}%
\pgfpathlineto{\pgfqpoint{3.650468in}{2.303340in}}%
\pgfpathlineto{\pgfqpoint{3.651519in}{2.303266in}}%
\pgfpathlineto{\pgfqpoint{3.651131in}{2.303825in}}%
\pgfpathlineto{\pgfqpoint{3.651578in}{2.303375in}}%
\pgfpathlineto{\pgfqpoint{3.652801in}{2.303136in}}%
\pgfpathlineto{\pgfqpoint{3.651584in}{2.303435in}}%
\pgfpathlineto{\pgfqpoint{3.652893in}{2.303242in}}%
\pgfpathlineto{\pgfqpoint{3.655831in}{2.303077in}}%
\pgfpathlineto{\pgfqpoint{3.652895in}{2.303251in}}%
\pgfpathlineto{\pgfqpoint{3.655844in}{2.303145in}}%
\pgfpathlineto{\pgfqpoint{3.656808in}{2.303366in}}%
\pgfpathlineto{\pgfqpoint{3.655992in}{2.303026in}}%
\pgfpathlineto{\pgfqpoint{3.656971in}{2.303227in}}%
\pgfpathlineto{\pgfqpoint{3.657185in}{2.303152in}}%
\pgfpathlineto{\pgfqpoint{3.657860in}{2.303540in}}%
\pgfpathlineto{\pgfqpoint{3.658073in}{2.303495in}}%
\pgfpathlineto{\pgfqpoint{3.658964in}{2.303187in}}%
\pgfpathlineto{\pgfqpoint{3.658405in}{2.303683in}}%
\pgfpathlineto{\pgfqpoint{3.659202in}{2.303345in}}%
\pgfpathlineto{\pgfqpoint{3.660311in}{2.303763in}}%
\pgfpathlineto{\pgfqpoint{3.659370in}{2.303226in}}%
\pgfpathlineto{\pgfqpoint{3.660343in}{2.303743in}}%
\pgfpathlineto{\pgfqpoint{3.660653in}{2.303451in}}%
\pgfpathlineto{\pgfqpoint{3.660486in}{2.303793in}}%
\pgfpathlineto{\pgfqpoint{3.661461in}{2.303750in}}%
\pgfpathlineto{\pgfqpoint{3.662484in}{2.304187in}}%
\pgfpathlineto{\pgfqpoint{3.661541in}{2.303603in}}%
\pgfpathlineto{\pgfqpoint{3.662631in}{2.303907in}}%
\pgfpathlineto{\pgfqpoint{3.663746in}{2.303740in}}%
\pgfpathlineto{\pgfqpoint{3.663047in}{2.304113in}}%
\pgfpathlineto{\pgfqpoint{3.663747in}{2.303752in}}%
\pgfpathlineto{\pgfqpoint{3.664934in}{2.302798in}}%
\pgfpathlineto{\pgfqpoint{3.663751in}{2.303773in}}%
\pgfpathlineto{\pgfqpoint{3.664944in}{2.302820in}}%
\pgfpathlineto{\pgfqpoint{3.666051in}{2.303440in}}%
\pgfpathlineto{\pgfqpoint{3.666098in}{2.303436in}}%
\pgfpathlineto{\pgfqpoint{3.667008in}{2.303368in}}%
\pgfpathlineto{\pgfqpoint{3.666259in}{2.303457in}}%
\pgfpathlineto{\pgfqpoint{3.667209in}{2.303398in}}%
\pgfpathlineto{\pgfqpoint{3.669086in}{2.303390in}}%
\pgfpathlineto{\pgfqpoint{3.667215in}{2.303431in}}%
\pgfpathlineto{\pgfqpoint{3.669090in}{2.303421in}}%
\pgfpathlineto{\pgfqpoint{3.670202in}{2.303449in}}%
\pgfpathlineto{\pgfqpoint{3.669525in}{2.303345in}}%
\pgfpathlineto{\pgfqpoint{3.670205in}{2.303432in}}%
\pgfpathlineto{\pgfqpoint{3.671255in}{2.303358in}}%
\pgfpathlineto{\pgfqpoint{3.670209in}{2.303446in}}%
\pgfpathlineto{\pgfqpoint{3.671318in}{2.303423in}}%
\pgfpathlineto{\pgfqpoint{3.672422in}{2.303456in}}%
\pgfpathlineto{\pgfqpoint{3.671390in}{2.303374in}}%
\pgfpathlineto{\pgfqpoint{3.672431in}{2.303372in}}%
\pgfpathlineto{\pgfqpoint{3.673609in}{2.303349in}}%
\pgfpathlineto{\pgfqpoint{3.672440in}{2.303444in}}%
\pgfpathlineto{\pgfqpoint{3.673640in}{2.303410in}}%
\pgfpathlineto{\pgfqpoint{3.674479in}{2.303455in}}%
\pgfpathlineto{\pgfqpoint{3.673723in}{2.303359in}}%
\pgfpathlineto{\pgfqpoint{3.674750in}{2.303398in}}%
\pgfpathlineto{\pgfqpoint{3.675779in}{2.303364in}}%
\pgfpathlineto{\pgfqpoint{3.674815in}{2.303446in}}%
\pgfpathlineto{\pgfqpoint{3.675861in}{2.303422in}}%
\pgfpathlineto{\pgfqpoint{3.676652in}{2.303453in}}%
\pgfpathlineto{\pgfqpoint{3.676024in}{2.303350in}}%
\pgfpathlineto{\pgfqpoint{3.676972in}{2.303393in}}%
\pgfpathlineto{\pgfqpoint{3.678066in}{2.303445in}}%
\pgfpathlineto{\pgfqpoint{3.677247in}{2.303320in}}%
\pgfpathlineto{\pgfqpoint{3.678087in}{2.303430in}}%
\pgfpathlineto{\pgfqpoint{3.679165in}{2.303377in}}%
\pgfpathlineto{\pgfqpoint{3.678455in}{2.303451in}}%
\pgfpathlineto{\pgfqpoint{3.679199in}{2.303399in}}%
\pgfpathlineto{\pgfqpoint{3.680259in}{2.303455in}}%
\pgfpathlineto{\pgfqpoint{3.679363in}{2.303367in}}%
\pgfpathlineto{\pgfqpoint{3.680327in}{2.303428in}}%
\pgfpathlineto{\pgfqpoint{3.681148in}{2.303352in}}%
\pgfpathlineto{\pgfqpoint{3.680348in}{2.303442in}}%
\pgfpathlineto{\pgfqpoint{3.681439in}{2.303414in}}%
\pgfpathlineto{\pgfqpoint{3.682537in}{2.303444in}}%
\pgfpathlineto{\pgfqpoint{3.681472in}{2.303378in}}%
\pgfpathlineto{\pgfqpoint{3.682556in}{2.303426in}}%
\pgfpathlineto{\pgfqpoint{3.683052in}{2.303343in}}%
\pgfpathlineto{\pgfqpoint{3.682618in}{2.303449in}}%
\pgfpathlineto{\pgfqpoint{3.683669in}{2.303402in}}%
\pgfpathlineto{\pgfqpoint{3.684655in}{2.303461in}}%
\pgfpathlineto{\pgfqpoint{3.684769in}{2.303252in}}%
\pgfpathlineto{\pgfqpoint{3.684774in}{2.303299in}}%
\pgfpathlineto{\pgfqpoint{3.685901in}{2.302524in}}%
\pgfpathlineto{\pgfqpoint{3.684776in}{2.303312in}}%
\pgfpathlineto{\pgfqpoint{3.685943in}{2.302630in}}%
\pgfpathlineto{\pgfqpoint{3.687065in}{2.303082in}}%
\pgfpathlineto{\pgfqpoint{3.686356in}{2.302449in}}%
\pgfpathlineto{\pgfqpoint{3.687077in}{2.303001in}}%
\pgfpathlineto{\pgfqpoint{3.687178in}{2.302844in}}%
\pgfpathlineto{\pgfqpoint{3.688142in}{2.303598in}}%
\pgfpathlineto{\pgfqpoint{3.688160in}{2.303497in}}%
\pgfpathlineto{\pgfqpoint{3.688969in}{2.304179in}}%
\pgfpathlineto{\pgfqpoint{3.688215in}{2.303279in}}%
\pgfpathlineto{\pgfqpoint{3.689285in}{2.303953in}}%
\pgfpathlineto{\pgfqpoint{3.690283in}{2.304005in}}%
\pgfpathlineto{\pgfqpoint{3.689883in}{2.304389in}}%
\pgfpathlineto{\pgfqpoint{3.690361in}{2.304216in}}%
\pgfpathlineto{\pgfqpoint{3.690506in}{2.304356in}}%
\pgfpathlineto{\pgfqpoint{3.691407in}{2.303838in}}%
\pgfpathlineto{\pgfqpoint{3.691463in}{2.303965in}}%
\pgfpathlineto{\pgfqpoint{3.692511in}{2.303703in}}%
\pgfpathlineto{\pgfqpoint{3.691987in}{2.304226in}}%
\pgfpathlineto{\pgfqpoint{3.692582in}{2.303844in}}%
\pgfpathlineto{\pgfqpoint{3.692601in}{2.303944in}}%
\pgfpathlineto{\pgfqpoint{3.693651in}{2.303496in}}%
\pgfpathlineto{\pgfqpoint{3.693678in}{2.303519in}}%
\pgfpathlineto{\pgfqpoint{3.694112in}{2.303307in}}%
\pgfpathlineto{\pgfqpoint{3.694674in}{2.303964in}}%
\pgfpathlineto{\pgfqpoint{3.694765in}{2.303804in}}%
\pgfpathlineto{\pgfqpoint{3.696002in}{2.303246in}}%
\pgfpathlineto{\pgfqpoint{3.694991in}{2.303915in}}%
\pgfpathlineto{\pgfqpoint{3.696017in}{2.303319in}}%
\pgfpathlineto{\pgfqpoint{3.700882in}{2.302787in}}%
\pgfpathlineto{\pgfqpoint{3.700887in}{2.302742in}}%
\pgfpathlineto{\pgfqpoint{3.701848in}{2.302672in}}%
\pgfpathlineto{\pgfqpoint{3.701215in}{2.302910in}}%
\pgfpathlineto{\pgfqpoint{3.701995in}{2.302823in}}%
\pgfpathlineto{\pgfqpoint{3.704745in}{2.302129in}}%
\pgfpathlineto{\pgfqpoint{3.702001in}{2.302855in}}%
\pgfpathlineto{\pgfqpoint{3.704889in}{2.302509in}}%
\pgfpathlineto{\pgfqpoint{3.705098in}{2.302813in}}%
\pgfpathlineto{\pgfqpoint{3.704941in}{2.302492in}}%
\pgfpathlineto{\pgfqpoint{3.705999in}{2.302493in}}%
\pgfpathlineto{\pgfqpoint{3.706987in}{2.301793in}}%
\pgfpathlineto{\pgfqpoint{3.706100in}{2.302580in}}%
\pgfpathlineto{\pgfqpoint{3.707133in}{2.301970in}}%
\pgfpathlineto{\pgfqpoint{3.707443in}{2.301833in}}%
\pgfpathlineto{\pgfqpoint{3.708207in}{2.302197in}}%
\pgfpathlineto{\pgfqpoint{3.708237in}{2.302096in}}%
\pgfpathlineto{\pgfqpoint{3.710266in}{2.302409in}}%
\pgfpathlineto{\pgfqpoint{3.708249in}{2.302045in}}%
\pgfpathlineto{\pgfqpoint{3.710292in}{2.302261in}}%
\pgfpathlineto{\pgfqpoint{3.710372in}{2.302194in}}%
\pgfpathlineto{\pgfqpoint{3.711281in}{2.302999in}}%
\pgfpathlineto{\pgfqpoint{3.711352in}{2.302870in}}%
\pgfpathlineto{\pgfqpoint{3.711636in}{2.303059in}}%
\pgfpathlineto{\pgfqpoint{3.712361in}{2.302612in}}%
\pgfpathlineto{\pgfqpoint{3.712456in}{2.302700in}}%
\pgfpathlineto{\pgfqpoint{3.712752in}{2.302510in}}%
\pgfpathlineto{\pgfqpoint{3.713481in}{2.303250in}}%
\pgfpathlineto{\pgfqpoint{3.713482in}{2.303245in}}%
\pgfpathlineto{\pgfqpoint{3.714059in}{2.303395in}}%
\pgfpathlineto{\pgfqpoint{3.714517in}{2.303050in}}%
\pgfpathlineto{\pgfqpoint{3.714591in}{2.303171in}}%
\pgfpathlineto{\pgfqpoint{3.715517in}{2.302384in}}%
\pgfpathlineto{\pgfqpoint{3.714862in}{2.303280in}}%
\pgfpathlineto{\pgfqpoint{3.715726in}{2.302570in}}%
\pgfpathlineto{\pgfqpoint{3.716548in}{2.303216in}}%
\pgfpathlineto{\pgfqpoint{3.715786in}{2.302411in}}%
\pgfpathlineto{\pgfqpoint{3.716869in}{2.303130in}}%
\pgfpathlineto{\pgfqpoint{3.717014in}{2.303070in}}%
\pgfpathlineto{\pgfqpoint{3.717747in}{2.303771in}}%
\pgfpathlineto{\pgfqpoint{3.717944in}{2.303543in}}%
\pgfpathlineto{\pgfqpoint{3.718945in}{2.303742in}}%
\pgfpathlineto{\pgfqpoint{3.718179in}{2.303140in}}%
\pgfpathlineto{\pgfqpoint{3.719068in}{2.303600in}}%
\pgfpathlineto{\pgfqpoint{3.719077in}{2.303578in}}%
\pgfpathlineto{\pgfqpoint{3.719549in}{2.303953in}}%
\pgfpathlineto{\pgfqpoint{3.720167in}{2.303794in}}%
\pgfpathlineto{\pgfqpoint{3.720413in}{2.303909in}}%
\pgfpathlineto{\pgfqpoint{3.721218in}{2.302796in}}%
\pgfpathlineto{\pgfqpoint{3.721955in}{2.302359in}}%
\pgfpathlineto{\pgfqpoint{3.722334in}{2.302731in}}%
\pgfpathlineto{\pgfqpoint{3.723479in}{2.303181in}}%
\pgfpathlineto{\pgfqpoint{3.722516in}{2.302690in}}%
\pgfpathlineto{\pgfqpoint{3.723492in}{2.303171in}}%
\pgfpathlineto{\pgfqpoint{3.724489in}{2.302938in}}%
\pgfpathlineto{\pgfqpoint{3.723794in}{2.303486in}}%
\pgfpathlineto{\pgfqpoint{3.724608in}{2.303053in}}%
\pgfpathlineto{\pgfqpoint{3.724898in}{2.303404in}}%
\pgfpathlineto{\pgfqpoint{3.725651in}{2.302608in}}%
\pgfpathlineto{\pgfqpoint{3.725691in}{2.302673in}}%
\pgfpathlineto{\pgfqpoint{3.725695in}{2.302650in}}%
\pgfpathlineto{\pgfqpoint{3.726762in}{2.303170in}}%
\pgfpathlineto{\pgfqpoint{3.726781in}{2.303157in}}%
\pgfpathlineto{\pgfqpoint{3.728902in}{2.302901in}}%
\pgfpathlineto{\pgfqpoint{3.726833in}{2.302921in}}%
\pgfpathlineto{\pgfqpoint{3.728913in}{2.302842in}}%
\pgfpathlineto{\pgfqpoint{3.729809in}{2.302304in}}%
\pgfpathlineto{\pgfqpoint{3.729122in}{2.303093in}}%
\pgfpathlineto{\pgfqpoint{3.730037in}{2.302456in}}%
\pgfpathlineto{\pgfqpoint{3.731143in}{2.302800in}}%
\pgfpathlineto{\pgfqpoint{3.730295in}{2.302171in}}%
\pgfpathlineto{\pgfqpoint{3.731200in}{2.302763in}}%
\pgfpathlineto{\pgfqpoint{3.732626in}{2.302191in}}%
\pgfpathlineto{\pgfqpoint{3.731204in}{2.302794in}}%
\pgfpathlineto{\pgfqpoint{3.732686in}{2.302372in}}%
\pgfpathlineto{\pgfqpoint{3.733326in}{2.302808in}}%
\pgfpathlineto{\pgfqpoint{3.732808in}{2.302264in}}%
\pgfpathlineto{\pgfqpoint{3.733831in}{2.302646in}}%
\pgfpathlineto{\pgfqpoint{3.733885in}{2.302584in}}%
\pgfpathlineto{\pgfqpoint{3.734753in}{2.303789in}}%
\pgfpathlineto{\pgfqpoint{3.734865in}{2.303533in}}%
\pgfpathlineto{\pgfqpoint{3.735458in}{2.303878in}}%
\pgfpathlineto{\pgfqpoint{3.734939in}{2.303408in}}%
\pgfpathlineto{\pgfqpoint{3.735983in}{2.303558in}}%
\pgfpathlineto{\pgfqpoint{3.736037in}{2.303319in}}%
\pgfpathlineto{\pgfqpoint{3.736437in}{2.303751in}}%
\pgfpathlineto{\pgfqpoint{3.737101in}{2.303522in}}%
\pgfpathlineto{\pgfqpoint{3.738492in}{2.303686in}}%
\pgfpathlineto{\pgfqpoint{3.737311in}{2.303148in}}%
\pgfpathlineto{\pgfqpoint{3.738540in}{2.303423in}}%
\pgfpathlineto{\pgfqpoint{3.739585in}{2.303272in}}%
\pgfpathlineto{\pgfqpoint{3.738671in}{2.303604in}}%
\pgfpathlineto{\pgfqpoint{3.739664in}{2.303345in}}%
\pgfpathlineto{\pgfqpoint{3.740339in}{2.303956in}}%
\pgfpathlineto{\pgfqpoint{3.739799in}{2.303148in}}%
\pgfpathlineto{\pgfqpoint{3.740856in}{2.303418in}}%
\pgfpathlineto{\pgfqpoint{3.743099in}{2.303399in}}%
\pgfpathlineto{\pgfqpoint{3.740858in}{2.303426in}}%
\pgfpathlineto{\pgfqpoint{3.743105in}{2.303420in}}%
\pgfpathlineto{\pgfqpoint{3.744108in}{2.303447in}}%
\pgfpathlineto{\pgfqpoint{3.743133in}{2.303372in}}%
\pgfpathlineto{\pgfqpoint{3.744216in}{2.303436in}}%
\pgfpathlineto{\pgfqpoint{3.745172in}{2.303364in}}%
\pgfpathlineto{\pgfqpoint{3.744253in}{2.303445in}}%
\pgfpathlineto{\pgfqpoint{3.745328in}{2.303412in}}%
\pgfpathlineto{\pgfqpoint{3.746822in}{2.303386in}}%
\pgfpathlineto{\pgfqpoint{3.745338in}{2.303464in}}%
\pgfpathlineto{\pgfqpoint{3.746832in}{2.303420in}}%
\pgfpathlineto{\pgfqpoint{3.747723in}{2.303451in}}%
\pgfpathlineto{\pgfqpoint{3.747567in}{2.303323in}}%
\pgfpathlineto{\pgfqpoint{3.747943in}{2.303421in}}%
\pgfpathlineto{\pgfqpoint{3.749056in}{2.303368in}}%
\pgfpathlineto{\pgfqpoint{3.747954in}{2.303436in}}%
\pgfpathlineto{\pgfqpoint{3.749088in}{2.303396in}}%
\pgfpathlineto{\pgfqpoint{3.750114in}{2.303447in}}%
\pgfpathlineto{\pgfqpoint{3.749308in}{2.303357in}}%
\pgfpathlineto{\pgfqpoint{3.750204in}{2.303425in}}%
\pgfpathlineto{\pgfqpoint{3.751121in}{2.303366in}}%
\pgfpathlineto{\pgfqpoint{3.750212in}{2.303441in}}%
\pgfpathlineto{\pgfqpoint{3.751318in}{2.303385in}}%
\pgfpathlineto{\pgfqpoint{3.752263in}{2.303450in}}%
\pgfpathlineto{\pgfqpoint{3.752038in}{2.303346in}}%
\pgfpathlineto{\pgfqpoint{3.752430in}{2.303423in}}%
\pgfpathlineto{\pgfqpoint{3.753487in}{2.303359in}}%
\pgfpathlineto{\pgfqpoint{3.752443in}{2.303448in}}%
\pgfpathlineto{\pgfqpoint{3.753543in}{2.303410in}}%
\pgfpathlineto{\pgfqpoint{3.754637in}{2.303455in}}%
\pgfpathlineto{\pgfqpoint{3.753605in}{2.303346in}}%
\pgfpathlineto{\pgfqpoint{3.754655in}{2.303436in}}%
\pgfpathlineto{\pgfqpoint{3.755637in}{2.303324in}}%
\pgfpathlineto{\pgfqpoint{3.754760in}{2.303456in}}%
\pgfpathlineto{\pgfqpoint{3.755767in}{2.303413in}}%
\pgfpathlineto{\pgfqpoint{3.755914in}{2.303356in}}%
\pgfpathlineto{\pgfqpoint{3.756886in}{2.303437in}}%
\pgfpathlineto{\pgfqpoint{3.757904in}{2.303360in}}%
\pgfpathlineto{\pgfqpoint{3.756947in}{2.303442in}}%
\pgfpathlineto{\pgfqpoint{3.757997in}{2.303417in}}%
\pgfpathlineto{\pgfqpoint{3.759065in}{2.303450in}}%
\pgfpathlineto{\pgfqpoint{3.758071in}{2.303376in}}%
\pgfpathlineto{\pgfqpoint{3.759109in}{2.303420in}}%
\pgfpathlineto{\pgfqpoint{3.759869in}{2.303222in}}%
\pgfpathlineto{\pgfqpoint{3.759115in}{2.303451in}}%
\pgfpathlineto{\pgfqpoint{3.760226in}{2.303453in}}%
\pgfpathlineto{\pgfqpoint{3.761154in}{2.303475in}}%
\pgfpathlineto{\pgfqpoint{3.760551in}{2.303048in}}%
\pgfpathlineto{\pgfqpoint{3.761332in}{2.303288in}}%
\pgfpathlineto{\pgfqpoint{3.761620in}{2.303070in}}%
\pgfpathlineto{\pgfqpoint{3.762177in}{2.303810in}}%
\pgfpathlineto{\pgfqpoint{3.762425in}{2.303564in}}%
\pgfpathlineto{\pgfqpoint{3.763389in}{2.303916in}}%
\pgfpathlineto{\pgfqpoint{3.762614in}{2.303414in}}%
\pgfpathlineto{\pgfqpoint{3.763559in}{2.303760in}}%
\pgfpathlineto{\pgfqpoint{3.764331in}{2.304043in}}%
\pgfpathlineto{\pgfqpoint{3.763848in}{2.303467in}}%
\pgfpathlineto{\pgfqpoint{3.764675in}{2.303673in}}%
\pgfpathlineto{\pgfqpoint{3.764677in}{2.303646in}}%
\pgfpathlineto{\pgfqpoint{3.765560in}{2.304140in}}%
\pgfpathlineto{\pgfqpoint{3.765778in}{2.303925in}}%
\pgfpathlineto{\pgfqpoint{3.766478in}{2.303603in}}%
\pgfpathlineto{\pgfqpoint{3.765801in}{2.303942in}}%
\pgfpathlineto{\pgfqpoint{3.766919in}{2.303899in}}%
\pgfpathlineto{\pgfqpoint{3.767972in}{2.304178in}}%
\pgfpathlineto{\pgfqpoint{3.767450in}{2.303430in}}%
\pgfpathlineto{\pgfqpoint{3.768037in}{2.303998in}}%
\pgfpathlineto{\pgfqpoint{3.769082in}{2.303942in}}%
\pgfpathlineto{\pgfqpoint{3.768417in}{2.304472in}}%
\pgfpathlineto{\pgfqpoint{3.769146in}{2.304057in}}%
\pgfpathlineto{\pgfqpoint{3.769151in}{2.304094in}}%
\pgfpathlineto{\pgfqpoint{3.770206in}{2.303133in}}%
\pgfpathlineto{\pgfqpoint{3.770452in}{2.303357in}}%
\pgfpathlineto{\pgfqpoint{3.770717in}{2.302959in}}%
\pgfpathlineto{\pgfqpoint{3.771315in}{2.303060in}}%
\pgfpathlineto{\pgfqpoint{3.772493in}{2.303314in}}%
\pgfpathlineto{\pgfqpoint{3.771345in}{2.303016in}}%
\pgfpathlineto{\pgfqpoint{3.772619in}{2.303113in}}%
\pgfpathlineto{\pgfqpoint{3.773725in}{2.302452in}}%
\pgfpathlineto{\pgfqpoint{3.772681in}{2.303171in}}%
\pgfpathlineto{\pgfqpoint{3.773759in}{2.302642in}}%
\pgfpathlineto{\pgfqpoint{3.773786in}{2.302524in}}%
\pgfpathlineto{\pgfqpoint{3.774699in}{2.303358in}}%
\pgfpathlineto{\pgfqpoint{3.774728in}{2.303297in}}%
\pgfpathlineto{\pgfqpoint{3.775539in}{2.303775in}}%
\pgfpathlineto{\pgfqpoint{3.774816in}{2.303199in}}%
\pgfpathlineto{\pgfqpoint{3.775859in}{2.303707in}}%
\pgfpathlineto{\pgfqpoint{3.776923in}{2.303383in}}%
\pgfpathlineto{\pgfqpoint{3.776112in}{2.303855in}}%
\pgfpathlineto{\pgfqpoint{3.776996in}{2.303543in}}%
\pgfpathlineto{\pgfqpoint{3.776998in}{2.303604in}}%
\pgfpathlineto{\pgfqpoint{3.777605in}{2.302755in}}%
\pgfpathlineto{\pgfqpoint{3.778070in}{2.302989in}}%
\pgfpathlineto{\pgfqpoint{3.779069in}{2.302842in}}%
\pgfpathlineto{\pgfqpoint{3.778581in}{2.303316in}}%
\pgfpathlineto{\pgfqpoint{3.779189in}{2.302914in}}%
\pgfpathlineto{\pgfqpoint{3.779757in}{2.303702in}}%
\pgfpathlineto{\pgfqpoint{3.780349in}{2.303451in}}%
\pgfpathlineto{\pgfqpoint{3.781227in}{2.303752in}}%
\pgfpathlineto{\pgfqpoint{3.780360in}{2.303508in}}%
\pgfpathlineto{\pgfqpoint{3.781398in}{2.303919in}}%
\pgfpathlineto{\pgfqpoint{3.782536in}{2.304341in}}%
\pgfpathlineto{\pgfqpoint{3.781559in}{2.303825in}}%
\pgfpathlineto{\pgfqpoint{3.782585in}{2.304251in}}%
\pgfpathlineto{\pgfqpoint{3.782602in}{2.304203in}}%
\pgfpathlineto{\pgfqpoint{3.783247in}{2.304716in}}%
\pgfpathlineto{\pgfqpoint{3.783679in}{2.304523in}}%
\pgfpathlineto{\pgfqpoint{3.784272in}{2.304925in}}%
\pgfpathlineto{\pgfqpoint{3.784548in}{2.304344in}}%
\pgfpathlineto{\pgfqpoint{3.784808in}{2.304682in}}%
\pgfpathlineto{\pgfqpoint{3.784890in}{2.304571in}}%
\pgfpathlineto{\pgfqpoint{3.785696in}{2.305065in}}%
\pgfpathlineto{\pgfqpoint{3.785912in}{2.304936in}}%
\pgfpathlineto{\pgfqpoint{3.786876in}{2.305362in}}%
\pgfpathlineto{\pgfqpoint{3.785954in}{2.304767in}}%
\pgfpathlineto{\pgfqpoint{3.787034in}{2.305031in}}%
\pgfpathlineto{\pgfqpoint{3.787249in}{2.304732in}}%
\pgfpathlineto{\pgfqpoint{3.787703in}{2.305228in}}%
\pgfpathlineto{\pgfqpoint{3.788159in}{2.304841in}}%
\pgfpathlineto{\pgfqpoint{3.789213in}{2.304961in}}%
\pgfpathlineto{\pgfqpoint{3.788899in}{2.304653in}}%
\pgfpathlineto{\pgfqpoint{3.789275in}{2.304890in}}%
\pgfpathlineto{\pgfqpoint{3.790372in}{2.304979in}}%
\pgfpathlineto{\pgfqpoint{3.790224in}{2.305261in}}%
\pgfpathlineto{\pgfqpoint{3.790381in}{2.305000in}}%
\pgfpathlineto{\pgfqpoint{3.790603in}{2.304813in}}%
\pgfpathlineto{\pgfqpoint{3.791376in}{2.305670in}}%
\pgfpathlineto{\pgfqpoint{3.791457in}{2.305316in}}%
\pgfpathlineto{\pgfqpoint{3.791476in}{2.305347in}}%
\pgfpathlineto{\pgfqpoint{3.792475in}{2.304818in}}%
\pgfpathlineto{\pgfqpoint{3.792543in}{2.305025in}}%
\pgfpathlineto{\pgfqpoint{3.793775in}{2.304338in}}%
\pgfpathlineto{\pgfqpoint{3.792655in}{2.305068in}}%
\pgfpathlineto{\pgfqpoint{3.793776in}{2.304353in}}%
\pgfpathlineto{\pgfqpoint{3.794011in}{2.304497in}}%
\pgfpathlineto{\pgfqpoint{3.794808in}{2.303798in}}%
\pgfpathlineto{\pgfqpoint{3.794860in}{2.303886in}}%
\pgfpathlineto{\pgfqpoint{3.796024in}{2.303026in}}%
\pgfpathlineto{\pgfqpoint{3.795038in}{2.303994in}}%
\pgfpathlineto{\pgfqpoint{3.796042in}{2.303079in}}%
\pgfpathlineto{\pgfqpoint{3.797082in}{2.303223in}}%
\pgfpathlineto{\pgfqpoint{3.796295in}{2.302915in}}%
\pgfpathlineto{\pgfqpoint{3.797232in}{2.302999in}}%
\pgfpathlineto{\pgfqpoint{3.797539in}{2.302858in}}%
\pgfpathlineto{\pgfqpoint{3.797872in}{2.303443in}}%
\pgfpathlineto{\pgfqpoint{3.798336in}{2.303151in}}%
\pgfpathlineto{\pgfqpoint{3.799137in}{2.303203in}}%
\pgfpathlineto{\pgfqpoint{3.798404in}{2.302994in}}%
\pgfpathlineto{\pgfqpoint{3.799466in}{2.302898in}}%
\pgfpathlineto{\pgfqpoint{3.800426in}{2.302400in}}%
\pgfpathlineto{\pgfqpoint{3.799885in}{2.303205in}}%
\pgfpathlineto{\pgfqpoint{3.800595in}{2.302643in}}%
\pgfpathlineto{\pgfqpoint{3.801500in}{2.303033in}}%
\pgfpathlineto{\pgfqpoint{3.800684in}{2.302521in}}%
\pgfpathlineto{\pgfqpoint{3.801729in}{2.302959in}}%
\pgfpathlineto{\pgfqpoint{3.802020in}{2.302692in}}%
\pgfpathlineto{\pgfqpoint{3.802679in}{2.303236in}}%
\pgfpathlineto{\pgfqpoint{3.802829in}{2.303103in}}%
\pgfpathlineto{\pgfqpoint{3.803337in}{2.303479in}}%
\pgfpathlineto{\pgfqpoint{3.803911in}{2.302816in}}%
\pgfpathlineto{\pgfqpoint{3.803933in}{2.302853in}}%
\pgfpathlineto{\pgfqpoint{3.805174in}{2.302393in}}%
\pgfpathlineto{\pgfqpoint{3.804149in}{2.302937in}}%
\pgfpathlineto{\pgfqpoint{3.805180in}{2.302432in}}%
\pgfpathlineto{\pgfqpoint{3.805876in}{2.302840in}}%
\pgfpathlineto{\pgfqpoint{3.806314in}{2.302776in}}%
\pgfpathlineto{\pgfqpoint{3.807180in}{2.302177in}}%
\pgfpathlineto{\pgfqpoint{3.807584in}{2.302496in}}%
\pgfpathlineto{\pgfqpoint{3.807599in}{2.302580in}}%
\pgfpathlineto{\pgfqpoint{3.808625in}{2.301527in}}%
\pgfpathlineto{\pgfqpoint{3.812606in}{2.303002in}}%
\pgfpathlineto{\pgfqpoint{3.808689in}{2.301355in}}%
\pgfpathlineto{\pgfqpoint{3.812673in}{2.302820in}}%
\pgfpathlineto{\pgfqpoint{3.812724in}{2.302723in}}%
\pgfpathlineto{\pgfqpoint{3.813491in}{2.303408in}}%
\pgfpathlineto{\pgfqpoint{3.813751in}{2.303275in}}%
\pgfpathlineto{\pgfqpoint{3.814361in}{2.303505in}}%
\pgfpathlineto{\pgfqpoint{3.814059in}{2.303079in}}%
\pgfpathlineto{\pgfqpoint{3.814858in}{2.303185in}}%
\pgfpathlineto{\pgfqpoint{3.815737in}{2.303144in}}%
\pgfpathlineto{\pgfqpoint{3.815401in}{2.303635in}}%
\pgfpathlineto{\pgfqpoint{3.815960in}{2.303418in}}%
\pgfpathlineto{\pgfqpoint{3.817069in}{2.303444in}}%
\pgfpathlineto{\pgfqpoint{3.816021in}{2.303356in}}%
\pgfpathlineto{\pgfqpoint{3.817072in}{2.303421in}}%
\pgfpathlineto{\pgfqpoint{3.818140in}{2.303374in}}%
\pgfpathlineto{\pgfqpoint{3.817081in}{2.303439in}}%
\pgfpathlineto{\pgfqpoint{3.818203in}{2.303418in}}%
\pgfpathlineto{\pgfqpoint{3.819270in}{2.303445in}}%
\pgfpathlineto{\pgfqpoint{3.818235in}{2.303359in}}%
\pgfpathlineto{\pgfqpoint{3.819313in}{2.303425in}}%
\pgfpathlineto{\pgfqpoint{3.820115in}{2.303365in}}%
\pgfpathlineto{\pgfqpoint{3.819653in}{2.303453in}}%
\pgfpathlineto{\pgfqpoint{3.820425in}{2.303416in}}%
\pgfpathlineto{\pgfqpoint{3.821531in}{2.303441in}}%
\pgfpathlineto{\pgfqpoint{3.820477in}{2.303366in}}%
\pgfpathlineto{\pgfqpoint{3.821545in}{2.303434in}}%
\pgfpathlineto{\pgfqpoint{3.822657in}{2.303356in}}%
\pgfpathlineto{\pgfqpoint{3.821554in}{2.303441in}}%
\pgfpathlineto{\pgfqpoint{3.822660in}{2.303372in}}%
\pgfpathlineto{\pgfqpoint{3.823777in}{2.303452in}}%
\pgfpathlineto{\pgfqpoint{3.823777in}{2.303437in}}%
\pgfpathlineto{\pgfqpoint{3.824726in}{2.303337in}}%
\pgfpathlineto{\pgfqpoint{3.824240in}{2.303455in}}%
\pgfpathlineto{\pgfqpoint{3.824889in}{2.303418in}}%
\pgfpathlineto{\pgfqpoint{3.825949in}{2.303446in}}%
\pgfpathlineto{\pgfqpoint{3.824917in}{2.303365in}}%
\pgfpathlineto{\pgfqpoint{3.826002in}{2.303428in}}%
\pgfpathlineto{\pgfqpoint{3.827082in}{2.303361in}}%
\pgfpathlineto{\pgfqpoint{3.826160in}{2.303455in}}%
\pgfpathlineto{\pgfqpoint{3.827113in}{2.303413in}}%
\pgfpathlineto{\pgfqpoint{3.828091in}{2.303449in}}%
\pgfpathlineto{\pgfqpoint{3.827314in}{2.303359in}}%
\pgfpathlineto{\pgfqpoint{3.828225in}{2.303417in}}%
\pgfpathlineto{\pgfqpoint{3.829339in}{2.303368in}}%
\pgfpathlineto{\pgfqpoint{3.828252in}{2.303446in}}%
\pgfpathlineto{\pgfqpoint{3.829339in}{2.303375in}}%
\pgfpathlineto{\pgfqpoint{3.830448in}{2.303437in}}%
\pgfpathlineto{\pgfqpoint{3.829375in}{2.303313in}}%
\pgfpathlineto{\pgfqpoint{3.830453in}{2.303410in}}%
\pgfpathlineto{\pgfqpoint{3.831552in}{2.303383in}}%
\pgfpathlineto{\pgfqpoint{3.830474in}{2.303446in}}%
\pgfpathlineto{\pgfqpoint{3.831593in}{2.303410in}}%
\pgfpathlineto{\pgfqpoint{3.832704in}{2.303439in}}%
\pgfpathlineto{\pgfqpoint{3.831679in}{2.303361in}}%
\pgfpathlineto{\pgfqpoint{3.832714in}{2.303424in}}%
\pgfpathlineto{\pgfqpoint{3.833826in}{2.303357in}}%
\pgfpathlineto{\pgfqpoint{3.832716in}{2.303438in}}%
\pgfpathlineto{\pgfqpoint{3.833829in}{2.303391in}}%
\pgfpathlineto{\pgfqpoint{3.834478in}{2.303460in}}%
\pgfpathlineto{\pgfqpoint{3.834747in}{2.303029in}}%
\pgfpathlineto{\pgfqpoint{3.834935in}{2.303281in}}%
\pgfpathlineto{\pgfqpoint{3.836026in}{2.302737in}}%
\pgfpathlineto{\pgfqpoint{3.834940in}{2.303307in}}%
\pgfpathlineto{\pgfqpoint{3.836251in}{2.303120in}}%
\pgfpathlineto{\pgfqpoint{3.836262in}{2.303186in}}%
\pgfpathlineto{\pgfqpoint{3.836770in}{2.302773in}}%
\pgfpathlineto{\pgfqpoint{3.837359in}{2.303016in}}%
\pgfpathlineto{\pgfqpoint{3.837835in}{2.303387in}}%
\pgfpathlineto{\pgfqpoint{3.838243in}{2.302735in}}%
\pgfpathlineto{\pgfqpoint{3.838480in}{2.303135in}}%
\pgfpathlineto{\pgfqpoint{3.839731in}{2.302484in}}%
\pgfpathlineto{\pgfqpoint{3.838712in}{2.303299in}}%
\pgfpathlineto{\pgfqpoint{3.839763in}{2.302605in}}%
\pgfpathlineto{\pgfqpoint{3.840873in}{2.303102in}}%
\pgfpathlineto{\pgfqpoint{3.839882in}{2.302426in}}%
\pgfpathlineto{\pgfqpoint{3.840912in}{2.303064in}}%
\pgfpathlineto{\pgfqpoint{3.842025in}{2.302949in}}%
\pgfpathlineto{\pgfqpoint{3.841491in}{2.303287in}}%
\pgfpathlineto{\pgfqpoint{3.842029in}{2.302963in}}%
\pgfpathlineto{\pgfqpoint{3.843110in}{2.302962in}}%
\pgfpathlineto{\pgfqpoint{3.842135in}{2.302721in}}%
\pgfpathlineto{\pgfqpoint{3.843165in}{2.302818in}}%
\pgfpathlineto{\pgfqpoint{3.844605in}{2.302133in}}%
\pgfpathlineto{\pgfqpoint{3.843166in}{2.302826in}}%
\pgfpathlineto{\pgfqpoint{3.844636in}{2.302212in}}%
\pgfpathlineto{\pgfqpoint{3.845558in}{2.302525in}}%
\pgfpathlineto{\pgfqpoint{3.845194in}{2.302083in}}%
\pgfpathlineto{\pgfqpoint{3.845754in}{2.302377in}}%
\pgfpathlineto{\pgfqpoint{3.845771in}{2.302501in}}%
\pgfpathlineto{\pgfqpoint{3.846781in}{2.301907in}}%
\pgfpathlineto{\pgfqpoint{3.846838in}{2.302046in}}%
\pgfpathlineto{\pgfqpoint{3.847953in}{2.301809in}}%
\pgfpathlineto{\pgfqpoint{3.847083in}{2.302294in}}%
\pgfpathlineto{\pgfqpoint{3.847955in}{2.301822in}}%
\pgfpathlineto{\pgfqpoint{3.848885in}{2.301895in}}%
\pgfpathlineto{\pgfqpoint{3.848331in}{2.301561in}}%
\pgfpathlineto{\pgfqpoint{3.849070in}{2.301833in}}%
\pgfpathlineto{\pgfqpoint{3.849109in}{2.301759in}}%
\pgfpathlineto{\pgfqpoint{3.850153in}{2.302287in}}%
\pgfpathlineto{\pgfqpoint{3.850155in}{2.302269in}}%
\pgfpathlineto{\pgfqpoint{3.850439in}{2.302082in}}%
\pgfpathlineto{\pgfqpoint{3.850320in}{2.302411in}}%
\pgfpathlineto{\pgfqpoint{3.851266in}{2.302322in}}%
\pgfpathlineto{\pgfqpoint{3.852026in}{2.302469in}}%
\pgfpathlineto{\pgfqpoint{3.851633in}{2.301877in}}%
\pgfpathlineto{\pgfqpoint{3.852365in}{2.302159in}}%
\pgfpathlineto{\pgfqpoint{3.857378in}{2.303252in}}%
\pgfpathlineto{\pgfqpoint{3.858455in}{2.303306in}}%
\pgfpathlineto{\pgfqpoint{3.857762in}{2.303023in}}%
\pgfpathlineto{\pgfqpoint{3.858491in}{2.303246in}}%
\pgfpathlineto{\pgfqpoint{3.859606in}{2.303133in}}%
\pgfpathlineto{\pgfqpoint{3.859032in}{2.303597in}}%
\pgfpathlineto{\pgfqpoint{3.859669in}{2.303244in}}%
\pgfpathlineto{\pgfqpoint{3.860045in}{2.303482in}}%
\pgfpathlineto{\pgfqpoint{3.860311in}{2.303082in}}%
\pgfpathlineto{\pgfqpoint{3.860776in}{2.303146in}}%
\pgfpathlineto{\pgfqpoint{3.860887in}{2.303103in}}%
\pgfpathlineto{\pgfqpoint{3.861576in}{2.303804in}}%
\pgfpathlineto{\pgfqpoint{3.861874in}{2.303621in}}%
\pgfpathlineto{\pgfqpoint{3.864617in}{2.303040in}}%
\pgfpathlineto{\pgfqpoint{3.861880in}{2.303654in}}%
\pgfpathlineto{\pgfqpoint{3.864645in}{2.303186in}}%
\pgfpathlineto{\pgfqpoint{3.865694in}{2.303171in}}%
\pgfpathlineto{\pgfqpoint{3.865320in}{2.302531in}}%
\pgfpathlineto{\pgfqpoint{3.865743in}{2.303007in}}%
\pgfpathlineto{\pgfqpoint{3.866203in}{2.302658in}}%
\pgfpathlineto{\pgfqpoint{3.865751in}{2.303021in}}%
\pgfpathlineto{\pgfqpoint{3.866859in}{2.302869in}}%
\pgfpathlineto{\pgfqpoint{3.866860in}{2.302876in}}%
\pgfpathlineto{\pgfqpoint{3.867444in}{2.301622in}}%
\pgfpathlineto{\pgfqpoint{3.867849in}{2.301670in}}%
\pgfpathlineto{\pgfqpoint{3.867922in}{2.301484in}}%
\pgfpathlineto{\pgfqpoint{3.868812in}{2.302150in}}%
\pgfpathlineto{\pgfqpoint{3.868947in}{2.302053in}}%
\pgfpathlineto{\pgfqpoint{3.869627in}{2.302660in}}%
\pgfpathlineto{\pgfqpoint{3.868986in}{2.301945in}}%
\pgfpathlineto{\pgfqpoint{3.870144in}{2.302062in}}%
\pgfpathlineto{\pgfqpoint{3.871207in}{2.301654in}}%
\pgfpathlineto{\pgfqpoint{3.870158in}{2.302113in}}%
\pgfpathlineto{\pgfqpoint{3.871292in}{2.301734in}}%
\pgfpathlineto{\pgfqpoint{3.872311in}{2.302299in}}%
\pgfpathlineto{\pgfqpoint{3.871348in}{2.301633in}}%
\pgfpathlineto{\pgfqpoint{3.872426in}{2.302123in}}%
\pgfpathlineto{\pgfqpoint{3.872455in}{2.302002in}}%
\pgfpathlineto{\pgfqpoint{3.873502in}{2.302581in}}%
\pgfpathlineto{\pgfqpoint{3.873505in}{2.302567in}}%
\pgfpathlineto{\pgfqpoint{3.873659in}{2.302778in}}%
\pgfpathlineto{\pgfqpoint{3.874444in}{2.302223in}}%
\pgfpathlineto{\pgfqpoint{3.874620in}{2.302616in}}%
\pgfpathlineto{\pgfqpoint{3.875697in}{2.302199in}}%
\pgfpathlineto{\pgfqpoint{3.875743in}{2.302294in}}%
\pgfpathlineto{\pgfqpoint{3.876873in}{2.302979in}}%
\pgfpathlineto{\pgfqpoint{3.875851in}{2.302278in}}%
\pgfpathlineto{\pgfqpoint{3.876888in}{2.302880in}}%
\pgfpathlineto{\pgfqpoint{3.876985in}{2.302884in}}%
\pgfpathlineto{\pgfqpoint{3.877896in}{2.302083in}}%
\pgfpathlineto{\pgfqpoint{3.877945in}{2.302185in}}%
\pgfpathlineto{\pgfqpoint{3.878906in}{2.302959in}}%
\pgfpathlineto{\pgfqpoint{3.877951in}{2.302153in}}%
\pgfpathlineto{\pgfqpoint{3.879170in}{2.302741in}}%
\pgfpathlineto{\pgfqpoint{3.880287in}{2.302038in}}%
\pgfpathlineto{\pgfqpoint{3.879194in}{2.302829in}}%
\pgfpathlineto{\pgfqpoint{3.880322in}{2.302093in}}%
\pgfpathlineto{\pgfqpoint{3.880438in}{2.302358in}}%
\pgfpathlineto{\pgfqpoint{3.881089in}{2.301855in}}%
\pgfpathlineto{\pgfqpoint{3.881443in}{2.302203in}}%
\pgfpathlineto{\pgfqpoint{3.882221in}{2.302325in}}%
\pgfpathlineto{\pgfqpoint{3.881454in}{2.302282in}}%
\pgfpathlineto{\pgfqpoint{3.882459in}{2.302705in}}%
\pgfpathlineto{\pgfqpoint{3.882465in}{2.302735in}}%
\pgfpathlineto{\pgfqpoint{3.883500in}{2.302162in}}%
\pgfpathlineto{\pgfqpoint{3.883540in}{2.302264in}}%
\pgfpathlineto{\pgfqpoint{3.883990in}{2.302088in}}%
\pgfpathlineto{\pgfqpoint{3.884611in}{2.302652in}}%
\pgfpathlineto{\pgfqpoint{3.884631in}{2.302639in}}%
\pgfpathlineto{\pgfqpoint{3.885842in}{2.303139in}}%
\pgfpathlineto{\pgfqpoint{3.884784in}{2.302462in}}%
\pgfpathlineto{\pgfqpoint{3.885885in}{2.303006in}}%
\pgfpathlineto{\pgfqpoint{3.885905in}{2.302962in}}%
\pgfpathlineto{\pgfqpoint{3.886914in}{2.303800in}}%
\pgfpathlineto{\pgfqpoint{3.886954in}{2.303693in}}%
\pgfpathlineto{\pgfqpoint{3.888138in}{2.303051in}}%
\pgfpathlineto{\pgfqpoint{3.887051in}{2.303771in}}%
\pgfpathlineto{\pgfqpoint{3.888182in}{2.303077in}}%
\pgfpathlineto{\pgfqpoint{3.888571in}{2.303510in}}%
\pgfpathlineto{\pgfqpoint{3.889149in}{2.302997in}}%
\pgfpathlineto{\pgfqpoint{3.889296in}{2.303117in}}%
\pgfpathlineto{\pgfqpoint{3.890154in}{2.303011in}}%
\pgfpathlineto{\pgfqpoint{3.889830in}{2.303475in}}%
\pgfpathlineto{\pgfqpoint{3.890408in}{2.303153in}}%
\pgfpathlineto{\pgfqpoint{3.891512in}{2.303444in}}%
\pgfpathlineto{\pgfqpoint{3.890431in}{2.303100in}}%
\pgfpathlineto{\pgfqpoint{3.891530in}{2.303406in}}%
\pgfpathlineto{\pgfqpoint{3.892615in}{2.303367in}}%
\pgfpathlineto{\pgfqpoint{3.891534in}{2.303441in}}%
\pgfpathlineto{\pgfqpoint{3.892646in}{2.303423in}}%
\pgfpathlineto{\pgfqpoint{3.893725in}{2.303447in}}%
\pgfpathlineto{\pgfqpoint{3.892840in}{2.303349in}}%
\pgfpathlineto{\pgfqpoint{3.893757in}{2.303417in}}%
\pgfpathlineto{\pgfqpoint{3.894704in}{2.303365in}}%
\pgfpathlineto{\pgfqpoint{3.893894in}{2.303451in}}%
\pgfpathlineto{\pgfqpoint{3.894869in}{2.303408in}}%
\pgfpathlineto{\pgfqpoint{3.894959in}{2.303359in}}%
\pgfpathlineto{\pgfqpoint{3.895989in}{2.303440in}}%
\pgfpathlineto{\pgfqpoint{3.897003in}{2.303359in}}%
\pgfpathlineto{\pgfqpoint{3.895992in}{2.303443in}}%
\pgfpathlineto{\pgfqpoint{3.897100in}{2.303400in}}%
\pgfpathlineto{\pgfqpoint{3.898210in}{2.303442in}}%
\pgfpathlineto{\pgfqpoint{3.897126in}{2.303360in}}%
\pgfpathlineto{\pgfqpoint{3.898212in}{2.303421in}}%
\pgfpathlineto{\pgfqpoint{3.899329in}{2.303440in}}%
\pgfpathlineto{\pgfqpoint{3.898269in}{2.303379in}}%
\pgfpathlineto{\pgfqpoint{3.899342in}{2.303427in}}%
\pgfpathlineto{\pgfqpoint{3.900077in}{2.303367in}}%
\pgfpathlineto{\pgfqpoint{3.899386in}{2.303448in}}%
\pgfpathlineto{\pgfqpoint{3.900453in}{2.303418in}}%
\pgfpathlineto{\pgfqpoint{3.901536in}{2.303442in}}%
\pgfpathlineto{\pgfqpoint{3.900990in}{2.303339in}}%
\pgfpathlineto{\pgfqpoint{3.901565in}{2.303407in}}%
\pgfpathlineto{\pgfqpoint{3.902672in}{2.303440in}}%
\pgfpathlineto{\pgfqpoint{3.901955in}{2.303333in}}%
\pgfpathlineto{\pgfqpoint{3.902685in}{2.303415in}}%
\pgfpathlineto{\pgfqpoint{3.902694in}{2.303393in}}%
\pgfpathlineto{\pgfqpoint{3.903840in}{2.303440in}}%
\pgfpathlineto{\pgfqpoint{3.904545in}{2.303360in}}%
\pgfpathlineto{\pgfqpoint{3.903846in}{2.303450in}}%
\pgfpathlineto{\pgfqpoint{3.904951in}{2.303412in}}%
\pgfpathlineto{\pgfqpoint{3.905817in}{2.303461in}}%
\pgfpathlineto{\pgfqpoint{3.905074in}{2.303311in}}%
\pgfpathlineto{\pgfqpoint{3.906064in}{2.303433in}}%
\pgfpathlineto{\pgfqpoint{3.907085in}{2.303372in}}%
\pgfpathlineto{\pgfqpoint{3.906269in}{2.303459in}}%
\pgfpathlineto{\pgfqpoint{3.907179in}{2.303422in}}%
\pgfpathlineto{\pgfqpoint{3.908280in}{2.303373in}}%
\pgfpathlineto{\pgfqpoint{3.907194in}{2.303442in}}%
\pgfpathlineto{\pgfqpoint{3.908307in}{2.303404in}}%
\pgfpathlineto{\pgfqpoint{3.909230in}{2.303452in}}%
\pgfpathlineto{\pgfqpoint{3.909408in}{2.303314in}}%
\pgfpathlineto{\pgfqpoint{3.909417in}{2.303370in}}%
\pgfpathlineto{\pgfqpoint{3.909462in}{2.303300in}}%
\pgfpathlineto{\pgfqpoint{3.910384in}{2.304402in}}%
\pgfpathlineto{\pgfqpoint{3.910471in}{2.304133in}}%
\pgfpathlineto{\pgfqpoint{3.914287in}{2.304125in}}%
\pgfpathlineto{\pgfqpoint{3.910475in}{2.304147in}}%
\pgfpathlineto{\pgfqpoint{3.914291in}{2.304161in}}%
\pgfpathlineto{\pgfqpoint{3.915045in}{2.304270in}}%
\pgfpathlineto{\pgfqpoint{3.914642in}{2.303839in}}%
\pgfpathlineto{\pgfqpoint{3.915403in}{2.304143in}}%
\pgfpathlineto{\pgfqpoint{3.915585in}{2.303975in}}%
\pgfpathlineto{\pgfqpoint{3.916438in}{2.304579in}}%
\pgfpathlineto{\pgfqpoint{3.916910in}{2.304869in}}%
\pgfpathlineto{\pgfqpoint{3.917432in}{2.304175in}}%
\pgfpathlineto{\pgfqpoint{3.917530in}{2.304182in}}%
\pgfpathlineto{\pgfqpoint{3.918436in}{2.303697in}}%
\pgfpathlineto{\pgfqpoint{3.917554in}{2.304124in}}%
\pgfpathlineto{\pgfqpoint{3.918485in}{2.303516in}}%
\pgfpathlineto{\pgfqpoint{3.919284in}{2.303042in}}%
\pgfpathlineto{\pgfqpoint{3.918611in}{2.303657in}}%
\pgfpathlineto{\pgfqpoint{3.919605in}{2.303290in}}%
\pgfpathlineto{\pgfqpoint{3.920050in}{2.303634in}}%
\pgfpathlineto{\pgfqpoint{3.920502in}{2.302889in}}%
\pgfpathlineto{\pgfqpoint{3.920660in}{2.302871in}}%
\pgfpathlineto{\pgfqpoint{3.920944in}{2.302621in}}%
\pgfpathlineto{\pgfqpoint{3.921757in}{2.303244in}}%
\pgfpathlineto{\pgfqpoint{3.922308in}{2.303576in}}%
\pgfpathlineto{\pgfqpoint{3.921795in}{2.303191in}}%
\pgfpathlineto{\pgfqpoint{3.922868in}{2.303221in}}%
\pgfpathlineto{\pgfqpoint{3.922932in}{2.303098in}}%
\pgfpathlineto{\pgfqpoint{3.923301in}{2.303621in}}%
\pgfpathlineto{\pgfqpoint{3.923968in}{2.303386in}}%
\pgfpathlineto{\pgfqpoint{3.924875in}{2.303469in}}%
\pgfpathlineto{\pgfqpoint{3.924400in}{2.303119in}}%
\pgfpathlineto{\pgfqpoint{3.925080in}{2.303362in}}%
\pgfpathlineto{\pgfqpoint{3.925291in}{2.303511in}}%
\pgfpathlineto{\pgfqpoint{3.925849in}{2.302989in}}%
\pgfpathlineto{\pgfqpoint{3.926166in}{2.303110in}}%
\pgfpathlineto{\pgfqpoint{3.927369in}{2.302670in}}%
\pgfpathlineto{\pgfqpoint{3.926259in}{2.303273in}}%
\pgfpathlineto{\pgfqpoint{3.927380in}{2.302754in}}%
\pgfpathlineto{\pgfqpoint{3.927482in}{2.302878in}}%
\pgfpathlineto{\pgfqpoint{3.928291in}{2.302362in}}%
\pgfpathlineto{\pgfqpoint{3.928493in}{2.302780in}}%
\pgfpathlineto{\pgfqpoint{3.929418in}{2.302503in}}%
\pgfpathlineto{\pgfqpoint{3.929605in}{2.302743in}}%
\pgfpathlineto{\pgfqpoint{3.930618in}{2.302814in}}%
\pgfpathlineto{\pgfqpoint{3.929836in}{2.302475in}}%
\pgfpathlineto{\pgfqpoint{3.930701in}{2.302522in}}%
\pgfpathlineto{\pgfqpoint{3.932809in}{2.303244in}}%
\pgfpathlineto{\pgfqpoint{3.932890in}{2.302984in}}%
\pgfpathlineto{\pgfqpoint{3.933956in}{2.302710in}}%
\pgfpathlineto{\pgfqpoint{3.933202in}{2.303192in}}%
\pgfpathlineto{\pgfqpoint{3.934009in}{2.302872in}}%
\pgfpathlineto{\pgfqpoint{3.934118in}{2.303044in}}%
\pgfpathlineto{\pgfqpoint{3.934765in}{2.302309in}}%
\pgfpathlineto{\pgfqpoint{3.935048in}{2.302439in}}%
\pgfpathlineto{\pgfqpoint{3.935115in}{2.302367in}}%
\pgfpathlineto{\pgfqpoint{3.936137in}{2.303045in}}%
\pgfpathlineto{\pgfqpoint{3.936142in}{2.303028in}}%
\pgfpathlineto{\pgfqpoint{3.945296in}{2.300907in}}%
\pgfpathlineto{\pgfqpoint{3.945324in}{2.300811in}}%
\pgfpathlineto{\pgfqpoint{3.945670in}{2.301365in}}%
\pgfpathlineto{\pgfqpoint{3.946401in}{2.301125in}}%
\pgfpathlineto{\pgfqpoint{3.947522in}{2.301266in}}%
\pgfpathlineto{\pgfqpoint{3.946773in}{2.300944in}}%
\pgfpathlineto{\pgfqpoint{3.947537in}{2.301250in}}%
\pgfpathlineto{\pgfqpoint{3.947602in}{2.300988in}}%
\pgfpathlineto{\pgfqpoint{3.948540in}{2.301450in}}%
\pgfpathlineto{\pgfqpoint{3.948651in}{2.301230in}}%
\pgfpathlineto{\pgfqpoint{3.949666in}{2.301432in}}%
\pgfpathlineto{\pgfqpoint{3.949424in}{2.300880in}}%
\pgfpathlineto{\pgfqpoint{3.949766in}{2.301307in}}%
\pgfpathlineto{\pgfqpoint{3.950018in}{2.300911in}}%
\pgfpathlineto{\pgfqpoint{3.950495in}{2.301497in}}%
\pgfpathlineto{\pgfqpoint{3.950875in}{2.301360in}}%
\pgfpathlineto{\pgfqpoint{3.951635in}{2.301648in}}%
\pgfpathlineto{\pgfqpoint{3.950975in}{2.301098in}}%
\pgfpathlineto{\pgfqpoint{3.951991in}{2.301432in}}%
\pgfpathlineto{\pgfqpoint{3.952845in}{2.302003in}}%
\pgfpathlineto{\pgfqpoint{3.952265in}{2.301232in}}%
\pgfpathlineto{\pgfqpoint{3.953124in}{2.301529in}}%
\pgfpathlineto{\pgfqpoint{3.953148in}{2.301466in}}%
\pgfpathlineto{\pgfqpoint{3.954194in}{2.301976in}}%
\pgfpathlineto{\pgfqpoint{3.954223in}{2.301891in}}%
\pgfpathlineto{\pgfqpoint{3.955025in}{2.301473in}}%
\pgfpathlineto{\pgfqpoint{3.954230in}{2.301952in}}%
\pgfpathlineto{\pgfqpoint{3.955439in}{2.301884in}}%
\pgfpathlineto{\pgfqpoint{3.956508in}{2.302018in}}%
\pgfpathlineto{\pgfqpoint{3.955557in}{2.301636in}}%
\pgfpathlineto{\pgfqpoint{3.956639in}{2.301891in}}%
\pgfpathlineto{\pgfqpoint{3.956729in}{2.301855in}}%
\pgfpathlineto{\pgfqpoint{3.957190in}{2.302387in}}%
\pgfpathlineto{\pgfqpoint{3.957729in}{2.302215in}}%
\pgfpathlineto{\pgfqpoint{3.958808in}{2.302993in}}%
\pgfpathlineto{\pgfqpoint{3.957733in}{2.302204in}}%
\pgfpathlineto{\pgfqpoint{3.959182in}{2.302799in}}%
\pgfpathlineto{\pgfqpoint{3.960014in}{2.301958in}}%
\pgfpathlineto{\pgfqpoint{3.959221in}{2.302946in}}%
\pgfpathlineto{\pgfqpoint{3.960313in}{2.302390in}}%
\pgfpathlineto{\pgfqpoint{3.960896in}{2.302685in}}%
\pgfpathlineto{\pgfqpoint{3.960334in}{2.302336in}}%
\pgfpathlineto{\pgfqpoint{3.961501in}{2.302092in}}%
\pgfpathlineto{\pgfqpoint{3.961909in}{2.301976in}}%
\pgfpathlineto{\pgfqpoint{3.962513in}{2.302547in}}%
\pgfpathlineto{\pgfqpoint{3.962600in}{2.302337in}}%
\pgfpathlineto{\pgfqpoint{3.962707in}{2.302243in}}%
\pgfpathlineto{\pgfqpoint{3.963573in}{2.302972in}}%
\pgfpathlineto{\pgfqpoint{3.963574in}{2.302958in}}%
\pgfpathlineto{\pgfqpoint{3.964262in}{2.303062in}}%
\pgfpathlineto{\pgfqpoint{3.963785in}{2.302718in}}%
\pgfpathlineto{\pgfqpoint{3.964681in}{2.302813in}}%
\pgfpathlineto{\pgfqpoint{3.965935in}{2.303446in}}%
\pgfpathlineto{\pgfqpoint{3.965272in}{2.302498in}}%
\pgfpathlineto{\pgfqpoint{3.965954in}{2.303429in}}%
\pgfpathlineto{\pgfqpoint{3.966577in}{2.303354in}}%
\pgfpathlineto{\pgfqpoint{3.966119in}{2.303455in}}%
\pgfpathlineto{\pgfqpoint{3.967067in}{2.303386in}}%
\pgfpathlineto{\pgfqpoint{3.968114in}{2.303444in}}%
\pgfpathlineto{\pgfqpoint{3.967567in}{2.303317in}}%
\pgfpathlineto{\pgfqpoint{3.968184in}{2.303428in}}%
\pgfpathlineto{\pgfqpoint{3.969285in}{2.303364in}}%
\pgfpathlineto{\pgfqpoint{3.968205in}{2.303437in}}%
\pgfpathlineto{\pgfqpoint{3.969296in}{2.303403in}}%
\pgfpathlineto{\pgfqpoint{3.970417in}{2.303447in}}%
\pgfpathlineto{\pgfqpoint{3.969433in}{2.303364in}}%
\pgfpathlineto{\pgfqpoint{3.970424in}{2.303410in}}%
\pgfpathlineto{\pgfqpoint{3.971649in}{2.303368in}}%
\pgfpathlineto{\pgfqpoint{3.970427in}{2.303432in}}%
\pgfpathlineto{\pgfqpoint{3.971661in}{2.303411in}}%
\pgfpathlineto{\pgfqpoint{3.972591in}{2.303452in}}%
\pgfpathlineto{\pgfqpoint{3.972127in}{2.303357in}}%
\pgfpathlineto{\pgfqpoint{3.972773in}{2.303427in}}%
\pgfpathlineto{\pgfqpoint{3.973682in}{2.303357in}}%
\pgfpathlineto{\pgfqpoint{3.972924in}{2.303452in}}%
\pgfpathlineto{\pgfqpoint{3.973886in}{2.303386in}}%
\pgfpathlineto{\pgfqpoint{3.974970in}{2.303445in}}%
\pgfpathlineto{\pgfqpoint{3.973940in}{2.303343in}}%
\pgfpathlineto{\pgfqpoint{3.975019in}{2.303426in}}%
\pgfpathlineto{\pgfqpoint{3.976078in}{2.303360in}}%
\pgfpathlineto{\pgfqpoint{3.975340in}{2.303451in}}%
\pgfpathlineto{\pgfqpoint{3.976129in}{2.303411in}}%
\pgfpathlineto{\pgfqpoint{3.977239in}{2.303442in}}%
\pgfpathlineto{\pgfqpoint{3.976448in}{2.303336in}}%
\pgfpathlineto{\pgfqpoint{3.977242in}{2.303411in}}%
\pgfpathlineto{\pgfqpoint{3.977249in}{2.303395in}}%
\pgfpathlineto{\pgfqpoint{3.979327in}{2.303431in}}%
\pgfpathlineto{\pgfqpoint{3.980158in}{2.303362in}}%
\pgfpathlineto{\pgfqpoint{3.979492in}{2.303459in}}%
\pgfpathlineto{\pgfqpoint{3.980438in}{2.303418in}}%
\pgfpathlineto{\pgfqpoint{3.981484in}{2.303444in}}%
\pgfpathlineto{\pgfqpoint{3.980441in}{2.303390in}}%
\pgfpathlineto{\pgfqpoint{3.981551in}{2.303383in}}%
\pgfpathlineto{\pgfqpoint{3.983567in}{2.303434in}}%
\pgfpathlineto{\pgfqpoint{3.983571in}{2.303387in}}%
\pgfpathlineto{\pgfqpoint{3.984383in}{2.302982in}}%
\pgfpathlineto{\pgfqpoint{3.983711in}{2.303452in}}%
\pgfpathlineto{\pgfqpoint{3.984692in}{2.303288in}}%
\pgfpathlineto{\pgfqpoint{3.984699in}{2.303340in}}%
\pgfpathlineto{\pgfqpoint{3.985443in}{2.302664in}}%
\pgfpathlineto{\pgfqpoint{3.985788in}{2.302954in}}%
\pgfpathlineto{\pgfqpoint{3.986805in}{2.302357in}}%
\pgfpathlineto{\pgfqpoint{3.986110in}{2.303118in}}%
\pgfpathlineto{\pgfqpoint{3.986913in}{2.302431in}}%
\pgfpathlineto{\pgfqpoint{3.987060in}{2.302546in}}%
\pgfpathlineto{\pgfqpoint{3.987157in}{2.302192in}}%
\pgfpathlineto{\pgfqpoint{3.988023in}{2.302394in}}%
\pgfpathlineto{\pgfqpoint{3.988859in}{2.302310in}}%
\pgfpathlineto{\pgfqpoint{3.988516in}{2.302730in}}%
\pgfpathlineto{\pgfqpoint{3.989123in}{2.302558in}}%
\pgfpathlineto{\pgfqpoint{3.989166in}{2.302641in}}%
\pgfpathlineto{\pgfqpoint{3.989820in}{2.301823in}}%
\pgfpathlineto{\pgfqpoint{3.990147in}{2.301984in}}%
\pgfpathlineto{\pgfqpoint{3.991260in}{2.301834in}}%
\pgfpathlineto{\pgfqpoint{3.990332in}{2.302318in}}%
\pgfpathlineto{\pgfqpoint{3.991277in}{2.301881in}}%
\pgfpathlineto{\pgfqpoint{3.992359in}{2.302629in}}%
\pgfpathlineto{\pgfqpoint{3.991328in}{2.301773in}}%
\pgfpathlineto{\pgfqpoint{3.992410in}{2.302471in}}%
\pgfpathlineto{\pgfqpoint{3.993494in}{2.302032in}}%
\pgfpathlineto{\pgfqpoint{3.992505in}{2.302643in}}%
\pgfpathlineto{\pgfqpoint{3.993533in}{2.302106in}}%
\pgfpathlineto{\pgfqpoint{3.994135in}{2.301921in}}%
\pgfpathlineto{\pgfqpoint{3.993664in}{2.302327in}}%
\pgfpathlineto{\pgfqpoint{3.994636in}{2.302229in}}%
\pgfpathlineto{\pgfqpoint{3.994930in}{2.302507in}}%
\pgfpathlineto{\pgfqpoint{3.994725in}{2.302210in}}%
\pgfpathlineto{\pgfqpoint{3.995758in}{2.302298in}}%
\pgfpathlineto{\pgfqpoint{3.996185in}{2.302036in}}%
\pgfpathlineto{\pgfqpoint{3.996641in}{2.302685in}}%
\pgfpathlineto{\pgfqpoint{3.996864in}{2.302396in}}%
\pgfpathlineto{\pgfqpoint{3.997030in}{2.302464in}}%
\pgfpathlineto{\pgfqpoint{3.997815in}{2.301853in}}%
\pgfpathlineto{\pgfqpoint{3.997955in}{2.301972in}}%
\pgfpathlineto{\pgfqpoint{3.999178in}{2.301307in}}%
\pgfpathlineto{\pgfqpoint{3.998224in}{2.302367in}}%
\pgfpathlineto{\pgfqpoint{3.999190in}{2.301320in}}%
\pgfpathlineto{\pgfqpoint{4.000213in}{2.301445in}}%
\pgfpathlineto{\pgfqpoint{3.999848in}{2.300983in}}%
\pgfpathlineto{\pgfqpoint{4.000300in}{2.301278in}}%
\pgfpathlineto{\pgfqpoint{4.000922in}{2.300969in}}%
\pgfpathlineto{\pgfqpoint{4.000434in}{2.301358in}}%
\pgfpathlineto{\pgfqpoint{4.001428in}{2.301176in}}%
\pgfpathlineto{\pgfqpoint{4.002993in}{2.301862in}}%
\pgfpathlineto{\pgfqpoint{4.001893in}{2.300756in}}%
\pgfpathlineto{\pgfqpoint{4.003046in}{2.301775in}}%
\pgfpathlineto{\pgfqpoint{4.003923in}{2.301704in}}%
\pgfpathlineto{\pgfqpoint{4.003558in}{2.302223in}}%
\pgfpathlineto{\pgfqpoint{4.004154in}{2.301904in}}%
\pgfpathlineto{\pgfqpoint{4.004409in}{2.301493in}}%
\pgfpathlineto{\pgfqpoint{4.004159in}{2.301906in}}%
\pgfpathlineto{\pgfqpoint{4.005270in}{2.301954in}}%
\pgfpathlineto{\pgfqpoint{4.006427in}{2.302835in}}%
\pgfpathlineto{\pgfqpoint{4.005372in}{2.301881in}}%
\pgfpathlineto{\pgfqpoint{4.006466in}{2.302796in}}%
\pgfpathlineto{\pgfqpoint{4.006629in}{2.302520in}}%
\pgfpathlineto{\pgfqpoint{4.007070in}{2.302872in}}%
\pgfpathlineto{\pgfqpoint{4.007577in}{2.302813in}}%
\pgfpathlineto{\pgfqpoint{4.008484in}{2.303249in}}%
\pgfpathlineto{\pgfqpoint{4.007616in}{2.302727in}}%
\pgfpathlineto{\pgfqpoint{4.008724in}{2.303038in}}%
\pgfpathlineto{\pgfqpoint{4.008876in}{2.302731in}}%
\pgfpathlineto{\pgfqpoint{4.009213in}{2.303187in}}%
\pgfpathlineto{\pgfqpoint{4.009844in}{2.302922in}}%
\pgfpathlineto{\pgfqpoint{4.011195in}{2.303585in}}%
\pgfpathlineto{\pgfqpoint{4.009902in}{2.302846in}}%
\pgfpathlineto{\pgfqpoint{4.011229in}{2.303435in}}%
\pgfpathlineto{\pgfqpoint{4.012326in}{2.302884in}}%
\pgfpathlineto{\pgfqpoint{4.011262in}{2.303486in}}%
\pgfpathlineto{\pgfqpoint{4.012366in}{2.302942in}}%
\pgfpathlineto{\pgfqpoint{4.014541in}{2.303578in}}%
\pgfpathlineto{\pgfqpoint{4.012381in}{2.302888in}}%
\pgfpathlineto{\pgfqpoint{4.014568in}{2.303497in}}%
\pgfpathlineto{\pgfqpoint{4.015086in}{2.303503in}}%
\pgfpathlineto{\pgfqpoint{4.015614in}{2.304182in}}%
\pgfpathlineto{\pgfqpoint{4.015622in}{2.304145in}}%
\pgfpathlineto{\pgfqpoint{4.016533in}{2.305513in}}%
\pgfpathlineto{\pgfqpoint{4.015633in}{2.304140in}}%
\pgfpathlineto{\pgfqpoint{4.016773in}{2.305195in}}%
\pgfpathlineto{\pgfqpoint{4.017859in}{2.304757in}}%
\pgfpathlineto{\pgfqpoint{4.017065in}{2.305410in}}%
\pgfpathlineto{\pgfqpoint{4.017907in}{2.304890in}}%
\pgfpathlineto{\pgfqpoint{4.018141in}{2.304911in}}%
\pgfpathlineto{\pgfqpoint{4.018775in}{2.305741in}}%
\pgfpathlineto{\pgfqpoint{4.018829in}{2.305673in}}%
\pgfpathlineto{\pgfqpoint{4.019533in}{2.306099in}}%
\pgfpathlineto{\pgfqpoint{4.019112in}{2.305558in}}%
\pgfpathlineto{\pgfqpoint{4.019941in}{2.305663in}}%
\pgfpathlineto{\pgfqpoint{4.020645in}{2.305325in}}%
\pgfpathlineto{\pgfqpoint{4.020109in}{2.305736in}}%
\pgfpathlineto{\pgfqpoint{4.021066in}{2.305427in}}%
\pgfpathlineto{\pgfqpoint{4.022564in}{2.304732in}}%
\pgfpathlineto{\pgfqpoint{4.021068in}{2.305446in}}%
\pgfpathlineto{\pgfqpoint{4.022656in}{2.304835in}}%
\pgfpathlineto{\pgfqpoint{4.022726in}{2.304905in}}%
\pgfpathlineto{\pgfqpoint{4.023736in}{2.304145in}}%
\pgfpathlineto{\pgfqpoint{4.023745in}{2.304149in}}%
\pgfpathlineto{\pgfqpoint{4.024766in}{2.304195in}}%
\pgfpathlineto{\pgfqpoint{4.024184in}{2.303956in}}%
\pgfpathlineto{\pgfqpoint{4.024850in}{2.304045in}}%
\pgfpathlineto{\pgfqpoint{4.024851in}{2.304027in}}%
\pgfpathlineto{\pgfqpoint{4.025883in}{2.305006in}}%
\pgfpathlineto{\pgfqpoint{4.025895in}{2.304983in}}%
\pgfpathlineto{\pgfqpoint{4.026865in}{2.305033in}}%
\pgfpathlineto{\pgfqpoint{4.026513in}{2.304749in}}%
\pgfpathlineto{\pgfqpoint{4.026994in}{2.304790in}}%
\pgfpathlineto{\pgfqpoint{4.026996in}{2.304761in}}%
\pgfpathlineto{\pgfqpoint{4.027398in}{2.305340in}}%
\pgfpathlineto{\pgfqpoint{4.028070in}{2.305232in}}%
\pgfpathlineto{\pgfqpoint{4.029412in}{2.305502in}}%
\pgfpathlineto{\pgfqpoint{4.028114in}{2.305156in}}%
\pgfpathlineto{\pgfqpoint{4.029444in}{2.305453in}}%
\pgfpathlineto{\pgfqpoint{4.030457in}{2.305310in}}%
\pgfpathlineto{\pgfqpoint{4.030003in}{2.305688in}}%
\pgfpathlineto{\pgfqpoint{4.030552in}{2.305557in}}%
\pgfpathlineto{\pgfqpoint{4.031893in}{2.306082in}}%
\pgfpathlineto{\pgfqpoint{4.030553in}{2.305549in}}%
\pgfpathlineto{\pgfqpoint{4.032032in}{2.305926in}}%
\pgfpathlineto{\pgfqpoint{4.032033in}{2.305915in}}%
\pgfpathlineto{\pgfqpoint{4.032779in}{2.306928in}}%
\pgfpathlineto{\pgfqpoint{4.033098in}{2.306636in}}%
\pgfpathlineto{\pgfqpoint{4.033130in}{2.306736in}}%
\pgfpathlineto{\pgfqpoint{4.033778in}{2.305907in}}%
\pgfpathlineto{\pgfqpoint{4.034171in}{2.306002in}}%
\pgfpathlineto{\pgfqpoint{4.035339in}{2.305948in}}%
\pgfpathlineto{\pgfqpoint{4.034315in}{2.306259in}}%
\pgfpathlineto{\pgfqpoint{4.035353in}{2.305985in}}%
\pgfpathlineto{\pgfqpoint{4.036027in}{2.306134in}}%
\pgfpathlineto{\pgfqpoint{4.036407in}{2.305699in}}%
\pgfpathlineto{\pgfqpoint{4.036456in}{2.305757in}}%
\pgfpathlineto{\pgfqpoint{4.037065in}{2.305689in}}%
\pgfpathlineto{\pgfqpoint{4.036902in}{2.306044in}}%
\pgfpathlineto{\pgfqpoint{4.037555in}{2.305903in}}%
\pgfpathlineto{\pgfqpoint{4.038671in}{2.306455in}}%
\pgfpathlineto{\pgfqpoint{4.037602in}{2.305814in}}%
\pgfpathlineto{\pgfqpoint{4.038686in}{2.306419in}}%
\pgfpathlineto{\pgfqpoint{4.039382in}{2.306855in}}%
\pgfpathlineto{\pgfqpoint{4.039768in}{2.306094in}}%
\pgfpathlineto{\pgfqpoint{4.039778in}{2.306126in}}%
\pgfpathlineto{\pgfqpoint{4.041088in}{2.303373in}}%
\pgfpathlineto{\pgfqpoint{4.040214in}{2.306259in}}%
\pgfpathlineto{\pgfqpoint{4.041102in}{2.303421in}}%
\pgfpathlineto{\pgfqpoint{4.042216in}{2.303440in}}%
\pgfpathlineto{\pgfqpoint{4.041125in}{2.303334in}}%
\pgfpathlineto{\pgfqpoint{4.042220in}{2.303427in}}%
\pgfpathlineto{\pgfqpoint{4.043240in}{2.303370in}}%
\pgfpathlineto{\pgfqpoint{4.042256in}{2.303448in}}%
\pgfpathlineto{\pgfqpoint{4.043332in}{2.303400in}}%
\pgfpathlineto{\pgfqpoint{4.044615in}{2.303440in}}%
\pgfpathlineto{\pgfqpoint{4.044622in}{2.303417in}}%
\pgfpathlineto{\pgfqpoint{4.045709in}{2.303371in}}%
\pgfpathlineto{\pgfqpoint{4.044737in}{2.303444in}}%
\pgfpathlineto{\pgfqpoint{4.045735in}{2.303402in}}%
\pgfpathlineto{\pgfqpoint{4.046870in}{2.303440in}}%
\pgfpathlineto{\pgfqpoint{4.045792in}{2.303375in}}%
\pgfpathlineto{\pgfqpoint{4.046875in}{2.303410in}}%
\pgfpathlineto{\pgfqpoint{4.047916in}{2.303373in}}%
\pgfpathlineto{\pgfqpoint{4.046922in}{2.303442in}}%
\pgfpathlineto{\pgfqpoint{4.047997in}{2.303413in}}%
\pgfpathlineto{\pgfqpoint{4.049092in}{2.303441in}}%
\pgfpathlineto{\pgfqpoint{4.048105in}{2.303369in}}%
\pgfpathlineto{\pgfqpoint{4.049110in}{2.303431in}}%
\pgfpathlineto{\pgfqpoint{4.049940in}{2.303349in}}%
\pgfpathlineto{\pgfqpoint{4.049133in}{2.303446in}}%
\pgfpathlineto{\pgfqpoint{4.050223in}{2.303424in}}%
\pgfpathlineto{\pgfqpoint{4.051326in}{2.303442in}}%
\pgfpathlineto{\pgfqpoint{4.050273in}{2.303385in}}%
\pgfpathlineto{\pgfqpoint{4.051337in}{2.303425in}}%
\pgfpathlineto{\pgfqpoint{4.052468in}{2.303384in}}%
\pgfpathlineto{\pgfqpoint{4.051363in}{2.303447in}}%
\pgfpathlineto{\pgfqpoint{4.052471in}{2.303406in}}%
\pgfpathlineto{\pgfqpoint{4.053635in}{2.303438in}}%
\pgfpathlineto{\pgfqpoint{4.052491in}{2.303375in}}%
\pgfpathlineto{\pgfqpoint{4.053638in}{2.303426in}}%
\pgfpathlineto{\pgfqpoint{4.054287in}{2.303354in}}%
\pgfpathlineto{\pgfqpoint{4.053807in}{2.303450in}}%
\pgfpathlineto{\pgfqpoint{4.054750in}{2.303396in}}%
\pgfpathlineto{\pgfqpoint{4.055914in}{2.303437in}}%
\pgfpathlineto{\pgfqpoint{4.054759in}{2.303390in}}%
\pgfpathlineto{\pgfqpoint{4.055935in}{2.303419in}}%
\pgfpathlineto{\pgfqpoint{4.057033in}{2.303378in}}%
\pgfpathlineto{\pgfqpoint{4.055951in}{2.303437in}}%
\pgfpathlineto{\pgfqpoint{4.057060in}{2.303413in}}%
\pgfpathlineto{\pgfqpoint{4.059909in}{2.303663in}}%
\pgfpathlineto{\pgfqpoint{4.057063in}{2.303433in}}%
\pgfpathlineto{\pgfqpoint{4.059942in}{2.303826in}}%
\pgfpathlineto{\pgfqpoint{4.060697in}{2.303967in}}%
\pgfpathlineto{\pgfqpoint{4.060370in}{2.303551in}}%
\pgfpathlineto{\pgfqpoint{4.061043in}{2.303599in}}%
\pgfpathlineto{\pgfqpoint{4.061168in}{2.303537in}}%
\pgfpathlineto{\pgfqpoint{4.061534in}{2.303922in}}%
\pgfpathlineto{\pgfqpoint{4.062145in}{2.303768in}}%
\pgfpathlineto{\pgfqpoint{4.062533in}{2.304009in}}%
\pgfpathlineto{\pgfqpoint{4.062943in}{2.303571in}}%
\pgfpathlineto{\pgfqpoint{4.063254in}{2.303715in}}%
\pgfpathlineto{\pgfqpoint{4.064490in}{2.302962in}}%
\pgfpathlineto{\pgfqpoint{4.063259in}{2.303734in}}%
\pgfpathlineto{\pgfqpoint{4.064526in}{2.303056in}}%
\pgfpathlineto{\pgfqpoint{4.065288in}{2.303396in}}%
\pgfpathlineto{\pgfqpoint{4.064789in}{2.302972in}}%
\pgfpathlineto{\pgfqpoint{4.065650in}{2.303284in}}%
\pgfpathlineto{\pgfqpoint{4.065707in}{2.303163in}}%
\pgfpathlineto{\pgfqpoint{4.066533in}{2.303580in}}%
\pgfpathlineto{\pgfqpoint{4.066759in}{2.303338in}}%
\pgfpathlineto{\pgfqpoint{4.067867in}{2.303524in}}%
\pgfpathlineto{\pgfqpoint{4.066953in}{2.303040in}}%
\pgfpathlineto{\pgfqpoint{4.067943in}{2.303417in}}%
\pgfpathlineto{\pgfqpoint{4.068755in}{2.302436in}}%
\pgfpathlineto{\pgfqpoint{4.069078in}{2.302775in}}%
\pgfpathlineto{\pgfqpoint{4.070202in}{2.302832in}}%
\pgfpathlineto{\pgfqpoint{4.069303in}{2.302658in}}%
\pgfpathlineto{\pgfqpoint{4.070216in}{2.302763in}}%
\pgfpathlineto{\pgfqpoint{4.070445in}{2.302602in}}%
\pgfpathlineto{\pgfqpoint{4.071098in}{2.303021in}}%
\pgfpathlineto{\pgfqpoint{4.071319in}{2.302931in}}%
\pgfpathlineto{\pgfqpoint{4.071568in}{2.302827in}}%
\pgfpathlineto{\pgfqpoint{4.071729in}{2.303263in}}%
\pgfpathlineto{\pgfqpoint{4.072389in}{2.303280in}}%
\pgfpathlineto{\pgfqpoint{4.072682in}{2.303497in}}%
\pgfpathlineto{\pgfqpoint{4.073436in}{2.302866in}}%
\pgfpathlineto{\pgfqpoint{4.073483in}{2.302884in}}%
\pgfpathlineto{\pgfqpoint{4.073655in}{2.303064in}}%
\pgfpathlineto{\pgfqpoint{4.074485in}{2.302347in}}%
\pgfpathlineto{\pgfqpoint{4.074957in}{2.302028in}}%
\pgfpathlineto{\pgfqpoint{4.074624in}{2.302467in}}%
\pgfpathlineto{\pgfqpoint{4.075599in}{2.302316in}}%
\pgfpathlineto{\pgfqpoint{4.076684in}{2.302344in}}%
\pgfpathlineto{\pgfqpoint{4.076209in}{2.301692in}}%
\pgfpathlineto{\pgfqpoint{4.076707in}{2.302259in}}%
\pgfpathlineto{\pgfqpoint{4.077068in}{2.301789in}}%
\pgfpathlineto{\pgfqpoint{4.076737in}{2.302328in}}%
\pgfpathlineto{\pgfqpoint{4.077819in}{2.302261in}}%
\pgfpathlineto{\pgfqpoint{4.078378in}{2.302712in}}%
\pgfpathlineto{\pgfqpoint{4.078899in}{2.301935in}}%
\pgfpathlineto{\pgfqpoint{4.078918in}{2.301973in}}%
\pgfpathlineto{\pgfqpoint{4.078918in}{2.301969in}}%
\pgfpathlineto{\pgfqpoint{4.079851in}{2.302777in}}%
\pgfpathlineto{\pgfqpoint{4.079920in}{2.302661in}}%
\pgfpathlineto{\pgfqpoint{4.080398in}{2.303050in}}%
\pgfpathlineto{\pgfqpoint{4.080030in}{2.302594in}}%
\pgfpathlineto{\pgfqpoint{4.081036in}{2.302809in}}%
\pgfpathlineto{\pgfqpoint{4.081071in}{2.302673in}}%
\pgfpathlineto{\pgfqpoint{4.081876in}{2.303307in}}%
\pgfpathlineto{\pgfqpoint{4.082132in}{2.303044in}}%
\pgfpathlineto{\pgfqpoint{4.083155in}{2.303265in}}%
\pgfpathlineto{\pgfqpoint{4.082356in}{2.302810in}}%
\pgfpathlineto{\pgfqpoint{4.083262in}{2.303176in}}%
\pgfpathlineto{\pgfqpoint{4.084240in}{2.302542in}}%
\pgfpathlineto{\pgfqpoint{4.083431in}{2.303249in}}%
\pgfpathlineto{\pgfqpoint{4.084398in}{2.302831in}}%
\pgfpathlineto{\pgfqpoint{4.084870in}{2.303355in}}%
\pgfpathlineto{\pgfqpoint{4.085323in}{2.302543in}}%
\pgfpathlineto{\pgfqpoint{4.085511in}{2.302848in}}%
\pgfpathlineto{\pgfqpoint{4.086380in}{2.302660in}}%
\pgfpathlineto{\pgfqpoint{4.085685in}{2.303063in}}%
\pgfpathlineto{\pgfqpoint{4.086623in}{2.302840in}}%
\pgfpathlineto{\pgfqpoint{4.087743in}{2.303304in}}%
\pgfpathlineto{\pgfqpoint{4.086712in}{2.302822in}}%
\pgfpathlineto{\pgfqpoint{4.087749in}{2.303266in}}%
\pgfpathlineto{\pgfqpoint{4.088154in}{2.303098in}}%
\pgfpathlineto{\pgfqpoint{4.088536in}{2.303520in}}%
\pgfpathlineto{\pgfqpoint{4.088851in}{2.303393in}}%
\pgfpathlineto{\pgfqpoint{4.089912in}{2.304017in}}%
\pgfpathlineto{\pgfqpoint{4.089102in}{2.303303in}}%
\pgfpathlineto{\pgfqpoint{4.089990in}{2.303823in}}%
\pgfpathlineto{\pgfqpoint{4.091050in}{2.303336in}}%
\pgfpathlineto{\pgfqpoint{4.089998in}{2.303837in}}%
\pgfpathlineto{\pgfqpoint{4.091143in}{2.303469in}}%
\pgfpathlineto{\pgfqpoint{4.092003in}{2.303604in}}%
\pgfpathlineto{\pgfqpoint{4.091830in}{2.304143in}}%
\pgfpathlineto{\pgfqpoint{4.092191in}{2.303835in}}%
\pgfpathlineto{\pgfqpoint{4.093036in}{2.303993in}}%
\pgfpathlineto{\pgfqpoint{4.092749in}{2.303659in}}%
\pgfpathlineto{\pgfqpoint{4.093303in}{2.303826in}}%
\pgfpathlineto{\pgfqpoint{4.093412in}{2.303750in}}%
\pgfpathlineto{\pgfqpoint{4.094196in}{2.304498in}}%
\pgfpathlineto{\pgfqpoint{4.094268in}{2.304483in}}%
\pgfpathlineto{\pgfqpoint{4.094743in}{2.304660in}}%
\pgfpathlineto{\pgfqpoint{4.095356in}{2.304224in}}%
\pgfpathlineto{\pgfqpoint{4.095949in}{2.303994in}}%
\pgfpathlineto{\pgfqpoint{4.095588in}{2.304589in}}%
\pgfpathlineto{\pgfqpoint{4.096472in}{2.304183in}}%
\pgfpathlineto{\pgfqpoint{4.097105in}{2.304230in}}%
\pgfpathlineto{\pgfqpoint{4.097480in}{2.303646in}}%
\pgfpathlineto{\pgfqpoint{4.098563in}{2.302716in}}%
\pgfpathlineto{\pgfqpoint{4.097488in}{2.303673in}}%
\pgfpathlineto{\pgfqpoint{4.098629in}{2.302847in}}%
\pgfpathlineto{\pgfqpoint{4.098630in}{2.302853in}}%
\pgfpathlineto{\pgfqpoint{4.099043in}{2.302312in}}%
\pgfpathlineto{\pgfqpoint{4.099712in}{2.302582in}}%
\pgfpathlineto{\pgfqpoint{4.100825in}{2.302292in}}%
\pgfpathlineto{\pgfqpoint{4.099908in}{2.302839in}}%
\pgfpathlineto{\pgfqpoint{4.100844in}{2.302304in}}%
\pgfpathlineto{\pgfqpoint{4.102067in}{2.302554in}}%
\pgfpathlineto{\pgfqpoint{4.100853in}{2.302382in}}%
\pgfpathlineto{\pgfqpoint{4.102083in}{2.302607in}}%
\pgfpathlineto{\pgfqpoint{4.103146in}{2.302924in}}%
\pgfpathlineto{\pgfqpoint{4.102197in}{2.302418in}}%
\pgfpathlineto{\pgfqpoint{4.103205in}{2.302825in}}%
\pgfpathlineto{\pgfqpoint{4.104306in}{2.302805in}}%
\pgfpathlineto{\pgfqpoint{4.103495in}{2.303129in}}%
\pgfpathlineto{\pgfqpoint{4.104319in}{2.302832in}}%
\pgfpathlineto{\pgfqpoint{4.105552in}{2.303519in}}%
\pgfpathlineto{\pgfqpoint{4.104345in}{2.302794in}}%
\pgfpathlineto{\pgfqpoint{4.105577in}{2.303444in}}%
\pgfpathlineto{\pgfqpoint{4.105840in}{2.303376in}}%
\pgfpathlineto{\pgfqpoint{4.105912in}{2.303694in}}%
\pgfpathlineto{\pgfqpoint{4.106679in}{2.303599in}}%
\pgfpathlineto{\pgfqpoint{4.107592in}{2.303824in}}%
\pgfpathlineto{\pgfqpoint{4.106924in}{2.303366in}}%
\pgfpathlineto{\pgfqpoint{4.107799in}{2.303603in}}%
\pgfpathlineto{\pgfqpoint{4.107929in}{2.303527in}}%
\pgfpathlineto{\pgfqpoint{4.108859in}{2.304252in}}%
\pgfpathlineto{\pgfqpoint{4.108860in}{2.304243in}}%
\pgfpathlineto{\pgfqpoint{4.110302in}{2.304510in}}%
\pgfpathlineto{\pgfqpoint{4.108947in}{2.304106in}}%
\pgfpathlineto{\pgfqpoint{4.110349in}{2.304366in}}%
\pgfpathlineto{\pgfqpoint{4.110472in}{2.304095in}}%
\pgfpathlineto{\pgfqpoint{4.111403in}{2.304867in}}%
\pgfpathlineto{\pgfqpoint{4.111404in}{2.304861in}}%
\pgfpathlineto{\pgfqpoint{4.112651in}{2.305688in}}%
\pgfpathlineto{\pgfqpoint{4.111471in}{2.304776in}}%
\pgfpathlineto{\pgfqpoint{4.112656in}{2.305676in}}%
\pgfpathlineto{\pgfqpoint{4.113370in}{2.305555in}}%
\pgfpathlineto{\pgfqpoint{4.113636in}{2.306021in}}%
\pgfpathlineto{\pgfqpoint{4.113757in}{2.305877in}}%
\pgfpathlineto{\pgfqpoint{4.114750in}{2.305920in}}%
\pgfpathlineto{\pgfqpoint{4.114032in}{2.305515in}}%
\pgfpathlineto{\pgfqpoint{4.114869in}{2.305875in}}%
\pgfpathlineto{\pgfqpoint{4.116297in}{2.303371in}}%
\pgfpathlineto{\pgfqpoint{4.114871in}{2.305888in}}%
\pgfpathlineto{\pgfqpoint{4.116332in}{2.303409in}}%
\pgfpathlineto{\pgfqpoint{4.117425in}{2.303451in}}%
\pgfpathlineto{\pgfqpoint{4.116391in}{2.303358in}}%
\pgfpathlineto{\pgfqpoint{4.117461in}{2.303425in}}%
\pgfpathlineto{\pgfqpoint{4.118114in}{2.303342in}}%
\pgfpathlineto{\pgfqpoint{4.117715in}{2.303459in}}%
\pgfpathlineto{\pgfqpoint{4.118573in}{2.303401in}}%
\pgfpathlineto{\pgfqpoint{4.119577in}{2.303448in}}%
\pgfpathlineto{\pgfqpoint{4.118605in}{2.303356in}}%
\pgfpathlineto{\pgfqpoint{4.119686in}{2.303395in}}%
\pgfpathlineto{\pgfqpoint{4.120766in}{2.303378in}}%
\pgfpathlineto{\pgfqpoint{4.119698in}{2.303446in}}%
\pgfpathlineto{\pgfqpoint{4.120798in}{2.303411in}}%
\pgfpathlineto{\pgfqpoint{4.121855in}{2.303444in}}%
\pgfpathlineto{\pgfqpoint{4.120905in}{2.303361in}}%
\pgfpathlineto{\pgfqpoint{4.121912in}{2.303424in}}%
\pgfpathlineto{\pgfqpoint{4.122951in}{2.303362in}}%
\pgfpathlineto{\pgfqpoint{4.122155in}{2.303456in}}%
\pgfpathlineto{\pgfqpoint{4.123024in}{2.303417in}}%
\pgfpathlineto{\pgfqpoint{4.124369in}{2.303383in}}%
\pgfpathlineto{\pgfqpoint{4.123026in}{2.303427in}}%
\pgfpathlineto{\pgfqpoint{4.124376in}{2.303419in}}%
\pgfpathlineto{\pgfqpoint{4.125477in}{2.303450in}}%
\pgfpathlineto{\pgfqpoint{4.124492in}{2.303374in}}%
\pgfpathlineto{\pgfqpoint{4.125496in}{2.303427in}}%
\pgfpathlineto{\pgfqpoint{4.126201in}{2.303352in}}%
\pgfpathlineto{\pgfqpoint{4.125632in}{2.303446in}}%
\pgfpathlineto{\pgfqpoint{4.126607in}{2.303393in}}%
\pgfpathlineto{\pgfqpoint{4.127678in}{2.303445in}}%
\pgfpathlineto{\pgfqpoint{4.126636in}{2.303379in}}%
\pgfpathlineto{\pgfqpoint{4.127728in}{2.303432in}}%
\pgfpathlineto{\pgfqpoint{4.128815in}{2.303366in}}%
\pgfpathlineto{\pgfqpoint{4.127875in}{2.303447in}}%
\pgfpathlineto{\pgfqpoint{4.128840in}{2.303416in}}%
\pgfpathlineto{\pgfqpoint{4.129948in}{2.303447in}}%
\pgfpathlineto{\pgfqpoint{4.129104in}{2.303364in}}%
\pgfpathlineto{\pgfqpoint{4.129953in}{2.303438in}}%
\pgfpathlineto{\pgfqpoint{4.131174in}{2.303437in}}%
\pgfpathlineto{\pgfqpoint{4.129996in}{2.303369in}}%
\pgfpathlineto{\pgfqpoint{4.131181in}{2.303426in}}%
\pgfpathlineto{\pgfqpoint{4.132115in}{2.303360in}}%
\pgfpathlineto{\pgfqpoint{4.131196in}{2.303445in}}%
\pgfpathlineto{\pgfqpoint{4.132293in}{2.303375in}}%
\pgfpathlineto{\pgfqpoint{4.133309in}{2.303445in}}%
\pgfpathlineto{\pgfqpoint{4.132411in}{2.303351in}}%
\pgfpathlineto{\pgfqpoint{4.133405in}{2.303401in}}%
\pgfpathlineto{\pgfqpoint{4.134480in}{2.303169in}}%
\pgfpathlineto{\pgfqpoint{4.134154in}{2.303533in}}%
\pgfpathlineto{\pgfqpoint{4.134530in}{2.303233in}}%
\pgfpathlineto{\pgfqpoint{4.135100in}{2.303315in}}%
\pgfpathlineto{\pgfqpoint{4.135623in}{2.302787in}}%
\pgfpathlineto{\pgfqpoint{4.135627in}{2.302816in}}%
\pgfpathlineto{\pgfqpoint{4.136415in}{2.302756in}}%
\pgfpathlineto{\pgfqpoint{4.136129in}{2.303131in}}%
\pgfpathlineto{\pgfqpoint{4.136690in}{2.303137in}}%
\pgfpathlineto{\pgfqpoint{4.137153in}{2.303538in}}%
\pgfpathlineto{\pgfqpoint{4.137595in}{2.302823in}}%
\pgfpathlineto{\pgfqpoint{4.137797in}{2.303029in}}%
\pgfpathlineto{\pgfqpoint{4.138953in}{2.303313in}}%
\pgfpathlineto{\pgfqpoint{4.137868in}{2.302853in}}%
\pgfpathlineto{\pgfqpoint{4.138960in}{2.303249in}}%
\pgfpathlineto{\pgfqpoint{4.139859in}{2.302707in}}%
\pgfpathlineto{\pgfqpoint{4.139020in}{2.303334in}}%
\pgfpathlineto{\pgfqpoint{4.140160in}{2.303015in}}%
\pgfpathlineto{\pgfqpoint{4.140821in}{2.303593in}}%
\pgfpathlineto{\pgfqpoint{4.141282in}{2.303224in}}%
\pgfpathlineto{\pgfqpoint{4.141333in}{2.303204in}}%
\pgfpathlineto{\pgfqpoint{4.142334in}{2.303716in}}%
\pgfpathlineto{\pgfqpoint{4.142362in}{2.303653in}}%
\pgfpathlineto{\pgfqpoint{4.142516in}{2.303879in}}%
\pgfpathlineto{\pgfqpoint{4.143345in}{2.303154in}}%
\pgfpathlineto{\pgfqpoint{4.143464in}{2.303396in}}%
\pgfpathlineto{\pgfqpoint{4.144361in}{2.302915in}}%
\pgfpathlineto{\pgfqpoint{4.143529in}{2.303190in}}%
\pgfpathlineto{\pgfqpoint{4.144396in}{2.302687in}}%
\pgfpathlineto{\pgfqpoint{4.145230in}{2.302553in}}%
\pgfpathlineto{\pgfqpoint{4.144866in}{2.302961in}}%
\pgfpathlineto{\pgfqpoint{4.145507in}{2.302722in}}%
\pgfpathlineto{\pgfqpoint{4.146377in}{2.302921in}}%
\pgfpathlineto{\pgfqpoint{4.145714in}{2.302471in}}%
\pgfpathlineto{\pgfqpoint{4.146615in}{2.302644in}}%
\pgfpathlineto{\pgfqpoint{4.147802in}{2.302102in}}%
\pgfpathlineto{\pgfqpoint{4.146811in}{2.302720in}}%
\pgfpathlineto{\pgfqpoint{4.147824in}{2.302236in}}%
\pgfpathlineto{\pgfqpoint{4.148296in}{2.302740in}}%
\pgfpathlineto{\pgfqpoint{4.148953in}{2.302475in}}%
\pgfpathlineto{\pgfqpoint{4.149299in}{2.302653in}}%
\pgfpathlineto{\pgfqpoint{4.149885in}{2.302120in}}%
\pgfpathlineto{\pgfqpoint{4.150050in}{2.302317in}}%
\pgfpathlineto{\pgfqpoint{4.151057in}{2.301947in}}%
\pgfpathlineto{\pgfqpoint{4.150497in}{2.302398in}}%
\pgfpathlineto{\pgfqpoint{4.151170in}{2.302121in}}%
\pgfpathlineto{\pgfqpoint{4.152182in}{2.302362in}}%
\pgfpathlineto{\pgfqpoint{4.151340in}{2.301958in}}%
\pgfpathlineto{\pgfqpoint{4.152332in}{2.302245in}}%
\pgfpathlineto{\pgfqpoint{4.152451in}{2.302213in}}%
\pgfpathlineto{\pgfqpoint{4.153138in}{2.302991in}}%
\pgfpathlineto{\pgfqpoint{4.153399in}{2.302908in}}%
\pgfpathlineto{\pgfqpoint{4.154518in}{2.302851in}}%
\pgfpathlineto{\pgfqpoint{4.154143in}{2.302536in}}%
\pgfpathlineto{\pgfqpoint{4.154524in}{2.302795in}}%
\pgfpathlineto{\pgfqpoint{4.155597in}{2.302629in}}%
\pgfpathlineto{\pgfqpoint{4.155015in}{2.303077in}}%
\pgfpathlineto{\pgfqpoint{4.155641in}{2.302705in}}%
\pgfpathlineto{\pgfqpoint{4.156614in}{2.302965in}}%
\pgfpathlineto{\pgfqpoint{4.155643in}{2.302701in}}%
\pgfpathlineto{\pgfqpoint{4.156834in}{2.302507in}}%
\pgfpathlineto{\pgfqpoint{4.157104in}{2.302423in}}%
\pgfpathlineto{\pgfqpoint{4.157881in}{2.303367in}}%
\pgfpathlineto{\pgfqpoint{4.157912in}{2.303279in}}%
\pgfpathlineto{\pgfqpoint{4.158246in}{2.303665in}}%
\pgfpathlineto{\pgfqpoint{4.158959in}{2.303025in}}%
\pgfpathlineto{\pgfqpoint{4.159017in}{2.303106in}}%
\pgfpathlineto{\pgfqpoint{4.160167in}{2.303360in}}%
\pgfpathlineto{\pgfqpoint{4.159300in}{2.302845in}}%
\pgfpathlineto{\pgfqpoint{4.160168in}{2.303344in}}%
\pgfpathlineto{\pgfqpoint{4.161277in}{2.303285in}}%
\pgfpathlineto{\pgfqpoint{4.160473in}{2.303570in}}%
\pgfpathlineto{\pgfqpoint{4.161284in}{2.303345in}}%
\pgfpathlineto{\pgfqpoint{4.161403in}{2.303495in}}%
\pgfpathlineto{\pgfqpoint{4.162224in}{2.302851in}}%
\pgfpathlineto{\pgfqpoint{4.162376in}{2.302956in}}%
\pgfpathlineto{\pgfqpoint{4.162418in}{2.302897in}}%
\pgfpathlineto{\pgfqpoint{4.163292in}{2.303507in}}%
\pgfpathlineto{\pgfqpoint{4.163464in}{2.303435in}}%
\pgfpathlineto{\pgfqpoint{4.163632in}{2.303550in}}%
\pgfpathlineto{\pgfqpoint{4.164249in}{2.302767in}}%
\pgfpathlineto{\pgfqpoint{4.164557in}{2.303085in}}%
\pgfpathlineto{\pgfqpoint{4.164912in}{2.302934in}}%
\pgfpathlineto{\pgfqpoint{4.165298in}{2.303389in}}%
\pgfpathlineto{\pgfqpoint{4.165659in}{2.303259in}}%
\pgfpathlineto{\pgfqpoint{4.166229in}{2.303888in}}%
\pgfpathlineto{\pgfqpoint{4.165718in}{2.303172in}}%
\pgfpathlineto{\pgfqpoint{4.166778in}{2.303414in}}%
\pgfpathlineto{\pgfqpoint{4.168491in}{2.303586in}}%
\pgfpathlineto{\pgfqpoint{4.166781in}{2.303378in}}%
\pgfpathlineto{\pgfqpoint{4.168521in}{2.303472in}}%
\pgfpathlineto{\pgfqpoint{4.168836in}{2.303334in}}%
\pgfpathlineto{\pgfqpoint{4.169534in}{2.304481in}}%
\pgfpathlineto{\pgfqpoint{4.169610in}{2.304311in}}%
\pgfpathlineto{\pgfqpoint{4.170740in}{2.304919in}}%
\pgfpathlineto{\pgfqpoint{4.170746in}{2.304880in}}%
\pgfpathlineto{\pgfqpoint{4.171907in}{2.304601in}}%
\pgfpathlineto{\pgfqpoint{4.171025in}{2.305015in}}%
\pgfpathlineto{\pgfqpoint{4.171908in}{2.304607in}}%
\pgfpathlineto{\pgfqpoint{4.172986in}{2.304830in}}%
\pgfpathlineto{\pgfqpoint{4.172278in}{2.304388in}}%
\pgfpathlineto{\pgfqpoint{4.173033in}{2.304761in}}%
\pgfpathlineto{\pgfqpoint{4.174024in}{2.304449in}}%
\pgfpathlineto{\pgfqpoint{4.173286in}{2.304994in}}%
\pgfpathlineto{\pgfqpoint{4.174163in}{2.304586in}}%
\pgfpathlineto{\pgfqpoint{4.175148in}{2.304858in}}%
\pgfpathlineto{\pgfqpoint{4.174260in}{2.304524in}}%
\pgfpathlineto{\pgfqpoint{4.175289in}{2.304672in}}%
\pgfpathlineto{\pgfqpoint{4.175371in}{2.304629in}}%
\pgfpathlineto{\pgfqpoint{4.175862in}{2.305122in}}%
\pgfpathlineto{\pgfqpoint{4.176384in}{2.304944in}}%
\pgfpathlineto{\pgfqpoint{4.176427in}{2.305139in}}%
\pgfpathlineto{\pgfqpoint{4.177058in}{2.304411in}}%
\pgfpathlineto{\pgfqpoint{4.177471in}{2.304561in}}%
\pgfpathlineto{\pgfqpoint{4.177473in}{2.304533in}}%
\pgfpathlineto{\pgfqpoint{4.177956in}{2.304936in}}%
\pgfpathlineto{\pgfqpoint{4.178574in}{2.304766in}}%
\pgfpathlineto{\pgfqpoint{4.179344in}{2.304873in}}%
\pgfpathlineto{\pgfqpoint{4.179563in}{2.304420in}}%
\pgfpathlineto{\pgfqpoint{4.179672in}{2.304599in}}%
\pgfpathlineto{\pgfqpoint{4.179803in}{2.304524in}}%
\pgfpathlineto{\pgfqpoint{4.180663in}{2.304980in}}%
\pgfpathlineto{\pgfqpoint{4.180768in}{2.304805in}}%
\pgfpathlineto{\pgfqpoint{4.180887in}{2.304871in}}%
\pgfpathlineto{\pgfqpoint{4.181728in}{2.304090in}}%
\pgfpathlineto{\pgfqpoint{4.181859in}{2.304259in}}%
\pgfpathlineto{\pgfqpoint{4.182841in}{2.303871in}}%
\pgfpathlineto{\pgfqpoint{4.182159in}{2.304505in}}%
\pgfpathlineto{\pgfqpoint{4.182997in}{2.303913in}}%
\pgfpathlineto{\pgfqpoint{4.183911in}{2.303750in}}%
\pgfpathlineto{\pgfqpoint{4.183640in}{2.304330in}}%
\pgfpathlineto{\pgfqpoint{4.184109in}{2.303962in}}%
\pgfpathlineto{\pgfqpoint{4.189226in}{2.303313in}}%
\pgfpathlineto{\pgfqpoint{4.189227in}{2.303301in}}%
\pgfpathlineto{\pgfqpoint{4.190223in}{2.302697in}}%
\pgfpathlineto{\pgfqpoint{4.189529in}{2.303582in}}%
\pgfpathlineto{\pgfqpoint{4.190340in}{2.303257in}}%
\pgfpathlineto{\pgfqpoint{4.193097in}{2.303429in}}%
\pgfpathlineto{\pgfqpoint{4.193100in}{2.303398in}}%
\pgfpathlineto{\pgfqpoint{4.194149in}{2.303359in}}%
\pgfpathlineto{\pgfqpoint{4.193110in}{2.303433in}}%
\pgfpathlineto{\pgfqpoint{4.194213in}{2.303419in}}%
\pgfpathlineto{\pgfqpoint{4.195294in}{2.303445in}}%
\pgfpathlineto{\pgfqpoint{4.194269in}{2.303362in}}%
\pgfpathlineto{\pgfqpoint{4.195324in}{2.303409in}}%
\pgfpathlineto{\pgfqpoint{4.196497in}{2.303367in}}%
\pgfpathlineto{\pgfqpoint{4.195327in}{2.303441in}}%
\pgfpathlineto{\pgfqpoint{4.196508in}{2.303418in}}%
\pgfpathlineto{\pgfqpoint{4.197603in}{2.303438in}}%
\pgfpathlineto{\pgfqpoint{4.196574in}{2.303357in}}%
\pgfpathlineto{\pgfqpoint{4.197619in}{2.303412in}}%
\pgfpathlineto{\pgfqpoint{4.198744in}{2.303361in}}%
\pgfpathlineto{\pgfqpoint{4.197636in}{2.303438in}}%
\pgfpathlineto{\pgfqpoint{4.198745in}{2.303366in}}%
\pgfpathlineto{\pgfqpoint{4.199851in}{2.303447in}}%
\pgfpathlineto{\pgfqpoint{4.198975in}{2.303346in}}%
\pgfpathlineto{\pgfqpoint{4.199859in}{2.303429in}}%
\pgfpathlineto{\pgfqpoint{4.200775in}{2.303345in}}%
\pgfpathlineto{\pgfqpoint{4.199912in}{2.303454in}}%
\pgfpathlineto{\pgfqpoint{4.200975in}{2.303415in}}%
\pgfpathlineto{\pgfqpoint{4.201977in}{2.303445in}}%
\pgfpathlineto{\pgfqpoint{4.201074in}{2.303351in}}%
\pgfpathlineto{\pgfqpoint{4.202086in}{2.303385in}}%
\pgfpathlineto{\pgfqpoint{4.203184in}{2.303437in}}%
\pgfpathlineto{\pgfqpoint{4.202532in}{2.303321in}}%
\pgfpathlineto{\pgfqpoint{4.203198in}{2.303368in}}%
\pgfpathlineto{\pgfqpoint{4.204359in}{2.303380in}}%
\pgfpathlineto{\pgfqpoint{4.203213in}{2.303440in}}%
\pgfpathlineto{\pgfqpoint{4.204365in}{2.303424in}}%
\pgfpathlineto{\pgfqpoint{4.205419in}{2.303443in}}%
\pgfpathlineto{\pgfqpoint{4.204388in}{2.303376in}}%
\pgfpathlineto{\pgfqpoint{4.205475in}{2.303427in}}%
\pgfpathlineto{\pgfqpoint{4.206553in}{2.303352in}}%
\pgfpathlineto{\pgfqpoint{4.205496in}{2.303444in}}%
\pgfpathlineto{\pgfqpoint{4.206586in}{2.303412in}}%
\pgfpathlineto{\pgfqpoint{4.207664in}{2.303449in}}%
\pgfpathlineto{\pgfqpoint{4.206615in}{2.303363in}}%
\pgfpathlineto{\pgfqpoint{4.207697in}{2.303386in}}%
\pgfpathlineto{\pgfqpoint{4.208797in}{2.303443in}}%
\pgfpathlineto{\pgfqpoint{4.207869in}{2.303338in}}%
\pgfpathlineto{\pgfqpoint{4.208819in}{2.303436in}}%
\pgfpathlineto{\pgfqpoint{4.209069in}{2.303328in}}%
\pgfpathlineto{\pgfqpoint{4.209857in}{2.303922in}}%
\pgfpathlineto{\pgfqpoint{4.209920in}{2.303808in}}%
\pgfpathlineto{\pgfqpoint{4.209984in}{2.304009in}}%
\pgfpathlineto{\pgfqpoint{4.210513in}{2.303399in}}%
\pgfpathlineto{\pgfqpoint{4.211034in}{2.303775in}}%
\pgfpathlineto{\pgfqpoint{4.211348in}{2.303572in}}%
\pgfpathlineto{\pgfqpoint{4.211861in}{2.304136in}}%
\pgfpathlineto{\pgfqpoint{4.212142in}{2.303867in}}%
\pgfpathlineto{\pgfqpoint{4.215839in}{2.303453in}}%
\pgfpathlineto{\pgfqpoint{4.215841in}{2.303494in}}%
\pgfpathlineto{\pgfqpoint{4.216949in}{2.304027in}}%
\pgfpathlineto{\pgfqpoint{4.215881in}{2.303384in}}%
\pgfpathlineto{\pgfqpoint{4.216986in}{2.303863in}}%
\pgfpathlineto{\pgfqpoint{4.217968in}{2.303696in}}%
\pgfpathlineto{\pgfqpoint{4.217559in}{2.304227in}}%
\pgfpathlineto{\pgfqpoint{4.218098in}{2.303847in}}%
\pgfpathlineto{\pgfqpoint{4.221171in}{2.303573in}}%
\pgfpathlineto{\pgfqpoint{4.221205in}{2.303725in}}%
\pgfpathlineto{\pgfqpoint{4.222053in}{2.304020in}}%
\pgfpathlineto{\pgfqpoint{4.221458in}{2.303443in}}%
\pgfpathlineto{\pgfqpoint{4.222326in}{2.303886in}}%
\pgfpathlineto{\pgfqpoint{4.222695in}{2.304131in}}%
\pgfpathlineto{\pgfqpoint{4.223371in}{2.303466in}}%
\pgfpathlineto{\pgfqpoint{4.223383in}{2.303500in}}%
\pgfpathlineto{\pgfqpoint{4.223901in}{2.303133in}}%
\pgfpathlineto{\pgfqpoint{4.223512in}{2.303660in}}%
\pgfpathlineto{\pgfqpoint{4.224498in}{2.303413in}}%
\pgfpathlineto{\pgfqpoint{4.226121in}{2.304279in}}%
\pgfpathlineto{\pgfqpoint{4.224526in}{2.303326in}}%
\pgfpathlineto{\pgfqpoint{4.226171in}{2.304125in}}%
\pgfpathlineto{\pgfqpoint{4.227101in}{2.303405in}}%
\pgfpathlineto{\pgfqpoint{4.227312in}{2.303639in}}%
\pgfpathlineto{\pgfqpoint{4.228047in}{2.304074in}}%
\pgfpathlineto{\pgfqpoint{4.227383in}{2.303486in}}%
\pgfpathlineto{\pgfqpoint{4.228432in}{2.303749in}}%
\pgfpathlineto{\pgfqpoint{4.229345in}{2.303624in}}%
\pgfpathlineto{\pgfqpoint{4.229011in}{2.303959in}}%
\pgfpathlineto{\pgfqpoint{4.229541in}{2.303809in}}%
\pgfpathlineto{\pgfqpoint{4.230408in}{2.304222in}}%
\pgfpathlineto{\pgfqpoint{4.229921in}{2.303609in}}%
\pgfpathlineto{\pgfqpoint{4.230667in}{2.303963in}}%
\pgfpathlineto{\pgfqpoint{4.231770in}{2.303673in}}%
\pgfpathlineto{\pgfqpoint{4.231828in}{2.303750in}}%
\pgfpathlineto{\pgfqpoint{4.232537in}{2.303961in}}%
\pgfpathlineto{\pgfqpoint{4.232392in}{2.303615in}}%
\pgfpathlineto{\pgfqpoint{4.232940in}{2.303769in}}%
\pgfpathlineto{\pgfqpoint{4.232986in}{2.303845in}}%
\pgfpathlineto{\pgfqpoint{4.234087in}{2.303200in}}%
\pgfpathlineto{\pgfqpoint{4.234543in}{2.303438in}}%
\pgfpathlineto{\pgfqpoint{4.235122in}{2.303027in}}%
\pgfpathlineto{\pgfqpoint{4.235196in}{2.303108in}}%
\pgfpathlineto{\pgfqpoint{4.235999in}{2.303467in}}%
\pgfpathlineto{\pgfqpoint{4.235479in}{2.303036in}}%
\pgfpathlineto{\pgfqpoint{4.236319in}{2.303237in}}%
\pgfpathlineto{\pgfqpoint{4.236451in}{2.303016in}}%
\pgfpathlineto{\pgfqpoint{4.237080in}{2.303750in}}%
\pgfpathlineto{\pgfqpoint{4.237417in}{2.303446in}}%
\pgfpathlineto{\pgfqpoint{4.237736in}{2.303828in}}%
\pgfpathlineto{\pgfqpoint{4.238188in}{2.303199in}}%
\pgfpathlineto{\pgfqpoint{4.238529in}{2.303417in}}%
\pgfpathlineto{\pgfqpoint{4.243403in}{2.303582in}}%
\pgfpathlineto{\pgfqpoint{4.243405in}{2.303606in}}%
\pgfpathlineto{\pgfqpoint{4.243740in}{2.304056in}}%
\pgfpathlineto{\pgfqpoint{4.243414in}{2.303585in}}%
\pgfpathlineto{\pgfqpoint{4.244513in}{2.303521in}}%
\pgfpathlineto{\pgfqpoint{4.246272in}{2.303588in}}%
\pgfpathlineto{\pgfqpoint{4.244567in}{2.303686in}}%
\pgfpathlineto{\pgfqpoint{4.246330in}{2.303791in}}%
\pgfpathlineto{\pgfqpoint{4.246603in}{2.303911in}}%
\pgfpathlineto{\pgfqpoint{4.247276in}{2.303475in}}%
\pgfpathlineto{\pgfqpoint{4.247429in}{2.303589in}}%
\pgfpathlineto{\pgfqpoint{4.247653in}{2.303216in}}%
\pgfpathlineto{\pgfqpoint{4.248472in}{2.304108in}}%
\pgfpathlineto{\pgfqpoint{4.248533in}{2.303930in}}%
\pgfpathlineto{\pgfqpoint{4.248562in}{2.304017in}}%
\pgfpathlineto{\pgfqpoint{4.248829in}{2.303560in}}%
\pgfpathlineto{\pgfqpoint{4.249607in}{2.303619in}}%
\pgfpathlineto{\pgfqpoint{4.250682in}{2.303364in}}%
\pgfpathlineto{\pgfqpoint{4.250141in}{2.303919in}}%
\pgfpathlineto{\pgfqpoint{4.250729in}{2.303402in}}%
\pgfpathlineto{\pgfqpoint{4.251788in}{2.303433in}}%
\pgfpathlineto{\pgfqpoint{4.251454in}{2.303109in}}%
\pgfpathlineto{\pgfqpoint{4.251846in}{2.303376in}}%
\pgfpathlineto{\pgfqpoint{4.252901in}{2.302991in}}%
\pgfpathlineto{\pgfqpoint{4.252098in}{2.303619in}}%
\pgfpathlineto{\pgfqpoint{4.252970in}{2.303156in}}%
\pgfpathlineto{\pgfqpoint{4.254099in}{2.303661in}}%
\pgfpathlineto{\pgfqpoint{4.252976in}{2.303112in}}%
\pgfpathlineto{\pgfqpoint{4.254113in}{2.303616in}}%
\pgfpathlineto{\pgfqpoint{4.257803in}{2.303189in}}%
\pgfpathlineto{\pgfqpoint{4.257811in}{2.303090in}}%
\pgfpathlineto{\pgfqpoint{4.258708in}{2.302803in}}%
\pgfpathlineto{\pgfqpoint{4.257859in}{2.303160in}}%
\pgfpathlineto{\pgfqpoint{4.258925in}{2.303042in}}%
\pgfpathlineto{\pgfqpoint{4.259190in}{2.302783in}}%
\pgfpathlineto{\pgfqpoint{4.260060in}{2.303326in}}%
\pgfpathlineto{\pgfqpoint{4.260651in}{2.303071in}}%
\pgfpathlineto{\pgfqpoint{4.261126in}{2.303752in}}%
\pgfpathlineto{\pgfqpoint{4.261155in}{2.303675in}}%
\pgfpathlineto{\pgfqpoint{4.261800in}{2.303886in}}%
\pgfpathlineto{\pgfqpoint{4.261238in}{2.303448in}}%
\pgfpathlineto{\pgfqpoint{4.262266in}{2.303683in}}%
\pgfpathlineto{\pgfqpoint{4.262537in}{2.303537in}}%
\pgfpathlineto{\pgfqpoint{4.263175in}{2.304268in}}%
\pgfpathlineto{\pgfqpoint{4.263354in}{2.304012in}}%
\pgfpathlineto{\pgfqpoint{4.263782in}{2.303905in}}%
\pgfpathlineto{\pgfqpoint{4.263960in}{2.304391in}}%
\pgfpathlineto{\pgfqpoint{4.264408in}{2.304358in}}%
\pgfpathlineto{\pgfqpoint{4.265148in}{2.304722in}}%
\pgfpathlineto{\pgfqpoint{4.265234in}{2.303375in}}%
\pgfpathlineto{\pgfqpoint{4.265428in}{2.303426in}}%
\pgfpathlineto{\pgfqpoint{4.266317in}{2.303363in}}%
\pgfpathlineto{\pgfqpoint{4.265554in}{2.303448in}}%
\pgfpathlineto{\pgfqpoint{4.266540in}{2.303421in}}%
\pgfpathlineto{\pgfqpoint{4.267671in}{2.303445in}}%
\pgfpathlineto{\pgfqpoint{4.266659in}{2.303351in}}%
\pgfpathlineto{\pgfqpoint{4.267690in}{2.303430in}}%
\pgfpathlineto{\pgfqpoint{4.268313in}{2.303348in}}%
\pgfpathlineto{\pgfqpoint{4.267730in}{2.303444in}}%
\pgfpathlineto{\pgfqpoint{4.268802in}{2.303405in}}%
\pgfpathlineto{\pgfqpoint{4.269803in}{2.303452in}}%
\pgfpathlineto{\pgfqpoint{4.268935in}{2.303368in}}%
\pgfpathlineto{\pgfqpoint{4.269916in}{2.303418in}}%
\pgfpathlineto{\pgfqpoint{4.271010in}{2.303354in}}%
\pgfpathlineto{\pgfqpoint{4.269955in}{2.303446in}}%
\pgfpathlineto{\pgfqpoint{4.271039in}{2.303420in}}%
\pgfpathlineto{\pgfqpoint{4.272145in}{2.303442in}}%
\pgfpathlineto{\pgfqpoint{4.271082in}{2.303352in}}%
\pgfpathlineto{\pgfqpoint{4.272156in}{2.303434in}}%
\pgfpathlineto{\pgfqpoint{4.272232in}{2.303335in}}%
\pgfpathlineto{\pgfqpoint{4.272250in}{2.303451in}}%
\pgfpathlineto{\pgfqpoint{4.273268in}{2.303385in}}%
\pgfpathlineto{\pgfqpoint{4.274315in}{2.303458in}}%
\pgfpathlineto{\pgfqpoint{4.273303in}{2.303356in}}%
\pgfpathlineto{\pgfqpoint{4.274382in}{2.303432in}}%
\pgfpathlineto{\pgfqpoint{4.275243in}{2.303350in}}%
\pgfpathlineto{\pgfqpoint{4.274623in}{2.303453in}}%
\pgfpathlineto{\pgfqpoint{4.275494in}{2.303398in}}%
\pgfpathlineto{\pgfqpoint{4.276585in}{2.303442in}}%
\pgfpathlineto{\pgfqpoint{4.275611in}{2.303354in}}%
\pgfpathlineto{\pgfqpoint{4.276617in}{2.303429in}}%
\pgfpathlineto{\pgfqpoint{4.277684in}{2.303382in}}%
\pgfpathlineto{\pgfqpoint{4.276705in}{2.303449in}}%
\pgfpathlineto{\pgfqpoint{4.277732in}{2.303415in}}%
\pgfpathlineto{\pgfqpoint{4.278919in}{2.303443in}}%
\pgfpathlineto{\pgfqpoint{4.277749in}{2.303350in}}%
\pgfpathlineto{\pgfqpoint{4.278919in}{2.303408in}}%
\pgfpathlineto{\pgfqpoint{4.279928in}{2.303378in}}%
\pgfpathlineto{\pgfqpoint{4.278924in}{2.303442in}}%
\pgfpathlineto{\pgfqpoint{4.280036in}{2.303420in}}%
\pgfpathlineto{\pgfqpoint{4.281133in}{2.303440in}}%
\pgfpathlineto{\pgfqpoint{4.280552in}{2.303346in}}%
\pgfpathlineto{\pgfqpoint{4.281147in}{2.303391in}}%
\pgfpathlineto{\pgfqpoint{4.281796in}{2.303340in}}%
\pgfpathlineto{\pgfqpoint{4.281152in}{2.303453in}}%
\pgfpathlineto{\pgfqpoint{4.282259in}{2.303390in}}%
\pgfpathlineto{\pgfqpoint{4.283318in}{2.303450in}}%
\pgfpathlineto{\pgfqpoint{4.282339in}{2.303340in}}%
\pgfpathlineto{\pgfqpoint{4.283374in}{2.303418in}}%
\pgfpathlineto{\pgfqpoint{4.284448in}{2.303197in}}%
\pgfpathlineto{\pgfqpoint{4.283934in}{2.303523in}}%
\pgfpathlineto{\pgfqpoint{4.284506in}{2.303264in}}%
\pgfpathlineto{\pgfqpoint{4.284539in}{2.303365in}}%
\pgfpathlineto{\pgfqpoint{4.285214in}{2.302364in}}%
\pgfpathlineto{\pgfqpoint{4.285591in}{2.302604in}}%
\pgfpathlineto{\pgfqpoint{4.286337in}{2.302250in}}%
\pgfpathlineto{\pgfqpoint{4.285592in}{2.302619in}}%
\pgfpathlineto{\pgfqpoint{4.286735in}{2.302560in}}%
\pgfpathlineto{\pgfqpoint{4.286745in}{2.302580in}}%
\pgfpathlineto{\pgfqpoint{4.287430in}{2.302203in}}%
\pgfpathlineto{\pgfqpoint{4.287834in}{2.302386in}}%
\pgfpathlineto{\pgfqpoint{4.289381in}{2.302073in}}%
\pgfpathlineto{\pgfqpoint{4.287849in}{2.302476in}}%
\pgfpathlineto{\pgfqpoint{4.289464in}{2.302332in}}%
\pgfpathlineto{\pgfqpoint{4.290038in}{2.302807in}}%
\pgfpathlineto{\pgfqpoint{4.290601in}{2.302750in}}%
\pgfpathlineto{\pgfqpoint{4.291337in}{2.302499in}}%
\pgfpathlineto{\pgfqpoint{4.291218in}{2.302910in}}%
\pgfpathlineto{\pgfqpoint{4.291732in}{2.302635in}}%
\pgfpathlineto{\pgfqpoint{4.292553in}{2.302820in}}%
\pgfpathlineto{\pgfqpoint{4.292032in}{2.302306in}}%
\pgfpathlineto{\pgfqpoint{4.292848in}{2.302683in}}%
\pgfpathlineto{\pgfqpoint{4.292924in}{2.302561in}}%
\pgfpathlineto{\pgfqpoint{4.293742in}{2.303119in}}%
\pgfpathlineto{\pgfqpoint{4.293959in}{2.302708in}}%
\pgfpathlineto{\pgfqpoint{4.297253in}{2.303527in}}%
\pgfpathlineto{\pgfqpoint{4.297256in}{2.303593in}}%
\pgfpathlineto{\pgfqpoint{4.297404in}{2.303747in}}%
\pgfpathlineto{\pgfqpoint{4.298169in}{2.303345in}}%
\pgfpathlineto{\pgfqpoint{4.298368in}{2.303584in}}%
\pgfpathlineto{\pgfqpoint{4.299203in}{2.303331in}}%
\pgfpathlineto{\pgfqpoint{4.298521in}{2.303680in}}%
\pgfpathlineto{\pgfqpoint{4.299480in}{2.303575in}}%
\pgfpathlineto{\pgfqpoint{4.300596in}{2.303886in}}%
\pgfpathlineto{\pgfqpoint{4.299591in}{2.303356in}}%
\pgfpathlineto{\pgfqpoint{4.300634in}{2.303827in}}%
\pgfpathlineto{\pgfqpoint{4.301759in}{2.303296in}}%
\pgfpathlineto{\pgfqpoint{4.301816in}{2.303391in}}%
\pgfpathlineto{\pgfqpoint{4.302927in}{2.303817in}}%
\pgfpathlineto{\pgfqpoint{4.301954in}{2.303165in}}%
\pgfpathlineto{\pgfqpoint{4.302953in}{2.303724in}}%
\pgfpathlineto{\pgfqpoint{4.303771in}{2.303500in}}%
\pgfpathlineto{\pgfqpoint{4.303323in}{2.303868in}}%
\pgfpathlineto{\pgfqpoint{4.304059in}{2.303838in}}%
\pgfpathlineto{\pgfqpoint{4.304967in}{2.303590in}}%
\pgfpathlineto{\pgfqpoint{4.304360in}{2.304077in}}%
\pgfpathlineto{\pgfqpoint{4.305183in}{2.303841in}}%
\pgfpathlineto{\pgfqpoint{4.305379in}{2.303733in}}%
\pgfpathlineto{\pgfqpoint{4.306334in}{2.304586in}}%
\pgfpathlineto{\pgfqpoint{4.307903in}{2.304796in}}%
\pgfpathlineto{\pgfqpoint{4.306352in}{2.304521in}}%
\pgfpathlineto{\pgfqpoint{4.307910in}{2.304757in}}%
\pgfpathlineto{\pgfqpoint{4.308094in}{2.304662in}}%
\pgfpathlineto{\pgfqpoint{4.308983in}{2.305360in}}%
\pgfpathlineto{\pgfqpoint{4.308990in}{2.305349in}}%
\pgfpathlineto{\pgfqpoint{4.309978in}{2.305693in}}%
\pgfpathlineto{\pgfqpoint{4.309171in}{2.305183in}}%
\pgfpathlineto{\pgfqpoint{4.310124in}{2.305613in}}%
\pgfpathlineto{\pgfqpoint{4.310155in}{2.305547in}}%
\pgfpathlineto{\pgfqpoint{4.311181in}{2.306314in}}%
\pgfpathlineto{\pgfqpoint{4.311200in}{2.306314in}}%
\pgfpathlineto{\pgfqpoint{4.312558in}{2.306467in}}%
\pgfpathlineto{\pgfqpoint{4.311207in}{2.306289in}}%
\pgfpathlineto{\pgfqpoint{4.312662in}{2.306281in}}%
\pgfpathlineto{\pgfqpoint{4.312665in}{2.306232in}}%
\pgfpathlineto{\pgfqpoint{4.313758in}{2.306799in}}%
\pgfpathlineto{\pgfqpoint{4.314452in}{2.306834in}}%
\pgfpathlineto{\pgfqpoint{4.313935in}{2.306625in}}%
\pgfpathlineto{\pgfqpoint{4.314842in}{2.306519in}}%
\pgfpathlineto{\pgfqpoint{4.315976in}{2.305997in}}%
\pgfpathlineto{\pgfqpoint{4.314858in}{2.306609in}}%
\pgfpathlineto{\pgfqpoint{4.315985in}{2.306058in}}%
\pgfpathlineto{\pgfqpoint{4.316024in}{2.306185in}}%
\pgfpathlineto{\pgfqpoint{4.316914in}{2.305321in}}%
\pgfpathlineto{\pgfqpoint{4.317027in}{2.305444in}}%
\pgfpathlineto{\pgfqpoint{4.318208in}{2.305309in}}%
\pgfpathlineto{\pgfqpoint{4.317086in}{2.305528in}}%
\pgfpathlineto{\pgfqpoint{4.318211in}{2.305317in}}%
\pgfpathlineto{\pgfqpoint{4.319379in}{2.305556in}}%
\pgfpathlineto{\pgfqpoint{4.318214in}{2.305304in}}%
\pgfpathlineto{\pgfqpoint{4.319407in}{2.305520in}}%
\pgfpathlineto{\pgfqpoint{4.320738in}{2.306284in}}%
\pgfpathlineto{\pgfqpoint{4.320866in}{2.305990in}}%
\pgfpathlineto{\pgfqpoint{4.321012in}{2.306152in}}%
\pgfpathlineto{\pgfqpoint{4.322012in}{2.305450in}}%
\pgfpathlineto{\pgfqpoint{4.322072in}{2.305552in}}%
\pgfpathlineto{\pgfqpoint{4.323099in}{2.305185in}}%
\pgfpathlineto{\pgfqpoint{4.323118in}{2.305304in}}%
\pgfpathlineto{\pgfqpoint{4.323185in}{2.305364in}}%
\pgfpathlineto{\pgfqpoint{4.323668in}{2.304725in}}%
\pgfpathlineto{\pgfqpoint{4.324195in}{2.305028in}}%
\pgfpathlineto{\pgfqpoint{4.324208in}{2.305066in}}%
\pgfpathlineto{\pgfqpoint{4.325344in}{2.304394in}}%
\pgfpathlineto{\pgfqpoint{4.326413in}{2.304604in}}%
\pgfpathlineto{\pgfqpoint{4.326110in}{2.304100in}}%
\pgfpathlineto{\pgfqpoint{4.326482in}{2.304591in}}%
\pgfpathlineto{\pgfqpoint{4.327490in}{2.303813in}}%
\pgfpathlineto{\pgfqpoint{4.326495in}{2.304625in}}%
\pgfpathlineto{\pgfqpoint{4.327653in}{2.303880in}}%
\pgfpathlineto{\pgfqpoint{4.327863in}{2.304061in}}%
\pgfpathlineto{\pgfqpoint{4.328414in}{2.303651in}}%
\pgfpathlineto{\pgfqpoint{4.328742in}{2.303673in}}%
\pgfpathlineto{\pgfqpoint{4.328963in}{2.303355in}}%
\pgfpathlineto{\pgfqpoint{4.329721in}{2.304081in}}%
\pgfpathlineto{\pgfqpoint{4.329837in}{2.303964in}}%
\pgfpathlineto{\pgfqpoint{4.330763in}{2.304293in}}%
\pgfpathlineto{\pgfqpoint{4.330106in}{2.303634in}}%
\pgfpathlineto{\pgfqpoint{4.330975in}{2.304266in}}%
\pgfpathlineto{\pgfqpoint{4.331284in}{2.304054in}}%
\pgfpathlineto{\pgfqpoint{4.331976in}{2.304771in}}%
\pgfpathlineto{\pgfqpoint{4.332060in}{2.304691in}}%
\pgfpathlineto{\pgfqpoint{4.333275in}{2.305322in}}%
\pgfpathlineto{\pgfqpoint{4.332113in}{2.304457in}}%
\pgfpathlineto{\pgfqpoint{4.333279in}{2.305293in}}%
\pgfpathlineto{\pgfqpoint{4.333817in}{2.305057in}}%
\pgfpathlineto{\pgfqpoint{4.333399in}{2.305392in}}%
\pgfpathlineto{\pgfqpoint{4.334412in}{2.305272in}}%
\pgfpathlineto{\pgfqpoint{4.334446in}{2.305379in}}%
\pgfpathlineto{\pgfqpoint{4.335101in}{2.304600in}}%
\pgfpathlineto{\pgfqpoint{4.335466in}{2.304730in}}%
\pgfpathlineto{\pgfqpoint{4.336595in}{2.304275in}}%
\pgfpathlineto{\pgfqpoint{4.335697in}{2.304889in}}%
\pgfpathlineto{\pgfqpoint{4.336600in}{2.304298in}}%
\pgfpathlineto{\pgfqpoint{4.337506in}{2.304603in}}%
\pgfpathlineto{\pgfqpoint{4.336909in}{2.304145in}}%
\pgfpathlineto{\pgfqpoint{4.337744in}{2.304412in}}%
\pgfpathlineto{\pgfqpoint{4.338795in}{2.304022in}}%
\pgfpathlineto{\pgfqpoint{4.338301in}{2.304554in}}%
\pgfpathlineto{\pgfqpoint{4.338872in}{2.304208in}}%
\pgfpathlineto{\pgfqpoint{4.339920in}{2.304217in}}%
\pgfpathlineto{\pgfqpoint{4.338999in}{2.303873in}}%
\pgfpathlineto{\pgfqpoint{4.340003in}{2.304088in}}%
\pgfpathlineto{\pgfqpoint{4.341137in}{2.303332in}}%
\pgfpathlineto{\pgfqpoint{4.340008in}{2.304108in}}%
\pgfpathlineto{\pgfqpoint{4.341248in}{2.303418in}}%
\pgfpathlineto{\pgfqpoint{4.341816in}{2.303455in}}%
\pgfpathlineto{\pgfqpoint{4.341496in}{2.303346in}}%
\pgfpathlineto{\pgfqpoint{4.342358in}{2.303426in}}%
\pgfpathlineto{\pgfqpoint{4.343383in}{2.303373in}}%
\pgfpathlineto{\pgfqpoint{4.342438in}{2.303456in}}%
\pgfpathlineto{\pgfqpoint{4.343470in}{2.303423in}}%
\pgfpathlineto{\pgfqpoint{4.344604in}{2.303443in}}%
\pgfpathlineto{\pgfqpoint{4.343536in}{2.303371in}}%
\pgfpathlineto{\pgfqpoint{4.344609in}{2.303427in}}%
\pgfpathlineto{\pgfqpoint{4.345340in}{2.303341in}}%
\pgfpathlineto{\pgfqpoint{4.344834in}{2.303448in}}%
\pgfpathlineto{\pgfqpoint{4.345721in}{2.303411in}}%
\pgfpathlineto{\pgfqpoint{4.346798in}{2.303448in}}%
\pgfpathlineto{\pgfqpoint{4.345769in}{2.303372in}}%
\pgfpathlineto{\pgfqpoint{4.346843in}{2.303412in}}%
\pgfpathlineto{\pgfqpoint{4.347994in}{2.303439in}}%
\pgfpathlineto{\pgfqpoint{4.346858in}{2.303361in}}%
\pgfpathlineto{\pgfqpoint{4.347999in}{2.303424in}}%
\pgfpathlineto{\pgfqpoint{4.349087in}{2.303336in}}%
\pgfpathlineto{\pgfqpoint{4.348025in}{2.303442in}}%
\pgfpathlineto{\pgfqpoint{4.349111in}{2.303421in}}%
\pgfpathlineto{\pgfqpoint{4.350205in}{2.303450in}}%
\pgfpathlineto{\pgfqpoint{4.349387in}{2.303358in}}%
\pgfpathlineto{\pgfqpoint{4.350224in}{2.303432in}}%
\pgfpathlineto{\pgfqpoint{4.351296in}{2.303374in}}%
\pgfpathlineto{\pgfqpoint{4.350282in}{2.303446in}}%
\pgfpathlineto{\pgfqpoint{4.351342in}{2.303413in}}%
\pgfpathlineto{\pgfqpoint{4.352136in}{2.303457in}}%
\pgfpathlineto{\pgfqpoint{4.351431in}{2.303377in}}%
\pgfpathlineto{\pgfqpoint{4.352453in}{2.303425in}}%
\pgfpathlineto{\pgfqpoint{4.353038in}{2.303357in}}%
\pgfpathlineto{\pgfqpoint{4.352459in}{2.303445in}}%
\pgfpathlineto{\pgfqpoint{4.353564in}{2.303416in}}%
\pgfpathlineto{\pgfqpoint{4.354661in}{2.303442in}}%
\pgfpathlineto{\pgfqpoint{4.353709in}{2.303357in}}%
\pgfpathlineto{\pgfqpoint{4.354684in}{2.303397in}}%
\pgfpathlineto{\pgfqpoint{4.356735in}{2.303435in}}%
\pgfpathlineto{\pgfqpoint{4.354688in}{2.303376in}}%
\pgfpathlineto{\pgfqpoint{4.356736in}{2.303410in}}%
\pgfpathlineto{\pgfqpoint{4.357793in}{2.303358in}}%
\pgfpathlineto{\pgfqpoint{4.357042in}{2.303459in}}%
\pgfpathlineto{\pgfqpoint{4.357848in}{2.303412in}}%
\pgfpathlineto{\pgfqpoint{4.358791in}{2.303493in}}%
\pgfpathlineto{\pgfqpoint{4.358940in}{2.303205in}}%
\pgfpathlineto{\pgfqpoint{4.358941in}{2.303213in}}%
\pgfpathlineto{\pgfqpoint{4.359826in}{2.302929in}}%
\pgfpathlineto{\pgfqpoint{4.359087in}{2.303453in}}%
\pgfpathlineto{\pgfqpoint{4.360055in}{2.303215in}}%
\pgfpathlineto{\pgfqpoint{4.360198in}{2.303128in}}%
\pgfpathlineto{\pgfqpoint{4.360662in}{2.303837in}}%
\pgfpathlineto{\pgfqpoint{4.361005in}{2.303794in}}%
\pgfpathlineto{\pgfqpoint{4.362114in}{2.304177in}}%
\pgfpathlineto{\pgfqpoint{4.361446in}{2.303717in}}%
\pgfpathlineto{\pgfqpoint{4.362136in}{2.304101in}}%
\pgfpathlineto{\pgfqpoint{4.362318in}{2.304242in}}%
\pgfpathlineto{\pgfqpoint{4.362900in}{2.303275in}}%
\pgfpathlineto{\pgfqpoint{4.363171in}{2.303314in}}%
\pgfpathlineto{\pgfqpoint{4.363742in}{2.302980in}}%
\pgfpathlineto{\pgfqpoint{4.364263in}{2.303506in}}%
\pgfpathlineto{\pgfqpoint{4.364271in}{2.303493in}}%
\pgfpathlineto{\pgfqpoint{4.364502in}{2.303634in}}%
\pgfpathlineto{\pgfqpoint{4.364879in}{2.303185in}}%
\pgfpathlineto{\pgfqpoint{4.365377in}{2.303393in}}%
\pgfpathlineto{\pgfqpoint{4.365391in}{2.303341in}}%
\pgfpathlineto{\pgfqpoint{4.366281in}{2.304242in}}%
\pgfpathlineto{\pgfqpoint{4.366400in}{2.304061in}}%
\pgfpathlineto{\pgfqpoint{4.366425in}{2.304090in}}%
\pgfpathlineto{\pgfqpoint{4.367295in}{2.303119in}}%
\pgfpathlineto{\pgfqpoint{4.367320in}{2.303161in}}%
\pgfpathlineto{\pgfqpoint{4.367668in}{2.302895in}}%
\pgfpathlineto{\pgfqpoint{4.368352in}{2.303536in}}%
\pgfpathlineto{\pgfqpoint{4.368425in}{2.303353in}}%
\pgfpathlineto{\pgfqpoint{4.368945in}{2.303466in}}%
\pgfpathlineto{\pgfqpoint{4.369434in}{2.302876in}}%
\pgfpathlineto{\pgfqpoint{4.369502in}{2.303001in}}%
\pgfpathlineto{\pgfqpoint{4.370548in}{2.302650in}}%
\pgfpathlineto{\pgfqpoint{4.369787in}{2.303138in}}%
\pgfpathlineto{\pgfqpoint{4.370632in}{2.302739in}}%
\pgfpathlineto{\pgfqpoint{4.371188in}{2.302877in}}%
\pgfpathlineto{\pgfqpoint{4.371624in}{2.302285in}}%
\pgfpathlineto{\pgfqpoint{4.371736in}{2.302445in}}%
\pgfpathlineto{\pgfqpoint{4.372641in}{2.301884in}}%
\pgfpathlineto{\pgfqpoint{4.371811in}{2.302572in}}%
\pgfpathlineto{\pgfqpoint{4.372864in}{2.302094in}}%
\pgfpathlineto{\pgfqpoint{4.374108in}{2.302714in}}%
\pgfpathlineto{\pgfqpoint{4.374133in}{2.302658in}}%
\pgfpathlineto{\pgfqpoint{4.374944in}{2.302091in}}%
\pgfpathlineto{\pgfqpoint{4.374145in}{2.302689in}}%
\pgfpathlineto{\pgfqpoint{4.375260in}{2.302186in}}%
\pgfpathlineto{\pgfqpoint{4.375381in}{2.302081in}}%
\pgfpathlineto{\pgfqpoint{4.376238in}{2.302859in}}%
\pgfpathlineto{\pgfqpoint{4.376292in}{2.302790in}}%
\pgfpathlineto{\pgfqpoint{4.376298in}{2.302859in}}%
\pgfpathlineto{\pgfqpoint{4.377250in}{2.301924in}}%
\pgfpathlineto{\pgfqpoint{4.377384in}{2.302007in}}%
\pgfpathlineto{\pgfqpoint{4.378533in}{2.302382in}}%
\pgfpathlineto{\pgfqpoint{4.377453in}{2.301954in}}%
\pgfpathlineto{\pgfqpoint{4.378538in}{2.302363in}}%
\pgfpathlineto{\pgfqpoint{4.379662in}{2.301584in}}%
\pgfpathlineto{\pgfqpoint{4.378619in}{2.302472in}}%
\pgfpathlineto{\pgfqpoint{4.379683in}{2.301634in}}%
\pgfpathlineto{\pgfqpoint{4.380739in}{2.301754in}}%
\pgfpathlineto{\pgfqpoint{4.379927in}{2.301346in}}%
\pgfpathlineto{\pgfqpoint{4.380799in}{2.301643in}}%
\pgfpathlineto{\pgfqpoint{4.381146in}{2.301510in}}%
\pgfpathlineto{\pgfqpoint{4.381562in}{2.301945in}}%
\pgfpathlineto{\pgfqpoint{4.381923in}{2.301680in}}%
\pgfpathlineto{\pgfqpoint{4.382595in}{2.301824in}}%
\pgfpathlineto{\pgfqpoint{4.382326in}{2.301510in}}%
\pgfpathlineto{\pgfqpoint{4.383037in}{2.301763in}}%
\pgfpathlineto{\pgfqpoint{4.384165in}{2.301591in}}%
\pgfpathlineto{\pgfqpoint{4.383454in}{2.301891in}}%
\pgfpathlineto{\pgfqpoint{4.384172in}{2.301615in}}%
\pgfpathlineto{\pgfqpoint{4.384916in}{2.302032in}}%
\pgfpathlineto{\pgfqpoint{4.384180in}{2.301573in}}%
\pgfpathlineto{\pgfqpoint{4.385298in}{2.301862in}}%
\pgfpathlineto{\pgfqpoint{4.386619in}{2.301669in}}%
\pgfpathlineto{\pgfqpoint{4.385359in}{2.301992in}}%
\pgfpathlineto{\pgfqpoint{4.386665in}{2.301901in}}%
\pgfpathlineto{\pgfqpoint{4.387590in}{2.302289in}}%
\pgfpathlineto{\pgfqpoint{4.387277in}{2.301828in}}%
\pgfpathlineto{\pgfqpoint{4.387786in}{2.302198in}}%
\pgfpathlineto{\pgfqpoint{4.388906in}{2.301832in}}%
\pgfpathlineto{\pgfqpoint{4.388053in}{2.302596in}}%
\pgfpathlineto{\pgfqpoint{4.388920in}{2.301850in}}%
\pgfpathlineto{\pgfqpoint{4.389954in}{2.302053in}}%
\pgfpathlineto{\pgfqpoint{4.389190in}{2.302200in}}%
\pgfpathlineto{\pgfqpoint{4.389999in}{2.302197in}}%
\pgfpathlineto{\pgfqpoint{4.390802in}{2.302457in}}%
\pgfpathlineto{\pgfqpoint{4.390123in}{2.302151in}}%
\pgfpathlineto{\pgfqpoint{4.391127in}{2.302389in}}%
\pgfpathlineto{\pgfqpoint{4.391368in}{2.302070in}}%
\pgfpathlineto{\pgfqpoint{4.391909in}{2.302663in}}%
\pgfpathlineto{\pgfqpoint{4.392227in}{2.302630in}}%
\pgfpathlineto{\pgfqpoint{4.393320in}{2.302713in}}%
\pgfpathlineto{\pgfqpoint{4.392630in}{2.302309in}}%
\pgfpathlineto{\pgfqpoint{4.393338in}{2.302605in}}%
\pgfpathlineto{\pgfqpoint{4.393433in}{2.302352in}}%
\pgfpathlineto{\pgfqpoint{4.394265in}{2.303021in}}%
\pgfpathlineto{\pgfqpoint{4.394435in}{2.302942in}}%
\pgfpathlineto{\pgfqpoint{4.395134in}{2.303397in}}%
\pgfpathlineto{\pgfqpoint{4.394589in}{2.302898in}}%
\pgfpathlineto{\pgfqpoint{4.395607in}{2.303147in}}%
\pgfpathlineto{\pgfqpoint{4.397024in}{2.303524in}}%
\pgfpathlineto{\pgfqpoint{4.395627in}{2.303105in}}%
\pgfpathlineto{\pgfqpoint{4.397067in}{2.303401in}}%
\pgfpathlineto{\pgfqpoint{4.397617in}{2.303247in}}%
\pgfpathlineto{\pgfqpoint{4.398138in}{2.303737in}}%
\pgfpathlineto{\pgfqpoint{4.398146in}{2.303706in}}%
\pgfpathlineto{\pgfqpoint{4.399222in}{2.303696in}}%
\pgfpathlineto{\pgfqpoint{4.398486in}{2.303443in}}%
\pgfpathlineto{\pgfqpoint{4.399262in}{2.303657in}}%
\pgfpathlineto{\pgfqpoint{4.400362in}{2.303436in}}%
\pgfpathlineto{\pgfqpoint{4.399932in}{2.304004in}}%
\pgfpathlineto{\pgfqpoint{4.400382in}{2.303466in}}%
\pgfpathlineto{\pgfqpoint{4.401317in}{2.303701in}}%
\pgfpathlineto{\pgfqpoint{4.401122in}{2.303401in}}%
\pgfpathlineto{\pgfqpoint{4.401507in}{2.303603in}}%
\pgfpathlineto{\pgfqpoint{4.401547in}{2.303451in}}%
\pgfpathlineto{\pgfqpoint{4.402451in}{2.304100in}}%
\pgfpathlineto{\pgfqpoint{4.402614in}{2.303735in}}%
\pgfpathlineto{\pgfqpoint{4.403148in}{2.303881in}}%
\pgfpathlineto{\pgfqpoint{4.403665in}{2.303361in}}%
\pgfpathlineto{\pgfqpoint{4.403682in}{2.303408in}}%
\pgfpathlineto{\pgfqpoint{4.404503in}{2.303086in}}%
\pgfpathlineto{\pgfqpoint{4.403793in}{2.303812in}}%
\pgfpathlineto{\pgfqpoint{4.404795in}{2.303376in}}%
\pgfpathlineto{\pgfqpoint{4.405133in}{2.303691in}}%
\pgfpathlineto{\pgfqpoint{4.405486in}{2.303116in}}%
\pgfpathlineto{\pgfqpoint{4.405900in}{2.303266in}}%
\pgfpathlineto{\pgfqpoint{4.407005in}{2.302611in}}%
\pgfpathlineto{\pgfqpoint{4.406315in}{2.303415in}}%
\pgfpathlineto{\pgfqpoint{4.407046in}{2.302719in}}%
\pgfpathlineto{\pgfqpoint{4.407133in}{2.302900in}}%
\pgfpathlineto{\pgfqpoint{4.408033in}{2.302007in}}%
\pgfpathlineto{\pgfqpoint{4.408098in}{2.302135in}}%
\pgfpathlineto{\pgfqpoint{4.408883in}{2.301953in}}%
\pgfpathlineto{\pgfqpoint{4.408378in}{2.302326in}}%
\pgfpathlineto{\pgfqpoint{4.409216in}{2.301994in}}%
\pgfpathlineto{\pgfqpoint{4.409664in}{2.302256in}}%
\pgfpathlineto{\pgfqpoint{4.409906in}{2.301841in}}%
\pgfpathlineto{\pgfqpoint{4.410327in}{2.301967in}}%
\pgfpathlineto{\pgfqpoint{4.410454in}{2.301760in}}%
\pgfpathlineto{\pgfqpoint{4.410967in}{2.302406in}}%
\pgfpathlineto{\pgfqpoint{4.411427in}{2.302148in}}%
\pgfpathlineto{\pgfqpoint{4.433199in}{2.303428in}}%
\pgfpathlineto{\pgfqpoint{4.433202in}{2.303385in}}%
\pgfpathlineto{\pgfqpoint{4.434111in}{2.303309in}}%
\pgfpathlineto{\pgfqpoint{4.433920in}{2.303569in}}%
\pgfpathlineto{\pgfqpoint{4.434317in}{2.303321in}}%
\pgfpathlineto{\pgfqpoint{4.434365in}{2.303494in}}%
\pgfpathlineto{\pgfqpoint{4.435282in}{2.302828in}}%
\pgfpathlineto{\pgfqpoint{4.435381in}{2.302888in}}%
\pgfpathlineto{\pgfqpoint{4.436553in}{2.301670in}}%
\pgfpathlineto{\pgfqpoint{4.435532in}{2.303022in}}%
\pgfpathlineto{\pgfqpoint{4.436615in}{2.301883in}}%
\pgfpathlineto{\pgfqpoint{4.436740in}{2.302335in}}%
\pgfpathlineto{\pgfqpoint{4.437582in}{2.301137in}}%
\pgfpathlineto{\pgfqpoint{4.437696in}{2.301286in}}%
\pgfpathlineto{\pgfqpoint{4.438139in}{2.300708in}}%
\pgfpathlineto{\pgfqpoint{4.437715in}{2.301337in}}%
\pgfpathlineto{\pgfqpoint{4.438815in}{2.301357in}}%
\pgfpathlineto{\pgfqpoint{4.439104in}{2.301456in}}%
\pgfpathlineto{\pgfqpoint{4.439676in}{2.301033in}}%
\pgfpathlineto{\pgfqpoint{4.439925in}{2.301305in}}%
\pgfpathlineto{\pgfqpoint{4.440823in}{2.300544in}}%
\pgfpathlineto{\pgfqpoint{4.441068in}{2.300812in}}%
\pgfpathlineto{\pgfqpoint{4.441959in}{2.301286in}}%
\pgfpathlineto{\pgfqpoint{4.441130in}{2.300669in}}%
\pgfpathlineto{\pgfqpoint{4.442196in}{2.301137in}}%
\pgfpathlineto{\pgfqpoint{4.444537in}{2.300372in}}%
\pgfpathlineto{\pgfqpoint{4.442199in}{2.301089in}}%
\pgfpathlineto{\pgfqpoint{4.444541in}{2.300322in}}%
\pgfpathlineto{\pgfqpoint{4.445323in}{2.300042in}}%
\pgfpathlineto{\pgfqpoint{4.444855in}{2.300591in}}%
\pgfpathlineto{\pgfqpoint{4.445662in}{2.300218in}}%
\pgfpathlineto{\pgfqpoint{4.446881in}{2.300460in}}%
\pgfpathlineto{\pgfqpoint{4.445687in}{2.300169in}}%
\pgfpathlineto{\pgfqpoint{4.446901in}{2.300371in}}%
\pgfpathlineto{\pgfqpoint{4.447937in}{2.299809in}}%
\pgfpathlineto{\pgfqpoint{4.448059in}{2.300134in}}%
\pgfpathlineto{\pgfqpoint{4.448095in}{2.300211in}}%
\pgfpathlineto{\pgfqpoint{4.448547in}{2.299582in}}%
\pgfpathlineto{\pgfqpoint{4.449154in}{2.299876in}}%
\pgfpathlineto{\pgfqpoint{4.450218in}{2.299140in}}%
\pgfpathlineto{\pgfqpoint{4.449167in}{2.299907in}}%
\pgfpathlineto{\pgfqpoint{4.450296in}{2.299250in}}%
\pgfpathlineto{\pgfqpoint{4.451160in}{2.299342in}}%
\pgfpathlineto{\pgfqpoint{4.450536in}{2.298930in}}%
\pgfpathlineto{\pgfqpoint{4.451405in}{2.299196in}}%
\pgfpathlineto{\pgfqpoint{4.452121in}{2.299033in}}%
\pgfpathlineto{\pgfqpoint{4.451635in}{2.299515in}}%
\pgfpathlineto{\pgfqpoint{4.452518in}{2.299178in}}%
\pgfpathlineto{\pgfqpoint{4.452974in}{2.299455in}}%
\pgfpathlineto{\pgfqpoint{4.452611in}{2.298987in}}%
\pgfpathlineto{\pgfqpoint{4.453633in}{2.299198in}}%
\pgfpathlineto{\pgfqpoint{4.454715in}{2.298811in}}%
\pgfpathlineto{\pgfqpoint{4.454009in}{2.299328in}}%
\pgfpathlineto{\pgfqpoint{4.454860in}{2.298886in}}%
\pgfpathlineto{\pgfqpoint{4.455770in}{2.299225in}}%
\pgfpathlineto{\pgfqpoint{4.454909in}{2.298839in}}%
\pgfpathlineto{\pgfqpoint{4.455979in}{2.299049in}}%
\pgfpathlineto{\pgfqpoint{4.457196in}{2.298325in}}%
\pgfpathlineto{\pgfqpoint{4.455985in}{2.299077in}}%
\pgfpathlineto{\pgfqpoint{4.457215in}{2.298367in}}%
\pgfpathlineto{\pgfqpoint{4.457650in}{2.298545in}}%
\pgfpathlineto{\pgfqpoint{4.458208in}{2.298134in}}%
\pgfpathlineto{\pgfqpoint{4.458325in}{2.298300in}}%
\pgfpathlineto{\pgfqpoint{4.459515in}{2.298978in}}%
\pgfpathlineto{\pgfqpoint{4.458671in}{2.298051in}}%
\pgfpathlineto{\pgfqpoint{4.459622in}{2.298806in}}%
\pgfpathlineto{\pgfqpoint{4.459903in}{2.298579in}}%
\pgfpathlineto{\pgfqpoint{4.460501in}{2.299422in}}%
\pgfpathlineto{\pgfqpoint{4.460711in}{2.299201in}}%
\pgfpathlineto{\pgfqpoint{4.461336in}{2.299643in}}%
\pgfpathlineto{\pgfqpoint{4.460786in}{2.299075in}}%
\pgfpathlineto{\pgfqpoint{4.461833in}{2.299228in}}%
\pgfpathlineto{\pgfqpoint{4.462381in}{2.298723in}}%
\pgfpathlineto{\pgfqpoint{4.462165in}{2.299251in}}%
\pgfpathlineto{\pgfqpoint{4.462961in}{2.298947in}}%
\pgfpathlineto{\pgfqpoint{4.463766in}{2.299082in}}%
\pgfpathlineto{\pgfqpoint{4.462989in}{2.298845in}}%
\pgfpathlineto{\pgfqpoint{4.464078in}{2.298883in}}%
\pgfpathlineto{\pgfqpoint{4.465240in}{2.298540in}}%
\pgfpathlineto{\pgfqpoint{4.464291in}{2.299067in}}%
\pgfpathlineto{\pgfqpoint{4.465253in}{2.298578in}}%
\pgfpathlineto{\pgfqpoint{4.465926in}{2.299188in}}%
\pgfpathlineto{\pgfqpoint{4.465479in}{2.298568in}}%
\pgfpathlineto{\pgfqpoint{4.466379in}{2.299043in}}%
\pgfpathlineto{\pgfqpoint{4.467098in}{2.298734in}}%
\pgfpathlineto{\pgfqpoint{4.467379in}{2.299151in}}%
\pgfpathlineto{\pgfqpoint{4.467489in}{2.299066in}}%
\pgfpathlineto{\pgfqpoint{4.467952in}{2.299208in}}%
\pgfpathlineto{\pgfqpoint{4.468359in}{2.298723in}}%
\pgfpathlineto{\pgfqpoint{4.468583in}{2.298878in}}%
\pgfpathlineto{\pgfqpoint{4.469007in}{2.299028in}}%
\pgfpathlineto{\pgfqpoint{4.469715in}{2.298414in}}%
\pgfpathlineto{\pgfqpoint{4.469717in}{2.298401in}}%
\pgfpathlineto{\pgfqpoint{4.470716in}{2.298915in}}%
\pgfpathlineto{\pgfqpoint{4.470798in}{2.298806in}}%
\pgfpathlineto{\pgfqpoint{4.471135in}{2.298828in}}%
\pgfpathlineto{\pgfqpoint{4.471755in}{2.298480in}}%
\pgfpathlineto{\pgfqpoint{4.471907in}{2.298726in}}%
\pgfpathlineto{\pgfqpoint{4.473003in}{2.299055in}}%
\pgfpathlineto{\pgfqpoint{4.471983in}{2.298620in}}%
\pgfpathlineto{\pgfqpoint{4.473037in}{2.298995in}}%
\pgfpathlineto{\pgfqpoint{4.474361in}{2.298648in}}%
\pgfpathlineto{\pgfqpoint{4.473069in}{2.299084in}}%
\pgfpathlineto{\pgfqpoint{4.474401in}{2.298758in}}%
\pgfpathlineto{\pgfqpoint{4.474913in}{2.298973in}}%
\pgfpathlineto{\pgfqpoint{4.475424in}{2.298295in}}%
\pgfpathlineto{\pgfqpoint{4.475494in}{2.298376in}}%
\pgfpathlineto{\pgfqpoint{4.476672in}{2.297943in}}%
\pgfpathlineto{\pgfqpoint{4.475599in}{2.298496in}}%
\pgfpathlineto{\pgfqpoint{4.476707in}{2.298116in}}%
\pgfpathlineto{\pgfqpoint{4.477333in}{2.298563in}}%
\pgfpathlineto{\pgfqpoint{4.477741in}{2.297933in}}%
\pgfpathlineto{\pgfqpoint{4.477822in}{2.298127in}}%
\pgfpathlineto{\pgfqpoint{4.477845in}{2.298004in}}%
\pgfpathlineto{\pgfqpoint{4.478820in}{2.298697in}}%
\pgfpathlineto{\pgfqpoint{4.478918in}{2.298412in}}%
\pgfpathlineto{\pgfqpoint{4.480766in}{2.298697in}}%
\pgfpathlineto{\pgfqpoint{4.478920in}{2.298429in}}%
\pgfpathlineto{\pgfqpoint{4.480809in}{2.298865in}}%
\pgfpathlineto{\pgfqpoint{4.481374in}{2.299128in}}%
\pgfpathlineto{\pgfqpoint{4.481895in}{2.298595in}}%
\pgfpathlineto{\pgfqpoint{4.482020in}{2.298440in}}%
\pgfpathlineto{\pgfqpoint{4.482858in}{2.299264in}}%
\pgfpathlineto{\pgfqpoint{4.482991in}{2.298855in}}%
\pgfpathlineto{\pgfqpoint{4.482994in}{2.298889in}}%
\pgfpathlineto{\pgfqpoint{4.483376in}{2.298320in}}%
\pgfpathlineto{\pgfqpoint{4.484072in}{2.298540in}}%
\pgfpathlineto{\pgfqpoint{4.484120in}{2.298450in}}%
\pgfpathlineto{\pgfqpoint{4.485108in}{2.299194in}}%
\pgfpathlineto{\pgfqpoint{4.485157in}{2.299021in}}%
\pgfpathlineto{\pgfqpoint{4.485593in}{2.299320in}}%
\pgfpathlineto{\pgfqpoint{4.486089in}{2.298853in}}%
\pgfpathlineto{\pgfqpoint{4.486288in}{2.299047in}}%
\pgfpathlineto{\pgfqpoint{4.487046in}{2.298396in}}%
\pgfpathlineto{\pgfqpoint{4.486333in}{2.299211in}}%
\pgfpathlineto{\pgfqpoint{4.487439in}{2.298758in}}%
\pgfpathlineto{\pgfqpoint{4.487681in}{2.299051in}}%
\pgfpathlineto{\pgfqpoint{4.488198in}{2.298500in}}%
\pgfpathlineto{\pgfqpoint{4.488540in}{2.298609in}}%
\pgfpathlineto{\pgfqpoint{4.489062in}{2.298452in}}%
\pgfpathlineto{\pgfqpoint{4.489573in}{2.299011in}}%
\pgfpathlineto{\pgfqpoint{4.489629in}{2.298850in}}%
\pgfpathlineto{\pgfqpoint{4.490973in}{2.303457in}}%
\pgfpathlineto{\pgfqpoint{4.489849in}{2.298836in}}%
\pgfpathlineto{\pgfqpoint{4.490991in}{2.303410in}}%
\pgfpathlineto{\pgfqpoint{4.492027in}{2.303369in}}%
\pgfpathlineto{\pgfqpoint{4.491049in}{2.303443in}}%
\pgfpathlineto{\pgfqpoint{4.492112in}{2.303407in}}%
\pgfpathlineto{\pgfqpoint{4.493467in}{2.303394in}}%
\pgfpathlineto{\pgfqpoint{4.492120in}{2.303431in}}%
\pgfpathlineto{\pgfqpoint{4.493468in}{2.303416in}}%
\pgfpathlineto{\pgfqpoint{4.494545in}{2.303442in}}%
\pgfpathlineto{\pgfqpoint{4.493482in}{2.303378in}}%
\pgfpathlineto{\pgfqpoint{4.494581in}{2.303421in}}%
\pgfpathlineto{\pgfqpoint{4.495237in}{2.303364in}}%
\pgfpathlineto{\pgfqpoint{4.494587in}{2.303444in}}%
\pgfpathlineto{\pgfqpoint{4.495693in}{2.303405in}}%
\pgfpathlineto{\pgfqpoint{4.496799in}{2.303440in}}%
\pgfpathlineto{\pgfqpoint{4.495760in}{2.303363in}}%
\pgfpathlineto{\pgfqpoint{4.496810in}{2.303409in}}%
\pgfpathlineto{\pgfqpoint{4.510731in}{2.303093in}}%
\pgfpathlineto{\pgfqpoint{4.510732in}{2.303078in}}%
\pgfpathlineto{\pgfqpoint{4.511520in}{2.302944in}}%
\pgfpathlineto{\pgfqpoint{4.511803in}{2.303320in}}%
\pgfpathlineto{\pgfqpoint{4.511838in}{2.303216in}}%
\pgfpathlineto{\pgfqpoint{4.516179in}{2.303744in}}%
\pgfpathlineto{\pgfqpoint{4.516207in}{2.303599in}}%
\pgfpathlineto{\pgfqpoint{4.517275in}{2.303110in}}%
\pgfpathlineto{\pgfqpoint{4.516275in}{2.303672in}}%
\pgfpathlineto{\pgfqpoint{4.517342in}{2.303220in}}%
\pgfpathlineto{\pgfqpoint{4.518191in}{2.303538in}}%
\pgfpathlineto{\pgfqpoint{4.517584in}{2.303074in}}%
\pgfpathlineto{\pgfqpoint{4.518462in}{2.303358in}}%
\pgfpathlineto{\pgfqpoint{4.518462in}{2.303340in}}%
\pgfpathlineto{\pgfqpoint{4.519450in}{2.304226in}}%
\pgfpathlineto{\pgfqpoint{4.519557in}{2.303762in}}%
\pgfpathlineto{\pgfqpoint{4.521727in}{2.303436in}}%
\pgfpathlineto{\pgfqpoint{4.519617in}{2.303926in}}%
\pgfpathlineto{\pgfqpoint{4.521870in}{2.303778in}}%
\pgfpathlineto{\pgfqpoint{4.522986in}{2.304069in}}%
\pgfpathlineto{\pgfqpoint{4.522055in}{2.303458in}}%
\pgfpathlineto{\pgfqpoint{4.523006in}{2.304022in}}%
\pgfpathlineto{\pgfqpoint{4.523019in}{2.303990in}}%
\pgfpathlineto{\pgfqpoint{4.523717in}{2.304533in}}%
\pgfpathlineto{\pgfqpoint{4.524101in}{2.304463in}}%
\pgfpathlineto{\pgfqpoint{4.524476in}{2.304556in}}%
\pgfpathlineto{\pgfqpoint{4.524950in}{2.303742in}}%
\pgfpathlineto{\pgfqpoint{4.525081in}{2.303903in}}%
\pgfpathlineto{\pgfqpoint{4.525084in}{2.303873in}}%
\pgfpathlineto{\pgfqpoint{4.526168in}{2.304424in}}%
\pgfpathlineto{\pgfqpoint{4.526177in}{2.304340in}}%
\pgfpathlineto{\pgfqpoint{4.526618in}{2.304242in}}%
\pgfpathlineto{\pgfqpoint{4.527069in}{2.304761in}}%
\pgfpathlineto{\pgfqpoint{4.527292in}{2.304481in}}%
\pgfpathlineto{\pgfqpoint{4.528314in}{2.304860in}}%
\pgfpathlineto{\pgfqpoint{4.527323in}{2.304446in}}%
\pgfpathlineto{\pgfqpoint{4.528426in}{2.304791in}}%
\pgfpathlineto{\pgfqpoint{4.528500in}{2.304704in}}%
\pgfpathlineto{\pgfqpoint{4.529486in}{2.305415in}}%
\pgfpathlineto{\pgfqpoint{4.529513in}{2.305313in}}%
\pgfpathlineto{\pgfqpoint{4.529641in}{2.305505in}}%
\pgfpathlineto{\pgfqpoint{4.530036in}{2.304938in}}%
\pgfpathlineto{\pgfqpoint{4.530530in}{2.304869in}}%
\pgfpathlineto{\pgfqpoint{4.531652in}{2.304531in}}%
\pgfpathlineto{\pgfqpoint{4.530699in}{2.305076in}}%
\pgfpathlineto{\pgfqpoint{4.531659in}{2.304561in}}%
\pgfpathlineto{\pgfqpoint{4.532258in}{2.304334in}}%
\pgfpathlineto{\pgfqpoint{4.532790in}{2.304654in}}%
\pgfpathlineto{\pgfqpoint{4.533572in}{2.304822in}}%
\pgfpathlineto{\pgfqpoint{4.533195in}{2.304443in}}%
\pgfpathlineto{\pgfqpoint{4.533911in}{2.304795in}}%
\pgfpathlineto{\pgfqpoint{4.535261in}{2.304102in}}%
\pgfpathlineto{\pgfqpoint{4.534129in}{2.304937in}}%
\pgfpathlineto{\pgfqpoint{4.535334in}{2.304207in}}%
\pgfpathlineto{\pgfqpoint{4.535567in}{2.304795in}}%
\pgfpathlineto{\pgfqpoint{4.536466in}{2.304628in}}%
\pgfpathlineto{\pgfqpoint{4.537302in}{2.305300in}}%
\pgfpathlineto{\pgfqpoint{4.536476in}{2.304602in}}%
\pgfpathlineto{\pgfqpoint{4.537735in}{2.305062in}}%
\pgfpathlineto{\pgfqpoint{4.538774in}{2.304428in}}%
\pgfpathlineto{\pgfqpoint{4.537801in}{2.305157in}}%
\pgfpathlineto{\pgfqpoint{4.538951in}{2.304618in}}%
\pgfpathlineto{\pgfqpoint{4.538960in}{2.304692in}}%
\pgfpathlineto{\pgfqpoint{4.539427in}{2.304146in}}%
\pgfpathlineto{\pgfqpoint{4.540049in}{2.304236in}}%
\pgfpathlineto{\pgfqpoint{4.541183in}{2.303602in}}%
\pgfpathlineto{\pgfqpoint{4.540110in}{2.304285in}}%
\pgfpathlineto{\pgfqpoint{4.541214in}{2.303669in}}%
\pgfpathlineto{\pgfqpoint{4.542325in}{2.303891in}}%
\pgfpathlineto{\pgfqpoint{4.541862in}{2.303473in}}%
\pgfpathlineto{\pgfqpoint{4.542331in}{2.303830in}}%
\pgfpathlineto{\pgfqpoint{4.542969in}{2.303469in}}%
\pgfpathlineto{\pgfqpoint{4.542744in}{2.304058in}}%
\pgfpathlineto{\pgfqpoint{4.543460in}{2.303588in}}%
\pgfpathlineto{\pgfqpoint{4.544206in}{2.303485in}}%
\pgfpathlineto{\pgfqpoint{4.543612in}{2.303742in}}%
\pgfpathlineto{\pgfqpoint{4.544552in}{2.303812in}}%
\pgfpathlineto{\pgfqpoint{4.545456in}{2.304185in}}%
\pgfpathlineto{\pgfqpoint{4.544878in}{2.303657in}}%
\pgfpathlineto{\pgfqpoint{4.545682in}{2.304013in}}%
\pgfpathlineto{\pgfqpoint{4.546772in}{2.303298in}}%
\pgfpathlineto{\pgfqpoint{4.545711in}{2.304069in}}%
\pgfpathlineto{\pgfqpoint{4.546878in}{2.303558in}}%
\pgfpathlineto{\pgfqpoint{4.547029in}{2.303757in}}%
\pgfpathlineto{\pgfqpoint{4.547342in}{2.303139in}}%
\pgfpathlineto{\pgfqpoint{4.547981in}{2.303387in}}%
\pgfpathlineto{\pgfqpoint{4.548596in}{2.303688in}}%
\pgfpathlineto{\pgfqpoint{4.548035in}{2.303190in}}%
\pgfpathlineto{\pgfqpoint{4.549118in}{2.302975in}}%
\pgfpathlineto{\pgfqpoint{4.549674in}{2.302586in}}%
\pgfpathlineto{\pgfqpoint{4.549229in}{2.303031in}}%
\pgfpathlineto{\pgfqpoint{4.550242in}{2.302766in}}%
\pgfpathlineto{\pgfqpoint{4.550250in}{2.302786in}}%
\pgfpathlineto{\pgfqpoint{4.551306in}{2.301931in}}%
\pgfpathlineto{\pgfqpoint{4.551308in}{2.301950in}}%
\pgfpathlineto{\pgfqpoint{4.551886in}{2.301536in}}%
\pgfpathlineto{\pgfqpoint{4.551609in}{2.301990in}}%
\pgfpathlineto{\pgfqpoint{4.552425in}{2.301852in}}%
\pgfpathlineto{\pgfqpoint{4.552872in}{2.302061in}}%
\pgfpathlineto{\pgfqpoint{4.552517in}{2.301710in}}%
\pgfpathlineto{\pgfqpoint{4.553540in}{2.301894in}}%
\pgfpathlineto{\pgfqpoint{4.554124in}{2.301734in}}%
\pgfpathlineto{\pgfqpoint{4.554440in}{2.302267in}}%
\pgfpathlineto{\pgfqpoint{4.554575in}{2.302293in}}%
\pgfpathlineto{\pgfqpoint{4.555408in}{2.302482in}}%
\pgfpathlineto{\pgfqpoint{4.554648in}{2.302111in}}%
\pgfpathlineto{\pgfqpoint{4.555689in}{2.302384in}}%
\pgfpathlineto{\pgfqpoint{4.556357in}{2.302348in}}%
\pgfpathlineto{\pgfqpoint{4.556018in}{2.302761in}}%
\pgfpathlineto{\pgfqpoint{4.556784in}{2.302638in}}%
\pgfpathlineto{\pgfqpoint{4.557855in}{2.302496in}}%
\pgfpathlineto{\pgfqpoint{4.556818in}{2.302725in}}%
\pgfpathlineto{\pgfqpoint{4.558039in}{2.302822in}}%
\pgfpathlineto{\pgfqpoint{4.558858in}{2.303071in}}%
\pgfpathlineto{\pgfqpoint{4.558284in}{2.302713in}}%
\pgfpathlineto{\pgfqpoint{4.559156in}{2.302892in}}%
\pgfpathlineto{\pgfqpoint{4.560248in}{2.303234in}}%
\pgfpathlineto{\pgfqpoint{4.559311in}{2.302512in}}%
\pgfpathlineto{\pgfqpoint{4.560300in}{2.303124in}}%
\pgfpathlineto{\pgfqpoint{4.560302in}{2.303090in}}%
\pgfpathlineto{\pgfqpoint{4.561368in}{2.303720in}}%
\pgfpathlineto{\pgfqpoint{4.561385in}{2.303672in}}%
\pgfpathlineto{\pgfqpoint{4.561899in}{2.303973in}}%
\pgfpathlineto{\pgfqpoint{4.561421in}{2.303514in}}%
\pgfpathlineto{\pgfqpoint{4.562504in}{2.303628in}}%
\pgfpathlineto{\pgfqpoint{4.563068in}{2.303204in}}%
\pgfpathlineto{\pgfqpoint{4.562519in}{2.303661in}}%
\pgfpathlineto{\pgfqpoint{4.563667in}{2.303406in}}%
\pgfpathlineto{\pgfqpoint{4.564452in}{2.303552in}}%
\pgfpathlineto{\pgfqpoint{4.563925in}{2.303226in}}%
\pgfpathlineto{\pgfqpoint{4.564778in}{2.303410in}}%
\pgfpathlineto{\pgfqpoint{4.565877in}{2.303355in}}%
\pgfpathlineto{\pgfqpoint{4.564799in}{2.303442in}}%
\pgfpathlineto{\pgfqpoint{4.565942in}{2.303416in}}%
\pgfpathlineto{\pgfqpoint{4.567037in}{2.303448in}}%
\pgfpathlineto{\pgfqpoint{4.566023in}{2.303319in}}%
\pgfpathlineto{\pgfqpoint{4.567056in}{2.303420in}}%
\pgfpathlineto{\pgfqpoint{4.567820in}{2.303333in}}%
\pgfpathlineto{\pgfqpoint{4.567236in}{2.303459in}}%
\pgfpathlineto{\pgfqpoint{4.568168in}{2.303418in}}%
\pgfpathlineto{\pgfqpoint{4.570076in}{2.303404in}}%
\pgfpathlineto{\pgfqpoint{4.568172in}{2.303430in}}%
\pgfpathlineto{\pgfqpoint{4.570078in}{2.303423in}}%
\pgfpathlineto{\pgfqpoint{4.571068in}{2.303450in}}%
\pgfpathlineto{\pgfqpoint{4.570218in}{2.303347in}}%
\pgfpathlineto{\pgfqpoint{4.571188in}{2.303423in}}%
\pgfpathlineto{\pgfqpoint{4.572107in}{2.303366in}}%
\pgfpathlineto{\pgfqpoint{4.571263in}{2.303447in}}%
\pgfpathlineto{\pgfqpoint{4.572301in}{2.303415in}}%
\pgfpathlineto{\pgfqpoint{4.573377in}{2.303340in}}%
\pgfpathlineto{\pgfqpoint{4.572318in}{2.303452in}}%
\pgfpathlineto{\pgfqpoint{4.573426in}{2.303419in}}%
\pgfpathlineto{\pgfqpoint{4.574312in}{2.303451in}}%
\pgfpathlineto{\pgfqpoint{4.573511in}{2.303364in}}%
\pgfpathlineto{\pgfqpoint{4.574539in}{2.303422in}}%
\pgfpathlineto{\pgfqpoint{4.575668in}{2.303447in}}%
\pgfpathlineto{\pgfqpoint{4.574619in}{2.303357in}}%
\pgfpathlineto{\pgfqpoint{4.575686in}{2.303395in}}%
\pgfpathlineto{\pgfqpoint{4.576786in}{2.303370in}}%
\pgfpathlineto{\pgfqpoint{4.575721in}{2.303456in}}%
\pgfpathlineto{\pgfqpoint{4.576798in}{2.303392in}}%
\pgfpathlineto{\pgfqpoint{4.578457in}{2.303403in}}%
\pgfpathlineto{\pgfqpoint{4.576804in}{2.303426in}}%
\pgfpathlineto{\pgfqpoint{4.578457in}{2.303412in}}%
\pgfpathlineto{\pgfqpoint{4.579476in}{2.303446in}}%
\pgfpathlineto{\pgfqpoint{4.578630in}{2.303366in}}%
\pgfpathlineto{\pgfqpoint{4.579570in}{2.303425in}}%
\pgfpathlineto{\pgfqpoint{4.580560in}{2.303361in}}%
\pgfpathlineto{\pgfqpoint{4.579638in}{2.303451in}}%
\pgfpathlineto{\pgfqpoint{4.580681in}{2.303418in}}%
\pgfpathlineto{\pgfqpoint{4.581757in}{2.303444in}}%
\pgfpathlineto{\pgfqpoint{4.580877in}{2.303362in}}%
\pgfpathlineto{\pgfqpoint{4.581794in}{2.303430in}}%
\pgfpathlineto{\pgfqpoint{4.582198in}{2.303319in}}%
\pgfpathlineto{\pgfqpoint{4.581946in}{2.303446in}}%
\pgfpathlineto{\pgfqpoint{4.582906in}{2.303407in}}%
\pgfpathlineto{\pgfqpoint{4.584031in}{2.303960in}}%
\pgfpathlineto{\pgfqpoint{4.583017in}{2.303362in}}%
\pgfpathlineto{\pgfqpoint{4.584034in}{2.303947in}}%
\pgfpathlineto{\pgfqpoint{4.585183in}{2.304740in}}%
\pgfpathlineto{\pgfqpoint{4.584059in}{2.303901in}}%
\pgfpathlineto{\pgfqpoint{4.585306in}{2.304395in}}%
\pgfpathlineto{\pgfqpoint{4.586136in}{2.304217in}}%
\pgfpathlineto{\pgfqpoint{4.586389in}{2.304529in}}%
\pgfpathlineto{\pgfqpoint{4.586415in}{2.304461in}}%
\pgfpathlineto{\pgfqpoint{4.586890in}{2.304763in}}%
\pgfpathlineto{\pgfqpoint{4.586526in}{2.304413in}}%
\pgfpathlineto{\pgfqpoint{4.587528in}{2.304506in}}%
\pgfpathlineto{\pgfqpoint{4.588672in}{2.305322in}}%
\pgfpathlineto{\pgfqpoint{4.587680in}{2.304465in}}%
\pgfpathlineto{\pgfqpoint{4.588673in}{2.305316in}}%
\pgfpathlineto{\pgfqpoint{4.589773in}{2.304855in}}%
\pgfpathlineto{\pgfqpoint{4.589026in}{2.305541in}}%
\pgfpathlineto{\pgfqpoint{4.589811in}{2.304911in}}%
\pgfpathlineto{\pgfqpoint{4.590980in}{2.305368in}}%
\pgfpathlineto{\pgfqpoint{4.591021in}{2.305179in}}%
\pgfpathlineto{\pgfqpoint{4.591228in}{2.305058in}}%
\pgfpathlineto{\pgfqpoint{4.592100in}{2.305882in}}%
\pgfpathlineto{\pgfqpoint{4.592109in}{2.305862in}}%
\pgfpathlineto{\pgfqpoint{4.593169in}{2.305762in}}%
\pgfpathlineto{\pgfqpoint{4.592535in}{2.305564in}}%
\pgfpathlineto{\pgfqpoint{4.593210in}{2.305654in}}%
\pgfpathlineto{\pgfqpoint{4.594538in}{2.305371in}}%
\pgfpathlineto{\pgfqpoint{4.593448in}{2.306055in}}%
\pgfpathlineto{\pgfqpoint{4.594550in}{2.305418in}}%
\pgfpathlineto{\pgfqpoint{4.595496in}{2.305957in}}%
\pgfpathlineto{\pgfqpoint{4.594710in}{2.305112in}}%
\pgfpathlineto{\pgfqpoint{4.595761in}{2.305769in}}%
\pgfpathlineto{\pgfqpoint{4.595776in}{2.305669in}}%
\pgfpathlineto{\pgfqpoint{4.596649in}{2.306187in}}%
\pgfpathlineto{\pgfqpoint{4.596866in}{2.305930in}}%
\pgfpathlineto{\pgfqpoint{4.597849in}{2.306252in}}%
\pgfpathlineto{\pgfqpoint{4.597109in}{2.305836in}}%
\pgfpathlineto{\pgfqpoint{4.597986in}{2.306008in}}%
\pgfpathlineto{\pgfqpoint{4.598755in}{2.306207in}}%
\pgfpathlineto{\pgfqpoint{4.598546in}{2.305865in}}%
\pgfpathlineto{\pgfqpoint{4.599071in}{2.305776in}}%
\pgfpathlineto{\pgfqpoint{4.599170in}{2.305720in}}%
\pgfpathlineto{\pgfqpoint{4.599831in}{2.306495in}}%
\pgfpathlineto{\pgfqpoint{4.600151in}{2.306370in}}%
\pgfpathlineto{\pgfqpoint{4.601097in}{2.306828in}}%
\pgfpathlineto{\pgfqpoint{4.600400in}{2.306028in}}%
\pgfpathlineto{\pgfqpoint{4.601282in}{2.306812in}}%
\pgfpathlineto{\pgfqpoint{4.601556in}{2.306362in}}%
\pgfpathlineto{\pgfqpoint{4.602240in}{2.307023in}}%
\pgfpathlineto{\pgfqpoint{4.602407in}{2.306750in}}%
\pgfpathlineto{\pgfqpoint{4.603746in}{2.307731in}}%
\pgfpathlineto{\pgfqpoint{4.602711in}{2.306460in}}%
\pgfpathlineto{\pgfqpoint{4.603755in}{2.307693in}}%
\pgfpathlineto{\pgfqpoint{4.604865in}{2.307489in}}%
\pgfpathlineto{\pgfqpoint{4.604155in}{2.307830in}}%
\pgfpathlineto{\pgfqpoint{4.604879in}{2.307554in}}%
\pgfpathlineto{\pgfqpoint{4.605881in}{2.307519in}}%
\pgfpathlineto{\pgfqpoint{4.605023in}{2.307280in}}%
\pgfpathlineto{\pgfqpoint{4.605991in}{2.307483in}}%
\pgfpathlineto{\pgfqpoint{4.606320in}{2.307468in}}%
\pgfpathlineto{\pgfqpoint{4.606768in}{2.306688in}}%
\pgfpathlineto{\pgfqpoint{4.606946in}{2.306911in}}%
\pgfpathlineto{\pgfqpoint{4.606959in}{2.306828in}}%
\pgfpathlineto{\pgfqpoint{4.607660in}{2.307228in}}%
\pgfpathlineto{\pgfqpoint{4.608051in}{2.307036in}}%
\pgfpathlineto{\pgfqpoint{4.608809in}{2.308341in}}%
\pgfpathlineto{\pgfqpoint{4.608066in}{2.306989in}}%
\pgfpathlineto{\pgfqpoint{4.609216in}{2.307634in}}%
\pgfpathlineto{\pgfqpoint{4.610156in}{2.307219in}}%
\pgfpathlineto{\pgfqpoint{4.609518in}{2.307721in}}%
\pgfpathlineto{\pgfqpoint{4.610368in}{2.307458in}}%
\pgfpathlineto{\pgfqpoint{4.610867in}{2.307716in}}%
\pgfpathlineto{\pgfqpoint{4.611339in}{2.307119in}}%
\pgfpathlineto{\pgfqpoint{4.611476in}{2.307387in}}%
\pgfpathlineto{\pgfqpoint{4.612698in}{2.307938in}}%
\pgfpathlineto{\pgfqpoint{4.611910in}{2.307018in}}%
\pgfpathlineto{\pgfqpoint{4.612765in}{2.307843in}}%
\pgfpathlineto{\pgfqpoint{4.613691in}{2.307630in}}%
\pgfpathlineto{\pgfqpoint{4.612967in}{2.308203in}}%
\pgfpathlineto{\pgfqpoint{4.613873in}{2.307923in}}%
\pgfpathlineto{\pgfqpoint{4.614535in}{2.308270in}}%
\pgfpathlineto{\pgfqpoint{4.613875in}{2.307913in}}%
\pgfpathlineto{\pgfqpoint{4.614926in}{2.307577in}}%
\pgfpathlineto{\pgfqpoint{4.615827in}{2.307615in}}%
\pgfpathlineto{\pgfqpoint{4.615713in}{2.308005in}}%
\pgfpathlineto{\pgfqpoint{4.616022in}{2.307808in}}%
\pgfpathlineto{\pgfqpoint{4.616604in}{2.308578in}}%
\pgfpathlineto{\pgfqpoint{4.617144in}{2.307923in}}%
\pgfpathlineto{\pgfqpoint{4.618089in}{2.307666in}}%
\pgfpathlineto{\pgfqpoint{4.617526in}{2.308106in}}%
\pgfpathlineto{\pgfqpoint{4.618262in}{2.307798in}}%
\pgfpathlineto{\pgfqpoint{4.619234in}{2.307918in}}%
\pgfpathlineto{\pgfqpoint{4.618791in}{2.307575in}}%
\pgfpathlineto{\pgfqpoint{4.619370in}{2.307733in}}%
\pgfpathlineto{\pgfqpoint{4.620477in}{2.307615in}}%
\pgfpathlineto{\pgfqpoint{4.620174in}{2.308029in}}%
\pgfpathlineto{\pgfqpoint{4.620492in}{2.307643in}}%
\pgfpathlineto{\pgfqpoint{4.621711in}{2.308239in}}%
\pgfpathlineto{\pgfqpoint{4.620623in}{2.307404in}}%
\pgfpathlineto{\pgfqpoint{4.621745in}{2.308145in}}%
\pgfpathlineto{\pgfqpoint{4.621965in}{2.307867in}}%
\pgfpathlineto{\pgfqpoint{4.622333in}{2.308410in}}%
\pgfpathlineto{\pgfqpoint{4.622867in}{2.308131in}}%
\pgfpathlineto{\pgfqpoint{4.623050in}{2.308069in}}%
\pgfpathlineto{\pgfqpoint{4.624008in}{2.308579in}}%
\pgfpathlineto{\pgfqpoint{4.624931in}{2.308283in}}%
\pgfpathlineto{\pgfqpoint{4.624532in}{2.308694in}}%
\pgfpathlineto{\pgfqpoint{4.625139in}{2.308419in}}%
\pgfpathlineto{\pgfqpoint{4.625191in}{2.308499in}}%
\pgfpathlineto{\pgfqpoint{4.625563in}{2.308053in}}%
\pgfpathlineto{\pgfqpoint{4.626244in}{2.308302in}}%
\pgfpathlineto{\pgfqpoint{4.626814in}{2.308656in}}%
\pgfpathlineto{\pgfqpoint{4.627369in}{2.307955in}}%
\pgfpathlineto{\pgfqpoint{4.631193in}{2.308163in}}%
\pgfpathlineto{\pgfqpoint{4.631197in}{2.308194in}}%
\pgfpathlineto{\pgfqpoint{4.631205in}{2.308249in}}%
\pgfpathlineto{\pgfqpoint{4.631744in}{2.307572in}}%
\pgfpathlineto{\pgfqpoint{4.632212in}{2.307517in}}%
\pgfpathlineto{\pgfqpoint{4.632518in}{2.307261in}}%
\pgfpathlineto{\pgfqpoint{4.633208in}{2.307774in}}%
\pgfpathlineto{\pgfqpoint{4.633318in}{2.307616in}}%
\pgfpathlineto{\pgfqpoint{4.633323in}{2.307639in}}%
\pgfpathlineto{\pgfqpoint{4.634313in}{2.306287in}}%
\pgfpathlineto{\pgfqpoint{4.635152in}{2.306257in}}%
\pgfpathlineto{\pgfqpoint{4.635395in}{2.306658in}}%
\pgfpathlineto{\pgfqpoint{4.635575in}{2.306871in}}%
\pgfpathlineto{\pgfqpoint{4.636391in}{2.306310in}}%
\pgfpathlineto{\pgfqpoint{4.636496in}{2.306482in}}%
\pgfpathlineto{\pgfqpoint{4.636748in}{2.306499in}}%
\pgfpathlineto{\pgfqpoint{4.637529in}{2.305700in}}%
\pgfpathlineto{\pgfqpoint{4.637540in}{2.305739in}}%
\pgfpathlineto{\pgfqpoint{4.638341in}{2.305393in}}%
\pgfpathlineto{\pgfqpoint{4.637565in}{2.305785in}}%
\pgfpathlineto{\pgfqpoint{4.638642in}{2.305879in}}%
\pgfpathlineto{\pgfqpoint{4.639646in}{2.306520in}}%
\pgfpathlineto{\pgfqpoint{4.639699in}{2.305257in}}%
\pgfpathlineto{\pgfqpoint{4.640338in}{2.303341in}}%
\pgfpathlineto{\pgfqpoint{4.640900in}{2.303370in}}%
\pgfpathlineto{\pgfqpoint{4.642000in}{2.303353in}}%
\pgfpathlineto{\pgfqpoint{4.640923in}{2.303443in}}%
\pgfpathlineto{\pgfqpoint{4.642025in}{2.303396in}}%
\pgfpathlineto{\pgfqpoint{4.643083in}{2.303448in}}%
\pgfpathlineto{\pgfqpoint{4.642396in}{2.303347in}}%
\pgfpathlineto{\pgfqpoint{4.643140in}{2.303437in}}%
\pgfpathlineto{\pgfqpoint{4.644193in}{2.303356in}}%
\pgfpathlineto{\pgfqpoint{4.643189in}{2.303443in}}%
\pgfpathlineto{\pgfqpoint{4.644253in}{2.303401in}}%
\pgfpathlineto{\pgfqpoint{4.645381in}{2.303443in}}%
\pgfpathlineto{\pgfqpoint{4.644264in}{2.303374in}}%
\pgfpathlineto{\pgfqpoint{4.645384in}{2.303426in}}%
\pgfpathlineto{\pgfqpoint{4.645759in}{2.303351in}}%
\pgfpathlineto{\pgfqpoint{4.645387in}{2.303443in}}%
\pgfpathlineto{\pgfqpoint{4.646497in}{2.303415in}}%
\pgfpathlineto{\pgfqpoint{4.647487in}{2.303458in}}%
\pgfpathlineto{\pgfqpoint{4.646609in}{2.303306in}}%
\pgfpathlineto{\pgfqpoint{4.647608in}{2.303402in}}%
\pgfpathlineto{\pgfqpoint{4.648366in}{2.303356in}}%
\pgfpathlineto{\pgfqpoint{4.647951in}{2.303450in}}%
\pgfpathlineto{\pgfqpoint{4.648719in}{2.303420in}}%
\pgfpathlineto{\pgfqpoint{4.649751in}{2.303449in}}%
\pgfpathlineto{\pgfqpoint{4.648804in}{2.303366in}}%
\pgfpathlineto{\pgfqpoint{4.649833in}{2.303416in}}%
\pgfpathlineto{\pgfqpoint{4.650978in}{2.303453in}}%
\pgfpathlineto{\pgfqpoint{4.649852in}{2.303366in}}%
\pgfpathlineto{\pgfqpoint{4.650979in}{2.303425in}}%
\pgfpathlineto{\pgfqpoint{4.651742in}{2.303357in}}%
\pgfpathlineto{\pgfqpoint{4.651242in}{2.303451in}}%
\pgfpathlineto{\pgfqpoint{4.652090in}{2.303419in}}%
\pgfpathlineto{\pgfqpoint{4.652909in}{2.303454in}}%
\pgfpathlineto{\pgfqpoint{4.652546in}{2.303358in}}%
\pgfpathlineto{\pgfqpoint{4.653202in}{2.303419in}}%
\pgfpathlineto{\pgfqpoint{4.654291in}{2.303443in}}%
\pgfpathlineto{\pgfqpoint{4.653307in}{2.303356in}}%
\pgfpathlineto{\pgfqpoint{4.654318in}{2.303435in}}%
\pgfpathlineto{\pgfqpoint{4.655174in}{2.303349in}}%
\pgfpathlineto{\pgfqpoint{4.654353in}{2.303444in}}%
\pgfpathlineto{\pgfqpoint{4.655430in}{2.303410in}}%
\pgfpathlineto{\pgfqpoint{4.656497in}{2.303451in}}%
\pgfpathlineto{\pgfqpoint{4.655524in}{2.303365in}}%
\pgfpathlineto{\pgfqpoint{4.656541in}{2.303409in}}%
\pgfpathlineto{\pgfqpoint{4.657719in}{2.303434in}}%
\pgfpathlineto{\pgfqpoint{4.656549in}{2.303391in}}%
\pgfpathlineto{\pgfqpoint{4.657720in}{2.303402in}}%
\pgfpathlineto{\pgfqpoint{4.658339in}{2.303374in}}%
\pgfpathlineto{\pgfqpoint{4.658566in}{2.303789in}}%
\pgfpathlineto{\pgfqpoint{4.658809in}{2.303620in}}%
\pgfpathlineto{\pgfqpoint{4.658981in}{2.303797in}}%
\pgfpathlineto{\pgfqpoint{4.659626in}{2.303181in}}%
\pgfpathlineto{\pgfqpoint{4.659898in}{2.303408in}}%
\pgfpathlineto{\pgfqpoint{4.660973in}{2.303263in}}%
\pgfpathlineto{\pgfqpoint{4.660460in}{2.303906in}}%
\pgfpathlineto{\pgfqpoint{4.661011in}{2.303348in}}%
\pgfpathlineto{\pgfqpoint{4.662017in}{2.304009in}}%
\pgfpathlineto{\pgfqpoint{4.661076in}{2.303250in}}%
\pgfpathlineto{\pgfqpoint{4.662140in}{2.303845in}}%
\pgfpathlineto{\pgfqpoint{4.663278in}{2.303108in}}%
\pgfpathlineto{\pgfqpoint{4.662223in}{2.303915in}}%
\pgfpathlineto{\pgfqpoint{4.663373in}{2.303237in}}%
\pgfpathlineto{\pgfqpoint{4.664547in}{2.304105in}}%
\pgfpathlineto{\pgfqpoint{4.663384in}{2.303190in}}%
\pgfpathlineto{\pgfqpoint{4.664550in}{2.304088in}}%
\pgfpathlineto{\pgfqpoint{4.664560in}{2.303975in}}%
\pgfpathlineto{\pgfqpoint{4.665603in}{2.304574in}}%
\pgfpathlineto{\pgfqpoint{4.665627in}{2.304481in}}%
\pgfpathlineto{\pgfqpoint{4.666291in}{2.304861in}}%
\pgfpathlineto{\pgfqpoint{4.665871in}{2.304438in}}%
\pgfpathlineto{\pgfqpoint{4.666795in}{2.304747in}}%
\pgfpathlineto{\pgfqpoint{4.667197in}{2.304720in}}%
\pgfpathlineto{\pgfqpoint{4.667783in}{2.305408in}}%
\pgfpathlineto{\pgfqpoint{4.667892in}{2.305071in}}%
\pgfpathlineto{\pgfqpoint{4.668794in}{2.305413in}}%
\pgfpathlineto{\pgfqpoint{4.667997in}{2.304942in}}%
\pgfpathlineto{\pgfqpoint{4.669076in}{2.305011in}}%
\pgfpathlineto{\pgfqpoint{4.670113in}{2.304369in}}%
\pgfpathlineto{\pgfqpoint{4.669085in}{2.305029in}}%
\pgfpathlineto{\pgfqpoint{4.670211in}{2.304595in}}%
\pgfpathlineto{\pgfqpoint{4.671290in}{2.304894in}}%
\pgfpathlineto{\pgfqpoint{4.670295in}{2.304537in}}%
\pgfpathlineto{\pgfqpoint{4.671345in}{2.304728in}}%
\pgfpathlineto{\pgfqpoint{4.672191in}{2.304478in}}%
\pgfpathlineto{\pgfqpoint{4.671417in}{2.304789in}}%
\pgfpathlineto{\pgfqpoint{4.672443in}{2.304900in}}%
\pgfpathlineto{\pgfqpoint{4.672658in}{2.305102in}}%
\pgfpathlineto{\pgfqpoint{4.673123in}{2.304524in}}%
\pgfpathlineto{\pgfqpoint{4.673548in}{2.304778in}}%
\pgfpathlineto{\pgfqpoint{4.673555in}{2.304748in}}%
\pgfpathlineto{\pgfqpoint{4.674598in}{2.305456in}}%
\pgfpathlineto{\pgfqpoint{4.674616in}{2.305385in}}%
\pgfpathlineto{\pgfqpoint{4.675726in}{2.305906in}}%
\pgfpathlineto{\pgfqpoint{4.674992in}{2.305239in}}%
\pgfpathlineto{\pgfqpoint{4.675748in}{2.305869in}}%
\pgfpathlineto{\pgfqpoint{4.676788in}{2.305873in}}%
\pgfpathlineto{\pgfqpoint{4.676179in}{2.305487in}}%
\pgfpathlineto{\pgfqpoint{4.676870in}{2.305819in}}%
\pgfpathlineto{\pgfqpoint{4.678262in}{2.304972in}}%
\pgfpathlineto{\pgfqpoint{4.677337in}{2.306102in}}%
\pgfpathlineto{\pgfqpoint{4.678286in}{2.305061in}}%
\pgfpathlineto{\pgfqpoint{4.679091in}{2.305215in}}%
\pgfpathlineto{\pgfqpoint{4.678587in}{2.304885in}}%
\pgfpathlineto{\pgfqpoint{4.679393in}{2.304970in}}%
\pgfpathlineto{\pgfqpoint{4.680112in}{2.304867in}}%
\pgfpathlineto{\pgfqpoint{4.680483in}{2.305226in}}%
\pgfpathlineto{\pgfqpoint{4.680487in}{2.305212in}}%
\pgfpathlineto{\pgfqpoint{4.681741in}{2.306005in}}%
\pgfpathlineto{\pgfqpoint{4.680537in}{2.305152in}}%
\pgfpathlineto{\pgfqpoint{4.681743in}{2.305987in}}%
\pgfpathlineto{\pgfqpoint{4.682838in}{2.305750in}}%
\pgfpathlineto{\pgfqpoint{4.682398in}{2.306121in}}%
\pgfpathlineto{\pgfqpoint{4.682864in}{2.305811in}}%
\pgfpathlineto{\pgfqpoint{4.685555in}{2.306404in}}%
\pgfpathlineto{\pgfqpoint{4.682866in}{2.305828in}}%
\pgfpathlineto{\pgfqpoint{4.685557in}{2.306429in}}%
\pgfpathlineto{\pgfqpoint{4.686442in}{2.306559in}}%
\pgfpathlineto{\pgfqpoint{4.685669in}{2.306176in}}%
\pgfpathlineto{\pgfqpoint{4.686654in}{2.306266in}}%
\pgfpathlineto{\pgfqpoint{4.687723in}{2.305201in}}%
\pgfpathlineto{\pgfqpoint{4.686657in}{2.306285in}}%
\pgfpathlineto{\pgfqpoint{4.687868in}{2.305390in}}%
\pgfpathlineto{\pgfqpoint{4.688358in}{2.305526in}}%
\pgfpathlineto{\pgfqpoint{4.687961in}{2.305113in}}%
\pgfpathlineto{\pgfqpoint{4.688975in}{2.305313in}}%
\pgfpathlineto{\pgfqpoint{4.692391in}{2.305488in}}%
\pgfpathlineto{\pgfqpoint{4.688977in}{2.305295in}}%
\pgfpathlineto{\pgfqpoint{4.692423in}{2.305325in}}%
\pgfpathlineto{\pgfqpoint{4.693546in}{2.305160in}}%
\pgfpathlineto{\pgfqpoint{4.693282in}{2.305732in}}%
\pgfpathlineto{\pgfqpoint{4.693547in}{2.305160in}}%
\pgfpathlineto{\pgfqpoint{4.694649in}{2.305534in}}%
\pgfpathlineto{\pgfqpoint{4.693849in}{2.305024in}}%
\pgfpathlineto{\pgfqpoint{4.694687in}{2.305467in}}%
\pgfpathlineto{\pgfqpoint{4.694720in}{2.305411in}}%
\pgfpathlineto{\pgfqpoint{4.695787in}{2.305828in}}%
\pgfpathlineto{\pgfqpoint{4.695788in}{2.305818in}}%
\pgfpathlineto{\pgfqpoint{4.696399in}{2.306078in}}%
\pgfpathlineto{\pgfqpoint{4.695959in}{2.305768in}}%
\pgfpathlineto{\pgfqpoint{4.696904in}{2.305811in}}%
\pgfpathlineto{\pgfqpoint{4.698365in}{2.306456in}}%
\pgfpathlineto{\pgfqpoint{4.697152in}{2.305634in}}%
\pgfpathlineto{\pgfqpoint{4.698380in}{2.306393in}}%
\pgfpathlineto{\pgfqpoint{4.699423in}{2.306171in}}%
\pgfpathlineto{\pgfqpoint{4.698862in}{2.306709in}}%
\pgfpathlineto{\pgfqpoint{4.699529in}{2.306265in}}%
\pgfpathlineto{\pgfqpoint{4.701040in}{2.306312in}}%
\pgfpathlineto{\pgfqpoint{4.699531in}{2.306286in}}%
\pgfpathlineto{\pgfqpoint{4.701085in}{2.306487in}}%
\pgfpathlineto{\pgfqpoint{4.701085in}{2.306519in}}%
\pgfpathlineto{\pgfqpoint{4.701769in}{2.306115in}}%
\pgfpathlineto{\pgfqpoint{4.702194in}{2.306428in}}%
\pgfpathlineto{\pgfqpoint{4.702902in}{2.306344in}}%
\pgfpathlineto{\pgfqpoint{4.703216in}{2.306853in}}%
\pgfpathlineto{\pgfqpoint{4.703268in}{2.306780in}}%
\pgfpathlineto{\pgfqpoint{4.704361in}{2.307222in}}%
\pgfpathlineto{\pgfqpoint{4.703651in}{2.306488in}}%
\pgfpathlineto{\pgfqpoint{4.704391in}{2.307195in}}%
\pgfpathlineto{\pgfqpoint{4.704993in}{2.306737in}}%
\pgfpathlineto{\pgfqpoint{4.704456in}{2.307268in}}%
\pgfpathlineto{\pgfqpoint{4.705525in}{2.307014in}}%
\pgfpathlineto{\pgfqpoint{4.706151in}{2.307401in}}%
\pgfpathlineto{\pgfqpoint{4.705597in}{2.306793in}}%
\pgfpathlineto{\pgfqpoint{4.706652in}{2.307072in}}%
\pgfpathlineto{\pgfqpoint{4.707731in}{2.307264in}}%
\pgfpathlineto{\pgfqpoint{4.706695in}{2.307121in}}%
\pgfpathlineto{\pgfqpoint{4.707750in}{2.307409in}}%
\pgfpathlineto{\pgfqpoint{4.708384in}{2.307611in}}%
\pgfpathlineto{\pgfqpoint{4.707858in}{2.307229in}}%
\pgfpathlineto{\pgfqpoint{4.708859in}{2.307371in}}%
\pgfpathlineto{\pgfqpoint{4.709544in}{2.306958in}}%
\pgfpathlineto{\pgfqpoint{4.709016in}{2.307484in}}%
\pgfpathlineto{\pgfqpoint{4.709980in}{2.307236in}}%
\pgfpathlineto{\pgfqpoint{4.710904in}{2.308012in}}%
\pgfpathlineto{\pgfqpoint{4.710016in}{2.307108in}}%
\pgfpathlineto{\pgfqpoint{4.711214in}{2.307934in}}%
\pgfpathlineto{\pgfqpoint{4.712240in}{2.307704in}}%
\pgfpathlineto{\pgfqpoint{4.711391in}{2.308193in}}%
\pgfpathlineto{\pgfqpoint{4.712341in}{2.307867in}}%
\pgfpathlineto{\pgfqpoint{4.712829in}{2.308092in}}%
\pgfpathlineto{\pgfqpoint{4.713053in}{2.307664in}}%
\pgfpathlineto{\pgfqpoint{4.713450in}{2.307788in}}%
\pgfpathlineto{\pgfqpoint{4.714205in}{2.307799in}}%
\pgfpathlineto{\pgfqpoint{4.713636in}{2.308198in}}%
\pgfpathlineto{\pgfqpoint{4.714537in}{2.308047in}}%
\pgfpathlineto{\pgfqpoint{4.714549in}{2.308142in}}%
\pgfpathlineto{\pgfqpoint{4.714923in}{2.303366in}}%
\pgfpathlineto{\pgfqpoint{4.715381in}{2.303436in}}%
\pgfpathlineto{\pgfqpoint{4.716368in}{2.303366in}}%
\pgfpathlineto{\pgfqpoint{4.715426in}{2.303448in}}%
\pgfpathlineto{\pgfqpoint{4.716494in}{2.303418in}}%
\pgfpathlineto{\pgfqpoint{4.717521in}{2.303450in}}%
\pgfpathlineto{\pgfqpoint{4.716545in}{2.303377in}}%
\pgfpathlineto{\pgfqpoint{4.717610in}{2.303376in}}%
\pgfpathlineto{\pgfqpoint{4.719108in}{2.303399in}}%
\pgfpathlineto{\pgfqpoint{4.717618in}{2.303424in}}%
\pgfpathlineto{\pgfqpoint{4.719111in}{2.303408in}}%
\pgfpathlineto{\pgfqpoint{4.719963in}{2.303459in}}%
\pgfpathlineto{\pgfqpoint{4.719195in}{2.303345in}}%
\pgfpathlineto{\pgfqpoint{4.720223in}{2.303414in}}%
\pgfpathlineto{\pgfqpoint{4.721233in}{2.303447in}}%
\pgfpathlineto{\pgfqpoint{4.720239in}{2.303363in}}%
\pgfpathlineto{\pgfqpoint{4.721334in}{2.303431in}}%
\pgfpathlineto{\pgfqpoint{4.722353in}{2.303373in}}%
\pgfpathlineto{\pgfqpoint{4.721344in}{2.303444in}}%
\pgfpathlineto{\pgfqpoint{4.722448in}{2.303388in}}%
\pgfpathlineto{\pgfqpoint{4.723538in}{2.303442in}}%
\pgfpathlineto{\pgfqpoint{4.722452in}{2.303378in}}%
\pgfpathlineto{\pgfqpoint{4.723562in}{2.303428in}}%
\pgfpathlineto{\pgfqpoint{4.724658in}{2.303358in}}%
\pgfpathlineto{\pgfqpoint{4.723713in}{2.303448in}}%
\pgfpathlineto{\pgfqpoint{4.724674in}{2.303424in}}%
\pgfpathlineto{\pgfqpoint{4.725359in}{2.303350in}}%
\pgfpathlineto{\pgfqpoint{4.724721in}{2.303443in}}%
\pgfpathlineto{\pgfqpoint{4.725785in}{2.303410in}}%
\pgfpathlineto{\pgfqpoint{4.726906in}{2.303432in}}%
\pgfpathlineto{\pgfqpoint{4.725827in}{2.303387in}}%
\pgfpathlineto{\pgfqpoint{4.726909in}{2.303426in}}%
\pgfpathlineto{\pgfqpoint{4.727837in}{2.303367in}}%
\pgfpathlineto{\pgfqpoint{4.727048in}{2.303444in}}%
\pgfpathlineto{\pgfqpoint{4.728020in}{2.303422in}}%
\pgfpathlineto{\pgfqpoint{4.729103in}{2.303445in}}%
\pgfpathlineto{\pgfqpoint{4.728098in}{2.303371in}}%
\pgfpathlineto{\pgfqpoint{4.729136in}{2.303397in}}%
\pgfpathlineto{\pgfqpoint{4.730195in}{2.303368in}}%
\pgfpathlineto{\pgfqpoint{4.729278in}{2.303451in}}%
\pgfpathlineto{\pgfqpoint{4.730248in}{2.303392in}}%
\pgfpathlineto{\pgfqpoint{4.731318in}{2.303447in}}%
\pgfpathlineto{\pgfqpoint{4.731150in}{2.303345in}}%
\pgfpathlineto{\pgfqpoint{4.731360in}{2.303410in}}%
\pgfpathlineto{\pgfqpoint{4.732341in}{2.303461in}}%
\pgfpathlineto{\pgfqpoint{4.731450in}{2.303351in}}%
\pgfpathlineto{\pgfqpoint{4.732474in}{2.303426in}}%
\pgfpathlineto{\pgfqpoint{4.733354in}{2.303267in}}%
\pgfpathlineto{\pgfqpoint{4.732684in}{2.303452in}}%
\pgfpathlineto{\pgfqpoint{4.733586in}{2.303395in}}%
\pgfpathlineto{\pgfqpoint{4.734544in}{2.304001in}}%
\pgfpathlineto{\pgfqpoint{4.733821in}{2.303377in}}%
\pgfpathlineto{\pgfqpoint{4.734716in}{2.303818in}}%
\pgfpathlineto{\pgfqpoint{4.735641in}{2.303772in}}%
\pgfpathlineto{\pgfqpoint{4.734907in}{2.303571in}}%
\pgfpathlineto{\pgfqpoint{4.735784in}{2.303436in}}%
\pgfpathlineto{\pgfqpoint{4.736915in}{2.302987in}}%
\pgfpathlineto{\pgfqpoint{4.735864in}{2.303481in}}%
\pgfpathlineto{\pgfqpoint{4.736920in}{2.303019in}}%
\pgfpathlineto{\pgfqpoint{4.737173in}{2.302815in}}%
\pgfpathlineto{\pgfqpoint{4.737578in}{2.303263in}}%
\pgfpathlineto{\pgfqpoint{4.738031in}{2.303103in}}%
\pgfpathlineto{\pgfqpoint{4.739112in}{2.303509in}}%
\pgfpathlineto{\pgfqpoint{4.738463in}{2.302908in}}%
\pgfpathlineto{\pgfqpoint{4.739160in}{2.303484in}}%
\pgfpathlineto{\pgfqpoint{4.740554in}{2.302573in}}%
\pgfpathlineto{\pgfqpoint{4.739214in}{2.303610in}}%
\pgfpathlineto{\pgfqpoint{4.740555in}{2.302587in}}%
\pgfpathlineto{\pgfqpoint{4.741563in}{2.303239in}}%
\pgfpathlineto{\pgfqpoint{4.740631in}{2.302360in}}%
\pgfpathlineto{\pgfqpoint{4.741698in}{2.303054in}}%
\pgfpathlineto{\pgfqpoint{4.742643in}{2.302374in}}%
\pgfpathlineto{\pgfqpoint{4.741992in}{2.303244in}}%
\pgfpathlineto{\pgfqpoint{4.742823in}{2.302961in}}%
\pgfpathlineto{\pgfqpoint{4.743270in}{2.303235in}}%
\pgfpathlineto{\pgfqpoint{4.743794in}{2.302745in}}%
\pgfpathlineto{\pgfqpoint{4.743929in}{2.302868in}}%
\pgfpathlineto{\pgfqpoint{4.743985in}{2.302755in}}%
\pgfpathlineto{\pgfqpoint{4.744834in}{2.303268in}}%
\pgfpathlineto{\pgfqpoint{4.745025in}{2.303143in}}%
\pgfpathlineto{\pgfqpoint{4.746075in}{2.303336in}}%
\pgfpathlineto{\pgfqpoint{4.745917in}{2.302917in}}%
\pgfpathlineto{\pgfqpoint{4.746157in}{2.303297in}}%
\pgfpathlineto{\pgfqpoint{4.747402in}{2.302686in}}%
\pgfpathlineto{\pgfqpoint{4.746159in}{2.303304in}}%
\pgfpathlineto{\pgfqpoint{4.747449in}{2.302927in}}%
\pgfpathlineto{\pgfqpoint{4.747609in}{2.303094in}}%
\pgfpathlineto{\pgfqpoint{4.748494in}{2.301894in}}%
\pgfpathlineto{\pgfqpoint{4.749647in}{2.301659in}}%
\pgfpathlineto{\pgfqpoint{4.748688in}{2.302129in}}%
\pgfpathlineto{\pgfqpoint{4.749647in}{2.301664in}}%
\pgfpathlineto{\pgfqpoint{4.750320in}{2.301943in}}%
\pgfpathlineto{\pgfqpoint{4.749764in}{2.301558in}}%
\pgfpathlineto{\pgfqpoint{4.750757in}{2.301614in}}%
\pgfpathlineto{\pgfqpoint{4.751275in}{2.301340in}}%
\pgfpathlineto{\pgfqpoint{4.751770in}{2.301832in}}%
\pgfpathlineto{\pgfqpoint{4.751862in}{2.301724in}}%
\pgfpathlineto{\pgfqpoint{4.752903in}{2.302118in}}%
\pgfpathlineto{\pgfqpoint{4.751974in}{2.301592in}}%
\pgfpathlineto{\pgfqpoint{4.752989in}{2.301945in}}%
\pgfpathlineto{\pgfqpoint{4.754097in}{2.301751in}}%
\pgfpathlineto{\pgfqpoint{4.753102in}{2.302244in}}%
\pgfpathlineto{\pgfqpoint{4.754124in}{2.301791in}}%
\pgfpathlineto{\pgfqpoint{4.754901in}{2.302508in}}%
\pgfpathlineto{\pgfqpoint{4.754171in}{2.301675in}}%
\pgfpathlineto{\pgfqpoint{4.755260in}{2.302161in}}%
\pgfpathlineto{\pgfqpoint{4.756151in}{2.301584in}}%
\pgfpathlineto{\pgfqpoint{4.755327in}{2.302238in}}%
\pgfpathlineto{\pgfqpoint{4.756391in}{2.301927in}}%
\pgfpathlineto{\pgfqpoint{4.756660in}{2.302076in}}%
\pgfpathlineto{\pgfqpoint{4.757369in}{2.301600in}}%
\pgfpathlineto{\pgfqpoint{4.757495in}{2.301797in}}%
\pgfpathlineto{\pgfqpoint{4.757513in}{2.301894in}}%
\pgfpathlineto{\pgfqpoint{4.757907in}{2.301317in}}%
\pgfpathlineto{\pgfqpoint{4.758582in}{2.301571in}}%
\pgfpathlineto{\pgfqpoint{4.758653in}{2.301388in}}%
\pgfpathlineto{\pgfqpoint{4.758946in}{2.302226in}}%
\pgfpathlineto{\pgfqpoint{4.759670in}{2.301865in}}%
\pgfpathlineto{\pgfqpoint{4.761288in}{2.302447in}}%
\pgfpathlineto{\pgfqpoint{4.759714in}{2.301725in}}%
\pgfpathlineto{\pgfqpoint{4.761312in}{2.302369in}}%
\pgfpathlineto{\pgfqpoint{4.761627in}{2.302129in}}%
\pgfpathlineto{\pgfqpoint{4.761454in}{2.302518in}}%
\pgfpathlineto{\pgfqpoint{4.762431in}{2.302291in}}%
\pgfpathlineto{\pgfqpoint{4.762685in}{2.302462in}}%
\pgfpathlineto{\pgfqpoint{4.763176in}{2.301971in}}%
\pgfpathlineto{\pgfqpoint{4.763538in}{2.302209in}}%
\pgfpathlineto{\pgfqpoint{4.764622in}{2.302363in}}%
\pgfpathlineto{\pgfqpoint{4.763586in}{2.302059in}}%
\pgfpathlineto{\pgfqpoint{4.764673in}{2.302259in}}%
\pgfpathlineto{\pgfqpoint{4.765551in}{2.301591in}}%
\pgfpathlineto{\pgfqpoint{4.764720in}{2.302398in}}%
\pgfpathlineto{\pgfqpoint{4.765805in}{2.301712in}}%
\pgfpathlineto{\pgfqpoint{4.766957in}{2.301597in}}%
\pgfpathlineto{\pgfqpoint{4.766426in}{2.302018in}}%
\pgfpathlineto{\pgfqpoint{4.766963in}{2.301647in}}%
\pgfpathlineto{\pgfqpoint{4.767793in}{2.301991in}}%
\pgfpathlineto{\pgfqpoint{4.768189in}{2.301650in}}%
\pgfpathlineto{\pgfqpoint{4.769202in}{2.301579in}}%
\pgfpathlineto{\pgfqpoint{4.768620in}{2.301932in}}%
\pgfpathlineto{\pgfqpoint{4.769297in}{2.301721in}}%
\pgfpathlineto{\pgfqpoint{4.773086in}{2.301058in}}%
\pgfpathlineto{\pgfqpoint{4.773090in}{2.301000in}}%
\pgfpathlineto{\pgfqpoint{4.773280in}{2.300906in}}%
\pgfpathlineto{\pgfqpoint{4.774025in}{2.301458in}}%
\pgfpathlineto{\pgfqpoint{4.774187in}{2.301420in}}%
\pgfpathlineto{\pgfqpoint{4.775019in}{2.301293in}}%
\pgfpathlineto{\pgfqpoint{4.774418in}{2.301817in}}%
\pgfpathlineto{\pgfqpoint{4.775283in}{2.301623in}}%
\pgfpathlineto{\pgfqpoint{4.775337in}{2.301832in}}%
\pgfpathlineto{\pgfqpoint{4.776150in}{2.301261in}}%
\pgfpathlineto{\pgfqpoint{4.776353in}{2.301328in}}%
\pgfpathlineto{\pgfqpoint{4.776466in}{2.301261in}}%
\pgfpathlineto{\pgfqpoint{4.777304in}{2.301780in}}%
\pgfpathlineto{\pgfqpoint{4.777452in}{2.301541in}}%
\pgfpathlineto{\pgfqpoint{4.778529in}{2.301652in}}%
\pgfpathlineto{\pgfqpoint{4.778011in}{2.301277in}}%
\pgfpathlineto{\pgfqpoint{4.778586in}{2.301525in}}%
\pgfpathlineto{\pgfqpoint{4.779210in}{2.301345in}}%
\pgfpathlineto{\pgfqpoint{4.779613in}{2.301771in}}%
\pgfpathlineto{\pgfqpoint{4.779694in}{2.301611in}}%
\pgfpathlineto{\pgfqpoint{4.780923in}{2.301006in}}%
\pgfpathlineto{\pgfqpoint{4.779742in}{2.301690in}}%
\pgfpathlineto{\pgfqpoint{4.781052in}{2.301130in}}%
\pgfpathlineto{\pgfqpoint{4.781147in}{2.301350in}}%
\pgfpathlineto{\pgfqpoint{4.782132in}{2.300806in}}%
\pgfpathlineto{\pgfqpoint{4.782145in}{2.300855in}}%
\pgfpathlineto{\pgfqpoint{4.783062in}{2.301334in}}%
\pgfpathlineto{\pgfqpoint{4.782913in}{2.301539in}}%
\pgfpathlineto{\pgfqpoint{4.783102in}{2.301432in}}%
\pgfpathlineto{\pgfqpoint{4.783559in}{2.301793in}}%
\pgfpathlineto{\pgfqpoint{4.784129in}{2.300785in}}%
\pgfpathlineto{\pgfqpoint{4.784187in}{2.300813in}}%
\pgfpathlineto{\pgfqpoint{4.784262in}{2.300726in}}%
\pgfpathlineto{\pgfqpoint{4.785152in}{2.301225in}}%
\pgfpathlineto{\pgfqpoint{4.785280in}{2.301202in}}%
\pgfpathlineto{\pgfqpoint{4.785367in}{2.301113in}}%
\pgfpathlineto{\pgfqpoint{4.786214in}{2.301969in}}%
\pgfpathlineto{\pgfqpoint{4.786235in}{2.301922in}}%
\pgfpathlineto{\pgfqpoint{4.786834in}{2.302092in}}%
\pgfpathlineto{\pgfqpoint{4.787117in}{2.301608in}}%
\pgfpathlineto{\pgfqpoint{4.787335in}{2.301722in}}%
\pgfpathlineto{\pgfqpoint{4.787839in}{2.301078in}}%
\pgfpathlineto{\pgfqpoint{4.787336in}{2.301725in}}%
\pgfpathlineto{\pgfqpoint{4.788491in}{2.301282in}}%
\pgfpathlineto{\pgfqpoint{4.790161in}{2.303451in}}%
\pgfpathlineto{\pgfqpoint{4.788557in}{2.301192in}}%
\pgfpathlineto{\pgfqpoint{4.790214in}{2.303368in}}%
\pgfpathlineto{\pgfqpoint{4.793346in}{2.303424in}}%
\pgfpathlineto{\pgfqpoint{4.794226in}{2.303356in}}%
\pgfpathlineto{\pgfqpoint{4.793438in}{2.303449in}}%
\pgfpathlineto{\pgfqpoint{4.794457in}{2.303445in}}%
\pgfpathlineto{\pgfqpoint{4.795564in}{2.303338in}}%
\pgfpathlineto{\pgfqpoint{4.794565in}{2.303451in}}%
\pgfpathlineto{\pgfqpoint{4.795570in}{2.303413in}}%
\pgfpathlineto{\pgfqpoint{4.796611in}{2.303448in}}%
\pgfpathlineto{\pgfqpoint{4.795697in}{2.303364in}}%
\pgfpathlineto{\pgfqpoint{4.796686in}{2.303420in}}%
\pgfpathlineto{\pgfqpoint{4.797815in}{2.303380in}}%
\pgfpathlineto{\pgfqpoint{4.796697in}{2.303434in}}%
\pgfpathlineto{\pgfqpoint{4.797867in}{2.303414in}}%
\pgfpathlineto{\pgfqpoint{4.798949in}{2.303451in}}%
\pgfpathlineto{\pgfqpoint{4.797935in}{2.303375in}}%
\pgfpathlineto{\pgfqpoint{4.798983in}{2.303434in}}%
\pgfpathlineto{\pgfqpoint{4.800054in}{2.303378in}}%
\pgfpathlineto{\pgfqpoint{4.799057in}{2.303462in}}%
\pgfpathlineto{\pgfqpoint{4.800100in}{2.303416in}}%
\pgfpathlineto{\pgfqpoint{4.801116in}{2.303450in}}%
\pgfpathlineto{\pgfqpoint{4.800182in}{2.303366in}}%
\pgfpathlineto{\pgfqpoint{4.801213in}{2.303404in}}%
\pgfpathlineto{\pgfqpoint{4.802340in}{2.303436in}}%
\pgfpathlineto{\pgfqpoint{4.801229in}{2.303385in}}%
\pgfpathlineto{\pgfqpoint{4.802358in}{2.303428in}}%
\pgfpathlineto{\pgfqpoint{4.802624in}{2.303337in}}%
\pgfpathlineto{\pgfqpoint{4.802603in}{2.303453in}}%
\pgfpathlineto{\pgfqpoint{4.803469in}{2.303401in}}%
\pgfpathlineto{\pgfqpoint{4.804579in}{2.303381in}}%
\pgfpathlineto{\pgfqpoint{4.803523in}{2.303444in}}%
\pgfpathlineto{\pgfqpoint{4.804587in}{2.303413in}}%
\pgfpathlineto{\pgfqpoint{4.808774in}{2.303464in}}%
\pgfpathlineto{\pgfqpoint{4.808783in}{2.303413in}}%
\pgfpathlineto{\pgfqpoint{4.808800in}{2.303376in}}%
\pgfpathlineto{\pgfqpoint{4.809799in}{2.304060in}}%
\pgfpathlineto{\pgfqpoint{4.809866in}{2.303941in}}%
\pgfpathlineto{\pgfqpoint{4.810686in}{2.304569in}}%
\pgfpathlineto{\pgfqpoint{4.811015in}{2.304099in}}%
\pgfpathlineto{\pgfqpoint{4.811775in}{2.303842in}}%
\pgfpathlineto{\pgfqpoint{4.811179in}{2.304222in}}%
\pgfpathlineto{\pgfqpoint{4.812113in}{2.304293in}}%
\pgfpathlineto{\pgfqpoint{4.820054in}{2.304160in}}%
\pgfpathlineto{\pgfqpoint{4.820056in}{2.304186in}}%
\pgfpathlineto{\pgfqpoint{4.821160in}{2.304445in}}%
\pgfpathlineto{\pgfqpoint{4.820234in}{2.304045in}}%
\pgfpathlineto{\pgfqpoint{4.821181in}{2.304396in}}%
\pgfpathlineto{\pgfqpoint{4.821512in}{2.304070in}}%
\pgfpathlineto{\pgfqpoint{4.822232in}{2.305010in}}%
\pgfpathlineto{\pgfqpoint{4.822256in}{2.304979in}}%
\pgfpathlineto{\pgfqpoint{4.823226in}{2.305751in}}%
\pgfpathlineto{\pgfqpoint{4.823412in}{2.305583in}}%
\pgfpathlineto{\pgfqpoint{4.824416in}{2.305388in}}%
\pgfpathlineto{\pgfqpoint{4.823728in}{2.305698in}}%
\pgfpathlineto{\pgfqpoint{4.824529in}{2.305479in}}%
\pgfpathlineto{\pgfqpoint{4.825640in}{2.305784in}}%
\pgfpathlineto{\pgfqpoint{4.824747in}{2.305350in}}%
\pgfpathlineto{\pgfqpoint{4.825656in}{2.305758in}}%
\pgfpathlineto{\pgfqpoint{4.826798in}{2.306188in}}%
\pgfpathlineto{\pgfqpoint{4.826803in}{2.306169in}}%
\pgfpathlineto{\pgfqpoint{4.827513in}{2.305500in}}%
\pgfpathlineto{\pgfqpoint{4.826815in}{2.306203in}}%
\pgfpathlineto{\pgfqpoint{4.827963in}{2.305738in}}%
\pgfpathlineto{\pgfqpoint{4.828793in}{2.306471in}}%
\pgfpathlineto{\pgfqpoint{4.828167in}{2.305667in}}%
\pgfpathlineto{\pgfqpoint{4.829101in}{2.306225in}}%
\pgfpathlineto{\pgfqpoint{4.829339in}{2.306026in}}%
\pgfpathlineto{\pgfqpoint{4.830069in}{2.306492in}}%
\pgfpathlineto{\pgfqpoint{4.830213in}{2.306221in}}%
\pgfpathlineto{\pgfqpoint{4.830633in}{2.305962in}}%
\pgfpathlineto{\pgfqpoint{4.831329in}{2.306347in}}%
\pgfpathlineto{\pgfqpoint{4.831456in}{2.306337in}}%
\pgfpathlineto{\pgfqpoint{4.832198in}{2.307042in}}%
\pgfpathlineto{\pgfqpoint{4.832373in}{2.306935in}}%
\pgfpathlineto{\pgfqpoint{4.832789in}{2.307134in}}%
\pgfpathlineto{\pgfqpoint{4.833238in}{2.306786in}}%
\pgfpathlineto{\pgfqpoint{4.833487in}{2.306959in}}%
\pgfpathlineto{\pgfqpoint{4.834668in}{2.306524in}}%
\pgfpathlineto{\pgfqpoint{4.833878in}{2.307387in}}%
\pgfpathlineto{\pgfqpoint{4.834715in}{2.306629in}}%
\pgfpathlineto{\pgfqpoint{4.834932in}{2.306731in}}%
\pgfpathlineto{\pgfqpoint{4.835487in}{2.306265in}}%
\pgfpathlineto{\pgfqpoint{4.835805in}{2.306389in}}%
\pgfpathlineto{\pgfqpoint{4.835806in}{2.306359in}}%
\pgfpathlineto{\pgfqpoint{4.836509in}{2.306775in}}%
\pgfpathlineto{\pgfqpoint{4.836907in}{2.306568in}}%
\pgfpathlineto{\pgfqpoint{4.837282in}{2.306917in}}%
\pgfpathlineto{\pgfqpoint{4.836974in}{2.306524in}}%
\pgfpathlineto{\pgfqpoint{4.838044in}{2.306611in}}%
\pgfpathlineto{\pgfqpoint{4.838397in}{2.306376in}}%
\pgfpathlineto{\pgfqpoint{4.839113in}{2.306993in}}%
\pgfpathlineto{\pgfqpoint{4.839139in}{2.306932in}}%
\pgfpathlineto{\pgfqpoint{4.840199in}{2.306549in}}%
\pgfpathlineto{\pgfqpoint{4.839277in}{2.307033in}}%
\pgfpathlineto{\pgfqpoint{4.840301in}{2.306799in}}%
\pgfpathlineto{\pgfqpoint{4.841277in}{2.307163in}}%
\pgfpathlineto{\pgfqpoint{4.840333in}{2.306724in}}%
\pgfpathlineto{\pgfqpoint{4.841445in}{2.307010in}}%
\pgfpathlineto{\pgfqpoint{4.841712in}{2.306809in}}%
\pgfpathlineto{\pgfqpoint{4.842398in}{2.307585in}}%
\pgfpathlineto{\pgfqpoint{4.842517in}{2.307396in}}%
\pgfpathlineto{\pgfqpoint{4.843266in}{2.307962in}}%
\pgfpathlineto{\pgfqpoint{4.842590in}{2.307318in}}%
\pgfpathlineto{\pgfqpoint{4.843640in}{2.307523in}}%
\pgfpathlineto{\pgfqpoint{4.844687in}{2.307282in}}%
\pgfpathlineto{\pgfqpoint{4.843934in}{2.307712in}}%
\pgfpathlineto{\pgfqpoint{4.844760in}{2.307499in}}%
\pgfpathlineto{\pgfqpoint{4.846012in}{2.307764in}}%
\pgfpathlineto{\pgfqpoint{4.844878in}{2.307739in}}%
\pgfpathlineto{\pgfqpoint{4.846090in}{2.307965in}}%
\pgfpathlineto{\pgfqpoint{4.847076in}{2.308301in}}%
\pgfpathlineto{\pgfqpoint{4.846128in}{2.307901in}}%
\pgfpathlineto{\pgfqpoint{4.847216in}{2.308078in}}%
\pgfpathlineto{\pgfqpoint{4.847392in}{2.307868in}}%
\pgfpathlineto{\pgfqpoint{4.847566in}{2.308384in}}%
\pgfpathlineto{\pgfqpoint{4.848237in}{2.308512in}}%
\pgfpathlineto{\pgfqpoint{4.848988in}{2.309240in}}%
\pgfpathlineto{\pgfqpoint{4.848597in}{2.308402in}}%
\pgfpathlineto{\pgfqpoint{4.849375in}{2.309112in}}%
\pgfpathlineto{\pgfqpoint{4.849460in}{2.309294in}}%
\pgfpathlineto{\pgfqpoint{4.850543in}{2.308549in}}%
\pgfpathlineto{\pgfqpoint{4.850800in}{2.308653in}}%
\pgfpathlineto{\pgfqpoint{4.851255in}{2.308140in}}%
\pgfpathlineto{\pgfqpoint{4.851641in}{2.308397in}}%
\pgfpathlineto{\pgfqpoint{4.851906in}{2.308293in}}%
\pgfpathlineto{\pgfqpoint{4.852725in}{2.308872in}}%
\pgfpathlineto{\pgfqpoint{4.852729in}{2.308868in}}%
\pgfpathlineto{\pgfqpoint{4.853597in}{2.309306in}}%
\pgfpathlineto{\pgfqpoint{4.852760in}{2.308772in}}%
\pgfpathlineto{\pgfqpoint{4.853937in}{2.308931in}}%
\pgfpathlineto{\pgfqpoint{4.854786in}{2.308462in}}%
\pgfpathlineto{\pgfqpoint{4.854326in}{2.309022in}}%
\pgfpathlineto{\pgfqpoint{4.855052in}{2.308852in}}%
\pgfpathlineto{\pgfqpoint{4.855593in}{2.309238in}}%
\pgfpathlineto{\pgfqpoint{4.855076in}{2.308826in}}%
\pgfpathlineto{\pgfqpoint{4.856175in}{2.308957in}}%
\pgfpathlineto{\pgfqpoint{4.857039in}{2.308560in}}%
\pgfpathlineto{\pgfqpoint{4.857359in}{2.308717in}}%
\pgfpathlineto{\pgfqpoint{4.858524in}{2.309227in}}%
\pgfpathlineto{\pgfqpoint{4.857539in}{2.308340in}}%
\pgfpathlineto{\pgfqpoint{4.858561in}{2.309135in}}%
\pgfpathlineto{\pgfqpoint{4.858948in}{2.308945in}}%
\pgfpathlineto{\pgfqpoint{4.859506in}{2.309440in}}%
\pgfpathlineto{\pgfqpoint{4.859664in}{2.309364in}}%
\pgfpathlineto{\pgfqpoint{4.860083in}{2.309648in}}%
\pgfpathlineto{\pgfqpoint{4.860746in}{2.309131in}}%
\pgfpathlineto{\pgfqpoint{4.860758in}{2.309141in}}%
\pgfpathlineto{\pgfqpoint{4.861086in}{2.309022in}}%
\pgfpathlineto{\pgfqpoint{4.861320in}{2.308368in}}%
\pgfpathlineto{\pgfqpoint{4.861721in}{2.308556in}}%
\pgfpathlineto{\pgfqpoint{4.861802in}{2.308394in}}%
\pgfpathlineto{\pgfqpoint{4.862359in}{2.309177in}}%
\pgfpathlineto{\pgfqpoint{4.862779in}{2.308916in}}%
\pgfpathlineto{\pgfqpoint{4.862836in}{2.309022in}}%
\pgfpathlineto{\pgfqpoint{4.863254in}{2.308539in}}%
\pgfpathlineto{\pgfqpoint{4.863879in}{2.308700in}}%
\pgfpathlineto{\pgfqpoint{4.863952in}{2.308740in}}%
\pgfpathlineto{\pgfqpoint{4.864442in}{2.306001in}}%
\pgfpathlineto{\pgfqpoint{4.865525in}{2.303354in}}%
\pgfpathlineto{\pgfqpoint{4.865725in}{2.303396in}}%
\pgfpathlineto{\pgfqpoint{4.866787in}{2.303454in}}%
\pgfpathlineto{\pgfqpoint{4.866441in}{2.303325in}}%
\pgfpathlineto{\pgfqpoint{4.866838in}{2.303393in}}%
\pgfpathlineto{\pgfqpoint{4.867837in}{2.303455in}}%
\pgfpathlineto{\pgfqpoint{4.866858in}{2.303378in}}%
\pgfpathlineto{\pgfqpoint{4.867951in}{2.303426in}}%
\pgfpathlineto{\pgfqpoint{4.867961in}{2.303342in}}%
\pgfpathlineto{\pgfqpoint{4.867985in}{2.303449in}}%
\pgfpathlineto{\pgfqpoint{4.869061in}{2.303415in}}%
\pgfpathlineto{\pgfqpoint{4.870175in}{2.303442in}}%
\pgfpathlineto{\pgfqpoint{4.869097in}{2.303375in}}%
\pgfpathlineto{\pgfqpoint{4.870178in}{2.303426in}}%
\pgfpathlineto{\pgfqpoint{4.871263in}{2.303363in}}%
\pgfpathlineto{\pgfqpoint{4.870343in}{2.303448in}}%
\pgfpathlineto{\pgfqpoint{4.871290in}{2.303423in}}%
\pgfpathlineto{\pgfqpoint{4.872413in}{2.303439in}}%
\pgfpathlineto{\pgfqpoint{4.871397in}{2.303346in}}%
\pgfpathlineto{\pgfqpoint{4.872414in}{2.303430in}}%
\pgfpathlineto{\pgfqpoint{4.873509in}{2.303360in}}%
\pgfpathlineto{\pgfqpoint{4.872587in}{2.303451in}}%
\pgfpathlineto{\pgfqpoint{4.873525in}{2.303405in}}%
\pgfpathlineto{\pgfqpoint{4.874551in}{2.303452in}}%
\pgfpathlineto{\pgfqpoint{4.873661in}{2.303360in}}%
\pgfpathlineto{\pgfqpoint{4.874647in}{2.303433in}}%
\pgfpathlineto{\pgfqpoint{4.875474in}{2.303341in}}%
\pgfpathlineto{\pgfqpoint{4.874689in}{2.303446in}}%
\pgfpathlineto{\pgfqpoint{4.875759in}{2.303394in}}%
\pgfpathlineto{\pgfqpoint{4.876891in}{2.303375in}}%
\pgfpathlineto{\pgfqpoint{4.875766in}{2.303444in}}%
\pgfpathlineto{\pgfqpoint{4.876900in}{2.303420in}}%
\pgfpathlineto{\pgfqpoint{4.877990in}{2.303442in}}%
\pgfpathlineto{\pgfqpoint{4.877124in}{2.303372in}}%
\pgfpathlineto{\pgfqpoint{4.878013in}{2.303428in}}%
\pgfpathlineto{\pgfqpoint{4.879064in}{2.303361in}}%
\pgfpathlineto{\pgfqpoint{4.878086in}{2.303447in}}%
\pgfpathlineto{\pgfqpoint{4.879125in}{2.303420in}}%
\pgfpathlineto{\pgfqpoint{4.880189in}{2.303449in}}%
\pgfpathlineto{\pgfqpoint{4.879136in}{2.303390in}}%
\pgfpathlineto{\pgfqpoint{4.880244in}{2.303414in}}%
\pgfpathlineto{\pgfqpoint{4.881346in}{2.303382in}}%
\pgfpathlineto{\pgfqpoint{4.880359in}{2.303454in}}%
\pgfpathlineto{\pgfqpoint{4.881361in}{2.303419in}}%
\pgfpathlineto{\pgfqpoint{4.882394in}{2.303447in}}%
\pgfpathlineto{\pgfqpoint{4.881393in}{2.303376in}}%
\pgfpathlineto{\pgfqpoint{4.882471in}{2.303393in}}%
\pgfpathlineto{\pgfqpoint{4.883240in}{2.303320in}}%
\pgfpathlineto{\pgfqpoint{4.883211in}{2.303528in}}%
\pgfpathlineto{\pgfqpoint{4.883581in}{2.303400in}}%
\pgfpathlineto{\pgfqpoint{4.884195in}{2.303814in}}%
\pgfpathlineto{\pgfqpoint{4.883635in}{2.303341in}}%
\pgfpathlineto{\pgfqpoint{4.884712in}{2.303573in}}%
\pgfpathlineto{\pgfqpoint{4.885670in}{2.303526in}}%
\pgfpathlineto{\pgfqpoint{4.885059in}{2.303929in}}%
\pgfpathlineto{\pgfqpoint{4.885822in}{2.303615in}}%
\pgfpathlineto{\pgfqpoint{4.886630in}{2.304400in}}%
\pgfpathlineto{\pgfqpoint{4.887002in}{2.304232in}}%
\pgfpathlineto{\pgfqpoint{4.887034in}{2.304165in}}%
\pgfpathlineto{\pgfqpoint{4.887293in}{2.304636in}}%
\pgfpathlineto{\pgfqpoint{4.888082in}{2.304584in}}%
\pgfpathlineto{\pgfqpoint{4.888484in}{2.304914in}}%
\pgfpathlineto{\pgfqpoint{4.888096in}{2.304522in}}%
\pgfpathlineto{\pgfqpoint{4.889185in}{2.304471in}}%
\pgfpathlineto{\pgfqpoint{4.889208in}{2.304396in}}%
\pgfpathlineto{\pgfqpoint{4.890041in}{2.305056in}}%
\pgfpathlineto{\pgfqpoint{4.890285in}{2.304740in}}%
\pgfpathlineto{\pgfqpoint{4.890769in}{2.304998in}}%
\pgfpathlineto{\pgfqpoint{4.890295in}{2.304707in}}%
\pgfpathlineto{\pgfqpoint{4.891420in}{2.304855in}}%
\pgfpathlineto{\pgfqpoint{4.894867in}{2.305830in}}%
\pgfpathlineto{\pgfqpoint{4.895691in}{2.306241in}}%
\pgfpathlineto{\pgfqpoint{4.894984in}{2.305782in}}%
\pgfpathlineto{\pgfqpoint{4.895983in}{2.306000in}}%
\pgfpathlineto{\pgfqpoint{4.896768in}{2.305709in}}%
\pgfpathlineto{\pgfqpoint{4.896043in}{2.306051in}}%
\pgfpathlineto{\pgfqpoint{4.897094in}{2.305999in}}%
\pgfpathlineto{\pgfqpoint{4.898156in}{2.306279in}}%
\pgfpathlineto{\pgfqpoint{4.897541in}{2.305798in}}%
\pgfpathlineto{\pgfqpoint{4.898220in}{2.306029in}}%
\pgfpathlineto{\pgfqpoint{4.898405in}{2.305782in}}%
\pgfpathlineto{\pgfqpoint{4.898660in}{2.306348in}}%
\pgfpathlineto{\pgfqpoint{4.899331in}{2.306052in}}%
\pgfpathlineto{\pgfqpoint{4.899558in}{2.306086in}}%
\pgfpathlineto{\pgfqpoint{4.900098in}{2.305358in}}%
\pgfpathlineto{\pgfqpoint{4.900382in}{2.305548in}}%
\pgfpathlineto{\pgfqpoint{4.900605in}{2.305387in}}%
\pgfpathlineto{\pgfqpoint{4.901423in}{2.306358in}}%
\pgfpathlineto{\pgfqpoint{4.901430in}{2.306332in}}%
\pgfpathlineto{\pgfqpoint{4.901876in}{2.306564in}}%
\pgfpathlineto{\pgfqpoint{4.901554in}{2.306129in}}%
\pgfpathlineto{\pgfqpoint{4.902539in}{2.306288in}}%
\pgfpathlineto{\pgfqpoint{4.902750in}{2.306048in}}%
\pgfpathlineto{\pgfqpoint{4.903521in}{2.306624in}}%
\pgfpathlineto{\pgfqpoint{4.903643in}{2.306417in}}%
\pgfpathlineto{\pgfqpoint{4.904777in}{2.306741in}}%
\pgfpathlineto{\pgfqpoint{4.904258in}{2.306056in}}%
\pgfpathlineto{\pgfqpoint{4.904798in}{2.306699in}}%
\pgfpathlineto{\pgfqpoint{4.905191in}{2.306509in}}%
\pgfpathlineto{\pgfqpoint{4.905798in}{2.306941in}}%
\pgfpathlineto{\pgfqpoint{4.905888in}{2.306936in}}%
\pgfpathlineto{\pgfqpoint{4.906910in}{2.307622in}}%
\pgfpathlineto{\pgfqpoint{4.905891in}{2.306933in}}%
\pgfpathlineto{\pgfqpoint{4.907036in}{2.307537in}}%
\pgfpathlineto{\pgfqpoint{4.907199in}{2.307686in}}%
\pgfpathlineto{\pgfqpoint{4.907633in}{2.307125in}}%
\pgfpathlineto{\pgfqpoint{4.908148in}{2.307441in}}%
\pgfpathlineto{\pgfqpoint{4.908762in}{2.306612in}}%
\pgfpathlineto{\pgfqpoint{4.908166in}{2.307539in}}%
\pgfpathlineto{\pgfqpoint{4.909288in}{2.307060in}}%
\pgfpathlineto{\pgfqpoint{4.909506in}{2.307232in}}%
\pgfpathlineto{\pgfqpoint{4.909943in}{2.306716in}}%
\pgfpathlineto{\pgfqpoint{4.910395in}{2.306949in}}%
\pgfpathlineto{\pgfqpoint{4.910444in}{2.306903in}}%
\pgfpathlineto{\pgfqpoint{4.910888in}{2.307426in}}%
\pgfpathlineto{\pgfqpoint{4.911501in}{2.307061in}}%
\pgfpathlineto{\pgfqpoint{4.918966in}{2.306841in}}%
\pgfpathlineto{\pgfqpoint{4.911507in}{2.307143in}}%
\pgfpathlineto{\pgfqpoint{4.918969in}{2.306877in}}%
\pgfpathlineto{\pgfqpoint{4.919345in}{2.307105in}}%
\pgfpathlineto{\pgfqpoint{4.919712in}{2.306576in}}%
\pgfpathlineto{\pgfqpoint{4.920065in}{2.306701in}}%
\pgfpathlineto{\pgfqpoint{4.921168in}{2.306495in}}%
\pgfpathlineto{\pgfqpoint{4.920341in}{2.307089in}}%
\pgfpathlineto{\pgfqpoint{4.921195in}{2.306611in}}%
\pgfpathlineto{\pgfqpoint{4.922266in}{2.306423in}}%
\pgfpathlineto{\pgfqpoint{4.921922in}{2.306014in}}%
\pgfpathlineto{\pgfqpoint{4.922278in}{2.306340in}}%
\pgfpathlineto{\pgfqpoint{4.922441in}{2.305993in}}%
\pgfpathlineto{\pgfqpoint{4.923325in}{2.306672in}}%
\pgfpathlineto{\pgfqpoint{4.923384in}{2.306536in}}%
\pgfpathlineto{\pgfqpoint{4.924447in}{2.306915in}}%
\pgfpathlineto{\pgfqpoint{4.923540in}{2.306220in}}%
\pgfpathlineto{\pgfqpoint{4.924505in}{2.306751in}}%
\pgfpathlineto{\pgfqpoint{4.925284in}{2.307209in}}%
\pgfpathlineto{\pgfqpoint{4.924720in}{2.306585in}}%
\pgfpathlineto{\pgfqpoint{4.925643in}{2.306995in}}%
\pgfpathlineto{\pgfqpoint{4.926643in}{2.306468in}}%
\pgfpathlineto{\pgfqpoint{4.925649in}{2.307013in}}%
\pgfpathlineto{\pgfqpoint{4.926788in}{2.306675in}}%
\pgfpathlineto{\pgfqpoint{4.927181in}{2.307158in}}%
\pgfpathlineto{\pgfqpoint{4.926883in}{2.306540in}}%
\pgfpathlineto{\pgfqpoint{4.927930in}{2.306882in}}%
\pgfpathlineto{\pgfqpoint{4.932226in}{2.306030in}}%
\pgfpathlineto{\pgfqpoint{4.933047in}{2.305497in}}%
\pgfpathlineto{\pgfqpoint{4.932278in}{2.306166in}}%
\pgfpathlineto{\pgfqpoint{4.933358in}{2.305824in}}%
\pgfpathlineto{\pgfqpoint{4.933364in}{2.305845in}}%
\pgfpathlineto{\pgfqpoint{4.934196in}{2.305116in}}%
\pgfpathlineto{\pgfqpoint{4.934421in}{2.305329in}}%
\pgfpathlineto{\pgfqpoint{4.934455in}{2.305267in}}%
\pgfpathlineto{\pgfqpoint{4.935222in}{2.306096in}}%
\pgfpathlineto{\pgfqpoint{4.935507in}{2.305820in}}%
\pgfpathlineto{\pgfqpoint{4.935802in}{2.305910in}}%
\pgfpathlineto{\pgfqpoint{4.936465in}{2.305322in}}%
\pgfpathlineto{\pgfqpoint{4.936566in}{2.305417in}}%
\pgfpathlineto{\pgfqpoint{4.937716in}{2.304547in}}%
\pgfpathlineto{\pgfqpoint{4.936698in}{2.305447in}}%
\pgfpathlineto{\pgfqpoint{4.937867in}{2.304607in}}%
\pgfpathlineto{\pgfqpoint{4.937876in}{2.304647in}}%
\pgfpathlineto{\pgfqpoint{4.938399in}{2.303913in}}%
\pgfpathlineto{\pgfqpoint{4.938968in}{2.304261in}}%
\pgfpathlineto{\pgfqpoint{4.940172in}{2.303360in}}%
\pgfpathlineto{\pgfqpoint{4.938972in}{2.304291in}}%
\pgfpathlineto{\pgfqpoint{4.940209in}{2.303414in}}%
\pgfpathlineto{\pgfqpoint{4.941200in}{2.303447in}}%
\pgfpathlineto{\pgfqpoint{4.940311in}{2.303355in}}%
\pgfpathlineto{\pgfqpoint{4.941320in}{2.303434in}}%
\pgfpathlineto{\pgfqpoint{4.942302in}{2.303355in}}%
\pgfpathlineto{\pgfqpoint{4.941481in}{2.303448in}}%
\pgfpathlineto{\pgfqpoint{4.942431in}{2.303411in}}%
\pgfpathlineto{\pgfqpoint{4.943538in}{2.303447in}}%
\pgfpathlineto{\pgfqpoint{4.942503in}{2.303366in}}%
\pgfpathlineto{\pgfqpoint{4.943559in}{2.303411in}}%
\pgfpathlineto{\pgfqpoint{4.944779in}{2.303384in}}%
\pgfpathlineto{\pgfqpoint{4.943563in}{2.303443in}}%
\pgfpathlineto{\pgfqpoint{4.944783in}{2.303410in}}%
\pgfpathlineto{\pgfqpoint{4.945784in}{2.303453in}}%
\pgfpathlineto{\pgfqpoint{4.944849in}{2.303374in}}%
\pgfpathlineto{\pgfqpoint{4.945894in}{2.303392in}}%
\pgfpathlineto{\pgfqpoint{4.947771in}{2.303451in}}%
\pgfpathlineto{\pgfqpoint{4.945905in}{2.303365in}}%
\pgfpathlineto{\pgfqpoint{4.947781in}{2.303379in}}%
\pgfpathlineto{\pgfqpoint{4.948759in}{2.303315in}}%
\pgfpathlineto{\pgfqpoint{4.947878in}{2.303444in}}%
\pgfpathlineto{\pgfqpoint{4.948892in}{2.303392in}}%
\pgfpathlineto{\pgfqpoint{4.949950in}{2.303326in}}%
\pgfpathlineto{\pgfqpoint{4.948941in}{2.303456in}}%
\pgfpathlineto{\pgfqpoint{4.950017in}{2.303408in}}%
\pgfpathlineto{\pgfqpoint{4.951128in}{2.303440in}}%
\pgfpathlineto{\pgfqpoint{4.950033in}{2.303373in}}%
\pgfpathlineto{\pgfqpoint{4.951130in}{2.303426in}}%
\pgfpathlineto{\pgfqpoint{4.952045in}{2.303366in}}%
\pgfpathlineto{\pgfqpoint{4.951323in}{2.303455in}}%
\pgfpathlineto{\pgfqpoint{4.952242in}{2.303400in}}%
\pgfpathlineto{\pgfqpoint{4.953217in}{2.303444in}}%
\pgfpathlineto{\pgfqpoint{4.952618in}{2.303364in}}%
\pgfpathlineto{\pgfqpoint{4.953354in}{2.303424in}}%
\pgfpathlineto{\pgfqpoint{4.953596in}{2.303351in}}%
\pgfpathlineto{\pgfqpoint{4.953453in}{2.303452in}}%
\pgfpathlineto{\pgfqpoint{4.954465in}{2.303403in}}%
\pgfpathlineto{\pgfqpoint{4.955566in}{2.303381in}}%
\pgfpathlineto{\pgfqpoint{4.954513in}{2.303443in}}%
\pgfpathlineto{\pgfqpoint{4.955608in}{2.303418in}}%
\pgfpathlineto{\pgfqpoint{4.956683in}{2.303447in}}%
\pgfpathlineto{\pgfqpoint{4.955689in}{2.303360in}}%
\pgfpathlineto{\pgfqpoint{4.956720in}{2.303414in}}%
\pgfpathlineto{\pgfqpoint{4.957811in}{2.303453in}}%
\pgfpathlineto{\pgfqpoint{4.956847in}{2.303349in}}%
\pgfpathlineto{\pgfqpoint{4.957847in}{2.303388in}}%
\pgfpathlineto{\pgfqpoint{4.958038in}{2.303310in}}%
\pgfpathlineto{\pgfqpoint{4.958618in}{2.303896in}}%
\pgfpathlineto{\pgfqpoint{4.958907in}{2.303802in}}%
\pgfpathlineto{\pgfqpoint{4.960053in}{2.304312in}}%
\pgfpathlineto{\pgfqpoint{4.959377in}{2.303668in}}%
\pgfpathlineto{\pgfqpoint{4.960057in}{2.304286in}}%
\pgfpathlineto{\pgfqpoint{4.960081in}{2.304254in}}%
\pgfpathlineto{\pgfqpoint{4.960834in}{2.305115in}}%
\pgfpathlineto{\pgfqpoint{4.961021in}{2.305041in}}%
\pgfpathlineto{\pgfqpoint{4.962098in}{2.305090in}}%
\pgfpathlineto{\pgfqpoint{4.961851in}{2.304684in}}%
\pgfpathlineto{\pgfqpoint{4.962133in}{2.305034in}}%
\pgfpathlineto{\pgfqpoint{4.963214in}{2.305085in}}%
\pgfpathlineto{\pgfqpoint{4.962916in}{2.304771in}}%
\pgfpathlineto{\pgfqpoint{4.963246in}{2.304990in}}%
\pgfpathlineto{\pgfqpoint{4.964094in}{2.305312in}}%
\pgfpathlineto{\pgfqpoint{4.963336in}{2.304846in}}%
\pgfpathlineto{\pgfqpoint{4.964480in}{2.305025in}}%
\pgfpathlineto{\pgfqpoint{4.964614in}{2.304813in}}%
\pgfpathlineto{\pgfqpoint{4.965500in}{2.305596in}}%
\pgfpathlineto{\pgfqpoint{4.965530in}{2.305526in}}%
\pgfpathlineto{\pgfqpoint{4.966051in}{2.305735in}}%
\pgfpathlineto{\pgfqpoint{4.966334in}{2.305118in}}%
\pgfpathlineto{\pgfqpoint{4.966614in}{2.305152in}}%
\pgfpathlineto{\pgfqpoint{4.967611in}{2.304732in}}%
\pgfpathlineto{\pgfqpoint{4.966825in}{2.305405in}}%
\pgfpathlineto{\pgfqpoint{4.967744in}{2.304819in}}%
\pgfpathlineto{\pgfqpoint{4.967756in}{2.304826in}}%
\pgfpathlineto{\pgfqpoint{4.968618in}{2.304182in}}%
\pgfpathlineto{\pgfqpoint{4.968788in}{2.304317in}}%
\pgfpathlineto{\pgfqpoint{4.969481in}{2.304061in}}%
\pgfpathlineto{\pgfqpoint{4.968949in}{2.304556in}}%
\pgfpathlineto{\pgfqpoint{4.969949in}{2.304228in}}%
\pgfpathlineto{\pgfqpoint{4.970356in}{2.304493in}}%
\pgfpathlineto{\pgfqpoint{4.970683in}{2.304015in}}%
\pgfpathlineto{\pgfqpoint{4.971068in}{2.304445in}}%
\pgfpathlineto{\pgfqpoint{4.971261in}{2.304109in}}%
\pgfpathlineto{\pgfqpoint{4.972177in}{2.304517in}}%
\pgfpathlineto{\pgfqpoint{4.972178in}{2.304514in}}%
\pgfpathlineto{\pgfqpoint{4.973017in}{2.303998in}}%
\pgfpathlineto{\pgfqpoint{4.972195in}{2.304519in}}%
\pgfpathlineto{\pgfqpoint{4.973321in}{2.304217in}}%
\pgfpathlineto{\pgfqpoint{4.973991in}{2.304271in}}%
\pgfpathlineto{\pgfqpoint{4.974398in}{2.303888in}}%
\pgfpathlineto{\pgfqpoint{4.974407in}{2.303923in}}%
\pgfpathlineto{\pgfqpoint{4.975396in}{2.303624in}}%
\pgfpathlineto{\pgfqpoint{4.974729in}{2.304031in}}%
\pgfpathlineto{\pgfqpoint{4.975532in}{2.303774in}}%
\pgfpathlineto{\pgfqpoint{4.975583in}{2.303913in}}%
\pgfpathlineto{\pgfqpoint{4.976116in}{2.303274in}}%
\pgfpathlineto{\pgfqpoint{4.976621in}{2.303463in}}%
\pgfpathlineto{\pgfqpoint{4.977249in}{2.303677in}}%
\pgfpathlineto{\pgfqpoint{4.976783in}{2.303090in}}%
\pgfpathlineto{\pgfqpoint{4.977728in}{2.303351in}}%
\pgfpathlineto{\pgfqpoint{4.978521in}{2.303084in}}%
\pgfpathlineto{\pgfqpoint{4.978148in}{2.303446in}}%
\pgfpathlineto{\pgfqpoint{4.978844in}{2.303231in}}%
\pgfpathlineto{\pgfqpoint{4.979935in}{2.303199in}}%
\pgfpathlineto{\pgfqpoint{4.979438in}{2.302689in}}%
\pgfpathlineto{\pgfqpoint{4.979955in}{2.303165in}}%
\pgfpathlineto{\pgfqpoint{4.980570in}{2.303065in}}%
\pgfpathlineto{\pgfqpoint{4.980271in}{2.303299in}}%
\pgfpathlineto{\pgfqpoint{4.981022in}{2.303474in}}%
\pgfpathlineto{\pgfqpoint{4.981327in}{2.303719in}}%
\pgfpathlineto{\pgfqpoint{4.981716in}{2.303258in}}%
\pgfpathlineto{\pgfqpoint{4.982119in}{2.303314in}}%
\pgfpathlineto{\pgfqpoint{4.982270in}{2.303156in}}%
\pgfpathlineto{\pgfqpoint{4.983127in}{2.303860in}}%
\pgfpathlineto{\pgfqpoint{4.983203in}{2.303831in}}%
\pgfpathlineto{\pgfqpoint{4.983728in}{2.304256in}}%
\pgfpathlineto{\pgfqpoint{4.983233in}{2.303796in}}%
\pgfpathlineto{\pgfqpoint{4.984338in}{2.304080in}}%
\pgfpathlineto{\pgfqpoint{4.984419in}{2.304054in}}%
\pgfpathlineto{\pgfqpoint{4.985110in}{2.304557in}}%
\pgfpathlineto{\pgfqpoint{4.985431in}{2.304335in}}%
\pgfpathlineto{\pgfqpoint{4.986134in}{2.304686in}}%
\pgfpathlineto{\pgfqpoint{4.985823in}{2.304226in}}%
\pgfpathlineto{\pgfqpoint{4.986544in}{2.304313in}}%
\pgfpathlineto{\pgfqpoint{4.986607in}{2.304246in}}%
\pgfpathlineto{\pgfqpoint{4.987531in}{2.304795in}}%
\pgfpathlineto{\pgfqpoint{4.987622in}{2.304654in}}%
\pgfpathlineto{\pgfqpoint{4.988691in}{2.304580in}}%
\pgfpathlineto{\pgfqpoint{4.988143in}{2.304150in}}%
\pgfpathlineto{\pgfqpoint{4.988720in}{2.304486in}}%
\pgfpathlineto{\pgfqpoint{4.990372in}{2.304446in}}%
\pgfpathlineto{\pgfqpoint{4.988725in}{2.304451in}}%
\pgfpathlineto{\pgfqpoint{4.990429in}{2.304292in}}%
\pgfpathlineto{\pgfqpoint{4.991054in}{2.304861in}}%
\pgfpathlineto{\pgfqpoint{4.991542in}{2.304252in}}%
\pgfpathlineto{\pgfqpoint{4.992530in}{2.304517in}}%
\pgfpathlineto{\pgfqpoint{4.991789in}{2.303819in}}%
\pgfpathlineto{\pgfqpoint{4.992659in}{2.304304in}}%
\pgfpathlineto{\pgfqpoint{4.992670in}{2.304261in}}%
\pgfpathlineto{\pgfqpoint{4.993643in}{2.304957in}}%
\pgfpathlineto{\pgfqpoint{4.995130in}{2.304929in}}%
\pgfpathlineto{\pgfqpoint{4.993747in}{2.304736in}}%
\pgfpathlineto{\pgfqpoint{4.995165in}{2.304832in}}%
\pgfpathlineto{\pgfqpoint{4.995741in}{2.304655in}}%
\pgfpathlineto{\pgfqpoint{4.996062in}{2.305171in}}%
\pgfpathlineto{\pgfqpoint{4.996254in}{2.305091in}}%
\pgfpathlineto{\pgfqpoint{4.997353in}{2.305647in}}%
\pgfpathlineto{\pgfqpoint{4.996479in}{2.304865in}}%
\pgfpathlineto{\pgfqpoint{4.997420in}{2.305625in}}%
\pgfpathlineto{\pgfqpoint{4.998557in}{2.305430in}}%
\pgfpathlineto{\pgfqpoint{4.997764in}{2.306053in}}%
\pgfpathlineto{\pgfqpoint{4.998561in}{2.305495in}}%
\pgfpathlineto{\pgfqpoint{4.998573in}{2.305534in}}%
\pgfpathlineto{\pgfqpoint{4.999422in}{2.304804in}}%
\pgfpathlineto{\pgfqpoint{4.999641in}{2.305182in}}%
\pgfpathlineto{\pgfqpoint{5.000479in}{2.305040in}}%
\pgfpathlineto{\pgfqpoint{5.000243in}{2.305460in}}%
\pgfpathlineto{\pgfqpoint{5.000755in}{2.305185in}}%
\pgfpathlineto{\pgfqpoint{5.001513in}{2.305681in}}%
\pgfpathlineto{\pgfqpoint{5.000930in}{2.304947in}}%
\pgfpathlineto{\pgfqpoint{5.001871in}{2.305272in}}%
\pgfpathlineto{\pgfqpoint{5.001899in}{2.305168in}}%
\pgfpathlineto{\pgfqpoint{5.002116in}{2.305535in}}%
\pgfpathlineto{\pgfqpoint{5.002983in}{2.305281in}}%
\pgfpathlineto{\pgfqpoint{5.003299in}{2.305254in}}%
\pgfpathlineto{\pgfqpoint{5.003433in}{2.305693in}}%
\pgfpathlineto{\pgfqpoint{5.003908in}{2.305905in}}%
\pgfpathlineto{\pgfqpoint{5.004089in}{2.305866in}}%
\pgfpathlineto{\pgfqpoint{5.005045in}{2.306248in}}%
\pgfpathlineto{\pgfqpoint{5.005952in}{2.306160in}}%
\pgfpathlineto{\pgfqpoint{5.005618in}{2.306715in}}%
\pgfpathlineto{\pgfqpoint{5.006142in}{2.306435in}}%
\pgfpathlineto{\pgfqpoint{5.006591in}{2.306956in}}%
\pgfpathlineto{\pgfqpoint{5.006150in}{2.306416in}}%
\pgfpathlineto{\pgfqpoint{5.007279in}{2.306417in}}%
\pgfpathlineto{\pgfqpoint{5.008149in}{2.306290in}}%
\pgfpathlineto{\pgfqpoint{5.007451in}{2.306756in}}%
\pgfpathlineto{\pgfqpoint{5.008393in}{2.306412in}}%
\pgfpathlineto{\pgfqpoint{5.009477in}{2.306235in}}%
\pgfpathlineto{\pgfqpoint{5.008444in}{2.306507in}}%
\pgfpathlineto{\pgfqpoint{5.009570in}{2.306345in}}%
\pgfpathlineto{\pgfqpoint{5.010702in}{2.306696in}}%
\pgfpathlineto{\pgfqpoint{5.010188in}{2.305925in}}%
\pgfpathlineto{\pgfqpoint{5.010703in}{2.306680in}}%
\pgfpathlineto{\pgfqpoint{5.011751in}{2.306483in}}%
\pgfpathlineto{\pgfqpoint{5.010763in}{2.306453in}}%
\pgfpathlineto{\pgfqpoint{5.011812in}{2.306360in}}%
\pgfpathlineto{\pgfqpoint{5.012172in}{2.306122in}}%
\pgfpathlineto{\pgfqpoint{5.012780in}{2.306534in}}%
\pgfpathlineto{\pgfqpoint{5.012920in}{2.306432in}}%
\pgfpathlineto{\pgfqpoint{5.014017in}{2.306515in}}%
\pgfpathlineto{\pgfqpoint{5.013693in}{2.306122in}}%
\pgfpathlineto{\pgfqpoint{5.014031in}{2.306407in}}%
\pgfpathlineto{\pgfqpoint{5.015355in}{2.303369in}}%
\pgfpathlineto{\pgfqpoint{5.014122in}{2.306487in}}%
\pgfpathlineto{\pgfqpoint{5.015522in}{2.303417in}}%
\pgfpathlineto{\pgfqpoint{5.016614in}{2.303447in}}%
\pgfpathlineto{\pgfqpoint{5.015534in}{2.303350in}}%
\pgfpathlineto{\pgfqpoint{5.016633in}{2.303435in}}%
\pgfpathlineto{\pgfqpoint{5.017718in}{2.303368in}}%
\pgfpathlineto{\pgfqpoint{5.016934in}{2.303455in}}%
\pgfpathlineto{\pgfqpoint{5.017745in}{2.303409in}}%
\pgfpathlineto{\pgfqpoint{5.018774in}{2.303372in}}%
\pgfpathlineto{\pgfqpoint{5.017764in}{2.303445in}}%
\pgfpathlineto{\pgfqpoint{5.018863in}{2.303401in}}%
\pgfpathlineto{\pgfqpoint{5.019933in}{2.303445in}}%
\pgfpathlineto{\pgfqpoint{5.018957in}{2.303374in}}%
\pgfpathlineto{\pgfqpoint{5.019974in}{2.303412in}}%
\pgfpathlineto{\pgfqpoint{5.020929in}{2.303373in}}%
\pgfpathlineto{\pgfqpoint{5.020046in}{2.303445in}}%
\pgfpathlineto{\pgfqpoint{5.021089in}{2.303415in}}%
\pgfpathlineto{\pgfqpoint{5.022172in}{2.303437in}}%
\pgfpathlineto{\pgfqpoint{5.021099in}{2.303374in}}%
\pgfpathlineto{\pgfqpoint{5.022203in}{2.303427in}}%
\pgfpathlineto{\pgfqpoint{5.023122in}{2.303362in}}%
\pgfpathlineto{\pgfqpoint{5.022253in}{2.303454in}}%
\pgfpathlineto{\pgfqpoint{5.023315in}{2.303417in}}%
\pgfpathlineto{\pgfqpoint{5.024453in}{2.303381in}}%
\pgfpathlineto{\pgfqpoint{5.023317in}{2.303437in}}%
\pgfpathlineto{\pgfqpoint{5.024470in}{2.303421in}}%
\pgfpathlineto{\pgfqpoint{5.025417in}{2.303458in}}%
\pgfpathlineto{\pgfqpoint{5.024515in}{2.303365in}}%
\pgfpathlineto{\pgfqpoint{5.025582in}{2.303428in}}%
\pgfpathlineto{\pgfqpoint{5.026631in}{2.303375in}}%
\pgfpathlineto{\pgfqpoint{5.025711in}{2.303453in}}%
\pgfpathlineto{\pgfqpoint{5.026694in}{2.303423in}}%
\pgfpathlineto{\pgfqpoint{5.027589in}{2.303449in}}%
\pgfpathlineto{\pgfqpoint{5.027488in}{2.303341in}}%
\pgfpathlineto{\pgfqpoint{5.027805in}{2.303430in}}%
\pgfpathlineto{\pgfqpoint{5.028139in}{2.303353in}}%
\pgfpathlineto{\pgfqpoint{5.027977in}{2.303461in}}%
\pgfpathlineto{\pgfqpoint{5.028917in}{2.303403in}}%
\pgfpathlineto{\pgfqpoint{5.029911in}{2.303444in}}%
\pgfpathlineto{\pgfqpoint{5.029018in}{2.303374in}}%
\pgfpathlineto{\pgfqpoint{5.030031in}{2.303419in}}%
\pgfpathlineto{\pgfqpoint{5.031162in}{2.303387in}}%
\pgfpathlineto{\pgfqpoint{5.030048in}{2.303443in}}%
\pgfpathlineto{\pgfqpoint{5.031191in}{2.303432in}}%
\pgfpathlineto{\pgfqpoint{5.032305in}{2.303440in}}%
\pgfpathlineto{\pgfqpoint{5.031212in}{2.303388in}}%
\pgfpathlineto{\pgfqpoint{5.032321in}{2.303429in}}%
\pgfpathlineto{\pgfqpoint{5.033434in}{2.303097in}}%
\pgfpathlineto{\pgfqpoint{5.032899in}{2.303474in}}%
\pgfpathlineto{\pgfqpoint{5.033444in}{2.303109in}}%
\pgfpathlineto{\pgfqpoint{5.034628in}{2.304034in}}%
\pgfpathlineto{\pgfqpoint{5.033527in}{2.302895in}}%
\pgfpathlineto{\pgfqpoint{5.034677in}{2.303844in}}%
\pgfpathlineto{\pgfqpoint{5.035442in}{2.303669in}}%
\pgfpathlineto{\pgfqpoint{5.035780in}{2.304148in}}%
\pgfpathlineto{\pgfqpoint{5.035781in}{2.304128in}}%
\pgfpathlineto{\pgfqpoint{5.035817in}{2.304071in}}%
\pgfpathlineto{\pgfqpoint{5.036551in}{2.304563in}}%
\pgfpathlineto{\pgfqpoint{5.036876in}{2.304319in}}%
\pgfpathlineto{\pgfqpoint{5.036965in}{2.304434in}}%
\pgfpathlineto{\pgfqpoint{5.037840in}{2.303243in}}%
\pgfpathlineto{\pgfqpoint{5.037943in}{2.303487in}}%
\pgfpathlineto{\pgfqpoint{5.038149in}{2.303670in}}%
\pgfpathlineto{\pgfqpoint{5.038944in}{2.302978in}}%
\pgfpathlineto{\pgfqpoint{5.038997in}{2.303041in}}%
\pgfpathlineto{\pgfqpoint{5.039000in}{2.303013in}}%
\pgfpathlineto{\pgfqpoint{5.039658in}{2.303562in}}%
\pgfpathlineto{\pgfqpoint{5.040098in}{2.303224in}}%
\pgfpathlineto{\pgfqpoint{5.041064in}{2.303735in}}%
\pgfpathlineto{\pgfqpoint{5.040148in}{2.303092in}}%
\pgfpathlineto{\pgfqpoint{5.041236in}{2.303562in}}%
\pgfpathlineto{\pgfqpoint{5.041534in}{2.303353in}}%
\pgfpathlineto{\pgfqpoint{5.042304in}{2.303866in}}%
\pgfpathlineto{\pgfqpoint{5.042328in}{2.303856in}}%
\pgfpathlineto{\pgfqpoint{5.043137in}{2.304491in}}%
\pgfpathlineto{\pgfqpoint{5.043468in}{2.304251in}}%
\pgfpathlineto{\pgfqpoint{5.044703in}{2.303731in}}%
\pgfpathlineto{\pgfqpoint{5.043554in}{2.304459in}}%
\pgfpathlineto{\pgfqpoint{5.044712in}{2.303763in}}%
\pgfpathlineto{\pgfqpoint{5.045266in}{2.303948in}}%
\pgfpathlineto{\pgfqpoint{5.045685in}{2.303475in}}%
\pgfpathlineto{\pgfqpoint{5.045820in}{2.303684in}}%
\pgfpathlineto{\pgfqpoint{5.045837in}{2.303663in}}%
\pgfpathlineto{\pgfqpoint{5.046694in}{2.304214in}}%
\pgfpathlineto{\pgfqpoint{5.046842in}{2.304211in}}%
\pgfpathlineto{\pgfqpoint{5.046857in}{2.304250in}}%
\pgfpathlineto{\pgfqpoint{5.047906in}{2.303591in}}%
\pgfpathlineto{\pgfqpoint{5.047920in}{2.303679in}}%
\pgfpathlineto{\pgfqpoint{5.047977in}{2.303785in}}%
\pgfpathlineto{\pgfqpoint{5.049066in}{2.302864in}}%
\pgfpathlineto{\pgfqpoint{5.050268in}{2.302364in}}%
\pgfpathlineto{\pgfqpoint{5.049075in}{2.302934in}}%
\pgfpathlineto{\pgfqpoint{5.050285in}{2.302432in}}%
\pgfpathlineto{\pgfqpoint{5.050961in}{2.302925in}}%
\pgfpathlineto{\pgfqpoint{5.050383in}{2.302314in}}%
\pgfpathlineto{\pgfqpoint{5.051399in}{2.302478in}}%
\pgfpathlineto{\pgfqpoint{5.055487in}{2.302913in}}%
\pgfpathlineto{\pgfqpoint{5.051401in}{2.302493in}}%
\pgfpathlineto{\pgfqpoint{5.055492in}{2.302949in}}%
\pgfpathlineto{\pgfqpoint{5.056449in}{2.302997in}}%
\pgfpathlineto{\pgfqpoint{5.055858in}{2.302725in}}%
\pgfpathlineto{\pgfqpoint{5.056595in}{2.302837in}}%
\pgfpathlineto{\pgfqpoint{5.056717in}{2.302641in}}%
\pgfpathlineto{\pgfqpoint{5.057468in}{2.303118in}}%
\pgfpathlineto{\pgfqpoint{5.057699in}{2.303099in}}%
\pgfpathlineto{\pgfqpoint{5.058830in}{2.302825in}}%
\pgfpathlineto{\pgfqpoint{5.057744in}{2.303178in}}%
\pgfpathlineto{\pgfqpoint{5.058831in}{2.302841in}}%
\pgfpathlineto{\pgfqpoint{5.059582in}{2.303240in}}%
\pgfpathlineto{\pgfqpoint{5.058985in}{2.302723in}}%
\pgfpathlineto{\pgfqpoint{5.059958in}{2.302930in}}%
\pgfpathlineto{\pgfqpoint{5.061002in}{2.302490in}}%
\pgfpathlineto{\pgfqpoint{5.060024in}{2.303042in}}%
\pgfpathlineto{\pgfqpoint{5.061140in}{2.302532in}}%
\pgfpathlineto{\pgfqpoint{5.061982in}{2.302686in}}%
\pgfpathlineto{\pgfqpoint{5.061420in}{2.302228in}}%
\pgfpathlineto{\pgfqpoint{5.062252in}{2.302552in}}%
\pgfpathlineto{\pgfqpoint{5.062575in}{2.302744in}}%
\pgfpathlineto{\pgfqpoint{5.062863in}{2.302383in}}%
\pgfpathlineto{\pgfqpoint{5.063375in}{2.302505in}}%
\pgfpathlineto{\pgfqpoint{5.063864in}{2.302375in}}%
\pgfpathlineto{\pgfqpoint{5.064465in}{2.302721in}}%
\pgfpathlineto{\pgfqpoint{5.064480in}{2.302656in}}%
\pgfpathlineto{\pgfqpoint{5.065421in}{2.302836in}}%
\pgfpathlineto{\pgfqpoint{5.064517in}{2.302539in}}%
\pgfpathlineto{\pgfqpoint{5.065617in}{2.302482in}}%
\pgfpathlineto{\pgfqpoint{5.066722in}{2.302351in}}%
\pgfpathlineto{\pgfqpoint{5.066788in}{2.302043in}}%
\pgfpathlineto{\pgfqpoint{5.067176in}{2.301939in}}%
\pgfpathlineto{\pgfqpoint{5.067336in}{2.302350in}}%
\pgfpathlineto{\pgfqpoint{5.067887in}{2.302216in}}%
\pgfpathlineto{\pgfqpoint{5.068558in}{2.302519in}}%
\pgfpathlineto{\pgfqpoint{5.067987in}{2.302127in}}%
\pgfpathlineto{\pgfqpoint{5.069016in}{2.302354in}}%
\pgfpathlineto{\pgfqpoint{5.069673in}{2.301984in}}%
\pgfpathlineto{\pgfqpoint{5.069216in}{2.302468in}}%
\pgfpathlineto{\pgfqpoint{5.070133in}{2.302147in}}%
\pgfpathlineto{\pgfqpoint{5.071093in}{2.302162in}}%
\pgfpathlineto{\pgfqpoint{5.070139in}{2.302125in}}%
\pgfpathlineto{\pgfqpoint{5.071219in}{2.301890in}}%
\pgfpathlineto{\pgfqpoint{5.071891in}{2.301162in}}%
\pgfpathlineto{\pgfqpoint{5.071221in}{2.301898in}}%
\pgfpathlineto{\pgfqpoint{5.072370in}{2.301456in}}%
\pgfpathlineto{\pgfqpoint{5.072830in}{2.301739in}}%
\pgfpathlineto{\pgfqpoint{5.073344in}{2.301083in}}%
\pgfpathlineto{\pgfqpoint{5.073468in}{2.301128in}}%
\pgfpathlineto{\pgfqpoint{5.073944in}{2.300939in}}%
\pgfpathlineto{\pgfqpoint{5.073490in}{2.301183in}}%
\pgfpathlineto{\pgfqpoint{5.074529in}{2.301476in}}%
\pgfpathlineto{\pgfqpoint{5.074550in}{2.301536in}}%
\pgfpathlineto{\pgfqpoint{5.074929in}{2.301035in}}%
\pgfpathlineto{\pgfqpoint{5.075631in}{2.301276in}}%
\pgfpathlineto{\pgfqpoint{5.076576in}{2.301523in}}%
\pgfpathlineto{\pgfqpoint{5.076760in}{2.301122in}}%
\pgfpathlineto{\pgfqpoint{5.077595in}{2.301259in}}%
\pgfpathlineto{\pgfqpoint{5.077414in}{2.300869in}}%
\pgfpathlineto{\pgfqpoint{5.077870in}{2.301071in}}%
\pgfpathlineto{\pgfqpoint{5.078126in}{2.300876in}}%
\pgfpathlineto{\pgfqpoint{5.078939in}{2.301494in}}%
\pgfpathlineto{\pgfqpoint{5.078949in}{2.301482in}}%
\pgfpathlineto{\pgfqpoint{5.079967in}{2.301512in}}%
\pgfpathlineto{\pgfqpoint{5.079422in}{2.301241in}}%
\pgfpathlineto{\pgfqpoint{5.080052in}{2.301364in}}%
\pgfpathlineto{\pgfqpoint{5.080164in}{2.301194in}}%
\pgfpathlineto{\pgfqpoint{5.080963in}{2.302020in}}%
\pgfpathlineto{\pgfqpoint{5.081143in}{2.301832in}}%
\pgfpathlineto{\pgfqpoint{5.081395in}{2.301984in}}%
\pgfpathlineto{\pgfqpoint{5.082099in}{2.301411in}}%
\pgfpathlineto{\pgfqpoint{5.082232in}{2.301475in}}%
\pgfpathlineto{\pgfqpoint{5.083294in}{2.300788in}}%
\pgfpathlineto{\pgfqpoint{5.082280in}{2.301484in}}%
\pgfpathlineto{\pgfqpoint{5.083359in}{2.300875in}}%
\pgfpathlineto{\pgfqpoint{5.083395in}{2.300906in}}%
\pgfpathlineto{\pgfqpoint{5.083692in}{2.300357in}}%
\pgfpathlineto{\pgfqpoint{5.084445in}{2.300490in}}%
\pgfpathlineto{\pgfqpoint{5.085624in}{2.300230in}}%
\pgfpathlineto{\pgfqpoint{5.084478in}{2.300262in}}%
\pgfpathlineto{\pgfqpoint{5.085669in}{2.300000in}}%
\pgfpathlineto{\pgfqpoint{5.086790in}{2.299776in}}%
\pgfpathlineto{\pgfqpoint{5.085935in}{2.300372in}}%
\pgfpathlineto{\pgfqpoint{5.086792in}{2.299801in}}%
\pgfpathlineto{\pgfqpoint{5.088023in}{2.300680in}}%
\pgfpathlineto{\pgfqpoint{5.086955in}{2.299599in}}%
\pgfpathlineto{\pgfqpoint{5.088028in}{2.300663in}}%
\pgfpathlineto{\pgfqpoint{5.089022in}{2.300106in}}%
\pgfpathlineto{\pgfqpoint{5.088037in}{2.300719in}}%
\pgfpathlineto{\pgfqpoint{5.089136in}{2.300723in}}%
\pgfpathlineto{\pgfqpoint{5.090309in}{2.303449in}}%
\pgfpathlineto{\pgfqpoint{5.090397in}{2.303412in}}%
\pgfpathlineto{\pgfqpoint{5.091488in}{2.303446in}}%
\pgfpathlineto{\pgfqpoint{5.090416in}{2.303387in}}%
\pgfpathlineto{\pgfqpoint{5.091521in}{2.303433in}}%
\pgfpathlineto{\pgfqpoint{5.092631in}{2.303347in}}%
\pgfpathlineto{\pgfqpoint{5.091572in}{2.303451in}}%
\pgfpathlineto{\pgfqpoint{5.092633in}{2.303367in}}%
\pgfpathlineto{\pgfqpoint{5.093716in}{2.303449in}}%
\pgfpathlineto{\pgfqpoint{5.093747in}{2.303415in}}%
\pgfpathlineto{\pgfqpoint{5.094729in}{2.303356in}}%
\pgfpathlineto{\pgfqpoint{5.093755in}{2.303434in}}%
\pgfpathlineto{\pgfqpoint{5.094867in}{2.303415in}}%
\pgfpathlineto{\pgfqpoint{5.095923in}{2.303455in}}%
\pgfpathlineto{\pgfqpoint{5.094903in}{2.303378in}}%
\pgfpathlineto{\pgfqpoint{5.095982in}{2.303433in}}%
\pgfpathlineto{\pgfqpoint{5.096994in}{2.303377in}}%
\pgfpathlineto{\pgfqpoint{5.096548in}{2.303462in}}%
\pgfpathlineto{\pgfqpoint{5.097094in}{2.303393in}}%
\pgfpathlineto{\pgfqpoint{5.098144in}{2.303370in}}%
\pgfpathlineto{\pgfqpoint{5.097124in}{2.303434in}}%
\pgfpathlineto{\pgfqpoint{5.098215in}{2.303417in}}%
\pgfpathlineto{\pgfqpoint{5.099367in}{2.303442in}}%
\pgfpathlineto{\pgfqpoint{5.098243in}{2.303394in}}%
\pgfpathlineto{\pgfqpoint{5.099374in}{2.303399in}}%
\pgfpathlineto{\pgfqpoint{5.100425in}{2.303376in}}%
\pgfpathlineto{\pgfqpoint{5.099640in}{2.303457in}}%
\pgfpathlineto{\pgfqpoint{5.100486in}{2.303415in}}%
\pgfpathlineto{\pgfqpoint{5.101532in}{2.303448in}}%
\pgfpathlineto{\pgfqpoint{5.100515in}{2.303359in}}%
\pgfpathlineto{\pgfqpoint{5.101598in}{2.303411in}}%
\pgfpathlineto{\pgfqpoint{5.102744in}{2.303375in}}%
\pgfpathlineto{\pgfqpoint{5.101631in}{2.303446in}}%
\pgfpathlineto{\pgfqpoint{5.102751in}{2.303413in}}%
\pgfpathlineto{\pgfqpoint{5.103809in}{2.303446in}}%
\pgfpathlineto{\pgfqpoint{5.103020in}{2.303359in}}%
\pgfpathlineto{\pgfqpoint{5.103863in}{2.303425in}}%
\pgfpathlineto{\pgfqpoint{5.104969in}{2.303369in}}%
\pgfpathlineto{\pgfqpoint{5.103920in}{2.303446in}}%
\pgfpathlineto{\pgfqpoint{5.104975in}{2.303413in}}%
\pgfpathlineto{\pgfqpoint{5.106177in}{2.303443in}}%
\pgfpathlineto{\pgfqpoint{5.105014in}{2.303382in}}%
\pgfpathlineto{\pgfqpoint{5.106180in}{2.303426in}}%
\pgfpathlineto{\pgfqpoint{5.106643in}{2.303358in}}%
\pgfpathlineto{\pgfqpoint{5.106375in}{2.303450in}}%
\pgfpathlineto{\pgfqpoint{5.107292in}{2.303409in}}%
\pgfpathlineto{\pgfqpoint{5.108417in}{2.303201in}}%
\pgfpathlineto{\pgfqpoint{5.107878in}{2.303621in}}%
\pgfpathlineto{\pgfqpoint{5.108419in}{2.303211in}}%
\pgfpathlineto{\pgfqpoint{5.108714in}{2.303367in}}%
\pgfpathlineto{\pgfqpoint{5.109327in}{2.302634in}}%
\pgfpathlineto{\pgfqpoint{5.109490in}{2.302804in}}%
\pgfpathlineto{\pgfqpoint{5.110723in}{2.302102in}}%
\pgfpathlineto{\pgfqpoint{5.109492in}{2.302810in}}%
\pgfpathlineto{\pgfqpoint{5.110780in}{2.302215in}}%
\pgfpathlineto{\pgfqpoint{5.111344in}{2.302449in}}%
\pgfpathlineto{\pgfqpoint{5.110946in}{2.302011in}}%
\pgfpathlineto{\pgfqpoint{5.111909in}{2.302377in}}%
\pgfpathlineto{\pgfqpoint{5.112065in}{2.302183in}}%
\pgfpathlineto{\pgfqpoint{5.112306in}{2.302751in}}%
\pgfpathlineto{\pgfqpoint{5.113002in}{2.302712in}}%
\pgfpathlineto{\pgfqpoint{5.114086in}{2.302757in}}%
\pgfpathlineto{\pgfqpoint{5.113007in}{2.302701in}}%
\pgfpathlineto{\pgfqpoint{5.114214in}{2.302493in}}%
\pgfpathlineto{\pgfqpoint{5.115112in}{2.301657in}}%
\pgfpathlineto{\pgfqpoint{5.114295in}{2.302589in}}%
\pgfpathlineto{\pgfqpoint{5.115384in}{2.301871in}}%
\pgfpathlineto{\pgfqpoint{5.116358in}{2.302553in}}%
\pgfpathlineto{\pgfqpoint{5.115623in}{2.301734in}}%
\pgfpathlineto{\pgfqpoint{5.116527in}{2.302175in}}%
\pgfpathlineto{\pgfqpoint{5.117262in}{2.301365in}}%
\pgfpathlineto{\pgfqpoint{5.116546in}{2.302234in}}%
\pgfpathlineto{\pgfqpoint{5.117688in}{2.301536in}}%
\pgfpathlineto{\pgfqpoint{5.118718in}{2.301536in}}%
\pgfpathlineto{\pgfqpoint{5.118239in}{2.301080in}}%
\pgfpathlineto{\pgfqpoint{5.118794in}{2.301410in}}%
\pgfpathlineto{\pgfqpoint{5.119021in}{2.301283in}}%
\pgfpathlineto{\pgfqpoint{5.119550in}{2.301727in}}%
\pgfpathlineto{\pgfqpoint{5.119893in}{2.301658in}}%
\pgfpathlineto{\pgfqpoint{5.123124in}{2.300543in}}%
\pgfpathlineto{\pgfqpoint{5.119896in}{2.301669in}}%
\pgfpathlineto{\pgfqpoint{5.123350in}{2.301018in}}%
\pgfpathlineto{\pgfqpoint{5.123940in}{2.301481in}}%
\pgfpathlineto{\pgfqpoint{5.123436in}{2.300900in}}%
\pgfpathlineto{\pgfqpoint{5.124491in}{2.301298in}}%
\pgfpathlineto{\pgfqpoint{5.124722in}{2.300966in}}%
\pgfpathlineto{\pgfqpoint{5.125432in}{2.301602in}}%
\pgfpathlineto{\pgfqpoint{5.125601in}{2.301444in}}%
\pgfpathlineto{\pgfqpoint{5.126716in}{2.302103in}}%
\pgfpathlineto{\pgfqpoint{5.125605in}{2.301438in}}%
\pgfpathlineto{\pgfqpoint{5.126744in}{2.301968in}}%
\pgfpathlineto{\pgfqpoint{5.127549in}{2.301600in}}%
\pgfpathlineto{\pgfqpoint{5.126759in}{2.302048in}}%
\pgfpathlineto{\pgfqpoint{5.127917in}{2.301961in}}%
\pgfpathlineto{\pgfqpoint{5.129106in}{2.302808in}}%
\pgfpathlineto{\pgfqpoint{5.129291in}{2.302705in}}%
\pgfpathlineto{\pgfqpoint{5.129841in}{2.302286in}}%
\pgfpathlineto{\pgfqpoint{5.129513in}{2.302775in}}%
\pgfpathlineto{\pgfqpoint{5.130417in}{2.302393in}}%
\pgfpathlineto{\pgfqpoint{5.131736in}{2.302408in}}%
\pgfpathlineto{\pgfqpoint{5.130462in}{2.302230in}}%
\pgfpathlineto{\pgfqpoint{5.131769in}{2.302236in}}%
\pgfpathlineto{\pgfqpoint{5.131971in}{2.301943in}}%
\pgfpathlineto{\pgfqpoint{5.131805in}{2.302244in}}%
\pgfpathlineto{\pgfqpoint{5.132877in}{2.302272in}}%
\pgfpathlineto{\pgfqpoint{5.133590in}{2.302519in}}%
\pgfpathlineto{\pgfqpoint{5.133020in}{2.302122in}}%
\pgfpathlineto{\pgfqpoint{5.133992in}{2.302400in}}%
\pgfpathlineto{\pgfqpoint{5.134291in}{2.302505in}}%
\pgfpathlineto{\pgfqpoint{5.134632in}{2.302156in}}%
\pgfpathlineto{\pgfqpoint{5.135078in}{2.302164in}}%
\pgfpathlineto{\pgfqpoint{5.136092in}{2.302038in}}%
\pgfpathlineto{\pgfqpoint{5.135662in}{2.302537in}}%
\pgfpathlineto{\pgfqpoint{5.136189in}{2.302177in}}%
\pgfpathlineto{\pgfqpoint{5.137445in}{2.302818in}}%
\pgfpathlineto{\pgfqpoint{5.136224in}{2.302114in}}%
\pgfpathlineto{\pgfqpoint{5.137454in}{2.302786in}}%
\pgfpathlineto{\pgfqpoint{5.138113in}{2.301982in}}%
\pgfpathlineto{\pgfqpoint{5.138594in}{2.302430in}}%
\pgfpathlineto{\pgfqpoint{5.139508in}{2.302710in}}%
\pgfpathlineto{\pgfqpoint{5.138625in}{2.302426in}}%
\pgfpathlineto{\pgfqpoint{5.139719in}{2.302529in}}%
\pgfpathlineto{\pgfqpoint{5.141540in}{2.301973in}}%
\pgfpathlineto{\pgfqpoint{5.141590in}{2.302075in}}%
\pgfpathlineto{\pgfqpoint{5.142112in}{2.302338in}}%
\pgfpathlineto{\pgfqpoint{5.141746in}{2.301896in}}%
\pgfpathlineto{\pgfqpoint{5.142700in}{2.302014in}}%
\pgfpathlineto{\pgfqpoint{5.144123in}{2.302290in}}%
\pgfpathlineto{\pgfqpoint{5.142720in}{2.301949in}}%
\pgfpathlineto{\pgfqpoint{5.144153in}{2.302208in}}%
\pgfpathlineto{\pgfqpoint{5.144687in}{2.302147in}}%
\pgfpathlineto{\pgfqpoint{5.144786in}{2.302577in}}%
\pgfpathlineto{\pgfqpoint{5.145254in}{2.302424in}}%
\pgfpathlineto{\pgfqpoint{5.145344in}{2.302516in}}%
\pgfpathlineto{\pgfqpoint{5.146314in}{2.301791in}}%
\pgfpathlineto{\pgfqpoint{5.146320in}{2.301824in}}%
\pgfpathlineto{\pgfqpoint{5.146346in}{2.301754in}}%
\pgfpathlineto{\pgfqpoint{5.146587in}{2.302180in}}%
\pgfpathlineto{\pgfqpoint{5.147420in}{2.302075in}}%
\pgfpathlineto{\pgfqpoint{5.148480in}{2.302255in}}%
\pgfpathlineto{\pgfqpoint{5.147513in}{2.301813in}}%
\pgfpathlineto{\pgfqpoint{5.148533in}{2.302044in}}%
\pgfpathlineto{\pgfqpoint{5.149424in}{2.301780in}}%
\pgfpathlineto{\pgfqpoint{5.148662in}{2.302224in}}%
\pgfpathlineto{\pgfqpoint{5.149669in}{2.301849in}}%
\pgfpathlineto{\pgfqpoint{5.150018in}{2.302105in}}%
\pgfpathlineto{\pgfqpoint{5.150558in}{2.301377in}}%
\pgfpathlineto{\pgfqpoint{5.150781in}{2.301795in}}%
\pgfpathlineto{\pgfqpoint{5.152065in}{2.300955in}}%
\pgfpathlineto{\pgfqpoint{5.150787in}{2.301812in}}%
\pgfpathlineto{\pgfqpoint{5.152096in}{2.301032in}}%
\pgfpathlineto{\pgfqpoint{5.153237in}{2.301200in}}%
\pgfpathlineto{\pgfqpoint{5.152669in}{2.300794in}}%
\pgfpathlineto{\pgfqpoint{5.153238in}{2.301194in}}%
\pgfpathlineto{\pgfqpoint{5.154400in}{2.300808in}}%
\pgfpathlineto{\pgfqpoint{5.154405in}{2.300813in}}%
\pgfpathlineto{\pgfqpoint{5.154759in}{2.301310in}}%
\pgfpathlineto{\pgfqpoint{5.155477in}{2.300648in}}%
\pgfpathlineto{\pgfqpoint{5.155512in}{2.300674in}}%
\pgfpathlineto{\pgfqpoint{5.156516in}{2.300072in}}%
\pgfpathlineto{\pgfqpoint{5.155526in}{2.300702in}}%
\pgfpathlineto{\pgfqpoint{5.156649in}{2.300192in}}%
\pgfpathlineto{\pgfqpoint{5.156977in}{2.300396in}}%
\pgfpathlineto{\pgfqpoint{5.157447in}{2.299993in}}%
\pgfpathlineto{\pgfqpoint{5.157762in}{2.300203in}}%
\pgfpathlineto{\pgfqpoint{5.158454in}{2.299684in}}%
\pgfpathlineto{\pgfqpoint{5.157814in}{2.300314in}}%
\pgfpathlineto{\pgfqpoint{5.158905in}{2.299882in}}%
\pgfpathlineto{\pgfqpoint{5.159033in}{2.299526in}}%
\pgfpathlineto{\pgfqpoint{5.160025in}{2.300109in}}%
\pgfpathlineto{\pgfqpoint{5.161478in}{2.299835in}}%
\pgfpathlineto{\pgfqpoint{5.160061in}{2.300164in}}%
\pgfpathlineto{\pgfqpoint{5.161488in}{2.299908in}}%
\pgfpathlineto{\pgfqpoint{5.162493in}{2.301110in}}%
\pgfpathlineto{\pgfqpoint{5.161702in}{2.299861in}}%
\pgfpathlineto{\pgfqpoint{5.162633in}{2.300983in}}%
\pgfpathlineto{\pgfqpoint{5.162924in}{2.300488in}}%
\pgfpathlineto{\pgfqpoint{5.162651in}{2.301003in}}%
\pgfpathlineto{\pgfqpoint{5.163761in}{2.300903in}}%
\pgfpathlineto{\pgfqpoint{5.165043in}{2.303438in}}%
\pgfpathlineto{\pgfqpoint{5.163802in}{2.300850in}}%
\pgfpathlineto{\pgfqpoint{5.165044in}{2.303433in}}%
\pgfpathlineto{\pgfqpoint{5.166034in}{2.303368in}}%
\pgfpathlineto{\pgfqpoint{5.165059in}{2.303449in}}%
\pgfpathlineto{\pgfqpoint{5.166155in}{2.303433in}}%
\pgfpathlineto{\pgfqpoint{5.166613in}{2.303317in}}%
\pgfpathlineto{\pgfqpoint{5.166246in}{2.303448in}}%
\pgfpathlineto{\pgfqpoint{5.167268in}{2.303412in}}%
\pgfpathlineto{\pgfqpoint{5.168367in}{2.303442in}}%
\pgfpathlineto{\pgfqpoint{5.167638in}{2.303360in}}%
\pgfpathlineto{\pgfqpoint{5.168382in}{2.303428in}}%
\pgfpathlineto{\pgfqpoint{5.169443in}{2.303353in}}%
\pgfpathlineto{\pgfqpoint{5.168620in}{2.303455in}}%
\pgfpathlineto{\pgfqpoint{5.169495in}{2.303387in}}%
\pgfpathlineto{\pgfqpoint{5.170775in}{2.303444in}}%
\pgfpathlineto{\pgfqpoint{5.170794in}{2.303400in}}%
\pgfpathlineto{\pgfqpoint{5.171584in}{2.303336in}}%
\pgfpathlineto{\pgfqpoint{5.170847in}{2.303448in}}%
\pgfpathlineto{\pgfqpoint{5.171908in}{2.303374in}}%
\pgfpathlineto{\pgfqpoint{5.173034in}{2.303452in}}%
\pgfpathlineto{\pgfqpoint{5.173046in}{2.303411in}}%
\pgfpathlineto{\pgfqpoint{5.174088in}{2.303350in}}%
\pgfpathlineto{\pgfqpoint{5.173057in}{2.303432in}}%
\pgfpathlineto{\pgfqpoint{5.174183in}{2.303396in}}%
\pgfpathlineto{\pgfqpoint{5.175129in}{2.303453in}}%
\pgfpathlineto{\pgfqpoint{5.174239in}{2.303372in}}%
\pgfpathlineto{\pgfqpoint{5.175295in}{2.303391in}}%
\pgfpathlineto{\pgfqpoint{5.176435in}{2.303369in}}%
\pgfpathlineto{\pgfqpoint{5.175310in}{2.303438in}}%
\pgfpathlineto{\pgfqpoint{5.176447in}{2.303422in}}%
\pgfpathlineto{\pgfqpoint{5.177548in}{2.303439in}}%
\pgfpathlineto{\pgfqpoint{5.176524in}{2.303376in}}%
\pgfpathlineto{\pgfqpoint{5.177562in}{2.303406in}}%
\pgfpathlineto{\pgfqpoint{5.178657in}{2.303350in}}%
\pgfpathlineto{\pgfqpoint{5.177575in}{2.303440in}}%
\pgfpathlineto{\pgfqpoint{5.178692in}{2.303414in}}%
\pgfpathlineto{\pgfqpoint{5.179790in}{2.303441in}}%
\pgfpathlineto{\pgfqpoint{5.178732in}{2.303384in}}%
\pgfpathlineto{\pgfqpoint{5.179813in}{2.303433in}}%
\pgfpathlineto{\pgfqpoint{5.180707in}{2.303372in}}%
\pgfpathlineto{\pgfqpoint{5.179989in}{2.303446in}}%
\pgfpathlineto{\pgfqpoint{5.180925in}{2.303419in}}%
\pgfpathlineto{\pgfqpoint{5.182024in}{2.303443in}}%
\pgfpathlineto{\pgfqpoint{5.180987in}{2.303370in}}%
\pgfpathlineto{\pgfqpoint{5.182037in}{2.303401in}}%
\pgfpathlineto{\pgfqpoint{5.183104in}{2.303268in}}%
\pgfpathlineto{\pgfqpoint{5.182718in}{2.303584in}}%
\pgfpathlineto{\pgfqpoint{5.183149in}{2.303389in}}%
\pgfpathlineto{\pgfqpoint{5.184271in}{2.303787in}}%
\pgfpathlineto{\pgfqpoint{5.183206in}{2.303370in}}%
\pgfpathlineto{\pgfqpoint{5.184301in}{2.303702in}}%
\pgfpathlineto{\pgfqpoint{5.185359in}{2.303149in}}%
\pgfpathlineto{\pgfqpoint{5.185073in}{2.303848in}}%
\pgfpathlineto{\pgfqpoint{5.185441in}{2.303299in}}%
\pgfpathlineto{\pgfqpoint{5.186696in}{2.304369in}}%
\pgfpathlineto{\pgfqpoint{5.185501in}{2.303150in}}%
\pgfpathlineto{\pgfqpoint{5.186705in}{2.304340in}}%
\pgfpathlineto{\pgfqpoint{5.187014in}{2.303924in}}%
\pgfpathlineto{\pgfqpoint{5.187270in}{2.304446in}}%
\pgfpathlineto{\pgfqpoint{5.187828in}{2.304086in}}%
\pgfpathlineto{\pgfqpoint{5.188765in}{2.304792in}}%
\pgfpathlineto{\pgfqpoint{5.187967in}{2.304043in}}%
\pgfpathlineto{\pgfqpoint{5.188965in}{2.304691in}}%
\pgfpathlineto{\pgfqpoint{5.189909in}{2.305008in}}%
\pgfpathlineto{\pgfqpoint{5.189016in}{2.304573in}}%
\pgfpathlineto{\pgfqpoint{5.190128in}{2.304789in}}%
\pgfpathlineto{\pgfqpoint{5.190130in}{2.304768in}}%
\pgfpathlineto{\pgfqpoint{5.191113in}{2.305319in}}%
\pgfpathlineto{\pgfqpoint{5.191212in}{2.305213in}}%
\pgfpathlineto{\pgfqpoint{5.192400in}{2.305883in}}%
\pgfpathlineto{\pgfqpoint{5.191424in}{2.304874in}}%
\pgfpathlineto{\pgfqpoint{5.192403in}{2.305864in}}%
\pgfpathlineto{\pgfqpoint{5.195526in}{2.305935in}}%
\pgfpathlineto{\pgfqpoint{5.195550in}{2.305833in}}%
\pgfpathlineto{\pgfqpoint{5.196161in}{2.305534in}}%
\pgfpathlineto{\pgfqpoint{5.195766in}{2.306084in}}%
\pgfpathlineto{\pgfqpoint{5.196658in}{2.305899in}}%
\pgfpathlineto{\pgfqpoint{5.196987in}{2.306244in}}%
\pgfpathlineto{\pgfqpoint{5.197481in}{2.305680in}}%
\pgfpathlineto{\pgfqpoint{5.197766in}{2.305814in}}%
\pgfpathlineto{\pgfqpoint{5.198636in}{2.305478in}}%
\pgfpathlineto{\pgfqpoint{5.198062in}{2.305959in}}%
\pgfpathlineto{\pgfqpoint{5.198883in}{2.305669in}}%
\pgfpathlineto{\pgfqpoint{5.198884in}{2.305684in}}%
\pgfpathlineto{\pgfqpoint{5.199854in}{2.304999in}}%
\pgfpathlineto{\pgfqpoint{5.199972in}{2.305227in}}%
\pgfpathlineto{\pgfqpoint{5.200615in}{2.304965in}}%
\pgfpathlineto{\pgfqpoint{5.200103in}{2.305364in}}%
\pgfpathlineto{\pgfqpoint{5.201074in}{2.305360in}}%
\pgfpathlineto{\pgfqpoint{5.202200in}{2.305560in}}%
\pgfpathlineto{\pgfqpoint{5.201191in}{2.305261in}}%
\pgfpathlineto{\pgfqpoint{5.202321in}{2.305316in}}%
\pgfpathlineto{\pgfqpoint{5.203336in}{2.304867in}}%
\pgfpathlineto{\pgfqpoint{5.202587in}{2.305574in}}%
\pgfpathlineto{\pgfqpoint{5.203456in}{2.304967in}}%
\pgfpathlineto{\pgfqpoint{5.203459in}{2.305003in}}%
\pgfpathlineto{\pgfqpoint{5.204357in}{2.304297in}}%
\pgfpathlineto{\pgfqpoint{5.204521in}{2.304460in}}%
\pgfpathlineto{\pgfqpoint{5.205261in}{2.304383in}}%
\pgfpathlineto{\pgfqpoint{5.204571in}{2.304609in}}%
\pgfpathlineto{\pgfqpoint{5.205630in}{2.304722in}}%
\pgfpathlineto{\pgfqpoint{5.205676in}{2.304817in}}%
\pgfpathlineto{\pgfqpoint{5.206505in}{2.303824in}}%
\pgfpathlineto{\pgfqpoint{5.206726in}{2.304003in}}%
\pgfpathlineto{\pgfqpoint{5.207756in}{2.304355in}}%
\pgfpathlineto{\pgfqpoint{5.206809in}{2.303904in}}%
\pgfpathlineto{\pgfqpoint{5.207880in}{2.304050in}}%
\pgfpathlineto{\pgfqpoint{5.208887in}{2.303617in}}%
\pgfpathlineto{\pgfqpoint{5.208251in}{2.304232in}}%
\pgfpathlineto{\pgfqpoint{5.209026in}{2.303682in}}%
\pgfpathlineto{\pgfqpoint{5.210432in}{2.303161in}}%
\pgfpathlineto{\pgfqpoint{5.209191in}{2.303834in}}%
\pgfpathlineto{\pgfqpoint{5.210474in}{2.303334in}}%
\pgfpathlineto{\pgfqpoint{5.210566in}{2.303527in}}%
\pgfpathlineto{\pgfqpoint{5.211543in}{2.302505in}}%
\pgfpathlineto{\pgfqpoint{5.212148in}{2.302977in}}%
\pgfpathlineto{\pgfqpoint{5.211742in}{2.302279in}}%
\pgfpathlineto{\pgfqpoint{5.212685in}{2.302732in}}%
\pgfpathlineto{\pgfqpoint{5.213845in}{2.303307in}}%
\pgfpathlineto{\pgfqpoint{5.212782in}{2.302671in}}%
\pgfpathlineto{\pgfqpoint{5.213918in}{2.303261in}}%
\pgfpathlineto{\pgfqpoint{5.214550in}{2.303430in}}%
\pgfpathlineto{\pgfqpoint{5.213975in}{2.303429in}}%
\pgfpathlineto{\pgfqpoint{5.214829in}{2.303889in}}%
\pgfpathlineto{\pgfqpoint{5.215765in}{2.304109in}}%
\pgfpathlineto{\pgfqpoint{5.215021in}{2.303687in}}%
\pgfpathlineto{\pgfqpoint{5.215946in}{2.303886in}}%
\pgfpathlineto{\pgfqpoint{5.216727in}{2.303945in}}%
\pgfpathlineto{\pgfqpoint{5.216075in}{2.304199in}}%
\pgfpathlineto{\pgfqpoint{5.217023in}{2.304147in}}%
\pgfpathlineto{\pgfqpoint{5.217334in}{2.303901in}}%
\pgfpathlineto{\pgfqpoint{5.218139in}{2.304194in}}%
\pgfpathlineto{\pgfqpoint{5.219167in}{2.304123in}}%
\pgfpathlineto{\pgfqpoint{5.218516in}{2.304434in}}%
\pgfpathlineto{\pgfqpoint{5.219247in}{2.304274in}}%
\pgfpathlineto{\pgfqpoint{5.220241in}{2.305055in}}%
\pgfpathlineto{\pgfqpoint{5.220460in}{2.304839in}}%
\pgfpathlineto{\pgfqpoint{5.220586in}{2.304700in}}%
\pgfpathlineto{\pgfqpoint{5.221384in}{2.305117in}}%
\pgfpathlineto{\pgfqpoint{5.221568in}{2.304937in}}%
\pgfpathlineto{\pgfqpoint{5.221799in}{2.305105in}}%
\pgfpathlineto{\pgfqpoint{5.222188in}{2.304675in}}%
\pgfpathlineto{\pgfqpoint{5.222683in}{2.304946in}}%
\pgfpathlineto{\pgfqpoint{5.223735in}{2.303976in}}%
\pgfpathlineto{\pgfqpoint{5.222691in}{2.304971in}}%
\pgfpathlineto{\pgfqpoint{5.223832in}{2.304060in}}%
\pgfpathlineto{\pgfqpoint{5.224761in}{2.304289in}}%
\pgfpathlineto{\pgfqpoint{5.223945in}{2.303729in}}%
\pgfpathlineto{\pgfqpoint{5.224953in}{2.303928in}}%
\pgfpathlineto{\pgfqpoint{5.225126in}{2.304181in}}%
\pgfpathlineto{\pgfqpoint{5.226092in}{2.303420in}}%
\pgfpathlineto{\pgfqpoint{5.226372in}{2.303619in}}%
\pgfpathlineto{\pgfqpoint{5.226734in}{2.303007in}}%
\pgfpathlineto{\pgfqpoint{5.227195in}{2.303143in}}%
\pgfpathlineto{\pgfqpoint{5.228169in}{2.303200in}}%
\pgfpathlineto{\pgfqpoint{5.227474in}{2.302944in}}%
\pgfpathlineto{\pgfqpoint{5.228298in}{2.303020in}}%
\pgfpathlineto{\pgfqpoint{5.228750in}{2.302785in}}%
\pgfpathlineto{\pgfqpoint{5.229372in}{2.303216in}}%
\pgfpathlineto{\pgfqpoint{5.229405in}{2.303154in}}%
\pgfpathlineto{\pgfqpoint{5.229644in}{2.303276in}}%
\pgfpathlineto{\pgfqpoint{5.230275in}{2.302766in}}%
\pgfpathlineto{\pgfqpoint{5.230513in}{2.303095in}}%
\pgfpathlineto{\pgfqpoint{5.231629in}{2.302889in}}%
\pgfpathlineto{\pgfqpoint{5.231056in}{2.303403in}}%
\pgfpathlineto{\pgfqpoint{5.231637in}{2.302948in}}%
\pgfpathlineto{\pgfqpoint{5.232763in}{2.302592in}}%
\pgfpathlineto{\pgfqpoint{5.231694in}{2.303001in}}%
\pgfpathlineto{\pgfqpoint{5.232816in}{2.302708in}}%
\pgfpathlineto{\pgfqpoint{5.232845in}{2.302796in}}%
\pgfpathlineto{\pgfqpoint{5.233452in}{2.302159in}}%
\pgfpathlineto{\pgfqpoint{5.233903in}{2.302178in}}%
\pgfpathlineto{\pgfqpoint{5.234005in}{2.302392in}}%
\pgfpathlineto{\pgfqpoint{5.234908in}{2.301604in}}%
\pgfpathlineto{\pgfqpoint{5.234982in}{2.301705in}}%
\pgfpathlineto{\pgfqpoint{5.236083in}{2.301371in}}%
\pgfpathlineto{\pgfqpoint{5.235448in}{2.302219in}}%
\pgfpathlineto{\pgfqpoint{5.236107in}{2.301476in}}%
\pgfpathlineto{\pgfqpoint{5.236737in}{2.301621in}}%
\pgfpathlineto{\pgfqpoint{5.237173in}{2.301186in}}%
\pgfpathlineto{\pgfqpoint{5.237213in}{2.301365in}}%
\pgfpathlineto{\pgfqpoint{5.237377in}{2.301303in}}%
\pgfpathlineto{\pgfqpoint{5.238000in}{2.301834in}}%
\pgfpathlineto{\pgfqpoint{5.238286in}{2.301785in}}%
\pgfpathlineto{\pgfqpoint{5.239458in}{2.303456in}}%
\pgfpathlineto{\pgfqpoint{5.238659in}{2.301635in}}%
\pgfpathlineto{\pgfqpoint{5.239468in}{2.303438in}}%
\pgfpathlineto{\pgfqpoint{5.240422in}{2.303341in}}%
\pgfpathlineto{\pgfqpoint{5.239663in}{2.303449in}}%
\pgfpathlineto{\pgfqpoint{5.240580in}{2.303418in}}%
\pgfpathlineto{\pgfqpoint{5.241693in}{2.303456in}}%
\pgfpathlineto{\pgfqpoint{5.240645in}{2.303372in}}%
\pgfpathlineto{\pgfqpoint{5.241693in}{2.303439in}}%
\pgfpathlineto{\pgfqpoint{5.242463in}{2.303361in}}%
\pgfpathlineto{\pgfqpoint{5.241744in}{2.303446in}}%
\pgfpathlineto{\pgfqpoint{5.242804in}{2.303420in}}%
\pgfpathlineto{\pgfqpoint{5.243730in}{2.303449in}}%
\pgfpathlineto{\pgfqpoint{5.243007in}{2.303350in}}%
\pgfpathlineto{\pgfqpoint{5.243917in}{2.303431in}}%
\pgfpathlineto{\pgfqpoint{5.244586in}{2.303341in}}%
\pgfpathlineto{\pgfqpoint{5.244192in}{2.303457in}}%
\pgfpathlineto{\pgfqpoint{5.245028in}{2.303416in}}%
\pgfpathlineto{\pgfqpoint{5.246098in}{2.303445in}}%
\pgfpathlineto{\pgfqpoint{5.245311in}{2.303319in}}%
\pgfpathlineto{\pgfqpoint{5.246143in}{2.303417in}}%
\pgfpathlineto{\pgfqpoint{5.247104in}{2.303377in}}%
\pgfpathlineto{\pgfqpoint{5.246188in}{2.303438in}}%
\pgfpathlineto{\pgfqpoint{5.247257in}{2.303415in}}%
\pgfpathlineto{\pgfqpoint{5.248205in}{2.303450in}}%
\pgfpathlineto{\pgfqpoint{5.247346in}{2.303368in}}%
\pgfpathlineto{\pgfqpoint{5.248369in}{2.303395in}}%
\pgfpathlineto{\pgfqpoint{5.249469in}{2.303436in}}%
\pgfpathlineto{\pgfqpoint{5.248455in}{2.303351in}}%
\pgfpathlineto{\pgfqpoint{5.249491in}{2.303399in}}%
\pgfpathlineto{\pgfqpoint{5.250956in}{2.303443in}}%
\pgfpathlineto{\pgfqpoint{5.250965in}{2.303429in}}%
\pgfpathlineto{\pgfqpoint{5.251959in}{2.303352in}}%
\pgfpathlineto{\pgfqpoint{5.251010in}{2.303449in}}%
\pgfpathlineto{\pgfqpoint{5.252077in}{2.303411in}}%
\pgfpathlineto{\pgfqpoint{5.253209in}{2.303379in}}%
\pgfpathlineto{\pgfqpoint{5.252089in}{2.303448in}}%
\pgfpathlineto{\pgfqpoint{5.253213in}{2.303411in}}%
\pgfpathlineto{\pgfqpoint{5.254251in}{2.303450in}}%
\pgfpathlineto{\pgfqpoint{5.253418in}{2.303351in}}%
\pgfpathlineto{\pgfqpoint{5.254324in}{2.303389in}}%
\pgfpathlineto{\pgfqpoint{5.255369in}{2.303370in}}%
\pgfpathlineto{\pgfqpoint{5.254389in}{2.303446in}}%
\pgfpathlineto{\pgfqpoint{5.255433in}{2.303423in}}%
\pgfpathlineto{\pgfqpoint{5.256349in}{2.303450in}}%
\pgfpathlineto{\pgfqpoint{5.255710in}{2.303360in}}%
\pgfpathlineto{\pgfqpoint{5.256543in}{2.303397in}}%
\pgfpathlineto{\pgfqpoint{5.257673in}{2.303360in}}%
\pgfpathlineto{\pgfqpoint{5.256547in}{2.303432in}}%
\pgfpathlineto{\pgfqpoint{5.257674in}{2.303374in}}%
\pgfpathlineto{\pgfqpoint{5.258795in}{2.303909in}}%
\pgfpathlineto{\pgfqpoint{5.257857in}{2.303326in}}%
\pgfpathlineto{\pgfqpoint{5.258810in}{2.303896in}}%
\pgfpathlineto{\pgfqpoint{5.260076in}{2.303675in}}%
\pgfpathlineto{\pgfqpoint{5.258868in}{2.303994in}}%
\pgfpathlineto{\pgfqpoint{5.260119in}{2.303789in}}%
\pgfpathlineto{\pgfqpoint{5.261231in}{2.304652in}}%
\pgfpathlineto{\pgfqpoint{5.260476in}{2.303572in}}%
\pgfpathlineto{\pgfqpoint{5.261245in}{2.304616in}}%
\pgfpathlineto{\pgfqpoint{5.261314in}{2.304553in}}%
\pgfpathlineto{\pgfqpoint{5.262072in}{2.305185in}}%
\pgfpathlineto{\pgfqpoint{5.262346in}{2.304911in}}%
\pgfpathlineto{\pgfqpoint{5.263234in}{2.305273in}}%
\pgfpathlineto{\pgfqpoint{5.262364in}{2.304856in}}%
\pgfpathlineto{\pgfqpoint{5.263498in}{2.305112in}}%
\pgfpathlineto{\pgfqpoint{5.264260in}{2.304883in}}%
\pgfpathlineto{\pgfqpoint{5.264553in}{2.305404in}}%
\pgfpathlineto{\pgfqpoint{5.264587in}{2.305331in}}%
\pgfpathlineto{\pgfqpoint{5.264912in}{2.305849in}}%
\pgfpathlineto{\pgfqpoint{5.264589in}{2.305329in}}%
\pgfpathlineto{\pgfqpoint{5.265711in}{2.305195in}}%
\pgfpathlineto{\pgfqpoint{5.266658in}{2.304769in}}%
\pgfpathlineto{\pgfqpoint{5.265906in}{2.305328in}}%
\pgfpathlineto{\pgfqpoint{5.266849in}{2.305068in}}%
\pgfpathlineto{\pgfqpoint{5.267180in}{2.305457in}}%
\pgfpathlineto{\pgfqpoint{5.267650in}{2.304933in}}%
\pgfpathlineto{\pgfqpoint{5.267970in}{2.305254in}}%
\pgfpathlineto{\pgfqpoint{5.268345in}{2.305122in}}%
\pgfpathlineto{\pgfqpoint{5.268787in}{2.305486in}}%
\pgfpathlineto{\pgfqpoint{5.269080in}{2.305286in}}%
\pgfpathlineto{\pgfqpoint{5.270053in}{2.305750in}}%
\pgfpathlineto{\pgfqpoint{5.269084in}{2.305268in}}%
\pgfpathlineto{\pgfqpoint{5.270227in}{2.305728in}}%
\pgfpathlineto{\pgfqpoint{5.271266in}{2.304860in}}%
\pgfpathlineto{\pgfqpoint{5.270230in}{2.305736in}}%
\pgfpathlineto{\pgfqpoint{5.271529in}{2.305063in}}%
\pgfpathlineto{\pgfqpoint{5.272472in}{2.305694in}}%
\pgfpathlineto{\pgfqpoint{5.271573in}{2.305214in}}%
\pgfpathlineto{\pgfqpoint{5.272476in}{2.305739in}}%
\pgfpathlineto{\pgfqpoint{5.273302in}{2.305864in}}%
\pgfpathlineto{\pgfqpoint{5.272905in}{2.305310in}}%
\pgfpathlineto{\pgfqpoint{5.273587in}{2.305679in}}%
\pgfpathlineto{\pgfqpoint{5.274695in}{2.305706in}}%
\pgfpathlineto{\pgfqpoint{5.273644in}{2.305946in}}%
\pgfpathlineto{\pgfqpoint{5.274770in}{2.305817in}}%
\pgfpathlineto{\pgfqpoint{5.275873in}{2.306478in}}%
\pgfpathlineto{\pgfqpoint{5.274780in}{2.305801in}}%
\pgfpathlineto{\pgfqpoint{5.275955in}{2.306435in}}%
\pgfpathlineto{\pgfqpoint{5.276279in}{2.306616in}}%
\pgfpathlineto{\pgfqpoint{5.276852in}{2.305899in}}%
\pgfpathlineto{\pgfqpoint{5.276971in}{2.305931in}}%
\pgfpathlineto{\pgfqpoint{5.277680in}{2.305624in}}%
\pgfpathlineto{\pgfqpoint{5.277275in}{2.306062in}}%
\pgfpathlineto{\pgfqpoint{5.278096in}{2.305658in}}%
\pgfpathlineto{\pgfqpoint{5.278328in}{2.305807in}}%
\pgfpathlineto{\pgfqpoint{5.279008in}{2.305160in}}%
\pgfpathlineto{\pgfqpoint{5.279170in}{2.305317in}}%
\pgfpathlineto{\pgfqpoint{5.279679in}{2.304950in}}%
\pgfpathlineto{\pgfqpoint{5.280289in}{2.305069in}}%
\pgfpathlineto{\pgfqpoint{5.281393in}{2.304842in}}%
\pgfpathlineto{\pgfqpoint{5.280482in}{2.305240in}}%
\pgfpathlineto{\pgfqpoint{5.281416in}{2.304865in}}%
\pgfpathlineto{\pgfqpoint{5.282566in}{2.305169in}}%
\pgfpathlineto{\pgfqpoint{5.282066in}{2.304603in}}%
\pgfpathlineto{\pgfqpoint{5.282589in}{2.305105in}}%
\pgfpathlineto{\pgfqpoint{5.282596in}{2.305002in}}%
\pgfpathlineto{\pgfqpoint{5.283166in}{2.305466in}}%
\pgfpathlineto{\pgfqpoint{5.283695in}{2.305226in}}%
\pgfpathlineto{\pgfqpoint{5.284058in}{2.304910in}}%
\pgfpathlineto{\pgfqpoint{5.283763in}{2.305365in}}%
\pgfpathlineto{\pgfqpoint{5.284811in}{2.305232in}}%
\pgfpathlineto{\pgfqpoint{5.284852in}{2.305353in}}%
\pgfpathlineto{\pgfqpoint{5.285721in}{2.304513in}}%
\pgfpathlineto{\pgfqpoint{5.285852in}{2.304559in}}%
\pgfpathlineto{\pgfqpoint{5.286784in}{2.304305in}}%
\pgfpathlineto{\pgfqpoint{5.286132in}{2.304730in}}%
\pgfpathlineto{\pgfqpoint{5.286974in}{2.304475in}}%
\pgfpathlineto{\pgfqpoint{5.287941in}{2.304972in}}%
\pgfpathlineto{\pgfqpoint{5.286989in}{2.304402in}}%
\pgfpathlineto{\pgfqpoint{5.288102in}{2.304931in}}%
\pgfpathlineto{\pgfqpoint{5.288918in}{2.305179in}}%
\pgfpathlineto{\pgfqpoint{5.288175in}{2.304806in}}%
\pgfpathlineto{\pgfqpoint{5.289215in}{2.304829in}}%
\pgfpathlineto{\pgfqpoint{5.290043in}{2.304777in}}%
\pgfpathlineto{\pgfqpoint{5.289845in}{2.305293in}}%
\pgfpathlineto{\pgfqpoint{5.290318in}{2.304985in}}%
\pgfpathlineto{\pgfqpoint{5.291106in}{2.305504in}}%
\pgfpathlineto{\pgfqpoint{5.290528in}{2.304662in}}%
\pgfpathlineto{\pgfqpoint{5.291448in}{2.305310in}}%
\pgfpathlineto{\pgfqpoint{5.291779in}{2.305512in}}%
\pgfpathlineto{\pgfqpoint{5.292142in}{2.304733in}}%
\pgfpathlineto{\pgfqpoint{5.292553in}{2.305197in}}%
\pgfpathlineto{\pgfqpoint{5.293111in}{2.304435in}}%
\pgfpathlineto{\pgfqpoint{5.293685in}{2.305036in}}%
\pgfpathlineto{\pgfqpoint{5.294775in}{2.305149in}}%
\pgfpathlineto{\pgfqpoint{5.293945in}{2.304655in}}%
\pgfpathlineto{\pgfqpoint{5.294804in}{2.304990in}}%
\pgfpathlineto{\pgfqpoint{5.295601in}{2.304619in}}%
\pgfpathlineto{\pgfqpoint{5.294820in}{2.305082in}}%
\pgfpathlineto{\pgfqpoint{5.295953in}{2.304758in}}%
\pgfpathlineto{\pgfqpoint{5.296907in}{2.305094in}}%
\pgfpathlineto{\pgfqpoint{5.296028in}{2.304690in}}%
\pgfpathlineto{\pgfqpoint{5.297125in}{2.305033in}}%
\pgfpathlineto{\pgfqpoint{5.298086in}{2.304703in}}%
\pgfpathlineto{\pgfqpoint{5.297230in}{2.305247in}}%
\pgfpathlineto{\pgfqpoint{5.298242in}{2.304851in}}%
\pgfpathlineto{\pgfqpoint{5.299456in}{2.304888in}}%
\pgfpathlineto{\pgfqpoint{5.298365in}{2.304690in}}%
\pgfpathlineto{\pgfqpoint{5.299511in}{2.304825in}}%
\pgfpathlineto{\pgfqpoint{5.299641in}{2.304773in}}%
\pgfpathlineto{\pgfqpoint{5.300317in}{2.305398in}}%
\pgfpathlineto{\pgfqpoint{5.300598in}{2.305098in}}%
\pgfpathlineto{\pgfqpoint{5.300643in}{2.305266in}}%
\pgfpathlineto{\pgfqpoint{5.301686in}{2.304595in}}%
\pgfpathlineto{\pgfqpoint{5.301695in}{2.304647in}}%
\pgfpathlineto{\pgfqpoint{5.302824in}{2.304638in}}%
\pgfpathlineto{\pgfqpoint{5.301927in}{2.305002in}}%
\pgfpathlineto{\pgfqpoint{5.302840in}{2.304681in}}%
\pgfpathlineto{\pgfqpoint{5.303798in}{2.304751in}}%
\pgfpathlineto{\pgfqpoint{5.303120in}{2.304394in}}%
\pgfpathlineto{\pgfqpoint{5.303944in}{2.304564in}}%
\pgfpathlineto{\pgfqpoint{5.305006in}{2.304431in}}%
\pgfpathlineto{\pgfqpoint{5.304707in}{2.304772in}}%
\pgfpathlineto{\pgfqpoint{5.305052in}{2.304648in}}%
\pgfpathlineto{\pgfqpoint{5.305971in}{2.305204in}}%
\pgfpathlineto{\pgfqpoint{5.305159in}{2.304487in}}%
\pgfpathlineto{\pgfqpoint{5.306194in}{2.304964in}}%
\pgfpathlineto{\pgfqpoint{5.306743in}{2.304732in}}%
\pgfpathlineto{\pgfqpoint{5.306972in}{2.305251in}}%
\pgfpathlineto{\pgfqpoint{5.307308in}{2.304928in}}%
\pgfpathlineto{\pgfqpoint{5.308344in}{2.305372in}}%
\pgfpathlineto{\pgfqpoint{5.307751in}{2.304783in}}%
\pgfpathlineto{\pgfqpoint{5.308440in}{2.305248in}}%
\pgfpathlineto{\pgfqpoint{5.308859in}{2.305002in}}%
\pgfpathlineto{\pgfqpoint{5.309305in}{2.305644in}}%
\pgfpathlineto{\pgfqpoint{5.309528in}{2.305465in}}%
\pgfpathlineto{\pgfqpoint{5.310207in}{2.305672in}}%
\pgfpathlineto{\pgfqpoint{5.310590in}{2.305132in}}%
\pgfpathlineto{\pgfqpoint{5.310608in}{2.305181in}}%
\pgfpathlineto{\pgfqpoint{5.311151in}{2.304915in}}%
\pgfpathlineto{\pgfqpoint{5.310618in}{2.305228in}}%
\pgfpathlineto{\pgfqpoint{5.311731in}{2.305027in}}%
\pgfpathlineto{\pgfqpoint{5.312800in}{2.305292in}}%
\pgfpathlineto{\pgfqpoint{5.312117in}{2.304833in}}%
\pgfpathlineto{\pgfqpoint{5.312852in}{2.305166in}}%
\pgfpathlineto{\pgfqpoint{5.314114in}{2.303360in}}%
\pgfpathlineto{\pgfqpoint{5.313330in}{2.305658in}}%
\pgfpathlineto{\pgfqpoint{5.314120in}{2.303415in}}%
\pgfpathlineto{\pgfqpoint{5.315216in}{2.303440in}}%
\pgfpathlineto{\pgfqpoint{5.314138in}{2.303379in}}%
\pgfpathlineto{\pgfqpoint{5.315242in}{2.303397in}}%
\pgfpathlineto{\pgfqpoint{5.316153in}{2.303343in}}%
\pgfpathlineto{\pgfqpoint{5.315406in}{2.303456in}}%
\pgfpathlineto{\pgfqpoint{5.316354in}{2.303422in}}%
\pgfpathlineto{\pgfqpoint{5.317492in}{2.303449in}}%
\pgfpathlineto{\pgfqpoint{5.316388in}{2.303338in}}%
\pgfpathlineto{\pgfqpoint{5.317497in}{2.303430in}}%
\pgfpathlineto{\pgfqpoint{5.318602in}{2.303362in}}%
\pgfpathlineto{\pgfqpoint{5.317531in}{2.303441in}}%
\pgfpathlineto{\pgfqpoint{5.318609in}{2.303406in}}%
\pgfpathlineto{\pgfqpoint{5.319655in}{2.303374in}}%
\pgfpathlineto{\pgfqpoint{5.318614in}{2.303439in}}%
\pgfpathlineto{\pgfqpoint{5.319733in}{2.303413in}}%
\pgfpathlineto{\pgfqpoint{5.320778in}{2.303446in}}%
\pgfpathlineto{\pgfqpoint{5.319912in}{2.303354in}}%
\pgfpathlineto{\pgfqpoint{5.320844in}{2.303392in}}%
\pgfpathlineto{\pgfqpoint{5.321911in}{2.303381in}}%
\pgfpathlineto{\pgfqpoint{5.320873in}{2.303435in}}%
\pgfpathlineto{\pgfqpoint{5.321958in}{2.303414in}}%
\pgfpathlineto{\pgfqpoint{5.323005in}{2.303449in}}%
\pgfpathlineto{\pgfqpoint{5.322177in}{2.303350in}}%
\pgfpathlineto{\pgfqpoint{5.323069in}{2.303406in}}%
\pgfpathlineto{\pgfqpoint{5.323576in}{2.303330in}}%
\pgfpathlineto{\pgfqpoint{5.323224in}{2.303446in}}%
\pgfpathlineto{\pgfqpoint{5.324181in}{2.303398in}}%
\pgfpathlineto{\pgfqpoint{5.325140in}{2.303450in}}%
\pgfpathlineto{\pgfqpoint{5.324586in}{2.303348in}}%
\pgfpathlineto{\pgfqpoint{5.325292in}{2.303424in}}%
\pgfpathlineto{\pgfqpoint{5.326269in}{2.303373in}}%
\pgfpathlineto{\pgfqpoint{5.325525in}{2.303460in}}%
\pgfpathlineto{\pgfqpoint{5.326405in}{2.303422in}}%
\pgfpathlineto{\pgfqpoint{5.327491in}{2.303442in}}%
\pgfpathlineto{\pgfqpoint{5.326676in}{2.303362in}}%
\pgfpathlineto{\pgfqpoint{5.327517in}{2.303430in}}%
\pgfpathlineto{\pgfqpoint{5.328628in}{2.303383in}}%
\pgfpathlineto{\pgfqpoint{5.327545in}{2.303454in}}%
\pgfpathlineto{\pgfqpoint{5.328633in}{2.303414in}}%
\pgfpathlineto{\pgfqpoint{5.329757in}{2.303441in}}%
\pgfpathlineto{\pgfqpoint{5.328778in}{2.303363in}}%
\pgfpathlineto{\pgfqpoint{5.329758in}{2.303437in}}%
\pgfpathlineto{\pgfqpoint{5.330492in}{2.303350in}}%
\pgfpathlineto{\pgfqpoint{5.329789in}{2.303448in}}%
\pgfpathlineto{\pgfqpoint{5.330870in}{2.303393in}}%
\pgfpathlineto{\pgfqpoint{5.331948in}{2.303376in}}%
\pgfpathlineto{\pgfqpoint{5.330916in}{2.303443in}}%
\pgfpathlineto{\pgfqpoint{5.331982in}{2.303423in}}%
\pgfpathlineto{\pgfqpoint{5.332907in}{2.303667in}}%
\pgfpathlineto{\pgfqpoint{5.332567in}{2.303265in}}%
\pgfpathlineto{\pgfqpoint{5.333095in}{2.303398in}}%
\pgfpathlineto{\pgfqpoint{5.333275in}{2.303652in}}%
\pgfpathlineto{\pgfqpoint{5.334030in}{2.303011in}}%
\pgfpathlineto{\pgfqpoint{5.334171in}{2.303074in}}%
\pgfpathlineto{\pgfqpoint{5.334656in}{2.302919in}}%
\pgfpathlineto{\pgfqpoint{5.334239in}{2.303232in}}%
\pgfpathlineto{\pgfqpoint{5.335257in}{2.303358in}}%
\pgfpathlineto{\pgfqpoint{5.336279in}{2.303817in}}%
\pgfpathlineto{\pgfqpoint{5.335650in}{2.303207in}}%
\pgfpathlineto{\pgfqpoint{5.336381in}{2.303658in}}%
\pgfpathlineto{\pgfqpoint{5.337500in}{2.304297in}}%
\pgfpathlineto{\pgfqpoint{5.336396in}{2.303652in}}%
\pgfpathlineto{\pgfqpoint{5.337607in}{2.304253in}}%
\pgfpathlineto{\pgfqpoint{5.338681in}{2.304128in}}%
\pgfpathlineto{\pgfqpoint{5.338347in}{2.304569in}}%
\pgfpathlineto{\pgfqpoint{5.338737in}{2.304216in}}%
\pgfpathlineto{\pgfqpoint{5.339866in}{2.304808in}}%
\pgfpathlineto{\pgfqpoint{5.339869in}{2.304794in}}%
\pgfpathlineto{\pgfqpoint{5.339947in}{2.304635in}}%
\pgfpathlineto{\pgfqpoint{5.340369in}{2.305422in}}%
\pgfpathlineto{\pgfqpoint{5.340887in}{2.305293in}}%
\pgfpathlineto{\pgfqpoint{5.341083in}{2.305482in}}%
\pgfpathlineto{\pgfqpoint{5.341710in}{2.304658in}}%
\pgfpathlineto{\pgfqpoint{5.341985in}{2.304941in}}%
\pgfpathlineto{\pgfqpoint{5.343065in}{2.304495in}}%
\pgfpathlineto{\pgfqpoint{5.342242in}{2.305267in}}%
\pgfpathlineto{\pgfqpoint{5.343107in}{2.304668in}}%
\pgfpathlineto{\pgfqpoint{5.344160in}{2.305328in}}%
\pgfpathlineto{\pgfqpoint{5.343114in}{2.304649in}}%
\pgfpathlineto{\pgfqpoint{5.344241in}{2.305245in}}%
\pgfpathlineto{\pgfqpoint{5.344780in}{2.304701in}}%
\pgfpathlineto{\pgfqpoint{5.344277in}{2.305352in}}%
\pgfpathlineto{\pgfqpoint{5.345362in}{2.305132in}}%
\pgfpathlineto{\pgfqpoint{5.346314in}{2.304895in}}%
\pgfpathlineto{\pgfqpoint{5.345599in}{2.305187in}}%
\pgfpathlineto{\pgfqpoint{5.346512in}{2.305005in}}%
\pgfpathlineto{\pgfqpoint{5.347454in}{2.305733in}}%
\pgfpathlineto{\pgfqpoint{5.346519in}{2.304996in}}%
\pgfpathlineto{\pgfqpoint{5.347661in}{2.305459in}}%
\pgfpathlineto{\pgfqpoint{5.347706in}{2.305544in}}%
\pgfpathlineto{\pgfqpoint{5.348210in}{2.304573in}}%
\pgfpathlineto{\pgfqpoint{5.348467in}{2.304665in}}%
\pgfpathlineto{\pgfqpoint{5.349272in}{2.304427in}}%
\pgfpathlineto{\pgfqpoint{5.348539in}{2.304846in}}%
\pgfpathlineto{\pgfqpoint{5.349593in}{2.304637in}}%
\pgfpathlineto{\pgfqpoint{5.350687in}{2.305175in}}%
\pgfpathlineto{\pgfqpoint{5.349603in}{2.304615in}}%
\pgfpathlineto{\pgfqpoint{5.350725in}{2.305096in}}%
\pgfpathlineto{\pgfqpoint{5.351766in}{2.305332in}}%
\pgfpathlineto{\pgfqpoint{5.350771in}{2.304959in}}%
\pgfpathlineto{\pgfqpoint{5.351849in}{2.305050in}}%
\pgfpathlineto{\pgfqpoint{5.352411in}{2.304577in}}%
\pgfpathlineto{\pgfqpoint{5.351985in}{2.305086in}}%
\pgfpathlineto{\pgfqpoint{5.352972in}{2.304900in}}%
\pgfpathlineto{\pgfqpoint{5.353521in}{2.305107in}}%
\pgfpathlineto{\pgfqpoint{5.354042in}{2.304572in}}%
\pgfpathlineto{\pgfqpoint{5.354056in}{2.304608in}}%
\pgfpathlineto{\pgfqpoint{5.354964in}{2.304798in}}%
\pgfpathlineto{\pgfqpoint{5.354540in}{2.304467in}}%
\pgfpathlineto{\pgfqpoint{5.355191in}{2.304623in}}%
\pgfpathlineto{\pgfqpoint{5.356732in}{2.303724in}}%
\pgfpathlineto{\pgfqpoint{5.355330in}{2.304783in}}%
\pgfpathlineto{\pgfqpoint{5.356809in}{2.303814in}}%
\pgfpathlineto{\pgfqpoint{5.357893in}{2.304148in}}%
\pgfpathlineto{\pgfqpoint{5.357950in}{2.304067in}}%
\pgfpathlineto{\pgfqpoint{5.358761in}{2.303850in}}%
\pgfpathlineto{\pgfqpoint{5.358167in}{2.304214in}}%
\pgfpathlineto{\pgfqpoint{5.359065in}{2.303939in}}%
\pgfpathlineto{\pgfqpoint{5.360015in}{2.303880in}}%
\pgfpathlineto{\pgfqpoint{5.359761in}{2.303530in}}%
\pgfpathlineto{\pgfqpoint{5.360134in}{2.303638in}}%
\pgfpathlineto{\pgfqpoint{5.360258in}{2.303417in}}%
\pgfpathlineto{\pgfqpoint{5.360731in}{2.304159in}}%
\pgfpathlineto{\pgfqpoint{5.361228in}{2.303838in}}%
\pgfpathlineto{\pgfqpoint{5.361617in}{2.303895in}}%
\pgfpathlineto{\pgfqpoint{5.361942in}{2.303385in}}%
\pgfpathlineto{\pgfqpoint{5.362298in}{2.303414in}}%
\pgfpathlineto{\pgfqpoint{5.363537in}{2.303643in}}%
\pgfpathlineto{\pgfqpoint{5.362307in}{2.303444in}}%
\pgfpathlineto{\pgfqpoint{5.363561in}{2.303719in}}%
\pgfpathlineto{\pgfqpoint{5.364087in}{2.303863in}}%
\pgfpathlineto{\pgfqpoint{5.364610in}{2.303352in}}%
\pgfpathlineto{\pgfqpoint{5.364653in}{2.303363in}}%
\pgfpathlineto{\pgfqpoint{5.365373in}{2.303106in}}%
\pgfpathlineto{\pgfqpoint{5.364720in}{2.303508in}}%
\pgfpathlineto{\pgfqpoint{5.365772in}{2.303295in}}%
\pgfpathlineto{\pgfqpoint{5.365781in}{2.303332in}}%
\pgfpathlineto{\pgfqpoint{5.366631in}{2.302519in}}%
\pgfpathlineto{\pgfqpoint{5.366815in}{2.302821in}}%
\pgfpathlineto{\pgfqpoint{5.367815in}{2.302354in}}%
\pgfpathlineto{\pgfqpoint{5.367127in}{2.303130in}}%
\pgfpathlineto{\pgfqpoint{5.367958in}{2.302414in}}%
\pgfpathlineto{\pgfqpoint{5.368825in}{2.302916in}}%
\pgfpathlineto{\pgfqpoint{5.368134in}{2.302360in}}%
\pgfpathlineto{\pgfqpoint{5.369097in}{2.302823in}}%
\pgfpathlineto{\pgfqpoint{5.370663in}{2.302608in}}%
\pgfpathlineto{\pgfqpoint{5.369125in}{2.302992in}}%
\pgfpathlineto{\pgfqpoint{5.370734in}{2.302821in}}%
\pgfpathlineto{\pgfqpoint{5.371241in}{2.303000in}}%
\pgfpathlineto{\pgfqpoint{5.371730in}{2.302426in}}%
\pgfpathlineto{\pgfqpoint{5.371836in}{2.302655in}}%
\pgfpathlineto{\pgfqpoint{5.372715in}{2.302219in}}%
\pgfpathlineto{\pgfqpoint{5.371854in}{2.302704in}}%
\pgfpathlineto{\pgfqpoint{5.372955in}{2.302464in}}%
\pgfpathlineto{\pgfqpoint{5.373861in}{2.302608in}}%
\pgfpathlineto{\pgfqpoint{5.373324in}{2.302152in}}%
\pgfpathlineto{\pgfqpoint{5.374064in}{2.302376in}}%
\pgfpathlineto{\pgfqpoint{5.375233in}{2.301338in}}%
\pgfpathlineto{\pgfqpoint{5.374068in}{2.302389in}}%
\pgfpathlineto{\pgfqpoint{5.375236in}{2.301380in}}%
\pgfpathlineto{\pgfqpoint{5.375988in}{2.301542in}}%
\pgfpathlineto{\pgfqpoint{5.375595in}{2.301074in}}%
\pgfpathlineto{\pgfqpoint{5.376335in}{2.301157in}}%
\pgfpathlineto{\pgfqpoint{5.376970in}{2.301215in}}%
\pgfpathlineto{\pgfqpoint{5.376625in}{2.300852in}}%
\pgfpathlineto{\pgfqpoint{5.377447in}{2.301105in}}%
\pgfpathlineto{\pgfqpoint{5.377473in}{2.301078in}}%
\pgfpathlineto{\pgfqpoint{5.378298in}{2.301621in}}%
\pgfpathlineto{\pgfqpoint{5.378549in}{2.301368in}}%
\pgfpathlineto{\pgfqpoint{5.379633in}{2.301509in}}%
\pgfpathlineto{\pgfqpoint{5.378712in}{2.301134in}}%
\pgfpathlineto{\pgfqpoint{5.379850in}{2.301225in}}%
\pgfpathlineto{\pgfqpoint{5.379852in}{2.301193in}}%
\pgfpathlineto{\pgfqpoint{5.380232in}{2.301767in}}%
\pgfpathlineto{\pgfqpoint{5.380946in}{2.301572in}}%
\pgfpathlineto{\pgfqpoint{5.380953in}{2.301590in}}%
\pgfpathlineto{\pgfqpoint{5.381427in}{2.301183in}}%
\pgfpathlineto{\pgfqpoint{5.382022in}{2.301225in}}%
\pgfpathlineto{\pgfqpoint{5.382294in}{2.301214in}}%
\pgfpathlineto{\pgfqpoint{5.383022in}{2.301965in}}%
\pgfpathlineto{\pgfqpoint{5.383106in}{2.301803in}}%
\pgfpathlineto{\pgfqpoint{5.383520in}{2.302046in}}%
\pgfpathlineto{\pgfqpoint{5.383738in}{2.301639in}}%
\pgfpathlineto{\pgfqpoint{5.384219in}{2.301811in}}%
\pgfpathlineto{\pgfqpoint{5.385361in}{2.300912in}}%
\pgfpathlineto{\pgfqpoint{5.384256in}{2.301854in}}%
\pgfpathlineto{\pgfqpoint{5.385363in}{2.300929in}}%
\pgfpathlineto{\pgfqpoint{5.385875in}{2.301191in}}%
\pgfpathlineto{\pgfqpoint{5.385442in}{2.300915in}}%
\pgfpathlineto{\pgfqpoint{5.386474in}{2.300942in}}%
\pgfpathlineto{\pgfqpoint{5.387094in}{2.300549in}}%
\pgfpathlineto{\pgfqpoint{5.386609in}{2.301035in}}%
\pgfpathlineto{\pgfqpoint{5.387585in}{2.300952in}}%
\pgfpathlineto{\pgfqpoint{5.387729in}{2.301023in}}%
\pgfpathlineto{\pgfqpoint{5.387957in}{2.300548in}}%
\pgfpathlineto{\pgfqpoint{5.388697in}{2.300949in}}%
\pgfpathlineto{\pgfqpoint{5.388702in}{2.300940in}}%
\pgfpathlineto{\pgfqpoint{5.388966in}{2.302879in}}%
\pgfpathlineto{\pgfqpoint{5.390117in}{2.303449in}}%
\pgfpathlineto{\pgfqpoint{5.390795in}{2.303357in}}%
\pgfpathlineto{\pgfqpoint{5.390225in}{2.303454in}}%
\pgfpathlineto{\pgfqpoint{5.391230in}{2.303380in}}%
\pgfpathlineto{\pgfqpoint{5.392284in}{2.303371in}}%
\pgfpathlineto{\pgfqpoint{5.391252in}{2.303437in}}%
\pgfpathlineto{\pgfqpoint{5.392340in}{2.303414in}}%
\pgfpathlineto{\pgfqpoint{5.393250in}{2.303455in}}%
\pgfpathlineto{\pgfqpoint{5.392912in}{2.303339in}}%
\pgfpathlineto{\pgfqpoint{5.393451in}{2.303406in}}%
\pgfpathlineto{\pgfqpoint{5.394412in}{2.303374in}}%
\pgfpathlineto{\pgfqpoint{5.393628in}{2.303448in}}%
\pgfpathlineto{\pgfqpoint{5.394564in}{2.303412in}}%
\pgfpathlineto{\pgfqpoint{5.395577in}{2.303457in}}%
\pgfpathlineto{\pgfqpoint{5.394584in}{2.303355in}}%
\pgfpathlineto{\pgfqpoint{5.395677in}{2.303429in}}%
\pgfpathlineto{\pgfqpoint{5.396349in}{2.303331in}}%
\pgfpathlineto{\pgfqpoint{5.395761in}{2.303444in}}%
\pgfpathlineto{\pgfqpoint{5.396789in}{2.303387in}}%
\pgfpathlineto{\pgfqpoint{5.397660in}{2.303452in}}%
\pgfpathlineto{\pgfqpoint{5.397700in}{2.303320in}}%
\pgfpathlineto{\pgfqpoint{5.397901in}{2.303413in}}%
\pgfpathlineto{\pgfqpoint{5.398993in}{2.303350in}}%
\pgfpathlineto{\pgfqpoint{5.397970in}{2.303446in}}%
\pgfpathlineto{\pgfqpoint{5.399016in}{2.303422in}}%
\pgfpathlineto{\pgfqpoint{5.400122in}{2.303451in}}%
\pgfpathlineto{\pgfqpoint{5.399381in}{2.303362in}}%
\pgfpathlineto{\pgfqpoint{5.400127in}{2.303409in}}%
\pgfpathlineto{\pgfqpoint{5.401225in}{2.303377in}}%
\pgfpathlineto{\pgfqpoint{5.400165in}{2.303447in}}%
\pgfpathlineto{\pgfqpoint{5.401255in}{2.303413in}}%
\pgfpathlineto{\pgfqpoint{5.402334in}{2.303453in}}%
\pgfpathlineto{\pgfqpoint{5.401393in}{2.303343in}}%
\pgfpathlineto{\pgfqpoint{5.402372in}{2.303404in}}%
\pgfpathlineto{\pgfqpoint{5.402825in}{2.303330in}}%
\pgfpathlineto{\pgfqpoint{5.402392in}{2.303445in}}%
\pgfpathlineto{\pgfqpoint{5.403484in}{2.303424in}}%
\pgfpathlineto{\pgfqpoint{5.404593in}{2.303440in}}%
\pgfpathlineto{\pgfqpoint{5.403525in}{2.303363in}}%
\pgfpathlineto{\pgfqpoint{5.404622in}{2.303409in}}%
\pgfpathlineto{\pgfqpoint{5.405710in}{2.303333in}}%
\pgfpathlineto{\pgfqpoint{5.404671in}{2.303445in}}%
\pgfpathlineto{\pgfqpoint{5.405737in}{2.303400in}}%
\pgfpathlineto{\pgfqpoint{5.406814in}{2.303442in}}%
\pgfpathlineto{\pgfqpoint{5.405900in}{2.303343in}}%
\pgfpathlineto{\pgfqpoint{5.406855in}{2.303435in}}%
\pgfpathlineto{\pgfqpoint{5.407966in}{2.303278in}}%
\pgfpathlineto{\pgfqpoint{5.407479in}{2.303549in}}%
\pgfpathlineto{\pgfqpoint{5.407970in}{2.303292in}}%
\pgfpathlineto{\pgfqpoint{5.409152in}{2.303658in}}%
\pgfpathlineto{\pgfqpoint{5.408028in}{2.303226in}}%
\pgfpathlineto{\pgfqpoint{5.409153in}{2.303653in}}%
\pgfpathlineto{\pgfqpoint{5.409174in}{2.303632in}}%
\pgfpathlineto{\pgfqpoint{5.409697in}{2.304342in}}%
\pgfpathlineto{\pgfqpoint{5.410225in}{2.303936in}}%
\pgfpathlineto{\pgfqpoint{5.411255in}{2.304713in}}%
\pgfpathlineto{\pgfqpoint{5.410252in}{2.303844in}}%
\pgfpathlineto{\pgfqpoint{5.411370in}{2.304646in}}%
\pgfpathlineto{\pgfqpoint{5.412279in}{2.304251in}}%
\pgfpathlineto{\pgfqpoint{5.411379in}{2.304662in}}%
\pgfpathlineto{\pgfqpoint{5.412519in}{2.304361in}}%
\pgfpathlineto{\pgfqpoint{5.413741in}{2.304569in}}%
\pgfpathlineto{\pgfqpoint{5.412827in}{2.304008in}}%
\pgfpathlineto{\pgfqpoint{5.413752in}{2.304541in}}%
\pgfpathlineto{\pgfqpoint{5.414865in}{2.303975in}}%
\pgfpathlineto{\pgfqpoint{5.414131in}{2.304775in}}%
\pgfpathlineto{\pgfqpoint{5.414890in}{2.304070in}}%
\pgfpathlineto{\pgfqpoint{5.415670in}{2.304540in}}%
\pgfpathlineto{\pgfqpoint{5.416020in}{2.304243in}}%
\pgfpathlineto{\pgfqpoint{5.416020in}{2.304239in}}%
\pgfpathlineto{\pgfqpoint{5.416904in}{2.305223in}}%
\pgfpathlineto{\pgfqpoint{5.416985in}{2.305036in}}%
\pgfpathlineto{\pgfqpoint{5.416996in}{2.305118in}}%
\pgfpathlineto{\pgfqpoint{5.417920in}{2.303918in}}%
\pgfpathlineto{\pgfqpoint{5.418049in}{2.304022in}}%
\pgfpathlineto{\pgfqpoint{5.419211in}{2.303618in}}%
\pgfpathlineto{\pgfqpoint{5.418111in}{2.304150in}}%
\pgfpathlineto{\pgfqpoint{5.419263in}{2.303692in}}%
\pgfpathlineto{\pgfqpoint{5.419349in}{2.303829in}}%
\pgfpathlineto{\pgfqpoint{5.419836in}{2.303426in}}%
\pgfpathlineto{\pgfqpoint{5.420371in}{2.303614in}}%
\pgfpathlineto{\pgfqpoint{5.421578in}{2.304345in}}%
\pgfpathlineto{\pgfqpoint{5.420420in}{2.303438in}}%
\pgfpathlineto{\pgfqpoint{5.421642in}{2.304206in}}%
\pgfpathlineto{\pgfqpoint{5.422832in}{2.303775in}}%
\pgfpathlineto{\pgfqpoint{5.421749in}{2.304278in}}%
\pgfpathlineto{\pgfqpoint{5.422835in}{2.303811in}}%
\pgfpathlineto{\pgfqpoint{5.423704in}{2.304345in}}%
\pgfpathlineto{\pgfqpoint{5.422902in}{2.303675in}}%
\pgfpathlineto{\pgfqpoint{5.424021in}{2.304160in}}%
\pgfpathlineto{\pgfqpoint{5.424097in}{2.304223in}}%
\pgfpathlineto{\pgfqpoint{5.425175in}{2.303441in}}%
\pgfpathlineto{\pgfqpoint{5.425338in}{2.303651in}}%
\pgfpathlineto{\pgfqpoint{5.425960in}{2.303115in}}%
\pgfpathlineto{\pgfqpoint{5.426269in}{2.303179in}}%
\pgfpathlineto{\pgfqpoint{5.427500in}{2.303162in}}%
\pgfpathlineto{\pgfqpoint{5.427535in}{2.303277in}}%
\pgfpathlineto{\pgfqpoint{5.428349in}{2.303457in}}%
\pgfpathlineto{\pgfqpoint{5.428472in}{2.303085in}}%
\pgfpathlineto{\pgfqpoint{5.428645in}{2.303231in}}%
\pgfpathlineto{\pgfqpoint{5.429033in}{2.303545in}}%
\pgfpathlineto{\pgfqpoint{5.428679in}{2.303093in}}%
\pgfpathlineto{\pgfqpoint{5.429770in}{2.303285in}}%
\pgfpathlineto{\pgfqpoint{5.430898in}{2.302701in}}%
\pgfpathlineto{\pgfqpoint{5.429967in}{2.303453in}}%
\pgfpathlineto{\pgfqpoint{5.430905in}{2.302739in}}%
\pgfpathlineto{\pgfqpoint{5.431387in}{2.302965in}}%
\pgfpathlineto{\pgfqpoint{5.431826in}{2.302378in}}%
\pgfpathlineto{\pgfqpoint{5.432001in}{2.302524in}}%
\pgfpathlineto{\pgfqpoint{5.432098in}{2.302325in}}%
\pgfpathlineto{\pgfqpoint{5.432841in}{2.303055in}}%
\pgfpathlineto{\pgfqpoint{5.433006in}{2.303067in}}%
\pgfpathlineto{\pgfqpoint{5.434111in}{2.303990in}}%
\pgfpathlineto{\pgfqpoint{5.433038in}{2.303046in}}%
\pgfpathlineto{\pgfqpoint{5.434188in}{2.303881in}}%
\pgfpathlineto{\pgfqpoint{5.434509in}{2.303665in}}%
\pgfpathlineto{\pgfqpoint{5.435234in}{2.304420in}}%
\pgfpathlineto{\pgfqpoint{5.436243in}{2.304531in}}%
\pgfpathlineto{\pgfqpoint{5.435356in}{2.304108in}}%
\pgfpathlineto{\pgfqpoint{5.436353in}{2.304395in}}%
\pgfpathlineto{\pgfqpoint{5.437595in}{2.303443in}}%
\pgfpathlineto{\pgfqpoint{5.436357in}{2.304425in}}%
\pgfpathlineto{\pgfqpoint{5.437628in}{2.303615in}}%
\pgfpathlineto{\pgfqpoint{5.438684in}{2.303784in}}%
\pgfpathlineto{\pgfqpoint{5.438295in}{2.303360in}}%
\pgfpathlineto{\pgfqpoint{5.438746in}{2.303722in}}%
\pgfpathlineto{\pgfqpoint{5.439029in}{2.303284in}}%
\pgfpathlineto{\pgfqpoint{5.439693in}{2.303979in}}%
\pgfpathlineto{\pgfqpoint{5.439857in}{2.303740in}}%
\pgfpathlineto{\pgfqpoint{5.440968in}{2.304348in}}%
\pgfpathlineto{\pgfqpoint{5.439985in}{2.303435in}}%
\pgfpathlineto{\pgfqpoint{5.441000in}{2.304218in}}%
\pgfpathlineto{\pgfqpoint{5.441984in}{2.304377in}}%
\pgfpathlineto{\pgfqpoint{5.442184in}{2.304117in}}%
\pgfpathlineto{\pgfqpoint{5.443236in}{2.303718in}}%
\pgfpathlineto{\pgfqpoint{5.442187in}{2.304123in}}%
\pgfpathlineto{\pgfqpoint{5.443318in}{2.303866in}}%
\pgfpathlineto{\pgfqpoint{5.444717in}{2.304793in}}%
\pgfpathlineto{\pgfqpoint{5.444821in}{2.304660in}}%
\pgfpathlineto{\pgfqpoint{5.444876in}{2.304491in}}%
\pgfpathlineto{\pgfqpoint{5.445843in}{2.305060in}}%
\pgfpathlineto{\pgfqpoint{5.445914in}{2.304990in}}%
\pgfpathlineto{\pgfqpoint{5.446723in}{2.305132in}}%
\pgfpathlineto{\pgfqpoint{5.446284in}{2.304754in}}%
\pgfpathlineto{\pgfqpoint{5.447022in}{2.304889in}}%
\pgfpathlineto{\pgfqpoint{5.447767in}{2.304612in}}%
\pgfpathlineto{\pgfqpoint{5.447265in}{2.304940in}}%
\pgfpathlineto{\pgfqpoint{5.448137in}{2.304868in}}%
\pgfpathlineto{\pgfqpoint{5.448925in}{2.305957in}}%
\pgfpathlineto{\pgfqpoint{5.449286in}{2.305868in}}%
\pgfpathlineto{\pgfqpoint{5.449307in}{2.305752in}}%
\pgfpathlineto{\pgfqpoint{5.449999in}{2.306846in}}%
\pgfpathlineto{\pgfqpoint{5.450289in}{2.306695in}}%
\pgfpathlineto{\pgfqpoint{5.450729in}{2.306959in}}%
\pgfpathlineto{\pgfqpoint{5.451126in}{2.306427in}}%
\pgfpathlineto{\pgfqpoint{5.451392in}{2.306508in}}%
\pgfpathlineto{\pgfqpoint{5.451393in}{2.306493in}}%
\pgfpathlineto{\pgfqpoint{5.451841in}{2.307074in}}%
\pgfpathlineto{\pgfqpoint{5.452485in}{2.306878in}}%
\pgfpathlineto{\pgfqpoint{5.453136in}{2.307163in}}%
\pgfpathlineto{\pgfqpoint{5.453482in}{2.306785in}}%
\pgfpathlineto{\pgfqpoint{5.453603in}{2.306952in}}%
\pgfpathlineto{\pgfqpoint{5.454616in}{2.306621in}}%
\pgfpathlineto{\pgfqpoint{5.453736in}{2.307122in}}%
\pgfpathlineto{\pgfqpoint{5.454723in}{2.306827in}}%
\pgfpathlineto{\pgfqpoint{5.455867in}{2.307442in}}%
\pgfpathlineto{\pgfqpoint{5.454776in}{2.306768in}}%
\pgfpathlineto{\pgfqpoint{5.455962in}{2.307270in}}%
\pgfpathlineto{\pgfqpoint{5.456397in}{2.307012in}}%
\pgfpathlineto{\pgfqpoint{5.456161in}{2.307343in}}%
\pgfpathlineto{\pgfqpoint{5.457095in}{2.307239in}}%
\pgfpathlineto{\pgfqpoint{5.457131in}{2.307379in}}%
\pgfpathlineto{\pgfqpoint{5.457879in}{2.306565in}}%
\pgfpathlineto{\pgfqpoint{5.458138in}{2.306780in}}%
\pgfpathlineto{\pgfqpoint{5.459278in}{2.306056in}}%
\pgfpathlineto{\pgfqpoint{5.459285in}{2.306063in}}%
\pgfpathlineto{\pgfqpoint{5.459560in}{2.305763in}}%
\pgfpathlineto{\pgfqpoint{5.459334in}{2.306194in}}%
\pgfpathlineto{\pgfqpoint{5.460408in}{2.306017in}}%
\pgfpathlineto{\pgfqpoint{5.461063in}{2.306055in}}%
\pgfpathlineto{\pgfqpoint{5.461315in}{2.305360in}}%
\pgfpathlineto{\pgfqpoint{5.461492in}{2.305561in}}%
\pgfpathlineto{\pgfqpoint{5.461784in}{2.304960in}}%
\pgfpathlineto{\pgfqpoint{5.461562in}{2.305609in}}%
\pgfpathlineto{\pgfqpoint{5.462612in}{2.305348in}}%
\pgfpathlineto{\pgfqpoint{5.463722in}{2.303360in}}%
\pgfpathlineto{\pgfqpoint{5.462968in}{2.305534in}}%
\pgfpathlineto{\pgfqpoint{5.463935in}{2.303396in}}%
\pgfpathlineto{\pgfqpoint{5.465452in}{2.303400in}}%
\pgfpathlineto{\pgfqpoint{5.463942in}{2.303441in}}%
\pgfpathlineto{\pgfqpoint{5.465455in}{2.303420in}}%
\pgfpathlineto{\pgfqpoint{5.466472in}{2.303448in}}%
\pgfpathlineto{\pgfqpoint{5.466134in}{2.303352in}}%
\pgfpathlineto{\pgfqpoint{5.466566in}{2.303409in}}%
\pgfpathlineto{\pgfqpoint{5.467616in}{2.303444in}}%
\pgfpathlineto{\pgfqpoint{5.466598in}{2.303358in}}%
\pgfpathlineto{\pgfqpoint{5.467679in}{2.303406in}}%
\pgfpathlineto{\pgfqpoint{5.468760in}{2.303365in}}%
\pgfpathlineto{\pgfqpoint{5.467730in}{2.303441in}}%
\pgfpathlineto{\pgfqpoint{5.468792in}{2.303393in}}%
\pgfpathlineto{\pgfqpoint{5.469808in}{2.303447in}}%
\pgfpathlineto{\pgfqpoint{5.468961in}{2.303349in}}%
\pgfpathlineto{\pgfqpoint{5.469906in}{2.303409in}}%
\pgfpathlineto{\pgfqpoint{5.471032in}{2.303387in}}%
\pgfpathlineto{\pgfqpoint{5.469918in}{2.303438in}}%
\pgfpathlineto{\pgfqpoint{5.471037in}{2.303403in}}%
\pgfpathlineto{\pgfqpoint{5.471978in}{2.303450in}}%
\pgfpathlineto{\pgfqpoint{5.471196in}{2.303303in}}%
\pgfpathlineto{\pgfqpoint{5.472149in}{2.303432in}}%
\pgfpathlineto{\pgfqpoint{5.473149in}{2.303358in}}%
\pgfpathlineto{\pgfqpoint{5.472279in}{2.303450in}}%
\pgfpathlineto{\pgfqpoint{5.473260in}{2.303409in}}%
\pgfpathlineto{\pgfqpoint{5.474368in}{2.303444in}}%
\pgfpathlineto{\pgfqpoint{5.473602in}{2.303356in}}%
\pgfpathlineto{\pgfqpoint{5.474373in}{2.303412in}}%
\pgfpathlineto{\pgfqpoint{5.477841in}{2.303409in}}%
\pgfpathlineto{\pgfqpoint{5.477842in}{2.303420in}}%
\pgfpathlineto{\pgfqpoint{5.478820in}{2.303451in}}%
\pgfpathlineto{\pgfqpoint{5.478047in}{2.303352in}}%
\pgfpathlineto{\pgfqpoint{5.478953in}{2.303424in}}%
\pgfpathlineto{\pgfqpoint{5.479803in}{2.303354in}}%
\pgfpathlineto{\pgfqpoint{5.479009in}{2.303444in}}%
\pgfpathlineto{\pgfqpoint{5.480064in}{2.303414in}}%
\pgfpathlineto{\pgfqpoint{5.480423in}{2.303333in}}%
\pgfpathlineto{\pgfqpoint{5.481177in}{2.303445in}}%
\pgfpathlineto{\pgfqpoint{5.482183in}{2.303367in}}%
\pgfpathlineto{\pgfqpoint{5.482290in}{2.303430in}}%
\pgfpathlineto{\pgfqpoint{5.482290in}{2.303435in}}%
\pgfpathlineto{\pgfqpoint{5.482892in}{2.302941in}}%
\pgfpathlineto{\pgfqpoint{5.483379in}{2.303210in}}%
\pgfpathlineto{\pgfqpoint{5.484197in}{2.302741in}}%
\pgfpathlineto{\pgfqpoint{5.483434in}{2.303280in}}%
\pgfpathlineto{\pgfqpoint{5.484512in}{2.302983in}}%
\pgfpathlineto{\pgfqpoint{5.485091in}{2.303313in}}%
\pgfpathlineto{\pgfqpoint{5.485614in}{2.302785in}}%
\pgfpathlineto{\pgfqpoint{5.485617in}{2.302799in}}%
\pgfpathlineto{\pgfqpoint{5.485793in}{2.302770in}}%
\pgfpathlineto{\pgfqpoint{5.486134in}{2.303262in}}%
\pgfpathlineto{\pgfqpoint{5.486709in}{2.303066in}}%
\pgfpathlineto{\pgfqpoint{5.487527in}{2.303641in}}%
\pgfpathlineto{\pgfqpoint{5.486770in}{2.302900in}}%
\pgfpathlineto{\pgfqpoint{5.487831in}{2.303197in}}%
\pgfpathlineto{\pgfqpoint{5.488670in}{2.303385in}}%
\pgfpathlineto{\pgfqpoint{5.488243in}{2.302968in}}%
\pgfpathlineto{\pgfqpoint{5.488954in}{2.303061in}}%
\pgfpathlineto{\pgfqpoint{5.490035in}{2.302456in}}%
\pgfpathlineto{\pgfqpoint{5.489326in}{2.303366in}}%
\pgfpathlineto{\pgfqpoint{5.490293in}{2.302634in}}%
\pgfpathlineto{\pgfqpoint{5.490562in}{2.302798in}}%
\pgfpathlineto{\pgfqpoint{5.491249in}{2.302301in}}%
\pgfpathlineto{\pgfqpoint{5.491394in}{2.302393in}}%
\pgfpathlineto{\pgfqpoint{5.492398in}{2.303032in}}%
\pgfpathlineto{\pgfqpoint{5.491610in}{2.302230in}}%
\pgfpathlineto{\pgfqpoint{5.492532in}{2.302846in}}%
\pgfpathlineto{\pgfqpoint{5.492873in}{2.302632in}}%
\pgfpathlineto{\pgfqpoint{5.493444in}{2.303160in}}%
\pgfpathlineto{\pgfqpoint{5.493636in}{2.302982in}}%
\pgfpathlineto{\pgfqpoint{5.494552in}{2.302235in}}%
\pgfpathlineto{\pgfqpoint{5.493811in}{2.303036in}}%
\pgfpathlineto{\pgfqpoint{5.494807in}{2.302413in}}%
\pgfpathlineto{\pgfqpoint{5.494964in}{2.302476in}}%
\pgfpathlineto{\pgfqpoint{5.495856in}{2.301918in}}%
\pgfpathlineto{\pgfqpoint{5.495900in}{2.302084in}}%
\pgfpathlineto{\pgfqpoint{5.496011in}{2.301980in}}%
\pgfpathlineto{\pgfqpoint{5.496906in}{2.302509in}}%
\pgfpathlineto{\pgfqpoint{5.496982in}{2.302378in}}%
\pgfpathlineto{\pgfqpoint{5.498153in}{2.303235in}}%
\pgfpathlineto{\pgfqpoint{5.498185in}{2.303157in}}%
\pgfpathlineto{\pgfqpoint{5.500061in}{2.303550in}}%
\pgfpathlineto{\pgfqpoint{5.498190in}{2.303182in}}%
\pgfpathlineto{\pgfqpoint{5.500090in}{2.303660in}}%
\pgfpathlineto{\pgfqpoint{5.500426in}{2.303685in}}%
\pgfpathlineto{\pgfqpoint{5.500741in}{2.303076in}}%
\pgfpathlineto{\pgfqpoint{5.501187in}{2.303283in}}%
\pgfpathlineto{\pgfqpoint{5.502470in}{2.303536in}}%
\pgfpathlineto{\pgfqpoint{5.501304in}{2.303046in}}%
\pgfpathlineto{\pgfqpoint{5.502495in}{2.303474in}}%
\pgfpathlineto{\pgfqpoint{5.503608in}{2.303360in}}%
\pgfpathlineto{\pgfqpoint{5.502806in}{2.303763in}}%
\pgfpathlineto{\pgfqpoint{5.503611in}{2.303370in}}%
\pgfpathlineto{\pgfqpoint{5.503660in}{2.303441in}}%
\pgfpathlineto{\pgfqpoint{5.504491in}{2.302683in}}%
\pgfpathlineto{\pgfqpoint{5.504685in}{2.302849in}}%
\pgfpathlineto{\pgfqpoint{5.505384in}{2.303467in}}%
\pgfpathlineto{\pgfqpoint{5.504933in}{2.302698in}}%
\pgfpathlineto{\pgfqpoint{5.505834in}{2.302994in}}%
\pgfpathlineto{\pgfqpoint{5.506376in}{2.302820in}}%
\pgfpathlineto{\pgfqpoint{5.506169in}{2.303322in}}%
\pgfpathlineto{\pgfqpoint{5.506947in}{2.302969in}}%
\pgfpathlineto{\pgfqpoint{5.507962in}{2.303345in}}%
\pgfpathlineto{\pgfqpoint{5.507088in}{2.302746in}}%
\pgfpathlineto{\pgfqpoint{5.508081in}{2.303281in}}%
\pgfpathlineto{\pgfqpoint{5.509195in}{2.303070in}}%
\pgfpathlineto{\pgfqpoint{5.508478in}{2.303779in}}%
\pgfpathlineto{\pgfqpoint{5.509215in}{2.303135in}}%
\pgfpathlineto{\pgfqpoint{5.509884in}{2.303693in}}%
\pgfpathlineto{\pgfqpoint{5.509335in}{2.303000in}}%
\pgfpathlineto{\pgfqpoint{5.510340in}{2.303365in}}%
\pgfpathlineto{\pgfqpoint{5.510387in}{2.303202in}}%
\pgfpathlineto{\pgfqpoint{5.510833in}{2.303840in}}%
\pgfpathlineto{\pgfqpoint{5.511363in}{2.303807in}}%
\pgfpathlineto{\pgfqpoint{5.512115in}{2.304337in}}%
\pgfpathlineto{\pgfqpoint{5.511492in}{2.303676in}}%
\pgfpathlineto{\pgfqpoint{5.512498in}{2.304025in}}%
\pgfpathlineto{\pgfqpoint{5.513053in}{2.303357in}}%
\pgfpathlineto{\pgfqpoint{5.512626in}{2.304191in}}%
\pgfpathlineto{\pgfqpoint{5.513633in}{2.303545in}}%
\pgfpathlineto{\pgfqpoint{5.514420in}{2.303614in}}%
\pgfpathlineto{\pgfqpoint{5.514095in}{2.303967in}}%
\pgfpathlineto{\pgfqpoint{5.514667in}{2.303956in}}%
\pgfpathlineto{\pgfqpoint{5.514855in}{2.304170in}}%
\pgfpathlineto{\pgfqpoint{5.515720in}{2.303419in}}%
\pgfpathlineto{\pgfqpoint{5.515746in}{2.303480in}}%
\pgfpathlineto{\pgfqpoint{5.516031in}{2.303418in}}%
\pgfpathlineto{\pgfqpoint{5.516534in}{2.303803in}}%
\pgfpathlineto{\pgfqpoint{5.516851in}{2.303610in}}%
\pgfpathlineto{\pgfqpoint{5.517911in}{2.303931in}}%
\pgfpathlineto{\pgfqpoint{5.517116in}{2.303417in}}%
\pgfpathlineto{\pgfqpoint{5.517969in}{2.303821in}}%
\pgfpathlineto{\pgfqpoint{5.518989in}{2.303592in}}%
\pgfpathlineto{\pgfqpoint{5.518675in}{2.304160in}}%
\pgfpathlineto{\pgfqpoint{5.519090in}{2.303711in}}%
\pgfpathlineto{\pgfqpoint{5.519683in}{2.303526in}}%
\pgfpathlineto{\pgfqpoint{5.519994in}{2.303923in}}%
\pgfpathlineto{\pgfqpoint{5.520215in}{2.303699in}}%
\pgfpathlineto{\pgfqpoint{5.520482in}{2.303933in}}%
\pgfpathlineto{\pgfqpoint{5.520806in}{2.303519in}}%
\pgfpathlineto{\pgfqpoint{5.521325in}{2.303663in}}%
\pgfpathlineto{\pgfqpoint{5.521803in}{2.303923in}}%
\pgfpathlineto{\pgfqpoint{5.521397in}{2.303576in}}%
\pgfpathlineto{\pgfqpoint{5.522450in}{2.303714in}}%
\pgfpathlineto{\pgfqpoint{5.523088in}{2.303700in}}%
\pgfpathlineto{\pgfqpoint{5.522725in}{2.303998in}}%
\pgfpathlineto{\pgfqpoint{5.523426in}{2.304238in}}%
\pgfpathlineto{\pgfqpoint{5.524165in}{2.304494in}}%
\pgfpathlineto{\pgfqpoint{5.523832in}{2.303932in}}%
\pgfpathlineto{\pgfqpoint{5.524532in}{2.304095in}}%
\pgfpathlineto{\pgfqpoint{5.525667in}{2.304390in}}%
\pgfpathlineto{\pgfqpoint{5.524822in}{2.303738in}}%
\pgfpathlineto{\pgfqpoint{5.525683in}{2.304369in}}%
\pgfpathlineto{\pgfqpoint{5.526660in}{2.303910in}}%
\pgfpathlineto{\pgfqpoint{5.526153in}{2.304500in}}%
\pgfpathlineto{\pgfqpoint{5.526818in}{2.304008in}}%
\pgfpathlineto{\pgfqpoint{5.527917in}{2.304348in}}%
\pgfpathlineto{\pgfqpoint{5.526872in}{2.303773in}}%
\pgfpathlineto{\pgfqpoint{5.528050in}{2.304181in}}%
\pgfpathlineto{\pgfqpoint{5.533089in}{2.306422in}}%
\pgfpathlineto{\pgfqpoint{5.533371in}{2.306647in}}%
\pgfpathlineto{\pgfqpoint{5.534077in}{2.306052in}}%
\pgfpathlineto{\pgfqpoint{5.534180in}{2.306145in}}%
\pgfpathlineto{\pgfqpoint{5.535308in}{2.305785in}}%
\pgfpathlineto{\pgfqpoint{5.534245in}{2.306267in}}%
\pgfpathlineto{\pgfqpoint{5.535310in}{2.305791in}}%
\pgfpathlineto{\pgfqpoint{5.535765in}{2.305653in}}%
\pgfpathlineto{\pgfqpoint{5.536362in}{2.306220in}}%
\pgfpathlineto{\pgfqpoint{5.536380in}{2.306192in}}%
\pgfpathlineto{\pgfqpoint{5.537654in}{2.306123in}}%
\pgfpathlineto{\pgfqpoint{5.536422in}{2.306317in}}%
\pgfpathlineto{\pgfqpoint{5.537669in}{2.306219in}}%
\pgfpathlineto{\pgfqpoint{5.538020in}{2.306655in}}%
\pgfpathlineto{\pgfqpoint{5.538591in}{2.303378in}}%
\pgfpathlineto{\pgfqpoint{5.538613in}{2.303422in}}%
\pgfpathlineto{\pgfqpoint{5.539807in}{2.303359in}}%
\pgfpathlineto{\pgfqpoint{5.538630in}{2.303436in}}%
\pgfpathlineto{\pgfqpoint{5.539811in}{2.303391in}}%
\pgfpathlineto{\pgfqpoint{5.540881in}{2.303449in}}%
\pgfpathlineto{\pgfqpoint{5.539956in}{2.303367in}}%
\pgfpathlineto{\pgfqpoint{5.540924in}{2.303409in}}%
\pgfpathlineto{\pgfqpoint{5.541942in}{2.303367in}}%
\pgfpathlineto{\pgfqpoint{5.540972in}{2.303440in}}%
\pgfpathlineto{\pgfqpoint{5.542035in}{2.303408in}}%
\pgfpathlineto{\pgfqpoint{5.543144in}{2.303440in}}%
\pgfpathlineto{\pgfqpoint{5.542300in}{2.303365in}}%
\pgfpathlineto{\pgfqpoint{5.543147in}{2.303406in}}%
\pgfpathlineto{\pgfqpoint{5.544225in}{2.303356in}}%
\pgfpathlineto{\pgfqpoint{5.543172in}{2.303450in}}%
\pgfpathlineto{\pgfqpoint{5.544265in}{2.303409in}}%
\pgfpathlineto{\pgfqpoint{5.545408in}{2.303440in}}%
\pgfpathlineto{\pgfqpoint{5.544283in}{2.303385in}}%
\pgfpathlineto{\pgfqpoint{5.545418in}{2.303400in}}%
\pgfpathlineto{\pgfqpoint{5.545927in}{2.303324in}}%
\pgfpathlineto{\pgfqpoint{5.545450in}{2.303443in}}%
\pgfpathlineto{\pgfqpoint{5.546531in}{2.303411in}}%
\pgfpathlineto{\pgfqpoint{5.547513in}{2.303445in}}%
\pgfpathlineto{\pgfqpoint{5.546725in}{2.303371in}}%
\pgfpathlineto{\pgfqpoint{5.547642in}{2.303394in}}%
\pgfpathlineto{\pgfqpoint{5.548620in}{2.303452in}}%
\pgfpathlineto{\pgfqpoint{5.548372in}{2.303325in}}%
\pgfpathlineto{\pgfqpoint{5.548755in}{2.303416in}}%
\pgfpathlineto{\pgfqpoint{5.549852in}{2.303447in}}%
\pgfpathlineto{\pgfqpoint{5.548764in}{2.303384in}}%
\pgfpathlineto{\pgfqpoint{5.549880in}{2.303428in}}%
\pgfpathlineto{\pgfqpoint{5.549974in}{2.303337in}}%
\pgfpathlineto{\pgfqpoint{5.550091in}{2.303444in}}%
\pgfpathlineto{\pgfqpoint{5.550992in}{2.303415in}}%
\pgfpathlineto{\pgfqpoint{5.551785in}{2.303367in}}%
\pgfpathlineto{\pgfqpoint{5.551020in}{2.303442in}}%
\pgfpathlineto{\pgfqpoint{5.552103in}{2.303420in}}%
\pgfpathlineto{\pgfqpoint{5.553183in}{2.303445in}}%
\pgfpathlineto{\pgfqpoint{5.552616in}{2.303337in}}%
\pgfpathlineto{\pgfqpoint{5.553215in}{2.303427in}}%
\pgfpathlineto{\pgfqpoint{5.554254in}{2.303359in}}%
\pgfpathlineto{\pgfqpoint{5.553382in}{2.303451in}}%
\pgfpathlineto{\pgfqpoint{5.554327in}{2.303418in}}%
\pgfpathlineto{\pgfqpoint{5.555568in}{2.303444in}}%
\pgfpathlineto{\pgfqpoint{5.554342in}{2.303334in}}%
\pgfpathlineto{\pgfqpoint{5.555575in}{2.303425in}}%
\pgfpathlineto{\pgfqpoint{5.556275in}{2.303352in}}%
\pgfpathlineto{\pgfqpoint{5.555631in}{2.303449in}}%
\pgfpathlineto{\pgfqpoint{5.556686in}{2.303411in}}%
\pgfpathlineto{\pgfqpoint{5.557133in}{2.303447in}}%
\pgfpathlineto{\pgfqpoint{5.557694in}{2.302909in}}%
\pgfpathlineto{\pgfqpoint{5.557762in}{2.303026in}}%
\pgfpathlineto{\pgfqpoint{5.558672in}{2.302833in}}%
\pgfpathlineto{\pgfqpoint{5.558235in}{2.303410in}}%
\pgfpathlineto{\pgfqpoint{5.558877in}{2.303019in}}%
\pgfpathlineto{\pgfqpoint{5.559839in}{2.303752in}}%
\pgfpathlineto{\pgfqpoint{5.558898in}{2.302997in}}%
\pgfpathlineto{\pgfqpoint{5.560084in}{2.303584in}}%
\pgfpathlineto{\pgfqpoint{5.560985in}{2.303469in}}%
\pgfpathlineto{\pgfqpoint{5.560514in}{2.304051in}}%
\pgfpathlineto{\pgfqpoint{5.561189in}{2.303698in}}%
\pgfpathlineto{\pgfqpoint{5.561201in}{2.303744in}}%
\pgfpathlineto{\pgfqpoint{5.561877in}{2.303208in}}%
\pgfpathlineto{\pgfqpoint{5.562279in}{2.303371in}}%
\pgfpathlineto{\pgfqpoint{5.563228in}{2.304030in}}%
\pgfpathlineto{\pgfqpoint{5.562530in}{2.303262in}}%
\pgfpathlineto{\pgfqpoint{5.563445in}{2.303980in}}%
\pgfpathlineto{\pgfqpoint{5.564271in}{2.303557in}}%
\pgfpathlineto{\pgfqpoint{5.563726in}{2.304178in}}%
\pgfpathlineto{\pgfqpoint{5.564562in}{2.303808in}}%
\pgfpathlineto{\pgfqpoint{5.565540in}{2.303966in}}%
\pgfpathlineto{\pgfqpoint{5.564617in}{2.303729in}}%
\pgfpathlineto{\pgfqpoint{5.565682in}{2.303857in}}%
\pgfpathlineto{\pgfqpoint{5.565838in}{2.303676in}}%
\pgfpathlineto{\pgfqpoint{5.566754in}{2.304346in}}%
\pgfpathlineto{\pgfqpoint{5.566778in}{2.304231in}}%
\pgfpathlineto{\pgfqpoint{5.567849in}{2.304415in}}%
\pgfpathlineto{\pgfqpoint{5.567443in}{2.303902in}}%
\pgfpathlineto{\pgfqpoint{5.567894in}{2.304342in}}%
\pgfpathlineto{\pgfqpoint{5.568600in}{2.304041in}}%
\pgfpathlineto{\pgfqpoint{5.567945in}{2.304475in}}%
\pgfpathlineto{\pgfqpoint{5.569013in}{2.304348in}}%
\pgfpathlineto{\pgfqpoint{5.569629in}{2.304833in}}%
\pgfpathlineto{\pgfqpoint{5.569215in}{2.304199in}}%
\pgfpathlineto{\pgfqpoint{5.570127in}{2.304393in}}%
\pgfpathlineto{\pgfqpoint{5.570612in}{2.304001in}}%
\pgfpathlineto{\pgfqpoint{5.570204in}{2.304432in}}%
\pgfpathlineto{\pgfqpoint{5.571263in}{2.304250in}}%
\pgfpathlineto{\pgfqpoint{5.572775in}{2.304121in}}%
\pgfpathlineto{\pgfqpoint{5.571282in}{2.304298in}}%
\pgfpathlineto{\pgfqpoint{5.572817in}{2.304342in}}%
\pgfpathlineto{\pgfqpoint{5.573708in}{2.304597in}}%
\pgfpathlineto{\pgfqpoint{5.573083in}{2.304204in}}%
\pgfpathlineto{\pgfqpoint{5.573961in}{2.304352in}}%
\pgfpathlineto{\pgfqpoint{5.574653in}{2.304096in}}%
\pgfpathlineto{\pgfqpoint{5.574454in}{2.304438in}}%
\pgfpathlineto{\pgfqpoint{5.575074in}{2.304324in}}%
\pgfpathlineto{\pgfqpoint{5.575612in}{2.303903in}}%
\pgfpathlineto{\pgfqpoint{5.576091in}{2.304751in}}%
\pgfpathlineto{\pgfqpoint{5.576139in}{2.304692in}}%
\pgfpathlineto{\pgfqpoint{5.576318in}{2.304788in}}%
\pgfpathlineto{\pgfqpoint{5.576600in}{2.304281in}}%
\pgfpathlineto{\pgfqpoint{5.577239in}{2.304479in}}%
\pgfpathlineto{\pgfqpoint{5.583178in}{2.304439in}}%
\pgfpathlineto{\pgfqpoint{5.583183in}{2.304394in}}%
\pgfpathlineto{\pgfqpoint{5.583326in}{2.304246in}}%
\pgfpathlineto{\pgfqpoint{5.584265in}{2.304877in}}%
\pgfpathlineto{\pgfqpoint{5.584266in}{2.304865in}}%
\pgfpathlineto{\pgfqpoint{5.584793in}{2.304479in}}%
\pgfpathlineto{\pgfqpoint{5.584277in}{2.304890in}}%
\pgfpathlineto{\pgfqpoint{5.585419in}{2.304618in}}%
\pgfpathlineto{\pgfqpoint{5.586650in}{2.305021in}}%
\pgfpathlineto{\pgfqpoint{5.585458in}{2.304489in}}%
\pgfpathlineto{\pgfqpoint{5.586665in}{2.304963in}}%
\pgfpathlineto{\pgfqpoint{5.587195in}{2.304723in}}%
\pgfpathlineto{\pgfqpoint{5.587753in}{2.305218in}}%
\pgfpathlineto{\pgfqpoint{5.587763in}{2.305167in}}%
\pgfpathlineto{\pgfqpoint{5.588886in}{2.304837in}}%
\pgfpathlineto{\pgfqpoint{5.587825in}{2.305346in}}%
\pgfpathlineto{\pgfqpoint{5.589063in}{2.305070in}}%
\pgfpathlineto{\pgfqpoint{5.589148in}{2.305202in}}%
\pgfpathlineto{\pgfqpoint{5.589488in}{2.304742in}}%
\pgfpathlineto{\pgfqpoint{5.590166in}{2.304940in}}%
\pgfpathlineto{\pgfqpoint{5.590846in}{2.304395in}}%
\pgfpathlineto{\pgfqpoint{5.590179in}{2.304969in}}%
\pgfpathlineto{\pgfqpoint{5.591308in}{2.304698in}}%
\pgfpathlineto{\pgfqpoint{5.592178in}{2.304518in}}%
\pgfpathlineto{\pgfqpoint{5.592453in}{2.305092in}}%
\pgfpathlineto{\pgfqpoint{5.593494in}{2.304720in}}%
\pgfpathlineto{\pgfqpoint{5.592522in}{2.305158in}}%
\pgfpathlineto{\pgfqpoint{5.593581in}{2.304863in}}%
\pgfpathlineto{\pgfqpoint{5.598072in}{2.304460in}}%
\pgfpathlineto{\pgfqpoint{5.593585in}{2.304889in}}%
\pgfpathlineto{\pgfqpoint{5.598083in}{2.304523in}}%
\pgfpathlineto{\pgfqpoint{5.598223in}{2.304730in}}%
\pgfpathlineto{\pgfqpoint{5.599029in}{2.304090in}}%
\pgfpathlineto{\pgfqpoint{5.599191in}{2.304442in}}%
\pgfpathlineto{\pgfqpoint{5.599313in}{2.304488in}}%
\pgfpathlineto{\pgfqpoint{5.599883in}{2.303830in}}%
\pgfpathlineto{\pgfqpoint{5.600256in}{2.303874in}}%
\pgfpathlineto{\pgfqpoint{5.600843in}{2.303610in}}%
\pgfpathlineto{\pgfqpoint{5.601225in}{2.304327in}}%
\pgfpathlineto{\pgfqpoint{5.601353in}{2.304194in}}%
\pgfpathlineto{\pgfqpoint{5.602319in}{2.304013in}}%
\pgfpathlineto{\pgfqpoint{5.601409in}{2.304397in}}%
\pgfpathlineto{\pgfqpoint{5.602470in}{2.304289in}}%
\pgfpathlineto{\pgfqpoint{5.607528in}{2.303567in}}%
\pgfpathlineto{\pgfqpoint{5.602482in}{2.304365in}}%
\pgfpathlineto{\pgfqpoint{5.607537in}{2.303662in}}%
\pgfpathlineto{\pgfqpoint{5.608034in}{2.303182in}}%
\pgfpathlineto{\pgfqpoint{5.608661in}{2.303724in}}%
\pgfpathlineto{\pgfqpoint{5.609184in}{2.303682in}}%
\pgfpathlineto{\pgfqpoint{5.609600in}{2.304298in}}%
\pgfpathlineto{\pgfqpoint{5.609656in}{2.304280in}}%
\pgfpathlineto{\pgfqpoint{5.609682in}{2.304301in}}%
\pgfpathlineto{\pgfqpoint{5.610295in}{2.303780in}}%
\pgfpathlineto{\pgfqpoint{5.610739in}{2.303922in}}%
\pgfpathlineto{\pgfqpoint{5.611192in}{2.303707in}}%
\pgfpathlineto{\pgfqpoint{5.611527in}{2.304382in}}%
\pgfpathlineto{\pgfqpoint{5.611821in}{2.304183in}}%
\pgfpathlineto{\pgfqpoint{5.612322in}{2.304498in}}%
\pgfpathlineto{\pgfqpoint{5.612765in}{2.303853in}}%
\pgfpathlineto{\pgfqpoint{5.612912in}{2.303917in}}%
\pgfpathlineto{\pgfqpoint{5.613908in}{2.303352in}}%
\pgfpathlineto{\pgfqpoint{5.613348in}{2.303989in}}%
\pgfpathlineto{\pgfqpoint{5.614085in}{2.303411in}}%
\pgfpathlineto{\pgfqpoint{5.615215in}{2.303432in}}%
\pgfpathlineto{\pgfqpoint{5.614099in}{2.303361in}}%
\pgfpathlineto{\pgfqpoint{5.615229in}{2.303407in}}%
\pgfpathlineto{\pgfqpoint{5.616540in}{2.303431in}}%
\pgfpathlineto{\pgfqpoint{5.615231in}{2.303395in}}%
\pgfpathlineto{\pgfqpoint{5.616547in}{2.303392in}}%
\pgfpathlineto{\pgfqpoint{5.617570in}{2.303377in}}%
\pgfpathlineto{\pgfqpoint{5.616561in}{2.303446in}}%
\pgfpathlineto{\pgfqpoint{5.617662in}{2.303409in}}%
\pgfpathlineto{\pgfqpoint{5.618772in}{2.303440in}}%
\pgfpathlineto{\pgfqpoint{5.617679in}{2.303344in}}%
\pgfpathlineto{\pgfqpoint{5.618775in}{2.303403in}}%
\pgfpathlineto{\pgfqpoint{5.619866in}{2.303389in}}%
\pgfpathlineto{\pgfqpoint{5.618792in}{2.303445in}}%
\pgfpathlineto{\pgfqpoint{5.619891in}{2.303412in}}%
\pgfpathlineto{\pgfqpoint{5.620919in}{2.303447in}}%
\pgfpathlineto{\pgfqpoint{5.619981in}{2.303362in}}%
\pgfpathlineto{\pgfqpoint{5.621002in}{2.303422in}}%
\pgfpathlineto{\pgfqpoint{5.621873in}{2.303454in}}%
\pgfpathlineto{\pgfqpoint{5.621261in}{2.303355in}}%
\pgfpathlineto{\pgfqpoint{5.622113in}{2.303436in}}%
\pgfpathlineto{\pgfqpoint{5.622851in}{2.303354in}}%
\pgfpathlineto{\pgfqpoint{5.622244in}{2.303444in}}%
\pgfpathlineto{\pgfqpoint{5.623224in}{2.303423in}}%
\pgfpathlineto{\pgfqpoint{5.624332in}{2.303445in}}%
\pgfpathlineto{\pgfqpoint{5.623229in}{2.303383in}}%
\pgfpathlineto{\pgfqpoint{5.624338in}{2.303421in}}%
\pgfpathlineto{\pgfqpoint{5.625460in}{2.303381in}}%
\pgfpathlineto{\pgfqpoint{5.624424in}{2.303455in}}%
\pgfpathlineto{\pgfqpoint{5.625472in}{2.303418in}}%
\pgfpathlineto{\pgfqpoint{5.626546in}{2.303448in}}%
\pgfpathlineto{\pgfqpoint{5.625566in}{2.303359in}}%
\pgfpathlineto{\pgfqpoint{5.626588in}{2.303403in}}%
\pgfpathlineto{\pgfqpoint{5.627970in}{2.303400in}}%
\pgfpathlineto{\pgfqpoint{5.626592in}{2.303431in}}%
\pgfpathlineto{\pgfqpoint{5.627975in}{2.303418in}}%
\pgfpathlineto{\pgfqpoint{5.628811in}{2.303459in}}%
\pgfpathlineto{\pgfqpoint{5.628491in}{2.303336in}}%
\pgfpathlineto{\pgfqpoint{5.629088in}{2.303421in}}%
\pgfpathlineto{\pgfqpoint{5.630640in}{2.303399in}}%
\pgfpathlineto{\pgfqpoint{5.629090in}{2.303430in}}%
\pgfpathlineto{\pgfqpoint{5.630640in}{2.303422in}}%
\pgfpathlineto{\pgfqpoint{5.631694in}{2.303446in}}%
\pgfpathlineto{\pgfqpoint{5.631068in}{2.303352in}}%
\pgfpathlineto{\pgfqpoint{5.631752in}{2.303407in}}%
\pgfpathlineto{\pgfqpoint{5.633328in}{2.303680in}}%
\pgfpathlineto{\pgfqpoint{5.631765in}{2.303439in}}%
\pgfpathlineto{\pgfqpoint{5.633336in}{2.303786in}}%
\pgfpathlineto{\pgfqpoint{5.633405in}{2.303819in}}%
\pgfpathlineto{\pgfqpoint{5.633956in}{2.302914in}}%
\pgfpathlineto{\pgfqpoint{5.634403in}{2.303014in}}%
\pgfpathlineto{\pgfqpoint{5.634512in}{2.302769in}}%
\pgfpathlineto{\pgfqpoint{5.635103in}{2.303265in}}%
\pgfpathlineto{\pgfqpoint{5.635510in}{2.303204in}}%
\pgfpathlineto{\pgfqpoint{5.635524in}{2.303295in}}%
\pgfpathlineto{\pgfqpoint{5.636549in}{2.302410in}}%
\pgfpathlineto{\pgfqpoint{5.636583in}{2.302488in}}%
\pgfpathlineto{\pgfqpoint{5.636588in}{2.302452in}}%
\pgfpathlineto{\pgfqpoint{5.637584in}{2.303159in}}%
\pgfpathlineto{\pgfqpoint{5.637641in}{2.303134in}}%
\pgfpathlineto{\pgfqpoint{5.638117in}{2.303600in}}%
\pgfpathlineto{\pgfqpoint{5.638760in}{2.303248in}}%
\pgfpathlineto{\pgfqpoint{5.639745in}{2.302938in}}%
\pgfpathlineto{\pgfqpoint{5.638939in}{2.303479in}}%
\pgfpathlineto{\pgfqpoint{5.639891in}{2.303038in}}%
\pgfpathlineto{\pgfqpoint{5.640985in}{2.303316in}}%
\pgfpathlineto{\pgfqpoint{5.640124in}{2.302778in}}%
\pgfpathlineto{\pgfqpoint{5.641020in}{2.303231in}}%
\pgfpathlineto{\pgfqpoint{5.642080in}{2.302337in}}%
\pgfpathlineto{\pgfqpoint{5.641263in}{2.303349in}}%
\pgfpathlineto{\pgfqpoint{5.642224in}{2.302408in}}%
\pgfpathlineto{\pgfqpoint{5.643322in}{2.302772in}}%
\pgfpathlineto{\pgfqpoint{5.642958in}{2.302175in}}%
\pgfpathlineto{\pgfqpoint{5.643352in}{2.302725in}}%
\pgfpathlineto{\pgfqpoint{5.643469in}{2.302620in}}%
\pgfpathlineto{\pgfqpoint{5.644208in}{2.303344in}}%
\pgfpathlineto{\pgfqpoint{5.644400in}{2.303208in}}%
\pgfpathlineto{\pgfqpoint{5.645087in}{2.303473in}}%
\pgfpathlineto{\pgfqpoint{5.644667in}{2.303154in}}%
\pgfpathlineto{\pgfqpoint{5.645513in}{2.303191in}}%
\pgfpathlineto{\pgfqpoint{5.646567in}{2.302615in}}%
\pgfpathlineto{\pgfqpoint{5.645777in}{2.303288in}}%
\pgfpathlineto{\pgfqpoint{5.646679in}{2.302881in}}%
\pgfpathlineto{\pgfqpoint{5.647164in}{2.303255in}}%
\pgfpathlineto{\pgfqpoint{5.647572in}{2.302585in}}%
\pgfpathlineto{\pgfqpoint{5.647792in}{2.302921in}}%
\pgfpathlineto{\pgfqpoint{5.647830in}{2.303029in}}%
\pgfpathlineto{\pgfqpoint{5.648795in}{2.302403in}}%
\pgfpathlineto{\pgfqpoint{5.648831in}{2.302401in}}%
\pgfpathlineto{\pgfqpoint{5.649677in}{2.301975in}}%
\pgfpathlineto{\pgfqpoint{5.648974in}{2.302624in}}%
\pgfpathlineto{\pgfqpoint{5.650014in}{2.302126in}}%
\pgfpathlineto{\pgfqpoint{5.650409in}{2.302158in}}%
\pgfpathlineto{\pgfqpoint{5.650816in}{2.301662in}}%
\pgfpathlineto{\pgfqpoint{5.651090in}{2.301849in}}%
\pgfpathlineto{\pgfqpoint{5.652176in}{2.301142in}}%
\pgfpathlineto{\pgfqpoint{5.651260in}{2.301931in}}%
\pgfpathlineto{\pgfqpoint{5.652249in}{2.301248in}}%
\pgfpathlineto{\pgfqpoint{5.657018in}{2.301672in}}%
\pgfpathlineto{\pgfqpoint{5.657030in}{2.301577in}}%
\pgfpathlineto{\pgfqpoint{5.657096in}{2.301377in}}%
\pgfpathlineto{\pgfqpoint{5.657720in}{2.301857in}}%
\pgfpathlineto{\pgfqpoint{5.658134in}{2.301752in}}%
\pgfpathlineto{\pgfqpoint{5.659055in}{2.302146in}}%
\pgfpathlineto{\pgfqpoint{5.658571in}{2.301551in}}%
\pgfpathlineto{\pgfqpoint{5.659281in}{2.301974in}}%
\pgfpathlineto{\pgfqpoint{5.660358in}{2.301835in}}%
\pgfpathlineto{\pgfqpoint{5.659575in}{2.302266in}}%
\pgfpathlineto{\pgfqpoint{5.660397in}{2.301977in}}%
\pgfpathlineto{\pgfqpoint{5.660875in}{2.301685in}}%
\pgfpathlineto{\pgfqpoint{5.661401in}{2.302435in}}%
\pgfpathlineto{\pgfqpoint{5.661448in}{2.302382in}}%
\pgfpathlineto{\pgfqpoint{5.662557in}{2.302699in}}%
\pgfpathlineto{\pgfqpoint{5.661691in}{2.302044in}}%
\pgfpathlineto{\pgfqpoint{5.662628in}{2.302530in}}%
\pgfpathlineto{\pgfqpoint{5.663644in}{2.302038in}}%
\pgfpathlineto{\pgfqpoint{5.662644in}{2.302567in}}%
\pgfpathlineto{\pgfqpoint{5.663782in}{2.302129in}}%
\pgfpathlineto{\pgfqpoint{5.665140in}{2.301175in}}%
\pgfpathlineto{\pgfqpoint{5.663820in}{2.302212in}}%
\pgfpathlineto{\pgfqpoint{5.665217in}{2.301357in}}%
\pgfpathlineto{\pgfqpoint{5.666335in}{2.301869in}}%
\pgfpathlineto{\pgfqpoint{5.665376in}{2.301200in}}%
\pgfpathlineto{\pgfqpoint{5.666345in}{2.301797in}}%
\pgfpathlineto{\pgfqpoint{5.667536in}{2.302109in}}%
\pgfpathlineto{\pgfqpoint{5.666354in}{2.301815in}}%
\pgfpathlineto{\pgfqpoint{5.667584in}{2.302242in}}%
\pgfpathlineto{\pgfqpoint{5.667872in}{2.302702in}}%
\pgfpathlineto{\pgfqpoint{5.667601in}{2.302209in}}%
\pgfpathlineto{\pgfqpoint{5.668708in}{2.302379in}}%
\pgfpathlineto{\pgfqpoint{5.670430in}{2.303019in}}%
\pgfpathlineto{\pgfqpoint{5.668717in}{2.302361in}}%
\pgfpathlineto{\pgfqpoint{5.670516in}{2.302776in}}%
\pgfpathlineto{\pgfqpoint{5.670535in}{2.302827in}}%
\pgfpathlineto{\pgfqpoint{5.671675in}{2.302078in}}%
\pgfpathlineto{\pgfqpoint{5.672804in}{2.302290in}}%
\pgfpathlineto{\pgfqpoint{5.671854in}{2.301875in}}%
\pgfpathlineto{\pgfqpoint{5.672842in}{2.302252in}}%
\pgfpathlineto{\pgfqpoint{5.672844in}{2.302224in}}%
\pgfpathlineto{\pgfqpoint{5.673393in}{2.303037in}}%
\pgfpathlineto{\pgfqpoint{5.673930in}{2.302817in}}%
\pgfpathlineto{\pgfqpoint{5.675066in}{2.302897in}}%
\pgfpathlineto{\pgfqpoint{5.674355in}{2.302490in}}%
\pgfpathlineto{\pgfqpoint{5.675073in}{2.302872in}}%
\pgfpathlineto{\pgfqpoint{5.675830in}{2.302719in}}%
\pgfpathlineto{\pgfqpoint{5.675481in}{2.303238in}}%
\pgfpathlineto{\pgfqpoint{5.676179in}{2.303028in}}%
\pgfpathlineto{\pgfqpoint{5.677199in}{2.303456in}}%
\pgfpathlineto{\pgfqpoint{5.676550in}{2.303001in}}%
\pgfpathlineto{\pgfqpoint{5.677307in}{2.303334in}}%
\pgfpathlineto{\pgfqpoint{5.678359in}{2.303253in}}%
\pgfpathlineto{\pgfqpoint{5.677489in}{2.303507in}}%
\pgfpathlineto{\pgfqpoint{5.678426in}{2.303398in}}%
\pgfpathlineto{\pgfqpoint{5.679274in}{2.303836in}}%
\pgfpathlineto{\pgfqpoint{5.679555in}{2.303569in}}%
\pgfpathlineto{\pgfqpoint{5.679753in}{2.303280in}}%
\pgfpathlineto{\pgfqpoint{5.680526in}{2.304245in}}%
\pgfpathlineto{\pgfqpoint{5.680603in}{2.304192in}}%
\pgfpathlineto{\pgfqpoint{5.681439in}{2.304590in}}%
\pgfpathlineto{\pgfqpoint{5.680880in}{2.303883in}}%
\pgfpathlineto{\pgfqpoint{5.681757in}{2.304442in}}%
\pgfpathlineto{\pgfqpoint{5.681914in}{2.304364in}}%
\pgfpathlineto{\pgfqpoint{5.682575in}{2.304996in}}%
\pgfpathlineto{\pgfqpoint{5.682855in}{2.304699in}}%
\pgfpathlineto{\pgfqpoint{5.683940in}{2.304744in}}%
\pgfpathlineto{\pgfqpoint{5.683013in}{2.304334in}}%
\pgfpathlineto{\pgfqpoint{5.683965in}{2.304630in}}%
\pgfpathlineto{\pgfqpoint{5.685046in}{2.304112in}}%
\pgfpathlineto{\pgfqpoint{5.684006in}{2.304728in}}%
\pgfpathlineto{\pgfqpoint{5.685117in}{2.304248in}}%
\pgfpathlineto{\pgfqpoint{5.686092in}{2.303846in}}%
\pgfpathlineto{\pgfqpoint{5.685186in}{2.304451in}}%
\pgfpathlineto{\pgfqpoint{5.686260in}{2.304066in}}%
\pgfpathlineto{\pgfqpoint{5.687458in}{2.304848in}}%
\pgfpathlineto{\pgfqpoint{5.686263in}{2.304061in}}%
\pgfpathlineto{\pgfqpoint{5.687468in}{2.304800in}}%
\pgfpathlineto{\pgfqpoint{5.689164in}{2.303394in}}%
\pgfpathlineto{\pgfqpoint{5.688198in}{2.305604in}}%
\pgfpathlineto{\pgfqpoint{5.689190in}{2.303417in}}%
\pgfpathlineto{\pgfqpoint{5.690301in}{2.303443in}}%
\pgfpathlineto{\pgfqpoint{5.689357in}{2.303374in}}%
\pgfpathlineto{\pgfqpoint{5.690302in}{2.303427in}}%
\pgfpathlineto{\pgfqpoint{5.691235in}{2.303337in}}%
\pgfpathlineto{\pgfqpoint{5.690370in}{2.303442in}}%
\pgfpathlineto{\pgfqpoint{5.691414in}{2.303408in}}%
\pgfpathlineto{\pgfqpoint{5.692505in}{2.303439in}}%
\pgfpathlineto{\pgfqpoint{5.691895in}{2.303339in}}%
\pgfpathlineto{\pgfqpoint{5.692526in}{2.303413in}}%
\pgfpathlineto{\pgfqpoint{5.693587in}{2.303380in}}%
\pgfpathlineto{\pgfqpoint{5.692548in}{2.303445in}}%
\pgfpathlineto{\pgfqpoint{5.693643in}{2.303415in}}%
\pgfpathlineto{\pgfqpoint{5.694655in}{2.303450in}}%
\pgfpathlineto{\pgfqpoint{5.693679in}{2.303373in}}%
\pgfpathlineto{\pgfqpoint{5.694754in}{2.303414in}}%
\pgfpathlineto{\pgfqpoint{5.695851in}{2.303445in}}%
\pgfpathlineto{\pgfqpoint{5.694767in}{2.303386in}}%
\pgfpathlineto{\pgfqpoint{5.695869in}{2.303417in}}%
\pgfpathlineto{\pgfqpoint{5.696959in}{2.303375in}}%
\pgfpathlineto{\pgfqpoint{5.695884in}{2.303448in}}%
\pgfpathlineto{\pgfqpoint{5.696997in}{2.303412in}}%
\pgfpathlineto{\pgfqpoint{5.698038in}{2.303455in}}%
\pgfpathlineto{\pgfqpoint{5.697198in}{2.303348in}}%
\pgfpathlineto{\pgfqpoint{5.698107in}{2.303394in}}%
\pgfpathlineto{\pgfqpoint{5.699062in}{2.303454in}}%
\pgfpathlineto{\pgfqpoint{5.698209in}{2.303364in}}%
\pgfpathlineto{\pgfqpoint{5.699219in}{2.303378in}}%
\pgfpathlineto{\pgfqpoint{5.700302in}{2.303445in}}%
\pgfpathlineto{\pgfqpoint{5.699220in}{2.303377in}}%
\pgfpathlineto{\pgfqpoint{5.700348in}{2.303395in}}%
\pgfpathlineto{\pgfqpoint{5.701403in}{2.303369in}}%
\pgfpathlineto{\pgfqpoint{5.700351in}{2.303446in}}%
\pgfpathlineto{\pgfqpoint{5.701479in}{2.303415in}}%
\pgfpathlineto{\pgfqpoint{5.702532in}{2.303448in}}%
\pgfpathlineto{\pgfqpoint{5.701556in}{2.303368in}}%
\pgfpathlineto{\pgfqpoint{5.702591in}{2.303416in}}%
\pgfpathlineto{\pgfqpoint{5.703621in}{2.303356in}}%
\pgfpathlineto{\pgfqpoint{5.702609in}{2.303440in}}%
\pgfpathlineto{\pgfqpoint{5.703708in}{2.303396in}}%
\pgfpathlineto{\pgfqpoint{5.704845in}{2.303443in}}%
\pgfpathlineto{\pgfqpoint{5.703729in}{2.303380in}}%
\pgfpathlineto{\pgfqpoint{5.704850in}{2.303426in}}%
\pgfpathlineto{\pgfqpoint{5.705465in}{2.303360in}}%
\pgfpathlineto{\pgfqpoint{5.704858in}{2.303456in}}%
\pgfpathlineto{\pgfqpoint{5.705961in}{2.303421in}}%
\pgfpathlineto{\pgfqpoint{5.706095in}{2.303362in}}%
\pgfpathlineto{\pgfqpoint{5.707090in}{2.303713in}}%
\pgfpathlineto{\pgfqpoint{5.707223in}{2.303416in}}%
\pgfpathlineto{\pgfqpoint{5.708121in}{2.304639in}}%
\pgfpathlineto{\pgfqpoint{5.708163in}{2.304419in}}%
\pgfpathlineto{\pgfqpoint{5.708718in}{2.304166in}}%
\pgfpathlineto{\pgfqpoint{5.709245in}{2.304646in}}%
\pgfpathlineto{\pgfqpoint{5.709259in}{2.304621in}}%
\pgfpathlineto{\pgfqpoint{5.709828in}{2.305175in}}%
\pgfpathlineto{\pgfqpoint{5.710394in}{2.305038in}}%
\pgfpathlineto{\pgfqpoint{5.711346in}{2.304822in}}%
\pgfpathlineto{\pgfqpoint{5.710530in}{2.305182in}}%
\pgfpathlineto{\pgfqpoint{5.711518in}{2.304940in}}%
\pgfpathlineto{\pgfqpoint{5.711857in}{2.305157in}}%
\pgfpathlineto{\pgfqpoint{5.712398in}{2.304640in}}%
\pgfpathlineto{\pgfqpoint{5.712632in}{2.304910in}}%
\pgfpathlineto{\pgfqpoint{5.712994in}{2.304619in}}%
\pgfpathlineto{\pgfqpoint{5.713543in}{2.305244in}}%
\pgfpathlineto{\pgfqpoint{5.713742in}{2.304953in}}%
\pgfpathlineto{\pgfqpoint{5.714767in}{2.305102in}}%
\pgfpathlineto{\pgfqpoint{5.714438in}{2.304550in}}%
\pgfpathlineto{\pgfqpoint{5.714872in}{2.304998in}}%
\pgfpathlineto{\pgfqpoint{5.715981in}{2.304404in}}%
\pgfpathlineto{\pgfqpoint{5.714976in}{2.305167in}}%
\pgfpathlineto{\pgfqpoint{5.716030in}{2.304482in}}%
\pgfpathlineto{\pgfqpoint{5.717200in}{2.305324in}}%
\pgfpathlineto{\pgfqpoint{5.716102in}{2.304360in}}%
\pgfpathlineto{\pgfqpoint{5.717242in}{2.305251in}}%
\pgfpathlineto{\pgfqpoint{5.718357in}{2.304750in}}%
\pgfpathlineto{\pgfqpoint{5.717273in}{2.305322in}}%
\pgfpathlineto{\pgfqpoint{5.718422in}{2.304852in}}%
\pgfpathlineto{\pgfqpoint{5.719225in}{2.305391in}}%
\pgfpathlineto{\pgfqpoint{5.718503in}{2.304777in}}%
\pgfpathlineto{\pgfqpoint{5.719577in}{2.305312in}}%
\pgfpathlineto{\pgfqpoint{5.719644in}{2.305171in}}%
\pgfpathlineto{\pgfqpoint{5.720207in}{2.305814in}}%
\pgfpathlineto{\pgfqpoint{5.720683in}{2.305438in}}%
\pgfpathlineto{\pgfqpoint{5.725207in}{2.305676in}}%
\pgfpathlineto{\pgfqpoint{5.720693in}{2.305382in}}%
\pgfpathlineto{\pgfqpoint{5.725212in}{2.305636in}}%
\pgfpathlineto{\pgfqpoint{5.726017in}{2.305410in}}%
\pgfpathlineto{\pgfqpoint{5.725360in}{2.305872in}}%
\pgfpathlineto{\pgfqpoint{5.726321in}{2.305705in}}%
\pgfpathlineto{\pgfqpoint{5.726847in}{2.306225in}}%
\pgfpathlineto{\pgfqpoint{5.726413in}{2.305555in}}%
\pgfpathlineto{\pgfqpoint{5.727443in}{2.305800in}}%
\pgfpathlineto{\pgfqpoint{5.727452in}{2.305770in}}%
\pgfpathlineto{\pgfqpoint{5.728222in}{2.306411in}}%
\pgfpathlineto{\pgfqpoint{5.728532in}{2.306219in}}%
\pgfpathlineto{\pgfqpoint{5.729546in}{2.306514in}}%
\pgfpathlineto{\pgfqpoint{5.729669in}{2.306176in}}%
\pgfpathlineto{\pgfqpoint{5.730795in}{2.305734in}}%
\pgfpathlineto{\pgfqpoint{5.729783in}{2.306253in}}%
\pgfpathlineto{\pgfqpoint{5.730803in}{2.305749in}}%
\pgfpathlineto{\pgfqpoint{5.731675in}{2.305962in}}%
\pgfpathlineto{\pgfqpoint{5.731346in}{2.305355in}}%
\pgfpathlineto{\pgfqpoint{5.731931in}{2.305890in}}%
\pgfpathlineto{\pgfqpoint{5.733408in}{2.305022in}}%
\pgfpathlineto{\pgfqpoint{5.732039in}{2.306082in}}%
\pgfpathlineto{\pgfqpoint{5.733496in}{2.305211in}}%
\pgfpathlineto{\pgfqpoint{5.734527in}{2.305628in}}%
\pgfpathlineto{\pgfqpoint{5.733620in}{2.305003in}}%
\pgfpathlineto{\pgfqpoint{5.734622in}{2.305509in}}%
\pgfpathlineto{\pgfqpoint{5.735070in}{2.305227in}}%
\pgfpathlineto{\pgfqpoint{5.735718in}{2.305888in}}%
\pgfpathlineto{\pgfqpoint{5.735719in}{2.305876in}}%
\pgfpathlineto{\pgfqpoint{5.736850in}{2.305049in}}%
\pgfpathlineto{\pgfqpoint{5.735817in}{2.306007in}}%
\pgfpathlineto{\pgfqpoint{5.736915in}{2.305209in}}%
\pgfpathlineto{\pgfqpoint{5.737976in}{2.305315in}}%
\pgfpathlineto{\pgfqpoint{5.737432in}{2.304825in}}%
\pgfpathlineto{\pgfqpoint{5.738029in}{2.305269in}}%
\pgfpathlineto{\pgfqpoint{5.739290in}{2.305871in}}%
\pgfpathlineto{\pgfqpoint{5.738031in}{2.305244in}}%
\pgfpathlineto{\pgfqpoint{5.739294in}{2.305853in}}%
\pgfpathlineto{\pgfqpoint{5.740042in}{2.305358in}}%
\pgfpathlineto{\pgfqpoint{5.739305in}{2.305859in}}%
\pgfpathlineto{\pgfqpoint{5.740433in}{2.305731in}}%
\pgfpathlineto{\pgfqpoint{5.741358in}{2.305833in}}%
\pgfpathlineto{\pgfqpoint{5.740434in}{2.305718in}}%
\pgfpathlineto{\pgfqpoint{5.741520in}{2.305406in}}%
\pgfpathlineto{\pgfqpoint{5.742463in}{2.305410in}}%
\pgfpathlineto{\pgfqpoint{5.741819in}{2.305754in}}%
\pgfpathlineto{\pgfqpoint{5.742621in}{2.305559in}}%
\pgfpathlineto{\pgfqpoint{5.742784in}{2.305719in}}%
\pgfpathlineto{\pgfqpoint{5.743691in}{2.304895in}}%
\pgfpathlineto{\pgfqpoint{5.743700in}{2.304950in}}%
\pgfpathlineto{\pgfqpoint{5.744522in}{2.304339in}}%
\pgfpathlineto{\pgfqpoint{5.743791in}{2.305068in}}%
\pgfpathlineto{\pgfqpoint{5.744830in}{2.304614in}}%
\pgfpathlineto{\pgfqpoint{5.744876in}{2.304740in}}%
\pgfpathlineto{\pgfqpoint{5.745133in}{2.304150in}}%
\pgfpathlineto{\pgfqpoint{5.745891in}{2.304114in}}%
\pgfpathlineto{\pgfqpoint{5.745973in}{2.304028in}}%
\pgfpathlineto{\pgfqpoint{5.746850in}{2.304861in}}%
\pgfpathlineto{\pgfqpoint{5.746851in}{2.304860in}}%
\pgfpathlineto{\pgfqpoint{5.746938in}{2.304979in}}%
\pgfpathlineto{\pgfqpoint{5.747410in}{2.304087in}}%
\pgfpathlineto{\pgfqpoint{5.747795in}{2.304068in}}%
\pgfpathlineto{\pgfqpoint{5.747880in}{2.304008in}}%
\pgfpathlineto{\pgfqpoint{5.748784in}{2.304986in}}%
\pgfpathlineto{\pgfqpoint{5.748886in}{2.304709in}}%
\pgfpathlineto{\pgfqpoint{5.749797in}{2.304901in}}%
\pgfpathlineto{\pgfqpoint{5.748985in}{2.304606in}}%
\pgfpathlineto{\pgfqpoint{5.749994in}{2.304628in}}%
\pgfpathlineto{\pgfqpoint{5.750755in}{2.304299in}}%
\pgfpathlineto{\pgfqpoint{5.751023in}{2.304773in}}%
\pgfpathlineto{\pgfqpoint{5.751106in}{2.304604in}}%
\pgfpathlineto{\pgfqpoint{5.751738in}{2.304843in}}%
\pgfpathlineto{\pgfqpoint{5.751385in}{2.304546in}}%
\pgfpathlineto{\pgfqpoint{5.752224in}{2.304629in}}%
\pgfpathlineto{\pgfqpoint{5.752244in}{2.304589in}}%
\pgfpathlineto{\pgfqpoint{5.752947in}{2.305381in}}%
\pgfpathlineto{\pgfqpoint{5.753272in}{2.305202in}}%
\pgfpathlineto{\pgfqpoint{5.753483in}{2.305389in}}%
\pgfpathlineto{\pgfqpoint{5.754270in}{2.304186in}}%
\pgfpathlineto{\pgfqpoint{5.754280in}{2.304219in}}%
\pgfpathlineto{\pgfqpoint{5.754690in}{2.303838in}}%
\pgfpathlineto{\pgfqpoint{5.755168in}{2.304307in}}%
\pgfpathlineto{\pgfqpoint{5.755403in}{2.304093in}}%
\pgfpathlineto{\pgfqpoint{5.756612in}{2.303455in}}%
\pgfpathlineto{\pgfqpoint{5.755415in}{2.304119in}}%
\pgfpathlineto{\pgfqpoint{5.756622in}{2.303490in}}%
\pgfpathlineto{\pgfqpoint{5.756689in}{2.303349in}}%
\pgfpathlineto{\pgfqpoint{5.757269in}{2.304165in}}%
\pgfpathlineto{\pgfqpoint{5.757416in}{2.304259in}}%
\pgfpathlineto{\pgfqpoint{5.758223in}{2.304612in}}%
\pgfpathlineto{\pgfqpoint{5.757446in}{2.304211in}}%
\pgfpathlineto{\pgfqpoint{5.758547in}{2.304448in}}%
\pgfpathlineto{\pgfqpoint{5.759230in}{2.304925in}}%
\pgfpathlineto{\pgfqpoint{5.758773in}{2.304364in}}%
\pgfpathlineto{\pgfqpoint{5.759675in}{2.304464in}}%
\pgfpathlineto{\pgfqpoint{5.759723in}{2.304407in}}%
\pgfpathlineto{\pgfqpoint{5.760430in}{2.305040in}}%
\pgfpathlineto{\pgfqpoint{5.760723in}{2.304841in}}%
\pgfpathlineto{\pgfqpoint{5.761858in}{2.305291in}}%
\pgfpathlineto{\pgfqpoint{5.760898in}{2.304603in}}%
\pgfpathlineto{\pgfqpoint{5.761862in}{2.305278in}}%
\pgfpathlineto{\pgfqpoint{5.761864in}{2.305251in}}%
\pgfpathlineto{\pgfqpoint{5.762935in}{2.305787in}}%
\pgfpathlineto{\pgfqpoint{5.762950in}{2.305774in}}%
\pgfpathlineto{\pgfqpoint{5.764317in}{2.303360in}}%
\pgfpathlineto{\pgfqpoint{5.763133in}{2.305852in}}%
\pgfpathlineto{\pgfqpoint{5.764337in}{2.303423in}}%
\pgfpathlineto{\pgfqpoint{5.765287in}{2.303450in}}%
\pgfpathlineto{\pgfqpoint{5.764412in}{2.303377in}}%
\pgfpathlineto{\pgfqpoint{5.765449in}{2.303410in}}%
\pgfpathlineto{\pgfqpoint{5.766556in}{2.303439in}}%
\pgfpathlineto{\pgfqpoint{5.765479in}{2.303343in}}%
\pgfpathlineto{\pgfqpoint{5.766569in}{2.303396in}}%
\pgfpathlineto{\pgfqpoint{5.767680in}{2.303385in}}%
\pgfpathlineto{\pgfqpoint{5.766633in}{2.303451in}}%
\pgfpathlineto{\pgfqpoint{5.767683in}{2.303405in}}%
\pgfpathlineto{\pgfqpoint{5.768786in}{2.303449in}}%
\pgfpathlineto{\pgfqpoint{5.767742in}{2.303363in}}%
\pgfpathlineto{\pgfqpoint{5.768794in}{2.303397in}}%
\pgfpathlineto{\pgfqpoint{5.769881in}{2.303443in}}%
\pgfpathlineto{\pgfqpoint{5.768797in}{2.303373in}}%
\pgfpathlineto{\pgfqpoint{5.769908in}{2.303432in}}%
\pgfpathlineto{\pgfqpoint{5.770194in}{2.303353in}}%
\pgfpathlineto{\pgfqpoint{5.769921in}{2.303451in}}%
\pgfpathlineto{\pgfqpoint{5.771020in}{2.303430in}}%
\pgfpathlineto{\pgfqpoint{5.771649in}{2.303350in}}%
\pgfpathlineto{\pgfqpoint{5.771045in}{2.303464in}}%
\pgfpathlineto{\pgfqpoint{5.772131in}{2.303416in}}%
\pgfpathlineto{\pgfqpoint{5.773202in}{2.303450in}}%
\pgfpathlineto{\pgfqpoint{5.772233in}{2.303357in}}%
\pgfpathlineto{\pgfqpoint{5.773247in}{2.303428in}}%
\pgfpathlineto{\pgfqpoint{5.774023in}{2.303356in}}%
\pgfpathlineto{\pgfqpoint{5.773269in}{2.303447in}}%
\pgfpathlineto{\pgfqpoint{5.774358in}{2.303418in}}%
\pgfpathlineto{\pgfqpoint{5.775455in}{2.303440in}}%
\pgfpathlineto{\pgfqpoint{5.774412in}{2.303367in}}%
\pgfpathlineto{\pgfqpoint{5.775472in}{2.303404in}}%
\pgfpathlineto{\pgfqpoint{5.776450in}{2.303340in}}%
\pgfpathlineto{\pgfqpoint{5.775508in}{2.303446in}}%
\pgfpathlineto{\pgfqpoint{5.776584in}{2.303417in}}%
\pgfpathlineto{\pgfqpoint{5.777568in}{2.303449in}}%
\pgfpathlineto{\pgfqpoint{5.776639in}{2.303372in}}%
\pgfpathlineto{\pgfqpoint{5.777695in}{2.303424in}}%
\pgfpathlineto{\pgfqpoint{5.778783in}{2.303382in}}%
\pgfpathlineto{\pgfqpoint{5.777748in}{2.303447in}}%
\pgfpathlineto{\pgfqpoint{5.778823in}{2.303422in}}%
\pgfpathlineto{\pgfqpoint{5.779930in}{2.303438in}}%
\pgfpathlineto{\pgfqpoint{5.778853in}{2.303367in}}%
\pgfpathlineto{\pgfqpoint{5.779936in}{2.303392in}}%
\pgfpathlineto{\pgfqpoint{5.780941in}{2.303373in}}%
\pgfpathlineto{\pgfqpoint{5.780054in}{2.303455in}}%
\pgfpathlineto{\pgfqpoint{5.781046in}{2.303420in}}%
\pgfpathlineto{\pgfqpoint{5.782063in}{2.303682in}}%
\pgfpathlineto{\pgfqpoint{5.781112in}{2.303375in}}%
\pgfpathlineto{\pgfqpoint{5.782192in}{2.303593in}}%
\pgfpathlineto{\pgfqpoint{5.782193in}{2.303574in}}%
\pgfpathlineto{\pgfqpoint{5.783019in}{2.304204in}}%
\pgfpathlineto{\pgfqpoint{5.783290in}{2.303815in}}%
\pgfpathlineto{\pgfqpoint{5.785134in}{2.304461in}}%
\pgfpathlineto{\pgfqpoint{5.785140in}{2.304512in}}%
\pgfpathlineto{\pgfqpoint{5.785280in}{2.304679in}}%
\pgfpathlineto{\pgfqpoint{5.786146in}{2.304004in}}%
\pgfpathlineto{\pgfqpoint{5.786169in}{2.304011in}}%
\pgfpathlineto{\pgfqpoint{5.786253in}{2.303841in}}%
\pgfpathlineto{\pgfqpoint{5.786662in}{2.304217in}}%
\pgfpathlineto{\pgfqpoint{5.787281in}{2.303977in}}%
\pgfpathlineto{\pgfqpoint{5.788190in}{2.303795in}}%
\pgfpathlineto{\pgfqpoint{5.787545in}{2.304247in}}%
\pgfpathlineto{\pgfqpoint{5.788392in}{2.304015in}}%
\pgfpathlineto{\pgfqpoint{5.788463in}{2.304120in}}%
\pgfpathlineto{\pgfqpoint{5.789159in}{2.303692in}}%
\pgfpathlineto{\pgfqpoint{5.789497in}{2.303872in}}%
\pgfpathlineto{\pgfqpoint{5.789567in}{2.303739in}}%
\pgfpathlineto{\pgfqpoint{5.790381in}{2.304302in}}%
\pgfpathlineto{\pgfqpoint{5.790586in}{2.304232in}}%
\pgfpathlineto{\pgfqpoint{5.790637in}{2.304077in}}%
\pgfpathlineto{\pgfqpoint{5.791071in}{2.304651in}}%
\pgfpathlineto{\pgfqpoint{5.791643in}{2.304558in}}%
\pgfpathlineto{\pgfqpoint{5.792761in}{2.305091in}}%
\pgfpathlineto{\pgfqpoint{5.791777in}{2.304465in}}%
\pgfpathlineto{\pgfqpoint{5.792775in}{2.305040in}}%
\pgfpathlineto{\pgfqpoint{5.793973in}{2.304694in}}%
\pgfpathlineto{\pgfqpoint{5.793540in}{2.305326in}}%
\pgfpathlineto{\pgfqpoint{5.793976in}{2.304719in}}%
\pgfpathlineto{\pgfqpoint{5.795012in}{2.304591in}}%
\pgfpathlineto{\pgfqpoint{5.794302in}{2.304384in}}%
\pgfpathlineto{\pgfqpoint{5.795130in}{2.304464in}}%
\pgfpathlineto{\pgfqpoint{5.795635in}{2.304612in}}%
\pgfpathlineto{\pgfqpoint{5.796265in}{2.304112in}}%
\pgfpathlineto{\pgfqpoint{5.797378in}{2.304410in}}%
\pgfpathlineto{\pgfqpoint{5.796361in}{2.304046in}}%
\pgfpathlineto{\pgfqpoint{5.797429in}{2.304350in}}%
\pgfpathlineto{\pgfqpoint{5.798155in}{2.304026in}}%
\pgfpathlineto{\pgfqpoint{5.797552in}{2.304428in}}%
\pgfpathlineto{\pgfqpoint{5.798547in}{2.304341in}}%
\pgfpathlineto{\pgfqpoint{5.799434in}{2.305285in}}%
\pgfpathlineto{\pgfqpoint{5.799695in}{2.304944in}}%
\pgfpathlineto{\pgfqpoint{5.800833in}{2.304478in}}%
\pgfpathlineto{\pgfqpoint{5.799944in}{2.305066in}}%
\pgfpathlineto{\pgfqpoint{5.800834in}{2.304481in}}%
\pgfpathlineto{\pgfqpoint{5.801624in}{2.304827in}}%
\pgfpathlineto{\pgfqpoint{5.800925in}{2.304407in}}%
\pgfpathlineto{\pgfqpoint{5.801982in}{2.304608in}}%
\pgfpathlineto{\pgfqpoint{5.802755in}{2.304500in}}%
\pgfpathlineto{\pgfqpoint{5.802340in}{2.304864in}}%
\pgfpathlineto{\pgfqpoint{5.803069in}{2.304832in}}%
\pgfpathlineto{\pgfqpoint{5.803552in}{2.304969in}}%
\pgfpathlineto{\pgfqpoint{5.804126in}{2.304368in}}%
\pgfpathlineto{\pgfqpoint{5.804158in}{2.304476in}}%
\pgfpathlineto{\pgfqpoint{5.804628in}{2.304063in}}%
\pgfpathlineto{\pgfqpoint{5.805230in}{2.304772in}}%
\pgfpathlineto{\pgfqpoint{5.805253in}{2.304727in}}%
\pgfpathlineto{\pgfqpoint{5.813250in}{2.307300in}}%
\pgfpathlineto{\pgfqpoint{5.805309in}{2.304534in}}%
\pgfpathlineto{\pgfqpoint{5.813251in}{2.307290in}}%
\pgfpathlineto{\pgfqpoint{5.814024in}{2.307019in}}%
\pgfpathlineto{\pgfqpoint{5.814193in}{2.307345in}}%
\pgfpathlineto{\pgfqpoint{5.814368in}{2.307206in}}%
\pgfpathlineto{\pgfqpoint{5.814369in}{2.307220in}}%
\pgfpathlineto{\pgfqpoint{5.814920in}{2.306669in}}%
\pgfpathlineto{\pgfqpoint{5.815453in}{2.306894in}}%
\pgfpathlineto{\pgfqpoint{5.816663in}{2.305948in}}%
\pgfpathlineto{\pgfqpoint{5.815523in}{2.307055in}}%
\pgfpathlineto{\pgfqpoint{5.816667in}{2.305979in}}%
\pgfpathlineto{\pgfqpoint{5.817753in}{2.306604in}}%
\pgfpathlineto{\pgfqpoint{5.817001in}{2.305603in}}%
\pgfpathlineto{\pgfqpoint{5.817794in}{2.306421in}}%
\pgfpathlineto{\pgfqpoint{5.817917in}{2.306349in}}%
\pgfpathlineto{\pgfqpoint{5.818726in}{2.307120in}}%
\pgfpathlineto{\pgfqpoint{5.818855in}{2.307073in}}%
\pgfpathlineto{\pgfqpoint{5.819306in}{2.307149in}}%
\pgfpathlineto{\pgfqpoint{5.819940in}{2.306762in}}%
\pgfpathlineto{\pgfqpoint{5.819955in}{2.306797in}}%
\pgfpathlineto{\pgfqpoint{5.819957in}{2.306805in}}%
\pgfpathlineto{\pgfqpoint{5.820976in}{2.306229in}}%
\pgfpathlineto{\pgfqpoint{5.820981in}{2.306248in}}%
\pgfpathlineto{\pgfqpoint{5.821850in}{2.306531in}}%
\pgfpathlineto{\pgfqpoint{5.821002in}{2.306160in}}%
\pgfpathlineto{\pgfqpoint{5.822247in}{2.306105in}}%
\pgfpathlineto{\pgfqpoint{5.822396in}{2.305747in}}%
\pgfpathlineto{\pgfqpoint{5.823130in}{2.306361in}}%
\pgfpathlineto{\pgfqpoint{5.823354in}{2.306191in}}%
\pgfpathlineto{\pgfqpoint{5.823628in}{2.306438in}}%
\pgfpathlineto{\pgfqpoint{5.823913in}{2.305972in}}%
\pgfpathlineto{\pgfqpoint{5.824460in}{2.306089in}}%
\pgfpathlineto{\pgfqpoint{5.824541in}{2.305913in}}%
\pgfpathlineto{\pgfqpoint{5.825383in}{2.306677in}}%
\pgfpathlineto{\pgfqpoint{5.825508in}{2.306589in}}%
\pgfpathlineto{\pgfqpoint{5.826324in}{2.306724in}}%
\pgfpathlineto{\pgfqpoint{5.825661in}{2.306271in}}%
\pgfpathlineto{\pgfqpoint{5.826627in}{2.306567in}}%
\pgfpathlineto{\pgfqpoint{5.827338in}{2.306161in}}%
\pgfpathlineto{\pgfqpoint{5.826818in}{2.306759in}}%
\pgfpathlineto{\pgfqpoint{5.827745in}{2.306417in}}%
\pgfpathlineto{\pgfqpoint{5.828791in}{2.306367in}}%
\pgfpathlineto{\pgfqpoint{5.828126in}{2.305895in}}%
\pgfpathlineto{\pgfqpoint{5.828866in}{2.306288in}}%
\pgfpathlineto{\pgfqpoint{5.829027in}{2.306449in}}%
\pgfpathlineto{\pgfqpoint{5.829999in}{2.305919in}}%
\pgfpathlineto{\pgfqpoint{5.830945in}{2.306436in}}%
\pgfpathlineto{\pgfqpoint{5.831137in}{2.306333in}}%
\pgfpathlineto{\pgfqpoint{5.832667in}{2.306773in}}%
\pgfpathlineto{\pgfqpoint{5.831196in}{2.306245in}}%
\pgfpathlineto{\pgfqpoint{5.832687in}{2.306706in}}%
\pgfpathlineto{\pgfqpoint{5.833037in}{2.306344in}}%
\pgfpathlineto{\pgfqpoint{5.833503in}{2.306842in}}%
\pgfpathlineto{\pgfqpoint{5.833807in}{2.306605in}}%
\pgfpathlineto{\pgfqpoint{5.834843in}{2.306189in}}%
\pgfpathlineto{\pgfqpoint{5.833834in}{2.306665in}}%
\pgfpathlineto{\pgfqpoint{5.834958in}{2.306322in}}%
\pgfpathlineto{\pgfqpoint{5.835332in}{2.306265in}}%
\pgfpathlineto{\pgfqpoint{5.835948in}{2.306861in}}%
\pgfpathlineto{\pgfqpoint{5.835964in}{2.306848in}}%
\pgfpathlineto{\pgfqpoint{5.836628in}{2.307088in}}%
\pgfpathlineto{\pgfqpoint{5.836257in}{2.306666in}}%
\pgfpathlineto{\pgfqpoint{5.837087in}{2.306907in}}%
\pgfpathlineto{\pgfqpoint{5.838269in}{2.303375in}}%
\pgfpathlineto{\pgfqpoint{5.837690in}{2.307352in}}%
\pgfpathlineto{\pgfqpoint{5.838396in}{2.303395in}}%
\pgfpathlineto{\pgfqpoint{5.839507in}{2.303446in}}%
\pgfpathlineto{\pgfqpoint{5.838399in}{2.303368in}}%
\pgfpathlineto{\pgfqpoint{5.839510in}{2.303418in}}%
\pgfpathlineto{\pgfqpoint{5.840385in}{2.303350in}}%
\pgfpathlineto{\pgfqpoint{5.839650in}{2.303452in}}%
\pgfpathlineto{\pgfqpoint{5.840621in}{2.303411in}}%
\pgfpathlineto{\pgfqpoint{5.841691in}{2.303443in}}%
\pgfpathlineto{\pgfqpoint{5.840704in}{2.303356in}}%
\pgfpathlineto{\pgfqpoint{5.841733in}{2.303408in}}%
\pgfpathlineto{\pgfqpoint{5.842772in}{2.303353in}}%
\pgfpathlineto{\pgfqpoint{5.841737in}{2.303441in}}%
\pgfpathlineto{\pgfqpoint{5.842867in}{2.303418in}}%
\pgfpathlineto{\pgfqpoint{5.843870in}{2.303450in}}%
\pgfpathlineto{\pgfqpoint{5.842882in}{2.303360in}}%
\pgfpathlineto{\pgfqpoint{5.843980in}{2.303419in}}%
\pgfpathlineto{\pgfqpoint{5.844699in}{2.303333in}}%
\pgfpathlineto{\pgfqpoint{5.844067in}{2.303450in}}%
\pgfpathlineto{\pgfqpoint{5.845092in}{2.303421in}}%
\pgfpathlineto{\pgfqpoint{5.846210in}{2.303439in}}%
\pgfpathlineto{\pgfqpoint{5.845144in}{2.303379in}}%
\pgfpathlineto{\pgfqpoint{5.846213in}{2.303428in}}%
\pgfpathlineto{\pgfqpoint{5.846975in}{2.303357in}}%
\pgfpathlineto{\pgfqpoint{5.846228in}{2.303453in}}%
\pgfpathlineto{\pgfqpoint{5.847325in}{2.303415in}}%
\pgfpathlineto{\pgfqpoint{5.848401in}{2.303445in}}%
\pgfpathlineto{\pgfqpoint{5.847435in}{2.303370in}}%
\pgfpathlineto{\pgfqpoint{5.848445in}{2.303424in}}%
\pgfpathlineto{\pgfqpoint{5.848727in}{2.303344in}}%
\pgfpathlineto{\pgfqpoint{5.848456in}{2.303448in}}%
\pgfpathlineto{\pgfqpoint{5.849558in}{2.303404in}}%
\pgfpathlineto{\pgfqpoint{5.850544in}{2.303343in}}%
\pgfpathlineto{\pgfqpoint{5.849638in}{2.303446in}}%
\pgfpathlineto{\pgfqpoint{5.850670in}{2.303422in}}%
\pgfpathlineto{\pgfqpoint{5.851773in}{2.303444in}}%
\pgfpathlineto{\pgfqpoint{5.850754in}{2.303368in}}%
\pgfpathlineto{\pgfqpoint{5.851787in}{2.303415in}}%
\pgfpathlineto{\pgfqpoint{5.853467in}{2.303381in}}%
\pgfpathlineto{\pgfqpoint{5.851790in}{2.303435in}}%
\pgfpathlineto{\pgfqpoint{5.853475in}{2.303402in}}%
\pgfpathlineto{\pgfqpoint{5.854580in}{2.303452in}}%
\pgfpathlineto{\pgfqpoint{5.853514in}{2.303351in}}%
\pgfpathlineto{\pgfqpoint{5.854587in}{2.303415in}}%
\pgfpathlineto{\pgfqpoint{5.855679in}{2.303442in}}%
\pgfpathlineto{\pgfqpoint{5.854776in}{2.303361in}}%
\pgfpathlineto{\pgfqpoint{5.855701in}{2.303428in}}%
\pgfpathlineto{\pgfqpoint{5.856538in}{2.303333in}}%
\pgfpathlineto{\pgfqpoint{5.855806in}{2.303443in}}%
\pgfpathlineto{\pgfqpoint{5.856812in}{2.303388in}}%
\pgfpathlineto{\pgfqpoint{5.859887in}{2.302464in}}%
\pgfpathlineto{\pgfqpoint{5.859888in}{2.302433in}}%
\pgfpathlineto{\pgfqpoint{5.860730in}{2.302310in}}%
\pgfpathlineto{\pgfqpoint{5.860383in}{2.302571in}}%
\pgfpathlineto{\pgfqpoint{5.860998in}{2.302454in}}%
\pgfpathlineto{\pgfqpoint{5.862006in}{2.302659in}}%
\pgfpathlineto{\pgfqpoint{5.861719in}{2.301903in}}%
\pgfpathlineto{\pgfqpoint{5.862131in}{2.302558in}}%
\pgfpathlineto{\pgfqpoint{5.862631in}{2.302547in}}%
\pgfpathlineto{\pgfqpoint{5.863163in}{2.303064in}}%
\pgfpathlineto{\pgfqpoint{5.863189in}{2.303044in}}%
\pgfpathlineto{\pgfqpoint{5.863270in}{2.303255in}}%
\pgfpathlineto{\pgfqpoint{5.864064in}{2.302556in}}%
\pgfpathlineto{\pgfqpoint{5.864287in}{2.302682in}}%
\pgfpathlineto{\pgfqpoint{5.864800in}{2.303046in}}%
\pgfpathlineto{\pgfqpoint{5.865311in}{2.302102in}}%
\pgfpathlineto{\pgfqpoint{5.865367in}{2.302165in}}%
\pgfpathlineto{\pgfqpoint{5.866461in}{2.302001in}}%
\pgfpathlineto{\pgfqpoint{5.866059in}{2.302348in}}%
\pgfpathlineto{\pgfqpoint{5.866489in}{2.302120in}}%
\pgfpathlineto{\pgfqpoint{5.867129in}{2.302224in}}%
\pgfpathlineto{\pgfqpoint{5.866772in}{2.301849in}}%
\pgfpathlineto{\pgfqpoint{5.867592in}{2.301979in}}%
\pgfpathlineto{\pgfqpoint{5.868427in}{2.302341in}}%
\pgfpathlineto{\pgfqpoint{5.867647in}{2.301862in}}%
\pgfpathlineto{\pgfqpoint{5.868732in}{2.302022in}}%
\pgfpathlineto{\pgfqpoint{5.869021in}{2.301689in}}%
\pgfpathlineto{\pgfqpoint{5.869586in}{2.302276in}}%
\pgfpathlineto{\pgfqpoint{5.869843in}{2.302016in}}%
\pgfpathlineto{\pgfqpoint{5.870989in}{2.302487in}}%
\pgfpathlineto{\pgfqpoint{5.869884in}{2.301953in}}%
\pgfpathlineto{\pgfqpoint{5.870992in}{2.302468in}}%
\pgfpathlineto{\pgfqpoint{5.872137in}{2.301691in}}%
\pgfpathlineto{\pgfqpoint{5.872139in}{2.301722in}}%
\pgfpathlineto{\pgfqpoint{5.872678in}{2.302167in}}%
\pgfpathlineto{\pgfqpoint{5.873272in}{2.301946in}}%
\pgfpathlineto{\pgfqpoint{5.873391in}{2.301809in}}%
\pgfpathlineto{\pgfqpoint{5.873941in}{2.302268in}}%
\pgfpathlineto{\pgfqpoint{5.874371in}{2.302134in}}%
\pgfpathlineto{\pgfqpoint{5.875241in}{2.302548in}}%
\pgfpathlineto{\pgfqpoint{5.874469in}{2.301979in}}%
\pgfpathlineto{\pgfqpoint{5.875540in}{2.302279in}}%
\pgfpathlineto{\pgfqpoint{5.876386in}{2.302185in}}%
\pgfpathlineto{\pgfqpoint{5.875694in}{2.302493in}}%
\pgfpathlineto{\pgfqpoint{5.876640in}{2.302437in}}%
\pgfpathlineto{\pgfqpoint{5.877425in}{2.302806in}}%
\pgfpathlineto{\pgfqpoint{5.876641in}{2.302431in}}%
\pgfpathlineto{\pgfqpoint{5.877781in}{2.302568in}}%
\pgfpathlineto{\pgfqpoint{5.878243in}{2.302462in}}%
\pgfpathlineto{\pgfqpoint{5.878810in}{2.303263in}}%
\pgfpathlineto{\pgfqpoint{5.878874in}{2.303222in}}%
\pgfpathlineto{\pgfqpoint{5.879868in}{2.303875in}}%
\pgfpathlineto{\pgfqpoint{5.878975in}{2.303079in}}%
\pgfpathlineto{\pgfqpoint{5.880015in}{2.303816in}}%
\pgfpathlineto{\pgfqpoint{5.880711in}{2.303914in}}%
\pgfpathlineto{\pgfqpoint{5.880455in}{2.303334in}}%
\pgfpathlineto{\pgfqpoint{5.881088in}{2.303550in}}%
\pgfpathlineto{\pgfqpoint{5.882086in}{2.303147in}}%
\pgfpathlineto{\pgfqpoint{5.881317in}{2.303658in}}%
\pgfpathlineto{\pgfqpoint{5.882209in}{2.303191in}}%
\pgfpathlineto{\pgfqpoint{5.883486in}{2.303733in}}%
\pgfpathlineto{\pgfqpoint{5.882328in}{2.303096in}}%
\pgfpathlineto{\pgfqpoint{5.883491in}{2.303683in}}%
\pgfpathlineto{\pgfqpoint{5.883652in}{2.303535in}}%
\pgfpathlineto{\pgfqpoint{5.884375in}{2.304264in}}%
\pgfpathlineto{\pgfqpoint{5.884596in}{2.303798in}}%
\pgfpathlineto{\pgfqpoint{5.885602in}{2.303977in}}%
\pgfpathlineto{\pgfqpoint{5.884859in}{2.303620in}}%
\pgfpathlineto{\pgfqpoint{5.885709in}{2.303839in}}%
\pgfpathlineto{\pgfqpoint{5.886340in}{2.303204in}}%
\pgfpathlineto{\pgfqpoint{5.885717in}{2.303847in}}%
\pgfpathlineto{\pgfqpoint{5.886833in}{2.303599in}}%
\pgfpathlineto{\pgfqpoint{5.887195in}{2.303821in}}%
\pgfpathlineto{\pgfqpoint{5.887571in}{2.303394in}}%
\pgfpathlineto{\pgfqpoint{5.887946in}{2.303533in}}%
\pgfpathlineto{\pgfqpoint{5.888101in}{2.303440in}}%
\pgfpathlineto{\pgfqpoint{5.888753in}{2.304115in}}%
\pgfpathlineto{\pgfqpoint{5.889025in}{2.304001in}}%
\pgfpathlineto{\pgfqpoint{5.889025in}{2.304012in}}%
\pgfpathlineto{\pgfqpoint{5.889670in}{2.303565in}}%
\pgfpathlineto{\pgfqpoint{5.890130in}{2.303888in}}%
\pgfpathlineto{\pgfqpoint{5.891170in}{2.303513in}}%
\pgfpathlineto{\pgfqpoint{5.890353in}{2.304264in}}%
\pgfpathlineto{\pgfqpoint{5.891252in}{2.303704in}}%
\pgfpathlineto{\pgfqpoint{5.892086in}{2.303838in}}%
\pgfpathlineto{\pgfqpoint{5.891329in}{2.303844in}}%
\pgfpathlineto{\pgfqpoint{5.892349in}{2.304185in}}%
\pgfpathlineto{\pgfqpoint{5.892698in}{2.304303in}}%
\pgfpathlineto{\pgfqpoint{5.892885in}{2.303734in}}%
\pgfpathlineto{\pgfqpoint{5.893452in}{2.303907in}}%
\pgfpathlineto{\pgfqpoint{5.894880in}{2.304120in}}%
\pgfpathlineto{\pgfqpoint{5.893462in}{2.304018in}}%
\pgfpathlineto{\pgfqpoint{5.894887in}{2.304142in}}%
\pgfpathlineto{\pgfqpoint{5.895644in}{2.304562in}}%
\pgfpathlineto{\pgfqpoint{5.895136in}{2.303918in}}%
\pgfpathlineto{\pgfqpoint{5.895999in}{2.304172in}}%
\pgfpathlineto{\pgfqpoint{5.896729in}{2.304432in}}%
\pgfpathlineto{\pgfqpoint{5.896190in}{2.303880in}}%
\pgfpathlineto{\pgfqpoint{5.897077in}{2.303813in}}%
\pgfpathlineto{\pgfqpoint{5.897132in}{2.303698in}}%
\pgfpathlineto{\pgfqpoint{5.897596in}{2.304403in}}%
\pgfpathlineto{\pgfqpoint{5.898080in}{2.304336in}}%
\pgfpathlineto{\pgfqpoint{5.898671in}{2.304932in}}%
\pgfpathlineto{\pgfqpoint{5.898165in}{2.304275in}}%
\pgfpathlineto{\pgfqpoint{5.899201in}{2.304467in}}%
\pgfpathlineto{\pgfqpoint{5.899729in}{2.304054in}}%
\pgfpathlineto{\pgfqpoint{5.899445in}{2.304495in}}%
\pgfpathlineto{\pgfqpoint{5.900335in}{2.304280in}}%
\pgfpathlineto{\pgfqpoint{5.901034in}{2.304484in}}%
\pgfpathlineto{\pgfqpoint{5.901421in}{2.303821in}}%
\pgfpathlineto{\pgfqpoint{5.901429in}{2.303852in}}%
\pgfpathlineto{\pgfqpoint{5.902501in}{2.304251in}}%
\pgfpathlineto{\pgfqpoint{5.901449in}{2.303758in}}%
\pgfpathlineto{\pgfqpoint{5.902598in}{2.304122in}}%
\pgfpathlineto{\pgfqpoint{5.903561in}{2.303743in}}%
\pgfpathlineto{\pgfqpoint{5.902682in}{2.304207in}}%
\pgfpathlineto{\pgfqpoint{5.903728in}{2.303935in}}%
\pgfpathlineto{\pgfqpoint{5.904369in}{2.304094in}}%
\pgfpathlineto{\pgfqpoint{5.904128in}{2.303602in}}%
\pgfpathlineto{\pgfqpoint{5.904841in}{2.303955in}}%
\pgfpathlineto{\pgfqpoint{5.905971in}{2.303646in}}%
\pgfpathlineto{\pgfqpoint{5.904854in}{2.304046in}}%
\pgfpathlineto{\pgfqpoint{5.906007in}{2.303666in}}%
\pgfpathlineto{\pgfqpoint{5.907138in}{2.303855in}}%
\pgfpathlineto{\pgfqpoint{5.906323in}{2.303443in}}%
\pgfpathlineto{\pgfqpoint{5.907150in}{2.303834in}}%
\pgfpathlineto{\pgfqpoint{5.907232in}{2.303702in}}%
\pgfpathlineto{\pgfqpoint{5.907775in}{2.304394in}}%
\pgfpathlineto{\pgfqpoint{5.908228in}{2.304182in}}%
\pgfpathlineto{\pgfqpoint{5.909334in}{2.304921in}}%
\pgfpathlineto{\pgfqpoint{5.908271in}{2.304058in}}%
\pgfpathlineto{\pgfqpoint{5.909367in}{2.304768in}}%
\pgfpathlineto{\pgfqpoint{5.909797in}{2.304755in}}%
\pgfpathlineto{\pgfqpoint{5.909920in}{2.304998in}}%
\pgfpathlineto{\pgfqpoint{5.910469in}{2.304903in}}%
\pgfpathlineto{\pgfqpoint{5.911608in}{2.305199in}}%
\pgfpathlineto{\pgfqpoint{5.910786in}{2.304695in}}%
\pgfpathlineto{\pgfqpoint{5.911609in}{2.305191in}}%
\pgfpathlineto{\pgfqpoint{5.912812in}{2.305971in}}%
\pgfpathlineto{\pgfqpoint{5.912990in}{2.304891in}}%
\pgfpathlineto{\pgfqpoint{5.914052in}{2.303379in}}%
\pgfpathlineto{\pgfqpoint{5.914210in}{2.303422in}}%
\pgfpathlineto{\pgfqpoint{5.915281in}{2.303449in}}%
\pgfpathlineto{\pgfqpoint{5.914732in}{2.303345in}}%
\pgfpathlineto{\pgfqpoint{5.915320in}{2.303386in}}%
\pgfpathlineto{\pgfqpoint{5.916135in}{2.303357in}}%
\pgfpathlineto{\pgfqpoint{5.915362in}{2.303443in}}%
\pgfpathlineto{\pgfqpoint{5.916432in}{2.303384in}}%
\pgfpathlineto{\pgfqpoint{5.917512in}{2.303452in}}%
\pgfpathlineto{\pgfqpoint{5.916439in}{2.303377in}}%
\pgfpathlineto{\pgfqpoint{5.917545in}{2.303413in}}%
\pgfpathlineto{\pgfqpoint{5.918654in}{2.303441in}}%
\pgfpathlineto{\pgfqpoint{5.917651in}{2.303335in}}%
\pgfpathlineto{\pgfqpoint{5.918660in}{2.303373in}}%
\pgfpathlineto{\pgfqpoint{5.919784in}{2.303440in}}%
\pgfpathlineto{\pgfqpoint{5.918668in}{2.303364in}}%
\pgfpathlineto{\pgfqpoint{5.919805in}{2.303428in}}%
\pgfpathlineto{\pgfqpoint{5.920274in}{2.303313in}}%
\pgfpathlineto{\pgfqpoint{5.920034in}{2.303453in}}%
\pgfpathlineto{\pgfqpoint{5.920915in}{2.303456in}}%
\pgfpathlineto{\pgfqpoint{5.921656in}{2.303357in}}%
\pgfpathlineto{\pgfqpoint{5.922030in}{2.303400in}}%
\pgfpathlineto{\pgfqpoint{5.923066in}{2.303452in}}%
\pgfpathlineto{\pgfqpoint{5.922143in}{2.303381in}}%
\pgfpathlineto{\pgfqpoint{5.923143in}{2.303402in}}%
\pgfpathlineto{\pgfqpoint{5.924276in}{2.303363in}}%
\pgfpathlineto{\pgfqpoint{5.923154in}{2.303442in}}%
\pgfpathlineto{\pgfqpoint{5.924281in}{2.303405in}}%
\pgfpathlineto{\pgfqpoint{5.925354in}{2.303447in}}%
\pgfpathlineto{\pgfqpoint{5.924366in}{2.303357in}}%
\pgfpathlineto{\pgfqpoint{5.925392in}{2.303430in}}%
\pgfpathlineto{\pgfqpoint{5.926469in}{2.303371in}}%
\pgfpathlineto{\pgfqpoint{5.925635in}{2.303450in}}%
\pgfpathlineto{\pgfqpoint{5.926504in}{2.303412in}}%
\pgfpathlineto{\pgfqpoint{5.927583in}{2.303448in}}%
\pgfpathlineto{\pgfqpoint{5.926782in}{2.303357in}}%
\pgfpathlineto{\pgfqpoint{5.927618in}{2.303429in}}%
\pgfpathlineto{\pgfqpoint{5.928709in}{2.303376in}}%
\pgfpathlineto{\pgfqpoint{5.927639in}{2.303443in}}%
\pgfpathlineto{\pgfqpoint{5.928733in}{2.303420in}}%
\pgfpathlineto{\pgfqpoint{5.929800in}{2.303447in}}%
\pgfpathlineto{\pgfqpoint{5.928776in}{2.303380in}}%
\pgfpathlineto{\pgfqpoint{5.929847in}{2.303394in}}%
\pgfpathlineto{\pgfqpoint{5.931023in}{2.303433in}}%
\pgfpathlineto{\pgfqpoint{5.929877in}{2.303376in}}%
\pgfpathlineto{\pgfqpoint{5.931028in}{2.303412in}}%
\pgfpathlineto{\pgfqpoint{5.932141in}{2.303053in}}%
\pgfpathlineto{\pgfqpoint{5.931053in}{2.303442in}}%
\pgfpathlineto{\pgfqpoint{5.932157in}{2.303152in}}%
\pgfpathlineto{\pgfqpoint{5.932323in}{2.303024in}}%
\pgfpathlineto{\pgfqpoint{5.932908in}{2.303633in}}%
\pgfpathlineto{\pgfqpoint{5.933225in}{2.303571in}}%
\pgfpathlineto{\pgfqpoint{5.933600in}{2.303796in}}%
\pgfpathlineto{\pgfqpoint{5.934093in}{2.302809in}}%
\pgfpathlineto{\pgfqpoint{5.934255in}{2.302880in}}%
\pgfpathlineto{\pgfqpoint{5.934972in}{2.302625in}}%
\pgfpathlineto{\pgfqpoint{5.935353in}{2.303024in}}%
\pgfpathlineto{\pgfqpoint{5.935362in}{2.302990in}}%
\pgfpathlineto{\pgfqpoint{5.936352in}{2.303099in}}%
\pgfpathlineto{\pgfqpoint{5.935486in}{2.303363in}}%
\pgfpathlineto{\pgfqpoint{5.936429in}{2.303322in}}%
\pgfpathlineto{\pgfqpoint{5.937043in}{2.303663in}}%
\pgfpathlineto{\pgfqpoint{5.936472in}{2.303258in}}%
\pgfpathlineto{\pgfqpoint{5.937516in}{2.303097in}}%
\pgfpathlineto{\pgfqpoint{5.937984in}{2.302991in}}%
\pgfpathlineto{\pgfqpoint{5.938540in}{2.303782in}}%
\pgfpathlineto{\pgfqpoint{5.938604in}{2.303722in}}%
\pgfpathlineto{\pgfqpoint{5.938605in}{2.303738in}}%
\pgfpathlineto{\pgfqpoint{5.939268in}{2.303136in}}%
\pgfpathlineto{\pgfqpoint{5.939707in}{2.303592in}}%
\pgfpathlineto{\pgfqpoint{5.940774in}{2.303480in}}%
\pgfpathlineto{\pgfqpoint{5.939763in}{2.303611in}}%
\pgfpathlineto{\pgfqpoint{5.940838in}{2.303616in}}%
\pgfpathlineto{\pgfqpoint{5.940880in}{2.303698in}}%
\pgfpathlineto{\pgfqpoint{5.941750in}{2.303098in}}%
\pgfpathlineto{\pgfqpoint{5.941938in}{2.303331in}}%
\pgfpathlineto{\pgfqpoint{5.943038in}{2.302893in}}%
\pgfpathlineto{\pgfqpoint{5.941950in}{2.303366in}}%
\pgfpathlineto{\pgfqpoint{5.943081in}{2.303000in}}%
\pgfpathlineto{\pgfqpoint{5.943227in}{2.303313in}}%
\pgfpathlineto{\pgfqpoint{5.944146in}{2.302445in}}%
\pgfpathlineto{\pgfqpoint{5.944147in}{2.302450in}}%
\pgfpathlineto{\pgfqpoint{5.944185in}{2.302223in}}%
\pgfpathlineto{\pgfqpoint{5.945200in}{2.303327in}}%
\pgfpathlineto{\pgfqpoint{5.945236in}{2.303222in}}%
\pgfpathlineto{\pgfqpoint{5.948299in}{2.303919in}}%
\pgfpathlineto{\pgfqpoint{5.945249in}{2.303175in}}%
\pgfpathlineto{\pgfqpoint{5.948300in}{2.303897in}}%
\pgfpathlineto{\pgfqpoint{5.948775in}{2.303519in}}%
\pgfpathlineto{\pgfqpoint{5.948363in}{2.304021in}}%
\pgfpathlineto{\pgfqpoint{5.949416in}{2.303844in}}%
\pgfpathlineto{\pgfqpoint{5.949580in}{2.304054in}}%
\pgfpathlineto{\pgfqpoint{5.950376in}{2.303527in}}%
\pgfpathlineto{\pgfqpoint{5.950517in}{2.303581in}}%
\pgfpathlineto{\pgfqpoint{5.951654in}{2.303151in}}%
\pgfpathlineto{\pgfqpoint{5.950530in}{2.303613in}}%
\pgfpathlineto{\pgfqpoint{5.951800in}{2.303420in}}%
\pgfpathlineto{\pgfqpoint{5.953523in}{2.303819in}}%
\pgfpathlineto{\pgfqpoint{5.951821in}{2.303370in}}%
\pgfpathlineto{\pgfqpoint{5.953543in}{2.303731in}}%
\pgfpathlineto{\pgfqpoint{5.953969in}{2.303218in}}%
\pgfpathlineto{\pgfqpoint{5.953576in}{2.303794in}}%
\pgfpathlineto{\pgfqpoint{5.954680in}{2.303492in}}%
\pgfpathlineto{\pgfqpoint{5.956163in}{2.303569in}}%
\pgfpathlineto{\pgfqpoint{5.954686in}{2.303435in}}%
\pgfpathlineto{\pgfqpoint{5.956208in}{2.303420in}}%
\pgfpathlineto{\pgfqpoint{5.956458in}{2.303414in}}%
\pgfpathlineto{\pgfqpoint{5.957300in}{2.303801in}}%
\pgfpathlineto{\pgfqpoint{5.957303in}{2.303788in}}%
\pgfpathlineto{\pgfqpoint{5.958379in}{2.303520in}}%
\pgfpathlineto{\pgfqpoint{5.957854in}{2.304024in}}%
\pgfpathlineto{\pgfqpoint{5.958428in}{2.303599in}}%
\pgfpathlineto{\pgfqpoint{5.959142in}{2.303869in}}%
\pgfpathlineto{\pgfqpoint{5.958686in}{2.303408in}}%
\pgfpathlineto{\pgfqpoint{5.959549in}{2.303809in}}%
\pgfpathlineto{\pgfqpoint{5.960599in}{2.303650in}}%
\pgfpathlineto{\pgfqpoint{5.960125in}{2.303955in}}%
\pgfpathlineto{\pgfqpoint{5.960666in}{2.303700in}}%
\pgfpathlineto{\pgfqpoint{5.961711in}{2.303244in}}%
\pgfpathlineto{\pgfqpoint{5.961811in}{2.303415in}}%
\pgfpathlineto{\pgfqpoint{5.962649in}{2.304268in}}%
\pgfpathlineto{\pgfqpoint{5.961859in}{2.303351in}}%
\pgfpathlineto{\pgfqpoint{5.962946in}{2.304019in}}%
\pgfpathlineto{\pgfqpoint{5.963557in}{2.303765in}}%
\pgfpathlineto{\pgfqpoint{5.962970in}{2.304047in}}%
\pgfpathlineto{\pgfqpoint{5.964108in}{2.304052in}}%
\pgfpathlineto{\pgfqpoint{5.965238in}{2.304223in}}%
\pgfpathlineto{\pgfqpoint{5.964736in}{2.303826in}}%
\pgfpathlineto{\pgfqpoint{5.965239in}{2.304211in}}%
\pgfpathlineto{\pgfqpoint{5.966354in}{2.303788in}}%
\pgfpathlineto{\pgfqpoint{5.965247in}{2.304234in}}%
\pgfpathlineto{\pgfqpoint{5.966404in}{2.303906in}}%
\pgfpathlineto{\pgfqpoint{5.967571in}{2.304250in}}%
\pgfpathlineto{\pgfqpoint{5.966642in}{2.303650in}}%
\pgfpathlineto{\pgfqpoint{5.967600in}{2.304201in}}%
\pgfpathlineto{\pgfqpoint{5.968550in}{2.303659in}}%
\pgfpathlineto{\pgfqpoint{5.967604in}{2.304209in}}%
\pgfpathlineto{\pgfqpoint{5.968740in}{2.303828in}}%
\pgfpathlineto{\pgfqpoint{5.969517in}{2.304012in}}%
\pgfpathlineto{\pgfqpoint{5.969185in}{2.303454in}}%
\pgfpathlineto{\pgfqpoint{5.969839in}{2.303627in}}%
\pgfpathlineto{\pgfqpoint{5.970364in}{2.303135in}}%
\pgfpathlineto{\pgfqpoint{5.969841in}{2.303637in}}%
\pgfpathlineto{\pgfqpoint{5.970976in}{2.303274in}}%
\pgfpathlineto{\pgfqpoint{5.972019in}{2.303399in}}%
\pgfpathlineto{\pgfqpoint{5.971214in}{2.302894in}}%
\pgfpathlineto{\pgfqpoint{5.972086in}{2.303245in}}%
\pgfpathlineto{\pgfqpoint{5.972611in}{2.302901in}}%
\pgfpathlineto{\pgfqpoint{5.973199in}{2.303248in}}%
\pgfpathlineto{\pgfqpoint{5.974266in}{2.303377in}}%
\pgfpathlineto{\pgfqpoint{5.973578in}{2.302903in}}%
\pgfpathlineto{\pgfqpoint{5.974318in}{2.303254in}}%
\pgfpathlineto{\pgfqpoint{5.975082in}{2.303283in}}%
\pgfpathlineto{\pgfqpoint{5.974461in}{2.303576in}}%
\pgfpathlineto{\pgfqpoint{5.975328in}{2.303745in}}%
\pgfpathlineto{\pgfqpoint{5.976487in}{2.304229in}}%
\pgfpathlineto{\pgfqpoint{5.975342in}{2.303717in}}%
\pgfpathlineto{\pgfqpoint{5.976489in}{2.304207in}}%
\pgfpathlineto{\pgfqpoint{5.977407in}{2.304114in}}%
\pgfpathlineto{\pgfqpoint{5.976715in}{2.304430in}}%
\pgfpathlineto{\pgfqpoint{5.977607in}{2.304219in}}%
\pgfpathlineto{\pgfqpoint{5.978675in}{2.304479in}}%
\pgfpathlineto{\pgfqpoint{5.977875in}{2.303845in}}%
\pgfpathlineto{\pgfqpoint{5.978726in}{2.304344in}}%
\pgfpathlineto{\pgfqpoint{5.979101in}{2.304133in}}%
\pgfpathlineto{\pgfqpoint{5.979642in}{2.304883in}}%
\pgfpathlineto{\pgfqpoint{5.979833in}{2.304441in}}%
\pgfpathlineto{\pgfqpoint{5.980130in}{2.304374in}}%
\pgfpathlineto{\pgfqpoint{5.980601in}{2.303579in}}%
\pgfpathlineto{\pgfqpoint{5.983713in}{2.302410in}}%
\pgfpathlineto{\pgfqpoint{5.980605in}{2.303603in}}%
\pgfpathlineto{\pgfqpoint{5.983745in}{2.302623in}}%
\pgfpathlineto{\pgfqpoint{5.984525in}{2.302825in}}%
\pgfpathlineto{\pgfqpoint{5.983944in}{2.302413in}}%
\pgfpathlineto{\pgfqpoint{5.984856in}{2.302600in}}%
\pgfpathlineto{\pgfqpoint{5.986645in}{2.303234in}}%
\pgfpathlineto{\pgfqpoint{5.984885in}{2.302483in}}%
\pgfpathlineto{\pgfqpoint{5.986699in}{2.303073in}}%
\pgfpathlineto{\pgfqpoint{5.987102in}{2.302763in}}%
\pgfpathlineto{\pgfqpoint{5.986980in}{2.303135in}}%
\pgfpathlineto{\pgfqpoint{5.987809in}{2.303072in}}%
\pgfpathlineto{\pgfqpoint{5.988958in}{2.303452in}}%
\pgfpathlineto{\pgfqpoint{5.987813in}{2.303070in}}%
\pgfpathlineto{\pgfqpoint{5.988987in}{2.303419in}}%
\pgfpathlineto{\pgfqpoint{5.990044in}{2.303372in}}%
\pgfpathlineto{\pgfqpoint{5.989047in}{2.303451in}}%
\pgfpathlineto{\pgfqpoint{5.990107in}{2.303394in}}%
\pgfpathlineto{\pgfqpoint{5.991176in}{2.303441in}}%
\pgfpathlineto{\pgfqpoint{5.990386in}{2.303358in}}%
\pgfpathlineto{\pgfqpoint{5.991219in}{2.303419in}}%
\pgfpathlineto{\pgfqpoint{5.992327in}{2.303439in}}%
\pgfpathlineto{\pgfqpoint{5.991304in}{2.303374in}}%
\pgfpathlineto{\pgfqpoint{5.992333in}{2.303425in}}%
\pgfpathlineto{\pgfqpoint{5.993016in}{2.303355in}}%
\pgfpathlineto{\pgfqpoint{5.992509in}{2.303448in}}%
\pgfpathlineto{\pgfqpoint{5.993444in}{2.303412in}}%
\pgfpathlineto{\pgfqpoint{5.994553in}{2.303442in}}%
\pgfpathlineto{\pgfqpoint{5.993655in}{2.303361in}}%
\pgfpathlineto{\pgfqpoint{5.994555in}{2.303391in}}%
\pgfpathlineto{\pgfqpoint{5.995502in}{2.303449in}}%
\pgfpathlineto{\pgfqpoint{5.994786in}{2.303367in}}%
\pgfpathlineto{\pgfqpoint{5.995669in}{2.303430in}}%
\pgfpathlineto{\pgfqpoint{5.996743in}{2.303376in}}%
\pgfpathlineto{\pgfqpoint{5.995675in}{2.303448in}}%
\pgfpathlineto{\pgfqpoint{5.996779in}{2.303415in}}%
\pgfpathlineto{\pgfqpoint{5.997806in}{2.303450in}}%
\pgfpathlineto{\pgfqpoint{5.996890in}{2.303349in}}%
\pgfpathlineto{\pgfqpoint{5.997891in}{2.303410in}}%
\pgfpathlineto{\pgfqpoint{5.998990in}{2.303446in}}%
\pgfpathlineto{\pgfqpoint{5.998115in}{2.303367in}}%
\pgfpathlineto{\pgfqpoint{5.999004in}{2.303416in}}%
\pgfpathlineto{\pgfqpoint{6.000119in}{2.303375in}}%
\pgfpathlineto{\pgfqpoint{5.999009in}{2.303435in}}%
\pgfpathlineto{\pgfqpoint{6.000135in}{2.303422in}}%
\pgfpathlineto{\pgfqpoint{6.001235in}{2.303448in}}%
\pgfpathlineto{\pgfqpoint{6.000306in}{2.303371in}}%
\pgfpathlineto{\pgfqpoint{6.001246in}{2.303431in}}%
\pgfpathlineto{\pgfqpoint{6.001267in}{2.303457in}}%
\pgfpathlineto{\pgfqpoint{6.002362in}{2.303366in}}%
\pgfpathlineto{\pgfqpoint{6.003511in}{2.303442in}}%
\pgfpathlineto{\pgfqpoint{6.002428in}{2.303345in}}%
\pgfpathlineto{\pgfqpoint{6.003511in}{2.303428in}}%
\pgfpathlineto{\pgfqpoint{6.003852in}{2.303314in}}%
\pgfpathlineto{\pgfqpoint{6.003597in}{2.303448in}}%
\pgfpathlineto{\pgfqpoint{6.004622in}{2.303418in}}%
\pgfpathlineto{\pgfqpoint{6.005722in}{2.303438in}}%
\pgfpathlineto{\pgfqpoint{6.004805in}{2.303356in}}%
\pgfpathlineto{\pgfqpoint{6.005735in}{2.303425in}}%
\pgfpathlineto{\pgfqpoint{6.006522in}{2.303355in}}%
\pgfpathlineto{\pgfqpoint{6.005782in}{2.303444in}}%
\pgfpathlineto{\pgfqpoint{6.006575in}{2.303394in}}%
\pgfusepath{stroke}%
\end{pgfscope}%
\begin{pgfscope}%
\pgfsetrectcap%
\pgfsetmiterjoin%
\pgfsetlinewidth{0.803000pt}%
\definecolor{currentstroke}{rgb}{0.000000,0.000000,0.000000}%
\pgfsetstrokecolor{currentstroke}%
\pgfsetdash{}{0pt}%
\pgfpathmoveto{\pgfqpoint{0.894648in}{0.771190in}}%
\pgfpathlineto{\pgfqpoint{0.894648in}{4.601775in}}%
\pgfusepath{stroke}%
\end{pgfscope}%
\begin{pgfscope}%
\pgfsetrectcap%
\pgfsetmiterjoin%
\pgfsetlinewidth{0.803000pt}%
\definecolor{currentstroke}{rgb}{0.000000,0.000000,0.000000}%
\pgfsetstrokecolor{currentstroke}%
\pgfsetdash{}{0pt}%
\pgfpathmoveto{\pgfqpoint{6.250000in}{0.771190in}}%
\pgfpathlineto{\pgfqpoint{6.250000in}{4.601775in}}%
\pgfusepath{stroke}%
\end{pgfscope}%
\begin{pgfscope}%
\pgfsetrectcap%
\pgfsetmiterjoin%
\pgfsetlinewidth{0.803000pt}%
\definecolor{currentstroke}{rgb}{0.000000,0.000000,0.000000}%
\pgfsetstrokecolor{currentstroke}%
\pgfsetdash{}{0pt}%
\pgfpathmoveto{\pgfqpoint{0.894648in}{0.771190in}}%
\pgfpathlineto{\pgfqpoint{6.250000in}{0.771190in}}%
\pgfusepath{stroke}%
\end{pgfscope}%
\begin{pgfscope}%
\pgfsetrectcap%
\pgfsetmiterjoin%
\pgfsetlinewidth{0.803000pt}%
\definecolor{currentstroke}{rgb}{0.000000,0.000000,0.000000}%
\pgfsetstrokecolor{currentstroke}%
\pgfsetdash{}{0pt}%
\pgfpathmoveto{\pgfqpoint{0.894648in}{4.601775in}}%
\pgfpathlineto{\pgfqpoint{6.250000in}{4.601775in}}%
\pgfusepath{stroke}%
\end{pgfscope}%
\end{pgfpicture}%
\makeatother%
\endgroup%
}
    \caption{Traveled path at $\Delta U = 80 eV$}
    \label{fig:particle_path_flashing}
\end{figure}



\begin{equation}
    U(x,t) = U_r(x)f(t),
\end{equation}

\begin{equation}
    U_r(x) = \begin{cases}
        \frac{x}{\alpha}\Delta U, & 0 \leq x < \alpha L \\
        \frac{L - x}{\alpha - 1}\Delta U, & \alpha L \leq x < L
    \end{cases}
\end{equation}

\begin{equation}
    f(t) = \begin{cases}
        1, & t < t_0 \\
        0, & t \geq t_0
    \end{cases}
\end{equation}

\begin{equation}
    x_{n+1} = x_n - \frac{1}{{\gamma}_i} \frac{\partial U}{\partial x}(x_n, t_n) \delta t + \sqrt{\frac{2 k_B T \delta t}{{\gamma}_i}} \hat{\xi}_n,
\end{equation}

\begin{equation}
    \left\lvert{x_n - x}\right\rvert << \alpha L
\end{equation}


Choosing {$\delta t$} such that it the equation 
\begin{equation}
    \frac{1}{{\gamma}_i} \max \left\lvert{\frac{\partial U}{\partial x}}\right\rvert \delta t << \alpha L
\end{equation} 
is satisfied, ensures that the particle does not move more than a fraction of the potential's spatial period in one time step. This is crucial for accurately capturing the dynamics of the system, as it allows for a more precise representation of the particle's motion within the potential landscape. By keeping the time step sufficiently small, we can ensure that the numerical integration of the particle's trajectory remains stable and accurate, preventing any significant overshooting that could lead to erroneous results in our simulations.

\subsection{Reduced units}
To simplify our calculations and reduce the number of parameters in our system, we can introduce reduced units. This involves choosing characteristic scales for length, energy, and time, which allows us to express all quantities in terms of these scales.
\begin{equation}
    \hat{x}_n = \hat{x} 
\end{equation}

\begin{gather}
    \hat{x} = \frac{x}{L}, \quad
    \hat{t} = \omega t, \quad
    \omega = \frac{\Delta U}{\gamma_i L^2}, \\
    \hat{U}(\hat{x},\hat{t}) = \frac{U(x,t)}{\Delta U}, \quad
    \hat{D} = \frac{k_B T}{\Delta U}
\end{gather}

\begin{equation}
    p(U) = \frac{e^{-\frac{U}{k_B T}}}{k_B T (1-e^{-\frac{\Delta U}{k_B T}})} 
\end{equation}

\begin{equation}
    n(x,t) = \frac{N}{\sqrt{4 \pi D t}} \exp\left(-\frac{x^2}{4Dt}\right)
\end{equation}


\begin{equation}
    m \frac{d^2 \bf{x}}{d t^2} = - \nabla U(\bf{x}, t) - \gamma \frac{d\bf{x}}{dt} + \xi (t)
\end{equation}
